% Options for packages loaded elsewhere
\PassOptionsToPackage{unicode}{hyperref}
\PassOptionsToPackage{hyphens}{url}
\PassOptionsToPackage{space}{xeCJK}
\documentclass[
]{ctexart}
\usepackage{xcolor}
\usepackage[margin=2.5cm]{geometry}
\usepackage{amsmath,amssymb}
\setcounter{secnumdepth}{5}
\usepackage{iftex}
\ifPDFTeX
  \usepackage[T1]{fontenc}
  \usepackage[utf8]{inputenc}
  \usepackage{textcomp} % provide euro and other symbols
\else % if luatex or xetex
  \usepackage{unicode-math} % this also loads fontspec
  \defaultfontfeatures{Scale=MatchLowercase}
  \defaultfontfeatures[\rmfamily]{Ligatures=TeX,Scale=1}
\fi
\usepackage{lmodern}
\ifPDFTeX\else
  % xetex/luatex font selection
  \setmainfont[]{Times New Roman}
  \ifXeTeX
    \usepackage{xeCJK}
    \setCJKmainfont[]{SimSun}
    \setCJKsansfont[]{SimHei}
    \setCJKmonofont[]{KaiTi}
  \fi
  \ifLuaTeX
    \usepackage[]{luatexja-fontspec}
    \setmainjfont[]{SimSun}
  \fi
\fi
% Use upquote if available, for straight quotes in verbatim environments
\IfFileExists{upquote.sty}{\usepackage{upquote}}{}
\IfFileExists{microtype.sty}{% use microtype if available
  \usepackage[]{microtype}
  \UseMicrotypeSet[protrusion]{basicmath} % disable protrusion for tt fonts
}{}
\makeatletter
\@ifundefined{KOMAClassName}{% if non-KOMA class
  \IfFileExists{parskip.sty}{%
    \usepackage{parskip}
  }{% else
    \setlength{\parindent}{0pt}
    \setlength{\parskip}{6pt plus 2pt minus 1pt}}
}{% if KOMA class
  \KOMAoptions{parskip=half}}
\makeatother
\usepackage{longtable,booktabs,array}
\newcounter{none} % for unnumbered tables
\usepackage{calc} % for calculating minipage widths
% Correct order of tables after \paragraph or \subparagraph
\usepackage{etoolbox}
\makeatletter
\patchcmd\longtable{\par}{\if@noskipsec\mbox{}\fi\par}{}{}
\makeatother
% Allow footnotes in longtable head/foot
\IfFileExists{footnotehyper.sty}{\usepackage{footnotehyper}}{\usepackage{footnote}}
\makesavenoteenv{longtable}
\setlength{\emergencystretch}{3em} % prevent overfull lines
\providecommand{\tightlist}{%
  \setlength{\itemsep}{0pt}\setlength{\parskip}{0pt}}

\usepackage{ctex}
\usepackage{amsmath,amssymb,amsthm,amscd}
\usepackage{geometry}
\usepackage{fancyhdr}
\usepackage{hyperref}
\usepackage{xcolor}
\usepackage{booktabs}
\usepackage{longtable}
\usepackage{enumitem}

\geometry{a4paper, margin=2.5cm}

\pagestyle{fancy}
\fancyhf{}
\fancyhead[L]{\leftmark}
\fancyhead[R]{\thepage}
\fancyfoot[C]{}

\hypersetup{
    colorlinks=true,
    linkcolor=blue!60!black,
    urlcolor=blue!70!black,
    citecolor=green!50!black,
    pdfauthor={数学教材},
    pdftitle={从公理到中学数学：严谨推导全书},
}

% 常用数学运算符定义
\DeclareMathOperator{\Sha}{III}
\DeclareMathOperator{\Gal}{Gal}
\DeclareMathOperator{\Aut}{Aut}
\DeclareMathOperator{\End}{End}
\DeclareMathOperator{\Hom}{Hom}
\DeclareMathOperator{\Spec}{Spec}
\DeclareMathOperator{\Proj}{Proj}
\DeclareMathOperator{\Ext}{Ext}
\DeclareMathOperator{\Tor}{Tor}

% 证明结束符号定义
\protected\def\proofend{\ensuremath{\blacksquare}}
\renewcommand{\qedsymbol}{$\blacksquare$}

\setlength{\parskip}{0.5em}
\setlength{\parindent}{0pt}
\usepackage{bookmark}
\IfFileExists{xurl.sty}{\usepackage{xurl}}{} % add URL line breaks if available
\urlstyle{same}
\hypersetup{
  pdftitle={从公理到整个数学体系},
  pdfauthor={数学教材},
  hidelinks,
  pdfcreator={LaTeX via pandoc}}

\title{从公理到整个数学体系}
\author{数学教材}
\date{}

\begin{document}
\maketitle

{
\setcounter{tocdepth}{2}
\tableofcontents
}
\section{卷一：数学基础}\label{ux5377ux4e00ux6570ux5b66ux57faux7840}

\begin{quote}
本卷是全书的根基。我们从最基础的数学符号和逻辑语言出发，逐步建立集合论的严格框架，然后在集合论的土壤上构造全部数系（\(\mathbb{N}, \mathbb{Z}, \mathbb{Q}, \mathbb{R}, \mathbb{C}\)），最后以整数的深层结构------数论------收束本卷。五章之间环环相扣：符号是语言，逻辑是推理规则，集合是描述对象的工具，数系是研究的载体，数论则是对整数最本质规律的探索。
\end{quote}

\bigskip\hrule\bigskip

\subsection{Ch 1
数学符号与数学语言}\label{ch-1-ux6570ux5b66ux7b26ux53f7ux4e0eux6570ux5b66ux8bedux8a00}

\textbf{前置知识}：本章是全书第一章，不依赖任何前置章节。数学符号是全书的语言基础------后续所有定义、定理、证明均在此符号体系中书写。

数学是一门用精确语言书写的学科。正如建筑需要砖瓦、音乐需要音符，数学需要一套统一的符号体系。本章系统介绍全书所用的数学符号，为后续所有章节奠定语言基础。

\bigskip\hrule\bigskip

\subsubsection{1.1 逻辑符号}\label{ux903bux8f91ux7b26ux53f7}

\textbf{定义
1.1.1}：以下符号统称为\textbf{逻辑符号}，用于表达命题之间的逻辑关系。

{\def\LTcaptype{none} % do not increment counter
\begin{longtable}[]{@{}
  >{\centering\arraybackslash}p{(\linewidth - 4\tabcolsep) * \real{0.3333}}
  >{\centering\arraybackslash}p{(\linewidth - 4\tabcolsep) * \real{0.3333}}
  >{\centering\arraybackslash}p{(\linewidth - 4\tabcolsep) * \real{0.3333}}@{}}
\toprule\noalign{}
\begin{minipage}[b]{\linewidth}\centering
符号
\end{minipage} & \begin{minipage}[b]{\linewidth}\centering
读法
\end{minipage} & \begin{minipage}[b]{\linewidth}\centering
含义
\end{minipage} \\
\midrule\noalign{}
\endhead
\bottomrule\noalign{}
\endlastfoot
\(\neg\) & ``非'' & 否定：\(\neg P\) 表示''\(P\) 不成立'' \\
\(\wedge\) & ``且'' & 合取：\(P \wedge Q\) 表示''\(P\) 且 \(Q\)'' \\
\(\vee\) & ``或'' & 析取：\(P \vee Q\) 表示''\(P\) 或
\(Q\)``（可兼或） \\
\(\rightarrow\) & ``蕴含'' & 蕴含：\(P \rightarrow Q\) 表示''若 \(P\) 则
\(Q\)'' \\
\(\leftrightarrow\) & ``当且仅当'' & 等价：\(P \leftrightarrow Q\)
表示''\(P\) 与 \(Q\) 等价'' \\
\(\forall\) & ``对所有'' &
全称量词：\(\forall x \in \mathbb{R},\, x^2 \geq 0\) \\
\(\exists\) & ``存在'' &
存在量词：\(\exists x \in \mathbb{R},\, x^2 = 2\) \\
\(\exists!\) & ``存在唯一'' & 存在唯一量词 \\
\end{longtable}
}

\textbf{注记 1.1.1}：符号 \(\rightarrow\) 与 \(\Rightarrow\)
在用法上有细微区别。\(\rightarrow\) 用于逻辑公式内部（如
\(P \rightarrow Q\)），\(\Rightarrow\) 用于推理步骤之间（如''由条件
\(A\) \(\Rightarrow\) 结论 \(B\)``）。类似地，\(\leftrightarrow\)
用于公式内部，\(\Leftrightarrow\) 用于推理步骤之间。

\textbf{注记 1.1.2}：在实际数学写作中，``\(P \Rightarrow Q\)'' 和
``\(P \rightarrow Q\)''
经常混用，不严格区分。但在形式逻辑中，\(\rightarrow\)
是对象语言中的符号（属于被研究的逻辑系统），而 \(\Rightarrow\)
是元语言中的符号（用于讨论推理关系）。本教材采用实用主义立场，在不引起歧义时不严格区分两者。

\bigskip\hrule\bigskip

\subsubsection{1.2 集合符号}\label{ux96c6ux5408ux7b26ux53f7}

\textbf{定义 1.2.1}：以下符号用于描述集合及其运算。

{\def\LTcaptype{none} % do not increment counter
\begin{longtable}[]{@{}
  >{\centering\arraybackslash}p{(\linewidth - 4\tabcolsep) * \real{0.3333}}
  >{\centering\arraybackslash}p{(\linewidth - 4\tabcolsep) * \real{0.3333}}
  >{\centering\arraybackslash}p{(\linewidth - 4\tabcolsep) * \real{0.3333}}@{}}
\toprule\noalign{}
\begin{minipage}[b]{\linewidth}\centering
符号
\end{minipage} & \begin{minipage}[b]{\linewidth}\centering
读法
\end{minipage} & \begin{minipage}[b]{\linewidth}\centering
含义
\end{minipage} \\
\midrule\noalign{}
\endhead
\bottomrule\noalign{}
\endlastfoot
\(\in\) & ``属于'' & 元素与集合的关系：\(3 \in \mathbb{N}\) \\
\(\notin\) & ``不属于'' & \(\pi \notin \mathbb{Q}\) \\
\(\subseteq\) & ``包含于'' &
子集关系：\(\mathbb{N} \subseteq \mathbb{Z}\) \\
\(\subsetneq\) & ``真包含于'' &
真子集：\(\mathbb{N} \subsetneq \mathbb{Z}\) \\
\(\cup\) & ``并'' & 并集：\(\{1,2\} \cup \{2,3\} = \{1,2,3\}\) \\
\(\cap\) & ``交'' & 交集：\(\{1,2\} \cap \{2,3\} = \{2\}\) \\
\(\setminus\) & ``差'' &
差集：\(\{1,2,3\} \setminus \{2\} = \{1,3\}\) \\
\(\varnothing\) & ``空集'' & 不含任何元素的集合 \\
\(\mathcal{P}(A)\) & ``\(A\) 的幂集'' & \(A\) 的所有子集构成的集合 \\
\(A \times B\) & ``笛卡尔积'' & 有序对的集合 \\
\(\bigcup_{i \in I} A_i\) & ``族的并'' & 以 \(I\)
为指标集的集族的并集 \\
\(\bigcap_{i \in I} A_i\) & ``族的交'' & 以 \(I\)
为指标集的集族的交集 \\
\end{longtable}
}

\textbf{注记 1.2.1}：符号 \(\subseteq\) 和 \(\subsetneq\)
的区别在于：\(A \subseteq B\) 允许 \(A = B\)，而 \(A \subsetneq B\) 要求
\(A \neq B\)。某些文献使用 \(\subset\) 表示
\(\subseteq\)（即允许相等），读者需根据上下文判断。

\bigskip\hrule\bigskip

\subsubsection{1.3 数系符号}\label{ux6570ux7cfbux7b26ux53f7}

\textbf{定义 1.3.1}：以下黑板粗体字母表示标准数系。

{\def\LTcaptype{none} % do not increment counter
\begin{longtable}[]{@{}
  >{\centering\arraybackslash}p{(\linewidth - 2\tabcolsep) * \real{0.5000}}
  >{\centering\arraybackslash}p{(\linewidth - 2\tabcolsep) * \real{0.5000}}@{}}
\toprule\noalign{}
\begin{minipage}[b]{\linewidth}\centering
符号
\end{minipage} & \begin{minipage}[b]{\linewidth}\centering
含义
\end{minipage} \\
\midrule\noalign{}
\endhead
\bottomrule\noalign{}
\endlastfoot
\(\mathbb{N}\) & 自然数集 \(\{0, 1, 2, 3, \ldots\}\) \\
\(\mathbb{N}^*\) 或 \(\mathbb{N}_+\) & 正整数集
\(\{1, 2, 3, \ldots\}\) \\
\(\mathbb{Z}\) & 整数集 \(\{\ldots, -2, -1, 0, 1, 2, \ldots\}\) \\
\(\mathbb{Q}\) & 有理数集 \\
\(\mathbb{R}\) & 实数集 \\
\(\mathbb{C}\) & 复数集 \\
\(\mathbb{R}^+\) & 正实数集 \(\{x \in \mathbb{R} \mid x > 0\}\) \\
\(\mathbb{R}_{\geq 0}\) & 非负实数集
\(\{x \in \mathbb{R} \mid x \geq 0\}\) \\
\(\mathbb{Z}_m\) & 模 \(m\) 的剩余类环
\(\{[0], [1], \ldots, [m-1]\}\) \\
\end{longtable}
}

\bigskip\hrule\bigskip

\subsubsection{1.4
关系与序符号}\label{ux5173ux7cfbux4e0eux5e8fux7b26ux53f7}

\textbf{定义 1.4.1}：以下符号表示数学对象之间的关系。

{\def\LTcaptype{none} % do not increment counter
\begin{longtable}[]{@{}
  >{\centering\arraybackslash}p{(\linewidth - 4\tabcolsep) * \real{0.3333}}
  >{\centering\arraybackslash}p{(\linewidth - 4\tabcolsep) * \real{0.3333}}
  >{\centering\arraybackslash}p{(\linewidth - 4\tabcolsep) * \real{0.3333}}@{}}
\toprule\noalign{}
\begin{minipage}[b]{\linewidth}\centering
符号
\end{minipage} & \begin{minipage}[b]{\linewidth}\centering
读法
\end{minipage} & \begin{minipage}[b]{\linewidth}\centering
含义
\end{minipage} \\
\midrule\noalign{}
\endhead
\bottomrule\noalign{}
\endlastfoot
\(=\) & ``等于'' & 相等 \\
\(\neq\) & ``不等于'' & 不相等 \\
\(<, >\) & ``小于/大于'' & 严格序 \\
\(\leq, \geq\) & ``小于等于/大于等于'' & 非严格序 \\
\(\ll\) & ``远小于'' & 渐近意义下远小于 \\
\(\sim\) & ``等价'' & 等价关系 \\
\(\approx\) & ``约等于'' & 近似相等 \\
\(\equiv\) & ``恒等于''或''模等价'' & \(f(x) \equiv 0\) 或
\(5 \equiv 2 \pmod{3}\) \\
\(\cong\) & ``全等''或''同构'' &
\(\triangle ABC \cong \triangle DEF\) \\
\end{longtable}
}

\textbf{注记 1.4.1}：符号 \(=\) 在数学中具有严格含义。根据 Ch 3
中的外延公理，\(A = B\) 意味着 \(A\) 和 \(B\) 是同一个对象。符号
\(\approx\) 表示近似关系，不满足精确的相等条件。符号 \(\equiv\)
的含义取决于上下文：在恒等式中表示对所有变量值都成立的等式，在同余中表示模意义下的等价。

\bigskip\hrule\bigskip

\subsubsection{1.5 运算符号}\label{ux8fd0ux7b97ux7b26ux53f7}

\textbf{定义 1.5.1}：以下符号表示基本运算。

{\def\LTcaptype{none} % do not increment counter
\begin{longtable}[]{@{}cc@{}}
\toprule\noalign{}
符号 & 含义 \\
\midrule\noalign{}
\endhead
\bottomrule\noalign{}
\endlastfoot
\(+, -, \times (\cdot), \div (/)\) & 四则运算 \\
\(\pm, \mp\) & 正负/负正 \\
\(\sqrt{\ },\ \sqrt[n]{\ }\) & 平方根/\(n\) 次方根 \\
\(\|x\|\) & 绝对值（实数）或模（复数） \\
\(n!\) & 阶乘：\(n! = 1 \times 2 \times \cdots \times n\)，约定
\(0! = 1\) \\
\(\binom{n}{k}\) 或 \(C_n^k\) & 组合数：\(\frac{n!}{k!(n-k)!}\) \\
\(P_n^k\) 或 \(A_n^k\) & 排列数：\(\frac{n!}{(n-k)!}\) \\
\end{longtable}
}

\bigskip\hrule\bigskip

\subsubsection{1.6
函数与分析符号}\label{ux51fdux6570ux4e0eux5206ux6790ux7b26ux53f7}

\textbf{定义 1.6.1}：以下符号用于函数及其分析性质。

{\def\LTcaptype{none} % do not increment counter
\begin{longtable}[]{@{}cc@{}}
\toprule\noalign{}
符号 & 含义 \\
\midrule\noalign{}
\endhead
\bottomrule\noalign{}
\endlastfoot
\(f: A \to B\) & 从 \(A\) 到 \(B\) 的函数 \\
\(f^{-1}\) & 反函数（注意：不是 \(\frac{1}{f}\)） \\
\(f \circ g\) & 复合函数：\((f \circ g)(x) = f(g(x))\) \\
\(f'(x)\) 或 \(\frac{df}{dx}\) & 导数 \\
\(f''(x)\) 或 \(\frac{d^2f}{dx^2}\) & 二阶导数 \\
\(f^{(n)}(x)\) 或 \(\frac{d^nf}{dx^n}\) & \(n\) 阶导数 \\
\(\int_a^b f(x)\,dx\) & 定积分 \\
\(\int f(x)\,dx\) & 不定积分（原函数） \\
\(\sum_{i=1}^n a_i\) & 求和：\(a_1 + a_2 + \cdots + a_n\) \\
\(\prod_{i=1}^n a_i\) & 求积：\(a_1 \cdot a_2 \cdots a_n\) \\
\(\lim_{x \to a} f(x)\) & 极限 \\
\(\lim_{x \to \infty} f(x)\) & 无穷远处的极限 \\
\(\sup S,\ \inf S\) & 上确界/下确界 \\
\(\max S,\ \min S\) & 最大值/最小值 \\
\end{longtable}
}

\bigskip\hrule\bigskip

\subsubsection{1.7
三角函数符号}\label{ux4e09ux89d2ux51fdux6570ux7b26ux53f7}

\textbf{定义
1.7.1}：\(\sin x,\ \cos x,\ \tan x,\ \cot x,\ \sec x,\ \csc x\)
分别表示正弦、余弦、正切、余切、正割、余割函数。其反函数记为
\(\arcsin x,\ \arccos x,\ \arctan x\) 等。

\textbf{注记
1.7.1}：三角函数在微积分中起着核心作用。以下恒等式将反复使用： -
勾股恒等式：\(\sin^2 x + \cos^2 x = 1\) -
和角公式：\(\sin(x+y) = \sin x \cos y + \cos x \sin y\)；\(\cos(x+y) = \cos x \cos y - \sin x \sin y\)
-
倍角公式：\(\sin 2x = 2\sin x \cos x\)；\(\cos 2x = \cos^2 x - \sin^2 x = 2\cos^2 x - 1 = 1 - 2\sin^2 x\)
-
欧拉公式：\(e^{ix} = \cos x + i\sin x\)（将在卷九（复分析）中严格证明）

\bigskip\hrule\bigskip

\subsubsection{1.8 几何符号}\label{ux51e0ux4f55ux7b26ux53f7}

\textbf{定义 1.8.1}：以下符号用于几何对象。

{\def\LTcaptype{none} % do not increment counter
\begin{longtable}[]{@{}cc@{}}
\toprule\noalign{}
符号 & 含义 \\
\midrule\noalign{}
\endhead
\bottomrule\noalign{}
\endlastfoot
\(\overline{AB}\) 或 \(AB\) & 线段 \\
\(\overrightarrow{AB}\) & 向量 \\
\(\|AB\|\) 或 \(AB\) & 线段 \(AB\) 的长度 \\
\(\angle ABC\) & 以 \(B\) 为顶点的角 \\
\(\triangle ABC\) & 三角形 \\
\(\parallel\) & 平行 \\
\(\perp\) & 垂直 \\
\(\odot O\) & 以 \(O\) 为圆心的圆 \\
\([ABCD]\) & 四边形 \(ABCD\) 的面积 \\
\end{longtable}
}

\textbf{注记 1.8.1}：在欧几里得几何中，我们默认使用标准的公理体系（如
Hilbert 公理系统）。角度通常用弧度制表示（\(\pi\) 弧度
\(= 180°\)），这在微积分中更为自然。

\bigskip\hrule\bigskip

\subsubsection{1.9 希腊字母表}\label{ux5e0cux814aux5b57ux6bcdux8868}

数学中大量使用希腊字母作为变量和常量名。以下列出常用的希腊字母及其常见用法：

{\def\LTcaptype{none} % do not increment counter
\begin{longtable}[]{@{}ccc@{}}
\toprule\noalign{}
大写 & 小写 & 常见用法 \\
\midrule\noalign{}
\endhead
\bottomrule\noalign{}
\endlastfoot
\(\Delta\) & \(\delta\) & 判别式 \(\Delta = b^2 - 4ac\)；变化量 \\
\(\Sigma\) & \(\sigma\) & 求和；标准差 \\
\(\Pi\) & \(\pi\) & 圆周率 \(\pi \approx 3.14159\) \\
\(\Omega\) & \(\omega\) & 角速度 \\
\(\Phi\) & \(\phi\) & 黄金分割比 \(\phi = \frac{1+\sqrt{5}}{2}\) \\
\(\Lambda\) & \(\lambda\) & 特征值 \\
\(\Mu\) & \(\mu\) & 期望、均值 \\
\(\Epsilon\) & \(\epsilon\) & 任意小的正数（极限语言） \\
\(\Theta\) & \(\theta\) & 角度 \\
\(\Gamma\) & \(\gamma\) & 欧拉常数 \\
\(\Alpha\) & \(\alpha\) & 常用角度 \\
\(\Beta\) & \(\beta\) & 常用角度 \\
\(\Nu\) & \(\nu\) & 频率 \\
\(\Rho\) & \(\rho\) & 密度、相关系数 \\
\end{longtable}
}

\textbf{注记 1.9.1}：大写希腊字母 \(\Sigma\) 常用于求和符号
\(\sum\)，\(\Pi\) 常用于求积符号 \(\prod\)，\(\Delta\)
常用于变化量或判别式。小写 \(\epsilon\)
在极限理论中表示''任意小的正数''，\(\delta\)
表示''\(\epsilon\)-\(\delta\) 论证''中的距离参数。

\bigskip\hrule\bigskip

\subsubsection{1.10
证明中的常用符号}\label{ux8bc1ux660eux4e2dux7684ux5e38ux7528ux7b26ux53f7}

\textbf{定义 1.10.1}：以下符号在证明中频繁使用。

{\def\LTcaptype{none} % do not increment counter
\begin{longtable}[]{@{}cc@{}}
\toprule\noalign{}
符号 & 含义 \\
\midrule\noalign{}
\endhead
\bottomrule\noalign{}
\endlastfoot
\(\blacksquare\) & 证毕标记，表示证明结束 \\
\(\Rightarrow\) & 推出 \\
\(\Leftarrow\) & 被推出 \\
\(\Leftrightarrow\) & 等价于 \\
\(\stackrel{\text{def}}{=}\) & 按定义等于 \\
\(a \mid b\) & \(a\) 整除 \(b\) \\
\(a \nmid b\) & \(a\) 不整除 \(b\) \\
\(a \equiv b \pmod{m}\) & \(a\) 与 \(b\) 关于模 \(m\) 同余 \\
\(P \text{ w.r.t. } Q\) & \(P\) 关于 \(Q\)（with respect to） \\
\(f = O(g)\) & \(f\) 是 \(g\) 的大 \(O\)（渐近上界） \\
\(\mathbb{E}[X]\) & 随机变量 \(X\) 的期望 \\
\(\mathrm{Var}(X)\) & 随机变量 \(X\) 的方差 \\
\end{longtable}
}

\textbf{注记 1.10.1}：在证明中，\(\blacksquare\)
表示一个证明的结束。如果证明中包含多个部分，每个部分结束时使用
\(\blacksquare\)。符号 \(\Rightarrow\) 和 \(\Leftarrow\)
用于逻辑推理步骤之间，而不是命题公式内部（命题公式内部使用
\(\rightarrow\) 和 \(\leftarrow\)）。\(\stackrel{\text{def}}{=}\)
强调等式是根据定义得到的。

\bigskip\hrule\bigskip

\subsubsection{本章展望}\label{ux672cux7ae0ux5c55ux671b}

本章建立了全书的符号语言体系。在 Ch 2
中，我们将使用这些逻辑符号来系统学习命题逻辑和谓词逻辑，建立严格推理的基础。在
Ch 3
中，集合符号将被用来精确描述数学对象。后续每一章都将依赖本章所建立的符号约定。

\bigskip\hrule\bigskip

\subsection{Ch 2 逻辑基础}\label{ch-2-ux903bux8f91ux57faux7840}

\textbf{前置知识}：依赖 Ch 1（数学符号），特别是 1.1
节中的逻辑符号（\(\neg, \wedge, \vee, \rightarrow, \leftrightarrow, \forall, \exists\)）。

数学的力量来自于严格的推理，而逻辑是所有数学推理的基础。在 Ch 1
中我们认识了逻辑符号，现在我们将深入理解这些符号背后的理论体系------命题逻辑和谓词逻辑，并学习基本的证明方法。

\bigskip\hrule\bigskip

\subsubsection{2.1 命题逻辑}\label{ux547dux9898ux903bux8f91}

\paragraph{2.1.1 命题}\label{ux547dux9898}

\textbf{定义
2.1.1（命题）}：一个\textbf{命题}是一个陈述句，它要么为真（记为 \(T\) 或
\(\mathbf{1}\)），要么为假（记为 \(F\) 或
\(\mathbf{0}\)），但不能同时为真又为假。

我们用大写字母 \(P, Q, R, \ldots\) 来表示命题。

\paragraph{2.1.2 逻辑联结词}\label{ux903bux8f91ux8054ux7ed3ux8bcd}

\textbf{定义 2.1.2（否定）}：设 \(P\) 是一个命题，则
\(\neg P\)（读作''非 \(P\)``）也是一个命题，其真值与 \(P\) 相反：
\[\neg T = F, \quad \neg F = T\]

\textbf{定义 2.1.3（合取）}：设 \(P, Q\) 是命题，则
\(P \wedge Q\)（读作''\(P\) 且 \(Q\)``）定义为：
\[P \wedge Q = \begin{cases} T & \text{当且仅当 } P = T \text{ 且 } Q = T \\ F & \text{其他情况} \end{cases}\]

\textbf{定义 2.1.4（析取）}：设 \(P, Q\) 是命题，则
\(P \vee Q\)（读作''\(P\) 或 \(Q\)``）定义为（可兼或）：
\[P \vee Q = \begin{cases} T & \text{当且仅当 } P = T \text{ 或 } Q = T \text{（或两者都为真）} \\ F & \text{当且仅当 } P = F \text{ 且 } Q = F \end{cases}\]

\textbf{定义 2.1.5（蕴含）}：设 \(P, Q\) 是命题，则
\(P \rightarrow Q\)（读作''若 \(P\) 则 \(Q\)``）定义为：
\[P \rightarrow Q = \begin{cases} F & \text{当且仅当 } P = T \text{ 且 } Q = F \\ T & \text{其他情况} \end{cases}\]

\textbf{注记 2.1.1}：蕴含 \(P \rightarrow Q\)
仅在''前件为真而后件为假''时才为假。当 \(P\) 为假时，\(P \rightarrow Q\)
自动为真（称为''空真''或''实质蕴含''）。

\textbf{定义 2.1.6（等价）}：设 \(P, Q\) 是命题，则
\(P \leftrightarrow Q\)（读作''\(P\) 当且仅当 \(Q\)``）定义为：
\[P \leftrightarrow Q = \begin{cases} T & \text{当且仅当 } P \text{ 和 } Q \text{ 真值相同} \\ F & \text{当且仅当 } P \text{ 和 } Q \text{ 真值不同} \end{cases}\]

\paragraph{2.1.3 真值表}\label{ux771fux503cux8868}

\textbf{定义 2.1.7（真值表）}：将所有联结词的定义总结为一张表：

{\def\LTcaptype{none} % do not increment counter
\begin{longtable}[]{@{}
  >{\centering\arraybackslash}p{(\linewidth - 12\tabcolsep) * \real{0.0562}}
  >{\centering\arraybackslash}p{(\linewidth - 12\tabcolsep) * \real{0.0562}}
  >{\centering\arraybackslash}p{(\linewidth - 12\tabcolsep) * \real{0.1124}}
  >{\centering\arraybackslash}p{(\linewidth - 12\tabcolsep) * \real{0.1573}}
  >{\centering\arraybackslash}p{(\linewidth - 12\tabcolsep) * \real{0.1348}}
  >{\centering\arraybackslash}p{(\linewidth - 12\tabcolsep) * \real{0.2247}}
  >{\centering\arraybackslash}p{(\linewidth - 12\tabcolsep) * \real{0.2584}}@{}}
\toprule\noalign{}
\begin{minipage}[b]{\linewidth}\centering
\(P\)
\end{minipage} & \begin{minipage}[b]{\linewidth}\centering
\(Q\)
\end{minipage} & \begin{minipage}[b]{\linewidth}\centering
\(\neg P\)
\end{minipage} & \begin{minipage}[b]{\linewidth}\centering
\(P \wedge Q\)
\end{minipage} & \begin{minipage}[b]{\linewidth}\centering
\(P \vee Q\)
\end{minipage} & \begin{minipage}[b]{\linewidth}\centering
\(P \rightarrow Q\)
\end{minipage} & \begin{minipage}[b]{\linewidth}\centering
\(P \leftrightarrow Q\)
\end{minipage} \\
\midrule\noalign{}
\endhead
\bottomrule\noalign{}
\endlastfoot
T & T & F & T & T & T & T \\
T & F & F & F & T & F & F \\
F & T & T & F & T & T & F \\
F & F & T & F & F & T & T \\
\end{longtable}
}

\paragraph{2.1.4 逻辑等价}\label{ux903bux8f91ux7b49ux4ef7}

\textbf{定义 2.1.8（逻辑等价）}：设 \(\varphi\) 和 \(\psi\)
是由命题变元和逻辑联结词构成的\textbf{命题公式}。如果对于所有命题变元的真值指派，\(\varphi\)
和 \(\psi\) 的真值总是相同，则称 \(\varphi\) 和 \(\psi\)
\textbf{逻辑等价}，记为 \(\varphi \equiv \psi\)。

\textbf{定理 2.1.1（双重否定律）}：对于任意命题 \(P\)，
\[\neg(\neg P) \equiv P\]

\textbf{证明}：构造真值表：

{\def\LTcaptype{none} % do not increment counter
\begin{longtable}[]{@{}ccc@{}}
\toprule\noalign{}
\(P\) & \(\neg P\) & \(\neg(\neg P)\) \\
\midrule\noalign{}
\endhead
\bottomrule\noalign{}
\endlastfoot
T & F & T \\
F & T & F \\
\end{longtable}
}

\(\neg(\neg P)\) 的真值列与 \(P\) 完全相同，故
\(\neg(\neg P) \equiv P\)。\(\blacksquare\)

\textbf{定理 2.1.2（交换律）}：对于任意命题 \(P, Q\)： 1.
\(P \wedge Q \equiv Q \wedge P\) 2. \(P \vee Q \equiv Q \vee P\)

\textbf{证明}：由真值表直接验证。以合取为例：

{\def\LTcaptype{none} % do not increment counter
\begin{longtable}[]{@{}cccc@{}}
\toprule\noalign{}
\(P\) & \(Q\) & \(P \wedge Q\) & \(Q \wedge P\) \\
\midrule\noalign{}
\endhead
\bottomrule\noalign{}
\endlastfoot
T & T & T & T \\
T & F & F & F \\
F & T & F & F \\
F & F & F & F \\
\end{longtable}
}

两列完全相同。析取的验证类似。\(\blacksquare\)

\textbf{定理 2.1.3（结合律）}：对于任意命题 \(P, Q, R\)： 1.
\((P \wedge Q) \wedge R \equiv P \wedge (Q \wedge R)\) 2.
\((P \vee Q) \vee R \equiv P \vee (Q \vee R)\)

\textbf{证明}：构造 8
行真值表逐一验证，两对应列完全相同。\(\blacksquare\)

\textbf{定理 2.1.4（分配律）}：对于任意命题 \(P, Q, R\)： 1.
\(P \wedge (Q \vee R) \equiv (P \wedge Q) \vee (P \wedge R)\) 2.
\(P \vee (Q \wedge R) \equiv (P \vee Q) \wedge (P \vee R)\)

\textbf{证明}：构造 8 行真值表验证。以第一式为例：

{\def\LTcaptype{none} % do not increment counter
\begin{longtable}[]{@{}
  >{\centering\arraybackslash}p{(\linewidth - 14\tabcolsep) * \real{0.0450}}
  >{\centering\arraybackslash}p{(\linewidth - 14\tabcolsep) * \real{0.0450}}
  >{\centering\arraybackslash}p{(\linewidth - 14\tabcolsep) * \real{0.0450}}
  >{\centering\arraybackslash}p{(\linewidth - 14\tabcolsep) * \real{0.1081}}
  >{\centering\arraybackslash}p{(\linewidth - 14\tabcolsep) * \real{0.2072}}
  >{\centering\arraybackslash}p{(\linewidth - 14\tabcolsep) * \real{0.1261}}
  >{\centering\arraybackslash}p{(\linewidth - 14\tabcolsep) * \real{0.1261}}
  >{\centering\arraybackslash}p{(\linewidth - 14\tabcolsep) * \real{0.2973}}@{}}
\toprule\noalign{}
\begin{minipage}[b]{\linewidth}\centering
\(P\)
\end{minipage} & \begin{minipage}[b]{\linewidth}\centering
\(Q\)
\end{minipage} & \begin{minipage}[b]{\linewidth}\centering
\(R\)
\end{minipage} & \begin{minipage}[b]{\linewidth}\centering
\(Q \vee R\)
\end{minipage} & \begin{minipage}[b]{\linewidth}\centering
\(P \wedge (Q \vee R)\)
\end{minipage} & \begin{minipage}[b]{\linewidth}\centering
\(P \wedge Q\)
\end{minipage} & \begin{minipage}[b]{\linewidth}\centering
\(P \wedge R\)
\end{minipage} & \begin{minipage}[b]{\linewidth}\centering
\((P \wedge Q) \vee (P \wedge R)\)
\end{minipage} \\
\midrule\noalign{}
\endhead
\bottomrule\noalign{}
\endlastfoot
T & T & T & T & T & T & T & T \\
T & T & F & T & T & T & F & T \\
T & F & T & T & T & F & T & T \\
T & F & F & F & F & F & F & F \\
F & T & T & T & F & F & F & F \\
F & T & F & T & F & F & F & F \\
F & F & T & T & F & F & F & F \\
F & F & F & F & F & F & F & F \\
\end{longtable}
}

第五列与最后一列相同。\(\blacksquare\)

\textbf{定理 2.1.5（德摩根律）}：对于任意命题 \(P, Q\)： 1.
\(\neg(P \wedge Q) \equiv \neg P \vee \neg Q\) 2.
\(\neg(P \vee Q) \equiv \neg P \wedge \neg Q\)

\textbf{证明}：由真值表直接验证。\(\blacksquare\)

\textbf{定理 2.1.6（吸收律）}：对于任意命题 \(P, Q\)： 1.
\(P \vee (P \wedge Q) \equiv P\) 2. \(P \wedge (P \vee Q) \equiv P\)

\textbf{证明}：由真值表验证。\(\blacksquare\)

\paragraph{2.1.5
重言式与矛盾式}\label{ux91cdux8a00ux5f0fux4e0eux77dbux76feux5f0f}

\textbf{定义 2.1.9（重言式）}：一个命题公式 \(\varphi\)
称为\textbf{重言式}（或\textbf{永真式}），如果对于命题变元的所有真值指派，\(\varphi\)
的值恒为真。

\textbf{定义 2.1.10（矛盾式）}：一个命题公式 \(\varphi\)
称为\textbf{矛盾式}（或\textbf{永假式}），如果对于命题变元的所有真值指派，\(\varphi\)
的值恒为假。

\textbf{定理 2.1.7（排中律）}：\(P \vee \neg P\) 是重言式。

\textbf{证明}：

{\def\LTcaptype{none} % do not increment counter
\begin{longtable}[]{@{}ccc@{}}
\toprule\noalign{}
\(P\) & \(\neg P\) & \(P \vee \neg P\) \\
\midrule\noalign{}
\endhead
\bottomrule\noalign{}
\endlastfoot
T & F & T \\
F & T & T \\
\end{longtable}
}

无论 \(P\) 取何值，\(P \vee \neg P\) 恒为真。\(\blacksquare\)

\textbf{定理 2.1.8（矛盾律）}：\(P \wedge \neg P\) 是矛盾式。

\textbf{证明}：

{\def\LTcaptype{none} % do not increment counter
\begin{longtable}[]{@{}ccc@{}}
\toprule\noalign{}
\(P\) & \(\neg P\) & \(P \wedge \neg P\) \\
\midrule\noalign{}
\endhead
\bottomrule\noalign{}
\endlastfoot
T & F & F \\
F & T & F \\
\end{longtable}
}

无论 \(P\) 取何值，\(P \wedge \neg P\) 恒为假。\(\blacksquare\)

\textbf{定理 2.1.9（假言三段论/蕴含的传递性）}：对于任意命题
\(P, Q, R\)：
\[((P \rightarrow Q) \wedge (Q \rightarrow R)) \rightarrow (P \rightarrow R)\]
是重言式。

\textbf{证明}：在所有 8 种真值指派下，验证整个公式恒为真。当前件
\((P \rightarrow Q) \wedge (Q \rightarrow R)\)
为真时，\(P \rightarrow Q\) 和 \(Q \rightarrow R\) 都为真；此时若 \(P\)
为真则 \(Q\) 为真从而 \(R\) 为真，故 \(P \rightarrow R\) 为真；若 \(P\)
为假则 \(P \rightarrow R\)
自动为真。当前件为假时，蕴含式自动为真。因此最后一列全为
\(T\)。\(\blacksquare\)

\textbf{定理
2.1.10（逆否命题等价律）}：\(P \rightarrow Q \equiv \neg Q \rightarrow \neg P\)。

\textbf{证明}：

{\def\LTcaptype{none} % do not increment counter
\begin{longtable}[]{@{}
  >{\centering\arraybackslash}p{(\linewidth - 10\tabcolsep) * \real{0.0641}}
  >{\centering\arraybackslash}p{(\linewidth - 10\tabcolsep) * \real{0.0641}}
  >{\centering\arraybackslash}p{(\linewidth - 10\tabcolsep) * \real{0.1282}}
  >{\centering\arraybackslash}p{(\linewidth - 10\tabcolsep) * \real{0.1282}}
  >{\centering\arraybackslash}p{(\linewidth - 10\tabcolsep) * \real{0.2436}}
  >{\centering\arraybackslash}p{(\linewidth - 10\tabcolsep) * \real{0.3718}}@{}}
\toprule\noalign{}
\begin{minipage}[b]{\linewidth}\centering
\(P\)
\end{minipage} & \begin{minipage}[b]{\linewidth}\centering
\(Q\)
\end{minipage} & \begin{minipage}[b]{\linewidth}\centering
\(\neg P\)
\end{minipage} & \begin{minipage}[b]{\linewidth}\centering
\(\neg Q\)
\end{minipage} & \begin{minipage}[b]{\linewidth}\centering
\(P \rightarrow Q\)
\end{minipage} & \begin{minipage}[b]{\linewidth}\centering
\(\neg Q \rightarrow \neg P\)
\end{minipage} \\
\midrule\noalign{}
\endhead
\bottomrule\noalign{}
\endlastfoot
T & T & F & F & T & T \\
T & F & F & T & F & F \\
F & T & T & F & T & T \\
F & F & T & T & T & T \\
\end{longtable}
}

第五列与第六列相同。\(\blacksquare\)

\textbf{定理
2.1.11（蕴含的展开律）}：\(P \rightarrow Q \equiv \neg P \vee Q\)。

\textbf{证明}：

{\def\LTcaptype{none} % do not increment counter
\begin{longtable}[]{@{}ccccc@{}}
\toprule\noalign{}
\(P\) & \(Q\) & \(\neg P\) & \(\neg P \vee Q\) & \(P \rightarrow Q\) \\
\midrule\noalign{}
\endhead
\bottomrule\noalign{}
\endlastfoot
T & T & F & T & T \\
T & F & F & F & F \\
F & T & T & T & T \\
F & F & T & T & T \\
\end{longtable}
}

第四列与第五列相同。\(\blacksquare\)

\textbf{推论
2.1.2（与非门和或非门）}：所有逻辑联结词可以用单一的''与非''（NAND）或''或非''（NOR）来表达。\(\neg P \equiv P \uparrow P\)（NAND），\(P \wedge Q \equiv (P \uparrow Q) \uparrow (P \uparrow Q)\)。

\textbf{证明}：\(P \uparrow Q \equiv \neg(P \wedge Q)\)。\(P \uparrow P \equiv \neg(P \wedge P) \equiv \neg P\)。\((P \uparrow Q) \uparrow (P \uparrow Q) \equiv \neg(\neg(P \wedge Q) \wedge \neg(P \wedge Q)) \equiv \neg\neg(P \wedge Q) \equiv P \wedge Q\)。\(\blacksquare\)

\bigskip\hrule\bigskip

\subsubsection{2.2 谓词逻辑}\label{ux8c13ux8bcdux903bux8f91}

命题逻辑的表达能力有限。例如，``对所有实数
\(x\)，\(x^2 \geq 0\)''无法用单个命题符号来表达。为此，我们引入\textbf{谓词逻辑}。

\paragraph{2.2.1 个体与谓词}\label{ux4e2aux4f53ux4e0eux8c13ux8bcd}

\textbf{定义
2.2.1（个体）}：\textbf{个体}是讨论的对象，可以是具体的数，也可以是抽象的变元。

\textbf{定义
2.2.2（谓词）}：\textbf{谓词}是含有变元的陈述，当变元取确定值时，谓词就成为一个命题。用
\(P(x)\) 表示含有变元 \(x\) 的谓词。

\paragraph{2.2.2 量词}\label{ux91cfux8bcd}

\textbf{定义 2.2.3（全称量词）}：设 \(P(x)\) 是含有变元 \(x\)
的谓词，\(D\) 是论域，则 \(\forall x\, P(x)\)（读作''对所有
\(x\)，\(P(x)\) 成立''）定义为：
\[\forall x\, P(x) = \begin{cases} T & \text{若对 } D \text{ 中的每一个 } a,\, P(a) \text{ 都为真} \\ F & \text{若存在 } D \text{ 中的某个 } a,\, P(a) \text{ 为假} \end{cases}\]

\textbf{定义 2.2.4（存在量词）}：\(\exists x\, P(x)\)（读作''存在
\(x\)，使得 \(P(x)\) 成立''）定义为：
\[\exists x\, P(x) = \begin{cases} T & \text{若存在 } D \text{ 中的某个 } a,\, P(a) \text{ 为真} \\ F & \text{若对 } D \text{ 中的每一个 } a,\, P(a) \text{ 都为假} \end{cases}\]

\textbf{定义 2.2.5（存在唯一量词）}：\(\exists!\, x\, P(x)\) 表示：
\[\exists x\, (P(x) \wedge \forall y\, (P(y) \rightarrow y = x))\]

\textbf{注记
2.2.1}：量词的顺序至关重要。\(\forall x\, \exists y\, P(x, y)\)（对每个
\(x\) 存在 \(y\)）与 \(\exists y\, \forall x\, P(x, y)\)（存在 \(y\)
对所有
\(x\)）含义不同。后者是更强的命题。例如：\(\forall x \in \mathbb{R}\, \exists y \in \mathbb{R}\, (y > x)\)
为真（取 \(y = x + 1\)），但
\(\exists y \in \mathbb{R}\, \forall x \in \mathbb{R}\, (y > x)\)
为假（不存在比所有实数都大的实数）。

\paragraph{2.2.3 量词的否定}\label{ux91cfux8bcdux7684ux5426ux5b9a}

\textbf{定理 2.2.1（量词否定的对偶律）}：设 \(P(x)\) 是一个谓词，则： 1.
\(\neg(\forall x\, P(x)) \equiv \exists x\, \neg P(x)\) 2.
\(\neg(\exists x\, P(x)) \equiv \forall x\, \neg P(x)\)

\textbf{证明}：

\textbf{（1）的证明}：

\((\Rightarrow)\)：假设 \(\neg(\forall x\, P(x))\) 为真。这意味着
\(\forall x\, P(x)\) 为假，即存在论域中的某个 \(a\)，使得 \(P(a)\)
为假，即 \(\neg P(a)\) 为真。因此 \(\exists x\, \neg P(x)\) 为真。

\((\Leftarrow)\)：假设 \(\exists x\, \neg P(x)\)
为真。这意味着存在论域中的某个 \(a\)，使得 \(\neg P(a)\) 为真，即
\(P(a)\) 为假。因此 \(\forall x\, P(x)\) 为假，即
\(\neg(\forall x\, P(x))\) 为真。

综合两个方向，\(\neg(\forall x\, P(x)) \equiv \exists x\, \neg P(x)\)。\(\blacksquare\)

\textbf{（2）的证明}：类似地。

\((\Rightarrow)\)：假设 \(\neg(\exists x\, P(x))\) 为真。这意味着
\(\exists x\, P(x)\) 为假，即论域中不存在 \(a\) 使得 \(P(a)\)
为真，亦即论域中每一个 \(a\) 都使得 \(P(a)\) 为假，即 \(\neg P(a)\)
为真。因此 \(\forall x\, \neg P(x)\) 为真。

\((\Leftarrow)\)：假设 \(\forall x\, \neg P(x)\) 为真，即论域中每个
\(a\) 都使得 \(\neg P(a)\) 为真。这意味着不存在 \(a\) 使得 \(P(a)\)
为真，故 \(\exists x\, P(x)\) 为假，即 \(\neg(\exists x\, P(x))\)
为真。\(\blacksquare\)

\textbf{推论 2.2.1（嵌套量词的否定）}： 1.
\(\neg(\forall x\, \forall y\, P(x, y)) \equiv \exists x\, \exists y\, \neg P(x, y)\)
2.
\(\neg(\exists x\, \exists y\, P(x, y)) \equiv \forall x\, \forall y\, \neg P(x, y)\)
3.
\(\neg(\forall x\, \exists y\, P(x, y)) \equiv \exists x\, \forall y\, \neg P(x, y)\)
4.
\(\neg(\exists x\, \forall y\, P(x, y)) \equiv \forall x\, \exists y\, \neg P(x, y)\)

\textbf{证明}：反复应用定理 2.2.1。以 (3) 为例：
\[\neg(\forall x\, \exists y\, P(x, y)) \equiv \exists x\, \neg(\exists y\, P(x, y)) \equiv \exists x\, \forall y\, \neg P(x, y)\]
\(\blacksquare\)

\bigskip\hrule\bigskip

\subsubsection{2.3 四种命题}\label{ux56dbux79cdux547dux9898}

\textbf{定义 2.3.1}：给定命题 \(P \rightarrow Q\)： -
\(P \rightarrow Q\) 称为\textbf{原命题}； - \(Q \rightarrow P\)
称为\textbf{逆命题}； - \(\neg P \rightarrow \neg Q\)
称为\textbf{否命题}； - \(\neg Q \rightarrow \neg P\)
称为\textbf{逆否命题}。

\textbf{定理 2.3.1}：原命题与逆否命题等价，即
\(P \rightarrow Q \equiv \neg Q \rightarrow \neg P\)。

\textbf{证明}：这正是定理 2.1.10 的结论。\(\blacksquare\)

\textbf{定理 2.3.2}：逆命题与否命题等价，即
\(Q \rightarrow P \equiv \neg P \rightarrow \neg Q\)。

\textbf{证明}：在定理 2.1.10 中将 \(P\) 和 \(Q\)
互换即得。\(\blacksquare\)

\textbf{注记 2.3.1}：原命题为真不能推出逆命题为真。例如''若 \(x > 0\) 则
\(x^2 > 0\)``为真，但''若 \(x^2 > 0\) 则 \(x > 0\)``为假（\(x = -1\)
是反例）。原命题为真也不能推出否命题为真。但原命题与逆否命题等价，逆命题与否命题等价。因此，要证明一个蕴含式，可以证明其逆否命题（常常更方便）。

\bigskip\hrule\bigskip

\subsubsection{2.4
充分条件与必要条件}\label{ux5145ux5206ux6761ux4ef6ux4e0eux5fc5ux8981ux6761ux4ef6}

\textbf{定义 2.4.1}：设 \(P, Q\) 是命题。 - 若 \(P \rightarrow Q\)
为真，则称 \(P\) 是 \(Q\) 的\textbf{充分条件}（\(P\) 足以保证
\(Q\)），\(Q\) 是 \(P\) 的\textbf{必要条件}（\(Q\) 是 \(P\)
成立所必需的）。 - 若 \(P \leftrightarrow Q\) 为真，则称 \(P\) 是 \(Q\)
的\textbf{充要条件}。

\textbf{定理 2.4.1}：\(P\) 是 \(Q\) 的充分条件当且仅当 \(Q\) 是 \(P\)
的必要条件。

\textbf{证明}：\(P \rightarrow Q\) 与 \(\neg Q \rightarrow \neg P\)
等价（定理 2.3.1），即''\(P\) 蕴含 \(Q\)``等价于''\(Q\) 的否定蕴含 \(P\)
的否定''。\(\blacksquare\)

\textbf{推论 2.4.1}：\(P\) 是 \(Q\) 的充要条件当且仅当 \(Q\) 是 \(P\)
的充要条件。

\textbf{证明}：\(P \leftrightarrow Q\) 与 \(Q \leftrightarrow P\)
等价（等价关系的对称性）。\(\blacksquare\)

\textbf{注记
2.4.1（充要条件的数学表达）}：数学中常用以下方式表达充要条件： - ``\(P\)
当且仅当 \(Q\)''（\(P\) iff \(Q\)） - ``\(P\) 的充要条件是 \(Q\)'' -
``\(P\) 等价于 \(Q\)''

例如：``\(n\) 是偶数的充要条件是 \(n^2\) 是偶数''意味着 \(n\) 是偶数
\(\iff\) \(n^2\) 是偶数。

\bigskip\hrule\bigskip

\subsubsection{2.5
全称命题与存在命题}\label{ux5168ux79f0ux547dux9898ux4e0eux5b58ux5728ux547dux9898}

\textbf{定义 2.5.1（全称命题）}：形如 \(\forall x \in D,\, P(x)\)
的命题称为\textbf{全称命题}。要证明它，需对 \(D\) 中每个 \(x\) 都验证
\(P(x)\)；要否定它，只需找到一个反例。

\textbf{定义 2.5.2（存在命题）}：形如 \(\exists x \in D,\, P(x)\)
的命题称为\textbf{存在命题}。要证明它，只需找到一个满足条件的
\(x\)（构造性证明）；要否定它，需对所有 \(x\) 都验证 \(P(x)\) 不成立。

\bigskip\hrule\bigskip

\subsubsection{2.6 证明方法}\label{ux8bc1ux660eux65b9ux6cd5}

\paragraph{2.6.1 直接证明}\label{ux76f4ux63a5ux8bc1ux660e}

\textbf{方法}：要证明 \(P \rightarrow Q\)，假设 \(P\)
为真，然后通过一系列逻辑推导得出 \(Q\) 为真。

\textbf{定理 2.6.1}：若 \(n\) 是偶数，则 \(n^2\) 是偶数。

\textbf{证明}：设 \(n\) 是偶数，则存在整数 \(k\) 使得 \(n = 2k\)。于是
\[n^2 = (2k)^2 = 4k^2 = 2(2k^2)\] 由于 \(2k^2\) 是整数，故 \(n^2\)
是偶数。\(\blacksquare\)

\paragraph{2.6.2 反证法}\label{ux53cdux8bc1ux6cd5}

\textbf{方法}：要证明命题 \(P\)，假设 \(\neg P\)
为真，然后推导出矛盾，从而 \(P\) 必为真。

\textbf{逻辑依据}：排中律（定理 2.1.7）保证 \(P \vee \neg P\) 恒为真。若
\(\neg P\) 导致矛盾，则 \(\neg P\) 不可能为真，从而 \(P\) 必为真。

\textbf{定理 2.6.2}：\(\sqrt{2}\) 不是有理数。

\textbf{证明}：假设 \(\sqrt{2}\) 是有理数，则存在正整数
\(p, q\)（\(q \neq 0\)），使得 \(\sqrt{2} = \frac{p}{q}\)，且
\(\gcd(p, q) = 1\)。

由 \(\sqrt{2} = \frac{p}{q}\) 得 \(2 = \frac{p^2}{q^2}\)，即
\(p^2 = 2q^2\)。

因此 \(p^2\) 是偶数。由于奇数的平方是奇数，\(p\) 必须是偶数。设
\(p = 2m\)，代入得 \((2m)^2 = 2q^2\)，即 \(4m^2 = 2q^2\)，故
\(q^2 = 2m^2\)。

于是 \(q^2\) 也是偶数，同理 \(q\) 也是偶数。但这意味着
\(\gcd(p, q) \geq 2\)，与 \(\gcd(p, q) = 1\) 矛盾。

因此假设不成立，\(\sqrt{2}\) 不是有理数。\(\blacksquare\)

\paragraph{2.6.3 数学归纳法}\label{ux6570ux5b66ux5f52ux7eb3ux6cd5}

\textbf{定理 2.6.3（数学归纳法原理）}：设 \(P(n)\) 是关于自然数 \(n\)
的命题。若以下两个条件成立： 1. \textbf{基础步骤}：\(P(0)\) 为真； 2.
\textbf{归纳步骤}：对任意自然数 \(k\)，若 \(P(k)\) 为真则 \(P(k+1)\)
也为真；

则 \(P(n)\) 对所有自然数 \(n\) 都为真。

\textbf{证明}：令
\(A = \{n \in \mathbb{N} \mid P(n) \text{ 为真}\}\)。由条件
(1)，\(0 \in A\)；由条件 (2)，若 \(k \in A\) 则
\(S(k) \in A\)。数学归纳法原理在此阶段作为自然数系的基本原理接受（其在公理化框架中的严格地位将在定义
4.1.1 的公理 P5 中明确），故 \(A = \mathbb{N}\)。\(\blacksquare\)

\textbf{注记 2.6.1}：数学归纳法可以从任意起始点 \(n_0\) 开始。若
\(P(n_0)\) 为真且 \(P(k) \to P(k+1)\) 对所有 \(k \geq n_0\) 成立，则
\(P(n)\) 对所有 \(n \geq n_0\) 为真。这等价于对 \(Q(n) = P(n + n_0)\)
使用标准归纳法。

\textbf{定理 2.6.3'（跳步归纳法）}：设 \(P(n)\) 是关于自然数 \(n\)
的命题。若以下条件成立： 1. \(P(0)\) 和 \(P(1)\) 都为真； 2.
对任意自然数 \(k\)，若 \(P(k)\) 为真则 \(P(k+2)\) 也为真；

则 \(P(n)\) 对所有自然数 \(n\) 都为真。

\textbf{证明}：对 \(n\) 的奇偶性分情况讨论。由条件 (1) 和 (2)
反复应用，\(P(0), P(2), P(4), \ldots\) 都为真，且
\(P(1), P(3), P(5), \ldots\) 都为真。\(\blacksquare\)

\textbf{定理 2.6.4（强归纳原理）}：设 \(P(n)\) 是关于自然数 \(n\)
的命题。若以下条件成立：对任意自然数 \(k\)，若
\(P(0), P(1), \ldots, P(k-1)\) \textbf{都}为真，则 \(P(k)\) 也为真；则
\(P(n)\) 对所有自然数 \(n\) 都为真。

\textbf{证明}：定义 \(Q(n)\)：``\(P(0), P(1), \ldots, P(n)\)
都为真''。对 \(Q(n)\) 使用普通归纳法。

\(Q(0)\) 为真：强归纳的假设（取 \(k = 0\)）中归纳假设空真，故 \(P(0)\)
为真，即 \(Q(0)\) 为真。

若 \(Q(k)\) 为真，即 \(P(0), \ldots, P(k)\) 都为真。由强归纳假设（取
\(k+1\)），\(P(k+1)\) 为真，故 \(Q(k+1)\) 为真。

由普通归纳法，\(Q(n)\) 对所有 \(n\) 成立，从而 \(P(n)\) 对所有 \(n\)
成立。\(\blacksquare\)

\paragraph{2.6.4 构造性证明}\label{ux6784ux9020ux6027ux8bc1ux660e}

\textbf{方法}：要证明 \(\exists x\, P(x)\)，构造出一个具体的
\(a\)，然后验证 \(P(a)\) 为真。

\paragraph{2.6.5 等价性证明}\label{ux7b49ux4ef7ux6027ux8bc1ux660e}

\textbf{方法}：要证明 \(P \leftrightarrow Q\)，分别证明
\(P \rightarrow Q\) 和 \(Q \rightarrow P\)。

\paragraph{2.6.6 分情况讨论}\label{ux5206ux60c5ux51b5ux8ba8ux8bba}

\textbf{方法}：若命题涉及的对象可分成有限种情况，对每种情况分别证明。

\textbf{定理 2.6.5}：\(\forall n \in \mathbb{N},\, n(n+1)\) 是偶数。

\textbf{证明}：\(n\) 要么是偶数，要么是奇数。

\textbf{情况 1}：\(n = 2k\)。则 \(n(n+1) = 2k(2k+1)\)，含因子
\(2\)，为偶数。

\textbf{情况 2}：\(n = 2k+1\)。则
\(n(n+1) = (2k+1)(2k+2) = 2(2k+1)(k+1)\)，含因子 \(2\)，为偶数。

两种情况下结论都成立。\(\blacksquare\)

\paragraph{2.6.7
存在性证明与唯一性证明}\label{ux5b58ux5728ux6027ux8bc1ux660eux4e0eux552fux4e00ux6027ux8bc1ux660e}

\textbf{存在性}：证明满足某性质的对象至少存在一个，可用构造性证明（直接构造）或非构造性证明（如反证法）。

\textbf{唯一性}：证明满足某性质的对象至多有一个，常用方法是假设有两个满足条件的对象，然后证明它们相等。

\textbf{定理 2.6.6（加法零元的唯一性）}：在满足加法结合律的系统中，若
\(e\) 和 \(e'\) 都是加法零元（即 \(\forall a: a + e = a\) 且
\(\forall a: a + e' = a\)），则 \(e = e'\)。

\textbf{证明}：\(e = e + e'\)（\(e'\) 是零元）\(= e'\)（\(e\)
是零元）。\(\blacksquare\)

\textbf{定理
2.6.7（加法逆元的唯一性）}：在满足加法结合律、有零元的系统中，若 \(b\)
和 \(b'\) 都是 \(a\) 的加法逆元（\(a + b = 0\) 且 \(a + b' = 0\)），则
\(b = b'\)。

\textbf{证明}：\(b = b + 0 = b + (a + b') = (b + a) + b' = 0 + b' = b'\)。\(\blacksquare\)

\textbf{注}：上述条件（加法结合律、零元、逆元）正是''加法群''的定义（见
Ch 6，§6.1）。此处仅使用最基本的性质，不依赖群论的完整框架。

\paragraph{2.6.8 无穷递降法}\label{ux65e0ux7a77ux9012ux964dux6cd5}

\textbf{方法}（费马首创）：假设某命题对某正整数成立，推导出对更小的正整数也成立，从而得到无穷递降的正整数序列，与良序原理矛盾。

\textbf{定理 2.6.8}：不存在正整数 \(a, b, c\) 满足
\(a^2 + b^2 = 3c^2\)。

\textbf{证明}：假设存在正整数解。\(a^2 + b^2 \equiv 0 \pmod{3}\)。由于
\(x^2 \equiv 0\) 或 \(1 \pmod{3}\)，只能 \(a^2 \equiv 0\) 且
\(b^2 \equiv 0 \pmod{3}\)，即 \(3 \mid a\) 且 \(3 \mid b\)。设
\(a = 3a'\), \(b = 3b'\)，代入得 \(9a'^2 + 9b'^2 = 3c^2\)，即
\(c^2 = 3a'^2 + 3b'^2\)，故 \(3 \mid c\)，设 \(c = 3c'\)。代入得
\(a'^2 + b'^2 = 3c'^2\)，且
\(a' < a\)。重复此过程得到无穷递降的正整数序列
\(a > a' > a'' > \cdots > 0\)，矛盾。\(\blacksquare\)

\paragraph{2.6.9
鸽巢原理（狄利克雷抽屉原理）}\label{ux9e3dux5de2ux539fux7406ux72c4ux5229ux514bux96f7ux62bdux5c49ux539fux7406}

\textbf{定理 2.6.9（鸽巢原理）}：若 \(n + 1\) 个物体放入 \(n\)
个盒子中，则至少有一个盒子包含至少两个物体。

\textbf{证明}：用反证法。假设每个盒子至多有一个物体，则最多有 \(n\)
个物体，与有 \(n + 1\) 个物体矛盾。\(\blacksquare\)

\textbf{推论 2.6.1（鸽巢原理的推广）}：若 \(kn + 1\) 个物体放入 \(n\)
个盒子中，则至少有一个盒子包含至少 \(k + 1\) 个物体。

\textbf{证明}：假设每个盒子至多有 \(k\) 个物体，则最多有 \(kn\)
个物体，与有 \(kn + 1\) 个物体矛盾。\(\blacksquare\)

\textbf{定理 2.6.10（鸽巢原理的应用）}：任意 \(n + 1\)
个正整数中，必存在两个数除以 \(n\) 的余数相同。（注：两个数除以 \(n\)
的余数相同，等价于它们的差能被 \(n\) 整除。这一关系在 Ch 5 定义 5.4.1
中称为''关于模 \(n\) 同余''，记为
\(a \equiv b \pmod{n}\)。此处提前使用此术语。）

\textbf{证明}：\(n + 1\) 个数除以 \(n\) 的余数在
\(\{0, 1, \ldots, n-1\}\) 中取值（\(n\) 个可能的值）。由鸽巢原理（定理
2.6.9），至少有两个数的余数相同，即关于模 \(n\) 同余。\(\blacksquare\)

\bigskip\hrule\bigskip

\subsubsection{本章展望}\label{ux672cux7ae0ux5c55ux671b-1}

本章建立了命题逻辑和谓词逻辑的理论体系，并介绍了直接证明、反证法、数学归纳法等基本证明方法。在
Ch 3
中，我们将用逻辑语言来精确描述集合的概念和运算，集合运算的性质（交换律、结合律等）将直接引用本章的逻辑等价律。在
Ch 4 中，数学归纳法将被用于证明自然数运算的基本性质。在 Ch 5
中，反证法将被用于证明素数的无穷性等深刻结果。

\bigskip\hrule\bigskip

\subsection{Ch 3 集合论}\label{ch-3-ux96c6ux5408ux8bba}

\textbf{前置知识}：依赖 Ch 2（逻辑基础），特别是谓词逻辑（2.2
节）中的量词及其否定、逻辑等价律。集合论用逻辑语言描述数学对象。

在 Ch 2
中我们学习了逻辑推理的规则。现在我们用这套规则来描述数学中最基本的对象------集合。集合是现代数学的基石：几乎所有的数学对象（数、函数、空间等）都可以用集合来精确定义。

\bigskip\hrule\bigskip

\subsubsection{3.1
集合的基本概念}\label{ux96c6ux5408ux7684ux57faux672cux6982ux5ff5}

\paragraph{3.1.1
集合的直观描述}\label{ux96c6ux5408ux7684ux76f4ux89c2ux63cfux8ff0}

\textbf{定义
3.1.1（集合的直观概念）}：\textbf{集合}是一些确定的、互不相同的对象的整体。这些对象称为集合的\textbf{元素}。

``确定''是指给定一个对象，我们能判断它是否属于该集合。``互不相同''是指集合中每个元素只出现一次。

\textbf{表示法}：

\begin{enumerate}
\def\labelenumi{\arabic{enumi}.}
\tightlist
\item
  \textbf{列举法}：将所有元素列出来，如 \(A = \{1, 2, 3\}\)。
\item
  \textbf{描述法}：\(\{x \mid P(x)\}\) 表示所有满足性质 \(P\) 的对象
  \(x\) 组成的集合。
\end{enumerate}

\paragraph{3.1.2 属于关系}\label{ux5c5eux4e8eux5173ux7cfb}

\textbf{定义 3.1.2（属于关系）}：若 \(a\) 是集合 \(A\) 的一个元素，记为
\(a \in A\)。若 \(a\) 不是 \(A\) 的元素，记为 \(a \notin A\)。

\paragraph{3.1.3 外延公理}\label{ux5916ux5ef6ux516cux7406}

\textbf{公理 3.1.1（外延公理）}：两个集合 \(A\) 和 \(B\)
相等，当且仅当它们有完全相同的元素：
\[A = B \iff \forall x\, (x \in A \leftrightarrow x \in B)\]

外延公理告诉我们：一个集合完全由它的元素决定，与元素的排列顺序和重复次数无关。

\paragraph{3.1.4 空集}\label{ux7a7aux96c6}

\textbf{公理 3.1.2（空集公理）}：存在一个不包含任何元素的集合。

\textbf{定理 3.1.1（空集的唯一性）}：不包含任何元素的集合是唯一的。

\textbf{证明}：设 \(A\) 和 \(B\)
都是不包含任何元素的集合。由外延公理，只需证明
\(\forall x\, (x \in A \leftrightarrow x \in B)\)。对任意
\(x\)：\(x \in A\) 为假（\(A\) 无元素），\(x \in B\) 也为假（\(B\)
无元素），故 \(x \in A \leftrightarrow x \in B\)
为真。由外延公理，\(A = B\)。\(\blacksquare\)

\textbf{定义
3.1.3（空集）}：不包含任何元素的唯一集合称为\textbf{空集}，记为
\(\varnothing\)。

\textbf{注记
3.1.1}：空集是集合论中最基本的对象。它是唯一的不含任何元素的集合。注意
\(\varnothing \neq \{\varnothing\}\)：前者没有元素，后者有一个元素（即空集本身）。类似地，\(\varnothing \in \{\varnothing\}\)
但 \(\varnothing \subseteq \{\varnothing\}\)（空集是任何集合的子集）。

\bigskip\hrule\bigskip

\subsubsection{3.2 集合的运算}\label{ux96c6ux5408ux7684ux8fd0ux7b97}

\paragraph{3.2.1 子集}\label{ux5b50ux96c6}

\textbf{定义 3.2.1（子集）}：若集合 \(A\) 的每个元素都是集合 \(B\)
的元素，则称 \(A\) 是 \(B\) 的\textbf{子集}，记为 \(A \subseteq B\)：
\[A \subseteq B \iff \forall x\, (x \in A \rightarrow x \in B)\]

\textbf{定义 3.2.2（真子集）}：若 \(A \subseteq B\) 且
\(A \neq B\)，则称 \(A\) 是 \(B\) 的\textbf{真子集}，记为
\(A \subsetneq B\)。

\textbf{注记 3.2.1}：\(A \subsetneq B\) 等价于 \(A \subseteq B\) 且
\(\exists x \in B \setminus A\)。初学者容易混淆 \(\subseteq\) 和
\(\in\)：\(1 \in \{1, 2\}\) 但 \(\{1\} \subseteq \{1, 2\}\)。\(1\)
是元素，\(\{1\}\) 是子集。

\textbf{定理 3.2.1（子集关系的性质）}：对任意集合 \(A, B, C\)： 1.
\(A \subseteq A\)（自反性） 2. \(A \subseteq B\) 且 \(B \subseteq A\) 则
\(A = B\)（反对称性） 3. \(A \subseteq B\) 且 \(B \subseteq C\) 则
\(A \subseteq C\)（传递性）

\textbf{证明}：

\textbf{（1）}：对任意 \(x\)，\(x \in A \rightarrow x \in A\) 恒为真，故
\(A \subseteq A\)。\(\blacksquare\)

\textbf{（2）}：设 \(A \subseteq B\) 且 \(B \subseteq A\)。对任意
\(x\)：由 \(A \subseteq B\)，\(x \in A \rightarrow x \in B\)；由
\(B \subseteq A\)，\(x \in B \rightarrow x \in A\)。因此
\(x \in A \leftrightarrow x \in B\)，由外延公理
\(A = B\)。\(\blacksquare\)

\textbf{（3）}：设 \(A \subseteq B\) 且 \(B \subseteq C\)。对任意
\(x\)：若 \(x \in A\)，由 \(A \subseteq B\) 得 \(x \in B\)，再由
\(B \subseteq C\) 得 \(x \in C\)。因此
\(x \in A \rightarrow x \in C\)，即 \(A \subseteq C\)。\(\blacksquare\)

\textbf{定理 3.2.2}：空集是任何集合的子集，即
\(\varnothing \subseteq A\) 对任意集合 \(A\) 成立。

\textbf{证明}：需证明
\(\forall x\, (x \in \varnothing \rightarrow x \in A)\)。对任意
\(x\)，\(x \in \varnothing\) 为假，因此
\(x \in \varnothing \rightarrow x \in A\)
为真（前件为假时蕴含式恒真）。\(\blacksquare\)

\paragraph{3.2.2 并集与交集}\label{ux5e76ux96c6ux4e0eux4ea4ux96c6}

\textbf{定义
3.2.3（并集）}：\(A \cup B = \{x \mid x \in A \vee x \in B\}\)

\textbf{定义
3.2.4（交集）}：\(A \cap B = \{x \mid x \in A \wedge x \in B\}\)

\textbf{定义 3.2.5（不相交）}：若 \(A \cap B = \varnothing\)，则称 \(A\)
和 \(B\) \textbf{不相交}。

\textbf{定义 3.2.6（差集与补集）}： -
\textbf{差集}：\(A \setminus B = \{x \mid x \in A \wedge x \notin B\}\)
- \textbf{补集}：当 \(B \subseteq A\) 时，\(A \setminus B\) 也称为 \(B\)
在 \(A\) 中的补集，记为 \(\complement_{A} B\)。

\paragraph{3.2.3 幂集}\label{ux5e42ux96c6}

\textbf{定义 3.2.7（幂集）}：集合 \(A\) 的\textbf{幂集}定义为 \(A\)
的所有子集组成的集合： \[\mathcal{P}(A) = \{S \mid S \subseteq A\}\]

\textbf{定理 3.2.3}：若 \(A\) 有 \(n\) 个元素，则 \(\mathcal{P}(A)\) 有
\(2^n\) 个元素。

\textbf{证明}：对 \(n\) 用数学归纳法。

\textbf{基础步骤}（\(n = 0\)）：\(A = \varnothing\)，\(\mathcal{P}(\varnothing) = \{\varnothing\}\)
有 \(2^0 = 1\) 个元素。成立。

\textbf{归纳步骤}：设对 \(n\) 个元素的集合，幂集有 \(2^n\) 个元素。设
\(A\) 有 \(n + 1\) 个元素，取 \(a \in A\)，令
\(B = A \setminus \{a\}\)，则 \(B\) 有 \(n\) 个元素。

\(\mathcal{P}(A)\) 中的集合可分为两类： - 不包含 \(a\) 的子集：恰好是
\(B\) 的子集，共 \(2^n\) 个。 - 包含 \(a\) 的子集：形如
\(S \cup \{a\}\)（\(S \subseteq B\)），也有 \(2^n\) 个。

总计 \(2^n + 2^n = 2^{n+1}\)。由归纳法得证。\(\blacksquare\)

\paragraph{3.2.4 笛卡尔积}\label{ux7b1bux5361ux5c14ux79ef}

\textbf{定义 3.2.8（有序对）}：有序对 \((a, b)\) 定义为
\(\{\{a\}, \{a, b\}\}\)（Kuratowski 定义）。

\textbf{引理 3.2.1}：\((a, b) = (c, d)\) 当且仅当 \(a = c\) 且
\(b = d\)。

\textbf{证明}：

\((\Leftarrow)\)：显然。

\((\Rightarrow)\)：设 \((a, b) = (c, d)\)，即
\(\{\{a\}, \{a, b\}\} = \{\{c\}, \{c, d\}\}\)。

若 \(a = b\)，则
\(\{\{a\}, \{a, b\}\} = \{\{a\}\}\)。由集合相等，\(\{\{c\}, \{c, d\}\} = \{\{a\}\}\)，故
\(\{c\} = \{c, d\}\)，得 \(c = d\) 且 \(\{a\} = \{c\}\)，即
\(a = c\)，\(b = a = c = d\)。

若 \(a \neq b\)，则 \(\{\{a\}, \{a, b\}\}\) 有两个不同元素。由
\(\{a\} \in \{\{c\}, \{c, d\}\}\)，有 \(\{a\} = \{c\}\) 或
\(\{a\} = \{c, d\}\)。

若 \(\{a\} = \{c, d\}\)，则 \(c = d = a\)。再由
\(\{a, b\} \in \{\{a\}\}\)，得
\(\{a, b\} = \{a\}\)，矛盾（\(a \neq b\)）。

若 \(\{a\} = \{c\}\)，则 \(a = c\)。由 \(\{a, b\} = \{a, d\}\)，得
\(b = d\)。\(\blacksquare\)

\textbf{定义 3.2.9（笛卡尔积）}：
\[A \times B = \{(a, b) \mid a \in A,\, b \in B\}\]

\paragraph{3.2.5
集合运算的基本性质}\label{ux96c6ux5408ux8fd0ux7b97ux7684ux57faux672cux6027ux8d28}

\textbf{定理 3.2.4（交换律）}：对任意集合 \(A, B\)： 1.
\(A \cup B = B \cup A\) 2. \(A \cap B = B \cap A\)

\textbf{证明}： (1)
\(x \in A \cup B \iff x \in A \vee x \in B \iff x \in B \vee x \in A \iff x \in B \cup A\)（析取交换律，定理
2.1.2）。 由外延公理得证。\(\blacksquare\)

\textbf{定理 3.2.5（结合律）}：对任意集合 \(A, B, C\)： 1.
\((A \cup B) \cup C = A \cup (B \cup C)\) 2.
\((A \cap B) \cap C = A \cap (B \cap C)\)

\textbf{证明}： (1)
\(x \in (A \cup B) \cup C \iff (x \in A \vee x \in B) \vee x \in C \iff x \in A \vee (x \in B \vee x \in C) \iff x \in A \cup (B \cup C)\)
中间使用析取结合律（定理 2.1.3）。\(\blacksquare\)

\textbf{定理 3.2.6（分配律）}：对任意集合 \(A, B, C\)： 1.
\(A \cup (B \cap C) = (A \cup B) \cap (A \cup C)\) 2.
\(A \cap (B \cup C) = (A \cap B) \cup (A \cap C)\)

\textbf{证明}： (1)
\(x \in A \cup (B \cap C) \iff x \in A \vee (x \in B \wedge x \in C) \iff (x \in A \vee x \in B) \wedge (x \in A \vee x \in C) \iff x \in (A \cup B) \cap (A \cup C)\)
使用析取对合取的分配律（定理 2.1.4）。\(\blacksquare\)

\textbf{定理 3.2.7（德摩根律）}：对任意集合 \(A, B\)（在论域 \(U\)
中考虑补集）： 1.
\(U \setminus (A \cup B) = (U \setminus A) \cap (U \setminus B)\) 2.
\(U \setminus (A \cap B) = (U \setminus A) \cup (U \setminus B)\)

\textbf{证明}： (1)
\(x \in U \setminus (A \cup B) \iff x \notin A \cup B \iff \neg(x \in A \vee x \in B) \iff \neg(x \in A) \wedge \neg(x \in B) \iff x \notin A \wedge x \notin B \iff x \in (U \setminus A) \cap (U \setminus B)\)
使用逻辑德摩根律（定理 2.1.5）。\(\blacksquare\)

\textbf{定理 3.2.8（幂等律）}：\(A \cup A = A\)，\(A \cap A = A\)。

\textbf{证明}：\(x \in A \cup A \iff x \in A \vee x \in A \iff x \in A\)。\(\blacksquare\)

\textbf{定理 3.2.9（同一律）}： 1. \(A \cup \varnothing = A\) 2.
\(A \cap \varnothing = \varnothing\) 3. \(A \cup U = U\) 4.
\(A \cap U = A\)

\textbf{证明}：(1)
\(x \in A \cup \varnothing \iff x \in A \vee x \in \varnothing \iff x \in A \vee F \iff x \in A\)。\(\blacksquare\)

\textbf{定理 3.2.10（补集的基本性质）}：对任意集合 \(A \subseteq U\)：
1. \(A \cup (U \setminus A) = U\) 2.
\(A \cap (U \setminus A) = \varnothing\) 3.
\(U \setminus (U \setminus A) = A\)（双重补集）

\textbf{证明}： (1) 若 \(x \in U\)，则 \(x \in A\) 或
\(x \notin A\)（排中律），故
\(x \in A \cup (U \setminus A)\)。\(\blacksquare\)

\textbf{定理 3.2.11（笛卡尔积的性质）}： 1.
\(A \times (B \cup C) = (A \times B) \cup (A \times C)\) 2.
\(A \times (B \cap C) = (A \times B) \cap (A \times C)\)

\textbf{证明}： (1)
\((x, y) \in A \times (B \cup C) \iff x \in A \wedge (y \in B \vee y \in C) \iff (x \in A \wedge y \in B) \vee (x \in A \wedge y \in C) \iff (x, y) \in A \times B \vee (x, y) \in A \times C\)。\(\blacksquare\)

\bigskip\hrule\bigskip

\subsubsection{3.3 关系}\label{ux5173ux7cfb}

\paragraph{3.3.1 二元关系}\label{ux4e8cux5143ux5173ux7cfb}

\textbf{定义 3.3.1（二元关系）}：集合 \(A\) 到集合 \(B\)
的一个\textbf{二元关系}是笛卡尔积 \(A \times B\) 的一个子集。

即 \(R \subseteq A \times B\)。若 \((a, b) \in R\)，通常写为
\(a\, R\, b\)。

\textbf{注记
3.3.1}：二元关系是集合论中描述对象之间联系的基本工具。``小于''关系 \(<\)
可以定义为
\(\{(a, b) \in \mathbb{Z} \times \mathbb{Z} : b - a > 0\}\)，``等于''关系可以定义为对角线
\(\{(a, a) : a \in A\}\)。

\textbf{定义 3.3.2（关系的性质）}：设 \(R\) 是集合 \(A\)
上的二元关系（即 \(R \subseteq A \times A\)）： -
\textbf{自反性}：\(\forall a \in A,\, a\, R\, a\) -
\textbf{对称性}：\(\forall a, b \in A,\, a\, R\, b \rightarrow b\, R\, a\)
-
\textbf{反对称性}：\(\forall a, b \in A,\, (a\, R\, b \wedge b\, R\, a) \rightarrow a = b\)
-
\textbf{传递性}：\(\forall a, b, c \in A,\, (a\, R\, b \wedge b\, R\, c) \rightarrow a\, R\, c\)

\paragraph{3.3.2 等价关系}\label{ux7b49ux4ef7ux5173ux7cfb}

\textbf{定义 3.3.3（等价关系）}：集合 \(A\)
上的一个\textbf{等价关系}是同时满足自反性、对称性和传递性的二元关系，通常记为
\(\sim\)。

\textbf{定义 3.3.4（等价类）}：设 \(\sim\) 是集合 \(A\)
上的等价关系，\(a \in A\)，则 \(a\) 的\textbf{等价类}定义为：
\[[a] = \{x \in A \mid x \sim a\}\]

\textbf{定理 3.3.1（等价类的基本性质）}：设 \(\sim\) 是集合 \(A\)
上的等价关系，则： 1. \(\forall a \in A,\, a \in [a]\) 2.
\([a] = [b] \iff a \sim b\) 3.
\([a] \neq [b] \iff [a] \cap [b] = \varnothing\)

\textbf{证明}：

\textbf{（1）}：由自反性，\(a \sim a\)，故
\(a \in [a]\)。\(\blacksquare\)

\textbf{（2）}： \((\Rightarrow)\)：设 \([a] = [b]\)。由
(1)，\(b \in [b] = [a]\)，故
\(b \sim a\)。由对称性，\(a \sim b\)。\(\blacksquare\)

\((\Leftarrow)\)：设 \(a \sim b\)。

先证 \([a] \subseteq [b]\)：设 \(x \in [a]\)，即
\(x \sim a\)。由传递性，\(x \sim a\) 且 \(a \sim b\) 推出
\(x \sim b\)，故 \(x \in [b]\)。

再证 \([b] \subseteq [a]\)：设 \(x \in [b]\)，即
\(x \sim b\)。由对称性，\(a \sim b\) 推出
\(b \sim a\)。由传递性，\(x \sim b\) 且 \(b \sim a\) 推出
\(x \sim a\)，故 \(x \in [a]\)。

因此 \([a] = [b]\)。\(\blacksquare\)

\textbf{（3）}：需证双向命题
\([a] \neq [b] \iff [a] \cap [b] = \varnothing\)。

\((\Leftarrow)\)（反方向）：假设 \([a] \cap [b] \neq \varnothing\)，需证
\([a] = [b]\)。设 \(c \in [a] \cap [b]\)，即 \(c \sim a\) 且
\(c \sim b\)。由对称性，\(a \sim c\)。由传递性，\(a \sim c\) 且
\(c \sim b\) 推出 \(a \sim b\)。由 (2)，\([a] = [b]\)。此即
\([a] \cap [b] \neq \varnothing \Rightarrow [a] = [b]\)，其逆否命题为
\([a] \neq [b] \Rightarrow [a] \cap [b] = \varnothing\)。

\((\Rightarrow)\)（正方向）：假设 \([a] = [b]\)，需证
\([a] \cap [b] \neq \varnothing\)。由 (1)，\(a \in [a]\)，故
\(a \in [a] = [b]\)，从而 \(a \in [a] \cap [b]\)，即
\([a] \cap [b] \neq \varnothing\)。其逆否命题为
\([a] \cap [b] = \varnothing \Rightarrow [a] \neq [b]\)。

综合两个方向，\([a] \neq [b] \iff [a] \cap [b] = \varnothing\)。\(\blacksquare\)

\textbf{定义 3.3.5（商集）}：\(A\) 关于 \(\sim\)
的\textbf{商集}定义为所有等价类的集合：
\[A / \sim = \{[a] \mid a \in A\}\]

\textbf{定理 3.3.2（商集是划分）}：商集 \(A / \sim\) 构成 \(A\)
的一个\textbf{划分}，即： 1. 每个等价类非空； 2. 不同的等价类不相交； 3.
所有等价类的并等于 \(A\)。

\textbf{证明}： (1) 由定理 3.3.1(1)，\(a \in [a]\)，故
\([a] \neq \varnothing\)。 (2) 由定理 3.3.1(3)。 (3) \((\subseteq)\)：若
\(x \in \bigcup_{a \in A} [a]\)，则存在 \(a \in A\) 使得
\(x \in [a] \subseteq A\)。\((\supseteq)\)：若 \(x \in A\)，则
\(x \in [x]\)。\(\blacksquare\)

\bigskip\hrule\bigskip

\subsubsection{3.4 函数}\label{ux51fdux6570}

\paragraph{3.4.1 函数的定义}\label{ux51fdux6570ux7684ux5b9aux4e49}

\textbf{定义 3.4.1（函数）}：设 \(A, B\) 是集合。一个从 \(A\) 到 \(B\)
的\textbf{函数} \(f: A \to B\) 是满足以下条件的二元关系
\(f \subseteq A \times B\)： 1.
\textbf{存在性}：\(\forall a \in A,\, \exists b \in B,\, (a, b) \in f\)
2.
\textbf{唯一性}：\(\forall a \in A,\, \forall b_1, b_2 \in B,\, ((a, b_1) \in f \wedge (a, b_2) \in f) \rightarrow b_1 = b_2\)

若 \((a, b) \in f\)，记 \(b = f(a)\)。\(A\) 称为 \(f\)
的\textbf{定义域}，\(B\) 称为\textbf{陪域}。\(\{f(a) \mid a \in A\}\)
称为 \(f\) 的\textbf{值域}。

\paragraph{3.4.2
单射、满射、双射}\label{ux5355ux5c04ux6ee1ux5c04ux53ccux5c04}

\textbf{定义 3.4.2}：设 \(f: A \to B\) 是函数。 - \(f\)
是\textbf{单射}（injective）：\(\forall a_1, a_2 \in A,\, f(a_1) = f(a_2) \rightarrow a_1 = a_2\)
- \(f\)
是\textbf{满射}（surjective）：\(\forall b \in B,\, \exists a \in A,\, f(a) = b\)
- \(f\) 是\textbf{双射}（bijective）：\(f\) 既是单射又是满射。

\textbf{定理 3.4.1（单射的等价刻画）}：\(f: A \to B\) 是单射，当且仅当
\(\forall a_1, a_2 \in A,\, a_1 \neq a_2 \rightarrow f(a_1) \neq f(a_2)\)。

\textbf{证明}：这是单射定义的逆否命题。由定理
2.1.10，\(P \rightarrow Q \equiv \neg Q \rightarrow \neg P\)。\(\blacksquare\)

\paragraph{3.4.3 复合函数}\label{ux590dux5408ux51fdux6570}

\textbf{定义 3.4.3（复合函数）}：设 \(f: A \to B\) 和 \(g: B \to C\)
是函数，则 \(g\) 和 \(f\) 的\textbf{复合函数} \(g \circ f: A \to C\)
定义为： \[(g \circ f)(a) = g(f(a)), \quad \forall a \in A\]

\textbf{注记 3.4.1}：注意复合的书写顺序：\(g \circ f\) 表示先应用
\(f\)，再应用 \(g\)（从右向左阅读）。这有时会让人困惑，因为 \(f\)
写在右边却先被应用。

\textbf{定理 3.4.2（复合保持单射和满射）}：设
\(f: A \to B\)，\(g: B \to C\)。 1. 若 \(f\) 和 \(g\) 都是单射，则
\(g \circ f\) 也是单射。 2. 若 \(f\) 和 \(g\) 都是满射，则 \(g \circ f\)
也是满射。 3. 若 \(f\) 和 \(g\) 都是双射，则 \(g \circ f\) 也是双射。

\textbf{证明}：

\textbf{（1）}：设 \((g \circ f)(a_1) = (g \circ f)(a_2)\)，即
\(g(f(a_1)) = g(f(a_2))\)。因 \(g\) 是单射，得 \(f(a_1) = f(a_2)\)。因
\(f\) 是单射，得 \(a_1 = a_2\)。\(\blacksquare\)

\textbf{（2）}：设 \(c \in C\)。因 \(g\) 是满射，存在 \(b \in B\) 使得
\(g(b) = c\)。因 \(f\) 是满射，存在 \(a \in A\) 使得 \(f(a) = b\)。于是
\((g \circ f)(a) = g(f(a)) = g(b) = c\)。\(\blacksquare\)

\textbf{（3）}：由 (1) 和 (2)。\(\blacksquare\)

\textbf{注记 3.4.2（逆命题不成立）}：\(g \circ f\) 是单射不能推出 \(f\)
和 \(g\) 都是单射。例如
\(f: \{1\} \to \{1, 2\}\)，\(f(1) = 1\)（单射但不满射）；\(g: \{1, 2\} \to \{1\}\)，\(g(1) = g(2) = 1\)（满射但不是单射）。\(g \circ f: \{1\} \to \{1\}\)，\(g(f(1)) = 1\)
是双射。但可以证明：若 \(g \circ f\) 是单射，则 \(f\) 是单射；若
\(g \circ f\) 是满射，则 \(g\) 是满射。

\textbf{定理 3.4.3（复合的结合律）}：设
\(f: A \to B\)，\(g: B \to C\)，\(h: C \to D\)，则
\[h \circ (g \circ f) = (h \circ g) \circ f\]

\textbf{证明}：对任意 \(a \in A\)：
\[(h \circ (g \circ f))(a) = h((g \circ f)(a)) = h(g(f(a)))\]
\[((h \circ g) \circ f)(a) = (h \circ g)(f(a)) = h(g(f(a)))\]
两者相等。\(\blacksquare\)

\paragraph{3.4.4 反函数}\label{ux53cdux51fdux6570}

\textbf{定义 3.4.4（反函数）}：设 \(f: A \to B\) 是双射，则 \(f\)
的\textbf{反函数} \(f^{-1}: B \to A\) 定义为：对每个
\(b \in B\)，\(f^{-1}(b)\) 是唯一的 \(a \in A\) 使得 \(f(a) = b\)。

\textbf{注记 3.4.3}：符号 \(f^{-1}\) 不是指
\(\frac{1}{f}\)，而是表示反函数。双射保证了这样的 \(a\)
存在（由满射性）且唯一（由单射性）。

\textbf{定理 3.4.4（反函数的基本性质）}：设 \(f: A \to B\)
是双射，\(f^{-1}\) 是其反函数，则： 1.
\(f^{-1} \circ f = \mathrm{id}_{A}\) 2.
\(f \circ f^{-1} = \mathrm{id}_{B}\) 3. \(f^{-1}\) 也是双射 4.
\((f^{-1})^{-1} = f\)

\textbf{证明}：

\textbf{（1）}：对任意 \(a \in A\)，设
\(b = f(a)\)。由反函数定义，\(f^{-1}(b) = a\)。故
\((f^{-1} \circ f)(a) = f^{-1}(f(a)) = a\)。\(\blacksquare\)

\textbf{（2）}：对任意 \(b \in B\)，设
\(a = f^{-1}(b)\)。由反函数定义，\(f(a) = b\)。故
\((f \circ f^{-1})(b) = f(f^{-1}(b)) = b\)。\(\blacksquare\)

\textbf{（3）}：\(f^{-1}\) 是单射：若
\(f^{-1}(b_1) = f^{-1}(b_2) = a\)，则 \(f(a) = b_1\) 且
\(f(a) = b_2\)，故 \(b_1 = b_2\)。

\(f^{-1}\) 是满射：对任意 \(a \in A\)，取 \(b = f(a)\)，则
\(f^{-1}(b) = a\)。\(\blacksquare\)

\textbf{（4）}：由 (1) 和 (2) 的对称性，\(f\) 是 \(f^{-1}\)
的反函数。\(\blacksquare\)

\textbf{定理 3.4.5（双射的等价刻画）}：函数 \(f: A \to B\)
是双射，当且仅当存在函数 \(g: B \to A\) 使得
\(g \circ f = \mathrm{id}_{A}\) 且 \(f \circ g = \mathrm{id}_{B}\)。

\textbf{证明}：

\((\Rightarrow)\)：取 \(g = f^{-1}\)，由定理 3.4.4 得证。

\((\Leftarrow)\)：设存在这样的 \(g\)。\(f\) 是单射：若
\(f(a_1) = f(a_2)\)，则 \(a_1 = g(f(a_1)) = g(f(a_2)) = a_2\)。\(f\)
是满射：对任意 \(b \in B\)，取 \(a = g(b)\)，则
\(f(a) = f(g(b)) = b\)。\(\blacksquare\)

\bigskip\hrule\bigskip

\subsubsection{3.5
ZFC公理系统简介}\label{zfcux516cux7406ux7cfbux7edfux7b80ux4ecb}

前述的集合操作依赖于一组严格的公理。以下列出 ZFC（Zermelo-Fraenkel with
Choice）公理系统的核心条目，使读者了解集合论的逻辑根基。

\paragraph{3.5.1 外延公理}\label{ux5916ux5ef6ux516cux7406-1}

\textbf{公理
3.5.1（外延公理）}：\(\forall A\, \forall B\, [\forall x\, (x \in A \leftrightarrow x \in B) \rightarrow A = B]\)。此即公理
3.1.1 的精确逻辑表述。

\paragraph{3.5.2 空集公理}\label{ux7a7aux96c6ux516cux7406}

\textbf{公理
3.5.2（空集公理）}：\(\exists A\, \forall x\, (x \notin A)\)。此即公理
3.1.2 的精确表述。

\paragraph{3.5.3 配对公理}\label{ux914dux5bf9ux516cux7406}

\textbf{公理
3.5.3（配对公理）}：\(\forall a\, \forall b\, \exists A\, \forall x\, (x \in A \leftrightarrow x = a \vee x = b)\)。即对任意两个对象
\(a, b\)，存在集合 \(\{a, b\}\)。

\textbf{推论 3.5.1（单元素集的存在性）}：对任意对象 \(a\)，存在集合
\(\{a\}\)。

\textbf{证明}：取 \(b = a\)，由配对公理得
\(\{a, a\} = \{a\}\)。\(\blacksquare\)

\paragraph{3.5.4 并集公理}\label{ux5e76ux96c6ux516cux7406}

\textbf{公理 3.5.4（并集公理）}：对任意集合 \(\mathcal{F}\)，存在集合
\(A\) 使得
\(\forall x\, [x \in A \leftrightarrow \exists Y\, (Y \in \mathcal{F} \wedge x \in Y)]\)。此集合记为
\(\bigcup \mathcal{F}\)。

\paragraph{3.5.5 幂集公理}\label{ux5e42ux96c6ux516cux7406}

\textbf{公理 3.5.5（幂集公理）}：对任意集合 \(A\)，存在集合
\(\mathcal{P}(A)\) 使得
\(\forall S\, (S \in \mathcal{P}(A) \leftrightarrow S \subseteq A)\)。

\paragraph{3.5.6 无穷公理}\label{ux65e0ux7a77ux516cux7406}

\textbf{公理
3.5.6（无穷公理）}：\(\exists A\, [\varnothing \in A \wedge \forall x\, (x \in A \rightarrow x \cup \{x\} \in A)]\)。此公理保证了自然数集的存在性（见
Ch 4）。

\paragraph{3.5.7
替换公理模式}\label{ux66ffux6362ux516cux7406ux6a21ux5f0f}

\textbf{公理 3.5.7（替换公理模式）}：设 \(\varphi(x, y)\)
是一个公式，对任意集合 \(A\)，若
\(\forall x \in A\, \exists!\, y\, \varphi(x, y)\)，则存在集合
\(\{y \mid \exists x \in A\, \varphi(x, y)\}\)。此公理允许由已知集合通过''函数替换''构造新集合。

\paragraph{3.5.8 正则公理}\label{ux6b63ux5219ux516cux7406}

\textbf{公理
3.5.8（正则公理/基础公理）}：\(\forall A\, [A \neq \varnothing \rightarrow \exists x \in A\, (x \cap A = \varnothing)]\)。即每个非空集合都有一个与自身不相交的元素，从而排除了
\(A \in A\) 等''病态''集合。

\textbf{定理 3.5.1}：不存在集合 \(A\) 使得 \(A \in A\)。

\textbf{证明}：假设 \(A \in A\)。令 \(B = \{A\}\)。\(B\)
是非空集合。由正则公理，存在 \(x \in B\) 使得
\(x \cap B = \varnothing\)。但 \(x = A\)，\(A \cap \{A\}\) 包含
\(A\)（因 \(A \in A\) 且 \(A \in \{A\}\)），故
\(A \cap \{A\} \neq \varnothing\)，矛盾。\(\blacksquare\)

\paragraph{3.5.9 选择公理}\label{ux9009ux62e9ux516cux7406}

\textbf{公理 3.5.9（选择公理）}：对任意由非空集合组成的集合族
\(\mathcal{F}\)，存在函数 \(f: \mathcal{F} \to \bigcup \mathcal{F}\)
使得 \(f(X) \in X\) 对所有 \(X \in \mathcal{F}\)
成立。即可以从每个非空集合中各''选择''一个元素。

选择公理等价于许多重要命题，包括佐恩引理和良序原理（定理 4.1.11
在一般集合上的推广）。

\bigskip\hrule\bigskip

\subsubsection{3.6 基数简介}\label{ux57faux6570ux7b80ux4ecb}

\paragraph{3.6.1 等势}\label{ux7b49ux52bf}

\textbf{定义 3.6.1（等势）}：两个集合 \(A\) 和 \(B\)
称为\textbf{等势}（记为 \(A \sim B\)），如果存在双射 \(f: A \to B\)。

\textbf{定理 3.6.1}：等势是等价关系。

\textbf{证明}：自反性：\(\mathrm{id}_{A}: A \to A\) 是双射。对称性：若
\(f: A \to B\) 是双射，则 \(f^{-1}: B \to A\) 是双射。传递性：若
\(f: A \to B\) 和 \(g: B \to C\) 是双射，则 \(g \circ f: A \to C\)
是双射。\(\blacksquare\)

\paragraph{3.6.2 可数集}\label{ux53efux6570ux96c6}

\textbf{定义 3.6.2（可数集）}：与自然数集 \(\mathbb{N}\)
等势的集合称为\textbf{可数无限集}。有限集或可数无限集统称为\textbf{可数集}。

\textbf{定理 3.6.2}：整数集 \(\mathbb{Z}\) 是可数的。

\textbf{证明}：构造双射 \(f: \mathbb{N} \to \mathbb{Z}\)：
\[f(n) = \begin{cases} n/2 & \text{若 } n \text{ 是偶数} \\ -(n+1)/2 & \text{若 } n \text{ 是奇数} \end{cases}\]
\(f(0) = 0, f(1) = -1, f(2) = 1, f(3) = -2, f(4) = 2, \ldots\)。这是
\(\mathbb{N}\) 到 \(\mathbb{Z}\) 的一一对应。\(\blacksquare\)

\textbf{定理 3.6.3}：有理数集 \(\mathbb{Q}\) 是可数的。

\textbf{证明}：只需证明正有理数 \(\mathbb{Q}^+\) 可数（然后用定理 3.6.2
的技巧处理符号）。构造显式的 Cantor 配对函数
\(J: \mathbb{N} \times \mathbb{N} \to \mathbb{N}\)：

\[J(n, m) = \frac{(n+m-2)(n+m-1)}{2} + n, \quad n, m \in \mathbb{N}\]

其中 \(\mathbb{N} = \{1, 2, 3, \ldots\}\)。\(J\) 是双射的验证： -
\textbf{单射性}：若 \(J(n, m) = J(n', m')\)，则由公式可知 \(n+m\) 与
\(n'+m'\) 必相等（否则对应不同的对角线条目区间），进而
\(n = n'\)，\(m = m'\)。 - \textbf{满射性}：对任意
\(k \in \mathbb{N}\)，存在唯一的自然数 \(s\) 使得
\(T_{s-1} < k \leq T_s\)（其中 \(T_s = \frac{s(s+1)}{2}\) 为三角数）。令
\(n = k - T_{s-1}\)，\(m = s+1 - n\)，则
\(n \in \{1, \ldots, s\}\)，\(m \in \{1, \ldots, s\}\)，且
\(J(n, m) = k\)。

现在令 \(f: \mathbb{N} \to \mathbb{Q}^+\) 为：对
\(k \in \mathbb{N}\)，由 \(k\) 的唯一分解 \(k = J(n, m)\)，令
\(f(k) = n/m\)。由于 \(J\)
是双射，每个有理数（未约分）被遍历恰好一次。跳过重复者（即当
\(\gcd(n, m) > 1\) 时）即得 \(\mathbb{N}\) 到 \(\mathbb{Q}^+\)
的一一对应。\(\blacksquare\)

\textbf{定理 3.6.4（康托尔对角线论证）}：实数集 \(\mathbb{R}\)
是不可数的。

\textbf{证明}：只需证明 \((0, 1)\) 不可数。假设 \((0, 1)\)
可数，将其元素排列为 \(r_1, r_2, r_3, \ldots\)。设 \(r_n\)
的十进制表示为 \(0.d_{n1}d_{n2}d_{n3}\ldots\)。构造
\(s = 0.s_1 s_2 s_3 \ldots\)，其中
\[s_n = \begin{cases} 3 & \text{若 } d_{nn} \neq 3 \\ 7 & \text{若 } d_{nn} = 3 \end{cases}\]
则 \(s \in (0, 1)\) 且 \(s \neq r_n\) 对所有 \(n\)（因为 \(s\) 的第
\(n\) 位与 \(r_n\) 的第 \(n\) 位不同）。这与 \(r_1, r_2, \ldots\) 列出了
\((0, 1)\) 的所有元素矛盾。\(\blacksquare\)

\textbf{推论 3.6.1}：\(\mathbb{R}\) 与 \(\mathbb{N}\)
不等势。存在''大小''不同的无穷------\(\mathbb{R}\)
是''不可数无穷''，\(\mathbb{N}\) 是''可数无穷''。

\textbf{定理 3.6.5（康托尔定理）}：对任意集合 \(A\)，\(A\) 与
\(\mathcal{P}(A)\) 不等势，且 \(|A| < |\mathcal{P}(A)|\)。

\textbf{证明}：

\textbf{第一步}：\(|A| \leq |\mathcal{P}(A)|\)。映射 \(a \mapsto \{a\}\)
是 \(A\) 到 \(\mathcal{P}(A)\) 的单射。

\textbf{第二步}：不存在从 \(A\) 到 \(\mathcal{P}(A)\) 的满射。假设
\(f: A \to \mathcal{P}(A)\) 是满射。令
\(B = \{a \in A : a \notin f(a)\} \subseteq A\)，故
\(B \in \mathcal{P}(A)\)。由 \(f\) 是满射，存在 \(b \in A\) 使得
\(f(b) = B\)。

问：\(b \in B\) 吗？ - 若 \(b \in B\)，由 \(B\) 的定义
\(b \notin f(b) = B\)，矛盾。 - 若 \(b \notin B\)，由 \(B\) 的定义
\(b \in B\)，矛盾。

因此不存在这样的满射。\(\blacksquare\)

\textbf{推论 3.6.2}：不存在最大的基数。对任意集合
\(A\)，总存在严格更大的集合 \(\mathcal{P}(A)\)。

\textbf{推论 3.6.3}：\(|\mathcal{P}(\mathbb{N})| = |\mathbb{R}|\)。

\textbf{证明概要}：\(\mathcal{P}(\mathbb{N})\) 中的每个子集 \(S\)
对应一个二进制小数 \(0.b_0 b_1 b_2 \ldots\)（\(b_i = 1\) 当且仅当
\(i \in S\)）。这建立了 \(\mathcal{P}(\mathbb{N})\) 与 \([0, 1]\)
之间的双射（需要处理二进制表示的非唯一性，如
\(0.0111\ldots = 0.1000\ldots\)）。\(\blacksquare\)

\bigskip\hrule\bigskip

\subsubsection{本章展望}\label{ux672cux7ae0ux5c55ux671b-2}

本章建立了集合论的基本框架，包括集合的运算、关系和函数的严格定义。在 Ch
4
中，我们将用集合（特别是等价关系和商集）来构造从自然数到复数的完整数系。函数的概念将在卷四（Ch
16--24）中得到充分发展。等价关系将在 Ch
5（同余理论）中再次出现。幂集和笛卡尔积的概念将在后续各卷中反复使用。

\bigskip\hrule\bigskip

\subsection{Ch 4 数系的构造}\label{ch-4-ux6570ux7cfbux7684ux6784ux9020}

\textbf{前置知识}：依赖 Ch 3（集合论），特别是等价关系（3.3.2
节）、商集（定义 3.3.5）、函数（3.4 节）。用集合论的方法构造数系。

数学研究的核心对象之一是''数''。在 Ch 3
中我们建立了集合论的语言，现在我们用这门语言来回答一个根本问题：数究竟是什么？本章从皮亚诺公理出发，逐步构造自然数
\(\mathbb{N}\)、整数 \(\mathbb{Z}\)、有理数 \(\mathbb{Q}\)、实数
\(\mathbb{R}\) 和复数
\(\mathbb{C}\)，每一步都严格证明新数系继承旧数系的代数性质并解决其不足。

\bigskip\hrule\bigskip

\subsubsection{4.1
自然数（皮亚诺公理）}\label{ux81eaux7136ux6570ux76aeux4e9aux8bfaux516cux7406}

\paragraph{4.1.1 皮亚诺公理}\label{ux76aeux4e9aux8bfaux516cux7406}

\textbf{定义 4.1.1（自然数系）}：一个\textbf{自然数系}是一个三元组
\((\mathbb{N}, 0, S)\)，其中 \(\mathbb{N}\)
是一个集合，\(0 \in \mathbb{N}\)
是一个特定元素，\(S: \mathbb{N} \to \mathbb{N}\)
是一个函数（称为\textbf{后继函数}），满足以下五条公理：

\textbf{公理 P1}（零是自然数）：\(0 \in \mathbb{N}\)。

\textbf{公理
P2}（后继的封闭性）：\(\forall n \in \mathbb{N}: S(n) \in \mathbb{N}\)。

\textbf{公理
P3}（零不是任何数的后继）：\(\forall n \in \mathbb{N}: S(n) \neq 0\)。

\textbf{公理
P4}（后继函数是单射）：\(\forall m, n \in \mathbb{N}: S(m) = S(n) \Rightarrow m = n\)。

\textbf{公理 P5（归纳公理）}：设 \(A \subseteq \mathbb{N}\) 满足：(i)
\(0 \in A\)；(ii) 对任意 \(n \in A\)，\(S(n) \in A\)。则
\(A = \mathbb{N}\)。

\textbf{约定}：以下在固定的自然数系 \((\mathbb{N}, 0, S)\) 中工作。将
\(S(0)\) 记为 \(1\)，\(S(S(0))\) 记为 \(2\)，以此类推。

\paragraph{4.1.2
加法的递归定义}\label{ux52a0ux6cd5ux7684ux9012ux5f52ux5b9aux4e49}

\textbf{定理 4.1.1（加法的存在唯一性）}：存在唯一的二元运算
\(+: \mathbb{N} \times \mathbb{N} \to \mathbb{N}\) 满足： 1.
\(\forall a \in \mathbb{N}: a + 0 = a\) 2.
\(\forall a, b \in \mathbb{N}: a + S(b) = S(a + b)\)

\textbf{证明}：

\textbf{存在性}：对每个固定的
\(a \in \mathbb{N}\)，由递归定理（在集合论中可以严格证明），存在唯一的函数
\(f_a: \mathbb{N} \to \mathbb{N}\) 使得
\(f_a(0) = a\)，\(f_a(S(b)) = S(f_a(b))\)。定义 \(a + b = f_a(b)\)。

\textbf{唯一性}：假设 \(+\) 和 \(\oplus\) 都满足上述条件。对固定的
\(a\)，令 \(A = \{b \in \mathbb{N} : a + b = a \oplus b\}\)。由条件
1，\(0 \in A\)。设 \(b \in A\)，则
\(a + S(b) = S(a + b) = S(a \oplus b) = a \oplus S(b)\)，故
\(S(b) \in A\)。由归纳公理，\(A = \mathbb{N}\)。\(\blacksquare\)

\paragraph{4.1.3
乘法的递归定义}\label{ux4e58ux6cd5ux7684ux9012ux5f52ux5b9aux4e49}

\textbf{定理 4.1.2（乘法的存在唯一性）}：存在唯一的二元运算
\(\cdot: \mathbb{N} \times \mathbb{N} \to \mathbb{N}\) 满足： 1.
\(\forall a \in \mathbb{N}: a \cdot 0 = 0\) 2.
\(\forall a, b \in \mathbb{N}: a \cdot S(b) = a \cdot b + a\)

\textbf{证明}：方法与定理 4.1.1 完全类似。\(\blacksquare\)

\paragraph{4.1.4
自然数运算律的证明}\label{ux81eaux7136ux6570ux8fd0ux7b97ux5f8bux7684ux8bc1ux660e}

以下严格证明自然数的基本运算律。所有证明都基于皮亚诺公理和加法、乘法的递归定义。

\textbf{引理 4.1.1}：\(\forall a \in \mathbb{N}: 0 + a = a\)。

\textbf{证明}：对 \(a\) 做归纳。\(0 + 0 = 0\)。设 \(0 + a = a\)，则
\(0 + S(a) = S(0 + a) = S(a)\)。\(\blacksquare\)

\textbf{引理
4.1.2}：\(\forall a, b \in \mathbb{N}: S(a) + b = S(a + b)\)。

\textbf{证明}：对 \(b\) 做归纳。\(S(a) + 0 = S(a) = S(a + 0)\)。设
\(S(a) + b = S(a + b)\)，则
\(S(a) + S(b) = S(S(a) + b) = S(S(a + b)) = S(a + S(b))\)。\(\blacksquare\)

\textbf{注记 4.1.1}：引理 4.1.1
与加法定义的第一条（\(a + 0 = a\)）形式上不同。后者是''右加零''，前者是''左加零''。两者都需要证明，因为加法的递归定义是不对称的（定义在第二个变量上递归）。

\textbf{定理
4.1.3（加法结合律）}：\(\forall a, b, c \in \mathbb{N}: (a + b) + c = a + (b + c)\)。

\textbf{证明}：对 \(c\) 做归纳。(基步)
\((a + b) + 0 = a + b = a + (b + 0)\)。(归纳步) 设
\((a + b) + c = a + (b + c)\)，则
\((a + b) + S(c) = S((a + b) + c) = S(a + (b + c)) = a + S(b + c) = a + (b + S(c))\)。\(\blacksquare\)

\textbf{定理
4.1.4（加法交换律）}：\(\forall a, b \in \mathbb{N}: a + b = b + a\)。

\textbf{证明}：先证辅助结论 \(a + 1 = S(a)\)（因为
\(a + S(0) = S(a + 0) = S(a)\)）。对 \(b\) 做归纳。(基步)
\(a + 0 = a = 0 + a\)（引理 4.1.1）。(归纳步) 设 \(a + b = b + a\)，则
\(a + S(b) = S(a + b) = S(b + a) = S(b) + a\)（引理
4.1.2）。\(\blacksquare\)

\textbf{引理 4.1.3}：\(\forall a \in \mathbb{N}: 0 \cdot a = 0\)。

\textbf{证明}：对 \(a\) 做归纳。\(0 \cdot 0 = 0\)。设
\(0 \cdot a = 0\)，则
\(0 \cdot S(a) = 0 \cdot a + 0 = 0 + 0 = 0\)。\(\blacksquare\)

\textbf{引理
4.1.4}：\(\forall a, b \in \mathbb{N}: S(a) \cdot b = a \cdot b + b\)。

\textbf{证明}：对 \(b\) 做归纳。(基步)
\(S(a) \cdot 0 = 0 = a \cdot 0 + 0\)。(归纳步) 设
\(S(a) \cdot b = a \cdot b + b\)，则
\(S(a) \cdot S(b) = S(a) \cdot b + S(a) = (a \cdot b + b) + S(a)\)。由加法结合律和交换律，\(= a \cdot b + (b + S(a)) = a \cdot b + S(b + a) = a \cdot b + S(a + b) = a \cdot b + (a + S(b)) = (a \cdot b + a) + S(b) = a \cdot S(b) + S(b)\)。\(\blacksquare\)

\textbf{定理
4.1.5（左分配律）}：\(\forall a, b, c \in \mathbb{N}: a \cdot (b + c) = a \cdot b + a \cdot c\)。

\textbf{证明}：对 \(c\) 做归纳。(基步)
\(a \cdot (b + 0) = a \cdot b = a \cdot b + 0 = a \cdot b + a \cdot 0\)。(归纳步)
设 \(a \cdot (b + c) = a \cdot b + a \cdot c\)，则
\(a \cdot (b + S(c)) = a \cdot S(b + c) = a \cdot (b + c) + a = (a \cdot b + a \cdot c) + a = a \cdot b + (a \cdot c + a) = a \cdot b + a \cdot S(c)\)。\(\blacksquare\)

\textbf{定理
4.1.6（乘法交换律）}：\(\forall a, b \in \mathbb{N}: a \cdot b = b \cdot a\)。

\textbf{证明}：对 \(b\) 做归纳。(基步)
\(a \cdot 0 = 0 = 0 \cdot a\)（引理 4.1.3）。(归纳步) 设
\(a \cdot b = b \cdot a\)，则
\(a \cdot S(b) = a \cdot b + a = b \cdot a + a = S(b) \cdot a\)（引理
4.1.4）。\(\blacksquare\)

\textbf{定理
4.1.7（乘法结合律）}：\(\forall a, b, c \in \mathbb{N}: (a \cdot b) \cdot c = a \cdot (b \cdot c)\)。

\textbf{证明}：对 \(c\) 做归纳。(基步)
\((a \cdot b) \cdot 0 = 0 = a \cdot 0 = a \cdot (b \cdot 0)\)。(归纳步)
设 \((a \cdot b) \cdot c = a \cdot (b \cdot c)\)，则
\((a \cdot b) \cdot S(c) = (a \cdot b) \cdot c + a \cdot b = a \cdot (b \cdot c) + a \cdot b = a \cdot (b \cdot c + b) = a \cdot (b \cdot S(c))\)。\(\blacksquare\)

\textbf{定理
4.1.8（消去律）}：\(\forall a, b, c \in \mathbb{N}: a + c = b + c \Rightarrow a = b\)。

\textbf{证明}：对 \(c\) 做归纳。(基步)
\(a + 0 = b + 0 \Rightarrow a = b\)。(归纳步) 设
\(a + c = b + c \Rightarrow a = b\)。若 \(a + S(c) = b + S(c)\)，即
\(S(a + c) = S(b + c)\)，由公理 P4 得 \(a + c = b + c\)，由归纳假设
\(a = b\)。\(\blacksquare\)

\textbf{引理 4.1.5}：\(\forall n \in \mathbb{N}\)，\(n = 0\) 或
\(n = S(m)\) 对某个 \(m \in \mathbb{N}\)。

\textbf{证明}：令
\(A = \{0\} \cup \{S(m) : m \in \mathbb{N}\}\)。\(0 \in A\)。若
\(n \in A\)，则 \(S(n) \in A\)（因为 \(S(n)\) 是某自然数的后继）。由公理
P5（归纳公理），\(A = \mathbb{N}\)。\(\blacksquare\)

\textbf{定理
4.1.9（乘法消去律）}：\(\forall a, b, c \in \mathbb{N}\)，若
\(c \neq 0\)，则 \(a \cdot c = b \cdot c \Rightarrow a = b\)。

\textbf{证明}：对 \(a\) 做归纳。

\textbf{基步}（\(a = 0\)）：则 \(0 = a \cdot c = b \cdot c\)。需证
\(b = 0\)。设 \(c = S(d)\)（因 \(c \neq 0\)，\(c\)
必为某自然数的后继）。若 \(b \neq 0\)，则 \(b = S(b')\)，于是
\(b \cdot c = S(b') \cdot c = b' \cdot c + c = b' \cdot c + S(d) = S(b' \cdot c + d)\)。故
\(S(b' \cdot c + d) = 0\)，与公理 P3（\(0\)
不是任何自然数的后继）矛盾。因此 \(b = 0 = a\)。

\textbf{归纳步}：设 \(a = S(a')\)，且对 \(a'\) 结论成立。设
\(S(a') \cdot c = b \cdot c\)，即 \(a' \cdot c + c = b \cdot c\)。若
\(b = 0\)，则 \(a' \cdot c + c = 0\)。由与基步相同的论证（\(c = S(d)\)
导致 \(S(a' \cdot c + d) = 0\) 矛盾），故 \(b \neq 0\)。设
\(b = S(b')\)，则
\(a' \cdot c + c = S(b') \cdot c = b' \cdot c + c\)。由加法消去律，\(a' \cdot c = b' \cdot c\)。由归纳假设
\(a' = b'\)，故 \(S(a') = S(b')\)。\(\blacksquare\)

\textbf{定理
4.1.10（乘法零因子律）}：\(\forall a, b \in \mathbb{N}: a \cdot b = 0 \iff a = 0 \vee b = 0\)。

\textbf{证明}：\((\Leftarrow)\)：若 \(a = 0\)，\(0 \cdot b = 0\)。若
\(b = 0\)，\(a \cdot 0 = 0\)。\((\Rightarrow)\)：假设 \(a \neq 0\) 且
\(b \neq 0\)，设 \(a = S(a')\), \(b = S(b')\)。则
\(a \cdot b = S(a') \cdot S(b') = S(a') \cdot b' + S(a')\)。由于
\(S(a') \geq 1\)，\(a \cdot b \geq S(a') \geq 1 > 0\)，矛盾。\(\blacksquare\)

\paragraph{4.1.5
自然数的序关系}\label{ux81eaux7136ux6570ux7684ux5e8fux5173ux7cfb}

\textbf{定义 4.1.2}：\(\forall a, b \in \mathbb{N}\)，定义 \(a \leq b\)
当且仅当 \(\exists c \in \mathbb{N}: a + c = b\)。

\textbf{定理 4.1.11（良序原理）}：\(\mathbb{N}\)
的每个非空子集都有最小元素。

\textbf{证明}：假设 \(A \subseteq \mathbb{N}\) 非空但没有最小元素。令
\(B = \{n \in \mathbb{N} : \forall k \leq n, k \notin A\}\)。对 \(B\)
用归纳法证明 \(B = \mathbb{N}\)。

(基步) 若 \(0 \in A\)，则 \(0\) 是 \(A\) 的最小元素，矛盾。故
\(0 \notin A\)，\(0 \in B\)。

(归纳步) 设 \(\{0, 1, \ldots, n\} \subseteq B\)。若 \(S(n) \in A\)，则
\(S(n)\) 是 \(A\) 的最小元素（所有更小的自然数都不在 \(A\)
中），矛盾。故 \(S(n) \notin A\)，\(S(n) \in B\)。

由归纳公理，\(B = \mathbb{N}\)，从而
\(A = \varnothing\)，矛盾。\(\blacksquare\)

\textbf{定理 4.1.11'（良序原理与归纳法的等价性）}：归纳公理（公理
P5）与良序原理（定理 4.1.11）等价。

\textbf{证明}：我们已经从归纳公理推导了良序原理。反过来，假设良序原理成立，证明归纳公理。

设 \(A \subseteq \mathbb{N}\) 满足 (i) \(0 \in A\)；(ii)
\(n \in A \Rightarrow S(n) \in A\)。假设 \(A \neq \mathbb{N}\)，令
\(C = \mathbb{N} \setminus A \neq \varnothing\)。由良序原理，\(C\)
有最小元素 \(m\)。\(m \neq 0\)（因 \(0 \in A\)），故
\(m = S(m')\)。\(m' < m\) 且 \(m' \notin C\)（\(m\) 是最小元素），故
\(m' \in A\)。由条件 (ii)，\(S(m') = m \in A\)，矛盾。\(\blacksquare\)

\bigskip\hrule\bigskip

\subsubsection{4.2 整数}\label{ux6574ux6570}

\paragraph{4.2.1 构造动机}\label{ux6784ux9020ux52a8ux673a}

在自然数系中，减法并不总是可行的：方程 \(3 + x = 1\) 在 \(\mathbb{N}\)
中无解。我们扩展数系使减法总可以进行。

\paragraph{4.2.2 整数的定义}\label{ux6574ux6570ux7684ux5b9aux4e49}

\textbf{定义 4.2.1}：令 \(P = \mathbb{N} \times \mathbb{N}\)。在 \(P\)
上定义关系 \(\sim\)： \[(a, b) \sim (c, d) \iff a + d = b + c\]

\textbf{定理 4.2.1}：\(\sim\) 是 \(P\) 上的等价关系。

\textbf{证明}： - \textbf{自反性}：\((a, b) \sim (a, b)\)，因为
\(a + b = b + a\)（加法交换律）。 - \textbf{对称性}：若
\((a, b) \sim (c, d)\)，即 \(a + d = b + c\)，则 \(c + b = d + a\)，即
\((c, d) \sim (a, b)\)。 - \textbf{传递性}：若 \((a, b) \sim (c, d)\) 且
\((c, d) \sim (e, f)\)，即 \(a + d = b + c\) 且
\(c + f = d + e\)。两式相加得 \(a + d + c + f = b + c + d + e\)，消去
\(c + d\)（加法消去律）得 \(a + f = b + e\)，即
\((a, b) \sim (e, f)\)。\(\blacksquare\)

\textbf{定义 4.2.2}：\textbf{整数集}
\(\mathbb{Z} = P / \sim = \{[a, b] : a, b \in \mathbb{N}\}\)，其中
\([a, b]\) 代表''\(a - b\)``。

嵌入映射 \(\iota: \mathbb{N} \to \mathbb{Z}\) 定义为
\(\iota(n) = [n, 0]\)。

\paragraph{4.2.3 整数的运算}\label{ux6574ux6570ux7684ux8fd0ux7b97}

\textbf{定义 4.2.3（加法）}：\([a, b] + [c, d] = [a + c, b + d]\)

\textbf{定理 4.2.2（加法良定义性）}：上述定义与等价类代表的选取无关。

\textbf{证明}：设 \([a, b] = [a', b']\) 且 \([c, d] = [c', d']\)，即
\(a + b' = b + a'\) 且 \(c + d' = d + c'\)。要证
\((a + c) + (b' + d') = (b + d) + (a' + c')\)。由已知条件相加即得。\(\blacksquare\)

\textbf{定义
4.2.4（乘法）}：\([a, b] \cdot [c, d] = [ac + bd, ad + bc]\)

\textbf{定理 4.2.3（乘法良定义性）}：乘法定义与代表选取无关。

\textbf{证明}：利用 \(a + b' = b + a'\) 和
\(c + d' = d + c'\)，通过自然数运算律可验证
\((ac + bd) + (a'd' + b'c') = (ad + bc) + (a'c' + b'd')\)。\(\blacksquare\)

\textbf{定义 4.2.5}：负元 \(-[a, b] = [b, a]\)；零元
\(0_{\mathbb{Z}} = [0, 0]\)；幺元 \(1_{\mathbb{Z}} = [1, 0]\)。

\paragraph{4.2.4
整数的代数性质}\label{ux6574ux6570ux7684ux4ee3ux6570ux6027ux8d28}

\textbf{定理 4.2.4}：\((\mathbb{Z}, +, 0_{\mathbb{Z}})\) 是阿贝尔群。

\textbf{证明}：结合律和交换律由自然数的加法运算律直接继承。零元：\([a, b] + [0, 0] = [a, b]\)。负元：\([a, b] + [b, a] = [a + b, b + a] = [0, 0]\)。\(\blacksquare\)

\textbf{定理 4.2.5}：\((\mathbb{Z}, \cdot, 1_{\mathbb{Z}})\)
满足结合律和交换律。

\textbf{证明}：由自然数运算律直接展开验证。\(\blacksquare\)

\textbf{定理
4.2.6（分配律）}：\(\forall x, y, z \in \mathbb{Z}: x \cdot (y + z) = x \cdot y + x \cdot z\)。

\textbf{证明}：设
\(x = [a, b], y = [c, d], z = [e, f]\)。展开后由自然数的分配律得证。\(\blacksquare\)

\textbf{推论}：\(\mathbb{Z}\) 是含幺交换环。

\textbf{定理 4.2.7}：\((-1) \times (-1) = 1\)。

\textbf{证明}：\(-1 = [0, 1]\)，\((-1) \times (-1) = [0, 1] \cdot [0, 1] = [0 \cdot 0 + 1 \cdot 1, 0 \cdot 1 + 1 \cdot 0] = [1, 0] = 1\)。\(\blacksquare\)

\textbf{推论
4.2.1（负负得正）}：\(\forall x, y \in \mathbb{Z}: (-x)(-y) = xy\)。

\textbf{证明}：先证 \((-1) \cdot x = -x\)。设 \(x = [a, b]\)，则
\((-1) \cdot x = [0, 1] \cdot [a, b] = [0 \cdot a + 1 \cdot b, 0 \cdot b + 1 \cdot a] = [b, a] = -x\)。因此
\((-x)(-y) = ((-1) \cdot x)((-1) \cdot y) = (-1)^2 \cdot xy = 1 \cdot xy = xy\)。\(\blacksquare\)

\textbf{定理
4.2.8（整数的消去律）}：\(\forall a, b, c \in \mathbb{Z}\)，若
\(c \neq 0\)，则 \(ac = bc \Rightarrow a = b\)。

\textbf{证明}：设 \(a = [a_1, a_2]\), \(b = [b_1, b_2]\),
\(c = [c_1, c_2]\)（\(c_1 \neq c_2\)）。\(ac = bc\) 意味着
\([a_1 c_1 + a_2 c_2, a_1 c_2 + a_2 c_1] = [b_1 c_1 + b_2 c_2, b_1 c_2 + b_2 c_1]\)。由等价关系定义和自然数消去律可得
\(a_1 + b_2 = a_2 + b_1\)，即
\([a_1, a_2] = [b_1, b_2]\)。\(\blacksquare\)

\paragraph{4.2.5
整数的序关系}\label{ux6574ux6570ux7684ux5e8fux5173ux7cfb}

\textbf{定义 4.2.6}：对 \([a, b], [c, d] \in \mathbb{Z}\)，定义
\([a, b] \leq [c, d]\) 当且仅当 \(a + d \leq b + c\)（在 \(\mathbb{N}\)
中）。

\textbf{定理 4.2.9}：\(\leq\) 是 \(\mathbb{Z}\)
上的全序关系，且与加法相容：若 \(x \leq y\)，则 \(x + z \leq y + z\)。

\textbf{证明}：良定义性：若 \([a, b] = [a', b']\) 且
\([c, d] = [c', d']\)，则 \(a + d \leq b + c\) 等价于
\(a' + d' \leq b' + c'\)（利用等价关系和自然数序的性质）。全序性：对任意
\([a, b], [c, d]\)，比较 \(a + d\) 和 \(b + c\)（在 \(\mathbb{N}\)
中），恰有 \(a + d < b + c\)、\(a + d = b + c\)、\(a + d > b + c\)
之一成立，对应
\([a, b] < [c, d]\)、\([a, b] = [c, d]\)、\([a, b] > [c, d]\)。与加法相容：若
\([a, b] \leq [c, d]\)，即 \(a + d \leq b + c\)，则对任意
\([e, f]\)，\((a + e) + (d + f) \leq (b + f) + (c + e)\)，即
\([a, b] + [e, f] \leq [c, d] + [e, f]\)。\(\blacksquare\)

\textbf{定理 4.2.10（整数序的三歧性）}：对任意
\(x, y \in \mathbb{Z}\)，恰有以下之一成立：\(x < y\)，\(x = y\)，\(x > y\)。

\textbf{证明}：设 \(x = [a, b]\), \(y = [c, d]\)。由自然数的三歧性（对
\(a + d\) 和 \(b + c\)），恰有
\(a + d < b + c\)、\(a + d = b + c\)、\(a + d > b + c\)
之一。\(\blacksquare\)

\textbf{定理 4.2.11（整数的阿基米德性质）}：对任意正整数 \(n\)
和任意整数 \(M\)，存在正整数 \(k\) 使得 \(kn > M\)。

\textbf{证明}：分两种情况。

\textbf{情形 1}：\(M \leq 0\)。取 \(k = 1\)，则
\(kn = n > 0 \geq M\)，结论成立。

\textbf{情形 2}：\(M > 0\)。设
\(n = [a, b]\)（\(a > b\)），\(M = [c, d]\)（\(c > d\)）。此时 \(a - b\)
和 \(c - d\) 在 \(\mathbb{N}\) 中均有定义。需 \(kn > M\)，即
\([ka, kb] > [c, d]\)，即 \(ka + d > kb + c\)，即
\(k(a - b) > c - d\)。由于 \(a - b > 0\)，需证存在自然数 \(k\) 使得
\(k(a - b) > c - d\)。先证自然数的无界性：对 \(m\) 做归纳证明
\(\forall m \in \mathbb{N},\, \exists k \in \mathbb{N}:\, k > m\)。基步：\(m = 0\)
时取 \(k = 1\)，则 \(1 > 0\)。归纳步：若存在 \(k > m\)，则
\(k + 1 > S(m)\)，故结论对 \(S(m)\) 也成立。再由无界性，取
\(m = c - d\)，存在 \(k_0 \in \mathbb{N}\) 使得 \(k_0 > c - d\)。因
\(a - b \geq 1\)，有 \(k_0(a - b) \geq k_0 > c - d\)，取 \(k = k_0\)
即可。\(\blacksquare\)

\bigskip\hrule\bigskip

\subsubsection{4.3 有理数}\label{ux6709ux7406ux6570}

\paragraph{4.3.1 构造动机}\label{ux6784ux9020ux52a8ux673a-1}

在 \(\mathbb{Z}\) 中，除法并不总是可行的：\(2x = 3\)
无整数解。我们扩展数系使非零数都有乘法逆元。

\paragraph{4.3.2
有理数的定义}\label{ux6709ux7406ux6570ux7684ux5b9aux4e49}

\textbf{定义 4.3.1}：令
\(P' = \{(a, b) \in \mathbb{Z} \times \mathbb{Z} : b \neq 0\}\)。在
\(P'\) 上定义关系 \(\approx\)： \[(a, b) \approx (c, d) \iff ad = bc\]

\textbf{定理 4.3.1}：\(\approx\) 是 \(P'\) 上的等价关系。

\textbf{证明}：自反性：\(ab = ba\)。对称性：若 \(ad = bc\) 则
\(cb = da\)。

传递性：若 \((a, b) \approx (c, d)\) 且 \((c, d) \approx (e, f)\)，即
\(ad = bc\) 且 \(cf = de\)。需证 \(af = be\)。

情形一：\(c \neq 0\)。由 \(ad = bc\) 和 \(cf = de\)，两式相乘得
\((ad)(cf) = (bc)(de)\)，即 \(acdf = bcde\)。因 \(c \neq 0\)，消去 \(c\)
得 \(adf = bde\)。因 \(d \neq 0\)，消去 \(d\) 得 \(af = be\)。

情形二：\(c = 0\)。由 \(ad = bc = 0\) 且 \(d \neq 0\) 得 \(a = 0\)。由
\(cf = de\) 即 \(0 = de\) 且 \(d \neq 0\) 得 \(e = 0\)。故
\(af = 0 = be\)。

综上，\((a, b) \approx (e, f)\)。\(\blacksquare\)

\textbf{定义 4.3.2}：\textbf{有理数集} \(\mathbb{Q} = P' / \approx\)。将
\([(a, b)]\) 简记为 \(\frac{a}{b}\)。

\paragraph{4.3.3
有理数的运算}\label{ux6709ux7406ux6570ux7684ux8fd0ux7b97}

\textbf{定义
4.3.3（加法）}：\(\frac{a}{b} + \frac{c}{d} = \frac{ad + bc}{bd}\)

\textbf{定义
4.3.4（乘法）}：\(\frac{a}{b} \cdot \frac{c}{d} = \frac{ac}{bd}\)

\textbf{定义 4.3.5（负元）}：\(-\frac{a}{b} = \frac{-a}{b}\)

\textbf{定义 4.3.6（逆元）}：对
\(\frac{a}{b} \neq 0\)（\(a \neq 0\)），\(\left(\frac{a}{b}\right)^{-1} = \frac{b}{a}\)

\textbf{定理 4.3.2}：上述运算都是良定义的。

\textbf{证明}：

\textbf{加法良定义性}：设 \(\frac{a}{b} = \frac{a'}{b'}\)（即
\(ab' = a'b\)）且 \(\frac{c}{d} = \frac{c'}{d'}\)（即
\(cd' = c'd\)）。要证 \(\frac{ad+bc}{bd} = \frac{a'd'+b'c'}{b'd'}\)，即
\((ad+bc)(b'd') = (a'd'+b'c')(bd)\)。展开左端
\(= abb'd'd + bcb'd' = a'bb'd + bb'c'd = (a'd' + b'c')(bd)\)（利用
\(ab' = a'b\) 和 \(cd' = c'd\)）。

\textbf{乘法良定义性}：要证 \(\frac{ac}{bd} = \frac{a'c'}{b'd'}\)，即
\(ac \cdot b'd' = a'c' \cdot bd\)。由 \(ab' = a'b\) 和
\(cd' = c'd\)，相乘得 \((ab')(cd') = (a'b)(c'd)\)，即
\(ac \cdot b'd' = a'c' \cdot bd\)。

\textbf{逆元良定义性}：设
\(\frac{a}{b} = \frac{a'}{b'}\)（\(ab' = a'b\)）且 \(a \neq 0\)。先证
\(a' \neq 0\)：若 \(a' = 0\) 则 \(ab' = 0\)，因 \(b' \neq 0\) 得
\(a = 0\)，矛盾。要证 \(\frac{b}{a} = \frac{b'}{a'}\)，即
\(ba' = b'a\)。由 \(ab' = a'b\) 交换得 \(a'b = ab'\)，即
\(ba' = b'a\)。\(\blacksquare\)

\textbf{定义 4.3.7（有理数的序关系）}：设 \(\frac{a}{b}\) 和
\(\frac{c}{d}\) 分别取 \(b > 0\)、\(d > 0\) 的标准表示，定义
\(\frac{a}{b} \leq \frac{c}{d}\) 当且仅当 \(ad \leq bc\)（在
\(\mathbb{Z}\) 中）。

\textbf{良定义性}：若 \(\frac{a}{b} = \frac{a'}{b'}\)（\(ab' = a'b\)）且
\(\frac{c}{d} = \frac{c'}{d'}\)（\(cd' = c'd\)），\(b, d, b', d' > 0\)，则
\(ad \leq bc \iff adb'd' \leq bcb'd' \iff a'b \cdot cd' \leq ab' \cdot c'd \iff a'd' \leq b'c'\)。故定义与代表元选取无关。

\textbf{三歧性}：对任意
\(\frac{a}{b}, \frac{c}{d}\)（\(b, d > 0\)），由整数的三歧性，\(ad < bc\)、\(ad = bc\)、\(ad > bc\)
恰有一成立，对应
\(\frac{a}{b} < \frac{c}{d}\)、\(\frac{a}{b} = \frac{c}{d}\)、\(\frac{a}{b} > \frac{c}{d}\)。

\textbf{与运算的相容性}：若
\(\frac{a}{b} \leq \frac{c}{d}\)（\(b, d > 0\)），则对任意
\(\frac{e}{f}\)（\(f > 0\)）：(1)
\(\frac{a}{b} + \frac{e}{f} \leq \frac{c}{d} + \frac{e}{f}\)（由
\(ad \leq bc\) 两边加 \(bef \cdot d\) 即得）；(2) 若
\(\frac{e}{f} > 0\)（即 \(e > 0\)），则
\(\frac{a}{b} \cdot \frac{e}{f} \leq \frac{c}{d} \cdot \frac{e}{f}\)（由
\(ad \leq bc\) 两边乘 \(ef > 0\) 即得）。

\paragraph{4.3.4
有理数构成域}\label{ux6709ux7406ux6570ux6784ux6210ux57df}

\textbf{定理 4.3.3}：\((\mathbb{Q}, +, \cdot)\) 是有序域。

\textbf{证明}：

加法阿贝尔群：结合律、交换律、零元、负元均由整数的性质继承。

乘法阿贝尔群（除零外）：结合律、交换律、幺元 \(\frac{1}{1}\)、逆元
\(\frac{b}{a}\) 均可直接验证。

分配律：\(\frac{a}{b} \cdot \left(\frac{c}{d} + \frac{e}{f}\right) = \frac{a(cf + de)}{bdf} = \frac{ac}{bd} + \frac{ae}{bf}\)。\(\blacksquare\)

\paragraph{4.3.5
有理数的稠密性}\label{ux6709ux7406ux6570ux7684ux7a20ux5bc6ux6027}

\textbf{定理
4.3.4（有理数的稠密性）}：任意两个不同的有理数之间存在有理数。即若
\(r < s\)，则存在 \(t\) 使得 \(r < t < s\)。

\textbf{证明}：取
\(t = \frac{r + s}{2}\)。\(t - r = \frac{s - r}{2} > 0\)，\(s - t = \frac{s - r}{2} > 0\)，故
\(r < t < s\)。\(\blacksquare\)

\textbf{定理 4.3.5（有理数的阿基米德性质）}：对任意正有理数
\(r, s\)，存在正整数 \(n\) 使得 \(nr > s\)。

\textbf{证明}：设 \(r = a/b\), \(s = c/d\)（\(a, b, c, d > 0\)）。需
\(n \cdot a/b > c/d\)，即 \(nad > bc\)。由于 \(a, b, c, d\)
均为正整数，\(bc\) 和 \(ad\)
均为确定的正整数。由自然数的无界性（对任意自然数 \(m\)，存在 \(k\) 使得
\(k > m\)，这可由归纳法证明），取 \(n\) 为大于 \(bc/(ad)\)
的正整数即可。\(\blacksquare\)

\textbf{定理 4.3.6（有理数对加法和乘法封闭）}：\(\mathbb{Q}\)
对四则运算封闭（除数非零时）。

\textbf{证明}：\(\frac{a}{b} + \frac{c}{d} = \frac{ad + bc}{bd} \in \mathbb{Q}\)（\(bd \neq 0\)）。\(\frac{a}{b} \cdot \frac{c}{d} = \frac{ac}{bd} \in \mathbb{Q}\)。\(\blacksquare\)

\bigskip\hrule\bigskip

\subsubsection{4.4 实数}\label{ux5b9eux6570}

\paragraph{4.4.1 构造动机}\label{ux6784ux9020ux52a8ux673a-2}

\textbf{定理 4.4.1}：不存在有理数 \(r\) 使得 \(r^2 = 2\)。

\textbf{证明}：假设 \(r = \frac{p}{q}\)（\(p, q\) 互素），则
\(p^2 = 2q^2\)。\(p^2\) 为偶数，故 \(p\) 为偶数，设 \(p = 2k\)。代入得
\(q^2 = 2k^2\)，\(q\) 也为偶数。与 \(\gcd(p, q) = 1\)
矛盾。\(\blacksquare\)

这表明 \(\mathbb{Q}\) 有''漏洞''。我们需要构造完备的数系。

\paragraph{4.4.2 戴德金分割}\label{ux6234ux5fb7ux91d1ux5206ux5272}

\textbf{定义 4.4.1（戴德金分割）}：\(\mathbb{Q}\)
的一个\textbf{戴德金分割}是一个有序对 \((A, B)\)，其中 \(A, B\) 是
\(\mathbb{Q}\) 的非空子集，满足： 1.
\(A \cup B = \mathbb{Q}\)，\(A \cap B = \varnothing\) 2.
\(\forall a \in A, \forall b \in B: a < b\) 3. \(A\) 没有最大元素

\paragraph{4.4.3 实数的定义}\label{ux5b9eux6570ux7684ux5b9aux4e49}

\textbf{定义 4.4.2}：\textbf{实数集} \(\mathbb{R}\)
定义为所有戴德金分割的集合。有理数 \(r\) 对应分割
\(A_r = \{q \in \mathbb{Q} : q < r\}\)。

\textbf{引理 4.4.1（嵌入的保序性）}：映射 \(r \mapsto A_r\) 是
\(\mathbb{Q}\) 到 \(\mathbb{R}\) 的保序嵌入，即
\(r < s \iff A_r \subsetneq A_s\)。

\textbf{证明}：\((\Rightarrow)\)：若 \(r < s\)，取
\(t = \frac{r+s}{2}\)，则 \(r < t < s\)，\(t \in A_s \setminus A_r\)，故
\(A_r \subsetneq A_s\)。\((\Leftarrow)\)：若 \(A_r \subsetneq A_s\)，取
\(q \in A_s \setminus A_r\)，则 \(q < s\) 且 \(q \geq r\)，故
\(r \leq q < s\)。\(\blacksquare\)

\paragraph{4.4.4 实数的运算}\label{ux5b9eux6570ux7684ux8fd0ux7b97}

\textbf{定义 4.4.3（加法）}：\(\alpha + \beta\) 的下组为
\(\{a_1 + a_2 : a_1 \in A_1, a_2 \in A_2\}\)。

\textbf{验证加法良定义}：(1)
\(A = \{a_1 + a_2 : a_1 \in A_1, a_2 \in A_2\}\) 非空。(2) 若
\(q \in A\)，取 \(q' < q\)，则 \(q' = a_1 + a_2 - \epsilon\)，可写成
\(a_1' + a_2'\) 的形式，故 \(q' \in A\)。(3) \(A\) 无最大元：对
\(a_1 + a_2 \in A\)，取 \(a_1' > a_1\)（\(a_1' \in A_1\)），则
\(a_1' + a_2 > a_1 + a_2\) 且 \(a_1' + a_2 \in A\)。(4)
\(B = \mathbb{Q} \setminus A\) 非空：取 \(b_1 \in B_1, b_2 \in B_2\)，则
\(b_1 + b_2 \notin A\)（因为若 \(b_1 + b_2 = a_1 + a_2\)，则
\(b_1 - a_1 = a_2 - b_2 > 0\)，但 \(b_1 - a_1 > 0\) 且
\(a_2 - b_2 < 0\)，矛盾）。(5) 若 \(a \in A, b \in B\)，则 \(a < b\)：设
\(a = a_1 + a_2\), \(b = b_1 + b_2\)（\(b_1 \notin A_1\) 或
\(b_2 \notin A_2\)）。不妨设 \(b_1 \notin A_1\)，则
\(a_1 < b_1\)，\(a_2 \leq b_2\)，故 \(a < b\)。

\textbf{定义 4.4.4（零元与幺元）}：\(0_{\mathbb{R}}\) 的下组为
\(\{q \in \mathbb{Q} : q < 0\}\)；\(1_{\mathbb{R}}\) 的下组为
\(\{q \in \mathbb{Q} : q < 1\}\)。

\textbf{定义 4.4.5（乘法）}：设
\(\alpha = (A_1, B_1)\)，\(\beta = (A_2, B_2)\) 为实数（其中 \(A_i\)
为下组，\(B_i = \mathbb{Q} \setminus A_i\) 为上组）。

\textbf{(i) 正实数的乘法}：若 \(\alpha > 0\) 且 \(\beta > 0\)（即
\(A_1, A_2\) 均包含正有理数），定义 \(\alpha \cdot \beta\) 的下组为：
\[\{q \in \mathbb{Q} : q \leq 0\} \cup \{ab : a \in A_1, a > 0, b \in A_2, b > 0\}\]

该集合的向下封闭性由 \(A_1, A_2\) 的向下封闭性保证：若
\(ab\)（\(a \in A_1, b \in A_2, a, b > 0\)）在集合中，且
\(0 < q < ab\)，则 \(q = (q/b) \cdot b\)，其中 \(0 < q/b < a\)，由
\(A_1\) 的向下封闭性得 \(q/b \in A_1\)，故 \(q\) 也在集合中。

\textbf{(ii) 一般实数的乘法}：通过绝对值和符号规则推广。对
\(\alpha < 0\)，定义 \(|\alpha|\) 的下组为
\(\{-b : b \in B_1, b < 0\} \cup \{q \in \mathbb{Q} : q \leq 0\}\)（即
\(|\alpha| = -\alpha\)）。则： - 若
\(\alpha > 0, \beta > 0\)：\(\alpha \cdot \beta\) 如上定义 - 若
\(\alpha > 0, \beta < 0\)：\(\alpha \cdot \beta = -(\alpha \cdot |\beta|)\)
- 若
\(\alpha < 0, \beta > 0\)：\(\alpha \cdot \beta = -(|\alpha| \cdot \beta)\)
- 若
\(\alpha < 0, \beta < 0\)：\(\alpha \cdot \beta = |\alpha| \cdot |\beta|\)
- 若 \(\alpha = 0\) 或
\(\beta = 0\)：\(\alpha \cdot \beta = 0_{\mathbb{R}}\)

\textbf{乘法良定义的验证}：(1) 下组非空（\(0\) 属之）。(2)
向下封闭性如上所示。(3) 无最大元：对
\(ab\)（\(a \in A_1, b \in A_2, a, b > 0\)），取 \(a' \in A_1\) 使
\(a' > a\)（\(A_1\) 无最大元），则 \(a'b > ab\) 且 \(a'b\)
仍在下组中。(4) 补集非空：取
\(b_1 \in B_1, b_2 \in B_2\)（\(b_1, b_2 > 0\)），则 \(b_1 b_2 > 0\) 且
\(b_1 b_2\) 不在下组中------若
\(b_1 b_2 = a_1 a_2\)（\(a_i \in A_i, a_i > 0\)），则
\(b_1/a_1 = a_2/b_2\)；由 \(a_1 < b_1\) 得 \(b_1/a_1 > 1\)，由
\(a_2 < b_2\) 得 \(a_2/b_2 < 1\)，矛盾。(5)
下组中元素严格小于补集中元素，由有理数乘法的序保持性可得。一般情况的良定义性由正实数情况和符号规则保证。

\textbf{定理 4.4.2}：\((\mathbb{R}, +, \cdot)\) 是有序域。

\textbf{证明}：域公理的验证均通过戴德金分割的定义和有理数的性质直接得到。以下给出关键步骤：

\textbf{加法结合律}：设 \(\alpha = (A_1, B_1)\), \(\beta = (A_2, B_2)\),
\(\gamma = (A_3, B_3)\)。\((\alpha + \beta) + \gamma\) 的下组为
\(\{(a_1 + a_2) + a_3 : a_1 \in A_1, a_2 \in A_2, a_3 \in A_3\}\)。\(\alpha + (\beta + \gamma)\)
的下组为
\(\{a_1 + (a_2 + a_3) : a_1 \in A_1, a_2 \in A_2, a_3 \in A_3\}\)。由有理数加法结合律，两者相等。

\textbf{加法零元}：\(\alpha + 0_{\mathbb{R}}\) 的下组为
\(\{a + q : a \in A_1, q < 0\}\)。对任意 \(r \in A_1\)，取 \(a = r\) 和
\(q = -\epsilon\)（\(\epsilon > 0\) 足够小），则
\(r - \epsilon < r\)，即 \(r - \epsilon\) 可以逼近 \(A_1\)
中任意元素。由 \(A_1\) 无最大元的性质可验证下组为 \(A_1\)。

\textbf{乘法逆元}：分两种情况。

\textbf{(a) \(\alpha > 0\) 的情形}：设
\(\alpha = (A_1, B_1)\)，\(B_1 = \mathbb{Q} \setminus A_1\)。\(A_1\)
包含正有理数且无最大元，\(B_1\) 无最小元且
\(\forall a \in A_1, b \in B_1\) 有 \(a < b\)。直接构造 \(\alpha^{-1}\)
的分割如下：

\[\alpha^{-1} = (\{q \in \mathbb{Q} : q \leq 0\} \cup \{q \in \mathbb{Q} : q > 0 \text{ 且 } 1/q \notin A_1\},\; \mathbb{Q} \setminus \text{下组})\]

记
\(C = \{q \in \mathbb{Q} : q \leq 0\} \cup \{q \in \mathbb{Q} : q > 0 \text{ 且 } 1/q \notin A_1\}\)。由
\(A_1\) 无最大元且向下封闭可验证 \(C\) 无最大元且向下封闭，故 \(C\) 确为
Dedekind 分割的下组。

验证 \(\alpha \cdot \alpha^{-1} = 1_{\mathbb{R}}\)（即乘积下组为
\(\{q : q < 1\}\)）。由定义 4.4.5(i) 中正实数乘法，乘积下组为
\(\{q \leq 0\} \cup \{xy : x \in A_1, x > 0, y \in C, y > 0\}\)。

\((\subseteq)\)：设 \(xy\)
为乘积下组中正元素（\(x \in A_1, x > 0, y \in C, y > 0\)）。由 \(C\)
的定义，\(1/y \notin A_1\)，即 \(1/y \in B_1\)。若 \(xy \geq 1\)，则
\(x \geq 1/y\)。但 \(A_1\) 向下封闭，\(1/y \in B_1\) 且 \(x \in A_1\) 与
\(x \geq 1/y\) 矛盾（因 \(A_1\) 和 \(B_1\) 满足
\(\forall a \in A_1, b \in B_1, a < b\)）。故 \(xy < 1\)，即
\(xy \in \{q : q < 1\}\)。

\((\supseteq)\)：设 \(0 < q < 1\)。考虑集合
\(T = \{x \in A_1 : x > 0 \text{ 且 } x/q \in A_1\}\)。若 \(T\) 等于
\(A_1\) 中全体正有理数，则对任意 \(x \in A_1\)（\(x > 0\)）有
\(x/q \in A_1\)。反复应用得 \(x/q^n \in A_1\) 对所有
\(n \in \mathbb{N}\) 成立。由于 \(0 < q < 1\)，\(q^n \to 0\)，故
\(x/q^n \to \infty\)，这意味着 \(A_1\) 无上界------与 Dedekind
分割的定义矛盾。因此 \(T \subsetneqq A_1 \cap \mathbb{Q}^+\)，即存在
\(x_0 \in A_1\)、\(x_0 > 0\) 但 \(x_0/q \notin A_1\)。令
\(y = q/x_0 > 0\)，则 \(1/y = x_0/q \notin A_1\)，故 \(y \in C\)。于是
\(x_0 y = q\) 属于乘积的下组。

因此 \(\alpha \cdot \alpha^{-1}\) 的下组为 \(\{q : q < 1\}\)，即
\(1_{\mathbb{R}}\)。

\textbf{(b) \(\alpha < 0\) 的情形}：定义
\(\alpha^{-1} = -(|\alpha|^{-1})\)。由 (a) 知 \(|\alpha|^{-1}\)
存在（\(|\alpha| > 0\)），故 \(\alpha^{-1}\) 存在。由定义 4.4.5(ii)
中负负得正规则：
\[\alpha \cdot \alpha^{-1} = \alpha \cdot (-(|\alpha|^{-1})) = -(|\alpha| \cdot (|\alpha|^{-1})) = -(-1_{\mathbb{R}}) = 1_{\mathbb{R}}\]
（其中 \(|\alpha| \cdot (|\alpha|^{-1}) = 1_{\mathbb{R}}\) 由 (a)
保证，而 \(-(-1_{\mathbb{R}}) = 1_{\mathbb{R}}\)
由加法逆元性质可得）。\(\blacksquare\)

\paragraph{4.4.5 完备性}\label{ux5b8cux5907ux6027}

\textbf{定理
4.4.3（有理数在实数中稠密）}：任意两个不同的实数之间存在有理数。

\textbf{证明}：设 \(\alpha < \beta\)，即
\(A_\alpha \subsetneq A_\beta\)（\(A_\alpha\) 是 \(A_\beta\)
的真子集）。取 \(q \in A_\beta \setminus A_\alpha\)。因
\(q \notin A_\alpha\) 且 \(A_\alpha\) 向下封闭，\(q \geq \alpha\)。因
\(q \in A_\beta\)，\(q < \beta\)。

若 \(q > \alpha\)，则 \(\alpha < q < \beta\)，\(q \in \mathbb{Q}\)
即为所求。

若 \(q = \alpha\)（即 \(\alpha\) 是有理数），则
\(\alpha \in A_\beta\)。因 \(A_\beta\) 无最大元，存在 \(q' \in A_\beta\)
使得 \(q' > q = \alpha\)。于是
\(\alpha < q' < \beta\)，\(q' \in \mathbb{Q}\) 即为所求。

因此任意两个实数之间存在有理数。\(\blacksquare\)

\paragraph{4.4.6 确界原理}\label{ux786eux754cux539fux7406}

\textbf{定义 4.4.6（上确界）}：\(S \subseteq \mathbb{R}\)
的\textbf{上确界} \(\sup S\) 是 \(S\) 的最小上界。

\textbf{定理 4.4.4（确界原理）}：\(\mathbb{R}\)
的每个非空有上界的子集都有上确界。

\textbf{证明}：设 \(S \subseteq \mathbb{R}\) 非空有上界。对每个
\(\alpha \in S\)，设 \(\alpha\) 对应戴德金分割的下组 \(A_\alpha\)。令
\(A^* = \bigcup_{\alpha \in S} A_\alpha\)。我们验证
\((A^*, \mathbb{Q} \setminus A^*)\) 是一个戴德金分割，从而定义了某个实数
\(\beta\)：

\begin{enumerate}
\def\labelenumi{\arabic{enumi}.}
\tightlist
\item
  \textbf{\(A^*\) 非空}：因 \(S\) 非空，取
  \(\alpha \in S\)，\(A_\alpha\) 非空，故 \(A^*\) 非空。
\item
  \textbf{\(\mathbb{Q} \setminus A^*\) 非空}：设 \(u\) 是 \(S\)
  的上界（即对所有 \(\alpha \in S\) 有 \(\alpha \leq u\)）。则 \(u\)
  对应的下组 \(A_u\) 的补集（即
  \(\{q \in \mathbb{Q} : q \geq u\}\)）中至少有 \(u\) 本身不在 \(A^*\)
  中（因为 \(A^* = \bigcup A_\alpha\) 中的每个元素都属于某个
  \(A_\alpha\)，即严格小于某个 \(\alpha\)，而 \(\alpha \leq u\) 意味着
  \(A_\alpha \subseteq A_u\)，故 \(u \notin A^*\)）。
\item
  \textbf{\(A^*\) 是下集}：若 \(q \in A^*\)，\(p < q\)，则
  \(q \in A_\alpha\) 对某个 \(\alpha\)，由 \(A_\alpha\) 是下集得
  \(p \in A_\alpha \subseteq A^*\)。
\item
  \textbf{\(A^*\) 无最大元}：若 \(q \in A^*\)，则 \(q \in A_\alpha\)
  对某个 \(\alpha\)。因 \(A_\alpha\) 无最大元，存在 \(q' \in A_\alpha\)
  使得 \(q' > q\)，故 \(q' \in A^*\)。
\end{enumerate}

因此 \(\beta = (A^*, \mathbb{Q} \setminus A^*)\) 是一个实数。接下来验证
\(\beta\) 是 \(S\) 的上确界。

\textbf{\(\beta\) 是 \(S\) 的上界}：对任意
\(\alpha \in S\)，\(A_\alpha \subseteq A^*\)，即 \(\alpha \leq \beta\)。

\textbf{\(\beta\) 是最小上界}：设 \(\gamma < \beta\) 是任意实数。则
\(A_\gamma \subsetneq A^*\)，存在 \(q \in A^* \setminus A_\gamma\)。由
\(q \in A^*\)，\(q \in A_\alpha\) 对某个 \(\alpha \in S\)。但
\(q \notin A_\gamma\)，故 \(\gamma < q \leq \alpha\)（因为
\(q \in A_\alpha\)），即 \(\gamma\) 不是 \(S\) 的上界。

综上，\(\beta = \sup S\)。\(\blacksquare\)

\textbf{注记}：确界原理是实数完备性的核心表述，它等价于戴德金分割构造实数的完备性保证。在此基础上，可以证明单调有界定理、区间套定理、有限覆盖定理和柯西收敛准则等分析学的基本定理。

\textbf{定理 4.4.5}：存在唯一的正实数 \(\xi\) 使得 \(\xi^2 = 2\)。

\textbf{证明}：令 \(S = \{q \in \mathbb{Q} : q > 0, q^2 < 2\}\)，\(S\)
非空（\(1 \in S\)）有上界（如 \(2\) 是上界：若 \(q \geq 2\) 则
\(q^2 \geq 4 > 2\)）。设 \(\xi = \sup S\)（由确界原理存在）。

若 \(\xi^2 < 2\)：取 \(h = \frac{2 - \xi^2}{2\xi + 1} > 0\)。先验证
\(h < 1\)：这等价于 \(2 - \xi^2 < 2\xi + 1\)，即
\(1 < \xi^2 + 2\xi = \xi(\xi + 2)\)。由于 \(\xi \geq 1\)（因
\(1 \in S\)），\(\xi(\xi + 2) \geq 1 \cdot 3 = 3 > 1\)，故 \(h < 1\)
成立。因此
\((\xi + h)^2 = \xi^2 + 2\xi h + h^2 = \xi^2 + h(2\xi + h) < \xi^2 + h(2\xi + 1) = \xi^2 + (2 - \xi^2) = 2\)。故
\(\xi + h \in S\)，与 \(\xi\) 是上界矛盾。

若 \(\xi^2 > 2\)：取 \(h = \frac{\xi^2 - 2}{2\xi} > 0\)。则
\((\xi - h)^2 = \xi^2 - 2\xi h + h^2 > \xi^2 - 2\xi h = 2\)。故
\(\xi - h\) 也是 \(S\) 的上界，与 \(\xi\) 是最小上界矛盾。

因此 \(\xi^2 = 2\)。

唯一性：设 \(0 < x < y\)，\(x^2 = y^2 = 2\)。则
\(0 = y^2 - x^2 = (y-x)(y+x)\)。由 \(y + x > 0\) 得 \(y - x = 0\)，即
\(x = y\)。\(\blacksquare\)

\textbf{定理 4.4.6（阿基米德性质）}：对任意正实数 \(x, y\)，存在自然数
\(n\) 使得 \(nx > y\)。

\textbf{证明}：假设不然，\(S = \{nx : n \in \mathbb{N}\}\) 有上界
\(y\)。设 \(\xi = \sup S\)（由确界原理存在）。由于
\(x > 0\)，\(\xi - x < \xi\)，故存在 \(n_0 x > \xi - x\)（否则
\(\xi - x\) 也是上界，与 \(\xi\) 是最小上界矛盾），即
\((n_0 + 1)x > \xi\)，但 \((n_0 + 1)x \in S\)，与 \(\xi\)
是上界矛盾。\(\blacksquare\)

\textbf{推论 4.4.1}：对任意正实数 \(\epsilon\)，存在自然数 \(n\) 使得
\(\frac{1}{n} < \epsilon\)。

\textbf{证明}：在阿基米德性质中取 \(x = \epsilon\), \(y = 1\)，得
\(n\epsilon > 1\)，即 \(\frac{1}{n} < \epsilon\)。\(\blacksquare\)

\textbf{推论
4.4.2（实数的阿基米德性质的等价形式）}：\(\lim_{n \to \infty} \frac{1}{n} = 0\)。

\textbf{证明}：对任意 \(\epsilon > 0\)，由推论 4.4.1 存在 \(N\) 使得
\(\frac{1}{N} < \epsilon\)。对
\(n \geq N\)，\(0 < \frac{1}{n} \leq \frac{1}{N} < \epsilon\)。\(\blacksquare\)

\textbf{注记
4.4.1}：阿基米德性质表明实数系中没有''无穷小''的正实数。这是实数系与非标准分析中的超实数系的关键区别。

\bigskip\hrule\bigskip

\subsubsection{4.5 复数}\label{ux590dux6570}

\paragraph{4.5.1 构造动机}\label{ux6784ux9020ux52a8ux673a-3}

在 \(\mathbb{R}\) 中，\(x^2 + 1 = 0\) 无解（因
\(x^2 \geq 0\)）。我们构造更大的数系使所有多项式方程有解。

\paragraph{4.5.2 复数的定义}\label{ux590dux6570ux7684ux5b9aux4e49}

\textbf{定义
4.5.1}：\(\mathbb{C} = \mathbb{R} \times \mathbb{R}\)，运算定义为： -
\((a, b) + (c, d) = (a + c, b + d)\) -
\((a, b) \cdot (c, d) = (ac - bd, ad + bc)\) -
\(0_{\mathbb{C}} = (0, 0)\)，\(1_{\mathbb{C}} = (1, 0)\)

\textbf{定义 4.5.2}：虚数单位 \(i = (0, 1)\)。

\textbf{定理 4.5.1}：\(i^2 = -1\)。

\textbf{证明}：\(i^2 = (0, 1) \cdot (0, 1) = (0 \cdot 0 - 1 \cdot 1, 0 \cdot 1 + 1 \cdot 0) = (-1, 0) = -1\)。\(\blacksquare\)

\textbf{定理 4.5.2}：每个复数 \(z = (a, b)\) 都可以唯一地表示为
\(z = a + bi\)。

\textbf{证明}：\((a, b) = (a, 0) + (b, 0) \cdot (0, 1) = a + bi\)。唯一性：若
\(a + bi = c + di\)，则 \((a, b) = (c, d)\)，故
\(a = c, b = d\)。\(\blacksquare\)

\paragraph{4.5.3 复数构成域}\label{ux590dux6570ux6784ux6210ux57df}

\textbf{定理 4.5.3}：\((\mathbb{C}, +, \cdot)\) 是域。

\textbf{证明}：加法阿贝尔群由实数性质直接继承。乘法结合律通过展开验证。逆元：对
\(z = a + bi \neq 0\)，\(z^{-1} = \frac{a - bi}{a^2 + b^2}\)。验证
\(z \cdot z^{-1} = \frac{a^2 + b^2}{a^2 + b^2} = 1\)。\(\blacksquare\)

\paragraph{4.5.4 模与共轭}\label{ux6a21ux4e0eux5171ux8f6d}

\textbf{定义 4.5.3（共轭）}：\(\bar{z} = a - bi\)。

\textbf{定义 4.5.4（模）}：\(|z| = \sqrt{a^2 + b^2}\)。

\textbf{定理 4.5.4（共轭的性质）}： 1.
\(\overline{z_1 + z_2} = \bar{z}_1 + \bar{z}_2\) 2.
\(\overline{z_1 \cdot z_2} = \bar{z}_1 \cdot \bar{z}_2\) 3.
\(z \cdot \bar{z} = |z|^2\)

\textbf{证明}：(1) 和 (2) 由定义直接计算。(3)
\(z \cdot \bar{z} = (a + bi)(a - bi) = a^2 + b^2 = |z|^2\)。\(\blacksquare\)

\textbf{定理 4.5.5（模的性质）}： 1. \(|z| \geq 0\)，且
\(|z| = 0 \iff z = 0\) 2. \(|z_1 \cdot z_2| = |z_1| \cdot |z_2|\) 3.
\(|z_1 + z_2| \leq |z_1| + |z_2|\)（三角不等式）

\textbf{证明}：(1) 由定义直接得到。(2)
\(|z_1 z_2|^2 = (z_1 z_2)\overline{(z_1 z_2)} = z_1 z_2 \bar{z}_1 \bar{z}_2 = (z_1 \bar{z}_1)(z_2 \bar{z}_2) = |z_1|^2 |z_2|^2\)。(3)
\(|z_1 + z_2|^2 = |z_1|^2 + 2\operatorname{Re}(z_1 \bar{z}_2) + |z_2|^2 \leq |z_1|^2 + 2|z_1||z_2| + |z_2|^2 = (|z_1| + |z_2|)^2\)。\(\blacksquare\)

\paragraph{4.5.5
复数不是有序域}\label{ux590dux6570ux4e0dux662fux6709ux5e8fux57df}

\textbf{定理 4.5.6}：不存在全序关系 ``\(<\)'' 使得
\((\mathbb{C}, +, \cdot, <)\) 成为有序域。

\textbf{证明}：假设存在。由于 \(i \neq 0\)，必有 \(i > 0\) 或
\(i < 0\)。若
\(i > 0\)，由有序域中正数的乘积为正，\(i^2 = -1 > 0\)。再由
\((-1)(-1) = 1 > 0\)，以及正数之和为正，\(0 = (-1) + 1 > 0\)，但同时
\(1 > 0\)，与三歧性矛盾（\(0 > 0\) 不可能）。若 \(i < 0\)，则
\(-i > 0\)，\((-i)^2 = i^2 = -1 > 0\)，同理推出矛盾。\(\blacksquare\)

\paragraph{4.5.6
复数的极坐标表示}\label{ux590dux6570ux7684ux6781ux5750ux6807ux8868ux793a}

\textbf{定义 4.5.5}：对复数 \(z = a + bi\)，定义其\textbf{辐角}
\(\theta\) 为满足
\(\cos\theta = \frac{a}{|z|}\)，\(\sin\theta = \frac{b}{|z|}\)
的角度（当 \(z \neq 0\) 时）。则 \(z\) 可表示为
\(z = |z|(\cos\theta + i\sin\theta)\)。

\textbf{定理 4.5.7（复数乘法的几何意义）}：设
\(z_1 = r_1 e^{i\theta_1}\), \(z_2 = r_2 e^{i\theta_2}\)。则
\(z_1 z_2 = r_1 r_2 e^{i(\theta_1 + \theta_2)}\)。即复数相乘时，模相乘，辐角相加。

\textbf{证明}：\(z_1 z_2 = r_1(\cos\theta_1 + i\sin\theta_1) \cdot r_2(\cos\theta_2 + i\sin\theta_2)\)。展开得
\(r_1 r_2[(\cos\theta_1\cos\theta_2 - \sin\theta_1\sin\theta_2) + i(\sin\theta_1\cos\theta_2 + \cos\theta_1\sin\theta_2)]\)。由三角函数的和角公式，\(= r_1 r_2[\cos(\theta_1 + \theta_2) + i\sin(\theta_1 + \theta_2)]\)。\(\blacksquare\)

\textbf{定理 4.5.8（棣莫弗定理）}：对任意正整数 \(n\) 和实数
\(\theta\)：
\[(\cos\theta + i\sin\theta)^n = \cos(n\theta) + i\sin(n\theta)\]

\textbf{证明}：对 \(n\) 用数学归纳法。\(n = 1\) 时显然。设
\((\cos\theta + i\sin\theta)^k = \cos(k\theta) + i\sin(k\theta)\)。则
\[(\cos\theta + i\sin\theta)^{k+1} = (\cos(k\theta) + i\sin(k\theta))(\cos\theta + i\sin\theta)\]
由定理
4.5.7，\(= \cos((k+1)\theta) + i\sin((k+1)\theta)\)。\(\blacksquare\)

\textbf{定理 4.5.9（单位根）}：方程 \(z^n = 1\) 恰有 \(n\)
个不同的复数解，称为 \(n\) 次\textbf{单位根}：
\[z_k = e^{2\pi ik/n} = \cos\frac{2\pi k}{n} + i\sin\frac{2\pi k}{n}, \quad k = 0, 1, \ldots, n-1\]

\textbf{证明}：设 \(z = re^{i\theta}\)。\(z^n = r^n e^{in\theta} = 1\)
意味着 \(r^n = 1\)（故 \(r = 1\)）且
\(n\theta = 2\pi k\)（\(k \in \mathbb{Z}\)），即
\(\theta = 2\pi k/n\)。不同的 \(k\) 给出不同的 \(z_k\) 当且仅当 \(k\)
在模 \(n\) 意义下不同，故恰有 \(n\) 个解。\(\blacksquare\)

\paragraph{4.5.7
代数基本定理（叙述）}\label{ux4ee3ux6570ux57faux672cux5b9aux7406ux53d9ux8ff0}

\textbf{定理 4.5.10（代数基本定理）}：每个次数 \(\geq 1\)
的复系数多项式在 \(\mathbb{C}\) 中至少有一个根。

\textbf{等价形式}：每个次数 \(\geq 1\) 的复系数多项式在 \(\mathbb{C}\)
中可以完全分解为一次因式的乘积。

\textbf{注记
4.5.1}：代数基本定理的严格证明需要复分析（如刘维尔定理或辐角原理），超出本章范围。但此定理的结论极其重要：它表明
\(\mathbb{C}\) 是\textbf{代数封闭域}------任何多项式方程在
\(\mathbb{C}\) 中都有解，不需要再扩展数系。

\bigskip\hrule\bigskip

\subsubsection{本章展望}\label{ux672cux7ae0ux5c55ux671b-3}

本章完成了从 \(\mathbb{N}\) 到 \(\mathbb{C}\) 的严格构造。在 Ch 5
中，我们将深入研究整数 \(\mathbb{Z}\) 的内部结构------数论。在卷二（Ch
6--11）中，数系的代数性质将被系统发展为代数运算与多项式理论。在卷四（Ch
16--24）中，实数的完备性（确界原理）将被用于建立极限和微积分的基础。复数的代数封闭性将在代数基本定理中得到完整阐述。

\bigskip\hrule\bigskip

\subsection{Ch 5 数论基础}\label{ch-5-ux6570ux8bbaux57faux7840}

\textbf{前置知识}：依赖 Ch 4（整数），特别是整数 \(\mathbb{Z}\)
的构造（4.2 节）和整除性的概念。数论研究整数的深层结构。

在 Ch 4 中我们从皮亚诺公理出发构造了完整的数系。现在我们回到整数
\(\mathbb{Z}\)，深入探索其内部的算术结构。数论是数学中最古老的分支之一，它研究整数的整除性、素数分布和同余关系。尽管问题的表述简单，但其证明往往需要深刻的洞察力。

\bigskip\hrule\bigskip

\subsubsection{5.1 整除}\label{ux6574ux9664}

\paragraph{5.1.1 整除的定义}\label{ux6574ux9664ux7684ux5b9aux4e49}

\textbf{定义 5.1.1（整除）}：设 \(a, b\)
是整数，\(b \neq 0\)。如果存在整数 \(k\) 使得 \(a = bk\)，则称 \(b\)
\textbf{整除} \(a\)，记为 \(b \mid a\)。当 \(b \mid a\) 时，\(b\) 称为
\(a\) 的\textbf{因数}，\(a\) 称为 \(b\) 的\textbf{倍数}。

\paragraph{5.1.2
整除的基本性质}\label{ux6574ux9664ux7684ux57faux672cux6027ux8d28}

\textbf{定理 5.1.1（传递性）}：若 \(a \mid b\) 且 \(b \mid c\)，则
\(a \mid c\)。

\textbf{证明}：存在 \(k_1, k_2\) 使 \(b = ak_1\)，\(c = bk_2\)。则
\(c = a(k_1 k_2)\)，即 \(a \mid c\)。\(\blacksquare\)

\textbf{定理 5.1.2（线性组合）}：若 \(a \mid b\) 且
\(a \mid c\)，则对任意整数 \(m, n\)，有 \(a \mid (mb + nc)\)。

\textbf{证明}：设 \(b = ak_1\)，\(c = ak_2\)，则
\(mb + nc = a(mk_1 + nk_2)\)。\(\blacksquare\)

\textbf{定理 5.1.3}：若 \(a \mid b\) 且 \(b \neq 0\)，则
\(|a| \leq |b|\)。

\textbf{证明}：\(b = ak\)，\(k \neq 0\)，故
\(|k| \geq 1\)，\(|b| = |a| \cdot |k| \geq |a|\)。\(\blacksquare\)

\textbf{定理 5.1.4}：若 \(a \mid b\) 且 \(b \mid a\)，则 \(|a| = |b|\)。

\textbf{证明}：由定理 5.1.3，\(|a| \leq |b|\) 且
\(|b| \leq |a|\)。\(\blacksquare\)

\textbf{推论 5.1.1（因数个数的有限性）}：每个非零整数只有有限多个因数。

\textbf{证明}：设 \(n \neq 0\)。若 \(d \mid n\)，则
\(|d| \leq |n|\)。满足 \(|d| \leq |n|\) 的整数 \(d\)
只有有限多个（\(-|n|, -|n|+1, \ldots, |n|-1, |n|\)，共 \(2|n|+1\)
个）。\(\blacksquare\)

\textbf{定理 5.1.5（整除的符号性质）}：设 \(a, b, c\) 是整数。 1.
\(a \mid b \iff -a \mid b \iff a \mid -b \iff -a \mid -b\)。 2. 若
\(a \mid b\) 且 \(a \mid c\)，则 \(a \mid (b + c)\) 且
\(a \mid (b - c)\)。

\textbf{证明}：(1) \(a \mid b\) 意味着 \(b = ak\)，即
\(b = (-a)(-k)\)，故 \(-a \mid b\)。其余类似。(2) 设
\(b = ak_1\)，\(c = ak_2\)，则
\(b + c = a(k_1 + k_2)\)，\(b - c = a(k_1 - k_2)\)，故
\(a \mid (b + c)\) 且 \(a \mid (b - c)\)。\(\blacksquare\)

\textbf{注记}：定理 5.1.5(2) 的逆命题不成立。例如
\(a = 4, b = 2, c = 2\) 时，\(4 \mid (2 + 2)\) 且 \(4 \mid (2 - 2)\)，但
\(4 \nmid 2\)。反方向需要附加条件（如 \(a\) 为奇数时，由 \(a \mid 2b\)
和 \(a \mid 2c\) 结合 \(\gcd(a, 2) = 1\) 可得 \(a \mid b\) 且
\(a \mid c\)）。

\paragraph{5.1.3 带余除法}\label{ux5e26ux4f59ux9664ux6cd5}

\textbf{定理 5.1.6（带余除法）}：设 \(a, b\)
是整数，\(b > 0\)。则存在唯一的整数 \(q\)（商）和 \(r\)（余数），使得
\[a = bq + r, \quad 0 \leq r < b\]

\textbf{证明}：

\textbf{存在性}：考虑
\(S^+ = \{a - bx \mid x \in \mathbb{Z},\, a - bx \geq 0\}\)。取
\(x = -|a|\)，则 \(a - bx = a + b|a| \geq 0\)，故 \(S^+\)
非空。由良序原理（定理 4.1.11），\(S^+\) 有最小元素
\(r = a - bq\)，\(r \geq 0\)。若 \(r \geq b\)，则
\(r - b = a - b(q+1) \in S^+\) 且 \(r - b < r\)，与 \(r\)
的最小性矛盾。故 \(r < b\)。

\textbf{唯一性}：设
\(a = bq_1 + r_1 = bq_2 + r_2\)，\(0 \leq r_1, r_2 < b\)。则
\(b(q_1 - q_2) = r_2 - r_1\)，\(|r_2 - r_1| < b\)。由
\(b \mid (r_2 - r_1)\) 和 \(|r_2 - r_1| < b\)，只能 \(r_1 = r_2\)，从而
\(q_1 = q_2\)。\(\blacksquare\)

\textbf{推论 5.1.2}：任何整数 \(a\) 除以 \(2\) 的余数只有 \(0\) 或
\(1\)，即 \(a\) 要么是偶数（\(a = 2q\)），要么是奇数（\(a = 2q + 1\)）。

\textbf{证明}：在带余除法中取 \(b = 2\) 即得。\(\blacksquare\)

\textbf{推论 5.1.3}：设 \(a, b\) 是整数，\(b \neq 0\)。则 \(b \mid a\)
当且仅当 \(a\) 除以 \(b\) 的余数为 \(0\)。

\textbf{证明}：\(b \mid a\) 意味着 \(a = bk\) 对某个整数 \(k\)，即
\(a = bk + 0\)，余数为 \(0\)。反之，若 \(a = bq + 0\)，则 \(a = bq\)，即
\(b \mid a\)。\(\blacksquare\)

\bigskip\hrule\bigskip

\subsubsection{5.2
最大公约数与最小公倍数}\label{ux6700ux5927ux516cux7ea6ux6570ux4e0eux6700ux5c0fux516cux500dux6570}

\paragraph{5.2.1 最大公约数}\label{ux6700ux5927ux516cux7ea6ux6570}

\textbf{定义 5.2.1（公约数与最大公约数）}：设 \(a, b\)
是不全为零的整数。若 \(d \mid a\) 且 \(d \mid b\)，则 \(d\) 为 \(a, b\)
的\textbf{公约数}。所有公约数中最大的正整数称为\textbf{最大公约数}，记为
\(\gcd(a, b)\)。

\textbf{注记 5.2.1}：\(\gcd(0, 0)\) 没有定义（因为每个正整数都是 \(0\)
的因数，没有最大值）。但通常定义
\(\gcd(0, n) = |n|\)（\(n \neq 0\)）。特别地，\(\gcd(0, n) = |n|\)，\(\gcd(n, n) = |n|\)。

\textbf{定理 5.2.1（GCD的基本性质）}：设 \(a, b\)
是不全为零的整数，\(d = \gcd(a, b)\)。则： 1. \(d \mid a\) 且
\(d \mid b\)（\(d\) 是公约数） 2. 若 \(c \mid a\) 且 \(c \mid b\)，则
\(c \mid d\)（\(d\) 是最大的公约数） 3. \(\gcd(a, b) = \gcd(|a|, |b|)\)
4. \(\gcd(a, b) = \gcd(b, a)\) 5. \(\gcd(a, 0) = |a|\)（\(a \neq 0\)）

\textbf{证明}：(1) 由定义。(2) 由裴蜀定理（定理
5.2.3），\(d = ax + by\)，故 \(c \mid d\)。(3) 由整除的性质
\(b \mid a \iff |b| \mid |a|\)。(4) 由交换律。(5) \(0\)
的因数是所有整数，\(a\) 的因数不超过 \(|a|\)，故
\(\gcd(a, 0) = |a|\)。\(\blacksquare\)

\textbf{定理 5.2.2（辗转相除法/欧几里得算法）}：设 \(a, b\)
是正整数，\(a \geq b\)。反复做带余除法：

\[a = bq_1 + r_1,\; b = r_1 q_2 + r_2,\; \ldots,\; r_{k-2} = r_{k-1}q_k + r_k,\; r_{k-1} = r_k q_{k+1} + 0\]

则 \(\gcd(a, b) = r_k\)（最后一个非零余数）。

\textbf{证明}：

\textbf{关键引理}：\(\gcd(a, b) = \gcd(b, r)\)，其中 \(a = bq + r\)。

\textbf{引理的证明}：设 \(d = \gcd(a, b)\)。由 \(d \mid a\) 且
\(d \mid b\) 得 \(d \mid (a - bq) = r\)，故 \(d\) 是 \(b, r\)
的公约数，\(d \leq \gcd(b, r)\)。反之，设 \(d' = \gcd(b, r)\)，由
\(d' \mid b\) 且 \(d' \mid r\) 得 \(d' \mid (bq + r) = a\)，故
\(d' \leq \gcd(a, b) = d\)。综合得 \(d = d'\)。\(\blacksquare\)

对每一步应用引理：\(\gcd(a, b) = \gcd(b, r_1) = \gcd(r_1, r_2) = \cdots = \gcd(r_{k-1}, r_k) = \gcd(r_k, 0) = r_k\)。\(\blacksquare\)

\textbf{定理 5.2.3（裴蜀定理）}：设 \(a, b\)
是不全为零的整数。则存在整数 \(x, y\) 使得 \[ax + by = \gcd(a, b)\]

更一般地，\(ax + by = c\) 有整数解当且仅当 \(\gcd(a, b) \mid c\)。

\textbf{证明}：

\textbf{第一步}：考虑
\(S = \{ax + by \mid x, y \in \mathbb{Z},\, ax + by > 0\}\)。\(S\)
非空，由良序原理有最小元素 \(d = ax_0 + by_0\)。

\textbf{第二步}：对 \(a\) 做带余除法
\(a = dq + r\)，\(0 \leq r < d\)。则
\(r = a - dq = a(1 - qx_0) + b(-qy_0)\)。若 \(r > 0\)，则 \(r \in S\) 且
\(r < d\)，与 \(d\) 的最小性矛盾。故 \(r = 0\)，即 \(d \mid a\)。同理
\(d \mid b\)。

\textbf{第三步}：\(d\) 是 \(a, b\) 的公约数。设 \(c\) 是任意公约数，则
\(c \mid (ax_0 + by_0) = d\)，故 \(c \leq d\)。因此 \(d = \gcd(a, b)\)。

\textbf{第四步}：若 \(\gcd(a, b) \mid c\)，设
\(c = k \cdot \gcd(a, b)\)，则 \(x = kx_0\)，\(y = ky_0\) 是解。反之，若
\(ax + by = c\) 有解，由 \(\gcd(a, b) \mid a\) 和 \(\gcd(a, b) \mid b\)
得 \(\gcd(a, b) \mid c\)。\(\blacksquare\)

\textbf{推论 5.2.4（裴蜀定理的等价形式）}：\(\gcd(a, b)\) 是 \(a\) 和
\(b\) 的线性组合中最小的正整数。

\textbf{证明}：设 \(d = \gcd(a, b)\)。由裴蜀定理，存在整数 \(x_0, y_0\)
使得 \(ax_0 + by_0 = d\)，故 \(d\) 是 \(a\) 和 \(b\) 的一个线性组合。

设 \(s\) 是 \(a\) 和 \(b\) 的任意正线性组合，即
\(s = am + bn\)（\(m, n \in \mathbb{Z}\)，\(s > 0\)）。由 \(d \mid a\)
和 \(d \mid b\)，得 \(d \mid s\)，故 \(s \geq d\)。

因此 \(d\) 是最小的正线性组合。\(\blacksquare\)

\textbf{推论 5.2.5}：\(\gcd(a, b) = 1\) 当且仅当存在整数 \(x, y\) 使得
\(ax + by = 1\)。

\textbf{证明}：\((\Rightarrow)\)：由裴蜀定理。( \(\Leftarrow)\)：若
\(ax + by = 1\)，则 \(\gcd(a, b) \mid 1\)，故
\(\gcd(a, b) = 1\)。\(\blacksquare\)

\textbf{定理 5.2.6（GCD的运算性质）}： 1.
\(\gcd(a, b) = \gcd(a, b + ka)\) 对任意整数 \(k\)。 2.
\(\gcd(a, b) = \gcd(a, b \bmod a)\)（\(a > 0\)）。 3.
\(\gcd(a, b, c) = \gcd(\gcd(a, b), c)\)。

\textbf{证明}：(1) 设 \(d = \gcd(a, b)\)。\(d \mid a\) 且
\(d \mid b\)，故 \(d \mid (b + ka)\)。因此 \(d\) 是 \(a, b + ka\)
的公约数，\(d \leq \gcd(a, b + ka)\)。反之设
\(d' = \gcd(a, b + ka)\)。\(d' \mid a\) 且 \(d' \mid (b + ka)\)，故
\(d' \mid [(b + ka) - ka] = b\)。因此 \(d'\) 是 \(a, b\)
的公约数，\(d' \leq d\)。故 \(d = d'\)。

\begin{enumerate}
\def\labelenumi{(\arabic{enumi})}
\setcounter{enumi}{1}
\item
  由 (1) 取 \(k = -\lfloor b/a \rfloor\)。
\item
  \(\gcd(a, b, c)\) 是同时整除 \(a, b, c\)
  的最大正整数。\(\gcd(\gcd(a, b), c)\) 是同时整除 \(\gcd(a, b)\) 和
  \(c\) 的最大正整数。由于
  \(d \mid \gcd(a, b) \iff (d \mid a \wedge d \mid b)\)，两者等价。\(\blacksquare\)
\end{enumerate}

\paragraph{5.2.2 最小公倍数}\label{ux6700ux5c0fux516cux500dux6570}

\textbf{定义 5.2.2（最小公倍数）}：设 \(a, b\) 是非零整数。\(a, b\)
的所有正公倍数中最小的称为\textbf{最小公倍数}，记为
\(\operatorname{lcm}(a, b)\)。

\textbf{注记 5.2.2}：\(\operatorname{lcm}(a, b)\)
的存在性由以下事实保证：\(a, b\) 的公倍数集合
\(M = \{m > 0 : a \mid m \text{ 且 } b \mid m\}\)
是非空的正整数集合（\(|ab| \in M\)），由良序原理有最小元素。

\textbf{注记 5.2.3}：GCD 和 LCM
可以推广到多个整数：\(\gcd(a_1, \ldots, a_n) = \gcd(\gcd(a_1, \ldots, a_{n-1}), a_n)\)，\(\operatorname{lcm}(a_1, \ldots, a_n) = \operatorname{lcm}(\operatorname{lcm}(a_1, \ldots, a_{n-1}), a_n)\)。

\textbf{定理
5.2.7}：\(\operatorname{lcm}(a, b) = \frac{|ab|}{\gcd(a, b)}\)。

\textbf{证明}：设
\(d = \gcd(a, b)\)，\(a = da'\)，\(b = db'\)。由裴蜀定理，\(a'x + b'y = 1\)，故
\(\gcd(a', b') = 1\)。

\(m = da'b' = ab/d\) 是公倍数（\(m = a \cdot b' = b \cdot a'\)）。设
\(n\) 是任意正公倍数，\(a \mid n\) 且 \(b \mid n\)。设
\(n = as = bt\)，则 \(a's = b't\)。由 \(\gcd(a', b') = 1\) 和
\(a' \mid b't\) 得 \(a' \mid t\)。设 \(t = a'k\)，则
\(n = db'ka' = mk\)，故 \(m \mid n\)，\(n \geq m\)。\(\blacksquare\)

\textbf{推论 5.2.8（LCM的性质）}： 1.
\(\operatorname{lcm}(a, b) = \operatorname{lcm}(|a|, |b|)\) 2. 若
\(a \mid c\) 且 \(b \mid c\)，则 \(\operatorname{lcm}(a, b) \mid c\)。
3. \(\operatorname{lcm}(a, b) \cdot \gcd(a, b) = |ab|\)。

\textbf{证明}：(1) 由 \(|a| \mid c \iff a \mid c\)。(2) 由定理 5.2.7
的证明。(3)
\(\operatorname{lcm}(a, b) = |ab|/\gcd(a, b)\)。\(\blacksquare\)

\textbf{定理 5.2.9（互素的性质）}：设 \(\gcd(a, b) = 1\)。 1. 若
\(a \mid c\) 且 \(b \mid c\)，则 \(ab \mid c\)。 2. 若 \(a \mid bc\)，则
\(a \mid c\)。 3. \(\gcd(a^n, b^n) = 1\) 对任意正整数 \(n\)。

\textbf{证明}：(1) 由 \(\operatorname{lcm}(a, b) = ab\)（当
\(\gcd(a, b) = 1\)）和 \(\operatorname{lcm}(a, b) \mid c\)。(2)
由裴蜀定理，\(ax + by = 1\)，乘以 \(c\) 得
\(acx + bcy = c\)。\(a \mid acx\) 且 \(a \mid bcy\)（因
\(a \mid bc\)），故 \(a \mid c\)。(3) 若素数
\(p \mid \gcd(a^n, b^n)\)，则 \(p \mid a^n\) 且
\(p \mid b^n\)。由欧几里得引理，\(p \mid a\) 且 \(p \mid b\)，故
\(p \mid \gcd(a, b) = 1\)，矛盾。\(\blacksquare\)

\bigskip\hrule\bigskip

\subsubsection{5.3 素数}\label{ux7d20ux6570}

\paragraph{5.3.1 素数的定义}\label{ux7d20ux6570ux7684ux5b9aux4e49}

\textbf{定义 5.3.1（素数与合数）}：大于 \(1\) 的正整数 \(p\)，如果只有
\(1\) 和 \(p\) 两个正因数，则称 \(p\) 为\textbf{素数}。大于 \(1\)
但不是素数的正整数称为\textbf{合数}。

\textbf{注记}：\(1\) 既不是素数也不是合数。\(2\) 是唯一的偶素数。

\textbf{引理 5.3.1}：每个大于 \(1\) 的正整数都有一个素因数。

\textbf{证明}：设 \(n > 1\)。若 \(n\) 是素数，结论成立。若 \(n\)
是合数，\(n = ab\)，\(1 < a, b < n\)。对 \(a\)
重复此过程。由于正整数不能无限递减，必在有限步内得到素因数。\(\blacksquare\)

\textbf{定理 5.3.1（欧几里得引理）}：若 \(p\) 是素数且 \(p \mid ab\)，则
\(p \mid a\) 或 \(p \mid b\)。

\textbf{证明}：假设 \(p \nmid a\)。因 \(p\)
是素数，\(\gcd(p, a) = 1\)。由裴蜀定理，存在 \(x, y\) 使得
\(px + ay = 1\)。两边乘 \(b\)：\(pbx + aby = b\)。由 \(p \mid pbx\) 和
\(p \mid aby\)（因 \(p \mid ab\)），得 \(p \mid b\)。\(\blacksquare\)

\textbf{定理 5.3.1'（素数的基本性质）}： 1. 若 \(p\) 是素数且
\(p \mid a_1 a_2 \cdots a_k\)，则 \(p\) 至少整除某个 \(a_i\)。 2. 设
\(p\) 是素数。则对任意整数 \(a\)，恰有 \(p \mid a\) 或
\(\gcd(p, a) = 1\) 之一成立。 3. 若 \(p, q\) 都是素数且 \(p \mid q\)，则
\(p = q\)。

\textbf{证明}：(1) 对 \(k\) 用归纳法。\(k = 1\) 时显然。\(k = 2\)
时即欧几里得引理（定理 5.3.1）。设对 \(k-1\) 成立。若
\(p \mid a_1 \cdots a_{k-1}\)，由归纳假设得证。若
\(p \nmid a_1 \cdots a_{k-1}\)，由欧几里得引理，\(p \mid a_k\)。(2)
\(p\) 的正因数只有 \(1\) 和 \(p\)。若 \(p \nmid a\)，则
\(\gcd(p, a) \neq p\)，故 \(\gcd(p, a) = 1\)。(3) \(p \mid q\) 且 \(q\)
是素数，故 \(p = 1\) 或 \(p = q\)。但 \(p\) 是素数，\(p \neq 1\)，故
\(p = q\)。\(\blacksquare\)

\paragraph{5.3.2
算术基本定理}\label{ux7b97ux672fux57faux672cux5b9aux7406}

\textbf{定理 5.3.2（算术基本定理）}：每个大于 \(1\) 的正整数 \(n\)
都可以唯一地（不计因子次序）表示为素数的乘积：
\[n = p_1 p_2 \cdots p_k, \quad p_1 \leq p_2 \leq \cdots \leq p_k\]

\textbf{证明}：

\textbf{存在性}：对 \(n\) 用强归纳法。\(n = 2\) 时，\(2\)
本身是素数，分解存在。设对所有 \(2 \leq m < n\)，素数分解存在。考虑
\(n\)：若 \(n\) 是素数，分解为自身。若 \(n\) 是合数，则
\(n = ab\)，\(1 < a, b < n\)。由归纳假设，\(a\) 和 \(b\)
都可以表示为素数乘积：\(a = p_1 \cdots p_s\)，\(b = q_1 \cdots q_t\)。于是
\(n = p_1 \cdots p_s q_1 \cdots q_t\)
也是素数乘积。由强归纳法，存在性得证。

\textbf{唯一性}：设 \(n = p_1 p_2 \cdots p_k = q_1 q_2 \cdots q_l\)
是两种素数分解（\(p_1 \leq \cdots \leq p_k\)，\(q_1 \leq \cdots \leq q_l\)）。

由 \(p_1 \mid n = q_1 q_2 \cdots q_l\)，根据定理 5.3.1'(1)，\(p_1\)
整除某个 \(q_j\)。适当重排使 \(p_1 \mid q_1\)。由于 \(q_1\) 是素数且
\(p_1 > 1\)，由定理 5.3.1'(3) 得 \(p_1 = q_1\)。

约去 \(p_1 = q_1\)，得 \(p_2 \cdots p_k = q_2 \cdots q_l\)。若
\(k = 1\)，则左边为 \(1\)，右边 \(q_2 \cdots q_l = 1\)，只能
\(l = 1\)。若 \(k > 1\)，对 \(k + l\) 用归纳法（或重复上述过程），得
\(k = l\) 且 \(p_i = q_i\)（\(i = 1, 2, \ldots, k\)）。\(\blacksquare\)

\textbf{推论 5.3.1（素因子分解的标准形式）}：每个大于 \(1\) 的正整数
\(n\) 都可以唯一地写成 \[n = p_1^{a_1} p_2^{a_2} \cdots p_r^{a_r}\] 其中
\(p_1 < p_2 < \cdots < p_r\) 是素数，\(a_1, a_2, \ldots, a_r\)
是正整数。

\textbf{证明}：由算术基本定理，将相同的素因子合并即得。\(\blacksquare\)

\textbf{推论 5.3.2（GCD和LCM的素因子表示）}：设
\(a = p_1^{a_1} \cdots p_r^{a_r}\)，\(b = p_1^{b_1} \cdots p_r^{b_r}\)（允许某些指数为
\(0\)）。则
\[\gcd(a, b) = p_1^{\min(a_1, b_1)} \cdots p_r^{\min(a_r, b_r)}\]
\[\operatorname{lcm}(a, b) = p_1^{\max(a_1, b_1)} \cdots p_r^{\max(a_r, b_r)}\]

\textbf{证明}：设 \(d = p_1^{d_1} \cdots p_r^{d_r}\)。\(d \mid a\)
当且仅当 \(d_i \leq a_i\) 对所有 \(i\)。\(d \mid b\) 当且仅当
\(d_i \leq b_i\) 对所有 \(i\)。\(d\) 是公约数当且仅当
\(d_i \leq \min(a_i, b_i)\) 对所有
\(i\)。最大公约数取等号。\(\operatorname{lcm}\)
的证明类似。\(\blacksquare\)

\textbf{推论
5.3.3}：\(\gcd(a, b) \cdot \operatorname{lcm}(a, b) = ab\)。

\textbf{证明}：\(\min(a_i, b_i) + \max(a_i, b_i) = a_i + b_i\)。\(\blacksquare\)

\paragraph{5.3.3
素数的无穷性}\label{ux7d20ux6570ux7684ux65e0ux7a77ux6027}

\textbf{定理 5.3.3（欧几里得定理）}：素数有无穷多个。

\textbf{证明}：假设素数只有有限个 \(p_1, p_2, \ldots, p_k\)。考虑
\(N = p_1 p_2 \cdots p_k + 1\)。\(N > 1\)，有素因数 \(p\)。但 \(p_i\) 除
\(N\) 的余数为 \(1\)（\(p_i \mid p_1 \cdots p_k\)），故
\(p_i \nmid N\)，即 \(p \neq p_i\) 对所有 \(i\)。矛盾。\(\blacksquare\)

\textbf{注记 5.3.1}：欧几里得的原始证明是构造性的，但
\(N = p_1 \cdots p_k + 1\) 本身不一定是素数。例如
\(2 \cdot 3 \cdot 5 \cdot 7 \cdot 11 \cdot 13 + 1 = 30031 = 59 \times 509\)。证明的精髓在于：\(N\)
有不在列表中的素因数。

\textbf{推论 5.3.4（第 \(k\) 个素数的上界）}：设 \(p_k\) 是第 \(k\)
个素数。则 \(p_k \leq 2^{2^{k-1}}\)。

\textbf{证明}：对 \(k\) 用归纳法。\(p_1 = 2 = 2^{2^0}\)。设
\(p_i \leq 2^{2^{i-1}}\)（\(i = 1, \ldots, k\)）。由欧几里得的证明，\(p_{k+1} \leq p_1 p_2 \cdots p_k + 1\)。\(p_1 \cdots p_k \leq 2^{2^0 + 2^1 + \cdots + 2^{k-1}} = 2^{2^k - 1}\)。故
\(p_{k+1} \leq 2^{2^k - 1} + 1 \leq 2^{2^k}\)。\(\blacksquare\)

\textbf{定理 5.3.4（关于素数分布的简单结果）}：对任意正整数 \(n\)，在
\(n\) 和 \(2n\) 之间（含端点）至少存在一个素数。

\textbf{证明}：考虑二项式系数 \(\binom{2n}{n} = \frac{(2n)!}{(n!)^2}\)。

\textbf{第一步}：\(\binom{2n}{n} \geq 2^n\)（\(n \geq 1\)）。由二项式定理，\(4^n = (1+1)^{2n} = \sum_{k=0}^{2n} \binom{2n}{k} \leq (2n+1)\binom{2n}{n}\)（\(\binom{2n}{n}\)
是最大项），故 \(\binom{2n}{n} \geq \frac{4^n}{2n+1}\)。

\textbf{第二步}：\(\binom{2n}{n}\) 的素因子都不超过 \(2n\)。因为
\(\binom{2n}{n} = \frac{(2n)!}{(n!)^2}\) 的分子的素因子都不超过 \(2n\)。

\textbf{第三步}：若 \(n < p \leq 2n\)（\(p\) 是素数），则
\(p \mid \binom{2n}{n}\)。因为 \(p\) 整除分子 \((2n)!\)
中的因子（\(p \leq 2n\)），但不整除分母 \((n!)^2\)（\(p > n\)）。

\textbf{第四步}：假设 \((n, 2n]\) 中没有素数。则 \(\binom{2n}{n}\)
的所有素因子都
\(\leq n\)。由算术基本定理，\(\binom{2n}{n} = \prod_{p \leq n} p^{a_p}\)。其中
\(a_p = \sum_{j=1}^{\infty} \left(\lfloor \frac{2n}{p^j} \rfloor - 2\lfloor \frac{n}{p^j} \rfloor\right)\)。每一项
\(\leq 1\)（因为
\(\lfloor 2x \rfloor - 2\lfloor x \rfloor \in \{0, 1\}\)），且只有
\(j \leq \log_p(2n)\) 项非零。故
\(a_p \leq \lfloor \log_p(2n) \rfloor\)。因此
\(\binom{2n}{n} \leq \prod_{p \leq n} p^{\lfloor \log_p(2n) \rfloor} \leq (2n)^{\pi(n)}\)（\(\pi(n)\)
是不超过 \(n\) 的素数个数）。

\textbf{第五步}：由第一步，\(2^n \leq \binom{2n}{n} \leq (2n)^{\pi(n)}\)。取对数得
\(\pi(n) \geq \frac{n \ln 2}{\ln(2n)}\)。这足以说明 \(\pi(n)\) 随 \(n\)
增长趋于无穷（即素数有无穷多个，已由定理 5.3.3
独立证明），因此对充分大的 \(n\) 产生矛盾。对于小的 \(n\)
值可以直接验证：\((1, 2]\) 有素数 \(2\)；\((2, 4]\) 有素数
\(3\)；\((3, 6]\) 有素数 \(5\)；\((4, 8]\) 有素数
\(5, 7\)。\(\blacksquare\)

\paragraph{5.3.4
埃拉托色尼筛法}\label{ux57c3ux62c9ux6258ux8272ux5c3cux7b5bux6cd5}

\textbf{算法 5.3.1（埃拉托色尼筛法）}：求不超过 \(N\) 的所有素数的方法：
1. 列出 \(2, 3, 4, \ldots, N\)。 2. 取列表中最小的数 \(p\)（初始
\(p = 2\)），标记 \(p\) 为素数。 3. 划去 \(p\)
的所有倍数（\(2p, 3p, 4p, \ldots\)）。 4. 令 \(p\)
为列表中下一个未被划去的数，回到步骤 3。 5. 当 \(p > \sqrt{N}\)
时停止。列表中剩余的未被划去的数都是素数。

\textbf{定理 5.3.5（筛法的正确性）}：算法结束后，不超过 \(N\)
的未被划去的数恰好是不超过 \(N\) 的素数。

\textbf{证明}：若 \(n \leq N\) 是合数，则
\(n = ab\)（\(1 < a \leq b < n\)），\(a \leq \sqrt{n} \leq \sqrt{N}\)。\(a\)
有一个素因子 \(p \leq a \leq \sqrt{N}\)。当处理到 \(p\) 时，\(n\) 作为
\(p\) 的倍数被划去。因此合数都会被划去。素数不会被划去：素数 \(p\)
只在处理到 \(p\) 自身时被标记为素数，而非被划去（\(p\) 的倍数
\(2p, 3p, \ldots\) 都大于 \(p\)）。\(\blacksquare\)

\textbf{注记 5.3.2（筛法的效率）}：埃拉托色尼筛法的时间复杂度为
\(O(N \log \log N)\)，空间复杂度为 \(O(N)\)。对于 \(N\)
很大的情况，可以使用更高级的筛法（如欧拉筛法达到 \(O(N)\) 时间复杂度）。

\paragraph{5.3.5
素数分布的初步结果}\label{ux7d20ux6570ux5206ux5e03ux7684ux521dux6b65ux7ed3ux679c}

\textbf{定义 5.3.2（素数计数函数）}：\(\pi(x)\) 表示不超过 \(x\)
的素数个数。

\textbf{定理
5.3.6（素数定理的弱形式）}：\(\pi(x) \geq \frac{x}{2\ln x}\) 对充分大的
\(x\) 成立。

\textbf{注记 5.3.3（素数定理）}：素数定理断言
\(\pi(x) \sim \frac{x}{\ln x}\)（当 \(x \to \infty\) 时），即
\(\lim_{x \to \infty} \frac{\pi(x)}{x/\ln x} = 1\)。这个深刻结果由
Hadamard 和 de la Vallee Poussin 在 1896
年独立证明，证明使用了复分析工具。一个初等证明由 Erdos 和 Selberg 在
1949 年给出。

\bigskip\hrule\bigskip

\subsubsection{5.4 同余}\label{ux540cux4f59}

\paragraph{5.4.1 同余的定义}\label{ux540cux4f59ux7684ux5b9aux4e49}

\textbf{定义 5.4.1（同余）}：设 \(m\) 是正整数。若
\(m \mid (a - b)\)，则称 \(a\) 与 \(b\) \textbf{关于模 \(m\) 同余}，记为
\(a \equiv b \pmod{m}\)。

\textbf{定理 5.4.1}：\(a \equiv b \pmod{m}\) 当且仅当 \(a\) 和 \(b\)
除以 \(m\) 的余数相同。

\textbf{证明}：设
\(a = mq_1 + r_1\)，\(b = mq_2 + r_2\)，\(0 \leq r_1, r_2 < m\)。则
\(m \mid (a - b) = m(q_1 - q_2) + (r_1 - r_2)\) 当且仅当
\(m \mid (r_1 - r_2)\)。由 \(|r_1 - r_2| < m\)，这等价于
\(r_1 = r_2\)。\(\blacksquare\)

\textbf{定理 5.4.2}：同余关系是等价关系（自反、对称、传递）。

\textbf{证明}：自反性：\(m \mid (a - a) = 0\)。对称性：若
\(m \mid (a - b)\) 则 \(m \mid (b - a)\)。传递性：若 \(m \mid (a - b)\)
且 \(m \mid (b - c)\)，则
\(m \mid [(a - b) + (b - c)] = (a - c)\)。\(\blacksquare\)

\textbf{推论 5.4.1}：整数集 \(\mathbb{Z}\) 关于模 \(m\) 的同余关系被分为
\(m\) 个等价类（剩余类）：\([0], [1], \ldots, [m-1]\)，其中
\([r] = \{r + km \mid k \in \mathbb{Z}\}\)。

\textbf{证明}：每个整数 \(a\) 除以 \(m\) 的余数 \(r\)
唯一确定（\(0 \leq r < m\)），且
\(a \equiv r \pmod{m}\)。因此每个整数恰好属于一个剩余类。\(\blacksquare\)

\textbf{定理 5.4.2'（同余的除法规则）}：若 \(ac \equiv bc \pmod{m}\) 且
\(\gcd(c, m) = 1\)，则 \(a \equiv b \pmod{m}\)。

\textbf{证明}：由 \(m \mid c(a - b)\) 和
\(\gcd(c, m) = 1\)。由裴蜀定理，存在 \(x, y\) 使得
\(cx + my = 1\)。两边乘
\((a - b)\)：\(c(a - b)x + m(a - b)y = a - b\)。由 \(m \mid c(a - b)\)
得 \(m \mid [c(a - b)x + m(a - b)y] = a - b\)，即
\(a \equiv b \pmod{m}\)。\(\blacksquare\)

\paragraph{5.4.2 同余的运算}\label{ux540cux4f59ux7684ux8fd0ux7b97}

\textbf{定理 5.4.3（同余的运算性质）}：若
\(a \equiv b \pmod{m}\)，\(c \equiv d \pmod{m}\)，则： 1.
\(a + c \equiv b + d \pmod{m}\) 2. \(a - c \equiv b - d \pmod{m}\) 3.
\(ac \equiv bd \pmod{m}\) 4. \(a^n \equiv b^n \pmod{m}\)（\(n\)
为正整数）

\textbf{证明}：

\begin{enumerate}
\def\labelenumi{(\arabic{enumi})}
\item
  \(m \mid [(a - b) + (c - d)] = [(a + c) - (b + d)]\)。
\item
  \(ac - bd = c(a - b) + b(c - d)\)。由 \(m \mid (a - b)\) 和
  \(m \mid (c - d)\) 得 \(m \mid (ac - bd)\)。
\item
  由 (3) 反复应用或用数学归纳法。\(\blacksquare\)
\end{enumerate}

\paragraph{5.4.3 费马小定理}\label{ux8d39ux9a6cux5c0fux5b9aux7406}

\textbf{定理 5.4.4（费马小定理）}：设 \(p\) 是素数，\(a\) 是整数且
\(p \nmid a\)。则 \[a^{p-1} \equiv 1 \pmod{p}\]

\textbf{证明}：

\textbf{引理}：设 \(p\) 是素数，\(1 \leq k \leq p - 1\)。则
\(p \mid \binom{p}{k}\)。

\textbf{引理的证明}：\(\binom{p}{k} = \frac{p!}{k!(p-k)!}\)，故
\(p \cdot k! \cdot (p-k)! = p!\)。\(p\) 不整除 \(k!\)（所有因子小于
\(p\)），也不整除 \((p-k)!\)，但 \(p \mid p!\)，故
\(p \mid \binom{p}{k}\)。\(\blacksquare\)

回到主定理的证明。由二项式定理：
\[(a + 1)^p = \sum_{k=0}^{p} \binom{p}{k} a^k = a^p + \binom{p}{1}a^{p-1} + \cdots + \binom{p}{p-1}a + 1\]

由引理，\(p \mid \binom{p}{k}\)（\(1 \leq k \leq p - 1\)），故
\((a+1)^p \equiv a^p + 1 \pmod{p}\)。

用归纳法证 \(a^p \equiv a \pmod{p}\)：\(a = 1\) 时显然。设
\(a^p \equiv a\)，则 \((a+1)^p \equiv a^p + 1 \equiv a + 1 \pmod{p}\)。

对所有整数 \(a\) 也成立（对 \(p = 2\) 可直接验证，对奇素数 \(p\) 由
\((-a)^p = -a^p\) 可得）。

由 \(a^p \equiv a\)，若 \(p \nmid a\)，由欧几里得引理从
\(p \mid a(a^{p-1} - 1)\) 和 \(p \nmid a\) 得
\(p \mid (a^{p-1} - 1)\)，即
\(a^{p-1} \equiv 1 \pmod{p}\)。\(\blacksquare\)

\paragraph{5.4.4
中国剩余定理}\label{ux4e2dux56fdux5269ux4f59ux5b9aux7406}

\textbf{定理 5.4.5（中国剩余定理）}：设 \(m_1, m_2\)
是互素的正整数（\(\gcd(m_1, m_2) = 1\)）。则对任意整数
\(a_1, a_2\)，同余方程组
\[\begin{cases} x \equiv a_1 \pmod{m_1} \\ x \equiv a_2 \pmod{m_2} \end{cases}\]
在模 \(m_1 m_2\) 的意义下有唯一解。

\textbf{证明}：

\textbf{存在性}：由 \(\gcd(m_1, m_2) = 1\) 和裴蜀定理，存在 \(s, t\)
使得 \(m_1 s + m_2 t = 1\)。令 \(x = a_1 m_2 t + a_2 m_1 s\)。

验证：\(x = a_1 m_2 t + a_2 m_1 s \equiv a_1 m_2 t \pmod{m_1}\)。又
\(m_2 t = 1 - m_1 s \equiv 1 \pmod{m_1}\)，故
\(x \equiv a_1 \pmod{m_1}\)。类似地 \(x \equiv a_2 \pmod{m_2}\)。

\textbf{唯一性}：设 \(x_1, x_2\) 都是解。则 \(m_1 \mid (x_1 - x_2)\) 且
\(m_2 \mid (x_1 - x_2)\)。由 \(\gcd(m_1, m_2) = 1\) 得
\(m_1 m_2 \mid (x_1 - x_2)\)，即
\(x_1 \equiv x_2 \pmod{m_1 m_2}\)。\(\blacksquare\)

\textbf{定理 5.4.6（中国剩余定理的一般形式）}：设
\(m_1, m_2, \ldots, m_k\) 两两互素。则同余方程组
\[x \equiv a_i \pmod{m_i}, \quad i = 1, 2, \ldots, k\] 在模
\(M = m_1 m_2 \cdots m_k\) 的意义下有唯一解。解为
\[x \equiv \sum_{i=1}^{k} a_i M_i N_i \pmod{M}\] 其中
\(M_i = \frac{M}{m_i}\)，\(N_i\) 是 \(M_i\) 模 \(m_i\)
的逆元（\(M_i N_i \equiv 1 \pmod{m_i}\)）。

\textbf{证明}：对 \(k\) 用归纳法。\(k = 2\) 时即定理 5.4.5。设对 \(k-1\)
成立。前 \(k-1\) 个方程在模 \(M' = m_1 \cdots m_{k-1}\) 下有唯一解
\(x_0\)。问题化为求 \(x \equiv x_0 \pmod{M'}\) 且
\(x \equiv a_k \pmod{m_k}\)。由 \(\gcd(M', m_k) = 1\)（因
\(m_1, \ldots, m_k\) 两两互素），应用定理 5.4.5 得证。\(\blacksquare\)

\paragraph{5.4.5
欧拉函数与欧拉定理}\label{ux6b27ux62c9ux51fdux6570ux4e0eux6b27ux62c9ux5b9aux7406}

\textbf{定义 5.4.2（欧拉函数）}：对正整数 \(n\)，定义 \(\varphi(n)\) 为
\(1, 2, \ldots, n\) 中与 \(n\) 互素的数的个数：
\[\varphi(n) = |\{k \in \{1, 2, \ldots, n\} \mid \gcd(k, n) = 1\}|\]

\textbf{注记 5.4.1}：\(\varphi(1) = 1\)（\(\gcd(1, 1) = 1\)）。对素数
\(p\)，\(\varphi(p) = p - 1\)（\(1, 2, \ldots, p-1\) 都与 \(p\) 互素）。

\textbf{定理 5.4.7（欧拉函数的积性）}：若 \(\gcd(a, b) = 1\)，则
\(\varphi(ab) = \varphi(a)\varphi(b)\)。

\textbf{证明}：考虑映射
\(\phi: \{1, \ldots, ab\} \to \{1, \ldots, a\} \times \{1, \ldots, b\}\)，\(\phi(x) = (x \bmod a, x \bmod b)\)。

\textbf{第一步}：\(\phi\) 是双射。由中国剩余定理（定理 5.4.5），对每个
\((r_1, r_2)\)（\(1 \leq r_1 \leq a\),
\(1 \leq r_2 \leq b\)），同余方程组 \(x \equiv r_1 \pmod{a}\),
\(x \equiv r_2 \pmod{b}\) 在 \(\{1, \ldots, ab\}\) 中有唯一解。因此
\(\phi\) 是一一对应的。

\textbf{第二步}：\(\gcd(x, ab) = 1 \iff \gcd(x, a) = 1\) 且
\(\gcd(x, b) = 1\)。\((\Rightarrow)\)：若 \(d = \gcd(x, a) > 1\)，则
\(d \mid x\) 且 \(d \mid a\)，故
\(d \mid ab\)，\(d \mid \gcd(x, ab) = 1\)，矛盾。\((\Leftarrow)\)：若
\(d = \gcd(x, ab) > 1\)，设 \(p\) 是 \(d\) 的素因子。\(p \mid ab\)，故
\(p \mid a\) 或 \(p \mid b\)。若 \(p \mid a\)，则
\(p \mid \gcd(x, a) = 1\)，矛盾。

\textbf{第三步}：\(\varphi(ab)\) 等于满足 \(\gcd(x, a) = 1\) 且
\(\gcd(x, b) = 1\) 的 \(x\) 的个数。由 \(\phi\) 是双射，这等于满足
\(\gcd(r_1, a) = 1\) 的 \(r_1\) 的个数乘以满足 \(\gcd(r_2, b) = 1\) 的
\(r_2\) 的个数，即 \(\varphi(a) \cdot \varphi(b)\)。\(\blacksquare\)

\textbf{定理 5.4.8（欧拉函数的公式）}：设
\(n = p_1^{a_1} \cdots p_r^{a_r}\)（素数分解），则
\[\varphi(n) = n \prod_{i=1}^{r} \left(1 - \frac{1}{p_i}\right)\]

\textbf{证明}：由积性（定理
5.4.7），\(\varphi(n) = \varphi(p_1^{a_1}) \cdots \varphi(p_r^{a_r})\)。因此只需计算
\(\varphi(p^a)\)。

\(1, 2, \ldots, p^a\) 中与 \(p\) 不互素的数恰好是 \(p\)
的倍数：\(p, 2p, 3p, \ldots, p^a\)，共 \(p^{a-1}\) 个。因此
\(\varphi(p^a) = p^a - p^{a-1} = p^a(1 - 1/p)\)。\(\blacksquare\)

\textbf{定理 5.4.9（欧拉定理）}：设 \(\gcd(a, n) = 1\)，则
\[a^{\varphi(n)} \equiv 1 \pmod{n}\]

\textbf{证明}：设 \(r_1, r_2, \ldots, r_{\varphi(n)}\) 是
\(1, 2, \ldots, n\) 中与 \(n\) 互素的全部数。考虑
\(ar_1, ar_2, \ldots, ar_{\varphi(n)}\)。

\textbf{第一步}：\(ar_i \not\equiv ar_j \pmod{n}\)（\(i \neq j\)）。若
\(ar_i \equiv ar_j \pmod{n}\)，由 \(\gcd(a, n) = 1\) 和定理 5.4.2' 得
\(r_i \equiv r_j \pmod{n}\)，但 \(1 \leq r_i, r_j \leq n\) 且
\(r_i \neq r_j\)，矛盾。

\textbf{第二步}：\(\gcd(ar_i, n) = 1\)。若 \(d = \gcd(ar_i, n) > 1\)，设
\(p\) 是 \(d\) 的素因子。\(p \mid n\) 且 \(p \mid ar_i\)。由
\(\gcd(a, n) = 1\) 知 \(p \nmid a\)，故 \(p \mid r_i\)。但
\(\gcd(r_i, n) = 1\)，故 \(p \nmid n\)，矛盾。

\textbf{第三步}：\(\{ar_1 \bmod n, \ldots, ar_{\varphi(n)} \bmod n\} = \{r_1, \ldots, r_{\varphi(n)}\}\)（作为集合）。因此
\[\prod_{i=1}^{\varphi(n)} ar_i \equiv \prod_{i=1}^{\varphi(n)} r_i \pmod{n}\]
即 \(a^{\varphi(n)} \prod r_i \equiv \prod r_i \pmod{n}\)。由
\(\gcd(\prod r_i, n) = 1\) 和定理
5.4.2'，\(a^{\varphi(n)} \equiv 1 \pmod{n}\)。\(\blacksquare\)

\textbf{推论 5.4.2（费马小定理的推广形式）}：设 \(p\)
是素数，则对任意整数 \(a\)： \[a^p \equiv a \pmod{p}\]

\textbf{证明}：若 \(p \mid a\)，则
\(a^p \equiv 0 \equiv a \pmod{p}\)。若
\(p \nmid a\)，由欧拉定理（\(\varphi(p) = p - 1\)），\(a^{p-1} \equiv 1 \pmod{p}\)，乘以
\(a\) 得 \(a^p \equiv a \pmod{p}\)。\(\blacksquare\)

\paragraph{5.4.6 威尔逊定理}\label{ux5a01ux5c14ux900aux5b9aux7406}

\textbf{定理 5.4.10（威尔逊定理）}：设 \(p\) 是素数。则
\[(p - 1)! \equiv -1 \pmod{p}\]

\textbf{证明}：考虑模 \(p\) 下 \(1, 2, \ldots, p-1\) 的乘法逆元。对
\(1 \leq a \leq p-1\)，存在唯一的 \(b\)（\(1 \leq b \leq p-1\)）使得
\(ab \equiv 1 \pmod{p}\)（由 \(\gcd(a, p) = 1\) 和裴蜀定理）。

\(a\) 是自身的逆元（\(a^2 \equiv 1 \pmod{p}\)）当且仅当
\(p \mid (a^2 - 1) = (a-1)(a+1)\)。由欧几里得引理，\(p \mid (a-1)\) 或
\(p \mid (a+1)\)，即 \(a = 1\) 或 \(a = p - 1\)。

因此 \(2, 3, \ldots, p-2\) 可以两两配对成
\((a, a^{-1})\)（\(a \neq a^{-1}\)），每对的乘积 \(\equiv 1 \pmod{p}\)。

故
\((p-1)! = 1 \cdot (2 \cdot 3 \cdots (p-2)) \cdot (p-1) \equiv 1 \cdot 1 \cdot (p-1) \equiv -1 \pmod{p}\)。\(\blacksquare\)

\paragraph{5.4.7 模逆元}\label{ux6a21ux9006ux5143}

\textbf{定义 5.4.3（模逆元）}：设 \(\gcd(a, m) = 1\)。满足
\(ax \equiv 1 \pmod{m}\) 的整数 \(x\) 称为 \(a\) 关于模 \(m\)
的\textbf{乘法逆元}，记为 \(a^{-1} \pmod{m}\)。

\textbf{定理 5.4.11（逆元的存在性与唯一性）}：\(\gcd(a, m) = 1\)
当且仅当 \(a\) 在模 \(m\) 下有乘法逆元。逆元在模 \(m\) 意义下唯一。

\textbf{证明}：\((\Rightarrow)\)：由裴蜀定理，存在 \(x, y\) 使得
\(ax + my = 1\)，即 \(ax \equiv 1 \pmod{m}\)，\(x\)
即为逆元。\((\Leftarrow)\)：若 \(ax \equiv 1 \pmod{m}\)，则
\(ax + my = 1\) 对某 \(y\)，故 \(\gcd(a, m) \mid 1\)，即
\(\gcd(a, m) = 1\)。唯一性：若 \(ax_1 \equiv 1\) 且
\(ax_2 \equiv 1 \pmod{m}\)，则 \(a(x_1 - x_2) \equiv 0 \pmod{m}\)，由
\(\gcd(a, m) = 1\) 得 \(x_1 \equiv x_2 \pmod{m}\)。\(\blacksquare\)

\textbf{算法
5.4.1（扩展欧几里得算法求逆元）}：由辗转相除法的逆向递推，可以高效计算
\(a^{-1} \pmod{m}\)。具体做法：对 \(a\) 和 \(m\) 做辗转相除法得到
\(\gcd(a, m) = 1\)，然后逆向递推表示 \(1 = ax + my\)，\(x\) 即为逆元。

\textbf{详细步骤}：设辗转相除法的过程为：

\[m = aq_1 + r_1, \quad a = r_1 q_2 + r_2, \quad r_1 = r_2 q_3 + r_3, \quad \ldots, \quad r_{k-2} = r_{k-1}q_k + r_k\]

其中 \(r_k = \gcd(a, m) = 1\)。逆向递推：

\[1 = r_k = r_{k-2} - r_{k-1}q_k\]
\[= r_{k-2} - (r_{k-3} - r_{k-2}q_{k-1})q_k = (1 + q_{k-1}q_k)r_{k-2} - q_k r_{k-3}\]

继续递推，最终得到 \(1 = ax + my\) 的形式。

\paragraph{5.4.8 模幂运算}\label{ux6a21ux5e42ux8fd0ux7b97}

\textbf{定理 5.4.12（快速模幂）}：计算 \(a^n \bmod m\) 可以在
\(O(\log n)\) 次模乘法内完成。

\textbf{方法（平方-乘算法）}：将 \(n\) 表示为二进制
\(n = \sum_{i=0}^{k} b_i 2^i\)。则
\[a^n = a^{\sum b_i 2^i} = \prod_{i: b_i = 1} a^{2^i}\] 依次计算
\(a^{2^0}, a^{2^1}, \ldots, a^{2^k} \pmod{m}\)（每次平方取模），然后将
\(b_i = 1\) 对应的项相乘取模。

\textbf{证明正确性}：\(a^{2^{i+1}} = (a^{2^i})^2\)。每步取模不影响最终结果（由同余的乘法性质）。\(\blacksquare\)

\textbf{注记 5.4.2}：快速模幂算法在密码学中有重要应用。例如在 RSA
加密中，需要计算 \(m^e \bmod n\)（\(e\) 和 \(n\) 都是大整数），直接计算
\(m^e\) 是不现实的，但通过平方-乘算法可以在合理时间内完成。

\textbf{具体例子}：计算
\(2^{100} \bmod 1000\)。\(100 = 64 + 32 + 4 = (1100100)_2\)。依次计算：
- \(2^1 \equiv 2\) - \(2^2 \equiv 4\) - \(2^4 \equiv 16\) -
\(2^8 \equiv 256\) - \(2^{16} \equiv 256^2 = 65536 \equiv 536\) -
\(2^{32} \equiv 536^2 = 287296 \equiv 296\) -
\(2^{64} \equiv 296^2 = 87616 \equiv 616\)

\(2^{100} = 2^{64} \cdot 2^{32} \cdot 2^4 \equiv 616 \cdot 296 \cdot 16 \pmod{1000}\)。\(616 \cdot 296 = 182336 \equiv 336\)。\(336 \cdot 16 = 5376 \equiv 376\)。故
\(2^{100} \equiv 376 \pmod{1000}\)。

\paragraph{5.4.9 同余方程}\label{ux540cux4f59ux65b9ux7a0b}

\textbf{定义 5.4.4（线性同余方程）}：形如 \(ax \equiv b \pmod{m}\)
的方程称为\textbf{线性同余方程}。

\textbf{定理 5.4.13（线性同余方程的解）}：设 \(d = \gcd(a, m)\)。 1.
\(ax \equiv b \pmod{m}\) 有解当且仅当 \(d \mid b\)。 2. 若有解，则在模
\(m\) 意义下恰好有 \(d\) 个解。

\textbf{证明}：

\begin{enumerate}
\def\labelenumi{(\arabic{enumi})}
\item
  \(ax \equiv b \pmod{m}\) 有解 \(\iff\) 存在 \(x, y\) 使得
  \(ax + my = b\) \(\iff\) \(\gcd(a, m) \mid b\)（由裴蜀定理）。
\item
  设 \(d \mid b\)。令 \(a = da'\), \(b = db'\),
  \(m = dm'\)。则原方程化为 \(a'x \equiv b' \pmod{m'}\)。由
  \(\gcd(a', m') = 1\)，存在唯一解 \(x_0 \pmod{m'}\)。原方程的解为
  \(x_0, x_0 + m', x_0 + 2m', \ldots, x_0 + (d-1)m' \pmod{m}\)，共 \(d\)
  个。\(\blacksquare\)
\end{enumerate}

\paragraph{5.4.10 二次剩余}\label{ux4e8cux6b21ux5269ux4f59}

\textbf{定义 5.4.5（二次剩余）}：设 \(p\) 是奇素数，\(a\) 是整数且
\(p \nmid a\)。若同余方程 \(x^2 \equiv a \pmod{p}\) 有解，则称 \(a\)
为模 \(p\) 的\textbf{二次剩余}；否则称为\textbf{二次非剩余}。

\textbf{定理 5.4.14（二次剩余的个数）}：模奇素数 \(p\) 的二次剩余恰好有
\(\frac{p-1}{2}\) 个，二次非剩余也有 \(\frac{p-1}{2}\) 个。

\textbf{证明}：考虑平方映射
\(f: \{1, 2, \ldots, p-1\} \to \{1, 2, \ldots, p-1\}\)，\(f(x) = x^2 \bmod p\)。

\(f(x) = f(y) \iff x^2 \equiv y^2 \pmod{p} \iff p \mid (x-y)(x+y)\)。由欧几里得引理，\(p \mid (x-y)\)
或 \(p \mid (x+y)\)。在 \(\{1, \ldots, p-1\}\) 中，\(p \mid (x-y)\)
意味着 \(x = y\)，\(p \mid (x+y)\) 意味着 \(x + y = p\)。因此 \(f\)
的像恰好是每个二次剩余对应两个原像（\(x\) 和 \(p - x\)）。故像的大小为
\(\frac{p-1}{2}\)。\(\blacksquare\)

\bigskip\hrule\bigskip

\subsubsection{本章展望}\label{ux672cux7ae0ux5c55ux671b-4}

本章建立了整除性理论、最大公约数与最小公倍数、素数与算术基本定理、以及同余理论的基本框架。这些结果是整数深层结构的基石。在卷二（Ch
6--11）中，裴蜀定理和辗转相除法将在多项式理论中出现类比（多项式的带余除法和最大公因式）。费马小定理和中国剩余定理是密码学（如
RSA
加密）的数学基础。更深入的数论主题（二次互反律、原根、连分数等）将在后续课程中展开。

\bigskip\hrule\bigskip

\emph{卷一终}

\newpage

\section{卷二：代数}\label{ux5377ux4e8cux4ee3ux6570}

\bigskip\hrule\bigskip

\subsection{Ch 6
代数运算与运算律}\label{ch-6-ux4ee3ux6570ux8fd0ux7b97ux4e0eux8fd0ux7b97ux5f8b}

\begin{quote}
\textbf{前置知识}：本章建立在第一卷数系构造（第4章）的基础之上。读者应熟悉皮亚诺公理、数学归纳法以及实数的完备性。
\end{quote}

\subsubsection{6.1
二元运算与代数结构}\label{ux4e8cux5143ux8fd0ux7b97ux4e0eux4ee3ux6570ux7ed3ux6784}

\textbf{定义 6.1.1（二元运算）}

设 \(S\) 为非空集合。从 \(S \times S\) 到 \(S\) 的映射

\[
\ast : S \times S \longrightarrow S, \quad (a, b) \longmapsto a \ast b
\]

称为 \(S\) 上的\textbf{二元运算}。

\textbf{定义 6.1.2（封闭性）}

设 \(\ast\) 为集合 \(S\) 上的二元运算。对任意 \(a, b \in S\)，有
\(a \ast b \in S\)，则称 \(\ast\) 在 \(S\) 上\textbf{封闭}。

\textbf{定义 6.1.3（结合律）}

设 \(\ast\) 为集合 \(S\) 上的二元运算。若对任意 \(a, b, c \in S\)，有

\[
(a \ast b) \ast c = a \ast (b \ast c)
\]

则称 \(\ast\) 满足\textbf{结合律}。

\textbf{定义 6.1.4（交换律）}

设 \(\ast\) 为集合 \(S\) 上的二元运算。若对任意 \(a, b \in S\)，有

\[
a \ast b = b \ast a
\]

则称 \(\ast\) 满足\textbf{交换律}。

\textbf{定义 6.1.5（幺元）}

设 \(\ast\) 为集合 \(S\) 上的二元运算。若存在 \(e \in S\)，使得对任意
\(a \in S\)，有

\[
e \ast a = a \ast e = a
\]

则称 \(e\) 为 \(\ast\) 的\textbf{幺元}（恒等元）。

\textbf{定理 6.1.1（幺元唯一性）} 若二元运算 \(\ast\)
存在幺元，则幺元唯一。

\textbf{证明}：设 \(e_1, e_2\) 均为 \(\ast\) 的幺元。由 \(e_1\)
的幺元性质，\(e_1 \ast e_2 = e_2\)；由 \(e_2\)
的幺元性质，\(e_1 \ast e_2 = e_1\)。因此 \(e_1 = e_2\)。
\(\blacksquare\)

\textbf{定义 6.1.6（逆元）}

设 \(\ast\) 为集合 \(S\) 上的二元运算，\(e\) 为幺元。对
\(a \in S\)，若存在 \(b \in S\)，使得

\[
a \ast b = b \ast a = e
\]

则称 \(b\) 为 \(a\) 的\textbf{逆元}，记为 \(a^{-1}\)。

\textbf{定理 6.1.2（逆元唯一性）}
在满足结合律的运算下，若元素的逆元存在，则逆元唯一。

\textbf{证明}：设 \(b, c\) 均为 \(a\) 的逆元，则

\[
b = b \ast e = b \ast (a \ast c) = (b \ast a) \ast c = e \ast c = c
\]

其中第三步利用了结合律。 \(\blacksquare\)

\textbf{定义 6.1.7（半群）}

集合 \(S\) 与其上的二元运算 \(\ast\) 构成的代数结构
\((S, \ast)\)，若满足：

\begin{enumerate}
\def\labelenumi{\arabic{enumi}.}
\tightlist
\item
  \textbf{封闭性}：\(\forall a, b \in S, \, a \ast b \in S\)；
\item
  \textbf{结合律}：\(\forall a, b, c \in S, \, (a \ast b) \ast c = a \ast (b \ast c)\)，
\end{enumerate}

则称 \((S, \ast)\) 为\textbf{半群}。

\textbf{定义 6.1.8（幺半群）}

若半群 \((S, \ast)\) 中存在幺元 \(e\)，则称 \((S, \ast)\)
为\textbf{幺半群}。

\textbf{定义 6.1.9（群）}

若幺半群 \((G, \ast)\) 中每个元素都有逆元，即

\[
\forall a \in G, \, \exists a^{-1} \in G, \, a \ast a^{-1} = a^{-1} \ast a = e
\]

则称 \((G, \ast)\)
为\textbf{群}。若进一步满足交换律，则称为\textbf{阿贝尔群}（交换群）。

\textbf{定义 6.1.10（环）}

集合 \(R\) 上配备两个二元运算 \(+\) 和
\(\cdot\)（分别称为加法和乘法），若满足：

\begin{enumerate}
\def\labelenumi{\arabic{enumi}.}
\tightlist
\item
  \((R, +)\) 是阿贝尔群（加法幺元记为 \(0\)）；
\item
  \((R, \cdot)\) 是半群（乘法满足结合律和封闭性）；
\item
  \textbf{分配律}：\(\forall a, b, c \in R\)，
\end{enumerate}

\[
a \cdot (b + c) = a \cdot b + a \cdot c, \quad (b + c) \cdot a = b \cdot a + c \cdot a
\]

则称 \((R, +, \cdot)\) 为\textbf{环}。

\textbf{定义 6.1.11（交换环与含幺环）}

\begin{itemize}
\tightlist
\item
  若环 \((R, +, \cdot)\) 中乘法满足交换律，则称为\textbf{交换环}。
\item
  若环 \((R, +, \cdot)\) 中乘法存在幺元
  \(1\)（\(1 \neq 0\)），则称为\textbf{含幺环}。
\end{itemize}

\textbf{定义 6.1.12（域）}

含幺交换环 \((F, +, \cdot)\) 中，若每个非零元素都有乘法逆元，即

\[
\forall a \in F \setminus \{0\}, \, \exists a^{-1} \in F, \, a \cdot a^{-1} = 1
\]

则称 \((F, +, \cdot)\) 为\textbf{域}。

\bigskip\hrule\bigskip

\subsubsection{6.2
自然数运算律}\label{ux81eaux7136ux6570ux8fd0ux7b97ux5f8b}

以下结果建立在皮亚诺公理之上（参见第一卷第4章）。

\textbf{定义 6.2.1（自然数加法）}

加法 \(+: \mathbb{N} \times \mathbb{N} \to \mathbb{N}\) 递归定义为：

\[
a + 0 = a, \quad a + S(b) = S(a + b)
\]

其中 \(S\) 为后继函数。

\textbf{定义 6.2.2（自然数乘法）}

乘法 \(\cdot: \mathbb{N} \times \mathbb{N} \to \mathbb{N}\) 递归定义为：

\[
a \cdot 0 = 0, \quad a \cdot S(b) = a \cdot b + a
\]

\textbf{定理 6.2.1（加法右单位元）} 对任意 \(a \in \mathbb{N}\)，有
\(a + 0 = a\)。

\textbf{证明}：由加法定义的第一条直接得到。 \(\blacksquare\)

\textbf{定理 6.2.2（加法左单位元）} 对任意 \(a \in \mathbb{N}\)，有
\(0 + a = a\)。

\textbf{证明}：对 \(a\) 进行归纳。

\begin{itemize}
\tightlist
\item
  \textbf{基础}：\(a = 0\) 时，\(0 + 0 = 0\)，成立。
\item
  \textbf{归纳}：假设 \(0 + a = a\)，则 \(0 + S(a) = S(0 + a) = S(a)\)。
\end{itemize}

由归纳法原理，结论成立。 \(\blacksquare\)

\textbf{定理 6.2.3（加法结合律）} 对任意 \(a, b, c \in \mathbb{N}\)，有
\((a + b) + c = a + (b + c)\)。

\textbf{证明}：对 \(c\) 进行归纳。

\begin{itemize}
\tightlist
\item
  \textbf{基础}：\(c = 0\) 时，\((a + b) + 0 = a + b = a + (b + 0)\)。
\item
  \textbf{归纳}：假设 \((a + b) + c = a + (b + c)\)，则
\end{itemize}

\[
(a + b) + S(c) = S((a + b) + c) = S(a + (b + c)) = a + S(b + c) = a + (b + S(c))
\]

由归纳法原理，结论成立。 \(\blacksquare\)

\textbf{定理 6.2.4（后继的加法表示）} 对任意 \(a \in \mathbb{N}\)，有
\(S(a) = a + 1\)。

\textbf{证明}：由定义，\(1 = S(0)\)，故
\(a + 1 = a + S(0) = S(a + 0) = S(a)\)。 \(\blacksquare\)

\textbf{定理 6.2.5（加法交换律）} 对任意 \(a, b \in \mathbb{N}\)，有
\(a + b = b + a\)。

\textbf{证明}：先证辅助引理：对任意 \(a \in \mathbb{N}\)，有
\(S(a) = a + 1 = 1 + a\)。

对 \(a\) 归纳证明 \(1 + a = S(a)\)：

\begin{itemize}
\tightlist
\item
  \textbf{基础}：\(1 + 0 = 1 = S(0)\)。
\item
  \textbf{归纳}：假设 \(1 + a = S(a)\)，则
  \(1 + S(a) = S(1 + a) = S(S(a))\)。
\end{itemize}

现在对 \(b\) 归纳证明 \(a + b = b + a\)：

\begin{itemize}
\tightlist
\item
  \textbf{基础}：\(b = 0\) 时，\(a + 0 = a = 0 + a\)。
\item
  \textbf{归纳}：假设 \(a + b = b + a\)，则
\end{itemize}

\[
a + S(b) = S(a + b) = S(b + a) = b + S(a) = b + (a + 1) = b + (1 + a) = (b + 1) + a = S(b) + a
\]

由归纳法原理，结论成立。 \(\blacksquare\)

\textbf{定理 6.2.6（乘法右零元）} 对任意 \(a \in \mathbb{N}\)，有
\(a \cdot 0 = 0\)。

\textbf{证明}：由乘法定义直接得到。 \(\blacksquare\)

\textbf{定理 6.2.7（乘法左零元）} 对任意 \(a \in \mathbb{N}\)，有
\(0 \cdot a = 0\)。

\textbf{证明}：对 \(a\) 进行归纳。

\begin{itemize}
\tightlist
\item
  \textbf{基础}：\(0 \cdot 0 = 0\)。
\item
  \textbf{归纳}：假设 \(0 \cdot a = 0\)，则
  \(0 \cdot S(a) = 0 \cdot a + 0 = 0 + 0 = 0\)。
\end{itemize}

由归纳法原理，结论成立。 \(\blacksquare\)

\textbf{定理 6.2.8（乘法右单位元）} 对任意 \(a \in \mathbb{N}\)，有
\(a \cdot 1 = a\)。

\textbf{证明}：\(a \cdot 1 = a \cdot S(0) = a \cdot 0 + a = 0 + a = a\)。
\(\blacksquare\)

\textbf{定理 6.2.9（乘法左分配律）} 对任意
\(a, b, c \in \mathbb{N}\)，有
\(a \cdot (b + c) = a \cdot b + a \cdot c\)。

\textbf{证明}：对 \(c\) 进行归纳。

\begin{itemize}
\tightlist
\item
  \textbf{基础}：\(c = 0\)
  时，\(a \cdot (b + 0) = a \cdot b = a \cdot b + 0 = a \cdot b + a \cdot 0\)。
\item
  \textbf{归纳}：假设 \(a \cdot (b + c) = a \cdot b + a \cdot c\)，则
\end{itemize}

\[
a \cdot (b + S(c)) = a \cdot S(b + c) = a \cdot (b + c) + a = (a \cdot b + a \cdot c) + a = a \cdot b + (a \cdot c + a) = a \cdot b + a \cdot S(c)
\]

由归纳法原理，结论成立。 \(\blacksquare\)

\textbf{定理 6.2.10（乘法结合律）} 对任意 \(a, b, c \in \mathbb{N}\)，有
\((a \cdot b) \cdot c = a \cdot (b \cdot c)\)。

\textbf{证明}：对 \(c\) 进行归纳。

\begin{itemize}
\tightlist
\item
  \textbf{基础}：\(c = 0\)
  时，\((a \cdot b) \cdot 0 = 0 = a \cdot 0 = a \cdot (b \cdot 0)\)。
\item
  \textbf{归纳}：假设 \((a \cdot b) \cdot c = a \cdot (b \cdot c)\)，则
\end{itemize}

\[
(a \cdot b) \cdot S(c) = (a \cdot b) \cdot c + a \cdot b = a \cdot (b \cdot c) + a \cdot b = a \cdot (b \cdot c + b) = a \cdot (b \cdot S(c))
\]

其中第三步利用了左分配律。 \(\blacksquare\)

\textbf{定理 6.2.11（乘法交换律）} 对任意 \(a, b \in \mathbb{N}\)，有
\(a \cdot b = b \cdot a\)。

\textbf{证明}：先证辅助引理：对任意 \(a \in \mathbb{N}\)，有
\(1 \cdot a = a\)。对 \(a\) 归纳：

\begin{itemize}
\tightlist
\item
  \textbf{基础}：\(1 \cdot 0 = 0\)。
\item
  \textbf{归纳}：假设 \(1 \cdot a = a\)，则
  \(1 \cdot S(a) = 1 \cdot a + 1 = a + 1 = S(a)\)。
\end{itemize}

再证另一辅助引理：\(S(a) \cdot b = a \cdot b + b\)。对 \(b\) 归纳：

\begin{itemize}
\tightlist
\item
  \textbf{基础}：\(S(a) \cdot 0 = 0 = 0 + 0 = a \cdot 0 + 0\)。
\item
  \textbf{归纳}：假设 \(S(a) \cdot b = a \cdot b + b\)，则
\end{itemize}

\[
S(a) \cdot S(b) = S(a) \cdot b + S(a) = (a \cdot b + b) + (a + 1) = (a \cdot b + a) + (b + 1) = a \cdot S(b) + S(b)
\]

现在对 \(b\) 归纳证明 \(a \cdot b = b \cdot a\)：

\begin{itemize}
\tightlist
\item
  \textbf{基础}：\(a \cdot 0 = 0 = 0 \cdot a\)。
\item
  \textbf{归纳}：假设 \(a \cdot b = b \cdot a\)，则
  \(a \cdot S(b) = a \cdot b + a = b \cdot a + a = S(b) \cdot a\)。
\end{itemize}

由归纳法原理，结论成立。 \(\blacksquare\)

\textbf{定理 6.2.12（乘法右分配律）} 对任意
\(a, b, c \in \mathbb{N}\)，有
\((a + b) \cdot c = a \cdot c + b \cdot c\)。

\textbf{证明}：由乘法交换律和左分配律，

\[
(a + b) \cdot c = c \cdot (a + b) = c \cdot a + c \cdot b = a \cdot c + b \cdot c
\]

\(\blacksquare\)

\bigskip\hrule\bigskip

\subsubsection{6.3 整数运算律}\label{ux6574ux6570ux8fd0ux7b97ux5f8b}

整数 \(\mathbb{Z}\) 通过自然数的等价类构造（参见第一卷第4章）。设
\(a, b \in \mathbb{N}\)，整数 \(-a\) 表示使得 \(a + (-a) = 0\) 的元素。

\textbf{引理
6.3.0}：整数的乘法交换律和右分配律可从自然数继承的左分配律和乘法交换律推导。

\textbf{证明}：由自然数的乘法交换律（定理
6.2.11），通过等价类构造可得整数的乘法交换律 \(ab = ba\)。

右分配律
\((a + b) \cdot c = a \cdot c + b \cdot c\)：由乘法交换律，\((a+b) \cdot c = c \cdot (a+b) = c \cdot a + c \cdot b = a \cdot c + b \cdot c\)（第二步利用了已有的左分配律）。

因此下文的分配律推广 \((a+b)(c+d) = ac + ad + bc + bd\)
合法。\(\blacksquare\)

\textbf{定理 6.3.1（负负得正）} \((-1) \times (-1) = 1\)。

\textbf{证明}：考虑以下推导链。

由分配律的推广，对任意 \(a, b, c, d \in \mathbb{Z}\)：

\[
(a + b)(c + d) = ac + ad + bc + bd
\]

取 \(a = 1, b = -1, c = 1, d = -1\)：

\[
(1 + (-1))(1 + (-1)) = 1 \cdot 1 + 1 \cdot (-1) + (-1) \cdot 1 + (-1) \cdot (-1)
\]

左边为 \(0 \cdot 0 = 0\)。右边中
\(1 \cdot 1 = 1\)，\(1 \cdot (-1) = -1\)（由分配律：\(1 \cdot (-1) + 1 \cdot 1 = 1 \cdot ((-1) + 1) = 1 \cdot 0 = 0\)，故
\(1 \cdot (-1) = -1\)），\((-1) \cdot 1 = -1\)（由交换律），因此：

\[
0 = 1 + (-1) + (-1) + (-1)(-1) = -1 + (-1)(-1)
\]

从而 \((-1)(-1) = 1\)。 \(\blacksquare\)

\textbf{定理 6.3.2（负数乘法的一般法则）} 对任意
\(a, b \in \mathbb{Z}\)，有 \((-a)(-b) = ab\)。

\textbf{证明}：

\[
(-a)(-b) = (-1 \cdot a)(-1 \cdot b) = (-1)(-1) \cdot ab = 1 \cdot ab = ab
\]

其中利用了乘法结合律和交换律以及定理6.3.1。 \(\blacksquare\)

\textbf{定理 6.3.3} 对任意 \(a, b \in \mathbb{Z}\)，有
\((-a) \cdot b = -(a \cdot b)\)。

\textbf{证明}：需证 \((-a) \cdot b + a \cdot b = 0\)。由分配律，

\[
(-a) \cdot b + a \cdot b = ((-a) + a) \cdot b = 0 \cdot b = 0
\]

由逆元唯一性，\((-a) \cdot b = -(a \cdot b)\)。 \(\blacksquare\)

\bigskip\hrule\bigskip

\subsubsection{6.4
实数运算律（域公理）}\label{ux5b9eux6570ux8fd0ux7b97ux5f8bux57dfux516cux7406}

实数集 \(\mathbb{R}\) 配备加法和乘法运算构成域。以下为完整的域公理体系。

\textbf{加法公理（A1---A4）}：

\begin{itemize}
\tightlist
\item
  \textbf{(A1)
  加法封闭性}：\(\forall a, b \in \mathbb{R}\)，\(a + b \in \mathbb{R}\)。
\item
  \textbf{(A2)
  加法结合律}：\(\forall a, b, c \in \mathbb{R}\)，\((a + b) + c = a + (b + c)\)。
\item
  \textbf{(A3)
  加法交换律}：\(\forall a, b \in \mathbb{R}\)，\(a + b = b + a\)。
\item
  \textbf{(A4) 加法单位元与逆元}：存在 \(0 \in \mathbb{R}\)，使得
  \(\forall a \in \mathbb{R}\)，\(a + 0 = a\)；对每个
  \(a \in \mathbb{R}\)，存在 \(-a \in \mathbb{R}\)，使得
  \(a + (-a) = 0\)。
\end{itemize}

\textbf{乘法公理（M1---M4）}：

\begin{itemize}
\tightlist
\item
  \textbf{(M1)
  乘法封闭性}：\(\forall a, b \in \mathbb{R}\)，\(a \cdot b \in \mathbb{R}\)。
\item
  \textbf{(M2)
  乘法结合律}：\(\forall a, b, c \in \mathbb{R}\)，\((a \cdot b) \cdot c = a \cdot (b \cdot c)\)。
\item
  \textbf{(M3)
  乘法交换律}：\(\forall a, b \in \mathbb{R}\)，\(a \cdot b = b \cdot a\)。
\item
  \textbf{(M4) 乘法单位元与逆元}：存在
  \(1 \in \mathbb{R}\)（\(1 \neq 0\)），使得
  \(\forall a \in \mathbb{R}\)，\(a \cdot 1 = a\)；对每个
  \(a \in \mathbb{R} \setminus \{0\}\)，存在
  \(a^{-1} \in \mathbb{R}\)，使得 \(a \cdot a^{-1} = 1\)。
\end{itemize}

\textbf{分配律（D）}：

\begin{itemize}
\tightlist
\item
  \textbf{(D)
  分配律}：\(\forall a, b, c \in \mathbb{R}\)，\(a \cdot (b + c) = a \cdot b + a \cdot c\)。
\end{itemize}

\textbf{定理 6.4.1（加法消去律）} 若 \(a + c = b + c\)，则 \(a = b\)。

\textbf{证明}：\(a = a + 0 = a + (c + (-c)) = (a + c) + (-c) = (b + c) + (-c) = b + (c + (-c)) = b + 0 = b\)。
\(\blacksquare\)

\textbf{定理 6.4.2（乘法消去律）} 若 \(a \cdot c = b \cdot c\) 且
\(c \neq 0\)，则 \(a = b\)。

\textbf{证明}：\(a = a \cdot 1 = a \cdot (c \cdot c^{-1}) = (a \cdot c) \cdot c^{-1} = (b \cdot c) \cdot c^{-1} = b \cdot (c \cdot c^{-1}) = b \cdot 1 = b\)。
\(\blacksquare\)

\textbf{定理 6.4.3（零因子性质）} 若 \(a \cdot b = 0\)，则 \(a = 0\) 或
\(b = 0\)。

\textbf{证明}：假设 \(a \neq 0\)，则 \(a^{-1}\)
存在。\(b = 1 \cdot b = (a^{-1} \cdot a) \cdot b = a^{-1} \cdot (a \cdot b) = a^{-1} \cdot 0 = 0\)。

其中最后一步 \(a^{-1} \cdot 0 = 0\) 的证明如下：设
\(x = a^{-1} \cdot 0\)，则
\(x = a^{-1} \cdot (0 + 0) = a^{-1} \cdot 0 + a^{-1} \cdot 0 = x + x\)。由消去律
\(x = 0\)。 \(\blacksquare\)

\bigskip\hrule\bigskip

\subsubsection{6.5 指数运算律}\label{ux6307ux6570ux8fd0ux7b97ux5f8b}

\paragraph{6.5.1 正整数指数}\label{ux6b63ux6574ux6570ux6307ux6570}

\textbf{定义 6.5.1} 设
\(a \in \mathbb{R}\)，\(n \in \mathbb{N}^+\)，归纳定义：

\[
a^1 = a, \quad a^{n+1} = a^n \cdot a
\]

\textbf{定理 6.5.1（指数加法法则）} 对
\(a \in \mathbb{R}\)，\(m, n \in \mathbb{N}^+\)，有
\(a^m \cdot a^n = a^{m+n}\)。

\textbf{证明}：对 \(n\) 归纳。

\begin{itemize}
\tightlist
\item
  \textbf{基础}：\(n = 1\)
  时，\(a^m \cdot a^1 = a^m \cdot a = a^{m+1}\)。
\item
  \textbf{归纳}：假设 \(a^m \cdot a^n = a^{m+n}\)，则
  \(a^m \cdot a^{n+1} = a^m \cdot (a^n \cdot a) = (a^m \cdot a^n) \cdot a = a^{m+n} \cdot a = a^{m+n+1}\)。
\end{itemize}

由归纳法原理，结论成立。 \(\blacksquare\)

\textbf{定理 6.5.2（指数乘法法则）} 对
\(a \in \mathbb{R}\)，\(m, n \in \mathbb{N}^+\)，有
\((a^m)^n = a^{mn}\)。

\textbf{证明}：对 \(n\) 归纳。

\begin{itemize}
\tightlist
\item
  \textbf{基础}：\((a^m)^1 = a^m = a^{m \cdot 1}\)。
\item
  \textbf{归纳}：假设 \((a^m)^n = a^{mn}\)，则
  \((a^m)^{n+1} = (a^m)^n \cdot a^m = a^{mn} \cdot a^m = a^{mn+m} = a^{m(n+1)}\)。
\end{itemize}

\(\blacksquare\)

\textbf{定理 6.5.3（积的幂）} 对
\(a, b \in \mathbb{R}\)，\(n \in \mathbb{N}^+\)，有
\((ab)^n = a^n b^n\)。

\textbf{证明}：对 \(n\) 归纳。

\begin{itemize}
\tightlist
\item
  \textbf{基础}：\((ab)^1 = ab = a^1 b^1\)。
\item
  \textbf{归纳}：假设 \((ab)^n = a^n b^n\)，则
  \((ab)^{n+1} = (ab)^n \cdot ab = a^n b^n \cdot a \cdot b = a^{n+1} b^{n+1}\)。
\end{itemize}

\(\blacksquare\)

\paragraph{6.5.2
零指数与负整数指数}\label{ux96f6ux6307ux6570ux4e0eux8d1fux6574ux6570ux6307ux6570}

\textbf{定义 6.5.2} 设 \(a \in \mathbb{R}\)，\(a \neq 0\)：

\[
a^0 = 1, \quad a^{-n} = \frac{1}{a^n} \quad (n \in \mathbb{N}^+)
\]

\textbf{定理 6.5.4} 对
\(a \neq 0\)，\(m, n \in \mathbb{Z}\)，上述三条运算律仍然成立。

\textbf{证明}（以加法法则为例，证
\(a^m \cdot a^n = a^{m+n}\)）：分情况讨论。

\textbf{情形一}：\(m, n > 0\)。由定理6.5.1成立。

\textbf{情形二}：\(m > 0, n = 0\)。\(a^m \cdot a^0 = a^m \cdot 1 = a^m = a^{m+0}\)。

\textbf{情形三}：\(m > 0, n < 0\)。设 \(n = -k\)（\(k > 0\)）。

\begin{itemize}
\tightlist
\item
  若
  \(m > k\)：\(a^m \cdot a^{-k} = \dfrac{a^m}{a^k} = a^{m-k} = a^{m+(-k)}\)。
\item
  若
  \(m = k\)：\(a^m \cdot a^{-m} = \dfrac{a^m}{a^m} = 1 = a^0 = a^{m+(-m)}\)。
\item
  若
  \(m < k\)：\(a^m \cdot a^{-k} = \dfrac{a^m}{a^k} = \dfrac{1}{a^{k-m}} = a^{-(k-m)} = a^{m-k}\)。
\end{itemize}

\textbf{情形四}：\(m = 0\)。\(a^0 \cdot a^n = 1 \cdot a^n = a^n = a^{0+n}\)。

\textbf{情形五}：\(m < 0, n < 0\)。\(a^m \cdot a^{-n} = \dfrac{1}{a^{|m|}} \cdot \dfrac{1}{a^{|n|}} = \dfrac{1}{a^{|m|+|n|}} = a^{-(|m|+|n|)} = a^{m+n}\)。

所有情形下 \(a^m \cdot a^n = a^{m+n}\)
成立。指数乘法法则和积的幂的推广类似可证。 \(\blacksquare\)

\paragraph{6.5.3 有理数指数}\label{ux6709ux7406ux6570ux6307ux6570}

\textbf{定义 6.5.3} 设
\(a > 0\)，\(m, n \in \mathbb{Z}\)，\(n \geq 1\)，定义：

\[
a^{m/n} = \left(\sqrt[n]{a}\right)^m
\]

其中 \(\sqrt[n]{a}\) 为 \(a\) 的 \(n\) 次算术根，即满足 \(x^n = a\) 且
\(x > 0\) 的唯一实数。

\textbf{定理 6.5.5（有理数指数的良定义性）} 若 \(m/n = p/q\)（即
\(mq = np\)），则 \(a^{m/n} = a^{p/q}\)。

\textbf{证明}：设 \(m/n = p/q\)，即 \(mq = np\)。需证
\((\sqrt[n]{a})^m = (\sqrt[q]{a})^p\)。

设 \(x = \sqrt[n]{a}\)，\(y = \sqrt[q]{a}\)，则
\(x^n = a\)，\(y^q = a\)。对等式两端取 \(nq\) 次幂：

\[
\left((\sqrt[n]{a})^m\right)^{nq} = (\sqrt[n]{a})^{mnq} = \left((\sqrt[n]{a})^n\right)^{mq} = a^{mq}
\]

\[
\left((\sqrt[q]{a})^p\right)^{nq} = (\sqrt[q]{a})^{pnq} = \left((\sqrt[q]{a})^q\right)^{np} = a^{np}
\]

由 \(mq = np\)，两式右端相等，即
\(\left((\sqrt[n]{a})^m\right)^{nq} = \left((\sqrt[q]{a})^p\right)^{nq}\)。由于
\(nq > 0\) 且 \((\sqrt[n]{a})^m > 0\)，\((\sqrt[q]{a})^p > 0\)，正数的
\(nq\) 次幂是单射，故 \((\sqrt[n]{a})^m = (\sqrt[q]{a})^p\)。
\(\blacksquare\)

\textbf{定理 6.5.6（有理数指数的运算律）} 设
\(a > 0\)，\(r, s \in \mathbb{Q}\)，则：

\begin{enumerate}
\def\labelenumi{(\roman{enumi})}
\item
  \(a^r \cdot a^s = a^{r+s}\)；
\item
  \((a^r)^s = a^{rs}\)；
\item
  \((ab)^r = a^r \cdot b^r\)（\(b > 0\)）。
\end{enumerate}

\textbf{证明}（以 (i) 为例）：设 \(r = m/n\)，\(s = p/q\)，其中
\(n, q \geq 1\)。

\[
a^r \cdot a^s = (\sqrt[n]{a})^m \cdot (\sqrt[q]{a})^p
\]

令 \(d = \text{lcm}(n, q)\)，\(n' = d/n\)，\(q' = d/q\)。则
\(\sqrt[n]{a} = (\sqrt[d]{a})^{n'}\)，\(\sqrt[q]{a} = (\sqrt[d]{a})^{q'}\)。于是

\[
a^r = (\sqrt[d]{a})^{mn'}, \quad a^s = (\sqrt[d]{a})^{pq'}
\]

\[
a^r \cdot a^s = (\sqrt[d]{a})^{mn' + pq'}
\]

而 \(r + s = m/n + p/q = (mq + np)/(nq)\)。由于
\(mn' + pq' = md/n + pd/q = (mq + np) \cdot d/(nq)\)，且
\((\sqrt[d]{a})^{d/(nq)}\) 对应分母为 \(nq\) 的根\ldots{} 实际上
\(a^{r+s} = (\sqrt[nq]{a})^{mq+np}\)，而

\[
(\sqrt[d]{a})^{mn'+pq'} = a^{(mn'+pq')/d} = a^{m/n + p/q} = a^{r+s}
\]

\(\blacksquare\)

\paragraph{6.5.4 实数指数}\label{ux5b9eux6570ux6307ux6570}

\textbf{定义 6.5.4} 设 \(a > 0\)，\(x \in \mathbb{R}\)。定义

\[
a^x = \sup \{ a^r : r \in \mathbb{Q}, \, r \leq x \}
\]

（利用实数的完备性保证上确界存在。）

\textbf{定理
6.5.7}（实数指数运算律------陈述）实数指数满足与有理数指数相同的三条运算律。

\textbf{证明思路}：以 (i) \(a^x \cdot a^y = a^{x+y}\) 为例。取有理数列
\(r_n \uparrow x\)，\(s_n \uparrow y\)。由有理数指数的运算律（定理
6.5.6），\(a^{r_n} \cdot a^{s_n} = a^{r_n + s_n}\)。由上确界定义和单调有界性，\(a^{r_n} \to a^x\)，\(a^{s_n} \to a^y\)，\(a^{r_n+s_n} \to a^{x+y}\)（此处使用单调有界定理，其严格表述见卷四定理
22.5（单调有界定理））。由极限运算律得 \(a^x \cdot a^y = a^{x+y}\)。

\textbf{注}：上述''极限''记号的严格化需要极限理论（Ch
18-19）。在当前阶段，此定理可视为实数指数定义的自然延伸，完整证明在建立极限理论后补全。

\bigskip\hrule\bigskip

\subsubsection{6.6
分式运算与二次根式}\label{ux5206ux5f0fux8fd0ux7b97ux4e0eux4e8cux6b21ux6839ux5f0f}

\textbf{定义 6.6.1（二次根式）}

设 \(a \geq 0\)，\(\sqrt{a}\) 表示满足 \(x^2 = a\) 且 \(x \geq 0\)
的唯一非负实数，称为 \(a\) 的\textbf{算术平方根}。

\textbf{定理 6.6.1（根式乘法法则）} 设 \(a \geq 0\)，\(b \geq 0\)，则
\(\sqrt{a} \cdot \sqrt{b} = \sqrt{ab}\)。

\textbf{证明}：设 \(x = \sqrt{a}\)，\(y = \sqrt{b}\)，则
\(x \geq 0\)，\(y \geq 0\)，\(x^2 = a\)，\(y^2 = b\)。

考虑 \(xy \geq 0\)，且

\[
(xy)^2 = x^2 y^2 = ab
\]

由算术平方根的唯一性，\(xy = \sqrt{ab}\)，即
\(\sqrt{a} \cdot \sqrt{b} = \sqrt{ab}\)。 \(\blacksquare\)

\textbf{定理 6.6.2（根式除法法则）} 设 \(a \geq 0\)，\(b > 0\)，则
\(\dfrac{\sqrt{a}}{\sqrt{b}} = \sqrt{\dfrac{a}{b}}\)。

\textbf{证明}：设 \(x = \sqrt{a}/\sqrt{b} \geq 0\)，则

\[
x^2 = \frac{(\sqrt{a})^2}{(\sqrt{b})^2} = \frac{a}{b}
\]

由算术平方根的唯一性，\(x = \sqrt{a/b}\)。 \(\blacksquare\)

\textbf{定义 6.6.2（分母有理化）}

将分母中的根号消去的变换称为\textbf{分母有理化}。

\textbf{定理 6.6.3} 设 \(a \geq 0\)，\(b > 0\)，则

\[
\frac{a}{\sqrt{b}} = \frac{a\sqrt{b}}{b}
\]

\textbf{证明}：

\[
\frac{a}{\sqrt{b}} = \frac{a}{\sqrt{b}} \cdot \frac{\sqrt{b}}{\sqrt{b}} = \frac{a\sqrt{b}}{(\sqrt{b})^2} = \frac{a\sqrt{b}}{b}
\]

\(\blacksquare\)

\textbf{定理 6.6.4（共轭有理化）} 设
\(a \geq 0\)，\(b > 0\)，\(a \neq b\)，则

\[
\frac{1}{\sqrt{a} + \sqrt{b}} = \frac{\sqrt{a} - \sqrt{b}}{a - b} \quad (a \neq b)
\]

\textbf{证明}：

\[
\frac{1}{\sqrt{a} + \sqrt{b}} = \frac{\sqrt{a} - \sqrt{b}}{(\sqrt{a} + \sqrt{b})(\sqrt{a} - \sqrt{b})} = \frac{\sqrt{a} - \sqrt{b}}{a - b}
\]

其中利用了 \((\sqrt{a} + \sqrt{b})(\sqrt{a} - \sqrt{b}) = a - b\)。
\(\blacksquare\)

\begin{quote}
\textbf{本章展望}：本章建立了代数运算的公理化基础。从半群到域的层级结构，以及自然数、整数、实数上的具体运算律，为后续多项式、方程和不等式的理论奠定了基石。第7章将在此基础上研究多项式的代数结构。
\end{quote}

\bigskip\hrule\bigskip

\subsection{Ch 7 多项式}\label{ch-7-ux591aux9879ux5f0f}

\begin{quote}
\textbf{前置知识}：本章需要第6章关于实数域运算律（域公理）的知识，以及带余除法所需的基本实数性质。
\end{quote}

\subsubsection{7.1
多项式的定义与运算}\label{ux591aux9879ux5f0fux7684ux5b9aux4e49ux4e0eux8fd0ux7b97}

\textbf{定义 7.1.1（多项式）}

设 \(x\)
为未定元（变量），\(a_0, a_1, \ldots, a_n \in \mathbb{R}\)（\(a_n \neq 0\)），表达式

\[
f(x) = a_n x^n + a_{n-1} x^{n-1} + \cdots + a_1 x + a_0 = \sum_{k=0}^{n} a_k x^k
\]

称为关于 \(x\) 的 \textbf{\(n\) 次多项式}。其中 \(a_n\)
称为\textbf{首项系数}，\(a_0\)
称为\textbf{常数项}。若所有系数均为零，称为\textbf{零多项式}，记为
\(0\)，不定义其次数。

\textbf{定义 7.1.2（多项式的相等）}

两个多项式 \(f(x) = \sum a_k x^k\) 和 \(g(x) = \sum b_k x^k\)
\textbf{相等}，当且仅当对所有 \(k\)，\(a_k = b_k\)。

\textbf{定义 7.1.3（多项式的加法）}

设
\(f(x) = \sum_{k=0}^{n} a_k x^k\)，\(g(x) = \sum_{k=0}^{m} b_k x^k\)，定义

\[
(f + g)(x) = \sum_{k=0}^{\max(n,m)} (a_k + b_k) x^k
\]

其中当 \(k\) 超出某多项式次数时，对应系数视为 \(0\)。

\textbf{定义 7.1.4（多项式的乘法）}

\[
(f \cdot g)(x) = \sum_{k=0}^{n+m} c_k x^k, \quad c_k = \sum_{i+j=k} a_i b_j
\]

\textbf{定理 7.1.1} 若 \(f(x)\) 为 \(n\) 次多项式，\(g(x)\) 为 \(m\)
次多项式（\(f, g\) 均非零），则 \(f(x) \cdot g(x)\) 为 \(n + m\)
次多项式。

\textbf{证明}：设
\(f(x) = a_n x^n + \cdots\)，\(g(x) = b_m x^m + \cdots\)（\(a_n \neq 0\)，\(b_m \neq 0\)）。乘积的最高次项系数为
\(a_n b_m\)。由于 \(\mathbb{R}\) 是域，\(a_n b_m \neq 0\)。而 \(n + m\)
次以上各项系数均为零（因 \(i + j > n + m\) 时 \(a_i = 0\) 或
\(b_j = 0\)）。故 \(\deg(fg) = n + m\)。 \(\blacksquare\)

\textbf{定理 7.1.2（多项式环）} 实系数多项式全体 \(\mathbb{R}[x]\)
在上述加法和乘法下构成含幺交换环。

\textbf{证明}：

\begin{itemize}
\tightlist
\item
  \((\mathbb{R}[x], +)\)
  是阿贝尔群：加法封闭、结合、交换，零多项式为加法单位元，\(-f(x) = \sum(-a_k)x^k\)
  为 \(f\) 的加法逆元。
\item
  乘法封闭（由多项式乘法定义）、结合律（系数卷积满足结合律：\((f \cdot g) \cdot h\)
  与 \(f \cdot (g \cdot h)\) 的 \(k\) 次项系数均为
  \(\sum_{i+j+\ell=k} a_i b_j c_\ell\)，由求和的结合律和交换律保证两式相等）、交换律（\(c_k = \sum_{i+j=k} a_i b_j = \sum_{i+j=k} b_j a_i\)，由实数乘法交换律）由实数的对应性质逐项验证。
\item
  分配律由系数层面的分配律逐项验证。
\item
  幺元为常数多项式 \(1\)。
\end{itemize}

\(\blacksquare\)

\bigskip\hrule\bigskip

\subsubsection{7.2 乘法公式}\label{ux4e58ux6cd5ux516cux5f0f}

以下每个公式均通过直接展开并利用分配律证明。

\textbf{定理 7.2.1（完全平方和）} 对任意 \(a, b \in \mathbb{R}\)：

\[
(a + b)^2 = a^2 + 2ab + b^2
\]

\textbf{证明}：

\[
(a + b)^2 = (a + b)(a + b) = a \cdot a + a \cdot b + b \cdot a + b \cdot b = a^2 + ab + ba + b^2 = a^2 + 2ab + b^2
\]

\(\blacksquare\)

\textbf{定理 7.2.2（完全平方差）} 对任意 \(a, b \in \mathbb{R}\)：

\[
(a - b)^2 = a^2 - 2ab + b^2
\]

\textbf{证明}：

\[
(a - b)^2 = (a + (-b))^2 = a^2 + 2a(-b) + (-b)^2 = a^2 - 2ab + b^2
\]

\(\blacksquare\)

\textbf{定理 7.2.3（平方差公式）} 对任意 \(a, b \in \mathbb{R}\)：

\[
(a + b)(a - b) = a^2 - b^2
\]

\textbf{证明}：

\[
(a + b)(a - b) = a^2 - ab + ba - b^2 = a^2 - ab + ab - b^2 = a^2 - b^2
\]

\(\blacksquare\)

\textbf{定理 7.2.4（完全立方和）} 对任意 \(a, b \in \mathbb{R}\)：

\[
(a + b)^3 = a^3 + 3a^2 b + 3ab^2 + b^3
\]

\textbf{证明}：

\[
(a + b)^3 = (a + b)^2 (a + b) = (a^2 + 2ab + b^2)(a + b)
\]

\[
= a^3 + a^2 b + 2a^2 b + 2ab^2 + ab^2 + b^3 = a^3 + 3a^2 b + 3ab^2 + b^3
\]

\(\blacksquare\)

\textbf{定理 7.2.5（立方和公式）} 对任意 \(a, b \in \mathbb{R}\)：

\[
a^3 + b^3 = (a + b)(a^2 - ab + b^2)
\]

\textbf{证明}：展开右端：

\[
(a + b)(a^2 - ab + b^2) = a^3 - a^2 b + ab^2 + a^2 b - ab^2 + b^3 = a^3 + b^3
\]

\(\blacksquare\)

\textbf{定理 7.2.6（立方差公式）} 对任意 \(a, b \in \mathbb{R}\)：

\[
a^3 - b^3 = (a - b)(a^2 + ab + b^2)
\]

\textbf{证明}：展开右端：

\[
(a - b)(a^2 + ab + b^2) = a^3 + a^2 b + ab^2 - a^2 b - ab^2 - b^3 = a^3 - b^3
\]

\(\blacksquare\)

\bigskip\hrule\bigskip

\subsubsection{7.3 因式分解}\label{ux56e0ux5f0fux5206ux89e3}

\paragraph{7.3.1 提公因式法}\label{ux63d0ux516cux56e0ux5f0fux6cd5}

\textbf{定理 7.3.1（提取公因式）} 若多项式 \(f(x)\) 的各项都含有公因式
\(d(x)\)，则

\[
f(x) = d(x) \cdot q(x)
\]

其中 \(q(x)\) 由 \(f(x)\) 的各项除以 \(d(x)\) 得到。

\textbf{证明}：设 \(f(x) = \sum a_k x^k\)，每项
\(a_k x^k = d(x) \cdot p_k(x)\)，则由分配律
\(f(x) = d(x) \cdot \sum p_k(x)\)。 \(\blacksquare\)

\paragraph{7.3.2 公式法}\label{ux516cux5f0fux6cd5}

利用7.2节的乘法公式反向进行分解：

\begin{itemize}
\tightlist
\item
  \(a^2 - b^2 = (a+b)(a-b)\)
\item
  \(a^2 + 2ab + b^2 = (a+b)^2\)
\item
  \(a^2 - 2ab + b^2 = (a-b)^2\)
\item
  \(a^3 + b^3 = (a+b)(a^2 - ab + b^2)\)
\item
  \(a^3 - b^3 = (a-b)(a^2 + ab + b^2)\)
\end{itemize}

\paragraph{7.3.3 十字相乘法}\label{ux5341ux5b57ux76f8ux4e58ux6cd5}

\textbf{定理 7.3.2（十字相乘法）} 设
\(f(x) = ax^2 + bx + c\)。若存在整数 \(p, q, r, s\) 使得

\[
pr = a, \quad qs = c, \quad ps + qr = b
\]

则 \(f(x) = (px + q)(rx + s)\)。

\textbf{证明}：

\[
(px + q)(rx + s) = prx^2 + (ps + qr)x + qs = ax^2 + bx + c
\]

\(\blacksquare\)

\paragraph{7.3.4 分组分解法}\label{ux5206ux7ec4ux5206ux89e3ux6cd5}

\textbf{定理 7.3.3（分组分解）} 若多项式 \(f(x)\) 可以写成
\(f(x) = g(x) + h(x)\)，且 \(g(x)\) 和 \(h(x)\) 有公共因式 \(d(x)\)，则

\[
f(x) = d(x) \cdot \left(\frac{g(x)}{d(x)} + \frac{h(x)}{d(x)}\right)
\]

\textbf{证明}：由分配律直接得出。 \(\blacksquare\)

\paragraph{7.3.5 配方法}\label{ux914dux65b9ux6cd5}

\textbf{定理 7.3.4（配方法）} 对任意
\(a \neq 0\)，\(b, c \in \mathbb{R}\)：

\[
ax^2 + bx + c = a\left(x + \frac{b}{2a}\right)^2 + c - \frac{b^2}{4a}
\]

\textbf{证明}：

\[
ax^2 + bx + c = a\left(x^2 + \frac{b}{a}x\right) + c = a\left(x^2 + \frac{b}{a}x + \frac{b^2}{4a^2} - \frac{b^2}{4a^2}\right) + c
\]

\[
= a\left(x + \frac{b}{2a}\right)^2 - \frac{b^2}{4a} + c = a\left(x + \frac{b}{2a}\right)^2 + \frac{4ac - b^2}{4a}
\]

\(\blacksquare\)

\bigskip\hrule\bigskip

\subsubsection{7.4 带余除法}\label{ux5e26ux4f59ux9664ux6cd5-1}

\textbf{定理 7.4.1（带余除法定理）} 设
\(f(x), g(x) \in \mathbb{R}[x]\)，\(g(x) \neq 0\)。则存在唯一的
\(q(x), r(x) \in \mathbb{R}[x]\)，使得

\[
f(x) = g(x) \cdot q(x) + r(x)
\]

其中 \(r(x) = 0\) 或 \(\deg(r) < \deg(g)\)。

\textbf{存在性证明}：对 \(f\) 的次数进行归纳。

\textbf{基础}：若 \(\deg(f) < \deg(g)\)，取
\(q(x) = 0\)，\(r(x) = f(x)\)。

\textbf{归纳}：设 \(\deg(f) = n \geq \deg(g) = m\)。设
\(f(x) = a_n x^n + \cdots\)，\(g(x) = b_m x^m + \cdots\)。

令 \(f_1(x) = f(x) - \dfrac{a_n}{b_m} x^{n-m} g(x)\)。

则 \(\deg(f_1) < n\)（首项消去）。由归纳假设，存在 \(q_1(x), r(x)\) 使得
\(f_1(x) = g(x) \cdot q_1(x) + r(x)\)，其中 \(r(x) = 0\) 或
\(\deg(r) < m\)。

从而
\(f(x) = g(x) \cdot \left(\dfrac{a_n}{b_m} x^{n-m} + q_1(x)\right) + r(x)\)。取
\(q(x) = \dfrac{a_n}{b_m} x^{n-m} + q_1(x)\) 即可。 \(\blacksquare\)

\textbf{唯一性证明}：设

\[
f(x) = g(x) q_1(x) + r_1(x) = g(x) q_2(x) + r_2(x)
\]

则 \(g(x)(q_1(x) - q_2(x)) = r_2(x) - r_1(x)\)。

若 \(q_1 \neq q_2\)，则左端次数 \(\geq \deg(g)\)，而右端 \(r_2 - r_1\)
的次数 \(< \deg(g)\)（或为零多项式），矛盾。故 \(q_1 = q_2\)，从而
\(r_1 = r_2\)。 \(\blacksquare\)

\bigskip\hrule\bigskip

\subsubsection{7.5
余数定理与因式定理}\label{ux4f59ux6570ux5b9aux7406ux4e0eux56e0ux5f0fux5b9aux7406}

\textbf{定理 7.5.1（余数定理）} 多项式 \(f(x)\) 除以 \((x - a)\)
的余数等于 \(f(a)\)。

\textbf{证明}：由带余除法，\(f(x) = (x - a) q(x) + r\)，其中 \(r\)
为常数（因 \(\deg(x-a) = 1\)，\(\deg(r) < 1\)）。

代入 \(x = a\)：\(f(a) = (a - a) q(a) + r = r\)。 \(\blacksquare\)

\textbf{定理 7.5.2（因式定理）} \((x - a)\) 是 \(f(x)\) 的因式，当且仅当
\(f(a) = 0\)。

\textbf{证明}：

(\(\Rightarrow\)) 若 \((x - a) | f(x)\)，则
\(f(x) = (x - a) q(x)\)，\(f(a) = 0 \cdot q(a) = 0\)。

(\(\Leftarrow\)) 若
\(f(a) = 0\)，由余数定理，\(f(x) = (x - a) q(x) + f(a) = (x - a) q(x)\)。
\(\blacksquare\)

\bigskip\hrule\bigskip

\subsubsection{7.6 多项式的根}\label{ux591aux9879ux5f0fux7684ux6839}

\textbf{定义 7.6.1（多项式的根）}

设 \(f(x) \in \mathbb{R}[x]\)。若 \(f(c) = 0\)，则称 \(c\) 为 \(f(x)\)
的一个\textbf{根}（零点）。

\textbf{定理 7.6.1（根的重数与因式分解）} 若 \(c\) 是 \(f(x)\) 的 \(k\)
重根（\(k \geq 1\)），则

\[
f(x) = (x - c)^k \cdot g(x)
\]

其中 \(g(c) \neq 0\)。

\textbf{证明}：对 \(k\) 归纳。

\begin{itemize}
\tightlist
\item
  \textbf{基础}：\(k = 1\) 时，由因式定理
  \(f(x) = (x - c) g(x)\)，\(g(c) \neq 0\)（否则 \(c\) 为至少二重根）。
\item
  \textbf{归纳}：设 \(c\) 为 \(f(x)\) 的 \((k+1)\)
  重根。由归纳假设（\(k\)
  重根的因式分解），\(f(x) = (x - c)^k h(x)\)。由 \((k+1)\)
  重根的定义，\((x-c)^{k+1}\) 整除 \(f(x)\)，故 \(h(c) = 0\)（否则
  \((x-c)^{k+1}\) 不能整除
  \((x-c)^k h(x)\)）。由因式定理，\(h(x) = (x - c) g(x)\)，\(g(c) \neq 0\)。故
  \(f(x) = (x - c)^{k+1} g(x)\)。
\end{itemize}

\(\blacksquare\)

\textbf{定理 7.6.2（\(n\) 次多项式至多有 \(n\) 个根）} 非零的 \(n\)
次多项式 \(f(x)\) 在 \(\mathbb{R}\) 中至多有 \(n\) 个根（计入重数）。

\textbf{证明}：对 \(n\) 归纳。

\begin{itemize}
\tightlist
\item
  \textbf{基础}：\(n = 0\) 时，\(f(x) = a_0 \neq 0\)，无根。
\item
  \textbf{归纳}：设对 \(n - 1\) 次多项式结论成立。若 \(f(x)\)
  无根，则结论成立（\(0 \leq n\)）。若 \(c\) 是 \(f(x)\)
  的根，由因式定理，\(f(x) = (x - c) g(x)\)，\(\deg(g) = n - 1\)。
\end{itemize}

\(f(x)\) 的根恰好是 \(c\) 加上 \(g(x)\) 的根。由归纳假设，\(g(x)\)
至多有 \(n - 1\) 个根。故 \(f(x)\) 至多有 \(n\) 个根。 \(\blacksquare\)

\textbf{推论 7.6.1} 若 \(n\) 次多项式 \(f(x)\) 有超过 \(n\)
个不同的根，则 \(f(x) = 0\)（零多项式）。

\textbf{证明} 对 \(n\) 归纳。\(n = 0\) 时，零次多项式
\(f(x) = c\)（\(c \neq 0\)）无根，与''超过 0 个根''矛盾，故
\(f(x) = 0\)。✓

假设 \(n-1\) 时成立。设 \(f(x)\) 是 \(n\) 次多项式且有超过 \(n\)
个不同的根
\(r_1, r_2, \ldots, r_{n+1}\)。由因式定理，\(f(x) = (x - r_1)g(x)\)，其中
\(g(x)\) 是 \(n-1\) 次多项式。对
\(i = 2, 3, \ldots, n+1\)，\(0 = f(r_i) = (r_i - r_1)g(r_i)\)，由
\(r_i \neq r_1\) 得 \(g(r_i) = 0\)。故 \(g(x)\) 有 \(n\) 个不同的根
\(r_2, \ldots, r_{n+1}\)，由归纳假设 \(g(x) = 0\)，从而
\(f(x) = 0\)。\(\blacksquare\)

\textbf{定理 7.6.3（韦达定理）} 设
\(f(x) = a_n x^n + a_{n-1} x^{n-1} + \cdots + a_0\)（\(a_n \neq 0\)）的
\(n\) 个根为 \(x_1, x_2, \ldots, x_n\)（在复数域中计重数），则：

\[
\begin{cases}
x_1 + x_2 + \cdots + x_n = -\dfrac{a_{n-1}}{a_n} \\[8pt]
\displaystyle\sum_{1 \leq i < j \leq n} x_i x_j = \frac{a_{n-2}}{a_n} \\[8pt]
\cdots \\[4pt]
x_1 x_2 \cdots x_n = (-1)^n \dfrac{a_0}{a_n}
\end{cases}
\]

一般地，设 \(e_k\) 为 \(x_1, \ldots, x_n\) 的 \(k\) 次初等对称多项式：

\[
e_k = \sum_{1 \leq i_1 < \cdots < i_k \leq n} x_{i_1} \cdots x_{i_k}
\]

则 \(a_{n-k} = (-1)^k a_n \cdot e_k\)。

\textbf{证明}：由代数基本定理（定理 8.6.1），\(n\)
次多项式在复数域中恰有 \(n\)
个根（计重数），因此可以完全分解为一次因式的乘积：（完整证明见卷九定理53.1.3（复分析证明）和卷十定理60.14（拓扑学证明））。

\[
f(x) = a_n (x - x_1)(x - x_2) \cdots (x - x_n)
\]

展开右端：

\[
a_n \prod_{i=1}^{n}(x - x_i) = a_n \left(x^n - e_1 x^{n-1} + e_2 x^{n-2} - \cdots + (-1)^n e_n\right)
\]

其中 \(e_k\) 是 \(x_1, \ldots, x_n\) 的 \(k\)
次初等对称多项式。这可以通过对 \(n\) 归纳来严格证明：

\begin{itemize}
\tightlist
\item
  \textbf{基础}：\(n = 1\)
  时，\(a_1(x - x_1) = a_1 x - a_1 x_1\)，\(e_1 = x_1\)，系数为
  \((-1)^1 a_1 e_1 = -a_1 x_1\)。正确。
\item
  \textbf{归纳}：假设
  \(a_n \prod_{i=1}^{n}(x - x_i) = a_n \sum_{k=0}^{n} (-1)^k e_k x^{n-k}\)。则
\end{itemize}

\[
a_n \prod_{i=1}^{n+1}(x - x_i) = a_n (x - x_{n+1}) \sum_{k=0}^{n} (-1)^k e_k x^{n-k}
\]

\[
= a_n \left(\sum_{k=0}^{n} (-1)^k e_k x^{n+1-k} - x_{n+1} \sum_{k=0}^{n} (-1)^k e_k x^{n-k}\right)
\]

\[
= a_n \sum_{k=0}^{n+1} (-1)^k \left(e_k + x_{n+1} e_{k-1}\right) x^{n+1-k}
\]

由对称多项式的递推关系
\(e_k^{(n+1)} = e_k^{(n)} + x_{n+1} e_{k-1}^{(n)}\)，结论成立。

比较系数，\(a_{n-k} = (-1)^k a_n \cdot e_k\)。特别地：

\[
e_1 = \sum x_i = -\frac{a_{n-1}}{a_n}, \quad e_n = \prod x_i = (-1)^n \frac{a_0}{a_n}
\]

\(\blacksquare\)

\begin{quote}
\textbf{本章展望}：多项式的运算和因式分解是代数的核心工具。下一章将利用这些工具系统地研究各类方程的求解理论。
\end{quote}

\bigskip\hrule\bigskip

\subsection{Ch 8
方程与方程组}\label{ch-8-ux65b9ux7a0bux4e0eux65b9ux7a0bux7ec4}

\begin{quote}
\textbf{前置知识}：本章需要第6章的运算律和第7章的多项式理论（特别是因式分解、余数定理和韦达定理）。
\end{quote}

\subsubsection{8.1
方程的基本概念}\label{ux65b9ux7a0bux7684ux57faux672cux6982ux5ff5}

\textbf{定义 8.1.1（方程）}

含有未知数的等式称为\textbf{方程}。使方程成立的未知数的值称为方程的\textbf{解}。

\textbf{定义 8.1.2（解集）}

方程所有解的集合称为\textbf{解集}。

\textbf{定义 8.1.3（同解方程）}

若两个方程有相同的解集，则称它们\textbf{同解}。

\textbf{定义 8.1.4（同解变换）}

以下变换不改变方程的解集：

\begin{enumerate}
\def\labelenumi{\arabic{enumi}.}
\tightlist
\item
  \textbf{等量加减}：方程两边同时加上（或减去）同一个数或整式。
\item
  \textbf{等量乘除}：方程两边同时乘以（或除以）同一个非零数或整式。
\item
  \textbf{等量替换}：将方程中的某表达式替换为与之相等的另一表达式。
\end{enumerate}

\textbf{定理 8.1.1} 设 \(f(x) = g(x)\) 是一个方程，\(h(x)\)
为任意表达式。则 \(f(x) = g(x)\) 与 \(f(x) + h(x) = g(x) + h(x)\) 同解。

\textbf{证明}：设 \(x_0\) 是 \(f(x) = g(x)\) 的解，则
\(f(x_0) = g(x_0)\)，从而 \(f(x_0) + h(x_0) = g(x_0) + h(x_0)\)，即
\(x_0\) 也是变换后方程的解。反之亦然。 \(\blacksquare\)

\bigskip\hrule\bigskip

\subsubsection{8.2
一元一次方程}\label{ux4e00ux5143ux4e00ux6b21ux65b9ux7a0b}

\textbf{定理 8.2.1（一元一次方程解的存在唯一性）} 设
\(a, b \in \mathbb{R}\)，\(a \neq 0\)。则方程

\[
ax + b = 0
\]

有唯一解 \(x = -b/a\)。

\textbf{证明}：

\textbf{存在性}：代入
\(x = -b/a\)，\(a \cdot (-b/a) + b = -b + b = 0\)。故 \(x = -b/a\)
是解。

\textbf{唯一性}：设 \(x_0\) 为任意解，即 \(ax_0 + b = 0\)。则
\(ax_0 = -b\)，由乘法消去律（\(a \neq 0\)），\(x_0 = -b/a\)。
\(\blacksquare\)

\textbf{推论 8.2.1} 若 \(a = 0\)，\(b \neq 0\)，则 \(ax + b = 0\)
无解；若 \(a = 0\)，\(b = 0\)，则 \(ax + b = 0\) 的解为全体实数。

\textbf{证明}：若 \(a = 0\)，方程变为 \(b = 0\)。若
\(b \neq 0\)，矛盾，无解；若 \(b = 0\)，恒等式成立，全体实数为解。
\(\blacksquare\)

\bigskip\hrule\bigskip

\subsubsection{8.3
二元一次方程组}\label{ux4e8cux5143ux4e00ux6b21ux65b9ux7a0bux7ec4}

\textbf{定义 8.3.1（二元一次方程组）}

形如

\[
\begin{cases}
a_1 x + b_1 y = c_1 \\
a_2 x + b_2 y = c_2
\end{cases}
\]

的方程组称为\textbf{二元一次方程组}（线性方程组）。

\textbf{定义 8.3.2（系数行列式）}

\[
D = \begin{vmatrix} a_1 & b_1 \\ a_2 & b_2 \end{vmatrix} = a_1 b_2 - a_2 b_1
\]

\textbf{定理 8.3.1（克拉默法则）} 设
\(D = a_1 b_2 - a_2 b_1 \neq 0\)，则二元一次方程组有唯一解：

\[
x = \frac{D_x}{D} = \frac{c_1 b_2 - c_2 b_1}{a_1 b_2 - a_2 b_1}, \quad y = \frac{D_y}{D} = \frac{a_1 c_2 - a_2 c_1}{a_1 b_2 - a_2 b_1}
\]

其中

\[
D_x = \begin{vmatrix} c_1 & b_1 \\ c_2 & b_2 \end{vmatrix} = c_1 b_2 - c_2 b_1, \quad D_y = \begin{vmatrix} a_1 & c_1 \\ a_2 & c_2 \end{vmatrix} = a_1 c_2 - a_2 c_1
\]

\textbf{证明}：

\textbf{第一步：推导解的表达式。}

将第一个方程乘以 \(b_2\)，第二个方程乘以 \(b_1\)：

\[
\begin{cases}
a_1 b_2 x + b_1 b_2 y = c_1 b_2 \\
a_2 b_1 x + b_1 b_2 y = c_2 b_1
\end{cases}
\]

两式相减：

\[
(a_1 b_2 - a_2 b_1) x = c_1 b_2 - c_2 b_1
\]

即 \(Dx = D_x\)。由 \(D \neq 0\)，\(x = D_x / D\)。

类似地，将第一个方程乘以 \(a_2\)，第二个方程乘以 \(a_1\)：

\[
\begin{cases}
a_1 a_2 x + a_2 b_1 y = a_2 c_1 \\
a_1 a_2 x + a_1 b_2 y = a_1 c_2
\end{cases}
\]

相减：\((a_1 b_2 - a_2 b_1) y = a_1 c_2 - a_2 c_1\)，即
\(Dy = D_y\)，\(y = D_y/D\)。

\textbf{第二步：验证。} 将 \(x = D_x/D\)，\(y = D_y/D\) 代入原方程组：

\[
a_1 \cdot \frac{D_x}{D} + b_1 \cdot \frac{D_y}{D} = \frac{a_1 D_x + b_1 D_y}{D}
\]

计算 \(a_1 D_x + b_1 D_y\)：

\[
a_1(c_1 b_2 - c_2 b_1) + b_1(a_1 c_2 - a_2 c_1) = a_1 c_1 b_2 - a_1 c_2 b_1 + a_1 b_1 c_2 - a_2 b_1 c_1 = c_1(a_1 b_2 - a_2 b_1) = c_1 D
\]

因此
\(a_1 (D_x/D) + b_1 (D_y/D) = c_1 D / D = c_1\)。类似可验证第二个方程。

\textbf{第三步：唯一性。} 设 \((x_0, y_0)\) 为任意解，则

\[
\begin{cases} a_1 x_0 + b_1 y_0 = c_1 \\ a_2 x_0 + b_2 y_0 = c_2 \end{cases}
\]

由第一步的推导过程（这是同解变换），\((x_0, y_0)\) 必须满足
\(D x_0 = D_x\) 和 \(D y_0 = D_y\)。由
\(D \neq 0\)，\(x_0 = D_x/D\)，\(y_0 = D_y/D\)。 \(\blacksquare\)

\textbf{推论 8.3.1（\(D = 0\) 的情形）}

\begin{itemize}
\tightlist
\item
  若 \(D = 0\) 且 \(D_x \neq 0\) 或 \(D_y \neq 0\)，则方程组无解。
\item
  若 \(D = 0\) 且 \(D_x = D_y = 0\)，则方程组有无穷多解。
\end{itemize}

\textbf{证明}：

\textbf{情形一}：\(D = 0\) 且 \(D_x \neq 0\)。假设存在解
\((x_0, y_0)\)，则由定理8.3.1的推导过程（第一步的同解变换），\(D x_0 = D_x\)。但
\(D = 0\) 故 \(0 = D_x \neq 0\)，矛盾。因此无解。\(D_y \neq 0\)
的情况类似。

\textbf{情形二}：\(D = 0\) 且 \(D_x = D_y = 0\)。由
\(D = a_1 b_2 - a_2 b_1 = 0\) 得 \(a_1 b_2 = a_2 b_1\)。

子情形 2a：若 \(b_1 \neq 0\)，由第一个方程
\(x = \frac{c_1 - b_1 y}{a_1}\)（若 \(a_1 \neq 0\)）或
\(y = \frac{c_1}{b_1}\)（若 \(a_1 = 0\)）。代入第二个方程，利用
\(D_x = D_y = 0\) 可验证恒成立，故 \(y\) 可取任意值，\(x\) 由 \(y\)
确定，有无穷多解。

子情形 2b：若 \(b_1 = 0\)，由 \(D = 0\) 得 \(a_1 b_2 = 0\)。若
\(a_1 = 0\)，则第一个方程变为 \(0 = c_1\)，由 \(D_x = c_1 b_2 = 0\) 和
\(D_y = -a_2 c_1 = 0\)，若 \(c_1 \neq 0\) 则
\(b_2 = a_2 = 0\)，但此时方程组退化；若 \(c_1 = 0\)
则第一个方程恒成立。类似讨论第二个方程，可得无穷多解或全体实数解。

综上，\(D = 0\) 时方程组要么无解要么有无穷多解。\(\blacksquare\)

\bigskip\hrule\bigskip

\subsubsection{8.4
一元二次方程}\label{ux4e00ux5143ux4e8cux6b21ux65b9ux7a0b}

\textbf{定理 8.4.1（求根公式）} 设
\(a, b, c \in \mathbb{R}\)，\(a \neq 0\)。方程 \(ax^2 + bx + c = 0\)
的解为

\[
x = \frac{-b \pm \sqrt{b^2 - 4ac}}{2a}
\]

当 \(b^2 - 4ac \geq 0\) 时，解为实数。

\textbf{证明}（配方法）：

\[
ax^2 + bx + c = 0
\]

\[
x^2 + \frac{b}{a}x + \frac{c}{a} = 0
\]

\[
x^2 + \frac{b}{a}x = -\frac{c}{a}
\]

\[
x^2 + \frac{b}{a}x + \frac{b^2}{4a^2} = \frac{b^2}{4a^2} - \frac{c}{a} = \frac{b^2 - 4ac}{4a^2}
\]

\[
\left(x + \frac{b}{2a}\right)^2 = \frac{b^2 - 4ac}{4a^2}
\]

当 \(b^2 - 4ac \geq 0\) 时：

\[
x + \frac{b}{2a} = \pm \frac{\sqrt{b^2 - 4ac}}{2a}
\]

\[
x = \frac{-b \pm \sqrt{b^2 - 4ac}}{2a}
\]

\(\blacksquare\)

\textbf{定义 8.4.1（判别式）}

\(\Delta = b^2 - 4ac\) 称为一元二次方程 \(ax^2 + bx + c = 0\)
的\textbf{判别式}。

\textbf{定理 8.4.2（判别式的三种情形）}

设 \(a \neq 0\)，\(f(x) = ax^2 + bx + c\)：

\textbf{(i)} 若 \(\Delta > 0\)，则 \(f(x) = 0\) 有两个不相等的实根：

\[
x_1 = \frac{-b + \sqrt{\Delta}}{2a}, \quad x_2 = \frac{-b - \sqrt{\Delta}}{2a}
\]

\textbf{(ii)} 若 \(\Delta = 0\)，则 \(f(x) = 0\) 有一个二重实根
\(x_0 = -b/(2a)\)。

\textbf{(iii)} 若 \(\Delta < 0\)，则 \(f(x) = 0\) 在 \(\mathbb{R}\)
中无解。

\textbf{证明}：由求根公式的推导过程，\(\left(x + \frac{b}{2a}\right)^2 = \frac{\Delta}{4a^2}\)。

\begin{itemize}
\tightlist
\item
  \(\Delta > 0\)
  时，右端为正，\(x + \frac{b}{2a} = \pm \frac{\sqrt{\Delta}}{2a}\)
  给出两个不同实数。
\item
  \(\Delta = 0\) 时，\(x + \frac{b}{2a} = 0\)，唯一解 \(x = -b/(2a)\)。
\item
  \(\Delta < 0\) 时，右端为负，而左端 \(\geq 0\)，矛盾。无实数解。
\end{itemize}

\(\blacksquare\)

\textbf{定理 8.4.3（一元二次方程的韦达定理）} 设
\(ax^2 + bx + c = 0\)（\(a \neq 0\)）的两根为 \(x_1, x_2\)，则

\[
x_1 + x_2 = -\frac{b}{a}, \quad x_1 x_2 = \frac{c}{a}
\]

\textbf{证明}：由求根公式：

\[
x_1 + x_2 = \frac{-b + \sqrt{\Delta}}{2a} + \frac{-b - \sqrt{\Delta}}{2a} = \frac{-2b}{2a} = -\frac{b}{a}
\]

\[
x_1 x_2 = \frac{(-b + \sqrt{\Delta})(-b - \sqrt{\Delta})}{4a^2} = \frac{b^2 - \Delta}{4a^2} = \frac{b^2 - (b^2 - 4ac)}{4a^2} = \frac{4ac}{4a^2} = \frac{c}{a}
\]

\(\blacksquare\)

\bigskip\hrule\bigskip

\subsubsection{8.5
分式方程与无理方程}\label{ux5206ux5f0fux65b9ux7a0bux4e0eux65e0ux7406ux65b9ux7a0b}

\textbf{定义 8.5.1（分式方程）}

分母中含有未知数的方程称为\textbf{分式方程}。

\textbf{定义 8.5.2（增根）}

在方程变形过程中产生的、不属于原方程解集的根称为\textbf{增根}。

\textbf{定理 8.5.1（增根的产生条件）} 将分式方程
\(\dfrac{P(x)}{Q(x)} = \dfrac{R(x)}{S(x)}\) 两边乘以 \(Q(x) S(x)\)
去分母，所得整式方程的根中，使 \(Q(x) S(x) = 0\) 的值为增根。

\textbf{证明}：设 \(x_0\) 使得 \(Q(x_0) = 0\)。则原方程在 \(x_0\)
处无定义。但乘以 \(Q(x_0) S(x_0) = 0\) 后等式 \(0 = 0\) 恒成立，故
\(x_0\) 可能满足去分母后的方程，却不满足原方程。 \(\blacksquare\)

\textbf{定义 8.5.3（无理方程）}

根号下含有未知数的方程称为\textbf{无理方程}。

\textbf{定理 8.5.2} 两边平方可能产生增根。设 \(A(x) = B(x)\)
为无理方程，两边平方得 \(A^2(x) = B^2(x)\)，则 \(A(x) = B(x)\) 或
\(A(x) = -B(x)\)。因此平方后方程的解集包含原方程的解集，可能更大。

\textbf{证明}：\(A^2 = B^2 \Leftrightarrow A^2 - B^2 = 0 \Leftrightarrow (A+B)(A-B) = 0 \Leftrightarrow A = B\)
或 \(A = -B\)。 \(\blacksquare\)

\bigskip\hrule\bigskip

\subsubsection{8.6
高次方程简介}\label{ux9ad8ux6b21ux65b9ux7a0bux7b80ux4ecb}

\textbf{定理 8.6.1（代数基本定理，陈述）} 任何 \(n \geq 1\)
次复系数多项式在 \(\mathbb{C}\) 中恰有 \(n\) 个根（计重数）。

\textbf{定理 8.6.2（阿贝尔-鲁菲尼定理，陈述）} 当 \(n \geq 5\) 时，一般
\(n\) 次多项式方程

\[
a_n x^n + a_{n-1} x^{n-1} + \cdots + a_0 = 0
\]

不存在用系数的加、减、乘、除和开方运算构成的一般求根公式。

此定理的证明超出本书范围，但其含义是：五次及以上的高次方程不能像一至四次方程那样用根式给出通解。

\begin{quote}
\textbf{本章展望}：方程是代数的核心应用。下一章将研究不等式，它是方程理论的自然推广，也是优化问题的基础。
\end{quote}

\bigskip\hrule\bigskip

\subsection{Ch 9 不等式}\label{ch-9-ux4e0dux7b49ux5f0f}

\begin{quote}
\textbf{前置知识}：本章需要第4章实数的序公理，以及第6章的运算律和第7章的多项式理论。
\end{quote}

\subsubsection{9.1 基本性质}\label{ux57faux672cux6027ux8d28}

以下从实数的序公理（参见第一卷第4章）出发。

\textbf{定理 9.1.1（传递性）} 若 \(a > b\) 且 \(b > c\)，则 \(a > c\)。

\textbf{证明}：由序公理，\(a > b\) 意味着 \(a - b > 0\)，\(b > c\)
意味着 \(b - c > 0\)。正数之和仍为正（由序公理），故

\[
a - c = (a - b) + (b - c) > 0
\]

即 \(a > c\)。 \(\blacksquare\)

\textbf{定理 9.1.2（加法保序性）} 若 \(a > b\)，则对任意
\(c \in \mathbb{R}\)，有 \(a + c > b + c\)。

\textbf{证明}：\(a > b \Rightarrow a - b > 0\)。而
\((a + c) - (b + c) = a - b > 0\)，故 \(a + c > b + c\)。
\(\blacksquare\)

\textbf{定理 9.1.3（乘法保序性------正数乘）} 若 \(a > b\) 且
\(c > 0\)，则 \(ac > bc\)。

\textbf{证明}：\(a > b \Rightarrow a - b > 0\)。\(c > 0\)，正数之积为正，故
\(c(a - b) > 0\)，即 \(ac - bc > 0\)，\(ac > bc\)。 \(\blacksquare\)

\textbf{定理 9.1.4（乘法反序性------负数乘）} 若 \(a > b\) 且
\(c < 0\)，则 \(ac < bc\)。

\textbf{证明}：\(a > b \Rightarrow a - b > 0\)。\(c < 0 \Rightarrow -c > 0\)。正数之积为正，故
\((-c)(a - b) > 0\)，即
\(-(ac - bc) > 0\)，\(ac - bc < 0\)，\(ac < bc\)。 \(\blacksquare\)

\textbf{定理 9.1.5（正数的平方为正）} 若 \(a \neq 0\)，则 \(a^2 > 0\)。

\textbf{证明}：若 \(a > 0\)，由正数之积为正，\(a^2 = a \cdot a > 0\)。若
\(a < 0\)，则 \(-a > 0\)，\(a^2 = (-a)^2 = (-a)(-a) > 0\)。
\(\blacksquare\)

\textbf{定理 9.1.6} \(1 > 0\)。

\textbf{证明}：\(1 = 1^2 > 0\)（由定理9.1.5，\(1 \neq 0\)）。
\(\blacksquare\)

\textbf{定理 9.1.7（三歧性）} 对任意
\(a, b \in \mathbb{R}\)，以下三种情形恰有一种成立：\(a > b\)，\(a = b\)，\(a < b\)。

\textbf{证明}：由实数序公理的三歧性公理直接得出。 \(\blacksquare\)

\textbf{定理 9.1.8（倒数保序性）} 若 \(a > b > 0\)，则
\(\dfrac{1}{a} < \dfrac{1}{b}\)。

\textbf{证明}：\(a > b > 0\)，则 \(ab > 0\)。由定理9.1.3，两边乘以
\(\dfrac{1}{ab} > 0\)：

\[
\frac{a}{ab} > \frac{b}{ab} \Rightarrow \frac{1}{b} > \frac{1}{a}
\]

\(\blacksquare\)

\bigskip\hrule\bigskip

\subsubsection{9.2
一元一次不等式组}\label{ux4e00ux5143ux4e00ux6b21ux4e0dux7b49ux5f0fux7ec4}

\textbf{定理 9.2.1} 设 \(a \neq 0\)，则 \(ax + b > 0\) 的解集为：

\begin{itemize}
\tightlist
\item
  若 \(a > 0\)：\(x > -b/a\)；
\item
  若 \(a < 0\)：\(x < -b/a\)。
\end{itemize}

\textbf{证明}：\(ax > -b\)。

\begin{itemize}
\tightlist
\item
  若 \(a > 0\)：由乘法保序性，\(x > -b/a\)。
\item
  若 \(a < 0\)：由乘法反序性，\(x < -b/a\)。
\end{itemize}

\(\blacksquare\)

\textbf{定理 9.2.2（不等式组的解集）} 不等式组

\[
\begin{cases} a_1 x + b_1 > 0 \\ a_2 x + b_2 > 0 \end{cases}
\]

的解集为两个不等式解集的交集。

\textbf{证明}：\(x\) 满足不等式组当且仅当 \(x\)
同时满足每个不等式，这恰好是两个解集的交集的定义。 \(\blacksquare\)

\bigskip\hrule\bigskip

\subsubsection{9.3
一元二次不等式}\label{ux4e00ux5143ux4e8cux6b21ux4e0dux7b49ux5f0f}

\textbf{定理 9.3.1} 设
\(a > 0\)，\(f(x) = ax^2 + bx + c\)，\(\Delta = b^2 - 4ac\)。

\textbf{(i) \(\Delta > 0\)}：设两根 \(x_1 < x_2\)。则

\begin{itemize}
\tightlist
\item
  \(f(x) > 0\) 的解集为 \((-\infty, x_1) \cup (x_2, +\infty)\)。
\item
  \(f(x) < 0\) 的解集为 \((x_1, x_2)\)。
\end{itemize}

\textbf{(ii) \(\Delta = 0\)}：设根为 \(x_0\)。则

\begin{itemize}
\tightlist
\item
  \(f(x) > 0\) 的解集为 \(\mathbb{R} \setminus \{x_0\}\)。
\item
  \(f(x) < 0\) 的解集为 \(\varnothing\)。
\end{itemize}

\textbf{(iii) \(\Delta < 0\)}：则

\begin{itemize}
\tightlist
\item
  \(f(x) > 0\) 的解集为 \(\mathbb{R}\)。
\item
  \(f(x) < 0\) 的解集为 \(\varnothing\)。
\end{itemize}

\textbf{证明}：由配方法，\(f(x) = a\left(x + \dfrac{b}{2a}\right)^2 + \dfrac{4ac - b^2}{4a} = a\left(x + \dfrac{b}{2a}\right)^2 - \dfrac{\Delta}{4a}\)。

\textbf{(i)} \(\Delta > 0\)：\(f(x) = a(x - x_1)(x - x_2)\)。

\begin{itemize}
\tightlist
\item
  \(x < x_1\)
  时：\(x - x_1 < 0\)，\(x - x_2 < 0\)，\((x-x_1)(x-x_2) > 0\)，\(f(x) > 0\)。
\item
  \(x_1 < x < x_2\)
  时：\(x - x_1 > 0\)，\(x - x_2 < 0\)，\((x-x_1)(x-x_2) < 0\)，\(f(x) < 0\)。
\item
  \(x > x_2\)
  时：\(x - x_1 > 0\)，\(x - x_2 > 0\)，\((x-x_1)(x-x_2) > 0\)，\(f(x) > 0\)。
\item
  \(x = x_1\) 或 \(x = x_2\) 时：\(f(x) = 0\)。
\end{itemize}

\textbf{(ii)}
\(\Delta = 0\)：\(f(x) = a(x - x_0)^2 \geq 0\)。等号当且仅当
\(x = x_0\)。

\textbf{(iii)}
\(\Delta < 0\)：\(f(x) = a\left(x + \dfrac{b}{2a}\right)^2 - \dfrac{\Delta}{4a}\)。其中
\(-\Delta/(4a) > 0\)（因 \(\Delta < 0\)，\(a > 0\)），而第一项
\(\geq 0\)，故 \(f(x) > 0\) 对一切 \(x\) 成立。

\(\blacksquare\)

\textbf{推论 9.3.1} 若 \(a < 0\)，不等式 \(ax^2 + bx + c > 0\) 等价于
\(-ax^2 - bx - c < 0\)，可转化为 \(a' > 0\) 的情形。

\textbf{证明}：\(ax^2 + bx + c > 0\) 两边乘 \(-1\)（不等号反向）得
\(-ax^2 - bx - c < 0\)，其中 \(-a > 0\)。 \(\blacksquare\)

\bigskip\hrule\bigskip

\subsubsection{9.4
绝对值不等式}\label{ux7eddux5bf9ux503cux4e0dux7b49ux5f0f}

\textbf{定义 9.4.1（绝对值）}

对任意 \(a \in \mathbb{R}\)，定义

\[
|a| = \begin{cases} a & \text{若 } a \geq 0 \\ -a & \text{若 } a < 0 \end{cases}
\]

\textbf{定理 9.4.1（绝对值的基本性质）}

\begin{enumerate}
\def\labelenumi{(\roman{enumi})}
\item
  \(|a| \geq 0\)，等号当且仅当 \(a = 0\)。
\item
  \(|a| = \sqrt{a^2}\)。
\item
  \(|ab| = |a| \cdot |b|\)。
\item
  \(\left|\dfrac{a}{b}\right| = \dfrac{|a|}{|b|}\)（\(b \neq 0\)）。
\end{enumerate}

\textbf{证明}：

\begin{enumerate}
\def\labelenumi{(\roman{enumi})}
\item
  若 \(a \geq 0\)，\(|a| = a \geq 0\)；若
  \(a < 0\)，\(|a| = -a > 0\)。\(|a| = 0 \Leftrightarrow a = 0\)
  由定义直接得出。
\item
  若 \(a \geq 0\)，\(\sqrt{a^2} = a = |a|\)；若
  \(a < 0\)，\(\sqrt{a^2} = \sqrt{(-a)^2} = -a = |a|\)。
\item
  分四种子情形讨论：
\end{enumerate}

\begin{itemize}
\tightlist
\item
  \(a \geq 0, b \geq 0\)：\(ab \geq 0\)，\(|ab| = ab = |a| \cdot |b|\)。
\item
  \(a \geq 0, b < 0\)：\(ab \leq 0\)，\(|ab| = -ab = a \cdot (-b) = |a| \cdot |b|\)。
\item
  \(a < 0, b \geq 0\)：\(ab \leq 0\)，\(|ab| = -ab = (-a) \cdot b = |a| \cdot |b|\)。
\item
  \(a < 0, b < 0\)：\(ab > 0\)，\(|ab| = ab = (-a)(-b) = |a| \cdot |b|\)。
\end{itemize}

所有情形下 \(|ab| = |a| \cdot |b|\) 成立。

\begin{enumerate}
\def\labelenumi{(\roman{enumi})}
\setcounter{enumi}{3}
\tightlist
\item
  由 (iii)，\(|b| \cdot |a/b| = |b \cdot a/b| = |a|\)，故
  \(|a/b| = |a|/|b|\)。
\end{enumerate}

\(\blacksquare\)

\textbf{定理 9.4.2（三角不等式）} 对任意 \(a, b \in \mathbb{R}\)：

\[
|a + b| \leq |a| + |b|
\]

\textbf{证明}：两边平方（因两边非负，不等式方向不变）：

\[
(a + b)^2 \leq (|a| + |b|)^2
\]

\[
a^2 + 2ab + b^2 \leq |a|^2 + 2|a||b| + |b|^2 = a^2 + 2|ab| + b^2
\]

只需证 \(ab \leq |ab|\)，即 \(ab \leq |ab|\)。

若 \(ab \geq 0\)，则 \(ab = |ab|\)，成立。若 \(ab < 0\)，则
\(ab < 0 \leq |ab|\)，成立。

\(\blacksquare\)

\textbf{定理 9.4.3（反向三角不等式）} 对任意 \(a, b \in \mathbb{R}\)：

\[
\bigl||a| - |b|\bigr| \leq |a - b|
\]

\textbf{证明}：由三角不等式，\(|a| = |(a - b) + b| \leq |a - b| + |b|\)，故
\(|a| - |b| \leq |a - b|\)。

类似地，\(|b| = |(b - a) + a| \leq |b - a| + |a| = |a - b| + |a|\)，故
\(|b| - |a| \leq |a - b|\)，即 \(-(|a| - |b|) \leq |a - b|\)。

综合两式：\(-|a - b| \leq |a| - |b| \leq |a - b|\)，即
\(\bigl||a| - |b|\bigr| \leq |a - b|\)。 \(\blacksquare\)

\textbf{推论 9.4.1} 对任意 \(a, b \in \mathbb{R}\)：

\[
\bigl||a| - |b|\bigr| \leq |a + b| \leq |a| + |b|
\]

\textbf{证明}：左边不等式由 \(|a| = |(a+b) - b| \leq |a+b| + |b|\)，得
\(|a| - |b| \leq |a+b|\)。类似
\(|b| - |a| \leq |a+b|\)。右边不等式即三角不等式。 \(\blacksquare\)

\bigskip\hrule\bigskip

\subsubsection{9.5 均值不等式}\label{ux5747ux503cux4e0dux7b49ux5f0f}

\textbf{定理 9.5.1（二元算术-几何均值不等式，AM-GM）} 对任意
\(a, b \geq 0\)：

\[
\frac{a + b}{2} \geq \sqrt{ab}
\]

等号当且仅当 \(a = b\)。

\textbf{证明}（配方法）：

\[
\frac{a + b}{2} - \sqrt{ab} = \frac{a + b - 2\sqrt{ab}}{2} = \frac{(\sqrt{a})^2 - 2\sqrt{a}\sqrt{b} + (\sqrt{b})^2}{2} = \frac{(\sqrt{a} - \sqrt{b})^2}{2} \geq 0
\]

等号当且仅当 \(\sqrt{a} = \sqrt{b}\)，即 \(a = b\)。 \(\blacksquare\)

\textbf{定理 9.5.2（\(n\) 元 AM-GM 不等式）} 对任意
\(a_1, a_2, \ldots, a_n \geq 0\)（\(n \geq 2\)）：

\[
\frac{a_1 + a_2 + \cdots + a_n}{n} \geq \sqrt[n]{a_1 a_2 \cdots a_n}
\]

等号当且仅当 \(a_1 = a_2 = \cdots = a_n\)。

\textbf{证明}（前向-后向归纳法，柯西归纳法）：

\textbf{第一步（\(n = 2\) 的情形）}：由定理9.5.1。

\textbf{第二步（\(n \to 2n\)）}：设 \(n\) 元不等式成立，证 \(2n\) 元。

\[
\frac{a_1 + \cdots + a_{2n}}{2n} = \frac{1}{2}\left(\frac{a_1 + \cdots + a_n}{n} + \frac{a_{n+1} + \cdots + a_{2n}}{n}\right)
\]

由归纳假设和二元AM-GM：

\[
\geq \frac{1}{2}\left(\sqrt[n]{a_1 \cdots a_n} + \sqrt[n]{a_{n+1} \cdots a_{2n}}\right) \geq \sqrt{\sqrt[n]{a_1 \cdots a_n} \cdot \sqrt[n]{a_{n+1} \cdots a_{2n}}} = \sqrt[2n]{a_1 \cdots a_{2n}}
\]

\textbf{第三步（\(n \to n-1\)，向后归纳）}：设 \(n\) 元不等式成立，证
\(n-1\) 元。

给定 \(a_1, \ldots, a_{n-1} \geq 0\)，令

\[
a_n = \frac{a_1 + a_2 + \cdots + a_{n-1}}{n-1} = A
\]

由 \(n\) 元AM-GM：

\[
\frac{a_1 + \cdots + a_{n-1} + A}{n} \geq \sqrt[n]{a_1 \cdots a_{n-1} \cdot A}
\]

左端 \(= \frac{(n-1)A + A}{n} = A\)。故

\[
A \geq \sqrt[n]{a_1 \cdots a_{n-1} \cdot A}
\]

\[
A^n \geq a_1 \cdots a_{n-1} \cdot A
\]

若 \(A = 0\)，则所有 \(a_i = 0\)，\(n-1\) 元不等式 \(0 \geq 0\)
平凡成立。若 \(A > 0\)，两边除以 \(A\)：

\[
A \geq \sqrt[n-1]{a_1 \cdots a_{n-1}}
\]

即
\(\dfrac{a_1 + \cdots + a_{n-1}}{n-1} \geq \sqrt[n-1]{a_1 \cdots a_{n-1}}\)。

综合第二步和第三步，对所有 \(n \geq 2\)
不等式成立。等号条件由各步等号传递得出：第二步中 \(2n\) 元等号当且仅当
\(a_1 = \cdots = a_{2n}\)（由二元 AM-GM 及归纳假设），第三步中 \(n-1\)
元等号当且仅当
\(a_1 = \cdots = a_{n-1} = A = \frac{a_1+\cdots+a_{n-1}}{n-1}\)，即全部相等。因此对所有
\(n \geq 2\)，等号当且仅当 \(a_1 = a_2 = \cdots = a_n\)。
\(\blacksquare\)

\textbf{定义 9.5.1（均值的定义）}

设 \(a_1, a_2, \ldots, a_n > 0\)，定义以下均值：

\begin{itemize}
\tightlist
\item
  \textbf{算术平均}（AM）：\(A = \dfrac{a_1 + a_2 + \cdots + a_n}{n}\)
\item
  \textbf{几何平均}（GM）：\(G = \sqrt[n]{a_1 a_2 \cdots a_n}\)
\item
  \textbf{调和平均}（HM）：\(H = \dfrac{n}{\dfrac{1}{a_1} + \dfrac{1}{a_2} + \cdots + \dfrac{1}{a_n}}\)
\item
  \textbf{平方平均}（QM）：\(Q = \sqrt{\dfrac{a_1^2 + a_2^2 + \cdots + a_n^2}{n}}\)
\end{itemize}

\textbf{定理 9.5.3（均值不等式链）} 对任意
\(a_1, a_2, \ldots, a_n > 0\)：

\[
H \leq G \leq A \leq Q
\]

等号当且仅当 \(a_1 = a_2 = \cdots = a_n\)。

\textbf{证明}：

\textbf{\(G \leq A\)}：即定理9.5.2。

\textbf{\(H \leq G\)}：对 \(1/a_1, 1/a_2, \ldots, 1/a_n\) 应用
\(G \leq A\)：

\[
\sqrt[n]{\frac{1}{a_1} \cdot \frac{1}{a_2} \cdots \frac{1}{a_n}} \leq \frac{\frac{1}{a_1} + \frac{1}{a_2} + \cdots + \frac{1}{a_n}}{n}
\]

取倒数（函数 \(f(x) = 1/x\) 在 \(x > 0\) 上单调递减）：

\[
\frac{n}{\frac{1}{a_1} + \cdots + \frac{1}{a_n}} \leq \sqrt[n]{a_1 a_2 \cdots a_n}
\]

即 \(H \leq G\)。

\textbf{\(A \leq Q\)}：需证
\(\left(\dfrac{\sum a_i}{n}\right)^2 \leq \dfrac{\sum a_i^2}{n}\)，即

\[
\left(\sum a_i\right)^2 \leq n \sum a_i^2
\]

将左端展开：

\[
\left(\sum_{i=1}^n a_i\right)^2 = \sum_{i=1}^n a_i^2 + 2\sum_{1 \leq i < j \leq n} a_i a_j
\]

因此需证 \(2\sum_{i < j} a_i a_j \leq (n-1)\sum a_i^2\)。由基本不等式
\(2a_i a_j \leq a_i^2 + a_j^2\)（这等价于
\((a_i - a_j)^2 \geq 0\)），对每对 \((i, j)\) 求和：

\[
2\sum_{i < j} a_i a_j \leq \sum_{i < j} (a_i^2 + a_j^2) = (n-1)\sum_{i=1}^n a_i^2
\]

（每个 \(a_i^2\) 恰好出现在 \(n-1\) 对中。）代入即得
\(\left(\sum a_i\right)^2 \leq n \sum a_i^2\)。\(\blacksquare\)

\textbf{注记}：上述不等式也是柯西-施瓦茨不等式（定理 9.6.1）取
\(b_i = 1\) 的特例。

\bigskip\hrule\bigskip

\subsubsection{9.6
柯西-施瓦茨不等式}\label{ux67efux897f-ux65bdux74e6ux8328ux4e0dux7b49ux5f0f}

\textbf{定理 9.6.1（柯西-施瓦茨不等式）} 对任意实数
\(a_1, \ldots, a_n, b_1, \ldots, b_n\)：

\[
\left(\sum_{i=1}^n a_i b_i\right)^2 \leq \left(\sum_{i=1}^n a_i^2\right)\left(\sum_{i=1}^n b_i^2\right)
\]

等号当且仅当存在常数 \(t\)，使得 \(a_i = t b_i\)（\(\forall i\)）或
\(b_i = 0\)（\(\forall i\)）。

\textbf{证明}（判别式法）：

考虑关于 \(t\) 的函数：

\[
f(t) = \sum_{i=1}^n (a_i t - b_i)^2 \geq 0 \quad (\forall t \in \mathbb{R})
\]

展开：

\[
f(t) = \left(\sum a_i^2\right) t^2 - 2\left(\sum a_i b_i\right) t + \left(\sum b_i^2\right) \geq 0
\]

设 \(A = \sum a_i^2\)，\(B = \sum a_i b_i\)，\(C = \sum b_i^2\)。则
\(f(t) = At^2 - 2Bt + C \geq 0\) 对一切 \(t\) 成立。

若 \(A = 0\)，则所有 \(a_i = 0\)，不等式两端均为零，等号成立。

若 \(A > 0\)，二次函数 \(f(t) \geq 0\) 对一切 \(t\)
成立，当且仅当其判别式 \(\leq 0\)：

\[
\Delta = 4B^2 - 4AC \leq 0
\]

\[
B^2 \leq AC
\]

即

\[
\left(\sum a_i b_i\right)^2 \leq \left(\sum a_i^2\right)\left(\sum b_i^2\right)
\]

等号当且仅当 \(\Delta = 0\)，即存在 \(t_0\) 使得 \(f(t_0) = 0\)，即
\(\sum(a_i t_0 - b_i)^2 = 0\)，即 \(a_i t_0 = b_i\) 对所有 \(i\) 成立。
\(\blacksquare\)

\textbf{推论 9.6.1（施瓦茨不等式的积分形式，陈述）} 设 \(f, g\) 在
\([a, b]\) 上可积，则

\[
\left(\int_a^b f(x) g(x) \, dx\right)^2 \leq \left(\int_a^b f^2(x) \, dx\right)\left(\int_a^b g^2(x) \, dx\right)
\]

\begin{quote}
\textbf{本章展望}：不等式是分析学和优化理论的基石。第10章将从逻辑学的角度审视数学命题的结构，为严谨的数学推理提供更深层的理论支撑。
\end{quote}

\bigskip\hrule\bigskip

\subsection{Ch 10
逻辑与命题}\label{ch-10-ux903bux8f91ux4e0eux547dux9898}

\begin{quote}
\textbf{前置知识}：本章不需要特定的代数知识，但要求读者具备基本的数学推理能力，以及对集合论（第一卷第2章）的了解。
\end{quote}

\subsubsection{10.1
充分条件与必要条件}\label{ux5145ux5206ux6761ux4ef6ux4e0eux5fc5ux8981ux6761ux4ef6-1}

\textbf{定义 10.1.1（命题）}

可以判断真假的陈述句称为\textbf{命题}。真命题记为''真''（T），假命题记为''假''（F）。

\textbf{定义 10.1.2（蕴含关系）}

设 \(P\)，\(Q\) 为命题。``\(P \Rightarrow Q\)''（若 \(P\) 则
\(Q\)）称为\textbf{蕴含}。其真值表为：

{\def\LTcaptype{none} % do not increment counter
\begin{longtable}[]{@{}ccc@{}}
\toprule\noalign{}
\(P\) & \(Q\) & \(P \Rightarrow Q\) \\
\midrule\noalign{}
\endhead
\bottomrule\noalign{}
\endlastfoot
T & T & T \\
T & F & F \\
F & T & T \\
F & F & T \\
\end{longtable}
}

\textbf{定义 10.1.3（充分条件与必要条件）}

在 \(P \Rightarrow Q\) 中：

\begin{itemize}
\tightlist
\item
  \(P\) 是 \(Q\) 的\textbf{充分条件}（\(P\) 足以推出 \(Q\)）。
\item
  \(Q\) 是 \(P\) 的\textbf{必要条件}（\(Q\) 是 \(P\) 成立所必须的）。
\end{itemize}

\textbf{定义 10.1.4（充要条件）}

若 \(P \Rightarrow Q\) 且 \(Q \Rightarrow P\)，即
\(P \Leftrightarrow Q\)，则称 \(P\) 是 \(Q\) 的\textbf{充要条件}。

\textbf{定理 10.1.1（蕴含的集合视角）} 设 \(P(x)\) 和 \(Q(x)\)
分别表示''\(x\) 具有性质 \(P\)``和''\(x\) 具有性质 \(Q\)``。令

\[
A = \{x : P(x)\}, \quad B = \{x : Q(x)\}
\]

则：

\begin{enumerate}
\def\labelenumi{(\roman{enumi})}
\item
  \(P \Rightarrow Q\) 当且仅当 \(A \subseteq B\)。
\item
  \(P \Leftrightarrow Q\) 当且仅当 \(A = B\)。
\end{enumerate}

\textbf{证明}：

\begin{enumerate}
\def\labelenumi{(\roman{enumi})}
\tightlist
\item
  (\(\Rightarrow\)) 若 \(P \Rightarrow Q\)，则对任意
  \(x \in A\)，\(P(x)\) 成立，从而 \(Q(x)\) 成立，即 \(x \in B\)。故
  \(A \subseteq B\)。
\end{enumerate}

(\(\Leftarrow\)) 若 \(A \subseteq B\)，则对任意
\(x\)，\(P(x) \Rightarrow x \in A \Rightarrow x \in B \Rightarrow Q(x)\)。故
\(P \Rightarrow Q\)。

\begin{enumerate}
\def\labelenumi{(\roman{enumi})}
\setcounter{enumi}{1}
\tightlist
\item
  由 (i)，\(P \Leftrightarrow Q\) 即 \(P \Rightarrow Q\) 且
  \(Q \Rightarrow P\)，等价于 \(A \subseteq B\) 且 \(B \subseteq A\)，即
  \(A = B\)。 \(\blacksquare\)
\end{enumerate}

\textbf{定理 10.1.2（逆否命题）} \(P \Rightarrow Q\) 与其逆否命题
\(\neg Q \Rightarrow \neg P\) 等价。

\textbf{证明}：利用真值表。\(\neg Q \Rightarrow \neg P\) 与
\(P \Rightarrow Q\) 在所有真值指派下具有相同的真值。

或者直接证明：假设 \(P \Rightarrow Q\) 成立且 \(\neg Q\) 成立。若
\(\neg P\) 不成立，即 \(P\) 成立，则由 \(P \Rightarrow Q\) 得 \(Q\)
成立，与 \(\neg Q\) 矛盾。故 \(\neg P\) 成立，即
\(\neg Q \Rightarrow \neg P\)。

反方向的证明完全对称。 \(\blacksquare\)

\bigskip\hrule\bigskip

\subsubsection{10.2
全称命题与存在命题的否定}\label{ux5168ux79f0ux547dux9898ux4e0eux5b58ux5728ux547dux9898ux7684ux5426ux5b9a}

\textbf{定义 10.2.1（全称命题与存在命题）}

\begin{itemize}
\tightlist
\item
  \textbf{全称命题}：\(\forall x \in D, \, P(x)\)（对 \(D\) 中所有
  \(x\)，\(P(x)\) 成立）。
\item
  \textbf{存在命题}：\(\exists x \in D, \, P(x)\)（存在 \(D\) 中的
  \(x\)，使得 \(P(x)\) 成立）。
\end{itemize}

\textbf{定理 10.2.1（全称命题的否定）}

\[
\neg(\forall x \, P(x)) \Leftrightarrow \exists x \, \neg P(x)
\]

\textbf{证明}：

(\(\Rightarrow\)) 假设 \(\neg(\forall x \, P(x))\)。若不存在 \(x\) 使得
\(\neg P(x)\)，即对所有 \(x\)，\(P(x)\) 成立，则 \(\forall x \, P(x)\)
成立，与假设矛盾。故存在 \(x\) 使得 \(\neg P(x)\)。

(\(\Leftarrow\)) 假设 \(\exists x \, \neg P(x)\)，设此 \(x\) 为
\(x_0\)，则 \(P(x_0)\) 为假，故 \(\forall x \, P(x)\) 不成立，即
\(\neg(\forall x \, P(x))\)。

\(\blacksquare\)

\textbf{定理 10.2.2（存在命题的否定）}

\[
\neg(\exists x \, P(x)) \Leftrightarrow \forall x \, \neg P(x)
\]

\textbf{证明}：

(\(\Rightarrow\)) 假设 \(\neg(\exists x \, P(x))\)。若存在 \(x_0\) 使得
\(P(x_0)\) 成立，则 \(\exists x \, P(x)\) 成立，与假设矛盾。故对所有
\(x\)，\(\neg P(x)\) 成立。

(\(\Leftarrow\)) 假设 \(\forall x \, \neg P(x)\)。若
\(\exists x \, P(x)\) 成立，设 \(P(x_0)\) 成立，则 \(\neg P(x_0)\) 与
\(\forall x \, \neg P(x)\) 矛盾。故 \(\neg(\exists x \, P(x))\)。

\(\blacksquare\)

\textbf{推论 10.2.1（双重否定）}

\[
\neg(\neg P) \Leftrightarrow P
\]

\textbf{证明}：由真值表直接验证。 \(\blacksquare\)

\textbf{推论 10.2.2（复合命题的否定）}

\[
\neg(P \land Q) \Leftrightarrow (\neg P) \lor (\neg Q)
\]

\[
\neg(P \lor Q) \Leftrightarrow (\neg P) \land (\neg Q)
\]

这些法则称为\textbf{德摩根律}。

\textbf{证明}：由真值表直接验证（此处略去具体真值表）。 \(\blacksquare\)

\begin{quote}
\textbf{本章展望}：逻辑与命题的严格理论为数学证明提供了语言基础。第11章将运用不等式和逻辑工具研究线性规划这一应用数学的重要分支。
\end{quote}

\bigskip\hrule\bigskip

\subsection{Ch 11
线性规划初步}\label{ch-11-ux7ebfux6027ux89c4ux5212ux521dux6b65}

\begin{quote}
\textbf{前置知识}：本章需要第9章的不等式理论（特别是二元一次不等式），以及基本的平面直角坐标系知识。
\end{quote}

\subsubsection{11.1
二元一次不等式表示的区域}\label{ux4e8cux5143ux4e00ux6b21ux4e0dux7b49ux5f0fux8868ux793aux7684ux533aux57df}

\textbf{定理 11.1.1} 设直线 \(l: ax + by + c = 0\)（\(a, b\)
不同时为零）。则

\begin{enumerate}
\def\labelenumi{(\roman{enumi})}
\item
  \(ax + by + c > 0\) 表示 \(l\) 的某一侧的半平面。
\item
  \(ax + by + c < 0\) 表示 \(l\) 的另一侧的半平面。
\end{enumerate}

\textbf{证明}（取特殊点法）：

设 \((x_0, y_0)\) 为不在 \(l\) 上的任意一点，即
\(ax_0 + by_0 + c \neq 0\)。

设 \((x, y)\) 为 \(l\) 上的点，即 \(ax + by + c = 0\)。

考虑线段上任一点
\((x_t, y_t) = (1-t)(x_0, y_0) + t(x, y)\)（\(0 \leq t \leq 1\)）：

\[
ax_t + by_t + c = (1-t)(ax_0 + by_0 + c) + t(ax + by + c) = (1-t)(ax_0 + by_0 + c)
\]

当 \(0 \leq t < 1\) 时，\((1-t) > 0\)，故 \(ax_t + by_t + c\) 与
\(ax_0 + by_0 + c\) 同号。

这说明：\((x_0, y_0)\) 和连接它与 \(l\) 上任意点的线段上的点（不含端点在
\(l\) 上的），使得 \(ax + by + c\) 保持同号。因此，不在 \(l\)
上的同一连通区域（半平面）内的点使 \(ax + by + c\) 同号。

为确定哪个半平面对应 \(> 0\)，只需取一个特殊点检验。例如，若
\(c \neq 0\)，取原点 \((0, 0)\)：

\begin{itemize}
\tightlist
\item
  若 \(c > 0\)，则原点在 \(ax + by + c > 0\) 的半平面内。
\item
  若 \(c < 0\)，则原点在 \(ax + by + c < 0\) 的半平面内。
\end{itemize}

\(\blacksquare\)

\textbf{推论 11.1.1} 二元一次不等式组

\[
\begin{cases}
a_1 x + b_1 y + c_1 \leq 0 \\
a_2 x + b_2 y + c_2 \leq 0 \\
\vdots \\
a_n x + b_n y + c_n \leq 0
\end{cases}
\]

的解集是各不等式对应半平面的交集。若此交集非空且有界，则为凸多边形。

\textbf{证明}：由定理11.1.1，每个二元一次不等式
\(a_i x + b_i y + c_i \leq 0\)
的解集为一个闭半平面。不等式组的解集为满足所有不等式的点的集合，即各半平面的交集。有限个半平面的交集为凸集（因半平面为凸集，凸集的交仍为凸集）。当此交集非空且有界时，凸集为凸多边形。\(\blacksquare\)

\bigskip\hrule\bigskip

\subsubsection{11.2
线性规划基本定理}\label{ux7ebfux6027ux89c4ux5212ux57faux672cux5b9aux7406}

\textbf{定义 11.2.1（线性规划问题）}

\textbf{线性规划问题}是求线性目标函数 \(z = ax + by\)
在由线性不等式组（约束条件）确定的可行域上的最大值或最小值。

\textbf{定义 11.2.2（可行域与顶点）}

\begin{itemize}
\tightlist
\item
  \textbf{可行域}：满足所有约束条件的点 \((x, y)\) 的集合。
\item
  \textbf{顶点}：可行域的凸多边形的顶点（角点）。
\end{itemize}

\textbf{定理 11.2.1（线性规划基本定理）} 设可行域 \(D\)
为非空有界的凸多边形，目标函数 \(z = ax + by\) 在 \(D\)
上取得最大值和最小值。最大值和最小值必在 \(D\) 的某个顶点处取得。

\textbf{证明}：

\textbf{第一步：最优值的存在性。} \(D\) 为非空有界闭集，\(z = ax + by\)
为连续函数。由闭区间上连续函数的极值定理的推广（多元版本），\(z\) 在
\(D\)
上取得最大值和最小值。（注：多元连续函数的极值定理将在后续多元分析章节中严格建立，此处仅作直观应用。）

\textbf{第二步：最优值在边界取得。} 若 \(a = b = 0\)，则 \(z \equiv 0\)
为常函数，\(D\) 的任意顶点均为最优点，结论成立。以下设
\((a,b) \neq (0,0)\)。

设 \((x_0, y_0) \in D\) 为 \(z\) 的最大值点。假设 \((x_0, y_0)\) 不在
\(D\) 的边界上，即 \((x_0, y_0)\) 是 \(D\) 的内点。则存在
\(\varepsilon > 0\) 使得以 \((x_0, y_0)\) 为圆心、\(\varepsilon\)
为半径的圆含于 \(D\) 内。考虑方向 \((a, b)\) 上的点
\((x_0 + ta, y_0 + tb)\)（\(t > 0\) 充分小）：

\[
z(x_0 + ta, y_0 + tb) = a(x_0 + ta) + b(y_0 + tb) = z(x_0, y_0) + t(a^2 + b^2) > z(x_0, y_0)
\]

这与 \((x_0, y_0)\) 是最大值点矛盾。故最大值点在边界上。

\textbf{第三步：最优值在顶点取得。} \(D\)
的边界由若干线段组成。设最大值在边界线段 \(\overline{PQ}\)（端点
\(P, Q\) 为顶点）的内点 \(M\) 处取得。设
\(M = (1-\lambda)P + \lambda Q\)（\(0 < \lambda < 1\)）。由线性：

\[
z(M) = (1-\lambda) z(P) + \lambda z(Q)
\]

由于 \(z(M)\) 是 \(z(P)\) 和 \(z(Q)\)
的凸组合，\(z(M) \leq \max(z(P), z(Q))\)。但 \(z(M)\) 是最大值，故
\(z(M) = z(P) = z(Q)\)。因此 \(P\) 和 \(Q\)
也是最大值点，即最大值在顶点处取得。

\(\blacksquare\)

\textbf{推论 11.2.1}
求解线性规划问题只需在可行域的所有顶点处计算目标函数的值，取最大（或最小）者即可。

\bigskip\hrule\bigskip

\subsubsection{11.3 图解法}\label{ux56feux89e3ux6cd5}

\textbf{定义 11.3.1（图解法）}

对于只有两个变量的线性规划问题，可以用平面图形求解，方法如下：

\begin{enumerate}
\def\labelenumi{\arabic{enumi}.}
\tightlist
\item
  在平面直角坐标系中画出各约束条件对应的直线。
\item
  确定每条直线对应的半平面，取其交集得到可行域。
\item
  画出目标函数 \(z = ax + by\) 的一组等值线（平行直线族
  \(ax + by = k\)）。
\item
  沿目标函数增大的方向平移等值线，最后离开可行域的点即为最优解。
\end{enumerate}

\textbf{定理 11.3.1（等值线的平行性）} 目标函数 \(z = ax + by\)
的等值线族 \(\{ax + by = k : k \in \mathbb{R}\}\) 中的直线互相平行。

\textbf{证明}：等值线 \(ax + by = k\) 的法向量为 \((a, b)\)，与 \(k\)
无关。故所有等值线法向量相同，互相平行。 \(\blacksquare\)

\textbf{定理 11.3.2} 若可行域 \(D\) 非空有界，且目标函数 \(z = ax + by\)
的等值线不与 \(D\) 的任一边平行，则最优解唯一，位于某一顶点处。

\textbf{证明}：由定理11.2.1，最优值在某顶点 \(V\) 处取得。假设另一点
\(W \in D\) 也取得最优值，则 \(z(W) = z(V)\)，即 \(W\) 在过 \(V\)
的等值线上。若 \(W\) 不是顶点，则 \(W\) 在某条边 \(\overline{PQ}\)
上。等值线方向向量为 \((-b, a)\)，边 \(\overline{PQ}\) 方向向量为
\(Q - P\)。若两者不平行，等值线与边至多一个交点，故 \(W\)
必与某一端点重合。 \(\blacksquare\)

\textbf{定理 11.3.3（无界可行域的情形）} 若可行域 \(D\)
无界，目标函数可能无最大值或无最小值（或两者皆无）。

\textbf{证明}：设可行域 \(D\) 包含射线
\(\{(x_0 + t d_1, y_0 + t d_2) : t \geq 0\}\)，其中
\((d_1, d_2) \neq (0, 0)\)。则

\[
z(x_0 + td_1, y_0 + td_2) = z(x_0, y_0) + t(ad_1 + bd_2)
\]

\begin{itemize}
\tightlist
\item
  若 \(ad_1 + bd_2 > 0\)，当 \(t \to +\infty\) 时
  \(z \to +\infty\)，无最大值。
\item
  若 \(ad_1 + bd_2 < 0\)，当 \(t \to +\infty\) 时
  \(z \to -\infty\)，无最小值。
\item
  若 \(ad_1 + bd_2 = 0\)，沿射线 \(z\) 为常数，需要进一步分析。
\end{itemize}

\(\blacksquare\)

\begin{quote}
\textbf{本章展望}：线性规划是运筹学的核心内容。本章仅涉及二元情形的图解法，更一般的单纯形法和对偶理论将在高等代数和运筹学课程中深入讨论。至此，卷二''代数''的内容全部完成，为读者进入高等数学的学习奠定了坚实的代数基础。
\end{quote}

\newpage

\section{卷三：初等几何}\label{ux5377ux4e09ux521dux7b49ux51e0ux4f55}

\begin{quote}
本卷从希尔伯特公理系统出发，以严格的逻辑推导建立欧氏几何的完整理论体系。我们首先在第12章确立几何的公理基础，然后在第13章系统推导平面几何的全部核心定理，在第14章将这些结果推广到三维空间，最后在第15章引入向量这一强有力的代数工具，为后续的解析几何和函数分析奠定基础。
\end{quote}

\bigskip\hrule\bigskip

\subsection{Ch 12
欧氏几何公理体系}\label{ch-12-ux6b27ux6c0fux51e0ux4f55ux516cux7406ux4f53ux7cfb}

\textbf{前置知识}：本章假定读者已掌握卷一 Ch 2
逻辑基础（命题逻辑、谓词逻辑、证明方法）和 Ch 3
集合论（集合的运算、关系与函数）。几何公理将使用逻辑语言精确表述几何直觉，每一条公理都是一个严格的形式化命题。

\bigskip\hrule\bigskip

几何学是数学中最古老的分支之一。公元前约300年，欧几里得在《几何原本》中首次尝试用公理化方法建立几何学体系。然而，欧几里得的公理系统在严谨性上存在明显不足------他的一些''公理''实际上是定理，某些证明中隐含了未经声明的假设。1899年，德国数学家大卫·希尔伯特在其著作《几何基础》中提出了一个完整的公理系统，将欧氏几何建立在严格的逻辑基础之上。本章详细阐述这一公理系统，并从中推导出若干最基本的几何命题。

希尔伯特公理系统由五组公理构成：关联公理（描述点、直线、平面之间的结合关系）、顺序公理（描述点在直线上的排列顺序）、合同公理（描述线段和角的全等关系）、平行公理（描述平行线的存在性和唯一性）以及连续公理（描述直线的连续性和完备性）。这五组公理共同刻画了欧氏几何的全部结构。

\subsubsection{12.1
基本概念与关联公理}\label{ux57faux672cux6982ux5ff5ux4e0eux5173ux8054ux516cux7406}

\paragraph{12.1.1
基本对象与基本关系}\label{ux57faux672cux5bf9ux8c61ux4e0eux57faux672cux5173ux7cfb}

希尔伯特公理系统建立在三个\textbf{基本对象}和三个\textbf{基本关系}之上：

\begin{itemize}
\tightlist
\item
  \textbf{基本对象}：点（point）、直线（line）、平面（plane）。
\item
  \textbf{基本关系}：关联（incidence，又称''结合''或''属于''）、介于（betweenness，又称''顺序''）、合同（congruence，又称''全等''）。
\end{itemize}

这些基本对象和基本关系本身不加定义，它们的意义完全由公理所规定的性质来确定。这正体现了公理化方法的精神：我们不追问''点是什么''，而只关心''点满足什么性质''。

我们采用如下记号： - 点用大写字母 \(A, B, C, \ldots\) 表示； -
直线用小写字母 \(a, b, c, \ldots\) 或 \(\overleftrightarrow{AB}\) 表示；
- 平面用希腊字母 \(\alpha, \beta, \gamma, \ldots\) 表示； - ``点 \(A\)
在直线 \(a\) 上''或''直线 \(a\) 经过点 \(A\)``记为 \(A \in a\)； -''点
\(A\) 在平面 \(\alpha\) 上''记为 \(A \in \alpha\)； - ``直线 \(a\)
在平面 \(\alpha\) 上''记为 \(a \subset \alpha\)。

\paragraph{12.1.2
关联公理（结合公理）I1--I8}\label{ux5173ux8054ux516cux7406ux7ed3ux5408ux516cux7406i1i8}

关联公理描述点、直线、平面之间的基本结合关系。

\textbf{公理 I1}：对于任意两个不同的点 \(A\) 和 \(B\)，存在一条直线
\(a\) 使得 \(A \in a\) 且 \(B \in a\)。

\textbf{公理 I2}：对于任意两个不同的点 \(A\) 和 \(B\)，至多存在一条直线
\(a\) 使得 \(A \in a\) 且 \(B \in a\)。

由公理 I1 和 I2
可知，任意两个不同的点\textbf{唯一确定}一条直线。我们将通过点 \(A\) 和
\(B\) 的唯一直线记为 \(\overleftrightarrow{AB}\)。

\textbf{公理 I3}：每条直线上至少存在两个不同的点。存在三个不共线的点。

\textbf{公理 I4}：对于任意三个不共线的点
\(A\)、\(B\)、\(C\)，存在一个平面 \(\alpha\) 使得
\(A \in \alpha\)、\(B \in \alpha\)、\(C \in \alpha\)。对于任意平面，其上至少存在一个点。

\textbf{公理 I5}：对于任意三个不共线的点
\(A\)、\(B\)、\(C\)，至多存在一个平面使得 \(A\)、\(B\)、\(C\) 都在其上。

由公理 I4 和 I5 可知，任意三个不共线的点\textbf{唯一确定}一个平面。

\textbf{公理 I6}：如果一条直线上的两个不同的点都在平面 \(\alpha\)
上，那么该直线上的所有点都在 \(\alpha\) 上（即直线含于 \(\alpha\)）。

\textbf{公理
I7}：如果两个不同的平面有一个公共点，那么它们至少还有一个公共点。

\textbf{公理 I8}：至少存在四个不共面的点。

公理 I8
保证了空间的''三维性''------排除了所有点都在同一个平面上的退化情况。

\textbf{定义
12.1}（共线）：如果若干个点都在同一条直线上，则称这些点是\textbf{共线的}。

\textbf{定义
12.2}（共面）：如果若干个点都在同一个平面上，则称这些点是\textbf{共面的}。

\textbf{命题 12.1}：两条不同的直线至多有一个公共点。

\textbf{证明}：假设直线 \(a\) 和直线 \(b\) 有两个不同的公共点 \(A\) 和
\(B\)。根据公理 I2，通过 \(A\) 和 \(B\) 至多有一条直线，因此
\(a = b\)，与 \(a\) 和 \(b\) 不同的假设矛盾。\(\blacksquare\)

\textbf{命题
12.2}：如果两个不同的平面有一个公共点，那么它们的交集是一条直线。

\textbf{证明}：设平面 \(\alpha\) 和平面 \(\beta\) 有一个公共点
\(P\)。根据公理 I7，它们至少还有一个公共点 \(Q\)（\(Q \neq P\)）。由公理
I2，直线 \(\overleftrightarrow{PQ}\) 唯一确定。

设 \(R\) 是 \(\alpha\) 和 \(\beta\) 的任意公共点。如果 \(R\) 不在
\(\overleftrightarrow{PQ}\) 上，则 \(P\)、\(Q\)、\(R\) 不共线，由公理
I5，它们唯一确定一个平面。但 \(P\)、\(Q\)、\(R\) 都在 \(\alpha\)
上，也都在 \(\beta\) 上，因此
\(\alpha = \beta\)，与两个平面不同的假设矛盾。

因此，\(\alpha\) 和 \(\beta\) 的所有公共点都在
\(\overleftrightarrow{PQ}\) 上。另一方面，由公理
I6，\(\overleftrightarrow{PQ}\) 上的所有点既在 \(\alpha\) 上也在
\(\beta\) 上。所以
\(\alpha \cap \beta = \overleftrightarrow{PQ}\)。\(\blacksquare\)

\subsubsection{12.2 顺序公理}\label{ux987aux5e8fux516cux7406}

顺序公理描述点在直线上的相对位置关系，引入''介于''的概念。这是建立线段、射线、角的内部等概念的基础。

\textbf{定义 12.3}（介于关系）：设 \(A\)、\(B\)、\(C\)
是一条直线上的三个不同的点。如果 \(B\) 在 \(A\) 和 \(C\) 之间，则记为
\(A * B * C\)。

\textbf{公理 O1}：如果 \(A * B * C\)，则 \(A\)、\(B\)、\(C\)
是同一条直线上三个不同的点，且 \(C * B * A\)。

\textbf{公理 O2}：对于任意两个不同的点 \(A\) 和 \(B\)，直线
\(\overleftrightarrow{AB}\) 上至少存在一个点 \(C\) 使得 \(A * B * C\)。

\textbf{公理
O3}：对于同一条直线上的任意三个不同的点，至多有一个点介于另外两个点之间。

\textbf{公理 O4}（帕施公理）：设 \(A\)、\(B\)、\(C\)
是三个不共线的点，\(a\) 是一条不经过 \(A\)、\(B\)、\(C\)
中任何一个点的直线。如果直线 \(a\) 经过线段 \(AB\)
的某个内点，那么它也经过线段 \(BC\) 或线段 \(AC\) 的某个内点。

\textbf{定义 12.4}（线段）：对于两个不同的点 \(A\) 和
\(B\)，\textbf{线段} \(AB\) 是指直线 \(\overleftrightarrow{AB}\)
上所有满足 \(A * X * B\) 的点 \(X\) 以及 \(A\)、\(B\)
两点本身的集合。\(A\) 和 \(B\) 称为线段 \(AB\) 的\textbf{端点}。线段
\(AB\) 记为 \(\overline{AB}\)。

\textbf{定义 12.5}（射线）：对于两个不同的点 \(A\) 和 \(B\)，以 \(A\)
为端点、经过 \(B\) 的\textbf{射线} \(\overrightarrow{AB}\) 是指直线
\(\overleftrightarrow{AB}\) 上所有满足 \(A * B * X\) 或 \(X = A\) 或
\(X = B\) 的点 \(X\) 的集合。

\textbf{定义 12.6}（开半直线）：直线 \(\overleftrightarrow{AB}\)
上除去线段 \(\overline{AB}\) 的两个部分（不包括 \(A\) 和
\(B\)），分别称为以 \(A\) 为分界点、关于 \(B\)
的\textbf{开半直线}及其\textbf{对侧}。

\textbf{引理 12.1}（介于关系的基本推论）：设 \(A\)、\(B\)、\(C\)、\(D\)
是同一条直线上的不同点。

\begin{enumerate}
\def\labelenumi{\arabic{enumi}.}
\tightlist
\item
  若 \(A * B * C\)，则 \(C * B * A\)。
\item
  若 \(A * B * D\) 且 \(C * B * D\)，则 \(A * B * C\)。
\end{enumerate}

\textbf{证明}：

性质 1 是公理 O1 的直接推论。

性质 2 的证明：由公理 O3，在 \(\{A, B, C\}\)
中至多一种''介于''关系成立。我们用反证法排除 \(A * C * B\) 和
\(B * A * C\) 两种可能：

\begin{itemize}
\item
  若 \(A * C * B\)，则 \(C\) 在 \(A\) 和 \(B\) 之间。由 \(A * B * D\) 知
  \(B\) 在 \(A\) 和 \(D\) 之间，故 \(C\) 介于 \(A\) 与 \(B\) 之间、而
  \(B\) 又介于 \(A\) 与 \(D\) 之间，因此 \(B\) 在 \(C\) 和 \(D\)
  之间，即 \(C * B * D\)。但这意味着 \(B\) 和 \(D\) 在 \(C\) 的两侧，而
  \(A * C * B\) 表明 \(B\) 和 \(A\) 也在 \(C\) 的两侧。综合知 \(A\) 和
  \(D\) 在 \(C\) 的同侧，从而 \(A * C * D\)。又 \(C * B * D\) 表明 \(D\)
  在 \(B\) 的远离 \(C\) 的一侧，而 \(A * C * D\) 表明 \(D\) 在 \(C\)
  的远离 \(A\) 的一侧，两者一致；但 \(A * C * B\) 表明 \(B\) 在 \(C\)
  的远离 \(A\) 的一侧，而 \(C * B * D\) 表明 \(B\) 在 \(C\) 和 \(D\)
  之间，即 \(B\) 在 \(C\) 的远离 \(A\) 的一侧、\(D\) 在 \(B\) 的远离
  \(C\) 的一侧，这与 \(A * C * D\) 中 \(D\) 在 \(C\) 的远离 \(A\)
  的一侧且 \(A\) 在 \(C\) 和 \(D\) 之间矛盾（\(A\) 不可能既在 \(C\)
  的远离 \(B\) 的一侧又在 \(C\) 和 \(D\) 之间）。因此 \(A * C * B\)
  不成立。
\item
  若 \(B * A * C\)，则 \(A\) 在 \(B\) 和 \(C\) 之间。由 \(A * B * D\) 知
  \(B\) 在 \(A\) 和 \(D\) 之间，故 \(B\) 和 \(D\) 在 \(A\) 的两侧；而
  \(B * A * C\) 表明 \(B\) 和 \(C\) 也在 \(A\) 的两侧，因此 \(C\) 和
  \(D\) 在 \(A\) 的同侧。由 \(C * B * D\)（性质 1 得
  \(D * B * C\)），\(B\) 在 \(C\) 和 \(D\) 之间，\(C\) 和 \(D\) 在 \(B\)
  的两侧。但 \(B * A * C\) 表明 \(C\) 在 \(A\) 的远离 \(B\) 的一侧，而
  \(A * B * D\) 表明 \(D\) 在 \(B\) 的远离 \(A\) 的一侧。\(C\) 和 \(D\)
  在 \(A\) 的同侧，但 \(B\) 在 \(C\) 和 \(D\) 之间要求 \(C\) 和 \(D\) 在
  \(B\) 的两侧，而 \(B\) 在 \(A\) 和 \(D\) 之间要求 \(D\) 在 \(A\)
  的远离 \(B\) 的一侧、\(C\) 在 \(A\) 的靠近 \(B\) 的一侧（由
  \(B * A * C\)），矛盾。因此 \(B * A * C\) 不成立。
\end{itemize}

排除以上两种可能后，必有 \(A * B * C\) 成立。\(\blacksquare\)''

\textbf{命题 12.3}：设 \(A\)、\(B\)、\(C\)
是同一条直线上的三个不同的点。则恰好有以下三种情况之一成立：

\begin{enumerate}
\def\labelenumi{\arabic{enumi}.}
\tightlist
\item
  \(A * B * C\)；
\item
  \(A * C * B\)；
\item
  \(B * A * C\)。
\end{enumerate}

\textbf{证明}：由公理
O3（唯一性），至多有一种''介于''关系成立。下面证明至少有一种成立，采用反证法。

假设 \(A\)、\(B\)、\(C\) 之间不存在任何''介于''关系，即
\(A * B * C\)、\(A * C * B\)、\(B * A * C\) 均不成立。我们分步导出矛盾。

\textbf{第一步}：由公理 O2（存在性），存在点 \(D\) 使得
\(A * B * D\)。若 \(D = C\)，则 \(A * B * C\) 成立，与假设矛盾。故
\(D \neq C\)，即 \(A\)、\(B\)、\(C\)、\(D\) 是直线上四个不同的点。

\textbf{第二步}：考虑 \(C\) 与 \(D\) 的位置。由公理 O3，在三个点
\(\{B, C, D\}\) 中恰好有一种''介于''关系成立。由 \(A * B * D\) 知
\(B \neq D\)，且
\(C \neq B\)、\(C \neq D\)，故以下三种情况恰有一种成立：

\textbf{(i)} \(B * C * D\)：此时考虑三个点 \(\{A, B, C\}\)。由公理
O4（帕施公理），构造辅助线：取不在直线 \(l\) 上的点 \(P\)，作三角形
\(\triangle ABD\)。点 \(C\) 满足 \(B * C * D\)，即 \(C\) 在线段 \(BD\)
上。由于 \(A * B * D\)，点 \(B\) 在线段 \(AD\) 上。在三角形
\(\triangle ACD\) 中，\(B\) 在线段 \(AD\) 上，\(B * C * D\) 表明 \(C\)
在线段 \(BD\) 上。由公理 O3，我们分析 \(\{A, B, C\}\)
的排列：\(A * B * C\) 和 \(B * A * C\) 互斥（公理 O3 唯一性）。若
\(A * B * C\)，则与假设矛盾。因此必有 \(B * A * C\) 或
\(A * C * B\)。但若 \(B * A * C\)，则 \(A\) 在 \(B\) 和 \(C\) 之间；而由
\(A * B * D\) 和 \(B * C * D\)，\(A\)、\(B\)、\(C\) 在 \(D\) 的同侧且
\(B\) 在 \(A\) 与 \(C\) 之间，这与 \(B * A * C\) 矛盾（公理
O3：\(A * B * D\) 和 \(B * A * C\) 不能同时成立，否则 \(B\) 既在 \(A\)
和 \(D\) 之间又在 \(A\) 的另一侧）。类似地，\(A * C * B\)
也会导致矛盾。因此在情形 (i) 下 \(A * B * C\) 必成立，与假设矛盾。

\textbf{(ii)} \(B * D * C\)：此时 \(D\) 在 \(B\) 和 \(C\) 之间。由
\(A * B * D\) 和 \(B * D * C\)，在三个点 \(\{A, B, C\}\) 中：由公理
O3，排列 \(A * B * C\)、\(A * C * B\)、\(B * A * C\)
恰有一种成立。我们排除后两种：若 \(B * A * C\)，则 \(A\) 在 \(B\) 和
\(C\) 之间；但 \(A * B * D\) 表明 \(B\) 在 \(A\) 和 \(D\)
之间，\(B * D * C\) 表明 \(D\) 在 \(B\) 和 \(C\) 之间，因此 \(A\) 和
\(C\) 在 \(B\) 的异侧，\(A\) 不可能在 \(B\) 和 \(C\) 之间，矛盾。若
\(A * C * B\)，则 \(C\) 在 \(A\) 和 \(B\) 之间；但 \(A * B * D\) 表明
\(B\) 在 \(A\) 和 \(D\) 之间，\(C\) 在 \(B\) 的远离 \(A\) 一侧（因为
\(B * D * C\)），矛盾。因此在情形 (ii) 下 \(A * B * C\)
必成立，与假设矛盾。

\textbf{(iii)} \(C * B * D\)：此时 \(B\) 在 \(C\) 和 \(D\) 之间。结合
\(A * B * D\)（\(B\) 在 \(A\) 和 \(D\) 之间），可知 \(A\) 和 \(C\) 在
\(B\) 的两侧（因为 \(B\) 将直线分为两个半直线，\(A\) 和 \(D\) 在 \(B\)
的两侧，而 \(C\) 和 \(D\) 也在 \(B\) 的两侧，故 \(A\) 和 \(C\) 在 \(B\)
的异侧）。因此 \(B * A * C\)（\(B\) 在 \(A\) 和 \(C\)
之间），与假设矛盾。

\textbf{第三步}：以上三种情况均与假设矛盾，故假设不成立。三个点之间至少有一种''介于''关系成立。结合至多一种成立（公理
O3），恰好有一种''介于''关系成立。\(\blacksquare\)

\textbf{定义 12.7}（三角形）：三个不共线的点 \(A\)、\(B\)、\(C\)
以及三条线段 \(\overline{AB}\)、\(\overline{BC}\)、\(\overline{CA}\)
组成的图形称为\textbf{三角形}，记为 \(\triangle ABC\)。点
\(A\)、\(B\)、\(C\) 称为三角形的\textbf{顶点}，线段
\(\overline{AB}\)、\(\overline{BC}\)、\(\overline{CA}\)
称为三角形的\textbf{边}。

\subsubsection{12.3 合同公理}\label{ux5408ux540cux516cux7406}

合同公理引入线段和角的''全等''（合同）关系，用 \(\cong\)
表示。合同关系刻画了''长度相等''和''角度相等''的几何直觉。

\paragraph{12.3.1 线段的合同}\label{ux7ebfux6bb5ux7684ux5408ux540c}

\textbf{公理 C1}：设 \(A\)、\(B\) 是直线 \(a\) 上的两个点，\(A'\) 是直线
\(a'\) 上的点，\(p'\) 是 \(a'\) 上以 \(A'\) 为端点的一条射线。则在射线
\(p'\) 上存在唯一的点 \(B'\) 使得
\(\overline{AB} \cong \overline{A'B'}\)。

\textbf{公理 C2}（合同的传递性）：如果
\(\overline{AB} \cong \overline{A'B'}\) 且
\(\overline{AB} \cong \overline{A''B''}\)，则
\(\overline{A'B'} \cong \overline{A''B''}\)。

\textbf{公理 C3}（合同的可加性）：设 \(B\) 是线段 \(\overline{AC}\)
上的点（即 \(A * B * C\)），\(B'\) 是线段 \(\overline{A'C'}\) 上的点（即
\(A' * B' * C'\)）。如果 \(\overline{AB} \cong \overline{A'B'}\) 且
\(\overline{BC} \cong \overline{B'C'}\)，则
\(\overline{AC} \cong \overline{A'C'}\)。

由公理 C1--C3 可知，线段的合同关系 \(\cong\)
是一个等价关系（自反性、对称性、传递性），并且合同关系在同一条直线上具有可加性。这使得''长度''的概念有了严格的公理基础。

\paragraph{12.3.2 角的合同}\label{ux89d2ux7684ux5408ux540c}

\textbf{定义 12.8}（角）：设 \(h\) 和 \(k\) 是从同一点 \(O\)
出发的两条不共线的射线。由这两条射线所确定的有界区域（不包含 \(O\)
的对顶角区域）称为角的\textbf{内部}。\(h\)、\(k\)
以及角内部的所有点组成的集合称为一个\textbf{角}，记为 \(\angle(h, k)\)
或 \(\angle hOk\)。射线 \(h\) 和 \(k\) 称为角的\textbf{边}，点 \(O\)
称为角的\textbf{顶点}。

\textbf{公理 C4}：设 \(\angle(h, k)\) 是一个角，\(h'\) 是直线 \(a'\)
上以 \(O'\) 为端点的射线，\(a'\) 有确定的一侧 \(H'\)。则在 \(H'\)
上存在且仅存在一条以 \(O'\) 为端点的射线 \(k'\)，使得
\(\angle(h, k) \cong \angle(h', k')\)。

\textbf{公理 C5}：角的合同关系具有对称性，即若
\(\angle(h, k) \cong \angle(h', k')\)，则
\(\angle(h', k') \cong \angle(h, k)\)。此外，\(\angle(h,k) \cong \angle(k, h)\)。

\textbf{公理 C6}（SAS 公理）：设 \(A\)、\(B\)、\(C\)
是三个不共线的点，\(A'\)、\(B'\)、\(C'\) 也是三个不共线的点。如果

\[\overline{AB} \cong \overline{A'B'}, \quad \overline{AC} \cong \overline{A'C'}, \quad \angle BAC \cong \angle B'A'C',\]

则 \(\angle ABC \cong \angle A'B'C'\) 且
\(\angle ACB \cong \angle A'C'B'\)。

公理 C6
是希尔伯特公理系统中最关键的公理之一------它实际上就是三角形全等的
SAS（边角边）判定准则。注意欧几里得在《几何原本》中将 SAS
作为定理来证明，但希尔伯特认识到其证明实际上隐含了未声明的假设，因此将它提升为公理。

\textbf{定义 12.9}（直角、垂直）：设 \(A\)、\(B\)、\(C\) 是直线 \(h\)
上的三个点（\(A * B * C\)），\(k\) 是以 \(B\) 为端点的射线。如果
\(\angle ABk \cong \angle CBk\)，则称 \(\angle ABk\)
为\textbf{直角}，射线 \(k\) 与直线 \(h\) \textbf{互相垂直}，记为
\(k \perp h\)。

\subsubsection{12.4 平行公理}\label{ux5e73ux884cux516cux7406}

\textbf{公理 P}（平行公理 / 第五公设）：设 \(a\) 是一条直线，\(P\)
是不在 \(a\) 上的一个点。则在 \(P\) 和 \(a\)
确定的平面上，至多存在一条直线经过 \(P\) 且与 \(a\) 不相交。

\textbf{定义
12.10}（平行）：在同一平面上，两条不相交的直线称为\textbf{平行}的，记为
\(a \parallel b\)。

平行公理是欧氏几何区别于非欧几何的关键所在。如果否定这条公理，就得到两种非欧几何：
-
\textbf{双曲几何}（罗巴切夫斯基几何）：过直线外一点有\textbf{无数条}平行线。
- \textbf{椭圆几何}（黎曼几何）：过直线外一点\textbf{不存在}平行线。

\textbf{命题 12.4}（平行线的存在性）：设 \(a\) 是一条直线，\(P\) 是不在
\(a\) 上的一个点。则在 \(P\) 和 \(a\) 确定的平面上，至少存在一条直线经过
\(P\) 且与 \(a\) 平行。

\textbf{证明}：在直线 \(a\) 上取两点 \(Q\)、\(R\)。在 \(P\) 点处，由公理
C4，可以构造射线 \(\overrightarrow{PS}\) 使得
\(\angle QPS \cong \angle PQR\)（同位角相等），且
\(\overrightarrow{PS}\) 与 \(a\) 在 \(\overleftrightarrow{PQ}\)
的同一侧。直线 \(\overleftrightarrow{PS}\) 与 \(a\) 被
\(\overleftrightarrow{PQ}\) 所截，同位角相等，故
\(\overleftrightarrow{PS} \parallel a\)。

\textbf{严格论证}：假设 \(\overleftrightarrow{PS}\) 与 \(a\) 相交于点
\(X\)（\(X \neq Q\)）。考虑 \(\triangle PQX\)，\(\angle QPS\) 是
\(\triangle PQX\) 在顶点 \(P\) 处的角。由于 \(X\) 在 \(a\) 上且
\(X \neq Q\)，\(\angle PQX\) 是 \(\triangle PQX\) 在顶点 \(Q\)
处的角。由外角定理（定理 12.3）（此定理将在后文定理 12.3
中正式陈述和证明，此处仅使用其结论），\(\triangle PQX\)
的外角严格大于不相邻内角。但 \(\angle QPS = \angle PQR\)（由构造），而
\(\angle PQR\) 与 \(\angle PQX\) 共享边 \(\overrightarrow{QP}\) 且 \(R\)
和 \(X\) 在 \(a\) 上（\(R\) 是 \(a\) 上的点使得 \(\angle PQR\)
为原角），这导致矛盾。故 \(\overleftrightarrow{PS}\) 不与 \(a\) 相交，即
\(\overleftrightarrow{PS} \parallel a\)。\(\blacksquare\)

结合平行公理 P 和命题
12.4，我们知道：经过直线外一点\textbf{恰好}存在一条直线与已知直线平行。

\subsubsection{12.5 连续公理}\label{ux8fdeux7eedux516cux7406}

\textbf{公理 D1}（阿基米德公理）：设 \(\overline{AB}\) 和
\(\overline{CD}\) 是两条线段。则存在正整数 \(n\)，使得将
\(\overline{CD}\) 沿射线 \(\overrightarrow{AB}\) 放置 \(n\)
次后，其长度超过 \(\overline{AB}\)。即存在点
\(A = A_0, A_1, A_2, \ldots, A_n\) 在射线 \(\overrightarrow{AB}\)
上，使得 \(A_{i-1}A_i \cong CD\)（\(i = 1, 2, \ldots, n\)），且
\(A * B * A_n\)。

\textbf{公理
D2}（直线完备性公理）：直线上的点组成的系统不能被扩充为一个仍然满足公理
I1--I2、O1--O4、C1--C3 和 D1 的更大的系统。

公理 D2
保证了直线上的点构成一个''完备''的系统，没有''空隙''。它等价于实数的完备性公理（确界原理），使得我们可以对线段进行任意精度的测量。

\subsubsection{12.6
从公理推导基本命题}\label{ux4eceux516cux7406ux63a8ux5bfcux57faux672cux547dux9898}

在建立了完整的公理系统之后，我们现在从中推导若干最基本的几何命题。这些命题是后续章节中全部平面几何和立体几何定理的基础。

\textbf{引理 12.0（SSS
全等判定）}：若两个三角形的三对边分别合同，则两个三角形全等。即：若
\(\overline{AB} \cong \overline{A'B'}\)，\(\overline{BC} \cong \overline{B'C'}\)，\(\overline{CA} \cong \overline{C'A'}\)，则
\(\triangle ABC \cong \triangle A'B'C'\)。

\textbf{证明}：将 \(\triangle A'B'C'\) 放置在平面上使得 \(A'B'\) 与
\(AB\) 重合，且 \(C'\) 与 \(C\) 在 \(AB\) 的同侧。若 \(C' = C\)
则已证。若 \(C' \neq C\)，由
\(\overline{AC} \cong \overline{A'C'} = \overline{AC'}\)，\(\overline{BC} \cong \overline{B'C'} = \overline{BC'}\)，知
\(C\) 和 \(C'\) 到 \(A\)、\(B\) 的距离相等。取 \(AB\) 的中点
\(M\)（由后文定理 12.1 或直接构造），连接 \(CC'\)。在 \(\triangle ACC'\)
中，\(AC = AC'\)，故 \(\triangle ACC'\) 是等腰三角形，\(M\)（\(AB\)
中点）到 \(C\) 和 \(C'\) 的距离相等。在 \(\triangle MCC'\)
中，\(MC = MC'\) 且 \(CC'\) 公共。若 \(C' \neq C\)，则
\(\triangle MCC'\) 退化，矛盾。故 \(C' = C\)，三角形全等。

\textbf{注} 此处的 SSS 证明不依赖定理 12.1（中点存在性），而是直接从 SAS
公理出发。证明思路：将 \(\triangle A'B'C'\) 反射到
\(\overleftrightarrow{AB}\) 的另一侧得到 \(\triangle A'B'C''\)，利用 SAS
证明 \(\triangle ABC \cong \triangle ABC''\)，再由唯一性（SAS
保证三角形由两边一夹角唯一确定）得 \(C = C''\)，从而 \(C' = C\)（因为
\(C'\) 和 \(C''\) 关于 \(\overleftrightarrow{AB}\) 对称且在同一侧时
\(C' = C'' = C\)）。\(\blacksquare\)

\paragraph{12.6.1
线段中点的存在性和唯一性}\label{ux7ebfux6bb5ux4e2dux70b9ux7684ux5b58ux5728ux6027ux548cux552fux4e00ux6027}

\textbf{定理 12.1}（线段中点的存在性和唯一性）：对于任意线段
\(\overline{AB}\)，存在唯一的点 \(M\) 使得 \(A * M * B\) 且
\(\overline{AM} \cong \overline{MB}\)。

\textbf{证明}：

\textbf{存在性}：在直线 \(\overleftrightarrow{AB}\) 外取一点
\(C\)（由公理 I8，存在不共面的四个点，故不共线的三个点必然存在，再由公理
I3，可以在 \(\overleftrightarrow{AB}\) 外取点）。在射线
\(\overrightarrow{AC}\) 上取点 \(P'\) 使得
\(\overline{AP'} \cong \overline{AB}\)（由公理 C1
保证这样的点存在）。由公理 C4，以 \(\overleftrightarrow{AB}\) 为界、在
\(C\) 所在的一侧作射线 \(\overrightarrow{BP'}\)，使得
\(\angle ABP' \cong \angle BAP'\)。在射线 \(\overrightarrow{BP'}\)
上取点 \(P\) 使得 \(\overline{BP} \cong \overline{AP'}\)（由公理 C1）。

在 \(\triangle ABP'\) 和 \(\triangle BAP\)
中：\(\overline{AP'} \cong \overline{AB}\)（由构造），\(\angle BAP' \cong \angle ABP\)（由构造），\(\overline{BP} \cong \overline{AP'}\)（由构造）。又
\(\overline{AB} \cong \overline{BA}\)（同一线段），故
\(\overline{AP'} \cong \overline{AB} \cong \overline{BA} \cong \overline{BP}\)。由
SAS（公理 C6），\(\triangle ABP' \cong \triangle BAP\)，因此
\(\overline{PA} \cong \overline{PB}\)。

类似地，在 \(\overleftrightarrow{AB}\) 的另一侧取点 \(Q\) 使得
\(\overline{QA} \cong \overline{QB}\)（同样的构造方法）。连接
\(\overline{PQ}\)，设 \(\overleftrightarrow{PQ}\) 与
\(\overleftrightarrow{AB}\) 的交点为 \(M\)。

在 \(\triangle PAQ\) 和 \(\triangle PBQ\) 中： -
\(\overline{PA} \cong \overline{PB}\)（由构造）， -
\(\overline{QA} \cong \overline{QB}\)（由构造）， -
\(\overline{PQ} \cong \overline{PQ}\)（公共边）。

由 SSS 全等条件（注：SSS 判定可从 SAS 公理严格推出，其证明见定理
13.3；此处的 SSS 应用可理解为：由下方 SAS 论证可得 \(\triangle PAB\)
的对称性，等价地推出
\(\triangle PAQ \cong \triangle PBQ\)），\(\triangle PAQ \cong \triangle PBQ\)，因此
\(\angle APQ \cong \angle BPQ\)。

在 \(\triangle PAM\) 和 \(\triangle PBM\) 中： -
\(\overline{PA} \cong \overline{PB}\)（已证）， -
\(\angle APM \cong \angle BPM\)（已证）， -
\(\overline{PM} \cong \overline{PM}\)（公共边）。

由 SAS（公理 C6），\(\triangle PAM \cong \triangle PBM\)。因此
\(\overline{AM} \cong \overline{BM}\)，即 \(M\) 是 \(\overline{AB}\)
的中点。

\textbf{唯一性}：假设存在两个不同的中点 \(M\) 和 \(M'\)，即
\(\overline{AM} \cong \overline{MB}\) 且
\(\overline{AM'} \cong \overline{M'B}\)，且
\(A * M * B\)，\(A * M' * B\)。

假设 \(M \neq M'\)。由命题 12.3（三点之间恰好有一种''介于''关系），\(M\)
和 \(M'\) 在 \(\overleftrightarrow{AB}\) 上的顺序只有以下几种可能：

\textbf{情况 (i)}：\(A * M * M' * B\)。此时 \(A * M * M'\) 且
\(M * M' * B\)。由公理 C3（可加性）：由
\(A * M * M'\)，\(\overline{AM} \cong \overline{AM}\)（自反性），\(\overline{MM'} \cong \overline{MM'}\)，可得
\(\overline{AM'} \cong \overline{AM} + \overline{MM'}\)（即
\(\overline{AM'}\) 由 \(\overline{AM}\) 和 \(\overline{MM'}\)
拼接而成，严格地说：存在辅助线段使得合同关系保持可加）。同理，由
\(M * M' * B\) 和
\(\overline{MM'} \cong \overline{MM'}\)，\(\overline{M'B} \cong \overline{M'B}\)，可得
\(\overline{MB} \cong \overline{MM'} + \overline{M'B}\)。

由中点条件：\(\overline{AM} \cong \overline{MB}\) 且
\(\overline{AM'} \cong \overline{M'B}\)。代入上述关系：
\[\overline{AM} \cong \overline{MB} \cong \overline{MM'} + \overline{M'B} \cong \overline{MM'} + \overline{AM'}\]
又 \(\overline{AM'} \cong \overline{AM} + \overline{MM'}\)（由
\(A * M * M'\) 和可加性），故：
\[\overline{AM} \cong \overline{MM'} + (\overline{AM} + \overline{MM'})\]

这意味着
\(\overline{AM} \cong \overline{AM} + 2\cdot\overline{MM'}\)。但根据公理
C3 和顺序公理，线段经可加拼接后严格变长（\(A * M * M'\) 保证
\(\overline{MM'}\) 非零------因为 \(M \neq M'\)，由公理 O3
不存在零长度线段），矛盾。

\textbf{情况 (ii)}：\(A * M' * M * B\)。对称地，交换 \(M\) 和 \(M'\)
的角色，得到同样的矛盾。

\textbf{情况 (iii)}：\(M * A * M'\) 或 \(M' * A * M\)。这与
\(A * M * B\) 和 \(A * M' * B\) 矛盾（因为 \(A\) 是端点，\(M\) 和 \(M'\)
都在 \(A\) 和 \(B\) 之间，不可能在 \(A\) 的外侧）。

综上，\(M = M'\)，中点唯一。\(\blacksquare\)

\paragraph{12.6.2
角平分线的存在性和唯一性}\label{ux89d2ux5e73ux5206ux7ebfux7684ux5b58ux5728ux6027ux548cux552fux4e00ux6027}

\textbf{定理 12.2}（角平分线的存在性和唯一性）：对于任意角
\(\angle BAC\)，存在唯一的以 \(A\) 为端点的射线
\(\overrightarrow{AD}\)（\(D\) 在 \(\angle BAC\) 的内部），使得
\(\angle BAD \cong \angle DAC\)。

\textbf{证明}：

\textbf{存在性}：在射线 \(\overrightarrow{AB}\) 上取点 \(B'\)，在射线
\(\overrightarrow{AC}\) 上取点 \(C'\) 使得
\(\overline{AB'} \cong \overline{AC'}\)（由公理 C1，这样的 \(C'\)
存在）。连接 \(\overline{B'C'}\)，取其中点 \(M\)（由定理
12.1，中点存在）。射线 \(\overrightarrow{AM}\) 即为角平分线。

在 \(\triangle AB'M\) 和 \(\triangle AC'M\) 中： -
\(\overline{AB'} \cong \overline{AC'}\)（由构造）， -
\(\overline{B'M} \cong \overline{C'M}\)（\(M\) 是 \(B'C'\) 的中点）， -
\(\overline{AM} \cong \overline{AM}\)（公共边）。

由 SSS 全等（由 SAS 可推出 SSS，见定理
13.3），\(\triangle AB'M \cong \triangle AC'M\)，所以
\(\angle B'AM \cong \angle C'AM\)。

\textbf{唯一性}：我们首先引入角的大小比较。设 \(\angle(h_1, k_1)\) 和
\(\angle(h_2, k_2)\) 是以同一点 \(O\)
为顶点、\(h_1 = h_2\)（即一条公共边）的两个角。若射线 \(k_1\) 在
\(\angle(h_2, k_2)\) 的内部，则称
\(\angle(h_1, k_1) < \angle(h_2, k_2)\)。由角的内部的定义（射线 \(k_1\)
在角的内部意味着它与角的两边相交但不与任一边共线），这个比较关系是良定义的。

现在证明唯一性。假设存在两条不同的角平分线 \(\overrightarrow{AD}\) 和
\(\overrightarrow{AD'}\)，即 \(\angle BAD \cong \angle DAC\) 且
\(\angle BAD' \cong \angle D'AC\)。不妨设 \(\overrightarrow{AD'}\) 在
\(\angle BAD\)
的内部。则由角的大小比较定义，\(\angle BAD' < \angle BAD\)。

但由角平分线条件，\(\angle BAD' = \frac{1}{2}\angle BAC = \angle BAD\)，与
\(\angle BAD' < \angle BAD\) 矛盾。

因此角平分线唯一。\(\blacksquare\)

\paragraph{12.6.3 外角定理}\label{ux5916ux89d2ux5b9aux7406}

\textbf{定理
12.3}（外角定理）：三角形的一个外角大于与它不相邻的任意一个内角。

\textbf{注}：外角定理是不依赖平行公理的纯绝对几何定理。它在合同公理和顺序公理的基础上即可证明。

\textbf{证明}：

设 \(\triangle ABC\)，延长 \(\overline{BC}\) 到 \(D\)（由公理 O2，这样的
\(D\) 存在），则 \(\angle ACD\) 是 \(\triangle ABC\) 的一个外角。要证明
\(\angle ACD > \angle BAC\) 且 \(\angle ACD > \angle ABC\)。

\textbf{证明 \(\angle ACD > \angle BAC\)}：

取 \(\overline{AC}\) 的中点 \(M\)（由定理 12.1）。连接 \(\overline{BM}\)
并延长到 \(E\) 使得 \(B * M * E\) 且
\(\overline{ME} \cong \overline{BM}\)（由公理 C1）。连接
\(\overline{CE}\)。

在 \(\triangle ABM\) 和 \(\triangle CEM\) 中： -
\(\overline{AM} \cong \overline{CM}\)（\(M\) 是 \(\overline{AC}\)
的中点）， - \(\angle AMB \cong \angle CME\)（对顶角）， -
\(\overline{BM} \cong \overline{ME}\)（由构造）。

由 SAS（公理 C6），\(\triangle ABM \cong \triangle CEM\)。因此
\(\angle BAM \cong \angle ECM\)。

现在证明 \(E\) 在 \(\angle ACD\) 的内部，从而
\(\angle ACD > \angle ECM = \angle BAM = \angle BAC\)。由于
\(B * M * E\) 且 \(M\) 是 \(AC\) 的中点（\(A * M * C\)），点 \(E\) 和
\(B\) 位于直线 \(AC\) 的两侧。又 \(D\) 是 \(BC\)
延长线上的点（\(B * C * D\)），故 \(D\) 和 \(B\) 位于 \(C\) 的两侧，从而
\(D\) 与 \(E\) 在直线 \(AC\) 的同一侧。此外，由
\(\triangle ABM \cong \triangle CEM\) 可得
\(\angle MEC = \angle MBA\)，且 \(\angle MEC\) 与 \(\angle MCD\)
的关系保证射线 \(CE\) 严格位于射线 \(CA\) 与射线 \(CD\) 之间（因为
\(\angle ACM < \angle ACD\)，且 \(E\) 在 \(B\) 的对侧使得
\(\angle MCE < \angle ACD - \angle ACM\)）。因此 \(\angle ECD\) 是
\(\angle ACD\) 的一部分，从而：

\[\angle ACD > \angle ECM = \angle BAM = \angle BAC\]

\textbf{证明 \(\angle ACD > \angle ABC\)}：

类似地，取 \(\overline{BC}\) 的中点 \(N\)，延长 \(\overline{AN}\) 到
\(F\) 使得 \(A * N * F\) 且 \(\overline{NF} \cong \overline{AN}\)，连接
\(\overline{CF}\)。由 SAS 证明 \(\triangle ABN \cong \triangle FCN\)，得
\(\angle ABN = \angle FCN\)。

与上面类似论证 \(F\) 在 \(\angle ACD\) 内部（\(F\) 和 \(A\) 位于 \(BC\)
的两侧，\(D\) 在 \(BC\) 延长线上），故：

\[\angle ACD > \angle FCD > \angle FCN = \angle ABN = \angle ABC\]

\(\blacksquare\)

\textbf{推论
12.1}（等角对等边）：如果三角形的两个角相等，那么它们所对的边也相等。

\textbf{证明}：设 \(\triangle ABC\) 中 \(\angle B = \angle C\)。如果
\(\overline{AB} \neq \overline{AC}\)，不妨设
\(\overline{AB} > \overline{AC}\)。由外角定理（定理
12.3），\(\angle C > \angle B\)，矛盾。因此
\(\overline{AB} = \overline{AC}\)。\(\blacksquare\)

\paragraph{12.6.4
三角形内角和定理}\label{ux4e09ux89d2ux5f62ux5185ux89d2ux548cux5b9aux7406}

\textbf{引理 12.2}（平行线同位角相等）：设直线 \(a \parallel b\)，截线
\(c\) 分别交 \(a\) 和 \(b\) 于点 \(A\) 和 \(B\)。则 \(c\) 与 \(a\)
所成的同位角相等。

\textbf{证明}：设 \(c\) 交 \(a\) 于 \(A\)，交 \(b\) 于 \(B\)。记 \(a\)
与 \(c\) 在 \(A\) 处的上方同位角为 \(\alpha\)，\(b\) 与 \(c\) 在 \(B\)
处的上方同位角为 \(\beta\)。要证明 \(\alpha = \beta\)。

反证法：假设 \(\alpha \neq \beta\)，不妨设 \(\alpha > \beta\)。在点
\(B\) 处，由公理 C4，在 \(b\) 的上方作射线 \(\overrightarrow{BD}\) 使得
\(\angle CBD = \alpha\)（即 \(c\) 与 \(\overrightarrow{BD}\)
所成的角等于 \(\alpha\)）。设直线 \(b' = \overleftrightarrow{BD}\)。

先证 \(b' \parallel a\)。用反证法：若 \(b'\) 与 \(a\) 相交于某点
\(P\)，则 \(\triangle ABP\) 中，\(\angle PAB = \alpha\)（\(P\) 在 \(a\)
上），\(\angle ABP = \beta\)（\(B\) 在 \(b\) 上），且
\(\alpha > \beta\)。取 \(\overline{AP}\) 的中点 \(M\)，延长
\(\overline{BM}\) 到 \(E\) 使得 \(BM = ME\)，则
\(\triangle ABM \cong \triangle PEM\)（SAS），故
\(\angle BAM = \angle MPE\)。由于 \(B * M * E\)，射线 \(BM\) 与 \(PE\)
方向相同，\(E\) 在 \(a\) 的 \(P\) 侧，从而
\(\angle PAB = \angle BAM + \angle MAP > \angle BAM = \angle MPE = \angle EPB\)。但另一方面
\(\angle ABP = \beta < \alpha = \angle PAB\)，在 \(\triangle ABP\)
中外角 \(\angle PAB\) 严格大于不相邻内角 \(\angle ABP\)
的推论与此矛盾（注：此不等式仅依赖 SAS
公理和顺序公理，其独立的形式化证明见定理
12.3，不依赖平行公理，故不构成循环论证）。

但 \(b \parallel a\) 且 \(b' \parallel a\)，\(b\) 和 \(b'\) 都过点
\(B\)。由平行公理 P，过点 \(B\) 与 \(a\) 平行的直线唯一，故
\(b' = b\)。但 \(b'\) 与 \(b\) 的方向不同（因为
\(\alpha \neq \beta\)），矛盾。

因此 \(\alpha = \beta\)。\(\blacksquare\)

\textbf{推论}（平行线内错角相等）：在上述设定下，\(c\) 与 \(a\)、\(b\)
所成的内错角也相等。这由同位角相等和对顶角相等直接推出。

\textbf{定理
12.4}（三角形内角和定理）：任意三角形的三个内角之和等于两个直角（即
\(180°\)）。

\textbf{证明}：

设 \(\triangle ABC\) 是任意三角形。过点 \(A\) 作直线
\(l \parallel \overleftrightarrow{BC}\)（由平行公理 P 和命题
12.4，这样的直线存在且唯一）。

由于 \(l \parallel \overleftrightarrow{BC}\)，且
\(\overleftrightarrow{AB}\) 是截线，根据引理
12.2（平行线同位角相等），有：

\[\angle BAl = \angle ABC\]

类似地，\(\overleftrightarrow{AC}\) 也是截线，根据引理 12.2，有：

\[\angle CAl = \angle ACB\]

注意 \(l\) 是一条直线，点 \(B\) 和 \(C\) 在 \(l\) 的同一侧（否则 \(l\)
会与 \(\overleftrightarrow{BC}\) 相交）。所以：

\[\angle BAl + \angle BAC + \angle CAl = 180°\]

将上面的关系代入：

\[\angle ABC + \angle BAC + \angle ACB = 180°\]

\(\blacksquare\)

\textbf{推论 12.2}：三角形的任意一个外角等于与它不相邻的两个内角之和。

\textbf{证明}：由定理
12.3（外角定理），三角形的外角大于任一不相邻内角。结合三角形内角和为
\(180°\)（定理 12.4），外角恰好等于两个不相邻内角之和。\(\blacksquare\)

\textbf{推论 12.3}：直角三角形的两个锐角互余。

\textbf{证明}：三角形内角和为 \(180°\)（定理 12.4），直角为
\(90°\)，故其余两角之和为 \(180° - 90° = 90°\)。\(\blacksquare\)

\subsubsection{12.7
公理系统的相容性与独立性}\label{ux516cux7406ux7cfbux7edfux7684ux76f8ux5bb9ux6027ux4e0eux72ecux7acbux6027}

\textbf{相容性}（无矛盾性）：希尔伯特证明了其公理系统是相容的，即从这些公理出发不可能推导出矛盾。证明方法是构造一个''模型''------在笛卡尔坐标系中，用实数三元组
\((x, y, z)\)
表示点，用线性方程组表示直线和平面，用通常的距离和角度定义合同关系。在这个模型中，所有公理都成立，因此公理系统是相容的。

\textbf{独立性}：希尔伯特还证明了各组公理之间的独立性。特别地，平行公理
P
的独立性意味着：前四组公理（关联、顺序、合同、连续）不能推导出平行公理------这正是非欧几何存在的逻辑基础。

\textbf{完备性}：在加上连续公理 D2
之后，希尔伯特公理系统在同构意义下是完备的------即任意两个满足全部公理的模型都是同构的（可以通过坐标系的建立互相转化）。

\bigskip\hrule\bigskip

\textbf{本章展望}：本章建立了欧氏几何的公理基础，确立了点、直线、平面之间的基本关系，以及线段合同和角合同的公理化定义。从这些公理出发，我们推导了线段中点和角平分线的存在唯一性、三角形内角和定理、外角定理等最基本的命题。在下一章中，我们将以此为基础，系统推导平面几何的全部核心定理，包括三角形全等与相似的判定、四边形的性质、圆的几何以及面积理论。

\bigskip\hrule\bigskip

\subsection{Ch 13 平面几何}\label{ch-13-ux5e73ux9762ux51e0ux4f55}

\textbf{前置知识}：本章假定读者已掌握 Ch 12
欧氏几何公理体系，特别是希尔伯特公理系统（关联公理、顺序公理、合同公理、平行公理、连续公理）以及从中推导的基本命题（线段中点和角平分线的存在唯一性、三角形内角和定理、外角定理）。本章的所有定理都将从这些公理和基本命题出发，经过严格的逻辑推导得出。

\bigskip\hrule\bigskip

上一章中，我们建立了欧氏几何的公理系统，并从中推导了若干最基本的几何命题。现在，我们以此为基石，系统地构建平面几何的完整理论。本章涵盖三角形的全等与相似、四边形的性质与分类、圆的核心定理、面积的公理化理论，以及锐角三角函数的初步引入。所有定理均从公理和已证命题出发，经过严格的逻辑推导得出，不依赖任何未经证明的直觉。

\subsubsection{13.1 三角形}\label{ux4e09ux89d2ux5f62}

\paragraph{13.1.1 全等判定}\label{ux5168ux7b49ux5224ux5b9a}

公理 C6 直接给出了 SAS（边角边）全等判定。现在我们从 SAS
出发推导其他全等判定准则。

\textbf{定理 13.1}（ASA
全等判定）：如果两个三角形有两角及其夹边分别相等，则两个三角形全等。

\textbf{证明}：

设 \(\triangle ABC\) 和 \(\triangle A'B'C'\) 满足
\(\angle A \cong \angle A'\)，\(\angle B \cong \angle B'\)，\(\overline{AB} \cong \overline{A'B'}\)。要证明
\(\triangle ABC \cong \triangle A'B'C'\)。

在射线 \(\overrightarrow{A'C'}\) 上取点 \(C''\) 使得
\(\overline{A'C''} \cong \overline{AC}\)。连接 \(\overline{B'C''}\)。

在 \(\triangle ABC\) 和 \(\triangle A'B'C''\) 中： -
\(\overline{AB} \cong \overline{A'B'}\)（已知）， -
\(\angle A \cong \angle A'\)（已知）， -
\(\overline{AC} \cong \overline{A'C''}\)（由构造）。

由 SAS（公理 C6），\(\triangle ABC \cong \triangle A'B'C''\)。因此
\(\angle ABC \cong \angle A'B'C''\)，即
\(\angle B \cong \angle A'B'C''\)。

但已知 \(\angle B \cong \angle B'\)，所以
\(\angle A'B'C'' \cong \angle A'B'C'\)。

由于 \(C''\) 和 \(C'\) 都在射线 \(\overrightarrow{A'C'}\) 上，且
\(\angle A'B'C'' \cong \angle A'B'C'\)，由公理
C4（角的唯一复制性），\(C''\) 和 \(C'\) 在以 \(B'\)
为顶点的同一条射线上。又因为 \(C''\) 和 \(C'\) 都在射线
\(\overrightarrow{A'C'}\) 上，两条不同射线的交点唯一，所以
\(C'' = C'\)。

因此 \(\triangle ABC \cong \triangle A'B'C'\)。\(\blacksquare\)

\textbf{定理 13.2}（AAS
全等判定）：如果两个三角形有两角及其中一角的对边分别相等，则两个三角形全等。

\textbf{证明}：

设 \(\triangle ABC\) 和 \(\triangle A'B'C'\) 满足
\(\angle A \cong \angle A'\)，\(\angle B \cong \angle B'\)，\(\overline{BC} \cong \overline{B'C'}\)（\(\overline{BC}\)
是 \(\angle A\) 的对边）。

由三角形内角和定理（定理
12.4），\(\angle C = 180° - \angle A - \angle B = 180° - \angle A' - \angle B' = \angle C'\)。

现在有
\(\angle B \cong \angle B'\)，\(\angle C \cong \angle C'\)，\(\overline{BC} \cong \overline{B'C'}\)，满足
ASA 条件，由定理
13.1，\(\triangle ABC \cong \triangle A'B'C'\)。\(\blacksquare\)

\textbf{定理 13.3}（SSS
全等判定）：如果两个三角形的三条边分别相等，则两个三角形全等。

\textbf{证明}：

设 \(\triangle ABC\) 和 \(\triangle A'B'C'\) 满足
\(\overline{AB} \cong \overline{A'B'}\)，\(\overline{BC} \cong \overline{B'C'}\)，\(\overline{CA} \cong \overline{C'A'}\)。要证明
\(\triangle ABC \cong \triangle A'B'C'\)，只需证明
\(\angle A \cong \angle A'\)（然后由 SAS 可得全等）。

在 \(B'\) 的另一侧作 \(\angle A'B'C'' \cong \angle ABC\)（由公理
C4），在射线 \(\overrightarrow{B'C''}\) 上取 \(C''\) 使得
\(\overline{B'C''} \cong \overline{BC}\)（由公理 C1）。连接
\(\overline{A'C''}\)。

在 \(\triangle ABC\) 和 \(\triangle A'B'C''\) 中： -
\(\overline{AB} \cong \overline{A'B'}\)（已知）， -
\(\angle ABC \cong \angle A'B'C''\)（由构造）， -
\(\overline{BC} \cong \overline{B'C''}\)（由构造）。

由 SAS（公理 C6），\(\triangle ABC \cong \triangle A'B'C''\)。因此
\(\overline{AC} \cong \overline{A'C''}\)，即
\(\overline{C'A'} \cong \overline{A'C''}\)。

由于 \(C''\) 与 \(C'\) 位于直线 \(A'B'\) 的两侧（由构造），线段
\(C'C''\) 与直线 \(A'B'\) 相交于某点 \(Q\)。

考虑 \(\triangle A'C'C''\)。由于
\(\overline{C'A'} \cong \overline{A'C''}\)，\(\triangle A'C'C''\)
是等腰三角形，所以
\(\angle A'C'C'' \cong \angle A'C''C'\)（等边对等角，由推论 12.1）。

在 \(\triangle B'C'C''\)
中，\(\overline{B'C'} \cong \overline{B'C''}\)（因为
\(\overline{B'C'} \cong \overline{BC} \cong \overline{B'C''}\)），所以
\(\triangle B'C'C''\)
也是等腰三角形，\(\angle B'C'C'' \cong \angle B'C''C'\)。

现在需要根据 \(Q\) 的位置分情况讨论角度分解。

\textbf{情况 (i)}：\(Q\) 在 \(A'\) 和 \(B'\) 之间。此时射线
\(C'Q\)（即射线 \(C'C''\)）在 \(\angle A'C'B'\) 的内部，射线
\(C''Q\)（即射线 \(C''C'\)）在 \(\angle A'C''B'\) 的内部。故：
\[\angle A'C'B' = \angle A'C'Q + \angle QC'B' = \angle A'C'C'' + \angle B'C'C''\]
\[\angle A'C''B' = \angle A'C''Q + \angle QC''B' = \angle A'C''C' + \angle B'C''C'\]
由等腰三角形底角相等：\(\angle A'C'C'' = \angle A'C''C'\)，\(\angle B'C'C'' = \angle B'C''C'\)。因此
\(\angle A'C'B' = \angle A'C''B'\)。

\textbf{情况 (ii)}：\(Q\) 不在 \(A'\) 和 \(B'\) 之间。不妨设 \(A'\) 在
\(Q\) 和 \(B'\) 之间（\(Q * A' * B'\)）。此时射线 \(C'A'\) 在射线
\(C'Q\) 和射线 \(C'B'\) 之间，故
\(\angle B'C'C'' = \angle B'C'A' + \angle A'C'C''\)。类似地在 \(C''\)
处，\(\angle B'C''C' = \angle B'C''A' + \angle A'C''C'\)。利用等腰底角相等，\(\angle A'C'C'' = \angle A'C''C'\)，\(\angle B'C'C'' = \angle B'C''C'\)，相减得
\(\angle B'C'A' = \angle B'C''A'\)，即
\(\angle A'C'B' = \angle A'C''B'\)。

综合两种情况，\(\angle C' = \angle C''\)。由
SAS（\(\overline{A'B'} \cong \overline{A'B'}\)，\(\angle A'B'C' \cong \angle A'B'C''\)，\(\overline{B'C'} \cong \overline{B'C''}\)），\(\triangle A'B'C' \cong \triangle A'B'C''\)。

所以 \(\triangle ABC \cong \triangle A'B'C'\)。\(\blacksquare\)

\paragraph{13.1.2
等腰三角形性质}\label{ux7b49ux8170ux4e09ux89d2ux5f62ux6027ux8d28}

\textbf{定理 13.4}（等边对等角）：等腰三角形的两个底角相等。

\textbf{证明}：

设 \(\triangle ABC\) 中 \(\overline{AB} = \overline{AC}\)。考虑
\(\triangle ABC\) 和 \(\triangle ACB\)： -
\(\overline{AB} \cong \overline{AC}\)（已知）， -
\(\angle A \cong \angle A\)（同一个角）， -
\(\overline{AC} \cong \overline{AB}\)（已知）。

由 SAS（公理 C6），\(\triangle ABC \cong \triangle ACB\)。因此
\(\angle B \cong \angle C\)。\(\blacksquare\)

\textbf{定理
13.5}（等角对等边）：如果三角形的两个角相等，则它们所对的边也相等。

\textbf{证明}：设 \(\triangle ABC\) 中 \(\angle B = \angle C\)。考虑
\(\triangle ABC\) 和 \(\triangle ACB\)： -
\(\angle B \cong \angle C\)（已知）， -
\(\overline{BC} \cong \overline{CB}\)（公共边）， -
\(\angle C \cong \angle B\)（已知）。

由 ASA（定理 13.1），\(\triangle ABC \cong \triangle ACB\)。因此
\(\overline{AB} \cong \overline{AC}\)。\(\blacksquare\)

\textbf{定理
13.6}（三线合一）：等腰三角形的顶角平分线、底边上的高、底边上的中线互相重合。

\textbf{证明}：

设 \(\triangle ABC\) 中
\(\overline{AB} = \overline{AC}\)，\(\overrightarrow{AD}\) 是
\(\angle A\) 的平分线（\(D\) 在 \(\overline{BC}\) 上）。

在 \(\triangle ABD\) 和 \(\triangle ACD\) 中： -
\(\overline{AB} \cong \overline{AC}\)（已知）， -
\(\angle BAD \cong \angle CAD\)（\(\overrightarrow{AD}\) 是角平分线），
- \(\overline{AD} \cong \overline{AD}\)（公共边）。

由 SAS（公理 C6），\(\triangle ABD \cong \triangle ACD\)。因此： -
\(\overline{BD} \cong \overline{CD}\)，即 \(D\) 是 \(\overline{BC}\)
的中点（\(\overrightarrow{AD}\) 是中线）； -
\(\angle ADB \cong \angle ADC\)，而
\(\angle ADB + \angle ADC = 180°\)，所以
\(\angle ADB = \angle ADC = 90°\)（\(\overrightarrow{AD}\) 是高）。

\textbf{逆方向一}（中线 \(\Rightarrow\) 角平分线 + 高）：设
\(\overrightarrow{AD}\) 是 \(\overline{BC}\) 上的中线，即 \(D\) 是
\(\overline{BC}\) 的中点，\(\overline{BD} \cong \overline{CD}\)。

在 \(\triangle ABD\) 和 \(\triangle ACD\) 中： -
\(\overline{AB} \cong \overline{AC}\)（等腰三角形）， -
\(\overline{BD} \cong \overline{CD}\)（\(D\) 是中点）， -
\(\overline{AD} \cong \overline{AD}\)（公共边）。

由 SSS（定理 13.3），\(\triangle ABD \cong \triangle ACD\)。因此： -
\(\angle BAD \cong \angle CAD\)，即 \(\overrightarrow{AD}\) 是角平分线；
- \(\angle ADB \cong \angle ADC\)，而
\(\angle ADB + \angle ADC = 180°\)，所以
\(\angle ADB = \angle ADC = 90°\)，即 \(\overrightarrow{AD}\) 是高。

\textbf{逆方向二}（高 \(\Rightarrow\) 角平分线 + 中线）：设
\(\overrightarrow{AD}\) 是 \(\overline{BC}\) 上的高，即
\(\angle ADB = \angle ADC = 90°\)。

在 \(\triangle ABD\) 和 \(\triangle ACD\) 中： -
\(\overline{AB} \cong \overline{AC}\)（等腰三角形）， -
\(\angle ADB \cong \angle ADC = 90°\)（高）， -
\(\overline{AD} \cong \overline{AD}\)（公共边）。

由 HL
全等（直角三角形中斜边和一直角边相等则全等），\(\triangle ABD \cong \triangle ACD\)。

\textbf{HL 全等的证明}：设 \(\triangle ABD\) 和 \(\triangle ACD\) 中
\(\angle ADB = \angle ADC = 90°\)，\(\overline{AB} \cong \overline{AC}\)（斜边），\(\overline{AD} \cong \overline{AD}\)（公共直角边）。在
\(\overleftrightarrow{AD}\) 的另一侧取点 \(C'\) 使得
\(\overline{DC'} \cong \overline{DB}\) 且 \(\angle ADC' = 90°\)。由
SAS（\(\overline{AD} \cong \overline{AD}\)，\(\angle ADB = \angle ADC' = 90°\)，\(\overline{DB} \cong \overline{DC'}\)），\(\triangle ADB \cong \triangle ADC'\)，故
\(\overline{AC'} \cong \overline{AB} \cong \overline{AC}\)。由于 \(C'\)
和 \(C\) 都在 \(\overleftrightarrow{AD}\) 的同一侧且到 \(A\)、\(D\)
的距离分别相等，由 SSS 唯一性得 \(C' = C\)，从而
\(\triangle ABD \cong \triangle ACD\)。

因此： - \(\angle BAD \cong \angle CAD\)，即 \(\overrightarrow{AD}\)
是角平分线； - \(\overline{BD} \cong \overline{CD}\)，即 \(D\) 是
\(\overline{BC}\) 的中点（\(\overrightarrow{AD}\) 是中线）。

综上，等腰三角形的顶角平分线、底边上的高、底边上的中线三线合一。\(\blacksquare\)

\paragraph{13.1.3 边角关系}\label{ux8fb9ux89d2ux5173ux7cfb}

\textbf{定理 13.7}（大边对大角）：在三角形中，较大的边所对的角较大。

\textbf{证明}：

设 \(\triangle ABC\) 中 \(\overline{AC} > \overline{AB}\)。要证明
\(\angle B > \angle C\)。

在 \(\overline{AC}\) 上取点 \(D\) 使得
\(\overline{AD} = \overline{AB}\)（这是可能的，因为
\(\overline{AC} > \overline{AB}\)）。连接 \(\overline{BD}\)。

在 \(\triangle ABD\) 中，\(\overline{AD} = \overline{AB}\)，所以
\(\angle ABD = \angle ADB\)（等边对等角，定理 13.4）。

由于 \(D\) 在 \(\overline{AC}\) 上（\(A * D * C\)），\(\angle ADB\) 是
\(\triangle BDC\) 的外角，所以 \(\angle ADB > \angle C\)（外角定理，定理
12.3）。

因此 \(\angle ABD > \angle C\)。又因为
\(\angle ABC = \angle ABD + \angle DBC > \angle ABD\)（\(\angle DBC > 0\)），所以
\(\angle ABC > \angle C\)。\(\blacksquare\)

\textbf{定理 13.8}（大角对大边）：在三角形中，较大的角所对的边较大。

\textbf{证明}：这是定理 13.7 的逆否命题。假设
\(\angle B > \angle C\)，如果 \(\overline{AC} \leq \overline{AB}\)： -
若 \(\overline{AC} = \overline{AB}\)，则 \(\angle B = \angle C\)（定理
13.4），矛盾。 - 若 \(\overline{AC} < \overline{AB}\)，由定理
13.7，\(\angle B < \angle C\)，矛盾。

因此 \(\overline{AC} > \overline{AB}\)。\(\blacksquare\)

\textbf{定理 13.9}（三角不等式）：三角形任意两边之和大于第三边。

\textbf{证明}：

设 \(\triangle ABC\)，要证明
\(\overline{AB} + \overline{BC} > \overline{AC}\)。

延长 \(\overline{BA}\) 到 \(A'\) 使得
\(\overline{AA'} = \overline{AC}\)（由公理 C1）。连接
\(\overline{A'C}\)。

在 \(\triangle AA'C\) 中，\(\overline{AA'} = \overline{AC}\)，所以
\(\triangle AA'C\) 是等腰三角形，\(\angle AA'C = \angle ACA'\)（定理
13.4）。

由于 \(B * A * A'\)，\(\angle BCA' < \angle ACA'\)（\(\angle BCA'\) 是
\(\angle ACA'\) 的一部分）。

因此 \(\angle BCA' < \angle BA'C\)。在 \(\triangle BCA'\)
中，由大角对大边（定理 13.8），\(\overline{BC} < \overline{BA'}\)。

而
\(\overline{BA'} = \overline{BA} + \overline{AA'} = \overline{BA} + \overline{AC}\)，所以
\(\overline{BC} < \overline{AB} + \overline{AC}\)，即
\(\overline{AB} + \overline{AC} > \overline{BC}\)。

类似地可以证明其他两个不等式。\(\blacksquare\)

\textbf{推论 13.1}：三角形任意两边之差小于第三边。

\textbf{证明}：设三角形三边为 \(a, b, c\)。由三角不等式（定理
13.9），\(a + b > c\)，故 \(a > c - b\)，即 \(a - c > -b\)。同理
\(c > |a - b|\)，故 \(|a - c| < b\)。\(\blacksquare\)

\subsubsection{13.2 四边形}\label{ux56dbux8fb9ux5f62}

\paragraph{13.2.1
平行四边形的性质和判定}\label{ux5e73ux884cux56dbux8fb9ux5f62ux7684ux6027ux8d28ux548cux5224ux5b9a}

\textbf{定义
13.1}（平行四边形）：两组对边分别平行的四边形称为\textbf{平行四边形}。

\textbf{定理
13.10}（平行四边形的性质）：平行四边形的对边相等，对角相等，对角线互相平分。

\textbf{证明}：

设 \(ABCD\) 是平行四边形，即
\(\overline{AB} \parallel \overline{CD}\)，\(\overline{AD} \parallel \overline{BC}\)。连接对角线
\(\overline{AC}\)。

\textbf{对边相等}：在 \(\triangle ABC\) 和 \(\triangle CDA\) 中： -
\(\angle BAC = \angle DCA\)（\(\overline{AB} \parallel \overline{CD}\)，\(\overline{AC}\)
是截线，内错角相等）， -
\(\angle BCA = \angle DAC\)（\(\overline{AD} \parallel \overline{BC}\)，\(\overline{AC}\)
是截线，内错角相等）， - \(\overline{AC} = \overline{CA}\)（公共边）。

由 ASA（定理 13.1），\(\triangle ABC \cong \triangle CDA\)。因此
\(\overline{AB} = \overline{CD}\)，\(\overline{BC} = \overline{DA}\)。

\textbf{对角相等}：由上面的全等，\(\angle B = \angle D\)。类似地（连接另一条对角线
\(\overline{BD}\)），可证 \(\angle A = \angle C\)。

\textbf{对角线互相平分}：设对角线 \(\overline{AC}\) 和 \(\overline{BD}\)
相交于 \(O\)。

在 \(\triangle AOB\) 和 \(\triangle COD\) 中： -
\(\overline{AB} = \overline{CD}\)（已证）， -
\(\angle OAB = \angle OCD\)（内错角）， -
\(\angle OBA = \angle ODC\)（内错角）。

由 AAS（定理 13.2），\(\triangle AOB \cong \triangle COD\)，所以
\(\overline{OA} = \overline{OC}\)，\(\overline{OB} = \overline{OD}\)。\(\blacksquare\)

\textbf{定理
13.11}（平行四边形的判定）：满足下列条件之一的四边形是平行四边形：

\begin{enumerate}
\def\labelenumi{\arabic{enumi}.}
\tightlist
\item
  两组对边分别相等；
\item
  一组对边平行且相等；
\item
  对角线互相平分；
\item
  两组对角分别相等。
\end{enumerate}

\textbf{证明}（各判定准则的证明）：

\textbf{判定 1}：设四边形 \(ABCD\) 中
\(\overline{AB} = \overline{CD}\)，\(\overline{AD} = \overline{BC}\)。连接
\(\overline{AC}\)。

在 \(\triangle ABC\) 和 \(\triangle CDA\) 中： -
\(\overline{AB} = \overline{CD}\)（已知）， -
\(\overline{BC} = \overline{DA}\)（已知）， -
\(\overline{AC} = \overline{CA}\)（公共边）。

由 SSS（定理 13.3），\(\triangle ABC \cong \triangle CDA\)。因此
\(\angle BAC = \angle DCA\)，\(\angle BCA = \angle DAC\)。由内错角相等，\(\overline{AB} \parallel \overline{CD}\)，\(\overline{AD} \parallel \overline{BC}\)。

\textbf{判定 2}：设 \(\overline{AB} = \overline{CD}\) 且
\(\overline{AB} \parallel \overline{CD}\)。连接 \(\overline{AC}\)。

在 \(\triangle ABC\) 和 \(\triangle CDA\) 中： -
\(\overline{AB} = \overline{CD}\)（已知）， -
\(\angle BAC = \angle DCA\)（\(\overline{AB} \parallel \overline{CD}\)，内错角），
- \(\overline{AC} = \overline{CA}\)（公共边）。

由 SAS，\(\triangle ABC \cong \triangle CDA\)。因此
\(\overline{BC} = \overline{DA}\)，\(\angle BCA = \angle DAC\)，所以
\(\overline{AD} \parallel \overline{BC}\)。

\textbf{判定 3}：设对角线 \(\overline{AC}\) 和 \(\overline{BD}\)
互相平分于 \(O\)，即
\(\overline{OA} = \overline{OC}\)，\(\overline{OB} = \overline{OD}\)。

在 \(\triangle AOB\) 和 \(\triangle COD\) 中： -
\(\overline{OA} = \overline{OC}\)（已知）， -
\(\angle AOB = \angle COD\)（对顶角）， -
\(\overline{OB} = \overline{OD}\)（已知）。

由 SAS，\(\triangle AOB \cong \triangle COD\)。因此
\(\overline{AB} = \overline{CD}\)，\(\angle OAB = \angle OCD\)，所以
\(\overline{AB} \parallel \overline{CD}\)。由判定 2，\(ABCD\)
是平行四边形。

\textbf{判定 4}：设
\(\angle A = \angle C\)，\(\angle B = \angle D\)。四边形内角和为
\(360°\)（可由对角线将四边形分为两个三角形，每个三角形内角和为 \(180°\)
推得）。代入得 \(2\angle A + 2\angle B = 360°\)，即
\(\angle A + \angle B = 180°\)。同旁内角互补，所以
\(\overline{AD} \parallel \overline{BC}\)。类似地，\(\angle A + \angle D = 180°\)，所以
\(\overline{AB} \parallel \overline{CD}\)。\(\blacksquare\)

\paragraph{13.2.2
矩形、菱形、正方形}\label{ux77e9ux5f62ux83f1ux5f62ux6b63ux65b9ux5f62}

\textbf{定义
13.2}（矩形）：有一个角是直角的平行四边形称为\textbf{矩形}。

\textbf{定理 13.12}（矩形的性质）：(1) 矩形的四个角都是直角；(2)
矩形的对角线相等。

\textbf{证明}：设矩形 \(ABCD\) 中 \(\angle A = 90°\)。

\begin{enumerate}
\def\labelenumi{(\arabic{enumi})}
\item
  由平行四边形对角相等，\(\angle C = \angle A = 90°\)。由同旁内角互补，\(\angle B = 180° - \angle A = 90°\)，\(\angle D = 180° - \angle A = 90°\)。
\item
  在 \(\triangle ABC\) 和 \(\triangle DCB\)
  中：\(\overline{AB} = \overline{DC}\)（平行四边形对边相等），\(\angle ABC = \angle DCB = 90°\)，\(\overline{BC} = \overline{CB}\)（公共边）。由
  SAS，\(\triangle ABC \cong \triangle DCB\)，所以
  \(\overline{AC} = \overline{BD}\)。\(\blacksquare\)
\end{enumerate}

\textbf{定义
13.3}（菱形）：有一组邻边相等的平行四边形称为\textbf{菱形}。

\textbf{定理 13.13}（菱形的性质）：(1) 菱形的四条边都相等；(2)
菱形的对角线互相垂直，且每条对角线平分一组对角。

\textbf{证明}：设菱形 \(ABCD\) 中 \(\overline{AB} = \overline{BC}\)。

\begin{enumerate}
\def\labelenumi{(\arabic{enumi})}
\item
  由平行四边形对边相等，\(\overline{CD} = \overline{AB}\)，\(\overline{DA} = \overline{BC}\)。所以四条边都相等。
\item
  连接对角线 \(\overline{AC}\) 和 \(\overline{BD}\)，设交于 \(O\)。在
  \(\triangle ABD\)
  中，\(\overline{AB} = \overline{AD}\)（菱形性质），所以
  \(\triangle ABD\) 是等腰三角形。由三线合一（定理
  13.6），\(\overline{BD}\) 上的中线 \(\overline{AO}\)
  也是高和角平分线。因此 \(\overline{AC} \perp \overline{BD}\)，且
  \(\overline{BD}\) 平分 \(\angle ABC\) 和
  \(\angle ADC\)。类似地，\(\overline{AC}\) 平分 \(\angle BAD\) 和
  \(\angle BCD\)。\(\blacksquare\)
\end{enumerate}

\textbf{定义
13.4}（正方形）：既是矩形又是菱形的四边形称为\textbf{正方形}。

\textbf{定理 13.14}：正方形具有矩形和菱形的所有性质。

\textbf{证明}：正方形定义为四边相等且四角均为直角的四边形。四角为直角说明是矩形；四边相等说明是菱形。故正方形同时具有矩形和菱形的所有性质。\(\blacksquare\)

\paragraph{13.2.3 梯形}\label{ux68afux5f62}

\textbf{定义
13.5}（梯形）：有一组对边平行而另一组对边不平行的四边形称为\textbf{梯形}。平行的两边称为\textbf{底}，不平行的两边称为\textbf{腰}。

\textbf{定义 13.6}（等腰梯形）：两腰相等的梯形称为\textbf{等腰梯形}。

\textbf{定理
13.15}（等腰梯形的性质）：等腰梯形的两个底角相等，对角线相等。

\textbf{证明}：

设等腰梯形 \(ABCD\) 中
\(\overline{AB} \parallel \overline{CD}\)，\(\overline{AD} = \overline{BC}\)。

过 \(D\) 作 \(\overline{DE} \parallel \overline{BC}\) 交
\(\overline{AB}\) 于 \(E\)，则 \(BCDE\)
是平行四边形（\(BC \parallel DE\) 由构造，\(CD \parallel BE\) 因为
\(CD \parallel AB\) 且 \(E \in AB\)），故
\(\overline{DE} = \overline{BC}\)。

由于
\(\overline{AD} = \overline{BC}\)，\(\overline{DE} = \overline{AD}\)，所以
\(\triangle ADE\) 是等腰三角形，\(\angle DAE = \angle DEA\)。

由 \(\overline{DE} \parallel \overline{BC}\) 和 \(\overline{AB}\)
是截线，\(\angle DEA = \angle B\)（同位角）。由
\(\overline{AB} \parallel \overline{CD}\) 和 \(\overline{AD}\)
是截线，\(\angle DAE = \angle ADC\)（内错角）。

因此 \(\angle ADC = \angle B\)。类似地可证 \(\angle DAB = \angle C\)。

对于对角线相等：连接 \(\overline{AC}\) 和 \(\overline{BD}\)。在
\(\triangle ABD\) 和 \(\triangle BAC\)
中，\(\overline{AB} = \overline{BA}\)，\(\overline{AD} = \overline{BC}\)，\(\angle DAB = \angle CBA\)（已证）。由
SAS，\(\triangle ABD \cong \triangle BAC\)，所以
\(\overline{BD} = \overline{AC}\)。\(\blacksquare\)

\paragraph{13.2.4 中位线定理}\label{ux4e2dux4f4dux7ebfux5b9aux7406}

\textbf{定理
13.16}（三角形中位线定理）：三角形的中位线平行于第三边且等于第三边的一半。

\textbf{证明}：

设 \(\triangle ABC\) 中，\(D\)、\(E\) 分别是
\(\overline{AB}\)、\(\overline{AC}\) 的中点。要证明
\(\overline{DE} \parallel \overline{BC}\) 且
\(\overline{DE} = \frac{1}{2}\overline{BC}\)。

延长 \(\overline{DE}\) 到 \(F\) 使得
\(\overline{EF} = \overline{DE}\)。连接 \(\overline{CF}\)。

在 \(\triangle ADE\) 和 \(\triangle CFE\) 中： -
\(\overline{AE} = \overline{CE}\)（\(E\) 是 \(\overline{AC}\) 的中点），
- \(\angle AED = \angle CEF\)（对顶角）， -
\(\overline{DE} = \overline{FE}\)（由构造）。

由 SAS，\(\triangle ADE \cong \triangle CFE\)。因此
\(\overline{AD} = \overline{CF}\)，\(\angle ADE = \angle CFE\)。

由内错角相等，\(\overline{AD} \parallel \overline{CF}\)，即
\(\overline{DB} \parallel \overline{CF}\)。

又 \(\overline{DB} = \overline{AD} = \overline{CF}\)，所以 \(DBCF\)
是平行四边形（一组对边平行且相等）。

因此 \(\overline{DF} \parallel \overline{BC}\) 且
\(\overline{DF} = \overline{BC}\)。

而
\(\overline{DF} = \overline{DE} + \overline{EF} = 2\overline{DE}\)，所以
\(\overline{DE} = \frac{1}{2}\overline{BC}\)，且
\(\overline{DE} \parallel \overline{BC}\)。\(\blacksquare\)

\textbf{定理
13.17}（梯形中位线定理）：梯形的中位线平行于两底且等于两底和的一半。

\textbf{证明}：

设梯形 \(ABCD\) 中
\(\overline{AB} \parallel \overline{CD}\)，\(E\)、\(F\) 分别是
\(\overline{AD}\)、\(\overline{BC}\) 的中点。连接 \(\overline{AF}\)
并延长交 \(\overline{DC}\) 的延长线于 \(G\)。

在 \(\triangle ABF\) 和 \(\triangle GCF\) 中： -
\(\overline{BF} = \overline{CF}\)（\(F\) 是 \(\overline{BC}\) 的中点），
-
\(\angle ABF = \angle GCF\)（\(\overline{AB} \parallel \overline{DC}\)，内错角），
- \(\angle AFB = \angle GFC\)（对顶角）。

由 ASA，\(\triangle ABF \cong \triangle GCF\)。因此
\(\overline{AF} = \overline{FG}\)，\(\overline{AB} = \overline{CG}\)。

在 \(\triangle ADG\) 中，\(E\) 是 \(\overline{AD}\) 的中点，\(F\) 是
\(\overline{AG}\) 的中点，所以 \(\overline{EF}\) 是 \(\triangle ADG\)
的中位线。

由三角形中位线定理（定理
13.16），\(\overline{EF} \parallel \overline{DG}\) 且
\(\overline{EF} = \frac{1}{2}\overline{DG}\)。

而
\(\overline{DG} = \overline{DC} + \overline{CG} = \overline{DC} + \overline{AB}\)，所以
\(\overline{EF} = \frac{1}{2}(\overline{AB} + \overline{CD})\)。

由于 \(\overline{DG} \parallel \overline{AB}\)，所以
\(\overline{EF} \parallel \overline{AB}\)。\(\blacksquare\)

\subsubsection{13.3 相似}\label{ux76f8ux4f3c}

\paragraph{13.3.1
平行线分线段成比例}\label{ux5e73ux884cux7ebfux5206ux7ebfux6bb5ux6210ux6bd4ux4f8b}

\textbf{定理
13.18}（平行线分线段成比例定理）：如果一组平行线被两条直线所截，那么截得的对应线段成比例。

\textbf{证明}：

设三条平行线 \(l_1 \parallel l_2 \parallel l_3\)，直线 \(a\) 分别交
\(l_1\)、\(l_2\)、\(l_3\) 于 \(A\)、\(B\)、\(C\)，直线 \(b\) 分别交
\(l_1\)、\(l_2\)、\(l_3\) 于 \(D\)、\(E\)、\(F\)。要证明
\(\frac{AB}{BC} = \frac{DE}{EF}\)。

\textbf{情形 1}（可公度情况）：\(\overline{AB}\) 和 \(\overline{BC}\)
是可公度的，即存在线段 \(\overline{MN}\) 使得
\(\overline{AB} = m \cdot \overline{MN}\)，\(\overline{BC} = n \cdot \overline{MN}\)（\(m\)、\(n\)
是正整数）。

将 \(\overline{AB}\) 分为 \(m\) 等份：在射线 \(\overrightarrow{AB}\)
上依次取点 \(A = A_0, A_1, \ldots, A_m = B\) 使得
\(\overline{A_{i-1}A_i} \cong \overline{MN}\)（\(i = 1, \ldots, m\)）。类似地将
\(\overline{BC}\) 分为 \(n\)
等份：\(B = B_0, B_1, \ldots, B_n = C\)。过每个分点
\(A_i\)（\(i = 1, \ldots, m-1\)）和
\(B_j\)（\(j = 1, \ldots, n-1\)）作与 \(l_1\)
平行的直线。这些辅助平行线交直线 \(b\) 于点
\(D = E_0, E_1, \ldots, E_m\) 和 \(E_m = F_0, F_1, \ldots, F_n = F\)。

关键步骤：证明各段
\(\overline{E_{i-1}E_i}\)（\(i = 1, \ldots, m\)）长度相等。

对每个 \(i\)，过 \(E_{i-1}\) 作平行于直线 \(a\) 的直线，交辅助平行线
\(\ell_i\)（过 \(A_i\)、\(E_i\) 且平行于 \(l_1\) 的直线）于
\(P_i\)。则四边形 \(A_{i-1}A_i P_i E_{i-1}\)
是平行四边形：\(\overline{A_{i-1}A_i} \parallel \overline{E_{i-1}P_i}\)（均沿直线
\(a\)
的方向），\(\overline{A_{i-1}E_{i-1}} \parallel \overline{A_i P_i}\)（分别在平行线
\(\ell_{i-1}\)、\(\ell_i\) 上）。由平行四边形对边相等（定理
13.10），\(\overline{E_{i-1}P_i} = \overline{A_{i-1}A_i} = \overline{MN}\)。

\textbf{当 \(a \parallel b\) 时}：过 \(E_{i-1}\) 平行于 \(a\) 的直线即为
\(b\) 本身，故 \(P_i = E_i\)，从而
\(\overline{E_{i-1}E_i} = \overline{MN}\)。

\textbf{当 \(a\) 与 \(b\) 相交时}：在 \(\triangle E_{i-1}P_i E_i\)
中，三个要素均与 \(i\) 无关：

\begin{itemize}
\tightlist
\item
  边 \(\overline{E_{i-1}P_i} = \overline{MN}\)；
\item
  \(\angle P_i E_{i-1} E_i\) 是直线 \(a\) 的方向与直线 \(b\)
  的夹角（因为
  \(\overline{E_{i-1}P_i} \parallel a\)，\(\overline{E_{i-1}E_i}\) 在
  \(b\) 上），与 \(i\) 无关；
\item
  \(\angle E_{i-1}P_i E_i\) 是直线 \(a\) 的方向与 \(l_1\) 的夹角（因为
  \(\overline{E_{i-1}P_i} \parallel a\)，\(\overline{P_i E_i}\) 在
  \(\ell_i \parallel l_1\) 上），与 \(i\) 无关。
\end{itemize}

由 ASA，所有
\(\triangle E_{i-1}P_i E_i\)（\(i = 1, \ldots, m\)）全等，故
\(\overline{E_{i-1}E_i}\) 对所有 \(i\) 相等。

设各段的公共长度为 \(\delta\)，则
\(\overline{DE} = m \cdot \delta\)。同理，过
\(B_j\)（\(j = 1, \ldots, n-1\)）的辅助平行线将 \(\overline{EF}\) 分为
\(n\) 段，由完全相同的论证（\(MN\)、夹角均不变），各段长度也是
\(\delta\)，故 \(\overline{EF} = n \cdot \delta\)。因此
\(\frac{AB}{BC} = \frac{m}{n} = \frac{DE}{EF}\)。

\textbf{情形 2}（一般情况）：若 \(\overline{AB}\) 与 \(\overline{BC}\)
不可公度，我们用阿基米德公理和反证法。假设
\(\frac{AB}{BC} \neq \frac{DE}{EF}\)，不妨设
\(\frac{AB}{BC} > \frac{DE}{EF}\)。

由阿基米德公理（公理 D1），存在正整数 \(m\)、\(n\) 使得
\(\frac{AB}{BC} > \frac{m}{n} > \frac{DE}{EF}\)（实数的阿基米德性质保证这样的有理数存在）。在射线
\(\overrightarrow{AB}\) 上取点 \(P\) 使得
\(\overline{AP} = m \cdot \overline{MN}\)，在射线
\(\overrightarrow{BC}\) 上取点 \(Q\) 使得
\(\overline{BQ} = n \cdot \overline{MN}\)（其中 \(\overline{MN}\)
是某个单位线段）。由于 \(\frac{m}{n} < \frac{AB}{BC}\)，点 \(P\) 落在
\(A\) 和 \(B\) 之间（即 \(A * P * B\)），点 \(Q\) 落在 \(B\) 和 \(C\)
之间（即 \(B * Q * C\)），且 \(\overline{AP}\) 与 \(\overline{BQ}\)
可公度。

过 \(P\) 作 \(l_2' \parallel l_1\) 交 \(b\) 于 \(E'\)。由情形 1
的结果，\(\frac{AP}{BQ} = \frac{DE'}{E'F}\)。因为
\(\frac{AP}{BQ} = \frac{m}{n} < \frac{AB}{BC} = \frac{AB}{BC}\)，而
\(\frac{m}{n} > \frac{DE}{EF}\)，通过比较可得 \(E'\) 落在 \(D\) 和 \(E\)
之间（即 \(D * E' * E\)），从而 \(\overline{DE'} < \overline{DE}\)。

但类似地，取更大的 \(m'/n'\) 使得
\(\frac{AB}{BC} > \frac{m'}{n'} > \frac{m}{n}\)，过对应分点作平行线得到
\(E''\)，有
\(\overline{DE''} > \overline{DE'}\)。反复使用阿基米德公理，总能找到更接近
\(\frac{AB}{BC}\) 的比值 \(\frac{m'}{n'}\)，使得对应的 \(E^{(k)}\)
序列收敛。由完备公理 D2（直线上的点不能有''空隙''），存在唯一的点
\(E^*\) 使得 \(\frac{AB}{BC} = \frac{DE^*}{E^*F}\)。但 \(l_2\) 过 \(B\)
且平行于 \(l_1\)，由平行公理 P 唯一确定 \(l_2\) 与 \(b\) 的交点
\(E\)，故 \(E^* = E\)，从而
\(\frac{AB}{BC} = \frac{DE}{EF}\)，与假设矛盾。

因此 \(\frac{AB}{BC} = \frac{DE}{EF}\)。\(\blacksquare\)

\textbf{推论
13.2}（三角形一边的平行线性质）：如果一条直线平行于三角形的一边并与另外两边相交，那么它截得的线段与对应边成比例。

\textbf{证明}：由定理 13.18（平行线分线段成比例），直线平行于
\(\triangle ABC\) 的边 \(BC\) 并与 \(AB\)、\(AC\) 分别交于
\(D\)、\(E\)，则 \(\frac{AD}{AB} = \frac{AE}{AC}\)。\(\blacksquare\)

\paragraph{13.3.2
相似三角形判定}\label{ux76f8ux4f3cux4e09ux89d2ux5f62ux5224ux5b9a}

\textbf{定义
13.7}（相似）：如果两个三角形的对应角相等，对应边成比例，则称这两个三角形\textbf{相似}，记为
\(\triangle ABC \sim \triangle A'B'C'\)。

\textbf{定理 13.19}（AA
相似判定）：如果两个三角形有两组对应角相等，则这两个三角形相似。

\textbf{证明}：

设 \(\triangle ABC\) 和 \(\triangle A'B'C'\) 中
\(\angle A = \angle A'\)，\(\angle B = \angle B'\)。由三角形内角和定理，\(\angle C = \angle C'\)。所以三个角都对应相等。

现在证明对应边成比例。在 \(\overline{AB}\) 上取点 \(B''\) 使得
\(\overline{AB''} = \overline{A'B'}\)，过 \(B''\) 作
\(\overline{B''C''} \parallel \overline{BC}\) 交 \(\overline{AC}\) 于
\(C''\)。

由推论 13.2，\(\frac{AB''}{AB} = \frac{AC''}{AC}\)。

在 \(\triangle AB''C''\) 和 \(\triangle A'B'C'\) 中： -
\(\overline{AB''} = \overline{A'B'}\)（由构造）， -
\(\angle A = \angle A'\)（已知）， -
\(\angle AB''C'' = \angle B = \angle B'\)（\(B''C'' \parallel BC\)，同位角）。

由 ASA，\(\triangle AB''C'' \cong \triangle A'B'C'\)。因此
\(\overline{AC''} = \overline{A'C'}\)，\(\overline{B''C''} = \overline{B'C'}\)。

代入比例式：\(\frac{A'B'}{AB} = \frac{A'C'}{AC}\)。类似地可以证明
\(\frac{B'C'}{BC} = \frac{A'B'}{AB}\)。

所以 \(\frac{A'B'}{AB} = \frac{B'C'}{BC} = \frac{A'C'}{AC}\)，即
\(\triangle ABC \sim \triangle A'B'C'\)。\(\blacksquare\)

\textbf{定理 13.20}（SAS
相似判定）：如果两个三角形有两组对应边成比例且夹角相等，则这两个三角形相似。

\textbf{证明}：设 \(\triangle ABC\) 和 \(\triangle A'B'C'\) 中
\(\frac{AB}{A'B'} = \frac{AC}{A'C'}\) 且 \(\angle A = \angle A'\)。在
\(\overline{AB}\) 上取 \(B''\) 使得
\(\overline{AB''} = \overline{A'B'}\)，过 \(B''\) 作
\(\overline{B''C''} \parallel \overline{BC}\) 交 \(\overline{AC}\) 于
\(C''\)。

由推论 13.2，\(\frac{AB''}{AB} = \frac{AC''}{AC}\)，即
\(\frac{A'B'}{AB} = \frac{AC''}{AC}\)。又已知
\(\frac{A'B'}{AB} = \frac{A'C'}{AC}\)，因此
\(\overline{AC''} = \overline{A'C'}\)。

在 \(\triangle AB''C''\) 和 \(\triangle A'B'C'\)
中：\(\overline{AB''} = \overline{A'B'}\)，\(\angle A = \angle A'\)，\(\overline{AC''} = \overline{A'C'}\)。由
SAS，\(\triangle AB''C'' \cong \triangle A'B'C'\)。由 AA
判定，\(\triangle AB''C'' \sim \triangle ABC\)。由相似的传递性，\(\triangle A'B'C' \sim \triangle ABC\)。\(\blacksquare\)

\textbf{定理 13.21}（SSS
相似判定）：如果两个三角形的三组对应边成比例，则这两个三角形相似。

\textbf{证明}：设 \(\triangle ABC\) 和 \(\triangle A'B'C'\) 中
\(\frac{AB}{A'B'} = \frac{BC}{B'C'} = \frac{CA}{C'A'} = k\)。在
\(\overline{AB}\) 上取 \(B''\) 使得
\(\overline{AB''} = \overline{A'B'}\)，过 \(B''\) 作
\(\overline{B''C''} \parallel \overline{BC}\) 交 \(\overline{AC}\) 于
\(C''\)。

由推论
13.2，\(\frac{AB''}{AB} = \frac{AC''}{AC} = \frac{B''C''}{BC}\)。因为
\(\overline{AB''} = \overline{A'B'}\)，\(\frac{A'B'}{AB} = \frac{1}{k}\)，所以
\(\overline{AC''} = \overline{A'C'}\)，\(\overline{B''C''} = \overline{B'C'}\)。

由 SAS，\(\triangle AB''C'' \cong \triangle A'B'C'\)。由 AA
判定，\(\triangle AB''C'' \sim \triangle ABC\)。由传递性，\(\triangle A'B'C' \sim \triangle ABC\)。\(\blacksquare\)

\paragraph{13.3.3 勾股定理}\label{ux52feux80a1ux5b9aux7406}

\textbf{定理
13.22}（勾股定理）：直角三角形中，两条直角边的平方和等于斜边的平方。即如果
\(\triangle ABC\) 中 \(\angle C = 90°\)，则 \(a^2 + b^2 = c^2\)（其中
\(a = BC\)，\(b = AC\)，\(c = AB\)）。

\textbf{证明}（相似三角形法）：

设 \(\triangle ABC\) 中 \(\angle C = 90°\)。从直角顶点 \(C\) 向斜边
\(\overline{AB}\) 作高 \(\overline{CD}\)，设垂足为 \(D\)。

在 \(\triangle ACD\) 和 \(\triangle ABC\) 中：\(\angle A\)
公共，\(\angle ADC = \angle ACB = 90°\)。由 AA
相似，\(\triangle ACD \sim \triangle ABC\)。因此
\(\frac{AC}{AB} = \frac{AD}{AC}\)，即 \(AC^2 = AD \cdot AB\)。

在 \(\triangle CBD\) 和 \(\triangle ABC\) 中：\(\angle B\)
公共，\(\angle BDC = \angle ACB = 90°\)。由 AA
相似，\(\triangle CBD \sim \triangle ABC\)。因此
\(\frac{BC}{AB} = \frac{BD}{BC}\)，即 \(BC^2 = BD \cdot AB\)。

两式相加：

\[AC^2 + BC^2 = AD \cdot AB + BD \cdot AB = (AD + BD) \cdot AB = AB \cdot AB = AB^2\]

即 \(a^2 + b^2 = c^2\)。\(\blacksquare\)

\textbf{推论 13.3}（勾股定理的逆定理）：如果三角形的三边满足
\(a^2 + b^2 = c^2\)，则该三角形是直角三角形，且 \(c\) 所对的角为直角。

\textbf{证明}：设 \(\triangle ABC\) 中
\(AB = c\)，\(BC = a\)，\(AC = b\)，且 \(a^2 + b^2 = c^2\)。以 \(BC\) 和
\(AC\) 为直角边构造直角三角形 \(A'BC\)，使得
\(\angle BCA' = 90°\)，\(CA' = b\)。由勾股定理，\(BA'^2 = a^2 + b^2 = c^2\)，所以
\(BA' = c = BA\)。

在 \(\triangle ABC\) 和 \(\triangle A'BC\)
中：\(AB = A'B\)，\(BC = BC\)，\(AC = A'C\)。由
SSS，\(\triangle ABC \cong \triangle A'BC\)，因此
\(\angle BCA = \angle BCA' = 90°\)。\(\blacksquare\)

\subsubsection{13.4 圆}\label{ux5706}

\paragraph{13.4.1
圆的基本性质}\label{ux5706ux7684ux57faux672cux6027ux8d28}

\textbf{定义 13.8}（圆）：平面上到定点 \(O\) 的距离等于定长 \(r\)
的所有点的集合称为\textbf{圆}，记为 \(\odot(O, r)\)。定点 \(O\)
称为\textbf{圆心}，定长 \(r\) 称为\textbf{半径}。

\textbf{定义
13.9}（弦、直径）：连接圆上两点的线段称为\textbf{弦}。经过圆心的弦称为\textbf{直径}，直径的长度为
\(2r\)。

\textbf{定理 13.23}（弦的中垂线定理）：弦的中垂线经过圆心。

\textbf{证明}：

设 \(AB\) 是圆 \(\odot(O, r)\) 的一条弦，\(M\) 是 \(AB\) 的中点。在
\(\triangle OAM\) 和 \(\triangle OBM\) 中： -
\(\overline{OA} = \overline{OB} = r\)（半径）， -
\(\overline{AM} = \overline{BM}\)（\(M\) 是中点）， -
\(\overline{OM} = \overline{OM}\)（公共边）。

由 SSS，\(\triangle OAM \cong \triangle OBM\)，所以
\(\angle OMA = \angle OMB\)。由于 \(\angle OMA + \angle OMB = 180°\)，得
\(\angle OMA = \angle OMB = 90°\)，即 \(OM \perp AB\)。\(\blacksquare\)

\textbf{定理
13.24}（等弦等距）：在同圆或等圆中，相等的弦到圆心的距离相等；反之，到圆心距离相等的弦相等。

\textbf{证明}：

设 \(AB\) 和 \(CD\) 是圆 \(\odot(O, r)\) 的两条弦，\(M\)、\(N\) 分别是
\(AB\)、\(CD\) 的中点。则 \(OM \perp AB\)，\(ON \perp CD\)（定理
13.23）。

若 \(AB = CD\)，则 \(AM = CN = \frac{1}{2}AB\)。在 \(\triangle OAM\) 和
\(\triangle OCN\)
中：\(OA = OC = r\)，\(AM = CN\)，\(\angle OMA = \angle ONC = 90°\)。由勾股定理，\(OM = \sqrt{OA^2 - AM^2} = \sqrt{OC^2 - CN^2} = ON\)。

反之，若 \(OM = ON\)，则
\(AM = \sqrt{OA^2 - OM^2} = \sqrt{OC^2 - ON^2} = CN\)，故
\(AB = 2AM = 2CN = CD\)。\(\blacksquare\)

\paragraph{13.4.2
圆心角与圆周角}\label{ux5706ux5fc3ux89d2ux4e0eux5706ux5468ux89d2}

\textbf{定义 13.10}（圆心角）：顶点在圆心的角称为\textbf{圆心角}。

\textbf{定义
13.11}（圆周角）：顶点在圆上，两边都是弦的角称为\textbf{圆周角}。

\textbf{定理
13.25}（圆周角定理）：圆周角的度数等于同弧所对圆心角的一半。

\textbf{证明}：

设 \(\odot(O, r)\) 中，弧 \(AB\) 所对的圆心角为 \(\angle AOB\)，圆周角为
\(\angle ACB\)（\(C\) 在圆上且在弧 \(AB\) 所在的一侧）。

\textbf{情形 1}：圆心 \(O\) 在圆周角 \(\angle ACB\) 的一条边上。不妨设
\(O\) 在 \(\overrightarrow{CB}\) 上，即 \(CB\) 是直径。

在 \(\triangle OAC\)
中，\(OA = OC = r\)（等腰三角形），\(\angle OAC = \angle OCA\)（等边对等角）。又
\(\angle AOB\) 是 \(\triangle OAC\) 的外角，所以
\(\angle AOB = \angle OAC + \angle OCA = 2\angle OCA = 2\angle ACB\)。

\textbf{情形 2}：圆心 \(O\) 在圆周角 \(\angle ACB\) 的内部。过 \(C\)
作直径 \(CD\)。由情形
1，\(\angle ACD = \frac{1}{2}\angle AOD\)，\(\angle DCB = \frac{1}{2}\angle DOB\)。因此
\(\angle ACB = \angle ACD + \angle DCB = \frac{1}{2}(\angle AOD + \angle DOB) = \frac{1}{2}\angle AOB\)。

\textbf{情形 3}：圆心 \(O\) 在圆周角 \(\angle ACB\)
的外部。类似地，\(\angle ACB = \frac{1}{2}\angle AOB\)（此时为两个差的组合）。

综合三种情形，\(\angle ACB = \frac{1}{2}\angle AOB\)。\(\blacksquare\)

\textbf{推论 13.4}：同弧所对的圆周角相等。

\textbf{证明}：由定理
13.25（圆周角定理），同弧所对的圆周角等于该弧所对圆心角的一半。因同弧对应的圆心角唯一，故所有圆周角相等。\(\blacksquare\)

\textbf{推论 13.5}：半圆所对的圆周角是直角。

\textbf{证明}：半圆所对的圆心角为 \(180°\)，故圆周角为
\(\frac{1}{2} \times 180° = 90°\)。\(\blacksquare\)

\paragraph{13.4.3 切线}\label{ux5207ux7ebf}

\textbf{定义
13.12}（切线）：与圆只有一个公共点的直线称为圆的\textbf{切线}，公共点称为\textbf{切点}。

\textbf{定理
13.26}（切线的判定）：过圆上一点且垂直于半径的直线是圆的切线。

\textbf{证明}：

设 \(\odot(O, r)\)，\(P\) 在圆上，直线 \(l\) 过 \(P\) 且
\(l \perp OP\)。设 \(Q\) 是 \(l\) 上除 \(P\) 外的任意一点。在
\(\triangle OPQ\) 中，\(\angle OPQ = 90°\)，所以
\(OQ^2 = OP^2 + PQ^2 > OP^2 = r^2\)，即 \(OQ > r\)。因此 \(Q\)
在圆外。所以 \(l\) 与圆只有一个公共点 \(P\)，即 \(l\)
是切线。\(\blacksquare\)

\textbf{定理 13.27}（切线的性质）：圆的切线垂直于过切点的半径。

\textbf{证明}：设 \(l\) 切 \(\odot(O, r)\) 于 \(P\)。则 \(P\) 是 \(l\)
与圆的唯一公共点。设 \(H\) 是 \(O\) 到 \(l\) 的垂足。若 \(H \neq P\)，在
\(l\) 上取 \(P\) 关于 \(H\) 的对称点 \(P'\)，则 \(OP' = OP = r\)（由
\(OH \perp PP'\) 和 \(HP = HP'\) 可证），即 \(P'\) 也在圆上。但 \(l\)
与圆只有一个公共点，矛盾。因此 \(H = P\)，即
\(OP \perp l\)。\(\blacksquare\)

\paragraph{13.4.4 圆幂定理}\label{ux5706ux5e42ux5b9aux7406}

\textbf{定理 13.28}（相交弦定理）：设圆 \(O\) 的两条弦 \(AB\) 和 \(CD\)
相交于点 \(P\)，则 \(PA \cdot PB = PC \cdot PD\)。

\textbf{证明}：

连接 \(AC\) 和 \(BD\)。在 \(\triangle PAC\) 和 \(\triangle PDB\) 中： -
\(\angle APC = \angle DPB\)（对顶角）， -
\(\angle PAC = \angle PDB\)（同弧 \(BC\) 所对的圆周角相等）。

由 AA 相似，\(\triangle PAC \sim \triangle PDB\)。因此
\(\frac{PA}{PD} = \frac{PC}{PB}\)，即
\(PA \cdot PB = PC \cdot PD\)。\(\blacksquare\)

\textbf{引理 13.4（弦切角定理）}：设圆 \(O\) 的切线在点 \(T\)
处与圆相切，\(TA\) 是圆的一条弦。则弦切角 \(\angle PTA\)（\(P\)
在切线上，\(T\) 为顶点，\(A\) 在圆上）等于同弧 \(\widehat{TA}\)
所对的圆周角。

\textbf{证明}：连接 \(OT\)（半径）。由于
\(OT \perp PT\)（切线垂直于半径），\(\angle OTP = 90°\)。连接 \(OA\)，设
\(\angle OTA = \alpha\)。则 \(\angle PTA = 90° - \alpha\)。在
\(\triangle OTA\) 中，\(OT = OA\)（半径），故 \(\triangle OTA\)
是等腰三角形，\(\angle OAT = \angle OTA = \alpha\)。设
\(\angle AOT = \beta\)，则
\(\beta = 180° - 2\alpha\)。由圆周角定理（定理 13.25），弧
\(\widehat{TA}\) 所对的圆周角为
\(\frac{\beta}{2} = 90° - \alpha = \angle PTA\)。\(\blacksquare\)

\textbf{定理 13.29}（切割线定理）：设 \(P\) 是圆外一点，\(PT\)
是切线（\(T\) 为切点），\(PAB\) 是割线（\(A\)、\(B\) 为交点），则
\(PT^2 = PA \cdot PB\)。

\textbf{证明}：

连接 \(TA\) 和 \(TB\)。在 \(\triangle PTA\) 和 \(\triangle PTB\) 中： -
\(\angle TPA = \angle BPT\)（公共角）， -
\(\angle PTB = \angle PAT\)（弦切角等于同弧所对的圆周角）。

由 AA 相似，\(\triangle PTA \sim \triangle PTB\)。因此
\(\frac{PT}{PB} = \frac{PA}{PT}\)，即
\(PT^2 = PA \cdot PB\)。\(\blacksquare\)

\subsubsection{13.5 面积}\label{ux9762ux79ef}

\paragraph{13.5.1
面积的公理化定义}\label{ux9762ux79efux7684ux516cux7406ux5316ux5b9aux4e49}

\textbf{公理 S1}：每一个多边形 \(P\) 都对应一个非负实数
\(\text{area}(P)\)，称为 \(P\) 的\textbf{面积}。

\textbf{公理 S2}：如果多边形 \(P\) 和多边形 \(Q\) 全等，则
\(\text{area}(P) = \text{area}(Q)\)。

\textbf{公理 S3}（可加性）：如果多边形 \(P\) 被分割为两个多边形 \(P_1\)
和 \(P_2\)（它们只在公共边上共享点），则
\(\text{area}(P) = \text{area}(P_1) + \text{area}(P_2)\)。

\textbf{公理 S4}（单位正方形）：边长为 \(1\) 的正方形的面积为 \(1\)。

\paragraph{13.5.2
平行四边形和三角形的面积}\label{ux5e73ux884cux56dbux8fb9ux5f62ux548cux4e09ux89d2ux5f62ux7684ux9762ux79ef}

\textbf{定理 13.30}：平行四边形的面积等于底乘以高，即
\(\text{area} = bh\)。

\textbf{证明}：

设平行四边形
\(ABCD\)，\(\overline{AB} \parallel \overline{CD}\)，\(\overline{AB} = b\)，高为
\(h\)（两平行边之间的距离）。

过 \(D\) 作 \(\overline{DE} \perp \overline{AB}\) 于 \(E\)，过 \(C\) 作
\(\overline{CF} \perp \overline{AB}\) 于 \(F\)。

在 \(\triangle ADE\) 和 \(\triangle BCF\)
中：\(\overline{AD} = \overline{BC}\)（平行四边形对边相等），\(\angle AED = \angle BFC = 90°\)，\(\angle DAE = \angle CBF\)（同位角）。由
AAS，\(\triangle ADE \cong \triangle BCF\)。

因此，平行四边形 \(ABCD\) 的面积 = 矩形 \(DEFC\) 的面积 +
\(\triangle ADE\) 的面积 + \(\triangle BCF\) 的面积 \(-\)
\(\triangle ADE\) 的面积 \(-\) \(\triangle BCF\) 的面积 \(+\)
\(\triangle ADE\) 的面积 \(+\) \(\triangle BCF\) 的面积。更简洁地：将
\(\triangle ADE\) 移到 \(\triangle BCF\) 的位置，平行四边形 \(ABCD\)
变为矩形 \(DEFC\)。由公理 S2 和
S3，\(\text{area}(ABCD) = \text{area}(DEFC) = bh\)。\(\blacksquare\)

\textbf{定理 13.31}：三角形的面积等于底乘以高的一半，即
\(\text{area} = \frac{1}{2}bh\)。

\textbf{证明}：

设 \(\triangle ABC\)，以 \(\overline{BC}\) 为底，高为 \(h\)。以
\(\overline{BC}\) 为一边，过 \(A\) 作 \(\overline{BC}\) 的平行线
\(l\)，过 \(C\) 作 \(\overline{AB}\) 的平行线交 \(l\) 于 \(D\)。则
\(ABCD\)
是平行四边形（\(\overline{AD} \parallel \overline{BC}\)，\(\overline{AB} \parallel \overline{CD}\)），其面积为
\(bh\)。

\(\triangle ABC\) 将平行四边形 \(ABCD\)
分为两个全等的三角形（\(\triangle ABC \cong \triangle CDA\)，由
SAS：\(\overline{AB} = \overline{CD}\)，\(\angle BAC = \angle DCA\)（内错角），\(\overline{AC}\)
公共）。因此
\(\text{area}(\triangle ABC) = \frac{1}{2} \text{area}(ABCD) = \frac{1}{2}bh\)。\(\blacksquare\)

\paragraph{13.5.3 海伦公式}\label{ux6d77ux4f26ux516cux5f0f}

\textbf{定理 13.32}（海伦公式）：设三角形三边为
\(a\)、\(b\)、\(c\)，半周长
\(s = \frac{a + b + c}{2}\)，则三角形的面积为

\[S = \sqrt{s(s-a)(s-b)(s-c)}\]

\textbf{证明}：

设 \(\triangle ABC\) 中，\(BC = a\)，\(AC = b\)，\(AB = c\)。以 \(BC\)
为底，高为 \(h\)，\(\overline{BC}\) 上的垂足为 \(D\)，设 \(BD = x\)，则
\(DC = a - x\)。

由勾股定理： \[c^2 - x^2 = h^2 = b^2 - (a - x)^2\]

展开得 \(c^2 - x^2 = b^2 - a^2 + 2ax - x^2\)，化简得
\(x = \frac{a^2 + c^2 - b^2}{2a}\)。

因此
\(h^2 = c^2 - x^2 = c^2 - \left(\frac{a^2 + c^2 - b^2}{2a}\right)^2\)。

\[S^2 = \frac{1}{4}a^2 h^2 = \frac{1}{4}\left[a^2 c^2 - \frac{(a^2 + c^2 - b^2)^2}{4}\right]\]

\[= \frac{1}{16}\left[4a^2 c^2 - (a^2 + c^2 - b^2)^2\right]\]

利用平方差公式
\(4a^2c^2 - (a^2 + c^2 - b^2)^2 = [2ac - (a^2 + c^2 - b^2)][2ac + (a^2 + c^2 - b^2)]\)：

\[= [b^2 - (a - c)^2][(a + c)^2 - b^2]\]

\[= (b - a + c)(b + a - c)(a + c - b)(a + c + b)\]

代入 \(s = \frac{a + b + c}{2}\)： - \(b + a - c = 2(s - c)\) -
\(b - a + c = 2(s - a)\) - \(a + c - b = 2(s - b)\) - \(a + c + b = 2s\)

\[S^2 = \frac{1}{16} \cdot 16 \cdot s(s - a)(s - b)(s - c) = s(s - a)(s - b)(s - c)\]

因此 \(S = \sqrt{s(s-a)(s-b)(s-c)}\)。\(\blacksquare\)

\subsubsection{13.6
锐角三角函数}\label{ux9510ux89d2ux4e09ux89d2ux51fdux6570}

\paragraph{13.6.1
三角函数的定义}\label{ux4e09ux89d2ux51fdux6570ux7684ux5b9aux4e49}

在建立了面积理论和相似理论之后，我们现在引入锐角三角函数的概念。三角函数将角度与线段比值联系起来，是连接几何与代数的重要桥梁。

\textbf{定义 13.13}（锐角三角函数）：设 \(\theta\)
是一个锐角（\(0° < \theta < 90°\)）。在一个包含角 \(\theta\)
的直角三角形中，设 \(\theta\) 的对边长为 \(a\)，邻边长为 \(b\)，斜边长为
\(c\)，定义：

\[\sin\theta = \frac{a}{c}, \quad \cos\theta = \frac{b}{c}, \quad \tan\theta = \frac{a}{b}\]

分别称为角 \(\theta\) 的\textbf{正弦}、\textbf{余弦}和\textbf{正切}。

\textbf{定义 13.13'}（三角函数向钝角的推广）：对于钝角
\(90° < \theta < 180°\)，利用其补角 \(180° - \theta\)（为锐角）定义：

\[\sin\theta = \sin(180° - \theta), \quad \cos\theta = -\cos(180° - \theta), \quad \tan\theta = -\tan(180° - \theta)\]

\textbf{几何意义}：若 \(\theta\) 是 \(\triangle ABC\)
中的钝角（\(\angle A = \theta > 90°\)），从 \(B\) 向 \(AC\) 的延长线作高
\(BH\)，则
\(\sin\theta = \frac{BH}{AB}\)，\(\cos\theta = -\frac{AH}{AB}\)（\(H\)
在 \(AC\) 的延长线上，故 \(AH\) 取负值）。

\textbf{性质}：对 \(0° < \theta < 180°\)，有
\(\sin\theta > 0\)；\(\cos\theta > 0\) 当且仅当
\(\theta < 90°\)；\(\cos\theta < 0\) 当且仅当 \(\theta > 90°\)。且
\(\sin^2\theta + \cos^2\theta = 1\) 仍然成立（由定义直接验证）。

\textbf{定理 13.33}（三角函数的良定义性）：三角函数的值只依赖于角
\(\theta\) 的大小，与所选取的直角三角形无关。

\textbf{证明}：

设 \(\triangle ABC\) 和 \(\triangle A'B'C'\) 都是以 \(\theta\)
为一个锐角的直角三角形，其中
\(\angle C = \angle C' = 90°\)，\(\angle A = \angle A' = \theta\)。

由 AA 相似（定理
13.19），\(\triangle ABC \sim \triangle A'B'C'\)。由相似三角形对应边成比例：

\[\frac{BC}{AB} = \frac{B'C'}{A'B'}, \quad \frac{AC}{AB} = \frac{A'C'}{A'B'}, \quad \frac{BC}{AC} = \frac{B'C'}{A'C'}\]

即 \(\sin\theta\)、\(\cos\theta\)、\(\tan\theta\)
的值在两个三角形中相同。\(\blacksquare\)

\paragraph{13.6.2
三角函数的基本关系}\label{ux4e09ux89d2ux51fdux6570ux7684ux57faux672cux5173ux7cfb}

\textbf{定理 13.34}（平方关系）：对于任意锐角
\(\theta\)，\(\sin^2\theta + \cos^2\theta = 1\)。

\textbf{证明}：

在含角 \(\theta\) 的直角三角形中，设对边为 \(a\)，邻边为 \(b\)，斜边为
\(c\)。由勾股定理，\(a^2 + b^2 = c^2\)。因此：

\[\sin^2\theta + \cos^2\theta = \frac{a^2}{c^2} + \frac{b^2}{c^2} = \frac{a^2 + b^2}{c^2} = \frac{c^2}{c^2} = 1\]

\(\blacksquare\)

\textbf{定理 13.35}（商关系）：对于任意锐角
\(\theta\)，\(\tan\theta = \frac{\sin\theta}{\cos\theta}\)。

\textbf{证明}：

\[\frac{\sin\theta}{\cos\theta} = \frac{a/c}{b/c} = \frac{a}{b} = \tan\theta\]

\(\blacksquare\)

\textbf{定理 13.36}（互余关系）：对于任意锐角
\(\theta\)，\(\sin(90° - \theta) = \cos\theta\)，\(\cos(90° - \theta) = \sin\theta\)。

\textbf{证明}：

在直角三角形中，设一个锐角为 \(\theta\)，另一个锐角为
\(90° - \theta\)。\(\theta\) 的对边是 \(90° - \theta\)
的邻边，\(\theta\) 的邻边是 \(90° - \theta\) 的对边。因此：

\[\sin(90° - \theta) = \frac{\theta\text{的邻边}}{\text{斜边}} = \cos\theta\]

\[\cos(90° - \theta) = \frac{\theta\text{的对边}}{\text{斜边}} = \sin\theta\]

\(\blacksquare\)

\paragraph{13.6.3
特殊角的三角函数值}\label{ux7279ux6b8aux89d2ux7684ux4e09ux89d2ux51fdux6570ux503c}

\textbf{定理
13.37}：\(\sin 30° = \frac{1}{2}\)，\(\cos 30° = \frac{\sqrt{3}}{2}\)，\(\tan 30° = \frac{\sqrt{3}}{3}\)。

\textbf{证明}：

考虑等边三角形 \(\triangle ABC\)（边长为 \(2\)）。作
\(\overline{AD} \perp \overline{BC}\) 于 \(D\)。由三线合一（定理
13.6），\(D\) 是 \(\overline{BC}\) 的中点，\(BD = 1\)。在
\(\triangle ABD\)
中，\(\angle ADB = 90°\)，\(\angle BAD = 30°\)（等边三角形每个角为
\(60°\)，高平分顶角）。由勾股定理，\(AD = \sqrt{AB^2 - BD^2} = \sqrt{4 - 1} = \sqrt{3}\)。

\[\sin 30° = \frac{BD}{AB} = \frac{1}{2}, \quad \cos 30° = \frac{AD}{AB} = \frac{\sqrt{3}}{2}, \quad \tan 30° = \frac{BD}{AD} = \frac{1}{\sqrt{3}} = \frac{\sqrt{3}}{3}\]

\(\blacksquare\)

\textbf{定理
13.38}：\(\sin 45° = \frac{\sqrt{2}}{2}\)，\(\cos 45° = \frac{\sqrt{2}}{2}\)，\(\tan 45° = 1\)。

\textbf{证明}：

考虑等腰直角三角形，两直角边长为 \(1\)，斜边长为
\(\sqrt{2}\)（由勾股定理）。两锐角均为 \(45°\)。

\[\sin 45° = \frac{1}{\sqrt{2}} = \frac{\sqrt{2}}{2}, \quad \cos 45° = \frac{1}{\sqrt{2}} = \frac{\sqrt{2}}{2}, \quad \tan 45° = \frac{1}{1} = 1\]

\(\blacksquare\)

\textbf{定理
13.39}：\(\sin 60° = \frac{\sqrt{3}}{2}\)，\(\cos 60° = \frac{1}{2}\)，\(\tan 60° = \sqrt{3}\)。

\textbf{证明}：由互余关系（定理
13.36），\(\sin 60° = \cos 30° = \frac{\sqrt{3}}{2}\)，\(\cos 60° = \sin 30° = \frac{1}{2}\)，\(\tan 60° = \frac{\sin 60°}{\cos 60°} = \frac{\sqrt{3}/2}{1/2} = \sqrt{3}\)。\(\blacksquare\)

\paragraph{13.6.4
正弦定理与余弦定理}\label{ux6b63ux5f26ux5b9aux7406ux4e0eux4f59ux5f26ux5b9aux7406}

\textbf{定理 13.40}（正弦定理）：在 \(\triangle ABC\) 中，设角
\(A\)、\(B\)、\(C\) 所对的边分别为 \(a\)、\(b\)、\(c\)，外接圆半径为
\(R\)，则

\[\frac{a}{\sin A} = \frac{b}{\sin B} = \frac{c}{\sin C} = 2R\]

\textbf{证明}：

设 \(\triangle ABC\) 的外接圆为 \(\odot(O, R)\)。

\textbf{情形 1}：\(\angle A\) 为锐角。连接 \(BO\) 并延长交圆于 \(D\)。则
\(BD = 2R\)（直径）。连接 \(DC\)。由推论
13.5，\(\angle BCD = 90°\)（半圆上的圆周角）。又
\(\angle BDC = \angle BAC = A\)（同弧所对的圆周角相等，推论 13.4）。在
\(\triangle BCD\) 中，\(\sin(\angle BDC) = \frac{BC}{BD}\)，即
\(\sin A = \frac{a}{2R}\)，故 \(\frac{a}{\sin A} = 2R\)。

\textbf{情形 2}：\(\angle A\) 为直角。则
\(a = 2R\)（直径），\(\sin A = \sin 90° = 1\)，\(\frac{a}{\sin A} = 2R\)。

\textbf{情形 3}：\(\angle A\) 为钝角。设 \(\triangle ABC\) 的外接圆为
\(\odot(O, R)\)。连接 \(BO\) 并延长交圆于 \(D\)，则
\(BD = 2R\)（直径）。连接 \(DC\)。

由推论 13.5，\(\angle BCD = 90°\)（半圆所对的圆周角是直角）。

现在确定 \(\angle BDC\)。弧 \(BC\) 有两条：劣弧 \(BC\) 和优弧
\(BAC\)（经过 \(A\) 的弧）。\(\angle A\) 是优弧 \(BAC\)
所对的圆周角，\(\angle BDC\) 是劣弧 \(BC\)
所对的圆周角。由圆周角定理（定理
13.25），\(\angle A = \frac{1}{2} \cdot (\text{优弧 } BAC)\)，\(\angle BDC = \frac{1}{2} \cdot (\text{劣弧 } BC)\)。由于优弧
\(BAC\) + 劣弧 \(BC = 360°\)，故 \(\angle A + \angle BDC = 180°\)，即
\(\angle BDC = 180° - \angle A\)。

在直角三角形 \(\triangle BCD\) 中（\(\angle BCD = 90°\)）：

\[\sin(\angle BDC) = \frac{BC}{BD} = \frac{a}{2R}\]

因此 \(\sin(180° - A) = \frac{a}{2R}\)。由于
\(\sin(180° - A) = \sin A\)（正弦函数的互补性质，可由定义 13.13
和互余关系推导：在含角 \(\angle BDC = 180° - A\)
的直角三角形中，\(\sin(180° - A) = \sin A\)），故：

\[\sin A = \frac{a}{2R}\]

即 \(\frac{a}{\sin A} = 2R\)。

同理，\(\frac{b}{\sin B} = 2R\)，\(\frac{c}{\sin C} = 2R\)。\(\blacksquare\)

\textbf{定理 13.41}（余弦定理）：在 \(\triangle ABC\)
中，\(a^2 = b^2 + c^2 - 2bc\cos A\)，\(b^2 = a^2 + c^2 - 2ac\cos B\)，\(c^2 = a^2 + b^2 - 2ab\cos C\)。

\textbf{证明}：

从顶点 \(C\) 向边 \(\overline{AB}\) 作高 \(\overline{CD}\)，设垂足为
\(D\)，\(AD = x\)，\(CD = h\)。

在 \(\triangle ACD\) 中（\(\angle ADC = 90°\)），由勾股定理：

\[b^2 = h^2 + x^2\]

在 \(\triangle BCD\) 中（\(\angle BDC = 90°\)），由勾股定理：

\[a^2 = h^2 + (c - x)^2\]

两式相减：

\[a^2 - b^2 = (c - x)^2 - x^2 = c^2 - 2cx\]

因此 \(a^2 = b^2 + c^2 - 2cx\)。

又在 \(\triangle ACD\) 中，\(\cos A = \frac{x}{b}\)，故
\(x = b\cos A\)。代入上式得：

\[a^2 = b^2 + c^2 - 2bc\cos A\]

同理可证其余两式。\(\blacksquare\)

\textbf{推论 13.6}：当 \(A = 90°\)
时，\(\cos A = 0\)，余弦定理退化为勾股定理 \(a^2 = b^2 + c^2\)。

\textbf{证明}：在余弦定理 \(a^2 = b^2 + c^2 - 2bc\cos A\) 中令
\(A = 90°\)，则 \(\cos 90° = 0\)，得
\(a^2 = b^2 + c^2\)。\(\blacksquare\)

\bigskip\hrule\bigskip

\textbf{本章展望}：本章从希尔伯特公理系统出发，系统推导了平面几何的全部核心定理：三角形全等与相似的判定准则、等腰三角形和边角关系、四边形的性质与分类、圆的核心定理（圆周角定理、切线定理、圆幂定理）、面积的公理化理论以及锐角三角函数。在下一章中，我们将把这些平面几何的结果推广到三维空间，建立立体几何的完整理论体系，包括线面关系的判定与性质、空间角与距离、以及简单几何体的体积与表面积。

\bigskip\hrule\bigskip

\subsection{Ch 14 立体几何}\label{ch-14-ux7acbux4f53ux51e0ux4f55}

\textbf{前置知识}：本章假定读者已掌握 Ch 12
欧氏几何公理体系（特别是关联公理 I1--I8 和平行公理 P）以及 Ch 13
平面几何的全部定理（三角形全等与相似、四边形性质、圆的定理、面积理论）。本章将平面几何的结果系统推广到三维空间，建立直线与平面的位置关系、空间角与距离、以及简单几何体的体积与表面积理论。

\bigskip\hrule\bigskip

上一章中，我们建立了平面几何的完整理论。然而，我们所生活的物理空间是三维的。从平面到空间的推广并非简单的''多加一个维度''------它引入了全新的概念（如异面直线、二面角）和全新的关系（如线面平行、面面垂直），这些在平面中没有对应物。本章严格建立立体几何的理论基础，从空间中直线与平面的基本关系出发，逐步推导空间角与距离的概念，最终建立各类简单几何体的体积与表面积公式。

\subsubsection{14.1
空间中的直线与平面关系}\label{ux7a7aux95f4ux4e2dux7684ux76f4ux7ebfux4e0eux5e73ux9762ux5173ux7cfb}

\paragraph{14.1.1
线面平行的判定与性质}\label{ux7ebfux9762ux5e73ux884cux7684ux5224ux5b9aux4e0eux6027ux8d28}

\textbf{定义
14.1}（直线与平面平行）：如果一条直线与一个平面没有公共点，则称该直线与该平面\textbf{平行}，记为
\(l \parallel \alpha\)。

\textbf{定义
14.2}（直线在平面内）：如果一条直线上的所有点都在一个平面内，则称该直线\textbf{在平面内}（或含于平面），记为
\(l \subset \alpha\)。

\textbf{定义
14.3}（直线与平面相交）：如果一条直线与一个平面有且仅有一个公共点，则称该直线与该平面\textbf{相交}，公共点称为\textbf{交点}。

空间中直线与平面的位置关系有且仅有三种：直线在平面内、直线与平面平行、直线与平面相交。

\textbf{定理
14.1}（线面平行的判定定理）：如果平面外一条直线与平面内一条直线平行，则该直线与此平面平行。

即：设 \(l \not\subset \alpha\)，\(m \subset \alpha\)，且
\(l \parallel m\)，则 \(l \parallel \alpha\)。

\textbf{证明}（反证法）：

假设 \(l\) 与 \(\alpha\) 不平行，则 \(l\) 与 \(\alpha\)
有公共点。设公共点为 \(P\)。

因为 \(l \parallel m\)，所以 \(l\) 与 \(m\) 不相交。但 \(P \in l\) 且
\(P \in \alpha\)，\(m \subset \alpha\)。

如果 \(P \in m\)，则 \(P\) 是 \(l\) 与 \(m\) 的公共点，与
\(l \parallel m\) 矛盾。

如果 \(P \notin m\)，则 \(P\) 和 \(m\) 在同一平面 \(\alpha\) 上，过
\(P\) 作直线 \(m' \parallel m\)。由平行公理（在平面 \(\alpha\)
上），这样的 \(m'\) 存在且唯一。

但 \(l\) 也过 \(P\) 且
\(l \parallel m\)。由空间中平行线的传递性（两条直线都平行于第三条直线，且过同一点，则它们重合），\(l = m'\)。

因此 \(l = m' \subset \alpha\)，与 \(l \not\subset \alpha\) 矛盾。

所以 \(l \parallel \alpha\)。\(\blacksquare\)

\textbf{定理
14.2}（线面平行的性质定理）：如果一条直线与一个平面平行，过该直线的平面与此平面相交，则该直线与交线平行。

即：设
\(l \parallel \alpha\)，\(l \subset \beta\)，\(\alpha \cap \beta = m\)，则
\(l \parallel m\)。

\textbf{证明}：

因为 \(l \parallel \alpha\)，所以 \(l\) 与 \(\alpha\) 没有公共点。

因为 \(m \subset \alpha\)，所以 \(l\) 与 \(m\) 没有公共点（\(l\)
上的点都不在 \(\alpha\) 上，自然不在 \(m\) 上）。

又因为 \(l \subset \beta\)，\(m \subset \beta\)，所以 \(l\) 和 \(m\)
在同一平面 \(\beta\) 上。

在同一平面上没有公共点的两条直线平行，因此
\(l \parallel m\)。\(\blacksquare\)

\textbf{定理
14.3}（线面平行性质的推论）：如果一条直线与一个平面平行，则过该直线的任意平面与该平面的交线都与该直线平行。

\textbf{证明}：设直线 \(l \parallel\) 平面 \(\alpha\)，\(\beta\) 为过
\(l\) 的平面，\(\beta \cap \alpha = m\)。若 \(l\) 与 \(m\)
相交，则交点在 \(\alpha\) 上（因 \(m \subset \alpha\)），与
\(l \parallel \alpha\) 矛盾。又 \(l, m \subset \beta\)，故
\(l \parallel m\)。\(\blacksquare\)

\textbf{推论
14.1}：如果两条平行线中的一条平行于一个平面，则另一条也平行于该平面（或在该平面内）。

\textbf{证明}：设 \(l_1 \parallel l_2\)，\(l_1 \parallel \alpha\)。在
\(\alpha\) 内取一点 \(P\)（\(P \notin l_2\)），过 \(P\) 和 \(l_1\)
作平面 \(\beta\)。设 \(\beta \cap \alpha = m\)。由定理
14.2，\(l_1 \parallel m\)。又
\(l_1 \parallel l_2\)，由平行线的传递性，\(l_2 \parallel m\)。由定理
14.1，\(l_2 \parallel \alpha\)（若
\(l_2 \not\subset \alpha\)）。\(\blacksquare\)

\paragraph{14.1.2
面面平行的判定与性质}\label{ux9762ux9762ux5e73ux884cux7684ux5224ux5b9aux4e0eux6027ux8d28}

\textbf{定义
14.4}（平面与平面平行）：如果两个平面没有公共点，则称这两个平面\textbf{平行}，记为
\(\alpha \parallel \beta\)。

\textbf{定理
14.4}（面面平行的判定定理）：如果一个平面内有两条相交直线分别平行于另一个平面，则这两个平面平行。

即：设
\(a \subset \alpha\)，\(b \subset \alpha\)，\(a \cap b = P\)，\(a \parallel \beta\)，\(b \parallel \beta\)，则
\(\alpha \parallel \beta\)。

\textbf{证明}（反证法）：

假设 \(\alpha\) 与 \(\beta\) 不平行，则它们相交。设
\(\alpha \cap \beta = l\)。

因为
\(a \parallel \beta\)，\(a \subset \alpha\)，由线面平行的性质定理（定理
14.2），\(a \parallel l\)。

因为
\(b \parallel \beta\)，\(b \subset \alpha\)，由线面平行的性质定理，\(b \parallel l\)。

所以 \(a \parallel l\) 且 \(b \parallel l\)，即
\(a \parallel b\)。但这与 \(a \cap b = P\) 矛盾。

因此 \(\alpha \parallel \beta\)。\(\blacksquare\)

\textbf{定理 14.5}（面面平行的性质定理）：

\begin{enumerate}
\def\labelenumi{\arabic{enumi}.}
\tightlist
\item
  如果两个平行平面都与第三个平面相交，则交线平行。
\item
  如果一条直线垂直于两个平行平面中的一个，则也垂直于另一个。
\end{enumerate}

\textbf{证明}：

\textbf{性质 1}：设
\(\alpha \parallel \beta\)，\(\gamma \cap \alpha = a\)，\(\gamma \cap \beta = b\)。

由于
\(a \subset \alpha\)，\(b \subset \beta\)，\(\alpha \parallel \beta\)，所以
\(a\) 与 \(b\) 没有公共点。

又 \(a \subset \gamma\)，\(b \subset \gamma\)，\(a\) 和 \(b\)
在同一平面上，所以 \(a \parallel b\)。

\textbf{性质 2}：设 \(\alpha \parallel \beta\)，\(l \perp \alpha\) 于
\(A\)。要证明 \(l \perp \beta\)。

过 \(l\) 作平面 \(\gamma\)，设
\(\gamma \cap \alpha = a\)，\(\gamma \cap \beta = b\)。由性质
1，\(a \parallel b\)。

因为 \(l \perp \alpha\)，所以 \(l \perp a\)。又 \(a \parallel b\)，所以
\(l \perp b\)。

对 \(\beta\) 内过垂足 \(H\)（\(l\) 与 \(\beta\) 的交点）的任意直线
\(c\)，取包含 \(l\) 和 \(c\) 的平面 \(\gamma\)。则
\(\gamma \cap \alpha = a_c\)（过 \(H\) 在 \(\alpha\)
内），\(\gamma \cap \beta = c\)。由性质 1，\(a_c \parallel c\)。由
\(l \perp \alpha\) 得 \(l \perp a_c\)。再由 \(a_c \parallel c\) 得
\(l \perp c\)。

因此 \(l\) 垂直于 \(\beta\) 上所有过 \(l\) 与 \(\beta\)
的交点的直线，所以 \(l \perp \beta\)。\(\blacksquare\)

\paragraph{14.1.3
线面垂直的判定与性质}\label{ux7ebfux9762ux5782ux76f4ux7684ux5224ux5b9aux4e0eux6027ux8d28}

\textbf{定义
14.5}（直线与平面垂直）：如果一条直线与一个平面内的\textbf{所有}直线都垂直，则称该直线与该平面\textbf{垂直}。该直线称为平面的\textbf{垂线}，垂线与平面的交点称为\textbf{垂足}。

\textbf{定义
14.6}（斜线与射影）：如果一条直线与平面相交但不垂直，则称该直线为平面的\textbf{斜线}，交点称为\textbf{斜足}。从斜线上一点向平面作垂线，垂足与斜足之间的线段称为斜线在平面上的\textbf{射影}。

\textbf{定理
14.6}（线面垂直的判定定理）：如果一条直线与一个平面内的两条相交直线都垂直，则该直线与此平面垂直。

即：设
\(a \subset \alpha\)，\(b \subset \alpha\)，\(a \cap b = P\)，\(l \perp a\)，\(l \perp b\)，则
\(l \perp \alpha\)。

\textbf{证明}：

设 \(l\) 过点 \(P\)（若不过 \(P\)，可作平移使过 \(P\) 的直线与 \(l\)
平行，保持垂直关系）。要证明 \(l\) 垂直于 \(\alpha\) 内过 \(P\)
的任意直线 \(c\)。

在 \(l\) 上取点 \(A\)（不在 \(\alpha\) 上），取 \(A'\) 使得 \(PA' = PA\)
且 \(A\)、\(A'\) 在 \(\alpha\) 的两侧（即 \(A*P*A'\)）。

\textbf{关键步骤}：证明 \(\alpha\) 上到 \(P\) 距离为 \(PA\) 的任意点到
\(A\) 和 \(A'\) 的距离相等。

在 \(a\) 上取点 \(B\) 使得 \(PB = PA\)。在 \(\triangle APB\) 和
\(\triangle A'PB\) 中： - \(PA = PA'\)（由构造） -
\(\angle APB = \angle A'PB = 90°\)（\(l \perp a\)，故
\(\angle APB = 90°\)；又
\(A*P*A'\)，\(\angle A'PB = 180° - \angle APB = 90°\)） -
\(PB = PB\)（公共边）

由 SAS，\(\triangle APB \cong \triangle A'PB\)，所以 \(AB = A'B\)。

同理，在 \(b\) 上取 \(C\) 使得 \(PC = PA\)，可证 \(AC = A'C\)。

现在，设 \(c\) 是 \(\alpha\) 内过 \(P\) 的任意直线。在包含 \(l\) 和
\(c\) 的平面 \(\beta\) 中，设 \(D\) 在 \(c\) 上且 \(PD = PA\)。

利用余弦定理计算 \(AD\) 和 \(A'D\)。设 \(\angle APD = \theta\)，则
\(\angle A'PD = 180° - \theta\)。

\[AD^2 = PA^2 + PD^2 - 2 \cdot PA \cdot PD \cos\theta\]
\[A'D^2 = PA'^2 + PD^2 - 2 \cdot PA' \cdot PD \cos(180° - \theta) = PA'^2 + PD^2 + 2 \cdot PA' \cdot PD \cos\theta\]

所以 \(AD^2 - A'D^2 = -4 \cdot PA \cdot PD \cos\theta\)。

\(AD = A'D\) 当且仅当 \(\cos\theta = 0\)，即 \(\theta = 90°\)。

现在需要证明 \(AD = A'D\)，从而由上式推出 \(\cos\theta = 0\)，即
\(\theta = 90°\)。

利用 \(a\) 和 \(b\) 提供的信息。在 \(\alpha\) 内，\(a\) 和 \(b\) 张成过
\(P\) 的整个平面。\(D\) 是 \(c\) 上满足 \(PD = PA\) 的点。在 \(c\) 上取
\(P\) 关于 \(D\) 的另一侧的点 \(D'\) 使得 \(PD' = PD = PA\)。在 \(l\)
上，\(A\) 和 \(A'\) 也关于 \(P\) 对称（\(PA = PA'\)，\(A * P * A'\)）。

由于 \(a\)、\(b\) 都是 \(\alpha\) 内过 \(P\) 且与 \(l\)
垂直的直线，\(a\) 和 \(b\) 上满足 \(PX = PA\) 的点 \(X\) 均满足
\(AX = A'X\)（已证）。现在，\(D\) 在 \(c\) 上，而 \(c\) 是 \(\alpha\)
内过 \(P\) 的另一条直线。设 \(\alpha\) 内过 \(P\) 且满足 \(PX = PA\)
的点 \(X\) 的集合为 \(S_P\)（即以 \(P\) 为圆心、\(PA\) 为半径的圆周在
\(\alpha\) 内的部分）。集合 \(S_P\) 是 \(\alpha\) 内的一个闭合曲线。

\(a\) 上的 \(B\) 和 \(b\) 上的 \(C\) 都在 \(S_P\) 上，且
\(AB = A'B\)，\(AC = A'C\)。现在 \(D\) 也在 \(S_P\)
上。由对称性论证（\(A\) 和 \(A'\) 关于 \(\alpha\) 内过 \(P\) 的平面
\(l\) 对称），\(S_P\) 上的每个点到 \(A\) 和 \(A'\) 的距离相等。这是因为
\(A\) 和 \(A'\) 关于由 \(l\) 和 \(\alpha\) 的法向量确定的对称面对称，且
\(S_P\) 在此对称面上是不变的。更严格地：取 \(S_P\) 上任意点 \(X\)，设
\(X\) 在 \(a\) 和 \(b\) 方向上的分量分别为 \(s\) 和 \(t\)（即
\(\overrightarrow{PX} = s \cdot \hat{a} + t \cdot \hat{b}\)，其中
\(\hat{a}\)、\(\hat{b}\) 是 \(a\)、\(b\)
方向的单位向量）。由余弦定理，\(AX^2 = PA^2 + PX^2 - 2 \cdot PA \cdot PX \cos\angle APX\)，\(A'X^2 = PA^2 + PX^2 + 2 \cdot PA \cdot PX \cos\angle APX\)。由
\(a\) 和 \(b\) 两个方向上的垂直关系以及 \(D\) 在 \(a\)、\(b\)
张成的平面内这一事实，可利用阿基米德公理和帕施公理论证
\(AD^2 = A'D^2\)（具体地，\(S_P\) 上到 \(A\) 和 \(A'\) 等距的点构成
\(S_P\) 的一个闭子集，包含 \(B\) 和 \(C\)；由 \(a \perp l\) 和
\(b \perp l\) 确保该子集非空且在 \(S_P\)
上稠密，结合完备性公理可知该子集等于 \(S_P\)）。

\textbf{严格补充}：对 \(\alpha\) 内过 \(P\) 的任意直线 \(c\)，可在
\(\alpha\) 内将 \(c\) 表示为沿 \(a\) 方向和 \(b\) 方向的线性组合。设
\(c\) 上距 \(P\) 为 \(d\) 的点为 \(D_c\)，其在 \(a\) 和 \(b\)
方向上的分量分别为 \(d_a\) 和 \(d_b\)（\(d_a^2 + d_b^2 = d^2\)）。由
\(l \perp a\) 和 \(l \perp b\)，利用勾股定理（在含 \(l\) 和 \(a\)
的平面内，\(AD^2 = AP^2 + PD^2\)；类似对 \(b\)），可得 \(A\) 到 \(D_c\)
的距离平方为 \(AP^2 + d^2\)，与 \(c\) 的方向无关。因此 \(l\) 垂直于
\(\alpha\) 内所有过垂足的直线，即 \(l \perp \alpha\)。

因此 \(AD = A'D\)，从而 \(\cos\theta = 0\)，\(\theta = 90°\)，即
\(l \perp c\)。由于 \(c\) 是 \(\alpha\) 内过 \(P\)
的任意直线，\(l \perp \alpha\)。\(\blacksquare\)

\textbf{定理 14.7}（线面垂直的性质定理）：

\begin{enumerate}
\def\labelenumi{\arabic{enumi}.}
\tightlist
\item
  如果两条直线都垂直于同一个平面，则这两条直线平行。
\item
  过一点与已知平面垂直的直线有且仅有一条。
\end{enumerate}

\textbf{证明}：

\textbf{性质 1}：设 \(a \perp \alpha\)，\(b \perp \alpha\)。假设 \(a\)
与 \(b\) 不平行，则 \(a\) 和 \(b\) 相交或异面。

如果 \(a\) 和 \(b\) 相交于 \(P\)，取 \(\alpha\) 中过 \(P\)
的任意一条直线 \(l\)。因为 \(a \perp \alpha\)，\(b \perp \alpha\)，所以
\(a \perp l\) 且 \(b \perp l\)。在由 \(a\) 和 \(l\) 确定的平面内，过
\(P\) 与 \(l\) 垂直的直线唯一（平面几何基本事实），故
\(a = b\)，与假设矛盾。

如果 \(a\) 和 \(b\) 异面，过 \(b\) 上一点 \(Q\) 作
\(a' \parallel a\)。由线面垂直的判定（\(\alpha\) 内两条相交直线分别与
\(a'\) 垂直，可由 \(a \perp \alpha\) 推出），\(a' \perp \alpha\)。但
\(b \perp \alpha\)，\(a'\) 和 \(b\) 都过 \(Q\) 且都垂直于
\(\alpha\)，由上述推理（相交情形）知 \(a' = b\)，故
\(a \parallel b\)，与 \(a\)、\(b\) 异面矛盾。

因此 \(a \parallel b\)。

\textbf{性质 2}：设 \(P\) 是空间中一点，\(l\) 和 \(l'\) 都过 \(P\)
且都垂直于平面 \(\alpha\)。由性质 1，\(l \parallel l'\)。又 \(l\) 和
\(l'\) 都过 \(P\)，故
\(l = l'\)（过同一点的平行线重合）。\(\blacksquare\)

\paragraph{14.1.4
面面垂直的判定与性质}\label{ux9762ux9762ux5782ux76f4ux7684ux5224ux5b9aux4e0eux6027ux8d28}

\textbf{定义
14.7}（二面角）：从一条直线出发的两个半平面所组成的图形称为\textbf{二面角}，这条直线称为\textbf{棱}，两个半平面称为\textbf{面}。

\textbf{定义
14.8}（二面角的平面角）：在二面角的棱上任取一点，过该点分别在两个半平面内作棱的垂线，这两条射线所成的角称为二面角的\textbf{平面角}。

\textbf{命题 14.1}：二面角的平面角与棱上点的选取无关。

\textbf{证明}：设在棱上取两点
\(P\)、\(Q\)，分别作棱的垂线。由这些垂线确定的平面角相等（因为对应的射线互相平行，异面直线所成角的定义保证了角的大小不变）。\(\blacksquare\)

\textbf{定义
14.9}（两个平面垂直）：如果两个平面所成的二面角是直角（即二面角的平面角为
\(90°\)），则称这两个平面\textbf{互相垂直}，记为
\(\alpha \perp \beta\)。

\textbf{定理
14.8}（面面垂直的判定定理）：如果一个平面经过另一个平面的一条垂线，则这两个平面互相垂直。

即：设 \(l \perp \alpha\)，\(l \subset \beta\)，则
\(\beta \perp \alpha\)。

\textbf{证明}：

设 \(\alpha \cap \beta = m\)。在 \(m\) 上取点 \(P\)，在 \(\beta\) 内作
\(PA \perp m\) 于 \(P\)（\(A \in l\)）。

因为 \(l \perp \alpha\)，\(m \subset \alpha\)，所以 \(l \perp m\)，即
\(\angle APl = 90°\)。

在 \(\beta\) 内，\(PA \perp m\)，\(l \perp m\)，且 \(PA\) 和 \(l\) 都在
\(\beta\) 内。所以 \(PA\) 与 \(l\) 的夹角（在 \(\beta\) 内度量）就是以
\(m\) 为棱的二面角的平面角。

由于 \(l \perp \alpha\)，\(\angle APl = 90°\)（\(l\) 垂直于 \(\alpha\)
内过 \(P\) 的直线 \(PA\)），所以二面角的平面角为 \(90°\)，即
\(\beta \perp \alpha\)。\(\blacksquare\)

\textbf{定理 14.9}（面面垂直的性质定理）：

\begin{enumerate}
\def\labelenumi{\arabic{enumi}.}
\tightlist
\item
  如果两个平面互相垂直，那么在一个平面内垂直于交线的直线垂直于另一个平面。
\item
  如果两个平面互相垂直，那么从一个平面内一点作另一个平面的垂线，垂足在交线上。
\end{enumerate}

\textbf{证明}：

\textbf{性质 1}：设
\(\alpha \perp \beta\)，\(\alpha \cap \beta = m\)，\(l \subset \alpha\)，\(l \perp m\)。

在 \(\beta\) 内作 \(n \perp m\) 于
\(P\)（\(P = l \cap m\)）。由二面角的定义，\(l\) 与 \(n\)
的夹角是二面角的平面角，等于 \(90°\)。

所以 \(l \perp n\)。又
\(l \perp m\)，\(n \cap m = P\)，\(n, m \subset \beta\)，由线面垂直的判定（定理
14.6），\(l \perp \beta\)。

\textbf{性质 2}：设
\(\alpha \perp \beta\)，\(A \in \alpha\)，\(AH \perp \beta\) 于
\(H\)。要证明 \(H \in m\)（\(m\) 是交线）。

过 \(A\) 作 \(AB \perp m\) 于 \(B\)，则 \(AB \subset \alpha\)。由性质
1，\(AB \perp \beta\)。

由于过一点到一个平面的垂线唯一（定理 14.7），\(AH\) 与 \(AB\) 重合，所以
\(H = B \in m\)。\(\blacksquare\)

\paragraph{14.1.5 三垂线定理}\label{ux4e09ux5782ux7ebfux5b9aux7406}

\textbf{定理
14.10}（三垂线定理）：如果平面内的一条直线与一条斜线在这个平面内的射影垂直，则它也与这条斜线垂直。

即：设 \(PO \perp \alpha\) 于 \(O\)，\(PA\)
是斜线（\(A \in \alpha\)），\(a \subset \alpha\)，\(a \perp OA\)（\(OA\)
是 \(PA\) 在 \(\alpha\) 上的射影），则 \(a \perp PA\)。

\textbf{证明}：

因为 \(PO \perp \alpha\)，所以 \(PO \perp a\)（\(a \subset \alpha\)）。

又 \(OA \perp a\)（已知），且 \(PO\) 和 \(OA\) 相交于 \(O\)。

所以 \(a \perp\) 平面 \(POA\)（\(a\) 垂直于平面 \(POA\)
内的两条相交直线，由定理 14.6）。

因此 \(a \perp PA\)（\(PA \subset\) 平面 \(POA\)）。\(\blacksquare\)

\textbf{定理
14.11}（三垂线定理的逆定理）：如果平面内的一条直线与一条斜线垂直，则它也与这条斜线在平面内的射影垂直。

即：设 \(PO \perp \alpha\) 于 \(O\)，\(PA\)
是斜线，\(a \subset \alpha\)，\(a \perp PA\)，则 \(a \perp OA\)。

\textbf{证明}：

因为 \(PO \perp \alpha\)，所以 \(PO \perp a\)。

又 \(a \perp PA\)（已知），且 \(PO\) 和 \(PA\) 相交于 \(P\)。

所以 \(a \perp\) 平面 \(POA\)。

因此 \(a \perp OA\)（\(OA \subset\) 平面 \(POA\)）。\(\blacksquare\)

\subsubsection{14.2
空间中的角与距离}\label{ux7a7aux95f4ux4e2dux7684ux89d2ux4e0eux8dddux79bb}

\paragraph{14.2.1
异面直线所成的角}\label{ux5f02ux9762ux76f4ux7ebfux6240ux6210ux7684ux89d2}

\textbf{定义
14.10}（异面直线）：不在同一平面内的两条直线称为\textbf{异面直线}。

\textbf{定义
14.11}（异面直线所成的角）：过空间中任意一点分别作与两条异面直线平行的直线，这两条相交直线所成的锐角（或直角）称为两条异面直线所成的角。

\textbf{定理 14.12}：异面直线所成的角的大小与所取点的位置无关。

\textbf{证明}：

设 \(a\)、\(b\) 是异面直线，\(P\)、\(Q\) 是空间中任意两点。

过 \(P\) 作 \(a' \parallel a\)，\(b' \parallel b\)，角
\(\alpha = \angle(a', b')\)。

过 \(Q\) 作 \(a'' \parallel a\)，\(b'' \parallel b\)，角
\(\beta = \angle(a'', b'')\)。

因为 \(a' \parallel a \parallel a''\)，所以 \(a' \parallel a''\)。

因为 \(b' \parallel b \parallel b''\)，所以 \(b' \parallel b''\)。

由平行线的性质（同位角或内错角相等），\(\alpha = \beta\)。\(\blacksquare\)

\paragraph{14.2.2
直线与平面所成的角}\label{ux76f4ux7ebfux4e0eux5e73ux9762ux6240ux6210ux7684ux89d2}

\textbf{定义
14.12}：一条直线与它在平面上的射影所成的锐角称为该直线与平面所成的角。

特别地，当直线垂直于平面时，所成的角为
\(90°\)；当直线平行于或在平面内时，所成的角为 \(0°\)。

\textbf{定理
14.13}：斜线与平面所成的角是斜线与平面内所有直线所成角中的最小值。

\textbf{证明}：

设斜线 \(PB\) 与平面 \(\alpha\) 所成的角为 \(\theta\)（\(P\) 在
\(\alpha\) 外，\(O\) 是 \(P\) 在 \(\alpha\) 上的射影，\(B\)
是斜足，\(\angle PBO = \theta\)）。

设 \(l\) 是 \(\alpha\) 上过 \(B\) 的任意直线，\(PB\) 与 \(l\) 所成的角为
\(\phi\)。

作 \(PH \perp l\) 于 \(H\)。在 \(\triangle PBH\)
中，\(\angle PHB = 90°\)，所以 \(PH = PB \sin\phi\)。

在 \(\triangle POB\) 中，\(\angle POB = 90°\)，所以
\(PO = PB \sin\theta\)。

因为 \(O\) 是 \(P\) 到 \(\alpha\)
的最近点（垂足），\(PO \leq PH\)（\(P\) 到平面的距离不超过 \(P\)
到平面内任意直线的距离，由垂线段最短定理）。

所以 \(PB \sin\theta \leq PB \sin\phi\)，即
\(\sin\theta \leq \sin\phi\)。

因为 \(\theta\) 和 \(\phi\) 都是锐角（或
\(\theta = \phi = 90°\)），\(\theta \leq \phi\)。

等号成立当且仅当 \(H = O\)，即 \(l \perp OB\)（\(l\) 垂直于 \(PB\) 在
\(\alpha\) 上的投影）。\(\blacksquare\)

\paragraph{14.2.3
点到平面的距离、异面直线的距离}\label{ux70b9ux5230ux5e73ux9762ux7684ux8dddux79bbux5f02ux9762ux76f4ux7ebfux7684ux8dddux79bb}

\textbf{定义
14.13}（点到平面的距离）：从一点到一个平面所作的垂线段的长度称为该点到该平面的\textbf{距离}。

\textbf{定理
14.14}（垂线段最短）：点到平面的垂线段是该点到平面上所有点的连线中最短的。

\textbf{证明}：

设 \(P\) 在平面 \(\alpha\) 外，\(PO \perp \alpha\) 于 \(O\)。设 \(A\) 是
\(\alpha\) 上不同于 \(O\) 的点。

在 \(\triangle POA\) 中，\(\angle POA = 90°\)，所以
\(PA^2 = PO^2 + OA^2 > PO^2\)，即 \(PA > PO\)。\(\blacksquare\)

\textbf{定义
14.14}（异面直线的距离）：与两条异面直线都垂直相交的线段的长度称为两条异面直线的\textbf{公垂线段}的长度，也称为两条异面直线的\textbf{距离}。

\textbf{定理
14.15}：两条异面直线的公垂线段存在且唯一，且其长度是两条异面直线上任意两点间距离的最小值。

\textbf{证明}（概要）：

设 \(a\)、\(b\) 是异面直线。

\textbf{(1) 过 \(a\) 作平面 \(\alpha \parallel b\) 的存在性论证}：在
\(a\) 上取一点 \(A\)。过 \(A\) 作直线
\(b' \parallel b\)（空间中过一点恰有一条直线与给定直线平行）。若
\(b' = a\)，则 \(a \parallel b\)，与 \(a\)、\(b\) 异面矛盾，故
\(b' \neq a\)。\(a\) 与 \(b'\) 确定唯一平面 \(\alpha\)。又
\(b \parallel b' \subset \alpha\) 且 \(b \not\subset \alpha\)（否则
\(a\) 与 \(b\) 共面于 \(\alpha\)，与异面矛盾），故
\(\alpha \parallel b\)。

\textbf{(2) 过 \(b\) 作平面 \(\beta \perp \alpha\) 的论证}：在 \(b\)
上取一点 \(B\)。由 \(B\) 向平面 \(\alpha\) 作垂线
\(BB'\)（\(B' \in \alpha\)），则 \(BB' \perp \alpha\)。过 \(b\) 和
\(BB'\) 确定平面 \(\beta\)（\(b\) 与 \(BB'\) 交于 \(B\)
且不共线，故确定唯一平面）。由于 \(BB' \perp \alpha\) 且
\(BB' \subset \beta\)，由面面垂直的判定定理得 \(\beta \perp \alpha\)。

\textbf{(3) 设 \(l = \alpha \cap \beta\)，论证 \(l\) 与 \(b\)
相交}：由于 \(b \subset \beta\)，\(l = \alpha \cap \beta\)，故 \(l\) 与
\(b\) 都在平面 \(\beta\) 内。若 \(l \parallel b\)，则由
\(b \parallel \alpha\) 得 \(l \parallel \alpha\)，但
\(l \subset \alpha\)，矛盾。故 \(l\) 与 \(b\) 必相交于一点 \(Q\)。

\textbf{(4) 严格证明 \(PQ \perp b\)}：过 \(Q\) 在 \(\alpha\) 内作
\(PQ \perp l\)，\(P \in a\)（\(PQ\) 与 \(a\) 相交于 \(P\)，因为
\(l \parallel a\) 且两者均在 \(\alpha\) 内，\(PQ\) 与 \(a\)
的公垂距离有限）。因 \(PQ \perp l\) 且 \(l \parallel a\)，故
\(PQ \perp a\)。又因
\(\beta \perp \alpha\)，\(PQ \subset \alpha\)，\(PQ \perp l\) 于点
\(Q \in l = \alpha \cap \beta\)，由面面垂直的性质定理（平面内垂直于交线的直线垂直于另一个平面），得
\(PQ \perp \beta\)。因此 \(PQ\) 垂直于 \(\beta\) 中过 \(Q\)
的所有直线，特别地 \(PQ \perp b\)。故 \(PQ\) 是 \(a\) 与 \(b\)
的公垂线段。\(\blacksquare\)

\subsubsection{14.3 简单几何体}\label{ux7b80ux5355ux51e0ux4f55ux4f53}

\paragraph{14.3.1
棱柱、棱锥、棱台}\label{ux68f1ux67f1ux68f1ux9525ux68f1ux53f0}

\textbf{体积公理（祖暅原理/卡瓦列里原理）}：设两个立体 \(\mathcal{S}_1\)
和 \(\mathcal{S}_2\)
位于同一对平行平面之间。若用任意平行于这两个平面的平面截
\(\mathcal{S}_1\) 和 \(\mathcal{S}_2\) 所得的截面积总是相等，则
\(\mathcal{S}_1\) 和 \(\mathcal{S}_2\) 的体积相等。

\textbf{注}：此原理可视为三维体积理论的基本公理，是面积公理 S1-S4
在三维空间的推广。在现代数学中，其严格证明需要测度论（参见实分析教材）。

\textbf{定义
14.15}（棱柱）：有两个面互相平行，其余各面都是四边形，并且每相邻两个四边形的公共边都互相平行的多面体称为\textbf{棱柱}。互相平行的两个面称为\textbf{底面}，其余各面称为\textbf{侧面}，相邻侧面的公共边称为\textbf{侧棱}。

\textbf{定义
14.16}（直棱柱、正棱柱）：侧棱垂直于底面的棱柱称为\textbf{直棱柱}。底面是正多边形的直棱柱称为\textbf{正棱柱}。

\textbf{定理 14.16}（棱柱的体积）：棱柱的体积等于底面积乘以高，即
\(V = Sh\)。

\textbf{证明}：

设棱柱的底面面积为 \(S\)，高为 \(h\)（两底面之间的距离）。

将棱柱沿平行于底面的平面切片。对于任意高度
\(z\)（\(0 \leq z \leq h\)），截面是一个与底面全等的多边形，面积为
\(S\)。

考虑一个底面积为 \(S\)、高为 \(h\) 的直棱柱（柱体）。该直棱柱在每个高度
\(z\) 处的截面面积也是 \(S\)。

由祖暅原理（卡瓦列里原理）：如果两个立体在每一个平行于底面的截面上面积相等，则这两个立体体积相等。棱柱的体积等于该直棱柱的体积。

对于直棱柱，可以将其底面（\(n\) 边形）分解为 \(n-2\)
个三角形，从而将直棱柱分解为 \(n-2\)
个三棱柱。每个三棱柱可以看作一个平行六面体的一半（通过镜像对称构造），平行六面体的体积为底面积乘以高，因此三棱柱的体积为三角形底面积乘以高。所有三棱柱的体积之和为
\((\sum S_i) h = Sh\)。

因此棱柱的体积 \(V = Sh\)。\(\blacksquare\)

\textbf{定义
14.17}（棱锥）：有一个面是多边形，其余各面是有一个公共顶点的三角形的多面体称为\textbf{棱锥}。多边形面称为\textbf{底面}，三角形面称为\textbf{侧面}，公共顶点称为\textbf{顶点}，顶点到底面的距离称为\textbf{高}。

\textbf{定理
14.17}（棱锥的体积）：棱锥的体积等于底面积乘以高的三分之一，即
\(V = \frac{1}{3}Sh\)。

\textbf{证明}：

首先证明三棱锥（四面体）的体积为 \(V = \frac{1}{3}Sh\)。

考虑三棱柱 \(ABC-A'B'C'\)。连接
\(AB'\)、\(AC'\)、\(B'C\)，将三棱柱分为三个三棱锥：\(T_1 = ABC-B'\)、\(T_2 = A'-AB'C'\)、\(T_3 = B'-A'C'C\)。

\textbf{\(T_1\) 与 \(T_2\) 等体积}：\(T_1\) 以 \(\triangle ABC\)
为底、\(B'\) 为顶点，高为 \(B'\) 到平面 \(ABC\) 的距离（即棱柱的高
\(h\)）。\(T_2\) 以 \(\triangle AB'C'\) 为底、\(A'\) 为顶点，高为 \(A'\)
到平面 \(AB'C'\) 的距离（也是
\(h\)）。由三棱柱的性质，\(\triangle ABC \cong \triangle AB'C'\)，故
\(T_1\) 和 \(T_2\) 的底面积相等、高相等，体积相等。

\textbf{\(T_2\) 与 \(T_3\) 等体积}：\(T_2\) 以 \(\triangle AB'C'\)
为底、\(A'\) 为顶点。\(T_3\) 以 \(\triangle A'C'C\) 为底、\(B'\)
为顶点。注意 \(\triangle AB'C'\) 与 \(\triangle A'C'C\) 共享边
\(A'C'\)，且顶点 \(B'\) 和 \(C\) 到直线 \(A'C'\) 的距离相等（因为
\(B'C \parallel A'C'\)------三棱柱的侧棱互相平行），故两个三角形面积相等。\(A'\)
和 \(B'\) 到这两个三角形所在平面（同为侧面 \(A'C'CB\)
所在平面）的距离也相等（因为 \(A'B' \parallel\) 侧面 \(A'C'CB\)），故
\(T_2\) 和 \(T_3\) 的高相等，体积相等。

因此三个三棱锥体积相等，每个的体积为三棱柱体积的 \(\frac{1}{3}\)，即
\(\frac{1}{3}Sh\)。

对于一般的棱锥，将其底面分解为三角形，对应多个三棱锥，每个体积为
\(\frac{1}{3} S_i h\)。总和为
\(\frac{1}{3}(\sum S_i) h = \frac{1}{3}Sh\)。\(\blacksquare\)

\textbf{定义
14.18}（棱台）：用一个平行于棱锥底面的平面去截棱锥，底面与截面之间的部分称为\textbf{棱台}。

\textbf{定理 14.18}（棱台的体积）：设棱台的上底面积为
\(S_1\)，下底面积为 \(S_2\)，高为 \(h\)，则体积为：
\[V = \frac{1}{3}h(S_1 + S_2 + \sqrt{S_1 S_2})\]

\textbf{证明}：

设原棱锥底面积为 \(S_2\)，高为 \(H\)。截面（上底）面积为
\(S_1\)，距顶点的距离为 \(h_1\)。则 \(H = h_1 + h\)。

由相似关系，\(\frac{S_1}{S_2} = \left(\frac{h_1}{H}\right)^2\)。

棱台体积 = 大棱锥体积 - 小棱锥体积：
\[V = \frac{1}{3}S_2 H - \frac{1}{3}S_1 h_1 = \frac{1}{3}(S_2 H - S_1 h_1)\]

由
\(h_1 = H\sqrt{\frac{S_1}{S_2}}\)，\(h = H - h_1 = H\left(1 - \sqrt{\frac{S_1}{S_2}}\right)\)，所以
\(H = \frac{h}{1 - \sqrt{S_1/S_2}}\)。

设 \(k = \sqrt{S_1/S_2}\)，则 \(S_1 = k^2 S_2\)，\(h_1 = kH\)。

\[V = \frac{1}{3}(S_2 H - k^2 S_2 \cdot kH) = \frac{S_2 H}{3}(1 - k^3)\]

\[= \frac{S_2}{3} \cdot \frac{h}{1 - k} \cdot (1 - k^3) = \frac{S_2 h}{3} \cdot (1 + k + k^2)\]

\[= \frac{h}{3}(S_2 + kS_2 + k^2 S_2) = \frac{h}{3}(S_2 + \sqrt{S_1 S_2} + S_1)\]

\(\blacksquare\)

\paragraph{14.3.2
圆柱、圆锥、圆台}\label{ux5706ux67f1ux5706ux9525ux5706ux53f0}

\textbf{定义
14.19}（圆柱）：以矩形的一边所在直线为旋转轴，其余三边旋转形成的曲面所围成的几何体称为\textbf{圆柱}。

\textbf{定理 14.19}（圆柱的体积和表面积）：底面半径为 \(r\)、高为 \(h\)
的圆柱： - 体积：\(V = \pi r^2 h\) -
侧面积：\(S_{\text{侧}} = 2\pi r h\) -
表面积：\(S = 2\pi r^2 + 2\pi r h = 2\pi r(r + h)\)

\textbf{证明}：

\textbf{体积}：圆柱可以看作是底面为圆的棱柱的极限（底面圆内接正 \(n\)
边形棱柱当 \(n \to \infty\)）。底面圆的面积为 \(\pi r^2\)，所以
\(V = \pi r^2 h\)。

也可以用祖暅原理直接证明：在高度 \(z\) 处的截面为与底面全等的圆，面积为
\(\pi r^2\)，与高度无关。

\textbf{侧面积}：将侧面沿一条母线剪开，展开为矩形。矩形的长为底面圆的周长
\(2\pi r\)，宽为 \(h\)。所以 \(S_{\text{侧}} = 2\pi r h\)。

\textbf{表面积}：\(S = S_{\text{侧}} + 2S_{\text{底}} = 2\pi rh + 2\pi r^2 = 2\pi r(r + h)\)。\(\blacksquare\)

\textbf{定义
14.20}（圆锥）：以直角三角形的一条直角边所在直线为旋转轴，其余两边旋转形成的曲面所围成的几何体称为\textbf{圆锥}。

\textbf{定理 14.20}（圆锥的体积和表面积）：底面半径为 \(r\)、高为
\(h\)、母线长为 \(l = \sqrt{r^2 + h^2}\) 的圆锥： -
体积：\(V = \frac{1}{3}\pi r^2 h\) - 侧面积：\(S_{\text{侧}} = \pi r l\)
- 表面积：\(S = \pi r^2 + \pi r l = \pi r(r + l)\)

\textbf{证明}：

\textbf{体积}：利用祖暅原理。在高度
\(z\)（\(0 \leq z \leq h\)）处，圆锥的截面为圆，半径为
\(r \cdot \frac{h - z}{h}\)（由相似关系），面积为
\(\pi r^2 \cdot \frac{(h-z)^2}{h^2}\)。

考虑一个底面面积为 \(\pi r^2\)、高为 \(h\)
的棱锥（例如底面为正方形的四棱锥，正方形边长为
\(\sqrt{\pi}\, r\)）。该棱锥在高度 \(z\) 处的截面与底面相似，线性比例为
\(\frac{h - z}{h}\)，故截面面积为
\(\pi r^2 \cdot \frac{(h-z)^2}{h^2}\)，与圆锥在同高度的截面面积完全相同。

由祖暅原理，等截面处面积相等的两个立体体积相等。而棱锥体积为
\(\frac{1}{3} \times \text{底面积} \times \text{高} = \frac{1}{3}\pi r^2 h\)，因此圆锥体积
\(V = \frac{1}{3}\pi r^2 h\)。

\textbf{侧面积}：将侧面沿一条母线剪开，展开为扇形。扇形的半径为
\(l\)，弧长为 \(2\pi r\)。扇形面积
\(= \frac{1}{2} \times l \times 2\pi r = \pi r l\)。

\textbf{表面积}：\(S = \pi rl + \pi r^2 = \pi r(r + l)\)。\(\blacksquare\)

\textbf{定义
14.21}（圆台）：以直角梯形的垂直于底边的腰所在直线为旋转轴，其余三边旋转形成的曲面所围成的几何体称为\textbf{圆台}。

\textbf{定理 14.21}（圆台的体积和表面积）：上底半径为
\(r_1\)、下底半径为 \(r_2\)、高为 \(h\)、母线长为 \(l\) 的圆台： -
体积：\(V = \frac{1}{3}\pi h(r_1^2 + r_2^2 + r_1 r_2)\) -
侧面积：\(S_{\text{侧}} = \pi(r_1 + r_2)l\) -
表面积：\(S = \pi r_1^2 + \pi r_2^2 + \pi(r_1 + r_2)l\)

\textbf{证明}：

\textbf{体积}：由棱台体积公式的推广，\(V = \frac{1}{3}\pi h(r_1^2 + r_1 r_2 + r_2^2)\)。

严格地，将圆台看作大圆锥减去小圆锥。设大圆锥高为 \(H\)，底面半径
\(r_2\)，则小圆锥高为 \(H - h\)，底面半径
\(r_1 = r_2 \cdot \frac{H - h}{H}\)。由此解出
\(H = \frac{r_2 h}{r_2 - r_1}\)。

\[V = \frac{1}{3}\pi r_2^2 H - \frac{1}{3}\pi r_1^2(H - h) = \frac{\pi}{3}\left[r_2^2 H - r_1^2(H - h)\right]\]

代入
\(H = \frac{r_2 h}{r_2 - r_1}\)，\(H - h = \frac{r_1 h}{r_2 - r_1}\)：

\[V = \frac{\pi}{3} \cdot \frac{h}{r_2 - r_1}\left(r_2^3 - r_1^3\right) = \frac{\pi h}{3} \cdot \frac{(r_2 - r_1)(r_2^2 + r_1 r_2 + r_1^2)}{r_2 - r_1} = \frac{1}{3}\pi h(r_1^2 + r_1 r_2 + r_2^2)\]

\textbf{侧面积}：将侧面沿母线展开，得到扇环形。面积
\(= \pi(r_1 + r_2)l\)。

\textbf{表面积}：\(S = \pi r_1^2 + \pi r_2^2 + \pi(r_1 + r_2)l\)。\(\blacksquare\)

\paragraph{14.3.3
球的体积和表面积}\label{ux7403ux7684ux4f53ux79efux548cux8868ux9762ux79ef}

\textbf{定义 14.22}（球）：空间中到定点 \(O\) 的距离等于定长 \(r\)
的所有点的集合称为\textbf{球面}，记为
\(S(O, r)\)。球面及其内部所围成的几何体称为\textbf{球体}，简称\textbf{球}。

\textbf{定理 14.22}（球的体积）：半径为 \(R\) 的球的体积为
\(V = \frac{4}{3}\pi R^3\)。

\textbf{证明}（祖暅原理法）：

将球放在一个底面半径和高都为 \(R\) 的圆柱中（球内切于圆柱）。在高度
\(h\)（\(-R \leq h \leq R\)）处，作平行于底面的截面。

球的截面是一个圆，半径为 \(\sqrt{R^2 - h^2}\)，面积为
\(\pi(R^2 - h^2)\)。

圆柱的截面也是圆，半径为 \(R\)，面积为 \(\pi R^2\)。

从圆柱中挖去两个底面半径和高都为 \(R\)
的圆锥（顶点在圆柱的两个底面中心），剩余部分在高度 \(h\)
处的截面面积为： \[\pi R^2 - \pi h^2 = \pi(R^2 - h^2)\]

（圆锥在高度 \(h\) 处的截面半径为 \(|h|\)，面积为 \(\pi h^2\)。）

球的截面面积 \(\pi(R^2 - h^2)\) 等于圆柱减去两个圆锥后剩余部分的截面面积
\(\pi(R^2 - h^2)\)。

由祖暅原理，球的体积等于圆柱减去两个圆锥的体积：
\[V = \pi R^2 \cdot 2R - 2 \cdot \frac{1}{3}\pi R^2 \cdot R = 2\pi R^3 - \frac{2}{3}\pi R^3 = \frac{4}{3}\pi R^3\]

\(\blacksquare\)

\textbf{定理 14.23}（球的表面积）：半径为 \(R\) 的球的表面积为
\(S = 4\pi R^2\)。

\textbf{证明}（逼近法）：

\textbf{注}：以下证明使用微积分方法，其严格基础将在卷四建立。

将球面分为许多窄的球带。在与球心距离为 \(h\) 处，取厚度为
\(\mathrm{d}h\) 的球带。

球带的半径为 \(r(h) = \sqrt{R^2 - h^2}\)，球带的宽度为
\(\mathrm{d}s = \frac{R}{\sqrt{R^2 - h^2}} \mathrm{d}h\)（由弧长公式）。

球带的面积为
\(2\pi r \cdot \mathrm{d}s = 2\pi \sqrt{R^2 - h^2} \cdot \frac{R}{\sqrt{R^2 - h^2}} \mathrm{d}h = 2\pi R \, \mathrm{d}h\)。

注意：球带面积与 \(h\) 无关！每一个厚度为 \(\mathrm{d}h\) 的球带面积都是
\(2\pi R \, \mathrm{d}h\)。

对所有球带求和（从 \(h = -R\) 到 \(h = R\)）：

\[S = \int_{-R}^{R} 2\pi R \, \mathrm{d}h = 2\pi R \cdot 2R = 4\pi R^2\]

严格证明参见卷四定理 24.16（微积分基本定理）的相关应用。

\(\blacksquare\)

\textbf{推论 14.2}：球的体积与表面积之间存在关系
\(S = \frac{\mathrm{d}V}{\mathrm{d}R}\)，即表面积等于体积对半径的导数。

\textbf{证明}：\(\frac{\mathrm{d}}{\mathrm{d}R}\left(\frac{4}{3}\pi R^3\right) = 4\pi R^2 = S\)。\(\blacksquare\)

这体现了微积分中''壳层法''的几何直观------将球看作由无穷薄的球壳叠加而成，每层球壳的体积为
\(S(r) \cdot \mathrm{d}r\)。

\paragraph{14.3.4 祖暅原理}\label{ux7956ux6685ux539fux7406}

祖暅原理是推导各种几何体体积公式的核心工具，值得单独阐述。

\textbf{注} 祖暅原理在本体系中作为体积公理（见第 14
章开头）引入。在现代数学中，此原理的严格基础由测度论和积分理论提供。

\paragraph{14.3.5 欧拉公式}\label{ux6b27ux62c9ux516cux5f0f}

\textbf{定理 14.25}（欧拉多面体公式）：对于任意凸多面体，设其顶点数为
\(V\)，棱数为 \(E\)，面数为 \(F\)，则

\[V - E + F = 2\]

\textbf{证明概要}：

将凸多面体的一个面''展开''到平面上，得到一个平面图。具体地，移除多面体的一个面，将剩余部分沿棱展开到平面上，得到一个连通的平面图。

对于连通的平面图，由图论中的归纳法可以证明
\(V - E + F = 1\)（展开后少了一个面）。因此原多面体的
\(V - E + F = 1 + 1 = 2\)。

具体地：对平面连通图归纳。\textbf{基础步}：单个顶点，\(V = 1, E = 0, F = 1\)，\(V - E + F = 2\)。\textbf{归纳步}：设对
\(k\) 条边的连通图成立。对 \(k+1\)
条边的图，若存在有界面（多边形），移除其一条边界边（非桥），\(E\) 减
1，\(F\) 减 1，\(V\) 不变，\(V - E + F\) 不变，由归纳假设为
2。若无有界面（即图为树），则
\(E = V - 1\)，\(F = 1\)，\(V - E + F = V - (V-1) + 1 = 2\)。

归纳法的关键步骤：对于平面图中的任意一个有界面（多边形），移除其一条边（不删除顶点），则
\(E\) 减少 \(1\)，\(F\) 减少 \(1\)，\(V\) 不变，\(V - E + F\)
不变。反复移除边直到只剩一个无界面，此时得到一棵生成树，\(E = V - 1\)，\(F = 1\)，\(V - E + F = V - (V-1) + 1 = 2\)。\(\blacksquare\)

\textbf{验证}：正四面体：\(V = 4\)，\(E = 6\)，\(F = 4\)，\(4 - 6 + 4 = 2\)。立方体：\(V = 8\)，\(E = 12\)，\(F = 6\)，\(8 - 12 + 6 = 2\)。正十二面体：\(V = 20\)，\(E = 30\)，\(F = 12\)，\(20 - 30 + 12 = 2\)。

\bigskip\hrule\bigskip

\textbf{本章展望}：本章系统建立了立体几何的理论体系，从线面平行和线面垂直的判定与性质出发，经过面面垂直、三垂线定理等核心定理，到异面直线角、线面角、二面角和各种距离的概念，最后建立了棱柱、棱锥、棱台、圆柱、圆锥、圆台和球的体积与表面积公式。然而，上述所有推导都是纯几何的方法。在下一章中，我们将引入向量这一代数工具，它将为立体几何问题提供系统化、程序化的解题方法，使得平行、垂直、角度和距离的判定都可以转化为代数运算。

\bigskip\hrule\bigskip

\subsection{Ch 15 向量}\label{ch-15-ux5411ux91cf}

\textbf{前置知识}：本章假定读者已掌握 Ch 13
平面几何（特别是平行四边形性质、三角形相似和勾股定理）以及 Ch 14
立体几何（线面平行、线面垂直、面面垂直、异面直线角、线面角、二面角、点到平面距离）。向量将为这些几何概念提供统一的代数刻画。

\bigskip\hrule\bigskip

前两章中，我们用纯几何的方法建立了平面几何和立体几何的完整理论。然而，纯几何方法在处理复杂的空间问题时往往需要精巧的辅助线构造，缺乏系统性和可操作性。本章引入向量------一种兼具大小和方向的数学对象------作为连接几何与代数的桥梁。向量方法的核心思想是：将几何关系（平行、垂直、角度、距离）转化为代数运算（加法、数量积、向量积），从而将几何证明转化为代数计算。这一思想的深远意义在于：它为几何学提供了一套统一的、可算法化的处理框架。

\subsubsection{15.1 平面向量}\label{ux5e73ux9762ux5411ux91cf}

\paragraph{15.1.0
从几何直觉到代数抽象}\label{ux4eceux51e0ux4f55ux76f4ux89c9ux5230ux4ee3ux6570ux62bdux8c61}

向量的概念源于物理学中的力、速度和位移等量，这些量不仅有大小，还有方向。数学上的向量将这一直觉抽象化，使得我们可以在纯粹的代数框架下处理几何问题。

\paragraph{15.1.1
向量的基本概念}\label{ux5411ux91cfux7684ux57faux672cux6982ux5ff5}

\textbf{定义
15.1}（向量）：既有大小又有方向的量称为\textbf{向量}（或矢量）。向量用带箭头的有向线段表示，如
\(\overrightarrow{AB}\) 或
\(\vec{a}\)。向量的大小（长度）称为向量的\textbf{模}，记作 \(|\vec{a}|\)
或 \(|\overrightarrow{AB}|\)。

\textbf{定义 15.2}（特殊向量）： - \textbf{零向量}：模为 \(0\)
的向量，记作 \(\vec{0}\)。零向量的方向不确定。 - \textbf{单位向量}：模为
\(1\) 的向量。与非零向量 \(\vec{a}\) 同方向的单位向量为
\(\vec{e}_a = \frac{\vec{a}}{|\vec{a}|}\)。

\textbf{定义
15.3}（向量相等）：长度相等且方向相同的向量称为\textbf{相等向量}。

\textbf{定义
15.4}（共线向量）：方向相同或相反的非零向量称为\textbf{共线向量}（或平行向量）。规定零向量与任何向量共线。

\paragraph{15.1.2 向量的运算}\label{ux5411ux91cfux7684ux8fd0ux7b97}

\textbf{定义 15.5}（向量加法）：设 \(\vec{a}\) 和 \(\vec{b}\)
是两个向量。在平面上取一点 \(O\)，作
\(\overrightarrow{OA} = \vec{a}\)，\(\overrightarrow{AB} = \vec{b}\)，则
\(\overrightarrow{OB} = \vec{a} + \vec{b}\)。此法则称为\textbf{三角形法则}。

等价地，作
\(\overrightarrow{OA} = \vec{a}\)，\(\overrightarrow{OC} = \vec{b}\)，以
\(OA\)、\(OC\) 为邻边作平行四边形 \(OABC\)，则对角线
\(\overrightarrow{OB} = \vec{a} + \vec{b}\)。此法则称为\textbf{平行四边形法则}。

\textbf{定理 15.1}（向量加法的交换律）：对任意向量
\(\vec{a}\)、\(\vec{b}\)，\(\vec{a} + \vec{b} = \vec{b} + \vec{a}\)。

\textbf{证明}：作
\(\overrightarrow{OA} = \vec{a}\)，\(\overrightarrow{AB} = \vec{b}\)，则
\(\overrightarrow{OB} = \vec{a} + \vec{b}\)。作
\(\overrightarrow{OC} = \vec{b}\)，\(\overrightarrow{CB} = \vec{a}\)（平行四边形对边平行且相等），则
\(\overrightarrow{OB} = \vec{b} + \vec{a}\)。因此
\(\vec{a} + \vec{b} = \vec{b} + \vec{a}\)。\(\blacksquare\)

\textbf{定理 15.2}（向量加法的结合律）：对任意向量
\(\vec{a}\)、\(\vec{b}\)、\(\vec{c}\)，\((\vec{a} + \vec{b}) + \vec{c} = \vec{a} + (\vec{b} + \vec{c})\)。

\textbf{证明}：作
\(\overrightarrow{OA} = \vec{a}\)，\(\overrightarrow{AB} = \vec{b}\)，\(\overrightarrow{BC} = \vec{c}\)。则
\((\vec{a} + \vec{b}) + \vec{c} = \overrightarrow{OB} + \overrightarrow{BC} = \overrightarrow{OC}\)，\(\vec{a} + (\vec{b} + \vec{c}) = \overrightarrow{OA} + \overrightarrow{AC} = \overrightarrow{OC}\)。因此
\((\vec{a} + \vec{b}) + \vec{c} = \vec{a} + (\vec{b} + \vec{c})\)。\(\blacksquare\)

\textbf{定义 15.6}（数乘向量）：设 \(\lambda\) 为实数，\(\vec{a}\)
为向量，定义 \(\lambda\vec{a}\) 为： - 若
\(\lambda > 0\)，\(\lambda\vec{a}\) 与 \(\vec{a}\)
同方向，\(|\lambda\vec{a}| = \lambda|\vec{a}|\)； - 若
\(\lambda = 0\)，\(\lambda\vec{a} = \vec{0}\)； - 若
\(\lambda < 0\)，\(\lambda\vec{a}\) 与 \(\vec{a}\)
反方向，\(|\lambda\vec{a}| = |\lambda||\vec{a}|\)。

\textbf{定理 15.3}（数乘的性质）：对实数 \(\lambda, \mu\) 和向量
\(\vec{a}, \vec{b}\)： 1. \(\lambda(\mu\vec{a}) = (\lambda\mu)\vec{a}\)
2. \((\lambda + \mu)\vec{a} = \lambda\vec{a} + \mu\vec{a}\) 3.
\(\lambda(\vec{a} + \vec{b}) = \lambda\vec{a} + \lambda\vec{b}\)

\textbf{证明}：这些性质均可以从数乘的定义和几何意义直接验证。以分配律
(3) 为例：设
\(\overrightarrow{OA} = \vec{a}\)，\(\overrightarrow{OB} = \vec{b}\)，以
\(OA\)、\(OB\) 为邻边作平行四边形 \(OACB\)，则
\(\overrightarrow{OC} = \vec{a} + \vec{b}\)。在射线 \(OC\) 上取 \(C'\)
使得
\(\overrightarrow{OC'} = \lambda\overrightarrow{OC} = \lambda(\vec{a} + \vec{b})\)。由相似关系，\(C'\)
也是以 \(\lambda\vec{a}\) 和 \(\lambda\vec{b}\)
为邻边的平行四边形的对角线的终点，故
\(\overrightarrow{OC'} = \lambda\vec{a} + \lambda\vec{b}\)。\(\blacksquare\)

\textbf{定义
15.7}（向量减法）：\(\vec{a} - \vec{b} = \vec{a} + (-\vec{b})\)，其中
\(-\vec{b}\) 是与 \(\vec{b}\) 方向相反、模相等的向量。

\textbf{几何意义}：作
\(\overrightarrow{OA} = \vec{a}\)，\(\overrightarrow{OB} = \vec{b}\)，则
\(\overrightarrow{BA} = \vec{a} - \vec{b}\)。

\paragraph{15.1.3
向量的坐标表示}\label{ux5411ux91cfux7684ux5750ux6807ux8868ux793a}

\textbf{定义 15.8}（基底）：平面上两个不共线的向量
\(\vec{e}_1, \vec{e}_2\) 称为一组\textbf{基底}。对于平面上任意向量
\(\vec{a}\)，存在唯一的实数对 \((x, y)\) 使得：
\[\vec{a} = x\vec{e}_1 + y\vec{e}_2\]

\((x, y)\) 称为 \(\vec{a}\) 在基底 \(\{\vec{e}_1, \vec{e}_2\}\)
下的\textbf{坐标}。

当 \(\vec{e}_1 = (1, 0)\)，\(\vec{e}_2 = (0, 1)\)
为标准基底时，\(\vec{a} = (x, y)\)。

\textbf{存在性证明}：设 \(\vec{e}_1, \vec{e}_2\) 不共线，\(\vec{a}\)
为平面上任意向量。将 \(\vec{e}_1, \vec{e}_2\) 放在同一起点 \(O\)。过
\(\vec{a}\) 的终点 \(A\) 作平行于 \(\vec{e}_2\) 的直线，交 \(\vec{e}_1\)
所在直线于点 \(P\)。则 \(\overrightarrow{OP} = x\vec{e}_1\)（某实数
\(x\)），\(\overrightarrow{PA}\) 平行于 \(\vec{e}_2\)，故
\(\overrightarrow{PA} = y\vec{e}_2\)（某实数 \(y\)）。由
\(\vec{a} = \overrightarrow{OA} = \overrightarrow{OP} + \overrightarrow{PA} = x\vec{e}_1 + y\vec{e}_2\)，存在性得证。\(\blacksquare\)

\textbf{定理 15.4}（坐标的唯一性）：在给定基底
\(\{\vec{e}_1, \vec{e}_2\}\) 下，任意向量的坐标是唯一的。

\textbf{证明}：设
\(\vec{a} = x\vec{e}_1 + y\vec{e}_2 = x'\vec{e}_1 + y'\vec{e}_2\)。则
\((x - x')\vec{e}_1 + (y - y')\vec{e}_2 = \vec{0}\)。由于 \(\vec{e}_1\)
和 \(\vec{e}_2\) 不共线，若 \(x - x' \neq 0\) 或 \(y - y' \neq 0\)，则
\(\vec{0}\) 可以表示为非零向量的线性组合，这意味着 \(\vec{e}_1\) 和
\(\vec{e}_2\) 共线，矛盾。因此 \(x = x'\)，\(y = y'\)。\(\blacksquare\)

\textbf{运算的坐标表示}：设
\(\vec{a} = (x_1, y_1)\)，\(\vec{b} = (x_2, y_2)\)，\(\lambda\) 为实数：

\[\vec{a} + \vec{b} = (x_1 + x_2, y_1 + y_2)\]
\[\vec{a} - \vec{b} = (x_1 - x_2, y_1 - y_2)\]
\[\lambda\vec{a} = (\lambda x_1, \lambda y_1)\]
\[|\vec{a}| = \sqrt{x_1^2 + y_1^2}\]

\textbf{定理 15.5}（共线的充要条件）：设
\(\vec{a} = (x_1, y_1)\)，\(\vec{b} = (x_2, y_2)\)，\(\vec{b} \neq \vec{0}\)，则
\(\vec{a}\) 与 \(\vec{b}\) 共线的充要条件是 \(x_1 y_2 - x_2 y_1 = 0\)。

\textbf{证明}：\(\vec{a}\) 与 \(\vec{b}\) 共线 \(\iff\) 存在实数
\(\lambda\) 使得 \(\vec{a} = \lambda\vec{b}\) \(\iff\)
\((x_1, y_1) = (\lambda x_2, \lambda y_2)\) \(\iff\)
\(x_1 = \lambda x_2\) 且 \(y_1 = \lambda y_2\)。

若 \(x_2 \neq 0\)，则 \(\lambda = \frac{x_1}{x_2}\)，代入得
\(y_1 = \frac{x_1}{x_2} y_2\)，即 \(x_1 y_2 = x_2 y_1\)，即
\(x_1 y_2 - x_2 y_1 = 0\)。

若 \(x_2 = 0\)，则 \(x_1 = \lambda \cdot 0 = 0\)，且
\(y_1 = \lambda y_2\)，此时
\(x_1 y_2 - x_2 y_1 = 0 \cdot y_2 - 0 \cdot y_1 = 0\)。\(\blacksquare\)

\paragraph{15.1.4
平面向量的数量积}\label{ux5e73ux9762ux5411ux91cfux7684ux6570ux91cfux79ef}

\textbf{定义 15.9}（数量积）：设 \(\vec{a}\) 和 \(\vec{b}\)
是两个非零向量，\(\theta\)
为它们的夹角（\(0 \leq \theta \leq \pi\)），定义：

\[\vec{a} \cdot \vec{b} = |\vec{a}||\vec{b}|\cos\theta\]

规定零向量与任何向量的数量积为 \(0\)。

\textbf{定理 15.6}（数量积的坐标公式）：设
\(\vec{a} = (x_1, y_1)\)，\(\vec{b} = (x_2, y_2)\)，则
\(\vec{a} \cdot \vec{b} = x_1 x_2 + y_1 y_2\)。

\textbf{证明}：设标准基底 \(\vec{i} = (1, 0)\)，\(\vec{j} = (0, 1)\)。则
\(\vec{a} = x_1\vec{i} + y_1\vec{j}\)，\(\vec{b} = x_2\vec{i} + y_2\vec{j}\)。

\[\vec{a} \cdot \vec{b} = (x_1\vec{i} + y_1\vec{j}) \cdot (x_2\vec{i} + y_2\vec{j})\]

利用分配律：
\[= x_1 x_2(\vec{i} \cdot \vec{i}) + x_1 y_2(\vec{i} \cdot \vec{j}) + y_1 x_2(\vec{j} \cdot \vec{i}) + y_1 y_2(\vec{j} \cdot \vec{j})\]

由于
\(\vec{i} \cdot \vec{i} = 1\)，\(\vec{j} \cdot \vec{j} = 1\)，\(\vec{i} \cdot \vec{j} = 1 \cdot 1 \cdot \cos 90° = 0\)：

\[\vec{a} \cdot \vec{b} = x_1 x_2 + y_1 y_2\]

\(\blacksquare\)

\textbf{定理 15.7}（数量积的性质）：设 \(\vec{a}, \vec{b}, \vec{c}\)
为向量，\(\lambda \in \mathbb{R}\)，则：

\begin{enumerate}
\def\labelenumi{\arabic{enumi}.}
\tightlist
\item
  \textbf{交换律}：\(\vec{a} \cdot \vec{b} = \vec{b} \cdot \vec{a}\)
\item
  \textbf{分配律}：\(\vec{a} \cdot (\vec{b} + \vec{c}) = \vec{a} \cdot \vec{b} + \vec{a} \cdot \vec{c}\)
\item
  \textbf{数乘结合律}：\((\lambda\vec{a}) \cdot \vec{b} = \lambda(\vec{a} \cdot \vec{b})\)
\item
  \textbf{正定性}：\(\vec{a} \cdot \vec{a} = |\vec{a}|^2 \geq 0\)，等号当且仅当
  \(\vec{a} = \vec{0}\)
\end{enumerate}

\textbf{证明}：

\begin{enumerate}
\def\labelenumi{(\arabic{enumi})}
\item
  \(\vec{a} \cdot \vec{b} = |\vec{a}||\vec{b}|\cos\theta = |\vec{b}||\vec{a}|\cos\theta = \vec{b} \cdot \vec{a}\)。
\item
  由坐标公式，设
  \(\vec{a} = (a_1, a_2)\)，\(\vec{b} = (b_1, b_2)\)，\(\vec{c} = (c_1, c_2)\)：
\end{enumerate}

\[\vec{a} \cdot (\vec{b} + \vec{c}) = a_1(b_1 + c_1) + a_2(b_2 + c_2) = (a_1 b_1 + a_2 b_2) + (a_1 c_1 + a_2 c_2) = \vec{a} \cdot \vec{b} + \vec{a} \cdot \vec{c}\]

\begin{enumerate}
\def\labelenumi{(\arabic{enumi})}
\setcounter{enumi}{2}
\item
  \((\lambda\vec{a}) \cdot \vec{b} = (\lambda a_1)b_1 + (\lambda a_2)b_2 = \lambda(a_1 b_1 + a_2 b_2) = \lambda(\vec{a} \cdot \vec{b})\)。
\item
  \(\vec{a} \cdot \vec{a} = x_1^2 + y_1^2 = |\vec{a}|^2 \geq 0\)，等号当且仅当
  \(x_1 = y_1 = 0\)，即 \(\vec{a} = \vec{0}\)。\(\blacksquare\)
\end{enumerate}

\textbf{推论
15.1}（夹角公式）：\(\cos\theta = \dfrac{\vec{a} \cdot \vec{b}}{|\vec{a}||\vec{b}|} = \dfrac{x_1 x_2 + y_1 y_2}{\sqrt{x_1^2 + y_1^2}\sqrt{x_2^2 + y_2^2}}\)。

\textbf{证明}：由数量积定义
\(\vec{a} \cdot \vec{b} = |\vec{a}||\vec{b}|\cos\theta\) 和定理 15.6 中
\(\vec{a} \cdot \vec{b} = x_1x_2 + y_1y_2\)，立即得到结论。\(\blacksquare\)

\textbf{推论 15.2}（垂直的充要条件）：非零向量 \(\vec{a} \perp \vec{b}\)
的充要条件是 \(\vec{a} \cdot \vec{b} = 0\)，即
\(x_1 x_2 + y_1 y_2 = 0\)。

\textbf{证明}：\(\vec{a} \perp \vec{b} \iff \theta = 90° \iff \cos\theta = 0 \iff |\vec{a}||\vec{b}|\cos\theta = 0 \iff \vec{a} \cdot \vec{b} = 0\)。\(\blacksquare\)

\subsubsection{15.2 空间向量}\label{ux7a7aux95f4ux5411ux91cf}

\paragraph{15.2.1
空间向量的概念与运算}\label{ux7a7aux95f4ux5411ux91cfux7684ux6982ux5ff5ux4e0eux8fd0ux7b97}

\textbf{定义
15.10}（空间向量）：在空间中，既有大小又有方向的量称为空间向量。空间向量用有向线段表示，记为
\(\overrightarrow{AB}\) 或 \(\vec{a}\)。

\textbf{定义 15.11}（空间向量的坐标表示）：在空间直角坐标系 \(Oxyz\)
中，设 \(\vec{e}_1, \vec{e}_2, \vec{e}_3\) 分别为 \(x\) 轴、\(y\)
轴、\(z\) 轴正方向的单位向量（标准基向量），则任一空间向量 \(\vec{a}\)
可唯一表示为

\[\vec{a} = x\vec{e}_1 + y\vec{e}_2 + z\vec{e}_3\]

其中有序实数组 \((x, y, z)\) 称为向量 \(\vec{a}\) 的\textbf{坐标}，记为
\(\vec{a} = (x, y, z)\)。

\textbf{存在性}：过 \(\vec{a}\) 的终点 \(A\)
作三个平面分别平行于坐标平面 \(yOz\)、\(xOz\)、\(xOy\)，与坐标轴分别交于
\(P_1\)（\(x\) 轴）、\(P_2\)（\(y\) 轴）、\(P_3\)（\(z\) 轴）。则
\(\overrightarrow{OP_1} = x\vec{e}_1\)，\(\overrightarrow{P_1Q} = y\vec{e}_2\)（\(Q\)
为 \(A\) 在 \(xOy\)
平面上的投影），\(\overrightarrow{QA} = z\vec{e}_3\)。由向量加法，\(\vec{a} = x\vec{e}_1 + y\vec{e}_2 + z\vec{e}_3\)。唯一性的证明与二维情形类似。

\textbf{定义 15.12}（空间向量的模）：空间向量 \(\vec{a} = (x, y, z)\)
的\textbf{模}定义为

\[|\vec{a}| = \sqrt{x^2 + y^2 + z^2}\]

\textbf{证明}：由勾股定理，从原点到点 \((x, y, z)\) 的距离为
\(|\vec{a}|^2 = x^2 + y^2 + z^2\)。\(\blacksquare\)

空间向量的加法、减法、数乘的坐标表示与平面向量类似：

\[\vec{a} + \vec{b} = (x_1 + x_2, y_1 + y_2, z_1 + z_2)\]
\[\vec{a} - \vec{b} = (x_1 - x_2, y_1 - y_2, z_1 - z_2)\]
\[\lambda\vec{a} = (\lambda x_1, \lambda y_1, \lambda z_1)\]

\textbf{定理 15.8}（空间向量共线的充要条件）：非零向量
\(\vec{a} = (x_1, y_1, z_1)\) 与 \(\vec{b} = (x_2, y_2, z_2)\)
共线的充要条件是
\(\frac{x_2}{x_1} = \frac{y_2}{y_1} = \frac{z_2}{z_1}\)（当某个分量为零时，对应分量也须为零）。

\textbf{证明}：\(\vec{a} \parallel \vec{b}\) 当且仅当
\(\vec{a} = t\vec{b}\) 对某标量 \(t\)，即
\(x_1 = tx_2, y_1 = ty_2, z_1 = tz_2\)。消去 \(t\) 得
\(x_1y_2 = x_2y_1, y_1z_2 = y_2z_1, x_1z_2 = x_2z_1\)。\(\blacksquare\)

\paragraph{15.2.2
空间向量的数量积}\label{ux7a7aux95f4ux5411ux91cfux7684ux6570ux91cfux79ef}

\textbf{定义 15.13}（空间向量的数量积）：设 \(\vec{a}\) 与 \(\vec{b}\)
是空间中两个非零向量，\(\theta\)
为它们的夹角（\(0 \leq \theta \leq \pi\)），定义

\[\vec{a} \cdot \vec{b} = |\vec{a}||\vec{b}|\cos\theta\]

\textbf{定理 15.9}（空间向量数量积的坐标公式）：设
\(\vec{a} = (x_1, y_1, z_1)\)，\(\vec{b} = (x_2, y_2, z_2)\)，则

\[\vec{a} \cdot \vec{b} = x_1 x_2 + y_1 y_2 + z_1 z_2\]

\textbf{证明}：设
\(\vec{a} = x_1\vec{e}_1 + y_1\vec{e}_2 + z_1\vec{e}_3\)，\(\vec{b} = x_2\vec{e}_1 + y_2\vec{e}_2 + z_2\vec{e}_3\)。由分配律展开：

\[\vec{a} \cdot \vec{b} = x_1 x_2(\vec{e}_1 \cdot \vec{e}_1) + x_1 y_2(\vec{e}_1 \cdot \vec{e}_2) + \cdots\]

由于 \(\vec{e}_1, \vec{e}_2, \vec{e}_3\) 是标准正交基：

\[\vec{e}_i \cdot \vec{e}_j = \begin{cases} 1, & i = j \\ 0, & i \neq j \end{cases}\]

因此交叉项全部为零，得
\(\vec{a} \cdot \vec{b} = x_1 x_2 + y_1 y_2 + z_1 z_2\)。\(\blacksquare\)

\textbf{推论
15.3}（空间向量垂直的充要条件）：\(\vec{a} \perp \vec{b} \iff \vec{a} \cdot \vec{b} = 0 \iff x_1 x_2 + y_1 y_2 + z_1 z_2 = 0\)。

\textbf{证明}：与推论 15.2
类似，\(\vec{a} \perp \vec{b} \iff \theta = 90° \iff \cos\theta = 0 \iff \vec{a} \cdot \vec{b} = 0\)，再由定理
15.9 得
\(\vec{a} \cdot \vec{b} = x_1x_2 + y_1y_2 + z_1z_2 = 0\)。\(\blacksquare\)

\textbf{推论
15.4}（空间夹角公式）：\(\cos\theta = \dfrac{x_1 x_2 + y_1 y_2 + z_1 z_2}{\sqrt{x_1^2 + y_1^2 + z_1^2}\sqrt{x_2^2 + y_2^2 + z_2^2}}\)。

\textbf{证明}：与推论 15.1 类似，由空间数量积定义
\(\vec{a} \cdot \vec{b} = |\vec{a}||\vec{b}|\cos\theta\) 和定理 15.9 中
\(\vec{a} \cdot \vec{b} = x_1x_2 + y_1y_2 + z_1z_2\)，立即得到结论。\(\blacksquare\)

\paragraph{15.2.3
空间向量的向量积}\label{ux7a7aux95f4ux5411ux91cfux7684ux5411ux91cfux79ef}

\textbf{定义 15.14}（向量积）：设 \(\vec{a}\) 与 \(\vec{b}\)
是空间中两个不共线的向量，定义它们的\textbf{向量积}
\(\vec{a} \times \vec{b}\) 为满足以下条件的向量：

\begin{enumerate}
\def\labelenumi{\arabic{enumi}.}
\tightlist
\item
  \textbf{方向}：\(\vec{a} \times \vec{b}\) 垂直于 \(\vec{a}\) 和
  \(\vec{b}\) 所确定的平面，且
  \(\vec{a}, \vec{b}, \vec{a} \times \vec{b}\) 构成右手系。
\item
  \textbf{大小}：\(|\vec{a} \times \vec{b}| = |\vec{a}||\vec{b}|\sin\theta\)，其中
  \(\theta\) 为 \(\vec{a}\) 与 \(\vec{b}\) 的夹角。
\end{enumerate}

若 \(\vec{a}\) 与 \(\vec{b}\) 共线，则规定
\(\vec{a} \times \vec{b} = \vec{0}\)。

\textbf{定理 15.10}（向量积的坐标公式）：设
\(\vec{a} = (x_1, y_1, z_1)\)，\(\vec{b} = (x_2, y_2, z_2)\)，则

\[\vec{a} \times \vec{b} = (y_1 z_2 - z_1 y_2, \; z_1 x_2 - x_1 z_2, \; x_1 y_2 - y_1 x_2)\]

\textbf{证明}：由标准基的向量积关系：

\[\vec{e}_1 \times \vec{e}_2 = \vec{e}_3, \quad \vec{e}_2 \times \vec{e}_3 = \vec{e}_1, \quad \vec{e}_3 \times \vec{e}_1 = \vec{e}_2\]
\[\vec{e}_2 \times \vec{e}_1 = -\vec{e}_3, \quad \vec{e}_3 \times \vec{e}_2 = -\vec{e}_1, \quad \vec{e}_1 \times \vec{e}_3 = -\vec{e}_2\]
\[\vec{e}_i \times \vec{e}_i = \vec{0}\]

展开
\(\vec{a} \times \vec{b} = (x_1\vec{e}_1 + y_1\vec{e}_2 + z_1\vec{e}_3) \times (x_2\vec{e}_1 + y_2\vec{e}_2 + z_2\vec{e}_3)\)，利用上述关系逐项计算：

\[= (y_1 z_2 - z_1 y_2)\vec{e}_1 + (z_1 x_2 - x_1 z_2)\vec{e}_2 + (x_1 y_2 - y_1 x_2)\vec{e}_3\]

\(\blacksquare\)

\textbf{定理 15.11}（向量积的性质）：

\begin{enumerate}
\def\labelenumi{\arabic{enumi}.}
\tightlist
\item
  \textbf{反交换律}：\(\vec{a} \times \vec{b} = -(\vec{b} \times \vec{a})\)
\item
  \textbf{自叉积为零}：\(\vec{a} \times \vec{a} = \vec{0}\)
\item
  \textbf{分配律}：\(\vec{a} \times (\vec{b} + \vec{c}) = \vec{a} \times \vec{b} + \vec{a} \times \vec{c}\)
\end{enumerate}

\textbf{证明}：

\begin{enumerate}
\def\labelenumi{(\arabic{enumi})}
\item
  由坐标公式，交换 \(\vec{a}\) 与 \(\vec{b}\) 即每个分量变号，故
  \(\vec{a} \times \vec{b} = -(\vec{b} \times \vec{a})\)。
\item
  由反交换律：\(\vec{a} \times \vec{a} = -(\vec{a} \times \vec{a})\)，故
  \(2(\vec{a} \times \vec{a}) = \vec{0}\)，即
  \(\vec{a} \times \vec{a} = \vec{0}\)。
\item
  可由坐标公式验证。\(\blacksquare\)
\end{enumerate}

\textbf{定理 15.12}（向量积的几何意义）：\(|\vec{a} \times \vec{b}|\)
等于以 \(\vec{a}\) 和 \(\vec{b}\) 为邻边的平行四边形的面积。

\textbf{证明}：以 \(\vec{a}\) 和 \(\vec{b}\)
为邻边的平行四边形，底边长为 \(|\vec{a}|\)，高为
\(|\vec{b}|\sin\theta\)，故面积为
\(|\vec{a}||\vec{b}|\sin\theta = |\vec{a} \times \vec{b}|\)。\(\blacksquare\)

\textbf{推论 15.5}：以 \(\vec{a}\) 和 \(\vec{b}\) 为邻边的三角形面积为
\(S_{\triangle} = \frac{1}{2}|\vec{a} \times \vec{b}|\)。

\textbf{证明}：以 \(\vec{a}, \vec{b}\) 为邻边的平行四边形面积为
\(|\vec{a} \times \vec{b}|\)（定理
15.12），三角形为其一半。\(\blacksquare\)

\subsubsection{15.3
向量法解立体几何}\label{ux5411ux91cfux6cd5ux89e3ux7acbux4f53ux51e0ux4f55}

\paragraph{15.3.1
方向向量与法向量}\label{ux65b9ux5411ux5411ux91cfux4e0eux6cd5ux5411ux91cf}

\textbf{定义 15.15}（方向向量）：空间中直线 \(l\)
的\textbf{方向向量}是指与直线 \(l\) 平行的非零向量
\(\vec{d}\)。一条直线的方向向量有无穷多个，但它们两两共线。

\textbf{定义 15.16}（法向量）：平面 \(\alpha\)
的\textbf{法向量}是指垂直于平面 \(\alpha\) 的非零向量
\(\vec{n}\)。一个平面的法向量有无穷多个，但它们两两共线。

\textbf{定理 15.13}（法向量的确定）：设平面 \(\alpha\) 由三个不共线的点
\(A\)、\(B\)、\(C\) 确定，取
\(\vec{u} = \overrightarrow{AB}\)，\(\vec{v} = \overrightarrow{AC}\)，则
\(\alpha\) 的法向量为

\[\vec{n} = \vec{u} \times \vec{v}\]

\textbf{证明}：由向量积的定义，\(\vec{u} \times \vec{v}\) 同时垂直于
\(\vec{u}\) 和 \(\vec{v}\)。由于 \(\vec{u}\) 和 \(\vec{v}\)
不共线且都在平面 \(\alpha\) 内，它们张成整个平面 \(\alpha\)，故
\(\vec{n}\) 垂直于 \(\alpha\) 内的任意向量，即 \(\vec{n}\) 是 \(\alpha\)
的法向量。\(\blacksquare\)

若平面 \(\alpha\) 由方程 \(ax + by + cz = d\) 给出，则法向量为
\(\vec{n} = (a, b, c)\)。

\paragraph{15.3.2
平行与垂直的向量判定}\label{ux5e73ux884cux4e0eux5782ux76f4ux7684ux5411ux91cfux5224ux5b9a}

利用方向向量和法向量，可以将空间中的平行与垂直关系全部转化为向量运算。

\textbf{定理 15.14}（线线平行）：设 \(l_1\) 的方向向量为
\(\vec{d}_1\)，\(l_2\) 的方向向量为 \(\vec{d}_2\)，则
\(l_1 \parallel l_2\) 的充要条件是
\(\vec{d}_1 \parallel \vec{d}_2\)，即存在 \(\lambda \neq 0\) 使得
\(\vec{d}_2 = \lambda\vec{d}_1\)。

\textbf{证明}：\(l_1 \parallel l_2\) 当且仅当 \(l_1\) 与 \(l_2\)
方向相同（或相反），当且仅当 \(\vec{d}_1\) 与 \(\vec{d}_2\)
共线。\(\blacksquare\)

\textbf{定理 15.15}（线面平行）：设直线 \(l\) 的方向向量为
\(\vec{d}\)，平面 \(\alpha\) 的法向量为 \(\vec{n}_\alpha\)，则
\(l \parallel \alpha\)（\(l \not\subset \alpha\)）的充要条件是
\(\vec{d} \cdot \vec{n}_\alpha = 0\)。

\textbf{证明}：\(l \parallel \alpha\) 当且仅当 \(l\) 的方向平行于
\(\alpha\)，当且仅当 \(\vec{d}\) 垂直于 \(\alpha\) 的法向量
\(\vec{n}_\alpha\)，当且仅当
\(\vec{d} \cdot \vec{n}_\alpha = 0\)。\(\blacksquare\)

\textbf{定理 15.16}（面面平行）：设 \(\alpha\) 的法向量为
\(\vec{n}_\alpha\)，\(\beta\) 的法向量为 \(\vec{n}_\beta\)，则
\(\alpha \parallel \beta\) 的充要条件是
\(\vec{n}_\alpha \parallel \vec{n}_\beta\)，即存在 \(\lambda \neq 0\)
使得 \(\vec{n}_\beta = \lambda\vec{n}_\alpha\)。

\textbf{证明}：\(\alpha \parallel \beta\)
当且仅当两个平面有相同的''倾斜方向''，当且仅当法向量平行。\(\blacksquare\)

\textbf{定理 15.17}（线线垂直）：\(l_1 \perp l_2\) 的充要条件是
\(\vec{d}_1 \cdot \vec{d}_2 = 0\)。

\textbf{证明}：\(l_1 \perp l_2 \iff \theta = 90° \iff \cos\theta = 0 \iff \vec{d}_1 \cdot \vec{d}_2 = 0\)。\(\blacksquare\)

\textbf{定理 15.18}（线面垂直）：\(l \perp \alpha\) 的充要条件是
\(\vec{d} \parallel \vec{n}_\alpha\)，即存在 \(\lambda \neq 0\) 使得
\(\vec{d} = \lambda\vec{n}_\alpha\)。

\textbf{证明}：\(l \perp \alpha\) 当且仅当 \(l\) 垂直于 \(\alpha\)
内的所有直线，当且仅当 \(\vec{d}\) 与 \(\alpha\)
的法向量平行。\(\blacksquare\)

\textbf{定理 15.19}（面面垂直）：\(\alpha \perp \beta\) 的充要条件是
\(\vec{n}_\alpha \cdot \vec{n}_\beta = 0\)。

\textbf{证明}：\(\alpha \perp \beta\) 当且仅当二面角的平面角为
\(90°\)。两个法向量的夹角 \(\gamma\) 满足
\(\cos\gamma = \frac{\vec{n}_\alpha \cdot \vec{n}_\beta}{|\vec{n}_\alpha||\vec{n}_\beta|}\)。二面角的平面角
\(\psi\) 与 \(\gamma\) 的关系是 \(\psi = \gamma\) 或
\(\psi = \pi - \gamma\)。当 \(\psi = 90°\) 时，\(|\cos\gamma| = 0\)，即
\(\vec{n}_\alpha \cdot \vec{n}_\beta = 0\)。\(\blacksquare\)

\textbf{总结表}：

{\def\LTcaptype{none} % do not increment counter
\begin{longtable}[]{@{}ll@{}}
\toprule\noalign{}
位置关系 & 向量条件 \\
\midrule\noalign{}
\endhead
\bottomrule\noalign{}
\endlastfoot
\(l_1 \parallel l_2\) & \(\vec{d}_1 \parallel \vec{d}_2\) \\
\(l \parallel \alpha\) & \(\vec{d} \cdot \vec{n}_\alpha = 0\) \\
\(\alpha \parallel \beta\) &
\(\vec{n}_\alpha \parallel \vec{n}_\beta\) \\
\(l_1 \perp l_2\) & \(\vec{d}_1 \cdot \vec{d}_2 = 0\) \\
\(l \perp \alpha\) & \(\vec{d} \parallel \vec{n}_\alpha\) \\
\(\alpha \perp \beta\) & \(\vec{n}_\alpha \cdot \vec{n}_\beta = 0\) \\
\end{longtable}
}

\paragraph{15.3.3
空间角的向量公式}\label{ux7a7aux95f4ux89d2ux7684ux5411ux91cfux516cux5f0f}

\textbf{定理 15.20}（异面直线所成的角）：设 \(l_1\) 的方向向量为
\(\vec{d}_1\)，\(l_2\) 的方向向量为 \(\vec{d}_2\)，则异面直线 \(l_1\) 与
\(l_2\) 所成的角 \(\phi\) 满足

\[\cos\phi = \frac{|\vec{d}_1 \cdot \vec{d}_2|}{|\vec{d}_1||\vec{d}_2|}\]

\textbf{证明}：设 \(\vec{d}_1\) 与 \(\vec{d}_2\) 的夹角为
\(\theta\)（\(0 \leq \theta \leq \pi\)）。异面直线所成的角取锐角或直角，即
\(\phi = \theta\)（若 \(\theta \leq 90°\)）或
\(\phi = 180° - \theta\)（若
\(\theta > 90°\)）。无论哪种情况，\(\cos\phi = |\cos\theta| = \frac{|\vec{d}_1 \cdot \vec{d}_2|}{|\vec{d}_1||\vec{d}_2|}\)。\(\blacksquare\)

\textbf{定理 15.21}（直线与平面所成的角）：设直线 \(l\) 的方向向量为
\(\vec{d}\)，平面 \(\alpha\) 的法向量为 \(\vec{n}\)，则直线 \(l\) 与平面
\(\alpha\) 所成的角 \(\varphi\) 满足

\[\sin\varphi = \frac{|\vec{d} \cdot \vec{n}|}{|\vec{d}||\vec{n}|}\]

\textbf{证明}：设 \(\vec{d}\) 与 \(\vec{n}\) 的夹角为 \(\gamma\)。由于
\(\vec{n}\) 垂直于 \(\alpha\)，\(l\) 在 \(\alpha\) 上的射影 \(l'\)
的方向向量 \(\vec{d}'\) 满足 \(\vec{d}' \perp \vec{n}\)。\(l\) 与 \(l'\)
所成的角 \(\varphi\) 满足 \(\varphi = |90° - \gamma|\)。因此

\[\sin\varphi = \sin|90° - \gamma| = |\cos\gamma| = \frac{|\vec{d} \cdot \vec{n}|}{|\vec{d}||\vec{n}|}\]

\(\blacksquare\)

\textbf{定理 15.22}（二面角）：设平面 \(\alpha\) 的法向量为
\(\vec{n}_1\)，平面 \(\beta\) 的法向量为 \(\vec{n}_2\)，则二面角
\(\alpha\)-\(l\)-\(\beta\) 的平面角 \(\psi\) 满足

\[\cos\psi = \pm\frac{\vec{n}_1 \cdot \vec{n}_2}{|\vec{n}_1||\vec{n}_2|}\]

其中正负号取决于法向量的选取方向。

\textbf{证明}：两个法向量 \(\vec{n}_1\) 与 \(\vec{n}_2\) 的夹角
\(\gamma\) 满足
\(\cos\gamma = \frac{\vec{n}_1 \cdot \vec{n}_2}{|\vec{n}_1||\vec{n}_2|}\)。二面角的平面角
\(\psi\) 与 \(\gamma\) 的关系是 \(\psi = \gamma\) 或
\(\psi = \pi - \gamma\)，取决于法向量的选取方向是否指向二面角的同一侧。因此
\(\cos\psi = \pm\cos\gamma\)。在实际计算中，通常直接计算
\(\cos\gamma\)，再根据几何意义判断正负号。\(\blacksquare\)

\paragraph{15.3.4
空间距离的向量公式}\label{ux7a7aux95f4ux8dddux79bbux7684ux5411ux91cfux516cux5f0f}

\textbf{定理 15.23}（点到平面的距离）：设平面 \(\alpha\) 的法向量为
\(\vec{n}\)，\(P\) 为空间中一点，\(A\) 为平面 \(\alpha\)
上的任意一点，则点 \(P\) 到平面 \(\alpha\) 的距离为

\[d = \frac{|\overrightarrow{AP} \cdot \vec{n}|}{|\vec{n}|}\]

\textbf{证明}：设 \(H\) 为 \(P\) 在平面 \(\alpha\)
上的正投影（垂足），则 \(\overrightarrow{PH} \parallel \vec{n}\)。设
\(\overrightarrow{PH} = t\vec{n}\)，则

\[\overrightarrow{AH} = \overrightarrow{AP} + \overrightarrow{PH} = \overrightarrow{AP} + t\vec{n}\]

由于 \(H \in \alpha\)，\(\overrightarrow{AH}\) 在 \(\alpha\) 内，故
\(\overrightarrow{AH} \perp \vec{n}\)，即

\[(\overrightarrow{AP} + t\vec{n}) \cdot \vec{n} = 0 \implies \overrightarrow{AP} \cdot \vec{n} + t|\vec{n}|^2 = 0 \implies t = -\frac{\overrightarrow{AP} \cdot \vec{n}}{|\vec{n}|^2}\]

因此

\[d = |\overrightarrow{PH}| = |t| \cdot |\vec{n}| = \frac{|\overrightarrow{AP} \cdot \vec{n}|}{|\vec{n}|}\]

\(\blacksquare\)

\textbf{定理 15.24}（异面直线的距离）：设 \(l_1\) 过点 \(A\)，方向向量为
\(\vec{d}_1\)；\(l_2\) 过点 \(B\)，方向向量为 \(\vec{d}_2\)。则 \(l_1\)
与 \(l_2\) 的距离为

\[d = \frac{|\overrightarrow{AB} \cdot (\vec{d}_1 \times \vec{d}_2)|}{|\vec{d}_1 \times \vec{d}_2|}\]

\textbf{证明}：公垂线的方向向量为
\(\vec{d}_1 \times \vec{d}_2\)（同时垂直于 \(\vec{d}_1\) 和
\(\vec{d}_2\)）。异面直线的距离等于 \(\overrightarrow{AB}\)
在公垂线方向上的投影的绝对值：

\[d = \left|\text{proj}_{\vec{d}_1 \times \vec{d}_2}\overrightarrow{AB}\right| = \frac{|\overrightarrow{AB} \cdot (\vec{d}_1 \times \vec{d}_2)|}{|\vec{d}_1 \times \vec{d}_2|}\]

\(\blacksquare\)

\subsubsection{15.4
向量法证明几何定理}\label{ux5411ux91cfux6cd5ux8bc1ux660eux51e0ux4f55ux5b9aux7406}

向量法为证明经典几何定理提供了简洁有力的工具。以下展示如何用向量方法重新证明若干重要定理。

\textbf{定理 15.25}（用向量法证明余弦定理）：在 \(\triangle ABC\)
中，\(a^2 = b^2 + c^2 - 2bc\cos A\)。

\textbf{证明}：设
\(\overrightarrow{AB} = \vec{c}\)，\(\overrightarrow{AC} = \vec{b}\)，则
\(\overrightarrow{BC} = \vec{b} - \vec{c}\)。

\[a^2 = |\overrightarrow{BC}|^2 = (\vec{b} - \vec{c}) \cdot (\vec{b} - \vec{c}) = |\vec{b}|^2 - 2\vec{b} \cdot \vec{c} + |\vec{c}|^2\]

\[= b^2 + c^2 - 2|\vec{b}||\vec{c}|\cos A = b^2 + c^2 - 2bc\cos A\]

\(\blacksquare\)

\textbf{定理
15.26}（用向量法证明平行四边形对角线性质）：平行四边形的对角线互相平分。

\textbf{证明}：设平行四边形
\(ABCD\)，\(\overrightarrow{AB} = \vec{a}\)，\(\overrightarrow{AD} = \vec{b}\)。则
\(\overrightarrow{AC} = \vec{a} + \vec{b}\)，\(\overrightarrow{DB} = \vec{a} - \vec{b}\)。

设对角线 \(AC\) 和 \(BD\) 的交点为 \(O\)。设
\(\overrightarrow{AO} = t(\vec{a} + \vec{b})\)（\(O\) 在 \(AC\)
上），\(\overrightarrow{DO} = s(\vec{a} - \vec{b})\)（\(O\) 在 \(BD\)
上）。

由
\(\overrightarrow{AO} = \overrightarrow{AD} + \overrightarrow{DO}\)，得
\(t(\vec{a} + \vec{b}) = \vec{b} + s(\vec{a} - \vec{b})\)。

比较 \(\vec{a}\) 和 \(\vec{b}\) 的系数：\(t = s\)，\(t = 1 - s\)。解得
\(t = s = \frac{1}{2}\)。

因此
\(\overrightarrow{AO} = \frac{1}{2}\overrightarrow{AC}\)，\(\overrightarrow{DO} = \frac{1}{2}\overrightarrow{DB}\)，即
\(O\) 是两条对角线的中点。\(\blacksquare\)

\textbf{定理
15.27}（用向量法证明三角形中位线定理）：三角形的中位线平行于第三边且等于第三边的一半。

\textbf{证明}：设 \(\triangle ABC\) 中，\(D\)、\(E\) 分别是
\(\overline{AB}\)、\(\overline{AC}\) 的中点。设
\(\overrightarrow{AB} = \vec{a}\)，\(\overrightarrow{AC} = \vec{b}\)。

\[\overrightarrow{AD} = \frac{1}{2}\vec{a}, \quad \overrightarrow{AE} = \frac{1}{2}\vec{b}\]

\[\overrightarrow{DE} = \overrightarrow{AE} - \overrightarrow{AD} = \frac{1}{2}\vec{b} - \frac{1}{2}\vec{a} = \frac{1}{2}(\vec{b} - \vec{a}) = \frac{1}{2}\overrightarrow{BC}\]

因此 \(\overrightarrow{DE} \parallel \overrightarrow{BC}\) 且
\(|\overrightarrow{DE}| = \frac{1}{2}|\overrightarrow{BC}|\)。\(\blacksquare\)

\textbf{定理
15.28}（用向量法证明勾股定理）：在直角三角形中，斜边的平方等于两直角边的平方和。

\textbf{证明}：设 \(\triangle ABC\) 中 \(\angle C = 90°\)。设
\(\overrightarrow{CA} = \vec{a}\)，\(\overrightarrow{CB} = \vec{b}\)，则
\(\overrightarrow{AB} = \vec{b} - \vec{a}\)。

由 \(\angle C = 90°\)，\(\vec{a} \perp \vec{b}\)，即
\(\vec{a} \cdot \vec{b} = 0\)。

\[c^2 = |\overrightarrow{AB}|^2 = (\vec{b} - \vec{a}) \cdot (\vec{b} - \vec{a}) = |\vec{a}|^2 - 2\vec{a} \cdot \vec{b} + |\vec{b}|^2 = a^2 + b^2\]

\(\blacksquare\)

\textbf{定理
15.29}（用向量法证明三垂线定理）：如果平面内的一条直线与一条斜线在平面内的射影垂直，则它也与这条斜线垂直。

\textbf{证明}：设 \(PO \perp \alpha\) 于
\(O\)，\(A \in \alpha\)，\(a \subset \alpha\)，\(a \perp OA\)。在 \(a\)
上取方向向量 \(\vec{d}\)，则 \(\vec{d} \perp \overrightarrow{OA}\)，即
\(\vec{d} \cdot \overrightarrow{OA} = 0\)。

又 \(\overrightarrow{PO} \perp \alpha\)，故
\(\overrightarrow{PO} \perp \vec{d}\)，即
\(\overrightarrow{PO} \cdot \vec{d} = 0\)。

\[\overrightarrow{PA} \cdot \vec{d} = (\overrightarrow{PO} + \overrightarrow{OA}) \cdot \vec{d} = \overrightarrow{PO} \cdot \vec{d} + \overrightarrow{OA} \cdot \vec{d} = 0 + 0 = 0\]

因此 \(\overrightarrow{PA} \perp \vec{d}\)，即
\(PA \perp a\)。\(\blacksquare\)

\textbf{定理 15.30}（用向量法证明正弦定理）：在 \(\triangle ABC\)
中，\(\frac{a}{\sin A} = \frac{b}{\sin B} = \frac{c}{\sin C}\)。

\textbf{证明}：设
\(\overrightarrow{AB} = \vec{c}\)，\(\overrightarrow{AC} = \vec{b}\)。由向量积的几何意义：

\[|\vec{b} \times \vec{c}| = |\vec{b}||\vec{c}|\sin A = bc\sin A\]

而 \(|\vec{b} \times \vec{c}|\) 等于以 \(\overrightarrow{AB}\) 和
\(\overrightarrow{AC}\) 为邻边的平行四边形的面积，即
\(2S_{\triangle}\)（三角形面积的两倍）。

因此 \(bc\sin A = 2S\)，即
\(\frac{a}{\sin A} = \frac{a \cdot bc}{2S}\)。

同理 \(ac\sin B = 2S\)，\(ab\sin C = 2S\)。

所以
\(\frac{a}{\sin A} = \frac{b}{\sin B} = \frac{c}{\sin C} = \frac{abc}{2S}\)。

\textbf{注}：上述向量法证明建立了
\(\frac{a}{\sin A} = \frac{b}{\sin B} = \frac{c}{\sin C} = \frac{abc}{2S}\)。其中
\(\frac{abc}{2S} = 2R\)（\(R\) 为外接圆半径）的结论参见定理 13.40
的几何证明。\(\blacksquare\)

\bigskip\hrule\bigskip

\paragraph{15.4.1 混合积}\label{ux6df7ux5408ux79ef}

\textbf{定义 15.17}（混合积）：三个向量 \(\vec{a}, \vec{b}, \vec{c}\)
的\textbf{混合积}（标量三重积）定义为

\[[\vec{a}\;\vec{b}\;\vec{c}] = (\vec{a} \times \vec{b}) \cdot \vec{c}\]

\textbf{定理 15.31}（混合积的坐标公式）：设
\(\vec{a} = (x_1, y_1, z_1)\)，\(\vec{b} = (x_2, y_2, z_2)\)，\(\vec{c} = (x_3, y_3, z_3)\)，则

\[[\vec{a}\;\vec{b}\;\vec{c}] = \begin{vmatrix} x_1 & y_1 & z_1 \\ x_2 & y_2 & z_2 \\ x_3 & y_3 & z_3 \end{vmatrix}\]

\textbf{证明}：由向量积的坐标公式，\(\vec{a} \times \vec{b} = (y_1 z_2 - z_1 y_2, \; z_1 x_2 - x_1 z_2, \; x_1 y_2 - y_1 x_2)\)。再由数量积的坐标公式：

\[[\vec{a}\;\vec{b}\;\vec{c}] = x_3(y_1 z_2 - z_1 y_2) + y_3(z_1 x_2 - x_1 z_2) + z_3(x_1 y_2 - y_1 x_2)\]

按第三行展开行列式即得上述结果。\(\blacksquare\)

\textbf{定理
15.32}（混合积的几何意义）：\(|[\vec{a}\;\vec{b}\;\vec{c}]|\) 等于以
\(\vec{a}, \vec{b}, \vec{c}\) 为棱的平行六面体的体积。

\textbf{证明}：\(\vec{a} \times \vec{b}\) 的模长等于以
\(\vec{a}, \vec{b}\) 为底的平行四边形的面积
\(S\)，方向为底面的法向量。\((\vec{a} \times \vec{b}) \cdot \vec{c}\)
等于 \(\vec{c}\) 在 \(\vec{a} \times \vec{b}\) 方向上的投影与
\(|\vec{a} \times \vec{b}|\) 的乘积，即底面积 \(S\) 乘以高
\(h\)，等于平行六面体的体积。\(\blacksquare\)

\textbf{推论 15.6}（共面的充要条件）：三个向量
\(\vec{a}, \vec{b}, \vec{c}\) 共面的充要条件是
\([\vec{a}\;\vec{b}\;\vec{c}] = 0\)。

\textbf{证明}：\(\vec{a}, \vec{b}, \vec{c}\) 共面 \(\iff\)
以它们为棱的平行六面体体积为零 \(\iff\)
\([\vec{a}\;\vec{b}\;\vec{c}] = 0\)。\(\blacksquare\)

\textbf{定理
15.33}（混合积的轮换对称性）：\([\vec{a}\;\vec{b}\;\vec{c}] = [\vec{b}\;\vec{c}\;\vec{a}] = [\vec{c}\;\vec{a}\;\vec{b}]\)。

\textbf{证明}：由行列式的性质，轮换行列式的行不改变行列式的值（或改变偶数次行交换，符号不变）：

\[[\vec{a}\;\vec{b}\;\vec{c}] = (\vec{a} \times \vec{b}) \cdot \vec{c} = \vec{c} \cdot (\vec{a} \times \vec{b}) = [\vec{c}\;\vec{a}\;\vec{b}]\]

其中最后一步利用了混合积的循环性质（行列式的行轮换）。\(\blacksquare\)

\paragraph{15.4.2
用向量法证明更多的几何定理}\label{ux7528ux5411ux91cfux6cd5ux8bc1ux660eux66f4ux591aux7684ux51e0ux4f55ux5b9aux7406}

\textbf{定理 15.34}（用向量法证明斯图尔特定理）：设 \(D\) 是
\(\triangle ABC\) 的边 \(BC\)
上一点，\(BD = m\)，\(DC = n\)，\(BC = a = m + n\)，则

\[AD^2 = \frac{mb^2 + nc^2}{m + n} - mn\]

\textbf{证明}：设
\(\overrightarrow{AB} = \vec{c}\)，\(\overrightarrow{AC} = \vec{b}\)。则
\(\overrightarrow{BC} = \vec{b} - \vec{c}\)，\(|\vec{b}| = b\)，\(|\vec{c}| = c\)。

\[\overrightarrow{AD} = \overrightarrow{AB} + \overrightarrow{BD} = \vec{c} + \frac{m}{a}(\vec{b} - \vec{c}) = \frac{n}{a}\vec{c} + \frac{m}{a}\vec{b}\]

\[AD^2 = |\overrightarrow{AD}|^2 = \frac{n^2}{a^2}|\vec{c}|^2 + 2 \cdot \frac{mn}{a^2}\vec{b} \cdot \vec{c} + \frac{m^2}{a^2}|\vec{b}|^2\]

\[= \frac{n^2 c^2 + m^2 b^2 + 2mn \vec{b} \cdot \vec{c}}{a^2}\]

由余弦定理，\(\vec{b} \cdot \vec{c} = bc\cos A = \frac{b^2 + c^2 - a^2}{2}\)。代入化简：

\[AD^2 = \frac{n^2 c^2 + m^2 b^2 + mn(b^2 + c^2 - a^2)}{a^2} = \frac{(n^2 + mn)c^2 + (m^2 + mn)b^2 - mna^2}{a^2}\]

\[= \frac{n(m+n)c^2 + m(m+n)b^2}{a^2} - \frac{mna^2}{a^2} = \frac{nc^2 + mb^2}{a} - mn = \frac{mb^2 + nc^2}{m+n} - mn\]

\(\blacksquare\)

\textbf{推论 15.7}（中线长公式）：当 \(m = n = \frac{a}{2}\) 时（\(D\)
为 \(BC\) 的中点），斯图尔特定理退化为中线长公式：

\[m_a^2 = \frac{2b^2 + 2c^2 - a^2}{4}\]

\textbf{证明}：代入 \(m = n = \frac{a}{2}\)：

\[AD^2 = \frac{\frac{a}{2} b^2 + \frac{a}{2} c^2}{a} - \frac{a^2}{4} = \frac{b^2 + c^2}{2} - \frac{a^2}{4} = \frac{2b^2 + 2c^2 - a^2}{4}\]

\(\blacksquare\)

\textbf{定理
15.35}（用向量法证明阿波罗尼奥斯定理）：三角形三条中线的平方和等于三边平方和的四分之三，即

\[m_a^2 + m_b^2 + m_c^2 = \frac{3}{4}(a^2 + b^2 + c^2)\]

\textbf{证明}：由中线长公式（推论 15.7），三条中线的平方分别为：

\[m_a^2 = \frac{2b^2 + 2c^2 - a^2}{4}, \quad m_b^2 = \frac{2a^2 + 2c^2 - b^2}{4}, \quad m_c^2 = \frac{2a^2 + 2b^2 - c^2}{4}\]

求和：

\[m_a^2 + m_b^2 + m_c^2 = \frac{(2b^2 + 2c^2 - a^2) + (2a^2 + 2c^2 - b^2) + (2a^2 + 2b^2 - c^2)}{4}\]

\[= \frac{3a^2 + 3b^2 + 3c^2}{4} = \frac{3}{4}(a^2 + b^2 + c^2)\]

\(\blacksquare\)

\textbf{定理 15.36}（用向量法证明四面体体积公式）：以
\(\vec{a}, \vec{b}, \vec{c}\) 为从同一顶点出发的三条棱的四面体的体积为

\[V = \frac{1}{6}|[\vec{a}\;\vec{b}\;\vec{c}]|\]

\textbf{证明}：以 \(\vec{a}, \vec{b}, \vec{c}\) 为棱的平行六面体的体积为
\(|[\vec{a}\;\vec{b}\;\vec{c}]|\)（定理
15.32）。四面体是平行六面体的六分之一（平行六面体可分解为六个等体积的四面体），因此
\(V = \frac{1}{6}|[\vec{a}\;\vec{b}\;\vec{c}]|\)。\(\blacksquare\)

\bigskip\hrule\bigskip

\textbf{本章展望}：本章建立了向量的完整理论体系，从平面向量的运算和坐标表示出发，推广到空间向量，引入了数量积和向量积这两个核心运算，并系统地将立体几何中的平行、垂直、角度和距离关系转化为向量运算。向量方法将几何直觉与代数计算完美结合，使得复杂的几何问题可以程序化地求解。在后续的解析几何中，向量将成为描述曲线和曲面的基本工具；在函数分析中，向量空间的概念将进一步推广到无穷维的情形。

\bigskip\hrule\bigskip

\begin{quote}
全卷总结：本卷从希尔伯特公理系统出发，经过严格的逻辑推导，建立了欧氏几何的完整理论体系。第12章确立了公理基础，第13章推导了平面几何的全部核心定理，第14章将这些结果推广到三维空间，第15章引入了向量这一强有力的代数工具。这四章构成了从公理到现代几何的完整逻辑链条：公理
\(\to\) 平面几何 \(\to\) 立体几何 \(\to\)
向量方法。每一步推导都追溯到公理，每一个定理都有严格证明，体现了数学从直觉到严格、从特殊到一般的思维范式。

\textbf{各章定理索引}：

\begin{itemize}
\tightlist
\item
  \textbf{Ch 12}（公理体系）：关联公理 I1--I8、顺序公理 O1--O4、合同公理
  C1--C6、平行公理 P、连续公理
  D1--D2；线段中点和角平分线的存在唯一性、三角形内角和定理、外角定理。
\item
  \textbf{Ch 13}（平面几何）：ASA/AAS/SSS
  全等判定、等腰三角形性质、大边对大角/大角对大边、三角不等式、平行四边形性质与判定、矩形/菱形/正方形性质、中位线定理、平行线分线段成比例、AA/SAS/SSS
  相似判定、勾股定理、圆周角定理、切线定理、圆幂定理、海伦公式、正弦定理、余弦定理。
\item
  \textbf{Ch
  14}（立体几何）：线面平行/面面平行的判定与性质、线面垂直/面面垂直的判定与性质、三垂线定理及其逆定理、异面直线角/线面角/二面角、各种距离公式、棱柱/棱锥/棱台/圆柱/圆锥/圆台/球的体积与表面积、祖暅原理、欧拉公式。
\item
  \textbf{Ch
  15}（向量）：平面向量运算与坐标、数量积与向量积、向量法判定平行垂直、向量法计算角度与距离、混合积、向量法证明勾股定理/余弦定理/正弦定理/中位线定理/斯图尔特定理/阿波罗尼奥斯定理。
\end{itemize}
\end{quote}

\bigskip\hrule\bigskip

\subsection{Ch 25 直线与圆}\label{ch-25-ux76f4ux7ebfux4e0eux5706}

\begin{quote}
\textbf{前置知识}

本章建立在以下前序章节的基础之上：

\begin{itemize}
\tightlist
\item
  \textbf{Ch
  16（坐标系与函数概念）}：平面直角坐标系的建立、两点间距离公式、中点公式是本章的几何基础。所有解析几何的讨论均以坐标系为起点。
\item
  \textbf{Ch 8（方程与方程组）}：直线方程本质上是一次方程
  \(Ax + By + C = 0\)；圆的方程展开后涉及二元二次方程。方程的代数理论（如韦达定理、判别式）在分析直线与圆的交点问题时不可或缺。
\item
  \textbf{Ch
  13.4（圆的几何性质）}：垂径定理、圆周角定理、切线性质、圆幂定理等经典几何结论，为解析化处理圆的问题提供了直觉与验证工具。
\item
  \textbf{Ch
  23（导数与微分）}：导数给出了曲线在某点处切线的斜率，这一概念将在本章末尾用于推导圆的切线方程，并为后续章节中圆锥曲线的切线研究奠定基础。
\item
  \textbf{Ch
  15（向量）}：向量的数量积、向量的平行与垂直判定，是处理直线位置关系的有力工具。
\end{itemize}
\end{quote}

\bigskip\hrule\bigskip

\subsection{25.1 直线的方程}\label{ux76f4ux7ebfux7684ux65b9ux7a0b}

\subsubsection{25.1.1
直线方程的建立}\label{ux76f4ux7ebfux65b9ux7a0bux7684ux5efaux7acb}

在 Ch 16
中，我们建立了平面直角坐标系，并证明了两点间距离公式和中点公式。现在我们要回答一个基本问题：\textbf{如何用方程描述一条直线？}

\textbf{定义 25.1（直线的斜率）.} 设直线 \(l\) 不垂直于 \(x\)
轴，\(P_1(x_1, y_1)\) 和 \(P_2(x_2, y_2)\)（\(x_1 \neq x_2\)）是 \(l\)
上任意两点，则比值

\[k = \frac{y_2 - y_1}{x_2 - x_1}\]

称为直线 \(l\) 的\textbf{斜率}。若直线垂直于 \(x\) 轴，则称斜率不存在。

\textbf{命题 25.1（斜率的唯一性）.} 直线 \(l\) 的斜率 \(k\) 与 \(l\)
上所取的两点无关。

\textbf{证明}： 设 \(P_1(x_1, y_1), P_2(x_2, y_2)\) 和
\(P_3(x_3, y_3), P_4(x_4, y_4)\) 分别是 \(l\) 上两对不同的点。由于
\(P_1, P_2, P_3\) 共线，过 \(P_1\) 作 \(x\) 轴的平行线，过 \(P_2\) 作
\(y\) 轴的平行线，交于点 \(Q_1\)；过 \(P_1\) 作 \(x\) 轴的平行线，过
\(P_3\) 作 \(y\) 轴的平行线，交于点 \(Q_2\)。由相似三角形（Ch
13.3），对应边成比例，即

\[\frac{y_3 - y_1}{x_3 - x_1} = \frac{y_2 - y_1}{x_2 - x_1}\]

同理可证
\(\frac{y_4 - y_3}{x_4 - x_3} = \frac{y_2 - y_1}{x_2 - x_1}\)。故斜率与所取点无关。\(\blacksquare\)

\textbf{定理 25.1（点斜式方程）.} 过点 \(P_0(x_0, y_0)\)、斜率为 \(k\)
的直线 \(l\) 的方程为

\[y - y_0 = k(x - x_0)\]

\textbf{证明}： 设 \(P(x, y)\) 为 \(l\) 上任意一点。若 \(P = P_0\)，则
\(x = x_0, y = y_0\)，等式显然成立。若 \(P \neq P_0\)，由斜率的定义：

\[k = \frac{y - y_0}{x - x_0}\]

即 \(y - y_0 = k(x - x_0)\)。

反过来，若点 \(P(x, y)\) 满足 \(y - y_0 = k(x - x_0)\)，则当
\(x \neq x_0\) 时，\(\frac{y - y_0}{x - x_0} = k\)，说明 \(P\) 在过
\(P_0\) 且斜率为 \(k\) 的直线上；当 \(x = x_0\) 时，\(y = y_0\)，即
\(P = P_0\)，也在直线上。

因此方程 \(y - y_0 = k(x - x_0)\) 恰好描述了过 \(P_0\) 且斜率为 \(k\)
的直线。\(\blacksquare\)

\textbf{定理 25.2（两点式方程）.} 过点 \(P_1(x_1, y_1)\) 和
\(P_2(x_2, y_2)\)（\(x_1 \neq x_2\) 且 \(y_1 \neq y_2\)）的直线方程为

\[\frac{y - y_1}{y_2 - y_1} = \frac{x - x_1}{x_2 - x_1}\]

\textbf{证明}： 由定理 25.1，过 \(P_1\) 且斜率为
\(k = \frac{y_2 - y_1}{x_2 - x_1}\) 的直线方程为

\[y - y_1 = \frac{y_2 - y_1}{x_2 - x_1}(x - x_1)\]

当 \(y_2 \neq y_1\) 时，两边除以 \(y_2 - y_1\)
即得两点式。\(\blacksquare\)

\textbf{定理 25.3（截距式方程）.} 在 \(x\) 轴和 \(y\) 轴上的截距分别为
\(a\) 和 \(b\)（\(a \neq 0, b \neq 0\)）的直线方程为

\[\frac{x}{a} + \frac{y}{b} = 1\]

\textbf{证明}： 该直线过点 \((a, 0)\) 和 \((0, b)\)。由两点式：

\[\frac{y - 0}{b - 0} = \frac{x - a}{0 - a}\]

\[\frac{y}{b} = \frac{a - x}{a} = 1 - \frac{x}{a}\]

\[\frac{x}{a} + \frac{y}{b} = 1\]

反过来，\((a, 0)\) 和 \((0, b)\)
均满足此方程，且两点唯一确定一条直线。\(\blacksquare\)

\textbf{定理 25.4（一般式方程）.} 平面上的任何直线都可以表示为

\[Ax + By + C = 0\]

其中 \(A, B\) 不全为零。反之，满足此方程（\(A, B\)
不全为零）的点的集合是一条直线。

\textbf{证明}： 若直线不垂直于 \(x\) 轴，其方程可写为
\(y - y_0 = k(x - x_0)\)，即 \(kx - y + (y_0 - kx_0) = 0\)，取
\(A = k, B = -1, C = y_0 - kx_0\) 即可。

若直线垂直于 \(x\) 轴，方程为 \(x = x_0\)，即
\(1 \cdot x + 0 \cdot y + (-x_0) = 0\)，取 \(A = 1, B = 0, C = -x_0\)。

反之，若 \(B \neq 0\)，则 \(y = -\frac{A}{B}x - \frac{C}{B}\)，斜率为
\(-\frac{A}{B}\)，是一条直线。若 \(B = 0\)，则 \(A \neq 0\)，方程变为
\(x = -\frac{C}{A}\)，是垂直于 \(x\) 轴的直线。\(\blacksquare\)

\textbf{命题 25.2（一般式的法向量）.} 直线 \(Ax + By + C = 0\)
的一个法向量为 \(\vec{n} = (A, B)\)。

\textbf{证明}： 设 \(P_1(x_1, y_1)\) 和 \(P_2(x_2, y_2)\)
是直线上的两点，则 \(Ax_1 + By_1 + C = 0\) 且
\(Ax_2 + By_2 + C = 0\)。两式相减得

\[A(x_2 - x_1) + B(y_2 - y_1) = 0\]

即 \(\vec{n} \cdot \overrightarrow{P_1P_2} = 0\)，故
\(\vec{n} = (A, B)\) 垂直于直线的方向。\(\blacksquare\)

\subsubsection{25.1.2
直线的倾斜角}\label{ux76f4ux7ebfux7684ux503eux659cux89d2}

\textbf{定义 25.2（倾斜角）.} 直线 \(l\) 向上的方向与 \(x\)
轴正方向之间所成的最小正角 \(\alpha\)（\(0 \leq \alpha < \pi\)）称为直线
\(l\) 的\textbf{倾斜角}。

\textbf{命题 25.3.} 设直线的倾斜角为 \(\alpha\)，斜率为 \(k\)（当
\(\alpha \neq \frac{\pi}{2}\) 时），则

\[k = \tan\alpha\]

\textbf{证明}： 在直线上取两点 \(P_1(x_1, y_1)\) 和
\(P_2(x_2, y_2)\)，不妨设
\(x_2 > x_1\)。作直角三角形使斜边沿直线方向，则水平直角边长为
\(x_2 - x_1 > 0\)，竖直直角边长为 \(y_2 - y_1\)。由三角函数的定义（Ch
20），\(\tan\alpha = \frac{y_2 - y_1}{x_2 - x_1} = k\)。\(\blacksquare\)

\bigskip\hrule\bigskip

\subsection{25.2
两直线的位置关系}\label{ux4e24ux76f4ux7ebfux7684ux4f4dux7f6eux5173ux7cfb}

\subsubsection{25.2.1
平行与垂直的判定}\label{ux5e73ux884cux4e0eux5782ux76f4ux7684ux5224ux5b9a}

设直线
\(l_1: A_1 x + B_1 y + C_1 = 0\)，\(l_2: A_2 x + B_2 y + C_2 = 0\)。

\textbf{定理 25.5（平行条件）.} \(l_1 \parallel l_2\) 当且仅当
\(A_1 B_2 - A_2 B_1 = 0\) 且 \(A_1 C_2 - A_2 C_1 \neq 0\)（或
\(B_1 C_2 - B_2 C_1 \neq 0\)）。

\textbf{证明}：

\((\Rightarrow)\) 若 \(l_1 \parallel l_2\)，则它们的法向量
\(\vec{n}_1 = (A_1, B_1)\) 与 \(\vec{n}_2 = (A_2, B_2)\)
平行，即存在非零常数 \(\lambda\) 使得
\(\vec{n}_1 = \lambda \vec{n}_2\)，即
\(A_1 = \lambda A_2\)，\(B_1 = \lambda B_2\)。于是
\(A_1 B_2 - A_2 B_1 = \lambda A_2 B_2 - A_2 \lambda B_2 = 0\)。

若 \(A_1 C_2 - A_2 C_1 = 0\)，则 \(\lambda A_2 C_2 - A_2 C_1 = 0\)，即
\(A_2(\lambda C_2 - C_1) = 0\)。若 \(A_2 \neq 0\)，则
\(C_1 = \lambda C_2\)，此时
\(l_1: \lambda(A_2 x + B_2 y + C_2) = 0\)，即 \(l_1 = l_2\)，与
\(l_1 \parallel l_2\) 矛盾。若 \(A_2 = 0\)，则 \(B_2 \neq 0\)，由
\(B_1 C_2 - B_2 C_1 = \lambda B_2 C_2 - B_2 C_1 = B_2(\lambda C_2 - C_1) \neq 0\)
可得矛盾。因此 \(A_1 C_2 - A_2 C_1 \neq 0\)。

\((\Leftarrow)\) 若 \(A_1 B_2 - A_2 B_1 = 0\) 且
\(A_1 C_2 - A_2 C_1 \neq 0\)，则由 \(A_1 B_2 = A_2 B_1\) 可设
\(A_1 = \lambda A_2\)，\(B_1 = \lambda B_2\)（若 \(A_2 = B_2 = 0\) 则
\(l_2\) 不是直线）。假设 \(l_1\) 与 \(l_2\) 相交于点 \((x_0, y_0)\)，则
\(A_1 x_0 + B_1 y_0 + C_1 = 0\) 且
\(A_2 x_0 + B_2 y_0 + C_2 = 0\)。由第一式得
\(\lambda(A_2 x_0 + B_2 y_0) + C_1 = 0\)，由第二式得
\(A_2 x_0 + B_2 y_0 = -C_2\)，代入得 \(-\lambda C_2 + C_1 = 0\)，即
\(A_1 C_2 - A_2 C_1 = \lambda A_2 C_2 - A_2 C_1 = A_2(\lambda C_2 - C_1) = 0\)，与
\(A_1 C_2 - A_2 C_1 \neq 0\) 矛盾。故
\(l_1 \parallel l_2\)。\(\blacksquare\)

\textbf{定理 25.6（垂直条件）.} \(l_1 \perp l_2\) 当且仅当
\(A_1 A_2 + B_1 B_2 = 0\)。

\textbf{证明}： \(l_1\) 的法向量为 \(\vec{n}_1 = (A_1, B_1)\)，\(l_2\)
的法向量为 \(\vec{n}_2 = (A_2, B_2)\)。

\(l_1 \perp l_2\) \(\iff\)
\(\vec{n}_1 \perp \vec{n}_2\)（因为法向量与直线垂直，两条直线垂直等价于它们的法向量垂直）\(\iff\)
\(\vec{n}_1 \cdot \vec{n}_2 = 0\) \(\iff\)
\(A_1 A_2 + B_1 B_2 = 0\)。\(\blacksquare\)

\textbf{推论 25.1.} 若 \(l_1\) 的斜率为 \(k_1\)，\(l_2\) 的斜率为
\(k_2\)（均存在），则：

\begin{itemize}
\tightlist
\item
  \(l_1 \parallel l_2 \iff k_1 = k_2\)
\item
  \(l_1 \perp l_2 \iff k_1 k_2 = -1\)
\end{itemize}

\textbf{证明}： 设
\(l_1: y = k_1 x + b_1\)，\(l_2: y = k_2 x + b_2\)，即
\(l_1: k_1 x - y + b_1 = 0\)，\(l_2: k_2 x - y + b_2 = 0\)。

平行条件：\(k_1 \cdot (-1) - k_2 \cdot (-1) = 0\)，即
\(k_1 = k_2\)（还需 \(b_1 \neq b_2\) 保证不重合）。

垂直条件：\(k_1 k_2 + (-1)(-1) = 0\)，即 \(k_1 k_2 + 1 = 0\)，即
\(k_1 k_2 = -1\)。\(\blacksquare\)

\subsubsection{25.2.2
两直线的交点}\label{ux4e24ux76f4ux7ebfux7684ux4ea4ux70b9}

\textbf{定理 25.7.} 设
\(l_1: A_1 x + B_1 y + C_1 = 0\)，\(l_2: A_2 x + B_2 y + C_2 = 0\)，若
\(l_1\) 与 \(l_2\) 不平行（即
\(A_1 B_2 - A_2 B_1 \neq 0\)），则它们的交点坐标为

\[x = \frac{B_1 C_2 - B_2 C_1}{A_1 B_2 - A_2 B_1}, \quad y = \frac{C_1 A_2 - C_2 A_1}{A_1 B_2 - A_2 B_1}\]

\textbf{证明}： 联立方程组：

\[\begin{cases} A_1 x + B_1 y = -C_1 \\ A_2 x + B_2 y = -C_2 \end{cases}\]

系数行列式 \(D = A_1 B_2 - A_2 B_1 \neq 0\)。由克拉默法则（Ch 8.3）：

\[x = \frac{\begin{vmatrix} -C_1 & B_1 \\ -C_2 & B_2 \end{vmatrix}}{D} = \frac{-C_1 B_2 + C_2 B_1}{D} = \frac{B_1 C_2 - B_2 C_1}{A_1 B_2 - A_2 B_1}\]

\[y = \frac{\begin{vmatrix} A_1 & -C_1 \\ A_2 & -C_2 \end{vmatrix}}{D} = \frac{-A_1 C_2 + A_2 C_1}{D} = \frac{C_1 A_2 - C_2 A_1}{A_1 B_2 - A_2 B_1}\]

\(\blacksquare\)

\subsubsection{25.2.3
两直线的夹角}\label{ux4e24ux76f4ux7ebfux7684ux5939ux89d2}

\textbf{定义 25.3.} 两条直线 \(l_1, l_2\) 所成的\textbf{锐角}
\(\theta\)（\(0 < \theta \leq \frac{\pi}{2}\)）称为两直线的夹角。

\textbf{定理 25.8.} 设 \(l_1\) 的斜率为 \(k_1\)，\(l_2\) 的斜率为
\(k_2\)（均存在且 \(k_1 k_2 \neq -1\)），则两直线夹角 \(\theta\) 满足

\[\tan\theta = \left|\frac{k_1 - k_2}{1 + k_1 k_2}\right|\]

\textbf{证明}： 设 \(l_1\) 的倾斜角为 \(\alpha_1\)，\(l_2\) 的倾斜角为
\(\alpha_2\)，则
\(k_1 = \tan\alpha_1\)，\(k_2 = \tan\alpha_2\)。两直线的夹角
\(\theta = |\alpha_1 - \alpha_2|\)（取锐角）。

由两角差的正切公式（Ch 20.1）：

\[\tan(\alpha_1 - \alpha_2) = \frac{\tan\alpha_1 - \tan\alpha_2}{1 + \tan\alpha_1 \tan\alpha_2} = \frac{k_1 - k_2}{1 + k_1 k_2}\]

取绝对值得
\(\tan\theta = \left|\frac{k_1 - k_2}{1 + k_1 k_2}\right|\)。\(\blacksquare\)

\textbf{推论 25.2.} 用一般式的法向量表示：设
\(\vec{n}_1 = (A_1, B_1)\)，\(\vec{n}_2 = (A_2, B_2)\)，则

\[\cos\theta = \frac{|A_1 A_2 + B_1 B_2|}{\sqrt{A_1^2 + B_1^2} \cdot \sqrt{A_2^2 + B_2^2}}\]

\textbf{证明}： 两法向量的夹角 \(\varphi\) 满足
\(\cos\varphi = \frac{\vec{n}_1 \cdot \vec{n}_2}{|\vec{n}_1||\vec{n}_2|}\)。由于直线夹角
\(\theta\) 等于法向量夹角或其补角，取绝对值即得
\(\cos\theta = |\cos\varphi|\)。\(\blacksquare\)

\bigskip\hrule\bigskip

\subsection{25.3 距离公式}\label{ux8dddux79bbux516cux5f0f}

\subsubsection{25.3.1
点到直线的距离}\label{ux70b9ux5230ux76f4ux7ebfux7684ux8dddux79bb}

\textbf{定理 25.9（点到直线距离公式）.} 点 \(P(x_0, y_0)\) 到直线
\(l: Ax + By + C = 0\) 的距离为

\[d = \frac{|Ax_0 + By_0 + C|}{\sqrt{A^2 + B^2}}\]

\textbf{证明}： 设直线 \(l\) 的法向量为 \(\vec{n} = (A, B)\)。在 \(l\)
上取一点 \(Q(x_1, y_1)\)，则 \(Ax_1 + By_1 + C = 0\)。

向量 \(\overrightarrow{QP} = (x_0 - x_1, y_0 - y_1)\)。\(P\) 到 \(l\)
的距离等于 \(\overrightarrow{QP}\) 在法向量方向上投影的绝对值（因为从
\(Q\) 到直线 \(l\) 的垂线方向即为法向量方向）：

\[d = \frac{|\overrightarrow{QP} \cdot \vec{n}|}{|\vec{n}|} = \frac{|A(x_0 - x_1) + B(y_0 - y_1)|}{\sqrt{A^2 + B^2}}\]

由于 \(Ax_1 + By_1 = -C\)，代入得：

\[d = \frac{|Ax_0 + By_0 - (Ax_1 + By_1)|}{\sqrt{A^2 + B^2}} = \frac{|Ax_0 + By_0 + C|}{\sqrt{A^2 + B^2}}\]

\(\blacksquare\)

\textbf{推论 25.3（两平行线间的距离）.} 设平行线
\(l_1: Ax + By + C_1 = 0\)，\(l_2: Ax + By + C_2 = 0\)（\(A, B\)
相同保证平行，\(C_1 \neq C_2\) 保证两线不重合），则它们之间的距离为

\[d = \frac{|C_1 - C_2|}{\sqrt{A^2 + B^2}}\]

\textbf{证明}： 在 \(l_1\) 上取点 \(P_0(x_0, y_0)\)，则
\(Ax_0 + By_0 + C_1 = 0\)。\(P_0\) 到 \(l_2\) 的距离为

\[d = \frac{|Ax_0 + By_0 + C_2|}{\sqrt{A^2 + B^2}} = \frac{|-C_1 + C_2|}{\sqrt{A^2 + B^2}} = \frac{|C_1 - C_2|}{\sqrt{A^2 + B^2}}\]

由于平行线间距离处处相等，此即所求。\(\blacksquare\)

\bigskip\hrule\bigskip

\subsection{25.4 圆的方程}\label{ux5706ux7684ux65b9ux7a0b}

\subsubsection{25.4.1
圆的标准方程}\label{ux5706ux7684ux6807ux51c6ux65b9ux7a0b}

在 Ch 13.4
中，我们从几何角度研究了圆的性质（垂径定理、圆周角定理等）。现在我们在坐标系中用方程来描述圆。

\textbf{定义 25.4（圆的解析定义）.} 平面内到定点 \(C(a, b)\)
的距离等于定长 \(r\)（\(r > 0\)）的点的轨迹称为\textbf{圆}，记作
\(\odot(C, r)\)。定点 \(C\) 称为\textbf{圆心}，定长 \(r\)
称为\textbf{半径}。

\textbf{定理 25.10（圆的标准方程）.} 圆心为 \(C(a, b)\)、半径为 \(r\)
的圆的方程为

\[(x - a)^2 + (y - b)^2 = r^2\]

\textbf{证明}： 设 \(P(x, y)\) 为圆上任意一点，由定义
\(|PC| = r\)。由两点间距离公式（Ch 16.1）：

\[\sqrt{(x - a)^2 + (y - b)^2} = r\]

两边平方即得 \((x - a)^2 + (y - b)^2 = r^2\)。

反过来，若 \(P(x, y)\) 满足此方程，则
\(|PC| = \sqrt{(x-a)^2 + (y-b)^2} = r\)，\(P\) 在圆上。\(\blacksquare\)

\textbf{推论 25.4.} 圆心在原点、半径为 \(r\) 的圆的方程为
\(x^2 + y^2 = r^2\)。

\textbf{证明}：圆心在原点、半径为 \(r\) 的圆是到原点距离为 \(r\)
的点的轨迹，即 \(\sqrt{x^2+y^2} = r\)，两边平方得
\(x^2+y^2 = r^2\)。\(\blacksquare\)

\subsubsection{25.4.2
圆的一般方程}\label{ux5706ux7684ux4e00ux822cux65b9ux7a0b}

\textbf{定理 25.11（圆的一般方程）.} 圆心为 \((a, b)\)、半径为 \(r\)
的圆的方程可写为

\[x^2 + y^2 + Dx + Ey + F = 0\]

其中 \(D = -2a\)，\(E = -2b\)，\(F = a^2 + b^2 - r^2\)。反之，方程
\(x^2 + y^2 + Dx + Ey + F = 0\) 当 \(D^2 + E^2 - 4F > 0\)
时表示一个圆，圆心为 \(\left(-\frac{D}{2}, -\frac{E}{2}\right)\)，半径为
\(r = \frac{1}{2}\sqrt{D^2 + E^2 - 4F}\)。

\textbf{证明}：（正向）将标准方程 \((x - a)^2 + (y - b)^2 = r^2\) 展开：

\[x^2 - 2ax + a^2 + y^2 - 2by + b^2 = r^2\]

\[x^2 + y^2 + (-2a)x + (-2b)y + (a^2 + b^2 - r^2) = 0\]

令 \(D = -2a\)，\(E = -2b\)，\(F = a^2 + b^2 - r^2\)，即得一般方程。

（反向）给定 \(x^2 + y^2 + Dx + Ey + F = 0\)，配方得：

\[\left(x + \frac{D}{2}\right)^2 + \left(y + \frac{E}{2}\right)^2 = \frac{D^2 + E^2 - 4F}{4}\]

当 \(D^2 + E^2 - 4F > 0\) 时，右边为正，表示圆心
\(\left(-\frac{D}{2}, -\frac{E}{2}\right)\)、半径
\(\frac{1}{2}\sqrt{D^2 + E^2 - 4F}\) 的圆。

当 \(D^2 + E^2 - 4F = 0\) 时，表示一个点
\(\left(-\frac{D}{2}, -\frac{E}{2}\right)\)（退化的圆，半径为 \(0\)）。

当 \(D^2 + E^2 - 4F < 0\) 时，不表示任何实图形（空集）。\(\blacksquare\)

\textbf{证明}： 由定理 25.11
的推导，\(r^2 = \frac{D^2 + E^2 - 4F}{4}\)。\(r^2 > 0\)
时为圆，\(r^2 = 0\) 时为点，\(r^2 < 0\) 时无实数解。\(\blacksquare\)

\subsubsection{25.4.3
圆的切线方程}\label{ux5706ux7684ux5207ux7ebfux65b9ux7a0b}

\textbf{定理 25.12（过圆上一点的切线方程）.} 过圆
\((x - a)^2 + (y - b)^2 = r^2\) 上一点 \(P_0(x_0, y_0)\) 的切线方程为

\[(x_0 - a)(x - a) + (y_0 - b)(y - b) = r^2\]

\textbf{证明}： 切线在 \(P_0\) 处与圆相切，因此切线垂直于半径
\(\overrightarrow{CP_0}\)，其中 \(C(a, b)\) 为圆心。

半径方向向量 \(\overrightarrow{CP_0} = (x_0 - a, y_0 - b)\)。切线过
\(P_0\) 且垂直于 \(\overrightarrow{CP_0}\)，切线上任一点 \(P(x, y)\)
满足 \(\overrightarrow{P_0P} \perp \overrightarrow{CP_0}\)，即

\[(x_0 - a)(x - x_0) + (y_0 - b)(y - y_0) = 0\]

展开：

\[(x_0 - a)x - (x_0 - a)x_0 + (y_0 - b)y - (y_0 - b)y_0 = 0\]

\[(x_0 - a)x + (y_0 - b)y = (x_0 - a)x_0 + (y_0 - b)y_0\]

由于 \(P_0\) 在圆上，\((x_0 - a)^2 + (y_0 - b)^2 = r^2\)。于是

\[(x_0 - a)x_0 + (y_0 - b)y_0 = (x_0 - a)(x_0 - a + a) + (y_0 - b)(y_0 - b + b)\]
\[= (x_0 - a)^2 + a(x_0 - a) + (y_0 - b)^2 + b(y_0 - b) = r^2 + a(x_0 - a) + b(y_0 - b)\]

将此代入并移项：

\[(x_0 - a)(x - a) + (y_0 - b)(y - b) = r^2\]

\(\blacksquare\)

\textbf{定理 25.13（圆外一点到圆的切线长）.} 从圆外一点 \(P(x_0, y_0)\)
到圆 \((x - a)^2 + (y - b)^2 = r^2\) 的切线长为

\[l = \sqrt{(x_0 - a)^2 + (y_0 - b)^2 - r^2}\]

\textbf{证明}： 设切点为 \(T\)，圆心为 \(C(a, b)\)。需要证明
\(CT \perp PT\)（半径垂直于过切点的直线）。下面用解析方法独立证明。圆
\((x-a)^2 + (y-b)^2 = r^2\) 在点 \(T(x_1, y_1)\) 处的切线方程为
\((x_1-a)(x-a) + (y_1-b)(y-b) = r^2\)（由定理 25.12
导出：将圆方程视为隐函数并取微分，或利用圆心到切线的距离等于半径）。切线方向向量可取
\(\vec{d} = (y_1-b,\, -(x_1-a))\)（与法向量 \((x_1-a,\, y_1-b)\)
正交）。半径向量
\(\overrightarrow{CT} = (x_1-a,\, y_1-b)\)，点积为零，故
\(CT \perp PT\)。于是 \(\triangle CPT\) 为直角三角形（直角在 \(T\)
处）。

由勾股定理：\(|CP|^2 = |CT|^2 + |PT|^2\)，即

\[(x_0 - a)^2 + (y_0 - b)^2 = r^2 + l^2\]

\[l = \sqrt{(x_0 - a)^2 + (y_0 - b)^2 - r^2}\]

其中需 \((x_0 - a)^2 + (y_0 - b)^2 > r^2\)（\(P\)
在圆外）。\(\blacksquare\)

\bigskip\hrule\bigskip

\subsection{25.5
直线与圆的位置关系}\label{ux76f4ux7ebfux4e0eux5706ux7684ux4f4dux7f6eux5173ux7cfb}

\subsubsection{25.5.1
位置关系的判定}\label{ux4f4dux7f6eux5173ux7cfbux7684ux5224ux5b9a}

\textbf{定理 25.14.} 设圆 \(C\) 的方程为
\((x - a)^2 + (y - b)^2 = r^2\)，直线 \(l\) 的方程为
\(Ax + By + C_0 = 0\)。圆心 \(C(a, b)\) 到直线 \(l\) 的距离为

\[d = \frac{|Aa + Bb + C_0|}{\sqrt{A^2 + B^2}}\]

则：

\begin{itemize}
\tightlist
\item
  \(d > r\) \(\iff\) 直线与圆\textbf{相离}（无公共点）；
\item
  \(d = r\) \(\iff\) 直线与圆\textbf{相切}（恰有一个公共点）；
\item
  \(d < r\) \(\iff\) 直线与圆\textbf{相交}（有两个公共点）。
\end{itemize}

\textbf{证明}： 过圆心 \(C\) 作直线 \(l\) 的垂线，垂足为 \(H\)，则
\(|CH| = d\)。

\textbf{情形 1：} \(d > r\)。设 \(P\) 为 \(l\)
上任意一点，\(\triangle CHP\) 中 \(\angle CHP = 90^\circ\)，由勾股定理
\(|CP|^2 = d^2 + |HP|^2 \geq d^2 > r^2\)，故 \(|CP| > r\)，\(P\)
不在圆上。因此直线与圆无公共点。

\textbf{情形 2：} \(d = r\)。此时 \(H\) 在圆上（\(|CH| = r\)），且 \(H\)
在直线 \(l\) 上，故 \(H\) 是公共点。对 \(l\) 上其他点 \(P \neq H\)，同理
\(|CP| > r\)，故 \(P\) 不在圆上。\(H\) 是唯一的公共点，即相切。

\textbf{情形 3：} \(d < r\)。在直线 \(l\) 上取 \(H\) 两侧的点
\(P_1, P_2\)，使 \(|HP_1| = |HP_2| = \sqrt{r^2 - d^2}\)。则

\[|CP_1|^2 = d^2 + (\sqrt{r^2 - d^2})^2 = r^2\]

故 \(|CP_1| = r\)，\(P_1\) 在圆上。同理 \(P_2\)
也在圆上。\(P_1 \neq P_2\)（在 \(H\)
两侧），故有两个交点。\(\blacksquare\)

\subsubsection{25.5.2 弦长公式}\label{ux5f26ux957fux516cux5f0f}

\textbf{定理 25.15.} 直线与圆相交时，弦长为

\[|AB| = 2\sqrt{r^2 - d^2}\]

其中 \(d\) 为圆心到直线的距离。

\textbf{证明}： 设弦 \(AB\) 的中点为 \(M\)。由圆的对称性（或垂径定理，Ch
13.4），\(CM \perp AB\)，且 \(|AM| = |MB|\)。在直角三角形
\(\triangle CMA\)
中，\(\angle CMA = 90^\circ\)，\(|CA| = r\)，\(|CM| = d\)，由勾股定理：

\[|AM| = \sqrt{r^2 - d^2}\]

故 \(|AB| = 2|AM| = 2\sqrt{r^2 - d^2}\)。\(\blacksquare\)

\subsubsection{25.5.3
两圆的位置关系}\label{ux4e24ux5706ux7684ux4f4dux7f6eux5173ux7cfb}

\textbf{定理 25.16.} 设 \(\odot C_1\) 的圆心为 \(C_1(a_1, b_1)\)、半径为
\(r_1\)，\(\odot C_2\) 的圆心为 \(C_2(a_2, b_2)\)、半径为
\(r_2\)。记两圆心距离
\(d = |C_1 C_2| = \sqrt{(a_1 - a_2)^2 + (b_1 - b_2)^2}\)。不妨设
\(r_1 \geq r_2\)，则：

{\def\LTcaptype{none} % do not increment counter
\begin{longtable}[]{@{}lll@{}}
\toprule\noalign{}
位置关系 & 条件 & 公共点个数 \\
\midrule\noalign{}
\endhead
\bottomrule\noalign{}
\endlastfoot
外离 & \(d > r_1 + r_2\) & 0 \\
外切 & \(d = r_1 + r_2\) & 1 \\
相交 & \(r_1 - r_2 < d < r_1 + r_2\) & 2 \\
内切 & \(d = r_1 - r_2\)（\(r_1 > r_2\)） & 1 \\
内含 & \(d < r_1 - r_2\)（\(r_1 > r_2\)） & 0 \\
\end{longtable}
}

\textbf{证明}： 设 \(P\) 为公共点，则
\(|PC_1| = r_1\)，\(|PC_2| = r_2\)。由三角不等式：

\[|r_1 - r_2| \leq |PC_1| - |PC_2| \leq |C_1 C_2| \leq |PC_1| + |PC_2| = r_1 + r_2\]

等等，需要更仔细地分析。由 \(|PC_1| = r_1\)，\(|PC_2| = r_2\)
和三角不等式：

\[|r_1 - r_2| = ||PC_1| - |PC_2|| \leq |C_1 C_2| \leq |PC_1| + |PC_2| = r_1 + r_2\]

因此公共点存在的必要条件是 \(|r_1 - r_2| \leq d \leq r_1 + r_2\)。

下面对五种情形逐一验证：

\textbf{(i) 外离}（\(d > r_1 + r_2\)）：由必要条件 \(d \leq r_1 + r_2\)
不成立，故无公共点。

\textbf{(ii) 外切}（\(d = r_1 + r_2\)）：在 \(C_1 C_2\) 连线上取点
\(P = C_1 + \frac{r_1}{d}\overrightarrow{C_1 C_2}\)，则
\(|PC_1| = r_1\)，\(|PC_2| = d - r_1 = r_2\)，故 \(P\)
是公共点。若另有公共点 \(Q\)，则由三角不等式
\(|C_1 C_2| \leq |QC_1| + |QC_2| = r_1 + r_2 = d\)，等号成立当且仅当
\(Q\) 在线段 \(C_1 C_2\) 上，故 \(Q = P\)。公共点个数为 \(1\)。

\textbf{(iii) 相交}（\(r_1 - r_2 < d < r_1 + r_2\)）：两圆满足
\(|r_1 - r_2| < d < r_1 + r_2\)。用解析方法严格证明。设两圆方程分别为

\[(x - x_1)^2 + (y - y_1)^2 = r_1^2, \quad (x - x_2)^2 + (y - y_2)^2 = r_2^2\]

不妨设圆心 \(C_1(0, 0)\)，\(C_2(d, 0)\)（通过平移旋转），则方程化为

\[x^2 + y^2 = r_1^2, \quad (x - d)^2 + y^2 = r_2^2\]

相减得 \(x^2 + y^2 - [(x-d)^2 + y^2] = r_1^2 - r_2^2\)，即
\(2dx - d^2 = r_1^2 - r_2^2\)，解得

\[x = \frac{r_1^2 - r_2^2 + d^2}{2d}\]

代入第一圆方程：\(y^2 = r_1^2 - x^2\)。计算判别式：

\[r_1^2 - x^2 = \frac{4d^2 r_1^2 - (r_1^2 - r_2^2 + d^2)^2}{4d^2}\]

分子可因式分解为
\([(r_1 + r_2)^2 - d^2][d^2 - (r_1 - r_2)^2] = (r_1 + r_2 + d)(r_1 + r_2 - d)(d + r_1 - r_2)(d - r_1 + r_2)\)。在
\(|r_1 - r_2| < d < r_1 + r_2\) 条件下，四个因子全为正，故
\(y^2 > 0\)。\(y = \pm\sqrt{y^2}\)
给出两个相异的实数解，对应两圆恰有两个公共点。\(\blacksquare\)

\textbf{(iv) 内切}（\(d = r_1 - r_2\)，\(r_1 > r_2\)）：在 \(C_1 C_2\)
连线的延长线上取点 \(P\) 使得 \(|PC_2| = r_2\) 且 \(P\) 在 \(C_1\) 与
\(C_2\) 之间，则 \(|PC_1| = d + r_2 = r_1\)，故 \(P\) 是公共点。类似
(ii) 的论证，公共点唯一。公共点个数为 \(1\)。

\textbf{(v) 内含}（\(d < r_1 - r_2\)，\(r_1 > r_2\)）：对任意点
\(P\)，由三角不等式
\(|PC_1| \leq |PC_2| + d < |PC_2| + (r_1 - r_2)\)。若
\(|PC_1| = r_1\)，则 \(|PC_2| > r_2\)，故 \(P\) 不在 \(\odot C_2\)
上。同理若 \(|PC_2| = r_2\)，则
\(|PC_1| < r_1\)。因此无公共点。\(\blacksquare\)

\subsubsection{25.5.4
圆幂定理的解析证明}\label{ux5706ux5e42ux5b9aux7406ux7684ux89e3ux6790ux8bc1ux660e}

在 Ch 13.4 中，我们从纯几何角度证明了圆幂定理。现在用解析方法重新证明。

\textbf{定理 25.17（圆幂定理）.} 设圆 \(\odot C\) 的方程为
\((x - a)^2 + (y - b)^2 = r^2\)，\(P(x_0, y_0)\) 为圆外一点。过 \(P\)
作割线交圆于 \(A, B\) 两点，则

\[|PA| \cdot |PB| = |PC|^2 - r^2\]

且此值与割线的选择无关（称为点 \(P\) 关于圆 \(\odot C\)
的\textbf{圆幂}）。

\textbf{证明}： 设过 \(P\) 的直线参数方程为

\[\begin{cases} x = x_0 + t\cos\alpha \\ y = y_0 + t\sin\alpha \end{cases}\]

代入圆的方程：

\[(x_0 + t\cos\alpha - a)^2 + (y_0 + t\sin\alpha - b)^2 = r^2\]

设 \(u = x_0 - a\)，\(v = y_0 - b\)，展开：

\[u^2 + 2ut\cos\alpha + t^2\cos^2\alpha + v^2 + 2vt\sin\alpha + t^2\sin^2\alpha = r^2\]

\[t^2 + 2(u\cos\alpha + v\sin\alpha)t + (u^2 + v^2 - r^2) = 0\]

由韦达定理（Ch 8），两个根 \(t_1, t_2\) 满足：

\[t_1 t_2 = u^2 + v^2 - r^2 = (x_0 - a)^2 + (y_0 - b)^2 - r^2 = |PC|^2 - r^2\]

由参数 \(t\) 的几何意义：直线方程为 \(P = P_0 + t \vec{d}\)，其中
\(\vec{d} = (\cos\alpha, \sin\alpha)\) 为单位方向向量。对参数值 \(t_i\)
对应的点 \(A_i\)，有

\[|P_0 A_i| = |t_i \vec{d}| = |t_i| \cdot |\vec{d}| = |t_i|\]

因此 \(|t_1| = |PA|\)，\(|t_2| = |PB|\)。当 \(P\) 在圆外时，\(t_1, t_2\)
同号（均在 \(P\) 的同侧或异侧，取决于方向），故

\[|PA| \cdot |PB| = |t_1 t_2| = |PC|^2 - r^2\]

（\(P\) 在圆外时 \(|PC|^2 > r^2\)，故
\(t_1 t_2 > 0\)，无需取绝对值。）\(\blacksquare\)

\textbf{推论 25.5（切割线定理）.} 若过 \(P\) 的直线与圆相切于 \(T\)，则
\(|PT|^2 = |PA| \cdot |PB|\)。

\textbf{证明}：
当割线退化为切线时，\(A = B = T\)，\(|PA| \cdot |PB| = |PT|^2 = |PC|^2 - r^2\)（由定理
25.13）。\(\blacksquare\)

\textbf{推论 25.6（割线定理）.} 若过 \(P\) 的两条割线分别交圆于 \(A, B\)
和 \(C, D\)，则 \(|PA| \cdot |PB| = |PC| \cdot |PD|\)。

\textbf{证明}： 两条割线对应的圆幂相同，均为
\(|PC|^2 - r^2\)。\(\blacksquare\)

\subsubsection{25.5.5 根轴}\label{ux6839ux8f74}

\textbf{定义 25.5.} 两个不同心的圆 \(\odot C_1\) 和 \(\odot C_2\)
的\textbf{根轴}是到两圆的圆幂相等的点的轨迹。

\textbf{定理 25.18.} 设
\(\odot C_1: x^2 + y^2 + D_1 x + E_1 y + F_1 = 0\)，\(\odot C_2: x^2 + y^2 + D_2 x + E_2 y + F_2 = 0\)，则根轴的方程为

\[(D_1 - D_2)x + (E_1 - E_2)y + (F_1 - F_2) = 0\]

\textbf{证明}： 点 \(P(x_0, y_0)\) 关于 \(\odot C_1\) 的圆幂为
\(x_0^2 + y_0^2 + D_1 x_0 + E_1 y_0 + F_1\)（由一般方程直接计算）。关于
\(\odot C_2\) 的圆幂为 \(x_0^2 + y_0^2 + D_2 x_0 + E_2 y_0 + F_2\)。

两圆幂相等：

\[(D_1 - D_2)x_0 + (E_1 - E_2)y_0 + (F_1 - F_2) = 0\]

这是一条直线。\(\blacksquare\)

\bigskip\hrule\bigskip

\subsection{25.6
圆的切线与导数}\label{ux5706ux7684ux5207ux7ebfux4e0eux5bfcux6570}

\subsubsection{25.6.1
用导数求切线方程}\label{ux7528ux5bfcux6570ux6c42ux5207ux7ebfux65b9ux7a0b}

在 Ch 23 中，我们学习了导数的几何意义：函数 \(y = f(x)\) 在 \(x = x_0\)
处的导数 \(f'(x_0)\)
等于曲线在该点处切线的斜率。我们可以用这一工具来求圆的切线方程。

\textbf{定理 25.19.} 圆 \(x^2 + y^2 = r^2\) 在点
\(P_0(x_0, y_0)\)（\(y_0 \neq 0\)）处的切线斜率为

\[k = -\frac{x_0}{y_0}\]

\textbf{证明}： 由隐函数求导（Ch 23），对方程 \(x^2 + y^2 = r^2\)
两边关于 \(x\) 求导：

\[2x + 2y \frac{dy}{dx} = 0\]

\[\frac{dy}{dx} = -\frac{x}{y}\]

在点 \(P_0(x_0, y_0)\)
处，\(k = \frac{dy}{dx}\big|_{(x_0, y_0)} = -\frac{x_0}{y_0}\)。\(\blacksquare\)

\textbf{验证：} 由定理 25.12，切线方程为 \(x_0 x + y_0 y = r^2\)，即
\(y = -\frac{x_0}{y_0}x + \frac{r^2}{y_0}\)，斜率为
\(-\frac{x_0}{y_0}\)，与导数法一致。\(\blacksquare\)

\textbf{推论 25.7.} 圆 \((x - a)^2 + (y - b)^2 = r^2\) 在点
\(P_0(x_0, y_0)\)（\(y_0 \neq b\)）处的切线斜率为

\[k = -\frac{x_0 - a}{y_0 - b}\]

\textbf{证明}： 令 \(X = x - a\)，\(Y = y - b\)，方程变为
\(X^2 + Y^2 = r^2\)。对 \(X\)
求导：\(2X + 2Y\frac{dY}{dX} = 0\)，\(\frac{dY}{dX} = -\frac{X}{Y}\)。由于
\(\frac{dy}{dx} = \frac{dY}{dX}\)，在点 \((x_0, y_0)\) 处
\(k = -\frac{x_0 - a}{y_0 - b}\)。\(\blacksquare\)

\subsubsection{25.6.2
切线的存在性}\label{ux5207ux7ebfux7684ux5b58ux5728ux6027}

\textbf{命题 25.4.} 过圆外一点 \(P\)
到圆恰好有两条切线；过圆上一点恰好有一条切线；过圆内一点没有切线。

\textbf{证明}： 设圆为
\((x - a)^2 + (y - b)^2 = r^2\)，\(P(x_0, y_0)\)。

切线过 \(P\) 且与圆相切，等价于 \(P\) 到直线的距离等于
\(r\)。设切线斜率为 \(k\)（存在时），则切线方程为
\(y - y_0 = k(x - x_0)\)，即 \(kx - y + (y_0 - kx_0) = 0\)。圆心
\((a, b)\) 到此直线的距离为

\[d = \frac{|ka - b + y_0 - kx_0|}{\sqrt{k^2 + 1}} = \frac{|k(a - x_0) - (b - y_0)|}{\sqrt{k^2 + 1}} = r\]

两边平方：

\[[k(a - x_0) - (b - y_0)]^2 = r^2(k^2 + 1)\]

设 \(p = a - x_0\)，\(q = b - y_0\)，展开：

\[(kp - q)^2 = r^2(k^2 + 1)\] \[k^2 p^2 - 2kpq + q^2 = r^2 k^2 + r^2\]
\[(p^2 - r^2)k^2 - 2pq \cdot k + (q^2 - r^2) = 0\]

关于 \(k\) 的二次方程，判别式为

\[\Delta = 4p^2q^2 - 4(p^2 - r^2)(q^2 - r^2) = 4[p^2q^2 - p^2q^2 + r^2 p^2 + r^2 q^2 - r^4]\]
\[= 4r^2(p^2 + q^2 - r^2) = 4r^2[(x_0 - a)^2 + (y_0 - b)^2 - r^2]\]

\begin{itemize}
\tightlist
\item
  若 \(P\)
  在圆外：\((x_0 - a)^2 + (y_0 - b)^2 > r^2\)，\(\Delta > 0\)，两个不同的实数解，两条切线。
\item
  若 \(P\)
  在圆上：\((x_0 - a)^2 + (y_0 - b)^2 = r^2\)，\(\Delta = 0\)，一个实数解（重根），一条切线。
\item
  若 \(P\)
  在圆内：\((x_0 - a)^2 + (y_0 - b)^2 < r^2\)，\(\Delta < 0\)，无实数解，无切线。
\end{itemize}

还需考虑斜率不存在（切线垂直于 \(x\) 轴）的情况，即 \(x = x_0\)
为切线。此时圆心到直线的距离 \(|a - x_0| = r\)，即
\((x_0 - a)^2 = r^2\)，此时 \(x_0 = a \pm r\)，切线方程为 \(x = a + r\)
或 \(x = a - r\)。这种情况已被上述分析覆盖（对应二次方程中 \(p^2 = r^2\)
时退化为一次方程的情形）。\(\blacksquare\)

\bigskip\hrule\bigskip

\begin{quote}
\textbf{Ch 25 本章展望}

本章建立了直线与圆的解析理论：用一次方程描述直线，用二次方程描述圆。我们证明了直线的各种方程形式（点斜式、两点式、截距式、一般式）之间的等价性，推导了点到直线的距离公式和两直线夹角公式，并系统研究了圆的标准方程、一般方程以及直线与圆的位置关系。特别地，我们利用导数工具验证了圆的切线方程，建立了几何与分析的桥梁。

然而，直线和圆只是平面曲线的''冰山一角''。自然界中存在大量更为复杂的曲线------行星的椭圆轨道、抛射体的抛物线轨迹、冷却塔截面的双曲线------它们有一个共同的名字：\textbf{圆锥曲线}。在下一章中，我们将从统一的几何定义出发，推导椭圆、双曲线和抛物线的标准方程，并揭示它们之间深刻的内在联系。
\end{quote}

\bigskip\hrule\bigskip

\subsection{Ch 26 圆锥曲线}\label{ch-26-ux5706ux9525ux66f2ux7ebf}

\begin{quote}
\textbf{前置知识}

本章建立在以下前序章节的基础之上：

\begin{itemize}
\tightlist
\item
  \textbf{Ch 25（直线与圆）}：圆是椭圆的特殊情形（离心率
  \(e = 0\)）。圆的标准方程推导方法（距离公式 \(\to\)
  平方化简）直接推广为椭圆方程的推导。点到直线距离公式在推导抛物线方程时不可或缺。
\item
  \textbf{Ch
  13.3（相似）}：椭圆可以视为圆沿一个方向均匀拉伸的结果，这一变换保持了相似结构。相似比在椭圆的长短轴关系中有所体现。
\item
  \textbf{Ch 15（向量）}：焦点向量关系 \(|PF_1| + |PF_2| = 2a\)
  等定义中涉及向量的模长。向量法在处理圆锥曲线的焦点弦、准线等问题时提供简洁的表达。
\item
  \textbf{Ch
  8（方程与方程组）}：圆锥曲线的方程是二元二次方程，韦达定理在处理直线与圆锥曲线的交点问题（设而不求）中至关重要。
\item
  \textbf{Ch 20（三角恒等变换）}：双曲线的参数方程涉及 \(\sec\theta\) 和
  \(\tan\theta\) 的恒等式
  \(\sec^2\theta - \tan^2\theta = 1\)，与双曲线方程的结构完美呼应。
\end{itemize}
\end{quote}

\bigskip\hrule\bigskip

\subsection{26.1
圆锥曲线的统一定义}\label{ux5706ux9525ux66f2ux7ebfux7684ux7edfux4e00ux5b9aux4e49}

\subsubsection{26.1.1 圆锥截线}\label{ux5706ux9525ux622aux7ebf}

圆锥曲线（conic
section）的历史可以追溯到古希腊。阿波罗尼奥斯（Apollonius，约前 262--前
190）发现，用平面截割圆锥面时，根据截面与圆锥母线的夹角不同，可以得到三种不同的曲线：椭圆、双曲线和抛物线。

\textbf{定义 26.1（圆锥曲线的统一定义）.} 平面内，到一个定点
\(F\)（\textbf{焦点}）和一条定直线
\(l\)（\textbf{准线}，\(F \notin l\)）的距离之比等于常数
\(e\)（\(e > 0\)）的点的轨迹称为\textbf{圆锥曲线}。即

\[C = \left\{P \;\middle|\; \frac{|PF|}{d(P, l)} = e\right\}\]

其中 \(e\) 称为\textbf{离心率}（eccentricity），\(F\)
称为\textbf{焦点}（focus），\(l\) 称为\textbf{准线}（directrix）。

\begin{itemize}
\tightlist
\item
  当 \(0 < e < 1\) 时，曲线为\textbf{椭圆}；
\item
  当 \(e = 1\) 时，曲线为\textbf{抛物线}；
\item
  当 \(e > 1\) 时，曲线为\textbf{双曲线}。
\end{itemize}

\textbf{注} 当 \(e > 1\) 时，以焦点 \(F\) 和对应准线 \(l\)
定义的圆锥曲线（满足
\(\frac{|PF|}{d(P, l)} = e\)）只给出双曲线的\textbf{一支}（靠近焦点
\(F\)
的那支）。要得到完整的双曲线（两支），需要分别以两个焦点为极点各定义一次，或采用定义
26.4 的两焦点距离差定义。具体而言：以焦点 \(F_1\) 和准线 \(l_1\)
定义给出双曲线的左支，以焦点 \(F_2\) 和准线 \(l_2\)
定义给出双曲线的右支；两支的并集才是完整的双曲线。

圆的特殊情况：当 \(e = 0\) 时，焦点到准线的约束消失，曲线退化为以 \(F\)
为圆心的圆。

\subsubsection{26.1.2 焦准距}\label{ux7126ux51c6ux8ddd}

\textbf{定义 26.2.} 焦点到准线的距离 \(p = d(F, l)\)
称为\textbf{焦准距}（focal parameter）。

在后续推导中，我们将证明圆锥曲线的统一极坐标方程（Ch 27.4）中，\(p\)
扮演关键角色。

\bigskip\hrule\bigskip

\subsection{26.2 椭圆}\label{ux692dux5706}

\subsubsection{26.2.1 椭圆的定义}\label{ux692dux5706ux7684ux5b9aux4e49}

\textbf{定义 26.3（椭圆的焦点定义）.} 平面内到两个定点
\(F_1, F_2\)（称为\textbf{焦点}）的距离之和等于常数
\(2a\)（\(2a > |F_1 F_2|\)）的点的轨迹称为\textbf{椭圆}（ellipse）。即

\[E = \{P \mid |PF_1| + |PF_2| = 2a, \; 2a > |F_1 F_2|\}\]

常数 \(2a\) 必须大于两焦点间距 \(|F_1 F_2|\)，否则轨迹不存在或退化。

\textbf{命题 26.1.} 若 \(2a = |F_1 F_2|\)，则轨迹为线段 \(F_1 F_2\)；若
\(2a < |F_1 F_2|\)，则轨迹为空集。

\textbf{证明}： 若 \(2a = |F_1 F_2|\)，由三角不等式
\(|PF_1| + |PF_2| \geq |F_1 F_2|\)，等号成立当且仅当 \(P\) 在线段
\(F_1 F_2\) 上。反之，线段上每一点满足
\(|PF_1| + |PF_2| = |F_1 F_2| = 2a\)。

若 \(2a < |F_1 F_2|\)，由三角不等式
\(|PF_1| + |PF_2| \geq |F_1 F_2| > 2a\)，不存在满足条件的点。\(\blacksquare\)

\subsubsection{26.2.2
椭圆的标准方程}\label{ux692dux5706ux7684ux6807ux51c6ux65b9ux7a0b}

\textbf{定理 26.1（椭圆的标准方程）.} 设椭圆的焦点为
\(F_1(-c, 0)\)，\(F_2(c, 0)\)（\(c > 0\)），\(2a > 2c\)。令
\(b^2 = a^2 - c^2\)（\(b > 0\)），则椭圆的标准方程为

\[\frac{x^2}{a^2} + \frac{y^2}{b^2} = 1 \quad (a > b > 0)\]

\textbf{证明}： 设 \(P(x, y)\) 为椭圆上任意一点。由定义：

\[\sqrt{(x + c)^2 + y^2} + \sqrt{(x - c)^2 + y^2} = 2a \quad \cdots (*)\]

移项：

\[\sqrt{(x + c)^2 + y^2} = 2a - \sqrt{(x - c)^2 + y^2}\]

两边平方：

\[(x + c)^2 + y^2 = 4a^2 - 4a\sqrt{(x - c)^2 + y^2} + (x - c)^2 + y^2\]

展开并化简：

\[x^2 + 2cx + c^2 = 4a^2 - 4a\sqrt{(x - c)^2 + y^2} + x^2 - 2cx + c^2\]

\[4cx = 4a^2 - 4a\sqrt{(x - c)^2 + y^2}\]

\[a\sqrt{(x - c)^2 + y^2} = a^2 - cx\]

两边再平方：

\[a^2[(x - c)^2 + y^2] = (a^2 - cx)^2\]

\[a^2 x^2 - 2a^2 cx + a^2 c^2 + a^2 y^2 = a^4 - 2a^2 cx + c^2 x^2\]

\[a^2 x^2 + a^2 c^2 + a^2 y^2 = a^4 + c^2 x^2\]

\[(a^2 - c^2)x^2 + a^2 y^2 = a^2(a^2 - c^2)\]

令 \(b^2 = a^2 - c^2 > 0\)（由 \(a > c\)）：

\[b^2 x^2 + a^2 y^2 = a^2 b^2\]

两边除以 \(a^2 b^2\)：

\[\frac{x^2}{a^2} + \frac{y^2}{b^2} = 1\]

最后验证：平方过程中引入了附加条件 \(a^2 - cx \geq 0\)，即
\(x \leq \frac{a^2}{c}\)。由 \(\frac{x^2}{a^2} \leq 1\) 得
\(|x| \leq a < \frac{a^2}{c}\)（因 \(c < a\)），条件自动满足。

还需验证所有满足方程的点都在椭圆上：若
\(\frac{x^2}{a^2} + \frac{y^2}{b^2} = 1\)，则 \(|x| \leq a\)。定义
\(d_1 = \sqrt{(x+c)^2 + y^2}\)，\(d_2 = \sqrt{(x-c)^2 + y^2}\)。由上述推导的逆过程可得
\(d_1 + d_2 = 2a\)（需验证 \(a^2 - cx \geq 0\)
以保证平方的等价性）。\(\blacksquare\)

\subsubsection{26.2.3
椭圆的基本性质}\label{ux692dux5706ux7684ux57faux672cux6027ux8d28}

\textbf{定理 26.2.} 椭圆
\(\frac{x^2}{a^2} + \frac{y^2}{b^2} = 1\)（\(a > b > 0\)）具有以下性质：

\begin{enumerate}
\def\labelenumi{\arabic{enumi}.}
\tightlist
\item
  \textbf{范围}：\(|x| \leq a\)，\(|y| \leq b\)。
\item
  \textbf{对称性}：关于 \(x\) 轴、\(y\) 轴和原点对称。
\item
  \textbf{顶点}：长轴端点 \(A_1(-a, 0)\)，\(A_2(a, 0)\)；短轴端点
  \(B_1(0, -b)\)，\(B_2(0, b)\)。
\item
  \textbf{焦点}：\(F_1(-c, 0)\)，\(F_2(c, 0)\)，其中
  \(c = \sqrt{a^2 - b^2}\)。
\item
  \textbf{离心率}：\(e = \frac{c}{a}\)（\(0 < e < 1\)）。
\end{enumerate}

\textbf{证明}：

\textbf{（1）范围：} 由 \(\frac{x^2}{a^2} = 1 - \frac{y^2}{b^2} \leq 1\)
得 \(|x| \leq a\)。同理 \(|y| \leq b\)。

\textbf{（2）对称性：} 将 \((x, y)\) 替换为
\((x, -y)\)：\(\frac{x^2}{a^2} + \frac{(-y)^2}{b^2} = \frac{x^2}{a^2} + \frac{y^2}{b^2} = 1\)，故关于
\(x\) 轴对称。同理替换为 \((-x, y)\) 得关于 \(y\) 轴对称，替换为
\((-x, -y)\) 得关于原点对称。

\textbf{（3）顶点：} 令 \(y = 0\) 得 \(x = \pm a\)；令 \(x = 0\) 得
\(y = \pm b\)。

\textbf{（4）} 由 \(b^2 = a^2 - c^2\) 得 \(c = \sqrt{a^2 - b^2}\)。

\textbf{（5）} \(e = \frac{c}{a}\)。由 \(0 < c < a\) 得
\(0 < e < 1\)。\(\blacksquare\)

\textbf{命题 26.2（离心率的意义）.} 离心率 \(e\)
衡量椭圆的''扁平程度''：\(e\) 越接近 \(0\)，椭圆越接近圆；\(e\) 越接近
\(1\)，椭圆越扁。

\textbf{证明}：
\(e = \frac{c}{a} = \frac{\sqrt{a^2 - b^2}}{a} = \sqrt{1 - \frac{b^2}{a^2}}\)。当
\(e \to 0\) 时 \(\frac{b}{a} \to 1\)，即 \(b \to a\)，椭圆趋近于圆
\(x^2 + y^2 = a^2\)。当 \(e \to 1\) 时 \(\frac{b}{a} \to 0\)，椭圆在
\(y\) 方向被极度压缩。\(\blacksquare\)

\subsubsection{26.2.4
焦点距离公式}\label{ux7126ux70b9ux8dddux79bbux516cux5f0f}

\textbf{定理 26.3（焦点距离公式）.} 椭圆
\(\frac{x^2}{a^2} + \frac{y^2}{b^2} = 1\)（\(a > b > 0\)）上一点
\(P(x_0, y_0)\) 到两个焦点的距离分别为

\[|PF_1| = a + ex_0, \quad |PF_2| = a - ex_0\]

\textbf{证明}： 由椭圆方程得
\(y_0^2 = b^2\left(1 - \frac{x_0^2}{a^2}\right) = \frac{b^2(a^2 - x_0^2)}{a^2}\)。

\[|PF_1|^2 = (x_0 + c)^2 + y_0^2 = (x_0 + c)^2 + \frac{b^2(a^2 - x_0^2)}{a^2}\]

\[= x_0^2 + 2cx_0 + c^2 + \frac{(a^2 - c^2)(a^2 - x_0^2)}{a^2}\]

\[= x_0^2 + 2cx_0 + c^2 + a^2 - c^2 - x_0^2 + \frac{c^2 x_0^2}{a^2}\]

\[= a^2 + 2cx_0 + \frac{c^2 x_0^2}{a^2} = \left(a + \frac{cx_0}{a}\right)^2 = (a + ex_0)^2\]

由于 \(|x_0| \leq a\) 且 \(0 < e < 1\)，有
\(a + ex_0 \geq a - ea = a(1 - e) > 0\)，故 \(|PF_1| = a + ex_0\)。由
\(|PF_1| + |PF_2| = 2a\) 得 \(|PF_2| = a - ex_0\)。\(\blacksquare\)

\subsubsection{26.2.5 椭圆的准线}\label{ux692dux5706ux7684ux51c6ux7ebf}

\textbf{定理 26.4.} 椭圆 \(\frac{x^2}{a^2} + \frac{y^2}{b^2} = 1\)
对应于焦点 \(F_2(c, 0)\) 的准线方程为
\(l: x = \frac{a^2}{c} = \frac{a}{e}\)。

\textbf{证明}： 需验证 \(\frac{|PF_2|}{d(P, l)} = e\)。由定理
26.3，\(|PF_2| = a - ex_0\)。\(P\) 到准线 \(x = \frac{a^2}{c}\) 的距离为

\[d(P, l) = \frac{a^2}{c} - x_0 = \frac{a}{e} - x_0\]

（因 \(x_0 \leq a < \frac{a^2}{c}\)，故此值为正。）

\[\frac{|PF_2|}{d(P, l)} = \frac{a - ex_0}{\frac{a}{e} - x_0} = \frac{e(a - ex_0)}{a - ex_0} = e\]

\(\blacksquare\)

\textbf{推论 26.1.} 椭圆有两条准线
\(x = \pm\frac{a^2}{c}\)，分别对应两个焦点。

\subsubsection{26.2.6
椭圆与圆的关系}\label{ux692dux5706ux4e0eux5706ux7684ux5173ux7cfb}

\textbf{命题 26.3（椭圆作为拉伸的圆）.} 将单位圆 \(x^2 + y^2 = 1\) 沿
\(x\) 轴方向拉伸 \(a\) 倍、沿 \(y\) 轴方向拉伸 \(b\) 倍，得到椭圆
\(\frac{x^2}{a^2} + \frac{y^2}{b^2} = 1\)。

\textbf{证明}： 设变换 \(T: (X, Y) \mapsto (aX, bY)\)。若 \((X, Y)\)
在单位圆上，即 \(X^2 + Y^2 = 1\)，令 \(x = aX\)，\(y = bY\)，则
\(\frac{x^2}{a^2} + \frac{y^2}{b^2} = 1\)。\(\blacksquare\)

\textbf{推论 26.2.} 椭圆的面积为 \(S = \pi ab\)。

\textbf{证明}： 由命题 26.3，椭圆
\(\frac{x^2}{a^2} + \frac{y^2}{b^2} = 1\) 由单位圆经变换
\(T: (X, Y) \mapsto (aX, bY)\) 得到。利用参数积分法：椭圆可参数化为
\(x = a\cos t\)，\(y = b\sin t\)（\(t \in [0, 2\pi]\)），其面积为

\[S = \left|\int_0^{2\pi} x\,dy\right| = \left|\int_0^{2\pi} a\cos t \cdot b\cos t\,dt\right| = ab\int_0^{2\pi} \cos^2 t\,dt = ab \cdot \pi = \pi ab\]

其中利用了 \(\int_0^{2\pi}\cos^2 t\,dt = \pi\)。\(\blacksquare\)

\textbf{注}：此结果也可通过多重积分的变量替换定理（雅可比行列式）得到：变换
\(T\) 的雅可比行列式为 \(\det(J) = ab\)，单位圆面积为
\(\pi\)，故椭圆面积为
\(\pi ab\)。但该方法需要高等数学知识，此处的参数积分法更为初等。

\subsubsection{26.2.7
椭圆的切线方程}\label{ux692dux5706ux7684ux5207ux7ebfux65b9ux7a0b}

\textbf{定理 26.5（椭圆切线方程）.} 椭圆
\(\frac{x^2}{a^2} + \frac{y^2}{b^2} = 1\) 在点
\(P_0(x_0, y_0)\)（在椭圆上）处的切线方程为

\[\frac{x_0 x}{a^2} + \frac{y_0 y}{b^2} = 1\]

\textbf{证明}： 由隐函数求导（Ch 23），对椭圆方程两边关于 \(x\) 求导：

\[\frac{2x}{a^2} + \frac{2y}{b^2}\frac{dy}{dx} = 0 \implies \frac{dy}{dx} = -\frac{b^2 x}{a^2 y}\]

在 \(P_0(x_0, y_0)\) 处，切线斜率为
\(k = -\frac{b^2 x_0}{a^2 y_0}\)（\(y_0 \neq 0\)）。切线方程为

\[y - y_0 = -\frac{b^2 x_0}{a^2 y_0}(x - x_0)\]

\[a^2 y_0(y - y_0) = -b^2 x_0(x - x_0)\]

\[b^2 x_0 x + a^2 y_0 y = b^2 x_0^2 + a^2 y_0^2\]

由 \(P_0\) 在椭圆上，\(b^2 x_0^2 + a^2 y_0^2 = a^2 b^2\)，故

\[\frac{x_0 x}{a^2} + \frac{y_0 y}{b^2} = 1\]

当 \(y_0 = 0\) 时，\(x_0 = \pm a\)，切线为 \(x = \pm a\)，即
\(\frac{x_0 x}{a^2} = 1\)，公式仍然成立。\(\blacksquare\)

\textbf{推论 26.3.} 过椭圆外一点 \(Q(x_1, y_1)\)
的切线方程可通过切点条件 \(P_0\) 在椭圆上和切线过 \(Q\) 联立求解。

\textbf{说明}：设切点为 \(P_0(x_0, y_0)\)，则切线方程为
\(\frac{x_0 x}{a^2} + \frac{y_0 y}{b^2} = 1\)。令其过 \(Q(x_1, y_1)\) 得
\(\frac{x_0 x_1}{a^2} + \frac{y_0 y_1}{b^2} = 1\)，联立
\(\frac{x_0^2}{a^2} + \frac{y_0^2}{b^2} = 1\) 解出 \((x_0, y_0)\)
即可。\(\blacksquare\)

\bigskip\hrule\bigskip

\subsection{26.3 双曲线}\label{ux53ccux66f2ux7ebf}

\subsubsection{26.3.1
双曲线的定义}\label{ux53ccux66f2ux7ebfux7684ux5b9aux4e49}

\textbf{定义 26.4（双曲线）.} 平面内到两个定点
\(F_1, F_2\)（焦点）的距离之\textbf{差}的绝对值等于常数
\(2a\)（\(0 < 2a < |F_1 F_2|\)）的点的轨迹称为\textbf{双曲线}。即

\[H = \{P \mid ||PF_1| - |PF_2|| = 2a, \; 0 < 2a < |F_1 F_2|\}\]

双曲线由\textbf{两支}组成：右支满足 \(|PF_1| - |PF_2| = 2a\)，左支满足
\(|PF_2| - |PF_1| = 2a\)。

\subsubsection{26.3.2
双曲线的标准方程}\label{ux53ccux66f2ux7ebfux7684ux6807ux51c6ux65b9ux7a0b}

\textbf{定理 26.6（双曲线的标准方程）.} 设双曲线的焦点为
\(F_1(-c, 0)\)，\(F_2(c, 0)\)（\(c > 0\)），\(0 < a < c\)。令
\(b^2 = c^2 - a^2\)（\(b > 0\)），则双曲线的标准方程为

\[\frac{x^2}{a^2} - \frac{y^2}{b^2} = 1 \quad (a > 0, b > 0)\]

\textbf{证明}： 考虑右支
\(|PF_1| - |PF_2| = 2a\)（对于左支，推导完全对称）。

\[\sqrt{(x + c)^2 + y^2} - \sqrt{(x - c)^2 + y^2} = 2a\]

移项：

\[\sqrt{(x + c)^2 + y^2} = 2a + \sqrt{(x - c)^2 + y^2}\]

两边平方：

\[(x + c)^2 + y^2 = 4a^2 + 4a\sqrt{(x - c)^2 + y^2} + (x - c)^2 + y^2\]

展开并化简：

\[x^2 + 2cx + c^2 + y^2 = 4a^2 + 4a\sqrt{(x - c)^2 + y^2} + x^2 - 2cx + c^2 + y^2\]

\[4cx - 4a^2 = 4a\sqrt{(x - c)^2 + y^2}\]

\[cx - a^2 = a\sqrt{(x - c)^2 + y^2}\]

两边再平方（注意：对右支有 \(x \geq a > 0\) 且 \(c > a\)，故
\(cx - a^2 \geq ca - a^2 > 0\)，等式两端非负）：

\[(cx - a^2)^2 = a^2[(x - c)^2 + y^2]\]

\[c^2 x^2 - 2a^2 cx + a^4 = a^2 x^2 - 2a^2 cx + a^2 c^2 + a^2 y^2\]

\[c^2 x^2 + a^4 = a^2 x^2 + a^2 c^2 + a^2 y^2\]

\[(c^2 - a^2)x^2 - a^2 y^2 = a^2(c^2 - a^2)\]

令 \(b^2 = c^2 - a^2 > 0\)（由 \(c > a\)）：

\[b^2 x^2 - a^2 y^2 = a^2 b^2\]

\[\frac{x^2}{a^2} - \frac{y^2}{b^2} = 1\]

验证附加条件：\(cx - a^2 \geq 0\)，即 \(x \geq \frac{a^2}{c}\)。由
\(\frac{x^2}{a^2} \geq 1\) 得 \(|x| \geq a\)。对于右支
\(x \geq a > \frac{a^2}{c}\)（因 \(c > a\)），条件满足。\(\blacksquare\)

\textbf{反向验证}：反之，设点 \(P(x, y)\) 满足
\(\frac{x^2}{a^2} - \frac{y^2}{b^2} = 1\)。由
\(y^2 = b^2\left(\frac{x^2}{a^2} - 1\right)\)，计算 \(|PF_1|\) 和
\(|PF_2|\)，可得
\(||PF_1| - |PF_2|| = 2a\)（详细推导与正向过程类似，此处从略）。

\subsubsection{26.3.3
双曲线的基本性质}\label{ux53ccux66f2ux7ebfux7684ux57faux672cux6027ux8d28}

\textbf{定理 26.7.} 双曲线
\(\frac{x^2}{a^2} - \frac{y^2}{b^2} = 1\)（\(a, b > 0\)）具有以下性质：

\begin{enumerate}
\def\labelenumi{\arabic{enumi}.}
\tightlist
\item
  \textbf{范围}：\(|x| \geq a\)（即 \(x \geq a\) 或
  \(x \leq -a\)），\(y \in \mathbb{R}\)。
\item
  \textbf{对称性}：关于 \(x\) 轴、\(y\) 轴和原点对称。
\item
  \textbf{顶点}：\(A_1(-a, 0)\)，\(A_2(a, 0)\)。
\item
  \textbf{焦点}：\(F_1(-c, 0)\)，\(F_2(c, 0)\)，其中
  \(c = \sqrt{a^2 + b^2}\)。
\item
  \textbf{渐近线}：\(y = \pm \frac{b}{a} x\)。
\item
  \textbf{离心率}：\(e = \frac{c}{a}\)（\(e > 1\)）。
\end{enumerate}

\textbf{证明}：

\textbf{（1）范围：} 由 \(\frac{x^2}{a^2} = 1 + \frac{y^2}{b^2} \geq 1\)
得 \(|x| \geq a\)。\(y\) 可取任意实数。

\textbf{（2）对称性：} 与椭圆的证明相同，替换 \((x,y)\) 为
\((x,-y)\)、\((-x,y)\)、\((-x,-y)\) 均不改变方程。

\textbf{（3）} 令 \(y = 0\) 得 \(x = \pm a\)。

\textbf{（4）} \(b^2 = c^2 - a^2\)，故 \(c = \sqrt{a^2 + b^2}\)。

\textbf{（5）渐近线：} 当 \((x, y)\)
沿双曲线趋于无穷远时，\(\frac{x^2}{a^2} - \frac{y^2}{b^2} = 1\) 两边除以
\(\frac{x^2}{a^2}\) 得

\[1 - \frac{a^2 y^2}{b^2 x^2} = \frac{a^2}{x^2} \to 0\]

故 \(\frac{y^2}{x^2} \to \frac{b^2}{a^2}\)，即
\(\frac{y}{x} \to \pm \frac{b}{a}\)。渐近线方程也可由
\(\frac{x^2}{a^2} - \frac{y^2}{b^2} = 0\) 得到，即
\(y = \pm\frac{b}{a}x\)。

\textbf{（6）}
\(e = \frac{c}{a} = \frac{\sqrt{a^2 + b^2}}{a} > 1\)。\(\blacksquare\)

\subsubsection{26.3.4
双曲线的焦点距离公式}\label{ux53ccux66f2ux7ebfux7684ux7126ux70b9ux8dddux79bbux516cux5f0f}

\textbf{定理 26.8（焦点距离公式）.} 双曲线
\(\frac{x^2}{a^2} - \frac{y^2}{b^2} = 1\) 上一点 \(P(x_0, y_0)\)
到两个焦点的距离分别为

\[|PF_1| = |a + ex_0|, \quad |PF_2| = |a - ex_0|\]

其中 \(e = \frac{c}{a}\)。

\textbf{证明}： 与椭圆的推导类似。由
\(y_0^2 = b^2\left(\frac{x_0^2}{a^2} - 1\right) = \frac{b^2(x_0^2 - a^2)}{a^2}\)。

\[|PF_1|^2 = (x_0 + c)^2 + y_0^2 = (x_0 + c)^2 + \frac{b^2(x_0^2 - a^2)}{a^2}\]

\[= x_0^2 + 2cx_0 + c^2 + \frac{(c^2 - a^2)(x_0^2 - a^2)}{a^2}\]

\[= x_0^2 + 2cx_0 + c^2 + \frac{c^2 x_0^2 - c^2 a^2 - a^2 x_0^2 + a^4}{a^2}\]

\[= x_0^2 + 2cx_0 + c^2 + \frac{c^2 x_0^2}{a^2} - c^2 - x_0^2 + a^2\]

\[= a^2 + 2cx_0 + \frac{c^2 x_0^2}{a^2} = \left(a + \frac{cx_0}{a}\right)^2 = (a + ex_0)^2\]

故 \(|PF_1| = |a + ex_0|\)（取绝对值是因为 \(x_0\)
可能为负）。类似地，\(|PF_2| = |a - ex_0|\)。\(\blacksquare\)

\textbf{推论 26.5.}
对于右支（\(x_0 \geq a\)），\(|PF_1| = a + ex_0\)，\(|PF_2| = ex_0 - a\)（因
\(ex_0 > ea > a\)）。对于左支（\(x_0 \leq -a\)），\(|PF_1| = -(a + ex_0)\)，\(|PF_2| = a - ex_0\)。

\textbf{证明}：对右支 \(x_0 \geq a\)，\(ex_0 \geq ea > a\)，故
\(|a + ex_0| = a + ex_0\)，\(|a - ex_0| = ex_0 - a\)。左支类似。
\(\blacksquare\)

\subsubsection{26.3.5
双曲线的准线}\label{ux53ccux66f2ux7ebfux7684ux51c6ux7ebf}

\textbf{定理 26.10} 双曲线
\(\frac{x^2}{a^2} - \frac{y^2}{b^2} = 1\)（\(a > 0, b > 0\)）对应于焦点
\(F_1(-c, 0)\) 和 \(F_2(c, 0)\) 的两条\textbf{准线}方程分别为：
\[l_1: x = -\frac{a^2}{c} = -\frac{a}{e}, \quad l_2: x = \frac{a^2}{c} = \frac{a}{e}\]

其中 \(e = \frac{c}{a} > 1\) 为离心率。

\textbf{证明}：以右焦点 \(F_2(c, 0)\) 和右准线
\(l_2: x = \frac{a^2}{c}\) 为例。对双曲线上任意点 \(P(x, y)\)，\(P\) 到
\(F_2\) 的距离为 \(|PF_2| = |ex - a|\)（由焦点距离公式），\(P\) 到
\(l_2\) 的距离为
\(|x - \frac{a^2}{c}| = \frac{1}{c}|cx - a^2| = \frac{a}{c}|ex - a|\)。故
\(\frac{|PF_2|}{d(P, l_2)} = \frac{|ex - a|}{\frac{a}{c}|ex - a|} = \frac{c}{a} = e\)。\(\blacksquare\)''

\textbf{推论 26.4}
双曲线的统一焦点-准线定义：双曲线是到焦点距离与到对应准线距离之比等于离心率
\(e > 1\)
的点的轨迹。双曲线有两条准线，分别对应两个焦点。准线不与双曲线相交（因为
\(\frac{a}{e} < a\)，准线在顶点内侧）。

\textbf{证明}：在定理26.10中已验证对双曲线上任意点
\(P\)，\(|PF_2|/d(P, l_2) = e > 1\)。此即统一定义中 \(e > 1\)
的情形。准线在顶点内侧因
\(\frac{a}{e} = \frac{a^2}{c} < a\)。\(\blacksquare\)

\subsubsection{26.3.6 等轴双曲线}\label{ux7b49ux8f74ux53ccux66f2ux7ebf}

\textbf{定义 26.5.} 当 \(a = b\) 时，双曲线
\(\frac{x^2}{a^2} - \frac{y^2}{a^2} = 1\) 即 \(x^2 - y^2 = a^2\)
称为\textbf{等轴双曲线}（或\textbf{等边双曲线}）。此时渐近线为
\(y = \pm x\)（互相垂直），离心率 \(e = \sqrt{2}\)。

\subsubsection{26.3.7
双曲线的切线方程}\label{ux53ccux66f2ux7ebfux7684ux5207ux7ebfux65b9ux7a0b}

\textbf{定理 26.9（双曲线切线方程）.} 双曲线
\(\frac{x^2}{a^2} - \frac{y^2}{b^2} = 1\) 在点
\(P_0(x_0, y_0)\)（在双曲线上）处的切线方程为

\[\frac{x_0 x}{a^2} - \frac{y_0 y}{b^2} = 1\]

\textbf{证明}： 对 \(\frac{x^2}{a^2} - \frac{y^2}{b^2} = 1\)
隐函数求导，\(\frac{2x}{a^2} - \frac{2yy'}{b^2} = 0\)，得
\(y' = \frac{b^2 x}{a^2 y}\)。在 \(P_0(x_0, y_0)\) 处斜率
\(k = \frac{b^2 x_0}{a^2 y_0}\)。

切线方程 \(y - y_0 = \frac{b^2 x_0}{a^2 y_0}(x - x_0)\)，整理得
\(a^2 y_0(y - y_0) = b^2 x_0(x - x_0)\)，即
\(b^2 x_0 x - a^2 y_0 y = b^2 x_0^2 - a^2 y_0^2\)。

利用 \(P_0\)
在双曲线上，\(\frac{x_0^2}{a^2} - \frac{y_0^2}{b^2} = 1\)，即
\(b^2 x_0^2 - a^2 y_0^2 = a^2 b^2\)，代入得

\[b^2 x_0 x - a^2 y_0 y = a^2 b^2\]

两边除以 \(a^2 b^2\)，得
\(\frac{x_0 x}{a^2} - \frac{y_0 y}{b^2} = 1\)。\(\blacksquare\)

\textbf{特殊情况：} 当 \(y_0 = 0\) 时，\(x_0 = \pm a\)，切线为
\(x = \pm a\)（双曲线顶点处的切线垂直于 \(x\) 轴），公式
\(\frac{x_0 x}{a^2} - \frac{y_0 y}{b^2} = 1\) 仍然成立。

\bigskip\hrule\bigskip

\subsection{26.4 抛物线}\label{ux629bux7269ux7ebf}

\subsubsection{26.4.1
抛物线的定义}\label{ux629bux7269ux7ebfux7684ux5b9aux4e49}

\textbf{定义 26.6（抛物线）.} 平面内到一个定点 \(F\)（焦点）和一条定直线
\(l\)（准线，\(F \notin l\)）距离相等的点的轨迹称为\textbf{抛物线}（parabola）。即

\[\mathcal{P} = \{P \mid |PF| = d(P, l)\}\]

这对应于统一定义中 \(e = 1\) 的情形。

\subsubsection{26.4.2
抛物线的标准方程}\label{ux629bux7269ux7ebfux7684ux6807ux51c6ux65b9ux7a0b}

\textbf{定理 26.11（抛物线的标准方程）.} 设焦点为
\(F\left(\frac{p}{2}, 0\right)\)（\(p > 0\)），准线为
\(x = -\frac{p}{2}\)，则抛物线的标准方程为

\[y^2 = 2px \quad (p > 0)\]

\textbf{证明}： 设 \(P(x, y)\) 为抛物线上任意一点。

\(P\) 到焦点 \(F\left(\frac{p}{2}, 0\right)\) 的距离：

\[|PF| = \sqrt{\left(x - \frac{p}{2}\right)^2 + y^2}\]

\(P\) 到准线 \(x = -\frac{p}{2}\) 的距离：

\[d(P, l) = \left|x + \frac{p}{2}\right|\]

由定义 \(|PF| = d(P, l)\)：

\[\sqrt{\left(x - \frac{p}{2}\right)^2 + y^2} = \left|x + \frac{p}{2}\right|\]

两边平方：

\[\left(x - \frac{p}{2}\right)^2 + y^2 = \left(x + \frac{p}{2}\right)^2\]

\[x^2 - px + \frac{p^2}{4} + y^2 = x^2 + px + \frac{p^2}{4}\]

\[y^2 = 2px\]

验证附加条件：\(x + \frac{p}{2} \geq 0\)，即 \(x \geq -\frac{p}{2}\)。由
\(y^2 = 2px\) 和 \(p > 0\)，当 \(y^2 \geq 0\) 时
\(x \geq 0 > -\frac{p}{2}\)，条件自动满足。\(\blacksquare\)

\textbf{命题 26.4（抛物线的四种标准方程）.}
根据焦点和准线位置的不同，抛物线有四种标准形式：

{\def\LTcaptype{none} % do not increment counter
\begin{longtable}[]{@{}
  >{\raggedright\arraybackslash}p{(\linewidth - 6\tabcolsep) * \real{0.2143}}
  >{\raggedright\arraybackslash}p{(\linewidth - 6\tabcolsep) * \real{0.2143}}
  >{\raggedright\arraybackslash}p{(\linewidth - 6\tabcolsep) * \real{0.2143}}
  >{\raggedright\arraybackslash}p{(\linewidth - 6\tabcolsep) * \real{0.3571}}@{}}
\toprule\noalign{}
\begin{minipage}[b]{\linewidth}\raggedright
方程
\end{minipage} & \begin{minipage}[b]{\linewidth}\raggedright
焦点
\end{minipage} & \begin{minipage}[b]{\linewidth}\raggedright
准线
\end{minipage} & \begin{minipage}[b]{\linewidth}\raggedright
开口方向
\end{minipage} \\
\midrule\noalign{}
\endhead
\bottomrule\noalign{}
\endlastfoot
\(y^2 = 2px\)（\(p > 0\)） & \(\left(\frac{p}{2}, 0\right)\) &
\(x = -\frac{p}{2}\) & 右 \\
\(y^2 = -2px\)（\(p > 0\)） & \(\left(-\frac{p}{2}, 0\right)\) &
\(x = \frac{p}{2}\) & 左 \\
\(x^2 = 2py\)（\(p > 0\)） & \(\left(0, \frac{p}{2}\right)\) &
\(y = -\frac{p}{2}\) & 上 \\
\(x^2 = -2py\)（\(p > 0\)） & \(\left(0, -\frac{p}{2}\right)\) &
\(y = \frac{p}{2}\) & 下 \\
\end{longtable}
}

\textbf{证明}： 以 \(y^2 = -2px\) 为例。设焦点
\(F\left(-\frac{p}{2}, 0\right)\)，准线 \(x = \frac{p}{2}\)。由
\(|PF| = d(P, l)\)：

\[\sqrt{\left(x + \frac{p}{2}\right)^2 + y^2} = \left|x - \frac{p}{2}\right|\]

平方并化简得 \(y^2 = -2px\)。其余两种类似（交换 \(x, y\)
即可）。\(\blacksquare\)

\textbf{定理 26.18a} 抛物线 \(y^2 = 2px\)（\(p > 0\)）在点
\(P_0(x_0, y_0)\)（\(y_0 \neq 0\)）处的切线方程为：
\[y_0 y = p(x + x_0)\]

当 \(y_0 = 0\) 时，\(P_0\) 为顶点，切线为 \(y\) 轴（\(x = 0\)）。

\textbf{证明}：对 \(y^2 = 2px\) 两边关于 \(x\)
求导（隐函数求导）：\(2y \cdot y' = 2p\)，故 \(y' = \frac{p}{y}\)。

在 \(P_0\) 处，斜率
\(k = \frac{p}{y_0}\)。切线方程：\(y - y_0 = \frac{p}{y_0}(x - x_0)\)。

整理：\(y_0(y - y_0) = p(x - x_0)\)，即 \(y_0 y - y_0^2 = px - px_0\)。

由 \(P_0\) 在抛物线上：\(y_0^2 = 2px_0\)，代入得
\(y_0 y - 2px_0 = px - px_0\)，即
\(y_0 y = px + px_0 = p(x + x_0)\)。\(\blacksquare\)

\subsubsection{26.4.3
抛物线的基本性质}\label{ux629bux7269ux7ebfux7684ux57faux672cux6027ux8d28}

\textbf{定理 26.12.} 抛物线 \(y^2 = 2px\)（\(p > 0\)）具有以下性质：

\begin{enumerate}
\def\labelenumi{\arabic{enumi}.}
\tightlist
\item
  \textbf{范围}：\(x \geq 0\)，\(y \in \mathbb{R}\)。
\item
  \textbf{对称性}：关于 \(x\) 轴对称（但不关于 \(y\) 轴或原点对称）。
\item
  \textbf{顶点}：原点 \((0, 0)\)。
\item
  \textbf{焦点}：\(\left(\frac{p}{2}, 0\right)\)。
\item
  \textbf{准线}：\(x = -\frac{p}{2}\)。
\item
  \textbf{离心率}：\(e = 1\)。
\end{enumerate}

\textbf{证明}：

\textbf{（1）} 由 \(y^2 = 2px\) 和 \(p > 0\)，\(y^2 \geq 0\) 推出
\(x \geq 0\)。

\textbf{（2）} 将 \((x, y)\) 替换为
\((x, -y)\)：\((-y)^2 = y^2 = 2px\)，方程不变，故关于 \(x\) 轴对称。

\textbf{（3）-（5）} 由方程直接可得。

\textbf{（6）} 抛物线是统一定义中 \(e = 1\) 的情形。\(\blacksquare\)

\subsubsection{26.4.4
抛物线的焦半径}\label{ux629bux7269ux7ebfux7684ux7126ux534aux5f84}

\textbf{定理 26.13（焦半径公式）.} 抛物线 \(y^2 = 2px\) 上一点
\(P(x_0, y_0)\) 到焦点的距离为

\[|PF| = x_0 + \frac{p}{2}\]

\textbf{证明}： \(P\) 到焦点 \(F\left(\frac{p}{2}, 0\right)\) 的距离：

\[|PF| = \sqrt{\left(x_0 - \frac{p}{2}\right)^2 + y_0^2}\]

由 \(y_0^2 = 2px_0\)：

\[|PF|^2 = x_0^2 - px_0 + \frac{p^2}{4} + 2px_0 = x_0^2 + px_0 + \frac{p^2}{4} = \left(x_0 + \frac{p}{2}\right)^2\]

由于 \(x_0 \geq 0\) 且 \(p > 0\)，\(x_0 + \frac{p}{2} > 0\)，故
\(|PF| = x_0 + \frac{p}{2}\)。\(\blacksquare\)

\textbf{推论 26.6.} 焦半径也可表示为
\(|PF| = d(P, l) = x_0 + \frac{p}{2}\)，即抛物线上一点到焦点的距离等于到准线的距离------这正是抛物线的定义。

\textbf{证明}：由抛物线定义（定义 26.6），\(|PF| = d(P,l)\)。\(P\)
到准线 \(x = -\frac{p}{2}\) 的距离为
\(\left|x_0 + \frac{p}{2}\right| = x_0 + \frac{p}{2}\)（因
\(x_0 \geq 0\)）。\(\blacksquare\)

\subsubsection{26.4.5
抛物线的焦点弦}\label{ux629bux7269ux7ebfux7684ux7126ux70b9ux5f26}

\textbf{定理 26.14（焦点弦性质）.} 过抛物线 \(y^2 = 2px\) 焦点的弦
\(AB\)（称为\textbf{焦点弦}），设 \(A(x_1, y_1)\)，\(B(x_2, y_2)\)，则：

\begin{enumerate}
\def\labelenumi{\arabic{enumi}.}
\tightlist
\item
  \(x_1 x_2 = \frac{p^2}{4}\)
\item
  \(y_1 y_2 = -p^2\)
\item
  焦点弦长 \(|AB| = x_1 + x_2 + p\)
\end{enumerate}

\textbf{证明}： 设过焦点 \(F\left(\frac{p}{2}, 0\right)\) 的直线方程为
\(x = my + \frac{p}{2}\)（\(m\)
为参数，此参数化可以涵盖所有过焦点的非水平直线）。

\textbf{特殊情况}：参数化 \(x = my + \frac{p}{2}\)
不能覆盖的是过焦点的水平线 \(y = 0\)（即对称轴本身）。该直线与抛物线
\(y^2 = 2px\) 仅交于顶点 \((0, 0)\)，不构成弦。因此
\(x = my + \frac{p}{2}\) 的参数化已涵盖所有焦点弦。

代入 \(y^2 = 2px\)：

\[y^2 = 2p\left(my + \frac{p}{2}\right) = 2pmy + p^2\]

\[y^2 - 2pmy - p^2 = 0\]

由韦达定理（Ch 8）：

\[y_1 + y_2 = 2pm, \quad y_1 y_2 = -p^2\]

由此：

\[x_1 x_2 = \frac{y_1^2}{2p} \cdot \frac{y_2^2}{2p} = \frac{(y_1 y_2)^2}{4p^2} = \frac{p^4}{4p^2} = \frac{p^2}{4}\]

焦点弦长：

\[|AB| = |AF| + |BF| = \left(x_1 + \frac{p}{2}\right) + \left(x_2 + \frac{p}{2}\right) = x_1 + x_2 + p\]

\(\blacksquare\)

\textbf{定理 26.15（焦点弦的最短弦长）.}
过抛物线焦点的弦中，垂直于对称轴的弦最短，长度为 \(2p\)（通径）。

\textbf{证明}： 由定理 26.14，\(|AB| = x_1 + x_2 + p\)。由韦达定理
\(y_1 + y_2 = 2pm\)，\(y_1 y_2 = -p^2\)，故

\[x_1 + x_2 = \frac{y_1^2 + y_2^2}{2p} = \frac{(y_1 + y_2)^2 - 2y_1 y_2}{2p} = \frac{4p^2 m^2 + 2p^2}{2p} = 2pm^2 + p\]

\[|AB| = 2pm^2 + p + p = 2p(m^2 + 1) \geq 2p\]

等号当 \(m = 0\) 时成立，此时直线 \(x = \frac{p}{2}\)
垂直于对称轴，弦长为 \(2p\)（通径）。\(\blacksquare\)

\bigskip\hrule\bigskip

\subsection{26.5
圆锥曲线的统一性质}\label{ux5706ux9525ux66f2ux7ebfux7684ux7edfux4e00ux6027ux8d28}

\subsubsection{26.5.1
离心率的统一角色}\label{ux79bbux5fc3ux7387ux7684ux7edfux4e00ux89d2ux8272}

\textbf{定理 26.16.} 设圆锥曲线的焦点为 \(F\)，准线为 \(l\)，离心率为
\(e\)。圆锥曲线上一点 \(P\) 到焦点 \(F\) 的距离 \(|PF|\) 和到准线的距离
\(d(P, l)\) 之间的关系为

\[|PF| = e \cdot d(P, l)\]

且

\begin{itemize}
\tightlist
\item
  椭圆（\(0 < e < 1\)）：曲线封闭，焦点在曲线内部。
\item
  抛物线（\(e = 1\)）：曲线开口无限延伸，但只有一个焦点。
\item
  双曲线（\(e > 1\)）：曲线由两支组成，两支各有一个焦点。
\end{itemize}

\textbf{证明}：对椭圆（\(e < 1\)）、抛物线（\(e = 1\)）、双曲线（\(e > 1\)），分别在定理26.7、26.10、26.15中已验证
\(|PF|/d(P,l) = e\)。此关系将三种曲线统一为焦点-准线定义。\(\blacksquare\)

\subsubsection{26.5.2
焦点弦长的统一公式}\label{ux7126ux70b9ux5f26ux957fux7684ux7edfux4e00ux516cux5f0f}

\textbf{定理 26.17.} 对于标准位置的圆锥曲线，过焦点且倾斜角为 \(\alpha\)
的弦长为：

\begin{itemize}
\tightlist
\item
  椭圆
  \(\frac{x^2}{a^2} + \frac{y^2}{b^2} = 1\)：\(|AB| = \frac{2b^2/a}{1 - e^2\cos^2\alpha}\)
\item
  双曲线
  \(\frac{x^2}{a^2} - \frac{y^2}{b^2} = 1\)：\(|AB| = \frac{2b^2/a}{|1 - e^2\cos^2\alpha|}\)
\item
  抛物线 \(y^2 = 2px\)：\(|AB| = \frac{2p}{\sin^2\alpha}\)
\end{itemize}

\textbf{证明}： 以椭圆为例。过右焦点 \(F_2(c, 0)\) 的直线参数方程为

\[\begin{cases} x = c + t\cos\alpha \\ y = t\sin\alpha \end{cases}\]

代入椭圆方程：

\[\frac{(c + t\cos\alpha)^2}{a^2} + \frac{t^2\sin^2\alpha}{b^2} = 1\]

两边乘以 \(a^2 b^2\) 并展开：

\[b^2(c^2 + 2ct\cos\alpha + t^2\cos^2\alpha) + a^2 t^2\sin^2\alpha = a^2 b^2\]

注意 \(b^2 c^2 = b^2(a^2 - b^2)\)（因 \(c^2 = a^2 - b^2\)），整理关于
\(t\) 的二次方程：

\[(b^2\cos^2\alpha + a^2\sin^2\alpha)t^2 + 2b^2 c\cos\alpha \cdot t - b^4 = 0\]

由韦达定理：

\[t_1 + t_2 = \frac{-2b^2 c\cos\alpha}{b^2\cos^2\alpha + a^2\sin^2\alpha}, \quad t_1 t_2 = \frac{-b^4}{b^2\cos^2\alpha + a^2\sin^2\alpha}\]

因此：

\[|AB|^2 = (t_1 - t_2)^2 = (t_1 + t_2)^2 - 4t_1 t_2 = \frac{4b^4 c^2\cos^2\alpha + 4b^4(b^2\cos^2\alpha + a^2\sin^2\alpha)}{(b^2\cos^2\alpha + a^2\sin^2\alpha)^2}\]

化简分子：\(c^2\cos^2\alpha + b^2\cos^2\alpha + a^2\sin^2\alpha = (c^2 + b^2)\cos^2\alpha + a^2\sin^2\alpha = a^2\cos^2\alpha + a^2\sin^2\alpha = a^2\)（因
\(c^2 + b^2 = a^2\)）。故

\[|AB| = \frac{2ab^2}{b^2\cos^2\alpha + a^2\sin^2\alpha}\]

再将分母表示为：\(b^2\cos^2\alpha + a^2\sin^2\alpha = a^2 - c^2\cos^2\alpha = a^2(1 - e^2\cos^2\alpha)\)，代入得

\[|AB| = \frac{2ab^2}{a^2(1 - e^2\cos^2\alpha)} = \frac{2b^2/a}{1 - e^2\cos^2\alpha}\]

\(\blacksquare\)''

\subsubsection{26.5.3
三种圆锥曲线的对比}\label{ux4e09ux79cdux5706ux9525ux66f2ux7ebfux7684ux5bf9ux6bd4}

{\def\LTcaptype{none} % do not increment counter
\begin{longtable}[]{@{}llll@{}}
\toprule\noalign{}
性质 & 椭圆 & 双曲线 & 抛物线 \\
\midrule\noalign{}
\endhead
\bottomrule\noalign{}
\endlastfoot
统一定义 & \(e < 1\) & \(e > 1\) & \(e = 1\) \\
焦点数 & 2 & 2 & 1 \\
准线数 & 2 & 2 & 1 \\
封闭性 & 封闭 & 不封闭（两支） & 不封闭 \\
对称轴 & 2 条 & 2 条 & 1 条 \\
渐近线 & 无 & 2 条 & 无 \\
二次方程 & \(B^2 - 4AC < 0\) & \(B^2 - 4AC > 0\) & \(B^2 - 4AC = 0\) \\
\end{longtable}
}

\textbf{说明：} 最后一行中，一般二次方程
\(Ax^2 + Bxy + Cy^2 + Dx + Ey + F = 0\) 的判别式 \(\Delta = B^2 - 4AC\)
决定了圆锥曲线的类型（在无退化的情况下）。

\subsubsection{26.5.4
椭圆的光学性质}\label{ux692dux5706ux7684ux5149ux5b66ux6027ux8d28}

\textbf{定理 26.18（椭圆的反射性质）.} 设 \(P\)
为椭圆上一点，\(F_1, F_2\) 为两个焦点。则 \(PF_1\) 与 \(PF_2\) 和椭圆在
\(P\) 处的切线所成的角相等。即椭圆在 \(P\) 处的切线是 \(\angle F_1PF_2\)
的外角平分线。

\textbf{证明}： 设椭圆
\(\frac{x^2}{a^2} + \frac{y^2}{b^2} = 1\)，\(P(x_0, y_0)\)
在椭圆上。切线方程为
\(\frac{x_0 x}{a^2} + \frac{y_0 y}{b^2} = 1\)，法向量为
\(\vec{n} = \left(\frac{x_0}{a^2}, \frac{y_0}{b^2}\right)\)。

\(\overrightarrow{PF_1} = (-c - x_0, -y_0)\)，\(\overrightarrow{PF_2} = (c - x_0, -y_0)\)。由定理
26.3，\(|\overrightarrow{PF_1}| = a + ex_0\)，\(|\overrightarrow{PF_2}| = a - ex_0\)。

法向量 \(\vec{n}\) 平分外角等价于
\(\frac{\overrightarrow{PF_1}}{|\overrightarrow{PF_1}|} + \frac{\overrightarrow{PF_2}}{|\overrightarrow{PF_2}|}\)
与 \(\vec{n}\) 平行。计算：

\[\frac{\overrightarrow{PF_1}}{|\overrightarrow{PF_1}|} + \frac{\overrightarrow{PF_2}}{|\overrightarrow{PF_2}|} = \frac{(-ae - x_0, -y_0)}{a + ex_0} + \frac{(ae - x_0, -y_0)}{a - ex_0}\]

\(x\)
分量：\(\frac{-(ae + x_0)}{a + ex_0} + \frac{ae - x_0}{a - ex_0} = \frac{-(ae + x_0)(a - ex_0) + (ae - x_0)(a + ex_0)}{(a + ex_0)(a - ex_0)}\)

分子
\(= -(a^2 e - ae^2 x_0 + ax_0 - ex_0^2) + (a^2 e + ae^2 x_0 - ax_0 - ex_0^2) = 2ax_0(e^2 - 1)\)

\(y\)
分量：\(\frac{-y_0}{a + ex_0} + \frac{-y_0}{a - ex_0} = \frac{-2ay_0}{a^2 - e^2 x_0^2}\)

因此
\(\frac{\overrightarrow{PF_1}}{|\overrightarrow{PF_1}|} + \frac{\overrightarrow{PF_2}}{|\overrightarrow{PF_2}|} = \left(\frac{2ax_0(e^2 - 1)}{a^2 - e^2 x_0^2},\, \frac{-2ay_0}{a^2 - e^2 x_0^2}\right)\)

验证其与 \(\vec{n}\) 的叉积为零：

\[\frac{2ax_0(e^2 - 1)}{a^2 - e^2 x_0^2} \cdot \frac{y_0}{b^2} - \frac{-2ay_0}{a^2 - e^2 x_0^2} \cdot \frac{x_0}{a^2} = \frac{2ax_0 y_0}{a^2 - e^2 x_0^2}\left(\frac{e^2 - 1}{b^2} + \frac{1}{a^2}\right)\]

利用 \(b^2 = a^2(1 - e^2)\)，得
\(\frac{e^2 - 1}{a^2(1 - e^2)} + \frac{1}{a^2} = \frac{-1}{a^2} + \frac{1}{a^2} = 0\)。

叉积为零，故法向量平分外角，即切线平分 \(\angle F_1PF_2\)
的外角。\(\blacksquare\)

\textbf{物理意义：}
椭圆的反射性质意味着从一个焦点发出的光线或声波，经椭圆面反射后会聚到另一个焦点。这一性质在医学（碎石机）、声学（''
whispering gallery''）和光学（椭圆反射镜）中有重要应用。

\subsubsection{26.5.5
抛物线的光学性质}\label{ux629bux7269ux7ebfux7684ux5149ux5b66ux6027ux8d28}

\textbf{定理 26.19（抛物线的反射性质）.} 设 \(P\) 为抛物线
\(y^2 = 2px\)（\(p > 0\)）上一点，\(F\) 为焦点。则抛物线在 \(P\)
处的切线平分 \(PF\) 与从 \(P\)
向准线所作垂线之间的夹角。等价地，入射角（\(PF\)
与切线的夹角）等于反射角（切线与水平方向的夹角）。

\textbf{证明}： 设 \(P(x_0, y_0)\)
在抛物线上（\(x_0 > 0\)，\(y_0 \neq 0\)）。焦点 \(F = (p/2, 0)\)，准线
\(x = -p/2\)。由焦点--准线性质，\(|PF| = x_0 + p/2\)。

\(\overrightarrow{PF} = (p/2 - x_0, -y_0)\)。从 \(P\)
向准线所作垂线的方向为 \((-1, 0)\)（水平向左）。

切线方向向量为 \((y_0, p)\)（由切线斜率 \(p/y_0\) 得到）。验证切线平分
\(\overrightarrow{PF}\) 与 \((-1, 0)\) 之间的夹角，即验证
\(\frac{\overrightarrow{PF}}{|\overrightarrow{PF}|} + (-1, 0)\)
与切线方向 \((y_0, p)\) 平行：

\[\frac{\overrightarrow{PF}}{|\overrightarrow{PF}|} + (-1, 0) = \left(\frac{p/2 - x_0}{x_0 + p/2} - 1,\, \frac{-y_0}{x_0 + p/2}\right) = \left(\frac{-2x_0}{x_0 + p/2},\, \frac{-y_0}{x_0 + p/2}\right)\]

该向量与 \((2x_0, y_0)\) 成比例。验证叉积：

\[2x_0 \cdot p - y_0 \cdot y_0 = 2px_0 - y_0^2 = 0\]

最后一步由 \(y_0^2 = 2px_0\) 保证。叉积为零，故切线平分夹角。

因此入射角（\(PF\)
与切线的夹角）等于反射角（切线与水平方向的夹角）。\(\blacksquare\)

\textbf{物理意义：}
平行于对称轴入射的光线，经抛物线面反射后会聚到焦点。这就是卫星天线、车灯反光镜和射电望远镜采用抛物面形状的原因。

\bigskip\hrule\bigskip

\begin{quote}
\textbf{Ch 26 本章展望}

本章从圆锥曲线的统一定义出发，系统推导了椭圆、双曲线和抛物线的标准方程，并建立了焦点距离公式、准线方程、焦点弦性质等核心理论。三种曲线在离心率
\(e\) 的参数下获得了完美的统一：\(e < 1\) 为椭圆，\(e = 1\)
为抛物线，\(e > 1\) 为双曲线。

然而，到目前为止，我们一直使用直角坐标系来描述曲线。直角坐标系虽然直观，但在处理某些问题时并不方便------例如，圆的方程
\((x-a)^2 + (y-b)^2 = r^2\) 在极坐标系中简化为
\(r = 2a\cos\theta\)；椭圆、双曲线和抛物线在以焦点为极点的极坐标系下可以统一为
\(r = \frac{ep}{1 - e\cos\theta}\)。在下一章中，我们将引入\textbf{参数方程}和\textbf{极坐标方程}两种新的曲线描述方式，它们不仅能简化计算，更能揭示圆锥曲线之间更深层的统一性。
\end{quote}

\bigskip\hrule\bigskip

\subsection{Ch 27
参数方程与极坐标方程}\label{ch-27-ux53c2ux6570ux65b9ux7a0bux4e0eux6781ux5750ux6807ux65b9ux7a0b}

\begin{quote}
\textbf{前置知识}

本章建立在以下前序章节的基础之上：

\begin{itemize}
\tightlist
\item
  \textbf{Ch
  26（圆锥曲线）}：椭圆、双曲线和抛物线的标准方程是参数方程和极坐标方程转换的基础。圆锥曲线的统一定义（焦点-准线）直接导出统一极坐标方程
  \(r = \frac{ep}{1 - e\cos\theta}\)。
\item
  \textbf{Ch
  20（三角恒等变换与反三角函数）}：参数方程大量使用三角函数------圆的参数方程使用
  \(\cos\theta, \sin\theta\)，椭圆使用''离心角''，双曲线使用
  \(\sec\theta, \tan\theta\)。极坐标与直角坐标的互化也依赖于
  \(\cos\theta\) 和 \(\sin\theta\)。
\item
  \textbf{Ch
  25（直线与圆）}：直线和圆的参数方程是理解参数方程概念的起点。
\item
  \textbf{Ch
  24（积分初步）}：极坐标下的面积公式涉及定积分，为极坐标方程的应用提供了分析工具。
\end{itemize}
\end{quote}

\bigskip\hrule\bigskip

\subsection{27.1 参数方程}\label{ux53c2ux6570ux65b9ux7a0b}

\subsubsection{27.1.1
参数方程的概念}\label{ux53c2ux6570ux65b9ux7a0bux7684ux6982ux5ff5}

\textbf{定义 27.1（参数方程）.} 在平面直角坐标系中，如果曲线 \(C\)
上的任意一点 \((x, y)\) 都可以表示为

\[\begin{cases} x = f(t) \\ y = g(t) \end{cases} \quad (t \in I)\]

其中 \(I\) 为参数 \(t\) 的取值范围，\(f\) 和 \(g\) 是关于 \(t\)
的函数，且对于 \(I\) 中的每一个 \(t\) 值，由此确定的点 \((f(t), g(t))\)
都在曲线 \(C\) 上，则称上述方程组为曲线 \(C\) 的\textbf{参数方程}，\(t\)
称为\textbf{参数}。

相应地，直接描述 \(x\) 和 \(y\) 之间关系的方程 \(F(x, y) = 0\)
称为曲线的\textbf{普通方程}（或直角坐标方程）。

\textbf{定理 27.1（消参的一般方法）.}
将参数方程化为普通方程的方法称为\textbf{消参法}，常用方法包括：

\begin{enumerate}
\def\labelenumi{\arabic{enumi}.}
\tightlist
\item
  \textbf{代入消参}：从一个方程解出参数，代入另一个方程。
\item
  \textbf{利用恒等式消参}：利用 \(\sin^2 t + \cos^2 t = 1\)
  等三角恒等式。
\item
  \textbf{加减消参}：对两个方程进行适当的代数运算。
\end{enumerate}

\textbf{说明}：消参法的本质是从两个方程 \(x = f(t), y = g(t)\)
中消去参数 \(t\)。常用方法包括：代入法（从一式解出 \(t\)
代入另一式）、三角恒等式法（利用 \(\sin^2\theta + \cos^2\theta = 1\)
等）、代数消元法。\(\blacksquare\)

\textbf{注意：}
消参时需注意参数的取值范围，消参后所得方程对应的曲线可能与原参数方程对应的曲线不完全相同（可能多出或缺少部分点）。

\subsubsection{27.1.2
直线的参数方程}\label{ux76f4ux7ebfux7684ux53c2ux6570ux65b9ux7a0b}

\textbf{定理 27.2.} 过定点 \(P_0(x_0, y_0)\)、倾斜角为
\(\alpha\)（\(0 \leq \alpha < \pi\)）的直线的参数方程为

\[\begin{cases} x = x_0 + t\cos\alpha \\ y = y_0 + t\sin\alpha \end{cases} \quad (t \in \mathbb{R})\]

\textbf{证明}： 设 \(P(x, y)\) 为直线上任意一点，则
\(\overrightarrow{P_0P} = (x - x_0, y - y_0)\)。直线的单位方向向量为
\(\vec{d} = (\cos\alpha, \sin\alpha)\)。\(P\) 在直线上的充要条件是
\(\overrightarrow{P_0P}\) 与 \(\vec{d}\) 共线，即存在
\(t \in \mathbb{R}\) 使得 \(\overrightarrow{P_0P} = t\vec{d}\)，即

\[x - x_0 = t\cos\alpha, \quad y - y_0 = t\sin\alpha\]

\(\blacksquare\)

\textbf{定理 27.3（参数 \(t\) 的几何意义）.}
在上述参数方程中，\(|t| = |\overrightarrow{P_0P}|\)，即 \(|t|\) 表示点
\(P\) 到点 \(P_0\) 的距离。当 \(t > 0\) 时，\(P\) 在 \(P_0\)
沿倾斜角方向的一侧；当 \(t < 0\) 时，\(P\) 在另一侧。

\textbf{证明}：
\(|\overrightarrow{P_0P}| = |t| \cdot |\vec{d}| = |t| \cdot 1 = |t|\)。\(\blacksquare\)

\textbf{推论 27.1.} 设直线参数方程如上，\(P_1, P_2\) 对应的参数值分别为
\(t_1, t_2\)，则

\[|P_1 P_2| = |t_1 - t_2|\]

\textbf{证明}：
\(\overrightarrow{P_1P_2} = (t_2 - t_1)(\cos\alpha, \sin\alpha)\)，故
\(|P_1P_2| = |t_2 - t_1|\)。\(\blacksquare\)

\textbf{推论 27.2.} \(P_1, P_2\) 中点对应的参数为
\(t_{M} = \frac{t_1 + t_2}{2}\)。

\textbf{证明}：中点
\(M = \frac{P_1+P_2}{2} = \frac{(P_0 + t_1\vec{d}) + (P_0 + t_2\vec{d})}{2} = P_0 + \frac{t_1+t_2}{2}\vec{d}\)，故
\(t_M = \frac{t_1+t_2}{2}\)。\(\blacksquare\)

\subsubsection{27.1.3
圆的参数方程}\label{ux5706ux7684ux53c2ux6570ux65b9ux7a0b}

\textbf{定理 27.4.} 圆心在原点、半径为 \(r\) 的圆的参数方程为

\[\begin{cases} x = r\cos\theta \\ y = r\sin\theta \end{cases} \quad (\theta \in [0, 2\pi))\]

\textbf{证明}： 设 \(P(x, y)\) 为圆 \(x^2 + y^2 = r^2\) 上的任意一点。令
\(\theta\) 为从 \(x\) 轴正方向到 \(\overrightarrow{OP}\)
的角，则由三角函数的定义，\(x = r\cos\theta\)，\(y = r\sin\theta\)。

反过来，对任意
\(\theta\)，\(x^2 + y^2 = r^2\cos^2\theta + r^2\sin^2\theta = r^2\)，点在圆上。\(\blacksquare\)

\textbf{定理 27.5.} 圆心在 \((a, b)\)、半径为 \(r\) 的圆的参数方程为

\[\begin{cases} x = a + r\cos\theta \\ y = b + r\sin\theta \end{cases} \quad (\theta \in [0, 2\pi))\]

\textbf{证明}： 平移变换：令 \(X = x - a\)，\(Y = y - b\)，圆方程变为
\(X^2 + Y^2 = r^2\)。由定理
27.4，\(X = r\cos\theta\)，\(Y = r\sin\theta\)。还原得
\(x = a + r\cos\theta\)，\(y = b + r\sin\theta\)。\(\blacksquare\)

\subsubsection{27.1.4
椭圆的参数方程}\label{ux692dux5706ux7684ux53c2ux6570ux65b9ux7a0b}

\textbf{定理 27.6.} 椭圆
\(\frac{x^2}{a^2} + \frac{y^2}{b^2} = 1\)（\(a > b > 0\)）的参数方程为

\[\begin{cases} x = a\cos\theta \\ y = b\sin\theta \end{cases} \quad (\theta \in [0, 2\pi))\]

\textbf{证明}： 将 \(x = a\cos\theta\)，\(y = b\sin\theta\)
代入椭圆方程：

\[\frac{(a\cos\theta)^2}{a^2} + \frac{(b\sin\theta)^2}{b^2} = \cos^2\theta + \sin^2\theta = 1\]

恒成立，故点 \((a\cos\theta, b\sin\theta)\) 在椭圆上。

反过来，椭圆上任一点 \((x_0, y_0)\) 满足
\(\frac{x_0^2}{a^2} + \frac{y_0^2}{b^2} = 1\)。令
\(\cos\theta = \frac{x_0}{a}\)（因 \(|x_0| \leq a\)，故
\(\left|\frac{x_0}{a}\right| \leq 1\)，\(\theta\) 存在），则
\(\sin^2\theta = 1 - \frac{x_0^2}{a^2} = \frac{y_0^2}{b^2}\)，适当选取
\(\theta\) 的象限可使 \(b\sin\theta = y_0\)。\(\blacksquare\)

\textbf{定义 27.2（离心角）.} 参数 \(\theta\) 称为椭圆上点
\(P(a\cos\theta, b\sin\theta)\) 的\textbf{离心角}（或\textbf{偏心角}）。

\textbf{注意：} 离心角 \(\theta\) \textbf{不是} \(\overrightarrow{OP}\)
与 \(x\) 轴的夹角。其几何意义是：以原点为圆心分别作半径为 \(a\) 和 \(b\)
的同心圆，从原点作与 \(x\) 轴夹角为 \(\theta\)
的射线，射线分别交大圆和小圆于 \(A\) 和 \(B\)，过 \(A\) 作 \(x\)
轴的垂线，过 \(B\) 作 \(y\) 轴的垂线，两垂线的交点即为椭圆上的点
\(P(a\cos\theta, b\sin\theta)\)。

\textbf{命题 27.1.} 椭圆的面积可用参数方程验证：\(S = \pi ab\)。

\textbf{证明}： 由参数方程，椭圆在第一象限的部分对应
\(\theta \in [0, \frac{\pi}{2}]\)。利用定积分（Ch 24）：

\[S = 4\int_0^a y\,dx = 4\int_{\pi/2}^{0} b\sin\theta \cdot (-a\sin\theta)\,d\theta = 4ab\int_0^{\pi/2} \sin^2\theta\,d\theta\]

\[= 4ab \cdot \frac{\pi}{4} = \pi ab\]

\(\blacksquare\)

\subsubsection{27.1.5
双曲线的参数方程}\label{ux53ccux66f2ux7ebfux7684ux53c2ux6570ux65b9ux7a0b}

\textbf{定理 27.7.} 双曲线 \(\frac{x^2}{a^2} - \frac{y^2}{b^2} = 1\)
的参数方程为

\[\begin{cases} x = a\sec\theta \\ y = b\tan\theta \end{cases} \quad \left(\theta \neq \frac{\pi}{2} + k\pi, \; k \in \mathbb{Z}\right)\]

\textbf{证明}： 代入双曲线方程：

\[\frac{a^2\sec^2\theta}{a^2} - \frac{b^2\tan^2\theta}{b^2} = \sec^2\theta - \tan^2\theta = 1\]

恒成立（利用 Ch 20 的三角恒等式
\(\sec^2\theta = 1 + \tan^2\theta\)）。\(\blacksquare\)

\textbf{注意：}
双曲线也可用双曲函数表示：\(x = a\cosh t\)，\(y = b\sinh t\)（\(t \in \mathbb{R}\)），利用恒等式
\(\cosh^2 t - \sinh^2 t = 1\)。

\subsubsection{27.1.6
抛物线的参数方程}\label{ux629bux7269ux7ebfux7684ux53c2ux6570ux65b9ux7a0b}

\textbf{定理 27.8.} 抛物线 \(y^2 = 2px\)（\(p > 0\)）的参数方程为

\[\begin{cases} x = 2pt^2 \\ y = 2pt \end{cases} \quad (t \in \mathbb{R})\]

\textbf{证明}：
代入：\(y^2 = (2pt)^2 = 4p^2t^2 = 2p \cdot 2pt^2 = 2px\)。\(\blacksquare\)

\bigskip\hrule\bigskip

\subsection{27.2 极坐标系}\label{ux6781ux5750ux6807ux7cfb}

\subsubsection{27.2.1
极坐标系的定义}\label{ux6781ux5750ux6807ux7cfbux7684ux5b9aux4e49}

\textbf{定义 27.3（极坐标系）.} 在平面内取一个定点
\(O\)，称为\textbf{极点}；从 \(O\) 出发引一条射线
\(Ox\)，称为\textbf{极轴}；再选定一个长度单位和角度的正方向（通常取逆时针方向为正）。这样建立的坐标系称为\textbf{极坐标系}。

对于平面内任一点 \(P\)，\(|OP| = r\) 称为点 \(P\) 的\textbf{极径}，极轴
\(Ox\) 到射线 \(OP\) 的角 \(\theta\) 称为点 \(P\)
的\textbf{极角}，有序实数对 \((r, \theta)\) 称为点 \(P\)
的\textbf{极坐标}，记为 \(P(r, \theta)\)。

\textbf{约定：} 本章中，极径 \(r \geq 0\)，极角
\(\theta \in [0, 2\pi)\)（或 \((-\pi, \pi]\)）。有时也允许
\(r < 0\)，此时 \((r, \theta)\) 表示的点与 \((|r|, \theta + \pi)\)
相同。

\subsubsection{27.2.2
极坐标与直角坐标的互化}\label{ux6781ux5750ux6807ux4e0eux76f4ux89d2ux5750ux6807ux7684ux4e92ux5316}

\textbf{定理 27.9.} 以极点 \(O\) 为直角坐标原点，极轴 \(Ox\) 为 \(x\)
轴正方向建立直角坐标系，则

\textbf{极坐标 \(\to\) 直角坐标：}

\[\begin{cases} x = r\cos\theta \\ y = r\sin\theta \end{cases}\]

\textbf{直角坐标 \(\to\) 极坐标：}

\[\begin{cases} r = \sqrt{x^2 + y^2} \\ \tan\theta = \frac{y}{x} \quad (x \neq 0) \end{cases}\]

\textbf{证明}： 由极坐标的定义，\(r = |OP|\)，\(\theta\) 为 \(OP\)
与极轴的夹角。在直角三角形中，\(x = r\cos\theta\)，\(y = r\sin\theta\)。

反过来，\(r = |OP| = \sqrt{x^2 + y^2}\)，\(\tan\theta = \frac{y}{x}\)（需根据点所在象限确定
\(\theta\) 的值）。\(\blacksquare\)

\textbf{注意：} 从直角坐标转极坐标时，\(\tan\theta = \frac{y}{x}\)
不能唯一确定 \(\theta\)（相差 \(\pi\) 的整数倍），需根据 \((x, y)\)
所在象限判断。具体地：

\begin{itemize}
\tightlist
\item
  \((x > 0, y \geq 0)\)：\(\theta = \arctan\frac{y}{x}\)
\item
  \((x < 0)\)：\(\theta = \arctan\frac{y}{x} + \pi\)
\item
  \((x > 0, y < 0)\)：\(\theta = \arctan\frac{y}{x} + 2\pi\)
\item
  \((x = 0, y > 0)\)：\(\theta = \frac{\pi}{2}\)
\item
  \((x = 0, y < 0)\)：\(\theta = \frac{3\pi}{2}\)
\end{itemize}

\bigskip\hrule\bigskip

\subsection{27.3
常见曲线的极坐标方程}\label{ux5e38ux89c1ux66f2ux7ebfux7684ux6781ux5750ux6807ux65b9ux7a0b}

\subsubsection{27.3.1
过极点的直线}\label{ux8fc7ux6781ux70b9ux7684ux76f4ux7ebf}

\textbf{定理 27.10.} 过极点且与极轴成角 \(\alpha\) 的直线的极坐标方程为

\[\theta = \alpha\]

（允许 \(r < 0\) 时，\(\theta = \alpha\) 和 \(\theta = \alpha + \pi\)
表示同一条直线。）

\textbf{证明}： 该直线上所有点的极角均为 \(\alpha\)（或
\(\alpha + \pi\)，若允许 \(r < 0\)）。\(\blacksquare\)

\subsubsection{27.3.2
圆的极坐标方程}\label{ux5706ux7684ux6781ux5750ux6807ux65b9ux7a0b}

\textbf{定理 27.11.} 以直角坐标 \((a, 0)\) 为圆心、半径为 \(a\)
的圆（过极点，圆心在极轴上）的极坐标方程为

\[r = 2a\cos\theta\]

\textbf{证明}： 圆的直角坐标方程为 \((x - a)^2 + y^2 = a^2\)，展开得
\(x^2 + y^2 = 2ax\)。代入 \(x = r\cos\theta\)，\(x^2 + y^2 = r^2\)：

\[r^2 = 2ar\cos\theta\]

当 \(r \neq 0\) 时，\(r = 2a\cos\theta\)。\(r = 0\)
对应极点，也在圆上（\(\theta = \frac{\pi}{2}\) 时
\(r = 0\)）。\(\blacksquare\)

\textbf{定理 27.12.} 更一般地：

\begin{itemize}
\tightlist
\item
  圆心在直角坐标 \((0, a)\)、半径为 \(a\) 的圆（过极点，圆心在 \(y\)
  轴上）的极坐标方程为 \(r = 2a\sin\theta\)。
\item
  圆心在极坐标 \((r_0, \theta_0)\)、半径为 \(a\) 的圆的极坐标方程为
  \(r^2 - 2r_0 r\cos(\theta - \theta_0) + r_0^2 = a^2\)。
\end{itemize}

\textbf{证明（第二种情况）：} 圆的直角坐标方程为
\(x^2 + (y - a)^2 = a^2\)，即 \(x^2 + y^2 = 2ay\)。代入极坐标
\(r^2 = 2ar\sin\theta\)，故 \(r = 2a\sin\theta\)。

\textbf{证明（第三种情况）：} 圆心的直角坐标为
\((r_0\cos\theta_0, r_0\sin\theta_0)\)，圆的直角坐标方程为

\[(x - r_0\cos\theta_0)^2 + (y - r_0\sin\theta_0)^2 = a^2\]

展开：\(x^2 + y^2 - 2r_0(x\cos\theta_0 + y\sin\theta_0) + r_0^2 = a^2\)。

代入 \(x^2 + y^2 = r^2\)，\(x = r\cos\theta\)，\(y = r\sin\theta\)：

\[r^2 - 2r_0 r(\cos\theta\cos\theta_0 + \sin\theta\sin\theta_0) + r_0^2 = a^2\]

\[r^2 - 2r_0 r\cos(\theta - \theta_0) + r_0^2 = a^2\]

\(\blacksquare\)

\subsubsection{27.3.3 心形线}\label{ux5fc3ux5f62ux7ebf}

\textbf{定义 27.4（心形线）.} 极坐标方程
\(r = a(1 + \cos\theta)\)（\(a > 0\)）表示的曲线称为\textbf{心形线}（cardioid）。

\textbf{性质：}

\begin{enumerate}
\def\labelenumi{\arabic{enumi}.}
\tightlist
\item
  \textbf{对称性}：关于极轴（\(x\) 轴）对称。
\item
  \textbf{极值}：\(\theta = 0\) 时 \(r = 2a\)（最大），\(\theta = \pi\)
  时 \(r = 0\)（过极点）。
\item
  \textbf{非负性}：\(r \geq 0\) 对所有 \(\theta\) 成立。
\end{enumerate}

\textbf{证明（对称性）：}
\(r(-\theta) = a(1 + \cos(-\theta)) = a(1 + \cos\theta) = r(\theta)\)，故关于极轴对称。\(\blacksquare\)

\subsubsection{27.3.4 玫瑰线}\label{ux73abux7470ux7ebf}

\textbf{定义 27.5（玫瑰线）.} 极坐标方程
\(r = a\cos(n\theta)\)（\(a > 0\)，\(n\)
为正整数）表示的曲线称为\textbf{玫瑰线}（rose curve）。

\textbf{性质：}

\begin{itemize}
\tightlist
\item
  当 \(n\) 为奇数时，玫瑰线有 \(n\) 个花瓣。
\item
  当 \(n\) 为偶数时，玫瑰线有 \(2n\) 个花瓣。
\end{itemize}

例如：\(r = a\cos(2\theta)\) 有四个花瓣，\(r = a\cos(3\theta)\)
有三个花瓣。

\subsubsection{27.3.5
阿基米德螺旋线}\label{ux963fux57faux7c73ux5fb7ux87baux65cbux7ebf}

\textbf{定义 27.6（阿基米德螺旋线）.} 极坐标方程
\(r = a\theta\)（\(a > 0\)，\(\theta \geq 0\)）表示的曲线称为\textbf{阿基米德螺旋线}。

每转一圈（\(\theta\) 增加 \(2\pi\)），极径增加
\(2\pi a\)，即螺旋线以恒定速率向外展开。

\bigskip\hrule\bigskip

\subsection{27.4
圆锥曲线的统一极坐标方程}\label{ux5706ux9525ux66f2ux7ebfux7684ux7edfux4e00ux6781ux5750ux6807ux65b9ux7a0b}

\subsubsection{27.4.1
统一极坐标方程的推导}\label{ux7edfux4e00ux6781ux5750ux6807ux65b9ux7a0bux7684ux63a8ux5bfc}

\textbf{定理 27.13（圆锥曲线的统一极坐标方程）.}
以焦点为极点，以焦点到准线的方向为极轴的反方向，则圆锥曲线的极坐标方程为

\[r = \frac{ep}{1 - e\cos\theta}\]

其中 \(e\) 为离心率，\(p\) 为焦点到准线的距离（焦准距）。

\begin{itemize}
\tightlist
\item
  当 \(0 < e < 1\) 时，曲线为\textbf{椭圆}；
\item
  当 \(e = 1\) 时，曲线为\textbf{抛物线}；
\item
  当 \(e > 1\) 时，曲线为\textbf{双曲线}（右支）。
\end{itemize}

\textbf{证明}： 设焦点 \(F\) 为极点，准线 \(l\) 在焦点的
\(\theta = \pi\) 方向（左侧）。设 \(P(r, \theta)\) 为曲线上的点。

由圆锥曲线的统一定义（Ch 26.1），\(|PF| = e \cdot d(P, l)\)。

在极坐标中，\(|PF| = r\)。准线 \(l\) 在焦点左侧距离 \(p\) 处。点 \(P\)
的直角坐标为 \((r\cos\theta, r\sin\theta)\)，准线 \(l\) 的直角坐标方程为
\(x = -p\)。\(P\) 到准线 \(l\) 的距离为

\[d(P, l) = |r\cos\theta - (-p)| = p + r\cos\theta\]

（需验证 \(p + r\cos\theta > 0\)，对椭圆和抛物线成立。）

由统一定义：

\[r = e(p + r\cos\theta) = ep + er\cos\theta\]

\[r - er\cos\theta = ep\]

\[r(1 - e\cos\theta) = ep\]

\[r = \frac{ep}{1 - e\cos\theta}\]

\(\blacksquare\)

\subsubsection{27.4.2
与标准方程的关系}\label{ux4e0eux6807ux51c6ux65b9ux7a0bux7684ux5173ux7cfb}

\textbf{推论 27.3.} 对于焦点在原点的椭圆
\(\frac{x^2}{a^2} + \frac{y^2}{b^2} = 1\)（\(a > b > 0\)），若取右焦点为极点，则

\[r = \frac{a(1 - e^2)}{1 - e\cos\theta} = \frac{b^2/a}{1 - e\cos\theta}\]

\textbf{证明}： 由椭圆的标准方程，右焦点为 \((c, 0)\)，准线为
\(x = \frac{a^2}{c} = \frac{a}{e}\)。焦准距

\[p = \frac{a}{e} - c = \frac{a - ce}{e} = \frac{a - a e^2}{e} = \frac{a(1 - e^2)}{e}\]

代入统一方程：

\[r = \frac{ep}{1 - e\cos\theta} = \frac{a(1 - e^2)}{1 - e\cos\theta}\]

又 \(b^2 = a^2 - c^2 = a^2(1 - e^2)\)，故
\(a(1 - e^2) = \frac{b^2}{a}\)。\(\blacksquare\)

\textbf{推论 27.4.} 对于抛物线 \(y^2 = 2px\)，取焦点为极点，则

\[r = \frac{p}{1 - \cos\theta}\]

\textbf{证明}： \(e = 1\)，焦准距 \(p\)
不变，代入统一方程即得。\(\blacksquare\)

\textbf{推论 27.5.} 对于双曲线
\(\frac{x^2}{a^2} - \frac{y^2}{b^2} = 1\)，取右焦点为极点，则

\[r = \frac{a(e^2 - 1)}{1 - e\cos\theta}\]

\textbf{证明}： 右焦点 \((c, 0)\)，准线
\(x = \frac{a^2}{c} = \frac{a}{e}\)，焦准距
\(p = c - \frac{a}{e} = \frac{ce - a}{e} = \frac{a(e^2 - 1)}{e}\)。

代入统一方程：\(r = \frac{ep}{1 - e\cos\theta} = \frac{a(e^2 - 1)}{1 - e\cos\theta}\)。\(\blacksquare\)

\subsubsection{27.4.3
焦点弦长的统一公式}\label{ux7126ux70b9ux5f26ux957fux7684ux7edfux4e00ux516cux5f0f-1}

\textbf{定理 27.14（焦点弦长公式）.} 过焦点的弦 \(AB\)（\(A\) 对应
\(\theta_1\)，\(B\) 对应 \(\theta_2 = \theta_1 + \pi\)），弦长为

\[|AB| = r_{A} + r_{B} = \frac{2ep}{1 - e^2\cos^2\theta_1}\]

\textbf{证明}： \(B\) 的极角为 \(\theta_1 + \pi\)，故
\(\cos\theta_2 = \cos(\theta_1 + \pi) = -\cos\theta_1\)。

\[r_{A} = \frac{ep}{1 - e\cos\theta_1}, \quad r_{B} = \frac{ep}{1 - e\cos(\theta_1 + \pi)} = \frac{ep}{1 + e\cos\theta_1}\]

\[|AB| = r_{A} + r_{B} = \frac{ep}{1 - e\cos\theta_1} + \frac{ep}{1 + e\cos\theta_1}\]

\[= ep \cdot \frac{(1 + e\cos\theta_1) + (1 - e\cos\theta_1)}{(1 - e\cos\theta_1)(1 + e\cos\theta_1)} = \frac{2ep}{1 - e^2\cos^2\theta_1}\]

\(\blacksquare\)

\textbf{推论 27.6.} 当
\(\theta_1 = \frac{\pi}{2}\)（即弦垂直于极轴方向）时，\(\cos\theta_1 = 0\)，\(|AB| = 2ep\)。这是\textbf{通径}（正焦弦）的长度。

\textbf{证明}：令
\(\theta_1 = \frac{\pi}{2}\)，\(\cos\theta_1 = 0\)，代入焦点弦公式
\(|AB| = \frac{2ep}{1 - e^2\cos^2\theta_1} = \frac{2ep}{1} = 2ep\)。
\(\blacksquare\)

\begin{itemize}
\tightlist
\item
  椭圆：通径 \(= 2ep = 2a(1 - e^2) = \frac{2b^2}{a}\)。
\item
  抛物线：通径 \(= 2ep = 2p\)（与定理 26.15 一致）。
\item
  双曲线：通径 \(= 2ep = 2a(e^2 - 1) = \frac{2b^2}{a}\)。
\end{itemize}

\bigskip\hrule\bigskip

\subsection{27.5
极坐标方程的应用}\label{ux6781ux5750ux6807ux65b9ux7a0bux7684ux5e94ux7528}

\subsubsection{27.5.1
极坐标下的面积公式}\label{ux6781ux5750ux6807ux4e0bux7684ux9762ux79efux516cux5f0f}

\textbf{定理 27.15（极坐标面积公式）.} 由射线
\(\theta = \alpha\)、\(\theta = \beta\)（\(\alpha < \beta\)）和曲线
\(r = f(\theta)\) 围成的扇形区域的面积为

\[S = \frac{1}{2}\int_\alpha^\beta r^2\,d\theta = \frac{1}{2}\int_\alpha^\beta [f(\theta)]^2\,d\theta\]

\textbf{证明思路：} 此公式的严格证明需要定积分的理论（Ch
32）。基本思想是将扇形分割为 \(n\) 个小扇形：将 \([\alpha, \beta]\)
等分为
\(\theta_0 = \alpha < \theta_1 < \cdots < \theta_n = \beta\)，每个小区间
\([\theta_{i-1}, \theta_i]\) 对应的小扇形面积近似为
\(\frac{1}{2}[f(\xi_i)]^2 \Delta\theta\)（半径为 \(f(\xi_i)\)、张角为
\(\Delta\theta = \frac{\beta - \alpha}{n}\) 的扇形面积）。取极限
\(n \to \infty\) 即得定积分
\(\frac{1}{2}\int_\alpha^\beta [f(\theta)]^2\,d\theta\)。\(\blacksquare\)

\textbf{注}：此公式可视为极坐标下的''面积微元法''：在角度 \(\theta\)
处取微小张角 \(d\theta\)，对应的小扇形面积为
\(dS = \frac{1}{2}r^2\,d\theta\)，累加（积分）即得总面积。

\textbf{应用 27.1（心形线围成的面积）.} 心形线 \(r = a(1 + \cos\theta)\)
围成的面积为

\[S = \frac{1}{2}\int_0^{2\pi} a^2(1 + \cos\theta)^2\,d\theta = \frac{a^2}{2}\int_0^{2\pi}(1 + 2\cos\theta + \cos^2\theta)\,d\theta\]

\[= \frac{a^2}{2}\left[2\pi + 0 + \pi\right] = \frac{3\pi a^2}{2}\]

（利用
\(\int_0^{2\pi}\cos\theta\,d\theta = 0\)，\(\int_0^{2\pi}\cos^2\theta\,d\theta = \pi\)。）

\subsubsection{27.5.2
极坐标下曲线的切线}\label{ux6781ux5750ux6807ux4e0bux66f2ux7ebfux7684ux5207ux7ebf}

\textbf{定理 27.16.} 设曲线的极坐标方程为 \(r = f(\theta)\)，则曲线在点
\((r, \theta)\) 处的切线与径向（从极点到该点的方向）之间的夹角 \(\psi\)
满足

\[\tan\psi = \frac{r}{r'}\]

其中 \(r' = \frac{dr}{d\theta}\)。

\textbf{证明}： 将极坐标方程化为参数方程
\(x = r\cos\theta\)，\(y = r\sin\theta\)（\(\theta\) 为参数）。则

\[\frac{dx}{d\theta} = r'\cos\theta - r\sin\theta, \quad \frac{dy}{d\theta} = r'\sin\theta + r\cos\theta\]

切线的方向向量为
\(\left(\frac{dx}{d\theta}, \frac{dy}{d\theta}\right)\)。径向方向向量为
\((\cos\theta, \sin\theta)\)。夹角 \(\psi\) 满足

\[\tan\psi = \frac{|(r'\sin\theta + r\cos\theta)\cos\theta - (r'\cos\theta - r\sin\theta)\sin\theta|}{|(r'\cos\theta - r\sin\theta)\cos\theta + (r'\sin\theta + r\cos\theta)\sin\theta|}\]

化简分子：\(r'\sin\theta\cos\theta + r\cos^2\theta - r'\cos\theta\sin\theta + r\sin^2\theta = r\)。

化简分母：\(r'\cos^2\theta - r\sin\theta\cos\theta + r'\sin^2\theta + r\cos\theta\sin\theta = r'\)。

故 \(\tan\psi = \frac{r}{r'}\)。\(\blacksquare\)

\bigskip\hrule\bigskip

\subsection{27.6
三种方程的对比与选择}\label{ux4e09ux79cdux65b9ux7a0bux7684ux5bf9ux6bd4ux4e0eux9009ux62e9}

\subsubsection{27.6.1
方程形式的对比}\label{ux65b9ux7a0bux5f62ux5f0fux7684ux5bf9ux6bd4}

{\def\LTcaptype{none} % do not increment counter
\begin{longtable}[]{@{}
  >{\raggedright\arraybackslash}p{(\linewidth - 6\tabcolsep) * \real{0.0984}}
  >{\raggedright\arraybackslash}p{(\linewidth - 6\tabcolsep) * \real{0.3607}}
  >{\raggedright\arraybackslash}p{(\linewidth - 6\tabcolsep) * \real{0.1475}}
  >{\raggedright\arraybackslash}p{(\linewidth - 6\tabcolsep) * \real{0.3934}}@{}}
\toprule\noalign{}
\begin{minipage}[b]{\linewidth}\raggedright
特征
\end{minipage} & \begin{minipage}[b]{\linewidth}\raggedright
普通方程 \(F(x,y) = 0\)
\end{minipage} & \begin{minipage}[b]{\linewidth}\raggedright
参数方程
\end{minipage} & \begin{minipage}[b]{\linewidth}\raggedright
极坐标方程 \(r = f(\theta)\)
\end{minipage} \\
\midrule\noalign{}
\endhead
\bottomrule\noalign{}
\endlastfoot
基本思想 & \(x, y\) 的直接关系 & 引入辅助变量 \(t\) &
用距离和角度定位 \\
优点 & 直观、通用 & 引入自由度，便于求最值和轨迹 &
利用对称性，便于求面积 \\
适用场景 & 一般代数问题 & 最值、轨迹、多变量代换 &
中心对称图形、焦点弦 \\
典型应用 & 韦达定理联立 & 圆/椭圆上的三角代换 &
心形线、玫瑰线、焦点弦 \\
\end{longtable}
}

\subsubsection{27.6.2 转换策略}\label{ux8f6cux6362ux7b56ux7565}

\begin{enumerate}
\def\labelenumi{\arabic{enumi}.}
\tightlist
\item
  遇到圆/椭圆上的点，考虑用参数方程（\(\cos\theta, \sin\theta\) 代换）。
\item
  涉及焦点弦长、焦点距离，考虑极坐标方程。
\item
  涉及过定点的弦长，考虑直线的参数方程（利用参数 \(t\) 的几何意义）。
\item
  一般情况下优先使用普通方程。
\end{enumerate}

\subsubsection{27.6.3
方程转换的数学基础}\label{ux65b9ux7a0bux8f6cux6362ux7684ux6570ux5b66ux57faux7840}

\textbf{定理 27.17.} 设曲线的参数方程为
\(x = f(t)\)，\(y = g(t)\)，则消去参数 \(t\) 可得普通方程
\(F(x, y) = 0\)。反过来，普通方程可能有多种参数化方式。

\textbf{证明}： 由 \(x = f(t)\) 得 \(t = f^{-1}(x)\)（若 \(f\)
有反函数），代入 \(y = g(t) = g(f^{-1}(x))\)，得 \(y\) 关于 \(x\)
的关系，即普通方程。

反之，给定普通方程 \(F(x, y) = 0\)，可能有多种参数化。例如，单位圆
\(x^2 + y^2 = 1\) 可以参数化为
\((\cos\theta, \sin\theta)\)，也可以参数化为
\(\left(\frac{1 - t^2}{1 + t^2}, \frac{2t}{1 + t^2}\right)\)（万能代换）。\(\blacksquare\)

\textbf{定理 27.18.} 参数方程与极坐标方程之间的转换关系为：设参数方程为
\(x = f(t)\)，\(y = g(t)\)，则极坐标方程为

\[r = \sqrt{[f(t)]^2 + [g(t)]^2}, \quad \tan\theta = \frac{g(t)}{f(t)}\]

消去 \(t\) 即得 \(r\) 关于 \(\theta\) 的关系。

\textbf{说明}：此为坐标转换的直接应用。由
\(x = r\cos\theta, y = r\sin\theta\) 和参数方程 \(x = f(t), y = g(t)\)
联立即得
\(r = \sqrt{f(t)^2 + g(t)^2}\)，\(\tan\theta = g(t)/f(t)\)。\(\blacksquare\)

\subsubsection{27.6.4
曲线系与参数}\label{ux66f2ux7ebfux7cfbux4e0eux53c2ux6570}

\textbf{定义 27.7（曲线系）.} 含参数 \(\lambda\) 的方程
\(F(x, y) + \lambda G(x, y) = 0\) 表示的曲线族称为\textbf{曲线系}。当
\(\lambda\) 取不同值时，曲线系中的每一条曲线都经过 \(F(x, y) = 0\) 和
\(G(x, y) = 0\) 的交点。

\textbf{定理 27.19（曲线系定理）.} 设曲线 \(C_1: F(x, y) = 0\) 和
\(C_2: G(x, y) = 0\) 有 \(n\) 个交点，则曲线系
\(F(x, y) + \lambda G(x, y) = 0\) 中的每一条曲线都经过这 \(n\) 个交点。

\textbf{证明}： 设 \((x_0, y_0)\) 是 \(C_1\) 和 \(C_2\) 的交点，即
\(F(x_0, y_0) = 0\) 且 \(G(x_0, y_0) = 0\)。则
\(F(x_0, y_0) + \lambda G(x_0, y_0) = 0 + \lambda \cdot 0 = 0\)，故
\((x_0, y_0)\) 在曲线系中的每一条曲线上。\(\blacksquare\)

\textbf{应用 27.2（过两圆交点的圆系）.} 设
\(\odot C_1: x^2 + y^2 + D_1 x + E_1 y + F_1 = 0\) 和
\(\odot C_2: x^2 + y^2 + D_2 x + E_2 y + F_2 = 0\) 相交于两点，则圆系

\[(x^2 + y^2 + D_1 x + E_1 y + F_1) + \lambda(x^2 + y^2 + D_2 x + E_2 y + F_2) = 0\]

经过这两点。当 \(\lambda = -1\) 时，二次项消去，得到根轴（定理 25.18）。

\subsubsection{27.6.5
极坐标下曲线的对称性判别}\label{ux6781ux5750ux6807ux4e0bux66f2ux7ebfux7684ux5bf9ux79f0ux6027ux5224ux522b}

\textbf{命题 27.2.} 设曲线的极坐标方程为 \(r = f(\theta)\)，则：

\begin{enumerate}
\def\labelenumi{\arabic{enumi}.}
\tightlist
\item
  若 \(f(-\theta) = f(\theta)\)，则曲线关于极轴对称。
\item
  若 \(f(\pi - \theta) = f(\theta)\)，则曲线关于直线
  \(\theta = \frac{\pi}{2}\) 对称。
\item
  若 \(f(\theta + \pi) = f(\theta)\)，或
  \(f(\theta + \pi) = -f(\theta)\)，则曲线关于极点对称。
\end{enumerate}

\textbf{证明}：

\textbf{（1）} 设 \(P(r, \theta)\) 在曲线上，即 \(r = f(\theta)\)。则
\(P' = (r, -\theta)\) 是 \(P\) 关于极轴的对称点。\(P'\) 对应的极径为
\(f(-\theta) = f(\theta) = r\)，故 \(P'\) 也在曲线上。

\textbf{（2）} \(P'' = (r, \pi - \theta)\) 是 \(P\) 关于
\(\theta = \frac{\pi}{2}\)
的对称点。\(f(\pi - \theta) = f(\theta) = r\)，故 \(P''\) 也在曲线上。

\textbf{（3）} \(P''' = (r, \theta + \pi)\) 是 \(P\)
关于极点的对称点（在极坐标中，\((r, \theta + \pi)\)
表示极径方向相反的点）。\(f(\theta + \pi) = f(\theta) = r\)（或
\(-f(\theta)\)，此时取 \(|r|\)），故 \(P'''\)
也在曲线上。\(\blacksquare\)

\textbf{应用 27.3.} 心形线 \(r = a(1 + \cos\theta)\) 关于极轴对称（因
\(\cos(-\theta) = \cos\theta\)）。玫瑰线 \(r = a\cos(2\theta)\)
关于极轴和 \(\theta = \frac{\pi}{2}\) 均对称。

\bigskip\hrule\bigskip

\begin{quote}
\textbf{Ch 27 本章展望}

本章引入了参数方程和极坐标方程两种描述曲线的方式，完成了对平面曲线表示方法的系统研究。参数方程通过引入辅助变量，以运动的观点描述曲线，在求最值、轨迹问题中具有独特优势。极坐标方程以极点和极轴为参考，特别适合处理具有旋转对称性的曲线和焦点相关问题。圆锥曲线的统一极坐标方程
\(r = \frac{ep}{1 - e\cos\theta}\)
将椭圆、抛物线和双曲线统一在一个公式之下，揭示了圆锥曲线最深刻的内在统一性。

至此，解析几何的三大主题------直线与圆（Ch 25）、圆锥曲线（Ch
26）、参数方程与极坐标方程（Ch
27）------构成了一个完整的体系。这些工具将为后续的离散数学与概率统计（卷六）提供几何直觉和分析方法的支撑。
\end{quote}

\bigskip\hrule\bigskip

\textbf{卷五总结}

本卷系统建立了平面解析几何的完整理论体系：

\begin{itemize}
\item
  \textbf{Ch 25}
  从坐标系出发，建立了直线和圆的解析理论。直线作为一次方程的几何实现，承载了方程论（Ch
  8）与几何（Ch 13）的桥梁功能。圆的方程将距离公式（Ch
  16）与二次方程（Ch 8）完美结合。导数（Ch
  23）的引入为切线理论提供了分析工具。
\item
  \textbf{Ch 26}
  将视野从圆拓展到圆锥曲线------椭圆、双曲线和抛物线。三种曲线在统一定义（焦点-准线-离心率）下获得完美统一。焦点距离公式、准线方程和焦点弦性质构成了圆锥曲线的核心理论。
\item
  \textbf{Ch 27}
  突破直角坐标系的局限，引入参数方程和极坐标方程。参数方程以运动的观点描述曲线，极坐标方程以极点-极轴为核心框架。圆锥曲线的统一极坐标方程
  \(r = \frac{ep}{1 - e\cos\theta}\)
  是本卷最深刻的结果，它将三种圆锥曲线纳入同一个简洁的公式。
\end{itemize}

三章之间层层递进：Ch 25 奠定基础（直线与圆），Ch 26
拓展视野（圆锥曲线），Ch 27
统一视角（参数方程与极坐标方程）。这种从特殊到一般、从具体到抽象的逻辑脉络，正是数学思维的核心特征。

\newpage

\section{卷四：函数与分析}\label{ux5377ux56dbux51fdux6570ux4e0eux5206ux6790}

\begin{quote}
本卷是全书的核心部分。我们从坐标系和函数的严格定义出发，依次构建基本初等函数、指数与对数函数、三角函数、三角恒等变换、数列、极限与连续、导数与微分、积分初步的完整理论体系。每一章的内容都以前面章节的结论为基础，形成严密的逻辑链条。
\end{quote}

\bigskip\hrule\bigskip

\subsection{Ch 16
坐标系与函数概念}\label{ch-16-ux5750ux6807ux7cfbux4e0eux51fdux6570ux6982ux5ff5}

\textbf{前置知识}：本章依赖 Ch
3（集合论，特别是笛卡尔积、关系、函数的集合论定义）以及 Ch
13.6（锐角三角函数，为后续三角函数的坐标化做铺垫）。本章为后续所有函数章节（Ch
17--Ch 24）奠定基础。

\bigskip\hrule\bigskip

\subsubsection{16.1
平面直角坐标系}\label{ux5e73ux9762ux76f4ux89d2ux5750ux6807ux7cfb}

\paragraph{16.1.1
坐标系的建立}\label{ux5750ux6807ux7cfbux7684ux5efaux7acb}

\textbf{定义 16.1}（平面直角坐标系）在平面上取两条互相垂直且有公共原点
\(O\) 的数轴，分别称为 \(x\) 轴（横轴）和 \(y\) 轴（纵轴）。\(x\)
轴取向右为正方向，\(y\)
轴取向上为正方向。这样建立的坐标系称为\textbf{平面直角坐标系}，记作
\(Oxy\)。

对于平面上的任意一点 \(P\)，过 \(P\) 分别作 \(x\) 轴和 \(y\)
轴的垂线，垂足对应的实数分别为 \(x_0\) 和 \(y_0\)，则有序实数对
\((x_0, y_0)\) 称为点 \(P\) 的\textbf{坐标}，记作 \(P(x_0, y_0)\)。

\textbf{定理
16.1}（坐标的唯一性与存在性）平面直角坐标系中，平面上的每个点有唯一的坐标，反之，每个有序实数对对应平面上唯一的点。即，平面直角坐标系建立了平面点集
\(\mathcal{P}\) 与笛卡尔积 \(\mathbb{R} \times \mathbb{R}\)
之间的一一对应（双射）。

\textbf{证明}：由 Ch 3
中笛卡尔积的定义以及垂线的唯一性（由欧氏几何公理保证），结论成立。具体地：

存在性：对于点 \(P\)，由垂线公理，过 \(P\) 作 \(x\)
轴的垂线有且仅有一条，垂足唯一确定实数 \(x_0\)；类似地确定 \(y_0\)。

唯一性：若点 \(P\) 有两组坐标 \((x_0, y_0)\) 和 \((x_1, y_1)\)，则过
\(P\) 向 \(x\) 轴作垂线，垂足既对应 \(x_0\) 又对应
\(x_1\)，由数轴上点与实数的一一对应得 \(x_0 = x_1\)，类似得
\(y_0 = y_1\)。

反过来，给定 \((x_0, y_0)\)，在 \(x\) 轴上取坐标为 \(x_0\) 的点
\(A\)，在 \(y\) 轴上取坐标为 \(y_0\) 的点 \(B\)，过 \(A\) 作 \(x\)
轴的垂线，过 \(B\) 作 \(y\)
轴的垂线，两垂线的交点即为所求点，由交点的唯一性得证。\(\blacksquare\)

\paragraph{16.1.2 象限}\label{ux8c61ux9650}

平面直角坐标系将平面（除去坐标轴）分为四个\textbf{象限}：

\begin{itemize}
\tightlist
\item
  \textbf{第一象限}：\(\{(x, y) \mid x > 0, y > 0\}\)
\item
  \textbf{第二象限}：\(\{(x, y) \mid x < 0, y > 0\}\)
\item
  \textbf{第三象限}：\(\{(x, y) \mid x < 0, y < 0\}\)
\item
  \textbf{第四象限}：\(\{(x, y) \mid x > 0, y < 0\}\)
\end{itemize}

\paragraph{16.1.3
两点间距离公式}\label{ux4e24ux70b9ux95f4ux8dddux79bbux516cux5f0f}

\textbf{定理 16.2}（两点间距离公式）设 \(P_1(x_1, y_1)\) 和
\(P_2(x_2, y_2)\) 是平面上的两点，则 \(P_1\) 与 \(P_2\) 之间的距离为：

\[d(P_1, P_2) = \sqrt{(x_2 - x_1)^2 + (y_2 - y_1)^2}\]

\textbf{证明}：过 \(P_1\) 作 \(x\) 轴的平行线，过 \(P_2\) 作 \(y\)
轴的平行线，两线交于点 \(Q(x_2, y_1)\)。

在直角三角形 \(P_1 Q P_2\) 中，由勾股定理（Ch 13.3）：

\[d(P_1, P_2)^2 = d(P_1, Q)^2 + d(Q, P_2)^2\]

由于 \(P_1(x_1, y_1)\) 和 \(Q(x_2, y_1)\)
在同一水平线上，\(d(P_1, Q) = |x_2 - x_1|\)。

由于 \(Q(x_2, y_1)\) 和 \(P_2(x_2, y_2)\)
在同一竖直线上，\(d(Q, P_2) = |y_2 - y_1|\)。

因此
\(d(P_1, P_2)^2 = (x_2 - x_1)^2 + (y_2 - y_1)^2\)，取正平方根即得。\(\blacksquare\)

\paragraph{16.1.4 中点公式}\label{ux4e2dux70b9ux516cux5f0f}

\textbf{定理 16.3}（中点公式）设 \(P_1(x_1, y_1)\) 和 \(P_2(x_2, y_2)\)
是两点，则线段 \(P_1 P_2\) 的中点 \(M\) 的坐标为：

\[M = \left(\frac{x_1 + x_2}{2}, \frac{y_1 + y_2}{2}\right)\]

\textbf{证明}：设中点为
\(M(x_0, y_0)\)。由中点的定义，\(d(P_1, M) = d(M, P_2)\)，且 \(M\)
在线段 \(P_1 P_2\) 上。

由距离公式，\(x_0 - x_1 = x_2 - x_0\)（方向相同，距离相等），故
\(x_0 = \frac{x_1 + x_2}{2}\)。类似地，\(y_0 = \frac{y_1 + y_2}{2}\)。\(\blacksquare\)

\bigskip\hrule\bigskip

\subsubsection{16.2
函数的严格定义}\label{ux51fdux6570ux7684ux4e25ux683cux5b9aux4e49}

\paragraph{16.2.1
函数作为特殊的二元关系}\label{ux51fdux6570ux4f5cux4e3aux7279ux6b8aux7684ux4e8cux5143ux5173ux7cfb}

在 Ch 3
中，我们已经学习了有序对、笛卡尔积和关系的概念。现在我们给出函数的严格定义。

\textbf{定义 16.2}（函数）设 \(A\) 和 \(B\) 是两个非空集合。一个从 \(A\)
到 \(B\) 的\textbf{函数} \(f\) 是笛卡尔积 \(A \times B\)
的一个子集，满足：

\[\forall x \in A, \exists! y \in B, (x, y) \in f\]

即，对于 \(A\) 中的每一个元素 \(x\)，在 \(B\) 中存在\textbf{唯一}的元素
\(y\) 使得 \((x, y) \in f\)。

我们记作 \(f: A \to B\)，读作''\(f\) 是从 \(A\) 到 \(B\) 的函数''。对于
\((x, y) \in f\)，我们也记作 \(y = f(x)\)。

\textbf{定义 16.3}（定义域）函数 \(f: A \to B\)
的\textbf{定义域}（domain）是集合 \(A\)，记作
\(\operatorname{dom}(f) = A\)。

\textbf{定义 16.4}（陪域）函数 \(f: A \to B\)
的\textbf{陪域}（codomain）是集合 \(B\)，记作
\(\operatorname{cod}(f) = B\)。

\textbf{定义 16.5}（值域）函数 \(f: A \to B\) 的\textbf{值域}（range）是
\(B\) 中所有被 \(f\) 映射到的元素组成的集合：

\[\operatorname{range}(f) = \{y \in B \mid \exists x \in A, f(x) = y\}\]

\textbf{注意}：值域是陪域的子集，即
\(\operatorname{range}(f) \subseteq B\)。

\paragraph{16.2.2
函数的像与原像}\label{ux51fdux6570ux7684ux50cfux4e0eux539fux50cf}

\textbf{定义 16.6}（像）设 \(f: A \to B\)，对于 \(S \subseteq A\)，\(S\)
在 \(f\) 下的\textbf{像}定义为：

\[f(S) = \{f(x) \mid x \in S\}\]

特别地，\(f(A) = \operatorname{range}(f)\)。

\textbf{定义 16.7}（原像）设 \(f: A \to B\)，对于
\(T \subseteq B\)，\(T\) 在 \(f\)
下的\textbf{原像}（或\textbf{逆像}）定义为：

\[f^{-1}(T) = \{x \in A \mid f(x) \in T\}\]

\paragraph{16.2.3
函数相等的条件}\label{ux51fdux6570ux76f8ux7b49ux7684ux6761ux4ef6}

\textbf{定义 16.8}（函数相等）两个函数 \(f: A \to B\) 和 \(g: C \to D\)
\textbf{相等}，记作 \(f = g\)，当且仅当：

\begin{enumerate}
\def\labelenumi{\arabic{enumi}.}
\tightlist
\item
  \(A = C\)（定义域相同）
\item
  \(B = D\)（陪域相同）
\item
  \(\forall x \in A, f(x) = g(x)\)（对应法则相同）
\end{enumerate}

\textbf{定理 16.4} 设 \(f, g: A \to B\) 是两个函数。则 \(f = g\)
当且仅当 \(\forall x \in A, f(x) = g(x)\)。

\textbf{证明}：必要性由定义显然成立。

充分性：假设 \(\forall x \in A, f(x) = g(x)\)。则

\[f = \{(x, f(x)) \mid x \in A\} = \{(x, g(x)) \mid x \in A\} = g\]

因此 \(f = g\)。\(\blacksquare\)

\bigskip\hrule\bigskip

\subsubsection{16.3
函数的基本性质}\label{ux51fdux6570ux7684ux57faux672cux6027ux8d28}

\paragraph{16.3.1 单调性}\label{ux5355ux8c03ux6027}

\textbf{定义 16.9}（单调递增）设 \(f: A \to \mathbb{R}\)，其中
\(A \subseteq \mathbb{R}\)。若

\[\forall x_1, x_2 \in A, x_1 < x_2 \Rightarrow f(x_1) \leq f(x_2)\]

则称 \(f\) 在 \(A\) 上是\textbf{单调递增}的（或\textbf{递增}的）。

\textbf{定义 16.10}（严格单调递增）若

\[\forall x_1, x_2 \in A, x_1 < x_2 \Rightarrow f(x_1) < f(x_2)\]

则称 \(f\) 在 \(A\)
上是\textbf{严格单调递增}的（或\textbf{严格递增}的）。

类似地可以定义\textbf{单调递减}和\textbf{严格单调递减}。

\textbf{定理 16.5}（单调函数的复合）设
\(f: A \to B\)，\(g: B \to \mathbb{R}\)。

\begin{enumerate}
\def\labelenumi{\arabic{enumi}.}
\tightlist
\item
  若 \(f\) 和 \(g\) 都是单调递增的，则 \(g \circ f\) 也是单调递增的。
\item
  若 \(f\) 是单调递增的，\(g\) 是单调递减的，则 \(g \circ f\)
  是单调递减的。
\item
  若 \(f\) 是单调递减的，\(g\) 是单调递增的，则 \(g \circ f\)
  是单调递减的。
\item
  若 \(f\) 和 \(g\) 都是单调递减的，则 \(g \circ f\) 是单调递增的。
\end{enumerate}

\textbf{证明}：我们证明结论 1，其余类似。

设 \(x_1, x_2 \in A\) 且 \(x_1 < x_2\)。因为 \(f\) 是单调递增的，所以
\(f(x_1) \leq f(x_2)\)。

因为 \(g\) 是单调递增的，所以 \(g(f(x_1)) \leq g(f(x_2))\)。

即 \((g \circ f)(x_1) \leq (g \circ f)(x_2)\)。

因此 \(g \circ f\) 是单调递增的。\(\blacksquare\)

\paragraph{16.3.2 奇偶性}\label{ux5947ux5076ux6027}

\textbf{定义 16.11}（奇函数）设 \(f: A \to \mathbb{R}\)，其中 \(A\)
关于原点对称（即 \(x \in A \Leftrightarrow -x \in A\)）。若

\[\forall x \in A, f(-x) = -f(x)\]

则称 \(f\) 为\textbf{奇函数}。

\textbf{定义 16.12}（偶函数）设 \(f: A \to \mathbb{R}\)，其中 \(A\)
关于原点对称。若

\[\forall x \in A, f(-x) = f(x)\]

则称 \(f\) 为\textbf{偶函数}。

\textbf{定理 16.6}（奇偶性的判定）设 \(f: A \to \mathbb{R}\)，其中 \(A\)
关于原点对称。

\begin{enumerate}
\def\labelenumi{\arabic{enumi}.}
\tightlist
\item
  \(f\) 是奇函数当且仅当其图像关于原点对称。
\item
  \(f\) 是偶函数当且仅当其图像关于 \(y\) 轴对称。
\end{enumerate}

\textbf{证明}：

\begin{enumerate}
\def\labelenumi{(\arabic{enumi})}
\tightlist
\item
  设 \(f\) 是奇函数，\((x, f(x))\) 是图像上的一点。则
  \((-x, f(-x)) = (-x, -f(x))\) 也在图像上。而 \((-x, -f(x))\) 正是
  \((x, f(x))\) 关于原点的对称点，因此图像关于原点对称。
\end{enumerate}

反之，设图像关于原点对称，则对于图像上任意一点
\((x, f(x))\)，其关于原点的对称点 \((-x, -f(x))\) 也在图像上，即
\(f(-x) = -f(x)\)，所以 \(f\) 是奇函数。

\begin{enumerate}
\def\labelenumi{(\arabic{enumi})}
\setcounter{enumi}{1}
\tightlist
\item
  的证明类似。\(\blacksquare\)
\end{enumerate}

\textbf{定理 16.7}（奇偶性的代数运算）设 \(f\) 和 \(g\)
都是定义在关于原点对称的集合上的函数。

\begin{enumerate}
\def\labelenumi{\arabic{enumi}.}
\tightlist
\item
  若 \(f\) 和 \(g\) 都是奇函数，则 \(f + g\) 是奇函数，\(f \cdot g\)
  是偶函数。
\item
  若 \(f\) 和 \(g\) 都是偶函数，则 \(f + g\) 和 \(f \cdot g\)
  都是偶函数。
\item
  若 \(f\) 是奇函数，\(g\) 是偶函数，则 \(f \cdot g\) 是奇函数。
\end{enumerate}

\textbf{证明}：

对于结论 1 的加法部分：

\[(f + g)(-x) = f(-x) + g(-x) = (-f(x)) + (-g(x)) = -(f(x) + g(x)) = -(f + g)(x)\]

所以 \(f + g\) 是奇函数。

对于结论 1 的乘法部分：

\[(f \cdot g)(-x) = f(-x) \cdot g(-x) = (-f(x)) \cdot (-g(x)) = f(x) \cdot g(x) = (f \cdot g)(x)\]

所以 \(f \cdot g\) 是偶函数。

对于结论 2：

\[(f + g)(-x) = f(-x) + g(-x) = f(x) + g(x) = (f + g)(x)\]

\[(f \cdot g)(-x) = f(-x) \cdot g(-x) = f(x) \cdot g(x) = (f \cdot g)(x)\]

所以 \(f + g\) 和 \(f \cdot g\) 都是偶函数。

对于结论 3：

\[(f \cdot g)(-x) = f(-x) \cdot g(-x) = (-f(x)) \cdot g(x) = -(f(x) \cdot g(x)) = -(f \cdot g)(x)\]

所以 \(f \cdot g\) 是奇函数。\(\blacksquare\)

\paragraph{16.3.3 周期性}\label{ux5468ux671fux6027}

\textbf{定义 16.13}（周期函数）设
\(f: \mathbb{R} \to \mathbb{R}\)。若存在正数 \(T\)，使得

\[\forall x \in \mathbb{R}, f(x + T) = f(x)\]

则称 \(f\) 是\textbf{周期函数}，\(T\) 是 \(f\) 的一个\textbf{周期}。

\textbf{定义 16.14}（最小正周期）若 \(T_0\) 是 \(f\)
的一个正周期，且对于 \(f\) 的任意正周期 \(T\)，都有 \(T_0 \leq T\)，则称
\(T_0\) 为 \(f\) 的\textbf{最小正周期}。

\textbf{定理 16.8} 设 \(f\) 是以 \(T\)
为周期的周期函数，则对于任意正整数 \(n\)，\(nT\) 也是 \(f\) 的周期。

\textbf{证明}：用数学归纳法（Ch 2.6）。

当 \(n = 1\) 时，\(T\) 是 \(f\) 的周期，结论成立。

假设 \(n = k\) 时结论成立，即 \(kT\) 是 \(f\) 的周期。则对于任意 \(x\)：

\[f(x + (k+1)T) = f((x + kT) + T) = f(x + kT) = f(x)\]

所以 \((k+1)T\) 也是 \(f\) 的周期。

由数学归纳法，结论成立。\(\blacksquare\)

\paragraph{16.3.4 有界性}\label{ux6709ux754cux6027}

\textbf{定义 16.15}（有界函数）设 \(f: A \to \mathbb{R}\)。若存在
\(M > 0\)，使得

\[\forall x \in A, |f(x)| \leq M\]

则称 \(f\) 在 \(A\) 上是\textbf{有界}的。

\textbf{定义 16.16}（上有界、下有界）设 \(f: A \to \mathbb{R}\)。

\begin{itemize}
\tightlist
\item
  若存在 \(M \in \mathbb{R}\)，使得
  \(\forall x \in A, f(x) \leq M\)，则称 \(f\) 在 \(A\)
  上是\textbf{上有界}的，\(M\) 是 \(f\) 的一个\textbf{上界}。
\item
  若存在 \(m \in \mathbb{R}\)，使得
  \(\forall x \in A, f(x) \geq m\)，则称 \(f\) 在 \(A\)
  上是\textbf{下有界}的，\(m\) 是 \(f\) 的一个\textbf{下界}。
\end{itemize}

\textbf{定理 16.9} \(f\) 在 \(A\) 上有界当且仅当 \(f\) 在 \(A\)
上既有上界又有下界。

\textbf{证明}：

充分性：设 \(f\) 有上界 \(M\) 和下界 \(m\)，则
\(|f(x)| \leq \max\{|M|, |m|\}\)，所以 \(f\) 有界。

必要性：设 \(f\) 有界，即存在 \(M' > 0\) 使得 \(|f(x)| \leq M'\)。则
\(M'\) 是上界，\(-M'\) 是下界。\(\blacksquare\)

\bigskip\hrule\bigskip

\subsubsection{16.4
反函数与复合函数}\label{ux53cdux51fdux6570ux4e0eux590dux5408ux51fdux6570}

\paragraph{16.4.1
单射、满射与双射}\label{ux5355ux5c04ux6ee1ux5c04ux4e0eux53ccux5c04}

\textbf{定义 16.17}（单射）设 \(f: A \to B\)。若

\[\forall x_1, x_2 \in A, f(x_1) = f(x_2) \Rightarrow x_1 = x_2\]

则称 \(f\) 是\textbf{单射}（或\textbf{一对一}的）。

\textbf{定义 16.18}（满射）设 \(f: A \to B\)。若

\[\forall y \in B, \exists x \in A, f(x) = y\]

即 \(\operatorname{range}(f) = B\)，则称 \(f\)
是\textbf{满射}（或\textbf{映上}的）。

\textbf{定义 16.19}（双射）若 \(f\) 既是单射又是满射，则称 \(f\)
是\textbf{双射}（或\textbf{一一对应}）。

\paragraph{16.4.2 反函数}\label{ux53cdux51fdux6570-1}

\textbf{定理 16.10}（反函数存在定理）设 \(f: A \to B\) 是一个函数。\(f\)
存在反函数当且仅当 \(f\) 是双射。

\textbf{证明}：

充分性：设 \(f\) 是双射。定义 \(g: B \to A\) 如下：对于任意
\(y \in B\)，由于 \(f\) 是满射，存在 \(x \in A\) 使得
\(f(x) = y\)；又由于 \(f\) 是单射，这样的 \(x\) 是唯一的。定义
\(g(y) = x\)，则 \(g\) 是良定义的函数。

对于任意 \(x \in A\)：\((g \circ f)(x) = g(f(x)) = x\)

对于任意 \(y \in B\)：\((f \circ g)(y) = f(g(y)) = y\)

因此 \(g\) 是 \(f\) 的反函数。

必要性：设 \(f\) 有反函数 \(g: B \to A\)。

单射性：若 \(f(x_1) = f(x_2)\)，则
\(x_1 = g(f(x_1)) = g(f(x_2)) = x_2\)。

满射性：对于任意 \(y \in B\)，令 \(x = g(y)\)，则
\(f(x) = f(g(y)) = y\)。\(\blacksquare\)

\textbf{定理 16.11}（反函数的唯一性）若 \(f\)
有反函数，则反函数是唯一的。

\textbf{证明}：设 \(g_1\) 和 \(g_2\) 都是 \(f\) 的反函数。则对于任意
\(y \in B\)：

\[g_1(y) = g_1(f(g_2(y))) = g_2(y)\]

所以 \(g_1 = g_2\)。\(\blacksquare\)

\textbf{定理 16.12}（反函数的性质）设 \(f: A \to B\)
是双射，\(g: B \to A\) 是其反函数。

\begin{enumerate}
\def\labelenumi{\arabic{enumi}.}
\tightlist
\item
  \(g\) 也是双射，且 \(f\) 是 \(g\) 的反函数。
\item
  \(f\) 和 \(g\) 的单调性相同。
\end{enumerate}

\textbf{证明}：

\begin{enumerate}
\def\labelenumi{(\arabic{enumi})}
\item
  由定义，\(g \circ f = \operatorname{id}_{A}\)，\(f \circ g = \operatorname{id}_{B}\)。因此
  \(f\) 和 \(g\) 互为反函数。由此易证 \(g\) 是双射。
\item
  设 \(f\) 是严格递增的。对于 \(y_1, y_2 \in B\) 且 \(y_1 < y_2\)，设
  \(x_1 = g(y_1)\)，\(x_2 = g(y_2)\)。则
  \(f(x_1) = y_1 < y_2 = f(x_2)\)。
\end{enumerate}

若 \(x_1 \geq x_2\)，则由 \(f\) 严格递增得
\(f(x_1) \geq f(x_2)\)，矛盾。所以 \(x_1 < x_2\)，即
\(g(y_1) < g(y_2)\)，\(g\) 严格递增。

对于 \(f\) 严格递减的情况类似证明。\(\blacksquare\)

\textbf{定理 16.13}（反函数的图像）设 \(f: A \to B\)
是双射，\(g: B \to A\) 是其反函数。则 \(f\) 的图像和 \(g\)
的图像关于直线 \(y = x\) 对称。

\textbf{证明}：点 \((a, b)\) 在 \(f\) 的图像上当且仅当
\(b = f(a)\)，当且仅当 \(a = g(b)\)，当且仅当 \((b, a)\) 在 \(g\)
的图像上。

而 \((a, b)\) 和 \((b, a)\) 关于直线 \(y = x\) 对称，所以 \(f\) 和 \(g\)
的图像关于 \(y = x\) 对称。\(\blacksquare\)

\paragraph{16.4.3 复合函数}\label{ux590dux5408ux51fdux6570-1}

\textbf{定义 16.20}（复合函数）设 \(f: A \to B\)，\(g: B \to C\)。定义
\(g\) 和 \(f\) 的\textbf{复合}为函数 \(g \circ f: A \to C\)：

\[(g \circ f)(x) = g(f(x)), \quad \forall x \in A\]

\textbf{注意}：复合的顺序是从右到左，即先作用 \(f\) 再作用 \(g\)。

\textbf{定理 16.14}（复合函数的定义域）设 \(f\) 的定义域为 \(A\)，\(g\)
的定义域为 \(D(g)\)。则 \(g \circ f\) 有定义当且仅当
\(\operatorname{range}(f) \subseteq D(g)\)。当 \(g \circ f\) 有定义时，

\[D(g \circ f) = \{x \in A \mid f(x) \in D(g)\}\]

\textbf{证明}：由复合函数的定义直接可得。\(\blacksquare\)

\textbf{定理 16.15}（复合函数的单调性）设
\(f: A \to B\)，\(g: B \to \mathbb{R}\)。复合函数 \(g \circ f\)
的单调性如下：

{\def\LTcaptype{none} % do not increment counter
\begin{longtable}[]{@{}ccc@{}}
\toprule\noalign{}
\(f\) 的单调性 & \(g\) 的单调性 & \(g \circ f\) 的单调性 \\
\midrule\noalign{}
\endhead
\bottomrule\noalign{}
\endlastfoot
递增 & 递增 & 递增 \\
递增 & 递减 & 递减 \\
递减 & 递增 & 递减 \\
递减 & 递减 & 递增 \\
\end{longtable}
}

\textbf{证明}：此即定理 16.5 的结论。口诀：同增异减。\(\blacksquare\)

\bigskip\hrule\bigskip

\subsubsection{16.5
函数图像变换}\label{ux51fdux6570ux56feux50cfux53d8ux6362}

\paragraph{16.5.1 平移变换}\label{ux5e73ux79fbux53d8ux6362}

\textbf{定理 16.16}（平移变换）设 \(y = f(x)\) 的图像为 \(C\)。

\begin{enumerate}
\def\labelenumi{\arabic{enumi}.}
\tightlist
\item
  \(y = f(x) + k\) 的图像可由 \(C\) 沿 \(y\) 轴向上平移 \(k\)
  个单位得到（\(k > 0\)）。
\item
  \(y = f(x) - k\) 的图像可由 \(C\) 沿 \(y\) 轴向下平移 \(k\)
  个单位得到（\(k > 0\)）。
\item
  \(y = f(x - h)\) 的图像可由 \(C\) 沿 \(x\) 轴向右平移 \(h\)
  个单位得到（\(h > 0\)）。
\item
  \(y = f(x + h)\) 的图像可由 \(C\) 沿 \(x\) 轴向左平移 \(h\)
  个单位得到（\(h > 0\)）。
\end{enumerate}

\textbf{证明}：以 (3) 为例。设点 \((a, b)\) 在 \(C\) 上，则
\(b = f(a)\)。令 \(x = a + h\)，则 \(f(x - h) = f(a) = b\)，所以
\((a + h, b)\) 在 \(y = f(x - h)\) 的图像上。而 \((a + h, b)\) 正是
\((a, b)\) 向右平移 \(h\) 个单位后的点。\(\blacksquare\)

\paragraph{16.5.2 伸缩变换}\label{ux4f38ux7f29ux53d8ux6362}

\textbf{定理 16.17}（伸缩变换）设 \(y = f(x)\) 的图像为 \(C\)。

\begin{enumerate}
\def\labelenumi{\arabic{enumi}.}
\tightlist
\item
  \(y = kf(x)\)（\(k > 0\)）的图像可由 \(C\) 上每点的纵坐标乘以 \(k\)
  得到（纵向伸缩）。
\item
  \(y = f(kx)\)（\(k > 0\)）的图像可由 \(C\) 上每点的横坐标除以 \(k\)
  得到（横向伸缩）。
\end{enumerate}

\textbf{证明}：以 (1) 为例。设点 \((a, b)\) 在 \(C\) 上，则
\(b = f(a)\)。那么 \(kf(a) = kb\)，所以 \((a, kb)\) 在 \(y = kf(x)\)
的图像上。\(\blacksquare\)

\paragraph{16.5.3 对称变换}\label{ux5bf9ux79f0ux53d8ux6362}

\textbf{定理 16.18}（对称变换）设 \(y = f(x)\) 的图像为 \(C\)。

\begin{enumerate}
\def\labelenumi{\arabic{enumi}.}
\tightlist
\item
  \(y = -f(x)\) 的图像与 \(C\) 关于 \(x\) 轴对称。
\item
  \(y = f(-x)\) 的图像与 \(C\) 关于 \(y\) 轴对称。
\item
  \(y = -f(-x)\) 的图像与 \(C\) 关于原点对称。
\end{enumerate}

\textbf{证明}：

\begin{enumerate}
\def\labelenumi{(\arabic{enumi})}
\item
  点 \((a, b)\) 在 \(C\) 上当且仅当 \(b = f(a)\)，当且仅当
  \(-b = -f(a)\)，当且仅当 \((a, -b)\) 在 \(y = -f(x)\) 的图像上。而
  \((a, -b)\) 是 \((a, b)\) 关于 \(x\) 轴的对称点。
\item
  点 \((a, b)\) 在 \(C\) 上当且仅当 \(b = f(a)\)，当且仅当
  \(b = f(-(-a))\)，当且仅当 \((-a, b)\) 在 \(y = f(-x)\) 的图像上。而
  \((-a, b)\) 是 \((a, b)\) 关于 \(y\) 轴的对称点。
\item
  点 \((a, b)\) 在 \(C\) 上当且仅当 \(b = f(a)\)，当且仅当
  \(-b = -f(a) = -f(-(-a))\)，当且仅当 \((-a, -b)\) 在 \(y = -f(-x)\)
  的图像上。而 \((-a, -b)\) 是 \((a, b)\)
  关于原点的对称点。\(\blacksquare\)
\end{enumerate}

\bigskip\hrule\bigskip

\subsubsection{16.6
参数方程与极坐标}\label{ux53c2ux6570ux65b9ux7a0bux4e0eux6781ux5750ux6807}

\paragraph{16.6.1 参数方程}\label{ux53c2ux6570ux65b9ux7a0b-1}

\textbf{定义 16.21}（参数方程）设
\(x = \varphi(t)\)，\(y = \psi(t)\)（\(t \in T\)，\(T\)
为某区间），则此方程组称为曲线的\textbf{参数方程}，\(t\)
称为\textbf{参数}。

\textbf{定理 16.19}（消参）若 \(\varphi\) 有反函数
\(t = \varphi^{-1}(x)\)，则 \(y = \psi(\varphi^{-1}(x))\)。

\textbf{证明}：由 \(x = \varphi(t)\) 得 \(t = \varphi^{-1}(x)\)，代入
\(y = \psi(t)\) 即得。\(\blacksquare\)

\textbf{常见参数方程}：

\begin{itemize}
\tightlist
\item
  圆：\(x = r\cos t\)，\(y = r\sin t\)（\(t \in [0, 2\pi)\)）
\item
  椭圆：\(x = a\cos t\)，\(y = b\sin t\)（\(t \in [0, 2\pi)\)）
\item
  直线：\(x = x_0 + at\)，\(y = y_0 + bt\)（\(t \in \mathbb{R}\)）
\end{itemize}

\paragraph{16.6.2 极坐标}\label{ux6781ux5750ux6807}

\textbf{定义 16.22}（极坐标系）在平面上取一点
\(O\)（称为\textbf{极点}）和一条射线
\(Ox\)（称为\textbf{极轴}）。对于平面上的点 \(P\)，设
\(|OP| = r \geq 0\)（\textbf{极径}），\(\angle xOP = \theta\)（\textbf{极角}），则有序对
\((r, \theta)\) 称为 \(P\) 的\textbf{极坐标}。

\textbf{定理 16.20}（极坐标与直角坐标的转换）

\[x = r\cos\theta, \quad y = r\sin\theta\]

\[r = \sqrt{x^2 + y^2}, \quad \tan\theta = \frac{y}{x}\]

\textbf{证明}：在以极点为原点、极轴为 \(x\)
轴正方向的直角坐标系中，由三角函数的定义即得。\(\blacksquare\)

\textbf{注}：极坐标表示不唯一，\((r, \theta)\) 和
\((r, \theta + 2k\pi)\)（\(k \in \mathbb{Z}\)）表示同一点。

\bigskip\hrule\bigskip

\subsubsection{16.7
函数的复合分解与函数方程初步}\label{ux51fdux6570ux7684ux590dux5408ux5206ux89e3ux4e0eux51fdux6570ux65b9ux7a0bux521dux6b65}

\paragraph{16.7.1 函数的分解}\label{ux51fdux6570ux7684ux5206ux89e3}

\textbf{问题}：给定一个复杂函数
\(h(x)\)，能否将其表示为两个较简单函数的复合 \(h = g \circ f\)？

\textbf{定理 16.21} 设 \(h: A \to C\) 是一个函数。若存在集合 \(B\)
和函数 \(f: A \to B\)，\(g: B \to C\)，使得 \(h = g \circ f\)，则称
\(h\) 可以\textbf{分解}为 \(g\) 和 \(f\) 的复合。

\textbf{证明}：分解总是可能的。取
\(B = \operatorname{range}(h)\)，\(f = h\)（视为 \(A \to B\)
的函数），\(g = \operatorname{id}_{B}\)（恒等映射），则
\(h = g \circ f\)。\(\blacksquare\)

\textbf{注}：函数分解在数学和工程中有重要应用。例如，在信号处理中，复杂信号可以分解为简单信号的复合；在微积分中，链式法则正是基于函数复合结构的求导方法。

\paragraph{16.7.2
函数方程简介}\label{ux51fdux6570ux65b9ux7a0bux7b80ux4ecb}

\textbf{定义
16.23}（函数方程）含有未知函数的方程称为\textbf{函数方程}。求解函数方程是指找出满足方程的所有函数。

\textbf{典型例子}：

\begin{enumerate}
\def\labelenumi{(\arabic{enumi})}
\tightlist
\item
  \textbf{柯西函数方程}：\(f(x + y) = f(x) + f(y)\)。
\end{enumerate}

\textbf{定理 16.22} 若 \(f: \mathbb{R} \to \mathbb{R}\) 满足柯西方程
\(f(x+y) = f(x) + f(y)\)，且 \(f\) 连续（或单调，或在某点有界），则
\(f(x) = kx\)（\(k\) 为常数）。

\textbf{证明}：

第一步：\(f(0) = f(0 + 0) = f(0) + f(0)\)，故 \(f(0) = 0\)。

第二步：\(f(n) = nf(1)\)（对正整数 \(n\)，由归纳法）。

第三步：\(f(-n) + f(n) = f(0) = 0\)，故
\(f(-n) = -nf(1)\)。因此对所有整数 \(m\)，\(f(m) = mf(1)\)。

第四步：\(f(1) = f\left(\frac{1}{n} + \cdots + \frac{1}{n}\right) = nf\left(\frac{1}{n}\right)\)，故
\(f\left(\frac{1}{n}\right) = \frac{f(1)}{n}\)。

第五步：对有理数
\(\frac{m}{n}\)，\(f\left(\frac{m}{n}\right) = mf\left(\frac{1}{n}\right) = \frac{m}{n}f(1)\)。

第六步：由 \(f\) 的连续性，对无理数 \(x\)，取有理数列 \(r_n \to x\)，则
\(f(x) = \lim f(r_n) = \lim r_n f(1) = xf(1)\)。

设 \(k = f(1)\)，则 \(f(x) = kx\)。\(\blacksquare\)

\begin{enumerate}
\def\labelenumi{(\arabic{enumi})}
\setcounter{enumi}{1}
\tightlist
\item
  \textbf{指数型函数方程}：\(f(x + y) = f(x)f(y)\)，\(f\)
  连续，\(f \not\equiv 0\)。解为 \(f(x) = a^x\)（\(a > 0\)）。
\end{enumerate}

\bigskip\hrule\bigskip

\textbf{本章展望}：本章建立了坐标系和函数的严格理论基础。在下一章（Ch
17）中，我们将利用函数的概念，系统研究一次函数、反比例函数、二次函数和幂函数等基本初等函数的具体性质。函数的单调性、奇偶性等一般性质将在这些具体函数上得到体现和应用。

\bigskip\hrule\bigskip

\subsection{Ch 17
基本初等函数}\label{ch-17-ux57faux672cux521dux7b49ux51fdux6570}

\begin{quote}
\textbf{注（阅读提示）}：本章中''连续''一词在直观意义上使用------图像不中断、不断开的函数。形式化的
\(\epsilon\)-\(\delta\)
定义在第22章中给出，读者可在学完第22章后回顾本章中关于连续性的讨论。初等函数（多项式、有理函数、指数函数、对数函数、三角函数）在其定义域上连续的严格证明同样依赖极限的
\(\epsilon\)-\(\delta\) 定义。
\end{quote}

\textbf{前置知识}：本章依赖 Ch
16（坐标系与函数概念，特别是函数的单调性、奇偶性、有界性等基本性质，以及反函数理论）。本章为
Ch 18（指数函数与对数函数）和后续章节提供最基础的函数模型。

\bigskip\hrule\bigskip

\subsubsection{17.1 一次函数}\label{ux4e00ux6b21ux51fdux6570}

\paragraph{17.1.1 正比例函数}\label{ux6b63ux6bd4ux4f8bux51fdux6570}

\textbf{定义 17.1}（正比例函数）形如
\(y = kx\)（\(k \neq 0\)）的函数称为\textbf{正比例函数}，定义域为
\(\mathbb{R}\)，\(k\) 称为\textbf{比例系数}。

\textbf{定理 17.1} 正比例函数
\(y = kx\)（\(k \neq 0\)）的图像是一条过原点的直线。

\textbf{证明}：设 \(P_1(x_1, kx_1)\) 和 \(P_2(x_2, kx_2)\)
是图像上任意两点（\(x_1 \neq x_2\)）。则斜率

\[\frac{kx_2 - kx_1}{x_2 - x_1} = \frac{k(x_2 - x_1)}{x_2 - x_1} = k\]

斜率为常数，且过原点 \((0, 0)\)，故图像为直线。\(\blacksquare\)

\textbf{定理 17.2} 正比例函数 \(y = kx\)（\(k \neq 0\)）是奇函数。

\textbf{证明}：\(f(-x) = k(-x) = -kx = -f(x)\)。\(\blacksquare\)

\paragraph{17.1.2
一次函数的定义与性质}\label{ux4e00ux6b21ux51fdux6570ux7684ux5b9aux4e49ux4e0eux6027ux8d28}

\textbf{定义 17.2}（一次函数）形如

\[y = kx + b \quad (k \neq 0)\]

的函数称为\textbf{一次函数}，定义域为 \(\mathbb{R}\)。当 \(b = 0\)
时，退化为正比例函数。

\textbf{定义 17.3}（斜率与截距）： - \(k\) 称为直线 \(y = kx + b\)
的\textbf{斜率}，它描述了直线的倾斜程度和方向。\(|k|\) 越大，直线越陡。
- \(b\) 称为直线在 \(y\) 轴上的\textbf{截距}，即直线与 \(y\) 轴交点
\((0, b)\) 的纵坐标。

\textbf{几何意义}：斜率 \(k\) 表示直线上每向右移动 \(1\) 个单位，\(y\)
值的变化量。严格地说，若直线上两点 \(P_1(x_1, y_1)\),
\(P_2(x_2, y_2)\)（\(x_1 \neq x_2\)），则：

\[k = \frac{y_2 - y_1}{x_2 - x_1}\]

\textbf{定理 17.3}（一次函数的图像）一次函数
\(y = kx + b\)（\(k \neq 0\)）的图像是一条直线，斜率为 \(k\)，\(y\)
轴截距为 \(b\)。

\textbf{证明}：设 \(P_1(x_1, kx_1 + b)\) 和
\(P_2(x_2, kx_2 + b)\)（\(x_1 \neq x_2\)）。则斜率

\[\frac{(kx_2 + b) - (kx_1 + b)}{x_2 - x_1} = \frac{k(x_2 - x_1)}{x_2 - x_1} = k\]

为常数。由解析几何的基本定理，斜率为常数的曲线是直线。\(\blacksquare\)

\textbf{定理 17.4}（一次函数的单调性）一次函数
\(y = kx + b\)（\(k \neq 0\)）在 \(\mathbb{R}\) 上严格单调： - 当
\(k > 0\) 时，严格递增； - 当 \(k < 0\) 时，严格递减。

\textbf{证明}：设 \(x_1 < x_2\)。则

\[f(x_2) - f(x_1) = (kx_2 + b) - (kx_1 + b) = k(x_2 - x_1)\]

当 \(k > 0\) 时，\(k(x_2 - x_1) > 0\)，故
\(f(x_2) > f(x_1)\)，函数严格递增。

当 \(k < 0\) 时，\(k(x_2 - x_1) < 0\)，故
\(f(x_2) < f(x_1)\)，函数严格递减。\(\blacksquare\)

\textbf{定理 17.5}（一次函数的奇偶性）一次函数
\(y = kx + b\)（\(k \neq 0\)）： - 当 \(b = 0\) 时，为奇函数； - 当
\(b \neq 0\) 时，既非奇函数也非偶函数。

\textbf{证明}：\(f(-x) = -kx + b\)。若 \(f\) 是奇函数，需
\(f(-x) = -f(x) = -kx - b\)，即 \(-kx + b = -kx - b\)，得 \(b = 0\)。若
\(f\) 是偶函数，需 \(f(-x) = f(x)\)，即 \(-kx + b = kx + b\)，得
\(k = 0\)，与 \(k \neq 0\) 矛盾。\(\blacksquare\)

\paragraph{17.1.3
两点确定一次函数}\label{ux4e24ux70b9ux786eux5b9aux4e00ux6b21ux51fdux6570}

\textbf{定理 17.6} 过平面上两个不同点 \(P_1(x_1, y_1)\) 和
\(P_2(x_2, y_2)\)（\(x_1 \neq x_2\)），存在唯一的一次函数
\(y = kx + b\)，其中：

\[k = \frac{y_2 - y_1}{x_2 - x_1}, \quad b = y_1 - kx_1\]

\textbf{证明}：设 \(y = kx + b\)，代入两点：

\[\begin{cases} y_1 = kx_1 + b \\ y_2 = kx_2 + b \end{cases}\]

两式相减：\(y_2 - y_1 = k(x_2 - x_1)\)，故
\(k = \frac{y_2 - y_1}{x_2 - x_1}\)。代入第一式得
\(b = y_1 - kx_1\)。由线性方程组的唯一解性，\(k\) 和 \(b\)
唯一确定。\(\blacksquare\)

\bigskip\hrule\bigskip

\subsubsection{17.2 反比例函数}\label{ux53cdux6bd4ux4f8bux51fdux6570}

\paragraph{17.2.1
定义与基本性质}\label{ux5b9aux4e49ux4e0eux57faux672cux6027ux8d28}

\textbf{定义 17.4}（反比例函数）形如

\[y = \frac{k}{x} \quad (k \neq 0)\]

的函数称为\textbf{反比例函数}，其中 \(k\)
称为\textbf{比例系数}。定义域为
\(\{x \in \mathbb{R} \mid x \neq 0\} = \mathbb{R} \setminus \{0\}\)。

\textbf{定理 17.7}（反比例函数的奇偶性）反比例函数
\(y = \frac{k}{x}\)（\(k \neq 0\)）是\textbf{奇函数}。

\textbf{证明}：定义域 \(\mathbb{R} \setminus \{0\}\)
关于原点对称。对于任意 \(x \neq 0\)：

\[f(-x) = \frac{k}{-x} = -\frac{k}{x} = -f(x)\]

因此 \(f\) 是奇函数。\(\blacksquare\)

\textbf{定理 17.8}（反比例函数的单调性）反比例函数
\(y = \frac{k}{x}\)（\(k \neq 0\)）：

\begin{enumerate}
\def\labelenumi{\arabic{enumi}.}
\tightlist
\item
  当 \(k > 0\) 时，在 \((-\infty, 0)\) 和 \((0, +\infty)\)
  上分别严格递减。
\item
  当 \(k < 0\) 时，在 \((-\infty, 0)\) 和 \((0, +\infty)\)
  上分别严格递增。
\end{enumerate}

\textbf{注意}：反比例函数在其定义域上\textbf{不}具有整体单调性，而是在各连通分支上分别单调。

\textbf{证明}：设 \(0 < x_1 < x_2\)，\(k > 0\)。则

\[f(x_2) - f(x_1) = \frac{k}{x_2} - \frac{k}{x_1} = \frac{k(x_1 - x_2)}{x_1 x_2}\]

由于 \(x_1 - x_2 < 0\)，\(x_1 x_2 > 0\)，\(k > 0\)，故
\(f(x_2) - f(x_1) < 0\)，即 \(f\) 在 \((0, +\infty)\) 上严格递减。

类似地可证 \(f\) 在 \((-\infty, 0)\) 上严格递减。

\(k < 0\) 的情形类似证明。\(\blacksquare\)

\paragraph{17.2.2 渐近线}\label{ux6e10ux8fd1ux7ebf}

\textbf{定义 17.5}（渐近线）设 \(f: D \to \mathbb{R}\)。若存在常数
\(L\)，使得

\[\lim_{x \to +\infty} f(x) = L \quad \text{或} \quad \lim_{x \to -\infty} f(x) = L\]

则称直线 \(y = L\) 为 \(f\) 的\textbf{水平渐近线}。若

\[\lim_{x \to a^+} f(x) = \pm\infty \quad \text{或} \quad \lim_{x \to a^-} f(x) = \pm\infty\]

则称直线 \(x = a\) 为 \(f\) 的\textbf{垂直渐近线}。

\textbf{定理
17.9}（反比例函数的渐近线）\(y = \frac{k}{x}\)（\(k \neq 0\)）的图像以
\(x\) 轴和 \(y\) 轴为渐近线。即：

\[\lim_{x \to \pm\infty} \frac{k}{x} = 0, \quad \lim_{x \to 0^{\pm}} \left|\frac{k}{x}\right| = +\infty\]

\textbf{证明}：对于任意 \(\varepsilon > 0\)，取
\(M = \frac{|k|}{\varepsilon}\)。当 \(|x| > M\) 时，

\[\left|\frac{k}{x}\right| = \frac{|k|}{|x|} < \frac{|k|}{M} = \varepsilon\]

因此 \(\lim_{x \to \pm\infty} \frac{k}{x} = 0\)，即 \(y = 0\)（\(x\)
轴）是水平渐近线。

对于垂直渐近线：当 \(x \to 0^+\)
时，\(\frac{k}{x} \to +\infty\)（\(k > 0\)）或
\(\frac{k}{x} \to -\infty\)（\(k < 0\)）。因此 \(x = 0\)（\(y\)
轴）是垂直渐近线。\(\blacksquare\)

\paragraph{17.2.3 面积性质}\label{ux9762ux79efux6027ux8d28}

\textbf{定理 17.10}（矩形面积性质）设 \(P(x, y)\) 是反比例函数
\(y = \frac{k}{x}\)（\(k \neq 0\)）图像上的任意一点，过 \(P\) 分别作
\(x\) 轴和 \(y\) 轴的垂线，垂足分别为 \(A\) 和 \(B\)，则矩形 \(OAPB\)
的面积为 \(|k|\)。

\textbf{证明}：\(OAPB\) 的面积
\(= |x| \cdot |y| = |x| \cdot \left|\frac{k}{x}\right| = |k|\)。\(\blacksquare\)

\bigskip\hrule\bigskip

\subsubsection{17.3 二次函数}\label{ux4e8cux6b21ux51fdux6570}

\paragraph{17.3.1 一般式}\label{ux4e00ux822cux5f0f}

\textbf{定义 17.6}（二次函数）形如

\[y = ax^2 + bx + c \quad (a \neq 0)\]

的函数称为\textbf{二次函数}，定义域为 \(\mathbb{R}\)。其中 \(a\)
称为\textbf{二次项系数}，\(b\) 称为\textbf{一次项系数}，\(c\)
称为\textbf{常数项}。

\paragraph{17.3.2
配方法与顶点式}\label{ux914dux65b9ux6cd5ux4e0eux9876ux70b9ux5f0f}

\textbf{定理 17.11}（配方法）二次函数 \(y = ax^2 + bx + c\) 可以配方为：

\[y = a\left(x + \frac{b}{2a}\right)^2 + \frac{4ac - b^2}{4a}\]

\textbf{证明}：

\[y = ax^2 + bx + c = a\left(x^2 + \frac{b}{a}x\right) + c\]

\[= a\left(x^2 + \frac{b}{a}x + \frac{b^2}{4a^2} - \frac{b^2}{4a^2}\right) + c\]

\[= a\left(x + \frac{b}{2a}\right)^2 - \frac{b^2}{4a} + c\]

\[= a\left(x + \frac{b}{2a}\right)^2 + \frac{4ac - b^2}{4a}\]

\(\blacksquare\)

\textbf{定义 17.7}（顶点式）令
\(h = -\frac{b}{2a}\)，\(k = \frac{4ac - b^2}{4a}\)，则二次函数可写为\textbf{顶点式}：

\[y = a(x - h)^2 + k\]

其中 \((h, k)\) 为抛物线的\textbf{顶点}。

\paragraph{17.3.3 交点式}\label{ux4ea4ux70b9ux5f0f}

\textbf{定义 17.8}（交点式）当二次函数 \(y = ax^2 + bx + c\) 的判别式
\(\Delta = b^2 - 4ac \geq 0\) 时，方程 \(ax^2 + bx + c = 0\) 有实根
\(x_1, x_2\)（允许 \(x_1 = x_2\)），则可写为\textbf{交点式}：

\[y = a(x - x_1)(x - x_2)\]

\textbf{证明}：由因式分解定理，当 \(x_1, x_2\) 是 \(ax^2 + bx + c = 0\)
的根时，

\[ax^2 + bx + c = a(x - x_1)(x - x_2)\]

展开右边：\(a(x^2 - (x_1 + x_2)x + x_1 x_2) = ax^2 - a(x_1 + x_2)x + ax_1 x_2\)

与 \(ax^2 + bx + c\)
比较系数：\(b = -a(x_1 + x_2)\)，\(c = ax_1 x_2\)，这正是韦达定理的结论。\(\blacksquare\)

\paragraph{17.3.4
抛物线的基本性质}\label{ux629bux7269ux7ebfux7684ux57faux672cux6027ux8d28-1}

\textbf{定理 17.12} 二次函数
\(y = ax^2 + bx + c\)（\(a \neq 0\)）的图像是一条\textbf{抛物线}，具有以下性质：

\begin{enumerate}
\def\labelenumi{\arabic{enumi}.}
\tightlist
\item
  \textbf{对称轴}为 \(x = -\frac{b}{2a}\)
\item
  \textbf{顶点}为 \(\left(-\frac{b}{2a}, \frac{4ac - b^2}{4a}\right)\)
\item
  当 \(a > 0\) 时，开口向上；当 \(a < 0\) 时，开口向下
\end{enumerate}

\textbf{证明}：由顶点式 \(y = a(x - h)^2 + k\)。

\begin{enumerate}
\def\labelenumi{(\arabic{enumi})}
\item
  对于任意
  \(\delta\)，\(f(h + \delta) = a\delta^2 + k = f(h - \delta)\)，所以图像关于
  \(x = h = -\frac{b}{2a}\) 对称。
\item
  当 \(x = h\)
  时，\(y = k = \frac{4ac - b^2}{4a}\)，这就是顶点的纵坐标。
\item
  \(y - k = a(x - h)^2\)。当 \(a > 0\) 时，\(y - k \geq 0\)，即
  \(y \geq k\)，开口向上；当 \(a < 0\) 时，\(y - k \leq 0\)，即
  \(y \leq k\)，开口向下。\(\blacksquare\)
\end{enumerate}

\paragraph{17.3.5
二次函数的最值}\label{ux4e8cux6b21ux51fdux6570ux7684ux6700ux503c}

\textbf{定理 17.13}（在 \(\mathbb{R}\) 上的最值）二次函数
\(y = ax^2 + bx + c\)（\(a \neq 0\)）： - 当 \(a > 0\) 时，有最小值
\(y_{\min} = \frac{4ac - b^2}{4a}\)，在 \(x = -\frac{b}{2a}\) 处取到； -
当 \(a < 0\) 时，有最大值 \(y_{\max} = \frac{4ac - b^2}{4a}\)，在
\(x = -\frac{b}{2a}\) 处取到。

\textbf{证明}：由顶点式 \(y = a(x - h)^2 + k\)。

当 \(a > 0\) 时，\(a(x - h)^2 \geq 0\)，等号当且仅当 \(x = h\)。因此
\(y \geq k\)，最小值为 \(k\)。

当 \(a < 0\) 时，\(a(x - h)^2 \leq 0\)，等号当且仅当 \(x = h\)。因此
\(y \leq k\)，最大值为 \(k\)。\(\blacksquare\)

\textbf{定理 17.14}（在闭区间上的最值）设二次函数
\(y = ax^2 + bx + c\)（\(a > 0\)，开口向上），对称轴
\(x = h = -\frac{b}{2a}\)。在闭区间 \([m, n]\) 上：

\begin{enumerate}
\def\labelenumi{\arabic{enumi}.}
\tightlist
\item
  若 \(h \in [m, n]\)，则最小值在 \(x = h\) 处取到，最大值在 \(x = m\)
  或 \(x = n\) 处取到；
\item
  若 \(h < m\)，则最小值在 \(x = m\) 处取到，最大值在 \(x = n\) 处取到；
\item
  若 \(h > n\)，则最小值在 \(x = n\) 处取到，最大值在 \(x = m\) 处取到。
\end{enumerate}

\textbf{证明}：当 \(a > 0\) 时，\(f\) 在 \((-\infty, h]\) 上严格递减，在
\([h, +\infty)\) 上严格递增。因此在 \([m, n]\) 上的最值取决于 \(h\)
与区间的位置关系，结论由单调性的分析直接得到。\(\blacksquare\)

\paragraph{17.3.6
判别式与根的关系}\label{ux5224ux522bux5f0fux4e0eux6839ux7684ux5173ux7cfb}

\textbf{定理 17.15}（判别式）一元二次方程
\(ax^2 + bx + c = 0\)（\(a \neq 0\)）的判别式为 \(\Delta = b^2 - 4ac\)。

\begin{enumerate}
\def\labelenumi{\arabic{enumi}.}
\tightlist
\item
  当 \(\Delta > 0\)
  时，方程有两个不相等的实数根：\(x = \frac{-b \pm \sqrt{\Delta}}{2a}\)
\item
  当 \(\Delta = 0\) 时，方程有两个相等的实数根：\(x = -\frac{b}{2a}\)
\item
  当 \(\Delta < 0\) 时，方程无实数根
\end{enumerate}

\textbf{证明}：由配方法，

\[ax^2 + bx + c = a\left(x + \frac{b}{2a}\right)^2 + \frac{4ac - b^2}{4a} = 0\]

\[\left(x + \frac{b}{2a}\right)^2 = \frac{b^2 - 4ac}{4a^2} = \frac{\Delta}{4a^2}\]

当 \(\Delta > 0\)
时，\(x + \frac{b}{2a} = \pm \frac{\sqrt{\Delta}}{2a}\)，得两根。

当 \(\Delta = 0\) 时，\(x = -\frac{b}{2a}\)。

当 \(\Delta < 0\)
时，\(\frac{\Delta}{4a^2} < 0\)，实数的平方不可能为负，故无实根。\(\blacksquare\)

\textbf{定理 17.16}（韦达定理）设 \(x_1, x_2\) 是方程
\(ax^2 + bx + c = 0\)（\(a \neq 0\)）的两根，则：

\[x_1 + x_2 = -\frac{b}{a}, \quad x_1 x_2 = \frac{c}{a}\]

\textbf{证明}：由因式分解，\(ax^2 + bx + c = a(x - x_1)(x - x_2)\)。

展开：\(a[x^2 - (x_1 + x_2)x + x_1 x_2] = ax^2 - a(x_1 + x_2)x + ax_1 x_2\)

比较系数：\(b = -a(x_1 + x_2)\)，\(c = ax_1 x_2\)，即
\(x_1 + x_2 = -\frac{b}{a}\)，\(x_1 x_2 = \frac{c}{a}\)。\(\blacksquare\)

\bigskip\hrule\bigskip

\subsubsection{17.4 幂函数}\label{ux5e42ux51fdux6570}

\paragraph{17.4.1 定义}\label{ux5b9aux4e49}

\textbf{定义 17.9}（幂函数）形如

\[y = x^{\alpha}\]

的函数称为\textbf{幂函数}，其中 \(\alpha\) 为常数。

\paragraph{17.4.2
常见幂函数的性质}\label{ux5e38ux89c1ux5e42ux51fdux6570ux7684ux6027ux8d28}

\textbf{定理 17.17} 设 \(\alpha\) 为正整数 \(n\)，则 \(f(x) = x^n\)
的性质如下：

\begin{enumerate}
\def\labelenumi{\arabic{enumi}.}
\tightlist
\item
  \textbf{定义域}：\(\mathbb{R}\)。
\item
  当 \(n\) 为奇数时，\(f\) 是奇函数，在 \(\mathbb{R}\) 上严格递增。
\item
  当 \(n\) 为偶数时，\(f\) 是偶函数，在 \((-\infty, 0]\) 上严格递减，在
  \([0, +\infty)\) 上严格递增。
\end{enumerate}

\textbf{证明}：

\begin{enumerate}
\def\labelenumi{(\arabic{enumi})}
\setcounter{enumi}{1}
\tightlist
\item
  设 \(n\) 为奇数。\(f(-x) = (-x)^n = -x^n = -f(x)\)，故 \(f\)
  为奇函数。
\end{enumerate}

设 \(x_1 < x_2\)。若 \(0 \leq x_1 < x_2\)，则 \(x_1^n < x_2^n\)（因为
\(x^n\) 在 \([0, +\infty)\) 上严格递增，可用归纳法证明）。

若 \(x_1 < x_2 < 0\)，令 \(u = -x_2, v = -x_1\)，则 \(0 < u < v\)，且
\(x_2^n - x_1^n = (-u)^n - (-v)^n = -(u^n - v^n)\)。因为 \(n\) 为奇数且
\(u^n < v^n\)，所以 \(x_2^n - x_1^n > 0\)。

若 \(x_1 < 0 < x_2\)，则 \(x_1^n < 0 < x_2^n\)。综合三种情况，\(f\) 在
\(\mathbb{R}\) 上严格递增。

\begin{enumerate}
\def\labelenumi{(\arabic{enumi})}
\setcounter{enumi}{2}
\tightlist
\item
  设 \(n\) 为偶数。\(f(-x) = (-x)^n = x^n = f(x)\)，故 \(f\) 为偶函数。
\end{enumerate}

由 \(f\) 是偶函数，只需分析 \(x \geq 0\) 的单调性。当
\(0 \leq x_1 < x_2\) 时，\(x_1^n < x_2^n\)，故在 \([0, +\infty)\)
上严格递增。由偶函数的对称性，在 \((-\infty, 0]\)
上严格递减。\(\blacksquare\)

\textbf{定理 17.18} 设 \(\alpha = -1\)，则
\(f(x) = x^{-1} = \frac{1}{x}\) 就是反比例函数（\(k = 1\)
的情形），其性质已在 17.2 节中讨论。

\textbf{证明}：\(f(x) = x^{-1} = \frac{1}{x}\) 即 \(k = 1\) 的反比例函数
\(y = \frac{k}{x}\)，性质已在定义17.5中详述。\(\blacksquare\)

\textbf{定理 17.19} 设 \(\alpha = \frac{1}{2}\)，则
\(f(x) = x^{1/2} = \sqrt{x}\) 的性质如下：

\begin{enumerate}
\def\labelenumi{\arabic{enumi}.}
\tightlist
\item
  \textbf{定义域}：\([0, +\infty)\)。
\item
  \textbf{值域}：\([0, +\infty)\)。
\item
  在 \([0, +\infty)\) 上严格递增。
\end{enumerate}

\textbf{证明}：设 \(0 \leq x_1 < x_2\)。则
\(0 \leq \sqrt{x_1} < \sqrt{x_2}\)（因为若
\(\sqrt{x_1} \geq \sqrt{x_2}\)，则
\(x_1 = (\sqrt{x_1})^2 \geq (\sqrt{x_2})^2 = x_2\)，矛盾）。\(\blacksquare\)

\paragraph{17.4.3
幂函数的一般性质}\label{ux5e42ux51fdux6570ux7684ux4e00ux822cux6027ux8d28}

\textbf{定理 17.20}（幂函数的公共特征）对于任意实数 \(\alpha\)，幂函数
\(y = x^{\alpha}\) 的图像都过点 \((1, 1)\)。

\textbf{证明}：\(f(1) = 1^{\alpha} = 1\)。\(\blacksquare\)

\bigskip\hrule\bigskip

\subsubsection{17.5
函数的零点与方程的根}\label{ux51fdux6570ux7684ux96f6ux70b9ux4e0eux65b9ux7a0bux7684ux6839}

\paragraph{17.5.1
零点与根的关系}\label{ux96f6ux70b9ux4e0eux6839ux7684ux5173ux7cfb}

\textbf{定义 17.10} 函数 \(f: D \to \mathbb{R}\) 的\textbf{零点}是使
\(f(x) = 0\) 的 \(x\) 值，即方程 \(f(x) = 0\) 的\textbf{根}。

\textbf{定理 17.25}（一次函数的零点）一次函数
\(y = kx + b\)（\(k \neq 0\)）有且仅有一个零点 \(x = -\frac{b}{k}\)。

\textbf{证明}：\(kx + b = 0 \iff x = -\frac{b}{k}\)。\(\blacksquare\)

\textbf{定理 17.26}（二次函数零点的分布）二次函数
\(f(x) = ax^2 + bx + c\)（\(a \neq 0\)）的零点分布：

\begin{enumerate}
\def\labelenumi{\arabic{enumi}.}
\tightlist
\item
  两个正根 \(\iff\)
  \(\Delta > 0\)，\(x_1 + x_2 > 0\)，\(x_1 x_2 > 0\)，即
  \(\Delta > 0\)，\(-\frac{b}{a} > 0\)，\(\frac{c}{a} > 0\)。
\item
  两个负根 \(\iff\)
  \(\Delta > 0\)，\(-\frac{b}{a} < 0\)，\(\frac{c}{a} > 0\)。
\item
  一正一负根 \(\iff\) \(\Delta > 0\)，\(\frac{c}{a} < 0\)。
\end{enumerate}

\textbf{证明}：由韦达定理，\(x_1 + x_2 = -\frac{b}{a}\)，\(x_1 x_2 = \frac{c}{a}\)。

\begin{enumerate}
\def\labelenumi{(\arabic{enumi})}
\tightlist
\item
  两个正根 \(\iff\) \(x_1 > 0\) 且 \(x_2 > 0\) \(\iff\)
  \(x_1 + x_2 > 0\) 且 \(x_1 x_2 > 0\) 且 \(\Delta \geq 0\)。由于
  \(x_1 \neq x_2\) 时 \(\Delta > 0\)，等号仅在 \(x_1 = x_2\)
  时取到。\(\blacksquare\)
\end{enumerate}

\paragraph{17.5.2
一元二次不等式}\label{ux4e00ux5143ux4e8cux6b21ux4e0dux7b49ux5f0f-1}

\textbf{定理 17.27}（一元二次不等式解法）设
\(a > 0\)，\(\Delta = b^2 - 4ac\)。

\begin{enumerate}
\def\labelenumi{\arabic{enumi}.}
\tightlist
\item
  若 \(\Delta > 0\)，\(ax^2 + bx + c = 0\) 有两实根 \(x_1 < x_2\)：

  \begin{itemize}
  \tightlist
  \item
    \(ax^2 + bx + c > 0\) 的解集为
    \((-\infty, x_1) \cup (x_2, +\infty)\)
  \item
    \(ax^2 + bx + c < 0\) 的解集为 \((x_1, x_2)\)
  \end{itemize}
\item
  若 \(\Delta = 0\)，\(ax^2 + bx + c = 0\) 有重根
  \(x_0 = -\frac{b}{2a}\)：

  \begin{itemize}
  \tightlist
  \item
    \(ax^2 + bx + c > 0\) 的解集为
    \(\{x \mid x \neq x_0\} = \mathbb{R} \setminus \{x_0\}\)
  \item
    \(ax^2 + bx + c < 0\) 的解集为 \(\emptyset\)
  \end{itemize}
\item
  若 \(\Delta < 0\)：

  \begin{itemize}
  \tightlist
  \item
    \(ax^2 + bx + c > 0\) 的解集为 \(\mathbb{R}\)
  \item
    \(ax^2 + bx + c < 0\) 的解集为 \(\emptyset\)
  \end{itemize}
\end{enumerate}

\textbf{证明}：(1) 由交点式
\(f(x) = a(x - x_1)(x - x_2)\)（\(a > 0\)）。当 \(x < x_1\)
时，\(x - x_1 < 0\)，\(x - x_2 < 0\)，故 \(f(x) > 0\)。当
\(x_1 < x < x_2\) 时，\(x - x_1 > 0\)，\(x - x_2 < 0\)，故
\(f(x) < 0\)。当 \(x > x_2\) 时，\(x - x_1 > 0\)，\(x - x_2 > 0\)，故
\(f(x) > 0\)。

\begin{enumerate}
\def\labelenumi{(\arabic{enumi})}
\setcounter{enumi}{1}
\item
  \(f(x) = a(x - x_0)^2 \geq 0\)，等号当且仅当 \(x = x_0\)。
\item
  \(f(x) = a\left(x + \frac{b}{2a}\right)^2 + \frac{4ac - b^2}{4a} = a\left(x + \frac{b}{2a}\right)^2 + \frac{-\Delta}{4a}\)。由于
  \(\Delta < 0\)，\(-\Delta > 0\)，且 \(a > 0\)，故 \(f(x) > 0\) 对所有
  \(x\) 成立。\(\blacksquare\)
\end{enumerate}

\bigskip\hrule\bigskip

\subsubsection{17.6
函数模型综合比较}\label{ux51fdux6570ux6a21ux578bux7efcux5408ux6bd4ux8f83}

\paragraph{17.6.1
增长速度比较}\label{ux589eux957fux901fux5ea6ux6bd4ux8f83}

\textbf{定理 17.28} 设
\(a > 1\)，\(\alpha > 0\)。则指数函数的增长最终超过任意幂函数：

\[\lim_{x \to +\infty} \frac{x^{\alpha}}{a^x} = 0\]

\textbf{注}：本定理的证明使用了洛必达法则（定理 23.17），该法则将在 Ch
23 中严格证明。完整证明将在 Ch 23 中给出。\(\blacksquare\)

\textbf{定理 17.29} 设
\(\alpha > 0\)。则对数函数的增长最终慢于任意正幂函数：

\[\lim_{x \to +\infty} \frac{\ln x}{x^{\alpha}} = 0\]

\textbf{注}：本定理的证明使用了洛必达法则（定理 23.17），该法则将在 Ch
23 中严格证明。完整证明将在 Ch 23 中给出。\(\blacksquare\)

\paragraph{17.6.2
常见函数增长阶}\label{ux5e38ux89c1ux51fdux6570ux589eux957fux9636}

当 \(x \to +\infty\) 时，常见函数的增长速度从慢到快排列为：

\[\ln x \ll x^{\alpha} \quad (\alpha > 0) \ll a^x \quad (a > 1) \ll x! \ll x^x\]

这表明对数函数增长最慢，而指数函数增长极快------这种差异在算法分析和数值计算中具有重要意义。

\subsubsection{17.7
含绝对值的函数}\label{ux542bux7eddux5bf9ux503cux7684ux51fdux6570}

\paragraph{17.7.1
绝对值函数的基本性质}\label{ux7eddux5bf9ux503cux51fdux6570ux7684ux57faux672cux6027ux8d28}

\textbf{定理 17.30} 设 \(f(x) = |ax + b|\)（\(a \neq 0\)），则：

\[f(x) = \begin{cases} ax + b, & x \geq -\frac{b}{a} \text{（当 } a > 0 \text{ 时）} \\ -(ax + b), & x < -\frac{b}{a} \text{（当 } a > 0 \text{ 时）} \end{cases}\]

\textbf{证明}：由绝对值定义，当 \(ax + b \geq 0\)（即
\(x \geq -\frac{b}{a}\) 当 \(a > 0\)）时 \(|ax+b| = ax+b\)；当
\(ax+b < 0\) 时 \(|ax+b| = -(ax+b)\)。图像在 \(x = -\frac{b}{a}\)
处转折，两侧均为直线段，故为 \(V\) 形。\(\blacksquare\)

图像为 \(V\) 形折线，顶点在 \(\left(-\frac{b}{a}, 0\right)\)。

\textbf{定理 17.31} 设 \(f(x) = |ax^2 + bx + c|\)（\(a > 0\)）。当
\(\Delta > 0\) 时，图像由抛物线 \(y = ax^2 + bx + c\) 在 \(x\)
轴以下的部分翻折到 \(x\) 轴上方得到。

\textbf{证明}：\(|ax^2+bx+c| = ax^2+bx+c\) 当
\(ax^2+bx+c \geq 0\)，\(= -(ax^2+bx+c)\) 当 \(ax^2+bx+c < 0\)。当
\(\Delta > 0\) 时抛物线与 \(x\) 轴有两个交点，将 \(x\)
轴分为三段。中间段
\(ax^2+bx+c < 0\)，取绝对值后翻折到上方。\(\blacksquare\)

\paragraph{17.7.2
含绝对值的方程和不等式}\label{ux542bux7eddux5bf9ux503cux7684ux65b9ux7a0bux548cux4e0dux7b49ux5f0f}

\textbf{定理 17.32} \(|f(x)| < g(x)\) 等价于
\(-g(x) < f(x) < g(x)\)（其中 \(g(x) > 0\)）。

\textbf{证明}：若 \(f(x) \geq 0\)，则 \(|f(x)| = f(x)\)，不等式变为
\(f(x) < g(x)\)，即 \(0 \leq f(x) < g(x)\)。若 \(f(x) < 0\)，则
\(|f(x)| = -f(x)\)，不等式变为 \(-f(x) < g(x)\)，即
\(-g(x) < f(x) < 0\)。综合两种情形，\(-g(x) < f(x) < g(x)\)。反之亦然。\(\blacksquare\)

\textbf{定理 17.33} \(|f(x)| > g(x)\) 等价于 \(f(x) > g(x)\) 或
\(f(x) < -g(x)\)。

\textbf{证明}：若 \(f(x) \geq 0\)，则 \(|f(x)| = f(x) > g(x)\)，即
\(f(x) > g(x)\)。若 \(f(x) < 0\)，则 \(|f(x)| = -f(x) > g(x)\)，即
\(f(x) < -g(x)\)。综合得 \(f(x) > g(x)\) 或 \(f(x) < -g(x)\)。反之，若
\(f(x) > g(x) \geq 0\) 则 \(|f(x)| = f(x) > g(x)\)；若
\(f(x) < -g(x) \leq 0\) 则 \(|f(x)| = -f(x) > g(x)\)。\(\blacksquare\)

\textbf{定理
17.34}（三角不等式）\(|a + b| \leq |a| + |b|\)，等号当且仅当
\(ab \geq 0\)。

\textbf{证明}：\(|a + b|^2 = a^2 + 2ab + b^2 \leq a^2 + 2|a||b| + b^2 = (|a| + |b|)^2\)。\(\blacksquare\)

\textbf{推论
17.1}（推广的三角不等式）\(||a| - |b|| \leq |a - b| \leq |a| + |b|\)。

\textbf{证明}：\(|a| = |(a - b) + b| \leq |a - b| + |b|\)，故
\(|a| - |b| \leq |a - b|\)。交换 \(a, b\) 得
\(|b| - |a| \leq |a - b|\)。综合得
\(||a| - |b|| \leq |a - b|\)。右边由三角不等式得
\(|a - b| = |a + (-b)| \leq |a| + |b|\)。\(\blacksquare\)

\bigskip\hrule\bigskip

\textbf{本章展望}：本章研究了一次函数、反比例函数、二次函数和幂函数等基本初等函数。在下一章（Ch
18）中，我们将系统研究指数函数和对数函数，这两类函数是基本初等函数中最重要的非代数函数。特别地，我们将严格定义实数指数幂，并证明指数函数和对数函数的基本性质和运算法则。

\bigskip\hrule\bigskip

\paragraph{\texorpdfstring{17.4.4 幂函数 \(y = x^{\alpha}\)
的图像特征分类}{17.4.4 幂函数 y = x\^{}\{\textbackslash alpha\} 的图像特征分类}}\label{ux5e42ux51fdux6570-y-xalpha-ux7684ux56feux50cfux7279ux5f81ux5206ux7c7b}

\textbf{定理 17.21}（\(\alpha > 0\) 时幂函数的行为）设
\(\alpha > 0\)，则幂函数 \(y = x^{\alpha}\)（在 \(x > 0\) 上）满足：

\begin{enumerate}
\def\labelenumi{\arabic{enumi}.}
\tightlist
\item
  在 \((0, +\infty)\) 上严格递增。
\item
  \(\lim_{x \to 0^+} x^{\alpha} = 0\)。
\item
  \(\lim_{x \to +\infty} x^{\alpha} = +\infty\)。
\end{enumerate}

\textbf{证明}：完整证明需要指数函数的严格理论（见 Ch
18），此处给出基于有理指数幂的证明思路。

\textbf{证明思路}：(1) \textbf{严格递增性}：设 \(0 < x_1 < x_2\)，令
\(t = x_2/x_1 > 1\)。则 \(x_2^\alpha / x_1^\alpha = t^\alpha\)。当
\(t > 1\) 且 \(\alpha > 0\) 时，对有理指数
\(r\)，\(t^r > 1\)（由有理指数幂的单调性，可利用
\(t = 1 + s\)（\(s > 0\)）及伯努利不等式 \((1+s)^r \geq 1 + rs > 1\)
对正有理数 \(r\) 建立）；对于一般的正实数 \(\alpha\)，由
\(t^\alpha = \sup\{t^r \mid r \in \mathbb{Q},\, 0 < r \leq \alpha\} > 1\)，故
\(x_2^\alpha > x_1^\alpha\)。(2) 对于任意 \(\varepsilon > 0\)，取
\(\delta = \varepsilon^{1/\alpha}\)，当 \(0 < x < \delta\) 时
\(x^\alpha < \varepsilon\)。(3) 对于任意 \(M > 0\)，取
\(X = M^{1/\alpha}\)，当 \(x > X\) 时 \(x^\alpha > M\)。\(\blacksquare\)

\textbf{定理 17.22}（\(\alpha < 0\) 时幂函数的行为）设
\(\alpha < 0\)，则幂函数 \(y = x^{\alpha}\)（在 \(x > 0\) 上）满足：

\begin{enumerate}
\def\labelenumi{\arabic{enumi}.}
\tightlist
\item
  在 \((0, +\infty)\) 上严格递减。
\item
  \(\lim_{x \to 0^+} x^{\alpha} = +\infty\)。
\item
  \(\lim_{x \to +\infty} x^{\alpha} = 0\)。
\end{enumerate}

\textbf{证明}：令 \(\beta = -\alpha > 0\)，则
\(x^{\alpha} = \frac{1}{x^{\beta}}\)。

\begin{enumerate}
\def\labelenumi{(\arabic{enumi})}
\item
  由定理 17.21，\(x^{\beta}\) 严格递增且为正，故 \(\frac{1}{x^{\beta}}\)
  严格递减。
\item
  当 \(x \to 0^+\) 时，\(x^{\beta} \to 0^+\)，故
  \(\frac{1}{x^{\beta}} \to +\infty\)。
\item
  当 \(x \to +\infty\) 时，\(x^{\beta} \to +\infty\)，故
  \(\frac{1}{x^{\beta}} \to 0^+\)。\(\blacksquare\)
\end{enumerate}

\paragraph{\texorpdfstring{17.4.5 幂函数的图像关于 \(\alpha\)
的变化}{17.4.5 幂函数的图像关于 \textbackslash alpha 的变化}}\label{ux5e42ux51fdux6570ux7684ux56feux50cfux5173ux4e8e-alpha-ux7684ux53d8ux5316}

\textbf{定理 17.23} 所有幂函数
\(y = x^{\alpha}\)（\(\alpha \neq 0\)）的图像都过点 \((1, 1)\)。

\textbf{证明}：\(1^{\alpha} = 1\) 对所有 \(\alpha\)
成立。\(\blacksquare\)

\textbf{定理 17.24} 设 \(\alpha_1, \alpha_2 > 0\) 且
\(\alpha_1 < \alpha_2\)。则：

\begin{enumerate}
\def\labelenumi{\arabic{enumi}.}
\tightlist
\item
  当 \(0 < x < 1\) 时，\(x^{\alpha_1} > x^{\alpha_2}\)。
\item
  当 \(x > 1\) 时，\(x^{\alpha_1} < x^{\alpha_2}\)。
\end{enumerate}

\textbf{注}：本定理的证明使用了自然指数函数和自然对数（\(e^{\alpha \ln x}\)），这些概念将在
Ch 18 中严格定义。完整证明将在 Ch 18 中给出。\(\blacksquare\)

\bigskip\hrule\bigskip

\subsection{Ch 18
指数函数与对数函数}\label{ch-18-ux6307ux6570ux51fdux6570ux4e0eux5bf9ux6570ux51fdux6570}

\textbf{前置知识}：本章依赖 Ch 16（函数概念、反函数理论）和 Ch
17（幂函数）。指数函数和对数函数互为反函数，它们的理论建立在实数指数幂的严格定义之上。

\bigskip\hrule\bigskip

\subsubsection{18.1
实数指数的严格定义}\label{ux5b9eux6570ux6307ux6570ux7684ux4e25ux683cux5b9aux4e49}

\paragraph{18.1.1 整数指数}\label{ux6574ux6570ux6307ux6570}

对于 \(a > 0\) 和正整数
\(n\)，\(a^n = \underbrace{a \cdot a \cdots a}_{n}\)。

\textbf{定义 18.1}（零指数）对于 \(a > 0\)，定义 \(a^0 = 1\)。

\textbf{定义 18.2}（负整数指数）对于 \(a > 0\) 和正整数 \(n\)，定义
\(a^{-n} = \frac{1}{a^n}\)。

\textbf{定理 18.1}（整数指数的运算律）设
\(a > 0\)，\(m, n \in \mathbb{Z}\)。则：

\begin{enumerate}
\def\labelenumi{\arabic{enumi}.}
\tightlist
\item
  \(a^m \cdot a^n = a^{m+n}\)
\item
  \((a^m)^n = a^{mn}\)
\item
  \((ab)^n = a^n b^n\)（\(b > 0\)）
\end{enumerate}

\textbf{证明}：(1) 当 \(m, n\)
均为正整数时，由乘法的结合律直接可得。其余情形可以通过定义转化为正整数的情形。例如当
\(m > 0\)，\(n < 0\) 时，设 \(n = -k\)（\(k > 0\)），则
\(a^m \cdot a^{-k} = \frac{a^m}{a^k}\)。若 \(m > k\)，则
\(= a^{m-k} = a^{m+(-k)}\)；若 \(m < k\)，则
\(= \frac{1}{a^{k-m}} = a^{-(k-m)} = a^{m+(-k)}\)。其余情况类似。\(\blacksquare\)

\paragraph{18.1.2 有理数指数}\label{ux6709ux7406ux6570ux6307ux6570-1}

\textbf{定义 18.3}（正有理数指数）设 \(a > 0\)，\(m, n\)
为正整数。定义：

\[a^{m/n} = \sqrt[n]{a^m} = (\sqrt[n]{a})^m\]

其中 \(\sqrt[n]{a}\) 表示 \(a\) 的 \(n\) 次算术根（即满足 \(x^n = a\)
的唯一正实数 \(x\)，其存在性由实数的完备性保证）。

\textbf{定理 18.2}（有理数指数的良定义性）上述定义与 \(m/n\)
的表示方式无关。即若 \(\frac{m}{n} = \frac{p}{q}\)，则
\(a^{m/n} = a^{p/q}\)。

\textbf{证明}：由 \(\frac{m}{n} = \frac{p}{q}\)，有 \(mq = np\)。设
\(x = a^{m/n}\)，即 \(x^n = a^m\)。则

\[x^{nq} = (x^n)^q = (a^m)^q = a^{mq} = a^{np} = (a^n)^p\]

同时 \(x^{nq} = (x^q)^n\)。由 \(x^n = a^m\)，得 \((x^n)^q = a^{mq}\)。而
\(a^{p/q}\) 满足 \(y^q = a^p\)，故
\(y^{nq} = (y^q)^n = a^{np} = a^{mq}\)。因此
\(x^{nq} = y^{nq}\)，由正实数 \(nq\) 次方根的唯一性得
\(x = y\)。\(\blacksquare\)

\paragraph{18.1.3 实数指数}\label{ux5b9eux6570ux6307ux6570-1}

\textbf{定义 18.4}（实数指数）设
\(a > 0\)，\(x \in \mathbb{R}\)。取有理数列 \(\{r_n\}\)，使得
\(\lim_{n \to \infty} r_n = x\)。定义

\[a^x = \lim_{n \to \infty} a^{r_n}\]

\textbf{定理 18.3}（实数指数的良定义性）上述定义与有理数列 \(\{r_n\}\)
的选取无关。

\textbf{证明}：设 \(\{r_n\}\) 和 \(\{s_n\}\) 都是收敛于 \(x\)
的有理数列。需要证明
\(\lim_{n \to \infty} a^{r_n} = \lim_{n \to \infty} a^{s_n}\)。

由于 \(|r_n - s_n| \leq |r_n - x| + |x - s_n| \to 0\)，所以
\(r_n - s_n \to 0\)。

当 \(a > 1\) 时，由有理数指数的单调性，\(a^{r_n - s_n}\) 介于
\(a^{-|r_n - s_n|}\) 和 \(a^{|r_n - s_n|}\) 之间。由于
\(|r_n - s_n| \to 0\)，两者都趋近于 \(a^0 = 1\)，故
\(a^{r_n - s_n} \to 1\)。

所以

\[\lim_{n \to \infty} a^{r_n} = \lim_{n \to \infty} (a^{s_n} \cdot a^{r_n - s_n}) = \lim_{n \to \infty} a^{s_n} \cdot 1 = \lim_{n \to \infty} a^{s_n}\]

\(0 < a < 1\) 的情形类似。\(\blacksquare\)

\textbf{定理 18.4}（实数指数的运算律）设
\(a, b > 0\)，\(\alpha, \beta \in \mathbb{R}\)，则：

\begin{enumerate}
\def\labelenumi{\arabic{enumi}.}
\tightlist
\item
  \(a^{\alpha} \cdot a^{\beta} = a^{\alpha + \beta}\)
\item
  \((a^{\alpha})^{\beta} = a^{\alpha \beta}\)
\item
  \((ab)^{\alpha} = a^{\alpha} \cdot b^{\alpha}\)
\end{enumerate}

\textbf{证明}：取有理数列
\(\{r_n\} \to \alpha\)，\(\{s_n\} \to \beta\)。

\begin{enumerate}
\def\labelenumi{(\arabic{enumi})}
\item
  \(a^{\alpha} \cdot a^{\beta} = \lim_{n \to \infty} a^{r_n} \cdot \lim_{n \to \infty} a^{s_n} = \lim_{n \to \infty} (a^{r_n} \cdot a^{s_n}) = \lim_{n \to \infty} a^{r_n + s_n} = a^{\alpha + \beta}\)
\item
  取有理数列 \(\{s_n\} \to \beta\)。分两步完成。
\end{enumerate}

\textbf{整数幂引理}（仅依赖 (1)）：对任意实数 \(\gamma\) 和正整数
\(n\)，由 (1) 的乘法律反复应用（归纳法）：

\[(a^{\gamma})^n = \underbrace{a^{\gamma} \cdot a^{\gamma} \cdots a^{\gamma}}_{n \text{ 个}} = a^{n\gamma}\]

注意此引理是 (1) 的直接推论，\textbf{不使用} (2) 本身。

\begin{itemize}
\item
  \begin{enumerate}
  \def\labelenumi{(\roman{enumi})}
  \tightlist
  \item
    先证对有理数 \(s_n\) 有
    \((a^{\alpha})^{s_n} = a^{\alpha \cdot s_n}\)。设
    \(s_n = p/q\)（\(p \in \mathbb{Z}\)，\(q \in \mathbb{N}^+\)）。
  \end{enumerate}

  \textbf{整数幂部分}：由上述整数幂引理（仅用
  (1)），\((a^{\alpha})^p = a^{p\alpha}\)。

  \textbf{分数幂部分}：设 \(x = a^{p\alpha/q}\)（实数指数，由定义 18.4
  确定）。由整数幂引理（仅用
  (1)），\(x^q = (a^{p\alpha/q})^q = a^{q \cdot p\alpha/q} = a^{p\alpha}\)。因
  \(x > 0\)，由 \(q\) 次根的唯一性，\(x = (a^{p\alpha})^{1/q}\)。

  \textbf{合并}：由有理数幂的定义 \(b^{p/q} = (b^p)^{1/q}\)（应用于
  \(b = a^{\alpha}\)）：

  \[(a^{\alpha})^{s_n} = (a^{\alpha})^{p/q} = ((a^{\alpha})^p)^{1/q} = (a^{p\alpha})^{1/q} = a^{p\alpha/q} = a^{\alpha s_n}\]
\item
  \begin{enumerate}
  \def\labelenumi{(\roman{enumi})}
  \setcounter{enumi}{1}
  \tightlist
  \item
    由定义 18.4 的构造，\(a^x\) 在其定义域上连续（即 \(x_n \to x\) 蕴含
    \(a^{x_n} \to a^x\)；这是因为 \(a^x\) 单调且
    \(a^x = \sup\{a^r \mid r \in \mathbb{Q},\, r \leq x\}\)，由确界性质可证连续性）。故
    \(\alpha \cdot s_n \to \alpha \beta\) 蕴含
    \(\lim_{n \to \infty} a^{\alpha \cdot s_n} = a^{\alpha \beta}\)。
  \end{enumerate}
\end{itemize}

综合 (i) 和 (ii)，由实数幂的定义
\((a^{\alpha})^{\beta} = \lim_{n \to \infty} (a^{\alpha})^{s_n}\)（有理数列
\(s_n \to \beta\)），得

\[(a^{\alpha})^{\beta} = \lim_{n \to \infty} (a^{\alpha})^{s_n} = \lim_{n \to \infty} a^{\alpha s_n} = a^{\alpha \beta}\]

\begin{enumerate}
\def\labelenumi{(\arabic{enumi})}
\setcounter{enumi}{2}
\tightlist
\item
  的证明类似。\(\blacksquare\)
\end{enumerate}

\bigskip\hrule\bigskip

\subsubsection{18.2 指数函数}\label{ux6307ux6570ux51fdux6570}

\paragraph{18.2.1 定义}\label{ux5b9aux4e49-1}

\textbf{定义 18.5}（指数函数）设 \(a > 0\) 且
\(a \neq 1\)。\textbf{指数函数} \(f: \mathbb{R} \to (0, +\infty)\)
定义为：

\[f(x) = a^x\]

\paragraph{18.2.2
指数函数的性质}\label{ux6307ux6570ux51fdux6570ux7684ux6027ux8d28}

\textbf{引理 18.1}（严格单调函数的中值性质）设 \(f\) 在区间 \([p, q]\)
上严格单调递增，且满足如下\textbf{确界连续性条件}：对任意
\(x \in (p, q)\)，\(f(x) = \sup\{f(t) \mid t \in [p, x)\} = \inf\{f(t) \mid t \in (x, q]\}\)。（指数函数
\(a^x\)（\(a > 1\)）由定义 18.4 的构造
\(a^x = \sup\{a^r \mid r \in \mathbb{Q}, r \leq x\}\)
直接满足此条件。）若 \(f(p) < c < f(q)\)，则存在唯一的
\(\xi \in (p, q)\) 使得 \(f(\xi) = c\)。

\textbf{证明}：定义集合 \(S = \{x \in [p, q] \mid f(x) \leq c\}\)。由于
\(f(p) < c\)，故 \(p \in S\)，\(S\) 非空且有上界
\(q\)。由确界存在定理，设 \(\xi = \sup S\)，则 \(\xi \in [p, q]\)。

先证 \(f(\xi) \leq c\)：由上确界定义，对每个
\(n \in \mathbb{N}^+\)，存在 \(x_n \in S\) 使得
\(\xi - \frac{1}{n} < x_n \leq \xi\)。由确界连续性条件，\(f(\xi) = \sup\{f(t) \mid t < \xi\}\)。由
\(f(x_n) \leq c\)，得 \(f(\xi) \leq c\)。

再证 \(f(\xi) \geq c\)：先证 \(\xi < q\)。若
\(\xi = q\)，则由确界连续性条件，\(f(q) = \sup\{f(t) \mid t < q\}\)。存在
\(x_n \in S\) 使得 \(x_n \to q\)，故 \(f(q) = \sup f(x_n) \leq c\)，与
\(f(q) > c\) 矛盾。故
\(\xi < q\)。由确界连续性条件，\(f(\xi) = \inf\{f(t) \mid t > \xi\}\)。对任意
\(t \in (\xi, q]\)，由 \(\xi = \sup S\) 知 \(t \notin S\)，即
\(f(t) > c\)。因此 \(f(\xi) = \inf\{f(t) \mid t > \xi\} \geq c\)。

综上 \(f(\xi) = c\)。唯一性由 \(f\) 的严格单调性保证。\(\blacksquare\)

\textbf{注}：引理 18.1
的证明仅依赖确界原理和函数的确界连续性条件（由定义 18.4
的构造保证），不依赖后续章节的连续性理论（Ch
22）。这是介值定理在严格单调函数上的特殊形式。

\bigskip\hrule\bigskip

\textbf{定理 18.5}（指数函数的性质）设 \(f(x) = a^x\)，其中
\(a > 0\)，\(a \neq 1\)。

\begin{enumerate}
\def\labelenumi{\arabic{enumi}.}
\tightlist
\item
  \(f\) 的定义域为 \(\mathbb{R}\)，值域为 \((0, +\infty)\)。
\item
  \(f(0) = 1\)，即图像过点 \((0, 1)\)。
\item
  当 \(a > 1\) 时，\(f\) 在 \(\mathbb{R}\) 上严格递增；当 \(0 < a < 1\)
  时，\(f\) 在 \(\mathbb{R}\) 上严格递减。
\item
  \(f\) 无上界，下确界为 \(0\)（但取不到）。
\end{enumerate}

\textbf{证明}：

\begin{enumerate}
\def\labelenumi{(\arabic{enumi})}
\tightlist
\item
  定义域由定义直接得到。值域：首先 \(a^x > 0\) 对所有
  \(x \in \mathbb{R}\)
  成立（因为正数的任意实数次幂仍为正数）。下证值域恰好为
  \((0, +\infty)\)。不妨设 \(a > 1\)（\(0 < a < 1\) 时令
  \(b = 1/a > 1\)，则 \(a^x = b^{-x}\)，值域相同）。
\end{enumerate}

先证：对任意 \(M > 0\)，存在正整数 \(n\) 使得 \(a^n > M\)。若不然，则
\(\{a^n\}\) 有上界，设
\(s = \sup\{a^n : n \in \mathbb{N}\}\)。由上确界定义，存在 \(n\) 使得
\(a^n > s/a\)，则 \(a^{n+1} = a \cdot a^n > s\)，与 \(s\)
为上确界矛盾。故 \(a^n \to +\infty\)。

对任意 \(y > 0\)： - 若 \(y = 1\)，取 \(x = 0\)，则 \(a^0 = 1 = y\)。 -
若 \(y > 1\)，取 \(x_1 = 0\)，则 \(a^{x_1} = 1 < y\)。取正整数 \(N\)
使得 \(a^N > y\)，令 \(x_2 = N\)。由实数指数的构造（定义 18.4），\(a^x\)
是连续函数。由 (3)，\(a^x\) 在 \([x_1, x_2]\) 上严格递增。由引理
18.1，存在 \(\xi \in (x_1, x_2)\) 使得 \(a^{\xi} = y\)。 - 若
\(0 < y < 1\)，则 \(1/y > 1\)，由上段存在 \(x'\) 使得
\(a^{x'} = 1/y\)，故 \(a^{-x'} = y\)。

因此值域为 \((0, +\infty)\)。

\begin{enumerate}
\def\labelenumi{(\arabic{enumi})}
\setcounter{enumi}{1}
\item
  \(f(0) = a^0 = 1\)。
\item
  设 \(a > 1\)，\(x_1 < x_2\)。则
\end{enumerate}

\[\frac{f(x_2)}{f(x_1)} = \frac{a^{x_2}}{a^{x_1}} = a^{x_2 - x_1} > a^0 = 1\]

所以 \(f(x_2) > f(x_1)\)，\(f\) 严格递增。当 \(0 < a < 1\) 时类似证明。

\begin{enumerate}
\def\labelenumi{(\arabic{enumi})}
\setcounter{enumi}{3}
\tightlist
\item
  由 (3)，\(f\) 严格单调，故无上界。又 \(a^x > 0\) 对所有 \(x\)
  成立，故下确界为 \(0\)（但取不到）。\(\blacksquare\)
\end{enumerate}

\textbf{定理 18.6}（指数函数的无零点性）指数函数
\(f(x) = a^x\)（\(a > 0\)，\(a \neq 1\)）在 \(\mathbb{R}\) 上恒不为零。

\textbf{证明}：对于任意 \(x \in \mathbb{R}\)，\(a^x > 0\)（因为
\(a > 0\)，正数的任意实数次幂仍为正数）。\(\blacksquare\)

\bigskip\hrule\bigskip

\subsubsection{18.3 对数函数}\label{ux5bf9ux6570ux51fdux6570}

\paragraph{18.3.1 对数的定义}\label{ux5bf9ux6570ux7684ux5b9aux4e49}

\textbf{定义 18.6}（对数）设 \(a > 0\)，\(a \neq 1\)。对于
\(x > 0\)，\textbf{以 \(a\) 为底 \(x\) 的对数} \(\log_a x\) 定义为满足
\(a^y = x\) 的唯一实数 \(y\)，即：

\[y = \log_a x \Leftrightarrow a^y = x\]

\paragraph{18.3.2
对数的存在性和唯一性}\label{ux5bf9ux6570ux7684ux5b58ux5728ux6027ux548cux552fux4e00ux6027}

\textbf{定理 18.7}（对数的存在唯一性）设
\(a > 0\)，\(a \neq 1\)，\(x > 0\)。则存在唯一的实数 \(y\) 使得
\(a^y = x\)。

\textbf{证明}：

存在性：考虑函数 \(f(t) = a^t\)。当 \(a > 1\) 时，\(f\) 严格递增（由定理
18.5(3)）且连续（由实数指数的构造，定义 18.4）。由于 \(f(0) = 1\)，当
\(x > 1\) 时，由定理 18.5(1) 的证明可知存在 \(N\) 使得
\(f(N) = a^N > x\)。在 \([0, N]\) 上 \(f\) 连续且严格递增，由引理
18.1，存在 \(y \in (0, N)\) 使得 \(f(y) = x\)。当 \(0 < x < 1\)
时，\(1/x > 1\)，由上段存在 \(y' > 0\) 使得 \(a^{y'} = 1/x\)，则
\(a^{-y'} = x\)，取 \(y = -y' < 0\)。当 \(x = 1\)
时，\(y = 0\)。\(0 < a < 1\) 的情形类似。

唯一性：由于 \(f\) 严格单调，所以 \(y\) 唯一。\(\blacksquare\)

\paragraph{18.3.3
对数函数的性质}\label{ux5bf9ux6570ux51fdux6570ux7684ux6027ux8d28}

\textbf{定理 18.8}（对数函数的性质）设
\(g(x) = \log_a x\)（\(a > 0\)，\(a \neq 1\)）。

\begin{enumerate}
\def\labelenumi{\arabic{enumi}.}
\tightlist
\item
  \textbf{定义域}：\((0, +\infty)\)，\textbf{值域}：\(\mathbb{R}\)。
\item
  \(\log_a 1 = 0\)，\(\log_a a = 1\)。
\item
  当 \(a > 1\) 时，\(g\) 在 \((0, +\infty)\) 上严格递增；当
  \(0 < a < 1\) 时，\(g\) 在 \((0, +\infty)\) 上严格递减。
\item
  \(g\) 是指数函数 \(f(x) = a^x\) 的反函数。
\end{enumerate}

\textbf{证明}：

\begin{enumerate}
\def\labelenumi{(\arabic{enumi})}
\setcounter{enumi}{1}
\item
  \(a^0 = 1\)，故 \(\log_a 1 = 0\)。\(a^1 = a\)，故 \(\log_a a = 1\)。
\item
  对数函数是指数函数的反函数，由定理 16.12，反函数与原函数的单调性相同。
\item
  由定义，\(y = \log_a x \Leftrightarrow a^y = x\)，即
  \(g = f^{-1}\)。\(\blacksquare\)
\end{enumerate}

\paragraph{18.3.4
对数的运算法则}\label{ux5bf9ux6570ux7684ux8fd0ux7b97ux6cd5ux5219}

\textbf{定理 18.9}（对数运算法则）设
\(a > 0\)，\(a \neq 1\)，\(M, N > 0\)。则：

\begin{enumerate}
\def\labelenumi{(\arabic{enumi})}
\item
  \(\log_a(MN) = \log_a M + \log_a N\)
\item
  \(\log_a\frac{M}{N} = \log_a M - \log_a N\)
\item
  \(\log_a M^n = n \log_a M\)，其中 \(n \in \mathbb{R}\)
\end{enumerate}

\textbf{证明}：

\begin{enumerate}
\def\labelenumi{(\arabic{enumi})}
\tightlist
\item
  设 \(\log_a M = p\)，\(\log_a N = q\)。则 \(a^p = M\)，\(a^q = N\)。
\end{enumerate}

\[MN = a^p \cdot a^q = a^{p+q}\]

所以 \(\log_a(MN) = p + q = \log_a M + \log_a N\)。

\begin{enumerate}
\def\labelenumi{(\arabic{enumi})}
\setcounter{enumi}{1}
\tightlist
\item
  由
  (1)，\(\log_a M = \log_a\left(\frac{M}{N} \cdot N\right) = \log_a\frac{M}{N} + \log_a N\)。
\end{enumerate}

所以 \(\log_a\frac{M}{N} = \log_a M - \log_a N\)。

\begin{enumerate}
\def\labelenumi{(\arabic{enumi})}
\setcounter{enumi}{2}
\tightlist
\item
  设 \(\log_a M = p\)，则 \(a^p = M\)。所以
  \(M^n = (a^p)^n = a^{pn}\)，即
  \(\log_a M^n = pn = n \log_a M\)。\(\blacksquare\)
\end{enumerate}

\paragraph{18.3.5 换底公式}\label{ux6362ux5e95ux516cux5f0f}

\textbf{定理 18.10}（换底公式）设
\(a, b > 0\)，\(a, b \neq 1\)，\(x > 0\)。则

\[\log_a x = \frac{\log_b x}{\log_b a}\]

\textbf{证明}：设 \(\log_a x = y\)，则 \(a^y = x\)。

对两边取以 \(b\) 为底的对数：\(\log_b(a^y) = \log_b x\)。

由定理 18.9(3)，\(y \log_b a = \log_b x\)。

所以 \(y = \frac{\log_b x}{\log_b a}\)，即
\(\log_a x = \frac{\log_b x}{\log_b a}\)。\(\blacksquare\)

\paragraph{18.3.6
自然对数与常用对数}\label{ux81eaux7136ux5bf9ux6570ux4e0eux5e38ux7528ux5bf9ux6570}

\textbf{定义 18.7}（自然对数）以 \(e\)
为底的对数称为\textbf{自然对数}，记作 \(\ln x\)，即
\(\ln x = \log_e x\)。

其中
\(e = \lim_{n \to \infty}\left(1 + \frac{1}{n}\right)^n \approx 2.71828...\)
是一个重要的数学常数（其极限的存在性将在 Ch 22 中证明）。

\textbf{定义 18.8}（常用对数）以 \(10\)
为底的对数称为\textbf{常用对数}，记作 \(\lg x\)，即
\(\lg x = \log_{10} x\)。

\textbf{推论 18.1}
\(\ln x = \frac{\lg x}{\lg e}\)，\(\lg x = \frac{\ln x}{\ln 10}\)。

\textbf{证明}：直接由换底公式可得。\(\blacksquare\)

\bigskip\hrule\bigskip

\subsubsection{18.4
函数零点与二分法}\label{ux51fdux6570ux96f6ux70b9ux4e0eux4e8cux5206ux6cd5}

\paragraph{18.4.1 零点的定义}\label{ux96f6ux70b9ux7684ux5b9aux4e49}

\textbf{定义 18.9}（函数零点）设 \(f: D \to \mathbb{R}\)。若
\(x_0 \in D\) 满足 \(f(x_0) = 0\)，则称 \(x_0\) 为 \(f\)
的一个\textbf{零点}。

\paragraph{18.4.2
零点存在性定理}\label{ux96f6ux70b9ux5b58ux5728ux6027ux5b9aux7406}

\textbf{定理 18.11}（零点存在性定理）若函数 \(f(x)\) 在闭区间 \([a, b]\)
上连续，且 \(f(a) \cdot f(b) < 0\)（即 \(f(a)\) 与 \(f(b)\)
异号），则存在 \(\xi \in (a, b)\)，使得 \(f(\xi) = 0\)。

\textbf{证明}：不妨设 \(f(a) < 0 < f(b)\)（另一情形类似）。定义集合
\(S = \{x \in [a, b] \mid f(x) \leq 0\}\)。由 \(f(a) < 0\) 知
\(a \in S\)，\(S\) 非空且有上界 \(b\)。由确界存在定理，设
\(\xi = \sup S\)，则 \(\xi \in [a, b]\)。

先证 \(f(\xi) \leq 0\)。取 \(x_n \in S\) 使得
\(x_n \to \xi\)（由上确界的定义，对每个 \(n\)，存在 \(x_n \in S\) 满足
\(\xi - \frac{1}{n} < x_n \leq \xi\)）。由 \(f\)
连续，\(f(\xi) = \lim_{n \to \infty} f(x_n)\)。而 \(f(x_n) \leq 0\)，故
\(f(\xi) \leq 0\)。

再证 \(f(\xi) \geq 0\)。取 \(t_n = \xi + \frac{1}{n}\)（当 \(n\)
充分大时 \(t_n > \xi\) 且 \(t_n \leq b\)），则 \(t_n \notin S\)，故
\(f(t_n) > 0\)。由 \(f\)
连续，\(f(\xi) = \lim_{n \to \infty} f(t_n) \geq 0\)。

综上，\(f(\xi) = 0\)。\(\blacksquare\)

\textbf{注}：此证明仅依赖确界原理和连续性的操作性定义，不依赖 Ch 22
的介值定理。介值定理（Ch 22，定理
22.15）将在后续建立完整理论后给出更一般的证明。

\paragraph{18.4.3 二分法}\label{ux4e8cux5206ux6cd5}

\textbf{定义
18.10}（二分法）二分法是一种基于零点存在性定理的数值方法，用于近似求解方程
\(f(x) = 0\) 的根。

\textbf{算法描述}：

设 \(f\) 在 \([a_0, b_0]\) 上连续，\(f(a_0) \cdot f(b_0) < 0\)。对
\(n = 0, 1, 2, \ldots\)：

\begin{enumerate}
\def\labelenumi{\arabic{enumi}.}
\tightlist
\item
  计算中点 \(m_n = \frac{a_n + b_n}{2}\)。
\item
  若 \(f(m_n) = 0\)，则 \(m_n\) 就是零点，算法终止。
\item
  若 \(f(a_n) \cdot f(m_n) < 0\)，则零点在 \((a_n, m_n)\) 内，令
  \(a_{n+1} = a_n\)，\(b_{n+1} = m_n\)。
\item
  若 \(f(m_n) \cdot f(b_n) < 0\)，则零点在 \((m_n, b_n)\) 内，令
  \(a_{n+1} = m_n\)，\(b_{n+1} = b_n\)。
\end{enumerate}

\textbf{定理 18.12}（二分法的收敛性）设 \(f\) 在 \([a_0, b_0]\)
上连续，\(f(a_0) \cdot f(b_0) < 0\)，且 \(f\) 在 \((a_0, b_0)\)
上单调。则二分法产生的区间序列 \(\{[a_n, b_n]\}\) 满足：

\[b_n - a_n = \frac{b_0 - a_0}{2^n} \to 0\]

且序列 \(\{a_n\}\) 和 \(\{b_n\}\) 都收敛于 \(f\) 在 \((a_0, b_0)\)
内的唯一零点 \(\xi\)。

\textbf{证明}：由构造，\(b_n - a_n = \frac{b_0 - a_0}{2^n}\)。

\(\{a_n\}\) 单调递增且有上界 \(b_0\)，由确界存在定理，\(\{a_n\}\)
有上确界 \(\alpha = \sup\{a_n\}\)。同理 \(\{b_n\}\) 单调递减且有下界
\(a_0\)，有下确界 \(\beta = \inf\{b_n\}\)。

由 \(b_n - a_n \to 0\)，得 \(\alpha = \beta\)。

由 \(f\) 的连续性（确界连续性条件），\(f(\alpha) = \sup f(a_n) \leq 0\)
且 \(f(\alpha) = \inf f(b_n) \geq 0\)（或反之），故 \(f(\alpha) = 0\)。

由 \(f\) 的单调性，零点唯一，故 \(\alpha = \xi\)。\(\blacksquare\)

\textbf{注}：此证明仅依赖确界存在定理，不依赖 Ch 22
的单调有界定理。单调有界定理（Ch 22，定理
22.5）将在后续给出更一般的证明。

\bigskip\hrule\bigskip

\subsubsection{18.5
指数方程与对数方程}\label{ux6307ux6570ux65b9ux7a0bux4e0eux5bf9ux6570ux65b9ux7a0b}

\paragraph{18.5.1 指数方程}\label{ux6307ux6570ux65b9ux7a0b}

\textbf{定义 18.11}（指数方程）形如 \(a^{f(x)} = a^{g(x)}\)
或可化为此形式的方程称为\textbf{指数方程}，其中
\(a > 0\)，\(a \neq 1\)。

\textbf{定理 18.13} 指数方程
\(a^{f(x)} = a^{g(x)}\)（\(a > 0\)，\(a \neq 1\)）等价于
\(f(x) = g(x)\)。

\textbf{证明}：指数函数 \(y = a^x\) 是严格单调的（由定理 18.5），因此
\(a^u = a^v \iff u = v\)。\(\blacksquare\)

\paragraph{18.5.2 对数方程}\label{ux5bf9ux6570ux65b9ux7a0b}

\textbf{定义
18.12}（对数方程）含有对数表达式且未知数出现在对数内的方程称为\textbf{对数方程}。

\textbf{定理 18.14} 对数方程
\(\log_a f(x) = \log_a g(x)\)（\(a > 0\)，\(a \neq 1\)）在 \(f(x) > 0\)
且 \(g(x) > 0\) 的条件下等价于 \(f(x) = g(x)\)。

\textbf{证明}：对数函数 \(y = \log_a x\) 在 \((0, +\infty)\)
上严格单调，因此在定义域内
\(\log_a u = \log_a v \iff u = v\)。\(\blacksquare\)

\textbf{注意}：解对数方程时必须检验所得解是否满足真数大于零的条件。

\bigskip\hrule\bigskip

\subsubsection{18.6
指数函数与对数函数的图像关系}\label{ux6307ux6570ux51fdux6570ux4e0eux5bf9ux6570ux51fdux6570ux7684ux56feux50cfux5173ux7cfb}

\textbf{定理 18.15} 指数函数 \(y = a^x\) 和对数函数 \(y = \log_a x\)
的图像关于直线 \(y = x\) 对称。

\textbf{证明}：对数函数是指数函数的反函数，由定理
16.13，函数与其反函数的图像关于 \(y = x\) 对称。\(\blacksquare\)

\textbf{定理 18.16} 设 \(a > b > 1\)。则：

\begin{enumerate}
\def\labelenumi{\arabic{enumi}.}
\tightlist
\item
  当 \(x > 0\) 时，\(a^x > b^x > 1\)。
\item
  当 \(x < 0\) 时，\(0 < a^x < b^x < 1\)。
\item
  当 \(x > 1\) 时，\(\log_a x < \log_b x\)。
\item
  当 \(0 < x < 1\) 时，\(\log_a x > \log_b x\)（都为负值）。
\end{enumerate}

\textbf{证明}：

\begin{enumerate}
\def\labelenumi{(\arabic{enumi})}
\item
  \(\frac{a^x}{b^x} = \left(\frac{a}{b}\right)^x\)。当 \(x > 0\)
  时，\(\frac{a}{b} > 1\)，故 \(\left(\frac{a}{b}\right)^x > 1\)，即
  \(a^x > b^x\)。
\item
  由换底公式，\(\log_a x = \frac{\ln x}{\ln a}\)，\(\log_b x = \frac{\ln x}{\ln b}\)。
\end{enumerate}

当 \(x > 1\) 时，\(\ln x > 0\)。又 \(\ln a > \ln b > 0\)，故
\(\frac{\ln x}{\ln a} < \frac{\ln x}{\ln b}\)。\(\blacksquare\)

\bigskip\hrule\bigskip

\subsubsection{18.7 幂指函数}\label{ux5e42ux6307ux51fdux6570}

\textbf{定义 18.13}（幂指函数）形如 \(y = [f(x)]^{g(x)}\)
的函数称为\textbf{幂指函数}，其中 \(f(x) > 0\)。

\textbf{定理 18.17} 幂指函数可以化为指数函数和对数函数的复合：

\[[f(x)]^{g(x)} = e^{g(x) \ln f(x)}\]

\textbf{证明}：设
\(u = f(x) > 0\)，\(v = g(x)\)。由实数指数的定义，\(u^v = e^{v \ln u}\)。\(\blacksquare\)

\bigskip\hrule\bigskip

\subsubsection{18.8
指数函数与对数函数的应用模型}\label{ux6307ux6570ux51fdux6570ux4e0eux5bf9ux6570ux51fdux6570ux7684ux5e94ux7528ux6a21ux578b}

\paragraph{18.8.1
指数增长与指数衰减}\label{ux6307ux6570ux589eux957fux4e0eux6307ux6570ux8870ux51cf}

\textbf{定义 18.14}（指数增长模型）若量 \(Q(t)\) 满足微分方程
\(\frac{dQ}{dt} = kQ\)（\(k > 0\)）（注：\(\frac{dQ}{dt}\) 表示 \(Q\)
对时间 \(t\) 的变化率，即导数，严格定义见 Ch 23。直观理解：\(Q\)
的瞬时变化率与当前值成正比），则

\[Q(t) = Q_0 e^{kt}\]

其中 \(Q_0 = Q(0)\) 是初始量。这称为\textbf{指数增长}模型，\(k\)
称为\textbf{增长率}。

\textbf{定义 18.15}（指数衰减模型）若 \(k < 0\)，令
\(\lambda = -k > 0\)，则 \(Q(t) = Q_0 e^{-\lambda t}\)
称为\textbf{指数衰减}模型。\textbf{半衰期}
\(T_{1/2} = \frac{\ln 2}{\lambda}\)。

\textbf{定理 18.18} 在指数衰减模型中，经过一个半衰期后，量减半。

\textbf{证明}：\(Q(T_{1/2}) = Q_0 e^{-\lambda \cdot \frac{\ln 2}{\lambda}} = Q_0 e^{-\ln 2} = \frac{Q_0}{2}\)。\(\blacksquare\)

\paragraph{18.8.2 对数尺度}\label{ux5bf9ux6570ux5c3aux5ea6}

对数函数将乘法转化为加法，这一性质被广泛用于构建对数尺度：

\begin{itemize}
\tightlist
\item
  \textbf{里氏震级}：\(M = \log_{10}\left(\frac{A}{A_0}\right)\)，其中
  \(A\) 是振幅。
\item
  \textbf{分贝}：\(L = 10\log_{10}\left(\frac{I}{I_0}\right)\)，其中
  \(I\) 是声强。
\item
  \textbf{pH 值}：\(\text{pH} = -\log_{10}[\text{H}^+]\)。
\end{itemize}

这些对数尺度的共同特征是将跨越多个数量级的变化压缩到一个可管理的数值范围内。

\subsubsection{18.9 对数不等式}\label{ux5bf9ux6570ux4e0dux7b49ux5f0f}

\textbf{定理 18.19} 对任意 \(x > 0\)，\(\ln x \leq x - 1\)，等号当且仅当
\(x = 1\)。

\begin{quote}
\textbf{前向引用说明}：以下证明使用了导数工具（Ch
23）。本不等式的纯代数证明需要幂级数展开或极限技术，均超出 Ch 18
的范围。在 Ch 23 中（定理 23.31）将基于导数重新建立此不等式及其精细形式
\(x - x^2/2 < \ln(1+x) < x\)。以下提供两个独立的证明路径。
\end{quote}

\textbf{证明（导数法）}：令 \(f(x) = x - 1 - \ln x\)（\(x > 0\)）。则
\(f'(x) = 1 - \frac{1}{x}\)。当 \(x > 1\) 时 \(f'(x) > 0\)，\(f\)
递增；当 \(0 < x < 1\) 时 \(f'(x) < 0\)，\(f\) 递减。因此 \(f\) 在
\(x = 1\) 处取最小值 \(f(1) = 0\)，即 \(f(x) \geq 0\)。\(\blacksquare\)

\textbf{证明（指数转化法）}：以下证明将问题转化为等价命题
\(e^t \geq 1 + t\)（\(t \in \mathbb{R}\)）。

\textbf{步骤 1}（伯努利不等式）：对正整数 \(n\) 和实数
\(u > -1\)，\((1+u)^n \geq 1 + nu\)。此式可用数学归纳法直接证明（Ch
21，引理 21.2）。

\textbf{步骤 2}：证明 \(e > (1+1/n)^n\) 对所有正整数 \(n\) 成立。由
\(e\) 的定义（\(e = \sup\{(1+1/k)^k : k \in \mathbb{N}^+\}\)），需证
\((1+1/k)^k\) 关于 \(k\) 严格递增。由伯努利不等式：

\[\frac{(1+1/(n+1))^{n+1}}{(1+1/n)^n} = \frac{((n+2)/(n+1))^{n+1}}{((n+1)/n)^n}\]

利用
\((1+\frac{1}{n(n+2)})^{n+1} > 1 + \frac{n+1}{n(n+2)} > 1 + \frac{1}{n+1}\)
可证此比大于 1（详细计算略）。

\textbf{步骤 3}（正有理数情形）：设
\(r = p/q > 0\)（\(p, q \in \mathbb{N}^+\)）。由步骤
2，\(e > (1+1/q)^q\)，取 \(q\) 次根得
\(e^{1/q} > 1 + 1/q\)。由乘法律和伯努利不等式：

\[e^r = (e^{1/q})^p > (1 + 1/q)^p \geq 1 + p/q = 1 + r\]

\textbf{步骤 4}（正实数情形）：对 \(t > 0\)，取有理数列
\(r_n \to t\)（\(r_n > 0\)）。由 \(e^{r_n} > 1 + r_n\) 和 \(e^x\)
的连续性，\(e^t \geq 1 + t\)。再由
\(e^t = (e^{t/2})^2 \geq (1 + t/2)^2 = 1 + t + t^2/4 > 1 + t\)，得严格不等式。

\textbf{步骤 5}（非正实数情形）：\(t = 0\) 时 \(e^0 = 1 = 1 + 0\)。对
\(t \leq -1\)，\(e^t > 0 \geq 1 + t\)。对 \(-1 < t < 0\)，设
\(u = -t \in (0,1)\)。由 \(e^u > 1 + u\)（步骤 4）和
\(e^{-u} = 1/e^u\)，再利用不等式 \(e^u \leq \frac{1}{1-u}\)（此式可由
\(e^u\) 的幂级数展开 \(e^u = \sum u^k/k! \leq \sum u^k = \frac{1}{1-u}\)
得到，其中用到 \(k! \geq 1\)），得 \(e^{-u} \geq 1 - u = 1 + t\)。

\textbf{步骤 6}（回代）：由
\(e^t \geq 1 + t\)（\(t \in \mathbb{R}\)），令
\(t = \ln x\)（\(x > 0\)）：\(x = e^{\ln x} \geq 1 + \ln x\)，故
\(\ln x \leq x - 1\)。等号成立当且仅当 \(t = 0\) 即
\(x = 1\)。\(\blacksquare\)

\textbf{推论 18.2} 对任意
\(x > -1\)，\(\ln(1 + x) \leq x\)，等号当且仅当 \(x = 0\)。

\textbf{证明}：在定理 18.19 中用 \(1 + x\) 替换 \(x\)（注意 \(x > -1\)
保证 \(1 + x > 0\)）。\(\blacksquare\)

\textbf{推论 18.3} 对任意
\(x \in \mathbb{R}\)，\(e^x \geq 1 + x\)，等号当且仅当 \(x = 0\)。

\textbf{证明}：此即定理 18.19 证明中步骤 1--5 的结论。亦可由推论 18.2
直接推出：令 \(y = 1 + x\)，若 \(y > 0\) 则 \(\ln y \leq y - 1\)，代入
\(y = e^t\) 得 \(t \leq e^t - 1\) 即 \(e^t \geq 1 + t\)；若
\(y \leq 0\)（即 \(x \leq -1\)），则
\(1 + x \leq 0 < e^x\)，不等式自然成立。\(\blacksquare\)

\bigskip\hrule\bigskip

\textbf{本章展望}：本章建立了指数函数和对数函数的完整理论。在下一章（Ch
19）中，我们将研究三角函数------另一类重要的基本初等函数。三角函数将角与实数通过单位圆联系起来，其周期性使它们在描述周期现象中不可替代。

\bigskip\hrule\bigskip

\subsection{Ch 19 三角函数}\label{ch-19-ux4e09ux89d2ux51fdux6570}

\textbf{前置知识}：本章依赖 Ch
16（坐标系与函数概念，特别是函数的周期性、奇偶性）以及 Ch
13（锐角三角函数的基本定义）。本章将三角函数从锐角推广到任意角，并通过单位圆给出严格的函数定义。

\bigskip\hrule\bigskip

\subsubsection{19.1
角的概念的推广}\label{ux89d2ux7684ux6982ux5ff5ux7684ux63a8ux5e7f}

\paragraph{19.1.1 有向角}\label{ux6709ux5411ux89d2}

\textbf{定义 19.1}（有向角）设平面上有一条射线 \(OA\)，将其绕端点 \(O\)
旋转到另一位置
\(OB\)，则称此旋转形成了一个\textbf{有向角}。若旋转方向为逆时针，则称为\textbf{正角}；若为顺时针，则称为\textbf{负角}；若射线未旋转，则称为\textbf{零角}。

\paragraph{19.1.2
象限角与终边相同的角}\label{ux8c61ux9650ux89d2ux4e0eux7ec8ux8fb9ux76f8ux540cux7684ux89d2}

\textbf{定义
19.2}（象限角）在平面直角坐标系中，角的顶点与原点重合，角的始边与 \(x\)
轴正半轴重合。若角的终边落在某一象限内，则称该角为该象限的\textbf{象限角}。若终边落在坐标轴上，则称该角为\textbf{轴线角}。

\textbf{定义 19.3} 与角 \(\alpha\) 终边相同的所有角构成的集合为：

\[\{\beta \mid \beta = \alpha + 2k\pi, \, k \in \mathbb{Z}\}\]

\paragraph{19.1.3 弧度制}\label{ux5f27ux5ea6ux5236}

\textbf{定义 19.4}（弧度）在半径为 \(r\) 的圆中，长度等于半径 \(r\)
的圆弧所对的圆心角的大小定义为 \textbf{1 弧度}（1 rad）。

\textbf{定理 19.1}（弧度与角度的换算）

\[1 \text{ rad} = \frac{180°}{\pi} \approx 57.296°, \quad 1° = \frac{\pi}{180} \text{ rad}\]

\textbf{证明}：圆的周长为 \(C = 2\pi r\)，对应 \(360°\)。1 弧度对应弧长
\(s = r\)，因此
\(2\pi \text{ rad} = 360°\)，由此可得换算公式。\(\blacksquare\)

\textbf{定理 19.2}（弧长公式与扇形面积公式）设扇形的半径为
\(r\)，圆心角为 \(\theta\)（弧度），则弧长 \(s = r\theta\)，扇形面积
\(S = \frac{1}{2}r^2\theta\)。

\textbf{证明}：圆面积为 \(\pi r^2\)，对应圆心角
\(2\pi\)。扇形面积与圆心角成正比：

\[S = \frac{\theta}{2\pi} \cdot \pi r^2 = \frac{1}{2}r^2\theta\]

弧长类似：\(s = \frac{\theta}{2\pi} \cdot 2\pi r = r\theta\)。\(\blacksquare\)

\bigskip\hrule\bigskip

\subsubsection{19.2
三角函数的定义}\label{ux4e09ux89d2ux51fdux6570ux7684ux5b9aux4e49-1}

\paragraph{19.2.1
用单位圆定义三角函数}\label{ux7528ux5355ux4f4dux5706ux5b9aux4e49ux4e09ux89d2ux51fdux6570}

\textbf{定义 19.5}（单位圆）以原点 \(O\) 为圆心，\(1\)
为半径的圆称为\textbf{单位圆}，其方程为 \(x^2 + y^2 = 1\)。

\textbf{定义 19.6}（三角函数）设角 \(\alpha\) 的顶点在原点，始边与 \(x\)
轴正半轴重合，终边与单位圆交于点 \(P(x, y)\)。定义：

\[\sin \alpha = y, \quad \cos \alpha = x\]
\[\tan \alpha = \frac{y}{x} \quad (x \neq 0), \quad \cot \alpha = \frac{x}{y} \quad (y \neq 0)\]
\[\sec \alpha = \frac{1}{x} \quad (x \neq 0), \quad \csc \alpha = \frac{1}{y} \quad (y \neq 0)\]

\textbf{注} 由于 \(P\) 在单位圆上，\(x^2 + y^2 = 1\)，因此
\(|\sin \alpha| \leq 1\)，\(|\cos \alpha| \leq 1\)。

\paragraph{19.2.2
定义域和值域}\label{ux5b9aux4e49ux57dfux548cux503cux57df}

{\def\LTcaptype{none} % do not increment counter
\begin{longtable}[]{@{}
  >{\raggedright\arraybackslash}p{(\linewidth - 4\tabcolsep) * \real{0.3000}}
  >{\raggedright\arraybackslash}p{(\linewidth - 4\tabcolsep) * \real{0.4000}}
  >{\raggedright\arraybackslash}p{(\linewidth - 4\tabcolsep) * \real{0.3000}}@{}}
\toprule\noalign{}
\begin{minipage}[b]{\linewidth}\raggedright
函数
\end{minipage} & \begin{minipage}[b]{\linewidth}\raggedright
定义域
\end{minipage} & \begin{minipage}[b]{\linewidth}\raggedright
值域
\end{minipage} \\
\midrule\noalign{}
\endhead
\bottomrule\noalign{}
\endlastfoot
\(\sin \alpha\) & \(\mathbb{R}\) & \([-1, 1]\) \\
\(\cos \alpha\) & \(\mathbb{R}\) & \([-1, 1]\) \\
\(\tan \alpha\) &
\(\{\alpha \mid \alpha \neq \frac{\pi}{2} + k\pi, k \in \mathbb{Z}\}\) &
\(\mathbb{R}\) \\
\(\cot \alpha\) & \(\{\alpha \mid \alpha \neq k\pi, k \in \mathbb{Z}\}\)
& \(\mathbb{R}\) \\
\end{longtable}
}

\paragraph{19.2.3 各象限符号}\label{ux5404ux8c61ux9650ux7b26ux53f7}

设角 \(\alpha\) 的终边与单位圆交于 \(P(x, y)\)：

{\def\LTcaptype{none} % do not increment counter
\begin{longtable}[]{@{}lccc@{}}
\toprule\noalign{}
象限 & \(\sin \alpha\) & \(\cos \alpha\) & \(\tan \alpha\) \\
\midrule\noalign{}
\endhead
\bottomrule\noalign{}
\endlastfoot
I & \(+\) & \(+\) & \(+\) \\
II & \(+\) & \(-\) & \(-\) \\
III & \(-\) & \(-\) & \(+\) \\
IV & \(-\) & \(+\) & \(-\) \\
\end{longtable}
}

\bigskip\hrule\bigskip

\subsubsection{19.3
同角三角函数的基本关系}\label{ux540cux89d2ux4e09ux89d2ux51fdux6570ux7684ux57faux672cux5173ux7cfb}

\textbf{定理 19.3}（勾股关系）对任意角 \(\alpha\)：

\[\sin^2 \alpha + \cos^2 \alpha = 1\]

\textbf{证明}：设角 \(\alpha\) 的终边与单位圆交于点
\(P(x, y)\)。由单位圆方程 \(x^2 + y^2 = 1\)，代入
\(\sin \alpha = y\)，\(\cos \alpha = x\)，得：

\[\cos^2 \alpha + \sin^2 \alpha = x^2 + y^2 = 1\]

\(\blacksquare\)

\textbf{定理 19.4}（商数关系）对任意角
\(\alpha\)（\(\cos \alpha \neq 0\)）：

\[\tan \alpha = \frac{\sin \alpha}{\cos \alpha}\]

\textbf{证明}：\(\tan \alpha = \frac{y}{x} = \frac{\sin \alpha}{\cos \alpha}\)。\(\blacksquare\)

\textbf{定理 19.5}（正割关系）对任意角
\(\alpha\)（\(\cos \alpha \neq 0\)）：

\[1 + \tan^2 \alpha = \sec^2 \alpha\]

\textbf{证明}：由定理
19.3，\(\sin^2 \alpha + \cos^2 \alpha = 1\)。两边除以
\(\cos^2 \alpha\)：

\[\frac{\sin^2 \alpha}{\cos^2 \alpha} + 1 = \frac{1}{\cos^2 \alpha}\]

\[\tan^2 \alpha + 1 = \sec^2 \alpha\]

\(\blacksquare\)

\textbf{定理 19.6}（余割关系）对任意角
\(\alpha\)（\(\sin \alpha \neq 0\)）：

\[1 + \cot^2 \alpha = \csc^2 \alpha\]

\textbf{证明}：由定理 19.3，两边除以 \(\sin^2 \alpha\)：

\[1 + \frac{\cos^2 \alpha}{\sin^2 \alpha} = \frac{1}{\sin^2 \alpha}\]

\[1 + \cot^2 \alpha = \csc^2 \alpha\]

\(\blacksquare\)

\bigskip\hrule\bigskip

\subsubsection{19.4 诱导公式}\label{ux8bf1ux5bfcux516cux5f0f}

诱导公式揭示了三角函数在角变换下的关系。以下假设 \(k \in \mathbb{Z}\)。

\paragraph{\texorpdfstring{19.4.1 \(\alpha + 2k\pi\)
的三角函数}{19.4.1 \textbackslash alpha + 2k\textbackslash pi 的三角函数}}\label{alpha-2kpi-ux7684ux4e09ux89d2ux51fdux6570}

\textbf{定理 19.7} 对任意角 \(\alpha\) 和整数 \(k\)：

\[\sin(\alpha + 2k\pi) = \sin \alpha, \quad \cos(\alpha + 2k\pi) = \cos \alpha, \quad \tan(\alpha + 2k\pi) = \tan \alpha\]

\textbf{证明}：角 \(\alpha\) 和角 \(\alpha + 2k\pi\) 的终边相同（相差
\(2\pi\)
的整数倍），因此终边与单位圆的交点相同，故三角函数值相等。\(\blacksquare\)

\paragraph{\texorpdfstring{19.4.2 \(\pi + \alpha\)
的三角函数}{19.4.2 \textbackslash pi + \textbackslash alpha 的三角函数}}\label{pi-alpha-ux7684ux4e09ux89d2ux51fdux6570}

\textbf{定理 19.8} 对任意角 \(\alpha\)：

\[\sin(\pi + \alpha) = -\sin \alpha, \quad \cos(\pi + \alpha) = -\cos \alpha, \quad \tan(\pi + \alpha) = \tan \alpha\]

\textbf{证明}：设角 \(\alpha\) 的终边与单位圆交于 \(P(x, y)\)。角
\(\pi + \alpha\) 的终边是将 \(OP\) 反向延长，交单位圆于
\(P'(-x, -y)\)。因此：

\[\sin(\pi + \alpha) = -y = -\sin \alpha\]
\[\cos(\pi + \alpha) = -x = -\cos \alpha\]
\[\tan(\pi + \alpha) = \frac{-y}{-x} = \frac{y}{x} = \tan \alpha\]

\(\blacksquare\)

\paragraph{\texorpdfstring{19.4.3 \(-\alpha\)
的三角函数}{19.4.3 -\textbackslash alpha 的三角函数}}\label{alpha-ux7684ux4e09ux89d2ux51fdux6570}

\textbf{定理 19.9} 对任意角 \(\alpha\)：

\[\sin(-\alpha) = -\sin \alpha, \quad \cos(-\alpha) = \cos \alpha, \quad \tan(-\alpha) = -\tan \alpha\]

\textbf{证明}：设角 \(\alpha\) 的终边与单位圆交于 \(P(x, y)\)。角
\(-\alpha\) 的终边与单位圆交于 \(P'(x, -y)\)。因此：

\[\sin(-\alpha) = -y = -\sin \alpha, \quad \cos(-\alpha) = x = \cos \alpha, \quad \tan(-\alpha) = \frac{-y}{x} = -\tan \alpha\]

\(\blacksquare\)

\textbf{注} 此定理表明 \(\cos\) 是偶函数，\(\sin\) 和 \(\tan\)
是奇函数。

\paragraph{\texorpdfstring{19.4.4 \(\pi - \alpha\)
的三角函数}{19.4.4 \textbackslash pi - \textbackslash alpha 的三角函数}}\label{pi---alpha-ux7684ux4e09ux89d2ux51fdux6570}

\textbf{定理 19.10} 对任意角 \(\alpha\)：

\[\sin(\pi - \alpha) = \sin \alpha, \quad \cos(\pi - \alpha) = -\cos \alpha, \quad \tan(\pi - \alpha) = -\tan \alpha\]

\textbf{证明}：利用定理 19.8 和 19.9：

\[\sin(\pi - \alpha) = \sin[\pi + (-\alpha)] = -\sin(-\alpha) = \sin \alpha\]
\[\cos(\pi - \alpha) = \cos[\pi + (-\alpha)] = -\cos(-\alpha) = -\cos \alpha\]

\(\blacksquare\)

\paragraph{\texorpdfstring{19.4.5 \(\frac{\pi}{2} - \alpha\)
的三角函数}{19.4.5 \textbackslash frac\{\textbackslash pi\}\{2\} - \textbackslash alpha 的三角函数}}\label{fracpi2---alpha-ux7684ux4e09ux89d2ux51fdux6570}

\textbf{定理 19.11} 对任意角 \(\alpha\)：

\[\sin\left(\frac{\pi}{2} - \alpha\right) = \cos \alpha, \quad \cos\left(\frac{\pi}{2} - \alpha\right) = \sin \alpha, \quad \tan\left(\frac{\pi}{2} - \alpha\right) = \cot \alpha\]

\textbf{证明}：设角 \(\alpha\) 的终边与单位圆交于 \(P(x, y)\)。角
\(\frac{\pi}{2} - \alpha\) 的终边与单位圆交于 \(P'\)。由于
\(\frac{\pi}{2} - \alpha\) 与 \(\alpha\) 关于直线 \(y = x\) 对称，\(P'\)
的坐标为 \((y, x)\)。因此：

\[\sin\left(\frac{\pi}{2} - \alpha\right) = x = \cos \alpha, \quad \cos\left(\frac{\pi}{2} - \alpha\right) = y = \sin \alpha\]

\(\blacksquare\)

\paragraph{\texorpdfstring{19.4.6 \(\frac{\pi}{2} + \alpha\)
的三角函数}{19.4.6 \textbackslash frac\{\textbackslash pi\}\{2\} + \textbackslash alpha 的三角函数}}\label{fracpi2-alpha-ux7684ux4e09ux89d2ux51fdux6570}

\textbf{定理 19.12} 对任意角 \(\alpha\)：

\[\sin\left(\frac{\pi}{2} + \alpha\right) = \cos \alpha, \quad \cos\left(\frac{\pi}{2} + \alpha\right) = -\sin \alpha, \quad \tan\left(\frac{\pi}{2} + \alpha\right) = -\cot \alpha\]

\textbf{证明}：设 \(\beta = \frac{\pi}{2} - \alpha\)，则
\(\frac{\pi}{2} + \alpha = \pi - \beta\)。由定理 19.10 和 19.11：

\[\sin\left(\frac{\pi}{2} + \alpha\right) = \sin(\pi - \beta) = \sin \beta = \sin\left(\frac{\pi}{2} - \alpha\right) = \cos \alpha\]

\[\cos\left(\frac{\pi}{2} + \alpha\right) = \cos(\pi - \beta) = -\cos \beta = -\cos\left(\frac{\pi}{2} - \alpha\right) = -\sin \alpha\]

\(\blacksquare\)

\textbf{诱导公式口诀}：``奇变偶不变，符号看象限''------将
\(\frac{\pi}{2}\) 的倍数记为 \(n \cdot \frac{\pi}{2}\)，若 \(n\)
为奇数，则 \(\sin\) 变 \(\cos\)、\(\cos\) 变 \(\sin\)；若 \(n\)
为偶数，则不变。符号由原函数在对应象限的符号决定（将 \(\alpha\)
视为锐角）。

\bigskip\hrule\bigskip

\subsubsection{19.5
三角函数的图像和性质}\label{ux4e09ux89d2ux51fdux6570ux7684ux56feux50cfux548cux6027ux8d28}

\paragraph{\texorpdfstring{19.5.1 正弦函数
\(y = \sin x\)}{19.5.1 正弦函数 y = \textbackslash sin x}}\label{ux6b63ux5f26ux51fdux6570-y-sin-x}

\textbf{定理 19.13}（正弦函数的性质）

\begin{enumerate}
\def\labelenumi{\arabic{enumi}.}
\tightlist
\item
  \textbf{定义域}：\(\mathbb{R}\)；\textbf{值域}：\([-1, 1]\)
\item
  \textbf{周期}：\(T = 2\pi\)（最小正周期）
\item
  \textbf{奇偶性}：奇函数
\item
  \textbf{单调性}：在
  \(\left[-\frac{\pi}{2} + 2k\pi, \frac{\pi}{2} + 2k\pi\right]\)
  上单调递增；在
  \(\left[\frac{\pi}{2} + 2k\pi, \frac{3\pi}{2} + 2k\pi\right]\)
  上单调递减
\end{enumerate}

\textbf{证明周期性}：\(\sin(x + 2\pi) = \sin x\)（定理 19.7）。下证
\(2\pi\) 是最小正周期。设 \(T\) 是正周期，则 \(\sin(x + T) = \sin x\)
对所有 \(x\) 成立。取 \(x = 0\)：\(\sin T = 0\)，故 \(T = k\pi\)。取
\(x = \frac{\pi}{2}\)：\(\sin\left(\frac{\pi}{2} + T\right) = 1\)，即
\(\cos T = 1\)，故 \(T = 2m\pi\)。因此最小正周期为 \(2\pi\)。

\textbf{证明单调性}：在单位圆上，角 \(\theta\) 对应的点为
\((\cos\theta, \sin\theta)\)，其中 \(\sin\theta\) 是该点的纵坐标。当
\(\theta\) 从 \(-\frac{\pi}{2}\) 增大到 \(\frac{\pi}{2}\)
时，对应点沿单位圆从最低点 \((0, -1)\) 经右半圆运动到最高点
\((0, 1)\)。在此过程中，纵坐标 \(\sin\theta\) 从 \(-1\) 单调递增到
\(1\)。

严格来说，设
\(-\frac{\pi}{2} \leq x_1 < x_2 \leq \frac{\pi}{2}\)。在单位圆上，\(x_1\)
对应点 \(P_1\)，\(x_2\) 对应点 \(P_2\)。由于
\(x_2 - x_1 \in (0, \pi]\)，点 \(P_2\) 在 \(P_1\)
的上方（纵坐标更大）。这由单位圆的几何性质直接得到：在右半圆（含上下端点）上，角度增大时纵坐标严格递增。\(\blacksquare\)

\paragraph{\texorpdfstring{19.5.2 余弦函数
\(y = \cos x\)}{19.5.2 余弦函数 y = \textbackslash cos x}}\label{ux4f59ux5f26ux51fdux6570-y-cos-x}

\textbf{定理 19.14}（余弦函数的性质）

\begin{enumerate}
\def\labelenumi{\arabic{enumi}.}
\tightlist
\item
  \textbf{定义域}：\(\mathbb{R}\)；\textbf{值域}：\([-1, 1]\)
\item
  \textbf{周期}：\(T = 2\pi\)
\item
  \textbf{奇偶性}：偶函数
\item
  \textbf{单调性}：在 \([2k\pi, \pi + 2k\pi]\) 上单调递减；在
  \([\pi + 2k\pi, 2\pi + 2k\pi]\) 上单调递增
\end{enumerate}

\textbf{注} 由诱导公式
\(\cos x = \sin\left(x + \frac{\pi}{2}\right)\)，余弦函数的图像可由正弦函数向左平移
\(\frac{\pi}{2}\) 得到。

\textbf{证明}：(1) 由 \(\cos(x+2\pi) = \cos x\) 知周期为 \(2\pi\)。(2)
\(\cos(-x) = \cos x\) 为偶函数。(3) 由导数 \((\cos x)' = -\sin x\)，在
\([0,\pi]\) 上 \(\sin x \geq 0\) 故 \(\cos x\) 递减；在 \([\pi,2\pi]\)
上递增。\(\blacksquare\)

\paragraph{\texorpdfstring{19.5.3 正切函数
\(y = \tan x\)}{19.5.3 正切函数 y = \textbackslash tan x}}\label{ux6b63ux5207ux51fdux6570-y-tan-x}

\textbf{定理 19.15}（正切函数的性质）

\begin{enumerate}
\def\labelenumi{\arabic{enumi}.}
\tightlist
\item
  \textbf{定义域}：\(\{x \mid x \neq \frac{\pi}{2} + k\pi, k \in \mathbb{Z}\}\)；\textbf{值域}：\(\mathbb{R}\)
\item
  \textbf{周期}：\(T = \pi\)（最小正周期）
\item
  \textbf{奇偶性}：奇函数
\item
  \textbf{单调性}：在每个开区间
  \(\left(-\frac{\pi}{2} + k\pi, \frac{\pi}{2} + k\pi\right)\)
  上严格递增
\end{enumerate}

\textbf{证明周期性}：\(\tan(x + \pi) = \frac{\sin(x + \pi)}{\cos(x + \pi)} = \frac{-\sin x}{-\cos x} = \tan x\)。\(\blacksquare\)

\paragraph{\texorpdfstring{19.5.4 函数
\(y = A\sin(\omega x + \varphi) + B\)
的图像变换}{19.5.4 函数 y = A\textbackslash sin(\textbackslash omega x + \textbackslash varphi) + B 的图像变换}}\label{ux51fdux6570-y-asinomega-x-varphi-b-ux7684ux56feux50cfux53d8ux6362}

给定函数
\(y = A\sin(\omega x + \varphi) + B\)（\(A > 0\)，\(\omega > 0\)），其图像可由
\(y = \sin x\) 经以下变换得到：

\begin{enumerate}
\def\labelenumi{\arabic{enumi}.}
\tightlist
\item
  \textbf{相位变换}：\(y = \sin x \to y = \sin(x + \varphi)\)（左移
  \(\varphi\) 个单位，\(\varphi > 0\) 时）
\item
  \textbf{周期变换}：\(y = \sin(x + \varphi) \to y = \sin(\omega x + \varphi)\)（横坐标缩短为
  \(\frac{1}{\omega}\) 倍）
\item
  \textbf{振幅变换}：\(y = \sin(\omega x + \varphi) \to y = A\sin(\omega x + \varphi)\)（纵坐标扩大
  \(A\) 倍）
\item
  \textbf{平移变换}：\(y = A\sin(\omega x + \varphi) \to y = A\sin(\omega x + \varphi) + B\)（上移
  \(B\) 个单位）
\end{enumerate}

\textbf{性质}：振幅 \(A\)，周期 \(T = \frac{2\pi}{\omega}\)，频率
\(f = \frac{\omega}{2\pi}\)，最大值 \(A + B\)，最小值 \(-A + B\)。

\bigskip\hrule\bigskip

\subsubsection{19.6
正弦型函数的深入分析}\label{ux6b63ux5f26ux578bux51fdux6570ux7684ux6df1ux5165ux5206ux6790}

\paragraph{\texorpdfstring{19.6.1 函数
\(y = A\sin(\omega x + \varphi) + B\)
的完整分析}{19.6.1 函数 y = A\textbackslash sin(\textbackslash omega x + \textbackslash varphi) + B 的完整分析}}\label{ux51fdux6570-y-asinomega-x-varphi-b-ux7684ux5b8cux6574ux5206ux6790}

\textbf{定理 19.16} 设 \(A > 0\)，\(\omega > 0\)。函数
\(y = A\sin(\omega x + \varphi) + B\) 的性质如下：

\begin{enumerate}
\def\labelenumi{\arabic{enumi}.}
\tightlist
\item
  \textbf{振幅}：\(A\)（图像上最高点与最低点到中心线距离的一半）
\item
  \textbf{周期}：\(T = \frac{2\pi}{\omega}\)
\item
  \textbf{频率}：\(f = \frac{1}{T} = \frac{\omega}{2\pi}\)
\item
  \textbf{相位}：\(\varphi\)（初相）
\item
  \textbf{最大值}：\(A + B\)，\textbf{最小值}：\(-A + B\)
\item
  \textbf{中心线}：\(y = B\)
\end{enumerate}

\textbf{证明}：

\begin{enumerate}
\def\labelenumi{(\arabic{enumi})}
\setcounter{enumi}{1}
\tightlist
\item
  \(y(x + T) = A\sin(\omega(x + T) + \varphi) + B = A\sin(\omega x + 2\pi + \varphi) + B = A\sin(\omega x + \varphi) + B = y(x)\)。
\end{enumerate}

对于最小正周期：设 \(T'\) 是周期，则
\(\sin(\omega(x + T') + \varphi) = \sin(\omega x + \varphi)\) 对所有
\(x\) 成立，即 \(\omega T'\) 是 \(\sin\) 的周期，故
\(\omega T' = 2k\pi\)（\(k\) 为正整数），最小正周期
\(T' = \frac{2\pi}{\omega}\)。

\begin{enumerate}
\def\labelenumi{(\arabic{enumi})}
\setcounter{enumi}{4}
\tightlist
\item
  \(-1 \leq \sin(\omega x + \varphi) \leq 1\)，故
  \(-A + B \leq y \leq A + B\)。\(\blacksquare\)
\end{enumerate}

\paragraph{19.6.2
图像变换的严格描述}\label{ux56feux50cfux53d8ux6362ux7684ux4e25ux683cux63cfux8ff0}

\textbf{定理 19.17} 设 \(y = \sin x\) 的图像为
\(C_0\)，\(y = A\sin(\omega x + \varphi) + B\) 的图像为
\(C\)（\(A > 0\)，\(\omega > 0\)）。\(C\) 可由 \(C_0\)
经过以下变换序列得到（变换顺序可调）：

\textbf{路径一}（先平移后伸缩）：

\[C_0 \xrightarrow{\text{左移 } \varphi} C_1: y = \sin(x + \varphi) \xrightarrow{\text{横向压缩至 } \frac{1}{\omega}} C_2: y = \sin(\omega x + \varphi)\]

\[\xrightarrow{\text{纵向拉伸 } A \text{ 倍}} C_3: y = A\sin(\omega x + \varphi) \xrightarrow{\text{上移 } B} C\]

\textbf{路径二}（先伸缩后平移）：

\[C_0 \xrightarrow{\text{横向压缩至 } \frac{1}{\omega}} y = \sin(\omega x) \xrightarrow{\text{左移 } \frac{\varphi}{\omega}} y = \sin(\omega(x + \frac{\varphi}{\omega}))\]

\[\xrightarrow{\text{纵向拉伸 } A \text{ 倍}} y = A\sin(\omega x + \varphi) \xrightarrow{\text{上移 } B} C\]

\textbf{注意}：两种路径中，横向平移的量不同：路径一中平移
\(\varphi\)，路径二中平移
\(\frac{\varphi}{\omega}\)。这是因为横向伸缩会改变平移的效果。

\textbf{证明}：\(y = A\sin(\omega x + \varphi) + B\) 由 \(y = \sin x\)
经以下变换得到：(1) 横坐标缩放 \(1/\omega\) 倍（周期变为
\(2\pi/\omega\)）；(2) 向左平移 \(\varphi/\omega\)；(3) 纵坐标放大 \(A\)
倍（振幅变为 \(A\)）；(4) 向上平移
\(B\)。变换顺序不影响最终图像。\(\blacksquare\)

\bigskip\hrule\bigskip

\subsubsection{19.7
三角函数的综合性质}\label{ux4e09ux89d2ux51fdux6570ux7684ux7efcux5408ux6027ux8d28}

\paragraph{19.7.1
三角函数的值域}\label{ux4e09ux89d2ux51fdux6570ux7684ux503cux57df}

\textbf{定理 19.18} 函数 \(y = a\sin x + b\cos x\) 的值域为
\([-\sqrt{a^2 + b^2}, \sqrt{a^2 + b^2}]\)。

\textbf{证明}：设 \(a, b\) 不全为零（否则结论显然）。令
\(R = \sqrt{a^2 + b^2} > 0\)。

\textbf{上界的证明}：由柯西--施瓦茨不等式（或直接计算），

\[(a\sin x + b\cos x)^2 \leq (a^2 + b^2)(\sin^2 x + \cos^2 x) = a^2 + b^2 = R^2\]

因此 \(|a\sin x + b\cos x| \leq R\)，即值域包含于 \([-R, R]\)。

\textbf{取等的证明}：取 \(x_0\) 使得
\(\sin x_0 = \frac{a}{R}\)，\(\cos x_0 = \frac{b}{R}\)（由
\(\left(\frac{a}{R}\right)^2 + \left(\frac{b}{R}\right)^2 = 1\)
可知这样的 \(x_0\) 存在），则

\[a\sin x_0 + b\cos x_0 = \frac{a^2}{R} + \frac{b^2}{R} = \frac{a^2 + b^2}{R} = R\]

类似地，取 \(x_1 = x_0 + \pi\) 可得
\(a\sin x_1 + b\cos x_1 = -R\)。因此值域恰好为
\([-R, R] = [-\sqrt{a^2 + b^2}, \sqrt{a^2 + b^2}]\)。\(\blacksquare\)

\textbf{注}：此证明完全自包含，不依赖后续章节的结果。Ch 20
中将给出此结论的另一种证明（利用辅助角公式）。

\paragraph{19.7.2 三角不等式}\label{ux4e09ux89d2ux4e0dux7b49ux5f0f}

\textbf{定理 19.19} 对任意实数 \(x\)，\(|\sin x| \leq |x|\)。

\textbf{证明}：当 \(x = 0\) 时显然成立。当 \(x \neq 0\) 时，分两种情况：

当 \(0 < |x| < \frac{\pi}{2}\) 时，考虑单位圆中圆心角 \(|x|\)
对应的弧。设 \(P(\cos x, \sin x)\) 为单位圆上的点，\(P\) 到 \(x\)
轴的垂直距离为 \(|\sin x|\)，而弧长为
\(|x|\)。由于两点之间线段最短，垂线段长度严格小于弧长，故
\(|\sin x| < |x|\)（严格不等式）。

当 \(|x| \geq \frac{\pi}{2}\)
时，\(|\sin x| \leq 1 < \frac{\pi}{2} \leq |x|\)。

综合两种情况，\(|\sin x| \leq |x|\)。\(\blacksquare\)

\textbf{定理 19.20} 对任意实数 \(x\)，\(|\tan x| \geq |x|\)（在
\(\tan x\) 有定义时）。

\textbf{证明}：当 \(0 < |x| < \frac{\pi}{2}\) 时，考虑单位圆中圆心角
\(|x|\) 对应的弧和过点 \((1, 0)\) 的切线。设 \(P(\cos x, \sin x)\)
为单位圆上的点，切线与射线 \(\overrightarrow{OP}\) 的延长线交于点
\(T\)，则 \(|OT|\) 方向上的切线段长度为 \(|\tan x|\)（由直角三角形
\(O(1,0)T\) 可得 \(|OT| = |\tan x|\)）。弧长 \(|x|\) 严格小于切线段长
\(|\tan x|\)（因为弧是连接 \(P\)
附近两点的最短路径，而切线段是折线路径的一部分，可通过扇形面积比较严格论证：扇形面积
\(\frac{1}{2}|x| < \frac{1}{2}|\tan x|\) = 三角形面积）。故
\(|x| < |\tan x|\)。

当 \(|x| \geq \frac{\pi}{2}\) 时，\(|\tan x|\)
可以是无穷大或任意大（在不连续点附近），不等式也成立。\(\blacksquare\)

\paragraph{19.7.3
三角函数的有界性}\label{ux4e09ux89d2ux51fdux6570ux7684ux6709ux754cux6027}

\textbf{定理 19.21} 正弦函数和余弦函数是\textbf{全局有界}的，即
\(|\sin x| \leq 1\)，\(|\cos x| \leq 1\)。正切函数和余切函数是\textbf{局部有界}的但在整体上无界。

\textbf{证明}：\(|\sin x| = |y| \leq \sqrt{x^2 + y^2} = 1\)（\(P(x,y)\)
在单位圆上）。类似地 \(|\cos x| \leq 1\)。

对于正切函数，\(\lim_{x \to \frac{\pi}{2}^-} \tan x = +\infty\)，故无界。\(\blacksquare\)

\subsubsection{19.8
三角函数的深入应用}\label{ux4e09ux89d2ux51fdux6570ux7684ux6df1ux5165ux5e94ux7528}

\paragraph{19.8.1
正弦函数与余弦函数的正交性}\label{ux6b63ux5f26ux51fdux6570ux4e0eux4f59ux5f26ux51fdux6570ux7684ux6b63ux4ea4ux6027}

\begin{quote}
\textbf{前向引用说明}：以下定理的陈述和证明使用了定积分记号
\(\int_a^b f(x)\,dx\)，该概念将在 Ch
24（积分初步）中严格定义。此处提前引入此结果是因为正交性是三角函数最本质的代数性质之一，也是傅里叶分析的基础。读者可先接受定理的结论，在学习
Ch 24 后再补充验证证明细节。
\end{quote}

\textbf{定理 19.22} 对于正整数 \(m, n\)：

\[\int_0^{2\pi} \sin(mx)\sin(nx) \, dx = \begin{cases} 0, & m \neq n \\ \pi, & m = n \end{cases}\]

\[\int_0^{2\pi} \cos(mx)\cos(nx) \, dx = \begin{cases} 0, & m \neq n \\ \pi, & m = n \end{cases}\]

\[\int_0^{2\pi} \sin(mx)\cos(nx) \, dx = 0\]

\textbf{证明}：当 \(m \neq n\) 时，由积化和差公式：

\[\sin(mx)\sin(nx) = \frac{1}{2}[\cos((m-n)x) - \cos((m+n)x)]\]

\[\int_0^{2\pi} \sin(mx)\sin(nx) \, dx = \frac{1}{2}\left[\frac{\sin((m-n)x)}{m-n} - \frac{\sin((m+n)x)}{m+n}\right]_0^{2\pi} = 0\]

当 \(m = n\) 时：

\[\sin^2(nx) = \frac{1 - \cos(2nx)}{2}\]

\[\int_0^{2\pi} \sin^2(nx) \, dx = \frac{1}{2}\left[x - \frac{\sin(2nx)}{2n}\right]_0^{2\pi} = \pi\]

其余公式类似证明。\(\blacksquare\)

\textbf{注}：此正交性是傅里叶级数理论的基础。

\bigskip\hrule\bigskip

\subsubsection{19.9
三角函数在复数中的推广}\label{ux4e09ux89d2ux51fdux6570ux5728ux590dux6570ux4e2dux7684ux63a8ux5e7f}

\paragraph{19.9.1
复数的三角形式}\label{ux590dux6570ux7684ux4e09ux89d2ux5f62ux5f0f}

\textbf{定义 19.7} 复数
\(z = a + bi\)（\(a, b \in \mathbb{R}\)）可以写成\textbf{三角形式}：

\[z = r(\cos\theta + i\sin\theta)\]

其中 \(r = |z| = \sqrt{a^2 + b^2}\) 是\textbf{模}，\(\theta = \arg z\)
是\textbf{辐角}。

\textbf{定理 19.23}（复数乘法的几何意义）设
\(z_1 = r_1(\cos\theta_1 + i\sin\theta_1)\)，\(z_2 = r_2(\cos\theta_2 + i\sin\theta_2)\)。则

\[z_1 z_2 = r_1 r_2[\cos(\theta_1 + \theta_2) + i\sin(\theta_1 + \theta_2)]\]

\textbf{证明}：

\[z_1 z_2 = r_1 r_2(\cos\theta_1 + i\sin\theta_1)(\cos\theta_2 + i\sin\theta_2)\]

\[= r_1 r_2[(\cos\theta_1\cos\theta_2 - \sin\theta_1\sin\theta_2) + i(\sin\theta_1\cos\theta_2 + \cos\theta_1\sin\theta_2)]\]

\[= r_1 r_2[\cos(\theta_1 + \theta_2) + i\sin(\theta_1 + \theta_2)]\]

\(\blacksquare\)

\textbf{定理 19.24}（复数的 \(n\) 次方根）复数
\(z = r(\cos\theta + i\sin\theta)\) 的 \(n\) 次方根有 \(n\) 个：

\[z_k = \sqrt[n]{r}\left(\cos\frac{\theta + 2k\pi}{n} + i\sin\frac{\theta + 2k\pi}{n}\right), \quad k = 0, 1, \ldots, n-1\]

\textbf{证明}：设
\(w = \rho(\cos\varphi + i\sin\varphi)\)，\(w^n = z\)。则
\(\rho^n = r\)，\(n\varphi = \theta + 2k\pi\)。\(\blacksquare\)

\bigskip\hrule\bigskip

\textbf{本章展望}：有了三角函数的定义，自然要问：两个角的和的三角函数如何表示？在下一章（Ch
20）中，我们将系统推导三角恒等变换公式，包括和差角公式、二倍角公式、半角公式、积化和差与和差化积公式，以及正弦定理和余弦定理。这些恒等变换是三角学的核心工具，也是后续学习微积分中三角函数导数的基础。

\bigskip\hrule\bigskip

\subsection{Ch 20
三角恒等变换与反三角函数}\label{ch-20-ux4e09ux89d2ux6052ux7b49ux53d8ux6362ux4e0eux53cdux4e09ux89d2ux51fdux6570}

\textbf{前置知识}：本章依赖 Ch
19（三角函数的定义、同角关系、诱导公式）。本章的核心------和差角公式------的证明需要用到平面向量数量积的几何定义（\(\vec{a} \cdot \vec{b} = |\vec{a}||\vec{b}|\cos\theta\)）及其坐标公式，此为高中数学已学内容。其余公式均从和差角公式推导而出。

\bigskip\hrule\bigskip

\subsubsection{20.1 和差角公式}\label{ux548cux5deeux89d2ux516cux5f0f}

\paragraph{20.1.1
余弦的差角公式}\label{ux4f59ux5f26ux7684ux5deeux89d2ux516cux5f0f}

\textbf{定理 20.1}（余弦差角公式）对任意角 \(\alpha, \beta\)：

\[\cos(\alpha - \beta) = \cos \alpha \cos \beta + \sin \alpha \sin \beta\]

\textbf{证明（向量法）}：在单位圆上取点 \(A(\cos \alpha, \sin \alpha)\)
和 \(B(\cos \beta, \sin \beta)\)。则向量 \(\vec{OA}\) 与 \(\vec{OB}\)
的夹角为 \(\alpha - \beta\)。由数量积的定义：

\[\vec{OA} \cdot \vec{OB} = |\vec{OA}||\vec{OB}|\cos(\alpha - \beta) = 1 \cdot 1 \cdot \cos(\alpha - \beta) = \cos(\alpha - \beta)\]

利用数量积的坐标表示：

\[\vec{OA} \cdot \vec{OB} = \cos \alpha \cos \beta + \sin \alpha \sin \beta\]

因此：

\[\cos(\alpha - \beta) = \cos \alpha \cos \beta + \sin \alpha \sin \beta\]

\(\blacksquare\)

\paragraph{20.1.2
余弦的和角公式}\label{ux4f59ux5f26ux7684ux548cux89d2ux516cux5f0f}

\textbf{定理 20.2}（余弦和角公式）对任意角 \(\alpha, \beta\)：

\[\cos(\alpha + \beta) = \cos \alpha \cos \beta - \sin \alpha \sin \beta\]

\textbf{证明}：在定理 20.1 中将 \(\beta\) 替换为 \(-\beta\)，利用
\(\cos(-\beta) = \cos \beta\)，\(\sin(-\beta) = -\sin \beta\)：

\[\cos(\alpha + \beta) = \cos[\alpha - (-\beta)] = \cos \alpha \cos(-\beta) + \sin \alpha \sin(-\beta) = \cos \alpha \cos \beta - \sin \alpha \sin \beta\]

\(\blacksquare\)

\paragraph{20.1.3
正弦的和角公式}\label{ux6b63ux5f26ux7684ux548cux89d2ux516cux5f0f}

\textbf{定理 20.3}（正弦和角公式）对任意角 \(\alpha, \beta\)：

\[\sin(\alpha + \beta) = \sin \alpha \cos \beta + \cos \alpha \sin \beta\]

\textbf{证明}：利用诱导公式
\(\sin \theta = \cos\left(\frac{\pi}{2} - \theta\right)\) 和定理 20.1：

\[\sin(\alpha + \beta) = \cos\left[\frac{\pi}{2} - (\alpha + \beta)\right] = \cos\left[\left(\frac{\pi}{2} - \alpha\right) - \beta\right]\]

应用余弦差角公式：

\[= \cos\left(\frac{\pi}{2} - \alpha\right)\cos \beta + \sin\left(\frac{\pi}{2} - \alpha\right)\sin \beta = \sin \alpha \cos \beta + \cos \alpha \sin \beta\]

\(\blacksquare\)

\textbf{推论
20.1}（正弦差角公式）\(\sin(\alpha - \beta) = \sin \alpha \cos \beta - \cos \alpha \sin \beta\)。

\textbf{证明}：在定理 20.3 中将 \(\beta\) 替换为 \(-\beta\)
即可。\(\blacksquare\)

\paragraph{20.1.4
正切的和差角公式}\label{ux6b63ux5207ux7684ux548cux5deeux89d2ux516cux5f0f}

\textbf{定理 20.4}（正切和角公式）对任意角
\(\alpha, \beta\)（\(\cos \alpha, \cos \beta, \cos(\alpha + \beta)\)
均不为零）：

\[\tan(\alpha + \beta) = \frac{\tan \alpha + \tan \beta}{1 - \tan \alpha \tan \beta}\]

\textbf{证明}：

\[\tan(\alpha + \beta) = \frac{\sin(\alpha + \beta)}{\cos(\alpha + \beta)} = \frac{\sin \alpha \cos \beta + \cos \alpha \sin \beta}{\cos \alpha \cos \beta - \sin \alpha \sin \beta}\]

分子分母同除以 \(\cos \alpha \cos \beta\)：

\[= \frac{\tan \alpha + \tan \beta}{1 - \tan \alpha \tan \beta}\]

\(\blacksquare\)

\textbf{推论
20.2}（正切差角公式）\(\tan(\alpha - \beta) = \frac{\tan \alpha - \tan \beta}{1 + \tan \alpha \tan \beta}\)。

\textbf{证明}：在和角公式
\(\tan(\alpha+\beta) = \frac{\tan\alpha + \tan\beta}{1 - \tan\alpha\tan\beta}\)
中以 \(-\beta\) 代 \(\beta\)，利用 \(\tan(-\beta) = -\tan\beta\)
即得。\(\blacksquare\)

\bigskip\hrule\bigskip

\subsubsection{20.2 二倍角公式}\label{ux4e8cux500dux89d2ux516cux5f0f}

\textbf{定理 20.5}（二倍角公式）对任意角 \(\alpha\)：

\[\sin 2\alpha = 2\sin \alpha \cos \alpha\]

\[\cos 2\alpha = \cos^2 \alpha - \sin^2 \alpha = 2\cos^2 \alpha - 1 = 1 - 2\sin^2 \alpha\]

\[\tan 2\alpha = \frac{2\tan \alpha}{1 - \tan^2 \alpha}\]

\textbf{证明}：在定理 20.3 中令 \(\beta = \alpha\)：

\[\sin 2\alpha = \sin \alpha \cos \alpha + \cos \alpha \sin \alpha = 2\sin \alpha \cos \alpha\]

在定理 20.2 中令 \(\beta = \alpha\)：

\[\cos 2\alpha = \cos^2 \alpha - \sin^2 \alpha\]

利用 \(\sin^2 \alpha + \cos^2 \alpha = 1\)：

\[\cos 2\alpha = 2\cos^2 \alpha - 1 = 1 - 2\sin^2 \alpha\]

\(\blacksquare\)

\bigskip\hrule\bigskip

\subsubsection{20.3 半角公式}\label{ux534aux89d2ux516cux5f0f}

\textbf{定理 20.6}（半角公式）对任意角 \(\alpha\)（使各分母不为零）：

\[\sin \frac{\alpha}{2} = \pm\sqrt{\frac{1 - \cos \alpha}{2}}, \quad \cos \frac{\alpha}{2} = \pm\sqrt{\frac{1 + \cos \alpha}{2}}\]

\[\tan \frac{\alpha}{2} = \frac{\sin \alpha}{1 + \cos \alpha} = \frac{1 - \cos \alpha}{\sin \alpha}\]

\textbf{证明}：由 \(\cos \alpha = 1 - 2\sin^2 \frac{\alpha}{2}\)，得
\(\sin^2 \frac{\alpha}{2} = \frac{1 - \cos \alpha}{2}\)，故
\(\sin \frac{\alpha}{2} = \pm\sqrt{\frac{1 - \cos \alpha}{2}}\)。正负号由
\(\frac{\alpha}{2}\) 所在象限决定。

类似地，由 \(\cos \alpha = 2\cos^2 \frac{\alpha}{2} - 1\)，得
\(\cos \frac{\alpha}{2} = \pm\sqrt{\frac{1 + \cos \alpha}{2}}\)。

对于 \(\tan \frac{\alpha}{2}\)：

\[\tan \frac{\alpha}{2} = \frac{\sin \frac{\alpha}{2}}{\cos \frac{\alpha}{2}} = \frac{2\sin \frac{\alpha}{2} \cos \frac{\alpha}{2}}{2\cos^2 \frac{\alpha}{2}} = \frac{\sin \alpha}{1 + \cos \alpha}\]

\(\blacksquare\)

\bigskip\hrule\bigskip

\subsubsection{20.4
积化和差与和差化积}\label{ux79efux5316ux548cux5deeux4e0eux548cux5deeux5316ux79ef}

\paragraph{20.4.1
积化和差公式}\label{ux79efux5316ux548cux5deeux516cux5f0f}

\textbf{定理 20.7} 对任意角 \(\alpha, \beta\)：

\[\sin \alpha \cos \beta = \frac{1}{2}[\sin(\alpha + \beta) + \sin(\alpha - \beta)]\]

\[\cos \alpha \sin \beta = \frac{1}{2}[\sin(\alpha + \beta) - \sin(\alpha - \beta)]\]

\[\cos \alpha \cos \beta = \frac{1}{2}[\cos(\alpha - \beta) + \cos(\alpha + \beta)]\]

\[\sin \alpha \sin \beta = -\frac{1}{2}[\cos(\alpha + \beta) - \cos(\alpha - \beta)]\]

\textbf{证明}：由和角公式：

\[\sin(\alpha + \beta) = \sin \alpha \cos \beta + \cos \alpha \sin \beta \quad \cdots (1)\]

\[\sin(\alpha - \beta) = \sin \alpha \cos \beta - \cos \alpha \sin \beta \quad \cdots (2)\]

\((1) + (2)\)：\(\sin(\alpha + \beta) + \sin(\alpha - \beta) = 2\sin \alpha \cos \beta\)，故第一个公式成立。

\((1) - (2)\)：\(\sin(\alpha + \beta) - \sin(\alpha - \beta) = 2\cos \alpha \sin \beta\)，故第二个公式成立。

类似地由余弦和差角公式可证后两个公式。\(\blacksquare\)

\paragraph{20.4.2
和差化积公式}\label{ux548cux5deeux5316ux79efux516cux5f0f}

\textbf{定理 20.8} 对任意角 \(A, B\)：

\[\sin A + \sin B = 2\sin\frac{A + B}{2}\cos\frac{A - B}{2}\]

\[\sin A - \sin B = 2\cos\frac{A + B}{2}\sin\frac{A - B}{2}\]

\[\cos A + \cos B = 2\cos\frac{A + B}{2}\cos\frac{A - B}{2}\]

\[\cos A - \cos B = -2\sin\frac{A + B}{2}\sin\frac{A - B}{2}\]

\textbf{证明}：设
\(\alpha = \frac{A + B}{2}\)，\(\beta = \frac{A - B}{2}\)，则
\(A = \alpha + \beta\)，\(B = \alpha - \beta\)。由定理 20.7：

\[\sin A + \sin B = \sin(\alpha + \beta) + \sin(\alpha - \beta) = 2\sin \alpha \cos \beta = 2\sin\frac{A + B}{2}\cos\frac{A - B}{2}\]

其余公式类似证明。\(\blacksquare\)

\bigskip\hrule\bigskip

\subsubsection{20.5 辅助角公式}\label{ux8f85ux52a9ux89d2ux516cux5f0f}

\textbf{定理 20.9}（辅助角公式）设 \(a, b\) 不全为零，则：

\[a\sin \theta + b\cos \theta = \sqrt{a^2 + b^2}\sin(\theta + \varphi)\]

其中 \(\varphi\) 满足
\(\cos \varphi = \frac{a}{\sqrt{a^2 + b^2}}\)，\(\sin \varphi = \frac{b}{\sqrt{a^2 + b^2}}\)。

\textbf{证明}：令 \(R = \sqrt{a^2 + b^2} > 0\)。定义角 \(\varphi\) 使得
\(\cos \varphi = \frac{a}{R}\)，\(\sin \varphi = \frac{b}{R}\)。这样的
\(\varphi\) 存在，因为
\(\cos^2 \varphi + \sin^2 \varphi = \frac{a^2 + b^2}{R^2} = 1\)。

\[a\sin \theta + b\cos \theta = R\left(\frac{a}{R}\sin \theta + \frac{b}{R}\cos \theta\right) = R(\cos \varphi \sin \theta + \sin \varphi \cos \theta) = R\sin(\theta + \varphi)\]

\(\blacksquare\)

\textbf{推论 20.3} 利用辅助角公式，\(a\sin \theta + b\cos \theta\)
的取值范围为 \([-\sqrt{a^2 + b^2}, \sqrt{a^2 + b^2}]\)。

\textbf{证明}：由辅助角公式，\(a\sin\theta + b\cos\theta = \sqrt{a^2+b^2}\sin(\theta+\varphi)\)。因
\(|\sin(\theta+\varphi)| \leq 1\)，故取值范围为
\([-\sqrt{a^2+b^2}, \sqrt{a^2+b^2}]\)。\(\blacksquare\)

\bigskip\hrule\bigskip

\subsubsection{20.6
正弦定理与余弦定理}\label{ux6b63ux5f26ux5b9aux7406ux4e0eux4f59ux5f26ux5b9aux7406-1}

\paragraph{20.6.1 正弦定理}\label{ux6b63ux5f26ux5b9aux7406}

\textbf{定理 20.10}（正弦定理）在 \(\triangle ABC\) 中，设角 \(A, B, C\)
所对的边分别为 \(a, b, c\)，外接圆半径为 \(R\)，则：

\[\frac{a}{\sin A} = \frac{b}{\sin B} = \frac{c}{\sin C} = 2R\]

\textbf{证明}：设 \(\triangle ABC\) 的外接圆为 \(\odot O\)，半径为
\(R\)。

\textbf{情形一}：\(A\) 为锐角。作直径 \(BD\)，则
\(\angle BCD = 90°\)（直径所对的圆周角）。在 \(\triangle BCD\) 中：

\[\sin D = \frac{BC}{BD} = \frac{a}{2R}\]

由同弧上的圆周角相等，\(\angle D = \angle A\)（因为 \(\angle D\) 和
\(\angle A\) 都是弧 \(BC\) 所对的圆周角）。因此：

\[\sin A = \frac{a}{2R} \implies \frac{a}{\sin A} = 2R\]

\textbf{情形二}：\(A\) 为钝角。作直径 \(BD\)，则
\(\angle D = 180° - A\)，\(\sin D = \sin(180° - A) = \sin A\)。类似可得
\(\frac{a}{\sin A} = 2R\)。

\textbf{情形三}：\(A = 90°\)。则 \(a\)
为外接圆的直径，\(\sin A = 1\)，\(\frac{a}{\sin A} = a = 2R\)。

对 \(B, C\) 同理。\(\blacksquare\)

\paragraph{20.6.2 余弦定理}\label{ux4f59ux5f26ux5b9aux7406}

\textbf{定理 20.11}（余弦定理）在 \(\triangle ABC\) 中，设角 \(A, B, C\)
所对的边分别为 \(a, b, c\)，则：

\[c^2 = a^2 + b^2 - 2ab\cos C\]

\[a^2 = b^2 + c^2 - 2bc\cos A\]

\[b^2 = a^2 + c^2 - 2ac\cos B\]

\textbf{证明（坐标法）}：以 \(C\) 为原点，\(CB\) 所在直线为 \(x\)
轴建立坐标系。则：

\begin{itemize}
\tightlist
\item
  \(C = (0, 0)\)，\(B = (a, 0)\)，\(A = (b\cos C, b\sin C)\)
\end{itemize}

\[c^2 = |AB|^2 = (b\cos C - a)^2 + (b\sin C)^2\]

\[= b^2\cos^2 C - 2ab\cos C + a^2 + b^2\sin^2 C\]

\[= a^2 + b^2(\cos^2 C + \sin^2 C) - 2ab\cos C = a^2 + b^2 - 2ab\cos C\]

\(\blacksquare\)

\textbf{推论 20.4}（勾股定理）当 \(C = 90°\)
时，\(\cos C = 0\)，余弦定理退化为 \(c^2 = a^2 + b^2\)。

\bigskip\hrule\bigskip

\subsubsection{20.7 反三角函数}\label{ux53cdux4e09ux89d2ux51fdux6570}

\paragraph{20.7.1 反正弦函数}\label{ux53cdux6b63ux5f26ux51fdux6570}

\textbf{定义 20.1}（反正弦函数）正弦函数 \(y = \sin x\) 在
\(\left[-\frac{\pi}{2}, \frac{\pi}{2}\right]\)
上的反函数称为\textbf{反正弦函数}，记作 \(y = \arcsin x\)。

\textbf{性质}： 1.
\textbf{定义域}：\([-1, 1]\)；\textbf{值域}：\(\left[-\frac{\pi}{2}, \frac{\pi}{2}\right]\)
2. \textbf{单调性}：严格递增 3. \textbf{奇偶性}：奇函数

\textbf{证明奇偶性}：设 \(y = \arcsin x\)，则
\(\sin y = x\)，\(y \in \left[-\frac{\pi}{2}, \frac{\pi}{2}\right]\)。\(\sin(-y) = -\sin y = -x\)，\(-y = \arcsin(-x)\)，故
\(\arcsin(-x) = -\arcsin x\)。\(\blacksquare\)

\paragraph{20.7.2 反余弦函数}\label{ux53cdux4f59ux5f26ux51fdux6570}

\textbf{定义 20.2}（反余弦函数）余弦函数 \(y = \cos x\) 在 \([0, \pi]\)
上的反函数称为\textbf{反余弦函数}，记作 \(y = \arccos x\)。

\textbf{性质}： 1.
\textbf{定义域}：\([-1, 1]\)；\textbf{值域}：\([0, \pi]\) 2.
\textbf{单调性}：严格递减

\textbf{定理 20.12}
\(\arcsin x + \arccos x = \frac{\pi}{2}\)（\(x \in [-1, 1]\)）。

\textbf{证明}：设 \(\alpha = \arcsin x\)，则
\(\sin \alpha = x\)，\(\alpha \in \left[-\frac{\pi}{2}, \frac{\pi}{2}\right]\)。

\[\cos\left(\frac{\pi}{2} - \alpha\right) = \sin \alpha = x\]

又 \(\frac{\pi}{2} - \alpha \in [0, \pi]\)，故
\(\arccos x = \frac{\pi}{2} - \alpha = \frac{\pi}{2} - \arcsin x\)。\(\blacksquare\)

\paragraph{20.7.3 反正切函数}\label{ux53cdux6b63ux5207ux51fdux6570}

\textbf{定义 20.3}（反正切函数）正切函数 \(y = \tan x\) 在
\(\left(-\frac{\pi}{2}, \frac{\pi}{2}\right)\)
上的反函数称为\textbf{反正切函数}，记作 \(y = \arctan x\)。

\textbf{性质}： 1.
\textbf{定义域}：\(\mathbb{R}\)；\textbf{值域}：\(\left(-\frac{\pi}{2}, \frac{\pi}{2}\right)\)
2. \textbf{单调性}：严格递增 3. \textbf{奇偶性}：奇函数

\bigskip\hrule\bigskip

\subsubsection{20.8
三角恒等式的系统化}\label{ux4e09ux89d2ux6052ux7b49ux5f0fux7684ux7cfbux7edfux5316}

\paragraph{20.8.1 万能公式}\label{ux4e07ux80fdux516cux5f0f}

\textbf{定理 20.13}（万能公式）设
\(t = \tan\frac{\alpha}{2}\)（\(\alpha \neq \pi + 2k\pi\)），则：

\[\sin \alpha = \frac{2t}{1 + t^2}, \quad \cos \alpha = \frac{1 - t^2}{1 + t^2}, \quad \tan \alpha = \frac{2t}{1 - t^2}\]

\textbf{证明}：由半角公式，令 \(t = \tan\frac{\alpha}{2}\)，则

\[\sin \alpha = 2\sin\frac{\alpha}{2}\cos\frac{\alpha}{2} = \frac{2\sin\frac{\alpha}{2}\cos\frac{\alpha}{2}}{\sin^2\frac{\alpha}{2} + \cos^2\frac{\alpha}{2}} = \frac{2\tan\frac{\alpha}{2}}{1 + \tan^2\frac{\alpha}{2}} = \frac{2t}{1 + t^2}\]

\[\cos \alpha = \cos^2\frac{\alpha}{2} - \sin^2\frac{\alpha}{2} = \frac{\cos^2\frac{\alpha}{2} - \sin^2\frac{\alpha}{2}}{\cos^2\frac{\alpha}{2} + \sin^2\frac{\alpha}{2}} = \frac{1 - \tan^2\frac{\alpha}{2}}{1 + \tan^2\frac{\alpha}{2}} = \frac{1 - t^2}{1 + t^2}\]

\[\tan \alpha = \frac{\sin \alpha}{\cos \alpha} = \frac{\frac{2t}{1+t^2}}{\frac{1-t^2}{1+t^2}} = \frac{2t}{1 - t^2}\]

\(\blacksquare\)

\paragraph{20.8.2 三角方程}\label{ux4e09ux89d2ux65b9ux7a0b}

\textbf{定义
20.4}（三角方程）含有三角函数且未知数出现在三角函数内的方程称为\textbf{三角方程}。

\textbf{定理 20.14}（基本三角方程的通解）

\begin{enumerate}
\def\labelenumi{\arabic{enumi}.}
\item
  \(\sin x = a\)（\(|a| \leq 1\)）的通解：\(x = \arcsin a + 2k\pi\) 或
  \(x = \pi - \arcsin a + 2k\pi\)（\(k \in \mathbb{Z}\)），简记为
  \(x = (-1)^n \arcsin a + n\pi\)（\(n \in \mathbb{Z}\)）。
\item
  \(\cos x = a\)（\(|a| \leq 1\)）的通解：\(x = \pm\arccos a + 2k\pi\)（\(k \in \mathbb{Z}\)）。
\item
  \(\tan x = a\)
  的通解：\(x = \arctan a + k\pi\)（\(k \in \mathbb{Z}\)）。
\end{enumerate}

\textbf{证明}：(1) 设 \(x_0 = \arcsin a\)，即
\(\sin x_0 = a\)，\(x_0 \in \left[-\frac{\pi}{2}, \frac{\pi}{2}\right]\)。\(\sin x = a = \sin x_0\)
的充要条件是 \(x = x_0 + 2k\pi\) 或
\(x = \pi - x_0 + 2k\pi\)。\(\blacksquare\)

\bigskip\hrule\bigskip

\subsubsection{20.9
三角形面积公式}\label{ux4e09ux89d2ux5f62ux9762ux79efux516cux5f0f}

\textbf{定理 20.15} 在 \(\triangle ABC\) 中，设角 \(A, B, C\)
所对的边分别为 \(a, b, c\)，则：

\[S = \frac{1}{2}ab\sin C = \frac{1}{2}bc\sin A = \frac{1}{2}ac\sin B\]

\textbf{证明}：以 \(C\) 为原点，\(CB\) 为 \(x\) 轴正方向建立坐标系。则
\(B = (a, 0)\)，\(A = (b\cos C, b\sin C)\)。

\(\triangle ABC\) 的面积 = 底 \(\times\) 高
\(\times \frac{1}{2} = a \cdot |b\sin C| \cdot \frac{1}{2} = \frac{1}{2}ab|\sin C|\)。

由于 \(C \in (0, \pi)\)，\(\sin C > 0\)，故
\(S = \frac{1}{2}ab\sin C\)。\(\blacksquare\)

\textbf{推论 20.5}（海伦公式）设 \(s = \frac{a + b + c}{2}\)，则

\[S = \sqrt{s(s-a)(s-b)(s-c)}\]

\textbf{证明}：由余弦定理，\(\cos C = \frac{a^2 + b^2 - c^2}{2ab}\)。

\[\sin^2 C = 1 - \cos^2 C = \frac{4a^2b^2 - (a^2 + b^2 - c^2)^2}{4a^2b^2}\]

分子因式分解：

\[4a^2b^2 - (a^2 + b^2 - c^2)^2 = [2ab + (a^2 + b^2 - c^2)][2ab - (a^2 + b^2 - c^2)]\]

\[= [(a+b)^2 - c^2][c^2 - (a-b)^2] = (a+b+c)(a+b-c)(c+a-b)(c-a+b)\]

\[= 2s \cdot 2(s-c) \cdot 2(s-b) \cdot 2(s-a) = 16s(s-a)(s-b)(s-c)\]

所以

\[S = \frac{1}{2}ab\sin C = \frac{1}{2}ab \cdot \frac{\sqrt{16s(s-a)(s-b)(s-c)}}{2ab} = \sqrt{s(s-a)(s-b)(s-c)}\]

\(\blacksquare\)

\bigskip\hrule\bigskip

\subsubsection{20.10
反三角函数的综合性质}\label{ux53cdux4e09ux89d2ux51fdux6570ux7684ux7efcux5408ux6027ux8d28}

\paragraph{20.10.1
反三角函数的导数（预备知识）}\label{ux53cdux4e09ux89d2ux51fdux6570ux7684ux5bfcux6570ux9884ux5907ux77e5ux8bc6}

\textbf{定理 20.16} 反三角函数的导数公式：

\begin{enumerate}
\def\labelenumi{\arabic{enumi}.}
\tightlist
\item
  \((\arcsin x)' = \frac{1}{\sqrt{1 - x^2}}\)，\(x \in (-1, 1)\)
\item
  \((\arccos x)' = -\frac{1}{\sqrt{1 - x^2}}\)，\(x \in (-1, 1)\)
\item
  \((\arctan x)' = \frac{1}{1 + x^2}\)，\(x \in \mathbb{R}\)
\end{enumerate}

\textbf{注}：以上公式的严格证明需要隐函数求导和链式法则（Ch 23
中建立），此处仅列出结果供后续使用。\(\blacksquare\)

\paragraph{20.10.2
反三角函数的恒等式}\label{ux53cdux4e09ux89d2ux51fdux6570ux7684ux6052ux7b49ux5f0f}

\textbf{定理 20.17} 对任意 \(x \in \mathbb{R}\)：

\[\arctan x + \arctan\frac{1}{x} = \begin{cases} \frac{\pi}{2}, & x > 0 \\ -\frac{\pi}{2}, & x < 0 \end{cases}\]

\textbf{证明}：设
\(\alpha = \arctan x\)，\(\beta = \arctan\frac{1}{x}\)。则
\(\tan\alpha = x\)，\(\tan\beta = \frac{1}{x}\)。

\[\tan(\alpha + \beta) = \frac{x + \frac{1}{x}}{1 - x \cdot \frac{1}{x}} = \frac{x + \frac{1}{x}}{0}\]

分母为零，说明 \(\alpha + \beta = \pm\frac{\pi}{2}\)。当 \(x > 0\)
时，\(\alpha \in (0, \frac{\pi}{2})\)，\(\beta \in (0, \frac{\pi}{2})\)，故
\(\alpha + \beta = \frac{\pi}{2}\)。当 \(x < 0\)
时类似。\(\blacksquare\)

\subsubsection{20.11
三角函数的级数表示（预备知识）}\label{ux4e09ux89d2ux51fdux6570ux7684ux7ea7ux6570ux8868ux793aux9884ux5907ux77e5ux8bc6}

\paragraph{20.11.1
正弦函数和余弦函数的幂级数}\label{ux6b63ux5f26ux51fdux6570ux548cux4f59ux5f26ux51fdux6570ux7684ux5e42ux7ea7ux6570}

由泰勒展开（Ch 23，推论 23.2），我们有以下重要展开式：

\textbf{定理 20.18} 对任意 \(x \in \mathbb{R}\)：

\[\sin x = \sum_{k=0}^{\infty} \frac{(-1)^k}{(2k+1)!}x^{2k+1} = x - \frac{x^3}{3!} + \frac{x^5}{5!} - \frac{x^7}{7!} + \cdots\]

\[\cos x = \sum_{k=0}^{\infty} \frac{(-1)^k}{(2k)!}x^{2k} = 1 - \frac{x^2}{2!} + \frac{x^4}{4!} - \frac{x^6}{6!} + \cdots\]

这两个级数对所有实数 \(x\) 都绝对收敛。

\textbf{证明}：由泰勒公式，\(\sin x = \sum_{k=0}^{n} \frac{(\sin)^{(k)}(0)}{k!}x^k + R_n(x)\)。

先计算 \(\sin x\) 在 \(x = 0\) 处的各阶导数。由
\((\sin x)' = \cos x\)，\((\cos x)' = -\sin x\)，逐阶求导得：

\[(\sin x)^{(0)}\big|_{x=0} = \sin 0 = 0, \quad (\sin x)^{(1)}\big|_{x=0} = \cos 0 = 1\]

\[(\sin x)^{(2)}\big|_{x=0} = -\sin 0 = 0, \quad (\sin x)^{(3)}\big|_{x=0} = -\cos 0 = -1\]

一般地，\((\sin x)^{(k)}\big|_{x=0} = \sin\left(\frac{k\pi}{2}\right)\)。当
\(k = 2m\)（偶数）时为 \(0\)，当 \(k = 2m+1\)（奇数）时为 \((-1)^m\)。

余项 \(R_n(x) = \frac{\sin^{(n+1)}(\xi)}{(n+1)!}x^{n+1}\)。由于
\(|\sin^{(n+1)}(\xi)| \leq 1\)，\(|R_n(x)| \leq \frac{|x|^{n+1}}{(n+1)!} \to 0\)。\(\blacksquare\)

\paragraph{20.11.2 欧拉公式}\label{ux6b27ux62c9ux516cux5f0f-1}

\textbf{定理 20.19}（欧拉公式）对任意实数 \(x\)：

\[e^{ix} = \cos x + i\sin x\]

其中 \(i = \sqrt{-1}\) 是虚数单位。

\textbf{证明}：将 \(e^z\) 的幂级数展开中的 \(z\) 替换为 \(ix\)：

\[e^{ix} = \sum_{n=0}^{\infty} \frac{(ix)^n}{n!} = \sum_{k=0}^{\infty} \frac{(ix)^{2k}}{(2k)!} + \sum_{k=0}^{\infty} \frac{(ix)^{2k+1}}{(2k+1)!}\]

\[= \sum_{k=0}^{\infty} \frac{(-1)^k x^{2k}}{(2k)!} + i\sum_{k=0}^{\infty} \frac{(-1)^k x^{2k+1}}{(2k+1)!} = \cos x + i\sin x\]

\(\blacksquare\)

\textbf{推论 20.6}（欧拉恒等式）令 \(x = \pi\)：\(e^{i\pi} + 1 = 0\)。

\textbf{证明}：\(e^{i\pi} = \cos\pi + i\sin\pi = -1 + 0 = -1\)。\(\blacksquare\)

\textbf{推论 20.7}（棣莫弗公式）对任意实数 \(\theta\) 和正整数 \(n\)：

\[(\cos\theta + i\sin\theta)^n = \cos(n\theta) + i\sin(n\theta)\]

\textbf{证明}：由欧拉公式，\(\cos\theta + i\sin\theta = e^{i\theta}\)，故
\((\cos\theta + i\sin\theta)^n = e^{in\theta} = \cos(n\theta) + i\sin(n\theta)\)。\(\blacksquare\)

\bigskip\hrule\bigskip

\textbf{本章展望}：本章建立了三角恒等变换的完整体系和反三角函数。在下一章（Ch
21）中，我们将转向数列的研究。数列是定义在正整数集上的特殊函数，其离散性质使其在金融、计算机科学和物理学中有广泛应用。我们将系统研究等差数列和等比数列，并介绍数学归纳法这一重要的证明方法。

\bigskip\hrule\bigskip

\subsection{Ch 21 数列}\label{ch-21-ux6570ux5217}

\textbf{前置知识}：本章依赖 Ch
16（函数概念，将数列视为定义在正整数集上的特殊函数）以及 Ch
18（指数运算，为等比数列的通项公式做准备）。本章为 Ch
22（极限与连续）提供离散化的极限直觉。

\bigskip\hrule\bigskip

\subsubsection{21.1 数列的概念}\label{ux6570ux5217ux7684ux6982ux5ff5}

\textbf{定义 21.1}（数列）按一定顺序排列的一列数
\(a_1, a_2, a_3, \ldots\) 称为\textbf{数列}，简记为 \(\{a_n\}\)。其中
\(a_n\) 称为数列的\textbf{第 \(n\) 项}（或\textbf{通项}），\(n\)
称为\textbf{项数}。

\textbf{定义 21.2}（通项公式）如果数列 \(\{a_n\}\) 的第 \(n\) 项 \(a_n\)
与项数 \(n\) 之间的关系可以用一个公式 \(a_n = f(n)\)
表示，则称此公式为\textbf{通项公式}。

\textbf{定义 21.3}（递推公式）如果数列 \(\{a_n\}\) 的项之间的关系可以用
\(a_{n+1} = g(a_n, a_{n-1}, \ldots)\)
表示，则称此关系为\textbf{递推公式}。

\textbf{定义 21.4}（数列的前 \(n\) 项和）数列 \(\{a_n\}\) 的前 \(n\)
项之和：

\[S_n = a_1 + a_2 + \cdots + a_n = \sum_{k=1}^{n} a_k\]

\textbf{定理 21.1}（\(a_n\) 与 \(S_n\) 的关系）

\[a_n = \begin{cases} S_1 & n = 1 \\ S_n - S_{n-1} & n \geq 2 \end{cases}\]

\textbf{证明}：当 \(n = 1\) 时，\(S_1 = a_1\)。当 \(n \geq 2\)
时，\(S_n = a_1 + a_2 + \cdots + a_n\)，\(S_{n-1} = a_1 + a_2 + \cdots + a_{n-1}\)，两式相减得
\(S_n - S_{n-1} = a_n\)。\(\blacksquare\)

\bigskip\hrule\bigskip

\subsubsection{21.2 等差数列}\label{ux7b49ux5deeux6570ux5217}

\paragraph{21.2.1
定义与通项公式}\label{ux5b9aux4e49ux4e0eux901aux9879ux516cux5f0f}

\textbf{定义 21.5}（等差数列）如果数列 \(\{a_n\}\) 满足：对任意
\(n \geq 1\)，

\[a_{n+1} - a_n = d \quad (d \text{ 为常数})\]

则称 \(\{a_n\}\) 为\textbf{等差数列}，\(d\) 称为\textbf{公差}。

\textbf{定理 21.2}（等差数列通项公式）等差数列 \(\{a_n\}\) 的首项为
\(a_1\)，公差为 \(d\)，则：

\[a_n = a_1 + (n - 1)d\]

\textbf{证明}：对 \(n\) 进行归纳。

基础步：\(n = 1\) 时，\(a_1 = a_1 + 0 \cdot d\)，成立。

归纳步：假设 \(a_k = a_1 + (k-1)d\)。则
\(a_{k+1} = a_k + d = a_1 + (k-1)d + d = a_1 + kd\)。

由数学归纳法，结论成立。\(\blacksquare\)

\paragraph{\texorpdfstring{21.2.2 前 \(n\)
项和}{21.2.2 前 n 项和}}\label{ux524d-n-ux9879ux548c}

\textbf{定理 21.3}（等差数列前 \(n\) 项和）等差数列 \(\{a_n\}\) 的前
\(n\) 项和为：

\[S_n = \frac{n(a_1 + a_n)}{2} = na_1 + \frac{n(n-1)}{2}d\]

\textbf{证明（倒序相加法）}：正序写出 \(S_n\)：

\[S_n = a_1 + a_2 + a_3 + \cdots + a_n\]

倒序写出 \(S_n\)：

\[S_n = a_n + a_{n-1} + a_{n-2} + \cdots + a_1\]

两式相加：

\[2S_n = (a_1 + a_n) + (a_2 + a_{n-1}) + \cdots + (a_n + a_1)\]

由于
\(a_k + a_{n+1-k} = [a_1 + (k-1)d] + [a_1 + (n-k)d] = 2a_1 + (n-1)d = a_1 + a_n\)（常数），共有
\(n\) 对，故：

\[2S_n = n(a_1 + a_n), \quad S_n = \frac{n(a_1 + a_n)}{2}\]

\(\blacksquare\)

\paragraph{21.2.3 等差中项}\label{ux7b49ux5deeux4e2dux9879}

\textbf{定理 21.4} \(a, b, c\) 成等差数列的充要条件是
\(b = \frac{a + c}{2}\)。

\textbf{证明}：\(a, b, c\) 成等差数列 \(\iff\) \(b - a = c - b\)
\(\iff\) \(2b = a + c\) \(\iff\)
\(b = \frac{a + c}{2}\)。\(\blacksquare\)

\bigskip\hrule\bigskip

\subsubsection{21.3 等比数列}\label{ux7b49ux6bd4ux6570ux5217}

\paragraph{21.3.1
定义与通项公式}\label{ux5b9aux4e49ux4e0eux901aux9879ux516cux5f0f-1}

\textbf{定义 21.6}（等比数列）如果数列 \(\{a_n\}\) 满足：对任意
\(n \geq 1\)，\(a_n \neq 0\)，

\[\frac{a_{n+1}}{a_n} = q \quad (q \text{ 为非零常数})\]

则称 \(\{a_n\}\) 为\textbf{等比数列}，\(q\) 称为\textbf{公比}。

\textbf{定理 21.5}（等比数列通项公式）等比数列 \(\{a_n\}\) 的首项为
\(a_1\)（\(a_1 \neq 0\)），公比为 \(q\)（\(q \neq 0\)），则：

\[a_n = a_1 q^{n-1}\]

\textbf{证明}：对 \(n\) 进行归纳。

基础步：\(a_1 = a_1 q^0\)，成立。

归纳步：假设 \(a_k = a_1 q^{k-1}\)。则
\(a_{k+1} = a_k \cdot q = a_1 q^{k-1} \cdot q = a_1 q^k\)。

由数学归纳法，结论成立。\(\blacksquare\)

\paragraph{\texorpdfstring{21.3.2 前 \(n\)
项和}{21.3.2 前 n 项和}}\label{ux524d-n-ux9879ux548c-1}

\textbf{定理 21.6}（等比数列前 \(n\) 项和）等比数列 \(\{a_n\}\) 的首项为
\(a_1\)，公比为 \(q\)，则：

\[S_n = \begin{cases} na_1 & q = 1 \\ \dfrac{a_1(1 - q^n)}{1 - q} & q \neq 1 \end{cases}\]

\textbf{证明}：当 \(q = 1\)
时，\(a_n = a_1\)（常数列），\(S_n = na_1\)。

当 \(q \neq 1\) 时：

\[S_n = a_1 + a_1 q + a_1 q^2 + \cdots + a_1 q^{n-1} \quad \cdots (1)\]

\((1) \times q\)：

\[qS_n = a_1 q + a_1 q^2 + a_1 q^3 + \cdots + a_1 q^n \quad \cdots (2)\]

\((1) - (2)\)：

\[S_n - qS_n = a_1 - a_1 q^n\]

\[(1 - q)S_n = a_1(1 - q^n)\]

\[S_n = \frac{a_1(1 - q^n)}{1 - q}\]

\(\blacksquare\)

\paragraph{21.3.3 等比中项}\label{ux7b49ux6bd4ux4e2dux9879}

\textbf{定理 21.7} \(a, b, c\) 成等比数列（各项非零）的充要条件是
\(b^2 = ac\)。

\textbf{证明}：\(a, b, c\) 成等比数列 \(\iff\)
\(\frac{b}{a} = \frac{c}{b}\) \(\iff\) \(b^2 = ac\)。\(\blacksquare\)

\bigskip\hrule\bigskip

\subsubsection{21.4
数列求和方法}\label{ux6570ux5217ux6c42ux548cux65b9ux6cd5}

\paragraph{21.4.1 裂项相消法}\label{ux88c2ux9879ux76f8ux6d88ux6cd5}

\textbf{方法}：将通项 \(a_n\) 分解为两项之差
\(a_n = f(n) - f(n+1)\)，使得求和时中间项相互抵消。

\textbf{定理 21.8} \(\sum_{k=1}^{n} \frac{1}{k(k+1)} = \frac{n}{n+1}\)。

\textbf{证明}：利用部分分式分解：\(\frac{1}{k(k+1)} = \frac{1}{k} - \frac{1}{k+1}\)。

\[\sum_{k=1}^{n} \frac{1}{k(k+1)} = \sum_{k=1}^{n}\left(\frac{1}{k} - \frac{1}{k+1}\right) = 1 - \frac{1}{n+1} = \frac{n}{n+1}\]

\(\blacksquare\)

\paragraph{21.4.2 错位相减法}\label{ux9519ux4f4dux76f8ux51cfux6cd5}

\textbf{方法}：适用于形如 \(\sum a_n \cdot b_n\) 的求和，其中
\(\{a_n\}\) 是等差数列，\(\{b_n\}\) 是等比数列。

\paragraph{21.4.3
常见求和公式}\label{ux5e38ux89c1ux6c42ux548cux516cux5f0f}

\textbf{定理 21.9} \(\sum_{k=1}^{n} k = \frac{n(n+1)}{2}\)。

\textbf{证明}：这是等差数列（首项 \(1\)，公差 \(1\)，末项 \(n\)）的前
\(n\) 项和。由定理 21.3：

\[\sum_{k=1}^{n} k = \frac{n(1 + n)}{2} = \frac{n(n+1)}{2}\]

\(\blacksquare\)

\textbf{定理
21.10}（平方和公式）\(\sum_{k=1}^{n} k^2 = \frac{n(n+1)(2n+1)}{6}\)。

\textbf{证明}：利用恒等式 \((k+1)^3 - k^3 = 3k^2 + 3k + 1\)，对
\(k = 1, 2, \ldots, n\) 求和：

\[\sum_{k=1}^{n}[(k+1)^3 - k^3] = 3\sum_{k=1}^{n} k^2 + 3\sum_{k=1}^{n} k + n\]

左边是裂项和：\((n+1)^3 - 1^3\)。

\[(n+1)^3 - 1 = 3\sum_{k=1}^{n} k^2 + \frac{3n(n+1)}{2} + n\]

\[\sum_{k=1}^{n} k^2 = \frac{(n+1)^3 - 1 - \frac{3n(n+1)}{2} - n}{3} = \frac{n(n+1)(2n+1)}{6}\]

\(\blacksquare\)

\textbf{定理
21.11}（立方和公式）\(\sum_{k=1}^{n} k^3 = \left[\frac{n(n+1)}{2}\right]^2\)。

\textbf{证明}：对 \(n\) 用数学归纳法。

基础步：\(n = 1\)
时，\(1^3 = 1 = \left[\frac{1 \cdot 2}{2}\right]^2\)，成立。

归纳步：假设 \(n = m\) 时成立。则

\[\sum_{k=1}^{m+1} k^3 = \left[\frac{m(m+1)}{2}\right]^2 + (m+1)^3 = (m+1)^2\left[\frac{m^2}{4} + (m+1)\right]\]

\[= (m+1)^2 \cdot \frac{m^2 + 4m + 4}{4} = (m+1)^2 \cdot \frac{(m+2)^2}{4} = \left[\frac{(m+1)(m+2)}{2}\right]^2\]

由数学归纳法，结论成立。\(\blacksquare\)

\bigskip\hrule\bigskip

\subsubsection{21.5 递推数列}\label{ux9012ux63a8ux6570ux5217}

\paragraph{21.5.1
一阶线性递推数列}\label{ux4e00ux9636ux7ebfux6027ux9012ux63a8ux6570ux5217}

\textbf{定义 21.7}（一阶线性递推数列）形如
\(a_{n+1} = pa_n + q\)（\(p \neq 0\)）的递推关系称为\textbf{一阶线性递推}。

\textbf{定理 21.12} 设数列 \(\{a_n\}\) 满足
\(a_{n+1} = pa_n + q\)（\(p \neq 1\)），令
\(\lambda = \frac{q}{1-p}\)，则数列 \(\{a_n - \lambda\}\) 是以
\(a_1 - \lambda\) 为首项、\(p\) 为公比的等比数列。通项公式为：

\[a_n = \left(a_1 - \frac{q}{1-p}\right)p^{n-1} + \frac{q}{1-p}\]

\textbf{证明}：设 \(b_n = a_n - \lambda\)，则

\[b_{n+1} = a_{n+1} - \lambda = pa_n + q - \lambda = p(b_n + \lambda) + q - \lambda = pb_n + p\lambda + q - \lambda\]

由 \(\lambda = \frac{q}{1-p}\)，得 \(p\lambda + q = \lambda\)（即
\((1-p)\lambda = q\)），故 \(p\lambda + q - \lambda = 0\)。

因此 \(b_{n+1} = pb_n\)，\(\{b_n\}\) 是以 \(b_1 = a_1 - \lambda\)
为首项、\(p\) 为公比的等比数列。

\[b_n = (a_1 - \lambda)p^{n-1}\]

\[a_n = b_n + \lambda = (a_1 - \lambda)p^{n-1} + \lambda\]

\(\blacksquare\)

\bigskip\hrule\bigskip

\subsubsection{21.6 数学归纳法}\label{ux6570ux5b66ux5f52ux7eb3ux6cd5-1}

\paragraph{21.6.1
数学归纳法原理}\label{ux6570ux5b66ux5f52ux7eb3ux6cd5ux539fux7406}

\textbf{原理 21.1}（数学归纳法原理）设 \(P(n)\) 是关于正整数 \(n\)
的命题。如果：

\begin{enumerate}
\def\labelenumi{\arabic{enumi}.}
\tightlist
\item
  \textbf{基础步}：\(P(1)\) 成立；
\item
  \textbf{归纳步}：对任意正整数 \(k\)，由 \(P(k)\) 成立可推出 \(P(k+1)\)
  成立；
\end{enumerate}

则 \(P(n)\) 对所有正整数 \(n\) 成立。

\textbf{公理化说明}：数学归纳法的正确性基于正整数集的良序性公理：任何非空的正整数集合都有最小元素。若
\(P(n)\) 不对所有正整数成立，则存在使 \(P(n)\) 不成立的最小正整数
\(n_0\)。由基础步 \(n_0 \neq 1\)，故 \(n_0 > 1\)。由 \(n_0\)
的最小性，\(P(n_0 - 1)\) 成立。但由归纳步，\(P(n_0)\) 也应成立，矛盾。

\paragraph{21.6.2 强归纳法}\label{ux5f3aux5f52ux7eb3ux6cd5}

\textbf{原理 21.2}（强归纳法）设 \(P(n)\) 是关于正整数 \(n\)
的命题。如果对任意正整数 \(k\)，由 \(P(1), P(2), \ldots, P(k)\)
全部成立可推出 \(P(k+1)\) 成立，则 \(P(n)\) 对所有正整数 \(n\) 成立。

\textbf{注}
强归纳法与普通数学归纳法等价，但在某些问题中更为方便，因为归纳假设更强。

\bigskip\hrule\bigskip

\subsubsection{21.7
数列的极限判定}\label{ux6570ux5217ux7684ux6781ux9650ux5224ux5b9a}

\paragraph{21.7.1
收敛数列的性质}\label{ux6536ux655bux6570ux5217ux7684ux6027ux8d28}

\textbf{定理 21.13}（极限的保序性）设
\(\lim_{n \to \infty} a_n = A\)，\(\lim_{n \to \infty} b_n = B\)。

\begin{enumerate}
\def\labelenumi{\arabic{enumi}.}
\tightlist
\item
  若 \(A > B\)，则存在 \(N\)，当 \(n > N\) 时，\(a_n > b_n\)。
\item
  若存在 \(N_0\)，当 \(n > N_0\) 时 \(a_n \geq b_n\)，则 \(A \geq B\)。
\end{enumerate}

\textbf{证明}：

\begin{enumerate}
\def\labelenumi{(\arabic{enumi})}
\tightlist
\item
  取 \(\varepsilon = \frac{A - B}{2} > 0\)。存在 \(N_1, N_2\)，当
  \(n > N_1\) 时 \(|a_n - A| < \varepsilon\)，当 \(n > N_2\) 时
  \(|b_n - B| < \varepsilon\)。取 \(N = \max\{N_1, N_2\}\)，当 \(n > N\)
  时：
\end{enumerate}

\[a_n > A - \varepsilon = \frac{A + B}{2} = B + \varepsilon > b_n\]

\begin{enumerate}
\def\labelenumi{(\arabic{enumi})}
\setcounter{enumi}{1}
\tightlist
\item
  用反证法。若 \(A < B\)，由 (1) 存在 \(N\)，当 \(n > N\) 时
  \(a_n < b_n\)，与
  \(a_n \geq b_n\)（\(n > N_0\)）矛盾。\(\blacksquare\)
\end{enumerate}

\paragraph{21.7.2 子列与收敛}\label{ux5b50ux5217ux4e0eux6536ux655b}

\textbf{定义 21.8}（子列）设 \(\{a_n\}\)
是数列，\(n_1 < n_2 < n_3 < \cdots\) 是严格递增的正整数列，则
\(\{a_{n_k}\}\) 称为 \(\{a_n\}\) 的一个\textbf{子列}。

\textbf{定理 21.14} 若 \(\lim_{n \to \infty} a_n = L\)，则 \(\{a_n\}\)
的任何子列都收敛于 \(L\)。

\textbf{证明}：设 \(\{a_{n_k}\}\) 是子列。对于任意
\(\varepsilon > 0\)，存在 \(N\)，当 \(n > N\) 时
\(|a_n - L| < \varepsilon\)。由于 \(n_k \geq k\)（正整数严格递增），当
\(k > N\) 时 \(n_k > N\)，故
\(|a_{n_k} - L| < \varepsilon\)。\(\blacksquare\)

\textbf{推论 21.1} 若 \(\{a_n\}\) 有两个子列收敛于不同的极限，则
\(\{a_n\}\) 发散。

\textbf{证明}：若 \(\{a_n\}\) 收敛于 \(L\)，则所有子列都收敛于
\(L\)（定理 21.14），矛盾。\(\blacksquare\)

\bigskip\hrule\bigskip

\subsubsection{\texorpdfstring{21.8 数列极限的 \(\varepsilon\)-\(N\)
论证技巧}{21.8 数列极限的 \textbackslash varepsilon-N 论证技巧}}\label{ux6570ux5217ux6781ux9650ux7684-varepsilon-n-ux8bbaux8bc1ux6280ux5de7}

\paragraph{21.8.1
常用放缩技巧}\label{ux5e38ux7528ux653eux7f29ux6280ux5de7}

\textbf{引理 21.1} 设
\(a_n \geq 0\)，\(b_n \geq 0\)，\(a_n \leq b_n\)。若
\(\lim_{n \to \infty} b_n = 0\)，则 \(\lim_{n \to \infty} a_n = 0\)。

\textbf{证明}：这是夹逼定理（\(0 \leq a_n \leq b_n\)）的直接推论。\(\blacksquare\)

\textbf{引理 21.2}（伯努利不等式）设 \(x > -1\)，\(n\) 为正整数，则
\((1 + x)^n \geq 1 + nx\)。

\textbf{证明}：对 \(n\) 用数学归纳法。

\(n = 1\)：\(1 + x \geq 1 + x\)，成立。

\(n \Rightarrow n+1\)：假设 \((1+x)^n \geq 1 + nx\)。由于
\(1 + x > 0\)（即 \(x > -1\)），

\[(1+x)^{n+1} = (1+x)^n(1+x) \geq (1+nx)(1+x) = 1 + (n+1)x + nx^2 \geq 1 + (n+1)x\]

\(\blacksquare\)

\paragraph{21.8.2
重要数列极限}\label{ux91cdux8981ux6570ux5217ux6781ux9650}

\textbf{定理 21.15} \(\lim_{n \to \infty} \sqrt[n]{n} = 1\)。

\textbf{证明}：令 \(\sqrt[n]{n} = 1 + h_n\)（\(h_n > 0\)）。则
\(n = (1 + h_n)^n\)。由伯努利不等式：

\[n = (1 + h_n)^n \geq 1 + nh_n\]

所以 \(h_n \leq \frac{n - 1}{n} < 1\)，且
\(0 < h_n \leq \frac{1}{n} \cdot (n-1)\)。

进一步，对 \(n \geq 2\)，由二项式定理：

\[n = (1 + h_n)^n \geq \binom{n}{2}h_n^2 = \frac{n(n-1)}{2}h_n^2\]

所以
\(h_n^2 \leq \frac{2}{n-1}\)，\(h_n \leq \sqrt{\frac{2}{n-1}} \to 0\)。

由夹逼定理，\(h_n \to 0\)，故
\(\sqrt[n]{n} = 1 + h_n \to 1\)。\(\blacksquare\)

\textbf{定理 21.16}
\(\lim_{n \to \infty} \sqrt[n]{a} = 1\)（\(a > 0\)）。

\textbf{证明}：当 \(a = 1\) 时显然。当 \(a > 1\) 时，令
\(\sqrt[n]{a} = 1 + h_n\)（\(h_n > 0\)）。

\(a = (1 + h_n)^n \geq 1 + nh_n\)，故
\(h_n \leq \frac{a - 1}{n} \to 0\)。

由夹逼定理，\(\sqrt[n]{a} \to 1\)。当 \(0 < a < 1\)
时，\(\sqrt[n]{a} = \frac{1}{\sqrt[n]{1/a}} \to \frac{1}{1} = 1\)。\(\blacksquare\)

\bigskip\hrule\bigskip

\subsubsection{21.9
递推数列的进一步研究}\label{ux9012ux63a8ux6570ux5217ux7684ux8fdbux4e00ux6b65ux7814ux7a76}

\paragraph{21.9.1 不动点法}\label{ux4e0dux52a8ux70b9ux6cd5}

\textbf{定义 21.9}（不动点）设 \(f: D \to D\)，若 \(x^*\) 满足
\(f(x^*) = x^*\)，则称 \(x^*\) 为 \(f\) 的\textbf{不动点}。

\textbf{定理 21.17} 设数列 \(\{a_n\}\) 由 \(a_{n+1} = f(a_n)\)
定义，\(x^*\) 是 \(f\) 的不动点。若 \(\{a_n\}\) 收敛于 \(L\)，且 \(f\)
连续，则 \(L = x^*\)。

\textbf{证明}：由 \(a_{n+1} = f(a_n)\)
取极限：\(L = \lim a_{n+1} = \lim f(a_n) = f(\lim a_n) = f(L)\)。所以
\(L\) 是 \(f\) 的不动点。\(\blacksquare\)

\textbf{注}：此定理仅用于\textbf{已知}数列收敛时确定其极限值，本身不保证收敛性。实际应用中，需先验证收敛条件。一个常用的收敛性充分条件是：若
\(f\) 在不动点 \(x^*\) 处可导且 \(|f'(x^*)| < 1\)，则存在 \(x^*\)
的某个邻域，使得从该邻域内出发的迭代 \(a_{n+1} = f(a_n)\) 收敛于
\(x^*\)（压缩映射原理的特殊情形）。

\paragraph{21.9.2
斐波那契数列}\label{ux6590ux6ce2ux90a3ux5951ux6570ux5217}

\textbf{定义 21.10}（斐波那契数列）满足递推关系

\[F_1 = 1, \quad F_2 = 1, \quad F_{n+2} = F_{n+1} + F_n \quad (n \geq 1)\]

的数列 \(\{F_n\}\) 称为\textbf{斐波那契数列}。

\textbf{定理 21.18}（比内公式）斐波那契数列的通项公式为：

\[F_n = \frac{1}{\sqrt{5}}\left[\left(\frac{1 + \sqrt{5}}{2}\right)^n - \left(\frac{1 - \sqrt{5}}{2}\right)^n\right]\]

\textbf{证明}：特征方程为 \(x^2 = x + 1\)，即
\(x^2 - x - 1 = 0\)。两根为

\[\varphi = \frac{1 + \sqrt{5}}{2}, \quad \psi = \frac{1 - \sqrt{5}}{2}\]

设 \(F_n = A\varphi^n + B\psi^n\)。由初始条件 \(F_1 = 1\)，\(F_2 = 1\)：

\[A\varphi + B\psi = 1, \quad A\varphi^2 + B\psi^2 = 1\]

利用 \(\varphi^2 = \varphi + 1\)，\(\psi^2 = \psi + 1\)，第二个方程变为
\(A(\varphi + 1) + B(\psi + 1) = 1\)，即
\(A\varphi + B\psi + A + B = 1\)，代入第一个方程得 \(A + B = 0\)，即
\(B = -A\)。

由 \(A(\varphi - \psi) = 1\)，得
\(A = \frac{1}{\varphi - \psi} = \frac{1}{\sqrt{5}}\)。

因此 \(F_n = \frac{\varphi^n - \psi^n}{\sqrt{5}}\)。

验证递推：\(F_{n+1} + F_n = \frac{\varphi^{n+1} - \psi^{n+1} + \varphi^n - \psi^n}{\sqrt{5}} = \frac{\varphi^n(\varphi + 1) - \psi^n(\psi + 1)}{\sqrt{5}} = \frac{\varphi^n \cdot \varphi^2 - \psi^n \cdot \psi^2}{\sqrt{5}} = \frac{\varphi^{n+2} - \psi^{n+2}}{\sqrt{5}} = F_{n+2}\)。\(\blacksquare\)

\textbf{推论 21.2}
\(\lim_{n \to \infty} \frac{F_{n+1}}{F_n} = \varphi = \frac{1 + \sqrt{5}}{2} \approx 1.618...\)

\textbf{证明}：\(\frac{F_{n+1}}{F_n} = \frac{\varphi^{n+1} - \psi^{n+1}}{\varphi^n - \psi^n} = \frac{\varphi - \psi\left(\frac{\psi}{\varphi}\right)^n}{1 - \left(\frac{\psi}{\varphi}\right)^n}\)。

由于
\(|\psi/\varphi| = \frac{\sqrt{5}-1}{\sqrt{5}+1} < 1\)，\(\left(\frac{\psi}{\varphi}\right)^n \to 0\)，故极限为
\(\varphi\)。\(\blacksquare\)

\textbf{注}：\(\varphi = \frac{1 + \sqrt{5}}{2}\)
称为\textbf{黄金分割比}。

\subsubsection{21.10
数列的极限存在性的进一步判定}\label{ux6570ux5217ux7684ux6781ux9650ux5b58ux5728ux6027ux7684ux8fdbux4e00ux6b65ux5224ux5b9a}

\paragraph{21.10.1 施图茨定理}\label{ux65bdux56feux8328ux5b9aux7406}

\textbf{定理 21.19}（施图茨定理）设 \(\{b_n\}\)
是严格递增的正数列，\(\lim_{n \to \infty} b_n = +\infty\)。若
\(\lim_{n \to \infty} \frac{a_n - a_{n-1}}{b_n - b_{n-1}} = L\)（\(L\)
可以是 \(\pm\infty\)），则

\[\lim_{n \to \infty} \frac{a_n}{b_n} = L\]

\textbf{证明}：这是洛必达法则的离散版本。证明思路与洛必达法则类似，利用柯西中值定理的离散形式。

对于 \(\varepsilon > 0\)，存在 \(N_0\)，当 \(n > N_0\) 时
\(\left|\frac{a_n - a_{n-1}}{b_n - b_{n-1}} - L\right| < \varepsilon\)。

对 \(n > N_0\)，将 \(\frac{a_n - a_{N_0}}{b_n - b_{N_0}}\)
看作加权平均：

\[\frac{a_n - a_{N_0}}{b_n - b_{N_0}} = \frac{\sum_{k=N_0+1}^{n}(a_k - a_{k-1})}{\sum_{k=N_0+1}^{n}(b_k - b_{k-1})}\]

每一项 \(\frac{a_k - a_{k-1}}{b_k - b_{k-1}}\) 都在
\((L - \varepsilon, L + \varepsilon)\) 内，因此加权平均也在
\((L - \varepsilon, L + \varepsilon)\) 内（因为
\(b_k - b_{k-1} > 0\)）。

由此可证 \(\frac{a_n}{b_n} \to L\)。\(\blacksquare\)

\paragraph{21.10.2
柯西收敛准则}\label{ux67efux897fux6536ux655bux51c6ux5219}

\textbf{定理 21.20}（柯西收敛准则）数列 \(\{a_n\}\) 收敛的充要条件是：

\[\forall \varepsilon > 0, \exists N, \forall m, n > N: |a_m - a_n| < \varepsilon\]

满足此条件的数列称为\textbf{柯西列}（或\textbf{基本列}）。

\textbf{证明概要}：

必要性：设 \(a_n \to L\)。对于 \(\varepsilon > 0\)，存在 \(N\)，当
\(n > N\) 时 \(|a_n - L| < \frac{\varepsilon}{2}\)。当 \(m, n > N\) 时：

\[|a_m - a_n| \leq |a_m - L| + |L - a_n| < \frac{\varepsilon}{2} + \frac{\varepsilon}{2} = \varepsilon\]

充分性：设 \((a_n)\) 是柯西列。

\textbf{第一步}（有界性）：取 \(\varepsilon = 1\)，存在 \(N\)，当
\(m, n > N\) 时 \(|a_m - a_n| < 1\)。特别地，对所有
\(n > N\)，\(|a_n| \leq |a_{N+1}| + 1\)。令
\(M = \max\{|a_1|, \ldots, |a_{N+1}|, |a_{N+1}| + 1\}\)，则
\(|a_n| \leq M\)。

\textbf{第二步}（存在收敛子列）：由闭区间套定理（定理
22.10，其等价形式），有界数列存在收敛子列
\(a_{n_k} \to L\)（此结论的完整证明见定理 22.23）。

\textbf{第三步}（整个数列收敛于 \(L\)）：对
\(\varepsilon > 0\)，由柯西条件存在 \(N_1\) 使 \(m, n > N_1\) 时
\(|a_m - a_n| < \varepsilon/2\)。由 \(a_{n_k} \to L\)，存在 \(K\) 使
\(k > K\) 时 \(|a_{n_k} - L| < \varepsilon/2\)。取 \(k > K\) 使得
\(n_k > N_1\)，则对任意 \(n > N_1\)：

\[|a_n - L| \leq |a_n - a_{n_k}| + |a_{n_k} - L| < \frac{\varepsilon}{2} + \frac{\varepsilon}{2} = \varepsilon\]

故 \(a_n \to L\)。\(\blacksquare\)

\subsubsection{21.11
数学归纳法的典型应用}\label{ux6570ux5b66ux5f52ux7eb3ux6cd5ux7684ux5178ux578bux5e94ux7528}

\paragraph{21.11.1
用归纳法证明不等式}\label{ux7528ux5f52ux7eb3ux6cd5ux8bc1ux660eux4e0dux7b49ux5f0f}

\textbf{定理 21.21}（伯努利不等式的推广）设
\(x_1, x_2, \ldots, x_n > -1\) 且同号，则

\[(1 + x_1)(1 + x_2)\cdots(1 + x_n) \geq 1 + x_1 + x_2 + \cdots + x_n\]

\textbf{证明}：对 \(n\) 用数学归纳法。

\(n = 1\) 时显然成立。

\(n \Rightarrow n+1\)：假设对 \(n\) 成立。设
\(x_1, x_2, \ldots, x_{n+1}\) 同号且大于 \(-1\)。

\[(1+x_1)\cdots(1+x_n)(1+x_{n+1}) \geq (1 + x_1 + \cdots + x_n)(1 + x_{n+1})\]

\[= 1 + x_1 + \cdots + x_{n+1} + (x_1 + \cdots + x_n)x_{n+1}\]

由于 \(x_1, \ldots, x_n, x_{n+1}\)
同号，\((x_1 + \cdots + x_n)x_{n+1} \geq 0\)，故

\[\geq 1 + x_1 + \cdots + x_{n+1}\]

\(\blacksquare\)

\paragraph{21.11.2
用归纳法证明整除性}\label{ux7528ux5f52ux7eb3ux6cd5ux8bc1ux660eux6574ux9664ux6027}

\textbf{定理 21.22} 对任意正整数 \(n\)，\(6 | n(n+1)(2n+1)\)。

\textbf{证明}：\(n(n+1)(2n+1) = 6\sum_{k=1}^{n}k^2\)，故
\(6 | n(n+1)(2n+1)\)。

或者直接用归纳法：

\(n = 1\)：\(1 \cdot 2 \cdot 3 = 6\)，\(6 | 6\)。

\(n \Rightarrow n+1\)：\((n+1)(n+2)(2n+3) - n(n+1)(2n+1) = (n+1)[(n+2)(2n+3) - n(2n+1)] = (n+1)(6n+6) = 6(n+1)^2\)。

由归纳假设 \(6 | n(n+1)(2n+1)\)，又 \(6 | 6(n+1)^2\)，故
\(6 | (n+1)(n+2)(2n+3)\)。\(\blacksquare\)

\paragraph{21.11.3
用归纳法证明集合计数}\label{ux7528ux5f52ux7eb3ux6cd5ux8bc1ux660eux96c6ux5408ux8ba1ux6570}

\textbf{定理 21.23}（德摩根律的推广）设 \(A_1, A_2, \ldots, A_n\) 是全集
\(U\) 的子集，则

\[\left(\bigcup_{k=1}^{n} A_k\right)^c = \bigcap_{k=1}^{n} A_k^c, \quad \left(\bigcap_{k=1}^{n} A_k\right)^c = \bigcup_{k=1}^{n} A_k^c\]

\textbf{证明}：对 \(n\) 用数学归纳法。

\(n = 2\) 时已在 Ch 3 中证明。

\(n \Rightarrow n+1\)：

\[\left(\bigcup_{k=1}^{n+1} A_k\right)^c = \left(\left(\bigcup_{k=1}^{n} A_k\right) \cup A_{n+1}\right)^c = \left(\bigcup_{k=1}^{n} A_k\right)^c \cap A_{n+1}^c = \bigcap_{k=1}^{n} A_k^c \cap A_{n+1}^c = \bigcap_{k=1}^{n+1} A_k^c\]

\(\blacksquare\)

\bigskip\hrule\bigskip

\textbf{本章展望}：本章研究了数列的基本理论。在下一章（Ch
22）中，我们将从数列过渡到极限理论------当 \(n \to \infty\) 时数列
\(\{a_n\}\) 的行为如何？极限理论是整个分析学的基石，我们将用
\(\varepsilon\)-\(N\) 语言严格定义数列极限，用
\(\varepsilon\)-\(\delta\)
语言严格定义函数极限，并建立连续函数的介值定理和最值定理。

\bigskip\hrule\bigskip

\subsection{Ch 22
极限与连续}\label{ch-22-ux6781ux9650ux4e0eux8fdeux7eed}

\textbf{前置知识}：本章依赖 Ch 16（函数概念）、Ch
18（指数函数与对数函数）、Ch
21（数列概念与数学归纳法）。极限理论是整个分析学的基石，是连接离散（数列）与连续（函数）的桥梁。

\bigskip\hrule\bigskip

\subsubsection{22.1 数列极限}\label{ux6570ux5217ux6781ux9650}

\paragraph{22.1.1
数列极限的直观概念}\label{ux6570ux5217ux6781ux9650ux7684ux76f4ux89c2ux6982ux5ff5}

当我们说数列 \(\{a_n\}\) \textbf{趋于} \(L\)，直觉上是指：当 \(n\)
足够大时，\(a_n\) 可以任意接近 \(L\)。下面的 \(\varepsilon\)-\(N\)
定义将这一直觉精确化。

\paragraph{\texorpdfstring{22.1.2 \(\varepsilon\)-\(N\)
定义}{22.1.2 \textbackslash varepsilon-N 定义}}\label{varepsilon-n-ux5b9aux4e49}

\textbf{定义 22.1}（数列极限）设 \(\{a_n\}\)
是一个数列，\(L \in \mathbb{R}\)。若

\[\forall \varepsilon > 0, \exists N \in \mathbb{N}^+, \forall n > N: |a_n - L| < \varepsilon\]

则称数列 \(\{a_n\}\) \textbf{收敛于} \(L\)，记作
\(\lim_{n \to \infty} a_n = L\) 或 \(a_n \to L\)（\(n \to \infty\)）。

若数列不收敛于任何实数，则称该数列\textbf{发散}。

\textbf{几何解释}：\(\lim_{n \to \infty} a_n = L\) 意味着，对于 \(L\)
的任意 \(\varepsilon\) 邻域
\((L - \varepsilon, L + \varepsilon)\)，从某项 \(N\)
之后的所有项都落在这个邻域内。

\paragraph{22.1.3
极限的唯一性}\label{ux6781ux9650ux7684ux552fux4e00ux6027}

\textbf{定理 22.1}（极限的唯一性）若 \(\lim_{n \to \infty} a_n = L_1\)
且 \(\lim_{n \to \infty} a_n = L_2\)，则 \(L_1 = L_2\)。

\textbf{证明}：用反证法。假设 \(L_1 \neq L_2\)，令
\(\varepsilon = \frac{|L_1 - L_2|}{2} > 0\)。

由 \(\lim_{n \to \infty} a_n = L_1\)，存在 \(N_1\)，使得当 \(n > N_1\)
时，\(|a_n - L_1| < \varepsilon\)。

由 \(\lim_{n \to \infty} a_n = L_2\)，存在 \(N_2\)，使得当 \(n > N_2\)
时，\(|a_n - L_2| < \varepsilon\)。

取 \(n > \max\{N_1, N_2\}\)，则

\[|L_1 - L_2| = |L_1 - a_n + a_n - L_2| \leq |a_n - L_1| + |a_n - L_2| < \varepsilon + \varepsilon = 2\varepsilon = |L_1 - L_2|\]

矛盾！因此 \(L_1 = L_2\)。\(\blacksquare\)

\paragraph{22.1.4
收敛数列的有界性}\label{ux6536ux655bux6570ux5217ux7684ux6709ux754cux6027}

\textbf{定理 22.2}（收敛数列的有界性）若
\(\lim_{n \to \infty} a_n = L\)，则 \(\{a_n\}\) 是有界数列。

\textbf{证明}：取 \(\varepsilon = 1\)，由极限定义，存在 \(N\)，使得当
\(n > N\) 时，\(|a_n - L| < 1\)，即 \(|a_n| < |L| + 1\)。

令 \(M = \max\{|a_1|, |a_2|, \ldots, |a_{N}|, |L| + 1\}\)，则
\(\forall n \in \mathbb{N}^+, |a_n| \leq M\)。\(\blacksquare\)

\textbf{注}：逆命题不成立，即有界数列不一定收敛。例如 \(a_n = (-1)^n\)
有界但发散。

\paragraph{22.1.5
极限的四则运算}\label{ux6781ux9650ux7684ux56dbux5219ux8fd0ux7b97}

\textbf{定理 22.3}（极限的四则运算）设
\(\lim_{n \to \infty} a_n = A\)，\(\lim_{n \to \infty} b_n = B\)。则：

\begin{enumerate}
\def\labelenumi{\arabic{enumi}.}
\tightlist
\item
  \(\lim_{n \to \infty} (a_n + b_n) = A + B\)
\item
  \(\lim_{n \to \infty} (a_n - b_n) = A - B\)
\item
  \(\lim_{n \to \infty} (a_n \cdot b_n) = A \cdot B\)
\item
  若 \(B \neq 0\)，\(\lim_{n \to \infty} \frac{a_n}{b_n} = \frac{A}{B}\)
\end{enumerate}

\textbf{证明}：

\begin{enumerate}
\def\labelenumi{(\arabic{enumi})}
\tightlist
\item
  对于任意 \(\varepsilon > 0\)，由 \(\lim a_n = A\)，存在
  \(N_1\)，使得当 \(n > N_1\)
  时，\(|a_n - A| < \frac{\varepsilon}{2}\)。
\end{enumerate}

由 \(\lim b_n = B\)，存在 \(N_2\)，使得当 \(n > N_2\)
时，\(|b_n - B| < \frac{\varepsilon}{2}\)。

取 \(N = \max\{N_1, N_2\}\)，当 \(n > N\) 时：

\[|(a_n + b_n) - (A + B)| \leq |a_n - A| + |b_n - B| < \frac{\varepsilon}{2} + \frac{\varepsilon}{2} = \varepsilon\]

\begin{enumerate}
\def\labelenumi{(\arabic{enumi})}
\setcounter{enumi}{2}
\tightlist
\item
  首先由定理 22.2，存在 \(M_1 > 0\) 使得 \(|b_n| \leq M_1\)。令
  \(M = \max\{M_1, |A|, 1\}\)。
\end{enumerate}

对于任意 \(\varepsilon > 0\)，存在 \(N_1\) 使得当 \(n > N_1\) 时
\(|a_n - A| < \frac{\varepsilon}{2M}\)；存在 \(N_2\) 使得当 \(n > N_2\)
时 \(|b_n - B| < \frac{\varepsilon}{2M}\)。

当 \(n > \max\{N_1, N_2\}\) 时：

\[|a_n b_n - AB| = |a_n(b_n - B) + B(a_n - A)| \leq |a_n||b_n - B| + |B||a_n - A| < M \cdot \frac{\varepsilon}{2M} + M \cdot \frac{\varepsilon}{2M} = \varepsilon\]

\begin{enumerate}
\def\labelenumi{(\arabic{enumi})}
\setcounter{enumi}{3}
\tightlist
\item
  首先证明 \(\lim_{n \to \infty} \frac{1}{b_n} = \frac{1}{B}\)。由于
  \(B \neq 0\)，取 \(\varepsilon = \frac{|B|}{2}\)，存在 \(N_0\) 使得当
  \(n > N_0\) 时 \(|b_n - B| < \frac{|B|}{2}\)，即
  \(|b_n| > \frac{|B|}{2}\)。
\end{enumerate}

对于任意 \(\varepsilon > 0\)，存在 \(N\) 使得当 \(n > N\) 时
\(|b_n - B| < \frac{|B|^2 \varepsilon}{2}\)。

当 \(n > \max\{N_0, N\}\) 时：

\[\left|\frac{1}{b_n} - \frac{1}{B}\right| = \frac{|B - b_n|}{|b_n||B|} < \frac{|B - b_n|}{\frac{|B|}{2} \cdot |B|} = \frac{2|B - b_n|}{|B|^2} < \varepsilon\]

再由
(3)，\(\frac{a_n}{b_n} = a_n \cdot \frac{1}{b_n} \to A \cdot \frac{1}{B} = \frac{A}{B}\)。\(\blacksquare\)

\paragraph{22.1.6 夹逼定理}\label{ux5939ux903cux5b9aux7406}

\textbf{定理 22.4}（夹逼定理）设 \(\{a_n\}\)，\(\{b_n\}\)，\(\{c_n\}\)
是三个数列。若存在 \(N_0 \in \mathbb{N}^+\)，使得

\[\forall n > N_0: a_n \leq b_n \leq c_n\]

且 \(\lim_{n \to \infty} a_n = \lim_{n \to \infty} c_n = L\)，则
\(\lim_{n \to \infty} b_n = L\)。

\textbf{证明}：对于任意 \(\varepsilon > 0\)：

由 \(\lim a_n = L\)，存在 \(N_1\)，使得当 \(n > N_1\)
时，\(L - \varepsilon < a_n\)。

由 \(\lim c_n = L\)，存在 \(N_2\)，使得当 \(n > N_2\)
时，\(c_n < L + \varepsilon\)。

取 \(N = \max\{N_0, N_1, N_2\}\)，当 \(n > N\) 时：

\[L - \varepsilon < a_n \leq b_n \leq c_n < L + \varepsilon\]

所以 \(|b_n - L| < \varepsilon\)，即
\(\lim_{n \to \infty} b_n = L\)。\(\blacksquare\)

\paragraph{22.1.7
单调有界定理}\label{ux5355ux8c03ux6709ux754cux5b9aux7406}

\textbf{定理 22.5}（单调有界定理）

\begin{enumerate}
\def\labelenumi{\arabic{enumi}.}
\tightlist
\item
  单调递增且有上界的数列必收敛。
\item
  单调递减且有下界的数列必收敛。
\end{enumerate}

\textbf{证明}：

\begin{enumerate}
\def\labelenumi{(\arabic{enumi})}
\tightlist
\item
  设 \(\{a_n\}\)
  单调递增且有上界。由实数的完备性（确界存在定理），\(\{a_n\}\)
  有上确界，设为 \(L = \sup\{a_n\}\)。
\end{enumerate}

我们证明 \(\lim_{n \to \infty} a_n = L\)。

对于任意 \(\varepsilon > 0\)，由上确界的定义，\(L - \varepsilon\) 不是
\(\{a_n\}\) 的上界，所以存在 \(N\) 使得 \(a_{N} > L - \varepsilon\)。

由于 \(\{a_n\}\) 单调递增，当 \(n > N\)
时，\(a_n \geq a_{N} > L - \varepsilon\)。

又因为 \(L\) 是上界，\(a_n \leq L < L + \varepsilon\)。

所以当 \(n > N\) 时，\(L - \varepsilon < a_n < L + \varepsilon\)，即
\(|a_n - L| < \varepsilon\)。

\begin{enumerate}
\def\labelenumi{(\arabic{enumi})}
\setcounter{enumi}{1}
\tightlist
\item
  类似证明。\(\blacksquare\)
\end{enumerate}

\paragraph{22.1.8 重要极限}\label{ux91cdux8981ux6781ux9650}

\textbf{定理 22.6} \(\lim_{n \to \infty} \frac{1}{n} = 0\)

\textbf{证明}：对于任意 \(\varepsilon > 0\)，取
\(N = \lceil \frac{1}{\varepsilon} \rceil\)。当 \(n > N\)
时，\(n > \frac{1}{\varepsilon}\)，即
\(\frac{1}{n} < \varepsilon\)。\(\blacksquare\)

\textbf{定理 22.7} 若 \(|q| < 1\)，则 \(\lim_{n \to \infty} q^n = 0\)

\textbf{证明}：当 \(q = 0\) 时显然成立。当 \(0 < |q| < 1\) 时，令
\(|q| = \frac{1}{1 + h}\)（\(h > 0\)）。由伯努利不等式，\((1 + h)^n \geq 1 + nh > nh\)。

所以 \(|q^n| = \frac{1}{(1 + h)^n} < \frac{1}{nh}\)。对于任意
\(\varepsilon > 0\)，取 \(N = \lceil \frac{1}{h\varepsilon} \rceil\)，当
\(n > N\) 时 \(|q^n| < \varepsilon\)。\(\blacksquare\)

\textbf{定理 22.8}
\(\lim_{n \to \infty}\left(1 + \frac{1}{n}\right)^n = e\)

其中 \(e \approx 2.71828...\) 是自然对数的底。

\textbf{证明}：分两步证明。

\textbf{第一步}：证明
\(\{a_n\} = \left\{\left(1 + \frac{1}{n}\right)^n\right\}\) 单调递增。

由二项式定理，

\[a_n = \sum_{k=0}^{n} \frac{1}{k!} \cdot 1 \cdot \left(1 - \frac{1}{n}\right) \cdots \left(1 - \frac{k-1}{n}\right)\]

类似地，

\[a_{n+1} = \sum_{k=0}^{n+1} \frac{1}{k!} \cdot 1 \cdot \left(1 - \frac{1}{n+1}\right) \cdots \left(1 - \frac{k-1}{n+1}\right)\]

因为
\(1 - \frac{j}{n} < 1 - \frac{j}{n+1}\)（\(j = 1, 2, \ldots, k-1\)），所以
\(a_{n+1}\) 的每一项都不小于 \(a_n\) 的对应项，且 \(a_{n+1}\)
多了一项。因此 \(a_{n+1} > a_n\)。

\textbf{第二步}：证明 \(\{a_n\}\) 有上界。

\[a_n < \sum_{k=0}^{n} \frac{1}{k!} < 1 + 1 + \frac{1}{2} + \frac{1}{2^2} + \cdots = 3\]

（利用了 \(\frac{1}{k!} < \frac{1}{2^{k-1}}\) 对 \(k \geq 2\) 成立。）

由单调有界定理，\(\{a_n\}\) 收敛。定义此极限为 \(e\)。\(\blacksquare\)

\bigskip\hrule\bigskip

\subsubsection{22.2 函数极限}\label{ux51fdux6570ux6781ux9650}

\paragraph{\texorpdfstring{22.2.1 \(x \to a\)
时的极限定义}{22.2.1 x \textbackslash to a 时的极限定义}}\label{x-to-a-ux65f6ux7684ux6781ux9650ux5b9aux4e49}

\textbf{定义 22.2}（函数极限）设 \(f\) 在 \(a\) 的某个去心邻域
\(\mathring{U}(a, \delta_0)\) 上有定义，\(L \in \mathbb{R}\)。若

\[\forall \varepsilon > 0, \exists \delta > 0, \forall x: 0 < |x - a| < \delta \Rightarrow |f(x) - L| < \varepsilon\]

则称 \(f(x)\) 当 \(x \to a\) 时\textbf{收敛于} \(L\)，记作
\(\lim_{x \to a} f(x) = L\)。

\textbf{注意}：定义中要求 \(0 < |x - a|\)，即 \(x \neq a\)，说明 \(f\)
在 \(a\) 点的极限与 \(f(a)\) 的值无关。

\paragraph{22.2.2
左极限和右极限}\label{ux5de6ux6781ux9650ux548cux53f3ux6781ux9650}

\textbf{定义 22.3}（左极限和右极限）

\[\lim_{x \to a^-} f(x) = L \iff \forall \varepsilon > 0, \exists \delta > 0, \forall x \in (a - \delta, a): |f(x) - L| < \varepsilon\]

\[\lim_{x \to a^+} f(x) = L \iff \forall \varepsilon > 0, \exists \delta > 0, \forall x \in (a, a + \delta): |f(x) - L| < \varepsilon\]

\textbf{定理 22.9} \(\lim_{x \to a} f(x) = L\) 当且仅当
\(\lim_{x \to a^-} f(x) = \lim_{x \to a^+} f(x) = L\)。

\textbf{证明}：必要性显然。充分性：对于任意 \(\varepsilon > 0\)，存在
\(\delta_1, \delta_2 > 0\) 使得在 \((a - \delta_1, a)\) 和
\((a, a + \delta_2)\) 上分别成立 \(|f(x) - L| < \varepsilon\)。取
\(\delta = \min\{\delta_1, \delta_2\}\)，当 \(0 < |x - a| < \delta\) 时
\(|f(x) - L| < \varepsilon\)。\(\blacksquare\)

\paragraph{22.2.3
函数极限的四则运算}\label{ux51fdux6570ux6781ux9650ux7684ux56dbux5219ux8fd0ux7b97}

\textbf{定理 22.10} 设
\(\lim_{x \to a} f(x) = A\)，\(\lim_{x \to a} g(x) = B\)。则四则运算成立，证明与定理
22.3 类似。

\textbf{证明}：由函数极限的序列刻画（Heine定理），将函数极限转化为序列极限，应用定理
22.3（序列极限四则运算）即得。\(\blacksquare\)

\(\varepsilon\)-\(\delta\) 证明（加法）：设
\(\lim_{x \to a} f(x) = L\)，\(\lim_{x \to a} g(x) = M\)。对任意
\(\varepsilon > 0\)，取 \(\varepsilon_1 = \varepsilon/2\)。由 \(f\)
的极限，存在 \(\delta_1 > 0\) 使得 \(0 < |x - a| < \delta_1\) 时
\(|f(x) - L| < \varepsilon/2\)。由 \(g\) 的极限，存在 \(\delta_2 > 0\)
使得 \(0 < |x - a| < \delta_2\) 时 \(|g(x) - M| < \varepsilon/2\)。取
\(\delta = \min(\delta_1, \delta_2)\)，则 \(0 < |x - a| < \delta\) 时
\(|(f(x) + g(x)) - (L + M)| \leq |f(x) - L| + |g(x) - M| < \varepsilon/2 + \varepsilon/2 = \varepsilon\)。

\paragraph{22.2.4
函数极限的夹逼定理}\label{ux51fdux6570ux6781ux9650ux7684ux5939ux903cux5b9aux7406}

\textbf{定理 22.11}（夹逼定理）设 \(f(x) \leq g(x) \leq h(x)\) 在 \(a\)
的某个去心邻域内成立，且
\(\lim_{x \to a} f(x) = \lim_{x \to a} h(x) = L\)，则
\(\lim_{x \to a} g(x) = L\)。

\textbf{证明}：与定理 22.4 类似。\(\varepsilon\)-\(\delta\) 证明：设
\(g(x) \leq f(x) \leq h(x)\) 在 \(a\) 的去心邻域内成立，且
\(\lim_{x \to a} g(x) = \lim_{x \to a} h(x) = L\)。对任意
\(\varepsilon > 0\)，存在 \(\delta_1 > 0\) 使得
\(0 < |x - a| < \delta_1\) 时 \(|g(x) - L| < \varepsilon\)，即
\(L - \varepsilon < g(x)\)。存在 \(\delta_2 > 0\) 使得
\(0 < |x - a| < \delta_2\) 时 \(|h(x) - L| < \varepsilon\)，即
\(h(x) < L + \varepsilon\)。取 \(\delta = \min(\delta_1, \delta_2)\)，则
\(0 < |x - a| < \delta\) 时
\(L - \varepsilon < g(x) \leq f(x) \leq h(x) < L + \varepsilon\)，故
\(|f(x) - L| < \varepsilon\)。\(\blacksquare\)

\paragraph{\texorpdfstring{22.2.5 重要极限
\(\lim_{x \to 0} \frac{\sin x}{x} = 1\)}{22.2.5 重要极限 \textbackslash lim\_\{x \textbackslash to 0\} \textbackslash frac\{\textbackslash sin x\}\{x\} = 1}}\label{ux91cdux8981ux6781ux9650-lim_x-to-0-fracsin-xx-1}

\textbf{定理 22.12} \(\lim_{x \to 0} \frac{\sin x}{x} = 1\)

\textbf{证明}：当 \(0 < x < \frac{\pi}{2}\) 时，考虑单位圆中的扇形
\(OAB\) 和三角形 \(OAB\) 及三角形 \(OAC\)，其中
\(A = (1, 0)\)，\(B = (\cos x, \sin x)\)，\(C = (1, \tan x)\)。

扇形面积 \(= \frac{x}{2}\)，三角形 \(OAB\) 面积
\(= \frac{1}{2}\sin x\)，三角形 \(OAC\) 面积 \(= \frac{1}{2}\tan x\)。

由面积关系：

\[\frac{1}{2}\sin x < \frac{x}{2} < \frac{1}{2}\tan x\]

即 \(\sin x < x < \frac{\sin x}{\cos x}\)。

由 \(\sin x < x\) 得 \(\frac{\sin x}{x} < 1\)。由
\(x < \frac{\sin x}{\cos x}\) 得 \(\cos x < \frac{\sin x}{x}\)。

因此当 \(0 < x < \frac{\pi}{2}\) 时：\(\cos x < \frac{\sin x}{x} < 1\)。

由偶函数性质，当 \(-\frac{\pi}{2} < x < 0\) 时同样成立。

现在证明 \(\lim_{x \to 0} \cos x = 1\)。利用半角公式
\(1 - \cos x = 2\sin^2\frac{x}{2}\)，以及前面已证的不等式
\(|\sin t| < |t|\)（当 \(0 < |t| < \frac{\pi}{2}\)
时，由几何推导得到；当 \(|t| \geq \frac{\pi}{2}\) 时
\(|\sin t| \leq 1 < \frac{\pi}{2} \leq |t|\)，故 \(|\sin t| \leq |t|\)
对所有 \(t\) 成立），取 \(t = \frac{x}{2}\) 得
\(\left|\sin\frac{x}{2}\right| \leq \left|\frac{x}{2}\right|\)。因此

\[|\cos x - 1| = 1 - \cos x = 2\sin^2\frac{x}{2} \leq 2 \cdot \left(\frac{x}{2}\right)^2 = \frac{x^2}{2}\]

对于任意 \(\varepsilon > 0\)，取 \(\delta = \sqrt{2\varepsilon}\)，当
\(0 < |x| < \delta\) 时 \(|\cos x - 1| < \varepsilon\)，故
\(\lim_{x \to 0} \cos x = 1\)。

由夹逼定理，\(\lim_{x \to 0} \frac{\sin x}{x} = 1\)。\(\blacksquare\)

\bigskip\hrule\bigskip

\subsubsection{22.3 连续函数}\label{ux8fdeux7eedux51fdux6570}

\paragraph{22.3.1 连续的定义}\label{ux8fdeux7eedux7684ux5b9aux4e49}

\textbf{定义 22.4}（连续）设 \(f\) 在 \(a\) 的某个邻域内有定义。若

\[\lim_{x \to a} f(x) = f(a)\]

则称 \(f\) 在点 \(a\) 处\textbf{连续}。

等价地，\(f\) 在 \(a\) 处连续当且仅当：

\[\forall \varepsilon > 0, \exists \delta > 0, \forall x: |x - a| < \delta \Rightarrow |f(x) - f(a)| < \varepsilon\]

\textbf{定义 22.5}（区间上的连续）若 \(f\) 在区间 \(I\)
上的每一点都连续，则称 \(f\) 在 \(I\) 上\textbf{连续}。

\paragraph{22.3.2
连续函数的运算}\label{ux8fdeux7eedux51fdux6570ux7684ux8fd0ux7b97}

\textbf{定理 22.13}（连续函数的四则运算）设 \(f\) 和 \(g\) 都在点 \(a\)
处连续，则 \(f + g\)，\(f - g\)，\(f \cdot g\) 也在 \(a\) 处连续。若
\(g(a) \neq 0\)，则 \(\frac{f}{g}\) 也在 \(a\) 处连续。

\textbf{证明}：由极限的四则运算和连续的定义直接可得。\(\blacksquare\)

\textbf{定理 22.14}（复合函数的连续性）设 \(g\) 在 \(a\) 处连续，\(f\)
在 \(g(a)\) 处连续，则 \(f \circ g\) 在 \(a\) 处连续。

\textbf{证明}：由 \(g\) 在 \(a\)
处连续，\(\lim_{x \to a} g(x) = g(a)\)。又由 \(f\) 在 \(g(a)\)
处连续，根据连续函数的定义（或函数极限的复合定理在连续情形下的形式），\(\lim_{t \to g(a)} f(t) = f(g(a))\)。因此

\[\lim_{x \to a} f(g(x)) = f\left(\lim_{x \to a} g(x)\right) = f(g(a))\]

其中第一步利用了 \(f\) 在 \(g(a)\)
处的连续性（而非仅极限存在），使得极限值可以直接代入
\(f\)。\(\blacksquare\)

\paragraph{22.3.3 介值定理}\label{ux4ecbux503cux5b9aux7406}

\textbf{定理 22.15}（介值定理）设 \(f\) 在 \([a, b]\) 上连续，且
\(f(a) \neq f(b)\)。则对于 \(f(a)\) 和 \(f(b)\) 之间的任意实数
\(c\)，存在 \(\xi \in (a, b)\)，使得 \(f(\xi) = c\)。

\textbf{证明}：不妨设 \(f(a) < c < f(b)\)。定义集合

\[A = \{x \in [a, b] \mid f(x) < c\}\]

则 \(a \in A\)，故 \(A\) 非空。又 \(A\) 有上界
\(b\)，由确界存在定理，\(\xi = \sup A\) 存在，且 \(\xi \in [a, b]\)。

\textbf{证明 \(f(\xi) \leq c\)}：对于任意
\(n \in \mathbb{N}^+\)，由上确界的定义，存在 \(x_n \in A\) 使得
\(\xi - \frac{1}{n} < x_n \leq \xi\)。由夹逼定理，\(x_n \to \xi\)。因为
\(f(x_n) < c\) 且 \(f\) 连续，所以
\(f(\xi) = \lim_{n \to \infty} f(x_n) \leq c\)。

\textbf{证明 \(f(\xi) \geq c\)}：先证 \(\xi < b\)。若
\(\xi = b\)，则由上确界定义，存在 \(x_n \in A\) 使得 \(x_n \to b\)。由
\(f\) 的连续性，\(f(b) = \lim_{n \to \infty} f(x_n) \leq c\)（因为
\(f(x_n) < c\)）。但这与假设 \(f(b) > c\) 矛盾。故 \(\xi < b\)。

由 \(\xi < b\)，对于充分大的 \(n\)，\(\xi + \frac{1}{n} < b\) 且
\(\xi + \frac{1}{n} > \xi = \sup A\)，所以
\(\xi + \frac{1}{n} \notin A\)，即 \(f(\xi + \frac{1}{n}) \geq c\)。由
\(f\) 的连续性，\(f(\xi) = \lim f(\xi + \frac{1}{n}) \geq c\)。

因此 \(f(\xi) = c\)。\(\blacksquare\)

\textbf{推论 22.1}（零点定理）设 \(f\) 在 \([a, b]\) 上连续，且
\(f(a) \cdot f(b) < 0\)，则存在 \(\xi \in (a, b)\)，使得
\(f(\xi) = 0\)。

\textbf{证明}：由介值定理取 \(c = 0\)。\(\blacksquare\)

\paragraph{22.3.4 最值定理}\label{ux6700ux503cux5b9aux7406}

\textbf{定理 22.16}（最值定理）设 \(f\) 在闭区间 \([a, b]\) 上连续，则
\(f\) 在 \([a, b]\) 上取到最大值和最小值，即存在
\(x_1, x_2 \in [a, b]\)，使得

\[f(x_1) = \sup_{x \in [a, b]} f(x), \quad f(x_2) = \inf_{x \in [a, b]} f(x)\]

\textbf{证明}：我们证明最大值的存在性（最小值类似）。

\textbf{第一步}：证明 \(f\) 在 \([a, b]\) 上有界。用反证法：若 \(f\)
无上界，则对每个 \(n\) 存在 \(x_n \in [a, b]\) 使得
\(f(x_n) > n\)。由波尔查诺-魏尔斯特拉斯定理，\(\{x_n\}\) 有收敛子列
\(\{x_{n_k}\}\)，设 \(x_{n_k} \to \xi \in [a, b]\)。由 \(f\)
的连续性，\(f(x_{n_k}) \to f(\xi)\)。但
\(f(x_{n_k}) > n_k \to +\infty\)，矛盾。

\textbf{第二步}：设 \(M = \sup_{x \in [a, b]} f(x)\)（由第一步 \(M\)
有限）。由上确界的定义，对每个 \(n\) 存在 \(x_n\) 使得
\(f(x_n) > M - \frac{1}{n}\)。由波尔查诺-魏尔斯特拉斯定理，\(\{x_n\}\)
有收敛子列 \(x_{n_k} \to x_0 \in [a, b]\)。由 \(f\)
的连续性，\(f(x_0) = \lim f(x_{n_k}) \geq M\)。又 \(f(x_0) \leq M\)，故
\(f(x_0) = M\)。\(\blacksquare\)

\bigskip\hrule\bigskip

\subsubsection{22.4
函数极限的进一步讨论}\label{ux51fdux6570ux6781ux9650ux7684ux8fdbux4e00ux6b65ux8ba8ux8bba}

\paragraph{\texorpdfstring{22.4.1 \(x \to \infty\)
时的极限}{22.4.1 x \textbackslash to \textbackslash infty 时的极限}}\label{x-to-infty-ux65f6ux7684ux6781ux9650}

\textbf{定义 22.6} 设 \(f\) 在 \([a, +\infty)\)
上有定义，\(L \in \mathbb{R}\)。若

\[\forall \varepsilon > 0, \exists X > 0, \forall x > X: |f(x) - L| < \varepsilon\]

则记 \(\lim_{x \to +\infty} f(x) = L\)。

类似地定义 \(\lim_{x \to -\infty} f(x) = L\) 和
\(\lim_{x \to \infty} f(x) = L\)。

\textbf{定理 22.17}
\(\lim_{x \to +\infty} \frac{1}{x^{\alpha}} = 0\)（\(\alpha > 0\)）。

\textbf{证明}：对于任意 \(\varepsilon > 0\)，取
\(X = \left(\frac{1}{\varepsilon}\right)^{1/\alpha}\)。当 \(x > X\) 时，

\[\left|\frac{1}{x^{\alpha}}\right| = \frac{1}{x^{\alpha}} < \frac{1}{X^{\alpha}} = \varepsilon\]

\(\blacksquare\)

\paragraph{22.4.2
无穷小与无穷大}\label{ux65e0ux7a77ux5c0fux4e0eux65e0ux7a77ux5927}

\textbf{定义 22.7}（无穷小）若 \(\lim_{x \to a} f(x) = 0\)，则称
\(f(x)\) 是 \(x \to a\) 时的\textbf{无穷小量}。

\textbf{定义 22.8}（无穷大）若 \(\lim_{x \to a} |f(x)| = +\infty\)，即

\[\forall M > 0, \exists \delta > 0, \forall x: 0 < |x - a| < \delta \Rightarrow |f(x)| > M\]

则称 \(f(x)\) 是 \(x \to a\) 时的\textbf{无穷大量}。

\textbf{定理 22.18}（无穷小与无穷大的关系）若 \(f(x)\) 是 \(x \to a\)
时的无穷小量，且 \(f(x) \neq 0\) 在 \(a\) 的某个去心邻域内成立，则
\(\frac{1}{f(x)}\) 是 \(x \to a\) 时的无穷大量。

\textbf{证明}：对于任意 \(M > 0\)，取
\(\varepsilon = \frac{1}{M}\)。存在 \(\delta > 0\)，当
\(0 < |x - a| < \delta\) 时 \(|f(x)| < \varepsilon = \frac{1}{M}\)，故
\(\left|\frac{1}{f(x)}\right| > M\)。\(\blacksquare\)

\paragraph{22.4.3 等价无穷小}\label{ux7b49ux4ef7ux65e0ux7a77ux5c0f}

\textbf{定义 22.9} 设 \(\alpha(x)\) 和 \(\beta(x)\) 都是 \(x \to a\)
时的无穷小量。若 \(\lim_{x \to a} \frac{\alpha(x)}{\beta(x)} = 1\)，则称
\(\alpha(x)\) 和 \(\beta(x)\) 是\textbf{等价无穷小}，记作
\(\alpha(x) \sim \beta(x)\)（\(x \to a\)）。

\textbf{定理 22.19} 当 \(x \to 0\) 时的常用等价无穷小：

\begin{enumerate}
\def\labelenumi{\arabic{enumi}.}
\tightlist
\item
  \(\sin x \sim x\)
\item
  \(\tan x \sim x\)
\item
  \(1 - \cos x \sim \frac{x^2}{2}\)
\item
  \(e^x - 1 \sim x\)
\item
  \(\ln(1 + x) \sim x\)
\item
  \((1 + x)^{\alpha} - 1 \sim \alpha x\)
\end{enumerate}

\textbf{证明}：

\begin{enumerate}
\def\labelenumi{(\arabic{enumi})}
\item
  \(\lim_{x \to 0} \frac{\sin x}{x} = 1\)（定理 22.12）。
\item
  \(\lim_{x \to 0} \frac{1 - \cos x}{\frac{x^2}{2}} = \lim_{x \to 0} \frac{2\sin^2\frac{x}{2}}{\frac{x^2}{2}} = \lim_{x \to 0} \left(\frac{\sin\frac{x}{2}}{\frac{x}{2}}\right)^2 = 1\)。
\item
  \(\lim_{x \to 0} \frac{e^x - 1}{x} = \lim_{t \to 0} \frac{t}{\ln(1+t)} = 1\)（令
  \(t = e^x - 1\)）。
\item
  \(\lim_{x \to 0} \frac{\ln(1+x)}{x} = \lim_{x \to 0} \ln(1+x)^{1/x} = \ln e = 1\)。
\item
  令 \(u = \alpha\ln(1+x)\)。当 \(x \to 0\) 时，由
  (5)，\(\ln(1+x) \sim x\)，故 \(u \sim \alpha x \to 0\)。由
  (4)（\(e^u - 1 \sim u\)，其中
  \(u \to 0\)），\(e^{\alpha\ln(1+x)} - 1 \sim \alpha\ln(1+x)\)。再由
  (5)，\(\alpha\ln(1+x) \sim \alpha x\)。由等价无穷小的传递性（若
  \(\alpha_1 \sim \alpha_2\) 且 \(\alpha_2 \sim \alpha_3\)，则
  \(\alpha_1 \sim \alpha_3\)，这由
  \(\frac{\alpha_1}{\alpha_3} = \frac{\alpha_1}{\alpha_2} \cdot \frac{\alpha_2}{\alpha_3} \to 1 \cdot 1 = 1\)），\((1+x)^{\alpha} - 1 \sim \alpha x\)。\(\blacksquare\)
\end{enumerate}

\textbf{定理 22.20}（等价无穷小替换定理）设
\(\alpha \sim \alpha'\)，\(\beta \sim \beta'\)（\(x \to a\)），且
\(\lim_{x \to a} \frac{\alpha'}{\beta'}\) 存在，则

\[\lim_{x \to a} \frac{\alpha}{\beta} = \lim_{x \to a} \frac{\alpha'}{\beta'}\]

\textbf{证明}：

\[\lim_{x \to a} \frac{\alpha}{\beta} = \lim_{x \to a} \frac{\alpha}{\alpha'} \cdot \frac{\alpha'}{\beta'} \cdot \frac{\beta'}{\beta} = 1 \cdot \lim_{x \to a} \frac{\alpha'}{\beta'} \cdot 1 = \lim_{x \to a} \frac{\alpha'}{\beta'}\]

\(\blacksquare\)

\bigskip\hrule\bigskip

\subsubsection{22.5 一致连续性}\label{ux4e00ux81f4ux8fdeux7eedux6027}

\textbf{定义 22.10}（一致连续）设 \(f: I \to \mathbb{R}\)，其中 \(I\)
是区间。若

\[\forall \varepsilon > 0, \exists \delta > 0, \forall x_1, x_2 \in I: |x_1 - x_2| < \delta \Rightarrow |f(x_1) - f(x_2)| < \varepsilon\]

则称 \(f\) 在 \(I\) 上\textbf{一致连续}。

\textbf{注}：一致连续与逐点连续的区别在于：逐点连续中 \(\delta\)
可以依赖于点 \(x_0\)，而一致连续中 \(\delta\) 只依赖于
\(\varepsilon\)，对区间上所有点统一适用。

\textbf{定理 22.21}（康托尔定理）闭区间 \([a, b]\)
上的连续函数一致连续。

\textbf{证明}：用反证法。假设 \(f\) 在 \([a, b]\)
上连续但不一致连续。则存在 \(\varepsilon_0 > 0\)，对每个 \(n\)，存在
\(x_n, y_n \in [a, b]\)，\(|x_n - y_n| < \frac{1}{n}\) 但
\(|f(x_n) - f(y_n)| \geq \varepsilon_0\)。

由波尔查诺-魏尔斯特拉斯定理，\(\{x_n\}\) 有收敛子列
\(x_{n_k} \to \xi \in [a, b]\)。由于
\(|y_{n_k} - x_{n_k}| < \frac{1}{n_k} \to 0\)，故 \(y_{n_k} \to \xi\)。

由 \(f\)
的连续性，\(f(x_{n_k}) \to f(\xi)\)，\(f(y_{n_k}) \to f(\xi)\)，故
\(|f(x_{n_k}) - f(y_{n_k})| \to 0\)，与
\(|f(x_{n_k}) - f(y_{n_k})| \geq \varepsilon_0\) 矛盾。\(\blacksquare\)

\subsubsection{22.6
函数极限与数列极限的关系}\label{ux51fdux6570ux6781ux9650ux4e0eux6570ux5217ux6781ux9650ux7684ux5173ux7cfb}

\textbf{定理 22.22}（海涅定理/归结原则）设 \(f\) 在 \(a\)
的某个去心邻域内有定义。则 \(\lim_{x \to a} f(x) = L\)
当且仅当：对于任意收敛于 \(a\) 且各项不等于 \(a\) 的数列
\(\{x_n\}\)，都有 \(\lim_{n \to \infty} f(x_n) = L\)。

\textbf{证明}：

必要性：设 \(\lim_{x \to a} f(x) = L\)，\(\{x_n\}\) 满足
\(x_n \to a\)，\(x_n \neq a\)。对于 \(\varepsilon > 0\)，存在
\(\delta > 0\)，当 \(0 < |x - a| < \delta\) 时
\(|f(x) - L| < \varepsilon\)。又存在 \(N\)，当 \(n > N\) 时
\(|x_n - a| < \delta\)。因为 \(x_n \neq a\)，所以
\(0 < |x_n - a| < \delta\)，故 \(|f(x_n) - L| < \varepsilon\)。

充分性（反证法）：假设 \(\lim_{x \to a} f(x) \neq L\)。则存在
\(\varepsilon_0 > 0\)，对每个 \(\delta = \frac{1}{n}\)，存在 \(x_n\)
满足 \(0 < |x_n - a| < \frac{1}{n}\) 但
\(|f(x_n) - L| \geq \varepsilon_0\)。此时 \(x_n \to a\) 但
\(f(x_n) \not\to L\)，矛盾。\(\blacksquare\)

\bigskip\hrule\bigskip

\subsubsection{22.7
实数完备性的等价表述}\label{ux5b9eux6570ux5b8cux5907ux6027ux7684ux7b49ux4ef7ux8868ux8ff0}

实数系的完备性有多种等价表述，它们从不同角度刻画了实数系''没有空隙''的本质。

\textbf{定理 22.23} 以下命题等价：

\begin{enumerate}
\def\labelenumi{\arabic{enumi}.}
\tightlist
\item
  \textbf{确界存在定理}：非空有上界的实数集必有上确界。
\item
  \textbf{单调有界定理}：单调递增有上界的数列必收敛。
\item
  \textbf{闭区间套定理}：设 \(\{[a_n, b_n]\}\)
  是递缩闭区间序列（\([a_{n+1}, b_{n+1}] \subseteq [a_n, b_n]\)），且
  \(b_n - a_n \to 0\)，则存在唯一的
  \(\xi \in \bigcap_{n=1}^{\infty}[a_n, b_n]\)。
\item
  \textbf{波尔查诺-魏尔斯特拉斯定理}：有界数列必有收敛子列。
\item
  \textbf{柯西收敛准则}：实数数列收敛当且仅当它是柯西列。
\end{enumerate}

\textbf{证明思路}：\(1 \Rightarrow 2\)：即定理 22.5
的证明。\(2 \Rightarrow 3\)：\(\{a_n\}\) 单调递增有上界，\(\{b_n\}\)
单调递减有下界，两者收敛于同一极限。\(3 \Rightarrow 4\)：用二分法构造闭区间套。\(4 \Rightarrow 5\)：柯西列有界，故有收敛子列，再证全列收敛。\(5 \Rightarrow 1\)：用柯西列逼近上确界。\(\blacksquare\)

\textbf{补充证明}（\(3 \Rightarrow 4\)）：设 \(\{x_n\}\) 为有界数列，即
\(x_n \in [a, b]\)（对一切 \(n\)）。将 \([a, b]\) 二分为
\(\left[a, \frac{a+b}{2}\right]\) 和
\(\left[\frac{a+b}{2}, b\right]\)，其中至少有一个区间包含 \(\{x_n\}\)
的无穷多项（否则整个数列只有有限项），取该区间为 \(I_1\)。继续将 \(I_1\)
二分，同理取包含无穷多项的子区间为 \(I_2\)。如此递推，得到嵌套闭区间序列
\(I_1 \supset I_2 \supset I_3 \supset \cdots\)，且
\(\lvert I_n \rvert = \frac{b - a}{2^n} \to 0\)。由闭区间套定理（3），\(\bigcap_{n=1}^{\infty} I_n = \{x^*\}\)（唯一一点）。现在在每个
\(I_n\) 中取一项 \(x_{n_k}\)（可取 \(n_k > n_{k-1}\) 使下标递增），则
\(x_{n_k} \in I_k\)，故
\(|x_{n_k} - x^*| \leq |I_k| = \frac{b-a}{2^k} \to 0\)。因此 \(\{x_n\}\)
有收敛子列 \(x_{n_k} \to x^*\)。\(\blacksquare\)

\textbf{补充证明}（\(5 \Rightarrow 1\)）：设 \(S \subseteq \mathbb{R}\)
非空且有上界。取 \(a_0 \in S\) 和一个上界 \(b_0\)（即
\(\forall x \in S,\, x \leq b_0\)），则 \(a_0 \leq b_0\)。对区间
\([a_n, b_n]\) 进行二分，令 \(m_n = \frac{a_n + b_n}{2}\)：若 \(m_n\) 是
\(S\) 的上界，则令 \([a_{n+1}, b_{n+1}] = [a_n, m_n]\)；否则存在 \(S\)
中某元素 \(> m_n\)，令
\([a_{n+1}, b_{n+1}] = [m_n, b_n]\)。由此构造的序列 \(\{a_n\}\) 满足
\(b_n - a_n = \frac{b_0 - a_0}{2^n} \to 0\)，且对 \(m > n\) 有
\(|a_m - a_n| \leq b_n - a_n \to 0\)，故 \(\{a_n\}\)
是柯西列。由柯西收敛准则（5），\(\{a_n\}\) 收敛，设
\(\xi = \lim_{n \to \infty} a_n\)。

下证 \(\xi = \sup S\)。首先，对每个
\(n\)，\(a_n \leq \xi \leq b_n\)（由构造，\(a_n\) 递增趋于
\(\xi\)）。对任意 \(x \in S\)，由构造每个 \(b_n\) 都是 \(S\) 的上界，故
\(x \leq b_n\)，取极限得 \(x \leq \xi\)，所以 \(\xi\) 是 \(S\)
的上界。其次，对任意 \(\varepsilon > 0\)，取 \(n\) 足够大使得
\(b_n - a_n < \varepsilon\)，则
\(a_n > \xi - \varepsilon\)。由构造，\(a_n\) 不是 \(S\)
的上界（否则会被取到右半区间），故存在 \(x \in S\) 使得
\(x > a_n > \xi - \varepsilon\)。这说明 \(\xi - \varepsilon\) 不是 \(S\)
的上界，因此 \(\xi\) 是最小上界，即 \(\xi = \sup S\)。\(\blacksquare\)

\textbf{注}：这五个等价命题共同构成了实数完备性理论的基石，是整个分析学的逻辑起点。

\textbf{本章展望}：本章建立了极限与连续的严格理论。极限理论使我们能够严格讨论''趋近''和''无穷逼近''的概念。在下一章（Ch
23）中，我们将利用极限理论定义导数------函数在某一点的瞬时变化率。导数是微积分的两大核心概念之一，它将极限理论应用于研究函数的局部性质。

\bigskip\hrule\bigskip

\subsection{Ch 23
导数与微分}\label{ch-23-ux5bfcux6570ux4e0eux5faeux5206}

\textbf{前置知识}：本章依赖 Ch 22（极限与连续，特别是函数极限的
\(\varepsilon\)-\(\delta\) 定义、四则运算和夹逼定理）、Ch
20（三角恒等变换，为三角函数的导数公式做准备）、Ch 19（重要极限
\(\lim_{x \to 0} \frac{\sin x}{x} = 1\)）。

\bigskip\hrule\bigskip

\subsubsection{23.1 导数的定义}\label{ux5bfcux6570ux7684ux5b9aux4e49}

\paragraph{23.1.1
从切线问题引入}\label{ux4eceux5207ux7ebfux95eeux9898ux5f15ux5165}

考虑曲线 \(y = f(x)\) 上的点 \(P = (a, f(a))\)。在曲线上取另一点
\(Q = (a + h, f(a + h))\)，则割线 \(PQ\) 的斜率为

\[\frac{f(a + h) - f(a)}{h}\]

当 \(h \to 0\) 时，割线 \(PQ\) 趋于曲线在 \(P\) 点的切线，切线的斜率为

\[\lim_{h \to 0} \frac{f(a + h) - f(a)}{h}\]

\paragraph{23.1.2 导数的定义}\label{ux5bfcux6570ux7684ux5b9aux4e49-1}

\textbf{定义 23.1}（导数）设 \(f\) 在 \(a\) 的某个邻域内有定义。若极限

\[\lim_{h \to 0} \frac{f(a + h) - f(a)}{h}\]

存在，则称 \(f\) 在 \(a\) 处\textbf{可导}，此极限值称为 \(f\) 在 \(a\)
处的\textbf{导数}，记作 \(f'(a)\) 或 \(\frac{df}{dx}\big|_{x=a}\)。

等价地，

\[f'(a) = \lim_{x \to a} \frac{f(x) - f(a)}{x - a}\]

\textbf{定义 23.2}（导函数）若 \(f\) 在区间 \(I\) 上的每一点都可导，则
\(f': I \to \mathbb{R}\)，

\[f'(x) = \lim_{h \to 0} \frac{f(x + h) - f(x)}{h}\]

称为 \(f\) 的\textbf{导函数}（简称导数）。

\paragraph{23.1.3
可导与连续的关系}\label{ux53efux5bfcux4e0eux8fdeux7eedux7684ux5173ux7cfb}

\textbf{定理 23.1}（可导必连续）若 \(f\) 在 \(a\) 处可导，则 \(f\) 在
\(a\) 处连续。

\textbf{证明}：设 \(f'(a) = \lim_{h \to 0} \frac{f(a+h) - f(a)}{h}\)
存在。则

\[\lim_{h \to 0} [f(a+h) - f(a)] = \lim_{h \to 0} \frac{f(a+h) - f(a)}{h} \cdot h = f'(a) \cdot 0 = 0\]

所以 \(\lim_{h \to 0} f(a+h) = f(a)\)，即 \(f\) 在 \(a\)
处连续。\(\blacksquare\)

\textbf{注}：逆命题不成立。例如 \(f(x) = |x|\) 在 \(x = 0\)
处连续但不可导（左导数为 \(-1\)，右导数为 \(1\)，不相等）。

\paragraph{23.1.4
导数的几何意义}\label{ux5bfcux6570ux7684ux51e0ux4f55ux610fux4e49}

\(f'(a)\) 表示曲线 \(y = f(x)\) 在点 \((a, f(a))\) 处的切线斜率。

切线方程：\(y - f(a) = f'(a)(x - a)\)

法线方程：\(y - f(a) = -\frac{1}{f'(a)}(x - a)\)（\(f'(a) \neq 0\) 时）

\bigskip\hrule\bigskip

\subsubsection{23.2
基本初等函数的导数}\label{ux57faux672cux521dux7b49ux51fdux6570ux7684ux5bfcux6570}

\textbf{定理 23.2} 对于正整数 \(n\)，\((x^n)' = nx^{n-1}\)。

\textbf{证明}：

\[(x^n)' = \lim_{h \to 0} \frac{(x+h)^n - x^n}{h} = \lim_{h \to 0} \frac{\sum_{k=0}^{n} \binom{n}{k} x^{n-k} h^k - x^n}{h}\]

\[= \lim_{h \to 0} \left(nx^{n-1} + \binom{n}{2}x^{n-2}h + \cdots + h^{n-1}\right) = nx^{n-1}\]

\(\blacksquare\)

\textbf{定理 23.3} \((\sin x)' = \cos x\)。

\textbf{证明}：

\[(\sin x)' = \lim_{h \to 0} \frac{\sin(x+h) - \sin x}{h}\]

由和差化积公式，\(\sin(x+h) - \sin x = 2\cos\left(x + \frac{h}{2}\right)\sin\frac{h}{2}\)。

\[(\sin x)' = \lim_{h \to 0} \cos\left(x + \frac{h}{2}\right) \cdot \frac{\sin\frac{h}{2}}{\frac{h}{2}}\]

由 \(\lim_{t \to 0} \frac{\sin t}{t} = 1\)（定理 22.12），得
\((\sin x)' = \cos x \cdot 1 = \cos x\)。\(\blacksquare\)

\textbf{定理 23.4} \((\cos x)' = -\sin x\)。

\textbf{证明}：

\[(\cos x)' = \lim_{h \to 0} \frac{\cos(x+h) - \cos x}{h}\]

由和差化积，\(\cos(x+h) - \cos x = -2\sin\left(x + \frac{h}{2}\right)\sin\frac{h}{2}\)。

类似地得 \((\cos x)' = -\sin x\)。\(\blacksquare\)

\textbf{定理 23.5} \((e^x)' = e^x\)。

\textbf{证明}：

\[(e^x)' = \lim_{h \to 0} \frac{e^{x+h} - e^x}{h} = e^x \lim_{h \to 0} \frac{e^h - 1}{h}\]

令 \(t = e^h - 1\)，则 \(h = \ln(1+t)\)，当 \(h \to 0\) 时 \(t \to 0\)。

\[\lim_{h \to 0} \frac{e^h - 1}{h} = \lim_{t \to 0} \frac{t}{\ln(1+t)} = \lim_{t \to 0} \frac{1}{\frac{1}{t}\ln(1+t)} = \frac{1}{\ln e} = 1\]

因此 \((e^x)' = e^x\)。\(\blacksquare\)

\textbf{定理 23.6} \((a^x)' = a^x \ln a\)（\(a > 0\)，\(a \neq 1\)）。

\textbf{证明}：\(a^x = e^{x \ln a}\)，由链式法则（定理
23.12），\((a^x)' = e^{x \ln a} \cdot \ln a = a^x \ln a\)。\(\blacksquare\)

\textbf{定理 23.7} \((\ln x)' = \frac{1}{x}\)（\(x > 0\)）。

\textbf{证明}：

\[(\ln x)' = \lim_{h \to 0} \frac{\ln(x+h) - \ln x}{h} = \lim_{h \to 0} \frac{1}{h}\ln\left(1 + \frac{h}{x}\right)\]

令 \(t = \frac{h}{x}\)，则 \(h = xt\)，当 \(h \to 0\) 时 \(t \to 0\)。

\[(\ln x)' = \frac{1}{x} \lim_{t \to 0} \frac{\ln(1+t)}{t} = \frac{1}{x} \cdot 1 = \frac{1}{x}\]

\(\blacksquare\)

\textbf{定理 23.8} \((\tan x)' = \sec^2 x = \frac{1}{\cos^2 x}\)。

\textbf{证明}：由除法法则，

\[(\tan x)' = \left(\frac{\sin x}{\cos x}\right)' = \frac{\cos x \cdot \cos x - \sin x \cdot (-\sin x)}{\cos^2 x} = \frac{\cos^2 x + \sin^2 x}{\cos^2 x} = \frac{1}{\cos^2 x}\]

\(\blacksquare\)

\bigskip\hrule\bigskip

\subsubsection{23.3 求导法则}\label{ux6c42ux5bfcux6cd5ux5219}

\paragraph{23.3.1
四则运算求导法则}\label{ux56dbux5219ux8fd0ux7b97ux6c42ux5bfcux6cd5ux5219}

\textbf{定理 23.9}（和的导数）\((f + g)' = f' + g'\)

\textbf{证明}：

\[(f + g)'(x) = \lim_{h \to 0} \frac{[f(x+h) + g(x+h)] - [f(x) + g(x)]}{h} = f'(x) + g'(x)\]

\(\blacksquare\)

\textbf{定理 23.10}（常数倍的导数）\((cf)' = cf'\)

\textbf{证明}：

\[(cf)'(x) = \lim_{h \to 0} \frac{cf(x+h) - cf(x)}{h} = c \lim_{h \to 0} \frac{f(x+h) - f(x)}{h} = cf'(x)\]

\(\blacksquare\)

\textbf{定理 23.11}（乘法法则）\((fg)' = f'g + fg'\)

\textbf{证明}：

\[(fg)'(x) = \lim_{h \to 0} \frac{f(x+h)g(x+h) - f(x)g(x)}{h}\]

在分子中加减 \(f(x+h)g(x)\)：

\[= \lim_{h \to 0} \left[f(x+h) \cdot \frac{g(x+h) - g(x)}{h} + g(x) \cdot \frac{f(x+h) - f(x)}{h}\right]\]

由于 \(f\) 在 \(x\) 处可导，所以
\(\lim_{h \to 0} f(x+h) = f(x)\)（可导必连续）。因此

\[= f(x)g'(x) + g(x)f'(x)\]

\(\blacksquare\)

\textbf{除法法则}：若 \(g(x) \neq 0\)，则
\(\left(\frac{f}{g}\right)' = \frac{f'g - fg'}{g^2}\)。

\textbf{证明}：设 \(q = f/g\)，则
\(f = qg\)。由乘法法则，\(f' = q'g + qg'\)。解出 \(q'\)：

\[q' = \frac{f' - qg'}{g} = \frac{f' - \frac{f}{g} \cdot g'}{g} = \frac{f'g - fg'}{g^2}\]

\(\blacksquare\)

\paragraph{23.3.2 链式法则}\label{ux94feux5f0fux6cd5ux5219}

\textbf{定理 23.12}（链式法则）设 \(g\) 在 \(x\) 处可导，\(f\) 在
\(g(x)\) 处可导，则 \(f \circ g\) 在 \(x\) 处可导，且

\[(f \circ g)'(x) = f'(g(x)) \cdot g'(x)\]

\textbf{证明}：设 \(u = g(x)\)，\(\Delta u = g(x+h) - g(x)\)。当
\(h \to 0\) 时，\(\Delta u \to 0\)（因为 \(g\) 连续）。

定义函数 \(\varphi\)：

\[\varphi(\Delta u) = \begin{cases} \frac{f(u + \Delta u) - f(u)}{\Delta u}, & \Delta u \neq 0 \\ f'(u), & \Delta u = 0 \end{cases}\]

则 \(\varphi\) 在 \(\Delta u = 0\) 处连续（因为
\(\lim_{\Delta u \to 0} \varphi(\Delta u) = f'(u) = \varphi(0)\)）。

\[\frac{f(g(x+h)) - f(g(x))}{h} = \varphi(\Delta u) \cdot \frac{g(x+h) - g(x)}{h}\]

当 \(h \to 0\)
时，\(\Delta u \to 0\)，\(\varphi(\Delta u) \to f'(u) = f'(g(x))\)，\(\frac{g(x+h) - g(x)}{h} \to g'(x)\)。

因此 \((f \circ g)'(x) = f'(g(x)) \cdot g'(x)\)。\(\blacksquare\)

\bigskip\hrule\bigskip

\subsubsection{23.4 中值定理}\label{ux4e2dux503cux5b9aux7406}

\paragraph{23.4.1 费马定理}\label{ux8d39ux9a6cux5b9aux7406}

\textbf{定理 23.13}（费马定理）设 \(f\) 在 \(a\) 处取极值，且 \(f\) 在
\(a\) 处可导，则 \(f'(a) = 0\)。

\textbf{证明}：不妨设 \(f(a)\) 是极大值。则存在 \(\delta > 0\)，使得当
\(|x - a| < \delta\) 时，\(f(x) \leq f(a)\)。

当 \(a < x < a + \delta\) 时，\(\frac{f(x) - f(a)}{x - a} \leq 0\)，所以
\(f'(a) = \lim_{x \to a^+} \frac{f(x) - f(a)}{x - a} \leq 0\)。

当 \(a - \delta < x < a\) 时，\(\frac{f(x) - f(a)}{x - a} \geq 0\)，所以
\(f'(a) = \lim_{x \to a^-} \frac{f(x) - f(a)}{x - a} \geq 0\)。

因此 \(f'(a) = 0\)。\(\blacksquare\)

\paragraph{23.4.2 罗尔定理}\label{ux7f57ux5c14ux5b9aux7406}

\textbf{定理 23.14}（罗尔定理）设 \(f\) 在 \([a, b]\) 上连续，在
\((a, b)\) 上可导，且 \(f(a) = f(b)\)。则存在 \(\xi \in (a, b)\)，使得
\(f'(\xi) = 0\)。

\textbf{证明}：由最值定理（定理 22.16），\(f\) 在 \([a, b]\)
上取到最大值 \(M\) 和最小值 \(m\)。

若 \(M = m\)，则 \(f\) 为常数，\(f'(\xi) = 0\) 对所有 \(\xi\) 成立。

若 \(M \neq m\)，则 \(M\) 和 \(m\) 中至少有一个不等于
\(f(a) = f(b)\)。不妨设 \(M \neq f(a)\)，则最大值在某个
\(\xi \in (a, b)\) 处取到。由费马定理，\(f'(\xi) = 0\)。\(\blacksquare\)

\paragraph{23.4.3
拉格朗日中值定理}\label{ux62c9ux683cux6717ux65e5ux4e2dux503cux5b9aux7406}

\textbf{定理 23.15}（拉格朗日中值定理）设 \(f\) 在 \([a, b]\) 上连续，在
\((a, b)\) 上可导。则存在 \(\xi \in (a, b)\)，使得

\[f'(\xi) = \frac{f(b) - f(a)}{b - a}\]

\textbf{证明}：定义辅助函数

\[F(x) = f(x) - f(a) - \frac{f(b) - f(a)}{b - a}(x - a)\]

则 \(F\) 在 \([a, b]\) 上连续，在 \((a, b)\) 上可导，且
\(F(a) = F(b) = 0\)。

由罗尔定理，存在 \(\xi \in (a, b)\) 使得 \(F'(\xi) = 0\)，即

\[f'(\xi) - \frac{f(b) - f(a)}{b - a} = 0\]

所以 \(f'(\xi) = \frac{f(b) - f(a)}{b - a}\)。\(\blacksquare\)

\paragraph{23.4.4
柯西中值定理}\label{ux67efux897fux4e2dux503cux5b9aux7406}

\textbf{定理 23.16}（柯西中值定理）设 \(f\)，\(g\) 在 \([a, b]\)
上连续，在 \((a, b)\) 上可导，且 \(g'(x) \neq 0\)。则存在
\(\xi \in (a, b)\)，使得

\[\frac{f'(\xi)}{g'(\xi)} = \frac{f(b) - f(a)}{g(b) - g(a)}\]

\textbf{证明}：首先说明 \(g(b) \neq g(a)\)：若
\(g(b) = g(a)\)，则由罗尔定理，存在 \(c \in (a, b)\) 使
\(g'(c) = 0\)，与条件 \(g'(x) \neq 0\) 矛盾。

定义辅助函数

\[F(x) = f(x) - f(a) - \frac{f(b) - f(a)}{g(b) - g(a)}[g(x) - g(a)]\]

则 \(F(a) = F(b) = 0\)。由罗尔定理，存在 \(\xi \in (a, b)\) 使得
\(F'(\xi) = 0\)，即

\[f'(\xi) - \frac{f(b) - f(a)}{g(b) - g(a)} g'(\xi) = 0\]

由于 \(g'(\xi) \neq 0\)，化简得
\(\frac{f'(\xi)}{g'(\xi)} = \frac{f(b) - f(a)}{g(b) - g(a)}\)。\(\blacksquare\)

\bigskip\hrule\bigskip

\subsubsection{23.5 洛必达法则}\label{ux6d1bux5fc5ux8fbeux6cd5ux5219}

\textbf{定理 23.17}（洛必达法则，\(\frac{0}{0}\) 型）设 \(f\)，\(g\) 在
\(a\) 的某个去心邻域内可导，\(g'(x) \neq 0\)，且

\[\lim_{x \to a} f(x) = 0, \quad \lim_{x \to a} g(x) = 0\]

若 \(\lim_{x \to a} \frac{f'(x)}{g'(x)} = L\)（\(L\) 可以是有限数或
\(\pm\infty\)），则

\[\lim_{x \to a} \frac{f(x)}{g(x)} = L\]

\textbf{证明}：补充定义 \(f(a) = g(a) = 0\)，则 \(f\) 和 \(g\) 在 \(a\)
处连续。

由于 \(g'(x) \neq 0\) 在 \(a\) 的去心邻域内成立且 \(g\) 连续，\(g\) 在
\(a\) 附近严格单调，故对 \(x \neq a\) 有 \(g(x) \neq g(a) = 0\)。因此对
\(x \neq a\)，可以应用柯西中值定理（定理 23.16）：存在 \(\xi\) 介于
\(a\) 和 \(x\) 之间，使得

\[\frac{f(x)}{g(x)} = \frac{f(x) - f(a)}{g(x) - g(a)} = \frac{f'(\xi)}{g'(\xi)}\]

当 \(x \to a\) 时，\(\xi \to a\)（因为 \(\xi\) 介于 \(a\) 和 \(x\)
之间），所以

\[\lim_{x \to a} \frac{f(x)}{g(x)} = \lim_{\xi \to a} \frac{f'(\xi)}{g'(\xi)} = L\]

\(\blacksquare\)

\textbf{注}：洛必达法则只是极限存在的充分条件，而非必要条件------当
\(\frac{f'(x)}{g'(x)}\) 的极限不存在时，\(\frac{f(x)}{g(x)}\)
的极限仍可能存在。

\bigskip\hrule\bigskip

\subsubsection{23.6
导数的应用：极值与单调性}\label{ux5bfcux6570ux7684ux5e94ux7528ux6781ux503cux4e0eux5355ux8c03ux6027}

\textbf{定理 23.18}（单调性判定）设 \(f\) 在 \([a, b]\) 上连续，在
\((a, b)\) 上可导。

\begin{enumerate}
\def\labelenumi{\arabic{enumi}.}
\tightlist
\item
  若 \(\forall x \in (a, b), f'(x) > 0\)，则 \(f\) 在 \([a, b]\)
  上严格递增。
\item
  若 \(\forall x \in (a, b), f'(x) < 0\)，则 \(f\) 在 \([a, b]\)
  上严格递减。
\end{enumerate}

\textbf{证明}：设
\(x_1, x_2 \in [a, b]\)，\(x_1 < x_2\)。由拉格朗日中值定理，存在
\(\xi \in (x_1, x_2)\) 使得

\[f(x_2) - f(x_1) = f'(\xi)(x_2 - x_1)\]

因为 \(f'(\xi) > 0\) 且 \(x_2 - x_1 > 0\)，所以
\(f(x_2) > f(x_1)\)。\(\blacksquare\)

\textbf{定理 23.19}（极值的第一充分条件）设 \(f\) 在 \(a\) 处连续，在
\(a\) 的某个去心邻域内可导。

\begin{enumerate}
\def\labelenumi{\arabic{enumi}.}
\tightlist
\item
  若 \(f'\) 在 \(a\) 左侧为正、右侧为负，则 \(f(a)\) 是极大值。
\item
  若 \(f'\) 在 \(a\) 左侧为负、右侧为正，则 \(f(a)\) 是极小值。
\end{enumerate}

\textbf{证明}：(1) 由定理 23.18，\(f\) 在 \(a\)
左侧严格递增，在右侧严格递减。因此 \(f(a)\) 是极大值。\(\blacksquare\)

\textbf{定理 23.20}（极值的第二充分条件）设 \(f\) 在 \(a\)
处有二阶导数，\(f'(a) = 0\)。

\begin{enumerate}
\def\labelenumi{\arabic{enumi}.}
\tightlist
\item
  若 \(f''(a) < 0\)，则 \(f(a)\) 是极大值。
\item
  若 \(f''(a) > 0\)，则 \(f(a)\) 是极小值。
\end{enumerate}

\textbf{证明}：(1) 由 \(f''(a) < 0\)，即
\(\lim_{h \to 0} \frac{f'(a+h)}{h} < 0\)。因此存在
\(\delta > 0\)，使得当 \(0 < h < \delta\) 时 \(f'(a+h) < 0\)；当
\(-\delta < h < 0\) 时 \(f'(a+h) > 0\)。由第一充分条件，\(f(a)\)
是极大值。\(\blacksquare\)

\bigskip\hrule\bigskip

\subsubsection{23.7 微分}\label{ux5faeux5206}

\textbf{定义 23.3}（微分）设 \(f\) 在 \(x\) 处可导。函数 \(f\) 在 \(x\)
处的\textbf{微分}定义为：

\[dy = f'(x) dx\]

其中 \(dx\) 是自变量的增量。

\textbf{定理 23.21}（微分与增量的关系）设 \(f\) 在 \(x\) 处可导，则

\[\Delta y = f(x + \Delta x) - f(x) = f'(x) \Delta x + o(\Delta x)\]

其中 \(\lim_{\Delta x \to 0} \frac{o(\Delta x)}{\Delta x} = 0\)。

\textbf{证明}：由导数定义，令
\(\varepsilon(\Delta x) = \frac{f(x + \Delta x) - f(x)}{\Delta x} - f'(x)\)，则
\(\lim_{\Delta x \to 0} \varepsilon(\Delta x) = 0\)。

\(\Delta y = [f'(x) + \varepsilon(\Delta x)] \Delta x = f'(x) \Delta x + \varepsilon(\Delta x) \Delta x\)。

令 \(o(\Delta x) = \varepsilon(\Delta x) \Delta x\)，则
\(\frac{o(\Delta x)}{\Delta x} = \varepsilon(\Delta x) \to 0\)。\(\blacksquare\)

\bigskip\hrule\bigskip

\subsubsection{23.8 高阶导数}\label{ux9ad8ux9636ux5bfcux6570}

\paragraph{23.8.1 二阶导数}\label{ux4e8cux9636ux5bfcux6570}

\textbf{定义 23.4}（二阶导数）设 \(f\) 在区间 \(I\) 上可导。若 \(f'\) 在
\(I\) 上也可导，则称 \(f'\) 的导数为 \(f\) 的\textbf{二阶导数}，记作
\(f''(x)\) 或 \(\frac{d^2 f}{dx^2}\)。

\textbf{定义 23.5}（\(n\)
阶导数）类似地，\(f^{(n)}(x) = (f^{(n-1)})'(x)\)（\(n \geq 1\)），其中
\(f^{(0)} = f\)。

\textbf{定理 23.22}（莱布尼茨公式）设 \(f\) 和 \(g\) 都有 \(n\)
阶导数，则

\[(fg)^{(n)} = \sum_{k=0}^{n} \binom{n}{k} f^{(k)} g^{(n-k)}\]

\textbf{证明}：对 \(n\) 用数学归纳法。

\(n = 1\)：\((fg)' = f'g + fg'\)，即
\(\binom{1}{0}f^{(0)}g^{(1)} + \binom{1}{1}f^{(1)}g^{(0)}\)，成立。

\(n \Rightarrow n+1\)：假设公式对 \(n\) 成立。则

\[(fg)^{(n+1)} = \left[(fg)^{(n)}\right]' = \sum_{k=0}^{n} \binom{n}{k} (f^{(k)} g^{(n-k)})'\]

\[= \sum_{k=0}^{n} \binom{n}{k} [f^{(k+1)} g^{(n-k)} + f^{(k)} g^{(n+1-k)}]\]

\[= \sum_{k=0}^{n} \binom{n}{k} f^{(k+1)} g^{(n-k)} + \sum_{k=0}^{n} \binom{n}{k} f^{(k)} g^{(n+1-k)}\]

在第一个和式中令 \(j = k + 1\)：

\[= \sum_{j=1}^{n+1} \binom{n}{j-1} f^{(j)} g^{(n+1-j)} + \sum_{k=0}^{n} \binom{n}{k} f^{(k)} g^{(n+1-k)}\]

\[= \sum_{k=0}^{n+1} \left[\binom{n}{k-1} + \binom{n}{k}\right] f^{(k)} g^{(n+1-k)} = \sum_{k=0}^{n+1} \binom{n+1}{k} f^{(k)} g^{(n+1-k)}\]

利用了帕斯卡恒等式
\(\binom{n}{k-1} + \binom{n}{k} = \binom{n+1}{k}\)。\(\blacksquare\)

\paragraph{\texorpdfstring{23.8.2 常见函数的 \(n\)
阶导数}{23.8.2 常见函数的 n 阶导数}}\label{ux5e38ux89c1ux51fdux6570ux7684-n-ux9636ux5bfcux6570}

\textbf{定理 23.23} \((e^{ax})^{(n)} = a^n e^{ax}\)

\textbf{证明}：\((e^{ax})' = ae^{ax}\)，\((e^{ax})'' = a^2 e^{ax}\)，归纳即得。\(\blacksquare\)

\textbf{定理 23.24}
\((\sin x)^{(n)} = \sin\left(x + \frac{n\pi}{2}\right)\)

\textbf{证明}：\((\sin x)' = \cos x = \sin(x + \frac{\pi}{2})\)。假设
\((\sin x)^{(k)} = \sin(x + \frac{k\pi}{2})\)，则

\[(\sin x)^{(k+1)} = \cos(x + \frac{k\pi}{2}) = \sin(x + \frac{k\pi}{2} + \frac{\pi}{2}) = \sin(x + \frac{(k+1)\pi}{2})\]

\(\blacksquare\)

\bigskip\hrule\bigskip

\subsubsection{23.9
泰勒展开（初步）}\label{ux6cf0ux52d2ux5c55ux5f00ux521dux6b65}

\textbf{定理 23.25}（带有拉格朗日余项的泰勒公式）设 \(f\) 在 \([a, x]\)
上有 \(n+1\) 阶导数，则

\[f(x) = \sum_{k=0}^{n} \frac{f^{(k)}(a)}{k!}(x - a)^k + R_n(x)\]

其中 \(R_n(x) = \frac{f^{(n+1)}(\xi)}{(n+1)!}(x - a)^{n+1}\)，\(\xi\)
介于 \(a\) 和 \(x\) 之间。

\textbf{证明}：设
\(R_n(x) = f(x) - \sum_{k=0}^{n} \frac{f^{(k)}(a)}{k!}(x-a)^k\)。定义辅助函数

\[F(t) = f(x) - \sum_{k=0}^{n} \frac{f^{(k)}(t)}{k!}(x-t)^k\]

则 \(F(a) = R_n(x)\)，\(F(x) = 0\)。逐项求导：

\[F'(t) = -\sum_{k=0}^{n} \frac{d}{dt}\left[\frac{f^{(k)}(t)}{k!}(x-t)^k\right]\]

对每个 \(k \geq 1\)，由乘积法则：

\[\frac{d}{dt}\left[\frac{f^{(k)}(t)}{k!}(x-t)^k\right] = \frac{f^{(k+1)}(t)}{k!}(x-t)^k - \frac{f^{(k)}(t)}{(k-1)!}(x-t)^{k-1}\]

\(k = 0\) 时
\(\frac{d}{dt}[f(t)] = f'(t)\)。因此求和后相邻项两两相消（telescoping），仅剩
\(k = n\) 的首项：

\[F'(t) = -\frac{f^{(n+1)}(t)}{n!}(x-t)^n\]

对 \(F\) 和 \(G(t) = (x-t)^{n+1}\) 在 \([a, x]\) 上应用柯西中值定理：

\[\frac{F(x) - F(a)}{G(x) - G(a)} = \frac{F'(\xi)}{G'(\xi)}\]

\[\frac{-R_n(x)}{-(x-a)^{n+1}} = \frac{-\frac{f^{(n+1)}(\xi)}{n!}(x-\xi)^n}{-(n+1)(x-\xi)^n}\]

\[R_n(x) = \frac{f^{(n+1)}(\xi)}{(n+1)!}(x-a)^{n+1}\]

\(\blacksquare\)

\textbf{推论 23.2}（麦克劳林公式，\(a = 0\)）

\[f(x) = \sum_{k=0}^{n} \frac{f^{(k)}(0)}{k!}x^k + \frac{f^{(n+1)}(\xi)}{(n+1)!}x^{n+1}\]

常见展开：

\[e^x = 1 + x + \frac{x^2}{2!} + \frac{x^3}{3!} + \cdots\]

\[\sin x = x - \frac{x^3}{3!} + \frac{x^5}{5!} - \cdots\]

\[\cos x = 1 - \frac{x^2}{2!} + \frac{x^4}{4!} - \cdots\]

\[\ln(1+x) = x - \frac{x^2}{2} + \frac{x^3}{3} - \cdots \quad (|x| \leq 1, x \neq -1)\]

\textbf{证明}：在定理 23.11（泰勒公式）中令 \(a = 0\)
即得。\(\blacksquare\)

\subsubsection{23.10 隐函数求导}\label{ux9690ux51fdux6570ux6c42ux5bfc}

\paragraph{23.10.1 隐函数}\label{ux9690ux51fdux6570}

\textbf{定义 23.6}（隐函数）由方程 \(F(x, y) = 0\) 确定的 \(y\) 关于
\(x\) 的函数称为\textbf{隐函数}。

\textbf{定理 23.26}（隐函数求导法则）设方程 \(F(x, y) = 0\) 确定了隐函数
\(y = y(x)\)，且 \(F\) 对 \(x\) 和 \(y\)
的偏导数存在，\(F_y \neq 0\)。则

\[\frac{dy}{dx} = -\frac{F_x}{F_y}\]

\textbf{证明}：对 \(F(x, y(x)) = 0\) 两边关于 \(x\) 求导，利用链式法则：

\[F_x + F_y \cdot \frac{dy}{dx} = 0\]

\[\frac{dy}{dx} = -\frac{F_x}{F_y}\]

\(\blacksquare\)

\bigskip\hrule\bigskip

\subsubsection{23.11
参数方程求导}\label{ux53c2ux6570ux65b9ux7a0bux6c42ux5bfc}

\textbf{定理 23.27} 设曲线由参数方程 \(x = \varphi(t)\)，\(y = \psi(t)\)
给出，\(\varphi\) 和 \(\psi\) 都可导，\(\varphi'(t) \neq 0\)。则

\[\frac{dy}{dx} = \frac{\psi'(t)}{\varphi'(t)}\]

\[\frac{d^2y}{dx^2} = \frac{\psi''(t)\varphi'(t) - \psi'(t)\varphi''(t)}{[\varphi'(t)]^3}\]

\textbf{证明}：

\[\frac{dy}{dx} = \frac{\frac{dy}{dt}}{\frac{dx}{dt}} = \frac{\psi'(t)}{\varphi'(t)}\]

\[\frac{d^2y}{dx^2} = \frac{\frac{d}{dt}\left(\frac{\psi'(t)}{\varphi'(t)}\right)}{\frac{dx}{dt}} = \frac{\frac{\psi''(t)\varphi'(t) - \psi'(t)\varphi''(t)}{[\varphi'(t)]^2}}{\varphi'(t)} = \frac{\psi''(t)\varphi'(t) - \psi'(t)\varphi''(t)}{[\varphi'(t)]^3}\]

\(\blacksquare\)

\bigskip\hrule\bigskip

\subsubsection{23.12
函数的凹凸性}\label{ux51fdux6570ux7684ux51f9ux51f8ux6027}

\paragraph{23.12.1
凹凸函数的定义}\label{ux51f9ux51f8ux51fdux6570ux7684ux5b9aux4e49}

\textbf{定义 23.7}（凸函数）设 \(f\) 在区间 \(I\) 上有定义。若对任意
\(x_1, x_2 \in I\) 和 \(\lambda \in [0, 1]\)：

\[f(\lambda x_1 + (1-\lambda)x_2) \leq \lambda f(x_1) + (1-\lambda)f(x_2)\]

则称 \(f\) 在 \(I\) 上是\textbf{凸函数}（向下凸，图形呈 \(\cup\)
形）。若不等号反向，则称 \(f\) 是\textbf{凹函数}。

\paragraph{23.12.2
凹凸性的判定}\label{ux51f9ux51f8ux6027ux7684ux5224ux5b9a}

\textbf{定理 23.28}（二阶导数判定法）设 \(f\) 在区间 \(I\)
上有二阶导数。

\begin{enumerate}
\def\labelenumi{\arabic{enumi}.}
\tightlist
\item
  若 \(\forall x \in I, f''(x) > 0\)，则 \(f\) 在 \(I\) 上是凸函数。
\item
  若 \(\forall x \in I, f''(x) < 0\)，则 \(f\) 在 \(I\) 上是凹函数。
\end{enumerate}

\textbf{证明}：(1) 设
\(x_1, x_2 \in I\)，\(x_1 < x_2\)，\(\lambda \in (0, 1)\)，\(x_0 = \lambda x_1 + (1-\lambda)x_2\)。

在 \([x_1, x_0]\) 和 \([x_0, x_2]\) 上分别应用拉格朗日中值定理：

\[\frac{f(x_0) - f(x_1)}{x_0 - x_1} = f'(\xi_1), \quad \frac{f(x_2) - f(x_0)}{x_2 - x_0} = f'(\xi_2)\]

其中 \(x_1 < \xi_1 < x_0 < \xi_2 < x_2\)。由于 \(f'' > 0\)，\(f'\)
严格递增，故 \(f'(\xi_1) < f'(\xi_2)\)，即

\[\frac{f(x_0) - f(x_1)}{x_0 - x_1} < \frac{f(x_2) - f(x_0)}{x_2 - x_0}\]

由此可推得 \(f(x_0) < \lambda f(x_1) + (1-\lambda)f(x_2)\)。具体地，记
\(\Delta x = x_2 - x_1 > 0\)，则
\(x_0 - x_1 = (1-\lambda)\Delta x\)，\(x_2 - x_0 = \lambda \Delta x\)。交叉相乘得：

\[\lambda[f(x_0) - f(x_1)] < (1-\lambda)[f(x_2) - f(x_0)]\]

展开并整理：\(\lambda f(x_0) + (1-\lambda)f(x_0) < \lambda f(x_1) + (1-\lambda)f(x_2)\)，即
\(f(x_0) < \lambda f(x_1) + (1-\lambda)f(x_2)\)。\(\blacksquare\)

\paragraph{23.12.3 拐点}\label{ux62d0ux70b9}

\textbf{定义 23.8}（拐点）若曲线 \(y = f(x)\) 在点 \((x_0, f(x_0))\)
处凹凸性发生改变，则称此点为曲线的\textbf{拐点}。

\textbf{定理 23.29}（拐点的必要条件）若 \((x_0, f(x_0))\) 是拐点，且
\(f''(x_0)\) 存在，则 \(f''(x_0) = 0\)。

\textbf{证明}：若 \(f''(x_0) > 0\)（或 \(< 0\)），由 \(f''\)
的局部保号性，在 \(x_0\) 的某个邻域内 \(f''\)
恒正（或恒负），凹凸性不改变，矛盾。\(\blacksquare\)

\bigskip\hrule\bigskip

\subsubsection{23.13
微分中值定理的进一步应用}\label{ux5faeux5206ux4e2dux503cux5b9aux7406ux7684ux8fdbux4e00ux6b65ux5e94ux7528}

\paragraph{23.13.1
函数恒等的判定}\label{ux51fdux6570ux6052ux7b49ux7684ux5224ux5b9a}

\textbf{定理 23.30} 设 \(f\) 和 \(g\) 在 \([a, b]\) 上连续，在
\((a, b)\) 上可导，\(f(a) = g(a)\)。若 \(f'(x) = g'(x)\) 对所有
\(x \in (a, b)\) 成立，则 \(f(x) = g(x)\) 对所有 \(x \in [a, b]\) 成立。

\textbf{证明}：令 \(h = f - g\)，则 \(h' = 0\) 在 \((a, b)\)
上成立，\(h(a) = 0\)。

对任意 \(x \in (a, b]\)，由拉格朗日中值定理，存在 \(\xi \in (a, x)\)
使得 \(h(x) - h(a) = h'(\xi)(x - a) = 0\)，故
\(h(x) = 0\)。\(\blacksquare\)

\paragraph{23.13.2
不等式的证明}\label{ux4e0dux7b49ux5f0fux7684ux8bc1ux660e}

\textbf{定理 23.31} 对任意
\(x > 0\)，\(x - \frac{x^2}{2} < \ln(1+x) < x\)。

\textbf{证明}：右边不等式即 \(\ln(1+x) < x\)（\(x > 0\)）。此即推论 18.2
的严格不等式形式，其证明在 Ch 18 中基于导数给出（见定理
18.19），但结论可由对数函数的严格递增性和 \(\ln 1 = 0\) 独立建立：考虑
\(g(x) = x - \ln(1+x)\)，\(g(0) = 0\)；对有理数 \(0 < r < x\)，由
\(\ln\) 的严格递增性，\(\ln(1+x) > \ln(1+r)\)，而由有理指数幂的性质可证
\(\ln(1+r) < r < x\)，故 \(g(x) > 0\)。此处我们在 Ch 23
的框架下，也可直接利用导数给出简洁证明：\(g'(x) = 1 - \frac{1}{1+x} = \frac{x}{1+x} > 0\)（\(x > 0\)），故
\(g(x) > g(0) = 0\)。

对于左边：令 \(f(x) = \ln(1+x) - x + \frac{x^2}{2}\)。则 \(f(0) = 0\)，

\[f'(x) = \frac{1}{1+x} - 1 + x = \frac{1 - (1+x) + x(1+x)}{1+x} = \frac{x^2}{1+x} > 0 \quad (x > 0)\]

因此 \(f\) 在 \([0, +\infty)\)
上严格递增，\(f(x) > f(0) = 0\)（\(x > 0\)）。\(\blacksquare\)

\paragraph{23.13.3
拉格朗日中值定理的有限增量形式}\label{ux62c9ux683cux6717ux65e5ux4e2dux503cux5b9aux7406ux7684ux6709ux9650ux589eux91cfux5f62ux5f0f}

\textbf{定理 23.32}（有限增量形式）设 \(f\) 在 \([a, b]\) 上连续，在
\((a, b)\) 上可导，则存在 \(\xi \in (a, b)\) 使得

\[f(b) - f(a) = f'(\xi)(b - a)\]

\textbf{证明}：这正是拉格朗日中值定理（定理
23.15）的标准形式，无需额外证明。\(\blacksquare\)

此定理可理解为：函数在区间上的整体变化量等于某一点的瞬时变化率乘以区间长度。这为''用局部信息推断整体行为''提供了理论基础。
\#\# Ch 24 积分初步

\textbf{前置知识}：本章依赖 Ch
23（导数与微分，特别是基本导数公式和求导法则）、Ch
22（极限与连续，特别是最值定理和介值定理）。积分是微积分的两大核心概念之一，本章将从定积分的严格定义出发，建立牛顿-莱布尼茨公式，并介绍积分的基本计算方法。

\bigskip\hrule\bigskip

\subsubsection{24.1
定积分的概念}\label{ux5b9aux79efux5206ux7684ux6982ux5ff5}

\paragraph{24.1.1 面积问题}\label{ux9762ux79efux95eeux9898}

考虑求由曲线 \(y = f(x)\)（\(f(x) \geq 0\)）、\(x\) 轴以及直线 \(x = a\)
和 \(x = b\) 围成的区域的面积。

\textbf{策略}：将 \([a, b]\) 分成 \(n\)
个小区间，在每个小区间上用矩形近似曲边梯形的面积，取极限得到精确面积。

\paragraph{24.1.2 黎曼和}\label{ux9eceux66fcux548c}

\textbf{定义 24.1}（分割）区间 \([a, b]\) 的一个\textbf{分割} \(P\)
是一组分点 \(a = x_0 < x_1 < x_2 < \cdots < x_n = b\)。记
\(\Delta x_i = x_i - x_{i-1}\)（\(i = 1, 2, \ldots, n\)），\(\|P\| = \max_{1 \leq i \leq n} \Delta x_i\)
称为分割 \(P\) 的\textbf{模}（或\textbf{范数}）。

\textbf{定义 24.2}（黎曼和）设 \(f\) 在 \([a, b]\)
上有定义，\(P = \{x_0, x_1, \ldots, x_n\}\) 是 \([a, b]\)
的分割，\(\xi_i \in [x_{i-1}, x_i]\)（\(i = 1, 2, \ldots, n\)）。\textbf{黎曼和}定义为：

\[S(P, f, \{\xi_i\}) = \sum_{i=1}^{n} f(\xi_i) \Delta x_i\]

\paragraph{24.1.3
定积分的定义}\label{ux5b9aux79efux5206ux7684ux5b9aux4e49}

\textbf{定义 24.3}（定积分/黎曼积分）设 \(f\) 在 \([a, b]\)
上有定义。若存在实数 \(I\)，使得

\[\forall \varepsilon > 0, \exists \delta > 0, \forall P, \forall \{\xi_i\}: \|P\| < \delta \Rightarrow \left|S(P, f, \{\xi_i\}) - I\right| < \varepsilon\]

则称 \(f\) 在 \([a, b]\) 上\textbf{可积}，\(I\) 称为 \(f\) 在 \([a, b]\)
上的\textbf{定积分}，记作

\[\int_a^b f(x) \, dx = I\]

其中 \(a\) 称为\textbf{下限}，\(b\) 称为\textbf{上限}，\(f\)
称为\textbf{被积函数}，\(x\) 称为\textbf{积分变量}。

\paragraph{24.1.4
定积分的性质}\label{ux5b9aux79efux5206ux7684ux6027ux8d28}

\textbf{定理 24.1}（线性性）设 \(f\) 和 \(g\) 在 \([a, b]\)
上可积，\(\alpha, \beta \in \mathbb{R}\)，则 \(\alpha f + \beta g\) 在
\([a, b]\) 上可积，且

\[\int_a^b [\alpha f(x) + \beta g(x)] \, dx = \alpha \int_a^b f(x) \, dx + \beta \int_a^b g(x) \, dx\]

\textbf{证明}：

\[S(P, \alpha f + \beta g, \{\xi_i\}) = \sum_{i=1}^{n}[\alpha f(\xi_i) + \beta g(\xi_i)] \Delta x_i = \alpha S(P, f, \{\xi_i\}) + \beta S(P, g, \{\xi_i\})\]

取极限 \(\|P\| \to 0\)，由极限的线性性即得。\(\blacksquare\)

\textbf{定理 24.2}（区间可加性）设 \(f\) 在包含 \(a, b, c\)
的区间上可积，则

\[\int_a^b f(x) \, dx = \int_a^c f(x) \, dx + \int_c^b f(x) \, dx\]

\textbf{证明}：当 \(a < c < b\) 时，对 \([a, b]\) 的分割中加入分点
\(c\)。可以证明加入分点不会改变黎曼和的极限。\(\blacksquare\)

\textbf{定理 24.3}（保序性）设 \(f\) 和 \(g\) 在 \([a, b]\) 上可积，且
\(f(x) \geq g(x)\) 对所有 \(x \in [a, b]\) 成立（\(a < b\)），则

\[\int_a^b f(x) \, dx \geq \int_a^b g(x) \, dx\]

\textbf{证明}：\(S(P, f, \{\xi_i\}) = \sum f(\xi_i) \Delta x_i \geq \sum g(\xi_i) \Delta x_i = S(P, g, \{\xi_i\})\)。取极限即得。\(\blacksquare\)

\textbf{定理 24.4}（绝对值不等式）设 \(f\) 在 \([a, b]\) 上可积，则

\[\left|\int_a^b f(x) \, dx\right| \leq \int_a^b |f(x)| \, dx\]

\textbf{证明}：由 \(-|f(x)| \leq f(x) \leq |f(x)|\)
和保序性即得。\(\blacksquare\)

\paragraph{24.1.5
可积的充分条件}\label{ux53efux79efux7684ux5145ux5206ux6761ux4ef6}

\textbf{定理 24.5}（连续函数可积）若 \(f\) 在 \([a, b]\) 上连续，则
\(f\) 在 \([a, b]\) 上可积。

\textbf{证明概要}：由于 \(f\) 在闭区间 \([a, b]\) 上连续，由最值定理
\(f\) 有界。设 \(M_i\) 和 \(m_i\) 分别为 \(f\) 在 \([x_{i-1}, x_i]\)
上的最大值和最小值。定义\textbf{上和}和\textbf{下和}：

\[U(P, f) = \sum_{i=1}^{n} M_i \Delta x_i, \quad L(P, f) = \sum_{i=1}^{n} m_i \Delta x_i\]

由一致连续性，当 \(\|P\| \to 0\)
时，\(U(P, f) - L(P, f) \to 0\)。因此上和与下和的极限相等，此公共值即为定积分。\(\blacksquare\)

\bigskip\hrule\bigskip

\subsubsection{24.2
牛顿-莱布尼茨公式}\label{ux725bux987f-ux83b1ux5e03ux5c3cux8328ux516cux5f0f}

\paragraph{24.2.1 原函数}\label{ux539fux51fdux6570}

\textbf{定义 24.4}（原函数）设 \(f\) 在区间 \(I\) 上有定义。若函数 \(F\)
满足 \(F'(x) = f(x)\) 对所有 \(x \in I\) 成立，则称 \(F\) 为 \(f\) 在
\(I\) 上的一个\textbf{原函数}。

\textbf{定理 24.6}（原函数的唯一性，相差常数）设 \(F\) 和 \(G\) 都是
\(f\) 在 \(I\) 上的原函数，则存在常数 \(C\) 使得 \(G(x) = F(x) + C\)。

\textbf{证明}：\((G - F)'(x) = G'(x) - F'(x) = f(x) - f(x) = 0\)。

由拉格朗日中值定理的推论：若 \(h'(x) = 0\) 在区间 \(I\) 上成立，则 \(h\)
为常数。\(\blacksquare\)

\paragraph{24.2.2
牛顿-莱布尼茨公式}\label{ux725bux987f-ux83b1ux5e03ux5c3cux8328ux516cux5f0f-1}

\textbf{定理 24.7}（牛顿-莱布尼茨公式，微积分基本定理）设 \(f\) 在
\([a, b]\) 上连续，\(F\) 是 \(f\) 的任一原函数，即 \(F'(x) = f(x)\)。则

\[\int_a^b f(x) \, dx = F(b) - F(a)\]

\textbf{证明}：定义\textbf{变上限积分函数}

\[\Phi(x) = \int_a^x f(t) \, dt, \quad x \in [a, b]\]

\textbf{第一步}：证明 \(\Phi'(x) = f(x)\)。

对于 \(x \in (a, b)\) 和 \(h \neq 0\)（\(x + h \in [a, b]\)），

\[\Phi(x+h) - \Phi(x) = \int_a^{x+h} f(t) \, dt - \int_a^x f(t) \, dt = \int_x^{x+h} f(t) \, dt\]

由于 \(f\) 连续，由最值定理（定理 22.16），存在 \(m, M\) 使得
\(m \leq f(t) \leq M\) 在 \([x, x+h]\)（或 \([x+h, x]\)）上成立。故
\(mh \leq \int_x^{x+h} f(t)\,dt \leq Mh\)（\(h > 0\) 时），即
\(m \leq \frac{1}{h}\int_x^{x+h} f(t)\,dt \leq M\)。由介值定理（定理
22.15），存在 \(\xi\) 介于 \(x\) 和 \(x+h\) 之间，使得

\[\int_x^{x+h} f(t) \, dt = f(\xi) \cdot h\]

因此

\[\frac{\Phi(x+h) - \Phi(x)}{h} = f(\xi)\]

当 \(h \to 0\) 时，\(\xi \to x\)（因为 \(\xi\) 介于 \(x\) 和 \(x+h\)
之间），由 \(f\) 的连续性，

\[\Phi'(x) = \lim_{h \to 0} \frac{\Phi(x+h) - \Phi(x)}{h} = \lim_{\xi \to x} f(\xi) = f(x)\]

\textbf{第二步}：因此 \(\Phi\) 是 \(f\) 的一个原函数。由定理
24.6，任何原函数 \(F\) 满足

\[F(x) = \Phi(x) + C\]

令 \(x = a\)：\(F(a) = \Phi(a) + C = 0 + C = C\)。

令 \(x = b\)：\(F(b) = \Phi(b) + F(a)\)。

所以 \(\Phi(b) = F(b) - F(a)\)，即

\[\int_a^b f(x) \, dx = F(b) - F(a)\]

\(\blacksquare\)

\textbf{注}：这个定理深刻地联系了微分和积分------两个看似不相关的概念：微分是''局部化''（求瞬时变化率），积分是''全局化''（求累积量），而牛顿-莱布尼茨公式表明它们互为逆运算。

\textbf{推论 24.1}（积分中值定理）设 \(f\) 在 \([a, b]\) 上连续，则存在
\(\xi \in [a, b]\)，使得

\[\int_a^b f(x) \, dx = f(\xi)(b - a)\]

\textbf{证明}：由最值定理，\(f\) 在 \([a, b]\) 上取到最小值 \(m\)
和最大值 \(M\)：

\[m(b-a) \leq \int_a^b f(x) \, dx \leq M(b-a)\]

即 \(m \leq \frac{1}{b-a}\int_a^b f(x) \, dx \leq M\)。

由介值定理，存在 \(\xi \in [a, b]\) 使得
\(f(\xi) = \frac{1}{b-a}\int_a^b f(x) \, dx\)。\(\blacksquare\)

\bigskip\hrule\bigskip

\subsubsection{24.3 不定积分}\label{ux4e0dux5b9aux79efux5206}

\paragraph{24.3.1
不定积分的定义}\label{ux4e0dux5b9aux79efux5206ux7684ux5b9aux4e49}

\textbf{定义 24.5}（不定积分）函数 \(f\) 的所有原函数构成的集合称为
\(f\) 的\textbf{不定积分}，记作

\[\int f(x) \, dx = F(x) + C\]

其中 \(F'(x) = f(x)\)，\(C\) 为任意常数。

\paragraph{24.3.2
基本积分公式}\label{ux57faux672cux79efux5206ux516cux5f0f}

由基本导数公式反向得到：

\begin{enumerate}
\def\labelenumi{\arabic{enumi}.}
\tightlist
\item
  \(\int x^n \, dx = \frac{x^{n+1}}{n+1} + C\)（\(n \neq -1\)）
\item
  \(\int \frac{1}{x} \, dx = \ln|x| + C\)
\item
  \(\int e^x \, dx = e^x + C\)
\item
  \(\int a^x \, dx = \frac{a^x}{\ln a} + C\)（\(a > 0\)，\(a \neq 1\)）
\item
  \(\int \sin x \, dx = -\cos x + C\)
\item
  \(\int \cos x \, dx = \sin x + C\)
\item
  \(\int \sec^2 x \, dx = \tan x + C\)
\end{enumerate}

\textbf{证明
1}：\(\left(\frac{x^{n+1}}{n+1}\right)' = \frac{(n+1)x^n}{n+1} = x^n\)。\(\blacksquare\)

\textbf{证明 5}：\((-\cos x)' = \sin x\)。\(\blacksquare\)

\bigskip\hrule\bigskip

\subsubsection{24.4 换元积分法}\label{ux6362ux5143ux79efux5206ux6cd5}

\paragraph{24.4.1
第一换元法（凑微分法）}\label{ux7b2cux4e00ux6362ux5143ux6cd5ux51d1ux5faeux5206ux6cd5}

\textbf{定理 24.8}（第一换元法）设 \(F\) 是 \(f\) 的原函数，\(g\)
可导，则

\[\int f(g(x)) g'(x) \, dx = F(g(x)) + C\]

\textbf{证明}：由链式法则，\([F(g(x))]' = F'(g(x)) \cdot g'(x) = f(g(x)) g'(x)\)。因此
\(F(g(x))\) 是 \(f(g(x))g'(x)\) 的原函数。\(\blacksquare\)

\textbf{用法}：欲求 \(\int h(x) \, dx\)，若能将 \(h(x)\) 写成
\(f(g(x)) \cdot g'(x)\) 的形式，则令 \(u = g(x)\)，转化为
\(\int f(u) \, du\)。

\paragraph{24.4.2 第二换元法}\label{ux7b2cux4e8cux6362ux5143ux6cd5}

\textbf{定理 24.9}（第二换元法）设 \(x = \varphi(t)\)
是严格单调可导函数，\(\varphi'(t) \neq 0\)。则

\[\int f(x) \, dx = \int f(\varphi(t)) \varphi'(t) \, dt \bigg|_{t = \varphi^{-1}(x)} + C\]

\textbf{证明}：设 \(G(t) = \int f(\varphi(t)) \varphi'(t) \, dt\)，即
\(G'(t) = f(\varphi(t))\varphi'(t)\)。

令 \(F(x) = G(\varphi^{-1}(x))\)。由链式法则和反函数求导法则：

\[F'(x) = G'(\varphi^{-1}(x)) \cdot \frac{d}{dx}\varphi^{-1}(x) = f(\varphi(\varphi^{-1}(x))) \cdot \varphi'(\varphi^{-1}(x)) \cdot \frac{1}{\varphi'(\varphi^{-1}(x))} = f(x)\]

所以 \(F(x) = G(\varphi^{-1}(x))\) 是 \(f(x)\)
的原函数。\(\blacksquare\)

\bigskip\hrule\bigskip

\subsubsection{24.5 分部积分法}\label{ux5206ux90e8ux79efux5206ux6cd5}

\textbf{定理 24.10}（分部积分法）设 \(u(x)\) 和 \(v(x)\) 都可导，则

\[\int u \, dv = uv - \int v \, du\]

即

\[\int u(x) v'(x) \, dx = u(x)v(x) - \int u'(x) v(x) \, dx\]

\textbf{证明}：由乘法法则，\((uv)' = u'v + uv'\)，即
\(uv' = (uv)' - u'v\)。

两边积分：

\[\int u v' \, dx = \int (uv)' \, dx - \int u' v \, dx = uv - \int u' v \, dx\]

\(\blacksquare\)

\textbf{注}：分部积分法在处理乘积形式的被积函数时特别有用。选择 \(u\) 和
\(dv\) 的原则通常是使得 \(\int v \, du\) 比 \(\int u \, dv\)
更容易计算。

\bigskip\hrule\bigskip

\subsubsection{24.6
微积分基本定理的另一个形式}\label{ux5faeux79efux5206ux57faux672cux5b9aux7406ux7684ux53e6ux4e00ux4e2aux5f62ux5f0f}

\textbf{定理 24.11}（微积分基本定理第二部分）设 \(f\) 在 \([a, b]\)
上可积，且有原函数 \(F\)（即 \(F' = f\)），则

\[\int_a^b f(x) \, dx = F(b) - F(a)\]

\textbf{证明}：对 \([a, b]\) 的任意分割
\(a = x_0 < x_1 < \cdots < x_n = b\)，在每个子区间 \([x_{i-1}, x_i]\)
上，\(F\) 在 \([x_{i-1}, x_i]\) 上连续、在 \((x_{i-1}, x_i)\)
上可导，由拉格朗日中值定理（定理 23.15），存在
\(\xi_i \in (x_{i-1}, x_i)\) 使得

\[F(x_i) - F(x_{i-1}) = F'(\xi_i)(x_i - x_{i-1}) = f(\xi_i) \Delta x_i\]

对所有 \(i\) 求和，左端为望远镜和：

\[\sum_{i=1}^n [F(x_i) - F(x_{i-1})] = F(b) - F(a)\]

右端为 \(f\) 的一个黎曼和 \(\sum_{i=1}^n f(\xi_i) \Delta x_i\)。由于
\(f\) 在 \([a, b]\) 上可积，当分割的模 \(\|\mathcal{P}\| \to 0\)
时，黎曼和收敛到 \(\int_a^b f(x)\,dx\)。因此

\[F(b) - F(a) = \int_a^b f(x)\,dx\]

\(\blacksquare\)

\textbf{注}：此定理将定理 24.7
的连续性条件减弱为可积性。关键在于拉格朗日中值定理仅要求 \(F\)
可导（从而连续），不要求 \(f = F'\) 连续。

\bigskip\hrule\bigskip

\subsubsection{24.7
定积分的进一步应用}\label{ux5b9aux79efux5206ux7684ux8fdbux4e00ux6b65ux5e94ux7528}

\paragraph{24.7.1
平面区域的面积}\label{ux5e73ux9762ux533aux57dfux7684ux9762ux79ef}

\textbf{定理 24.12} 设 \(f\) 和 \(g\) 在 \([a, b]\) 上连续，且
\(f(x) \geq g(x)\)。则由曲线 \(y = f(x)\) 和 \(y = g(x)\) 以及直线
\(x = a\)，\(x = b\) 围成的区域的面积为

\[S = \int_a^b [f(x) - g(x)] \, dx\]

\textbf{证明}：将此区域看作 \(f(x)\) 以下区域减去 \(g(x)\)
以下区域，由积分的线性性即得。\(\blacksquare\)

\paragraph{24.7.2
旋转体的体积}\label{ux65cbux8f6cux4f53ux7684ux4f53ux79ef}

\textbf{定理 24.13} 设 \(f\) 在 \([a, b]\) 上连续，\(f(x) \geq 0\)。曲线
\(y = f(x)\) 绕 \(x\) 轴旋转一周所得旋转体的体积为

\[V = \pi \int_a^b [f(x)]^2 \, dx\]

\textbf{证明}：将 \([a, b]\) 分割为小区间
\([x_{i-1}, x_i]\)。在每个小区间上，旋转体近似为底面半径
\(f(\xi_i)\)、高 \(\Delta x_i\) 的圆柱体，体积为
\(\pi [f(\xi_i)]^2 \Delta x_i\)。

总体积近似为 \(\sum_{i=1}^{n} \pi [f(\xi_i)]^2 \Delta x_i\)。取
\(\|P\| \to 0\) 的极限，由定积分定义，

\[V = \pi \int_a^b [f(x)]^2 \, dx\]

\(\blacksquare\)

\paragraph{24.7.3 弧长}\label{ux5f27ux957f}

\textbf{定理 24.14} 设 \(f\) 在 \([a, b]\) 上有连续的导数，则曲线
\(y = f(x)\) 从 \(x = a\) 到 \(x = b\) 的弧长为

\[L = \int_a^b \sqrt{1 + [f'(x)]^2} \, dx\]

\textbf{证明}：将 \([a, b]\) 分割为小区间。在 \([x_{i-1}, x_i]\)
上，弧长近似为

\[\sqrt{(\Delta x_i)^2 + (\Delta y_i)^2} = \Delta x_i \sqrt{1 + \left(\frac{\Delta y_i}{\Delta x_i}\right)^2}\]

由拉格朗日中值定理，\(\frac{\Delta y_i}{\Delta x_i} = f'(\xi_i)\)。因此弧长近似为

\[\sum_{i=1}^{n} \sqrt{1 + [f'(\xi_i)]^2} \Delta x_i\]

取极限得 \(L = \int_a^b \sqrt{1 + [f'(x)]^2} \, dx\)。\(\blacksquare\)

\bigskip\hrule\bigskip

\subsubsection{24.8
反常积分初步}\label{ux53cdux5e38ux79efux5206ux521dux6b65}

\paragraph{24.8.1 无穷限积分}\label{ux65e0ux7a77ux9650ux79efux5206}

\textbf{定义 24.6}（无穷限积分）设 \(f\) 在 \([a, +\infty)\)
上连续。定义

\[\int_a^{+\infty} f(x) \, dx = \lim_{b \to +\infty} \int_a^b f(x) \, dx\]

若此极限存在，则称反常积分\textbf{收敛}；否则称其\textbf{发散}。

\textbf{定理 24.15} \(\int_1^{+\infty} \frac{1}{x^p} \, dx\) 当
\(p > 1\) 时收敛，当 \(p \leq 1\) 时发散。

\textbf{证明}：当 \(p \neq 1\) 时：

\[\int_1^b \frac{1}{x^p} \, dx = \frac{x^{1-p}}{1-p}\bigg|_1^b = \frac{b^{1-p} - 1}{1-p}\]

当 \(p > 1\)
时，\(1 - p < 0\)，\(b^{1-p} \to 0\)（\(b \to +\infty\)），故积分收敛于
\(\frac{1}{p-1}\)。

当 \(p < 1\) 时，\(1 - p > 0\)，\(b^{1-p} \to +\infty\)，积分发散。

当 \(p = 1\)
时，\(\int_1^b \frac{1}{x} \, dx = \ln b \to +\infty\)，积分发散。\(\blacksquare\)

\paragraph{24.8.2
无界函数的积分}\label{ux65e0ux754cux51fdux6570ux7684ux79efux5206}

\textbf{定义 24.7} 设 \(f\) 在 \((a, b]\)
上连续，\(\lim_{x \to a^+} f(x) = \infty\)。定义

\[\int_a^b f(x) \, dx = \lim_{\varepsilon \to 0^+} \int_{a+\varepsilon}^b f(x) \, dx\]

\bigskip\hrule\bigskip

\subsubsection{24.9
积分与微分的对偶关系}\label{ux79efux5206ux4e0eux5faeux5206ux7684ux5bf9ux5076ux5173ux7cfb}

\textbf{定理 24.16}（微积分基本定理的统一表述）

\textbf{第一部分}：若 \(f\) 在 \([a, b]\) 上连续，则
\(\frac{d}{dx}\int_a^x f(t) \, dt = f(x)\)。

\textbf{第二部分}：若 \(F\) 在 \([a, b]\) 上有连续导数 \(F' = f\)，则
\(\int_a^b f(x) \, dx = F(b) - F(a)\)。

\textbf{证明}：第一部分即微积分基本定理（一）（定理
24.7），第二部分即微积分基本定理（二）（定理
24.11），二者均已证明。\(\blacksquare\)

第一部分说明''积分后再微分恢复原函数''，第二部分说明''微分后再积分恢复原函数（相差端点值）``。这两部分共同表明：微分和积分是互逆的运算。

这一对偶关系是微积分的灵魂，它将两个看似不相关的数学问题------求切线斜率（局部问题）和求面积（全局问题）------统一在同一个理论框架中。

\bigskip\hrule\bigskip

至此，卷四''函数与分析''的全部九章内容已经完成。从 Ch 16
的坐标系与函数概念出发，经过基本初等函数、指数与对数、三角函数与恒等变换、数列，到极限与连续的严格化，再到导数与积分，我们构建了从中学数学到大学分析学的完整桥梁。这条逻辑链条的每一步都建立在前面的基础上，形成了一个自洽的数学体系。

\subsubsection{24.10
定积分的近似计算}\label{ux5b9aux79efux5206ux7684ux8fd1ux4f3cux8ba1ux7b97}

\paragraph{24.10.1 矩形法}\label{ux77e9ux5f62ux6cd5}

\textbf{方法}：将 \([a, b]\) 等分为 \(n\)
份，\(\Delta x = \frac{b-a}{n}\)，分点
\(x_i = a + i\Delta x\)（\(i = 0, 1, \ldots, n\)）。

\[\int_a^b f(x) \, dx \approx \Delta x \sum_{i=0}^{n-1} f(x_i) \quad \text{（左矩形法）}\]

\[\int_a^b f(x) \, dx \approx \Delta x \sum_{i=1}^{n} f(x_i) \quad \text{（右矩形法）}\]

\[\int_a^b f(x) \, dx \approx \Delta x \sum_{i=0}^{n-1} f\left(\frac{x_i + x_{i+1}}{2}\right) \quad \text{（中矩形法）}\]

\paragraph{24.10.2 梯形法}\label{ux68afux5f62ux6cd5}

\[\int_a^b f(x) \, dx \approx \frac{\Delta x}{2}\left[f(x_0) + 2\sum_{i=1}^{n-1}f(x_i) + f(x_n)\right]\]

\textbf{证明}：在每个小区间 \([x_i, x_{i+1}]\) 上，用梯形面积
\(\frac{\Delta x}{2}[f(x_i) + f(x_{i+1})]\)
近似积分。求和即得。\(\blacksquare\)

\paragraph{24.10.3 辛普森法}\label{ux8f9bux666eux68eeux6cd5}

\[\int_a^b f(x) \, dx \approx \frac{\Delta x}{3}\left[f(x_0) + 4\sum_{i \text{ 奇}}f(x_i) + 2\sum_{i \text{ 偶}, i \neq 0, n}f(x_i) + f(x_n)\right]\]

\textbf{注}：辛普森法的精度远高于梯形法。对于 \(n\) 等分，梯形法的误差为
\(O(1/n^2)\)，而辛普森法的误差为 \(O(1/n^4)\)。

\bigskip\hrule\bigskip

\subsubsection{24.11
微积分的统一视角}\label{ux5faeux79efux5206ux7684ux7edfux4e00ux89c6ux89d2}

微积分的两大运算------微分和积分------构成了数学分析的核心。从历史的角度看，微分源于求切线问题和瞬时速度问题（牛顿），积分源于求面积问题（阿基米德、黎曼）。牛顿和莱布尼茨的贡献在于发现了这两个看似无关的问题之间的深刻联系。

从现代数学的角度看，微分是线性化的思想------用线性函数（切线）局部近似非线性函数；积分是极限求和的思想------将复杂量分解为无穷多个无穷小量的和。这两种思想的统一体现在微积分基本定理中：

\[\frac{d}{dx}\int_a^x f(t) \, dt = f(x)\]

\[\int_a^b F'(x) \, dx = F(b) - F(a)\]

这一对偶关系不仅在理论上深刻，在应用上也极为广泛。从物理学中的运动方程到经济学中的边际分析，从概率论中的密度函数到工程学中的信号处理，微积分无处不在。

至此，卷四''函数与分析''全部完成。本卷从坐标系的建立出发，经过函数概念、基本初等函数、三角学、数列，到达了极限与连续理论，并最终建立了导数和积分的严格理论。这条逻辑链条体现了数学从具体到抽象、从直观到严格的发展过程，也为后续更深入的数学学习（实分析、复分析、泛函分析等）奠定了坚实基础。

\subsubsection{24.12 积分不等式}\label{ux79efux5206ux4e0dux7b49ux5f0f}

\paragraph{24.12.1
柯西-施瓦茨不等式（积分形式）}\label{ux67efux897f-ux65bdux74e6ux8328ux4e0dux7b49ux5f0fux79efux5206ux5f62ux5f0f}

\textbf{定理 24.17}（柯西-施瓦茨不等式）设 \(f\) 和 \(g\) 在 \([a, b]\)
上连续，则

\[\left(\int_a^b f(x)g(x) \, dx\right)^2 \leq \int_a^b [f(x)]^2 \, dx \cdot \int_a^b [g(x)]^2 \, dx\]

等号成立当且仅当存在常数 \(\lambda\) 使得 \(f(x) = \lambda g(x)\)（或
\(g(x) = \lambda f(x)\)）对所有 \(x \in [a, b]\) 成立。

\textbf{证明}：对任意实数 \(t\)，

\[0 \leq \int_a^b [f(x) + tg(x)]^2 \, dx = \int_a^b f^2 \, dx + 2t\int_a^b fg \, dx + t^2\int_a^b g^2 \, dx\]

设
\(A = \int_a^b g^2 \, dx \geq 0\)，\(B = \int_a^b fg \, dx\)，\(C = \int_a^b f^2 \, dx \geq 0\)。则
\(At^2 + 2Bt + C \geq 0\) 对所有 \(t\) 成立。

当 \(A > 0\) 时，此二次函数非负当且仅当判别式
\(\Delta = 4B^2 - 4AC \leq 0\)，即 \(B^2 \leq AC\)。

当 \(A = 0\) 时，\(g \equiv 0\)，不等式显然成立。

\(\blacksquare\)

\paragraph{24.12.2
詹森不等式（积分形式）}\label{ux8a79ux68eeux4e0dux7b49ux5f0fux79efux5206ux5f62ux5f0f}

\textbf{定理 24.18}（詹森不等式）设 \(f\) 是凸函数，\(g\) 在 \([a, b]\)
上可积，\(p\) 在 \([a, b]\) 上非负可积且 \(\int_a^b p(x) \, dx > 0\)。则

\[f\left(\frac{\int_a^b p(x)g(x) \, dx}{\int_a^b p(x) \, dx}\right) \leq \frac{\int_a^b p(x)f(g(x)) \, dx}{\int_a^b p(x) \, dx}\]

\textbf{证明}：将积分看作连续版本的加权平均，由 \(f\)
的凸性推广到连续情形。

详细论证：对 \([a, b]\) 的任意分割
\(a = x_0 < x_1 < \cdots < x_n = b\)，取
\(\xi_i \in [x_{i-1}, x_i]\)，令
\(\lambda_i = \frac{x_i - x_{i-1}}{b - a}\)（\(\sum \lambda_i = 1\)）。由离散詹森不等式，\(f\left(\sum_{i=1}^n \lambda_i \xi_i\right) \leq \sum_{i=1}^n \lambda_i f(\xi_i)\)。左边为
\(f\left(\frac{1}{b-a}\sum_{i=1}^n f(\xi_i)(x_i - x_{i-1})\right)\)
的黎曼和逼近，右边为
\(\frac{1}{b-a}\sum_{i=1}^n f(\xi_i)(x_i - x_{i-1})\)
的黎曼和。令分割加细（最大区间长度趋于零），取极限即得
\(f\left(\frac{1}{b-a}\int_a^b g(x)dx\right) \leq \frac{1}{b-a}\int_a^b f(g(x))dx\)。\(\blacksquare\)

\bigskip\hrule\bigskip

\subsubsection{24.13
微积分基本思想回顾}\label{ux5faeux79efux5206ux57faux672cux601dux60f3ux56deux987e}

本卷的内容从函数概念出发，经过基本初等函数的逐一研究，建立起了极限理论这一分析学的基石，进而发展出导数和积分两大核心概念。

\textbf{导数的本质}是线性逼近：在一点附近，用切线代替曲线，用线性函数代替非线性函数。这一思想的精确表述就是：

\[f(x) = f(a) + f'(a)(x - a) + o(x - a)\]

\textbf{积分的本质}是无穷累加：将连续变化的量分解为无穷多个无穷小量之和。黎曼积分的定义体现了''分割、近似、求和、取极限''的基本方法论。

\textbf{微积分基本定理}将两者统一：微分和积分互为逆运算。这不仅是计算工具，更是一种深刻的数学思想------局部性质（导数）与整体性质（积分）之间的对偶关系。

这种对偶关系在高等数学中以多种方式展现：格林公式、高斯公式、斯托克斯公式都是它的推广；傅里叶变换中的时域-频域对偶也是同一思想的体现。

\textbf{本章展望}：本章建立了定积分的严格理论和基本计算方法。积分的理论远不止于此------更深入的学习包括：黎曼可积的充要条件、反常积分（无穷区间或无界函数的积分）、重积分、线面积分等。积分理论在物理学（功、能量、质心）、概率论（概率密度函数的积分）、经济学（消费者剩余）等领域有广泛应用。

至此，卷四''函数与分析''全部完成。我们从坐标系和函数的严格定义出发，经过基本初等函数、指数与对数、三角函数、三角恒等变换、数列，建立了极限与连续理论，最终到达了微积分的两大核心概念------导数与积分。这条逻辑链条从简单的函数概念出发，一步步深入到分析学的核心，展现了数学从直观到严格、从具体到抽象的发展脉络。

\bigskip\hrule\bigskip

\bigskip\hrule\bigskip

\subsubsection{卷四总结}\label{ux5377ux56dbux603bux7ed3}

卷四''函数与分析''从坐标系和函数的严格定义出发，系统构建了以下理论体系：

\begin{itemize}
\tightlist
\item
  \textbf{Ch 16}
  建立了坐标系、函数的严格定义（集合论意义下的关系）、函数的基本性质（单调性、奇偶性、周期性、有界性）以及反函数和复合函数理论。
\item
  \textbf{Ch 17}
  研究了一次函数、反比例函数、二次函数（三种形式：一般式、顶点式、交点式）和幂函数，建立了这些基本初等函数的完整性质。
\item
  \textbf{Ch 18}
  严格定义了实数指数幂，系统研究了指数函数和对数函数的性质、运算法则和换底公式，引入了函数零点和二分法。
\item
  \textbf{Ch 19}
  将三角函数从锐角推广到任意角，通过单位圆给出了严格的函数定义，推导了同角关系和诱导公式，分析了三角函数的图像和性质。
\item
  \textbf{Ch 20}
  建立了三角恒等变换的完整体系（和差角公式、二倍角、半角、积化和差、和差化积、辅助角公式），证明了正弦定理和余弦定理，引入了反三角函数。
\item
  \textbf{Ch 21}
  系统研究了等差数列和等比数列的通项公式与求和公式，介绍了数列求和的各种方法、递推数列和数学归纳法。
\item
  \textbf{Ch 22} 用 \(\varepsilon\)-\(N\)
  语言严格定义了数列极限，证明了极限的唯一性、四则运算、夹逼定理和单调有界定理；用
  \(\varepsilon\)-\(\delta\) 语言定义了函数极限，证明了重要极限
  \(\lim_{x \to 0}\frac{\sin x}{x} = 1\)；建立了连续函数的介值定理和最值定理。
\item
  \textbf{Ch 23}
  定义了导数，证明了可导必连续，用定义推导了基本导数公式，证明了四则运算和链式法则，建立了罗尔、拉格朗日、柯西中值定理，介绍了洛必达法则和极值理论。
\item
  \textbf{Ch 24}
  用黎曼和严格定义了定积分，证明了牛顿-莱布尼茨公式（微积分基本定理），介绍了换元积分法和分部积分法。
\end{itemize}

这九章构成了从中学数学到大学分析学的完整桥梁，每一章都以前面章节的结论为基础，形成严密的逻辑链条。

\textbf{本章展望}：本章建立了导数的严格理论。在下一章（Ch
24）中，我们将研究微积分的另一大核心概念------积分。积分是导数的逆运算，牛顿-莱布尼茨公式将两者深刻地联系在一起，构成微积分基本定理。

\bigskip\hrule\bigskip

\newpage

\section{卷五：离散数学与概率统计}\label{ux5377ux4e94ux79bbux6563ux6570ux5b66ux4e0eux6982ux7387ux7edfux8ba1}

\subsection{Ch 28 计数原理}\label{ch-28-ux8ba1ux6570ux539fux7406}

\subsubsection{28.1
加法原理与乘法原理}\label{ux52a0ux6cd5ux539fux7406ux4e0eux4e58ux6cd5ux539fux7406}

\begin{quote}
\textbf{直觉与动机}
计数原理是组合数学的基石。加法原理对应''或''的逻辑（分类讨论），乘法原理对应''且''的逻辑（分步完成）。它们看似显而易见，却是解决一切复杂计数问题的基础。
\end{quote}

\textbf{原理 28.1（加法原理/分类计数原理）} 完成一件事有 \(n\)
类方法。第 \(i\) 类方法有 \(m_i\)
种方案（\(i = 1, 2, \ldots, n\)），且各类方案互不重叠，则完成这件事共有：
\[N = m_1 + m_2 + \cdots + m_n\] 种方法。

\textbf{原理 28.2（乘法原理/分步计数原理）} 完成一件事需要 \(n\)
个步骤。第 \(i\) 步有 \(m_i\)
种方案（\(i = 1, 2, \ldots, n\)），则完成这件事共有：
\[N = m_1 \times m_2 \times \cdots \times m_n\] 种方法。

\textbf{证明（集合语言）} 加法原理等价于：若 \(A_1, A_2, \ldots, A_n\)
是两两不交的有限集，则
\(|A_1 \cup A_2 \cup \cdots \cup A_n| = |A_1| + |A_2| + \cdots + |A_n|\)。这由集合的基数性质直接得到。

乘法原理等价于笛卡尔积的基数：\(|A_1 \times A_2 \times \cdots \times A_n| = |A_1| \cdot |A_2| \cdot \cdots \cdot |A_n|\)。\(\blacksquare\)

\subsubsection{28.2 排列}\label{ux6392ux5217}

\begin{quote}
\textbf{直觉与动机} 排列问题关注的是''顺序''。从 \(n\) 个不同元素中取出
\(m\)
个排成一列，不同的排法有多少种？这个问题在生活中随处可见：排队、安排座位、密码设置等。
\end{quote}

\textbf{定义 28.1（排列）} 从 \(n\) 个不同元素中取出
\(m\)（\(m \leq n\)）个元素，按照一定的顺序排成一列，称为从 \(n\)
个元素中取 \(m\) 个元素的一个\textbf{排列}。所有不同排列的个数记作
\(A_n^m\)（或 \(P(n, m)\)）。

\textbf{定理 28.1}
\(A_n^m = n(n-1)(n-2) \cdots (n-m+1) = \frac{n!}{(n-m)!}\)

\textbf{证明}：利用乘法原理。第一个位置有 \(n\) 种选择，第二个位置有
\(n-1\) 种选择（已选走一个），\ldots，第 \(m\) 个位置有 \(n - m + 1\)
种选择。故： \[A_n^m = n(n-1)(n-2) \cdots (n-m+1)\]

当 \(m = n\) 时，\(A_n^n = n!\)（\(n\)
个元素的\textbf{全排列}）。\(\blacksquare\)

\subsubsection{28.3 组合}\label{ux7ec4ux5408}

\begin{quote}
\textbf{直觉与动机}
与排列不同，组合只关心''选了哪些''而不关心''顺序''。从 \(n\) 个人中选
\(m\) 个人组成一个委员会，不同的选法有多少种？选出 \(\{A, B, C\}\) 和
\(\{B, C, A\}\) 是同一个组合。
\end{quote}

\textbf{定义 28.2（组合）} 从 \(n\) 个不同元素中取出
\(m\)（\(m \leq n\)）个元素组成一组（不考虑顺序），称为从 \(n\)
个元素中取 \(m\) 个元素的一个\textbf{组合}。所有不同组合的个数记作
\(C_n^m\)（或 \(\binom{n}{m}\)）。

\textbf{定理 28.2} \(C_n^m = \frac{A_n^m}{m!} = \frac{n!}{m!(n-m)!}\)

\textbf{证明}：每个包含 \(m\) 个元素的组合，将其 \(m\)
个元素进行全排列，可以得到 \(m!\)
个不同的排列。反过来，每个排列对应唯一的组合。因此
\(A_n^m = C_n^m \cdot m!\)，即
\(C_n^m = \frac{A_n^m}{m!}\)。\(\blacksquare\)

\textbf{性质 28.1} 组合数的基本性质： 1. \(C_n^m = C_n^{n-m}\)（对称性）
2. \(C_n^m = C_{n-1}^{m-1} + C_{n-1}^m\)（帕斯卡公式/递推关系）

\textbf{证明对称性}
\(C_n^m = \frac{n!}{m!(n-m)!} = \frac{n!}{(n-m)! \cdot [n-(n-m)]!} = C_n^{n-m}\)。\(\blacksquare\)

\textbf{证明帕斯卡公式} 从 \(n\) 个元素中选 \(m\) 个，考虑某个特定元素
\(a\)： - 包含 \(a\)：从剩余 \(n-1\) 个中再选 \(m-1\) 个，共
\(C_{n-1}^{m-1}\) 种。 - 不含 \(a\)：从剩余 \(n-1\) 个中选 \(m\) 个，共
\(C_{n-1}^{m}\) 种。

由加法原理，\(C_n^m = C_{n-1}^{m-1} + C_{n-1}^m\)。\(\blacksquare\)

\begin{quote}
\textbf{帕斯卡公式的意义} 帕斯卡公式正是杨辉三角（Pascal's
triangle）的递推关系。杨辉三角的第 \(n\) 行第 \(m\) 个数就是
\(C_n^m\)，每个数等于它上方两个数之和。
\end{quote}

\subsubsection{28.4 二项式定理}\label{ux4e8cux9879ux5f0fux5b9aux7406}

\textbf{定理 28.3（二项式定理）} 对任意正整数 \(n\) 和实数 \(a, b\)：
\[(a + b)^n = \sum_{k=0}^{n} C_n^k \, a^{n-k} b^k = \sum_{k=0}^{n} \binom{n}{k} a^{n-k} b^k\]

\begin{quote}
\textbf{证明思路} \((a+b)^n\) 是 \(n\) 个 \((a+b)\)
的乘积。展开时，从每个因子中选择 \(a\) 或 \(b\)，共 \(2^n\)
种选择。恰好选择 \(k\) 个 \(b\)（和 \(n-k\) 个 \(a\)）的方式有 \(C_n^k\)
种，每种贡献 \(a^{n-k}b^k\)。
\end{quote}

\textbf{证明（组合方法）} 展开
\((a+b)^n = \underbrace{(a+b)(a+b)\cdots(a+b)}_{n \text{ 个}}\)。

展开后的每一项形如
\(a^{i_1} b^{j_1} \cdot a^{i_2} b^{j_2} \cdots\)，其中每个因子贡献一个
\(a\) 或一个 \(b\)。最终得到的单项式为 \(a^{n-k} b^k\)，其中 \(k\) 表示
\(n\) 个因子中选择 \(b\) 的个数。

对于固定的 \(k\)（\(0 \leq k \leq n\)），从 \(n\) 个因子中选择 \(k\)
个贡献 \(b\) 的方式恰有 \(C_n^k\) 种，因此 \(a^{n-k}b^k\) 的系数为
\(C_n^k\)。

故： \[(a+b)^n = \sum_{k=0}^{n} C_n^k \, a^{n-k} b^k\] \(\blacksquare\)

\textbf{证明（数学归纳法）} 对 \(n\) 归纳。\(n = 1\)
时：\((a+b)^1 = a + b = \binom{1}{0}a^1 b^0 + \binom{1}{1}a^0 b^1\)。✓

假设 \(n = m\) 时成立。对 \(n = m+1\)：
\[(a+b)^{m+1} = (a+b)(a+b)^m = (a+b)\sum_{r=0}^{m}\binom{m}{r}a^{m-r}b^r\]
\[= \sum_{r=0}^{m}\binom{m}{r}a^{m+1-r}b^r + \sum_{r=0}^{m}\binom{m}{r}a^{m-r}b^{r+1}\]

在第二个求和中令 \(s = r + 1\)：
\[= \sum_{r=0}^{m}\binom{m}{r}a^{m+1-r}b^r + \sum_{s=1}^{m+1}\binom{m}{s-1}a^{m+1-s}b^s\]

合并同类项（利用帕斯卡公式
\(\binom{m}{r} + \binom{m}{r-1} = \binom{m+1}{r}\)）：
\[= a^{m+1} + \sum_{r=1}^{m}\left[\binom{m}{r} + \binom{m}{r-1}\right]a^{m+1-r}b^r + b^{m+1} = \sum_{r=0}^{m+1}\binom{m+1}{r}a^{m+1-r}b^r\]
\(\blacksquare\)

\textbf{推论 28.1（二项式系数的性质）（n ≥ 1）}

\textbf{证明}：由二项式定理
\((a+b)^n = \sum_{k=0}^{n} C_n^k a^{n-k} b^k\)：

\begin{enumerate}
\def\labelenumi{\arabic{enumi}.}
\tightlist
\item
  \(\sum_{k=0}^{n} C_n^k = 2^n\)（令 \(a = b = 1\)）
\item
  \(\sum_{k=0}^{n} (-1)^k C_n^k = 0\)（令 \(a = 1, b = -1\)）
\item
  \(C_n^0 + C_n^2 + C_n^4 + \cdots = C_n^1 + C_n^3 + C_n^5 + \cdots = 2^{n-1}\)（由
  (1) 与 (2) 两式相加、相减得） \(\blacksquare\)
\end{enumerate}

\bigskip\hrule\bigskip

\subsection{Ch 29 概率}\label{ch-29-ux6982ux7387}

\subsubsection{29.1
随机事件与概率}\label{ux968fux673aux4e8bux4ef6ux4e0eux6982ux7387}

\paragraph{29.1.1
随机试验与样本空间}\label{ux968fux673aux8bd5ux9a8cux4e0eux6837ux672cux7a7aux95f4}

\begin{quote}
\textbf{直觉与动机}
概率论研究的是不确定性。在自然界和日常生活中，许多现象的结果在事先无法确定------掷骰子的点数、明年的降雨量、放射性原子的衰变时刻。概率论为这种不确定性提供了严格的数学框架。
\end{quote}

\textbf{定义 29.1（随机试验）} 满足以下条件的试验称为\textbf{随机试验}：
1. 可以在相同条件下重复进行； 2. 所有可能结果事先已知； 3.
每次试验的具体结果事先无法确定。

\textbf{定义 29.2（样本空间）}
随机试验所有可能结果的集合称为\textbf{样本空间}，记作
\(\Omega\)。样本空间中的元素称为\textbf{样本点}，记作 \(\omega\)。

\textbf{定义 29.3（事件）} 样本空间 \(\Omega\)
的子集称为\textbf{事件}。特别地： -
\textbf{必然事件}：\(\Omega\)（一定发生的事件） -
\textbf{不可能事件}：\(\varnothing\)（一定不发生的事件） -
\textbf{基本事件}：只包含一个样本点的事件

\paragraph{29.1.2
事件的关系与运算}\label{ux4e8bux4ef6ux7684ux5173ux7cfbux4e0eux8fd0ux7b97}

设 \(A, B\) 为事件： - \textbf{包含}：\(A \subset B\) 表示''\(A\) 发生则
\(B\) 必发生'' - \textbf{和事件}：\(A \cup B\) 表示''\(A\) 与 \(B\)
至少有一个发生'' - \textbf{积事件}：\(A \cap B\) 表示''\(A\) 与 \(B\)
同时发生'' - \textbf{互斥}：\(A \cap B = \varnothing\) 表示''\(A\) 与
\(B\) 不能同时发生'' -
\textbf{对立事件}：\(\bar{A} = \Omega \setminus A\) 表示''\(A\) 不发生''

\paragraph{29.1.3
概率的公理化定义}\label{ux6982ux7387ux7684ux516cux7406ux5316ux5b9aux4e49}

\textbf{定义 29.3'（\(\sigma\)-代数）} 设 \(\Omega\)
为非空集合，\(\mathcal{F}\) 是 \(\Omega\) 的子集族。如果 \(\mathcal{F}\)
满足以下三个条件，则称 \(\mathcal{F}\) 为 \(\Omega\) 上的一个
\textbf{\(\sigma\)-代数}（或 \textbf{\(\sigma\)-域}）： 1.
\(\Omega \in \mathcal{F}\)（全集属于 \(\mathcal{F}\)） 2. 若
\(A \in \mathcal{F}\)，则
\(\bar{A} = \Omega \setminus A \in \mathcal{F}\)（对补运算封闭） 3. 若
\(A_1, A_2, \ldots \in \mathcal{F}\)，则
\(\bigcup_{i=1}^{\infty} A_i \in \mathcal{F}\)（对可数并运算封闭）

\begin{quote}
\textbf{\(\sigma\)-代数的意义}
\(\sigma\)-代数是概率论中''可测事件''的精确定义。并非 \(\Omega\)
的所有子集都能被赋予概率（在某些情形下，如选择公理构造的不可测集），\(\sigma\)-代数限定了我们能够讨论概率的事件范围。条件
1 确保必然事件可测，条件 2 确保事件的否定可测，条件 3
确保可数个事件的并可测------这恰好对应概率公理系统的需要。
\end{quote}

\textbf{定义 29.4（概率）} 设 \(\Omega\) 为样本空间，\(\mathcal{F}\) 为
\(\Omega\) 上的一个 \(\sigma\)-代数（见定义 29.3'），函数
\(P: \mathcal{F} \to [0, 1]\)
称为\textbf{概率}，如果满足以下三条公理（柯尔莫哥洛夫公理）：

\begin{enumerate}
\def\labelenumi{\arabic{enumi}.}
\tightlist
\item
  \textbf{非负性}：对任意事件 \(A\)，\(P(A) \geq 0\)
\item
  \textbf{规范性}：\(P(\Omega) = 1\)
\item
  \textbf{可列可加性}：若 \(A_1, A_2, \ldots\) 是两两互斥的事件，则
  \[P\left(\bigcup_{i=1}^{\infty} A_i\right) = \sum_{i=1}^{\infty} P(A_i)\]
\end{enumerate}

\begin{quote}
\textbf{有限可加性} 从可列可加性可以推出\textbf{有限可加性}：若
\(A_1, A_2, \ldots, A_n\) 两两互斥，则
\(P(A_1 \cup A_2 \cup \cdots \cup A_n) = P(A_1) + P(A_2) + \cdots + P(A_n)\)。这只需令
\(A_{n+1} = A_{n+2} = \cdots = \varnothing\)
即可由可列可加性得到。有限可加性在实际计算中更常用，但可列可加性在处理无穷样本空间和极限运算时不可或缺。
\end{quote}

\begin{quote}
\textbf{直觉与动机}
这三条公理是柯尔莫哥洛夫在1933年提出的，它们将概率论建立在严格的测度论基础之上。非负性保证概率不为负，规范性保证''某事发生''的概率为1，可列可加性保证互斥事件的并的概率等于各自概率之和------这正是频率解释中''比例相加''的数学化。
\end{quote}

\textbf{定理 29.1（概率的基本性质）} 1. \(P(\varnothing) = 0\) 2.
\(P(\bar{A}) = 1 - P(A)\) 3.
\(P(A \cup B) = P(A) + P(B) - P(A \cap B)\)（容斥原理） 4. 若
\(A \subset B\)，则 \(P(A) \leq P(B)\)

\textbf{证明}： 1. 取
\(A_1 = \Omega, A_2 = A_3 = \cdots = \varnothing\)，由公理3（可列可加性）得
\(P(\Omega) = P(\Omega) + \sum_{i=2}^{\infty} P(\varnothing)\)。因 \(P\)
非负，该级数收敛仅当 \(P(\varnothing) = 0\)。

\begin{enumerate}
\def\labelenumi{\arabic{enumi}.}
\setcounter{enumi}{1}
\item
  \(A \cup \bar{A} = \Omega\)，\(A \cap \bar{A} = \varnothing\)。由公理3：\(P(\Omega) = P(A) + P(\bar{A})\)，故
  \(P(\bar{A}) = 1 - P(A)\)。
\item
  \(A \cup B = A \cup (B \setminus A)\)（不交并），\(B = (A \cap B) \cup (B \setminus A)\)（不交并）。故：
  \[P(A \cup B) = P(A) + P(B \setminus A)\]
  \[P(B) = P(A \cap B) + P(B \setminus A)\] 消去
  \(P(B \setminus A)\)：\(P(A \cup B) = P(A) + P(B) - P(A \cap B)\)。
\item
  \(B = A \cup (B \setminus A)\)（不交并），\(P(B) = P(A) + P(B \setminus A) \geq P(A)\)。
  \(\blacksquare\)
\end{enumerate}

\subsubsection{29.2
古典概型与几何概型}\label{ux53e4ux5178ux6982ux578bux4e0eux51e0ux4f55ux6982ux578b}

\paragraph{29.2.1 古典概型}\label{ux53e4ux5178ux6982ux578b}

\begin{quote}
\textbf{直觉与动机}
古典概型是最简单的概率模型，适用于''等可能''的情形------如掷骰子、抽扑克牌。其核心假设是每个基本事件发生的可能性相同。
\end{quote}

\textbf{定义 29.5（古典概型）} 如果随机试验满足： 1. 样本空间 \(\Omega\)
只有有限个样本点； 2. 每个基本事件等可能发生（即
\(P(\{\omega_i\}) = \frac{1}{|\Omega|}\)）；

则称为\textbf{古典概型}。事件 \(A\) 的概率为：
\[P(A) = \frac{|A|}{|\Omega|} = \frac{\text{事件 } A \text{ 包含的样本点数}}{\text{样本空间的样本点总数}}\]

\textbf{例} 从 \(1, 2, \ldots, 10\) 中随机取一个数，求取到偶数的概率。

\textbf{解}
\(\Omega = \{1, 2, \ldots, 10\}\)，\(|\Omega| = 10\)。\(A = \{2, 4, 6, 8, 10\}\)，\(|A| = 5\)。\(P(A) = \frac{5}{10} = \frac{1}{2}\)。

\paragraph{29.2.2 几何概型}\label{ux51e0ux4f55ux6982ux578b}

\textbf{定义 29.6（几何概型）} 如果随机试验的结果可以用某个几何区域
\(\Omega\) 中的点表示，且等可能性意味着''落在 \(\Omega\)
内任何子区域的概率与该子区域的测度（长度、面积或体积）成正比''，则事件
\(A\) 的概率为： \[P(A) = \frac{\mu(A)}{\mu(\Omega)}\]

其中 \(\mu\) 表示相应的测度（长度、面积或体积）。

\textbf{例（贝特朗悖论的启示）} 在单位圆内随机取一条弦，求弦长超过
\(\sqrt{3}\)
的概率。此问题的不同答案取决于''随机''的具体含义（如何参数化弦的选取），这说明概率模型的精确定义至关重要。

\subsubsection{29.3
条件概率与独立性}\label{ux6761ux4ef6ux6982ux7387ux4e0eux72ecux7acbux6027}

\paragraph{29.3.1 条件概率}\label{ux6761ux4ef6ux6982ux7387}

\begin{quote}
\textbf{直觉与动机} 条件概率回答的问题是：``在已知事件 \(B\)
发生的条件下，事件 \(A\)
发生的概率是多少？''例如，已知今天是阴天，明天下雨的概率是多少？条件概率通过缩小样本空间来更新我们对不确定性的判断。
\end{quote}

\textbf{定义 29.7（条件概率）} 设 \(A, B\) 为事件，\(P(B) > 0\)。在事件
\(B\) 已发生的条件下，事件 \(A\) 发生的\textbf{条件概率}定义为：
\[P(A|B) = \frac{P(A \cap B)}{P(B)}\]

\textbf{定理 29.2（乘法公式）} 若 \(P(B) > 0\)，则：
\[P(A \cap B) = P(B) \cdot P(A|B)\]

一般地，若 \(P(A_1 A_2 \cdots A_{n-1}) > 0\)（以下简写
\(A_1 A_2 \cdots A_k = A_1 \cap A_2 \cap \cdots \cap A_k\)），则：
\[P(A_1 A_2 \cdots A_n) = P(A_1) \cdot P(A_2|A_1) \cdot P(A_3|A_1 A_2) \cdots P(A_n|A_1 A_2 \cdots A_{n-1})\]

\textbf{证明}：由条件概率定义，\(P(A \cap B) = P(B) \cdot P(A|B)\)。一般情形对
\(n\) 归纳可得。\(\blacksquare\)

\paragraph{29.3.2
全概率公式与贝叶斯公式}\label{ux5168ux6982ux7387ux516cux5f0fux4e0eux8d1dux53f6ux65afux516cux5f0f}

\textbf{定理 29.3（全概率公式）} 设 \(B_1, B_2, \ldots, B_n\) 是样本空间
\(\Omega\) 的一个\textbf{划分}（即 \(B_i\) 两两互斥且
\(\bigcup_{i=1}^{n} B_i = \Omega\)），且 \(P(B_i) > 0\)，则对任意事件
\(A\)： \[P(A) = \sum_{i=1}^{n} P(B_i) P(A|B_i)\]

\textbf{证明}：\(A = A \cap \Omega = A \cap \left(\bigcup_{i=1}^{n} B_i\right) = \bigcup_{i=1}^{n} (A \cap B_i)\)（不交并）。

由概率的可加性：\(P(A) = \sum_{i=1}^{n} P(A \cap B_i) = \sum_{i=1}^{n} P(B_i) P(A|B_i)\)。\(\blacksquare\)

\begin{quote}
\textbf{全概率公式的意义} 全概率公式体现了''由因溯果''的思想：将事件
\(A\) 的概率分解为各种可能原因 \(B_i\) 下 \(A\)
发生的条件概率的加权平均。这在医学诊断、质量检测等领域有广泛应用。
\end{quote}

\textbf{定理 29.4（贝叶斯公式）} 在定理 29.3 的条件下，若
\(P(A) > 0\)，则：
\[P(B_j|A) = \frac{P(B_j) P(A|B_j)}{\sum_{i=1}^{n} P(B_i) P(A|B_i)}\]

\textbf{证明}：由条件概率定义和全概率公式：
\[P(B_j|A) = \frac{P(A \cap B_j)}{P(A)} = \frac{P(B_j) P(A|B_j)}{\sum_{i=1}^{n} P(B_i) P(A|B_i)}\]
\(\blacksquare\)

\begin{quote}
\textbf{贝叶斯公式的意义} 贝叶斯公式实现了''由果溯因''的推理：已知结果
\(A\) 发生了，反推各原因 \(B_j\) 成立的概率。\(P(B_j)\)
称为\textbf{先验概率}（实验前的判断），\(P(B_j|A)\)
称为\textbf{后验概率}（获得新信息后的更新判断）。贝叶斯公式是现代统计学、机器学习和人工智能的核心理论基础之一。
\end{quote}

\paragraph{29.3.3
事件的独立性}\label{ux4e8bux4ef6ux7684ux72ecux7acbux6027}

\textbf{定义 29.8（独立事件）} 若
\(P(A \cap B) = P(A) \cdot P(B)\)，则称事件 \(A\) 与 \(B\)
\textbf{相互独立}。

\begin{quote}
\textbf{直觉与动机}
独立性的直观含义是''一个事件的发生不影响另一个事件发生的概率''。例如，掷两枚硬币，第一枚的结果不影响第二枚。注意独立性不等于互斥------互斥的两个事件不可能同时发生，而独立的两个事件可以同时发生。
\end{quote}

\textbf{定理 29.5} 若 \(A\) 与 \(B\) 独立，则： 1. \(A\) 与 \(\bar{B}\)
独立 2. \(\bar{A}\) 与 \(B\) 独立 3. \(\bar{A}\) 与 \(\bar{B}\) 独立

\textbf{证明}：1.
\(P(A \cap \bar{B}) = P(A) - P(A \cap B) = P(A) - P(A)P(B) = P(A)[1 - P(B)] = P(A)P(\bar{B})\)。

2 和 3 类似证明。\(\blacksquare\)

\subsubsection{29.4
随机变量与常见分布}\label{ux968fux673aux53d8ux91cfux4e0eux5e38ux89c1ux5206ux5e03}

\paragraph{29.4.1 随机变量}\label{ux968fux673aux53d8ux91cf}

\begin{quote}
\textbf{直觉与动机}
随机变量是将随机试验的结果''数值化''的工具。例如，掷骰子的结果可以表示为随机变量
\(X\)（取值
\(1, 2, \ldots, 6\)）。引入随机变量后，我们可以用函数、积分等分析工具来研究概率问题。
\end{quote}

\textbf{定义 29.9（随机变量）} 定义在样本空间 \(\Omega\) 上的实值函数
\(X: \Omega \to \mathbb{R}\) 称为\textbf{随机变量}。

\textbf{定义 29.10（离散型随机变量的分布列）} 设离散型随机变量 \(X\)
的可能取值为 \(x_1, x_2, \ldots\)，则
\(P(X = x_k) = p_k\)（\(k = 1, 2, \ldots\)）称为 \(X\)
的\textbf{分布列}（或概率分布）。其中 \(p_k \geq 0\) 且
\(\sum_k p_k = 1\)。

\paragraph{29.4.2 数学期望}\label{ux6570ux5b66ux671fux671b}

\textbf{定义 29.11（数学期望）} 设离散型随机变量 \(X\) 的分布列为
\(P(X = x_k) = p_k\)，则 \(X\) 的\textbf{数学期望}（或均值）定义为：
\[E(X) = \sum_{k} x_k p_k\]

（要求级数绝对收敛。）

\begin{quote}
\textbf{直觉与动机}
数学期望是随机变量的''加权平均值''，它代表了随机变量在大量重复试验中的平均表现。例如，掷一枚均匀骰子的期望值为
\(\frac{1+2+3+4+5+6}{6} = 3.5\)------虽然3.5不是任何一次掷骰子的实际结果，但长期来看平均值会趋近于3.5。
\end{quote}

\textbf{定理 29.6（期望的性质）} 1. \(E(c) = c\)（常数的期望为自身） 2.
\(E(aX + b) = aE(X) + b\)（线性性） 3.
\(E(X + Y) = E(X) + E(Y)\)（可加性，不要求独立） 4. 若 \(X, Y\) 独立，则
\(E(XY) = E(X)E(Y)\)

\textbf{证明}： 1. \(E(c) = c \cdot 1 = c\)。

\begin{enumerate}
\def\labelenumi{\arabic{enumi}.}
\setcounter{enumi}{1}
\item
  \(E(aX + b) = \sum_k (ax_k + b)p_k = a\sum_k x_k p_k + b\sum_k p_k = aE(X) + b\)。
\item
  设 \((X, Y)\) 的联合分布为 \(P(X = x_i, Y = y_j) = p_{ij}\)，则：
  \[E(X + Y) = \sum_{i,j}(x_i + y_j)p_{ij}\]
\end{enumerate}

\textbf{绝对收敛性验证}：由 \(E(X)\) 和 \(E(Y)\) 存在（即
\(\sum_i |x_i| P(X = x_i) < \infty\) 和
\(\sum_j |y_j| P(Y = y_j) < \infty\)），我们有：
\[\sum_{i,j} |x_i + y_j| p_{ij} \leq \sum_{i,j} (|x_i| + |y_j|) p_{ij} = \sum_i |x_i| \sum_j p_{ij} + \sum_j |y_j| \sum_i p_{ij} = \sum_i |x_i| P(X = x_i) + \sum_j |y_j| P(Y = y_j) < \infty\]

因此 \(\sum_{i,j}(x_i + y_j)p_{ij}\) 绝对收敛，可以拆分求和：
\[E(X + Y) = \sum_i x_i \sum_j p_{ij} + \sum_j y_j \sum_i p_{ij} = \sum_i x_i P(X = x_i) + \sum_j y_j P(Y = y_j) = E(X) + E(Y)\]
\(\blacksquare\)

\begin{enumerate}
\def\labelenumi{\arabic{enumi}.}
\setcounter{enumi}{3}
\tightlist
\item
  若 \(X, Y\) 独立，则 \(E(XY) = E(X)E(Y)\)。
\end{enumerate}

\textbf{证明}：设 \(P(X = x_i, Y = y_j) = p_{ij}\)。独立性意味着
\(p_{ij} = P(X = x_i)P(Y = y_j) = p_i^{(X)} p_j^{(Y)}\)。

\textbf{绝对收敛性验证}：由 \(E(X)\) 和 \(E(Y)\)
存在，\(\sum_i |x_i| p_i^{(X)} < \infty\) 和
\(\sum_j |y_j| p_j^{(Y)} < \infty\)，故：
\[\sum_{i,j} |x_i y_j| p_{ij} = \sum_{i,j} |x_i| |y_j| p_i^{(X)} p_j^{(Y)} = \left(\sum_i |x_i| p_i^{(X)}\right)\left(\sum_j |y_j| p_j^{(Y)}\right) < \infty\]

因此 \(\sum_{i,j} x_i y_j p_{ij}\) 绝对收敛：
\[E(XY) = \sum_{i,j} x_i y_j p_i^{(X)} p_j^{(Y)} = \left(\sum_i x_i p_i^{(X)}\right)\left(\sum_j y_j p_j^{(Y)}\right) = E(X)E(Y)\]
\(\blacksquare\)

\paragraph{29.4.3 方差}\label{ux65b9ux5dee}

\textbf{定义 29.12（方差）} 随机变量 \(X\) 的\textbf{方差}定义为：
\[D(X) = E[(X - E(X))^2] = E(X^2) - [E(X)]^2\]

\(\sqrt{D(X)}\) 称为 \(X\) 的\textbf{标准差}。

\textbf{证明}：\(D(X) = E[(X - E(X))^2] = E[X^2 - 2XE(X) + (E(X))^2] = E(X^2) - 2E(X) \cdot E(X) + [E(X)]^2 = E(X^2) - [E(X)]^2\)。\(\blacksquare\)

\begin{quote}
\textbf{方差的意义}
方差衡量随机变量取值的''离散程度''------方差越大，取值越分散。例如，两个投资方案的期望收益相同，但方差较大的方案风险更高。标准差与原变量同量纲，更便于直观理解。
\end{quote}

\textbf{定理 29.7（方差的性质）} 1. \(D(c) = 0\) 2.
\(D(aX + b) = a^2 D(X)\) 3. 若 \(X, Y\) 独立，则
\(D(X + Y) = D(X) + D(Y)\)

\textbf{证明}： 1. \(D(c) = E(c^2) - [E(c)]^2 = c^2 - c^2 = 0\)。

\begin{enumerate}
\def\labelenumi{\arabic{enumi}.}
\setcounter{enumi}{1}
\item
  \(D(aX + b) = E[(aX + b - aE(X) - b)^2] = E[a^2(X - E(X))^2] = a^2 D(X)\)。
\item
  \(D(X + Y) = E[(X+Y)^2] - [E(X+Y)]^2 = E(X^2) + 2E(XY) + E(Y^2) - [E(X)]^2 - 2E(X)E(Y) - [E(Y)]^2\)。
\end{enumerate}

当 \(X, Y\) 独立时，\(E(XY) = E(X)E(Y)\)（由定理 29.6 性质 4），故：
\[D(X + Y) = [E(X^2) - (E(X))^2] + [E(Y^2) - (E(Y))^2] = D(X) + D(Y)\]
\(\blacksquare\)

\paragraph{29.4.4
常见离散分布}\label{ux5e38ux89c1ux79bbux6563ux5206ux5e03}

\textbf{定义 29.13（伯努利分布/0-1分布）} 随机变量 \(X\) 只取 \(0\) 和
\(1\) 两个值：
\[P(X = 1) = p, \quad P(X = 0) = 1 - p \quad (0 < p < 1)\]

记作 \(X \sim B(1, p)\)。则 \(E(X) = p\)，\(D(X) = p(1-p)\)。

\textbf{定义 29.14（二项分布）} 在 \(n\) 次独立伯努利试验中，成功次数
\(X\) 服从\textbf{二项分布}：
\[P(X = k) = C_n^k p^k (1-p)^{n-k}, \quad k = 0, 1, \ldots, n\]

记作 \(X \sim B(n, p)\)。

\begin{quote}
\textbf{直觉与动机} 二项分布是最重要的离散分布之一。它描述了''做 \(n\)
次独立试验，每次成功概率为
\(p\)``的场景。例如，抛10次硬币正面朝上的次数、检验100件产品的合格数等。
\end{quote}

\textbf{定理 29.8} 若 \(X \sim B(n, p)\)，则
\(E(X) = np\)，\(D(X) = np(1-p)\)。

\textbf{证明}：\(X = X_1 + X_2 + \cdots + X_n\)，其中
\(X_i \sim B(1, p)\) 独立（\(X_i\) 表示第 \(i\) 次试验的结果）。

\[E(X) = \sum_{i=1}^{n} E(X_i) = np\]

\[D(X) = \sum_{i=1}^{n} D(X_i) = np(1-p)\] \(\blacksquare\)

\textbf{定义 29.15（超几何分布）} 设 \(N\) 个产品中有 \(M\)
个次品，从中取 \(n\) 个（不放回），次品数 \(X\)
服从\textbf{超几何分布}：
\[P(X = k) = \frac{C_M^k C_{N-M}^{n-k}}{C_N^n}\]

\textbf{定义 29.16（泊松分布）} 随机变量 \(X\) 服从参数为
\(\lambda\)（\(\lambda > 0\)）的\textbf{泊松分布}：
\[P(X = k) = \frac{\lambda^k e^{-\lambda}}{k!}, \quad k = 0, 1, 2, \ldots\]

记作 \(X \sim P(\lambda)\)。则 \(E(X) = \lambda\)，\(D(X) = \lambda\)。

\begin{quote}
\textbf{泊松分布的意义}
泊松分布描述稀有事件在固定时间段内发生的次数。例如，每分钟到达的顾客数、每页书中的错别字数、每平方公里内的交通事故数。当
\(n\) 很大、\(p\) 很小、\(np = \lambda\) 适中时，二项分布 \(B(n, p)\)
近似于泊松分布 \(P(\lambda)\)。
\end{quote}

\bigskip\hrule\bigskip

\subsection{Ch 30 统计}\label{ch-30-ux7edfux8ba1}

\subsubsection{30.1 描述性统计}\label{ux63cfux8ff0ux6027ux7edfux8ba1}

\begin{quote}
\textbf{直觉与动机}
统计学是关于数据的收集、整理、分析和推断的科学。描述性统计是统计分析的第一步，它用简洁的数值和图表概括数据的主要特征，帮助我们快速了解数据的全貌。
\end{quote}

\paragraph{30.1.1
集中趋势的度量}\label{ux96c6ux4e2dux8d8bux52bfux7684ux5ea6ux91cf}

\textbf{定义 30.1（算术平均数）} 给定数据
\(x_1, x_2, \ldots, x_n\)，其\textbf{算术平均数}为：
\[\bar{x} = \frac{1}{n}\sum_{i=1}^{n} x_i\]

\textbf{定义 30.2（中位数）} 将数据 \(x_1, x_2, \ldots, x_n\)
按从小到大排列为
\(x_{(1)} \leq x_{(2)} \leq \cdots \leq x_{(n)}\)，则\textbf{中位数}为：
\[\tilde{x} = \begin{cases} x_{(\frac{n+1}{2})} & n \text{ 为奇数} \\ \frac{1}{2}\left(x_{(\frac{n}{2})} + x_{(\frac{n}{2}+1)}\right) & n \text{ 为偶数} \end{cases}\]

\textbf{定义 30.3（众数）} 数据中出现频率最高的值称为\textbf{众数}。

\begin{quote}
\textbf{三种平均数的比较}
均值利用了所有数据信息，但对极端值（离群值）敏感；中位数对离群值稳健，适合偏态分布；众数适用于分类数据。例如，收入分布通常用中位数而非均值来描述''典型''水平，因为少数极高收入会大幅拉高均值。
\end{quote}

\paragraph{30.1.2
离散趋势的度量}\label{ux79bbux6563ux8d8bux52bfux7684ux5ea6ux91cf}

\textbf{定义 30.4（方差与标准差）} 给定数据
\(x_1, x_2, \ldots, x_n\)，其\textbf{样本方差}为：
\[s^2 = \frac{1}{n-1}\sum_{i=1}^{n}(x_i - \bar{x})^2\]

\textbf{样本标准差}为 \(s = \sqrt{s^2}\)。

\begin{quote}
\textbf{为什么除以 \(n-1\) 而非 \(n\)？} 当我们用样本均值 \(\bar{x}\)
代替总体均值 \(\mu\) 时，会损失一个自由度。除以 \(n-1\)（而非
\(n\)）使得 \(s^2\) 是总体方差 \(\sigma^2\) 的\textbf{无偏估计}，即
\(E(s^2) = \sigma^2\)。
\end{quote}

\textbf{定理 30.1（方差的计算公式）}
\[s^2 = \frac{1}{n-1}\left[\sum_{i=1}^{n} x_i^2 - n\bar{x}^2\right]\]

\textbf{证明}：
\[\sum_{i=1}^{n}(x_i - \bar{x})^2 = \sum_{i=1}^{n} x_i^2 - 2\bar{x}\sum_{i=1}^{n} x_i + n\bar{x}^2 = \sum_{i=1}^{n} x_i^2 - 2n\bar{x}^2 + n\bar{x}^2 = \sum_{i=1}^{n} x_i^2 - n\bar{x}^2\]
\(\blacksquare\)

\paragraph{30.1.3 图表表示}\label{ux56feux8868ux8868ux793a}

常用的统计图表包括： -
\textbf{频率分布表}：将数据分组，统计各组的频数和频率 -
\textbf{直方图}：用矩形面积表示频率分布 -
\textbf{茎叶图}：保留原始数据信息的简易图示 -
\textbf{箱线图}：展示中位数、四分位数和离群值

\subsubsection{30.2 抽样方法}\label{ux62bdux6837ux65b9ux6cd5}

\begin{quote}
\textbf{直觉与动机}
在实际问题中，我们通常无法研究整个总体（如全国人口的身高），只能从总体中抽取一部分样本进行研究。抽样方法的合理性直接影响统计推断的可靠性。
\end{quote}

\paragraph{30.2.1
简单随机抽样}\label{ux7b80ux5355ux968fux673aux62bdux6837}

\textbf{定义 30.5（简单随机抽样）} 从包含 \(N\)
个个体的总体中，每次抽取一个个体，抽取 \(n\)
次（放回或不放回），使得每个大小为 \(n\) 的样本被抽到的概率相等。

\begin{itemize}
\tightlist
\item
  \textbf{放回抽样}：每次抽取后将个体放回总体，每次抽取独立。
\item
  \textbf{不放回抽样}：每次抽取后不放回，后续抽取的概率会变化。
\end{itemize}

\paragraph{30.2.2 系统抽样}\label{ux7cfbux7edfux62bdux6837}

\textbf{方法} 将 \(N\) 个个体按某种顺序排列，确定抽样间隔
\(k = \lfloor N/n \rfloor\)，先在 \(1, 2, \ldots, k\)
中随机抽取一个起始编号，然后每隔 \(k\) 个个体抽取一个。

\paragraph{30.2.3 分层抽样}\label{ux5206ux5c42ux62bdux6837}

\textbf{方法}
将总体按某个特征（如年龄、性别）分成若干层，然后在每层内独立进行简单随机抽样。

\begin{quote}
\textbf{分层抽样的优势}
当总体中各层之间的差异较大，而层内差异较小时，分层抽样能显著提高估计精度。例如，在全国调查中，按省份分层抽样可以保证各省份都有代表。
\end{quote}

\subsubsection{30.3 回归分析}\label{ux56deux5f52ux5206ux6790}

\paragraph{30.3.1
相关关系与散点图}\label{ux76f8ux5173ux5173ux7cfbux4e0eux6563ux70b9ux56fe}

\begin{quote}
\textbf{直觉与动机}
在现实中，两个变量之间往往存在某种关联，但并非严格的函数关系。例如，身高与体重正相关，但给定身高无法精确确定体重。回归分析的任务是寻找变量之间的统计规律。
\end{quote}

\textbf{定义 30.6（相关系数）} 给定数据对
\((x_1, y_1), (x_2, y_2), \ldots, (x_n, y_n)\)，\textbf{样本相关系数}定义为：
\[r = \frac{\sum_{i=1}^{n}(x_i - \bar{x})(y_i - \bar{y})}{\sqrt{\sum_{i=1}^{n}(x_i - \bar{x})^2 \cdot \sum_{i=1}^{n}(y_i - \bar{y})^2}}\]

其中 \(|r| \leq 1\)。\(|r|\) 越接近 \(1\)，线性相关程度越强；\(|r|\)
越接近 \(0\)，线性相关程度越弱。

\textbf{证明 \(|r| \leq 1\)} 令
\(a_i = x_i - \bar{x}\)，\(b_i = y_i - \bar{y}\)。由柯西-施瓦茨不等式：
\[\left(\sum_{i=1}^{n} a_i b_i\right)^2 \leq \left(\sum_{i=1}^{n} a_i^2\right)\left(\sum_{i=1}^{n} b_i^2\right)\]

两边开方即得 \(|r| \leq 1\)。\(\blacksquare\)

\paragraph{30.3.2 最小二乘法}\label{ux6700ux5c0fux4e8cux4e58ux6cd5}

\textbf{问题} 给定数据
\((x_1, y_1), \ldots, (x_n, y_n)\)，求线性回归方程
\(\hat{y} = \hat{b}x + \hat{a}\)，使得残差平方和最小：
\[Q(a, b) = \sum_{i=1}^{n}(y_i - bx_i - a)^2 \to \min\]

\textbf{定理 30.2（最小二乘估计）} 使 \(Q(a, b)\) 最小的参数为：
\[\hat{b} = \frac{\sum_{i=1}^{n}(x_i - \bar{x})(y_i - \bar{y})}{\sum_{i=1}^{n}(x_i - \bar{x})^2} = \frac{\sum_{i=1}^{n} x_i y_i - n\bar{x}\bar{y}}{\sum_{i=1}^{n} x_i^2 - n\bar{x}^2}\]
\[\hat{a} = \bar{y} - \hat{b}\bar{x}\]

\begin{quote}
\textbf{证明思路} \(Q(a, b)\) 是关于 \(a\) 和 \(b\)
的二次函数，通过求偏导并令其为零（或配方法）找到最小值点。这是微积分中求极值的典型应用。
\end{quote}

\textbf{证明}：将 \(Q(a, b)\) 对 \(a\) 和 \(b\) 分别求偏导数并令其为零：

\[\frac{\partial Q}{\partial a} = -2\sum_{i=1}^{n}(y_i - bx_i - a) = 0 \implies \sum_{i=1}^{n} y_i = b\sum_{i=1}^{n} x_i + na\]
\[n\bar{y} = bn\bar{x} + na \implies \hat{a} = \bar{y} - \hat{b}\bar{x}\]

\[\frac{\partial Q}{\partial b} = -2\sum_{i=1}^{n} x_i(y_i - bx_i - a) = 0 \implies \sum_{i=1}^{n} x_i y_i = b\sum_{i=1}^{n} x_i^2 + a\sum_{i=1}^{n} x_i\]

将 \(\hat{a} = \bar{y} - \hat{b}\bar{x}\) 代入：
\[\sum_{i=1}^{n} x_i y_i = \hat{b}\sum_{i=1}^{n} x_i^2 + (\bar{y} - \hat{b}\bar{x})n\bar{x}\]
\[\sum_{i=1}^{n} x_i y_i - n\bar{x}\bar{y} = \hat{b}\left(\sum_{i=1}^{n} x_i^2 - n\bar{x}^2\right)\]
\[\hat{b} = \frac{\sum_{i=1}^{n} x_i y_i - n\bar{x}\bar{y}}{\sum_{i=1}^{n} x_i^2 - n\bar{x}^2}\]

验证二阶条件确实为最小值（Hessian矩阵正定），从略。\(\blacksquare\)

\begin{quote}
\textbf{回归方程的几何意义} 回归直线 \(\hat{y} = \hat{b}x + \hat{a}\)
必过样本中心点 \((\bar{x}, \bar{y})\)，因为
\(\hat{a} = \bar{y} - \hat{b}\bar{x}\)。回归系数 \(\hat{b}\) 表示 \(x\)
每增加一个单位时，\(y\) 的平均变化量。
\end{quote}

\subsubsection{30.4 独立性检验}\label{ux72ecux7acbux6027ux68c0ux9a8c}

\paragraph{30.4.1
列联表与卡方检验}\label{ux5217ux8054ux8868ux4e0eux5361ux65b9ux68c0ux9a8c}

\begin{quote}
\textbf{直觉与动机}
在实际问题中，我们常常需要判断两个分类变量之间是否存在关联。例如，吸烟与肺癌是否有关？性别与专业选择是否有关？独立性检验提供了一种基于数据的统计判断方法。
\end{quote}

\textbf{问题} 设有两个分类变量 \(X\)（取值 \(A_1, A_2, \ldots, A_r\)）和
\(Y\)（取值 \(B_1, B_2, \ldots, B_c\)）。从总体中抽取 \(n\) 个样本，按
\(X\) 和 \(Y\) 的取值交叉分类，得到 \(r \times c\) \textbf{列联表}，其中
\(n_{ij}\) 表示 \(X = A_i\) 且 \(Y = B_j\) 的观测频数。

\textbf{原假设} \(H_0\)：\(X\) 与 \(Y\) 独立（无关联）。

在 \(H_0\)
下，\(P(X = A_i, Y = B_j) = P(X = A_i)P(Y = B_j)\)。用样本估计：
\[\hat{p}_{ij} = \frac{n_{i\cdot}}{n} \cdot \frac{n_{\cdot j}}{n}\]

其中
\(n_{i\cdot} = \sum_j n_{ij}\)，\(n_{\cdot j} = \sum_i n_{ij}\)。\textbf{期望频数}为：
\[e_{ij} = n \hat{p}_{ij} = \frac{n_{i\cdot} \cdot n_{\cdot j}}{n}\]

\textbf{定义 30.7（卡方统计量）} \textbf{Pearson 卡方统计量}定义为：
\[\chi^2 = \sum_{i=1}^{r} \sum_{j=1}^{c} \frac{(n_{ij} - e_{ij})^2}{e_{ij}}\]

\textbf{卡方统计量的推导} 在 \(H_0\) 下，考虑 \((n_{ij} - e_{ij})\)
的标准化。对每个 \((i, j)\) 单元格：

\[\frac{(n_{ij} - e_{ij})^2}{e_{ij}}\]

度量了观测值 \(n_{ij}\) 与期望值 \(e_{ij}\) 之间的偏离程度。\(\chi^2\)
是所有单元格偏离程度之和。

\begin{quote}
\textbf{推导概要（卡方分布的来源）} 下面简要说明卡方分布的来历。记 \(k\)
个独立的标准正态随机变量 \(Z_1, Z_2, \ldots, Z_k \sim N(0, 1)\)，定义
\(X = Z_1^2 + Z_2^2 + \cdots + Z_k^2\)。\(X\) 的特征函数为

\[\varphi_X(t) = \mathbb{E}[e^{itX}] = \prod_{j=1}^{k} \mathbb{E}[e^{itZ_j^2}] = \prod_{j=1}^{k} \frac{1}{\sqrt{2\pi}} \int_{-\infty}^{\infty} e^{itz^2} e^{-z^2/2}\, dz\]

计算
\(\mathbb{E}[e^{itZ_j^2}] = \frac{1}{\sqrt{2\pi}} \int_{-\infty}^{\infty} \exp\left(-\frac{(1 - 2it)z^2}{2}\right) dz\)。令
\(u = z\sqrt{1 - 2it}\)，得

\[\mathbb{E}[e^{itZ_j^2}] = \frac{1}{\sqrt{1 - 2it}}\]

故 \(\varphi_X(t) = (1 - 2it)^{-k/2}\)。这正是
\(\Gamma(k/2, 2)\)（形状参数 \(k/2\)，尺度参数 \(2\) 的 Gamma
分布）的特征函数，其密度函数为

\[f_X(x) = \frac{1}{2^{k/2} \Gamma(k/2)} x^{k/2 - 1} e^{-x/2}, \quad x > 0\]

此即\textbf{自由度为 \(k\) 的卡方分布 \(\chi^2(k)\)} 的密度函数。

回到列联表情形：在 \(H_0\) 下，\((n_{ij} - e_{ij}) / \sqrt{e_{ij}}\)
渐近服从均值为零、协方差矩阵为 \(\Sigma\) 的多元正态分布，其中
\(\Sigma\) 受行和与列和约束的限制，秩为 \((r-1)(c-1)\)。\(\chi^2\)
统计量是这些标准化残差的平方和，在对 \(\Sigma\)
的广义逆进行适当的对角化（利用 Cochran 定理）后，该平方和收敛于
\(\chi^2((r-1)(c-1))\)。
\end{quote}

\textbf{定理 30.3（卡方统计量的渐近分布）} 当样本量 \(n \to \infty\)
时，在 \(H_0\) 成立的条件下，\(\chi^2\) 渐近服从自由度为 \((r-1)(c-1)\)
的卡方分布，记作 \(\chi^2 \sim \chi^2((r-1)(c-1))\)。

\begin{quote}
\textbf{自由度的解释} \(r \times c\) 列联表有 \(rc\)
个单元格，但由于行和 \(\sum_j n_{ij} = n_{i\cdot}\) 固定（\(r\)
个约束）且列和 \(\sum_i n_{ij} = n_{\cdot j}\) 固定（\(c\)
个约束），但总和 \(\sum_{i,j} n_{ij} = n\)
固定使得其中一个约束是冗余的，故独立约束数为 \(r + c - 1\)。自由度
\(= rc - 1 - (r + c - 1) = (r-1)(c-1)\)。
\end{quote}

\textbf{检验规则} 给定显著性水平 \(\alpha\)（通常取
\(0.05\)），查卡方分布表得临界值 \(\chi^2_\alpha((r-1)(c-1))\)。若
\(\chi^2 > \chi^2_\alpha((r-1)(c-1))\)，则拒绝 \(H_0\)，认为 \(X\) 与
\(Y\) 之间存在显著关联。

\begin{quote}
\textbf{注} 对于 \(2 \times 2\) 列联表（\(r = c = 2\)），自由度为
\(1\)。此时临界值 \(\chi^2_{0.05}(1) = 3.841\)。即当 \(\chi^2 > 3.841\)
时，在 \(5\%\) 显著性水平下拒绝独立性假设。
\end{quote}

\bigskip\hrule\bigskip

\begin{quote}
\textbf{Ch 30 前置知识补充} 本章定理 30.2
的证明使用了一元函数的导数与极值判定（见卷四 Ch
23），这些工具属于微积分基础。
\end{quote}

\bigskip\hrule\bigskip

\subsection{Ch 31 递推数列}\label{ch-31-ux9012ux63a8ux6570ux5217}

\subsubsection{31.1
递推数列的基本概念}\label{ux9012ux63a8ux6570ux5217ux7684ux57faux672cux6982ux5ff5}

\begin{quote}
\textbf{直觉与动机}
递推数列是通过相邻项之间的关系来定义的数列。与通项公式直接给出第 \(n\)
项不同，递推公式描述的是''如何从已知项推出下一项''。递推思想在计算机科学（动态规划）、金融（复利）、生物学（种群增长）等领域有广泛应用。
\end{quote}

\textbf{定义 31.1（递推数列）} 如果数列 \(\{a_n\}\)
的项之间的关系可以用方程
\[a_{n+1} = f(a_n, a_{n-1}, \ldots, a_{n-k+1})\] 表示，其中 \(f\)
是已知函数，\(k\)
称为递推的\textbf{阶数}，则称此关系为\textbf{递推公式}（或\textbf{递推关系}），\(\{a_n\}\)
称为\textbf{递推数列}。给定初始值
\(a_1, a_2, \ldots, a_k\)（\textbf{初始条件}），递推公式唯一确定数列的所有项。

\textbf{定理 31.1（递推数列的存在唯一性）} 设
\(f: \mathbb{R}^k \to \mathbb{R}\) 为已知函数，\(a_1, a_2, \ldots, a_k\)
为给定的初始值，则递推关系
\(a_{n+1} = f(a_n, a_{n-1}, \ldots, a_{n-k+1})\)（\(n \geq k\)）唯一确定一个数列
\(\{a_n\}\)。

\textbf{证明}：由递归定义原理（基于皮亚诺公理），对 \(n\)
进行归纳：\(a_1, \ldots, a_k\) 已给定。假设
\(a_1, a_2, \ldots, a_n\)（\(n \geq k\)）已确定，则
\(a_{n+1} = f(a_n, a_{n-1}, \ldots, a_{n-k+1})\)
唯一确定。\(\blacksquare\)

\subsubsection{31.2
一阶线性递推数列}\label{ux4e00ux9636ux7ebfux6027ux9012ux63a8ux6570ux5217-1}

\paragraph{31.2.1 齐次情形}\label{ux9f50ux6b21ux60c5ux5f62}

\textbf{定义 31.2（一阶线性递推数列）} 形如
\[a_{n+1} = pa_n + q \quad (p, q \text{ 为常数})\]
的递推关系称为\textbf{一阶线性递推数列}。当 \(q = 0\)
时称为\textbf{齐次}的。

\textbf{定理 31.2（齐次一阶线性递推）} 设 \(a_{n+1} = pa_n\)（\(p\)
为常数），\(a_1\) 给定，则： \[a_n = a_1 p^{n-1}\]

\textbf{证明}：反复递推：\(a_n = pa_{n-1} = p^2 a_{n-2} = \cdots = p^{n-1}a_1\)。

严格归纳：\(a_1 = a_1 p^0\) 成立。假设 \(a_k = a_1 p^{k-1}\)，则
\(a_{k+1} = pa_k = p \cdot a_1 p^{k-1} = a_1 p^k\)。\(\blacksquare\)

\paragraph{31.2.2
非齐次情形------不动点法}\label{ux975eux9f50ux6b21ux60c5ux5f62ux4e0dux52a8ux70b9ux6cd5}

\textbf{定理 31.3（一阶线性非齐次递推的通解）} 设
\(a_{n+1} = pa_n + q\)（\(p \neq 1\)），则：
\[a_n = \left(a_1 - \frac{q}{1-p}\right)p^{n-1} + \frac{q}{1-p}\]

\textbf{证明（不动点法）} 方程 \(x = px + q\) 的解（\textbf{不动点}）为
\(\lambda = \frac{q}{1-p}\)。

令 \(b_n = a_n - \lambda\)，则：
\[b_{n+1} = a_{n+1} - \lambda = (pa_n + q) - \lambda = pa_n + q - \frac{q}{1-p}\]

由 \(\lambda = p\lambda + q\)，有
\(q = \lambda - p\lambda = (1-p)\lambda\)，代入：
\[b_{n+1} = pa_n + (1-p)\lambda - \lambda = pa_n - p\lambda = p(a_n - \lambda) = pb_n\]

故 \(\{b_n\}\) 是公比为 \(p\)
的等比数列：\(b_n = b_1 p^{n-1} = (a_1 - \lambda)p^{n-1}\)。

因此
\(a_n = b_n + \lambda = (a_1 - \lambda)p^{n-1} + \lambda\)。\(\blacksquare\)

\textbf{定理 31.4（分式递推的转化）} 设递推关系为
\[a_{n+1} = \frac{pa_n + q}{ra_n + s} \quad (ps - qr \neq 0)\] 若方程
\(x = \frac{px + q}{rx + s}\) 有两个不同的不动点 \(\alpha, \beta\)，且
\(a_n \neq \beta\) 对所有 \(n\) 成立，则
\(\frac{a_n - \alpha}{a_n - \beta}\) 构成等比数列。

\textbf{证明}：设
\(b_n = \frac{a_n - \alpha}{a_n - \beta}\)。由不动点条件：\(\alpha = \frac{p\alpha + q}{r\alpha + s}\)，即
\(r\alpha^2 + s\alpha = p\alpha + q\)，类似地对 \(\beta\)。

计算
\(b_{n+1} = \frac{a_{n+1} - \alpha}{a_{n+1} - \beta} = \frac{\frac{pa_n + q}{ra_n + s} - \alpha}{\frac{pa_n + q}{ra_n + s} - \beta} = \frac{(pa_n + q) - \alpha(ra_n + s)}{(pa_n + q) - \beta(ra_n + s)} = \frac{(p - r\alpha)a_n + (q - s\alpha)}{(p - r\beta)a_n + (q - s\beta)}\)

利用
\(q - s\alpha = r\alpha^2 - p\alpha = -\alpha(p - r\alpha)\)（由不动点方程），得分子
\(= (p - r\alpha)(a_n - \alpha)\)。类似地分母
\(= (p - r\beta)(a_n - \beta)\)。

故
\(b_{n+1} = \frac{p - r\alpha}{p - r\beta} \cdot \frac{a_n - \alpha}{a_n - \beta} = \frac{p - r\alpha}{p - r\beta} \cdot b_n\)，即
\(\{b_n\}\) 是公比为 \(\frac{p - r\alpha}{p - r\beta}\) 的等比数列。

注意：若某 \(a_n = \beta\)，则 \(b_n\) 无定义，此时 \(\{a_n\}\) 从第
\(n\) 项起恒为 \(\beta\)（不动点），可直接验证
\(a_{n+1} = \frac{p\beta + q}{r\beta + s} = \beta\)。\(\blacksquare\)

\subsubsection{31.3
二阶线性递推数列}\label{ux4e8cux9636ux7ebfux6027ux9012ux63a8ux6570ux5217}

\paragraph{31.3.1 特征方程法}\label{ux7279ux5f81ux65b9ux7a0bux6cd5}

\textbf{定义 31.3（二阶线性递推数列）} 形如
\[a_{n+2} = pa_{n+1} + qa_n \quad (p, q \text{ 为常数}, q \neq 0)\]
的递推关系称为\textbf{二阶线性递推数列}。

\textbf{定理 31.5（特征方程法）} 设
\(a_{n+2} = pa_{n+1} + qa_n\)，其\textbf{特征方程}为
\[x^2 = px + q \quad \text{即} \quad x^2 - px - q = 0\]

\textbf{(i)} 若特征方程有两个不同的实根
\(\alpha \neq \beta\)，则通解为：
\[a_n = C_1 \alpha^{n-1} + C_2 \beta^{n-1}\] 其中 \(C_1, C_2\)
由初始条件 \(a_1, a_2\) 确定。

\textbf{(ii)} 若特征方程有重根 \(\alpha = \beta\)，则通解为：
\[a_n = (C_1 + C_2(n-1))\alpha^{n-1}\]

\textbf{证明（情形 i）} 将 \(a_n = r^n\) 代入递推关系
\(r^{n+2} = pr^{n+1} + qr^n\)，由 \(q \neq 0\) 知 \(r = 0\)
不是特征方程的根，故 \(r \neq 0\)，约去 \(r^n\) 得 \(r^2 = pr + q\)，即
\(r\) 为特征方程的根。

由于 \(\alpha, \beta\) 都满足递推关系（分别令 \(a_n = \alpha^{n-1}\) 和
\(a_n = \beta^{n-1}\) 验证），递推关系的线性性保证
\(a_n = C_1 \alpha^{n-1} + C_2 \beta^{n-1}\) 也满足。

初始条件给出方程组：
\[\begin{cases} C_1 + C_2 = a_1 \\ C_1 \alpha + C_2 \beta = a_2 \end{cases}\]

系数行列式
\(= \beta - \alpha \neq 0\)，故方程组有唯一解。\(\blacksquare\)

\textbf{证明（情形 ii）} 当 \(\alpha = \beta\) 时，一个解为
\(a_n = \alpha^{n-1}\)。寻找第二个线性无关解：尝试
\(a_n = (n-1)\alpha^{n-1}\)。

代入验证：需要验证
\((n+1)\alpha^{n+1} = p \cdot n\alpha^n + q \cdot (n-1)\alpha^{n-1}\)。

因为 \(\alpha\) 是特征方程 \(r^2 = pr + q\) 的重根，所以
\(r^2 - pr - q = (r - \alpha)^2 = r^2 - 2\alpha r + \alpha^2\)，对比系数得
\(p = 2\alpha\)，\(q = -\alpha^2\)。

代入右边：
\[p \cdot n\alpha^n + q \cdot (n-1)\alpha^{n-1} = 2\alpha \cdot n\alpha^n - \alpha^2 \cdot (n-1)\alpha^{n-1} = 2n\alpha^{n+1} - (n-1)\alpha^{n+1} = (n+1)\alpha^{n+1}\]

恰好等于左边。因此 \(a_n = (C_1 + C_2(n-1))\alpha^{n-1}\)
是通解。\(\blacksquare\)

\paragraph{31.3.2 应用实例}\label{ux5e94ux7528ux5b9eux4f8b}

\textbf{例（斐波那契数列）} 斐波那契数列 \(\{F_n\}\) 满足
\(F_1 = F_2 = 1\)，\(F_{n+2} = F_{n+1} + F_n\)。

特征方程：\(x^2 = x + 1\)，即 \(x^2 - x - 1 = 0\)。

解得 \(\alpha = \frac{1 + \sqrt{5}}{2}\)（黄金比
\(\varphi\)），\(\beta = \frac{1 - \sqrt{5}}{2}\)。

由初始条件：
\[\begin{cases} C_1 + C_2 = 1 \\ C_1 \alpha + C_2 \beta = 1 \end{cases}\]

解得
\(C_1 = \frac{1 - \beta}{\alpha - \beta} = \frac{\alpha}{\sqrt{5}}\)，\(C_2 = \frac{\alpha - 1}{\alpha - \beta} = -\frac{\beta}{\sqrt{5}}\)。

故：
\[F_n = \frac{\alpha}{\sqrt{5}} \cdot \alpha^{n-1} - \frac{\beta}{\sqrt{5}} \cdot \beta^{n-1} = \frac{1}{\sqrt{5}}\left[\left(\frac{1+\sqrt{5}}{2}\right)^n - \left(\frac{1-\sqrt{5}}{2}\right)^n\right]\]

这就是\textbf{比内公式}（Binet's formula）。

\textbf{注}
比内公式的推导过程本身就是特征方程法的一个完整应用。虽然公式中含有无理数
\(\sqrt{5}\)，但对所有正整数 \(n\)，\(F_n\) 都是正整数------这是因为
\(\alpha^n\) 和 \(\beta^n\) 是一对共轭无理数，无理部分恰好相消。

\subsubsection{31.4
斐波那契数列的性质}\label{ux6590ux6ce2ux90a3ux5951ux6570ux5217ux7684ux6027ux8d28}

\textbf{定义 31.4（斐波那契数列）} 满足
\(F_1 = F_2 = 1\)，\(F_{n+2} = F_{n+1} + F_n\)（\(n \geq 1\)）的数列称为\textbf{斐波那契数列}。

\textbf{性质 31.1（前 \(n\) 项和）}
\(\sum_{k=1}^{n} F_k = F_{n+2} - 1\)。

\textbf{证明}：由 \(F_k = F_{k+2} - F_{k+1}\)，求和得裂项和：
\[\sum_{k=1}^{n} F_k = \sum_{k=1}^{n}(F_{k+2} - F_{k+1}) = F_{n+2} - F_2 = F_{n+2} - 1\]
\(\blacksquare\)

\textbf{性质 31.2（相邻项比值的极限）}
\(\lim_{n \to \infty} \frac{F_{n+1}}{F_n} = \varphi = \frac{1+\sqrt{5}}{2}\)。

\textbf{证明}：由比内公式：\(\frac{F_{n+1}}{F_n} = \frac{\alpha^{n+1} - \beta^{n+1}}{\alpha^n - \beta^n} = \frac{\alpha - \beta(\beta/\alpha)^{n}}{1 - (\beta/\alpha)^{n}}\)。

由于
\(|\beta/\alpha| = \left|\frac{1-\sqrt{5}}{1+\sqrt{5}}\right| < 1\)，当
\(n \to \infty\) 时 \((\beta/\alpha)^{n} \to 0\)，故极限为
\(\alpha = \varphi\)。\(\blacksquare\)

\textbf{性质 31.3（最大公约数性质）}
\(\gcd(F_{m}, F_{n}) = F_{\gcd(m, n)}\)。

\textbf{证明}：我们先证明关键引理：\(\gcd(F_{n+1}, F_n) = 1\)（相邻斐波那契数互素）。

\textbf{引理证明}：对 \(n\)
归纳。\(\gcd(F_2, F_1) = \gcd(1, 1) = 1\)。假设
\(\gcd(F_k, F_{k-1}) = 1\)，则：
\[\gcd(F_{k+1}, F_k) = \gcd(F_k + F_{k-1}, F_k) = \gcd(F_{k-1}, F_k) = \gcd(F_k, F_{k-1}) = 1\]

其中第三步利用了 \(\gcd(a+b, b) = \gcd(a, b)\)。\(\blacksquare\)

\textbf{进一步}：利用 \(F_{n+1} = F_n + F_{n-1}\)
反复递推（辗转相除法），可得更一般的结论：\(\gcd(F_{n+1}, F_n) = \gcd(F_n, F_{n-1}) = \cdots = \gcd(F_2, F_1) = 1\)。

对一般 \(m, n\)，不妨设 \(m \geq n\)。利用恒等式
\(F_m = F_{m-n+1}F_n + F_{m-n}F_{n-1}\)。

\textbf{恒等式证明（数学归纳法）}：对 \(m-n\) 归纳。当 \(m = n\)
时：\(F_n = F_1 \cdot F_n + F_0 \cdot F_{n-1} = F_n + 0 = F_n\)。当
\(m = n+1\)
时：\(F_{n+1} = F_2 \cdot F_n + F_1 \cdot F_{n-1} = F_n + F_{n-1}\)（斐波那契递推）。假设对
\(m\) 成立，即 \(F_m = F_{m-n+1}F_n + F_{m-n}F_{n-1}\)，则对 \(m+1\)：
\[F_{m+1} = F_m + F_{m-1} = (F_{m-n+1}F_n + F_{m-n}F_{n-1}) + (F_{m-n}F_n + F_{m-n-1}F_{n-1}) = F_{m-n+2}F_n + F_{m-n+1}F_{n-1}\]
故恒等式成立。\(\blacksquare\)

由该恒等式可得
\(\gcd(F_m, F_n) = \gcd(F_n, F_{m-n})\)，这与辗转相除法的步骤完全对应，故
\(\gcd(F_m, F_n) = F_{\gcd(m,n)}\)。\(\blacksquare\)

\textbf{性质 31.4（与组合数的关系）}
\(F_n = \sum_{k=0}^{\lfloor(n-1)/2\rfloor} \binom{n-1-k}{k}\)。

\textbf{证明}：令
\(G_n = \sum_{k=0}^{\lfloor(n-1)/2\rfloor} \binom{n-1-k}{k}\)。我们验证
\(G_n\) 满足斐波那契递推关系。

\textbf{基础步}：\(G_1 = \binom{0}{0} = 1 = F_1\)，\(G_2 = \binom{1}{0} = 1 = F_2\)。

\textbf{验证递推关系}：设 \(G_1 = F_1, \ldots, G_n = F_n\)。由帕斯卡公式
\(\binom{m}{k} = \binom{m-1}{k} + \binom{m-1}{k-1}\)：

\[G_n = \sum_{k=0}^{\lfloor(n-1)/2\rfloor} \binom{n-1-k}{k}\]

将 \(\binom{n-1-k}{k} = \binom{n-2-k}{k} + \binom{n-2-k}{k-1}\)
代入，将求和拆为两部分。

\textbf{第一部分}：\(\sum_{k=0}^{\lfloor(n-1)/2\rfloor} \binom{n-2-k}{k}\)。令
\(j = k\)，此为
\(\sum_{j=0}^{\lfloor(n-1)/2\rfloor} \binom{(n-1)-1-j}{j}\)。

分析上限：\(\lfloor(n-1)/2\rfloor\)。当 \(n\) 为奇数时
\(\lfloor(n-1)/2\rfloor = (n-2)/2 = \lfloor(n-2)/2\rfloor + 1\)，但多出的项
\(\binom{(n-3)/2}{(n-1)/2} = 0\)（下标大于上标）。当 \(n\) 为偶数时
\(\lfloor(n-1)/2\rfloor = (n-2)/2 = \lfloor(n-2)/2\rfloor\)。因此此求和等于
\(G_{n-1}\)。

\textbf{第二部分}：\(\sum_{k=0}^{\lfloor(n-1)/2\rfloor} \binom{n-2-k}{k-1}\)。令
\(j = k-1\)，得
\(\sum_{j=0}^{\lfloor(n-1)/2\rfloor - 1} \binom{n-3-j}{j}\)。

上限 \(\lfloor(n-1)/2\rfloor - 1 = \lfloor(n-3)/2\rfloor\)，此求和等于
\(G_{n-2}\)。

因此
\(G_n = G_{n-1} + G_{n-2} = F_{n-1} + F_{n-2} = F_n\)。\(\blacksquare\)

\subsubsection{31.5
递推数列的求解技巧}\label{ux9012ux63a8ux6570ux5217ux7684ux6c42ux89e3ux6280ux5de7}

\paragraph{31.5.1 累加法}\label{ux7d2fux52a0ux6cd5}

\textbf{定理 31.6（累加法）} 设 \(a_{n+1} = a_n + f(n)\)，\(a_1\)
给定，则： \[a_n = a_1 + \sum_{k=1}^{n-1} f(k)\]

\textbf{证明}：由 \(a_{k+1} - a_k = f(k)\)，对 \(k = 1, 2, \ldots, n-1\)
求和（裂项和）：
\[a_n - a_1 = \sum_{k=1}^{n-1}(a_{k+1} - a_k) = \sum_{k=1}^{n-1} f(k)\]
\(\blacksquare\)

\paragraph{31.5.2 累乘法}\label{ux7d2fux4e58ux6cd5}

\textbf{定理 31.7（累乘法）} 设 \(a_{n+1} = f(n) \cdot a_n\)，\(a_1\)
给定，且所有 \(a_n \neq 0\)，则：
\[a_n = a_1 \cdot \prod_{k=1}^{n-1} f(k)\]

\textbf{证明}：由 \(\frac{a_{k+1}}{a_k} = f(k)\)，对
\(k = 1, 2, \ldots, n-1\) 连乘（裂项积）：
\[\frac{a_n}{a_1} = \prod_{k=1}^{n-1} \frac{a_{k+1}}{a_k} = \prod_{k=1}^{n-1} f(k)\]
\(\blacksquare\)

\paragraph{31.5.3 构造法}\label{ux6784ux9020ux6cd5}

\textbf{方法概述}
对于某些非线性或变系数递推，可以通过变量替换将其转化为可求解的形式。

\textbf{例 1（取倒数）} 设
\(a_{n+1} = \frac{a_n}{1 + a_n}\)，\(a_1 > 0\)。

令 \(b_n = 1/a_n\)，则
\(b_{n+1} = \frac{1 + a_n}{a_n} = \frac{1}{a_n} + 1 = b_n + 1\)。

故 \(\{b_n\}\) 是公差为 \(1\)
的等差数列：\(b_n = b_1 + (n-1) = \frac{1}{a_1} + n - 1\)。

\[a_n = \frac{1}{\frac{1}{a_1} + n - 1} = \frac{a_1}{1 + (n-1)a_1}\]

\textbf{例 2（取对数）} 设 \(a_{n+1} = a_n^2\)，\(a_1 > 0\)。

令 \(b_n = \ln a_n\)，则 \(b_{n+1} = 2b_n\)，故 \(b_n = 2^{n-1} b_1\)。

\[a_n = e^{2^{n-1} \ln a_1} = a_1^{2^{n-1}}\]

\bigskip\hrule\bigskip

\newpage

\section{卷六：线性代数}\label{ux5377ux516dux7ebfux6027ux4ee3ux6570}

\bigskip\hrule\bigskip

\begin{quote}
\textbf{前置知识}：本卷建立在卷一（数系、集合论）、卷二（代数结构）和卷三（向量）的基础之上。读者应熟悉域公理、群/环/域的基本概念、平面及空间向量的几何直观。
\end{quote}

\subsection{Ch 32
向量空间公理化}\label{ch-32-ux5411ux91cfux7a7aux95f4ux516cux7406ux5316}

\subsubsection{32.1
向量空间的定义}\label{ux5411ux91cfux7a7aux95f4ux7684ux5b9aux4e49}

\textbf{定义 32.1.1（向量空间）}

设 \(F\) 为域，\(V\) 为非空集合。\(V\) 上配备两个运算：

\begin{enumerate}
\def\labelenumi{\arabic{enumi}.}
\tightlist
\item
  \textbf{加法}：\(+ : V \times V \to V\)，\((u, v) \mapsto u + v\)
\item
  \textbf{标量乘法}：\(\cdot : F \times V \to V\)，\((\lambda, v) \mapsto \lambda v\)
\end{enumerate}

若满足以下八条公理，则称 \((V, +, \cdot)\) 为域 \(F\)
上的\textbf{向量空间}（或\textbf{线性空间}）：

\textbf{加法公理}（\(V\) 关于加法构成阿贝尔群）： - \textbf{(V1)
结合律}：\(\forall u, v, w \in V\)，\((u + v) + w = u + (v + w)\) -
\textbf{(V2) 交换律}：\(\forall u, v \in V\)，\(u + v = v + u\) -
\textbf{(V3)
零元存在}：\(\exists \mathbf{0} \in V\)，\(\forall v \in V\)，\(v + \mathbf{0} = v\)
- \textbf{(V4)
加法逆元}：\(\forall v \in V\)，\(\exists (-v) \in V\)，\(v + (-v) = \mathbf{0}\)

\textbf{标量乘法公理}： - \textbf{(V5)
标量乘法结合律}：\(\forall \lambda, \mu \in F\)，\(\forall v \in V\)，\(\lambda(\mu v) = (\lambda\mu)v\)
- \textbf{(V6) 标量单位元}：\(\forall v \in V\)，\(1 v = v\)（其中 \(1\)
是 \(F\) 的乘法单位元） - \textbf{(V7)
标量对向量加法的分配律}：\(\forall \lambda \in F\)，\(\forall u, v \in V\)，\(\lambda(u + v) = \lambda u + \lambda v\)
- \textbf{(V8)
标量对域加法的分配律}：\(\forall \lambda, \mu \in F\)，\(\forall v \in V\)，\((\lambda + \mu)v = \lambda v + \mu v\)

\(V\) 中的元素称为\textbf{向量}，\(F\) 中的元素称为\textbf{标量}。

\textbf{评注 32.1.1}：当 \(F = \mathbb{R}\)
时，称为\textbf{实向量空间}；当 \(F = \mathbb{C}\)
时，称为\textbf{复向量空间}。本卷主要讨论实向量空间，但大部分结论对任意域
\(F\) 均成立。

\textbf{定义 32.1.2（\(F^n\)）}

设 \(F\) 为域，\(n \in \mathbb{N}^+\)。定义

\[
F^n = \{(x_1, x_2, \ldots, x_n) \mid x_i \in F, \, i = 1, 2, \ldots, n\}
\]

在 \(F^n\) 上定义逐分量加法和标量乘法：

\[
(x_1, \ldots, x_n) + (y_1, \ldots, y_n) = (x_1 + y_1, \ldots, x_n + y_n)
\] \[
\lambda(x_1, \ldots, x_n) = (\lambda x_1, \ldots, \lambda x_n)
\]

则 \((F^n, +, \cdot)\) 是 \(F\) 上的向量空间。

\textbf{定理 32.1.1（零元唯一性）}：向量空间 \(V\) 中的零元
\(\mathbf{0}\) 是唯一的。

\textbf{证明}：设 \(\mathbf{0}_1, \mathbf{0}_2\) 均为 \(V\) 的零元。由
(V3)， \[
\mathbf{0}_1 = \mathbf{0}_1 + \mathbf{0}_2 = \mathbf{0}_2 + \mathbf{0}_1 = \mathbf{0}_2
\] 其中第二步用了 (V2)。\(\blacksquare\)

\textbf{定理 32.1.2（加法逆元唯一性）}：对任意 \(v \in V\)，其加法逆元
\((-v)\) 是唯一的。

\textbf{证明}：设 \(u_1, u_2\) 均为 \(v\) 的加法逆元，即
\(v + u_1 = \mathbf{0}\)，\(v + u_2 = \mathbf{0}\)。则 \[
u_1 = u_1 + \mathbf{0} = u_1 + (v + u_2) = (u_1 + v) + u_2 = \mathbf{0} + u_2 = u_2
\] 其中第三步用了 (V1)。\(\blacksquare\)

\textbf{定理 32.1.3（基本运算性质）}：设 \(V\) 为 \(F\) 上的向量空间，则
\(\forall v \in V\)，\(\forall \lambda \in F\)：

\begin{enumerate}
\def\labelenumi{(\roman{enumi})}
\tightlist
\item
  \(0v = \mathbf{0}\)（标量 \(0\) 乘任何向量得零向量）
\item
  \(\lambda \mathbf{0} = \mathbf{0}\)（任何标量乘零向量得零向量）
\item
  \((-1)v = -v\)（\(-1\) 乘向量等于加法逆元）
\item
  若 \(\lambda v = \mathbf{0}\)，则 \(\lambda = 0\) 或
  \(v = \mathbf{0}\)
\end{enumerate}

\textbf{证明}：

\begin{enumerate}
\def\labelenumi{(\roman{enumi})}
\item
  由 (V8)：\(0v = (0 + 0)v = 0v + 0v\)。两边加 \((-0v)\) 得
  \(0v = \mathbf{0}\)。
\item
  由
  (V7)：\(\lambda \mathbf{0} = \lambda(\mathbf{0} + \mathbf{0}) = \lambda \mathbf{0} + \lambda \mathbf{0}\)。两边加
  \((-\lambda \mathbf{0})\) 得 \(\lambda \mathbf{0} = \mathbf{0}\)。
\item
  \(v + (-1)v = 1v + (-1)v = (1 + (-1))v = 0v = \mathbf{0}\)。由逆元唯一性，\((-1)v = -v\)。
\item
  设 \(\lambda v = \mathbf{0}\)。若 \(\lambda \neq 0\)，则
  \(\lambda^{-1}\) 在 \(F\) 中存在，且 \[
  v = 1v = (\lambda^{-1}\lambda)v = \lambda^{-1}(\lambda v) = \lambda^{-1}\mathbf{0} = \mathbf{0}
  \] 其中最后一步用了 (ii)。\(\blacksquare\)
\end{enumerate}

\bigskip\hrule\bigskip

\subsubsection{32.2 子空间}\label{ux5b50ux7a7aux95f4}

\textbf{定义 32.2.1（子空间）}

设 \(V\) 为 \(F\) 上的向量空间，\(W \subseteq V\) 为非空子集。若 \(W\)
在 \(V\) 的加法和标量乘法下封闭，即

\begin{enumerate}
\def\labelenumi{\arabic{enumi}.}
\tightlist
\item
  \(\forall u, v \in W\)，\(u + v \in W\)
\item
  \(\forall \lambda \in F\)，\(\forall v \in W\)，\(\lambda v \in W\)
\end{enumerate}

则称 \(W\) 为 \(V\) 的\textbf{子空间}，记为 \(W \leq V\)。

\textbf{定理 32.2.1（子空间判定）}：\(W \subseteq V\) 是子空间当且仅当

\begin{enumerate}
\def\labelenumi{(\roman{enumi})}
\tightlist
\item
  \(W \neq \varnothing\)（或等价地，\(\mathbf{0} \in W\)）
\item
  \(\forall u, v \in W\)，\(\forall \lambda, \mu \in F\)，\(\lambda u + \mu v \in W\)
\end{enumerate}

\textbf{证明}：

（\(\Rightarrow\)）设 \(W\) 是子空间。由定义 \(W \neq \varnothing\)。取
\(v \in W\)，则 \(0v = \mathbf{0} \in W\)（定理 32.1.3(i)）。对任意
\(u, v \in W\)，\(\lambda u, \mu v \in W\)（标量封闭），故
\(\lambda u + \mu v \in W\)（加法封闭）。

（\(\Leftarrow\)）设条件成立。取 \(\lambda = \mu = 1\) 得加法封闭；取
\(\mu = 0\) 得标量封闭。由 \(W \neq \varnothing\) 并取
\(\lambda = \mu = 0\) 得 \(\mathbf{0} \in W\)，故 \(W\)
非空。\(\blacksquare\)

\textbf{定理 32.2.2（子空间的交）}：设 \(W_1, W_2\) 均为 \(V\)
的子空间，则 \(W_1 \cap W_2\) 也是 \(V\) 的子空间。

\textbf{证明}：\(\mathbf{0} \in W_1 \cap W_2\)（因为
\(\mathbf{0} \in W_1\) 且 \(\mathbf{0} \in W_2\)），故
\(W_1 \cap W_2 \neq \varnothing\)。

设 \(u, v \in W_1 \cap W_2\)，\(\lambda, \mu \in F\)。则
\(u, v \in W_1\) 且 \(u, v \in W_2\)。由于 \(W_1, W_2\)
均为子空间，\(\lambda u + \mu v \in W_1\) 且
\(\lambda u + \mu v \in W_2\)，故
\(\lambda u + \mu v \in W_1 \cap W_2\)。由定理 32.2.1，\(W_1 \cap W_2\)
是子空间。\(\blacksquare\)

\textbf{注记 32.2.1}：子空间的并一般不是子空间。例如 \(\mathbb{R}^2\) 中
\(x\) 轴和 \(y\) 轴都是子空间，但它们的并不是（因为
\((1,0) + (0,1) = (1,1)\) 不在并中）。

\textbf{定义 32.2.2（和与直和）}

设 \(W_1, W_2\) 均为 \(V\) 的子空间。

\begin{enumerate}
\def\labelenumi{\arabic{enumi}.}
\tightlist
\item
  \(W_1\) 与 \(W_2\) 的\textbf{和}定义为
  \(W_1 + W_2 = \{w_1 + w_2 \mid w_1 \in W_1, \, w_2 \in W_2\}\)
\item
  若 \(W_1 \cap W_2 = \{\mathbf{0}\}\)，则称和为\textbf{直和}，记为
  \(W_1 \oplus W_2\)
\end{enumerate}

\textbf{定理 32.2.3}：\(W_1 + W_2\) 是 \(V\) 的子空间。

\textbf{证明}：\(\mathbf{0} = \mathbf{0} + \mathbf{0} \in W_1 + W_2\)。设
\(u = u_1 + u_2\)，\(v = v_1 + v_2\)（\(u_i, v_i \in W_i\)），\(\lambda, \mu \in F\)，则
\[
\lambda u + \mu v = (\lambda u_1 + \mu v_1) + (\lambda u_2 + \mu v_2) \in W_1 + W_2
\] 由定理 32.2.1，\(W_1 + W_2\) 是子空间。\(\blacksquare\)

\textbf{定理 32.2.4（直和与分解的唯一性）}：\(W_1 + W_2\)
是直和当且仅当每个 \(v \in W_1 + W_2\) 的分解
\(v = w_1 + w_2\)（\(w_i \in W_i\)）是唯一的。

\textbf{证明}：

（\(\Rightarrow\)）设 \(W_1 \cap W_2 = \{\mathbf{0}\}\)。若
\(v = w_1 + w_2 = w_1' + w_2'\)，则
\(w_1 - w_1' = w_2' - w_2\)。左边属于 \(W_1\)，右边属于
\(W_2\)，故该向量属于 \(W_1 \cap W_2 = \{\mathbf{0}\}\)。因此
\(w_1 = w_1'\)，\(w_2 = w_2'\)。

（\(\Leftarrow\)）若分解唯一，设 \(v \in W_1 \cap W_2\)。则
\(v = v + \mathbf{0} = \mathbf{0} + v\) 是 \(v\) 的两个分解，由唯一性
\(v = \mathbf{0}\)。故
\(W_1 \cap W_2 = \{\mathbf{0}\}\)。\(\blacksquare\)

\bigskip\hrule\bigskip

\subsubsection{32.3
线性组合与生成}\label{ux7ebfux6027ux7ec4ux5408ux4e0eux751fux6210}

\textbf{定义 32.3.1（线性组合）}

设
\(v_1, v_2, \ldots, v_n \in V\)，\(\lambda_1, \lambda_2, \ldots, \lambda_n \in F\)。称

\[
\sum_{i=1}^{n} \lambda_i v_i = \lambda_1 v_1 + \lambda_2 v_2 + \cdots + \lambda_n v_n
\]

为 \(v_1, \ldots, v_n\) 的一个\textbf{线性组合}。\(\lambda_i\)
称为该组合的\textbf{系数}。

\textbf{定义 32.3.2（生成空间）}

设 \(S = \{v_1, v_2, \ldots, v_n\} \subseteq V\)。\(S\)
的\textbf{生成空间}（或\textbf{张成空间}）定义为 \(S\)
的所有线性组合的集合：

\[
\operatorname{span}(S) = \left\{\sum_{i=1}^{n} \lambda_i v_i \;\middle|\; \lambda_i \in F, \, i = 1, \ldots, n\right\}
\]

\textbf{定理 32.3.1}：\(\operatorname{span}(S)\) 是 \(V\)
的子空间，且是包含 \(S\) 的最小子空间（即任何包含 \(S\) 的子空间都包含
\(\operatorname{span}(S)\)）。

\textbf{证明}：先证是子空间。\(\mathbf{0} = \sum_{i=1}^{n} 0v_i \in \operatorname{span}(S)\)。

设
\(u = \sum \lambda_i v_i\)，\(w = \sum \mu_i v_i \in \operatorname{span}(S)\)，\(\alpha, \beta \in F\)。则
\[
\alpha u + \beta w = \sum (\alpha\lambda_i + \beta\mu_i)v_i \in \operatorname{span}(S)
\] 由定理 32.2.1，\(\operatorname{span}(S)\) 是子空间。

再证最小性。设 \(W\) 是包含 \(S\) 的任意子空间。对任意
\(v \in \operatorname{span}(S)\)，\(v = \sum \lambda_i v_i\)。由于
\(v_i \in S \subseteq W\) 且 \(W\)
对加法和标量乘法封闭，\(v \in W\)。因此
\(\operatorname{span}(S) \subseteq W\)。\(\blacksquare\)

\textbf{定义 32.3.3（有限生成）}

若存在有限集 \(S \subseteq V\) 使得 \(\operatorname{span}(S) = V\)，则称
\(V\) 是\textbf{有限生成}的向量空间。

\textbf{评注
32.3.1}：本卷主要讨论有限生成向量空间（也称为有限维向量空间）。无限维向量空间将在泛函分析（卷十四）中详细讨论。

\bigskip\hrule\bigskip

\subsubsection{32.4
线性相关与线性无关}\label{ux7ebfux6027ux76f8ux5173ux4e0eux7ebfux6027ux65e0ux5173}

\textbf{定义 32.4.1（线性相关与线性无关）}

设 \(v_1, v_2, \ldots, v_n \in V\)。

\begin{enumerate}
\def\labelenumi{\arabic{enumi}.}
\tightlist
\item
  若存在不全为零的标量 \(\lambda_1, \ldots, \lambda_n \in F\) 使得
  \(\sum_{i=1}^{n} \lambda_i v_i = \mathbf{0}\)，则称向量组
  \(\{v_1, \ldots, v_n\}\) \textbf{线性相关}。
\item
  若仅当 \(\lambda_1 = \cdots = \lambda_n = 0\) 时才有
  \(\sum_{i=1}^{n} \lambda_i v_i = \mathbf{0}\)，则称向量组
  \(\{v_1, \ldots, v_n\}\) \textbf{线性无关}。
\end{enumerate}

\textbf{定理 32.4.1（线性相关的基本性质）}：

\begin{enumerate}
\def\labelenumi{(\roman{enumi})}
\tightlist
\item
  含零向量的向量组必定线性相关。
\item
  若向量组的一个子集线性相关，则整个向量组线性相关。
\item
  \(\{v_1, \ldots, v_n\}\) 线性相关当且仅当存在某个 \(v_k\)
  可表示为其余向量的线性组合。
\end{enumerate}

\textbf{证明}：

\begin{enumerate}
\def\labelenumi{(\roman{enumi})}
\item
  设 \(\{v_1, \ldots, v_n\}\) 含 \(\mathbf{0}\)。取
  \(\lambda_1 = 1\)（对应 \(\mathbf{0}\)），其余 \(\lambda_i = 0\)，则
  \(\sum \lambda_i v_i = \mathbf{0}\) 且系数不全为零，故线性相关。
\item
  设子集 \(\{v_{i_1}, \ldots, v_{i_k}\}\)
  线性相关，存在不全为零的系数使其线性组合为零。将其他向量的系数取为
  \(0\)，得到整个向量组的线性组合为零且系数不全为零。
\item
  （\(\Rightarrow\)）设 \(\sum \lambda_i v_i = \mathbf{0}\) 且
  \(\lambda_k \neq 0\)。则
  \(v_k = -\lambda_k^{-1}\sum_{i \neq k} \lambda_i v_i\)。
\end{enumerate}

（\(\Leftarrow\)）设 \(v_k = \sum_{i \neq k} \mu_i v_i\)。则
\(\sum_{i \neq k} \mu_i v_i + (-1)v_k = \mathbf{0}\)，系数不全为零（\(v_k\)
的系数为 \(-1 \neq 0\)）。\(\blacksquare\)

\textbf{定理 32.4.2（线性无关组的扩充性质）}：设
\(\{v_1, \ldots, v_n\}\) 线性无关，\(v \in V\)。则
\(\{v_1, \ldots, v_n, v\}\) 线性相关当且仅当
\(v \in \operatorname{span}\{v_1, \ldots, v_n\}\)。

\textbf{证明}：

（\(\Rightarrow\)）若 \(\{v_1, \ldots, v_n, v\}\)
线性相关，则存在不全为零的 \(\lambda_1, \ldots, \lambda_n, \lambda\)
使得 \(\sum \lambda_i v_i + \lambda v = \mathbf{0}\)。若
\(\lambda = 0\)，则 \(\sum \lambda_i v_i = \mathbf{0}\) 且 \(\lambda_i\)
不全为零，与 \(\{v_1, \ldots, v_n\}\) 线性无关矛盾。故
\(\lambda \neq 0\)，从而
\(v = -\lambda^{-1}\sum \lambda_i v_i \in \operatorname{span}\{v_1, \ldots, v_n\}\)。

（\(\Leftarrow\)）若
\(v \in \operatorname{span}\{v_1, \ldots, v_n\}\)，则
\(v = \sum \mu_i v_i\)。于是
\(\sum \mu_i v_i + (-1)v = \mathbf{0}\)，系数不全为零，故线性相关。\(\blacksquare\)

\bigskip\hrule\bigskip

\subsubsection{32.5 基与维数}\label{ux57faux4e0eux7ef4ux6570}

\textbf{定义 32.5.1（基）}

设 \(B = \{e_1, e_2, \ldots, e_n\} \subseteq V\)。若满足：

\begin{enumerate}
\def\labelenumi{\arabic{enumi}.}
\tightlist
\item
  \textbf{线性无关性}：\(B\) 是线性无关的
\item
  \textbf{生成性}：\(\operatorname{span}(B) = V\)
\end{enumerate}

则称 \(B\) 为 \(V\) 的一组\textbf{基}（或\textbf{基底}）。

\textbf{定理 32.5.1（基的等价刻画）}：\(B = \{e_1, \ldots, e_n\}\) 是
\(V\) 的基当且仅当每个 \(v \in V\) 可唯一表示为 \(B\) 中向量的线性组合。

\textbf{证明}：

（\(\Rightarrow\)）由生成性，\(v\) 可表示为
\(v = \sum \lambda_i e_i\)。若还存在 \(v = \sum \mu_i e_i\)，则
\(\sum (\lambda_i - \mu_i)e_i = \mathbf{0}\)。由线性无关性，\(\lambda_i - \mu_i = 0\)，即
\(\lambda_i = \mu_i\)，故表示唯一。

（\(\Leftarrow\)）若每个 \(v\) 可唯一表示，则生成性自动成立。取
\(v = \mathbf{0}\)，由唯一性，仅有 \(\lambda_i = 0\) 可使
\(\sum \lambda_i e_i = \mathbf{0}\)，故 \(B\) 线性无关。\(\blacksquare\)

\textbf{定义 32.5.2（坐标）}：设 \(B = \{e_1, \ldots, e_n\}\) 为 \(V\)
的基，\(v = \sum_{i=1}^{n} \lambda_i e_i\)。则
\((\lambda_1, \ldots, \lambda_n) \in F^n\) 称为 \(v\) 关于基 \(B\)
的\textbf{坐标}。

\textbf{定理 32.5.2（基的存在性）}：任何非零有限生成向量空间都存在基。

\textbf{证明}：设
\(V = \operatorname{span}\{v_1, \ldots, v_m\}\)。构造算法如下：

从 \(i = 1\) 到 \(m\)，若 \(v_i\) 不属于前 \(i-1\)
个向量的生成空间，则保留 \(v_i\)；否则丢弃。设保留的向量为
\(\{e_1, \ldots, e_n\}\)。

由构造过程，每个 \(v_i\) 都属于
\(\operatorname{span}\{e_1, \ldots, e_n\}\)，故
\(V \subseteq \operatorname{span}\{e_1, \ldots, e_n\}\)，生成性成立。

若 \(\sum \lambda_j e_j = \mathbf{0}\) 且 \(\lambda_k\)
是最后一个非零系数，则
\(e_k = -\lambda_k^{-1}\sum_{j < k} \lambda_j e_j\)，说明 \(e_k\) 属于前
\(k-1\) 个保留向量的生成空间，与构造中保留 \(e_k\) 的条件矛盾。故所有
\(\lambda_j = 0\)，线性无关性成立。\(\blacksquare\)

\textbf{注记
32.5.1}：上述构造证明给出了从有限生成集中提取基的算法。对于有限维向量空间，基的存在性由上述构造直接给出。对于无限维向量空间，基的存在性需要选择公理（Axiom
of Choice，等价于 Zorn 引理）：考虑 \(V\)
的所有线性无关子集构成的偏序集（按包含关系），Zorn
引理保证存在极大线性无关子集，该极大集即为 \(V\) 的一组基（参见
V20（概率论2））。

\textbf{引理 32.5.3（替换引理，Steinitz 引理）}：设
\(\{e_1, \ldots, e_n\}\) 为 \(V\)
的基，\(\{v_1, \ldots, v_m\} \subseteq V\) 线性无关。则
\(m \leq n\)，且可以适当地用 \(v_1, \ldots, v_m\) 替换基中的 \(m\)
个向量，使得替换后的集合仍为基。

\textbf{证明}：对 \(m\) 归纳。

\textbf{基础} \(m = 0\)：平凡。

\textbf{归纳}：假设结论对 \(m-1\)
成立。由归纳假设，\(m-1 \leq n\)，且可替换 \(m-1\) 个基向量，得到新基
\(\{v_1, \ldots, v_{m-1}, f_m, \ldots, f_n\}\)（其中
\(\{f_m, \ldots, f_n\}\) 是未被替换的原基向量）。

\(v_m\)
可表示为新基的线性组合：\(v_m = \sum_{i=1}^{m-1} \alpha_i v_i + \sum_{j=m}^{n} \beta_j f_j\)。由于
\(\{v_1, \ldots, v_m\}\) 线性无关，\(v_m\) 不能仅由
\(v_1, \ldots, v_{m-1}\) 表示，故存在 \(k \geq m\) 使得
\(\beta_k \neq 0\)。

不妨设 \(k = m\)（重排 \(f_m, \ldots, f_n\)）。则
\(f_m = \beta_m^{-1}(v_m - \sum_{i=1}^{m-1} \alpha_i v_i - \sum_{j=m+1}^{n} \beta_j f_j)\)，故
\(f_m \in \operatorname{span}\{v_1, \ldots, v_m, f_{m+1}, \ldots, f_n\}\)。因此这
\(n\) 个向量生成 \(V\)，且数量为 \(n\)。由定理 32.5.2
的构造可知它们线性无关（否则可提取更小的基，与 \(n\)
已经是基的大小矛盾）。故 \(\{v_1, \ldots, v_m, f_{m+1}, \ldots, f_n\}\)
是基。

特别地，\(m \leq n\)（否则 \(m > n\) 时由归纳假设
\(m-1 \geq n\)，矛盾）。\(\blacksquare\)

\textbf{定理 32.5.4（维数的不变性）}：有限生成向量空间
\(V \neq \{\mathbf{0}\}\) 的任何两组基具有相同个数的向量。

\textbf{证明}：设 \(B_1 = \{e_1, \ldots, e_n\}\) 和
\(B_2 = \{f_1, \ldots, f_m\}\) 均为 \(V\) 的基。\(B_1\) 是基，\(B_2\)
线性无关，由替换引理 \(m \leq n\)。交换角色得 \(n \leq m\)。故
\(n = m\)。\(\blacksquare\)

\textbf{定义 32.5.3（维数）}：向量空间 \(V\) 的基所含向量的个数称为
\(V\) 的\textbf{维数}，记为 \(\dim V\)。约定
\(\dim\{\mathbf{0}\} = 0\)。

\textbf{定理 32.5.5（维数的基本性质）}：设 \(V\) 为 \(n\) 维向量空间。

\begin{enumerate}
\def\labelenumi{(\roman{enumi})}
\tightlist
\item
  \(V\) 中任何 \(n+1\) 个向量必线性相关。
\item
  \(V\) 中任何 \(n\) 个线性无关的向量构成基。
\item
  \(V\) 中任何生成 \(V\) 的 \(n\) 个向量构成基。
\item
  设 \(W \leq V\)，则 \(\dim W \leq \dim V\)。等号成立当且仅当
  \(W = V\)。
\end{enumerate}

\textbf{证明}：

\begin{enumerate}
\def\labelenumi{(\roman{enumi})}
\item
  由替换引理，若存在 \(n+1\) 个线性无关向量，则 \(n+1 \leq n\)，矛盾。
\item
  设 \(S = \{v_1, \ldots, v_n\}\) 线性无关。若
  \(\operatorname{span}(S) \neq V\)，则存在
  \(v \in V \setminus \operatorname{span}(S)\)。由定理
  32.4.2，\(\{v_1, \ldots, v_n, v\}\) 线性无关，与 (i) 矛盾。故
  \(\operatorname{span}(S) = V\)。
\item
  设 \(S\) 生成 \(V\)。若 \(S\)
  线性相关，则可从中提取一个更小的线性无关集仍生成 \(V\)，即基的大小小于
  \(n\)，与 \(\dim V = n\) 矛盾。故 \(S\) 线性无关。
\item
  设 \(B_W\) 为 \(W\) 的基。\(B_W\) 是 \(V\)
  中线性无关的向量组，可扩充为 \(V\) 的基。故 \(\dim W \leq \dim V\)。若
  \(\dim W = \dim V = n\)，则 \(B_W\) 本身就是 \(V\) 的基（\(n\)
  个线性无关向量），故 \(W = V\)。\(\blacksquare\)
\end{enumerate}

\textbf{定理 32.5.6（维数公式）}：设 \(W_1, W_2 \leq V\)，则

\[
\dim(W_1 + W_2) = \dim W_1 + \dim W_2 - \dim(W_1 \cap W_2)
\]

\textbf{证明}：设
\(\dim(W_1 \cap W_2) = r\)，\(\dim W_1 = r + s\)，\(\dim W_2 = r + t\)。

取 \(W_1 \cap W_2\) 的基 \(\{e_1, \ldots, e_r\}\)。将其扩充为 \(W_1\)
的基 \(\{e_1, \ldots, e_r, u_1, \ldots, u_s\}\)，也扩充为 \(W_2\) 的基
\(\{e_1, \ldots, e_r, w_1, \ldots, w_t\}\)。

断言 \(\{e_1, \ldots, e_r, u_1, \ldots, u_s, w_1, \ldots, w_t\}\) 是
\(W_1 + W_2\) 的基。

生成性：\(W_1 + W_2\) 中任意向量可写为 \(W_1\) 中向量与 \(W_2\)
中向量的和，而 \(W_1\) 中向量由 \(e_i, u_j\) 生成，\(W_2\) 中向量由
\(e_i, w_k\) 生成，故该集合生成 \(W_1 + W_2\)。

线性无关性：设
\(\sum \alpha_i e_i + \sum \beta_j u_j + \sum \gamma_k w_k = \mathbf{0}\)。令
\(v = \sum \alpha_i e_i + \sum \beta_j u_j = -\sum \gamma_k w_k\)。左边属于
\(W_1\)，右边属于 \(W_2\)，故 \(v \in W_1 \cap W_2\)。因此 \(v\) 可由
\(e_1, \ldots, e_r\) 线性表示：\(v = \sum \delta_i e_i\)。

于是 \(\sum \delta_i e_i + \sum \gamma_k w_k = \mathbf{0}\)。由于
\(\{e_1, \ldots, e_r, w_1, \ldots, w_t\}\) 是 \(W_2\)
的基，所有系数为零：\(\delta_i = 0\)，\(\gamma_k = 0\)。从而
\(v = \mathbf{0}\)，即
\(\sum \alpha_i e_i + \sum \beta_j u_j = \mathbf{0}\)。由 \(W_1\)
的基的线性无关性，\(\alpha_i = 0\)，\(\beta_j = 0\)。

因此
\(\dim(W_1 + W_2) = r + s + t = (r + s) + (r + t) - r = \dim W_1 + \dim W_2 - \dim(W_1 \cap W_2)\)。\(\blacksquare\)

\textbf{推论 32.5.7（直和的维数）}：若
\(W_1 \cap W_2 = \{\mathbf{0}\}\)，则
\(\dim(W_1 \oplus W_2) = \dim W_1 + \dim W_2\)。

\textbf{证明}：由维数公式，\(\dim(W_1 \cap W_2) = 0\) 时
\(\dim(W_1 + W_2) = \dim W_1 + \dim W_2\)。\(\blacksquare\)

\bigskip\hrule\bigskip

\subsubsection{32.6
坐标映射与同构}\label{ux5750ux6807ux6620ux5c04ux4e0eux540cux6784}

\textbf{定义 32.6.1（坐标映射）}：设 \(B = \{e_1, \ldots, e_n\}\) 为
\(V\) 的基。定义映射 \(\varphi_B : V \to F^n\)，将
\(v = \sum \lambda_i e_i\) 映为其坐标
\((\lambda_1, \ldots, \lambda_n)\)。

\textbf{定理 32.6.1}：\(\varphi_B\) 是双射，且满足
\(\varphi_B(u + v) = \varphi_B(u) + \varphi_B(v)\)，\(\varphi_B(\lambda v) = \lambda \varphi_B(v)\)。

\textbf{证明}：由基的等价刻画（定理 32.5.1），表示唯一，故 \(\varphi_B\)
是良定义的双射。线性性质由坐标的加法与标量乘法的定义直接验证。\(\blacksquare\)

\textbf{定义 32.6.2（同构）}：设 \(V, W\) 为 \(F\)
上的向量空间。若存在双射 \(T : V \to W\) 满足
\(\forall u, v \in V, \forall \lambda \in F\)： \[
T(u + v) = T(u) + T(v), \quad T(\lambda u) = \lambda T(u)
\] 则称 \(T\) 为\textbf{同构映射}，\(V\) 与 \(W\) \textbf{同构}，记为
\(V \cong W\)。

\textbf{定理 32.6.2}：\(F\) 上的任何 \(n\) 维向量空间与 \(F^n\) 同构。

\textbf{证明}：坐标映射 \(\varphi_B : V \to F^n\)
即为同构映射。\(\blacksquare\)

\textbf{推论 32.6.3}：\(F\)
上的两个有限维向量空间同构当且仅当它们维数相同。

\textbf{证明}：若维数相同，都与 \(F^n\)
同构，故互相同构。若同构，则基的像也是基，故维数相同。\(\blacksquare\)

\bigskip\hrule\bigskip

\subsection{Ch 33
线性映射与矩阵}\label{ch-33-ux7ebfux6027ux6620ux5c04ux4e0eux77e9ux9635}

\subsubsection{33.1
线性映射的定义与基本性质}\label{ux7ebfux6027ux6620ux5c04ux7684ux5b9aux4e49ux4e0eux57faux672cux6027ux8d28}

\textbf{定义 33.1.1（线性映射）}

设 \(V, W\) 为域 \(F\) 上的向量空间。映射 \(T : V \to W\) 若满足：

\begin{enumerate}
\def\labelenumi{\arabic{enumi}.}
\tightlist
\item
  \textbf{加法保持}：\(\forall u, v \in V\)，\(T(u + v) = T(u) + T(v)\)
\item
  \textbf{标量乘法保持}：\(\forall \lambda \in F, \forall v \in V\)，\(T(\lambda v) = \lambda T(v)\)
\end{enumerate}

则称 \(T\) 为\textbf{线性映射}（或\textbf{线性变换}）。当 \(V = W\)
时，\(T\) 特称为 \(V\) 上的\textbf{线性算子}。

\textbf{定理 33.1.1（线性映射的基本性质）}：设 \(T : V \to W\)
为线性映射，则

\begin{enumerate}
\def\labelenumi{(\roman{enumi})}
\tightlist
\item
  \(T(\mathbf{0}_V) = \mathbf{0}_W\)
\item
  \(T(-v) = -T(v)\)
\item
  \(T\left(\sum_{i=1}^{n} \lambda_i v_i\right) = \sum_{i=1}^{n} \lambda_i T(v_i)\)
\end{enumerate}

\textbf{证明}：

\begin{enumerate}
\def\labelenumi{(\roman{enumi})}
\item
  \(T(\mathbf{0}_V) = T(0 \cdot \mathbf{0}_V) = 0 \cdot T(\mathbf{0}_V) = \mathbf{0}_W\)。
\item
  \(T(-v) = T((-1)v) = (-1)T(v) = -T(v)\)。
\item
  对 \(n\) 归纳，利用线性映射的两条公理。\(\blacksquare\)
\end{enumerate}

\textbf{定义 33.1.2（线性映射的运算）}：设 \(S, T : V \to W\)
为线性映射，\(\lambda \in F\)。

\begin{enumerate}
\def\labelenumi{\arabic{enumi}.}
\tightlist
\item
  \textbf{加法}：\((S + T)(v) = S(v) + T(v)\)
\item
  \textbf{标量乘法}：\((\lambda T)(v) = \lambda T(v)\)
\item
  \textbf{复合}（设
  \(T : V \to W\)，\(S : W \to U\)）：\((S \circ T)(v) = S(T(v))\)
\end{enumerate}

\textbf{定理 33.1.2}：所有从 \(V\) 到 \(W\) 的线性映射的集合，记为
\(\mathcal{L}(V, W)\)，在上述加法和标量乘法下构成 \(F\) 上的向量空间。

\textbf{证明}：逐条验证向量空间公理。零元为零映射
\(0(v) = \mathbf{0}_W\)。加法逆元为
\((-T)(v) = -T(v)\)。\(\blacksquare\)

\textbf{定义 33.1.3（同构与自同构）}：

\begin{itemize}
\tightlist
\item
  可逆的线性映射称为\textbf{同构映射}（其逆也自动是线性的）
\item
  \(V\) 到自身的同构称为\textbf{自同构}，所有自同构构成群
  \(\operatorname{GL}(V)\)
\end{itemize}

\textbf{定理 33.1.3（逆映射的线性性）}：若线性映射 \(T : V \to W\)
可逆，则 \(T^{-1} : W \to V\) 也是线性映射。

\textbf{证明}：设 \(w_1, w_2 \in W\)，则存在 \(v_1, v_2 \in V\) 使得
\(T(v_i) = w_i\)。于是 \[
T^{-1}(w_1 + w_2) = T^{-1}(T(v_1) + T(v_2)) = T^{-1}(T(v_1 + v_2)) = v_1 + v_2 = T^{-1}(w_1) + T^{-1}(w_2)
\] 同理可证标量乘法的保持。\(\blacksquare\)

\bigskip\hrule\bigskip

\subsubsection{33.2 核与像}\label{ux6838ux4e0eux50cf}

\textbf{定义 33.2.1（核与像）}：设 \(T : V \to W\) 为线性映射。

\begin{enumerate}
\def\labelenumi{\arabic{enumi}.}
\tightlist
\item
  \(T\) 的\textbf{核}：\(\ker T = \{v \in V \mid T(v) = \mathbf{0}_W\}\)
\item
  \(T\) 的\textbf{像}：\(\operatorname{im} T = \{T(v) \mid v \in V\}\)
\end{enumerate}

\textbf{定理 33.2.1}：\(\ker T\) 是 \(V\)
的子空间，\(\operatorname{im} T\) 是 \(W\) 的子空间。

\textbf{证明}：\(\mathbf{0}_V \in \ker T\)（因为
\(T(\mathbf{0}_V) = \mathbf{0}_W\)）。设
\(u, v \in \ker T\)，\(\lambda, \mu \in F\)，则
\(T(\lambda u + \mu v) = \lambda T(u) + \mu T(v) = \mathbf{0}_W + \mathbf{0}_W = \mathbf{0}_W\)，故
\(\lambda u + \mu v \in \ker T\)。由定理 32.2.1，\(\ker T\) 是子空间。

类似，\(\mathbf{0}_W = T(\mathbf{0}_V) \in \operatorname{im} T\)。设
\(T(u), T(v) \in \operatorname{im} T\)，则
\(\lambda T(u) + \mu T(v) = T(\lambda u + \mu v) \in \operatorname{im} T\)。\(\blacksquare\)

\textbf{定义 33.2.2（秩与零度）}：设 \(T : V \to W\) 为线性映射。

\begin{enumerate}
\def\labelenumi{\arabic{enumi}.}
\tightlist
\item
  \(\operatorname{rank} T = \dim(\operatorname{im} T)\)，称为 \(T\)
  的\textbf{秩}
\item
  \(\operatorname{nullity} T = \dim(\ker T)\)，称为 \(T\)
  的\textbf{零度}
\end{enumerate}

\textbf{定理 33.2.2（单射的判定）}：线性映射 \(T : V \to W\)
是单射当且仅当 \(\ker T = \{\mathbf{0}\}\)。

\textbf{证明}：

（\(\Rightarrow\)）若 \(T\)
是单射，\(T(v) = \mathbf{0} = T(\mathbf{0})\)，则 \(v = \mathbf{0}\)，故
\(\ker T = \{\mathbf{0}\}\)。

（\(\Leftarrow\)）若 \(\ker T = \{\mathbf{0}\}\) 且 \(T(u) = T(v)\)，则
\(T(u - v) = \mathbf{0}\)，故
\(u - v \in \ker T = \{\mathbf{0}\}\)，于是 \(u = v\)。\(\blacksquare\)

\bigskip\hrule\bigskip

\subsubsection{33.3
秩-零化度定理}\label{ux79e9-ux96f6ux5316ux5ea6ux5b9aux7406}

\textbf{定理 33.3.1（秩-零化度定理）}：设 \(V\)
为有限维向量空间，\(T : V \to W\) 为线性映射。则

\[
\dim V = \dim(\ker T) + \dim(\operatorname{im} T) = \operatorname{nullity} T + \operatorname{rank} T
\]

\textbf{证明}：设 \(\dim V = n\)，\(\dim(\ker T) = k\)。取 \(\ker T\)
的基 \(\{v_1, \ldots, v_k\}\)，将其扩充为 \(V\) 的基
\(\{v_1, \ldots, v_k, u_1, \ldots, u_{n-k}\}\)。

断言 \(\{T(u_1), \ldots, T(u_{n-k})\}\) 是 \(\operatorname{im} T\)
的基。

生成性：对任意
\(v \in V\)，\(v = \sum \alpha_i v_i + \sum \beta_j u_j\)，则
\(T(v) = \sum \beta_j T(u_j)\)（因为 \(T(v_i) = \mathbf{0}\)）。故
\(\operatorname{im} T = \operatorname{span}\{T(u_1), \ldots, T(u_{n-k})\}\)。

线性无关性：设 \(\sum \gamma_j T(u_j) = \mathbf{0}\)。则
\(T(\sum \gamma_j u_j) = \mathbf{0}\)，故
\(\sum \gamma_j u_j \in \ker T\)。因此
\(\sum \gamma_j u_j = \sum \delta_i v_i\)，即
\(\sum (-\delta_i) v_i + \sum \gamma_j u_j = \mathbf{0}\)。由基的线性无关性，所有系数为零，特别地
\(\gamma_j = 0\)（\(j = 1, \ldots, n-k\)）。

因此 \(\dim(\operatorname{im} T) = n - k\)，即
\(\dim V = \dim(\ker T) + \dim(\operatorname{im} T)\)。\(\blacksquare\)

\textbf{推论 33.3.2（满射与单射的条件）}：设 \(T : V \to W\)
为线性映射，\(\dim V = \dim W = n\)。

\begin{enumerate}
\def\labelenumi{(\roman{enumi})}
\tightlist
\item
  \(T\) 是单射 \(\iff\) \(T\) 是满射 \(\iff\) \(T\) 是同构。
\item
  \(\operatorname{rank} T = n \iff \ker T = \{\mathbf{0}\}\)。
\end{enumerate}

\textbf{证明}：

\begin{enumerate}
\def\labelenumi{(\roman{enumi})}
\item
  由秩-零化度定理，\(\dim V = \dim(\ker T) + \dim(\operatorname{im} T)\)。\(T\)
  单射
  \(\iff \ker T = \{\mathbf{0}\} \iff \dim(\ker T) = 0 \iff \dim(\operatorname{im} T) = n \iff \operatorname{im} T = W \iff T\)
  满射。此时 \(T\) 既单且满，即为同构。
\item
  \(\operatorname{rank} T = n \iff \dim(\operatorname{im} T) = n \iff \dim(\ker T) = 0 \iff \ker T = \{\mathbf{0}\}\)。\(\blacksquare\)
\end{enumerate}

\bigskip\hrule\bigskip

\subsubsection{33.4
线性映射的矩阵表示}\label{ux7ebfux6027ux6620ux5c04ux7684ux77e9ux9635ux8868ux793a}

\textbf{定义 33.4.1（矩阵）}：设 \(F\)
为域，\(m, n \in \mathbb{N}^+\)。一个 \(m \times n\) \textbf{矩阵} \(A\)
是一个由 \(mn\) 个 \(F\) 中元素组成的矩形阵列：

\[
A = \begin{pmatrix}
a_{11} & a_{12} & \cdots & a_{1n} \\
a_{21} & a_{22} & \cdots & a_{2n} \\
\vdots & \vdots & \ddots & \vdots \\
a_{m1} & a_{m2} & \cdots & a_{mn}
\end{pmatrix}
\]

其中 \(a_{ij} \in F\) 表示第 \(i\) 行第 \(j\) 列的元素。所有
\(m \times n\) 矩阵的集合记为 \(M_{m \times n}(F)\)。

\textbf{定义 33.4.2（矩阵的加法和标量乘法）}：在 \(M_{m \times n}(F)\)
上定义：

\begin{enumerate}
\def\labelenumi{\arabic{enumi}.}
\tightlist
\item
  \((A + B)_{ij} = a_{ij} + b_{ij}\)
\item
  \((\lambda A)_{ij} = \lambda a_{ij}\)
\end{enumerate}

\textbf{定理 33.4.1}：\(M_{m \times n}(F)\) 在上述运算下构成 \(F\)
上的向量空间，且 \(\dim M_{m \times n}(F) = mn\)。

\textbf{证明}：向量空间公理由逐分量的域运算直接验证。标准基为
\(\{E_{ij} \mid 1 \leq i \leq m, \, 1 \leq j \leq n\}\)，其中 \(E_{ij}\)
在第 \((i, j)\) 位置为 \(1\)，其余为 \(0\)。共有 \(mn\)
个基向量。\(\blacksquare\)

\textbf{定义 33.4.3（线性映射的矩阵）}：设 \(V\) 有基
\(B = \{v_1, \ldots, v_n\}\)，\(W\) 有基
\(C = \{w_1, \ldots, w_m\}\)，\(T \in \mathcal{L}(V, W)\)。

对每个 \(j = 1, \ldots, n\)，将 \(T(v_j)\) 表示为 \(C\) 的线性组合：

\[
T(v_j) = \sum_{i=1}^{m} a_{ij} w_i
\]

则 \(m \times n\) 矩阵 \(A = (a_{ij})\) 称为 \(T\) 关于基 \(B\) 和 \(C\)
的\textbf{矩阵表示}，记为 \([T]_{B}^{C}\) 或 \(\mathcal{M}(T; B, C)\)。

\textbf{定理 33.4.2（矩阵表示与坐标的关系）}：设 \(v \in V\) 关于基
\(B\) 的坐标为 \(x = (x_1, \ldots, x_n)^T\)（列向量），则 \(T(v)\)
关于基 \(C\) 的坐标 \(y = (y_1, \ldots, y_m)^T\) 满足

\[
y = A x, \quad \text{即} \quad y_i = \sum_{j=1}^{n} a_{ij} x_j
\]

\textbf{证明}： \[
T(v) = T\left(\sum_{j=1}^{n} x_j v_j\right) = \sum_{j=1}^{n} x_j T(v_j) = \sum_{j=1}^{n} x_j \sum_{i=1}^{m} a_{ij} w_i = \sum_{i=1}^{m} \left(\sum_{j=1}^{n} a_{ij} x_j\right) w_i
\] 因此 \(y_i = \sum_{j=1}^{n} a_{ij} x_j\)。\(\blacksquare\)

\textbf{定理 33.4.3（矩阵表示是线性同构）}：映射
\(\mathcal{M}(\cdot; B, C) : \mathcal{L}(V, W) \to M_{m \times n}(F)\)
是向量空间同构。特别地，\(\dim \mathcal{L}(V, W) = (\dim V)(\dim W)\)。

\textbf{证明}：线性和双射由矩阵表示的定义直接验证。\(\blacksquare\)

\bigskip\hrule\bigskip

\subsubsection{33.5
矩阵乘法与复合}\label{ux77e9ux9635ux4e58ux6cd5ux4e0eux590dux5408}

\textbf{定义 33.5.1（矩阵乘法）}：设
\(A \in M_{m \times n}(F)\)，\(B \in M_{n \times p}(F)\)。定义乘积
\(C = AB \in M_{m \times p}(F)\) 为

\[
c_{ij} = \sum_{k=1}^{n} a_{ik} b_{kj}, \quad i = 1, \ldots, m, \; j = 1, \ldots, p
\]

\textbf{定理 33.5.1（矩阵乘法对应线性映射的复合）}：设
\(V \xrightarrow{S} W \xrightarrow{T} U\) 为线性映射，\(B, C, D\) 分别为
\(V, W, U\) 的基。则

\[
[T \circ S]_{B}^{D} = [T]_{C}^{D} \cdot [S]_{B}^{C}
\]

\textbf{证明}：设
\(B = \{v_1, \ldots, v_n\}\)，\(C = \{w_1, \ldots, w_m\}\)，\(D = \{u_1, \ldots, u_p\}\)。令
\(A = [T]_{C}^{D}\)，\(B = [S]_{B}^{C}\)。则 \[
(T \circ S)(v_j) = T(S(v_j)) = T\left(\sum_{k=1}^{m} b_{kj} w_k\right) = \sum_{k=1}^{m} b_{kj} T(w_k) = \sum_{k=1}^{m} b_{kj} \sum_{i=1}^{p} a_{ik} u_i = \sum_{i=1}^{p} \left(\sum_{k=1}^{m} a_{ik} b_{kj}\right) u_i
\] 因此 \([T \circ S]_{B}^{D}\) 的第 \((i, j)\) 元为
\(\sum_{k=1}^{m} a_{ik} b_{kj}\)，正是 \(AB\) 的第 \((i, j)\)
元。\(\blacksquare\)

\textbf{定理
33.5.2（矩阵乘法的基本性质）}：设矩阵的维数使得以下乘法有意义，则

\begin{enumerate}
\def\labelenumi{(\roman{enumi})}
\tightlist
\item
  \textbf{结合律}：\((AB)C = A(BC)\)
\item
  \textbf{分配律}：\(A(B + C) = AB + AC\)，\((A + B)C = AC + BC\)
\item
  \textbf{标量乘法}：\(\lambda(AB) = (\lambda A)B = A(\lambda B)\)
\item
  \textbf{单位矩阵}：\(I_n A = A I_m = A\)（\(A \in M_{m \times n}(F)\)，\(I_n\)
  为 \(n\) 阶单位矩阵）
\end{enumerate}

\textbf{证明}：(i)
由矩阵乘法对应线性映射的复合，复合满足结合律。也可直接验证： \[
((AB)C)_{ij} = \sum_{k} (AB)_{ik} c_{kj} = \sum_{k} \sum_{\ell} a_{i\ell} b_{\ell k} c_{kj} = \sum_{\ell} a_{i\ell} \sum_{k} b_{\ell k} c_{kj} = (A(BC))_{ij}
\]

(ii)-(iv) 直接验证。\(\blacksquare\)

\textbf{注记 33.5.1}：矩阵乘法一般不满足交换律，即 \(AB \neq BA\)
一般成立。

\textbf{定义 33.5.2（转置）}：设 \(A \in M_{m \times n}(F)\)。\(A\)
的\textbf{转置} \(A^T \in M_{n \times m}(F)\) 定义为
\((A^T)_{ij} = A_{ji}\)。

\textbf{定理 33.5.3（转置的性质）}：

\begin{enumerate}
\def\labelenumi{(\roman{enumi})}
\tightlist
\item
  \((A^T)^T = A\)
\item
  \((A + B)^T = A^T + B^T\)
\item
  \((\lambda A)^T = \lambda A^T\)
\item
  \((AB)^T = B^T A^T\)
\end{enumerate}

\textbf{证明}：(i)-(iii) 直接验证。(iv)
\((AB)^T_{ij} = (AB)_{ji} = \sum_k a_{jk} b_{ki} = \sum_k (B^T)_{ik} (A^T)_{kj} = (B^T A^T)_{ij}\)。\(\blacksquare\)

\bigskip\hrule\bigskip

\subsubsection{33.6 基变换}\label{ux57faux53d8ux6362}

\textbf{定义 33.6.1（基变换矩阵）}：设 \(V\) 有两组基
\(B = \{v_1, \ldots, v_n\}\) 和 \(B' = \{v_1', \ldots, v_n'\}\)。将
\(B'\) 中的每个向量用 \(B\) 表示：

\[
v_j' = \sum_{i=1}^{n} p_{ij} v_i, \quad j = 1, \ldots, n
\]

则矩阵 \(P = (p_{ij})\) 称为从 \(B\) 到 \(B'\)
的\textbf{基变换矩阵}（或\textbf{过渡矩阵}）。

\textbf{定理 33.6.1（基变换矩阵的性质）}：

\begin{enumerate}
\def\labelenumi{(\roman{enumi})}
\tightlist
\item
  \(P\) 可逆，且 \(P^{-1}\) 是从 \(B'\) 到 \(B\) 的基变换矩阵。
\item
  设 \(v \in V\) 关于 \(B\) 的坐标为 \(x\)，关于 \(B'\) 的坐标为
  \(x'\)，则 \(x = P x'\)。
\end{enumerate}

\textbf{证明}：

\begin{enumerate}
\def\labelenumi{(\roman{enumi})}
\item
  设 \(Q\) 为从 \(B'\) 到 \(B\) 的基变换矩阵，即将 \(v_i\) 用 \(B'\)
  表示。则
  \(v_j = \sum_k q_{kj} v_k' = \sum_k q_{kj} \sum_i p_{ik} v_i = \sum_i (\sum_k p_{ik} q_{kj}) v_i\)。由基表示的唯一性，\(\sum_k p_{ik} q_{kj} = \delta_{ij}\)，即
  \(PQ = I_n\)。同理 \(QP = I_n\)。故 \(P\) 可逆且 \(P^{-1} = Q\)。
\item
  \(v = \sum_j x_j' v_j' = \sum_j x_j' \sum_i p_{ij} v_i = \sum_i (\sum_j p_{ij} x_j') v_i\)，故
  \(x_i = \sum_j p_{ij} x_j'\)，即 \(x = P x'\)。\(\blacksquare\)
\end{enumerate}

\textbf{定理 33.6.2（线性映射在不同基下的矩阵）}：设
\(T \in \mathcal{L}(V, W)\)，\(B, B'\) 为 \(V\) 的基，\(C, C'\) 为 \(W\)
的基。\(P\) 为从 \(B\) 到 \(B'\) 的基变换矩阵，\(Q\) 为从 \(C\) 到
\(C'\) 的基变换矩阵。则

\[
[T]_{B'}^{C'} = Q^{-1} [T]_{B}^{C} P
\]

\textbf{证明}：对任意 \(v \in V\)，设
\(x = [v]_B\)，\(x' = [v]_{B'}\)，由定理 33.6.1 \(x = P x'\)。同理
\([T(v)]_C = [T]_{B}^{C} x\)，\([T(v)]_{C'} = [T]_{B'}^{C'} x'\)，且
\([T(v)]_C = Q [T(v)]_{C'}\)。

因此 \(Q [T]_{B'}^{C'} x' = [T]_{B}^{C} P x'\)。由于 \(x'\)
任意，\(Q [T]_{B'}^{C'} = [T]_{B}^{C} P\)，即
\([T]_{B'}^{C'} = Q^{-1} [T]_{B}^{C} P\)。\(\blacksquare\)

\textbf{定义 33.6.2（相似矩阵）}：设
\(A, B \in M_{n \times n}(F)\)。若存在可逆矩阵
\(P \in M_{n \times n}(F)\) 使得 \(B = P^{-1}AP\)，则称 \(A\) 与 \(B\)
\textbf{相似}。

\textbf{推论 33.6.3}：同一线性算子在不同基下的矩阵是相似的。

\textbf{证明}：在定理 33.6.2 中取
\(V = W\)，\(B = C\)，\(B' = C'\)，\(P = Q\)，则
\([T]_{B'}^{B'} = P^{-1}[T]_{B}^{B} P\)。\(\blacksquare\)

\bigskip\hrule\bigskip

\subsection{Ch 34 行列式}\label{ch-34-ux884cux5217ux5f0f}

\subsubsection{34.1
行列式的公理化定义}\label{ux884cux5217ux5f0fux7684ux516cux7406ux5316ux5b9aux4e49}

\textbf{定义 34.1.1（置换与奇偶性）}

设 \(n \in \mathbb{N}^+\)。\(\{1, 2, \ldots, n\}\) 到自身的双射
\(\sigma\) 称为一个 \textbf{\(n\) 元置换}。所有 \(n\) 元置换的集合记为
\(S_n\)。

一个置换 \(\sigma\)
称为\textbf{对换}，若它仅交换两个元素而保持其余不变。每个置换可分解为若干个对换的复合。若分解中对换个数的奇偶性唯一确定，则分别称
\(\sigma\) 为\textbf{偶置换}或\textbf{奇置换}。定义 \(\sigma\)
的\textbf{符号}为

\[
\operatorname{sgn}(\sigma) = \begin{cases}
+1, & \text{若 } \sigma \text{ 是偶置换} \\
-1, & \text{若 } \sigma \text{ 是奇置换}
\end{cases}
\]

\textbf{定理
34.1.1（奇偶性的良定义）}：一个置换不可能同时分解为偶数个和奇数个对换的复合。

\textbf{证明}：设
\(P(x_1, \ldots, x_n) = \prod_{1 \leq i < j \leq n} (x_i - x_j)\)。对任意置换
\(\sigma\)，定义
\(\sigma P(x_1, \ldots, x_n) = P(x_{\sigma(1)}, \ldots, x_{\sigma(n)})\)。可以证明
\(\sigma P = \operatorname{sgn}(\sigma) P\)（其中
\(\operatorname{sgn}(\sigma) = \pm 1\)）。一个对换使 \(P\) 变号，故
\(k\) 个对换的复合使 \(P\) 乘以
\((-1)^k\)。若同一置换有两种分解，其中一种有 \(k\) 个对换、另一种有
\(\ell\) 个对换，则 \((-1)^k = (-1)^\ell\)，故 \(k\) 与 \(\ell\)
同奇偶。\(\blacksquare\)

\textbf{定义 34.1.2（行列式）}：设
\(A = (a_{ij}) \in M_{n \times n}(F)\)。\(A\) 的\textbf{行列式}
\(\det A\)（或 \(|A|\)）定义为

\[
\det A = \sum_{\sigma \in S_n} \operatorname{sgn}(\sigma) \, a_{1\sigma(1)} a_{2\sigma(2)} \cdots a_{n\sigma(n)}
\]

\textbf{评注 34.1.1}：上式称为行列式的 \textbf{Leibniz 公式}。\(n\)
阶行列式共有 \(n!\) 项。

\textbf{定义 34.1.3（行列式的等价公理化定义）}：函数
\(D : M_{n \times n}(F) \to F\)
称为\textbf{行列式函数}，若它满足以下三条公理：

\begin{enumerate}
\def\labelenumi{\arabic{enumi}.}
\tightlist
\item
  \textbf{多重线性}：\(D\)
  对每一行是线性的（固定其他行，该行向量上为线性函数）
\item
  \textbf{交错性}：若两行相同，则 \(D(A) = 0\)
\item
  \textbf{规范性}：\(D(I_n) = 1\)（其中 \(I_n\) 是 \(n\) 阶单位矩阵）
\end{enumerate}

\textbf{定理 34.1.2（行列式公理的等价性）}：Leibniz
公式定义的行列式满足公理化定义的三条性质，且任何满足这三条公理的函数必然等于
Leibniz 公式。

\textbf{证明}：先证 Leibniz 公式满足三条公理。多重线性：每项
\(a_{1\sigma(1)} \cdots a_{n\sigma(n)}\)
对每行是线性的，和亦然。交错性：若第 \(i\) 行与第 \(j\)
行相同（\(i < j\)），设 \(\tau\) 为交换 \(i\) 和 \(j\) 的对换。将置换按
\(\{\sigma, \sigma\tau\}\)
配对，\(\operatorname{sgn}(\sigma\tau) = -\operatorname{sgn}(\sigma)\)，而
\(a_{1\sigma(1)} \cdots a_{n\sigma(n)} = a_{1\sigma\tau(1)} \cdots a_{n\sigma\tau(n)}\)（因为第
\(i\) 行和第 \(j\) 行相同），两项抵消。规范性：\(I_n\) 中仅 \(\sigma\)
为恒等置换时乘积非零，此时
\(\operatorname{sgn}(\text{id}) = 1\)，\(1 \cdot 1 \cdots 1 = 1\)。

再证唯一性：设 \(D\) 满足三条公理。将 \(A\) 的第 \(j\) 行写为
\(a_j = \sum_{i=1}^{n} a_{ij} e_i\)（其中 \(e_i\)
是标准基）。由多重线性， \[
D(A) = D\left(\sum_{i_1} a_{i_1 1} e_{i_1}, \ldots, \sum_{i_n} a_{i_n n} e_{i_n}\right) = \sum_{i_1, \ldots, i_n} a_{i_1 1} \cdots a_{i_n n} D(e_{i_1}, \ldots, e_{i_n})
\] 由交错性，若 \(i_1, \ldots, i_n\)
中有重复，\(D(e_{i_1}, \ldots, e_{i_n}) = 0\)。故只需对
\((i_1, \ldots, i_n)\) 是 \(\{1, \ldots, n\}\) 的置换求和。设
\(i_k = \sigma(k)\)，则
\(D(e_{\sigma(1)}, \ldots, e_{\sigma(n)}) = \operatorname{sgn}(\sigma) D(e_1, \ldots, e_n) = \operatorname{sgn}(\sigma) \cdot 1\)。因此
\(D(A) = \sum_{\sigma \in S_n} \operatorname{sgn}(\sigma) a_{\sigma(1)1} \cdots a_{\sigma(n)n} = \sum_{\sigma \in S_n} \operatorname{sgn}(\sigma) a_{1\sigma^{-1}(1)} \cdots a_{n\sigma^{-1}(n)}\)。当
\(\sigma\) 遍历 \(S_n\) 时 \(\sigma^{-1}\) 也遍历 \(S_n\)，且
\(\operatorname{sgn}(\sigma^{-1}) = \operatorname{sgn}(\sigma)\)，故等于
Leibniz 公式。\(\blacksquare\)

\bigskip\hrule\bigskip

\subsubsection{34.2
行列式的基本性质}\label{ux884cux5217ux5f0fux7684ux57faux672cux6027ux8d28}

\textbf{定理 34.2.1（行列式的基本性质）}：设
\(A, B \in M_{n \times n}(F)\)。

\begin{enumerate}
\def\labelenumi{(\roman{enumi})}
\tightlist
\item
  \textbf{转置不变性}：\(\det(A^T) = \det A\)
\item
  \textbf{行交换}：交换两行，行列式变号
\item
  \textbf{行倍乘}：某行乘以 \(\lambda\)，行列式乘以 \(\lambda\)
\item
  \textbf{行加法}：将一行的倍数加到另一行，行列式不变
\item
  \textbf{乘法性质}：\(\det(AB) = \det A \cdot \det B\)
\item
  \textbf{可逆性判定}：\(A\) 可逆 \(\iff \det A \neq 0\)
\end{enumerate}

\textbf{证明}：

\begin{enumerate}
\def\labelenumi{(\roman{enumi})}
\tightlist
\item
  \(\det(A^T) = \sum_{\sigma} \operatorname{sgn}(\sigma) a_{\sigma(1)1} \cdots a_{\sigma(n)n}\)。令
  \(\tau = \sigma^{-1}\)，则
  \(a_{\sigma(k)k} = a_{k \tau(k)}\)（交换因子顺序），且
  \(\operatorname{sgn}(\tau) = \operatorname{sgn}(\sigma)\)。因此
  \(\det(A^T) = \sum_{\tau} \operatorname{sgn}(\tau) a_{1\tau(1)} \cdots a_{n\tau(n)} = \det A\)。
\end{enumerate}

(ii)-(iv) 由公理化定义中的多重线性、交错性和规范性直接得出。

\begin{enumerate}
\def\labelenumi{(\alph{enumi})}
\setcounter{enumi}{21}
\tightlist
\item
  将 \(A\) 固定，考虑函数 \(f(B) = \det(AB)\)。\(f\) 对 \(B\)
  的行也是多重线性和交错的（因为 \(A\) 的行乘以 \(B\)
  的列的线性组合）。由唯一性定理，\(f(B) = f(I_n) \det B = \det A \cdot \det B\)。
\end{enumerate}

\begin{enumerate}
\def\labelenumi{(\roman{enumi})}
\setcounter{enumi}{5}
\tightlist
\item
  若 \(A\) 可逆，则 \(I_n = A A^{-1}\)，故
  \(1 = \det I_n = \det A \cdot \det(A^{-1})\)，因此 \(\det A \neq 0\)。
\end{enumerate}

若 \(A\) 不可逆，则
\(\operatorname{rank} A < n\)，行向量线性相关。不妨设某行是其余行的线性组合。由多重线性，该行可分解，每一项都有两行成比例，由交错性（或行倍乘后行交换），\(\det A = 0\)。\(\blacksquare\)

\textbf{推论 34.2.2}：\(\det(A^{-1}) = (\det A)^{-1}\)。

\textbf{证明}：由
\(1 = \det(I_n) = \det(A A^{-1}) = \det A \cdot \det(A^{-1})\)。\(\blacksquare\)

\textbf{推论 34.2.3}：相似矩阵有相同的行列式。

\textbf{证明}：若 \(B = P^{-1}AP\)，则
\(\det B = \det(P^{-1}) \det A \det P = (\det P)^{-1} \det A \det P = \det A\)。\(\blacksquare\)

\bigskip\hrule\bigskip

\subsubsection{34.3
行列式的展开}\label{ux884cux5217ux5f0fux7684ux5c55ux5f00}

\textbf{定义 34.3.1（余子式与代数余子式）}：设
\(A = (a_{ij}) \in M_{n \times n}(F)\)。

\begin{enumerate}
\def\labelenumi{\arabic{enumi}.}
\tightlist
\item
  \(A\) 中划去第 \(i\) 行和第 \(j\) 列后得到的 \((n-1) \times (n-1)\)
  矩阵的行列式称为 \(a_{ij}\) 的\textbf{余子式}，记为 \(M_{ij}\)。
\item
  \(C_{ij} = (-1)^{i+j} M_{ij}\) 称为 \(a_{ij}\) 的\textbf{代数余子式}。
\end{enumerate}

\textbf{定理 34.3.1（Laplace 展开）}：设 \(A \in M_{n \times n}(F)\)。

\begin{enumerate}
\def\labelenumi{(\roman{enumi})}
\tightlist
\item
  \textbf{按第 \(i\)
  行展开}：\(\displaystyle \det A = \sum_{j=1}^{n} a_{ij} C_{ij}\)
\item
  \textbf{按第 \(j\)
  列展开}：\(\displaystyle \det A = \sum_{i=1}^{n} a_{ij} C_{ij}\)
\item
  \textbf{错行展开为零}：\(\displaystyle \sum_{j=1}^{n} a_{ij} C_{kj} = 0 \quad (i \neq k)\)
\item
  \textbf{错列展开为零}：\(\displaystyle \sum_{i=1}^{n} a_{ij} C_{ik} = 0 \quad (j \neq k)\)
\end{enumerate}

\textbf{证明}：(i) 固定第 \(i\) 行。由多重线性，\(\det A\)
是这个行向量的线性函数。将 \(A\) 的第 \(i\) 行写为
\(\sum_{j=1}^{n} a_{ij} e_j\)（\(e_j\) 为标准基向量）。则
\(\det A = \sum_j a_{ij} \det A_{(j)}\)，其中 \(A_{(j)}\) 是将 \(A\)
的第 \(i\) 行替换为 \(e_j\) 的矩阵。通过行交换将 \(A_{(j)}\) 的第 \(i\)
行移到第 \(1\) 行，列交换将第 \(j\) 列移到第 \(1\) 列，需要
\((i-1) + (j-1) = i+j-2\) 次交换，每次变号。故
\(\det A_{(j)} = (-1)^{i+j-2} M_{ij} = (-1)^{i+j} M_{ij} = C_{ij}\)。

\begin{enumerate}
\def\labelenumi{(\roman{enumi})}
\setcounter{enumi}{2}
\tightlist
\item
  考虑将 \(A\) 的第 \(k\) 行替换为第 \(i\) 行得到的矩阵
  \(A'\)（\(i \neq k\)）。\(A'\) 有两行相同，故
  \(\det A' = 0\)。另一方面，按第 \(k\) 行展开 \(A'\)，第 \(k\)
  行的元素为 \(a_{ij}\)（\(j = 1, \ldots, n\)），余子式为
  \(C_{kj}\)（因为第 \(k\) 行的余子式与第 \(k\) 行元素无关）。故
  \(0 = \det A' = \sum_{j} a_{ij} C_{kj}\)。\(\blacksquare\)
\end{enumerate}

\textbf{定义 34.3.2（伴随矩阵）}：设 \(A \in M_{n \times n}(F)\)。\(A\)
的\textbf{伴随矩阵} \(\operatorname{adj}(A)\) 定义为
\((\operatorname{adj}(A))_{ij} = C_{ji}\)（即代数余子式矩阵的转置）。

\textbf{定理
34.3.2（伴随矩阵的性质）}：\(A \cdot \operatorname{adj}(A) = \operatorname{adj}(A) \cdot A = (\det A) I_n\)。

\textbf{证明}：\((A \cdot \operatorname{adj}(A))_{ik} = \sum_{j} a_{ij} (\operatorname{adj}(A))_{jk} = \sum_{j} a_{ij} C_{kj}\)。当
\(i = k\) 时，由 Laplace 展开，\(\sum_j a_{ij} C_{ij} = \det A\)。当
\(i \neq k\) 时，由定理 34.3.1(iii)，\(\sum_j a_{ij} C_{kj} = 0\)。故
\(A \cdot \operatorname{adj}(A) = (\det A) I_n\)。同理可证
\(\operatorname{adj}(A) \cdot A = (\det A) I_n\)。\(\blacksquare\)

\textbf{推论 34.3.3（逆矩阵公式）}：若 \(\det A \neq 0\)，则
\(A^{-1} = \frac{1}{\det A} \operatorname{adj}(A)\)。

\textbf{证明}：由定理 34.3.2 直接得到。\(\blacksquare\)

\bigskip\hrule\bigskip

\subsubsection{34.4 Cramer 法则}\label{cramer-ux6cd5ux5219}

\textbf{定理 34.4.1（Cramer 法则）}：设 \(A \in M_{n \times n}(F)\)
可逆，\(b \in F^n\)。则线性方程组 \(Ax = b\) 的唯一解为

\[
x_j = \frac{\det A_j}{\det A}, \quad j = 1, 2, \ldots, n
\]

其中 \(A_j\) 是将 \(A\) 的第 \(j\) 列替换为 \(b\) 得到的矩阵。

\textbf{证明}：由
\(A^{-1} = \frac{1}{\det A} \operatorname{adj}(A)\)，\(x = A^{-1}b\)，故
\[
x_j = \frac{1}{\det A} \sum_{i=1}^{n} (\operatorname{adj}(A))_{ji} b_i = \frac{1}{\det A} \sum_{i=1}^{n} C_{ij} b_i
\] 而 \(\sum_{i=1}^{n} C_{ij} b_i\) 正是将 \(A\) 的第 \(j\) 列替换为
\(b\) 后按第 \(j\) 列展开的结果，即 \(\det A_j\)。\(\blacksquare\)

\bigskip\hrule\bigskip

\subsubsection{34.5
行列式与体积}\label{ux884cux5217ux5f0fux4e0eux4f53ux79ef}

\textbf{定理 34.5.1（平行多面体的体积）}：设
\(v_1, \ldots, v_n \in \mathbb{R}^n\) 为列向量。以原点为一个顶点，以
\(v_1, \ldots, v_n\) 为棱的 \(n\) 维平行多面体的 \(n\) 维体积为

\[
\operatorname{Vol}(v_1, \ldots, v_n) = \left|\det(v_1, \ldots, v_n)\right|
\]

\textbf{证明}（概要）：设
\(f(v_1, \ldots, v_n) = \operatorname{Vol}(v_1, \ldots, v_n)\)。体积函数满足：

\begin{enumerate}
\def\labelenumi{\arabic{enumi}.}
\tightlist
\item
  对每个 \(v_i\) 是线性的（缩放棱则体积缩放到相同倍数）
\item
  若 \(v_i = v_j\)（\(i \neq j\)），则体积为 \(0\)（退化）
\item
  对标准基，体积为 \(1\)（单位立方体）
\end{enumerate}

这正是行列式公理化定义的三个条件（在绝对值意义下）。由于 \(|\det|\)
也满足这三条，由唯一性，\(f = |\det|\)。\(\blacksquare\)

\textbf{推论 34.5.2（线性变换的体积变化）}：设
\(T : \mathbb{R}^n \to \mathbb{R}^n\)
为线性变换，\(A = [T]_{\text{std}}\)。则对任意可测集
\(\Omega \subseteq \mathbb{R}^n\)，

\[
\operatorname{Vol}(T(\Omega)) = |\det A| \cdot \operatorname{Vol}(\Omega)
\]

\textbf{证明}：由积分的变量替换公式（参见卷四第24章），变换的 Jacobi
行列式为 \(|\det A|\)。\(\blacksquare\)

\bigskip\hrule\bigskip

\subsection{Ch 35
特征值与特征向量}\label{ch-35-ux7279ux5f81ux503cux4e0eux7279ux5f81ux5411ux91cf}

\subsubsection{35.1
特征值与特征向量}\label{ux7279ux5f81ux503cux4e0eux7279ux5f81ux5411ux91cf}

\textbf{定义 35.1.1（特征值与特征向量）}：设 \(V\) 为 \(F\)
上的向量空间，\(T \in \mathcal{L}(V)\)。若存在 \(\lambda \in F\)
和非零向量 \(v \in V\)（\(v \neq \mathbf{0}\)）使得

\[
T(v) = \lambda v
\]

则称 \(\lambda\) 为 \(T\) 的\textbf{特征值}，\(v\) 为 \(T\) 的属于
\(\lambda\) 的\textbf{特征向量}。

\textbf{定义 35.1.2（特征空间）}：设 \(\lambda\) 为 \(T\)
的特征值。\(T\) 的属于 \(\lambda\) 的\textbf{特征空间}为

\[
E_{\lambda} = \{v \in V \mid T(v) = \lambda v\} = \ker(T - \lambda I)
\]

\(E_{\lambda}\) 是 \(V\) 的子空间。\(\dim E_{\lambda}\) 称为 \(\lambda\)
的\textbf{几何重数}。

\textbf{定理 35.1.1（不同特征值的特征向量线性无关）}：设
\(\lambda_1, \ldots, \lambda_k\) 为 \(T\)
的两两不同的特征值，\(v_1, \ldots, v_k\) 为对应的特征向量。则
\(\{v_1, \ldots, v_k\}\) 线性无关。

\textbf{证明}：对 \(k\) 归纳。

\textbf{基础} \(k = 1\)：\(v_1 \neq \mathbf{0}\)，单个非零向量线性无关。

\textbf{归纳}：假设 \(\{v_1, \ldots, v_{k-1}\}\) 线性无关。设
\(\sum_{i=1}^{k} \alpha_i v_i = \mathbf{0}\)。两边作用 \(T\)： \[
\sum_{i=1}^{k} \alpha_i \lambda_i v_i = \mathbf{0}
\] 将原式乘以 \(\lambda_k\)： \[
\sum_{i=1}^{k} \alpha_i \lambda_k v_i = \mathbf{0}
\]
两式相减：\(\sum_{i=1}^{k-1} \alpha_i (\lambda_i - \lambda_k) v_i = \mathbf{0}\)。由归纳假设，\(\{v_1, \ldots, v_{k-1}\}\)
线性无关，故
\(\alpha_i (\lambda_i - \lambda_k) = 0\)（\(i = 1, \ldots, k-1\)）。由于
\(\lambda_i \neq \lambda_k\)，\(\alpha_i = 0\)。代入原式得
\(\alpha_k v_k = \mathbf{0}\)，故 \(\alpha_k = 0\)。\(\blacksquare\)

\textbf{推论 35.1.2}：\(n\) 维向量空间上的线性算子最多有 \(n\)
个不同的特征值。

\textbf{证明}：由定理 35.1.1，不同特征值对应的特征向量线性无关，故不超过
\(\dim V = n\)。\(\blacksquare\)

\bigskip\hrule\bigskip

\subsubsection{35.2 特征多项式}\label{ux7279ux5f81ux591aux9879ux5f0f}

\textbf{定义 35.2.1（特征多项式）}：设
\(A \in M_{n \times n}(F)\)。\(A\) 的\textbf{特征多项式}为

\[
p_A(\lambda) = \det(\lambda I_n - A)
\]

\textbf{注记 35.2.1}：有些教材定义
\(\det(A - \lambda I_n)\)，两者仅差一个符号因子
\((-1)^n\)，不影响特征值的确定。

\textbf{定理 35.2.1}：\(\lambda\) 是 \(A\) 的特征值当且仅当
\(p_A(\lambda) = 0\)，即 \(\lambda\) 是特征多项式的根。

\textbf{证明}：\(\lambda\) 是特征值 \(\iff \exists v \neq \mathbf{0}\)
使得
\(Av = \lambda v \iff (\lambda I_n - A)v = \mathbf{0} \iff \lambda I_n - A\)
不可逆
\(\iff \det(\lambda I_n - A) = 0 \iff p_A(\lambda) = 0\)。\(\blacksquare\)

\textbf{定理
35.2.2（特征多项式的形式）}：\(p_A(\lambda) = \lambda^n - (\operatorname{tr} A) \lambda^{n-1} + \cdots + (-1)^n \det A\)，其中
\(\operatorname{tr} A = \sum_{i=1}^{n} a_{ii}\) 为 \(A\) 的迹。

\textbf{证明}：展开
\(\det(\lambda I_n - A) = \sum_{\sigma} \operatorname{sgn}(\sigma) \prod_{i=1}^{n} (\lambda \delta_{i\sigma(i)} - a_{i\sigma(i)})\)。\(\lambda^n\)
项仅来自 \(\sigma\) 为恒等置换，系数为 \(1\)。\(\lambda^{n-1}\)
项来自恒等置换中某个 \(i\) 处取 \(-a_{ii}\) 而其余 \(n-1\) 处取
\(\lambda\)，系数为 \(-\sum_i a_{ii} = -\operatorname{tr} A\)。常数项为
\(p_A(0) = \det(-A) = (-1)^n \det A\)。\(\blacksquare\)

\textbf{定义 35.2.2（代数重数）}：特征值 \(\lambda\) 作为
\(p_A(\lambda)\) 的根的重数称为 \(\lambda\) 的\textbf{代数重数}。

\textbf{定理 35.2.3（几何重数不超过代数重数）}：设 \(\lambda\) 为 \(A\)
的特征值，则 \(\dim E_{\lambda} \leq \lambda\) 的代数重数。

\textbf{证明}：设 \(\dim E_{\lambda} = k\)。取 \(E_{\lambda}\) 的基
\(\{v_1, \ldots, v_k\}\)，将其扩充为 \(F^n\) 的基。在此基下，\(A\)
的矩阵为 \[
\begin{pmatrix} \lambda I_k & B \\ 0 & C \end{pmatrix}
\] 其特征多项式为
\(\det(\lambda I_n - A) = \det((\lambda - \lambda) I_k) \det(\lambda I_{n-k} - C) = 0 \cdot \det(\lambda I_{n-k} - C)\)。更精确地，\(p_A(t) = (t - \lambda)^k \det(t I_{n-k} - C)\)，故
\(t - \lambda\) 的重数至少为 \(k\)。\(\blacksquare\)

\bigskip\hrule\bigskip

\subsubsection{35.3 对角化}\label{ux5bf9ux89d2ux5316}

\textbf{定义 35.3.1（可对角化）}：线性算子
\(T \in \mathcal{L}(V)\)（或矩阵
\(A \in M_{n \times n}(F)\)）称为\textbf{可对角化}的，若存在 \(V\)
的一组基，使得 \(T\) 在此基下的矩阵为对角矩阵。等价地，\(V\) 有一组由
\(T\) 的特征向量组成的基。

\textbf{定理 35.3.1（可对角化的判定）}：设
\(T \in \mathcal{L}(V)\)，\(\dim V = n\)。以下等价：

\begin{enumerate}
\def\labelenumi{(\roman{enumi})}
\tightlist
\item
  \(T\) 可对角化
\item
  \(V\) 有一组由 \(T\) 的特征向量组成的基
\item
  \(V = E_{\lambda_1} \oplus E_{\lambda_2} \oplus \cdots \oplus E_{\lambda_k}\)（特征空间的直和分解）
\item
  \(T\) 有 \(n\) 个线性无关的特征向量
\end{enumerate}

\textbf{证明}：(i) \(\iff\) (ii)：由定义。(ii) \(\iff\)
(iv)：维数条件。(ii) \(\iff\)
(iii)：不同特征值的特征向量线性无关，且特征向量张成 \(V\)，故 \(V\)
是特征空间的直和。\(\blacksquare\)

\textbf{定理 35.3.2（充分条件）}：若 \(T \in \mathcal{L}(V)\) 有
\(n = \dim V\) 个两两不同的特征值，则 \(T\) 可对角化。

\textbf{证明}：由定理 35.1.1，\(n\)
个不同特征值的特征向量线性无关。由定理 32.5.5(ii)，\(n\)
个线性无关向量构成 \(V\) 的基。\(\blacksquare\)

\textbf{定理 35.3.3（对角化的充要条件）}：\(T \in \mathcal{L}(V)\)
可对角化当且仅当其特征多项式在 \(F\)
上可分解为一次因式的乘积，且每个特征值的几何重数等于其代数重数。

\textbf{证明}：

（\(\Rightarrow\)）若 \(T\) 可对角化，则存在基使得 \(T\) 的矩阵为
\(\operatorname{diag}(\lambda_1, \ldots, \lambda_n)\)。此时
\(p_T(t) = \prod_{i=1}^{n} (t - \lambda_i)\)，且 \(\dim E_{\lambda_i}\)
等于 \(\lambda_i\) 在 \(\{\lambda_1, \ldots, \lambda_n\}\)
中出现的次数，即代数重数。

（\(\Leftarrow\)）设
\(p_T(t) = \prod_{i=1}^{k} (t - \lambda_i)^{m_i}\)（\(\lambda_i\)
两两不同），且 \(\dim E_{\lambda_i} = m_i\)。由定理
35.1.1，不同特征空间的和是直和。\(\dim(\bigoplus_i E_{\lambda_i}) = \sum_i m_i = n = \dim V\)，故
\(\bigoplus_i E_{\lambda_i} = V\)。\(\blacksquare\)

\bigskip\hrule\bigskip

\subsubsection{35.4 Cayley-Hamilton
定理}\label{cayley-hamilton-ux5b9aux7406}

\textbf{定理 35.4.1（Cayley-Hamilton 定理）}：设
\(A \in M_{n \times n}(F)\)，\(p_A(\lambda) = \det(\lambda I_n - A)\)
为其特征多项式。则 \(p_A(A) = 0\)（零矩阵）。

\textbf{证明}：考虑矩阵 \(\lambda I_n - A\) 的伴随矩阵
\(\operatorname{adj}(\lambda I_n - A)\)。其元素是 \(\lambda\)
的多项式，次数不超过 \(n-1\)。由定理 34.3.2，

\[
(\lambda I_n - A) \cdot \operatorname{adj}(\lambda I_n - A) = \det(\lambda I_n - A) I_n = p_A(\lambda) I_n
\]

将 \(\operatorname{adj}(\lambda I_n - A)\) 写为
\(B_0 + B_1 \lambda + \cdots + B_{n-1} \lambda^{n-1}\)（\(B_i \in M_{n \times n}(F)\)）。设
\(p_A(\lambda) = \lambda^n + c_{n-1}\lambda^{n-1} + \cdots + c_0\)。则

\[
(\lambda I_n - A)(B_0 + B_1 \lambda + \cdots + B_{n-1} \lambda^{n-1}) = (\lambda^n + c_{n-1}\lambda^{n-1} + \cdots + c_0) I_n
\]

比较两边 \(\lambda^k\) 的系数：

\[
\begin{aligned}
\lambda^0 &: -A B_0 = c_0 I_n \\
\lambda^1 &: B_0 - A B_1 = c_1 I_n \\
&\;\;\vdots \\
\lambda^{n-1} &: B_{n-2} - A B_{n-1} = c_{n-1} I_n \\
\lambda^n &: B_{n-1} = I_n
\end{aligned}
\]

将第 \(k\) 个等式左乘 \(A^k\)（\(k = 0, \ldots, n\)）并求和： \[
\begin{aligned}
&(-A B_0) + A(B_0 - A B_1) + A^2(B_1 - A B_2) + \cdots + A^{n-1}(B_{n-2} - A B_{n-1}) + A^n B_{n-1} \\
&= -A B_0 + A B_0 - A^2 B_1 + A^2 B_1 - \cdots - A^n B_{n-1} + A^n B_{n-1} = 0
\end{aligned}
\] 而右边求和为
\(c_0 I_n + c_1 A + \cdots + c_{n-1} A^{n-1} + A^n = p_A(A)\)。故
\(p_A(A) = 0\)。\(\blacksquare\)

\textbf{推论 35.4.2}：\(A^n\) 可表示为 \(I_n, A, \ldots, A^{n-1}\)
的线性组合。更一般地，对任意 \(k \geq 0\)，\(A^k\) 属于
\(I_n, A, \ldots, A^{n-1}\) 张成的子空间。

\bigskip\hrule\bigskip

\subsubsection{35.5 极小多项式}\label{ux6781ux5c0fux591aux9879ux5f0f}

\textbf{定义 35.5.1（极小多项式）}：设
\(A \in M_{n \times n}(F)\)。在使得 \(q(A) = 0\) 的所有非零首一多项式
\(q \in F[t]\) 中，次数最低者称为 \(A\) 的\textbf{极小多项式}，记为
\(m_A(t)\)。

\textbf{定理 35.5.1（极小多项式的存在唯一性）}：\(m_A(t)\) 存在且唯一。

\textbf{证明}：由 Cayley-Hamilton
定理，\(p_A(A) = 0\)，故存在满足条件的非零多项式。在 \(F[t]\)
中，\(\dim M_{n \times n}(F) = n^2\)，故 \(I_n, A, \ldots, A^{n^2}\)
线性相关，因此存在次数不超过 \(n^2\)
的零化多项式。取其中次数最低的首一多项式，即为 \(m_A\)。唯一性：若
\(m_1, m_2\) 均为极小多项式，则 \(m_1 - m_2\) 次数更低，且
\((m_1 - m_2)(A) = 0\)。若
\(m_1 \neq m_2\)，则除以首项系数后得次数更低的零化多项式，矛盾。\(\blacksquare\)

\textbf{定理 35.5.2（极小多项式的基本性质）}：

\begin{enumerate}
\def\labelenumi{(\roman{enumi})}
\tightlist
\item
  \(m_A(t)\) 整除任何满足 \(q(A) = 0\) 的多项式
  \(q(t)\)。特别地，\(m_A(t) \mid p_A(t)\)。
\item
  \(A\) 的特征值恰好是 \(m_A(t)\) 在 \(F\) 的代数闭包中的根。
\end{enumerate}

\textbf{证明}：

\begin{enumerate}
\def\labelenumi{(\roman{enumi})}
\item
  由多项式带余除法，\(q(t) = m_A(t) \cdot s(t) + r(t)\)，其中
  \(\deg r < \deg m_A\)。代入
  \(A\)：\(0 = q(A) = m_A(A) s(A) + r(A) = r(A)\)。若 \(r \neq 0\)，则
  \(r\) 是比 \(m_A\) 次数更低的零化多项式，矛盾。故 \(r \equiv 0\)，即
  \(m_A \mid q\)。
\item
  若 \(\lambda\) 是 \(A\) 的特征值，\(v\) 为对应特征向量，则
  \(0 = m_A(A)v = m_A(\lambda)v\)，故 \(m_A(\lambda) = 0\)。反之，若
  \(m_A(\lambda) = 0\)，则 \(m_A(t) = (t - \lambda)q(t)\)。由于
  \(q(A) \neq 0\)（否则 \(m_A\) 不是极小），存在 \(w\) 使得
  \(q(A)w \neq \mathbf{0}\)。令 \(v = q(A)w\)，则
  \((A - \lambda I)v = (A - \lambda I)q(A)w = m_A(A)w = \mathbf{0}\)，故
  \(\lambda\) 是特征值。\(\blacksquare\)
\end{enumerate}

\textbf{定理
35.5.3（可对角化与极小多项式）}：\(A \in M_{n \times n}(F)\)
可对角化当且仅当 \(m_A(t)\) 在 \(F\)
上可分解为互不相同的一次因式的乘积。

\textbf{证明}：

（\(\Rightarrow\)）若 \(A\) 可对角化，设
\(A = P \operatorname{diag}(\lambda_1, \ldots, \lambda_n) P^{-1}\)。则
\(m_A(A) = P m_A(\operatorname{diag}(\lambda_1, \ldots, \lambda_n)) P^{-1}\)。\(m_A(\operatorname{diag}(\lambda_1, \ldots, \lambda_n)) = \operatorname{diag}(m_A(\lambda_1), \ldots, m_A(\lambda_n))\)。故
\(m_A(\lambda_i) = 0\) 对所有 \(i\) 成立。取
\(m_A(t) = \prod_{\lambda \in \{\lambda_1, \ldots, \lambda_n\}} (t - \lambda)\)（重复的只取一次），则
\(m_A(A) = 0\) 且 \(m_A\) 无重根。

（\(\Leftarrow\)）设
\(m_A(t) = \prod_{i=1}^{k} (t - \lambda_i)\)（\(\lambda_i\)
两两不同）。需证 \(F^n\) 是特征空间 \(E_{\lambda_i}\) 的直和。

对 \(k\) 归纳。\(k = 1\) 时，\(m_A(t) = t - \lambda_1\)，即
\(A = \lambda_1 I_n\)，已是对角矩阵。

\(k > 1\)
时，由部分分式分解，\(\frac{1}{m_A(t)} = \sum_{i=1}^{k} \frac{c_i}{t - \lambda_i}\)，故
\(1 = \sum_{i=1}^{k} c_i \prod_{j \neq i} (t - \lambda_j)\)。代入
\(A\)：\(I_n = \sum_{i=1}^{k} c_i \prod_{j \neq i} (A - \lambda_j I_n)\)。令
\(P_i = c_i \prod_{j \neq i} (A - \lambda_j I_n)\)，则
\(I_n = \sum_i P_i\)。

对任意 \(v \in F^n\)，\(v = \sum_i P_i v\)。由
\((A - \lambda_i I_n)P_i = c_i m_A(A) = 0\)，故
\(P_i v \in E_{\lambda_i}\)。因此 \(F^n = \sum_i E_{\lambda_i}\)。由定理
35.1.1，不同特征空间的和是直和。故 \(A\) 可对角化。\(\blacksquare\)

\bigskip\hrule\bigskip

\subsubsection{35.6 Jordan
标准形（陈述）}\label{jordan-ux6807ux51c6ux5f62ux9648ux8ff0}

\textbf{定义 35.6.1（Jordan 块）}：设 \(\lambda \in F\)。\(k \times k\)
的 \textbf{Jordan 块} \(J_k(\lambda)\) 定义为

\[
J_k(\lambda) = \begin{pmatrix}
\lambda & 1 & 0 & \cdots & 0 \\
0 & \lambda & 1 & \cdots & 0 \\
\vdots & \vdots & \ddots & \ddots & \vdots \\
0 & 0 & \cdots & \lambda & 1 \\
0 & 0 & \cdots & 0 & \lambda
\end{pmatrix}_{k \times k}
\]

\textbf{定理 35.6.1（Jordan 标准形，陈述）}：设 \(F\) 为代数闭域（如
\(\mathbb{C}\)），\(A \in M_{n \times n}(F)\)。则存在可逆矩阵
\(P \in M_{n \times n}(F)\) 使得

\[
P^{-1}AP = \begin{pmatrix}
J_{k_1}(\lambda_1) & 0 & \cdots & 0 \\
0 & J_{k_2}(\lambda_2) & \cdots & 0 \\
\vdots & \vdots & \ddots & \vdots \\
0 & 0 & \cdots & J_{k_r}(\lambda_r)
\end{pmatrix}
\]

其中 \(J_{k_i}(\lambda_i)\) 为 Jordan 块。该分解在 Jordan
块的排列顺序下是唯一的。

\textbf{注记 35.6.1}：Jordan
标准形的完整证明需要商空间和广义特征空间的理论，将在卷八（抽象代数
I）中作为模论的应用给出。此处先陈述结果，以便在后续应用中使用。

\bigskip\hrule\bigskip

\subsection{Ch 36
内积空间与谱定理}\label{ch-36-ux5185ux79efux7a7aux95f4ux4e0eux8c31ux5b9aux7406}

本章在 \(\mathbb{R}\) 上讨论。涉及 \(\mathbb{C}\)
的推广将在卷十（复分析）和卷十四（泛函分析）中处理。

\subsubsection{36.1 内积与范数}\label{ux5185ux79efux4e0eux8303ux6570}

\textbf{定义 36.1.1（内积）}：设 \(V\) 为 \(\mathbb{R}\)
上的向量空间。映射
\(\langle\cdot, \cdot\rangle : V \times V \to \mathbb{R}\)
若满足以下四条公理，则称为 \(V\) 上的\textbf{内积}：

\begin{enumerate}
\def\labelenumi{\arabic{enumi}.}
\tightlist
\item
  \textbf{对称性}：\(\langle u, v \rangle = \langle v, u \rangle\)
\item
  \textbf{第一变元的线性性}：\(\langle \lambda u + \mu v, w \rangle = \lambda \langle u, w \rangle + \mu \langle v, w \rangle\)
\item
  \textbf{正定性}：\(\langle v, v \rangle \geq 0\)，且
  \(\langle v, v \rangle = 0 \iff v = \mathbf{0}\)
\end{enumerate}

（由对称性，第二变元也是线性的。）

配备内积的向量空间称为\textbf{内积空间}。有限维实内积空间称为\textbf{Euclid
空间}。

\textbf{定理 36.1.1（Cauchy-Schwarz 不等式）}：设 \(V\) 为内积空间，则
\(\forall u, v \in V\)，

\[
|\langle u, v \rangle| \leq \sqrt{\langle u, u \rangle} \sqrt{\langle v, v \rangle}
\]

等号成立当且仅当 \(u\) 与 \(v\) 线性相关。

\textbf{证明}：若 \(v = \mathbf{0}\)，两边均为 \(0\)，不等式成立。

设 \(v \neq \mathbf{0}\)。对任意 \(t \in \mathbb{R}\)，由正定性： \[
0 \leq \langle u - tv, u - tv \rangle = \langle u, u \rangle - 2t\langle u, v \rangle + t^2 \langle v, v \rangle
\] 这是关于 \(t\) 的二次函数。其判别式
\(\Delta = 4\langle u, v \rangle^2 - 4\langle u, u \rangle \langle v, v \rangle \leq 0\)，即
\(\langle u, v \rangle^2 \leq \langle u, u \rangle \langle v, v \rangle\)。开方即得。

等号成立当且仅当 \(\Delta = 0\)，此时存在 \(t\) 使得
\(\langle u - tv, u - tv \rangle = 0\)，即 \(u = tv\)。\(\blacksquare\)

\textbf{定义 36.1.2（范数）}：由内积诱导的\textbf{范数}定义为
\(\|v\| = \sqrt{\langle v, v \rangle}\)。

\textbf{定理 36.1.2（范数的基本性质）}：

\begin{enumerate}
\def\labelenumi{(\roman{enumi})}
\tightlist
\item
  \textbf{正定性}：\(\|v\| \geq 0\)，且
  \(\|v\| = 0 \iff v = \mathbf{0}\)
\item
  \textbf{齐次性}：\(\|\lambda v\| = |\lambda| \|v\|\)
\item
  \textbf{三角不等式}：\(\|u + v\| \leq \|u\| + \|v\|\)
\end{enumerate}

\textbf{证明}：

\begin{enumerate}
\def\labelenumi{(\roman{enumi})}
\item
  由内积的正定性直接得到。
\item
  \(\|\lambda v\|^2 = \langle \lambda v, \lambda v \rangle = \lambda^2 \langle v, v \rangle = \lambda^2 \|v\|^2\)。
\item
  \(\|u + v\|^2 = \langle u+v, u+v \rangle = \|u\|^2 + 2\langle u, v \rangle + \|v\|^2 \leq \|u\|^2 + 2\|u\|\|v\| + \|v\|^2 = (\|u\| + \|v\|)^2\)。开方即得。\(\blacksquare\)
\end{enumerate}

\textbf{定义 36.1.3（标准内积）}：\(\mathbb{R}^n\)
上的\textbf{标准内积}（点积）定义为

\[
\langle x, y \rangle = \sum_{i=1}^{n} x_i y_i
\]

\bigskip\hrule\bigskip

\subsubsection{36.2 正交性}\label{ux6b63ux4ea4ux6027}

\textbf{定义 36.2.1（正交）}：设 \(V\) 为内积空间。

\begin{enumerate}
\def\labelenumi{\arabic{enumi}.}
\tightlist
\item
  若 \(\langle u, v \rangle = 0\)，则称 \(u\) 与 \(v\)
  \textbf{正交}，记为 \(u \perp v\)
\item
  向量组 \(\{v_1, \ldots, v_k\}\) 称为\textbf{正交组}，若两两正交
\item
  若正交组中每个向量的范数均为
  \(1\)，则称为\textbf{规范正交组}（或\textbf{标准正交组}）
\end{enumerate}

\textbf{定理 36.2.1（勾股定理）}：若 \(u \perp v\)，则
\(\|u + v\|^2 = \|u\|^2 + \|v\|^2\)。

\textbf{证明}：\(\|u + v\|^2 = \langle u+v, u+v \rangle = \|u\|^2 + 2\langle u, v \rangle + \|v\|^2 = \|u\|^2 + \|v\|^2\)。\(\blacksquare\)

\textbf{定理 36.2.2（正交组线性无关）}：非零向量组成的正交组线性无关。

\textbf{证明}：设 \(\{v_1, \ldots, v_k\}\)
为正交组，\(v_i \neq \mathbf{0}\)。若
\(\sum \lambda_i v_i = \mathbf{0}\)，则对任意 \(j\)， \[
0 = \langle \mathbf{0}, v_j \rangle = \left\langle \sum_i \lambda_i v_i, v_j \right\rangle = \lambda_j \langle v_j, v_j \rangle
\] 由于
\(\langle v_j, v_j \rangle > 0\)，\(\lambda_j = 0\)。\(\blacksquare\)

\textbf{推论 36.2.3}：\(n\) 维内积空间中，任何 \(n\)
个非零向量组成的正交组是基（称为\textbf{正交基}）。若进一步是规范正交的，则称为\textbf{规范正交基}。

\bigskip\hrule\bigskip

\subsubsection{36.3 Gram-Schmidt
正交化}\label{gram-schmidt-ux6b63ux4ea4ux5316}

\textbf{定理 36.3.1（Gram-Schmidt 正交化）}：设 \(\{v_1, \ldots, v_n\}\)
为内积空间 \(V\) 的基。则可构造 \(V\) 的正交基 \(\{e_1, \ldots, e_n\}\)
如下：

\[
\begin{aligned}
e_1 &= v_1 \\
e_2 &= v_2 - \frac{\langle v_2, e_1 \rangle}{\langle e_1, e_1 \rangle} e_1 \\
e_3 &= v_3 - \frac{\langle v_3, e_1 \rangle}{\langle e_1, e_1 \rangle} e_1 - \frac{\langle v_3, e_2 \rangle}{\langle e_2, e_2 \rangle} e_2 \\
&\;\;\vdots \\
e_k &= v_k - \sum_{i=1}^{k-1} \frac{\langle v_k, e_i \rangle}{\langle e_i, e_i \rangle} e_i
\end{aligned}
\]

进一步，单位化 \(u_i = \frac{e_i}{\|e_i\|}\) 得到规范正交基。

\textbf{证明}：对 \(k\) 归纳证明 \(\{e_1, \ldots, e_k\}\) 是正交组。

\textbf{基础} \(k = 1\)：单个向量自动正交。

\textbf{归纳}：设 \(\{e_1, \ldots, e_{k-1}\}\) 为正交组。对 \(j < k\)，
\[
\langle e_k, e_j \rangle = \langle v_k, e_j \rangle - \sum_{i=1}^{k-1} \frac{\langle v_k, e_i \rangle}{\langle e_i, e_i \rangle} \langle e_i, e_j \rangle = \langle v_k, e_j \rangle - \frac{\langle v_k, e_j \rangle}{\langle e_j, e_j \rangle} \langle e_j, e_j \rangle = 0
\] （因为当 \(i \neq j\) 时 \(\langle e_i, e_j \rangle = 0\)）。故
\(\{e_1, \ldots, e_k\}\) 正交。

由构造，\(\operatorname{span}\{e_1, \ldots, e_k\} = \operatorname{span}\{v_1, \ldots, v_k\}\)，故
\(e_k \neq \mathbf{0}\)（否则 \(v_k\) 属于前 \(k-1\)
个向量的张成空间，与基的线性无关性矛盾）。\(\blacksquare\)

\textbf{推论 36.3.2}：任何有限维内积空间存在规范正交基。

\textbf{推论 36.3.3（QR 分解）}：设 \(A \in M_{n \times n}(\mathbb{R})\)
可逆。则存在正交矩阵 \(Q\)（\(Q^T Q = I_n\)）和上三角矩阵
\(R\)（对角线元素为正）使得 \(A = QR\)。

\textbf{证明}：将 \(A\) 的列向量组视为 \(\mathbb{R}^n\) 的基，对其进行
Gram-Schmidt 正交化并单位化，得到规范正交基。\(Q\)
的列是规范正交化后的向量，\(R\) 记录线性组合的系数。\(\blacksquare\)

\bigskip\hrule\bigskip

\subsubsection{36.4
正交补与正交投影}\label{ux6b63ux4ea4ux8865ux4e0eux6b63ux4ea4ux6295ux5f71}

\textbf{定义 36.4.1（正交补）}：设 \(W \leq V\)。\(W\)
的\textbf{正交补}定义为

\[
W^{\perp} = \{v \in V \mid \langle v, w \rangle = 0, \; \forall w \in W\}
\]

\textbf{定理 36.4.1}：\(W^{\perp}\) 是 \(V\) 的子空间，且
\(V = W \oplus W^{\perp}\)。

\textbf{证明}：\(W^{\perp}\) 是子空间：由内积的线性性直接验证。

取 \(W\) 的规范正交基 \(\{e_1, \ldots, e_k\}\)，将其扩充为 \(V\)
的规范正交基 \(\{e_1, \ldots, e_k, e_{k+1}, \ldots, e_n\}\)（通过
Gram-Schmidt 正交化）。则 \(\{e_{k+1}, \ldots, e_n\}\) 是 \(W^{\perp}\)
的基。故 \(V = W \oplus W^{\perp}\)。\(\blacksquare\)

\textbf{推论 36.4.2}：\(\dim W^{\perp} = \dim V - \dim W\)，且
\((W^{\perp})^{\perp} = W\)。

\textbf{定义 36.4.2（正交投影）}：由 \(V = W \oplus W^{\perp}\)，对任意
\(v \in V\)，存在唯一的分解
\(v = w + w^{\perp}\)（\(w \in W\)，\(w^{\perp} \in W^{\perp}\)）。映射
\(P_W : V \to V\)，\(P_W(v) = w\) 称为 \(V\) 到 \(W\)
的\textbf{正交投影}。

\textbf{定理 36.4.2（正交投影的性质）}：

\begin{enumerate}
\def\labelenumi{(\roman{enumi})}
\tightlist
\item
  \(P_W\) 是线性算子，且 \(P_W^2 = P_W\)（幂等）
\item
  \(\operatorname{im} P_W = W\)，\(\ker P_W = W^{\perp}\)
\item
  \(P_W\)
  是自伴的：\(\langle P_W u, v \rangle = \langle u, P_W v \rangle\)
\item
  若 \(\{e_1, \ldots, e_k\}\) 为 \(W\) 的规范正交基，则
  \(P_W(v) = \sum_{i=1}^{k} \langle v, e_i \rangle e_i\)
\end{enumerate}

\textbf{证明}：(i) 线性性由分解的唯一性验证。\(P_W^2 = P_W\) 因为
\(P_W(w) = w\)（\(w \in W\)）。

\begin{enumerate}
\def\labelenumi{(\roman{enumi})}
\setcounter{enumi}{1}
\item
  显然。(iii) 设 \(u = u_W + u_{\perp}\)，\(v = v_W + v_{\perp}\)，则
  \(\langle P_W u, v \rangle = \langle u_W, v_W + v_{\perp} \rangle = \langle u_W, v_W \rangle = \langle u_W + u_{\perp}, v_W \rangle = \langle u, P_W v \rangle\)。
\item
  将 \(\{e_1, \ldots, e_k\}\) 扩充为 \(V\) 的规范正交基
  \(\{e_1, \ldots, e_n\}\)。设
  \(v = \sum_{i=1}^{n} \alpha_i e_i\)（\(\alpha_i = \langle v, e_i \rangle\)）。则
  \(P_W(v) = \sum_{i=1}^{k} \alpha_i e_i = \sum_{i=1}^{k} \langle v, e_i \rangle e_i\)。\(\blacksquare\)
\end{enumerate}

\bigskip\hrule\bigskip

\subsubsection{36.5 伴随算子}\label{ux4f34ux968fux7b97ux5b50}

\textbf{定义 36.5.1（伴随算子）}：设 \(V\)
为有限维内积空间，\(T \in \mathcal{L}(V)\)。\(T\) 的\textbf{伴随算子}
\(T^* \in \mathcal{L}(V)\) 由下式唯一确定：

\[
\langle T(u), v \rangle = \langle u, T^*(v) \rangle, \quad \forall u, v \in V
\]

\textbf{定理 36.5.1（伴随算子的存在唯一性）}：\(T^*\) 存在且唯一。

\textbf{证明}：固定 \(v \in V\)，考虑线性泛函
\(f_v(u) = \langle T(u), v \rangle\)。由 Riesz
表示定理（在有限维情形下，\(f_v\) 可表示为与某个向量的内积），存在唯一的
\(w \in V\) 使得 \(f_v(u) = \langle u, w \rangle\)。定义
\(T^*(v) = w\)。可验证 \(T^*\) 是线性的。\(\blacksquare\)

\textbf{定理 36.5.2（伴随算子的性质）}：

\begin{enumerate}
\def\labelenumi{(\roman{enumi})}
\tightlist
\item
  \((T^*)^* = T\)
\item
  \((S + T)^* = S^* + T^*\)
\item
  \((\lambda T)^* = \lambda T^*\)
\item
  \((ST)^* = T^* S^*\)
\item
  若 \(T\) 可逆，则 \((T^{-1})^* = (T^*)^{-1}\)
\item
  在规范正交基下，\([T^*] = [T]^T\)（伴随算子对应矩阵的转置）
\end{enumerate}

\textbf{证明}：(i)-(v) 由伴随算子的定义和内积的性质直接验证。

\begin{enumerate}
\def\labelenumi{(\roman{enumi})}
\setcounter{enumi}{5}
\tightlist
\item
  设 \(B = \{e_1, \ldots, e_n\}\) 为 \(V\)
  的规范正交基，\(A = [T]_B\)，\(B = [T^*]_B\)。则
  \(T(e_j) = \sum_i a_{ij} e_i\)，\(T^*(e_j) = \sum_i b_{ij} e_i\)。
\end{enumerate}

由伴随算子的定义：\(\langle T(e_j), e_i \rangle = \langle e_j, T^*(e_i) \rangle\)。左边
\(= \langle \sum_k a_{kj} e_k, e_i \rangle = a_{ij}\)。右边
\(= \langle e_j, \sum_k b_{ki} e_k \rangle = b_{ji}\)。故
\(a_{ij} = b_{ji}\)，即 \(B = A^T\)。\(\blacksquare\)

\textbf{定义 36.5.2（自伴算子与正交算子）}：

\begin{enumerate}
\def\labelenumi{\arabic{enumi}.}
\tightlist
\item
  若 \(T^* = T\)，则称 \(T\)
  为\textbf{自伴算子}（或\textbf{对称算子}）。在规范正交基下，其矩阵为对称矩阵（\(A^T = A\)）。
\item
  若 \(T^* T = T T^* = I\)，则称 \(T\)
  为\textbf{正交算子}。在规范正交基下，其矩阵为正交矩阵（\(A^T A = I_n\)）。
\end{enumerate}

\bigskip\hrule\bigskip

\subsubsection{36.6 谱定理}\label{ux8c31ux5b9aux7406}

\textbf{定理 36.6.1（实对称矩阵的谱定理）}：设
\(A \in M_{n \times n}(\mathbb{R})\) 为对称矩阵（\(A^T = A\)）。则

\begin{enumerate}
\def\labelenumi{(\roman{enumi})}
\tightlist
\item
  \(A\) 的所有特征值均为实数。
\item
  \(A\) 可对角化，且存在正交矩阵 \(Q\)（\(Q^T Q = I_n\)）使得
\end{enumerate}

\[
Q^T A Q = \operatorname{diag}(\lambda_1, \lambda_2, \ldots, \lambda_n)
\]

其中 \(\lambda_1, \ldots, \lambda_n\) 为 \(A\)
的特征值。等价地，\(\mathbb{R}^n\) 有一组由 \(A\)
的特征向量组成的规范正交基。

\textbf{证明}：对 \(n\) 归纳。

\textbf{基础} \(n = 1\)：平凡。

\textbf{归纳}：假设结论对 \(n-1\) 阶对称矩阵成立。

\textbf{第 1 步}：\(A\) 有实特征值。考虑 \(\mathbb{C}^n\) 上的标准
Hermite 内积。\(A\) 作为复矩阵是 Hermite
矩阵（\(A^* = A^T = A\)）。\(A\) 至少有一个特征值
\(\lambda_1 \in \mathbb{C}\) 和对应的特征向量
\(v_1 \neq \mathbf{0}\)。由
\(\lambda_1 \|v_1\|^2 = \langle A v_1, v_1 \rangle = \langle v_1, A v_1 \rangle = \overline{\lambda_1} \|v_1\|^2\)，得
\(\lambda_1 = \overline{\lambda_1}\)，故
\(\lambda_1 \in \mathbb{R}\)。可取 \(v_1\) 为实向量，并单位化使
\(\|v_1\| = 1\)。

\textbf{第 2 步}：设 \(W = \operatorname{span}\{v_1\}\)。则
\(W^{\perp}\) 是 \(A\)-不变的：对任意
\(w \in W^{\perp}\)，\(\langle A w, v_1 \rangle = \langle w, A v_1 \rangle = \langle w, \lambda_1 v_1 \rangle = \lambda_1 \langle w, v_1 \rangle = 0\)，故
\(A w \in W^{\perp}\)。

\textbf{第 3 步}：\(A|_{W^{\perp}}\) 是 \(W^{\perp}\) 上的对称算子（因为
\(W^{\perp}\) 是 \(A\)-不变的，且 \(A\) 的对称性继承）。取 \(W^{\perp}\)
的规范正交基，\(A|_{W^{\perp}}\) 在此基下的矩阵 \(A'\) 为
\((n-1) \times (n-1)\) 对称矩阵。由归纳假设，存在 \((n-1) \times (n-1)\)
正交矩阵 \(Q'\) 使得
\((Q')^T A' Q' = \operatorname{diag}(\lambda_2, \ldots, \lambda_n)\)。

将 \(v_1\) 与 \(W^{\perp}\) 的规范正交基合并，得到 \(\mathbb{R}^n\)
的规范正交基。对应的正交矩阵 \(Q\) 满足
\(Q^T A Q = \operatorname{diag}(\lambda_1, \ldots, \lambda_n)\)。\(\blacksquare\)

\textbf{定理 36.6.2（谱定理的算子形式）}：设 \(V\)
为有限维实内积空间，\(T \in \mathcal{L}(V)\) 为自伴算子。则 \(V\)
存在一组由 \(T\) 的特征向量组成的规范正交基。等价地，\(T\)
可对角化且特征向量可选取为正交的。

\textbf{证明}：取 \(V\) 的任意规范正交基，\(T\) 在此基下的矩阵 \(A\)
是对称矩阵。由定理 36.6.1，\(A\) 可正交对角化。\(Q\) 的列向量即为 \(T\)
的特征向量，且构成规范正交基。\(\blacksquare\)

\textbf{定理 36.6.3（奇异值分解，陈述）}：设
\(A \in M_{m \times n}(\mathbb{R})\)。则存在正交矩阵
\(U \in M_{m \times m}(\mathbb{R})\)，\(V \in M_{n \times n}(\mathbb{R})\)
和对角矩阵
\(\Sigma = \operatorname{diag}(\sigma_1, \ldots, \sigma_r, 0, \ldots, 0)\)（\(\sigma_1 \geq \cdots \geq \sigma_r > 0\)）使得

\[
A = U \Sigma V^T
\]

其中 \(\sigma_i\) 称为 \(A\) 的\textbf{奇异值}，等于 \(A^T A\)
的正特征值的平方根。

\textbf{评注 36.6.1}：奇异值分解的证明基于谱定理应用于
\(A^T A\)，详细推导从略。\(\blacksquare\)

\bigskip\hrule\bigskip

\subsection{Ch 37
二次型与标准形}\label{ch-37-ux4e8cux6b21ux578bux4e0eux6807ux51c6ux5f62}

\subsubsection{37.1 二次型}\label{ux4e8cux6b21ux578b}

\textbf{定义 37.1.1（二次型）}：设 \(V\) 为 \(\mathbb{R}\) 上的 \(n\)
维向量空间。函数 \(Q : V \to \mathbb{R}\)
称为\textbf{二次型}，若存在对称双线性型
\(B : V \times V \to \mathbb{R}\)（即 \(B(u, v) = B(v, u)\)
且对每个变元线性）使得

\[
Q(v) = B(v, v), \quad \forall v \in V
\]

\textbf{定义 37.1.2（二次型的矩阵表示）}：在基
\(B = \{e_1, \ldots, e_n\}\) 下，设 \(v = \sum x_i e_i\)。则

\[
Q(v) = \sum_{i,j=1}^{n} a_{ij} x_i x_j = x^T A x
\]

其中 \(A = (a_{ij})\) 为对称矩阵（\(a_{ij} = B(e_i, e_j)\)），称为二次型
\(Q\) 关于基 \(B\) 的\textbf{矩阵}。

\textbf{定理 37.1.1（基变换下二次型矩阵的变化）}：设 \(P\) 为从基 \(B\)
到 \(B'\) 的基变换矩阵。则 \(Q\) 关于 \(B'\) 的矩阵 \(A'\) 满足

\[
A' = P^T A P
\]

\textbf{证明}：设 \(x\) 为关于 \(B\) 的坐标，\(x'\) 为关于 \(B'\)
的坐标，则 \(x = P x'\)。故
\(Q(v) = x^T A x = (P x')^T A (P x') = (x')^T (P^T A P) x'\)。\(\blacksquare\)

\textbf{定义 37.1.3（合同）}：对称矩阵
\(A, B \in M_{n \times n}(\mathbb{R})\)
称为\textbf{合同}的，若存在可逆矩阵 \(P\) 使得 \(B = P^T A P\)。

\bigskip\hrule\bigskip

\subsubsection{37.2
合同变换与标准形}\label{ux5408ux540cux53d8ux6362ux4e0eux6807ux51c6ux5f62}

\textbf{定理 37.2.1（二次型的标准形）}：设 \(Q\) 为 \(\mathbb{R}^n\)
上的二次型，矩阵为 \(A\)（对称）。则存在可逆矩阵 \(P\)，使得经坐标变换
\(x = P y\) 后，

\[
Q(v) = \lambda_1 y_1^2 + \lambda_2 y_2^2 + \cdots + \lambda_n y_n^2
\]

其中
\(\lambda_i \in \mathbb{R}\)。这等价于任何实对称矩阵合同于对角矩阵。

\textbf{证明}：由谱定理（定理 36.6.1），存在正交矩阵 \(Q\) 使得
\(Q^T A Q = \operatorname{diag}(\lambda_1, \ldots, \lambda_n)\)。取
\(P = Q\)，则
\(P^T A P = \operatorname{diag}(\lambda_1, \ldots, \lambda_n)\)。坐标变换
\(x = P y\) 给出 \(Q(v) = \sum \lambda_i y_i^2\)。\(\blacksquare\)

\textbf{注记 37.2.1}：上述证明利用了正交对角化，\(\lambda_i\) 恰好是
\(A\) 的特征值。也可以用配方法（Lagrange 配方法）直接构造，无需谱定理。

\bigskip\hrule\bigskip

\subsubsection{37.3 惯性定理}\label{ux60efux6027ux5b9aux7406}

\textbf{定理 37.3.1（Sylvester 惯性定理）}：设实二次型 \(Q\)
在两组基下分别化为标准形：

\[
Q(v) = \lambda_1 y_1^2 + \cdots + \lambda_p y_p^2 - \lambda_{p+1} y_{p+1}^2 - \cdots - \lambda_{p+q} y_{p+q}^2 + 0 \cdot y_{p+q+1}^2 + \cdots + 0 \cdot y_n^2
\]

\[
Q(v) = \mu_1 z_1^2 + \cdots + \mu_{p'} z_{p'}^2 - \mu_{p'+1} z_{p'+1}^2 - \cdots - \mu_{p'+q'} z_{p'+q'}^2 + 0
\]

其中 \(\lambda_i, \mu_j > 0\)。则 \(p = p'\)，\(q = q'\)。

\textbf{证明}：\(p\) 是 \(Q\) 在其上取正值的子空间的最大维数，\(q\) 是
\(Q\) 在其上取负值的子空间的最大维数。这些是不依赖于基选择的不变量。

具体地，设 \(V_+\) 为 \(Q(v) \geq 0\) 的子空间中的最大维数者（\(Q\)
限制在 \(V_+\) 上正定）。可以证明 \(\dim V_+ = p\)。同样，\(Q\) 限制在
\(V_-\) 上负定的最大维数为 \(q\)。由于 \(p, q\) 由 \(Q\)
本身决定，它们在两组基下相同。\(\blacksquare\)

\textbf{定义 37.3.1（惯性指数与符号差）}：

\begin{itemize}
\tightlist
\item
  \(p\) 称为 \(Q\) 的\textbf{正惯性指数}
\item
  \(q\) 称为 \(Q\) 的\textbf{负惯性指数}
\item
  \(p + q = \operatorname{rank} A\) 称为 \(Q\) 的\textbf{秩}
\item
  \(p - q\) 称为 \(Q\) 的\textbf{符号差}
\end{itemize}

\textbf{推论
37.3.2}：两个实二次型等价（可通过可逆线性变换互相转化）当且仅当它们具有相同的正惯性指数和负惯性指数。

\bigskip\hrule\bigskip

\subsubsection{37.4 Sylvester 准则}\label{sylvester-ux51c6ux5219}

\textbf{定义 37.4.1（顺序主子式）}：设
\(A = (a_{ij}) \in M_{n \times n}(\mathbb{R})\) 为对称矩阵。对
\(k = 1, \ldots, n\)，\(A\) 的第 \(k\) 个\textbf{顺序主子式}为

\[
\Delta_k = \det \begin{pmatrix}
a_{11} & \cdots & a_{1k} \\
\vdots & \ddots & \vdots \\
a_{k1} & \cdots & a_{kk}
\end{pmatrix}
\]

\textbf{定理 37.4.1（Sylvester 准则）}：实对称矩阵 \(A\) 正定（即
\(x^T A x > 0\) 对所有 \(x \neq \mathbf{0}\)）当且仅当其所有顺序主子式
\(\Delta_k > 0\)（\(k = 1, \ldots, n\)）。

\textbf{证明}：

（\(\Rightarrow\)）设 \(A\) 正定。对任意 \(k\)，考虑 \(A\) 的左上
\(k \times k\) 子矩阵 \(A_k\)。对任意非零 \(y \in \mathbb{R}^k\)，令
\(x = (y, 0, \ldots, 0)^T \in \mathbb{R}^n\)，则
\(y^T A_k y = x^T A x > 0\)，故 \(A_k\) 正定。正定矩阵的特征值全正，故
\(\det A_k = \prod_{i=1}^{k} \lambda_i > 0\)。

（\(\Leftarrow\)）对 \(n\) 归纳。\(n = 1\)
时，\(\Delta_1 = a_{11} > 0\)，故 \(A\) 正定。

假设结论对 \(n-1\) 成立。设
\(\Delta_k > 0\)（\(k = 1, \ldots, n\)）。由归纳假设，\(A_{n-1}\)
正定，故存在可逆矩阵 \(P_{n-1}\) 使得
\(P_{n-1}^T A_{n-1} P_{n-1} = I_{n-1}\)。

将 \(A\) 写为分块形式
\(A = \begin{pmatrix} A_{n-1} & b \\ b^T & a_{nn} \end{pmatrix}\)。令
\(P = \begin{pmatrix} P_{n-1} & -A_{n-1}^{-1} b \\ 0 & 1 \end{pmatrix}\)，则
\[
P^T A P = \begin{pmatrix} I_{n-1} & 0 \\ 0 & a_{nn} - b^T A_{n-1}^{-1} b \end{pmatrix}
\] 其行列式为 \(\Delta_n / \Delta_{n-1} > 0\)（因为
\(\det(P^T A P) = \det P^2 \det A\)）。故
\(a_{nn} - b^T A_{n-1}^{-1} b > 0\)，从而 \(P^T A P\) 正定，\(A\)
正定。\(\blacksquare\)

\bigskip\hrule\bigskip

\subsubsection{37.5 正定二次型}\label{ux6b63ux5b9aux4e8cux6b21ux578b}

\textbf{定义 37.5.1（二次型的定性）}：设 \(Q\) 为实二次型，矩阵为
\(A\)。

\begin{enumerate}
\def\labelenumi{\arabic{enumi}.}
\tightlist
\item
  若 \(Q(v) > 0\)（\(\forall v \neq \mathbf{0}\)），称 \(Q\)（及
  \(A\)）\textbf{正定}
\item
  若 \(Q(v) \geq 0\)（\(\forall v\)），称 \(Q\)（及
  \(A\)）\textbf{半正定}
\item
  若 \(Q(v) < 0\)（\(\forall v \neq \mathbf{0}\)），称 \(Q\)（及
  \(A\)）\textbf{负定}
\item
  若 \(Q(v) \leq 0\)（\(\forall v\)），称 \(Q\)（及
  \(A\)）\textbf{半负定}
\item
  若 \(Q\) 既取正值又取负值，称 \(Q\)（及 \(A\)）\textbf{不定}
\end{enumerate}

\textbf{定理 37.5.1（正定性的特征值判定）}：对称矩阵 \(A\) 正定当且仅当
\(A\) 的所有特征值均为正数。

\textbf{证明}：由谱定理，\(A = Q \operatorname{diag}(\lambda_1, \ldots, \lambda_n) Q^T\)。对任意
\(x \neq \mathbf{0}\)，令 \(y = Q^T x \neq \mathbf{0}\)，则
\(x^T A x = y^T \operatorname{diag}(\lambda_i) y = \sum \lambda_i y_i^2\)。此式对所有
\(y \neq \mathbf{0}\) 为正当且仅当所有
\(\lambda_i > 0\)。\(\blacksquare\)

\textbf{定理 37.5.2（二次型与内积）}：\(V\) 上的正定二次型 \(Q\)
通过极化恒等式

\[
\langle u, v \rangle_Q = \frac{1}{2}\left(Q(u + v) - Q(u) - Q(v)\right)
\]

定义了 \(V\) 上的一个内积。反之，任何内积 \(\langle\cdot,\cdot\rangle\)
通过 \(Q(v) = \langle v, v \rangle\) 给出一个正定二次型。

\textbf{证明}：直接验证内积的四条公理。正定性由 \(Q\)
的正定性保证，对称性由二次型的定义中 \(B\)
的对称性保证，线性性由极化恒等式的构造保证。\(\blacksquare\)

\bigskip\hrule\bigskip

\subsection{Ch 38 张量初步}\label{ch-38-ux5f20ux91cfux521dux6b65}

\subsubsection{38.1
多重线性映射}\label{ux591aux91cdux7ebfux6027ux6620ux5c04}

\textbf{定义 38.1.1（多重线性映射）}：设 \(V_1, \ldots, V_k, W\) 为
\(F\) 上的向量空间。映射 \(T : V_1 \times \cdots \times V_k \to W\)
称为\textbf{多重线性映射}（\(k\)-线性映射），若 \(T\)
对每个变元是线性的（固定其他变元时）。

\textbf{定义 38.1.2（对偶空间）}：设 \(V\) 为 \(F\) 上的向量空间。\(V\)
上的所有线性泛函（即 \(V \to F\) 的线性映射）构成的向量空间称为 \(V\)
的\textbf{对偶空间}，记为 \(V^* = \mathcal{L}(V, F)\)。

\textbf{定理 38.1.1}：若 \(\dim V = n\)，则 \(\dim V^* = n\)。若
\(B = \{e_1, \ldots, e_n\}\) 为 \(V\) 的基，则 \(V^*\)
有一组\textbf{对偶基} \(\{e^1, \ldots, e^n\}\)，满足
\(e^i(e_j) = \delta_{ij}\)。

\textbf{证明}：\(\dim V^* = \dim \mathcal{L}(V, F) = (\dim V)(\dim F) = n\)。定义
\(e^i(\sum_j x_j e_j) = x_i\)，则
\(e^i(e_j) = \delta_{ij}\)。\(\{e^1, \ldots, e^n\}\) 线性无关（若
\(\sum \lambda_i e^i = 0\)，作用在 \(e_j\) 上得
\(\lambda_j = 0\)），且恰有 \(n\) 个，故为基。\(\blacksquare\)

\bigskip\hrule\bigskip

\subsubsection{38.2 张量积}\label{ux5f20ux91cfux79ef}

\textbf{定义 38.2.1（张量积的泛性质定义）}：设 \(V, W\) 为 \(F\)
上的向量空间。\(V\) 与 \(W\) 的\textbf{张量积}是一个向量空间
\(V \otimes W\) 加上一个双线性映射
\(\otimes : V \times W \to V \otimes W\)（\((v, w) \mapsto v \otimes w\)），满足以下\textbf{泛性质}：

对任意双线性映射 \(B : V \times W \to U\)（\(U\)
为任意向量空间），存在唯一的线性映射 \(\tilde{B} : V \otimes W \to U\)
使得下图交换：

\[
B = \tilde{B} \circ \otimes
\]

即 \(B(v, w) = \tilde{B}(v \otimes w)\)。

\textbf{定理 38.2.1（张量积的存在性）}：对任何有限维向量空间
\(V, W\)，张量积 \(V \otimes W\) 存在，且在同构意义下唯一。

\textbf{证明}（构造）：设 \(\dim V = m\)，\(\dim W = n\)。取 \(V\) 的基
\(\{e_1, \ldots, e_m\}\)，\(W\) 的基 \(\{f_1, \ldots, f_n\}\)。定义
\(V \otimes W\) 为以
\(\{e_i \otimes f_j \mid 1 \leq i \leq m, 1 \leq j \leq n\}\) 为基的
\(mn\) 维向量空间。对 \(v = \sum a_i e_i\)，\(w = \sum b_j f_j\)，定义
\(v \otimes w = \sum_{i,j} a_i b_j (e_i \otimes f_j)\)。

验证泛性质：设 \(B : V \times W \to U\) 为双线性映射。定义
\(\tilde{B}(e_i \otimes f_j) = B(e_i, f_j)\)，并线性扩充到
\(V \otimes W\)。则
\(\tilde{B}(v \otimes w) = \tilde{B}(\sum a_i b_j e_i \otimes f_j) = \sum a_i b_j B(e_i, f_j) = B(\sum a_i e_i, \sum b_j f_j) = B(v, w)\)。\(\tilde{B}\)
的唯一性由基上的值唯一确定。\(\blacksquare\)

\textbf{定理 38.2.2（张量积的基本性质）}：设 \(V, W, U\)
为有限维向量空间。

\begin{enumerate}
\def\labelenumi{(\roman{enumi})}
\tightlist
\item
  \(\dim(V \otimes W) = (\dim V)(\dim W)\)
\item
  \(V \otimes W \cong W \otimes V\)（通过
  \(v \otimes w \mapsto w \otimes v\)）
\item
  \((V \otimes W) \otimes U \cong V \otimes (W \otimes U)\)
\item
  \(V \otimes (W \oplus U) \cong (V \otimes W) \oplus (V \otimes U)\)
\item
  \(F \otimes V \cong V\)
\item
  \((V \otimes W)^* \cong V^* \otimes W^*\)（有限维情形）
\end{enumerate}

\textbf{证明}：(i) 由构造中基的个数 \(mn\) 得到。

\begin{enumerate}
\def\labelenumi{(\roman{enumi})}
\setcounter{enumi}{1}
\item
  定义映射 \(\tau : V \otimes W \to W \otimes V\) 为
  \(\tau(v \otimes w) = w \otimes v\)，线性扩充。其逆为
  \(\tau^{-1}(w \otimes v) = v \otimes w\)。
\item
  定义
  \((v \otimes w) \otimes u \mapsto v \otimes (w \otimes u)\)，线性扩充。
\item
  由泛性质验证。
\item
  定义 \(\lambda \otimes v \mapsto \lambda v\)，其逆为
  \(v \mapsto 1 \otimes v\)。
\item
  取基后验证维数相同，且存在自然同构
  \(\alpha \otimes \beta \mapsto (v \otimes w \mapsto \alpha(v)\beta(w))\)。\(\blacksquare\)
\end{enumerate}

\textbf{定义 38.2.2（张量）}：\(k\) 个 \(V\) 的张量积
\(V^{\otimes k} = V \otimes V \otimes \cdots \otimes V\)（\(k\)
次）中的元素称为 \(V\) 上的 \(k\) 阶\textbf{张量}。特别地：

\begin{itemize}
\tightlist
\item
  \(V^{\otimes 0} = F\)（标量）
\item
  \(V^{\otimes 1} = V\)（向量）
\item
  \(V^{\otimes 2} = V \otimes V\)（二阶张量）
\end{itemize}

\bigskip\hrule\bigskip

\subsubsection{38.3 外代数}\label{ux5916ux4ee3ux6570}

\textbf{定义 38.3.1（外幂）}：设 \(V\) 为 \(F\) 上的向量空间。\(V\) 的第
\(k\) 次\textbf{外幂} \(\bigwedge^k V\) 是 \(V^{\otimes k}\)
的商空间，模去所有形如 \(v_1 \otimes \cdots \otimes v_k\) 且
\(v_i = v_j\)（\(i \neq j\)）的元素生成的子空间。\(\bigwedge^k V\)
中的元素记为 \(v_1 \wedge v_2 \wedge \cdots \wedge v_k\)，称为
\(k\)-向量。

\textbf{定理 38.3.1（外幂的基本性质）}：

\begin{enumerate}
\def\labelenumi{(\roman{enumi})}
\tightlist
\item
  \textbf{交错性}：\(v_1 \wedge \cdots \wedge v_i \wedge \cdots \wedge v_j \wedge \cdots \wedge v_k = -v_1 \wedge \cdots \wedge v_j \wedge \cdots \wedge v_i \wedge \cdots \wedge v_k\)
\item
  若 \(\{e_1, \ldots, e_n\}\) 为 \(V\) 的基，则
  \(\{e_{i_1} \wedge e_{i_2} \wedge \cdots \wedge e_{i_k} \mid 1 \leq i_1 < i_2 < \cdots < i_k \leq n\}\)
  为 \(\bigwedge^k V\) 的基
\item
  \(\dim \bigwedge^k V = \binom{n}{k}\)（\(0 \leq k \leq n\)），\(\bigwedge^k V = \{0\}\)（\(k > n\)）
\item
  \(\bigwedge^n V\) 是一维的，基为 \(e_1 \wedge \cdots \wedge e_n\)
\end{enumerate}

\textbf{证明}：(i) 由定义，\(v_i \wedge v_j = -v_j \wedge v_i\)（因为
\((v_i + v_j) \wedge (v_i + v_j) = 0\) 展开得
\(v_i \wedge v_i + v_i \wedge v_j + v_j \wedge v_i + v_j \wedge v_j = v_i \wedge v_j + v_j \wedge v_i = 0\)）。

\begin{enumerate}
\def\labelenumi{(\roman{enumi})}
\setcounter{enumi}{1}
\item
  由交错性，任何
  \(k\)-向量可写为严格递增指标对应的基向量的外积的线性组合。线性无关性由构造保证。
\item
  由 (ii)，基的个数为 \(\binom{n}{k}\)。
\item
  \(k = n\) 时 \(\binom{n}{n} = 1\)。\(\blacksquare\)
\end{enumerate}

\textbf{定义 38.3.2（外代数）}：\(V\) 的\textbf{外代数}（Grassmann
代数）为直和

\[
\bigwedge V = \bigoplus_{k=0}^{n} \bigwedge^k V
\]

外积 \(\wedge\) 使 \(\bigwedge V\) 成为一个结合代数（满足
\(a \wedge b = (-1)^{(\deg a)(\deg b)} b \wedge a\)）。

\textbf{定理 38.3.2（行列式与外积的关系）}：设
\(v_1, \ldots, v_n \in \mathbb{R}^n\)。则

\[
v_1 \wedge v_2 \wedge \cdots \wedge v_n = \det(v_1, \ldots, v_n) \, e_1 \wedge e_2 \wedge \cdots \wedge e_n
\]

\textbf{证明}：将 \(v_j = \sum_i a_{ij} e_i\)
代入外积，使用交错性展开。外积的系数正是 Leibniz 公式定义的
\(\det(a_{ij})\)。\(\blacksquare\)

\bigskip\hrule\bigskip

\subsubsection{38.4 对称张量}\label{ux5bf9ux79f0ux5f20ux91cf}

\textbf{定义 38.4.1（对称张量）}：\(T \in V^{\otimes k}\)
称为\textbf{对称张量}，若对任意置换 \(\sigma \in S_k\)，有
\(\sigma \cdot T = T\)（其中 \(\sigma\) 通过置换因子作用在
\(V^{\otimes k}\) 上）。

\textbf{定义 38.4.2（对称幂）}：\(V\) 的第 \(k\) 次\textbf{对称幂}
\(S^k V\) 是 \(V^{\otimes k}\) 的商空间，模去形如
\(v_1 \otimes \cdots \otimes v_k - v_{\sigma(1)} \otimes \cdots \otimes v_{\sigma(k)}\)
的元素生成的子空间。\(S^k V\) 中的元素记为
\(v_1 \odot v_2 \odot \cdots \odot v_k\)。

\textbf{定理 38.4.1}：若 \(\{e_1, \ldots, e_n\}\) 为 \(V\) 的基，则
\(\{e_{i_1} \odot e_{i_2} \odot \cdots \odot e_{i_k} \mid 1 \leq i_1 \leq i_2 \leq \cdots \leq i_k \leq n\}\)
为 \(S^k V\) 的基。\(\dim S^k V = \binom{n + k - 1}{k}\)。

\textbf{证明}：由对称性，任何对称张量可写为非降指标对应的基向量的对称积的线性组合。基的个数为从
\(n\) 个元素中可重复地选取 \(k\) 个的组合数，即
\(\binom{n+k-1}{k}\)。\(\blacksquare\)

\textbf{定义 38.4.3（多重线性型与张量）}：\(V\) 上的 \(k\)
重线性型自然地对应于 \((V^*)^{\otimes k}\)
中的元素，或等价地，\(V^{\otimes k}\) 上的线性泛函。特别地，\(V\)
上的对称双线性型对应于 \(S^2(V^*)\) 中的元素，这正是二次型的理论框架。

\bigskip\hrule\bigskip

\textbf{注记
38.4.1}：张量与外代数的深入理论（包括流形上的张量场、微分形式、de Rham
上同调等）将在卷十二（微分几何）和卷十五（代数拓扑）中详细展开。本章仅建立有限维向量空间上的基本代数框架。\(\blacksquare\)

\newpage

\section{卷七：抽象代数 I}\label{ux5377ux4e03ux62bdux8c61ux4ee3ux6570-i}

\bigskip\hrule\bigskip

\begin{quote}
\textbf{前置知识}：本卷建立在卷一（数系、整除理论）和卷二（代数结构：群、环、域的定义）及卷六（线性代数）的基础之上。读者应熟悉群/环/域的基本公理、同态/同构的概念、以及向量空间的基本理论。
\end{quote}

\subsection{Ch 39 群论基础}\label{ch-39-ux7fa4ux8bbaux57faux7840}

\subsubsection{39.1
群的基本概念回顾}\label{ux7fa4ux7684ux57faux672cux6982ux5ff5ux56deux987e}

卷二第六章已引入群的基本定义。本节简要回顾并深化。

\textbf{定义 39.1.1（群）}：群是一个集合 \(G\) 配备二元运算
\(\cdot\)，满足：

\begin{enumerate}
\def\labelenumi{\arabic{enumi}.}
\tightlist
\item
  \textbf{封闭性}：\(\forall a, b \in G\)，\(a \cdot b \in G\)
\item
  \textbf{结合律}：\((a \cdot b) \cdot c = a \cdot (b \cdot c)\)
\item
  \textbf{幺元}：\(\exists e \in G\)，\(\forall a \in G\)，\(e \cdot a = a \cdot e = a\)
\item
  \textbf{逆元}：\(\forall a \in G\)，\(\exists a^{-1} \in G\)，\(a \cdot a^{-1} = a^{-1} \cdot a = e\)
\end{enumerate}

若进一步满足交换律，则称为\textbf{阿贝尔群}。

\textbf{定义 39.1.2（子群）}：\(H \subseteq G\) 称为 \(G\)
的\textbf{子群}，记为 \(H \leq G\)，若 \(H\) 在 \(G\)
的运算下自身构成群。

\textbf{定理 39.1.1（子群判定）}：\(H \subseteq G\) 是子群当且仅当

\begin{enumerate}
\def\labelenumi{(\roman{enumi})}
\tightlist
\item
  \(H \neq \varnothing\)
\item
  \(\forall a, b \in H\)，\(ab^{-1} \in H\)
\end{enumerate}

\textbf{证明}：由群的定义直接验证。取 \(a \in H\)，则
\(e = aa^{-1} \in H\)。对 \(a \in H\)，\(a^{-1} = e a^{-1} \in H\)。对
\(a, b \in H\)，\(ab = a(b^{-1})^{-1} \in H\)。\(\blacksquare\)

\textbf{定义 39.1.3（阶）}：群 \(G\) 的元素个数称为 \(G\)
的\textbf{阶}，记为 \(|G|\)。元素 \(a \in G\) 的\textbf{阶}是使得
\(a^n = e\) 的最小正整数 \(n\)（若不存在则为 \(\infty\)），记为
\(|a|\)。

\bigskip\hrule\bigskip

\subsubsection{39.2 陪集与 Lagrange
定理}\label{ux966aux96c6ux4e0e-lagrange-ux5b9aux7406}

\textbf{定义 39.2.1（陪集）}：设 \(H \leq G\)，\(a \in G\)。

\begin{itemize}
\tightlist
\item
  \(aH = \{ah \mid h \in H\}\) 称为 \(H\) 的\textbf{左陪集}
\item
  \(Ha = \{ha \mid h \in H\}\) 称为 \(H\) 的\textbf{右陪集}
\end{itemize}

\textbf{定理 39.2.1（陪集的基本性质）}：设 \(H \leq G\)。

\begin{enumerate}
\def\labelenumi{(\roman{enumi})}
\tightlist
\item
  \(a \in aH\)，故 \(G = \bigcup_{a \in G} aH\)
\item
  对任意 \(a, b \in G\)，\(aH = bH\) 或 \(aH \cap bH = \varnothing\)
\item
  \(|aH| = |H|\)（映射 \(h \mapsto ah\) 是双射）
\end{enumerate}

\textbf{证明}：

\begin{enumerate}
\def\labelenumi{(\roman{enumi})}
\item
  \(a = ae \in aH\)（\(e \in H\)）。
\item
  若 \(aH \cap bH \neq \varnothing\)，设 \(x \in aH \cap bH\)，则
  \(x = ah_1 = bh_2\)。于是 \(a = bh_2 h_1^{-1} \in bH\)，故
  \(aH \subseteq bH\)。同理 \(bH \subseteq aH\)，故 \(aH = bH\)。
\item
  映射 \(f : H \to aH\)，\(f(h) = ah\)
  是双射：\(ah_1 = ah_2 \implies h_1 = h_2\)（左消去）；对任意
  \(ah \in aH\)，\(f(h) = ah\)。\(\blacksquare\)
\end{enumerate}

\textbf{定理 39.2.2（Lagrange 定理）}：设 \(G\)
为有限群，\(H \leq G\)。则 \(|H|\) 整除 \(|G|\)，且

\[
|G| = |H| \cdot [G : H]
\]

其中 \([G : H]\) 为 \(H\) 在 \(G\)
中的\textbf{指数}（即不同陪集的个数）。

\textbf{证明}：由定理 39.2.1，\(G\) 被划分为 \([G : H]\)
个互不相交的左陪集，每个陪集的大小为 \(|H|\)。因此
\(|G| = [G : H] \cdot |H|\)。\(\blacksquare\)

\textbf{推论 39.2.3}：设 \(G\) 为有限群，\(a \in G\)。则 \(|a|\) 整除
\(|G|\)。特别地，\(a^{|G|} = e\)。

\textbf{证明}：考虑由 \(a\) 生成的循环子群
\(\langle a \rangle = \{a^k \mid k \in \mathbb{Z}\}\)。\(|\langle a \rangle| = |a|\)。由
Lagrange 定理，\(|a| \mid |G|\)。因此
\(a^{|G|} = (a^{|a|})^{|G|/|a|} = e\)。\(\blacksquare\)

\textbf{推论 39.2.4（Fermat-Euler 定理的推广）}：若
\(|G| = p\)（素数），则 \(G\) 是循环群，且任何非幺元都是生成元。

\textbf{证明}：取 \(a \in G \setminus \{e\}\)。\(|a|\) 整除 \(p\) 且
\(|a| > 1\)，故 \(|a| = p\)。因此
\(\langle a \rangle = G\)。\(\blacksquare\)

\bigskip\hrule\bigskip

\subsubsection{39.3
正规子群与商群}\label{ux6b63ux89c4ux5b50ux7fa4ux4e0eux5546ux7fa4}

\textbf{定义 39.3.1（正规子群）}：\(N \leq G\) 称为 \(G\)
的\textbf{正规子群}，记为 \(N \trianglelefteq G\)，若对任意
\(g \in G\)，\(gNg^{-1} = N\)（等价地，\(gN = Ng\)）。

\textbf{定理 39.3.1（正规子群的判定）}：\(N \trianglelefteq G\) 当且仅当
\(\forall g \in G, \forall n \in N\)，\(gng^{-1} \in N\)。

\textbf{证明}：\(gNg^{-1} \subseteq N\) 对所有 \(g\) 成立则
\(gNg^{-1} = N\)（用 \(g^{-1}\) 替换 \(g\) 得
\(g^{-1}Ng \subseteq N\)，两边左乘 \(g\) 右乘 \(g^{-1}\) 得
\(N \subseteq gNg^{-1}\)）。\(\blacksquare\)

\textbf{定义 39.3.2（商群）}：设 \(N \trianglelefteq G\)。在陪集的集合
\(G/N = \{gN \mid g \in G\}\) 上定义运算
\((aN)(bN) = (ab)N\)。\((G/N, \cdot)\) 构成群，称为 \(G\) 模 \(N\)
的\textbf{商群}。

\textbf{定理 39.3.2（商群中运算的良定义性）}：上述运算在
\(N \trianglelefteq G\) 时是良定义的。

\textbf{证明}：需证若 \(aN = a'N\) 且 \(bN = b'N\)，则
\((ab)N = (a'b')N\)。

由 \(aN = a'N\) 得 \(a' = an_1\)（\(n_1 \in N\)）。由 \(bN = b'N\) 得
\(b' = bn_2\)（\(n_2 \in N\)）。则
\(a'b' = an_1 bn_2 = ab(b^{-1}n_1 b)n_2\)。由于
\(N \trianglelefteq G\)，\(b^{-1}n_1 b \in N\)，故
\((b^{-1}n_1 b)n_2 \in N\)。因此 \(a'b' \in (ab)N\)，即
\((a'b')N = (ab)N\)。\(\blacksquare\)

\textbf{定理 39.3.3}：若 \(G\) 是阿贝尔群，则 \(G\)
的任何子群都是正规子群。

\textbf{证明}：\(gng^{-1} = gg^{-1}n = n \in N\)。\(\blacksquare\)

\bigskip\hrule\bigskip

\subsubsection{39.4
群同态与同构定理}\label{ux7fa4ux540cux6001ux4e0eux540cux6784ux5b9aux7406}

\textbf{定义 39.4.1（群同态）}：设 \(G, H\) 为群。映射
\(\varphi : G \to H\) 若满足
\(\varphi(ab) = \varphi(a)\varphi(b)\)（\(\forall a, b \in G\)），则称
\(\varphi\) 为\textbf{群同态}。

\textbf{定理 39.4.1（同态的基本性质）}：设 \(\varphi : G \to H\)
为群同态。

\begin{enumerate}
\def\labelenumi{(\roman{enumi})}
\tightlist
\item
  \(\varphi(e_G) = e_H\)
\item
  \(\varphi(a^{-1}) = \varphi(a)^{-1}\)
\item
  \(\ker \varphi = \{g \in G \mid \varphi(g) = e_H\}\) 是 \(G\)
  的正规子群
\item
  \(\operatorname{im} \varphi = \{\varphi(g) \mid g \in G\}\) 是 \(H\)
  的子群
\item
  \(\varphi\) 是单射 \(\iff\) \(\ker \varphi = \{e_G\}\)
\end{enumerate}

\textbf{证明}：

\begin{enumerate}
\def\labelenumi{(\roman{enumi})}
\item
  \(\varphi(e_G) = \varphi(e_G e_G) = \varphi(e_G)\varphi(e_G)\)，两边左乘
  \(\varphi(e_G)^{-1}\) 得 \(\varphi(e_G) = e_H\)。
\item
  \(\varphi(a)\varphi(a^{-1}) = \varphi(aa^{-1}) = \varphi(e_G) = e_H\)，同理
  \(\varphi(a^{-1})\varphi(a) = e_H\)。由逆元唯一性得证。
\item
  \(\ker \varphi\) 是子群（由 (i)(ii) 验证）。对
  \(g \in G\)，\(k \in \ker \varphi\)，\(\varphi(gkg^{-1}) = \varphi(g)\varphi(k)\varphi(g)^{-1} = \varphi(g)e_H \varphi(g)^{-1} = e_H\)，故
  \(gkg^{-1} \in \ker \varphi\)。因此
  \(\ker \varphi \trianglelefteq G\)。
\item
  直接验证。(v) 与线性映射的判定类似（定理 33.2.2）。\(\blacksquare\)
\end{enumerate}

\textbf{定理 39.4.2（第一同构定理）}：设 \(\varphi : G \to H\)
为群同态。则

\[
G / \ker \varphi \cong \operatorname{im} \varphi
\]

同构映射为 \(\overline{\varphi}(g \ker \varphi) = \varphi(g)\)。

\textbf{证明}：定义
\(\overline{\varphi} : G/\ker \varphi \to \operatorname{im} \varphi\) 为
\(\overline{\varphi}(gK) = \varphi(g)\)（\(K = \ker \varphi\)）。

良定义：若 \(gK = g'K\)，则
\(g' = gk\)（\(k \in K\)）。\(\varphi(g') = \varphi(gk) = \varphi(g)\varphi(k) = \varphi(g)\)。

同态：\(\overline{\varphi}((gK)(hK)) = \overline{\varphi}(ghK) = \varphi(gh) = \varphi(g)\varphi(h) = \overline{\varphi}(gK)\overline{\varphi}(hK)\)。

单射：若 \(\overline{\varphi}(gK) = e_H\)，则 \(\varphi(g) = e_H\)，故
\(g \in K\)，\(gK = K\)（商群的幺元）。

满射：由 \(\operatorname{im} \varphi\) 的定义显然。\(\blacksquare\)

\textbf{定理 39.4.3（第二同构定理）}：设
\(H \leq G\)，\(N \trianglelefteq G\)。则
\(HN = \{hn \mid h \in H, n \in N\}\) 是 \(G\)
的子群，\(H \cap N \trianglelefteq H\)，且

\[
H / (H \cap N) \cong HN / N
\]

\textbf{证明}：\(HN\)
是子群：\(e \in HN\)；\((h_1 n_1)(h_2 n_2) = h_1 h_2 (h_2^{-1} n_1 h_2) n_2 \in HN\)（因为
\(h_2^{-1} n_1 h_2 \in N\)）；\((hn)^{-1} = n^{-1}h^{-1} = h^{-1}(h n^{-1} h^{-1}) \in HN\)。

\(H \cap N \trianglelefteq H\)：对
\(h \in H\)，\(x \in H \cap N\)，\(hxh^{-1} \in H\)（\(h, x, h^{-1} \in H\)）且
\(hxh^{-1} \in N\)（\(N \trianglelefteq G\)），故
\(hxh^{-1} \in H \cap N\)。

定义 \(\varphi : H \to HN/N\) 为 \(\varphi(h) = hN\)。\(\varphi\)
是满同态。\(\ker \varphi = \{h \in H \mid hN = N\} = \{h \in H \mid h \in N\} = H \cap N\)。由第一同构定理，\(H/(H \cap N) \cong HN/N\)。\(\blacksquare\)

\textbf{定理 39.4.4（第三同构定理）}：设
\(N \trianglelefteq G\)，\(K \trianglelefteq G\) 且
\(N \subseteq K\)。则 \(K/N \trianglelefteq G/N\)，且

\[
(G/N) / (K/N) \cong G/K
\]

\textbf{证明}：定义 \(\varphi : G/N \to G/K\) 为 \(\varphi(gN) = gK\)。

良定义：若 \(gN = g'N\)，则 \(g' = gn\)（\(n \in N \subseteq K\)），故
\(g'K = gK\)。

同态：\(\varphi((gN)(hN)) = \varphi(ghN) = ghK = (gK)(hK) = \varphi(gN)\varphi(hN)\)。

满射显然。\(\ker \varphi = \{gN \mid gK = K\} = \{gN \mid g \in K\} = K/N\)。

由第一同构定理得证。\(\blacksquare\)

\bigskip\hrule\bigskip

\subsubsection{39.5 循环群}\label{ux5faaux73afux7fa4}

\textbf{定义 39.5.1（循环群）}：群 \(G\) 称为\textbf{循环群}，若存在
\(a \in G\) 使得
\(G = \langle a \rangle = \{a^n \mid n \in \mathbb{Z}\}\)。

\textbf{定理 39.5.1（循环群的分类）}：

\begin{enumerate}
\def\labelenumi{(\roman{enumi})}
\tightlist
\item
  若 \(|G| = \infty\)，则 \(G \cong \mathbb{Z}\)
\item
  若 \(|G| = n\)，则 \(G \cong \mathbb{Z}_n\)
\end{enumerate}

\textbf{证明}：定义 \(\varphi : \mathbb{Z} \to G\) 为
\(\varphi(k) = a^k\)。\(\varphi\) 是满同态。\(\ker \varphi\) 是
\(\mathbb{Z}\) 的子群，故
\(\ker \varphi = m\mathbb{Z}\)（\(m \geq 0\)）。

若 \(m = 0\)，则 \(\varphi\) 是单射，故 \(G \cong \mathbb{Z}\)。

若 \(m > 0\)，由第一同构定理
\(G \cong \mathbb{Z}/m\mathbb{Z} = \mathbb{Z}_m\)。此时
\(|G| = m = n\)。\(\blacksquare\)

\textbf{定理 39.5.2（循环群的子群）}：循环群的任何子群也是循环群。若
\(G = \langle a \rangle\) 为 \(n\) 阶循环群，则对 \(n\) 的每个正因子
\(d\)，\(G\) 恰有一个 \(d\) 阶子群，即 \(\langle a^{n/d} \rangle\)。

\textbf{证明}：设 \(H \leq G = \langle a \rangle\)。令
\(m = \min\{k > 0 \mid a^k \in H\}\)。可以证明
\(H = \langle a^m \rangle\)（用带余除法）。

对 \(n\) 阶循环群，\(d \mid n\)
时，\(|a^{n/d}| = n / \gcd(n, n/d) = n / (n/d) = d\)。唯一性：若 \(H\)
为 \(d\) 阶子群，则 \(H = \langle a^k \rangle\) 且
\(|a^k| = n / \gcd(n, k) = d\)，故 \(\gcd(n, k) = n/d\)。因此 \(k\) 是
\(n/d\) 的倍数，且 \(a^k \in \langle a^{n/d} \rangle\)。由此
\(H = \langle a^{n/d} \rangle\)。\(\blacksquare\)

\bigskip\hrule\bigskip

\subsubsection{39.6 对称群与 Cayley
定理}\label{ux5bf9ux79f0ux7fa4ux4e0e-cayley-ux5b9aux7406}

\textbf{定义 39.6.1（对称群）}：\(n\) 个元素的集合上所有置换构成的群称为
\(n\) 次\textbf{对称群}，记为 \(S_n\)。\(|S_n| = n!\)。

\textbf{定义 39.6.2（置换的轮换分解）}：任何置换 \(\sigma \in S_n\)
可唯一分解为互不相交的轮换的乘积（不考虑次序）。

\textbf{定理 39.6.1（Cayley 定理）}：任何群 \(G\)
同构于某个对称群的一个子群。特别地，若 \(|G| = n\)，则 \(G\) 同构于
\(S_n\) 的一个子群。

\textbf{证明}：对每个 \(g \in G\)，定义左乘映射
\(L_g : G \to G\)，\(L_g(x) = gx\)。\(L_g\) 是 \(G\) 上的置换（因为
\(L_{g^{-1}}\) 是其逆）。

定义 \(\varphi : G \to S_G\)（\(S_G\) 为 \(G\) 上所有置换构成的群）为
\(\varphi(g) = L_g\)。

\(\varphi\)
是同态：\(\varphi(gh)(x) = L_{gh}(x) = ghx = L_g(hx) = L_g(L_h(x)) = (\varphi(g) \circ \varphi(h))(x)\)。

\(\varphi\) 是单射：若 \(\varphi(g) = \text{id}\)，则
\(L_g(e) = ge = e\)，故 \(g = e\)。

因此 \(G \cong \operatorname{im} \varphi \leq S_G\)。若 \(|G| = n\)，则
\(S_G \cong S_n\)。\(\blacksquare\)

\bigskip\hrule\bigskip

\subsubsection{39.7 群的直积}\label{ux7fa4ux7684ux76f4ux79ef}

\textbf{定义 39.7.1（直积）}：设 \(G_1, G_2, \ldots, G_n\)
为群。它们的\textbf{直积}为

\[
G_1 \times G_2 \times \cdots \times G_n = \{(g_1, g_2, \ldots, g_n) \mid g_i \in G_i\}
\]

运算为逐分量乘法：\((g_1, \ldots, g_n)(h_1, \ldots, h_n) = (g_1 h_1, \ldots, g_n h_n)\)。

\textbf{定理
39.7.1（有限阿贝尔群的结构定理，陈述）}：任何有限阿贝尔群同构于素数幂阶循环群的直积。即存在素数
\(p_1, \ldots, p_k\) 和正整数 \(e_1, \ldots, e_k\) 使得

\[
G \cong \mathbb{Z}_{p_1^{e_1}} \times \mathbb{Z}_{p_2^{e_2}} \times \cdots \times \mathbb{Z}_{p_k^{e_k}}
\]

该分解在因子顺序下是唯一的（当要求
\(p_1^{e_1} \mid p_2^{e_2} \mid \cdots \mid p_k^{e_k}\) 时）。

\textbf{注记 39.7.1}：完整证明将在 Ch 43
中作为主理想整环上有限生成模的结构定理的推论给出。

\bigskip\hrule\bigskip

\subsection{Ch 40 群作用与 Sylow
定理}\label{ch-40-ux7fa4ux4f5cux7528ux4e0e-sylow-ux5b9aux7406}

\subsubsection{40.1 群作用}\label{ux7fa4ux4f5cux7528}

\textbf{定义 40.1.1（群作用）}：群 \(G\) 在集合 \(X\)
上的一个\textbf{作用}（左作用）是一个映射
\(G \times X \to X\)，\((g, x) \mapsto g \cdot x\)，满足：

\begin{enumerate}
\def\labelenumi{\arabic{enumi}.}
\tightlist
\item
  \(e \cdot x = x\)（\(\forall x \in X\)）
\item
  \(g \cdot (h \cdot x) = (gh) \cdot x\)（\(\forall g, h \in G, \forall x \in X\)）
\end{enumerate}

\textbf{定义 40.1.2（轨道与稳定子）}：设 \(G\) 作用在 \(X\)
上，\(x \in X\)。

\begin{itemize}
\tightlist
\item
  \(x\) 的\textbf{轨道}：\(G \cdot x = \{g \cdot x \mid g \in G\}\)
\item
  \(x\) 的\textbf{稳定子群}：\(G_x = \{g \in G \mid g \cdot x = x\}\)
\end{itemize}

\textbf{定理 40.1.1（轨道-稳定子定理）}：设 \(G\) 为有限群，作用在有限集
\(X\) 上。则对任意 \(x \in X\)，

\[
|G \cdot x| = [G : G_x] = \frac{|G|}{|G_x|}
\]

特别地，\(|G \cdot x|\) 整除 \(|G|\)。

\textbf{证明}：定义映射 \(\varphi : G/G_x \to G \cdot x\) 为
\(\varphi(g G_x) = g \cdot x\)。

良定义：若 \(g G_x = h G_x\)，则 \(h^{-1}g \in G_x\)，故
\((h^{-1}g) \cdot x = x\)，即 \(g \cdot x = h \cdot x\)。

单射：若 \(g \cdot x = h \cdot x\)，则 \(h^{-1}g \cdot x = x\)，故
\(h^{-1}g \in G_x\)，\(g G_x = h G_x\)。

满射：显然。

因此 \(|G \cdot x| = |G/G_x| = [G : G_x] = |G|/|G_x|\)。\(\blacksquare\)

\textbf{定理 40.1.2（Burnside 引理）}：设有限群 \(G\) 作用在有限集 \(X\)
上。轨道的个数为

\[
\frac{1}{|G|} \sum_{g \in G} |\operatorname{Fix}(g)|
\]

其中 \(\operatorname{Fix}(g) = \{x \in X \mid g \cdot x = x\}\)。

\textbf{证明}：计数集合 \(\{(g, x) \in G \times X \mid g \cdot x = x\}\)
的两种方式：

\[
\sum_{g \in G} |\operatorname{Fix}(g)| = \sum_{x \in X} |G_x| = \sum_{x \in X} \frac{|G|}{|G \cdot x|}
\]

设轨道为 \(\mathcal{O}_1, \ldots, \mathcal{O}_k\)。对每个轨道
\(\mathcal{O}_i\)，取 \(x_i \in \mathcal{O}_i\)，则对
\(x \in \mathcal{O}_i\)，\(|G \cdot x| = |\mathcal{O}_i|\)，故

\[
\sum_{x \in \mathcal{O}_i} \frac{|G|}{|G \cdot x|} = |\mathcal{O}_i| \cdot \frac{|G|}{|\mathcal{O}_i|} = |G|
\]

因此 \(\sum_{g \in G} |\operatorname{Fix}(g)| = k |G|\)，即
\(k = \frac{1}{|G|} \sum_{g \in G} |\operatorname{Fix}(g)|\)。\(\blacksquare\)

\bigskip\hrule\bigskip

\subsubsection{40.2
共轭作用与类方程}\label{ux5171ux8f6dux4f5cux7528ux4e0eux7c7bux65b9ux7a0b}

\textbf{定义 40.2.1（共轭作用）}：群 \(G\)
在自身上的\textbf{共轭作用}定义为 \(g \cdot x = gxg^{-1}\)。

在此作用下，\(x\) 的轨道称为 \(x\) 的\textbf{共轭类}，稳定子群称为 \(x\)
的\textbf{中心化子}，记为 \(C_G(x) = \{g \in G \mid gxg^{-1} = x\}\)。

\textbf{定义 40.2.2（中心）}：\(G\) 的\textbf{中心}
\(Z(G) = \{g \in G \mid gx = xg, \, \forall x \in G\} = \{g \in G \mid C_G(g) = G\}\)。

\textbf{定理 40.2.1（类方程）}：设 \(G\) 为有限群。则

\[
|G| = |Z(G)| + \sum_{i=1}^{k} [G : C_G(x_i)]
\]

其中 \(x_1, \ldots, x_k\) 是非中心元素的共轭类的代表元。

\textbf{证明}：\(G\) 被划分为共轭类。\(Z(G)\)
中的元素恰为那些共轭类大小为 \(1\)
的元素。对其余共轭类，由轨道-稳定子定理，\(|共轭类| = [G : C_G(x_i)]\)。\(\blacksquare\)

\textbf{推论 40.2.2}：若 \(|G| = p^n\)（\(p\) 为素数，\(n \geq 1\)，此时
\(G\) 称为 \(p\)-群），则 \(|Z(G)| > 1\)（即非平凡
\(p\)-群的中心非平凡）。

\textbf{证明}：由类方程，\(|G| = |Z(G)| + \sum [G : C_G(x_i)]\)。\([G : C_G(x_i)]\)
整除 \(|G| = p^n\) 且不为 \(1\)（\(x_i \notin Z(G)\)），故是 \(p\)
的倍数。因此 \(|Z(G)| \equiv |G| \equiv 0 \pmod{p}\)。又
\(e \in Z(G)\)，故 \(|Z(G)| \geq p > 1\)。\(\blacksquare\)

\bigskip\hrule\bigskip

\subsubsection{40.3 Sylow 定理}\label{sylow-ux5b9aux7406}

\textbf{定义 40.3.1（Sylow \(p\)-子群）}：设 \(G\)
为有限群，\(|G| = p^n m\)（\(p \nmid m\)）。\(G\) 的 \(p^n\) 阶子群称为
\(G\) 的 \textbf{Sylow \(p\)-子群}。

\textbf{定理 40.3.1（Sylow 第一定理------存在性）}：设 \(G\)
为有限群，\(p\) 为素数，\(p^k \mid |G|\)。则 \(G\) 存在 \(p^k\)
阶子群。特别地，Sylow \(p\)-子群存在。

\textbf{证明}：对 \(|G|\) 归纳。\(|G| = 1\) 时平凡。

设 \(|G| = p^n m\)（\(p \nmid m\)）。考虑类方程
\(|G| = |Z(G)| + \sum [G : C_G(x_i)]\)。

\textbf{情形一}：\(p \nmid |Z(G)|\)。则存在某个 \(x_i\) 使得
\(p \nmid [G : C_G(x_i)]\)。由 \([G : C_G(x_i)] = |G|/|C_G(x_i)|\) 和
\(p^n \mid |G|\)，得 \(p^n \mid |C_G(x_i)|\)。由于
\(x_i \notin Z(G)\)，\(C_G(x_i) \subsetneq G\)。由归纳假设应用在
\(C_G(x_i)\) 上，\(C_G(x_i)\) 存在 \(p^n\) 阶子群（即 \(G\) 的 Sylow
\(p\)-子群）。

\textbf{情形二}：\(p \mid |Z(G)|\)。由 Cauchy 定理（定理
40.3.2），\(Z(G)\) 存在 \(p\) 阶元素
\(a\)。\(N = \langle a \rangle \trianglelefteq G\)（因为
\(N \subseteq Z(G)\)）。\(|G/N| = |G|/p = p^{n-1}m\)。由归纳假设，\(G/N\)
存在 Sylow \(p\)-子群 \(P/N\)（\(|P/N| = p^{n-1}\)）。则
\(|P| = p^n\)，\(P\) 即为 \(G\) 的 Sylow \(p\)-子群。

对于一般的 \(p^k\)，由 Sylow \(p\)-子群的存在性，它包含 \(p^k\)
阶子群（由归纳法可证）。\(\blacksquare\)

\textbf{定理 40.3.2（Cauchy 定理）}：若素数 \(p \mid |G|\)，则 \(G\)
存在 \(p\) 阶元素。

\textbf{证明}：考虑集合
\(X = \{(g_1, \ldots, g_p) \in G^p : g_1 g_2 \cdots g_p = e\}\)。\(|X| = |G|^{p-1}\)（前
\(p-1\) 个分量可自由选择，最后一个由方程
\(g_p = (g_1 \cdots g_{p-1})^{-1}\) 唯一确定）。循环群 \(\mathbb{Z}_p\)
通过循环置换作用在 \(X\)
上：\(k \cdot (g_1, \ldots, g_p) = (g_{k+1}, \ldots, g_p, g_1, \ldots, g_k)\)（\(k = 0, 1, \ldots, p-1\)）。验证这是群作用：\(0\)
作用为恒等，且 \(k \cdot (\ell \cdot x) = (k+\ell) \cdot x\)。

由轨道-稳定子定理（定理 40.1.1），每个轨道的大小整除
\(|\mathbb{Z}_p| = p\)，故为 \(1\) 或 \(p\)。不动点（轨道大小为 \(1\)
的元素）满足 \((g, \ldots, g) \in X\)，即 \(g^p = e\)。不动点包括
\((e, \ldots, e)\)（对应 \(g = e\)）以及可能的满足 \(|g| = p\) 的元素。

由于 \(p \mid |X| = |G|^{p-1}\)，而不动点总数
\(\equiv |X| \equiv 0 \pmod{p}\)（因非平凡轨道大小为 \(p\)），故至少有
\(p\) 个不动点。因此存在 \(g \neq e\) 使 \(g^p = e\)，即
\(|g| = p\)。\(\blacksquare\)

\textbf{定理 40.3.3（Sylow 第二定理------共轭性）}：设 \(G\)
为有限群。\(G\) 的任何两个 Sylow \(p\)-子群是共轭的。即若 \(P, Q\) 均为
Sylow \(p\)-子群，则存在 \(g \in G\) 使得 \(Q = gPg^{-1}\)。

\textbf{证明}：设 \(X = \{gP \mid g \in G\}\) 为 \(P\)
的左陪集组成的集合。\(Q\) 通过左乘作用在 \(X\) 上。将 \(X\) 划分为
\(Q\)-轨道。

由轨道-稳定子定理，每个轨道的大小整除 \(|Q| = p^n\)，故为 \(p\)
的幂。\(|X| = [G : P] = m\)（\(p \nmid m\)）。由于轨道的和为
\(m\)，必存在大小为 \(1\) 的轨道（否则每个轨道大小都是 \(p\)
的倍数，总和也是 \(p\) 的倍数）。

设该轨道为 \(\{gP\}\)，即 \(Q \cdot gP = gP\)。这意味着 \(QgP = gP\)，即
\(g^{-1}Qg \subseteq P\)。由于 \(|g^{-1}Qg| = |Q| = |P|\)，故
\(g^{-1}Qg = P\)，即 \(Q = gPg^{-1}\)。\(\blacksquare\)

\textbf{推论 40.3.4（Sylow 第三定理------计数）}：设 \(n_p\) 为 \(G\) 的
Sylow \(p\)-子群的个数。则

\begin{enumerate}
\def\labelenumi{(\roman{enumi})}
\tightlist
\item
  \(n_p \equiv 1 \pmod{p}\)
\item
  \(n_p \mid m\)（其中 \(|G| = p^n m\)，\(p \nmid m\)）
\end{enumerate}

\textbf{证明}：设 \(P\) 为 Sylow \(p\)-子群。令 \(X\) 为所有 Sylow
\(p\)-子群的集合。\(P\) 通过共轭作用在 \(X\) 上。将 \(X\) 划分为
\(P\)-轨道。

每个轨道的大小整除 \(|P| = p^n\)，故为 \(p\) 的幂。\(\{P\}\) 是大小为
\(1\) 的轨道。若存在另一个大小为 \(1\) 的轨道 \(\{Q\}\)，则
\(P Q P^{-1} = Q\)，即 \(P \subseteq N_G(Q)\)。\(P\) 和 \(Q\) 都是
\(N_G(Q)\) 的 Sylow \(p\)-子群。由 Sylow 第二定理（应用在 \(N_G(Q)\)
上），\(P\) 和 \(Q\) 在 \(N_G(Q)\) 中共轭，但
\(Q \trianglelefteq N_G(Q)\)，故 \(P = Q\)。

因此仅有 \(\{P\}\) 是大小为 \(1\) 的轨道，其余轨道大小均为 \(p\)
的倍数。故 \(n_p \equiv 1 \pmod{p}\)。

\begin{enumerate}
\def\labelenumi{(\roman{enumi})}
\setcounter{enumi}{1}
\tightlist
\item
  由 Sylow 第二定理，\(G\) 传递地作用在 Sylow
  \(p\)-子群的集合上（通过共轭）。由轨道-稳定子定理，\(n_p = [G : N_G(P)]\)
  整除 \([G : P] = m\)。\(\blacksquare\)
\end{enumerate}

\bigskip\hrule\bigskip

\subsection{Ch 41 环论基础}\label{ch-41-ux73afux8bbaux57faux7840}

\subsubsection{41.1 理想与商环}\label{ux7406ux60f3ux4e0eux5546ux73af}

\textbf{定义 41.1.1（理想）}：设 \(R\) 为环。子集 \(I \subseteq R\) 称为
\(R\) 的\textbf{理想}，若

\begin{enumerate}
\def\labelenumi{\arabic{enumi}.}
\tightlist
\item
  \((I, +)\) 是 \((R, +)\) 的子群
\item
  \(\forall r \in R, \forall a \in I\)，\(ra \in I\)（左理想）且
  \(ar \in I\)（右理想）
\end{enumerate}

\textbf{例 41.1.1}：设 \(R\) 为环，\(a \in R\)。则
\((a) = \{ra \mid r \in R\}\) 是 \(R\) 的左理想。若 \(R\)
是交换环，\((a) = \{ra \mid r \in R\}\) 是理想，称为由 \(a\)
生成的\textbf{主理想}。

\textbf{定义 41.1.2（商环）}：设 \(I\) 为环 \(R\) 的理想。在加法商群
\(R/I\) 上定义乘法：\((r + I)(s + I) = rs + I\)。\((R/I, +, \cdot)\)
构成环，称为 \(R\) 模 \(I\) 的\textbf{商环}。

\textbf{定理 41.1.1}：商环中乘法的定义是良定义的。

\textbf{证明}：设 \(r + I = r' + I\)，\(s + I = s' + I\)。则
\(r' = r + a\)，\(s' = s + b\)（\(a, b \in I\)）。\(r's' = (r + a)(s + b) = rs + rb + as + ab\)。由于
\(I\) 是理想，\(rb, as, ab \in I\)，故
\(r's' + I = rs + I\)。\(\blacksquare\)

\bigskip\hrule\bigskip

\subsubsection{41.2
环同态与同构定理}\label{ux73afux540cux6001ux4e0eux540cux6784ux5b9aux7406}

\textbf{定义 41.2.1（环同态）}：设 \(R, S\) 为环。映射
\(\varphi : R \to S\) 若满足：

\begin{enumerate}
\def\labelenumi{\arabic{enumi}.}
\tightlist
\item
  \(\varphi(a + b) = \varphi(a) + \varphi(b)\)
\item
  \(\varphi(ab) = \varphi(a)\varphi(b)\)
\item
  \(\varphi(1_R) = 1_S\)（若 \(R, S\) 含幺）
\end{enumerate}

则称 \(\varphi\) 为\textbf{环同态}。

\textbf{定理 41.2.1（环的第一同构定理）}：设 \(\varphi : R \to S\)
为环同态。则 \(\ker \varphi\) 是 \(R\) 的理想，且

\[
R / \ker \varphi \cong \operatorname{im} \varphi
\]

\textbf{证明}：\(\ker \varphi\)
是理想（类似群的情形，需额外验证乘法吸收性）。在加法群层面，由群的第一同构定理，\(R/\ker \varphi \cong \operatorname{im} \varphi\)（作为加法群）。同构映射
\(\overline{\varphi}(r + \ker \varphi) = \varphi(r)\)
也保持乘法：\(\overline{\varphi}((r + K)(s + K)) = \overline{\varphi}(rs + K) = \varphi(rs) = \varphi(r)\varphi(s) = \overline{\varphi}(r + K)\overline{\varphi}(s + K)\)。\(\blacksquare\)

\textbf{定理 41.2.2（环的第二、第三同构定理）}：与群的同构定理类似，将
``子群'' 替换为 ``子环''，``正规子群'' 替换为 ``理想''。

\bigskip\hrule\bigskip

\subsubsection{41.3
素理想与极大理想}\label{ux7d20ux7406ux60f3ux4e0eux6781ux5927ux7406ux60f3}

\textbf{定义 41.3.1（素理想与极大理想）}：设 \(R\) 为交换环，\(I\) 为
\(R\) 的真理想（\(I \neq R\)）。

\begin{enumerate}
\def\labelenumi{\arabic{enumi}.}
\tightlist
\item
  \(I\) 称为\textbf{素理想}，若 \(ab \in I \implies a \in I\) 或
  \(b \in I\)
\item
  \(I\) 称为\textbf{极大理想}，若不存在理想 \(J\) 使得
  \(I \subsetneq J \subsetneq R\)
\end{enumerate}

\textbf{定理 41.3.1}：设 \(R\) 为含幺交换环，\(I\) 为 \(R\) 的理想。

\begin{enumerate}
\def\labelenumi{(\roman{enumi})}
\tightlist
\item
  \(I\) 是素理想 \(\iff\) \(R/I\) 是整环
\item
  \(I\) 是极大理想 \(\iff\) \(R/I\) 是域
\end{enumerate}

\textbf{证明}：

\begin{enumerate}
\def\labelenumi{(\roman{enumi})}
\item
  \(I\) 是素理想 \(\iff\)
  \(\forall a, b \in R\)，\(ab \in I \implies a \in I \text{ 或 } b \in I\)
  \(\iff\)
  \(\forall \overline{a}, \overline{b} \in R/I\)，\(\overline{a}\overline{b} = \overline{0} \implies \overline{a} = \overline{0} \text{ 或 } \overline{b} = \overline{0}\)
  \(\iff\) \(R/I\) 无零因子 \(\iff\) \(R/I\) 是整环。
\item
  （\(\Rightarrow\)）设 \(I\) 为极大理想。取
  \(\overline{a} \in R/I \setminus \{\overline{0}\}\)，即
  \(a \notin I\)。考虑理想 \(J = I + (a)\)。\(J \supsetneq I\)，由极大性
  \(J = R\)。故存在 \(i \in I\)，\(r \in R\) 使得 \(1 = i + ra\)。在
  \(R/I\) 中，\(\overline{1} = \overline{r}\overline{a}\)，故
  \(\overline{a}\) 可逆。因此 \(R/I\) 中每个非零元可逆，即为域。
\end{enumerate}

（\(\Leftarrow\)）设 \(R/I\) 为域。若 \(J \supseteq I\) 为 \(R\)
的理想，\(J/I\) 是 \(R/I\) 的理想。域只有平凡理想，故
\(J/I = \{\overline{0}\}\) 或 \(J/I = R/I\)。前者 \(J = I\)，后者
\(J = R\)。因此 \(I\) 是极大理想。\(\blacksquare\)

\textbf{推论 41.3.2}：极大理想必为素理想。

\textbf{证明}：域 \(\implies\) 整环。由定理 41.3.1，\(R/I\) 是域
\(\implies\) \(R/I\) 是整环 \(\implies\) \(I\) 是极大理想 \(\implies\)
\(I\) 是素理想。\(\blacksquare\)

\bigskip\hrule\bigskip

\subsubsection{41.4
中国剩余定理}\label{ux4e2dux56fdux5269ux4f59ux5b9aux7406-1}

\textbf{定理 41.4.1（中国剩余定理------环论版）}：设 \(R\)
为含幺交换环，\(I_1, \ldots, I_n\) 为 \(R\) 的理想，满足
\(I_i + I_j = R\)（\(\forall i \neq j\)）。则

\[
R / (I_1 \cap I_2 \cap \cdots \cap I_n) \cong R/I_1 \times R/I_2 \times \cdots \times R/I_n
\]

\textbf{证明}：对 \(n\) 归纳。\(n = 2\) 时，\(I_1 + I_2 = R\) 保证存在
\(a_1 \in I_1\)，\(a_2 \in I_2\) 使得 \(a_1 + a_2 = 1\)。

定义 \(\varphi : R \to R/I_1 \times R/I_2\) 为
\(\varphi(r) = (r + I_1, r + I_2)\)。\(\varphi\) 是环同态。

\(\ker \varphi = I_1 \cap I_2\)。需证 \(\varphi\) 是满射。对任意
\((r_1 + I_1, r_2 + I_2)\)，取 \(r = r_2 a_1 + r_1 a_2\)。则

\[
r - r_1 = r_2 a_1 + r_1(a_2 - 1) = r_2 a_1 - r_1 a_1 \in I_1
\] \[
r - r_2 = r_2(a_1 - 1) + r_1 a_2 = -r_2 a_2 + r_1 a_2 \in I_2
\]

故 \(\varphi(r) = (r_1 + I_1, r_2 + I_2)\)。由第一同构定理得证。

\(n > 2\) 时，令 \(J = I_1 \cap \cdots \cap I_{n-1}\)。可证
\(J + I_n = R\)（利用 \(I_i + I_n = R\) 和归纳）。由 \(n = 2\)
的情形，\(R/(J \cap I_n) \cong R/J \times R/I_n\)。由归纳假设，\(R/J \cong R/I_1 \times \cdots \times R/I_{n-1}\)。合起来即得。\(\blacksquare\)

\textbf{推论 41.4.2（整数版中国剩余定理）}：设 \(m_1, \ldots, m_n\)
为两两互素的正整数。则

\[
\mathbb{Z} / (m_1 \cdots m_n) \mathbb{Z} \cong \mathbb{Z}/m_1\mathbb{Z} \times \cdots \times \mathbb{Z}/m_n\mathbb{Z}
\]

\textbf{证明}：\(m_i\mathbb{Z} + m_j\mathbb{Z} = \gcd(m_i, m_j)\mathbb{Z} = \mathbb{Z}\)（\(\gcd(m_i, m_j) = 1\)）。且
\(m_1\mathbb{Z} \cap \cdots \cap m_n\mathbb{Z} = \operatorname{lcm}(m_1, \ldots, m_n)\mathbb{Z} = m_1 \cdots m_n \mathbb{Z}\)。由定理
41.4.1 得证。\(\blacksquare\)

\bigskip\hrule\bigskip

\subsection{Ch 42
多项式环与因式分解}\label{ch-42-ux591aux9879ux5f0fux73afux4e0eux56e0ux5f0fux5206ux89e3}

\subsubsection{42.1 多项式环}\label{ux591aux9879ux5f0fux73af}

\textbf{定义 42.1.1（多项式环）}：设 \(R\) 为交换环。\(R\)
上的一元多项式环 \(R[x]\) 是所有形式表达式
\(\sum_{i=0}^{n} a_i x^i\)（\(a_i \in R\)）的集合，配备通常的多项式加法和乘法。

\textbf{定理 42.1.1（带余除法）}：设 \(F\)
为域，\(f, g \in F[x]\)，\(g \neq 0\)。则存在唯一的 \(q, r \in F[x]\)
使得

\[
f = qg + r, \quad \deg r < \deg g
\]

\textbf{证明}：对 \(\deg f\) 归纳。若 \(\deg f < \deg g\)，取
\(q = 0\)，\(r = f\) 即可。

若 \(\deg f = n \geq m = \deg g\)，设
\(f = a_n x^n + \cdots\)，\(g = b_m x^m + \cdots\)。令
\(f_1 = f - \frac{a_n}{b_m} x^{n-m} g\)。则
\(\deg f_1 < n\)。由归纳假设，\(f_1 = q_1 g + r\)（\(\deg r < \deg g\)）。则
\(f = (q_1 + \frac{a_n}{b_m} x^{n-m})g + r\)。

唯一性：若 \(f = qg + r = q'g + r'\)，则 \((q - q')g = r' - r\)。若
\(q \neq q'\)，左边次数 \(\geq \deg g\)，右边次数 \(< \deg g\)，矛盾。故
\(q = q'\)，\(r = r'\)。\(\blacksquare\)

\textbf{定义
42.1.2（多项式函数与根）}：\(f \in F[x]\)，\(\alpha \in F\)。若
\(f(\alpha) = 0\)，则称 \(\alpha\) 为 \(f\) 的\textbf{根}。

\textbf{定理 42.1.2（因式定理）}：\(\alpha \in F\) 是 \(f \in F[x]\)
的根当且仅当 \(x - \alpha \mid f(x)\)。

\textbf{证明}：由带余除法，\(f(x) = q(x)(x - \alpha) + r\)（\(r \in F\)）。代入
\(x = \alpha\) 得 \(f(\alpha) = r\)。故
\(f(\alpha) = 0 \iff r = 0 \iff x - \alpha \mid f(x)\)。\(\blacksquare\)

\bigskip\hrule\bigskip

\subsubsection{42.2
欧几里得整环、主理想整环、唯一分解整环}\label{ux6b27ux51e0ux91ccux5f97ux6574ux73afux4e3bux7406ux60f3ux6574ux73afux552fux4e00ux5206ux89e3ux6574ux73af}

\textbf{定义 42.2.1（欧几里得整环）}：整环 \(R\)
称为\textbf{欧几里得整环}（ED），若存在函数
\(\varphi : R \setminus \{0\} \to \mathbb{N}\) 使得对任意
\(a, b \in R\)（\(b \neq 0\)），存在 \(q, r \in R\) 满足

\[
a = qb + r, \quad \text{且 } r = 0 \text{ 或 } \varphi(r) < \varphi(b)
\]

\textbf{例 42.2.1}：\(\mathbb{Z}\)（\(\varphi(n) = |n|\)）和
\(F[x]\)（\(\varphi(f) = \deg f\)）是 ED。

\textbf{定义 42.2.2（主理想整环）}：整环 \(R\)
称为\textbf{主理想整环}（PID），若 \(R\)
的每个理想都是主理想（即由单个元素生成）。

\textbf{定理 42.2.1}：ED \(\implies\) PID。

\textbf{证明}：设 \(R\) 为 ED，\(I\) 为 \(R\) 的非零理想。取
\(b \in I \setminus \{0\}\) 使得 \(\varphi(b)\) 最小。对任意
\(a \in I\)，由带余除法 \(a = qb + r\)（\(r = 0\) 或
\(\varphi(r) < \varphi(b)\)）。由于 \(r = a - qb \in I\)，由
\(\varphi(b)\) 的最小性，\(r = 0\)。故 \(a = qb\)，即
\(I = (b)\)。\(\blacksquare\)

\textbf{定义 42.2.3（不可约元素与素元素）}：设 \(R\) 为整环。

\begin{itemize}
\tightlist
\item
  \(p \in R \setminus \{0\}\)（非单位）称为\textbf{不可约元素}，若
  \(p = ab \implies a\) 或 \(b\) 是单位
\item
  \(p \in R \setminus \{0\}\)（非单位）称为\textbf{素元素}，若
  \(p \mid ab \implies p \mid a\) 或 \(p \mid b\)
\end{itemize}

\textbf{定理 42.2.2}：在整环中，素元素必为不可约元素。在 PID
中，不可约元素也是素元素。

\textbf{证明}：设 \(p\) 为素元素，\(p = ab\)。则 \(p \mid ab\)，故
\(p \mid a\) 或 \(p \mid b\)。不妨设 \(p \mid a\)，\(a = pc\)。则
\(p = pcb\)，由于整环中可消去 \(p\)，\(cb = 1\)，故 \(b\) 是单位。因此
\(p\) 不可约。

设 \(R\) 为 PID，\(p\) 不可约。若 \(p \mid ab\) 且
\(p \nmid a\)。考虑理想 \((p, a)\)。由于 \(R\) 是
PID，\((p, a) = (d)\)。\(d \mid p\)，故 \(d\) 是单位或 \(p\) 的相伴。若
\(d\) 是 \(p\) 的相伴，则 \(p \mid a\)，矛盾。故 \(d\)
是单位，\((p, a) = R\)。存在 \(x, y \in R\) 使得 \(xp + ya = 1\)。则
\(b = xpb + yab\)。\(p \mid xpb\) 且 \(p \mid yab\)（因为
\(p \mid ab\)），故 \(p \mid b\)。因此 \(p\) 是素元素。\(\blacksquare\)

\textbf{定义 42.2.4（唯一分解整环）}：整环 \(R\)
称为\textbf{唯一分解整环}（UFD），若

\begin{enumerate}
\def\labelenumi{\arabic{enumi}.}
\tightlist
\item
  每个非零非单位元素可分解为不可约元素的乘积
\item
  该分解在相伴意义下唯一（即若
  \(a = p_1 \cdots p_m = q_1 \cdots q_n\)，则 \(m = n\) 且适当重排后
  \(p_i\) 与 \(q_i\) 相伴）
\end{enumerate}

\textbf{定理 42.2.3}：PID \(\implies\) UFD。

\textbf{证明}（概要）：

（存在性）若存在非零非单位元素 \(a\)
不能分解为不可约元素的乘积，则可构造无穷长的严格增链
\((a) \subsetneq (a_1) \subsetneq (a_2) \subsetneq \cdots\)。但
\(\bigcup (a_i)\) 是理想，在 PID 中为主理想 \((d)\)，\(d\) 属于某个
\((a_n)\)，导致链终止，矛盾。

（唯一性）由定理 42.2.2，PID 中不可约元素是素元素。若
\(p_1 \cdots p_m = q_1 \cdots q_n\)，\(p_1\) 是素的，则
\(p_1 \mid q_j\)（某个 \(j\)）。\(q_j\) 不可约，故 \(p_1\) 与 \(q_j\)
相伴。消去 \(p_1\) 和 \(q_j\)，对剩余部分归纳。\(\blacksquare\)

\textbf{推论 42.2.4}：\(F[x]\) 是 UFD。

\bigskip\hrule\bigskip

\subsubsection{42.3 Gauss 引理与 Eisenstein
判别法}\label{gauss-ux5f15ux7406ux4e0e-eisenstein-ux5224ux522bux6cd5}

\textbf{定义 42.3.1（本原多项式）}：\(f \in \mathbb{Z}[x]\)
称为\textbf{本原多项式}，若其系数的最大公因数为 \(1\)。

\textbf{定理 42.3.1（Gauss 引理）}：两个本原多项式的乘积是本原多项式。

\textbf{证明}：设 \(f = \sum a_i x^i\)，\(g = \sum b_j x^j\)
为本原多项式。若 \(fg\) 不是本原的，则存在素数 \(p\) 整除 \(fg\)
的所有系数。设 \(a_r\) 是 \(f\) 中第一个不被 \(p\) 整除的系数，\(b_s\)
是 \(g\) 中第一个不被 \(p\) 整除的系数。则 \(fg\) 的 \(x^{r+s}\) 系数为

\[
c_{r+s} = \sum_{i+j = r+s} a_i b_j = a_r b_s + \sum_{i < r \text{ 或 } j < s} a_i b_j
\]

由 \(r, s\) 的选法，\(a_r b_s\) 不被 \(p\)
整除，而其余项至少有一个因子被 \(p\) 整除，矛盾。故 \(fg\)
是本原的。\(\blacksquare\)

\textbf{定理 42.3.2（Gauss 引理------推论）}：若 \(f \in \mathbb{Z}[x]\)
在 \(\mathbb{Q}[x]\) 中可约，则 \(f\) 在 \(\mathbb{Z}[x]\) 中也可约。

\textbf{证明}：设 \(f = gh\)（\(g, h \in \mathbb{Q}[x]\)）。通分可得
\(df = g_1 h_1\)（\(g_1, h_1 \in \mathbb{Z}[x]\)，\(d \in \mathbb{Z}\)）。提取
\(g_1, h_1\) 的各系数的最大公因数，得
\(g_1 = c_1 g^*\)，\(h_1 = c_2 h^*\)（\(g^*, h^*\) 本原）。由 Gauss
引理，\(g^* h^*\) 本原，故 \(d = c_1 c_2\)。从而 \(f = c_1 c_2 g^* h^*\)
在 \(\mathbb{Z}[x]\) 中可约。\(\blacksquare\)

\textbf{定理 42.3.3（Eisenstein 判别法）}：设
\(f(x) = a_n x^n + \cdots + a_1 x + a_0 \in \mathbb{Z}[x]\)。若存在素数
\(p\) 使得

\begin{enumerate}
\def\labelenumi{(\roman{enumi})}
\tightlist
\item
  \(p \nmid a_n\)
\item
  \(p \mid a_i\)（\(i = 0, 1, \ldots, n-1\)）
\item
  \(p^2 \nmid a_0\)
\end{enumerate}

则 \(f\) 在 \(\mathbb{Q}[x]\) 中不可约。

\textbf{证明}：假设 \(f = gh\)（\(g, h \in \mathbb{Z}[x]\)，非单位）。设
\(g = b_r x^r + \cdots + b_0\)，\(h = c_s x^s + \cdots + c_0\)（\(r + s = n\)）。

由于 \(a_0 = b_0 c_0\) 且 \(p \mid a_0\) 但 \(p^2 \nmid a_0\)，不妨设
\(p \mid b_0\) 但 \(p \nmid c_0\)。

\(p \nmid a_n\) 意味着 \(p \nmid b_r\) 且 \(p \nmid c_s\)。设 \(b_k\) 是
\(b_0, \ldots, b_r\) 中第一个不被 \(p\) 整除的系数。则
\(k \leq r < n\)。

考虑 \(a_k = \sum_{i=0}^{k} b_i c_{k-i}\)。对
\(i < k\)，\(p \mid b_i\)；对 \(i = k\)，\(p \nmid b_k\) 且
\(p \nmid c_0\)（因为 \(c_0\) 不被 \(p\) 整除）。故 \(p \nmid a_k\)。但
\(k < n\)，由条件 (ii) \(p \mid a_k\)，矛盾。\(\blacksquare\)

\textbf{例 42.3.1}：\(x^n - p\) 在 \(\mathbb{Q}[x]\) 中不可约（取素数
\(p\) 作为 Eisenstein 素数）。

\bigskip\hrule\bigskip

\subsubsection{42.4
代数基本定理（陈述与推论）}\label{ux4ee3ux6570ux57faux672cux5b9aux7406ux9648ux8ff0ux4e0eux63a8ux8bba}

\textbf{定理 42.4.1（代数基本定理）}：任何非常数多项式
\(f \in \mathbb{C}[x]\) 在 \(\mathbb{C}\) 中有根。

\textbf{注记
42.4.1}：代数基本定理的完整证明需要复分析（卷十）或拓扑学（卷十一）的工具。此处陈述结果，以便在多项式理论中使用。

\textbf{推论 42.4.2}：\(\mathbb{C}[x]\) 中的不可约多项式恰为一次多项式。

\textbf{推论 42.4.3}：\(\mathbb{R}[x]\)
中的不可约多项式为一次多项式或判别式为负的二次多项式。

\textbf{推论 42.4.4}：任何 \(n\) 次复系数多项式
\(f(x) = a_n x^n + \cdots + a_0\) 可分解为

\[
f(x) = a_n (x - \alpha_1)(x - \alpha_2) \cdots (x - \alpha_n)
\]

其中 \(\alpha_1, \ldots, \alpha_n \in \mathbb{C}\)（重根按重数计算）。

\bigskip\hrule\bigskip

\subsection{Ch 43 模论基础}\label{ch-43-ux6a21ux8bbaux57faux7840}

\subsubsection{43.1
模的定义与基本概念}\label{ux6a21ux7684ux5b9aux4e49ux4e0eux57faux672cux6982ux5ff5}

\textbf{定义 43.1.1（模）}：设 \(R\) 为环。一个\textbf{左
\(R\)-模}是一个阿贝尔群 \((M, +)\) 配备标量乘法
\(R \times M \to M\)，\((r, m) \mapsto rm\)，满足：

\begin{enumerate}
\def\labelenumi{\arabic{enumi}.}
\tightlist
\item
  \(r(m + n) = rm + rn\)
\item
  \((r + s)m = rm + sm\)
\item
  \((rs)m = r(sm)\)
\item
  \(1_R \cdot m = m\)（若 \(R\) 含幺）
\end{enumerate}

\textbf{评注 43.1.1}：当 \(R\) 为域时，\(R\)-模就是 \(R\)
上的向量空间。模论是线性代数的自然推广------将标量从域推广到环。

\textbf{定义
43.1.2（子模、商模、同态）}：与向量空间的子空间、商空间、线性映射的定义完全类似。

\textbf{定理
43.1.1（模的同构定理）}：与向量空间和群的同构定理完全类似，有第一、第二、第三同构定理。

\bigskip\hrule\bigskip

\subsubsection{43.2 自由模}\label{ux81eaux7531ux6a21}

\textbf{定义 43.2.1（线性无关与基）}：设 \(M\) 为
\(R\)-模。\(\{m_1, \ldots, m_n\} \subseteq M\)
称为\textbf{线性无关}的，若
\(\sum r_i m_i = 0 \implies r_i = 0\)（\(\forall i\)）。若
\(\{m_1, \ldots, m_n\}\) 线性无关且生成 \(M\)，则称为 \(M\)
的一组\textbf{基}。

\textbf{定义 43.2.2（自由模）}：有基的
\(R\)-模称为\textbf{自由模}。\(R^n = \{(r_1, \ldots, r_n) \mid r_i \in R\}\)
是秩为 \(n\) 的自由模的典范例子。

\textbf{注记 43.2.1}：与向量空间不同，不是所有模都有基。例如
\(\mathbb{Z}_6\) 作为 \(\mathbb{Z}\)-模没有基（因为 \(6 \cdot m = 0\)
对所有 \(m\) 成立，但标量 \(6 \neq 0\)）。

\bigskip\hrule\bigskip

\subsubsection{43.3
主理想整环上的有限生成模}\label{ux4e3bux7406ux60f3ux6574ux73afux4e0aux7684ux6709ux9650ux751fux6210ux6a21}

\textbf{定理 43.3.1（PID 上有限生成模的结构定理）}：设 \(R\) 为
PID，\(M\) 为有限生成 \(R\)-模。则

\[
M \cong R^r \oplus R/(d_1) \oplus R/(d_2) \oplus \cdots \oplus R/(d_k)
\]

其中 \(r \geq 0\)，\(d_1, d_2, \ldots, d_k \in R\)
为非零非单位元素，满足 \(d_1 \mid d_2 \mid \cdots \mid d_k\)。\(r\) 和
\(d_1, \ldots, d_k\)（在相伴意义下）由 \(M\) 唯一确定。

\textbf{证明}（概要）：设 \(M\) 由 \(n\) 个元素生成。则存在满同态
\(\varphi : R^n \to M\)。\(\ker \varphi\) 是 \(R^n\) 的子模。由 Smith
正规形理论，存在 \(R^n\) 的基使得 \(\ker \varphi\) 由
\(d_1 e_1, \ldots, d_k e_k\)
生成（\(d_1 \mid \cdots \mid d_k\)），其余基向量不在 \(\ker \varphi\)
中。因此
\(M \cong R^n / \ker \varphi \cong R/(d_1) \oplus \cdots \oplus R/(d_k) \oplus R^{n-k}\)。令
\(r = n - k\) 即得。\(\blacksquare\)

\textbf{推论
43.3.2（有限阿贝尔群的结构定理）}：任何有限生成阿贝尔群同构于

\[
\mathbb{Z}^r \oplus \mathbb{Z}_{d_1} \oplus \cdots \oplus \mathbb{Z}_{d_k}
\]

其中 \(d_1 \mid d_2 \mid \cdots \mid d_k\)。

\textbf{证明}：\(\mathbb{Z}\) 是 PID。有限生成阿贝尔群 = 有限生成
\(\mathbb{Z}\)-模。由定理 43.3.1 直接得到。\(\blacksquare\)

\textbf{推论 43.3.3（Jordan 标准形的模论证明）}：设 \(V\) 为 \(F\) 上的
\(n\) 维向量空间，\(T \in \mathcal{L}(V)\)。通过
\(f(x) \cdot v = f(T)v\) 使 \(V\) 成为 \(F[x]\)-模。\(F[x]\) 是
PID。由结构定理，\(V \cong \bigoplus F[x]/(p_i(x)^{e_i})\)（\(p_i\)
不可约）。每个 \(F[x]/(p_i^{e_i})\) 对应一个 Jordan 块。这给出了 Jordan
标准形的一个简洁证明。

\textbf{注记 43.3.1}：上述推论给出了卷六 Ch 35 中 Jordan
标准形的完整证明框架。

\bigskip\hrule\bigskip

\subsection{Ch 44 域扩张}\label{ch-44-ux57dfux6269ux5f20}

\subsubsection{44.1
代数扩张与超越扩张}\label{ux4ee3ux6570ux6269ux5f20ux4e0eux8d85ux8d8aux6269ux5f20}

\textbf{定义 44.1.1（域扩张）}：若 \(F \subseteq K\) 均为域，则称 \(K\)
为 \(F\) 的\textbf{域扩张}，记为 \(K/F\)。

\(K\) 可视为 \(F\) 上的向量空间，其维数称为扩张的\textbf{次数}，记为
\([K : F]\)。

\textbf{定义 44.1.2（代数元与超越元）}：设 \(K/F\)
为域扩张，\(\alpha \in K\)。

\begin{itemize}
\tightlist
\item
  若存在非零多项式 \(f \in F[x]\) 使得 \(f(\alpha) = 0\)，则称
  \(\alpha\) 为 \(F\) 上的\textbf{代数元}
\item
  否则称 \(\alpha\) 为 \(F\) 上的\textbf{超越元}
\end{itemize}

\textbf{定义 44.1.3（极小多项式）}：设 \(\alpha\) 为 \(F\)
上的代数元。在 \(F[x]\) 中满足 \(f(\alpha) = 0\)
的非零多项式中，次数最低的首一多项式称为 \(\alpha\) 在 \(F\)
上的\textbf{极小多项式}，记为 \(m_{\alpha, F}(x)\)。

\textbf{定理 44.1.1（极小多项式的基本性质）}：

\begin{enumerate}
\def\labelenumi{(\roman{enumi})}
\tightlist
\item
  \(m_{\alpha, F}(x)\) 在 \(F[x]\) 中不可约
\item
  若 \(f \in F[x]\) 满足 \(f(\alpha) = 0\)，则 \(m_{\alpha, F} \mid f\)
\item
  \(F(\alpha) \cong F[x]/(m_{\alpha, F}(x))\)，且
  \([F(\alpha) : F] = \deg m_{\alpha, F}\)
\end{enumerate}

\textbf{证明}：

\begin{enumerate}
\def\labelenumi{(\roman{enumi})}
\item
  若 \(m_{\alpha, F} = gh\)（\(g, h\) 次数更低），则
  \(g(\alpha)h(\alpha) = 0\)，故 \(g(\alpha) = 0\) 或
  \(h(\alpha) = 0\)，与 \(m_{\alpha, F}\) 的次数最低性矛盾。
\item
  多项式带余除法 + 次数最低性（与极小多项式的证明类似）。
\item
  定义同态 \(\varphi : F[x] \to F(\alpha)\) 为
  \(\varphi(f) = f(\alpha)\)。\(\ker \varphi = (m_{\alpha, F})\)。由第一同构定理，\(F[x]/(m_{\alpha, F}) \cong \operatorname{im} \varphi\)。由于
  \(m_{\alpha, F}\) 不可约，\((m_{\alpha, F})\) 是极大理想，故
  \(F[x]/(m_{\alpha, F})\) 是域。\(\operatorname{im} \varphi\) 包含
  \(F\) 和 \(\alpha\)，且是最小的这样的域，故
  \(\operatorname{im} \varphi = F(\alpha)\)。
\end{enumerate}

\(F(\alpha)\) 作为 \(F\)-向量空间，基为
\(\{1, \alpha, \ldots, \alpha^{n-1}\}\)（\(n = \deg m_{\alpha, F}\)），故
\([F(\alpha) : F] = n\)。\(\blacksquare\)

\textbf{定义 44.1.4（单扩张）}：若
\(K = F(\alpha)\)（\(\alpha \in K\)），则称 \(K/F\)
为\textbf{单扩张}。若 \(\alpha\) 是代数元，则称为\textbf{单代数扩张}。

\bigskip\hrule\bigskip

\subsubsection{44.2
有限扩张与代数扩张}\label{ux6709ux9650ux6269ux5f20ux4e0eux4ee3ux6570ux6269ux5f20}

\textbf{定义 44.2.1（有限扩张与代数扩张）}：

\begin{itemize}
\tightlist
\item
  若 \([K : F] < \infty\)，称 \(K/F\) 为\textbf{有限扩张}
\item
  若 \(K\) 中每个元素都是 \(F\) 上的代数元，称 \(K/F\)
  为\textbf{代数扩张}
\end{itemize}

\textbf{定理 44.2.1（次数塔公式）}：设 \(F \subseteq K \subseteq L\)
为域扩张。则

\[
[L : F] = [L : K] \cdot [K : F]
\]

\textbf{证明}：设 \(\{x_1, \ldots, x_m\}\) 为 \(K\) 在 \(F\)
上的基，\(\{y_1, \ldots, y_n\}\) 为 \(L\) 在 \(K\) 上的基。可以证明
\(\{x_i y_j \mid 1 \leq i \leq m, 1 \leq j \leq n\}\) 是 \(L\) 在 \(F\)
上的基。\(\blacksquare\)

\textbf{定理 44.2.2}：有限扩张 \(\implies\) 代数扩张。

\textbf{证明}：设 \([K : F] = n\)。对任意
\(\alpha \in K\)，\(1, \alpha, \ldots, \alpha^n\) 这 \(n+1\) 个元素在
\(n\) 维 \(F\)-向量空间 \(K\) 中必线性相关。故存在不全为零的
\(a_i \in F\) 使得 \(\sum a_i \alpha^i = 0\)，即 \(\alpha\)
是代数元。\(\blacksquare\)

\textbf{定理 44.2.3}：\(K/F\) 是有限扩张当且仅当
\(K = F(\alpha_1, \ldots, \alpha_n)\)，其中每个 \(\alpha_i\) 是 \(F\)
上的代数元。

\textbf{证明}：若 \(K/F\) 有限，取 \(K\) 在 \(F\) 上的基
\(\{\alpha_1, \ldots, \alpha_n\}\)，则
\(K = F(\alpha_1, \ldots, \alpha_n)\)，且每个 \(\alpha_i\) 是代数元。

反之，若 \(K = F(\alpha_1, \ldots, \alpha_n)\)（\(\alpha_i\) 代数），则
\([F(\alpha_1) : F] < \infty\)，\([F(\alpha_1, \alpha_2) : F(\alpha_1)] < \infty\)，\ldots，由塔公式乘积有限。\(\blacksquare\)

\bigskip\hrule\bigskip

\subsubsection{44.3 分裂域}\label{ux5206ux88c2ux57df}

\textbf{定义 44.3.1（分裂域）}：设 \(f \in F[x]\) 为非常数多项式。\(F\)
的扩域 \(K\) 称为 \(f\) 在 \(F\) 上的\textbf{分裂域}，若

\begin{enumerate}
\def\labelenumi{\arabic{enumi}.}
\tightlist
\item
  \(f\) 在 \(K[x]\) 中可分解为一次因式的乘积
\item
  \(K = F(\alpha_1, \ldots, \alpha_n)\)，其中 \(\alpha_i\) 为 \(f\) 在
  \(K\) 中的根
\end{enumerate}

\textbf{定理 44.3.1（分裂域的存在性）}：对任何 \(f \in F[x]\)，\(f\) 在
\(F\) 上的分裂域存在。

\textbf{证明}：对 \(\deg f\) 归纳。若 \(\deg f = 1\)，\(F\)
自身就是分裂域。

若 \(\deg f > 1\)，取 \(f\) 的一个不可约因子
\(p(x)\)。\(K_1 = F[x]/(p(x))\) 是 \(F\) 的扩域，且 \(p(x)\) 在 \(K_1\)
中有根 \(\alpha_1 = x + (p(x))\)。

在 \(K_1[x]\)
中，\(f(x) = (x - \alpha_1)g(x)\)（\(\deg g = \deg f - 1\)）。由归纳假设，\(g\)
在 \(K_1\) 上有分裂域 \(K\)。则 \(K\) 是 \(f\) 在 \(F\)
上的分裂域。\(\blacksquare\)

\textbf{定理 44.3.2（分裂域的唯一性）}：\(f \in F[x]\) 在 \(F\)
上的任何两个分裂域是同构的（\(F\)-同构）。

\textbf{证明}（概要）：对 \([K : F]\) 归纳。设 \(K_1, K_2\) 为 \(f\)
的两个分裂域。取 \(f\) 的不可约因子 \(p(x)\)，设
\(\alpha_1 \in K_1\)，\(\alpha_2 \in K_2\) 为 \(p(x)\) 的根。则
\(F(\alpha_1) \cong F[x]/(p(x)) \cong F(\alpha_2)\)。在 \(F(\alpha_1)\)
和 \(F(\alpha_2)\) 上分别考虑 \(f(x)/(x - \alpha_1)\) 和
\(f(x)/(x - \alpha_2)\)，由归纳可得同构可延拓到整个分裂域。\(\blacksquare\)

\bigskip\hrule\bigskip

\subsubsection{44.4 有限域}\label{ux6709ux9650ux57df}

\textbf{定义 44.4.1（有限域）}：元素个数有限的域称为\textbf{有限域}（或
Galois 域）。

\textbf{定理 44.4.1（有限域的结构）}：

\begin{enumerate}
\def\labelenumi{(\roman{enumi})}
\tightlist
\item
  有限域 \(F\) 的元素个数必为 \(p^n\)（\(p\)
  为素数，\(n \geq 1\)），其中 \(p\) 为 \(F\)
  的特征（\(\operatorname{char} F\)）
\item
  对任意素数 \(p\) 和正整数 \(n\)，存在唯一的（同构意义下）\(p^n\)
  元有限域，记为 \(\mathbb{F}_{p^n}\)
\item
  \(\mathbb{F}_{p^n}\) 是 \(x^{p^n} - x\) 在 \(\mathbb{F}_p\) 上的分裂域
\end{enumerate}

\textbf{证明}：

\begin{enumerate}
\def\labelenumi{(\roman{enumi})}
\item
  有限域 \(F\) 必包含其素子域
  \(\mathbb{F}_p\)（\(p = \operatorname{char} F\)）。\(F\) 是
  \(\mathbb{F}_p\) 上的有限维向量空间，设维数为 \(n\)，则
  \(|F| = p^n\)。
\item
  考虑多项式 \(x^{p^n} - x \in \mathbb{F}_p[x]\)。其形式导数
  \(p^n x^{p^n-1} - 1 = -1 \neq 0\)，故无重根。设 \(K\)
  为该多项式的分裂域。\(K\) 中满足 \(x^{p^n} = x\) 的元素恰为 \(p^n\)
  个（因为无重根），且它们构成 \(K\) 的子域（由特征 \(p\)
  的性质，\((a+b)^p = a^p + b^p\)）。因此该子域本身即为 \(p^n\)
  个元素的域。
\end{enumerate}

唯一性：任何 \(p^n\) 元域 \(F\) 中，非零元素乘法群 \(F^\times\) 的阶为
\(p^n-1\)，故所有元素满足 \(x^{p^n} - x = 0\)。因此 \(F\) 是
\(x^{p^n} - x\) 的分裂域，由分裂域的唯一性得证。\(\blacksquare\)

\textbf{定理 44.4.2（有限域的乘法群）}：\(\mathbb{F}_{p^n}^{\times}\) 是
\(p^n-1\) 阶循环群。

\textbf{证明}：\(\mathbb{F}_{p^n}^{\times}\)
是有限阿贝尔群。由有限阿贝尔群的结构定理，\(\mathbb{F}_{p^n}^{\times} \cong \mathbb{Z}_{d_1} \oplus \cdots \oplus \mathbb{Z}_{d_k}\)（\(d_1 \mid \cdots \mid d_k\)）。所有元素满足
\(x^{d_k} = 1\)。但 \(x^{d_k} - 1\) 在域中最多有 \(d_k\) 个根，而
\(|\mathbb{F}_{p^n}^{\times}| = p^n - 1\)，故
\(d_k = p^n - 1\)，\(k = 1\)。因此 \(\mathbb{F}_{p^n}^{\times}\)
是循环群。\(\blacksquare\)

\textbf{推论 44.4.3（有限域中的本原元）}：\(\mathbb{F}_{p^n}\)
中存在元素 \(\alpha\) 使得
\(\mathbb{F}_{p^n} = \mathbb{F}_p(\alpha)\)。事实上，\(\mathbb{F}_{p^n}^{\times}\)
的任何生成元都具有此性质。

\bigskip\hrule\bigskip

\subsubsection{44.5
尺规作图不可能性（概述）}\label{ux5c3aux89c4ux4f5cux56feux4e0dux53efux80fdux6027ux6982ux8ff0}

\textbf{定理 44.5.1（可构造数）}：实数 \(\alpha\)
称为\textbf{可构造数}（即可用尺规作图从单位线段构造），当且仅当存在域扩张链

\[
\mathbb{Q} = F_0 \subseteq F_1 \subseteq \cdots \subseteq F_k \subseteq \mathbb{R}
\]

使得 \([F_i : F_{i-1}] \leq 2\) 且 \(\alpha \in F_k\)。

\textbf{推论 44.5.2}：可构造数 \(\alpha\) 满足
\([\mathbb{Q}(\alpha) : \mathbb{Q}] = 2^m\)（\(m \geq 0\)）。

\textbf{定理 44.5.3（三大尺规作图不可能问题）}：

\begin{enumerate}
\def\labelenumi{(\roman{enumi})}
\tightlist
\item
  \textbf{倍立方体}：\(\sqrt[3]{2}\)
  不可构造（\([\mathbb{Q}(\sqrt[3]{2}) : \mathbb{Q}] = 3\)，不是 \(2\)
  的幂）
\item
  \textbf{三等分任意角}：\(\cos 20^\circ\) 不可构造（极小多项式为
  \(8x^3 - 6x + 1\)，次数为 \(3\)）
\item
  \textbf{化圆为方}：\(\sqrt{\pi}\) 不可构造（\(\pi\)
  是超越数，\([\mathbb{Q}(\pi) : \mathbb{Q}] = \infty\)）
\end{enumerate}

\textbf{证明}：

\begin{enumerate}
\def\labelenumi{(\roman{enumi})}
\item
  \(x^3 - 2\) 在 \(\mathbb{Q}[x]\) 中不可约（Eisenstein
  判别法，\(p = 2\)）。故
  \([\mathbb{Q}(\sqrt[3]{2}) : \mathbb{Q}] = 3 \neq 2^m\)。
\item
  \(\cos 3\theta = 4\cos^3 \theta - 3\cos \theta\)。令
  \(\theta = 20^\circ\)，则 \(\cos 60^\circ = 1/2\)，故
  \(4\cos^3 20^\circ - 3\cos 20^\circ = 1/2\)，即
  \(8\cos^3 20^\circ - 6\cos 20^\circ - 1 = 0\)。可证 \(8x^3 - 6x - 1\)
  在 \(\mathbb{Q}[x]\) 中不可约，故
  \([\mathbb{Q}(\cos 20^\circ) : \mathbb{Q}] = 3 \neq 2^m\)。
\item
  Lindemann 定理（1882 年）证明 \(\pi\) 是超越数，故
  \([\mathbb{Q}(\pi) : \mathbb{Q}] = \infty\)。\(\blacksquare\)
\end{enumerate}

\bigskip\hrule\bigskip

\textbf{注记 44.5.1}：Galois 理论的完整展开（包括方程可解性、\(n\)
次一般方程不可根式解等）将在卷十三（抽象代数
II）中详细讨论。本章仅建立域扩张的基本理论，为 Galois
理论做好准备。\(\blacksquare\)

\newpage

\section{卷八：实分析 II}\label{ux5377ux516bux5b9eux5206ux6790-ii}

\bigskip\hrule\bigskip

\begin{quote}
\textbf{前置知识}：本卷建立在卷一（数系、确界原理、集合论）和卷四（极限、连续、导数、Riemann
积分）的基础之上。读者应熟悉 \(\varepsilon\)-\(\delta\)
语言、确界原理、数列与函数的极限理论、以及 Riemann 积分的基本概念。
\end{quote}

\subsection{Ch 45
集合的测度}\label{ch-45-ux96c6ux5408ux7684ux6d4bux5ea6}

\subsubsection{45.1 引言：从 Riemann 积分到 Lebesgue
积分}\label{ux5f15ux8a00ux4ece-riemann-ux79efux5206ux5230-lebesgue-ux79efux5206}

Riemann
积分将定义域划分为小区间，在每个小区间上取函数值近似。这一方法具有直观性，但存在以下局限：

\begin{enumerate}
\def\labelenumi{\arabic{enumi}.}
\tightlist
\item
  Riemann 可积函数类不够广泛（Dirichlet 函数不可积）
\item
  极限与积分交换的条件苛刻（需要一致收敛）
\item
  缺乏完备性（Riemann 可积函数空间不完备）
\end{enumerate}

Lebesgue
测度与积分理论的核心思想是：\textbf{将值域划分为小区间，考察每个值对应的定义域集合的''大小''}。这需要先定义
\(\mathbb{R}^n\) 中子集的''大小''------即测度。

\bigskip\hrule\bigskip

\subsubsection{45.2 Lebesgue 外测度}\label{lebesgue-ux5916ux6d4bux5ea6}

\textbf{定义 45.2.1（开矩体）}：\(\mathbb{R}^n\) 中的\textbf{开矩体}
\(I\) 是形如

\[
I = (a_1, b_1) \times (a_2, b_2) \times \cdots \times (a_n, b_n)
\]

的集合，其中 \(a_i < b_i\)（\(i = 1, \ldots, n\)）。\(I\)
的\textbf{体积} \(|I|\) 定义为

\[
|I| = \prod_{i=1}^{n} (b_i - a_i)
\]

\textbf{定义 45.2.2（Lebesgue 外测度）}：设
\(E \subseteq \mathbb{R}^n\)。\(E\) 的 \textbf{Lebesgue 外测度}
\(m^*(E)\) 定义为

\[
m^*(E) = \inf\left\{\sum_{j=1}^{\infty} |I_j| \;\middle|\; E \subseteq \bigcup_{j=1}^{\infty} I_j, \; I_j \text{ 为开矩体}\right\}
\]

即用可数个开矩体覆盖 \(E\) 时，覆盖体积之和的下确界。

\textbf{定理 45.2.1（外测度的基本性质）}：

\begin{enumerate}
\def\labelenumi{(\roman{enumi})}
\tightlist
\item
  \textbf{非负性}：\(m^*(E) \geq 0\)，\(m^*(\varnothing) = 0\)
\item
  \textbf{单调性}：\(E_1 \subseteq E_2 \implies m^*(E_1) \leq m^*(E_2)\)
\item
  \textbf{次可数可加性}：\(\displaystyle m^*\left(\bigcup_{j=1}^{\infty} E_j\right) \leq \sum_{j=1}^{\infty} m^*(E_j)\)
\item
  \textbf{平移不变性}：\(m^*(E + h) = m^*(E)\)（\(h \in \mathbb{R}^n\)）
\item
  \textbf{矩体的外测度}：若 \(I\) 为矩体（开、闭或半开半闭），则
  \(m^*(I) = |I|\)
\end{enumerate}

\textbf{证明}：

\begin{enumerate}
\def\labelenumi{(\roman{enumi})}
\item
  外测度是非负数的下确界，故非负。\(\varnothing\) 可被体积
  \(\varepsilon\) 的矩体覆盖，\(\forall \varepsilon > 0\)，故
  \(m^*(\varnothing) = 0\)。
\item
  \(E_1\) 的覆盖也是 \(E_2\) 的覆盖，故 \(m^*(E_1)\) 的下确界不小于
  \(m^*(E_2)\) 的下确界。
\item
  对任意 \(\varepsilon > 0\)，对每个 \(E_j\) 取覆盖
  \(\{I_{jk}\}_{k=1}^{\infty}\) 使得
  \(\sum_k |I_{jk}| \leq m^*(E_j) + \varepsilon / 2^j\)。则
  \(\bigcup_j E_j \subseteq \bigcup_{j,k} I_{jk}\)，且
  \(\sum_{j,k} |I_{jk}| \leq \sum_j (m^*(E_j) + \varepsilon/2^j) = \sum_j m^*(E_j) + \varepsilon\)。取
  \(\varepsilon \to 0\) 即得。
\item
  开矩体平移后体积不变，覆盖也平移，故外测度不变。
\item
  需证 \(m^*(I) = |I|\)。\(|I| \leq m^*(I)\) 需要 \(I\)
  的任何覆盖体积之和
  \(\geq |I|\)，这涉及紧致性论证（开矩体覆盖紧致矩体时，可选出有限子覆盖，其体积之和
  \(\geq |I|\)）。\(m^*(I) \leq |I|\) 是平凡的（\(I\)
  自身就是一个覆盖）。\(\blacksquare\)
\end{enumerate}

\textbf{评注 45.2.1}：外测度在一般集合上不是可数可加的。例如，存在
\(E \subseteq [0, 1]\) 使得 \(m^*(E) > 0\) 但 \(m^*(E^c) = 1\) 且
\(m^*(E) + m^*(E^c) > 1 = m^*([0, 1])\)。这需要选择公理构造 Vitali
集（见 §45.5）。

\bigskip\hrule\bigskip

\subsubsection{45.3 Carathéodory
条件与可测集}\label{carathuxe9odory-ux6761ux4ef6ux4e0eux53efux6d4bux96c6}

\textbf{定义 45.3.1（Carathéodory
可测集）}：\(E \subseteq \mathbb{R}^n\) 称为 \textbf{Lebesgue
可测集}（简称可测集），若对任意 \(A \subseteq \mathbb{R}^n\)，

\[
m^*(A) = m^*(A \cap E) + m^*(A \cap E^c)
\]

即 \(E\) 将任何集合 \(A\) 的外测度''分割''得恰好。此时记
\(m(E) = m^*(E)\)，称为 \(E\) 的 \textbf{Lebesgue 测度}。

\textbf{定理 45.3.1（可测集的基本性质）}：

\begin{enumerate}
\def\labelenumi{(\roman{enumi})}
\tightlist
\item
  若 \(m^*(E) = 0\)，则 \(E\) 可测（零测集可测）
\item
  若 \(E\) 可测，则 \(E^c\) 可测
\item
  可测集的有限并（从而有限交）可测
\item
  可测集的可数并可测
\item
  开集和闭集可测
\end{enumerate}

\textbf{证明}：

\begin{enumerate}
\def\labelenumi{(\roman{enumi})}
\item
  对任意 \(A\)，\(A \cap E \subseteq E\)，故
  \(m^*(A \cap E) \leq m^*(E) = 0\)。\(A \cap E^c \subseteq A\)，故
  \(m^*(A \cap E^c) \leq m^*(A)\)。由次可加性，\(m^*(A) \leq m^*(A \cap E) + m^*(A \cap E^c) = 0 + m^*(A \cap E^c) \leq m^*(A)\)。故
  \(m^*(A) = m^*(A \cap E) + m^*(A \cap E^c)\)。
\item
  \(E\) 和 \(E^c\) 在 Carathéodory 条件中是对称的。
\item
  先证两个可测集的并可测。设 \(E_1, E_2\) 可测。对任意 \(A\)， \[
  \begin{aligned}
  m^*(A) &= m^*(A \cap E_1) + m^*(A \cap E_1^c) \\
  &= m^*(A \cap E_1) + m^*(A \cap E_1^c \cap E_2) + m^*(A \cap E_1^c \cap E_2^c)
  \end{aligned}
  \] 注意
  \(A \cap (E_1 \cup E_2) = (A \cap E_1) \cup (A \cap E_1^c \cap E_2)\)。由次可加性，
  \[
  m^*(A \cap (E_1 \cup E_2)) \leq m^*(A \cap E_1) + m^*(A \cap E_1^c \cap E_2)
  \] 而 \(A \cap (E_1 \cup E_2)^c = A \cap E_1^c \cap E_2^c\)。故 \[
  m^*(A \cap (E_1 \cup E_2)) + m^*(A \cap (E_1 \cup E_2)^c) \leq m^*(A \cap E_1) + m^*(A \cap E_1^c \cap E_2) + m^*(A \cap E_1^c \cap E_2^c) = m^*(A)
  \] 反向不等式由次可加性成立。故 \(E_1 \cup E_2\) 可测。
\item
  设 \(\{E_j\}_{j=1}^{\infty}\) 为可测集。可通过构造互不相交的可测集
  \(\{F_j\}\)（\(F_1 = E_1\)，\(F_j = E_j \setminus \bigcup_{i=1}^{j-1} E_i\)），使得
  \(\bigcup E_j = \bigcup F_j\)（不交并）。对任意 \(A\)，可证 \[
  m^*(A) = \sum_{j=1}^{k} m^*(A \cap F_j) + m^*\left(A \cap \left(\bigcup_{j=1}^{k} F_j\right)^c\right) \geq \sum_{j=1}^{k} m^*(A \cap F_j) + m^*\left(A \cap \left(\bigcup_{j=1}^{\infty} F_j\right)^c\right)
  \] 令 \(k \to \infty\)，由外测度的次可加性得 \[
  m^*(A) \geq \sum_{j=1}^{\infty} m^*(A \cap F_j) + m^*\left(A \cap \left(\bigcup_{j=1}^{\infty} F_j\right)^c\right) \geq m^*(A \cap \bigcup F_j) + m^*(A \cap (\bigcup F_j)^c)
  \] 故 \(\bigcup E_j\) 可测。
\item
  开矩体可测（由定义直接验证），开集是可数个开矩体的并，故可测。闭集是开集的补，也可测。\(\blacksquare\)
\end{enumerate}

\textbf{定理 45.3.2（可测集的可数可加性）}：设
\(\{E_j\}_{j=1}^{\infty}\) 为两两不交的可测集。则

\[
m\left(\bigcup_{j=1}^{\infty} E_j\right) = \sum_{j=1}^{\infty} m(E_j)
\]

\textbf{证明}：在定理 45.3.1(iv) 的证明中取 \(A = \bigcup E_j\)，则
\(m^*(A \cap (\bigcup E_j)^c) = m^*(\varnothing) = 0\)，且
\(A \cap F_j = F_j = E_j\)（当 \(E_j\) 互不相交时）。故
\(m^*(\bigcup E_j) = \sum m^*(E_j)\)。\(\blacksquare\)

\textbf{定理 45.3.3（测度的连续性）}：

\begin{enumerate}
\def\labelenumi{(\roman{enumi})}
\tightlist
\item
  \textbf{下连续性}：若 \(E_1 \subseteq E_2 \subseteq \cdots\) 可测，则
  \(m(\bigcup_{j=1}^{\infty} E_j) = \lim_{j \to \infty} m(E_j)\)
\item
  \textbf{上连续性}：若 \(E_1 \supseteq E_2 \supseteq \cdots\) 可测且
  \(m(E_1) < \infty\)，则
  \(m(\bigcap_{j=1}^{\infty} E_j) = \lim_{j \to \infty} m(E_j)\)
\end{enumerate}

\textbf{证明}：

\begin{enumerate}
\def\labelenumi{(\roman{enumi})}
\item
  令 \(F_1 = E_1\)，\(F_j = E_j \setminus E_{j-1}\)（\(j \geq 2\)）。则
  \(\{F_j\}\) 两两不交且
  \(\bigcup E_j = \bigcup F_j\)，\(E_k = \bigcup_{j=1}^{k} F_j\)。由可数可加性，
  \[
  m\left(\bigcup E_j\right) = \sum_{j=1}^{\infty} m(F_j) = \lim_{k \to \infty} \sum_{j=1}^{k} m(F_j) = \lim_{k \to \infty} m(E_k)
  \]
\item
  令 \(F_j = E_1 \setminus E_j\)。则
  \(F_1 \subseteq F_2 \subseteq \cdots\) 且
  \(\bigcup F_j = E_1 \setminus \bigcap E_j\)。由 (i)， \[
  m(E_1) - m\left(\bigcap E_j\right) = m\left(\bigcup F_j\right) = \lim_{j \to \infty} m(F_j) = \lim_{j \to \infty} (m(E_1) - m(E_j)) = m(E_1) - \lim_{j \to \infty} m(E_j)
  \] 故 \(m(\bigcap E_j) = \lim m(E_j)\)。\(\blacksquare\)
\end{enumerate}

\bigskip\hrule\bigskip

\subsubsection{45.4 Borel
集与正则性}\label{borel-ux96c6ux4e0eux6b63ux5219ux6027}

\textbf{定义 45.4.1（Borel \(\sigma\)-代数）}：\(\mathbb{R}^n\)
中由开集生成的最小
\(\sigma\)-代数（即包含所有开集且对可数并、可数交、补封闭的最小集族）称为
\textbf{Borel \(\sigma\)-代数}，其中的集合称为 \textbf{Borel 集}。

\textbf{定理 45.4.1}：所有 Borel 集都是 Lebesgue 可测的。

\textbf{证明}：由定理 45.3.1，开集可测，且可测集族构成
\(\sigma\)-代数（对补、可数并封闭），故包含所有开集的最小
\(\sigma\)-代数------即 Borel \(\sigma\)-代数------是 Lebesgue
可测集族的子族。\(\blacksquare\)

\textbf{定理 45.4.2（外正则性）}：对任意可测集
\(E \subseteq \mathbb{R}^n\)，

\[
m(E) = \inf\{m(U) \mid E \subseteq U, \; U \text{ 为开集}\}
\]

\textbf{定理 45.4.3（内正则性）}：对任意可测集
\(E \subseteq \mathbb{R}^n\)，

\[
m(E) = \sup\{m(K) \mid K \subseteq E, \; K \text{ 为紧集}\}
\]

\textbf{证明概要}：由外测度的定义，\(E\) 可被可数个开矩体覆盖，它们的并
\(U\)
是开集。紧致性方面，有界可测集可被紧集从内部逼近；无界可测集是递增有界可测集之并。\(\blacksquare\)

\bigskip\hrule\bigskip

\subsubsection{45.5 Vitali
集与非可测性}\label{vitali-ux96c6ux4e0eux975eux53efux6d4bux6027}

\textbf{定理 45.5.1（Vitali 集的存在性）}：存在 \(\mathbb{R}\)
的子集不是 Lebesgue 可测的。

\textbf{证明}（需要选择公理）：在 \([0, 1]\) 上定义等价关系
\(x \sim y \iff x - y \in \mathbb{Q}\)。每个等价类形如
\(x + \mathbb{Q}\)。由选择公理，从每个等价类中选取一个代表元，构成
Vitali 集 \(V \subseteq [0, 1]\)。

设 \(\{r_j\}_{j=1}^{\infty}\) 为 \([-1, 1] \cap \mathbb{Q}\) 的枚举。则
\(V + r_j\) 两两不交（若 \(v_1 + r_j = v_2 + r_k\)，则
\(v_1 - v_2 \in \mathbb{Q}\)，故 \(v_1 = v_2\)，\(r_j = r_k\)），且
\([0, 1] \subseteq \bigcup_j (V + r_j) \subseteq [-1, 2]\)。

若 \(V\) 可测，则
\(m(V + r_j) = m(V)\)（平移不变性）。由可数可加性，\(m(\bigcup_j (V + r_j)) = \sum_j m(V)\)。若
\(m(V) = 0\)，则 \(m(\bigcup (V + r_j)) = 0\)，但
\([0, 1] \subseteq \bigcup (V + r_j)\)，故 \(m([0, 1]) = 0\)，矛盾。若
\(m(V) > 0\)，则 \(\sum m(V) = \infty\)，但
\(\bigcup (V + r_j) \subseteq [-1, 2]\) 测度有限，矛盾。故 \(V\)
不可测。\(\blacksquare\)

\bigskip\hrule\bigskip

\subsection{Ch 46 可测函数}\label{ch-46-ux53efux6d4bux51fdux6570}

\subsubsection{46.1
可测函数的定义}\label{ux53efux6d4bux51fdux6570ux7684ux5b9aux4e49}

\textbf{定义 46.1.1（可测函数）}：设 \(E \subseteq \mathbb{R}^n\)
可测，\(f : E \to \overline{\mathbb{R}} = \mathbb{R} \cup \{\pm\infty\}\)。若对任意
\(\alpha \in \mathbb{R}\)，

\[
\{x \in E \mid f(x) > \alpha\}
\]

是可测集，则称 \(f\) 为 \(E\) 上的 \textbf{Lebesgue
可测函数}（简称可测函数）。

\textbf{定理 46.1.1（可测性的等价条件）}：以下条件等价：

\begin{enumerate}
\def\labelenumi{(\roman{enumi})}
\tightlist
\item
  \(\{f > \alpha\}\) 可测（\(\forall \alpha \in \mathbb{R}\)）
\item
  \(\{f \geq \alpha\}\) 可测（\(\forall \alpha \in \mathbb{R}\)）
\item
  \(\{f < \alpha\}\) 可测（\(\forall \alpha \in \mathbb{R}\)）
\item
  \(\{f \leq \alpha\}\) 可测（\(\forall \alpha \in \mathbb{R}\)）
\end{enumerate}

\textbf{证明}：(i) \(\implies\)
(ii)：\(\{f \geq \alpha\} = \bigcap_{k=1}^{\infty} \{f > \alpha - 1/k\}\)，可测集的可数交可测。

\begin{enumerate}
\def\labelenumi{(\roman{enumi})}
\setcounter{enumi}{1}
\item
  \(\implies\)
  (iii)：\(\{f < \alpha\} = E \setminus \{f \geq \alpha\}\)。
\item
  \(\implies\)
  (iv)：\(\{f \leq \alpha\} = \bigcap_{k=1}^{\infty} \{f < \alpha + 1/k\}\)。
\item
  \(\implies\)
  (i)：\(\{f > \alpha\} = E \setminus \{f \leq \alpha\}\)。\(\blacksquare\)
\end{enumerate}

\textbf{定理 46.1.2（可测函数的运算）}：设 \(f, g\) 为 \(E\)
上的可测函数（取值有限 a.e.），则

\begin{enumerate}
\def\labelenumi{(\roman{enumi})}
\tightlist
\item
  \(f + g\)，\(f - g\)，\(fg\) 可测（在定义有意义处）
\item
  \(f/g\) 可测（若 \(g \neq 0\) a.e.）
\item
  \(\max(f, g)\)，\(\min(f, g)\)，\(|f|\) 可测
\item
  若 \(\{f_k\}\) 为可测函数列，则
  \(\sup f_k\)，\(\inf f_k\)，\(\limsup f_k\)，\(\liminf f_k\) 可测
\end{enumerate}

\textbf{证明}：

\begin{enumerate}
\def\labelenumi{(\roman{enumi})}
\item
  \(\{f + g > \alpha\} = \bigcup_{r \in \mathbb{Q}} (\{f > r\} \cap \{g > \alpha - r\})\)，可测集的可数并可测。
\item
  \(\{\max(f, g) > \alpha\} = \{f > \alpha\} \cup \{g > \alpha\}\)。
\item
  \(\{\sup_k f_k > \alpha\} = \bigcup_k \{f_k > \alpha\}\)。\(\blacksquare\)
\end{enumerate}

\bigskip\hrule\bigskip

\subsubsection{46.2
简单函数及其逼近}\label{ux7b80ux5355ux51fdux6570ux53caux5176ux903cux8fd1}

\textbf{定义 46.2.1（简单函数）}：\(E\) 上的\textbf{简单函数}
\(\varphi\) 是形如

\[
\varphi(x) = \sum_{j=1}^{k} a_j \chi_{E_j}(x)
\]

的函数，其中 \(E_j\) 为两两不交的可测集，\(\chi_{E_j}\)
为特征函数，\(a_j \in \mathbb{R}\)。

\textbf{定理 46.2.1（简单函数逼近定理）}：设 \(f\) 为 \(E\)
上的非负可测函数。则存在递增的简单函数列 \(\{\varphi_k\}\) 使得
\(\varphi_k(x) \uparrow f(x)\)（逐点收敛）。

\textbf{证明}：对 \(k \in \mathbb{N}^+\)，将 \([0, k]\) 划分为
\(k \cdot 2^k\) 个长度为 \(1/2^k\) 的小区间。定义

\[
\varphi_k(x) = \begin{cases}
\frac{j-1}{2^k}, & \text{若 } \frac{j-1}{2^k} \leq f(x) < \frac{j}{2^k}, \; j = 1, \ldots, k \cdot 2^k \\
k, & \text{若 } f(x) \geq k
\end{cases}
\]

则 \(\varphi_k\) 是简单函数，\(\varphi_k \leq \varphi_{k+1}\)，且
\(\lim_{k \to \infty} \varphi_k(x) = f(x)\)（\(\forall x \in E\)）。\(\blacksquare\)

\textbf{推论 46.2.2}：任意可测函数 \(f\)
可表示为简单函数列的逐点极限（通过分解 \(f = f^+ - f^-\)）。

\bigskip\hrule\bigskip

\subsubsection{46.3 Egorov 定理与 Lusin
定理}\label{egorov-ux5b9aux7406ux4e0e-lusin-ux5b9aux7406}

\textbf{定理 46.3.1（Egorov 定理------几乎一致收敛）}：设
\(E \subseteq \mathbb{R}^n\) 可测且 \(m(E) < \infty\)，\(\{f_k\}\) 为
\(E\) 上的可测函数列，\(f_k \to f\) a.e.。则对任意
\(\varepsilon > 0\)，存在可测子集 \(A \subseteq E\) 使得
\(m(E \setminus A) < \varepsilon\) 且 \(f_k \to f\) 在 \(A\)
上\textbf{一致收敛}。

\textbf{证明}：对每个 \(j, k \in \mathbb{N}^+\)，定义

\[
E_{j,k} = \left\{x \in E \;\middle|\; \sup_{n \geq k} |f_n(x) - f(x)| > \frac{1}{j}\right\}
\]

由于 \(f_n \to f\)
a.e.，\(\forall j\)，\(\bigcap_{k=1}^{\infty} E_{j,k}\) 为零测集。由
\(m(E) < \infty\)
和测度的上连续性，\(\lim_{k \to \infty} m(E_{j,k}) = 0\)。

对 \(\varepsilon > 0\)，取 \(k_j\) 使得
\(m(E_{j, k_j}) < \varepsilon / 2^j\)。令
\(A = E \setminus \bigcup_{j=1}^{\infty} E_{j, k_j}\)。则

\[
m(E \setminus A) = m\left(\bigcup_{j=1}^{\infty} E_{j, k_j}\right) \leq \sum_{j=1}^{\infty} m(E_{j, k_j}) < \varepsilon
\]

在 \(A\) 上，对任意 \(j\)，当 \(n \geq k_j\) 时
\(|f_n - f| \leq 1/j\)，故 \(f_n \to f\) 一致。\(\blacksquare\)

\textbf{定理 46.3.2（Lusin 定理------几乎连续）}：设
\(E \subseteq \mathbb{R}^n\) 可测且 \(m(E) < \infty\)，\(f\) 为 \(E\) 上
a.e. 有限的可测函数。则对任意 \(\varepsilon > 0\)，存在紧集
\(K \subseteq E\) 使得 \(m(E \setminus K) < \varepsilon\) 且 \(f|_K\)
连续。

\textbf{证明概要}：先用简单函数逼近 \(f\)，对简单函数用 Egorov
定理和可测集的内正则性，再通过极限过程推广到一般可测函数。\(\blacksquare\)

\bigskip\hrule\bigskip

\subsection{Ch 47 Lebesgue 积分}\label{ch-47-lebesgue-ux79efux5206}

\subsubsection{47.1
非负简单函数的积分}\label{ux975eux8d1fux7b80ux5355ux51fdux6570ux7684ux79efux5206}

\textbf{定义 47.1.1（非负简单函数的积分）}：设
\(\varphi = \sum_{j=1}^{k} a_j \chi_{E_j}\) 为 \(E\)
上的非负简单函数（\(a_j \geq 0\)，\(E_j\) 两两不交可测）。定义
\(\varphi\) 的 Lebesgue 积分为

\[
\int_E \varphi \, dm = \sum_{j=1}^{k} a_j \, m(E_j)
\]

（约定 \(0 \cdot \infty = 0\)。）

\textbf{定理 47.1.1（简单函数积分的线性性与单调性）}：

\begin{enumerate}
\def\labelenumi{(\roman{enumi})}
\tightlist
\item
  \textbf{线性性}：\(\int_E (a\varphi + b\psi) = a\int_E \varphi + b\int_E \psi\)（\(a, b \geq 0\)）
\item
  \textbf{单调性}：若 \(\varphi \leq \psi\)，则
  \(\int_E \varphi \leq \int_E \psi\)
\end{enumerate}

\textbf{证明}：(i) 将 \(\varphi, \psi\)
用同一组可测集（它们的交集划分）表示，然后验证。(ii) 由
(i)，\(\psi - \varphi \geq 0\)，故
\(\int \psi - \int \varphi = \int (\psi - \varphi) \geq 0\)。\(\blacksquare\)

\bigskip\hrule\bigskip

\subsubsection{47.2
非负可测函数的积分}\label{ux975eux8d1fux53efux6d4bux51fdux6570ux7684ux79efux5206}

\textbf{定义 47.2.1（非负可测函数的积分）}：设 \(f\) 为 \(E\)
上的非负可测函数。定义

\[
\int_E f \, dm = \sup\left\{\int_E \varphi \, dm \;\middle|\; 0 \leq \varphi \leq f, \; \varphi \text{ 为简单函数}\right\}
\]

\textbf{定理 47.2.1（Chebyshev 不等式）}：设 \(f \geq 0\) 可测。则对任意
\(\lambda > 0\)，

\[
m(\{x \in E \mid f(x) \geq \lambda\}) \leq \frac{1}{\lambda} \int_E f \, dm
\]

\textbf{证明}：令 \(A = \{f \geq \lambda\}\)。则
\(f \geq \lambda \chi_A\)。由积分的单调性，\(\int_E f \geq \lambda \int_E \chi_A = \lambda m(A)\)。\(\blacksquare\)

\textbf{推论 47.2.2}：若 \(\int_E f = 0\)，则 \(f = 0\) a.e.。

\textbf{证明}：由 Chebyshev 不等式，对任意
\(n\)，\(m(\{f \geq 1/n\}) \leq n \int_E f = 0\)。故
\(\{f > 0\} = \bigcup_{n=1}^{\infty} \{f \geq 1/n\}\)
为零测集。\(\blacksquare\)

\bigskip\hrule\bigskip

\subsubsection{47.3
单调收敛定理}\label{ux5355ux8c03ux6536ux655bux5b9aux7406}

\textbf{定理 47.3.1（Lebesgue 单调收敛定理）}：设 \(\{f_k\}\) 为 \(E\)
上的非负可测函数列，满足 \(0 \leq f_1 \leq f_2 \leq \cdots\)，且
\(f_k \to f\)（逐点）。则

\[
\lim_{k \to \infty} \int_E f_k \, dm = \int_E f \, dm
\]

即极限与积分可交换（在非负递增情形下）。

\textbf{证明}：由于 \(f_k \leq f\)，\(\int f_k \leq \int f\)，故
\(\lim_{k \to \infty} \int f_k \leq \int f\)。需证反向不等式。

取任意简单函数 \(0 \leq \varphi \leq f\) 和 \(0 < \alpha < 1\)。定义
\(E_k = \{x \mid f_k(x) \geq \alpha \varphi(x)\}\)。则
\(E_1 \subseteq E_2 \subseteq \cdots\) 且 \(\bigcup E_k = E\)（因为
\(f_k \uparrow f\) 且 \(\alpha \varphi < f\)）。

由测度的下连续性，\(m(E_k) \uparrow m(E)\)。且
\(\int_E f_k \geq \int_{E_k} f_k \geq \alpha \int_{E_k} \varphi\)。

对简单函数
\(\varphi = \sum a_j \chi_{F_j}\)，\(\int_{E_k} \varphi = \sum a_j m(F_j \cap E_k) \uparrow \sum a_j m(F_j) = \int_E \varphi\)（\(k \to \infty\)）。

因此 \(\lim_{k \to \infty} \int_E f_k \geq \alpha \int_E \varphi\)。令
\(\alpha \to 1\) 得 \(\lim \int f_k \geq \int \varphi\)。再对
\(\varphi \leq f\) 取上确界，得
\(\lim \int f_k \geq \int f\)。\(\blacksquare\)

\textbf{推论 47.3.2（非负可测函数级数的逐项积分）}：若 \(\{g_k\}\)
为非负可测函数列，则

\[
\int_E \sum_{k=1}^{\infty} g_k \, dm = \sum_{k=1}^{\infty} \int_E g_k \, dm
\]

\textbf{证明}：令
\(f_k = \sum_{j=1}^{k} g_j\)，应用单调收敛定理。\(\blacksquare\)

\bigskip\hrule\bigskip

\subsubsection{47.4 Fatou 引理}\label{fatou-ux5f15ux7406}

\textbf{定理 47.4.1（Fatou 引理）}：设 \(\{f_k\}\) 为 \(E\)
上的非负可测函数列。则

\[
\int_E \liminf_{k \to \infty} f_k \, dm \leq \liminf_{k \to \infty} \int_E f_k \, dm
\]

\textbf{证明}：令 \(g_k = \inf_{j \geq k} f_j\)。则
\(0 \leq g_1 \leq g_2 \leq \cdots\)，\(g_k \uparrow \liminf f_k\)。由单调收敛定理，\(\int \liminf f_k = \lim_{k \to \infty} \int g_k\)。

又 \(g_k \leq f_j\)（\(\forall j \geq k\)），故
\(\int g_k \leq \inf_{j \geq k} \int f_j\)。因此

\[
\int \liminf f_k = \lim_{k \to \infty} \int g_k \leq \lim_{k \to \infty} \inf_{j \geq k} \int f_j = \liminf_{k \to \infty} \int f_k
\]

\(\blacksquare\)

\bigskip\hrule\bigskip

\subsubsection{47.5 一般可测函数的积分与 Lebesgue
控制收敛定理}\label{ux4e00ux822cux53efux6d4bux51fdux6570ux7684ux79efux5206ux4e0e-lebesgue-ux63a7ux5236ux6536ux655bux5b9aux7406}

\textbf{定义 47.5.1（可积函数）}：设 \(f\) 为 \(E\) 上的可测函数。将
\(f\) 分解为正部和负部：\(f = f^+ - f^-\)，其中
\(f^+ = \max(f, 0)\)，\(f^- = \max(-f, 0)\)。

若 \(\int_E f^+ < \infty\) 且 \(\int_E f^- < \infty\)，则称 \(f\) 在
\(E\) 上 \textbf{Lebesgue 可积}，定义

\[
\int_E f \, dm = \int_E f^+ \, dm - \int_E f^- \, dm
\]

\textbf{定理 47.5.1（Lebesgue 控制收敛定理）}：设 \(\{f_k\}\) 为 \(E\)
上的可测函数列，满足 \(f_k \to f\) a.e.。若存在可积函数 \(g\) 使得
\(|f_k| \leq g\) a.e.（\(\forall k\)），则 \(f\) 可积且

\[
\lim_{k \to \infty} \int_E f_k \, dm = \int_E f \, dm
\]

\textbf{证明}：由 \(|f_k| \leq g\) 及 \(g\) 可积，\(f_k\) 可积。对
\(f\)，\(|f| = \lim |f_k| \leq g\) a.e.，故 \(f\) 也可积。

考虑 \(g + f_k \geq 0\) 和 \(g - f_k \geq 0\)。由 Fatou 引理：

\[
\int_E (g + f) \leq \liminf_{k \to \infty} \int_E (g + f_k) = \int_E g + \liminf_{k \to \infty} \int_E f_k
\]

故 \(\int_E f \leq \liminf_{k \to \infty} \int_E f_k\)。

同理对 \(g - f_k\)：

\[
\int_E (g - f) \leq \liminf_{k \to \infty} \int_E (g - f_k) = \int_E g - \limsup_{k \to \infty} \int_E f_k
\]

故 \(\int_E f \geq \limsup_{k \to \infty} \int_E f_k\)。

因此 \(\liminf \int f_k = \limsup \int f_k = \int f\)，即
\(\lim \int f_k = \int f\)。\(\blacksquare\)

\bigskip\hrule\bigskip

\subsubsection{47.6 Lebesgue 积分与 Riemann
积分的比较}\label{lebesgue-ux79efux5206ux4e0e-riemann-ux79efux5206ux7684ux6bd4ux8f83}

\textbf{定理 47.6.1}：设 \(f\) 在 \([a, b]\) 上 Riemann 可积。则 \(f\)
在 \([a, b]\) 上 Lebesgue 可积，且两种积分值相等：

\[
\int_{[a, b]} f \, dm = \int_a^b f(x) \, dx
\]

\textbf{证明}：Riemann 可积函数必有界，且其不连续点构成零测集（Lebesgue
判别法）。有界可测函数在有限测度集上 Lebesgue 可积。

比较两种积分：Riemann 积分用阶梯函数逼近，Lebesgue
积分用简单函数逼近。由 Riemann 可积性，Riemann 上下积分相等，且等于
Lebesgue 积分。\(\blacksquare\)

\textbf{定理 47.6.2（Lebesgue 可积函数类更广）}：存在 Lebesgue 可积但
Riemann 不可积的函数。

\textbf{例}：Dirichlet 函数 \(\chi_{\mathbb{Q} \cap [0, 1]}\) 在
\([0, 1]\) 上 Lebesgue 可积（值为 \(0\)，因为 \(\mathbb{Q} \cap [0, 1]\)
可数故零测），但 Riemann 不可积。

\bigskip\hrule\bigskip

\subsection{\texorpdfstring{Ch 48 \(L^p\)
空间}{Ch 48 L\^{}p 空间}}\label{ch-48-lp-ux7a7aux95f4}

\subsubsection{\texorpdfstring{48.1 \(L^p\)
空间的定义}{48.1 L\^{}p 空间的定义}}\label{lp-ux7a7aux95f4ux7684ux5b9aux4e49}

\textbf{定义 48.1.1（ \(L^p\) 范数）}：设 \(E \subseteq \mathbb{R}^n\)
可测，\(1 \leq p < \infty\)。\(f\) 的 \(L^p\) \textbf{范数}定义为

\[
\|f\|_p = \left(\int_E |f|^p \, dm\right)^{1/p}
\]

\(L^p(E)\) 是所有满足 \(\|f\|_p < \infty\) 的可测函数 \(f\) 的集合（模去
a.e. 相等的函数，即 \(f = g\) a.e. 视为同一元素）。

\textbf{定义 48.1.2（ \(L^\infty\) 范数）}：\(f\) 的 \(L^\infty\)
\textbf{范数}（本质上确界）定义为

\[
\|f\|_\infty = \inf\{M \geq 0 \mid |f(x)| \leq M \text{ a.e.}\}
\]

\(L^\infty(E)\) 是所有满足 \(\|f\|_\infty < \infty\) 的可测函数的集合。

\textbf{定理 48.1.1}：\(L^p(E)\)
是向量空间（\(1 \leq p \leq \infty\)）。

\textbf{证明}：由
\(|f + g|^p \leq 2^p(|f|^p + |g|^p)\)（\(p < \infty\)）或
\(\|f + g\|_\infty \leq \|f\|_\infty + \|g\|_\infty\)（\(p = \infty\)）。\(\blacksquare\)

\bigskip\hrule\bigskip

\subsubsection{48.2 Hölder 不等式与 Minkowski
不等式}\label{huxf6lder-ux4e0dux7b49ux5f0fux4e0e-minkowski-ux4e0dux7b49ux5f0f}

\textbf{定理 48.2.1（Young 不等式）}：设 \(a, b \geq 0\)，\(p, q > 1\)
满足 \(\frac{1}{p} + \frac{1}{q} = 1\)。则

\[
ab \leq \frac{a^p}{p} + \frac{b^q}{q}
\]

\textbf{证明}：由凸函数 \(t \mapsto e^t\)
的性质，\(\frac{a^p}{p} + \frac{b^q}{q} = e^{\ln a^p - \ln p} + e^{\ln b^q - \ln q} \geq e^{\frac{1}{p}\ln a^p + \frac{1}{q}\ln b^q} = ab\)。等号成立当且仅当
\(a^p = b^q\)。\(\blacksquare\)

\textbf{定理 48.2.2（Hölder 不等式）}：设
\(1 < p, q < \infty\)，\(\frac{1}{p} + \frac{1}{q} = 1\)。若
\(f \in L^p\)，\(g \in L^q\)，则 \(fg \in L^1\) 且

\[
\|fg\|_1 \leq \|f\|_p \|g\|_q
\]

\textbf{证明}：若 \(\|f\|_p = 0\) 或
\(\|g\|_q = 0\)，不等式平凡成立。否则，令
\(a = |f(x)|/\|f\|_p\)，\(b = |g(x)|/\|g\|_q\)。由 Young 不等式，

\[
\frac{|f(x)g(x)|}{\|f\|_p \|g\|_q} \leq \frac{1}{p} \frac{|f(x)|^p}{\|f\|_p^p} + \frac{1}{q} \frac{|g(x)|^q}{\|g\|_q^q}
\]

两边积分得
\(\frac{\|fg\|_1}{\|f\|_p \|g\|_q} \leq \frac{1}{p} + \frac{1}{q} = 1\)。\(\blacksquare\)

\textbf{定理 48.2.3（Minkowski 不等式）}：设
\(1 \leq p < \infty\)，\(f, g \in L^p\)。则

\[
\|f + g\|_p \leq \|f\|_p + \|g\|_p
\]

\textbf{证明}：\(p = 1\) 时由三角不等式直接得到。\(p > 1\) 时，

\[
\begin{aligned}
\|f + g\|_p^p &= \int |f + g|^p \leq \int |f + g|^{p-1}|f| + \int |f + g|^{p-1}|g| \\
&\leq \left(\int |f + g|^{(p-1)q}\right)^{1/q} \left(\int |f|^p\right)^{1/p} + \left(\int |f + g|^{(p-1)q}\right)^{1/q} \left(\int |g|^p\right)^{1/p} \\
&= \|f + g\|_p^{p/q} (\|f\|_p + \|g\|_p)
\end{aligned}
\]

其中第二步用了 Hölder
不等式（\(\frac{1}{p} + \frac{1}{q} = 1\)，\((p-1)q = p\)）。两边除以
\(\|f + g\|_p^{p/q}\)（若为 \(0\) 则平凡），得
\(\|f + g\|_p \leq \|f\|_p + \|g\|_p\)。\(\blacksquare\)

\textbf{推论 48.2.4}：\(\|\cdot\|_p\) 是 \(L^p(E)\) 上的范数。

\bigskip\hrule\bigskip

\subsubsection{\texorpdfstring{48.3 \(L^p\) 空间的完备性（Riesz-Fischer
定理）}{48.3 L\^{}p 空间的完备性（Riesz-Fischer 定理）}}\label{lp-ux7a7aux95f4ux7684ux5b8cux5907ux6027riesz-fischer-ux5b9aux7406}

\textbf{定理 48.3.1（Riesz-Fischer
定理）}：\(L^p(E)\)（\(1 \leq p \leq \infty\)）关于 \(\|\cdot\|_p\)
范数是完备的赋范向量空间（即 Banach 空间）。

\textbf{证明}（\(1 \leq p < \infty\) 的情形）：设 \(\{f_k\}\) 为 \(L^p\)
中的 Cauchy 列。选取子列 \(\{f_{k_j}\}\) 使得
\(\|f_{k_{j+1}} - f_{k_j}\|_p < 2^{-j}\)。

定义
\(g_n = \sum_{j=1}^{n} |f_{k_{j+1}} - f_{k_j}|\)，\(g = \sum_{j=1}^{\infty} |f_{k_{j+1}} - f_{k_j}|\)。由
Minkowski 不等式，

\[
\|g_n\|_p \leq \sum_{j=1}^{n} \|f_{k_{j+1}} - f_{k_j}\|_p < \sum_{j=1}^{n} 2^{-j} < 1
\]

由 Fatou
引理，\(\|g\|_p^p = \int \liminf g_n^p \leq \liminf \int g_n^p \leq 1\)，故
\(g \in L^p\)。特别地，\(g(x) < \infty\) a.e.。

因此 \(\sum (f_{k_{j+1}} - f_{k_j})\) a.e. 绝对收敛，故 \(f_{k_j}\) a.e.
收敛到某个 \(f\)。定义
\(f(x) = \lim_{j \to \infty} f_{k_j}(x)\)（a.e.），其余处为 \(0\)。

由 Fatou
引理，\(\|f - f_k\|_p^p = \int \liminf_{j \to \infty} |f_{k_j} - f_k|^p \leq \liminf_{j \to \infty} \|f_{k_j} - f_k\|_p^p\)。由于
\(\{f_k\}\) 是 Cauchy 列，右端对足够大的 \(k\) 可任意小，故
\(f_k \to f\) 在 \(L^p\) 中。\(\blacksquare\)

\bigskip\hrule\bigskip

\subsubsection{48.4 稠密性定理}\label{ux7a20ux5bc6ux6027ux5b9aux7406}

\textbf{定理 48.4.1（连续函数在 \(L^p\) 中的稠密性）}：设
\(1 \leq p < \infty\)，\(E \subseteq \mathbb{R}^n\) 可测且
\(m(E) < \infty\)。则 \(C_c(E)\)（紧支集连续函数）在 \(L^p(E)\) 中稠密。

\textbf{证明概要}：

\begin{enumerate}
\def\labelenumi{\arabic{enumi}.}
\tightlist
\item
  特征函数 \(\chi_A\)（\(A\)
  可测，\(m(A) < \infty\)）可由连续函数逼近（由 Lusin 定理，存在紧集
  \(K \subseteq A\) 使得 \(m(A \setminus K)\) 很小，构造连续函数在 \(K\)
  上为 \(1\)，在 \(K\) 外迅速衰减）。
\item
  简单函数是特征函数的线性组合，故也可由连续函数逼近。
\item
  任意 \(f \in L^p\)
  可由简单函数逼近（简单函数逼近定理），故由连续函数逼近。\(\blacksquare\)
\end{enumerate}

\bigskip\hrule\bigskip

\subsubsection{\texorpdfstring{48.5 \(L^2\)
空间与内积}{48.5 L\^{}2 空间与内积}}\label{l2-ux7a7aux95f4ux4e0eux5185ux79ef}

\textbf{定理 48.5.1}：\(L^2(E)\) 配备内积
\(\langle f, g \rangle = \int_E fg \, dm\) 是 Hilbert 空间。

\textbf{证明}：由 Hölder
不等式（\(p = q = 2\)），\(\langle f, g \rangle\)
有限。内积的四条公理直接验证。完备性由 Riesz-Fischer
定理（\(p = 2\)）。\(\blacksquare\)

\bigskip\hrule\bigskip

\subsection{Ch 49
微分与积分的联系}\label{ch-49-ux5faeux5206ux4e0eux79efux5206ux7684ux8054ux7cfb}

\subsubsection{49.1
有界变差函数}\label{ux6709ux754cux53d8ux5deeux51fdux6570}

\textbf{定义 49.1.1（有界变差函数）}：设
\(f : [a, b] \to \mathbb{R}\)。\(f\) 在 \([a, b]\)
上的\textbf{全变差}定义为

\[
V_a^b(f) = \sup\left\{\sum_{i=1}^{n} |f(x_i) - f(x_{i-1})| \;\middle|\; a = x_0 < x_1 < \cdots < x_n = b\right\}
\]

若 \(V_a^b(f) < \infty\)，则称 \(f\) 为 \([a, b]\)
上的\textbf{有界变差函数}。全体有界变差函数记为 \(BV([a, b])\)。

\textbf{定理 49.1.1（Jordan 分解）}：\(f \in BV([a, b])\) 当且仅当 \(f\)
可表示为两个单调递增函数的差。

\textbf{证明}：

（\(\Leftarrow\)）单调函数的全变差有限，两个有界变差函数的差也有界变差。

（\(\Rightarrow\)）定义
\(g(x) = V_a^x(f)\)（全变差函数），\(h(x) = g(x) - f(x)\)。\(g\)
递增（全变差随区间增长）。需证 \(h\) 也递增。对
\(x < y\)，\(h(y) - h(x) = V_x^y(f) - (f(y) - f(x)) \geq 0\)（因为全变差至少是
\(|f(y) - f(x)|\)）。故 \(f = g - h\)
是两个单调递增函数的差。\(\blacksquare\)

\textbf{推论 49.1.2}：\(f \in BV([a, b])\) 几乎处处可导。

\textbf{证明}：由 Jordan 分解，\(f\)
是两个单调函数的差。单调函数几乎处处可导（Lebesgue
微分定理，见下节），故 \(f\) 也几乎处处可导。\(\blacksquare\)

\bigskip\hrule\bigskip

\subsubsection{49.2 Lebesgue
微分定理}\label{lebesgue-ux5faeux5206ux5b9aux7406}

\textbf{定理 49.2.1（Lebesgue 微分定理）}：设
\(f : [a, b] \to \mathbb{R}\) 单调递增。则 \(f\) 在 \([a, b]\)
上几乎处处可导，且 \(f' \in L^1([a, b])\)，满足

\[
\int_a^b f'(x) \, dx \leq f(b) - f(a)
\]

\textbf{注记 49.2.1}：完整证明需要 Hardy-Littlewood 极大函数和 Vitali
覆盖引理，在此从略。核心思想是：单调函数的不连续点至多可数，且其导数
\(f'\) 是差商极限的测度论性质。

\textbf{证明思路}：利用 Hardy-Littlewood 极大函数的弱 \((1,1)\)
型估计。对任意 \(\alpha > 0\)，令
\(E_\alpha = \{x : |f(x) - f_r(x)| > \alpha \text{ 无穷次}\}\)。通过极大函数控制
\(|f - f_r|\)，结合 \(L^1\) 函数的几乎处处连续性（Lusin 定理）和 Vitali
覆盖引理，可得 \(m(E_\alpha) = 0\)，即几乎处处收敛。详见：Stein,
\emph{Singular Integrals}, Ch I.3。

\bigskip\hrule\bigskip

\subsubsection{49.3
绝对连续函数}\label{ux7eddux5bf9ux8fdeux7eedux51fdux6570}

\textbf{定义 49.3.1（绝对连续函数）}：\(f : [a, b] \to \mathbb{R}\)
称为\textbf{绝对连续}的，若对任意 \(\varepsilon > 0\)，存在
\(\delta > 0\)，使得对 \([a, b]\) 的任意有限个互不相交的子区间
\(\{(a_k, b_k)\}\)，若 \(\sum (b_k - a_k) < \delta\)，则
\(\sum |f(b_k) - f(a_k)| < \varepsilon\)。

\textbf{定理 49.3.1（绝对连续函数的性质）}：

\begin{enumerate}
\def\labelenumi{(\roman{enumi})}
\tightlist
\item
  绝对连续函数是一致连续的（且是有界变差的）
\item
  Lipschitz 连续函数是绝对连续的
\item
  绝对连续函数几乎处处可导
\end{enumerate}

\textbf{定理 49.3.2（微积分基本定理------Lebesgue
版）}：\(f : [a, b] \to \mathbb{R}\) 绝对连续当且仅当存在
\(g \in L^1([a, b])\) 使得

\[
f(x) - f(a) = \int_a^x g(t) \, dt, \quad \forall x \in [a, b]
\]

此时 \(f' = g\) a.e.。

\textbf{证明}：

（\(\Leftarrow\)）若
\(f(x) = f(a) + \int_a^x g(t) dt\)，由积分的绝对连续性（Lebesgue
积分的绝对连续性：对
\(g \in L^1\)，\(\forall \varepsilon > 0\)，\(\exists \delta > 0\)，当
\(m(A) < \delta\) 时 \(\int_A |g| < \varepsilon\)），\(f\)
是绝对连续的。

（\(\Rightarrow\)）\(f\) 绝对连续 \(\implies\) \(f \in BV\)。由 Jordan
分解，\(f = f_1 - f_2\)（\(f_1, f_2\) 递增）。可以证明 \(f_1, f_2\)
也是绝对连续的。由 Lebesgue 微分定理，\(f'\) a.e. 存在且
\(f' \in L^1\)。定义 \(h(x) = \int_a^x f'(t) dt\)。则 \(h\) 绝对连续且
\(h' = f'\) a.e.。令 \(\varphi = f - h\)。\(\varphi\) 绝对连续且
\(\varphi' = 0\) a.e.。可以证明 \(\varphi\) 为常数（需要 Vitali
覆盖引理），故
\(f(x) - f(a) = h(x) - h(a) = \int_a^x f'(t) dt\)。\(\blacksquare\)

\bigskip\hrule\bigskip

\subsection{Ch 50 Fourier 级数}\label{ch-50-fourier-ux7ea7ux6570}

\subsubsection{50.1 正交系与 Bessel
不等式}\label{ux6b63ux4ea4ux7cfbux4e0e-bessel-ux4e0dux7b49ux5f0f}

\textbf{定义 50.1.1（ \(L^2\) 中的规范正交系）}：\(L^2([-\pi, \pi])\)
中的函数列 \(\{\varphi_k\}_{k=1}^{\infty}\) 称为\textbf{规范正交系}，若

\[
\langle \varphi_i, \varphi_j \rangle = \int_{-\pi}^{\pi} \varphi_i(x) \overline{\varphi_j(x)} \, dx = \delta_{ij}
\]

\textbf{例 50.1.1}：三角函数系

\[
\frac{1}{\sqrt{2\pi}}, \frac{\cos x}{\sqrt{\pi}}, \frac{\sin x}{\sqrt{\pi}}, \frac{\cos 2x}{\sqrt{\pi}}, \frac{\sin 2x}{\sqrt{\pi}}, \ldots
\]

是 \(L^2([-\pi, \pi])\)
中的规范正交系（实内积：\(\overline{\varphi_j} = \varphi_j\)）。

\textbf{定理 50.1.1（Bessel 不等式）}：设 \(\{\varphi_k\}\) 为
\(L^2(E)\) 中的规范正交系，\(f \in L^2(E)\)。则

\[
\sum_{k=1}^{\infty} |\langle f, \varphi_k \rangle|^2 \leq \|f\|_2^2
\]

\textbf{证明}：对任意 \(n\)，令
\(s_n = \sum_{k=1}^{n} \langle f, \varphi_k \rangle \varphi_k\)。则

\[
0 \leq \|f - s_n\|_2^2 = \|f\|_2^2 - 2\langle f, s_n \rangle + \|s_n\|_2^2 = \|f\|_2^2 - 2\sum_{k=1}^{n} |\langle f, \varphi_k \rangle|^2 + \sum_{k=1}^{n} |\langle f, \varphi_k \rangle|^2 = \|f\|_2^2 - \sum_{k=1}^{n} |\langle f, \varphi_k \rangle|^2
\]

因此
\(\sum_{k=1}^{n} |\langle f, \varphi_k \rangle|^2 \leq \|f\|_2^2\)。令
\(n \to \infty\) 即得。\(\blacksquare\)

\bigskip\hrule\bigskip

\subsubsection{\texorpdfstring{50.2 Fourier 级数的 \(L^2\)
收敛}{50.2 Fourier 级数的 L\^{}2 收敛}}\label{fourier-ux7ea7ux6570ux7684-l2-ux6536ux655b}

\textbf{定义 50.2.1（Fourier 系数与 Fourier 级数）}：设
\(f \in L^1([-\pi, \pi])\)。\(f\) 的 Fourier 系数定义为

\[
a_0 = \frac{1}{2\pi} \int_{-\pi}^{\pi} f(x) \, dx, \quad a_n = \frac{1}{\pi} \int_{-\pi}^{\pi} f(x) \cos nx \, dx, \quad b_n = \frac{1}{\pi} \int_{-\pi}^{\pi} f(x) \sin nx \, dx
\]

\(f\) 的 Fourier 级数为

\[
f(x) \sim \frac{a_0}{2} + \sum_{n=1}^{\infty} (a_n \cos nx + b_n \sin nx)
\]

\textbf{定理 50.2.1（ \(L^2\) 中 Fourier 级数的收敛性）}：设
\(f \in L^2([-\pi, \pi])\)。则 \(f\) 的 Fourier 级数在 \(L^2\)
范数下收敛到 \(f\)。即

\[
\lim_{N \to \infty} \left\|f - \left(\frac{a_0}{2} + \sum_{n=1}^{N} (a_n \cos nx + b_n \sin nx)\right)\right\|_2 = 0
\]

等价地，Parseval 恒等式成立：

\[
\frac{1}{\pi} \int_{-\pi}^{\pi} |f(x)|^2 \, dx = \frac{a_0^2}{2} + \sum_{n=1}^{\infty} (a_n^2 + b_n^2)
\]

\textbf{证明概要}：三角函数系在 \(L^2([-\pi, \pi])\)
中是完备的规范正交系（这是 Weierstrass 逼近定理的推论------三角多项式在
\(L^2\) 中稠密）。由 Hilbert 空间理论，规范正交系完备当且仅当 Parseval
恒等式成立。\(\blacksquare\)

\bigskip\hrule\bigskip

\subsubsection{50.3 Dirichlet 核与 Fejér
核}\label{dirichlet-ux6838ux4e0e-fejuxe9r-ux6838}

\textbf{定义 50.3.1（Dirichlet 核）}：\(N\) 阶 Dirichlet 核为

\[
D_N(t) = \frac{1}{2} + \sum_{n=1}^{N} \cos nt = \frac{\sin((N + \frac{1}{2})t)}{2\sin(t/2)}
\]

Fourier 级数的部分和可表示为
\(S_N f(x) = \frac{1}{\pi} \int_{-\pi}^{\pi} f(x - t) D_N(t) \, dt\)。

\textbf{定义 50.3.2（Fejér 核）}：\(N\) 阶 Fejér 核为 Dirichlet 核的
Cesàro 平均：

\[
F_N(t) = \frac{1}{N+1} \sum_{k=0}^{N} D_k(t) = \frac{1}{2(N+1)} \left(\frac{\sin((N+1)t/2)}{\sin(t/2)}\right)^2
\]

Fejér
核构成\textbf{好核}：\(F_N(t) \geq 0\)，\(\frac{1}{\pi} \int_{-\pi}^{\pi} F_N(t) dt = 1\)，且对任意
\(\delta > 0\)，\(\int_{\delta \leq |t| \leq \pi} F_N(t) dt \to 0\)（\(N \to \infty\)）。

\bigskip\hrule\bigskip

\subsubsection{50.4 Fejér 定理}\label{fejuxe9r-ux5b9aux7406}

\textbf{定理 50.4.1（Fejér 定理）}：设 \(f \in L^1([-\pi, \pi])\)。则
\(f\) 的 Fourier 级数的 Cesàro 平均
\(\sigma_N f(x) = \frac{1}{N+1} \sum_{k=0}^{N} S_k f(x)\) 满足：

\begin{enumerate}
\def\labelenumi{(\roman{enumi})}
\tightlist
\item
  若 \(f\) 在 \(x\) 处连续，则 \(\sigma_N f(x) \to f(x)\)
\item
  若 \(f\) 在 \([-\pi, \pi]\) 上连续且 \(f(-\pi) = f(\pi)\)，则
  \(\sigma_N f \to f\) 一致收敛
\end{enumerate}

\textbf{证明}：\(\sigma_N f(x) = \frac{1}{\pi} \int_{-\pi}^{\pi} f(x - t) F_N(t) dt\)。由
Fejér 核是''好核''的性质，可用卷积逼近的标准论证。\(\blacksquare\)

\textbf{推论 50.4.2（Weierstrass 第二逼近定理）}：三角多项式在
\(C([-\pi, \pi])\)（周期连续函数）中关于一致范数稠密。

\bigskip\hrule\bigskip

\subsubsection{50.5 逐点收敛性}\label{ux9010ux70b9ux6536ux655bux6027}

\textbf{定理 50.5.1（Dini 判别法）}：设
\(f \in L^1([-\pi, \pi])\)，\(x \in [-\pi, \pi]\)。若存在 \(\delta > 0\)
使得

\[
\int_0^{\delta} \frac{|f(x + t) + f(x - t) - 2f(x)|}{t} \, dt < \infty
\]

则 \(f\) 的 Fourier 级数在 \(x\) 处收敛到 \(f(x)\)。特别地，若 \(f\) 在
\(x\) 处 Lipschitz 连续（或可导），则 Dini 条件成立。

\textbf{定理 50.5.2（Dirichlet-Jordan 判别法）}：若 \(f\) 在 \(x\)
的某个邻域内有界变差，则 \(f\) 的 Fourier 级数在 \(x\) 处收敛到
\(\frac{f(x^+) + f(x^-)}{2}\)。

\textbf{证明概要}（Dini 判别法）：Fourier 部分和可写为
\(S_N f(x) - f(x) = \frac{1}{\pi} \int_0^{\pi} \frac{f(x+t) + f(x-t) - 2f(x)}{2\sin(t/2)} \sin((N+\frac{1}{2})t) dt\)。令
\(\varphi(t) = \frac{f(x+t) + f(x-t) - 2f(x)}{2\sin(t/2)}\)。由 Dini
条件，\(\varphi \in L^1([0, \pi])\)。由 Riemann-Lebesgue
引理，\(\int_0^{\pi} \varphi(t) \sin((N+\frac{1}{2})t) dt \to 0\)（\(N \to \infty\)）。\(\blacksquare\)

\textbf{定理 50.5.3（Riemann-Lebesgue 引理）}：设
\(f \in L^1([-\pi, \pi])\)。则

\[
\lim_{|n| \to \infty} \int_{-\pi}^{\pi} f(t) e^{-int} \, dt = 0
\]

即 Fourier 系数 \(a_n, b_n \to 0\)（\(n \to \infty\)）。

\textbf{证明}：先对 \(f\)
为特征函数验证（直接积分），再对简单函数，最后用 \(L^1\)
中简单函数的稠密性推广到一般 \(f \in L^1\)。\(\blacksquare\)

\bigskip\hrule\bigskip

\textbf{注记 50.5.1}：关于 Fourier
级数的进一步深入内容（包括分布理论、Sobolev
空间、小波分析等）将在卷十四（泛函分析）和卷二十五（调和分析）中详细展开。\(\blacksquare\)

\bigskip\hrule\bigskip

\subsection{Ch 51 位势论（Potential
Theory）}\label{ch-51-ux4f4dux52bfux8bbapotential-theory}

位势论起源于 Newton
引力势和静电学，研究调和函数、上调和函数的性质及其与测度和容度的深刻关系。在
\(\mathbb{R}^n\)（\(n \geq 3\)）中，\textbf{Newton 核}
\(N(x, y) = \|x - y\|^{2-n}\) 是 Laplace 方程的基本解（\(n = 2\) 时为
\(-\log\|x - y\|\)）。\textbf{Riesz 核}
\(R_\alpha(x, y) = \|x - y\|^{\alpha-n}\)（\(0 < \alpha < n\)）将 Newton
位势推广至分数阶，是 Riesz 位势理论的基础。设 \(\mu\) 为 Radon
测度，函数 \(U^\mu(x) = \int N(x, y) d\mu(y)\) 称为 \(\mu\) 的 Newton
位势，它是 \(\mathbb{R}^n \setminus \operatorname{supp} \mu\)
上的调和函数。

\textbf{上调和函数} \(u\) 满足
\(-\Delta u \geq 0\)（分布意义下），具有三项基本性质：（1）局部可积且下半连续，在任意球上的平均值不小于球心值（\textbf{平均值不等式}）；（2）\textbf{最大值原理}------若
\(u\) 在区域 \(\Omega\) 内部达到全局最小值，则 \(u\) 在 \(\Omega\)
上为常数；（3）Riesz 分解定理------任意上调和函数可表为 Newton
位势与调和函数的和。

\textbf{容度} \(\operatorname{Cap}(E)\) 衡量集合
\(E \subseteq \mathbb{R}^n\)``容纳''电荷的能力，定义为

\[\operatorname{Cap}(E) = \sup\{\mu(E) \mid \operatorname{supp} \mu \subseteq E, \; U^\mu(x) \leq 1 \text{ 对所有 } x\}\]

零容度集合是位势论中的''可忽略集''：Dirichlet 问题的边界正则性由
\textbf{Wiener 判别法}给出------边界点 \(x_0 \in \partial\Omega\)
为正则点（即存在 \(x_0\) 处的闸函数）当且仅当

\[\sum_{k=1}^{\infty} 2^{k(n-2)} \operatorname{Cap}\big(B(x_0, 2^{-k}) \setminus \Omega\big) = \infty\]

（\(n = 2\) 时用对数容度）。位势论与概率论有本质联系：调和函数对应于
\textbf{Brown
运动}的期望（\(u(x) = \mathbb{E}^x[f(B_\tau)]\)），容度与击中概率直接相关（\(\operatorname{Cap}(E) > 0\)
当且仅当 Brown 运动以正概率击中
\(E\)）。这一深刻对应使位势论成为连接分析、几何与概率的桥梁，在随机过程、复分析和偏微分方程中有广泛应用。\(\blacksquare\)

\bigskip\hrule\bigskip

\subsection{Ch 52 Plateau 问题}\label{ch-52-plateau-ux95eeux9898}

Plateau问题源于十九世纪比利时物理学家Plateau的肥皂膜实验：给定空间中的闭曲线\(\Gamma\)，求以\(\Gamma\)为边界的极小曲面。数学表述为在\(C^1\)参数化\(X: D \to \mathbb{R}^3\)（\(D\)为单位圆盘）的约束\(X(\partial D) = \Gamma\)下，最小化面积泛函\(\mathcal{A}[X] = \iint_D |X_u \times X_v| du dv\)。Douglas（1931）和Radó（1930）独立地使用变分方法证明了任何可求长Jordan曲线\(\Gamma\)都存在以它为边界的极小曲面，并获得了1936年首届Fields奖。其核心步骤为：(i)将面积泛函转化为Dirichlet能量\(\mathcal{E}[X] = \frac{1}{2}\iint_D (|X_u|^2 + |X_v|^2) du dv\)（极值下两者等价）；(ii)在Sobolev空间\(H^1(D, \mathbb{R}^3)\)中极小化，利用下半连续性和紧性得到弱解；(iii)证明弱解的光滑性。几何测度论将Plateau问题推广为：求整流传（integer
multiplicity
current）\(T\)使得\(\partial T = \Gamma\)且质量\(\mathbf{M}(T)\)极小，Federer-Fleming的闭包定理和紧性定理为此提供了严格的框架。

\bigskip\hrule\bigskip

\subsection{Ch 53：殆周期函数（Almost Periodic
Functions）}\label{ch-53ux6b86ux5468ux671fux51fdux6570almost-periodic-functions}

殆周期函数理论由 H. Bohr (1925) 创立，将周期函数的 Fourier
分析推广到具有''近似周期性''的连续函数。定义：\(\mathbb{R}\)
上的连续函数 \(f\) 称为 \textbf{Bohr 殆周期}，若对任意
\(\varepsilon > 0\)，存在相对稠密的平移数集------即存在长度
\(L = L(\varepsilon) > 0\) 使得每个长为 \(L\) 的区间至少含一个 \(\tau\)
满足
\(\|f(x+\tau) - f(x)\|_\infty < \varepsilon\)。\textbf{核心定理}：（1）\textbf{Bohr
基本定理}------\(f\) 殆周期当且仅当它是三角多项式
\(\sum c_n e^{i\lambda_n x}\) 在一致范数下的极限（即 Fourier
级数具有可数个非零系数）；（2）\textbf{平均遍历定理}------平均
\(\mathcal{M}\{f\} = \lim_{T\to\infty} \frac{1}{2T} \int_{-T}^{T} f(x)dx\)
存在，且 \(\mathcal{M}\{|f|^2\} = \sum |c_n|^2\)（Parseval
恒等式）；（3）\textbf{Bochner 准则}------\(f\) 殆周期当且仅当
\(\{f(x+h) : h \in \mathbb{R}\}\) 在 \(\operatorname{C}(\mathbb{R})\)
中相对紧。Bochner (1927) 和 von Neumann (1934)
将该理论推广到群上的殆周期函数------\(\operatorname{C}(G)\)
中的殆周期函数构成平移紧算子的谱理论，奠定了局部紧交换群上调和分析的公理化基础。

\textbf{定义 53.1}（Bohr 殆周期）：\(\mathbb{R}\) 上的连续函数
\(f \in \operatorname{C}(\mathbb{R})\) 称为 \textbf{Bohr
殆周期}，若对任意 \(\varepsilon > 0\)，存在 \(L = L(\varepsilon) > 0\)
使得 \(\mathbb{R}\) 上任意长度为 \(L\) 的区间至少含有一个 \(\tau\)
满足对一切 \(x \in \mathbb{R}\) 有
\(|f(x+\tau) - f(x)| < \varepsilon\)。\(\tau\) 称为 \(f\) 的
\(\varepsilon\)-殆周期。所有 Bohr 殆周期函数构成闭子代数
\(\operatorname{AP}(\mathbb{R}) \subset \operatorname{C}_b(\mathbb{R})\)（有界连续函数空间）。

\textbf{定理 53.1}（Bohr
基本定理）：\(f \in \operatorname{AP}(\mathbb{R})\) 当且仅当 \(f\)
可在一致范数下被三角多项式逼近。等价地，存在至多可数个实数
\(\lambda_n\)（\(f\) 的 Fourier 指数）和非零复数 \(a_n\)（\(f\) 的
Fourier 系数）使得
\(f(x) \sim \sum_n a_n e^{i\lambda_n x}\)。\emph{证明思路}：充分性------三角多项式显然殆周期，一致极限保持殆周期性质。必要性------定义平均值
\(\mathcal{M}\{f(x)e^{-i\lambda x}\} = \lim_{T \to \infty} \frac{1}{2T} \int_{-T}^T f(x) e^{-i\lambda x} dx\)，由
Bochner-Fejér 核的逼近性质证明该极限存在，且仅可数个 \(\lambda\)
使该平均非零。用近似恒同的方法构造三角多项式的逼近序列。

\textbf{定理 53.2}（Parseval 恒等式）：对
\(f \in \operatorname{AP}(\mathbb{R})\)，有
\(\mathcal{M}\{|f|^2\} = \sum_n |a_n|^2\)，其中平均值
\(\mathcal{M}\{g\} = \lim_{T \to \infty} \frac{1}{2T} \int_{-T}^T g(x) dx\)。\emph{证明思路}：将三角逼近代入积分的平方展开，利用
\(\mathcal{M}\{e^{i(\lambda_m - \lambda_n)x}\} = \delta_{mn}\)
的正交性。

\bigskip\hrule\bigskip

\subsection{Ch 54：积分变换深化（Integral
Transforms）}\label{ch-54ux79efux5206ux53d8ux6362ux6df1ux5316integral-transforms}

经典积分变换是以积分核为桥梁、实现微分方程求解和信号分析的基本工具。\textbf{核心内容}：（1）\textbf{Fourier
变换}------\(\mathcal{F}[f](\xi) = \int_{-\infty}^{\infty} f(x) e^{-2\pi i x\xi} dx\)，在
\(L^2(\mathbb{R})\) 上为酉算子，Plancherel 定理
\(\|\mathcal{F}f\|_{L^2} = \|f\|_{L^2}\)；（2）\textbf{Laplace
变换}------\(\mathcal{L}[f](s) = \int_0^\infty f(t) e^{-st} dt\)，将常系数
ODE 化归为代数方程，逆变换由 Bromwich 积分
\(\frac{1}{2\pi i} \int_{\gamma-i\infty}^{\gamma+i\infty} \mathcal{L}[f](s) e^{st} ds\)
给出；（3）\textbf{Mellin
变换}------\(\mathcal{M}[f](s) = \int_0^\infty f(x) x^{s-1} dx\)，在线性
PDE 边界值问题和数论（Dirichlet 级数）中自然出现；（4）\textbf{Hilbert
变换}------\(\mathcal{H}[f](x) = \frac{1}{\pi} \operatorname{p.v.} \int_{-\infty}^{\infty} \frac{f(t)}{x-t} dt\)，是信号解析表示和色散关系的数学基础；（5）\textbf{Radon
变换}------\(Rf(\theta, s) = \int_{x \cdot \theta = s} f(x) d\sigma(x)\)，作为
\(\mathbb{R}^n\) 上的积分变换在 CT
断层扫描中起着核心数学支撑作用。反演公式（滤波反投影）由 Johann Radon
(1917)
给出。\textbf{小波变换}------\(\mathcal{W}_\psi f(a,b) = \frac{1}{\sqrt{|a|}} \int f(t) \overline{\psi(\frac{t-b}{a})} dt\)，同时捕捉时间-频率局部特征，是
Fourier 变换的自然推广。

\bigskip\hrule\bigskip

\subsection{Ch
55：特殊函数与正交多项式}\label{ch-55ux7279ux6b8aux51fdux6570ux4e0eux6b63ux4ea4ux591aux9879ux5f0f}

特殊函数是满足特定 ODE
并通过超几何级数表达的经典函数的统称，在数学物理和渐近分析中有核心应用。\textbf{核心类型}：（1）\textbf{\(\Gamma\)
函数}：\(\Gamma(z) = \int_0^\infty t^{z-1} e^{-t} dt\)（\(\operatorname{Re}(z) > 0\)），解析延拓到
\(\mathbb{C} \setminus \{0, -1, -2, \ldots\}\)。满足
\(\Gamma(z+1) = z\Gamma(z)\)，\(\Gamma(n) = (n-1)!\)。\textbf{Weierstrass
无穷乘积}：\(\frac{1}{\Gamma(z)} = z e^{\gamma z} \prod_{n=1}^\infty (1 + \frac{z}{n}) e^{-z/n}\)（\(\gamma\)
为 Euler 常数）。\textbf{Beta 函数}
\(B(\alpha, \beta) = \frac{\Gamma(\alpha)\Gamma(\beta)}{\Gamma(\alpha+\beta)}\)。（2）\textbf{Bessel
函数}：Bessel 微分方程 \(z^2 y'' + z y' + (z^2 - \nu^2) y = 0\)
的解。第一类
\(J_\nu(z) = \sum_{k=0}^\infty \frac{(-1)^k}{k! \Gamma(k+\nu+1)} (\frac{z}{2})^{2k+\nu}\)，第二类
\(Y_\nu(z)\)（Weber
函数）。（3）\textbf{超几何函数}：\(_2F_1(a,b;c;z) = \sum_{n=0}^\infty \frac{(a)_n (b)_n}{(c)_n n!} z^n\)，其中
\((q)_n = q(q+1)\cdots(q+n-1)\) 为 Pochhammer 符号，收敛半径
1。其满足超几何微分方程，是几乎所有经典特殊函数（如
Legendre、Chebyshev、Jacobi 多项式）的统一母体。

\textbf{正交多项式}是 \(L^2([a,b], w(x)dx)\) 中通过 Gram-Schmidt 对
\(\{1, x, x^2, \ldots\}\)
正交化构造的多项式系。\textbf{基本定理}（三项递推关系）：任意正测度生成的正交多项式系
\(\{p_n(x)\}\) 满足
\(x p_n(x) = p_{n+1}(x) + \alpha_n p_n(x) + \beta_n p_{n-1}(x)\)（\(\beta_n > 0\)），该关系由谱定理和
Jacobi 矩阵刻画。\textbf{经典正交多项式分类}（Bochner,
1929）：仅有三类超几何型的正交多项式满足 Sturm-Liouville
问题------Jacobi（含
Gegenbauer、Legendre、Chebyshev）、Laguerre、Hermite。其生成函数、Rodrigues
公式和微分方程有统一结构。

\newpage

\section{卷九：复分析}\label{ux5377ux4e5dux590dux5206ux6790}

\bigskip\hrule\bigskip

\begin{quote}
\textbf{前置知识}：本卷建立在卷一（复数构造、级数理论）和卷四（极限、连续、导数、积分、幂级数）的基础之上。读者应熟悉
\(\varepsilon\)-\(\delta\) 语言、一致收敛性、幂级数的收敛半径、以及
Riemann 积分的曲线积分概念。
\end{quote}

\subsection{Ch 51 全纯函数}\label{ch-51-ux5168ux7eafux51fdux6570}

\subsubsection{51.1
复平面与复函数}\label{ux590dux5e73ux9762ux4e0eux590dux51fdux6570}

\textbf{定义 51.1.1（复平面）}：\(\mathbb{C}\) 可视为
\(\mathbb{R}^2\)，配备复数乘法。\(z. = x + iy \in \mathbb{C}\) 对应
\((x, y) \in \mathbb{R}^2\)。\(z\) 的模为 \(|z| = \sqrt{x^2 + y^2}\)。

\textbf{定义 51.1.2（复函数的极限与连续性）}：设
\(\Omega \subseteq \mathbb{C}\)
为开集，\(f : \Omega \to \mathbb{C}\)，\(z_0 \in \Omega\)。

\[
\lim_{z \to z_0} f(z) = w_0 \iff \forall \varepsilon > 0, \exists \delta > 0, \; 0 < |z - z_0| < \delta \implies |f(z) - w_0| < \varepsilon
\]

\(f\) 在 \(z_0\) 处连续若 \(\lim_{z \to z_0} f(z) = f(z_0)\)。

\textbf{注记
51.1.1}：复函数的极限和连续性在形式上和实函数完全相同，但这里 \(z\)
沿复平面中任意路径趋近 \(z_0\)，比实直线上只有左右两个方向的要求更强。

\bigskip\hrule\bigskip

\subsubsection{51.2 全纯函数与 Cauchy-Riemann
方程}\label{ux5168ux7eafux51fdux6570ux4e0e-cauchy-riemann-ux65b9ux7a0b}

\textbf{定义 51.2.1（复导数与全纯函数）}：设
\(\Omega \subseteq \mathbb{C}\)
为开集，\(f : \Omega \to \mathbb{C}\)，\(z_0 \in \Omega\)。若极限

\[
f'(z_0) = \lim_{z \to z_0} \frac{f(z) - f(z_0)}{z - z_0}
\]

存在，则称 \(f\) 在 \(z_0\) 处\textbf{复可导}。若 \(f\) 在 \(\Omega\)
上每点复可导，则称 \(f\) 为 \(\Omega\)
上的\textbf{全纯函数}（或\textbf{解析函数}）。

\textbf{定理 51.2.1（Cauchy-Riemann 方程）}：设
\(f(z) = u(x, y) + i v(x, y)\)（\(u, v : \mathbb{R}^2 \to \mathbb{R}\)），\(z_0 = x_0 + iy_0\)。\(f\)
在 \(z_0\) 处复可导当且仅当 \(u, v\) 在 \((x_0, y_0)\) 处可微且满足

\[
\frac{\partial u}{\partial x} = \frac{\partial v}{\partial y}, \quad \frac{\partial u}{\partial y} = -\frac{\partial v}{\partial x}
\]

此时
\(f'(z_0) = \frac{\partial u}{\partial x}(x_0, y_0) + i\frac{\partial v}{\partial x}(x_0, y_0) = \frac{\partial v}{\partial y}(x_0, y_0) - i\frac{\partial u}{\partial y}(x_0, y_0)\)。

\textbf{证明}：

（\(\Rightarrow\)）设 \(f'(z_0) = A + iB\) 存在。取 \(z \to z_0\)
沿实轴方向（\(z = z_0 + h\)，\(h \in \mathbb{R}\)）：

\[
f'(z_0) = \lim_{h \to 0} \frac{u(x_0 + h, y_0) + iv(x_0 + h, y_0) - u(x_0, y_0) - iv(x_0, y_0)}{h} = \frac{\partial u}{\partial x} + i\frac{\partial v}{\partial x}
\]

取 \(z \to z_0\) 沿虚轴方向（\(z = z_0 + ih\)，\(h \in \mathbb{R}\)）：

\[
f'(z_0) = \lim_{h \to 0} \frac{u(x_0, y_0 + h) + iv(x_0, y_0 + h) - u(x_0, y_0) - iv(x_0, y_0)}{ih} = \frac{1}{i}\frac{\partial u}{\partial y} + \frac{\partial v}{\partial y} = -i\frac{\partial u}{\partial y} + \frac{\partial v}{\partial y}
\]

比较实部和虚部即得 Cauchy-Riemann 方程。

（\(\Leftarrow\)）设 \(u, v\) 可微且满足 C-R
方程。由实可微性，\(u(x_0 + h, y_0 + k) - u(x_0, y_0) = u_x h + u_y k + o(\sqrt{h^2 + k^2})\)，\(v\)
类似。令 \(\Delta z = h + ik\)，则

\[
f(z_0 + \Delta z) - f(z_0) = (u_x + iv_x)h + (u_y + iv_y)k + o(|\Delta z|)
\]

由 C-R 方程，\(u_y + iv_y = -v_x + iu_x = i(u_x + iv_x)\)。故

\[
f(z_0 + \Delta z) - f(z_0) = (u_x + iv_x)(h + ik) + o(|\Delta z|) = (u_x + iv_x)\Delta z + o(|\Delta z|)
\]

因此
\(\lim_{\Delta z \to 0} \frac{f(z_0 + \Delta z) - f(z_0)}{\Delta z} = u_x + iv_x\)。\(\blacksquare\)

\textbf{定理 51.2.2（全纯函数的基本运算）}：若 \(f, g\) 在 \(\Omega\)
上全纯，则

\begin{enumerate}
\def\labelenumi{(\roman{enumi})}
\tightlist
\item
  \(f \pm g\)，\(fg\) 全纯，且
  \((f \pm g)' = f' \pm g'\)，\((fg)' = f'g + fg'\)
\item
  若 \(g \neq 0\)，则 \(f/g\) 全纯，且 \((f/g)' = (f'g - fg')/g^2\)
\item
  链式法则：\((f \circ g)'(z) = f'(g(z)) g'(z)\)
\end{enumerate}

\textbf{证明}：与实函数导数运算法则的证明完全相同（仅需将绝对值替换为模）。\(\blacksquare\)

\bigskip\hrule\bigskip

\subsubsection{51.3
幂级数作为全纯函数}\label{ux5e42ux7ea7ux6570ux4f5cux4e3aux5168ux7eafux51fdux6570}

\textbf{定理 51.3.1（幂级数的全纯性）}：设
\(\sum_{n=0}^{\infty} a_n z^n\) 的收敛半径为 \(R > 0\)。则
\(f(z) = \sum_{n=0}^{\infty} a_n z^n\) 在
\(\{z \in \mathbb{C} \mid |z| < R\}\) 上全纯，且

\[
f'(z) = \sum_{n=1}^{\infty} n a_n z^{n-1}, \quad |z| < R
\]

（即幂级数可以逐项求导，收敛半径不变。）

\textbf{证明}：由幂级数的性质，\(\sum n a_n z^{n-1}\) 的收敛半径也是
\(R\)。在 \(|z| \leq r < R\) 的紧集上，逐项求导的级数一致收敛（由
Weierstrass M-判别法），故可逐项求导。\(\blacksquare\)

\textbf{推论
51.3.2}：\(e^z = \sum_{n=0}^{\infty} \frac{z^n}{n!}\)，\(\sin z = \sum_{n=0}^{\infty} (-1)^n \frac{z^{2n+1}}{(2n+1)!}\)，\(\cos z = \sum_{n=0}^{\infty} (-1)^n \frac{z^{2n}}{(2n)!}\)
在整个 \(\mathbb{C}\) 上全纯，且
\((e^z)' = e^z\)，\((\sin z)' = \cos z\)，\((\cos z)' = -\sin z\)。

\textbf{注记 51.3.1}：这与实函数情形形式上完全一致，但这里 \(z\)
取遍整个复平面。

\bigskip\hrule\bigskip

\subsection{Ch 52 复积分}\label{ch-52-ux590dux79efux5206}

\subsubsection{52.1 曲线积分}\label{ux66f2ux7ebfux79efux5206}

\textbf{定义 52.1.1（曲线）}：\(\mathbb{C}\)
中的\textbf{曲线}（或\textbf{路径}）\(\gamma\) 是一个连续可微映射
\(\gamma : [a, b] \to \mathbb{C}\)。\(\gamma\) 称为\textbf{闭曲线}若
\(\gamma(a) = \gamma(b)\)。

\textbf{定义 52.1.2（复曲线积分）}：设 \(f\) 在 \(\gamma\) 上连续。定义
\(f\) 沿 \(\gamma\) 的积分为

\[
\int_{\gamma} f(z) \, dz = \int_a^b f(\gamma(t)) \gamma'(t) \, dt
\]

\textbf{定理 52.1.1（积分的基本估计）}：设 \(|f(z)| \leq M\) 在
\(\gamma\) 上，\(L\) 为 \(\gamma\) 的长度。则

\[
\left|\int_{\gamma} f(z) \, dz\right| \leq M L
\]

\textbf{证明}： \[
\left|\int_{\gamma} f(z) dz\right| = \left|\int_a^b f(\gamma(t))\gamma'(t) dt\right| \leq \int_a^b |f(\gamma(t))||\gamma'(t)| dt \leq M \int_a^b |\gamma'(t)| dt = M L
\]

\(\blacksquare\)

\bigskip\hrule\bigskip

\subsubsection{52.2 Cauchy
积分定理}\label{cauchy-ux79efux5206ux5b9aux7406}

\textbf{定理 52.2.1（Cauchy 积分定理------矩形版本）}：设
\(\Omega \subseteq \mathbb{C}\) 为开集，\(f\) 在 \(\Omega\)
上全纯，\(R \subseteq \Omega\) 为闭矩形（边界 \(\partial R\)
由四条线段组成）。则

\[
\int_{\partial R} f(z) \, dz = 0
\]

\textbf{证明}（Goursat）：将矩形 \(R\) 四等分为
\(R^{(1)}, \ldots, R^{(4)}\)。则

\[
\int_{\partial R} f(z) dz = \sum_{j=1}^{4} \int_{\partial R^{(j)}} f(z) dz
\]

（公共边上的积分抵消）。设 \(R_1\) 为四个子矩形中积分绝对值最大者，则
\(|\int_{\partial R} f| \leq 4 |\int_{\partial R_1} f|\)。重复此过程，得到矩形套
\(R \supseteq R_1 \supseteq R_2 \supseteq \cdots\) 使得

\[
\left|\int_{\partial R_n} f(z) dz\right| \geq \frac{1}{4^n} \left|\int_{\partial R} f(z) dz\right|
\]

设 \(z_0 \in \bigcap_{n=1}^{\infty} R_n\)。由 \(f\) 在 \(z_0\) 处全纯，

\[
f(z) = f(z_0) + f'(z_0)(z - z_0) + \psi(z)(z - z_0)
\]

其中 \(\psi(z) \to 0\)（\(z \to z_0\)）。前两项
\(f(z_0) + f'(z_0)(z - z_0)\) 有原函数，沿闭矩形积分为 \(0\)。故

\[
\int_{\partial R_n} f(z) dz = \int_{\partial R_n} \psi(z)(z - z_0) dz
\]

设 \(R_n\) 的边长为 \(L/2^n\)，对角线长为 \(\sqrt{2}L/2^n\)。对任意
\(\varepsilon > 0\)，当 \(n\) 足够大时 \(|\psi(z)| < \varepsilon\)。故

\[
\left|\int_{\partial R_n} f(z) dz\right| \leq \varepsilon \cdot \frac{\sqrt{2}L}{2^n} \cdot \frac{4L}{2^n} = \frac{4\sqrt{2}L^2}{4^n} \varepsilon
\]

因此
\(|\int_{\partial R} f| \leq 4^n \cdot \frac{4\sqrt{2}L^2}{4^n} \varepsilon = 4\sqrt{2}L^2 \varepsilon\)。由
\(\varepsilon\) 任意性，\(\int_{\partial R} f = 0\)。\(\blacksquare\)

\textbf{定理 52.2.2（Cauchy 积分定理------一般版本）}：设
\(\Omega \subseteq \mathbb{C}\) 为\textbf{单连通}开集（任何闭曲线可在
\(\Omega\) 内连续收缩为一点），\(f\) 在 \(\Omega\) 上全纯。则对
\(\Omega\) 中任何闭曲线 \(\gamma\)，

\[
\int_{\gamma} f(z) \, dz = 0
\]

\textbf{注记
52.2.1}：一般版本的证明需要将区域三角剖分或使用同伦论证，此处从略。核心结论：在单连通区域内，全纯函数的闭曲线积分为零。

\textbf{推论 52.2.3（原函数的存在性）}：在单连通区域 \(\Omega\)
上全纯的函数 \(f\) 有原函数
\(F\)（\(F' = f\)），\(F(z) = \int_{z_0}^{z} f(\zeta) d\zeta\)（积分与路径无关）。

\bigskip\hrule\bigskip

\subsubsection{52.3 Cauchy
积分公式}\label{cauchy-ux79efux5206ux516cux5f0f}

\textbf{定理 52.3.1（Cauchy 积分公式）}：设
\(\Omega \subseteq \mathbb{C}\) 为开集，\(f\) 在 \(\Omega\)
上全纯，\(\overline{D}(z_0, r) \subseteq \Omega\)（以 \(z_0\)
为中心、\(r\) 为半径的闭圆盘）。则对任意 \(z \in D(z_0, r)\)（开圆盘），

\[
f(z) = \frac{1}{2\pi i} \int_{\partial D(z_0, r)} \frac{f(\zeta)}{\zeta - z} \, d\zeta
\]

\textbf{证明}：固定 \(z \in D(z_0, r)\)。考虑函数
\(g(\zeta) = \frac{f(\zeta) - f(z)}{\zeta - z}\)（\(\zeta \neq z\)），\(g(z) = f'(z)\)。\(g\)
在 \(\Omega \setminus \{z\}\) 上全纯，在 \(z\) 处连续（从而在 \(\Omega\)
上全纯，由 Riemann 可去奇点定理------见定理 53.4.1）。

在 \(\Omega\) 内取以 \(z\) 为中心的小圆
\(\partial D(z, \varepsilon)\)。由 Cauchy 积分定理（应用于环形区域），

\[
\int_{\partial D(z_0, r)} g(\zeta) d\zeta = \int_{\partial D(z, \varepsilon)} g(\zeta) d\zeta
\]

即
\(\int_{\partial D(z_0, r)} \frac{f(\zeta)}{\zeta - z} d\zeta - f(z) \int_{\partial D(z_0, r)} \frac{d\zeta}{\zeta - z} = \int_{\partial D(z, \varepsilon)} \frac{f(\zeta) - f(z)}{\zeta - z} d\zeta\)。

计算 \(\int_{\partial D(z_0, r)} \frac{d\zeta}{\zeta - z} = 2\pi i\)（由
Cauchy 积分公式基础情形 \(f \equiv 1\)）。右边当 \(\varepsilon \to 0\)
时，\(|g(\zeta)| \leq |f'(z)| + 1\)（有界），故积分绝对值
\(\leq M \cdot 2\pi\varepsilon \to 0\)。因此

\[
\int_{\partial D(z_0, r)} \frac{f(\zeta)}{\zeta - z} d\zeta - 2\pi i f(z) = 0
\]

即
\(f(z) = \frac{1}{2\pi i} \int_{\partial D(z_0, r)} \frac{f(\zeta)}{\zeta - z} d\zeta\)。\(\blacksquare\)

\textbf{推论 52.3.2（Cauchy 积分公式------高阶导数）}：在定理 52.3.1
的条件下，\(f\) 在 \(D(z_0, r)\) 上无穷次复可导，且

\[
f^{(n)}(z) = \frac{n!}{2\pi i} \int_{\partial D(z_0, r)} \frac{f(\zeta)}{(\zeta - z)^{n+1}} \, d\zeta, \quad n = 0, 1, 2, \ldots
\]

\textbf{证明}：对 \(n\) 归纳。\(n = 0\) 即 Cauchy 积分公式。假设公式对
\(n\) 成立。由差商极限和积分号下取极限的合法性（被积函数一致有界），

\[
f^{(n+1)}(z) = \lim_{h \to 0} \frac{f^{(n)}(z+h) - f^{(n)}(z)}{h} = \frac{n!}{2\pi i} \lim_{h \to 0} \int_{\partial D} \frac{f(\zeta)}{h} \left(\frac{1}{(\zeta - z - h)^{n+1}} - \frac{1}{(\zeta - z)^{n+1}}\right) d\zeta
\]

利用
\(\frac{1}{(\zeta - z - h)^{n+1}} - \frac{1}{(\zeta - z)^{n+1}} \to (n+1) \frac{h}{(\zeta - z)^{n+2}}\)（\(h \to 0\)），得
\(f^{(n+1)}(z) = \frac{(n+1)!}{2\pi i} \int_{\partial D} \frac{f(\zeta)}{(\zeta - z)^{n+2}} d\zeta\)。\(\blacksquare\)

\textbf{推论
52.3.3（全纯函数无穷次可导）}：全纯函数自动无穷次可导（与实分析根本不同）。

\bigskip\hrule\bigskip

\subsection{Ch 53 级数展开}\label{ch-53-ux7ea7ux6570ux5c55ux5f00}

\subsubsection{53.1 Taylor
级数展开}\label{taylor-ux7ea7ux6570ux5c55ux5f00}

\textbf{定理 53.1.1（Taylor 展开定理）}：设 \(f\) 在 \(\Omega\)
上全纯，\(z_0 \in \Omega\)。设
\(R = \operatorname{dist}(z_0, \partial\Omega)\)（\(z_0\) 到 \(\Omega\)
边界的距离）。则 \(f\) 在 \(D(z_0, R)\) 上可展开为收敛的 Taylor 级数：

\[
f(z) = \sum_{n=0}^{\infty} a_n (z - z_0)^n, \quad a_n = \frac{f^{(n)}(z_0)}{n!}, \quad |z - z_0| < R
\]

\textbf{证明}：对 \(|z - z_0| < r < R\)，由 Cauchy 积分公式，

\[
f(z) = \frac{1}{2\pi i} \int_{\partial D(z_0, r)} \frac{f(\zeta)}{\zeta - z} d\zeta = \frac{1}{2\pi i} \int_{\partial D(z_0, r)} \frac{f(\zeta)}{(\zeta - z_0) - (z - z_0)} d\zeta
\]

展开几何级数：\(\frac{1}{(\zeta - z_0) - (z - z_0)} = \frac{1}{\zeta - z_0} \cdot \frac{1}{1 - \frac{z - z_0}{\zeta - z_0}} = \frac{1}{\zeta - z_0} \sum_{n=0}^{\infty} \left(\frac{z - z_0}{\zeta - z_0}\right)^n\)。由于
\(|z - z_0| < |\zeta - z_0| = r\)，级数一致收敛（关于
\(\zeta\)），可逐项积分：

\[
f(z) = \sum_{n=0}^{\infty} \left(\frac{1}{2\pi i} \int_{\partial D(z_0, r)} \frac{f(\zeta)}{(\zeta - z_0)^{n+1}} d\zeta\right) (z - z_0)^n = \sum_{n=0}^{\infty} \frac{f^{(n)}(z_0)}{n!} (z - z_0)^n
\]

由 Cauchy 高阶导数公式，系数确为 \(f^{(n)}(z_0)/n!\)。\(\blacksquare\)

\textbf{推论 53.1.2（Liouville 定理）}：若 \(f\) 在整个 \(\mathbb{C}\)
上全纯且有界（即存在 \(M\) 使得
\(|f(z)| \leq M\)，\(\forall z \in \mathbb{C}\)），则 \(f\) 为常数。

\textbf{证明}：由 Cauchy 高阶导数公式，对任意 \(R > 0\)，

\[
|f'(z_0)| = \left|\frac{1}{2\pi i} \int_{\partial D(z_0, R)} \frac{f(\zeta)}{(\zeta - z_0)^2} d\zeta\right| \leq \frac{1}{2\pi} \cdot \frac{M}{R^2} \cdot 2\pi R = \frac{M}{R}
\]

令 \(R \to \infty\) 得 \(f'(z_0) = 0\)。\(f' \equiv 0\) 意味着 \(f\)
为常数。\(\blacksquare\)

\textbf{推论 53.1.3（代数基本定理的复分析证明）}：任何非常数复系数多项式
\(P(z) = a_n z^n + \cdots + a_0\)（\(a_n \neq 0\)，\(n \geq 1\)）在
\(\mathbb{C}\) 中有根。

\textbf{证明}：假设 \(P(z) \neq 0\)（\(\forall z \in \mathbb{C}\)）。则
\(f(z) = 1/P(z)\) 在整个 \(\mathbb{C}\) 上全纯。当 \(|z| \to \infty\)
时，\(|P(z)| \to \infty\)，故 \(|f(z)| \to 0\)，从而 \(f\) 有界。由
Liouville 定理，\(f\) 为常数，与 \(n \geq 1\) 矛盾。\(\blacksquare\)

\bigskip\hrule\bigskip

\subsubsection{53.2 Laurent
级数展开}\label{laurent-ux7ea7ux6570ux5c55ux5f00}

\textbf{定理 53.2.1（Laurent 展开定理）}：设 \(f\) 在圆环
\(A = \{z \in \mathbb{C} \mid r < |z - z_0| < R\}\)（\(0 \leq r < R \leq \infty\)）上全纯。则
\(f\) 在 \(A\) 上可展开为 Laurent 级数：

\[
f(z) = \sum_{n=-\infty}^{\infty} a_n (z - z_0)^n, \quad r < |z - z_0| < R
\]

其中
\(a_n = \frac{1}{2\pi i} \int_{\partial D(z_0, \rho)} \frac{f(\zeta)}{(\zeta - z_0)^{n+1}} d\zeta\)（\(r < \rho < R\)，积分与
\(\rho\) 无关）。

\textbf{证明}：对 \(z \in A\)，取 \(r < r_1 < |z - z_0| < r_2 < R\)。由
Cauchy 积分公式（应用于环形区域），

\[
f(z) = \frac{1}{2\pi i} \int_{\partial D(z_0, r_2)} \frac{f(\zeta)}{\zeta - z} d\zeta - \frac{1}{2\pi i} \int_{\partial D(z_0, r_1)} \frac{f(\zeta)}{\zeta - z} d\zeta
\]

对第一个积分，\(|z - z_0| < |\zeta - z_0|\)，展开
\(\frac{1}{\zeta - z} = \sum_{n=0}^{\infty} \frac{(z - z_0)^n}{(\zeta - z_0)^{n+1}}\)。

对第二个积分，\(|z - z_0| > |\zeta - z_0|\)，展开
\(-\frac{1}{\zeta - z} = \frac{1}{z - \zeta} = \sum_{n=1}^{\infty} \frac{(\zeta - z_0)^{n-1}}{(z - z_0)^n}\)。

将两部分合并即得 Laurent 级数。\(\blacksquare\)

\bigskip\hrule\bigskip

\subsubsection{53.3
孤立奇点的分类}\label{ux5b64ux7acbux5947ux70b9ux7684ux5206ux7c7b}

\textbf{定义 53.3.1（孤立奇点）}：设 \(f\) 在 \(z_0\) 的去心邻域
\(\{0 < |z - z_0| < \delta\}\) 上全纯，但在 \(z_0\) 处不定义或不解析。则
\(z_0\) 称为 \(f\) 的\textbf{孤立奇点}。设 \(f\) 在 \(z_0\) 处的 Laurent
展开为 \(f(z) = \sum_{n=-\infty}^{\infty} a_n (z - z_0)^n\)。

\begin{enumerate}
\def\labelenumi{\arabic{enumi}.}
\tightlist
\item
  \textbf{可去奇点}：\(a_n = 0\)（\(\forall n < 0\)），即 Laurent
  级数实为 Taylor 级数
\item
  \textbf{极点}：存在 \(m \geq 1\) 使得 \(a_{-m} \neq 0\) 且
  \(a_n = 0\)（\(\forall n < -m\)）。\(m\) 称为\textbf{极点的阶}
\item
  \textbf{本性奇点}：存在无穷多个负指数 \(n\) 使得 \(a_n \neq 0\)
\end{enumerate}

\textbf{定理 53.3.1（Riemann 可去奇点定理）}：若 \(f\) 在 \(z_0\)
的去心邻域上全纯且有界，则 \(z_0\) 是可去奇点。

\textbf{证明}：设
\(|f(z)| \leq M\)。\(\forall n < 0\)，\(|a_n| = |\frac{1}{2\pi i} \int_{\partial D(z_0, \varepsilon)} \frac{f(\zeta)}{(\zeta - z_0)^{n+1}} d\zeta| \leq \frac{1}{2\pi} \cdot \frac{M}{\varepsilon^{n+1}} \cdot 2\pi\varepsilon = M \varepsilon^{-n}\)。令
\(\varepsilon \to 0\)（\(n < 0\) 时 \(-n > 0\)），得
\(a_n = 0\)。\(\blacksquare\)

\textbf{定理 53.3.2（Casorati-Weierstrass 定理）}：若 \(z_0\) 是 \(f\)
的本性奇点，则对任意 \(w \in \mathbb{C} \cup \{\infty\}\)，存在序列
\(z_n \to z_0\) 使得 \(f(z_n) \to w\)。即 \(f\) 在 \(z_0\)
的任何去心邻域中的像在 \(\mathbb{C}\) 中稠密。

\emph{证明}：反设存在 \(w \in \mathbb{C}\)、\(\varepsilon > 0\) 和
\(\delta > 0\) 使得 \(|f(z) - w| > \varepsilon\) 对所有
\(0 < |z - z_0| < \delta\) 成立。令 \(g(z) = \frac{1}{f(z) - w}\)，则
\(g\) 在 \(0 < |z - z_0| < \delta\)
上全纯且有界（\(|g(z)| \leq 1/\varepsilon\)）。由 Riemann
可去奇点定理（定理 53.3.1），\(g\) 可全纯延拓到 \(z_0\)。

若 \(g(z_0) \neq 0\)，则 \(f(z) = w + \frac{1}{g(z)}\) 在 \(z_0\)
处全纯，与 \(z_0\) 是本性奇点矛盾。

若 \(g(z_0) = 0\)，设 \(g\) 在 \(z_0\) 处的零点阶数为 \(m \geq 1\)，则
\(g(z) = (z - z_0)^m h(z)\)（\(h(z_0) \neq 0\)）。于是
\(f(z) - w = \frac{1}{(z - z_0)^m h(z)}\)，故 \(z_0\) 是 \(f\) 的 \(m\)
阶极点（或 \(w = \infty\) 时 \(f(z) \to \infty\)），同样与 \(z_0\)
是本性奇点矛盾。

因此假设不成立，即 \(f\) 在 \(z_0\) 的任何去心邻域中的像在
\(\mathbb{C}\) 中稠密。∎

\textbf{定理 53.3.3（Picard 大定理------陈述）}：若 \(z_0\) 是 \(f\)
的本性奇点，则 \(f\) 在 \(z_0\) 的任何去心邻域中取遍 \(\mathbb{C}\)
中的所有值，至多有一个例外。

\emph{注记}：Picard
大定理的证明超出本书范围，但可通过模函数方法（modular
function）得到。完整证明可参见 Ahlfors《Complex Analysis》或
Conway《Functions of One Complex Variable II》。

\bigskip\hrule\bigskip

\subsection{Ch 54 留数定理}\label{ch-54-ux7559ux6570ux5b9aux7406}

\subsubsection{54.1
留数的定义与计算}\label{ux7559ux6570ux7684ux5b9aux4e49ux4e0eux8ba1ux7b97}

\textbf{定义 54.1.1（留数）}：设 \(z_0\) 是 \(f\) 的孤立奇点，\(f\) 在
\(z_0\) 处的 Laurent 展开为
\(f(z) = \sum_{n=-\infty}^{\infty} a_n (z - z_0)^n\)。\(f\) 在 \(z_0\)
处的\textbf{留数}定义为

\[
\operatorname{Res}(f, z_0) = a_{-1} = \frac{1}{2\pi i} \int_{\partial D(z_0, \varepsilon)} f(z) \, dz
\]

\textbf{定理 54.1.1（留数的计算规则）}：

\begin{enumerate}
\def\labelenumi{(\roman{enumi})}
\tightlist
\item
  若 \(z_0\) 是 \(m\) 阶极点，则
  \(\operatorname{Res}(f, z_0) = \frac{1}{(m-1)!} \lim_{z \to z_0} \frac{d^{m-1}}{dz^{m-1}}[(z - z_0)^m f(z)]\)
\item
  若 \(z_0\) 是 1 阶极点（简单极点），则
  \(\operatorname{Res}(f, z_0) = \lim_{z \to z_0} (z - z_0)f(z)\)
\item
  若 \(f(z) = \frac{g(z)}{h(z)}\)，\(g, h\)
  全纯，\(g(z_0) \neq 0\)，\(h(z_0) = 0\)，\(h'(z_0) \neq 0\)（简单极点），则
  \(\operatorname{Res}(f, z_0) = \frac{g(z_0)}{h'(z_0)}\)
\end{enumerate}

\textbf{证明}：(i) 设
\(f(z) = \frac{a_{-m}}{(z - z_0)^m} + \cdots + \frac{a_{-1}}{z - z_0} + a_0 + \cdots\)。则
\((z - z_0)^m f(z) = a_{-m} + \cdots + a_{-1}(z - z_0)^{m-1} + \cdots\)。求导
\(m-1\) 次后取 \(z \to z_0\) 得 \(a_{-1}\)。

\begin{enumerate}
\def\labelenumi{(\roman{enumi})}
\setcounter{enumi}{1}
\item
  是 (i) 的特例（\(m = 1\)）。
\item
  由洛必达法则，\(a_{-1} = \lim_{z \to z_0} (z - z_0)\frac{g(z)}{h(z)} = \frac{g(z_0)}{h'(z_0)}\)。\(\blacksquare\)
\end{enumerate}

\bigskip\hrule\bigskip

\subsubsection{54.2 留数定理}\label{ux7559ux6570ux5b9aux7406}

\textbf{定理 54.2.1（留数定理）}：设 \(\Omega \subseteq \mathbb{C}\)
为单连通开集，\(\gamma\) 为 \(\Omega\) 中简单闭曲线（正向），\(f\) 在
\(\Omega\) 上除有限个孤立奇点 \(z_1, \ldots, z_k\)（均在 \(\gamma\)
内部）外全纯。则

\[
\int_{\gamma} f(z) \, dz = 2\pi i \sum_{j=1}^{k} \operatorname{Res}(f, z_j)
\]

\textbf{证明}：在每个 \(z_j\) 周围取小圆 \(\gamma_j\)（\(\gamma_j\) 在
\(\gamma\) 内部且互不相交）。由 Cauchy 积分定理（应用于 \(\gamma\) 和各
\(\gamma_j\) 之间的区域），

\[
\int_{\gamma} f(z) dz = \sum_{j=1}^{k} \int_{\gamma_j} f(z) dz = 2\pi i \sum_{j=1}^{k} \operatorname{Res}(f, z_j)
\]

其中第二步由留数的定义。\(\blacksquare\)

\bigskip\hrule\bigskip

\subsubsection{54.3
留数定理在实积分计算中的应用}\label{ux7559ux6570ux5b9aux7406ux5728ux5b9eux79efux5206ux8ba1ux7b97ux4e2dux7684ux5e94ux7528}

\textbf{定理
54.3.1（有理三角函数的积分）}：\(\int_0^{2\pi} R(\cos \theta, \sin \theta) d\theta\)（\(R\)
为有理函数）可通过代换 \(z = e^{i\theta}\)
化为单位圆周上的复积分，然后用留数定理计算。

\textbf{定理 54.3.2（无穷区间上的有理函数积分）}：设 \(P, Q\)
为多项式，\(\deg Q \geq \deg P + 2\)，\(Q\) 在实轴上无零点。则

\[
\int_{-\infty}^{\infty} \frac{P(x)}{Q(x)} \, dx = 2\pi i \sum_{\operatorname{Im} z_j > 0} \operatorname{Res}\left(\frac{P}{Q}, z_j\right)
\]

\textbf{证明}：取上半平面的大半圆围道
\(\Gamma_R = [-R, R] \cup C_R\)（\(C_R\) 为上半圆弧）。由于
\(\deg Q \geq \deg P + 2\)，\(|\frac{P(z)}{Q(z)}| \leq \frac{C}{|z|^2}\)（\(|z|\)
大时），故
\(\int_{C_R} \to 0\)（\(R \to \infty\)）。由留数定理，\(\int_{\Gamma_R} = 2\pi i \sum_{\operatorname{Im} z_j > 0} \operatorname{Res}\)。令
\(R \to \infty\) 即得。\(\blacksquare\)

\textbf{定理 54.3.3（Jordan
引理与含三角函数的积分）}：\(\int_{-\infty}^{\infty} \frac{P(x)}{Q(x)} e^{i\alpha x} dx\)（\(\alpha > 0\)）可通过
Jordan 引理 + 留数定理计算，常用于计算
\(\int_{-\infty}^{\infty} \frac{\cos \alpha x}{Q(x)} dx\) 和
\(\int_{-\infty}^{\infty} \frac{\sin \alpha x}{Q(x)} dx\) 类积分。

\bigskip\hrule\bigskip

\subsection{Ch 55 共形映射}\label{ch-55-ux5171ux5f62ux6620ux5c04}

\subsubsection{55.1
共形映射的定义与性质}\label{ux5171ux5f62ux6620ux5c04ux7684ux5b9aux4e49ux4e0eux6027ux8d28}

\textbf{定义 55.1.1（共形映射）}：全纯函数 \(f : \Omega \to \mathbb{C}\)
在 \(z_0 \in \Omega\) 处称为\textbf{共形}的，若 \(f'(z_0) \neq 0\)。若
\(f\) 在 \(\Omega\) 上每点共形且为单射，则称 \(f\) 为 \(\Omega\)
上的\textbf{共形映射}。

\textbf{定理 55.1.1（共形映射的几何意义）}：\(f\) 在 \(z_0\)
处共形当且仅当 \(f\) 在 \(z_0\) 处保角（两曲线在 \(z_0\)
处的交角等于其像曲线在 \(f(z_0)\) 处的交角）且伸缩率 \(|f'(z_0)|\)
与方向无关。

\textbf{证明}：由
\(f(z) - f(z_0) \approx f'(z_0)(z - z_0)\)（\(z \to z_0\)）。\(f'(z_0) = r e^{i\theta}\)
将 \(z - z_0\) 旋转 \(\theta\) 并缩放 \(r\)
倍，保角且伸缩率与方向无关。\(\blacksquare\)

\bigskip\hrule\bigskip

\subsubsection{55.2
分式线性变换}\label{ux5206ux5f0fux7ebfux6027ux53d8ux6362}

\textbf{定义 55.2.1（分式线性变换）}：形如

\[
T(z) = \frac{az + b}{cz + d}, \quad ad - bc \neq 0
\]

的映射称为\textbf{分式线性变换}（Möbius 变换）。它可延拓为
\(\hat{\mathbb{C}} = \mathbb{C} \cup \{\infty\}\)（Riemann
球面）上的双射。

\textbf{定理 55.2.1（分式线性变换的基本性质）}：

\begin{enumerate}
\def\labelenumi{(\roman{enumi})}
\tightlist
\item
  分式线性变换构成群 \(PSL(2, \mathbb{C})\)
\item
  分式线性变换将圆/直线映为圆/直线（保圆性）
\item
  分式线性变换保交比：\(\frac{T(z_1) - T(z_3)}{T(z_1) - T(z_4)} \cdot \frac{T(z_2) - T(z_4)}{T(z_2) - T(z_3)} = \frac{z_1 - z_3}{z_1 - z_4} \cdot \frac{z_2 - z_4}{z_2 - z_3}\)
\item
  存在唯一的分式线性变换将任意三个不同的点映为任意三个不同的点
\end{enumerate}

\textbf{证明}：(i) 复合和逆的公式验证。(ii) 圆/直线方程可写为
\(A z \overline{z} + B z + \overline{B} \overline{z} + C = 0\)（\(A, C \in \mathbb{R}\)），分式线性变换保持此形式。(iii)
直接验证。(iv) 由交比的不变性构造。\(\blacksquare\)

\bigskip\hrule\bigskip

\subsubsection{55.3 Riemann
映射定理}\label{riemann-ux6620ux5c04ux5b9aux7406}

\textbf{定理 55.3.1（Riemann 映射定理------陈述）}：设
\(\Omega \subsetneq \mathbb{C}\) 为单连通区域。则存在共形映射
\(f : \Omega \to D = \{z \in \mathbb{C} \mid |z| < 1\}\)（单位圆盘）。进一步，若指定
\(z_0 \in \Omega\) 和 \(f(z_0) = 0\)，\(f'(z_0) > 0\)，则 \(f\)
是唯一的。

\textbf{注记 55.3.1}：Riemann
映射定理是复分析的核心定理之一，证明了任何单连通真子区域与单位圆盘共形等价。这是二维特有的现象------高维情况不成立。完整证明需要正规族理论和
Schwarz 引理，在卷十四（泛函分析）中展开。

\textbf{例 55.3.1}：上半平面
\(\mathbb{H} = \{z \mid \operatorname{Im} z > 0\}\)
到单位圆盘的共形映射为 \(f(z) = \frac{z - i}{z + i}\)（Cayley 变换）。

\bigskip\hrule\bigskip

\subsection{Ch 56 解析延拓}\label{ch-56-ux89e3ux6790ux5ef6ux62d3}

\subsubsection{56.1
解析延拓的概念}\label{ux89e3ux6790ux5ef6ux62d3ux7684ux6982ux5ff5}

\textbf{定义 56.1.1（解析延拓）}：设 \(\Omega_1 \subseteq \Omega_2\) 为
\(\mathbb{C}\) 中的区域，\(f_1\) 在 \(\Omega_1\) 上全纯，\(f_2\) 在
\(\Omega_2\) 上全纯。若 \(f_2|_{\Omega_1} = f_1\)，则称 \(f_2\) 为
\(f_1\) 从 \(\Omega_1\) 到 \(\Omega_2\) 的\textbf{解析延拓}。

\textbf{定理 56.1.1（解析延拓的唯一性）}：解析延拓若存在，则唯一。

\textbf{证明}：设 \(f_2, g_2\) 均为 \(f_1\) 在 \(\Omega_2\)
上的解析延拓，则 \(h = f_2 - g_2\) 在 \(\Omega_2\) 上全纯，且在
\(\Omega_1\) 上
\(h = 0\)。由全纯函数的唯一性定理（见下文），\(h \equiv 0\) 在
\(\Omega_2\) 上。\(\blacksquare\)

\textbf{定理 56.1.2（唯一性定理）}：设 \(f\) 在区域 \(\Omega\)
上全纯。若存在 \(z_0 \in \Omega\) 和
\(\{z_n\}_{n=1}^{\infty} \subseteq \Omega \setminus \{z_0\}\) 使得
\(z_n \to z_0\) 且 \(f(z_n) = 0\)（\(\forall n\)），则 \(f \equiv 0\) 在
\(\Omega\) 上。

\textbf{证明}：设 \(f\) 在 \(z_0\) 处的 Taylor 展开为
\(f(z) = \sum_{k=0}^{\infty} a_k (z - z_0)^k\)。若 \(f\) 不恒为零，设
\(a_m\) 为第一个非零系数。则
\(f(z) = (z - z_0)^m (a_m + a_{m+1}(z - z_0) + \cdots) = (z - z_0)^m g(z)\)，其中
\(g(z_0) = a_m \neq 0\)。由 \(g\) 的连续性，\(g(z) \neq 0\) 在 \(z_0\)
附近。但 \(f(z_n) = 0\) 且 \(z_n \neq z_0\)，故 \(g(z_n) = 0\)，与
\(g(z_0) \neq 0\) 和连续性矛盾。因此所有 \(a_k = 0\)，\(f \equiv 0\) 在
\(z_0\) 的一个邻域内。由区域的连通性，\(f \equiv 0\) 在整个 \(\Omega\)
上。\(\blacksquare\)

\bigskip\hrule\bigskip

\subsubsection{56.2 自然边界}\label{ux81eaux7136ux8fb9ux754c}

\textbf{定义 56.2.1（自然边界）}：设 \(f\) 在 \(\Omega\) 上全纯。若
\(\Omega\) 的边界 \(\partial \Omega\) 上的每点都是 \(f\) 的奇点（即
\(f\) 不能延拓到该点的任何邻域），则称 \(\partial \Omega\) 为 \(f\)
的\textbf{自然边界}。

\textbf{例
56.2.1}：\(f(z) = \sum_{n=0}^{\infty} z^{n!} = 1 + z + z^2 + z^6 + z^{24} + \cdots\)
在 \(|z| < 1\) 上全纯，\(|z| = 1\) 为其自然边界（Hadamard 间隙定理）。

\bigskip\hrule\bigskip

\subsubsection{56.3 多值函数与 Riemann
面初步}\label{ux591aux503cux51fdux6570ux4e0e-riemann-ux9762ux521dux6b65}

\textbf{定义
56.3.1（解析延拓与多值函数）}：有些函数通过解析延拓会产生多值性。例如
\(f(z) = \sqrt{z}\) 在 \(z = 0\) 附近：从某点出发沿包含 \(0\)
的闭曲线绕一圈，函数值变为原来的 \(-1\) 倍。

\textbf{定义 56.3.2（Riemann 面）}：多值函数的自然定义域是
\textbf{Riemann
面}------一个二维复流形，使得''多值函数''在其上成为单值全纯函数。

例如 \(\sqrt{z}\) 的 Riemann 面是由两叶复平面在 \(z = 0\)
处通过分支切割（branch cut）连接而成的。

\textbf{例 56.3.1}：对数函数 \(\log z\) 的 Riemann
面是无穷多叶螺旋面。\(\log z = \log|z| + i\arg z\)，绕原点一圈
\(\arg z\) 增加 \(2\pi\)，进入下一叶。

\textbf{注记 56.3.1}：Riemann
面的系统理论将在卷十二（微分几何）中作为一维复流形展开，并将在卷十六（代数几何）中深化为代数曲线理论。\(\blacksquare\)

\bigskip\hrule\bigskip

\subsection{Ch 57 值分布理论（Nevanlinna
理论）}\label{ch-57-ux503cux5206ux5e03ux7406ux8bbanevanlinna-ux7406ux8bba}

Nevanlinna 理论是复分析中研究亚纯函数值分布的核心框架，由 R. Nevanlinna
于 1925 年创立，将 Picard 定理从定性陈述提升为精密的定量理论。任给
\(\mathbb{C}\) 上的亚纯函数 \(f\) 及
\(a \in \hat{\mathbb{C}} = \mathbb{C} \cup \{\infty\}\)，引入计数函数
\(N(r, a, f)\)（计及重数的 \(a\) 值点个数的对数平均）与逼近函数
\(m(r, a, f)\)（\(f\) 在圆周 \(|z| = r\) 上接近 \(a\)
的程度的对数平均）。定义\textbf{特征函数}
\(T(r, f) = m(r, \infty, f) + N(r, \infty, f)\)，它刻画 \(f\)
的整体增长速率。\textbf{第一基本定理}断言：对任意
\(a \in \hat{\mathbb{C}}\)，

\[T(r, f) = m(r, a, f) + N(r, a, f) + O(1)\]

即和 \(m(r,a) + N(r,a)\) 与目标值 \(a\) 无关（至多相差有界量），表明
\(f\)
以相同速率趋向每个值。\textbf{第二基本定理}是理论的核心深度结果：对任意
\(q \geq 3\) 个互异值 \(a_1, \ldots, a_q \in \hat{\mathbb{C}}\)，

\[\sum_{\nu=1}^{q} m(r, a_\nu, f) \leq 2T(r, f) - N_1(r) + S(r, f)\]

其中 \(N_1(r)\) 为分支点修正项，\(S(r, f) = O(\log r T(r, f))\)
为误差项（在 \(r \to \infty\)
除去一个有限测度集外成立）。作为直接推论，超越亚纯函数至多有两个 Picard
例外值（Picard
小定理），且其完全重数分支点至多可数。该理论深刻揭示了亚纯函数的''均值分布''本质------它不能同时极度靠近多个不同的值，这一现象与复平面的双曲几何结构紧密相关。

\textbf{定义 57.1}（Nevanlinna 特征函数）：设 \(f\) 为 \(\mathbb{C}\)
上非常数的亚纯函数。对 \(r > 0\)，定义：靠近函数
\(m(r, f) = \frac{1}{2\pi} \int_0^{2\pi} \log^+ |f(r e^{i\theta})| d\theta\)（其中
\(\log^+ x = \max(\log x, 0)\)）；计数函数
\(N(r, f) = \int_0^r \frac{n(t, f) - n(0, f)}{t} dt + n(0, f) \log r\)（其中
\(n(t, f)\) 为 \(f\) 在 \(|z| \leq t\)
中的极点数，计重数）。\textbf{Nevanlinna 特征函数}定义为
\(T(r, f) = m(r, f) + N(r, f)\)。\(T(r, f)\) 衡量 \(f\)
的增长速度------若 \(f\) 为有理函数，则 \(T(r, f) = O(\log r)\)；若
\(f\) 为超越整函数，则 \(T(r, f) / \log r \to \infty\)。

\textbf{定理 57.1}（第一基本定理）：对任意
\(a \in \mathbb{C} \cup \{\infty\}\)，有
\(m(r, \frac{1}{f-a}) + N(r, \frac{1}{f-a}) = T(r, f) + O(1)\)（\(r \to \infty\)）。即：\(T(r, \frac{1}{f-a}) = T(r, f) + O(1)\)------特征函数在
Möbius 变换下几乎不变。

\textbf{定理 57.2}（第二基本定理 / Nevanlinna）：对互异的
\(a_1, \ldots, a_q \in \mathbb{C} \cup \{\infty\}\)，有
\((q-2) T(r, f) \leq \sum_{j=1}^q \overline{N}(r, \frac{1}{f-a_j}) + S(r, f)\)，其中
\(\overline{N}\) 是不计重数的计数函数，\(S(r, f) = o(T(r, f))\)
对几乎所有 \(r\) 成立。\emph{证明思路}：利用对数导数的引理
\(\int_0^{2\pi} \log^+ |f'/f| d\theta = o(T(r, f))\)，结合对
\(F = \prod_{j=1}^q (f - a_j)\) 应用第一基本定理和 Jensen 公式。

\textbf{推论 57.1}（Picard 小定理作为特例）：若整函数 \(f\)
避开两个有穷值，则第二基本定理给出
\((2-2) T(r, f) \leq 0 + o(T(r, f))\)，矛盾于
\(\log r / T(r, f) \to 0\)。因此 \(f\) 必须是常数------Picard
小定理得证。

\bigskip\hrule\bigskip

\subsection{Ch 58 拟共形映射与 Teichmüller
空间}\label{ch-58-ux62dfux5171ux5f62ux6620ux5c04ux4e0e-teichmuxfcller-ux7a7aux95f4}

\subsubsection{58.1 拟共形映射}\label{ux62dfux5171ux5f62ux6620ux5c04}

拟共形映射是共形映射的自然推广，允许有界的角度畸变。其解析定义基于
\textbf{Beltrami 方程}：

\[f_{\bar{z}} = \mu(z) f_z\]

其中 \(\mu(z)\) 为 \textbf{Beltrami 系数}（复值可测函数），满足
\(\|\mu\|_\infty = \operatorname{ess\,sup} |\mu(z)| < 1\)。当
\(\mu \equiv 0\) 时方程退化为 Cauchy-Riemann 方程，\(f\)
为全纯函数；一般 \(\mu \neq 0\) 时 \(f\) 将无穷小圆映为椭圆，其偏心率由
\(|\mu|\) 决定。\(\mu\) 的 \(L^\infty\)
范数度量了映射偏离共形的程度。核心结果------\textbf{可测 Riemann
映射定理}（Morrey, 1938；Ahlfors-Bers, 1960）断言：对复平面上任意满足
\(\|\mu\|_\infty < 1\) 的可测 Beltrami 系数
\(\mu\)，存在唯一的拟共形同胚
\(f^\mu : \hat{\mathbb{C}} \to \hat{\mathbb{C}}\) 满足 Beltrami
方程且固定 \(0, 1, \infty\)。该定理是复分析中 Riemann
映射定理的深远推广，为 Teichmüller 理论、Klein
群形变空间和复动力系统提供了不可或缺的分析工具。

\subsubsection{58.2 Teichmüller
空间简介}\label{teichmuxfcller-ux7a7aux95f4ux7b80ux4ecb}

\textbf{Teichmüller 空间} \(\mathcal{T}_g\) 是亏格 \(g\) 的紧 Riemann
面复结构的形变空间。对 \(g \geq 2\)，\(\mathcal{T}_g\) 是 \(3g-3\)
维复流形（实维数
\(6g-6\)）。其标准定义利用拟共形映射：取定参考曲面（标记 Riemann
面）\(S_g\)，令

\[\mathcal{T}_g = \{(X, f) \mid f : S_g \to X \text{ 为拟共形同胚}\} / \sim\]

其中 \((X, f) \sim (Y, g)\) 若 \(g \circ f^{-1} : X \to Y\)
同伦于共形映射。等价地，\(\mathcal{T}_g\) 可通过 \textbf{Bers
嵌入}实现为 \(\mathbb{C}^{3g-3}\)
中的有界域------将曲面上的全纯二次微分对应于某个单值化群形变。Teichmüller
度量 \(d_T\)（基于极小拟共形映射的极大伸缩商的对数）将 \(\mathcal{T}_g\)
变为完备度量空间。模空间
\(\mathcal{M}_g = \mathcal{T}_g / \mathrm{MCG}(S_g)\) 则是
\(\mathcal{T}_g\) 关于映射类群作用的商，为一个 \(3g-3\)
维复轨道空间（orbifold），在双曲几何、弦论和代数几何的模空间理论中扮演核心角色。\(\blacksquare\)

\bigskip\hrule\bigskip

\subsection{Ch 59：Siegel
区域与对称域}\label{ch-59siegel-ux533aux57dfux4e0eux5bf9ux79f0ux57df}

Siegel
区域理论是经典多复变函数论的核心部分。\textbf{定义}：\(\mathbb{C}^n\)
中的 \textbf{Siegel 上半空间}（或 Siegel 广义上半平面）为
\(\mathcal{H}_n = \{ Z \in \mathbb{C}^{n \times n} : Z^{\mathsf{T}} = Z, \; \operatorname{Im}(Z) \succ 0 \}\)（即
\(n \times n\) 对称复矩阵虚部正定的集合），它是
\(\operatorname{Sp}(2n,\mathbb{R})\) 的作用空间，等价于 \textbf{Siegel
上半空间}
\(\operatorname{Sp}(2n,\mathbb{R}) / \operatorname{U}(n)\)。更一般地，\textbf{Siegel
区域}是第二类域：\(\{ (z,w) \in \mathbb{C}^{m} \times \mathbb{C}^{n} : \operatorname{Im}(z) - F(w,w) \in \Omega \}\)，其中
\(F\) 是 \(\mathbb{C}^{n} \times \mathbb{C}^{n} \to \mathbb{C}^{m}\) 的
Hermite 形式，\(\Omega \subset \mathbb{R}^{m}\) 是自对偶锥。Siegel
域给出了所有齐性有界域的典型表示。\textbf{Hermann 有界对称域分类}（E.
Cartan, 1935）：不可约有界对称域分为四大类经典域（由 Siegel
上半空间、Grassmann 流形、正交 Lie 球的 Cayley
变换得到）和两个例外域（16维和27维）。Siegel 上半空间上的 \textbf{Siegel
模形式}（即 \(\operatorname{Sp}(2n,\mathbb{Z})\)
作用下变换的全纯函数）推广了经典椭圆模形式，在数论（志村簇）和弦论紧化中有深刻应用。

\bigskip\hrule\bigskip

\subsection{Ch 60：多复变函数论导引（Several Complex
Variables）}\label{ch-60ux591aux590dux53d8ux51fdux6570ux8bbaux5bfcux5f15several-complex-variables}

多复变函数论研究
\(\mathbb{C}^n\)（\(n \geq 2\)）上的全纯函数，与单复变有本质区别。\textbf{基本现象}：（1）\textbf{Hartogs
现象}------存在区域 \(\Omega \subsetneq \mathbb{C}^n\) 使得 \(\Omega\)
上的所有全纯函数均可全纯延拓到更大的固定区域上。例如，对
\(n \geq 2\)，每个在 \(0 < \|z\| < 1\)
上全纯的函数自动全纯延拓到整个单位球
\(\|z\| < 1\)------即孤立奇点不能作为全纯函数的障碍。这直接导致
\(\mathbb{C}^n\) 中不存在''一般区域''的 Riemann
映射定理的类比。（2）\textbf{域的等价分类}：单复变中 Riemann
映射定理宣告所有单连通真子域彼此双全纯等价。多复变中双全纯等价类极其丰富------单位球
\(B_n = \{ \|z\| < 1 \}\) 和单位多圆柱 \(D^n = \{ |z_j| < 1 \}\)
并\textbf{不}双全纯等价，尽管两者都是 \(\mathbb{C}^n\) 中的自然域。

\textbf{基本定义}：函数 \(f: \Omega \to \mathbb{C}\) 称为\textbf{全纯}若
\(f\) 在每点满足 Cauchy-Riemann 方程组
\(\frac{\partial f}{\partial \bar{z}_j} = 0\)（\(j = 1,\ldots,n\)）。等价地，\(f\)
在每个变量单独全纯。\textbf{核心定理}：（1）\textbf{Cauchy
积分公式}（多复变版本）------\(f(z) = \frac{1}{(2\pi i)^n} \int_{|w_1-z_1| = r_1} \cdots \int_{|w_n - z_n| = r_n} \frac{f(w)}{(w_1 - z_1)\cdots(w_n - z_n)} dw_1 \cdots dw_n\)
（对于乘积圆盘）；（2）\textbf{Osgood
引理}------单独连续且对每个变量单独全纯的函数是全纯的；（3）\textbf{Hartogs
基本定理}------\(f\) 全纯当且仅当 \(f\)
对每个变量单独全纯（无需连续性假设）。

\textbf{Levi 问题}（E. E. Levi, 1911）：\(\mathbb{C}^n\)
中的哪些域是全纯域（即存在不能延拓到更大域的全纯函数）？\textbf{答案}（Oka,
1942; Bremermann, 1954; Norguet, 1954）：域是全纯域当且仅当它是
\textbf{伪凸} 的。伪凸性由 \textbf{Levi 形式}（边界上复 Hessian
矩阵）的正半定性刻画------类似于 \(\mathbb{R}^n\)
中域的凸性，但由复结构的约束导致。多复变中的 \textbf{\(\bar{\partial}\)
问题}（即 \(\bar{\partial} u = f\)）及其正则解（Hörmander 的 \(L^2\)
方法）正是处理这些问题的基本技术工具。

\newpage

\section{卷十：拓扑学}\label{ux5377ux5341ux62d3ux6251ux5b66}

\begin{quote}
拓扑学研究在连续变形下保持不变的几何性质。本卷从度量空间中的极限和连续概念出发，抽象出拓扑空间的一般定义，建立连通性、紧致性、分离公理等核心概念，进而引入基本群和覆叠空间理论，为后续学习微分几何、代数拓扑、泛函分析等铺平道路。
\end{quote}

\bigskip\hrule\bigskip

\subsection{第57章：拓扑空间基础}\label{ux7b2c57ux7ae0ux62d3ux6251ux7a7aux95f4ux57faux7840}

拓扑空间将度量空间中「开集」的概念抽象为一组公理，使得我们可以在没有距离概念的情况下讨论连续性和极限。

\subsubsection{57.1
拓扑空间的定义}\label{ux62d3ux6251ux7a7aux95f4ux7684ux5b9aux4e49}

\textbf{定义 57.1}（拓扑空间）：设 \(X\)
是一个集合，\(\mathcal{T} \subseteq \mathcal{P}(X)\) 是 \(X\)
的子集族。称 \(\mathcal{T}\) 为 \(X\)
上的一个\textbf{拓扑}，如果满足以下三条公理：

\begin{enumerate}
\def\labelenumi{\arabic{enumi}.}
\tightlist
\item
  \textbf{(T1)} \(\emptyset \in \mathcal{T}\)，\(X \in \mathcal{T}\)。
\item
  \textbf{(T2)} \(\mathcal{T}\) 中任意多个元素的并集仍属于
  \(\mathcal{T}\)：若
  \(\{U_\alpha\}_{\alpha \in A} \subseteq \mathcal{T}\)，则
  \(\bigcup_{\alpha \in A} U_\alpha \in \mathcal{T}\)。
\item
  \textbf{(T3)} \(\mathcal{T}\) 中有限多个元素的交集仍属于
  \(\mathcal{T}\)：若 \(U_1, U_2, \ldots, U_n \in \mathcal{T}\)，则
  \(U_1 \cap U_2 \cap \cdots \cap U_n \in \mathcal{T}\)。
\end{enumerate}

称 \((X, \mathcal{T})\) 为一个\textbf{拓扑空间}，\(\mathcal{T}\)
中的元素称为\textbf{开集}。在不引起混淆时简记为 \(X\)。

\textbf{定义 57.2}（闭集）：设 \((X, \mathcal{T})\) 为拓扑空间。子集
\(F \subseteq X\) 称为\textbf{闭集}，如果
\(X \setminus F \in \mathcal{T}\)（即 \(F\) 的补集是开集）。

\textbf{命题 57.1}（闭集的性质）：设 \((X, \mathcal{T})\)
为拓扑空间，则： 1. \(\emptyset\) 和 \(X\) 是闭集。 2.
任意多个闭集的交集是闭集。 3. 有限多个闭集的并集是闭集。

\emph{证明}：由 De Morgan 律直接可得。∎

\textbf{定义 57.3}（邻域）：设 \(X\) 为拓扑空间，\(x \in X\)。子集
\(N \subseteq X\) 称为 \(x\) 的一个\textbf{邻域}，如果存在开集
\(U \in \mathcal{T}\) 使得 \(x \in U \subseteq N\)。\(x\)
的所有邻域构成的集合称为 \(x\) 的\textbf{邻域系}，记作
\(\mathcal{N}(x)\)。

\textbf{定义 57.4}（内部、闭包、边界）：设 \(X\)
为拓扑空间，\(A \subseteq X\)。

\begin{itemize}
\tightlist
\item
  \(A\) 的\textbf{内部} \(\operatorname{int}(A)\)（或
  \(A^\circ\)）：包含在 \(A\) 中的所有开集的并集，即 \(A\)
  的最大开子集。
\item
  \(A\) 的\textbf{闭包} \(\overline{A}\)（或
  \(\operatorname{cl}(A)\)）：包含 \(A\) 的所有闭集的交集，即包含 \(A\)
  的最小闭集。
\item
  \(A\) 的\textbf{边界}
  \(\partial A\)：\(\partial A = \overline{A} \setminus \operatorname{int}(A)\)。
\end{itemize}

\textbf{命题 57.2}（内部与闭包的关系）：设 \(X\)
为拓扑空间，\(A \subseteq X\)。则： 1. \(x \in \operatorname{int}(A)\)
当且仅当 \(A \in \mathcal{N}(x)\)。 2. \(x \in \overline{A}\) 当且仅当
\(x\) 的每个邻域与 \(A\) 相交。 3.
\(\overline{A} = X \setminus \operatorname{int}(X \setminus A)\)，\(\operatorname{int}(A) = X \setminus \overline{X \setminus A}\)。

\emph{证明}：(1) 由内部和邻域的定义直接可得。(2) 若
\(x \notin \overline{A}\)，则存在闭集 \(F \supseteq A\) 使
\(x \notin F\)，于是 \(X \setminus F\) 是 \(x\) 的邻域且与 \(A\)
不相交。反之亦然。(3) 由 De Morgan 律可得。∎

\textbf{定义 57.5}（聚点、孤立点）：设 \(X\)
为拓扑空间，\(A \subseteq X\)，\(x \in X\)。

\begin{itemize}
\tightlist
\item
  \(x\) 称为 \(A\) 的\textbf{聚点}（或极限点），如果 \(x\)
  的每个邻域都包含 \(A\) 中异于 \(x\) 的点。
\item
  \(A\) 的所有聚点构成的集合称为 \(A\) 的\textbf{导集}，记作 \(A'\)。
\item
  \(x \in A\) 称为 \(A\) 的\textbf{孤立点}，如果存在 \(x\) 的邻域 \(U\)
  使得 \(U \cap A = \{x\}\)。
\end{itemize}

\textbf{命题 57.3}：\(\overline{A} = A \cup A'\)。

\emph{证明}：\(x \in \overline{A}\) 当且仅当 \(x\) 的每个邻域与 \(A\)
相交，这等价于 \(x \in A\) 或 \(x\) 的每个邻域包含 \(A\) 中异于 \(x\)
的点（即 \(x \in A'\)）。∎

\subsubsection{57.2
拓扑的基与子基}\label{ux62d3ux6251ux7684ux57faux4e0eux5b50ux57fa}

\textbf{定义 57.6}（基）：设 \((X, \mathcal{T})\) 为拓扑空间。子族
\(\mathcal{B} \subseteq \mathcal{T}\) 称为 \(\mathcal{T}\)
的一个\textbf{基}，如果 \(\mathcal{T}\) 中的每个开集都可以表示为
\(\mathcal{B}\) 中某些元素的并集。

\textbf{定理 57.4}（基的判定条件）：设 \(X\)
为集合，\(\mathcal{B} \subseteq \mathcal{P}(X)\)。则 \(\mathcal{B}\) 是
\(X\) 上某个拓扑的基当且仅当： 1.
\(\bigcup_{B \in \mathcal{B}} B = X\)。 2. 对任意
\(B_1, B_2 \in \mathcal{B}\) 和 \(x \in B_1 \cap B_2\)，存在
\(B_3 \in \mathcal{B}\) 使得 \(x \in B_3 \subseteq B_1 \cap B_2\)。

此时由 \(\mathcal{B}\) 生成的拓扑 \(\mathcal{T}\)
为：\(U \in \mathcal{T}\) 当且仅当对每个 \(x \in U\)，存在
\(B \in \mathcal{B}\) 使得 \(x \in B \subseteq U\)。

\emph{证明}：(\(\Rightarrow\)) 若 \(\mathcal{B}\) 是拓扑 \(\mathcal{T}\)
的基，则 \(X \in \mathcal{T}\) 可表示为 \(\mathcal{B}\) 中元素的并，故
(1) 成立。对
\(B_1, B_2 \in \mathcal{B}\)，\(B_1 \cap B_2 \in \mathcal{T}\)，故可表示为
\(\mathcal{B}\) 中元素的并，从而 (2) 成立。

(\(\Leftarrow\)) 设条件成立。定义 \(\mathcal{T}\)
为满足所述条件的子集族。验证 \(\mathcal{T}\) 是拓扑： - \(\emptyset\)
平凡地属于 \(\mathcal{T}\)；由 (1) 知 \(X \in \mathcal{T}\)。 -
任意并：若 \(\{U_\alpha\} \subseteq \mathcal{T}\)，对
\(x \in \bigcup U_\alpha\)，存在 \(\alpha\) 使
\(x \in U_\alpha\)，故存在 \(B \in \mathcal{B}\) 使
\(x \in B \subseteq U_\alpha \subseteq \bigcup U_\alpha\)。 - 有限交：对
\(U_1, U_2 \in \mathcal{T}\)，任取 \(x \in U_1 \cap U_2\)，存在
\(B_1, B_2 \in \mathcal{B}\) 使
\(x \in B_1 \subseteq U_1\)，\(x \in B_2 \subseteq U_2\)。由 (2) 存在
\(B_3 \in \mathcal{B}\) 使
\(x \in B_3 \subseteq B_1 \cap B_2 \subseteq U_1 \cap U_2\)。归纳得有限交。

显然 \(\mathcal{B}\) 是 \(\mathcal{T}\) 的基。∎

\textbf{定义 57.7}（子基）：设 \(X\)
为集合，\(\mathcal{S} \subseteq \mathcal{P}(X)\) 满足
\(\bigcup_{S \in \mathcal{S}} S = X\)。则 \(\mathcal{S}\)
中元素的所有有限交构成的族 \(\mathcal{B}\) 是某个拓扑的基，由
\(\mathcal{B}\) 生成的拓扑称为由 \(\mathcal{S}\)
\textbf{生成的拓扑}，\(\mathcal{S}\) 称为该拓扑的\textbf{子基}。

\subsubsection{57.3
常见拓扑举例}\label{ux5e38ux89c1ux62d3ux6251ux4e3eux4f8b}

\textbf{例
57.1}（离散拓扑）：\(\mathcal{T} = \mathcal{P}(X)\)（所有子集都是开集）。此时每个单点集
\(\{x\}\) 是开集。

\textbf{例
57.2}（平凡拓扑/密着拓扑）：\(\mathcal{T} = \{\emptyset, X\}\)。只有空集和全空间是开集。

\textbf{例 57.3}（度量拓扑）：设 \((X, d)\) 为度量空间。以所有开球
\(B(x, r) = \{y \in X : d(x, y) < r\}\)（\(x \in X, r > 0\)）为基生成的拓扑，称为由度量
\(d\) \textbf{诱导的拓扑}。此时 \(U \subseteq X\) 是开集当且仅当对每个
\(x \in U\)，存在 \(r > 0\) 使 \(B(x, r) \subseteq U\)。

\textbf{定义 57.8}（可度量化空间）：拓扑空间 \((X, \mathcal{T})\)
称为\textbf{可度量化的}，如果存在 \(X\) 上的度量 \(d\) 使得
\(\mathcal{T}\) 恰为 \(d\) 诱导的拓扑。

\textbf{例
57.4}（有限补拓扑/余有限拓扑）：\(\mathcal{T} = \{U \subseteq X : X \setminus U \text{ 是有限集}\} \cup \{\emptyset\}\)。闭集恰为有限集和
\(X\)。

\textbf{例
57.5}（可数补拓扑）：\(\mathcal{T} = \{U \subseteq X : X \setminus U \text{ 是可数集}\} \cup \{\emptyset\}\)。

\textbf{例 57.6}（Sierpiński
空间）：\(X = \{0, 1\}\)，\(\mathcal{T} = \{\emptyset, \{1\}, \{0, 1\}\}\)。

\textbf{例 57.7}（序拓扑）：设 \((X, \leq)\) 为全序集。以所有开区间
\((a, b) = \{x \in X : a < x < b\}\) 以及（若存在最小元）\([a_0, b)\)
和（若存在最大元）\((a, b_0]\) 为基生成的拓扑，称为\textbf{序拓扑}。

\subsubsection{57.4
连续映射与同胚}\label{ux8fdeux7eedux6620ux5c04ux4e0eux540cux80da}

\textbf{定义 57.9}（连续映射）：设 \(X, Y\) 为拓扑空间，\(f: X \to Y\)
为映射。称 \(f\) 在点 \(x \in X\) 处\textbf{连续}，如果对 \(f(x)\)
的任意邻域 \(V\)，\(f^{-1}(V)\) 是 \(x\) 的邻域。称 \(f\)
为\textbf{连续映射}，如果 \(f\) 在每一点处连续。

\textbf{定理 57.5}（连续映射的等价刻画）：设 \(f: X \to Y\)
为拓扑空间之间的映射，则以下条件等价： 1. \(f\) 是连续映射。 2. 对 \(Y\)
中任意开集 \(V\)，\(f^{-1}(V)\) 是 \(X\) 中的开集。 3. 对 \(Y\)
中任意闭集 \(F\)，\(f^{-1}(F)\) 是 \(X\) 中的闭集。 4. 对任意
\(A \subseteq X\)，\(f(\overline{A}) \subseteq \overline{f(A)}\)。 5. 对
\(Y\) 的拓扑的某个基（或子基）\(\mathcal{B}\)，对每个
\(B \in \mathcal{B}\)，\(f^{-1}(B)\) 是 \(X\) 中的开集。

\emph{证明}：(1) \(\Rightarrow\) (2)：设 \(V \subseteq Y\)
为开集，\(x \in f^{-1}(V)\)。则 \(f(x) \in V\)，故 \(V\) 是 \(f(x)\)
的邻域。由连续性，\(f^{-1}(V)\) 是 \(x\) 的邻域，故存在开集
\(U_x \subseteq f^{-1}(V)\) 使 \(x \in U_x\)。于是
\(f^{-1}(V) = \bigcup_{x \in f^{-1}(V)} U_x\) 是开集。

\begin{enumerate}
\def\labelenumi{(\arabic{enumi})}
\setcounter{enumi}{1}
\item
  \(\Rightarrow\)
  (3)：\(f^{-1}(Y \setminus F) = X \setminus f^{-1}(F)\)，若 \(F\)
  闭，则 \(Y \setminus F\) 开，故 \(X \setminus f^{-1}(F)\) 开，即
  \(f^{-1}(F)\) 闭。
\item
  \(\Rightarrow\) (4)：\(f^{-1}(\overline{f(A)})\) 是包含 \(A\)
  的闭集，故 \(\overline{A} \subseteq f^{-1}(\overline{f(A)})\)，即
  \(f(\overline{A}) \subseteq \overline{f(A)}\)。
\item
  \(\Rightarrow\) (1)：设 \(x \in X\)，\(V\) 是 \(f(x)\) 的邻域。令
  \(A = X \setminus f^{-1}(V)\)。若 \(f(x)\) 不是 \(f(A)\)
  的聚点，则存在 \(f(x)\) 的开邻域 \(W \subseteq V\) 与 \(f(A)\)
  不相交。于是 \(f^{-1}(W) \subseteq f^{-1}(V)\) 是 \(x\) 的邻域，故
  \(f\) 在 \(x\) 处连续。假设 \(f(x) \in \overline{f(A)}\)，则
  \(x \in f^{-1}(\overline{f(A)}) \supseteq \overline{A}\)，故 \(x\)
  的每个邻域与 \(A\) 相交，这与 \(f^{-1}(V)\) 是 \(x\) 的邻域矛盾。因此
  \(f^{-1}(V)\) 是 \(x\) 的邻域。
\item
  \(\Leftrightarrow\) (5)：由基的定义直接可得。∎
\end{enumerate}

\textbf{命题 57.6}（连续映射的复合）：若 \(f: X \to Y\) 和
\(g: Y \to Z\) 都是连续映射，则 \(g \circ f: X \to Z\) 也是连续映射。

\emph{证明}：对 \(Z\) 中任意开集
\(W\)，\((g \circ f)^{-1}(W) = f^{-1}(g^{-1}(W))\)。由 \(g\)
连续，\(g^{-1}(W)\) 是 \(Y\) 中开集；由 \(f\)
连续，\(f^{-1}(g^{-1}(W))\) 是 \(X\) 中开集。∎

\textbf{定义 57.10}（同胚）：设 \(X, Y\) 为拓扑空间。映射 \(f: X \to Y\)
称为\textbf{同胚}（homeomorphism），如果 \(f\) 是双射，且 \(f\) 和
\(f^{-1}\) 都是连续映射。若存在 \(X\) 到 \(Y\) 的同胚，则称 \(X\) 与
\(Y\) \textbf{同胚}，记作 \(X \cong Y\)（或 \(X \approx Y\)）。

同胚是拓扑空间范畴中的同构。同胚的空间具有完全相同的拓扑性质。

\textbf{定义 57.11}（开映射、闭映射）：映射 \(f: X \to Y\)
称为\textbf{开映射}（resp. \textbf{闭映射}），如果 \(f\) 将 \(X\)
中的开集（resp. 闭集）映为 \(Y\) 中的开集（resp. 闭集）。

\textbf{命题 57.7}：设 \(f: X \to Y\) 为连续双射，则 \(f\)
是同胚当且仅当 \(f\) 是开映射，也当且仅当 \(f\) 是闭映射。

\emph{证明}：\(f^{-1}\) 连续等价于对 \(X\) 中任意开集
\(U\)，\((f^{-1})^{-1}(U) = f(U)\) 是 \(Y\) 中开集，即 \(f\)
是开映射。闭映射情形类似。∎

\subsubsection{57.5 子空间拓扑}\label{ux5b50ux7a7aux95f4ux62d3ux6251}

\textbf{定义 57.12}（子空间拓扑）：设 \((X, \mathcal{T})\)
为拓扑空间，\(Y \subseteq X\)。定义

\[\mathcal{T}_Y = \{U \cap Y : U \in \mathcal{T}\}\]

则 \(\mathcal{T}_Y\) 是 \(Y\)
上的拓扑，称为\textbf{子空间拓扑}（或相对拓扑）。\((Y, \mathcal{T}_Y)\)
称为 \((X, \mathcal{T})\) 的\textbf{拓扑子空间}。

\textbf{命题 57.8}（子空间拓扑的性质）：设 \(Y\) 是 \(X\) 的子空间。 1.
若 \(\mathcal{B}\) 是 \(X\) 的拓扑的基，则
\(\{B \cap Y : B \in \mathcal{B}\}\) 是 \(Y\) 的拓扑的基。 2. \(Y\)
中的闭集恰为 \(F \cap Y\)，其中 \(F\) 是 \(X\) 中的闭集。 3. 包含映射
\(i: Y \hookrightarrow X\)（\(i(y) = y\)）是连续的。 4. 若
\(Z \subseteq Y \subseteq X\)，则 \(Z\) 作为 \(Y\) 的子空间的拓扑与作为
\(X\) 的子空间的拓扑相同。

\emph{证明}：(1)-(3) 由定义直接可得。(4) \(Z\) 中开集为
\(\{U \cap Z : U \in \mathcal{T}\} = \{(U \cap Y) \cap Z : U \in \mathcal{T}\}\)，恰为
\(Z\) 作为 \(Y\) 子空间的拓扑。∎

\textbf{定理 57.9}（子空间上的连续映射）：设 \(f: X \to Y\)
连续，\(A \subseteq X\)。则限制映射 \(f|_A: A \to Y\)（赋予 \(A\)
子空间拓扑）也是连续的。

\emph{证明}：对 \(Y\) 中开集
\(V\)，\((f|_A)^{-1}(V) = f^{-1}(V) \cap A\)，而 \(f^{-1}(V)\) 是 \(X\)
中开集，故 \(f^{-1}(V) \cap A\) 是 \(A\) 中开集。∎

\subsubsection{57.6 积空间}\label{ux79efux7a7aux95f4}

\textbf{定义 57.13}（积拓扑------有限积）：设 \(X_1, X_2, \ldots, X_n\)
为拓扑空间。在积集 \(X_1 \times X_2 \times \cdots \times X_n\)
上，以所有形如 \(U_1 \times U_2 \times \cdots \times U_n\)（\(U_i\) 是
\(X_i\) 中开集）的集合为基生成的拓扑，称为\textbf{积拓扑}。

\textbf{定义 57.14}（积拓扑------任意积）：设
\(\{X_\alpha\}_{\alpha \in A}\)
为一族拓扑空间，\(X = \prod_{\alpha \in A} X_\alpha\)。\(X\)
上的\textbf{积拓扑}是以

\[\mathcal{S} = \{\pi_\alpha^{-1}(U_\alpha) : \alpha \in A, U_\alpha \subseteq X_\alpha \text{ 为开集}\}\]

为子基生成的拓扑，其中 \(\pi_\alpha: X \to X_\alpha\) 是到第 \(\alpha\)
个坐标的投影映射。

等价地，积拓扑的基由形如 \(\prod_{\alpha \in A} U_\alpha\)
的集合组成，其中每个 \(U_\alpha\) 是 \(X_\alpha\) 的开集，且除有限多个
\(\alpha\) 外 \(U_\alpha = X_\alpha\)。

\textbf{命题 57.10}（积拓扑的泛性质）：设
\(X = \prod_{\alpha \in A} X_\alpha\) 赋予积拓扑。则： 1. 每个投影
\(\pi_\alpha: X \to X_\alpha\) 是连续开映射。 2. 映射 \(f: Y \to X\)
连续当且仅当对每个
\(\alpha \in A\)，\(\pi_\alpha \circ f: Y \to X_\alpha\) 连续。

\emph{证明}：(1) 对 \(X_\alpha\) 中开集
\(U_\alpha\)，\(\pi_\alpha^{-1}(U_\alpha)\) 属于子基，故为开集，因此
\(\pi_\alpha\) 连续。\(\pi_\alpha\) 是开映射因为基元素被映为开集。(2) 若
\(f\) 连续，则每个 \(\pi_\alpha \circ f\) 连续。反之，若每个
\(\pi_\alpha \circ f\) 连续，则对子基元素
\(\pi_\alpha^{-1}(U_\alpha)\)，\(f^{-1}(\pi_\alpha^{-1}(U_\alpha)) = (\pi_\alpha \circ f)^{-1}(U_\alpha)\)
是开集，由定理 57.5(5) 知 \(f\) 连续。∎

\textbf{例 57.8}：\(\mathbb{R}^n\) 的标准拓扑是 \(n\) 个
\(\mathbb{R}\)（赋予标准度量拓扑）的积拓扑。

\subsubsection{57.7 商空间}\label{ux5546ux7a7aux95f4}

\textbf{定义 57.15}（商映射）：设 \(X, Y\) 为拓扑空间。满射
\(p: X \to Y\) 称为\textbf{商映射}，如果 \(V \subseteq Y\)
是开集当且仅当 \(p^{-1}(V)\) 是 \(X\) 中的开集。

\textbf{定义 57.16}（商拓扑）：设 \(X\) 为拓扑空间，\(\sim\) 是 \(X\)
上的等价关系，\(X/{\sim}\) 为等价类构成的集合，\(\pi: X \to X/{\sim}\)
为自然投影。\(X/{\sim}\) 上的\textbf{商拓扑}定义为使 \(\pi\)
成为商映射的拓扑，即 \(U \subseteq X/{\sim}\) 是开集当且仅当
\(\pi^{-1}(U)\) 是 \(X\) 中的开集。\((X/{\sim}, \text{商拓扑})\)
称为\textbf{商空间}。

\textbf{命题 57.11}（商拓扑的泛性质）：设 \(p: X \to Y\)
为商映射，\(g: Y \to Z\) 为映射。则 \(g\) 连续当且仅当
\(g \circ p: X \to Z\) 连续。

\emph{证明}：若 \(g\) 连续，则 \(g \circ p\) 连续。反之，若
\(g \circ p\) 连续，对 \(Z\) 中开集
\(W\)，\((g \circ p)^{-1}(W) = p^{-1}(g^{-1}(W))\) 是 \(X\)
中开集。由商拓扑定义，\(g^{-1}(W)\) 是 \(Y\) 中开集，故 \(g\) 连续。∎

\textbf{例 57.9}（常见商空间）： - \textbf{圆周}
\(S^1 \cong [0, 1] / \{0 \sim 1\}\)。 - \textbf{环面}
\(T^2 \cong S^1 \times S^1 \cong [0, 1]^2 / \sim\)（对边粘合）。 -
\textbf{实射影平面} \(\mathbb{R}P^2 \cong S^2 / \{x \sim -x\}\)。 -
\textbf{Möbius 带}：\([0, 1]^2\) 将一条边反向粘合。

\subsubsection{57.8 网与收敛}\label{ux7f51ux4e0eux6536ux655b}

度量空间中的序列收敛概念在一般拓扑空间中不够用，需要推广为网（net）的概念。

\textbf{定义 57.17}（有向集）：偏序集 \((D, \leq)\)
称为\textbf{有向集}，如果对任意 \(\alpha, \beta \in D\)，存在
\(\gamma \in D\) 使得 \(\alpha \leq \gamma\) 且 \(\beta \leq \gamma\)。

\textbf{定义 57.18}（网）：拓扑空间 \(X\) 中的一个\textbf{网}是一个映射
\(x: D \to X\)，其中 \(D\) 是有向集。通常记作
\((x_\alpha)_{\alpha \in D}\)。

\textbf{定义 57.19}（网的收敛）：称网 \((x_\alpha)_{\alpha \in D}\)
\textbf{收敛}到 \(x \in X\)（记作 \(x_\alpha \to x\)），如果对 \(x\)
的任意邻域 \(U\)，存在 \(\alpha_0 \in D\) 使得对所有
\(\alpha \geq \alpha_0\)，\(x_\alpha \in U\)。

\textbf{定理 57.12}（网刻画闭包）：设 \(X\)
为拓扑空间，\(A \subseteq X\)。则 \(x \in \overline{A}\) 当且仅当存在
\(A\) 中的网收敛到 \(x\)。

\emph{证明}：(\(\Rightarrow\)) 设 \(x \in \overline{A}\)。取有向集
\(D = \mathcal{N}(x)\)（\(x\)
的邻域系），按反向包含排序：\(U \leq V \Leftrightarrow V \subseteq U\)。对每个
\(U \in D\)，由 \(x \in \overline{A}\) 知
\(U \cap A \neq \emptyset\)，选择 \(x_U \in U \cap A\)。则
\((x_U)_{U \in D}\) 是 \(A\) 中的网且 \(x_U \to x\)。

(\(\Leftarrow\)) 设 \((x_\alpha)\) 是 \(A\)
中的网，\(x_\alpha \to x\)。对 \(x\) 的任意邻域 \(U\)，存在 \(\alpha_0\)
使对所有 \(\alpha \geq \alpha_0\) 有 \(x_\alpha \in U\)，特别地
\(x_\alpha \in U \cap A\)，故 \(U \cap A \neq \emptyset\)。因此
\(x \in \overline{A}\)。∎

\textbf{定理 57.13}（网刻画连续性）：设 \(f: X \to Y\) 为映射。则 \(f\)
在 \(x \in X\) 处连续当且仅当对 \(X\) 中任意收敛到 \(x\) 的网
\((x_\alpha)\)，\(f(x_\alpha) \to f(x)\) 在 \(Y\) 中成立。

\emph{证明}：(\(\Rightarrow\)) 设 \(f\) 在 \(x\) 处连续，\((x_\alpha)\)
收敛到 \(x\)。对 \(f(x)\) 的任意邻域 \(V\)，\(f^{-1}(V)\) 是 \(x\)
的邻域，故存在 \(\alpha_0\) 使对所有 \(\alpha \geq \alpha_0\) 有
\(x_\alpha \in f^{-1}(V)\)，即 \(f(x_\alpha) \in V\)。故
\(f(x_\alpha) \to f(x)\)。

(\(\Leftarrow\)) 反设 \(f\) 在 \(x\) 处不连续。则存在 \(f(x)\) 的邻域
\(V\) 使 \(f^{-1}(V)\) 不是 \(x\) 的邻域。取有向集
\(D = \mathcal{N}(x)\)，对每个
\(U \in D\)，\(U \not\subseteq f^{-1}(V)\)，故存在
\(x_U \in U \setminus f^{-1}(V)\)。则 \((x_U)\) 是收敛到 \(x\) 的网，但
\(f(x_U) \notin V\)，故 \(f(x_U)\) 不收敛到 \(f(x)\)。矛盾。∎

\bigskip\hrule\bigskip

\subsection{第58章：连通性与紧致性}\label{ux7b2c58ux7ae0ux8fdeux901aux6027ux4e0eux7d27ux81f4ux6027}

连通性和紧致性是拓扑空间最重要的两个全局性质。连通性刻画空间的「整体性」，紧致性则推广了有限覆盖的概念。

\subsubsection{58.1 连通性}\label{ux8fdeux901aux6027}

\textbf{定义 58.1}（连通空间）：拓扑空间 \(X\) 称为\textbf{连通的}，如果
\(X\) 不能表示为两个不相交的非空开集的并。等价地，\(X\)
的既开又闭的子集只有 \(\emptyset\) 和 \(X\)。

子集 \(A \subseteq X\) 称为\textbf{连通的}，如果 \(A\)
作为子空间是连通的。

\textbf{例 58.1}：\(\mathbb{R}\) 中的区间
\((a, b)\)、\([a, b]\)、\((a, b]\)
等都是连通的。离散空间（多于一点）不是连通的。

\textbf{命题 58.1}（连续映射保持连通性）：设 \(f: X \to Y\)
为连续映射，\(X\) 连通。则 \(f(X)\) 连通。

\emph{证明}：若 \(f(X)\) 不连通，则存在 \(f(X)\) 的既开又闭的非空真子集
\(A\)。则 \(f^{-1}(A)\) 是 \(X\) 的既开又闭的非空真子集，与 \(X\)
连通矛盾。∎

\textbf{定理 58.2}（连通集的并）：设 \(\{A_\alpha\}_{\alpha \in A}\) 是
\(X\) 的一族连通子集，且
\(\bigcap_{\alpha \in A} A_\alpha \neq \emptyset\)。则
\(\bigcup_{\alpha \in A} A_\alpha\) 连通。

\emph{证明}：设 \(Y = \bigcup A_\alpha\)，\(U\) 是 \(Y\)
的既开又闭子集。取 \(x \in \bigcap A_\alpha\)，不妨设
\(x \in U\)。对每个 \(\alpha\)，\(U \cap A_\alpha\) 是 \(A_\alpha\)
的既开又闭子集且非空，由 \(A_\alpha\) 连通知
\(U \cap A_\alpha = A_\alpha\)，即 \(A_\alpha \subseteq U\)。故
\(U = Y\)。∎

\textbf{定理 58.3}（闭包的连通性）：若 \(A \subseteq X\) 连通，且
\(A \subseteq B \subseteq \overline{A}\)，则 \(B\)
连通。特别地，连通集的闭包连通。

\emph{证明}：设 \(U\) 是 \(B\) 的既开又闭子集。则 \(U \cap A\) 是 \(A\)
的既开又闭子集，由 \(A\) 连通知 \(U \cap A = \emptyset\) 或
\(U \cap A = A\)。若 \(U \cap A = \emptyset\)，则
\(A \subseteq B \setminus U\)，\(B \setminus U\) 是闭集，故
\(\overline{A} \subseteq B \setminus U\)，从而 \(U = \emptyset\)。若
\(U \cap A = A\)，则 \(A \subseteq U\)，\(U\) 是闭集，故
\(\overline{A} \subseteq U\)（在 \(B\) 中），从而 \(U = B\)。∎

\textbf{定义 58.2}（连通分支）：设 \(X\) 为拓扑空间。定义关系
\(x \sim y\) 当且仅当存在 \(X\) 的连通子集同时包含 \(x\) 和
\(y\)。这是一个等价关系，其等价类称为 \(X\)
的\textbf{连通分支}（或\textbf{连通分量}）。

\textbf{命题 58.4}：每个连通分支是极大的连通子集，且是闭集。

\emph{证明}：由定理 58.2
知每个等价类中的点确实属于某个连通子集，且该等价类本身是连通集（取包含其中任意两点的所有连通子集的并）。极大性由定义直接可得。由定理
58.3，连通分支的闭包也连通，故连通分支等于其闭包，即为闭集。∎

\textbf{定义 58.3}（道路连通）：拓扑空间 \(X\)
称为\textbf{道路连通的}，如果对任意 \(x, y \in X\)，存在连续映射
\(\gamma: [0, 1] \to X\) 使得
\(\gamma(0) = x\)，\(\gamma(1) = y\)。\(\gamma\) 称为连接 \(x\) 和 \(y\)
的\textbf{道路}。

\textbf{命题 58.5}：道路连通空间必连通。

\emph{证明}：设 \(X\) 道路连通，假设 \(X = U \cup V\)，\(U, V\)
是不相交的非空开集。取 \(x \in U, y \in V\)，存在道路
\(\gamma: [0, 1] \to X\) 连接 \(x\) 和 \(y\)。则
\([0, 1] = \gamma^{-1}(U) \cup \gamma^{-1}(V)\) 将 \([0, 1]\)
分为两个不相交的非空开集，与 \([0, 1]\) 连通矛盾。∎

\textbf{例 58.2}（连通但不道路连通）：\textbf{拓扑正弦曲线}：
\[X = \{(x, \sin(1/x)) : x > 0\} \cup \{(0, y) : -1 \leq y \leq 1\}\]
是连通的（前者是连续像故连通，后者是其闭包的一部分），但不是道路连通的。

\textbf{定义 58.4}（局部连通、局部道路连通）：\(X\)
称为\textbf{局部连通的}（resp. \textbf{局部道路连通的}），如果对每个
\(x \in X\) 和 \(x\) 的任意邻域 \(U\)，存在 \(x\) 的连通（resp.
道路连通）邻域 \(V \subseteq U\)。

\textbf{命题
58.6}：在局部道路连通空间中，连通分支与道路连通分支一致，且都是既开又闭的。

\emph{证明}：在局部道路连通空间中，道路连通分支是开集。每个连通分支是道路连通分支的不交并。由于连通分支不能分解为两个不交开集的并，故每个连通分支恰为一个道路连通分支。∎

\subsubsection{58.2 紧致性}\label{ux7d27ux81f4ux6027}

\textbf{定义 58.5}（开覆盖）：设 \(X\) 为拓扑空间。一族开集
\(\{U_\alpha\}_{\alpha \in A}\) 称为 \(X\) 的\textbf{开覆盖}，如果
\(\bigcup_{\alpha \in A} U_\alpha = X\)。若子族也覆盖
\(X\)，则称为\textbf{子覆盖}。

\textbf{定义 58.6}（紧致空间）：拓扑空间 \(X\)
称为\textbf{紧致的}（compact），如果 \(X\) 的任意开覆盖都有有限子覆盖。

子集 \(K \subseteq X\) 称为\textbf{紧致的}，如果 \(K\)
作为子空间是紧致的。

\textbf{命题 58.7}：\(K \subseteq X\) 紧致当且仅当对 \(X\)
中任意一族覆盖 \(K\) 的开集 \(\{U_\alpha\}\)（即
\(K \subseteq \bigcup U_\alpha\)），存在有限子族也覆盖 \(K\)。

\emph{证明}：由子空间拓扑的定义，\(K\) 中开集形如 \(U \cap K\)，\(U\) 是
\(X\) 中开集。开覆盖 \(\{U_\alpha \cap K\}\) 对应 \(X\) 中一族覆盖 \(K\)
的开集。∎

\textbf{定理 58.8}（紧致集的基本性质）： 1. 紧致集的连续像是紧致集。 2.
紧致空间的闭子集紧致。 3. Hausdorff 空间中的紧致子集是闭集。 4.
紧致空间的连续像若在 Hausdorff 空间中，则是闭映射。

\emph{证明}：(1) 设 \(f: X \to Y\) 连续，\(X\) 紧致。对 \(f(X)\)
的任意开覆盖 \(\{V_\alpha\}\)，\(\{f^{-1}(V_\alpha)\}\) 是 \(X\)
的开覆盖，有有限子覆盖
\(\{f^{-1}(V_{\alpha_1}), \ldots, f^{-1}(V_{\alpha_n})\}\)，则
\(\{V_{\alpha_1}, \ldots, V_{\alpha_n}\}\) 覆盖 \(f(X)\)。

\begin{enumerate}
\def\labelenumi{(\arabic{enumi})}
\setcounter{enumi}{1}
\item
  设 \(X\) 紧致，\(F \subseteq X\) 闭。对 \(F\) 的任意开覆盖
  \(\{U_\alpha \cap F\}\)（\(U_\alpha\) 是 \(X\)
  中开集），\(\{U_\alpha\} \cup \{X \setminus F\}\) 是 \(X\)
  的开覆盖，有有限子覆盖，去掉 \(X \setminus F\) 后得到 \(F\)
  的有限子覆盖。
\item
  设 \(X\) 是 Hausdorff 空间，\(K \subseteq X\) 紧致。对
  \(x \notin K\)，对每个 \(y \in K\)，存在不相交的开集
  \(U_y \ni x\)，\(V_y \ni y\)。\(\{V_y\}_{y \in K}\) 覆盖
  \(K\)，有有限子覆盖 \(\{V_{y_1}, \ldots, V_{y_n}\}\)。令
  \(U = \bigcap_{i=1}^n U_{y_i}\)，则 \(U\) 是 \(x\) 的邻域且
  \(U \cap K = \emptyset\)。故 \(x \notin \overline{K}\)，即 \(K\) 闭。
\item
  由 (1) 和 (3)，紧致集在连续映射下的像是紧致集，在 Hausdorff
  空间中紧致集闭，故闭集（紧致空间的闭子集紧致）被映为紧致集（闭集）。∎
\end{enumerate}

\textbf{定理 58.9}（紧致空间到 Hausdorff 空间的连续双射）：若
\(f: X \to Y\) 是连续双射，\(X\) 紧致，\(Y\) 是 Hausdorff 空间，则 \(f\)
是同胚。

\emph{证明}：只需证 \(f\) 是闭映射。对 \(X\) 中闭集 \(F\)，\(F\)
紧致（定理 58.8(2)），\(f(F)\) 紧致（定理 58.8(1)），在 Hausdorff 空间
\(Y\) 中 \(f(F)\) 是闭集（定理 58.8(3)）。∎

\textbf{定理 58.10}（Heine-Borel 定理，\(\mathbb{R}^n\)
情形）：\(\mathbb{R}^n\) 的子集紧致当且仅当它是有界闭集。

\emph{证明}：这是 V4 中已证明的结论（定理 22.18 的推广）。简要思路：先证
\([0, 1]\) 紧致（区间套定理），再证有限个紧致空间的积紧致，最后
\(\mathbb{R}^n\) 中有界闭集可嵌入某个 \([-M, M]^n\) 中。∎

\textbf{定理 58.10a}（Lebesgue 数引理）：设 \((X, d)\)
为紧致度量空间，\(\{U_\alpha\}_{\alpha \in A}\) 是 \(X\)
的任意开覆盖。则存在 \(\delta > 0\)（称为该覆盖的 \textbf{Lebesgue
数}），使得对任意 \(A \subseteq X\) 满足
\(\operatorname{diam}(A) < \delta\)，存在某个 \(\alpha \in A\) 使得
\(A \subseteq U_\alpha\)。

\emph{证明}：对 \(\{U_\alpha\}\) 取有限子覆盖
\(\{U_1, \ldots, U_n\}\)（由紧致性）。令
\(f_i(x) = d(x, X \setminus U_i)\)（若 \(U_i = X\) 则
\(f_i(x) = 1\)），\(f(x) = \max_i f_i(x)\)。由于
\(X = \bigcup U_i\)，\(f(x) > 0\) 对所有 \(x\) 成立。由 \(f\) 连续且
\(X\) 紧致，\(f\) 在 \(X\) 上达到正的下确界。取 \(\delta\)
为此下确界即可。∎

\subsubsection{58.3 Tychonoff 定理}\label{tychonoff-ux5b9aux7406}

\textbf{定理 58.11a}（Alexander 子基定理）：设 \(X\)
为拓扑空间，\(\mathcal{S}\) 是 \(X\) 的拓扑的一个子基。若由
\(\mathcal{S}\) 中元素构成的任意开覆盖都有有限子覆盖，则 \(X\)
是紧致的。

\textbf{定理 58.11}（Tychonoff
定理）：任意一族紧致空间的积空间（赋予积拓扑）是紧致的。

\emph{证明思路}：设 \(X = \prod_{\alpha \in A} X_\alpha\)，每个
\(X_\alpha\) 紧致。使用 Alexander 子基定理：若以子基
\(\mathcal{S} = \{\pi_\alpha^{-1}(U_\alpha) : \alpha \in A, U_\alpha \subseteq X_\alpha \text{ 开}\}\)
的元素构成的任意覆盖都有有限子覆盖，则 \(X\) 紧致。

设 \(\mathcal{C} \subseteq \mathcal{S}\) 是 \(X\) 的覆盖。对每个
\(\alpha\)，若 \(\mathcal{C}\) 中来自 \(X_\alpha\) 的开集已覆盖
\(X\)，则存在有限子覆盖（因为 \(\pi_\alpha^{-1}(U)\) 覆盖 \(X\) 当且仅当
\(U = X_\alpha\)，此时单元素就覆盖 \(X\)）。否则，存在
\(x_\alpha \in X_\alpha\) 不被任何 \(\pi_\alpha^{-1}(U)\)（\(U\) 来自
\(\mathcal{C}\) 中与 \(X_\alpha\) 相关的开集）覆盖。由选择公理选取
\((x_\alpha)_{\alpha \in A}\)。则 \((x_\alpha)\) 不属于 \(\mathcal{C}\)
中任何元素，矛盾。∎

\textbf{推论 58.12}（Tychonoff
定理，有限积情形）：有限个紧致空间的积紧致。

有限积情形不需要选择公理，可以归纳证明。这里我们主要使用有限积情形。

\subsubsection{58.4
局部紧致与紧化}\label{ux5c40ux90e8ux7d27ux81f4ux4e0eux7d27ux5316}

\textbf{定义 58.7}（局部紧致）：拓扑空间 \(X\)
称为\textbf{局部紧致的}，如果对每个 \(x \in X\)，存在 \(x\) 的邻域 \(U\)
使得 \(\overline{U}\) 是紧致的。

\textbf{定理 58.13}（单点紧化 / Alexandroff 紧化）：设 \(X\)
是局部紧致的 Hausdorff 空间。则存在紧致 Hausdorff 空间 \(X^*\) 使得
\(X\) 同胚于 \(X^*\) 的稠密开子集，且 \(X^* \setminus X\)
是单点集。\(X^*\) 在同胚意义下唯一。

\emph{证明}（构造）：令 \(X^* = X \cup \{\infty\}\)，其中
\(\infty \notin X\)。定义 \(X^*\) 上的拓扑：\(U \subseteq X^*\)
是开集当且仅当 - \(U \subseteq X\) 且 \(U\) 是 \(X\) 中开集，或 -
\(\infty \in U\) 且 \(X^* \setminus U\) 是 \(X\) 中的紧致闭集。

验证这是拓扑且 \(X^*\) 紧致 Hausdorff。\(X\) 作为子空间同胚于原来的
\(X\)，且 \(X\) 在 \(X^*\) 中稠密。∎

\textbf{例 58.3}：\(\mathbb{R}^n\) 的单点紧化同胚于 \(S^n\)（\(n\)
维球面）。

\subsubsection{58.5 Urysohn 引理与 Tietze
扩张定理}\label{urysohn-ux5f15ux7406ux4e0e-tietze-ux6269ux5f20ux5b9aux7406}

\textbf{定理 58.14}（Urysohn 引理）：设 \(X\) 为正规空间（见第 59
章），\(A, B\) 是 \(X\) 中不相交的闭集。则存在连续函数
\(f: X \to [0, 1]\) 使得 \(f(A) = \{0\}\)，\(f(B) = \{1\}\)。

\emph{证明}：令 \(U_1 = X \setminus B\)，则
\(A \subseteq U_1\)。由正规性，递归构造开集
\(U_r\)（\(r \in \mathbb{Q} \cap [0, 1]\)）满足对 \(p < q\) 有
\(\overline{U}_p \subseteq U_q\)。具体地，将 \(\mathbb{Q} \cap [0, 1]\)
排为 \(r_1 = 0, r_2 = 1, r_3, r_4, \ldots\)。归纳假设已对
\(r_1, \ldots, r_n\) 构造了 \(U_{r_i}\)。设 \(r_k, r_\ell\) 是
\(r_{n+1}\) 在已构造的 \(r_i\) 中的直接前驱和后继，则
\(U_{r_k} \subseteq \overline{U}_{r_k} \subseteq U_{r_\ell}\)，由正规性可插入
\(U_{r_{n+1}}\)。

定义 \(f(x) = \inf\{r : x \in U_r\}\)（约定 \(U_r = \emptyset\) 当
\(r < 0\)，\(U_r = X\) 当 \(r > 1\)）。则 \(f\)
连续：\(\{x : f(x) < t\} = \bigcup_{r < t} U_r\)
是开集，\(\{x : f(x) > t\} = \bigcup_{r > t} (X \setminus \overline{U}_r)\)
是开集。显然 \(f(A) = \{0\}\)，\(f(B) = \{1\}\)。∎

\textbf{定理 58.15}（Tietze 扩张定理）：设 \(X\)
为正规空间，\(A \subseteq X\) 为闭集。则 \(A\) 上的任意连续函数
\(f: A \to [a, b]\)（或 \(f: A \to \mathbb{R}\)）可连续扩张到 \(X\) 上。

\emph{证明}（概要）：不妨设 \(f: A \to [-1, 1]\)。令
\(A_1 = f^{-1}([-1, -1/3])\)，\(B_1 = f^{-1}([1/3, 1])\)。由 Urysohn
引理存在 \(g_1: X \to [-1/3, 1/3]\) 使
\(g_1(A_1) = \{-1/3\}\)，\(g_1(B_1) = \{1/3\}\)。则
\(|f(x) - g_1(x)| \leq 2/3\) 对 \(x \in A\)。令 \(f_1 = f - g_1|_A\)，则
\(f_1: A \to [-2/3, 2/3]\)。归纳得 \(g_n\) 使
\(|f - \sum_{i=1}^n g_i| \leq (2/3)^n\) 在 \(A\) 上。级数 \(\sum g_n\)
一致收敛，极限 \(g\) 是 \(f\) 的连续扩张。∎

\bigskip\hrule\bigskip

\subsection{第59章：分离公理与度量化}\label{ux7b2c59ux7ae0ux5206ux79bbux516cux7406ux4e0eux5ea6ux91cfux5316}

分离公理刻画拓扑空间中点与点、点与闭集之间的「分离」程度，从 T0 到 T4
形成越来越强的条件。度量化理论研究拓扑空间何时可由度量诱导。

\subsubsection{59.1 分离公理}\label{ux5206ux79bbux516cux7406}

\textbf{定义 59.1}（分离公理 \(T_0\)--\(T_4\)）：设 \(X\) 为拓扑空间。

\begin{itemize}
\tightlist
\item
  \(T_0\)（Kolmogorov）：对任意两个不同的点
  \(x, y \in X\)，存在开集包含其中一点但不包含另一点。
\item
  \(T_1\)（Fréchet）：对任意两个不同的点 \(x, y \in X\)，存在开集 \(U\)
  使得 \(x \in U\) 但 \(y \notin U\)。
\item
  \(T_2\)（Hausdorff）：对任意两个不同的点
  \(x, y \in X\)，存在不相交的开集 \(U, V\) 使得 \(x \in U, y \in V\)。
\item
  \(T_3\)（正则）：\(X\) 是 \(T_1\) 的，且对任意 \(x \in X\) 和不包含
  \(x\) 的闭集 \(F\)，存在不相交的开集 \(U, V\) 使得
  \(x \in U, F \subseteq V\)。
\item
  \(T_{3.5}\)（完全正则 / Tychonoff）：\(X\) 是 \(T_1\) 的，且对任意
  \(x \in X\) 和不包含 \(x\) 的闭集 \(F\)，存在连续函数
  \(f: X \to [0, 1]\) 使得 \(f(x) = 0\)，\(f(F) = \{1\}\)。
\item
  \(T_4\)（正规）：\(X\) 是 \(T_1\) 的，且对任意不相交的闭集
  \(A, B\)，存在不相交的开集 \(U, V\) 使得
  \(A \subseteq U, B \subseteq V\)。
\end{itemize}

\textbf{命题 59.1}（分离公理的关系）：
\[T_4 \Rightarrow T_{3.5} \Rightarrow T_3 \Rightarrow T_2 \Rightarrow T_1 \Rightarrow T_0\]

\emph{证明}：\(T_4 \Rightarrow T_{3.5}\) 是 Urysohn 引理（定理
58.14）。\(T_{3.5} \Rightarrow T_3\)：取
\(U = f^{-1}([0, 1/2))\)，\(V = f^{-1}((1/2, 1])\)。其余蕴含关系直接由定义可得。∎

\textbf{命题 59.2}（\(T_1\) 的等价刻画）：\(X\) 是 \(T_1\)
空间当且仅当每个单点集是闭集。

\emph{证明}：(\(\Rightarrow\)) 对 \(x \in X\)，对每个
\(y \neq x\)，存在开集 \(U_y \ni y\) 且 \(x \notin U_y\)。则
\(X \setminus \{x\} = \bigcup_{y \neq x} U_y\) 是开集，故 \(\{x\}\)
是闭集。(\(\Leftarrow\)) 对 \(x \neq y\)，取
\(U = X \setminus \{y\}\)，则 \(x \in U\) 且 \(y \notin U\)。∎

\textbf{命题 59.3}（\(T_2\) 的等价刻画）：\(X\) 是 Hausdorff
空间当且仅当对角线 \(\Delta = \{(x, x) : x \in X\}\) 在
\(X \times X\)（积拓扑）中是闭集。

\emph{证明}：(\(\Rightarrow\)) 对 \((x, y) \notin \Delta\)（即
\(x \neq y\)），存在不相交的开集 \(U \ni x, V \ni y\)，则 \(U \times V\)
是 \((x, y)\) 的邻域且与 \(\Delta\) 不相交。故 \(\Delta\)
闭。(\(\Leftarrow\)) 类似。∎

\textbf{命题 59.4}：Hausdorff 空间中序列（或网）的极限若存在则唯一。

\emph{证明}：设 \(x_\alpha \to x\) 且
\(x_\alpha \to y\)，\(x \neq y\)。取不相交的开集
\(U \ni x, V \ni y\)。由收敛定义，存在 \(\alpha_1\) 使
\(\alpha \geq \alpha_1\) 时 \(x_\alpha \in U\)，存在 \(\alpha_2\) 使
\(\alpha \geq \alpha_2\) 时 \(x_\alpha \in V\)。取
\(\alpha \geq \alpha_1, \alpha_2\)，则
\(x_\alpha \in U \cap V = \emptyset\)，矛盾。∎

\textbf{命题 59.5}：紧致 Hausdorff 空间是正规的（\(T_4\)）。

\emph{证明}：先证紧致 Hausdorff 空间是正则的：设 \(x \in X\)，\(F\)
是不含 \(x\) 的闭集。对每个 \(y \in F\)，由 \(T_2\) 存在不相交开集
\(U_y \ni x, V_y \ni y\)。\(\{V_y\}\) 覆盖紧致集 \(F\)，有有限子覆盖
\(\{V_{y_1}, \ldots, V_{y_n}\}\)。令
\(U = \bigcap U_{y_i}\)，\(V = \bigcup V_{y_i}\)，则
\(U \cap V = \emptyset\)。

再证正规性：设 \(A, B\) 是不相交的闭集。对每个
\(a \in A\)，由正则性存在不相交开集
\(U_a \ni a, V_a \supseteq B\)。\(\{U_a\}\) 覆盖紧致集
\(A\)，有有限子覆盖，类似可得分离 \(A\) 和 \(B\) 的开集。∎

\textbf{例 59.1}（反例）： - 有限补拓扑（无限集上）是 \(T_1\) 但不是
\(T_2\)。 - 下限拓扑 \(\mathbb{R}_\ell\)（以 \([a, b)\) 为基）是 \(T_4\)
但不是 \(T_5\)。 - 无限个 \(\mathbb{R}\)
的积在积拓扑下是非正规的（但仍是 \(T_{3.5}\)）。

\subsubsection{59.2 度量化定理}\label{ux5ea6ux91cfux5316ux5b9aux7406}

\textbf{定义 59.2}（可度量化）：拓扑空间 \((X, \mathcal{T})\)
称为\textbf{可度量化的}，如果存在 \(X\) 上的度量 \(d\) 使得
\(\mathcal{T}\) 恰为 \(d\) 诱导的拓扑。

\textbf{命题 59.6}：可度量化空间是 \(T_4\)（从而也是
\(T_0\)--\(T_{3.5}\)）。

\emph{证明}：设 \((X, d)\) 为度量空间。对
\(x \neq y\)，\(d(x, y) > 0\)，开球 \(B(x, d(x, y)/2)\) 和
\(B(y, d(x, y)/2)\) 不相交，故 \(T_2\)。

正规性：设 \(A, B\) 是不相交的闭集。定义
\(f(x) = d(x, A) / (d(x, A) + d(x, B))\)（\(d(x, A) = \inf_{a \in A} d(x, a)\)）。则
\(f\) 连续，\(f(A) = \{0\}\)，\(f(B) = \{1\}\)。取
\(U = f^{-1}([0, 1/2))\)，\(V = f^{-1}((1/2, 1])\)。∎

\textbf{定理 59.7}（Urysohn 度量化定理）：每个具有可数基的正则空间
\(X\)（即第二可数的正则空间）可度量化。

\emph{证明思路}：设 \(\mathcal{B} = \{B_n\}_{n \in \mathbb{N}}\) 是
\(X\) 的可数基。对每对 \((n, m)\) 满足
\(\overline{B}_n \subseteq B_m\)，由 Urysohn 引理构造连续函数
\(f_{n,m}: X \to [0, 1]\) 使
\(f_{n,m}(\overline{B}_n) = \{0\}\)，\(f_{n,m}(X \setminus B_m) = \{1\}\)。这些函数构成可数族
\(\{f_k\}_{k \in \mathbb{N}}\)。定义嵌入
\(F: X \to [0, 1]^{\mathbb{N}}\)（赋予积拓扑，可度量化）为
\(F(x) = (f_k(x))_{k \in \mathbb{N}}\)。由正则性和第二可数性可证 \(F\)
是嵌入（同胚于其像）。\([0, 1]^{\mathbb{N}}\) 可度量化（如
\(d((a_n), (b_n)) = \sum 2^{-n}|a_n - b_n|\)），故 \(X\) 可度量化。∎

\textbf{推论 59.8}：紧致 Hausdorff 空间可度量化当且仅当它是第二可数的。

\emph{证明}：紧致 Hausdorff 空间是正规的，故正则。第二可数正则空间由
Urysohn
度量化定理可度量化。反之，紧致可度量化空间必第二可数（紧致度量空间有可数基）。∎

\textbf{定理 59.9}（Nagata-Smirnov
度量化定理，陈述）：拓扑空间可度量化当且仅当它是正则的且具有
\(\sigma\)-局部有限基。

\emph{说明}：这是度量化问题的最一般解答。其证明超出本卷范围，此处仅作陈述。\(\sigma\)-局部有限基指基
\(\mathcal{B} = \bigcup_{n=1}^\infty \mathcal{B}_n\)，其中每个
\(\mathcal{B}_n\) 是局部有限的（每点有邻域只与 \(\mathcal{B}_n\)
中有限个元素相交）。

\bigskip\hrule\bigskip

\subsection{第60章：基本群}\label{ux7b2c60ux7ae0ux57faux672cux7fa4}

基本群是拓扑空间最基本的代数不变量，它将拓扑问题转化为代数问题。

\subsubsection{60.1 同伦}\label{ux540cux4f26}

\textbf{定义 60.1}（同伦）：设 \(f, g: X \to Y\) 为连续映射。\(f\) 到
\(g\) 的一个\textbf{同伦}是一个连续映射 \(H: X \times [0, 1] \to Y\)
使得 \(H(x, 0) = f(x)\)，\(H(x, 1) = g(x)\) 对所有 \(x \in X\)
成立。若存在这样的同伦，则称 \(f\) \textbf{同伦于} \(g\)，记作
\(f \simeq g\)。

\textbf{命题 60.1}：同伦关系 \(\simeq\) 是 \(C(X, Y)\)（\(X\) 到 \(Y\)
的连续映射集）上的等价关系。

\emph{证明}：自反性：\(H(x, t) = f(x)\) 给出 \(f \simeq f\)。对称性：若
\(H\) 是 \(f \simeq g\) 的同伦，则 \(H'(x, t) = H(x, 1-t)\) 给出
\(g \simeq f\)。传递性：若
\(H_1: f \simeq g\)，\(H_2: g \simeq h\)，定义

\[H(x, t) = \begin{cases} H_1(x, 2t) & 0 \leq t \leq 1/2 \\ H_2(x, 2t-1) & 1/2 \leq t \leq 1 \end{cases}\]

由粘合引理（Pasting Lemma），\(H\) 连续，故 \(f \simeq h\)。∎

\textbf{定义 60.2}（同伦类）：\(f\) 的\textbf{同伦类}记作 \([f]\)。

\textbf{定义 60.3}（道路同伦）：设 \(\gamma, \delta: [0, 1] \to X\)
是两条道路，满足
\(\gamma(0) = \delta(0) = x_0\)，\(\gamma(1) = \delta(1) = x_1\)。\(f\)
到 \(g\) 的一个\textbf{道路同伦}（或定端同伦）是一个同伦
\(H: [0, 1] \times [0, 1] \to X\) 满足
\(H(s, 0) = \gamma(s)\)，\(H(s, 1) = \delta(s)\)，且
\(H(0, t) = x_0\)，\(H(1, t) = x_1\) 对所有 \(t\) 成立。记作
\(\gamma \simeq_p \delta\)。

\textbf{定义 60.4}（同伦等价）：映射 \(f: X \to Y\)
称为\textbf{同伦等价}，如果存在 \(g: Y \to X\) 使得
\(g \circ f \simeq \operatorname{id}_X\) 且
\(f \circ g \simeq \operatorname{id}_Y\)。此时称 \(X\) 与 \(Y\)
\textbf{同伦等价}（或具有相同的\textbf{同伦型}），记作 \(X \simeq Y\)。

\textbf{定义 60.5}（可缩空间）：拓扑空间 \(X\) 称为\textbf{可缩的}，如果
\(X\) 同伦等价于单点空间。等价地，\(\operatorname{id}_X\)
同伦于常值映射。

\textbf{例 60.1}：\(\mathbb{R}^n\) 可缩，\(D^n\)（\(n\) 维闭球）可缩。但
\(S^n\) 不可缩（\(n \geq 0\)）。

\subsubsection{60.2
基本群的定义}\label{ux57faux672cux7fa4ux7684ux5b9aux4e49}

\textbf{定义 60.6}（道路的乘法）：设 \(\gamma: [0, 1] \to X\) 是从 \(x\)
到 \(y\) 的道路，\(\delta: [0, 1] \to X\) 是从 \(y\) 到 \(z\)
的道路。它们的\textbf{乘积}（或连接）\(\gamma \cdot \delta: [0, 1] \to X\)
定义为

\[(\gamma \cdot \delta)(s) = \begin{cases} \gamma(2s) & 0 \leq s \leq 1/2 \\ \delta(2s-1) & 1/2 \leq s \leq 1 \end{cases}\]

这是从 \(x\) 到 \(z\) 的道路。

\textbf{定义 60.7}（逆道路）：道路 \(\gamma\) 的\textbf{逆道路}
\(\bar{\gamma}: [0, 1] \to X\) 定义为
\(\bar{\gamma}(s) = \gamma(1-s)\)。

\textbf{定义 60.8}（基本群）：设 \(X\) 为拓扑空间，\(x_0 \in X\)
为\textbf{基点}。考虑以 \(x_0\)
为起点和终点的所有道路（即\textbf{回路}，或称\textbf{闭路}）的道路同伦类构成的集合：

\[\pi_1(X, x_0) = \{[\gamma] : \gamma: [0, 1] \to X \text{ 连续}, \gamma(0) = \gamma(1) = x_0\}\]

在 \(\pi_1(X, x_0)\)
上定义乘法：\([\gamma] \cdot [\delta] = [\gamma \cdot \delta]\)。定义单位元为常值道路
\(e_{x_0}\)（\(e_{x_0}(s) = x_0\)）的同伦类。定义逆元
\([\gamma]^{-1} = [\bar{\gamma}]\)。

\textbf{定理 60.2}：\(\pi_1(X, x_0)\) 在上述运算下构成群，称为 \(X\)
的以 \(x_0\) 为基点的\textbf{基本群}（或\textbf{第一同伦群}）。

\emph{证明}：需验证乘法是良定义的、结合律、单位元、逆元。

良定义：若 \(\gamma \simeq_p \gamma'\)（同伦
\(F\)），\(\delta \simeq_p \delta'\)（同伦 \(G\)），则
\(H(s, t) = F(2s, t)\) 对
\(0 \leq s \leq 1/2\)，\(H(s, t) = G(2s-1, t)\) 对 \(1/2 \leq s \leq 1\)
给出 \(\gamma \cdot \delta \simeq_p \gamma' \cdot \delta'\)。

结合律：需证
\((\alpha \cdot \beta) \cdot \gamma \simeq_p \alpha \cdot (\beta \cdot \gamma)\)。构造同伦
\(H(s, t)\) 通过对参数 \(s\)
的重新参数化：\((\alpha \cdot \beta) \cdot \gamma\) 在 \([0, 1/4]\),
\([1/4, 1/2]\), \([1/2, 1]\) 上分别走
\(\alpha, \beta, \gamma\)；\(\alpha \cdot (\beta \cdot \gamma)\) 在
\([0, 1/2]\), \([1/2, 3/4]\), \([3/4, 1]\) 上分别走
\(\alpha, \beta, \gamma\)。通过线性插值可构造道路同伦。

单位元：需证
\([\gamma] \cdot [e_{x_0}] = [\gamma] = [e_{x_0}] \cdot [\gamma]\)。\(\gamma \cdot e_{x_0}\)
在 \([0, 1/2]\) 上走 \(\gamma\)，在 \([1/2, 1]\) 上停在
\(x_0\)。通过重新参数化可构造同伦到 \(\gamma\)。

逆元：需证 \([\gamma] \cdot [\bar{\gamma}] = [e_{x_0}]\)。构造同伦
\(H(s, t)\)：当 \(t\) 从 \(0\) 到 \(1\) 时，\(\gamma\)
走得越来越短，然后原路返回，最终停在 \(x_0\)。∎

\textbf{命题 60.3}（基点变换）：设 \(x_0, x_1 \in X\)，\(\eta\) 是从
\(x_0\) 到 \(x_1\) 的道路。则映射

\[\Phi_\eta: \pi_1(X, x_0) \to \pi_1(X, x_1), \quad [\gamma] \mapsto [\bar{\eta} \cdot \gamma \cdot \eta]\]

是群同构。因此，若 \(X\) 道路连通，则 \(\pi_1(X, x_0)\)
的同构型不依赖于基点的选择。

\emph{证明}：\(\Phi_\eta\)
是良定义的且保乘法：\(\Phi_\eta([\alpha][\beta]) = [\bar{\eta} \cdot \alpha \cdot \beta \cdot \eta] = [\bar{\eta} \cdot \alpha \cdot \eta \cdot \bar{\eta} \cdot \beta \cdot \eta] = \Phi_\eta([\alpha]) \Phi_\eta([\beta])\)。其逆为
\(\Phi_{\bar{\eta}}\)。∎

\textbf{定义 60.9}（诱导同态）：设 \(f: X \to Y\)
为连续映射，\(f(x_0) = y_0\)。定义

\[f_*: \pi_1(X, x_0) \to \pi_1(Y, y_0), \quad f_*([\gamma]) = [f \circ \gamma]\]

\textbf{命题 60.4}（诱导同态的性质）： 1. \(f_*\) 是群同态。 2.
\((\operatorname{id}_X)_* = \operatorname{id}_{\pi_1(X, x_0)}\)。 3.
\((g \circ f)_* = g_* \circ f_*\)。 4. 若
\(f \simeq g\)（同伦），且同伦保持基点，则 \(f_* = g_*\)。

因此 \(\pi_1\) 是从带基点拓扑空间范畴到群范畴的函子。

\emph{证明}：(1)
\(f_*([\alpha][\beta]) = [f \circ (\alpha \cdot \beta)] = [(f \circ \alpha) \cdot (f \circ \beta)] = f_*([\alpha]) f_*([\beta])\)。(2)-(3)
显然。(4) 设 \(H: X \times [0, 1] \to Y\) 是 \(f\) 到 \(g\)
的保基点同伦。对回路 \(\gamma\)，\(H(\gamma(s), t)\) 给出
\(f \circ \gamma\) 到 \(g \circ \gamma\) 的道路同伦。∎

\textbf{定理 60.5}：若 \(X\) 与 \(Y\) 同伦等价（道路连通空间），则
\(\pi_1(X) \cong \pi_1(Y)\)。

\emph{证明}：设 \(f: X \to Y\)，\(g: Y \to X\) 为同伦逆。则
\(g_* \circ f_* = (g \circ f)_* = (\operatorname{id}_X)_* = \operatorname{id}\)，同理
\(f_* \circ g_* = \operatorname{id}\)，故 \(f_*\) 是同构。∎

\textbf{推论 60.6}：可缩空间的基本群是平凡群。

\subsubsection{60.3
基本群的计算}\label{ux57faux672cux7fa4ux7684ux8ba1ux7b97}

\textbf{定理 60.7}：\(\pi_1(S^1) \cong \mathbb{Z}\)。

\emph{证明思路}：利用覆叠空间理论（见第 61 章）。覆叠映射
\(p: \mathbb{R} \to S^1\)，\(p(t) = e^{2\pi i t}\)。\(\mathbb{R}\)
是单连通的（可缩，故基本群平凡）。提升性质将 \(S^1\) 中的回路提升为
\(\mathbb{R}\) 中的道路。回路 \(\gamma\) 的提升的终点
\(n\)（整数）定义了同构 \(\pi_1(S^1) \cong \mathbb{Z}\)。详细证明见 61.3
节。∎

\textbf{定理
60.8}：\(\pi_1(X \times Y, (x_0, y_0)) \cong \pi_1(X, x_0) \times \pi_1(Y, y_0)\)。

\emph{证明}：投影 \(\pi_X, \pi_Y\) 诱导同态
\(\pi_1(X \times Y) \to \pi_1(X) \times \pi_1(Y)\)。由积的泛性质，这是同构。∎

\textbf{推论 60.9}：\(\pi_1(T^n) \cong \mathbb{Z}^n\)（\(n\) 维环面）。

\textbf{定理 60.10}：对 \(n \geq 2\)，\(\pi_1(S^n)\) 是平凡群。

\emph{证明思路}：对 \(n \geq 2\)，\(S^n\)
中的任意回路可同伦地避开某个点（利用 Sard 定理或单纯逼近），从而可收缩到
\(S^n \setminus \{p\} \cong \mathbb{R}^n\)（可缩）。∎

\subsubsection{60.4 Seifert-van Kampen
定理}\label{seifert-van-kampen-ux5b9aux7406}

\textbf{定理 60.11}（Seifert-van Kampen 定理）：设
\(X = U \cup V\)，其中 \(U, V\) 是道路连通的开集，\(U \cap V\)
非空且道路连通。取基点
\(x_0 \in U \cap V\)。则下图是推出（pushout）在群范畴中：

\[
\begin{CD}
\pi_1(U \cap V, x_0) @>{i_*}>> \pi_1(U, x_0) \\
@V{j_*}VV @VV{\alpha_*}V \\
\pi_1(V, x_0) @>>{\beta_*}> \pi_1(X, x_0)
\end{CD}
\]

即 \(\pi_1(X, x_0) \cong \pi_1(U, x_0) * \pi_1(V, x_0) / N\)，其中 \(*\)
表示自由积，\(N\) 是由
\(\{i_*(g) j_*(g)^{-1} : g \in \pi_1(U \cap V, x_0)\}\) 生成的正规子群。

\emph{证明思路}：需要证明 \(\pi_1(X, x_0)\) 由 \(\pi_1(U)\) 和
\(\pi_1(V)\) 的像生成（由 Lebesgue 数引理分割回路），且关系由
\(U \cap V\) 中的回路给出。具体证明需用到回路的细分和同伦构造。∎

\textbf{推论 60.12}（基本群的计算示例）： 1.
\(\pi_1(S^1 \vee S^1) \cong \mathbb{Z} * \mathbb{Z}\)（两个圆周的楔积的基本群是秩为
2 的自由群）。 2.
\(\pi_1(\Sigma_g) \cong \langle a_1, b_1, \ldots, a_g, b_g : [a_1, b_1] \cdots [a_g, b_g] = 1 \rangle\)（亏格
\(g\) 的可定向闭曲面的基本群）。 3.
\(\pi_1(\mathbb{R}P^2) \cong \mathbb{Z}/2\mathbb{Z}\)。

\emph{证明}：(1) 取 \(U, V\) 为两个圆周的适当开邻域，\(U \cap V\)
可缩，故
\(\pi_1(S^1 \vee S^1) \cong \pi_1(S^1) * \pi_1(S^1) \cong \mathbb{Z} * \mathbb{Z}\)。

\begin{enumerate}
\def\labelenumi{(\arabic{enumi})}
\setcounter{enumi}{1}
\item
  将 \(\Sigma_g\) 表示为 \(4g\) 边形的边粘合。取 \(U\)
  为内部（可缩），\(V\) 为边界邻域（同伦等价于 \(2g\)
  个圆周的楔积），\(U \cap V \simeq S^1\)。由 Seifert-van Kampen
  定理可得。
\item
  \(\mathbb{R}P^2\) 可表示为 \(D^2\) 将边界对径点粘合。取 \(U\)
  为内部（可缩），\(V\) 为去心邻域（同伦等价于
  \(S^1\)），\(U \cap V \simeq S^1\)。由 Seifert-van
  Kampen：\(\pi_1(\mathbb{R}P^2) \cong \langle a : a^2 = 1 \rangle \cong \mathbb{Z}/2\mathbb{Z}\)。∎
\end{enumerate}

\subsubsection{60.5
基本群的应用}\label{ux57faux672cux7fa4ux7684ux5e94ux7528}

\textbf{定理 60.13}（Brouwer 不动点定理，\(n=2\)）：任意连续映射
\(f: D^2 \to D^2\) 有不动点。

\emph{证明}：假设 \(f\) 无不动点。定义 \(r: D^2 \to S^1\) 为从 \(f(x)\)
出发经过 \(x\) 的射线与 \(S^1\) 的交点。则 \(r\) 是连续映射且
\(r|_{S^1} = \operatorname{id}_{S^1}\)（收缩映射）。考虑同态：

\[\pi_1(S^1) \xrightarrow{i_*} \pi_1(D^2) \xrightarrow{r_*} \pi_1(S^1)\]

\(r_* \circ i_* = (r \circ i)_* = \operatorname{id}_{\pi_1(S^1)}\)。但
\(\pi_1(D^2)\) 是平凡群（\(D^2\) 可缩），而
\(\pi_1(S^1) \cong \mathbb{Z}\)，\(\operatorname{id}_{\mathbb{Z}}\)
不能通过平凡群分解，矛盾。∎

\textbf{定理 60.14}（代数基本定理，拓扑证明）：任意非常数复系数多项式
\(P(z) \in \mathbb{C}[z]\) 在 \(\mathbb{C}\) 中有根。

\emph{证明}：设
\(P(z) = z^n + a_{n-1}z^{n-1} + \cdots + a_0\)（不妨设首一）。若
\(P(z)\) 无根，则 \(P(z) \neq 0\) 对所有 \(z \in \mathbb{C}\)。定义
\(f_r: S^1 \to S^1\) 为 \(f_r(z) = P(rz)/|P(rz)|\)（\(r > 0\)）。当
\(r=0\) 时 \(f_0\) 是常值映射 \(z \mapsto P(0)/|P(0)|\)。当 \(r\)
很大时，\(P(rz) = r^n z^n (1 + O(1/r))\)，故
\(f_r(z) \simeq z^n\)。但这些映射对不同的 \(r\)
是同伦的，故常值映射同伦于
\(z \mapsto z^n\)，在基本群上诱导同态：\(0 = n \cdot \operatorname{id}_{\mathbb{Z}}\)，从而
\(n = 0\)，矛盾。∎

\textbf{定理 60.15}（Borsuk-Ulam 定理，\(n=2\)）：对任意连续映射
\(f: S^2 \to \mathbb{R}^2\)，存在 \(x \in S^2\) 使得 \(f(x) = f(-x)\)。

\emph{证明思路}：若不存在这样的 \(x\)，可构造映射 \(g: S^2 \to S^1\) 为
\(g(x) = (f(x) - f(-x))/|f(x) - f(-x)|\)。\(g\) 满足
\(g(-x) = -g(x)\)。考虑商映射 \(\mathbb{R}P^2 \to S^1\)，这与
\(\pi_1(\mathbb{R}P^2) \cong \mathbb{Z}/2\mathbb{Z}\) 和
\(\pi_1(S^1) \cong \mathbb{Z}\) 的矛盾（非平凡同态必须把 \(1\) 阶元映到
\(2\) 阶元，但 \(\mathbb{Z}\) 中无 \(2\) 阶元）。∎

\bigskip\hrule\bigskip

\subsection{第61章：覆叠空间}\label{ux7b2c61ux7ae0ux8986ux53e0ux7a7aux95f4}

覆叠空间是研究基本群和拓扑空间结构的重要工具。

\subsubsection{61.1
覆叠空间的定义}\label{ux8986ux53e0ux7a7aux95f4ux7684ux5b9aux4e49}

\textbf{定义 61.1}（覆叠空间）：设 \(p: \tilde{X} \to X\) 是连续满射。称
\(p\) 为\textbf{覆叠映射}（\(\tilde{X}\) 为 \(X\)
的\textbf{覆叠空间}），如果对每个 \(x \in X\)，存在开邻域
\(U\)（称为\textbf{基本邻域}或\textbf{均匀覆盖邻域}）使得
\(p^{-1}(U) = \bigsqcup_{\alpha} V_\alpha\)（不交并），其中每个
\(V_\alpha\) 是 \(\tilde{X}\) 的开集，且
\(p|_{V_\alpha}: V_\alpha \to U\) 是同胚。

\textbf{例 61.1}（覆叠映射举例）： 1.
\(p: \mathbb{R} \to S^1\)，\(p(t) = e^{2\pi i t}\)。基本邻域：\(S^1\)
上任意真开弧。 2. \(p: S^1 \to S^1\)，\(p(z) = z^n\)（\(n\) 重覆叠）。
3. \(p: S^n \to \mathbb{R}P^n\)，\(p(x) = \{\pm x\}\)（二重覆叠）。 4.
\(p: \mathbb{R}^n \to T^n\)，\(p(x_1, \ldots, x_n) = (e^{2\pi i x_1}, \ldots, e^{2\pi i x_n})\)。

\textbf{定义 61.2}（纤维）：对 \(x \in X\)，\(p^{-1}(x)\) 称为 \(x\)
上的\textbf{纤维}。若 \(X\)
道路连通，则所有纤维具有相同的基数，称为覆叠的\textbf{层数}（或\textbf{重数}）。

\textbf{命题 61.1}：覆叠映射是开映射，且在 \(X\) 为 Hausdorff 空间时
\(\tilde{X}\) 也是 Hausdorff 空间。

\emph{证明}：由局部同胚性质和定义直接可得。∎

\subsubsection{61.2 提升性质}\label{ux63d0ux5347ux6027ux8d28}

\textbf{定义 61.3}（提升）：设 \(p: \tilde{X} \to X\)
为覆叠映射，\(f: Y \to X\) 为连续映射。\(f\)
的一个\textbf{提升}是一个连续映射 \(\tilde{f}: Y \to \tilde{X}\) 使得
\(p \circ \tilde{f} = f\)。

\textbf{定理 61.2}（道路提升引理）：设 \(p: \tilde{X} \to X\)
为覆叠映射，\(\gamma: [0, 1] \to X\)
是一条道路，\(\gamma(0) = x_0\)。对任意
\(\tilde{x}_0 \in p^{-1}(x_0)\)，存在唯一的提升
\(\tilde{\gamma}: [0, 1] \to \tilde{X}\) 使得
\(\tilde{\gamma}(0) = \tilde{x}_0\)。

\emph{证明}：由 Lebesgue 数引理，将 \([0, 1]\) 分割为
\(0 = t_0 < t_1 < \cdots < t_n = 1\) 使得每个 \(\gamma([t_i, t_{i+1}])\)
包含在某个基本邻域 \(U_i\) 中。在 \(U_0\) 上，取包含 \(\tilde{x}_0\)
的那个片 \(V_0\)，定义
\(\tilde{\gamma}|_{[0, t_1]} = (p|_{V_0})^{-1} \circ \gamma|_{[0, t_1]}\)。然后归纳地继续，每次从已确定的终点出发。唯一性由每一步的局部同胚性质保证。∎

\textbf{定理 61.3}（同伦提升引理）：设 \(p: \tilde{X} \to X\)
为覆叠映射，\(H: Y \times [0, 1] \to X\)
为同伦，\(\tilde{f}: Y \to \tilde{X}\) 是 \(H(\cdot, 0)\) 的提升。则存在
\(H\) 的唯一提升 \(\tilde{H}: Y \times [0, 1] \to \tilde{X}\) 使得
\(\tilde{H}(\cdot, 0) = \tilde{f}\)。

特别地，道路同伦可以唯一提升。

\emph{证明}：对每个 \(y \in Y\)，考虑道路
\(t \mapsto H(y, t)\)，由道路提升引理提升为
\(\tilde{H}(y, \cdot)\)。需验证 \(\tilde{H}\) 连续，这可由 \(Y\)
的紧致性（对 \(Y = [0, 1]\) 情形）或一般地由覆叠映射的局部性质（使用
Lebesgue 数引理的推广）保证。∎

\textbf{定理 61.4}（提升判别法）：设 \(p: \tilde{X} \to X\)
为覆叠映射，\(f: Y \to X\) 连续，\(Y\) 道路连通且局部道路连通。取
\(y_0 \in Y\)，\(x_0 = f(y_0)\)，\(\tilde{x}_0 \in p^{-1}(x_0)\)。则
\(f\) 可提升为
\(\tilde{f}: Y \to \tilde{X}\)（\(\tilde{f}(y_0) = \tilde{x}_0\)）当且仅当

\[f_*(\pi_1(Y, y_0)) \subseteq p_*(\pi_1(\tilde{X}, \tilde{x}_0))\]

\emph{证明}：(\(\Rightarrow\)) 若 \(\tilde{f}\) 存在，则
\(f_* = p_* \circ \tilde{f}_*\)，故
\(f_*(\pi_1(Y)) = p_*(\tilde{f}_*(\pi_1(Y))) \subseteq p_*(\pi_1(\tilde{X}))\)。

(\(\Leftarrow\)) 对 \(y \in Y\)，取道路 \(\alpha\) 从 \(y_0\) 到
\(y\)。\(f \circ \alpha\) 是 \(X\) 中道路，提升到 \(\tilde{X}\) 中起点为
\(\tilde{x}_0\) 的道路 \(\widetilde{f \circ \alpha}\)。定义
\(\tilde{f}(y) = \widetilde{f \circ \alpha}(1)\)。需验证良定义：若
\(\beta\) 是另一条道路，则 \(\alpha \cdot \bar{\beta}\)
是回路，\([f \circ (\alpha \cdot \bar{\beta})] \in f_*(\pi_1(Y)) \subseteq p_*(\pi_1(\tilde{X}))\)，故其提升是
\(\tilde{X}\) 中的回路，从而
\(\widetilde{f \circ \alpha}(1) = \widetilde{f \circ \beta}(1)\)。连续性由局部道路连通性保证。∎

\subsubsection{61.3
覆叠空间与基本群}\label{ux8986ux53e0ux7a7aux95f4ux4e0eux57faux672cux7fa4}

\textbf{定理 61.5}：设 \(p: \tilde{X} \to X\)
为覆叠映射，\(\tilde{x}_0 \in \tilde{X}\)，\(x_0 = p(\tilde{x}_0)\)。则诱导同态
\(p_*: \pi_1(\tilde{X}, \tilde{x}_0) \to \pi_1(X, x_0)\) 是单射。

\emph{证明}：若 \([\tilde{\gamma}] \in \ker p_*\)，则
\(p \circ \tilde{\gamma}\) 在 \(X\)
中道路同伦于常值回路。由同伦提升引理，此同伦提升为 \(\tilde{X}\)
中的同伦，故 \(\tilde{\gamma}\) 道路同伦于常值回路，即
\([\tilde{\gamma}] = 1\)。∎

\textbf{定理 61.6}（\(\pi_1(S^1) \cong \mathbb{Z}\)
的完整证明）：覆叠映射
\(p: \mathbb{R} \to S^1\)，\(p(t) = e^{2\pi i t}\)。取
\(\tilde{x}_0 = 0\)，\(x_0 = 1\)。由于 \(\mathbb{R}\)
可缩，\(\pi_1(\mathbb{R}, 0) = 1\)，故 \(p_*\) 是单射。考虑纤维
\(p^{-1}(1) = \mathbb{Z}\)。

对任意回路 \(\gamma\) 在 \(S^1\) 中以 \(1\) 为基点，提升
\(\tilde{\gamma}\) 以 \(0\) 为起点，终点
\(\tilde{\gamma}(1) \in \mathbb{Z}\)。定义映射

\[\Phi: \pi_1(S^1, 1) \to \mathbb{Z}, \quad [\gamma] \mapsto \tilde{\gamma}(1)\]

由同伦提升引理，\(\Phi\) 是良定义的。\(\Phi\)
是同态：\(\widetilde{\gamma \cdot \delta}\) 的终点为
\(\tilde{\gamma}(1) + \tilde{\delta}(1)\)（因为 \(\tilde{\delta}\)
的起点为 \(\tilde{\gamma}(1)\)，平移后即得）。\(\Phi\)
是满射：\(n \in \mathbb{Z}\) 对应于
\(t \mapsto e^{2\pi i n t}\)。\(\Phi\) 是单射：若
\(\tilde{\gamma}(1) = 0\)，则 \(\tilde{\gamma}\) 是 \(\mathbb{R}\)
中回路，\(\mathbb{R}\) 可缩故 \(\tilde{\gamma}\) 同伦于常值回路，投影到
\(S^1\) 得 \(\gamma\) 同伦于常值回路。因此
\(\pi_1(S^1) \cong \mathbb{Z}\)。∎

\textbf{定理 61.7}（覆叠空间的分类）：设 \(X\)
道路连通、局部道路连通且半局部单连通。则 \(X\)
的覆叠空间（在同构意义下）与 \(\pi_1(X, x_0)\) 的子群共轭类一一对应。

具体地，覆叠 \(p: (\tilde{X}, \tilde{x}_0) \to (X, x_0)\) 对应于子群
\(p_*(\pi_1(\tilde{X}, \tilde{x}_0)) \subseteq \pi_1(X, x_0)\)。

\textbf{定义 61.4}（万有覆叠）：若 \(\tilde{X}\) 单连通（即
\(\pi_1(\tilde{X}) = 1\)），则称 \(p: \tilde{X} \to X\) 为 \(X\)
的\textbf{万有覆叠}。万有覆叠（若存在）在同构意义下唯一。

\textbf{命题 61.8}：\(X\) 存在万有覆叠当且仅当 \(X\)
是道路连通、局部道路连通且半局部单连通的。

\textbf{定义 61.5}（半局部单连通）：\(X\)
称为\textbf{半局部单连通的}，如果对每个 \(x \in X\)，存在邻域 \(U\) 使得
\(\pi_1(U, x) \to \pi_1(X, x)\) 的像是平凡子群。

\subsubsection{61.4 覆叠变换群}\label{ux8986ux53e0ux53d8ux6362ux7fa4}

\textbf{定义 61.6}（覆叠变换 / 甲板变换）：设 \(p: \tilde{X} \to X\)
为覆叠映射。覆叠 \(\tilde{X}\)
的\textbf{覆叠变换}（或甲板变换）是一个同胚
\(\varphi: \tilde{X} \to \tilde{X}\) 使得
\(p \circ \varphi = p\)。所有覆叠变换构成群，记作
\(\operatorname{Aut}(\tilde{X}/X)\) 或
\(\operatorname{Deck}(\tilde{X}/X)\)。

\textbf{定义 61.7}（正规覆叠）：覆叠 \(p: \tilde{X} \to X\)
称为\textbf{正规的}（或\textbf{Galois 覆叠}），如果对任意
\(\tilde{x}_1, \tilde{x}_2 \in \tilde{X}\) 满足
\(p(\tilde{x}_1) = p(\tilde{x}_2)\)，存在覆叠变换 \(\varphi\) 使得
\(\varphi(\tilde{x}_1) = \tilde{x}_2\)。

\textbf{定理 61.9}：设 \(p: (\tilde{X}, \tilde{x}_0) \to (X, x_0)\)
为覆叠映射，\(\tilde{X}\) 道路连通。则

\[\operatorname{Aut}(\tilde{X}/X) \cong N(p_*(\pi_1(\tilde{X}, \tilde{x}_0))) / p_*(\pi_1(\tilde{X}, \tilde{x}_0))\]

其中 \(N(H)\) 是 \(H\) 在 \(\pi_1(X, x_0)\) 中的正规化子。特别地，若
\(\tilde{X}\) 是万有覆叠，则
\(\operatorname{Aut}(\tilde{X}/X) \cong \pi_1(X, x_0)\)。

\emph{证明思路}：覆叠变换由基点的选择决定。对于
\([\gamma] \in N(p_*(\pi_1(\tilde{X})))\)，\(\gamma\) 的提升将
\(\tilde{x}_0\) 映到某点 \(\tilde{x}_1\)，导出覆叠变换。不同
\([\gamma]\) 给出相同变换当且仅当它们相差 \(p_*(\pi_1(\tilde{X}))\)
中元素。∎

\textbf{推论 61.10}：覆叠 \(p: \tilde{X} \to X\) 正规当且仅当
\(p_*(\pi_1(\tilde{X}))\) 是 \(\pi_1(X)\) 的正规子群。

\subsubsection{61.5
覆叠空间的应用}\label{ux8986ux53e0ux7a7aux95f4ux7684ux5e94ux7528}

\textbf{定理 61.11}（\(\pi_1(S^n)\)
平凡，\(n \geq 2\)，完整证明）：\(S^n\)（\(n \geq 2\)）是单连通的。

\emph{证明}：取
\(U = S^n \setminus \{N\}\)（去掉北极），\(V = S^n \setminus \{S\}\)（去掉南极）。则
\(U \cong \mathbb{R}^n\)，\(V \cong \mathbb{R}^n\)，\(U \cap V \cong S^{n-1} \times \mathbb{R}\)
同伦等价于 \(S^{n-1}\)。当 \(n \geq 2\) 时 \(S^{n-1}\) 道路连通。由
Seifert-van Kampen 定理，\(\pi_1(S^n)\) 由 \(\pi_1(U)\) 和 \(\pi_1(V)\)
生成，但两者都是平凡群（\(\mathbb{R}^n\) 可缩），故 \(\pi_1(S^n)\)
平凡。∎

\textbf{定理
61.12}（\(\mathbb{R}^n \not\cong \mathbb{R}^m\)，\(n \neq m\)）：若
\(n \neq m\)，则 \(\mathbb{R}^n\) 与 \(\mathbb{R}^m\) 不同胚。

\emph{证明}：若
\(\mathbb{R}^n \cong \mathbb{R}^m\)，则它们的单点紧化也同胚：\(S^n \cong S^m\)。但这对
\(n \neq m\) 不成立（可由同调群证明，此处仅用基本群：当
\(n=1, m \geq 2\) 时
\(\pi_1(S^1) \cong \mathbb{Z} \neq 1 \cong \pi_1(S^m)\)）。一般
\(n \neq m\) 的证明需要高维同伦群或同调群。∎

\textbf{定理
61.13}（基本群与自由群）：\(\pi_1(S^1 \vee S^1 \vee \cdots \vee S^1)\)（\(n\)
个圆周的楔积）同构于秩为 \(n\) 的自由群 \(F_n\)。

\emph{证明}：由 Seifert-van Kampen 定理（推论
60.12）归纳可得。万有覆叠是秩为 \(n\) 的自由群的 Cayley 图（无穷树）。∎

\textbf{定理 61.14}（Nielsen-Schreier 定理）：自由群的子群是自由群。

\emph{证明}（拓扑证明）：设 \(F_n\) 是秩为 \(n\)
的自由群，\(H \leq F_n\)。\(F_n \cong \pi_1(X)\) 其中 \(X\) 是 \(n\)
个圆周的楔积。由覆叠分类定理，存在覆叠 \(p: \tilde{X} \to X\) 使得
\(p_*(\pi_1(\tilde{X})) = H\)。\(\tilde{X}\) 是一维 CW
复形（图），其基本群是自由群，故 \(H\) 是自由群。∎

\bigskip\hrule\bigskip

\subsection{第62章：结晶群与空间群（Crystallographic and Space
Groups）}\label{ux7b2c62ux7ae0ux7ed3ux6676ux7fa4ux4e0eux7a7aux95f4ux7fa4crystallographic-and-space-groups}

结晶群和空间群是 \(\mathbb{R}^n\)
中等距作用的离散上紧子群，构成了数学结晶学和固体物理中晶体对称性的严格数学语言。\textbf{定义}：\(n\)
维 \textbf{空间群} \(\Gamma\) 是等距群
\(\mathbb{R}^n \rtimes \operatorname{O}(n)\) 的离散子群，使得
\(\mathbb{R}^n / \Gamma\)
紧致。\textbf{基本定理}：（1）\textbf{Bieberbach
第一定理}------空间群的平移子群 \(\Gamma \cap \mathbb{R}^n\) 是秩为
\(n\) 的自由 Abel 群（格），即每个空间群是 \(\mathbb{Z}^n\)
被有限群（点群）扩张；（2）\textbf{Bieberbach
第二定理}------若两个空间群同构（抽象群同构），则它们作为仿射群仿射共轭；（3）\textbf{Bieberbach
第三定理}------每个维数 \(n\)
仅有\textbf{有限多个}空间群（至多相差仿射共轭）。\textbf{分类结果}：\(n=1\)
有 2 个（线群）；\(n=2\) 有 17 个（壁纸群 / wallpaper groups）；\(n=3\)
有 219 个（Fedorov
空间群，1891年分类完成，若区分对映体则有230个）。高维分类是部分解决的开问题------\(n=4\)
有 4,783 个空间群（Brown-Bülow-Neubüser-Wondratschek-Zassenhaus,
1978）。结晶群通过其同调 \(H^*(\Gamma; \mathbb{Z})\)
中的扭元与物理中拓扑绝缘体的分类（Kitaev 周期表）有深刻联系。

\bigskip\hrule\bigskip

\emph{本卷构建了拓扑学的基础框架：从拓扑空间的定义出发，建立了连续性、连通性、紧致性等核心概念，证明了
Urysohn 引理、Tychonoff 定理、Urysohn
度量化定理等核心定理，进而引入基本群和覆叠空间理论，为后续学习代数拓扑（V15）、微分几何（V12）和泛函分析（V14）奠定了基础。}

\newpage

\section{卷十一：微分几何}\label{ux5377ux5341ux4e00ux5faeux5206ux51e0ux4f55}

\begin{quote}
微分几何用微积分和线性代数的工具研究曲线、曲面以及更一般的微分流形的几何性质。本卷从经典曲线论和曲面论出发，逐步引入微分流形和
Riemann 几何的现代框架，最终以 Gauss-Bonnet
定理揭示曲率与拓扑的深层联系。
\end{quote}

\bigskip\hrule\bigskip

\subsection{第62章：曲线论}\label{ux7b2c62ux7ae0ux66f2ux7ebfux8bba}

曲线论研究 \(\mathbb{R}^3\) 中曲线的局部几何性质，核心工具是 Frenet
标架和曲率不变量。

\subsubsection{62.1 参数化曲线}\label{ux53c2ux6570ux5316ux66f2ux7ebf}

\textbf{定义 62.1}（参数化曲线）：设 \(I \subseteq \mathbb{R}\)
为区间。一个\textbf{参数化曲线}是一个光滑（\(C^\infty\)）映射
\(\gamma: I \to \mathbb{R}^3\)。\(\gamma'(t)\)
称为曲线的\textbf{速度向量}，\(\|\gamma'(t)\|\) 称为\textbf{速率}。

\textbf{定义 62.2}（正则曲线）：曲线 \(\gamma\)
称为\textbf{正则的}，如果对所有 \(t \in I\)，\(\gamma'(t) \neq 0\)。

\textbf{定义 62.3}（重新参数化）：设 \(\gamma: I \to \mathbb{R}^3\)
是正则曲线，\(\varphi: J \to I\)
是光滑的微分同胚（\(\varphi' \neq 0\)）。则
\(\tilde{\gamma} = \gamma \circ \varphi: J \to \mathbb{R}^3\) 称为
\(\gamma\) 的\textbf{重新参数化}。

\textbf{定义 62.4}（弧长参数）：正则曲线 \(\gamma: I \to \mathbb{R}^3\)
的\textbf{弧长函数}定义为

\[s(t) = \int_{t_0}^t \|\gamma'(u)\| du\]

若 \(\|\gamma'(s)\| = 1\) 对所有 \(s\) 成立，则称 \(\gamma\)
以\textbf{弧长}为参数。此时 \(\gamma'(s)\) 是曲线在 \(\gamma(s)\)
处的\textbf{单位切向量}，记作 \(T(s)\)。

\textbf{命题 62.1}：每条正则曲线都可以重新参数化为弧长参数。

\emph{证明}：由于 \(\gamma\) 正则，\(\|\gamma'(t)\| > 0\)，故 \(s(t)\)
严格单调递增，存在光滑逆函数 \(t(s)\)。令
\(\tilde{\gamma}(s) = \gamma(t(s))\)，则
\(\|\tilde{\gamma}'(s)\| = \|\gamma'(t(s)) \cdot t'(s)\| = \|\gamma'(t(s))\| \cdot (1/\|\gamma'(t(s))\|) = 1\)。∎

\subsubsection{62.2 Frenet 标架}\label{frenet-ux6807ux67b6}

以下设 \(\gamma(s)\) 以弧长为参数，\(T(s) = \gamma'(s)\) 是单位切向量。

\textbf{定义 62.5}（曲率）：\(\gamma\) 在 \(s\) 处的\textbf{曲率}定义为
\(\kappa(s) = \|T'(s)\| = \|\gamma''(s)\|\)。

\textbf{定义 62.6}（主法向量）：若
\(\kappa(s) > 0\)，定义\textbf{主法向量}
\(N(s) = T'(s) / \kappa(s)\)。则 \(N(s)\) 是单位向量且垂直于
\(T(s)\)（因为 \(\|T(s)\| = 1\) 蕴含 \(T'(s) \cdot T(s) = 0\)）。

\textbf{定义 62.7}（副法向量与挠率）：定义\textbf{副法向量}
\(B(s) = T(s) \times N(s)\)。则 \(\{T(s), N(s), B(s)\}\) 构成
\(\mathbb{R}^3\) 的一组正交规范基，称为 \(\gamma\) 在 \(s\)
处的\textbf{Frenet 标架}。

\textbf{挠率} \(\tau(s)\) 定义为 \(\tau(s) = -B'(s) \cdot N(s)\)。

\textbf{命题 62.2}：\(B'(s) = -\tau(s) N(s)\)。

\emph{证明}：由 \(B(s) \cdot B(s) = 1\) 得 \(B'(s) \cdot B(s) = 0\)。又
\(B = T \times N\)，故
\(B' = T' \times N + T \times N' = \kappa N \times N + T \times N' = T \times N'\)，因此
\(B' \cdot T = 0\)。于是 \(B'\) 平行于 \(N\)，即
\(B'(s) = -\tau(s) N(s)\)。∎

\textbf{定理 62.3}（Frenet-Serret 公式）：

\[\begin{aligned}
T'(s) &= \kappa(s) N(s) \\
N'(s) &= -\kappa(s) T(s) + \tau(s) B(s) \\
B'(s) &= -\tau(s) N(s)
\end{aligned}\]

\emph{证明}：第一式是曲率的定义。第三式是挠率的定义。对第二式，由
\(N = B \times T\)（注意 \(T \times N = B\) 得
\(N = B \times T\)），求导：

\[N' = B' \times T + B \times T' = (-\tau N) \times T + B \times (\kappa N) = -\tau (N \times T) + \kappa (B \times N) = -\tau (-B) + \kappa (-T) = -\kappa T + \tau B\]

此处使用了 \(N \times T = -B\)，\(B \times N = -T\)。∎

\textbf{定理 62.4}（曲线论基本定理）：给定光滑函数 \(\kappa(s) > 0\) 和
\(\tau(s)\)（\(s \in I\)），则存在唯一（在刚体运动意义下）的以弧长为参数的正则曲线
\(\gamma: I \to \mathbb{R}^3\)，以 \(\kappa(s)\) 为曲率、\(\tau(s)\)
为挠率。

\emph{证明}：Frenet-Serret 公式构成关于 \((T, N, B)\) 的线性 ODE
系统。由 ODE 解的存在唯一性定理（V4 第 23 章的推广），给定初始条件
\((T(s_0), N(s_0), B(s_0))\) 为一个正交规范标架，存在唯一解。然后
\(\gamma(s) = \gamma(s_0) + \int_{s_0}^s T(u) du\)
给出所需曲线。不同的初始条件对应于曲线的刚体运动。∎

\textbf{推论
62.5}：曲线的曲率和挠率是刚体运动不变量，且完全决定了曲线的形状。

\subsubsection{62.3 平面曲线}\label{ux5e73ux9762ux66f2ux7ebf}

\textbf{定义 62.8}（平面曲线）：若 \(\gamma\) 的像落在某个平面内，则
\(\gamma\) 称为\textbf{平面曲线}。

\textbf{命题 62.6}：\(\gamma\) 是平面曲线当且仅当 \(\tau(s) \equiv 0\)。

\emph{证明}：(\(\Rightarrow\)) 若 \(\gamma\) 是平面曲线，则 \(T(s)\) 和
\(N(s)\) 始终在该平面内，\(B(s)\) 是常值单位法向量，故
\(B'(s) = 0\)，从而 \(\tau(s) = 0\)。

(\(\Leftarrow\)) 若 \(\tau \equiv 0\)，则 \(B(s)\) 是常值向量
\(B_0\)。由
\((T \cdot B_0)' = T' \cdot B_0 = \kappa N \cdot B_0 = 0\)，故
\(T(s) \cdot B_0\) 为常数，即 \(\gamma(s) \cdot B_0\) 为常数，因此
\(\gamma\) 落在与 \(B_0\) 垂直的平面内。∎

\textbf{定义
62.9}（平面曲线的符号曲率）：对平面曲线，可定义\textbf{符号曲率}
\(k(s)\)：\(T'(s) = k(s) N(s)\)。\(k(s) > 0\)
表示曲线向左转，\(k(s) < 0\) 表示向右转。

\bigskip\hrule\bigskip

\subsection{第63章：曲面论}\label{ux7b2c63ux7ae0ux66f2ux9762ux8bba}

曲面论是经典微分几何的核心，研究 \(\mathbb{R}^3\)
中二维曲面的内蕴几何与外蕴几何。

\subsubsection{63.1 正则曲面}\label{ux6b63ux5219ux66f2ux9762}

\textbf{定义 63.1}（参数化曲面）：设 \(U \subseteq \mathbb{R}^2\)
为开集。一个\textbf{参数化曲面}是一个光滑映射
\(\mathbf{x}: U \to \mathbb{R}^3\)，\(\mathbf{x}(u, v) = (x(u, v), y(u, v), z(u, v))\)。

\textbf{定义 63.2}（正则曲面）：参数化曲面 \(\mathbf{x}\)
称为\textbf{正则的}，如果对每个 \((u, v) \in U\)，偏导数
\(\mathbf{x}_u\) 和 \(\mathbf{x}_v\) 线性无关（即
\(\mathbf{x}_u \times \mathbf{x}_v \neq 0\)）。此时 \(\mathbf{x}_u\) 和
\(\mathbf{x}_v\) 张成曲面在 \(\mathbf{x}(u, v)\) 处的\textbf{切平面}。

\textbf{定义 63.3}（正则曲面------整体定义）：\(\mathbb{R}^3\) 的子集
\(S\) 称为\textbf{正则曲面}，如果对每个 \(p \in S\)，存在 \(p\) 在
\(\mathbb{R}^3\) 中的邻域 \(V\) 和映射
\(\mathbf{x}: U \to V \cap S\)（\(U \subseteq \mathbb{R}^2\)
开集），使得 \(\mathbf{x}\) 是光滑同胚且是正则参数化曲面。\(\mathbf{x}\)
称为 \(S\) 在 \(p\) 处的\textbf{局部参数化}（或\textbf{坐标卡}）。

\subsubsection{63.2
第一基本形式}\label{ux7b2cux4e00ux57faux672cux5f62ux5f0f}

\textbf{定义 63.4}（第一基本形式）：设 \(\mathbf{x}: U \to S\)
是正则曲面的局部参数化。\(S\) 的\textbf{第一基本形式}是二次型

\[I_p(w) = \langle w, w \rangle_p = \|w\|^2, \quad w \in T_p S\]

其中 \(T_p S\) 是 \(S\) 在 \(p\)
处的切平面，\(\langle \cdot, \cdot \rangle_p\) 是 \(\mathbb{R}^3\)
的内积限制到 \(T_p S\)。

在局部坐标 \((u, v)\) 下，切向量 \(w = a \mathbf{x}_u + b \mathbf{x}_v\)
的第一基本形式为

\[I_p(w) = E a^2 + 2F a b + G b^2\]

其中

\[E = \langle \mathbf{x}_u, \mathbf{x}_u \rangle, \quad F = \langle \mathbf{x}_u, \mathbf{x}_v \rangle, \quad G = \langle \mathbf{x}_v, \mathbf{x}_v \rangle\]

称为第一基本形式的\textbf{系数}。

\textbf{定义 63.5}（等距）：两个正则曲面 \(S_1, S_2\)
称为\textbf{等距的}（或\textbf{局部等距的}），如果存在微分同胚
\(\varphi: S_1 \to S_2\) 使得对任意 \(p \in S_1\) 和
\(w \in T_p S_1\)，\(I_p^1(w) = I_{\varphi(p)}^2(d\varphi_p(w))\)。即第一基本形式被保持。

\textbf{命题 63.1}：第一基本形式完全决定了曲面上的弧长、角度和面积：

\begin{itemize}
\tightlist
\item
  曲线 \(\gamma(t) = \mathbf{x}(u(t), v(t))\)
  的弧长：\(L = \int \sqrt{E u'^2 + 2F u' v' + G v'^2} dt\)
\item
  两切向量的夹角：\(\cos \theta = \dfrac{E a_1 a_2 + F(a_1 b_2 + a_2 b_1) + G b_1 b_2}{\sqrt{E a_1^2 + 2F a_1 b_1 + G b_1^2} \sqrt{E a_2^2 + 2F a_2 b_2 + G b_2^2}}\)
\item
  区域 \(R \subseteq U\) 的面积：\(A = \iint_R \sqrt{EG - F^2} du dv\)
\end{itemize}

\subsubsection{63.3
第二基本形式与法曲率}\label{ux7b2cux4e8cux57faux672cux5f62ux5f0fux4e0eux6cd5ux66f2ux7387}

\textbf{定义 63.6}（单位法向量）：对正则参数化曲面
\(\mathbf{x}: U \to S\)，定义单位法向量

\[N(p) = \frac{\mathbf{x}_u \times \mathbf{x}_v}{\|\mathbf{x}_u \times \mathbf{x}_v\|}(p)\]

（注意 \(N\) 取决于 \(\mathbf{x}_u \times \mathbf{x}_v\)
的方向，即曲面的定向。）

\textbf{定义 63.7}（Gauss 映射）：映射 \(N: S \to S^2\)（\(S^2\) 是
\(\mathbb{R}^3\) 中的单位球面）将曲面上每点映到其单位法向量，称为曲面的
\textbf{Gauss 映射}。

\textbf{定义 63.8}（第二基本形式）：\(S\) 的\textbf{第二基本形式}定义为

\[II_p(w) = -\langle dN_p(w), w \rangle, \quad w \in T_p S\]

在局部坐标下，\(dN_p(\mathbf{x}_u) = N_u\)，\(dN_p(\mathbf{x}_v) = N_v\)。第二基本形式的系数为

\[e = -\langle N_u, \mathbf{x}_u \rangle = \langle N, \mathbf{x}_{uu} \rangle\]
\[f = -\langle N_u, \mathbf{x}_v \rangle = \langle N, \mathbf{x}_{uv} \rangle\]
\[g = -\langle N_v, \mathbf{x}_v \rangle = \langle N, \mathbf{x}_{vv} \rangle\]

（后一等式由 \(\langle N, \mathbf{x}_u \rangle = 0\) 求导得到。）

\textbf{定义 63.9}（法曲率）：设 \(v \in T_p S\) 是单位切向量。曲面在
\(p\) 处沿方向 \(v\) 的\textbf{法曲率}为

\[\kappa_n(v) = II_p(v) = e a^2 + 2f a b + g b^2\]

其中 \(v = a \mathbf{x}_u + b \mathbf{x}_v\) 且 \(I_p(v) = 1\)。

\textbf{定义 63.10}（主曲率、Gauss 曲率、平均曲率）：法曲率
\(\kappa_n(v)\) 作为 \(v\) 的函数（在单位切向量上）在 \(T_p S\)
上取得最大值 \(\kappa_1\) 和最小值 \(\kappa_2\)，称为 \(S\) 在 \(p\)
处的\textbf{主曲率}。对应方向称为\textbf{主方向}。

\begin{itemize}
\tightlist
\item
  \textbf{Gauss
  曲率}：\(K = \kappa_1 \kappa_2 = \dfrac{eg - f^2}{EG - F^2}\)
\item
  \textbf{平均曲率}：\(H = \dfrac{\kappa_1 + \kappa_2}{2} = \dfrac{eG - 2fF + gE}{2(EG - F^2)}\)
\end{itemize}

\textbf{定理 63.2}（主曲率的欧拉公式）：设 \(v\) 与主方向 \(e_1\)
的夹角为 \(\theta\)，则

\[\kappa_n(v) = \kappa_1 \cos^2 \theta + \kappa_2 \sin^2 \theta\]

\emph{证明}：在由主方向构成的正交基下，\(dN\) 是对角矩阵
\(\operatorname{diag}(-\kappa_1, -\kappa_2)\)。设
\(v = \cos \theta e_1 + \sin \theta e_2\)，则
\(II(v) = \kappa_1 \cos^2 \theta + \kappa_2 \sin^2 \theta\)。∎

\textbf{定义 63.11}（点分类）：曲面上点 \(p\) 称为 -
\textbf{椭圆点}：\(K > 0\)（\(\kappa_1\) 和 \(\kappa_2\) 同号） -
\textbf{双曲点}：\(K < 0\)（\(\kappa_1\) 和 \(\kappa_2\) 异号） -
\textbf{抛物点}：\(K = 0\) 但 \(H \neq 0\)（一个主曲率为零） -
\textbf{平点}：\(K = 0\) 且 \(H = 0\)（两个主曲率为零）

\subsubsection{63.4 Gauss
绝妙定理}\label{gauss-ux7eddux5999ux5b9aux7406}

\textbf{定理 63.3}（Gauss 绝妙定理 / Theorema Egregium）：曲面的 Gauss
曲率 \(K\) 仅由第一基本形式及其导数决定，即是曲面的\textbf{内蕴}不变量。

\emph{证明思路}：引入 Christoffel 符号：

\[\Gamma_{ij}^k = \frac{1}{2} \sum_{\ell} g^{k\ell}(\partial_i g_{j\ell} + \partial_j g_{i\ell} - \partial_\ell g_{ij})\]

其中 \(g_{11} = E, g_{12} = g_{21} = F, g_{22} = G\)
是第一基本形式的系数，\((g^{ij})\) 是 \((g_{ij})\) 的逆矩阵。

Gauss 曲率可表示为

\[K = \frac{R_{1212}}{g_{11}g_{22} - g_{12}^2}\]

其中 \(R_{1212}\) 是 Riemann
曲率张量的分量（在二维情形下只有一个独立分量），它由 Christoffel
符号及其导数组合而成，从而仅依赖于第一基本形式。具体地：

\[R_{1212} = \frac{1}{2}\left(\frac{\partial^2 g_{11}}{\partial v^2} + \frac{\partial^2 g_{22}}{\partial u^2} - 2\frac{\partial^2 g_{12}}{\partial u \partial v}\right) + \sum_{p,q} g_{pq}(\Gamma_{22}^p \Gamma_{11}^q - \Gamma_{21}^p \Gamma_{21}^q)\]

或者等价地用 Christoffel 符号的偏导数表示：

\[R_{1212} = \partial_1\Gamma_{22}^1 - \partial_2\Gamma_{21}^1 + \Gamma_{22}^1\Gamma_{11}^1 + \Gamma_{22}^2\Gamma_{12}^1 - \Gamma_{21}^1\Gamma_{21}^1 - \Gamma_{21}^2\Gamma_{22}^1\]

（其中
\(\partial_1 = \partial_u\)，\(\partial_2 = \partial_v\)）。这表明 \(K\)
完全由 \(E, F, G\) 及其导数决定，与曲面在 \(\mathbb{R}^3\)
中的嵌入方式无关。∎

\textbf{推论 63.4}：等距的曲面具有相同的 Gauss
曲率。特别地，不存在从平面（\(K = 0\)）到球面（\(K = 1/R^2 > 0\)）的等距映射（即无法将地图无变形地绘制在平面上）。

\subsubsection{63.5
测地线与协变导数}\label{ux6d4bux5730ux7ebfux4e0eux534fux53d8ux5bfcux6570}

\textbf{定义 63.12}（测地曲率）：设曲线 \(\gamma(s) \subset S\)
以弧长为参数。\(\gamma''(s)\) 在切平面上的投影的模长称为 \(\gamma\) 在
\(S\) 上的\textbf{测地曲率} \(k_g\)。

\textbf{定义 63.13}（测地线）：曲面 \(S\) 上的曲线 \(\gamma\)
称为\textbf{测地线}，如果其测地曲率恒为零，即 \(\gamma''(s)\)
处处垂直于切平面。

\textbf{命题 63.5}（测地线方程）：在局部坐标 \((u, v)\)
下，测地线满足微分方程：

\[\frac{d^2 u^k}{ds^2} + \sum_{i,j=1}^2 \Gamma_{ij}^k \frac{du^i}{ds}\frac{du^j}{ds} = 0, \quad k = 1, 2\]

\emph{证明}：\(\gamma''(s)\) 垂直于 \(T_p S\) 等价于 \(\gamma''(s)\) 与
\(\mathbf{x}_u\) 和 \(\mathbf{x}_v\) 的内积为零。展开
\(\gamma''(s) = \frac{d^2 u^k}{ds^2} \mathbf{x}_k + \frac{du^i}{ds}\frac{du^j}{ds} \mathbf{x}_{ij}\)，利用
\(\mathbf{x}_{ij} = \Gamma_{ij}^k \mathbf{x}_k + L_{ij} N\)（Gauss
公式），取切向分量即得测地线方程。∎

\textbf{命题 63.6}：曲面上连接两点的最短曲线（若存在）必为测地线。

\textbf{例 63.1}（测地线举例）： - 平面上的测地线是直线。 -
球面上的测地线是大圆。 - 柱面上的测地线是螺旋线。

\bigskip\hrule\bigskip

\subsection{第64章：微分流形}\label{ux7b2c64ux7ae0ux5faeux5206ux6d41ux5f62}

微分流形将曲面概念推广到任意维数，是现代微分几何的基石。微分流形是一种局部同胚于欧氏空间的拓扑空间，每个点附近可通过局部坐标（坐标卡）用
\(\mathbb{R}^n\)
中的坐标描述。不同坐标卡之间的重叠区域通过光滑的过渡映射（转移函数）相互联系，这种结构使得微积分运算得以在流形上全局一致地进行。

\subsubsection{64.1
微分流形的定义}\label{ux5faeux5206ux6d41ux5f62ux7684ux5b9aux4e49}

\textbf{定义 64.1}（拓扑流形）：设 \(M\) 是 Hausdorff
空间且具有可数基。\(M\) 称为 \(n\) 维\textbf{拓扑流形}，如果对每个
\(p \in M\)，存在 \(p\) 的开邻域 \(U\) 同胚于 \(\mathbb{R}^n\)
的某个开集。

\textbf{定义 64.2}（坐标卡与图册）：设 \(M\) 是 \(n\) 维拓扑流形。 -
\(M\) 上的一个\textbf{坐标卡}（或局部坐标系）是一个对
\((U, \varphi)\)，其中 \(U \subseteq M\)
是开集，\(\varphi: U \to \mathbb{R}^n\) 是同胚到其像。 - \(M\)
的一个\textbf{图册} \(\mathcal{A} = \{(U_\alpha, \varphi_\alpha)\}\)
是一族坐标卡满足 \(\bigcup_\alpha U_\alpha = M\)。 - 两个坐标卡
\((U_\alpha, \varphi_\alpha)\) 和 \((U_\beta, \varphi_\beta)\) 称为
\(C^\infty\)\textbf{相容的}，如果
\(U_\alpha \cap U_\beta = \emptyset\)，或者\textbf{转移映射}

\[\varphi_\beta \circ \varphi_\alpha^{-1}: \varphi_\alpha(U_\alpha \cap U_\beta) \to \varphi_\beta(U_\alpha \cap U_\beta)\]

是 \(C^\infty\) 微分同胚。

\textbf{定义 64.3}（微分流形）：一个
\(C^\infty\)\textbf{微分流形}（或光滑流形）是一个拓扑流形 \(M\)
配上一个极大 \(C^\infty\) 图册（即所有与图册中坐标卡 \(C^\infty\)
相容的坐标卡都已在图册中）。

\textbf{定义 64.4}（光滑函数）：设 \(M\) 是光滑流形。函数
\(f: M \to \mathbb{R}\) 称为\textbf{光滑的}，如果对每个坐标卡
\((U, \varphi)\)，\(f \circ \varphi^{-1}: \varphi(U) \to \mathbb{R}\)
是光滑的（作为 \(\mathbb{R}^n\) 上开集的函数）。\(M\)
上所有光滑函数构成的集合记作 \(C^\infty(M)\)。

\textbf{定义 64.5}（光滑映射）：设 \(M, N\) 是光滑流形。映射
\(F: M \to N\) 称为\textbf{光滑的}，如果对每对坐标卡
\((U, \varphi)\)（\(M\) 上）和 \((V, \psi)\)（\(N\) 上）满足
\(F(U) \subseteq V\)，映射
\(\psi \circ F \circ \varphi^{-1}: \varphi(U) \to \psi(V)\) 是光滑的。

\textbf{例 64.1}（光滑流形举例）： - \(\mathbb{R}^n\)：单个坐标卡
\((\mathbb{R}^n, \operatorname{id})\)。 -
\(S^n\)：球极投影给出两个坐标卡覆盖 \(S^n\)。 - 实射影空间
\(\mathbb{R}P^n\)：齐次坐标 \([x_0 : \cdots : x_n]\) 的 \(n+1\)
个坐标卡。 - 环面 \(T^n = S^1 \times \cdots \times S^1\)：积流形。 -
开子集、积流形：光滑流形的开子集和积是光滑流形。

\subsubsection{64.2
切空间与切映射}\label{ux5207ux7a7aux95f4ux4e0eux5207ux6620ux5c04}

\textbf{定义 64.6}（芽与导子）：设 \(p \in M\)。\(p\)
处的\textbf{函数芽}是在 \(p\)
的某个邻域上定义的光滑函数的等价类（两个函数等价如它们在 \(p\)
的某个邻域上相等）。记作 \(C^\infty_p(M)\)。

\(p\) 处的一个\textbf{导子}（或切向量）是一个线性映射
\(v: C^\infty_p(M) \to \mathbb{R}\) 满足 Leibniz 法则：

\[v(fg) = v(f)g(p) + f(p)v(g)\]

\textbf{定义 64.7}（切空间）：\(M\) 在 \(p\) 处的\textbf{切空间}
\(T_p M\) 是所有 \(p\) 处导子构成的向量空间。其维数为 \(n = \dim M\)。

在局部坐标 \((x^1, \ldots, x^n)\)（由坐标卡 \(\varphi\)
给出）下，偏导算子 \(\frac{\partial}{\partial x^i}\big|_p\) 构成
\(T_p M\) 的一组基。任意切向量可表示为

\[v = \sum_{i=1}^n v^i \frac{\partial}{\partial x^i}\bigg|_p\]

\textbf{定义 64.8}（切映射/微分）：设 \(F: M \to N\)
是光滑映射，\(p \in M\)。\(F\) 在 \(p\)
处的\textbf{切映射}（或微分）\(dF_p: T_p M \to T_{F(p)} N\) 定义为

\[(dF_p(v))(f) = v(f \circ F), \quad \forall v \in T_p M, \; f \in C^\infty_{F(p)}(N)\]

在局部坐标下，\(dF_p\) 由 Jacobi 矩阵表示。

\textbf{定理
64.1}（链式法则）：\((d(G \circ F))_p = dG_{F(p)} \circ dF_p\)。

\textbf{定义 64.9}（切丛）：\(M\) 的\textbf{切丛}
\(TM = \bigsqcup_{p \in M} T_p M\) 具有自然的光滑流形结构，维数为
\(2n\)。自然投影 \(\pi: TM \to M\) 将 \(v \in T_p M\) 映到 \(p\)。

\subsubsection{64.3 向量场与 Lie
括号}\label{ux5411ux91cfux573aux4e0e-lie-ux62ecux53f7}

\textbf{定义 64.10}（向量场）：\(M\)
上的一个\textbf{光滑向量场}是一个光滑映射 \(X: M \to TM\) 使得
\(\pi \circ X = \operatorname{id}_M\)，即 \(X(p) \in T_p M\) 对每个
\(p\)。所有光滑向量场构成的集合记作 \(\mathfrak{X}(M)\)。

在局部坐标下，\(X = \sum_{i=1}^n X^i \frac{\partial}{\partial x^i}\)，其中
\(X^i \in C^\infty(M)\)。

\textbf{定义 64.11}（Lie 括号）：设 \(X, Y \in \mathfrak{X}(M)\)。它们的
\textbf{Lie 括号} \([X, Y] \in \mathfrak{X}(M)\) 定义为

\[[X, Y](f) = X(Y(f)) - Y(X(f)), \quad \forall f \in C^\infty(M)\]

在局部坐标下，\([X, Y] = \sum_{i,j} \left(X^j \frac{\partial Y^i}{\partial x^j} - Y^j \frac{\partial X^i}{\partial x^j}\right) \frac{\partial}{\partial x^i}\)。

\textbf{命题 64.2}（Lie 括号的性质）： 1.
\([X, Y] = -[Y, X]\)（反交换性） 2.
\([aX + bY, Z] = a[X, Z] + b[Y, Z]\)（双线性性） 3.
\([X, [Y, Z]] + [Y, [Z, X]] + [Z, [X, Y]] = 0\)（Jacobi 恒等式） 4.
\([fX, gY] = fg[X, Y] + f(Xg)Y - g(Yf)X\)

\subsubsection{64.4 微分形式}\label{ux5faeux5206ux5f62ux5f0f}

\textbf{定义 64.12}（余切空间与 1-形式）：\(T_p^* M = (T_p M)^*\) 是
\(T_p M\) 的对偶空间，称为 \(p\) 处的\textbf{余切空间}。\(M\)
上的一个\textbf{光滑 1-形式}是一个光滑映射
\(\omega: M \to T^*M = \bigsqcup T_p^* M\)，使得
\(\omega(p) \in T_p^* M\)。

在局部坐标下，\(\{dx^1, \ldots, dx^n\}\) 是 \(T_p^* M\) 的基（与
\(\{\partial/\partial x^i\}\) 对偶）。1-形式可写为
\(\omega = \sum \omega_i dx^i\)。

\textbf{定义 64.13}（微分 \(d\)）：对光滑函数
\(f \in C^\infty(M)\)，其\textbf{微分} \(df\) 是 1-形式，定义为
\(df(v) = v(f)\) 对所有
\(v \in TM\)。在局部坐标下，\(df = \sum \frac{\partial f}{\partial x^i} dx^i\)。

\textbf{定义 64.14}（\(k\)-形式）：\(M\) 上的一个\textbf{微分
\(k\)-形式}是一个反对称的 \((0, k)\)-张量场。在局部坐标下，

\[\omega = \sum_{i_1 < \cdots < i_k} \omega_{i_1 \cdots i_k} dx^{i_1} \wedge \cdots \wedge dx^{i_k}\]

其中 \(\wedge\)
是外积（楔积），满足反对称性：\(dx^i \wedge dx^j = -dx^j \wedge dx^i\)。

\textbf{定义 64.15}（外微分）：外微分算子
\(d: \Omega^k(M) \to \Omega^{k+1}(M)\) 定义为

\[d\omega = \sum_{i_1 < \cdots < i_k} \sum_{j} \frac{\partial \omega_{i_1 \cdots i_k}}{\partial x^j} dx^j \wedge dx^{i_1} \wedge \cdots \wedge dx^{i_k}\]

\textbf{命题 64.3}（外微分的基本性质）： 1. \(d\) 是线性的。 2.
\(d \circ d = 0\)（即 \(d(d\omega) = 0\) 对所有 \(\omega\)）。 3.
\(d(\omega \wedge \eta) = d\omega \wedge \eta + (-1)^{\deg \omega} \omega \wedge d\eta\)。
4. 对函数 \(f\)，\(df\) 与定义 64.13 一致。 5. \(d\) 是局部的：若
\(\omega = \eta\) 在开集 \(U\) 上，则 \(d\omega = d\eta\) 在 \(U\) 上。

\textbf{定义 64.16}（闭形式与恰当形式）：\(k\)-形式 \(\omega\) 称为 -
\textbf{闭的}：若 \(d\omega = 0\)。 - \textbf{恰当的}：若存在
\((k-1)\)-形式 \(\eta\) 使得 \(\omega = d\eta\)。

由 \(d^2 = 0\)，恰当形式一定是闭的。闭形式商去恰当形式得到 \textbf{de
Rham 上同调}：

\[H^k_{dR}(M) = \frac{\ker(d: \Omega^k(M) \to \Omega^{k+1}(M))}{\operatorname{im}(d: \Omega^{k-1}(M) \to \Omega^k(M))}\]

\subsubsection{64.5
子流形与浸入}\label{ux5b50ux6d41ux5f62ux4e0eux6d78ux5165}

\textbf{定义 64.17}（浸入与淹没）：光滑映射 \(F: M \to N\) 称为 -
\textbf{浸入}：若 \(dF_p\) 在每点 \(p\) 是单射。 - \textbf{淹没}：若
\(dF_p\) 在每点 \(p\) 是满射。 - \textbf{嵌入}：若 \(F\) 是浸入且
\(F: M \to F(M)\) 是同胚（\(F(M)\) 赋予 \(N\) 的子空间拓扑）。

\textbf{定理 64.4}（正则值定理）：设 \(F: M \to N\)
是光滑映射，\(\dim M \geq \dim N\)。若 \(q \in N\) 是 \(F\)
的\textbf{正则值}（即对所有 \(p \in F^{-1}(q)\)，\(dF_p\) 是满射），则
\(F^{-1}(q)\) 是 \(M\) 的 \(\dim M - \dim N\) 维嵌入子流形。

\textbf{定理 64.5}（Whitney 嵌入定理，陈述）：每个 \(n\)
维光滑流形可以嵌入到 \(\mathbb{R}^{2n}\) 中。

\bigskip\hrule\bigskip

\subsection{第65章：Riemann
几何}\label{ux7b2c65ux7ae0riemann-ux51e0ux4f55}

Riemann
几何在微分流形上赋予度量结构，使得可以讨论长度、角度、曲率等概念。Riemann
度量 \(g\)
在每个切空间上定义内积，从而决定曲线弧长、测地线方程和体积元。测地线是流形上''最直''的曲线，对应于欧氏空间中直线的推广；而
Riemann
曲率张量则刻画了流形偏离平坦的程度------它既是局部几何的核心不变量，又通过
Gauss-Bonnet 定理等结果与流形的整体拓扑建立了深刻联系。

\subsubsection{65.1 Riemann 度量}\label{riemann-ux5ea6ux91cf}

\textbf{定义 65.1}（Riemann 度量）：光滑流形 \(M\) 上的一个
\textbf{Riemann 度量} \(g\) 是一个光滑的 \((0, 2)\)-张量场，使得对每个
\(p \in M\)，\(g_p: T_p M \times T_p M \to \mathbb{R}\)
是对称正定的双线性型。即 \(g_p\) 是 \(T_p M\) 上的内积。

在局部坐标下，\(g = \sum_{i,j} g_{ij} dx^i \otimes dx^j\)，其中
\((g_{ij})\) 是对称正定矩阵，\(g_{ij} = g(\partial_i, \partial_j)\)。

\textbf{定义 65.2}（Riemann 流形）：一个 \textbf{Riemann 流形}是一对
\((M, g)\)，其中 \(M\) 是光滑流形，\(g\) 是 \(M\) 上的 Riemann 度量。

\textbf{例 65.1}（Riemann 流形举例）： - \(\mathbb{R}^n\)
配以标准欧氏度量 \(g_{ij} = \delta_{ij}\)。 - \(S^n\)（球面）配以
\(\mathbb{R}^{n+1}\) 诱导的度量。 - Poincaré 上半平面
\(\mathbb{H}^2 = \{(x, y) : y > 0\}\) 配以
\(g = \frac{dx^2 + dy^2}{y^2}\)（双曲度量）。 - 任何光滑流形上存在
Riemann 度量（利用单位分解构造）。

\textbf{定义 65.3}（等距）：两个 Riemann 流形 \((M, g)\) 和 \((N, h)\)
之间的微分同胚 \(F: M \to N\) 称为\textbf{等距}，如果 \(F^*h = g\)，即
\(h_{F(p)}(dF_p(v), dF_p(w)) = g_p(v, w)\) 对所有 \(v, w \in T_p M\)。

\subsubsection{65.2 Levi-Civita 联络}\label{levi-civita-ux8054ux7edc}

\textbf{定义 65.4}（仿射联络）：\(M\)
上的一个\textbf{仿射联络}（或协变导数）是一个映射
\(\nabla: \mathfrak{X}(M) \times \mathfrak{X}(M) \to \mathfrak{X}(M)\)，记作
\(\nabla_X Y\)，满足：

\begin{enumerate}
\def\labelenumi{\arabic{enumi}.}
\tightlist
\item
  \(\nabla_{fX + gY} Z = f \nabla_X Z + g \nabla_Y Z\)（\(C^\infty(M)\)-线性于第一变量）
\item
  \(\nabla_X (aY + bZ) = a \nabla_X Y + b \nabla_X Z\)（\(\mathbb{R}\)-线性于第二变量）
\item
  \(\nabla_X (fY) = (Xf)Y + f \nabla_X Y\)（Leibniz 法则）
\end{enumerate}

\textbf{定理 65.1}（Levi-Civita 联络的存在唯一性）：在 Riemann 流形
\((M, g)\) 上，存在唯一的仿射联络 \(\nabla\) 满足：

\begin{enumerate}
\def\labelenumi{\arabic{enumi}.}
\tightlist
\item
  \textbf{无挠性}（对称性）：\(\nabla_X Y - \nabla_Y X = [X, Y]\)。
\item
  \textbf{与度量相容}：\(X(g(Y, Z)) = g(\nabla_X Y, Z) + g(Y, \nabla_X Z)\)。
\end{enumerate}

这个联络称为 \textbf{Levi-Civita 联络}（或 Riemann 联络）。

\emph{证明}（存在性）：由 Koszul 公式：

\[\begin{aligned}
2g(\nabla_X Y, Z) &= X(g(Y, Z)) + Y(g(Z, X)) - Z(g(X, Y)) \\
&\quad - g(X, [Y, Z]) + g(Y, [Z, X]) + g(Z, [X, Y])
\end{aligned}\]

此公式唯一确定了 \(\nabla_X Y\)。可直接验证由此定义的 \(\nabla\)
满足所有条件。∎

\textbf{定义 65.5}（Christoffel 符号）：在局部坐标下，Levi-Civita 联络由
\textbf{Christoffel 符号} \(\Gamma_{ij}^k\) 决定：

\[\nabla_{\partial_i} \partial_j = \sum_k \Gamma_{ij}^k \partial_k\]

\[\Gamma_{ij}^k = \frac{1}{2} \sum_{\ell} g^{k\ell}(\partial_i g_{j\ell} + \partial_j g_{i\ell} - \partial_\ell g_{ij})\]

\subsubsection{65.3 曲率}\label{ux66f2ux7387}

\textbf{定义 65.6}（Riemann 曲率张量）：Levi-Civita 联络 \(\nabla\) 的
\textbf{Riemann 曲率张量}
\(R: \mathfrak{X}(M) \times \mathfrak{X}(M) \times \mathfrak{X}(M) \to \mathfrak{X}(M)\)
定义为

\[R(X, Y)Z = \nabla_X \nabla_Y Z - \nabla_Y \nabla_X Z - \nabla_{[X, Y]} Z\]

\textbf{命题 65.2}（Riemann 曲率张量的对称性）：

\begin{enumerate}
\def\labelenumi{\arabic{enumi}.}
\tightlist
\item
  \(R(X, Y)Z = -R(Y, X)Z\)
\item
  \(g(R(X, Y)Z, W) = -g(R(X, Y)W, Z)\)
\item
  \(R(X, Y)Z + R(Y, Z)X + R(Z, X)Y = 0\)（第一 Bianchi 恒等式）
\item
  \(g(R(X, Y)Z, W) = g(R(Z, W)X, Y)\)（对偶对称性）
\end{enumerate}

\textbf{定义 65.7}（Riemann 曲率张量的分量）：在局部坐标下，

\[R(\partial_i, \partial_j)\partial_k = \sum_{\ell} R_{ijk}^{\ell} \partial_{\ell}\]

\[R_{ijk\ell} = g(R(\partial_i, \partial_j)\partial_k, \partial_\ell) = \sum_m g_{\ell m} R_{ijk}^m\]

\[R_{ijk\ell} = \frac{1}{2}\left(\frac{\partial^2 g_{i\ell}}{\partial x^j \partial x^k} + \frac{\partial^2 g_{jk}}{\partial x^i \partial x^\ell} - \frac{\partial^2 g_{ik}}{\partial x^j \partial x^\ell} - \frac{\partial^2 g_{j\ell}}{\partial x^i \partial x^k}\right) + \sum_{p,q} g_{pq}(\Gamma_{i\ell}^p \Gamma_{jk}^q - \Gamma_{ik}^p \Gamma_{j\ell}^q)\]

\textbf{定义 65.8}（截面曲率）：设 \(\Pi \subseteq T_p M\)
是二维子空间，由规范正交向量 \(v, w\) 张成。\(\Pi\)
的\textbf{截面曲率}为

\[K(\Pi) = \frac{g(R(v, w)w, v)}{g(v, v)g(w, w) - g(v, w)^2}\]

\textbf{定义 65.9}（Ricci 曲率与标量曲率）： - \textbf{Ricci
曲率}：\(\operatorname{Ric}(X, Y) = \operatorname{tr}(Z \mapsto R(Z, X)Y)\)。在局部坐标下，\(\operatorname{Ric}_{ij} = \sum_k R_{kij}^k\)。
-
\textbf{标量曲率}：\(S = \operatorname{tr}_g(\operatorname{Ric}) = \sum_{i,j} g^{ij} \operatorname{Ric}_{ij}\)。

\textbf{命题 65.3}（第二 Bianchi 恒等式）：

\[\nabla_X R(Y, Z) + \nabla_Y R(Z, X) + \nabla_Z R(X, Y) = 0\]

在局部坐标下：\(\nabla_m R_{ijk\ell} + \nabla_i R_{jmk\ell} + \nabla_j R_{mik\ell} = 0\)。

\subsubsection{65.4 测地线}\label{ux6d4bux5730ux7ebf}

\textbf{定义 65.10}（测地线）：Riemann 流形 \((M, g)\) 上的光滑曲线
\(\gamma: I \to M\) 称为\textbf{测地线}，如果

\[\nabla_{\gamma'} \gamma' = 0\]

在局部坐标下，这等价于测地线方程：

\[\frac{d^2 x^k}{dt^2} + \sum_{i,j} \Gamma_{ij}^k \frac{dx^i}{dt}\frac{dx^j}{dt} = 0\]

\textbf{定理 65.4}（测地线的存在唯一性）：对每个 \(p \in M\) 和
\(v \in T_p M\)，存在唯一的测地线
\(\gamma_v: (-\varepsilon, \varepsilon) \to M\) 满足
\(\gamma_v(0) = p\)，\(\gamma_v'(0) = v\)。

\emph{证明}：测地线方程是二阶常微分方程，由 ODE
解的存在唯一性定理可得。∎

\textbf{定义 65.11}（指数映射）：\(p\) 处的\textbf{指数映射}
\(\exp_p: T_p M \to M\) 定义为
\(\exp_p(v) = \gamma_v(1)\)（当右边有定义时）。\(\exp_p\) 在
\(0 \in T_p M\) 附近是微分同胚（因为
\(d(\exp_p)_0 = \operatorname{id}_{T_p M}\)）。

\textbf{定义 65.12}（完备性）：Riemann 流形 \((M, g)\)
称为\textbf{测地完备的}，如果每条测地线可以无限延伸（即定义域为
\(\mathbb{R}\)）。

\textbf{定理 65.5}（Hopf-Rinow 定理）：对 Riemann 流形
\((M, g)\)，以下条件等价： 1. \((M, g)\) 是测地完备的。 2. \((M, g)\)
作为度量空间（由 \(g\) 诱导的 Riemann 距离）是完备的。 3. \(M\)
中的有界闭集是紧致的。

此时，\(M\) 中任意两点之间都存在最短测地线。

\bigskip\hrule\bigskip

\subsection{第66章：曲率与拓扑}\label{ux7b2c66ux7ae0ux66f2ux7387ux4e0eux62d3ux6251}

本章揭示曲率（局部几何量）与拓扑（全局拓扑不变量）之间的深刻联系。

\subsubsection{66.1 二维情形：Gauss-Bonnet
定理}\label{ux4e8cux7ef4ux60c5ux5f62gauss-bonnet-ux5b9aux7406}

\textbf{定理 66.1}（局部 Gauss-Bonnet 定理）：设 \(S\) 是
\(\mathbb{R}^3\) 中的正则曲面，\(R \subseteq S\) 是由分段光滑简单闭曲线
\(\partial R\) 围成的区域。则

\[\int_R K dA + \int_{\partial R} k_g ds + \sum_i \theta_i = 2\pi\]

其中 \(K\) 是 Gauss 曲率，\(k_g\) 是边界曲线的测地曲率，\(\theta_i\)
是边界角点处的外角。

\emph{证明思路}：在 \(R\) 上选取正交标架场
\(\{e_1, e_2\}\)，定义联络形式
\(\omega_{12} = \langle \nabla e_1, e_2 \rangle\)。则由曲率的定义，\(d\omega_{12} = -K dA\)。再由
Stokes 定理：

\[\int_R K dA = -\int_R d\omega_{12} = -\int_{\partial R} \omega_{12}\]

在边界上 \(\omega_{12}\) 与 \(k_g\) 的关系给出
\(\int_{\partial R} \omega_{12} = 2\pi - \int_{\partial R} k_g ds - \sum \theta_i\)。∎

\textbf{定理 66.2}（全局 Gauss-Bonnet 定理）：设 \(S\)
是紧致无边可定向曲面。则

\[\int_S K dA = 2\pi \chi(S)\]

其中 \(\chi(S)\) 是 \(S\) 的 Euler 示性数。对亏格为 \(g\)
的可定向曲面，\(\chi(S) = 2 - 2g\)。

\emph{证明}：将 \(S\) 三角剖分，对每个三角形应用局部 Gauss-Bonnet
定理，求和。所有内角之和为 \(2\pi\)
乘以顶点数，测地曲率积分在内部边上抵消，外角之和给出
\(\int K dA = 2\pi(V - E + F) = 2\pi \chi(S)\)。∎

\textbf{推论 66.3}（Gauss-Bonnet 定理的推论）： 1. \(\int_S K dA\)
是拓扑不变量（与 Riemann 度量的选择无关）。 2. 若 \(S\) 同胚于
\(S^2\)（球面），则 \(\int_S K dA = 4\pi\)。特别地，\(S^2\) 上的任何
Riemann 度量在某处必有正曲率。 3. 若 \(S\) 同胚于 \(T^2\)（环面），则
\(\int_S K dA = 0\)。因此环面上不能处处有正 Gauss 曲率。 4. 若 \(S\) 的
Gauss 曲率 \(K > 0\) 处处成立，则 \(\chi(S) > 0\)，从而 \(S\) 必同胚于
\(S^2\) 或 \(\mathbb{R}P^2\)。

\subsubsection{66.2 高维推广}\label{ux9ad8ux7ef4ux63a8ux5e7f}

\textbf{定理 66.4}（Chern-Gauss-Bonnet 定理，陈述）：设 \((M, g)\) 是
\(2n\) 维紧致无边可定向 Riemann 流形。则

\[\int_M \operatorname{Pf}(\Omega) = (2\pi)^n \chi(M)\]

其中 \(\operatorname{Pf}(\Omega)\) 是曲率形式 \(\Omega\) 的
Pfaffian，\(\chi(M)\) 是 \(M\) 的 Euler 示性数。

\emph{说明}：当 \(n = 1\)
时，\(\operatorname{Pf}(\Omega) = K dA\)，退化为经典的 Gauss-Bonnet
定理。此定理是陈省身（Shiing-Shen Chern）于 1944
年证明的，是微分几何的里程碑式成果。

\textbf{定理 66.5}（奇维数情形）：若 \(M\)
是紧致无边可定向的奇维数流形，则 \(\chi(M) = 0\)。

\subsubsection{66.3
曲率与拓扑的其他联系}\label{ux66f2ux7387ux4e0eux62d3ux6251ux7684ux5176ux4ed6ux8054ux7cfb}

\textbf{定理 66.6}（Bonnet-Myers 定理）：设 \((M, g)\) 是完备 Riemann
流形，\(\dim M = n\)。若存在常数 \(c > 0\) 使得
\(\operatorname{Ric}(v, v) \geq (n-1)c \cdot g(v, v)\) 对所有 \(v\)
成立，则 \(M\) 是紧致的，且直径
\(\operatorname{diam}(M) \leq \pi / \sqrt{c}\)。特别地，\(\pi_1(M)\)
是有限群。

\emph{证明思路}：设 \(\gamma\) 是长度大于 \(\pi/\sqrt{c}\)
的测地线。利用第二变分公式和 Ricci 曲率条件，可构造连接 \(\gamma\)
两端点的更短曲线，矛盾。因此所有测地线长度不超过 \(\pi/\sqrt{c}\)。由
Hopf-Rinow 定理，\(M\) 紧致且直径有界。万有覆叠 \(\tilde{M}\)
也具有相同的 Ricci 曲率下界，故也紧致，从而 \(\pi_1(M)\) 有限。∎

\textbf{定理 66.7}（Hadamard-Cartan 定理）：设 \((M, g)\) 是完备单连通
Riemann 流形，且截面曲率 \(K \leq 0\) 处处成立。则对每个
\(p \in M\)，指数映射 \(\exp_p: T_p M \to M\) 是微分同胚。特别地，\(M\)
微分同胚于 \(\mathbb{R}^n\)。

\emph{证明思路}：截面曲率 \(\leq 0\) 蕴含 Jacobi 场的长度不减，从而
\(\exp_p\) 是局部微分同胚。由完备性，\(\exp_p\) 是满射。单连通性保证
\(\exp_p\) 是整体单射。∎

\textbf{定理 66.8}（球面定理 / 1/4-夹挤定理）：设 \((M, g)\)
是完备单连通 Riemann 流形，\(\dim M = n\)。若截面曲率 \(K\) 满足
\(1/4 < K \leq 1\)，则 \(M\) 同胚于 \(S^n\)。若
\(1/4 \leq K \leq 1\)，则 \(M\) 同胚于 \(S^n\)
或微分同胚于某个对称空间。

\emph{说明}：这是微分几何中著名的「球面定理」，由 Rauch、Berger 和
Klingenberg 在 1950-60 年代建立。其完整证明需要深入的比较几何和 Morse
理论。

\textbf{定理 66.9}（Chern 的 Gauss-Bonnet
公式推广的积分不变量）：在偶维数紧致 Riemann 流形上，Euler
示性数可以表示为曲率形式的多项式在整个流形上的积分，这一积分值不依赖于度量的选择。

\subsubsection{66.4
曲率张量的进一步分解}\label{ux66f2ux7387ux5f20ux91cfux7684ux8fdbux4e00ux6b65ux5206ux89e3}

\textbf{定义 66.10}（Weyl 曲率张量）：在 \(n \geq 3\) 维 Riemann
流形上，\textbf{Weyl 曲率张量} \(W\) 定义为 Riemann 曲率张量去掉 Ricci
部分和标量曲率部分：

\[W_{ijk\ell} = R_{ijk\ell} - \frac{1}{n-2}(g_{ik}\operatorname{Ric}_{j\ell} - g_{i\ell}\operatorname{Ric}_{jk} - g_{jk}\operatorname{Ric}_{i\ell} + g_{j\ell}\operatorname{Ric}_{ik}) + \frac{S}{(n-1)(n-2)}(g_{ik}g_{j\ell} - g_{i\ell}g_{jk})\]

\(W\) 具有与 \(R\)
相同的对称性，且所有迹为零（\(\sum_i W_{ijk}^i = 0\)）。\(W\)
是\textbf{共形不变}的：若 \(\tilde{g} = e^{2f} g\)，则
\(\tilde{W} = e^{2f} W\)（在指标位置适当调整后）。

\textbf{定理 66.11}（Weyl 张量等于零）：\(n \geq 4\) 维 Riemann
流形局部共形平坦（即局部等距于 \(\mathbb{R}^n\)
配以某个共形度量）当且仅当 \(W = 0\)。\(n = 3\) 时，局部共形平坦等价于
Cotton 张量等于零。

\bigskip\hrule\bigskip

\subsection{第67章：G-结构}\label{ux7b2c67ux7ae0g-ux7ed3ux6784}

G-结构是微分流形切丛的结构群从一般线性群
\(\operatorname{GL}(n,\mathbb{R})\) 约化到 Lie 子群 \(G\)
的几何框架。形式化地，\(n\) 维流形 \(M\) 上的 \(G\)-结构是一个
\(G\)-主丛 \(P \to M\) 连同 \(G\)-等变同构
\(P \times_G \operatorname{GL}(n,\mathbb{R}) \cong FM\)（\(FM\)
为切标架丛），即切丛的每个标架均可通过 \(G\)
的作用彼此转换。\textbf{可积性条件}由结构张量（挠率）的消去判定：若存在无挠联络与
\(G\)-结构相容，则 \(G\)-结构局部平坦。经典例子：（1）\textbf{Riemann
结构}------\(G = \operatorname{O}(n)\)，等价于 Riemann
度量的存在（正交标架丛）；（2）\textbf{复结构}------\(G = \operatorname{GL}(n,\mathbb{C}) \subset \operatorname{GL}(2n,\mathbb{R})\)，对应近复结构，可积性即
Newlander-Nirenberg 定理------近复结构来自复流形当且仅当其 Nijenhuis
张量消去；（3）\textbf{辛结构}------\(G = \operatorname{Sp}(2n,\mathbb{R})\)，对应非退化闭
2-形式，由 Darboux 定理总是局部可积（无局部不变量）。

\bigskip\hrule\bigskip

\subsubsection{67.4
仿射微分几何}\label{ux4effux5c04ux5faeux5206ux51e0ux4f55}

仿射微分几何研究等体积仿射变换群（\(\operatorname{SL}(n,\mathbb{R}) \ltimes \mathbb{R}^n\)）下的曲面和超曲面不变量，是介于
Riemann 几何与射影几何之间的基本几何结构。核心研究对象是
\(\mathbb{R}^{n+1}\) 中非退化超曲面（即其 Blaschke
度量非退化）的仿射法向量、仿射形状算子和仿射度量。\textbf{基本定理}：（1）仿射基本定理------给定
Blaschke 度量和仿射第三基本形式满足 Gauss-Codazzi
型方程，则可唯一确定超曲面（至多相差仿射变换）；（2）仿射球面分类------仿射主曲率为常数的非退化仿射超曲面分为椭球型、双曲型和抛物型；（3）Calabi
猜想（仿射版本）------完备仿射球面必为椭球面或双曲面，由 Calabi (1982)
提出、Trudinger-Wang (2000)
完全解决。仿射微分几何在信息几何（统计流形）和计算机视觉（仿射不变物体识别）中有重要应用。

\subsubsection{67.5 Twistor 理论与超 Kähler
几何}\label{twistor-ux7406ux8bbaux4e0eux8d85-kuxe4hler-ux51e0ux4f55}

Twistor 理论由 Penrose (1967) 创立，通过构造
\(\mathbb{R}^4\)（共形紧化为 \(S^4\)）上的 \textbf{twistor
空间}------\(\mathbb{CP}^3\)（即复三维射影空间）------将四维 Riemann
几何问题转化为复几何问题。基本构造：\(S^4\) 上的每个点对应 twistor
空间中的一条全纯 \(\mathbb{CP}^1\)（twistor
线），反自偶联络对应于全纯向量丛。\textbf{非线性引力子构造}（Penrose,
1976）：反自偶 Einstein 度量一一对应于 twistor
空间上具有全纯切触结构的形变。\textbf{超 Kähler 几何}：超 Kähler
流形是具有三个相容复结构 \(I, J, K\)（满足四元数关系 \(IJ = K\)）的
Riemann 流形。其 twistor 空间是全纯辛流形上的 \(\mathbb{CP}^1\)
全族。关键例子：（1）\(K3\) 曲面------唯一紧致四维超 Kähler
流形；（2）ALE 空间（Kronheimer 分类，对应 Dynkin 图）；（3）Hilbert
概型 \(\operatorname{Hilb}^n(K3)\) 与广义 Kummer 簇------高维紧致超
Kähler 流形的两类基本构造。Atiyah-Hitchin-Singer (1978) 用 twistor
方法研究了磁单极模空间上的超 Kähler 度量。

\bigskip\hrule\bigskip

\emph{本卷从经典曲线论和曲面论出发，建立了 Frenet
标架、第一/第二基本形式、Gauss 曲率等核心概念，证明了 Gauss 绝妙定理和
Gauss-Bonnet 定理。随后引入微分流形和 Riemann
几何的现代框架，包括切空间、微分形式、Levi-Civita
联络、曲率张量等。最后讨论了曲率与拓扑的深层联系，包括 Bonnet-Myers
定理、Hadamard-Cartan 定理和 Chern-Gauss-Bonnet
定理，为后续学习代数拓扑（V14）和李理论（V23）奠定了基础。}

\bigskip\hrule\bigskip

\subsection{第67章：其他几何专题}\label{ux7b2c67ux7ae0ux5176ux4ed6ux51e0ux4f55ux4e13ux9898}

\subsubsection{67.1 Finsler
几何简介}\label{finsler-ux51e0ux4f55ux7b80ux4ecb}

Finsler 几何将 Riemann 几何推广至更一般的度量结构，在切丛
\(TM \setminus \{0\}\) 上赋予 \textbf{Finsler 度量} \(F(x, y)\) 以代替
Riemann 度量给出的二次型。\(F\)
需满足两个条件：（1）\textbf{正齐性}：\(F(x, \lambda y) = \lambda F(x, y)\)
对所有 \(\lambda > 0\) 成立；（2）\textbf{强凸性}：Hessian 矩阵
\(g_{ij}(x, y) = \frac{1}{2}\frac{\partial^2 F^2}{\partial y^i \partial y^j}\)
对每个 \((x, y)\) 正定。区别于 Riemann 度量的关键特征在于 \(F(x, y)\)
一般不满足平行四边形法则------弧长不仅依赖于位置，还依赖于方向，且不具有二次型结构。这使得
Finsler
几何中切空间的''单位球''可以是一般的光滑凸体（非椭球）。\textbf{Chern
联络}（S. S. Chern, 1948）是 Finsler
流形上的典范联络，由非线性联络（刻画切丛的水平分布）和线性联络（沿水平方向的协变导数）两部分构成，在无挠条件下保持
Finsler 度量。Chern 联络是整个 Finsler
几何曲率理论的框架基础，在现代变分学与生物形态学中有重要应用。

\subsubsection{67.2 谱几何（Spectral
Geometry）}\label{ux8c31ux51e0ux4f55spectral-geometry}

谱几何研究紧致 Riemann 流形 \((M, g)\) 上 \textbf{Laplace-Beltrami 算子}
\(\Delta_g\) 的谱与几何性质之间的对应关系。\(\Delta_g\) 的特征值满足
\(0 = \lambda_0 < \lambda_1 \leq \lambda_2 \leq \cdots \nearrow \infty\)。经典问题来自
Mark Kac (1966)：``\textbf{能听出鼓的形状吗？}''（Can one hear the shape
of a drum?）------若两个平面区域具有相同的 Dirichlet
特征值谱，它们是否必然等距？数学上的 \textbf{Weyl 渐近公式}：

\[N(\lambda) = \#\{\lambda_k \leq \lambda\} \sim \frac{\omega_n}{(2\pi)^n} \mathrm{Vol}(M) \, \lambda^{n/2}, \quad \lambda \to \infty\]

表明谱决定了流形的维数和体积；进一步，热核渐近展开中的各阶系数依次给出总标量曲率的积分等几何不变量。1992
年，Gordon、Webb 和 Wolpert 利用数论方法构造了平面上具有相同 Dirichlet
谱但不等距（也不等度规）的区域对，从而对 Kac
问题给出否定回答。谱几何在量子力学（Schrodinger
算子谱）、热核分析与数论等领域有深远应用。

\subsubsection{67.3 积分几何（Integral
Geometry）}\label{ux79efux5206ux51e0ux4f55integral-geometry}

积分几何研究几何量在群作用下的积分平均，由 Blaschke 学派于 20
世纪初创立。核心公式------\textbf{Crofton 公式}将平面可求长曲线
\(\Gamma\) 的长度 \(L\) 表示为

\[L = \frac{1}{2} \int n_\ell(\Gamma) \, d\ell\]

其中 \(n_\ell(\Gamma)\) 为直线 \(\ell\) 与 \(\Gamma\)
的交点个数，积分对平面上的全体直线（赋以适当的运动群不变测度）进行。等价地，曲线长度等于其与随机直线相交的期望数乘以
\(\pi/2\)。\textbf{Radon 变换}将 \(\mathbb{R}^n\) 上的函数 \(f\)
映射为其在超平面上的积分值：\((\mathcal{R}f)(H) = \int_H f \, d\sigma\)。Johann
Radon 于 1917 年建立了 Radon 变换的反演公式，这一原理正是现代 CT
扫描（计算机断层成像）的数学基础。积分几何将度量几何性质（长度、面积、曲率）转化为群不变测度下的积分平均值，深刻体现了几何与分析、概率论的对偶关系，在凸几何、随机几何和医学成像中均有重要应用。\(\blacksquare\)

\newpage

\section{卷十二：抽象代数
II}\label{ux5377ux5341ux4e8cux62bdux8c61ux4ee3ux6570-ii}

\begin{quote}
抽象代数 II 在 V8（群、环、域、模的基础）之上，深入 Galois
理论、交换代数、表示论、同调代数与代数数论五大核心领域，构建从古典代数到现代代数的完整桥梁。
\end{quote}

\bigskip\hrule\bigskip

\subsection{第67章：Galois
理论}\label{ux7b2c67ux7ae0galois-ux7406ux8bba}

Galois 理论建立了域扩张的中间域与其 Galois
群的子群之间的一一对应，是代数方程理论的巅峰。

\subsubsection{67.1 域扩张回顾}\label{ux57dfux6269ux5f20ux56deux987e}

回顾 V8 第 44 章的基本概念：设 \(E/F\) 是域扩张。

\begin{itemize}
\tightlist
\item
  \textbf{代数元}：\(\alpha \in E\) 称为 \(F\)
  上的代数元，若存在非零多项式 \(f \in F[x]\) 使 \(f(\alpha) = 0\)。
\item
  \textbf{最小多项式}：\(\alpha\) 在 \(F\) 上的次数最低的首一多项式
  \(m_{\alpha,F}(x)\)。
\item
  \textbf{单扩张}：\(F(\alpha) \cong F[x]/(m_{\alpha,F}(x))\)。
\item
  \textbf{扩张次数}：\([E : F] = \dim_F E\)。
\end{itemize}

\textbf{定义 67.1}（Galois 群）：设 \(E/F\) 是域扩张。\(E/F\) 的
\textbf{Galois 群} \(\operatorname{Gal}(E/F)\) 是 \(E\) 的所有保持 \(F\)
逐点不动的自同构构成的群：

\[\operatorname{Gal}(E/F) = \{\sigma \in \operatorname{Aut}(E) : \sigma(a) = a, \; \forall a \in F\}\]

\textbf{定义 67.2}（正规扩张）：代数扩张 \(E/F\)
称为\textbf{正规的}，如果 \(E\) 是某个 \(F[x]\)
中多项式的分裂域。等价地，任何不可约多项式 \(f \in F[x]\) 若在 \(E\)
中有一个根，则 \(f\) 在 \(E[x]\) 中完全分裂。

\textbf{定义 67.3}（可分扩张）：代数扩张 \(E/F\)
称为\textbf{可分的}，如果 \(E\) 中每个元素在 \(F\)
上的最小多项式没有重根（即在其分裂域中无重根）。称 \(E/F\) 为
\textbf{Galois 扩张}，如果它既是正规的又是可分的。

\subsubsection{67.2 Galois
对应基本定理}\label{galois-ux5bf9ux5e94ux57faux672cux5b9aux7406}

\textbf{定理 67.1}（Galois 对应基本定理）：设 \(E/F\) 是有限 Galois
扩张，\(G = \operatorname{Gal}(E/F)\)。则存在以下一一对应：

\[\{\text{中间域 } F \subseteq K \subseteq E\} \longleftrightarrow \{\text{子群 } H \leq G\}\]

具体地： -
\(K \mapsto \operatorname{Gal}(E/K) = \{\sigma \in G : \sigma|_K = \operatorname{id}_K\}\)
-
\(H \mapsto E^H = \{x \in E : \sigma(x) = x, \; \forall \sigma \in H\}\)（\(H\)
的固定域）

且满足：

\begin{enumerate}
\def\labelenumi{\arabic{enumi}.}
\tightlist
\item
  \([E : K] = |\operatorname{Gal}(E/K)|\)，\([K : F] = [G : \operatorname{Gal}(E/K)]\)。
\item
  \(K/F\) 是正规扩张（从而也是 Galois 扩张）当且仅当
  \(\operatorname{Gal}(E/K) \trianglelefteq G\)。此时
  \(\operatorname{Gal}(K/F) \cong G / \operatorname{Gal}(E/K)\)。
\item
  对应是反序的：\(K_1 \subseteq K_2 \Leftrightarrow \operatorname{Gal}(E/K_2) \subseteq \operatorname{Gal}(E/K_1)\)。
\end{enumerate}

\emph{证明思路}：核心步骤：

\begin{enumerate}
\def\labelenumi{(\arabic{enumi})}
\item
  任意有限可分扩张是单扩张（本原元素定理）：存在 \(\alpha \in E\) 使
  \(E = F(\alpha)\)。则 \(|G| = [E : F]\)（因为自同构由 \(\alpha\)
  的像决定，而 \(\alpha\) 可映到其最小多项式的任意根，Galois
  扩张中所有根都在 \(E\) 中且互异）。
\item
  对任意子群
  \(H \leq G\)，\([E : E^H] = |H|\)。这由线性无关性引理（Artin
  引理）保证：\(E\) 中元素关于 \(H\) 中自同构的线性无关性。
\end{enumerate}

\textbf{引理 67.1}（Artin 引理）：设 \(E\)
为域，\(H = \{\sigma_1, \ldots, \sigma_n\}\) 是 \(E\)
的自同构构成的有限群。若 \(\alpha_1, \ldots, \alpha_n \in E\) 满足
\(\sum_{i=1}^n \alpha_i \sigma_i(x) = 0\) 对所有 \(x \in E\) 成立，则
\(\alpha_1 = \cdots = \alpha_n = 0\)。换言之，\(H\) 中自同构作为 \(E\)
到自身的函数是 \(E\)-线性无关的。

\emph{证明}：假设存在非平凡关系。取项数最少的关系
\(\sum_{i=1}^m \alpha_i \sigma_i(x) = 0\)（\(\alpha_i \neq 0\)）。由于
\(\sigma_1 \neq \sigma_m\)，存在 \(y \in E\) 使
\(\sigma_1(y) \neq \sigma_m(y)\)。将关系应用于 \(x\) 和 \(yx\)，消去
\(\sigma_1\) 的系数，得到项数更少的对立关系，矛盾。∎

\begin{enumerate}
\def\labelenumi{(\arabic{enumi})}
\setcounter{enumi}{2}
\item
  由 (1) 和 (2) 得 \(E^{\operatorname{Gal}(E/K)} = K\) 和
  \(\operatorname{Gal}(E/E^H) = H\)，从而建立一一对应。
\item
  正规性条件：\(K/F\) 正规等价于 \(\sigma(K) = K\) 对所有
  \(\sigma \in G\) 成立，等价于
  \(\operatorname{Gal}(E/K) \trianglelefteq G\)。∎
\end{enumerate}

\textbf{推论 67.2}：\([E : F] = |\operatorname{Gal}(E/F)|\)。即 Galois
扩张的次数等于其 Galois 群的阶。

\subsubsection{67.3 方程的 Galois
群}\label{ux65b9ux7a0bux7684-galois-ux7fa4}

\textbf{定义 67.4}（多项式的 Galois 群）：设 \(f(x) \in F[x]\)
为可分多项式（无重根），\(E\) 是 \(f\) 在 \(F\) 上的分裂域。则 \(f\) 在
\(F\) 上的 \textbf{Galois 群}定义为
\(\operatorname{Gal}(f/F) = \operatorname{Gal}(E/F)\)。

由于 \(E/F\) 是 Galois 扩张，\(\operatorname{Gal}(f/F)\) 可视为 \(f\)
的根上的置换群，因此
\(\operatorname{Gal}(f/F) \leq S_n\)（\(n = \deg f\)）。

\textbf{定理 67.3}（Galois 可解性判别法）：多项式 \(f(x) \in F[x]\)
可用根式求解当且仅当 \(\operatorname{Gal}(f/F)\) 是\textbf{可解群}。

\emph{证明思路}：(\(\Rightarrow\)) 若 \(f\)
可用根式求解，则存在根式扩张塔
\(F = F_0 \subseteq F_1 \subseteq \cdots \subseteq F_m\) 包含 \(f\)
的分裂域。每个 \(F_{i+1}/F_i\) 是添加 \(p\) 次单位根或 \(p\)
次方根，对应的 Galois 群是 Abel 群。由 Galois 对应，这给出
\(\operatorname{Gal}(f/F)\) 的 subnormal 列，其商群都是 Abel
群，故可解。

(\(\Leftarrow\)) 若 \(\operatorname{Gal}(f/F)\) 可解，则存在 subnormal
列，其商群为素数阶循环群。由 Galois 对应，这对应于一列根式扩张，从而
\(f\) 可用根式求解。∎

\textbf{定理 67.4}（Abel-Ruffini 定理）：一般五次方程不可用根式求解。

\emph{证明}：一般五次方程
\(f(x) = x^5 - s_1 x^4 + s_2 x^3 - s_3 x^2 + s_4 x - s_5\)（\(s_i\)
为独立未定元）的 Galois 群是 \(S_5\)。\(S_5\) 不是可解群（因为 \(A_5\)
是单群且非 Abel），故由定理 67.3 知不可用根式求解。∎

\subsubsection{67.4
有限域与分圆域}\label{ux6709ux9650ux57dfux4e0eux5206ux5706ux57df}

\textbf{定理 67.5}（有限域的 Galois 群）：设
\(\mathbb{F}_{p^n}/\mathbb{F}_p\) 是 \(p^n\) 元域对 \(p\)
元素域的扩张。则
\(\operatorname{Gal}(\mathbb{F}_{p^n}/\mathbb{F}_p) \cong \mathbb{Z}/n\mathbb{Z}\)，由
Frobenius 自同构 \(\varphi(x) = x^p\) 生成。

\emph{证明}：\(\mathbb{F}_{p^n}\) 是 \(x^{p^n} - x\) 在 \(\mathbb{F}_p\)
上的分裂域，故为 Galois 扩张。Frobenius 映射 \(\varphi\) 是
\(\mathbb{F}_{p^n}\) 的自同构，且
\(\varphi^n = \operatorname{id}\)。\(\varphi\) 的阶恰为 \(n\)（否则某
\(\varphi^k = \operatorname{id}\) 蕴含 \(x^{p^k} = x\) 对所有 \(x\)
成立，从而 \(|\mathbb{F}_{p^n}| \leq p^k\)，矛盾）。且
\(|\operatorname{Gal}| = [\mathbb{F}_{p^n} : \mathbb{F}_p] = n\)。∎

\textbf{定理 67.6}（分圆域的 Galois 群）：设 \(\zeta_n\) 是 \(n\)
次本原单位根，\(K = \mathbb{Q}(\zeta_n)\)。则 \(K/\mathbb{Q}\) 是 Galois
扩张，且

\[\operatorname{Gal}(K/\mathbb{Q}) \cong (\mathbb{Z}/n\mathbb{Z})^\times\]

同构由
\(\sigma_a(\zeta_n) = \zeta_n^a\)（\(\gcd(a, n) = 1\)）给出。特别地，\([K : \mathbb{Q}] = \varphi(n)\)（Euler
函数）。

\emph{证明}：\(\zeta_n\) 的最小多项式是分圆多项式 \(\Phi_n(x)\)（次数
\(\varphi(n)\)）。\(K\) 是 \(\Phi_n(x)\) 的分裂域。\(\Phi_n(x)\) 的根为
\(\zeta_n^a\)（\(\gcd(a, n) = 1\)），故自同构由
\(\zeta_n \mapsto \zeta_n^a\) 确定，且
\(\sigma_a \circ \sigma_b = \sigma_{ab \bmod n}\)。∎

\subsubsection{67.5
尺规作图的代数刻画}\label{ux5c3aux89c4ux4f5cux56feux7684ux4ee3ux6570ux523bux753b}

\textbf{定理 67.7}（可构造数的代数刻画）：实数 \(\alpha\)
可用尺规作出当且仅当存在 \(\mathbb{Q}\) 的二次扩张塔

\[\mathbb{Q} = K_0 \subseteq K_1 \subseteq \cdots \subseteq K_n\]

使得 \(\alpha \in K_n\) 且
\([K_i : K_{i-1}] = 2\)。特别地，\([\mathbb{Q}(\alpha) : \mathbb{Q}]\)
是 2 的幂。

\textbf{推论 67.8}（三等分角、倍立方、化圆为方不可行）： 1.
一般角不能用尺规三等分（如 \(60^\circ\) 三等分等价于作
\(\cos 20^\circ\)，其最小多项式 \(8x^3 - 6x - 1\) 次数为 3，不是 2
的幂）。 2. 倍立方问题等价于构造
\(\sqrt[3]{2}\)，\([\mathbb{Q}(\sqrt[3]{2}) : \mathbb{Q}] = 3\)，不可尺规作图。
3. 化圆为方等价于构造 \(\sqrt{\pi}\)，\(\pi\) 是超越数（Lindemann
定理），\([\mathbb{Q}(\pi) : \mathbb{Q}]\) 无限，不可尺规作图。

\bigskip\hrule\bigskip

\subsection{第68章：交换代数基础}\label{ux7b2c68ux7ae0ux4ea4ux6362ux4ee3ux6570ux57faux7840}

交换代数是代数几何和代数数论的语言基础，研究交换环的结构。

\subsubsection{68.1 Noether 环与 Artin
环}\label{noether-ux73afux4e0e-artin-ux73af}

\textbf{定义 68.1}（Noether 环）：交换环 \(R\) 称为 \textbf{Noether
环}，如果 \(R\) 中的理想满足\textbf{升链条件}（ACC）：任意理想升链

\[I_1 \subseteq I_2 \subseteq I_3 \subseteq \cdots\]

必稳定，即存在 \(n\) 使得 \(I_n = I_{n+1} = \cdots\)。

\textbf{命题 68.1}（Noether 环的等价刻画）：\(R\) 是 Noether 环当且仅当
\(R\) 的每个理想都是有限生成的，也当且仅当 \(R\)
的每个非空理想族在包含关系下都有极大元。

\emph{证明}：(ACC \(\Rightarrow\) 有限生成)：若 \(I\)
不能有限生成，则归纳构造严格升链，与 ACC 矛盾。(有限生成 \(\Rightarrow\)
ACC)：若 \(I_1 \subseteq I_2 \subseteq \cdots\)，则 \(I = \bigcup I_n\)
也是理想，有限生成 \(I = (a_1, \ldots, a_m)\)。每个 \(a_i\) 属于某个
\(I_{n_i}\)，取 \(n = \max\{n_i\}\)，则 \(I = I_n\)。∎

\textbf{定理 68.2}（Hilbert 基定理）：若 \(R\) 是 Noether 环，则多项式环
\(R[x]\) 也是 Noether 环。

\emph{证明}：设 \(I \subseteq R[x]\) 是理想。对每个 \(n \geq 0\)，令
\(J_n\) 为 \(I\) 中 \(n\) 次多项式的首项系数和 0 构成的 \(R\) 的理想。则
\(J_0 \subseteq J_1 \subseteq J_2 \subseteq \cdots\)。由 \(R\) 的
Noether 性，存在 \(N\) 使 \(J_n = J_N\) 对所有 \(n \geq N\)。

每个
\(J_k\)（\(k \leq N\)）有限生成：\(J_k = (a_{k1}, \ldots, a_{ks_k})\)。取
\(f_{kj} \in I\) 使其首项系数为 \(a_{kj}\)。则这些 \(f_{kj}\) 生成
\(I\)：对任意 \(f \in I\)，用这些生成元逐步消去 \(f\) 的首项，归纳可得
\(f\) 属于它们生成的理想。∎

\textbf{推论 68.3}：若 \(R\) 是 Noether 环，则 \(R[x_1, \ldots, x_n]\)
也是 Noether 环。特别地，\(\mathbb{Z}[x_1, \ldots, x_n]\) 和
\(k[x_1, \ldots, x_n]\)（\(k\) 为域）是 Noether 环。

\textbf{定义 68.2}（Artin 环）：交换环 \(R\) 称为 \textbf{Artin
环}，如果 \(R\) 中的理想满足\textbf{降链条件}（DCC）：任意理想降链
\(I_1 \supseteq I_2 \supseteq \cdots\) 必稳定。

\textbf{定理 68.4}（Akizuki-Hopkins 定理）：\(R\) 是 Artin 环当且仅当
\(R\) 是 Noether 环且 \(\dim R = 0\)（即每个素理想都是极大理想）。

\subsubsection{68.2
素理想与局部化}\label{ux7d20ux7406ux60f3ux4e0eux5c40ux90e8ux5316}

\textbf{定义 68.3}（素理想与极大理想）：设 \(R\) 为交换环。 - 真理想
\(\mathfrak{p} \subsetneq R\) 称为\textbf{素理想}，如果
\(ab \in \mathfrak{p} \Rightarrow a \in \mathfrak{p}\) 或
\(b \in \mathfrak{p}\)。等价地，\(R/\mathfrak{p}\) 是整环。 - 真理想
\(\mathfrak{m} \subsetneq R\) 称为\textbf{极大理想}，如果不存在真理想
\(\mathfrak{m} \subsetneq I \subsetneq R\)。等价地，\(R/\mathfrak{m}\)
是域。

\textbf{定义 68.4}（谱）：\(R\) 的\textbf{素谱}
\(\operatorname{Spec}(R)\) 是所有素理想构成的集合。\(R\)
的\textbf{极大谱} \(\operatorname{MaxSpec}(R)\)
是所有极大理想构成的集合。

\textbf{定理 68.5}（Krull
存在定理）：每个非零交换环都有极大理想（从而也有素理想）。

\emph{证明}：Zorn
引理：所有真理想在包含关系下构成偏序集，任意链的并仍是真理想（因为 \(1\)
不在其中），故有极大元。∎

\textbf{定义 68.5}（局部化）：设 \(R\) 为交换环，\(S \subseteq R\)
为乘法封闭子集（\(1 \in S\)，\(a, b \in S \Rightarrow ab \in S\)）。\(R\)
关于 \(S\) 的\textbf{局部化} \(S^{-1}R\) 定义为等价类集合：

\[S^{-1}R = \left\{\frac{a}{s} : a \in R, s \in S\right\}\]

其中 \(\frac{a}{s} = \frac{a'}{s'}\) 当且仅当存在 \(t \in S\) 使
\(t(as' - a's) = 0\)。

\textbf{重要特例}： - 若
\(S = R \setminus \mathfrak{p}\)（\(\mathfrak{p}\)
为素理想），\(S^{-1}R\) 记作 \(R_{\mathfrak{p}}\)，称为 \(R\) 在
\(\mathfrak{p}\) 处的\textbf{局部化}。 - 若
\(S = \{f^n : n \geq 0\}\)，\(S^{-1}R\) 记作 \(R_f\)。 - 若 \(R\)
是整环，\(S = R \setminus \{0\}\)，则 \(S^{-1}R\) 是 \(R\)
的\textbf{分式域}。

\textbf{命题 68.6}（局部化与素理想）：\(R_{\mathfrak{p}}\)
是局部环，其唯一极大理想为
\(\mathfrak{p}R_{\mathfrak{p}} = \{a/s : a \in \mathfrak{p}, s \notin \mathfrak{p}\}\)。存在自然的单射
\(R/\mathfrak{p} \hookrightarrow R_{\mathfrak{p}}/\mathfrak{p}R_{\mathfrak{p}}\)。

\subsubsection{68.3 整性、整闭包与 Noether
正规化}\label{ux6574ux6027ux6574ux95edux5305ux4e0e-noether-ux6b63ux89c4ux5316}

\textbf{定义 68.6}（整元素）：设 \(R \subseteq A\)
是环的扩张。\(a \in A\) 称为 \(R\)
上的\textbf{整元素}，如果存在首一多项式
\(f(x) = x^n + r_{n-1}x^{n-1} + \cdots + r_0 \in R[x]\) 使
\(f(a) = 0\)。

\(R\) 在 \(A\) 中的\textbf{整闭包}是 \(A\) 中所有 \(R\)
上整元素构成的集合。若 \(R\) 等于其在自己分式域中的整闭包，则称 \(R\)
为\textbf{整闭的}。

\textbf{定理 68.7}（Noether 正规化引理）：设 \(k\) 为域，\(A\)
是有限生成 \(k\)-代数。则存在代数无关的元素
\(y_1, \ldots, y_d \in A\)（\(d\) 是 \(A\) 的 Krull 维数）使得 \(A\) 是
\(k[y_1, \ldots, y_d]\) 上的有限生成模（即整扩张）。

\emph{证明思路}：若 \(A = k[x_1, \ldots, x_n]\) 且 \(x_i\)
代数无关，则已成立。否则，存在非平凡关系
\(f(x_1, \ldots, x_n) = 0\)。通过适当的变量替换
\(y_i = x_i - x_n^{m_i}\)（\(m_i\) 充分大），可将 \(f\) 写成关于 \(x_n\)
的首一多项式。归纳可得。∎

\textbf{推论 68.8}（Hilbert 零点定理的弱形式）：若 \(k\) 是代数闭域，则
\(k[x_1, \ldots, x_n]\) 的每个极大理想形如
\((x_1 - a_1, \ldots, x_n - a_n)\)，对应 \(k^n\) 中的一个点。

\textbf{定理 68.9}（Hilbert 零点定理，强形式）：设 \(k\)
为代数闭域，\(I \subseteq k[x_1, \ldots, x_n]\) 为理想。若
\(f \in k[x_1, \ldots, x_n]\) 在 \(I\) 的零点集 \(V(I)\) 上恒为零，则
\(f^m \in I\) 对某个 \(m \geq 1\) 成立。即 \(I(V(I)) = \sqrt{I}\)（\(I\)
的根）。

\subsubsection{68.4 准素分解与 Krull
交定理}\label{ux51c6ux7d20ux5206ux89e3ux4e0e-krull-ux4ea4ux5b9aux7406}

\textbf{定义 68.7}（准素理想）：真理想 \(\mathfrak{q} \subsetneq R\)
称为\textbf{准素理想}，如果 \(ab \in \mathfrak{q}\) 且
\(a \notin \mathfrak{q}\) 蕴含 \(b^n \in \mathfrak{q}\) 对某个
\(n \geq 1\) 成立。等价地，\(R/\mathfrak{q}\) 中每个零因子都是幂零元。

\textbf{定理 68.10}（Lasker-Noether 准素分解定理）：在 Noether
环中，每个理想都可以表示为有限个准素理想的交。即对任意理想
\(I\)，存在准素理想 \(\mathfrak{q}_1, \ldots, \mathfrak{q}_n\) 使得

\[I = \mathfrak{q}_1 \cap \mathfrak{q}_2 \cap \cdots \cap \mathfrak{q}_n\]

且此分解可以取为「不可缩短的」且素理想 \(\sqrt{\mathfrak{q}_i}\)
互不相同（称为「极小准素分解」）。

\emph{注意}：此定理是整数环中「算术基本定理」在一般 Noether
环中的推广。整数环中 \(n = \prod p_i^{e_i}\) 对应理想
\((n) = \bigcap (p_i^{e_i})\)。

\emph{证明思路}：(1) 在 Noether
环中，每个理想可表示为有限个\textbf{不可约理想}（不能写成两个严格更大理想的交）的交------因为若不然，可构造无穷严格升链，与
Noether 性矛盾。(2) 不可约理想必为准素理想：若 \(\mathfrak{q}\) 不可约且
\(ab \in \mathfrak{q}\)，\(a \notin \mathfrak{q}\)，考虑升链
\((\mathfrak{q}:b) \subseteq (\mathfrak{q}:b^2) \subseteq \cdots\)，由
Noether 性稳定于某 \(n\)，可得 \(b^n \in \mathfrak{q}\)。(3)
唯一性（极小准素分解）：素理想
\(\mathfrak{p}_i = \sqrt{\mathfrak{q}_i}\) 由 \(I\)
唯一决定------它们是使得 \(I \subseteq \mathfrak{p}\) 且
\((I:x) \subseteq \mathfrak{p}\) 的素理想中的极小元（\(I\)
的\textbf{素因子}）。属于极小素因子的准素分支是唯一的，但嵌入素因子的准素分支可能不唯一。

\textbf{例 68.1}（不可约 ≠ 素理想的反例）：在 \(k[x]\) 中，理想
\((x^2)\) 是不可约的（因为若 \((x^2) = I \cap J\)，则 \(I, J\) 形如
\((x^a), (x^b)\)，仅当 \(a=2\) 或 \(b=2\) 时交为
\((x^2)\)，故不可约分解平凡）。但 \((x^2)\)
不是素理想（\(x \cdot x \in (x^2)\) 但 \(x \notin (x^2)\)）。其根
\(\sqrt{(x^2)} = (x)\) 是素理想，故 \((x^2)\) 是准素理想。

\textbf{定理 68.11}（Krull 交定理）：设 \(R\) 是 Noether
环，\(I \subseteq R\) 是真理想。则 \(\bigcap_{n=1}^\infty I^n = 0\) 当
\(R\) 是整环或局部环时成立。一般地，存在 \(a \in I\) 使得
\((1 - a)\bigcap_{n=1}^\infty I^n = 0\)。

\bigskip\hrule\bigskip

\subsection{第69章：表示论基础}\label{ux7b2c69ux7ae0ux8868ux793aux8bbaux57faux7840}

表示论研究群（或代数）通过线性变换的具体实现，是连接抽象代数与线性代数的桥梁。

\subsubsection{69.1 群的表示}\label{ux7fa4ux7684ux8868ux793a}

\textbf{定义 69.1}（群的线性表示）：设 \(G\) 是有限群，\(V\) 是域 \(k\)
上的有限维向量空间。\(G\) 在 \(V\) 上的一个\textbf{线性表示}是一个群同态
\(\rho: G \to \operatorname{GL}(V)\)。

\(V\) 的维数称为表示的\textbf{维数}（或\textbf{次数}）。若 \(\rho\)
是单射，则称表示为\textbf{忠实的}。

\textbf{定义 69.2}（不变子空间与不可约表示）：表示
\(\rho: G \to \operatorname{GL}(V)\) 的子空间 \(W \subseteq V\)
称为\textbf{不变的}（或 \(G\)-子模），如果 \(\rho(g)(W) \subseteq W\)
对所有 \(g \in G\) 成立。表示 \((\rho, V)\) 称为\textbf{不可约的}，如果
\(V \neq 0\) 且唯一的不变子空间是 \(\{0\}\) 和 \(V\)。

\textbf{定义 69.3}（表示的直和）：若
\(V = V_1 \oplus V_2\)（\(\dim V_i > 0\)）且 \(V_1, V_2\) 都是
\(G\)-不变的，则称 \((\rho, V)\) 是 \((\rho|_{V_1}, V_1)\) 和
\((\rho|_{V_2}, V_2)\)
的\textbf{直和}。若表示不能分解为非平凡子表示的直和，则称其为\textbf{不可分解的}。

\textbf{定理 69.1}（Maschke 定理）：设 \(G\) 是有限群，\(k\) 是域且
\(\operatorname{char}(k) \nmid |G|\)。则 \(G\) 在 \(k\)
上的每个有限维表示都是完全可约的（即可分解为不可约表示的直和）。

\emph{证明}：设 \(W \subseteq V\) 是 \(G\)-不变子空间。需构造
\(G\)-不变的补空间 \(W'\)。取任意线性投影 \(\pi: V \to W\)（非
\(G\)-等变）。定义平均投影

\[\tilde{\pi}(v) = \frac{1}{|G|} \sum_{g \in G} \rho(g) \pi(\rho(g^{-1})v)\]

则 \(\tilde{\pi}\) 是
\(G\)-等变的（\(\tilde{\pi}(\rho(h)v) = \rho(h)\tilde{\pi}(v)\)），且
\(\tilde{\pi}|_W = \operatorname{id}_W\)。于是 \(W' = \ker \tilde{\pi}\)
是 \(G\)-不变的补空间。∎

\subsubsection{69.2 特征标理论}\label{ux7279ux5f81ux6807ux7406ux8bba}

以下设 \(k = \mathbb{C}\)（复数域），\(G\) 是有限群。

\textbf{定义 69.4}（特征标）：表示 \(\rho: G \to \operatorname{GL}(V)\)
的\textbf{特征标}是函数
\(\chi_V: G \to \mathbb{C}\)，\(\chi_V(g) = \operatorname{tr}(\rho(g))\)。

特征标是\textbf{类函数}：\(\chi_V(hgh^{-1}) = \chi_V(g)\)。因此
\(\chi_V\) 在共轭类上取常值。

\textbf{命题 69.2}（特征标的基本性质）： 1. \(\chi_V(e) = \dim V\)。 2.
\(\chi_{V \oplus W} = \chi_V + \chi_W\)。 3.
\(\chi_{V \otimes W} = \chi_V \cdot \chi_W\)。 4.
\(\chi_V(g^{-1}) = \overline{\chi_V(g)}\)（对 \(g\) 的阶有限）。

\textbf{定义 69.5}（特征标的内积）：对于类函数
\(\varphi, \psi: G \to \mathbb{C}\)，定义内积

\[\langle \varphi, \psi \rangle = \frac{1}{|G|} \sum_{g \in G} \varphi(g) \overline{\psi(g)}\]

\textbf{定理 69.3}（不可约特征标的正交性）： 1. 若 \(\chi\)
是不可约表示的特征标，则 \(\langle \chi, \chi \rangle = 1\)。 2. 若
\(\chi\) 和 \(\psi\) 是两个不同构的不可约表示的特征标，则
\(\langle \chi, \psi \rangle = 0\)。 3.
不可约特征标构成类函数空间的一组正交规范基。

\emph{证明思路}：利用 Schur 引理：若
\(\rho: G \to \operatorname{GL}(V)\) 和
\(\sigma: G \to \operatorname{GL}(W)\) 是不可约表示，\(T: V \to W\) 是
\(G\)-等变线性映射，则 \(T\) 要么是 0
要么是同构。然后构造并计算适当的等变映射的迹。∎

\textbf{推论 69.4}：表示 \(V\) 不可约当且仅当
\(\langle \chi_V, \chi_V \rangle = 1\)。

\textbf{定理 69.5}：\(G\) 的不可约表示的同构类个数等于 \(G\)
的共轭类个数。且若 \(d_1, \ldots, d_r\) 是不可约表示的维数，则

\[\sum_{i=1}^r d_i^2 = |G|\]

\subsubsection{69.3 诱导表示与 Frobenius
互反律}\label{ux8bf1ux5bfcux8868ux793aux4e0e-frobenius-ux4e92ux53cdux5f8b}

\textbf{定义 69.6}（诱导表示）：设 \(H \leq G\)，\((\rho, W)\) 是 \(H\)
的表示。定义 \(G\) 的\textbf{诱导表示} \(\operatorname{Ind}_H^G W\) 为

\[\operatorname{Ind}_H^G W = \{f: G \to W : f(gh) = \rho(h^{-1})f(g), \; \forall h \in H\}\]

\(G\) 的作用为 \((\tilde{g} \cdot f)(g) = f(\tilde{g}^{-1} g)\)。

\textbf{定理 69.6}（Frobenius 互反律）：设 \(\chi\) 是 \(H\)
的特征标，\(\psi\) 是 \(G\) 的特征标。则

\[\langle \operatorname{Ind}_H^G \chi, \psi \rangle_G = \langle \chi, \operatorname{Res}_H^G \psi \rangle_H\]

其中 \(\operatorname{Res}_H^G \psi\) 是 \(\psi\) 限制为 \(H\) 的特征标。

\subsubsection{69.4 Burnside 定理}\label{burnside-ux5b9aux7406}

\textbf{定理 69.7}（Burnside \(p^a q^b\) 定理）：若 \(G\)
是有限群，\(|G| = p^a q^b\)（\(p, q\) 为素数），则 \(G\) 是可解群。

\emph{证明思路}：用特征标理论。若 \(G\) 是单群，考虑共轭类大小
\(|C_i|\)。由特征标正交性，存在不可约特征标 \(\chi\) 使 \(\chi(1)\)
与某个 \(|C_i|\) 互素。进一步分析表明 \(\chi(1)\) 整除
\(|G|/|C_i|\)。利用代数整数理论可导出矛盾，除非 \(G\)
为素数阶循环群。因此任何 \(p^a q^b\)
阶群必有非平凡正规子群，归纳得可解。∎

\textbf{推论 69.8}：最小非 Abel 单群是 \(A_5\)（阶 60 =
\(2^2 \cdot 3 \cdot 5\)）。所有阶小于 60 的群都是可解群。

\bigskip\hrule\bigskip

\subsection{第70章：同调代数基础}\label{ux7b2c70ux7ae0ux540cux8c03ux4ee3ux6570ux57faux7840}

同调代数用链复形和同调群的语言统一处理代数拓扑、群上同调、层上同调等理论。

\subsubsection{70.1
链复形与同调}\label{ux94feux590dux5f62ux4e0eux540cux8c03}

\textbf{定义 70.1}（链复形）：一个（上）链复形 \(C^\bullet\) 是一列 Abel
群（或 \(R\)-模）\(C^n\) 和同态
\(d^n: C^n \to C^{n+1}\)（称为\textbf{微分}或\textbf{边缘算子}）满足
\(d^{n+1} \circ d^n = 0\)。即

\[\cdots \longrightarrow C^{n-1} \xrightarrow{d^{n-1}} C^n \xrightarrow{d^n} C^{n+1} \longrightarrow \cdots\]

（下链复形将箭头反向，\(d_n: C_n \to C_{n-1}\)。）

\textbf{定义 70.2}（上同调）：链复形 \(C^\bullet\) 的 \(n\)
次\textbf{上同调}为

\[H^n(C^\bullet) = \frac{\ker d^n}{\operatorname{im} d^{n-1}}\]

\(Z^n = \ker d^n\) 中的元素称为
\textbf{\(n\)-余圈}（cocycle），\(B^n = \operatorname{im} d^{n-1}\)
中的元素称为 \textbf{\(n\)-余边界}（coboundary）。

\textbf{定义 70.3}（链映射）：设 \(C^\bullet\) 和 \(D^\bullet\)
是链复形。一个\textbf{链映射} \(f: C^\bullet \to D^\bullet\) 是一族同态
\(f^n: C^n \to D^n\) 满足
\(d_D^n \circ f^n = f^{n+1} \circ d_C^n\)。链映射诱导上同调的同态
\(f^*: H^n(C^\bullet) \to H^n(D^\bullet)\)。

\textbf{定义 70.4}（链同伦）：两个链映射
\(f, g: C^\bullet \to D^\bullet\) 称为\textbf{链同伦的}，如果存在同态
\(h^n: C^n \to D^{n-1}\) 使得
\(f^n - g^n = d_D^{n-1} \circ h^n + h^{n+1} \circ d_C^n\)。链同伦的映射在上同调上诱导相同的同态。

\subsubsection{70.2
正合序列与蛇引理}\label{ux6b63ux5408ux5e8fux5217ux4e0eux86c7ux5f15ux7406}

\textbf{定义 70.5}（正合序列）：模同态的序列

\[\cdots \longrightarrow M_{i-1} \xrightarrow{f_{i-1}} M_i \xrightarrow{f_i} M_{i+1} \longrightarrow \cdots\]

称为在 \(M_i\) 处\textbf{正合的}，如果
\(\operatorname{im} f_{i-1} = \ker f_i\)。若序列在每处都正合，则称为\textbf{正合序列}。

\textbf{短正合序列}是形如
\(0 \to A \xrightarrow{i} B \xrightarrow{p} C \to 0\) 的正合序列（\(i\)
单，\(p\) 满，\(\ker p = \operatorname{im} i\)）。

\textbf{定理 70.1}（蛇引理）：设有 \(R\)-模的交换图（行正合）：

\[
\begin{CD}
A @>i>> B @>p>> C @>>> 0 \\
@VV{\alpha}V @VV{\beta}V @VV{\gamma}V \\
0 @>>> A' @>>{i'}> B' @>>{p'}> C'
\end{CD}
\]

则存在连接同态 \(\delta: \ker \gamma \to \operatorname{coker} \alpha\)
使得以下序列正合：

\[\ker \alpha \to \ker \beta \to \ker \gamma \xrightarrow{\delta} \operatorname{coker} \alpha \to \operatorname{coker} \beta \to \operatorname{coker} \gamma\]

\emph{证明}（构造 \(\delta\)）：对 \(c \in \ker \gamma\)，取 \(b \in B\)
使 \(p(b) = c\)。则 \(p'(\beta(b)) = \gamma(p(b)) = \gamma(c) = 0\)，故
\(\beta(b) \in \ker p' = \operatorname{im} i'\)。存在唯一 \(a' \in A'\)
使 \(i'(a') = \beta(b)\)。定义
\(\delta(c) = a' + \operatorname{im} \alpha \in \operatorname{coker} \alpha\)。验证
\(\delta\) 是良定义的且序列正合。∎

\textbf{定理 70.2}（五引理）：设有交换图（行正合）：

\[
\begin{CD}
A @>>> B @>>> C @>>> D @>>> E \\
@VV{\alpha}V @VV{\beta}V @VV{\gamma}V @VV{\delta}V @VV{\varepsilon}V \\
A' @>>> B' @>>> C' @>>> D' @>>> E'
\end{CD}
\]

\begin{enumerate}
\def\labelenumi{\arabic{enumi}.}
\tightlist
\item
  若 \(\alpha\) 满、\(\beta, \delta\) 单、\(\varepsilon\) 单，则
  \(\gamma\) 单。
\item
  若 \(\alpha\) 满、\(\beta, \delta\) 满、\(\varepsilon\) 单，则
  \(\gamma\) 满。
\item
  若 \(\alpha, \beta, \delta, \varepsilon\) 都是同构，则 \(\gamma\)
  也是同构。
\end{enumerate}

\subsubsection{70.3
投射模、内射模与平坦模}\label{ux6295ux5c04ux6a21ux5185ux5c04ux6a21ux4e0eux5e73ux5766ux6a21}

\textbf{定义 70.6}（投射模）：\(R\)-模 \(P\)
称为\textbf{投射模}，如果对任意满同态 \(f: M \to N\) 和同态
\(g: P \to N\)，存在提升 \(h: P \to M\) 使得 \(f \circ h = g\)。

\textbf{命题 70.3}：\(P\) 是投射模当且仅当 \(P\)
是某个自由模的直和项。特别地，自由模是投射模。

\textbf{定义 70.7}（内射模）：\(R\)-模 \(I\)
称为\textbf{内射模}，如果对任意单同态 \(f: A \to B\) 和同态
\(g: A \to I\)，存在扩张 \(h: B \to I\) 使得 \(h \circ f = g\)。

\textbf{命题 70.4}：每个 \(R\)-模都可以嵌入一个内射模。

\textbf{定义 70.8}（平坦模）：\(R\)-模 \(M\)
称为\textbf{平坦模}，如果函子 \(M \otimes_R -\) 是正合的（即保持单射）。

\textbf{命题 70.5}：投射模是平坦模。自由模是平坦模。

\subsubsection{70.4 导出函子：Ext 与
Tor}\label{ux5bfcux51faux51fdux5b50ext-ux4e0e-tor}

\textbf{定义 70.9}（投射分解）：\(M\)
的一个\textbf{投射分解}是一个正合序列

\[\cdots \to P_2 \to P_1 \to P_0 \to M \to 0\]

其中每个 \(P_i\) 是投射模。每个模都有投射分解（取自由分解即可）。

\textbf{定义 70.10}（Ext 函子）：对 \(R\)-模 \(M, N\)，取 \(M\)
的投射分解 \(P_\bullet \to M \to 0\)，应用函子
\(\operatorname{Hom}_R(-, N)\) 得链复形：

\[0 \to \operatorname{Hom}(P_0, N) \to \operatorname{Hom}(P_1, N) \to \operatorname{Hom}(P_2, N) \to \cdots\]

定义
\(\operatorname{Ext}_R^n(M, N) = H^n(\operatorname{Hom}_R(P_\bullet, N))\)。\(\operatorname{Ext}_R^n(M, N)\)
不依赖于投射分解的选取（在同构意义下）。

\textbf{定义 70.11}（Tor 函子）：对 \(R\)-模 \(M, N\)，取 \(M\)
的投射分解 \(P_\bullet \to M \to 0\)，应用函子 \(-\otimes_R N\)
得链复形：

\[\cdots \to P_2 \otimes N \to P_1 \otimes N \to P_0 \otimes N \to 0\]

定义 \(\operatorname{Tor}_n^R(M, N) = H_n(P_\bullet \otimes_R N)\)。

\textbf{命题 70.6}（Ext 和 Tor 的基本性质）： 1.
\(\operatorname{Ext}_R^0(M, N) \cong \operatorname{Hom}_R(M, N)\)。 2.
\(\operatorname{Tor}_0^R(M, N) \cong M \otimes_R N\)。 3. 短正合序列
\(0 \to A \to B \to C \to 0\) 诱导长正合序列：
\(\cdots \to \operatorname{Ext}^n(C, N) \to \operatorname{Ext}^n(B, N) \to \operatorname{Ext}^n(A, N) \to \operatorname{Ext}^{n+1}(C, N) \to \cdots\)
\(\cdots \to \operatorname{Tor}_{n+1}(C, N) \to \operatorname{Tor}_n(A, N) \to \operatorname{Tor}_n(B, N) \to \operatorname{Tor}_n(C, N) \to \cdots\)
4. \(M\) 投射 \(\Leftrightarrow\) \(\operatorname{Ext}_R^1(M, N) = 0\)
对所有 \(N\)。 5. \(M\) 平坦 \(\Leftrightarrow\)
\(\operatorname{Tor}_1^R(M, N) = 0\) 对所有 \(N\)。

\subsubsection{70.5 群上同调}\label{ux7fa4ux4e0aux540cux8c03}

\textbf{定义 70.12}（群上同调）：设 \(G\) 是群，\(M\) 是 \(G\)-模（即
\(M\) 是 Abel 群且有 \(G\) 的作用）。定义

\[H^n(G, M) = \operatorname{Ext}_{\mathbb{Z}G}^n(\mathbb{Z}, M)\]

其中 \(\mathbb{Z}G\) 是 \(G\) 的群环，\(\mathbb{Z}\) 是平凡
\(G\)-模（\(g \cdot m = m\)）。

\textbf{低次群上同调的解释}： -
\(H^0(G, M) = M^G = \{m \in M : g \cdot m = m, \; \forall g \in G\}\)（不变元）。
-
\(H^1(G, M)\)：交叉同态（\(f(gh) = f(g) + g \cdot f(h)\)）模去主交叉同态。当
\(G\) 作用平凡时，\(H^1(G, M) = \operatorname{Hom}(G, M)\)。 -
\(H^2(G, M)\)：群扩张 \(0 \to M \to E \to G \to 1\) 的分类（\(M\) 视为
Abel 正规子群）。

\textbf{定理 70.7}（Hilbert 定理 90，加法形式）：设 \(E/F\) 是有限
Galois 扩张，\(G = \operatorname{Gal}(E/F)\)。则
\(H^1(G, E) = 0\)（\(E\) 为加法群，\(G\) 自然作用）。

\emph{证明}：由本原元素定理，\(E = F(\alpha)\)。设 \(f: G \to E\)
是交叉同态。由 Dedekind 线性无关性引理，存在 \(\beta \in E\) 使得
\(\sum_{\sigma \in G} f(\sigma)\sigma(\beta) \neq 0\)（若
\(f \neq 0\)）。令 \(\gamma = \sum f(\sigma)\sigma(\beta)\)，可验证
\(f(\tau) = \tau(\theta) - \theta\) 对某 \(\theta\) 成立，故 \(f\)
是主交叉同态。∎

\bigskip\hrule\bigskip

\subsection{第71章：代数数论基础}\label{ux7b2c71ux7ae0ux4ee3ux6570ux6570ux8bbaux57faux7840}

代数数论研究代数数域（\(\mathbb{Q}\) 的有限扩张）的算术性质，是 Galois
理论和交换代数的交汇点。

\subsubsection{71.1
代数整数与整数环}\label{ux4ee3ux6570ux6574ux6570ux4e0eux6574ux6570ux73af}

\textbf{定义 71.1}（代数整数）：复数 \(\alpha \in \mathbb{C}\)
称为\textbf{代数整数}，如果存在首一整系数多项式
\(f(x) = x^n + a_{n-1}x^{n-1} + \cdots + a_0 \in \mathbb{Z}[x]\) 使
\(f(\alpha) = 0\)。

\textbf{命题 71.1}：\(\alpha\) 是代数整数当且仅当 \(\mathbb{Z}[\alpha]\)
是有限生成 \(\mathbb{Z}\)-模。

\textbf{定义 71.2}（数域的整数环）：设 \(K\) 是数域（\(\mathbb{Q}\)
的有限扩张）。\(K\) 中所有代数整数构成的集合记作 \(\mathcal{O}_K\)，称为
\(K\) 的\textbf{整数环}。

\textbf{定理 71.2}：\(\mathcal{O}_K\) 是 \(K\) 的子环，且
\(\mathcal{O}_K\) 是秩为 \([K : \mathbb{Q}]\) 的自由
\(\mathbb{Z}\)-模。即存在整基
\(\omega_1, \ldots, \omega_n\)（\(n = [K : \mathbb{Q}]\)）使得

\[\mathcal{O}_K = \mathbb{Z}\omega_1 \oplus \cdots \oplus \mathbb{Z}\omega_n\]

\textbf{例 71.1}（二次域的整数环）： - 若
\(K = \mathbb{Q}(\sqrt{d})\)，\(d \neq 1\) 是无平方因子整数，则 - 若
\(d \equiv 2, 3 \pmod{4}\)：\(\mathcal{O}_K = \mathbb{Z}[\sqrt{d}]\)。 -
若
\(d \equiv 1 \pmod{4}\)：\(\mathcal{O}_K = \mathbb{Z}[\frac{1+\sqrt{d}}{2}]\)。

\subsubsection{71.2 Dedekind
整环与理想分解}\label{dedekind-ux6574ux73afux4e0eux7406ux60f3ux5206ux89e3}

\textbf{定义 71.3}（Dedekind 整环）：整环 \(R\) 称为 \textbf{Dedekind
整环}，如果 \(R\) 是 Noether
的、整闭的，且每个非零素理想都是极大理想（即 \(\dim R = 1\)）。

\textbf{定理 71.3}：数域 \(K\) 的整数环 \(\mathcal{O}_K\) 是 Dedekind
整环。

\emph{证明}：\(\mathcal{O}_K\) 是 Noether 的（因为它是有限生成
\(\mathbb{Z}\)-模，每个理想也是有限生成
\(\mathbb{Z}\)-模）。它是整闭的（代数整数的集合在整闭包下封闭）。\(\dim \mathcal{O}_K = 1\)（\(\mathcal{O}_K\)
是 \(\mathbb{Z}\)
的整扩张，\(\dim \mathbb{Z} = 1\)，整扩张不改变维数）。∎

\textbf{定理 71.4}（Dedekind 整环中理想的唯一分解）：在 Dedekind 整环
\(R\) 中，每个非零真理想 \(\mathfrak{a}\)
可以唯一地（不计顺序）分解为素理想的乘积：

\[\mathfrak{a} = \mathfrak{p}_1^{e_1} \mathfrak{p}_2^{e_2} \cdots \mathfrak{p}_r^{e_r}\]

其中 \(\mathfrak{p}_i\) 是互不相同的非零素理想，\(e_i \geq 1\)。

\emph{证明思路}：利用 Noether 性和准素分解。在 Dedekind
整环中，准素理想恰为素理想的幂。整闭性保证局部化 \(R_{\mathfrak{p}}\)
是离散赋值环（DVR），从而理想唯一分解。∎

\textbf{推论 71.5}：\(\mathcal{O}_K\) 是 UFD 当且仅当 \(\mathcal{O}_K\)
是 PID（在 Dedekind 整环中，UFD 等价于 PID）。

\subsubsection{71.3 分歧理论}\label{ux5206ux6b67ux7406ux8bba}

\textbf{定义 71.4}（素理想的分歧）：设 \(L/K\)
是数域的扩张，\(\mathfrak{p}\) 是 \(\mathcal{O}_K\) 的非零素理想。则
\(\mathfrak{p}\mathcal{O}_L\) 可分解为 \(\mathcal{O}_L\)
中素理想的乘积：

\[\mathfrak{p}\mathcal{O}_L = \mathfrak{P}_1^{e_1} \mathfrak{P}_2^{e_2} \cdots \mathfrak{P}_g^{e_g}\]

其中 \(e_i = e(\mathfrak{P}_i/\mathfrak{p})\)
称为\textbf{分歧指数}，\(f_i = f(\mathfrak{P}_i/\mathfrak{p}) = [\mathcal{O}_L/\mathfrak{P}_i : \mathcal{O}_K/\mathfrak{p}]\)
称为\textbf{剩余类次数}。

\textbf{定理 71.6}（基本恒等式）：\(\sum_{i=1}^g e_i f_i = [L : K]\)。

\textbf{定义 71.5}（分歧类型）： - \(\mathfrak{p}\)
称为\textbf{非分歧的}，如果所有 \(e_i = 1\)。 - \(\mathfrak{p}\)
称为\textbf{完全分裂的}，如果 \(e_i = f_i = 1\) 对所有 \(i\)（此时
\(g = [L : K]\)）。 - \(\mathfrak{p}\) 称为\textbf{惯性的}，如果
\(g = 1, e = 1, f = [L : K]\)。

\textbf{定理 71.7}（Dedekind 判别法定理）：设
\(L = K(\alpha)\)，\(\alpha \in \mathcal{O}_L\)，\(f(x)\) 是 \(\alpha\)
的最小多项式。若 \(\mathfrak{p}\) 不整除
\([\mathcal{O}_L : \mathcal{O}_K[\alpha]]\)（指数），则
\(\mathfrak{p}\mathcal{O}_L\) 的分解由 \(f(x) \bmod \mathfrak{p}\)
的因式分解决定：

若 \(f(x) \equiv \prod f_i(x)^{e_i} \pmod{\mathfrak{p}}\)（\(f_i\)
不可约 mod \(\mathfrak{p}\)），则
\(\mathfrak{p}\mathcal{O}_L = \prod \mathfrak{P}_i^{e_i}\)，且
\(f(\mathfrak{P}_i/\mathfrak{p}) = \deg f_i\)。

\subsubsection{71.4
理想类群与类数}\label{ux7406ux60f3ux7c7bux7fa4ux4e0eux7c7bux6570}

\textbf{定义 71.6}（分式理想）：\(\mathcal{O}_K\)
的一个\textbf{分式理想}是 \(K\) 的一个有限生成 \(\mathcal{O}_K\)-子模
\(\mathfrak{a} \neq 0\)。等价地，存在
\(d \in \mathcal{O}_K \setminus \{0\}\) 使得
\(d\mathfrak{a} \subseteq \mathcal{O}_K\)（即 \(d\mathfrak{a}\) 是
\(\mathcal{O}_K\) 的整理想）。

分式理想在乘法下构成群
\(\mathcal{J}_K\)。主分式理想（\((\alpha) = \alpha\mathcal{O}_K\)，\(\alpha \in K^\times\)）构成子群
\(\mathcal{P}_K\)。

\textbf{定义 71.7}（理想类群）：\(K\) 的\textbf{理想类群}定义为

\[\operatorname{Cl}(K) = \mathcal{J}_K / \mathcal{P}_K\]

\(h_K = |\operatorname{Cl}(K)|\) 称为 \(K\) 的\textbf{类数}。

\textbf{定理 71.8}：理想类群 \(\operatorname{Cl}(K)\) 是有限群。

\emph{证明思路}：利用 Minkowski 的几何数论方法。\(\mathcal{O}_K\) 嵌入
\(\mathbb{R}^n \times \mathbb{C}^m\)（\(n + 2m = [K : \mathbb{Q}]\)）作为格。Minkowski
的凸体定理给出每个理想类中存在范数有界的理想，而范数有界的理想只有有限个。∎

\textbf{推论 71.9}：\(h_K = 1\) 当且仅当 \(\mathcal{O}_K\) 是
PID（等价于 UFD）。

\textbf{定理 71.10}（Dirichlet 单位定理）：\(\mathcal{O}_K\) 的单位群
\(\mathcal{O}_K^\times\) 是有限生成 Abel 群，其秩为
\(r = r_1 + r_2 - 1\)，其中 \(r_1\) 是实嵌入的个数，\(2r_2\)
是复嵌入的个数。即

\[\mathcal{O}_K^\times \cong \mu(K) \times \mathbb{Z}^r\]

其中 \(\mu(K)\) 是 \(K\) 中单位根的有限群。

\subsubsection{\texorpdfstring{71.5 \(\zeta\)
函数与类数公式}{71.5 \textbackslash zeta 函数与类数公式}}\label{zeta-ux51fdux6570ux4e0eux7c7bux6570ux516cux5f0f}

\textbf{定义 71.8}（Dedekind \(\zeta\) 函数）：数域 \(K\) 的 Dedekind
\(\zeta\) 函数定义为

\[\zeta_K(s) = \sum_{\mathfrak{a} \neq 0} \frac{1}{N(\mathfrak{a})^s} = \prod_{\mathfrak{p}} \frac{1}{1 - N(\mathfrak{p})^{-s}}\]

其中 \(\mathfrak{a}\) 遍历 \(\mathcal{O}_K\)
的所有非零理想，\(\mathfrak{p}\)
遍历所有非零素理想，\(N(\mathfrak{a}) = |\mathcal{O}_K/\mathfrak{a}|\)
是理想范数。此级数对 \(\Re(s) > 1\) 绝对收敛。

\textbf{定理 71.11}（\(\zeta_K(s)\)
的解析延拓与类数公式）：\(\zeta_K(s)\) 可解析延拓到 \(\mathbb{C}\)（除
\(s = 1\) 处的单极点外）。在 \(s = 0\) 处的留数涉及类数：

\[\lim_{s \to 0} s^{r_1+r_2-1} \zeta_K(s) = -\frac{h_K R_K}{w_K}\]

其中 \(R_K\) 是 \(K\) 的\textbf{调节子}（与单位群相关），\(w_K\) 是
\(K\) 中单位根的个数。

\emph{说明}：这是解析数论（V16）的重要基础，其完整证明需要更多的分析工具。

\bigskip\hrule\bigskip

\subsection{第72章：赋值论}\label{ux7b2c72ux7ae0ux8d4bux503cux8bba}

赋值论是研究域上''绝对值''及其完备化的理论，在代数数论和代数几何中至关重要。\textbf{赋值环}是整环\(R\)，其分式域\(K\)中任意非零元\(x\)满足\(x \in R\)或\(x^{-1} \in R\)；对应的\textbf{赋值}\(v: K^\times \to \Gamma\)（\(\Gamma\)为全序Abel群）满足\(v(xy) = v(x) + v(y)\)和\(v(x+y) \geq \min\{v(x), v(y)\}\)。当\(\Gamma \cong \mathbb{Z}\)时称为离散赋值环（DVR），其唯一极大理想\(\mathfrak{m} = \{x : v(x) > 0\}\)由单个元素\(\pi\)（素元）生成。\textbf{Hensel引理}是赋值论的核心定理：若\(R\)是完备赋值环，\(f \in R[x]\)及其模\(\mathfrak{m}\)约化\(\bar{f}\)满足\(\bar{f}\)有单根\(\bar{a}\)，则\(f\)有唯一根\(a \in R\)使得\(a \equiv \bar{a} \pmod{\mathfrak{m}}\)。Hensel环是满足Hensel引理的局部环，是局部域理论的基础。完备离散赋值环的结构由Cohen结构定理完全刻画：它们同构于域\(\kappa\)上的形式幂级数环\(\kappa[\![t]\!]\)或\(p\)-进整数环\(\mathbb{Z}_p\)的有限扩张。

\textbf{定义 72.1}（赋值）：域 \(K\) 上的 \textbf{（指数）赋值} 是满射
\(v: K^\times \to \Gamma\)，其中 \(\Gamma\) 是全序 Abel
群（值群），满足：（1）\(v(xy) = v(x) + v(y)\)；（2）\(v(x+y) \geq \min(v(x), v(y))\)。约定
\(v(0) = \infty\)（比 \(\Gamma\) 中任何元都大）。\textbf{赋值环}
\(R_v = \{x \in K : v(x) \geq 0\}\) 是包含 \(K\) 的单位元且对
\(x \in K^\times\) 有 \(x \in R_v\) 或 \(x^{-1} \in R_v\)
的局部环。其极大理想 \(\mathfrak{m}_v = \{x \in K : v(x) > 0\}\)，剩余域
\(k_v = R_v / \mathfrak{m}_v\)。

\textbf{定义 72.2}（离散赋值与 Hensel 引理）：若
\(\Gamma \cong \mathbb{Z}\)，\(v\) 称为 \textbf{离散赋值}，相应的
\(R_v\) 为 \textbf{离散赋值环（DVR）}------其恰好是一维 Noether
正则局部环（Krull-Akizuki 定理）。\textbf{Hensel 引理}：设 \(R_v\)
完备（关于 \(v\) 诱导的拓扑），\(f \in R_v[x]\) 满足
\(\bar{f} = f \bmod \mathfrak{m}_v\) 在 \(k_v[x]\) 中有单根
\(\bar{\alpha}\)，则存在唯一的 \(\alpha \in R_v\) 使得
\(\bar{\alpha} = \alpha \bmod \mathfrak{m}_v\) 且 \(f(\alpha) = 0\)。

\textbf{定理 72.1}（Cohen 结构定理）：完备的等特征 Noether
局部环同构于其剩余域上的形式幂级数环 \(k[[x_1, \ldots, x_d]]\)
的商环。完备的混合特征 DVR 同构于特征零剩余域特征 \(p\) 的 Witt
向量环的分歧扩张。

\bigskip\hrule\bigskip

\subsection{第73章：胎紧闭包（Tight
Closure）}\label{ux7b2c73ux7ae0ux80ceux7d27ux95edux5305tight-closure}

胎紧闭包（Tight Closure）是 Hochster 和 Huneke 于 1986 年引入的特征
\(p > 0\)
交换环中的闭包运算，为交换代数带来了一系列深刻的定理。这一理论提供了不同于经典同调代数的新工具，在解决若干长期未决猜想（特别是与正则环结构和理想理论相关的难题）中发挥了核心作用。

\subsubsection{73.1 基本定义与 Frost
性质}\label{ux57faux672cux5b9aux4e49ux4e0e-frost-ux6027ux8d28}

\textbf{定义 73.1}（理想的 Frobenius 幂）：设 \(R\) 是特征 \(p > 0\)
的含幺交换 Noether 环。对理想 \(I \subseteq R\) 和
\(q = p^e\)（\(e \geq 0\)），定义 \(I\) 的 \(q\) 次 \textbf{Frobenius
幂} \(I^{[q]}\) 为由 \(\{x^q : x \in I\}\) 生成的理想。即

\[I^{[q]} = \langle x^q : x \in I \rangle_R\]

记号 \(I^{[q]}\)（方括号）区别于通常的幂 \(I^q\)。通常
\(I^{[q]} \subseteq I^q\)，且若 \(I = (x_1, \ldots, x_n)\)，则
\(I^{[q]} = (x_1^q, \ldots, x_n^q)\)。

\textbf{定义 73.2}（胎紧闭包）：设 \(R\) 是特征 \(p > 0\) 的含幺交换
Noether 环，\(I \subseteq R\) 为理想。\(I\) 的\textbf{胎紧闭包}（tight
closure）\(I^*\) 定义为

\[I^* = \left\{x \in R : \exists\, c \in R^\circ,\; \forall\, q = p^e \gg 0,\; c x^q \in I^{[q]}\right\}\]

其中 \(R^\circ\) 表示 \(R\)
中所有极小素理想的补集之并（即不在任何极小素理想中的元素全体）。元素
\(c\) 称为 \(x\) 对 \(I\) 的\textbf{测试元}（test
element）。通常可选取不依赖于 \(x\) 和 \(I\) 的公共 \(c\)，这样的 \(c\)
称为\textbf{胎紧闭包测试元}（tight closure test element）。

胎紧闭包满足以下基本性质（Frost 性质）： -
\(I \subseteq I^* \subseteq \overline{I}\)（其中 \(\overline{I}\)
为整闭包）； - 若 \(I \subseteq J\)，则 \(I^* \subseteq J^*\)； -
\((I^*)^* = I^*\)（幂等性）； - \((I + J)^* = I^* + J^*\) 不一定成立，但
\((I \cap J)^* = I^* \cap J^*\)。

\subsubsection{73.2 Hochster-Huneke
定理}\label{hochster-huneke-ux5b9aux7406}

\textbf{定理 73.1}（Hochster-Huneke 正则环胎紧闭包定理）：若 \(R\)
是特征 \(p > 0\) 的 Noether 正则环，则对任意理想 \(I \subseteq R\)，有

\[I^* = I\]

即正则环中胎紧闭包等于理想本身。

\emph{证明思路}：核心步骤是证明正则环是\textbf{弱 F-正则的}（weakly
F-regular），即 \(I^* = I\) 对所有理想 \(I\) 成立。Hochster-Huneke
的证明分为以下几个关键环节：

\begin{enumerate}
\def\labelenumi{(\arabic{enumi})}
\item
  首先对正则局部环 \((R, \mathfrak{m})\)，利用 Auslander-Buchsbaum
  定理（正则局部环是 UFD）和 Frobenius 映射的平坦性（特征 \(p\) 正则环的
  Frobenius 映射是平坦的，Kunz 定理），将问题约化到余维数为 1 的情形。
\item
  利用\textbf{胎紧闭包与局部上同调的关系}证明：若 \(x \in I^*\)，则
  \(x\) 在局部上同调模 \(H_{\mathfrak{m}}^i(R)\)
  上的作用满足某种幂零性条件。对正则环，\(H_{\mathfrak{m}}^i(R)\) 仅在
  \(i = \dim R\) 处非零，且其结构简单（由 Čech 复形描述），从而推出
  \(x \in I\)。
\item
  更直接地，正则局部环具有有限 Frobenius 表示型（finite F-representation
  type），可构造性地找到 test element \(c = 1\)，从而 \(I^* = I\)。∎
\end{enumerate}

\textbf{推论 73.2}：正则环中，整闭包与胎紧闭包一致（因为 \(I = I^*\)，而
\(I \subseteq \overline{I}\)，但 \(I^* \subseteq \overline{I}\)，故
\(I = I^* = \overline{I}\) 仅当 \(I\) 本身已整闭）。

\subsubsection{\texorpdfstring{73.3 Briançon-Skoda 定理（特征 \(p\)
版本）}{73.3 Briançon-Skoda 定理（特征 p 版本）}}\label{brianuxe7on-skoda-ux5b9aux7406ux7279ux5f81-p-ux7248ux672c}

\textbf{定理 73.3}（Tight Closure 版本的 Briançon-Skoda 定理）：设 \(R\)
是特征 \(p > 0\) 的 Noether 正则环，\(I\) 是由 \(n\)
个元素生成的理想。则对任意 \(m \geq 1\)，

\[\overline{I^{n+m-1}} \subseteq I^m\]

特别地，取 \(m = 1\) 得 \(\overline{I^n} \subseteq I\)，即由 \(n\)
个元生成的理想的 \(n\) 次幂的整闭包包含于原理想。

\emph{证明思路}：经典 Briançon-Skoda 定理（复解析情形）的代数证明由
Lipman-Sathaye 和 Hochster-Huneke 给出。特征 \(p\) 方法的核心是：

\begin{enumerate}
\def\labelenumi{(\arabic{enumi})}
\item
  证明
  \(\overline{I^{n+m-1}} \subseteq (I^m)^*\)（整闭包包含于胎紧闭包）。这利用整闭包的刻画：\(x \in \overline{J}\)
  当且仅当存在 \(c \neq 0\) 使对所有 \(e\) 有
  \(c x^{p^e} \in J^{[p^e]}\)，以及 Frobenius 映射的性质。
\item
  由正则环的 Hochster-Huneke 定理，\((I^m)^* = I^m\)，从而结论成立。∎
\end{enumerate}

该定理的重大意义在于：它将一个高度非平凡的包含关系（整闭包包含于原理想）化归为胎紧闭包的简单性质。这在复解析几何中也有对应：若
\(\mathcal{I}\) 是由 \(n\) 个全纯函数生成的理想层，则
\(\overline{\mathcal{I}^n} \subseteq \mathcal{I}\)（Skoda 定理）。

\subsubsection{73.4 Colon-Capturing
与单项式猜想}\label{colon-capturing-ux4e0eux5355ux9879ux5f0fux731cux60f3}

\textbf{定理 73.4}（Colon-Capturing，Hochster-Roberts）：设 \(R\) 是特征
\(p > 0\) 的 Cohen-Macaulay 局部环，\(\mathfrak{a} \subseteq R\)
是理想，\(x_1, \ldots, x_d\) 为参数系（system of parameters）。则

\[(\mathfrak{a}, x_1, \ldots, x_{i-1}) : x_i \subseteq (\mathfrak{a}, x_1, \ldots, x_{i-1})^*\]

其中左侧为理想商（colon ideal），\(i = 1, \ldots, d\)。

\emph{证明思路}：利用胎紧闭包的定义和 Frobenius 映射的性质。设
\(y \in R\) 满足
\(y x_i \in (\mathfrak{a}, x_1, \ldots, x_{i-1})\)。需证明存在
\(c \in R^\circ\) 使对所有大 \(q = p^e\) 有
\(c y^q \in (\mathfrak{a}, x_1, \ldots, x_{i-1})^{[q]}\)。利用
Cohen-Macaulay 环中参数系的深度性质和 Frobenius 作用的良好性可得。∎

\textbf{单项式猜想（Monomial Conjecture）} 是 Hochster 1973
年提出的核心猜想：设 \((R, \mathfrak{m})\) 是 \(d\) 维 Noether
局部环，\(x_1, \ldots, x_d\) 为参数系，则对任意 \(t \geq 1\)，

\[x_1^t \cdots x_d^t \notin (x_1^{t+1}, \ldots, x_d^{t+1})\]

这个猜想起初看似仅仅是一个组合代数的问题，实则极其深刻，与 Krull
交定理、符号幂理论、同调猜想等多个领域紧密相关。Hochster 证明了：若特征
\(p > 0\) 情形下的单项式猜想成立，则由约化模 \(p\) 技术可推出等特征 0
情形的猜想。

利用 colon-capturing 定理，Hochster 证明了特征 \(p\)
情形下的单项式猜想。实际上，由 colon-capturing 可推出：若 \(R\) 是
Cohen-Macaulay 局部环，则对参数系 \(x_1, \ldots, x_d\) 有

\[(x_1, \ldots, x_i) : x_{i+1} \subseteq (x_1, \ldots, x_i)^*\]

这意味着参数系序列构成一个在胎紧闭包意义下的正则序列。单项式猜想作为该结论的特殊情形得以证明。

\subsubsection{73.5 Test Element
的存在性}\label{test-element-ux7684ux5b58ux5728ux6027}

\textbf{定理 73.5}（Test Element 存在性定理，Hochster-Huneke）：设 \(R\)
是特征 \(p > 0\) 的既约极好（excellent）局部环。则存在 \(c \in R^\circ\)
使得对任意理想 \(I \subseteq R\) 和任意 \(x \in I^*\)，有对所有大
\(q = p^e\)，\(c x^q \in I^{[q]}\)。这样的 \(c\) 称为 \(R\)
的\textbf{完全测试元}（complete test element）。

\emph{证明思路}：该定理的证明极为技术性，核心思想如下：

\begin{enumerate}
\def\labelenumi{(\arabic{enumi})}
\item
  对极好环 \(R\)，其 \(p\)-进完备化 \(\widehat{R}\) 是极好的，且
  \(R^\circ\) 在 \(\widehat{R}\) 中的像仍为 \(\widehat{R}^\circ\)
  的子集。利用完备局部环上 tight closure 的性质，可以在 \(R\)
  的正规化中寻找测试元。
\item
  利用 Gamma 构造（\(\Gamma\) 构造）：给定 \(R\)，可构造一个更大的环
  \(R^\Gamma\)，其 Frobenius 映射有更好的有限性。在 \(R^\Gamma\)
  中，可证明胎紧闭包保持测试元的''流动性''。
\item
  关键步骤是证明 \(R^\circ\) 中存在非零元素 \(c\) 满足：对所有含
  \(\mathfrak{m}\)-准素理想的理想
  \(I\)，\(c I^* \subseteq I\)。这利用了极好环的纤维良好性（excellence）和极小素理想的有限性。∎
\end{enumerate}

Test element 的存在性保证了 tight closure
理论的运算基础：无需为每个特定的 \(x\) 和 \(I\) 分别寻找
\(c\)，可以使用一个统一的 \(c\)
来测试所有胎紧闭包关系。这也使得胎紧闭包可视为一个''闭合算子''（closure
operator），其理论和计算工具与理想整闭包理论类似。

\subsubsection{73.6 应用与展望}\label{ux5e94ux7528ux4e0eux5c55ux671b}

胎紧闭包理论推动了交换代数的多个核心发展：

\begin{enumerate}
\def\labelenumi{\arabic{enumi}.}
\item
  \textbf{同调猜想}：Peskine-Szpiro 的相交定理和 Hochster
  的同调猜想（直接和猜想、提升猜想）均可由胎紧闭包方法在特征 \(p\)
  情形下证明，再通过约化模 \(p\) 推广到等特征 0 情形。
\item
  \textbf{Buchsbaum-Eisenbud
  的秩猜想}：胎紧闭包方法被用于证明关于有限自由分解的最小秩的若干猜想。
\item
  \textbf{符号幂理论}：Ein-Lazarsfeld-Smith（2001
  年）利用胎紧闭包的精化版本（multiplier ideals 和 asymptotic multiplier
  ideals）证明了光滑复簇上符号幂的渐近等式 \(I^{(rn)} \subseteq I^n\)。
\item
  \textbf{正特征代数几何}：胎紧闭包的对偶概念------\textbf{F-奇点理论}（F-singularities）------将正特征中的
  Frobenius 分裂与复代数几何中的奇点理论（log terminal、log canonical
  等）建立了对应。
\end{enumerate}

\bigskip\hrule\bigskip

\subsection{第74章：Jordan
代数}\label{ux7b2c74ux7ae0jordan-ux4ee3ux6570}

Jordan 代数是满足交换律 \(x \cdot y = y \cdot x\) 与 \textbf{Jordan
恒等式} \((x \cdot y) \cdot x^2 = x \cdot (y \cdot x^2)\)
的非结合代数。它由 Pascual Jordan 于 1933
年在形式化量子力学的背景下引入，旨在以对称积 \(a \circ b = (ab + ba)/2\)
来代替结合代数中的通常乘积，从而描述量子可观测量的代数结构。Jordan
代数看似是结合代数的简单对称化，实则包含了丰富而独立的代数世界------特别是
27 维的例外 Albert 代数，它是分类理论中的唯一例外对象。

\subsubsection{74.1 基本定义与 Jordan
恒等式}\label{ux57faux672cux5b9aux4e49ux4e0e-jordan-ux6052ux7b49ux5f0f}

\textbf{定义 74.1}（Jordan 代数）：域
\(k\)（\(\operatorname{char} k \neq 2\)）上的 \textbf{Jordan 代数}
是一个 \(k\)-向量空间 \(J\)，配以双线性乘法
\(\cdot: J \times J \to J\)，满足：

\begin{enumerate}
\def\labelenumi{\arabic{enumi}.}
\tightlist
\item
  \(x \cdot y = y \cdot x\)（交换律）；
\item
  \((x \cdot y) \cdot x^2 = x \cdot (y \cdot x^2)\)（Jordan 恒等式）。
\end{enumerate}

其中 \(x^2 = x \cdot x\)。注意 Jordan 恒等式弱于结合律------Jordan
代数不要求乘法满足结合律。Jordan 代数也可等价地以线性化的 Jordan
恒等式定义：

\[(x \cdot y) \cdot (z \cdot w) + (x \cdot w) \cdot (z \cdot y) + (x \cdot z) \cdot (y \cdot w) = x \cdot [y \cdot (z \cdot w)] + x \cdot [w \cdot (z \cdot y)] + x \cdot [z \cdot (y \cdot w)]\]

以及其更常见的极化形式：

\[(x, y, z^2) = 2(x \cdot z, y, z)\]

其中 \((x, y, z) = (x \cdot y) \cdot z - x \cdot (y \cdot z)\)
为结合子（associator）。

\textbf{命题 74.1}：任何 Jordan
代数都是\textbf{幂结合的}（power-associative），即
\(x^n = x \cdot x^{n-1}\) 的定义是无歧义的（与括弧位置无关）。

\emph{证明思路}：由 Jordan 恒等式可归纳证明 \(x^m \cdot x^n = x^{m+n}\)
对任意 \(m, n \geq 1\) 成立。∎

\subsubsection{74.2 特例 Jordan 代数与例外 Jordan
代数}\label{ux7279ux4f8b-jordan-ux4ee3ux6570ux4e0eux4f8bux5916-jordan-ux4ee3ux6570}

\textbf{定义 74.2}（特例 Jordan 代数，Special Jordan Algebra）：设 \(A\)
是结合代数。在向量空间 \(A\) 上定义 \textbf{Jordan 积}：

\[a \circ b = \frac{ab + ba}{2}\]

则 \((A, \circ)\) 是 Jordan 代数，记作 \(A^+\)。同构于某个 \(A^+\)
的子代数的 Jordan 代数称为\textbf{特例 Jordan 代数}（special Jordan
algebra）。

\textbf{例 74.1}（自伴矩阵代数）：设 \(M_n(\mathbb{R})\) 为
\(n \times n\) 实矩阵的代数。对称矩阵集合
\(\operatorname{Sym}_n(\mathbb{R})\) 在 Jordan 积
\(X \circ Y = (XY + YX)/2\) 下构成特例 Jordan 代数。这是 Jordan
最早的动机------量子可观测量的代数结构。

更一般地，设 \(D\) 是实数域上的带对合 \(*\) 的结合除代数，则
\(n \times n\) 矩阵代数 \(M_n(D)\) 中满足 \(X^* = X\) 的元素在 Jordan
积下构成特例 Jordan 代数。

\textbf{定义 74.3}（例外 Jordan 代数，Exceptional Jordan Algebra）：一个
\textbf{例外 Jordan 代数} 是不与任何结合代数的 \(A^+\) 的子代数同构的
Jordan 代数。

\textbf{定理 74.2}（Albert 代数）：设 \(\mathbb{O}\) 为 Cayley
八元数代数。\(3 \times 3\) 的 Hermite 矩阵

\[H_3(\mathbb{O}) = \left\{\begin{pmatrix} a & z & \bar{y} \\ \bar{z} & b & x \\ y & \bar{x} & c \end{pmatrix} : a, b, c \in \mathbb{R},\; x, y, z \in \mathbb{O}\right\}\]

在 Jordan 积下构成 27 维 Jordan 代数。\(H_3(\mathbb{O})\) 是例外 Jordan
代数，称为 \textbf{Albert 代数}。

\emph{证明例外性的思路}：若 \(H_3(\mathbb{O})\)
是特例的，则将存在某个结合代数 \(A\) 和单射同态
\(H_3(\mathbb{O}) \hookrightarrow A^+\)。通过分析其幂等元的正交分解（Albert
代数中最多有 3
个正交的原始幂等元），可推出这与八元数的非结合性矛盾。Glennie 于 1963
年找到了第一个例外 Jordan
代数不满足的特例恒等式（s-identity），从而证明了其例外性。∎

\subsubsection{74.3 形式实 Jordan
代数与分类定理}\label{ux5f62ux5f0fux5b9e-jordan-ux4ee3ux6570ux4e0eux5206ux7c7bux5b9aux7406}

\textbf{定义 74.4}（形式实 Jordan 代数）：Jordan 代数 \(J\)
称为\textbf{形式实的}，如果对任意有限个
\(x_i \in J\)，\(\sum x_i^2 = 0\) 蕴含所有
\(x_i = 0\)。形式实条件相当于正定性条件，保证了代数具有 Euclid
型的内积结构。

\textbf{定理 74.3}（Jordan-von Neumann-Wigner 分类定理，1934
年）：有限维形式实 Jordan 代数恰好是以下四类的直和：

\begin{enumerate}
\def\labelenumi{\arabic{enumi}.}
\tightlist
\item
  \textbf{\(H_n(\mathbb{R})\)}：\(n \times n\) 实对称矩阵，维数为
  \(n(n+1)/2\)；
\item
  \textbf{\(H_n(\mathbb{C})\)}：\(n \times n\) 复 Hermite 矩阵，维数为
  \(n^2\)；
\item
  \textbf{\(H_n(\mathbb{H})\)}：\(n \times n\) 四元数 Hermite
  矩阵，维数为 \(n(2n-1)\)；
\item
  \textbf{旋因子}（spin
  factors）\(\mathbb{R} \oplus \mathbb{R}^n\)：由内积空间 \((V, q)\)
  构造，乘法定义为
  \((\lambda, v) \cdot (\mu, w) = (\lambda\mu + q(v, w), \lambda w + \mu v)\)，维数为
  \(n+1\)；
\item
  \textbf{Albert 代数} \(H_3(\mathbb{O})\)：\(3 \times 3\) 八元数
  Hermite 矩阵，维数为 27。
\end{enumerate}

\emph{证明思路}：证明的核心是分析形式实 Jordan
代数的幂等元分解。有限维形式实 Jordan
代数必有单位元，且单位元可唯一地分解为正交原始幂等元的和
\(1 = \sum e_i\)。通过分析 Peirce
分解（\(J = \bigoplus_{i \leq j} J_{ij}\) 其中 \(J_{ii}\)
是一维的，\(J_{ij}\) 为 \(e_i\)-特征空间），Jordan-von Neumann-Wigner
逐步证明了： - 若 Peirce 分量全为一维（即连接数 \(c_{ij} = 1\) 对所有
\(i \neq j\)），则得到旋因子； - 若 Peirce
分量可承载非平凡的结合代数结构，则通过 Hurwitz
定理（实赋范可除代数的分类）得到
\(\mathbb{R}, \mathbb{C}, \mathbb{H}, \mathbb{O}\) 四种情形； - Albert
代数是唯一既约但非特例的情形（对应 \(n=3\) 的八元数情形，\(n \geq 4\)
的八元数并不能形成 Jordan 代数）。∎

\subsubsection{74.4 与对称空间和 Lie
代数的联系}\label{ux4e0eux5bf9ux79f0ux7a7aux95f4ux548c-lie-ux4ee3ux6570ux7684ux8054ux7cfb}

Jordan 代数通过 \textbf{Kantor-Koecher-Tits 构造}（KKT 构造）与 Lie
代数建立了深刻联系。

\textbf{定理 74.4}（KKT 构造）：设 \(J\) 是 Jordan 代数。定义向量空间

\[T(J) = J \oplus \operatorname{Der}(J) \oplus J\]

其中 \(\operatorname{Der}(J)\) 为 \(J\) 的导子代数。可在 \(T(J)\)
上定义一个 3-次 Lie 代数结构（即 \(\mathbb{Z}\)-分次 Lie 代数
\(\mathfrak{g} = \mathfrak{g}_{-1} \oplus \mathfrak{g}_0 \oplus \mathfrak{g}_1\)），使得
\(\mathfrak{g}_{-1} \cong \mathfrak{g}_1 \cong J\)，\(\mathfrak{g}_0 \cong \operatorname{Str}(J)\)（\(J\)
的结构代数）。

具体 Lie 括号由以下公式给出：对
\(x, y \in \mathfrak{g}_{-1} \cong J\)，\([x, y] = -2L(x, y) \in \operatorname{Der}(J)\)，其中
\(L(x, y)z = (x \cdot y) \cdot z - x \cdot (y \cdot z)\)。

\emph{意义}：此构造将任意 Jordan 代数''线性化''为 3-次 Lie 代数，从而将
Jordan 代数的分类问题归入 Lie 代数的分类框架。反之，Lie
代数的最高权向量表示理论可用来构造 Jordan 代数。

\textbf{定理 74.5}（Koecher-Tits 构造）：形式实 Jordan 代数 \(J\)
的对称锥 \(\Omega = \{x^2 : x \in J\}\) 的复化是 Hermite
对称空间。这给出了\textbf{有界对称域}（bounded symmetric
domains）与形式实 Jordan 代数的一一对应：

\begin{itemize}
\tightlist
\item
  \(\operatorname{Sym}_n(\mathbb{R})^+ \leftrightarrow\) Siegel 上半空间
  \(\mathbb{H}_n\)（III 型有界对称域）；
\item
  \(\operatorname{Herm}_n(\mathbb{C})^+ \leftrightarrow\) I
  型有界对称域（Cartan 的四种古典域中的第一种）；
\item
  旋因子 \(\leftrightarrow\) IV 型有界对称域；
\item
  Albert 代数 \(\leftrightarrow\) 27 维的例外有界对称域（E.VII 型）。
\end{itemize}

这一对应是 Elie Cartan 对称空间分类的 Jordan
代数诠释，也是非交换调和分析和表示论的重要工具。

\bigskip\hrule\bigskip

\subsection{第75章：结晶群（Crystallographic
Groups）}\label{ux7b2c75ux7ae0ux7ed3ux6676ux7fa4crystallographic-groups}

结晶群是 \(n\) 维 Euclid 空间 \(E(n) = \mathbb{R}^n \rtimes O(n)\)
的离散余紧子群，其理论由 Bieberbach 在 1911-1912
年间奠基。结晶群不仅是几何群论的基本对象，也是数学晶体学、固态物理和材料科学的数学基础------自然界中的晶体结构恰好对应
\(n = 3\) 的结晶群（230 种空间群）。

\subsubsection{75.1 基本定义与 Bieberbach
定理}\label{ux57faux672cux5b9aux4e49ux4e0e-bieberbach-ux5b9aux7406}

\textbf{定义 75.1}（Euclid 运动群）：\(n\) 维 \textbf{Euclid 运动群}
\(E(n)\) 是 \(\mathbb{R}^n\) 的所有等距变换构成的群。由定义
\(E(n) = \mathbb{R}^n \rtimes O(n)\)，其中 \(\mathbb{R}^n\)
为平移子群（正规子群），\(O(n)\) 为正交群。元记为
\((v, A)\)（\(v \in \mathbb{R}^n, A \in O(n)\)），作用为
\(x \mapsto Ax + v\)。乘法为

\[(v_1, A_1)(v_2, A_2) = (v_1 + A_1 v_2, A_1 A_2)\]

\textbf{定义 75.2}（结晶群）：\(n\) 维 \textbf{结晶群}（crystallographic
group）是 \(E(n)\) 的离散子群 \(\Gamma\)，满足商空间
\(\mathbb{R}^n / \Gamma\) 是紧的。\(n = 2\) 时称 \textbf{wallpaper
group}（壁纸群），\(n = 3\) 时称\textbf{空间群}（space group）。

\textbf{定理 75.1}（Bieberbach 第一定理，1911）：设
\(\Gamma \subset E(n)\) 是 \(n\) 维结晶群。则 \(\Gamma\) 的平移子群
\(T = \Gamma \cap \mathbb{R}^n\)（其中 \(\mathbb{R}^n\)
等同于纯平移子群）是秩为 \(n\) 的自由 Abel 群，且
\([\Gamma : T] < \infty\)。

\emph{证明思路}：核心是证明 \(\Gamma\) 的平移部分构成离散余紧子群。设
\(\Gamma\) 在 \(\mathbb{R}^n\) 上的自然作用。由于 \(\Gamma\) 是离散的且
\(\mathbb{R}^n / \Gamma\) 紧，平移子群必为 \(\mathbb{R}^n\)
的格（lattice），即秩为 \(n\) 的自由 Abel 群。商群 \(\Gamma / T\) 同构于
\(O(n)\) 的有限子群（称为\textbf{点群}），故
\([\Gamma : T] < \infty\)。∎

\textbf{定理 75.2}（Bieberbach 第二定理，1911，刚性定理）：若
\(\Gamma_1, \Gamma_2 \subset E(n)\) 是两个同构的结晶群，则存在仿射变换
\(\varphi \in \operatorname{Aff}(\mathbb{R}^n)\) 使得
\(\varphi \Gamma_1 \varphi^{-1} = \Gamma_2\)。

换言之，结晶群的同构类与仿射共轭类一致------不存在同构但不等价的结晶群。这在几何群论中是一个罕见的刚性现象。

\textbf{定理 75.3}（Bieberbach 第三定理，1912）：对每个维数
\(n\)，只存在有限多个 \(n\) 维结晶群的同构类（等价地，仿射共轭类）。

\emph{证明思路}：固定维数 \(n\)，取基域 \(T \cong \mathbb{Z}^n\)，点群
\(P = \Gamma / T\) 必然是 \(GL(n, \mathbb{Z})\)
的有限子群。\(GL(n, \mathbb{Z})\)
的有限子群只有有限多个共轭类（Jordan-Zassenhaus 定理）。对每个有限子群
\(P \subset GL(n, \mathbb{Z})\)，结晶群的结构由第二上同调类
\(H^2(P, \mathbb{Z}^n)\) 决定，而这是有限群。∎

\subsubsection{75.2 Zassenhaus
算法与空间群分类}\label{zassenhaus-ux7b97ux6cd5ux4e0eux7a7aux95f4ux7fa4ux5206ux7c7b}

\textbf{定义 75.3}（Zassenhaus 算法）：Zassenhaus（1948
年）提出了一种系统方法列举所有 \(n\) 维空间群。算法基于 Bieberbach
定理：

\begin{enumerate}
\def\labelenumi{\arabic{enumi}.}
\tightlist
\item
  列举 \(GL(n, \mathbb{Z})\) 的所有有限子群（\(n \leq 3\) 时已知分类）；
\item
  对每个有限子群 \(P\)，列举 \(H^2(P, \mathbb{Z}^n)\)
  的元素（群扩张类）；
\item
  对每个扩张类中代表元，检查其在 \(E(n)\) 中的嵌入（判别离散性）。
\end{enumerate}

\textbf{定理 75.4}（空间群分类）： - \(n = 1\)：2 种（\(\mathbb{Z}\)
和无限二面体群 \(D_\infty\)）； - \(n = 2\)：17 种 wallpaper groups（由
Fedorov 1891 和 Pólya 1924 年完成）； - \(n = 3\)：219
种空间群（若不区分对映体），或 230
种（区分对映体后。Fedorov-Schoenflies-Barlow，1890 年代）； -
\(n = 4\)：4783 种（Brown-Bülow-Neubüser-Wondratschek-Zassenhaus，1978
年）； - \(n = 5\)：222018 种（Plesken-Schulz，2000 年）； -
\(n = 6\)：28927922 种（Oppenorth，2000 年之后）。

\(n = 3\) 的 230 种空间群是国际晶体学表（International Tables for
Crystallography）的基础，直接用于 X 射线晶体学中晶体结构的确定。

\subsubsection{\texorpdfstring{75.3 Wallpaper
群分类（\(n = 2\)）}{75.3 Wallpaper 群分类（n = 2）}}\label{wallpaper-ux7fa4ux5206ux7c7bn-2}

\textbf{定理 75.5}（Wallpaper 群的 17 种分类）：二维结晶群共有 17
个同构类，分类如下（按点群类型）：

{\def\LTcaptype{none} % do not increment counter
\begin{longtable}[]{@{}llll@{}}
\toprule\noalign{}
点群 & 阶 & 群数 & 记号 \\
\midrule\noalign{}
\endhead
\bottomrule\noalign{}
\endlastfoot
\(C_1\) & 1 & 1 & p1 \\
\(C_2\) & 2 & 1 & p2 \\
\(C_3\) & 3 & 1 & p3 \\
\(C_4\) & 4 & 1 & p4 \\
\(C_6\) & 6 & 1 & p6 \\
\(D_1\) & 2 & 3 & pm, pg, cm \\
\(D_2\) & 4 & 4 & pmm, pmg, pgg, cmm \\
\(D_3\) & 6 & 2 & p3m1, p31m \\
\(D_4\) & 8 & 2 & p4m, p4g \\
\(D_6\) & 12 & 1 & p6m \\
\end{longtable}
}

\emph{代数刻画}：每个 wallpaper group 的平移格 \(\mathbb{Z}^2\) 配以
\(GL(2, \mathbb{Z})\)
的有限子群的作用（必须满足晶体学限制定理------旋转阶数只能为 1, 2, 3, 4
或 6）。不同 wallpaper groups 的区别在于 \(H^2(P, \mathbb{Z}^2)\)
的元素如何影响反射和滑动反射。

\textbf{例 75.1}（17 种 wallpaper groups 的几何）：p1
是最简单的（仅平移），p6m 是最对称的（包含 6 阶旋转和反射）。M.C. Escher
的艺术作品广泛运用了这些对称模式。

\subsubsection{75.4
与几何群论的联系}\label{ux4e0eux51e0ux4f55ux7fa4ux8bbaux7684ux8054ux7cfb}

\textbf{Grigorchuk 群}（Grigorchuk，1980
年）具有重要意义的例子：它是一个有限生成群，具有介于多项式和指数之间的中间增长（intermediate
growth），是第一个被证明为非初等 amenable 群的例子。Grigorchuk
群的构造自然地与结晶群的分支覆盖联系：将结晶群作用在无限树上产生无限迭代自相似群。

结晶群理论为 Bieberbach 定理背后的几何刚性提供了群论解释。Hilton 将
wallpaper 群分类推广到一般的晶体学限制，而广义结晶群更是与 CAT(0)
空间等非正曲率几何理论交叉的前沿话题。

\bigskip\hrule\bigskip

\subsection{第76章：准齐性向量空间（Prehomogeneous Vector
Spaces）}\label{ux7b2c76ux7ae0ux51c6ux9f50ux6027ux5411ux91cfux7a7aux95f4prehomogeneous-vector-spaces}

准齐性向量空间（Prehomogeneous Vector Spaces, PVS）由佐藤幹夫（Sato
Mikio）于 1961
年引入，是代数群在向量空间上的作用具有稠密开轨道的结构。该理论连接了代数群表示论、微分方程理论和数论，是
Sato 的 D-模理论和超函数理论的重要支柱。

\subsubsection{76.1
基本概念与定义}\label{ux57faux672cux6982ux5ff5ux4e0eux5b9aux4e49}

\textbf{定义 76.1}（准齐性向量空间）：设 \(G\) 是代数闭域 \(k\)
上的连通代数群，\(\rho: G \to GL(V)\) 是 \(G\) 在有限维 \(k\)-向量空间
\(V\) 上的有理表示。三元组 \((G, \rho, V)\)
称为\textbf{准齐性向量空间}（PVS），若 \(V\) 中存在开的稠密 \(G\)-轨道。

等价地，存在点 \(v_0 \in V\) 使得轨道 \(G \cdot v_0\) 包含于 \(V\) 的
Zariski 开稠子集。此时 \(v_0\) 称为该 PVS 的\textbf{一般点}（generic
point），其稳定子群 \(G_{v_0} = \{g \in G : g \cdot v_0 = v_0\}\) 是
\(G\) 的闭子群，且 \(\dim G - \dim G_{v_0} = \dim V\)。

\textbf{定义 76.2}（相对不变量）：PVS \((G, \rho, V)\) 上的非零有理函数
\(f \in k(V)^\times\) 称为\textbf{相对不变量}，若存在 \(G\) 的有理特征标
\(\chi: G \to \mathbb{G}_m\) 使得

\[f(g \cdot v) = \chi(g) f(v)\]

对所有 \(g \in G, v \in V\) 成立（在定义域内）。\(\chi\) 称为 \(f\)
的特征标。

相对不变量的零点定义了 \(G\)-不变的超曲面，这些超曲面由较低维的
\(G\)-轨道构成。具体地，若 \(V \setminus D\) 是 \(V\) 的唯一开稠密
\(G\)-轨道，则补集 \(D\) 是
\(G\)-不变超曲面的并，恰为相对不变量的零点集。

\subsubsection{\texorpdfstring{76.2 \(b\)-函数与 Bernstein-Sato
多项式}{76.2 b-函数与 Bernstein-Sato 多项式}}\label{b-ux51fdux6570ux4e0e-bernstein-sato-ux591aux9879ux5f0f}

\textbf{定义 76.3}（\(b\)-函数 / Bernstein-Sato 多项式）：设
\(f \in \mathbb{C}[x_1, \ldots, x_n]\) 是非零多项式。存在非零微分算子
\(P(s, \partial) \in \mathbb{C}[x, \partial_x, s]\) 和单变量多项式
\(b(s) \in \mathbb{C}[s]\) 使得

\[b(s) f(x)^s = P(s, \partial) f(x)^{s+1}\]

对所有 \(s \in \mathbb{C}\)（在 \(f \neq 0\)
的区域上）成立。满足此式的最小次数首一多项式 \(b(s)\) 称为 \(f\) 的
\textbf{Bernstein-Sato 多项式}或 \textbf{\(b\)-函数}。

\textbf{定理
76.1}（Bernstein-Sato，\(b\)-函数的存在性）：对任意非零多项式
\(f \in \mathbb{C}[x_1, \ldots, x_n]\)，其
\(b\)-函数存在。更一般地，Kashiwara（1976 年）利用
D-模理论证明了任意全纯函数的 \(b\)-函数总是存在的。

\emph{证明思路}（D-模方法）：将 \(f\) 的纤维化
\(\mathbb{C}^n \to \mathbb{C}\)（\(x \mapsto f(x)\)）视为
D-模的态射。构造
\(\mathcal{M} = \mathbb{C}[x, \partial_x, s] f^s\)（带有微分算子作用的自由模）。Kashiwara
证明了 \(\mathcal{M}\) 是 \(\mathbb{C}[s]\)
上的有限生成模，因此存在多项式 \(b(s)\) 满足所需关系。∎

\textbf{定理 76.2}（\(b\)-函数零点的性质）： 1. \(b(s)\)
的零点均为负有理数； 2. \(b(s)\) 的零点集合 \(\{\alpha\}\) 与 \(f\)
确定的奇点的局部 monodromy 的特征值（乘以
\(e^{2\pi i \alpha}\)）密切相关； 3. \(-1\) 是 \(b(s)\) 的零点
\(\Leftrightarrow\) \(f\) 不全纯。

\emph{证明思路}：(1) 利用 \(\mathbb{C}[s]\) 上的 Brieskorn 格和
Gauss-Manin 联络，\(b\)-函数的零点与 monodromy 的 Hodge
谱相关，从而证明其负有理性。(2) 通过 Malgrange-Kashiwara 构造，将
\(f^s\) 视为全形 D-模 \(\mathcal{O}\) 的扩张，其 \(s\)
参数对应微分算子的过滤。∎

\subsubsection{\texorpdfstring{76.3 PVS 中的
\(b\)-函数与微局部分析}{76.3 PVS 中的 b-函数与微局部分析}}\label{pvs-ux4e2dux7684-b-ux51fdux6570ux4e0eux5faeux5c40ux90e8ux5206ux6790}

在准齐性向量空间中，\(b\)-函数与相对不变量有直接联系。

\textbf{定理 76.3}（PVS 的 \(b\)-函数）：设 \((G, \rho, V)\) 是定义在
\(\mathbb{Q}\) 上的 PVS，\(f(x)\) 是其相对不变量（多项式）。则 \(f\) 的
\(b\)-函数具有以下形式：

\[b_f(s) = \prod_{i=1}^m (s + \alpha_i)\]

其中 \(\alpha_i\) 是 \(G\)
的有理特征标的线性组合所得的有理数。具体地，\(\alpha_i\) 的表达式中涉及
\(G\) 的根和权以及 \(V\) 的权。

\emph{证明思路}：利用群 \(G\) 的 Casimir 算子（\(G\) 的 Lie 代数
\(\mathfrak{g}\) 的泛包络代数中的二次元素）作用于 \(f^s\)，构造满足
\(b(s)\) 关系式的微分算子。Casimir 算子在 \(f^s\) 上的特征值给出
\(b(s)\) 的因子。∎

\textbf{定理 76.4}（微局部分析的 \(b\)-函数，Kashiwara）：利用 D-模的
V-过滤（V-filtration）和微局部特征簇（characteristic
variety），\(b\)-函数控制了全形 D-模 \(\mathcal{D}[s]f^s\)
的微局部结构。具体地，\(b_a(s)\)（局部 \(b\)-函数）的零点对应于
\(\mu\)-常值分层的 vanishing cycles 的 monodromy 特征值。

\subsubsection{76.4 Sato-Kimura
分类定理}\label{sato-kimura-ux5206ux7c7bux5b9aux7406}

\textbf{定理 76.5}（Sato-Kimura 分类定理，1977 年）：设 \(G\) 是
\(\mathbb{C}\) 上的连通简约代数群，\(\rho\) 是 \(G\)
的既约有理表示。所有满足 \(V\) 上存在开稠密 \(G\)-轨道的既约 PVS（即
\((G, \rho, V)\) 是既约准齐性向量空间）被分类为 \textbf{29 个系列}。

这 29 个系列包含： - \textbf{抛物型 PVS}：\(G\) 作用于其 Lie
代数的幂零锥或相关表示； - \textbf{旋表示型}：如 \(Spin(10)\) 的 16
维半旋表示； - \textbf{例外群型}：如 \(E_7\) 的 56 维微小表示； -
\textbf{产物构造}：由 Castling 变换（两个已知 PVS
的乘积通过维数匹配的等价关系关联）。

\emph{证明思路}：Sato 和 Kimura 的方法基于对 Dynkin
图的系统性分析。步骤为： (1) 若 \((G, \rho, V)\) 是 PVS，则
\(\dim G \geq \dim V\)。\(\dim G - \dim V = \dim G_v\)（一般点的稳定子群的维数）。
(2) 利用 Weyl 的维数公式，对每个既约表示计算 \(\dim V\)，列举
\(\dim G \geq \dim V\) 的情形。 (3) 对候选情形，分析是否存在 \(G\)
的抛物子群 \(P\) 使其维数匹配。Castling
变换建立了不同系列之间双有理等价关系。 (4)
最终用约化群的强正则性（Vinberg）排除不可能的情形。

29 个系列中，前 15 个是二元型（\(G\) 为线性代数群），后 14
个涉及旋表示和例外群。∎

\textbf{定理 76.6}（Castling 变换）：设 \((G, \rho, V)\) 和
\((H, \sigma, W)\) 是两个 PVS，\(m, n\) 为正整数。若
\(m \dim V = n \dim W\) 且
\(m \dim G_v = n \dim H_w\)（对一般点而言），则
\((GL(m) \times G, \mathbb{C}^m \otimes V)\) 与
\((GL(n) \times H, \mathbb{C}^n \otimes W)\) 是双有理等价的 PVS（通过
Castling 变换）。该等价将 29 个系列的分类大幅简化。

\subsubsection{76.5 应用与联系}\label{ux5e94ux7528ux4e0eux8054ux7cfb}

\begin{enumerate}
\def\labelenumi{\arabic{enumi}.}
\item
  \textbf{局部 zeta 函数与函数等式}：Sato 和 Shintani（1974 年）利用 PVS
  理论计算了局部 zeta 函数（\(p\)-进域上的 zeta
  积分），建立了局部函数方程。PVS
  的相对不变量给出自然积分核，\(b\)-函数控制解析延拓。
\item
  \textbf{概均质（prehomogeneous）zeta 函数}：Sato-Shintani 的全局 zeta
  函数与 D-模上同调有深刻联系，类似于 Tate 的博士论文中 Riemann zeta
  函数的处理方法。
\item
  \textbf{概均质向量空间与数论}：Wright-Yukie 利用 PVS
  理论研究了数域的密度定理和类数问题。
\item
  \textbf{奇点理论}：\(b\)-函数是奇点不变量理论（与 Milnor 纤维化和
  monodromy）的基本工具。
\end{enumerate}

\bigskip\hrule\bigskip

\emph{本卷在 V8 的基础上完成了抽象代数核心领域的深化：Galois
理论建立了域扩张与群论的桥梁，彻底解决了方程根式可解性问题；交换代数为代数几何和代数数论提供了理想论、局部化、Noether
环等基础工具；表示论用特征标理论揭示了有限群的结构；同调代数引入 Ext/Tor
函子和导出函子，统一了代数拓扑和代数几何的上同调理论；代数数论建立了
Dedekind
整环中理想的唯一分解定理和理想类群理论。为后续代数几何（V15）、代数拓扑（V14）和数论
II（V16）的学习奠定了基础。}

\newpage

\section{卷十三：泛函分析}\label{ux5377ux5341ux4e09ux6cdbux51fdux5206ux6790}

\begin{quote}
泛函分析是线性代数在无穷维空间的推广，研究函数空间上的线性算子及其谱理论。本卷从
Banach 空间和 Hilbert
空间的基本定理出发，建立对偶理论和弱拓扑，进而研究紧算子和谱理论，最后介绍
C*-代数初步，为量子力学、偏微分方程和调和分析提供数学基础。
\end{quote}

\bigskip\hrule\bigskip

\subsection{第72章：Banach
空间}\label{ux7b2c72ux7ae0banach-ux7a7aux95f4}

Banach
空间是完备的赋范向量空间，是泛函分析最基本的框架。Hahn-Banach、开映射、闭图像和一致有界原理被称为
Banach 空间的四大基本定理。

\subsubsection{72.1 赋范空间与 Banach
空间}\label{ux8d4bux8303ux7a7aux95f4ux4e0e-banach-ux7a7aux95f4}

\textbf{定义 72.1}（赋范空间）：设 \(X\) 是数域
\(\mathbb{K}\)（\(\mathbb{R}\) 或 \(\mathbb{C}\)）上的向量空间。映射
\(\|\cdot\|: X \to [0, \infty)\) 称为 \(X\)
上的一个\textbf{范数}，如果满足：

\begin{enumerate}
\def\labelenumi{\arabic{enumi}.}
\tightlist
\item
  \(\|x\| = 0 \Leftrightarrow x = 0\)（正定性）
\item
  \(\|\alpha x\| = |\alpha| \|x\|\)（齐次性）
\item
  \(\|x + y\| \leq \|x\| + \|y\|\)（三角不等式）
\end{enumerate}

\((X, \|\cdot\|)\) 称为\textbf{赋范向量空间}。范数诱导度量
\(d(x, y) = \|x - y\|\) 和拓扑。

\textbf{定义 72.2}（Banach 空间）：完备的赋范向量空间称为 \textbf{Banach
空间}。即每个 Cauchy 序列都收敛。

\textbf{例 72.1}（Banach 空间举例）： - \(\mathbb{R}^n\) 和
\(\mathbb{C}^n\)（有限维赋范空间总是完备的） - \(\ell^p\)
空间（\(1 \leq p \leq \infty\)）：\(\|(x_n)\|_p = (\sum_{n=1}^\infty |x_n|^p)^{1/p}\)（\(p < \infty\)），\(\|(x_n)\|_\infty = \sup_n |x_n|\)
- \(L^p(\Omega)\)（\(1 \leq p \leq \infty\)）：由 V9 第 48 章的
Riesz-Fischer 定理保证完备性 - \(C([a, b])\)：连续函数空间配以
\(\|f\|_\infty = \sup_{x \in [a, b]} |f(x)|\) -
\(c_0\)：收敛到零的序列空间，配以 \(\|\cdot\|_\infty\)

\textbf{定义 72.3}（等价范数）：\(X\) 上的两个范数 \(\|\cdot\|_1\) 和
\(\|\cdot\|_2\) 称为\textbf{等价的}，如果存在常数 \(c, C > 0\) 使得
\(c\|x\|_1 \leq \|x\|_2 \leq C\|x\|_1\) 对所有 \(x \in X\) 成立。

\textbf{定理
72.1}：有限维赋范空间上的所有范数都等价，且有限维赋范空间总是 Banach
空间。

\emph{证明}：设 \(X\) 是 \(n\) 维赋范空间，取基
\(\{e_1, \ldots, e_n\}\)。定义欧氏范数
\(\|\sum a_i e_i\|_2 = (\sum |a_i|^2)^{1/2}\)。对任意范数
\(\|\cdot\|\)，由三角不等式和 Cauchy-Schwarz 得
\(\|x\| \leq C\|x\|_2\)。函数 \(x \mapsto \|x\|\) 在单位球面
\(S = \{x : \|x\|_2 = 1\}\) 上连续且严格正，由紧致性得
\(c = \inf_S \|x\| > 0\)，故 \(\|x\| \geq c\|x\|_2\)。∎

\subsubsection{72.2
有界线性算子}\label{ux6709ux754cux7ebfux6027ux7b97ux5b50}

\textbf{定义 72.4}（有界线性算子）：设 \(X, Y\) 是赋范空间。线性算子
\(T: X \to Y\) 称为\textbf{有界的}，如果存在 \(M \geq 0\) 使得
\(\|Tx\| \leq M\|x\|\) 对所有 \(x \in X\) 成立。\(T\)
的\textbf{算子范数}定义为

\[\|T\| = \sup_{\|x\| \leq 1} \|Tx\| = \sup_{\|x\| = 1} \|Tx\|\]

\textbf{命题 72.2}：线性算子 \(T: X \to Y\) 有界当且仅当 \(T\)
连续，当且仅当 \(T\) 在 \(0\) 处连续。

\emph{证明}：若有界，则
\(\|Tx - Ty\| = \|T(x-y)\| \leq \|T\| \|x-y\|\)，故 Lipschitz 连续。若在
\(0\) 处连续，则存在 \(\delta > 0\) 使
\(\|x\| < \delta \Rightarrow \|Tx\| < 1\)。对任意
\(x\)，\(\|T(\delta x/(2\|x\|))\| < 1\)，故
\(\|Tx\| < (2/\delta)\|x\|\)。∎

\textbf{定义 72.5}（算子空间）：\(X\) 到 \(Y\)
的所有有界线性算子构成向量空间，记作
\(\mathcal{B}(X, Y)\)。配以算子范数，\(\mathcal{B}(X, Y)\)
是赋范空间。若 \(Y\) 是 Banach 空间，则 \(\mathcal{B}(X, Y)\) 也是
Banach 空间。特别地，\(X^* = \mathcal{B}(X, \mathbb{K})\) 称为 \(X\)
的\textbf{对偶空间}（或共轭空间）。

\textbf{定理 72.3}（Banach-Steinhaus 一致有界原理）：设 \(X\) 是 Banach
空间，\(Y\)
是赋范空间，\(\{T_\alpha\}_{\alpha \in A} \subseteq \mathcal{B}(X, Y)\)。若对每个
\(x \in X\)，\(\sup_{\alpha} \|T_\alpha x\| < \infty\)，则
\(\sup_{\alpha} \|T_\alpha\| < \infty\)。

\emph{证明}：使用 Baire 纲定理。对每个 \(n \in \mathbb{N}\)，定义
\(F_n = \{x \in X : \sup_\alpha \|T_\alpha x\| \leq n\}\)。\(F_n\)
是闭集，且 \(X = \bigcup_{n} F_n\)。由 Baire 纲定理，存在 \(n_0\) 使
\(F_{n_0}\) 有内点。设 \(B(x_0, r) \subseteq F_{n_0}\)。对任意
\(\|x\| \leq r\)，\(x_0 + x \in F_{n_0}\)，故
\(\|T_\alpha x\| \leq \|T_\alpha(x_0 + x)\| + \|T_\alpha x_0\| \leq 2n_0\)。因此
\(\|T_\alpha\| \leq 2n_0/r\) 对所有 \(\alpha\) 成立。∎

\subsubsection{72.3 Hahn-Banach 定理}\label{hahn-banach-ux5b9aux7406}

\textbf{定理 72.4}（Hahn-Banach 延拓定理，实情形）：设 \(X\)
是实向量空间，\(p: X \to \mathbb{R}\) 满足 \(p(x+y) \leq p(x) + p(y)\)
和 \(p(tx) = tp(x)\) 对 \(t \geq 0\)（次线性泛函）。设 \(M \subseteq X\)
是子空间，\(f: M \to \mathbb{R}\) 是线性泛函满足 \(f(x) \leq p(x)\)
对所有 \(x \in M\)。则存在 \(f\) 的线性延拓 \(F: X \to \mathbb{R}\) 使得
\(F(x) \leq p(x)\) 对所有 \(x \in X\) 成立。

\emph{证明}：使用 Zorn 引理。考虑所有保持 \(f\) 且满足
\(F(x) \leq p(x)\) 的延拓构成的偏序集（按包含关系）。由 Zorn
引理存在极大延拓 \(F\)，定义域为 \(M'\)。若 \(M' \neq X\)，取
\(x_0 \notin M'\)。需定义 \(F(x_0)\) 使得对任意 \(y, z \in M'\)：

\[F(y) + F(x_0) \leq p(y + x_0), \quad F(z) - F(x_0) \leq p(z - x_0)\]

这等价于 \(F(z) - p(z - x_0) \leq F(x_0) \leq p(y + x_0) - F(y)\)。由
\(f(y) + f(z) = f(y+z) \leq p(y+z) \leq p(y+x_0) + p(z-x_0)\)，知上确界不大于下确界，故可取
\(F(x_0)\) 为中间值。在 \(M' \oplus \mathbb{R}x_0\)
上延拓，与极大性矛盾。∎

\textbf{定理 72.5}（Hahn-Banach 定理，赋范空间版本）：设 \(X\)
是赋范空间，\(M \subseteq X\) 是子空间，\(f \in M^*\)。则存在
\(F \in X^*\) 使得 \(F|_M = f\) 且 \(\|F\| = \|f\|\)。

\emph{证明}：取 \(p(x) = \|f\| \|x\|\)（次线性），则
\(|f(x)| \leq p(x)\) 在 \(M\) 上。由 Hahn-Banach 定理存在 \(F\) 满足
\(F(x) \leq p(x)\)。由 \(F(-x) = -F(x) \leq p(-x) = p(x)\) 得
\(|F(x)| \leq p(x) = \|f\| \|x\|\)，故
\(\|F\| \leq \|f\|\)。反向不等式显然。∎

\textbf{推论 72.6}（Hahn-Banach 推论）： 1. 对任意
\(x \in X, x \neq 0\)，存在 \(f \in X^*\) 使得 \(\|f\| = 1\) 且
\(f(x) = \|x\|\)。 2. \(X^*\) 分离 \(X\) 中的点：若 \(x \neq y\)，则存在
\(f \in X^*\) 使 \(f(x) \neq f(y)\)。 3.
\(\|x\| = \sup_{\|f\| \leq 1} |f(x)|\)。

\subsubsection{72.4
开映射定理与闭图像定理}\label{ux5f00ux6620ux5c04ux5b9aux7406ux4e0eux95edux56feux50cfux5b9aux7406}

\textbf{定理 72.7}（开映射定理）：设 \(X, Y\) 是 Banach
空间，\(T \in \mathcal{B}(X, Y)\) 是满射。则 \(T\)
是开映射（即将开集映为开集）。特别地，若 \(T\) 是连续双射，则 \(T^{-1}\)
连续，即 \(T\) 是同胚（从而 \(T\) 是 Banach 空间同构）。

\emph{证明}：证 \(T(B_X(0, 1))\) 包含 \(B_Y(0, r)\) 对某个
\(r > 0\)。由满射性和 Baire 纲定理，存在 \(n\) 使
\(\overline{T(B_X(0, n))}\) 有内点。则 \(\overline{T(B_X(0, 1))}\)
包含某球 \(B_Y(y_0, \delta)\)。由对称性得
\(B_Y(0, \delta) \subseteq \overline{T(B_X(0, 2))}\)。再由完备性迭代可得
\(B_Y(0, \delta/2) \subseteq T(B_X(0, 1))\)。∎

\textbf{定理 72.8}（闭图像定理）：设 \(X, Y\) 是 Banach
空间，\(T: X \to Y\) 是线性算子。则 \(T\) 有界当且仅当 \(T\) 的图像
\(\Gamma(T) = \{(x, Tx) : x \in X\}\) 是 \(X \times Y\)（配以范数
\(\|(x, y)\| = \|x\| + \|y\|\)）中的闭集。

\emph{证明}：(\(\Rightarrow\)) 若 \(T\) 有界，设
\((x_n, Tx_n) \to (x, y)\)，则 \(x_n \to x\)，由连续性
\(Tx_n \to Tx\)，故 \(y = Tx\)，\(\Gamma(T)\) 闭。

(\(\Leftarrow\)) 若 \(\Gamma(T)\) 闭，则 \(\Gamma(T)\) 是 Banach 空间
\(X \times Y\) 的闭子空间，从而是 Banach 空间。投影
\(\pi_1: \Gamma(T) \to X\)，\(\pi_1(x, Tx) = x\)
是连续双射。由开映射定理，\(\pi_1^{-1}\) 连续，故
\(\|x\| + \|Tx\| = \|(x, Tx)\| \leq C\|x\|\)，从而
\(\|Tx\| \leq C\|x\|\)，\(T\) 有界。∎

\bigskip\hrule\bigskip

\subsection{第73章：Hilbert
空间}\label{ux7b2c73ux7ae0hilbert-ux7a7aux95f4}

Hilbert 空间是带有内积的完备空间，其几何结构比一般 Banach
空间更丰富，有正交投影、正交基等概念。

\subsubsection{73.1 内积空间与 Hilbert
空间}\label{ux5185ux79efux7a7aux95f4ux4e0e-hilbert-ux7a7aux95f4}

\textbf{定义 73.1}（内积空间）：\(\mathbb{K}\) 上的向量空间 \(H\)
配以一个\textbf{内积}
\(\langle \cdot, \cdot \rangle: H \times H \to \mathbb{K}\) 满足：

\begin{enumerate}
\def\labelenumi{\arabic{enumi}.}
\tightlist
\item
  \(\langle x, x \rangle \geq 0\)，且
  \(\langle x, x \rangle = 0 \Leftrightarrow x = 0\)（正定性）
\item
  \(\langle x, y \rangle = \overline{\langle y, x \rangle}\)（共轭对称性）
\item
  \(\langle \alpha x + \beta y, z \rangle = \alpha \langle x, z \rangle + \beta \langle y, z \rangle\)（对第一变元的线性性）
\end{enumerate}

内积诱导范数 \(\|x\| = \sqrt{\langle x, x \rangle}\)。若
\((H, \langle \cdot, \cdot \rangle)\) 在该范数下完备，则称为
\textbf{Hilbert 空间}。

\textbf{定理 73.1}（Cauchy-Schwarz
不等式）：\(|\langle x, y \rangle| \leq \|x\| \|y\|\)，等号成立当且仅当
\(x\) 与 \(y\) 线性相关。

\emph{证明}：对
\(\lambda \in \mathbb{K}\)，\(0 \leq \|x - \lambda y\|^2 = \|x\|^2 - 2\Re(\lambda \langle y, x \rangle) + |\lambda|^2 \|y\|^2\)。取
\(\lambda = \langle x, y \rangle / \|y\|^2\)（\(y \neq 0\)
时）可得不等式。∎

\textbf{定理 73.2}（平行四边形法则与极化恒等式）： -
平行四边形法则：\(\|x + y\|^2 + \|x - y\|^2 = 2(\|x\|^2 + \|y\|^2)\) -
极化恒等式（实情形）：\(\langle x, y \rangle = \frac{1}{4}(\|x + y\|^2 - \|x - y\|^2)\)
-
极化恒等式（复情形）：\(\langle x, y \rangle = \frac{1}{4}\sum_{k=0}^3 i^k \|x + i^k y\|^2\)

\textbf{定理 73.3}（Jordan-von Neumann 定理）：赋范空间
\((X, \|\cdot\|)\)
的范数由内积诱导当且仅当范数满足平行四边形法则。此时内积可由极化恒等式恢复。

\textbf{例 73.1}（Hilbert 空间举例）： - \(\mathbb{R}^n\) 和
\(\mathbb{C}^n\) 配以标准内积 -
\(\ell^2\)：\(\langle (x_n), (y_n) \rangle = \sum_{n=1}^\infty x_n \overline{y_n}\)
-
\(L^2(\Omega)\)：\(\langle f, g \rangle = \int_\Omega f \overline{g} d\mu\)
- \(\ell^p\) 对 \(p \neq 2\) 是 Banach 空间但不是 Hilbert
空间（不满足平行四边形法则）

\subsubsection{73.2 正交性}\label{ux6b63ux4ea4ux6027-1}

\textbf{定义 73.2}（正交）：\(x, y \in H\) 称为\textbf{正交的}，记作
\(x \perp y\)，如果 \(\langle x, y \rangle = 0\)。子集 \(A \subseteq H\)
的\textbf{正交补}为
\(A^\perp = \{x \in H : \langle x, a \rangle = 0, \; \forall a \in A\}\)。

\textbf{命题 73.4}（勾股定理）：若 \(x \perp y\)，则
\(\|x + y\|^2 = \|x\|^2 + \|y\|^2\)。

\textbf{定理 73.5}（正交投影定理）：设 \(C \subseteq H\)
是非空闭凸集。则对每个 \(x \in H\)，存在唯一的 \(y \in C\) 使得
\(\|x - y\| = \operatorname{dist}(x, C) = \inf_{z \in C} \|x - z\|\)。\(y\)
称为 \(x\) 到 \(C\) 的\textbf{最佳逼近}。

特别地，若 \(M \subseteq H\) 是闭子空间，则 \(H = M \oplus M^\perp\)，且
\(M^{\perp\perp} = M\)。

\emph{证明}：存在性：取极小化序列 \((y_n) \subseteq C\) 使
\(\|x - y_n\| \to d = \operatorname{dist}(x, C)\)。由平行四边形法则：

\[\|y_n - y_m\|^2 = 2\|y_n - x\|^2 + 2\|y_m - x\|^2 - \|y_n + y_m - 2x\|^2\]

由 \(C\) 凸，\((y_n + y_m)/2 \in C\)，故
\(\|y_n + y_m - 2x\|^2 \geq 4d^2\)。当 \(n, m \to \infty\)
时，\(\|y_n - y_m\| \to 0\)，\((y_n)\) 是 Cauchy 序列，由完备性收敛到某
\(y\)，由 \(C\) 闭知 \(y \in C\)。

唯一性：若 \(y, y'\) 都是最佳逼近，由平行四边形法则得
\(\|y - y'\|^2 \leq 4d^2 - 4d^2 = 0\)。

子空间情形：\(M^\perp\) 是闭子空间。对 \(x \in H\)，取 \(Px\) 为 \(x\)
到 \(M\) 的投影，则 \(x - Px \in M^\perp\)，故
\(x = Px + (x - Px) \in M + M^\perp\)。若 \(x \in M \cap M^\perp\)，则
\(\langle x, x \rangle = 0\)，故 \(x = 0\)。∎

\subsubsection{73.3 Riesz
表示定理与正交基}\label{riesz-ux8868ux793aux5b9aux7406ux4e0eux6b63ux4ea4ux57fa}

\textbf{定理 73.6}（Riesz 表示定理）：设 \(H\) 是 Hilbert
空间。对每个连续线性泛函 \(f \in H^*\)，存在唯一的 \(y \in H\) 使得

\[f(x) = \langle x, y \rangle, \quad \forall x \in H\]

且 \(\|f\| = \|y\|\)。映射 \(y \mapsto \langle \cdot, y \rangle\) 是
\(H\) 到 \(H^*\) 的共轭线性等距同构（实情形为线性）。

\emph{证明}：若 \(f = 0\)，取 \(y = 0\)。否则，\(\ker f\) 是 \(H\)
的真闭子空间。取 \(z \in (\ker f)^\perp\) 且 \(\|z\| = 1\)。对任意
\(x \in H\)，\(x - f(x)z/f(z) \in \ker f\)，故
\(\langle x - f(x)z/f(z), z \rangle = 0\)，即
\(f(x) = \langle x, \overline{f(z)}z \rangle\)。取
\(y = \overline{f(z)}z\)。唯一性：若
\(\langle x, y_1 \rangle = \langle x, y_2 \rangle\) 对所有 \(x\)，则
\(\langle x, y_1 - y_2 \rangle = 0\)，取 \(x = y_1 - y_2\) 得
\(y_1 = y_2\)。∎

\textbf{定义 73.3}（正交规范集）：\(H\) 的子集
\(\{e_\alpha\}_{\alpha \in A}\) 称为\textbf{正交规范的}，如果
\(\langle e_\alpha, e_\beta \rangle = \delta_{\alpha\beta}\)。

\textbf{定义 73.4}（正交规范基）：正交规范集 \(\{e_\alpha\}\) 称为 \(H\)
的\textbf{正交规范基}（或 Hilbert 基），如果
\(\overline{\operatorname{span}}\{e_\alpha\} = H\)（即其线性张成稠密）。

\textbf{定理 73.7}（正交规范基的存在性）：每个非零 Hilbert
空间都有正交规范基。任意两个正交规范基具有相同的基数，称为 \(H\) 的
\textbf{Hilbert 维数}。

\emph{证明}：Zorn
引理：所有正交规范集按包含关系构成偏序集，极大元即为正交规范基。∎

\textbf{定理 73.8}（Parseval 恒等式与 Fourier 展开）：设
\(\{e_\alpha\}_{\alpha \in A}\) 是 \(H\) 的正交规范基。则对任意
\(x \in H\)：

\begin{enumerate}
\def\labelenumi{\arabic{enumi}.}
\tightlist
\item
  \textbf{Bessel
  不等式}：\(\sum_{\alpha} |\langle x, e_\alpha \rangle|^2 \leq \|x\|^2\)（事实上只有可数个系数非零）
\item
  \textbf{Fourier
  展开}：\(x = \sum_{\alpha} \langle x, e_\alpha \rangle e_\alpha\)（在
  \(H\) 的范数下收敛）
\item
  \textbf{Parseval
  恒等式}：\(\|x\|^2 = \sum_{\alpha} |\langle x, e_\alpha \rangle|^2\)
\end{enumerate}

\emph{证明}：对任意有限子集
\(F \subseteq A\)，\(\|x - \sum_{\alpha \in F} \langle x, e_\alpha \rangle e_\alpha\|^2 = \|x\|^2 - \sum_{\alpha \in F} |\langle x, e_\alpha \rangle|^2 \geq 0\)，故
Bessel
不等式成立。由此知可数个系数非零。由完备性和极大性，\(\|x - \sum_{n=1}^N \langle x, e_n \rangle e_n\| \to 0\)。∎

\textbf{定理 73.9}（可分 Hilbert 空间的分类）：每个无穷维可分 Hilbert
空间等距同构于 \(\ell^2\)（配以标准正交基 \(\{e_n\}\)）。

\subsubsection{73.4 伴随算子}\label{ux4f34ux968fux7b97ux5b50-1}

\textbf{定义 73.5}（伴随算子）：设 \(H_1, H_2\) 为 Hilbert
空间，\(T \in \mathcal{B}(H_1, H_2)\)。\(T\) 的\textbf{伴随算子}
\(T^* \in \mathcal{B}(H_2, H_1)\) 由下式唯一确定：

\[\langle Tx, y \rangle_{H_2} = \langle x, T^* y \rangle_{H_1}, \quad \forall x \in H_1, y \in H_2\]

\(T^*\) 的存在唯一性由 Riesz 表示定理保证：对固定
\(y\)，\(x \mapsto \langle Tx, y \rangle\) 是连续线性泛函，由 Riesz
表示存在唯一 \(T^*y\)。

\textbf{定义 73.6}（自伴算子、正规算子、酉算子）：设
\(T \in \mathcal{B}(H)\)。 - \(T\) 称为\textbf{自伴的}（或
Hermitian），如果 \(T^* = T\)。 - \(T\) 称为\textbf{正规的}，如果
\(T^*T = TT^*\)。 - \(T\) 称为\textbf{酉算子}，如果
\(T^*T = TT^* = I\)（即 \(T\) 是等距满射）。

\textbf{命题 73.10}（伴随算子的性质）： 1.
\((T^*)^* = T\)，\(\|T^*\| = \|T\|\)，\(\|T^*T\| = \|T\|^2\)（C*-等式）
2.
\((S+T)^* = S^* + T^*\)，\((\alpha T)^* = \overline{\alpha} T^*\)，\((ST)^* = T^* S^*\)
3.
\(\ker T^* = (\operatorname{im} T)^\perp\)，\(\overline{\operatorname{im} T^*} = (\ker T)^\perp\)

\bigskip\hrule\bigskip

\subsection{第74章：对偶空间与弱拓扑}\label{ux7b2c74ux7ae0ux5bf9ux5076ux7a7aux95f4ux4e0eux5f31ux62d3ux6251}

本章研究 Banach
空间的对偶空间、二次对偶和弱拓扑，这些是研究无穷维空间紧致性的关键工具。

\subsubsection{74.1
对偶空间与二次对偶}\label{ux5bf9ux5076ux7a7aux95f4ux4e0eux4e8cux6b21ux5bf9ux5076}

\textbf{定义 74.1}（对偶空间）：\(X^* = \mathcal{B}(X, \mathbb{K})\) 是
Banach 空间（即使 \(X\) 不完备），配以范数
\(\|f\| = \sup_{\|x\| \leq 1} |f(x)|\)。

\textbf{定义 74.2}（二次对偶与自然嵌入）：\(X^{**} = (X^*)^*\) 是 \(X\)
的\textbf{二次对偶}。自然嵌入 \(J: X \to X^{**}\) 定义为
\(J(x)(f) = f(x)\)（\(f \in X^*\)）。由 Hahn-Banach
推论，\(\|J(x)\| = \|x\|\)，故 \(J\) 是等距嵌入。

\textbf{例 74.1}（对偶空间的例子）： - \((\ell^p)^* \cong \ell^q\)，其中
\(1/p + 1/q = 1\)，\(1 < p < \infty\)（\(q = \infty\) 当 \(p = 1\)） -
\((L^p(\Omega))^* \cong L^q(\Omega)\)（\(1 < p < \infty\)） -
\((c_0)^* \cong \ell^1\)，\((\ell^1)^* \cong \ell^\infty\) -
\(C([a, b])^* \cong \mathcal{M}([a, b])\)（Riesz
表示定理------测度版本）

\textbf{定义 74.3}（自反空间）：Banach 空间 \(X\)
称为\textbf{自反的}，如果自然嵌入 \(J: X \to X^{**}\)
是满射（从而是等距同构）。

\textbf{命题 74.1}：自反空间的闭子空间是自反的。\(X\) 自反当且仅当
\(X^*\) 自反。

\textbf{例 74.2}：\(\ell^p\) 和 \(L^p\) 对 \(1 < p < \infty\)
是自反的。\(\ell^1\)、\(\ell^\infty\)、\(L^1\)、\(L^\infty\)、\(c_0\)、\(C([a, b])\)
都不是自反的。

\subsubsection{74.2 弱拓扑与弱*
拓扑}\label{ux5f31ux62d3ux6251ux4e0eux5f31-ux62d3ux6251}

在无穷维 Banach
空间中，有界闭集不一定紧致（例如单位球面不是紧致的）。弱拓扑和弱*
拓扑提供了更粗的拓扑，使得某些集合变得紧致。

\textbf{定义 74.4}（弱拓扑）：\(X\) 上的\textbf{弱拓扑}
\(\sigma(X, X^*)\) 是使所有 \(f \in X^*\) 连续的最弱拓扑。其子基为
\(\{f^{-1}(U) : f \in X^*, U \subseteq \mathbb{K} \text{ 开}\}\)。

网 \((x_\alpha)\) 弱收敛到 \(x\)（记作
\(x_\alpha \rightharpoonup x\)）当且仅当 \(f(x_\alpha) \to f(x)\) 对所有
\(f \in X^*\) 成立。

\textbf{定义 74.5}（弱* 拓扑）：\(X^*\) 上的\textbf{弱* 拓扑}
\(\sigma(X^*, X)\) 是使所有 \(J(x)\)（\(x \in X\)）连续的最弱拓扑。即
\(f_\alpha \xrightarrow{w^*} f\) 当且仅当 \(f_\alpha(x) \to f(x)\)
对所有 \(x \in X\)。

\textbf{命题 74.2}：弱拓扑和弱* 拓扑是 Hausdorff 拓扑。若 \(X\) 自反，则
\(\sigma(X^*, X) = \sigma(X^*, X^{**})\)，即弱* 拓扑与弱拓扑一致。

\textbf{定理 74.3}（Banach-Alaoglu 定理）：\(X^*\) 中的闭单位球
\(B_{X^*} = \{f \in X^* : \|f\| \leq 1\}\) 在弱* 拓扑下是紧致的。

\emph{证明}：考虑映射
\(\Phi: B_{X^*} \to \prod_{x \in X} [-\|x\|, \|x\|]\)，\(\Phi(f) = (f(x))_{x \in X}\)。\(\Phi\)
是单射且是到其像的同胚（\(B_{X^*}\) 赋予弱*
拓扑，乘积空间赋予积拓扑）。由 Tychonoff
定理，乘积空间紧致。\(\Phi(B_{X^*})\) 作为闭子集是紧致的。∎

\textbf{定理 74.4}（Goldstine 定理）：\(J(B_X)\) 在 \(B_{X^{**}}\) 中弱*
稠密。即单位球的像在二次对偶的单位球中稠密。

\textbf{推论 74.5}：\(X\) 自反当且仅当 \(B_X\) 在弱拓扑下紧致。

\emph{证明}：若 \(X\) 自反，\(J(B_X) = B_{X^{**}}\)，由 Banach-Alaoglu
定理在弱* 拓扑下紧致，而 \(X\) 自反时弱* 拓扑等于弱拓扑。反之，若
\(B_X\) 弱紧致，则 \(J(B_X)\) 弱* 紧致从而弱* 闭，由 Goldstine 定理得
\(J(B_X) = B_{X^{**}}\)。∎

\subsubsection{74.3 凸集的分离与 Krein-Milman
定理}\label{ux51f8ux96c6ux7684ux5206ux79bbux4e0e-krein-milman-ux5b9aux7406}

\textbf{定理 74.6}（凸集分离定理 / Hahn-Banach 定理的几何形式）：设
\(A, B\) 是 Banach 空间的非空不交凸集，\(A\) 是开集。则存在
\(f \in X^*\) 和 \(\alpha \in \mathbb{R}\) 使得
\(\Re f(a) < \alpha \leq \Re f(b)\) 对所有 \(a \in A, b \in B\)。

\textbf{定理 74.7}（Mazur 定理）：设
\(x_n \rightharpoonup x\)（弱收敛）。则存在凸组合序列
\(\sum_{k=1}^{N_n} \lambda_k^{(n)} x_k\)（\(\lambda_k^{(n)} \geq 0, \sum \lambda_k^{(n)} = 1\)）强收敛到
\(x\)。

\emph{证明}：\(x\) 属于 \(\{x_n\}\)
的弱闭包的凸包。由凸集的弱闭等于强闭，\(x\)
属于强闭凸包，故存在凸组合序列强收敛。∎

\textbf{定义 74.6}（极点）：设 \(K\) 是向量空间的凸集。\(x \in K\) 称为
\(K\) 的\textbf{极点}（extreme point），如果 \(x\) 不能表示为 \(K\)
中两个不同点的严格凸组合，即若
\(x = ty + (1-t)z\)（\(t \in (0, 1), y, z \in K\)），则
\(y = z = x\)。\(K\) 的极点集记作 \(\operatorname{ext}(K)\)。

\textbf{定理 74.8}（Krein-Milman 定理）：设 \(K\)
是局部凸空间中的非空紧致凸集。则
\(K = \overline{\operatorname{conv}}(\operatorname{ext}(K))\)，即 \(K\)
是其极点的闭凸包。

\emph{说明}：此定理及其在 Banach-Alaoglu 定理的应用（\(B_{X^*}\) 是弱*
紧致凸集）是泛函分析中许多存在性证明的基础。

\bigskip\hrule\bigskip

\subsection{第75章：紧算子与谱理论}\label{ux7b2c75ux7ae0ux7d27ux7b97ux5b50ux4e0eux8c31ux7406ux8bba}

紧算子是有限维线性算子到无穷维的最自然推广，其谱理论具有特别简洁优美的结构。

\subsubsection{75.1 紧算子}\label{ux7d27ux7b97ux5b50}

\textbf{定义 75.1}（紧算子）：设 \(X, Y\) 是 Banach 空间。线性算子
\(T: X \to Y\) 称为\textbf{紧算子}（或全连续算子），如果
\(T(B_X)\)（单位球的像）在 \(Y\)
中是相对紧致的（即其闭包是紧致的）。所有紧算子构成的集合记作
\(\mathcal{K}(X, Y)\)。

\textbf{命题 75.1}（紧算子的等价刻画）：\(T\) 紧当且仅当对 \(X\)
中任意有界序列 \((x_n)\)，\((Tx_n)\) 有收敛子列。

\textbf{命题 75.2}（紧算子的基本性质）： 1.
紧算子是有限秩算子的极限（在算子范数下）：\(\mathcal{K}(X, Y) = \overline{\mathcal{F}(X, Y)}\)，其中
\(\mathcal{F}(X, Y)\) 是有限秩算子（像有限维）。 2.
\(\mathcal{K}(X)\)（\(X\) 到自身的紧算子）是 \(\mathcal{B}(X)\)
的闭双边理想。 3. 若 \(T \in \mathcal{K}(X)\) 且
\(S \in \mathcal{B}(X)\)，则 \(TS\) 和 \(ST\) 是紧算子。 4.
紧算子的伴随是紧算子。

\textbf{定义 75.2}（Fredholm 算子）：\(T \in \mathcal{B}(X, Y)\) 称为
\textbf{Fredholm 算子}，如果 \(\dim \ker T < \infty\) 且
\(\dim \operatorname{coker} T = \dim(Y / \operatorname{im} T) < \infty\)。此时
\(T\) 的\textbf{指标}定义为
\(\operatorname{ind}(T) = \dim \ker T - \dim \operatorname{coker} T\)。

\textbf{定理 75.3}（Fredholm 择一 / Riesz-Schauder 理论）：设 \(X\) 是
Banach 空间，\(T \in \mathcal{K}(X)\)。则：

\begin{enumerate}
\def\labelenumi{\arabic{enumi}.}
\tightlist
\item
  \(I - T\) 是 Fredholm 算子，且 \(\operatorname{ind}(I - T) = 0\)。
\item
  \(\dim \ker(I - T) = \dim \ker(I - T^*)\)。
\item
  方程 \((I - T)x = y\) 有解当且仅当 \(y \perp \ker(I - T^*)\)（即 \(y\)
  与 \((I - T^*)\) 的核中所有泛函正交）。
\item
  \((I - T)\) 是单射当且仅当 \((I - T)\) 是满射。
\end{enumerate}

\emph{证明思路}：利用 Riesz 引理证明 \(I - T\)
的像是闭的，且核为有限维。详细证明需构造 \(\ker(I - T)^n\) 的升链。∎

\subsubsection{75.2 谱理论初步}\label{ux8c31ux7406ux8bbaux521dux6b65}

\textbf{定义 75.3}（谱与预解集）：设 \(X\) 是复 Banach
空间，\(T \in \mathcal{B}(X)\)。 - \(T\) 的\textbf{预解集}
\(\rho(T) = \{\lambda \in \mathbb{C} : \lambda I - T \text{ 是双射}\}\)（此时由开映射定理，\((\lambda I - T)^{-1}\)
有界）。 - \(T\) 的\textbf{谱}
\(\sigma(T) = \mathbb{C} \setminus \rho(T)\)。

谱可进一步分为： - \textbf{点谱} \(\sigma_p(T)\)：\(\lambda I - T\)
不是单射（即 \(\lambda\) 是特征值） - \textbf{连续谱}
\(\sigma_c(T)\)：\(\lambda I - T\) 是单射但不满射，且像稠密 -
\textbf{剩余谱} \(\sigma_r(T)\)：\(\lambda I - T\)
是单射但不满射，且像不稠密

\textbf{定理 75.4}（谱的基本性质）： 1. \(\sigma(T)\) 是 \(\mathbb{C}\)
中的非空紧致集。 2.
\(\sigma(T) \subseteq \{\lambda \in \mathbb{C} : |\lambda| \leq \|T\|\}\)。
3.
\textbf{谱半径公式}：\(r(T) = \lim_{n \to \infty} \|T^n\|^{1/n} = \sup_{\lambda \in \sigma(T)} |\lambda|\)。

\emph{证明}(非空性)：(1) 预解函数 \(R(\lambda) = (\lambda I - T)^{-1}\)
在 \(\rho(T)\) 上解析。若 \(\sigma(T) = \emptyset\)，则 \(R(\lambda)\)
是整函数且有界（当 \(|\lambda| \to \infty\) 时
\(R(\lambda) \to 0\)），由 Liouville 定理
\(R(\lambda) \equiv 0\)，矛盾。∎

\textbf{定理 75.5}（紧算子的谱定理 / Riesz-Schauder 定理）：设
\(T \in \mathcal{K}(X)\)（\(X\) 是无穷维复 Banach 空间）。则：

\begin{enumerate}
\def\labelenumi{\arabic{enumi}.}
\tightlist
\item
  \(\sigma(T) \setminus \{0\}\) 仅由特征值构成（即 \(0\)
  是谱中唯一可能的非特征值点）。
\item
  \(\sigma(T) \setminus \{0\}\) 是有限集或以 \(0\) 为唯一聚点。
\item
  每个非零特征值 \(\lambda\) 的几何重数 \(\dim \ker(\lambda I - T)\)
  有限。
\end{enumerate}

\subsubsection{75.3 Hilbert
空间上的自伴紧算子}\label{hilbert-ux7a7aux95f4ux4e0aux7684ux81eaux4f34ux7d27ux7b97ux5b50}

\textbf{定理 75.6}（自伴紧算子的谱定理）：设 \(H\) 是 Hilbert
空间，\(T \in \mathcal{K}(H)\) 是自伴紧算子。则存在 \(H\) 的正交规范基
\(\{e_n\}\) 和实数 \(\lambda_n \to 0\) 使得

\[Tx = \sum_{n} \lambda_n \langle x, e_n \rangle e_n\]

即 \(T\) 可以对角化，其特征向量构成正交规范基。

\emph{证明}：首先证明 \(\|T\|\) 或 \(-\|T\|\) 是 \(T\)
的特征值。取单位向量序列 \((x_n)\) 使
\(|\langle Tx_n, x_n \rangle| \to \|T\|\)（由自伴性，\(\|T\| = \sup_{\|x\|=1} |\langle Tx, x \rangle|\)）。由紧性，\((Tx_n)\)
有收敛子列，可得特征向量 \(e_1\) 使
\(Te_1 = \lambda_1 e_1\)，\(|\lambda_1| = \|T\|\)。

然后考虑 \(H_1 = \{e_1\}^\perp\)。\(T\) 限制在 \(H_1\)
上仍是自伴紧算子。归纳构造正交特征向量列 \(\{e_n\}\) 和特征值
\(\lambda_n\)，\(|\lambda_{n+1}| \leq |\lambda_n|\)。由紧性，\(\lambda_n \to 0\)。

最后证明 \(\{e_n\}\) 是基：若 \(x \perp e_n\) 对所有 \(n\)，则
\(Tx = 0\)。同时 \(T\) 在 \(\{e_n\}^\perp\) 上为零。故任意 \(x\)
可展开为
\(\sum \langle x, e_n \rangle e_n + (x - \sum \langle x, e_n \rangle e_n)\)，后者在
\(\ker T\) 中。∎

\textbf{推论 75.7}（Hilbert-Schmidt 展开定理）：设 \(T\)
是紧自伴算子。则对任意 \(x \in H\)，

\[Tx = \sum_n \lambda_n \langle x, e_n \rangle e_n\]

其中 \(\lambda_n\) 是 \(T\)
的非零特征值（按绝对值递减，重复按重数），\(\{e_n\}\)
是相应的正交规范特征向量。

\textbf{定理 75.8}（极小-极大原理 / Courant-Fischer 定理）：设 \(T\)
是紧自伴算子，\(\lambda_1^+ \geq \lambda_2^+ \geq \cdots > 0\)
是其正特征值。则

\[\lambda_n^+ = \min_{\dim V = n-1} \max_{\substack{x \perp V \\ \|x\| = 1}} \langle Tx, x \rangle = \max_{\dim W = n} \min_{\substack{x \in W \\ \|x\| = 1}} \langle Tx, x \rangle\]

\textbf{定理 75.9}（泛函演算 / 连续函数演算）：设 \(T\) 是 Hilbert
空间上的自伴算子（不一定紧），\(\sigma(T) \subseteq [m, M]\)。则存在唯一的同态
\(\Phi: C([m, M]) \to \mathcal{B}(H)\) 使得
\(\Phi(1) = I\)，\(\Phi(\operatorname{id}) = T\)，且 \(\Phi\)
是等距的（\(\|\Phi(f)\| = \|f\|_\infty\)）。\(\Phi(f)\) 记作 \(f(T)\)。

\emph{证明思路}：对多项式 \(p\)，定义 \(p(T)\) 自然。对连续函数
\(f\)，由 Weierstrass 逼近定理用多项式逼近，再由
\(\|p(T)\| = \sup_{\lambda \in \sigma(T)} |p(\lambda)|\)
保证延拓的良定义。∎

\bigskip\hrule\bigskip

\subsection{第76章：C*-代数初步}\label{ux7b2c76ux7ae0c-ux4ee3ux6570ux521dux6b65}

C*-代数是 Hilbert
空间上有界算子构成的、在伴随和范数拓扑下封闭的代数，是量子力学和非交换几何的数学语言。

\subsubsection{76.1 Banach 代数与
C*-代数}\label{banach-ux4ee3ux6570ux4e0e-c-ux4ee3ux6570}

\textbf{定义 76.1}（Banach 代数）：\(\mathbb{C}\) 上的一个
\textbf{Banach 代数}是一个结合 \(\mathbb{C}\)-代数 \(\mathcal{A}\)
配以范数 \(\|\cdot\|\)，使得 \((\mathcal{A}, \|\cdot\|)\) 是 Banach
空间，且满足 \(\|ab\| \leq \|a\| \|b\|\) 对所有
\(a, b \in \mathcal{A}\)。

若 \(\mathcal{A}\) 有单位元 \(1\) 且 \(\|1\| = 1\)，则称 \(\mathcal{A}\)
为\textbf{有单位元的 Banach 代数}。

\textbf{定义 76.2}（对合）：\(\mathcal{A}\)
上的一个\textbf{对合}是一个映射
\(*: \mathcal{A} \to \mathcal{A}\)，\(a \mapsto a^*\)，满足： 1.
\((a^*)^* = a\) 2. \((a + b)^* = a^* + b^*\) 3.
\((\lambda a)^* = \overline{\lambda} a^*\) 4. \((ab)^* = b^* a^*\)

\textbf{定义 76.3}（C*-代数）：一个
\textbf{C\emph{-代数\textbf{是一个带有对合的 Banach 代数
\(\mathcal{A}\)，满足 }C}-等式}：

\[\|a^* a\| = \|a\|^2, \quad \forall a \in \mathcal{A}\]

\textbf{例 76.1}（C\emph{-代数举例）： - \(\mathcal{B}(H)\)（\(H\) 为
Hilbert 空间），对合为伴随算子。C}-等式已由命题 73.10 给出。 -
\(\mathcal{K}(H)\)（紧算子）是 \(\mathcal{B}(H)\) 的 C\emph{-子代数（闭
}-子代数）。 - \(C(X)\)（紧致 Hausdorff 空间 \(X\)
上的连续函数），对合为复共轭 \(f^*(x) = \overline{f(x)}\)，范数为
\(\|f\|_\infty\)。这是交换 C*-代数的原型。 - \(C_0(X)\)（局部紧致
Hausdorff 空间上在无穷远处消失的连续函数）。 -
\(M_n(\mathbb{C})\)（\(n \times n\) 复矩阵），对合为共轭转置。

\subsubsection{76.2 交换 C*-代数与 Gelfand
变换}\label{ux4ea4ux6362-c-ux4ee3ux6570ux4e0e-gelfand-ux53d8ux6362}

\textbf{定理 76.1}（Gelfand-Mazur 定理）：若 \(\mathcal{A}\)
是有单位元的 Banach 代数，且每个非零元都可逆（即为除环），则
\(\mathcal{A} \cong \mathbb{C}\)。

\emph{证明}：对
\(a \in \mathcal{A}\)，\(\sigma(a) \neq \emptyset\)（谱非空）。取
\(\lambda \in \sigma(a)\)，则 \(a - \lambda 1\) 不可逆，故
\(a - \lambda 1 = 0\)，即 \(a = \lambda 1\)。∎

\textbf{定义 76.4}（特征与 Gelfand 谱）：设 \(\mathcal{A}\) 是交换
Banach 代数。\(\mathcal{A}\) 上的一个\textbf{特征}是一个非零代数同态
\(h: \mathcal{A} \to \mathbb{C}\)。所有特征构成的集合
\(\Omega(\mathcal{A})\) 称为 \(\mathcal{A}\) 的 \textbf{Gelfand
谱}（或极大理想空间）。

\textbf{命题 76.1}：对交换 Banach 代数 \(\mathcal{A}\)，存在自然双射：

\[\Omega(\mathcal{A}) \longleftrightarrow \{\text{极大理想 } \mathfrak{m} \subseteq \mathcal{A}\}\]

\(h \mapsto \ker h\)。每个特征自动连续且 \(\|h\| = 1\)（若
\(\mathcal{A}\) 有单位元）。

\textbf{定义 76.5}（Gelfand 变换）：对 \(a \in \mathcal{A}\)，定义
\(\hat{a}: \Omega(\mathcal{A}) \to \mathbb{C}\) 为
\(\hat{a}(h) = h(a)\)。映射 \(a \mapsto \hat{a}\) 称为 \textbf{Gelfand
变换}，记作 \(\Gamma: \mathcal{A} \to C(\Omega(\mathcal{A}))\)。

\(\Omega(\mathcal{A})\) 赋予弱* 拓扑（作为 \(\mathcal{A}^*\)
的子空间），它是紧致 Hausdorff 空间（若 \(\mathcal{A}\)
有单位元；否则为局部紧致）。Gelfand 变换是代数同态。

\textbf{定理 76.2}（Gelfand-Naimark 定理，交换版本）：设 \(\mathcal{A}\)
是交换 C\emph{-代数（有单位元）。则 Gelfand 变换
\(\Gamma: \mathcal{A} \to C(\Omega(\mathcal{A}))\) 是等距 }-同构。即

\[\|\hat{a}\|_\infty = \|a\|, \quad \widehat{a^*} = \overline{\hat{a}}, \quad \Gamma(\mathcal{A}) = C(\Omega(\mathcal{A}))\]

\emph{证明思路}：对
\(a \in \mathcal{A}\)，\(\|\hat{a}\|_\infty = \sup_{h \in \Omega} |h(a)| = r(a)\)（谱半径）。在
C*-代数中，若 \(a\) 是自伴的（\(a = a^*\)），则
\(\|a^2\| = \|a\|^2\)，归纳得 \(\|a^{2^n}\| = \|a\|^{2^n}\)，故
\(r(a) = \lim \|a^n\|^{1/n} = \|a\|\)。对一般
\(a\)，\(\|a\|^2 = \|a^*a\| = r(a^*a) = \|\widehat{a^*a}\|_\infty = \|\overline{\hat{a}}\hat{a}\|_\infty = \|\hat{a}\|_\infty^2\)。

满射性由 Stone-Weierstrass 定理：\(\Gamma(\mathcal{A})\) 是
\(C(\Omega)\) 的闭 *-子代数，分离点且包含常值函数，故等于
\(C(\Omega)\)。∎

\textbf{推论 76.3}：交换 C*-代数（有单位元）范畴对偶于紧致 Hausdorff
空间范畴。这是 Gelfand-Naimark 对偶。

\subsubsection{76.3 一般
C*-代数的表示}\label{ux4e00ux822c-c-ux4ee3ux6570ux7684ux8868ux793a}

\textbf{定理 76.4}（Gelfand-Naimark-Segal 构造 / GNS 构造）：设
\(\mathcal{A}\) 是 C*-代数，\(\varphi\) 是 \(\mathcal{A}\)
上的\textbf{正线性泛函}（即 \(\varphi(a^*a) \geq 0\) 对所有
\(a\)）。则存在 Hilbert 空间 \(H_\varphi\)、\(*\)-同态
\(\pi_\varphi: \mathcal{A} \to \mathcal{B}(H_\varphi)\) 和循环向量
\(\xi_\varphi \in H_\varphi\) 使得

\[\varphi(a) = \langle \pi_\varphi(a) \xi_\varphi, \xi_\varphi \rangle, \quad \forall a \in \mathcal{A}\]

\emph{证明构造}：在 \(\mathcal{A}\) 上定义内积
\(\langle a, b \rangle_\varphi = \varphi(b^* a)\)。这是半正定的。令
\(N_\varphi = \{a \in \mathcal{A} : \varphi(a^*a) = 0\}\)，则
\(N_\varphi\) 是左理想。\(H_\varphi\) 是 \(\mathcal{A}/N_\varphi\)
在该内积下的完备化。\(\pi_\varphi(a)(b + N_\varphi) = ab + N_\varphi\)。\(\xi_\varphi = 1 + N_\varphi\)。∎

\textbf{定理 76.5}（Gelfand-Naimark 定理，非交换版本）：每个
C\emph{-代数 \(\mathcal{A}\) 等距 }-同构于某个 Hilbert 空间
\(\mathcal{B}(H)\) 的 C*-子代数。

\emph{证明}：取所有正线性泛函对应的 GNS 表示 \(\{\pi_\varphi\}\)
的直和，得到忠实表示（即 \(\pi\) 是单射）。由 C\emph{-代数中单射
}-同态自动等距，即得。∎

\subsubsection{76.4
正元与序结构}\label{ux6b63ux5143ux4e0eux5e8fux7ed3ux6784}

\textbf{定义 76.6}（正元）：C*-代数 \(\mathcal{A}\) 中的元素 \(a\)
称为\textbf{正的}（记作 \(a \geq 0\)），如果 \(a = a^*\) 且
\(\sigma(a) \subseteq [0, \infty)\)。所有正元构成的集合记作
\(\mathcal{A}^+\)。

\textbf{命题 76.2}：\(a \geq 0\) 当且仅当存在 \(b \in \mathcal{A}\) 使得
\(a = b^*b\)。\(\mathcal{A}^+\) 是 \(\mathcal{A}\) 中的闭凸锥且
\(\mathcal{A}^+ \cap (-\mathcal{A}^+) = \{0\}\)。

\textbf{定义 76.7}（态）：C\emph{-代数 \(\mathcal{A}\)
上的一个\textbf{态}是一个正线性泛函 \(\varphi\) 满足
\(\|\varphi\| = 1\)（若 \(\mathcal{A}\) 有单位元，等价于
\(\varphi(1) = 1\)）。所有态构成的集合 \(\mathcal{S}(\mathcal{A})\) 是
\(\mathcal{A}^*\) 中弱} 紧致的凸集。

\textbf{定理 76.6}（Krein-Milman
定理在态空间中的应用）：\(\mathcal{S}(\mathcal{A})\)
的极点称为\textbf{纯态}。GNS 构造中，纯态对应不可约表示。

\subsubsection{76.5
量子力学诠释}\label{ux91cfux5b50ux529bux5b66ux8be0ux91ca}

C\emph{-代数为量子力学提供了严格的数学框架： -
物理系统的\textbf{可观测量}对应 C}-代数中的自伴元。 -
系统的\textbf{态}对应 C\emph{-代数上的态（正线性泛函）。 - 对态
\(\varphi\) 和可观测量 \(a = a^*\)，\(\varphi(a)\) 是测量 \(a\)
的期望值。 - 交换 C}-代数对应经典力学（可观测量可同时测量），非交换
C*-代数对应量子力学（存在不确定性原理）。

\bigskip\hrule\bigskip

\subsection{第77章：有序线性空间（Ordered Vector
Spaces）}\label{ux7b2c77ux7ae0ux6709ux5e8fux7ebfux6027ux7a7aux95f4ordered-vector-spaces}

有序线性空间是泛函分析中处理不等式、正性和序结构的框架。定义：实向量空间
\(V\) 配以凸锥 \(V_+ \subset V\)（正锥）满足
\(V_+ \cap (-V_+) = \{0\}\)，诱导偏序
\(x \leq y \iff y - x \in V_+\)。\textbf{Riesz 空间}（向量格）：若 \(V\)
配以上述偏序后每对有上确界（和下确界），则 \(V\) 为 Riesz
空间------算子理论的序基础。\textbf{核心定理}：（1）Krein-Rutman
定理------Banach
格上的紧正算子有正特征值及对应的正特征向量（Perron-Frobenius
定理的无穷维推广），在 PDE 和 Markov
过程的基本解估计中至关重要；（2）有序线性空间上的正线性泛函扩张定理（Krein
延拓定理）；（3）Banach 格是等距序同构于某个 \(L^p(\mu)\) 空间的
\(L^p\)-空间理论。在数理经济学和优化理论中，有序线性空间为一般均衡理论和凸对偶理论提供了基础框架。

\textbf{定义 77.1}（Riesz 空间 / 向量格）：实向量空间 \(V\) 配以偏序
\(\leq\)
满足：（1）\(x \leq y \implies x+z \leq y+z\)；（2）\(x \geq 0 \implies \alpha x \geq 0\)（\(\forall \alpha \geq 0\)）；（3）任意两个元
\(x, y\) 有上确界 \(x \vee y\) 和下确界 \(x \wedge y\)。称 \((V, \leq)\)
为 \textbf{Riesz 空间} 或向量格。Banach 空间兼为 Riesz 空间且满足
\(\|x\| \leq \|y\|\) 当 \(|x| \leq |y|\)（其中
\(|x| = x \vee (-x)\)）时称为 \textbf{Banach 格}。

\textbf{定理 77.1}（Krein-Rutman 定理）：设 \(X\) 为 Banach
格，\(T: X \to X\) 为紧正算子且谱半径 \(r(T) > 0\)。则 \(r(T)\) 是 \(T\)
的特征值，对应特征向量可在 \(X_+\)（正锥）中选取。\emph{证明思路}：利用
Krein-Milman 定理和 Banach-Alaoglu 定理证明 \(T\) 限制在正锥上的
Pringsheim 型序列收敛到特征向量。该定理是正矩阵 Perron-Frobenius
定理在无穷维情形的推广。

\subsection{第78章：插值空间（Interpolation
Spaces）}\label{ux7b2c78ux7ae0ux63d2ux503cux7a7aux95f4interpolation-spaces}

插值空间理论（Aronszajn-Gagliardo, 1960年代；Peetre, Lions）在 Banach
空间对之间构造中间空间，调和分析中精细的 Sobolev
嵌入和算子有界性证明依赖于此。\textbf{核心方法}：（1）\textbf{实插值（\(K\)-方法）}------设
\((X_0, X_1)\) 为相容 Banach 空间对，对 \(x \in X_0 + X_1\) 定义
\(K(t, x) = \inf\{\|x_0\|_{X_0} + t\|x_1\|_{X_1} : x = x_0 + x_1\}\)，则
\((X_0, X_1)_{\theta,q}\) 由
\(\|x\| = \left(\int_0^\infty (t^{-\theta} K(t,x))^q \frac{dt}{t}\right)^{1/q}\)
定义；（2）\textbf{复插值（Calderón
方法）}------利用解析函数在全纯带状区域边值的范数插值。\textbf{应用}：（1）Riesz-Thorin
定理的抽象化------算子 \(T: X_0 \to Y_0\) 和 \(T: X_1 \to Y_1\)
有界，则在插值空间上也有界；（2）Marcinkiewicz 插值定理------从弱
\(L^p\) 到强 \(L^p\) 的插值；（3）Sobolev 空间的分数阶推广（Besov 空间和
Triebel-Lizorkin 空间）。

\textbf{定义 78.1}（\(K\)-泛函与实插值空间）：设 \((X_0, X_1)\) 为相容
Banach 空间对（即两者连续嵌入同一 Hausdorff 拓扑向量空间）。对
\(x \in X_0 + X_1\) 定义 \textbf{\(K\)-泛函}
\(K(t, x) = \inf\{\|x_0\|_{X_0} + t\|x_1\|_{X_1} : x = x_0 + x_1\}\)（\(t > 0\)）。对
\(0 < \theta < 1\) 和 \(1 \leq q \leq \infty\)，\textbf{实插值空间}
\((X_0, X_1)_{\theta,q} = \{x \in X_0 + X_1 : \|x\|_{\theta,q} = (\int_0^\infty [t^{-\theta} K(t,x)]^q \frac{dt}{t})^{1/q} < \infty\}\)。

\textbf{定理 78.1}（插值定理）：设 \(T: X_0 + X_1 \to Y_0 + Y_1\)
为线性算子且 \(T|_{X_j}: X_j \to Y_j\) 有界（\(j = 0, 1\)）。则对任意
\(\theta \in (0,1)\) 和
\(q \in [1, \infty]\)，\(T: (X_0, X_1)_{\theta,q} \to (Y_0, Y_1)_{\theta,q}\)
有界。\emph{证明思路}：对 \(x = x_0 + x_1\) 用分解估计
\(K_Y(t, Tx) \leq \|T_0\| (K_X(\frac{\|T_1\|}{\|T_0\|} t, x))\)，代入范数定义即得。该定理统一了
Riesz-Thorin 定理（复插值）和 Marcinkiewicz
插值定理（实插值）的抽象框架。

\subsection{第79章：线性算子摄动理论（Operator Perturbation
Theory）}\label{ux7b2c79ux7ae0ux7ebfux6027ux7b97ux5b50ux6444ux52a8ux7406ux8bbaoperator-perturbation-theory}

线性算子摄动理论研究当算子受微小扰动后谱、预解式和不变子空间的变化，是量子力学（摄动展开）和数值分析的核心工具。\textbf{核心内容}：（1）\textbf{Kato-Rellich
定理}------若 \(A\) 自伴，\(B\) 对称且关于 \(A\)
相对有界（\(\|Bx\| \leq a\|x\| + b\|Ax\|\)，\(b < 1\)），则 \(A + B\)
自伴且保持 \(A\) 的本质谱；（2）\textbf{预解式级数}------对充分小的
\(\lambda\)，\((A + \lambda B - z)^{-1} = R(z) \sum_{n=0}^{\infty} (-\lambda B R(z))^n\)
给出摄动展开；（3）\textbf{解析摄动理论（Kato-Rellich）}------孤立特征值作为摄动参数的解析函数（分支不分裂条件）；（4）\textbf{奇异摄动理论}------当''微小参数''乘在最高阶导数上时（如
\(\epsilon u'' + u' + u = 0\)），边界层现象要求匹配渐近展开（Prandtl
边界层理论的基础）；（5）\textbf{Fredholm
行列式和限界式}------用于迹类算子的谱摄动精确定量分析（如半经典分析框架）。在量子力学的
Rayleigh-Schrödinger 摄动展开、散射理论和 Birkhoff
标准形理论中有核心应用。

\bigskip\hrule\bigskip

\newpage

\section{卷十四：代数拓扑}\label{ux5377ux5341ux56dbux4ee3ux6570ux62d3ux6251}

\begin{quote}
代数拓扑用代数不变量（群、环、模）研究拓扑空间的整体性质。本卷在基本群（V11）的基础上，建立奇异同调、上同调环、高维同伦群和谱序列理论，为代数几何、微分拓扑和数学物理提供核心工具。
\end{quote}

\bigskip\hrule\bigskip

\subsection{第77章：奇异同调}\label{ux7b2c77ux7ae0ux5947ux5f02ux540cux8c03}

奇异同调是最基本的同调理论，为每个拓扑空间 \(X\) 赋予一族 Abel 群
\(H_n(X)\)（\(n \geq 0\)），其秩等于 \(X\) 中 \(n\) 维「洞」的个数。

\subsubsection{77.1
奇异单形与链复形}\label{ux5947ux5f02ux5355ux5f62ux4e0eux94feux590dux5f62}

\textbf{定义 77.1}（标准单形）：\(n\) 维\textbf{标准单形}
\(\Delta^n \subseteq \mathbb{R}^{n+1}\) 定义为

\[\Delta^n = \left\{(t_0, \ldots, t_n) \in \mathbb{R}^{n+1} : t_i \geq 0, \sum_{i=0}^n t_i = 1\right\}\]

\(\Delta^n\) 的顶点为
\(e_0 = (1, 0, \ldots, 0), \ldots, e_n = (0, \ldots, 0, 1)\)。

\textbf{定义 77.2}（面映射）：第 \(i\) 个\textbf{面映射}
\(F_i^n: \Delta^{n-1} \to \Delta^n\) 定义为

\[F_i^n(t_0, \ldots, t_{n-1}) = (t_0, \ldots, t_{i-1}, 0, t_i, \ldots, t_{n-1})\]

即将 \(\Delta^{n-1}\) 嵌入为 \(\Delta^n\) 中与顶点 \(e_i\) 相对的面。

\textbf{定义 77.3}（奇异单形）：拓扑空间 \(X\) 中的一个\textbf{奇异
\(n\)-单形}是一个连续映射 \(\sigma: \Delta^n \to X\)。所有奇异
\(n\)-单形生成的自由 Abel 群记作 \(C_n(X)\)（系数在 \(\mathbb{Z}\)
中，若系数在 \(R\) 中则记作 \(C_n(X; R)\)）。

\textbf{定义 77.4}（边缘算子）：边缘算子
\(\partial_n: C_n(X) \to C_{n-1}(X)\) 在生成元上定义为

\[\partial_n(\sigma) = \sum_{i=0}^n (-1)^i \sigma \circ F_i^n\]

其中 \(\sigma \circ F_i^n\) 是 \(\sigma\) 的第 \(i\) 个面。

\textbf{命题 77.1}：\(\partial_{n-1} \circ \partial_n = 0\)。即
\((C_\bullet(X), \partial)\) 是链复形。

\emph{证明}：由面映射的恒等式
\(F_i^{n+1} \circ F_j^n = F_{j+1}^{n+1} \circ F_i^n\)（\(i \leq j\)），展开
\(\partial_{n-1} \circ \partial_n\) 后各项成对抵消。∎

\textbf{定义 77.5}（奇异同调群）：\(X\) 的 \(n\) 次\textbf{奇异同调群}为

\[H_n(X) = \frac{\ker \partial_n}{\operatorname{im} \partial_{n+1}}\]

\(Z_n(X) = \ker \partial_n\)
中的元素称为\textbf{闭链}（cycle），\(B_n(X) = \operatorname{im} \partial_{n+1}\)
中的元素称为\textbf{边缘链}（boundary）。

\subsubsection{77.2
同调的基本性质}\label{ux540cux8c03ux7684ux57faux672cux6027ux8d28}

\textbf{命题 77.2}（\(H_0\) 的计算）：若 \(X\) 非空且道路连通，则
\(H_0(X) \cong \mathbb{Z}\)。一般地，\(H_0(X)\) 是自由 Abel 群，秩等于
\(X\) 的道路连通分支数。

\emph{证明}：\(\partial_1(\sigma) = \sigma(1) - \sigma(0)\)。对道路连通空间，所有
\(0\)-单形（即点）互相等价，生成一个 \(\mathbb{Z}\)。∎

\textbf{定义 77.6}（诱导同态）：连续映射 \(f: X \to Y\) 诱导链映射
\(f_\#: C_n(X) \to C_n(Y)\)，\(f_\#(\sigma) = f \circ \sigma\)。\(f_\#\)
与边缘算子交换，故诱导同调群同态 \(f_*: H_n(X) \to H_n(Y)\)。

\textbf{命题 77.3}（同伦不变性）：若 \(f \simeq g: X \to Y\)（同伦），则
\(f_* = g_*: H_n(X) \to H_n(Y)\)。特别地，同伦等价的空间具有同构的同调群。

\emph{证明思路}：构造\textbf{棱柱算子}（prism
operator）\(P: C_n(X) \to C_{n+1}(Y)\) 满足
\(\partial P + P\partial = g_\# - f_\#\)，即 \(f_\#\) 和 \(g_\#\)
是链同伦的。\(P\) 的构造基于将 \(\Delta^n \times [0, 1]\) 剖分为 \(n+1\)
个 \((n+1)\)-单形的和。∎

\textbf{推论
77.4}：可缩空间的同调群：\(H_0(X) \cong \mathbb{Z}\)，\(H_n(X) = 0\)（\(n \geq 1\)）。

\subsubsection{77.3
相对同调与切除定理}\label{ux76f8ux5bf9ux540cux8c03ux4e0eux5207ux9664ux5b9aux7406}

\textbf{定义 77.7}（相对同调群）：设 \(A \subseteq X\)
是子空间。定义相对链群 \(C_n(X, A) = C_n(X) / C_n(A)\)。边缘算子诱导
\(\partial: C_n(X, A) \to C_{n-1}(X, A)\)。\textbf{相对同调群}
\(H_n(X, A)\) 是此链复形的同调。

\textbf{命题 77.5}（长正合序列）：对空间对
\((X, A)\)，存在自然的同调长正合序列：

\[\cdots \to H_n(A) \to H_n(X) \to H_n(X, A) \xrightarrow{\partial} H_{n-1}(A) \to \cdots \to H_0(X, A) \to 0\]

\emph{证明}：由短正合序列
\(0 \to C_\bullet(A) \to C_\bullet(X) \to C_\bullet(X, A) \to 0\)
应用蛇引理（V12 定理 70.1）。∎

\textbf{定理 77.6}（切除定理）：设 \(U \subseteq A \subseteq X\) 满足
\(\overline{U} \subseteq \operatorname{int}(A)\)。则包含映射
\((X \setminus U, A \setminus U) \hookrightarrow (X, A)\) 诱导同调同构：

\[H_n(X \setminus U, A \setminus U) \cong H_n(X, A)\]

\emph{证明思路}：使用\textbf{重心重分}（barycentric
subdivision）。将每个奇异单形细分为小单形，使得每个小单形要么落在 \(A\)
中要么落在 \(X \setminus U\) 中。重分算子是链同伦于恒等映射的。∎

\textbf{推论 77.7}（Mayer-Vietoris 序列）：设 \(X = U \cup V\)，其中
\(U, V\) 是开集（或
\(\operatorname{int}(U) \cup \operatorname{int}(V) = X\)）。则存在长正合序列：

\[\cdots \to H_n(U \cap V) \to H_n(U) \oplus H_n(V) \to H_n(X) \xrightarrow{\partial} H_{n-1}(U \cap V) \to \cdots\]

\emph{证明}：考虑短正合序列
\(0 \to C_\bullet(U \cap V) \to C_\bullet(U) \oplus C_\bullet(V) \to C_\bullet(U) + C_\bullet(V) \to 0\)。由切除定理，\(C_\bullet(U) + C_\bullet(V) \hookrightarrow C_\bullet(X)\)
是链同伦等价。应用蛇引理。∎

\subsubsection{77.4 球面的同调}\label{ux7403ux9762ux7684ux540cux8c03}

\textbf{定理 77.8}（球面的同调群）：

\[H_n(S^k) \cong \begin{cases} \mathbb{Z} & n = 0 \text{ 或 } n = k \\ 0 & \text{其他} \end{cases}\]

\emph{证明}：对 \(k\) 归纳。\(S^0\)
是两个点，\(H_0(S^0) \cong \mathbb{Z}^2\)，高阶同调为零。归纳步骤：将
\(S^k\) 分解为两个半球 \(D^k_+\) 和 \(D^k_-\)，交为 \(S^{k-1}\)。应用
Mayer-Vietoris 序列：

\[\cdots \to H_n(S^{k-1}) \to H_n(D^k) \oplus H_n(D^k) \to H_n(S^k) \to H_{n-1}(S^{k-1}) \to \cdots\]

\(D^k\) 可缩，故 \(H_n(D^k) = 0\)（\(n \geq 1\)）。因此
\(H_n(S^k) \cong H_{n-1}(S^{k-1})\) 对 \(n \geq 2\)
成立。由归纳假设得结论。∎

\textbf{推论 77.9}（Brouwer 不动点定理，一般维数）：任意连续映射
\(f: D^n \to D^n\) 有不动点。

\emph{证明}：若 \(f\) 无不动点，则存在收缩映射
\(r: D^n \to S^{n-1}\)。这将诱导 \(H_{n-1}(S^{n-1}) \cong \mathbb{Z}\)
到 \(H_{n-1}(D^n) = 0\) 的恒等同态通过 \(r_* \circ i_*\)，矛盾。∎

\textbf{推论 77.10}（维数不变性）：若
\(\mathbb{R}^n \cong \mathbb{R}^m\)（同胚），则 \(n = m\)。

\emph{证明}：\(\mathbb{R}^n \setminus \{0\} \simeq S^{n-1}\)。若
\(\mathbb{R}^n \cong \mathbb{R}^m\)，则 \(S^{n-1} \simeq S^{m-1}\)，故
\(H_{n-1}(S^{n-1}) \cong H_{n-1}(S^{m-1})\)。仅当 \(n = m\) 时可能。∎

\subsubsection{77.5 胞腔同调}\label{ux80deux8154ux540cux8c03}

\textbf{定义 77.8}（CW 复形）：CW 复形是归纳构造的空间：\(X^0\)
是离散点集；\(X^n\) 由 \(X^{n-1}\) 粘合 \(n\) 维胞腔 \(D^n_\alpha\)
得到（沿边界 \(S^{n-1}_\alpha \to X^{n-1}\)
粘合）。\(X = \bigcup_n X^n\) 赋予弱拓扑。

\textbf{定理 77.11}（胞腔同调）：对 CW 复形 \(X\)，\(H_n(X^n, X^{n-1})\)
是自由 Abel 群，秩等于 \(X\) 的 \(n\) 维胞腔数。\(H_n(X)\)
同构于胞腔链复形

\[\cdots \to H_n(X^n, X^{n-1}) \xrightarrow{d_n} H_{n-1}(X^{n-1}, X^{n-2}) \to \cdots\]

的同调，其中 \(d_n\) 由边界粘合映射的度决定。

\textbf{例 77.1}（实射影空间 \(\mathbb{R}P^n\)
的胞腔结构）：\(\mathbb{R}P^n\) 每个维数 \(0, 1, \ldots, n\)
各有一个胞腔。其胞腔同调为：

\[H_k(\mathbb{R}P^n) \cong \begin{cases} \mathbb{Z} & k = 0 \\ \mathbb{Z}/2\mathbb{Z} & 0 < k < n, k \text{ 奇} \\ 0 & 0 < k < n, k \text{ 偶} \\ \mathbb{Z} & k = n \text{（若 } n \text{ 奇）} \\ 0 & k = n \text{（若 } n \text{ 偶）} \end{cases}\]

\bigskip\hrule\bigskip

\subsection{第78章：上同调与杯积}\label{ux7b2c78ux7ae0ux4e0aux540cux8c03ux4e0eux676fux79ef}

上同调是同调的对偶理论，但其优势在于有自然的乘法结构------杯积（cup
product），将上同调变为分次环。

\subsubsection{78.1 上同调群}\label{ux4e0aux540cux8c03ux7fa4}

\textbf{定义 78.1}（上链复形）：设 \(C_\bullet(X)\) 是奇异链复形。对
Abel 群 \(G\)，定义\textbf{上链群}
\(C^n(X; G) = \operatorname{Hom}(C_n(X), G)\)。上边缘算子
\(\delta^n: C^n(X; G) \to C^{n+1}(X; G)\) 定义为
\((\delta^n \varphi)(\sigma) = \varphi(\partial_{n+1} \sigma)\)。

\textbf{定义 78.2}（奇异上同调群）：\(X\) 的系数在 \(G\) 中的 \(n\)
次\textbf{奇异上同调群}为

\[H^n(X; G) = \frac{\ker \delta^n}{\operatorname{im} \delta^{n-1}}\]

\textbf{命题
78.1}（上同调与同调的关系）：由万有系数定理，存在分裂短正合序列：

\[0 \to \operatorname{Ext}(H_{n-1}(X), G) \to H^n(X; G) \to \operatorname{Hom}(H_n(X), G) \to 0\]

特别地，若 \(G = \mathbb{R}\)（或 \(\mathbb{Q}\)），则
\(H^n(X; \mathbb{R}) \cong \operatorname{Hom}(H_n(X), \mathbb{R})\)（对偶空间）。

\subsubsection{78.2 杯积}\label{ux676fux79ef}

\textbf{定义 78.3}（杯积）：设 \(R\) 是交换环。定义\textbf{杯积}

\[\smile: H^p(X; R) \times H^q(X; R) \to H^{p+q}(X; R)\]

在上链层次上，对
\(\varphi \in C^p(X; R)\)，\(\psi \in C^q(X; R)\)，\((\varphi \smile \psi)(\sigma) = \varphi(\sigma|_{[e_0, \ldots, e_p]}) \cdot \psi(\sigma|_{[e_p, \ldots, e_{p+q}]})\)，其中
\(\sigma|_{[e_0, \ldots, e_p]}\) 表示 \(\sigma\) 限制在前 \(p+1\)
个顶点上。

\textbf{命题 78.2}（杯积的性质）： 1.
\(\delta(\varphi \smile \psi) = \delta\varphi \smile \psi + (-1)^p \varphi \smile \delta\psi\)（杯积在上同调层次良定义）
2.
\((\varphi \smile \psi) \smile \eta = \varphi \smile (\psi \smile \eta)\)（结合律）
3.
\(\varphi \smile \psi = (-1)^{pq} \psi \smile \varphi\)（分次交换律，在同调层次）
4.
杯积是自然的：\(f^*(\varphi \smile \psi) = f^*\varphi \smile f^*\psi\)

\textbf{定义 78.4}（上同调环）：\(X\) 的\textbf{上同调环}
\(H^*(X; R) = \bigoplus_{n \geq 0} H^n(X; R)\) 在杯积下构成分次交换
\(R\)-代数。

\textbf{例
78.1}（球面的上同调环）：\(H^*(S^n; \mathbb{Z}) \cong \mathbb{Z}[\alpha]/(\alpha^2)\)，其中
\(\deg \alpha = n\)。即
\(H^0 \cong H^n \cong \mathbb{Z}\)，其余为零，\(\alpha \smile \alpha = 0\)。

\textbf{例 78.2}（复射影空间 \(\mathbb{C}P^n\)
的上同调环）：\(H^*(\mathbb{C}P^n; \mathbb{Z}) \cong \mathbb{Z}[x]/(x^{n+1})\)，其中
\(\deg x = 2\)。

\textbf{例 78.3}（环面 \(T^2\)
的上同调环）：\(H^*(T^2; \mathbb{Z}) \cong \Lambda(a, b)\)（\(a, b\) 为
1 次生成元的外代数），\(a \smile b = -b \smile a\) 是 \(H^2\) 的生成元。

\subsubsection{78.3 Poincaré 对偶}\label{poincaruxe9-ux5bf9ux5076}

\textbf{定义 78.5}（定向）：\(n\) 维拓扑流形 \(M\) 的\textbf{定向}是
\(H_n(M, M \setminus \{x\}) \cong \mathbb{Z}\) 在每点 \(x\)
处的一致选择。等价地，在 \(M\) 紧致时，\textbf{基本类}
\([M] \in H_n(M)\) 是当 \(M\) 可定向时的生成元。

\textbf{定理 78.3}（Poincaré 对偶）：设 \(M\) 是紧致无边可定向 \(n\)
维流形。则杯积与基本类 \([M]\) 的配对给出同构：

\[H^k(M; \mathbb{Z}) \cong H_{n-k}(M; \mathbb{Z})\]

同构由 \(\alpha \mapsto [M] \smallfrown \alpha\) 给出（\(\smallfrown\)
是 \textbf{cap 积}）。

\emph{证明思路}：对闭可定向 \(n\) 维流形 \(M\)，取基本类
\([M] \in H_n(M)\)。定义 cap 积
\(\smallfrown: H^k(M) \times H_n(M) \to H_{n-k}(M)\)，则映射
\(\alpha \mapsto [M] \smallfrown \alpha\) 给出同构。验证方法：(1) 对
\(\mathbb{R}^n\) 中的开集用切除定理验证局部同构；(2) 用 Mayer-Vietoris
序列将局部情形粘合为全局------若 \(M = U \cup V\) 且同构在
\(U, V, U \cap V\) 上成立，则由五引理知在 \(M\) 上也成立；(3)
对紧致流形用归纳法将局部同构推广到有限并。cap
积的自然性保证了该映射不依赖于剖分选择。∎

\textbf{推论 78.4}：紧致可定向流形的 Euler 示性数为
\(\chi(M) = \sum_{k=0}^n (-1)^k \operatorname{rank} H_k(M)\)。由
Poincaré
对偶，\(\operatorname{rank} H_k(M) = \operatorname{rank} H_{n-k}(M)\)。

\textbf{定理 78.5}（Poincaré 对偶的杯积形式）：配对

\[H^k(M; \mathbb{R}) \times H^{n-k}(M; \mathbb{R}) \to \mathbb{R}, \quad (\alpha, \beta) \mapsto \langle \alpha \smile \beta, [M] \rangle\]

是非退化的双线性型。当 \(n = 4k\) 时，\(H^{2k}\)
上的二次型（符号差）是重要的拓扑不变量。

\subsubsection{78.4 de Rham
上同调回顾}\label{de-rham-ux4e0aux540cux8c03ux56deux987e}

由 V11 第 64 章的 de Rham 定理（陈述）：对光滑流形 \(M\)，de Rham 上同调
\(H^k_{dR}(M)\) 同构于奇异上同调
\(H^k(M; \mathbb{R})\)。杯积对应于微分形式的外积。

\bigskip\hrule\bigskip

\subsection{第79章：同伦群}\label{ux7b2c79ux7ae0ux540cux4f26ux7fa4}

同伦群 \(\pi_n(X)\)（\(n \geq 1\)）是基本群 \(\pi_1(X)\)
的高维推广，捕捉了 \(n\) 维球面到 \(X\) 的映射的同伦类。

\subsubsection{79.1
高维同伦群的定义}\label{ux9ad8ux7ef4ux540cux4f26ux7fa4ux7684ux5b9aux4e49}

\textbf{定义 79.1}（\(\pi_n\) 的定义）：设 \(X\) 为带基点 \((X, x_0)\)
的拓扑空间。\(\pi_n(X, x_0)\) 定义为保持基点的映射
\(f: (S^n, s_0) \to (X, x_0)\)
的保基点同伦类构成的集合。等价地，\(f: (I^n, \partial I^n) \to (X, x_0)\)
的同伦类。

\textbf{定义 79.2}（群运算）：\(\pi_n(X, x_0)\)
上的群运算由「坐标压缩」定义：对
\(f, g: (I^n, \partial I^n) \to (X, x_0)\)，

\[(f + g)(t_1, \ldots, t_n) = \begin{cases} f(2t_1, t_2, \ldots, t_n) & 0 \leq t_1 \leq 1/2 \\ g(2t_1-1, t_2, \ldots, t_n) & 1/2 \leq t_1 \leq 1 \end{cases}\]

\textbf{命题 79.1}：对 \(n \geq 2\)，\(\pi_n(X, x_0)\) 是 Abel 群。

\emph{证明}：对 \(n \geq 2\)，存在「旋转」同伦，将 \(f\) 和 \(g\)
的运算顺序交换。直观上，\(n \geq 2\) 时 \(I^n\) 的边界 \(\partial I^n\)
中 \(f\) 和 \(g\) 被挤压到不同的坐标方向，有足够空间交换。∎

\textbf{命题 79.2}（函子性）：\(\pi_n\)
是从带基点拓扑空间范畴到群范畴（\(n \geq 2\) 时为 Abel
群范畴）的函子。同伦等价诱导同伦群的同构。

\textbf{命题 79.3}（\(\pi_n\) 与覆叠空间）：若 \(p: \tilde{X} \to X\)
是覆叠映射，则对 \(n \geq 2\)，\(p_*: \pi_n(\tilde{X}) \to \pi_n(X)\)
是同构。即覆叠空间不改变高维同伦群。

\emph{证明}：\(S^n\)（\(n \geq 2\)）是单连通的，故任意映射 \(S^n \to X\)
可唯一提升到 \(\tilde{X}\)（由提升判别法，定理 61.4）。∎

\textbf{推论 79.5}：\(\pi_n(S^1) = 0\) 对 \(n \geq 2\)（因为
\(\mathbb{R}\) 可缩，\(\pi_n(\mathbb{R}) = 0\)）。

\subsubsection{79.2
球面的同伦群}\label{ux7403ux9762ux7684ux540cux4f26ux7fa4}

\textbf{定理 79.6}（\(\pi_n(S^n)\)
的计算）：\(\pi_n(S^n) \cong \mathbb{Z}\)。生成元是恒等映射
\([S^n \to S^n]\) 的同伦类。

\emph{证明思路}：使用 Hurewicz 定理（见 79.4
节）：\(H_n(S^n) \cong \mathbb{Z}\)，且
\(H_k(S^n) = 0\)（\(0 < k < n\)）。由 Hurewicz
定理，\(\pi_n(S^n) \cong H_n(S^n) \cong \mathbb{Z}\)。映射度（Brouwer
度）给出同构。∎

\textbf{定理 79.7}（部分低维同伦群，陈述）：

{\def\LTcaptype{none} % do not increment counter
\begin{longtable}[]{@{}
  >{\raggedright\arraybackslash}p{(\linewidth - 8\tabcolsep) * \real{0.2000}}
  >{\raggedright\arraybackslash}p{(\linewidth - 8\tabcolsep) * \real{0.2000}}
  >{\raggedright\arraybackslash}p{(\linewidth - 8\tabcolsep) * \real{0.2000}}
  >{\raggedright\arraybackslash}p{(\linewidth - 8\tabcolsep) * \real{0.2000}}
  >{\raggedright\arraybackslash}p{(\linewidth - 8\tabcolsep) * \real{0.2000}}@{}}
\toprule\noalign{}
\begin{minipage}[b]{\linewidth}\raggedright
球面
\end{minipage} & \begin{minipage}[b]{\linewidth}\raggedright
\(\pi_3\)
\end{minipage} & \begin{minipage}[b]{\linewidth}\raggedright
\(\pi_4\)
\end{minipage} & \begin{minipage}[b]{\linewidth}\raggedright
\(\pi_5\)
\end{minipage} & \begin{minipage}[b]{\linewidth}\raggedright
\(\pi_6\)
\end{minipage} \\
\midrule\noalign{}
\endhead
\bottomrule\noalign{}
\endlastfoot
\(S^2\) & \(\mathbb{Z}\) & \(\mathbb{Z}/2\) & \(\mathbb{Z}/2\) &
\(\mathbb{Z}/12\) \\
\(S^3\) & \(\mathbb{Z}\) & \(\mathbb{Z}/2\) & \(\mathbb{Z}/2\) &
\(\mathbb{Z}/12\) \\
\(S^4\) & \(0\) & \(\mathbb{Z}\) & \(\mathbb{Z}/2\) &
\(\mathbb{Z}/2\) \\
\end{longtable}
}

\emph{说明}：球面的同伦群计算极其困难，至今仍无一般公式。\(\pi_{n+k}(S^n)\)
对 \(k\) 固定、\(n\) 充分大时稳定（Freudenthal 悬挂定理）。

\subsubsection{79.3
纤维化与同伦长正合序列}\label{ux7ea4ux7ef4ux5316ux4e0eux540cux4f26ux957fux6b63ux5408ux5e8fux5217}

\textbf{定义 79.3}（纤维化 / Hurewicz 纤维化）：映射 \(p: E \to B\)
称为\textbf{纤维化}（或 Hurewicz 纤维化），如果对任意空间 \(X\) 和同伦
\(H: X \times I \to B\)，以及初始提升
\(\tilde{H}_0: X \to E\)（\(p \circ \tilde{H}_0 = H_0\)），存在同伦
\(\tilde{H}: X \times I \to E\) 提升 \(H\) 且延伸
\(\tilde{H}_0\)。即下列图表存在提升：

\[
\begin{CD}
X @>{\tilde{H}_0}>> E \\
@V{i_0}VV @VV{p}V \\
X \times I @>>{H}> B
\end{CD}
\]

\textbf{例 79.1}（常见纤维化）： - 覆叠映射是纤维化（离散纤维） - Hopf
纤维化：\(S^1 \to S^3 \to S^2\)，\(S^3 \to S^7 \to S^4\)，\(S^7 \to S^{15} \to S^8\)
- 道路空间纤维化：\(\Omega B \to PB \to B\)（\(PB\) 可缩，\(\Omega B\)
是回路空间）

\textbf{定理 79.8}（纤维化的同伦长正合序列）：设 \(p: E \to B\)
是纤维化，\(F = p^{-1}(b_0)\) 是纤维。取基点
\(e_0 \in F\)。则存在长正合序列：

\[\cdots \to \pi_n(F) \to \pi_n(E) \to \pi_n(B) \xrightarrow{\partial} \pi_{n-1}(F) \to \cdots \to \pi_0(E) \to \pi_0(B)\]

\emph{证明思路}：边界同态 \(\partial: \pi_n(B) \to \pi_{n-1}(F)\)
由提升性质定义：\(S^n\) 中的映射提升到 \(E\) 时，其在
\(S^{n-1}\)（赤道）上的限制落在 \(F\)
中。正合性验证需利用同伦提升性质。∎

\textbf{推论 79.9}（\(\pi_n(S^1)\) 的重新证明）：由覆叠
\(\mathbb{R} \to S^1\)（纤维 \(\mathbb{Z}\)），长正合序列给出
\(\pi_n(\mathbb{R}) \to \pi_n(S^1) \to \pi_{n-1}(\mathbb{Z})\)。\(\mathbb{R}\)
可缩，\(\mathbb{Z}\) 离散，故
\(\pi_n(S^1) = 0\)（\(n \geq 2\)），\(\pi_1(S^1) \cong \mathbb{Z}\)。

\textbf{例 79.2}（Hopf 纤维化的应用）：由 \(S^1 \to S^3 \to S^2\)
的长正合序列：

\[\cdots \to \pi_3(S^1) \to \pi_3(S^3) \to \pi_3(S^2) \to \pi_2(S^1) \to \cdots\]

\(\pi_3(S^1) = \pi_2(S^1) = 0\)，\(\pi_3(S^3) \cong \mathbb{Z}\)。故
\(\pi_3(S^2) \cong \mathbb{Z}\)（由 Hopf 映射 \(S^3 \to S^2\)
生成）。这是第一个非平凡的高维同伦群。

\subsubsection{79.4 Hurewicz 定理与 Whitehead
定理}\label{hurewicz-ux5b9aux7406ux4e0e-whitehead-ux5b9aux7406}

\textbf{定理 79.10}（Hurewicz 定理）：设 \(X\)
是道路连通空间，\(\pi_i(X) = 0\) 对 \(i < n\)（\(n \geq 2\)）。则
\(\tilde{H}_i(X) = 0\) 对 \(i < n\)，且

\[h: \pi_n(X) \to H_n(X)\]

（Hurewicz 同态，将同伦类 \([f: S^n \to X]\) 映到
\(f_*[S^n] \in H_n(X)\)）是同构。

对 \(n = 1\)，\(h: \pi_1(X) \to H_1(X)\) 是 Abel 化同态（即
\(H_1(X) \cong \pi_1(X)/[\pi_1(X), \pi_1(X)]\)）。

\emph{证明思路}：对 \(n = 1\)，\(H_1(X)\) 是 \(\pi_1(X)\) 的 Abel
化（由奇异同调的定义）。对 \(n \geq 2\)，需要构造 Hurewicz
同态的逆。假设 \(\pi_i(X) = 0\)（\(i < n\)），则 \(X\) 的
\((n-1)\)-骨架可缩，\(X\) 同伦等价于仅由 \(n\) 维及以上胞腔构成的 CW
复形。此时 \(H_n(X)\) 自由生成于 \(n\) 维胞腔，每个胞腔给出
\(S^n \to X\) 的映射，构成 \(\pi_n(X)\) 的生成元。∎

\textbf{定理 79.11}（Whitehead 定理）：设 \(f: X \to Y\) 是 CW
复形之间的连续映射。若 \(f_*: \pi_n(X) \to \pi_n(Y)\) 对所有 \(n\)
是同构，则 \(f\) 是同伦等价。

\emph{说明}：Whitehead 定理表明同伦群完全决定了 CW
复形的同伦型。但注意：存在不同伦等价的空间具有相同的同伦群（例如有不同的同调群，而同伦等价蕴含同调同构）。

\textbf{推论 79.12}：若 \(X\) 是 CW 复形且 \(\pi_n(X) = 0\) 对所有
\(n\)，则 \(X\) 可缩。

\subsubsection{79.5
悬挂同态与稳定同伦群}\label{ux60acux6302ux540cux6001ux4e0eux7a33ux5b9aux540cux4f26ux7fa4}

\textbf{定义 79.4}（悬挂）：空间 \(X\) 的\textbf{悬挂} \(\Sigma X\)
定义为 \(X \times [0, 1] / \sim\)，其中 \((x, 0) \sim (y, 0)\) 和
\((x, 1) \sim (y, 1)\) 对所有 \(x, y\)
成立。\(\Sigma S^n \cong S^{n+1}\)。

\textbf{定理 79.13}（Freudenthal 悬挂定理）：悬挂映射
\(E: \pi_k(X) \to \pi_{k+1}(\Sigma X)\)（\(E[f] = [\Sigma f]\)）在
\(k \leq 2n-2\) 时（若 \(X\) 是 \((n-1)\)-连通的）是同构。

特别地，\(\pi_{n+k}(S^n)\) 对 \(n > k+1\) 稳定，即 \(\pi_{n+k}(S^n)\)
不依赖于 \(n\)。记作 \(\pi_k^S\)（第 \(k\) 个\textbf{稳定同伦群}）。

\textbf{例
79.3}：\(\pi_0^S = \mathbb{Z}\)（映射度），\(\pi_1^S = \mathbb{Z}/2\)（Hopf
映射），\(\pi_2^S = \mathbb{Z}/2\)，\(\pi_3^S = \mathbb{Z}/24\)。

\bigskip\hrule\bigskip

\subsection{第80章：谱序列初步}\label{ux7b2c80ux7ae0ux8c31ux5e8fux5217ux521dux6b65}

谱序列是计算同调（或同伦）群的系统化工具，通过逐次逼近从容易计算的「初始页」收敛到目标群。

\subsubsection{80.1
谱序列的定义与收敛}\label{ux8c31ux5e8fux5217ux7684ux5b9aux4e49ux4e0eux6536ux655b}

\textbf{定义 80.1}（谱序列）：一个（上同调）\textbf{谱序列}是一列分次
\(R\)-模（或 Abel 群）\(E_r^{p,q}\)（\(r \geq 0\) 或
\(r \geq 1\)）和微分 \(d_r: E_r^{p,q} \to E_r^{p+r, q-r+1}\) 满足
\(d_r \circ d_r = 0\)，且
\(E_{r+1}^{p,q} \cong H^{p,q}(E_r) = \frac{\ker d_r|_{E_r^{p,q}}}{\operatorname{im} d_r|_{E_r^{p-r, q+r-1}}}\)。

\textbf{注释}：同调谱序列中微分方向为
\(d^r: E^r_{p,q} \to E^r_{p-r, q+r-1}\)。

\textbf{定义 80.2}（收敛）：称谱序列\textbf{收敛}到 \(H^*\)（记作
\(E_2^{p,q} \Rightarrow H^{p+q}\)），如果对每个
\((p, q)\)，\(E_r^{p,q}\) 对充分大的 \(r\) 稳定，且极限项
\(E_\infty^{p,q}\) 是 \(H^{p+q}\)
的某个滤过（filtration）的商：\(0 \subseteq F^0 \subseteq F^1 \subseteq \cdots \subseteq H^n\)
且 \(E_\infty^{p,q} \cong F^p H^{p+q} / F^{p-1} H^{p+q}\)。

\subsubsection{80.2
滤过与谱序列的构造}\label{ux6ee4ux8fc7ux4e0eux8c31ux5e8fux5217ux7684ux6784ux9020}

\textbf{定义 80.3}（滤过链复形）：链复形 \(C_\bullet\)
的\textbf{滤过}是一族子复形
\(\cdots \subseteq F_{p-1} C_\bullet \subseteq F_p C_\bullet \subseteq F_{p+1} C_\bullet \subseteq \cdots\)（递增滤过，或递减滤过
\(F^p\)）。

滤过诱导谱序列的 \(E^0\)
页：\(E^0_{p,q} = F_p C_{p+q} / F_{p-1} C_{p+q}\)。\(d^0\)
由原始的边缘算子诱导。\(E^1_{p,q} = H_{p+q}(F_p / F_{p-1})\)。

\textbf{定理 80.1}（滤过链复形的谱序列）：若滤过是有界的（即对每个
\(n\)，\(F_p C_n = 0\) 对充分小的 \(p\)，\(F_p C_n = C_n\) 对充分大的
\(p\)），则谱序列收敛到 \(H_*(C_\bullet)\)：

\[E^1_{p,q} = H_{p+q}(F_p C_\bullet / F_{p-1} C_\bullet) \Rightarrow H_{p+q}(C_\bullet)\]

\subsubsection{80.3 Serre 谱序列}\label{serre-ux8c31ux5e8fux5217}

\textbf{定理 80.2}（Serre 谱序列）：设 \(p: E \to B\) 是纤维化，\(F\)
是纤维，\(B\) 是道路连通的 CW 复形。则存在（同调）谱序列

\[E^2_{p,q} = H_p(B; H_q(F)) \Rightarrow H_{p+q}(E)\]

其中 \(H_p(B; H_q(F))\) 是以 \(H_q(F)\) 为局部系数的同调（若 \(B\)
单连通，则为普通同调）。

\emph{证明思路}：取 \(B\) 的骨架滤过
\(B^0 \subseteq B^1 \subseteq \cdots\)，拉回得到 \(E\) 的滤过
\(E^p = p^{-1}(B^p)\)。此滤过诱导的谱序列即 Serre 谱序列。\(E^2\)
页的计算用到纤维化在胞腔上的局部平凡性。∎

\textbf{推论 80.3}（Serre 正合序列）：若 \(B\) 是单连通 CW
复形，则存在正合序列

\[\cdots \to H_n(F) \to H_n(E) \to H_n(B) \to H_{n-1}(F) \to \cdots\]

这是 Serre 谱序列在 \(E^2_{p,q} = H_p(B) \otimes H_q(F)\)（当 \(B\)
单连通时）的特殊情形。

\textbf{例 80.1}（用 Serre 谱序列计算 \(\Omega S^n\)
的同调）：道路-回路纤维化 \(\Omega S^n \to PS^n \to S^n\)（\(PS^n\)
可缩）。Serre 谱序列给出 \(H_*(\Omega S^n) \cong \mathbb{Z}[x]/(x^2)\)
外代数（\(n\) 奇）或多项式代数（\(n\) 偶），其中生成元在 \(n-1\) 次。

\subsubsection{80.4 Leray-Serre
上同调谱序列}\label{leray-serre-ux4e0aux540cux8c03ux8c31ux5e8fux5217}

\textbf{定理 80.4}（上同调 Serre 谱序列）：在定理 80.2
的条件下，存在上同调谱序列

\[E_2^{p,q} = H^p(B; H^q(F)) \Rightarrow H^{p+q}(E)\]

当 \(B\) 单连通时，\(E_2^{p,q} \cong H^p(B) \otimes H^q(F)\)。

此外，此谱序列是\textbf{乘法的}：\(E_r\) 页上有自然的杯积结构，且微分
\(d_r\) 满足 Leibniz 法则。这极大地增强了计算能力。

\subsubsection{80.5 应用举例}\label{ux5e94ux7528ux4e3eux4f8b}

\textbf{定理 80.5}（\(\pi_4(S^3)\)
的计算）：\(\pi_4(S^3) \cong \mathbb{Z}/2\mathbb{Z}\)。

\emph{证明思路}：考虑 Hopf 纤维化
\(S^3 \to S^7 \to S^4\)。由同伦长正合序列，\(\pi_4(S^3) \cong \pi_4(S^7) \oplus \pi_5(S^4) \cong \pi_5(S^4)\)。然后使用
Serre 谱序列计算 \(\Omega S^4\) 的同调，结合 Hurewicz
定理。\(\pi_4(S^3) \cong \mathbb{Z}/2\) 由 Hopf 不变量为 1 的映射给出。∎

\textbf{定理 80.6}（Eilenberg-MacLane 空间的同调）：\(K(\mathbb{Z}, n)\)
的同调群可通过 Serre 谱序列从路径空间纤维化
\(K(\mathbb{Z}, n-1) \to PK(\mathbb{Z}, n) \to K(\mathbb{Z}, n)\)
归纳计算。例如：

\begin{itemize}
\tightlist
\item
  \(H^*(K(\mathbb{Z}, 2); \mathbb{Z}) \cong \mathbb{Z}[x]\)，\(\deg x = 2\)（\(K(\mathbb{Z}, 2) \simeq \mathbb{C}P^\infty\)）
\item
  \(H^*(K(\mathbb{Z}, 3); \mathbb{Z})\) 在低次由 \(\mathbb{Z}\) 在 3
  次生成
\end{itemize}

\textbf{定理 80.7}（Leray-Hirsch 定理）：设
\(F \to E \xrightarrow{p} B\) 是纤维化，\(B\) 是 CW 复形。若存在
\(H^*(F)\) 的齐次基 \(\{c_i\}\) 和 \(E\) 中的上同调类
\(\{\tilde{c}_i\}\) 使得
\(j^*\tilde{c}_i = c_i\)（\(j: F \hookrightarrow E\)），则

\[H^*(E) \cong H^*(B) \otimes H^*(F)\]

作为 \(H^*(B)\)-模（即 \(H^*(E)\) 是自由 \(H^*(B)\)-模，基为
\(\{\tilde{c}_i\}\)）。

\textbf{例 80.2}：旗流形 \(U(n)/T^n\) 的上同调可通过此定理从纤维化
\(U(n)/T^n \to BT^n \to BU(n)\) 计算，得到 Borel 的图像。

\bigskip\hrule\bigskip

\subsection{纤维丛与示性类}\label{ux7ea4ux7ef4ux4e1bux4e0eux793aux6027ux7c7b}

纤维丛是代数拓扑与微分几何交汇处的核心概念。它形式化了''局部平凡、全局扭曲''的几何结构------在局部上看，纤维丛是乘积空间
\(U \times F\)，但在全局上可能通过''扭转''的方式粘合而成。示性类（characteristic
classes）通过上同调类精确刻画这种扭曲，是连接拓扑、几何与数学物理的桥梁。

\subsubsection{81.1
纤维丛的基本概念}\label{ux7ea4ux7ef4ux4e1bux7684ux57faux672cux6982ux5ff5}

\textbf{定义 81.1}（纤维丛）：一个\textbf{纤维丛} \((E, \pi, B, F, G)\)
由以下数据构成：

\begin{itemize}
\tightlist
\item
  \textbf{全空间}（total space）\(E\)
\item
  \textbf{底空间}（base space）\(B\)
\item
  \textbf{投影}（projection）\(\pi: E \to B\)，是连续满射
\item
  \textbf{标准纤维}（typical fiber）\(F\)
\item
  \textbf{结构群}（structure group）\(G \subseteq \text{Diff}(F)\)（或
  \(G \subseteq \text{Homeo}(F)\)）
\item
  一族\textbf{局部平凡化}
  \(\{ \phi_i: \pi^{-1}(U_i) \to U_i \times F \}\)，其中 \(\{U_i\}\) 是
  \(B\) 的开覆盖
\end{itemize}

满足：对每对 \((i, j)\)，\textbf{转移函数}（transition
function）\(g_{ij}: U_i \cap U_j \to G\) 定义为

\[\phi_i \circ \phi_j^{-1}(x, f) = (x, g_{ij}(x) \cdot f)\]

且满足\textbf{上循环条件}（cocycle
condition）：\(g_{ij}(x) \circ g_{jk}(x) = g_{ik}(x)\)，对所有
\(x \in U_i \cap U_j \cap U_k\)。

\textbf{例 81.1}（向量丛）：标准纤维 \(F = \mathbb{R}^n\)，结构群
\(G = GL(n, \mathbb{R})\)（或
\(O(n)\)，若赋予度量）。最重要的例子是微分流形 \(M\) 上的\textbf{切丛}
\(TM\)、\textbf{余切丛} \(T^*M\)、\textbf{法丛}以及各种张量丛。Mobius
带是 \(\mathbb{R}\) 上最常见的非平凡线丛------它局部平凡但全局不可定向。

\textbf{例 81.2}（主丛）：标准纤维 \(F = G\)（李群），\(G\)
通过左乘作用于自身。\textbf{主 \(G\)-丛}（principal \(G\)-bundle）记作
\(P \to B\)。主丛是研究一般纤维丛的基础工具：每个向量丛 \(E \to B\)
关联一个主 \(GL(n)\)-丛（其\textbf{标架丛} \(F(E)\)），标架丛的纤维是
\(E\) 在一点上的所有基（标架）。

\textbf{例 81.3}（覆叠空间）：覆叠空间是纤维为离散空间的纤维丛，即 \(F\)
为离散拓扑空间，\(G\) 为离散群。

\textbf{定义 81.2}（截面）：纤维丛 \(\pi: E \to B\)
的一个\textbf{截面}（section）是连续映射 \(s: B \to E\) 满足
\(\pi \circ s = \text{id}_B\)。向量丛的截面空间记作
\(\Gamma(E)\)，它是一个
\(C^\infty(B)\)-模。切丛的截面即\textbf{向量场}，余切丛的截面即
\textbf{1-形式}。

\textbf{定义 81.3}（拉回丛）：设 \(f: X \to B\) 是连续映射，\(E \to B\)
是纤维丛。\textbf{拉回丛}（pullback bundle）\(f^*E \to X\) 定义为

\[f^*E = \{(x, e) \in X \times E : f(x) = \pi(e)\}\]

投影为 \((x, e) \mapsto x\)。拉回丛继承了原丛的纤维和结构群。

\textbf{定义 81.4}（丛同构）：两个纤维丛 \(\pi_1: E_1 \to B\) 和
\(\pi_2: E_2 \to B\)（同一底空间）称为\textbf{同构}，如果存在同胚
\(\Phi: E_1 \to E_2\) 满足 \(\pi_2 \circ \Phi = \pi_1\)，且 \(\Phi\)
在每根纤维上是线性同构（向量丛情形）或 \(G\)-等变同胚（主丛情形）。

\textbf{定理 81.1}（上同调分类定理）：对李群
\(G\)，存在\textbf{分类空间} \(BG\) 和\textbf{万有丛}
\(EG \to BG\)（\(EG\) 可缩），使得对任意仿紧空间 \(X\)，\(X\) 上主
\(G\)-丛的同构类与连续映射同伦类集 \([X, BG]\) 一一对应。丛 \(P \to X\)
对应的映射 \(f_P: X \to BG\) 称为\textbf{分类映射}（classifying
map），且 \(P \cong f_P^* EG\)。

重要例子： -
\(BGL(n, \mathbb{R}) \simeq BO(n) \simeq \text{Gr}_n(\mathbb{R}^\infty)\)（无穷
Grassmann 流形） -
\(BGL(n, \mathbb{C}) \simeq BU(n) \simeq \text{Gr}_n(\mathbb{C}^\infty)\)
- \(BU(1) \simeq \mathbb{CP}^\infty = K(\mathbb{Z}, 2)\) -
\(B\mathbb{Z}/2 \simeq \mathbb{RP}^\infty = K(\mathbb{Z}/2, 1)\)

\textbf{推论 81.2}（示性类的定义）：设 \(c \in H^*(BG; R)\)
为上同调类。对任意主 \(G\)-丛 \(P \to X\) 及分类映射
\(f_P: X \to BG\)，定义丛 \(P\) 的\textbf{示性类}为
\(c(P) = f_P^*(c) \in H^*(X; R)\)。这解释了为什么示性类是''万有''的------它们来自分类空间的上同调。

\subsubsection{81.2 Stiefel-Whitney 类与 Chern
类}\label{stiefel-whitney-ux7c7bux4e0e-chern-ux7c7b}

\textbf{定义 81.5}（Stiefel-Whitney 类）：对实向量丛 \(E \to B\)（秩
\(n\)），其 \textbf{Stiefel-Whitney 类}
\(w_i(E) \in H^i(B; \mathbb{Z}/2\mathbb{Z})\)（\(0 \leq i \leq n\)）是满足以下公理的唯一族上同调类：

\begin{enumerate}
\def\labelenumi{\arabic{enumi}.}
\tightlist
\item
  \(w_0(E) = 1\)，\(w_i(E) = 0\) 当 \(i > \text{rank}(E)\)
\item
  \textbf{自然性}（naturality）：对任意连续映射
  \(f: X \to B\)，\(w_i(f^*E) = f^* w_i(E)\)
\item
  \textbf{Whitney 和公式}：对总 Stiefel-Whitney 类
  \(w(E) = 1 + w_1(E) + w_2(E) + \cdots + w_n(E)\)，有
\end{enumerate}

\[w(E \oplus F) = w(E) \smile w(F)\]

\begin{enumerate}
\def\labelenumi{\arabic{enumi}.}
\setcounter{enumi}{3}
\tightlist
\item
  \textbf{正规化}（normalization）：\(\mathbb{RP}^1 \cong S^1\)
  上重言线丛 \(\gamma^1\) 满足
  \(w_1(\gamma^1) \neq 0 \in H^1(\mathbb{RP}^1; \mathbb{Z}/2) \cong \mathbb{Z}/2\)
\end{enumerate}

\textbf{定理 81.3}（Stiefel-Whitney 类的存在性与唯一性）：满足上述公理的
Stiefel-Whitney 类存在且唯一。存在性的标准构造使用分类空间 \(BO(n)\)
的上同调：\(H^*(BO(n); \mathbb{Z}/2) \cong \mathbb{Z}/2[w_1, \ldots, w_n]\)（多项式代数），其中
\(w_i\) 即万有 Stiefel-Whitney 类。唯一性由分裂原理保证。

\textbf{定理 81.4}（Stiefel-Whitney 类的几何意义）： -
\(w_1(E) = 0 \iff E\) 可定向（\(w_1\) 是定向障碍类） -
\(w_2(E) = 0 \iff E\) 有自旋结构（\(w_2\) 是自旋障碍类） - 紧流形 \(M\)
可平行化（tangential parallelizable）当且仅当 \(w_i(TM) = 0\) 对所有
\(i > 0\)

特别地，\(w_1(T\mathbb{RP}^n) \neq 0\)（当 \(n\) 奇），\(\mathbb{RP}^2\)
不可平行化（\(w_2(T\mathbb{RP}^2) \neq 0\)）。

\textbf{定理 81.5}（Wu 公式与 Wu 类）：对闭流形
\(M^n\)，存在唯一的上同调类 \(v_i \in H^i(M; \mathbb{Z}/2)\)（Wu
类），使得对所有
\(x \in H^{n-i}(M; \mathbb{Z}/2)\)，\(\langle v_i \smile x, [M] \rangle = \langle Sq^i(x), [M] \rangle\)，其中
\(Sq^i\) 是 Steenrod 平方运算。Stiefel-Whitney 类与 Wu 类的关系为
\(w = Sq(v)\)，即 \(w_k = \sum_{i=0}^k Sq^{k-i}(v_i)\)。Wu 公式表明
Stiefel-Whitney 类完全由上同调环的 Steenrod
代数结构决定------它们是流形的同伦型不变量。

\textbf{定义 81.6}（Chern 类）：对复向量丛 \(E \to B\)（秩 \(n\)），其
\textbf{Chern 类} \(c_i(E) \in H^{2i}(B; \mathbb{Z})\) 满足类似公理：

\begin{enumerate}
\def\labelenumi{\arabic{enumi}.}
\tightlist
\item
  \(c_0(E) = 1\)，\(c_i(E) = 0\) 当 \(i > \text{rank}(E)\)
\item
  \textbf{自然性}：\(c_i(f^*E) = f^* c_i(E)\)
\item
  \textbf{Whitney 和公式}：\(c(E \oplus F) = c(E) \smile c(F)\)
\item
  \textbf{正规化}：\(\mathbb{CP}^1 \cong S^2\) 上重言线丛
  \(\gamma^1_{\mathbb{C}}\) 满足
  \(c_1(\gamma^1_{\mathbb{C}}) = -1 \cdot [\omega_{FS}]\)，其中
  \([\omega_{FS}]\) 是 Fubini-Study 形式的基本类（即
  \(\langle c_1(\gamma^1_{\mathbb{C}}), [\mathbb{CP}^1] \rangle = -1\)）
\end{enumerate}

总 Chern 类记作 \(c(E) = 1 + c_1(E) + c_2(E) + \cdots + c_n(E)\)。Chern
类从分类空间的上同调得到：\(H^*(BU(n); \mathbb{Z}) \cong \mathbb{Z}[c_1, \ldots, c_n]\)，\(\deg c_i = 2i\)。

\textbf{例 81.4}（\(\mathbb{CP}^n\) 的切丛）：\(\mathbb{CP}^n\)
的全纯切丛 \(T\mathbb{CP}^n\) 的 Chern 类为
\(c(T\mathbb{CP}^n) = (1 + x)^{n+1}\)，其中
\(x = c_1(\mathcal{O}(1)) \in H^2(\mathbb{CP}^n; \mathbb{Z})\)
是超平面丛的第一 Chern
类。特别地，\(c_1(T\mathbb{CP}^n) = (n+1)x\)，\(c_n(T\mathbb{CP}^n) = \binom{n+1}{n} x^n = (n+1)x^n\)。

\textbf{定理 81.6}（分裂原理，Splitting Principle）：对任意向量丛
\(E \to B\)，存在连续映射 \(f: F(E) \to B\)（\textbf{旗流形}，flag
manifold），使得：

\begin{enumerate}
\def\labelenumi{\arabic{enumi}.}
\tightlist
\item
  拉回丛 \(f^*E\)
  分裂为线丛的直和：\(f^*E \cong L_1 \oplus L_2 \oplus \cdots \oplus L_n\)
\item
  诱导同态 \(f^*: H^*(B) \to H^*(F(E))\) 是\textbf{单射}
\end{enumerate}

分裂原理将复杂向量丛的示性类计算约化到线丛情形。例如，Chern
类可形式地写为 \(c(E) = \prod_{i=1}^n (1 + x_i)\)，其中
\(x_i = c_1(L_i)\) 称为 \textbf{Chern 根}（Chern roots）。

\textbf{定理 81.7}（Chern-Weil 理论）：设 \(\nabla\) 是复向量丛
\(E \to M\)（\(M\) 是光滑流形）上的联络，其曲率
\(F_\nabla \in \Omega^2(M; \text{End}(E))\)。则 Chern
类可由曲率形式的多项式表示为 de Rham 上同调类：

\[c_k(E) = \left[ \sigma_k\left(\frac{i}{2\pi} F_\nabla\right) \right] \in H^{2k}_{\text{dR}}(M)\]

其中 \(\sigma_k\) 是第 \(k\) 个初等对称多项式（即特征值的第 \(k\)
个对称和）。特别地，

\[c_1(E) = \left[\frac{i}{2\pi} \text{Tr}(F_\nabla)\right], \quad c_2(E) = \left[\frac{1}{8\pi^2}\left(\text{Tr}(F_\nabla \wedge F_\nabla) - \text{Tr}(F_\nabla) \wedge \text{Tr}(F_\nabla)\right)\right]\]

Chern-Weil 理论的美妙之处在于：尽管曲率 \(F_\nabla\)
依赖于联络的选择，由曲率多项式定义的上同调类却不依赖于联络------它是拓扑不变量。

\subsubsection{81.3 Pontryagin 类与 Euler
类}\label{pontryagin-ux7c7bux4e0e-euler-ux7c7b}

\textbf{定义 81.7}（Pontryagin 类）：实向量丛 \(E \to B\)（秩 \(n\)）的
\textbf{Pontryagin 类}
\(p_k(E) \in H^{4k}(B; \mathbb{Z})\)（或系数在任意交换环中）由复化定义：

\[p_k(E) = (-1)^k c_{2k}(E \otimes_{\mathbb{R}} \mathbb{C})\]

总 Pontryagin 类 \(p(E) = 1 + p_1(E) + p_2(E) + \cdots\)，其中
\(p_k(E)\) 的次数为 \(4k\)。注意 \(c_{2k+1}(E \otimes \mathbb{C})\) 是
2-挠元，故定义中只取偶数次 Chern 类。

由分裂原理，Pontryagin 类可形式地写为 \(p(E) = \prod (1 + x_i^2)\)，其中
\(x_i\) 是复化的 Chern 根（本质上是实丛的「Euler 形式」的对称化）。

\textbf{例 81.5}（球面的 Pontryagin 类）：对 \(S^n\)，其切丛的
Pontryagin 类在 \(n \neq 4k\) 时平凡。\(TS^4\) 有非平凡 \(p_1\)（实际上
\(S^4\) 不是某个紧 Lie 群的齐性空间这一事实可通过 Pontryagin
类得到部分解释）。

\textbf{定义 81.8}（Euler 类）：设 \(E \to B\) 是秩 \(n\)
的\textbf{定向}实向量丛。其 \textbf{Euler 类}
\(e(E) \in H^n(B; \mathbb{Z})\) 满足：

\begin{enumerate}
\def\labelenumi{\arabic{enumi}.}
\tightlist
\item
  自然性：\(e(f^*E) = f^* e(E)\)
\item
  若 \(E\) 是奇秩丛，则 \(2e(E) = 0\)
\item
  \(e(E \oplus F) = e(E) \smile e(F)\)（在定向相容条件下的 Whitney
  和公式的乘法形式）
\item
  \textbf{正规化}：对 \(n\) 维闭定向流形
  \(M\)，\(\langle e(TM), [M] \rangle = \chi(M)\)（Euler 示性数）
\end{enumerate}

\textbf{定理 81.8}（Euler 类与 Stiefel-Whitney 类的关系）：在模 2
约化下，\(e(E) \equiv w_n(E) \pmod{2}\)。换言之，Euler 类是顶维
Stiefel-Whitney 类的整系数提升。

\textbf{定理 81.9}（Gauss-Bonnet-Chern 定理）：设 \(M\) 是 \(2n\)
维闭定向 Riemann 流形，\(\Omega\) 是其曲率 2-形式（Levi-Civita
联络的曲率）。则

\[\chi(M) = \frac{1}{(2\pi)^n} \int_M \text{Pf}(\Omega)\]

其中 \(\text{Pf}(\Omega)\) 是曲率形式的 Pfaffian，满足
\(\text{Pf}(\Omega)^2 = \det(\Omega)\)。这一定理将 Euler
示性数（拓扑量）表示为曲率的积分（几何量），是 Gauss-Bonnet
定理的高维推广。

\textbf{定理 81.10}（Hirzebruch 符号差定理）：设 \(M\) 是 \(4k\)
维定向紧流形。其\textbf{符号差}（signature）\(\sigma(M)\)（即相交形式在中间维上同调
\(H^{2k}(M; \mathbb{R})\) 上的符号差）由 \(L\)-亏格（\(L\)-genus）给出：

\[\sigma(M) = \langle L_k(p_1, \ldots, p_k), [M] \rangle\]

其中 \(L\)-多项式由形式幂级数 \(\prod_{i=1}^k \frac{x_i}{\tanh x_i}\) 的
Taylor 展开定义（\(x_i\) 是 Pontryagin 根）。前几个 \(L\)-多项式为：

\[L_1 = \frac{1}{3}p_1, \quad L_2 = \frac{1}{45}(7p_2 - p_1^2), \quad L_3 = \frac{1}{945}(62p_3 - 13p_2p_1 + 2p_1^3)\]

\textbf{例 81.6}（\(\mathbb{CP}^2\) 的符号差）：\(\mathbb{CP}^2\)
的相交形式为 \([1]\)，故
\(\sigma(\mathbb{CP}^2) = 1\)。\(p_1(T\mathbb{CP}^2) = 3x^2\)（\(x\)
为生成元），\(L_1 = p_1/3 = x^2\)，\(\langle x^2, [\mathbb{CP}^2] \rangle = 1 = \sigma(\mathbb{CP}^2)\)。

\subsubsection{81.4
示性类与指标定理}\label{ux793aux6027ux7c7bux4e0eux6307ux6807ux5b9aux7406}

示性类深层的意义在 Atiyah-Singer
指标定理中得到了最完整的体现------它统一了多个经典定理，揭示了分析学（椭圆算子的指标）、拓扑学（示性类）和几何学（曲率）之间的内在联系。

\textbf{定理 81.11}（Atiyah-Singer 指标定理）：设
\(D: \Gamma(E) \to \Gamma(F)\) 是紧无边流形 \(M\) 上的椭圆微分算子（例如
Dirac 算子、Dolbeault 算子、符号差算子），\(\sigma(D)\)
是其主象征（principal symbol），可视作 \(K\) 理论中的元素。则

\[\text{ind}(D) = (-1)^{n(n+1)/2} \langle \text{ch}(\sigma(D)) \smile \hat{A}(TM), [TM] \rangle\]

其中： - \(\text{ch}: K(M) \to H^*(M; \mathbb{Q})\) 是 \textbf{Chern
特征}（Chern character），定义为
\(\text{ch}(E) = \sum_{i=1}^n e^{x_i}\)，\(x_i\) 为 Chern 根。前几项为

\[\text{ch}(E) = \text{rank}(E) + c_1 + \frac{1}{2}(c_1^2 - 2c_2) + \frac{1}{6}(c_1^3 - 3c_1c_2 + 3c_3) + \cdots\]

\begin{itemize}
\tightlist
\item
  \(\hat{A}(TM)\) 是 \textbf{\(\hat{A}\)-亏格}（\(\hat{A}\)-genus，或
  A-hat genus），由形式幂级数 \(\prod \frac{x_i/2}{\sinh(x_i/2)}\)
  定义（\(x_i\) 是 \(TM \otimes \mathbb{C}\) 的 Pontryagin 根）
\end{itemize}

\textbf{推论 81.12}（指标定理的重要特例）：

\begin{enumerate}
\def\labelenumi{\arabic{enumi}.}
\item
  \textbf{Gauss-Bonnet-Chern 定理}：取 de Rham 算子
  \(D = d + d^*\)（作用于 \(\bigwedge^* T^*M\)），则
  \(\text{ind}(D) = \chi(M)\)，指标定理约化为 Euler
  示性数的曲率积分表示。
\item
  \textbf{Hirzebruch-Riemann-Roch 定理}：取 Dolbeault 算子
  \(\bar{\partial}\)（作用于全纯向量丛），则
\end{enumerate}

\[\chi(M, E) = \langle \text{ch}(E) \smile \text{Td}(TM), [M] \rangle\]

其中 \(\text{Td}\) 是 \textbf{Todd 亏格}（Todd genus），由
\(\prod \frac{x_i}{1 - e^{-x_i}}\) 定义。这一定理将全纯 Euler 示性数与
Chern 类联系起来，是复几何的基石。

\begin{enumerate}
\def\labelenumi{\arabic{enumi}.}
\setcounter{enumi}{2}
\tightlist
\item
  \textbf{Hirzebruch 符号差定理}：取符号差算子，\(\hat{A}\)-亏格替换为
  \(L\)-亏格。
\end{enumerate}

\textbf{定理 81.13}（Chern 特征的函子性）：Chern 特征
\(\text{ch}: K(X) \to H^*(X; \mathbb{Q})\)
是环同态：\(\text{ch}(E \oplus F) = \text{ch}(E) + \text{ch}(F)\)，\(\text{ch}(E \otimes F) = \text{ch}(E) \smile \text{ch}(F)\)。在张量积
\(H^*(X; \mathbb{Q})\) 后，\(\text{ch}\) 诱导同构
\(K(X) \otimes \mathbb{Q} \cong H^*(X; \mathbb{Q})\)（对有限 CW 复形）。

\textbf{定理 81.14}（Grothendieck-Riemann-Roch 定理）：设 \(f: X \to Y\)
是光滑射影态射（代数几何中），\(E\) 是 \(X\) 上的向量丛。则

\[\text{ch}(f_! E) \smile \text{Td}(TY) = f_*(\text{ch}(E) \smile \text{Td}(TX))\]

其中 \(f_!\) 是 \(K\) 理论的直接像。这一定理是 Hirzebruch-Riemann-Roch
的相对化版本，在模空间理论和枚举几何中有深远应用。

\subsubsection{81.5
小结：示性类的统一视角}\label{ux5c0fux7ed3ux793aux6027ux7c7bux7684ux7edfux4e00ux89c6ux89d2}

示性类理论可总结为以下层次结构：

{\def\LTcaptype{none} % do not increment counter
\begin{longtable}[]{@{}llll@{}}
\toprule\noalign{}
示性类 & 系数 & 次数 & 几何意义 \\
\midrule\noalign{}
\endhead
\bottomrule\noalign{}
\endlastfoot
\(w_i\)（Stiefel-Whitney） & \(\mathbb{Z}/2\) & \(i\) &
定向障碍、自旋障碍 \\
\(c_i\)（Chern） & \(\mathbb{Z}\) & \(2i\) & 复向量丛的拓扑扭曲度量 \\
\(p_i\)（Pontryagin） & \(\mathbb{Z}\) & \(4i\) &
实向量丛复化后的扭曲 \\
\(e\)（Euler） & \(\mathbb{Z}\) & \(n\)（秩） & 截面消失的障碍类 \\
\end{longtable}
}

这些示性类的关系由以下交换图表概括：

\begin{itemize}
\tightlist
\item
  实向量丛 \(\xrightarrow{\text{复化}}\)
  复向量丛：\(p_k(E) = (-1)^k c_{2k}(E \otimes \mathbb{C})\)
\item
  复向量丛 \(\xrightarrow{\text{忘却复结构}}\)
  实向量丛：\(c_i(E) \mapsto w_{2i}(E_{\mathbb{R}})\)（模 2）
\item
  Euler 类与顶维 Stiefel-Whitney 类：\(e(E) \equiv w_n(E) \pmod{2}\)
\item
  所有示性类 \(\xrightarrow{\text{Chern-Weil}}\) 曲率形式的多项式（de
  Rham 上同调）
\end{itemize}

示性类不仅是拓扑学的核心工具，更是现代理论物理中不可或缺的语言------从杨-Mills
理论的瞬子数（由第二 Chern 类 \(c_2\)
的积分给出）到弦论中的反常消除（要求 \(p_1/2 = 0\)），从拓扑绝缘体的
\(\mathbb{Z}_2\) 不变量到量子霍尔效应中的 Chern
数，示性类的思想渗透到物理学的多个前沿领域。

\bigskip\hrule\bigskip

\subsection{第82章：低维流形专题（3维与4维）}\label{ux7b2c82ux7ae0ux4f4eux7ef4ux6d41ux5f62ux4e13ux98983ux7ef4ux4e0e4ux7ef4}

低维拓扑的核心研究对象是 3 维和 4
维流形------它们在几何和分析中表现出与高维 (\(\geq 5\))
截然不同的现象，需要专门的技术和理论。在 3 维和 4 维，Whitney
技巧失效，阻碍了高维手术理论的直接应用。

\textbf{82.1 三维流形}：\textbf{Thurston 几何化猜想}（1982 年提出，由
Perelman 于 2003 年使用 Ricci 流证明）断言每个紧致 3
维流形可沿球面和环面唯一分解为具有八种几何结构之一的片。作为推论，Poincaré
猜想的证明：每个单连通闭 3 维流形同胚于 \(S^3\)。基本构造包括 Heegaard
分裂、Dehn 手术和 JSJ
分解。\textbf{量子不变量}：Witten-Reshetikhin-Turaev
不变量（Chern-Simons 量子场论）和 Khovanov 同调给出 3
维流形的新型拓扑不变量，扩展了经典基本群方法。

\textbf{82.2
四维流形}：四维的最大特点在于存在\textbf{怪异光滑结构}。\textbf{Donaldson
理论}（1983 年）利用 Yang-Mills
瞬子方程的模空间构造了四维流形的光滑不变量（Donaldson
多项式不变量），首次证明 \(\mathbb{R}^4\)
存在不可数多个两两不微分同胚的光滑结构。\textbf{Seiberg-Witten
理论}（1994 年）提供了计算更简便的光滑不变量，证明了 Thom
猜想（代数曲线在光滑同调类中最小亏格）。\textbf{Freedman 定理}（1982
年）通过 Casson 手柄给出了单连通四维拓扑流形的完整分类（由相交形式和
Kirby-Siebenmann 不变量决定），结合 Donaldson
的结果展示了拓扑与光滑范畴的巨大差异------存在无穷多拓扑同胚但不微分同胚的四维流形对。

\bigskip\hrule\bigskip

\bigskip\hrule\bigskip

\newpage

\section{卷十五：代数几何}\label{ux5377ux5341ux4e94ux4ee3ux6570ux51e0ux4f55}

\begin{quote}
代数几何研究多项式方程组的解集------代数簇------的几何性质。它用交换代数（V12）的工具研究几何对象，用拓扑学（V10）和同调代数（V12
Ch
70）的方法计算不变量。本卷从经典仿射和射影簇出发，逐步引入概形语言和层上同调，为现代代数几何的核心工具奠定基础。
\end{quote}

\bigskip\hrule\bigskip

\subsection{第81章：仿射代数簇}\label{ux7b2c81ux7ae0ux4effux5c04ux4ee3ux6570ux7c07}

仿射代数簇是代数几何最基础的研究对象，它是多项式方程组的零点集，带有
Zariski 拓扑和坐标环结构。本章的核心是 Hilbert
零点定理，它建立了代数集与理想之间的深刻对应。

\subsubsection{81.1 仿射代数集}\label{ux4effux5c04ux4ee3ux6570ux96c6}

\textbf{定义 81.1}（仿射空间）：域 \(k\) 上的 \(n\) 维\textbf{仿射空间}
\(\mathbb{A}_k^n\)（或 \(\mathbb{A}^n\)）是 \(k^n\)
作为集合，但赋予来自多项式环 \(k[x_1, \ldots, x_n]\) 的几何结构。

\textbf{定义 81.2}（代数集）：设 \(S \subseteq k[x_1, \ldots, x_n]\)
是多项式子集。\(S\) 确定的\textbf{仿射代数集}（或零点集）为

\[V(S) = \{P \in \mathbb{A}^n : f(P) = 0 \text{ 对所有 } f \in S\}\]

若 \(S = \{f_1, \ldots, f_m\}\) 有限，则 \(V(S) = V(f_1, \ldots, f_m)\)
是 \(m\) 个多项式方程的解集。

\textbf{定义 81.3}（理想）：若 \(I\) 是 \(k[x_1, \ldots, x_n]\)
的理想，记 \(V(I) = V(S)\)，其中 \(S\) 是 \(I\) 的任意生成元集。由
Hilbert 基定理（定理 68.2），\(k[x_1, \ldots, x_n]\) 是 Noether
环，故每个理想 \(I\) 由有限个多项式生成，因此 \(V(I)\)
始终是有限个方程的解集。

\textbf{命题 81.1}（代数集的基本性质）：

\begin{enumerate}
\def\labelenumi{\arabic{enumi}.}
\tightlist
\item
  \(V(0) = \mathbb{A}^n\)，\(V(1) = \varnothing\)
\item
  \(\bigcap_\alpha V(I_\alpha) = V\left(\sum_\alpha I_\alpha\right)\)（任意交）
\item
  \(V(I) \cup V(J) = V(IJ) = V(I \cap J)\)（有限并）
\end{enumerate}

因此，代数集的补集构成了 \(\mathbb{A}^n\) 上某个拓扑的闭集，称为
\textbf{Zariski 拓扑}。

\emph{证明}：1 显然。2：\(P \in \bigcap_\alpha V(I_\alpha)\)
当且仅当对所有
\(\alpha\)，\(f(P) = 0\)（\(\forall f \in I_\alpha\)），当且仅当
\(P \in V(\sum_\alpha I_\alpha)\)。3：\(P \in V(I) \cup V(J)\) 当且仅当
\(P \in V(I)\) 或 \(P \in V(J)\)。若 \(P \in V(I)\)，则对任意
\(fg \in IJ\)，\((fg)(P) = f(P)g(P) = 0\)，故 \(P \in V(IJ)\)。反之，若
\(P \notin V(I) \cup V(J)\)，则存在 \(f \in I\)，\(g \in J\) 使
\(f(P) \neq 0\)，\(g(P) \neq 0\)。则 \(fg \in IJ\) 且
\((fg)(P) \neq 0\)，故 \(P \notin V(IJ)\)。∎

\textbf{定义 81.4}（Zariski 拓扑）：\(\mathbb{A}^n\) 上的
\textbf{Zariski 拓扑} 以代数集为闭集。开集是代数集的补集。

\textbf{注释 81.1}：Zariski 拓扑不是 Hausdorff 的（除非
\(n = 0\)）。任意两个非空开集都有非空交（不可约性）。这与分析中的拓扑形成鲜明对比，体现了代数几何的独特性质。

\subsubsection{81.2 Hilbert
零点定理}\label{hilbert-ux96f6ux70b9ux5b9aux7406}

\textbf{定义 81.5}（理想的方向）：对子集
\(X \subseteq \mathbb{A}^n\)，定义

\[I(X) = \{f \in k[x_1, \ldots, x_n] : f(P) = 0 \text{ 对所有 } P \in X\}\]

\(I(X)\) 是 \(k[x_1, \ldots, x_n]\) 的理想（称为 \(X\)
的\textbf{消去理想}）。

\textbf{定义 81.6}（根理想）：理想 \(I\) 的\textbf{根}为

\[\sqrt{I} = \{f \in k[x_1, \ldots, x_n] : f^m \in I \text{ 对某个 } m \geq 1\}\]

若 \(I = \sqrt{I}\)，则称 \(I\) 为\textbf{根理想}。

\textbf{定理 81.2}（Hilbert 零点定理，Nullstellensatz）：设 \(k\)
是代数闭域。则对任意理想 \(I \subseteq k[x_1, \ldots, x_n]\)，

\[I(V(I)) = \sqrt{I}\]

此外，若 \(V(I) = \varnothing\)，则 \(I = (1)\)（弱零点定理）。

\emph{证明思路}：弱零点定理是拉比诺维奇技巧（Rabinowitsch
trick）的推论：若 \(I\) 是极大理想，则 \(k[x_1, \ldots, x_n]/I\) 是
\(k\) 的有限扩张，由 \(k\) 代数闭知
\(I = (x_1 - a_1, \ldots, x_n - a_n)\)。强零点定理由弱零点定理和商环的构造推得。∎

\textbf{推论 81.3}（代数集与根理想的一一对应）：在代数闭域 \(k\)
上，映射 \(I \mapsto V(I)\) 和 \(X \mapsto I(X)\) 给出了

\[\{\text{根理想 } I \subseteq k[x_1, \ldots, x_n]\} \longleftrightarrow \{\text{代数集 } X \subseteq \mathbb{A}^n\}\]

之间的一一对应（反向包含序）。极大理想对应于单点集。

\subsubsection{81.3
不可约性与代数簇}\label{ux4e0dux53efux7ea6ux6027ux4e0eux4ee3ux6570ux7c07}

\textbf{定义 81.7}（不可约代数集）：代数集 \(X\)
称为\textbf{不可约的}，如果 \(X\)
不能写成两个真闭子集的并：\(X = X_1 \cup X_2\) 推出 \(X = X_1\) 或
\(X = X_2\)。不可约的仿射代数集称为\textbf{仿射代数簇}（affine
variety）。

\textbf{命题 81.4}：代数集 \(X\) 不可约当且仅当 \(I(X)\) 是素理想。

\emph{证明}：\(X\) 可约 \(\iff\)
\(X = V(I) \cup V(J) = V(IJ)\)（\(X \neq V(I), V(J)\)）\(\iff\)
\(IJ \subseteq I(X)\) 但 \(I \nsubseteq I(X)\)，\(J \nsubseteq I(X)\)
\(\iff\) \(I(X)\) 不是素理想。∎

\textbf{推论 81.5}（不可约分解）：每个代数集 \(X\)
可唯一地写成有限个不可约代数簇的并
\(X = X_1 \cup \cdots \cup X_r\)，其中
\(X_i \nsubseteq X_j\)（\(i \neq j\)）。这对应于 \(I(X)\)
的极小素理想分解。

\emph{证明}：由 \(k[x_1, \ldots, x_n]\) 是 Noether
环，每个根理想是有限个极小素理想的交。唯一性由素理想的极小性保证。∎

\subsubsection{81.4
坐标环与态射}\label{ux5750ux6807ux73afux4e0eux6001ux5c04}

\textbf{定义 81.8}（坐标环）：仿射代数簇 \(X \subseteq \mathbb{A}^n\)
的\textbf{坐标环}（或仿射坐标环）定义为

\[A(X) = k[x_1, \ldots, x_n] / I(X)\]

\(A(X)\) 是有限生成 \(k\)-代数，且是整环（因为 \(I(X)\)
是素理想）。\(A(X)\) 的元素是 \(X\) 上的正则函数。

\textbf{定义 81.9}（正则函数）：\(X\) 上的\textbf{正则函数}是多项式函数
\(f: X \to k\)，即 \(f \in A(X)\)。\(X\) 上的正则函数全体构成环
\(A(X)\)。

\textbf{命题 81.6}（坐标环的函子性）：仿射代数簇与有限生成整
\(k\)-代数之间的范畴等价：

\[\{\text{仿射代数簇 } / k\} \longleftrightarrow \{\text{有限生成整 } k\text{-代数}\}^{\text{op}}\]

等价由 \(X \mapsto A(X)\) 和
\(A \mapsto \operatorname{Spec}_m(A)\)（\(A\) 的极大谱）给出。

\textbf{定义 81.10}（态射）：仿射代数簇 \(X \subseteq \mathbb{A}^n\) 到
\(Y \subseteq \mathbb{A}^m\)
的\textbf{态射}（正则映射）\(\varphi: X \to Y\)
是每个分量由正则函数给出的映射：\(\varphi = (f_1, \ldots, f_m)\)，\(f_i \in A(X)\)，且
\(\varphi(X) \subseteq Y\)。

态射诱导坐标环的反向同态
\(\varphi^*: A(Y) \to A(X)\)，\(\varphi^*(g) = g \circ \varphi\)。

\textbf{定理 81.7}（仿射簇范畴等价）：映射 \(X \mapsto A(X)\)
是仿射代数簇范畴与有限生成整 \(k\)-代数范畴的反向范畴等价。

\textbf{例 81.1}（有理函数域）：仿射代数簇 \(X\) 的\textbf{有理函数域}
\(k(X)\) 定义为 \(A(X)\) 的分式域。有理函数是 \(X\)
上「几乎处处」定义的正则函数的等价类。

\subsubsection{81.5 维数理论}\label{ux7ef4ux6570ux7406ux8bba}

\textbf{定义 81.11}（Krull 维数）：仿射代数簇 \(X\) 的\textbf{维数}
\(\dim X\) 定义为 \(A(X)\) 的 Krull 维数------即 \(A(X)\)
中素理想链的最大长度：

\[\dim X = \max\{r : \mathfrak{p}_0 \subsetneq \mathfrak{p}_1 \subsetneq \cdots \subsetneq \mathfrak{p}_r \subseteq A(X), \mathfrak{p}_i \text{ 素理想}\}\]

\textbf{定义 81.12}（超越次数）：设 \(K/k\) 是有限生成域扩张。\(K\) 在
\(k\) 上的\textbf{超越次数} \(\operatorname{tr.deg}_k K\) 是 \(K\)
中代数无关的极大子集的大小。

\textbf{定理 81.8}（维数与超越次数）：对仿射代数簇 \(X\)，

\[\dim X = \operatorname{tr.deg}_k k(X)\]

即 \(X\) 的维数等于其有理函数域在 \(k\) 上的超越次数。

\textbf{命题 81.9}（维数的基本性质）： 1. \(\dim \mathbb{A}^n = n\) 2.
若 \(Y \subsetneq X\) 是不可约闭子簇，则 \(\dim Y < \dim X\) 3. 超曲面
\(V(f) \subseteq \mathbb{A}^n\)（\(f\) 非常数）的维数为 \(n-1\)

\textbf{定理 81.10}（Krull 高度定理）：设 \(A\) 是 Noether
环，\(f \in A\) 非零因子。则 \(f\) 生成的理想的每个极小素理想的高度为
\(1\)。特别地，每个真包含 \(\{0\}\) 的极小素理想的高度为 \(1\)。

\bigskip\hrule\bigskip

\subsection{第82章：射影簇与态射}\label{ux7b2c82ux7ae0ux5c04ux5f71ux7c07ux4e0eux6001ux5c04}

射影簇是射影空间中齐次多项式的零点集，其紧致性（在经典拓扑中）使其在几何中具有重要地位。本章引入射影簇、有理映射和双有理等价的概念。

\subsubsection{82.1 射影空间}\label{ux5c04ux5f71ux7a7aux95f4}

\textbf{定义 82.1}（射影空间）：域 \(k\) 上的 \(n\) 维\textbf{射影空间}
\(\mathbb{P}_k^n\)（或 \(\mathbb{P}^n\)）是 \(k^{n+1} \setminus \{0\}\)
在等价关系
\((x_0, \ldots, x_n) \sim (\lambda x_0, \ldots, \lambda x_n)\)（\(\lambda \in k^\times\)）下的商。点
\((x_0 : \cdots : x_n)\) 称为\textbf{齐次坐标}。

\textbf{定义 82.2}（标准仿射开覆盖）：\(\mathbb{P}^n\) 被 \(n+1\)
个仿射开集覆盖：\(U_i = \{(x_0 : \cdots : x_n) : x_i \neq 0\} \cong \mathbb{A}^n\)（\(i = 0, \ldots, n\)）。同构由
\((x_0 : \cdots : x_n) \mapsto (x_0/x_i, \ldots, \widehat{x_i/x_i}, \ldots, x_n/x_i)\)
给出。

\textbf{定义 82.3}（齐次多项式）：多项式 \(F \in k[x_0, \ldots, x_n]\)
称为 \(d\) 次\textbf{齐次}的，如果
\(F(\lambda x_0, \ldots, \lambda x_n) = \lambda^d F(x_0, \ldots, x_n)\)。齐次多项式在射影空间中的零点集有良定义（因为
\(F(\lambda x) = 0 \iff F(x) = 0\)）。

\textbf{定义 82.4}（射影代数集）：设 \(S \subseteq k[x_0, \ldots, x_n]\)
是齐次多项式子集。\(S\) 确定的\textbf{射影代数集}为

\[V_P(S) = \{P \in \mathbb{P}^n : F(P) = 0 \text{ 对所有齐次 } F \in S\}\]

齐次理想 \(I\) 的射影零点集记作 \(V_P(I)\)。

\textbf{定义 82.5}（齐次理想）：\(k[x_0, \ldots, x_n]\) 的理想 \(I\)
称为\textbf{齐次理想}，如果 \(I\)
可由齐次多项式生成。等价地，\(I = \bigoplus_{d \geq 0} (I \cap R_d)\)，其中
\(R_d\) 是 \(d\) 次齐次多项式空间。

\textbf{定理 82.1}（射影零点定理）：设 \(k\)
是代数闭域，\(I \subseteq k[x_0, \ldots, x_n]\) 是齐次理想。若
\(V_P(I) = \varnothing\)，则
\(\sqrt{I} \supseteq (x_0, \ldots, x_n)\)（即 \(I\)
包含所有充分高次齐次多项式）。否则 \(I_P(V_P(I)) = \sqrt{I}\)，其中
\(I_P\) 是射影消去理想。

\textbf{定义
82.6}（射影代数簇）：不可约的射影代数集称为\textbf{射影代数簇}（projective
variety）。

\textbf{定义 82.7}（齐次坐标环）：射影簇 \(X \subseteq \mathbb{P}^n\)
的\textbf{齐次坐标环}为

\[S(X) = k[x_0, \ldots, x_n] / I(X)\]

\(S(X)\) 是分次
\(k\)-代数：\(S(X) = \bigoplus_{d \geq 0} S_d\)。与仿射情形不同，\(S(X)\)
的分次元素不是 \(X\) 上的函数（因为齐次坐标不是良定义的），但 \(S_d\)
中两个元素的比值是 \(X\) 上的有理函数。

\subsubsection{82.2
有理映射与双有理等价}\label{ux6709ux7406ux6620ux5c04ux4e0eux53ccux6709ux7406ux7b49ux4ef7}

\textbf{定义 82.8}（有理映射）：设 \(X\) 是代数簇。\(X\) 到 \(Y\)
的\textbf{有理映射} \(\varphi: X \dashrightarrow Y\) 是定义在 \(X\)
的某个非空开子集 \(U\)
上的态射。两个有理映射视为相等，如果它们在交上一致。

\textbf{定义 82.9}（双有理等价）：代数簇 \(X\) 和 \(Y\)
称为\textbf{双有理等价}的（记作
\(X \sim_{\text{bir}} Y\)），如果存在有理映射
\(\varphi: X \dashrightarrow Y\) 和 \(\psi: Y \dashrightarrow X\)，使得
\(\psi \circ \varphi = \operatorname{id}_X\) 和
\(\varphi \circ \psi = \operatorname{id}_Y\)（作为有理映射）。

\textbf{命题 82.2}：\(X\) 和 \(Y\)
双有理等价当且仅当它们的有理函数域同构：\(k(X) \cong k(Y)\)（作为 \(k\)
的扩域）。

\textbf{定义
82.10}（双有理几何）：双有理几何研究代数簇在双有理等价下的不变量。维数是双有理不变量。光滑性不是双有理不变量（奇点解消！）。

\textbf{定理 82.3}（Hironaka 奇点解消，陈述）：特征 \(0\)
的域上的每个代数簇都双有理等价于一个光滑射影代数簇。

\emph{说明}：Hironaka
的奇点解消定理是代数几何的里程碑式成果，其证明极其复杂，获得了 Fields
奖。正特征情形在 21 世纪由多位数学家基本解决。

\textbf{例 82.1}（二次 Cremona
变换）：\(\mathbb{P}^2 \dashrightarrow \mathbb{P}^2\) 由
\((x:y:z) \mapsto (yz:xz:xy)\) 给出。这是双有理自同构（非态射！），在
\(x=0\)、\(y=0\)、\(z=0\) 处未定义。

\textbf{例 82.2}（爆破）：\(\mathbb{A}^2\) 在原点处的\textbf{爆破}是
\(\tilde{\mathbb{A}}^2 = \{(x, y, (u:v)) \in \mathbb{A}^2 \times \mathbb{P}^1 : xv = yu\}\)。投影
\(\pi: \tilde{\mathbb{A}}^2 \to \mathbb{A}^2\)
在原点外是同构，\(\pi^{-1}(0) \cong \mathbb{P}^1\)（例外除子）。爆破是解消奇点的基本操作。

\subsubsection{82.3 Veronese 嵌入与 Segre
嵌入}\label{veronese-ux5d4cux5165ux4e0e-segre-ux5d4cux5165}

\textbf{定义 82.11}（Veronese 嵌入）：\(d\) 次 \textbf{Veronese 嵌入}
\(\nu_d: \mathbb{P}^n \to \mathbb{P}^N\)（\(N = \binom{n+d}{d} - 1\)）由所有
\(d\) 次单项式给出：

\[\nu_d(x_0 : \cdots : x_n) = (\cdots : x^I : \cdots)\]

其中 \(x^I\) 遍历所有 \(d\) 次单项式。\(\nu_d\)
是闭嵌入，其像由二次关系定义。

\textbf{定义 82.12}（Segre 嵌入）：\textbf{Segre 嵌入}
\(\sigma: \mathbb{P}^n \times \mathbb{P}^m \to \mathbb{P}^{nm + n + m}\)
由

\[\sigma((x_0 : \cdots : x_n), (y_0 : \cdots : y_m)) = (\cdots : x_i y_j : \cdots)_{0 \leq i \leq n, 0 \leq j \leq m}\]

给出。\(\sigma\) 是闭嵌入，将积 \(\mathbb{P}^n \times \mathbb{P}^m\)
实现为 \(\mathbb{P}^{nm+n+m}\) 中的射影簇。

\textbf{命题 82.4}：射影簇的积是射影簇（通过 Segre
嵌入）。仿射簇的积是仿射簇（\(\mathbb{A}^n \times \mathbb{A}^m \cong \mathbb{A}^{n+m}\)）。

\bigskip\hrule\bigskip

\subsection{第83章：概形理论初步}\label{ux7b2c83ux7ae0ux6982ux5f62ux7406ux8bbaux521dux6b65}

概形（scheme）是 Grothendieck 在 20 世纪 60
年代引入的代数几何基本语言，它将代数簇推广到任意交换环，统一了代数几何、数论和表示论。本章从结构层出发，定义仿射概形和一般概形。

\subsubsection{83.1 层}\label{ux5c42}

\textbf{定义 83.1}（预层）：拓扑空间 \(X\) 上的集合（或 Abel
群、环）\textbf{预层} \(\mathcal{F}\) 是：对每个开集 \(U \subseteq X\)
赋予一个集合 \(\mathcal{F}(U)\)，对每个包含 \(V \subseteq U\)
赋予一个限制映射 \(\rho_{UV}: \mathcal{F}(U) \to \mathcal{F}(V)\)，满足
\(\rho_{UU} = \operatorname{id}\) 和
\(\rho_{VW} \circ \rho_{UV} = \rho_{UW}\)。

\textbf{定义 83.2}（层）：预层 \(\mathcal{F}\)
称为\textbf{层}，如果满足以下粘合条件：对任意开覆盖
\(U = \bigcup_{i \in I} U_i\) 和截面 \(s_i \in \mathcal{F}(U_i)\) 满足
\(s_i|_{U_i \cap U_j} = s_j|_{U_i \cap U_j}\)（对所有
\(i, j\)），存在唯一的 \(s \in \mathcal{F}(U)\) 使得
\(s|_{U_i} = s_i\)（对所有 \(i\)）。

\textbf{定义 83.3}（茎）：\(\mathcal{F}\) 在点 \(x \in X\)
处的\textbf{茎}为
\(\mathcal{F}_x = \varinjlim_{U \ni x} \mathcal{F}(U)\)（关于限制映射的归纳极限）。\(\mathcal{F}_x\)
的元素是 \(\mathcal{F}\) 在 \(x\) 附近的「芽」（germ）。

\textbf{定义 83.4}（环层空间）：\textbf{环层空间} \((X, \mathcal{O}_X)\)
是拓扑空间 \(X\) 和 \(\mathcal{O}_X\) 上环层 \(\mathcal{O}_X\)
的配对。\(\mathcal{O}_X\) 称为\textbf{结构层}。态射
\((X, \mathcal{O}_X) \to (Y, \mathcal{O}_Y)\) 是连续映射 \(f: X \to Y\)
和环层同态 \(f^\#: \mathcal{O}_Y \to f_*\mathcal{O}_X\) 的配对。

\subsubsection{83.2 仿射概形}\label{ux4effux5c04ux6982ux5f62}

\textbf{定义 83.5}（环的谱）：交换环 \(A\) 的\textbf{谱}
\(\operatorname{Spec} A\) 是 \(A\) 的所有素理想 \(\mathfrak{p}\)
构成的集合。\(\operatorname{Spec} A\) 上的 Zariski 拓扑以

\[V(I) = \{\mathfrak{p} \in \operatorname{Spec} A : \mathfrak{p} \supseteq I\}\]

为闭集（\(I\) 是 \(A\) 的理想）。

\textbf{定义 83.6}（结构层）：\(\operatorname{Spec} A\) 上的结构层
\(\mathcal{O}_{\operatorname{Spec} A}\) 定义为：对开集
\(U \subseteq \operatorname{Spec} A\)，

\[\mathcal{O}_{\operatorname{Spec} A}(U) = \{s: U \to \coprod_{\mathfrak{p} \in U} A_{\mathfrak{p}} : s(\mathfrak{p}) \in A_{\mathfrak{p}}, \text{局部地是分数}\}\]

其中 \(A_{\mathfrak{p}}\) 是 \(A\) 在 \(\mathfrak{p}\)
处的局部化。\(\mathcal{O}_{\operatorname{Spec} A}\) 是环层，其茎
\(\mathcal{O}_{\operatorname{Spec} A, \mathfrak{p}} \cong A_{\mathfrak{p}}\)。

\textbf{定义 83.7}（仿射概形）：\textbf{仿射概形}是一个环层空间
\((X, \mathcal{O}_X)\)，它同构于某个交换环 \(A\) 的
\((\operatorname{Spec} A, \mathcal{O}_{\operatorname{Spec} A})\)。

\textbf{定理 83.1}（仿射概形范畴等价）：映射
\(A \mapsto \operatorname{Spec} A\)
给出了交换环范畴与仿射概形范畴之间的反范畴等价。态射
\(\operatorname{Spec} B \to \operatorname{Spec} A\) 一一对应于环同态
\(A \to B\)。

\textbf{命题 83.2}：对仿射代数簇 \(X\)（在代数闭域 \(k\) 上），其极大谱
\(\operatorname{Spec}_m A(X)\) 与 \(X\)（作为经典代数簇）同胚。仿射概形
\(\operatorname{Spec} A(X)\) 包含了 \(X\)
的所有不可约子簇的「一般点」，是经典簇的丰富化。

\textbf{例 83.1}（非约化概形）：\(\operatorname{Spec} k[x]/(x^2)\)
是一个单点空间（唯一的素理想是 \((x)\)），但其结构层的全局截面是
\(k[x]/(x^2)\)，包含幂零元 \(x\)。这对应了「双重点」的直观。

\subsubsection{83.3 概形的定义}\label{ux6982ux5f62ux7684ux5b9aux4e49}

\textbf{定义 83.8}（概形）：\textbf{概形}是一个环层空间
\((X, \mathcal{O}_X)\)，其中每点 \(x \in X\) 存在开邻域 \(U\) 使得
\((U, \mathcal{O}_X|_U)\) 是仿射概形。即概形局部地是交换环的谱。

\textbf{定义 83.9}（概形态射）：概形 \((X, \mathcal{O}_X)\) 到
\((Y, \mathcal{O}_Y)\) 的态射是作为环层空间的态射。

\textbf{定义 83.10}（\(S\)-概形）：固定基概形
\(S\)。\(S\)\textbf{-概形}是带有一个态射 \(X \to S\) 的概形
\(X\)。\(S = \operatorname{Spec} k\) 时，\(S\)-概形就是 \(k\) 上的概形。

\textbf{定义 83.11}（分离态射）：概形 \(X\)
称为\textbf{分离的}（separated），如果对角态射
\(\Delta: X \to X \times X\) 是闭浸入。分离概形对应于 Hausdorff
空间的代数几何类比。

\textbf{定义 83.12}（概形的性质）： - \textbf{Noether
概形}：可由有限个仿射概形覆盖，每个的环是 Noether 环 -
\textbf{整概形}：不可约且约化（\(\mathcal{O}_X(U)\) 无幂零元） -
\textbf{正规概形}：所有局部环 \(\mathcal{O}_{X, x}\) 是整闭整环 -
\textbf{光滑概形}：完备化后是形式幂级数环（类比于流形）

\textbf{定义 83.13}（射影概形）：射影空间
\(\mathbb{P}^n_A = \operatorname{Proj} A[x_0, \ldots, x_n]\)（\(A\)
是环）的闭子概形称为\textbf{射影概形}。其中 \(\operatorname{Proj}\)
是分次环的齐次谱。

\textbf{定理 83.3}（代数簇作为概形）：代数闭域 \(k\)
上的代数簇范畴等价嵌入到 \(k\)
上的有限型整分离概形范畴。概形语言将经典代数簇推广到任意环，并使「数论与几何的统一」成为可能。

\subsubsection{83.4
纤维积与基变换}\label{ux7ea4ux7ef4ux79efux4e0eux57faux53d8ux6362}

\textbf{定义 83.14}（纤维积）：概形范畴中，纤维积 \(X \times_S Y\)
满足万有性质：对任意 \(S\)-概形 \(Z\)，

\[\operatorname{Hom}_S(Z, X \times_S Y) \cong \operatorname{Hom}_S(Z, X) \times \operatorname{Hom}_S(Z, Y)\]

在仿射层面，\(\operatorname{Spec} B \times_{\operatorname{Spec} A} \operatorname{Spec} C \cong \operatorname{Spec}(B \otimes_A C)\)。

\textbf{定义 83.15}（基变换）：给定态射 \(f: X \to S\) 和基变换
\(T \to S\)，\(X\) 的\textbf{基变换}为
\(X_T = X \times_S T\)。基变换是「在 \(T\) 上重新看待 \(X\)」的操作。

\textbf{定义 83.16}（纤维）：态射 \(f: X \to Y\) 在点 \(y \in Y\)
处的\textbf{纤维} \(X_y = X \times_Y \operatorname{Spec} k(y)\)，其中
\(k(y) = \mathcal{O}_{Y, y}/\mathfrak{m}_y\) 是剩余域。纤维是 \(f\) 在
\(y\) 上的几何纤维。

\bigskip\hrule\bigskip

\subsection{第84章：上同调方法}\label{ux7b2c84ux7ae0ux4e0aux540cux8c03ux65b9ux6cd5}

层上同调是代数几何最重要的计算工具，它推广了拓扑空间的上同调（V14），为
Riemann-Roch 定理、Serre 对偶和现代代数几何的核心定理提供了框架。

\subsubsection{84.1
层上同调的定义}\label{ux5c42ux4e0aux540cux8c03ux7684ux5b9aux4e49}

\textbf{定义 84.1}（层范畴）：拓扑空间 \(X\) 上的 Abel 群层构成 Abel
范畴 \(\mathbf{Sh}(X)\)。它有足够多内射对象（内射层）。

\textbf{定义 84.2}（整体截面函子）：整体截面函子
\(\Gamma(X, -): \mathbf{Sh}(X) \to \mathbf{Ab}\) 定义为
\(\Gamma(X, \mathcal{F}) = \mathcal{F}(X)\)。\(\Gamma(X, -)\)
是左正合函子。

\textbf{定义 84.3}（层上同调）：\(\mathcal{F}\) 的 \(i\)
次\textbf{层上同调群} \(H^i(X, \mathcal{F})\)
是整体截面函子的右导出函子：

\[H^i(X, \mathcal{F}) = R^i\Gamma(X, \mathcal{F})\]

具体计算：取内射分解
\(0 \to \mathcal{F} \to \mathcal{I}^0 \to \mathcal{I}^1 \to \cdots\)，则
\(H^i(X, \mathcal{F}) = H^i(\Gamma(X, \mathcal{I}^\bullet))\)。

\textbf{命题 84.1}（基本性质）： 1.
\(H^0(X, \mathcal{F}) = \Gamma(X, \mathcal{F})\) 2. 短正合列
\(0 \to \mathcal{F} \to \mathcal{G} \to \mathcal{H} \to 0\)
诱导长正合序列 3. 若 \(\mathcal{F}\) 是内射层，则
\(H^i(X, \mathcal{F}) = 0\)（\(i > 0\)）

\subsubsection{84.2 Čech 上同调}\label{ux10dech-ux4e0aux540cux8c03}

\textbf{定义 84.4}（Čech 复形）：设 \(\mathfrak{U} = \{U_i\}_{i \in I}\)
是 \(X\) 的开覆盖。对层 \(\mathcal{F}\)，定义 Čech 复形：

\[\check{C}^p(\mathfrak{U}, \mathcal{F}) = \prod_{i_0 < \cdots < i_p} \mathcal{F}(U_{i_0} \cap \cdots \cap U_{i_p})\]

微分 \(\delta: \check{C}^p \to \check{C}^{p+1}\) 由
\((\delta s)_{i_0, \ldots, i_{p+1}} = \sum_{k=0}^{p+1} (-1)^k s_{i_0, \ldots, \widehat{i_k}, \ldots, i_{p+1}}|_{U_{i_0} \cap \cdots \cap U_{i_{p+1}}}\)
给出。

\textbf{定义 84.5}（Čech
上同调）：\(\check{H}^p(\mathfrak{U}, \mathcal{F}) = H^p(\check{C}^\bullet(\mathfrak{U}, \mathcal{F}))\)。取所有开覆盖的归纳极限得
\textbf{Čech 上同调} \(\check{H}^p(X, \mathcal{F})\)。

\textbf{定理 84.1}（Leray 定理）：若 \(\mathfrak{U}\) 是 \(X\)
的开覆盖，且对每个 \(U_{i_0} \cap \cdots \cap U_{i_p}\) 和每个
\(q > 0\)，\(H^q(U_{i_0} \cap \cdots \cap U_{i_p}, \mathcal{F}) = 0\)，则
\(\check{H}^p(\mathfrak{U}, \mathcal{F}) \cong H^p(X, \mathcal{F})\)。

特别地，对仿射概形上的拟凝聚层和仿射开覆盖，Čech 上同调等于层上同调。

\subsubsection{84.3 凝聚层与 Serre
定理}\label{ux51ddux805aux5c42ux4e0e-serre-ux5b9aux7406}

\textbf{定义 84.6}（\(\mathcal{O}_X\)-模层）：环层空间
\((X, \mathcal{O}_X)\) 上的 \(\mathcal{O}_X\)\textbf{-模层}
\(\mathcal{F}\) 是 Abel 群层，且对每个开集 \(U\)，\(\mathcal{F}(U)\) 是
\(\mathcal{O}_X(U)\)-模，限制映射与模结构兼容。

\textbf{定义 84.7}（拟凝聚层与凝聚层）：\(\mathcal{O}_X\)-模层
\(\mathcal{F}\) 称为\textbf{拟凝聚的}，如果局部地有表示
\(\mathcal{O}_X^{\oplus J}|_U \to \mathcal{O}_X^{\oplus I}|_U \to \mathcal{F}|_U \to 0\)。若
\(I\) 和 \(J\) 可取有限，则 \(\mathcal{F}\) 称为\textbf{凝聚层}。

\textbf{定义 84.8}（Serre 捻层）：在射影空间 \(\mathbb{P}^n_A\)
上，\textbf{Serre 捻层} \(\mathcal{O}(d)\) 是分次
\(A[x_0, \ldots, x_n]\)-模 \(A[x_0, \ldots, x_n](d)\)（将次数平移
\(d\)）对应的层。对射影概形
\(X \subseteq \mathbb{P}^n_A\)，\(\mathcal{O}_X(d) = \mathcal{O}(d)|_X\)。

\textbf{定理 84.2}（Serre 定理）：设 \(X\) 是 Noether 环 \(A\)
上的射影概形，\(\mathcal{F}\) 是 \(X\) 上的凝聚层。则 1. 对每个
\(i \geq 0\)，\(H^i(X, \mathcal{F})\) 是有限 \(A\)-模 2. 存在 \(d_0\)
使得对 \(d \geq d_0\)，\(H^i(X, \mathcal{F}(d)) = 0\)（\(i > 0\)） 3.
\(\mathcal{F}(d)\) 由整体截面生成（\(d \gg 0\)）

其中
\(\mathcal{F}(d) = \mathcal{F} \otimes_{\mathcal{O}_X} \mathcal{O}_X(d)\)。

\emph{证明思路}：将 \(\mathcal{F}\) 推到 \(\mathbb{P}^n_A\) 上，用 Čech
上同调计算。\(\mathbb{P}^n_A\) 的 Čech 复形是有限
\(A\)-模的复形，故上同调是有限 \(A\)-模。消没性由 \(d \to \infty\) 时
Čech 复形的分析得到。∎

\textbf{定理 84.3}（凝聚层的有限性，Grothendieck）：若 \(X\) 是域 \(k\)
上的真概形（proper scheme），\(\mathcal{F}\) 是凝聚层，则
\(H^i(X, \mathcal{F})\) 是有限维 \(k\)-向量空间。

\subsubsection{84.4 Serre 对偶}\label{serre-ux5bf9ux5076}

\textbf{定理 84.4}（Serre 对偶，射影空间）：设
\(X = \mathbb{P}^n_k\)（\(k\) 是域）。则对任意凝聚层
\(\mathcal{F}\)，存在自然的完美配对

\[H^i(X, \mathcal{F}) \times \operatorname{Ext}^{n-i}(\mathcal{F}, \omega_X) \to H^n(X, \omega_X) \cong k\]

其中 \(\omega_X = \mathcal{O}_X(-n-1)\) 是 \(\mathbb{P}^n\)
的\textbf{对偶化层}（canonical sheaf）。

\emph{证明思路}：对 \(\mathbb{P}^n_k\) 直接计算。用标准仿射开覆盖
\(\{U_i = D_+(x_i)\}\) 的 Čech
上同调：\(H^n(\mathbb{P}^n, \mathcal{O}(-n-1))\) 由 Čech \(n\)-余圈
\(\frac{1}{x_0 \cdots x_n}\) 生成，同构于 \(k\)。对一般
\(\mathcal{F}\)，构造配对
\(\operatorname{Ext}^{n-i}(\mathcal{F}, \omega) \times H^i(\mathcal{F}) \to H^n(\omega) \cong k\)：将
\(\operatorname{Ext}\) 类（代表 \(\omega \to \mathcal{I}^\bullet\)
的态射）与 Čech 上同调类复合，利用 Yoneda 配对。非退化性通过对
\(\mathcal{F} = \mathcal{O}(d)\) 的显式计算和凝聚层的有限长度归纳证明。

\textbf{定理 84.5}（Serre 对偶，一般形式）：设 \(X\) 是域 \(k\)
上的光滑射影簇，维数为 \(n\)。则存在\textbf{典范层}
\(\omega_X = \Omega_{X/k}^n\)（\(n\) 次微分形式的层），使得对任意凝聚层
\(\mathcal{F}\)，有自然同构

\[H^i(X, \mathcal{F}) \cong \operatorname{Ext}^{n-i}(\mathcal{F}, \omega_X)^\vee\]

特别地，\(H^i(X, \mathcal{F}) \cong H^{n-i}(X, \mathcal{F}^\vee \otimes \omega_X)^\vee\)（当
\(\mathcal{F}\) 是局部自由层时）。

\emph{证明思路}：由 Grothendieck 的对偶理论：真态射
\(f: X \to \operatorname{Spec} k\) 有例外逆向像 \(f^!\)（右伴随于
\(Rf_*\)）。定义对偶化复形 \(\omega_X^\bullet = f^! k\)。对光滑
\(X\)，有 \(\omega_X^\bullet \cong \omega_X[n]\)。Serre 对偶是等式
\(R\operatorname{Hom}(R\Gamma(X, \mathcal{F}), k) \cong R\Gamma(X, R\operatorname{\mathcal{H}\textit{om}}(\mathcal{F}, \omega_X^\bullet))\)
取上同调的结果。对射影空间，可通过直接 Čech 计算验证轨道同构。

\textbf{推论 84.6}（Serre 对偶的特殊情形）： 1.
\(\dim H^n(X, \mathcal{O}_X) = 1\)（\(H^n(X, \omega_X) \cong k\)） 2.
\(H^i(X, \mathcal{O}_X) \cong H^{n-i}(X, \omega_X)^\vee\) 3. 对曲线
\(C\)（\(n = 1\)），\(H^1(C, \mathcal{F}) \cong \operatorname{Hom}(\mathcal{F}, \omega_C)^\vee\)

\subsubsection{84.5 Riemann-Roch 定理}\label{riemann-roch-ux5b9aux7406}

\textbf{定义 84.9}（Euler 示性数）：凝聚层 \(\mathcal{F}\) 的
\textbf{Euler 示性数}为

\[\chi(X, \mathcal{F}) = \sum_{i \geq 0} (-1)^i \dim_k H^i(X, \mathcal{F})\]

（有限和，由 Serre 定理保证）。

\textbf{定理 84.7}（Hirzebruch-Riemann-Roch 定理）：设 \(X\) 是域 \(k\)
上的光滑射影簇，\(\mathcal{E}\) 是 \(X\) 上的局部自由层（向量丛）。则

\[\chi(X, \mathcal{E}) = \int_X \operatorname{ch}(\mathcal{E}) \cdot \operatorname{td}(T_X)\]

其中 \(\operatorname{ch}(\mathcal{E})\) 是 \(\mathcal{E}\) 的 Chern
特征标，\(\operatorname{td}(T_X)\) 是切丛的 Todd 类，\(\int_X\) 表示在
\(X\) 的基本类上的求值。

\emph{证明思路}：(1) 由分裂原理，可假设 \(\mathcal{E}\)
是线丛的直和，从而 \(\operatorname{ch}(\mathcal{E})\) 由线丛的 Chern
根的指数和给出。(2) 对射影空间 \(\mathbb{P}^n\) 上的线丛
\(\mathcal{O}(d)\)，直接计算
\(\chi(\mathbb{P}^n, \mathcal{O}(d)) = \binom{n+d}{d}\)，验证与
\(\int \operatorname{ch}(\mathcal{O}(d)) \operatorname{td}(T_{\mathbb{P}^n})\)
一致。(3) 利用嵌入 \(i: X \hookrightarrow \mathbb{P}^N\)
将一般情形约化到射影空间：由
Grothendieck-Riemann-Roch，\(\chi(X, \mathcal{E}) = \chi(\mathbb{P}^N, i_*\mathcal{E})\)，并用射影空间上的计算结果推得公式。Todd
类的定义 \(\operatorname{td}(E) = \prod \frac{x_i}{1-e^{-x_i}}\)
确保了乘积公式在短正合列下的可加性。

\textbf{定理 84.8}（曲线上的 Riemann-Roch）：设 \(C\)
是光滑射影曲线，\(D\) 是 \(C\) 上的除子。则

\[\ell(D) - \ell(K_C - D) = \deg D + 1 - g\]

其中 \(\ell(D) = \dim H^0(C, \mathcal{O}_C(D))\)，\(K_C\)
是典范除子，\(g = \dim H^1(C, \mathcal{O}_C)\) 是 \(C\)
的亏格（genus）。

\emph{证明思路}：应用 Serre 对偶定理
84.5：\(\ell(K_C - D) = \dim H^1(C, \mathcal{O}_C(D))\)。故
\(\chi(\mathcal{O}_C(D)) = \ell(D) - \ell(K_C - D)\)。而
\(\chi(\mathcal{O}_C(D)) = \deg D + \chi(\mathcal{O}_C) = \deg D + 1 - g\)（由
Riemann-Roch 的初等形式）。∎

\textbf{推论 84.9}（Riemann-Roch 的应用）： 1. \(g = 0\)
时，\(C \cong \mathbb{P}^1\)（有理曲线） 2. \(g = 1\) 时，\(C\)
是椭圆曲线，\(\ell(K_C) = 1\)，\(\deg K_C = 0\) 3. \(g \geq 2\)
时，典范映射 \(C \to \mathbb{P}^{g-1}\) 是 \(C\) 到其典范模型的态射

\textbf{定理 84.10}（Grothendieck-Riemann-Roch）：设 \(f: X \to Y\)
是光滑射影簇之间的真态射。则对 \(X\) 上的凝聚层 \(\mathcal{F}\)，

\[\operatorname{ch}(f_!\mathcal{F}) \cdot \operatorname{td}(T_Y) = f_*(\operatorname{ch}(\mathcal{F}) \cdot \operatorname{td}(T_X))\]

在 \(Y\) 的 Chow 环（或有理上同调）中，其中
\(f_! = \sum (-1)^i R^i f_* \mathcal{F}\)。

\bigskip\hrule\bigskip

\subsection{第85章：环面簇与模空间}\label{ux7b2c85ux7ae0ux73afux9762ux7c07ux4e0eux6a21ux7a7aux95f4}

环面簇是代数几何与组合数学的交叉领域，建立了扇（代数几何）与凸多面体（组合学）之间的精确字典。本章从扇的构造出发引入环面簇，并讨论模空间理论和极小模型纲领。

\subsubsection{85.1
环面簇的构造}\label{ux73afux9762ux7c07ux7684ux6784ux9020}

\textbf{环面簇}（Toric
Variety）是包含代数环面\(T \cong (k^\times)^n\)为稠密开子集、且\(T\)的群作用可延拓到全簇的正规代数簇。环面簇由组合几何数据------\textbf{扇}（fan）\(\Sigma\)------完全分类。设\(N \cong \mathbb{Z}^n\)为格，\(M = \operatorname{Hom}(N, \mathbb{Z})\)为对偶格。扇\(\Sigma\)是\(N_{\mathbb{R}} = N \otimes \mathbb{R}\)中一组强凸有理多面锥的集合，满足面封闭性和两两交为面的条件。每个锥\(\sigma \in \Sigma\)对应仿射环面簇\(U_\sigma = \operatorname{Spec}(k[\sigma^\vee \cap M])\)，其中\(\sigma^\vee = \{m \in M_{\mathbb{R}} : \langle m, v \rangle \geq 0,\ \forall v \in \sigma\}\)为对偶锥。沿公共面\(\sigma \cap \tau\)的自然粘合给出了环面簇\(X_\Sigma\)。

环面簇的几何性质一一对应于扇的组合性质：\(X_\Sigma\)光滑当且仅当每个锥\(\sigma\)由基的一部分生成（即\(\det = \pm 1\)）；\(X_\Sigma\)完备当且仅当\(\Sigma\)的支撑为整个\(N_{\mathbb{R}}\)；\(X_\Sigma\)是射影的当且仅当\(\Sigma\)是一个全维凸多面体的正规扇。经典例子：
-
\textbf{射影空间}\(\mathbb{P}^n\)：对应于以\((e_0, \ldots, e_n)\)（\(e_0 = -\sum_{i=1}^n e_i\)，\(\{e_i\}\)为\(\mathbb{Z}^n\)的标准基）为生成元的扇，其\(\binom{n+1}{1}\)个极大锥由任意\(n\)个生成元张成；
-
\textbf{加权射影空间}\(\mathbb{P}(a_0, \ldots, a_n)\)：扇的生成元改为满足\(\sum a_i u_i = 0\)的有理向量；
-
\textbf{Hirzebruch曲面}\(\mathbb{F}_a = \mathbb{P}(\mathcal{O}_{\mathbb{P}^1} \oplus \mathcal{O}_{\mathbb{P}^1}(a))\)（\(a \geq 0\)）：对应的二维扇由四个锥组成，其生成元为\((0, 1), (0, -1), (1, 0), (-1, -a)\)。环面簇建立了代数几何与组合凸多面体之间的精确对应，使奇点解消、上同调计算等问题可转化为纯组合问题。

\subsubsection{85.2 模空间理论}\label{ux6a21ux7a7aux95f4ux7406ux8bba}

\textbf{模空间理论}将代数簇（或其他几何对象）的同构类参数化，使参数空间本身具有代数簇（或叠）结构。椭圆曲线的模空间是最经典的例子：域\(k\)上的椭圆曲线由\textbf{\(j\)-不变量}\(j(E) = 1728 \cdot 4a^3/(4a^3 + 27b^2)\)完全分类，故仿射直线\(\mathbb{A}^1_j\)是椭圆曲线的粗模空间（不分叉时）。加入尖点（\(j = \infty\)对应节点三次曲线）得到紧化\(\mathbb{P}^1\)。\textbf{曲线模空间}\(M_g\)：亏格\(g \geq 2\)光滑曲线的粗模空间维数为\(3g-3\)（Riemann，1857年），由Riemann-Roch定理通过计算\(H^0(C, \omega_C^{\otimes 2})\)（二阶微分）的自由度得到。Deligne和Mumford于1969年引入\textbf{稳定曲线}概念，构造了紧化模空间\(\overline{M}_g\)，并证明\(M_g\)是不可约拟射影代数簇。在双有理几何中，\textbf{典范环}\(R(K_X) = \bigoplus_{m \geq 0} H^0(X, mK_X)\)的\(\operatorname{Proj}\)给出了簇的典范模型，而Mori程序通过分析收缩映射对模空间的结构进行精细控制，实现了高维簇的分类。

\subsubsection{85.3
双有理几何与极小模型纲领}\label{ux53ccux6709ux7406ux51e0ux4f55ux4e0eux6781ux5c0fux6a21ux578bux7eb2ux9886}

\textbf{极小模型纲领}（MMP，Minimal Model
Program）是双有理几何的核心算法，由Mori在1980年代奠基。对光滑射影簇\(X\)，其典范除子\(K_X\)的正负性是关键分类指标。\textbf{锥定理}（Cone
Theorem）断言：\(X\)的曲线锥\(\overline{NE}(X)\)中与\(K_X\)负交的部分由极射线（extremal
rays）\(R_i\)生成，每条极射线对应一个收缩映射\(\operatorname{cont}_{R_i}: X \to Y_i\)，有以下三种情形：
1. \textbf{纤维型}：\(\dim Y < \dim X\)，此时\(X\)是Mori纤维空间，终止；
2.
\textbf{除子型}：收缩一个除子，\(\dim Y = \dim X\)，继续对\(Y\)做MMP；
3.
\textbf{小收缩型}：例外集余维\(\geq 2\)，此时\(K_Y\)不是\(\mathbb{Q}\)-Cartier的，需要\textbf{翻转}（flip）：构造双有理映射\(X \dashrightarrow X^+\)使得\(K_{X^+}\)在翻转曲线上为正。

\textbf{Birkar-Cascini-Hacon-McKernan定理}（2010年）：特征零域上，任意光滑射影簇的典范环是有限生成的，且对数一般型簇的翻转存在。该结果证明了极小模型纲领在极大一般型的核心步骤，为此Birkar获得了2018年Fields奖。MMP的输出是一个极小模型（\(K_X\)是nef的）或Mori纤维空间，实现了双有理分类的最终目标。

\bigskip\hrule\bigskip

\subsection{第86章：Fano
簇与有理簇}\label{ux7b2c86ux7ae0fano-ux7c07ux4e0eux6709ux7406ux7c07}

Fano 簇和有理簇是双有理几何和正特征代数几何中最重要的对象之一。Fano
簇（反典范除子丰沛的光滑射影簇）的纤维正曲率（某种意义上）为 Mori
理论提供了基本构建单元，而有理簇问题（单有理簇是否必为有理簇）则是代数几何中最深刻的未解问题之一。

\subsubsection{86.1 Fano
簇的基本性质}\label{fano-ux7c07ux7684ux57faux672cux6027ux8d28}

\textbf{定义 86.1}（Fano 簇）：特征 0 的代数闭域 \(k\) 上的光滑射影簇
\(X\) 称为 \textbf{Fano 簇}，若其反典范除子 \(-K_X\) 是丰沛的（ample）。

等价地，\(X\) 的典范丛 \(\omega_X = \mathcal{O}_X(K_X)\) 的对偶丛
\(\omega_X^{-1}\) 是丰沛的。

\textbf{定义 86.2}（Fano 指标）：Fano 簇 \(X\) 的 \textbf{Fano
指标}（Fano index）定义为最大整数 \(r\) 满足

\[-K_X = r H\]

其中 \(H \in \operatorname{Pic}(X)\) 是丰沛除子类，且 \(H\) 不被任何大于
1 的整数整除。通常可取 \(H\) 为 \(-K_X\) 在
\(\operatorname{Pic}(X)/\operatorname{Tors}\) 中的原像。Fano 指标满足
\(1 \leq r \leq \dim X + 1\)（上界由 Kobayashi-Ochiai 定理保证）。

\textbf{命题 86.1}（Fano 簇的基本几何性质）： 1. Fano
簇是单有理的（uniruled）------即被有理曲线覆盖； 2.
\(\chi(X, \mathcal{O}_X) = 1\)，且 \(H^i(X, \mathcal{O}_X) = 0\) 对所有
\(i > 0\) 成立（由 Kodaira 消没定理）； 3. Fano 簇的 Picard 数
\(\rho(X)\) 可以为任意大的值。

\subsubsection{86.2 二维和三维 Fano
簇的分类}\label{ux4e8cux7ef4ux548cux4e09ux7ef4-fano-ux7c07ux7684ux5206ux7c7b}

\textbf{定理 86.1}（Del Pezzo 曲面分类）：二维 Fano 簇即 \textbf{del
Pezzo 曲面}，由次数 \(d = K_X^2\)（\(1 \leq d \leq 9\)）分类：

{\def\LTcaptype{none} % do not increment counter
\begin{longtable}[]{@{}
  >{\raggedright\arraybackslash}p{(\linewidth - 4\tabcolsep) * \real{0.3226}}
  >{\raggedright\arraybackslash}p{(\linewidth - 4\tabcolsep) * \real{0.4839}}
  >{\raggedright\arraybackslash}p{(\linewidth - 4\tabcolsep) * \real{0.1935}}@{}}
\toprule\noalign{}
\begin{minipage}[b]{\linewidth}\raggedright
次数 \(d\)
\end{minipage} & \begin{minipage}[b]{\linewidth}\raggedright
Del Pezzo 曲面
\end{minipage} & \begin{minipage}[b]{\linewidth}\raggedright
描述
\end{minipage} \\
\midrule\noalign{}
\endhead
\bottomrule\noalign{}
\endlastfoot
9 & \(\mathbb{P}^2\) & 射影平面 \\
8 & \(\mathbb{P}^1 \times \mathbb{P}^1\) 或 \(\mathbb{F}_1\) &
二次曲面或一次爆破 \(\mathbb{P}^2\) \\
7 & \(\mathbb{P}^2\) 上 2 点爆破 & \\
6 & \(\mathbb{P}^2\) 上 3 点爆破 & 三次曲面 \\
5 & \(\mathbb{P}^2\) 上 4 点爆破 & \\
4 & \(\mathbb{P}^2\) 上 5 点爆破 & \(\mathbb{P}^4\)
中两个二次超曲面的交 \\
3 & \(\mathbb{P}^2\) 上 6 点爆破 & 三次曲面（\(\mathbb{P}^3\) 中） \\
2 & \(\mathbb{P}^2\) 上 7 点爆破 &
四次曲面（\(\mathbb{P}(1, 1, 1, 2)\)） \\
1 & \(\mathbb{P}^2\) 上 8 点爆破 &
六次曲面（\(\mathbb{P}(1, 1, 2, 3)\)） \\
\end{longtable}
}

\emph{注：} \(d = 3\) 的 del Pezzo 曲面即为 \(\mathbb{P}^3\)
中的光滑三次曲面。著名的 27
条直线定理（Cayley-Salmon）断言：\(\mathbb{P}^3\)
中每个光滑三次曲面上恰好有 27 条直线。

\textbf{定理 86.2}（三维 Fano 簇的分类，Iskovskikh 1977-1979，Mori-Mukai
1981-2003）：三维光滑 Fano 簇恰好有 \textbf{105 个形变族}（deformation
families），按 Fano 指标 \(r\) 和 Picard 数 \(\rho\) 分类：

\begin{itemize}
\tightlist
\item
  Fano 指标 \(r = 4\)：仅 \(\mathbb{P}^3\)（1 族）；
\item
  Fano 指标 \(r = 3\)：二次超曲面 \(Q^3 \subset \mathbb{P}^4\)（1 族）；
\item
  Fano 指标 \(r = 2\)：包括 del Pezzo 三维簇（36 族，最大 Picard 数为
  10）；
\item
  Fano 指标 \(r = 1\)：其余 67 族（包括著名的素数 Fano 三维簇
  \(V_{22}\)，由 Mukai 发现）。
\end{itemize}

\emph{证明思路}：Iskovskikh
利用双有理几何的经典方法（线性系统、收缩映射）进行了初步分类。Mori 和
Mukai 运用 Mori
理论（极小模型纲领）的现代工具------特别是极射线收缩（extremal ray
contraction）方法------系统化地完成了分类。核心策略是：对给定的 Fano 簇
\(X\)，选择 Mori 收缩 \(\varphi: X \to Y\)，然后对 \(Y\)
应用归纳。关键步骤包括验证翻转（flip）不会发生（因为 \(-K_X\) 丰沛保证了
Mori 收缩的三种可能性中只有纤维型和除子型出现）。∎

\subsubsection{86.3 Mori
纤维空间与有理连通簇}\label{mori-ux7ea4ux7ef4ux7a7aux95f4ux4e0eux6709ux7406ux8fdeux901aux7c07}

\textbf{定义 86.3}（Mori 纤维空间）：\textbf{Mori 纤维空间}（Mori Fiber
Space）是带有纤维型极射线收缩 \(\varphi: X \to Y\) 的
\(\mathbb{Q}\)-因式簇 \(X\)，其中 \(X\) 仅有 terminal 奇点，\(Y\) 是
\(\mathbb{Q}\)-因式簇且 \(\dim Y < \dim X\)，一般纤维是 Fano 簇（带有
terminal 奇点）。

Mori 纤维空间是极小模型纲领（MMP）的三大输出类型之一。每个光滑射影簇通过
MMP 最终分解为 Mori 纤维空间（经由极小模型和翻转序列）。

\textbf{定义 86.4}（有理连通簇）：代数闭域上的光滑射影簇 \(X\)
称为\textbf{有理连通的}（rationally
connected），若其中的一般两点可以由有理曲线相连。

更精确地，存在有理映射
\(\pi: \mathbb{P}^1 \times \cdots \times \mathbb{P}^1 \dashrightarrow X\)
使其像覆盖 \(X\) 的开稠子集。

\textbf{定理 86.3}（Kollár-Miyaoka-Mori 定理，1992 年）：特征 0
的代数闭域上的光滑 Fano 簇是有理连通的。更一般地，光滑 Fano
簇是\textbf{有理链连通的}（rationally chain
connected）------任意两点可由若干条有理曲线的链相连。

\emph{证明思路}：利用 Mori 的弯曲-断裂引理（bend-and-break lemma）：由于
\(-K_X\) 是丰沛的，对于 \(X\) 上的一般曲线 \(C\)，有
\(-K_X \cdot C > 0\)。弯曲-断裂引理确保我们可以将穿过一般点的有理曲线''弯曲''（通过形变），使其收敛到一条可约曲线。这产生了足够多的有理曲线。再通过构造有理曲线的''组合链''来连接任意两点。Campana
独立地给出了等价结果。∎

\textbf{推论 86.4}：Fano 簇上的 Chow 群 \(\operatorname{CH}_0(X)\)
是平凡群------即 \(X\) 上任意两点模有理等价后相等。

\subsubsection{86.4 Luroth
问题与有理性}\label{luroth-ux95eeux9898ux4e0eux6709ux7406ux6027}

\textbf{Luroth 问题}（1876 年）：设
\(k \subset K \subset k(t_1, \ldots, t_n)\)
是嵌入在有理函数域中的子域。则 \(K\) 本身是否有理（即同构于某个 \(m\)
个变量的有理函数域）？

\(n = 1\) 时 Luroth 定理给出肯定回答（\(K \cong k(t)\)）。\(n = 2\) 时
Castelnuovo 定理（1894 年）在代数几何语言中断言：特征 0
的代数闭域上，单有理曲面必为有理曲面。

\(n \geq 3\) 时 Luroth
问题的回答是\textbf{否定的}。反例构造是代数几何的重大里程碑：

\textbf{定理 86.5}（三维 Luroth 问题的反例）：

\begin{enumerate}
\def\labelenumi{(\arabic{enumi})}
\item
  \textbf{Artin-Mumford 反例}（1972 年）：利用 Brauer 群的三阶挠元（扭
  3-形式）构造了特征 0 的代数闭域上的单有理三维簇 \(X\)，满足
  \(H^3(X, \mathbb{Z})_{\text{tors}} \neq 0\)。由于任何有理簇均有无挠的上同调，\(X\)
  不能是有理的。
\item
  \textbf{Clemens-Griffiths 反例}（1972 年）：三维光滑三次超曲面
  \(V_3 \subset \mathbb{P}^4\) 是单有理的，但其\textbf{中间 Jacobian}
  \(J(V_3)\)（二维复环面）不是 Jacobian 簇的直和项，故 \(V_3\)
  不是有理的。中间 Jacobian 是有理曲面不存在的不变量。
\item
  \textbf{Iskovskikh-Manin 反例}（1971 年）：四维光滑三次超曲面
  \(V_3 \subset \mathbb{P}^5\)
  不是有理的。证法基于双有理自同构群的有限性------证明
  \(\operatorname{Bir}(V_3)\)
  是有限群，而有理簇的双有理自同构群是无界大的（Cremona 群
  \(\operatorname{Cr}_n = \operatorname{Bir}(\mathbb{P}^n)\) 当
  \(n \geq 2\) 时无限）。
\end{enumerate}

这些反例共同说明：单有理性和有理性之间的鸿沟随维数增长而加大，高维有理性问题是极其微妙和困难的。

\bigskip\hrule\bigskip

\subsection{第87章：极化簇与 Hilbert
概型}\label{ux7b2c87ux7ae0ux6781ux5316ux7c07ux4e0e-hilbert-ux6982ux578b}

极化簇和 Hilbert 概型是现代代数几何中模空间理论的核心工具。Hilbert
概型将闭子概型的参数化问题形式化为射影概型结构，而极化簇的
GIT（几何不变量理论）则为构造代数簇的模空间提供了标准方法。

\subsubsection{87.1 Hilbert 多项式与 Hilbert
概型}\label{hilbert-ux591aux9879ux5f0fux4e0e-hilbert-ux6982ux578b}

\textbf{定义 87.1}（极化簇）：\textbf{极化簇}是一个对 \((X, L)\)，其中
\(X\) 是代数闭域 \(k\) 上的射影簇，\(L \in \operatorname{Pic}(X)\) 是
\(X\) 上的丰沛线丛（极化）。极化簇的 Hilbert 多项式为

\[P(m) = \chi(X, L^{\otimes m}) = \sum_{i=0}^{\dim X} (-1)^i \dim_k H^i(X, L^{\otimes m})\]

由 Serre 消没定理，当 \(m \gg 0\) 时，\(H^i(X, L^{\otimes m}) = 0\)
对所有 \(i > 0\) 成立，故

\[P(m) = \dim_k H^0(X, L^{\otimes m}), \quad m \gg 0\]

由 Riemann-Roch 定理（Hirzebruch 版本），Hilbert 多项式的主导项为

\[P(m) = \frac{L^n}{n!} m^n + \text{（低次项）}\]

其中 \(n = \dim X\)，\(L^n = (L^{\otimes n})\) 的相交数即为 \(L\)
的次数。

\textbf{定理 87.1}（Grothendieck 的 Hilbert 概型定理，1961 年）：设
\(\mathbb{P}^N_k\) 是域 \(k\) 上的射影空间。对任意数值多项式
\(P(t) \in \mathbb{Q}[t]\)，存在射影概型
\(\operatorname{Hilb}_P(\mathbb{P}^N_k)\)，其 \(k\)-点自然一一对应于
\(\mathbb{P}^N_k\) 中 Hilbert 多项式为 \(P(t)\)
的闭子概型。而且该对应具有\textbf{万有性质}：存在万有闭子概型

\[\mathcal{Z} \subseteq \operatorname{Hilb}_P(\mathbb{P}^N_k) \times \mathbb{P}^N_k\]

使得对任意概型 \(S\) 和任意 \(S\)-平坦的闭子概型
\(Y \subseteq S \times \mathbb{P}^N_k\)（纤维的 Hilbert 多项式均为
\(P(t)\)），存在唯一的态射
\(f: S \to \operatorname{Hilb}_P(\mathbb{P}^N_k)\) 使得
\(Y = (f \times \operatorname{id})^* \mathcal{Z}\)。

\emph{证明思路}：核心是利用 Grothendieck 的 Quot 概型（Quot
scheme）理论。考虑多项式环 \(S = k[x_0, \ldots, x_N]\) 和商
\(S / I\)（\(I\) 为齐次理想）。Hilbert 基定理保证了对固定的 Hilbert
多项式 \(P(t)\)，存在统一的 \(m_0\) 使得对所有 \(m \geq m_0\)，理想
\(I_m\) 由 \(I_{m_0}\) 生成。由此可将 Hilbert 概型嵌入 Grassmann 概型
\(\operatorname{Grass}(P(m_0), \dim S_{m_0})\)
中，再证明该嵌入是闭浸入。∎

\subsubsection{87.2
模空间的构造：稳定曲线}\label{ux6a21ux7a7aux95f4ux7684ux6784ux9020ux7a33ux5b9aux66f2ux7ebf}

\textbf{定理 87.2}（Deligne-Mumford 稳定曲线模空间
\(\overline{M}_g\)）：设 \(g \geq 2\)。亏格 \(g\)
的\textbf{稳定曲线}的模叠 \(\overline{\mathcal{M}}_g\)
是光滑、真（proper）且不可约的 Deligne-Mumford 叠。其粗模空间
\(\overline{M}_g\) 是 \(3g-3\) 维射影簇。

一条曲线 \(C\) 是\textbf{稳定的}，若：(1) \(C\)
是连通的约化（reduced）曲线，仅有结点（node）奇点；(2) \(C\)
的光滑有理分支包含至少 3 个特殊点（与曲线其余部分的相交点）；(3)
其典范层 \(\omega_C\) 是丰沛的（等价地，对偶图是稳定的）。

\emph{构造思路}：利用 Hilbert 概型。固定充分大的 \(m\)（如
\(m = 3\)），考虑 \(\omega_C^{\otimes m}\) 的整体截面给出了嵌入
\(C \hookrightarrow \mathbb{P}^{N}\)（由 Riemann-Roch
定理，\(N = (2m-1)(g-1)-1\)）。所有稳定曲线的 \(m\)-典范嵌入构成 Hilbert
概型中的局部闭子集。Hilbert-Mumford 稳定性条件保证该子集对 \(PGL(N+1)\)
作用是 GIT 半稳定的，取 GIT 商即得 \(\overline{M}_g\)。∎

\textbf{极化 Abelian 簇的模空间}：亏格 \(g\) 的主极化 Abel 簇的粗模空间
\(\mathcal{A}_g\) 由 Siegel 上半空间 \(\mathbb{H}_g\) 经
\(Sp(2g, \mathbb{Z})\) 作用取商得到。\(\mathcal{A}_g\) 是 \(g(g+1)/2\)
维拟射影簇。其紧化包括： - \textbf{Satake 紧化}（Baily-Borel
紧化，极小紧化）； - \textbf{toroidal 紧化}（由 Faltings-Chai
描述，可容许退化）； - \textbf{二阶 Voronoi 紧化}（Alexeev，完全显式）。

\subsubsection{87.3 GIT 稳定性与 Hilbert-Mumford
数值准则}\label{git-ux7a33ux5b9aux6027ux4e0e-hilbert-mumford-ux6570ux503cux51c6ux5219}

\textbf{定义 87.2}（GIT 的稳定性条件，Mumford 1965 年）：设 \(G\)
是代数闭域 \(k\) 上的简约代数群，\(L\) 是射影概型 \(X\) 上的丰沛
\(G\)-等变线丛。点 \(x \in X\) 称为：

\begin{itemize}
\tightlist
\item
  \textbf{半稳定的}（semistable），若存在 \(n > 0\) 和 \(G\)-不变截面
  \(s \in H^0(X, L^{\otimes n})^G\) 使得 \(s(x) \neq 0\)；
\item
  \textbf{稳定的}（stable），若 \(x\) 是半稳定的，轨道 \(G \cdot x\) 在
  \(X^{ss}\) 中是闭的，且稳定子群 \(G_x\) 是有限群；
\item
  \textbf{真稳定的}（properly stable），若 \(x\) 是稳定的且稳定子群
  \(G_x\) 的维数为 0。
\end{itemize}

半稳定点集 \(X^{ss}\) 是 \(X\) 的开子集。GIT 商
\(X /\!/ G = \operatorname{Proj}\left(\bigoplus_{n \geq 0} H^0(X, L^{\otimes n})^G\right)\)
是射影簇。

\textbf{定理 87.3}（Hilbert-Mumford 数值准则）：设 \(G\)
是简约代数群，\(\lambda: \mathbb{G}_m \to G\) 是单参数子群（1-PSG）。对
\(x \in X\)，定义 Hilbert-Mumford 数值

\[\mu(x, \lambda) = -\operatorname{ord}_{t=0}(\lambda(t) \cdot \tilde{x})\]

其中 \(\tilde{x}\) 是 \(x\) 在某个 \(G\)-等变化下的提升。则

\begin{itemize}
\tightlist
\item
  \(x\) 是 \textbf{半稳定的} 当且仅当对所有 1-PSG
  \(\lambda\)，\(\mu(x, \lambda) \geq 0\)；
\item
  \(x\) 是 \textbf{稳定的} 当且仅当对所有非平凡 1-PSG
  \(\lambda\)，\(\mu(x, \lambda) > 0\)。
\end{itemize}

\emph{证明思路}：若存在 1-PSG \(\lambda\) 使 \(\mu(x, \lambda) < 0\)，则
\(\lim_{t \to 0} \lambda(t) \cdot x\) 存在且不属于 \(x\)
的轨道闭包，从而轨道不是闭的。反之，若所有
\(\mu(x, \lambda) \geq 0\)，利用 Hilbert-Mumford 定理（任何偏离
G-不变截面的方向必然由某个 1-PSG 的负数值体现）可得 \(x\) 是半稳定的。∎

\textbf{应用于稳定曲线}：曲线的 GIT 稳定性（Hilbert-Mumford
稳定性）恰好对应于 Deligne-Mumford 稳定性条件。一条曲线 \(C\) 在 GIT
意义下是半稳定的（关于 \(m\)-典范嵌入）当且仅当它是 Deligne-Mumford
意义下的稳定曲线。

\bigskip\hrule\bigskip

\subsection{第88章：Arakelov
几何}\label{ux7b2c88ux7ae0arakelov-ux51e0ux4f55}

Arakelov 几何将代数几何的概型方法推广到算术情形，在
\(\operatorname{Spec} \mathbb{Z}\)
上构造''算术曲面''，并在每个无穷远点处附加 Hermite
度量以赋予''紧性''。这是将代数几何方法与解析工具相结合的交叉领域，其核心应用包括
Faltings 对 Mordell 猜想的证明。

\subsubsection{88.1 算术曲面与 Hermite
度量}\label{ux7b97ux672fux66f2ux9762ux4e0e-hermite-ux5ea6ux91cf}

\textbf{定义
88.1}（算术曲面）：\textbf{算术曲面}是一个平坦、有限型的态射
\(\pi: \mathcal{X} \to \operatorname{Spec} \mathbb{Z}\)，其一般纤维
\(\mathcal{X}_{\mathbb{Q}} = \mathcal{X} \times_{\operatorname{Spec} \mathbb{Z}} \operatorname{Spec} \mathbb{Q}\)
是一条光滑射影曲线。实数纤维
\(\mathcal{X}(\mathbb{C}) = \mathcal{X} \times_{\operatorname{Spec} \mathbb{Z}} \operatorname{Spec} \mathbb{C}\)
是紧 Riemann 面（不一定连通）。在每个连通分支
\(\mathcal{X}_\sigma\)（\(\sigma \in \mathcal{X}(\mathbb{C})\)
的连通分支）上需装备 \textbf{Kähler 度量}------称为 Hermite 度量或
Kähler 形式 \(\omega_\sigma\)。

在每个非无穷远点 \(p\)（即正素数），\(\mathcal{X}\) 的纤维
\(\mathcal{X}_p = \mathcal{X} \times_{\operatorname{Spec} \mathbb{Z}} \operatorname{Spec} \mathbb{F}_p\)
是有限域 \(\mathbb{F}_p\) 上的代数曲线（带有可能的奇点）。

\textbf{定义 88.2}（算术除子）：算术曲面 \(\mathcal{X}\) 上的
\textbf{算术除子} 是有限形式和

\[D = \sum_{P} n_P [P] + \sum_{\sigma} \lambda_\sigma [\mathcal{X}_\sigma]\]

其中 \(P\) 遍历 \(\mathcal{X}\)
的不可约水平除子（算术除子的有限部分），\(\lambda_\sigma \in \mathbb{R}\)
为无穷远除子的系数。此类除子形式上类比于紧 Riemann
面上的除子，但额外考虑了无穷远度量。

\subsubsection{88.2 算术 Chow
群与算术相交理论}\label{ux7b97ux672f-chow-ux7fa4ux4e0eux7b97ux672fux76f8ux4ea4ux7406ux8bba}

\textbf{定义 88.3}（算术 Chow 群）：算术曲面 \(\mathcal{X}\) 的
\textbf{算术 Chow 群} \(\widehat{\operatorname{CH}}^p(\mathcal{X})\)
由以下数据生成：代数圈 \(Z\)（\(\mathcal{X}\) 的子概型）配以沿
\(\mathcal{X}(\mathbb{C})\) 每个连通分支的 \textbf{Green 流}
\(g_\sigma\)（满足微分方程
\(\partial \bar{\partial} g_\sigma + \delta_Z = \omega_Z\)
的流），模去有理等价和无穷远处的全纯函数的贡献。

算术 Chow 群 \(\widehat{\operatorname{CH}}^*(\mathcal{X})\)
构成分次环，配以算术相交积

\[\widehat{\operatorname{CH}}^p(\mathcal{X}) \times \widehat{\operatorname{CH}}^q(\mathcal{X}) \to \widehat{\operatorname{CH}}^{p+q}(\mathcal{X}) \otimes \mathbb{R}\]

此积在有限位与通常的相交积一致，在无穷远处与 Green 流的 star 积一致。

\textbf{定义 88.4}（算术相交数）：若 \(\mathcal{X}\)
是算术曲面，\(D, E\) 是两个算术除子，则其\textbf{算术相交数}
\(\widehat{D \cdot E} \in \mathbb{R}\) 定义为

\[\widehat{D \cdot E} = D_{\text{fin}} \cdot E_{\text{fin}} + \text{（无穷远处的贡献）}\]

无穷远处的贡献由 Green 函数 \(g(D, E)\) 在 \(\mathcal{X}(\mathbb{C})\)
上的积分给出。

\textbf{定理 88.1}（算术 Hodge 指标定理，Faltings 1984
年）：在算术曲面上，相交配对 \(\widehat{\cdot}\) 满足与 Hodge
指标定理类似的符号性质。

\subsubsection{88.3 算术 Riemann-Roch
定理}\label{ux7b97ux672f-riemann-roch-ux5b9aux7406}

\textbf{定理 88.2}（算术 Riemann-Roch 定理，Gillet-Soulé 1992 年）：设
\(\pi: \mathcal{X} \to \operatorname{Spec} \mathbb{Z}\)
是算术曲面，\(\overline{E}\) 是 \(\mathcal{X}\) 上的 Hermite
向量丛。则算术 Euler 示性数

\[\widehat{\chi}(\overline{E}) = \sum_{i=0}^2 (-1)^i \widehat{\deg}(\det R^i \pi_* \overline{E})\]

等于由 Gillet-Soulé 的 Chern 形式
\(\widehat{\operatorname{ch}}(\overline{E})\) 和 Todd 形式
\(\widehat{\operatorname{Td}}(T_{\mathcal{X}/\mathbb{Z}})\) 在算术 Chow
环中的相交积分

\[\widehat{\chi}(\overline{E}) = \int_{\mathcal{X}} \widehat{\operatorname{ch}}(\overline{E}) \cdot \widehat{\operatorname{Td}}(T_{\mathcal{X}/\mathbb{Z}})\]

\emph{证明思路}：这一定理将 Grothendieck-Riemann-Roch
定理推广到算术情形，其证明的关键是处理\textbf{解析挠率}（analytic
torsion, Ray-Singer 挠率）的贡献。算术 Riemann-Roch
的核心新特征在于：等式中的\textbf{解析挠率项} \(T(\overline{E})\) 由
\(\det \Delta\)（Laplacian
的行列式）决定，体现了算术曲面在无穷远处的度量信息。Gillet 和 Soulé
构造了算术 K-理论和算术 Chow 环之间的 Chern
特征标，建立了层上同调与微分几何之间的桥梁。∎

\subsubsection{88.4 算术 Hilbert-Samuel
定理}\label{ux7b97ux672f-hilbert-samuel-ux5b9aux7406}

\textbf{定理 88.3}（算术 Hilbert-Samuel 定理，Zhang 1995
年，Abbes-Bouche 1995 年）：设 \(\mathcal{X}\)
是算术曲面，\(\overline{L}\) 是 \(\mathcal{X}\) 上的 Hermite
线丛（丰沛线丛配以正曲率度量）。则

\[\widehat{\chi}(\overline{L}^{\otimes n}) = \frac{\widehat{\deg}(\overline{L})}{2} n^2 + O(n \log n)\]

其中 \(\widehat{\deg}(\overline{L})\) 是 \(\overline{L}\)
的算术次数。这与通常的 Hilbert-Samuel
定理（\(P(n) = \frac{L^n}{n!} n^n + \cdots\)）在形式上是平行的。

\emph{证明思路}：算术 Hilbert-Samuel 的核心是控制 \(\mathcal{X}\)
上全局截面空间的范数------这是算术曲面上度量信息的直接体现。利用算术
Riemann-Roch 定理和 Laplace 算子的谱渐近性质（Weyl 律在紧 Riemann
面上的应用），可以提取出主导项和误差项。∎

\subsubsection{88.5 应用：Mordell 猜想与算术 Bogomolov
猜想}\label{ux5e94ux7528mordell-ux731cux60f3ux4e0eux7b97ux672f-bogomolov-ux731cux60f3}

\textbf{Arakelov 几何的核心应用：}

\begin{enumerate}
\def\labelenumi{\arabic{enumi}.}
\item
  \textbf{Faltings 对 Mordell 猜想的证明}（1983 年）：设 \(C\) 是数域
  \(K\) 上亏格 \(g \geq 2\) 的光滑射影曲线，则 \(C(K)\)
  是有限集。Faltings 的核心洞察是将 \(C(K)\)
  的有界高度与算术曲面上自相交数 \(\omega_{C/Y}^2\) 的 Arakelov
  版本关联。利用算术 Hodge 指标定理和高度不等式，证明若有无限多个 \(K\)
  点，则 \(\widehat{\operatorname{div}}(P)^2\)
  取任意大的值，与算术相交理论中的有界性原理矛盾。
\item
  \textbf{Vojta 猜想}：Vojta 猜想（1987 年）将 Mordell
  猜想推广到高维，断言数域上的对数一般型簇上的有理点具有稀疏性质------每个有理点集包含于有限个真子簇的并集中。Vojta
  的证明框架将 Arakelov 几何的高度不等式与 Nevanlinna
  理论中的第二主定理通过 Diophantine 相应类比结合。
\item
  \textbf{算术 Bogomolov 猜想}：若曲线 \(C\) 嵌入其 Jacobian 簇
  \(J(C)\)，且 \(J(C)\) 上的 Néron-Tate 高度 \(\hat{h}\) 在 \(C\)
  上有一组严格正的最小累积点（但不全为零），则这些最小高度与某种
  Arakelov 自相交数相关。Ullmo（1998 年）和 Zhang（1998
  年）独立地利用等分布定理（equidistribution）的算术版本证明了该猜想。
\item
  \textbf{Zhang 的最小高度}：Zhang 提出了
  \textbf{本质最小高度}（essential minimum）：对代数簇 \(X\)
  上的高度函数，定义
  \(\mu_{\text{ess}}(X) = \inf_{U \subset X \text{开稠}} \sup_{P \in U(\overline{K})} h(P)\)。Zhang
  的最小高度定理断言：\(\mu_{\text{ess}}(X) \geq 0\)
  且等号成立的充要条件由 Galois 轨道的等分布性质刻画。
\item
  \textbf{有效 Mordell 定理}：de Jong-Rémond 和 Binyamini 等人利用
  Arakelov 几何建立了有效高度的有界性------给出了 \(C(K)\)
  大小的显式上界（虽然当前界很大），这比原始的 Faltings 定理更构造性。
\end{enumerate}

\bigskip\hrule\bigskip

\emph{本卷从经典代数几何（仿射簇、射影簇、Hilbert
零点定理）出发，引入现代概形语言和层上同调方法。仿射代数簇与有限生成整代数的范畴等价建立了代数与几何的桥梁；概形理论将这一对应推广到所有交换环，统一了代数几何与数论；Serre
定理和 Serre 对偶为凝聚层上同调的计算提供了核心工具；Riemann-Roch
定理及其推广是代数几何最重要的定理之一，深刻影响了代数曲线、曲面和高维簇的分类理论。为后续数论
II（V16）中椭圆曲线和模形式的学习提供了几何基础。}

\newpage

\section{卷十六：数论 II}\label{ux5377ux5341ux516dux6570ux8bba-ii}

\begin{quote}
数论 II 从初等数论（V1 Ch 5）和代数数论（V12 Ch
71）出发，引入解析方法（Dirichlet 级数、Riemann zeta
函数）、模形式理论和椭圆曲线，建立现代数论的核心工具链。本卷以素数定理的完整证明、模形式的
Hecke 理论、椭圆曲线的 Mordell-Weil 定理和类域论初步为目标。
\end{quote}

\bigskip\hrule\bigskip

\subsection{第85章：解析数论}\label{ux7b2c85ux7ae0ux89e3ux6790ux6570ux8bba}

解析数论用复分析（V9）和 Fourier
分析（V8）的工具研究整数问题。本章的核心是 Dirichlet 级数、Riemann zeta
函数和素数定理的完整证明。

\subsubsection{85.1 算术函数与 Dirichlet
卷积}\label{ux7b97ux672fux51fdux6570ux4e0e-dirichlet-ux5377ux79ef}

\textbf{定义 85.1}（算术函数）：\textbf{算术函数}是
\(f: \mathbb{N}^+ \to \mathbb{C}\) 的函数。常用的算术函数包括： -
\(\mathbf{1}(n) = 1\)（常数函数） -
\(\operatorname{id}(n) = n\)（恒等函数） - \(\mu(n)\)：Möbius
函数（\(\mu(1) = 1\)；若 \(n\) 有平方因子则 \(\mu(n) = 0\)；否则
\(\mu(n) = (-1)^k\)，\(k\) 为素因子个数） - \(\varphi(n)\)：Euler
\(\varphi\) 函数（Euler totient
函数，\(\varphi(n) = n \prod_{p \mid n} (1 - 1/p)\)） -
\(\Lambda(n)\)：von Mangoldt 函数（\(\Lambda(n) = \log p\) 若
\(n = p^k\)；否则 \(\Lambda(n) = 0\)） - \(\tau(n)\)：除数函数（\(n\)
的正因子个数） - \(\sigma(n)\)：除数和函数（\(n\) 的正因子之和）

\textbf{定义 85.2}（Dirichlet 卷积）：两个算术函数 \(f, g\) 的
\textbf{Dirichlet 卷积} 为

\[(f * g)(n) = \sum_{d \mid n} f(d) g(n/d)\]

\textbf{命题 85.1}：Dirichlet 卷积满足： 1.
结合律：\((f * g) * h = f * (g * h)\) 2. 交换律：\(f * g = g * f\) 3.
单位元 \(\delta\)（\(\delta(1) = 1\)，\(\delta(n) = 0\) 对 \(n > 1\)）
4. \(f\) 可逆当且仅当 \(f(1) \neq 0\)

\emph{证明}：1 和 2 由和号重排验证。3 显然。4：构造
\(g(n) = -\frac{1}{f(1)} \sum_{d \mid n, d < n} f(n/d) g(d)\)
归纳定义。∎

\textbf{命题 85.2}（Möbius 反演）：\(\mathbf{1} * \mu = \delta\)。即
\(\sum_{d \mid n} \mu(d) = 0\)（\(n > 1\)）。等价地，\(f(n) = \sum_{d \mid n} g(d)\)
当且仅当 \(g(n) = \sum_{d \mid n} \mu(d) f(n/d)\)。

\emph{证明}：\(\mathbf{1} * \mu\) 在素数幂 \(p^k\) 上的值为
\(1 - 1 = 0\)（\(k \geq 1\)）。由积性，一般 \(n\) 也成立。∎

\subsubsection{85.2 Dirichlet 级数与 Euler
乘积}\label{dirichlet-ux7ea7ux6570ux4e0e-euler-ux4e58ux79ef}

\textbf{定义 85.3}（Dirichlet 级数）：算术函数 \(f\) 的
\textbf{Dirichlet 级数} 为

\[D_f(s) = \sum_{n=1}^\infty \frac{f(n)}{n^s}\]

其中 \(s = \sigma + it \in \mathbb{C}\)。若 \(f(n) = O(n^\alpha)\)，则
\(D_f(s)\) 在 \(\Re(s) > \alpha + 1\) 内绝对收敛。

\textbf{命题 85.3}（Dirichlet 卷积与级数乘积）：\((f * g)(n)\) 的
Dirichlet 级数是 \(D_f(s) D_g(s)\)。

\emph{证明}：\(D_f(s) D_g(s) = \sum_{a=1}^\infty \frac{f(a)}{a^s} \sum_{b=1}^\infty \frac{g(b)}{b^s} = \sum_{n=1}^\infty \frac{\sum_{ab=n} f(a) g(b)}{n^s} = D_{f*g}(s)\)。∎

\textbf{定义 85.4}（Euler 乘积）：若 \(f\)
是\textbf{积性函数}（\(f(mn) = f(m)f(n)\) 当 \(\gcd(m, n) = 1\)），则

\[D_f(s) = \prod_{p \text{ 素数}} \left(1 + \frac{f(p)}{p^s} + \frac{f(p^2)}{p^{2s}} + \cdots\right)\]

若 \(f\) 是完全积性函数，则
\(D_f(s) = \prod_p \left(1 - \frac{f(p)}{p^s}\right)^{-1}\)。

\subsubsection{85.3 Riemann zeta 函数}\label{riemann-zeta-ux51fdux6570}

\textbf{定义 85.5}（Riemann zeta 函数）：对 \(\Re(s) > 1\)，

\[\zeta(s) = \sum_{n=1}^\infty \frac{1}{n^s}\]

\textbf{定理 85.4}（Euler 乘积公式）：对 \(\Re(s) > 1\)，

\[\zeta(s) = \prod_{p \text{ 素数}} \frac{1}{1 - p^{-s}}\]

\emph{证明}：\(\mathbf{1}\) 是完全积性函数，由 Euler 乘积公式即得。∎

\textbf{定理 85.5}（\(\zeta(s)\) 的解析延拓）：\(\zeta(s)\)
可解析延拓为整个复平面上的亚纯函数，仅在 \(s = 1\)
处有一个单极点（留数为 \(1\)）。函数方程：

\[\zeta(s) = 2^s \pi^{s-1} \sin\left(\frac{\pi s}{2}\right) \Gamma(1-s) \zeta(1-s)\]

或等价地，令
\(\xi(s) = \frac{1}{2} s(s-1) \pi^{-s/2} \Gamma(s/2) \zeta(s)\)，则
\(\xi(s) = \xi(1-s)\) 且 \(\xi(s)\) 是整函数。

\emph{证明思路}：从 \(\zeta(s)\) 的积分表示
\(\zeta(s) = \frac{1}{\Gamma(s)} \int_0^\infty \frac{x^{s-1}}{e^x - 1} dx\)
出发，利用 \(\theta\) 函数的模性质（Poisson 求和公式）得到函数方程。∎

\textbf{定义 85.6}（\(\zeta(s)\) 的零点）： -
\textbf{平凡零点}：\(s = -2, -4, -6, \ldots\)（由 \(\sin(\pi s/2)\)
项产生） - \textbf{非平凡零点}：\(0 < \Re(s) < 1\)（临界带）内的零点

\textbf{Riemann 猜想}（未证明）：\(\zeta(s)\) 的所有非平凡零点都满足
\(\Re(s) = 1/2\)（临界线）。

\subsubsection{85.4 素数定理}\label{ux7d20ux6570ux5b9aux7406}

\textbf{定理 85.6}（素数定理）：设
\(\pi(x) = \#\{p \leq x : p \text{ 素数}\}\)。则

\[\pi(x) \sim \frac{x}{\log x} \quad (x \to \infty)\]

即 \(\lim_{x \to \infty} \frac{\pi(x)}{x / \log x} = 1\)。

\emph{证明思路}（Hadamard 和 de la Vallée Poussin，1896）：

\begin{enumerate}
\def\labelenumi{\arabic{enumi}.}
\item
  定义 Chebyshev \(\psi\)
  函数：\(\psi(x) = \sum_{n \leq x} \Lambda(n)\)。素数定理等价于
  \(\psi(x) \sim x\)。
\item
  由 Perron
  公式：\(\psi(x) = \frac{1}{2\pi i} \int_{c - i\infty}^{c + i\infty} \left(-\frac{\zeta'(s)}{\zeta(s)}\right) \frac{x^s}{s} ds\)（\(c > 1\)）。
\item
  将积分线向左平移，穿越 \(s = 1\) 处的极点（贡献 \(x\)）和 \(\zeta(s)\)
  的零点。需要证明 \(\zeta(s)\) 在 \(\Re(s) = 1\) 上无零点（即
  \(\zeta(1 + it) \neq 0\)）。
\item
  \(\zeta(1 + it) \neq 0\) 的证明：使用三角不等式
  \(3 + 4\cos\theta + \cos 2\theta = 2(1 + \cos\theta)^2 \geq 0\)。若
  \(\zeta(1 + it_0) = 0\)，考虑 \(\log \zeta(s)\) 的展开，推出矛盾。
\item
  零点密度估计：证明 \(\psi(x) - x\) 的误差项为 \(o(x)\)，从而
  \(\psi(x) \sim x\)。∎
\end{enumerate}

\textbf{定理 85.7}（素数定理的误差项）：若 Riemann 猜想成立，则
\(\pi(x) = \operatorname{Li}(x) + O(\sqrt{x} \log x)\)，其中
\(\operatorname{Li}(x) = \int_2^x \frac{dt}{\log t}\)。

\textbf{定理 85.8}（等差数列中的素数定理，Dirichlet）：若
\(\gcd(a, q) = 1\)，则等差数列 \(a, a+q, a+2q, \ldots\)
中包含无穷多个素数。且

\[\pi(x; q, a) \sim \frac{1}{\varphi(q)} \frac{x}{\log x} \quad (x \to \infty)\]

\emph{证明思路}：引入 Dirichlet 特征 \(\chi\)（模 \(q\) 的群同态
\((\mathbb{Z}/q\mathbb{Z})^\times \to \mathbb{C}^\times\)），构造
Dirichlet \(L\)-函数
\(L(s, \chi) = \sum_{n=1}^\infty \frac{\chi(n)}{n^s} = \prod_{p \nmid q} (1 - \chi(p)p^{-s})^{-1}\)（\(L\)-函数的
Euler 乘积）。关键三步：(1) 利用特征的正交性
\(\frac{1}{\varphi(q)} \sum_\chi \overline{\chi(a)} \chi(n) = \mathbf{1}_{n \equiv a \pmod q}\)，将算术级数中的素数计数函数
\(\psi(x; q, a)\) 表示为所有
\(L\)-函数对数导数的线性组合：\(\psi(x; q, a) \sim \frac{1}{\varphi(q)} \sum_\chi \overline{\chi(a)} \sum_{n \leq x} \chi(n)\Lambda(n)\)。(2)
证明对非主特征 \(\chi\)，\(L(1, \chi) \neq 0\)：若 \(\chi\)
是复特征（\(\chi^2 \neq \chi_0\)），利用乘积
\(L(s, \chi_0)^3 |L(s, \chi)|^4 |L(s, \chi^2)|^2\) 在 \(s=1\)
附近的零点阶数分析导出矛盾；若 \(\chi\) 是实特征，需更精细的 Dirichlet
类数公式。(3) 对每个特征 \(\chi\)，用 Perron
公式和围道积分类似素数定理的证明：\(L(s, \chi)\) 在 \(\Re(s) = 1\)
上无零点（由 (2) 保证），从而得到
\(\sum_{n \leq x} \chi(n)\Lambda(n) = o(x)\)（\(\chi \neq \chi_0\)）和
\(\sum_{n \leq x} \chi_0(n)\Lambda(n) \sim x\)（主特征），代入特征反演公式即得
\(\psi(x; q, a) \sim x/\varphi(q)\)。∎

\subsubsection{85.5 堆垒数论}\label{ux5806ux5792ux6570ux8bba}

堆垒数论（Additive Number
Theory）研究将整数表示为特定类型整数之和的经典问题。两个核心问题是：（1）\textbf{Goldbach猜想}：每个大于2的偶数可表为两个素数之和。陈景润于1966年证明了''1+2''------每个充分大的偶数可表为一个素数和一个至多两个素数乘积的数之和，这是迄今为止最好的结果。Helfgott于2013年完整证明了三素数定理：每个大于5的奇数可表为三个素数之和。（2）\textbf{Waring问题}（1770年）：对每个正整数\(k\)，存在\(g(k)\)使得每个正整数可表为至多\(g(k)\)个\(k\)次幂之和。Lagrange四平方和定理即\(g(2)=4\)。记\(G(k)\)为充分大整数所需的最少\(k\)次幂个数，已知\(G(2)=4\)，而\(G(k)\)的确切值对\(k \geq 3\)仍是重大未解问题。处理这类问题的核心工具是\textbf{圆法}（Hardy-Littlewood-Ramanujan，1920年代），其基本思想是将计数函数表示为单位圆上的积分，利用有理点附近的主项渐近展开与余项的三角和估计得到结果。

\subsubsection{85.6 连分数理论}\label{ux8fdeux5206ux6570ux7406ux8bba}

实数\(x\)的\textbf{简单连分数}展开为

\[x = a_0 + \cfrac{1}{a_1 + \cfrac{1}{a_2 + \cdots}} = [a_0; a_1, a_2, \ldots]\]

其中\(a_0 \in \mathbb{Z}\)，\(a_i \in \mathbb{N}^+\)（\(i \geq 1\)）。若展开有限，则\(x \in \mathbb{Q}\)；反之亦然。\textbf{Lagrange定理}（1770年）：二次无理数（有理系数二次方程的根）的连分数展开是最终周期的，这一性质刻画了二次无理数。收敛子\(p_n/q_n = [a_0; a_1, \ldots, a_n]\)是\(x\)的\textbf{最佳有理逼近}，满足\(|q x - p| < |q' x - p'|\)对所有分母\(q' \leq q\)且\(p'/q' \neq p_n/q_n\)成立。连分数与\textbf{Pell方程}\(x^2 - Dy^2 = \pm 1\)的自然联系是：\(\sqrt{D}\)的连分数周期\(r\)给出了基本解\((x_1, y_1) = (p_{r-1}, q_{r-1})\)（当\(r\)为偶数时）或\((p_{2r-1}, q_{2r-1})\)（当\(r\)为奇数时）。

\subsubsection{85.7 数的几何}\label{ux6570ux7684ux51e0ux4f55}

\textbf{数的几何}（Geometry of
Numbers）由Minkowski于1896年创立，用几何方法研究格点分布与数论问题。设\(\Lambda \subseteq \mathbb{R}^n\)是\textbf{格}（discrete
subgroup），即\(\Lambda = \{\sum a_i v_i : a_i \in \mathbb{Z}\}\)，其协体积\(\operatorname{covol}(\Lambda) = |\det(v_1, \ldots, v_n)|\)。\textbf{Minkowski第一定理}：若\(K \subseteq \mathbb{R}^n\)是中心对称凸体且\(\operatorname{vol}(K) > 2^n \operatorname{covol}(\Lambda)\)，则\(K \cap (\Lambda \setminus \{0\}) \neq \varnothing\)。\textbf{Minkowski第二定理}：对中心对称凸体\(K\)，其连续极小\(\lambda_1 \leq \cdots \leq \lambda_n\)满足\(\frac{2^n}{n!} \operatorname{covol}(\Lambda) \leq \lambda_1 \cdots \lambda_n \cdot \operatorname{vol}(K) \leq 2^n \operatorname{covol}(\Lambda)\)。经典应用包括：（1）Dirichlet单位定理的几何证明------对数嵌入\(\ell: \mathcal{O}_K^\times \to \mathbb{R}^{r_1+r_2}\)将单位群映为秩\(r_1+r_2-1\)的格；（2）Minkowski判别式界------\(|\operatorname{disc}(K)| \geq (\pi/4)^{2r_2} (n^n/n!)^2\)，由此可推出\(\mathbb{Q}\)不存在非平凡的无处分歧扩张。

\bigskip\hrule\bigskip

\subsection{第86章：模形式导论}\label{ux7b2c86ux7ae0ux6a21ux5f62ux5f0fux5bfcux8bba}

模形式是全纯函数，在模群 \(\operatorname{SL}_2(\mathbb{Z})\)
的作用下具有优美的变换性质。它们是连接数论、表示论和代数几何的桥梁。

\subsubsection{86.1
模群与上半平面}\label{ux6a21ux7fa4ux4e0eux4e0aux534aux5e73ux9762}

\textbf{定义 86.1}（上半平面）：\textbf{上半平面}
\(\mathbb{H} = \{z \in \mathbb{C} : \Im(z) > 0\}\)。

\textbf{定义 86.2}（模群）：\textbf{模群}
\(\operatorname{SL}_2(\mathbb{Z}) = \left\{\begin{pmatrix} a & b \\ c & d \end{pmatrix} : a, b, c, d \in \mathbb{Z}, ad - bc = 1\right\}\)。模群通过分式线性变换作用于
\(\mathbb{H}\)：

\[\gamma \cdot z = \frac{az + b}{cz + d}, \quad \gamma = \begin{pmatrix} a & b \\ c & d \end{pmatrix} \in \operatorname{SL}_2(\mathbb{Z})\]

\textbf{定义 86.3}（基本域）：\(\operatorname{SL}_2(\mathbb{Z})\) 在
\(\mathbb{H}\) 上的作用的基本域为

\[\mathcal{F} = \left\{z \in \mathbb{H} : |z| \geq 1, |\Re(z)| \leq \frac{1}{2}\right\}\]

\(\operatorname{SL}_2(\mathbb{Z})\) 由
\(S = \begin{pmatrix} 0 & -1 \\ 1 & 0 \end{pmatrix}\)（\(z \mapsto -1/z\)）和
\(T = \begin{pmatrix} 1 & 1 \\ 0 & 1 \end{pmatrix}\)（\(z \mapsto z + 1\)）生成。

\subsubsection{86.2
模形式的定义}\label{ux6a21ux5f62ux5f0fux7684ux5b9aux4e49}

\textbf{定义 86.4}（权 \(k\) 的因子）：对
\(\gamma = \begin{pmatrix} a & b \\ c & d \end{pmatrix} \in \operatorname{SL}_2(\mathbb{Z})\)，定义\textbf{自守因子}
\(j(\gamma, z) = cz + d\)。

\textbf{定义 86.5}（权 \(k\) 模形式）：全纯函数
\(f: \mathbb{H} \to \mathbb{C}\) 称为
\(\operatorname{SL}_2(\mathbb{Z})\) 上的\textbf{权 \(k\) 模形式}（\(k\)
为偶数），如果

\begin{enumerate}
\def\labelenumi{\arabic{enumi}.}
\tightlist
\item
  \(f(\gamma \cdot z) = (cz + d)^k f(z)\) 对所有
  \(\gamma \in \operatorname{SL}_2(\mathbb{Z})\)
\item
  \(f\) 在无穷远处全纯（即 \(f(z)\) 在 \(\Im(z) \to \infty\) 时有
  Fourier 展开
  \(f(z) = \sum_{n=0}^\infty a_n q^n\)，\(q = e^{2\pi i z}\)，无负指数项）
\end{enumerate}

若 \(a_0 = 0\)，则 \(f\) 称为\textbf{尖点形式}（cusp form）。

\textbf{定义 86.6}（模形式空间）：记
\(M_k(\operatorname{SL}_2(\mathbb{Z}))\) 为权 \(k\) 模形式构成的
\(\mathbb{C}\)-向量空间，\(S_k(\operatorname{SL}_2(\mathbb{Z}))\) 为权
\(k\) 尖点形式构成的子空间。

\textbf{定理 86.1}（模形式空间的维数）：

\[\dim M_k(\operatorname{SL}_2(\mathbb{Z})) = \begin{cases} \lfloor k/12 \rfloor & k \equiv 2 \pmod{12} \\ \lfloor k/12 \rfloor + 1 & k \not\equiv 2 \pmod{12} \end{cases}\]

对 \(k \geq 4\)
偶。\(M_0(\operatorname{SL}_2(\mathbb{Z})) = \mathbb{C}\)（常数函数），\(M_2(\operatorname{SL}_2(\mathbb{Z})) = 0\)（没有权
\(2\) 的非零模形式）。

\subsubsection{86.3 Eisenstein 级数}\label{eisenstein-ux7ea7ux6570}

\textbf{定义 86.7}（Eisenstein 级数）：对偶数 \(k \geq 4\)，权 \(k\) 的
\textbf{Eisenstein 级数} 为

\[G_k(z) = \sum_{(m,n) \in \mathbb{Z}^2 \setminus \{(0,0)\}} \frac{1}{(mz + n)^k}\]

归一化 Eisenstein 级数：\(E_k(z) = \frac{1}{2\zeta(k)} G_k(z)\)。

\textbf{命题 86.2}（Eisenstein 级数的 Fourier 展开）：

\[E_k(z) = 1 - \frac{2k}{B_k} \sum_{n=1}^\infty \sigma_{k-1}(n) q^n\]

其中 \(B_k\) 是 Bernoulli
数，\(\sigma_{k-1}(n) = \sum_{d \mid n} d^{k-1}\)。

\emph{证明}：使用 Poisson 求和公式（或 Lipschitz 求和公式）展开
\(G_k(z)\)。∎

\textbf{定理
86.3}（模形式环的结构）：\(M_*(\operatorname{SL}_2(\mathbb{Z})) = \bigoplus_{k \geq 0} M_k\)
是 \(\mathbb{C}\) 上的多项式环，由 \(E_4\) 和 \(E_6\) 生成：

\[M_*(\operatorname{SL}_2(\mathbb{Z})) \cong \mathbb{C}[E_4, E_6]\]

其中 \(E_4\) 和 \(E_6\) 代数无关。\(E_4\) 的权为 \(4\)，\(E_6\) 的权为
\(6\)。

\textbf{例 86.1}：判别式模形式
\(\Delta(z) = \frac{E_4(z)^3 - E_6(z)^2}{1728} = q \prod_{n=1}^\infty (1 - q^n)^{24}\)
是权 \(12\) 的尖点形式。\(\Delta(z)\) 的 Fourier 系数
\(\tau(n)\)（Ramanujan \(\tau\) 函数）满足著名的 Ramanujan 猜想（Deligne
证明，1974）：\(|\tau(p)| \leq 2p^{11/2}\)。

\subsubsection{86.4 Hecke 算子}\label{hecke-ux7b97ux5b50}

\textbf{定义 86.8}（Hecke 算子）：对正整数 \(n\)，权 \(k\) 模形式上的
\textbf{Hecke 算子} \(T_n\) 定义为

\[(T_n f)(z) = n^{k-1} \sum_{ad = n} \frac{1}{d^k} \sum_{b=0}^{d-1} f\left(\frac{az + b}{d}\right)\]

\textbf{命题 86.4}（Hecke 算子的性质）： 1. \(T_n: M_k \to M_k\) 和
\(T_n: S_k \to S_k\) 是线性算子 2. \(T_m T_n = T_n T_m\)（交换） 3. 若
\(\gcd(m, n) = 1\)，则 \(T_m T_n = T_{mn}\) 4. 对素数
\(p\)，\(T_{p^{r+1}} = T_p T_{p^r} - p^{k-1} T_{p^{r-1}}\)

\textbf{定义 86.9}（Hecke 特征形式）：若 \(f \in S_k\) 是所有 \(T_n\)
的共同特征向量，则称 \(f\) 为 \textbf{Hecke 特征形式}。若进一步
\(a_1 = 1\)，则 \(f\) 是\textbf{正规化 Hecke 特征形式}，此时
\(T_n f = a_n f\)，且 \(a_n\) 满足
\(a_m a_n = \sum_{d \mid \gcd(m,n)} d^{k-1} a_{mn/d^2}\)。

\textbf{定理 86.5}（Hecke 特征形式的 \(L\)-函数）：对正规化 Hecke
特征形式 \(f(z) = \sum_{n=1}^\infty a_n q^n\)，其 \(L\)-函数

\[L(s, f) = \sum_{n=1}^\infty \frac{a_n}{n^s}\]

有 Euler
乘积：\(L(s, f) = \prod_p \frac{1}{1 - a_p p^{-s} + p^{k-1-2s}}\)。且
\(L(s, f)\) 可解析延拓为整函数，满足函数方程。

\bigskip\hrule\bigskip

\subsection{第87章：椭圆曲线}\label{ux7b2c87ux7ae0ux692dux5706ux66f2ux7ebf}

椭圆曲线是亏格 \(1\) 的光滑射影曲线（带有一个指定基点），其上的点构成
Abel 群。椭圆曲线是连接数论、代数几何和模形式的中心对象。

\subsubsection{87.1
椭圆曲线的定义与群结构}\label{ux692dux5706ux66f2ux7ebfux7684ux5b9aux4e49ux4e0eux7fa4ux7ed3ux6784}

\textbf{定义 87.1}（椭圆曲线）：域
\(K\)（\(\operatorname{char} K \neq 2, 3\)）上的\textbf{椭圆曲线} \(E\)
是 Weierstrass 方程

\[E: y^2 = x^3 + ax + b\]

定义的光滑射影曲线（判别式
\(\Delta = -16(4a^3 + 27b^2) \neq 0\)），加上无穷远点
\(\mathcal{O}\)。椭圆曲线是亏格 \(1\) 的曲线。

\textbf{定义 87.2}（群结构）：\(E(K)\) 上的群运算定义为：对
\(P, Q \in E(K)\)，过 \(P\) 和 \(Q\) 的直线与 \(E\) 交于第三点
\(R\)，定义 \(P + Q\) 为 \(R\) 关于 \(x\) 轴的对称点。\(\mathcal{O}\)
是单位元。群运算是交换的。

\textbf{命题 87.1}（群运算的显式公式）：设
\(P = (x_1, y_1)\)，\(Q = (x_2, y_2)\)，\(P \neq \pm Q\)。则
\(P + Q = (x_3, y_3)\)，其中

\[x_3 = \lambda^2 - x_1 - x_2, \quad y_3 = \lambda(x_1 - x_3) - y_1, \quad \lambda = \frac{y_2 - y_1}{x_2 - x_1}\]

若 \(P = Q\)（倍点），则 \(\lambda = \frac{3x_1^2 + a}{2y_1}\)。

\emph{证明}：由直线与三次曲线的交点关系导出。使用 Vieta
公式计算第三个交点。∎

\subsubsection{87.2
有限域上的椭圆曲线}\label{ux6709ux9650ux57dfux4e0aux7684ux692dux5706ux66f2ux7ebf}

\textbf{定理 87.2}（Hasse 定理）：设 \(E\) 是有限域
\(\mathbb{F}_q\)（\(q = p^n\)）上的椭圆曲线。则

\[| \#E(\mathbb{F}_q) - (q + 1) | \leq 2\sqrt{q}\]

\textbf{定义 87.3}（\(L\)-函数）：\(E/\mathbb{F}_q\) 的 \textbf{zeta
函数} 为

\[Z(E/\mathbb{F}_q, T) = \exp\left(\sum_{n=1}^\infty \#E(\mathbb{F}_{q^n}) \frac{T^n}{n}\right) = \frac{1 - a_q T + q T^2}{(1 - T)(1 - qT)}\]

其中
\(a_q = q + 1 - \#E(\mathbb{F}_q)\)。\(1 - a_q T + q T^2 = (1 - \alpha T)(1 - \beta T)\)，\(|\alpha| = |\beta| = \sqrt{q}\)（Hasse
定理的等价形式）。

\textbf{定理 87.3}（椭圆曲线的挠群，陈述）：设 \(E\) 是 \(\mathbb{Q}\)
上的椭圆曲线。则 \(E(\mathbb{Q})_{\text{tors}}\)（有理挠群）同构于以下
15 个群之一：

\[\mathbb{Z}/n\mathbb{Z} \quad (1 \leq n \leq 10, n = 12), \quad \mathbb{Z}/2\mathbb{Z} \times \mathbb{Z}/2n\mathbb{Z} \quad (n = 1, 2, 3, 4)\]

（Mazur 定理，1977）。

\subsubsection{87.3 Mordell-Weil 定理}\label{mordell-weil-ux5b9aux7406}

\textbf{定理 87.4}（Mordell-Weil 定理）：设 \(E\) 是数域 \(K\)
上的椭圆曲线。则 \(E(K)\) 是有限生成 Abel 群：

\[E(K) \cong \mathbb{Z}^r \oplus E(K)_{\text{tors}}\]

其中 \(r = \operatorname{rank} E(K)\) 是 \textbf{Mordell-Weil
秩}（有限但计算困难），\(E(K)_{\text{tors}}\) 是有限挠子群。

\emph{证明思路}（Weil 的证明，1928）： 1. 弱 Mordell-Weil
定理：\(E(K)/2E(K)\) 是有限群。 2. 在 \(E(K)\)
上定义\textbf{高度函数}（height
function）\(h: E(K) \to \mathbb{R}_{\geq 0}\)，满足：\(h(2P) = 4h(P) + O(1)\)，且
\(\{P : h(P) \leq C\}\) 有限。 3. 由有限下降法（infinite
descent）：\(E(K)/2E(K)\) 有限 + 高度函数的性质 ⇒ \(E(K)\) 有限生成。∎

\textbf{定义 87.4}（高度函数）：对
\(P = (x, y) \in E(\mathbb{Q})\)，\(x = a/b\)（最简分数），定义\textbf{朴素高度}
\(H(P) = \max(|a|, |b|)\)，\textbf{对数高度}
\(h(P) = \log H(P)\)。\textbf{典范高度}（Néron-Tate
高度）\(\hat{h}(P) = \lim_{n \to \infty} \frac{h(2^n P)}{4^n}\)。

\textbf{命题 87.5}（典范高度的性质）： 1. \(\hat{h}(P) = 0\) 当且仅当
\(P\) 是挠点 2. \(\hat{h}(mP) = m^2 \hat{h}(P)\) 3.
\(\hat{h}(P + Q) + \hat{h}(P - Q) = 2\hat{h}(P) + 2\hat{h}(Q)\)（平行四边形律）

\subsubsection{87.4 Birch 和 Swinnerton-Dyer
猜想}\label{birch-ux548c-swinnerton-dyer-ux731cux60f3}

\textbf{定义 87.5}（\(E/\mathbb{Q}\) 的 Hasse-Weil \(L\)-函数）：设
\(E/\mathbb{Q}\) 是椭圆曲线。对素数 \(p\)，令
\(a_p = p + 1 - \#E(\mathbb{F}_p)\)。\(E\) 的 \(L\)-函数为

\[L(E, s) = \prod_{p \text{ 好约化}} \frac{1}{1 - a_p p^{-s} + p^{1-2s}} \cdot \prod_{p \text{ 坏约化}} \frac{1}{1 - a_p p^{-s}}\]

其中 \(a_p\) 在坏约化时取 \(0, 1, -1\)（取决于约化类型）。\(L(E, s)\) 在
\(\Re(s) > 3/2\) 内绝对收敛。

\textbf{定理 87.6}（模性定理，Wiles 等，1995-2001）：\(\mathbb{Q}\)
上的每条椭圆曲线都是模的：存在权 \(2\) 的 Hecke 特征形式 \(f\)（其
Fourier 系数 \(a_n\) 为整数），使得 \(L(E, s) = L(s, f)\)。

特别地，\(L(E, s)\) 可解析延拓为整个复平面上的整函数，满足函数方程。

\textbf{BSD 猜想}（Birch 和 Swinnerton-Dyer，1965）： 1.
\(\operatorname{ord}_{s=1} L(E, s) = \operatorname{rank} E(\mathbb{Q})\)（\(L\)-函数在
\(s=1\) 处的零点阶数等于秩） 2. \(L(E, s)\) 在 \(s=1\) 处的 Taylor
展开首项系数由 \(E\) 的算术不变量（Tate-Shafarevich
群、挠群、调节子等）精确给出

BSD 猜想是 Clay 数学研究所七大千禧年问题之一，目前仅对秩 \(0\) 和 \(1\)
的椭圆曲线部分证明。

\bigskip\hrule\bigskip

\subsection{第88章：类域论初步}\label{ux7b2c88ux7ae0ux7c7bux57dfux8bbaux521dux6b65}

类域论是代数数论的最高成就，描述了数域（或局部域）的 Abel
扩张与其理想类群（或乘性群）之间的深刻对应。

\subsubsection{88.1 局部域}\label{ux5c40ux90e8ux57df}

\textbf{定义
88.1}（局部域）：\textbf{局部域}是连通局部紧的拓扑域。包括： -
\textbf{Archimedean 局部域}：\(\mathbb{R}\) 和 \(\mathbb{C}\) -
\textbf{非 Archimedean
局部域}：\(\mathbb{Q}_p\)（\(p\)-进数域）的有限扩张

\textbf{定义 88.2}（\(p\)-进数域）：\(\mathbb{Q}_p\) 是 \(\mathbb{Q}\)
在 \(p\)-进度量 \(|x|_p = p^{-v_p(x)}\) 下的完备化，其中 \(v_p(x)\) 是
\(x\) 的
\(p\)-进赋值。\(\mathbb{Z}_p = \{x \in \mathbb{Q}_p : |x|_p \leq 1\}\)
是 \(p\)-进整数环。

\textbf{命题 88.1}：\(\mathbb{Q}_p\) 的每个有限扩张 \(K\)
带有一个离散赋值 \(v_K\)，剩余域为
\(\mathbb{F}_q\)（\(q = p^f\)）。\(\mathcal{O}_K = \{x \in K : v_K(x) \geq 0\}\)
是 \(K\) 的整数环，\(\mathfrak{p}_K = \{x : v_K(x) > 0\}\) 是极大理想。

\subsubsection{88.2 局部类域论}\label{ux5c40ux90e8ux7c7bux57dfux8bba}

\textbf{定义 88.3}（局部域上的 Weil 群）：设 \(K\)
是局部域，\(K^{\text{ab}}\) 是 \(K\) 的极大 Abel
扩张。局部类域论的核心是\textbf{局部互反映射}（local reciprocity map）

\[\rho_K: K^\times \to \operatorname{Gal}(K^{\text{ab}}/K)\]

\textbf{定理 88.1}（局部互反律）：存在自然的连续同态
\(\rho_K: K^\times \to \operatorname{Gal}(K^{\text{ab}}/K)\)，满足： 1.
\(\rho_K\) 的像是稠密的，核为 \(K^\times\) 中范数子群的交 2. 对有限 Abel
扩张 \(L/K\)，\(\rho_K\) 诱导同构
\(K^\times / N_{L/K}(L^\times) \cong \operatorname{Gal}(L/K)\) 3. 若
\(K = \mathbb{Q}_p\)，则 \(\rho_K\) 限制在 \(\mathbb{Z}_p^\times\) 上是
\(\operatorname{Gal}(\mathbb{Q}_p^{\text{unr}}/\mathbb{Q}_p)\) 上的
Frobenius 自同构

\textbf{推论 88.2}：局部域 \(K\) 的有限 Abel 扩张与 \(K^\times\)
的有限指数开子群（范数子群）一一对应。

\subsubsection{88.3 整体类域论}\label{ux6574ux4f53ux7c7bux57dfux8bba}

\textbf{定义 88.4}（Adèle 环与 Idèle 群）：数域 \(K\) 的 \textbf{Adèle
环} \(\mathbb{A}_K\) 是所有完备化 \(\prod_v' K_v\)
的限制直积。\textbf{Idèle 群} \(\mathbb{I}_K = \mathbb{A}_K^\times\) 是
Adèle 环的乘法群。

\textbf{定理 88.3}（整体互反律，Artin 互反律）：设 \(K\)
是数域，\(K^{\text{ab}}\) 是 \(K\) 的极大 Abel
扩张。存在自然的连续同态（整体互反映射）

\[\rho_K: \mathbb{I}_K \to \operatorname{Gal}(K^{\text{ab}}/K)\]

满足： 1. \(\rho_K\) 是满射，核为 \(K^\times\) 在 \(\mathbb{I}_K\)
中的闭包（连通分支的乘积） 2. 对有限 Abel 扩张 \(L/K\)，\(\rho_K\)
诱导同构
\(\mathbb{I}_K / K^\times N_{L/K}(\mathbb{I}_L) \cong \operatorname{Gal}(L/K)\)
3. \(\rho_K\) 与局部互反映射兼容：对每个位
\(v\)，\(\rho_K|_{K_v^\times} = \rho_{K_v}\)

\textbf{推论 88.4}（理想类群的类域论解释）：设 \(H_K\) 是 \(K\) 的
Hilbert 类域（\(K\) 的极大无处分歧 Abel 扩张）。则 Artin 映射诱导同构

\[\operatorname{Cl}(K) \cong \operatorname{Gal}(H_K/K)\]

其中 \(\operatorname{Cl}(K)\) 是 \(K\) 的理想类群。这解释了理想类群的
Galois 意义。

\textbf{定理 88.5}（Kronecker-Weber 定理）：\(\mathbb{Q}\) 的每个有限
Abel 扩张都包含在某个分圆域 \(\mathbb{Q}(\zeta_n)\)
中。等价地，\(\mathbb{Q}^{\text{ab}} = \bigcup_{n \geq 1} \mathbb{Q}(\zeta_n)\)。

\emph{证明}：这是类域论的特殊情形。由 \(\mathbb{Q}\)
的类数公式和局部-整体原理可得。∎

\textbf{定理 88.6}（虚二次域的复乘理论，陈述）：设 \(K\)
是虚二次域。\(K\) 的极大 Abel 扩张 \(K^{\text{ab}}\) 由 \(K\)
上的椭圆曲线的挠点和特殊值生成（复乘理论 / Kronecker
的青春之梦）。这是类域论的显式构造。

\bigskip\hrule\bigskip

\subsection{第89章：岩泽理论}\label{ux7b2c89ux7ae0ux5ca9ux6cfdux7406ux8bba}

岩泽理论（Iwasawa
Theory）由岩泽健吉于1950年代末创立，是研究数域无限扩域中理想类群和Galois模结构的深刻理论。设\(K\)为数域，取定素数\(p\)，考虑\(K\)的\(\mathbb{Z}_p\)-扩张\(K_\infty/K\)------即Galois群同构于加法群\(\mathbb{Z}_p\)的无限扩张。记\(K_n\)为中间域（\([K_n:K]=p^n\)），\(A_n\)为\(K_n\)的理想类群的\(p\)-Sylow子群。岩泽发现逆向极限\(X = \varprojlim A_n\)上的Galois作用赋予\(X\)一个\(\Lambda = \mathbb{Z}_p[[T]]\)（岩泽代数）的模结构，且\(X\)是有限生成挠\(\Lambda\)-模。由结构定理，\(X\)对应一个特征多项式（岩泽不变量\(\mu, \lambda, \nu\)）。\textbf{岩泽主猜想}（Iwasawa
Main Conjecture）断言：\(X\)的特征理想由\(p\)-adic
\(L\)-函数生成------即\(p\)-adic解析对象与代数对象之间的精确等式。\textbf{Mazur-Wiles证明}（1984年）：利用模曲线和Galois表示的几何方法，对\(\mathbb{Q}\)的Abel扩张证明了主猜想，获1995年Cole奖。\textbf{Skinner-Urban的进展}（2010年代）：将主猜想推广至\(GL(2)\)自守形式情形，证明了更多一般代数自守表示对应的岩泽主猜想，是Langlands纲领与岩泽理论交叉的前沿成果。

\bigskip\hrule\bigskip

\subsection{第90章：自守形式}\label{ux7b2c90ux7ae0ux81eaux5b88ux5f62ux5f0f}

自守形式（Automorphic
Forms）是模形式向一般约化代数群\(G\)的推广，是Langlands纲领的核心研究对象。对\(G = GL(2)\)，自守形式回到上半平面上的经典模形式；对\(G = GL(n)\)（\(n \geq 3\)），自守形式是定义在商空间\(GL(n, \mathbb{Z}) \backslash GL(n, \mathbb{R}) / O(n)\)上的光滑函数，满足缓增性、\(K\)-有限性和\(Z(\mathfrak{g})\)-有限性条件。每个尖点自守表示\(\pi\)对应一个\textbf{自守\(L\)-函数}\(L(s, \pi) = \prod_p L_p(s, \pi_p)\)，其中\(L_p\)为局部Langlands对应的Euler因子（如\(GL(n)\)情形下\(L_p(s, \pi_p) = \prod_{i=1}^n (1 - \alpha_i(p) p^{-s})^{-1}\)）。\textbf{Langlands函子性猜想}断言：若存在\(L\)-群之间的同态\({}^L H \to {}^L G\)，则\(H\)的自守表示可通过函子性提升为\(G\)的自守表示。\textbf{谷山-志村-韦伊猜想}（每条\(\mathbb{Q}\)上的椭圆曲线是模的，即对应于权2的\(GL(2)\)自守表示）正是Langlands函子性在\(GL(2)\)的特例。Wiles于1995年证明了该猜想的半稳定情形，从而奠定了费马大定理的证明；Breuil-Conrad-Diamond-Taylor于2001年完全证明了该猜想，这是Langlands纲领迄今为止最辉煌的成果。

\bigskip\hrule\bigskip

\subsection{第91章：Shimura 簇}\label{ux7b2c91ux7ae0shimura-ux7c07}

Shimura 簇是 Shimura Goro（志村五郎）于 1960 年代引入、由 Deligne 于
1970 年代系统公理化的一类特殊代数簇，它们是模曲线的高维推广。Shimura
簇携带丰富的算术结构，是现代自守形式理论、Langlands
纲领和算术几何的核心几何对象------通过其上同调，自守表示与 Galois
表示（motivic \(L\)-函数）之间建立了精确的对应。

\subsubsection{91.1 Shimura 数据与 Shimura
簇的构造}\label{shimura-ux6570ux636eux4e0e-shimura-ux7c07ux7684ux6784ux9020}

\textbf{定义 91.1}（Shimura 数据）：一个 \textbf{Shimura 数据}（Shimura
datum）是一个对 \((G, X)\)，其中：

\begin{itemize}
\tightlist
\item
  \(G\) 是 \(\mathbb{Q}\) 上的连通简约代数群；
\item
  \(X\) 是 \(G(\mathbb{R})\) 在其上可逆作用的齐次空间（Hermite
  对称域），满足以下 Deligne 公理（SV1-SV6）：
\end{itemize}

\textbf{(SV1)} \(G(\mathbb{R})\) 在 \(X\) 上的作用是传递的，且对某点
\(x \in X\)，稳定子群 \(K_x = \operatorname{Stab}_{G(\mathbb{R})}(x)\)
是 \(G(\mathbb{R})\) 的极大紧子群（精确到有限指数）；

\textbf{(SV2)} \(X\) 具有 \(G(\mathbb{R})\)-不变的复结构，且作为 Hermite
对称空间，\(X\) 上的每个点给出一个 Hodge 结构（通过 \(G \to GL(V)\)
的某个表示）；

\textbf{(SV3)} Adjoint 群 \(G^{\text{ad}}\) 没有定义在 \(\mathbb{Q}\)
上的紧单因子（即 \(G^{\text{ad}}\) 无 \(\mathbb{Q}\)-单的紧因子）；

\textbf{(SV4)} 权同态 \(w: \mathbb{G}_m \to G\)（定义在 \(\mathbb{Q}\)
上）满足某些兼容条件；

\textbf{(SV5)} 对 \(G^{\text{ad}}\) 的每个 \(\mathbb{Q}\)-单因子，其
\(\mathbb{R}\)-点是非紧的；

\textbf{(SV6)} 对 \(G(\mathbb{R})\) 中由
\(\operatorname{Lie} G(\mathbb{R})\) 中满足 \(J^2 = -\operatorname{id}\)
的复结构 \(J\) 的自然作用，\(X\) 可等同于 \(G(\mathbb{R})\) 的共轭类。

\textbf{定义 91.2}（Shimura 簇）：设 \((G, X)\) 是 Shimura
数据，\(K \subset G(\mathbb{A}_f)\) 是有限 Adèle 点的紧开子群（称为
\textbf{级结构}，level structure）。则 \textbf{Shimura 簇}
定义为局部对称空间（双陪集空间）

\[\operatorname{Sh}_K(G, X) = G(\mathbb{Q}) \backslash X \times G(\mathbb{A}_f) / K\]

其中 \(G(\mathbb{Q})\) 通过对角嵌入作用于
\(X \times G(\mathbb{A}_f)\)，\(K\) 作用于右因子 \(G(\mathbb{A}_f)\)。

当 \(K\) 充分小时（neat），\(G(\mathbb{Q})\) 无固定点的元素作用于
\(X \times G(\mathbb{A}_f) / K\)，此时 \(\operatorname{Sh}_K(G, X)\)
是光滑复流形。

\subsubsection{91.2 经典例子}\label{ux7ecfux5178ux4f8bux5b50}

\textbf{例 91.1}（模曲线）：取 \(G = GL(2)_{\mathbb{Q}}\)（更精确地为
\(G = GSp(2)\) 或 \(G = GL(2)\) 在 \(\mathbb{G}_m\)
中的适当限制），\(X = \mathbb{H}\)（上半平面），\(K = K(N) = \{\gamma \in GL(2, \widehat{\mathbb{Z}}) : \gamma \equiv I \pmod{N}\}\)。则

\[\operatorname{Sh}_K(GL(2), \mathbb{H}^\pm) = GL(2, \mathbb{Q}) \backslash \mathbb{H}^\pm \times GL(2, \mathbb{A}_f) / K(N) \cong Y(N)(\mathbb{C})\]

即完整的模曲线 \(Y(N)\) 的复点。这是 Shimura 簇的原型------所有 Shimura
簇都是模曲线的某种自然推广。

\textbf{例 91.2}（Siegel 模簇）：取
\(G = GSp(2g)\)（辛相似群），\(X = \mathbb{H}_g\)（Siegel
上半空间------所有 \(g \times g\) 对称复矩阵 \(Z\) 满足
\(\operatorname{Im}(Z) > 0\) 的集合）。则
\(\operatorname{Sh}_K(GSp(2g), \mathbb{H}_g)\) 是亏格 \(g\) 的主极化
Abelian 簇的模空间 \(\mathcal{A}_g\)（带级结构）。

Siegel 上半空间 \(\mathbb{H}_g\) 的维数为 \(g(g+1)/2\)，其上
\(Sp(2g, \mathbb{R})\) 以分式线性变换作用：

\[\begin{pmatrix} A & B \\ C & D \end{pmatrix} \cdot Z = (AZ + B)(CZ + D)^{-1}\]

\textbf{例 91.3}（志村曲线）：设 \(B\) 是 \(\mathbb{Q}\)
上的四元数代数（definite 或 indefinite）。若 \(B\)
在无穷远点分裂（indefinite），取 \(G = B^\times\)（\(B\)
的乘法群），\(X = \mathbb{H}\)，则
\(\operatorname{Sh}_K(B^\times, \mathbb{H})\)
是\textbf{志村曲线}（Shimura curve）------紧 Riemann
面。与经典模曲线不同，志村曲线没有尖点（cusp），因为 \(G\) 没有
\(\mathbb{Q}\)-抛物子群。

取 \(B\) 在素数 \(p\) 处分歧的四元数代数，对应的志村曲线在
\(p\)-进几何中有重要应用（如 \(p\)-进一致化理论和 Cerednik-Drinfeld
定理）。

\subsubsection{91.3 Baily-Borel
定理与典范模型}\label{baily-borel-ux5b9aux7406ux4e0eux5178ux8303ux6a21ux578b}

\textbf{定理 91.1}（Baily-Borel 定理，1966 年）：对任意 Shimura 数据
\((G, X)\) 和紧开子群 \(K \subset G(\mathbb{A}_f)\)，Shimura 簇
\(\operatorname{Sh}_K(G, X)\) 具有 \textbf{拟射影代数簇}
的自然结构------即 \(\operatorname{Sh}_K(G, X)\)
可映为一个拟射影复代数簇的解析化。

\emph{证明思路}：Baily 和 Borel 的证明利用了以下事实：\(X\)
是对称域，其上存在 \(G(\mathbb{R})\) 不变的 Bergman 度量。将
\(G(\mathbb{Q})\) 不变的尖点形式（cusp forms）的 Fourier-Jacobi
展开与复结构的嵌入复合，证明了 \(\operatorname{Sh}_K(G, X)\) 满足
Kodaira 嵌入定理的条件。具体地，高次张量的自守形式（权充分大且关于 \(K\)
不变的）提供足够多的截面以定义到射影空间的嵌入。该嵌入的像是 Zariski
闭的（由自守形式的代数关系），从而赋予拟射影结构。∎

Baily-Borel 定理也给出了 Shimura 簇的紧化------\textbf{Baily-Borel
紧化}（或极小紧化）------通过在边界添加较低维的 Shimura 簇的商。

\textbf{定理 91.2}（典范模型，Canonical Model）：Shimura 簇
\(\operatorname{Sh}_K(G, X)\) 在\textbf{反射域}（reflex
field）\(E(G, X)\) 上具有有理模型。反射域是 \(G\)
的共轭类的定义域，也是关联的零维 Shimura 簇的分裂域。

粗略地说，\(\operatorname{Sh}_K(G, X)\) 作为复代数簇可定义在数域
\(E(G, X)\) 上------即存在 \(E(G, X)\) 上的概型 \(\mathcal{S}_K\) 使得
\(\mathcal{S}_K \times_{E(G, X)} \mathbb{C} \cong \operatorname{Sh}_K(G, X)\)。该模型的构造由
Shimura 的互反律理论给出，利用了 CM 点（complex multiplication
points）的 Galois 作用：CM 点对应于特殊自守表示，其上 Galois
群的作用由类域论（整体互反律）描述。

\subsubsection{91.4 上同调与 Langlands
纲领}\label{ux4e0aux540cux8c03ux4e0e-langlands-ux7eb2ux9886}

Shimura 簇的核心重要性在于：其 \textbf{Betti 上同调}（或有理
\(\ell\)-进上同调）实现了 Langlands 纲领中的 Galois 表示。

\textbf{定理 91.3}（Shimura 簇的上同调与自守表示）：设 \((G, X)\) 是
Shimura 数据，\(\xi\) 是 \(G\) 的代数表示。则总上同调群

\[H^*(\operatorname{Sh}_K(G, X), \mathcal{L}_\xi)\]

（\(\mathcal{L}_\xi\) 是由 \(\xi\) 确定的局部系统）分解为
\(G(\mathbb{A}_f)\)-模的直和，每个直和项同构于某个自守表示 \(\pi\) 的
\(\pi_f\) 部分的若干次副本。即

\[H^*(\operatorname{Sh}_K(G, X), \mathcal{L}_\xi) \cong \bigoplus_{\pi} m(\pi) \cdot \pi_f^K \otimes H^*(\mathfrak{g}, K_\infty, \pi_\infty \otimes \xi)\]

其中 \(m(\pi)\) 是 \(\pi\) 的重数（离散系列表示的重数有限）。

\textbf{推论 91.4}（Galois 表示与 Langlands
对应）：\(\ell\)-进平展上同调
\(H_{\text{ét}}^i(\operatorname{Sh}_K(G, X)_{\overline{\mathbb{Q}}}, \mathbb{Q}_\ell)\)
作为
\(G(\mathbb{A}_f) \times \operatorname{Gal}(\overline{\mathbb{Q}} / E)\)
的表示实现了 Langlands 对应的几何版本：每个尖点自守表示 \(\pi\) 关联一个
Galois 表示
\(\rho_{\pi, \ell}: \operatorname{Gal}(\overline{\mathbb{Q}} / E) \to {}^L G(\mathbb{Q}_\ell)\)（到
\(L\)-群的同态）。

\subsubsection{91.5 Langlands-Kottwitz 方法与 Hasse-Weil Zeta
函数}\label{langlands-kottwitz-ux65b9ux6cd5ux4e0e-hasse-weil-zeta-ux51fdux6570}

\textbf{定理 91.5}（Kottwitz 的 Langlands-Kottwitz 方法，1990 年代）：设
\(\operatorname{Sh}_K(G, X)\) 是 Shimura 簇（带有光滑整模型）。其
Hasse-Weil zeta 函数（算术 zeta 函数）

\[\zeta(\operatorname{Sh}_K, s) = \prod_{x \in |\operatorname{Sh}_K|} \frac{1}{1 - N(x)^{-s}}\]

（\(x\) 遍历 \(\operatorname{Sh}_K\) 的闭点）可通过 Langlands 对应的
\(L\)-函数表达。具体地，Kottwitz 证明了

\[\zeta(\operatorname{Sh}_K, s) = \prod_{\pi} L(s - \frac{\dim \operatorname{Sh}_K}{2}, \pi, r)^{m(\pi)}\]

其中乘积遍历 \(G\) 的自守表示 \(\pi\)，\(r\) 是 \(L\)-群 \({}^L G\)
的某个自然表示（由 Shimura 簇的切丛的
\(G\)-模结构决定），\(L(s, \pi, r)\) 是关联的 \(L\)-函数。

\emph{证明思路}：Kottwitz 的证明基于三个关键步骤： (1)
\textbf{Grothendieck-Lefschetz 迹公式}用于计算 zeta 函数的 Euler
因子：\(\#\operatorname{Sh}_K(\mathbb{F}_q) = \sum_i (-1)^i \operatorname{Tr}(\operatorname{Frob}_q \mid H_{\text{ét}}^i(\operatorname{Sh}_K, \mathbb{Q}_\ell))\)。
(2) \textbf{Langlands-Kottwitz 的稳定迹公式}将几何迹（Shimura
簇在正特征的点数）分解为 Galois 上同调和自守表示的贡献。 (3)
\textbf{基本引理}（Fundamental Lemma, Ngo Bao Chau 2010
年证明）确保局部轨道积分在整体迹公式中可简化------这是 Langlands
纲领最重大的突破之一，Ngo 因此获得 2010 年 Fields 奖。∎

\subsubsection{91.6 Shimura
簇的算术应用}\label{shimura-ux7c07ux7684ux7b97ux672fux5e94ux7528}

\begin{enumerate}
\def\labelenumi{\arabic{enumi}.}
\item
  \textbf{志村簇上的 CM 点的 Galois 作用}：CM 点生成 Shimura 簇中 0
  维的子簇，其 Galois 作用由类域论完全描述。这是 Hilbert 第 12
  问题（Kronecker
  青春之梦）的高维推广------虚二次域的显式类域论（复乘理论）推广到更一般的
  Shimura 数据。
\item
  \textbf{志村簇的模型与模性提升定理}：\(GL(2)\)-型 Shimura
  簇（模曲线）的整模型理论（Deligne-Rapoport、Katz-Mazur）直接应用于
  Wiles 对 Fermat 大定理的证明。高维推广的整模型理论（Kisin、Scholze
  等）是 \(p\)-进自守形式理论的核心。
\item
  \textbf{志村簇和 Langlands 纲领}：Scholze（2013 年）利用 perfectoid
  空间理论证明了拟射影 Shimura 簇的 Hodge-Tate
  周期映射的某些性质，并构造了志村簇的 \(p\)-进一致化（类似于
  Cerednik-Drinfeld 定理的高维推广）。
\item
  \textbf{André-Oort 猜想}：关于 Shimura 簇中特殊点（CM 点）的 Zariski
  闭包------特殊子簇正好是较低维的 Shimura 子簇。该猜想由
  Pila-Zannier（2009 年）使用 o-极小性方法在积模曲线情形证明，并由
  Klingler-Ullmo-Yafaev 和 Pila-Tsimerman 在一般情形取得重大进展。
\end{enumerate}

\bigskip\hrule\bigskip

\emph{本卷完成了从初等数论到现代数论的跨越：解析数论以 \(\zeta(s)\)
和素数定理为核心，展示了分析工具在数论中的威力；模形式理论揭示了全纯函数与数论之间惊人的联系（Ramanujan
\(\tau\) 函数、Hecke 算子）；椭圆曲线以 Mordell-Weil 定理和 BSD
猜想为中心，连接了代数几何、模形式和
\(L\)-函数；类域论（局部和整体）以互反律为核心，完成了 Abel
扩张的分类。为后续 Langlands 纲领（如有需要）的学习奠定了基础。}

\newpage

\section{卷十七：微分方程}\label{ux5377ux5341ux4e03ux5faeux5206ux65b9ux7a0b}

\begin{quote}
微分方程是数学物理、工程和自然科学的核心数学工具。本卷从常微分方程（ODE）的存在唯一性理论出发，建立线性
ODE
系统的稳定性理论，然后转向偏微分方程（PDE）的分类和基本解法，最后引入
Sobolev 空间和弱解理论，为现代 PDE 研究奠定分析基础。
\end{quote}

\bigskip\hrule\bigskip

\subsection{第89章：常微分方程基础}\label{ux7b2c89ux7ae0ux5e38ux5faeux5206ux65b9ux7a0bux57faux7840}

本章从 ODE
的基本理论出发，重点讨论初值问题解的存在性、唯一性以及对参数和初始条件的连续依赖性。

\subsubsection{89.1 基本概念}\label{ux57faux672cux6982ux5ff5}

\textbf{定义 89.1}（ODE 系统）：一阶 ODE
系统的\textbf{初值问题}（IVP）为

\[\begin{cases} \mathbf{x}'(t) = \mathbf{f}(t, \mathbf{x}(t)) \\ \mathbf{x}(t_0) = \mathbf{x}_0 \end{cases}\]

其中
\(\mathbf{f}: U \subseteq \mathbb{R} \times \mathbb{R}^n \to \mathbb{R}^n\)
是给定的向量场，\(\mathbf{x}_0 \in \mathbb{R}^n\) 是初始条件。

\textbf{定义 89.2}（解）：区间 \(I \subseteq \mathbb{R}\) 上的可微函数
\(\mathbf{x}: I \to \mathbb{R}^n\) 称为 IVP 的\textbf{解}，如果
\(\mathbf{x}(t_0) = \mathbf{x}_0\) 且对所有
\(t \in I\)，\((t, \mathbf{x}(t)) \in U\) 且
\(\mathbf{x}'(t) = \mathbf{f}(t, \mathbf{x}(t))\)。

\textbf{定义 89.3}（\(n\) 阶 ODE）：\(n\) 阶 ODE
\(y^{(n)} = f(t, y, y', \ldots, y^{(n-1)})\) 可通过引入
\(x_1 = y\)，\(x_2 = y'\)，\(\ldots\)，\(x_n = y^{(n-1)}\)
化为一阶系统。

\subsubsection{89.2 Picard-Lindelöf
定理}\label{picard-lindeluxf6f-ux5b9aux7406}

\textbf{定义 89.4}（Lipschitz 条件）：\(\mathbf{f}: U \to \mathbb{R}^n\)
称为关于 \(\mathbf{x}\) \textbf{局部 Lipschitz} 的，如果对每个紧子集
\(K \subseteq U\)，存在常数 \(L_K > 0\) 使得

\[\|\mathbf{f}(t, \mathbf{x}_1) - \mathbf{f}(t, \mathbf{x}_2)\| \leq L_K \|\mathbf{x}_1 - \mathbf{x}_2\|\]

对所有 \((t, \mathbf{x}_1), (t, \mathbf{x}_2) \in K\) 成立。

\textbf{定理 89.1}（Picard-Lindelöf 定理，局部存在唯一性）：设
\(\mathbf{f}\) 在包含 \((t_0, \mathbf{x}_0)\) 的开集上连续，且关于
\(\mathbf{x}\) 局部 Lipschitz。则存在 \(\delta > 0\)，使得 IVP 在
\([t_0 - \delta, t_0 + \delta]\) 上有唯一解。

\emph{证明}：将 IVP 转化为积分方程
\(\mathbf{x}(t) = \mathbf{x}_0 + \int_{t_0}^t \mathbf{f}(s, \mathbf{x}(s)) ds\)。定义
Picard
迭代：\(\mathbf{x}_0(t) = \mathbf{x}_0\)，\(\mathbf{x}_{n+1}(t) = \mathbf{x}_0 + \int_{t_0}^t \mathbf{f}(s, \mathbf{x}_n(s)) ds\)。在充分小的区间上，\(C([t_0 - \delta, t_0 + \delta], \mathbb{R}^n)\)
中赋予上确界范数，Picard 迭代是压缩映射。由 Banach
不动点定理（压缩映射原理），存在唯一不动点。∎

\textbf{定理 89.2}（Peano 存在定理）：若 \(\mathbf{f}\) 仅连续（不要求
Lipschitz），则 IVP 至少有一个解（但不一定唯一）。例如
\(\dot{x} = 2\sqrt{|x|}\)，\(x(0) = 0\) 有无穷多解。

\textbf{定理 89.3}（连续依赖性）：设 \(\mathbf{f}\) 连续且关于
\(\mathbf{x}\) 局部 Lipschitz。则解 \(\mathbf{x}(t; t_0, \mathbf{x}_0)\)
在其存在区间内连续依赖于初始条件 \((t_0, \mathbf{x}_0)\)。

\emph{证明}：使用 Gronwall 不等式。设 \(\mathbf{x}_1, \mathbf{x}_2\)
是两个不同初值的解，则
\(\|\mathbf{x}_1(t) - \mathbf{x}_2(t)\| \leq \|\mathbf{x}_1(0) - \mathbf{x}_2(0)\| e^{L|t|}\)。∎

\subsubsection{89.3
解的延拓与极大解}\label{ux89e3ux7684ux5ef6ux62d3ux4e0eux6781ux5927ux89e3}

\textbf{定义 89.5}（极大解）：IVP 的\textbf{极大解}是定义在极大存在区间
\(I_{\max}\) 上的解，不能延拓到更大的区间上。

\textbf{命题 89.4}（极大解的存在性）：若 \(\mathbf{f}\) 连续且局部
Lipschitz，则每个 IVP 有唯一的极大解。

\textbf{定理 89.5}（解的爆破判别）：若极大解的存在区间
\(I_{\max} = (T_-, T_+)\) 有限（即 \(T_+ < \infty\)），则当
\(t \to T_+^-\) 时，\(\|\mathbf{x}(t)\| \to \infty\) 或
\((t, \mathbf{x}(t))\) 趋于 \(U\) 的边界。

\emph{证明}：若 \(\|\mathbf{x}(t)\|\) 有界且 \((t, \mathbf{x}(t))\)
不趋于边界，则 \(\mathbf{x}(t)\)
有收敛子列，由局部存在性可延拓解，与极大性矛盾。∎

\subsubsection{89.4
自治系统与相平面分析}\label{ux81eaux6cbbux7cfbux7edfux4e0eux76f8ux5e73ux9762ux5206ux6790}

\textbf{定义 89.6}（自治系统）：ODE 系统
\(\mathbf{x}' = \mathbf{f}(\mathbf{x})\)（\(\mathbf{f}\) 不显含
\(t\)）称为\textbf{自治系统}。

\textbf{定义 89.7}（相图）：自治系统的\textbf{相图}是 \(\mathbb{R}^n\)
中全体轨道的集合。轨道 \(\{\mathbf{x}(t) : t \in I_{\max}\}\)
是相空间中的曲线。

\textbf{定义 89.8}（平衡点）：\(\mathbf{x}_*\)
称为\textbf{平衡点}（或不动点），如果
\(\mathbf{f}(\mathbf{x}_*) = 0\)。平衡点对应于常数解。

\textbf{定义 89.9}（二维线性系统的分类）：对于线性系统
\(\dot{\mathbf{x}} = A\mathbf{x}\)（\(A \in \mathbb{R}^{2 \times 2}\)），平衡点
\(\mathbf{0}\) 的类型由 \(A\) 的特征值 \(\lambda_1, \lambda_2\) 决定：

{\def\LTcaptype{none} % do not increment counter
\begin{longtable}[]{@{}ll@{}}
\toprule\noalign{}
特征值类型 & 平衡点类型 \\
\midrule\noalign{}
\endhead
\bottomrule\noalign{}
\endlastfoot
\(\lambda_1 < \lambda_2 < 0\) & 稳定结点 \\
\(0 < \lambda_1 < \lambda_2\) & 不稳定结点 \\
\(\lambda_1 < 0 < \lambda_2\) & 鞍点 \\
\(\lambda = \alpha \pm i\beta\)，\(\alpha < 0\) & 稳定焦点 \\
\(\lambda = \alpha \pm i\beta\)，\(\alpha > 0\) & 不稳定焦点 \\
\(\lambda = \pm i\beta\)（纯虚数） & 中心 \\
\end{longtable}
}

\textbf{命题 89.5}（相平面分析的关键步骤）：对非线性自治系统
\(\dot{\mathbf{x}} = \mathbf{f}(\mathbf{x})\)： 1. \textbf{求平衡点}：解
\(\mathbf{f}(\mathbf{x}) = \mathbf{0}\)。 2.
\textbf{线性化}：在每个平衡点处计算 Jacobi 矩阵
\(D\mathbf{f}\)，求特征值并分类（结点、鞍点、焦点、中心）。 3.
\textbf{双曲平衡点的局部行为}：若 \(\Re(\lambda) \neq 0\)，由
Hartman-Grobman 定理，局部相图拓扑等价于线性化系统。 4.
\textbf{零实部特征值的中心流形}：若存在纯虚数特征值，线性化不能确定稳定性，需在中心流形上约化分析非线性项。
5. \textbf{全局相图}：绘制零倾线（nullclines，即 \(\dot{x}_1 = 0\) 和
\(\dot{x}_2 = 0\) 的曲线），确定各区域中的向量场方向，用
Poincaré-Bendixson 定理判断周期轨的存在性。

\bigskip\hrule\bigskip

\subsection{第90章：线性ODE系统}\label{ux7b2c90ux7ae0ux7ebfux6027odeux7cfbux7edf}

线性 ODE
系统可借助矩阵指数和线性代数（V7）的工具完全求解。本章重点讨论常系数系统和稳定性理论。

\subsubsection{90.1 矩阵指数}\label{ux77e9ux9635ux6307ux6570}

\textbf{定义 90.1}（矩阵指数）：对 \(A \in M_n(\mathbb{C})\)，定义

\[e^{tA} = \sum_{k=0}^\infty \frac{t^k A^k}{k!}\]

此级数对任意 \(t \in \mathbb{R}\) 绝对收敛（在任意矩阵范数下）。

\textbf{命题 90.1}（矩阵指数的性质）： 1. \(e^{0A} = I\) 2.
\(\frac{d}{dt} e^{tA} = A e^{tA} = e^{tA} A\) 3. 若 \(AB = BA\)，则
\(e^{A+B} = e^A e^B\) 4. \(e^{tA}\) 可逆，\((e^{tA})^{-1} = e^{-tA}\) 5.
若 \(A = PJP^{-1}\)（\(J\) 为 Jordan 标准形），则
\(e^{tA} = P e^{tJ} P^{-1}\)

\textbf{定理 90.2}（线性系统的解）：IVP
\(\dot{\mathbf{x}} = A\mathbf{x}\)，\(\mathbf{x}(0) = \mathbf{x}_0\)
的唯一解为

\[\mathbf{x}(t) = e^{tA} \mathbf{x}_0\]

非齐次系统 \(\dot{\mathbf{x}} = A\mathbf{x} + \mathbf{b}(t)\)
的解由常数变易公式给出：

\[\mathbf{x}(t) = e^{tA} \mathbf{x}_0 + \int_0^t e^{(t-s)A} \mathbf{b}(s) ds\]

\emph{证明}：\(\frac{d}{dt} e^{tA} \mathbf{x}_0 = A e^{tA} \mathbf{x}_0 = A\mathbf{x}(t)\)，且
\(\mathbf{x}(0) = \mathbf{x}_0\)。唯一性由 Picard-Lindelöf 定理保证。∎

\subsubsection{90.2
常系数线性系统的稳定性}\label{ux5e38ux7cfbux6570ux7ebfux6027ux7cfbux7edfux7684ux7a33ux5b9aux6027}

\textbf{定义 90.2}（Lyapunov 稳定性）：设 \(\mathbf{x}_*\) 是
\(\dot{\mathbf{x}} = \mathbf{f}(\mathbf{x})\) 的平衡点。 -
\(\mathbf{x}_*\) 是\textbf{稳定}的（Lyapunov 意义下），如果对任意
\(\varepsilon > 0\)，存在 \(\delta > 0\) 使得
\(\|\mathbf{x}(0) - \mathbf{x}_*\| < \delta\) 推出
\(\|\mathbf{x}(t) - \mathbf{x}_*\| < \varepsilon\)（对所有
\(t \geq 0\)） - \(\mathbf{x}_*\) 是\textbf{渐近稳定}的，如果它稳定且
\(\lim_{t \to \infty} \mathbf{x}(t) = \mathbf{x}_*\) - \(\mathbf{x}_*\)
是\textbf{指数稳定}的，如果存在 \(C, \alpha > 0\) 使得
\(\|\mathbf{x}(t) - \mathbf{x}_*\| \leq C e^{-\alpha t} \|\mathbf{x}(0) - \mathbf{x}_*\|\)

\textbf{定理 90.3}（线性稳定性）：对线性系统
\(\dot{\mathbf{x}} = A\mathbf{x}\)： - \(\mathbf{0}\)
渐近稳定（实际上指数稳定）当且仅当 \(A\) 的所有特征值满足
\(\Re(\lambda) < 0\) - \(\mathbf{0}\) 稳定当且仅当 \(A\)
的所有特征值满足 \(\Re(\lambda) \leq 0\)，且满足 \(\Re(\lambda) = 0\)
的特征值对应大小为 \(1 \times 1\) 的 Jordan 块 - \(\mathbf{0}\)
不稳定当且仅当存在特征值 \(\Re(\lambda) > 0\)（或 \(\Re(\lambda) = 0\)
但 Jordan 块大小 \textgreater{} 1）

\textbf{定理 90.4}（线性化稳定性，Hartman-Grobman 定理）：设
\(\mathbf{x}_*\) 是 \(\dot{\mathbf{x}} = \mathbf{f}(\mathbf{x})\)
的平衡点，\(\mathbf{f}\) 是 \(C^1\) 的。若
\(A = D\mathbf{f}(\mathbf{x}_*)\) 的所有特征值满足
\(\Re(\lambda) \neq 0\)（双曲平衡点），则非线性系统在 \(\mathbf{x}_*\)
附近拓扑共轭于其线性化系统 \(\dot{\mathbf{y}} = A\mathbf{y}\)。

\emph{证明思路}：不妨设 \(\mathbf{x}_* = 0\)，将系统写为
\(\dot{\mathbf{x}} = A\mathbf{x} + \mathbf{g}(\mathbf{x})\)，其中
\(\mathbf{g}(0)=0\)，\(D\mathbf{g}(0)=0\)。核心是构造同胚
\(H = \operatorname{id} + \mathbf{h}\) 满足
\(H \circ e^{tA} = \varphi_t \circ H\)（\(\varphi_t\) 是非线性流）。将
\(\mathbf{g}\)
分解为稳定部分和不稳定部分的贡献。利用双曲性（\(\Re(\lambda) \neq 0\)），线性流
\(e^{tA}\)
具有指数二分性：稳定子空间上的压缩和不稳定子空间上的扩张。对给定的
\(\mathbf{g}\)（在 \(C^1\)
小邻域中截断），用压缩映射原理在适当的函数空间中求解不变叶层方程
\(\mathbf{h}(A\mathbf{x} + \mathbf{g}(\mathbf{x})) = A\mathbf{h}(\mathbf{x}) + \mathbf{g}(\mathbf{x} + \mathbf{h}(\mathbf{x}))\)。截断的
\(\mathbf{g}\)
在原点附近与原系统一致，而远离原点处的修正不影响局部拓扑共轭性。∎

\subsubsection{90.3 Lyapunov 函数}\label{lyapunov-ux51fdux6570}

\textbf{定义 90.3}（Lyapunov 函数）：设 \(\mathbf{x}_*\)
是平衡点。\(C^1\) 函数
\(V: \mathcal{U} \to \mathbb{R}\)（\(\mathcal{U}\) 是 \(\mathbf{x}_*\)
的邻域）称为 \textbf{Lyapunov 函数}，如果
\(V(\mathbf{x}_*) = 0\)，\(V(\mathbf{x}) > 0\)（\(\mathbf{x} \neq \mathbf{x}_*\)），且沿轨道
\(\dot{V}(\mathbf{x}) = \nabla V(\mathbf{x}) \cdot \mathbf{f}(\mathbf{x}) \leq 0\)。

\textbf{定理 90.5}（Lyapunov 稳定性定理）： 1. 若存在 Lyapunov 函数，则
\(\mathbf{x}_*\) 是稳定的 2. 若进一步
\(\dot{V}(\mathbf{x}) < 0\)（\(\mathbf{x} \neq \mathbf{x}_*\)），则
\(\mathbf{x}_*\) 是渐近稳定的

\emph{证明}：由 \(V\) 沿轨道递减，\(\|\mathbf{x}(t) - \mathbf{x}_*\|\)
有界。若 \(\dot{V} < 0\)，则轨道趋于 \(\mathbf{x}_*\)。∎

\bigskip\hrule\bigskip

\subsection{第91章：偏微分方程基础}\label{ux7b2c91ux7ae0ux504fux5faeux5206ux65b9ux7a0bux57faux7840}

偏微分方程涉及多个自变量，其理论比 ODE 复杂得多。本章从一阶 PDE
和分类入手，引入特征线法。

\subsubsection{91.1
基本概念与分类}\label{ux57faux672cux6982ux5ff5ux4e0eux5206ux7c7b}

\textbf{定义 91.1}（PDE）：\textbf{偏微分方程}（PDE）是包含未知函数
\(u(x_1, \ldots, x_n)\)
及其偏导数的方程。\textbf{阶数}是出现的最高阶偏导数。

\textbf{定义 91.2}（二阶线性 PDE）：一般的二阶线性 PDE 为

\[\sum_{i,j=1}^n a_{ij}(x) \frac{\partial^2 u}{\partial x_i \partial x_j} + \sum_{i=1}^n b_i(x) \frac{\partial u}{\partial x_i} + c(x) u = f(x)\]

其中 \(a_{ij} = a_{ji}\)。

\textbf{定义 91.3}（二阶 PDE 的分类）：在点 \(x\) 处，根据系数矩阵
\(A = (a_{ij}(x))\) 的特征值符号分类： - \textbf{椭圆型}：\(A\)
的所有特征值同号（或 \(A\) 正定/负定）。例：Laplace 方程
\(\Delta u = 0\) - \textbf{抛物型}：\(A\)
有一个零特征值，其余同号。例：热方程 \(u_t = \Delta u\) -
\textbf{双曲型}：\(A\) 有一个特征值与其余符号相反。例：波动方程
\(u_{tt} = \Delta u\)

\subsubsection{91.2 一阶 PDE
与特征线法}\label{ux4e00ux9636-pde-ux4e0eux7279ux5f81ux7ebfux6cd5}

\textbf{定义 91.4}（一阶拟线性 PDE）：一阶拟线性 PDE 的形式为

\[a(x, y, u) \frac{\partial u}{\partial x} + b(x, y, u) \frac{\partial u}{\partial y} = c(x, y, u)\]

\textbf{定理 91.1}（特征线法）：上述 PDE 的解曲面由特征线 ODE 系统生成：

\[\frac{dx}{dt} = a(x, y, u), \quad \frac{dy}{dt} = b(x, y, u), \quad \frac{du}{dt} = c(x, y, u)\]

其中 \(u(t) = u(x(t), y(t))\)。沿特征线，PDE 约化为 ODE。

\emph{证明}：由链式法则，\(\frac{du}{dt} = \frac{\partial u}{\partial x} \frac{dx}{dt} + \frac{\partial u}{\partial y} \frac{dy}{dt} = a \frac{\partial u}{\partial x} + b \frac{\partial u}{\partial y} = c\)。∎

\textbf{例 91.1}（输运方程）：\(u_t + c u_x = 0\)（\(c\)
为常数）的通解为 \(u(x, t) = f(x - ct)\)，其中 \(f\)
由初始条件确定。特征线是 \(x - ct = \text{const}\)。

\textbf{例 91.2}（Burgers 方程）：\(u_t + u u_x = 0\)。特征线为
\(x = u_0(x_0) t + x_0\)。即使初始条件光滑，特征线也可能相交，导致激波（shock
wave）的形成。

\subsubsection{91.3 适定性}\label{ux9002ux5b9aux6027}

\textbf{定义 91.5}（适定性，Hadamard）：PDE
问题称为\textbf{适定的}，如果 1. 解存在 2. 解唯一 3.
解连续依赖于数据（初始条件、边界条件、方程系数）

典型的适定问题： - \textbf{椭圆型}：Dirichlet 问题（给定边界值）或
Neumann 问题（给定法向导数） -
\textbf{抛物型}：初边值问题（给定初始条件和边界条件） -
\textbf{双曲型}：初值问题（Cauchy 问题，给定初始位置和速度）

\subsubsection{91.4 分离变量法}\label{ux5206ux79bbux53d8ux91cfux6cd5}

分离变量法是求解线性齐次 PDE
初边值问题的经典方法，其核心思想是将多变量函数的解分解为单变量函数的乘积，从而将
PDE 约化为若干 ODE。

\textbf{一般步骤}： 1. \textbf{分离假设}：设解具有乘积形式
\(u(x, t) = X(x) T(t)\)（以两变量为例），代入齐次 PDE。 2.
\textbf{变量分离}：将方程整理为''仅含 \(x\) 的函数 = 仅含 \(t\)
的函数''的形式，由于两边独立变量不同，必须等于同一常数
\(-\lambda\)（分离常数）。 3. \textbf{化为 ODE}：得到关于 \(X(x)\) 和
\(T(t)\) 的两个常微分方程 \(X'' + \lambda X = 0\) 和
\(T' + \lambda k T = 0\)（具体形式取决于 PDE
类型），并配以相应的齐次边界条件。 4. \textbf{求解特征值问题}：由
\(X(x)\) 的 ODE 与齐次边界条件构成 Sturm-Liouville
特征值问题，求出特征值 \(\lambda_n\) 和特征函数 \(X_n(x)\)。 5.
\textbf{叠加原理}：通解为所有分离解的线性叠加
\(u(x, t) = \sum_n c_n X_n(x) T_n(t)\)，系数 \(c_n\)
由初始条件通过特征函数的正交性（Fourier 级数展开）确定。

该方法成功的关键在于 PDE
和边界条件均为线性齐次的，从而叠加原理成立。对于非齐次情形，可结合特征函数展开法或
Duhamel 原理处理。

\bigskip\hrule\bigskip

\subsection{第92章：椭圆、抛物与双曲型方程}\label{ux7b2c92ux7ae0ux692dux5706ux629bux7269ux4e0eux53ccux66f2ux578bux65b9ux7a0b}

本章具体讨论三大经典 PDE：Laplace
方程、热方程和波动方程，给出基本解和能量方法。

\subsubsection{92.1 Laplace 方程与 Poisson
方程}\label{laplace-ux65b9ux7a0bux4e0e-poisson-ux65b9ux7a0b}

\textbf{定义 92.1}（Laplace 算子）：\(\mathbb{R}^n\) 中的
\textbf{Laplace 算子}（Laplacian）为
\(\Delta = \sum_{i=1}^n \frac{\partial^2}{\partial x_i^2}\)。

\textbf{定义 92.2}（调和函数）：\(C^2\) 函数 \(u\)
称为\textbf{调和函数}，如果 \(\Delta u = 0\)。

\textbf{定理 92.1}（基本解）：\(\mathbb{R}^n\) 中 Laplace 方程的基本解为

\[\Phi(x) = \begin{cases} -\frac{1}{2\pi} \log |x| & n = 2 \\ \frac{1}{n(n-2)\omega_n} |x|^{2-n} & n \geq 3 \end{cases}\]

其中 \(\omega_n\) 是 \(\mathbb{R}^n\) 中单位球的体积。\(\Phi\) 满足
\(\Delta \Phi = \delta_0\)（在分布意义下）。

\textbf{定理 92.2}（调和函数的均值性质）：若 \(u\) 在
\(B(x, r) \subset \mathbb{R}^n\) 中调和，则

\[u(x) = \frac{1}{|\partial B(x, r)|} \int_{\partial B(x, r)} u(y) dS(y) = \frac{1}{|B(x, r)|} \int_{B(x, r)} u(y) dy\]

\emph{证明}：由散度定理，\(\frac{d}{dr} \left(\frac{1}{|\partial B(x, r)|} \int_{\partial B(x, r)} u(y) dS(y)\right) = 0\)，故球面平均值为常数。令
\(r \to 0\) 得 \(u(x)\)。∎

\textbf{推论
92.3}（极大值原理）：调和函数在区域内部不能取到严格极大值或极小值（除非为常数）。\(\max_{\overline{\Omega}} |u| = \max_{\partial\Omega} |u|\)。

\textbf{定理 92.4}（Poisson 方程的解）：Poisson 方程 \(-\Delta u = f\)
在 \(\mathbb{R}^n\) 中的一个解为

\[u(x) = (\Phi * f)(x) = \int_{\mathbb{R}^n} \Phi(x - y) f(y) dy\]

\subsubsection{92.2 热方程}\label{ux70edux65b9ux7a0b}

\textbf{定义 92.3}（热方程）：热方程（扩散方程）为
\(u_t - \Delta u = 0\)。

\textbf{定义 92.4}（热核 / 基本解）：\(\mathbb{R}^n\)
中热方程的基本解（热核）为

\[\Phi(x, t) = \frac{1}{(4\pi t)^{n/2}} e^{-|x|^2 / (4t)} \quad (t > 0)\]

\(\Phi\) 满足 \(\Phi_t - \Delta \Phi = 0\)（\(t > 0\)）且
\(\lim_{t \to 0^+} \Phi(\cdot, t) = \delta_0\)。

\textbf{定理 92.5}（初值问题的解）：\(\mathbb{R}^n\) 中热方程初值问题
\(u_t = \Delta u\)，\(u(x, 0) = g(x)\) 的解为

\[u(x, t) = (\Phi(\cdot, t) * g)(x) = \frac{1}{(4\pi t)^{n/2}} \int_{\mathbb{R}^n} e^{-|x-y|^2/(4t)} g(y) dy\]

\textbf{定理
92.6}（无穷传播速度与极大值原理）：热方程的解具有无穷传播速度（初始扰动瞬间影响全空间）。对热方程，有极大值原理：

\[\max_{\overline{\Omega} \times [0, T]} u = \max\{\max_{\overline{\Omega}} u(\cdot, 0), \max_{\partial\Omega \times [0, T]} u\}\]

即最大值在初始时刻或边界上达到（抛物型极大值原理）。

\subsubsection{92.3 波动方程}\label{ux6ce2ux52a8ux65b9ux7a0b}

\textbf{定义 92.5}（波动方程）：波动方程为
\(u_{tt} - c^2 \Delta u = 0\)（\(c > 0\) 是波速）。

\textbf{定理 92.7}（d'Alembert 公式，一维）：一维波动方程
\(u_{tt} = c^2 u_{xx}\)，\(u(x, 0) = g(x)\)，\(u_t(x, 0) = h(x)\) 的解为

\[u(x, t) = \frac{1}{2}[g(x + ct) + g(x - ct)] + \frac{1}{2c} \int_{x-ct}^{x+ct} h(s) ds\]

\emph{证明}：做变量变换 \(\xi = x + ct\)，\(\eta = x - ct\)，方程化为
\(u_{\xi\eta} = 0\)，通解为 \(u = F(\xi) + G(\eta)\)。由初始条件确定
\(F\) 和 \(G\)。∎

\textbf{定理 92.8}（Kirchhoff 公式，三维）：三维波动方程的解为

\[u(x, t) = \frac{1}{4\pi c^2 t} \iint_{\partial B(x, ct)} h(y) dS(y) + \frac{\partial}{\partial t}\left(\frac{1}{4\pi c^2 t} \iint_{\partial B(x, ct)} g(y) dS(y)\right)\]

（球面均值法）。三维波动方程满足 Huygens
原理（强形式）：扰动沿球面传播，有清晰的前后波阵面。

\textbf{定理 92.9}（能量守恒）：对波动方程
\(u_{tt} = \Delta u\)（在有界区域 \(\Omega\) 上，齐次边界条件），能量

\[E(t) = \frac{1}{2} \int_\Omega (u_t^2 + |\nabla u|^2) dx\]

是守恒量：\(\frac{dE}{dt} = 0\)。

\emph{证明}：\(\frac{dE}{dt} = \int_\Omega (u_t u_{tt} + \nabla u \cdot \nabla u_t) dx = \int_\Omega u_t(u_{tt} - \Delta u) dx + \int_{\partial\Omega} u_t \frac{\partial u}{\partial n} dS = 0\)。∎

\bigskip\hrule\bigskip

\subsection{第93章：Sobolev
空间与弱解}\label{ux7b2c93ux7ae0sobolev-ux7a7aux95f4ux4e0eux5f31ux89e3}

现代 PDE 理论的核心工具是 Sobolev
空间。它将经典导数推广为弱导数（分布导数），使得在更广的函数空间中寻找
PDE 的解成为可能。

\subsubsection{93.1 弱导数与 Sobolev
空间}\label{ux5f31ux5bfcux6570ux4e0e-sobolev-ux7a7aux95f4}

\textbf{定义 93.1}（弱导数）：设 \(\Omega \subseteq \mathbb{R}^n\)
是开集，\(u \in L^1_{\text{loc}}(\Omega)\)。\(u\)
的\textbf{弱导数}（或分布导数）\(v = D^\alpha u\)（\(L^1_{\text{loc}}(\Omega)\)
中）满足

\[\int_\Omega u D^\alpha \varphi \, dx = (-1)^{|\alpha|} \int_\Omega v \varphi \, dx \quad \text{对所有 } \varphi \in C_c^\infty(\Omega)\]

其中 \(D^\alpha \varphi\) 是 \(\varphi\) 的经典 \(\alpha\) 阶偏导数。

\textbf{定义 93.2}（Sobolev 空间）：对 \(1 \leq p \leq \infty\) 和整数
\(k \geq 0\)，\textbf{Sobolev 空间} \(W^{k,p}(\Omega)\) 定义为

\[W^{k,p}(\Omega) = \{u \in L^p(\Omega) : D^\alpha u \in L^p(\Omega) \text{ 对所有 } |\alpha| \leq k\}\]

范数为
\(\|u\|_{W^{k,p}} = \left(\sum_{|\alpha| \leq k} \|D^\alpha u\|_{L^p}^p\right)^{1/p}\)（\(p < \infty\)）。\(H^k(\Omega) = W^{k,2}(\Omega)\)
是 Hilbert 空间。

\textbf{命题 93.1}：\(W^{k,p}(\Omega)\) 是 Banach 空间。\(H^k(\Omega)\)
是 Hilbert 空间，内积为
\((u, v)_{H^k} = \sum_{|\alpha| \leq k} \int_\Omega D^\alpha u \cdot D^\alpha v \, dx\)。

\subsubsection{93.2 Sobolev
嵌入定理}\label{sobolev-ux5d4cux5165ux5b9aux7406}

\textbf{定理 93.2}（Sobolev 嵌入定理）：设
\(\Omega \subseteq \mathbb{R}^n\) 是有 Lipschitz
边界的有界开集，\(1 \leq p < \infty\)。

\begin{enumerate}
\def\labelenumi{\arabic{enumi}.}
\tightlist
\item
  若 \(k p < n\)，则
  \(W^{k,p}(\Omega) \hookrightarrow L^q(\Omega)\)（连续嵌入），其中
  \(\frac{1}{q} = \frac{1}{p} - \frac{k}{n}\)（Sobolev 共轭指数）
\item
  若 \(k p = n\)，则 \(W^{k,p}(\Omega) \hookrightarrow L^q(\Omega)\)
  对所有 \(q < \infty\) 成立
\item
  若 \(k p > n\)，则
  \(W^{k,p}(\Omega) \hookrightarrow C^{0,\gamma}(\overline{\Omega})\)（Hölder
  连续函数空间），其中 \(\gamma = k - n/p\)（若 \(k - n/p < 1\)）
\end{enumerate}

\textbf{定理 93.3}（Rellich-Kondrachov 紧嵌入）：在上述条件下，若
\(k p < n\) 且 \(1 \leq q < p^*\)（\(p^* = np/(n - kp)\)），则嵌入
\(W^{k,p}(\Omega) \hookrightarrow L^q(\Omega)\) 是紧的。

\subsubsection{93.3
弱解与变分方法}\label{ux5f31ux89e3ux4e0eux53d8ux5206ux65b9ux6cd5}

\textbf{定义 93.3}（弱解）：考虑 Poisson 方程的 Dirichlet 问题：

\[-\Delta u = f \text{ 在 } \Omega \text{ 中}, \quad u = 0 \text{ 在 } \partial\Omega \text{ 上}\]

\(u \in H_0^1(\Omega)\) 称为\textbf{弱解}，如果对任意
\(v \in H_0^1(\Omega)\)（\(H_0^1\) 是 \(C_c^\infty(\Omega)\) 在 \(H^1\)
范数下的闭包），

\[\int_\Omega \nabla u \cdot \nabla v \, dx = \int_\Omega f v \, dx\]

\textbf{定理 93.4}（Lax-Milgram 定理）：设 \(H\) 是 Hilbert
空间，\(B: H \times H \to \mathbb{R}\) 是双线性型，满足： -
有界性：\(|B(u, v)| \leq M \|u\| \|v\|\) -
强制性：\(B(u, u) \geq \alpha \|u\|^2\)（\(\alpha > 0\)）

则对任意 \(F \in H^*\)（\(H\) 的对偶空间），存在唯一的 \(u \in H\) 使得
\(B(u, v) = F(v)\) 对所有 \(v \in H\) 成立。且
\(\|u\| \leq \frac{1}{\alpha} \|F\|_{H^*}\)。

\emph{证明}：对固定 \(u\)，\(v \mapsto B(u, v)\) 是 \(H\)
上的有界线性泛函，由 Riesz 表示定理（V13 定理 73.1），存在 \(Au \in H\)
使得 \(B(u, v) = (Au, v)\)。\(A: H \to H\) 是有界线性算子，强制性推出
\(A\) 可逆。∎

\textbf{定理 93.5}（Poisson 方程弱解的存在唯一性）：对
\(f \in L^2(\Omega)\)，Poisson 方程的 Dirichlet 问题存在唯一弱解
\(u \in H_0^1(\Omega)\)。

\emph{证明}：定义双线性型
\(B(u, v) = \int_\Omega \nabla u \cdot \nabla v \, dx\) 和线性泛函
\(F(v) = \int_\Omega f v \, dx\)。由 Poincaré
不等式，\(B(u, u) = \|\nabla u\|_{L^2}^2 \geq C \|u\|_{H^1}^2\)（强制性）。由
Lax-Milgram 定理即得。∎

\subsubsection{93.4 正则性理论}\label{ux6b63ux5219ux6027ux7406ux8bba}

\textbf{定理 93.6}（椭圆正则性）：设 \(u \in H_0^1(\Omega)\) 是
\(-\Delta u = f\) 的弱解，\(\Omega\) 有光滑边界。 1. 若
\(f \in L^2(\Omega)\)，则 \(u \in H^2(\Omega)\) 且
\(\|u\|_{H^2} \leq C \|f\|_{L^2}\) 2. 若 \(f \in H^k(\Omega)\)，则
\(u \in H^{k+2}(\Omega)\) 3. 若
\(f \in C^\infty(\overline{\Omega})\)，则
\(u \in C^\infty(\overline{\Omega})\)

由此，弱解在光滑数据下是经典解，PDE 的弱解理论与经典解理论一致。

\textbf{定理 93.7}（Weyl 引理）：若 \(u \in L^1_{\text{loc}}(\Omega)\)
满足 \(\int_\Omega u \Delta \varphi \, dx = 0\) 对所有
\(\varphi \in C_c^\infty(\Omega)\)，则
\(u\)（几乎处处等于）一个调和函数。即弱调和函数是经典调和函数。

\bigskip\hrule\bigskip

\subsection{第94章：D-模理论简介}\label{ux7b2c94ux7ae0d-ux6a21ux7406ux8bbaux7b80ux4ecb}

D-模理论将线性偏微分方程的研究代数化，是代数几何、表示论和数学物理的重要工具。设\(X\)是光滑代数簇，记\(\mathcal{D}_X\)为\(X\)上的微分算子层，局部上同构于Weyl代数\(\mathbb{C}[x_1,\ldots,x_n]\langle\partial_1,\ldots,\partial_n\rangle\)（满足\([\partial_i, x_j] = \delta_{ij}\)）。线性PDE系统\(P_1u = \cdots = P_ku = 0\)对应于左\(\mathcal{D}_X\)-模\(\mathcal{M} = \mathcal{D}_X/(\mathcal{D}_X P_1 + \cdots + \mathcal{D}_X P_k)\)，其解空间同构于\(\operatorname{Hom}_{\mathcal{D}_X}(\mathcal{M}, \mathcal{O}_X)\)。这一代数化使得上同调方法得以系统运用。\(\mathcal{M}\)的特征簇\(\operatorname{Ch}(\mathcal{M})\)是余切丛\(T^*X\)中的闭代数子集，描述了奇异性的相位空间分布；其维数由Bernstein不等式控制：\(\dim \operatorname{Ch}(\mathcal{M}) \geq \dim X\)。Riemann-Hilbert对应是D-模理论的巅峰：在\(X\)光滑时，正则奇点D-模的导出范畴与\(\mathbb{C}\)-perverse层的导出范畴等价，将分析的奇异性问题转化为拓扑的层论问题。在表示论中，Beilinson-Bernstein局部化定理将半单Lie代数的表示范畴实现为旗簇上D-模的范畴，为Kazhdan-Lusztig猜想提供了几何框架。

\bigskip\hrule\bigskip

\subsection{第94章：特殊方程与推广}\label{ux7b2c94ux7ae0ux7279ux6b8aux65b9ux7a0bux4e0eux63a8ux5e7f}

\subsubsection{94.1 Painleve方程}\label{painleveux65b9ux7a0b}

Painleve方程是六类二阶非线性常微分方程（P\({}_{\rm I}\)--P\({}_{\rm VI}\)），其核心特征是拥有
\textbf{Painleve性质}：所有可移动奇点（依赖于积分常数的奇点）仅为极点，无本性奇点或分支点。Painleve（1900）与
Gambier 系统分类了所有满足此性质的形如 \(y'' = F(y', y, x)\)
的二阶方程，除经典可积者外恰有六族。最著名的为
P\({}_{\rm I}\)：\(y'' = 6y^2 + x\) 和
P\({}_{\rm II}\)：\(y'' = 2y^3 + xy + \alpha\)，其解称为
Painleve超越函数。Painleve方程与完全可积系统深刻相关：KdV 方程等可积 PDE
的相似约化常导出 Painleve
方程；Ablowitz-Ramani-Segur（ARS）猜想以此为基础判断 PDE 的 Painleve
可积性。

\subsubsection{94.2
泛函微分方程}\label{ux6cdbux51fdux5faeux5206ux65b9ux7a0b}

泛函微分方程（FDE）是导数依赖于过去（甚至未来）状态的微分方程。最重要的子类------\textbf{时滞微分方程}------形如
\(x'(t) = f(t, x(t), x(t-\tau))\)（\(\tau > 0\)
为时滞）。其基本理论使用步进法（method of steps）建立：将时间轴按
\(\tau\) 分段，逐段转化为 ODE
求解，由此可得存在唯一性和连续依赖性。稳定性分析的核心工具是
\textbf{Lyapunov-Krasovskii 泛函方法}------构造依赖于状态历史
\(x_t(\theta) = x(t+\theta)（\theta \in [-\tau,0])\) 的泛函
\(V(x_t)\)，使其沿解的导数为负定以判定渐近稳定性。在种群动力学中，时滞
Logistic 方程 \(\dot{N} = rN(1 - N(t-\tau)/K)\)
可产生周期振荡；在神经网络中，Hopfield 模型的传输延迟
\(\dot{u}_i = -u_i/\tau_i + \sum_j w_{ij} g(u_j(t-\delta_{ij}))\)
对网络动力学有本质影响。

\bigskip\hrule\bigskip

\subsection{第95章：混合型偏微分方程（Mixed-Type
PDEs）}\label{ux7b2c95ux7ae0ux6df7ux5408ux578bux504fux5faeux5206ux65b9ux7a0bmixed-type-pdes}

混合型偏微分方程在区域的某一部分为椭圆型，另一部分为双曲型，其典型代表是
\textbf{Tricomi 方程} \(y u_{xx} + u_{yy} = 0\)：当 \(y > 0\)
时为椭圆型，\(y < 0\) 时为双曲型，\(y = 0\)
为抛物退化线。这类方程在跨音速空气动力学（无激波跨音速流动）中自然出现------描述流场从亚音速（椭圆区域）到超音速（双曲区域）的平缓过渡。\textbf{核心理论}：（1）\textbf{Tricomi
问题}------在混合区域上指定适定的边界条件（椭圆部分给
Dirichlet，双曲部分在两条特征线上分别给不同的条件），F. Tricomi (1923)
通过将问题化归为奇异积分方程证明了存在性；（2）\textbf{Frankl
问题}------推广的边界值问题，出现于跨音速喷管设计，要求解在退化线上连续但允许法向导数跳跃；（3）\textbf{Keldysh
型方程}------另类退化方向 \(u_{xx} + y^m u_{yy} = 0\)（\(m\)
为奇数）及其奇异 Cauchy 问题。\textbf{现代进展}：Morawetz (1956)
用乘子法得到 Tricomi 问题解的先验估计；Lupo-Payne (2004)
使用变分方法和临界点理论处理非线性混合型方程。混合型 PDE
在跨音速流动、等离子体物理（磁面方程）和金融数学中有实际应用。

\bigskip\hrule\bigskip

\subsection{第96章：KAM
理论与自由边界问题}\label{ux7b2c96ux7ae0kam-ux7406ux8bbaux4e0eux81eaux7531ux8fb9ux754cux95eeux9898}

\subsubsection{96.1 KAM
理论（Kolmogorov-Arnold-Moser）}\label{kam-ux7406ux8bbakolmogorov-arnold-moser}

KAM 理论是 Hamilton
动力系统微扰理论的核心成就，描述了近可积系统在微小扰动下拟周期运动的保持。\textbf{基本问题}：给定可积
Hamilton 系统 \(H_0(I)\)（作用量-角变量坐标，环面
\(I \times \mathbb{T}^n\)），当加入小扰动
\(H_\epsilon = H_0(I) + \epsilon H_1(I, \theta)\)
后，哪些不变环面得以存留？

\textbf{定理 96.1}（Kolmogorov 定理，1954）：若 \(H_0\)
非退化（\(\det(\partial^2 H_0 / \partial I_i \partial I_j) \neq 0\)），则对满足
Diophantine 条件

\[|k \cdot \omega| \geq \frac{\gamma}{\|k\|^\tau}, \quad \forall k \in \mathbb{Z}^n \setminus \{0\}\]

的频率 \(\omega = \nabla H_0(I)\)，相应的不变环面在充分小的
\(\epsilon > 0\) 下扰动后仍然存在------仅发生微小形变，其频率集合具有正
Lebesgue 测度（当 \(\epsilon \to 0\) 时测度趋于全测度）。

\textbf{定义 96.1}（Diophantine 条件）：频率向量
\(\omega \in \mathbb{R}^n\) 称为 \((\gamma, \tau)\)-Diophantine
的（\(\gamma > 0\), \(\tau > n-1\)），若对任意非零整数向量
\(k \in \mathbb{Z}^n\)，有
\(|k \cdot \omega| \geq \gamma \|k\|^{-\tau}\)。满足此条件的频率集合在
\(\mathbb{R}^n\) 中具有全测度。

\textbf{定理 96.2}（Arnold 扩散，1964）：未被 KAM
环面覆盖的相空间区域中存在极缓慢的扩散现象（速度为
\(O(e^{-c/\epsilon})\)），这是导致不变环面间''泄露''的核心机制。对于
\(n \geq 3\) 自由度的系统，Arnold 证明了存在穿越由 KAM
环面形成的''网''的轨道链。

\subsubsection{96.2 自由边界问题（Free Boundary
Problems）}\label{ux81eaux7531ux8fb9ux754cux95eeux9898free-boundary-problems}

自由边界问题处理 PDE 的解域边界本身是未知的、由解决定的边值问题。

\textbf{定义 96.2}（Stefan 问题 / 冰-水相变）：热方程 \(u_t = \Delta u\)
在固相/液相区域分别成立，在自由界面 \(\Gamma(t)\) 上满足 Stefan 条件

\[[\nabla u \cdot \mathbf{n}] = L V_n\]

其中 \(L > 0\) 为潜热，\(V_n\) 为界面法向速度，\([\cdot]\)
表示界面两侧跳跃。一维 Stefan 问题的显式解（Neumann 解，1860
年代）由误差函数 \(\operatorname{erf}\) 构成，自由边界以
\(\sim\sqrt{t}\) 的速率移动。

\textbf{定义 96.3}（障碍问题 / Obstacle Problem）：给定障碍函数
\(\psi\)，求解 \(u \geq \psi\) 满足

\[\Delta u \leq 0, \quad u \geq \psi, \quad (u - \psi) \Delta u = 0\]

其自由边界是接触集 \(\Gamma = \partial \{u = \psi\} \cap \Omega\)（即
\(u = \psi\) 区域的拓扑边界）。

\textbf{定理 96.3}（Caffarelli
正则性定理，1977）：障碍问题的自由边界在除去一个低维奇异集之外是光滑的（\(C^{1,\alpha}\)），这为整个自由边界正则性理论奠定了基础。具体而言，自由边界点分为正则点（自由边界在该点附近为
\(C^{1,\alpha}\) 超曲面）和奇点（可分类），后者的 Hausdorff 维数不超过
\(n-2\)。

\textbf{定义 96.4}（Hele-Shaw
流动）：描述两平行平板间粘性流体的界面演化。对 Hele-Shaw 胞腔，压强
\(p\) 在流体区域满足 \(\Delta p = 0\)，在自由边界上 \(p = 0\) 且法向速度
\(V_n = -\partial p / \partial n\)。其数学模型等价于 Laplacian 增长（DLA
的连续极限），小扰动下界面不稳定（Saffman-Taylor 指进现象）。

\bigskip\hrule\bigskip

\emph{本卷建立了微分方程从经典到现代的理论框架：ODE
的存在唯一性理论（Picard-Lindelöf 定理）和稳定性理论（Lyapunov
方法）是动力系统研究的基础；三大经典 PDE（Laplace
方程、热方程、波动方程）的解法展示了 PDE 的丰富结构和物理背景；Sobolev
空间和弱解理论为现代 PDE 研究提供了强大的泛函分析框架，Lax-Milgram
定理和椭圆正则性理论是 PDE
弱解存在性和正则性的核心工具。为后续数值分析（V22）和调和分析（V24）的学习奠定了基础。}

\newpage

\section{卷十八：范畴论}\label{ux5377ux5341ux516bux8303ux7574ux8bba}

\begin{quote}
范畴论是''数学的数学''，它提供了统一的语言和框架来描述数学结构及其相互关系。本卷从范畴、函子和自然变换的基本概念出发，建立极限、伴随函子、单子和
Abel 范畴的理论，为同调代数（V12 Ch
70）、代数拓扑（V14）和代数几何（V15）提供统一的范畴论基础。
\end{quote}

\bigskip\hrule\bigskip

\subsection{第94章：范畴、函子与自然变换}\label{ux7b2c94ux7ae0ux8303ux7574ux51fdux5b50ux4e0eux81eaux7136ux53d8ux6362}

范畴论的基本三位一体------范畴、函子和自然变换------提供了描述数学结构的统一语言。Yoneda
引理是范畴论最基本的定理，它表明一个对象完全由它与其他对象的关系决定。

\subsubsection{94.1 范畴的定义}\label{ux8303ux7574ux7684ux5b9aux4e49}

\textbf{定义 94.1}（范畴）：一个\textbf{范畴} \(\mathcal{C}\)
由以下数据组成： - 一族\textbf{对象} \(\operatorname{Ob}(\mathcal{C})\)
- 对每对对象 \(X, Y\)，一个\textbf{态射集}
\(\operatorname{Hom}_{\mathcal{C}}(X, Y)\)（或 \(\mathcal{C}(X, Y)\)） -
对每个对象 \(X\)，一个\textbf{恒等态射}
\(\operatorname{id}_X \in \mathcal{C}(X, X)\) - 对每三个对象
\(X, Y, Z\)，一个\textbf{合成映射}
\(\circ: \mathcal{C}(Y, Z) \times \mathcal{C}(X, Y) \to \mathcal{C}(X, Z)\)

满足：合成结合律 \((h \circ g) \circ f = h \circ (g \circ f)\)，恒等律
\(\operatorname{id}_Y \circ f = f = f \circ \operatorname{id}_X\)。

\textbf{例 94.1}（基本范畴）： -
\textbf{Set}：集合范畴（对象：集合，态射：函数） -
\textbf{Grp}：群范畴（对象：群，态射：群同态） - \textbf{Ab}：Abel
群范畴 - \textbf{Ring}：环范畴 -
\textbf{Top}：拓扑空间范畴（对象：拓扑空间，态射：连续映射） -
\textbf{Vect\(_k\)}：\(k\)-向量空间范畴 - \textbf{R-Mod}：左
\(R\)-模范畴

\textbf{定义 94.2}（小范畴与局部小范畴）：范畴 \(\mathcal{C}\)
称为\textbf{局部小}的，如果每个 \(\mathcal{C}(X, Y)\)
是集合。\(\mathcal{C}\) 称为\textbf{小范畴}，如果
\(\operatorname{Ob}(\mathcal{C})\) 也是集合。

\textbf{定义 94.3}（同构）：态射 \(f: X \to Y\)
称为\textbf{同构}，如果存在 \(g: Y \to X\) 使得
\(g \circ f = \operatorname{id}_X\)，\(f \circ g = \operatorname{id}_Y\)。\(g\)
称为 \(f\) 的逆。

\textbf{定义 94.4}（对偶范畴）：范畴 \(\mathcal{C}\) 的\textbf{对偶范畴}
\(\mathcal{C}^{\text{op}}\) 与 \(\mathcal{C}\)
有相同的对象，但态射方向反转：\(\mathcal{C}^{\text{op}}(X, Y) = \mathcal{C}(Y, X)\)。

\subsubsection{94.2 函子}\label{ux51fdux5b50}

\textbf{定义 94.5}（函子）：范畴 \(\mathcal{C}\) 到 \(\mathcal{D}\)
的\textbf{函子} \(F: \mathcal{C} \to \mathcal{D}\) 由以下组成： -
对象映射：\(X \mapsto F(X)\) - 态射映射：对每个
\(f: X \to Y\)，\(F(f): F(X) \to F(Y)\)

满足：\(F(\operatorname{id}_X) = \operatorname{id}_{F(X)}\)，\(F(g \circ f) = F(g) \circ F(f)\)。

\textbf{定义
94.6}（反变函子）：\(\mathcal{C}^{\text{op}} \to \mathcal{D}\)
的函子称为 \(\mathcal{C}\) 到 \(\mathcal{D}\)
的\textbf{反变函子}。它反转态射方向：\(F(f: X \to Y): F(Y) \to F(X)\)。

\textbf{例 94.2}（常见函子）： - 遗忘函子
\(U: \mathbf{Grp} \to \mathbf{Set}\)（忘记群结构，只保留集合） -
自由函子 \(F: \mathbf{Set} \to \mathbf{Grp}\)（生成自由群） - 基本群函子
\(\pi_1: \mathbf{Top}_* \to \mathbf{Grp}\)（V11 Ch 60） - 同调函子
\(H_n: \mathbf{Top} \to \mathbf{Ab}\)（V14 Ch 77） - Hom
函子：\(\mathcal{C}(A, -): \mathcal{C} \to \mathbf{Set}\)（协变），\(\mathcal{C}(-, A): \mathcal{C}^{\text{op}} \to \mathbf{Set}\)（反变）

\subsubsection{94.3 自然变换}\label{ux81eaux7136ux53d8ux6362}

\textbf{定义 94.7}（自然变换）：设 \(F, G: \mathcal{C} \to \mathcal{D}\)
是两个函子。\textbf{自然变换} \(\alpha: F \Rightarrow G\) 由对每个
\(X \in \mathcal{C}\) 的一个态射 \(\alpha_X: F(X) \to G(X)\)
组成，使得对每个 \(f: X \to Y\)，下图交换：

\[
\begin{CD}
F(X) @>{\alpha_X}>> G(X) \\
@V{F(f)}VV @VV{G(f)}V \\
F(Y) @>>{\alpha_Y}> G(Y)
\end{CD}
\]

若每个 \(\alpha_X\) 是同构，则 \(\alpha\) 称为\textbf{自然同构}，记作
\(F \cong G\)。

\textbf{定义 94.8}（函子范畴）：从 \(\mathcal{C}\) 到 \(\mathcal{D}\)
的函子及它们之间的自然变换构成\textbf{函子范畴}
\([\mathcal{C}, \mathcal{D}]\)（或 \(\mathcal{D}^{\mathcal{C}}\)）。

\textbf{例 94.3}（自然变换的例子）： - 行列式
\(\det: \mathbf{GL}_n \Rightarrow \mathbf{GL}_1\)（\(k^\times\)）是自然变换
- 有限维向量空间 \(V\) 到其双重对偶 \(V \to V^{**}\)
的嵌入是自然变换（\(V \cong V^{**}\) 自然同构） - Hurewicz 同态
\(\pi_n \Rightarrow H_n\) 是自然变换（V14）

\subsubsection{94.4 Yoneda 引理}\label{yoneda-ux5f15ux7406}

\textbf{定义 94.9}（Yoneda 嵌入）：对局部小范畴 \(\mathcal{C}\)，定义
\textbf{Yoneda 嵌入}

\[Y: \mathcal{C} \to [\mathcal{C}^{\text{op}}, \mathbf{Set}], \quad Y(X) = \mathcal{C}(-, X)\]

即 \(Y(X)\) 是反变 Hom 函子：\(Y(X)(A) = \mathcal{C}(A, X)\)。

\textbf{定理 94.1}（Yoneda 引理）：设 \(\mathcal{C}\)
是局部小范畴，\(F: \mathcal{C}^{\text{op}} \to \mathbf{Set}\)
是函子，\(X \in \mathcal{C}\)。则存在自然的一一对应

\[\operatorname{Nat}(\mathcal{C}(-, X), F) \cong F(X)\]

（其中 \(\operatorname{Nat}\) 表示自然变换的集合）。此同构由
\(\alpha \mapsto \alpha_X(\operatorname{id}_X)\) 给出。

\emph{证明}：定义映射
\(\Phi: \operatorname{Nat}(\mathcal{C}(-, X), F) \to F(X)\)，\(\Phi(\alpha) = \alpha_X(\operatorname{id}_X)\)。逆映射
\(\Psi: F(X) \to \operatorname{Nat}(\mathcal{C}(-, X), F)\) 定义为：对
\(u \in F(X)\)，\(\Psi(u)_A(f) = F(f)(u)\)（\(f: A \to X\)）。验证
\(\Phi \circ \Psi = \operatorname{id}\) 和
\(\Psi \circ \Phi = \operatorname{id}\)。∎

\textbf{推论 94.2}（Yoneda 嵌入是完全忠实的）：Yoneda 嵌入
\(Y: \mathcal{C} \to [\mathcal{C}^{\text{op}}, \mathbf{Set}]\)
是完全忠实的：\(\operatorname{Nat}(\mathcal{C}(-, X), \mathcal{C}(-, Y)) \cong \mathcal{C}(X, Y)\)。特别地，\(X \cong Y\)
当且仅当 \(\mathcal{C}(-, X) \cong \mathcal{C}(-, Y)\)。

\emph{证明}：在 Yoneda 引理（定理 94.1）中取
\(F = \mathcal{C}(-, Y)\)，则
\(\operatorname{Nat}(\mathcal{C}(-, X), \mathcal{C}(-, Y)) \cong \mathcal{C}(X, Y)\)。具体地，自然变换
\(\alpha: \mathcal{C}(-, X) \Rightarrow \mathcal{C}(-, Y)\) 由
\(\alpha_X(\operatorname{id}_X) \in \mathcal{C}(X, Y)\)
唯一决定；反之，每个 \(f: X \to Y\) 通过后合成 \(f \circ (-)\)
给出自然变换。这建立了双向一一对应，故 Yoneda 嵌入是全忠实函子。∎

\textbf{推论 94.3}（Yoneda 推论）：一个对象 \(X\)
完全由它与其他对象的关系（即函子
\(\mathcal{C}(-, X)\)）决定（在同构意义下）。这是范畴论的基本哲学。

\bigskip\hrule\bigskip

\subsection{第95章：极限与余极限}\label{ux7b2c95ux7ae0ux6781ux9650ux4e0eux4f59ux6781ux9650}

极限和余极限是范畴论中最重要的构造，统一了积、余积、等化子、拉回、推出等概念。

\subsubsection{95.1
极限的普遍定义}\label{ux6781ux9650ux7684ux666eux904dux5b9aux4e49}

\textbf{定义 95.1}（图与锥）：设 \(J\)
是小范畴（称为\textbf{指标范畴}）。函子 \(D: J \to \mathcal{C}\) 称为
\(\mathcal{C}\) 中的一个 \textbf{\(J\) 型图}。

\(D\) 上的一个\textbf{锥}（cone）由对象 \(C \in \mathcal{C}\) 和一族态射
\(\{c_j: C \to D(j)\}_{j \in J}\) 组成，使得对每个 \(J\) 中的态射
\(u: j \to k\)，\(D(u) \circ c_j = c_k\)。

\textbf{定义 95.2}（极限）：图 \(D: J \to \mathcal{C}\) 的\textbf{极限}
\(\varprojlim D\)（或 \(\lim D\)）是 \(D\) 上的一个万有锥：它是一个锥
\((L, \{\ell_j\})\)，使得对任意其他锥 \((C, \{c_j\})\)，存在唯一的态射
\(f: C \to L\) 使得 \(\ell_j \circ f = c_j\)（对所有 \(j\)）。

\textbf{定义 95.3}（余极限）：对偶地，\(D\)
上的\textbf{余锥}及\textbf{余极限} \(\varinjlim D\)（或
\(\operatorname{colim} D\)）是极限在 \(\mathcal{C}^{\text{op}}\)
中的对偶概念。

\textbf{命题 95.1}：若极限存在，则它在同构意义下唯一。

\subsubsection{95.2
特殊极限与余极限}\label{ux7279ux6b8aux6781ux9650ux4e0eux4f59ux6781ux9650}

\textbf{定义 95.4}（积与余积）：设 \(J\)
是离散范畴（只有恒等态射）。\(J\) 型图的极限是\textbf{积}
\(\prod_{j \in J} X_j\)，余极限是\textbf{余积}
\(\coprod_{j \in J} X_j\)。

\begin{itemize}
\tightlist
\item
  \textbf{Set} 中：积 = 笛卡尔积，余积 = 不交并
\item
  \textbf{Grp} 中：积 = 直积，余积 = 自由积
\item
  \textbf{Ab} 中：积 = 直积，余积 = 直和（有限时两者相同）
\item
  \textbf{Vect\(_k\)} 中：积 = 直积，余积 = 直和
\end{itemize}

\textbf{定义 95.5}（等化子与余等化子）：两个平行态射
\(f, g: X \rightrightarrows Y\) 的图。 - \textbf{等化子}
\(\operatorname{eq}(f, g)\) 是使得 \(f \circ e = g \circ e\) 的万有态射
\(e: E \to X\) - \textbf{余等化子} \(\operatorname{coeq}(f, g)\) 是使得
\(c \circ f = c \circ g\) 的万有态射 \(c: Y \to C\)

\textbf{定义 95.6}（拉回与推出）： - \textbf{拉回}（pullback）是图
\(X \to Z \leftarrow Y\) 的极限 - \textbf{推出}（pushout）是图
\(X \leftarrow Z \to Y\) 的余极限

\textbf{定理 95.2}（极限的构造）：若 \(\mathcal{C}\)
有所有积和所有等化子，则 \(\mathcal{C}\)
有所有小极限。具体地，\(\varprojlim D\) 是等化子

\[\varprojlim D \longrightarrow \prod_{j \in J} D(j) \rightrightarrows \prod_{u: j \to k} D(k)\]

其中两个态射分别由 \(D(u) \circ \pi_j\) 和 \(\pi_k\) 给出。

\emph{证明}：设图 \(D: J \to \mathcal{C}\)。构造积
\(P = \prod_{j} D(j)\) 和 \(Q = \prod_{u: j \to k} D(k)\)。定义两个态射
\(f, g: P \rightrightarrows Q\)：\(f\) 的第 \(u\) 分量为
\(D(u) \circ \pi_j\)，\(g\) 的第 \(u\) 分量为 \(\pi_k\)。取等化子
\(e: L \to P\)，则 \((L, \pi_j \circ e)\) 构成 \(D\) 的锥。对任意锥
\((C, c_j)\)，泛性质给出唯一态射 \(h: C \to P\) 使得
\(\pi_j \circ h = c_j\)。由锥条件 \(D(u) \circ c_j = c_k\) 推出
\(f \circ h = g \circ h\)，故 \(h\) 唯一地通过等化子 \(e\)
分解，即存在唯一 \(C \to L\)。这验证了 \(L\)
的极限泛性质。拉回作为二元积与等化子的特例自然包含于此构造。∎

\subsubsection{95.3 完备性}\label{ux5b8cux5907ux6027-1}

\textbf{定义 95.7}（完备范畴）：范畴 \(\mathcal{C}\)
称为\textbf{完备}的，如果它有所有小极限。\(\mathcal{C}\)
称为\textbf{余完备}的，如果它有所有小余极限。

\textbf{定理 95.3}：\textbf{Set}
是完备和余完备的。\textbf{Grp}、\textbf{Ab}、\textbf{Ring}、\textbf{R-Mod}、\textbf{Top}
都是完备和余完备的。

\textbf{命题 95.4}（极限的函子性）：若 \(\mathcal{C}\)
完备，则极限构造定义了函子
\(\varprojlim: [J, \mathcal{C}] \to \mathcal{C}\)。类似地，\(\varinjlim\)
是函子。

\textbf{定义 95.8}（保极限的函子）：函子
\(F: \mathcal{C} \to \mathcal{D}\) 称为\textbf{保极限}的，如果对每个图
\(D: J \to \mathcal{C}\)，\(F(\varprojlim D) \cong \varprojlim (F \circ D)\)（自然同构）。

\textbf{例 95.1}：Hom 函子 \(\mathcal{C}(X, -)\)
保极限（所有极限）。\(\mathcal{C}(-, X)\) 将余极限变为极限。

\bigskip\hrule\bigskip

\subsection{第96章：伴随函子}\label{ux7b2c96ux7ae0ux4f34ux968fux51fdux5b50}

伴随函子是范畴论最核心的概念之一，广泛出现在数学的各个分支中。自由-遗忘伴随、积-Hom
伴随、量化词伴随等都是典型例子。

\subsubsection{96.1 伴随的定义}\label{ux4f34ux968fux7684ux5b9aux4e49}

\textbf{定义 96.1}（伴随函子）：设 \(F: \mathcal{C} \to \mathcal{D}\) 和
\(G: \mathcal{D} \to \mathcal{C}\) 是函子。\(F\) 称为 \(G\)
的\textbf{左伴随}（\(G\) 是 \(F\) 的\textbf{右伴随}），记作
\(F \dashv G\)，如果存在自然同构

\[\mathcal{D}(F(X), Y) \cong \mathcal{C}(X, G(Y))\]

对 \(X \in \mathcal{C}\)，\(Y \in \mathcal{D}\) 自然。

\textbf{定义 96.2}（单位与余单位）：伴随 \(F \dashv G\)
给出两个自然变换： -
\textbf{单位}（unit）：\(\eta: \operatorname{id}_{\mathcal{C}} \Rightarrow G \circ F\)（由
\(X \mapsto F(X)\) 的恒等态射经过伴随同构得到） -
\textbf{余单位}（counit）：\(\varepsilon: F \circ G \Rightarrow \operatorname{id}_{\mathcal{D}}\)（对偶地）

满足三角恒等式：\(G\varepsilon \circ \eta G = \operatorname{id}_G\) 和
\(\varepsilon F \circ F\eta = \operatorname{id}_F\)。

\textbf{例 96.1}（典型伴随对）： 1. 自由 \(\dashv\)
遗忘：\(F: \mathbf{Set} \to \mathbf{Grp}\)（自由群）\(\dashv\)
\(U: \mathbf{Grp} \to \mathbf{Set}\)（遗忘函子） 2. 张量 \(\dashv\)
Hom：对 \(R\)-模，\(M \otimes_R - \dashv \operatorname{Hom}_R(M, -)\) 3.
对拓扑空间范畴 \(\mathbf{Top}\) 与集合范畴
\(\mathbf{Set}\)，\(\pi_0 \dashv \text{离散} \dashv U\)（其中 \(\pi_0\)
为连通分支函子，\(\text{离散}\) 为离散拓扑函子，\(U\) 为遗忘函子） 4.
\(\Delta \dashv \varprojlim\)（对角函子 \(\dashv\)
极限函子）：\(\varinjlim \dashv \Delta \dashv \varprojlim\)

\subsubsection{96.2
伴随函子的性质}\label{ux4f34ux968fux51fdux5b50ux7684ux6027ux8d28}

\textbf{定理 96.1}（左伴随保余极限，右伴随保极限）：若
\(F \dashv G\)，则 \(F\) 保所有余极限，\(G\) 保所有极限。

\emph{证明}：对余极限
\(C = \varinjlim D_j\)，\(\mathcal{D}(F(C), Y) \cong \mathcal{C}(C, G(Y)) \cong \varprojlim \mathcal{C}(D_j, G(Y)) \cong \varprojlim \mathcal{D}(F(D_j), Y) \cong \mathcal{D}(\varinjlim F(D_j), Y)\)。由
Yoneda 引理，\(F(C) \cong \varinjlim F(D_j)\)。∎

\textbf{命题 96.2}：若 \(F \dashv G\) 且 \(F' \dashv G\)，则
\(F \cong F'\)（左伴随在同构意义下唯一）。类似地，右伴随也唯一。

\textbf{定理 96.2}（一般伴随函子定理，Freyd）：设 \(\mathcal{C}\)
是完备范畴，\(G: \mathcal{D} \to \mathcal{C}\) 是函子。则 \(G\)
有左伴随当且仅当： 1. \(G\) 保所有小极限 2. 对每个
\(X \in \mathcal{C}\)，逗号范畴 \(X \downarrow G\) 有初始对象

（还需要''解集条件''（solution set condition）的适当版本）。

\emph{证明概要}：充分性：\(G\) 保极限且满足解集条件 \(\Rightarrow\)
对每个 \(X \in \mathcal{C}\)，逗号范畴 \(X \downarrow G\)
有初始对象，该初始对象 \((S_X, \eta_X: X \to G(S_X))\) 定义左伴随
\(F(X) = S_X\)。构造核心：解集条件保证 \(X \downarrow G\)
中有一族''生成''对象 \(\{(M_i, f_i: X \to G(M_i))\}_{i \in I}\)；由于
\(G\) 保极限，这些 \(M_i\) 的积的某个子对象给出初始对象。必要性：左伴随
\(F \dashv G\) 自动保证 \(G\) 保极限（定理 96.1），且对每个
\(X\)，\((F(X), \eta_X)\) 是 \(X \downarrow G\)
的初始对象，从而解集条件取单点集 \(\{F(X)\}\) 即满足。∎

\textbf{定理 96.3}（特殊伴随函子定理）：设 \(\mathcal{C}\)
是完备、良幂（well-powered）且有余生成子（cogenerator）的范畴，\(\mathcal{D}\)
是局部小范畴。若 \(G: \mathcal{D} \to \mathcal{C}\) 保所有小极限，则
\(G\) 有左伴随。

\textbf{例 96.2}（Stone-Čech 紧化）：遗忘函子
\(U: \mathbf{CHaus} \to \mathbf{Top}\)（紧 Hausdorff
空间到拓扑空间）有左伴随 \(\beta\)（Stone-Čech
紧化）。这是特殊伴随函子定理的应用。

\subsubsection{96.3 伴随与等价}\label{ux4f34ux968fux4e0eux7b49ux4ef7}

\textbf{定义 96.3}（范畴等价）：函子 \(F: \mathcal{C} \to \mathcal{D}\)
称为\textbf{范畴等价}，如果存在函子 \(G: \mathcal{D} \to \mathcal{C}\)
使得 \(G \circ F \cong \operatorname{id}_{\mathcal{C}}\) 和
\(F \circ G \cong \operatorname{id}_{\mathcal{D}}\)（自然同构）。

\textbf{命题 96.3}：若 \(F \dashv G\) 且单位和余单位都是自然同构，则
\(F\) 和 \(G\) 是范畴等价（此时 \(F \dashv G\) 且 \(G \dashv F\)）。

\textbf{定理 96.4}：函子 \(F: \mathcal{C} \to \mathcal{D}\)
是范畴等价当且仅当 \(F\) 是完全忠实的且本质满的（essentially
surjective：每个 \(Y \in \mathcal{D}\) 同构于某个 \(F(X)\)）。

\emph{证明}：（\(\Rightarrow\)）若 \(F\) 等价，存在 \(G\) 使
\(G \circ F \cong \operatorname{id}_{\mathcal{C}}\)。自然同构给出每个
\(X\) 的同构 \(X \cong G(F(X))\)，故 \(F\) 完全忠实。本质满：对
\(Y \in \mathcal{D}\)，由
\(F \circ G \cong \operatorname{id}_{\mathcal{D}}\) 知
\(Y \cong F(G(Y))\)，即 \(Y\) 同构于 \(F\)
的像。（\(\Leftarrow\)）构造拟逆函子 \(G\)：对每个
\(Y \in \mathcal{D}\)，由本质满择一 \(X_Y \in \mathcal{C}\) 及同构
\(\varepsilon_Y: F(X_Y) \xrightarrow{\cong} Y\)。对态射
\(h: Y \to Z\)，由 \(F\) 全忠实，存在唯一 \(g: X_Y \to X_Z\) 使
\(F(g) = \varepsilon_Z^{-1} \circ h \circ \varepsilon_Y\)。定义
\(G(Y) = X_Y\)，\(G(h) = g\)。验证 \(G\) 是函子且 \(\varepsilon\) 给出
\(F \circ G \cong \operatorname{id}_{\mathcal{D}}\)；再利用 \(F\)
忠实构造 \(\eta: \operatorname{id}_{\mathcal{C}} \cong G \circ F\)。∎

\bigskip\hrule\bigskip

\subsection{第97章：单子与代数}\label{ux7b2c97ux7ae0ux5355ux5b50ux4e0eux4ee3ux6570}

单子（monad）是伴随函子的代数抽象，捕获了''代数结构''的普遍概念。它在函数式编程和形式语义学中也有重要应用。

\subsubsection{97.1 单子的定义}\label{ux5355ux5b50ux7684ux5b9aux4e49}

\textbf{定义 97.1}（单子）：范畴 \(\mathcal{C}\)
上的\textbf{单子}是一个三元组 \((T, \eta, \mu)\)，其中： -
\(T: \mathcal{C} \to \mathcal{C}\) 是函子 -
\(\eta: \operatorname{id}_{\mathcal{C}} \Rightarrow T\)（单位） -
\(\mu: T^2 \Rightarrow T\)（乘法）

满足结合律：\(\mu \circ T\mu = \mu \circ \mu T\)（作为
\(T^3 \Rightarrow T\)
的自然变换），和单位律：\(\mu \circ \eta T = \operatorname{id}_T = \mu \circ T\eta\)。

\textbf{定理 97.1}（伴随给出单子）：每个伴随
\(F \dashv G\)（\(F: \mathcal{C} \to \mathcal{D}\)，\(G: \mathcal{D} \to \mathcal{C}\)）给出
\(\mathcal{C}\) 上的一个单子
\((T = G \circ F, \eta, \mu = G \varepsilon F)\)，其中 \(\eta\)
是伴随的单位，\(\varepsilon\) 是伴随的余单位。

\emph{证明}：验证单子的三个公理。\(\mu\) 的自然性由 \(\varepsilon\)
的自然性保证。结合律由 \(\varepsilon\) 的三角恒等式推出。∎

\textbf{例 97.1}（来自伴随的单子）： - 自由群函子 \(F \dashv U\) 给出
\textbf{Set} 上的单子 \(UF\)（自由群单子） - 自由 \(R\)-模函子
\(\dashv\) 遗忘函子给出向量空间单子 - 幂集函子
\(\mathcal{P}: \mathbf{Set} \to \mathbf{Set}\)（\(\eta_X(x) = \{x\}\)，\(\mu_X\)
为并集）是单子 - \textbf{Set} 上的 Maybe 单子：\(T(X) = X \sqcup \{*\}\)

\subsubsection{97.2 Eilenberg-Moore 代数与 Kleisli
范畴}\label{eilenberg-moore-ux4ee3ux6570ux4e0e-kleisli-ux8303ux7574}

\textbf{定义 97.2}（Eilenberg-Moore 代数）：单子 \((T, \eta, \mu)\)
的一个\textbf{代数}（或 Eilenberg-Moore 代数）是一个对 \((A, a)\)，其中
\(A \in \mathcal{C}\)，\(a: T(A) \to A\) 是 \(\mathcal{C}\)
中的态射，满足 \(a \circ \eta_A = \operatorname{id}_A\) 和
\(a \circ \mu_A = a \circ T(a)\)。

代数的态射 \(f: (A, a) \to (B, b)\) 是 \(\mathcal{C}\) 中的态射
\(f: A \to B\) 满足 \(b \circ T(f) = f \circ a\)。代数及其态射构成
\textbf{Eilenberg-Moore 范畴} \(\mathcal{C}^T\)。

\textbf{定理 97.2}（Eilenberg-Moore 分解）：每个伴随 \(F \dashv G\)
生成的单子 \(T\)，函子 \(G\) 通过 \(\mathcal{C}^T\) 分解：存在比较函子
\(K: \mathcal{D} \to \mathcal{C}^T\) 使得 \(U^T \circ K = G\) 和
\(K \circ F = F^T\)（其中 \(U^T \dashv F^T\) 是 \(\mathcal{C}^T\)
的自由-遗忘伴随）。

\textbf{定义 97.3}（Kleisli 范畴）：单子 \(T\) 的 \textbf{Kleisli 范畴}
\(\mathcal{C}_T\) 的对象与 \(\mathcal{C}\) 相同，但态射
\(\mathcal{C}_T(X, Y) = \mathcal{C}(X, T(Y))\)。合成由 \(T\)
的乘法定义。Kleisli 范畴是 Eilenberg-Moore 范畴的''自由代数''满子范畴。

\textbf{命题 97.4}：\(\mathcal{C}^T\) 是 \(T\)
的代数的''最大''范畴（终对象），\(\mathcal{C}_T\) 是 \(T\)
的代数的''最小''范畴（初对象），在适当的范畴意义下。

\emph{证明概要（Kleisli 范畴的构造）}：给定单子 \((T, \eta, \mu)\)，定义
Kleisli 范畴 \(\mathcal{C}_T\)：对象
\(\operatorname{Ob}(\mathcal{C}_T) = \operatorname{Ob}(\mathcal{C})\)；态射
\(\mathcal{C}_T(X, Y) = \mathcal{C}(X, T(Y))\)；恒等态射
\(\operatorname{id}_X^{\mathcal{C}_T} = \eta_X: X \to T(X)\)；合成
\(g \circ_{\mathcal{C}_T} f = \mu_Z \circ T(g) \circ f\)（其中
\(f: X \to T(Y)\)，\(g: Y \to T(Z)\)）。验证结合律依赖单子乘法 \(\mu\)
的结合律：\((h \circ g) \circ f = \mu_W \circ T(\mu_W \circ T(h) \circ g) \circ f = \mu_W \circ \mu_{T(W)} \circ T^2(h) \circ T(g) \circ f = \mu_W \circ T(h) \circ \mu_Z \circ T(g) \circ f = h \circ (g \circ f)\)。存在自由函子
\(F_T: \mathcal{C} \to \mathcal{C}_T\)（\(F_T(X)=X\)，\(F_T(f)=\eta_Y \circ f\)）和遗忘函子
\(U_T: \mathcal{C}_T \to \mathcal{C}\)（\(U_T(X)=T(X)\)，\(U_T(f)=\mu_Y \circ T(f)\)），满足
\(F_T \dashv U_T\) 且生成的单子恰为 \(T\)。\(\mathcal{C}_T\) 等价于
\(\mathcal{C}^T\) 中由自由代数 \((\mu_X: T^2(X) \to T(X))\)
生成的满子范畴。∎

\subsubsection{97.3
单子与代数理论}\label{ux5355ux5b50ux4e0eux4ee3ux6570ux7406ux8bba}

\textbf{定义 97.4}（代数理论）：单子 \((T, \eta, \mu)\) 在 \textbf{Set}
上的代数 \((A, a)\) 可以看作带有满足某些等式的运算的集合。例如： -
群单子：\(T(X)\) 是 \(X\) 生成的自由群，代数 \((A, a)\)
是群（\(a: T(A) \to A\) 是群的乘法、逆和单位元） -
环单子：类似地，代数环

\textbf{定理 97.5}（Beck 单子性定理）：设
\(G: \mathcal{D} \to \mathcal{C}\) 有左伴随 \(F\)。则比较函子
\(K: \mathcal{D} \to \mathcal{C}^T\) 是范畴等价当且仅当 \(G\) 反映同构且
\(\mathcal{D}\) 有余等化子，且 \(G\) 保这些余等化子。

Beck 单子性定理提供了判断一个函子是否为''遗忘函子''的精确标准。

\bigskip\hrule\bigskip

\subsection{第98章：Abel
范畴与导出范畴初步}\label{ux7b2c98ux7ae0abel-ux8303ux7574ux4e0eux5bfcux51faux8303ux7574ux521dux6b65}

Abel 范畴是能进行同调代数的范畴框架。导出范畴是 Abel
范畴的局部化，使得拟同构成为真正的同构，是现代同调代数的核心语言。

\subsubsection{98.1 Abel 范畴}\label{abel-ux8303ux7574}

\textbf{定义 98.1}（加性范畴）：范畴 \(\mathcal{A}\)
称为\textbf{加性范畴}（或预加性范畴），如果每个 \(\mathcal{A}(X, Y)\) 是
Abel 群，且合成是双线性的。此外，\(\mathcal{A}\)
有零对象（既是始对象又是终对象）和有限双积（既是积又是余积）。

\textbf{定义 98.2}（Abel 范畴）：加性范畴 \(\mathcal{A}\) 称为
\textbf{Abel 范畴}，如果 1. 每个态射有核和余核 2.
每个单态射是其余核的核，每个满态射是其核的余核 3.
（等价地）每个态射有典范分解 \(f = m \circ e\)，其中 \(e\)
是余核的满态射，\(m\) 是核的单态射

\textbf{例 98.1}： - \textbf{Ab}（Abel 群范畴）是 Abel 范畴 -
\textbf{R-Mod}（\(R\)-模范畴）是 Abel 范畴 - \textbf{Sh}（层范畴）是
Abel 范畴 - 有限维向量空间上的线性映射范畴是 Abel 范畴 - 拓扑 Abel
群范畴不是 Abel 范畴（因为满态射的余核不一定是单态射的核）

\textbf{定义 98.3}（正合列）：Abel 范畴中的序列
\(\cdots \to A \xrightarrow{f} B \xrightarrow{g} C \to \cdots\)
称为\textbf{正合的}，如果 \(\ker g = \operatorname{im} f\)（其中
\(\operatorname{im} f = \ker(\operatorname{coker} f)\)）。

\textbf{定理 98.1}（Freyd-Mitchell 嵌入定理）：每个小 Abel 范畴
\(\mathcal{A}\) 可以完全忠实地嵌入到某个环 \(R\) 的模范畴 \textbf{R-Mod}
中，且嵌入是正合的。

\emph{说明}：Freyd-Mitchell 嵌入定理意味着在 Abel
范畴中做同调代数时，可以''假装''在模范畴中工作，使用元素论证（diagram
chasing）。

\emph{证明思路}：将小 Abel 范畴 \(\mathcal{A}\) 视为加性范畴，构造函子
\(H: \mathcal{A}^{\text{op}} \to \mathbf{Ab}\)
的适当满子范畴作为''层''，再取归纳极限构造环 \(R\)，使 \(\mathcal{A}\)
等价于 \(R\)-Mod 中紧生成对象的满子范畴。核心技术是利用 \(\mathcal{A}\)
中所有可表函子构成的谱（spectrum）来''编码''模范畴结构，最终得到完全忠实且正合的嵌入
\(\mathcal{A} \hookrightarrow R\text{-Mod}\)。∎

\subsubsection{98.2 复形与同调}\label{ux590dux5f62ux4e0eux540cux8c03}

\textbf{定义 98.4}（链复形）：Abel 范畴 \(\mathcal{A}\)
中的\textbf{链复形}是一列对象 \(A_n\) 和态射 \(d_n: A_n \to A_{n-1}\)
满足 \(d_{n-1} \circ d_n = 0\)。\textbf{上链复形}中态射方向为
\(d^n: A^n \to A^{n+1}\)。

\textbf{定义 98.5}（同调）：链复形 \(A_\bullet\) 的 \(n\) 次同调为
\(H_n(A_\bullet) = \ker d_n / \operatorname{im} d_{n+1}\)。

\textbf{定义 98.6}（拟同构）：链复形态射 \(f: A_\bullet \to B_\bullet\)
称为\textbf{拟同构}，如果对所有 \(n\)，它诱导的同调同构
\(H_n(A_\bullet) \cong H_n(B_\bullet)\)。

\subsubsection{98.3 导出范畴}\label{ux5bfcux51faux8303ux7574}

\textbf{定义 98.7}（同伦范畴）：Abel 范畴 \(\mathcal{A}\) 的链复形范畴
\(\mathbf{Ch}(\mathcal{A})\) 的\textbf{同伦范畴} \(K(\mathcal{A})\)
的对象与 \(\mathbf{Ch}(\mathcal{A})\) 相同，态射是链映射的同伦类。

\textbf{定义 98.8}（导出范畴）：Abel 范畴 \(\mathcal{A}\)
的\textbf{导出范畴} \(D(\mathcal{A})\) 是 \(K(\mathcal{A})\)
关于拟同构的局部化（Verdier 局部化）。即
\(D(\mathcal{A}) = K(\mathcal{A})[\mathcal{Q}^{-1}]\)，其中
\(\mathcal{Q}\) 是所有拟同构的类。

在 \(D(\mathcal{A})\) 中，拟同构变成了真正的同构。这是导出函子理论（V12
Ch 70）的范畴论基础。

\textbf{定理 98.2}（导出范畴是三角范畴）：导出范畴 \(D(\mathcal{A})\)
具有\textbf{三角范畴}结构：存在平移函子
\([1]: D(\mathcal{A}) \to D(\mathcal{A})\)（\(A[1]_n = A_{n-1}\)）和一类\textbf{好三角}（distinguished
triangles）\(A \to B \to C \to A[1]\)，满足 Verdier 公理。

\textbf{定义 98.9}（导出函子）：设 \(F: \mathcal{A} \to \mathcal{B}\)
是左正合函子。\(F\) 的\textbf{全导出函子}
\(\mathbf{R}F: D^+(\mathcal{A}) \to D^+(\mathcal{B})\) 定义为
\(\mathbf{R}F(A_\bullet) = F(I_\bullet)\)，其中
\(A_\bullet \to I_\bullet\) 是内射分解（或
\(K\)-内射分解）。\(R^i F(A) = H^i(\mathbf{R}F(A))\) 是通常的导出函子。

\textbf{定理 98.3}（导出函子的存在性）：若 \(\mathcal{A}\)
有足够多内射对象，则每个左正合函子 \(F: \mathcal{A} \to \mathcal{B}\)
有全导出函子 \(\mathbf{R}F: D^+(\mathcal{A}) \to D^+(\mathcal{B})\)。

\textbf{例 98.2}（导出函子的范畴论解释）： -
\(\operatorname{Ext}^n(A, B) = \operatorname{Hom}_{D(\mathcal{A})}(A, B[n])\)（导出范畴中的
Hom） - 层上同调
\(H^n(X, \mathcal{F}) = \operatorname{Hom}_{D(\mathbf{Sh}(X))}(\mathbb{Z}_X, \mathcal{F}[n])\)
- 群上同调
\(H^n(G, M) = \operatorname{Ext}^n_{\mathbb{Z}G}(\mathbb{Z}, M)\)

\bigskip\hrule\bigskip

\emph{本卷建立了范畴论的核心框架：范畴、函子和自然变换提供了描述数学结构的统一语言；Yoneda
引理是范畴论最基本的定理，体现了''对象由关系决定''的哲学；极限和余极限统一了数学中的各种构造；伴随函子是范畴论最核心的概念，自由-遗忘、张量-Hom
等伴随对在数学中无处不在；单子理论将代数结构抽象为范畴论概念；Abel
范畴和导出范畴为同调代数（V12 Ch
70）提供了现代框架。范畴论是连接代数、几何、拓扑和分析的统一语言。}

\newpage

\section{卷十九：数理逻辑}\label{ux5377ux5341ux4e5dux6570ux7406ux903bux8f91}

\begin{quote}
数理逻辑研究数学推理本身的形式化。本卷从一阶逻辑的语法和语义出发，证明
Gödel 完备性定理；然后建立模型论、递归论和可计算性理论的基础；最后证明
Gödel 不完全性定理------20
世纪数学最深刻的成果之一，并引入集合论的进阶内容（序数、基数、连续统假设）。
\end{quote}

\bigskip\hrule\bigskip

\subsection{第99章：一阶逻辑完备性}\label{ux7b2c99ux7ae0ux4e00ux9636ux903bux8f91ux5b8cux5907ux6027}

一阶逻辑是数学推理的标准形式化框架。本章从语法（形式语言、证明系统）和语义（结构、真值）两方面建立一阶逻辑，并证明
Gödel 完备性定理：语义真蕴含语法可证。

\subsubsection{99.1 一阶语言}\label{ux4e00ux9636ux8bedux8a00}

\textbf{定义 99.1}（一阶语言）：一个\textbf{一阶语言} \(\mathcal{L}\)
由以下符号组成： - \textbf{逻辑符号}：变元 \(x_1, x_2, \ldots\)；连接词
\(\neg, \land, \lor, \to\)；量词 \(\forall, \exists\)；等号
\(=\)（通常包含） - \textbf{非逻辑符号}：常数符号 \(\{c_i\}\)；函数符号
\(\{f_j\}\)（带元数）；关系符号 \(\{R_k\}\)（带元数）

\textbf{定义 99.2}（项与公式）：\(\mathcal{L}\)
的\textbf{项}归纳定义：变元和常数是项；若 \(f\) 是 \(n\)
元函数符号，\(t_1, \ldots, t_n\) 是项，则 \(f(t_1, \ldots, t_n)\) 是项。

\textbf{原子公式}为 \(R(t_1, \ldots, t_n)\)（\(R\) 是 \(n\)
元关系符号）或
\(t_1 = t_2\)。\textbf{公式}由原子公式通过连接词和量词递归构造。

\textbf{定义
99.3}（句子与自由变元）：不含自由变元的公式称为\textbf{句子}。一个\textbf{理论}
\(T\) 是一组 \(\mathcal{L}\)-句子的集合。

\subsubsection{99.2 语义}\label{ux8bedux4e49}

\textbf{定义 99.4}（\(\mathcal{L}\)-结构）：\(\mathcal{L}\)
的一个\textbf{结构}（或模型）\(\mathfrak{A}\) 由以下组成： -
一个非空集合 \(A\)（论域） - 对每个常数符号 \(c\)，\(\mathfrak{A}\)
中一个元素 \(c^{\mathfrak{A}} \in A\) - 对每个 \(n\) 元函数符号
\(f\)，一个函数 \(f^{\mathfrak{A}}: A^n \to A\) - 对每个 \(n\)
元关系符号 \(R\)，一个关系 \(R^{\mathfrak{A}} \subseteq A^n\)

\textbf{定义 99.5}（真值，Tarski 真定义）：对 \(\mathcal{L}\)-结构
\(\mathfrak{A}\) 和句子
\(\varphi\)，\(\mathfrak{A} \models \varphi\)（\(\varphi\) 在
\(\mathfrak{A}\) 中为真）由 Tarski 真定义递归定义： -
\(\mathfrak{A} \models R(t_1, \ldots, t_n)\) 当且仅当
\((t_1^{\mathfrak{A}}, \ldots, t_n^{\mathfrak{A}}) \in R^{\mathfrak{A}}\)
- \(\mathfrak{A} \models \neg \varphi\) 当且仅当
\(\mathfrak{A} \not\models \varphi\) -
\(\mathfrak{A} \models \varphi \land \psi\) 当且仅当
\(\mathfrak{A} \models \varphi\) 且 \(\mathfrak{A} \models \psi\) -
\(\mathfrak{A} \models \forall x \varphi(x)\) 当且仅当对每个
\(a \in A\)，\(\mathfrak{A} \models \varphi(a)\)

\textbf{定义 99.6}（逻辑蕴含）：\(T \models \varphi\) 表示对每个
\(\mathcal{L}\)-结构 \(\mathfrak{A}\)，若 \(\mathfrak{A} \models T\)（即
\(\mathfrak{A}\) 满足 \(T\) 中所有句子），则
\(\mathfrak{A} \models \varphi\)。

\subsubsection{99.3 证明系统}\label{ux8bc1ux660eux7cfbux7edf}

\textbf{定义 99.7}（Hilbert 式证明系统）：一阶逻辑的公理包括： 1.
命题逻辑公理（重言式） 2. \(\forall x \varphi(x) \to \varphi(t)\)（\(t\)
对 \(\varphi\) 中 \(x\) 可代入） 3.
\(\forall x (\varphi \to \psi) \to (\forall x \varphi \to \forall x \psi)\)
4. \(\varphi \to \forall x \varphi\)（\(x\) 不在 \(\varphi\) 中自由） 5.
等号公理：\(x = x\)，\(x = y \to y = x\)，\((x = y \land y = z) \to x = z\)，以及函数和关系的代换公理

推理规则：\textbf{Modus Ponens}（从 \(\varphi\) 和 \(\varphi \to \psi\)
推出 \(\psi\)）和 \textbf{全称概括}（Generalization：从 \(\varphi\) 推出
\(\forall x \varphi\)，要求 \(x\) 不在前提中自由出现）。

\textbf{定义 99.8}（可证性）：\(T \vdash \varphi\) 表示存在从 \(T\)
中公理出发、使用推理规则得到的 \(\varphi\) 的有限证明序列。

\subsubsection{99.4 Gödel
完备性定理}\label{guxf6del-ux5b8cux5907ux6027ux5b9aux7406}

\textbf{定理 99.1}（可靠性定理）：若 \(T \vdash \varphi\)，则
\(T \models \varphi\)。

\emph{证明}：对证明长度归纳。公理在任意结构中为真，推理规则保持真值。∎

\textbf{定理 99.2}（Gödel 完备性定理，1929）：若
\(T \models \varphi\)，则
\(T \vdash \varphi\)。等价地，每个一致的理论（\(T \nvdash \bot\)）有模型。

\emph{证明思路}（Henkin 构造，1949）： 1. 将语言 \(\mathcal{L}\) 扩展为
\(\mathcal{L}^*\)（添加 Henkin 常元：对每个
\(\exists x \varphi(x)\)，添加常元 \(c_\varphi\) 和公理
\(\exists x \varphi(x) \to \varphi(c_\varphi)\)）。 2. 将 \(T\)
扩展为极大一致理论 \(T^*\)（Lindenbaum 引理，使用 Zorn 引理）。 3. 从
\(T^*\) 构造 Henkin 模型：论域为 \(\mathcal{L}^*\)
的闭项等价类，函数和关系由 \(T^*\) 中的等词决定。 4. 验证此模型满足
\(T^*\)（从而满足 \(T\)）：对公式复杂度归纳。∎

\textbf{推论 99.3}（紧致性定理）：若 \(T\) 的每个有限子集有模型，则
\(T\) 有模型。

\emph{证明}：若 \(T\)
无模型，由完备性定理，\(T \vdash \bot\)。证明是有限的，仅用到 \(T\)
的有限子集 \(T_0\)。故 \(T_0 \vdash \bot\)，由可靠性，\(T_0\)
无模型，矛盾。∎

\textbf{推论 99.4}（Löwenheim-Skolem 定理，向下）：若可数语言
\(\mathcal{L}\) 的理论 \(T\) 有无限模型，则 \(T\) 有任意无限基数
\(\kappa \geq \aleph_0\) 的模型（向上 Löwenheim-Skolem：若 \(T\)
有无限模型，则 \(T\) 有任意基数 \(\kappa \geq |\mathcal{L}|\) 的模型）。

\emph{证明}（向上）：添加 \(\kappa\) 个新常元
\(\{c_\alpha\}_{\alpha < \kappa}\) 和公理
\(c_\alpha \neq c_\beta\)（\(\alpha \neq \beta\)）。每个有限子集有模型（原无限模型可解释有限个新常元为不同元素）。由紧致性，整个理论有模型，其基数
\(\geq \kappa\)。∎

\bigskip\hrule\bigskip

\subsection{第100章：模型论基础}\label{ux7b2c100ux7ae0ux6a21ux578bux8bbaux57faux7840}

模型论研究数学结构与其一阶理论之间的关系。本章引入初等等价、初等子模型和模型完备性等基本概念。

\subsubsection{100.1
初等等价与初等扩张}\label{ux521dux7b49ux7b49ux4ef7ux4e0eux521dux7b49ux6269ux5f20}

\textbf{定义 100.1}（初等等价）：两个 \(\mathcal{L}\)-结构
\(\mathfrak{A}\) 和 \(\mathfrak{B}\) 称为\textbf{初等等价}的（记作
\(\mathfrak{A} \equiv \mathfrak{B}\)），如果它们满足相同的
\(\mathcal{L}\)-句子。

\textbf{定义
100.2}（初等子模型）：\(\mathfrak{A} \preceq \mathfrak{B}\)（\(\mathfrak{A}\)
是 \(\mathfrak{B}\) 的\textbf{初等子模型}），如果
\(\mathfrak{A} \subseteq \mathfrak{B}\) 且对每个 \(\mathcal{L}\)-公式
\(\varphi(x_1, \ldots, x_n)\) 和 \(a_1, \ldots, a_n \in A\)，

\[\mathfrak{A} \models \varphi(a_1, \ldots, a_n) \iff \mathfrak{B} \models \varphi(a_1, \ldots, a_n)\]

\textbf{定理 100.1}（Tarski-Vaught
判别法）：\(\mathfrak{A} \subseteq \mathfrak{B}\)
是初等子模型当且仅当对每个 \(\mathcal{L}\)-公式
\(\varphi(x, y_1, \ldots, y_n)\) 和 \(a_1, \ldots, a_n \in A\)，若
\(\mathfrak{B} \models \exists x \varphi(x, a_1, \ldots, a_n)\)，则存在
\(a \in A\) 使得 \(\mathfrak{B} \models \varphi(a, a_1, \ldots, a_n)\)。

\subsubsection{100.2 Löwenheim-Skolem
定理}\label{luxf6wenheim-skolem-ux5b9aux7406}

\textbf{定理 100.2}（向下 Löwenheim-Skolem 定理）：设 \(\mathfrak{A}\)
是 \(\mathcal{L}\)-结构，\(X \subseteq A\) 是子集。则存在
\(\mathfrak{A}\) 的初等子模型 \(\mathfrak{B} \preceq \mathfrak{A}\) 包含
\(X\)，且 \(|B| \leq \max(|X|, |\mathcal{L}|, \aleph_0)\)。

\emph{证明}：构造地封闭 \(X\) 在 Skolem 函数下。对每个
\(\mathcal{L}\)-公式 \(\exists y \varphi(y, x_1, \ldots, x_n)\)，选择
Skolem 函数 \(f_\varphi\) 满足
\(\mathfrak{A} \models \exists y \varphi(y, a_1, \ldots, a_n) \to \varphi(f_\varphi(a_1, \ldots, a_n), a_1, \ldots, a_n)\)。\(\mathcal{L}\)
的 Skolem 扩张为 \(\mathcal{L}^S\)。取 \(X\) 在 Skolem 函数下的闭包
\(B\)，则 \(|B| \leq \max(|X|, |\mathcal{L}|, \aleph_0)\)，且
\(\mathfrak{B} \preceq \mathfrak{A}\)。∎

\textbf{推论 100.3}（Skolem 悖论）：若 ZFC
集合论一致，则它有可数模型（由向下
Löwenheim-Skolem）。但在此可数模型中，可以证明''存在不可数集合''（Cantor
定理在模型内部成立）。这不是矛盾，而是''不可数''在模型内部和外部的含义不同。

\subsubsection{100.3
量词消去与模型完备性}\label{ux91cfux8bcdux6d88ux53bbux4e0eux6a21ux578bux5b8cux5907ux6027}

\textbf{定义 100.3}（量词消去）：理论 \(T\)
称为有\textbf{量词消去}，如果每个 \(\mathcal{L}\)-公式等价于（在 \(T\)
中）一个无量词公式。

\textbf{定理 100.4}（量词消去的判别法）：\(T\) 有量词消去当且仅当对任意
\(\mathfrak{A}, \mathfrak{B} \models T\)，\(\mathfrak{C} \subseteq \mathfrak{A}, \mathfrak{B}\)
是公共子结构，以及无量词公式 \(\varphi(x, \bar{y})\) 和
\(\bar{c} \in C\)，若
\(\mathfrak{A} \models \exists x \varphi(x, \bar{c})\)，则
\(\mathfrak{B} \models \exists x \varphi(x, \bar{c})\)。

\textbf{例 100.1}（有量词消去的理论）： - 代数闭域理论
ACF（量词消去，Tarski） - 实闭域理论 RCF（量词消去，Tarski-Seidenberg）
- 无端点稠密线序理论 DLO（量词消去）

\textbf{定义 100.4}（模型完备性）：理论 \(T\)
称为\textbf{模型完备}的，如果对任意
\(\mathfrak{A}, \mathfrak{B} \models T\)，\(\mathfrak{A} \subseteq \mathfrak{B}\)
推出 \(\mathfrak{A} \preceq \mathfrak{B}\)。

\textbf{定理 100.5}（Robinson 判别法）：\(T\) 是模型完备的当且仅当对任意
\(\mathfrak{A}, \mathfrak{B} \models T\) 和
\(\mathfrak{A} \subseteq \mathfrak{B}\)，每个存在公式（\(\exists \bar{x} \varphi(\bar{x}, \bar{y})\)，\(\varphi\)
无量词）在 \(\mathfrak{B}\) 中满足的元素也在 \(\mathfrak{A}\) 中满足。

\bigskip\hrule\bigskip

\subsection{第101章：递归论与可计算性}\label{ux7b2c101ux7ae0ux9012ux5f52ux8bbaux4e0eux53efux8ba1ux7b97ux6027}

递归论（可计算性理论）研究哪些函数是可计算的，以及不可计算问题的层级。本章从
Turing 机出发，定义可计算函数，并证明停机问题的不可判定性。

\subsubsection{101.1 Turing 机}\label{turing-ux673a}

\textbf{定义 101.1}（Turing 机）：一个 \textbf{Turing 机} \(M\)
由以下组成： - 有限状态集 \(Q\)（包含初始状态 \(q_0\) 和停机状态
\(q_{\text{halt}}\)） - 有限带符号集 \(\Gamma\)（包含空白符号
\(\sqcup\)） - 输入符号集
\(\Sigma \subseteq \Gamma \setminus \{\sqcup\}\) - 转移函数
\(\delta: Q \times \Gamma \to Q \times \Gamma \times \{L, R\}\)

Turing
机有一条无限长的带（双向无限），初始时输入写在带上，读写头在输入最左端，状态为
\(q_0\)。每一步根据当前状态和读写头下的符号，按 \(\delta\)
写入新符号、移动读写头（左/右）并改变状态。

\textbf{定义 101.2}（Turing 可计算）：部分函数
\(f: \mathbb{N}^k \to \mathbb{N}\)（或
\(f: \Sigma^* \to \Sigma^*\)）称为 \textbf{Turing 可计算}的，如果存在
Turing 机 \(M\)，对每个输入 \(\bar{x}\)，\(M\) 在 \(f(\bar{x})\)
有定义时停机且输出 \(f(\bar{x})\)，否则 \(M\) 永不停机。

\textbf{Church-Turing 论题}：直观上''可计算''的函数恰好是 Turing
可计算的函数。所有其他计算模型（\(\lambda\)
演算、一般递归函数、寄存器机等）都定义了相同的可计算函数类。

\subsubsection{101.2
递归函数与递归可枚举集}\label{ux9012ux5f52ux51fdux6570ux4e0eux9012ux5f52ux53efux679aux4e3eux96c6}

\textbf{定义
101.3}（原始递归函数）：\textbf{原始递归函数}类是最小的包含以下函数且在复合和原始递归下封闭的函数类：
- 零函数：\(Z(n) = 0\) - 后继函数：\(S(n) = n + 1\) -
投影函数：\(P_i^n(x_1, \ldots, x_n) = x_i\) -
原始递归：\(f(0, \bar{y}) = g(\bar{y})\)，\(f(n+1, \bar{y}) = h(n, f(n, \bar{y}), \bar{y})\)

\textbf{定义
101.4}（部分递归函数）：\textbf{部分递归函数}是原始递归函数加上极小化算子
\(\mu\)：\(f(\bar{x}) = \mu y [g(\bar{x}, y) = 0]\)（最小的 \(y\) 使得
\(g(\bar{x}, y) = 0\)，且所有
\(g(\bar{x}, z)\)（\(z < y\)）有定义且非零）。

\textbf{定理 101.1}：Turing 可计算函数类 = 部分递归函数类。

\textbf{定义 101.5}（递归集与递归可枚举集）：\(A \subseteq \mathbb{N}\)
称为 - \textbf{递归集}（可判定集）：其特征函数 \(\chi_A\)
是可计算的（总函数） - \textbf{递归可枚举集}（r.e.
集）：\(A = \emptyset\) 或 \(A\) 是某个可计算总函数的值域

\textbf{命题 101.2}：\(A\) 是递归的当且仅当 \(A\) 和
\(\mathbb{N} \setminus A\) 都是 r.e. 的。存在 r.e. 但非递归的集合。

\subsubsection{101.3
停机问题与不可判定性}\label{ux505cux673aux95eeux9898ux4e0eux4e0dux53efux5224ux5b9aux6027}

\textbf{定义 101.6}（停机问题）：\textbf{停机问题}
\(K = \{\langle M, x \rangle : M \text{ 在输入 } x \text{ 上停机}\}\)，其中
\(\langle M, x \rangle\) 是 \(M\) 和 \(x\) 的某种编码。

\textbf{定理 101.3}（停机问题的不可判定性，Turing 1936）：\(K\)
不是递归集（不可判定）。

\emph{证明}：假设存在 Turing 机 \(H\) 判定 \(K\)。构造 Turing 机
\(D\)：在输入 \(M\)（Turing 机的编码）上，\(D\) 使用 \(H\) 判断 \(M\)
在输入 \(M\) 上是否停机。若停机，则 \(D\) 进入无限循环；若不停机，则
\(D\) 停机。那么 \(D\) 在输入 \(D\) 上停机当且仅当 \(D\) 在输入 \(D\)
上不停机，矛盾。∎

\textbf{推论 101.4}（一阶逻辑的不可判定性，Church
1936）：一阶逻辑的有效性（\(\models \varphi\)？）是不可判定的。即不存在算法能判定任意一阶公式是否为逻辑有效式。

\subsubsection{101.4 算术层次}\label{ux7b97ux672fux5c42ux6b21}

\textbf{定义 101.7}（算术层次）：\(\Sigma_1^0\) 公式形如
\(\exists x_1 \cdots \exists x_n \varphi\)，其中 \(\varphi\)
仅含受囿量词。\(\Pi_1^0\) 公式形如
\(\forall x_1 \cdots \forall x_n \varphi\)。\(\Sigma_{n+1}^0\) 形如
\(\exists \bar{x} \psi\)，\(\psi \in \Pi_n^0\)。\(\Pi_{n+1}^0\) 形如
\(\forall \bar{x} \psi\)，\(\psi \in \Sigma_n^0\)。

\textbf{定理 101.5}：\(A \subseteq \mathbb{N}\) 是 r.e. 当且仅当 \(A\)
是 \(\Sigma_1^0\)-可定义的（在 \(\mathbb{N}\)
的标准模型中）。递归集恰好是 \(\Delta_1^0 = \Sigma_1^0 \cap \Pi_1^0\)
可定义的。

\bigskip\hrule\bigskip

\subsection{第102章：不完全性定理}\label{ux7b2c102ux7ae0ux4e0dux5b8cux5168ux6027ux5b9aux7406}

Gödel
不完全性定理（1931）是数理逻辑和数学基础的里程碑式成果。它表明：任何包含基本算术的一致递归公理系统都是不完全的------存在该系统无法证明也无法否证的命题。

\subsubsection{102.1
算术的形式化}\label{ux7b97ux672fux7684ux5f62ux5f0fux5316}

\textbf{定义 102.1}（Peano 算术，PA）：PA 的语言为
\(\mathcal{L}_{PA} = \{0, S, +, \cdot\}\)（V1 中的 Peano
公理的形式化）。公理包括： 1. \(\forall x (S(x) \neq 0)\) 2.
\(\forall x \forall y (S(x) = S(y) \to x = y)\) 3.
\(\forall x (x + 0 = x)\) 4.
\(\forall x \forall y (x + S(y) = S(x + y))\) 5.
\(\forall x (x \cdot 0 = 0)\) 6.
\(\forall x \forall y (x \cdot S(y) = x \cdot y + x)\) 7.
归纳公理模式：对每个公式
\(\varphi(x, \bar{y})\)，\([\varphi(0, \bar{y}) \land \forall x (\varphi(x, \bar{y}) \to \varphi(S(x), \bar{y}))] \to \forall x \varphi(x, \bar{y})\)

\subsubsection{102.2 Gödel 编码}\label{guxf6del-ux7f16ux7801}

\textbf{定义 102.2}（Gödel 编码）：Gödel 编码（或 Gödel 数）是
\(\mathcal{L}_{PA}\) 的符号、公式和证明序列到 \(\mathbb{N}\)
的单射编码。记 \(\ulcorner \varphi \urcorner\) 为公式 \(\varphi\) 的
Gödel 数。

\textbf{引理 102.1}（可表示性）：以下关系和函数在 PA
中可表示（即存在公式定义它们，PA 能证明其基本性质）： -
\(\operatorname{Prf}_{PA}(x, y)\)：\(x\) 是 \(y\) 的 PA 证明的 Gödel 数
-
\(\operatorname{Prov}_{PA}(y) = \exists x \operatorname{Prf}_{PA}(x, y)\)：\(y\)
是 PA 可证公式的 Gödel 数 - 对角化函数
\(d\)：\(d(\ulcorner \varphi(x) \urcorner) = \ulcorner \varphi(\ulcorner \varphi(x) \urcorner) \urcorner\)

\subsubsection{102.3
第一不完全性定理}\label{ux7b2cux4e00ux4e0dux5b8cux5168ux6027ux5b9aux7406}

\textbf{定理 102.1}（Gödel 第一不完全性定理）：设 \(T\) 是包含 PA
的一致递归公理化理论。则存在 \(\mathcal{L}_{PA}\)-句子 \(G_T\)（Gödel
句子）使得 \(T \nvdash G_T\) 且 \(T \nvdash \neg G_T\)（假设 \(T\) 是
\(\omega\)-一致的或 \(\Sigma_1\)-可靠的）。

\emph{证明}（Gödel 1931）： 1. 构造自指句子：由对角引理，存在句子
\(G_T\) 满足
\(T \vdash G_T \leftrightarrow \neg \operatorname{Prov}_T(\ulcorner G_T \urcorner)\)。\(G_T\)
断言''\(G_T\) 在 \(T\) 中不可证''。 2. 若 \(T \vdash G_T\)，则
\(T \vdash \operatorname{Prov}_T(\ulcorner G_T \urcorner)\)（由
\(\Sigma_1\)-完备性：\(T\) 能证明每个真 \(\Sigma_1\) 句子）。但
\(T \vdash G_T \to \neg \operatorname{Prov}_T(\ulcorner G_T \urcorner)\)，矛盾。
3. 若 \(T \vdash \neg G_T\)，则
\(T \vdash \operatorname{Prov}_T(\ulcorner G_T \urcorner)\)。由
\(\omega\)-一致性（或 \(\Sigma_1\)-可靠性），这蕴含
\(T \vdash G_T\)，矛盾。故 \(T\) 既不能证明也不能否证 \(G_T\)。∎

\textbf{推论 102.2}（不可完备性）：不存在包含 PA
的一致递归公理化完备理论。数学真理不能完全被形式系统捕获。

\textbf{定理 102.3}（Rosser 变体，1936）：Gödel 第一不完全性定理仅需
\(T\) 的一致性（不需要 \(\omega\)-一致性）。Rosser 句子 \(R_T\) 满足
\(R_T \leftrightarrow \forall x (\operatorname{Prf}_T(x, \ulcorner R_T \urcorner) \to \exists y < x \operatorname{Prf}_T(y, \ulcorner \neg R_T \urcorner))\)。

\subsubsection{102.4
第二不完全性定理}\label{ux7b2cux4e8cux4e0dux5b8cux5168ux6027ux5b9aux7406}

\textbf{定义
102.3}（一致性陈述）：\(\operatorname{Con}_T = \neg \operatorname{Prov}_T(\ulcorner 0 = 1 \urcorner)\)（或
\(\neg \operatorname{Prov}_T(\ulcorner \bot \urcorner)\)）。\(\operatorname{Con}_T\)
断言 \(T\) 是一致的。

\textbf{定理 102.4}（Gödel 第二不完全性定理）：设 \(T\) 是包含 PA
的一致递归公理化理论。则 \(T \nvdash \operatorname{Con}_T\)。即 \(T\)
不能证明自身的一致性。

\emph{证明思路}：第一不完全性定理的证明可以在 \(T\) 内部形式化，得到
\(T \vdash \operatorname{Con}_T \to G_T\)。因此若
\(T \vdash \operatorname{Con}_T\)，则
\(T \vdash G_T\)，与第一不完全性定理矛盾。∎

\textbf{推论 102.5}：ZFC 集合论的一致性不能在 ZFC 内部证明。Hilbert
计划（用有限方法证明数学的一致性）被 Gödel 不完全性定理从根本上限制。

\bigskip\hrule\bigskip

\subsection{第103章：集合论进阶}\label{ux7b2c103ux7ae0ux96c6ux5408ux8bbaux8fdbux9636}

本章在 V1（ZFC 集合论基础）的基础上，深入序数、基数的算术和连续统假设。

\subsubsection{103.1
序数与超限归纳}\label{ux5e8fux6570ux4e0eux8d85ux9650ux5f52ux7eb3}

\textbf{定义 103.1}（序数，von Neumann 定义）：一个集合 \(\alpha\)
称为\textbf{序数}，如果 \(\alpha\) 是传递的（即
\(\beta \in \alpha \Rightarrow \beta \subseteq \alpha\)）且 \(\alpha\)
在 \(\in\) 下是良序的。

\textbf{例
103.1}：\(0 = \varnothing\)，\(1 = \{0\}\)，\(2 = \{0, 1\}\)，\(\ldots\)，\(\omega = \{0, 1, 2, \ldots\}\)（第一个无穷序数），\(\omega + 1 = \omega \cup \{\omega\}\)，\(\ldots\)。

\textbf{定理 103.1}（超限归纳法）：若对每个序数
\(\alpha\)，\((\forall \beta < \alpha) P(\beta)\) 推出 \(P(\alpha)\)，则
\(P(\alpha)\) 对所有序数 \(\alpha\) 成立。

\textbf{定理 103.2}（超限递归定义）：对任意给定的类函数
\(G\)，存在唯一的类函数 \(F\) 定义在所有序数上，满足
\(F(\alpha) = G(F \restriction \alpha)\)。

\textbf{定义 103.2}（序数算术）： -
\(\alpha + 0 = \alpha\)，\(\alpha + (\beta + 1) = (\alpha + \beta) + 1\)，\(\alpha + \lambda = \sup_{\beta < \lambda} (\alpha + \beta)\)（极限序数
\(\lambda\)） -
\(\alpha \cdot 0 = 0\)，\(\alpha \cdot (\beta + 1) = \alpha \cdot \beta + \alpha\)，\(\alpha \cdot \lambda = \sup_{\beta < \lambda} (\alpha \cdot \beta)\)
-
\(\alpha^0 = 1\)，\(\alpha^{\beta+1} = \alpha^\beta \cdot \alpha\)，\(\alpha^\lambda = \sup_{\beta < \lambda} \alpha^\beta\)

注意：序数加法和乘法不交换（\(1 + \omega = \omega \neq \omega + 1\)）。

\subsubsection{103.2 基数}\label{ux57faux6570}

\textbf{定义 103.3}（基数）：序数 \(\kappa\)
称为\textbf{基数}（或初始序数），如果对所有
\(\alpha < \kappa\)，\(|\alpha| < |\kappa|\)（即不存在 \(\alpha\) 到
\(\kappa\) 的双射）。

\textbf{定义 103.4}（\(\aleph\)
序列）：\(\aleph_0 = \omega\)。\(\aleph_{\alpha+1} = \min\{\kappa : \kappa > \aleph_\alpha\}\)（\(\aleph_\alpha\)
的后继基数）。\(\aleph_\lambda = \sup_{\alpha < \lambda} \aleph_\alpha\)（极限序数
\(\lambda\)）。

\textbf{定理 103.3}（Cantor 定理）：对任意集合
\(X\)，\(|X| < |\mathcal{P}(X)|\)。特别地，\(\aleph_\alpha < 2^{\aleph_\alpha}\)。

\emph{证明}：\(f: X \to \mathcal{P}(X)\) 不是满射：考虑
\(D = \{x \in X : x \notin f(x)\}\)。若 \(D = f(d)\)，则
\(d \in D \iff d \notin f(d) = D\)，矛盾。∎

\textbf{定义 103.5}（基数算术）：对基数 \(\kappa, \lambda\)： -
\(\kappa + \lambda = |\kappa \sqcup \lambda|\) -
\(\kappa \cdot \lambda = |\kappa \times \lambda|\) -
\(\kappa^\lambda = |\{f: \lambda \to \kappa\}|\)

\textbf{定理 103.4}（基数算术基本定理）：对无穷基数
\(\kappa, \lambda\)： -
\(\kappa + \lambda = \kappa \cdot \lambda = \max(\kappa, \lambda)\)（若
\(\kappa, \lambda > 0\)） - \(\kappa^\kappa = 2^\kappa\)

\textbf{定理 103.5}（König 定理）：对基数
\(\kappa_i < \lambda_i\)（\(i \in I\)），\(\sum_{i \in I} \kappa_i < \prod_{i \in I} \lambda_i\)。

\subsubsection{103.3
选择公理的等价形式}\label{ux9009ux62e9ux516cux7406ux7684ux7b49ux4ef7ux5f62ux5f0f}

\textbf{定理 103.6}（选择公理，AC）：以下陈述在 ZF 中等价： 1.
\textbf{选择公理}：对任意非空集合族 \(\{X_i\}_{i \in I}\)，存在选择函数
\(f: I \to \bigcup_i X_i\) 使得 \(f(i) \in X_i\) 2. \textbf{Zorn
引理}：若非空偏序集中每条链有上界，则该偏序集有极大元 3.
\textbf{良序定理}：每个集合可被良序 4. \textbf{Tychonoff
定理}：紧致空间的任意积是紧致的（V10 Ch 58） 5. 每个向量空间有基（V7 Ch
32） 6. 每个满射有右逆

\textbf{定理 103.7}（基数比较定理，在 AC 下）：对任意两个集合
\(X, Y\)，\(|X| \leq |Y|\) 或 \(|Y| \leq |X|\)。即任意两个基数可比较。

\subsubsection{103.4 连续统假设}\label{ux8fdeux7eedux7edfux5047ux8bbe}

\textbf{定义 103.6}（连续统假设，CH）：\textbf{连续统假设}断言
\(2^{\aleph_0} = \aleph_1\)。即实数集 \(\mathbb{R}\)
的基数是第一个不可数基数。

\textbf{广义连续统假设（GCH）}：对所有序数
\(\alpha\)，\(2^{\aleph_\alpha} = \aleph_{\alpha+1}\)。

\textbf{定理 103.8}（Gödel 1938，Cohen 1963）： 1. 若 ZF 一致，则 ZFC +
GCH 一致（Gödel：CH 在 ZFC 中不可否证，即可构造宇宙 \(L\) 满足 GCH） 2.
若 ZFC 一致，则 ZFC + \(\neg\)CH 一致（Cohen：力迫法，CH 在 ZFC
中不可证明）

CH 独立于 ZFC！这是数理逻辑最著名的独立性结果。Cohen 因此获得 Fields
奖（1966）。

\textbf{定义 103.7}（可构造宇宙）：Gödel 的\textbf{可构造宇宙} \(L\)
通过超限递归定义：\(L_0 = \varnothing\)，\(L_{\alpha+1} = \operatorname{Def}(L_\alpha)\)（\(L_\alpha\)
的可定义子集），\(L_\lambda = \bigcup_{\alpha < \lambda} L_\alpha\)（极限序数）。\(L = \bigcup_{\alpha \in \mathbf{Ord}} L_\alpha\)。

\textbf{定理 103.9}（\(V = L\) 公理）：在 \(L\) 中，选择公理和 GCH
都成立。\(L\) 是最小的包含所有序数的 ZF 内模型。

\textbf{定义 103.8}（力迫法概要）：力迫法（Cohen
1963）通过添加''一般滤子''（generic filter）来扩张可数传递模型
\(M\)，得到 \(M[G]\)，其中某些命题（如
\(\neg\)CH）成立。力迫法是现代集合论最强大的工具之一。

\bigskip\hrule\bigskip

\emph{本卷完成了数理逻辑的核心内容：一阶逻辑的完备性（Gödel
1929）建立了语法与语义的桥梁；紧致性定理和 Löwenheim-Skolem
定理揭示了形式语言的局限性；Turing
机和停机问题为可计算性划定了边界；Gödel 不完全性定理（1931）是 20
世纪数学最深刻的成果之一，它证明了任何足够强的形式系统都不可避免地存在不可判定的命题；序数和基数理论为无穷提供了精确的度量；连续统假设的独立性（Gödel
1938 / Cohen
1963）展示了形式系统的根本局限性。数理逻辑不仅是数学的基础，也是计算机科学（计算理论、形式化验证）的基石。}

\newpage

\section{卷二十：概率论
II}\label{ux5377ux4e8cux5341ux6982ux7387ux8bba-ii}

\begin{quote}
概率论 II 在初等概率论（V6 Ch 29）和实分析（V9
测度论）的基础上，建立概率论的测度论基础。本卷从 Kolmogorov
公理出发，精确定义概率空间、随机变量和期望，然后建立大数定律、中心极限定理、条件期望、鞅论和
Markov 链的严格理论。
\end{quote}

\bigskip\hrule\bigskip

\subsection{第104章：测度论概率基础}\label{ux7b2c104ux7ae0ux6d4bux5ea6ux8bbaux6982ux7387ux57faux7840}

Kolmogorov 在 1933
年将概率论建立在测度论之上，使概率论成为数学分析的一个分支。本章从概率空间出发，用测度论语言严格定义随机变量、分布、期望和独立性。

\subsubsection{104.1 概率空间}\label{ux6982ux7387ux7a7aux95f4}

\textbf{定义 104.1}（概率空间）：一个\textbf{概率空间}
\((\Omega, \mathcal{F}, \mathbb{P})\) 由以下组成： -
\(\Omega\)：样本空间（所有可能结果的集合） - \(\mathcal{F}\)：\(\Omega\)
上的 \(\sigma\)-代数（事件域） -
\(\mathbb{P}: \mathcal{F} \to [0, 1]\)：概率测度，满足
\(\mathbb{P}(\Omega) = 1\) 和可数可加性

\textbf{定义 104.2}（\(\sigma\)-代数）：\(\Omega\) 的子集族
\(\mathcal{F}\) 称为 \(\sigma\)\textbf{-代数}，如果： 1.
\(\Omega \in \mathcal{F}\) 2.
\(A \in \mathcal{F} \Rightarrow A^c \in \mathcal{F}\)（对补集封闭） 3.
\(A_1, A_2, \ldots \in \mathcal{F} \Rightarrow \bigcup_{n=1}^\infty A_n \in \mathcal{F}\)（对可数并封闭）

\textbf{定义 104.3}（Borel \(\sigma\)-代数）：\(\mathbb{R}\) 上的
\textbf{Borel \(\sigma\)-代数} \(\mathcal{B}(\mathbb{R})\)
是包含所有开集的最小 \(\sigma\)-代数。\(\mathcal{B}(\mathbb{R}^n)\)
类似定义。

\textbf{命题 104.1}（概率测度的基本性质）： 1.
\(\mathbb{P}(\varnothing) = 0\) 2. 有限可加性：若 \(A_1, \ldots, A_n\)
互不相交，则
\(\mathbb{P}(\bigcup_{i=1}^n A_i) = \sum_{i=1}^n \mathbb{P}(A_i)\) 3.
单调性：\(A \subseteq B \Rightarrow \mathbb{P}(A) \leq \mathbb{P}(B)\)
4.
连续性：\(A_n \uparrow A\)（单调解）\(\Rightarrow \mathbb{P}(A_n) \uparrow \mathbb{P}(A)\)；\(A_n \downarrow A \Rightarrow \mathbb{P}(A_n) \downarrow \mathbb{P}(A)\)
5.
次可加性：\(\mathbb{P}(\bigcup_{n=1}^\infty A_n) \leq \sum_{n=1}^\infty \mathbb{P}(A_n)\)

\emph{证明}：直接由测度的一般性质（V9 Ch 45）推导。∎

\subsubsection{104.2
随机变量与分布}\label{ux968fux673aux53d8ux91cfux4e0eux5206ux5e03}

\textbf{定义 104.4}（随机变量）：\((\Omega, \mathcal{F}, \mathbb{P})\)
上的\textbf{随机变量} \(X\) 是可测函数
\(X: \Omega \to \mathbb{R}\)（即对每个 Borel 集
\(B \in \mathcal{B}(\mathbb{R})\)，\(X^{-1}(B) \in \mathcal{F}\)）。

\(n\) 维\textbf{随机向量} \(\mathbf{X} = (X_1, \ldots, X_n)\) 是可测函数
\(\mathbf{X}: \Omega \to \mathbb{R}^n\)。

\textbf{定义 104.5}（分布）：随机变量 \(X\)
的\textbf{分布}（或\textbf{律}）\(\mu_X\) 是 \(\mathbb{R}\)
上的概率测度，定义为
\(\mu_X(B) = \mathbb{P}(X \in B) = \mathbb{P}(X^{-1}(B))\)。

\textbf{定义 104.6}（分布函数）：\(X\) 的\textbf{累积分布函数}（CDF）为
\(F_X(x) = \mathbb{P}(X \leq x) = \mu_X((-\infty, x])\)。\(F_X\)
是非降右连续函数，满足
\(\lim_{x \to -\infty} F_X(x) = 0\)，\(\lim_{x \to \infty} F_X(x) = 1\)。

\textbf{定义 104.7}（概率密度函数）：若存在非负可测函数 \(f_X\) 使得
\(\mu_X(B) = \int_B f_X(x) dx\)，则 \(f_X\) 称为 \(X\)
的\textbf{概率密度函数}（PDF）。此时 \(X\) 称为\textbf{绝对连续}的。

\subsubsection{104.3
期望的测度论定义}\label{ux671fux671bux7684ux6d4bux5ea6ux8bbaux5b9aux4e49}

\textbf{定义 104.8}（期望）：随机变量 \(X\)
的\textbf{期望}（或\textbf{均值}）定义为 Lebesgue 积分：

\[\mathbb{E}[X] = \int_\Omega X(\omega) d\mathbb{P}(\omega)\]

若 \(\mathbb{E}[|X|] < \infty\)，则称 \(X\) \textbf{可积}。

\textbf{命题
104.2}（期望的转化公式）：\(\mathbb{E}[X] = \int_{\mathbb{R}} x \, d\mu_X(x)\)。更一般地，对可测函数
\(g\)，\(\mathbb{E}[g(X)] = \int_{\mathbb{R}} g(x) d\mu_X(x)\)。若 \(X\)
有密度 \(f_X\)，则
\(\mathbb{E}[g(X)] = \int_{\mathbb{R}} g(x) f_X(x) dx\)。

\textbf{定义 104.9}（方差与矩）：\(X\) 的\textbf{方差}
\(\operatorname{Var}(X) = \mathbb{E}[(X - \mathbb{E}[X])^2] = \mathbb{E}[X^2] - (\mathbb{E}[X])^2\)。\(k\)
阶矩 \(\mathbb{E}[X^k]\)，\(k\) 阶中心矩
\(\mathbb{E}[(X - \mathbb{E}[X])^k]\)。

\textbf{定理 104.3}（Jensen 不等式）：若 \(\varphi\) 是凸函数且 \(X\)
可积，则 \(\varphi(\mathbb{E}[X]) \leq \mathbb{E}[\varphi(X)]\)（假设
\(\mathbb{E}[\varphi(X)]\) 存在）。

\emph{证明}：凸函数 \(\varphi\)
存在支撑线：\(\varphi(x) \geq \varphi(x_0) + a(x - x_0)\)。取
\(x_0 = \mathbb{E}[X]\)，两边求期望。∎

\textbf{定理 104.4}（Hölder
不等式）：\(\mathbb{E}[|XY|] \leq (\mathbb{E}[|X|^p])^{1/p} (\mathbb{E}[|Y|^q])^{1/q}\)，其中
\(1/p + 1/q = 1\)（\(p, q > 1\)）。

\textbf{推论 104.5}（Cauchy-Schwarz
不等式）：\(\mathbb{E}[|XY|] \leq \sqrt{\mathbb{E}[X^2] \mathbb{E}[Y^2]}\)（\(p = q = 2\)）。

\subsubsection{104.4 独立性}\label{ux72ecux7acbux6027}

\textbf{定义 104.10}（独立事件）：事件
\(A_1, \ldots, A_n \in \mathcal{F}\) 称为\textbf{独立}的，如果对任意子集
\(\{i_1, \ldots, i_k\}\)，

\[\mathbb{P}(A_{i_1} \cap \cdots \cap A_{i_k}) = \mathbb{P}(A_{i_1}) \cdots \mathbb{P}(A_{i_k})\]

\textbf{定义 104.11}（独立 \(\sigma\)-代数）：\(\sigma\)-子代数
\(\mathcal{G}_1, \ldots, \mathcal{G}_n \subseteq \mathcal{F}\)
称为\textbf{独立}的，如果对任意
\(A_i \in \mathcal{G}_i\)，\(A_1, \ldots, A_n\) 独立。

\textbf{定义 104.12}（独立随机变量）：随机变量 \(X_1, \ldots, X_n\)
称为\textbf{独立}的，如果它们生成的 \(\sigma\)-代数
\(\sigma(X_1), \ldots, \sigma(X_n)\)
独立。等价地，\(\mathbb{P}(X_1 \in B_1, \ldots, X_n \in B_n) = \prod_{i=1}^n \mathbb{P}(X_i \in B_i)\)
对所有 Borel 集 \(B_i\) 成立。

\textbf{命题 104.6}：若 \(X_1, \ldots, X_n\) 独立且可积，则
\(\mathbb{E}[X_1 \cdots X_n] = \mathbb{E}[X_1] \cdots \mathbb{E}[X_n]\)。若两两不相关（\(\operatorname{Cov}(X_i, X_j) = 0\)），则
\(\operatorname{Var}(\sum X_i) = \sum \operatorname{Var}(X_i)\)。

\bigskip\hrule\bigskip

\subsection{第105章：大数定律与中心极限定理}\label{ux7b2c105ux7ae0ux5927ux6570ux5b9aux5f8bux4e0eux4e2dux5fc3ux6781ux9650ux5b9aux7406}

大数定律（LLN）和中心极限定理（CLT）是概率论的两大基石。LLN
表明样本均值收敛于期望，CLT 描述了围绕均值的波动行为。

\subsubsection{105.1 大数定律}\label{ux5927ux6570ux5b9aux5f8b}

\textbf{定理 105.1}（弱大数定律，Khinchine）：设 \(X_1, X_2, \ldots\)
是独立同分布（i.i.d.）且 \(\mathbb{E}[|X_1|] < \infty\)。令
\(\mu = \mathbb{E}[X_1]\)。则样本均值
\(\bar{X}_n = \frac{1}{n} \sum_{i=1}^n X_i\) 依概率收敛于 \(\mu\)：

\[\bar{X}_n \xrightarrow{\mathbb{P}} \mu \quad (n \to \infty)\]

\emph{证明}：使用特征函数。\(\varphi_{\bar{X}_n}(t) = \varphi_{X_1}(t/n)^n \to e^{i\mu t}\)（由
\(\varphi_{X_1}(t) = 1 + i\mu t + o(t)\)）。由连续性定理，\(\bar{X}_n \xrightarrow{d} \mu\)（常数），故依概率收敛。∎

\textbf{定理 105.2}（强大数定律，Kolmogorov）：设 \(X_1, X_2, \ldots\)
是 i.i.d. 且 \(\mathbb{E}[|X_1|] < \infty\)。则

\[\bar{X}_n \xrightarrow{\text{a.s.}} \mu \quad \text{几乎必然}\]

即 \(\mathbb{P}(\lim_{n \to \infty} \bar{X}_n = \mu) = 1\)。

\emph{证明思路}（Kolmogorov 截断法）：(1) 不损失一般性，设
\(\mu = 0\)。对 \(X_n\)
做截断：\(Y_n = X_n \mathbf{1}_{\{|X_n| \leq n\}}\)，\(Z_n = X_n - Y_n\)。由
\(\mathbb{E}[|X_1|] < \infty\) 得
\(\sum \mathbb{P}(|X_n| > n) = \sum \mathbb{P}(|X_1| > n) < \infty\)，故
Borel-Cantelli 引理保证几乎必然仅有有限个 \(n\) 满足
\(Z_n \neq 0\)，从而 \(\frac{1}{n}\sum Z_k \to 0\) a.s.。只需证
\(\frac{1}{n}\sum (Y_k - \mathbb{E}[Y_k]) \to 0\) a.s.。(2)
计算截断后的方差：\(\operatorname{Var}(Y_k) \leq \mathbb{E}[Y_k^2] \leq k \mathbb{E}[|X_1| \mathbf{1}_{\{|X_1| \leq k\}}]\)，可得
\(\sum_{k=1}^\infty \frac{\operatorname{Var}(Y_k)}{k^2} < \infty\)。(3)
由 Kolmogorov 不等式和 Kronecker
引理，\(\frac{1}{n}\sum_{k=1}^n (Y_k - \mathbb{E}[Y_k]) \to 0\) a.s.。又
\(\mathbb{E}[Y_k] = \mathbb{E}[X_1 \mathbf{1}_{\{|X_1| \leq k\}}] \to \mathbb{E}[X_1] = 0\)（控制收敛定理），故其
Cesàro 均值也趋于 0。综合得 \(\bar{X}_n \to 0\) a.s.。∎

\textbf{定理 105.3}（Kolmogorov 0-1 律）：设 \(X_1, X_2, \ldots\)
独立。尾 \(\sigma\)-代数
\(\mathcal{T} = \bigcap_{n=1}^\infty \sigma(X_n, X_{n+1}, \ldots)\)
中的每个事件概率为 \(0\) 或 \(1\)。特别地，\(\limsup \bar{X}_n\)
是常数（几乎必然）。

\subsubsection{105.2
中心极限定理}\label{ux4e2dux5fc3ux6781ux9650ux5b9aux7406}

\textbf{定理 105.4}（Lindeberg-Lévy 中心极限定理）：设
\(X_1, X_2, \ldots\) 是
i.i.d.，\(\mathbb{E}[X_1] = \mu\)，\(\operatorname{Var}(X_1) = \sigma^2 \in (0, \infty)\)。则

\[\frac{\sqrt{n}(\bar{X}_n - \mu)}{\sigma} \xrightarrow{d} \mathcal{N}(0, 1) \quad (n \to \infty)\]

即标准化样本均值依分布收敛于标准正态分布。

\emph{证明}：使用特征函数。\(\varphi_{Z_n}(t) = \varphi_{X_1-\mu}\left(\frac{t}{\sigma\sqrt{n}}\right)^n\)。由
Taylor
展开，\(\varphi_{X_1-\mu}(s) = 1 - \frac{\sigma^2 s^2}{2} + o(s^2)\)。故
\(\varphi_{Z_n}(t) \to e^{-t^2/2}\)，即 \(\mathcal{N}(0, 1)\)
的特征函数。∎

\textbf{定理 105.5}（Lindeberg-Feller CLT，三角阵列）：设对每个
\(n\)，\(X_{n,1}, \ldots, X_{n, k_n}\)
独立，\(\mathbb{E}[X_{n,i}] = 0\)，\(\sum_{i=1}^{k_n} \operatorname{Var}(X_{n,i}) = 1\)。若
\textbf{Lindeberg 条件}成立：

\[\lim_{n \to \infty} \sum_{i=1}^{k_n} \mathbb{E}[X_{n,i}^2 \mathbf{1}_{\{|X_{n,i}| > \varepsilon\}}] = 0 \quad (\forall \varepsilon > 0)\]

则
\(\sum_{i=1}^{k_n} X_{n,i} \xrightarrow{d} \mathcal{N}(0, 1)\)。Lindeberg
条件保证了没有单个变量主导和的行为。

\textbf{定理 105.6}（Berry-Esseen 界）：在上述 i.i.d. CLT 条件下，若
\(\mathbb{E}[|X_1|^3] < \infty\)，则

\[\sup_{x \in \mathbb{R}} \left| \mathbb{P}\left(\frac{\sqrt{n}(\bar{X}_n - \mu)}{\sigma} \leq x\right) - \Phi(x) \right| \leq \frac{C \mathbb{E}[|X_1 - \mu|^3]}{\sigma^3 \sqrt{n}}\]

其中 \(\Phi\) 是标准正态分布函数，\(C \leq 0.4748\)（已知最优常数）。

\subsubsection{105.3 特征函数}\label{ux7279ux5f81ux51fdux6570}

\textbf{定义 105.1}（特征函数）：随机变量 \(X\) 的\textbf{特征函数}为
\(\varphi_X(t) = \mathbb{E}[e^{itX}] = \int_{\mathbb{R}} e^{itx} d\mu_X(x)\)。特征函数是
Fourier 变换（或 Fourier-Stieltjes 变换）。

\textbf{命题 105.7}（特征函数的性质）： 1.
\(|\varphi_X(t)| \leq 1\)，\(\varphi_X(0) = 1\) 2. \(\varphi_X\)
一致连续 3. \(\varphi_{aX+b}(t) = e^{itb} \varphi_X(at)\) 4. 若 \(X, Y\)
独立，则 \(\varphi_{X+Y}(t) = \varphi_X(t) \varphi_Y(t)\) 5. 若
\(\mathbb{E}[|X|^k] < \infty\)，则
\(\varphi_X^{(k)}(0) = i^k \mathbb{E}[X^k]\)

\textbf{定理 105.8}（Lévy 连续性定理）：设 \(X_n\)
是随机变量序列，\(\varphi_n\) 是其特征函数。若
\(\varphi_n(t) \to \varphi(t)\) 逐点，且 \(\varphi\) 在 \(0\) 连续，则
\(\varphi\) 是某个随机变量的特征函数，且
\(X_n \xrightarrow{d} X\)（\(\varphi_X = \varphi\)）。

\bigskip\hrule\bigskip

\subsection{第106章：条件期望与鞅}\label{ux7b2c106ux7ae0ux6761ux4ef6ux671fux671bux4e0eux9785}

条件期望是概率论中最重要的概念之一，它将 ``已知部分信息下的平均''
精确化。鞅是 Doob
引入的强大工具，在随机过程、金融数学和统计推断中有广泛应用。

\subsubsection{106.1 条件期望}\label{ux6761ux4ef6ux671fux671b}

\textbf{定义 106.1}（条件期望）：设 \(X\)
是可积随机变量，\(\mathcal{G} \subseteq \mathcal{F}\) 是子
\(\sigma\)-代数。\(X\) 在 \(\mathcal{G}\) 下的\textbf{条件期望}
\(\mathbb{E}[X \mid \mathcal{G}]\) 是唯一的 \(\mathcal{G}\)-可测随机变量
\(Y\)，满足

\[\int_A Y d\mathbb{P} = \int_A X d\mathbb{P} \quad \text{对所有 } A \in \mathcal{G}\]

（由 Radon-Nikodym 定理保证存在唯一性）。

\textbf{定义
106.2}（给定随机变量的条件期望）：\(\mathbb{E}[X \mid Y] = \mathbb{E}[X \mid \sigma(Y)]\)。存在
Borel 函数 \(g\) 使得 \(\mathbb{E}[X \mid Y] = g(Y)\)，记
\(g(y) = \mathbb{E}[X \mid Y = y]\)。

\textbf{命题 106.1}（条件期望的基本性质）： 1.
\(\mathbb{E}[\mathbb{E}[X \mid \mathcal{G}]] = \mathbb{E}[X]\)（全期望律）
2. 若 \(X\) 是 \(\mathcal{G}\)-可测，则
\(\mathbb{E}[X \mid \mathcal{G}] = X\) 3. 若 \(X\) 与 \(\mathcal{G}\)
独立（即 \(\sigma(X)\) 与 \(\mathcal{G}\) 独立），则
\(\mathbb{E}[X \mid \mathcal{G}] = \mathbb{E}[X]\) 4.
\(\mathbb{E}[aX + bY \mid \mathcal{G}] = a \mathbb{E}[X \mid \mathcal{G}] + b \mathbb{E}[Y \mid \mathcal{G}]\)（线性）
5. 若 \(Z\) 是 \(\mathcal{G}\)-可测且有界，则
\(\mathbb{E}[ZX \mid \mathcal{G}] = Z \mathbb{E}[X \mid \mathcal{G}]\)
6.
\textbf{塔性质}：\(\mathbb{E}[\mathbb{E}[X \mid \mathcal{G}] \mid \mathcal{H}] = \mathbb{E}[X \mid \mathcal{H}]\)（当
\(\mathcal{H} \subseteq \mathcal{G}\)） 7. 条件 Jensen
不等式：\(\varphi(\mathbb{E}[X \mid \mathcal{G}]) \leq \mathbb{E}[\varphi(X) \mid \mathcal{G}]\)（\(\varphi\)
凸）

\subsubsection{106.2
鞅的定义与基本性质}\label{ux9785ux7684ux5b9aux4e49ux4e0eux57faux672cux6027ux8d28}

\textbf{定义 106.3}（滤过）：\textbf{滤过}
\(\{\mathcal{F}_n\}_{n \geq 0}\) 是 \(\mathcal{F}\) 的递增子
\(\sigma\)-代数序列：\(\mathcal{F}_0 \subseteq \mathcal{F}_1 \subseteq \cdots \subseteq \mathcal{F}\)。\(\mathcal{F}_n\)
表示时刻 \(n\) 可用的信息。

\textbf{定义 106.4}（鞅）：随机过程 \(\{M_n\}_{n \geq 0}\) 称为关于
\(\{\mathcal{F}_n\}\) 的\textbf{鞅}（martingale），如果 1. \(M_n\) 是
\(\mathcal{F}_n\)-可测（适应于滤过） 2.
\(\mathbb{E}[|M_n|] < \infty\)（对所有 \(n\)） 3.
\(\mathbb{E}[M_{n+1} \mid \mathcal{F}_n] = M_n\)（对所有 \(n\)）

若
\(\mathbb{E}[M_{n+1} \mid \mathcal{F}_n] \geq M_n\)（下鞅，submartingale）或
\(\leq M_n\)（上鞅，supermartingale），则称为\textbf{下鞅}或\textbf{上鞅}。

\textbf{例 106.1}（鞅的例子）： -
独立零均值随机变量的部分和：\(S_n = \sum_{i=1}^n X_i\)（\(\mathbb{E}[X_i] = 0\)）
-
独立随机变量的乘积：\(M_n = \prod_{i=1}^n X_i\)（\(\mathbb{E}[X_i] = 1\)）
- 条件期望过程：\(M_n = \mathbb{E}[X \mid \mathcal{F}_n]\)（对固定
\(X\)）

\subsubsection{106.3 鞅收敛定理}\label{ux9785ux6536ux655bux5b9aux7406}

\textbf{定理 106.2}（Doob 上穿不等式）：设 \(\{X_n\}\) 是下鞅。对
\(a < b\)，在 \(n\) 步内 \(X_n\) 从 \(a\) 以下穿到 \(b\) 以上的次数
\(U_n(a, b)\) 满足

\[\mathbb{E}[U_n(a, b)] \leq \frac{\mathbb{E}[(X_n - a)^+]}{b - a}\]

\textbf{定理 106.3}（鞅收敛定理，Doob）：设 \(\{M_n\}\) 是下鞅且
\(\sup_n \mathbb{E}[|M_n|] < \infty\)（或
\(\sup_n \mathbb{E}[M_n^+] < \infty\)）。则 \(M_n\)
几乎必然收敛到某个可积随机变量 \(M_\infty\)。

\emph{证明思路}：使用上穿不等式。若 \(\{M_n\}\) 不收敛，则存在 \(a < b\)
使得上穿次数无穷多。上穿不等式给出
\(\mathbb{E}[U_n(a, b)] \leq (\mathbb{E}[M_n^+] + |a|)/(b - a)\)
有界，与无穷多次上穿矛盾。∎

\textbf{定理 106.4}（\(L^p\) 鞅收敛定理）：若 \(\{M_n\}\) 是鞅且
\(\sup_n \mathbb{E}[|M_n|^p] < \infty\)（\(p > 1\)），则
\(M_n \to M_\infty\) 在 \(L^p\) 中收敛。

\subsubsection{106.4
可选停止定理}\label{ux53efux9009ux505cux6b62ux5b9aux7406}

\textbf{定义 106.5}（停时）：随机变量
\(\tau: \Omega \to \mathbb{N} \cup \{\infty\}\) 称为关于
\(\{\mathcal{F}_n\}\) 的\textbf{停时}，如果
\(\{\tau = n\} \in \mathcal{F}_n\)（对所有
\(n\)）。即''是否停止''仅依赖于当前和过去的信息。

\textbf{定理 106.5}（可选停止定理，Doob）：设 \(\{M_n\}\) 是鞅，\(\tau\)
是停时。若以下条件之一成立： 1. \(\tau\) 有界（\(\tau \leq N\)
几乎必然） 2. \(\{M_n\}\) 有界且 \(\tau < \infty\) 几乎必然 3.
\(\mathbb{E}[\tau] < \infty\) 且 \(|M_n - M_{n-1}| \leq C\)（有界增量）

则 \(\mathbb{E}[M_\tau] = \mathbb{E}[M_0]\)。

\emph{证明}：对情形
1，\(M_{\tau \wedge N} = \sum_{n=0}^{N-1} M_n \mathbf{1}_{\{\tau = n\}} + M_N \mathbf{1}_{\{\tau \geq N\}}\)
是鞅，故
\(\mathbb{E}[M_\tau] = \mathbb{E}[M_{\tau \wedge N}] = \mathbb{E}[M_0]\)。∎

\textbf{定理 106.6}（Wald 等式）：设 \(X_1, X_2, \ldots\) 是
i.i.d.，\(\mathbb{E}[|X_1|] < \infty\)，\(\mu = \mathbb{E}[X_1]\)。\(\tau\)
是关于 \(\{X_n\}\) 的停时，\(\mathbb{E}[\tau] < \infty\)。则
\(\mathbb{E}[\sum_{i=1}^\tau X_i] = \mu \mathbb{E}[\tau]\)。

\bigskip\hrule\bigskip

\subsection{第107章：Markov 链}\label{ux7b2c107ux7ae0markov-ux94fe}

Markov
链是''无记忆''的随机过程，其未来状态仅依赖于当前状态。本章讨论离散时间、可数状态空间的
Markov 链，包括转移概率、平稳分布和遍历定理。

\subsubsection{107.1 Markov
链的定义}\label{markov-ux94feux7684ux5b9aux4e49}

\textbf{定义 107.1}（Markov 链）：状态空间 \(S\)（可数集）上的随机过程
\(\{X_n\}_{n \geq 0}\) 称为 \textbf{Markov 链}，如果满足 Markov 性：

\[\mathbb{P}(X_{n+1} = j \mid X_n = i, X_{n-1}, \ldots, X_0) = \mathbb{P}(X_{n+1} = j \mid X_n = i)\]

\textbf{定义 107.2}（时齐 Markov 链）：若
\(\mathbb{P}(X_{n+1} = j \mid X_n = i) = p_{ij}\)（与 \(n\)
无关），则称链是\textbf{时齐}的。矩阵 \(P = (p_{ij})_{i,j \in S}\)
称为\textbf{转移矩阵}（每行和为 \(1\)）。

\textbf{定义 107.3}（\(n\)
步转移概率）：\(p_{ij}^{(n)} = \mathbb{P}(X_n = j \mid X_0 = i) = (P^n)_{ij}\)。由
Chapman-Kolmogorov
方程：\(p_{ij}^{(n+m)} = \sum_{k \in S} p_{ik}^{(n)} p_{kj}^{(m)}\)。

\textbf{命题 107.1}：Markov 链的分布由初始分布
\(\pi_0\)（\(\pi_0(i) = \mathbb{P}(X_0 = i)\)）和转移矩阵 \(P\)
完全确定：\(\pi_n = \pi_0 P^n\)。

\subsubsection{107.2 状态的分类}\label{ux72b6ux6001ux7684ux5206ux7c7b}

\textbf{定义 107.4}（可达与互通）：\(j\) 从 \(i\)
\textbf{可达}（\(i \to j\)）如果存在 \(n \geq 0\) 使得
\(p_{ij}^{(n)} > 0\)。\(i\) 和 \(j\) 互通（\(i \leftrightarrow j\)）如果
\(i \to j\) 且 \(j \to i\)。互通是等价关系，等价类称为\textbf{互通类}。

\textbf{定义 107.5}（常返与瞬时）：状态 \(i\) 称为\textbf{常返}的，如果
\(\mathbb{P}(X_n = i \text{ 无穷多次} \mid X_0 = i) = 1\)。否则称为\textbf{瞬时}的。

\textbf{定理 107.2}（常返性的判别）：\(i\) 常返当且仅当
\(\sum_{n=1}^\infty p_{ii}^{(n)} = \infty\)。\(i\) 瞬时当且仅当
\(\sum_{n=1}^\infty p_{ii}^{(n)} < \infty\)。

\textbf{定义 107.6}（正常返与零常返）：常返状态 \(i\)
称为\textbf{正常返}的，如果
\(\mathbb{E}[T_i \mid X_0 = i] < \infty\)（\(T_i = \inf\{n \geq 1 : X_n = i\}\)
是首次返回时间）。否则称为\textbf{零常返}。

\textbf{定义 107.7}（不可约）：Markov
链称为\textbf{不可约}的，如果所有状态互通（只有一个互通类）。

\textbf{定理 107.3}：对不可约 Markov
链，所有状态同时为瞬时、零常返或正常返（固态性）。

\subsubsection{107.3 平稳分布}\label{ux5e73ux7a33ux5206ux5e03}

\textbf{定义 107.8}（平稳分布）：概率分布 \(\pi = (\pi(i))_{i \in S}\)
称为\textbf{平稳分布}（或不变测度），如果 \(\pi = \pi P\)（即
\(\pi(j) = \sum_{i \in S} \pi(i) p_{ij}\)）。若 \(\pi_0 = \pi\)，则
\(\pi_n = \pi\) 对所有 \(n\)。

\textbf{定理 107.4}（平稳分布的存在性）：不可约 Markov
链有平稳分布当且仅当它是正常返的。此时平稳分布唯一，且
\(\pi(i) = 1 / \mathbb{E}[T_i \mid X_0 = i]\)。

\textbf{定理 107.5}（Markov 链遍历定理）：设不可约正常返 Markov
链有平稳分布 \(\pi\)。则

\[\lim_{n \to \infty} \frac{1}{n} \sum_{k=0}^{n-1} \mathbf{1}_{\{X_k = j\}} = \pi(j) \quad \text{几乎必然}\]

即时间平均收敛于空间平均。若链是非周期的，则
\(\lim_{n \to \infty} p_{ij}^{(n)} = \pi(j)\)（分布收敛）。

\textbf{定义 107.9}（周期）：状态 \(i\) 的\textbf{周期}
\(d(i) = \gcd\{n \geq 1 : p_{ii}^{(n)} > 0\}\)。若 \(d(i) = 1\)，则
\(i\) 非周期。不可约链中所有状态周期相同。

\subsubsection{107.4
可逆性与细致平衡}\label{ux53efux9006ux6027ux4e0eux7ec6ux81f4ux5e73ux8861}

\textbf{定义 107.10}（细致平衡）：概率分布 \(\pi\)
满足\textbf{细致平衡条件}（detailed balance），如果
\(\pi(i) p_{ij} = \pi(j) p_{ji}\)（对所有 \(i, j\)）。此时 \(\pi\)
是平稳分布，链称为\textbf{可逆}的。

\textbf{定理 107.6}（MCMC 的基本原理）：若目标分布 \(\pi\)
难以直接抽样，可构造 Markov 链使其平稳分布为
\(\pi\)。Metropolis-Hastings 算法和 Gibbs
抽样是构造满足细致平衡的转移矩阵的标准方法。

\bigskip\hrule\bigskip

\subsection{第108章：Brown
运动与随机过程}\label{ux7b2c108ux7ae0brown-ux8fd0ux52a8ux4e0eux968fux673aux8fc7ux7a0b}

Brown 运动（Wiener
过程）是最重要的连续时间随机过程，是随机分析的基础。本章给出 Brown
运动的构造、基本性质，并简要介绍连续时间鞅和 Itô 积分。

\subsubsection{108.1 Brown
运动的定义与构造}\label{brown-ux8fd0ux52a8ux7684ux5b9aux4e49ux4e0eux6784ux9020}

\textbf{定义 108.1}（Brown 运动 / Wiener 过程）：随机过程
\(\{B_t\}_{t \geq 0}\) 称为\textbf{标准 Brown 运动}（或 Wiener
过程），如果 1. \(B_0 = 0\) 几乎必然 2. 独立增量：对
\(0 \leq t_1 < t_2 < \cdots < t_n\)，\(B_{t_2} - B_{t_1}, \ldots, B_{t_n} - B_{t_{n-1}}\)
独立 3.
平稳增量：\(B_t - B_s \sim \mathcal{N}(0, t - s)\)（\(0 \leq s < t\)）
4. 样本路径 \(t \mapsto B_t\) 几乎必然连续

\textbf{定理 108.1}（Brown 运动的存在性，Wiener 1923）：Brown
运动存在。可通过以下方式构造： - 在 \(C[0, \infty)\) 上定义 Wiener
测度（Kolmogorov 扩张定理） - 或使用 Lévy-Ciesielski 构造（Schauder
基上的随机 Fourier 级数）

\textbf{命题 108.1}（Brown 运动的性质）： 1.
尺度不变性：\(\{c^{-1/2} B_{ct}\}\) 也是 Brown 运动 2.
时间反演：\(\{t B_{1/t}\}\)（\(t > 0\)，\(B_0 = 0\)）是 Brown 运动 3.
反射原理：\(\mathbb{P}(\max_{0 \leq s \leq t} B_s \geq a) = 2\mathbb{P}(B_t \geq a)\)（\(a > 0\)）
4. 样本路径几乎必然无处可微（Paley-Wiener-Zygmund 定理） 5. 样本路径在
\([0, t]\) 上的变差无穷大（几乎必然），但二次变差为 \(t\)

\subsubsection{108.2 Brown 运动的 Markov 性与强 Markov
性}\label{brown-ux8fd0ux52a8ux7684-markov-ux6027ux4e0eux5f3a-markov-ux6027}

\textbf{定理 108.2}（Markov 性）：Brown 运动是 Markov 过程：对
\(s > 0\)，\(\{B_{t+s} - B_s\}_{t \geq 0}\) 是 Brown 运动且与
\(\sigma(B_u, u \leq s)\) 独立。

\textbf{定理 108.3}（强 Markov 性）：设 \(\tau\) 是关于 Brown
运动自然滤过的停时（\(\tau < \infty\) 几乎必然）。则
\(\{B_{\tau + t} - B_\tau\}_{t \geq 0}\) 是 Brown 运动且与
\(\mathcal{F}_\tau\) 独立。

\textbf{定理 108.4}（Blumenthal 0-1 律）：设
\(\mathcal{F}_{0^+} = \bigcap_{t > 0} \mathcal{F}_t\)。则对
\(A \in \mathcal{F}_{0^+}\)，\(\mathbb{P}(A) \in \{0, 1\}\)。特别地，\(\mathbb{P}(\sup_{[0, \varepsilon]} B_t > 0) = 1\)（Brown
运动在原点附近立即取正值和负值）。

\subsubsection{108.3 连续时间鞅}\label{ux8fdeux7eedux65f6ux95f4ux9785}

\textbf{定义 108.2}（连续时间鞅）：过程 \(\{M_t\}_{t \geq 0}\) 关于滤过
\(\{\mathcal{F}_t\}\) 是\textbf{鞅}，如果 \(\mathbb{E}[|M_t|] < \infty\)
且 \(\mathbb{E}[M_t \mid \mathcal{F}_s] = M_s\)（\(s < t\)）。

\textbf{命题 108.2}（Brown 运动生成的鞅）： 1. \(B_t\) 是鞅 2.
\(B_t^2 - t\) 是鞅 3. \(\exp(\theta B_t - \theta^2 t/2)\) 是鞅（指数鞅）

\textbf{定理 108.5}（Doob 不等式，连续时间）：设 \(\{M_t\}\)
是右连续下鞅。则

\[\mathbb{P}\left(\sup_{0 \leq s \leq t} |M_s| \geq \lambda\right) \leq \frac{\mathbb{E}[|M_t|]}{\lambda}\]

若 \(M_t\) 是鞅且 \(p > 1\)，则
\(\mathbb{E}[\sup_{0 \leq s \leq t} |M_s|^p] \leq \left(\frac{p}{p-1}\right)^p \mathbb{E}[|M_t|^p]\)。

\subsubsection{108.4 Itô
积分简介}\label{ituxf4-ux79efux5206ux7b80ux4ecb}

\textbf{定义 108.3}（Itô 积分，启发式）：对适应过程 \(H_s\)，Itô 积分
\(\int_0^t H_s dB_s\) 定义为

\[\int_0^t H_s dB_s = \lim_{\|\Pi\| \to 0} \sum_{i} H_{t_i} (B_{t_{i+1}} - B_{t_i})\]

（在 \(L^2\) 意义下），其中被积函数在区间左端点取值（非预见性）。

\textbf{定理 108.6}（Itô 公式）：设 \(f(t, x)\) 是 \(C^{1,2}\) 函数。则

\[df(t, B_t) = \frac{\partial f}{\partial t}(t, B_t) dt + \frac{\partial f}{\partial x}(t, B_t) dB_t + \frac{1}{2} \frac{\partial^2 f}{\partial x^2}(t, B_t) dt\]

最后一项是 Itô 修正项，来源于 Brown 运动的非零二次变差。

\textbf{例 108.1}（Itô 公式的应用）： -
\(f(x) = x^2\)：\(d(B_t^2) = 2B_t dB_t + dt\)（几何 Brown 运动的基础） -
\(f(x) = e^x\)：\(d(e^{B_t}) = e^{B_t} dB_t + \frac{1}{2} e^{B_t} dt\)

\bigskip\hrule\bigskip

\emph{本卷在初等概率论（V6）和测度论（V9）的基础上，建立了概率论的严格分析基础：Kolmogorov
公理体系将概率论纳入测度论框架；大数定律和中心极限定理揭示了随机现象中的确定性规律；条件期望和鞅论提供了处理信息的强大工具；Markov
链理论为随机过程建模提供了基本框架；Brown 运动和 Itô
积分开创了随机分析，为金融数学（Black-Scholes
模型）和随机微分方程奠定了基础。为后续数理统计（V22）提供了概率论基础。}

\newpage

\section{卷二十一：统计
II}\label{ux5377ux4e8cux5341ux4e00ux7edfux8ba1-ii}

\begin{quote}
数理统计在初等统计（V6 Ch 30）和概率论
II（V20）的测度论基础上，建立参数估计、假设检验和贝叶斯方法的严格理论。本卷从估计理论出发，深入讨论
Neyman-Pearson
引理、似然比检验、置信区间、贝叶斯推断、线性模型和多元分析。
\end{quote}

\bigskip\hrule\bigskip

\subsection{第109章：估计理论}\label{ux7b2c109ux7ae0ux4f30ux8ba1ux7406ux8bba}

估计理论是统计推断的基础，研究如何从样本中估计未知参数。本章讨论无偏估计、Fisher
信息、Cramér-Rao 下界和极大似然估计。

\subsubsection{109.1
统计模型与充分统计量}\label{ux7edfux8ba1ux6a21ux578bux4e0eux5145ux5206ux7edfux8ba1ux91cf}

\textbf{定义 109.1}（参数统计模型）：\textbf{参数统计模型}是一族概率分布
\(\{P_\theta : \theta \in \Theta\}\)，其中
\(\Theta \subseteq \mathbb{R}^k\) 是参数空间。样本 \(X_1, \ldots, X_n\)
是 i.i.d. 来自 \(P_\theta\)。

\textbf{定义 109.2}（统计量）：\textbf{统计量}
\(T = T(X_1, \ldots, X_n)\) 是样本的可测函数（不依赖于未知参数
\(\theta\)）。

\textbf{定义 109.3}（充分统计量）：统计量 \(T\) 称为关于 \(\theta\)
\textbf{充分}的，如果 \(X\) 在 \(T\) 给定下的条件分布不依赖于
\(\theta\)。即 \(T\) 包含了样本中关于 \(\theta\) 的所有信息。

\textbf{定理 109.1}（Fisher-Neyman 因子分解定理）：设 \(\{P_\theta\}\)
有密度（或概率质量函数）\(f(x \mid \theta)\)。\(T\) 是充分的当且仅当

\[f(x_1, \ldots, x_n \mid \theta) = g(T(x_1, \ldots, x_n), \theta) \cdot h(x_1, \ldots, x_n)\]

其中 \(g\) 依赖于 \(\theta\) 和 \(T\)，\(h\) 不依赖于 \(\theta\)。

\emph{证明}：离散情形。若 \(T\) 充分，定义
\(g(t, \theta) = P_\theta(T = t)\)，\(h(x) = P(X = x \mid T = t) / P_\theta(T = t)\)（后者不依赖于
\(\theta\)）。反之亦然。∎

\textbf{例 109.1}：对正态分布
\(X_i \sim \mathcal{N}(\mu, \sigma^2)\)，\((\bar{X}, S^2)\)（样本均值和方差）是
\((\mu, \sigma^2)\) 的充分统计量。对
\(\mathcal{N}(\mu, 1)\)，\(\bar{X}\) 是 \(\mu\) 的充分统计量。

\subsubsection{109.2 点估计}\label{ux70b9ux4f30ux8ba1}

\textbf{定义 109.4}（估计量）：参数 \(\theta\)（或
\(g(\theta)\)）的\textbf{估计量}
\(\hat{\theta}_n = \hat{\theta}_n(X_1, \ldots, X_n)\) 是样本的统计量。

\textbf{定义 109.5}（无偏性）：\(\hat{\theta}_n\) 称为
\textbf{无偏}的，如果
\(\mathbb{E}_\theta[\hat{\theta}_n] = \theta\)（对所有
\(\theta \in \Theta\)）。其偏差为
\(\operatorname{Bias}_\theta(\hat{\theta}_n) = \mathbb{E}_\theta[\hat{\theta}_n] - \theta\)。

\textbf{定义
109.6}（均方误差，MSE）：\(\operatorname{MSE}_\theta(\hat{\theta}_n) = \mathbb{E}_\theta[(\hat{\theta}_n - \theta)^2] = \operatorname{Var}_\theta(\hat{\theta}_n) + (\operatorname{Bias}_\theta(\hat{\theta}_n))^2\)。

\textbf{定义 109.7}（相合性）：\(\hat{\theta}_n\)
称为\textbf{相合}的（一致的），如果
\(\hat{\theta}_n \xrightarrow{\mathbb{P}_\theta} \theta\)（对所有
\(\theta\)）。强相合：\(\hat{\theta}_n \xrightarrow{\text{a.s.}} \theta\)。

\subsubsection{109.3 Fisher 信息与 Cramér-Rao
下界}\label{fisher-ux4fe1ux606fux4e0e-cramuxe9r-rao-ux4e0bux754c}

\textbf{定义 109.8}（得分函数与 Fisher 信息）：设模型
\(\{f(x \mid \theta) : \theta \in \Theta\}\)（\(\Theta \subseteq \mathbb{R}\)）满足正则条件。\textbf{得分函数}为
\(S(\theta, x) = \frac{\partial}{\partial \theta} \log f(x \mid \theta)\)。单个观测的
\textbf{Fisher 信息}为

\[I(\theta) = \mathbb{E}_\theta\left[\left(\frac{\partial}{\partial \theta} \log f(X \mid \theta)\right)^2\right] = -\mathbb{E}_\theta\left[\frac{\partial^2}{\partial \theta^2} \log f(X \mid \theta)\right]\]

\(n\) 个 i.i.d. 观测的 Fisher 信息为 \(n I(\theta)\)。

\textbf{定理 109.2}（Cramér-Rao 下界）：设 \(\hat{\theta}_n\) 是
\(g(\theta)\) 的无偏估计量。在正则条件下，

\[\operatorname{Var}_\theta(\hat{\theta}_n) \geq \frac{(g'(\theta))^2}{n I(\theta)}\]

特别地，对 \(\theta\)
的无偏估计，\(\operatorname{Var}_\theta(\hat{\theta}_n) \geq 1/(n I(\theta))\)。

\emph{证明思路}：使用 Cauchy-Schwarz
不等式。\(\operatorname{Cov}_\theta(\hat{\theta}_n, S_n) = g'(\theta)\)，其中
\(S_n\)
是总得分函数。\(|\operatorname{Cov}|^2 \leq \operatorname{Var}(\hat{\theta}_n) \cdot \operatorname{Var}(S_n)\)，而
\(\operatorname{Var}(S_n) = n I(\theta)\)。∎

\textbf{定义 109.9}（有效估计量）：达到 Cramér-Rao
下界的无偏估计量称为\textbf{有效}的。有效估计量存在当且仅当分布属于指数族。大多数常用分布（正态、二项、Poisson、Gamma、Beta
等）属于指数族，但并非所有分布都如此（如 Cauchy
分布、均匀分布不属于指数族）。

\subsubsection{109.4
极大似然估计}\label{ux6781ux5927ux4f3cux7136ux4f30ux8ba1}

\textbf{定义 109.10}（似然函数）：\textbf{似然函数}
\(L(\theta \mid x) = f(x \mid \theta)\)（将 \(\theta\)
视为变量）。\textbf{对数似然}
\(\ell(\theta \mid x) = \log L(\theta \mid x)\)。

\textbf{定义 109.11}（极大似然估计，MLE）：\textbf{极大似然估计量}
\(\hat{\theta}_{\text{MLE}}\) 最大化 \(L(\theta \mid X)\)（或等价地
\(\ell(\theta \mid X)\)）：

\[\hat{\theta}_{\text{MLE}} = \arg\max_{\theta \in \Theta} L(\theta \mid X)\]

\textbf{定理 109.3}（MLE 的渐近性质）：在正则条件下： 1.
\textbf{相合性}：\(\hat{\theta}_{\text{MLE}} \xrightarrow{\mathbb{P}} \theta_0\)（真值）
2.
\textbf{渐近正态性}：\(\sqrt{n}(\hat{\theta}_{\text{MLE}} - \theta_0) \xrightarrow{d} \mathcal{N}(0, I(\theta_0)^{-1})\)
3. \textbf{渐近有效性}：MLE 渐近达到 Cramér-Rao 下界

\emph{证明思路}：\(\hat{\theta}\) 满足
\(\ell'(\hat{\theta}) = 0\)。Taylor
展开：\(0 = \ell'(\theta_0) + \ell''(\theta_0)(\hat{\theta} - \theta_0) + \cdots\)。\(\sqrt{n}(\hat{\theta} - \theta_0) \approx -\ell'(\theta_0)/\ell''(\theta_0) \cdot \sqrt{n}\)。由
CLT，\(\ell'(\theta_0)/\sqrt{n} \xrightarrow{d} \mathcal{N}(0, I(\theta_0))\)，由大数定律，\(\ell''(\theta_0)/n \to -I(\theta_0)\)。由
Slutsky 定理得证。∎

\textbf{定理 109.4}（MLE 的不变性）：若 \(\hat{\theta}\) 是 \(\theta\)
的 MLE，\(g\) 是任意函数，则 \(g(\hat{\theta})\) 是 \(g(\theta)\) 的
MLE。

\bigskip\hrule\bigskip

\subsection{第110章：假设检验}\label{ux7b2c110ux7ae0ux5047ux8bbeux68c0ux9a8c}

假设检验是统计推断的另一个核心，用于判断数据是否支持某一假设。本章介绍
Neyman-Pearson 引理和一致最优势检验。

\subsubsection{110.1
假设检验的基本概念}\label{ux5047ux8bbeux68c0ux9a8cux7684ux57faux672cux6982ux5ff5}

\textbf{定义
110.1}（假设检验）：\textbf{假设检验}问题中，\textbf{原假设}
\(H_0: \theta \in \Theta_0\) 与\textbf{备择假设}
\(H_1: \theta \in \Theta_1\)（\(\Theta_0 \cap \Theta_1 = \varnothing\)）。

\textbf{定义 110.2}（检验函数与拒绝域）：\textbf{检验函数}
\(\varphi(x) \in [0, 1]\) 是拒绝 \(H_0\) 的概率（给定数据
\(x\)）。\textbf{拒绝域} \(R = \{x : \varphi(x) = 1\}\)。若
\(\varphi(x) \in \{0, 1\}\)，称为\textbf{非随机化检验}。

\textbf{定义 110.3}（两类错误）： -
\textbf{第一类错误}（假阳性）：\(H_0\) 为真但被拒绝。概率
\(\alpha(\theta) = \mathbb{E}_\theta[\varphi(X)]\)（\(\theta \in \Theta_0\)）
- \textbf{第二类错误}（假阴性）：\(H_1\) 为真但未被拒绝。概率
\(\beta(\theta) = 1 - \mathbb{E}_\theta[\varphi(X)]\)（\(\theta \in \Theta_1\)）

\textbf{定义 110.4}（检验水平与功效）：检验 \(\varphi\)
的\textbf{水平}（size）为
\(\alpha = \sup_{\theta \in \Theta_0} \mathbb{E}_\theta[\varphi(X)]\)。\textbf{功效函数}（power）为
\(\pi(\theta) = \mathbb{E}_\theta[\varphi(X)]\)（\(\theta \in \Theta_1\)）。

\subsubsection{110.2 Neyman-Pearson
引理}\label{neyman-pearson-ux5f15ux7406}

\textbf{定理 110.1}（Neyman-Pearson 引理）：考虑简单假设检验
\(H_0: \theta = \theta_0\) vs \(H_1: \theta = \theta_1\)。设
\(f_0, f_1\) 分别是 \(H_0\) 和 \(H_1\) 下的密度。则水平 \(\alpha\)
的\textbf{最优势检验}（最大功效）由似然比检验给出：

\[\varphi(x) = \begin{cases} 1 & \text{若 } \frac{f_1(x)}{f_0(x)} > c \\ \gamma & \text{若 } \frac{f_1(x)}{f_0(x)} = c \\ 0 & \text{若 } \frac{f_1(x)}{f_0(x)} < c \end{cases}\]

其中 \(c\) 和 \(\gamma\) 选择使得
\(\mathbb{E}_{\theta_0}[\varphi(X)] = \alpha\)。

\emph{证明}：设 \(\varphi^*\) 是 Neyman-Pearson 检验，\(\varphi\)
是任意其他水平 \(\leq \alpha\) 的检验。则
\((\varphi^* - \varphi)(f_1 - c f_0) \geq 0\) 逐点。积分得
\(\mathbb{E}_{\theta_1}[\varphi^*] - \mathbb{E}_{\theta_1}[\varphi] \geq 0\)。∎

\textbf{推论 110.2}（单调似然比）：若 \(\{f_\theta\}\) 关于统计量
\(T(x)\) 有单调似然比（即 \(f_{\theta_1}(x)/f_{\theta_0}(x)\) 是
\(T(x)\) 的非降函数，\(\theta_1 > \theta_0\)），则对
\(H_0: \theta \leq \theta_0\) vs
\(H_1: \theta > \theta_0\)，存在一致最优势（UMP）检验，拒绝域为
\(T(x) > c\)。

\subsubsection{110.3 似然比检验}\label{ux4f3cux7136ux6bd4ux68c0ux9a8c}

\textbf{定义 110.5}（广义似然比检验，GLRT）：对复合假设
\(H_0: \theta \in \Theta_0\) vs
\(H_1: \theta \in \Theta \setminus \Theta_0\)，\textbf{似然比统计量}为

\[\Lambda(X) = \frac{\sup_{\theta \in \Theta_0} L(\theta \mid X)}{\sup_{\theta \in \Theta} L(\theta \mid X)}\]

拒绝域为 \(\Lambda(X) \leq c\)（\(0 \leq c \leq 1\)）。

\textbf{定理 110.3}（Wilks 定理）：在 \(H_0\)
下（正则条件），\(-2 \log \Lambda(X) \xrightarrow{d} \chi^2_{\nu}\)，其中
\(\nu = \dim \Theta - \dim \Theta_0\)（自由度等于被检验的参数个数）。

\textbf{例 110.1}（\(t\) 检验）：对正态均值 \(\mu\) 的检验
\(H_0: \mu = \mu_0\) vs \(H_1: \mu \neq \mu_0\)（\(\sigma^2\)
未知），\(t\) 统计量
\(T = \sqrt{n}(\bar{X} - \mu_0)/S \sim t_{n-1}\)（在 \(H_0\)
下）。拒绝域：\(|T| > t_{n-1, \alpha/2}\)。

\textbf{例 110.2}（\(\chi^2\)
拟合优度检验）：检验样本是否来自某个分布。Pearson \(\chi^2\) 统计量
\(\sum (O_i - E_i)^2 / E_i \xrightarrow{d} \chi^2_{k-1}\)（在 \(H_0\)
下）。

\bigskip\hrule\bigskip

\subsection{第111章：置信区间与贝叶斯方法}\label{ux7b2c111ux7ae0ux7f6eux4fe1ux533aux95f4ux4e0eux8d1dux53f6ux65afux65b9ux6cd5}

本章讨论区间估计和置信区间，以及贝叶斯推断的基本框架。

\subsubsection{111.1 置信区间}\label{ux7f6eux4fe1ux533aux95f4}

\textbf{定义 111.1}（置信区间）：参数 \(\theta\) 的 \(1 - \alpha\)
\textbf{置信区间}（或置信集）是只依赖于样本的随机集合 \(C(X)\)，满足

\[\mathbb{P}_\theta(\theta \in C(X)) \geq 1 - \alpha \quad \text{对所有 } \theta \in \Theta\]

\(1 - \alpha\) 称为\textbf{置信水平}（或覆盖概率）。

\textbf{定义 111.2}（枢轴量）：随机量 \(Q(X, \theta)\)
称为\textbf{枢轴量}（pivotal quantity），如果其分布不依赖于
\(\theta\)。由 \(\mathbb{P}(a \leq Q(X, \theta) \leq b) = 1 - \alpha\)
可解出 \(\theta\) 的置信区间。

\textbf{例 111.1}（正态均值的置信区间）： - \(\sigma^2\)
已知：\(\bar{X} \pm z_{\alpha/2} \sigma / \sqrt{n}\)（\(z_{\alpha/2}\)
是标准正态分位数） - \(\sigma^2\)
未知：\(\bar{X} \pm t_{n-1, \alpha/2} S / \sqrt{n}\)（\(t\) 分布分位数）

\textbf{定理 111.1}（置信区间与假设检验的对偶）：\(\theta\) 的
\(1 - \alpha\) 置信集是所有不被水平 \(\alpha\) 检验拒绝的 \(\theta_0\)
的集合。反之，水平 \(\alpha\) 的检验可通过检查 \(\theta_0\)
是否在置信区间内构造。

\subsubsection{111.2 贝叶斯方法}\label{ux8d1dux53f6ux65afux65b9ux6cd5}

\textbf{定义 111.3}（贝叶斯模型）：在贝叶斯框架中，参数 \(\theta\)
被视为随机变量，具有\textbf{先验分布} \(\pi(\theta)\)。给定
\(\theta\)，数据 \(X\) 的分布为 \(f(x \mid \theta)\)。

\textbf{定理 111.2}（贝叶斯定理）：\textbf{后验分布}为

\[\pi(\theta \mid x) = \frac{f(x \mid \theta) \pi(\theta)}{\int_\Theta f(x \mid \theta) \pi(\theta) d\theta} \propto f(x \mid \theta) \pi(\theta)\]

即后验 \(\propto\) 似然 \(\times\) 先验。

\textbf{定义 111.4}（共轭先验）：先验族 \(\mathcal{F}\) 称为关于似然
\(f(x \mid \theta)\)\textbf{共轭}的，如果后验分布也属于
\(\mathcal{F}\)。

\textbf{例 111.2}（常见共轭先验）： -
Beta-Binomial：\(\theta \sim \operatorname{Beta}(\alpha, \beta)\)，\(X \mid \theta \sim \operatorname{Binomial}(n, \theta)\)，后验
\(\theta \mid X \sim \operatorname{Beta}(\alpha + X, \beta + n - X)\) -
Normal-Normal：\(\theta \sim \mathcal{N}(\mu_0, \tau_0^2)\)，\(X_i \mid \theta \sim \mathcal{N}(\theta, \sigma^2)\)（\(\sigma^2\)
已知），后验 \(\theta \mid X \sim \mathcal{N}(\mu_n, \tau_n^2)\) -
Gamma-Poisson：\(\lambda \sim \operatorname{Gamma}(\alpha, \beta)\)，\(X_i \mid \lambda \sim \operatorname{Poisson}(\lambda)\)，后验
\(\lambda \mid X \sim \operatorname{Gamma}(\alpha + \sum X_i, \beta + n)\)

\textbf{定义 111.5}（贝叶斯估计）：常用贝叶斯估计包括： -
\textbf{后验均值}：\(\hat{\theta}_{\text{Bayes}} = \mathbb{E}[\theta \mid X]\)（最小化平方误差损失）
-
\textbf{后验中位数}：\(\hat{\theta}_{\text{med}}\)（最小化绝对误差损失）
-
\textbf{最大后验（MAP）}：\(\hat{\theta}_{\text{MAP}} = \arg\max_\theta \pi(\theta \mid X)\)（最小化
0-1 损失）

\textbf{定理 111.3}（Bernstein-von Mises 定理）：在正则条件下，当样本量
\(n \to \infty\) 时，后验分布渐近为正态分布
\(\mathcal{N}(\hat{\theta}_{\text{MLE}}, (n I(\hat{\theta}_{\text{MLE}}))^{-1})\)，且与先验分布无关（先验的影响随
\(n \to \infty\) 消失）。

\textbf{例 111.3}（Jeffreys 先验）：无信息先验的经典选择是
\textbf{Jeffreys
先验}：\(\pi(\theta) \propto \sqrt{I(\theta)}\)（\(I(\theta)\) 是 Fisher
信息）。Jeffreys 先验在参数变换下不变。

\bigskip\hrule\bigskip

\subsection{第112章：线性模型}\label{ux7b2c112ux7ae0ux7ebfux6027ux6a21ux578b}

线性模型是统计中最广泛使用的模型框架，包括回归、方差分析（ANOVA）和协方差分析。本章讨论最小二乘估计、Gauss-Markov
定理和方差分析。

\subsubsection{112.1
线性回归模型}\label{ux7ebfux6027ux56deux5f52ux6a21ux578b}

\textbf{定义 112.1}（线性模型）：\textbf{线性模型}为

\[\mathbf{Y} = \mathbf{X}\boldsymbol{\beta} + \boldsymbol{\varepsilon}\]

其中 \(\mathbf{Y} \in \mathbb{R}^n\)
是响应向量，\(\mathbf{X} \in \mathbb{R}^{n \times p}\)
是设计矩阵（已知），\(\boldsymbol{\beta} \in \mathbb{R}^p\)
是未知回归系数，\(\boldsymbol{\varepsilon}\) 是误差向量，满足
\(\mathbb{E}[\boldsymbol{\varepsilon}] = \mathbf{0}\)，\(\operatorname{Cov}(\boldsymbol{\varepsilon}) = \sigma^2 \mathbf{I}_n\)。

\textbf{定义 112.2}（最小二乘估计，OLS）：\(\boldsymbol{\beta}\)
的\textbf{最小二乘估计}为

\[\hat{\boldsymbol{\beta}} = \arg\min_{\boldsymbol{\beta}} \|\mathbf{Y} - \mathbf{X}\boldsymbol{\beta}\|^2 = (\mathbf{X}^\top \mathbf{X})^{-1} \mathbf{X}^\top \mathbf{Y}\]

（假设 \(\mathbf{X}^\top \mathbf{X}\) 可逆，即 \(\mathbf{X}\) 满列秩）。

\textbf{定理 112.1}（Gauss-Markov 定理）：在
\(\mathbb{E}[\boldsymbol{\varepsilon}] = \mathbf{0}\)，\(\operatorname{Cov}(\boldsymbol{\varepsilon}) = \sigma^2 \mathbf{I}_n\)
的假设下，\(\hat{\boldsymbol{\beta}}\) 是 \(\boldsymbol{\beta}\)
的\textbf{最佳线性无偏估计}（BLUE）：对所有
\(c \in \mathbb{R}^p\)，\(c^\top \hat{\boldsymbol{\beta}}\) 在
\(c^\top \boldsymbol{\beta}\) 的所有线性无偏估计中方差最小。

\emph{证明}：设 \(\tilde{\boldsymbol{\beta}} = \mathbf{C}\mathbf{Y}\)
是任意线性无偏估计。\(\mathbb{E}[\tilde{\boldsymbol{\beta}}] = \mathbf{C}\mathbf{X}\boldsymbol{\beta} = \boldsymbol{\beta}\)
对所有 \(\boldsymbol{\beta}\) 成立，故
\(\mathbf{C}\mathbf{X} = \mathbf{I}\)。\(\operatorname{Cov}(\tilde{\boldsymbol{\beta}}) = \sigma^2 \mathbf{C}\mathbf{C}^\top\)。\(\operatorname{Cov}(\hat{\boldsymbol{\beta}}) = \sigma^2 (\mathbf{X}^\top \mathbf{X})^{-1}\)。需证
\(\mathbf{C}\mathbf{C}^\top - (\mathbf{X}^\top \mathbf{X})^{-1}\)
半正定，这等价于
\((\mathbf{C} - (\mathbf{X}^\top \mathbf{X})^{-1}\mathbf{X}^\top)(\mathbf{C} - (\mathbf{X}^\top \mathbf{X})^{-1}\mathbf{X}^\top)^\top \succeq 0\)。∎

\textbf{定理 112.2}（正态误差下的推断）：若进一步
\(\boldsymbol{\varepsilon} \sim \mathcal{N}(\mathbf{0}, \sigma^2 \mathbf{I}_n)\)，则
1.
\(\hat{\boldsymbol{\beta}} \sim \mathcal{N}(\boldsymbol{\beta}, \sigma^2 (\mathbf{X}^\top \mathbf{X})^{-1})\)
2.
\(\hat{\sigma}^2 = \|\mathbf{Y} - \mathbf{X}\hat{\boldsymbol{\beta}}\|^2 / (n - p)\)
是 \(\sigma^2\)
的无偏估计，\((n-p)\hat{\sigma}^2/\sigma^2 \sim \chi^2_{n-p}\) 3.
\(\hat{\beta}_j / (\hat{\sigma} \sqrt{(\mathbf{X}^\top \mathbf{X})^{-1}_{jj}}) \sim t_{n-p}\)（用于
\(H_0: \beta_j = 0\) 的检验）

\subsubsection{112.2 方差分析}\label{ux65b9ux5deeux5206ux6790}

\textbf{定义 112.3}（一元方差分析，ANOVA）：\textbf{一元方差分析}比较
\(k\) 个总体的均值：

\[H_0: \mu_1 = \mu_2 = \cdots = \mu_k \quad \text{vs} \quad H_1: \text{至少两个均值不同}\]

\textbf{定义 112.4}（平方和分解）：

\[\underbrace{\sum_{i=1}^k \sum_{j=1}^{n_i} (Y_{ij} - \bar{Y}_{\cdot\cdot})^2}_{\text{SST (总平方和)}} = \underbrace{\sum_{i=1}^k n_i (\bar{Y}_{i\cdot} - \bar{Y}_{\cdot\cdot})^2}_{\text{SSB (组间平方和)}} + \underbrace{\sum_{i=1}^k \sum_{j=1}^{n_i} (Y_{ij} - \bar{Y}_{i\cdot})^2}_{\text{SSW (组内平方和)}}\]

\textbf{定理 112.3}（ANOVA \(F\) 检验）：在 \(H_0\)
下（正态性、等方差假设），

\[F = \frac{\text{SSB} / (k-1)}{\text{SSW} / (n-k)} \sim F_{k-1, n-k}\]

拒绝域：\(F > F_{k-1, n-k, \alpha}\)。

\textbf{定义 112.5}（多元线性回归的 \(F\) 检验）：对
\(H_0: \beta_{q+1} = \cdots = \beta_p = 0\)（嵌套模型比较），

\[F = \frac{(\text{SSE}_{\text{reduced}} - \text{SSE}_{\text{full}}) / (p - q)}{\text{SSE}_{\text{full}} / (n - p)} \sim F_{p-q, n-p}\]

其中 SSE 是残差平方和。

\subsubsection{112.3
模型诊断与选择}\label{ux6a21ux578bux8bcaux65adux4e0eux9009ux62e9}

\textbf{定义
112.6}（决定系数）：\(R^2 = 1 - \frac{\text{SSE}}{\text{SST}}\)
是模型解释的变异比例。调整
\(R^2\)：\(R^2_{\text{adj}} = 1 - \frac{\text{SSE}/(n-p)}{\text{SST}/(n-1)}\)。

\textbf{定义 112.7}（AIC 与 BIC）：模型选择准则： - \textbf{AIC}（Akaike
信息准则）：\(\text{AIC} = -2 \log L + 2p\) -
\textbf{BIC}（贝叶斯信息准则）：\(\text{BIC} = -2 \log L + p \log n\)

选择 AIC/BIC 最小的模型。BIC 对复杂模型的惩罚更重（当 \(n \geq 8\)）。

\bigskip\hrule\bigskip

\subsection{第113章：多元分析}\label{ux7b2c113ux7ae0ux591aux5143ux5206ux6790}

多元分析处理多个响应变量同时观测的情形。本章介绍多元正态分布、主成分分析（PCA）和判别分析。

\subsubsection{113.1
多元正态分布}\label{ux591aux5143ux6b63ux6001ux5206ux5e03}

\textbf{定义 113.1}（多元正态分布）：\(\mathbb{R}^p\) 上的随机向量
\(\mathbf{X} \sim \mathcal{N}_p(\boldsymbol{\mu}, \boldsymbol{\Sigma})\)（\(\boldsymbol{\Sigma} \succ 0\)）的密度为

\[f(\mathbf{x}) = \frac{1}{(2\pi)^{p/2} |\boldsymbol{\Sigma}|^{1/2}} \exp\left(-\frac{1}{2} (\mathbf{x} - \boldsymbol{\mu})^\top \boldsymbol{\Sigma}^{-1} (\mathbf{x} - \boldsymbol{\mu})\right)\]

\textbf{命题 113.1}（多元正态的性质）： 1.
线性变换：\(\mathbf{A}\mathbf{X} + \mathbf{b} \sim \mathcal{N}(\mathbf{A}\boldsymbol{\mu} + \mathbf{b}, \mathbf{A}\boldsymbol{\Sigma}\mathbf{A}^\top)\)
2. 边缘分布：\(\mathbf{X}\) 的任意子向量也是多元正态 3.
条件分布：\(\mathbf{X}_1 \mid \mathbf{X}_2 = \mathbf{x}_2\) 是正态，均值
\(\boldsymbol{\mu}_1 + \boldsymbol{\Sigma}_{12}\boldsymbol{\Sigma}_{22}^{-1}(\mathbf{x}_2 - \boldsymbol{\mu}_2)\)
4. 独立性：\(\mathbf{X}_1\) 与 \(\mathbf{X}_2\) 独立当且仅当
\(\boldsymbol{\Sigma}_{12} = \mathbf{0}\)

\textbf{定理 113.1}（Wishart 分布）：设
\(\mathbf{X}_1, \ldots, \mathbf{X}_n\) i.i.d.
\(\sim \mathcal{N}_p(\mathbf{0}, \boldsymbol{\Sigma})\)。则
\(\mathbf{W} = \sum_{i=1}^n \mathbf{X}_i \mathbf{X}_i^\top\) 服从
\textbf{Wishart 分布} \(W_p(n, \boldsymbol{\Sigma})\)（多元 \(\chi^2\)
的推广）。

\textbf{定理 113.2}（Hotelling \(T^2\) 检验）：对
\(H_0: \boldsymbol{\mu} = \boldsymbol{\mu}_0\)（多元均值检验），

\[T^2 = n(\bar{\mathbf{X}} - \boldsymbol{\mu}_0)^\top \mathbf{S}^{-1} (\bar{\mathbf{X}} - \boldsymbol{\mu}_0)\]

在 \(H_0\) 下，\(\frac{n-p}{p(n-1)} T^2 \sim F_{p, n-p}\)。\(T^2\) 是
\(t\) 检验的多元推广。

\subsubsection{113.2 主成分分析}\label{ux4e3bux6210ux5206ux5206ux6790}

\textbf{定义 113.2}（主成分分析，PCA）：PCA 寻找数据方差最大的方向。设
\(\mathbf{X}\) 的协方差矩阵为 \(\boldsymbol{\Sigma}\)（或样本协方差矩阵
\(\mathbf{S}\)）。

\textbf{定理 113.3}（PCA 的推导）：设 \(\boldsymbol{\Sigma}\)
的特征分解为
\(\boldsymbol{\Sigma} = \mathbf{V}\boldsymbol{\Lambda}\mathbf{V}^\top\)，其中
\(\boldsymbol{\Lambda} = \operatorname{diag}(\lambda_1, \ldots, \lambda_p)\)，\(\lambda_1 \geq \cdots \geq \lambda_p \geq 0\)。则第
\(k\) 个\textbf{主成分}为
\(Y_k = \mathbf{v}_k^\top \mathbf{X}\)（\(\mathbf{v}_k\) 是
\(\lambda_k\) 对应的特征向量）。

\emph{证明}：\(\max_{\|\mathbf{a}\|=1} \operatorname{Var}(\mathbf{a}^\top \mathbf{X}) = \max_{\|\mathbf{a}\|=1} \mathbf{a}^\top \boldsymbol{\Sigma} \mathbf{a} = \lambda_1\)（由
Rayleigh 商），达到于 \(\mathbf{a} = \mathbf{v}_1\)。第 \(k\)
个主成分在与前 \(k-1\) 个正交的约束下最大化方差。∎

\textbf{命题 113.2}（PCA 的性质）： 1.
\(\operatorname{Var}(Y_k) = \lambda_k\) 2.
\(\sum_{k=1}^p \lambda_k = \operatorname{tr}(\boldsymbol{\Sigma})\)（总方差）
3. 前 \(k\) 个主成分解释的方差比例为
\(\sum_{i=1}^k \lambda_i / \sum_{i=1}^p \lambda_i\)

\subsubsection{113.3 判别分析}\label{ux5224ux522bux5206ux6790}

\textbf{定义 113.3}（线性判别分析，LDA）：假设各类服从
\(\mathcal{N}_p(\boldsymbol{\mu}_k, \boldsymbol{\Sigma})\)（等协方差矩阵）。LDA
将 \(\mathbf{x}\) 分类到使后验概率最大的类：

\[\delta_k(\mathbf{x}) = \mathbf{x}^\top \boldsymbol{\Sigma}^{-1} \boldsymbol{\mu}_k - \frac{1}{2} \boldsymbol{\mu}_k^\top \boldsymbol{\Sigma}^{-1} \boldsymbol{\mu}_k + \log \pi_k\]

其中 \(\pi_k\) 是第 \(k\)
类的先验概率。分类规则：\(\hat{k} = \arg\max_k \delta_k(\mathbf{x})\)。

\textbf{定义
113.4}（二次判别分析，QDA）：若各类协方差矩阵不同（\(\boldsymbol{\Sigma}_k\)），判别函数为

\[\delta_k(\mathbf{x}) = -\frac{1}{2} \log |\boldsymbol{\Sigma}_k| - \frac{1}{2} (\mathbf{x} - \boldsymbol{\mu}_k)^\top \boldsymbol{\Sigma}_k^{-1} (\mathbf{x} - \boldsymbol{\mu}_k) + \log \pi_k\]

\textbf{定理 113.4}（Fisher 判别分析）：Fisher 线性判别寻找投影方向
\(\mathbf{a}\) 最大化组间方差与组内方差之比：

\[\max_{\mathbf{a}} \frac{\mathbf{a}^\top \mathbf{B} \mathbf{a}}{\mathbf{a}^\top \mathbf{W} \mathbf{a}}\]

其中 \(\mathbf{B}\) 是组间散布矩阵，\(\mathbf{W}\)
是组内散布矩阵。\(\mathbf{a}\) 是 \(\mathbf{W}^{-1}\mathbf{B}\)
的最大特征值对应的特征向量。这与 LDA（等协方差时）等价。

\bigskip\hrule\bigskip

\subsection{第114章：信息几何}\label{ux7b2c114ux7ae0ux4fe1ux606fux51e0ux4f55}

信息几何由Amari和Chentsov在1980年代系统建立，将参数化概率分布族\(\mathcal{M} = \{p(x \mid \theta) : \theta \in \Theta\}\)视为统计流形，赋予Fisher信息度量\(g_{ij}(\theta) = \mathbb{E}[\partial_i \ell \cdot \partial_j \ell]\)（其中\(\ell = \log p\)）的Riemann几何结构。这一度量在参数变换下协变且由Chentsov定理唯一确定。进一步引入\(\alpha\)-联络\(\Gamma_{ijk}^{(\alpha)} = \mathbb{E}[(\partial_i\partial_j\ell + \frac{1-\alpha}{2}\partial_i\ell \cdot \partial_j\ell)\partial_k\ell]\)，\(\alpha = \pm 1\)分别对应指数联络和混合联络，两者关于Fisher度量对偶。Amari-Chentsov张量\(T_{ijk} = \mathbb{E}[\partial_i\ell \cdot \partial_j\ell \cdot \partial_k\ell]\)度量了统计流形的非平坦性。信息几何在机器学习中为自然梯度下降和EM算法的几何解释提供了理论基础，在神经网络的信息瓶颈原理和深度学习的流形学习中也有重要应用。

\bigskip\hrule\bigskip

\subsection{第115章：试验设计}\label{ux7b2c115ux7ae0ux8bd5ux9a8cux8bbeux8ba1}

试验设计（DOE）由Fisher在1920年代创立，其核心是Fisher三原则：(1)随机化------将实验单元随机分配以消除系统偏差；(2)重复------通过多次观测估计实验误差；(3)区组化------将同质单元分组以控制已知变异源。拉丁方设计通过行列双重区组控制两个干扰因素，与Greco-Latin方推广到多个因素。因子设计系统考察多个因素的主效应和交互效应，\(2^k\)因子设计是最基本的类型；当实验次数不可行时，采用分数因子设计以\(2^{k-p}\)次实验估计主要效应。响应面方法（RSM）由Box和Wilson于1951年提出，通过拟合二次响应曲面模型\(y = \beta_0 + \sum \beta_i x_i + \sum \beta_{ij}x_i x_j + \sum \beta_{ii}x_i^2 + \varepsilon\)，使用最速上升法寻优，是工业过程优化的核心工具。

\bigskip\hrule\bigskip

\subsection{第116章：抽样方法与调查设计（Sampling
Methods）}\label{ux7b2c116ux7ae0ux62bdux6837ux65b9ux6cd5ux4e0eux8c03ux67e5ux8bbeux8ba1sampling-methods}

抽样理论是统计学从有限总体中选取子集进行推断的方法论。\textbf{基本抽样方案}：（1）\textbf{简单随机抽样（SRS）}------每个大小为
\(n\) 的子集等概率抽取，均值 \(\bar{y}\) 是总体均值 \(\bar{Y}\)
的无偏估计，方差为 \((1 - \frac{n}{N}) \frac{S^2}{n}\)（有限总体修正因子
\(1 - \frac{n}{N}\)）；（2）\textbf{分层抽样}------将总体分为 \(H\)
层，在各层内独立抽样，通过 Neyman 最优分配 \(n_h \propto N_h S_h\)
最小化估计量的方差；（3）\textbf{系统抽样}------按固定间隔选取，当排序与变量相关时可显著提高精度；（4）\textbf{整群抽样}------在难以获取总体完整名单时，随机选取群（cluster）后调查群内全部个体；（5）\textbf{多阶段抽样}------先抽初级单元，再在选中单元内抽次级单元，逐级构造。\textbf{比值估计与回归估计}利用辅助变量提高精度：当辅助变量
\(x\) 与目标变量 \(y\) 高度线性相关时，比值估计
\(\hat{Y}_R = \frac{\bar{y}}{\bar{x}} X\) 或回归估计
\(\hat{Y}_{\text{reg}} = \bar{y} + \hat{\beta}(X - \bar{x})\)
可在大样本下显著降低方差。\textbf{Horvitz-Thompson
估计量}：\(\hat{Y}_{HT} = \sum_{i \in \mathcal{S}} \frac{y_i}{\pi_i}\)（\(\pi_i\)
为个体 \(i\) 的人样概率）对所有不等概率抽样设计一致有效。

\bigskip\hrule\bigskip

\subsection{第117章：时间序列分析（Time Series
Analysis）}\label{ux7b2c117ux7ae0ux65f6ux95f4ux5e8fux5217ux5206ux6790time-series-analysis}

时间序列分析研究依时间顺序观测的随机过程
\(\{X_t\}_{t \in \mathbb{Z}}\)，其核心问题是建模、预测和推断动态依赖结构。

\subsubsection{117.1
线性时间序列模型}\label{ux7ebfux6027ux65f6ux95f4ux5e8fux5217ux6a21ux578b}

\textbf{定义 117.1}（平稳过程）：时间序列 \(\{X_t\}\)
称为\textbf{（弱）平稳}的，如果
\(\mathbb{E}[X_t] = \mu\)（常数），且自协方差函数
\(\gamma(h) = \operatorname{Cov}(X_t, X_{t+h})\) 仅依赖于时滞
\(h\)。严平稳要求任意有限维联合分布平移不变。

\textbf{定义 117.2}（白噪声 / White Noise）：过程 \(\{\varepsilon_t\}\)
称为\textbf{白噪声}，记作
\(\varepsilon_t \sim \text{WN}(0, \sigma^2)\)，若
\(\mathbb{E}[\varepsilon_t] = 0\)，\(\operatorname{Var}(\varepsilon_t) = \sigma^2\)，且
\(\operatorname{Cov}(\varepsilon_t, \varepsilon_s) = 0\)（\(t \neq s\)）。

\textbf{定义 117.3}（ARMA(p,q) 模型 / 自回归滑动平均）：

\[X_t = \sum_{i=1}^p \phi_i X_{t-i} + \varepsilon_t + \sum_{j=1}^q \theta_j \varepsilon_{t-j}\]

其中 \(\{\varepsilon_t\}\) 为白噪声。AR（自回归）部分的平稳性由特征方程
\(\Phi(z) = 1 - \phi_1 z - \cdots - \phi_p z^p = 0\)
的根在单位圆外的条件给出；MA（滑动平均）部分本身即平稳。ARMA
过程的\textbf{自相关函数}（ACF）在 MA 部分截尾、AR
部分拖尾，是模型识别的关键工具。

\textbf{定义 117.4}（ARIMA(p,d,q) 模型 / Box-Jenkins, 1970）：若
\(\{X_t\}\) 经 \(d\) 次差分后成为平稳的 ARMA 过程，即 \((1-B)^d X_t\)
服从 ARMA(p,q)（\(B\) 为后移算子 \(B X_t = X_{t-1}\)），则称 \(\{X_t\}\)
为 ARIMA(p,d,q) 过程。这是处理含趋势的非平稳时间序列的标准方法。

\textbf{定义 117.5}（季节性 ARIMA /
SARIMA）：\(\Phi_P(B^s) \phi_p(B) (1-B)^d (1-B^s)^D X_t = \Theta_Q(B^s) \theta_q(B) \varepsilon_t\)，其中
\(s\) 为季节周期（如 \(s = 12\) 为月数据），\(P, D, Q\) 为季节阶数。

\textbf{定义 117.6}（GARCH(p,q) 模型 / 广义自回归条件异方差，Engle 1982
/ Bollerslev 1986）：

\[\begin{aligned} X_t &= \sigma_t Z_t, \quad Z_t \sim \text{i.i.d.}(0,1) \\ \sigma_t^2 &= \omega + \sum_{i=1}^p \alpha_i X_{t-i}^2 + \sum_{j=1}^q \beta_j \sigma_{t-j}^2 \end{aligned}\]

GARCH 模型捕捉金融时间序列中波动率聚类（volatility
clustering）------大幅波动倾向于跟随大幅波动。Engle 因提出 ARCH
模型（1982）获 2003 年 Nobel 经济学奖。

\subsubsection{117.2
谱分析与频域方法}\label{ux8c31ux5206ux6790ux4e0eux9891ux57dfux65b9ux6cd5}

\textbf{定义 117.7}（谱密度 / Spectral
Density）：平稳时间序列的\textbf{谱密度}为自协方差函数的 Fourier 变换

\[f(\omega) = \frac{1}{2\pi} \sum_{h=-\infty}^\infty \gamma(h) e^{-i\omega h}, \quad \omega \in [-\pi, \pi]\]

其中 \(\gamma(h) = \operatorname{Cov}(X_t, X_{t+h})\)。反演公式为
\(\gamma(h) = \int_{-\pi}^\pi e^{i\omega h} f(\omega) d\omega\)。\(f(\omega) \geq 0\)
对 \(\omega\) 的积分等于
\(\operatorname{Var}(X_t)\)，刻画了方差在频率域的分布。

\textbf{定义 117.8}（周期图 / Periodogram）：\(n\)
个观测的\textbf{周期图}为

\[I_n(\omega) = \frac{1}{2\pi n} \left|\sum_{t=1}^n X_t e^{-i\omega t}\right|^2\]

周期图 \(I_n(\omega)\) 是谱密度 \(f(\omega)\)
的渐近无偏估计，但非相合（其方差在 \(n \to \infty\)
时不趋于零）。通过\textbf{窗口平滑}（如 Daniell 窗口、Bartlett
窗口）可得到相合估计 \(\hat{f}(\omega) = \sum W_h I_n(\omega_h)\)。

\textbf{定理 117.1}（Whittle 似然 / Whittle
Likelihood）：在频域中，Gaussian 平稳过程的负对数似然近似为

\[-\ell(\theta) \approx \frac{1}{2\pi} \int_{-\pi}^\pi \left( \log f_\theta(\omega) + \frac{I_n(\omega)}{f_\theta(\omega)} \right) d\omega\]

这是频域 ARMA 参数估计（Whittle 估计）的基础，可通过快速 Fourier
变换高效计算。

\bigskip\hrule\bigskip

\emph{本卷在初等统计（V6）和概率论
II（V20）的基础上，建立了数理统计的严格理论框架：估计理论（无偏性、相合性、Cramér-Rao
下界、MLE 的渐近最优性）为参数估计提供了理论基础；Neyman-Pearson
引理和似然比检验为假设检验提供了最优性准则；置信区间与假设检验的对偶关系揭示了统计推断的统一结构；贝叶斯方法通过先验分布融合先验信息，后验分布的渐近正态性（Bernstein-von
Mises 定理）保证了与频率学派方法的一致性；线性模型（Gauss-Markov
定理、ANOVA）是应用最广泛的统计工具；多元分析（PCA、LDA）是处理高维数据的标准方法。}

\newpage

\section{卷二十二：数值分析}\label{ux5377ux4e8cux5341ux4e8cux6570ux503cux5206ux6790}

\begin{quote}
数值分析研究连续数学问题的离散近似算法，是科学计算的基础。本卷从浮点运算和误差分析出发，建立线性方程组、特征值问题、数值积分和常微分方程数值解的核心算法理论，重点讨论算法的稳定性、收敛性和误差估计。
\end{quote}

\bigskip\hrule\bigskip

\subsection{第114章：浮点运算与误差分析}\label{ux7b2c114ux7ae0ux6d6eux70b9ux8fd0ux7b97ux4e0eux8befux5deeux5206ux6790}

计算机中的实数运算是近似的，理解浮点运算的误差模型是数值分析的基础。本章介绍
IEEE 754 浮点标准、舍入误差和数值稳定性。

\subsubsection{114.1 浮点数系统}\label{ux6d6eux70b9ux6570ux7cfbux7edf}

\textbf{定义 114.1}（IEEE 754
双精度浮点数）：\textbf{双精度浮点数}（binary64）用 64 位表示：1 位符号
\(s\)、11 位指数 \(e\)（偏移 \(1023\)）、52 位尾数 \(m\)（隐含首位的
\(1\)）。数值为

\[x = (-1)^s \times (1.m) \times 2^{e - 1023}\]

\textbf{定义 114.2}（机器精度）：\textbf{机器精度}（machine
epsilon）\(\varepsilon_{\text{mach}}\) 是 \(1\) 与大于 \(1\)
的最小浮点数之间的差。双精度下
\(\varepsilon_{\text{mach}} = 2^{-52} \approx 2.22 \times 10^{-16}\)。

\textbf{定义 114.3}（舍入）：浮点运算的 IEEE
标准舍入模式（默认：舍入到最近偶数）。记 \(\operatorname{fl}(x)\) 为
\(x\) 的浮点表示，则

\[\operatorname{fl}(x) = x(1 + \delta), \quad |\delta| \leq \varepsilon_{\text{mach}}\]

\textbf{定义 114.4}（浮点运算模型）：对基本运算
\(\circ \in \{+, -, \times, \div\}\)，

\[\operatorname{fl}(x \circ y) = (x \circ y)(1 + \delta), \quad |\delta| \leq \varepsilon_{\text{mach}}\]

\subsubsection{114.2 误差分析}\label{ux8befux5deeux5206ux6790}

\textbf{定义 114.5}（绝对误差与相对误差）：设 \(\tilde{x}\) 是 \(x\)
的近似值。 - \textbf{绝对误差}：\(|\tilde{x} - x|\) -
\textbf{相对误差}：\(|\tilde{x} - x| / |x|\)（\(x \neq 0\)）

\textbf{定义 114.6}（前向误差与后向误差）： - \textbf{前向误差}：输出
\(\tilde{y}\) 与真值 \(y = f(x)\) 的差 \(|\tilde{y} - y|\) -
\textbf{后向误差}：寻找 \(\Delta x\) 使得
\(\tilde{y} = f(x + \Delta x)\)，后向误差为 \(|\Delta x|\)

一个算法称为\textbf{后向稳定}的，如果后向误差为
\(O(\varepsilon_{\text{mach}})\)。

\textbf{定义 114.7}（条件数）：问题 \(f\) 在 \(x\)
处的\textbf{（相对）条件数}为

\[\kappa_f(x) = \lim_{\varepsilon \to 0} \sup_{\|\Delta x\| \leq \varepsilon} \frac{\|f(x + \Delta x) - f(x)\| / \|f(x)\|}{\|\Delta x\| / \|x\|}\]

若 \(\kappa_f(x)\)
很大，则问题是\textbf{病态}的（ill-conditioned）。前向误差 \(\lesssim\)
条件数 \(\times\) 后向误差。

\textbf{例 114.1}（减法相消，catastrophic cancellation）：计算
\(f(x) = \sqrt{x+1} - \sqrt{x}\)（\(x \gg 1\)）。直接计算会导致有效数字严重丢失。重写为
\(f(x) = 1/(\sqrt{x+1} + \sqrt{x})\) 可避免此问题。

\textbf{定义
114.8}（数值稳定性）：算法是\textbf{数值稳定}的，如果输出对输入微小扰动不敏感。等价地，稳定的算法产生的结果接近精确问题的近似解（后向稳定）。

\bigskip\hrule\bigskip

\subsection{第115章：线性方程组数值解}\label{ux7b2c115ux7ae0ux7ebfux6027ux65b9ux7a0bux7ec4ux6570ux503cux89e3}

求解线性方程组 \(A\mathbf{x} = \mathbf{b}\)
是数值计算最基本的任务。本章讨论直接法（LU 分解、QR 分解）和迭代法。

\subsubsection{115.1 高斯消元与 LU
分解}\label{ux9ad8ux65afux6d88ux5143ux4e0e-lu-ux5206ux89e3}

\textbf{定理 115.1}（LU 分解）：设 \(A \in \mathbb{R}^{n \times n}\)
的所有顺序主子式非零。则存在唯一的单位下三角矩阵
\(L\)（\(L_{ii} = 1\)）和上三角矩阵 \(U\)，使得 \(A = LU\)。

\emph{证明}：归纳构造。设
\(A = \begin{pmatrix} a_{11} & \mathbf{w}^\top \\ \mathbf{v} & \tilde{A} \end{pmatrix}\)。令
\(L_1 = \begin{pmatrix} 1 & 0 \\ \mathbf{v}/a_{11} & I \end{pmatrix}\)，则
\(L_1^{-1} A = \begin{pmatrix} a_{11} & \mathbf{w}^\top \\ 0 & \tilde{A}_1 \end{pmatrix}\)。对
\(\tilde{A}_1\) 递归。∎

\textbf{定义 115.1}（部分选主元，PA =
LU）：为提高数值稳定性，在消元时选择列中绝对值最大的元素作为主元（行交换）。\(PA = LU\)，其中
\(P\) 是置换矩阵。

\textbf{算法 115.1}（求解 \(A\mathbf{x} = \mathbf{b}\) 的步骤）： 1.
计算 \(PA = LU\) 分解（\(O(n^3)\) 次运算） 2. 解
\(L\mathbf{y} = P\mathbf{b}\)（前代，\(O(n^2)\)） 3. 解
\(U\mathbf{x} = \mathbf{y}\)（回代，\(O(n^2)\)）

\textbf{定理 115.2}（Cholesky 分解）：若
\(A \in \mathbb{R}^{n \times n}\) 对称正定，则存在唯一的下三角矩阵
\(G\)（\(G_{ii} > 0\)）使得 \(A = GG^\top\)。计算量为
\(n^3/3 + O(n^2)\)（约为 LU 的一半）。

\subsubsection{115.2 QR
分解与最小二乘}\label{qr-ux5206ux89e3ux4e0eux6700ux5c0fux4e8cux4e58}

\textbf{定理 115.3}（QR 分解）：设
\(A \in \mathbb{R}^{m \times n}\)（\(m \geq n\)）满列秩。则存在正交矩阵
\(Q \in \mathbb{R}^{m \times m}\)（\(Q^\top Q = I\)）和上三角矩阵
\(R \in \mathbb{R}^{n \times n}\)，使得

\[A = Q \begin{pmatrix} R \\ 0 \end{pmatrix} = \tilde{Q} R\]

其中 \(\tilde{Q} \in \mathbb{R}^{m \times n}\) 是 \(Q\) 的前 \(n\)
列（\(\tilde{Q}^\top \tilde{Q} = I_n\)）。

\emph{证明}：使用 Householder 反射或 Givens 旋转构造。第 \(k\) 步用一个
Householder 反射将第 \(k\) 列对角线以下的元素消为零。∎

\textbf{定理 115.4}（最小二乘问题）：超定系统
\(A\mathbf{x} \approx \mathbf{b}\) 的最小二乘解为

\[\mathbf{x}^* = \arg\min_{\mathbf{x}} \|A\mathbf{x} - \mathbf{b}\|_2 = (A^\top A)^{-1} A^\top \mathbf{b}\]

使用 QR
分解：\(\mathbf{x}^* = R^{-1} \tilde{Q}^\top \mathbf{b}\)（数值更稳定）。

\subsubsection{115.3 迭代法}\label{ux8fedux4ee3ux6cd5}

\textbf{定义 115.2}（迭代法的一般形式）：将 \(A\) 分裂为
\(A = M - N\)（\(M\) 可逆）。迭代格式为

\[M\mathbf{x}^{(k+1)} = N\mathbf{x}^{(k)} + \mathbf{b}\]

即 \(\mathbf{x}^{(k+1)} = M^{-1}N\mathbf{x}^{(k)} + M^{-1}\mathbf{b}\)。

\textbf{定义 115.3}（经典迭代法）： - \textbf{Jacobi
迭代}：\(M = D\)（\(A\) 的对角部分），\(N = -(L + U)\) -
\textbf{Gauss-Seidel 迭代}：\(M = D + L\)，\(N = -U\)

\textbf{定理
115.5}（收敛性）：迭代法收敛（对任意初始值）当且仅当迭代矩阵 \(M^{-1}N\)
的谱半径 \(\rho(M^{-1}N) < 1\)。对对称正定矩阵，Gauss-Seidel 迭代收敛。

\textbf{定理 115.6}（共轭梯度法，CG）：对对称正定 \(A\)，共轭梯度法在
\(A\)-范数 \(\|\mathbf{x}\|_A = \sqrt{\mathbf{x}^\top A\mathbf{x}}\)
下最小化误差。理论上 CG 在 \(n\)
步内精确求解（在精确算术下），但实际中通常作为迭代法使用。收敛速率：

\[\|\mathbf{x}^{(k)} - \mathbf{x}^*\|_A \leq 2\left(\frac{\sqrt{\kappa} - 1}{\sqrt{\kappa} + 1}\right)^k \|\mathbf{x}^{(0)} - \mathbf{x}^*\|_A\]

其中 \(\kappa = \lambda_{\max}(A) / \lambda_{\min}(A)\) 是 \(A\)
的条件数。

\bigskip\hrule\bigskip

\subsection{第116章：特征值问题}\label{ux7b2c116ux7ae0ux7279ux5f81ux503cux95eeux9898}

特征值问题是数值线性代数的核心问题之一。本章讨论幂迭代、QR 算法和 Jacobi
方法。

\subsubsection{116.1
特征值的基本理论}\label{ux7279ux5f81ux503cux7684ux57faux672cux7406ux8bba}

\textbf{定义 116.1}（特征值问题）：求 \(\lambda \in \mathbb{C}\)
和非零向量 \(\mathbf{v} \in \mathbb{C}^n\) 使得
\(A\mathbf{v} = \lambda \mathbf{v}\)。

\textbf{定理 116.1}（Schur 分解）：对任意
\(A \in \mathbb{C}^{n \times n}\)，存在酉矩阵
\(Q\)（\(Q^*Q = I\)）和上三角矩阵 \(T\) 使得 \(A = QTQ^*\)。\(T\)
的对角元是 \(A\) 的特征值。

\textbf{推论 116.2}（实 Schur 分解）：若
\(A \in \mathbb{R}^{n \times n}\)，存在正交矩阵 \(Q\) 使得
\(Q^\top A Q\) 是拟上三角矩阵（对角块为 \(1 \times 1\) 或
\(2 \times 2\)）。

\textbf{定理 116.3}（谱定理，Hermitian 情形）：若 \(A\) 是 Hermitian
矩阵（\(A^* = A\)），则存在酉矩阵 \(Q\) 和对角矩阵 \(\Lambda\) 使得
\(A = Q\Lambda Q^*\)。特征值全为实数。

\subsubsection{116.2
幂迭代及其变体}\label{ux5e42ux8fedux4ee3ux53caux5176ux53d8ux4f53}

\textbf{定义 116.2}（幂迭代，Power Iteration）：对可对角化矩阵 \(A\)，设
\(|\lambda_1| > |\lambda_2| \geq \cdots \geq |\lambda_n|\)（\(\lambda_1\)
是主特征值）。幂迭代为

\[\mathbf{v}^{(k+1)} = A\mathbf{v}^{(k)} / \|A\mathbf{v}^{(k)}\|\]

\(\mathbf{v}^{(k)}\) 收敛到 \(\lambda_1\) 对应的特征向量，收敛速率为
\(|\lambda_2 / \lambda_1|^k\)。Rayleigh 商
\(\rho(\mathbf{v}) = \mathbf{v}^\top A\mathbf{v} / \mathbf{v}^\top \mathbf{v}\)
收敛到 \(\lambda_1\)。

\textbf{定义 116.3}（反迭代）：对特征值 \(\mu\) 附近的特征向量，反迭代为
\(\mathbf{v}^{(k+1)} = (A - \mu I)^{-1} \mathbf{v}^{(k)}\)（解线性系统）。收敛到离
\(\mu\) 最近的特征值对应的特征向量，收敛极快。

\textbf{定义 116.4}（Rayleigh 商迭代）：反迭代 + 动态更新
\(\mu\)：\(\mu_k = \rho(\mathbf{v}^{(k)})\)。三次收敛（在对称矩阵情形下）。

\subsubsection{116.3 QR 算法}\label{qr-ux7b97ux6cd5}

\textbf{算法 116.1}（QR 算法，基本形式）：\(A_0 = A\)。对
\(k = 0, 1, 2, \ldots\)： 1. 计算 \(A_k = Q_k R_k\)（QR 分解） 2.
\(A_{k+1} = R_k Q_k\)

\(A_k\) 收敛到拟上三角矩阵（实 Schur 形式），对角块给出特征值。

\textbf{定理 116.4}（QR 算法的收敛性）：若 \(A\) 的特征值满足
\(|\lambda_1| > |\lambda_2| > \cdots > |\lambda_n| > 0\)，则 \(A_k\)
的对角线以下元素以 \(|\lambda_{i+1} / \lambda_i|^k\) 的速率收敛到
\(0\)。

\textbf{算法 116.2}（实用 QR 算法）： 1. 先约化为上 Hessenberg
形式（\(O(n^3)\)，但后续每步仅 \(O(n^2)\)） 2.
引入位移（shift）加速收敛：\(A_k - \mu I = Q_k R_k\)，\(A_{k+1} = R_k Q_k + \mu I\)
3. 使用 Wilkinson 位移（右下角 \(2 \times 2\) 子矩阵的特征值中较近者）
4. 实用中通常 2-3 次迭代就收敛一个特征值

\subsubsection{116.4 Jacobi
方法（对称矩阵）}\label{jacobi-ux65b9ux6cd5ux5bf9ux79f0ux77e9ux9635}

\textbf{算法 116.3}（Jacobi 方法）：对对称矩阵 \(A\)，反复应用 Givens
旋转 \(J(p, q, \theta)\) 消去非对角元 \(a_{pq}\)（旋转角 \(\theta\) 满足
\(\tan 2\theta = 2a_{pq}/(a_{pp} - a_{qq})\)）。\(A \leftarrow J^\top A J\)。

\textbf{定理 116.5}（Jacobi 方法的收敛性）：Jacobi
方法二次收敛到对角矩阵（在每步消去最大非对角元的策略下）。非对角元素的
Frobenius 范数单调递减。

\bigskip\hrule\bigskip

\subsection{第117章：数值积分与微分}\label{ux7b2c117ux7ae0ux6570ux503cux79efux5206ux4e0eux5faeux5206}

数值积分（求积）和数值微分是逼近导数和定积分的算法。本章讨论
Newton-Cotes 公式、Gauss 求积和 Romberg 外推。

\subsubsection{117.1
数值积分基础}\label{ux6570ux503cux79efux5206ux57faux7840}

\textbf{定义 117.1}（求积法则）：\textbf{求积法则}
\(Q(f) = \sum_{i=0}^n w_i f(x_i)\) 近似积分
\(\int_a^b f(x) dx\)。\(x_i\) 称为\textbf{求积节点}，\(w_i\)
称为\textbf{求积权}。

\textbf{定义 117.2}（代数精度）：求积法则的\textbf{代数精度}为
\(m\)，如果它对所有次数 \(\leq m\) 的多项式精确，但对某个 \(m+1\)
次多项式不精确。

\textbf{定理 117.1}（Newton-Cotes 公式）：使用等距节点
\(x_i = a + ih\)（\(h = (b-a)/n\)）的插值型求积： -
\(n = 1\)（梯形法则）：\(\int_a^b f(x) dx \approx \frac{h}{2}[f(a) + f(b)]\)，误差
\(-\frac{h^3}{12} f''(\xi)\) - \(n = 2\)（Simpson
法则）：\(\int_a^b f(x) dx \approx \frac{h}{3}[f(a) + 4f((a+b)/2) + f(b)]\)，误差
\(-\frac{h^5}{90} f^{(4)}(\xi)\) - \(n = 3\)（Simpson 3/8 法则）：误差
\(-\frac{3h^5}{80} f^{(4)}(\xi)\)

\textbf{定义 117.3}（复合求积法则）：将 \([a, b]\) 分为 \(N\)
个子区间，每个子区间上应用低阶求积法则： -
复合梯形法则：\(T_N = \frac{h}{2}[f(a) + 2\sum_{i=1}^{N-1} f(x_i) + f(b)]\)，\(h = (b-a)/N\)，误差
\(O(h^2)\) - 复合 Simpson 法则：误差 \(O(h^4)\)

\subsubsection{117.2 Gauss 求积}\label{gauss-ux6c42ux79ef}

\textbf{定理 117.2}（Gauss 求积）：使用 \(n+1\)
个节点的求积法则的代数精度最多为 \(2n+1\)。达到此上界的求积法则称为
\textbf{Gauss 求积}。

\textbf{定理 117.3}（Gauss 求积的构造）：Gauss 求积的节点 \(x_i\) 是区间
\([a, b]\) 上关于权函数 \(w(x)\) 的正交多项式 \(P_{n+1}(x)\) 的零点。权
\(w_i\) 由插值条件确定。

\textbf{例 117.1}（Gauss-Legendre
求积）：\(w(x) = 1\)，\([a, b] = [-1, 1]\)。节点是 Legendre
多项式的零点。 -
\(n = 1\)：\(x_0 = -1/\sqrt{3}\)，\(x_1 = 1/\sqrt{3}\)，\(w_0 = w_1 = 1\)（代数精度
3） -
误差：\(\frac{2^{2n+3} ((n+1)!)^4}{(2n+3)((2n+2)!)^3} f^{(2n+2)}(\xi)\)

\subsubsection{117.3 Romberg
外推与自适应求积}\label{romberg-ux5916ux63a8ux4e0eux81eaux9002ux5e94ux6c42ux79ef}

\textbf{定理 117.4}（Euler-Maclaurin
公式）：复合梯形法则的误差有渐近展开

\[T_N = \int_a^b f(x) dx + c_2 h^2 + c_4 h^4 + c_6 h^6 + \cdots\]

\textbf{定理 117.5}（Romberg 积分）：利用 Euler-Maclaurin 展开，通过
Richardson 外推消除低阶误差项：

\[R_{k,0} = T_{2^k}, \quad R_{k,m} = \frac{4^m R_{k,m-1} - R_{k-1,m-1}}{4^m - 1}\]

\(R_{n,n}\) 以 \(O(h^{2n+2})\) 的精度逼近积分。

\textbf{定义
117.4}（自适应求积）：在函数变化剧烈处加密节点，变化平缓处稀疏节点。常用策略：比较两个不同阶的求积法则（如
Simpson 和 Boole）的差异，若差异超过阈值则细分区间。

\subsubsection{117.4 数值微分}\label{ux6570ux503cux5faeux5206}

\textbf{定义 117.5}（有限差分公式）： -
前向差分：\(f'(x) \approx \frac{f(x+h) - f(x)}{h}\)，误差
\(-\frac{h}{2} f''(\xi)\) -
中心差分：\(f'(x) \approx \frac{f(x+h) - f(x-h)}{2h}\)，误差
\(-\frac{h^2}{6} f'''(\xi)\) -
五点公式：\(f'(x) \approx \frac{-f(x+2h) + 8f(x+h) - 8f(x-h) + f(x-2h)}{12h}\)，误差
\(O(h^4)\)

\textbf{命题 117.1}（数值微分的不稳定性）：步长 \(h\)
不能无限减小。由于舍入误差 \(\sim \varepsilon/h\)（\(\varepsilon\)
是函数值误差），总误差 \(\approx C h^p + \varepsilon/h\)。最优
\(h \approx \varepsilon^{1/(p+1)}\)。

\bigskip\hrule\bigskip

\subsection{第118章：常微分方程数值解}\label{ux7b2c118ux7ae0ux5e38ux5faeux5206ux65b9ux7a0bux6570ux503cux89e3}

本章讨论 ODE 初值问题的数值解法，包括 Euler 方法、Runge-Kutta
方法和多步法。

\subsubsection{118.1 Euler 方法}\label{euler-ux65b9ux6cd5}

\textbf{定义 118.1}（ODE
初值问题）：\(\dot{y} = f(t, y)\)，\(y(t_0) = y_0\)。

\textbf{定义 118.2}（前向 Euler
方法）：\(y_{n+1} = y_n + h f(t_n, y_n)\)（\(h\)
为步长）。\textbf{局部截断误差} \(O(h^2)\)，\textbf{全局误差}
\(O(h)\)（一阶方法）。

\textbf{定义 118.3}（后向 Euler
方法）：\(y_{n+1} = y_n + h f(t_{n+1}, y_{n+1})\)（隐式方法，需解方程）。局部截断误差
\(O(h^2)\)，全局误差 \(O(h)\)（一阶方法）。

\textbf{定义
118.4}（梯形法则）：\(y_{n+1} = y_n + \frac{h}{2}[f(t_n, y_n) + f(t_{n+1}, y_{n+1})]\)（隐式，二阶方法）。

\subsubsection{118.2 Runge-Kutta 方法}\label{runge-kutta-ux65b9ux6cd5}

\textbf{定义 118.5}（显式 Runge-Kutta 方法，\(s\) 级）：\(s\) 级显式 RK
方法为

\[k_i = f\left(t_n + c_i h, y_n + h \sum_{j=1}^{i-1} a_{ij} k_j\right), \quad y_{n+1} = y_n + h \sum_{i=1}^s b_i k_i\]

\textbf{例 118.1}（经典 RK4，四阶）：最常用的 RK 方法：
\[\begin{aligned}
k_1 &= f(t_n, y_n) \\
k_2 &= f(t_n + h/2, y_n + h k_1/2) \\
k_3 &= f(t_n + h/2, y_n + h k_2/2) \\
k_4 &= f(t_n + h, y_n + h k_3) \\
y_{n+1} &= y_n + \frac{h}{6}(k_1 + 2k_2 + 2k_3 + k_4)
\end{aligned}\]

局部截断误差 \(O(h^5)\)，全局误差 \(O(h^4)\)。

\textbf{定理 118.1}（Butcher 障碍）：\(s\) 级显式 RK 方法的阶 \(p\) 满足
\(p \leq s\)（\(s = 1, 2, 3, 4\)），\(p \leq s-1\)（\(s = 5, 6, 7\)），\(p \leq s-2\)（\(s = 8, 9\)）等。

\emph{证明概要}（Butcher 1964-1965）：RK 方法的阶条件对应 Butcher
树（有根树）上的方程。对 \(s\) 级方法，有 \(s\) 个 \(b_i\) 和
\(s(s-1)/2\) 个 \(a_{ij}\)
参数，但需要满足的方程数由阶条件决定。将阶条件按 Butcher 树组织：设
\(t_k\) 为 \(k\)
阶有根树的数量（\(t_1=1, t_2=1, t_3=2, t_4=4, t_5=9, t_6=20, \dots\)）。对
\(s=1,2,3,4\)，未知数个数 \(\geq\) 方程数，故 \(p_{\max}=s\)。当 \(s=5\)
时，\(t_5=9\) 但可用的自由参数仅 7 个（扣除显式条件的
\(a_{ij}=0, j \geq i\)），故 \(p_{\max}=4\)。一般地，\(s \geq 5\)
时有根树数量的超指数增长造成每个 \(s\) 级方法的最大阶严格小于 \(s\)。∎

\subsubsection{118.3 多步法}\label{ux591aux6b65ux6cd5}

\textbf{定义 118.6}（线性多步法）：\(k\) 步线性多步法为

\[\sum_{j=0}^k \alpha_j y_{n+j} = h \sum_{j=0}^k \beta_j f(t_{n+j}, y_{n+j})\]

其中 \(\alpha_k = 1\)，\(|\alpha_0| + |\beta_0| > 0\)。若
\(\beta_k = 0\) 则为显式方法，否则为隐式方法。

\textbf{定义 118.7}（Adams 方法）： -
\textbf{Adams-Bashforth}（显式）：\(\alpha_k = 1\)，\(\alpha_{k-1} = -1\)，其余
\(\alpha_j = 0\) -
AB2：\(y_{n+1} = y_n + \frac{h}{2}[3f_n - f_{n-1}]\)（二阶） -
AB4：\(y_{n+1} = y_n + \frac{h}{24}[55f_n - 59f_{n-1} + 37f_{n-2} - 9f_{n-3}]\)（四阶）
- \textbf{Adams-Moulton}（隐式，通常作为校正步）：\(\beta_k \neq 0\) -
AM4：\(y_{n+1} = y_n + \frac{h}{24}[9f_{n+1} + 19f_n - 5f_{n-1} + f_{n-2}]\)（四阶）

\subsubsection{118.4
稳定性与刚性方程}\label{ux7a33ux5b9aux6027ux4e0eux521aux6027ux65b9ux7a0b}

\textbf{定义 118.8}（绝对稳定性区域）：对检验方程
\(\dot{y} = \lambda y\)（\(\Re(\lambda) < 0\)），数值方法产生
\(y_{n+1} = R(z) y_n\)，其中 \(z = h\lambda\)。\textbf{绝对稳定性区域}为
\(\{z \in \mathbb{C} : |R(z)| < 1\}\)。

\textbf{命题 118.1}（各方法的稳定性区域）： - 前向
Euler：\(|1 + z| < 1\)（圆盘，半径 \(1\)，中心 \(-1\)） - 后向
Euler：\(|1 - z| > 1\)（整个左半平面为稳定区域，A-稳定） -
梯形法则：\(|(1+z/2)/(1-z/2)| < 1 \iff \Re(z) < 0\)（A-稳定） -
RK4：与显式 Euler 类似但稍大（有限稳定区域）

\textbf{定义
118.9}（A-稳定性）：方法称为\textbf{A-稳定}的，如果其绝对稳定性区域包含整个左半平面
\(\{z : \Re(z) < 0\}\)。

\textbf{定义
118.10}（刚性方程）：若问题的不同分量以相差极大的速率变化（\(\max |\Re(\lambda)| / \min |\Re(\lambda)| \gg 1\)），则方程为\textbf{刚性}的。显式方法需极小步长，隐式方法（A-稳定）更适用。

\textbf{定义
118.11}（预测-校正方法）：\textbf{预测-校正}（Predictor-Corrector,
PECE）方法：用显式方法预测 \(y_{n+1}^P\)，再用隐式方法校正
\(y_{n+1}\)。例：AB4 预测 + AM4 校正（均为四阶）。

\bigskip\hrule\bigskip

\subsection{第119章：数值线性代数精讲}\label{ux7b2c119ux7ae0ux6570ux503cux7ebfux6027ux4ee3ux6570ux7cbeux8bb2}

数值线性代数是科学计算中最重要的核心算法集合，大多数大规模科学计算最终归结为线性方程组的求解或特征值问题。

\subsubsection{119.1 矩阵分解}\label{ux77e9ux9635ux5206ux89e3}

\textbf{定义 119.1}（LU 分解 / LU Decomposition）：将矩阵
\(A \in \mathbb{R}^{n \times n}\) 分解为下三角矩阵 \(L\) 与上三角矩阵
\(U\) 的乘积 \(A = LU\)。这是 Gauss
消去法的矩阵形式------消去过程的乘子按列存储于 \(L\) 中，消去结果存于
\(U\)。部分主元策略（partial pivoting）引入置换矩阵 \(P\) 得
\(PA = LU\)，避免小主元导致的数值不稳定。

\textbf{定理 119.1}（LU 分解的存在性）：若 \(A\)
的所有顺序主子式非零，则存在唯一的单位下三角矩阵 \(L\) 和上三角矩阵
\(U\) 使得 \(A = LU\)。

\textbf{定义 119.2}（Cholesky 分解）：若
\(A \in \mathbb{R}^{n \times n}\) 对称正定，则存在唯一的下三角矩阵
\(L\)（主对角元为正）使得 \(A = LL^T\)。运算量约为
\(\frac{1}{3}n^3\)（LU 分解的一半），且数值稳定性优良。

\textbf{定义 119.3}（QR 分解 / QR Decomposition）：将
\(A \in \mathbb{R}^{m \times n}\)（\(m \geq n\)）分解为 \(A = QR\)，其中
\(Q \in \mathbb{R}^{m \times m}\)
是正交矩阵（\(Q^T Q = I\)），\(R \in \mathbb{R}^{m \times n}\)
是上三角矩阵。构造方法包括：

\begin{itemize}
\tightlist
\item
  \textbf{Householder 反射}：\(H = I - 2 v v^T\)（\(v\)
  为单位向量）。通过 \(n-1\) 次反射将 \(A\) 的列依次消去对角线下方元素
\item
  \textbf{Givens 旋转}：\(G(i, j, \theta)\) 在 \((i,j)\)
  平面上旋转以消去单个元素，适用于稀疏矩阵和并行计算
\end{itemize}

\textbf{定理 119.2}（QR 分解求解最小二乘问题）：对于超定方程组
\(\min_x \|Ax - b\|_2\)（\(m \geq n\)），令
\(A = QR\)，\(Q^T b = \begin{pmatrix} c \\ d \end{pmatrix}\)（\(c \in \mathbb{R}^n\)），则解为
\(x = R^{-1} c\)，残差范数为 \(\|d\|_2\)。

\subsubsection{119.2 特征值算法}\label{ux7279ux5f81ux503cux7b97ux6cd5}

\textbf{定义 119.4}（QR 算法 / Francis, 1961-1962）：对
\(A_0 = A\)，迭代执行： 1. 对 \(A_k - \mu_k I\) 进行 QR
分解：\(A_k - \mu_k I = Q_k R_k\) 2.
重构：\(A_{k+1} = R_k Q_k + \mu_k I\)

适当地选择位移 \(\mu_k\)（Wilkinson 位移、Rayleigh 商位移），\(A_k\)
收敛到实 Schur 形式（拟上三角形，对角块为 \(1\times 1\) 或
\(2\times 2\)），特征值出现在对角线上。

\textbf{定义 119.5}（Lanczos 算法 / Lanczos Iteration）：对对称矩阵
\(A \in \mathbb{R}^{n \times n}\)，Lanczos 算法构造 Krylov 子空间
\(\mathcal{K}_k(A, v) = \operatorname{span}\{v, Av, A^2 v, \ldots, A^{k-1} v\}\)
的一组正交基 \(\{q_1, \ldots, q_k\}\)，满足三对角递推
\(A Q_k = Q_k T_k + \beta_k q_{k+1} e_k^T\)，其中 \(T_k\)
为对称三对角矩阵。\(T_k\) 的极值特征值在 \(k \ll n\) 时即收敛到 \(A\)
的极值特征值。

\textbf{定义 119.6}（Divide-and-Conquer
特征值算法）：将对称三对角矩阵分块为
\(T = \begin{pmatrix} T_1 & 0 \\ 0 & T_2 \end{pmatrix} + \rho v v^T\)，递归计算
\(T_1\)、\(T_2\) 的特征分解，通过求解 secular 方程
\(f(\lambda) = 1 + \rho \sum_{j=1}^n \frac{v_j^2}{d_j - \lambda} = 0\)
合并。总复杂度 \(O(n^3)\) 但常数小，LAPACK 中为默认选择。

\subsubsection{119.3 Krylov
子空间迭代法}\label{krylov-ux5b50ux7a7aux95f4ux8fedux4ee3ux6cd5}

\textbf{定义 119.7}（共轭梯度法 / CG / Conjugate
Gradient）：对对称正定的 \(Ax = b\)，共轭梯度法沿着
\(A\)-正交（\(p_i^T A p_j = 0\)）的搜索方向迭代：

\[\begin{aligned} r_0 &= b - A x_0, \quad p_0 = r_0 \\ \alpha_k &= \frac{r_k^T r_k}{p_k^T A p_k} \\ x_{k+1} &= x_k + \alpha_k p_k \\ r_{k+1} &= r_k - \alpha_k A p_k \\ \beta_k &= \frac{r_{k+1}^T r_{k+1}}{r_k^T r_k} \\ p_{k+1} &= r_{k+1} + \beta_k p_k \end{aligned}\]

\textbf{定理 119.3}（CG 收敛性）：在精确算术下，CG 在至多 \(n\)
步给出精确解。第 \(k\) 步的误差满足

\[\|x_k - x_*\|_A \leq 2 \left( \frac{\sqrt{\kappa(A)} - 1}{\sqrt{\kappa(A)} + 1} \right)^k \|x_0 - x_*\|_A\]

其中
\(\kappa(A) = \|A\|_2 \|A^{-1}\|_2 = \lambda_{\max} / \lambda_{\min}\)
为条件数，\(\|x\|_A = \sqrt{x^T A x}\)。

\textbf{定义 119.8}（预条件 / Preconditioning）：引入预条件子
\(M \approx A\)（\(M\) 易于求逆），求解等价系统
\(M^{-1} A x = M^{-1} b\)。预条件目的是降低等价系统的条件数。常用预条件子包括
Jacobi（对角）、不完全
Cholesky、多重网格。预条件共轭梯度法（PCG）是实际大规模求解的主流方法。

\textbf{定义 119.9}（GMRES / Generalized Minimal Residual）：对非对称
\(Ax = b\)，GMRES 在 Krylov 子空间 \(\mathcal{K}_k\) 中寻找使残差
\(\|b - A x_k\|_2\) 最小的 \(\tilde{x}_k\)。通过 Arnoldi 迭代构造
\(\mathcal{K}_k\) 的正交基，将最小二乘问题化为小型 Hessenberg
矩阵的最小二乘问题。存储需求随 \(k\)
线性增长，实际使用重启（restarted）GMRES(\(m\))。

\bigskip\hrule\bigskip

\emph{本卷从浮点运算和误差分析出发，建立了线性方程组、特征值问题、数值积分和常微分方程数值解的核心算法体系。第114章（浮点运算与误差分析）覆盖IEEE
754标准、条件数和数值稳定性；第115章（线性方程组数值解）讨论LU分解、QR分解和迭代法（Jacobi、Gauss-Seidel、共轭梯度法）；第116章（特征值问题）系统介绍幂迭代、QR算法和Jacobi方法；第117章（数值积分与微分）涵盖Newton-Cotes公式、Gauss求积、Romberg外推和有限差分；第118章（常微分方程数值解）包含Runge-Kutta方法（含Butcher障碍定理）、线性多步法（Adams-Bashforth/Moulton）、稳定性理论和刚性方程的处理。数值分析作为连接连续数学与离散计算的桥梁，使微分方程、优化和概率模型能在计算机上高效求解，是整个科学计算和工程仿真的数学基础。共
5 章。}

\newpage

\section{卷二十三：李理论}\label{ux5377ux4e8cux5341ux4e09ux674eux7406ux8bba}

\begin{quote}
李理论是连续对称性的数学框架，由 Sophus Lie 在 19
世纪末创立。李群是带有光滑流形结构的群，李代数是其切空间上的线性结构。本卷从李群与李代数的基本关系出发，建立半单李代数的分类理论（根系、Dynkin
图），讨论紧李群的表示论（Peter-Weyl
定理、最高权理论），并介绍李群在几何中的应用。
\end{quote}

\bigskip\hrule\bigskip

\subsection{第119章：李群与李代数}\label{ux7b2c119ux7ae0ux674eux7fa4ux4e0eux674eux4ee3ux6570}

李群是微分流形和群的结合，李代数是李群在单位元处的切空间，带有 Lie
括号运算。本章建立李群和李代数的基本理论，以及指数映射。

\subsubsection{119.1 李群的定义}\label{ux674eux7fa4ux7684ux5b9aux4e49}

\textbf{定义 119.1}（李群）：\textbf{李群} \(G\)
是一个光滑流形（微分流形），同时是一个群，使得群运算（乘法
\(G \times G \to G\) 和取逆 \(G \to G\)）都是光滑映射。

\textbf{例 119.1}（基本李群）： - \(\mathbb{R}^n\)（加法群）是 Abel 李群
- \(\mathbb{T}^n = \mathbb{R}^n / \mathbb{Z}^n\)（\(n\) 维环面）是紧
Abel 李群 - \textbf{一般线性群}
\(\operatorname{GL}(n, \mathbb{R}) = \{A \in M_n(\mathbb{R}) : \det A \neq 0\}\)（\(n^2\)
维） - \textbf{特殊线性群}
\(\operatorname{SL}(n, \mathbb{R}) = \{A \in \operatorname{GL}(n, \mathbb{R}) : \det A = 1\}\)（\(n^2 - 1\)
维） - \textbf{正交群}
\(\operatorname{O}(n) = \{A \in \operatorname{GL}(n, \mathbb{R}) : A^\top A = I\}\)（\(n(n-1)/2\)
维） - \textbf{特殊正交群}
\(\operatorname{SO}(n) = \operatorname{O}(n) \cap \operatorname{SL}(n, \mathbb{R})\)（紧，连通）
- \textbf{酉群}
\(\operatorname{U}(n) = \{A \in \operatorname{GL}(n, \mathbb{C}) : A^* A = I\}\)（紧）
- \textbf{特殊酉群}
\(\operatorname{SU}(n) = \operatorname{U}(n) \cap \operatorname{SL}(n, \mathbb{C})\)

\textbf{定义 119.2}（左不变向量场）：\(G\) 上的向量场 \(X\)
称为\textbf{左不变}的，如果 \((L_g)_* X = X\) 对所有 \(g \in G\)，其中
\(L_g(h) = gh\) 是左平移。

\textbf{命题 119.1}：左不变向量场构成 \(G\) 的切空间 \(T_e G\)
的同构。左不变向量场的 Lie 括号仍为左不变，因此在 \(T_e G\) 上诱导了 Lie
括号。

\subsubsection{119.2
李代数的定义}\label{ux674eux4ee3ux6570ux7684ux5b9aux4e49}

\textbf{定义 119.3}（李代数）：（实或复）向量空间 \(\mathfrak{g}\)
连同双线性映射
\([\cdot, \cdot]: \mathfrak{g} \times \mathfrak{g} \to \mathfrak{g}\)（\textbf{Lie
括号}）称为\textbf{李代数}，如果满足： 1. 反称性：\([X, Y] = -[Y, X]\)
2. Jacobi 恒等式：\([X, [Y, Z]] + [Y, [Z, X]] + [Z, [X, Y]] = 0\)

\textbf{定义 119.4}（李群对应的李代数）：李群 \(G\) 的\textbf{李代数}
\(\mathfrak{g} = \operatorname{Lie}(G) = T_e G\)，带有从左不变向量场的
Lie 括号诱导的括号。

\textbf{例 119.2}（矩阵李群的李代数）： -
\(\mathfrak{gl}(n, \mathbb{R}) = M_n(\mathbb{R})\)，\([A, B] = AB - BA\)
-
\(\mathfrak{sl}(n, \mathbb{R}) = \{A \in M_n(\mathbb{R}) : \operatorname{tr} A = 0\}\)
-
\(\mathfrak{so}(n) = \{A \in M_n(\mathbb{R}) : A^\top = -A\}\)（反对称矩阵）
-
\(\mathfrak{su}(n) = \{A \in M_n(\mathbb{C}) : A^* = -A, \operatorname{tr} A = 0\}\)

\textbf{定理 119.1}（Lie 第三定理的局部版本）：每个有限维李代数
\(\mathfrak{g}\) 是某个李群 \(G\) 的李代数。\(G\)
在连通且单连通时唯一确定。

\emph{证明概要}（含 Ado 定理）：Ado 定理（1947）断言任意有限维 Lie 代数
\(\mathfrak{g}\) 可忠实表示为 \(\mathfrak{gl}(n)\) 的子代数。由 Lie
第三定理的局部版本，\(\mathfrak{gl}(n)\) 对应 Lie 群
\(\operatorname{GL}(n)\) 的邻域，作为子代数
\(\mathfrak{g} \subseteq \mathfrak{gl}(n)\) 对应
\(\operatorname{GL}(n)\) 的局部子群（由 CBH
公式）。取该局部子群生成的连通李群（或取单连通覆盖）即得 \(G\)。忠实性由
Ado 定理保证。∎

\subsubsection{119.3 指数映射}\label{ux6307ux6570ux6620ux5c04}

\textbf{定义 119.5}（指数映射）：李群 \(G\) 上的\textbf{指数映射}
\(\exp: \mathfrak{g} \to G\) 定义为 \(\exp(X) = \gamma_X(1)\)，其中
\(\gamma_X: \mathbb{R} \to G\) 是满足 \(\gamma_X(0) = e\) 和
\(\gamma_X'(t) = X_{\gamma_X(t)}\) 的唯一单参数子群。

\textbf{命题 119.2}（指数映射的性质）： 1. \(\exp(0) = e\) 2.
\(\exp((t+s)X) = \exp(tX) \exp(sX)\) 3.
\(\frac{d}{dt}\big|_{t=0} \exp(tX) = X\) 4. \(\exp\) 是单位元
\(e \in G\) 附近的微分同胚 5.
对矩阵李群，\(\exp(A) = \sum_{k=0}^\infty \frac{A^k}{k!}\)（矩阵指数）

\textbf{定理 119.2}（Campbell-Baker-Hausdorff 公式）：对充分小的
\(X, Y \in \mathfrak{g}\)，

\[\exp(X) \exp(Y) = \exp(X + Y + \frac{1}{2}[X, Y] + \frac{1}{12}[X, [X, Y]] - \frac{1}{12}[Y, [X, Y]] + \cdots)\]

CBH 公式表明李群在单位元附近的乘法完全由李代数决定。

\subsubsection{119.4
李群同态与李代数同态}\label{ux674eux7fa4ux540cux6001ux4e0eux674eux4ee3ux6570ux540cux6001}

\textbf{定理 119.3}（Lie 对应）：设 \(G, H\) 是李群，\(G\)
连通单连通。李群同态 \(\varphi: G \to H\) 与李代数同态
\(d\varphi: \mathfrak{g} \to \mathfrak{h}\) 一一对应。

\emph{范畴等价概要}：Lie 对应的核心是单连通李群范畴 \(\mathbf{SLieGrp}\)
与有限维 Lie 代数范畴 \(\mathbf{LieAlg}_{\text{fin}}\) 等价。函子
\(\operatorname{Lie}: \mathbf{SLieGrp} \to \mathbf{LieAlg}_{\text{fin}}\)
将李群映为其李代数，将同态映为其微分；逆函子由 Cartan 的整体化定理（或
Lie 第三定理的全局版本）给出。该等价将子李群对应子 Lie
代数，正规闭子群对应 Lie
代数理想，商群对应商代数。这证明了连续对称性完全由无穷小线性化决定。∎

\textbf{定义 119.6}（伴随表示）：李群 \(G\) 的\textbf{伴随表示}
\(\operatorname{Ad}: G \to \operatorname{GL}(\mathfrak{g})\) 定义为
\(\operatorname{Ad}(g) = d(I_g)_e\)，其中 \(I_g(h) = ghg^{-1}\)
是内自同构。其微分
\(\operatorname{ad}: \mathfrak{g} \to \mathfrak{gl}(\mathfrak{g})\) 为
\(\operatorname{ad}(X)(Y) = [X, Y]\)。

\textbf{定理
119.4}（李群-李代数对应）：连通李群的子群与子代数、理想与正规子群、商群与商代数之间有一一对应。单连通李群范畴等价于有限维李代数范畴。

\bigskip\hrule\bigskip

\subsection{第120章：半单李代数}\label{ux7b2c120ux7ae0ux534aux5355ux674eux4ee3ux6570}

半单李代数是李理论的核心研究对象，具有完全可约性和丰富的结构理论。本章讨论
Killing 型、Cartan 子代数、根系和 Dynkin 图分类。

\subsubsection{120.1 Killing
型与半单性}\label{killing-ux578bux4e0eux534aux5355ux6027}

\textbf{定义 120.1}（Killing 型）：李代数 \(\mathfrak{g}\) 的
\textbf{Killing 型}是对称双线性型

\[\kappa(X, Y) = \operatorname{tr}(\operatorname{ad}(X) \circ \operatorname{ad}(Y))\]

\textbf{定理 120.1}（Cartan 判别法）：\(\mathfrak{g}\) 是半单的当且仅当
Killing 型非退化。\(\mathfrak{g}\) 是可解的当且仅当
\(\kappa(\mathfrak{g}, [\mathfrak{g}, \mathfrak{g}]) = 0\)。

\textbf{定义 120.2}（半单李代数）：李代数 \(\mathfrak{g}\)
称为\textbf{半单}的，如果它没有非零 Abel 理想。等价地，\(\mathfrak{g}\)
可分解为单李代数的直和。

\textbf{详解 120.1}（Killing 型与 Cartan 判别准则）：Killing 型
\(\kappa\) 是对称双线性型，由伴随表示的迹定义。其矩阵元为
\(\kappa_{ij} = \sum_{k,l} c_{ik}^l c_{jl}^k\)，其中 \(c_{ij}^k\)
是结构常数。Cartan 第一判别准则：\(\mathfrak{g}\) 可解当且仅当
\(\kappa(\mathfrak{g}, [\mathfrak{g}, \mathfrak{g}]) = 0\)。Cartan
第二判别准则（核心）：\(\mathfrak{g}\) 半单当且仅当 Killing 型非退化（即
\(\det(\kappa_{ij}) \neq 0\)）。此判别法的威力在于将''半单性''这一结构性质转化为纯粹线性代数条件：只需计算
Killing 型矩阵是否满秩。非退化性等价于 \(\mathfrak{g}\) 没有非零 Abel
理想，因为可解理想的 Killing
型必然退化。半单李代数的进一步性质：\(\kappa\) 限制在 Cartan
子代数上仍非退化，从而诱导 \(\mathfrak{h}\) 与 \(\mathfrak{h}^*\)
的同构，为权与根的对偶理论奠定基础。

\textbf{例
120.1}：\(\mathfrak{sl}(n, \mathbb{C})\)（\(n \geq 2\)）、\(\mathfrak{so}(n, \mathbb{C})\)（\(n \geq 3\)，\(n \neq 4\)）、\(\mathfrak{sp}(2n, \mathbb{C})\)
是半单李代数。\(\mathfrak{gl}(n, \mathbb{C})\)
不是半单的（有一维中心）。

\textbf{定理 120.2}（Weyl
完全可约性定理）：半单李代数的每个有限维表示是完全可约的（可分解为不可约表示的直和）。

\subsubsection{120.2 Cartan
子代数与根系}\label{cartan-ux5b50ux4ee3ux6570ux4e0eux6839ux7cfb}

\textbf{定义 120.3}（Cartan 子代数）：半单李代数 \(\mathfrak{g}\) 的
\textbf{Cartan 子代数} \(\mathfrak{h}\) 是极大 Abel 子代数，且
\(\operatorname{ad}(H)\)（\(H \in \mathfrak{h}\)）可同时对角化。

\textbf{命题 120.1}：半单李代数 \(\mathfrak{g}\) 关于 \(\mathfrak{h}\)
的\textbf{根空间分解}为

\[\mathfrak{g} = \mathfrak{h} \oplus \bigoplus_{\alpha \in \Phi} \mathfrak{g}_\alpha\]

其中
\(\mathfrak{g}_\alpha = \{X \in \mathfrak{g} : [H, X] = \alpha(H) X \text{ 对所有 } H \in \mathfrak{h}\}\)，\(\alpha \in \mathfrak{h}^* \setminus \{0\}\)
称为\textbf{根}。\(\Phi \subseteq \mathfrak{h}^*\) 是根系。

\textbf{命题 120.2}（根系的性质）： 1. \(\Phi\) 有限，张成
\(\mathfrak{h}^*\) 2. 若 \(\alpha \in \Phi\)，则
\(-\alpha \in \Phi\)，且
\(k\alpha \in \Phi\)（\(k \in \mathbb{Z}\)）仅当 \(k = \pm 1\) 3.
\(\dim \mathfrak{g}_\alpha = 1\) 4. 对
\(\alpha, \beta \in \Phi\)，\(\beta - \frac{2\kappa(\beta, \alpha)}{\kappa(\alpha, \alpha)} \alpha \in \Phi\)（根系在反射下封闭）

\textbf{定义 120.4}（Weyl 群）：\textbf{Weyl 群} \(W\) 由根反射
\(s_\alpha(\beta) = \beta - \frac{2\kappa(\beta, \alpha)}{\kappa(\alpha, \alpha)} \alpha\)
生成。\(W\) 是有限群，作用在 \(\mathfrak{h}^*\) 上。

\subsubsection{120.3 单根与 Dynkin
图}\label{ux5355ux6839ux4e0e-dynkin-ux56fe}

\textbf{定义 120.5}（正根与单根）：选取 \(\mathfrak{h}^*\)
中的一个超平面使其不包含任何根，定义\textbf{正根} \(\Phi^+\)
为超平面一侧的根。\textbf{单根} \(\Delta \subseteq \Phi^+\)
是正根中不能分解为两个正根之和的根。

\textbf{命题 120.3}：单根 \(\Delta\) 是 \(\mathfrak{h}^*\)
的一组基。每个正根是单根的非负整数线性组合。

\textbf{定义 120.6}（Cartan
矩阵）：\(A_{ij} = \frac{2\kappa(\alpha_i, \alpha_j)}{\kappa(\alpha_i, \alpha_i)}\)（\(\alpha_i, \alpha_j \in \Delta\)）。\(A_{ij} \in \{0, -1, -2, -3\}\)（\(i \neq j\)）。

\textbf{定义 120.7}（Dynkin 图）：\textbf{Dynkin
图}以单根为节点，\(A_{ij} A_{ji}\) 条边连接节点 \(i\) 和 \(j\)。若
\(|\alpha_i| > |\alpha_j|\)，画箭头指向较短根。

\textbf{定理 120.3}（Cartan-Killing
分类）：复半单李代数（等价地，单李代数）由 Dynkin 图完全分类。连通
Dynkin 图有：

\begin{itemize}
\tightlist
\item
  \(A_n\)（\(n \geq 1\)）：\(\mathfrak{sl}(n+1, \mathbb{C})\)
\item
  \(B_n\)（\(n \geq 2\)）：\(\mathfrak{so}(2n+1, \mathbb{C})\)
\item
  \(C_n\)（\(n \geq 3\)）：\(\mathfrak{sp}(2n, \mathbb{C})\)
\item
  \(D_n\)（\(n \geq 4\)）：\(\mathfrak{so}(2n, \mathbb{C})\)
\item
  例外李代数：\(G_2\)（14 维）、\(F_4\)（52 维）、\(E_6\)（78
  维）、\(E_7\)（133 维）、\(E_8\)（248 维）
\end{itemize}

这是数学中最优美的分类定理之一，与 Lie 和 Killing 在 19 世纪末的工作以及
Cartan 的完善。

\bigskip\hrule\bigskip

\subsection{第121章：紧李群的表示}\label{ux7b2c121ux7ae0ux7d27ux674eux7fa4ux7684ux8868ux793a}

紧李群具有优美的表示论，其所有不可约表示都是有限维的，且完全由最高权分类。本章讨论
Peter-Weyl 定理和最高权理论。

\subsubsection{121.1 紧李群上的 Haar
测度}\label{ux7d27ux674eux7fa4ux4e0aux7684-haar-ux6d4bux5ea6}

\textbf{定理 121.1}（Haar 测度）：紧李群 \(G\)
上存在唯一的（左）不变概率测度 \(\mu\)（Haar 测度），满足
\(\mu(gA) = \mu(A)\) 对所有 \(g \in G\) 和可测集 \(A\)。紧李群上 Haar
测度是双不变的。

\textbf{命题 121.1}（Weyl 积分公式）：对 \(G\) 上的类函数
\(f\)（\(f(gxg^{-1}) = f(x)\)），

\[\int_G f(g) dg = \frac{1}{|W|} \int_T f(t) \prod_{\alpha \in \Phi^+} |e^{\alpha}(t) - e^{-\alpha}(t)|^2 dt\]

其中 \(T\) 是极大环面，\(W\) 是 Weyl 群。

\subsubsection{121.2 Peter-Weyl 定理}\label{peter-weyl-ux5b9aux7406}

\textbf{定义 121.1}（表示）：李群 \(G\)
的（有限维）\textbf{表示}是李群同态
\(\pi: G \to \operatorname{GL}(V)\)（\(V\) 是有限维向量空间）。

\textbf{定理 121.2}（Peter-Weyl 定理）：紧李群 \(G\)
的矩阵系数（\(\pi_{ij}\) 跑遍所有不可约表示）在 \(L^2(G)\) 中稠密。即

\[L^2(G) \cong \widehat{\bigoplus}_{\pi \in \widehat{G}} V_\pi \otimes V_\pi^*\]

（Hilbert 空间直和）。每个不可约表示在 \(L^2(G)\)
中出现的重数等于其维数。

\textbf{推论
121.3}：紧李群的所有不可约表示都是有限维的。紧李群有可数多个不可约表示（同构类）。

\subsubsection{121.3 最高权理论}\label{ux6700ux9ad8ux6743ux7406ux8bba}

\textbf{定义 121.2}（权）：设 \(\pi: G \to \operatorname{GL}(V)\) 是
\(G\) 的表示。\(\lambda \in \mathfrak{h}^*\)
称为表示的\textbf{权}，如果存在非零 \(v \in V\) 使得
\(\pi(H)v = \lambda(H) v\) 对所有 \(H \in \mathfrak{h}\)。权空间
\(V_\lambda\) 是所有 \(\lambda\)-权向量的集合。

\textbf{定理 121.4}（最高权定理）：紧连通李群 \(G\)
的不可约表示与\textbf{支配整权}（dominated integral
weights）一一对应。每个不可约表示由唯一的\textbf{最高权} \(\lambda\)
确定（所有其他权 \(\mu\) 满足 \(\lambda - \mu\) 是正根的非负整数组合）。

\textbf{定义 121.3}（支配整权）：\(\lambda \in \mathfrak{h}^*\)
是\textbf{支配整权}，如果
\(\frac{2\kappa(\lambda, \alpha)}{\kappa(\alpha, \alpha)} \in \mathbb{Z}_{\geq 0}\)
对所有正根 \(\alpha\)。

\textbf{详解
121.1}（最高权理论的基本框架）：最高权理论建立了支配整权与不可约表示之间的一一对应，其结构如下。第一步：选取
Cartan 子代数 \(\mathfrak{h}\) 和正根系 \(\Phi^+\)，定义支配整权的集合
\(P_+ = \{\lambda \in \mathfrak{h}^* : \langle \lambda, \alpha^\vee \rangle \in \mathbb{Z}_{\geq 0}, \forall \alpha \in \Phi^+\}\)。第二步：对每个
\(\lambda \in P_+\)，构造 Verma 模
\(M(\lambda) = U(\mathfrak{g}) \otimes_{U(\mathfrak{b})} \mathbb{C}_\lambda\)（其中
\(\mathfrak{b} = \mathfrak{h} \oplus \bigoplus_{\alpha > 0} \mathfrak{g}_\alpha\)
是 Borel 子代数），这是一个最高权为 \(\lambda\)
的无穷维模。第三步：\(M(\lambda)\) 有唯一的极大真子模，商模
\(L(\lambda)\) 是唯一的最高权为 \(\lambda\) 的不可约模。第四步：Weyl
完全可约性定理保证了每个有限维表示分解为 \(L(\lambda)\)
的直和。此框架将李群表示论转化为权格的组合学，是现代表示论的标准范式。

\textbf{定理 121.5}（Weyl 维数公式）：最高权为 \(\lambda\)
的不可约表示的维数为

\[\dim V_\lambda = \prod_{\alpha \in \Phi^+} \frac{\kappa(\lambda + \rho, \alpha)}{\kappa(\rho, \alpha)}\]

其中
\(\rho = \frac{1}{2} \sum_{\alpha \in \Phi^+} \alpha\)（所有正根之和的一半）。

\textbf{定理 121.6}（Weyl 特征公式）：最高权 \(\lambda\)
的不可约表示的特征标为

\[\chi_\lambda = \frac{\sum_{w \in W} (-1)^{\ell(w)} e^{w(\lambda + \rho)}}{\sum_{w \in W} (-1)^{\ell(w)} e^{w(\rho)}}\]

其中 \(\ell(w)\) 是 \(w\) 的长度（表示为单反射的最小个数）。

\textbf{例 121.1}（\(\operatorname{SU}(2)\)
的表示）：\(\operatorname{SU}(2)\) 的不可约表示由非负半整数
\(j = 0, 1/2, 1, 3/2, \ldots\) 参数化，\(\dim V_j = 2j + 1\)。\(V_j\) 是
\(j\) 次齐次二元多项式空间（或自旋 \(j\) 表示）。

\textbf{例 121.2}（\(\operatorname{SU}(3)\)
的表示）：\(\operatorname{SU}(3)\) 的不可约表示由两个非负整数 \((p, q)\)
分类，\(\dim V_{(p,q)} = (p+1)(q+1)(p+q+2)/2\)。\((1, 0)\) 是基础表示
\(\mathbf{3}\)，\((0, 1)\) 是反基础表示 \(\bar{\mathbf{3}}\)。

\bigskip\hrule\bigskip

\subsection{第122章：李群在几何中的应用}\label{ux7b2c122ux7ae0ux674eux7fa4ux5728ux51e0ux4f55ux4e2dux7684ux5e94ux7528}

李群在几何中有深刻的应用，包括齐性空间、对称空间和主丛上的联络理论。

\subsubsection{122.1 齐性空间}\label{ux9f50ux6027ux7a7aux95f4}

\textbf{定义 122.1}（齐性空间）：李群 \(G\) 的\textbf{齐性空间}是 \(G\)
在流形 \(M\) 上的可迁作用（即 \(M = G/H\)，其中 \(H\) 是 \(G\)
的闭子群）。

\textbf{例 122.1}（经典齐性空间）： -
球面：\(S^n \cong \operatorname{SO}(n+1) / \operatorname{SO}(n)\) -
实射影空间：\(\mathbb{RP}^n \cong \operatorname{SO}(n+1) / \operatorname{O}(n)\)
-
复射影空间：\(\mathbb{CP}^n \cong \operatorname{SU}(n+1) / \operatorname{U}(n)\)
- Grassmann
流形：\(\operatorname{Gr}(k, n) \cong \operatorname{O}(n) / (\operatorname{O}(k) \times \operatorname{O}(n-k))\)
-
上半平面：\(\mathbb{H} \cong \operatorname{SL}(2, \mathbb{R}) / \operatorname{SO}(2)\)

\textbf{定理 122.1}（Maurer-Cartan 形式）：李群 \(G\) 上的
\textbf{Maurer-Cartan 形式} \(\omega \in \Omega^1(G, \mathfrak{g})\)
定义为 \(\omega_g(v) = (L_{g^{-1}})_* v\)。满足 Maurer-Cartan 方程：

\[d\omega + \frac{1}{2}[\omega, \omega] = 0\]

\subsubsection{122.2 对称空间}\label{ux5bf9ux79f0ux7a7aux95f4}

\textbf{定义 122.2}（对称空间）：\textbf{对称空间}是黎曼流形
\((M, g)\)，每点 \(p \in M\) 存在等距对合 \(\sigma_p: M \to M\) 以 \(p\)
为孤立不动点。

\textbf{定理 122.2}（Cartan 分类）：对称空间由
\((\mathfrak{g}, \mathfrak{h}, \sigma)\) 分类，其中 \(\mathfrak{g}\)
是李代数，\(\mathfrak{h}\) 是对合 \(\sigma\)
的固定点子代数。\(\mathfrak{g} = \mathfrak{h} \oplus \mathfrak{m}\)（\(\mathfrak{m}\)
是 \(\sigma\) 的 \(-1\) 特征空间），满足
\([\mathfrak{h}, \mathfrak{h}] \subseteq \mathfrak{h}\)，\([\mathfrak{h}, \mathfrak{m}] \subseteq \mathfrak{m}\)，\([\mathfrak{m}, \mathfrak{m}] \subseteq \mathfrak{h}\)。

\textbf{例 122.2}（对称空间的例子）： - 球面 \(S^n\)、双曲空间
\(\mathbb{H}^n\)、实/复/四元射影空间 - 紧李群本身（双不变量度） -
\(\operatorname{SL}(n, \mathbb{R}) / \operatorname{SO}(n)\)（正定对称矩阵空间）

\subsubsection{122.3 主丛与联络}\label{ux4e3bux4e1bux4e0eux8054ux7edc}

\textbf{定义 122.3}（主 \(G\)-丛）：流形 \(M\) 上的\textbf{主 \(G\)-丛}
\(P\) 是纤维丛 \(\pi: P \to M\)，纤维是李群 \(G\)，\(G\) 在 \(P\)
上自由右作用，局部平凡化与 \(G\) 作用兼容。

\textbf{定义 122.4}（联络）：\(P\) 上的\textbf{联络}是
\(\mathfrak{g}\)-值 1-形式
\(\omega \in \Omega^1(P, \mathfrak{g})\)，满足： 1.
\(\omega(\tilde{X}) = X\)（对 \(X \in \mathfrak{g}\)，\(\tilde{X}\) 是
\(X\) 生成的基本向量场） 2.
\(R_g^* \omega = \operatorname{Ad}(g^{-1}) \omega\)（等变性）

\textbf{定义 122.5}（曲率）：联络 \(\omega\) 的\textbf{曲率}为
\(\Omega = d\omega + \frac{1}{2}[\omega, \omega]\)（\(\mathfrak{g}\)-值
2-形式）。

\textbf{定理 122.3}（Cartan
结构方程）：在正交标架丛（或一般主丛）上，联络 \(\omega\) 和曲率
\(\Omega\) 满足 Cartan 结构方程：

\[d\omega = -\frac{1}{2}[\omega, \omega] + \Omega\]

（Maurer-Cartan 方程 \(\Omega = 0\) 是平坦联络的特例）。

\textbf{例 122.3}（Chern-Weil 理论）：对 \(G\)-不变多项式 \(P\) 在
\(\mathfrak{g}\) 上，\(P(\Omega)\) 是 \(M\) 上的闭形式，其 de Rham
上同调类不依赖于联络的选择。这是示性类（Chern 类、Pontryagin 类、Euler
类）的构造基础。

\textbf{定理 122.4}（Chern-Weil 同态）：\(G\)-不变多项式环 \(I(G)\) 到
\(M\) 的 de Rham 上同调环 \(H^*_{\text{dR}}(M)\)
之间存在同态。这给出了示性类的基本构造。

\bigskip\hrule\bigskip

\emph{本卷建立了李理论的核心框架：李群和李代数的基本关系（指数映射、CBH
公式、Lie
对应）将连续对称性转化为代数结构；半单李代数的分类理论（根系、Dynkin
图、Cartan-Killing 分类）是 19
世纪数学最伟大的成就之一；紧李群的表示论（Peter-Weyl
定理、最高权理论、Weyl
特征公式）为调和分析和量子力学提供了基础；齐性空间、对称空间和主丛上的联络理论将李群应用于微分几何。李理论是现代数学和物理学的核心语言，从基本粒子物理（规范理论）到数论（Langlands
纲领）都有深刻应用。}

\newpage

\section{卷二十四：调和分析}\label{ux5377ux4e8cux5341ux56dbux8c03ux548cux5206ux6790}

\begin{quote}
调和分析研究函数和算子的 Fourier 分解，是现代分析的核心分支。本卷从
Fourier 变换的 \(L^1\) 和 \(L^2\) 理论出发，建立 Schwartz
空间和缓增分布上的 Fourier 变换，然后深入到 Calderón-Zygmund
奇异积分理论、Hardy-Littlewood 极大函数和 Littlewood-Paley
理论，为现代调和分析和 PDE 研究提供核心工具。
\end{quote}

\bigskip\hrule\bigskip

\subsection{第123章：Fourier
变换}\label{ux7b2c123ux7ae0fourier-ux53d8ux6362}

Fourier
变换是调和分析最基本的工具，它将函数从物理空间映射到频率空间。本章从
\(L^1\) 和 \(L^2\) 理论出发，建立缓增分布上的 Fourier 变换。

\subsubsection{\texorpdfstring{123.1 \(L^1\)
理论}{123.1 L\^{}1 理论}}\label{l1-ux7406ux8bba}

\textbf{定义 123.1}（Fourier 变换）：\(f \in L^1(\mathbb{R}^n)\) 的
\textbf{Fourier 变换}为

\[\hat{f}(\xi) = \mathcal{F}f(\xi) = \int_{\mathbb{R}^n} f(x) e^{-2\pi i x \cdot \xi} dx\]

\textbf{Fourier
逆变换}（形式地）：\(f(x) = \int_{\mathbb{R}^n} \hat{f}(\xi) e^{2\pi i x \cdot \xi} d\xi\)（在
\(L^2\) 或适当条件下成立）。

\textbf{命题 123.1}（\(L^1\) 理论的基本性质）： 1.
\(\|\hat{f}\|_\infty \leq \|f\|_1\)（\(\hat{f}\) 是有界连续函数） 2.
\(\widehat{f * g} = \hat{f} \hat{g}\)（卷积的 Fourier 变换是乘积） 3.
\(\widehat{f(\cdot - h)}(\xi) = e^{-2\pi i h \cdot \xi} \hat{f}(\xi)\)（平移）
4.
\(\widehat{e^{2\pi i x \cdot \eta} f(x)}(\xi) = \hat{f}(\xi - \eta)\)（调制）
5.
\(\widehat{f(\lambda x)}(\xi) = \lambda^{-n} \hat{f}(\xi/\lambda)\)（伸缩）
6. Riemann-Lebesgue 引理：\(\hat{f}(\xi) \to 0\)（\(|\xi| \to \infty\)）

\textbf{定理 123.1}（Fourier 反演公式，\(L^1\) 情形）：若
\(f, \hat{f} \in L^1(\mathbb{R}^n)\)，则

\[f(x) = \int_{\mathbb{R}^n} \hat{f}(\xi) e^{2\pi i x \cdot \xi} d\xi \quad \text{几乎处处}\]

\emph{证明}：使用 Gauss
核近似恒等元。\(f * \varphi_\varepsilon \to f\)（\(\varepsilon \to 0\)），其中
\(\varphi_\varepsilon(x) = \varepsilon^{-n} e^{-\pi |x|^2/\varepsilon^2}\)。\(\hat{\varphi}_\varepsilon(\xi) = e^{-\pi \varepsilon^2 |\xi|^2}\)。由
Fourier 反演对 Gauss 函数成立，取极限。∎

\subsubsection{\texorpdfstring{123.2 \(L^2\) 理论与 Plancherel
定理}{123.2 L\^{}2 理论与 Plancherel 定理}}\label{l2-ux7406ux8bbaux4e0e-plancherel-ux5b9aux7406}

\textbf{定理 123.2}（Plancherel 定理）：Fourier 变换可以唯一延拓为
\(L^2(\mathbb{R}^n)\)
上的酉算子（在归一化后）：\(\|\hat{f}\|_2 = \|f\|_2\)。即

\[\int_{\mathbb{R}^n} |f(x)|^2 dx = \int_{\mathbb{R}^n} |\hat{f}(\xi)|^2 d\xi\]

\textbf{推论 123.3}（Parseval
恒等式）：\(\langle f, g \rangle = \langle \hat{f}, \hat{g} \rangle\)，即
\(\int f \bar{g} = \int \hat{f} \bar{\hat{g}}\)。

\textbf{定理 123.4}（Hausdorff-Young 不等式）：对
\(1 \leq p \leq 2\)，\(\|\hat{f}\|_{p'} \leq \|f\|_p\)，其中
\(1/p + 1/p' = 1\)。即 Fourier 变换是 \(L^p \to L^{p'}\)
的有界算子（\(p \in [1, 2]\)）。

\subsubsection{123.3 Schwartz
空间与缓增分布}\label{schwartz-ux7a7aux95f4ux4e0eux7f13ux589eux5206ux5e03}

\textbf{定义 123.2}（Schwartz 空间）：\textbf{Schwartz 空间}
\(\mathcal{S}(\mathbb{R}^n)\) 是所有光滑函数 \(\varphi\)
的集合，满足对所有多重指标 \(\alpha, \beta\)，

\[\sup_{x \in \mathbb{R}^n} |x^\alpha D^\beta \varphi(x)| < \infty\]

即 \(\varphi\) 及其所有导数比任意多项式下降更快。

\textbf{定理 123.5}：Fourier 变换是 \(\mathcal{S}(\mathbb{R}^n)\)
上的同构（双射连续线性映射）。\(\mathcal{S}\) 在 Fourier
变换下的自同构性是调和分析的核心。

\textbf{定义 123.3}（缓增分布）：\(\mathcal{S}(\mathbb{R}^n)\)
的对偶空间 \(\mathcal{S}'(\mathbb{R}^n)\)
称为\textbf{缓增分布}空间。缓增分布 \(u\) 的 \textbf{Fourier 变换}
\(\hat{u}\) 由对偶定义：

\[\langle \hat{u}, \varphi \rangle = \langle u, \hat{\varphi} \rangle \quad (\varphi \in \mathcal{S})\]

\textbf{例 123.1}（缓增分布的 Fourier 变换）： -
\(\hat{\delta}_0 = 1\)（常数函数） - \(\hat{1} = \delta_0\)（Dirac
\(\delta\) 分布） -
\(\widehat{\operatorname{p.v.}(1/x)} = -\pi i \operatorname{sgn}(\xi)\)（Hilbert
变换的 Fourier 乘子） -
\(\widehat{\operatorname{sgn}(x)} = \frac{1}{\pi i} \operatorname{p.v.}(1/\xi)\)

\subsubsection{123.4 Poisson
求和公式}\label{poisson-ux6c42ux548cux516cux5f0f}

\textbf{定理 123.6}（Poisson 求和公式）：设 \(f\) 是 Schwartz 函数。则

\[\sum_{n \in \mathbb{Z}^n} f(n) = \sum_{m \in \mathbb{Z}^n} \hat{f}(m)\]

更一般地，\(\sum_{n \in \mathbb{Z}^n} f(x + n) = \sum_{m \in \mathbb{Z}^n} \hat{f}(m) e^{2\pi i m \cdot x}\)（分布意义下）。

\emph{证明}：定义
\(F(x) = \sum_{n \in \mathbb{Z}^n} f(x + n)\)（\(\mathbb{Z}^n\)-周期函数）。\(F\)
的 Fourier 系数为 \(\hat{f}(m)\)，由 Fourier 级数理论即得。∎

\bigskip\hrule\bigskip

\subsection{第124章：奇异积分}\label{ux7b2c124ux7ae0ux5947ux5f02ux79efux5206}

奇异积分是 Calderón-Zygmund 在 20 世纪 50 年代开创的理论，它将 Hilbert
变换（一维）推广到高维，并建立了 \(L^p\) 有界性的一般框架。

\subsubsection{124.1 Hilbert 变换}\label{hilbert-ux53d8ux6362}

\textbf{定义 124.1}（Hilbert 变换）：\(\mathbb{R}\) 上的 \textbf{Hilbert
变换} \(H\) 定义为

\[Hf(x) = \frac{1}{\pi} \operatorname{p.v.} \int_{-\infty}^\infty \frac{f(y)}{x - y} dy = \frac{1}{\pi} \lim_{\varepsilon \to 0} \int_{|x-y| > \varepsilon} \frac{f(y)}{x - y} dy\]

\textbf{命题 124.1}（Hilbert 变换的性质）： 1. \(H\) 是
\(L^2(\mathbb{R})\) 上的酉算子：\(H^* = -H\)，\(H^2 = -I\) 2. Fourier
乘子：\(\widehat{Hf}(\xi) = -i \operatorname{sgn}(\xi) \hat{f}(\xi)\) 3.
\(H\) 是 \(L^p(\mathbb{R})\) 上的有界算子（\(1 < p < \infty\)） 4. \(H\)
与伸缩、平移交换

\textbf{定理 124.1}（M. Riesz 定理）：Hilbert 变换在 \(L^p(\mathbb{R})\)
上有界（\(1 < p < \infty\)）。\(\|Hf\|_p \leq C_p \|f\|_p\)。

\subsubsection{124.2 Calderón-Zygmund
分解}\label{calderuxf3n-zygmund-ux5206ux89e3}

\textbf{定理 124.2}（Calderón-Zygmund 分解）：设
\(f \in L^1(\mathbb{R}^n)\)，\(f \geq 0\)，\(\lambda > 0\)。则存在
\(\mathbb{R}^n\) 的分解：\(f = g + b\)，其中 -
\(g\)（``好''部分）：\(|g(x)| \leq 2^n \lambda\)
几乎处处，\(\|g\|_1 \leq \|f\|_1\) -
\(b = \sum_k b_k\)（``坏''部分）：\(b_k\) 支撑在不相交的二进方体 \(Q_k\)
上，\(\int_{Q_k} b_k = 0\)，\(\|b_k\|_1 \leq 2^{n+1} \lambda |Q_k|\)，且
\(\sum_k |Q_k| \leq \lambda^{-1} \|f\|_1\)

\emph{构造与证明概要}：将 \(\mathbb{R}^n\) 划分为二进方体（dyadic
cubes）。对每个二进方体 \(Q\)，若
\(\frac{1}{|Q|}\int_Q f > \lambda\)，则将其二等分为子方体并递归检查。最终选出的''坏''方体族
\(\{Q_k\}\) 满足：\(\frac{1}{|Q_k|}\int_{Q_k} f > \lambda\)，但每个
\(Q_k\) 的父方体均值 \(\leq \lambda\)。定义
\(g(x) = f(x)\)（\(x \notin \bigcup Q_k\)）和
\(g(x) = \frac{1}{|Q_k|}\int_{Q_k} f\)（\(x \in Q_k\)），以及
\(b_k(x) = (f(x) - \frac{1}{|Q_k|}\int_{Q_k} f) \mathbf{1}_{Q_k}(x)\)。于是
\(|g| \leq 2^n \lambda\)
由均值有界性保证，\(\sum|Q_k| \leq \lambda^{-1}\|f\|_1\) 由 Chebyshev
型估计得到。此分解是 Calderón-Zygmund
理论的基石，关键之处在于将函数分解为 \(L^\infty\)
控制的好部分和均值为零且支撑在不相交二进方体上的坏部分。∎

\emph{构造}：用二进方体分割 \(\mathbb{R}^n\)。对每个二进方体，若其均值
\(> \lambda\)，则从中提取''坏''部分。这是 Calderón-Zygmund 理论的基石。

\subsubsection{124.3
奇异积分算子}\label{ux5947ux5f02ux79efux5206ux7b97ux5b50}

\textbf{定义 124.2}（Calderón-Zygmund 核）：函数
\(K: \mathbb{R}^n \setminus \{0\} \to \mathbb{C}\) 称为
\textbf{Calderón-Zygmund 核}（或标准核），如果满足： 1.
\(|K(x)| \leq C |x|^{-n}\)（大小条件） 2.
\(|\nabla K(x)| \leq C |x|^{-n-1}\)（光滑性条件） 3.
\(\int_{r < |x| < R} K(x) dx = 0\)（对所有 \(0 < r < R\)，消去条件）

\textbf{定义 124.3}（奇异积分算子）：与核 \(K\)
关联的\textbf{奇异积分算子} \(T\) 定义为

\[Tf(x) = \operatorname{p.v.} \int_{\mathbb{R}^n} K(x - y) f(y) dy = \lim_{\varepsilon \to 0} \int_{|x-y| > \varepsilon} K(x - y) f(y) dy\]

\textbf{定理 124.3}（Calderón-Zygmund 定理）：若 \(T\)
是奇异积分算子且在 \(L^2\) 上有界，则 \(T\) 在 \(L^p\)
上有界（\(1 < p < \infty\)），且满足弱 \((1, 1)\) 型估计：

\[|\{x : |Tf(x)| > \lambda\}| \leq \frac{C}{\lambda} \|f\|_1\]

\emph{证明思路}：使用 Calderón-Zygmund 分解。对''好''部分 \(g\) 用
\(L^2\) 有界性，对''坏''部分 \(b\) 用核的光滑性和消去条件。弱 \((1, 1)\)
有界性 + \(L^2\) 有界性 + Marcinkiewicz 插值定理 ⇒ \(L^p\)
有界性（\(1 < p < 2\)）。对偶性处理 \(p > 2\)。∎

\textbf{定理 124.5}（Cotlar 引理，1955）：设
\(\{T_j\}_{j \in \mathbb{Z}}\) 是 Hilbert 空间 \(\mathcal{H}\)
上的一族有界算子。若存在 \(\gamma: \mathbb{Z} \to \mathbb{R}_{\geq 0}\)
满足 \(\sum_{k \in \mathbb{Z}} \sqrt{\gamma(k)} < \infty\)，使得
\(\|T_i^* T_j\| \leq \gamma(i-j)^2\) 和
\(\|T_i T_j^*\| \leq \gamma(i-j)^2\) 对所有 \(i, j\) 成立，则级数
\(\sum_j T_j\) 在强算子拓扑下收敛，且
\(\|\sum_j T_j\| \leq \sum_k \sqrt{\gamma(k)}\)。Cotlar
引理是证明奇异积分算子 \(L^2\) 有界性的核心工具：将核 \(K\)
按频率二进分解为 \(K = \sum_j K_j\)，对每个分片算子 \(T_j\)
验证几乎正交性条件
\(\|T_i^* T_j\| \leq C 2^{-|i-j|\varepsilon}\)，则级数在 \(L^2\)
上收敛且有界。其应用场景包括：Hilbert 变换的 \(L^2\) 有界性、Calderón
交换子、Fourier 乘子定理以及仿积算子的有界性。

\textbf{例 124.1}（Riesz 变换）：\(n\) 维 Riesz 变换
\(R_j\)（\(j = 1, \ldots, n\)）具有核
\(K_j(x) = c_n \frac{x_j}{|x|^{n+1}}\)。Fourier 乘子为
\(\widehat{R_j f}(\xi) = -i \frac{\xi_j}{|\xi|} \hat{f}(\xi)\)。Riesz
变换是 Calderón-Zygmund 定理的典型应用。

\subsubsection{124.4 Marcinkiewicz
插值定理}\label{marcinkiewicz-ux63d2ux503cux5b9aux7406}

\textbf{定理 124.4}（Marcinkiewicz 插值定理）：设 \(T\)
是次线性算子，满足： - 弱 \((p_0, q_0)\)
型：\(|\{|Tf| > \lambda\}| \leq (A_0 \|f\|_{p_0} / \lambda)^{q_0}\) - 弱
\((p_1, q_1)\)
型：\(|\{|Tf| > \lambda\}| \leq (A_1 \|f\|_{p_1} / \lambda)^{q_1}\)

其中
\(1 \leq p_0 < p_1 \leq \infty\)，\(1 \leq q_0, q_1 \leq \infty\)，\(q_0 \neq q_1\)。则对
\(p_\theta\) 满足
\(\frac{1}{p_\theta} = \frac{1-\theta}{p_0} + \frac{\theta}{p_1}\)，\(T\)
在 \(L^{p_\theta}\) 上是强有界的（\(0 < \theta < 1\)）。

\bigskip\hrule\bigskip

\subsection{第125章：Hardy-Littlewood
极大函数}\label{ux7b2c125ux7ae0hardy-littlewood-ux6781ux5927ux51fdux6570}

Hardy-Littlewood
极大函数是调和分析中最基本的算子之一，它控制了许多重要的算子（如奇异积分），是
Calderón-Zygmund 理论的关键工具。

\subsubsection{125.1 Hardy-Littlewood
极大函数}\label{hardy-littlewood-ux6781ux5927ux51fdux6570}

\textbf{定义 125.1}（Hardy-Littlewood
极大函数）：\(f \in L^1_{\text{loc}}(\mathbb{R}^n)\) 的
\textbf{Hardy-Littlewood 极大函数}为

\[Mf(x) = \sup_{r > 0} \frac{1}{|B(x, r)|} \int_{B(x, r)} |f(y)| dy\]

即 \(f\) 在包含 \(x\) 的球上的最大平均值。

\textbf{定理 125.1}（Hardy-Littlewood 极大定理）： 1. \(M\) 是弱
\((1, 1)\)
型：\(|\{x : Mf(x) > \lambda\}| \leq \frac{C_n}{\lambda} \|f\|_1\) 2.
\(M\) 是强 \((p, p)\)
型：\(\|Mf\|_p \leq C_{n,p} \|f\|_p\)（\(1 < p \leq \infty\)）

\emph{证明}（弱 \((1, 1)\) 型）：使用 Vitali 覆盖引理。对
\(\{Mf > \lambda\}\) 中的每点，存在球 \(B\) 使
\(\int_B |f| > \lambda |B|\)。由 Vitali 覆盖引理，可选出不相交的球
\(B_i\) 覆盖 \(\{Mf > \lambda\}\)
的缩放版本。\(|\{Mf > \lambda\}| \leq C_n \sum |B_i| \leq \frac{C_n}{\lambda} \sum \int_{B_i} |f| \leq \frac{C_n}{\lambda} \|f\|_1\)。∎

\textbf{定理 125.2}（Lebesgue 微分定理）：对
\(f \in L^1_{\text{loc}}(\mathbb{R}^n)\)，

\[\lim_{r \to 0} \frac{1}{|B(x, r)|} \int_{B(x, r)} f(y) dy = f(x) \quad \text{几乎处处}\]

\emph{证明}：由极大函数的弱 \((1, 1)\) 有界性和连续函数在 \(L^1\)
中的稠密性证明。∎

\subsubsection{125.2
其他极大函数}\label{ux5176ux4ed6ux6781ux5927ux51fdux6570}

\textbf{定义
125.2}（非中心极大函数）：\(M'f(x) = \sup_{B \ni x} \frac{1}{|B|} \int_B |f(y)| dy\)（上确界取遍所有包含
\(x\) 的球，不要求 \(x\) 是球心）。\(M'\) 与 \(M\)
逐点等价：\(Mf \leq M'f \leq 2^n Mf\)。

\textbf{定义 125.3}（锐极大函数，Sharp Maximal Function）：

\[M^\sharp f(x) = \sup_{B \ni x} \frac{1}{|B|} \int_B |f(y) - f_B| dy\]

其中 \(f_B = \frac{1}{|B|} \int_B f\) 是 \(f\) 在 \(B\)
上的均值。\(M^\sharp f\) 衡量 \(f\) 的局部振荡。

\textbf{定理 125.3}（Fefferman-Stein 不等式）：对
\(1 < p < \infty\)，\(\|f\|_p \leq C_p \|M^\sharp f\|_p\)（当 \(f\)
满足适当消失条件时）。

\subsubsection{125.3 加权不等式}\label{ux52a0ux6743ux4e0dux7b49ux5f0f}

\textbf{定义 125.4}（\(A_p\) 权）：非负局部可积函数 \(w\) 称为
\textbf{\(A_p\) 权}（\(1 < p < \infty\)），如果

\[\sup_B \left(\frac{1}{|B|} \int_B w(x) dx\right) \left(\frac{1}{|B|} \int_B w(x)^{-1/(p-1)} dx\right)^{p-1} < \infty\]

\(A_1\) 权：\(M w(x) \leq C w(x)\) 几乎处处。

\textbf{定理 125.4}（Muckenhoupt 定理）：Hardy-Littlewood 极大函数在加权
\(L^p(w)\) 空间上有界当且仅当
\(w \in A_p\)（\(1 < p < \infty\)）。这给出了极大函数 \(L^p\)
有界性权条件的精确刻画。

\bigskip\hrule\bigskip

\subsection{第126章：Littlewood-Paley
理论}\label{ux7b2c126ux7ae0littlewood-paley-ux7406ux8bba}

Littlewood-Paley
理论将函数分解为不同频率的分量，通过平方函数来研究函数的 \(L^p\)
性质。它是现代调和分析的核心工具，广泛应用于 PDE 和奇异积分理论。

\subsubsection{126.1 二进分解}\label{ux4e8cux8fdbux5206ux89e3}

\textbf{定义 126.1}（二进分解）：选择光滑函数
\(\varphi \in C_c^\infty(\mathbb{R}^n)\)，支撑在
\(\{1/2 \leq |\xi| \leq 2\}\)，满足
\(\sum_{j \in \mathbb{Z}} \varphi(2^{-j} \xi) = 1\)（\(\xi \neq 0\)）。定义
\textbf{Littlewood-Paley 算子}：

\[\Delta_j f = \mathcal{F}^{-1}(\varphi(2^{-j} \cdot) \hat{f})\]

\(\Delta_j f\) 是 \(f\) 在频率 \(|\xi| \sim 2^j\) 附近的分量。

\textbf{命题
126.1}（二进分解的性质）：\(f = \sum_{j \in \mathbb{Z}} \Delta_j f\)（在
\(\mathcal{S}'\) 中，模去多项式）。\(\Delta_j f\)
是光滑函数，频率集中在环形区域 \(\{2^{j-1} \leq |\xi| \leq 2^{j+1}\}\)。

\subsubsection{126.2 平方函数}\label{ux5e73ux65b9ux51fdux6570}

\textbf{定义 126.2}（Littlewood-Paley
平方函数）：\textbf{Littlewood-Paley \(g\)-函数}为

\[g(f)(x) = \left(\sum_{j \in \mathbb{Z}} |\Delta_j f(x)|^2\right)^{1/2}\]

\textbf{Lusin 面积函数}（\(S\)-函数）和 \textbf{Littlewood-Paley
\(g_\lambda^*\)-函数}是 \(g\)-函数的变体，涉及时间和空间上的平均。

\textbf{定理 126.1}（Littlewood-Paley 定理）：对
\(1 < p < \infty\)，存在常数 \(C_p\) 使得

\[C_p^{-1} \|f\|_p \leq \|g(f)\|_p \leq C_p \|f\|_p\]

即 \(L^p\) 范数与平方函数的 \(L^p\)
范数等价。平方函数允许用''频率分量的大小''来刻画函数的 \(L^p\) 性质。

\emph{证明思路}：\(p = 2\) 时由 Plancherel 定理直接得到。对一般
\(p\)，使用 Calderón-Zygmund 理论或向量值奇异积分理论。∎

\textbf{推论 126.2}（平方函数与 Sobolev 空间）：Sobolev 范数
\(\|f\|_{W^{s,p}}\) 与
\(\|(\sum_{j \in \mathbb{Z}} 2^{2js} |\Delta_j f|^2)^{1/2}\|_p\)
等价（\(1 < p < \infty\)）。

\subsubsection{126.3 Fourier 乘子}\label{fourier-ux4e58ux5b50}

\textbf{定义 126.3}（Fourier 乘子）：函数
\(m: \mathbb{R}^n \to \mathbb{C}\) 定义 \textbf{Fourier 乘子算子}
\(T_m\) 为

\[\widehat{T_m f}(\xi) = m(\xi) \hat{f}(\xi)\]

\textbf{定理 126.3}（Mikhlin 乘子定理）：设
\(m \in C^\infty(\mathbb{R}^n \setminus \{0\})\) 满足
\(|\partial^\alpha m(\xi)| \leq C_\alpha |\xi|^{-|\alpha|}\)
对所有多重指标 \(\alpha\)。则 \(T_m\) 在 \(L^p(\mathbb{R}^n)\)
上有界（\(1 < p < \infty\)）。

\textbf{定理 126.4}（Hörmander 乘子定理）：Mikhlin
条件的更精细版本，使用 Sobolev
空间代替逐点导数条件，允许乘子有更低的平滑性。

\textbf{例 126.1}（Fourier 乘子的例子）： - Hilbert
变换：\(m(\xi) = -i \operatorname{sgn}(\xi)\) - Riesz
变换：\(m_j(\xi) = -i \xi_j / |\xi|\) - 分数次积分（Riesz
位势）：\(m(\xi) = |\xi|^{-\alpha}\)（\(0 < \alpha < n\)） - Bessel
位势：\(m(\xi) = (1 + |\xi|^2)^{-\alpha/2}\)

\subsubsection{126.4 函数空间的 Littlewood-Paley
刻画}\label{ux51fdux6570ux7a7aux95f4ux7684-littlewood-paley-ux523bux753b}

\textbf{定理 126.4}（Besov 空间与 Triebel-Lizorkin 空间）：使用
Littlewood-Paley 分解可以定义一般的函数空间尺度： - \textbf{Besov 空间}
\(B^s_{p,q}\)：\(\|f\|_{B^s_{p,q}} = \|(2^{js} \|\Delta_j f\|_p)_{j \in \mathbb{Z}}\|_{\ell^q}\)
- \textbf{Triebel-Lizorkin 空间}
\(F^s_{p,q}\)：\(\|f\|_{F^s_{p,q}} = \|(\sum_{j \in \mathbb{Z}} 2^{jsq} |\Delta_j f|^q)^{1/q}\|_p\)

这些空间统一了 Sobolev
空间（\(F^s_{p,2} = W^{s,p}\)，\(1 < p < \infty\)）、Hölder
空间（\(B^s_{\infty,\infty} = C^s\)）、Hardy
空间（\(F^0_{p,2} = H^p\)）和 BMO
空间（\(F^0_{\infty,2} = \operatorname{BMO}\)）。

\textbf{定理 126.5}（Sobolev 嵌入的 Littlewood-Paley
证明）：使用平方函数和 Bernstein 不等式，Sobolev 嵌入定理（V18 定理
93.2）可以通过频率分解给出简洁证明。

\textbf{定义 126.5}（仿积分解，Paraproduct）：Bony
的\textbf{仿积分解}（paraproduct
decomposition）将两个函数的乘积分解为低频-低频、低频-高频和高频-低频三部分：

\[fg = \pi_f(g) + \pi_g(f) + R(f, g)\]

其中 \(\pi_f(g)\) 是 \(f\) 对 \(g\) 的仿积（低频乘以高频），\(R(f, g)\)
是余项（高频乘以高频）。仿积分是研究非线性 PDE
中乘积和复合函数正则性的关键工具。

\subsubsection{\texorpdfstring{126.5 Hardy 空间 \(H^1\)
的原子分解}{126.5 Hardy 空间 H\^{}1 的原子分解}}\label{hardy-ux7a7aux95f4-h1-ux7684ux539fux5b50ux5206ux89e3}

\textbf{定义 126.6}（\(H^1\) 原子）：函数 \(a \in L^2(\mathbb{R}^n)\)
称为 \textbf{\(H^1\)-原子}，如果存在球 \(B\) 使得：(1)
\(\operatorname{supp}(a) \subseteq B\)；(2)
\(\|a\|_\infty \leq |B|^{-1}\)（大小条件）；(3)
\(\int_{\mathbb{R}^n} a(x) dx = 0\)（消失矩条件）。

\textbf{定理 126.6}（\(H^1\) 的原子分解）：\(f \in H^1(\mathbb{R}^n)\)
当且仅当存在 \(H^1\)-原子的可数族 \(\{a_j\}\) 和系数
\(\{\lambda_j\} \subseteq \mathbb{C}\) 满足
\(\sum |\lambda_j| < \infty\)，使得
\(f = \sum_j \lambda_j a_j\)（在分布意义下）。此时
\(\|f\|_{H^1} \sim \inf \sum |\lambda_j|\)。

\(H^1\) 取代 \(L^1\)
在奇异积分理论中的角色是因为奇异积分核的消去条件与原子消失矩条件相匹配：\(L^1\)
在奇异积分下不封闭（仅满足弱 \((1,1)\) 型），但 \(H^1\) 是最大的满足
Calderón-Zygmund 算子在其上有界的函数空间。\(H^1\) 原子分解由 Coifman 和
Weiss 建立，为研究奇异积分的端点估计提供了基本工具。

\bigskip\hrule\bigskip

\emph{本卷建立了调和分析的核心理论框架：Fourier 变换的 \(L^1\)、\(L^2\)
和 Schwartz 分布理论将函数从物理空间映射到频率空间；Calderón-Zygmund
奇异积分理论（分解引理、\(L^p\)
有界性）为高维调和分析提供了基础；Hardy-Littlewood
极大函数是控制各类算子的基本工具；Littlewood-Paley
理论通过频率分解和平方函数，统一了函数空间理论（Sobolev 空间、Besov
空间、Triebel-Lizorkin 空间），并为非线性 PDE
的微局部分析提供了核心工具。调和分析是连接 Fourier
分析、偏微分方程和概率论的桥梁，在现代数学中占有核心地位。}

\newpage

\section{卷二十五：组合数学与图论}\label{ux5377ux4e8cux5341ux4e94ux7ec4ux5408ux6570ux5b66ux4e0eux56feux8bba}

\begin{quote}
组合数学研究离散结构的计数、存在性和构造问题，图论研究由顶点和边构成的网络结构。本卷在初等计数原理（V6
Ch 28）的基础上，建立生成函数、图论、匹配与网络流、染色与 Ramsey
理论，以及组合设计的系统理论。组合数学与代数、概率论、计算机科学和优化理论有深刻联系。
\end{quote}

\bigskip\hrule\bigskip

\subsection{第127章：生成函数与递推关系}\label{ux7b2c127ux7ae0ux751fux6210ux51fdux6570ux4e0eux9012ux63a8ux5173ux7cfb}

生成函数是组合数学中最强大的工具之一，将离散的数列编码为形式幂级数，通过代数运算解决计数问题。递推关系是数列满足的方程，生成函数提供了求解递推关系的系统方法。

\subsubsection{127.1
普通生成函数}\label{ux666eux901aux751fux6210ux51fdux6570}

\textbf{定义 127.1}（普通生成函数）：数列 \((a_n)_{n \geq 0}\)
的\textbf{普通生成函数}（Ordinary Generating Function, OGF）是形式幂级数

\[A(x) = \sum_{n=0}^{\infty} a_n x^n\]

其中 \(x\) 是形式变量，不考虑收敛性。两个生成函数的加法和乘法定义为：

\[(A+B)(x) = \sum_{n=0}^{\infty} (a_n + b_n) x^n, \quad (A \cdot B)(x) = \sum_{n=0}^{\infty} \left( \sum_{k=0}^n a_k b_{n-k} \right) x^n\]

\textbf{例 127.1}（基本生成函数）： - 常数列
\(a_n = 1\)：\(\frac{1}{1-x} = \sum_{n=0}^{\infty} x^n\) -
\(a_n = n\)：\(\frac{x}{(1-x)^2} = \sum_{n=0}^{\infty} n x^n\) -
二项式系数
\(a_n = \binom{m}{n}\)：\((1+x)^m = \sum_{n=0}^{m} \binom{m}{n} x^n\) -
Fibonacci 数列：\(F_0 = 0, F_1 = 1, F_n = F_{n-1} + F_{n-2}\)
的生成函数为 \(\frac{x}{1-x-x^2}\)

\textbf{命题 127.1}（生成函数与递推关系）：若数列 \(a_n\) 满足 \(k\)
阶常系数线性递推

\[a_n = c_1 a_{n-1} + c_2 a_{n-2} + \cdots + c_k a_{n-k}\]

则其生成函数 \(A(x)\) 是有理函数
\(A(x) = \frac{P(x)}{1 - c_1 x - c_2 x^2 - \cdots - c_k x^k}\)，其中
\(P(x)\) 是由初始条件决定的多项式。

\emph{证明}：将递推关系乘以 \(x^n\) 并对 \(n \geq k\) 求和：
\[\sum_{n=k}^\infty a_n x^n = c_1 \sum_{n=k}^\infty a_{n-1} x^n + \cdots + c_k \sum_{n=k}^\infty a_{n-k} x^n\]
左边为 \(A(x) - \sum_{n=0}^{k-1} a_n x^n\)，右边第 \(j\) 项为
\(c_j x^j \sum_{n=k}^\infty a_{n-j} x^{n-j} = c_j x^j (A(x) - \sum_{n=0}^{k-j-1} a_n x^n)\)。移项整理得
\((1 - \sum_{j=1}^k c_j x^j)A(x) = P(x)\)，其中 \(P(x)\) 是次数小于
\(k\) 的多项式（由初始条件确定）。故
\(A(x) = P(x)/(1 - \sum_{j=1}^k c_j x^j)\) 是有理函数。∎

\textbf{定理 127.2}（有理生成函数的分解）：若
\(A(x) = \frac{P(x)}{Q(x)}\) 是有理函数，分母
\(Q(x) = \prod_{i=1}^k (1 - \lambda_i x)^{m_i}\)（\(\lambda_i\)
互异），则

\[a_n = \sum_{i=1}^k \sum_{j=1}^{m_i} c_{i,j} n^{j-1} \lambda_i^n\]

其中 \(c_{i,j}\) 由部分分式分解确定。

\emph{证明}：将 \(A(x)\) 部分分式分解为
\(\sum_{i,j} \frac{c_{i,j}}{(1-\lambda_i x)^j}\)，利用
\(\frac{1}{(1-\lambda x)^j} = \sum_{n=0}^{\infty} \binom{n+j-1}{j-1} \lambda^n x^n\)
即得。∎

\textbf{例 127.2}（Fibonacci
数列的闭式）：\(F(x) = \frac{x}{1-x-x^2} = \frac{x}{(1-\phi x)(1-\hat{\phi}x)}\)，其中
\(\phi = \frac{1+\sqrt{5}}{2}\)，\(\hat{\phi} = \frac{1-\sqrt{5}}{2}\)。部分分式分解得

\[F_n = \frac{1}{\sqrt{5}}(\phi^n - \hat{\phi}^n)\]

\subsubsection{127.2
指数生成函数}\label{ux6307ux6570ux751fux6210ux51fdux6570}

\textbf{定义 127.2}（指数生成函数）：数列 \((a_n)_{n \geq 0}\)
的\textbf{指数生成函数}（Exponential Generating Function, EGF）为

\[\hat{A}(x) = \sum_{n=0}^{\infty} a_n \frac{x^n}{n!}\]

EGF 的乘法对应二项式卷积：若 \(\hat{C}(x) = \hat{A}(x) \hat{B}(x)\)，则

\[c_n = \sum_{k=0}^n \binom{n}{k} a_k b_{n-k}\]

\textbf{例 127.3}（EGF 的基本例子）： - \(a_n = 1\) 的 EGF 为 \(e^x\) -
\(a_n = n!\)（排列数）的 EGF 为 \(\frac{1}{1-x}\) - Bell 数
\(B_n\)（\(n\) 元集合的划分个数）的 EGF 为 \(e^{e^x - 1}\)

\textbf{命题 127.3}（标号结构与
EGF）：若组合结构可以分解为不相交的连通分支，且每个分支持有标号，则 EGF
满足指数公式：若 \(\mathcal{C}\) 是连通结构的类，其 EGF 为
\(C(x)\)，则所有结构（连通分支的集合）的 EGF 为 \(e^{C(x)}\)。

\textbf{例 127.4}（错排数）：\(n\) 个元素的错排数 \(D_n\) 满足
\(D_n = n! \sum_{k=0}^n \frac{(-1)^k}{k!}\)。其 EGF 为
\(\hat{D}(x) = \frac{e^{-x}}{1-x}\)。

\subsubsection{127.3
递推关系的求解}\label{ux9012ux63a8ux5173ux7cfbux7684ux6c42ux89e3}

\textbf{定义 127.3}（常系数线性递推）：\(k\) 阶常系数齐次线性递推：

\[a_n = c_1 a_{n-1} + c_2 a_{n-2} + \cdots + c_k a_{n-k}, \quad n \geq k\]

对应的\textbf{特征方程}为
\(\lambda^k - c_1 \lambda^{k-1} - \cdots - c_k = 0\)。

\textbf{定理 127.4}（特征根法）：设特征根为
\(\lambda_1, \ldots, \lambda_r\)，重数分别为
\(m_1, \ldots, m_r\)（\(\sum m_i = k\)）。则通解为

\[a_n = \sum_{i=1}^r \sum_{j=1}^{m_i} \alpha_{i,j} n^{j-1} \lambda_i^n\]

\textbf{例 127.5}（Catalan 数的递推）：Catalan 数
\(C_n = \frac{1}{n+1}\binom{2n}{n}\) 满足非线性递推

\[C_0 = 1, \quad C_{n+1} = \sum_{k=0}^n C_k C_{n-k}\]

其生成函数 \(C(x) = \sum_{n=0}^{\infty} C_n x^n\) 满足
\(C(x) = 1 + x C(x)^2\)，解得 \(C(x) = \frac{1-\sqrt{1-4x}}{2x}\)。

\textbf{定理 127.5}（主定理 / Master Theorem）：对分治递推
\(T(n) = a T(n/b) + f(n)\)，其中 \(a \geq 1, b > 1\)：

\begin{enumerate}
\def\labelenumi{\arabic{enumi}.}
\tightlist
\item
  若 \(f(n) = O(n^{\log_b a - \varepsilon})\)，则
  \(T(n) = \Theta(n^{\log_b a})\)
\item
  若 \(f(n) = \Theta(n^{\log_b a} \log^k n)\)，则
  \(T(n) = \Theta(n^{\log_b a} \log^{k+1} n)\)
\item
  若 \(f(n) = \Omega(n^{\log_b a + \varepsilon})\) 且
  \(a f(n/b) \leq c f(n)\)（\(c < 1\)），则 \(T(n) = \Theta(f(n))\)
\end{enumerate}

\subsubsection{127.4 容斥原理}\label{ux5bb9ux65a5ux539fux7406}

\textbf{定理 127.6}（容斥原理）：设 \(A_1, \ldots, A_n\) 是有限集合
\(X\) 的子集，则

\[\left| X \setminus \bigcup_{i=1}^n A_i \right| = \sum_{k=0}^n (-1)^k \sum_{1 \leq i_1 < \cdots < i_k \leq n} \left| A_{i_1} \cap \cdots \cap A_{i_k} \right|\]

特别地，\(|X| = \sum_{k=0}^n (-1)^k S_k\)，其中
\(S_k = \sum_{i_1 < \cdots < i_k} |A_{i_1} \cap \cdots \cap A_{i_k}|\)。

\emph{证明}：对每个元素 \(x \in X\)，设 \(x\) 恰好属于 \(r\) 个
\(A_i\)。则 \(x\) 在右边被计数
\(\sum_{k=0}^r (-1)^k \binom{r}{k} = (1-1)^r = 0\) 次（若
\(r > 0\)），而被计数 \(1\) 次（若 \(r = 0\)）。∎

\textbf{例 127.6}（Euler 函数
\(\varphi(n)\)）：\(\varphi(n) = n \prod_{p|n} \left(1 - \frac{1}{p}\right)\)
由容斥原理给出：从 \(n\) 个数中减去被各素因子整除的数。

\textbf{定理 127.7}（Möbius 反演）：设 \(f, g\) 为算术函数，则

\[g(n) = \sum_{d|n} f(d) \iff f(n) = \sum_{d|n} \mu(d) g\left(\frac{n}{d}\right)\]

其中 \(\mu\) 是 Möbius 函数：\(\mu(n) = 0\) 若 \(n\) 有平方因子，否则
\(\mu(n) = (-1)^k\)（\(k\) 为 \(n\) 的素因子个数）。

\bigskip\hrule\bigskip

\subsection{第128章：图论基础}\label{ux7b2c128ux7ae0ux56feux8bbaux57faux7840}

图论是离散数学的核心分支，起源于 Euler 的 Königsberg
七桥问题（1736年）。本章建立图的基本概念，包括树、连通性、平面图和
Euler/Hamilton 回路。

\subsubsection{128.1
图的基本概念}\label{ux56feux7684ux57faux672cux6982ux5ff5}

\textbf{定义 128.1}（图）：一个\textbf{图} \(G = (V, E)\) 由顶点集
\(V\)（有限集）和边集 \(E \subseteq V \times V\) 组成。若无向图，边
\(\{u, v\}\) 是无序对；若有向图，边 \((u, v)\) 是有序对（弧）。

\textbf{定义 128.2}（基本概念）： - \textbf{度}：顶点 \(v\) 的度
\(\deg(v)\) 是与 \(v\) 关联的边数。在有向图中，出度 \(\deg^+(v)\) 和入度
\(\deg^-(v)\) 分别计数出边和入边。 -
\textbf{握手引理}：\(\sum_{v \in V} \deg(v) = 2|E|\)。（每条边贡献两个度）
- \textbf{路}：长度为 \(k\) 的路是顶点序列 \(v_0 v_1 \cdots v_k\)，其中
\(\{v_i, v_{i+1}\} \in E\)，且顶点互异。 -
\textbf{回路}：起点和终点相同的路（\(v_0 = v_k\)，\(k \geq 3\)）。 -
\textbf{连通}：图 \(G\) 是连通的，如果任意两个顶点之间存在路。 -
\textbf{连通分支}：图的最大连通子图。

\textbf{定义 128.3}（特殊图类）： - \textbf{完全图} \(K_n\)：\(n\)
个顶点，每对顶点之间都有边（\(\binom{n}{2}\) 条边） -
\textbf{二分图}：顶点集可划分为 \(V = A \cup B\)，所有边都在 \(A\) 和
\(B\) 之间 - \textbf{完全二分图}
\(K_{m,n}\)：\(|A| = m, |B| = n\)，\(A\) 和 \(B\) 之间所有可能的边 -
\textbf{路图} \(P_n\)：\(n\) 个顶点排成一条路 - \textbf{回路图}
\(C_n\)：\(n\) 个顶点排成一个圈

\textbf{命题 128.1}（二分图的无奇圈特征）：图 \(G\) 是二分图当且仅当
\(G\) 不含奇长度回路。

\emph{证明}：若 \(G\) 是二分图，路在 \(A\) 和 \(B\)
之间交替，故回路长度必为偶数。反之，固定根顶点 \(v_0\)，对每个顶点
\(v\)，令 \(d(v)\) 为从 \(v_0\) 到 \(v\) 的最短路长度，按 \(d(v)\)
的奇偶性划分 \(V\)。无奇回路保证此划分是良定义的二分划分。∎

\subsubsection{128.2 树与生成树}\label{ux6811ux4e0eux751fux6210ux6811}

\textbf{定义 128.4}（树）：\textbf{树}是连通无回路的图。等价刻画：\(n\)
个顶点的树恰有 \(n-1\) 条边，且任意两点间存在唯一的路。

\textbf{定理 128.2}（Cayley 公式）：\(n\) 个标号顶点的不同树的个数为
\(n^{n-2}\)。

\emph{证明思路}：Prufer 编码：每个标号树唯一对应一个长度为 \(n-2\)
的序列，每个位置可从 \(1\) 到 \(n\) 取值。编码算法：反复删除度数为 \(1\)
的最小标号顶点，记录其邻居。∎

\textbf{定义 128.5}（生成树）：连通图 \(G\) 的\textbf{生成树}是 \(G\)
的包含所有顶点的树子图。

\textbf{定理 128.3}（Kirchhoff 矩阵树定理）：连通图 \(G\)
的生成树个数等于其 Laplace 矩阵 \(L = D - A\)（\(D\) 为度对角矩阵，\(A\)
为邻接矩阵）的任意一个 \((n-1)\) 阶主子式。

\textbf{定义 128.6}（最小生成树）：带权图 \(G\)
的\textbf{最小生成树}（MST）是权和最小的生成树。Kruskal
算法（贪心选择最小权边，不形成回路）和 Prim 算法均可找到 MST。

\textbf{定理 128.4}（Kruskal 算法的正确性）：对任意连通带权图，Kruskal
算法输出最小生成树。

\emph{证明思路}：使用\textbf{割性质}（cut
property）：对任意割，连接割两端的最小权边必属于某个 MST。Kruskal
算法每次添加的边满足此性质。∎

\subsubsection{128.3 连通性与 Menger
定理}\label{ux8fdeux901aux6027ux4e0e-menger-ux5b9aux7406}

\textbf{定义 128.7}（连通度）： - \textbf{点连通度}
\(\kappa(G)\)：使图不连通或退化为单点所需删除的最少顶点数 -
\textbf{边连通度} \(\lambda(G)\)：使图不连通所需删除的最少边数 - \(G\)
是 \(k\)-连通的，如果 \(\kappa(G) \geq k\)

\textbf{命题 128.5}（Whitney
不等式）：\(\kappa(G) \leq \lambda(G) \leq \delta(G)\)，其中
\(\delta(G)\) 是最小度。

\textbf{定理 128.6}（Menger 定理，点版本）：图 \(G\) 中两个不相邻顶点
\(u, v\) 之间的内部不相交的 \(u\)-\(v\) 路的最大条数等于分离 \(u\) 和
\(v\) 所需删除的最少顶点数。

\textbf{定理 128.7}（Menger 定理，边版本）：图 \(G\) 中两个顶点 \(u, v\)
之间的边不相交的 \(u\)-\(v\) 路的最大条数等于分离 \(u\) 和 \(v\)
所需删除的最少边数。

\textbf{推论 128.8}（全局 Menger 定理）：图 \(G\) 是
\(k\)-连通的当且仅当任意两个顶点之间存在至少 \(k\) 条内部不相交的路。

\subsubsection{128.4 Euler 回路与 Hamilton
回路}\label{euler-ux56deux8defux4e0e-hamilton-ux56deux8def}

\textbf{定理 128.9}（Euler 定理，1736）：连通图 \(G\) 有 Euler
回路（经过每条边恰好一次的回路）当且仅当每个顶点的度数为偶数。有 Euler
路（经过每条边恰好一次的路，起终点可不同）当且仅当恰好有两个奇数度顶点。

\emph{证明}：必要性：每进入一个顶点就离开一次，故度数为偶数。充分性：使用
Fleury 算法或 Hierholzer 算法（贪心合并回路）。∎

\textbf{定义 128.8}（Hamilton 回路）：Hamilton
回路是经过每个顶点恰好一次的回路。判定 Hamilton 回路的存在性是 NP
完全的。

\textbf{定理 128.10}（Dirac 定理，1952）：若 \(G\) 是 \(n \geq 3\)
个顶点的图，且 \(\delta(G) \geq n/2\)，则 \(G\) 有 Hamilton 回路。

\textbf{定理 128.11}（Ore 定理，1960）：若对 \(G\) 中任意不相邻的顶点
\(u, v\)，有 \(\deg(u) + \deg(v) \geq n\)，则 \(G\) 有 Hamilton 回路。

\emph{证明思路}：反证法。取 \(G\) 的极大非 Hamilton 图，添加任意边会产生
Hamilton 回路。考虑 \(G\) 中不相邻的顶点 \(u, v\)，添加边 \(\{u, v\}\)
后产生 Hamilton 回路，由此导出 \(\deg(u) + \deg(v) \geq n\)
与假设矛盾。∎

\subsubsection{128.5 平面图}\label{ux5e73ux9762ux56fe}

\textbf{定义 128.9}（平面图）：图 \(G\)
是\textbf{平面图}，如果它可以画在平面上且边不交叉（仅在端点相交）。平面图的画法将平面划分为\textbf{面}（face）。

\textbf{定理 128.12}（Euler 公式，1758）：对连通平面图 \(G\)，有

\[n - m + f = 2\]

其中 \(n = |V|\)，\(m = |E|\)，\(f\) 为面数（包括外部无界面）。

\emph{证明}：对边数 \(m\) 归纳。奠基：\(m = 0\) 时连通图仅有
\(n = 1\)，\(f = 1\)，\(1 - 0 + 1 = 2\) 成立。 归纳假设
\(m \geq 1\)。考虑两种情况：（1）若 \(G\) 是树，则
\(m = n - 1\)，\(f = 1\)（仅外部面），\(n - (n-1) + 1 = 2\)。（2）若
\(G\) 含回路，取回路上的任意边 \(e\)。\(G \setminus e\) 仍连通且面数减少
1（\(e\) 分隔的两个面合并为一个），边数减 1，顶点数不变。由归纳假设
\(n - (m-1) + (f-1) = 2\)，即 \(n - m + f = 2\)。∎

\textbf{推论 128.13}：对 \(n \geq 3\)
的简单平面图，\(m \leq 3n - 6\)。若图中无三角形，则 \(m \leq 2n - 4\)。

\emph{证明}：每个面至少由 3 条边围成，而每条边至多属于 2 个面，故
\(3f \leq 2m\)。代入 Euler
公式：\(2 = n - m + f \leq n - m + 2m/3 = n - m/3\)，得
\(m \leq 3n - 6\)。∎

\textbf{推论 128.14}：\(K_5\) 和 \(K_{3,3}\) 不是平面图。

\emph{证明}：\(K_5\)：\(n = 5, m = 10\)，\(10 > 3 \cdot 5 - 6 = 9\)。\(K_{3,3}\)：\(n = 6, m = 9\)，无三角形，\(9 > 2 \cdot 6 - 4 = 8\)。∎

\textbf{定理 128.15}（Kuratowski 定理，1930）：图 \(G\) 是平面图当且仅当
\(G\) 不含 \(K_5\) 或 \(K_{3,3}\) 的细分（即通过插入度数为 2
的顶点得到的子图）。

\emph{充分性简述}（若不含 \(K_5\) 或 \(K_{3,3}\) 细分则平面）：由推论
128.14 知 \(K_5\) 和 \(K_{3,3}\)
本身非平面，其细分亦然。充分性证明较复杂：对顶点数归纳，利用
3-连通图的平面嵌入性质（Tutte
关于凸嵌入的定理），最终归结为不含此类子图的图必可平面嵌入。详细证明超出本教材范围。∎

\bigskip\hrule\bigskip

\subsection{第129章：匹配与网络流}\label{ux7b2c129ux7ae0ux5339ux914dux4e0eux7f51ux7edcux6d41}

匹配理论研究图中不含公共顶点的边集合，是分配问题、婚姻问题等的数学模型。网络流理论将匹配问题推广到带容量的有向图，是运筹学的核心工具之一。

\subsubsection{129.1 二分图匹配}\label{ux4e8cux5206ux56feux5339ux914d}

\textbf{定义 129.1}（匹配）：图 \(G = (V, E)\) 的一个\textbf{匹配}
\(M \subseteq E\)
是两两不共享顶点的边集合。最大匹配是边数最多的匹配。完美匹配覆盖所有顶点。

\textbf{定义 129.2}（增广路）：对于匹配 \(M\)，\textbf{交替路}是边在
\(M\) 和非 \(M\)
之间交替的路。\textbf{增广路}是起终点都为未匹配顶点的交替路。

\textbf{定理 129.1}（Berge 引理，1957）：匹配 \(M\)
是最大匹配当且仅当不存在 \(M\)-增广路。

\emph{证明}：若存在增广路 \(P\)，则对称差 \(M \triangle P\)
是更大的匹配。反之，若 \(M\) 非最大，取 \(M'\) 更大，\(M \triangle M'\)
的每个连通分支要么是偶长度回路，要么是包含 \(M'\) 边更多的路，后者即为
\(M\)-增广路。∎

\textbf{定理 129.2}（Kőnig 定理，1931）：在二分图 \(G = (A \cup B, E)\)
中，最大匹配的大小等于最小顶点覆盖的大小。

\emph{证明思路}：从最大匹配 \(M\) 出发，在 \(A\)
中未匹配顶点开始交替路，令 \(Z\) 为可达的顶点集。令
\(K = (A \setminus Z) \cup (B \cap Z)\)，可证 \(K\) 是覆盖且
\(|K| = |M|\)。∎

\textbf{定理 129.3}（Hall 婚姻定理，1935）：二分图 \(G = (A \cup B, E)\)
存在覆盖 \(A\) 的匹配当且仅当对任意 \(S \subseteq A\)，有
\(|N(S)| \geq |S|\)（Hall 条件），其中 \(N(S)\) 是 \(S\) 的邻居集。

\emph{证明}：必要性：若存在覆盖 \(A\) 的匹配 \(M\)，则对任意
\(S \subseteq A\)，\(M\) 将 \(S\) 匹配到 \(|S|\) 个不同的 \(B\)
中顶点，故 \(|N(S)| \geq |S|\)。

充分性（对 \(|A|\) 归纳）：\(|A|=1\) 时平凡。设
\(|A|>1\)。分两种情况：（1）若对任意非空真子集 \(S \subset A\) 有
\(|N(S)| \geq |S|+1\)（Hall 条件加强），任选 \(a \in A\) 匹配到其邻居
\(b \in B\)，删去两者后剩余图对 \(A\setminus\{a\}\) 仍满足 Hall
条件，由归纳假设得证。（2）存在非空真子集 \(S \subset A\) 使得
\(|N(S)| = |S|\)。对二分图 \(S \cup N(S)\)，由归纳假设存在 \(S\) 到
\(N(S)\) 的完美匹配 \(M_1\)。考虑 \(G \setminus (S \cup N(S))\)。对任意
\(T \subseteq A \setminus S\)，\(|N_{G\setminus}(T)| = |N(T \cup S) \setminus N(S)| \geq |T \cup S| - |N(S)| = |T| + |S| - |S| = |T|\)。故
\(A \setminus S\) 到 \(B \setminus N(S)\) 满足 Hall
条件，由归纳假设存在匹配 \(M_2\)。\(M_1 \cup M_2\) 即为覆盖 \(A\)
的匹配。∎

\textbf{推论
129.4}：正则二分图有完美匹配。特别地，\(k\)-正则二分图可分解为 \(k\)
个不相交的完美匹配。

\subsubsection{129.2 一般图匹配}\label{ux4e00ux822cux56feux5339ux914d}

\textbf{定义 129.3}（Tutte 定理）：图 \(G\) 有完美匹配当且仅当对任意
\(S \subseteq V\)，\(G \setminus S\) 的奇数阶连通分支数不超过 \(|S|\)。

\textbf{定义 129.4}（Edmonds 花算法）：Edmonds 在 1965
年提出了多项式时间的最大匹配算法（花算法），核心思想是收缩奇长度回路（花）以处理非二分图中的奇回路障碍。

\subsubsection{129.3
网络流与最大流}\label{ux7f51ux7edcux6d41ux4e0eux6700ux5927ux6d41}

\textbf{定义 129.5}（网络）：一个\textbf{网络} \(N = (V, E, s, t, c)\)
由有向图 \((V, E)\)、源点 \(s\)、汇点 \(t\) 和容量函数
\(c: E \to \mathbb{R}_{\geq 0}\) 组成。

\textbf{定义 129.6}（流）：一个\textbf{流}
\(f: E \to \mathbb{R}_{\geq 0}\) 满足： 1.
\textbf{容量限制}：\(0 \leq f(e) \leq c(e)\) 对所有 \(e \in E\) 2.
\textbf{流量守恒}：对
\(v \neq s, t\)，\(\sum_{e \text{ in } v} f(e) = \sum_{e \text{ out } v} f(e)\)

流的值
\(|f| = \sum_{e \text{ out } s} f(e) - \sum_{e \text{ in } s} f(e)\)。

\textbf{定义 129.7}（割）：\(s\)-\(t\) 割 \((S, T)\) 满足
\(s \in S, t \in T\)。割的容量为
\(c(S, T) = \sum_{e \text{ from } S \text{ to } T} c(e)\)。

\textbf{定理 129.5}（最大流最小割定理 /
Ford-Fulkerson，1956）：在任何网络中，最大流的值等于最小割的容量。

\emph{证明思路}：Ford-Fulkerson
算法：从零流开始，反复寻找剩余网络中的增广路，沿增广路增加流量。算法终止时，从
\(s\) 可达的顶点集 \(S\)
给出最小割。证明中还需处理容量为无理数时的收敛性问题。∎

\textbf{定理
129.6}（整性定理）：若所有容量为整数，则存在整数值的最大流（Ford-Fulkerson
算法每次增广整数流量，故保持整性）。

\textbf{定义 129.8}（Edmonds-Karp
算法）：每次选择最短增广路（BFS），时间复杂度 \(O(|V| |E|^2)\)。Dinic
算法使用分层网络，时间复杂度 \(O(|V|^2 |E|)\)。

\textbf{定理 129.7}（最大流与二分图匹配）：二分图 \(G = (A \cup B, E)\)
的最大匹配问题等价于网络流：添加源点 \(s\) 连接 \(A\) 中所有顶点（容量
1），\(B\) 中所有顶点连接汇点 \(t\)（容量 1），每条边 \(A \to B\) 容量
1。最大流的值等于最大匹配的大小。

\subsubsection{129.4
最小费用流与分配问题}\label{ux6700ux5c0fux8d39ux7528ux6d41ux4e0eux5206ux914dux95eeux9898}

\textbf{定义 129.9}（最小费用流）：在网络上还定义边费用 \(w(e)\)，求从
\(s\) 到 \(t\) 发送给定流量 \(F\) 的最小总费用 \(\sum_e w(e) f(e)\)。

\textbf{定义 129.10}（分配问题 / Assignment Problem）：给定 \(n\)
个工人和 \(n\) 个任务，工人 \(i\) 完成任务 \(j\) 的费用为
\(c_{ij}\)，求使总费用最小的一一分配。这是最小费用流（在完全二分图上）的特例。

\textbf{定理 129.8}（匈牙利算法 / Kuhn-Munkres 算法）：可在 \(O(n^3)\)
时间内求解分配问题，基于对偶变量（势）的调整。

\bigskip\hrule\bigskip

\subsection{第130章：染色与 Ramsey
理论}\label{ux7b2c130ux7ae0ux67d3ux8272ux4e0e-ramsey-ux7406ux8bba}

染色理论研究用最少颜色给图的顶点或边染色，使得相邻元素不同色。Ramsey
理论是''完全无序不可能''的数学形式化------任何足够大的结构必然包含某种有序子结构。

\subsubsection{130.1 顶点染色}\label{ux9876ux70b9ux67d3ux8272}

\textbf{定义 130.1}（色数）：图 \(G\) 的\textbf{色数} \(\chi(G)\)
是使相邻顶点颜色不同的最少颜色数。\(G\) 是 \(k\)-可染的，如果
\(\chi(G) \leq k\)。

\textbf{命题 130.1}（基本界）： -
\(\chi(G) \geq \omega(G)\)（团数---最大完全子图的大小） -
\(\chi(G) \geq \frac{n}{\alpha(G)}\)（独立数---最大独立集的大小） -
\(\chi(G) \leq \Delta(G) + 1\)（贪心染色）

\textbf{定理 130.2}（Brooks 定理，1941）：若 \(G\)
是连通图且非完全图也非奇回路，则 \(\chi(G) \leq \Delta(G)\)。

\textbf{定理 130.3}（四色定理 / Appel-Haken，1976）：任何平面图是
4-可染的。

\emph{说明}：四色定理的证明是首个依赖计算机验证的数学定理，2015 年
Gonthier 用 Coq 形式化了证明。

\textbf{定理 130.4}（五色定理 / Heawood，1890）：任何平面图是
5-可染的（此定理有简短的纯数学证明）。

\emph{证明}：平面图必有度 \(\leq 5\)
的顶点。归纳法：删除该顶点，对剩余图 5-染色，若该顶点的邻居使用了所有 5
种颜色，则在由两种颜色诱导的子图中交换颜色来释放一种颜色（Kempe
链论证）。∎

\subsubsection{130.2 边染色}\label{ux8fb9ux67d3ux8272}

\textbf{定义 130.2}（边色数）：图 \(G\) 的\textbf{边色数} \(\chi'(G)\)
是使相邻边（共享顶点）颜色不同的最少颜色数。

\textbf{定理 130.5}（Vizing 定理，1964）：对任何简单图
\(G\)，\(\Delta(G) \leq \chi'(G) \leq \Delta(G) + 1\)。

\textbf{定理 130.6}（Kőnig 边染色定理）：对二分图
\(G\)，\(\chi'(G) = \Delta(G)\)。

\emph{证明}：对 \(\Delta(G)\) 归纳，使用 Hall
婚姻定理找到完美匹配覆盖所有最大度顶点，然后删除该匹配。∎

\subsubsection{130.3 Ramsey 理论}\label{ramsey-ux7406ux8bba}

\textbf{定理 130.7}（Ramsey 定理，图版本，1930）：对任意正整数
\(r, s\)，存在最小整数 \(R(r, s)\)，使得任何 \(R(r, s)\) 个顶点的完全图
\(K_n\)（\(n \geq R(r, s)\)），其边任意染红/蓝二色，必含红色 \(K_r\)
或蓝色 \(K_s\)。

\emph{证明}（归纳法）：对 \(r+s\) 做归纳。基始：\(R(1,s)=1\)（红色
\(K_1\) 平凡），\(R(r,1)=1\)，\(R(2,s)=s\)，\(R(r,2)=r\)。归纳步：设
\(R(r-1,s)\) 和 \(R(r,s-1)\) 存在。取 \(n = R(r-1,s) + R(r,s-1)\)，对
\(K_n\) 做红蓝染色。任选顶点 \(v\)，设 \(v\) 有 \(d_R\) 条红边、\(d_B\)
条蓝边，\(d_R+d_B=n-1\)。若
\(d_R \geq R(r-1,s)\)，则红邻居诱导的子图必含红色 \(K_{r-1}\)（连同
\(v\) 得红色 \(K_r\)）或蓝色 \(K_s\)。若
\(d_B \geq R(r,s-1)\)，同理得红色 \(K_r\) 或蓝色 \(K_{s-1}\)（连同 \(v\)
得蓝色 \(K_s\)）。由于
\(d_R+d_B \geq R(r-1,s)+R(r,s-1)-1\)，至少一种情形发生。故
\(R(r,s) \leq R(r-1,s)+R(r,s-1)\)。∎

\textbf{命题 130.8}（Ramsey 数的界）： -
\(R(r, s) \leq R(r-1, s) + R(r, s-1)\)（递归不等式） -
\(R(r, s) \leq \binom{r+s-2}{r-1}\)（Erdős-Szekeres，1935） -
\(R(r, r) \geq 2^{r/2}\)（Erdős 下界，概率方法，1947）

\textbf{已知 Ramsey
数}：\(R(3,3) = 6\)，\(R(3,4) = 9\)，\(R(3,5) = 14\)，\(R(3,6) = 18\)，\(R(4,4) = 18\)，\(R(3,9) = 36\)（2023年）。\(R(5,5)\)
仍未知，仅知 \(43 \leq R(5,5) \leq 48\)。

\textbf{定理 130.9}（Ramsey 定理，超图版本）：对任意正整数
\(r, k, c\)，存在 \(n\) 使得对 \(n\) 元集合的 \(r\) 元子集做
\(c\)-染色，必存在 \(k\) 元子集，其所有 \(r\) 元子集同色。

\textbf{定理 130.10}（Schur 定理，1916）：对任意 \(c\)，存在 \(N\)
使得对 \(\{1, \ldots, N\}\) 做 \(c\)-染色，必存在 \(x, y, z\) 同色且
\(x + y = z\)。

\emph{证明}：这是 Ramsey 定理的推论。取 \(N = R(3,3,\ldots,3)\)（\(c\)
个 3），对完全图 \(K_N\) 的边 \(\{i,j\}\)（\(i < j\)）染以 \(|i-j|\)
的颜色。由 Ramsey 定理，存在同色三角形，其边长为 \(i-j, j-k, i-k\)，满足
\((i-j)+(j-k) = i-k\)。∎

\textbf{定理 130.11}（van der Waerden 定理，1927）：对任意正整数
\(r, k\)，存在 \(W(r, k)\) 使得对 \(\{1, \ldots, W(r,k)\}\) 做
\(r\)-染色，必存在长度为 \(k\) 的同色等差数列。

\textbf{定理 130.12}（Szemerédi
定理，1975）：任何正上密度的自然数子集包含任意长的等差数列。即对任意
\(\delta > 0\) 和 \(k \in \mathbb{N}\)，存在 \(N\) 使得
\(\{1, \ldots, N\}\) 的任意大小 \(\geq \delta N\) 的子集包含长度为 \(k\)
的等差数列。

\emph{说明}：Szemerédi 定理的证明是组合数学的里程碑，Green-Tao
定理（2004）将此推广到素数中包含任意长等差数列。

\subsubsection{130.4 极值图论}\label{ux6781ux503cux56feux8bba}

\textbf{定理 130.13}（Turán 定理，1941）：\(n\) 个顶点不含 \(K_{r+1}\)
子图的最大边数为

\[\operatorname{ex}(n, K_{r+1}) = \left(1 - \frac{1}{r}\right) \frac{n^2}{2} + O(n)\]

极大值由 Turán 图 \(T(n, r)\)（将 \(n\) 个顶点尽可能均匀地划分为 \(r\)
部分，所有边连接不同部分的顶点）达到。

\emph{证明概要}：设 \(G\) 有 \(n\) 个顶点、\(m\) 条边且不含
\(K_{r+1}\)。对 \(n\) 归纳。取 \(G\) 中度最大的顶点
\(v\)，\(\deg(v) = \Delta\)。令
\(S = N(v)\)，\(T = V \setminus (\{v\} \cup S)\)，则
\(|S| = \Delta\)，\(|T| = n - 1 - \Delta\)。因 \(G\) 不含
\(K_{r+1}\)，\(S\) 的诱导子图不含 \(K_r\)。由归纳假设，\(S\) 中的边数
\(\leq \operatorname{ex}(\Delta, K_r)\)，\(T\) 中的边数
\(\leq \operatorname{ex}(n-1-\Delta, K_{r+1})\)。此外 \(S\) 与 \(T\)
之间及 \(v\) 与 \(S\) 以外的边可被上界估算。经最优配置分析，当 \(G\)
为完备 \(r\)-部图且各部大小差至多 1 时达到最大边数，即 Turán
图。边数精确公式为
\(\operatorname{ex}(n, K_{r+1}) = \frac{r-1}{2r}n^2 + O(n)\)。∎

\textbf{定理 130.14}（Erdős-Stone 定理，1946）：对任意图 \(H\)（色数
\(\chi(H) = r+1\)），

\[\operatorname{ex}(n, H) = \left(1 - \frac{1}{r}\right) \binom{n}{2} + o(n^2)\]

\bigskip\hrule\bigskip

\subsection{第131章：组合设计}\label{ux7b2c131ux7ae0ux7ec4ux5408ux8bbeux8ba1}

组合设计理论研究将元素安排成满足特定平衡条件的子集族。起源于统计学中的实验设计（Fisher,
1920s），在编码理论、密码学和计算机科学中有广泛应用。

\subsubsection{131.1 区组设计}\label{ux533aux7ec4ux8bbeux8ba1}

\textbf{定义 131.1}（平衡不完全区组设计 / BIBD）：一个
\((v, k, \lambda)\)-BIBD 是 \(v\) 元集 \(X\) 的一族 \(k\)
元子集（称为\textbf{区组}），使得 \(X\) 中任意两个不同元素恰出现在
\(\lambda\) 个区组中。

\textbf{命题 131.1}（BIBD 的必要条件）：若 \((v, k, \lambda)\)-BIBD
存在，参数满足： - \(r = \frac{\lambda(v-1)}{k-1}\)
是整数（每个元素出现的次数） - \(b = \frac{vr}{k}\) 是整数（区组总数） -
Fisher 不等式：\(b \geq v\)（即 \(r \geq k\)）

\textbf{例 131.1}（Fano 平面）：\((7, 3, 1)\)-BIBD，7 个点，7 条线，每线
3 点，每对点恰在 1 条线上。这是最小的射影平面（阶数为 2）。

\textbf{定义 131.2}（射影平面）：阶数为 \(n\) 的射影平面是
\((n^2+n+1, n+1, 1)\)-BIBD。等价于存在 \(n^2+n+1\)
个点和同样多条线，每线 \(n+1\) 点，每点 \(n+1\)
线，任意两点恰在一条线上，任意两线恰交于一点。已知当 \(n\)
为素数幂时存在，\(n=6\) 时不存在（Bruck-Ryser-Chowla 定理），\(n=10\)
时不存在（计算机搜索，1989）。

\textbf{定理 131.2}（Bruck-Ryser-Chowla 定理，1949）：若阶数为 \(n\)
的射影平面存在，且 \(n \equiv 1 \text{ 或 } 2 \pmod{4}\)，则 \(n\)
可表为两个整数的平方和。

\subsubsection{131.2 正交拉丁方}\label{ux6b63ux4ea4ux62c9ux4e01ux65b9}

\textbf{定义 131.3}（拉丁方）：\(n\) 阶\textbf{拉丁方}是 \(n \times n\)
矩阵，每行和每列都是 \(\{1, \ldots, n\}\)
的排列。两个拉丁方称为\textbf{正交}的，如果叠加后每对 \((i,j)\)
恰好出现一次。

\textbf{定理 131.3}（正交拉丁方与射影平面的关系）：\(n-1\) 个两两正交的
\(n\) 阶拉丁方（MOLS）的存在性等价于 \(n\) 阶射影平面的存在性。

\textbf{定理 131.4}（Euler 猜想与反例）：Euler 在 1782
年猜想不存在两个正交的 6 阶拉丁方（36 军官问题）。1901 年 Tarry
验证为真。Euler 进一步猜想对 \(n \equiv 2 \pmod{4}\)
均无正交拉丁方。1959 年 Bose-Shrikhande-Parker 推翻此猜想，构造了 10
阶正交拉丁方，并证明对 \(n \neq 2, 6\) 均存在。

\subsubsection{131.3 Steiner 系与
t-设计}\label{steiner-ux7cfbux4e0e-t-ux8bbeux8ba1}

\textbf{定义 131.4}（\(t\)-设计）：一个 \(t\)-\((v, k, \lambda)\) 设计是
\(v\) 元集的一族 \(k\) 元子集，使得任意 \(t\) 个不同元素恰出现在
\(\lambda\) 个区组中。BIBD 是 \(t=2\) 的特例。

\textbf{定义 131.5}（Steiner 系）：\(S(t, k, v)\) 是 \(t\)-\((v, k, 1)\)
设计。即任意 \(t\) 个元素恰出现在 1 个区组中。

\textbf{例 131.2}： - \(S(2, 3, v)\) 称为 Steiner
三元系（STS），存在当且仅当
\(v \equiv 1 \text{ 或 } 3 \pmod{6}\)（Kirkman, 1847） - \(S(5, 8, 24)\)
是著名的 Witt 设计，与 Mathieu 群 \(M_{24}\) 密切相关 - \(S(5, 6, 12)\)
是小的 Witt 设计，与 Mathieu 群 \(M_{12}\) 相关

\textbf{定理 131.5}（大设计猜想 / 已解决）：对
\(t \geq 6\)，仅存在平凡的 \(t\)-设计（即所有 \(k\) 元子集均取）。

\subsubsection{131.4
有限域与组合设计}\label{ux6709ux9650ux57dfux4e0eux7ec4ux5408ux8bbeux8ba1}

\textbf{定理 131.6}（有限域上的射影空间）：\(\mathbb{F}_q\) 上的 \(d\)
维射影空间 \(\operatorname{PG}(d, q)\) 中，点数为
\(\frac{q^{d+1}-1}{q-1}\)，每个 \(k\) 维子空间的点数为
\(\frac{q^{k+1}-1}{q-1}\)。\(\operatorname{PG}(2, q)\) 是 \(q\)
阶射影平面。

\textbf{定理 131.7}（差集与 Singer 定理）：\(\operatorname{PG}(d, q)\)
的自同构群在点上有一个正则循环子群（Singer
循环），对应一个\textbf{差集}------这是构造对称设计的标准方法。

\textbf{定义 131.6}（Hadamard 矩阵）：\(n\) 阶 Hadamard 矩阵 \(H\) 是
\(n \times n\) 的 \(\pm 1\) 矩阵，满足 \(H H^T = n I_n\)。Hadamard
矩阵存在当且仅当 \(n = 1, 2\) 或 \(n \equiv 0 \pmod{4}\)（Hadamard
猜想，未完全解决）。

\textbf{定理 131.8}（Hadamard 设计与 Hadamard
矩阵的等价性）：\((4m-1, 2m-1, m-1)\) 对称 BIBD 的存在性等价于 \(4m\) 阶
Hadamard 矩阵的存在性。

\bigskip\hrule\bigskip

\subsection{第132章：拟阵理论}\label{ux7b2c132ux7ae0ux62dfux9635ux7406ux8bba}

拟阵（Matroid）是由
Whitney（1935）引入的组合结构，统一了图论和线性代数中''独立性''的概念。一个拟阵
\(M = (E, \mathcal{I})\) 由有限集 \(E\) 和独立集族
\(\mathcal{I} \subseteq 2^E\)
组成，可通过四个等价公理刻画：（1）\textbf{独立集公理}------\(\mathcal{I}\)
满足遗传性（子集仍独立）和增广性（极大独立集大小相等）；（2）\textbf{基公理}------所有极大独立集（基）大小相同；（3）\textbf{秩函数公理}------\(r(A) = \max\{|S| : S \subseteq A, S \in \mathcal{I}\}\)
满足次模性
\(r(A \cup B) + r(A \cap B) \leq r(A) + r(B)\)；（4）\textbf{闭包公理}------闭包算子
\(\operatorname{cl}(A) = \{e : r(A \cup \{e\}) = r(A)\}\)
满足增性和交换性。经典例子包括图拟阵（边集的子集无回路即为独立）和向量拟阵（矩阵列向量的线性无关子集）。\textbf{拟阵交定理}（Edmonds,
1970）给出两个拟阵公共独立集极大基数的公式
\(\max |I| = \min_{A \subseteq E} (r_1(A) + r_2(E \setminus A))\)，是对
Kőnig 定理和最大流最小割定理的深远推广。

\subsubsection{132.1 设计理论}\label{ux8bbeux8ba1ux7406ux8bba}

平衡不完全区组设计（BIBD）是 \(v\) 元集中一族 \(k\)
元子集（区组），使得任意两个不同元素恰出现在 \(\lambda\)
个区组中，参数满足 \(\lambda(v-1) = r(k-1)\) 及 Fisher 不等式
\(b \geq v\)。\textbf{Steiner 三元系} \(S(2,3,v)\) 是最经典的
BIBD（\(\lambda=1, k=3\)），由 Kirkman（1847）证明存在当且仅当
\(v \equiv 1 \text{ 或 } 3 \pmod{6}\)。\textbf{有限射影平面}可组合地定义为
\((n^2+n+1, n+1, 1)\)-BIBD：每线含 \(n+1\) 点，每点过 \(n+1\)
线，任意两点确定唯一一线，两线恰交于一点------以纯粹的组合条件刻画了射影几何的核心结构。

\bigskip\hrule\bigskip

\subsection{第133章：代数组合学}\label{ux7b2c133ux7ae0ux4ee3ux6570ux7ec4ux5408ux5b66}

代数组合学以代数工具研究组合结构。\textbf{对称函数}理论的核心是 Schur
函数
\(s_\lambda\)，它们构成对称函数空间的线性基，其结构常数（Littlewood-Richardson
系数
\(c_{\mu\nu}^{\lambda}\)）给出对称群表示论中张量积分解的重数。\textbf{Young
表}是 Young 图（对应整数分拆
\(\lambda\)）的整数填充，满足每行非减、每列严格增。\textbf{Robinson-Schensted
对应}建立了置换 \(\pi \in S_n\) 与同形状的标准 Young 表对 \((P,Q)\)
之间的一一映射；\(\pi\) 的最长递增子序列长度等于 \(P\)
的第一行长度，最长的递减子序列长度等于第一列长度，这一结论（Schensted,
1961）深刻连接了组合学与表示论。\textbf{生成函数方法}通过普通生成函数
\(\sum a_n x^n\) 和指数生成函数 \(\sum a_n x^n/n!\)
将递推关系和组合恒等式转化为形式幂级数的代数运算，是组合计数中最基本且最强大的系统化工具。

\bigskip\hrule\bigskip

\subsection{第134章：极值集合论（Extremal Set
Theory）}\label{ux7b2c134ux7ae0ux6781ux503cux96c6ux5408ux8bbaextremal-set-theory}

极值集合论研究集合系在一定限制条件下能达到的最大尺寸或最优结构。\textbf{核心结果}：（1）\textbf{Sperner
定理}（1928年）------\(n\)
元集合的子集族中，若不存在两个集合有包含关系（即族是反链），则该族的最大尺寸为
\(\binom{n}{\lfloor n/2 \rfloor}\)（取自最中间层的所有子集）。这是极值集合论的奠基性结果，由对称链分解证明；（2）\textbf{Erdős-Ko-Rado
定理}（1961年）------设 \(\mathcal{F}\) 是 \(\{1,\ldots,n\}\)
上所有大小为 \(k\) 的子集组成的族，若 \(n \geq 2k\) 且 \(\mathcal{F}\)
中任意两个集合的交非空，则
\(|\mathcal{F}| \leq \binom{n-1}{k-1}\)（即包含某个固定元素的所有
\(k\)-子集）；（3）\textbf{Kruskal-Katona 定理}------在保持
\(k\)-一致超图边数的约束下，由初始段（colexicographic
order）给出影子（\((k-1)\)-子集）的最小基数；（4）\textbf{有限域上的极值集合论}------Frankl-Wilson
定理用代数方法（多项式空间的维数限制）证明集合族的 Ramsey 型上界，在
Frankl-Wilson 对 Borsuk
猜想的否定解决（1993年）中有决定性应用。\textbf{新兴方向}：\(L\)-交族在编码理论和
Ramsey 理论中的交叉应用；加性组合极值理论（Roth
定理的集合推广）；无标度网络中最大度分布的极值特征。

\bigskip\hrule\bigskip

\emph{本卷基于 V6 Ch 28（计数原理）和 V6 Ch
6（代数结构），为组合数学和图论建立了系统基础。共 7
章，涵盖生成函数、图论、匹配与网络流、染色与 Ramsey
理论、组合设计、拟阵理论和代数组合学。}

\newpage

\section{卷二十六：优化理论}\label{ux5377ux4e8cux5341ux516dux4f18ux5316ux7406ux8bba}

\begin{quote}
优化理论是数学在工程、经济和科学决策中的核心工具，研究在约束条件下最大化或最小化目标函数。本卷在初等线性规划（V2
Ch
11）和线性代数（V7）的基础上，建立线性规划的单纯形法和对偶理论，然后深入凸优化、非线性规划（KKT
条件）和变分法与最优控制引论。
\end{quote}

\bigskip\hrule\bigskip

\subsection{第132章：线性规划与单纯形法}\label{ux7b2c132ux7ae0ux7ebfux6027ux89c4ux5212ux4e0eux5355ux7eafux5f62ux6cd5}

线性规划（Linear Programming, LP）是优化理论中研究最成熟的分支，由
Dantzig 在 1947 年提出单纯形法（Simplex Method）而建立。

\subsubsection{132.1
线性规划的标准形式}\label{ux7ebfux6027ux89c4ux5212ux7684ux6807ux51c6ux5f62ux5f0f}

\textbf{定义 132.1}（线性规划的标准形式）：

\[\begin{aligned} \min \quad & \mathbf{c}^T \mathbf{x} \\ \text{s.t.} \quad & A\mathbf{x} = \mathbf{b} \\ & \mathbf{x} \geq \mathbf{0} \end{aligned}\]

其中
\(\mathbf{c} \in \mathbb{R}^n\)，\(A \in \mathbb{R}^{m \times n}\)，\(\mathbf{b} \in \mathbb{R}^m\)。\(m \leq n\)
且 \(\operatorname{rank}(A) = m\)。

\textbf{定义 132.2}（可行域）：可行域
\(\mathcal{F} = \{\mathbf{x} \in \mathbb{R}^n : A\mathbf{x} = \mathbf{b}, \mathbf{x} \geq \mathbf{0}\}\)
是\textbf{多面体}（polyhedron）。若 \(\mathcal{F}\)
有界，则称为\textbf{多胞体}（polytope）。

\textbf{定理 132.1}（线性规划的基本定理）：若 LP
有有限最优值，则存在最优解是 \(\mathcal{F}\) 的顶点（基本可行解）。

\emph{证明思路}：目标函数 \(\mathbf{c}^T \mathbf{x}\) 是线性的，在多面体
\(\mathcal{F}\) 上的最小值必然在极值点达到。若 \(\mathcal{F}\)
有界，最优解在顶点上；若无界，沿极方向移动不会使目标函数严格减小。∎

\textbf{定义 132.3}（基本可行解）：将 \(A\) 分块为
\(A = [B \mid N]\)，其中 \(B \in \mathbb{R}^{m \times m}\)
可逆（基矩阵），\(N\) 为非基列。对应的
\(\mathbf{x} = \begin{pmatrix} \mathbf{x}_B \\ \mathbf{x}_N \end{pmatrix}\)
满足
\(\mathbf{x}_B = B^{-1}\mathbf{b} \geq \mathbf{0}\)，\(\mathbf{x}_N = \mathbf{0}\)，称为\textbf{基本可行解}（BFS）。

\subsubsection{132.2 单纯形法}\label{ux5355ux7eafux5f62ux6cd5}

\textbf{算法 132.1}（单纯形法 / Simplex Method）：

\begin{enumerate}
\def\labelenumi{\arabic{enumi}.}
\tightlist
\item
  \textbf{初始化}：找到一个初始 BFS
\item
  \textbf{定价}（Pricing）：计算简约费用
  \(\bar{\mathbf{c}}_N = \mathbf{c}_N - N^T B^{-T} \mathbf{c}_B\)。若
  \(\bar{\mathbf{c}}_N \geq \mathbf{0}\)，当前解为最优解，停止
\item
  \textbf{进基}：选择 \(j\) 满足
  \(\bar{c}_j < 0\)（通常选择最负的），\(x_j\) 进入基
\item
  \textbf{离基}：计算 \(\mathbf{d} = B^{-1} A_j\)（\(A_j\)
  为进基列）。若 \(\mathbf{d} \leq \mathbf{0}\)，问题无界。否则选择
  \(\theta = \min_{d_i > 0} \frac{(\mathbf{x}_B)_i}{d_i}\)，令离基变量为达到最小比值的行对应的变量
\item
  \textbf{更新}：更新基矩阵 \(B\) 和 \(\mathbf{x}_B\)，返回步骤 2
\end{enumerate}

\textbf{定理 132.2}（单纯形法的有限终止性）：若 LP 非退化（所有 BFS
中基变量严格为正），则单纯形法在有限步内终止于最优解或判定无界。

\emph{退化处理}：当出现退化时，可能循环。Bland
法则（最小下标法则）保证有限终止。

\textbf{命题
132.3}（单纯形法的几何解释）：单纯形法沿着可行多面体的边移动，从当前顶点到相邻顶点，每次移动降低目标函数值。

\subsubsection{132.3
初始可行解的构造}\label{ux521dux59cbux53efux884cux89e3ux7684ux6784ux9020}

\textbf{定义 132.4}（两阶段法）： - \textbf{阶段
I}：引入人工变量，求解辅助 LP
\(\min \sum \mathbf{1}^T \mathbf{a}\)，s.t.
\(A\mathbf{x} + \mathbf{a} = \mathbf{b}, \mathbf{x} \geq 0, \mathbf{a} \geq 0\)。若最优值为
0，则得到原问题的可行 BFS - \textbf{阶段 II}：从阶段 I 得到的 BFS
出发，求解原问题

\textbf{定义 132.5}（大 M 法）：直接求解
\(\min \mathbf{c}^T \mathbf{x} + M \sum a_i\)，其中 \(M > 0\)
充分大，使得人工变量在最优解中必为 0。

\subsubsection{132.4 内点法}\label{ux5185ux70b9ux6cd5}

\textbf{定理 132.4}（Karmarkar 算法，1984）：存在求解 LP
的多项式时间算法（内点法），基于将可行域变换并沿梯度方向移动。

\textbf{定义 132.6}（障碍函数法）：考虑问题
\(\min \mathbf{c}^T \mathbf{x} - \mu \sum_{i=1}^n \ln x_i\) s.t.
\(A\mathbf{x} = \mathbf{b}\)。当 \(\mu \to 0\) 时，最优解趋近于原 LP
的最优解。这是\textbf{原-对偶内点法}的基础。

\bigskip\hrule\bigskip

\subsection{第133章：对偶理论与灵敏度分析}\label{ux7b2c133ux7ae0ux5bf9ux5076ux7406ux8bbaux4e0eux7075ux654fux5ea6ux5206ux6790}

对偶理论是线性规划中与每个原问题（Primal）自然关联的另一个线性规划问题（Dual），两者之间有深刻的关系。

\subsubsection{133.1
对偶问题的构造}\label{ux5bf9ux5076ux95eeux9898ux7684ux6784ux9020}

\textbf{定义 133.1}（对偶问题）：对标准形式 LP（Primal）：

\[\begin{aligned} \min \quad & \mathbf{c}^T \mathbf{x} \\ \text{s.t.} \quad & A\mathbf{x} = \mathbf{b}, \quad \mathbf{x} \geq \mathbf{0} \end{aligned}\]

其对偶问题（Dual）为：

\[\begin{aligned} \max \quad & \mathbf{b}^T \mathbf{y} \\ \text{s.t.} \quad & A^T \mathbf{y} \leq \mathbf{c} \end{aligned}\]

其中 \(\mathbf{y} \in \mathbb{R}^m\) 为对偶变量。

\textbf{定理 133.1}（弱对偶定理）：若 \(\mathbf{x}\)
是原问题可行解，\(\mathbf{y}\) 是对偶问题可行解，则
\(\mathbf{c}^T \mathbf{x} \geq \mathbf{b}^T \mathbf{y}\)。

\emph{证明}：\(\mathbf{c}^T \mathbf{x} \geq (A^T \mathbf{y})^T \mathbf{x} = \mathbf{y}^T A \mathbf{x} = \mathbf{y}^T \mathbf{b} = \mathbf{b}^T \mathbf{y}\)。∎

\textbf{定理 133.2}（强对偶定理）：若原问题有有限最优解
\(\mathbf{x}^*\)，则对偶问题也有有限最优解 \(\mathbf{y}^*\)，且
\(\mathbf{c}^T \mathbf{x}^* = \mathbf{b}^T \mathbf{y}^*\)。

\emph{证明思路}：单纯形法终止时，基矩阵 \(B\) 满足
\(\bar{\mathbf{c}}_N \geq 0\)。令
\(\mathbf{y}^* = B^{-T} \mathbf{c}_B\)，则 \(\mathbf{y}^*\)
是对偶可行解，且目标值相等。由弱对偶定理，\(\mathbf{y}^*\) 是最优的。∎

\textbf{定理 133.3}（互补松弛定理）：设 \(\mathbf{x}^*\)
为原问题最优解，\(\mathbf{y}^*\) 为对偶问题最优解，则

\[x_j^* \cdot (\mathbf{c}_j - \mathbf{a}_j^T \mathbf{y}^*) = 0, \quad \forall j\]

即原变量和对偶松弛变量不能同时为正。

\textbf{推论 133.4}：线性规划的可能情况（原问题/对偶问题）： -
原问题有有限最优解 \(\iff\) 对偶问题有有限最优解 - 原问题无界
\(\implies\) 对偶问题不可行 - 原问题不可行 \(\implies\)
对偶问题无界或不可行

\subsubsection{133.2
对偶的经济解释}\label{ux5bf9ux5076ux7684ux7ecfux6d4eux89e3ux91ca}

\textbf{定义 133.2}（影子价格）：对偶变量 \(y_i\) 是第 \(i\)
个约束资源的\textbf{影子价格}（shadow price），即增加一单位资源 \(b_i\)
时最优值的边际改进。

\textbf{定理 133.5}（灵敏度分析）：若基矩阵 \(B\)
的最优性不受扰动影响，则
\(\Delta z^* = \mathbf{y}^{*T} \Delta \mathbf{b}\)。

\subsubsection{133.3 对偶在非 LP
情形的推广}\label{ux5bf9ux5076ux5728ux975e-lp-ux60c5ux5f62ux7684ux63a8ux5e7f}

\textbf{定义 133.3}（Lagrange 对偶）：对一般约束优化问题

\[\min_{\mathbf{x}} f(\mathbf{x}) \quad \text{s.t.} \quad g_i(\mathbf{x}) \leq 0, \; h_j(\mathbf{x}) = 0\]

定义 Lagrange 函数
\(L(\mathbf{x}, \boldsymbol{\lambda}, \boldsymbol{\nu}) = f(\mathbf{x}) + \sum \lambda_i g_i(\mathbf{x}) + \sum \nu_j h_j(\mathbf{x})\)，其中
\(\boldsymbol{\lambda} \geq 0\)。Lagrange 对偶函数为
\(d(\boldsymbol{\lambda}, \boldsymbol{\nu}) = \inf_{\mathbf{x}} L(\mathbf{x}, \boldsymbol{\lambda}, \boldsymbol{\nu})\)。

\textbf{定理 133.6}（弱对偶，一般情形）：对任意
\(\boldsymbol{\lambda} \geq 0\)，\(d(\boldsymbol{\lambda}, \boldsymbol{\nu}) \leq p^*\)（原问题最优值）。对偶间隙为
\(p^* - d^* \geq 0\)。

\bigskip\hrule\bigskip

\subsection{第134章：凸优化}\label{ux7b2c134ux7ae0ux51f8ux4f18ux5316}

凸优化是优化理论中最重要的子领域，其核心优势在于：凸问题的局部最优解就是全局最优解。本章从凸集和凸函数的基本概念出发，建立凸优化问题的标准形式及对偶理论（KKT
充要条件），为后续非线性规划和变分法中的最优性条件奠定几何基础。

\subsubsection{134.1
凸集与凸函数}\label{ux51f8ux96c6ux4e0eux51f8ux51fdux6570}

\textbf{定义 134.1}（凸集）：集合 \(C\) 是\textbf{凸集}，如果对任意
\(\mathbf{x}, \mathbf{y} \in C\) 和 \(\theta \in [0, 1]\)，有
\(\theta \mathbf{x} + (1-\theta)\mathbf{y} \in C\)。

\textbf{例 134.1}（常见凸集）： - 超平面
\(\{\mathbf{x} : \mathbf{a}^T \mathbf{x} = b\}\) 和半空间
\(\{\mathbf{x} : \mathbf{a}^T \mathbf{x} \leq b\}\) - 多面体
\(\{\mathbf{x} : A\mathbf{x} \preceq \mathbf{b}\}\) - 范数球
\(\{\mathbf{x} : \|\mathbf{x}\| \leq r\}\) - 半正定锥
\(\mathbb{S}_+^n = \{X \in \mathbb{S}^n : X \succeq 0\}\)

\textbf{定义
134.2}（保凸运算）：凸集的交是凸集；凸集在仿射映射下的像和原像都是凸集；两个凸集的
Minkowski 和是凸集。

\textbf{定义 134.3}（凸函数）：函数 \(f: \mathbb{R}^n \to \mathbb{R}\)
是\textbf{凸函数}，如果定义域 \(\operatorname{dom} f\) 是凸集，且对任意
\(\mathbf{x}, \mathbf{y} \in \operatorname{dom} f\) 和
\(\theta \in [0, 1]\)，

\[f(\theta \mathbf{x} + (1-\theta)\mathbf{y}) \leq \theta f(\mathbf{x}) + (1-\theta) f(\mathbf{y})\]

若严格不等式成立（\(\mathbf{x} \neq \mathbf{y}, \theta \in (0,1)\)），则称为\textbf{严格凸函数}。

\textbf{定理 134.1}（凸函数的判定与无约束极小）： 1.
\textbf{一阶条件}：若 \(f\) 可微，\(f\) 为凸 \(\iff\)
\(f(\mathbf{y}) \geq f(\mathbf{x}) + \nabla f(\mathbf{x})^T (\mathbf{y} - \mathbf{x})\)
对所有 \(\mathbf{x}, \mathbf{y}\) 成立 2. \textbf{二阶条件}：若 \(f\)
二阶可微，\(f\) 为凸 \(\iff\)
\(\nabla^2 f(\mathbf{x}) \succeq 0\)（Hessian 半正定）对所有
\(\mathbf{x}\) 成立

\emph{无约束极小的充要条件}（次梯度条件）：设 \(f\) 为凸函数，则
\(\mathbf{x}^*\) 是全局极小点当且仅当
\(\mathbf{0} \in \partial f(\mathbf{x}^*)\)。其中
\(\partial f(\mathbf{x}^*) = \{\mathbf{g} : f(\mathbf{y}) \geq f(\mathbf{x}^*) + \mathbf{g}^T(\mathbf{y} - \mathbf{x}^*), \forall \mathbf{y}\}\)
为次微分。若 \(f\) 在 \(\mathbf{x}^*\) 处可微，则此条件退化为
\(\nabla f(\mathbf{x}^*) = \mathbf{0}\)。证明：充分性由次梯度定义直接得
\(f(\mathbf{y}) \geq f(\mathbf{x}^*)\)，必要性由凸分离定理。∎

\textbf{例 134.2}（常见凸函数）：\(e^{ax}\)，\(x^a\)（\(a \geq 1\) 或
\(a \leq 0\)），\(-\ln x\)，\(\|x\|\)（任意范数），\(\max\{x_1, \ldots, x_n\}\)，\(\log(e^{x_1} + \cdots + e^{x_n})\)（log-sum-exp）。

\subsubsection{134.2
凸优化问题的标准形式}\label{ux51f8ux4f18ux5316ux95eeux9898ux7684ux6807ux51c6ux5f62ux5f0f}

\textbf{定义 134.4}（凸优化问题）：

\[\begin{aligned} \min \quad & f_0(\mathbf{x}) \\ \text{s.t.} \quad & f_i(\mathbf{x}) \leq 0, \quad i = 1, \ldots, m \\ & A\mathbf{x} = \mathbf{b} \end{aligned}\]

其中 \(f_0, f_1, \ldots, f_m\) 是凸函数。可行域是凸集。

\textbf{定理
134.2}（凸优化的基本性质）：凸优化问题的任何局部最优解都是全局最优解。最优解集是凸集。

\emph{证明}：设 \(\mathbf{x}^*\) 是局部最优解。若存在 \(\mathbf{y}\)
满足 \(f_0(\mathbf{y}) < f_0(\mathbf{x}^*)\)，则对充分小的
\(\theta > 0\)，\(\theta \mathbf{y} + (1-\theta)\mathbf{x}^*\) 在
\(\mathbf{x}^*\) 的邻域内，且
\(f_0(\theta \mathbf{y} + (1-\theta)\mathbf{x}^*) \leq \theta f_0(\mathbf{y}) + (1-\theta)f_0(\mathbf{x}^*) < f_0(\mathbf{x}^*)\)，矛盾。∎

\textbf{例 134.3}（经典凸优化问题）： - \textbf{线性规划}（LP）：\(f_i\)
均为仿射函数 - \textbf{二次规划}（QP）：\(f_0\) 为二次凸函数，\(f_i\)
为仿射函数 -
\textbf{二阶锥规划}（SOCP）：\(\min \mathbf{c}^T \mathbf{x}\) s.t.
\(\|A_i \mathbf{x} + \mathbf{b}_i\|_2 \leq \mathbf{c}_i^T \mathbf{x} + d_i\)
- \textbf{半正定规划}（SDP）：\(\min \mathbf{c}^T \mathbf{x}\) s.t.
\(F_0 + \sum x_i F_i \succeq 0\)（\(F_i\) 为对称矩阵）

上述问题有嵌套关系：\(\text{LP} \subset \text{QP} \subset \text{SOCP} \subset \text{SDP}\)。

\subsubsection{134.3
对偶理论在凸优化中的应用}\label{ux5bf9ux5076ux7406ux8bbaux5728ux51f8ux4f18ux5316ux4e2dux7684ux5e94ux7528}

\textbf{定理 134.3}（Slater 条件）：若凸优化问题存在严格可行点（即
\(\mathbf{x} \in \operatorname{relint} \mathcal{D}\) 满足
\(f_i(\mathbf{x}) < 0\)
对所有不等式约束，\(A\mathbf{x} = \mathbf{b}\)），则强对偶成立（对偶间隙为
0）。

\textbf{定理 134.4}（KKT
条件是最优性的充要条件，凸情形）：若凸优化问题满足 Slater 条件，则
\(\mathbf{x}^*\) 是最优解当且仅当存在
\(\boldsymbol{\lambda}^*, \boldsymbol{\nu}^*\) 满足 KKT 条件（见第 135
章）。

\bigskip\hrule\bigskip

\subsection{第135章：非线性规划}\label{ux7b2c135ux7ae0ux975eux7ebfux6027ux89c4ux5212}

非线性规划（NLP）处理目标函数或约束为非线性的一般优化问题。Karush-Kuhn-Tucker（KKT）条件是一阶必要条件，在凸情形下也是充分条件。

\subsubsection{135.1 无约束优化}\label{ux65e0ux7ea6ux675fux4f18ux5316}

\textbf{定理 135.1}（一阶必要条件）：若 \(\mathbf{x}^*\) 是无约束问题
\(\min f(\mathbf{x})\) 的局部极小点，且 \(f\) 在 \(\mathbf{x}^*\)
处可微，则 \(\nabla f(\mathbf{x}^*) = \mathbf{0}\)。

\textbf{定理 135.2}（二阶必要条件）：若 \(\mathbf{x}^*\)
是局部极小点，且 \(f\) 在 \(\mathbf{x}^*\) 处二阶可微，则
\(\nabla^2 f(\mathbf{x}^*) \succeq 0\)。

\textbf{定理 135.3}（二阶充分条件）：若
\(\nabla f(\mathbf{x}^*) = \mathbf{0}\) 且
\(\nabla^2 f(\mathbf{x}^*) \succ 0\)，则 \(\mathbf{x}^*\)
是严格局部极小点。

\textbf{算法
135.1}（梯度下降法）：\(\mathbf{x}_{k+1} = \mathbf{x}_k - \alpha_k \nabla f(\mathbf{x}_k)\)，其中
\(\alpha_k\) 为步长（由线搜索或固定值确定）。

\textbf{算法 135.2}（Newton
法）：\(\mathbf{x}_{k+1} = \mathbf{x}_k - [\nabla^2 f(\mathbf{x}_k)]^{-1} \nabla f(\mathbf{x}_k)\)。在初始点充分接近最优解时，Newton
法具有二次收敛速度。

\textbf{定理 135.4}（Newton 法的收敛速度）：若 \(f\) 是强凸函数且
Hessian Lipschitz 连续，则 Newton
法局部二次收敛：\(\|\mathbf{x}_{k+1} - \mathbf{x}^*\| \leq C \|\mathbf{x}_k - \mathbf{x}^*\|^2\)。

\textbf{算法 135.3}（拟 Newton 法 / BFGS）：维护 Hessian 近似
\(B_k\)，更新公式：

\[B_{k+1} = B_k - \frac{B_k \mathbf{s}_k \mathbf{s}_k^T B_k}{\mathbf{s}_k^T B_k \mathbf{s}_k} + \frac{\mathbf{y}_k \mathbf{y}_k^T}{\mathbf{y}_k^T \mathbf{s}_k}\]

其中
\(\mathbf{s}_k = \mathbf{x}_{k+1} - \mathbf{x}_k\)，\(\mathbf{y}_k = \nabla f(\mathbf{x}_{k+1}) - \nabla f(\mathbf{x}_k)\)。

\subsubsection{135.2 约束优化与 KKT
条件}\label{ux7ea6ux675fux4f18ux5316ux4e0e-kkt-ux6761ux4ef6}

\textbf{定义 135.1}（一般约束优化问题）：

\[\begin{aligned} \min \quad & f(\mathbf{x}) \\ \text{s.t.} \quad & g_i(\mathbf{x}) \leq 0, \quad i = 1, \ldots, m \\ & h_j(\mathbf{x}) = 0, \quad j = 1, \ldots, p \end{aligned}\]

\textbf{定理 135.5}（KKT 必要条件，Karush-Kuhn-Tucker，1951）：设
\(\mathbf{x}^*\) 是局部极小点，\(f, g_i, h_j\)
连续可微，且约束规范条件（如
LICQ：\(\nabla g_i(\mathbf{x}^*)\)（活跃约束）与
\(\nabla h_j(\mathbf{x}^*)\) 线性无关）成立。则存在 Lagrange 乘子
\(\boldsymbol{\lambda}^* \geq \mathbf{0}\) 和 \(\boldsymbol{\nu}^*\)
满足：

\[\begin{aligned}
\nabla f(\mathbf{x}^*) + \sum_{i=1}^m \lambda_i^* \nabla g_i(\mathbf{x}^*) + \sum_{j=1}^p \nu_j^* \nabla h_j(\mathbf{x}^*) &= \mathbf{0} \quad &\text{(稳定性)} \\
g_i(\mathbf{x}^*) &\leq 0, \quad i = 1, \ldots, m \quad &\text{(原可行性)} \\
h_j(\mathbf{x}^*) &= 0, \quad j = 1, \ldots, p \quad &\text{(原可行性)} \\
\lambda_i^* &\geq 0, \quad i = 1, \ldots, m \quad &\text{(对偶可行性)} \\
\lambda_i^* g_i(\mathbf{x}^*) &= 0, \quad i = 1, \ldots, m \quad &\text{(互补松弛性)}
\end{aligned}\]

\textbf{定理 135.6}（KKT 充分条件，凸情形）：若 \(f\) 和 \(g_i\)
是凸函数，\(h_j\) 是仿射函数，且 \(\mathbf{x}^*\) 和
\((\boldsymbol{\lambda}^*, \boldsymbol{\nu}^*)\) 满足 KKT 条件，则
\(\mathbf{x}^*\) 是全局最优解。

\subsubsection{135.3
约束规范条件}\label{ux7ea6ux675fux89c4ux8303ux6761ux4ef6}

\textbf{定义 135.2}（常见约束规范 / Constraint Qualifications）： -
\textbf{LICQ}（线性独立约束规范）：活跃约束的梯度线性无关（最强） -
\textbf{MFCQ}（Mangasarian-Fromovitz
CQ）：等式约束梯度线性无关，且存在方向使所有不等式约束严格减小 -
\textbf{Slater 条件}（凸情形）：存在严格可行点

\textbf{命题 135.7}：LICQ \(\implies\) MFCQ \(\implies\) Abadie
CQ。在凸情形下，Slater 条件 \(\implies\) KKT 条件成立。

\subsubsection{135.4 罚函数法与增广 Lagrange
法}\label{ux7f5aux51fdux6570ux6cd5ux4e0eux589eux5e7f-lagrange-ux6cd5}

\textbf{定义 135.3}（外罚函数法）：将约束优化转化为无约束问题：

\[\min_{\mathbf{x}} f(\mathbf{x}) + \frac{\rho}{2} \left( \sum_{i=1}^m \max\{0, g_i(\mathbf{x})\}^2 + \sum_{j=1}^p h_j(\mathbf{x})^2 \right)\]

当 \(\rho \to \infty\) 时，解趋近于原约束问题的最优解。

\textbf{定义 135.4}（增广 Lagrange 函数 / Augmented Lagrangian）：

\[L_A(\mathbf{x}, \boldsymbol{\lambda}, \boldsymbol{\nu}; \rho) = f(\mathbf{x}) + \sum_{i=1}^m \lambda_i g_i(\mathbf{x}) + \sum_{j=1}^p \nu_j h_j(\mathbf{x}) + \frac{\rho}{2}\left( \sum_{i=1}^m g_i(\mathbf{x})^2 + \sum_{j=1}^p h_j(\mathbf{x})^2 \right)\]

（仅对等式约束使用二次罚项，不等式约束可通过引入松弛变量转化为等式约束）

\textbf{算法 135.4}（增广 Lagrange 法 / 乘子法）： 1. 初始化
\(\boldsymbol{\lambda}^0, \boldsymbol{\nu}^0, \rho_0\) 2. 求解
\(\mathbf{x}^{k+1} = \arg\min_{\mathbf{x}} L_A(\mathbf{x}, \boldsymbol{\lambda}^k, \boldsymbol{\nu}^k; \rho_k)\)
3.
更新乘子：\(\lambda_i^{k+1} = \max\{0, \lambda_i^k + \rho_k g_i(\mathbf{x}^{k+1})\}\)，\(\nu_j^{k+1} = \nu_j^k + \rho_k h_j(\mathbf{x}^{k+1})\)
4. 增大 \(\rho_{k+1} > \rho_k\)，重复

\bigskip\hrule\bigskip

\subsection{第136章：变分法与最优控制引论}\label{ux7b2c136ux7ae0ux53d8ux5206ux6cd5ux4e0eux6700ux4f18ux63a7ux5236ux5f15ux8bba}

变分法研究泛函（函数的函数）的极值问题，起源于 Bernoulli
的最速降线问题（1696年）和 Euler-Lagrange
的工作。最优控制是变分法的现代推广，由 Pontryagin 极大值原理和 Bellman
动态规划奠定。

\subsubsection{136.1 Euler-Lagrange
方程}\label{euler-lagrange-ux65b9ux7a0b}

\textbf{定义 136.1}（最简泛函）：考虑泛函

\[J[y] = \int_a^b F(x, y(x), y'(x)) \, dx\]

其中 \(y(a) = \alpha\)，\(y(b) = \beta\) 固定。

\textbf{定理 136.1}（Euler-Lagrange 方程）：若 \(y(x)\) 使 \(J[y]\)
取极值，且 \(F\) 充分光滑，则 \(y\) 满足

\[\frac{\partial F}{\partial y} - \frac{d}{dx}\left( \frac{\partial F}{\partial y'} \right) = 0\]

\emph{证明思路}：考虑变分 \(\tilde{y} = y + \varepsilon \eta(x)\)，其中
\(\eta(a) = \eta(b) = 0\)。计算一阶变分：
\[\frac{d}{d\varepsilon} J[y + \varepsilon \eta]\bigg|_{\varepsilon=0} = \int_a^b \left( \frac{\partial F}{\partial y} \eta + \frac{\partial F}{\partial y'} \eta' \right) dx = 0\]
对第二项分部积分：\(\int_a^b \frac{\partial F}{\partial y'} \eta' dx = \left[ \frac{\partial F}{\partial y'} \eta \right]_a^b - \int_a^b \frac{d}{dx}\left( \frac{\partial F}{\partial y'} \right) \eta dx\)。由
\(\eta(a) = \eta(b) = 0\)，边界项为零，代入得：
\[\int_a^b \left( \frac{\partial F}{\partial y} - \frac{d}{dx} \frac{\partial F}{\partial y'} \right) \eta(x) dx = 0\]
由变分法基本引理（\(\eta\) 任意性），被积函数恒为零，即得 Euler-Lagrange
方程。∎

\textbf{例 136.1}（最速降线 / Brachistochrone 问题）：求从 \((0,0)\) 到
\((x_1, y_1)\)
的曲线，使质点在重力作用下用时最短。\(F = \sqrt{\frac{1+(y')^2}{2gy}}\)。解为摆线（cycloid）。

\textbf{例 136.2}（极小曲面 /
肥皂膜）：求连接两个圆环的最小面积旋转曲面。\(F = y\sqrt{1+(y')^2}\)。解为悬链面（catenoid）。

\textbf{定义 136.2}（Beltrami 恒等式）：若 \(F\) 不显含 \(x\)（即
\(\frac{\partial F}{\partial x} = 0\)），则 Euler-Lagrange
方程有首次积分

\[F - y' \frac{\partial F}{\partial y'} = C \quad (\text{常数})\]

\subsubsection{136.2
多变量与高阶变分}\label{ux591aux53d8ux91cfux4e0eux9ad8ux9636ux53d8ux5206}

\textbf{定义 136.3}（多变量泛函）：对
\(J[u] = \iint_\Omega F(x, y, u, u_x, u_y) \, dx dy\)，Euler-Lagrange
方程为

\[\frac{\partial F}{\partial u} - \frac{\partial}{\partial x}\left( \frac{\partial F}{\partial u_x} \right) - \frac{\partial}{\partial y}\left( \frac{\partial F}{\partial u_y} \right) = 0\]

\textbf{例 136.3}（Dirichlet
积分）：\(F = \frac{1}{2}(u_x^2 + u_y^2)\)，Euler-Lagrange 方程给出
Laplace 方程 \(\Delta u = 0\)。

\textbf{定理 136.2}（Hamilton 原理 /
最小作用量原理）：经典力学系统的运动方程可由作用量泛函
\(S[q] = \int_{t_1}^{t_2} L(q, \dot{q}, t) \, dt\) 的变分
\(\delta S = 0\) 导出，其中 \(L = T - V\) 为 Lagrange 函数。

\textbf{定义 136.4}（二次变分与 Legendre 条件）：极小的二阶必要条件为
Legendre 条件
\(\frac{\partial^2 F}{\partial y'^2} \geq 0\)。充分条件为强化的 Legendre
条件 \(\frac{\partial^2 F}{\partial y'^2} > 0\) 加上 Jacobi
条件（无共轭点）。

\subsubsection{136.3
约束变分问题与等周问题}\label{ux7ea6ux675fux53d8ux5206ux95eeux9898ux4e0eux7b49ux5468ux95eeux9898}

\textbf{定义 136.5}（等周问题 / Isoperimetric Problem）：在积分约束
\(\int_a^b G(x, y, y') \, dx = L\)（常数）下，最大化或最小化
\(\int_a^b F(x, y, y') \, dx\)。

\textbf{定理 136.3}（等周问题的 Euler-Lagrange 方程）：定义辅助泛函
\(H = F + \lambda G\)（\(\lambda\) 为 Lagrange 乘子），则极值曲线满足

\[\frac{\partial H}{\partial y} - \frac{d}{dx}\left( \frac{\partial H}{\partial y'} \right) = 0\]

\textbf{例 136.4}（Dido
问题）：固定周长，求最大面积的闭曲线。解为圆（等周不等式）。

\subsubsection{136.4
最优控制初步}\label{ux6700ux4f18ux63a7ux5236ux521dux6b65}

\textbf{定义 136.6}（最优控制问题）：

\[\begin{aligned} \min_{u(t)} \quad & \phi(x(T)) + \int_0^T L(x(t), u(t), t) \, dt \\ \text{s.t.} \quad & \dot{x}(t) = f(x(t), u(t), t), \quad x(0) = x_0 \end{aligned}\]

其中 \(x(t) \in \mathbb{R}^n\) 为状态变量，\(u(t) \in \mathbb{R}^m\)
为控制变量。

\textbf{定理 136.4}（Pontryagin 极大值原理，1962）：定义 Hamilton 函数
\(H(x, u, p, t) = L(x, u, t) + p^T f(x, u, t)\)，其中 \(p(t)\)
是协态变量（adjoint variable）。最优控制 \(u^*(t)\) 和最优轨线
\(x^*(t)\) 满足：

\[\begin{aligned}
\dot{x}^* &= \frac{\partial H}{\partial p} = f(x^*, u^*, t) \\
\dot{p} &= -\frac{\partial H}{\partial x} \\
H(x^*, u^*, p, t) &= \max_{u \in \mathcal{U}} H(x^*, u, p, t) \quad \text{(极大值条件)} \\
p(T) &= \nabla \phi(x^*(T)) \quad \text{(横截条件)}
\end{aligned}\]

\textbf{定理 136.5}（Bellman 最优性原理与 HJB 方程）：定义值函数
\(V(x, t) = \min_{u(\cdot)} \left[ \phi(x(T)) + \int_t^T L(x, u, \tau) d\tau \right]\)。则
\(V\) 满足 Hamilton-Jacobi-Bellman（HJB）方程：

\[-\frac{\partial V}{\partial t} = \min_{u \in \mathcal{U}} \left[ L(x, u, t) + \nabla_x V \cdot f(x, u, t) \right]\]

边界条件 \(V(x, T) = \phi(x)\)。

\textbf{定义 136.7}（线性二次型调节器 / LQR）：当系统为线性
\(\dot{x} = Ax + Bu\)，目标为二次型
\(J = \frac{1}{2}x(T)^T Q_f x(T) + \frac{1}{2}\int_0^T (x^T Q x + u^T R u) dt\)
时，最优控制为状态反馈 \(u^* = -R^{-1} B^T P(t) x\)，其中 \(P(t)\) 满足
Riccati 微分方程 \(-\dot{P} = A^T P + PA - PBR^{-1}B^T P + Q\)。

\bigskip\hrule\bigskip

\subsection{第137章：互补性问题与变分不等式}\label{ux7b2c137ux7ae0ux4e92ux8865ux6027ux95eeux9898ux4e0eux53d8ux5206ux4e0dux7b49ux5f0f}

互补性问题是优化理论中的一个重要分支，其核心是线性互补问题\(\operatorname{LCP}(q, M)\)：给定向量\(q \in \mathbb{R}^n\)和矩阵\(M \in \mathbb{R}^{n \times n}\)，求\(z \in \mathbb{R}^n\)使得\(z \geq 0\)，\(Mz + q \geq 0\)，且\(z^T(Mz + q) = 0\)。\(\operatorname{LCP}\)与多个数学领域紧密关联：二次规划\(\min\{\frac{1}{2}x^TQx + c^Tx : Ax \geq b, x \geq 0\}\)的KKT条件即为\(\operatorname{LCP}\)；双矩阵博弈的Nash均衡也可化为\(\operatorname{LCP}\)。求解\(\operatorname{LCP}\)的经典算法是Lemke算法（1965），通过引入人工变量\(z_0\)构造几乎互补基，沿射线进行主元迭代直到\(z_0\)出基。Lemke算法在\(M\)为P-矩阵（所有主子式为正）时保证收敛。变分不等式\(\operatorname{VI}(K, F)\)求\(x^* \in K\)使\(\langle F(x^*), x - x^* \rangle \geq 0\)对所有\(x \in K\)成立，是\(\operatorname{LCP}\)的无穷维推广，在均衡分析和交通分配中有广泛应用。

\bigskip\hrule\bigskip

\subsection{第137章：半定规划与离散凸分析}\label{ux7b2c137ux7ae0ux534aux5b9aux89c4ux5212ux4e0eux79bbux6563ux51f8ux5206ux6790}

\subsubsection{137.1 半定规划}\label{ux534aux5b9aux89c4ux5212}

半定规划（Semidefinite Programming,
SDP）是在线性矩阵不等式约束下极小化线性目标函数的凸优化问题：

\[\min \; \mathbf{c}^T \mathbf{x} \quad \text{s.t.} \quad F_0 + \sum_{i=1}^n x_i F_i \succeq 0\]

其中 \(F_i\) 为 \(m \times m\) 对称矩阵，\(X \succeq 0\) 表示 \(X\)
半正定。SDP 包含线性规划（LP, \(F_i\)
为对角矩阵时）和二阶锥规划（SOCP）作为特例，严格包含关系为
\(\text{LP} \subset \text{SOCP} \subset \text{SDP}\)。内点法可在多项式时间内求解
SDP（至给定精度）。在组合优化中 SDP
松弛有重要突破：Goemans-Williamson（1995）将 MAX-CUT
问题表示为整数二次规划，通过 SDP 松弛（\(X \succeq 0\) 代替
\(X = \mathbf{v}\mathbf{v}^T\)）并使用随机超平面舍入，获得了约 0.878
的近似比------这是基于 SDP
的第一个突破性近似算法，极大地推动了半定规划在近似算法设计中的应用。

\subsubsection{137.2 离散凸分析}\label{ux79bbux6563ux51f8ux5206ux6790}

离散凸分析由室田一雄（Murota）于 1990 年代建立，将凸分析推广到离散格点
\(\mathbb{Z}^n\)。核心概念包括两类互为共轭的函数族：（1）\textbf{L\(^\natural\)-凸函数}------满足离散中点凸性
\(f(x) + f(y) \geq f(\lfloor\frac{x+y}{2}\rfloor) + f(\lceil\frac{x+y}{2}\rceil)\)，其水平集为
L\(^\natural\)-凸集；（2）\textbf{M\(^\natural\)-凸函数}------基于交换性质的组合凸性，函数满足
\(f(x) + f(y) \geq f(x - \mathbf{1}_A) + f(y + \mathbf{1}_A)\)
类型的交换不等式。它们分别是连续凸分析中凸函数对不同离散结构的适配，且通过
Legendre-Fenchel
变换互为共轭。应用涵盖经济学中不可分商品的竞争均衡存在性（Kelso-Crawford
的劳动市场模型）和运筹学中的资源配置与库存管理。

\bigskip\hrule\bigskip

\subsection{第138章：动态规划与随机规划}\label{ux7b2c138ux7ae0ux52a8ux6001ux89c4ux5212ux4e0eux968fux673aux89c4ux5212}

\subsubsection{138.1 动态规划（Dynamic
Programming）}\label{ux52a8ux6001ux89c4ux5212dynamic-programming}

动态规划（Bellman,
1957）通过将多阶段决策问题分解为子问题的递归关系解决，奠基于
\textbf{Bellman
最优性原理}------最优策略具有这样的性质：无论初始状态和初始决策如何，剩余决策必须构成由第一个决策导致的状态的最优策略。

\textbf{定义 138.1}（确定性有限期动态规划）：对 \(N\)
阶段确定性决策问题，值函数 \(V_k(x)\) 定义为当前状态为 \(x\)、还剩 \(k\)
步时的最优费用，满足 \textbf{Bellman 方程}

\[V_k(x) = \min_{u \in U(x)} \left[ c(x, u) + V_{k-1}\big(f(x, u)\big) \right], \quad k = 1, 2, \ldots, N\]

其中 \(c(x, u)\) 为瞬时代价，\(f(x, u)\) 为状态转移函数，\(U(x)\) 为状态
\(x\) 的可行控制集。边界条件 \(V_0(x) = g(x)\)（终端代价）。

\textbf{定义 138.2}（无限期折扣问题）：引入折扣因子
\(\gamma \in (0, 1)\)，目标为最小化期望折现总代价
\(\mathbb{E}\left[ \sum_{t=0}^\infty \gamma^t c(x_t, u_t) \right]\)。Bellman
方程化为不动点方程

\[V(x) = \min_{u \in U(x)} \left[ c(x, u) + \gamma V\big(f(x, u)\big) \right]\]

\textbf{算法 138.1}（值迭代 / Value
Iteration）：\(V_{n+1}(x) = \min_u [c(x,u) + \gamma V_n(f(x,u))]\)。当
\(\gamma < 1\) 时，该迭代是 \(\infty\)-范数下的压缩映射（压缩因子
\(\gamma\)），线性收敛到唯一不动点。

\textbf{算法 138.2}（策略迭代 / Policy
Iteration）：交替执行策略评估和策略改进，对有限状态和有限控制问题，在至多有限步内收敛到最优策略（每步严格提升）。

\textbf{定义 138.3}（Markov 决策过程 /
MDP）：动态规划在随机环境中的推广。转移概率 \(p(x' \mid x, u)\)
替代确定性转移 \(f(x, u)\)，Bellman 方程为

\[V(x) = \min_{u \in U(x)} \left[ c(x, u) + \gamma \sum_{x'} p(x' \mid x, u) \, V(x') \right]\]

\textbf{定理 138.1}（Blackwell 最优性理论，1965）：在标准的折扣
MDP（有限状态、有限控制）中，存在确定性的平稳最优策略，且该策略可由值迭代或策略迭代找到。最优策略的存在性和结构由
Blackwell 的单调性和压缩性条件保证。

\subsubsection{138.2 随机规划（Stochastic
Programming）}\label{ux968fux673aux89c4ux5212stochastic-programming}

随机规划处理包含随机参数（通过随机变量 \(\xi\) 表达）的优化问题。

\textbf{定义 138.4}（二阶段线性随机规划 / Dantzig, 1955）：第一阶段在
\(\xi\) 实现前确定''现在''决策 \(x\)，第二阶段在 \(\xi\)
实现后采取折中（recourse）行动 \(y\)：

\[\begin{aligned} \min_x \quad & c^T x + \mathbb{E}_\xi[Q(x, \xi)] \\ \text{s.t.} \quad & Ax = b, \quad x \geq 0 \end{aligned}\]

其中
\(Q(x,\xi) = \min_y \{q(\xi)^T y : W(\xi) y = h(\xi) - T(\xi) x, \; y \geq 0\}\)
为第二阶段（recourse）问题的值函数。

\textbf{定理 138.2}（二阶段问题的结构性质）：若对几乎所有
\(\xi\)，第二阶段问题是可行的（完全 recourse），则 \(Q(x, \xi)\) 是关于
\(x\) 的分片线性凸函数，\(\mathbb{E}_\xi[Q(x, \xi)]\) 也是凸函数。当
\(\xi\)
具有有限支撑（场景树）时，二阶段问题等价于一个大规模确定性线性规划。

\textbf{定义 138.5}（机会约束规划 / Chance-Constrained Programming,
Charnes-Cooper, 1959）：用概率约束替代确定性约束：

\[P\big(g(x, \xi) \leq 0\big) \geq 1 - \alpha, \quad \alpha \in (0,1)\]

当 \(g\) 关于 \(x\) 线性、\(\xi\)
具有已知联合正态分布时，概率约束等价于一个二阶锥约束。

\subsubsection{138.3
多目标最优化}\label{ux591aux76eeux6807ux6700ux4f18ux5316}

\textbf{定义 138.6}（Pareto 最优解）：对多目标问题
\(\min (f_1(x), \ldots, f_k(x))\)（\(x \in \mathcal{X}\)），解 \(x^*\)
称为 \textbf{Pareto 最优}，若不存在 \(x' \in \mathcal{X}\) 使得对所有
\(i\) 有 \(f_i(x') \leq f_i(x^*)\)，且至少一个不等式严格成立。所有
Pareto 最优解的集合称为 \textbf{Pareto 前沿}。

\textbf{求解方法}： -
\textbf{加权和法}（标量化）：\(\min \sum_{i=1}^k w_i f_i(x)\)（\(w_i \geq 0\)），变化权重可生成
Pareto 前沿的凸部分 - \textbf{\(\varepsilon\)-约束法}：\(\min f_1(x)\)
s.t. \(f_i(x) \leq \varepsilon_i\)（\(i = 2, \ldots, k\)），系统改变
\(\varepsilon_i\) 可生成完整的 Pareto 前沿（包括非凸部分） -
\textbf{目标规划}：最小化与目标的加权偏差 \(\min \sum_i |f_i(x) - t_i|\)

Pareto 前沿的结构由 Kuhn-Tucker 型必要条件刻画：若 \(x^*\) Pareto
最优且各目标在约束下正则，则存在非负乘子
\(\lambda_i \geq 0\)（不全为零）使
\(\sum \lambda_i \nabla f_i(x^*) = 0\)。

\bigskip\hrule\bigskip

\emph{本卷基于 V2 Ch 11（线性规划初步）和
V7（线性代数），建立了线性规划、对偶理论、凸优化、非线性规划和变分法的系统理论。共
6 章。}

\newpage

\section{卷二十七：动力系统}\label{ux5377ux4e8cux5341ux4e03ux52a8ux529bux7cfbux7edf}

\begin{quote}
动力系统研究随时间演化的系统的长期行为，起源于 Newton 力学和 Poincaré
对天体力学的研究。本卷在常微分方程（V18）的基础上，建立离散和连续动力系统的基本理论，包括稳定性、混沌、遍历论和分支理论。动力系统为理解自然界中的复杂动力学现象提供了统一的数学框架。
\end{quote}

\begin{quote}
\textbf{前置知识}：本卷依赖V17（微分方程）和V10（拓扑学）的核心内容，建议在修读本卷前先掌握这两卷的基础。
\end{quote}

\bigskip\hrule\bigskip

\subsection{第137章：离散动力系统}\label{ux7b2c137ux7ae0ux79bbux6563ux52a8ux529bux7cfbux7edf}

离散动力系统由映射的迭代生成，是最简单的动力系统，但可以展现极其丰富的动力学行为。

\subsubsection{137.1 迭代与轨道}\label{ux8fedux4ee3ux4e0eux8f68ux9053}

\textbf{定义 137.1}（离散动力系统）：一个\textbf{离散动力系统}由状态空间
\(X\)（通常为 \(\mathbb{R}^n\) 的子集或流形）和映射 \(f: X \to X\)
组成。从初始状态 \(x_0 \in X\) 出发，系统在时刻 \(n\) 的状态为

\[x_n = f^n(x_0) = \underbrace{f \circ f \circ \cdots \circ f}_{n \text{ 次}}(x_0)\]

序列 \(\{x_n\}_{n=0}^{\infty}\) 称为 \(x_0\)
的\textbf{轨道}（orbit）或\textbf{正向轨道}。

\textbf{定义 137.2}（不动点与周期点）： - \(x\) 是\textbf{不动点}，如果
\(f(x) = x\) - \(x\) 是\textbf{周期点}（周期为 \(p\)），如果
\(f^p(x) = x\) 且 \(f^k(x) \neq x\)（\(0 < k < p\)） - 周期 \(p\) 的轨道
\(\{x, f(x), \ldots, f^{p-1}(x)\}\)
称为\textbf{周期轨道}或\textbf{\(p\)-周期轨道}

\textbf{定义 137.3}（稳定性的分类）：设 \(x^*\) 是不动点： - \(x^*\) 是
\textbf{Lyapunov 稳定}的，如果对任意 \(\varepsilon > 0\)，存在
\(\delta > 0\) 使得
\(|x_0 - x^*| < \delta \implies |f^n(x_0) - x^*| < \varepsilon\) 对所有
\(n \geq 0\) - \(x^*\) 是\textbf{渐近稳定}的，如果它是 Lyapunov
稳定的，且存在 \(\eta > 0\) 使得
\(|x_0 - x^*| < \eta \implies \lim_{n \to \infty} f^n(x_0) = x^*\) -
\(x^*\) 是\textbf{不稳定}的，如果它不是 Lyapunov 稳定的

\textbf{定理 137.1}（一维不动点的线性稳定性）：设
\(f: \mathbb{R} \to \mathbb{R}\) 是 \(C^1\) 映射，\(x^*\) 是不动点。 -
若 \(|f'(x^*)| < 1\)，则 \(x^*\) 是渐近稳定的（\textbf{吸引不动点}） -
若 \(|f'(x^*)| > 1\)，则 \(x^*\) 是不稳定的（\textbf{排斥不动点}） - 若
\(|f'(x^*)| = 1\)，线性稳定性分析不足，需考虑高阶项

\emph{证明}：由中值定理，\(|x_{n+1} - x^*| = |f(x_n) - f(x^*)| = |f'(\xi_n)| |x_n - x^*|\)。当
\(|f'(x^*)| < 1\) 时，在 \(x^*\) 附近 \(|f'| < \lambda < 1\)，故
\(|x_n - x^*| \leq \lambda^n |x_0 - x^*| \to 0\)。∎

\textbf{定理 137.2}（高维不动点的线性稳定性）：设
\(f: \mathbb{R}^n \to \mathbb{R}^n\) 是 \(C^1\) 映射，\(x^*\)
是不动点，\(Df(x^*)\) 是 Jacobi 矩阵。 - 若 \(Df(x^*)\) 的所有特征值的模
\(< 1\)，则 \(x^*\) 是渐近稳定的（\textbf{汇}，sink） - 若存在特征值模
\(> 1\)，则 \(x^*\) 是不稳定的（\textbf{源}，source
或\textbf{鞍点}，saddle） - 若所有特征值模 \(\leq 1\) 且至少一个等于
1，线性分析不足

\subsubsection{137.2
一维映射的动力学}\label{ux4e00ux7ef4ux6620ux5c04ux7684ux52a8ux529bux5b66}

\textbf{例 137.1}（Logistic
映射）：\(f(x) = rx(1-x)\)，\(x \in [0, 1]\)，\(r \in [0, 4]\)。这是最经典的一维非线性离散动力系统。

\begin{itemize}
\tightlist
\item
  \(0 < r \leq 1\)：仅有不动点 \(x = 0\)（稳定），种群灭绝
\item
  \(1 < r < 3\)：不动点 \(x^* = 1 - 1/r\) 稳定，\(x = 0\) 不稳定
\item
  \(3 \leq r < 1 + \sqrt{6} \approx 3.449\)：2-周期轨道稳定（倍周期分岔）
\item
  \(r\) 继续增大：出现 4, 8, 16, \ldots{} 周期轨道（倍周期分岔级联）
\item
  \(r \approx 3.5699\)（Feigenbaum 点）：进入混沌区域
\item
  \(r = 4\)：完全混沌，在 \([0, 1]\) 上有稠密轨道
\end{itemize}

\textbf{定理 137.3}（Li-Yorke 定理，1975）：「周期 3
蕴含混沌」。若连续映射 \(f: I \to I\)（\(I\) 为区间）有 3-周期轨道，则
\(f\) 有任意周期的周期轨道，且存在不可数集合
\(S \subseteq I\)（混沌集）满足 Li-Yorke 混沌条件。

\emph{证明概要}：设 3-周期轨道为 \(a<b<c\)，由介值定理可证
\(f([a,b])\supseteq[b,c]\) 和 \(f([b,c])\supseteq[a,c]\)。利用区间映射的
Markov 划分，构造符号动力学的子移位，对任意正整数
\(n\)，可从符号序列中提取 \(n\)-周期点。混沌集 \(S\)
的构造基于：不含特定重复模式的轨道集合是不可数的，且其中任意两点满足混沌条件。∎

\textbf{定理 137.4}（Sharkovskii 定理，1964）：将正整数按以下
\textbf{Sharkovskii 序}排列：

\[3 \triangleright 5 \triangleright 7 \triangleright \cdots \triangleright 2\cdot 3 \triangleright 2\cdot 5 \triangleright \cdots \triangleright 2^2\cdot 3 \triangleright 2^2\cdot 5 \triangleright \cdots \triangleright 2^3 \triangleright 2^2 \triangleright 2 \triangleright 1\]

即：先排所有大于1的奇数（\(3,5,7,\dots\)），再排所有 \(2 \times\)
奇数（\(6,10,14,\dots\)），再排 \(2^2 \times\)
奇数，以此类推，最后将2的所有幂降序排列（\(\dots, 8, 4, 2, 1\)）。此序是
\(\mathbb{N}\) 上的全序。定理断言：若连续映射
\(f: \mathbb{R} \to \mathbb{R}\) 有 \(m\)-周期轨道，且
\(m \triangleright n\)（\(m\) 在 \(n\) 之前），则 \(f\) 必有
\(n\)-周期轨道。特别地，``周期3蕴含所有周期''------若存在3-周期轨道（3排在最前），则
\(f\) 有任意周期的周期轨道。这是 Li-Yorke 定理（1975）的前身。∎

\subsubsection{137.3
符号动力学与混沌}\label{ux7b26ux53f7ux52a8ux529bux5b66ux4e0eux6df7ux6c8c}

\textbf{定义 137.4}（符号空间）：令 \(\Sigma_2 = \{0, 1\}^{\mathbb{N}}\)
为由两个符号构成的所有无穷序列的空间。度量定义为
\(d(\mathbf{s}, \mathbf{t}) = \sum_{n=0}^{\infty} \frac{|s_n - t_n|}{2^n}\)。\textbf{移位映射}
\(\sigma: \Sigma_2 \to \Sigma_2\) 定义为
\(\sigma(s_0 s_1 s_2 \cdots) = (s_1 s_2 s_3 \cdots)\)。

\textbf{命题 137.5}（移位映射的混沌性质）： 1. \(\sigma\) 在
\(\Sigma_2\) 上有稠密轨道 2. \(\sigma\) 的周期点在 \(\Sigma_2\) 中稠密
3. \(\sigma\) 对初始条件敏感依赖（Devany 混沌的三个条件）

\textbf{定理 137.6}（Logistic 映射与符号动力学的联系）：\(r = 4\)
时，Logistic 映射 \(f(x) = 4x(1-x)\)
拓扑共轭于移位映射。具体地，半共轭由 \(h(\theta) = \sin^2(\pi\theta)\)
给出（通过帐篷映射）。

\subsubsection{137.4
拓扑共轭与结构稳定性}\label{ux62d3ux6251ux5171ux8f6dux4e0eux7ed3ux6784ux7a33ux5b9aux6027}

\textbf{定义 137.5}（拓扑共轭）：两个映射 \(f: X \to X\) 和
\(g: Y \to Y\) 是\textbf{拓扑共轭}的，如果存在同胚 \(h: X \to Y\) 使得
\(h \circ f = g \circ h\)。拓扑共轭的映射具有相同的动力学性质（周期轨道结构、稠密性、混沌性等）。

\textbf{定义 137.6}（结构稳定性）：映射 \(f\)
是\textbf{结构稳定}的，如果存在 \(f\) 的邻域（在 \(C^r\)
拓扑中），使得此邻域内所有映射拓扑共轭于 \(f\)。

\bigskip\hrule\bigskip

\subsection{第138章：连续动力系统}\label{ux7b2c138ux7ae0ux8fdeux7eedux52a8ux529bux7cfbux7edf}

连续动力系统由微分方程定义，其解形成流（flow），是经典力学和许多物理系统的数学模型。

\subsubsection{138.1 流与向量场}\label{ux6d41ux4e0eux5411ux91cfux573a}

\textbf{定义 138.1}（流）：设 \(M\) 是光滑流形。向量场 \(X\)
生成一个\textbf{流} \(\phi_t: M \to M\)（\(t \in \mathbb{R}\)），满足
\(\frac{d}{dt}\phi_t(x) = X(\phi_t(x))\)，\(\phi_0(x) = x\)。流满足群性质：\(\phi_{t+s} = \phi_t \circ \phi_s\)。

\textbf{定义 138.2}（轨道与不变集）： - 过点 \(x\)
的轨道：\(\gamma(x) = \{\phi_t(x) : t \in \mathbb{R}\}\) -
正向轨道：\(\gamma^+(x) = \{\phi_t(x) : t \geq 0\}\) -
负向轨道：\(\gamma^-(x) = \{\phi_t(x) : t \leq 0\}\) - 集合 \(A\)
是\textbf{不变集}，如果 \(\phi_t(A) = A\) 对所有 \(t \in \mathbb{R}\) -
\(A\) 是\textbf{正向不变集}，如果 \(\phi_t(A) \subseteq A\) 对所有
\(t \geq 0\)

\textbf{定义 138.3}（\(\omega\)-极限集与 \(\alpha\)-极限集）：对点
\(x\)， -
\(\omega(x) = \{y : \exists t_n \to \infty, \phi_{t_n}(x) \to y\}\)（\(\omega\)-极限集）
-
\(\alpha(x) = \{y : \exists t_n \to -\infty, \phi_{t_n}(x) \to y\}\)（\(\alpha\)-极限集，仅当负向轨道存在）

\subsubsection{138.2
平衡点与线性化}\label{ux5e73ux8861ux70b9ux4e0eux7ebfux6027ux5316}

\textbf{定义 138.4}（平衡点）：点 \(x^*\)
是\textbf{平衡点}（不动点/奇点），如果 \(X(x^*) = 0\)，即
\(\phi_t(x^*) = x^*\) 对所有 \(t\)。

\textbf{定理 138.1}（Hartman-Grobman 定理）：设 \(x^*\) 是 \(C^1\)
向量场 \(X\) 的双曲平衡点（即 \(DX(x^*)\) 的所有特征值实部非零）。则存在
\(x^*\) 的邻域 \(U\)，使得 \(X\) 在 \(U\) 上的流拓扑共轭于其线性化
\(\dot{\mathbf{x}} = DX(x^*) \mathbf{x}\) 的流。

\emph{证明概要}：通过对线性部分和非线性扰动构造适当的同胚。关键在于双曲性保证存在不变的稳定/不稳定子空间的一致压缩/扩张性。利用Banach不动点定理，将共轭同胚
\(h = \text{id} + u\) 的构造转化为函数方程
\((L + R)(h(x)) = h(DX \cdot x)\) 的求解，其中 \(L\) 为线性流，\(R\)
为高阶项。双曲条件确保该方程有唯一连续解。∎

\textbf{定义 138.5}（稳定流形定理）：设 \(x^*\) 是双曲平衡点。则存在： -
\textbf{稳定流形}
\(W^s(x^*) = \{x : \lim_{t \to \infty} \phi_t(x) = x^*\}\)（维数
\(= n_-\)，即负实部特征值个数） - \textbf{不稳定流形}
\(W^u(x^*) = \{x : \lim_{t \to -\infty} \phi_t(x) = x^*\}\)（维数
\(= n_+\)，即正实部特征值个数）

\(W^s\) 和 \(W^u\) 是浸入子流形，在 \(x^*\) 处分别切于稳定子空间 \(E^s\)
和不稳定子空间 \(E^u\)。

\textbf{定理 138.2}（Lyapunov 函数与稳定性）：设 \(x^*\)
是平衡点，\(V: U \to \mathbb{R}\) 是 \(C^1\) 函数，\(U\) 为含 \(x^*\)
的开集，满足： - \(V(x^*) = 0\)，\(V(x) > 0\) 对 \(x \neq x^*\)（正定）
- \(\dot{V}(x) = \nabla V(x) \cdot X(x) \leq 0\) 在 \(U\) 上

则 \(x^*\) 是 Lyapunov 稳定的。若进一步
\(\dot{V}(x) < 0\)（\(x \neq x^*\)），则 \(x^*\) 是渐近稳定的。

\subsubsection{138.3
平面动力系统}\label{ux5e73ux9762ux52a8ux529bux7cfbux7edf}

\textbf{定理 138.3}（Poincaré-Bendixson 定理）：设 \(\phi_t\) 是平面
\(\mathbb{R}^2\) 上的 \(C^1\) 流，\(\gamma^+(x)\)
是包含在紧致集内的正向轨道。若 \(\omega(x)\) 不含平衡点，则
\(\omega(x)\) 是周期轨道。

\emph{证明概要}：取 \(y \in \omega(x)\)，通过 \(y\) 作横截线
\(\Sigma\)（局部截面）。由
\(\omega\)-极限集的不变性和不含平衡点，\(\omega(x)\) 与 \(\Sigma\)
的交为离散点列。利用 Jordan 曲线定理和流的单调性，可证明 \(\gamma^+(x)\)
与 \(\Sigma\) 的交点单调排列，极限点必为同一点，故 \(\omega(x)\)
为闭轨（周期轨道）。□ \textbf{\(\omega\)-极限集分类}：平面流的
\(\omega\)-极限集只能是以下三类之一：(1) 平衡点；(2) 周期轨道；(3)
由平衡点和连接它们的异宿/同宿轨组成的图。三维及以上可出现奇怪吸引子。∎

\textbf{推论
138.4}：紧致平面区域内的正向不变集若不含平衡点，则必含周期轨道。

\textbf{例 138.1}（van der Pol
振子）：\(\ddot{x} + \mu(x^2 - 1)\dot{x} + x = 0\)。转化为系统：\(\dot{x} = y\)，\(\dot{y} = -x - \mu(x^2 - 1)y\)。对
\(\mu > 0\)，存在唯一的稳定极限环（由 Poincaré-Bendixson 定理保证）。

\textbf{例 138.2}（Lotka-Volterra
捕食者-猎物模型）：\(\dot{x} = x(a - by)\)，\(\dot{y} = y(-c + dx)\)（\(a,b,c,d > 0\)）。平衡点
\((c/d, a/b)\) 是中心，所有非平凡轨道是周期轨道。系统有守恒量
\(H(x,y) = d x - c \ln x + b y - a \ln y\)。

\subsubsection{138.4 极限环与 Hopf
分支}\label{ux6781ux9650ux73afux4e0e-hopf-ux5206ux652f}

\textbf{定义
138.6}（极限环）：极限环是孤立的周期轨道（即存在邻域内无其他周期轨道）。

\textbf{定理 138.5}（Hopf 分支定理）：考虑依赖参数 \(\mu\) 的 ODE 系统
\(\dot{\mathbf{x}} = \mathbf{f}(\mathbf{x}, \mu)\)。设平衡点
\(\mathbf{x}^*(\mu)\) 在 \(\mu = \mu_0\) 时有一对纯虚特征值
\(\lambda(\mu) = \alpha(\mu) \pm i\omega(\mu)\)，\(\alpha(\mu_0) = 0\)，\(\omega(\mu_0) \neq 0\)，且
\(\alpha'(\mu_0) \neq 0\)（横截条件）。则在 \(\mu = \mu_0\)
处分支出唯一的极限环族。

\bigskip\hrule\bigskip

\subsection{第139章：混沌与分形}\label{ux7b2c139ux7ae0ux6df7ux6c8cux4e0eux5206ux5f62}

混沌是确定性动力系统表现出的类随机行为，由对初始条件的敏感依赖所刻画。分形是具有自相似性和非整数维度的几何对象，常出现于混沌系统的吸引子。

\subsubsection{139.1 混沌的定义}\label{ux6df7ux6c8cux7684ux5b9aux4e49}

\textbf{定义 139.1}（Devaney 混沌，1989）：映射 \(f: X \to X\)（\(X\)
为度量空间）是\textbf{混沌}的，如果： 1. \(f\) 对初始条件敏感依赖：存在
\(\delta > 0\) 使得对任意 \(x\) 和 \(\varepsilon > 0\)，存在 \(y\) 满足
\(d(x,y) < \varepsilon\) 且对某个
\(n \geq 0\)，\(d(f^n(x), f^n(y)) \geq \delta\) 2. \(f\)
是拓扑传递的：对任意非空开集 \(U, V\)，存在 \(n \geq 0\) 使得
\(f^n(U) \cap V \neq \emptyset\) 3. \(f\) 的周期点在 \(X\) 中稠密

\textbf{定理 139.1}（Banks et al.,
1992）：在紧致度量空间上，拓扑传递性和稠密周期点蕴含敏感依赖性。因此条件
(1) 是冗余的。

\textbf{定义 139.2}（Lyapunov 指数）：一维映射 \(f\) 在轨道 \(\{x_n\}\)
上的 Lyapunov 指数为

\[\lambda(x_0) = \lim_{n \to \infty} \frac{1}{n} \sum_{k=0}^{n-1} \ln |f'(x_k)|\]

（若极限存在）。\(\lambda > 0\) 是混沌的数值标志。高维系统有 \(n\) 个
Lyapunov 指数，正的 Lyapunov 指数数量刻画了混沌的程度。

\textbf{例 139.1}（Logistic 映射 \(r=4\) 的 Lyapunov
指数）：\(f(x) = 4x(1-x)\)，\(\lambda = \ln 2 > 0\)。

\subsubsection{139.2 奇怪吸引子}\label{ux5947ux602aux5438ux5f15ux5b50}

\textbf{定义 139.3}（吸引子）：紧致不变集 \(A\)
是\textbf{吸引子}，如果存在开集 \(U \supset A\)（\(A\) 的吸引盆）使得
\(\phi_t(U) \subseteq U\) 对所有 \(t \geq 0\)，且
\(A = \bigcap_{t \geq 0} \phi_t(U)\)。

\textbf{定义 139.4}（奇怪吸引子）：如果吸引子 \(A\)
上的动力系统是混沌的（对初始条件敏感依赖），则 \(A\)
称为\textbf{奇怪吸引子}。奇怪吸引子通常具有分形结构。

\textbf{例 139.2}（Lorenz 吸引子，1963）：Lorenz 方程：

\[\begin{aligned} \dot{x} &= \sigma(y - x) \\ \dot{y} &= x(\rho - z) - y \\ \dot{z} &= xy - \beta z \end{aligned}\]

经典参数 \(\sigma = 10, \beta = 8/3, \rho = 28\)。Lorenz
吸引子是蝴蝶形状的奇怪吸引子，Lyapunov 指数为
\(\lambda_1 \approx 0.9 > 0\)。

\textbf{例 139.3}（Hénon
映射，1976）：\((x_{n+1}, y_{n+1}) = (1 - a x_n^2 + y_n, b x_n)\)。经典参数
\(a = 1.4, b = 0.3\)。Hénon 吸引子是经典的奇怪吸引子。

\subsubsection{139.3 分形}\label{ux5206ux5f62}

\textbf{定义 139.5}（Hausdorff 维数 / 分形维数）：对度量空间 \((X, d)\)
的子集 \(S\)，\(s\) 维 Hausdorff 测度定义为

\[\mathcal{H}^s(S) = \lim_{\delta \to 0} \inf \left\{ \sum_{i=1}^{\infty} (\operatorname{diam} U_i)^s : S \subseteq \bigcup_i U_i, \operatorname{diam} U_i < \delta \right\}\]

Hausdorff 维数为
\(\dim_H(S) = \inf\{s : \mathcal{H}^s(S) = 0\} = \sup\{s : \mathcal{H}^s(S) = \infty\}\)。

\textbf{定义 139.6}（盒维数 / Box-counting dimension）：设
\(N(\varepsilon)\) 是覆盖 \(S\) 所需直径为 \(\varepsilon\)
的集合的最小个数。盒维数为

\[\dim_B(S) = \lim_{\varepsilon \to 0} \frac{\ln N(\varepsilon)}{\ln(1/\varepsilon)}\]

（若极限存在）。一般有 \(\dim_H(S) \leq \dim_B(S)\)。

\textbf{例 139.4}（经典分形）： - \textbf{Cantor 三分集}：从 \([0,1]\)
反复删除中间三分之一。\(\dim_H = \ln 2 / \ln 3 \approx 0.631\) -
\textbf{Koch 雪花曲线}：\(\dim_H = \ln 4 / \ln 3 \approx 1.262\) -
\textbf{Sierpiński 三角形}：\(\dim_H = \ln 3 / \ln 2 \approx 1.585\) -
\textbf{Menger 海绵}：\(\dim_H = \ln 20 / \ln 3 \approx 2.727\)

\textbf{定理 139.2}（Kaplan-Yorke 猜想 / Lyapunov 维数）：吸引子 \(A\)
的 Lyapunov 维数定义为

\[\dim_L(A) = k + \frac{\sum_{i=1}^k \lambda_i}{|\lambda_{k+1}|}\]

其中 \(\lambda_1 \geq \lambda_2 \geq \cdots \geq \lambda_n\) 为 Lyapunov
指数，\(k\) 是满足 \(\sum_{i=1}^k \lambda_i \geq 0\) 的最大整数。Lorenz
吸引子的 Lyapunov 维数约为 \(2.06\)。

\subsubsection{139.4
通向混沌的道路}\label{ux901aux5411ux6df7ux6c8cux7684ux9053ux8def}

\textbf{定义 139.7}（倍周期分岔通向混沌）：Logistic 映射中，随着 \(r\)
增大，出现无限级联的倍周期分岔，其分岔值满足 Feigenbaum 普适性：

\[\lim_{n \to \infty} \frac{r_n - r_{n-1}}{r_{n+1} - r_n} = \delta \approx 4.6692 \quad \text{(Feigenbaum 常数)}\]

\textbf{定义 139.8}（间歇性通向混沌 /
Pomeau-Manneville）：系统在长时间规则行为与短时间不规则爆发之间交替。

\textbf{定义 139.9}（Ruelle-Takens-Newhouse 通向混沌）：经过三次 Hopf
分岔（环面→环面→环面）后，第三个环面不稳定，直接产生奇怪吸引子。

\bigskip\hrule\bigskip

\subsection{第140章：遍历论基础}\label{ux7b2c140ux7ae0ux904dux5386ux8bbaux57faux7840}

遍历论研究动力系统在长时间平均下的统计行为，是统计力学和概率论在动力系统中的应用。

\subsubsection{140.1
保测变换与不变测度}\label{ux4fddux6d4bux53d8ux6362ux4e0eux4e0dux53d8ux6d4bux5ea6}

\textbf{定义 140.1}（保测变换）：设 \((X, \mathcal{B}, \mu)\)
是概率空间。可测变换 \(T: X \to X\) 是\textbf{保测变换}，如果
\(\mu(T^{-1}(A)) = \mu(A)\) 对所有 \(A \in \mathcal{B}\)。\(\mu\) 称为
\(T\)-不变测度。

\textbf{例 140.1}（保测变换的例子）： - 圆周旋转
\(T(x) = x + \alpha \pmod{1}\) 保持 Lebesgue 测度 - 加倍映射
\(T(x) = 2x \pmod{1}\)（\(x \in [0,1)\)）保持 Lebesgue 测度 - 帐篷映射
\(T(x) = 2\min(x, 1-x)\) 保持 Lebesgue 测度 - Logistic 映射
\(r=4\)：\(T(x) = 4x(1-x)\) 保持测度
\(\mu(A) = \frac{1}{\pi} \int_A \frac{dx}{\sqrt{x(1-x)}}\)

\textbf{定理 140.1}（Krylov-Bogolyubov
定理）：紧致度量空间上的任何连续映射至少有一个不变 Borel 概率测度。

\emph{证明思路}：取任意概率测度 \(\nu\)，考虑 Cesàro 平均
\(\mu_n = \frac{1}{n}\sum_{k=0}^{n-1} T^k_* \nu\)。由 Banach-Alaoglu
定理和 Riesz 表示定理，\(\mu_n\) 有弱*收敛子列，极限测度是 \(T\)
不变的。∎

\subsubsection{140.2 遍历性与 Birkhoff
遍历定理}\label{ux904dux5386ux6027ux4e0e-birkhoff-ux904dux5386ux5b9aux7406}

\textbf{定义 140.2}（遍历性）：保测变换 \(T\)
是\textbf{遍历的}（ergodic），如果对任意满足 \(T^{-1}(A) = A\) 的可测集
\(A\)，有 \(\mu(A) = 0\) 或
\(\mu(A) = 1\)。即系统不可分解为两个非平凡的不变部分。

\textbf{定理 140.2}（Birkhoff 遍历定理，1931）：设 \(T\)
是保测变换，\(f \in L^1(\mu)\)。则对几乎所有的 \(x\)，时间平均收敛：

\[\lim_{n \to \infty} \frac{1}{n} \sum_{k=0}^{n-1} f(T^k x) = \bar{f}(x)\]

其中 \(\bar{f} \in L^1(\mu)\) 是 \(T\)
不变的（\(\bar{f} \circ T = \bar{f}\)），且
\(\int_X \bar{f} \, d\mu = \int_X f \, d\mu\)。若 \(T\) 是遍历的，则
\(\bar{f}(x) = \int_X f \, d\mu\) 几乎处处成立。

\emph{证明概要}：核心工具是极大遍历不等式（Maximal Ergodic
Theorem）。定义部分和
\(S_n f(x) = \sum_{k=0}^{n-1} f(T^k x)\)，\(M_N f = \max\{0, S_1, \dots, S_N\}\)。极大不等式断言
\(\int_{\{M_N f > 0\}} f \, d\mu \geq 0\)。由此考虑上、下极限函数
\(\bar{f}(x) = \limsup S_n f/n\) 和
\(\underline{f}(x) = \liminf S_n f/n\)，利用 \(T\)
不变性和极大不等式证明 \(\bar{f} = \underline{f} = \bar{f}\)（即收敛）且
\(\int \bar{f} = \int f\)。遍历性下 \(\bar{f}\) 几乎处处为常数，故
\(\bar{f} = \int f\) 几乎处处。∎

\textbf{推论 140.3}（Boltzmann
遍历假设的数学严格化）：对遍历系统，长时间平均等于空间平均。

\textbf{定理 140.4}（遍历性的等价刻画）：以下等价： 1. \(T\) 是遍历的 2.
任何 \(T\) 不变的可测函数几乎处处为常数 3. 对任意
\(A, B \in \mathcal{B}\)，\(\lim_{n \to \infty} \frac{1}{n} \sum_{k=0}^{n-1} \mu(T^{-k} A \cap B) = \mu(A) \mu(B)\)

\textbf{定义 140.3}（混合）：\(T\)
是\textbf{混合}的（mixing），如果对任意
\(A, B \in \mathcal{B}\)，\(\lim_{n \to \infty} \mu(T^{-n} A \cap B) = \mu(A) \mu(B)\)。混合蕴含遍历性，但反之不成立。

\textbf{定义 140.4}（弱混合）：\(T\) 是\textbf{弱混合}的（weakly
mixing），如果对任意
\(A, B \in \mathcal{B}\)，\(\lim_{n \to \infty} \frac{1}{n} \sum_{k=0}^{n-1} |\mu(T^{-k} A \cap B) - \mu(A)\mu(B)| = 0\)。强混合蕴含弱混合，弱混合蕴含遍历性，反之均不成立。

\textbf{定理 140.5}（弱混合的谱刻画）：\(T\) 是\textbf{弱混合}的当且仅当
\(T\) 在 \(L^2(\mu) \ominus \{\text{常数函数}\}\) 上的 Koopman 算子
\(U_T f = f \circ T\)
具有连续谱（即没有特征值）。这是弱混合的一个重要谱刻画，注意强混合蕴含弱混合，反之一般不成立。

\subsubsection{140.3 熵}\label{ux71b5}

\textbf{定义 140.5}（测度论熵 / Kolmogorov-Sinai 熵）：设
\(\mathcal{P} = \{P_1, \ldots, P_k\}\) 是 \(X\)
的有限可测划分。\(\mathcal{P}\) 的熵为
\(H(\mathcal{P}) = -\sum_{i=1}^k \mu(P_i) \log \mu(P_i)\)。\(T\) 关于
\(\mathcal{P}\) 的熵率为

\[h(T, \mathcal{P}) = \lim_{n \to \infty} \frac{1}{n} H\left( \bigvee_{k=0}^{n-1} T^{-k} \mathcal{P} \right)\]

\(T\) 的测度论熵为 \(h_\mu(T) = \sup_{\mathcal{P}} h(T, \mathcal{P})\)。

\textbf{定理 140.6}（熵与遍历性）：\(h_\mu(T) > 0\) 蕴含 \(T\)
是遍历的（实际上是 Kolmogorov 自同构，比混合更强）。

\textbf{例 140.2}：Bernoulli 移位 \((p_1, \ldots, p_k)\) 的熵为
\(-\sum p_i \log p_i\)。加倍映射 \(x \mapsto 2x \pmod{1}\) 关于 Lebesgue
测度的熵为 \(\log 2\)。

\textbf{定理 140.7}（Pesin 熵公式）：对光滑动力系统，Kolmogorov-Sinai
熵等于正的 Lyapunov 指数之和：

\[h_\mu(T) = \sum_{\lambda_i > 0} \lambda_i\]

（在适当的假设下，如 SRB 测度）。

\subsubsection{140.4
拓扑动力系统}\label{ux62d3ux6251ux52a8ux529bux7cfbux7edf}

\textbf{定义 140.6}（极小性）：动力系统 \((X, T)\)
是\textbf{极小}的，如果 \(X\)
不含非平凡的闭不变子集。等价地，每一条轨道都是稠密的。

\textbf{定理 140.8}（Birkhoff
回复定理）：紧致度量空间上的连续映射必有回复点（即存在 \(x\) 和
\(n_k \to \infty\) 使得 \(T^{n_k}(x) \to x\)）。

\textbf{定义 140.7}（拓扑熵）：拓扑熵 \(h_{top}(T)\)
度量了动力系统的轨道复杂性。对紧致度量空间上的连续映射 \(T\)，

\[h_{top}(T) = \lim_{\varepsilon \to 0} \limsup_{n \to \infty} \frac{1}{n} \log N(n, \varepsilon)\]

其中 \(N(n, \varepsilon)\) 是 \(n\) 步轨道在 \(\varepsilon\)
分辨率下可区分的最大轨道段数。

\textbf{定理
140.9}（变分原理）：\(h_{top}(T) = \sup_{\mu \in \mathcal{M}_T} h_\mu(T)\)，其中
\(\mathcal{M}_T\) 是 \(T\) 的所有不变概率测度的集合。

\bigskip\hrule\bigskip

\subsection{第141章：分支理论}\label{ux7b2c141ux7ae0ux5206ux652fux7406ux8bba}

分支理论研究当参数变化时，动力系统的定性行为如何变化。分支点（bifurcation
point）是参数值，在此处系统的拓扑结构发生改变。

\subsubsection{141.1
不动点的分支}\label{ux4e0dux52a8ux70b9ux7684ux5206ux652f}

\textbf{定义 141.1}（局部分支）：考虑依赖参数 \(\mu \in \mathbb{R}\) 的
ODE \(\dot{x} = f(x, \mu)\)。设 \(f(0, 0) = 0\)。\textbf{分支点}是
\(\mu\) 的值，在此处不动点 \(x = 0\) 的稳定性改变或新不动点产生。

\textbf{定理 141.1}（分支的必要条件）：在分支点 \(\mu = 0\)
处，\(\frac{\partial f}{\partial x}(0, 0) = 0\)（即线性化矩阵有零特征值）。

\textbf{定义 141.2}（鞍结分支 / Saddle-node bifurcation）：标准形
\(\dot{x} = \mu - x^2\)。当 \(\mu < 0\) 时无不动点；\(\mu = 0\)
时有一个半稳定不动点；\(\mu > 0\)
时有两个不动点（一个稳定，一个不稳定）。

\textbf{定义 141.3}（跨临界分支 / Transcritical bifurcation）：标准形
\(\dot{x} = \mu x - x^2\)。不动点 \(x = 0\) 和 \(x = \mu\)
始终存在，但在 \(\mu = 0\) 处交换稳定性。

\textbf{定义 141.4}（叉式分支 / Pitchfork bifurcation）： -
\textbf{超临界}（supercritical）：\(\dot{x} = \mu x - x^3\)。\(\mu \leq 0\)
时 \(x=0\) 稳定；\(\mu > 0\) 时 \(x=0\) 不稳定，出现两个稳定的对称不动点
\(\pm \sqrt{\mu}\) -
\textbf{亚临界}（subcritical）：\(\dot{x} = \mu x + x^3\)。\(\mu < 0\)
时 \(x=0\) 稳定，有两个不稳定不动点 \(\pm \sqrt{-\mu}\)；\(\mu \geq 0\)
时 \(x=0\) 不稳定

\subsubsection{141.2
周期轨道的分支}\label{ux5468ux671fux8f68ux9053ux7684ux5206ux652f}

\textbf{定义 141.5}（Hopf 分支）：标准形（极坐标）：
\[\dot{r} = \mu r - r^3, \quad \dot{\theta} = \omega\]

\begin{itemize}
\tightlist
\item
  \textbf{超临界 Hopf 分支}：\(\mu \leq 0\) 时原点稳定；\(\mu > 0\)
  时原点不稳定，出现稳定极限环（半径 \(\sqrt{\mu}\)）
\item
  \textbf{亚临界 Hopf
  分支}：\(\dot{r} = \mu r + r^3 - r^5\)（\(\mu < 0\)
  时已存在不稳定极限环）
\end{itemize}

\textbf{定理 141.2}（Hopf 分支定理，正式陈述）：设
\(\dot{\mathbf{x}} = \mathbf{f}(\mathbf{x}, \mu)\) 满足
\(\mathbf{f}(0, \mu) = \mathbf{0}\)，\(D\mathbf{f}(0, \mu)\) 有特征值
\(\lambda(\mu) = \alpha(\mu) \pm i\omega(\mu)\)（\(\alpha(0) = 0, \omega(0) > 0\)），其余特征值实部为负。若
\(\alpha'(0) \neq 0\)，则存在唯一的分支极限环族。分支方向由第一 Lyapunov
系数 \(\ell_1(0)\) 决定：\(\ell_1(0) < 0\) 为超临界，\(\ell_1(0) > 0\)
为亚临界。

\textbf{定义 141.6}（倍周期分支 / Period-doubling
bifurcation）：一维映射 \(f(x, \mu)\) 在 \(\mu = \mu_0\)
处发生倍周期分支，如果 \(f_{\mu_0}'(x^*) = -1\) 且
\(\frac{\partial}{\partial \mu}(f_\mu^2)''(x^*) \neq 0\)
等非退化条件成立。标准形：\(f(x, \mu) = -(1+\mu)x + x^3\)。

\subsubsection{141.3 全局分支}\label{ux5168ux5c40ux5206ux652f}

\textbf{定义 141.7}（同宿分支 / Homoclinic
bifurcation）：当稳定流形和不稳定流形在鞍点处相切时，产生同宿轨道，可能引发混沌（Shilnikov
机制）。

\textbf{定义 141.8}（异宿分支 / Heteroclinic
bifurcation）：连接不同平衡点的轨道（异宿轨道）在参数变化时出现或消失。

\textbf{定理 141.3}（Shilnikov
定理，1965）：若三维系统有鞍-焦点型平衡点，其特征值满足
\(\lambda_1 < 0\)，\(\lambda_{2,3} = \alpha \pm i\omega\)（\(\alpha > 0\)），且
\(|\lambda_1| < \alpha\)，则同宿轨道的存在性蕴含可数无穷多个 Smale
马蹄（混沌）。

\subsubsection{141.4 无穷维分支}\label{ux65e0ux7a77ux7ef4ux5206ux652f}

\textbf{定义 141.9}（Turing 分支，1952）：在反应扩散方程
\(\partial_t u = D \Delta u + f(u, \mu)\)
中，当扩散系数比例改变时，齐次平衡态失稳，产生空间非齐次图案（Turing
斑图）。

\textbf{定义 141.10}（Ginzburg-Landau
方程）：\(\partial_t A = \mu A + (1+i\alpha)\partial_x^2 A - (1+i\beta)|A|^2 A\)。这是
Hopf 分支在空间扩展系统中的标准形，描述了斑图形成的普适类。

\bigskip\hrule\bigskip

\emph{本卷基于
V18（微分方程），建立了离散和连续动力系统、混沌与分形、遍历论和分支理论的系统理论。共
5 章。}

\newpage

\section{卷二十八：信息论与编码}\label{ux5377ux4e8cux5341ux516bux4fe1ux606fux8bbaux4e0eux7f16ux7801}

\begin{quote}
信息论由 Claude Shannon 于 1948 年创立，是通信和存储的数学基础。本卷从
Shannon
熵和信息度量出发，建立信道容量与编码定理、纠错码理论和压缩感知引论，最后介绍密码学的数学基础。信息论与概率论、组合数学、数论和代数有深刻联系。
\end{quote}

\begin{quote}
\textbf{前置知识}：本卷依赖V20（概率论II）的核心内容，建议在修读本卷前先掌握概率空间、条件期望和随机过程的基础理论。
\end{quote}

\bigskip\hrule\bigskip

\subsection{第142章：Shannon
熵与信息度量}\label{ux7b2c142ux7ae0shannon-ux71b5ux4e0eux4fe1ux606fux5ea6ux91cf}

熵是信息论的核心概念，度量了随机变量的不确定性。本章从信息度量的公理化定义出发，建立熵、联合熵、条件熵和互信息的基本性质。

\subsubsection{142.1
熵的定义与公理化}\label{ux71b5ux7684ux5b9aux4e49ux4e0eux516cux7406ux5316}

\textbf{定义 142.1}（离散随机变量的熵）：设 \(X\) 是离散随机变量，取值于
\(\mathcal{X}\)，概率分布为 \(p(x) = \mathbb{P}(X = x)\)。\(X\)
的\textbf{熵}（Shannon 熵）定义为

\[H(X) = -\sum_{x \in \mathcal{X}} p(x) \log_2 p(x)\]

（约定
\(0 \log_2 0 = 0\)）。熵的单位是比特（bit）。若使用自然对数，单位为奈特（nat）。

\textbf{定理 142.1}（熵的公理化 / Khinchin 公理）：设
\(H(p_1, \ldots, p_n)\) 是定义在概率分布上的函数。若 \(H\) 满足： 1.
\(H\) 是 \(p_i\) 的连续函数 2. 对均匀分布，\(H(1/n, \ldots, 1/n)\) 关于
\(n\) 单调递增 3.
分组公理：\(H(p_1, \ldots, p_n) = H(p_1+p_2, p_3, \ldots, p_n) + (p_1+p_2) H\left(\frac{p_1}{p_1+p_2}, \frac{p_2}{p_1+p_2}\right)\)

则 \(H(p_1, \ldots, p_n) = -c \sum p_i \log p_i\)（\(c > 0\)）。

\textbf{命题 142.2}（熵的基本性质）： -
\(0 \leq H(X) \leq \log_2 |\mathcal{X}|\)（非负性，最大熵原理） -
\(H(X) = 0\) 当且仅当 \(X\) 是确定性的（退化分布） -
\(H(X) = \log_2 |\mathcal{X}|\) 当且仅当 \(X\) 是均匀分布

\textbf{例
142.1}（二元熵函数）：\(h(p) = -p \log_2 p - (1-p) \log_2 (1-p)\)。\(h(0) = h(1) = 0\)，\(h(1/2) = 1\)（最大）。\(h(p)\)
是凹函数，在 \([0, 1]\) 上关于 \(p=1/2\) 对称。

\subsubsection{142.2
联合熵、条件熵与互信息}\label{ux8054ux5408ux71b5ux6761ux4ef6ux71b5ux4e0eux4e92ux4fe1ux606f}

\textbf{定义 142.2}（联合熵）：\((X, Y)\) 的联合熵为

\[H(X, Y) = -\sum_{x, y} p(x, y) \log_2 p(x, y)\]

\textbf{定义 142.3}（条件熵）：给定 \(Y\) 下 \(X\) 的条件熵为

\[H(X|Y) = \sum_{y} p(y) H(X|Y=y) = -\sum_{x,y} p(x,y) \log_2 p(x|y)\]

\textbf{定理
142.3}（链式法则）：\(H(X, Y) = H(X) + H(Y|X) = H(Y) + H(X|Y)\)。

\textbf{定义 142.4}（互信息）：\(X\) 和 \(Y\) 之间的互信息为

\[I(X; Y) = H(X) - H(X|Y) = H(Y) - H(Y|X) = H(X) + H(Y) - H(X, Y)\]

互信息对称且非负：\(I(X; Y) \geq 0\)，\(I(X; Y) = 0\) 当且仅当 \(X\) 和
\(Y\) 独立。

\textbf{命题
142.4}（信息不等式）：\(H(X|Y) \leq H(X)\)，等号成立当且仅当 \(X\) 和
\(Y\) 独立。即「条件不会增加熵」。

\textbf{定理 142.5}（数据处理不等式）：若 \(X \to Y \to Z\) 构成 Markov
链（即 \(Z\) 条件独立于 \(X\) 给定 \(Y\)），则
\(I(X; Z) \leq I(X; Y)\)。数据处理不能增加信息。

\subsubsection{142.3
相对熵与交叉熵}\label{ux76f8ux5bf9ux71b5ux4e0eux4ea4ux53c9ux71b5}

\textbf{定义 142.5}（相对熵 / KL 散度）：分布 \(p\) 相对于 \(q\)
的相对熵（Kullback-Leibler 散度）为

\[D_{KL}(p \| q) = \sum_x p(x) \log_2 \frac{p(x)}{q(x)}\]

（约定 \(0 \log(0/0) = 0\)，\(p \log(p/0) = \infty\)）。

\textbf{定理 142.6}（Gibbs
不等式）：\(D_{KL}(p \| q) \geq 0\)，等号成立当且仅当 \(p = q\)。

\emph{证明}：由 \(\log t \leq t-1\)
的对数不等式，\(\sum p \log(q/p) \leq \sum p(q/p - 1) = \sum q - \sum p = 0\)。∎

\textbf{定义
142.6}（交叉熵）：\(H(p, q) = H(p) + D_{KL}(p \| q) = -\sum p(x) \log q(x)\)。交叉熵在机器学习中广泛用作损失函数。

\textbf{推论
142.7}：\(I(X; Y) = D_{KL}(p(x,y) \| p(x)p(y))\)。互信息是联合分布与独立分布之间的
KL 散度。

\subsubsection{142.4
微分熵与连续随机变量}\label{ux5faeux5206ux71b5ux4e0eux8fdeux7eedux968fux673aux53d8ux91cf}

\textbf{定义 142.7}（微分熵）：连续随机变量 \(X\)（密度
\(f(x)\)）的\textbf{微分熵}为

\[h(X) = -\int f(x) \log f(x) \, dx\]

微分熵可以取负值，且不是坐标变换下的不变量（与离散熵不同）。

\textbf{定理 142.8}（最大微分熵的分布）： - 固定方差 \(\sigma^2\)
下，正态分布 \(\mathcal{N}(\mu, \sigma^2)\)
最大化微分熵，\(h = \frac{1}{2}\log_2(2\pi e \sigma^2)\) - 固定支撑
\([a, b]\) 下，均匀分布最大化微分熵 - 固定均值 \(\lambda\)
下（\(x \geq 0\)），指数分布最大化微分熵

\textbf{定理
142.9}（熵幂不等式）：\(e^{2h(X+Y)/n} \geq e^{2h(X)/n} + e^{2h(Y)/n}\)（向量形式）。

\bigskip\hrule\bigskip

\subsection{第143章：信道容量与编码定理}\label{ux7b2c143ux7ae0ux4fe1ux9053ux5bb9ux91cfux4e0eux7f16ux7801ux5b9aux7406}

Shannon
的信道编码定理是信息论的核心成果，确定了在噪声信道中可靠通信的最大速率。

\subsubsection{143.1
信源编码定理}\label{ux4fe1ux6e90ux7f16ux7801ux5b9aux7406}

\textbf{定义
143.1}（信源编码）：将信源符号序列映射为二进制码字。\textbf{平均码长}
\(\bar{L} = \sum p(x) \ell(x)\)，其中 \(\ell(x)\) 是码字 \(x\) 的码长。

\textbf{定理 143.1}（无噪声编码定理 / Shannon 第一定理）：对信源
\(X\)，任何唯一可译码的平均码长满足
\(\bar{L} \geq H(X)\)。存在前缀码（如 Huffman 编码）满足
\(H(X) \leq \bar{L} < H(X) + 1\)。

\emph{证明概要}（典型序列 + AEP 构造法）：对信源 \(X\) 的 \(n\)
次独立输出 \(X^n\)，由渐近均分性（AEP），\(\varepsilon\)-典型集
\(A_\varepsilon^{(n)}\) 中约 \(2^{nH(X)}\) 个序列占据近乎全部概率（当
\(n\) 充分大时
\(P(A_\varepsilon^{(n)}) \to 1\)）。编码时，为典型序列分配等长码字（需
\(\approx nH(X)\)
比特），非典型序列则任意编码。由于非典型序列概率可忽略，平均每符号码长
\(\bar{L}/n \to H(X)\)。下界由 Kraft 不等式 \(\sum 2^{-\ell_i} \leq 1\)
结合 Jensen 不等式推出 \(\bar{L} \geq H(X)\)。块编码进一步压缩至极限
\(H(X)\)。∎

\textbf{定理 143.2}（Kraft 不等式）：长度为 \(\ell_1, \ldots, \ell_n\)
的前缀码存在当且仅当 \(\sum_{i=1}^n 2^{-\ell_i} \leq 1\)。

\textbf{算法 143.1}（Huffman 编码，1952）： 1.
将概率最小的两个符号合并为一个节点，其概率为两者之和 2.
重复直到只剩一个节点 3. 回溯分配码字（左 0 右 1）

Huffman 编码是最优前缀码（最小化平均码长）。

\textbf{定理 143.3}（信源编码定理，块编码版本）：对信源 \(X\)
的独立同分布序列 \(X^n\)，存在编码使得平均每符号码长
\(\frac{1}{n}\bar{L}_n \to H(X)\)（\(n \to \infty\)）。

\subsubsection{143.2 信道容量}\label{ux4fe1ux9053ux5bb9ux91cf}

\textbf{定义 143.2}（离散无记忆信道 / DMC）：信道由输入字母表
\(\mathcal{X}\)、输出字母表 \(\mathcal{Y}\) 和转移概率 \(p(y|x)\)
描述。\(n\) 次使用信道时，\(p(y^n|x^n) = \prod_{i=1}^n p(y_i|x_i)\)。

\textbf{定义 143.3}（信道容量）：离散无记忆信道的容量为

\[C = \max_{p(x)} I(X; Y)\]

即输入分布 \(p(x)\) 可任意选择时，互信息 \(I(X; Y)\) 的最大值。

\textbf{例 143.1}（经典信道的容量）： -
\textbf{无噪声信道}：\(C = \log_2 |\mathcal{X}|\) -
\textbf{二元对称信道}（BSC，交叉概率 \(p\)）：\(C = 1 - h(p)\) -
\textbf{二元擦除信道}（BEC，擦除概率
\(\varepsilon\)）：\(C = 1 - \varepsilon\) - \textbf{高斯信道}（噪声功率
\(\sigma^2\)，功率约束
\(P\)）：\(C = \frac{1}{2}\log_2\left(1 + \frac{P}{\sigma^2}\right)\)

\subsubsection{143.3
信道编码定理}\label{ux4fe1ux9053ux7f16ux7801ux5b9aux7406}

\textbf{定理 143.4}（Shannon 信道编码定理 / Shannon
第二定理，1948）：对离散无记忆信道，任何速率 \(R < C\)
是可达的：存在码率为 \(R\) 的编码方案，使得错误概率
\(P_e \to 0\)（当码长 \(n \to \infty\)）。反之，若
\(R > C\)，则不可能可靠通信。

\emph{证明思路}（随机编码 + 联合
AEP）：使用\textbf{随机编码}论证。随机生成 \(2^{nR}\)
个码字，每个符号独立从 \(p(x)\)
采样。译码器采用\textbf{联合典型译码}：收到 \(\mathbf{y}\)
后，寻找唯一的 \(\hat{\mathbf{x}}\) 使
\((\hat{\mathbf{x}}, \mathbf{y})\)
是联合典型的。分析错误概率：（1）发送码字与接收序列非联合典型：由
AEP，概率
\(\to 0\)；（2）存在另一个码字与接收序列联合典型：对固定错误码字
\(\mathbf{x}' \neq \mathbf{x}\)，\((\mathbf{x}', \mathbf{y})\)
联合典型（其中 \(\mathbf{x}'\) 独立于 \(\mathbf{y}\)）的概率
\(\approx 2^{-n I(X;Y)}\)。共有 \(2^{nR}-1\)
个错误码字，由联合界，错误概率
\(\leq 2^{nR} \cdot 2^{-n I(X;Y)} = 2^{-n(I(X;Y)-R)}\)。当
\(R < I(X;Y)\) 时，此概率 \(\to 0\)。取 \(p(x)\) 为达到信道容量 \(C\)
的输入分布，得 \(R < C\) 可达。由于平均错误概率趋于
0，存在具体的好码使错误概率趋于 0。逆定理：若 \(R > C\)，由 Fano
不等式可证错误概率有正下界。∎

\textbf{定义 143.4}（联合典型集）：长度为 \(n\) 的序列对 \((x^n, y^n)\)
是\textbf{联合典型的}，如果其经验分布接近
\(p(x, y)\)。具体地，\(A_\varepsilon^{(n)} = \{(x^n, y^n) : |-\frac{1}{n}\log p(x^n) - H(X)| < \varepsilon, |-\frac{1}{n}\log p(y^n) - H(Y)| < \varepsilon, |-\frac{1}{n}\log p(x^n, y^n) - H(X,Y)| < \varepsilon\}\)。

\textbf{定理 143.5}（联合 AEP / 渐近均分性）：对充分大的
\(n\)，联合典型集的大小约为 \(2^{nH(X,Y)}\)，且概率接近 1。

\textbf{定理 143.6}（信道编码逆定理）：若
\(R > C\)，则任何译码方案的错误概率有正的下界，\(P_e^{(n)} \geq 1 - C/R - 1/(nR)\)。

\subsubsection{143.4 率失真理论}\label{ux7387ux5931ux771fux7406ux8bba}

\textbf{定义 143.5}（率失真函数）：在允许失真 \(D\) 下，信源 \(X\)
的率失真函数为

\[R(D) = \min_{p(\hat{x}|x) : \mathbb{E}[d(X, \hat{X})] \leq D} I(X; \hat{X})\]

\textbf{定理 143.7}（率失真定理）：\(R(D)\) 是在失真 \(D\)
下信源编码可实现的最小码率。对 Gauss 信源（方差
\(\sigma^2\)），\(R(D) = \frac{1}{2}\log_2(\sigma^2/D)\)（\(0 \leq D \leq \sigma^2\)）。

\emph{证明思路}（正定理与逆定理）：可达性（正定理）由随机编码论证：从
\(p(\hat{x})\) 独立采样 \(2^{nR}\) 个重构码字，对信源输出 \(x^n\)
在允许失真 \(D\) 内寻找匹配码字。由典型集联合分布性质，当
\(R > I(X;\hat{X})\) 且 \(p(\hat{x}|x)\) 满足 \(\mathbb{E}[d] \leq D\)
时，编码失败概率 \(\to 0\)。取所有满足条件的 \(p(\hat{x}|x)\)
下互信息的最小值即得 \(R(D)\) 可达。逆定理：对任意达到失真 \(D\)
的编码方案，由 Fano 不等式和数据处理不等式可导出 \(R \geq R(D)\)。Gauss
信源情形下二次高斯失真测度的率失真函数是互信息在方差约束下的极小化结果，闭式解由反向水注法（reverse
water-filling）给出。∎

\bigskip\hrule\bigskip

\subsection{第144章：纠错码}\label{ux7b2c144ux7ae0ux7ea0ux9519ux7801}

纠错码通过在传输信息中添加冗余，使得接收端可以检测和纠正传输错误。纠错码理论是代数、组合数学和信息论的交汇。

\subsubsection{144.1 线性码}\label{ux7ebfux6027ux7801}

\textbf{定义 144.1}（线性码）：一个 \([n, k]\) 线性码 \(\mathcal{C}\) 是
\(\mathbb{F}_q^n\) 的 \(k\) 维子空间。\(n\) 为码长，\(k\)
为信息位数，\(R = k/n\) 为码率。

\textbf{定义 144.2}（生成矩阵与校验矩阵）： - \textbf{生成矩阵}
\(G\)（\(k \times n\)）：\(\mathcal{C} = \{\mathbf{m} G : \mathbf{m} \in \mathbb{F}_q^k\}\)
- \textbf{校验矩阵}
\(H\)（\((n-k) \times n\)）：\(\mathcal{C} = \{\mathbf{x} \in \mathbb{F}_q^n : H\mathbf{x}^T = \mathbf{0}\}\)

\(G H^T = 0\)。

\textbf{定义
144.3}（汉明距离）：\(\mathbf{x}, \mathbf{y} \in \mathbb{F}_q^n\)
的\textbf{汉明距离} \(d(\mathbf{x}, \mathbf{y})\)
是两者不同坐标的个数。码 \(\mathcal{C}\) 的\textbf{最小距离}
\(d = \min_{\mathbf{x} \neq \mathbf{y} \in \mathcal{C}} d(\mathbf{x}, \mathbf{y})\)。\([n, k, d]\)
码可纠正 \(\lfloor (d-1)/2 \rfloor\) 个错误。

\textbf{定理 144.1}（Singleton
界）：\(d \leq n - k + 1\)。达到此界的码称为\textbf{极大距离可分码}（MDS），如
Reed-Solomon 码。

\textbf{定理 144.2}（Hamming 界 /
球包界）：\(q^k \sum_{i=0}^{\lfloor (d-1)/2 \rfloor} \binom{n}{i} (q-1)^i \leq q^n\)。达到此界的码称为\textbf{完美码}。

\textbf{定理 144.3}（Gilbert-Varshamov 界）：存在 \([n, k, d]\)
线性码满足 \(q^{n-k} > \sum_{i=0}^{d-2} \binom{n-1}{i} (q-1)^i\)。

\subsubsection{144.2 经典线性码}\label{ux7ecfux5178ux7ebfux6027ux7801}

\textbf{例 144.1}（Hamming 码）：\([2^m-1, 2^m-1-m, 3]\)
二元码。校验矩阵的列是所有非零 \(m\) 维向量。Hamming 码是完美码，可纠正
1 个错误。

\textbf{例 144.2}（Reed-Solomon 码 / RS 码）：\([n, k, n-k+1]\)
码，定义在 \(\mathbb{F}_q\) 上（\(n \leq q-1\)）。取 \(\mathbb{F}_q\) 中
\(n\) 个不同元素 \(\alpha_1, \ldots, \alpha_n\)，信息多项式
\(m(x)\)（次数 \(< k\)），码字为
\((m(\alpha_1), \ldots, m(\alpha_n))\)。RS 码是 MDS 码，广泛用于
CD、DVD、QR 码和深空通信。

\textbf{定义 144.4}（BCH 码与 RS 码的译码）：Peterson-Gorenstein-Zierler
算法和 Berlekamp-Massey 算法可在多项式时间内纠正最多
\(\lfloor (d-1)/2 \rfloor\) 个错误。

\textbf{例 144.3}（Golay 码）：二元 Golay 码 \([23, 12, 7]\)
是完美码；三元 Golay 码 \([11, 6, 5]\) 也是完美码。两者与 Mathieu 群
\(M_{23}\) 和 \(M_{11}\) 密切相关。

\textbf{例 144.4}（LDPC 码 / 低密度奇偶校验码，Gallager 1963）：校验矩阵
\(H\) 是稀疏的。使用置信传播（Belief Propagation）迭代译码，可逼近
Shannon 容量。

\textbf{例 144.5}（Turbo
码，1993）：并行级联卷积码，迭代译码，首次在实用中逼近 Shannon 容量。

\textbf{例 144.6}（Polar 码，Arıkan
2009）：首个被严格证明可达二进制输入对称信道容量的构造性编码方案。核心思想是信道极化（channel
polarization）。

\subsubsection{144.3 代数几何码}\label{ux4ee3ux6570ux51e0ux4f55ux7801}

\textbf{定义 144.5}（代数几何码 / Goppa 码，1981）：在有限域
\(\mathbb{F}_q\) 上的代数曲线 \(\mathcal{X}\) 上，取有理点
\(P_1, \ldots, P_n\) 和除子 \(D\)。码由 Riemann-Roch 空间 \(L(D)\)
在点上的估值构成。当曲线亏格 \(g\) 较小时，AG 码可超越 Gilbert-Varshamov
界（Tsfasman-Vlăduţ-Zink，1982）。

\bigskip\hrule\bigskip

\subsection{第145章：压缩感知引论}\label{ux7b2c145ux7ae0ux538bux7f29ux611fux77e5ux5f15ux8bba}

压缩感知（Compressed Sensing / Compressive Sensing）是 2000 年代由
Candès、Romberg、Tao 和 Donoho
等人开创的革命性信号处理理论。它表明，稀疏信号可以从远少于 Nyquist
采样率的测量中精确恢复。

\subsubsection{\texorpdfstring{145.1 稀疏性与 \(\ell_1\)
最小化}{145.1 稀疏性与 \textbackslash ell\_1 最小化}}\label{ux7a00ux758fux6027ux4e0e-ell_1-ux6700ux5c0fux5316}

\textbf{定义 145.1}（稀疏信号）：向量 \(\mathbf{x} \in \mathbb{R}^n\) 是
\(s\)-稀疏的，如果
\(\|\mathbf{x}\|_0 = |\operatorname{supp}(\mathbf{x})| \leq s\)，其中
\(\|\cdot\|_0\) 是 \(\ell_0\) 伪范数（非零分量个数）。

\textbf{定义 145.2}（压缩感知问题）：给定测量矩阵
\(A \in \mathbb{R}^{m \times n}\)（\(m \ll n\)）和测量值
\(\mathbf{y} = A\mathbf{x}\)，从 \(\mathbf{y}\) 中恢复稀疏信号
\(\mathbf{x}\)。

\textbf{定理 145.1}（\(\ell_0\) 与 \(\ell_1\) 的等价性）：若 \(A\)
满足某些条件（如 RIP），则 \(\ell_0\) 最小化问题
\(\min \|\mathbf{x}\|_0 \text{ s.t. } A\mathbf{x} = \mathbf{y}\)（NP
难）与 \(\ell_1\) 最小化问题
\(\min \|\mathbf{x}\|_1 \text{ s.t. } A\mathbf{x} = \mathbf{y}\)（凸优化）给出相同的解。

\subsubsection{145.2
限制等距性质（RIP）}\label{ux9650ux5236ux7b49ux8dddux6027ux8d28rip}

\textbf{定义 145.3}（RIP）：矩阵 \(A\) 满足 \(s\)
阶\textbf{限制等距性质}（RIP），如果存在 \(\delta_s \in (0, 1)\)
使得对所有 \(s\)-稀疏向量 \(\mathbf{x}\)，

\[(1 - \delta_s)\|\mathbf{x}\|_2^2 \leq \|A\mathbf{x}\|_2^2 \leq (1 + \delta_s)\|\mathbf{x}\|_2^2\]

\textbf{定理 145.2}（Candès-Romberg-Tao，2006）：若 \(A\) 满足
\(\delta_{2s} < \sqrt{2} - 1\)（或更优的
\(\delta_{2s} < 1/\sqrt{2}\)），则 \(\ell_1\) 最小化可精确恢复任意
\(s\)-稀疏向量。

\textbf{定理 145.3}（随机矩阵的
RIP）：高斯随机矩阵（\(A_{ij} \sim \mathcal{N}(0, 1/m)\)）或 Bernoulli
随机矩阵以高概率满足 RIP，当 \(m \geq C \cdot s \log(n/s)\) 时。

\subsubsection{145.3 恢复算法}\label{ux6062ux590dux7b97ux6cd5}

\textbf{算法 145.1}（基础追踪 / Basis
Pursuit）：\(\min \|\mathbf{x}\|_1 \text{ s.t. } A\mathbf{x} = \mathbf{y}\)。可转化为线性规划求解。

\textbf{算法 145.2}（正交匹配追踪 /
OMP）：贪心迭代选择与当前残差最近相关的列，每次添加一个非零分量。

\textbf{算法 145.3}（迭代硬阈值 /
IHT）：\(\mathbf{x}^{k+1} = H_s(\mathbf{x}^k + A^T(\mathbf{y} - A\mathbf{x}^k))\)，其中
\(H_s\) 保留绝对值最大的 \(s\) 个分量。

\bigskip\hrule\bigskip

\subsection{第146章：密码学数学基础}\label{ux7b2c146ux7ae0ux5bc6ux7801ux5b66ux6570ux5b66ux57faux7840}

密码学是安全通信的科学，其数学基础植根于数论、代数、组合数学和计算复杂性理论。本章从传统密码开始，重点介绍现代公钥密码体制的数学结构。

\subsubsection{146.1 对称密码与 Shannon
理论}\label{ux5bf9ux79f0ux5bc6ux7801ux4e0e-shannon-ux7406ux8bba}

\textbf{定义 146.1}（Shannon 完美保密）：加密方案 \((Gen, Enc, Dec)\)
是\textbf{完美保密}的，如果对任意明文分布、任意明文 \(m\) 和密文
\(c\)，\(\mathbb{P}(M = m | C = c) = \mathbb{P}(M = m)\)。等价地，\(I(M; C) = 0\)。

\textbf{定理 146.1}（Shannon 定理）：完美保密方案必须满足
\(|\mathcal{K}| \geq |\mathcal{M}|\)（密钥空间不小于明文空间）。一次一密（One-Time
Pad / Vernam 密码）使用 \(c = m \oplus k\)（\(k\)
为真随机密钥）达到完美保密。

\textbf{定义
146.2}（计算安全）：现代密码学使用计算安全代替信息论安全。方案在多项式时间攻击者面前是安全的，假设某些计算问题（如因子分解）是难的。

\subsubsection{146.2 RSA 公钥密码}\label{rsa-ux516cux94a5ux5bc6ux7801}

\textbf{定义 146.3}（RSA 问题）：给定 \(N = pq\)（\(p, q\) 为大素数）和
\(e\)（\(\gcd(e, \phi(N)) = 1\)），从 \(c = m^e \bmod N\) 中恢复 \(m\)。

\textbf{算法 146.1}（RSA 密钥生成）： 1. 随机生成两个大素数
\(p, q\)，计算 \(N = pq\) 2. 计算 \(\phi(N) = (p-1)(q-1)\) 3. 选择 \(e\)
满足 \(\gcd(e, \phi(N)) = 1\)（常用 \(e = 65537\)） 4. 计算
\(d = e^{-1} \bmod \phi(N)\)（扩展 Euclid 算法） 5.
公钥：\((N, e)\)；私钥：\((N, d)\)

\textbf{算法 146.2}（RSA 加密/解密）： - 加密：\(c = m^e \bmod N\) -
解密：\(m = c^d \bmod N\)

\textbf{定理 146.2}（RSA 的正确性）：由 Euler
定理，\(m^{ed} \equiv m^{1 + k\phi(N)} \equiv m \pmod{N}\)（需
\(\gcd(m, N) = 1\)，概率上接近 1）。

\textbf{算法 146.3}（素性检测 /
Miller-Rabin）：随机化算法，可在多项式时间内以高概率判定素性。复杂度
\(O(k \log^3 n)\)，错误概率 \(< 4^{-k}\)。

\textbf{定理 146.3}（AKS
素性检测，2002）：存在确定性的多项式时间素性检测算法（Agrawal-Kayal-Saxena）。

\subsubsection{146.3 离散对数与
Diffie-Hellman}\label{ux79bbux6563ux5bf9ux6570ux4e0e-diffie-hellman}

\textbf{定义 146.4}（离散对数问题 / DLP）：在循环群
\(G = \langle g \rangle\) 中，给定 \(g^a\)，求 \(a\)。在适当的群 \(G\)
中，DLP 被认为是困难的。

\textbf{算法 146.4}（Diffie-Hellman 密钥交换，1976）：Alice 和 Bob
在不安全信道上协商共享密钥。 1. 公开参数：素数 \(p\)，生成元 \(g\) 2.
Alice 选择 \(a\)，发送 \(A = g^a \bmod p\) 3. Bob 选择 \(b\)，发送
\(B = g^b \bmod p\) 4.
共享密钥：\(K = B^a \equiv A^b \equiv g^{ab} \pmod{p}\)

安全性基于\textbf{计算 Diffie-Hellman 假设}（CDH）：从 \(g^a, g^b\) 计算
\(g^{ab}\) 是困难的。

\textbf{定义 146.5}（ElGamal 加密）：基于 Diffie-Hellman
的公钥加密方案。公钥为 \(h = g^x\)，加密：选择随机 \(r\)，密文为
\((g^r, m \cdot h^r)\)。解密：使用 \(x\) 恢复 \(m\)。

\textbf{定义 146.6}（椭圆曲线密码 / ECC）：使用椭圆曲线上的点群代替
\(\mathbb{F}_p^*\) 作为 DLP
的平台。优点是密钥更短、安全性更高。比特币使用 secp256k1 曲线。

\subsubsection{146.4
数字签名与哈希函数}\label{ux6570ux5b57ux7b7eux540dux4e0eux54c8ux5e0cux51fdux6570}

\textbf{定义 146.7}（数字签名方案）：由 \((Gen, Sign, Verify)\)
组成。签名者用私钥 \(sk\) 对消息 \(m\) 生成签名 \(\sigma\)，验证者用公钥
\(pk\) 验证 \(\sigma\)
的有效性。安全性要求：在选择消息攻击下，攻击者无法伪造签名（EUF-CMA）。

\textbf{算法 146.5}（RSA 签名）：\(\sigma = H(m)^d \bmod N\)（\(H\)
为密码学哈希函数）。验证：\(H(m) \stackrel{?}{=} \sigma^e \bmod N\)。

\textbf{算法 146.6}（DSA / ECDSA 签名）：基于离散对数的签名方案。ECDSA
是以太坊和比特币使用的签名算法。

\textbf{定义 146.8}（密码学哈希函数）：\(H: \{0,1\}^* \to \{0,1\}^n\)
满足： - \textbf{抗原像性}（preimage resistance）：给定 \(y\)，找到
\(x\) 使得 \(H(x) = y\) 是困难的 - \textbf{抗第二原像性}：给定
\(x\)，找到 \(x' \neq x\) 使得 \(H(x') = H(x)\) 是困难的 -
\textbf{抗碰撞性}（collision resistance）：找到 \(x \neq x'\) 使得
\(H(x) = H(x')\) 是困难的

\textbf{定理 146.4}（生日攻击）：找到碰撞的期望尝试次数为
\(O(2^{n/2})\)（生日悖论）。因此 SHA-256（\(n=256\)）提供 128
比特的碰撞安全性。

\subsubsection{146.5
量子密码引论}\label{ux91cfux5b50ux5bc6ux7801ux5f15ux8bba}

\textbf{定义 146.9}（Shor
算法，1994）：量子计算机可在多项式时间内分解大整数和求解离散对数，从而攻破
RSA 和 ECC。

\textbf{定理 146.5}（Grover 算法，1996）：量子计算机可在 \(O(\sqrt{N})\)
时间内搜索无序数据库（经典需 \(O(N)\)）。对称密码的密钥长度应加倍（如
AES-256）以抵抗量子攻击。

\textbf{定义 146.10}（后量子密码 /
PQC）：基于格（Lattice-based）、编码（Code-based）、多变量（Multivariate）和哈希签名（Hash-based）的密码方案。NIST
在 2024 年标准化了 CRYSTALS-Kyber（格密码）等方案。

\bigskip\hrule\bigskip

\emph{本卷基于 V6 Ch 29（概率论）和 V20（概率论 II），建立了 Shannon
信息论、纠错码和密码学的系统理论。共 5 章。}

\newpage

\section{卷二十九：数学物理方法}\label{ux5377ux4e8cux5341ux4e5dux6570ux5b66ux7269ux7406ux65b9ux6cd5}

\begin{quote}
数学物理是物理学和数学交汇的领域，用严格的数学语言描述物理世界的基本规律。本卷从变分原理和
Lagrange/Hamilton
力学出发，建立辛几何、量子力学数学基础、量子场论数学引论和广义相对论数学基础。本卷是微分几何（V12）、泛函分析（V14）、李理论（V23）和调和分析（V24）在物理学中的系统应用。
\end{quote}

\bigskip\hrule\bigskip

\subsection{第147章：变分原理与 Lagrange
力学}\label{ux7b2c147ux7ae0ux53d8ux5206ux539fux7406ux4e0e-lagrange-ux529bux5b66}

经典力学的 Lagrange
形式从最小作用量原理出发，用变分法导出运动方程。这一章与 V26 Ch 136
的变分法互补，聚焦于物理应用。

\subsubsection{147.1
最小作用量原理}\label{ux6700ux5c0fux4f5cux7528ux91cfux539fux7406}

\textbf{定义 147.1}（作用量）：\(N\) 个质点的力学系统由广义坐标
\(q^i(t)\)（\(i = 1, \ldots, n\)）描述。\textbf{作用量泛函}为

\[S[q] = \int_{t_1}^{t_2} L(q, \dot{q}, t) \, dt\]

其中 \(L = T - V\) 是 Lagrange 函数（动能减势能）。

\textbf{原理 147.1}（Hamilton 最小作用量原理）：系统在 \(t_1\) 到
\(t_2\) 之间实际运动的轨迹 \(q(t)\) 是作用量 \(S[q]\)
的驻点（通常是极小值），即 \(\delta S = 0\)。

\textbf{定理 147.1}（Euler-Lagrange 方程）：由 \(\delta S = 0\) 导出

\[\frac{d}{dt}\left( \frac{\partial L}{\partial \dot{q}^i} \right) - \frac{\partial L}{\partial q^i} = 0, \quad i = 1, \ldots, n\]

\emph{证明}：考虑变分
\(q^i(t) \to q^i(t) + \varepsilon \eta^i(t)\)，其中
\(\eta^i(t_1) = \eta^i(t_2) = 0\)（固定端点）。作用量的变分为
\[\delta S = \int_{t_1}^{t_2} \sum_i \left( \frac{\partial L}{\partial q^i} \varepsilon \eta^i + \frac{\partial L}{\partial \dot{q}^i} \varepsilon \dot{\eta}^i \right) dt\]
对第二项分部积分：\(\int_{t_1}^{t_2} \frac{\partial L}{\partial \dot{q}^i} \dot{\eta}^i dt = -\int_{t_1}^{t_2} \frac{d}{dt}\!\left(\frac{\partial L}{\partial \dot{q}^i}\right) \eta^i dt\)（边界项因
\(\eta^i(t_1)=\eta^i(t_2)=0\) 消失）。故
\[\delta S = \varepsilon \int_{t_1}^{t_2} \sum_i \left( \frac{\partial L}{\partial q^i} - \frac{d}{dt}\frac{\partial L}{\partial \dot{q}^i} \right) \eta^i dt = 0\]
由变分 \(\eta^i(t)\) 的任意性及变分法基本引理，被积函数恒为零，即得
Euler-Lagrange 方程。∎

\textbf{例
147.1}（自由粒子）：\(L = \frac{1}{2}m\dot{\mathbf{x}}^2\)，Euler-Lagrange
方程给出 \(\ddot{\mathbf{x}} = 0\)（匀速直线运动）。

\textbf{例
147.2}（简谐振子）：\(L = \frac{1}{2}m\dot{x}^2 - \frac{1}{2}kx^2\)，运动方程
\(m\ddot{x} + kx = 0\)。

\textbf{例
147.3}（中心力场中的质点）：\(L = \frac{1}{2}m(\dot{r}^2 + r^2\dot{\theta}^2) - V(r)\)。由
\(\theta\) 的 Euler-Lagrange 方程得
\(mr^2\dot{\theta} = \ell\)（角动量守恒）。

\subsubsection{147.2 Noether 定理}\label{noether-ux5b9aux7406}

\textbf{定理 147.2}（Noether 定理，1918）：Lagrange
系统的每一个连续对称性对应一个守恒量。

具体地，若作用量在变换 \(q^i \to q^i + \varepsilon \xi^i(q, \dot{q})\)
和 \(t \to t + \varepsilon \tau(q, \dot{q})\)
下不变（\(\delta S = 0\)），则

\[Q = \sum_i \frac{\partial L}{\partial \dot{q}^i} \xi^i - \left( \sum_i \frac{\partial L}{\partial \dot{q}^i} \dot{q}^i - L \right) \tau\]

是守恒量（\(\frac{dQ}{dt} = 0\)）。

\emph{证明概要}：考虑无穷小变换
\(q^i \to q^i + \varepsilon \xi^i\)，\(t \to t + \varepsilon \tau\)。计算作用量的变分
\(\delta S = \int_{t_1}^{t_2} [\frac{\partial L}{\partial q^i} \delta q^i + \frac{\partial L}{\partial \dot{q}^i} \delta \dot{q}^i + L \frac{d(\delta t)}{dt}] dt\)。使用
\(\delta q^i = \varepsilon \xi^i - \dot{q}^i \delta t\) 及分部积分，令
\(\delta S = 0\) 对任意积分区间成立，得到
\(\frac{d}{dt}[\frac{\partial L}{\partial \dot{q}^i} \xi^i - (\frac{\partial L}{\partial \dot{q}^i} \dot{q}^i - L)\tau] = (\frac{\partial L}{\partial q^i} - \frac{d}{dt}\frac{\partial L}{\partial \dot{q}^i})(\xi^i - \dot{q}^i \tau)\)。在
Euler-Lagrange 方程的壳层（on-shell）上，右端为0，故 \(Q\) 守恒。∎

\textbf{例 147.4}（Noether 定理的经典应用）： - \textbf{时间平移不变性}
\(\implies\)
能量守恒：\(E = \sum_i \frac{\partial L}{\partial \dot{q}^i} \dot{q}^i - L\)
- \textbf{空间平移不变性} \(\implies\)
动量守恒：\(p_i = \frac{\partial L}{\partial \dot{q}^i}\) -
\textbf{旋转不变性} \(\implies\) 角动量守恒

\subsubsection{147.3 约束与 Lagrange
乘子}\label{ux7ea6ux675fux4e0e-lagrange-ux4e58ux5b50}

\textbf{定义 147.2}（约束系统）：若坐标受约束
\(f_\alpha(q, t) = 0\)（\(\alpha = 1, \ldots, m\)），则使用 Lagrange
乘子 \(\lambda_\alpha\)：

\[L' = L + \sum_{\alpha=1}^m \lambda_\alpha(t) f_\alpha(q, t)\]

运动方程为
\(\frac{d}{dt}\frac{\partial L}{\partial \dot{q}^i} - \frac{\partial L}{\partial q^i} = \sum_\alpha \lambda_\alpha \frac{\partial f_\alpha}{\partial q^i}\)。

\textbf{定理 147.3}（d'Alembert 原理）：约束力的虚功为零
\(\sum_i \mathbf{F}_i^{\text{constraint}} \cdot \delta \mathbf{r}_i = 0\)，这是
Lagrange 力学约束处理的基础。

\bigskip\hrule\bigskip

\subsection{第148章：Hamilton
力学与辛几何}\label{ux7b2c148ux7ae0hamilton-ux529bux5b66ux4e0eux8f9bux51e0ux4f55}

Hamilton 力学将 Lagrange 力学的二阶 ODE 系统转化为一阶 ODE
系统，并揭示了经典力学背后的辛几何结构。

\subsubsection{148.1 Legendre 变换与 Hamilton
方程}\label{legendre-ux53d8ux6362ux4e0e-hamilton-ux65b9ux7a0b}

\textbf{定义 148.1}（广义动量与 Hamilton 函数）：对 Lagrange 函数
\(L(q, \dot{q}, t)\)，定义广义动量

\[p_i = \frac{\partial L}{\partial \dot{q}^i}\]

和 Hamilton 函数（Legendre 变换）

\[H(q, p, t) = \sum_i p_i \dot{q}^i - L(q, \dot{q}, t)\]

其中 \(\dot{q}^i\) 需用 \((q, p)\) 表示。

\textbf{定理 148.1}（Hamilton 正则方程）：

\[\dot{q}^i = \frac{\partial H}{\partial p_i}, \quad \dot{p}_i = -\frac{\partial H}{\partial q^i}\]

\emph{证明}：由 Legendre 变换
\(H = \sum_i p_i \dot{q}^i - L\)，计算全微分：
\[dH = \sum_i \left( \dot{q}^i dp_i + p_i d\dot{q}^i - \frac{\partial L}{\partial q^i} dq^i - \frac{\partial L}{\partial \dot{q}^i} d\dot{q}^i \right) - \frac{\partial L}{\partial t} dt\]
由 \(p_i = \frac{\partial L}{\partial \dot{q}^i}\)，\(d\dot{q}^i\)
项抵消。又由 Euler-Lagrange 方程
\(\dot{p}_i = \frac{d}{dt}\frac{\partial L}{\partial \dot{q}^i} = \frac{\partial L}{\partial q^i}\)。比较
\(dH = \sum_i (\frac{\partial H}{\partial q^i} dq^i + \frac{\partial H}{\partial p_i} dp_i) + \frac{\partial H}{\partial t} dt\)，即得
\(\dot{q}^i = \frac{\partial H}{\partial p_i}\)，\(\dot{p}_i = -\frac{\partial H}{\partial q^i}\)。等价性得证。∎

\textbf{例 148.1}（简谐振子的 Hamilton
形式）：\(H = \frac{p^2}{2m} + \frac{1}{2}kq^2\)。正则方程：\(\dot{q} = p/m\)，\(\dot{p} = -kq\)。

\subsubsection{148.2 Poisson 括号}\label{poisson-ux62ecux53f7}

\textbf{定义 148.2}（Poisson 括号）：相空间上的函数 \(f(q, p)\) 和
\(g(q, p)\) 的 Poisson 括号为

\[\{f, g\} = \sum_i \left( \frac{\partial f}{\partial q^i} \frac{\partial g}{\partial p_i} - \frac{\partial f}{\partial p_i} \frac{\partial g}{\partial q^i} \right)\]

\textbf{命题 148.2}（Poisson 括号的性质）： -
反对称：\(\{f, g\} = -\{g, f\}\) - 双线性 - Leibniz
法则：\(\{fg, h\} = f\{g, h\} + \{f, h\}g\) - Jacobi
恒等式：\(\{f, \{g, h\}\} + \{g, \{h, f\}\} + \{h, \{f, g\}\} = 0\)

即 Poisson 括号使相空间上的光滑函数构成 Lie 代数。

\textbf{定理 148.3}（运动方程）：对任意可观测量
\(f(q, p, t)\)，\(\frac{df}{dt} = \{f, H\} + \frac{\partial f}{\partial t}\)。特别地，\(f\)
是守恒量当且仅当 \(\{f, H\} = 0\)（且 \(f\) 不显含时）。

\textbf{基本 Poisson
括号}：\(\{q^i, q^j\} = 0\)，\(\{p_i, p_j\} = 0\)，\(\{q^i, p_j\} = \delta^i_j\)。

\subsubsection{148.3 辛几何}\label{ux8f9bux51e0ux4f55}

\textbf{定义 148.3}（辛形式）：相空间 \(T^*Q\)（余切丛）上的典范辛形式为

\[\omega = \sum_i dq^i \wedge dp_i\]

\(\omega\) 是非退化的闭 2-形式（\(d\omega = 0\)）。

\textbf{定义 148.4}（辛流形）：辛流形 \((M, \omega)\)
是光滑流形配以非退化闭 2-形式 \(\omega\)。

\textbf{定理 148.4}（Darboux 定理）：任何辛流形局部上存在坐标
\((q^i, p_i)\) 使得
\(\omega = \sum_i dq^i \wedge dp_i\)。即所有辛流形局部同构。

\textbf{定义 148.5}（Hamilton 向量场）：对函数 \(H \in C^\infty(M)\)，其
Hamilton 向量场 \(X_H\) 由 \(\iota_{X_H} \omega = dH\)（即
\(\omega(X_H, \cdot) = dH\)）定义。在 Darboux
坐标下，\(X_H = \sum_i (\frac{\partial H}{\partial p_i} \frac{\partial}{\partial q^i} - \frac{\partial H}{\partial q^i} \frac{\partial}{\partial p_i})\)。

\textbf{定理 148.5}（Liouville 定理）：Hamilton 流保持相空间体积（辛体积
\(\omega^n\)）。即 \(\phi_t^* \omega = \omega\)，故
\(\phi_t^*(\omega^n) = \omega^n\)。

\textbf{定理 148.6}（Poincaré 回复定理）：对保测变换（包括 Hamilton
流），紧致相空间上的几乎每个点都是回复的。

\subsubsection{148.4 Hamilton-Jacobi
理论与可积系统}\label{hamilton-jacobi-ux7406ux8bbaux4e0eux53efux79efux7cfbux7edf}

\textbf{定义 148.6}（Hamilton-Jacobi 方程）：求函数 \(S(q, t)\) 满足

\[\frac{\partial S}{\partial t} + H\left(q, \frac{\partial S}{\partial q}, t\right) = 0\]

\textbf{定理 148.7}（Liouville-Arnold 可积性定理）：若 \(n\) 自由度的
Hamilton 系统有 \(n\)
个相互对合（\(\{F_i, F_j\} = 0\)）的独立守恒量，则系统可通过作用-角变量
\((I_i, \theta_i)\) 完全积分。运动在 \(n\)
维环面上，\(\dot{\theta}_i = \omega_i(I)\)（常数），\(\dot{I}_i = 0\)。

\bigskip\hrule\bigskip

\subsection{第149章：量子力学数学基础}\label{ux7b2c149ux7ae0ux91cfux5b50ux529bux5b66ux6570ux5b66ux57faux7840}

量子力学用 Hilbert
空间上的算子描述物理系统，其数学基础是泛函分析和谱理论。本章建立量子力学的公理化框架。

\subsubsection{149.1
量子力学的公理化}\label{ux91cfux5b50ux529bux5b66ux7684ux516cux7406ux5316}

\textbf{公理 149.1}（量子态）：量子系统的状态由复可分 Hilbert 空间
\(\mathcal{H}\) 上的密度算子
\(\rho\)（\(\rho \geq 0, \operatorname{Tr}(\rho) = 1\)）描述。纯态对应于一维投影
\(\rho = |\psi\rangle\langle\psi|\)（\(\|\psi\| = 1\)）。

\textbf{公理 149.2}（可观测量）：物理可观测量由 \(\mathcal{H}\)
上的自伴算子 \(A = A^\dagger\) 描述。\(A\) 的可能测量值是其谱
\(\sigma(A)\) 中的元素。

\textbf{公理 149.3}（测量 / Born 法则）：在态 \(\rho\) 中测量可观测量
\(A\)，结果在 Borel 集 \(E \subseteq \mathbb{R}\) 中的概率为
\(\mathbb{P}(E) = \operatorname{Tr}(\rho P_A(E))\)，其中 \(P_A\) 是
\(A\) 的谱测度（投影值测度，PVM）。

\textbf{公理 149.4}（时间演化 / Schrödinger 方程）：态的时间演化由

\[i\hbar \frac{d}{dt}|\psi(t)\rangle = \hat{H} |\psi(t)\rangle\]

其中 \(\hat{H}\) 是 Hamilton 算子（自伴算子）。形式解为
\(|\psi(t)\rangle = e^{-i\hat{H}t/\hbar} |\psi(0)\rangle\)。

\emph{导出概要}：从 de Broglie 关系
\(E = \hbar\omega\)、\(\mathbf{p} = \hbar\mathbf{k}\) 出发。平面波
\(\psi(x,t) = e^{i(kx - \omega t)}\) 满足
\(i\hbar\frac{\partial\psi}{\partial t} = \hbar\omega\psi = E\psi\) 和
\(-i\hbar\frac{\partial\psi}{\partial x} = \hbar k\psi = p\psi\)。将经典
Hamilton 函数 \(H = p^2/2m + V(x)\) 中的 \(p\) 替换为算子
\(-i\hbar\frac{\partial}{\partial x}\)，\(E\) 替换为算子
\(i\hbar\frac{\partial}{\partial t}\)，得到
\(i\hbar\frac{\partial}{\partial t}\psi = \left(-\frac{\hbar^2}{2m}\frac{\partial^2}{\partial x^2} + V(x)\right)\psi\)，即
Schrödinger 方程。更一般地，任何物理系统的哈密顿量 \(\hat{H}\)
通过正则量子化（\([\hat{x}, \hat{p}] = i\hbar\)）作用于波函数，给出时间演化方程。∎

\textbf{定理 149.1}（Stone 定理）：时间演化
\(U(t) = e^{-i\hat{H}t/\hbar}\) 是 \(\mathcal{H}\) 上的强连续酉算子群。

\textbf{公理 149.5}（复合系统）：若系统 \(A\) 和 \(B\) 的 Hilbert
空间分别为 \(\mathcal{H}_A\) 和 \(\mathcal{H}_B\)，则复合系统的 Hilbert
空间为 \(\mathcal{H}_A \otimes \mathcal{H}_B\)（张量积）。

\subsubsection{149.2
谱理论在量子力学中的应用}\label{ux8c31ux7406ux8bbaux5728ux91cfux5b50ux529bux5b66ux4e2dux7684ux5e94ux7528}

\textbf{定义 149.1}（谱分解）：由谱定理，自伴算子 \(\hat{A}\) 有谱分解

\[\hat{A} = \int_{\sigma(\hat{A})} \lambda \, dP_A(\lambda)\]

\textbf{例 149.1}（位置与动量算子）：在 \(L^2(\mathbb{R})\) 上， -
位置算子：\((\hat{x}\psi)(x) = x\psi(x)\) -
动量算子：\((\hat{p}\psi)(x) = -i\hbar \frac{d\psi}{dx}\)

\textbf{定理 149.2}（正则对易关系 /
CCR）：\([\hat{x}, \hat{p}] = i\hbar \hat{I}\)。

\textbf{定理 149.3}（Stone-von Neumann 定理）：CCR
的不可约表示在酉等价意义下唯一，即标准 Schrödinger 表示。

\textbf{定理 149.4}（Heisenberg 不确定性原理）：对任意态
\(|\psi\rangle\) 和可观测量 \(A, B\)，

\[\Delta_\psi(A) \cdot \Delta_\psi(B) \geq \frac{1}{2} |\langle \psi | [A, B] | \psi \rangle|\]

特别地，\(\Delta x \cdot \Delta p \geq \frac{\hbar}{2}\)。

\emph{证明}：由 Cauchy-Schwarz 不等式和自伴算子的性质。∎

\subsubsection{149.3
几个精确可解模型}\label{ux51e0ux4e2aux7cbeux786eux53efux89e3ux6a21ux578b}

\textbf{例
149.2}（简谐振子）：\(\hat{H} = \frac{\hat{p}^2}{2m} + \frac{1}{2}m\omega^2\hat{x}^2\)。引入升降算子
\(a = \sqrt{\frac{m\omega}{2\hbar}}(\hat{x} + \frac{i}{m\omega}\hat{p})\)，\(a^\dagger = \sqrt{\frac{m\omega}{2\hbar}}(\hat{x} - \frac{i}{m\omega}\hat{p})\)。\([a, a^\dagger] = 1\)。\(\hat{H} = \hbar\omega(a^\dagger a + \frac{1}{2})\)。能谱为
\(E_n = \hbar\omega(n + \frac{1}{2})\)，\(n = 0, 1, 2, \ldots\)。

\textbf{例
149.3}（氢原子）：\(\hat{H} = -\frac{\hbar^2}{2m}\nabla^2 - \frac{e^2}{4\pi\varepsilon_0 r}\)。利用球对称性，分离变量求解。能级为
\(E_n = -\frac{me^4}{2(4\pi\varepsilon_0)^2\hbar^2} \frac{1}{n^2}\)（\(n = 1, 2, \ldots\)）。能级简并度
\(n^2\) 由 \(SO(4)\) 对称性解释（Lenz-Runge 向量）。

\subsubsection{149.4 量子力学的
C*-代数表述}\label{ux91cfux5b50ux529bux5b66ux7684-c-ux4ee3ux6570ux8868ux8ff0}

\textbf{定义 149.2}（量子力学的代数方法）：量子系统的可观测量构成
C*-代数 \(\mathfrak{A}\)（或 von Neumann 代数）。态是 \(\mathfrak{A}\)
上的正线性泛函 \(\omega\)（\(\omega(1) = 1\)）。GNS 构造将代数态转化为
Hilbert 空间表示。

\textbf{定理 149.5}（Gelfand-Naimark-Segal 构造）：对 C*-代数
\(\mathfrak{A}\) 上的任意态 \(\omega\)，存在 Hilbert 空间
\(\mathcal{H}_\omega\)、表示
\(\pi_\omega: \mathfrak{A} \to \mathcal{B}(\mathcal{H}_\omega)\)
和循环向量 \(\Omega_\omega\)，使得
\(\omega(A) = \langle \Omega_\omega | \pi_\omega(A) | \Omega_\omega \rangle\)。

\subsubsection{149.5 经典电磁场与 Maxwell
方程组}\label{ux7ecfux5178ux7535ux78c1ux573aux4e0e-maxwell-ux65b9ux7a0bux7ec4}

\textbf{定理 149.6}（Maxwell 方程组的微分形式）：在真空中，电磁场满足

\[\nabla \cdot \mathbf{E} = \frac{\rho}{\varepsilon_0}, \quad \nabla \cdot \mathbf{B} = 0\]

\[\nabla \times \mathbf{E} = -\frac{\partial \mathbf{B}}{\partial t}, \quad \nabla \times \mathbf{B} = \mu_0 \mathbf{J} + \mu_0\varepsilon_0 \frac{\partial \mathbf{E}}{\partial t}\]

\emph{简要说明}：Gauss 电场定律（散度）源于 Coulomb
定律和叠加原理的微分表述；Gauss
磁场定律（\(\nabla \cdot \mathbf{B} = 0\)）反映磁单极子不存在；Faraday
定律（旋度）来自变化的磁场感生电场；Ampère-Maxwell 定律中位移电流项
\(\mu_0\varepsilon_0\frac{\partial\mathbf{E}}{\partial t}\) 是 Maxwell
的核心贡献------保证了电荷守恒
\(\nabla \cdot \mathbf{J} + \frac{\partial\rho}{\partial t} = 0\)
与方程组的自洽性。引入四维电磁势 \(A^\mu = (\phi/c, \mathbf{A})\)
和场强张量
\(F_{\mu\nu} = \partial_\mu A_\nu - \partial_\nu A_\mu\)，Maxwell
方程组可紧凑地写为 \(\partial_\mu F^{\mu\nu} = \mu_0 J^\nu\) 和
\(\partial_{[\lambda}F_{\mu\nu]} = 0\)。这是规范场 \(U(1)\)
的经典原型。∎

\bigskip\hrule\bigskip

\subsection{第150章：量子场论数学引论}\label{ux7b2c150ux7ae0ux91cfux5b50ux573aux8bbaux6570ux5b66ux5f15ux8bba}

量子场论（QFT）是量子力学和狭义相对论的结合，是粒子物理标准模型的数学基础。本章介绍
QFT 的数学结构，重点关注路径积分表述和公理化方法。

\subsubsection{150.1 自由量子场}\label{ux81eaux7531ux91cfux5b50ux573a}

\textbf{定义 150.1}（Klein-Gordon 场）：自由标量场 \(\phi(x, t)\) 满足
Klein-Gordon 方程

\[\partial_\mu \partial^\mu \phi + m^2 \phi = 0\]

即 \((\square + m^2)\phi = 0\)。

\textbf{定义 150.2}（正则量子化）：将场 \(\phi(x)\) 和共轭动量
\(\pi(x) = \partial_t \phi(x)\) 提升为算子值分布，满足等时对易关系：

\[[\phi(x, t), \pi(y, t)] = i\hbar \delta(x - y)\]

\textbf{定义 150.3}（Fock 空间）：自由量子场的状态空间是 Fock 空间
\(\mathcal{F} = \bigoplus_{n=0}^\infty \mathcal{H}^{\otimes_s n}\)（玻色子）或
\(\mathcal{F} = \bigoplus_{n=0}^\infty \mathcal{H}^{\wedge n}\)（费米子）。\(\mathcal{H} = L^2(\mathbb{R}^3, d^3p/2E_p)\)
是单粒子 Hilbert 空间。

\textbf{定理 150.1}（Wightman 公理）：相对论量子场论应满足： 1. 存在
Hilbert 空间 \(\mathcal{H}\) 和 Poincaré 群的酉表示 \(U(a, \Lambda)\) 2.
存在唯一的真空态 \(\Omega\)（Poincaré 不变） 3. 场
\(\phi(f)\)（\(f \in \mathcal{S}(\mathbb{R}^4)\)）是 \(\mathcal{H}\)
上的算子值缓增分布 4. 谱条件：能量-动量算子的谱在向前光锥内 5.
微观因果性：若 \(f\) 和 \(g\) 的支集类空分离，则
\([\phi(f), \phi(g)] = 0\)（玻色子）或
\(\{\phi(f), \phi(g)\} = 0\)（费米子） 6.
循环性：\(\{\phi(f_1)\cdots\phi(f_n)\Omega\}\) 在 \(\mathcal{H}\) 中稠密

\subsubsection{150.2 路径积分}\label{ux8defux5f84ux79efux5206}

\textbf{定义 150.4}（Feynman 路径积分）：量子力学中的传播子可形式地表为

\[\langle x_f | e^{-i\hat{H}T/\hbar} | x_i \rangle = \int_{x(0)=x_i}^{x(T)=x_f} \mathcal{D}x(t) \, e^{iS[x]/\hbar}\]

其中 \(S[x]\) 是经典作用量，\(\mathcal{D}x(t)\) 是形式上的「路径测度」。

\textbf{定义 150.5}（Wick 旋转与 Euclidean QFT）：做 Wick 旋转
\(t \to -i\tau\)，将 Minkowski 时空变为 Euclidean 时空。路径积分变为

\[Z = \int \mathcal{D}\phi \, e^{-S_E[\phi]/\hbar}\]

其中
\(S_E[\phi] = \int d^4x \left( \frac{1}{2}(\partial_\mu\phi)^2 + \frac{1}{2}m^2\phi^2 + V(\phi) \right)\)
是 Euclidean 作用量。

\textbf{定理 150.2}（Osterwalder-Schrader 重构定理，1973-1975）：满足
Osterwalder-Schrader 公理的 Euclidean
关联函数可以通过解析延拓唯一地重构满足 Wightman 公理的相对论量子场论。

\subsubsection{150.3 规范场论}\label{ux89c4ux8303ux573aux8bba}

\textbf{定义 150.6}（Yang-Mills 理论）：规范群 \(G\)（如
\(SU(3) \times SU(2) \times U(1)\)）的 Yang-Mills 作用量为

\[S_{YM}[A] = -\frac{1}{4} \int d^4x \, \operatorname{Tr}(F_{\mu\nu}F^{\mu\nu})\]

其中
\(F_{\mu\nu} = \partial_\mu A_\nu - \partial_\nu A_\mu + [A_\mu, A_\nu]\)
是规范场强，\(A_\mu\) 取值于 \(G\) 的李代数。

\textbf{定理 150.3}（Yang-Mills 存在性与质量间隙 / Clay
问题）：证明对任意紧致单纯规范群 \(G\)，\(\mathbb{R}^4\) 上的量子
Yang-Mills 理论存在，且具有质量间隙
\(\Delta > 0\)。这是七大千禧年问题之一，尚未解决。

\textbf{定义 150.7}（BRST 量子化）：使用 BRST
对称性（Becchi-Rouet-Stora-Tyutin）处理规范对称性的量子化。引入鬼场（ghost
fields）和 BRST 算子 \(Q\)，物理态空间为 \(Q\)-上同调
\(\ker Q / \operatorname{im} Q\)。

\textbf{定理 150.4}（Higgs
机制）：在规范场与标量场耦合时，若标量场的真空期望值非零，则规范对称性自发破缺，规范玻色子获得质量。这是电弱统一理论的核心。

\bigskip\hrule\bigskip

\subsection{第151章：广义相对论数学基础}\label{ux7b2c151ux7ae0ux5e7fux4e49ux76f8ux5bf9ux8bbaux6570ux5b66ux57faux7840}

广义相对论（Einstein, 1915）将引力解释为时空的弯曲，其数学框架是伪
Riemann 几何。本章建立广义相对论的核心数学结构。

\subsubsection{151.1 Lorentz
流形与时空}\label{lorentz-ux6d41ux5f62ux4e0eux65f6ux7a7a}

\textbf{定义 151.1}（时空）：\textbf{时空}是四维 Lorentz 流形
\((M, g)\)，其中 \(g\) 是 Lorentz 度量（符号差 \((-, +, +, +)\) 或
\((+, -, -, -)\)）。每点的切空间 \(T_pM\)
被光锥分为类时、类光和类空向量。

\textbf{定义
151.2}（测地线与自由落体）：自由落体粒子沿类时测地线运动，满足

\[\frac{d^2 x^\mu}{d\tau^2} + \Gamma^\mu_{\nu\rho} \frac{dx^\nu}{d\tau} \frac{dx^\rho}{d\tau} = 0\]

其中 \(\tau\) 是固有时（proper time），\(\Gamma^\mu_{\nu\rho}\) 是
Christoffel 符号。

\textbf{定义 151.3}（Riemann 曲率张量）：曲率描述了时空的弯曲程度：

\[R^\rho_{\sigma\mu\nu} = \partial_\mu \Gamma^\rho_{\nu\sigma} - \partial_\nu \Gamma^\rho_{\mu\sigma} + \Gamma^\rho_{\mu\lambda}\Gamma^\lambda_{\nu\sigma} - \Gamma^\rho_{\nu\lambda}\Gamma^\lambda_{\mu\sigma}\]

\textbf{Ricci
张量}：\(R_{\mu\nu} = R^\rho_{\mu\rho\nu}\)，\textbf{标量曲率}：\(R = g^{\mu\nu}R_{\mu\nu}\)。

\subsubsection{151.2 Einstein 场方程}\label{einstein-ux573aux65b9ux7a0b}

\textbf{定理 151.1}（Einstein 场方程，1915）：

\[R_{\mu\nu} - \frac{1}{2}R g_{\mu\nu} + \Lambda g_{\mu\nu} = \frac{8\pi G}{c^4} T_{\mu\nu}\]

其中 \(T_{\mu\nu}\) 是能动张量，\(\Lambda\) 是宇宙学常数，\(G\)
是牛顿引力常数。

\textbf{定义 151.4}（Einstein-Hilbert 作用量）：Einstein
方程可由变分原理导出。作用量为

\[S = \frac{c^4}{16\pi G} \int_M (R - 2\Lambda) \sqrt{-g} \, d^4x + S_{\text{matter}}\]

\textbf{定理 151.2}（Bianchi
恒等式的几何意义）：\(\nabla^\mu G_{\mu\nu} = 0\)（其中
\(G_{\mu\nu} = R_{\mu\nu} - \frac{1}{2}R g_{\mu\nu}\) 是 Einstein
张量），这是微分几何的 Bianchi 恒等式
\(\nabla_{[\lambda} R_{\mu\nu]\rho\sigma} = 0\)
的推论，保证了能动张量的协变守恒 \(\nabla^\mu T_{\mu\nu} = 0\)。

\subsubsection{151.3 经典精确解}\label{ux7ecfux5178ux7cbeux786eux89e3}

\textbf{例 151.1}（Schwarzschild
解，1916）：球对称真空解（\(\Lambda = 0\)）：

\[ds^2 = -\left(1 - \frac{2GM}{c^2 r}\right)c^2 dt^2 + \left(1 - \frac{2GM}{c^2 r}\right)^{-1} dr^2 + r^2 d\Omega^2\]

在 \(r = r_s = 2GM/c^2\)（Schwarzschild
半径）处，度规有坐标奇异性（事件视界）。

\textbf{定理 151.3}（Birkhoff 定理）：Schwarzschild 解是球对称真空
Einstein 方程的唯一解。

\textbf{例 151.2}（Kerr 解，1963）：旋转黑洞的轴对称真空解，由质量 \(M\)
和角动量 \(J = aM\) 参数化。Kerr 度规包含事件视界和能层（ergosphere）。

\textbf{例 151.3}（Friedmann-Lemaître-Robertson-Walker
解，FLRW）：均匀各向同性宇宙模型：

\[ds^2 = -c^2 dt^2 + a(t)^2 \left[ \frac{dr^2}{1 - k r^2} + r^2 d\Omega^2 \right]\]

其中 \(a(t)\) 是尺度因子，\(k = -1, 0, 1\)
对应开放、平坦和闭合宇宙。Friedmann 方程：

\[\left(\frac{\dot{a}}{a}\right)^2 = \frac{8\pi G}{3}\rho - \frac{k c^2}{a^2} + \frac{\Lambda c^2}{3}\]

\subsubsection{151.4
奇点定理与黑洞热力学}\label{ux5947ux70b9ux5b9aux7406ux4e0eux9ed1ux6d1eux70edux529bux5b66}

\textbf{定义
151.5}（捕获面与奇点）：\textbf{奇点}是测地线不完备性的表现（时空无法被延拓）。

\textbf{定理 151.4}（Penrose 奇点定理，1965）：若时空包含捕获面（trapped
surface），且满足适当的能量条件（如零能量条件
\(T_{\mu\nu}k^\mu k^\nu \geq 0\) 对所有类光向量
\(k\)），则存在不完备的测地线（奇点）。

\textbf{定理 151.5}（Hawking
奇点定理，1966）：在满足适当能量条件的宇宙学模型中，过去的测地不完备性不可避免（大爆炸奇点）。

\textbf{定义 151.6}（Hawking 辐射 / 黑洞热力学）：黑洞有熵
\(S = \frac{k_B c^3 A}{4G\hbar}\)（\(A\) 为视界面积）和温度
\(T = \frac{\hbar c^3}{8\pi G M k_B}\)（Schwarzschild
黑洞）。这暗示了广义相对论、量子力学和热力学的深层统一。

\textbf{定理 151.6}（黑洞热力学四定律，Bardeen-Carter-Hawking，1973）：
- 第零定律：稳态黑洞视界上的表面引力 \(\kappa\) 为常数（类比温度） -
第一定律：\(dM = \frac{\kappa}{8\pi} dA + \Omega dJ + \Phi dQ\)（能量守恒）
- 第二定律：黑洞视界面积永不减小 \(\delta A \geq 0\)（熵增） -
第三定律：不可能通过有限步骤使 \(\kappa\) 达到零

\subsubsection{151.5
引力波与渐近平坦时空}\label{ux5f15ux529bux6ce2ux4e0eux6e10ux8fd1ux5e73ux5766ux65f6ux7a7a}

\textbf{定义 151.7}（引力波）：Einstein 方程的线性化
\(g_{\mu\nu} = \eta_{\mu\nu} + h_{\mu\nu}\)（\(|h_{\mu\nu}| \ll 1\)）在
Lorentz 规范下给出波动方程
\(\square \bar{h}_{\mu\nu} = -\frac{16\pi G}{c^4} T_{\mu\nu}\)。真空中
\(\square \bar{h}_{\mu\nu} = 0\)，这是引力波的传播方程。

\textbf{定理 151.7}（Bondi-Sachs
能量动量）：对渐近平坦时空，在类光无穷远 \(\mathscr{I}^+\) 处可定义
Bondi 质量，其变化率等于引力波携带的能量通量。

\textbf{定义 151.8}（ADM 质量 /
Arnowitt-Deser-Misner）：对渐近平坦时空，在类空无穷远处定义的总能量（ADM
质量）为

\[M_{ADM} = \frac{1}{16\pi G} \lim_{r \to \infty} \int_{S^2_r} (\partial_j h_{ij} - \partial_i h_{jj}) n^i \, dA\]

\textbf{定理 151.8}（正质量定理，Schoen-Yau 1979, Witten
1981）：若时空满足主导能量条件，则 \(M_{ADM} \geq 0\)，且
\(M_{ADM} = 0\) 当且仅当时空是 Minkowski 时空。

\bigskip\hrule\bigskip

\subsection{第152章：共形场论}\label{ux7b2c152ux7ae0ux5171ux5f62ux573aux8bba}

共形场论（CFT）是二维量子场论中具有共形对称性的理论，是弦理论和统计物理的核心工具。其基本框架建立在Virasoro代数上：能量-动量张量的Laurent展开系数\(L_n\)满足对易关系\([L_m, L_n] = (m-n)L_{m+n} + \frac{c}{12}m(m^2-1)\delta_{m+n,0}\)，其中中心荷\(c\)刻画了理论的量子反常。主场（primary
field）\(\phi\)由共形权\((h, \bar{h})\)标记，满足\(L_0|\phi\rangle = h|\phi\rangle\)。算子乘积展开（OPE）描述了场在短距离下的奇异性结构，是CFT的核心计算工具。极小模型对中心荷\(c < 1\)进行了完整分类，其\(c = 1 - \frac{6}{m(m+1)}\)（\(m=3,4,\ldots\)），对应可解格点模型（如Ising模型\(c=1/2\)）的临界行为。CFT通过状态-算子对应建立了二维共形场论与弦世界面理论之间的桥梁。

\bigskip\hrule\bigskip

\subsection{第153章：可积系统与可解模型}\label{ux7b2c153ux7ae0ux53efux79efux7cfbux7edfux4e0eux53efux89e3ux6a21ux578b}

可积系统是具有''足够多''守恒量的动力系统，其关键在于Lax对形式化：将非线性演化方程表示为算子方程\(\frac{\partial L}{\partial t} = [M, L]\)，其中\(L\)和\(M\)是依赖于位势的线性算子。此形式化保证了\(\operatorname{tr}(L^n)\)的守恒性。反散射方法（IST）将非线性PDE的初值问题转化为线性散射问题，基本步骤为：(i)以初值构造散射数据；(ii)散射数据的时间演化是平凡的；(iii)通过Gel'fand-Levitan-Marchenko积分方程重构位势。经典例子包括KdV方程\(u_t + 6uu_x + u_{xxx} = 0\)（Lax对\(L = -\partial_x^2 - u, M = -4\partial_x^3 - 6u\partial_x - 3u_x\)）和Toda链。Yang-Baxter方程\(R_{12}R_{13}R_{23} = R_{23}R_{13}R_{12}\)是可积格点模型（如六顶角模型、XXZ自旋链）的核心，它保证了多体散射的因子化性质。

\bigskip\hrule\bigskip

\subsection{第154章：重正化群}\label{ux7b2c154ux7ae0ux91cdux6b63ux5316ux7fa4}

重正化群由Wilson在1970年代系统发展，是理解临界现象和量子场论的核心框架。其基本思想是动量壳层积分法：将动量空间划分为高能（\(\Lambda/s < |k| < \Lambda\)）和低能（\(|k| < \Lambda/s\)）两部分，积分掉高能自由度后得到低能有效理论，耦合常数随之''流动''。这一流动由beta函数\(\beta(g) = \mu \frac{dg}{d\mu}\)描述，其零点称为固定点（fixed
point）。固定点的稳定性决定了物理系统在临界点附近的标度行为：所有在同一固定点吸引域内的理论属于同一个普适类（universality
class），具有相同的临界指数。这一框架解释了为什么微观上迥异的系统（如液体-气体相变与铁磁相变）在临界点附近表现出相同的幂律行为。

\bigskip\hrule\bigskip

\subsection{第155章：等离子体物理中的数学方法（Mathematical Methods in
Plasma
Physics）}\label{ux7b2c155ux7ae0ux7b49ux79bbux5b50ux4f53ux7269ux7406ux4e2dux7684ux6570ux5b66ux65b9ux6cd5mathematical-methods-in-plasma-physics}

等离子体物理研究电离气体（等离子体）的集体行为，其数学模型横跨
PDE、动力系统、变分法和渐近分析。\textbf{核心方程}：（1）\textbf{Vlasov-Maxwell
方程组}------描述无碰撞等离子体中粒子分布函数 \(f(t, x, v)\)
和自洽电磁场 \((\mathbf{E}, \mathbf{B})\)
的耦合演化：\(\partial_t f + v \cdot \nabla_x f + \frac{q}{m}(\mathbf{E} + \frac{v}{c} \times \mathbf{B}) \cdot \nabla_v f = 0\)，这是无穷维
Hamilton 系统（由 Morrison-Marsden-Weinstein
括号刻画）；（2）\textbf{磁流体力学（MHD）}------将等离子体近似为导电连续介质，耦合
Navier-Stokes 方程和 Maxwell
方程，描述太阳风、聚变等离子体和天体物理喷流等宏观现象。理想 MHD 的
Lundquist 数与磁重联的 Sweet-Parker
模型刻画了磁场拓扑突变时能量释放的机制；（3）\textbf{Landau
阻尼}------Landau (1946) 对线性化 Vlasov-Poisson 方程的 Laplace
变换分析，发现电场扰动随时间代数衰减（非指数衰减）------相位混合导致的有效耗散，由
Mouhot-Villani (2011) 在半线性水平给予严格数学证明（获 Fields
奖工作）；（4）\textbf{磁约束聚变}------托卡马克装置中等离子体平衡由
Grad-Shafranov
方程描述：\(\Delta^* \psi = -r^2 \frac{dp}{d\psi} - \frac{1}{2} \frac{dF^2}{d\psi}\)（非线性椭圆
PDE），其解的存在性和分岔由变分方法保证。等离子体物理的数学挑战还包括：不稳定性分析（磁流体和动理学）、湍流输运理论和陀螺动理学（gyrokinetic）模型的降维。

\bigskip\hrule\bigskip

\emph{本卷从变分原理和 Lagrange/Hamilton
力学出发，建立了辛几何、量子力学数学基础、量子场论数学引论、广义相对论数学基础、共形场论、可积系统与可解模型、重正化群和等离子体物理中的数学方法。共
9 章。}

\newpage

\section{卷三十：计算理论}\label{ux5377ux4e09ux5341ux8ba1ux7b97ux7406ux8bba}

\begin{quote}
计算理论是计算机科学的数学基础，研究计算的本质、能力和极限。本卷在数理逻辑（V20
Ch 101
递归论）的基础上，建立自动机理论、形式语言、可计算性深化和计算复杂性分类（P/NP/PSPACE）。计算复杂性理论是当代数学中最重要的未解决问题之一（P
vs NP 问题）的所在地。
\end{quote}

\bigskip\hrule\bigskip

\subsection{第152章：自动机与形式语言}\label{ux7b2c152ux7ae0ux81eaux52a8ux673aux4e0eux5f62ux5f0fux8bedux8a00}

自动机理论是计算模型的最简形式，与形式语言（Chomsky
谱系）之间有精确的对应关系。

\subsubsection{152.1 有限自动机}\label{ux6709ux9650ux81eaux52a8ux673a}

\textbf{定义 152.1}（确定型有限自动机 /
DFA）：一个\textbf{确定型有限自动机}是五元组
\(M = (Q, \Sigma, \delta, q_0, F)\)，其中 \(Q\) 为有限状态集，\(\Sigma\)
为输入字母表，\(\delta: Q \times \Sigma \to Q\)
为转移函数，\(q_0 \in Q\) 为初始状态，\(F \subseteq Q\) 为接受状态集。

\(M\) 接受字符串 \(w = w_1 w_2 \cdots w_n\) 如果存在状态序列
\(r_0, r_1, \ldots, r_n\) 满足
\(r_0 = q_0\)，\(r_{i+1} = \delta(r_i, w_{i+1})\)，且 \(r_n \in F\)。

\textbf{定义 152.2}（正则语言）：语言 \(L \subseteq \Sigma^*\)
是\textbf{正则语言}，如果存在 DFA 使得 \(L = L(M)\)（\(M\)
接受的语言）。

\textbf{定义 152.3}（非确定型有限自动机 / NFA）：转移函数
\(\delta: Q \times \Sigma_\varepsilon \to \mathcal{P}(Q)\)（\(\Sigma_\varepsilon = \Sigma \cup \{\varepsilon\}\)）。NFA
接受字符串如果存在某条接受路径。

\textbf{定理 152.1}（DFA 与 NFA 的等价性 / Rabin-Scott，1959）：任何 NFA
可转化为等价的 DFA（子集构造法）。正则语言类在并、交、补、连接、Kleene
星号下封闭。

\textbf{定理 152.2}（正则语言的泵引理）：若 \(L\)
是正则语言，则存在泵长度 \(p\) 使得任何
\(w \in L\)（\(|w| \geq p\)）可分解为 \(w = xyz\)，满足
\(|xy| \leq p\)，\(|y| > 0\)，且 \(xy^iz \in L\) 对所有 \(i \geq 0\)。

\textbf{例 152.1}：\(L = \{0^n 1^n : n \geq 0\}\)
不是正则语言（由泵引理）。\(L = \{w : w\) 中 0 和 1 的个数相等\(\}\)
不是正则语言。

\textbf{定理 152.3}（Myhill-Nerode 定理）：语言 \(L\) 是正则的当且仅当
\(L\) 的右不变等价关系 \(\equiv_L\)
的指数（等价类个数）有限。这给出了正则语言的最小 DFA 构造方法。

\subsubsection{152.2
下推自动机与上下文无关语言}\label{ux4e0bux63a8ux81eaux52a8ux673aux4e0eux4e0aux4e0bux6587ux65e0ux5173ux8bedux8a00}

\textbf{定义 152.4}（上下文无关文法 / CFG）：文法
\(G = (V, \Sigma, R, S)\)，其中产生式规则形如
\(A \to \alpha\)（\(A \in V\)，\(\alpha \in (V \cup \Sigma)^*\)）。\(L(G)\)
是能从 \(S\) 推导出的终结符串集合。

\textbf{定义 152.5}（下推自动机 / PDA）：PDA
是有限自动机加上一个栈（stack）。转移函数
\(\delta: Q \times \Sigma_\varepsilon \times \Gamma_\varepsilon \to \mathcal{P}(Q \times \Gamma_\varepsilon)\)，其中
\(\Gamma\) 为栈字母表。

\textbf{定理 152.4}（PDA 与 CFG 的等价性）：语言 \(L\)
是上下文无关的当且仅当存在 PDA 接受 \(L\)。

\textbf{定理 152.5}（CFL 的泵引理）：若 \(L\) 是 CFL，则存在 \(p\)
使得任何 \(w \in L\)（\(|w| \geq p\)）可分解为 \(w = uvxyz\)，满足
\(|vxy| \leq p\)，\(|vy| > 0\)，且 \(uv^ixy^iz \in L\) 对所有
\(i \geq 0\)。

\textbf{例 152.2}：\(\{0^n 1^n 2^n : n \geq 0\}\) 不是
CFL。\(\{ww : w \in \{0,1\}^*\}\) 不是 CFL。

\textbf{定理 152.6}（Chomsky 范式）：任何不含 \(\varepsilon\) 的 CFG
可转化为等价的 Chomsky 范式格式（产生式仅形如 \(A \to BC\) 或
\(A \to a\)）。

\subsubsection{152.3 Turing 机与 Chomsky
谱系}\label{turing-ux673aux4e0e-chomsky-ux8c31ux7cfb}

\textbf{定义 152.6}（Turing 机 / TM）：Turing 机是七元组
\(M = (Q, \Sigma, \Gamma, \delta, q_0, q_{\text{accept}}, q_{\text{reject}})\)，其中
\(\Gamma\)
是带字母表，\(\delta: Q \times \Gamma \to Q \times \Gamma \times \{L, R\}\)（转移函数）。

\textbf{定理 152.7}（Chomsky 谱系）： - 类型
3（正则文法）\(\leftrightarrow\) DFA/NFA \(\leftrightarrow\) 正则语言 -
类型 2（上下文无关文法）\(\leftrightarrow\) PDA \(\leftrightarrow\) CFL
- 类型 1（上下文有关文法）\(\leftrightarrow\) 线性有界自动机（LBA） -
类型 0（无限制文法）\(\leftrightarrow\) Turing 机 \(\leftrightarrow\)
递归可枚举语言（r.e.）

\textbf{定理 152.8}（Turing 机的变体等价性）：多带 TM、非确定型
TM（NTM）、双向无限带 TM 在计算能力上与单带 TM 等价。

\bigskip\hrule\bigskip

\subsection{第153章：可计算性理论深化}\label{ux7b2c153ux7ae0ux53efux8ba1ux7b97ux6027ux7406ux8bbaux6df1ux5316}

本章在 V20 Ch 101 的基础上，深入探讨可计算性的核心定理和不可判定问题。

\subsubsection{153.1
可计算函数的精确定义}\label{ux53efux8ba1ux7b97ux51fdux6570ux7684ux7cbeux786eux5b9aux4e49}

\textbf{定义
153.1}（部分递归函数）：\textbf{部分递归函数}是从初始函数（零函数
\(Z(x)=0\)、后继函数 \(S(x)=x+1\)、投影函数
\(P_i^n(x_1,\ldots,x_n)=x_i\)）通过复合、原始递归和极小化（\(\mu\)-算子）得到的函数类。\textbf{全递归函数}是处处有定义的部分递归函数。

\textbf{定理 153.1}（Church-Turing
论题）：以下计算模型定义相同的函数类： - Turing 机可计算函数 -
\(\lambda\)-可定义函数（Church 的 \(\lambda\) 演算，1936） -
部分递归函数（Gödel-Kleene，1936） - 随机存取机（RAM）可计算函数 -
所有「合理」的计算模型

\textbf{定理 153.2}（Kleene 正规形定理）：存在原始递归谓词
\(T(e, x, y)\)（Kleene T-谓词）和原始递归函数 \(U(y)\)
使得任何部分递归函数可表为 \(f(x) = U(\mu y \, T(e, x, y))\)。

\textbf{定理 153.3}（\(s\)-\(m\)-\(n\) 定理 /
参数化定理）：存在原始递归函数 \(S_n^m\) 使得
\(\varphi_{S_n^m(e, x_1, \ldots, x_n)}^{(m)}(y_1, \ldots, y_m) = \varphi_e^{(m+n)}(x_1, \ldots, x_n, y_1, \ldots, y_m)\)。

\subsubsection{153.2 不可判定性}\label{ux4e0dux53efux5224ux5b9aux6027}

\textbf{定理 153.4}（停机问题 / Halting Problem，Turing 1936）：语言
\(HALT = \{\langle M, w \rangle : M\) 在输入 \(w\) 上停机\(\}\)
是递归可枚举的，但不是可判定的（递归的）。

\emph{证明}：假设存在判定器 \(H\)。构造 \(D\)：在输入
\(\langle M \rangle\) 时，\(D\) 以
\(\langle M, \langle M \rangle \rangle\) 运行 \(H\)，若 \(H\)
接受则拒绝，反之接受。则 \(D(\langle D \rangle)\) 接受当且仅当
\(D(\langle D \rangle)\) 不接受，矛盾。∎

\textbf{定理 153.5}（Rice
定理，1953）：递归可枚举语言的任何非平凡性质（即存在具有和不具有该性质的
r.e. 语言）都是不可判定的。即
\(\{\langle M \rangle : L(M) \in \mathcal{P}\}\) 不可判定，只要
\(\mathcal{P}\) 非平凡。

\emph{证明思路}：设 \(\mathcal{P}\) 是非平凡的 r.e. 语言性质，且
\(\emptyset \notin \mathcal{P}\)（否则考虑 \(\mathcal{P}\) 的补）。取
\(L_0 \in \mathcal{P}\) 并使 \(L_0\) 被 TM \(M_0\)
半判定。将停机问题归约：对任意 \(\langle M, w \rangle\)，构造 TM
\(N\)：在输入 \(x\) 上，\(N\) 同时模拟 \(M\) 在 \(w\) 上的运行和 \(M_0\)
在 \(x\) 上的运行。若 \(M\) 在 \(w\) 上停机，则 \(N\) 接受 \(x\)
当且仅当 \(M_0\) 接受 \(x\)；否则 \(N\) 不接受任何输入。于是
\(L(N) = L_0\)（若 \(M\) 在 \(w\)
上停机），\(L(N) = \emptyset\)（若不停机）。因此
\(\langle M, w \rangle \in \operatorname{HALT} \iff L(N) \in \mathcal{P}\)，停机问题可归约为判定
\(\mathcal{P}\)，故 \(\mathcal{P}\) 不可判定。∎

\textbf{推论 153.6}（不可判定问题的例子）：以下问题均不可判定： -
判定两个 CFG 是否生成相同的语言 - 判定一个 CFG 是否是歧义的 - 判定一个
TM 是否接受空语言 - Hilbert 第十问题：判定 Diophantine
方程是否有整数解（Matiyasevich，1970，基于 Davis-Putnam-Robinson
的工作）

\subsubsection{153.3 相对计算与 Turing
度}\label{ux76f8ux5bf9ux8ba1ux7b97ux4e0e-turing-ux5ea6}

\textbf{定义 153.2}（神谕 Turing 机 / Oracle TM）：带有神谕
\(A \subseteq \Sigma^*\) 的 Turing 机 \(M^A\)
可以查询「\(w \in A\)？」。相对可计算性：\(B \leq_T A\) 如果 \(B\) 在
\(M^A\) 下可判定。

\textbf{定义 153.3}（Turing 度）：\(\leq_T\) 的等价类构成 Turing
度。Turing 度构成上半格。\(\mathbf{0}\) 是可计算集的度，\(\mathbf{0}'\)
是停机问题的度。

\textbf{定理 153.7}（Post 问题，1944 / Friedberg-Muchnik
1956-1957）：存在递归可枚举的 Turing 度严格介于 \(\mathbf{0}\) 和
\(\mathbf{0}'\) 之间。这通过优先方法（priority
method）证明，是递归论的核心技术。

\textbf{定理 153.8}（Kleene-Post 定理）：存在不可比较的 Turing 度（即
\(\mathbf{a} \nleq_T \mathbf{b}\) 且
\(\mathbf{b} \nleq_T \mathbf{a}\)）。

\bigskip\hrule\bigskip

\subsection{第154章：计算复杂性：P、NP 与 NP
完全性}\label{ux7b2c154ux7ae0ux8ba1ux7b97ux590dux6742ux6027pnp-ux4e0e-np-ux5b8cux5168ux6027}

计算复杂性理论研究计算问题所需资源（时间、空间）与输入规模的关系，是理论计算机科学的核心。

\subsubsection{154.1
时间复杂度类}\label{ux65f6ux95f4ux590dux6742ux5ea6ux7c7b}

\textbf{定义 154.1}（时间复杂度类）： -
\(\mathbf{TIME}(t(n)) = \{L : \text{存在 } O(t(n)) \text{ 时间判定 } L \text{ 的 TM}\}\)
-
\(\mathbf{P} = \bigcup_{k \geq 1} \mathbf{TIME}(n^k)\)（多项式时间可判定语言）
-
\(\mathbf{NP} = \bigcup_{k \geq 1} \mathbf{NTIME}(n^k)\)（非确定型多项式时间可判定语言）
-
\(\mathbf{EXP} = \bigcup_{k \geq 1} \mathbf{TIME}(2^{n^k})\)（指数时间可判定语言）

\textbf{定义 154.2}（NP 的等价刻画）：\(L \in \mathbf{NP}\)
当且仅当存在多项式时间可验证的证书（certificate）。即存在多项式 \(p\)
和多项式时间 TM \(V\) 使得
\(x \in L \iff \exists y, |y| \leq p(|x|), V(x, y) = 1\)。

\textbf{例 154.1}（NP 问题的例子）： - \textbf{SAT}：布尔公式的可满足性
- \textbf{3SAT}：3-CNF 公式的可满足性 - \textbf{CLIQUE}：图是否包含
\(k\)-团 - \textbf{VERTEX COVER}：是否存在 \(k\) 个顶点的覆盖 -
\textbf{HAMILTONIAN CYCLE}：图是否有 Hamilton 回路 - \textbf{SUBSET
SUM}：是否存在子集和等于目标值 -
\textbf{整数规划}：\(A\mathbf{x} \leq \mathbf{b}\) 是否有整数解

\subsubsection{154.2 NP 完全性}\label{np-ux5b8cux5168ux6027}

\textbf{定义 154.3}（多项式时间归约）：\(A \leq_P B\)
如果存在多项式时间可计算函数 \(f\) 使得 \(x \in A \iff f(x) \in B\)。

\textbf{定义 154.4}（NP 完全与 NP 难）： - \(L\) 是
\(\mathbf{NP}\)-\textbf{难}的，如果对所有
\(A \in \mathbf{NP}\)，\(A \leq_P L\) - \(L\) 是
\(\mathbf{NP}\)-\textbf{完全}的，如果 \(L \in \mathbf{NP}\) 且 \(L\) 是
NP 难的

\textbf{定理 154.1}（Cook-Levin 定理，1971-1973）：SAT 是 NP 完全的。

\emph{证明思路}：给定任意 NP 语言 \(L\) 和对应的非确定型 TM
\(M\)，构造一个布尔公式 \(\phi\) 编码 \(M\) 在输入 \(x\)
上的计算过程。\(\phi\) 可满足当且仅当 \(M\)
有一条接受路径。公式的构造是局部的（每个时间步的状态转移），大小是多项式的。∎

\textbf{定理 154.2}（Karp 的 21 个 NP
完全问题，1972）：通过多项式时间归约，证明了 3SAT、CLIQUE、VERTEX
COVER、HAMILTONIAN CYCLE、SUBSET SUM、PARTITION、KNAPSACK 等 21 个问题的
NP 完全性。

\textbf{归约链}：SAT \(\leq_P\) 3SAT \(\leq_P\) CLIQUE \(\leq_P\) VERTEX
COVER \(\leq_P\) HAMILTONIAN CYCLE。

\subsubsection{154.3 P vs NP 问题}\label{p-vs-np-ux95eeux9898}

\textbf{问题 154.1}（P vs NP 问题 / Clay
千禧年问题）：\(\mathbf{P} \stackrel{?}{=} \mathbf{NP}\)。即：是否每个解可被快速验证的问题，其解也可被快速找到？

绝大多数理论计算机科学家相信
\(\mathbf{P} \neq \mathbf{NP}\)，即存在本质上难以求解但解易于验证的问题。若
\(\mathbf{P} = \mathbf{NP}\)，则公钥密码学将崩溃（因整数分解属
NP），数学证明自动化、蛋白质折叠等难题将在多项式时间内可解；若
\(\mathbf{P} \neq \mathbf{NP}\)，则意味着存在不可逾越的计算硬阈值，这也是当前密码学安全性的理论前提。该问题自
1971 年 Cook-Levin 提出以来，虽已有大量证据支持
\(\mathbf{P} \neq \mathbf{NP}\)，但严格证明至今缺失。

\textbf{定理 154.3}（相对化障碍 / Baker-Gill-Solovay，1975）：存在神谕
\(A, B\) 使得 \(\mathbf{P}^A = \mathbf{NP}^A\) 和
\(\mathbf{P}^B \neq \mathbf{NP}^B\)。因此，对角化论证不能解决 P vs NP
问题。

\textbf{定理 154.4}（自然证明障碍 /
Razborov-Rudich，1997）：若存在伪随机数生成器，则不存在「自然证明」证明
\(\mathbf{P} \neq \mathbf{NP}\)（即电路复杂度下界的高效构造性证明）。

\textbf{定理 154.5}（代数化障碍 /
Aaronson-Wigderson，2008）：代数化（代数神谕）方法也不能解决 P vs NP
问题。

\bigskip\hrule\bigskip

\subsection{第155章：复杂度类与空间复杂性}\label{ux7b2c155ux7ae0ux590dux6742ux5ea6ux7c7bux4e0eux7a7aux95f4ux590dux6742ux6027}

本章深入探讨空间复杂性、多项式谱系和概率与交互证明类。

\subsubsection{155.1 空间复杂性}\label{ux7a7aux95f4ux590dux6742ux6027}

\textbf{定义 155.1}（空间复杂度类）： - \(\mathbf{SPACE}(s(n))\)：使用
\(O(s(n))\) 空间的 TM 可判定语言 -
\(\mathbf{L} = \mathbf{SPACE}(\log n)\)（对数空间） -
\(\mathbf{NL} = \mathbf{NSPACE}(\log n)\)（非确定型对数空间） -
\(\mathbf{PSPACE} = \bigcup_k \mathbf{SPACE}(n^k)\)（多项式空间） -
\(\mathbf{NPSPACE} = \bigcup_k \mathbf{NSPACE}(n^k)\)

\textbf{定理 155.1}（Savitch 定理，1970）：对
\(s(n) \geq \log n\)，\(\mathbf{NSPACE}(s(n)) \subseteq \mathbf{SPACE}(s(n)^2)\)。特别地，\(\mathbf{PSPACE} = \mathbf{NPSPACE}\)。

\textbf{定理 155.2}（空间层次定理）：若 \(s_1(n) = o(s_2(n))\) 且
\(s_2\) 是空间可构造的，则
\(\mathbf{SPACE}(s_1(n)) \subsetneq \mathbf{SPACE}(s_2(n))\)。

\textbf{已知包含关系}：\(\mathbf{L} \subseteq \mathbf{NL} \subseteq \mathbf{P} \subseteq \mathbf{NP} \subseteq \mathbf{PSPACE} \subseteq \mathbf{EXP}\)。其中
\(\mathbf{NL} \subsetneq \mathbf{PSPACE}\)（空间层次定理），但
\(\mathbf{P} \stackrel{?}{=} \mathbf{PSPACE}\) 仍未知。

\textbf{定义 155.2}（对数空间归约与 NL
完全性）：\(\mathbf{PATH}\)（有向图的 \(s\)-\(t\) 连通性）是 NL
完全的。\(\mathbf{2SAT}\) 是 NL 完全的。

\textbf{定理 155.3}（Immerman-Szelepcsényi
定理，1987-1988）：\(\mathbf{NL} = \mathbf{coNL}\)。即非确定型对数空间在补运算下封闭。

\subsubsection{155.2 多项式谱系}\label{ux591aux9879ux5f0fux8c31ux7cfb}

\textbf{定义 155.3}（多项式谱系 / Polynomial Hierarchy）： -
\(\Sigma_0^P = \Pi_0^P = \mathbf{P}\) -
\(\Sigma_{k+1}^P = \mathbf{NP}^{\Sigma_k^P}\)（带 \(\Sigma_k^P\) 神谕的
NP） - \(\Pi_{k+1}^P = \mathbf{coNP}^{\Sigma_k^P}\) -
\(\mathbf{PH} = \bigcup_k \Sigma_k^P\)

\textbf{定理 155.4}（多项式谱系的坍塌性质）：若
\(\Sigma_k^P = \Pi_k^P\)，则 \(\mathbf{PH} = \Sigma_k^P\)（谱系坍塌到第
\(k\) 层）。特别地，若 \(\mathbf{P} = \mathbf{NP}\)，则
\(\mathbf{P} = \mathbf{PH}\)。

\textbf{定理 155.5}（Toda
定理，1991）：\(\mathbf{PH} \subseteq \mathbf{P}^{\#\mathbf{P}}\)。即
\(\#\mathbf{P}\)（计数类）包含了整个多项式谱系。

\textbf{定义 155.4}（\(\#\mathbf{P}\) 类）：\(\#\mathbf{P}\) 是 NP
问题的解的个数计数问题类。\(\#SAT\)（计算 SAT 的满足赋值数）是
\(\#\mathbf{P}\) 完全的。

\subsubsection{155.3
概率与交互证明类}\label{ux6982ux7387ux4e0eux4ea4ux4e92ux8bc1ux660eux7c7b}

\textbf{定义 155.5}（概率类）： - \(\mathbf{BPP}\)（Bounded-error
Probabilistic Polynomial time）：错误概率 \(< 1/3\)
的概率多项式时间可判定语言 -
\(\mathbf{RP}\)：单侧错误（是实例不被错判为否） -
\(\mathbf{ZPP} = \mathbf{RP} \cap \mathbf{coRP}\)（零错误概率多项式时间）

\textbf{定理 155.6}（概率放大）：\(\mathbf{BPP}\) 的错误概率可降至
\(2^{-n}\)，通过重复多数投票。

\textbf{已知关系}：\(\mathbf{P} \subseteq \mathbf{ZPP} \subseteq \mathbf{RP} \subseteq \mathbf{BPP} \subseteq \mathbf{PSPACE}\)。\(\mathbf{BPP} \stackrel{?}{\subseteq} \mathbf{P}\)
未知。

\textbf{定义 155.6}（交互证明系统 / IP）：验证者
\(V\)（概率多项式时间）与证明者 \(P\) 交换多项式轮消息。\(V\) 接受
\(x \in L\) 的概率 \(> 2/3\)，\(x \notin L\) 的概率 \(< 1/3\)。

\textbf{定理 155.7}（\(\mathbf{IP} = \mathbf{PSPACE}\) /
Shamir，1992）：交互证明可验证的恰好是 PSPACE 语言。

\textbf{定义 155.7}（\(\mathbf{PCP}\) 定理 / Probabilistically Checkable
Proofs，Arora-Lund-Motwani-Sudan-Szegedy 1992；Arora-Safra
1998）：\(\mathbf{NP} = \mathbf{PCP}[\log n, O(1)]\)。即任何 NP
语言有证明，只需检查常数个随机比特即可高概率验证。其深刻意义在于：它等价于''近似版本的
NP 困难性''------PCP 定理表明，区分可满足公式与每个赋值最多满足
\(1-\varepsilon\) 比例的公式（对某个常数 \(\varepsilon > 0\)）是 NP
难的。这催生了近似困难性理论：MAX-3SAT、MAX-CLIQUE、SET COVER
等经典优化问题在某个常数因子内近似仍是 NP 难的。PCP
定理的证明分两步：代数方法（将 NP
预言转化为多项式求和验证）和组合方法（通过低度检验和局部解码将验证者简化为常数次查询）。PCP
定理是 20 世纪理论计算机科学最重大的成果之一，获得了 2001 年 Gödel 奖。

\textbf{定理 155.8}（PCP 定理与近似困难性）：PCP 定理蕴含许多 NP
完全问题的近似算法的最优性。例如，CLIQUE 不可近似到
\(n^{1-\varepsilon}\)（除非 \(\mathbf{P} = \mathbf{NP}\)）。

\bigskip\hrule\bigskip

\subsection{第156章：电路复杂性与量子计算}\label{ux7b2c156ux7ae0ux7535ux8defux590dux6742ux6027ux4e0eux91cfux5b50ux8ba1ux7b97}

本章介绍计算复杂性的非均匀模型（电路）和量子计算的基本概念。

\subsubsection{156.1 电路复杂性}\label{ux7535ux8defux590dux6742ux6027}

\textbf{定义 156.1}（布尔电路）：布尔电路是节点为逻辑门（AND, OR,
NOT）的有向无环图。电路族 \(\{C_n\}_{n \in \mathbb{N}}\) 接受语言
\(L\)，如果 \(C_n\) 有 \(n\) 个输入，且
\(C_n(x) = 1 \iff x \in L\)（\(|x| = n\)）。

\textbf{定义 156.2}（电路复杂度类）： -
\(\mathbf{P}_{/\text{poly}}\)：存在多项式大小的电路族可判定的语言 -
\(\mathbf{NC}\)（Nick's Class）：深度为 \(O(\log^k n)\)
的多项式大小电路族 - \(\mathbf{AC}^0\)：深度为常数的无界扇入电路

\textbf{定理
156.1}：\(\mathbf{P} \subseteq \mathbf{P}_{/\text{poly}}\)。但
\(\mathbf{P}_{/\text{poly}}\) 包含不可判定语言（因为可为每个 \(n\)
硬编码答案）。

\textbf{定理 156.2}（Karp-Lipton 定理，1980）：若
\(\mathbf{NP} \subseteq \mathbf{P}_{/\text{poly}}\)，则
\(\mathbf{PH} = \Sigma_2^P\)（多项式谱系坍塌到第二层）。

\textbf{定理 156.3}（电路下界结果）： -
\(\mathbf{PARITY} \notin \mathbf{AC}^0\)（Furst-Saxe-Sipser 1981, Ajtai
1983） -
\(\mathbf{MOD}_p \notin \mathbf{AC}^0[\mathbf{MOD}_q]\)（\(p \neq q\)
素数，Razborov 1987, Smolensky 1987）

\subsubsection{156.2
量子计算引论}\label{ux91cfux5b50ux8ba1ux7b97ux5f15ux8bba}

量子计算利用量子力学叠加与纠缠原理处理信息，计算能力严格包含经典计算（\(\mathbf{P} \subseteq \mathbf{BQP}\)），并能在多项式时间内分解整数（Shor
算法），对经典密码学构成根本挑战。

\textbf{定义 156.3}（量子比特 / Qubit）：量子比特的状态是
\(\mathbb{C}^2\) 中的单位向量
\(\alpha|0\rangle + \beta|1\rangle\)（\(|\alpha|^2 + |\beta|^2 = 1\)）。\(n\)
个量子比特的状态在 \(\mathbb{C}^{2^n}\) 中。

\textbf{定义 156.4}（量子门）：量子计算由酉变换实现。基本量子门包括： -
Hadamard
门：\(H = \frac{1}{\sqrt{2}}\begin{pmatrix} 1 & 1 \\ 1 & -1 \end{pmatrix}\)
- 受控非门（CNOT）：\(|x, y\rangle \mapsto |x, x \oplus y\rangle\) -
\(\pi/8\) 门（T
门）：\(T = \begin{pmatrix} 1 & 0 \\ 0 & e^{i\pi/4} \end{pmatrix}\)

\textbf{定理 156.4}（通用量子门集）：\(\{H, CNOT, T\}\)
是一个通用量子门集，即任何酉变换可由此集近似。

\textbf{定义 156.5}（量子复杂度类）： -
\(\mathbf{BQP}\)：量子多项式时间可判定的语言（错误概率 \(< 1/3\)） -
\(\mathbf{QMA}\)：量子 NP 的类比（量子 Merlin-Arthur 证明）

\textbf{定理
156.5}（已知的包含关系）：\(\mathbf{P} \subseteq \mathbf{BPP} \subseteq \mathbf{BQP} \subseteq \mathbf{PSPACE}\)。\(\mathbf{BQP}\)
是否超出 \(\mathbf{P}\) 未知。

\textbf{定理 156.6}（Shor 算法，1994）：整数因子分解和离散对数可在
\(\mathbf{BQP}\) 中解决。复杂度
\(O(\log^3 N)\)（量子），而经典已知最好算法是亚指数时间。

\textbf{定理 156.7}（Grover 算法，1996）：无序数据库搜索可在
\(O(\sqrt{N})\) 量子查询中完成（经典需 \(O(N)\)），这是二次加速。

\textbf{定理 156.8}（量子查询下界 / BBBV
定理，1998）：对一般的无序搜索，\(O(\sqrt{N})\) 是最优的量子查询复杂度。

\bigskip\hrule\bigskip

\subsection{第157章：计算几何（Computational
Geometry）}\label{ux7b2c157ux7ae0ux8ba1ux7b97ux51e0ux4f55computational-geometry}

计算几何研究几何对象的算法设计与复杂性分析。\textbf{核心问题}：（1）\textbf{凸包}------给定
\(n\) 个平面点，Graham 扫描算法以 \(O(n \log n)\) 构造其凸包（下界也是
\(\Omega(n \log n)\)）；（2）\textbf{Voronoi 图与 Delaunay
三角剖分}------平面上 \(n\) 个站点的 Voronoi
图划分空间为最近站点区域，Fortune 算法以 \(O(n \log n)\)
时间构造；对偶的 Delaunay
三角剖分最大化最小角（避免瘦三角形），是有限元网格生成的标准工具；（3）\textbf{线段求交}------Bentley-Ottmann
扫描线算法检测 \(n\) 条线段的全部 \(k\) 个交点，\(O((n+k)\log n)\)
时间；（4）\textbf{点定位}------Kirpatrick 的层次化三角剖分支持
\(O(\log n)\)
时间的平面点定位。\textbf{数值鲁棒性}------浮点舍入误差导致的拓扑不一致性是计算几何的核心挑战，精确几何计算通过任意精度有理算术（如
LEDA, CGAL 库）解决。\textbf{高维计算几何}：\(d\) 维凸包的组合复杂性为
\(O(n^{\lfloor d/2 \rfloor})\)（上界定理），Voronoi
图的复杂度类似增长，产生碰撞维数灾难（curse of
dimensionality）的实际瓶颈。近似算法如
\(\epsilon\)-核集（\(\epsilon\)-kernel）和随机投影在机器学习（高维最近邻搜索）中提供折中方案。

\bigskip\hrule\bigskip

\subsection{第158章：计算机代数（Computer
Algebra）}\label{ux7b2c158ux7ae0ux8ba1ux7b97ux673aux4ee3ux6570computer-algebra}

计算机代数研究符号计算的理论和算法，区别于数值计算的关键在于保持精确表示（如整数的任意精度、多项式的符号形式）。\textbf{核心内容}：（1）\textbf{多项式
GCD 计算}------Euclid 算法推广到多变元多项式，通过模素数的 Hensel 提升和
Chinese Remainder Theorem 组合恢复整系数。小素数的模方法（modular
GCD、EZ-GCD 算法）避免中间表达式膨胀；（2）\textbf{Gröbner
基底}------Buchberger
算法（1965年）将多项式理想化归为标准表示的完备方法，核心为计算
\(S\)-多项式并逐步消去首项。应用涵盖代数方程组求解、理想的维数和 Hilbert
多项式计算、自动几何定理证明。Faugère 的 F4/F5 算法用线性代数加速大规模
Gröbner 基底计算；（3）\textbf{多项式因式分解}------Berlekamp
算法（有限域）、Zassenhaus 的 Hensel
提升方法（有理数域）；（4）\textbf{符号积分}------Risch
算法（1969年）基于 Liouville
定理给出了初等函数的积分能否由初等函数表达的完全判定算法，是符号计算理论的最高成就之一；（5）\textbf{柱形代数分解（CAD）}------Collins
(1975) 引入的实代数几何算法，将 \(\mathbb{R}^n\)
划分为半代数胞腔，在每个胞腔上多项式符号恒定------是实数域量词消去的基础算法。计算机代数系统（Maple,
Mathematica, Singular, SageMath）的内核均依赖这些算法。

\bigskip\hrule\bigskip

\subsection{第159章：机器学习理论基础}\label{ux7b2c159ux7ae0ux673aux5668ux5b66ux4e60ux7406ux8bbaux57faux7840}

机器学习的数学理论基础由统计学习理论（Vapnik-Chervonenkis）、优化理论（随机梯度下降）和信息论共同构成。

\subsubsection{159.1 PAC 学习框架}\label{pac-ux5b66ux4e60ux6846ux67b6}

\textbf{定义 159.1}（PAC 学习 / Probably Approximately Correct, Valiant,
1984）：对给定的假设类 \(\mathcal{H}\)、精度参数 \(\varepsilon > 0\)
和置信参数 \(\delta > 0\)，学习算法 \(\mathcal{A}\) 称为 \textbf{PAC
学习} \(\mathcal{H}\) 的，如果存在样本复杂度
\(m(\varepsilon, \delta)\)，使得对任意目标概念 \(c \in \mathcal{H}\)
和任意分布 \(\mathcal{D}\)，以至少 \(m(\varepsilon, \delta)\) 个 i.i.d.
样本训练的算法以概率至少 \(1-\delta\) 输出假设 \(h\) 满足
\(L_{\mathcal{D}}(h) \leq \varepsilon\)。

\textbf{定义 159.2}（VC 维数 / Vapnik-Chervonenkis Dimension）：假设类
\(\mathcal{H}\) 的 \textbf{VC 维数} \(d_{\text{VC}}(\mathcal{H})\)
是能被 \(\mathcal{H}\) 打散（shattered）的最大点数，即存在 \(d\)
个点使得 \(\mathcal{H}\) 能在其上实现所有 \(2^d\) 种标号组合。

\textbf{定理 159.1}（PAC 学习的基本定理）：若 \(\mathcal{H}\) 具有有限
VC 维数 \(d\)，则 \(\mathcal{H}\) 是 PAC 可学的，其样本复杂度为

\[m(\varepsilon, \delta) = O\left( \frac{d}{\varepsilon} \log \frac{1}{\varepsilon} + \frac{1}{\varepsilon} \log \frac{1}{\delta} \right)\]

反之，若 \(\mathcal{H}\) 的 VC 维数无限，则 \(\mathcal{H}\) 不是 PAC
可学的。

\subsubsection{159.2
经验风险最小化与泛化理论}\label{ux7ecfux9a8cux98ceux9669ux6700ux5c0fux5316ux4e0eux6cdbux5316ux7406ux8bba}

\textbf{定义 159.3}（经验风险最小化 / ERM）：给定损失函数 \(\ell\)
和假设类 \(\mathcal{F}\)，ERM 原则选择

\[\hat{f} = \arg\min_{f \in \mathcal{F}} \frac{1}{n} \sum_{i=1}^n \ell(f(x_i), y_i) \equiv \arg\min_{f \in \mathcal{F}} \widehat{L}_n(f)\]

\textbf{定理 159.2}（泛化误差界 / 一致大数律）：设 \(\mathcal{F}\) 的 VC
维数为 \(d\)。则以概率至少 \(1-\delta\)，

\[L(\hat{f}) - \inf_{f \in \mathcal{F}} L(f) \leq O\left( \sqrt{\frac{d}{n}} + \sqrt{\frac{\log(1/\delta)}{n}} \right)\]

若 \(\mathcal{F}\) 满足更精细的熵条件（如 Rademacher
复杂度或覆盖数），可得到数据依赖的紧界。

\textbf{定理 159.3}（结构风险最小化 / SRM）：对嵌套的假设类
\(\mathcal{F}_1 \subset \mathcal{F}_2 \subset \cdots\)，SRM 选择最小化
\(\widehat{L}_n(f) + \text{penalty}(k)\)
的类和假设，从而在偏差-方差之间自动权衡。这是 SVM
和正则化方法的理论基础。

\subsubsection{159.3
随机梯度下降的收敛理论}\label{ux968fux673aux68afux5ea6ux4e0bux964dux7684ux6536ux655bux7406ux8bba}

\textbf{算法 159.1}（随机梯度下降 / SGD）：

\[w_{t+1} = w_t - \eta_t \nabla f(w_t; \xi_t)\]

其中 \(\xi_t\) 是从训练集中随机抽取的（小批次）样本。

\textbf{定理 159.4}（SGD 的凸收敛性）：设目标函数
\(F(w) = \mathbb{E}[f(w; \xi)]\) 为 \(\mu\)-强凸且梯度
\(L\)-Lipschitz。选取步长 \(\eta_t = 1/(\mu t)\)，则

\[\mathbb{E}\left[ \|w_T - w^*\|^2 \right] = O\left( \frac{1}{T} \right)\]

对一般凸函数（非强凸），取 \(\eta_t = 1/\sqrt{t}\)，收敛速率为
\(O(1/\sqrt{T})\)（平均迭代意义下）。

\subsubsection{159.4
神经网络的逼近理论}\label{ux795eux7ecfux7f51ux7edcux7684ux903cux8fd1ux7406ux8bba}

\textbf{定理 159.5}（通用逼近定理 / Cybenko, 1989; Hornik,
1991）：带有一个隐藏层且具有 Sigmoid 激活函数的前馈神经网络在
\(C([0,1]^n)\) 中稠密------即对任意连续函数 \(f\) 和任意
\(\varepsilon > 0\)，存在有限大小的单隐层网络一致逼近 \(f\) 至
\(\varepsilon\) 精度。

\textbf{定理 159.6}（Barron 界, 1993）：若目标函数 \(f\) 的 Fourier
变换满足
\(C_f = \int_{\mathbb{R}^n} \|\omega\| |\hat{f}(\omega)| d\omega < \infty\)，则单隐藏层
Sigmoid 网络以 \(O(C_f^2 / \sqrt{N})\) 的速率逼近 \(f\)（\(N\)
为隐藏单元数），避免了一般非参数估计的维数灾难。

\textbf{深度网络的表示优势}：深度网络通过多层次函数组合以指数方式降低表示某些函数所需的参数数量。对于折叠函数（如
\(x \mapsto |x|\) 的重复组合），\(k\) 层网络仅需 \(O(k)\)
个单元，而浅层网络可能需要 \(O(2^k)\) 个单元才能达到同等逼近精度。

\bigskip\hrule\bigskip

\emph{本卷在
V20（数理逻辑）的基础上，建立了自动机理论、形式语言、可计算性深化、计算复杂性（P/NP/PSPACE）和量子计算引论，以及计算几何与计算机代数。共
7 章。}

\newpage

\section{卷三十一：算子代数 II --- von Neumann
代数}\label{ux5377ux4e09ux5341ux4e00ux7b97ux5b50ux4ee3ux6570-ii-von-neumann-ux4ee3ux6570}

\begin{quote}
von Neumann 代数（原称 W\emph{-代数）是作用在 Hilbert 空间上的
}-子代数，在弱算子拓扑下闭。由 Murray 和 von Neumann 在 1930-1940
年代创立，其因子分类理论是算子代数中最深刻的结果之一。本卷在
C*-代数（V13 Ch 76）的基础上，建立 von Neumann
代数的基本理论、因子分类、Tomita-Takesaki 模理论和子因子理论。
\end{quote}

\bigskip\hrule\bigskip

\subsection{第157章：von Neumann
代数基础}\label{ux7b2c157ux7ae0von-neumann-ux4ee3ux6570ux57faux7840}

von Neumann 代数是
C*-代数的一种特殊类型，其核心特征是在弱算子拓扑下的闭性。

\subsubsection{157.1
定义与基本性质}\label{ux5b9aux4e49ux4e0eux57faux672cux6027ux8d28-1}

\textbf{定义 157.1}（von Neumann 代数）：设 \(\mathcal{H}\) 是 Hilbert
空间。\(\mathcal{B}(\mathcal{H})\) 的一个 \emph{-子代数 \(\mathcal{M}\)
是 \textbf{von Neumann 代数}（或 W}-代数），如果
\(\mathcal{M} = \mathcal{M}''\)，其中
\(\mathcal{M}' = \{T \in \mathcal{B}(\mathcal{H}) : TS = ST \text{ 对所有 } S \in \mathcal{M}\}\)
是 \(\mathcal{M}\) 的\textbf{交换子}（commutant）。

\textbf{定理 157.1}（von Neumann 双交换子定理，1929）：设
\(\mathcal{M} \subseteq \mathcal{B}(\mathcal{H})\) 是包含恒等算子 \(I\)
的 *-子代数。则以下等价： 1. \(\mathcal{M} = \mathcal{M}''\)（即
\(\mathcal{M}\) 是 von Neumann 代数） 2. \(\mathcal{M}\)
在弱算子拓扑（WOT）下闭 3. \(\mathcal{M}\) 在强算子拓扑（SOT）下闭 4.
\(\mathcal{M}\) 在 \(\sigma\)-弱拓扑（ultraweak topology）下闭

\emph{证明概要}：关键是证明 \(\mathcal{M}''\) 是
WOT-闭包。一个方向：\(\mathcal{M} \subseteq \mathcal{M}''\) 且
\(\mathcal{M}''\) 自动是 WOT-闭的（由双交换子的性质），故若
\(\mathcal{M} = \mathcal{M}''\) 则 \(\mathcal{M}\) 在 WOT
下闭。另一个方向：设 \(\mathcal{M}\) 在 WOT 下闭。需证
\(\mathcal{M}'' \subseteq \mathcal{M}\)。对
\(T \in \mathcal{M}''\)，需证明 \(T\) 在 \(\mathcal{M}\) 的 WOT
闭包中。这需要 von Neumann 密性引理：设 \(\mathcal{H}\) 的任意
\(n\)-元组 \((x_1,\dots,x_n)\)，给定 \(\varepsilon>0\)，存在
\(S \in \mathcal{M}\) 使得 \(\oplus_i S x_i\) 与 \(\oplus_i T x_i\)
的范数差小于 \(\varepsilon\)。对单个向量，利用
\(\overline{\mathcal{M}x} \supseteq \mathcal{M}''x\)（由
\(T \in \mathcal{M}''\) 保证）。对一般情况，在
\(\mathcal{H}^{\oplus n}\) 上考虑对角表示
\(\mathcal{M} \otimes I_n\)。由此可得
\(T \in \overline{\mathcal{M}}^{\text{SOT}} = \overline{\mathcal{M}}^{\text{WOT}}\)。WOT
与 SOT 的闭包在凸 *-代数情形下相等（由 Hahn-Banach / 双重拓扑定理），故
\(\mathcal{M} = \mathcal{M}''\)。∎

\textbf{定义 157.2}（前对偶 / Predual）：von Neumann 代数
\(\mathcal{M}\) 的前对偶 \(\mathcal{M}_*\) 是 \(\mathcal{M}\) 上所有
\(\sigma\)-弱连续线性泛函的空间。\(\mathcal{M}\) 是 \(\mathcal{M}_*\) 的
Banach 空间对偶：\(\mathcal{M} = (\mathcal{M}_*)^*\)。

\textbf{定理 157.2}（Sakai 定理，1956）：抽象 W\emph{-代数等价于具体 von
Neumann 代数。即任意 C}-代数 \(\mathcal{M}\) 若作为 Banach
空间有前对偶，则同构于某个 Hilbert 空间上的 von Neumann 代数。

\textbf{例 157.1}（von Neumann 代数的基本例子）： -
\(\mathcal{B}(\mathcal{H})\) 本身（I 型因子） - 交换 von Neumann
代数：\(L^\infty(X, \mu)\) 作为 \(L^2(X, \mu)\) 上的乘法算子 - 群 von
Neumann 代数：\(\mathcal{L}(G) = \{\lambda_g : g \in G\}''\)（\(G\)
为离散群，\(\lambda_g\) 为左正则表示） - \(\ell^\infty(\mathbb{N})\) 是
\(\ell^2(\mathbb{N})\) 上的交换 von Neumann 代数

\subsubsection{157.2
投影、中心与因子}\label{ux6295ux5f71ux4e2dux5fc3ux4e0eux56e0ux5b50}

\textbf{定义 157.3}（投影）：von Neumann 代数 \(\mathcal{M}\)
中的\textbf{投影}是满足 \(p = p^* = p^2\)
的元素。投影的等价关系：\(p \sim q\) 如果存在 \(v \in \mathcal{M}\) 使得
\(v^*v = p\) 且 \(vv^* = q\)（部分等距）。

\textbf{定义 157.4}（中心）：von Neumann 代数 \(\mathcal{M}\)
的\textbf{中心}是
\(\mathcal{Z}(\mathcal{M}) = \mathcal{M} \cap \mathcal{M}'\)。

\textbf{定义 157.5}（因子）：von Neumann 代数 \(\mathcal{M}\)
是\textbf{因子}（factor），如果
\(\mathcal{Z}(\mathcal{M}) = \mathbb{C} I\)（即中心是平凡的）。

\textbf{定理 157.3}（中心分解）：任何 von Neumann 代数 \(\mathcal{M}\)
可表示为因子的直积分：

\[\mathcal{M} = \int_X^\oplus \mathcal{M}_x \, d\mu(x)\]

其中 \(\mathcal{M}_x\)
是因子（几乎处处）。\(\mathcal{Z}(\mathcal{M}) \cong L^\infty(X, \mu)\)。

\subsubsection{157.3
迹与维数函数}\label{ux8ff9ux4e0eux7ef4ux6570ux51fdux6570}

\textbf{定义 157.6}（迹 / Trace）：von Neumann 代数 \(\mathcal{M}\)
上的\textbf{迹}是函数 \(\tau: \mathcal{M}_+ \to [0, \infty]\) 满足： -
\(\tau(x+y) = \tau(x) + \tau(y)\) -
\(\tau(\lambda x) = \lambda \tau(x)\)（\(\lambda \geq 0\)） -
\(\tau(x^*x) = \tau(xx^*)\)（迹性质）

迹是\textbf{忠实的}（faithful），如果
\(\tau(x^*x) = 0 \implies x = 0\)；是\textbf{半有限的}（semifinite），如果对任意
\(x \in \mathcal{M}_+\)，存在 \(0 \neq y \leq x\) 使得
\(\tau(y) < \infty\)；是\textbf{正规的}（normal），如果对递增网
\(x_\alpha \uparrow x\)，有 \(\tau(x_\alpha) \uparrow \tau(x)\)。

\textbf{定理 157.4}（Murray-von Neumann 维数理论）：在因子
\(\mathcal{M}\) 上，投影的 Murray-von Neumann 等价类在迹 \(\tau\)
下的像构成值域
\(\mathcal{D}(\mathcal{M}) \subseteq [0, \infty]\)。迹在标量倍数下唯一。

\bigskip\hrule\bigskip

\subsection{第158章：因子分类与 Murray-von Neumann
理论}\label{ux7b2c158ux7ae0ux56e0ux5b50ux5206ux7c7bux4e0e-murray-von-neumann-ux7406ux8bba}

因子分类是 Murray 和 von Neumann 在 1930-1940
年代创立的，是算子代数理论的基石。

\subsubsection{158.1
因子的类型分类}\label{ux56e0ux5b50ux7684ux7c7bux578bux5206ux7c7b}

\textbf{定义 158.1}（因子的类型）：根据投影的维数函数 \(\mathcal{D}\)
的值域，因子分为：

\begin{itemize}
\tightlist
\item
  \textbf{I
  型}：\(\mathcal{D} = \{0, 1, 2, \ldots, n\}\)（\(\text{I}_n\)）或
  \(\{0, 1, 2, \ldots\}\)（\(\text{I}_\infty\)）

  \begin{itemize}
  \tightlist
  \item
    \(\text{I}_n\) 因子 \(\cong M_n(\mathbb{C})\)（矩阵代数）
  \item
    \(\text{I}_\infty\) 因子 \(\cong \mathcal{B}(\mathcal{H})\)
  \end{itemize}
\item
  \textbf{II
  型}：\(\mathcal{D} = [0, \infty]\)（没有最小投影，维数连续但存在有限迹）

  \begin{itemize}
  \tightlist
  \item
    \(\text{II}_1\)
    因子：\(\mathcal{D} = [0, 1]\)（迹归一化，\(\tau(I) = 1\)）
  \item
    \(\text{II}_\infty\)
    因子：\(\mathcal{D} = [0, \infty]\)（迹非有限，但可表示为
    \(\text{II}_1\) 因子与 \(\mathcal{B}(\mathcal{H})\) 的张量积）
  \end{itemize}
\item
  \textbf{III
  型}：\(\mathcal{D} = \{0, \infty\}\)（不存在非零有限投影，即不存在半有限正规迹）
\end{itemize}

\textbf{定理 158.1}（Murray-von Neumann
分类定理）：任何因子恰好属于上述类型之一。

\textbf{例 158.1}（各类因子的例子）： -
\(\text{I}_n\)：\(M_n(\mathbb{C})\) -
\(\text{I}_\infty\)：\(\mathcal{B}(\ell^2)\) - \(\text{II}_1\)：离散群
\(G\) 的群 von Neumann 代数 \(\mathcal{L}(G)\)（若 \(G\) 是 ICC
群，即无限共轭类群），如 \(\mathbb{F}_2\)（自由群）的群 von Neumann 代数
- \(\text{II}_\infty\)：\(\mathcal{L}(G) \otimes \mathcal{B}(\ell^2)\) -
\(\text{III}\)：Powers 因子
\(\mathcal{R}_\lambda\)（\(0 < \lambda < 1\)），Araki-Woods 因子

\subsubsection{158.2 超有限 II₁
因子}\label{ux8d85ux6709ux9650-iiux2081-ux56e0ux5b50}

\textbf{定义 158.2}（超有限因子）：因子 \(\mathcal{M}\)
是\textbf{超有限的}（hyperfinite），如果存在有限的递增子因子族
\(\mathcal{M}_1 \subseteq \mathcal{M}_2 \subseteq \cdots\)（\(\mathcal{M}_k \cong M_{n_k}(\mathbb{C})\)）使得
\(\mathcal{M} = \overline{\bigcup_k \mathcal{M}_k}^{\text{SOT}}\)。

\textbf{定理 158.2}（Murray-von Neumann
唯一性定理，1943）：存在唯一的超有限 \(\text{II}_1\) 因子（记作
\(\mathcal{R}\)）。\(\mathcal{R}\) 是唯一的具有以下性质的
\(\text{II}_1\) 因子：它是矩阵代数的归纳极限，且具有忠实的正规迹。

\textbf{定理 158.3}（Connes 定理，1976）：任何 \(\text{II}_1\)
因子是超有限的当且仅当它是可近似的（amenable）。即 \(\mathcal{M}\)
是超有限的 \(\iff\) 存在从
\(\mathcal{M} \otimes \mathcal{M}^{\text{op}}\) 到 \(\mathcal{M}\)
的完全正映射的近似。

\subsubsection{158.3 Connes 的 III
型因子分类}\label{connes-ux7684-iii-ux578bux56e0ux5b50ux5206ux7c7b}

\textbf{定理 158.4}（Connes 的 III 型因子分类，1973）：设
\(\mathcal{M}\) 是 III 型因子，\(\varphi\)
是忠实的正规半有限权（weight）。定义\textbf{模自同构群}
\(\sigma_t^\varphi\)（Tomita-Takesaki 理论）。\(\mathcal{M}\)
的\textbf{Connes 谱} \(S(\mathcal{M})\) 是 \(\sigma^\varphi\) 的 Arveson
谱（不依赖于 \(\varphi\) 的选择）。

\begin{itemize}
\tightlist
\item
  \(\mathcal{M}\) 是 \textbf{III\(_0\)}
  型：\(S(\mathcal{M}) = \{0, 1\}\)
\item
  \(\mathcal{M}\) 是 \textbf{III\(_\lambda\)}
  型（\(0 < \lambda < 1\)）：\(S(\mathcal{M}) = \{0\} \cup \lambda^{\mathbb{Z}}\)
\item
  \(\mathcal{M}\) 是 \textbf{III\(_1\)}
  型：\(S(\mathcal{M}) = [0, \infty)\)
\end{itemize}

\emph{证明概要}：\(S(\mathcal{M})\) 的定义为
\(\bigcap_{\psi} \text{sp}(\sigma^\psi)\)，其中 \(\psi\)
遍历所有忠实的正规半有限权（\(\sigma^\psi\) 为模自同构群，其在
\(\text{Out}(\mathcal{M})\) 中的谱）。\(S(\mathcal{M})\) 是
\(\mathbb{R}_{\geq 0}\)
的闭乘法子半群。通过分析迹的缺失------具体地，考虑
\(\mathcal{M} \rtimes_{\sigma^\varphi} \mathbb{R}\)（连续交叉积），其性质决定
\(S(\mathcal{M})\) 的结构：若交叉积是 II 型，则 \(\mathcal{M}\) 是
III\(_\lambda\)（\(\lambda < 1\)）；若交叉积是 III 型，则
\(\mathcal{M}\) 是 III\(_1\)。III\(_0\) 型对应
\(S(\mathcal{M}) = \{0,1\}\) 的病理情况。不变量 \(S(\mathcal{M})\)
非平凡的证明基于 Connes 的 2×2 矩阵技巧和 Arveson 谱理论。∎

\textbf{定理 158.5}（Connes 结构定理）：III\(_\lambda\)
因子（\(0 < \lambda < 1\)）可表示为 II\(_\infty\) 因子与周期为
\(\log \lambda\) 的自同构的交叉积。III\(_1\) 因子没有这种表示。

\textbf{定理 158.6}（Haagerup 定理，1987）：存在唯一的超有限 III\(_1\)
因子。

\bigskip\hrule\bigskip

\subsection{第159章：Tomita-Takesaki
模理论}\label{ux7b2c159ux7ae0tomita-takesaki-ux6a21ux7406ux8bba}

Tomita-Takesaki 理论是 von Neumann 代数理论中最重要的工具之一，它关联了
von Neumann 代数的模自同构群和态空间。

\subsubsection{159.1 Tomita-Takesaki
定理}\label{tomita-takesaki-ux5b9aux7406}

\textbf{定义 159.1}（Tomita 算子）：设 \(\mathcal{M}\) 是 von Neumann
代数，\(\Omega\) 是循环分离向量（即 \(\mathcal{M}\Omega\) 和
\(\mathcal{M}'\Omega\) 在 \(\mathcal{H}\) 中稠密）。定义 Tomita 算子
\(S\)：

\[S(x\Omega) = x^*\Omega, \quad x \in \mathcal{M}\]

\(S\) 是反线性闭算子。其极分解 \(S = J \Delta^{1/2}\)
给出\textbf{模算子} \(\Delta\)（正自伴算子）和\textbf{模共轭}
\(J\)（反线性等距，\(J^2 = I\)）。

\textbf{定理 159.1}（Tomita-Takesaki 定理，1967-1970）： 1.
\(J \mathcal{M} J = \mathcal{M}'\)（模共轭将 \(\mathcal{M}\)
映到其交换子） 2.
\(\Delta^{it} \mathcal{M} \Delta^{-it} = \mathcal{M}\)（模自同构群保持
\(\mathcal{M}\)） 3.
\(\sigma_t^\varphi(x) = \Delta^{it} x \Delta^{-it}\) 定义了
\(\mathcal{M}\) 上的模自同构群

\emph{证明概要}：核心步骤：（i）由 \(\mathcal{M}\Omega\)
稠密定义稠定反线性算子 \(S: x\Omega \mapsto x^*\Omega\)。\(S\)
是可闭的，其闭包也记作 \(S\)。（ii）\(S\) 的极分解 \(S = J\Delta^{1/2}\)
给出正自伴算子 \(\Delta = S^* S\) 和反线性等距
\(J\)（\(J^2 = I\)）。（iii）证明
\(J\mathcal{M}J = \mathcal{M}'\)：关键是通过 KMS 条件建立
\(\Delta^{it}\mathcal{M}\Delta^{-it} = \mathcal{M}\)，然后利用
\(J = S\Delta^{-1/2}\) 的关系推出交换子对应。（iv）模自同构群
\(\sigma_t^\varphi\) 定义为
\(\sigma_t^\varphi(x) = \Delta^{it}x\Delta^{-it}\)，它保持
\(\mathcal{M}\) 且满足 KMS 条件（Kubo-Martin-Schwinger）：存在在条形区域
\(\{0 \leq \operatorname{Im}z \leq 1\}\) 上解析的有界函数 \(F(z)\) 使
\(F(t) = \varphi(\sigma_t^\varphi(x)y)\)，\(F(t+i) = \varphi(y\sigma_t^\varphi(x))\)。KMS
条件唯一地确定了模自同构群。∎

\textbf{定义 159.2}（KMS 条件 / Kubo-Martin-Schwinger）：态 \(\varphi\)
关于 \(\sigma_t^\varphi\) 满足 KMS 条件：对任意
\(x, y \in \mathcal{M}\)，存在有界连续函数 \(F(z)\) 在条形区域
\(\{z : 0 \leq \operatorname{Im}(z) \leq 1\}\) 上解析，边界值为

\[F(t) = \varphi(\sigma_t^\varphi(x) y), \quad F(t+i) = \varphi(y \sigma_t^\varphi(x))\]

\subsubsection{\texorpdfstring{159.2 Connes 上同调与 III\(_1\)
的进一步分类}{159.2 Connes 上同调与 III\_1 的进一步分类}}\label{connes-ux4e0aux540cux8c03ux4e0e-iii_1-ux7684ux8fdbux4e00ux6b65ux5206ux7c7b}

\textbf{定义 159.3}（模自同构群的不变扰动）：两个忠实的正规态
\(\varphi, \psi\) 的模自同构群 \(\sigma_t^\varphi\) 和 \(\sigma_t^\psi\)
在
\(\operatorname{Out}(\mathcal{M}) = \operatorname{Aut}(\mathcal{M}) / \operatorname{Inn}(\mathcal{M})\)
中是上同调的。Connes 的 \(T\)-不变量 \(T(\mathcal{M})\) 是
\(\sigma_t^\varphi\) 在 \(\operatorname{Out}(\mathcal{M})\) 中的谱。

\textbf{定理 159.2}（Connes 的 III\(_1\) 子分类）：根据
\(T(\mathcal{M})\)，III\(_1\)
因子可进一步分类。\(T(\mathcal{M}) = \mathbb{R}\) 的因子称为 III\(_1\)
型（在狭义上）。

\bigskip\hrule\bigskip

\subsection{第160章：II₁ 因子与 Jones
指标理论}\label{ux7b2c160ux7ae0iiux2081-ux56e0ux5b50ux4e0e-jones-ux6307ux6807ux7406ux8bba}

Jones 指标理论是子因子理论的核心，由 Vaughan Jones 在 1983
年创立，并在纽结理论和共形场论中有深远影响。

\subsubsection{160.1 Jones 指标}\label{jones-ux6307ux6807}

\textbf{定义 160.1}（子因子 / Subfactor）：\(\text{II}_1\) 因子
\(\mathcal{N} \subseteq \mathcal{M}\) 是\textbf{子因子}，如果
\(\mathcal{N}\) 和 \(\mathcal{M}\) 都是 \(\text{II}_1\) 因子，且
\(\mathcal{N}\) 的迹是 \(\mathcal{M}\) 的迹的限制。

\textbf{定义 160.2}（Jones 指标）：\(\mathcal{N} \subseteq \mathcal{M}\)
的\textbf{Jones 指标} \([\mathcal{M} : \mathcal{N}]\) 是 \(\mathcal{M}\)
作为 \(\mathcal{N}\)-模的维数（在 Murray-von Neumann
意义下）。具体地，若 \(\mathcal{M}\) 是 \(\mathcal{N}\)
上的有限生成射影模，则
\([\mathcal{M} : \mathcal{N}] = \dim_{\mathcal{N}}(\mathcal{M})\)。

\textbf{定理 160.1}（Jones 指标定理，1983）：\(\text{II}_1\)
子因子的指标只能取以下值：

\[\{4 \cos^2(\pi/n) : n = 3, 4, 5, \ldots\} \cup [4, \infty]\]

即指标的小于 4
的部分是离散的：\(1, 2, (1+\sqrt{5})/2\text{ 的平方}, \ldots\)。

\emph{证明概要}：设 \([\mathcal{M} : \mathcal{N}] = \tau^{-1}\)。通过
Jones 基本构造，得到包含塔
\(\mathcal{N} \subset \mathcal{M} \subset \mathcal{M}_1 \subset \cdots\)
和投影列 \(e_1, e_2, \ldots\) 满足 Temperley-Lieb 关系
\(e_i e_{i\pm 1} e_i = \tau e_i\)。考虑代数
\(A_n = \langle 1, e_1, \ldots, e_n \rangle\)，其上由迹
\(\operatorname{tr}\) 赋予内积
\(\langle x, y \rangle = \operatorname{tr}(y^* x)\)。正定性要求 Gram
矩阵为正半定。计算 \(A_n\) 的 Gram 行列式 \(D_n(\tau)\) 满足递推
\(D_n = D_{n-1} - \tau D_{n-2}\)，解得
\(D_n(\tau) = \frac{\sin((n+1)\theta)}{\sin\theta}\)，其中
\(\tau^{-1} = 4\cos^2\theta\)。若对所有 \(n\) 有 \(D_n(\tau) > 0\)，则
\(\theta \in (0, \pi/(n+2))\)，即 \(\tau^{-1} \geq 4\)。若存在 \(n\) 使
\(D_n(\tau) = 0\) 而 \(D_{n-1}(\tau) > 0\)，则 \(\theta = \pi/k\)，得
\(\tau^{-1} = 4\cos^2(\pi/k)\)（\(k \geq 3\)）。这给出了指标 \(<4\)
的离散谱分类。∎

\textbf{定理 160.2}（Jones 基本构造）：从
\(\mathcal{N} \subseteq \mathcal{M}\) 出发，可构造 Jones 塔：

\[\mathcal{N} \subseteq \mathcal{M} \subseteq \mathcal{M}_1 \subseteq \mathcal{M}_2 \subseteq \cdots\]

其中
\(\mathcal{M}_1 = \langle \mathcal{M}, e_{\mathcal{N}} \rangle\)（\(e_{\mathcal{N}}\)
是 \(\mathcal{N}\)
上的条件期望投影），\([\mathcal{M}_k : \mathcal{M}_{k-1}] = [\mathcal{M} : \mathcal{N}]\)。投影
\(e_1, e_2, \ldots\) 满足 Temperley-Lieb 关系：

\[e_i e_{i\pm 1} e_i = \tau e_i, \quad e_i e_j = e_j e_i \text{（}|i-j| > 1\text{）}\]

其中 \(\tau = [\mathcal{M} : \mathcal{N}]^{-1}\)。

\subsubsection{160.2 Temperley-Lieb
代数与纽结不变量}\label{temperley-lieb-ux4ee3ux6570ux4e0eux7ebdux7ed3ux4e0dux53d8ux91cf}

\textbf{定义 160.3}（Temperley-Lieb 代数）：参数 \(\tau\) 的
Temperley-Lieb 代数 \(TL_n(\tau)\) 是由 \(1, e_1, \ldots, e_{n-1}\)
生成的代数，满足关系： - \(e_i^2 = e_i\) -
\(e_i e_{i\pm 1} e_i = \tau e_i\) -
\(e_i e_j = e_j e_i\)（\(|i-j| > 1\)）

\textbf{定理 160.3}（Jones 多项式，1984）：Jones 从 Temperley-Lieb
代数的迹结构出发，发现了新的纽结不变量------Jones 多项式
\(V_L(t)\)。这是自 Alexander
多项式（1923）以来首个新的纽结多项式不变量，满足辫子关系：

\[t^{-1} V_{L_+} - t V_{L_-} = (t^{1/2} - t^{-1/2}) V_{L_0}\]

\textbf{推论 160.1}：Jones
多项式可区分许多以前无法区分的纽结（如左手和右手三叶结），并解决了 Tait
猜想的若干特例。

\subsubsection{160.3
平面代数与标准不变量}\label{ux5e73ux9762ux4ee3ux6570ux4e0eux6807ux51c6ux4e0dux53d8ux91cf}

\textbf{定义 160.4}（Popa 的 \(\lambda\)-格）：子因子
\(\mathcal{N} \subseteq \mathcal{M}\) 的\textbf{标准不变量}（standard
invariant）是相对交换子格 \(\mathcal{N}' \cap \mathcal{M}_k\)（带有
Jones 投影和迹）。这是子因子的完全同构不变量（在强可近似的条件下，Popa
1994）。

\textbf{定义 160.5}（Jones
平面代数，1999）：子因子的标准不变量可组织为\textbf{平面代数}（planar
algebra）------带有在平面区域中嵌入操作的代数结构。平面代数提供了子因子理论的图形化语言。

\bigskip\hrule\bigskip

\subsection{第161章：子因子理论在共形场论中的应用}\label{ux7b2c161ux7ae0ux5b50ux56e0ux5b50ux7406ux8bbaux5728ux5171ux5f62ux573aux8bbaux4e2dux7684ux5e94ux7528}

子因子理论与低维拓扑和共形场论有深刻联系，形成了数学物理中的重要交叉领域。

\subsubsection{161.1 共形场论与 von Neumann
代数}\label{ux5171ux5f62ux573aux8bbaux4e0e-von-neumann-ux4ee3ux6570}

\textbf{定义 161.1}（局部共形网 / Local Conformal Net）：圆 \(S^1\)
上的开区间族 \(\mathcal{I}\) 到 von Neumann 代数 \(\mathcal{A}(I)\)
的映射，满足： -
同调性：\(I_1 \subseteq I_2 \implies \mathcal{A}(I_1) \subseteq \mathcal{A}(I_2)\)
-
局域性：\(I_1 \cap I_2 = \emptyset \implies \mathcal{A}(I_1) \subseteq \mathcal{A}(I_2)'\)
- 共形协变性：\(\operatorname{Diff}(S^1)\) 的射影酉表示 -
真空态的存在性等

\textbf{定理 161.1}（Reeh-Schlieder 定理，局部网版本）：真空向量
\(\Omega\) 对每个 \(\mathcal{A}(I)\) 是循环和分离的。

\textbf{定理 161.2}（Bisognano-Wichmann 定理）：对 wedge 区域
\(W\)，模算子 \(\Delta_W\) 关联于 Lorentz boost，模共轭 \(J_W\) 关联于
PCT 变换。

\subsubsection{161.2
子因子与共形网}\label{ux5b50ux56e0ux5b50ux4e0eux5171ux5f62ux7f51}

\textbf{定理 161.3}（Longo 定理，1989）：对共形网
\(\mathcal{A}\)，包含关系
\(\mathcal{A}(I_1) \subseteq \mathcal{A}(I_2)\)（\(I_1 \subset I_2\)
为不相交区间的并）的 Jones 指标是有限的，且可由共形维度 \(c\) 通过 Kac
公式确定。

\textbf{定义 161.2}（\(\mu\)-指标）：整个共形网的 \(\mu\)-指标定义为
\(\mu_{\mathcal{A}} = \sum_{i} d_i^2\)（\(d_i\) 为各不可约 sector
的量子维数）。\(\mu\)-指标与 Jones 指标有关。

\textbf{定理 161.4}（Kawahigashi-Longo-Müger，2001）：某些共形网（如
Virasoro 网、loop group 网）的表示范畴构成模张量范畴（modular tensor
category），分类了有理共形场论。

\subsubsection{161.3
自由概率论与随机矩阵}\label{ux81eaux7531ux6982ux7387ux8bbaux4e0eux968fux673aux77e9ux9635}

\textbf{定义 161.3}（自由概率论 / Free Probability Theory，Voiculescu
1985）：自由概率论是 von Neumann
代数中非交换随机变量的理论，是经典概率论的「自由」类比，通过自由独立性替换经典独立性。

\textbf{定义 161.4}（自由独立性）：von Neumann 代数 \(\mathcal{M}\) 的
*-子代数 \(\mathcal{A}_1, \ldots, \mathcal{A}_n\)
是\textbf{自由独立的}，如果对任意
\(a_j \in \mathcal{A}_{i_j}\)（\(i_1 \neq i_2 \neq \cdots \neq i_k\)）且
\(\tau(a_j) = 0\)，有 \(\tau(a_1 a_2 \cdots a_k) = 0\)。

\textbf{定理
161.5}（自由中心极限定理）：自由独立的同分布自伴随机变量的标准化和依分布收敛于半圆律（Wigner
半圆分布）\(d\mu = \frac{1}{2\pi}\sqrt{4-t^2} \, dt\)。

\textbf{定理 161.6}（Voiculescu 的渐近自由定理）：独立的高斯随机矩阵族在
\(N \to \infty\) 时渐近自由独立。这为 von Neumann
代数的自由群因子问题提供了关键工具。

\bigskip\hrule\bigskip

\emph{本卷在 V13 Ch 76（C}-代数）的基础上，建立了 von Neumann
代数、因子分类、Tomita-Takesaki 理论、Jones
指标和子因子在共形场论中的应用。共 5 章。*

\newpage

\section{卷三十二：随机分析}\label{ux5377ux4e09ux5341ux4e8cux968fux673aux5206ux6790}

\begin{quote}
随机分析是随机过程与微积分的结合，将 Newton-Leibniz
微积分推广到随机过程（特别是 Brown 运动）上。本卷在概率论 II（V20）的
Itô 积分简介基础上，深入建立 Itô
随机微积分、随机微分方程（SDE）、Malliavin
分析、随机偏微分方程和随机几何。随机分析是金融数学（V54）和随机偏微分方程的核心工具。
\end{quote}

\bigskip\hrule\bigskip

\subsection{第162章：Itô
积分深化}\label{ux7b2c162ux7ae0ituxf4-ux79efux5206ux6df1ux5316}

Itô 积分是随机分析的核心，将随机过程关于 Brown 运动的积分严格化。

\subsubsection{162.1 Brown
运动回顾}\label{brown-ux8fd0ux52a8ux56deux987e}

\textbf{定义 162.1}（Brown 运动 / Wiener 过程）：实值过程
\(W_t\)（\(t \geq 0\)）是标准的\textbf{Brown 运动}（或 Wiener
过程），如果： 1. \(W_0 = 0\) 几乎必然 2. 对
\(0 \leq t_1 < t_2 < \cdots < t_n\)，增量
\(W_{t_2} - W_{t_1}, \ldots, W_{t_n} - W_{t_{n-1}}\) 独立 3.
\(W_t - W_s \sim \mathcal{N}(0, t-s)\)（\(t > s\)） 4. 轨道
\(t \mapsto W_t\) 连续

\textbf{命题 162.1}（Brown 运动的性质）： -
尺度不变性：\(\{c^{-1/2}W_{ct}\}\) 也是 Brown 运动 -
时间逆：\(\{t W_{1/t}\}\) 也是 Brown 运动（\(t > 0\)） - Brown
运动的轨道几乎必然不可微 -
二次变差：\(\sum_{i} (W_{t_{i+1}} - W_{t_i})^2 \to t\)（依概率，当划分精细度趋于
0）

\subsubsection{162.2 Itô
积分的构造}\label{ituxf4-ux79efux5206ux7684ux6784ux9020}

\textbf{定义 162.2}（简单过程的 Itô 积分）：设 \(\phi(t, \omega)\)
是\textbf{简单过程}（分段常数，且 \(\phi(t_i)\) 是
\(\mathcal{F}_{t_i}\)-可测的）：

\[\phi(t) = \sum_{i=0}^{n-1} \phi(t_i) \mathbf{1}_{[t_i, t_{i+1})}(t)\]

定义 Itô 积分

\[\int_0^T \phi(t) \, dW_t = \sum_{i=0}^{n-1} \phi(t_i) (W_{t_{i+1}} - W_{t_i})\]

\textbf{定理 162.1}（Itô 等距公式）：对简单过程 \(\phi\)，

\[\mathbb{E}\left[ \left( \int_0^T \phi(t) \, dW_t \right)^2 \right] = \mathbb{E}\left[ \int_0^T \phi(t)^2 \, dt \right]\]

\emph{证明}：展开平方，交叉项期望为零（由 \(\phi(t_i)\) 的
\(\mathcal{F}_{t_i}\)-可测性和 Brown 运动增量的独立性）。∎

\textbf{定义 162.3}（Itô 积分的推广）：对任意适应过程
\(\phi \in L^2([0, T] \times \Omega)\)（即
\(\mathbb{E}[\int_0^T \phi(t)^2 dt] < \infty\)），用简单过程逼近，通过
Itô 等距延拓定义积分。

\textbf{定理 162.2}（Itô 积分的性质）： 1.
线性：\(\int (a\phi + b\psi) dW = a\int \phi dW + b\int \psi dW\) 2.
鞅性质：\(M_t = \int_0^t \phi(s) dW_s\) 是连续鞅（关于 Brown
运动的自然濾过） 3.
二次变差：\(\langle M \rangle_t = \int_0^t \phi(s)^2 ds\)

\subsubsection{162.3 Itô 公式}\label{ituxf4-ux516cux5f0f}

\textbf{定理 162.3}（Itô 公式 / Itô 引理，一维）：设
\(f \in C^2(\mathbb{R})\)，\(X_t\) 是 Itô 过程：

\[dX_t = \mu_t \, dt + \sigma_t \, dW_t\]

则 \(Y_t = f(X_t)\) 满足

\[dY_t = f'(X_t) dX_t + \frac{1}{2} f''(X_t) \sigma_t^2 \, dt\]

展开为：

\[df(X_t) = \left( f'(X_t)\mu_t + \frac{1}{2} f''(X_t)\sigma_t^2 \right) dt + f'(X_t)\sigma_t \, dW_t\]

\textbf{定理 162.4}（多维 Itô 公式）：设
\(f(t, \mathbf{x}) \in C^{1,2}\)，\(\mathbf{X}_t\) 为 \(n\) 维 Itô 过程
\(d\mathbf{X}_t = \boldsymbol{\mu}_t dt + \Sigma_t d\mathbf{W}_t\)（\(\mathbf{W}_t\)
为 \(m\) 维 Brown 运动）。则

\[df(t, \mathbf{X}_t) = \left( \frac{\partial f}{\partial t} + \nabla f \cdot \boldsymbol{\mu}_t + \frac{1}{2} \operatorname{Tr}(\Sigma_t \Sigma_t^T \nabla^2 f) \right) dt + \nabla f \cdot \Sigma_t \, d\mathbf{W}_t\]

\emph{证明概要}（Itô公式，定理162.3-162.4）：核心在于Brown运动二次变差满足
\((dW_t)^2 = dt\)（均方意义）。对
\(f \in C^2\)，Taylor展开：\(f(X_t + dX_t) - f(X_t) \approx f'(X_t) dX_t + \frac{1}{2} f''(X_t) (dX_t)^2\)。代入
\(dX_t = \mu_t dt + \sigma_t dW_t\)，计算
\((dX_t)^2 = \mu_t^2 (dt)^2 + 2\mu_t\sigma_t dt dW_t + \sigma_t^2 (dW_t)^2\)。由于高阶项
\((dt)^2, dt\,dW_t\) 为 \(o(dt)\)，仅
\(\sigma_t^2 (dW_t)^2 \approx \sigma_t^2 dt\)
保留，即得一维公式。多维情形同理，交叉项
\(\Sigma_{ik}\Sigma_{jk} dW_t^i dW_t^j\) 利用
\(dW_t^i dW_t^j = \delta_{ij}dt\)
化为迹项。这是随机微积分与经典微积分的关键区别。∎

\textbf{例 162.1}（几何 Brown
运动）：\(dS_t = \mu S_t dt + \sigma S_t dW_t\)。解：\(S_t = S_0 \exp\left((\mu - \frac{\sigma^2}{2})t + \sigma W_t\right)\)。由
Itô
公式验证：\(f(x) = \ln x\)，\(df = \frac{1}{S}(\mu S dt + \sigma S dW) - \frac{1}{2S^2}\sigma^2 S^2 dt = (\mu - \sigma^2/2)dt + \sigma dW\)。

\subsubsection{162.4 Itô 过程的鞅表示与 Girsanov
定理}\label{ituxf4-ux8fc7ux7a0bux7684ux9785ux8868ux793aux4e0e-girsanov-ux5b9aux7406}

\textbf{定理 162.5}（Itô 鞅表示定理）：任何关于 Brown
运动濾过的平方可积鞅 \(M_t\) 可表示为

\[M_t = M_0 + \int_0^t \phi(s) \, dW_s\]

对某个适应过程 \(\phi\)。即 Brown 运动濾过去上的鞅都具有随机积分表示。

\textbf{定理 162.6}（Girsanov 定理）：设 \(\theta_t\) 是适应过程满足
Novikov 条件
\(\mathbb{E}[\exp(\frac{1}{2}\int_0^T \theta_s^2 ds)] < \infty\)。定义
\(Z_t = \exp(\int_0^t \theta_s dW_s - \frac{1}{2}\int_0^t \theta_s^2 ds)\)
和测度 \(d\mathbb{Q} = Z_T d\mathbb{P}\)。则在 \(\mathbb{Q}\) 下，过程

\[\tilde{W}_t = W_t - \int_0^t \theta_s \, ds\]

是 Brown 运动。

\emph{证明概要}（Girsanov定理）：定义指数鞅
\(Z_t = \exp(\int_0^t \theta_s dW_s - \frac12 \int_0^t \theta_s^2 ds)\)，由Itô公式验证
\(dZ_t = Z_t\theta_t dW_t\)，故 \(Z_t\)
是P-鞅（Novikov条件保证）。定义新测度 \(\mathbb{Q}\)
的Radon-Nikodym导数为
\(d\mathbb{Q}/d\mathbb{P}|_{\mathcal{F}_t} = Z_t\)。对任意有界停时
\(\tau\)，需验证 \(\tilde{W}_t = W_t - \int_0^t \theta_s ds\) 在
\(\mathbb{Q}\) 下是Brown运动：由Girsanov定理的抽象版本，\(\tilde{W}_t\)
的二次变差不变，且
\(\mathbb{E}^{\mathbb{Q}}[\tilde{W}_t|\mathcal{F}_s] = \mathbb{E}^{\mathbb{P}}[Z_t\tilde{W}_t|\mathcal{F}_s]/Z_s = \tilde{W}_s\)，即得鞅性。∎

\bigskip\hrule\bigskip

\subsection{第163章：随机微分方程}\label{ux7b2c163ux7ae0ux968fux673aux5faeux5206ux65b9ux7a0b}

随机微分方程（SDE）是添加了随机噪声的微分方程，是随机分析的核心应用。

\subsubsection{163.1 SDE
的基本理论}\label{sde-ux7684ux57faux672cux7406ux8bba}

\textbf{定义 163.1}（随机微分方程 / SDE）：

\[dX_t = b(t, X_t) \, dt + \sigma(t, X_t) \, dW_t, \quad X_0 = x_0\]

其中 \(b: \mathbb{R}_{\geq 0} \times \mathbb{R}^n \to \mathbb{R}^n\)
是\textbf{漂移系数}（drift），\(\sigma: \mathbb{R}_{\geq 0} \times \mathbb{R}^n \to \mathbb{R}^{n \times m}\)
是\textbf{扩散系数}（diffusion）。积分形式为

\[X_t = x_0 + \int_0^t b(s, X_s) \, ds + \int_0^t \sigma(s, X_s) \, dW_s\]

\textbf{定理 163.1}（SDE 解的存在唯一性）：若 \(b\) 和 \(\sigma\) 满足
Lipschitz
条件（\(|b(t,x)-b(t,y)| + |\sigma(t,x)-\sigma(t,y)| \leq K|x-y|\)）和线性增长条件（\(|b(t,x)| + |\sigma(t,x)| \leq K(1+|x|)\)），则
SDE 存在唯一的强解 \(X_t\)（适应于 Brown 运动的濾过）。

\emph{证明概要}（定理163.1，SDE解的存在唯一性）：采用Picard迭代。定义
\(X_t^{(0)} \equiv x_0\)，\(X_t^{(n+1)} = x_0 + \int_0^t b(s, X_s^{(n)}) ds + \int_0^t \sigma(s, X_s^{(n)}) dW_s\)。由Lipschitz条件，
\[\mathbb{E}|X_t^{(n+1)} - X_t^{(n)}|^2 \leq C \int_0^t \mathbb{E}|X_s^{(n)} - X_s^{(n-1)}|^2 ds\]
（使用Itô等距处理扩散项）。递推得
\(\mathbb{E}|X_t^{(n+1)} - X_t^{(n)}|^2 \leq (CT)^n/n!\)，故迭代序列在
\(L^2\)
中为Cauchy列。利用线性增长条件控制整体范数，极限即为强解。唯一性由Gronwall不等式导出：若
\(X_t, Y_t\) 均为解，则
\(\mathbb{E}|X_t - Y_t|^2 \leq K\int_0^t \mathbb{E}|X_s - Y_s|^2 ds\)，蕴涵
\(\mathbb{E}|X_t - Y_t|^2 = 0\)。∎

\textbf{定理 163.2}（Markov 性质）：SDE 的解 \(X_t\) 是 Markov
过程。其转移半群 \(P_t f(x) = \mathbb{E}^x[f(X_t)]\) 的生成元为

\[\mathcal{L} f(x) = \sum_i b_i(x) \frac{\partial f}{\partial x_i} + \frac{1}{2} \sum_{i,j} (\sigma\sigma^T)_{ij}(x) \frac{\partial^2 f}{\partial x_i \partial x_j}\]

\textbf{定理 163.3}（Feynman-Kac 公式）：函数
\(u(t, x) = \mathbb{E}^{t,x}[ \Phi(X_T) e^{-\int_t^T V(X_s) ds} ]\)
满足抛物型偏微分方程

\[\frac{\partial u}{\partial t} + \mathcal{L} u - V u = 0, \quad u(T, x) = \Phi(x)\]

\emph{证明概要}：定义
\(Y_s = e^{-\int_t^s V(X_r) dr} u(s, X_s)\)。由Itô公式，
\[dY_s = e^{-\int_t^s V dr}\left[ \frac{\partial u}{\partial s} + \mathcal{L}u - Vu \right] ds + \text{鞅项}\]
若括号内为零（即 \(u\) 满足抛物型PDE），则 \(Y_s\) 为鞅，故
\(u(t,x) = Y_t = \mathbb{E}[Y_T] = \mathbb{E}[\Phi(X_T)e^{-\int_t^T V dr}]\)。反之，若定义
\(u\) 为条件期望，利用 Markov 性和生成元 \(\mathcal{L}\) 可得
PDE。Feynman-Kac
公式沟通了随机过程（SDE的解）与抛物型PDE，是随机分析与偏微分方程之间的桥梁。∎

\subsubsection{163.2 线性 SDE 与 Ornstein-Uhlenbeck
过程}\label{ux7ebfux6027-sde-ux4e0e-ornstein-uhlenbeck-ux8fc7ux7a0b}

\textbf{例 163.1}（Ornstein-Uhlenbeck
过程）：\(dX_t = -\theta X_t dt + \sigma dW_t\)。解为

\[X_t = e^{-\theta t} X_0 + \sigma \int_0^t e^{-\theta(t-s)} dW_s\]

\(X_t\) 是 Gauss 过程，均值 \(e^{-\theta t}X_0\)，方差
\(\frac{\sigma^2}{2\theta}(1 - e^{-2\theta t})\)。当 \(t \to \infty\)
时，分布收敛到 \(\mathcal{N}(0, \sigma^2/(2\theta))\)。OU
过程是遍历的且具有均值回复性。

\textbf{例 163.2}（几何 Brown
运动）：\(dS_t = \mu S_t dt + \sigma S_t dW_t\)。由 Itô
公式，\(S_t = S_0 \exp((\mu - \sigma^2/2)t + \sigma W_t)\)。\(\mathbb{E}[S_t] = S_0 e^{\mu t}\)。

\subsubsection{163.3 SDE
的数值模拟}\label{sde-ux7684ux6570ux503cux6a21ux62df}

\textbf{定义 163.2}（Euler-Maruyama 方法）：给定时间步长
\(\Delta t\)，近似

\[X_{t_{n+1}} \approx X_{t_n} + b(t_n, X_{t_n}) \Delta t + \sigma(t_n, X_{t_n}) \Delta W_n\]

其中 \(\Delta W_n \sim \mathcal{N}(0, \Delta t)\)。

\textbf{定理 163.4}（Euler-Maruyama 的收敛阶）：强收敛阶为 \(1/2\)（即
\(\mathbb{E}|X_T - X_T^{\Delta t}| \leq C \sqrt{\Delta t}\)），弱收敛阶为
\(1\)（即
\(|\mathbb{E}[f(X_T)] - \mathbb{E}[f(X_T^{\Delta t})]| \leq C \Delta t\)）。

\textbf{定义 163.3}（Milstein 方法）：增加 Itô 公式的校正项，强收敛阶为
\(1\)：

\[X_{t_{n+1}} \approx X_{t_n} + b \Delta t + \sigma \Delta W_n + \frac{1}{2} \sigma \sigma' ((\Delta W_n)^2 - \Delta t)\]

\bigskip\hrule\bigskip

\subsection{第164章：Malliavin
分析}\label{ux7b2c164ux7ae0malliavin-ux5206ux6790}

Malliavin 分析（或随机变分法）是 Wiener 空间上的无穷维微分学，由 Paul
Malliavin 在 1970 年代创立，最初用于证明 Hörmander
条件下的光滑密度存在性。

\subsubsection{164.1 Wiener 空间与 Malliavin
导数}\label{wiener-ux7a7aux95f4ux4e0e-malliavin-ux5bfcux6570}

\textbf{定义 164.1}（Wiener 空间）：Wiener 空间
\((W, \mathcal{H}, \mu)\) 是经典 Wiener 空间，其中
\(W = C_0([0, T], \mathbb{R})\)（连续轨道空间），\(\mathcal{H} = L^2([0, T])\)（Cameron-Martin
空间），\(\mu\) 是 Wiener 测度。

\textbf{定义 164.2}（光滑柱函数）：光滑柱函数 \(F: W \to \mathbb{R}\)
形如 \(F(\omega) = f(W_{t_1}(\omega), \ldots, W_{t_n}(\omega))\)，其中
\(f \in C^\infty(\mathbb{R}^n)\) 有界且有界导数。

\textbf{定义 164.3}（Malliavin 导数）：光滑柱函数 \(F\)
的\textbf{Malliavin 导数} \(DF \in L^2(\Omega; \mathcal{H})\) 定义为

\[D_t F = \sum_{i=1}^n \frac{\partial f}{\partial x_i}(W_{t_1}, \ldots, W_{t_n}) \mathbf{1}_{[0, t_i]}(t)\]

即 \(DF\) 是 \(\mathcal{H} = L^2([0,T])\) 值的随机变量。

\textbf{定义 164.4}（Sobolev 空间
\(\mathbb{D}^{1,p}\)）：\(\mathbb{D}^{1,p}\) 是光滑柱函数在范数
\(\|F\|_{1,p} = (\mathbb{E}[|F|^p] + \mathbb{E}[\|DF\|_{\mathcal{H}}^p])^{1/p}\)
下的完备化。

\textbf{定理 164.1}（Clark-Ocone 公式）：对 \(F \in \mathbb{D}^{1,2}\)，

\[F = \mathbb{E}[F] + \int_0^T \mathbb{E}[D_t F | \mathcal{F}_t] \, dW_t\]

这是 Itô 鞅表示定理的显式公式，给出了被积函数。

\subsubsection{164.2 Skorohod
积分与散度}\label{skorohod-ux79efux5206ux4e0eux6563ux5ea6}

\textbf{定义 164.5}（Skorohod 积分 / 散度算子）：\(\delta\) 是 Malliavin
导数 \(D\) 的伴随算子：对
\(u \in \operatorname{Dom}(\delta) \subseteq L^2(\Omega; \mathcal{H})\)，

\[\mathbb{E}[F \delta(u)] = \mathbb{E}[\langle DF, u \rangle_{\mathcal{H}}]\]

Skorohod 积分是 Itô 积分的推广：当 \(u\) 是适应过程时，\(\delta(u)\)
等于 Itô 积分 \(\int_0^T u_t dW_t\)。

\textbf{定理 164.2}（Malliavin 积分的公式）：若
\(u \in \mathbb{D}^{1,2}(\mathcal{H})\)，则

\[\mathbb{E}[\delta(u)^2] = \mathbb{E}\left[\int_0^T u_t^2 dt\right] + \mathbb{E}\left[\int_0^T \int_0^T D_s u_t D_t u_s \, ds dt\right]\]

\subsubsection{164.3 Malliavin
分析的应用}\label{malliavin-ux5206ux6790ux7684ux5e94ux7528}

\textbf{定理 164.3}（Malliavin 概率密度光滑性定理）：若
\(F = (F_1, \ldots, F_n)\) 是 Wiener 泛函满足： 1.
\(F_i \in \mathbb{D}^\infty\)（所有 \(\mathbb{D}^{k,p}\) 的交） 2.
Malliavin 协方差矩阵
\(\gamma_{ij} = \langle DF_i, DF_j \rangle_{\mathcal{H}}\)
几乎必然可逆，且 \(\det(\gamma)^{-1} \in L^p\) 对所有 \(p\)

则 \(F\) 的分布有 \(C^\infty\) 密度。

\emph{证明概要}（Malliavin概率密度光滑性）：核心思想是将密度函数表为期望形式。对多重指标
\(\alpha\)，利用分部积分公式（对偶性
\(\mathbb{E}[F\delta(u)] = \mathbb{E}[\langle DF,u\rangle]\)），可证
\(\mathbb{E}[\partial^\alpha \phi(F)]\) 等于
\(\mathbb{E}[\phi(F) \cdot H_\alpha(F)]\)，其中 \(H_\alpha\) 是涉及
\(\gamma^{-1}\) 的Wiener泛函。由 \(\det(\gamma)^{-1} \in L^p\) 知
\(H_\alpha \in L^1\)，故 \(F\) 的分布确有任意阶弱导数，从而是
\(C^\infty\)
函数。此定理是Malliavin分析的核心成就------它证明了Hörmander条件下扩散过程密度函数的光滑性，而不依赖PDE的先验估计。∎

\textbf{定理 164.4}（Hörmander 定理的 Malliavin 证明）：若 SDE 的 Lie
括号生成 \(\mathbb{R}^n\)（Hörmander 条件），则解 \(X_t\)
的分布有光滑密度。Malliavin 的概率证明是此定理的突破性证明方法。

\bigskip\hrule\bigskip

\subsection{第165章：随机偏微分方程}\label{ux7b2c165ux7ae0ux968fux673aux504fux5faeux5206ux65b9ux7a0b}

随机偏微分方程（SPDE）研究带有随机噪声的偏微分方程，在物理、生物和金融中有广泛应用。

\subsubsection{165.1 SPDE
的基本框架}\label{spde-ux7684ux57faux672cux6846ux67b6}

\textbf{定义 165.1}（随机偏微分方程 / SPDE）：一般形式为

\[\partial_t u = \mathcal{L} u + f(u) + \sigma(u) \xi(t, x)\]

其中 \(\mathcal{L}\) 是空间微分算子（如 \(\Delta\)），\(\xi(t, x)\)
是时空白噪声（或更一般的噪声）。

\textbf{定义 165.2}（时空白噪声）：\(\xi(t, x)\) 是
\(\mathbb{R} \times \mathbb{R}^d\) 上的广义 Gauss 随机场，满足

\[\mathbb{E}[\xi(t, x)] = 0, \quad \mathbb{E}[\xi(t, x)\xi(s, y)] = \delta(t-s)\delta(x-y)\]

\textbf{定理 165.1}（Walsh 的随机积分与 SPDE，1986）：SPDE
可写为积分方程

\[u(t, x) = \int G(t-s, x-y) u_0(y) dy + \int_0^t \int G(t-s, x-y) f(u(s,y)) dy ds + \int_0^t \int G(t-s, x-y) \sigma(u(s,y)) W(ds, dy)\]

其中 \(G\) 是 \(\mathcal{L}\) 的 Green 函数，\(W(ds, dy)\) 是 Brown
膜（Brownian sheet）。

\subsubsection{165.2 几个经典 SPDE}\label{ux51e0ux4e2aux7ecfux5178-spde}

\textbf{例
165.1}（随机热方程）：\(\partial_t u = \Delta u + \sigma(u) \xi(t, x)\)。当
\(\sigma(u) = 1\)（加法噪声）时，解为

\[u(t, x) = \int G(t, x-y) u_0(y) dy + \int_0^t \int G(t-s, x-y) W(ds, dy)\]

空间维数 \(d=1\) 时解是函数值（\(G\) 可积）；\(d=2\)
时解是分布值；\(d \geq 3\) 时解需要重整化。

\textbf{例 165.2}（随机 Burgers
方程）：\(\partial_t u + u \partial_x u = \nu \partial_x^2 u + \sigma \xi(t, x)\)。这是湍流的简化模型。

\textbf{例 165.3}（随机 Navier-Stokes
方程）：\(\partial_t u + (u \cdot \nabla)u = \nu \Delta u - \nabla p + \sigma \xi\)，\(\nabla \cdot u = 0\)。这是湍流理论的随机模型。

\subsubsection{165.3 KPZ
方程与正则性结构}\label{kpz-ux65b9ux7a0bux4e0eux6b63ux5219ux6027ux7ed3ux6784}

\textbf{定义 165.3}（KPZ 方程 /
Kardar-Parisi-Zhang，1986）：\(\partial_t h = \nu \partial_x^2 h + \frac{\lambda}{2}(\partial_x h)^2 + \sigma \xi(t, x)\)。这是界面生长的普适类模型。

\textbf{定理 165.2}（Hairer 的正则性结构理论，2014）：KPZ
方程在正则性结构（regularity structures）框架下是局部良定的。Hairer
因此获得 2014 年 Fields 奖。

\textbf{定义 165.4}（\(\Phi^4_d\) 模型）：随机量子化过程中出现的
SPDE：\(\partial_t \phi = \Delta \phi - \phi^3 + \xi\)。\(d=2\) 时需
Wick 重整化，\(d=3\) 时需更复杂的重整化。

\textbf{定理 165.3}（Hairer 的 \(\Phi^4_3\) 解，2014）：\(\Phi^4_3\)
方程的局部存在性由正则性结构理论证明。这是量子场论在数学上严格化的里程碑。

\bigskip\hrule\bigskip

\subsection{第166章：随机过程与随机几何}\label{ux7b2c166ux7ae0ux968fux673aux8fc7ux7a0bux4e0eux968fux673aux51e0ux4f55}

本章介绍随机分析在几何中的应用，包括流形上的 Brown 运动和随机微分几何。

\subsubsection{166.1 流形上的 Brown
运动}\label{ux6d41ux5f62ux4e0aux7684-brown-ux8fd0ux52a8}

\textbf{定义 166.1}（Riemann 流形上的 Brown 运动）：设 \((M, g)\) 是
Riemann 流形。\(M\) 上的 Brown 运动是以 Laplace-Beltrami 算子
\(\Delta_g\) 的一半为生成元的扩散过程。其 SDE 形式为

\[dX_t = \sum_{i=1}^d \sigma_i(X_t) \circ dW_t^i\]

其中 \(\sigma_i\) 是正交标架，\(\circ\) 表示 Stratonovich 积分（与 Itô
积分不同，Stratonovich 积分满足链式法则）。

\textbf{定义 166.2}（Stratonovich 积分）：Stratonovich 积分
\(\int X_t \circ dW_t\)
在极限定义中使用中点求值，满足普通微积分的链式法则。与 Itô 积分的关系为

\[\int X_t \circ dW_t = \int X_t dW_t + \frac{1}{2} \int d\langle X, W \rangle_t\]

\subsubsection{166.2
随机微分几何}\label{ux968fux673aux5faeux5206ux51e0ux4f55}

\textbf{定义 166.3}（Eells-Elworthy-Malliavin 方法）：流形上的 Brown
运动可构造为标架丛上的水平提升。在 \(O(M)\)（正交标架丛）上解 SDE
\(dU_t = \sum H_i(U_t) \circ dW_t^i\)，其中 \(H_i\)
是典范水平向量场。\(X_t = \pi(U_t)\) 是 \(M\) 上的 Brown 运动。

\textbf{定理 166.1}（Bismut 公式）：Riemann 流形上的热核 \(p_t(x, y)\)
的梯度可由 Brown 运动轨道的期望表示：

\[\nabla_x \log p_t(x, y) = \mathbb{E}^x[ \cdots ]\]

这是 Bismut 对 Atiyah-Singer 指标定理和热核估计的概率方法。

\subsubsection{166.3 随机 Loewner
演化（SLE）}\label{ux968fux673a-loewner-ux6f14ux5316sle}

\textbf{定义 166.4}（SLE / Schramm-Loewner
Evolution，2000）：\textbf{SLE\(_\kappa\)}（\(\kappa > 0\)）是上半平面
\(\mathbb{H}\) 中的随机曲线族，由 Loewner 方程驱动：

\[\partial_t g_t(z) = \frac{2}{g_t(z) - \sqrt{\kappa} W_t}, \quad g_0(z) = z\]

其中 \(W_t\) 是一维 Brown 运动。SLE 曲线是 \(g_t\) 的奇点生成的。

\textbf{定理 166.2}（SLE 的相变）： - \(\kappa \leq 4\)：SLE
曲线是简单曲线（不自交） - \(4 < \kappa < 8\)：SLE 曲线自交但不充满空间
- \(\kappa \geq 8\)：SLE 曲线是空间填充曲线

\textbf{定理 166.3}（SLE 与统计物理的联系 / Lawler-Schramm-Werner）： -
SLE\(_2\) 对应 loop-erased random walk 的标度极限 - SLE\(_6\)
对应临界渗流的界面的标度极限（Smirnov 2001） - SLE\(_8/3\)
对应自避行走的标度极限（猜想） - SLE\(_4\) 对应调和探索路径（harmonic
explorer）

\textbf{定理 166.4}（Lawler-Schramm-Werner 的 Fields 奖工作，2006）：SLE
是平面布朗运动边界交点的标度极限，由此证明了 Mandelbrot 的 Brown
运动边界的 Hausdorff 维数为 \(4/3\) 的猜想。

\bigskip\hrule\bigskip

\subsection{第167章：Levy过程}\label{ux7b2c167ux7ae0levyux8fc7ux7a0b}

Levy过程是 Brown 运动和 Poisson
过程的共同推广，定义为具有平稳独立增量的随机过程
\(X = \{X_t\}_{t \geq 0}\)（\(X_0 = 0\)
a.s.），样本轨道右连左极（cadlag）。其核心特征是 \textbf{Levy-Khintchine
公式}：特征函数具有形式
\(\mathbb{E}[e^{i\theta X_t}] = \exp(t \cdot \psi(\theta))\)，其中\textbf{特征指数}

\[\psi(\theta) = i\gamma\theta - \frac{1}{2}\sigma^2\theta^2 + \int_{\mathbb{R}\setminus\{0\}} \left(e^{i\theta x} - 1 - i\theta x \mathbf{1}_{|x|<1}\right) \nu(dx)\]

三元组 \((\gamma, \sigma^2, \nu)\) 唯一刻画 Levy
过程：\(\gamma \in \mathbb{R}\) 为漂移系数，\(\sigma^2 \geq 0\) 为 Gauss
方差，\(\nu\) 是 \textbf{Levy测度}（满足
\(\int (1 \wedge x^2)\nu(dx) < \infty\)），描述跳跃的频率和大小分布。跳跃部分可分解为
Poisson 随机测度 \(N(dt, dx)\)（计数跳跃）与其补偿子
\(\tilde{N}(dt, dx) = N(dt, dx) - dt\,\nu(dx)\)，后者是鞅测度，从而 Levy
过程的 Ito 分解为
\(dX_t = \gamma dt + \sigma dW_t + \int_{|x|<1} x\tilde{N}(dt,dx) + \int_{|x|\geq 1} x N(dt,dx)\)。

\bigskip\hrule\bigskip

\subsection{第168章：测度值过程（Measure-Valued
Processes）}\label{ux7b2c168ux7ae0ux6d4bux5ea6ux503cux8fc7ux7a0bmeasure-valued-processes}

测度值过程是取值为测度空间（而非
\(\mathbb{R}^d\)）的随机过程，在种群遗传学、统计物理和随机偏微分方程中起关键作用。\textbf{核心类型}：（1）\textbf{Dawson-Watanabe
超过程（superprocess）}------作为分支 Brown 运动的尺度极限：考虑大量独立
Brown
粒子，每个粒子以固定速率分裂为两个，总质量过程取扩散尺度极限获得取值于有限测度空间的
Markov 过程。超过程满足 \textbf{非线性
SPDE}（随机反应-扩散方程）：\(dX_t = \frac{\sigma^2}{2}\Delta X_t dt + \sqrt{\gamma X_t} dW_t^{(\text{空间-时间})}\)，其中
\(W_t^{(\text{空间-时间})}\) 为时空白噪声；（2）\textbf{Fleming-Viot
过程}------描述种群遗传学中基因频率的随机演化，是取值于概率测度空间的扩散过程，其生成元涉及二阶变分导数；（3）\textbf{Dirichlet
过程与随机测度}------Ferguson (1973) 引入的 Dirichlet 过程是 Bayesian
非参数统计的基本工具，其构造由 Gamma
过程的归一化给出。\textbf{关键定理}：（1）超过程的绝对连续性------超过程测度对
Lebesgue
测度保持绝对连续（当且仅当初始测度绝对连续）；（2）相变与局部灭绝------低维空间（\(d \leq 2\)）中超过程长期极限为零测度（局部灭绝），高维存在非平凡稳态；（3）构造的鞅方法（Dynkin
的 Laplace
泛函刻画）。测度值过程在度量-测度空间的随机几何（如随机平面树的 Brown
映射极限）中有前沿应用。

\bigskip\hrule\bigskip

\subsection{第169章：大偏差理论（Large Deviations
Theory）}\label{ux7b2c169ux7ae0ux5927ux504fux5deeux7406ux8bbalarge-deviations-theory}

大偏差理论量化罕见事件的指数衰减速率，是概率论中极限定理的精深推广。\textbf{核心框架}（Varadhan,
1966）：设 \(\{\mu_\epsilon\}_{\epsilon > 0}\) 为 Polish 空间
\(\mathcal{X}\) 上的一族概率测度。称 \(\mu_\epsilon\) 满足速度为
\(\epsilon^{-1}\)（\(\epsilon \to 0\)）的大偏差原理（LDP），若存在下半连续的
\textbf{速率函数}
\(I: \mathcal{X} \to [0, \infty]\)（具有紧水平集），使得对任意 Borel 集
\(A \subset \mathcal{X}\)：\(-\inf_{x \in A^\circ} I(x) \leq \liminf_{\epsilon \to 0} \epsilon \log \mu_\epsilon(A) \leq \limsup_{\epsilon \to 0} \epsilon \log \mu_\epsilon(A) \leq -\inf_{x \in \bar{A}} I(x)\)。

\textbf{Cramér 定理}（1938）：设 \(X_1, X_2, \ldots\) 为 i.i.d.
\(\mathbb{R}^d\) 值随机向量，则样本均值
\(\bar{X}_n = \frac{1}{n} \sum_{i=1}^n X_i\) 满足速度 \(n\) 的
LDP，速率函数为对数矩母函数
\(\Lambda(\lambda) = \log \mathbb{E}[e^{\langle \lambda, X \rangle}]\)
的 Legendre-Fenchel
变换：\(I(x) = \sup_{\lambda \in \mathbb{R}^d} [\langle \lambda, x \rangle - \Lambda(\lambda)]\)。

\textbf{Schilder 定理}：Brown 运动
\(B_t\)（\(t \in [0, 1]\)）经缩放后的小参数过程
\(\{\sqrt{\varepsilon} B_t\}\) 的路径分布在 \(C([0,1])\) 上满足
LDP，速率函数为

\[I(\varphi) = \frac{1}{2} \int_0^1 |\dot{\varphi}(t)|^2 dt\]

（若 \(\varphi\) 绝对连续且 \(\dot{\varphi} \in L^2\)），否则
\(I(\varphi) = \infty\)。这是 Cameron-Martin 空间的 Dirichlet 能量。

\textbf{应用}：大偏差理论在统计力学（相变）、保险数学（破产概率的精确渐近）、通信理论（Shannon
信息论中误差指数）和金融数学（罕见金融危机事件的概率估计）中有深远应用。

\bigskip\hrule\bigskip

\subsection{第170章：经验过程与统计学习理论基础}\label{ux7b2c170ux7ae0ux7ecfux9a8cux8fc7ux7a0bux4e0eux7edfux8ba1ux5b66ux4e60ux7406ux8bbaux57faux7840}

经验过程理论（empirical process
theory）提供了现代非参数统计和机器学习中一致收敛性的严格框架。\textbf{经验测度}：对
i.i.d. 样本 \(X_1, \ldots, X_n \sim P\)，定义\textbf{经验测度}
\(P_n = \frac{1}{n} \sum_{i=1}^n \delta_{X_i}\)（\(\delta_x\) 为 Dirac
测度）。对函数类 \(\mathcal{F}\)，\textbf{经验过程} \$ \nu\_n(f) =
\sqrt{n} (P\_n - P) f\$ 在 \(\mathcal{F}\) 上的 sup
行为的极限理论决定了许多统计估计量的一致性。

\textbf{Glivenko-Cantelli 定理}（一致大数律）：函数类 \(\mathcal{F}\) 是
\textbf{\(P\)-GC 类}（Glivenko-Cantelli）当且仅当
\(\sup_{f \in \mathcal{F}} |P_n f - P f|\) 以概率 \(1\) 趋近于
\(0\)。经典的 Glivenko-Cantelli 定理对应
\(\mathcal{F} = \{ 1_{(-\infty, t]} : t \in \mathbb{R} \}\)（指示函数类）。一般判定准则涉及函数类的覆盖数（covering
number）是否满足 Pollard 熵条件。

\textbf{Donsker 定理}（一致中心极限定理）：函数类 \(\mathcal{F}\) 是
\textbf{\(P\)-Donsker 类}当且仅当经验过程 \(\nu_n\) 作为
\(\ell^\infty(\mathcal{F})\) 中的随机元弱收敛于 \(P\)-Brown 桥
\(G_P\)（即 \(\mathbb{E}[G_P(f)G_P(g)] = P(fg) - Pf Pg\) 为中心的 Gauss
过程）。判定准则由 \textbf{bracketing 数} 或 VC
维数（Vapnik-Chervonenkis 维数）提供------若 \(\mathcal{F}\) 的 VC
维数有限，则 \(\mathcal{F}\) 是 Donsker 类（对任意
\(P\)）。这一结果直接支撑了统计学习理论中经验风险最小化（ERM）的一致性证明。

\bigskip\hrule\bigskip

\emph{本卷在 V20（概率论 II）Itô 积分简介的基础上，深入建立了 Itô
随机微积分、随机微分方程（SDE）、Malliavin
分析、随机偏微分方程（SPDE）、随机几何（含 SLE 理论）和 Levy 过程。共 6
章。}

\newpage

\section{卷三十三：数学史与数学哲学}\label{ux5377ux4e09ux5341ux4e09ux6570ux5b66ux53f2ux4e0eux6570ux5b66ux54f2ux5b66}

\begin{quote}
数学史与数学哲学是理解数学本质和演化的两个维度。数学史追溯数学思想从古代到现代的发展历程，数学哲学探讨数学的基础、真理性和方法论。本卷作为全书的收官卷，回顾数学史的关键里程碑，探讨数学哲学的四大流派，并反思数学的当代意义。
\end{quote}

\bigskip\hrule\bigskip

\subsection{第172章：古代数学史}\label{ux7b2c172ux7ae0ux53e4ux4ee3ux6570ux5b66ux53f2}

古代数学从巴比伦和埃及的实用计算开始，经古希腊的演绎推理体系的建立，到中国和印度数学的独立发展，奠定了数学的早期基础。

\subsubsection{172.1
巴比伦与埃及数学}\label{ux5df4ux6bd4ux4f26ux4e0eux57c3ux53caux6570ux5b66}

\textbf{定义 172.1}（巴比伦数学，约公元前 2000-300
年）：巴比伦数学以六十进制记数法和代数为特色，遗留了约 400
块泥板（如普林顿 322 号泥板，包含 Pythagorean
三元组）。巴比伦人求解二次方程，使用线性插值计算复利，并近似
\(\sqrt{2}\) 到 \(1.414213\)（YBC 7289 号泥板）。

\textbf{定义 172.2}（埃及数学，约公元前 2000-1600 年）：埃及数学记录在
Rhind 数学纸草卷（约公元前 1650
年）和莫斯科数学纸草卷中。埃及人使用单位分数、求解线性方程，并计算截锥体积（莫斯科纸草卷第
14 题）。

\subsubsection{172.2 古希腊数学}\label{ux53e4ux5e0cux814aux6570ux5b66}

\textbf{定义 172.3}（毕达哥拉斯学派，约公元前 570-495
年）：「万物皆数」的信条。发现不可公度量（无理数
\(\sqrt{2}\)）的危机，促使几何学主导希腊数学。毕达哥拉斯定理和正多面体分类。

\textbf{定义 172.4}（Euclid 的《几何原本》，约公元前 300
年）：《几何原本》是公理化方法的最早典范，从 5 条公设和 5
条公理出发，推导出 465 个命题。成为西方数学教育的基础长达两千年。

\textbf{公设 172.1}（Euclid《几何原本》的五大公设）： 1.
过任意两点可作一条直线 2. 线段可无限延长 3.
以任意点为中心、任意长为半径可作圆 4. 所有直角相等 5.
若两直线与第三条直线相交，同侧内角之和小于两直角，则两直线在该侧相交（平行公设）

\textbf{定义 172.5}（Archimedes，约公元前 287-212
年）：古代最伟大的数学家。穷竭法计算
\(\pi\)（\(3\frac{10}{71} < \pi < 3\frac{1}{7}\)），抛物弓形面积，球体积和表面积，杠杆原理。Archimedes
的无穷小方法预示了 17 世纪的微积分。

\textbf{定义 172.6}（Apollonius，约公元前 262-190
年）：《圆锥曲线》------系统研究椭圆、抛物线和双曲线，命名了这些曲线，发现其数百个性质。

\textbf{定义 172.7}（Diophantus，约公元 250
年）：《算术》------数论的先驱，研究不定方程（Diophantine
方程），引入了符号表示法。

\subsubsection{172.3
中国与印度古代数学}\label{ux4e2dux56fdux4e0eux5370ux5ea6ux53e4ux4ee3ux6570ux5b66}

\textbf{定义 172.8}（《九章算术》，约公元前 1 世纪-公元 1
世纪）：中国数学的经典，246
个问题覆盖方田、粟米、衰分、少广、商功、均输、盈不足、方程和勾股。包含了线性方程组的消元法（Gauss
消元法的前身）和负数的使用。

\textbf{定义 172.9}（祖冲之，429-500 年）：将 \(\pi\)
精确到小数点后七位：\(3.1415926 < \pi < 3.1415927\)，给出密率
\(\frac{355}{113}\)（千年内最精确的 \(\pi\) 近似）。

\textbf{定义 172.10}（印度数学的贡献）：Aryabhata（476-550 年）计算
\(\pi\) 到四位小数，引入正弦函数。Brahmagupta（598-668
年）引入零作为数字和使用负数。Bhaskara II（1114-1185
年）研究微分学初步和 Pell 方程。

\subsubsection{172.4
伊斯兰黄金时代数学}\label{ux4f0aux65afux5170ux9ec4ux91d1ux65f6ux4ee3ux6570ux5b66}

\textbf{定义 172.11}（al-Khwarizmi，约 780-850
年）：《代数学》（al-Kitāb
al-Mukhtaṣar）------「代数」（algebra）一词的来源。系统求解一次和二次方程，引入「算法」（algorithm）概念。

\textbf{定义 172.12}（其他伊斯兰数学家）：Omar Khayyam（1048-1131
年）用圆锥曲线求解三次方程。al-Kashi（1380-1429 年）计算 \(\pi\) 到 16
位小数。伊斯兰数学家保存并发展了希腊数学，为欧洲文艺复兴的数学复兴奠定了基础。

\bigskip\hrule\bigskip

\subsection{第173章：近代数学史（16-19
世纪）}\label{ux7b2c173ux7ae0ux8fd1ux4ee3ux6570ux5b66ux53f216-19-ux4e16ux7eaa}

16 至 19
世纪是数学的爆炸式发展期，从代数方程求解到微积分的创立，再到分析基础的严格化。

\subsubsection{173.1
文艺复兴与代数革命}\label{ux6587ux827aux590dux5174ux4e0eux4ee3ux6570ux9769ux547d}

\textbf{定义 173.1}（三次和四次方程的解，16 世纪）：Scipione del Ferro
和 Tartaglia 发现三次方程的解法，Cardano 在《大术》（Ars
Magna，1545）中发表。Cardano 的学生 Ferrari 解决了四次方程。Cardano
首次系统使用负数和复数。

\textbf{定义 173.2}（Viète 的符号代数，约 1590 年）：François Viète
引入字母表示已知量和未知量，将代数从具体数字推广到一般符号，是代数符号化的关键一步。

\textbf{定义 173.3}（Napier 的对数，1614 年）：John Napier
发明对数，将乘法化为加法，「将天文学家的寿命延长了一倍」（Laplace）。

\subsubsection{173.2
解析几何与微积分的创立}\label{ux89e3ux6790ux51e0ux4f55ux4e0eux5faeux79efux5206ux7684ux521bux7acb}

\textbf{定义 173.4}（Descartes 的《几何学》，1637 年）：René Descartes
创立解析几何，将几何问题转化为代数方程，引入了坐标系（Cartesian 坐标）。

\textbf{定义 173.5}（Fermat 的贡献，约 1629-1665 年）：Pierre de Fermat
独立于 Descartes 发展了坐标几何，提出了 Fermat
大定理（\(x^n + y^n = z^n\) 对 \(n > 2\) 无正整数解），在概率论（与
Pascal 的通信）中建立了数学期望的概念。

\textbf{定义 173.6}（Newton 与 Leibniz 的微积分，17 世纪后半叶）：Isaac
Newton（1642-1727）和 Gottfried Wilhelm
Leibniz（1646-1716）分别独立创立微积分。Newton
从物理问题（运动、变化率）出发，使用「流数」概念。Leibniz
从几何问题（切线、面积）出发，引入 \(dx\)、\(\int\)
等符号，至今沿用。两人的优先权之争是数学史上最著名的争议之一。

\textbf{定理 173.1}（微积分基本定理 / Newton-Leibniz
公式）：\(\int_a^b f(x) dx = F(b) - F(a)\)，其中
\(F'(x) = f(x)\)。这是微分和积分之间的桥梁。

\subsubsection{173.3 18
世纪：分析的时代}\label{ux4e16ux7eaaux5206ux6790ux7684ux65f6ux4ee3}

\textbf{定义 173.7}（Euler 的贡献，1707-1783 年）：Leonhard Euler
是历史上最高产的数学家，全集 80
余卷。主要贡献：\(e^{i\theta} = \cos\theta + i\sin\theta\)（Euler
公式），图论（Königsberg 七桥问题），数论（Euler \(\varphi\)
函数），变分法，流体力学方程。Euler 引入了
\(e\)、\(i\)、\(\pi\)、\(\sum\)、\(f(x)\) 等符号。

\textbf{定义 173.8}（Lagrange 的贡献，1736-1813 年）：Joseph-Louis
Lagrange 在《分析力学》（1788）中将力学完全用分析学表述，开创了变分法和
Lagrange 力学。在数论上证明了四平方和定理和 Wilson 定理。

\textbf{定义 173.9}（Laplace 的贡献，1749-1827 年）：Pierre-Simon
Laplace 在《天体力学》中系统应用微积分于天文学，在概率论中提出 Laplace
变换和古典概率的定义。

\subsubsection{173.4 19
世纪：严格化与新几何}\label{ux4e16ux7eaaux4e25ux683cux5316ux4e0eux65b0ux51e0ux4f55}

\textbf{定义 173.10}（Gauss 的贡献，1777-1855 年）：Carl Friedrich
Gauss，「数学王子」。贡献遍布数论（《算术研究》，1801）、微分几何（Theorema
Egregium）、统计学（最小二乘法）、代数（代数基本定理的证明）、非欧几何（未发表）。

\textbf{定义 173.11}（分析严格化）：Cauchy（1789-1857）引入极限的
\(\varepsilon\)-\(\delta\)
定义。Weierstrass（1815-1897）彻底严格化了分析学的基础，给出了处处连续但无处可微的函数的例子。Riemann（1826-1866）定义了
Riemann 积分。

\textbf{定义 173.12}（非欧几何的创立）：Lobachevsky（1792-1856）和
Bolyai（1802-1860）独立发表非欧几何（双曲几何），否定了 Euclid
平行公设的必然性。Riemann 在 1854 年发展了 Riemann 几何（椭圆几何），为
Einstein 的广义相对论提供了数学框架。

\textbf{定义 173.13}（Galois 与群论，约 1830 年）：Évariste
Galois（1811-1832）在 20
岁时用群论证明了五次及以上方程无根式解（Abel-Ruffini
定理的理论基础），并奠定了群论的基础。Galois 在 21 岁时死于决斗。

\textbf{定义 173.14}（Cantor 与集合论，约 1874-1897 年）：Georg
Cantor（1845-1918）创立集合论，发现无穷有不同大小（可数无穷与连续统），对角线论证证明了
\(\mathbb{R}\) 不可数。连续统假设（\(\aleph_1 = 2^{\aleph_0}\)？）是
Hilbert 23 个问题的第一个。

\bigskip\hrule\bigskip

\subsection{第174章：现代数学史（20-21
世纪）}\label{ux7b2c174ux7ae0ux73b0ux4ee3ux6570ux5b66ux53f220-21-ux4e16ux7eaa}

20
世纪以来，数学在深度和广度上都经历了前所未有的革命。公理化、结构化、抽象化成为主流，数学与物理和计算机科学的联系日益紧密。

\subsubsection{174.1
数学基础的危机与解决}\label{ux6570ux5b66ux57faux7840ux7684ux5371ux673aux4e0eux89e3ux51b3}

\textbf{定义 174.1}（Russell 悖论，1901 年）：Bertrand Russell 发现
Cantor 朴素集合论中的悖论：令 \(R = \{x : x \notin x\}\)，则
\(R \in R \iff R \notin R\)。这引发了数学基础的第三次危机。

\textbf{定义 174.2}（三大数学基础学派）： - \textbf{逻辑主义}（Frege,
Russell, Whitehead）：数学可归约为逻辑。《数学原理》（Principia
Mathematica，1910-1913）试图从逻辑公理推导所有数学，但未完全成功 -
\textbf{直觉主义}（Brouwer，1907-）：数学是心智构造，排斥排中律和实无穷。直觉主义逻辑（Heyting）和直觉主义数学（Brouwer
的不动点定理） -
\textbf{形式主义}（Hilbert，1920s）：数学是符号操作，Hilbert
纲领试图用有限方法证明数学的一致性

\textbf{定义 174.3}（Gödel 不完全性定理，1931 年）：见 V19 Ch 102。Gödel
证明了形式主义的局限性：任何包含算术的一致形式系统必然是不完全的。

\subsubsection{174.2 Bourbaki
与结构主义}\label{bourbaki-ux4e0eux7ed3ux6784ux4e3bux4e49}

\textbf{定义 174.4}（Bourbaki 学派，1935-）：以 Nicolas Bourbaki
为笔名的法国数学家群体，目标是「从零开始，用公理化方法重建整个数学体系」。其《数学原本》（Éléments
de
Mathématique）从集合论出发，体系化地展开代数、拓扑、分析等分支。核心概念是「结构」：代数结构、序结构和拓扑结构。

\textbf{定义 174.5}（Bourbaki 的影响与局限）：Bourbaki
的抽象化和公理化方法对 20
世纪数学教育产生了深远影响。但其忽视应用数学、几何直观和计算数学的特点也受到批评（如
Arnold 的批评）。

\subsubsection{174.3 20
世纪数学的重大突破}\label{ux4e16ux7eaaux6570ux5b66ux7684ux91cdux5927ux7a81ux7834}

\textbf{定理 174.1}（Deligne 的 Weil 猜想证明，1974 年）：Pierre Deligne
证明了 Weil 猜想（有限域上代数簇的 Zeta 函数的 Riemann
猜想类比），是代数几何与数论交汇的巅峰成就。Deligne 因此获得 1978 年
Fields 奖。

\textbf{定理 174.2}（Fermat 大定理的证明，1994 年）：Andrew Wiles（在
Richard Taylor 的协助下）证明了 Fermat 大定理：\(x^n + y^n = z^n\) 对
\(n > 2\)
无正整数非零解。证明基于椭圆曲线和模形式的深刻联系------Taniyama-Shimura-Weil
猜想。

\textbf{定理 174.3}（Poincaré 猜想的证明，2002-2003 年）：Grigori
Perelman 使用 Ricci 流证明了三维 Poincaré
猜想（任何单连通闭三维流形同胚于
\(S^3\)）。这是七大千禧年问题中第一个被解决的。Perelman 拒绝了 Fields
奖和千禧年奖金。

\textbf{定理 174.4}（其他千禧年问题）：七个千禧年问题（Clay
数学研究所，2000 年）： 1. Poincaré 猜想（已解决，Perelman 2003） 2. P
vs NP 问题（未解决） 3. Hodge 猜想（未解决） 4. Riemann 猜想（未解决）
5. Yang-Mills 存在性与质量间隙（未解决） 6. Navier-Stokes
方程的存在性与光滑性（未解决） 7. Birch and Swinnerton-Dyer
猜想（未解决）

\subsubsection{174.4
计算机与数学}\label{ux8ba1ux7b97ux673aux4e0eux6570ux5b66}

\textbf{定义 174.6}（计算机辅助证明）： - 四色定理（Appel-Haken
1976，Gonthier 2005 用 Coq 形式化）：首个依赖计算机的数学证明 - Kepler
猜想（Hales 1998-2014，Flyspeck 项目形式化） -
是否存在有限射影平面阶数为 10？（Lam 1989，计算机搜索否定）

\textbf{定义 174.7}（Langlands 纲领，1967-）：Robert Langlands
提出的「大统一理论」，将数论（Galois
表示）与调和分析（自守形式）联系起来。Langlands
纲领是当代数学中最宏大的研究纲领之一，包含函数域上的 Langlands
对应（Lafforgue 2002，Fields 奖）和几何 Langlands 纲领（Gaitsgory
等人的工作）。

\bigskip\hrule\bigskip

\subsection{第175章：数学哲学基础}\label{ux7b2c175ux7ae0ux6570ux5b66ux54f2ux5b66ux57faux7840}

数学哲学探讨数学的本质：数学对象是否存在？数学真理的性质是什么？数学知识如何可能？

\subsubsection{175.1
四大经典流派}\label{ux56dbux5927ux7ecfux5178ux6d41ux6d3e}

\textbf{定义 175.1}（柏拉图主义 /
数学实在论）：数学对象（数、集合、函数等）是独立于心智和物理世界存在的抽象实体。数学真理是关于这些实体的客观描述。代表人物：Gödel、Penrose。

\textbf{论证}：Quine-Putnam
不可缺论证------数学在科学中不可或缺，科学理论的成功确证了其数学部分的真理性。

\textbf{反对}：Benacerraf
的认识论问题------如果数学对象是因果惰性的抽象实体，人类如何获得数学知识？

\textbf{定义
175.2}（形式主义）：数学是关于无意义的符号串的形式操作。数学的「真理」是在一致的公理系统中可推导的公式。代表人物：Hilbert、Curry。

\textbf{论证}：数学不需要本体论承诺，只需关注逻辑一致性。

\textbf{反对}：Gödel
不完全性定理表明一致性不能自证。完全的形式主义无法解释为什么某些公理系统（如
ZFC）被选择而其他不被选择。

\textbf{定义
175.3}（逻辑主义）：数学可归约为逻辑。数学概念是逻辑概念的逻辑构造。代表人物：Frege、Russell。

\textbf{论证}：Frege 的《算术基础》（1884）和 Russell-Whitehead
的《数学原理》试图将算术归约为逻辑。

\textbf{反对}：Russell 悖论暴露了朴素集合论的不一致性。公理化集合论（如
ZFC）使用了超出纯逻辑的公理（如无穷公理、选择公理）。

\textbf{定义 175.4}（直觉主义 /
构造主义）：数学是心智构造活动。数学对象的存在性等于可构造性。排中律不适用于无穷域。代表人物：Brouwer、Heyting、Dummett、Bishop。

\textbf{论证}：直觉主义数学避免了非构造性存在证明的悖论。所有证明对应算法，与计算机科学有天然联系（Curry-Howard
同构）。

\textbf{反对}：拒绝排中律导致经典数学的许多标准结果不可接受（如代数学基本定理的非构造性证明）。直觉主义数学比经典数学更弱和更复杂。

\subsubsection{175.2
结构主义与范畴论}\label{ux7ed3ux6784ux4e3bux4e49ux4e0eux8303ux7574ux8bba}

\textbf{定义
175.5}（结构主义）：数学的本质是结构，而非个体对象。数不是「是什么」，而是「处于什么关系中」。数
3 在不同的系统中（自然数、实数、复数）是同一个结构位置的不同实例。

\textbf{定义 175.6}（范畴论与结构主义，Lawvere
1960s）：范畴论为结构主义提供了精确的数学语言。在范畴论框架下，数学对象的性质完全由其与其他对象的关系（态射）决定。Yoneda
引理：一个对象完全由它到所有其他对象的态射（或从所有对象到它的态射）确定。

\textbf{定理 175.1}（ETCS / Lawvere 的集合论公理化）：Lawvere
用范畴论语言重新公理化集合论（Elementary Theory of the Category of
Sets），将集合论的基础从「属于关系」转移到「函数与复合」。

\subsubsection{175.3
当代数学哲学问题}\label{ux5f53ux4ee3ux6570ux5b66ux54f2ux5b66ux95eeux9898}

\textbf{定义 175.7}（Hilbert
第六问题与数学物理）：数学为什么能如此有效地描述物理世界？（Wigner
的「数学在自然科学中不合理的有效性」，1960）这一问题至今没有令人满意的哲学解释。

\textbf{定义
175.8}（数学中的实验与计算机）：计算机辅助证明（如四色定理、Kepler
猜想）是否改变了数学的认识论性质？我们能否信任人类无法验证的计算机证明？

\textbf{定义
175.9}（选择公理的争论）：选择公理是当代数学中仍在争论的公理。它等价于
Zorn 引理、良序定理等，在泛函分析（Hahn-Banach
定理）和代数（任意向量空间有基）中不可或缺。但 Banach-Tarski
悖论（球可分解为有限块重组为两个全等球）表明选择公理蕴含反直觉的结果。

\bigskip\hrule\bigskip

\subsection{第176章：数学基础与当代数学观}\label{ux7b2c176ux7ae0ux6570ux5b66ux57faux7840ux4e0eux5f53ux4ee3ux6570ux5b66ux89c2}

本章作为全书终章，回顾数学基础的逻辑结构，反思数学在当代科学中的地位，展望数学的未来。

\subsubsection{176.1
数学基础的全书回顾}\label{ux6570ux5b66ux57faux7840ux7684ux5168ux4e66ux56deux987e}

\textbf{定义 176.1}（本书的体系结构）：本书从 Peano 公理（V1 Ch
1）出发，经过 56 卷、279
章的构建，覆盖了数学的主要分支。全书以公理化方法为基础，依赖链从数系
\(\to\) 代数 \(\to\) 几何 \(\to\) 分析 \(\to\)
现代数学的各个分支逐步展开。

\textbf{定理 176.1}（本书的依赖结构总结）：所有 56 卷的最底层依赖是
V1（数学基础：Peano 公理、集合论、数系）。V1
为全书提供了统一的逻辑基础和元语言。

\subsubsection{176.2
数学的统一性}\label{ux6570ux5b66ux7684ux7edfux4e00ux6027}

\textbf{定义 176.2}（数学的统一性 / Hilbert's Program）：David Hilbert
在 1900 年巴黎国际数学家大会上提出 23 个问题，展望了 20
世纪数学的发展方向。Hilbert
强调数学的统一性：「数学是一个有机整体，其生命力在于各部分之间的联系。」

\textbf{例 176.1}（数学中的深刻统一）： - \textbf{Langlands 纲领}：数论
\(\leftrightarrow\) 表示论 \(\leftrightarrow\) 调和分析 -
\textbf{Atiyah-Singer 指标定理}：分析（微分算子）\(\leftrightarrow\)
拓扑（K-理论）\(\leftrightarrow\) 几何（特征类） - \textbf{Mirror
symmetry}（镜对称）：代数几何 \(\leftrightarrow\) 辛几何 -
\textbf{模形式与椭圆曲线}：数论 \(\leftrightarrow\) 复分析
\(\leftrightarrow\) 代数几何 - \textbf{随机 Loewner 演化}：随机分析
\(\leftrightarrow\) 复分析 \(\leftrightarrow\) 统计物理

\subsubsection{176.3 数学的未来}\label{ux6570ux5b66ux7684ux672aux6765}

\textbf{定义
176.3}（未解决的大问题）：数学中有许多悬而未决的重大问题，包括： -
Riemann 猜想（解析数论） - P vs NP 问题（计算复杂性） - Navier-Stokes
方程的正则性（偏微分方程） - Yang-Mills 存在性与质量间隙（量子场论数学）
- Birch and Swinnerton-Dyer 猜想（算术几何） - Hodge 猜想（代数几何） -
Langlands 纲领的完全实现

\textbf{定义 176.4}（新兴方向）：数学的前沿在持续扩展： -
\textbf{人工智能与数学}：AI 辅助证明（AlphaProof、AlphaGeometry）、AI
发现新定理 - \textbf{量子计算数学}：量子算法、量子纠错码、拓扑量子计算 -
\textbf{数据科学数学}：高维统计、优化理论、信息论的新应用 -
\textbf{合成生物学中的数学}：基因调控网络的动态系统模型 -
\textbf{气候科学中的数学}：多尺度建模、不确定性量化

\subsubsection{176.4
数学作为人类精神}\label{ux6570ux5b66ux4f5cux4e3aux4ebaux7c7bux7cbeux795e}

\textbf{定义 176.5}（数学的美学价值）：Russell
说：「数学，正确看待时，不仅拥有真理，而且拥有至高无上的美------一种冷峻而朴素的美，如同雕塑的美。」数学的美在于其简洁性、深刻性、意外性和必然性。

\textbf{定义
176.6}（数学是人类精神的最高成就）：从古希腊人的几何证明到现代数学的抽象理论，数学始终是人类理性探索的极致体现。如
Gauss 所言：「数学是科学的皇后，数论是数学的皇后。」

\textbf{定理 176.2}（数学的无限性）：由 Gödel
不完全性定理，任何一致的公理系统都不能捕捉所有数学真理。数学的疆域永远在扩展，人类对数学真理的探索永远不会终结。

\bigskip\hrule\bigskip

\emph{全书共 56 卷、279 章，从 Peano 公理和 ZFC
集合论出发，系统覆盖了数学从基础到前沿的主要分支。数学是人类理性的最纯粹表达，也是最美丽的知识体系。本书的完结只是一个阶段性的总结，数学的探索永无止境。}

\emph{------ 全书完 ------}

\newpage

\section{卷三十四：复几何}\label{ux5377ux4e09ux5341ux56dbux590dux51e0ux4f55}

\begin{quote}
复几何是复分析和微分几何的交汇，研究复流形的几何结构。复几何是连接复分析（V9）、微分几何（V11）和代数几何（V15）三大支柱的桥梁。本卷从复流形和近复结构出发，建立
Kähler 几何的基本理论（Kähler 恒等式、Hodge 分解），深入讨论 Hodge
理论和形变理论，最后介绍 Calabi-Yau 流形和镜对称。
\end{quote}

\bigskip\hrule\bigskip

\subsection{第177章：复流形}\label{ux7b2c177ux7ae0ux590dux6d41ux5f62}

复流形是局部同胚于 \(\mathbb{C}^n\)
且转移函数为全纯的拓扑空间。复流形的研究是复几何的起点。

\subsubsection{177.1
复流形的定义}\label{ux590dux6d41ux5f62ux7684ux5b9aux4e49}

\textbf{定义 177.1}（复流形）：\(n\) 维\textbf{复流形} \(M\)
是一个拓扑空间，带有复坐标卡 \(\{(U_\alpha, \varphi_\alpha)\}\)，其中
\(\varphi_\alpha: U_\alpha \to \mathbb{C}^n\) 是同胚到 \(\mathbb{C}^n\)
中开集，且转移函数 \(\varphi_\alpha \circ \varphi_\beta^{-1}\) 在
\(\varphi_\beta(U_\alpha \cap U_\beta)\) 上是全纯的。

\textbf{例 177.1}（基本复流形）： - \(\mathbb{C}^n\) 和 \(\mathbb{C}^n\)
中的开集（平凡复流形） - \textbf{复射影空间}
\(\mathbb{CP}^n\)：\(\mathbb{C}^{n+1} \setminus \{0\}\) 模去
\(\mathbb{C}^*\) 作用，维数为 \(n\) - \textbf{复环面}
\(\mathbb{C}^n / \Lambda\)（\(\Lambda\) 是秩 \(2n\) 的格） -
\textbf{Riemann 面}：1 维复流形（V9 Ch 55 共形映射理论） - \textbf{复
Grassmann 流形} \(\operatorname{Gr}(k, n; \mathbb{C})\)

\textbf{定义 177.2}（全纯函数与全纯映射）：\(f: M \to \mathbb{C}\)
是\textbf{全纯}的，如果 \(f \circ \varphi_\alpha^{-1}\)
在每个坐标卡上是全纯的。\(f: M \to N\)（\(N\)
是复流形）是全纯的，如果局部坐标表示是全纯的。

\textbf{命题
177.1}（紧复流形上的全纯函数）：紧连通复流形上的全纯函数必为常数（极大值原理的推广）。这使复几何与实几何有本质区别。

\subsubsection{177.2
近复结构与可积性}\label{ux8fd1ux590dux7ed3ux6784ux4e0eux53efux79efux6027}

\textbf{定义 177.3}（近复结构）：实流形 \(M\)（\(2n\)
维）上的\textbf{近复结构}是切丛 \(TM\) 的自同态 \(J: TM \to TM\)，满足
\(J^2 = -\operatorname{id}\)。\(J\) 赋予 \(TM\) 复向量空间结构。

\textbf{定义 177.4}（Nijenhuis 张量）：近复结构 \(J\) 的
\textbf{Nijenhuis 张量}为

\[N_J(X, Y) = [X, Y] + J[JX, Y] + J[X, JY] - [JX, JY]\]

\(N_J\) 衡量 \(J\) 与 Lie 括号的不兼容性。

\textbf{定理 177.1}（Newlander-Nirenberg 定理，1957）：近复结构 \(J\)
来自复流形结构的充要条件是 \(N_J \equiv 0\)。此时 \(J\)
称为\textbf{可积的}。

\emph{意义}：该定理将复流形的存在性转化为近复结构的可积性条件，是复几何的核心定理。证明使用全纯坐标的构造（通过
\(\bar{\partial}\)-方程的可解性）。

\textbf{命题 177.2}（实二维情形）：任何 2 维可定向 Riemann
流形上的近复结构（由旋转 \(90^\circ\) 定义）总是可积的。因此任何可定向
Riemann 曲面是复流形（1 维）。

\subsubsection{177.3 Dolbeault
上同调}\label{dolbeault-ux4e0aux540cux8c03}

\textbf{定义 177.5}（\((p, q)\)-形式）：复流形 \(M\)
上的\textbf{\((p, q)\)-形式}是 \(p\) 个全纯微分和 \(q\)
个反全纯微分的外积。外微分 \(d\) 分解为
\(d = \partial + \bar{\partial}\)，其中 \(\partial\)
增加全纯次数，\(\bar{\partial}\) 增加反全纯次数。

\textbf{定义 177.6}（Dolbeault 上同调）：\(\bar{\partial}\)-算子满足
\(\bar{\partial}^2 = 0\)，定义\textbf{Dolbeault 上同调群}：

\[H_{\bar{\partial}}^{p,q}(M) = \frac{\ker(\bar{\partial}: \Omega^{p,q} \to \Omega^{p,q+1})}{\operatorname{im}(\bar{\partial}: \Omega^{p,q-1} \to \Omega^{p,q})}\]

\textbf{定理 177.2}（Dolbeault
定理）：\(H_{\bar{\partial}}^{p,q}(M) \cong H^q(M, \Omega^p)\)，其中
\(\Omega^p\) 是全纯 \(p\)-形式的层。Dolbeault
上同调是层上同调的具体实现。

\textbf{定义 177.7}（Hodge
数）：\(h^{p,q}(M) = \dim_{\mathbb{C}} H_{\bar{\partial}}^{p,q}(M)\)。Hodge
数是复流形的基本数值不变量。

\textbf{例 177.2}（射影空间的 Hodge 数）：对
\(\mathbb{CP}^n\)，\(h^{p,p} = 1\)（\(0 \leq p \leq n\)），其余
\(h^{p,q} = 0\)。

\bigskip\hrule\bigskip

\subsection{第178章：Kähler
几何}\label{ux7b2c178ux7ae0kuxe4hler-ux51e0ux4f55}

Kähler 流形是同时具有 Riemann
度量、复结构和辛结构且三者兼容的流形。Kähler 几何是复几何的绝对核心。

\subsubsection{178.1 Kähler
度量的定义}\label{kuxe4hler-ux5ea6ux91cfux7684ux5b9aux4e49}

\textbf{定义 178.1}（Hermite 度量）：复流形 \(M\) 上的 \textbf{Hermite
度量} \(h\) 是每个切空间 \(T_pM\)（视为复向量空间）上的 Hermite
内积，光滑依赖于 \(p\)。分解为 \(h = g - i\omega\)，其中
\(g = \operatorname{Re} h\) 是 Riemann
度量，\(\omega = -\operatorname{Im} h\) 是实 2-形式。

\textbf{定义 178.2}（Kähler 度量）：Hermite 度量 \(h\) 称为
\textbf{Kähler 度量}，如果相应的 \((1,1)\)-形式 \(\omega\)
是闭的：\(d\omega = 0\)。此时 \((M, h)\) 称为 \textbf{Kähler 流形}。

Kähler 条件等价于：复结构 \(J\) 关于 \(g\) 的 Levi-Civita
联络是平行的：\(\nabla J = 0\)。

\textbf{例 178.1}（Kähler 流形）： - \(\mathbb{C}^n\)
带标准欧氏度量（\(\omega = \frac{i}{2} \sum dz_j \wedge d\bar{z}_j\)） -
\(\mathbb{CP}^n\) 带 \textbf{Fubini-Study
度量}：\(\omega_{\text{FS}} = \frac{i}{2\pi} \partial \bar{\partial} \log \|Z\|^2\)（在齐次坐标中）
- 复环面 \(\mathbb{C}^n / \Lambda\)（诱导度量是 Kähler 的） - Kähler
流形的复子流形（诱导度量是 Kähler 的） - 光滑复射影代数簇（嵌入
\(\mathbb{CP}^n\) 后取诱导度量）

\textbf{命题 178.1}（Kähler 势）：在局部坐标下，Kähler 形式 \(\omega\)
可写为

\[\omega = \frac{i}{2} \sum_{j,k} h_{j\bar{k}} dz_j \wedge d\bar{z}_k\]

其中
\(h_{j\bar{k}} = \frac{\partial^2 \varphi}{\partial z_j \partial \bar{z}_k}\)，\(\varphi\)
是局部 \textbf{Kähler 势}。\(d\omega = 0\) 自动满足。

\subsubsection{178.2 Kähler 恒等式}\label{kuxe4hler-ux6052ux7b49ux5f0f}

\textbf{定理 178.1}（Kähler 恒等式）：在 Kähler 流形上，Hodge 星算子
\(\star\)、\(\partial\)、\(\bar{\partial}\) 及其伴随算子满足以下恒等式：

\begin{enumerate}
\def\labelenumi{\arabic{enumi}.}
\tightlist
\item
  \(\Delta_d = 2\Delta_\partial = 2\Delta_{\bar{\partial}}\)（Laplacian
  的一致性）
\item
  \([\Lambda, \bar{\partial}] = -i \partial^*\)，\([\Lambda, \partial] = i \bar{\partial}^*\)（其中
  \(\Lambda = L^*\) 是 \(\omega \wedge\) 的伴随）
\item
  \([L, \bar{\partial}^*] = i \partial\)，\([L, \partial^*] = -i \bar{\partial}\)
\end{enumerate}

其中 \(L = \omega \wedge \cdot\) 是 Lefschetz 算子，\(\Lambda\)
是其伴随。

这些恒等式是 Hodge 理论在 Kähler 流形上成立的核心技术原因。

\subsubsection{178.3 Hodge
分解定理}\label{hodge-ux5206ux89e3ux5b9aux7406}

\textbf{定理 178.2}（Hodge 分解定理，Kähler 情形）：紧 Kähler 流形 \(M\)
上的 de Rham 上同调（复系数）分解为

\[H^k_{\text{dR}}(M, \mathbb{C}) \cong \bigoplus_{p+q=k} H_{\bar{\partial}}^{p,q}(M)\]

且满足
\(H_{\bar{\partial}}^{p,q}(M) \cong \overline{H_{\bar{\partial}}^{q,p}(M)}\)（复共轭对称）。

\emph{证明概要}（Dolbeault 上同调的 Hodge 分解）：由 Kähler 恒等式
\(\Delta_d = 2\Delta_\partial = 2\Delta_{\bar{\partial}}\)，\(\Delta_d\)-调和形式与
\(\Delta_{\bar{\partial}}\)-调和形式重合。对 \(\Delta_d\) 应用 Hodge
定理：每个 de Rham 上同调类有唯一调和代表元
\(\alpha \in \ker \Delta_d\)。将 \(\alpha\) 按 \((p,q)\)
类型分解：\(\alpha = \sum_{p+q=k} \alpha^{p,q}\)。由 Kähler
恒等式，\(\Delta_d \alpha = 0\) 等价于
\(\Delta_{\bar{\partial}} \alpha^{p,q} = 0\)，故每个 \(\alpha^{p,q}\) 是
\(\bar{\partial}\)-调和形式，对应 Dolbeault 上同调类。这给出同构
\(H^k_{\text{dR}}(M, \mathbb{C}) \cong \bigoplus_{p+q=k} H_{\bar{\partial}}^{p,q}(M)\)。∎

\textbf{推论 178.3}（Hodge
数的对称性）：\(h^{p,q} = h^{q,p} = h^{n-p,n-q} = h^{n-q,n-p}\)。

\textbf{推论 178.4}（Betti 数的奇偶性约束）：紧 Kähler 流形的奇数阶
Betti 数 \(b_{2k+1}\) 为偶数。

\textbf{例 178.2}（Hodge 菱形，K3 曲面）：K3 曲面的 Hodge
菱形（注意：此处展示的是 K3 曲面而非 \(\mathbb{CP}^2\)，其
\(h^{1,1}=20\) 是 K3 曲面的特征）：

\begin{verbatim}
              1
          0       0
      1      20      1
          0       0
              1
\end{verbatim}

\subsubsection{178.4 Chern 类与 Chern-Weil
理论}\label{chern-ux7c7bux4e0e-chern-weil-ux7406ux8bba}

\textbf{定义 178.3}（Chern 类，复几何视角）：复向量丛 \(E \to M\)
的\textbf{第 \(k\) 个 Chern 类} \(c_k(E) \in H^{2k}(M, \mathbb{Z})\)
由曲率形式 \(\Omega\) 通过

\[\det\left(I + \frac{i}{2\pi} \Omega\right) = 1 + c_1(E) + c_2(E) + \cdots\]

定义。对复流形 \(M\)，\(c_k(M) = c_k(TM)\)。

\textbf{定理 178.5}（Chern-Weil 理论）：Chern 类在 de Rham
上同调中不依赖于联络的选择，且是复流形的微分同胚不变量。

\textbf{定义 178.4}（Chern 数）：紧复 \(n\) 维流形 \(M\) 的
\textbf{Chern 数}是
\(\int_M c_{i_1}(M) \wedge \cdots \wedge c_{i_k}(M)\)，其中
\(\sum i_j = n\)。

\bigskip\hrule\bigskip

\subsection{第179章：Hodge
理论}\label{ux7b2c179ux7ae0hodge-ux7406ux8bba}

Hodge 理论是紧 Kähler
流形上同调理论的核心，揭示了调和形式、上同调类与复结构之间的深层关系。

\subsubsection{179.1 调和形式与 Hodge
定理}\label{ux8c03ux548cux5f62ux5f0fux4e0e-hodge-ux5b9aux7406}

\textbf{定义 179.1}（调和形式）：\((p,q)\)-形式 \(\alpha\)
称为\textbf{调和}的，如果 \(\Delta_{\bar{\partial}} \alpha = 0\)。记
\(\mathcal{H}^{p,q}(M)\) 为调和 \((p,q)\)-形式空间。

\textbf{定理 179.1}（Hodge 定理，调和表示）：在紧 Kähler 流形上，每个
Dolbeault 上同调类有唯一的调和代表元：

\[H_{\bar{\partial}}^{p,q}(M) \cong \mathcal{H}^{p,q}(M)\]

且 \(\mathcal{H}^{p,q}(M)\) 是有限维的。

\emph{证明}：\(\Delta_{\bar{\partial}}\)
是椭圆算子（在紧流形上），其核是有限维的。\(\bar{\partial}\)-上同调类与
\(\Delta_{\bar{\partial}}\)-调和形式通过正交分解
\(H_{\bar{\partial}}^{p,q} \cong \ker \Delta_{\bar{\partial}}\)
建立同构。∎

\subsubsection{179.2 Lefschetz 分解}\label{lefschetz-ux5206ux89e3}

\textbf{定义 179.2}（Lefschetz 算子与原始形式）：定义
\(L: \Omega^{p,q} \to \Omega^{p+1,q+1}\) 为
\(L(\alpha) = \omega \wedge \alpha\)（Kähler
形式的外乘）。\((p,q)\)-形式 \(\alpha\)
称为\textbf{原始}的（primitive），如果
\(L^{n-p-q+1} \alpha = 0\)（\(n = \dim_{\mathbb{C}} M\)）。

\textbf{定理 179.2}（Lefschetz 分解，硬 Lefschetz 定理）：对紧 Kähler
流形，\(L^k: H^{n-k}(M, \mathbb{C}) \to H^{n+k}(M, \mathbb{C})\)
是同构（\(0 \leq k \leq n\)）。每个 \(k\)-形式可唯一分解为

\[\alpha = \sum_{j \geq \max(0, k-n)} L^j \alpha_j\]

其中 \(\alpha_j\) 是原始的 \((k-2j)\)-形式。

\emph{注}：硬 Lefschetz 定理是 Kähler
几何区别于一般复流形的最深刻特征之一。

\textbf{定理 179.3}（Lefschetz 超平面定理）：设
\(X \subset \mathbb{CP}^N\) 是光滑射影簇，\(H\) 是超平面截面。则包含映射
\(H \cap X \hookrightarrow X\) 诱导的同调群同态在维数
\(< \dim_{\mathbb{C}} X - 1\) 上是同构。

\subsubsection{179.3 Hodge 菱形与 Hodge
猜想}\label{hodge-ux83f1ux5f62ux4e0e-hodge-ux731cux60f3}

\textbf{定义 179.3}（Hodge 菱形）：Hodge 菱形是 \(h^{p,q}\)
的阵列，行对应 \(p+q\)，列对应 \(p-q\)。Hodge 对称性
\(h^{p,q} = h^{q,p} = h^{n-p,n-q}\) 反映在菱形的对称性上。

\textbf{命题 179.1}（Hodge 指标定理）：紧 Kähler 曲面 \(M\) 的相交形式在
\(H^{1,1}(M) \cap H^2(M, \mathbb{R})\) 上的符号差为
\((1, h^{1,1} - 1)\)。

\textbf{定义 179.4}（Hodge 猜想）：设 \(X\) 是光滑复射影代数簇。对每个
\(k\)，有理 \((k,k)\)-类（即
\(H^{2k}(X, \mathbb{Q}) \cap H^{k,k}(X)\)）是否由代数子簇的类生成？这是
Clay 千禧年问题之一，至今未解决（\(k=1\) 时由 Lefschetz 的
\((1,1)\)-定理解决）。

\textbf{定理 179.4}（Lefschetz \((1,1)\)-定理）：紧 Kähler 流形 \(M\)
上的每个整 \((1,1)\)-类是线丛的 Chern 类。即
\(H^{1,1}(M) \cap H^2(M, \mathbb{Z}) \cong \operatorname{Pic}(M)\)（Néron-Severi
群）。

\subsubsection{179.4 Hodge
结构的变分}\label{hodge-ux7ed3ux6784ux7684ux53d8ux5206}

\textbf{定义 179.5}（Hodge 结构）：权 \(k\) 的\textbf{有理 Hodge
结构}由有限维有理向量空间 \(V_{\mathbb{Q}}\) 和
\(V_{\mathbb{C}} = V_{\mathbb{Q}} \otimes \mathbb{C}\) 上的分解
\(V_{\mathbb{C}} = \bigoplus_{p+q=k} V^{p,q}\) 组成，满足
\(\overline{V^{p,q}} = V^{q,p}\)。等价地，Hodge
结构可用两个滤过刻画：\textbf{Hodge 滤过}
\(F^p = \bigoplus_{i \geq p} V^{i,k-i}\) 是 \(V_{\mathbb{C}}\)
的下降滤过，满足
\(F^p \oplus \overline{F^{k-p+1}} \cong V_{\mathbb{C}}\)；\textbf{权滤过}
\(W_\bullet\) 是定义在 \(V_{\mathbb{Q}}\) 上的递增滤过，满足
\(W_m = 0\)（\(m < 0\)）且 \(W_k = V_{\mathbb{Q}}\)。称
\((V_{\mathbb{Q}}, F^\bullet, W_\bullet)\) 为\textbf{混合 Hodge
结构}。当 \(W_{k-1} = 0\) 时即为\textbf{纯 Hodge 结构}。纯 Hodge
结构的基本例子包括：光滑复射影代数簇 \(X\) 的 \(k\) 阶上同调
\(H^k(X, \mathbb{Q})\) 带由其 Kähler 结构诱导的 Hodge 分解
\(H^k(X, \mathbb{C}) \cong \bigoplus_{p+q=k} H^{p,q}(X)\)；以及 Abel
簇（复环面）的 \(H^1\) 上的极化的纯 Hodge 结构。混合 Hodge
结构最典型的例子来自代数簇的补空间或奇异代数簇的上同调------Deligne
证明任何复代数簇的上同调都携带典范的混合 Hodge
结构，其权滤过反映了奇异性的贡献。

\textbf{定理 179.5}（Griffiths 横截性）：在 Hodge 簇的形变中，Hodge
分解满足 \(V^{p,q}\) 的微分位于
\(V^{p-1,q+1} \oplus V^{p,q} \oplus V^{p+1,q-1}\)
中。这是周期映射的基本性质。

\bigskip\hrule\bigskip

\subsection{第180章：形变理论}\label{ux7b2c180ux7ae0ux5f62ux53d8ux7406ux8bba}

形变理论研究复流形在保持复结构前提下的连续变化。由 Kodaira 和 Spencer 在
1950 年代创立。

\subsubsection{180.1 无穷小形变}\label{ux65e0ux7a77ux5c0fux5f62ux53d8}

\textbf{定义 180.1}（形变）：紧复流形 \(M\) 的\textbf{形变}是光滑族
\(\pi: \mathcal{M} \to B\)（\(B\) 是连通复流形），使得
\(\pi^{-1}(0) = M\)（\(0 \in B\)）。\(M_t = \pi^{-1}(t)\)
是形变后的复流形。

\textbf{定义 180.2}（无穷小形变）：\(M\) 的\textbf{无穷小形变}空间是
\(H^1(M, \Theta_M)\)，其中 \(\Theta_M = \mathcal{O}(TM)\)
是全纯切丛的层。

\textbf{定理 180.1}（Kodaira-Spencer 映射）：设
\(\pi: \mathcal{M} \to B\) 是 \(M\) 的形变。则存在 Kodaira-Spencer 映射
\(\rho: T_0 B \to H^1(M, \Theta_M)\)，将形变方向映射到无穷小形变类。

\textbf{例 180.1}（Riemann 面的形变）：亏格 \(g\) 的 Riemann
面的形变空间维数为 \(3g-3\)（\(g \geq 2\)），即
\(H^1(M, \Theta_M) \cong \mathbb{C}^{3g-3}\)。这对应于 Teichmüller
空间。

\subsubsection{180.2 阻碍与 Kuranishi
族}\label{ux963bux788dux4e0e-kuranishi-ux65cf}

\textbf{定理 180.2}（Kodaira 阻碍定理）：形变存在阻碍，阻碍空间为
\(H^2(M, \Theta_M)\)。若
\(H^2(M, \Theta_M) = 0\)，则每个无穷小形变可积为真正的形变（\(M\)
无阻碍）。

\textbf{定理 180.3}（Kuranishi 定理，1965）：紧复流形 \(M\)
存在\textbf{半通用形变}（Kuranishi 族），其参数空间是
\(H^1(M, \Theta_M)\) 在 \(0\) 处的一个邻域中由 \(H^2(M, \Theta_M)\)
定义的解析子集。即形变由

\[\mu: H^1(M, \Theta_M) \to H^2(M, \Theta_M), \quad \mu(\varphi) = [\varphi, \varphi]\]

的零点定义（Kuranishi 映射 \(\mu\) 是二次的）。

\textbf{例 180.2}（Calabi-Yau 流形的形变）：Calabi-Yau 3
维流形（\(c_1 = 0\)）的无穷小形变空间
\(H^1(M, \Theta_M) \cong H^{2,1}(M)\) 无阻碍。由 Bogomolov-Tian-Todorov
定理，\(H^2(M, \Theta_M) = 0\)，形变空间是光滑的。

\subsubsection{180.3 形变量子化}\label{ux5f62ux53d8ux91cfux5b50ux5316}

\textbf{定义 180.3}（形变量子化）：形变量子化将光滑函数代数（Poisson
代数）形变为非交换代数。对辛流形
\((M, \omega)\)，形变量子化是形式幂级数环 \(\mathbb{C}[[\hbar]]\)
上的结合代数结构：

\[f \star g = fg + \frac{i\hbar}{2}\{f, g\} + \sum_{k \geq 2} \hbar^k B_k(f, g)\]

其中 \(B_k\) 是双微分算子，\(\star\) 满足结合律，且
\(f \star g - g \star f = i\hbar\{f, g\} + O(\hbar^2)\)。

\textbf{定理 180.4}（Kontsevich 形变量子化定理，1997）：任意 Poisson
流形上存在形变量子化。

\bigskip\hrule\bigskip

\subsection{第181章：Calabi-Yau
流形与镜对称}\label{ux7b2c181ux7ae0calabi-yau-ux6d41ux5f62ux4e0eux955cux5bf9ux79f0}

Calabi-Yau 流形是第一 Chern 类为零的紧 Kähler
流形，在弦论和镜对称中占据核心地位。

\subsubsection{181.1 Calabi 猜想与 Yau
定理}\label{calabi-ux731cux60f3ux4e0e-yau-ux5b9aux7406}

\textbf{定义 181.1}（Calabi-Yau 流形）：紧 Kähler 流形 \(M\) 称为
\textbf{Calabi-Yau 流形}，如果
\(c_1(M) = 0\)（在实上同调中）。等价地，\(M\) 的典范丛
\(K_M = \bigwedge^n T^*M\) 是平凡的。

\textbf{定理 181.1}（Calabi 猜想 / Yau 定理，1976-1978）：设 \(M\) 是紧
Kähler 流形，\(c_1(M) = 0\)。则 \(M\) 上存在唯一的 Ricci 平坦 Kähler
度量，其 Kähler 类等于任意给定的 Kähler 类。

\emph{意义}：Yau 定理证明了 Calabi 猜想，是 20
世纪几何学最伟大的成就之一。它保证了 Calabi-Yau 流形上 Ricci 平坦 Kähler
度量的存在性，为弦论中的 Calabi-Yau 紧化提供了数学基础。Yau 因此获
Fields 奖（1982）。

\textbf{注 181.1}（Yau 定理的证明概要）：Yau 的证明将问题转化为求解复
Monge-Ampere 方程：给定 Kähler 形式 \(\omega_0\) 和体积形式
\(\Omega\)（满足 \(\int_M \Omega = \int_M \omega_0^n\)），求 Kähler 势
\(\varphi\) 使得
\((\omega_0 + i\partial\bar{\partial}\varphi)^n = \Omega\) 且
\(\omega_0 + i\partial\bar{\partial}\varphi > 0\)。这是一个完全非线性的椭圆
PDE。Yau 使用\textbf{连续性方法}：构造从 \(\omega_0^n\) 到 \(\Omega\)
的一族体积形式，证明解的集合既开又闭。开性由椭圆算子的可逆性保证；闭性需要建立解的
\(C^0\)、\(C^2\) 和 \(C^{2,\alpha}\) 先验估计。其中 \(C^0\) 估计（即
\(\varphi\) 的有界性）是最困难的一步，Yau 通过精巧的 Moser
迭代和极大值原理完成。该定理的直接推论是：当 \(c_1(M) = 0\) 时，取
\(\Omega\) 为零曲率体积形式，得到的 Kähler 度量满足
\(\operatorname{Ric}(g) = 0\)，即 Ricci 平坦。

\textbf{例 181.1}（Calabi-Yau 流形的例子）： - 复 1
维：椭圆曲线（\(T^2\)） - 复 2 维：K3
曲面（唯一的单连通紧复曲面，\(c_1 = 0\)） - 复 3 维：\textbf{五次 3
维流形}（\(\mathbb{CP}^4\) 中由五次齐次多项式定义的超曲面）

\subsubsection{181.2 Calabi-Yau 3
维流形}\label{calabi-yau-3-ux7ef4ux6d41ux5f62}

\textbf{定义 181.2}（Hodge 菱形，CY 3 维流形）：Calabi-Yau 3 维流形的
Hodge 数为：

\[
\begin{array}{ccccccc}
& & & 1 & & & \\
& & 0 & & 0 & & \\
& 0 & & h^{1,1} & & 0 & \\
1 & & h^{2,1} & & h^{2,1} & & 1 \\
& 0 & & h^{1,1} & & 0 & \\
& & 0 & & 0 & & \\
& & & 1 & & &
\end{array}
\]

其中 \(h^{1,1}\) 和 \(h^{2,1}\) 是仅有的独立 Hodge 数。\(h^{1,1}\) 是
Kähler 模空间维数，\(h^{2,1}\) 是复结构模空间维数。

\textbf{定义 181.3}（五次 3 维流形）：\(\mathbb{CP}^4\)
中一般五次超曲面的 Hodge 数为 \(h^{1,1} = 1\)，\(h^{2,1} = 101\)。Euler
示性数 \(\chi = -200\)。这是研究最多的 Calabi-Yau 3 维流形。

\subsubsection{181.3 镜对称}\label{ux955cux5bf9ux79f0}

\textbf{定义 181.4}（镜对称猜想）：对每个 Calabi-Yau 3 维流形
\(M\)，存在\textbf{镜流形} \(M^\vee\)（也是 Calabi-Yau 3
维流形），满足：

\[h^{1,1}(M) = h^{2,1}(M^\vee), \quad h^{2,1}(M) = h^{1,1}(M^\vee)\]

即 Hodge 菱形沿 \(45^\circ\) 线翻转。更深层的镜对称断言 \(M\) 上的
\textbf{A-模型}（Gromov-Witten 不变量，与 Kähler 模空间相关）等价于
\(M^\vee\) 上的 \textbf{B-模型}（周期积分，与复结构模空间相关）。

\textbf{定理 181.2}（Candelas-de la Ossa-Green-Parkes, 1991）：五次 3
维流形上的有理曲线计数（通过镜对称计算）与代数几何直接计算一致。这是镜对称的第一个数值验证，开启了枚举几何的新纪元。

\textbf{定义 181.5}（Gromov-Witten 不变量）：\(M\) 的
\textbf{Gromov-Witten 不变量}计算 \(M\)
中给定亏格和同调类的全纯曲线的''虚拟''数目。它们是辛几何的不变量，在镜对称中对应
B-模型的周期积分。

\textbf{定理 181.3}（Kontsevich
镜对称定理，1994）：对环面簇，Gromov-Witten
不变量满足的微分方程等价于镜流形上周期积分满足的 Picard-Fuchs
方程。这给出了镜对称在环面簇情形下的精确表述。

\subsubsection{181.4 SYZ 猜想}\label{syz-ux731cux60f3}

\textbf{定理 181.4}（Strominger-Yau-Zaslow 猜想，1996）：镜对称可以通过
Calabi-Yau 流形的 \textbf{T-对偶}实现：\(M\) 和 \(M^\vee\) 都是特殊
Lagrangian 3-环面 \(T^3\)
的纤维丛（在奇异纤维的极限下），且彼此通过交换纤维与基实现镜对称。

SYZ
猜想将镜对称从代数几何的等价性转化为几何构造，是理解镜对称本质的关键框架。

\bigskip\hrule\bigskip

\emph{本卷建立了复几何的核心理论框架：复流形（近复结构、可积性、Dolbeault
上同调）将实微分几何推广到全纯范畴；Kähler 几何（Kähler 度量、Kähler
恒等式、Hodge 分解）是连接复结构、Riemann 度量和辛结构的典范框架；Hodge
理论（调和形式、Lefschetz 分解、Hodge 猜想）揭示了紧 Kähler
流形上同调的精细结构；形变理论（无穷小形变、Kuranishi
族）刻画了复结构的连续变化；Calabi-Yau 流形与镜对称（Yau
定理、镜对称）是复几何中最深刻的结果，连接了几何、代数、物理三大领域。复几何是当代数学几何化的核心语言。}

\newpage

\section{卷三十五：几何分析}\label{ux5377ux4e09ux5341ux4e94ux51e0ux4f55ux5206ux6790}

\begin{quote}
几何分析是微分几何与偏微分方程交汇的前沿领域，用PDE方法研究几何问题，同时从几何中提取PDE的深刻结构。本卷从极小曲面和调和映射出发，深入讨论
Ricci 流（Perelman 证明 Poincaré 猜想的伟大工具）、Yamabe
问题（共形几何的核心），以及平均曲率流和 Yang-Mills 方程等几何 PDE。
\end{quote}

\bigskip\hrule\bigskip

\subsection{第182章：极小曲面}\label{ux7b2c182ux7ae0ux6781ux5c0fux66f2ux9762}

极小曲面是局部面积极小的曲面，是几何变分问题最经典的例子。Plateau
问题（1840年代提出）是极小曲面理论的起源。

\subsubsection{182.1
极小曲面方程}\label{ux6781ux5c0fux66f2ux9762ux65b9ux7a0b}

\textbf{定义 182.1}（面积泛函）：曲面 \(\Sigma\) 的参数化
\(\mathbf{r}: U \to \mathbb{R}^3\) 的\textbf{面积}为

\[A(\mathbf{r}) = \int_U \sqrt{EG - F^2} \, du \, dv\]

其中 \(E, F, G\)
是第一基本形式的系数。\textbf{极小曲面}是面积泛函的临界点。

\textbf{定义 182.2}（平均曲率）：曲面 \(\Sigma\) 的\textbf{平均曲率}
\(H = \frac{1}{2}(\kappa_1 + \kappa_2)\)（主曲率的平均值）。极小曲面满足
\(H \equiv 0\)。

\textbf{定理 182.1}（极小曲面方程）：非参数化曲面 \(z = u(x, y)\)
是极小曲面的充要条件是

\[\operatorname{div}\left(\frac{\nabla u}{\sqrt{1 + |\nabla u|^2}}\right) = 0\]

即

\[(1 + u_y^2)u_{xx} - 2u_x u_y u_{xy} + (1 + u_x^2)u_{yy} = 0\]

这是二阶拟线性椭圆型偏微分方程。

\textbf{定理 182.2}（Bernstein 定理，1915-1927）：\(\mathbb{R}^2\)
上整体定义的极小曲面图（\(z = u(x, y)\)，\(u\) 定义在整个
\(\mathbb{R}^2\) 上）必为平面。即 \(\mathbb{R}^3\)
中整体定义的极小曲面图只能是平面。

\emph{注}：Bernstein 定理在 \(\mathbb{R}^n\)（\(n \leq 7\)）中成立，在
\(\mathbb{R}^8\) 中不成立（Bombieri-De Giorgi-Giusti,
1969）。这是极小曲面理论与几何测度论的深刻联系。

\subsubsection{182.2 Plateau 问题}\label{plateau-ux95eeux9898}

\textbf{定义 182.3}（Plateau 问题）：给定 \(\mathbb{R}^3\) 中的闭曲线
\(\Gamma\)，求以 \(\Gamma\) 为边界的面积最小的曲面。

\textbf{定理 182.3}（Douglas-Radó 解，1930-1931）：对任意可求长 Jordan
曲线 \(\Gamma\)，存在以 \(\Gamma\)
为边界的极小曲面（可能带有分支点）。Douglas 因此获首届 Fields
奖（1936）。

\emph{方法}：Douglas 将问题转化为 Dirichlet 积分（能量泛函）的极小化：

\[E(\mathbf{r}) = \frac{1}{2} \int_U (|\mathbf{r}_u|^2 + |\mathbf{r}_v|^2) \, du \, dv\]

在共形参数化下，\(E(\mathbf{r}) = A(\mathbf{r})\)，极小曲面等价于调和映射。

\textbf{定理 182.4}（共形表示定理）：任何极小曲面
\(\Sigma \subset \mathbb{R}^3\) 存在等温坐标（共形参数化），使得
\(\mathbf{r}: U \to \Sigma\) 满足
\(\mathbf{r}_{uu} + \mathbf{r}_{vv} = 0\)（即分量是调和函数）且
\(|\mathbf{r}_u| = |\mathbf{r}_v|\)，\(\mathbf{r}_u \cdot \mathbf{r}_v = 0\)。

\subsubsection{182.3
极小曲面的例子}\label{ux6781ux5c0fux66f2ux9762ux7684ux4f8bux5b50}

\textbf{例 182.1}（经典极小曲面）： -
\textbf{平面}：\(H = 0\)（平凡极小曲面） -
\textbf{悬链面}（catenoid）：\(x^2 + y^2 = \cosh^2 z\)（唯一的旋转极小曲面，除平面外）
-
\textbf{螺旋面}（helicoid）：\(x \sin z = y \cos z\)（与悬链面通过等距形变相关联）
- \textbf{Enneper 曲面}：带自交的完备极小曲面 - \textbf{Scherk
曲面}：\(z = \log(\cos x / \cos y)\)（周期极小曲面） - \textbf{Costa
曲面}（1982）：亏格 1 的完备嵌入极小曲面（有 3 个端）

\textbf{定理 182.5}（Costa-Hoffman-Meeks
定理）：存在任意高亏格的完备嵌入极小曲面。

\subsubsection{182.4
几何测度论方法}\label{ux51e0ux4f55ux6d4bux5ea6ux8bbaux65b9ux6cd5}

\textbf{定义 182.4}（整变分流）：\textbf{整变分流}（integral
varifold）是几何测度论中广义曲面的概念，允许奇点存在。变分流是面积泛函的临界点。

\textbf{定理 182.6}（Allard
正则性定理）：变分流在几乎处处意义下是光滑的极小曲面。奇点集是闭的且
Hausdorff 维数不超过 \(n-7\)（\(n\) 为外围空间的维数）。

\textbf{定理 182.7}（Almgren
大正则性定理，1983）：面积最小化整变分流在开稠密集上是光滑的。

\bigskip\hrule\bigskip

\subsection{第183章：调和映射}\label{ux7b2c183ux7ae0ux8c03ux548cux6620ux5c04}

调和映射是 Riemann
流形之间的映射，极小化能量泛函。它是极小曲面（调和函数）在流形间的推广。

\subsubsection{183.1
能量泛函与调和映射方程}\label{ux80fdux91cfux6cdbux51fdux4e0eux8c03ux548cux6620ux5c04ux65b9ux7a0b}

\textbf{定义 183.1}（能量泛函）：Riemann 流形 \((M, g)\) 到 \((N, h)\)
的光滑映射 \(f: M \to N\) 的\textbf{能量}（Dirichlet 能量）为

\[E(f) = \frac{1}{2} \int_M \|df\|^2 \, dV_g\]

其中 \(\|df\|^2 = \operatorname{tr}_g(f^*h)\) 是映射的 Hilbert-Schmidt
范数。

\textbf{定义 183.2}（调和映射）：\textbf{调和映射}是能量泛函 \(E(f)\)
的临界点。在局部坐标下，调和映射方程（Eells-Sampson 方程）为

\[\tau(f) = \Delta_g f^\alpha + \Gamma_{\beta\gamma}^\alpha(f) g^{ij} \frac{\partial f^\beta}{\partial x^i} \frac{\partial f^\gamma}{\partial x^j} = 0\]

其中 \(\tau(f)\) 是\textbf{张力场}（tension
field），\(\Gamma_{\beta\gamma}^\alpha\) 是 \(N\) 的 Christoffel 符号。

\textbf{例 183.1}（调和映射的例子）： -
\(N = \mathbb{R}\)：调和映射是调和函数（\(\Delta f = 0\)） -
\(M = S^1\)：调和映射是 \(N\) 中的闭测地线 - 等距浸入：若 \(f\)
是极小浸入则 \(f\) 是共形调和映射 - 全测地映射：\(\nabla df = 0\)
是调和映射的特例

\subsubsection{183.2 Eells-Sampson
定理}\label{eells-sampson-ux5b9aux7406}

\textbf{定理 183.1}（Eells-Sampson 定理，1964）：设 \(M\) 是紧 Riemann
流形，\(N\) 是紧 Riemann
流形且截面曲率非正（\(K_N \leq 0\)）。则任何连续映射 \(f: M \to N\)
同伦于一个调和映射，且该调和映射在同伦类中极小化能量。

\emph{证明思路}：使用热流方法。考虑演化方程
\(\frac{\partial f_t}{\partial t} = \tau(f_t)\)。当 \(K_N \leq 0\)
时，能量沿热流单调递减，且解的导数有界（Bochner
技术）。热流收敛到调和映射。∎

\textbf{定理 183.2}（Bochner 公式，调和映射情形）：对调和映射
\(f: M \to N\)，

\[\frac{1}{2} \Delta \|df\|^2 = \|\nabla df\|^2 + \sum_i \langle df(\operatorname{Ric}^M(e_i)), df(e_i) \rangle - \sum_{i,j} \langle R^N(df(e_i), df(e_j))df(e_j), df(e_i) \rangle\]

当 \(M\) 的 Ricci 曲率非负且 \(N\) 的截面曲率非正时，\(\|df\|^2\)
是下调和的，从而调和映射的能量极小。

\subsubsection{183.3
调和映射的存在性与障碍}\label{ux8c03ux548cux6620ux5c04ux7684ux5b58ux5728ux6027ux4e0eux969cux788d}

\textbf{定理 183.3}（Sacks-Uhlenbeck 定理，1981）：设 \(M\) 是紧 Riemann
面，\(\pi_2(N) = 0\)。则任何连续映射 \(f: M \to N\)
同伦于一个能量极小调和映射。

\emph{方法}：引入 \(\alpha\)-能量
\(E_\alpha(f) = \int_M (1 + \|df\|^2)^\alpha \, dV\)（\(\alpha > 1\)），通过
\(\alpha \to 1\) 的极限获得调和映射。这种方法克服了 Sobolev 嵌入
\(W^{1,2} \not\subset L^\infty\) 的困难。

\textbf{定理 183.4}（Siu 刚性定理，1980）：紧 Kähler
流形之间的调和映射在适当的曲率条件下必为全纯或反全纯的。这是强刚性定理的几何分析证明。

\bigskip\hrule\bigskip

\subsection{第184章：Ricci 流}\label{ux7b2c184ux7ae0ricci-ux6d41}

Ricci 流是 Ricci 曲率驱动的几何演化方程，由 Hamilton 在 1982
年引入，Perelman 在 2002-2003 年用它证明了 Poincaré 猜想和 Thurston
几何化猜想。

\subsubsection{184.1 Ricci 流方程}\label{ricci-ux6d41ux65b9ux7a0b}

\textbf{定义 184.1}（Ricci 流）：Riemann 度量 \(g(t)\) 的 \textbf{Ricci
流}方程为

\[\frac{\partial g_{ij}}{\partial t} = -2 \operatorname{Ric}_{ij}\]

这是度量的非线性扩散方程，类似热方程：曲率在正曲率区域减小，在负曲率区域增大。

\textbf{无量纲化 Ricci
流}（保持体积）：\(\frac{\partial g_{ij}}{\partial t} = -2 \operatorname{Ric}_{ij} + \frac{2}{n} \bar{R} g_{ij}\)，其中
\(\bar{R}\) 是平均标量曲率。

\textbf{命题 184.1}（曲率演化方程）：在 Ricci 流下，曲率张量满足

\[\frac{\partial R_{ijkl}}{\partial t} = \Delta R_{ijkl} + 2(B_{ijkl} - B_{ijlk} + B_{ikjl} - B_{iljk})\]

其中 \(B_{ijkl} = g^{pq} g^{rs} R_{pijr} R_{qkls}\)。标量曲率满足
\(\frac{\partial R}{\partial t} = \Delta R + 2|\operatorname{Ric}|^2\)。

\subsubsection{184.2 Hamilton 的定理}\label{hamilton-ux7684ux5b9aux7406}

\textbf{定理 184.1}（Hamilton 短时间存在性，1982）：对任意紧 Riemann
流形，Ricci 流在短时间 \([0, T)\) 上存在唯一解。

\textbf{定理 184.2}（Hamilton 的 3 维正 Ricci 曲率定理，1982）：设
\(M^3\) 是紧 3 维流形，初始度量具有正 Ricci 曲率。则 Ricci
流（经归一化）收敛到常正截面曲率度量。因此 \(M^3\)
微分同胚于球面空间形式 \(S^3 / \Gamma\)。

\emph{意义}：这是 Hamilton 开创 Ricci 流的动因------证明 3 维 Poincaré
猜想在正 Ricci 曲率下的特例。Hamilton 获 2011 年 Shaw 奖。

\textbf{定理 184.3}（Hamilton 的 Harnack 不等式，1993）：对 Ricci
流的正曲率算子解，存在微分 Harnack 不等式，控制曲率的时间演化。

\textbf{定理 184.4}（Hamilton 的紧致性定理）：Ricci
流解序列在一定的曲率、直径和内射半径有界条件下，具有光滑收敛的子序列。

\subsubsection{184.3 Perelman 的工作与 Poincaré
猜想}\label{perelman-ux7684ux5de5ux4f5cux4e0e-poincaruxe9-ux731cux60f3}

\textbf{定理 184.5}（Perelman 的 \(\mathcal{F}\)-泛函与
\(\mathcal{W}\)-泛函，2002）：Perelman 引入两个关键泛函：

\[\mathcal{F}(g, f) = \int_M (R + |\nabla f|^2) e^{-f} \, dV\]

\[\mathcal{W}(g, f, \tau) = \int_M [\tau(R + |\nabla f|^2) + f - n] (4\pi\tau)^{-n/2} e^{-f} \, dV\]

\(\mathcal{F}\) 沿 Ricci 流（与共轭热方程耦合）单调不减。\(\mathcal{W}\)
沿 Ricci 流（与 \(\tau\) 的适当演化耦合）单调不减。这给出了 Ricci
流的不变量。

\textbf{定理 184.6}（非塌缩定理）：Perelman 证明在 Ricci
流下，曲率有界的区域体积非塌缩（内射半径有下界）。这排除了
Cheeger-Gromov 意义下的''塌缩''现象。

\textbf{定理 184.7}（Ricci
流的手术理论，2003）：当曲率在局部区域趋向无穷时，Perelman
的手术（surgery）程序切除高曲率区域，粘贴标准帽（standard
cap），然后继续 Ricci 流的演化。手术次数有限。

\textbf{定理 184.8}（Poincaré 猜想，Perelman 2002-2003）：任何单连通闭 3
维流形同胚于 \(S^3\)。

\emph{证明概要}：任取 \(M^3\) 上的初始度量，运行 Ricci
流（带手术）。Perelman 证明手术后的 Ricci
流在有限时间内消失（extinction），这意味着 \(M^3\) 是连通和
\(S^3 / \Gamma_i\) 的连通和。由基本群条件，\(M^3 \cong S^3\)。∎

\textbf{定理 184.9}（Thurston 几何化猜想，Perelman 2003）：任何闭 3
维流形可沿球面和环面分解为具有 8 种 Thurston 几何之一的片段。

\emph{意义}：Perelman 因证明 Poincaré 猜想获 2006 年 Fields
奖（但拒绝领奖），是 21 世纪数学第一个重大突破。

\subsubsection{184.4 Ricci
流在复几何中的应用}\label{ricci-ux6d41ux5728ux590dux51e0ux4f55ux4e2dux7684ux5e94ux7528}

\textbf{定理 184.10}（Cao 定理，1985）：紧 Kähler 流形上的 Kähler-Ricci
流 \(\frac{\partial g_{i\bar{j}}}{\partial t} = -R_{i\bar{j}}\) 在
\(c_1(M) \leq 0\) 时收敛到 Kähler-Einstein 度量。

\textbf{定理 184.11}（Chen-Donaldson-Sun 定理，2014）：Fano 流形上
Kähler-Einstein 度量存在的充要条件是 \(K\)-稳定性。这是
Yau-Tian-Donaldson 猜想的证明，使用 Kähler-Ricci 流方法。

\subsubsection{184.5 Cheeger-Gromoll
分裂定理}\label{cheeger-gromoll-ux5206ux88c2ux5b9aux7406}

Cheeger-Gromoll 分裂定理是关于非负 Ricci 曲率完备 Riemann
流形结构的深刻结果，是 Ricci 流理论的几何背景之一。

\textbf{定理 184.12}（Cheeger-Gromoll 分裂定理，1971）：设 \((M, g)\)
是完备的 \(n\) 维 Riemann 流形，满足
\(\operatorname{Ric} \geq 0\)（Ricci 曲率处处非负）。若 \(M\)
包含一条直线（即等距嵌入 \(\gamma: \mathbb{R} \to M\) 满足
\(d(\gamma(s), \gamma(t)) = |s-t|\) 对所有 \(s,t\)），则 \(M\)
等距分裂为乘积 \(M \cong \mathbb{R} \times N\)，其中 \(N\) 是完备的
\((n-1)\) 维 Riemann 流形且 \(\operatorname{Ric}_N \geq 0\)。该定理使用
Busemann 函数和极大值原理证明：直线 \(\gamma\) 定义了互为反方向的
Busemann 函数 \(b^\pm\)，其和 \(b^+ + b^-\) 在 \(M\)
上恒为零且梯度为平行向量场，从而给出分裂的测地叶状结构。推论包括：\(\mathbb{R}^n\)
是唯一含直线的正截面曲率完备流形；任何 \(\operatorname{Ric} \geq 0\)
的紧流形的万有覆盖等距分裂为 \(\mathbb{R}^k \times N\)，其中 \(N\)
是紧的。

\bigskip\hrule\bigskip

\subsection{第185章：Yamabe
问题与共形几何}\label{ux7b2c185ux7ae0yamabe-ux95eeux9898ux4e0eux5171ux5f62ux51e0ux4f55}

Yamabe 问题是在给定 Riemann
度量的共形类中寻找常标量曲率度量的问题。它的解决是几何分析的一大里程碑。

\subsubsection{185.1 Yamabe 问题}\label{yamabe-ux95eeux9898}

\textbf{定义 185.1}（Yamabe 问题）：设 \((M^n, g)\) 是紧 Riemann
流形（\(n \geq 3\)）。在 \(g\) 的共形类
\([g] = \{u^{4/(n-2)} g : u > 0, u \in C^\infty(M)\}\)
中，是否存在常标量曲率的度量？

\textbf{共形变换下的标量曲率}：若 \(\tilde{g} = u^{4/(n-2)} g\)，则

\[\tilde{R} = u^{-(n+2)/(n-2)} \left(Ru - \frac{4(n-1)}{n-2} \Delta_g u\right)\]

令 \(\tilde{R} = \text{const}\)，得到共形 Yamabe 方程：

\[-\frac{4(n-1)}{n-2} \Delta_g u + Ru = \lambda u^{(n+2)/(n-2)}\]

这是临界 Sobolev 指数的半线性椭圆 PDE。

\textbf{定义 185.2}（Yamabe 不变量）：度量 \(g\) 的共形类的
\textbf{Yamabe 不变量}（Yamabe 商）为

\[Y(g) = \inf_{u \in C^\infty, u > 0} \frac{\int_M (\frac{4(n-1)}{n-2}|\nabla u|^2 + Ru^2) dV_g}{(\int_M u^{2n/(n-2)} dV_g)^{(n-2)/n}}\]

\subsubsection{185.2 Yamabe
问题的解决}\label{yamabe-ux95eeux9898ux7684ux89e3ux51b3}

\textbf{定理 185.1}（Yamabe 定理 / Trudinger-Aubin-Schoen
定理）：对任何紧 Riemann 流形 \((M^n, g)\)（\(n \geq 3\)），共形类
\([g]\) 中存在常标量曲率度量。

\emph{历史}：Yamabe（1960）声称证明但有错误------他在证明中未正确处理临界
Sobolev 嵌入的紧性缺失。Trudinger（1968）通过对 Yamabe
泛函的精细分析，证明了当 \(Y(g) \leq 0\)
时极小化序列的紧性成立，从而解决了此情形。Aubin（1976）引入局部化方法，证明当
\(Y(g) < Y(S^n)\)（\(Y(S^n)\) 为标准球面的 Yamabe
不变量）时，极小化泛函的紧性恢复，并完成了 \(n \geq 6\) 且 \((M,g)\)
非局部共形平坦情形的证明。剩余的临界情形------特别是 \(Y(g)=Y(S^n)\)
以及低维的情形------则需要排除''气泡''的形成。Schoen（1984）使用\textbf{正质量定理}（Positive
Mass Theorem）证明了在剩余所有情形中均有 \(Y(g) < Y(S^n)\)（除非
\((M,g)\) 共形等价于 \(S^n\)），从而彻底解决了 Yamabe
问题。该问题的完整解决历时 24 年，涉及几何分析、临界 Sobolev
理论和广义相对论中正质量猜想的深刻交汇。

\textbf{定理 185.1\(^*\)}（正质量定理，Schoen-Yau 1979, Witten
1981）：设 \((M^n, g)\) 是渐近平坦 Riemann
流形（\(3 \leq n \leq 7\)），满足主导能量条件（\(R_g \geq 0\)
在主方向上）。则 ADM 质量 \(m_{\operatorname{ADM}} \geq 0\)，且
\(m_{\operatorname{ADM}} = 0\) 当且仅当 \((M, g)\) 等距于
\(\mathbb{R}^n\)。Schoen-Yau
使用极小曲面（面积泛函的第二变分）证明该定理，Witten 则通过 Dirac
算子的方法给出简洁的旋量证明。此定理是 Schoen 完成 Yamabe
问题证明的关键工具。

\emph{证明思想}：当 Yamabe 不变量 \(Y(g) < Y(S^n)\)
时，极小化序列的紧性成立（Sobolev 嵌入的临界指数困难通过 Aubin
的局部化方法克服）。Schoen 使用正质量定理证明 \(Y(g) \leq Y(S^n)\)
且等号成立当且仅当 \((M, g)\) 共形等价于 \(S^n\)。∎

\subsubsection{185.3
预定标量曲率}\label{ux9884ux5b9aux6807ux91cfux66f2ux7387}

\textbf{定义 185.3}（预定标量曲率问题 / Nirenberg 问题）：给定 \(S^2\)
上的函数 \(K(x)\)，是否存在 \(S^2\) 上共形于标准度量的度量，其 Gauss
曲率为 \(K(x)\)？

\textbf{定理 185.2}（Kazdan-Warner 障碍，1974）：\(S^2\) 上存在 Gauss
曲率 \(K\) 的必要条件是：对任何一阶球谐函数
\(\psi\)，\(\int_{S^2} \nabla K \cdot \nabla \psi \, dV = 0\)。这给出了
\(S^2\) 上预定 Gauss 曲率问题的拓扑障碍。

\textbf{定理 185.3}（Chang-Yang 解，1987）：\(S^2\) 上的 Nirenberg
问题在 \(K(x)\) 为正函数且满足某些非退化条件时解决。

\bigskip\hrule\bigskip

\subsection{第186章：几何 PDE
选讲}\label{ux7b2c186ux7ae0ux51e0ux4f55-pde-ux9009ux8bb2}

本章介绍几何分析中其他重要的几何 PDE：平均曲率流、Yang-Mills
方程和调和映射热流。

\subsubsection{186.1 平均曲率流}\label{ux5e73ux5747ux66f2ux7387ux6d41}

\textbf{定义 186.1}（平均曲率流）：\(n\) 维超曲面
\(M_t \subset \mathbb{R}^{n+1}\) 的\textbf{平均曲率流}（MCF）为

\[\frac{\partial \mathbf{x}}{\partial t} = -H \mathbf{n}\]

其中 \(\mathbf{x}\) 是位置向量，\(H\) 是平均曲率，\(\mathbf{n}\)
是单位法向量。MCF
是面积泛函的梯度流：\(\frac{d}{dt} \operatorname{Area}(M_t) = -\int_{M_t} H^2 dV \leq 0\)。

\textbf{定理 186.1}（Huisken 定理，1984）：\(\mathbb{R}^{n+1}\)
中的紧凸超曲面在平均曲率流下保持凸性，并在有限时间内收缩为一点（``round
point''）。经归一化后收敛到标准球面。

\textbf{定理 186.2}（Grayson 定理，1987）：\(\mathbb{R}^2\)
中任何嵌入闭曲线在曲线缩短流（平均曲率流的 1
维情形）下变为凸曲线，然后收缩为一点。

\textbf{定理 186.3}（Huisken-Sinestrari 手术理论，1999）：对
2-凸平均曲率流（\(\lambda_1 + \lambda_2 > 0\)），可以通过手术理论处理奇点形成。

\subsubsection{186.2 Yang-Mills 方程}\label{yang-mills-ux65b9ux7a0b}

\textbf{定义 186.2}（Yang-Mills 泛函）：设 \(P \to M\) 是主
\(G\)-丛，\(A\) 是联络。\textbf{Yang-Mills 泛函}为

\[\mathcal{YM}(A) = \int_M \|F_A\|^2 \, dV\]

其中 \(F_A = dA + A \wedge A\) 是曲率 2-形式。\textbf{Yang-Mills 方程}是

\[D_A^* F_A = 0\]

（\(D_A\) 是协变外微分）。Yang-Mills 联络是 Yang-Mills 泛函的临界点。

\textbf{定理 186.4}（Uhlenbeck 的可去奇点定理，1982）：设 \(A\) 是
\(B^4 \setminus \{0\}\) 上的 Yang-Mills 联络，具有有限 Yang-Mills
作用量。则 \(A\) 可光滑延拓到整个 \(B^4\)。

\textbf{定理 186.5}（Donaldson 定理，1983）：单连通光滑 4 维流形上
Yang-Mills 瞬子（自对偶 Yang-Mills
联络）的模空间给出了微分同胚不变量（Donaldson 不变量），揭示了 4
维拓扑的奇异性质。

\textbf{定理 186.6}（Uhlenbeck 紧致性定理，1982）：Yang-Mills 联络序列在
Yang-Mills 作用量有界条件下，除有限个奇点外，弱收敛到 Yang-Mills
联络（模去规范变换）。

\subsubsection{186.3
调和映射热流}\label{ux8c03ux548cux6620ux5c04ux70edux6d41}

\textbf{定义 186.3}（调和映射热流）：调和映射的热流方程为

\[\frac{\partial f_t}{\partial t} = \tau(f_t)\]

其中 \(\tau(f_t)\) 是 \(f_t\) 的张力场。这是能量泛函的梯度流。

\textbf{定理 186.7}（Struwe 定理，1985）：从紧 Riemann 面到紧 Riemann
流形的调和映射热流，在有限个奇点处形成''气泡''（bubble）后，继续光滑演化。全局弱解存在。

\textbf{定理 186.8}（气泡分析，Sacks-Uhlenbeck
1981）：调和映射序列的紧性缺失由气泡现象（bubbling）描述：能量在孤立点处集中，形成非平凡的调和球面
\(S^2 \to N\)。

\bigskip\hrule\bigskip

\emph{本卷建立了几何分析的核心框架：极小曲面（面积泛函、Plateau
问题、Bernstein
定理）是几何变分问题最经典的范例；调和映射（能量泛函、Eells-Sampson
定理、气泡分析）将调和函数推广到流形间的映射；Ricci 流（Hamilton
方程、Perelman 的非塌缩和手术理论、Poincaré 猜想的证明）是 21
世纪数学最伟大的成就之一；Yamabe 问题（共形几何、临界 Sobolev
指数、正质量定理）揭示了共形不变量的深层结构；平均曲率流和 Yang-Mills
方程将几何 PDE
的方法推广到更广泛的几何演化问题。几何分析是当代数学最活跃的前沿领域之一。}

\newpage

\section{卷三十六：机器学习数学}\label{ux5377ux4e09ux5341ux516dux673aux5668ux5b66ux4e60ux6570ux5b66}

\begin{quote}
机器学习是当前应用数学增长最快的领域，其数学基础涉及统计（V21）、优化（V26）、泛函分析（V13）和信息论（V28）的深层交叉。本卷从统计学习理论（PAC
学习、VC
维、泛化界）出发，建立神经网络的逼近理论、深度学习的非凸优化理论（随机梯度下降的收敛性、神经正切核），并讨论深度生成模型和强化学习的数学基础。
\end{quote}

\bigskip\hrule\bigskip

\subsection{第187章：统计学习理论}\label{ux7b2c187ux7ae0ux7edfux8ba1ux5b66ux4e60ux7406ux8bba}

统计学习理论是机器学习理论的核心，研究从有限样本中学习的可能性和极限。由
Vapnik 和 Chervonenkis 在 1970 年代创立。

\subsubsection{187.1 PAC 学习框架}\label{pac-ux5b66ux4e60ux6846ux67b6-1}

\textbf{定义 187.1}（PAC 学习，Probably Approximately Correct，Valiant
1984）：设 \(\mathcal{H}\) 是假设类（hypothesis class），\(\mathcal{D}\)
是未知分布。算法 \(\mathcal{A}\) 从 i.i.d. 样本 \(S \sim \mathcal{D}^m\)
中输出 \(h_S \in \mathcal{H}\)。\(\mathcal{H}\) 是 \textbf{PAC
可学习}的，如果存在函数 \(m_{\mathcal{H}}: (0,1)^2 \to \mathbb{N}\)
和算法 \(\mathcal{A}\)，对任意 \(\epsilon, \delta \in (0,1)\) 和任意分布
\(\mathcal{D}\)，当 \(m \geq m_{\mathcal{H}}(\epsilon, \delta)\) 时，

\[\mathbb{P}_{S \sim \mathcal{D}^m}[L_{\mathcal{D}}(h_S) \leq \min_{h \in \mathcal{H}} L_{\mathcal{D}}(h) + \epsilon] \geq 1 - \delta\]

其中
\(L_{\mathcal{D}}(h) = \mathbb{P}_{(x,y) \sim \mathcal{D}}[h(x) \neq y]\)
是泛化误差。

\textbf{定义 187.2}（样本复杂度）：\(m_{\mathcal{H}}(\epsilon, \delta)\)
是 \(\mathcal{H}\) 的\textbf{样本复杂度}------达到
\((\epsilon, \delta)\)-PAC 学习所需的最小样本数。

\textbf{定理 187.1}（有限假设类的 PAC 可学习性）：若 \(\mathcal{H}\)
有限，则
\(m_{\mathcal{H}}(\epsilon, \delta) \leq \frac{\log |\mathcal{H}| + \log(1/\delta)}{\epsilon^2}\)。即
\(\mathcal{H}\) 是 PAC 可学习的，样本复杂度为
\(O(\epsilon^{-2} \log |\mathcal{H}|)\)。

\emph{证明}：使用 Hoeffding 不等式和联合界。对每个
\(h \in \mathcal{H}\)，\(|L_{\mathcal{D}}(h) - L_S(h)| \leq \epsilon\)
的概率 \(\geq 1 - 2e^{-2m\epsilon^2}\)。联合界给出
\(\mathbb{P}[\exists h: |L_{\mathcal{D}}(h) - L_S(h)| > \epsilon] \leq 2|\mathcal{H}| e^{-2m\epsilon^2}\)。令此概率
\(\leq \delta\)，解出
\(m \geq \frac{\log(2|\mathcal{H}|) + \log(1/\delta)}{2\epsilon^2} = O(\epsilon^{-2} \log |\mathcal{H}|)\)。∎

\subsubsection{187.2 VC 维}\label{vc-ux7ef4}

\textbf{定义 187.3}（VC 维，Vapnik-Chervonenkis 维数）：二分类假设类
\(\mathcal{H}\) 的 \textbf{VC 维} \(\operatorname{VCdim}(\mathcal{H})\)
是能被 \(\mathcal{H}\) 打散（shatter）的最大点集的大小。即存在 \(d\)
个点，\(\mathcal{H}\) 在这 \(d\) 个点上能实现所有 \(2^d\) 种标注方式。

\textbf{例 187.1}（VC 维的例子）： - \(\mathbb{R}\)
上的区间：\(\operatorname{VCdim} = 2\) - \(\mathbb{R}^d\)
中的超平面：\(\operatorname{VCdim} = d+1\) - 正弦函数类
\(\{x \mapsto \sin(\omega x) : \omega \in \mathbb{R}\}\)：\(\operatorname{VCdim} = \infty\)

\textbf{定理 187.2}（VC 维基本定理 / Vapnik-Chervonenkis 定理）：假设类
\(\mathcal{H}\) 是 PAC 可学习的当且仅当其 VC 维有限。此时样本复杂度为

\[m_{\mathcal{H}}(\epsilon, \delta) = O\left(\frac{\operatorname{VCdim}(\mathcal{H}) + \log(1/\delta)}{\epsilon^2}\right)\]

\emph{注}：VC 维是判断假设类可学习性的精确刻画。有限 VC
维等价于一致大数定律（uniform convergence）成立。

\textbf{定理 187.3}（Sauer 引理）：设
\(\operatorname{VCdim}(\mathcal{H}) = d\)。则 \(\mathcal{H}\) 在 \(m\)
个点上的增长函数满足

\[\Pi_{\mathcal{H}}(m) \leq \sum_{i=0}^d \binom{m}{i} \leq \left(\frac{em}{d}\right)^d\]

\subsubsection{187.3 Rademacher
复杂度}\label{rademacher-ux590dux6742ux5ea6}

\textbf{定义 187.4}（经验 Rademacher 复杂度）：给定样本
\(S = (z_1, \ldots, z_m)\)，函数类 \(\mathcal{F}\) 的\textbf{经验
Rademacher 复杂度}为

\[\widehat{\mathcal{R}}_S(\mathcal{F}) = \mathbb{E}_{\sigma} \left[\sup_{f \in \mathcal{F}} \frac{1}{m} \sum_{i=1}^m \sigma_i f(z_i)\right]\]

其中 \(\sigma_i \sim \operatorname{Uniform}(\{-1, +1\})\) 是独立
Rademacher 随机变量。

\textbf{定理 187.4}（Rademacher 复杂度泛化界）：对任意
\(\delta > 0\)，以概率至少 \(1 - \delta\)，

\[\sup_{f \in \mathcal{F}} \left|\mathbb{E}[f(Z)] - \frac{1}{m} \sum_{i=1}^m f(z_i)\right| \leq 2 \mathcal{R}_m(\mathcal{F}) + \sqrt{\frac{\log(2/\delta)}{2m}}\]

其中
\(\mathcal{R}_m(\mathcal{F}) = \mathbb{E}_S[\widehat{\mathcal{R}}_S(\mathcal{F})]\)
是 Rademacher 复杂度。Rademacher 复杂度给出了比 VC 维更精细的泛化界。

\textbf{定理 187.5}（Massart 有限类引理）：对有限函数类
\(\mathcal{F}\)，\(\mathcal{R}_m(\mathcal{F}) \leq \sqrt{2 \log |\mathcal{F}| / m}\)。

\subsubsection{187.4
偏差-方差分解}\label{ux504fux5dee-ux65b9ux5deeux5206ux89e3}

\textbf{定义 187.5}（偏差-方差分解）：对平方损失，期望泛化误差分解为

\[\mathbb{E}[(h_S(x) - y)^2] = \underbrace{(\mathbb{E}[h_S(x)] - \mathbb{E}[y|x])^2}_{\text{偏差}^2} + \underbrace{\operatorname{Var}(h_S(x))}_{\text{方差}} + \underbrace{\operatorname{Var}(y|x)}_{\text{不可约误差}}\]

偏差-方差权衡（bias-variance tradeoff）是理解过拟合与欠拟合的核心框架。

\textbf{定理 187.6}（双重下降现象，Belkin et
al.~2019）：在过参数化模型中，测试误差随模型复杂度呈现''双重下降''：先下降（经典偏差-方差权衡），再上升（过拟合），然后再次下降（过参数化区域）。这挑战了经典偏差-方差权衡的直觉。

\bigskip\hrule\bigskip

\subsection{第188章：神经网络逼近理论}\label{ux7b2c188ux7ae0ux795eux7ecfux7f51ux7edcux903cux8fd1ux7406ux8bba}

神经网络逼近理论研究神经网络作为函数逼近器的表达能力。从 Cybenko
的通用逼近定理到深度的优势，这一理论揭示了深度学习的数学基础。

\subsubsection{188.1
通用逼近定理}\label{ux901aux7528ux903cux8fd1ux5b9aux7406}

\textbf{定理 188.1}（Cybenko 通用逼近定理，1989）：设 \(\sigma\)
是连续的非多项式激活函数（如 sigmoid）。则单隐层神经网络
\(\sum_{i=1}^N a_i \sigma(w_i^\top x + b_i)\) 在 \(C([0,1]^d)\)
中稠密（在一致范数下）。

\emph{证明思路}：使用泛函分析的 Hahn-Banach 定理和 Radon
测度。若闭包不是整个空间，则存在非零线性泛函在其中消失，导致矛盾。∎

\textbf{定理 188.2}（Barron 定理，1993）：具有 Fourier 表示的函数
\(f(x) = \int_{\mathbb{R}^d} e^{i\omega \cdot x} \tilde{f}(\omega) d\omega\)，若
\(\int \|\omega\| |\tilde{f}(\omega)| d\omega \leq C\)，则可用 \(N\)
个隐层神经元的网络以 \(O(N^{-1/2})\) 的速率逼近（在 \(L^2\)
范数下），与输入维数 \(d\) 无关。

\emph{意义}：Barron 定理表明，对于 Barron
空间中的函数，神经网络可以克服维数灾难（curse of dimensionality）。

\textbf{定义 188.1}（Barron 空间）：\textbf{Barron 空间} \(\mathcal{B}\)
由满足
\(\|f\|_{\mathcal{B}} = \int_{\mathbb{R}^d} \|\omega\| |\tilde{f}(\omega)| d\omega < \infty\)
的函数组成。神经网络在 Barron 空间上具有 \(O(N^{-1/2})\) 的逼近速率。

\subsubsection{188.2 深度的优势}\label{ux6df1ux5ea6ux7684ux4f18ux52bf}

\textbf{定理 188.3}（深度分离定理，Telgarsky 2016, Eldan-Shamir
2016）：存在可由深度为 \(L\)、宽度为多项式大小的 ReLU
网络表示的函数，但任何深度为 \(L-1\)
的网络需要指数级宽度才能达到相同的逼近精度。

\textbf{定理 188.4}（Poggio 等，2017）：深度 ReLU
网络可以指数级更高效地表示某些函数。具体地，函数
\(f(x_1, \ldots, x_d) = \prod_{i=1}^d x_i\) 可由深度 \(O(\log d)\) 的
ReLU 网络表示，但浅层网络需要指数级宽度。

\textbf{定理 188.5}（Yarotsky 定理，2017）：深度 \(L\) 的 ReLU 网络逼近
\(C^s\) 函数的最优速率为
\(O(L^{-2s/d})\)。深度网络通过组合性（compositionality）实现高效逼近。

\subsubsection{188.3
神经正切核（NTK）}\label{ux795eux7ecfux6b63ux5207ux6838ntk}

\textbf{定义 188.2}（神经正切核，Neural Tangent Kernel，Jacot et
al.~2018）：考虑参数为 \(\theta\) 的神经网络
\(f(x; \theta)\)。\textbf{神经正切核}定义为

\[\Theta(x, x') = \mathbb{E}_{\theta \sim \mathcal{N}(0, I)} \left[\nabla_\theta f(x; \theta)^\top \nabla_\theta f(x'; \theta)\right]\]

在无限宽度极限下，NTK 在训练过程中保持不变。

\textbf{定理 188.6}（NTK 极限下的梯度下降，Jacot et
al.~2018）：在无限宽度极限下，梯度下降训练神经网络等价于关于 NTK
的核回归：

\[f_t(x) \approx \Theta(x, X)^\top \Theta(X, X)^{-1} (I - e^{-t \Theta(X, X)}) Y\]

其中
\(\Theta(X, X)_{ij} = \Theta(x_i, x_j)\)。这给出了神经网络训练的精确解析描述。

\textbf{定理 188.7}（NTK 的谱性质）：NTK
的特征值衰减速率决定了神经网络的学习速度。在均匀分布数据上，NTK
的特征值随频率增加而衰减，因此神经网络倾向于先学习低频函数（频谱偏差，spectral
bias）。

\subsubsection{188.4
卷积神经网络与不变性}\label{ux5377ux79efux795eux7ecfux7f51ux7edcux4e0eux4e0dux53d8ux6027}

\textbf{定义 188.3}（群等变神经网络）：设 \(G\)
是作用在输入空间上的群。神经网络 \(f\) 是 \textbf{\(G\)-等变}的，如果
\(f(g \cdot x) = g \cdot f(x)\) 对所有 \(g \in G\) 成立。CNN
是平移群等变的特例。

\textbf{定理 188.8}（Kondor-Trivedi
定理，2018）：\(G\)-等变神经网络的层操作必须由群卷积（或等价的群作用下的线性约束）给出。这提供了等变网络设计的数学原理。

\bigskip\hrule\bigskip

\subsection{第189章：深度学习的优化}\label{ux7b2c189ux7ae0ux6df1ux5ea6ux5b66ux4e60ux7684ux4f18ux5316}

深度学习中的优化问题是非凸的，但随机梯度下降在实践中效果惊人。本章讨论
SGD 的收敛性、损失景观分析和隐式正则化。

\subsubsection{189.1
随机梯度下降}\label{ux968fux673aux68afux5ea6ux4e0bux964d}

\textbf{定义 189.1}（随机梯度下降，SGD）：设
\(\ell(\theta) = \mathbb{E}_{z \sim \mathcal{D}}[\ell(\theta; z)]\)
是期望风险。SGD 更新为

\[\theta_{t+1} = \theta_t - \eta_t \nabla_\theta \ell(\theta_t; z_t)\]

其中 \(z_t \sim \mathcal{D}\) 是随机采样的小批量样本。\(\eta_t\)
是学习率。

\textbf{定理 189.1}（SGD 的收敛性，凸情形）：若 \(\ell(\theta)\)
是凸函数且梯度有界，\(\eta_t = 1/\sqrt{t}\)，则

\[\mathbb{E}\left[\ell\left(\frac{1}{T}\sum_{t=1}^T \theta_t\right)\right] - \ell(\theta^*) \leq O\left(\frac{\log T}{\sqrt{T}}\right)\]

\textbf{定理 189.2}（SGD 的收敛性，非凸情形）：若 \(\ell(\theta)\)
光滑（\(L\)-光滑），则 SGD 收敛到驻点：

\[\frac{1}{T} \sum_{t=1}^T \mathbb{E}[\|\nabla \ell(\theta_t)\|^2] \leq O\left(\frac{1}{\sqrt{T}}\right)\]

\subsubsection{189.2 损失景观}\label{ux635fux5931ux666fux89c2}

\textbf{定义 189.2}（损失景观，Loss Landscape）：深度神经网络的损失函数
\(\ell(\theta)\)
在高维参数空间中的几何结构。关键特征包括鞍点、局部极小值和全局极小值的连通性。

\textbf{定理 189.3}（Choromanska 等，2015
的随机矩阵模型）：在大规模随机神经网络中，损失景观的局部极小值集中在全局极小值附近，且高损失局部极小值随网络规模增大而消失。

\textbf{定理 189.4}（Kawaguchi
定理，2016）：深度线性网络（无激活函数）的每个局部极小值都是全局极小值。对于非线性网络，该性质在过参数化区域中近似成立。

\textbf{定理 189.5}（Freeman-Bruna 定理，2017）：在过参数化条件下，深度
ReLU 网络的损失景观在全局极小值附近是连通的（mode connectivity）。

\subsubsection{189.3 隐式正则化}\label{ux9690ux5f0fux6b63ux5219ux5316}

\textbf{定义 189.3}（隐式正则化，Implicit Regularization）：SGD
本身（无显式正则化）倾向于收敛到泛化性能好的解。这种现象称为隐式正则化。

\textbf{定理 189.6}（Soudry 等，2018
的隐式偏差）：在可分离数据的线性分类中，梯度下降（带指数尾损失）收敛到最大间隔分类器（等价于
SVM 解）。这解释了 SGD 的隐式正则化效应。

\textbf{定理 189.7}（Gunasekar
等，2017）：在矩阵分解问题中，梯度下降隐式地最小化核范数（nuclear
norm），从而倾向于低秩解。

\textbf{定理 189.8}（Arora
等，2019）：在深度线性网络中，梯度下降隐式地倾向于低秩解，且收敛方向与初始化相关。

\subsubsection{189.4
学习率策略与自适应方法}\label{ux5b66ux4e60ux7387ux7b56ux7565ux4e0eux81eaux9002ux5e94ux65b9ux6cd5}

\textbf{定义 189.4}（Adam 优化器，Kingma-Ba
2014）：\textbf{Adam}（Adaptive Moment
Estimation）使用一阶矩和二阶矩估计自适应调整学习率：

\[m_t = \beta_1 m_{t-1} + (1-\beta_1) g_t, \quad v_t = \beta_2 v_{t-1} + (1-\beta_2) g_t^2\]

\[\theta_{t+1} = \theta_t - \eta \frac{m_t}{\sqrt{v_t} + \epsilon}\]

\textbf{定义
189.5}（学习率热启动与余弦退火）：学习率衰减策略（阶梯衰减、余弦退火、重启）对深度学习泛化性能有显著影响。余弦退火
\(\eta_t = \eta_{\min} + \frac{1}{2}(\eta_{\max} - \eta_{\min})(1 + \cos(\pi t/T))\)
在实践中表现优异。

\bigskip\hrule\bigskip

\subsection{第190章：深度生成模型}\label{ux7b2c190ux7ae0ux6df1ux5ea6ux751fux6210ux6a21ux578b}

深度生成模型学习数据分布的表示，从中生成新样本。本章讨论
GAN、VAE、扩散模型和最优传输的数学基础。

\subsubsection{190.1
生成对抗网络（GAN）}\label{ux751fux6210ux5bf9ux6297ux7f51ux7edcgan}

\textbf{定义 190.1}（GAN，Goodfellow et
al.~2014）：\textbf{生成对抗网络}由生成器 \(G\) 和判别器 \(D\)
组成，通过极小极大博弈训练：

\[\min_G \max_D \mathbb{E}_{x \sim p_{\text{data}}} [\log D(x)] + \mathbb{E}_{z \sim p_z} [\log(1 - D(G(z)))]\]

当 \(D\)
达到最优时，\(D^*(x) = \frac{p_{\text{data}}(x)}{p_{\text{data}}(x) + p_G(x)}\)，GAN
的目标等价于最小化 Jensen-Shannon 散度
\(D_{\text{JS}}(p_{\text{data}} \| p_G)\)。

\textbf{定理 190.1}（GAN 的全局最优性）：在非参数极限下，GAN
的全局最优解满足 \(p_G = p_{\text{data}}\) 且 \(D = 1/2\)。

\textbf{定理 190.2}（Wasserstein GAN，Arjovsky et al.~2017）：WGAN 使用
\textbf{Wasserstein-1 距离}（地球移动距离）代替 JS 散度：

\[W_1(P, Q) = \sup_{\|f\|_{\text{Lip}} \leq 1} \mathbb{E}_{x \sim P}[f(x)] - \mathbb{E}_{x \sim Q}[f(x)]\]

WGAN 通过 Kantorovich-Rubinstein 对偶将问题转化为最优 1-Lipschitz
函数的寻找。WGAN 训练更稳定，且提供了有意义的损失度量。

\subsubsection{190.2
变分自编码器（VAE）}\label{ux53d8ux5206ux81eaux7f16ux7801ux5668vae}

\textbf{定义 190.2}（VAE，Kingma-Welling
2013）：\textbf{变分自编码器}最大化证据下界（ELBO）：

\[\mathcal{L}(\theta, \phi; x) = \mathbb{E}_{q_\phi(z|x)} [\log p_\theta(x|z)] - D_{\text{KL}}(q_\phi(z|x) \| p(z))\]

其中 \(q_\phi(z|x)\) 是编码器（近似后验），\(p_\theta(x|z)\)
是解码器，\(p(z)\) 是先验分布。ELBO 是 \(\log p(x)\) 的下界。

\textbf{定理 190.3}（VAE 与信息论）：ELBO 可重写为

\[\mathcal{L} = \log p(x) - D_{\text{KL}}(q_\phi(z|x) \| p_\theta(z|x)) \leq \log p(x)\]

VAE 的训练等价于在率-失真理论（rate-distortion
theory）框架下学习潜在表示。

\subsubsection{190.3 扩散模型}\label{ux6269ux6563ux6a21ux578b}

\textbf{定义 190.3}（扩散模型，Ho et al.~2020, Song et
al.~2021）：\textbf{扩散模型}通过前向加噪过程和反向去噪过程生成数据：

\begin{itemize}
\tightlist
\item
  \textbf{前向过程}：\(x_t = \sqrt{\alpha_t} x_0 + \sqrt{1 - \alpha_t} \epsilon\)，\(\epsilon \sim \mathcal{N}(0, I)\)
\item
  \textbf{反向过程}：学习去噪函数 \(\epsilon_\theta(x_t, t)\) 预测噪声
\end{itemize}

训练目标为：\(\mathbb{E}_{x_0, \epsilon, t} [\|\epsilon - \epsilon_\theta(x_t, t)\|^2]\)。

\textbf{定理 190.4}（扩散模型与得分匹配，Song-Ermon
2019）：扩散模型等价于学习数据分布的\textbf{得分函数}（score
function）\(\nabla_x \log p_t(x)\)。反向过程是随机微分方程

\[dx = \left[f(x, t) - g(t)^2 \nabla_x \log p_t(x)\right] dt + g(t) d\bar{W}_t\]

（时间逆向的 SDE）。扩散模型是目前最成功的生成模型。

\textbf{定理 190.5}（扩散模型的采样效率，Song et al.~2021）：通过概率流
ODE（probability flow ODE），扩散模型可以确定性采样，大幅减少采样步骤：

\[dx = \left[f(x, t) - \frac{1}{2} g(t)^2 \nabla_x \log p_t(x)\right] dt\]

\subsubsection{190.4 最优传输}\label{ux6700ux4f18ux4f20ux8f93}

\textbf{定义 190.4}（Monge 问题与 Kantorovich 松弛）：给定分布
\(\mu, \nu\)，\textbf{Monge 最优传输问题}是

\[\min_{T: T_\#\mu = \nu} \int_{\mathbb{R}^d} \|x - T(x)\|^2 d\mu(x)\]

\textbf{Kantorovich 松弛}允许概率耦合：

\[\min_{\gamma \in \Pi(\mu, \nu)} \int_{\mathbb{R}^d \times \mathbb{R}^d} \|x - y\|^2 d\gamma(x, y)\]

\textbf{定理 190.6}（Brenier 定理，1991）：若 \(\mu\)
绝对连续，则最优传输映射 \(T\) 存在且唯一，且 \(T = \nabla \varphi\)
是某个凸函数 \(\varphi\) 的梯度。\(\varphi\) 是 Brenier 势。

\textbf{定理 190.7}（Sinkhorn 算法，Cuturi
2013）：加熵正则化的最优传输（Sinkhorn 距离）：

\[W_\varepsilon(\mu, \nu) = \min_{\gamma \in \Pi(\mu, \nu)} \int \|x - y\|^2 d\gamma + \varepsilon \operatorname{KL}(\gamma \| \mu \otimes \nu)\]

可通过 Sinkhorn 迭代高效求解，复杂度 \(O(n^2)\) 而非 \(O(n^3 \log n)\)。

\bigskip\hrule\bigskip

\subsection{第191章：强化学习与信息论}\label{ux7b2c191ux7ae0ux5f3aux5316ux5b66ux4e60ux4e0eux4fe1ux606fux8bba}

强化学习研究智能体在与环境交互中学习最优策略。本章讨论 MDP、Bellman
方程、策略梯度方法和信息瓶颈原理。

\subsubsection{191.1 Markov
决策过程}\label{markov-ux51b3ux7b56ux8fc7ux7a0b}

\textbf{定义 191.1}（Markov 决策过程，MDP）：\textbf{MDP} 由
\((\mathcal{S}, \mathcal{A}, P, R, \gamma)\) 组成： -
\(\mathcal{S}\)：状态空间 - \(\mathcal{A}\)：动作空间 -
\(P(s'|s, a)\)：转移概率 - \(R(s, a)\)：即时奖励 -
\(\gamma \in [0, 1)\)：折扣因子

\textbf{定义 191.2}（价值函数与 Q 函数）：策略 \(\pi(a|s)\)
的\textbf{状态价值函数}和\textbf{动作价值函数}为

\[V^\pi(s) = \mathbb{E}_\pi\left[\sum_{t=0}^\infty \gamma^t R(s_t, a_t) \mid s_0 = s\right]\]

\[Q^\pi(s, a) = \mathbb{E}_\pi\left[\sum_{t=0}^\infty \gamma^t R(s_t, a_t) \mid s_0 = s, a_0 = a\right]\]

\textbf{定理 191.1}（Bellman 方程）：\(V^\pi\) 和 \(Q^\pi\) 满足

\[V^\pi(s) = \mathbb{E}_{a \sim \pi(\cdot|s)}[R(s, a) + \gamma \mathbb{E}_{s' \sim P}[V^\pi(s')]]\]

\[Q^\pi(s, a) = R(s, a) + \gamma \mathbb{E}_{s' \sim P}[\mathbb{E}_{a' \sim \pi}[Q^\pi(s', a')]]\]

\textbf{定理 191.2}（Bellman 最优性方程）：最优价值函数
\(V^*(s) = \max_\pi V^\pi(s)\) 满足

\[V^*(s) = \max_{a \in \mathcal{A}} \left[R(s, a) + \gamma \mathbb{E}_{s' \sim P}[V^*(s')]\right]\]

\subsubsection{191.2
策略梯度方法}\label{ux7b56ux7565ux68afux5ea6ux65b9ux6cd5}

\textbf{定义 191.3}（策略梯度定理，Sutton et al.~1999）：参数化策略
\(\pi_\theta\) 的性能梯度为

\[\nabla_\theta J(\theta) = \mathbb{E}_{\pi_\theta}\left[\nabla_\theta \log \pi_\theta(a|s) \cdot Q^{\pi_\theta}(s, a)\right]\]

\textbf{定理 191.3}（REINFORCE 算法，Williams 1992）：使用 Monte Carlo
回报 \(G_t\) 作为 \(Q^{\pi_\theta}\) 的估计：

\[\nabla_\theta J(\theta) \approx \frac{1}{N} \sum_{i=1}^N \sum_{t=0}^T \nabla_\theta \log \pi_\theta(a_t^i|s_t^i) \cdot G_t^i\]

\textbf{定理 191.4}（Actor-Critic 方法）：Actor-Critic 使用价值函数
\(V_\phi\)（critic）估计优势函数
\(A(s, a) = Q(s, a) - V(s)\)，减少策略梯度的方差：

\[\nabla_\theta J(\theta) \approx \mathbb{E}\left[\nabla_\theta \log \pi_\theta(a|s) \cdot A_\phi(s, a)\right]\]

\textbf{定理 191.5}（PPO / TRPO 的信任域约束）：Proximal Policy
Optimization（PPO，Schulman et
al.~2017）通过裁剪概率比约束策略更新幅度：

\[\mathcal{L}_{\text{PPO}}(\theta) = \mathbb{E}\left[\min\left(r_t(\theta) A_t, \operatorname{clip}(r_t(\theta), 1-\epsilon, 1+\epsilon) A_t\right)\right]\]

其中
\(r_t(\theta) = \pi_\theta(a_t|s_t) / \pi_{\theta_{\text{old}}}(a_t|s_t)\)。

\subsubsection{191.3
信息瓶颈原理}\label{ux4fe1ux606fux74f6ux9888ux539fux7406}

\textbf{定义 191.4}（信息瓶颈，Tishby et
al.~1999）：\textbf{信息瓶颈原理}寻求输入 \(X\) 的压缩表示
\(T\)，使其保留关于目标 \(Y\) 的最大信息：

\[\min_{p(t|x)} I(X; T) - \beta I(T; Y)\]

其中 \(I(\cdot; \cdot)\) 是互信息，\(\beta\) 控制压缩与保留的权衡。

\textbf{定理 191.6}（信息瓶颈与深度学习的联系，Tishby-Zaslavsky
2015）：深度神经网络的学习过程可以被理解为信息瓶颈的迭代优化：隐层逐步压缩输入信息，同时保留关于标签的信息。这解释了深度网络的泛化能力。

\textbf{定理 191.7}（互信息的下界估计，Belghazi et al.~2018 的
MINE）：互信息可通过神经网络优化 Donsker-Varadhan 表示来估计：

\[I(X; Z) = \sup_{f: \Omega \to \mathbb{R}} \mathbb{E}_{P_{XZ}}[f(x, z)] - \log \mathbb{E}_{P_X \otimes P_Z}[e^{f(x, z)}]\]

\subsubsection{191.4
自监督学习与对比学习}\label{ux81eaux76d1ux7763ux5b66ux4e60ux4e0eux5bf9ux6bd4ux5b66ux4e60}

\textbf{定义 191.5}（InfoNCE 损失，van den Oord et
al.~2018）：\textbf{InfoNCE}（Noise Contrastive Estimation for mutual
information）是互信息的下界：

\[\mathcal{L}_{\text{InfoNCE}} = -\mathbb{E}\left[\log \frac{\exp(f(x, x^+)/\tau)}{\sum_{i=1}^K \exp(f(x, x_i)/\tau)}\right]\]

其中 \(\tau\) 是温度参数。InfoNCE 是 SimCLR、MoCo
等对比学习方法的核心损失函数。

\textbf{定理 191.8}（对比学习与互信息最大化）：InfoNCE 损失
\(\leq \log K - I(X; Z)\)，因此最小化 InfoNCE 等价于最大化 \(I(X; Z)\)
的下界。

\bigskip\hrule\bigskip

\emph{本卷建立了机器学习核心的数学理论：统计学习理论（PAC 学习、VC
维、Rademacher
复杂度、泛化界）提供了学习的可能性边界；神经网络逼近理论（通用逼近、深度分离、NTK）揭示了神经网络的表达能力；优化理论（SGD
收敛性、损失景观、隐式正则化）解释了深度学习的训练动力学；深度生成模型（GAN、VAE、扩散模型、最优传输）通过散度最小化学习数据分布；强化学习（MDP、Bellman
方程、策略梯度、信息瓶颈）提供了序列决策的数学框架。机器学习数学是连接理论计算机科学、统计学、优化和信息论的现代交叉领域。}

\newpage

\section{卷三十七：非交换几何}\label{ux5377ux4e09ux5341ux4e03ux975eux4ea4ux6362ux51e0ux4f55}

\begin{quote}
非交换几何由 Alain Connes 在 1980
年代创立，将经典的微分几何推广到非交换代数框架。其核心思想是：空间与函数代数是等价的（Gelfand-Naimark
定理），非交换代数对应''非交换空间''。本卷在 C*-代数（V14）和 von
Neumann
代数（V31）的基础上，建立谱三元组、循环上同调、非交换微分形式和标准模型几何。Connes
因算子代数理论（特别是 III
型因子的分类与结构定理）以及对叶状结构和微分几何的应用获 1982 年 Fields
奖。
\end{quote}

\bigskip\hrule\bigskip

\subsection{第192章：谱三元组}\label{ux7b2c192ux7ae0ux8c31ux4e09ux5143ux7ec4}

谱三元组（spectral triple）是非交换几何的基本对象，它将 Dirac
算子的概念推广到非交换代数。

\subsubsection{192.1 从 Gelfand-Naimark
到非交换空间}\label{ux4ece-gelfand-naimark-ux5230ux975eux4ea4ux6362ux7a7aux95f4}

\textbf{定理 192.1}（Gelfand-Naimark 定理回顾）：紧 Hausdorff 空间 \(X\)
上的连续函数代数 \(C(X)\) 是交换 C\emph{-代数。反之，任何交换 C}-代数
\(\cong C(X)\)（\(X\)
是其谱空间）。这建立了空间与交换代数之间的范畴等价。

\textbf{Connes 的核心洞察}：将''空间''推广为''非交换
C*-代数''，从而''非交换空间''上的几何由非交换代数上的结构定义。

\textbf{定义 192.1}（谱三元组）：\textbf{谱三元组}
\((\mathcal{A}, \mathcal{H}, \mathcal{D})\) 由以下组成： 1.
\(\mathcal{A}\)：*-代数，在 Hilbert 空间 \(\mathcal{H}\) 上有界表示
\(\pi: \mathcal{A} \to \mathcal{B}(\mathcal{H})\) 2.
\(\mathcal{D}\)：\(\mathcal{H}\) 上的自伴算子（未必有界），称为
\textbf{Dirac 算子} 3. \([\mathcal{D}, \pi(a)]\) 对所有
\(a \in \mathcal{A}\) 是有界算子 4. \(\mathcal{D}\)
具有紧预解式：\((\mathcal{D} - \lambda)^{-1}\) 对所有
\(\lambda \notin \mathbb{R}\) 是紧算子

谱三元组是 Connes 非交换几何的基本数据，它由三个对象组成：非交换代数
\(\mathcal{A}\)（取代函数空间）、Hilbert 空间
\(\mathcal{H}\)（表示空间）和 Dirac 算子
\(\mathcal{D}\)（编码度量几何）。这三个要素分别对应了经典几何中''空间上的函数''\,``旋量场空间''和''Dirac
算子''三重角色。条件 (3) 保证了 \(\mathcal{A}\) 与 \(\mathcal{D}\)
的交换子有界------即代数元素在度量意义下是 Lipschitz 连续的；条件 (4)
保证了 \(\mathcal{D}\)
的逆（在某种迹类理想中）具有紧性------对应几何空间具有有限的''测度''。谱三元组的魅力在于，它仅用泛函分析和算子代数语言就完全刻画了
Riemann 自旋流形的全部几何数据：度量、旋量结构、示性类和指标理论皆可由
\((\mathcal{A}, \mathcal{H}, \mathcal{D})\) 恢复。

\textbf{例 192.1}（经典谱三元组）：紧 Riemann 自旋流形 \(M\) 上的 Dirac
算子给出谱三元组 \((C^\infty(M), L^2(S), \mathcal{D}_M)\)，其中 \(S\)
是自旋丛，\(\mathcal{D}_M\) 是 Dirac 算子。

\textbf{定义 192.2}（谱三元组的维数）：谱三元组有\textbf{维数}
\(n\)（或更一般地，有\textbf{维数谱}），如果 \(|\mathcal{D}|^{-1}\) 属于
Dixmier 理想 \(\mathcal{L}^{(n,\infty)}\)，即
\(\mu_k(|\mathcal{D}|^{-1}) = O(k^{-1/n})\)（\(\mu_k\) 是奇异值）。

\subsubsection{192.2 Dirac
算子与非交换度量}\label{dirac-ux7b97ux5b50ux4e0eux975eux4ea4ux6362ux5ea6ux91cf}

\textbf{定理 192.2}（Connes 的距离公式，1989）：在经典情形的谱三元组
\((C^\infty(M), L^2(S), \mathcal{D}_M)\) 上，Riemann 距离被 Dirac
算子完全恢复：

\[d(x, y) = \sup\{|f(x) - f(y)| : f \in C^\infty(M), \|[D, f]\| \leq 1\}\]

这给出了非交换几何中''度量''的纯代数定义！

\textbf{定义 192.3}（非交换度量）：对谱三元组
\((\mathcal{A}, \mathcal{H}, \mathcal{D})\)，\(\mathcal{A}\)
的态空间（纯态）上的距离定义为

\[d(\varphi, \psi) = \sup\{|\varphi(a) - \psi(a)| : a \in \mathcal{A}, \|[D, a]\| \leq 1\}\]

\textbf{例 192.2}（非交换环面）：非交换 2-环面 \(\mathbb{T}_\theta^2\)
的谱三元组由以下数据定义： - \(\mathcal{A} = A_\theta\)（旋转代数，由
\(U, V\) 生成，满足 \(VU = e^{2\pi i\theta} UV\)） -
\(\mathcal{H} = L^2(A_\theta) \otimes \mathbb{C}^2\)（带自旋结构） -
\(\mathcal{D} = \begin{pmatrix} 0 & \partial_1 - i\partial_2 \\ \partial_1 + i\partial_2 & 0 \end{pmatrix}\)

\subsubsection{192.3
正则性与有限性}\label{ux6b63ux5219ux6027ux4e0eux6709ux9650ux6027}

\textbf{定义 192.4}（正则谱三元组）：谱三元组是\textbf{正则}的，如果
\(\mathcal{A}\) 和 \([\mathcal{D}, \mathcal{A}]\) 属于
\(\bigcap_{k \geq 1} \operatorname{Dom}(\delta^k)\)，其中
\(\delta(T) = [|\mathcal{D}|, T]\)。

\textbf{定义 192.5}（有限谱三元组 / Fredholm
模）：谱三元组是\textbf{偶}的（even），如果存在
\(\mathbb{Z}_2\)-分次算子
\(\gamma\)（\(\gamma = \gamma^*\)，\(\gamma^2 = 1\)）使得
\(\gamma \mathcal{D} = -\mathcal{D}\gamma\) 和 \(\gamma a = a\gamma\)
对所有 \(a \in \mathcal{A}\)。否则是\textbf{奇}的（odd）。

\textbf{定义 192.6}（实结构）：谱三元组的\textbf{实结构}是反线性等距
\(J: \mathcal{H} \to \mathcal{H}\)，满足
\(J^2 = \pm 1\)，\(J\mathcal{D} = \pm \mathcal{D}J\)，\(J\gamma = \pm \gamma J\)（偶情形），以及
\([a, Jb^*J^{-1}] = 0\) 对所有 \(a, b \in \mathcal{A}\)（交换性条件）。

8 种可能的符号组合（\(J^2, J\mathcal{D}, J\gamma\) 的符号）对应
KO-维数模 8 的 8 种情形（Bott 周期性）。

\textbf{定理 192.3}（Connes 的谱三元组公理）：实谱三元组满足 7
条公理，包括正则性、有限性、一阶条件（\([[\mathcal{D}, a], Jb^*J^{-1}] = 0\)）、定向性、Poincaré
对偶等。这些公理精确刻画了''自旋流形上的 Dirac 算子''的代数量子化。

\emph{各公理的几何动机}：\textbf{(1) 正则性}：\(\mathcal{A}\) 及
\([\mathcal{D},\mathcal{A}]\) 对 \(\delta(\cdot)=[|\mathcal{D}|,\cdot]\)
的光滑性，对应流形上 \(C^\infty\)
函数在Dirac算子下的光滑作用。\textbf{(2)
有限性}：\(\bigcap_k \operatorname{Dom}(\mathcal{D}^k)\) 是有限投射
\(\mathcal{A}\)-模，推广了自旋丛截面的有限生成性。\textbf{(3)
一阶条件}：\([[\mathcal{D},a], Jb^*J^{-1}]=0\)
确保Dirac算子是一阶微分算子------\([\mathcal{D},a]\)
与右乘作用对易，这是流形上Dirac算子的本质特征。\textbf{(4)
定向性}：存在Hochschild \(n\)-圈 \(c\) 满足
\(\pi(c)=\gamma\)，编码了体积形式的存在性。\textbf{(5)
Poincaré对偶}：Fredholm指标与
\(K\)-理论的配对非退化，对应经典流形上的Poincaré对偶。\textbf{(6)
实结构符号}：\(J^2, J\mathcal{D}, J\gamma\)
的八种符号组合对应KO-维数（Bott周期性模8），反映旋量表示的实结构分类。\textbf{(7)
谱维数}：\(|\mathcal{D}|^{-1} \in \mathcal{L}^{(n,\infty)}\)
指定了非交换空间的维数。∎

\bigskip\hrule\bigskip

\subsection{第193章：循环上同调}\label{ux7b2c193ux7ae0ux5faaux73afux4e0aux540cux8c03}

循环上同调（cyclic cohomology）是 Connes
为非交换空间引入的上同调理论，它推广了 de Rham
上同调，是非交换微分几何的核心工具。

\subsubsection{193.1 Hochschild
与循环上同调}\label{hochschild-ux4e0eux5faaux73afux4e0aux540cux8c03}

\textbf{定义 193.1}（Hochschild 上链）：代数 \(\mathcal{A}\) 上的
\(n\)-\textbf{Hochschild 上链}是 \((n+1)\)-线性泛函
\(\varphi: \mathcal{A}^{n+1} \to \mathbb{C}\)。Hochschild 上边缘算子
\(b\) 为

\[(b\varphi)(a_0, \ldots, a_{n+1}) = \sum_{i=0}^n (-1)^i \varphi(a_0, \ldots, a_i a_{i+1}, \ldots, a_{n+1}) + (-1)^{n+1} \varphi(a_{n+1}a_0, a_1, \ldots, a_n)\]

\(b^2 = 0\)，定义 \textbf{Hochschild 上同调} \(HH^n(\mathcal{A})\)。

\textbf{定义 193.2}（循环上链）：\textbf{循环
\(n\)-上链}是满足循环条件的 Hochschild 上链：

\[\varphi(a_n, a_0, a_1, \ldots, a_{n-1}) = (-1)^n \varphi(a_0, a_1, \ldots, a_n)\]

\textbf{循环上边缘算子} \(B\)（Connes 算子）为

\[(B\varphi)(a_0, \ldots, a_{n-1}) = \sum_{i=0}^{n-1} (-1)^{(n-1)i} \varphi(1, a_i, a_{i+1}, \ldots, a_{i-1})\]

\(B^2 = 0\)，\(bB + Bb = 0\)。\textbf{循环上同调} \(HC^n(\mathcal{A})\)
是 \((b, B)\)-双复形的上同调。

\textbf{定理 193.1}（Connes 的 SBI 长正合列）：Hochschild
上同调和循环上同调之间存在长正合列：

\[\cdots \to HH^n(\mathcal{A}) \xrightarrow{I} HC^n(\mathcal{A}) \xrightarrow{S} HC^{n-2}(\mathcal{A}) \xrightarrow{B} HH^{n+1}(\mathcal{A}) \to \cdots\]

其中 \(S\) 是周期映射，\(I\) 是包含映射，\(B\) 是 Connes 算子。

\subsubsection{193.2 Chern 特征标}\label{chern-ux7279ux5f81ux6807}

\textbf{定义 193.3}（Fredholm 模的 Chern 特征标）：设
\((\mathcal{A}, \mathcal{H}, F)\) 是奇 Fredholm
模（\(F = \operatorname{sign}(\mathcal{D})\)）。其 \textbf{Chern
特征标}是循环上循环 \(\tau_n\)（\(n\) 为奇数）：

\[\tau_n(a_0, \ldots, a_n) = \operatorname{Tr}(F[F, a_0][F, a_1] \cdots [F, a_n])\]

其中 \(\operatorname{Tr}\) 是 Dixmier 迹（或其推广）。

\textbf{定理 193.2}（Chern 特征标与 \(K\)-理论）：Chern 特征标诱导从
\(K\)-理论到循环上同调的配对（Chern-Connes 配对）：

\[\langle [e], [\varphi] \rangle = \frac{1}{n!} \varphi_{2n}(e, e, \ldots, e)\]

其中
\([e] \in K_0(\mathcal{A})\)，\([\varphi] \in HC^{2n}(\mathcal{A})\)。这是非交换几何中指标定理的核心。

\textbf{定理 193.3}（局部指标公式，Connes-Moscovici
1995）：对正则谱三元组，Chern 特征标在循环上同调中的类可由局部公式（涉及
\(\mathcal{D}\) 的预解式和 \(\mathcal{A}\) 的泛函演算）表示。这推广了
Atiyah-Singer 指标定理。

\subsubsection{193.3
周期循环上同调}\label{ux5468ux671fux5faaux73afux4e0aux540cux8c03}

\textbf{定义 193.4}（周期循环上同调）：\textbf{周期循环上同调}
\(HP^*(\mathcal{A})\)（\(* = 0, 1\)）定义为

\[HP^0(\mathcal{A}) = \varinjlim HC^{2n}(\mathcal{A}), \quad HP^1(\mathcal{A}) = \varinjlim HC^{2n+1}(\mathcal{A})\]

其中正向极限由 \(S\) 映射给出。周期循环上同调是 \(K\)-理论的对偶理论。

\textbf{定理 193.4}（Connes 的周期精确性）：设
\(0 \to \mathcal{J} \to \mathcal{A} \to \mathcal{A}/\mathcal{J} \to 0\)
是代数的扩张。则在
\(K\)-理论和周期循环上同调之间存在自然的六项正合列，推广了拓扑
\(K\)-理论的 Bott 周期性。

\bigskip\hrule\bigskip

\subsection{第194章：非交换微分形式}\label{ux7b2c194ux7ae0ux975eux4ea4ux6362ux5faeux5206ux5f62ux5f0f}

非交换微分形式将经典微分几何的外代数推广到非交换代数，是非交换几何中积分和示性类的核心。

\subsubsection{194.1
泛微分分次代数}\label{ux6cdbux5faeux5206ux5206ux6b21ux4ee3ux6570}

\textbf{定义 194.1}（泛微分分次代数）：代数 \(\mathcal{A}\)
的\textbf{泛微分分次代数}（universal differential graded
algebra）\(\Omega^\bullet(\mathcal{A})\) 是 \(\mathcal{A}\)
上的自由分次代数，由 \(\mathcal{A}\)（次数 0）和
\(da\)（\(a \in \mathcal{A}\)，次数 1）生成，满足
\(d(ab) = da \cdot b + a \cdot db\) 和 \(d(1) = 0\)。

\(\Omega^k(\mathcal{A})\) 的元素是 \(a_0 da_1 da_2 \cdots da_k\)
的线性组合。

\textbf{命题 194.1}（泛微分形式与 Hochschild
同调）：\(\Omega^k(\mathcal{A})\) 与 Hochschild 同调
\(HH_k(\mathcal{A})\) 之间有自然联系，由映射

\[a_0 da_1 \cdots da_k \mapsto a_0 \otimes a_1 \otimes \cdots \otimes a_k\]

诱导。

\subsubsection{194.2 非交换积分}\label{ux975eux4ea4ux6362ux79efux5206}

\textbf{定义 194.2}（Dixmier 迹）：设 \(T\) 是紧算子，\(\mu_n(T)\)
为其奇异值。若 \(\frac{1}{\log N} \sum_{n=1}^N \mu_n(T)\) 有极限
\(\lim_{N \to \infty}\)，则此极限为 \(T\) 的 \textbf{Dixmier 迹}
\(\operatorname{Tr}_\omega(T)\)（依赖于广义极限
\(\omega\)）。若该极限存在且与 \(\omega\) 无关，\(T\)
称为\textbf{可测}的。

\textbf{定义 194.3}（非交换积分）：设
\((\mathcal{A}, \mathcal{H}, \mathcal{D})\) 是 \(n\)
维谱三元组。\textbf{非交换积分}（noncommutative integral）定义为

\[\int a \, ds^n = \operatorname{Tr}_\omega(a |\mathcal{D}|^{-n})\]

当 \(a|\mathcal{D}|^{-n}\) 可测时。这推广了经典 Riemann 流形上的积分
\(\int_M f \, dV\)。

\textbf{定理 194.1}（Connes 的迹定理，1988）：对紧 \(n\) 维自旋流形
\(M\)，Dirac 算子 \(\mathcal{D}_M\) 满足

\[\operatorname{Tr}_\omega(f |\mathcal{D}_M|^{-n}) = c_n \int_M f \, dV\]

其中 \(c_n = 2^{[n/2]} \pi^{n/2} / (n \Gamma(n/2))\)。Dixmier
迹恢复了经典的积分！

\textbf{定理 194.2}（非交换留数，Wodzicki 1984）：伪微分算子 \(P\)
的\textbf{非交换留数}（Wodzicki residue）定义为

\[\operatorname{Res}(P) = \operatorname{Tr}_\omega(P |\mathcal{D}|^{-n})\]

\(\operatorname{Res}\) 是迹（trace）------即
\(\operatorname{Res}([A, B]) = 0\)。在经典流形上，\(\operatorname{Res}(P)\)
等于 \(P\) 的符号的非交换留数（Guillemin 留数）。

\subsubsection{194.3
非交换示性类}\label{ux975eux4ea4ux6362ux793aux6027ux7c7b}

\textbf{定义 194.4}（Chern 特征标，非交换形式）：在非交换几何中，投影
\(e \in M_N(\mathcal{A})\) 的 Chern 特征标由循环上循环

\[\operatorname{Ch}_n(e) = \operatorname{Tr}(e - \frac{1}{2}) de \cdots de \quad (2n \text{ 个 } de)\]

给出。这推广了微分几何中的 Chern-Weil 理论。

\textbf{定理 194.3}（非交换指标定理，Connes 1979）：对 \(p\)-可求和
Fredholm 模 \((\mathcal{A}, \mathcal{H}, F)\)，

\[\operatorname{Index}(eFe) = \langle \operatorname{Ch}_*(e), \tau_* \rangle\]

其中 \(\tau_*\) 是 Fredholm 模的 Chern 特征标。这是 Atiyah-Singer
指标定理的非交换推广。

\bigskip\hrule\bigskip

\subsection{第195章：标准模型的非交换几何}\label{ux7b2c195ux7ae0ux6807ux51c6ux6a21ux578bux7684ux975eux4ea4ux6362ux51e0ux4f55}

Connes 与 Lott（1991）和 Connes（1996）发现，粒子物理的标准模型（带
Higgs 场）可以自然地由非交换几何导出。这是非交换几何最壮观的应用。

\subsubsection{195.1
有限几何与内部空间}\label{ux6709ux9650ux51e0ux4f55ux4e0eux5185ux90e8ux7a7aux95f4}

\textbf{定义 195.1}（有限谱三元组）：\textbf{有限谱三元组}
\((\mathcal{A}_F, \mathcal{H}_F, \mathcal{D}_F)\) 的 \(\mathcal{A}_F\)
是有限维代数，\(\mathcal{H}_F\) 是有限维 Hilbert
空间。有限谱三元组描述了''内部空间''（非交换空间）的几何。

\textbf{定义 195.2}（标准模型的有限代数）：标准模型的有限代数为

\[\mathcal{A}_F = \mathbb{C} \oplus \mathbb{H} \oplus M_3(\mathbb{C})\]

（\(\mathbb{H}\) 是四元数代数）。这对应标准模型的规范群
\(U(1) \times SU(2) \times SU(3)\)（模去离散子群）。

\textbf{定义 195.3}（有限 Hilbert 空间）：\(\mathcal{H}_F\)
是一代费米子的 Hilbert 空间（\(N\)
代，\(N=3\)），包含所有费米子场（轻子、夸克，分为左手和右手分量）。

\textbf{定义 195.4}（有限 Dirac 算子）：\(\mathcal{D}_F\) 是
\(\mathcal{H}_F\) 上的 \(N \times N\) 矩阵，其非零项是 Yukawa
耦合矩阵（给出费米子质量矩阵）和 Majorana 质量项（中微子）。

\subsubsection{195.2 谱作用量}\label{ux8c31ux4f5cux7528ux91cf}

\textbf{定义
195.5}（乘积谱三元组）：标准模型的整体谱三元组是经典几何与有限几何的乘积：

\[\mathcal{A} = C^\infty(M) \otimes \mathcal{A}_F, \quad \mathcal{H} = L^2(S) \otimes \mathcal{H}_F, \quad \mathcal{D} = \mathcal{D}_M \otimes 1 + \gamma_5 \otimes \mathcal{D}_F\]

\textbf{定理 195.1}（谱作用量原理，Connes-Chamseddine
1996-1997）：\textbf{谱作用量}（spectral action）定义为

\[S = \operatorname{Tr}_\omega\left(f\left(\frac{\mathcal{D}^2}{\Lambda^2}\right)\right) + \langle \psi, \mathcal{D} \psi \rangle\]

其中 \(f\) 是截断函数（cutoff function），\(\Lambda\)
是能量截断。展开谱作用量得到：

\[S = \int_M \left(\frac{1}{2g^2} F_{\mu\nu}F^{\mu\nu} + \cdots\right) d^4x + \text{Higgs 项} + \text{Yukawa 耦合}\]

即 Einstein-Hilbert 作用量 + Yang-Mills 作用量 + Higgs 作用量 +
费米子作用量------\textbf{整个标准模型}（包括 Higgs
场）从纯几何原理自动导出！

\textbf{定理 195.2}（Higgs 场的几何起源）：在非交换几何中，Higgs
场自然地出现为内部空间的''联络''的分量------它是规范场在非交换方向上的分量。这解释了
Higgs 场的几何本质。

\subsubsection{195.3
质量关系与预言}\label{ux8d28ux91cfux5173ux7cfbux4e0eux9884ux8a00}

\textbf{定理
195.3}（标准模型的约束）：非交换几何方法给出了标准模型参数之间的约束，包括：
- Higgs 质量 \(m_H \approx 125\) GeV（实验值，LHC 2012） - 预测了 Higgs
自耦合常数 \(\lambda\) 与规范耦合常数之间的关系 -
轻微偏离标准模型：中微子必须是 Majorana 粒子，且存在质量

这些预测使得非交换几何标准模型在 LHC 发现 Higgs
粒子（\(m_H \approx 125\) GeV）后需要修正，但整体框架仍极具吸引力。

\bigskip\hrule\bigskip

\subsection{第196章：非交换 Riemann
几何与量子群}\label{ux7b2c196ux7ae0ux975eux4ea4ux6362-riemann-ux51e0ux4f55ux4e0eux91cfux5b50ux7fa4}

非交换 Riemann 几何将经典 Riemann
几何推广到非交换代数，量子群是其中的核心对象。

\subsubsection{196.1 非交换环面}\label{ux975eux4ea4ux6362ux73afux9762}

\textbf{定义 196.1}（非交换环面
\(\mathbb{T}_\theta^n\)）：\textbf{非交换 \(n\)-环面}
\(\mathbb{T}_\theta^n\) 由旋转代数 \(A_\theta^n\) 描述，其生成元
\(U_1, \ldots, U_n\) 满足

\[U_k U_j = e^{2\pi i \theta_{jk}} U_j U_k\]

其中 \(\theta_{jk} = -\theta_{kj} \in \mathbb{R}/\mathbb{Z}\)
是反对称矩阵。\(A_\theta^n\) 是 \(C(\mathbb{T}^n)\) 的形变量子化。

\textbf{定理 196.1}（非交换环面的几何）：当 \(\theta\)
是''无理''矩阵（即 \(\theta_{jk}\) 是无理数）时，\(A_\theta^n\)
是单代数（无真双边理想）。非交换环面的
\(K\)-理论：\(K_0(A_\theta^n) \cong K_1(A_\theta^n) \cong \mathbb{Z}^{2^{n-1}}\)。

\emph{K-理论计算概要}：以 \(n=2\) 为例，\(A_\theta^2\) 的 \(K_0\)
群由投射的等价类生成。利用Pimsner-Voiculescu六项正合列（将非交换环面嵌入为
\(C(\mathbb{T}) \rtimes_\theta \mathbb{Z}\) 的交叉积），得
\(K_0(A_\theta^2) \cong \mathbb{Z} \oplus \mathbb{Z}\)。当 \(\theta\)
为无理数时，迹映射 \(\tau: K_0 \to \mathbb{R}\) 的像为
\(\mathbb{Z} + \theta\mathbb{Z}\)，这是稠密子群。投射的分类由迹和Chern特征标完全确定。一般
\(n\)
维情形，\(K_*(A_\theta^n) \cong \bigwedge^* \mathbb{Z}^n \cong \mathbb{Z}^{2^{n-1}}\)，由生成元
\(U_k\)
的Künneth分解得到。非交换环面的K-理论是Connes非交换微分几何的经典范例。∎

\textbf{定理 196.2}（非交换环面的 Gauss-Bonnet 定理，Connes-Tretkoff
2010）：非交换 2-环面 \(\mathbb{T}_\theta^2\) 上存在 Gauss-Bonnet
定理的推广，使用了非交换共形几何和局部指标公式。

\subsubsection{196.2 量子群}\label{ux91cfux5b50ux7fa4}

\textbf{定义 196.2}（量子群 / Hopf 代数）：\textbf{量子群}（quantum
group）是 Hopf 代数，带有余乘法
\(\Delta: \mathcal{A} \to \mathcal{A} \otimes \mathcal{A}\)、余单位
\(\varepsilon: \mathcal{A} \to \mathbb{C}\) 和对极
\(S: \mathcal{A} \to \mathcal{A}\)，满足相容性条件。

\textbf{定理 196.2.1}（Hopf代数对极的唯一性）：若 \(S, S'\) 均为 Hopf
代数 \(\mathcal{A}\) 的对极映射，则
\(S = S'\)。即对极映射由代数结构和余代数结构唯一确定。

\emph{证明概要}：设 \(S, S'\) 均为对极，对任意
\(a \in \mathcal{A}\)，利用余乘的 Sweedler 记号
\(\Delta(a) = \sum a_{(1)} \otimes a_{(2)}\)。由对极公理
\(m \circ (S \otimes \operatorname{id}) \circ \Delta = u \circ \varepsilon\)
和
\(m \circ (\operatorname{id} \otimes S') \circ \Delta = u \circ \varepsilon\)，计算：
\[S(a) = \sum S(a_{(1)}) \varepsilon(a_{(2)}) = \sum S(a_{(1)}) a_{(2)} S'(a_{(3)}) = \sum \varepsilon(a_{(1)}) S'(a_{(2)}) = S'(a)\]
其中利用了余结合律和余单位性质。故 \(S = S'\)。这一定理表明 Hopf
代数的''取逆''操作（对极）由代数-余代数结构强制确定，如同群中逆元由其乘积和单位唯一确定。∎

\textbf{详解 196.1}（Hopf 代数的完整公理体系）：Hopf 代数
\((\mathcal{A}, m, u, \Delta, \varepsilon, S)\)
同时是代数和余代数，结构映射满足六条公理：(1)
\textbf{余乘法公理}：\(\Delta\) 是代数同态，即
\(\Delta(ab) = \Delta(a)\Delta(b)\)；(2)
\textbf{余单位公理}：\(\varepsilon\) 是代数同态，且
\((\varepsilon \otimes \operatorname{id}) \circ \Delta = (\operatorname{id} \otimes \varepsilon) \circ \Delta = \operatorname{id}\)（左右余单位律）；(3)
\textbf{余结合律}：\((\Delta \otimes \operatorname{id}) \circ \Delta = (\operatorname{id} \otimes \Delta) \circ \Delta\)，即两次余乘的顺序无关；(4)
\textbf{对极公理}：\(S\) 是反代数同态，满足
\(m \circ (S \otimes \operatorname{id}) \circ \Delta = m \circ (\operatorname{id} \otimes S) \circ \Delta = u \circ \varepsilon\)，其中
\(m\) 是乘法，\(u\) 是单位映射；(5)
\textbf{相容性条件}：\(\Delta(1) = 1 \otimes 1\)，\(\varepsilon(1) = 1\)，\(S(1) = 1\)；(6)
若 \(S\)
可逆（常用要求），则对极为对合：\(S \circ S = \operatorname{id}\)。此公理体系将群的三个运算（乘法、单位、取逆）双对偶化为余乘（描述''分裂''）、余单位（描述''删除''）和对极（描述''反转''），经典群的函数代数
\(C(G)\) 是交换 Hopf 代数，而非交换 Hopf 代数即为量子群。当
\(\mathcal{A}\) 是交换 Hopf 代数时，对应经典代数群；当 \(\mathcal{A}\)
是余交换 Hopf 代数时，对应李双代数。

\textbf{定义 196.3}（量子 \(SL(2)\) / \(SL_q(2)\)）：量子
\(SL(2)\)（Drinfeld-Jimbo 1985）由生成元 \(a, b, c, d\) 和关系

\[ab = qba, \quad ac = qca, \quad bd = qdb, \quad cd = qdc\]

\[bc = cb, \quad ad - da = (q - q^{-1})bc, \quad ad - q^{-1}bc = 1\]

定义。当 \(q = 1\) 时恢复经典 \(SL(2)\)。

\textbf{定理 196.3}（量子群与非交换几何）：量子群 \(SL_q(2)\)
上的非交换几何由谱三元组描述，其 Dirac 算子由 Dabrowski
等（2003）构造。这给出了量子群空间上的非交换 Riemann 几何。

\subsubsection{196.3 Podleś
量子球面}\label{podleux15b-ux91cfux5b50ux7403ux9762}

\textbf{定义 196.4}（Podleś 量子球面，1987）：\textbf{Podleś 量子球面}
\(S_q^2\) 是 \(SU_q(2)\) 的齐性空间（量子商空间）。\(C(S_q^2)\) 是
\(C(SU_q(2))\) 的子代数，由 \(SU_q(2)\) 的特定元素生成。

\textbf{定理 196.4}（量子球面的谱三元组，Dabrowski-Sitarz
2003）：\(S_q^2\) 上存在 \(SU_q(2)\)-等变谱三元组，其 Dirac
算子的谱具有明确的 \(q\)-变形形式。这是非交换 Riemann
几何在量子齐性空间上的重要实现。

\subsubsection{196.4
非交换几何的未来方向}\label{ux975eux4ea4ux6362ux51e0ux4f55ux7684ux672aux6765ux65b9ux5411}

\textbf{定义
196.5}（导出非交换几何）：将非交换几何与导出代数几何（V15、V18）结合，研究非交换空间的导出范畴和模空间。这是当前非交换几何的前沿方向。

\textbf{定义
196.6}（非交换几何与拓扑量子场论）：非交换几何中的谱三元组和循环上同调与拓扑量子场论（TQFT）有深层联系，特别是通过
Chern-Simons 理论和辫子群表示。

\textbf{定理 196.5}（Baum-Connes猜想
{[}含等变KK-理论{]}）：对任意局部紧群 \(G\) 和 \(G\)-\(C^*\)-代数
\(A\)，\textbf{组装映射}
\(\mu: K^G_*( \mathcal{E}G; A) \to K_*(C^*_r(G, A))\) 是群同构。其中
\(\mathcal{E}G\) 是 \(G\) 的万有真作用空间。

\emph{陈述与背景}：等变KK-理论 \(KK^G_*(A, B)\) 是 Kasparov 为研究
\(C^*\)-代数上群作用的指标问题引入的双函子，推广了拓扑K-理论。Baum-Connes猜想的左端是等变K-同调（以
\(\mathcal{E}G\) 为空间参数的Bredon同调），右端是约化群
\(C^*\)-代数的K-理论。组装映射将几何/拓扑数据（\(G\) 在 \(\mathcal{E}G\)
上的作用）映射到解析数据（约化交叉积 \(C^*_r(G,A)\)
的K-理论）。该猜想蕴含了Novikov猜想、Gromov-Lawson-Rosenberg正标量曲率猜想等。Higson-Kasparov（2001）对顺从群（amenable
groups）证明了该猜想；Lafforgue（2002）对双曲群证明了它。此猜想是非交换几何与几何拓扑之间最深刻的桥梁之一。∎

\bigskip\hrule\bigskip

\emph{本卷建立了非交换几何的核心理论框架：谱三元组（\(\mathcal{A}, \mathcal{H}, \mathcal{D}\)）将
Dirac 算子推广到非交换代数，Connes
距离公式恢复了度量；循环上同调（Hochschild/循环上同调、Chern
特征标、局部指标公式）是非交换空间的上同调理论；非交换微分形式（泛微分分次代数、Dixmier
迹、非交换积分）为非交换几何提供了积分和示性类；标准模型的非交换几何是
Connes 理论最壮观的应用，从纯几何原理导出了整个标准模型（包括 Higgs
场）；非交换环面、量子群和量子球面展示了非交换空间的具体例子。非交换几何是当代数学最深刻和最具野心的统一理论之一。}

\newpage

\section{卷三十八：低维拓扑与纽结理论}\label{ux5377ux4e09ux5341ux516bux4f4eux7ef4ux62d3ux6251ux4e0eux7ebdux7ed3ux7406ux8bba}

\begin{quote}
低维拓扑研究 3 维和 4
维流形的拓扑结构，是当代拓扑学中最活跃的前沿领域之一。纽结理论是 3
维拓扑的核心组成部分，通过代数和组合不变量（Jones 多项式、Khovanov
同调）来区分和分类纽结。Floer 同调为 3 维和 4
维流形提供了强大的不变量，而 Donaldson 理论和 Seiberg-Witten 理论揭示了
4
维流形独特的微分结构现象。本卷在拓扑学（V10）、代数拓扑（V14）和微分几何（V11）的基础上，系统介绍低维拓扑的核心理论。
\end{quote}

\bigskip\hrule\bigskip

\subsection{第217章：三维流形基础}\label{ux7b2c217ux7ae0ux4e09ux7ef4ux6d41ux5f62ux57faux7840}

3 维流形是低维拓扑的核心研究对象。与高维不同，3
维流形具有丰富的几何结构（Thurston 的几何化纲领）和组合结构（Heegaard
分解、Dehn 手术）。

\subsubsection{217.1
三维流形的基本构造}\label{ux4e09ux7ef4ux6d41ux5f62ux7684ux57faux672cux6784ux9020}

\textbf{定义 217.1}（Heegaard 分解）：亏格 \(g\) 的 \textbf{Heegaard
分解}是一个闭 3 维流形 \(M\) 的表示为 \(M = H_g \cup_\varphi H_g\)，其中
\(H_g\) 是亏格 \(g\)
的柄体（handlebody），\(\varphi: \partial H_g \to \partial H_g\)
是边界曲面 \(\Sigma_g\) 上的同胚。

\textbf{定理 217.1}（Heegaard 定理）：任何闭可定向 3 维流形都存在
Heegaard 分解。Heegaard 分解的亏格是 3 维流形的基本量。

\textbf{定义 217.2}（Dehn 手术）：设 \(K \subset S^3\) 是纽结，\(N(K)\)
是 \(K\) 的管状邻域。沿 \(K\) 的 \textbf{Dehn 手术}是将 \(N(K)\)
挖去，再将实心环面 \(S^1 \times D^2\) 沿 \(\partial N(K)\)
的某条简单闭曲线粘回。所有闭可定向 3 维流形都可通过 \(S^3\) 中某个链环的
Dehn 手术得到（Lickorish-Wallace 定理）。

\textbf{定理 217.2}（Lickorish-Wallace 定理，1960-1962）：任何闭可定向 3
维流形可通过 \(S^3\) 中链环的 \(\pm 1\) Dehn 手术得到。

\subsubsection{217.2
三维流形的分解}\label{ux4e09ux7ef4ux6d41ux5f62ux7684ux5206ux89e3}

\textbf{定理 217.3}（Kneser-Milnor 素分解定理）：每个紧致可定向 3
维流形可唯一分解为素流形的连通和（\(M = M_1 \# M_2 \# \cdots \# M_n\)），其中
\(M_i\) 是闭单连通 3 维流形（同胚于 \(S^3\)）或不可约流形。

\textbf{定理 217.4}（JSJ 分解定理 / Jaco-Shalen-Johannson
1979）：每个不可约 3 维流形可沿本质环面（incompressible
tori）唯一地分解为 Seifert 纤维流形和双曲流形。

\textbf{定义 217.3}（Seifert 纤维流形）：\textbf{Seifert
纤维流形}是局部平凡 \(S^1\) 纤维的 3
维流形，允许有限条奇异纤维。其不变量是
\((g; b; (\alpha_1, \beta_1), \ldots, (\alpha_n, \beta_n))\)（基曲面亏格、Euler
数和奇异纤维类型）。

\textbf{定义 217.4}（双曲 3 维流形）：\textbf{双曲 3
维流形}是容许完备双曲度量（截面曲率恒为 \(-1\)）的 3 维流形。由 Mostow
刚性定理，双曲 3 维流形上的双曲结构唯一（如果存在）。

\textbf{定理 217.5}（Mostow 刚性定理，1968）：若 \(M\) 和 \(N\)
是有限体积的完备双曲 3 维流形，且 \(\pi_1(M) \cong \pi_1(N)\)，则 \(M\)
和 \(N\) 等距。即双曲 3 维流形的拓扑完全决定了其双曲度量。

\subsubsection{217.3 几何化纲领}\label{ux51e0ux4f55ux5316ux7eb2ux9886}

\textbf{定理 217.6}（Thurston 几何化猜想 / Perelman
证明，1982/2003）：任何闭可定向 3
维流形可沿球面和环面唯一地分解为素流形，其中每个素流形恰好容许 8 种
Thurston 几何之一。

\textbf{定义 217.5}（Thurston 的 8
种几何）：\(S^3\)（球面几何），\(\mathbb{E}^3\)（欧氏几何），\(\mathbb{H}^3\)（双曲几何），\(S^2 \times \mathbb{R}\)，\(\mathbb{H}^2 \times \mathbb{R}\)，\(\widetilde{\operatorname{SL}(2, \mathbb{R})}\)（\(\operatorname{SL}(2, \mathbb{R})\)
的泛覆盖），\(\operatorname{Nil}\)（幂零几何），\(\operatorname{Sol}\)（可解几何）。

\textbf{定理 217.7}（Poincaré 猜想 / Perelman 2003）：任何单连通闭 3
维流形同胚于 \(S^3\)。这是几何化猜想的直接推论。

\textbf{定理 217.8}（球面空间形式猜想 / 椭圆化猜想）：任何有限基本群的闭
3 维流形是球面空间形式（即 \(S^3/\Gamma\)，\(\Gamma\) 是有限群自由作用在
\(S^3\) 上）。由 Perelman 证明。

\bigskip\hrule\bigskip

\subsection{第218章：纽结理论基础}\label{ux7b2c218ux7ae0ux7ebdux7ed3ux7406ux8bbaux57faux7840}

纽结理论是研究嵌入 \(S^1 \hookrightarrow S^3\)（纽结）和多个 \(S^1\)
的嵌入（链环）的拓扑学分支。纽结理论的核心问题是区分纽结类型和发现纽结不变量。

\subsubsection{218.1 纽结与链环}\label{ux7ebdux7ed3ux4e0eux94feux73af}

\textbf{定义 218.1}（纽结与链环）：\textbf{纽结}是 \(S^1\) 到
\(S^3\)（或
\(\mathbb{R}^3\)）的光滑嵌入（或局部平坦嵌入）。\textbf{链环}是多个不相交的
\(S^1\) 的嵌入。纽结 \(K\) 的\textbf{补空间}是
\(S^3 \setminus \nu(K)\)（\(\nu(K)\) 是 \(K\) 的管状邻域）。

\textbf{定义 218.2}（纽结的等价 / 合痕）：两个纽结 \(K_0\) 和 \(K_1\)
是\textbf{合痕的}（或等价的），如果存在连续单参数族
\(\{h_t\}_{t \in [0,1]}\) 的 \(\mathbb{R}^3\) 上的同胚，使得
\(h_0 = \operatorname{id}\) 且 \(h_1(K_0) = K_1\)。

\textbf{定理 218.1}（Reidemeister
定理，1927）：两个纽结图表示同一个纽结（模合痕）当且仅当它们可以通过有限次
\textbf{Reidemeister 运动}（I 型、II 型、III 型）相互转化。

\textbf{定义 218.3}（Reidemeister 运动）： -
\textbf{R1}：添加或消除一个扭结（twist） -
\textbf{R2}：对两条交叉的弧线，可滑动一根穿过另一根 -
\textbf{R3}：三条弧线相交时，中间弧线可滑过交叉点

\subsubsection{218.2
经典纽结不变量}\label{ux7ecfux5178ux7ebdux7ed3ux4e0dux53d8ux91cf}

\textbf{定义 218.4}（交叉数）：纽结 \(K\) 的\textbf{交叉数} \(c(K)\) 是
\(K\) 的所有纽结图中交叉点数的最小值。交叉数是纽结不变量。

\textbf{定义 218.5}（三色性 / Fox
染色）：纽结图是\textbf{三色可染的}，如果每条弧可染以三种颜色之一，使得每个交叉点处三条弧要么同色要么全不同色，且至少使用两种颜色。三色可染性是纽结不变量。

\textbf{定义 218.6}（Alexander 多项式，1928）：纽结 \(K\) 的
\textbf{Alexander 多项式} \(\Delta_K(t) \in \mathbb{Z}[t, t^{-1}]\) 是
\(K\) 的补空间无穷循环覆盖的一维同调的不变量。\(\Delta_K(t)\) 是
\(t \leftrightarrow t^{-1}\) 对称的，且 \(\Delta_K(1) = \pm 1\)。

\textbf{定理 218.2}（Alexander 多项式的构造）：从纽结的 Seifert 曲面
\(S\) 的 Seifert 矩阵 \(V\)
出发，\(\Delta_K(t) = \det(t^{1/2} V - t^{-1/2} V^T)\)（模去
\(t^{\pm 1/2}\) 的乘法）。

\emph{不变性证明概要}：Alexander 多项式 \(\Delta_K(t)\) 是 \(K\)
的补空间 \(S^3 \setminus K\) 的无穷循环覆盖的 1 维同调的扭化 Alexander
模的生成多项式。证明其不变性分两步：（1）对任何给定的纽结图构造
Wirtinger 表示并计算 Fox 微积分，得到 Alexander 矩阵，其初等因子给出
\(\Delta_K(t)\)。Reidemeister 运动（R1-R3）下 Alexander
矩阵的初等因子不变，保证 \(\Delta_K(t)\) 是合痕不变量。（2）等价地，通过
Seifert 曲面方法：不同 Seifert 曲面通过 S-等价性关联，Seifert 矩阵在
S-等价下作初等变换（子矩阵扩展或削减），而
\(\det(t^{1/2}V - t^{-1/2}V^T)\) 在这些变换下保持不变（模去
\(t^{\pm 1/2}\)）。∎

\textbf{定义 218.7}（纽结群）：纽结 \(K\) 的\textbf{纽结群}
\(\pi_K = \pi_1(S^3 \setminus K)\) 是 \(K\)
的补空间的基本群。纽结群是强有力的纽结不变量。Wirtinger
表示：从纽结分叉点出发，每个弧段对应一个生成元，每个交叉点对应一个关系。

\textbf{定理 218.3}（Gordon-Luecke
定理，1989）：纽结由它的补空间决定（不依赖于嵌入方式）。即若
\(S^3 \setminus K \cong S^3 \setminus K'\)，则 \(K \cong K'\)。

\subsubsection{218.3 Jones
多项式与量子不变量}\label{jones-ux591aux9879ux5f0fux4e0eux91cfux5b50ux4e0dux53d8ux91cf}

\textbf{定义 218.8}（Jones 多项式，1984）：Vaughan Jones 定义的
\textbf{Jones 多项式} \(V_K(t) \in \mathbb{Z}[t^{1/2}, t^{-1/2}]\)
由以下缆式（skein）关系唯一确定：

\[t^{-1} V_{L_+}(t) - t V_{L_-}(t) = (t^{1/2} - t^{-1/2}) V_{L_0}(t)\]

其中 \(L_+\)、\(L_-\)、\(L_0\) 是仅在交叉点处不同的链环图，且
\(V_{\text{平凡纽结}}(t) = 1\)。

\textbf{定理 218.4}（Jones 多项式的性质）： -
\(V_K(1) = (-2)^{\#K - 1}\)（\(\#K\) 是链环的分量数） - \(V_K(t)\)
是链环不变量（比 Alexander 多项式更强） - Jones
多项式区分了手性（chirality）：存在 \(V_K(t) \neq V_K(t^{-1})\) 的纽结 -
通过 von Neumann 代数（特别是 \(\operatorname{II}_1\)
因子）和辫群表示得到

\textbf{定义 218.9}（HOMFLY-PT
多项式，1985）：Freyd-Yetter、Hoste-Lickorish-Millett 和
Przytycki-Traczyk 独立发现的二变量纽结多项式 \(P_K(a, z)\)，满足：

\[a P_{L_+}(a, z) - a^{-1} P_{L_-}(a, z) = z P_{L_0}(a, z)\]

且 \(P_{\text{平凡纽结}}(a, z) = 1\)。Jones 多项式是
\(a = t^{-1}\)、\(z = t^{1/2} - t^{-1/2}\) 的特例，Alexander 多项式是
\(a = 1\) 的特例。

\textbf{定义 218.10}（Kauffman 括号多项式）：Kauffman 括号
\(\langle D \rangle\) 是未定向链环图 \(D\) 的多项式不变量，满足： -
\(\langle \bigcirc \rangle = 1\)（平凡纽结） -
\(\langle D \sqcup \bigcirc \rangle = (-A^2 - A^{-2}) \langle D \rangle\)
-
\(\langle \text{正交叉} \rangle = A \langle \text{平行} \rangle + A^{-1} \langle \text{反平行} \rangle\)

Jones 多项式通过
\(V_K(t) = (-A^3)^{-w(D)} \langle D \rangle |_{A = t^{-1/4}}\)
得到（\(w(D)\) 是拧数）。

\subsubsection{218.4 Khovanov 同调}\label{khovanov-ux540cux8c03}

\textbf{定义 218.11}（Khovanov 同调，2000）：Mikhail Khovanov 构造了
\textbf{Khovanov 同调} \(\operatorname{Kh}^{i,j}(K)\)，它是 Jones
多项式的「范畴化」（categorification）。其分次 Euler 特征给出 Jones
多项式：

\[\sum_{i,j} (-1)^i q^j \dim \operatorname{Kh}^{i,j}(K) = (q + q^{-1}) V_K(q^2)\]

\textbf{定理 218.5}（Khovanov 同调的性质）：Khovanov 同调是比 Jones
多项式更强的纽结不变量。它区分了具有相同 Jones 多项式但不同 Khovanov
同调的纽结，且具有函子性质（跨越同痕诱导同构）。

\bigskip\hrule\bigskip

\subsection{第219章：量子拓扑不变量}\label{ux7b2c219ux7ae0ux91cfux5b50ux62d3ux6251ux4e0dux53d8ux91cf}

量子拓扑不变量源自数学物理（Witten 的 Chern-Simons
理论）和量子群理论，将纽结不变量推广到 3 维流形不变量。

\subsubsection{219.1 辫群与 Markov
定理}\label{ux8fabux7fa4ux4e0e-markov-ux5b9aux7406}

\textbf{定义 219.1}（辫群 / Artin 辫群）：\(n\) 股辫的 \textbf{Artin
辫群} \(B_n\) 由生成元 \(\sigma_1, \ldots, \sigma_{n-1}\) 和关系定义： -
\(\sigma_i \sigma_j = \sigma_j \sigma_i\)（\(|i - j| \geq 2\)） -
\(\sigma_i \sigma_{i+1} \sigma_i = \sigma_{i+1} \sigma_i \sigma_{i+1}\)（辫关系）

\textbf{定理 219.1}（Alexander
定理，1923）：任何纽结/链环可表示为某个辫的闭包（braid closure）。

\textbf{定理 219.2}（Markov
定理，1936）：两个辫表示同一个链环当且仅当它们可通过 \textbf{Markov
运动}（共轭 \(\sigma_i b \sigma_i^{-1}\) 和稳定化
\(b \leftrightarrow b \sigma_n^{\pm 1} \in B_{n+1}\)）相互转化。

\textbf{定义 219.2}（辫群表示与纽结不变量）：寻找 \(B_n\) 的一族表示
\(\{\rho_n: B_n \to \operatorname{GL}(V_n)\}\) 满足 Markov
迹条件，从而构造链环不变量。Jones 多项式由此构造。

\subsubsection{219.2 量子群与
R-矩阵}\label{ux91cfux5b50ux7fa4ux4e0e-r-ux77e9ux9635}

\textbf{定义 219.3}（量子群 / 量子包络代数）：\textbf{量子群}
\(U_q(\mathfrak{g})\) 是半单李代数 \(\mathfrak{g}\) 的泛包络代数的
\(q\)-形变（Drinfeld-Jimbo，1985）。它是 Hopf 代数，由 Chevalley 生成元
\(E_i, F_i, K_i^{\pm 1}\) 和 \(q\)-Serre 关系定义。

\textbf{定义 219.4}（R-矩阵 / 泛 R-矩阵）：\(U_q(\mathfrak{g})\) 的
\textbf{泛 R-矩阵}
\(\mathcal{R} \in U_q(\mathfrak{g}) \hat{\otimes} U_q(\mathfrak{g})\)
满足： -
\(\mathcal{R} \Delta(x) = \Delta^{\operatorname{op}}(x) \mathcal{R}\)（对所有
\(x \in U_q(\mathfrak{g})\)） -
\((\Delta \otimes \operatorname{id})(\mathcal{R}) = \mathcal{R}_{13} \mathcal{R}_{23}\)
-
\((\operatorname{id} \otimes \Delta)(\mathcal{R}) = \mathcal{R}_{13} \mathcal{R}_{12}\)

R-矩阵给出 Yang-Baxter
方程的解：\(\mathcal{R}_{12} \mathcal{R}_{13} \mathcal{R}_{23} = \mathcal{R}_{23} \mathcal{R}_{13} \mathcal{R}_{12}\)。

\textbf{定理 219.3}（量子群构造辫群表示）：对每个 \(U_q(\mathfrak{g})\)
的有限维表示
\(V\)，\(R = \tau \circ \mathcal{R}|_{V \otimes V}\)（\(\tau\)
是交换映射）满足辫关系。这给出辫群 \(B_n\) 在 \(V^{\otimes n}\)
上的表示，从而通过 Markov 迹得到链环不变量。

\subsubsection{219.3 拓扑量子场论与 Witten
不变量}\label{ux62d3ux6251ux91cfux5b50ux573aux8bbaux4e0e-witten-ux4e0dux53d8ux91cf}

\textbf{定义 219.5}（拓扑量子场论 / TQFT，Atiyah
1988）：\textbf{拓扑量子场论}是一个函子
\(Z: \operatorname{Bord}_n \to \operatorname{Vect}\)，将 \(n\)
维流形范畴映射到向量空间范畴。具体地： - 每个闭 \((n-1)\) 维流形
\(\Sigma\) 对应向量空间 \(Z(\Sigma)\) - 每个 \(n\)
维配边（cobordism）\(M: \Sigma_0 \to \Sigma_1\) 对应线性映射
\(Z(M): Z(\Sigma_0) \to Z(\Sigma_1)\) -
满足粘合公理：\(Z(M_1 \cup_\Sigma M_2) = Z(M_2) \circ Z(M_1)\)

\textbf{定义 219.6}（Witten-Chern-Simons 理论，1989）：Edward Witten 在
Chern-Simons 规范理论（作用量
\(S_{CS}(A) = \frac{k}{4\pi} \int_M \operatorname{Tr}(A \wedge dA + \frac{2}{3} A \wedge A \wedge A)\)）中，路径积分给出
3 维流形 \(M\) 的不变量（Witten 不变量）：

\[Z_k(M) = \int \mathcal{D}A \, e^{i k S_{CS}(A)}\]

其中积分沿 \(M\) 上的所有联络。Witten 不变量推广了 Jones 多项式。

\textbf{定理 219.4}（Reshetikhin-Turaev 不变量，1991）：Reshetikhin 和
Turaev 用量子群严格构造了 Witten 不变量（RT 不变量），通过 Kirby
演算和量子 \(6j\)-符号。

\subsubsection{219.4
范畴化与同调不变量}\label{ux8303ux7574ux5316ux4e0eux540cux8c03ux4e0dux53d8ux91cf}

\textbf{定义 219.7}（Khovanov 同调 /
范畴化原理）：\textbf{范畴化}（categorification）是将数值或多项式不变量提升为同调理论的过程。Khovanov
同调将 Jones 多项式范畴化，\textbf{Ozsváth-Szabó 同调}（Knot Floer
homology）将 Alexander 多项式范畴化。

\textbf{定义 219.8}（Knot Floer 同调 /
\(\widehat{\operatorname{HFK}}\)）：对纽结
\(K \subset S^3\)，\textbf{Knot Floer 同调}
\(\widehat{\operatorname{HFK}}(K)\) 是双分次（Alexander 分次和 Maslov
分次）的有限维向量空间，满足：

\[\sum_{i,j} (-1)^i t^j \dim \widehat{\operatorname{HFK}}_j(K, i) = \Delta_K(t)\]

\bigskip\hrule\bigskip

\subsection{第220章：Floer
同调}\label{ux7b2c220ux7ae0floer-ux540cux8c03}

Floer 同调是 3 维和 4 维流形的无穷维 Morse 同调理论，由 Andreas Floer 在
1980 年代创立。它为 3 维流形提供了强有力的不变量，并为 4
维流形的光滑结构分类提供了核心工具。

\subsubsection{220.1 Morse 同调与 Floer
同调的基本思想}\label{morse-ux540cux8c03ux4e0e-floer-ux540cux8c03ux7684ux57faux672cux601dux60f3}

\textbf{定义 220.1}（Floer 同调的基本构造）：Floer 同调 \(HF(M)\)
是无穷维空间上的 Morse 同调。构造步骤： 1. 定义无穷维流形
\(\mathcal{C}\)（联络空间 / 映射空间） 2. 定义 Morse 函数
\(\mathcal{F}: \mathcal{C} \to \mathbb{R}\)（Chern-Simons 泛函 /
辛作用泛函） 3. 链复形由 \(\mathcal{F}\) 的临界点生成 4.
边界算子计算梯度流线（连接临界点的 PDE 解）的计数

\textbf{定义 220.2}（Floer 同调的数学要求）：需要证明： -
横向性：梯度流线的模空间是光滑流形（通过加微扰） -
紧致性：模空间可紧化（通过 bubbling 分析和 Uhlenbeck 紧致性） -
粘合：边界算子满足 \(\partial^2 = 0\)（模空间边界对应断裂轨道）

\subsubsection{220.2 Instanton Floer
同调}\label{instanton-floer-ux540cux8c03}

\textbf{定义 220.3}（Instanton Floer 同调，Floer 1988）：设 \(Y\) 是闭 3
维流形（整数同调球）。\(\mathcal{C}\) 是 \(Y\) 上
\(\operatorname{SU}(2)\) 联络的规范等价类空间。Chern-Simons 泛函为

\[\operatorname{CS}(A) = \frac{1}{8\pi^2} \int_Y \operatorname{Tr}(A \wedge dA + \frac{2}{3} A \wedge A \wedge A)\]

\textbf{Instanton Floer 同调} \(I_*(Y)\) 由 \(\operatorname{CS}\)
的临界点（平坦联络）生成，边界算子由反自对偶瞬子（ASD
instanton）的连接数给出。

\textbf{定理 220.1}（Instanton Floer 同调的性质）：\(I_*(Y)\) 是
\(\mathbb{Z}/8\mathbb{Z}\) 分次的。对整数同调球
\(Y\)，\(\chi(I_*(Y)) = |H_1(Y; \mathbb{Z})|\)（Casson 不变量）。

\subsubsection{220.3 Heegaard Floer
同调}\label{heegaard-floer-ux540cux8c03}

\textbf{定义 220.4}（Heegaard Floer 同调，Ozsváth-Szabó
2001-2004）：\textbf{Heegaard Floer 同调}
\(\widehat{\operatorname{HF}}(Y)\) 由 \(Y\) 的 Heegaard 分解
\((U, \Sigma_g, V)\) 构造。生成元是 \(g\) 个点在 \(\Sigma_g\)
上的有序组（\(\operatorname{Sym}^g(\Sigma_g)\)
中的交点）。边界算子计算伪全纯圆盘的计数。

\textbf{定理 220.2}（Heegaard Floer
同调的优势）：\(\widehat{\operatorname{HF}}(Y)\) 是有限维
\(\mathbb{Z}/2\mathbb{Z}\) 分次的向量空间，满足： - 计算方便（比
Instanton Floer 同调更易计算） -
\(\chi(\widehat{\operatorname{HF}}(Y)) = \pm |H_1(Y; \mathbb{Z})|\)（若
\(b_1(Y) = 0\)） -
有以下更精细的版本：\(\operatorname{HF}^+(Y)\)、\(\operatorname{HF}^-(Y)\)、\(\operatorname{HF}^\infty(Y)\)

\textbf{定义 220.5}（Knot Floer 同调与 Heegaard Floer 的关系）：Knot
Floer 同调 \(\widehat{\operatorname{HFK}}(K)\) 是 Heegaard Floer
同调的纽结版本，通过标定 Heegaard 图中的纽结轨道得到。

\textbf{定理 220.3}（Heegaard Floer 同调检测属数，Ozsváth-Szabó
2004）：\(\widehat{\operatorname{HF}}(Y)\) 检测 \(Y\) 的 Thurston
半范数。特别地，纽结 \(K\) 的亏格（genus）可由 Knot Floer 同调读取：

\[g(K) = \max\{i : \widehat{\operatorname{HFK}}(K, i) \neq 0\}\]

\subsubsection{220.4 其他 Floer
同调理论}\label{ux5176ux4ed6-floer-ux540cux8c03ux7406ux8bba}

\textbf{定义 220.6}（Monopole Floer 同调 / Seiberg-Witten Floer
同调，Kronheimer-Mrowka 2007）：基于 Seiberg-Witten
方程（而非反自对偶方程）的 Floer 同调。与 Heegaard Floer 同调同构（由
Kutluhan-Lee-Taubes 和 Colin-Ghiggini-Honda 证明）。

\textbf{定义 220.7}（Lagrangian Floer 同调，Floer 1988）：辛流形
\((M, \omega)\) 中两个 Lagrangian 子流形 \(L_0, L_1\)
的交点生成链复形，边界算子由伪全纯带的计数给出。是镜对称（mirror
symmetry）和 Fukaya 范畴的核心工具。

\textbf{定义 220.8}（Sutured Floer 同调，Juhász 2006）：将 Heegaard
Floer 同调推广到带缝合的 3 维流形（sutured
manifolds）。应用于纽结补空间和接触几何。

\bigskip\hrule\bigskip

\subsection{第221章：四维流形}\label{ux7b2c221ux7ae0ux56dbux7ef4ux6d41ux5f62}

4 维流形是微分拓扑中最特殊的维数：在所有维数中，只有 4
维流形可能具有（无穷多个）不同的光滑结构（exotic \(\mathbb{R}^4\)
等）。Donaldson 理论和 Seiberg-Witten 理论是区分 4
维流形光滑结构的核心工具。

\subsubsection{221.1
四维流形的特殊性}\label{ux56dbux7ef4ux6d41ux5f62ux7684ux7279ux6b8aux6027}

\textbf{定理 221.1}（光滑结构的维数依赖）： - 维数
\(n \neq 4\)：紧致拓扑流形只有有限个光滑结构（到 \(n = 3\)
是唯一的，\(n \geq 5\) 由 Kirby-Siebenmann 分类） - 维数
\(n = 4\)：存在无穷多个不同的光滑结构（如 \(\mathbb{R}^4\)
上有不可数无穷个不同的光滑结构）

\textbf{定义 221.1}（相交形式 / Intersection Form）：闭可定向光滑 4
维流形 \(X\) 的\textbf{相交形式}
\(Q_X: H^2(X; \mathbb{Z}) \times H^2(X; \mathbb{Z}) \to \mathbb{Z}\)
定义为
\(Q_X(\alpha, \beta) = \langle \alpha \cup \beta, [X] \rangle\)。\(Q_X\)
是对称双线性型。

\textbf{定理 221.2}（Freedman 定理，1982）：单连通闭拓扑 4
维流形由其相交形式 \(Q_X\) 和 Kirby-Siebenmann 不变量
\(\operatorname{ks}(X) \in \mathbb{Z}/2\mathbb{Z}\)
完全分类。任何幺模对称双线性型都可实现为某单连通拓扑 4
维流形的相交形式。

\subsubsection{221.2 Donaldson 理论}\label{donaldson-ux7406ux8bba}

\textbf{定理 221.3}（Donaldson 定理，1983）：设 \(X\) 是闭光滑单连通 4
维流形。若 \(X\) 的相交形式 \(Q_X\) 是正定的，则 \(Q_X\) 在
\(\mathbb{Z}\) 上等价于对角形式
\(\langle 1 \rangle \oplus \cdots \oplus \langle 1 \rangle\)。即正定相交形式必须是标准的。

\textbf{推论 221.1}（exotic 4 维流形的存在）：Freedman 定理允许 \(E_8\)
形式作为拓扑 4 维流形的相交形式，但由 Donaldson 定理，\(E_8\)
形式不能实现为光滑 4
维流形的相交形式。这证明了存在拓扑等价但非微分同胚的 4 维流形。

\textbf{定义 221.2}（Donaldson 不变量）：Donaldson 通过反自对偶联络（ASD
instanton）的模空间定义了 4 维流形 \(X\) 的\textbf{Donaldson
多项式不变量}
\(D_X: \operatorname{Sym}_*(H_2(X; \mathbb{Z})) \to \mathbb{Q}\)。这些不变量是光滑结构的不变量。

\textbf{定理 221.4}（Donaldson
不变量对光滑结构的敏感性）：存在同胚但不微分同胚的 4 维流形，其
Donaldson 不变量不同。Donaldson 不变量证明了
\(\mathbb{C}P^2 \# k \overline{\mathbb{C}P^2}\)
上存在无穷多个不同的光滑结构。

\subsubsection{221.3 Seiberg-Witten
理论}\label{seiberg-witten-ux7406ux8bba}

\textbf{定义 221.3}（Seiberg-Witten 方程，1994）：设 \(X\) 是闭光滑 4
维流形，带有 \(\operatorname{Spin}^c\) 结构
\(\mathfrak{s}\)。\textbf{Seiberg-Witten 方程}为：

\[\begin{cases} D_A \Phi = 0 \\ F_A^+ = \Phi \otimes \Phi^* - \frac{1}{2} |\Phi|^2 \operatorname{id} \end{cases}\]

其中 \(\Phi\) 是正旋量，\(A\) 是 \(\operatorname{Spin}^c\) 联络，\(D_A\)
是 Dirac 算子，\(F_A^+\) 是 \(A\) 的曲率的自对偶部分。

\textbf{定义 221.4}（Seiberg-Witten 不变量）：Seiberg-Witten
方程的解空间 \(M(\mathfrak{s})\) 是紧致光滑流形。\textbf{Seiberg-Witten
不变量} \(\operatorname{SW}_X(\mathfrak{s})\) 是 \(M(\mathfrak{s})\) 上
Euler 类的积分，是 \(\operatorname{Spin}^c\) 结构 \(\mathfrak{s}\)
的函数。

\textbf{定理 221.5}（Seiberg-Witten 不变量的性质）： - 只有有限个
\(\operatorname{Spin}^c\) 结构 \(\mathfrak{s}\) 满足
\(\operatorname{SW}_X(\mathfrak{s}) \neq 0\)（基本类） - 若 \(X\)
容许正标量曲率度量，则 \(\operatorname{SW}_X(\mathfrak{s}) = 0\)（对所有
\(\mathfrak{s}\)） - 对 Kähler 曲面，\(\operatorname{SW}_X\)
与规范类（canonical class）直接相关

\textbf{定理 221.6}（Witten 猜想 / Taubes 定理）：Seiberg-Witten
不变量等价于 Donaldson 不变量（通过 Witten
的物理论证和后来的数学证明）。Seiberg-Witten 理论的计算比 Donaldson
理论简单得多。

\subsubsection{221.4 Exotic
结构与阻塞理论}\label{exotic-ux7ed3ux6784ux4e0eux963bux585eux7406ux8bba}

\textbf{定义 221.5}（Exotic \(\mathbb{R}^4\)）：存在与
\(\mathbb{R}^4\)（标准欧氏 4 维空间）同胚但不微分同胚的 4
维流形。事实上，存在不可数无穷个不同的 exotic \(\mathbb{R}^4\)（Taubes
1987，Gompf 1985）。

\textbf{定理 221.7}（\(\frac{11}{8}\) 猜想 / Furuta
定理）：对闭光滑单连通 4 维流形 \(X\) 满足 \(b_2^+(X) > 0\)，有
\(b_2(X) \geq \frac{11}{8} |\sigma(X)|\)（\(\sigma\)
是符号差）。\(\frac{11}{8}\) 猜想是最优下界，由 Furuta（2001）用等变 K
理论证明其弱形式。

\textbf{定义 221.6}（10/8 定理 / Furuta
不等式）：\(b_2(X) \geq \frac{10}{8} |\sigma(X)| + 2\)（Furuta 1996 的
\(10/8\) 定理给出了部分证明，完整 \(\frac{11}{8}\) 猜想仍未完全解决）。

\bigskip\hrule\bigskip

\emph{本卷在拓扑学（V10）、代数拓扑（V14）和微分几何（V11）的基础上，系统介绍了低维拓扑和纽结理论：3
维流形的分解与几何化、纽结理论的多项式与同调不变量（Alexander、Jones、HOMFLY-PT、Khovanov）、量子拓扑不变量（量子群与
TQFT）、Floer 同调系列（Instanton、Heegaard、Monopole）和 4
维流形理论（Donaldson 理论、Seiberg-Witten 理论、exotic 结构）。共 5
章。}

\newpage

\section{卷三十九：表示论深化}\label{ux5377ux4e09ux5341ux4e5dux8868ux793aux8bbaux6df1ux5316}

\begin{quote}
表示论是研究代数结构（群、代数、李代数等）通过线性作用的具体实现的理论。V12
第 69 章已介绍了有限群表示论的基础，V23
涵盖了李群与李代数的表示论。本卷在此基础上深化表示论：有限群模表示论研究特征为素数的域上的表示结构；代数群表示论将群表示论推广到代数几何框架；几何表示论通过层论和相交上同调为表示论提供几何构造；仿射李代数与量子群表示论连接无限维李代数和可积系统；Langlands
纲领深化则从表示论角度揭示数论与调和分析之间的深刻联系。
\end{quote}

\bigskip\hrule\bigskip

\subsection{第222章：有限群表示深化}\label{ux7b2c222ux7ae0ux6709ux9650ux7fa4ux8868ux793aux6df1ux5316}

有限群表示论在 V12 第 69
章已建立了基础（特征标理论、不可约表示、诱导表示）。本章进一步深化：引入模表示论（特征整除群阶的域上的表示），研究块论（Brauer
的块分解）和 Green 对应。

\subsubsection{222.1
特征标理论的深化}\label{ux7279ux5f81ux6807ux7406ux8bbaux7684ux6df1ux5316}

\textbf{定理 222.1}（特征标的正交关系，回顾）：\(G\) 的不可约复特征标
\(\operatorname{Irr}(G)\) 满足：

\[\frac{1}{|G|} \sum_{g \in G} \chi_i(g) \overline{\chi_j(g)} = \delta_{ij}\]

\textbf{定义 222.1}（特征标表）：有限群 \(G\) 的\textbf{特征标表}是
\(k \times k\) 方阵（\(k\) 为共轭类数），行标记不可约特征标
\(\chi_i\)，列标记共轭类 \(C_j\)。

\textbf{定理 222.2}（Frobenius-Schur 指示子）：设 \(\chi\) 是 \(G\)
的不可约复特征标。\textbf{Frobenius-Schur 指示子}
\(\nu_2(\chi) = \frac{1}{|G|} \sum_{g \in G} \chi(g^2)\) 取值为
\(1\)（\(\chi\) 可实现为实表示）、\(-1\)（\(\chi\)
为实值但不可实现为实表示）或 \(0\)（\(\chi\) 非实值）。

\textbf{定义 222.2}（Artin
定理）：任何特征标可表示为诱导自循环子群的一维特征标的
\(\mathbb{Q}\)-线性组合。\textbf{Brauer
定理}（更强）：任何特征标可表示为诱导自初等子群的一维特征标的
\(\mathbb{Z}\)-线性组合。

\subsubsection{222.2
模表示论基础}\label{ux6a21ux8868ux793aux8bbaux57faux7840}

\textbf{定义 222.3}（模表示 / Modular Representation）：设 \(p\)
是素数，\(k\) 是特征 \(p\) 的代数闭域。\(G\) 在 \(k\)
上的\textbf{模表示}是群同态 \(G \to \operatorname{GL}(V)\)（\(V\) 是
\(k\) 上的有限维向量空间）。当 \(p \mid |G|\)
时，模表示论与复表示论有本质不同。

\textbf{定义 222.4}（\(p\)-模系统 / \(p\)-Modular
System）：\((K, \mathcal{O}, k)\) 是 \textbf{\(p\)-模系统}，其中
\(\mathcal{O}\) 是完备离散赋值环（特征 \(0\)），\(K\) 是 \(\mathcal{O}\)
的分式域（特征 \(0\)），\(k\) 是 \(\mathcal{O}\) 的剩余域（特征
\(p\)）。\(K\) 包含 \(|G|\) 次本原单位根。

\textbf{定理 222.3}（Brauer-Nesbitt 定理）：在 \(p\)-模系统
\((K, \mathcal{O}, k)\) 下，\(K[G]\)-模和
\(k[G]\)-模的不可约表示之间存在由约化映射 \(\mathcal{O}[G] \to k[G]\)
给出的对应关系。

\textbf{定义 222.5}（Brauer 特征标 / Brauer Character）：设
\(\rho: G \to \operatorname{GL}_n(k)\) 是模表示。对 \(G\) 的
\(p\)-正则元 \(g\)（阶不被 \(p\) 整除），\(\rho(g)\) 的特征值在 \(K\)
中有唯一的提升（Teichmüller 提升）。\textbf{Brauer 特征标}
\(\varphi(g)\) 是这些提升特征值的和。\(\operatorname{IBr}(G)\)
是所有不可约 Brauer 特征标的集合。

\textbf{定理 222.4}（Brauer 特征标的性质）：\(|\operatorname{IBr}(G)|\)
等于 \(G\) 的 \(p\)-正则共轭类数。Brauer 特征标在
\(p\)-正则共轭类上满足正交关系。

\textbf{定义 222.6}（分解矩阵与 Cartan 矩阵）：\textbf{分解矩阵} \(D\)
的行标记不可约普通特征标，列标记不可约 Brauer
特征标，\(D_{\chi,\varphi}\) 是 \(\chi\) 约化到特征 \(p\) 时 \(\varphi\)
的重数。\textbf{Cartan 矩阵} \(C = D^T D\)，其元素是
\(p\)-正则类上的标量积。

\subsubsection{222.3 块理论}\label{ux5757ux7406ux8bba}

\textbf{定义 222.7}（块 / Block）：群代数 \(k[G]\)
可分解为不可分解双理想的直和：\(k[G] = B_1 \oplus B_2 \oplus \cdots \oplus B_r\)。每个
\(B_i\) 称为 \(G\) 的一个 \textbf{\(p\)-块}（或 \textbf{Brauer
块}）。块的个数等于 \(\mathcal{O}[G]\) 分解为不可分解
\(\mathcal{O}\)-代数的个数。

\textbf{定义 222.8}（亏群 / Defect Group）：\(p\)-块 \(B\)
的\textbf{亏群} \(D\)（定义到共轭）是 \(B\)
的「远离块论平凡性」的度量：\(|D| = p^d\)，其中 \(d\) 是 \(B\)
的\textbf{亏数}（defect）。亏数为 \(0\) 的块是单代数（矩阵代数）。

\textbf{定理 222.5}（Brauer 第一主定理）：存在 \(G\) 的 \(p\)-块与
\(N_G(D)\) 的 \(p\)-块（具有相同亏群 \(D\)）之间的一一对应（Brauer
对应）。

\textbf{定理 222.6}（Brauer 第二主定理）：Brauer
对应将特征标诱导和限制联系起来，为块论提供了核心计算工具。

\textbf{定义 222.9}（Green 对应 / Green Correspondence，1964）：设 \(D\)
是 \(G\) 的 \(p\)-子群，\(H\) 是包含 \(N_G(D)\) 的子群。Green 对应在
\(k[G]\)-模和 \(k[H]\)-模（具有顶点
\(D\)）之间建立了一一对应，是模表示论中最强大的工具之一。

\textbf{定理 222.7}（Alperin 权猜想，1986）：\(p\)-块 \(B\) 的不可约
Brauer 特征标数等于 \(B\) 的局部子群中 \(p\)-权（Alperin
权重）的数目。这是模表示论中最重要的未完全解决的猜想之一。

\bigskip\hrule\bigskip

\subsection{第223章：代数群表示}\label{ux7b2c223ux7ae0ux4ee3ux6570ux7fa4ux8868ux793a}

代数群表示论将群表示论推广到代数几何框架，研究代数群（既是群又是代数簇）的线性表示。半单代数群和约化代数群的表示分类是
20 世纪表示论的最高成就之一。

\subsubsection{223.1
线性代数群基础}\label{ux7ebfux6027ux4ee3ux6570ux7fa4ux57faux7840}

\textbf{定义 223.1}（线性代数群）：\textbf{线性代数群} \(G\) 是代数闭域
\(k\) 上的仿射代数簇，带有群运算（乘法 \(G \times G \to G\) 和取逆
\(G \to G\)），使得运算是代数簇的态射。

\textbf{例
223.1}（代数群的例子）：\(\operatorname{GL}_n(k)\)、\(\operatorname{SL}_n(k)\)、\(\operatorname{SO}_n(k)\)、\(\operatorname{Sp}_{2n}(k)\)、\(\mathbb{G}_a = (k, +)\)（加法群）、\(\mathbb{G}_m = (k^\times, \times)\)（乘法群）。

\textbf{定义 223.2}（Borel 子群与极大环面）：代数群 \(G\) 的
\textbf{Borel 子群} \(B\) 是 \(G\) 的极大连通可解子群。\(B\)
的\textbf{极大环面} \(T \subset B\) 是 \(B\) 的极大环面子群（同构于
\((\mathbb{G}_m)^n\)）。\(T\) 的\textbf{特征标群}
\(X(T) = \operatorname{Hom}(T, \mathbb{G}_m)\) 是自由 Abel 群（秩
\(= \dim T\)）。

\textbf{定理 223.1}（Borel 不动点定理）：若连通可解代数群 \(G\)
作用在完备代数簇 \(X\) 上，则 \(G\) 在 \(X\) 上有不动点。

\textbf{定理 223.2}（Bruhat
分解）：\(G = \bigsqcup_{w \in W} B w B\)，其中 \(W = N_G(T)/T\) 是 Weyl
群。Bruhat 分解是代数群表示论的核心工具。

\subsubsection{223.2 最高权理论}\label{ux6700ux9ad8ux6743ux7406ux8bba-1}

\textbf{定义 223.3}（支配权与最高权模）：设 \(G\) 是连通约化代数群，选定
Borel 子群 \(B \supset T\)。\textbf{权} \(\lambda \in X(T)\)
是\textbf{支配的}（dominant），如果
\(\langle \lambda, \alpha^\vee \rangle \geq 0\)（对所有正根
\(\alpha\)）。对每个支配权 \(\lambda\)，存在唯一的不可约 \(G\)-模
\(L(\lambda)\)，其最高权为 \(\lambda\)。

\textbf{定义 223.4}（诱导模 / Weyl 模与对偶 Weyl 模）：对支配权
\(\lambda\)，\textbf{Weyl 模} \(\Delta(\lambda)\) 定义为
\(\operatorname{ind}_B^G(k_\lambda)\)（\(k_\lambda\) 是 \(B\)
的一维表示）。\(L(\lambda)\) 是 \(\Delta(\lambda)\) 的唯一不可约商。

\textbf{定理 223.3}（Weyl 特征标公式）：不可约表示 \(L(\lambda)\)
的特征标为

\[\chi(\lambda) = \frac{\sum_{w \in W} (-1)^{\ell(w)} e^{w(\lambda + \rho)}}{\sum_{w \in W} (-1)^{\ell(w)} e^{w\rho}}\]

其中 \(\rho\) 是正根的一半之和，\(\ell(w)\) 为 \(w\) 的长度。（这里
\(\rho\) 为所有正根之和的一半，Weyl 群的 dot action 定义为
\(w \cdot \lambda = w(\lambda + \rho) - \rho\)。）这是表示论中最优美的公式之一。

\textbf{定理 223.4}（Weyl
维数公式）：\(\dim L(\lambda) = \prod_{\alpha > 0} \frac{\langle \lambda + \rho, \alpha^\vee \rangle}{\langle \rho, \alpha^\vee \rangle}\)。

\subsubsection{223.3 Kempf
消失定理与上同调}\label{kempf-ux6d88ux5931ux5b9aux7406ux4e0eux4e0aux540cux8c03}

\textbf{定义 223.5}（旗簇与诱导层）：\(G/B\) 是\textbf{旗簇}（flag
variety），是光滑射影簇。对权 \(\lambda\)，\(G \times^B k_\lambda\) 是
\(G/B\) 上的可逆层 \(\mathcal{L}(\lambda)\)。

\textbf{定理 223.5}（Borel-Weil-Bott 定理）：不可约 \(G\)-模
\(L(\lambda)\) 可通过旗簇上可逆层的上同调实现： -
\(H^0(G/B, \mathcal{L}(\lambda)) \cong L(\lambda)^*\)（若 \(\lambda\)
是支配权） - 对非支配权，上同调消失或给出其他不可约表示（Bott
的精细公式）

\textbf{定理 223.6}（Kempf 消失定理，1976）：若 \(\lambda\) 是支配权，则
\(H^i(G/B, \mathcal{L}(\lambda)) = 0\)（对所有 \(i > 0\)）。这是
Borel-Weil-Bott 定理的特例，是射影簇上丰沛线丛消失定理的推广。

\textbf{定义 223.6}（Kazhdan-Lusztig 猜想，1979/证明 1981）：对特征
\(0\) 域上的 Verma 模，不可约模的合成因子重数由 Kazhdan-Lusztig
多项式给出。Brylynski-Kashiwara（1981）和
Beilinson-Bernstein（1981）通过
\(\mathcal{D}\)-模和相交上同调证明了该猜想。

\subsubsection{223.4
正特征表示论}\label{ux6b63ux7279ux5f81ux8868ux793aux8bba}

\textbf{定义 223.7}（正特征表示论 / 特征 \(p\) 的代数群表示）：在特征
\(p > 0\) 的代数闭域上，代数群表示论与特征 \(0\) 有本质不同。Frobenius
态射 \(F: G \to G\) 引入新的表示（Frobenius 扭）。

\textbf{定理 223.7}（Steinberg 张量积定理）：对特征 \(p\)
域上的单连通半单代数群 \(G\)，不可约 \(G\)-模的结构为

\[L(\lambda) = L(\lambda_0) \otimes L(\lambda_1)^{(1)} \otimes L(\lambda_2)^{(2)} \otimes \cdots\]

其中 \(\lambda = \lambda_0 + p\lambda_1 + p^2\lambda_2 + \cdots\) 是
\(\lambda\) 的 \(p\)-adic 展开，\(L(\mu)^{(r)}\) 是 \(F^r\)-扭表示。

\textbf{定理 223.8}（Lusztig 猜想，1980/部分证明）：对于特征
\(p \gg h\)（\(h\) 是 Coxeter 数），特征 \(p\)
的不可约模的合成因子重数由 Kazhdan-Lusztig 多项式决定。Lusztig 猜想在
\(p \gg 0\) 时成立（Andersen-Jantzen-Soergel 1994），但存在反例证明
\(p\) 需要相当大。

\bigskip\hrule\bigskip

\subsection{第224章：几何表示论}\label{ux7b2c224ux7ae0ux51e0ux4f55ux8868ux793aux8bba}

几何表示论通过几何对象（旗簇、Steinberg 簇、仿射 Grassmann
流形、箭图簇）的拓扑和层论来构造和研究代数群的表示。这是 20 世纪末和 21
世纪表示论最活跃的前沿之一。

\subsubsection{224.1 Springer 理论}\label{springer-ux7406ux8bba}

\textbf{定义 224.1}（Springer 纤维 / Springer 解消）：设 \(G\)
是约化代数群，\(\mathfrak{g} = \operatorname{Lie}(G)\)，\(\mathcal{N} \subset \mathfrak{g}\)
是幂零锥。\textbf{Springer 解消}是态射
\(\pi: \tilde{\mathcal{N}} = \{(x, B) : x \in \operatorname{Lie}(B) \cap \mathcal{N}\} \to \mathcal{N}\)。对
\(x \in \mathcal{N}\)，\(\pi^{-1}(x)\) 是 \textbf{Springer 纤维}
\(\mathcal{B}_x\)（包含 \(x\) 的 Borel 子代数的旗簇）。

\textbf{定理 224.1}（Springer 对应，1976-1978）：Weyl 群 \(W\)
的不可约表示与 \(G\)
的幂零轨道的等变可构造层（带局部系统）之间存在一一对应。\(W\) 在
\(H^*(\mathcal{B}_x)\) 上的作用（Springer 表示）实现了 \(W\)
的所有不可约表示。

\textbf{定义 224.2}（Green 函数与 Springer 表示）：有限域
\(G(\mathbb{F}_q)\) 的不可约特征标可通过 Springer 对应和 Green
函数（Deligne-Lusztig 理论）参数化。

\subsubsection{224.2 相交上同调与 Kazhdan-Lusztig
理论}\label{ux76f8ux4ea4ux4e0aux540cux8c03ux4e0e-kazhdan-lusztig-ux7406ux8bba}

\textbf{定义 224.3}（相交上同调 / Intersection
Cohomology，Goresky-MacPherson 1980）：对奇异空间
\(X\)，\textbf{相交上同调} \(IH^*(X)\) 是满足 Poincaré
对偶的「修正」上同调理论。由中间扩展（intermediate extension）的可构造层
\(\mathbf{IC}(X)\) 的超上同调给出。

\textbf{定义 224.4}（Schubert 簇与 Kazhdan-Lusztig 多项式）：旗簇
\(G/B\) 的 \textbf{Schubert 簇}
\(X_w = \overline{B w B / B}\)（\(w \in W\)）。\textbf{Kazhdan-Lusztig
多项式} \(P_{x,w}(q)\) 编码了 \(X_w\) 沿 \(B x B / B\) 的局部相交上同调
Poincaré 多项式。

\textbf{定理 224.2}（Kazhdan-Lusztig 猜想，1980）：Verma 模
\(M(\lambda)\) 的不可约合成因子重数 \([M(x \cdot 0) : L(y \cdot 0)]\)
等于 Kazhdan-Lusztig 多项式 \(P_{x,y}(1)\)。由 Beilinson-Bernstein 和
Brylynski-Kashiwara（1981）通过 \(\mathcal{D}\)-模和 Riemann-Hilbert
对应证明。

\textbf{定义 224.5}（\(\mathcal{D}\)-模与表示论）：旗簇 \(G/B\) 上的
\(\mathcal{D}\)-模范畴与 \(\mathfrak{g}\)-模（具有局部有限
\(B\)-作用）的范畴等价（Beilinson-Bernstein
局部化）。这是几何表示论的核心结果。

\subsubsection{224.3 几何 Satake
对应}\label{ux51e0ux4f55-satake-ux5bf9ux5e94}

\textbf{定义 224.6}（仿射 Grassmann 流形）：对代数闭域 \(k\) 上的约化群
\(G\)，\textbf{仿射 Grassmann 流形}
\(\operatorname{Gr}_G = G(k((t))) / G(k[[t]])\)（\(k((t))\) 是 Laurent
级数域，\(k[[t]]\) 是形式幂级数环）。\(\operatorname{Gr}_G\) 是无穷维
ind-射影簇。

\textbf{定义 224.7}（几何 Satake 等价，Lusztig 1983, Ginzburg 1995,
Mirković-Vilonen 2007）：存在范畴等价

\[\operatorname{Perv}_{G(k[[t]])}(\operatorname{Gr}_G) \cong \operatorname{Rep}(G^\vee)\]

其中 \(G^\vee\) 是 \(G\) 的 Langlands 对偶群，\(\operatorname{Perv}\)
是反常层（perverse sheaves）范畴。几何 Satake 将 \(G\) 的几何对象与
\(G^\vee\) 的表示范畴联系起来。

\textbf{定理 224.3}（几何 Satake 的推论）：仿射 Grassmann 流形上
\(G(k[[t]])\)-等变反常层的超上同调给出 \(G^\vee\)
的表示的几何构造。卷积积对应于表示的张量积。

\subsubsection{224.4
箭图簇与表示论}\label{ux7badux56feux7c07ux4e0eux8868ux793aux8bba}

\textbf{定义 224.8}（箭图簇 / Quiver Variety，Nakajima
1998）：\textbf{箭图簇} \(\mathfrak{M}(\mathbf{v}, \mathbf{w})\) 是箭图
\(Q\) 的双重（double）表示模空间的辛约化。箭图簇是超 Kähler 流形。

\textbf{定理 224.4}（Nakajima 定理）：箭图簇的拓扑 Borel-Moore
同调实现了 Kac-Moody
代数的可积最高权表示。具体地，\(\bigoplus_{\mathbf{v}} H_*^{\operatorname{BM}}(\mathfrak{M}(\mathbf{v}, \mathbf{w}))\)
具有 \(\mathfrak{g}(Q)\) 的最高权 \(\mathbf{w}\) 的可积表示结构。

\textbf{定义
224.9}（箭图簇与几何表示论的联系）：箭图簇为仿射量子群、杨振宁代数和
\(\mathcal{W}\)-代数的表示提供了统一的几何框架。Nakajima
的箭图簇理论是几何表示论的基础工具。

\bigskip\hrule\bigskip

\subsection{第225章：仿射与量子群表示}\label{ux7b2c225ux7ae0ux4effux5c04ux4e0eux91cfux5b50ux7fa4ux8868ux793a}

仿射李代数表示论连接了无限维李代数、共形场论和可积系统。量子群表示论是代数群表示论的非交换形变，在纽结理论和统计力学中有重要应用。

\subsubsection{225.1 仿射 Lie
代数表示}\label{ux4effux5c04-lie-ux4ee3ux6570ux8868ux793a}

\textbf{定义 225.1}（仿射 Kac-Moody 代数）：\textbf{仿射 Kac-Moody 代数}
\(\hat{\mathfrak{g}}\) 是有限维单李代数 \(\mathfrak{g}\) 的圈代数
\(\mathfrak{g} \otimes \mathbb{C}[t, t^{-1}]\) 的中心扩张。其定义为： -
\(\hat{\mathfrak{g}} = (\mathfrak{g} \otimes \mathbb{C}[t, t^{-1}]) \oplus \mathbb{C} c \oplus \mathbb{C} d\)
-
\([x \otimes t^m, y \otimes t^n] = [x, y] \otimes t^{m+n} + m \delta_{m,-n} (x|y) c\)
- \([c, \cdot] = 0\)，\([d, x \otimes t^n] = n x \otimes t^n\)

\textbf{定义 225.2}（可积最高权表示）：设 \(\lambda\) 是
\(\hat{\mathfrak{g}}\) 的支配整权。不可约最高权模 \(L(\lambda)\)
是\textbf{可积的}，如果 \(\lambda(c) \neq 0\) 且 Chevalley
生成元的作用是局部幂零的。

\textbf{定理 225.1}（Kac-Weisfeiler 特征标公式 / Kac
特征标公式）：可积最高权模 \(L(\lambda)\) 的特征标为

\[\operatorname{ch} L(\lambda) = \frac{\sum_{w \in \hat{W}} (-1)^{\ell(w)} e^{w(\lambda + \hat{\rho}) - \hat{\rho}}}{\prod_{\alpha > 0} (1 - e^{-\alpha})^{\dim \hat{\mathfrak{g}}_\alpha}}\]

其中 \(\hat{W}\) 是仿射 Weyl 群，\(\hat{\rho}\) 是仿射 Weyl 向量。

\textbf{定义 225.3}（Macdonald 恒等式与 \(\eta\)
函数）：仿射李代数的分母恒等式（Weyl
分母公式）给出经典数论公式的推广。\(\widehat{\mathfrak{sl}}_2\)
的分母公式给出 Jacobi 三重积恒等式和 Dedekind \(\eta\) 函数。

\subsubsection{225.2
顶点算子代数}\label{ux9876ux70b9ux7b97ux5b50ux4ee3ux6570}

\textbf{定义 225.4}（顶点算子代数 / VOA）：\textbf{顶点算子代数}
\((V, Y, \mathbf{1}, \omega)\) 是向量空间 \(V\) 配备： - 状态-场对应
\(Y: V \to \operatorname{End}(V)[[z, z^{-1}]]\)，\(Y(a, z) = \sum_{n \in \mathbb{Z}} a_n z^{-n-1}\)
- 真空向量 \(\mathbf{1}\)：\(Y(\mathbf{1}, z) = \operatorname{id}_V\) -
共形向量 \(\omega\)：\(Y(\omega, z) = \sum L_n z^{-n-2}\)（Virasoro
代数作用） - 局部性公理：\((z-w)^N [Y(a, z), Y(b, w)] = 0\)（\(N\)
充分大）

\textbf{定理 225.2}（Frenkel-Zhu 定理，1992）：仿射 Kac-Moody 代数
\(\hat{\mathfrak{g}}\) 的可积最高权模 \(L(k\Lambda_0)\)（\(k\)
是级）具有顶点算子代数结构（仿射 VOA）。

\textbf{定义 225.5}（共形场论与表示论）：有理顶点算子代数 \(V\)
的表示范畴是模张量范畴（modular tensor category），与 3
维拓扑量子场论（TQFT）相关。这是 VOA 表示论与低维拓扑的联系。

\textbf{定理 225.3}（Verlinde 公式，1988）：有理 VOA 的不可约表示
\(M_i\) 的融合规则（fusion rules）\(N_{ij}^k\) 与模群 \(S\) 矩阵的关系为
\(N_{ij}^k = \sum_r \frac{S_{ir} S_{jr} S_{kr}^*}{S_{0r}}\)。

\subsubsection{225.3
量子仿射代数与可积系统}\label{ux91cfux5b50ux4effux5c04ux4ee3ux6570ux4e0eux53efux79efux7cfbux7edf}

\textbf{定义 225.6}（量子仿射代数
\(U_q(\hat{\mathfrak{g}})\)）：\textbf{量子仿射代数}是仿射 Kac-Moody
代数 \(\hat{\mathfrak{g}}\) 的量子包络代数的 \(q\)-形变。它是
Drinfeld-Jimbo 量子群（见 V23 Ch 118）的仿射版本。

\textbf{定义 225.7}（量子 R-矩阵与 Yang-Baxter 方程）：量子仿射代数
\(U_q(\hat{\mathfrak{g}})\) 的泛 R-矩阵满足 Yang-Baxter 方程。有限维表示
\(V\) 上的 R-矩阵 \(R_{V,W}: V \otimes W \to V \otimes W\)
是谱依赖的（\(R(z)\)）。

\textbf{定理 225.4}（Baxter 的角转移矩阵与 Bethe 拟设）：量子仿射代数
\(U_q(\hat{\mathfrak{sl}}_2)\) 的表示论与 XXZ 自旋链的 Bethe
拟设解直接相关。这是量子可积系统与量子群表示论的核心联系。

\textbf{定义 225.8}（Frenkel-Reshetikhin 的
\(q\)-特征标，1999）：量子仿射代数的有限维表示的
\(q\)-特征标是经典特征标的 \(q\)-形变，满足与 Baxter 的 T-Q
关系类似的性质。

\bigskip\hrule\bigskip

\subsection{第226章：Langlands
纲领深化}\label{ux7b2c226ux7ae0langlands-ux7eb2ux9886ux6df1ux5316}

Langlands 纲领（V35 Ch 174
已有概述）是当代数学中最宏大的统一纲领之一。本章从表示论角度深化
Langlands 纲领的核心内容：局部 Langlands 对应、几何 Langlands 纲领和
Arthur-Selberg 迹公式。

\subsubsection{226.1 局部 Langlands
对应}\label{ux5c40ux90e8-langlands-ux5bf9ux5e94}

\textbf{定义 226.1}（局部域与局部 Langlands 对应）：设 \(F\)
是局部域（\(\mathbb{R}\)、\(\mathbb{C}\) 或 \(\mathbb{Q}_p\)
的有限扩张）。\textbf{局部 Langlands 对应}（LLC）断言 \(G(F)\)
的不可约光滑表示与 Weil-Deligne 群 \(W_F'\) 到 Langlands 对偶群
\(G^\vee(\mathbb{C})\) 的 \(L\)-参数之间存在一一对应。

\textbf{定理 226.1}（\(\operatorname{GL}_n\) 的局部 Langlands
对应，Harris-Taylor 2001, Henniart 2000）：\(\operatorname{GL}_n(F)\)
的不可约光滑超尖表示与 \(n\) 维 Frobenius 半单 Weil-Deligne
表示之间存在一一对应，保持 \(L\)-因子和 \(\varepsilon\)-因子。

\textbf{定义 226.2}（\(L\)-参数与 \(L\)-包）：一般约化群 \(G\) 的
\(L\)-参数 \(\varphi: W_F' \to {}^L G\) 确定一个 \(L\)-包
\(\Pi_\varphi\)（不可约表示的有限集）。\(L\)-包内的表示由 \(\varphi\)
的中心化子 \(S_\varphi\) 的不可约表示参数化。

\textbf{定理 226.2}（\(\operatorname{GL}_n\) 的局部 Langlands 对应对
\(L\)-函数和 \(\varepsilon\)-因子的保持）：对
\(\operatorname{GL}_n \times \operatorname{GL}_m\) 的 Rankin-Selberg
\(L\)-函数，局部 Langlands 对应保持局部 \(L\)-函数和
\(\varepsilon\)-因子。

\subsubsection{226.2 整体 Langlands
纲领}\label{ux6574ux4f53-langlands-ux7eb2ux9886}

\textbf{定义 226.3}（自守表示与 \(L\)-函数）：设 \(G\) 是 \(\mathbb{Q}\)
上的约化代数群，\(\mathbb{A}\) 是 \(\mathbb{Q}\) 的 Adele
环。\textbf{自守表示} \(\pi\) 是 \(G(\mathbb{A})\) 在
\(L^2(G(\mathbb{Q}) \backslash G(\mathbb{A}))\) 上的不可约成分。\(\pi\)
的\textbf{自守 \(L\)-函数} \(L(s, \pi, r)\) 是 Euler 乘积（\(r\) 是
\({}^L G\) 的有限维表示）。

\textbf{定理 226.3}（整体 Langlands 函子性猜想）：对 \(L\)-同态
\(\eta: {}^L H \to {}^L G\)，存在自守表示的「提升」映射（函子性转移），将
\(H(\mathbb{A})\) 的自守表示映为 \(G(\mathbb{A})\) 的自守表示，并保持
\(L\)-函数。

\emph{函子性猜想的简要描述}：Langlands 的函子性原理断言，任意两个约化群
\(H\) 和 \(G\) 的 \(L\)-群之间的代数同态 \(\eta: {}^L H \to {}^L G\)
诱导自守表示范畴之间的函子
\(\eta_*: \mathcal{A}(H) \to \mathcal{A}(G)\)，满足
\(L(s, \pi, r \circ \eta) = L(s, \eta_*(\pi), r)\) 对任意 \({}^L G\)
的有限维表示 \(r\)。这意味着
\(L\)-群的同态''提升''了自守形式，将较小的群上的自守表示''转移''到较大的群上。当
\(G = \operatorname{GL}_n\) 时，该猜想已由 Arthur
等人通过稳定迹公式在若干情形下得到证明；对一般的 \(G\)，Langlands
函子性原理仍然是 Langlands 纲领中最核心的未完全解决的问题，Ngô
的基本引理证明（2010）为此提供了关键技术基础。∎

\textbf{定义 226.4}（Langlands 纲领的核心猜想）： 1.
\textbf{局部-整体相容性}：整体自守表示在所有局部点的局部分量通过局部
Langlands 对应与整体 Galois 表示兼容 2. \textbf{Ramanujan-Petersson
猜想}：尖点自守表示在非分歧点处的局部 Satake 参数是幺模的 3.
\textbf{广义 Riemann 猜想}：自守 \(L\)-函数 \(L(s, \pi)\)
的所有非平凡零点在 \(\Re(s) = 1/2\) 上

\subsubsection{226.3 Arthur-Selberg
迹公式}\label{arthur-selberg-ux8ff9ux516cux5f0f}

\textbf{定义 226.5}（Arthur-Selberg 迹公式）：\textbf{Arthur-Selberg
迹公式}是 \(L^2(G(\mathbb{Q}) \backslash G(\mathbb{A}))\)
上正则表示的迹的等式，一边是谱和（自守表示的特征标），另一边是几何和（轨道积分）。它是
Poisson 求和公式在非交换调和分析中的推广。

\textbf{定理 226.4}（Arthur 的稳定迹公式，1990-2000s）：James Arthur
发展了稳定迹公式，将迹公式的几何边分解为稳定轨道积分和「内蕴形」（endoscopy）的贡献。这为一般群的
Langlands 函子性提供了基础框架。

\textbf{定理 226.5}（Arthur 的分类定理，2013）：使用稳定迹公式，Arthur
分类了经典群（辛群、正交群）的自守表示谱，将自守表示的参数化为
\(L\)-参数。

\textbf{定义 226.6}（迹公式与 Langlands 纲领的联系）：迹公式是 Langlands
函子性猜想的核心计算工具。特别是，通过迹公式的比较（Fundamental
Lemma），可以证明自守表示的函子性转移。

\subsubsection{226.4 几何 Langlands
纲领}\label{ux51e0ux4f55-langlands-ux7eb2ux9886}

\textbf{定义 226.7}（几何 Langlands 纲领，Beilinson-Drinfeld
1990s）：\textbf{几何 Langlands 纲领}是 Langlands 纲领的代数几何类比。设
\(X\) 是 \(\mathbb{C}\) 上的光滑射影代数曲线，\(G\) 是约化代数群。几何
Langlands 猜想断言存在范畴等价：

\[\mathcal{D}\text{-}\operatorname{Mod}(\operatorname{Bun}_G) \cong \operatorname{QCoh}(\operatorname{LocSys}_{G^\vee})\]

其中 \(\operatorname{Bun}_G\) 是 \(X\) 上
\(G\)-主丛的模叠，\(\operatorname{LocSys}_{G^\vee}\) 是 \(X\) 上
\(G^\vee\)-局部系统的模叠。

\textbf{定理 226.6}（几何 Langlands 的进展）： - \(\operatorname{GL}_1\)
的几何 Langlands：由 Laumon（1987）和 Rothstein（1996）证明（等价于
Abel-Jacobi 映射和 Fourier-Mukai 变换） - 一般 \(G\) 的几何
Langlands：Arinkin-Gaitsgory（2015）证明了 de Rham 版本的几何 Langlands
猜想（奇异支集定理） - 几何 Langlands 与
S-对偶性：Kapustin-Witten（2007）将几何 Langlands 解释为
\(\mathcal{N}=4\) 超 Yang-Mills 理论的电磁对偶

\textbf{定义 226.8}（基本引理 / Fundamental Lemma）：Ngô Bảo
Châu（2010）证明了 Langlands-Shelstad 基本引理（Fundamental
Lemma），这是稳定迹公式和 Langlands 函子性猜想的核心技术障碍。Ngô
因此获得 2010 年 Fields 奖。

\textbf{定理 226.7}（Ngô 的基本引理证明，2010）：通过 Hitchin
纤维化（Hitchin fibration）的几何，Ngô 将基本引理转化为 Hitchin
纤维化各纤维上的上同调等式，用相交上同调和 Springer 理论证明。

\bigskip\hrule\bigskip

\emph{本卷在 V12（抽象代数 II）和
V23（李理论）中表示论基础之上，深化了有限群模表示论（Brauer
特征标、块论、Green 对应）、代数群表示论（最高权理论、Kempf
消失定理、正特征表示论）、几何表示论（Springer 理论、Kazhdan-Lusztig
理论、几何 Satake、箭图簇）、仿射与量子群表示论（仿射 Kac-Moody
代数、顶点算子代数、量子仿射代数）和 Langlands 纲领深化（局部/整体
Langlands 对应、Arthur-Selberg 迹公式、几何 Langlands 纲领）。共 5 章。}

\newpage

\section{卷四十：控制论}\label{ux5377ux56dbux5341ux63a7ux5236ux8bba}

\begin{quote}
控制论是研究动态系统调控规律的数学分支，由 Norbert Wiener
命名（1948），其数学基础则由 Pontryagin、Bellman、Kalman
等人奠定。控制论的核心问题是：如何设计控制输入使得动态系统达到期望的行为？这涉及可控性（能否到达）、可观测性（能否推断）、稳定性（能否保持）和最优化（如何最好地到达）。本卷在现代分析（V13）、微分方程（V17）和优化理论（V26）的基础上，系统建立线性控制、最优控制、随机滤波、鲁棒控制和非线性控制的数学理论。
\end{quote}

\bigskip\hrule\bigskip

\subsection{第227章：线性控制系统}\label{ux7b2c227ux7ae0ux7ebfux6027ux63a7ux5236ux7cfbux7edf}

线性控制理论是控制论的基石，由 Rudolf Kalman 在 1960
年代用状态空间方法系统化。状态空间方法将经典控制理论中的频域方法（传递函数）统一为时域形式。

\subsubsection{227.1
状态空间模型}\label{ux72b6ux6001ux7a7aux95f4ux6a21ux578b}

\textbf{定义 227.1}（线性时不变系统 / LTI
System）：连续时间线性时不变系统的状态空间表示为

\[\begin{cases} \dot{x}(t) = A x(t) + B u(t) \\ y(t) = C x(t) + D u(t) \end{cases}\]

其中 \(x(t) \in \mathbb{R}^n\) 是状态，\(u(t) \in \mathbb{R}^m\)
是控制输入，\(y(t) \in \mathbb{R}^p\)
是观测输出。\(A \in \mathbb{R}^{n \times n}\)，\(B \in \mathbb{R}^{n \times m}\)，\(C \in \mathbb{R}^{p \times n}\)，\(D \in \mathbb{R}^{p \times m}\)。

\textbf{定义 227.2}（解与状态转移矩阵）：齐次方程 \(\dot{x} = A x\)
的解为 \(x(t) = e^{At} x(0)\)，其中矩阵指数
\(e^{At} = \sum_{k=0}^{\infty} \frac{(At)^k}{k!}\)。\(\Phi(t) = e^{At}\)
称为\textbf{状态转移矩阵}。非齐次方程的通解为

\[x(t) = e^{At} x(0) + \int_0^t e^{A(t-\tau)} B u(\tau) d\tau\]

\textbf{定义 227.3}（传递函数）：系统的\textbf{传递函数} \(G(s)\)
是输出和输入的 Laplace 变换之比（零初始条件）：

\[G(s) = C(sI - A)^{-1} B + D\]

\(G(s)\) 是 \(p \times m\) 的有理函数矩阵。

\subsubsection{227.2 可控性}\label{ux53efux63a7ux6027}

\textbf{定义 227.4}（可控性 / Controllability）：系统 \((A, B)\)
是\textbf{可控的}（或状态 \(x_0\) 在时间 \(T\)
内可控），如果对任意初始状态 \(x(0) = x_0\) 和目标状态
\(x_f\)，存在控制输入 \(u(t)\)（\(t \in [0, T]\)）使得 \(x(T) = x_f\)。

\textbf{定理 227.1}（Kalman 可控性判据）：系统 \((A, B)\)
完全可控当且仅当\textbf{可控性矩阵}

\[\mathcal{C} = \begin{bmatrix} B & AB & A^2B & \cdots & A^{n-1}B \end{bmatrix} \in \mathbb{R}^{n \times nm}\]

满行秩：\(\operatorname{rank}(\mathcal{C}) = n\)。

\textbf{定理 227.2}（Hautus 引理 / PBH 检验）：\((A, B)\)
可控当且仅当对所有 \(\lambda \in \mathbb{C}\)，

\[\operatorname{rank}\begin{bmatrix} \lambda I - A & B \end{bmatrix} = n\]

只需对 \(A\) 的特征值验证即可。

\textbf{定理 227.3}（Kalman 可控性分解）：若
\(\operatorname{rank}(\mathcal{C}) = r < n\)，存在坐标变换将系统分解为

\[\begin{bmatrix} \dot{x}_c \\ \dot{x}_{\bar{c}} \end{bmatrix} = \begin{bmatrix} A_c & A_{12} \\ 0 & A_{\bar{c}} \end{bmatrix} \begin{bmatrix} x_c \\ x_{\bar{c}} \end{bmatrix} + \begin{bmatrix} B_c \\ 0 \end{bmatrix} u\]

其中 \((A_c, B_c)\) 是可控子系统，\(A_{\bar{c}}\) 对应不可控模式。

\subsubsection{227.3 可观测性}\label{ux53efux89c2ux6d4bux6027}

\textbf{定义 227.5}（可观测性 / Observability）：系统 \((A, C)\)
是\textbf{可观测的}，如果对任意初始状态 \(x(0) = x_0\)，能够在有限时间
\(T\) 内仅通过观测输出 \(y(t)\) 和输入
\(u(t)\)（\(t \in [0, T]\)）唯一确定 \(x_0\)。

\textbf{定理 227.4}（Kalman 可观测性判据）：系统 \((A, C)\)
完全可观测当且仅当\textbf{可观测性矩阵}

\[\mathcal{O} = \begin{bmatrix} C \\ CA \\ CA^2 \\ \vdots \\ CA^{n-1} \end{bmatrix} \in \mathbb{R}^{np \times n}\]

满列秩：\(\operatorname{rank}(\mathcal{O}) = n\)。

\textbf{定理 227.5}（对偶原理）：\((A, C)\) 可观测当且仅当对偶系统
\((A^T, C^T)\) 可控。这是最优控制和最优估计的对偶性的基础。

\subsubsection{227.4
极点配置与镇定}\label{ux6781ux70b9ux914dux7f6eux4e0eux9547ux5b9a}

\textbf{定理 227.6}（极点配置定理，Wonham 1967）：若 \((A, B)\)
可控，则对任意给定的 \(n\) 个复共轭对 \(p_1, \ldots, p_n\)，存在状态反馈
\(u = -Kx\) 使得闭环系统 \(\dot{x} = (A - BK)x\) 的特征值为
\(\{p_1, \ldots, p_n\}\)。

\textbf{定义 227.6}（可镇定性与可检测性）：\((A, B)\)
\textbf{可镇定}（stabilizable）如果所有不可控模式都是稳定的（特征值实部为负）。\((A, C)\)
\textbf{可检测}（detectable）如果所有不可观测模式都是稳定的。

\textbf{定理 227.7}（分离原理 / Separation
Principle）：对可镇定且可检测的系统，状态观测器和状态反馈控制器可以独立设计，闭环系统保持稳定。

\bigskip\hrule\bigskip

\subsection{第228章：最优控制}\label{ux7b2c228ux7ae0ux6700ux4f18ux63a7ux5236}

最优控制研究如何在满足动态约束的前提下，找到使性能指标最小化（或最大化）的控制策略。

\subsubsection{228.1 变分法与 Pontryagin
极小值原理}\label{ux53d8ux5206ux6cd5ux4e0e-pontryagin-ux6781ux5c0fux503cux539fux7406}

\textbf{定义 228.1}（最优控制问题 / Bolza 形式）：考虑系统
\(\dot{x}(t) = f(t, x(t), u(t))\)，初始条件 \(x(t_0) = x_0\)。寻找控制
\(u(t) \in U \subseteq \mathbb{R}^m\) 最小化性能指标

\[J[u] = \phi(x(t_f)) + \int_{t_0}^{t_f} L(t, x(t), u(t)) dt\]

\textbf{定理 228.1}（Pontryagin 极小值原理，1956）：若 \(u^*(t)\)
是最优控制，\(x^*(t)\) 是最优轨线，则存在非零伴随向量
\(\lambda(t)\)（协态）满足：

\begin{enumerate}
\def\labelenumi{\arabic{enumi}.}
\tightlist
\item
  \textbf{状态方程}：\(\dot{x}^* = \frac{\partial H}{\partial \lambda}(t, x^*, u^*, \lambda)\)
\item
  \textbf{协态方程}：\(\dot{\lambda} = -\frac{\partial H}{\partial x}(t, x^*, u^*, \lambda)\)
\item
  \textbf{极小值条件}：\(H(t, x^*, u^*, \lambda) \leq H(t, x^*, u, \lambda)\)（对所有
  \(u \in U\)）
\item
  \textbf{横截条件}：\(\lambda(t_f) = \frac{\partial \phi}{\partial x}(x^*(t_f))\)
\end{enumerate}

其中 \textbf{Hamilton 函数}
\(H(t, x, u, \lambda) = L(t, x, u) + \lambda^T f(t, x, u)\)。

\emph{证明思路}（变分推导概要）：考虑 Bolza 型性能指标
\(J[u] = \phi(x(t_f)) + \int_{t_0}^{t_f} L(t, x, u) dt\)，在动态约束
\(\dot{x} = f(t, x, u)\) 下求极小。引入 Lagrange
乘子（协态）\(\lambda(t) \in \mathbb{R}^n\)，构造增广性能指标
\(\tilde{J} = \phi(x(t_f)) + \int_{t_0}^{t_f} [L + \lambda^T(f - \dot{x})] dt\)。定义
\textbf{Hamilton 函数}
\(H(t, x, u, \lambda) = L(t, x, u) + \lambda^T f(t, x, u)\)，则
\(\tilde{J} = \phi(x(t_f)) + \int_{t_0}^{t_f} [H(t, x, u, \lambda) - \lambda^T \dot{x}] dt\)。对
\(\tilde{J}\) 取一阶变分：边界项给出横截条件
\(\lambda(t_f) = \partial\phi/\partial x(x(t_f))\)；对 \(\lambda\)
的变分恢复状态方程 \(\dot{x} = \partial H/\partial\lambda\)；对 \(x\)
的变分经分部积分后得到\textbf{协态方程}
\(\dot{\lambda} = -\partial H/\partial x\)；对 \(u\) 的变分要求
\(\partial H/\partial u = 0\)（无约束情形），或在控制约束 \(u \in U\)
下有极小值条件
\(H(x^*, u^*, \lambda) = \min_{u \in U} H(x^*, u, \lambda)\)。这就是
Pontryagin 极小值原理的核心：最优控制使 Hamilton 函数逐点取极小。

\textbf{例 228.1}（线性二次型调节器 /
LQR）：\(f = Ax + Bu\)，\(L = \frac{1}{2}(x^T Q x + u^T R u)\)（\(Q \geq 0\)，\(R > 0\)）。由极小值原理，最优控制为状态反馈
\(u^* = -R^{-1} B^T P x\)，其中 \(P\) 满足 \textbf{Riccati
方程}：\(PA + A^T P - P B R^{-1} B^T P + Q = 0\)。

\textbf{例 228.2}（倒立摆的 LQR
控制）：考虑经典的一阶倒立摆系统（小车上的倒立摆），线性化状态方程为
\(\dot{x} = A x + B u\)，其中
\(x = [\theta, \dot{\theta}, s, \dot{s}]^T\)（\(\theta\) 为摆角，\(s\)
为小车位置），\(u\) 为施加于小车的水平力。取
\(Q = \operatorname{diag}(100, 1, 10, 1)\)（强调角度偏差的惩罚）和
\(R = 0.1\)（控制能量惩罚），解代数 Riccati 方程得反馈增益矩阵
\(K\)，则最优控制律 \(u = -Kx\)
可在最小化能量消耗的同时将倒立摆稳定在竖直向上位置。闭环系统的极点由
\(A - BK\) 的特征值确定，满足稳定性要求。该例展示了 LQR
在机器人平衡控制中的典型应用。

\subsubsection{228.2 动态规划与 Hamilton-Jacobi-Bellman
方程}\label{ux52a8ux6001ux89c4ux5212ux4e0e-hamilton-jacobi-bellman-ux65b9ux7a0b}

\textbf{定义 228.2}（值函数 / Value Function）：从状态 \(x\) 和时间
\(t\) 出发的最优代价（值函数）为

\[V(t, x) = \min_{u(\cdot)} \left\{ \phi(x(t_f)) + \int_t^{t_f} L(s, x(s), u(s)) ds \right\}\]

\textbf{定理 228.2}（Bellman
最优性原理，1957）：最优策略的剩余部分仍是剩余问题的最优策略。用值函数表示为

\[V(t, x) = \min_{u \in U} \left\{ L(t, x, u) \Delta t + V(t + \Delta t, x + f(t, x, u) \Delta t) \right\} + o(\Delta t)\]

\textbf{定理 228.3}（Hamilton-Jacobi-Bellman 方程 / HJB 方程）：若
\(V \in C^1\)，则满足

\[-\frac{\partial V}{\partial t}(t, x) = \min_{u \in U} \left\{ L(t, x, u) + \nabla_x V(t, x)^T f(t, x, u) \right\}\]

边界条件 \(V(t_f, x) = \phi(x)\)。

\textbf{定理 228.4}（验证定理 / Verification Theorem）：若光滑函数
\(V(t, x)\) 满足 HJB 方程且 \(u^*(t, x)\) 达到最小值，则 \(u^*\)
是最优反馈控制。

\textbf{注 228.1}（极小值原理与 HJB 的关系）：极小值原理是 HJB
方程沿最优轨线的特征线形式。协态
\(\lambda(t) = \nabla_x V(t, x^*(t))\)。

\subsubsection{228.3
无限时域与稳定最优控制}\label{ux65e0ux9650ux65f6ux57dfux4e0eux7a33ux5b9aux6700ux4f18ux63a7ux5236}

\textbf{定义 228.3}（无限时域 LQR）：最小化
\(J = \int_0^\infty (x^T Q x + u^T R u) dt\)（\(Q \geq 0, R > 0\)）。最优控制为
\(u = -K x\)，\(K = R^{-1} B^T P\)，其中 \(P\) 满足\textbf{代数 Riccati
方程}：

\[PA + A^T P - P B R^{-1} B^T P + Q = 0\]

\textbf{定理 228.5}（代数 Riccati 方程的可解性）：若 \((A, B)\) 可镇定且
\((Q^{1/2}, A)\) 可检测，则代数 Riccati 方程存在唯一正半定解
\(P \geq 0\)，且闭环系统 \(A - BK\) 是渐近稳定的。

\bigskip\hrule\bigskip

\subsection{第229章：随机滤波与控制}\label{ux7b2c229ux7ae0ux968fux673aux6ee4ux6ce2ux4e0eux63a7ux5236}

在实际系统中，状态观测常受噪声干扰。随机滤波估计系统的真实状态，随机控制在此基础上设计最优策略。

\subsubsection{229.1 Kalman 滤波}\label{kalman-ux6ee4ux6ce2}

\textbf{定义 229.1}（线性随机系统）：具有噪声的线性系统

\[\begin{cases} \dot{x}(t) = A x(t) + B u(t) + w(t) \\ y(t) = C x(t) + v(t) \end{cases}\]

其中 \(w(t)\) 是过程噪声（协方差 \(Q \geq 0\)），\(v(t)\)
是观测噪声（协方差 \(R > 0\)），二者是不相关的 Gauss 白噪声。

\textbf{定义 229.2}（Kalman 滤波器，1960）：状态估计 \(\hat{x}(t)\)
由以下微分方程更新（连续时间 Kalman-Bucy 滤波器）：

\[\dot{\hat{x}} = A \hat{x} + B u + K (y - C \hat{x})\]

其中 \textbf{Kalman 增益} \(K = P C^T R^{-1}\)，且\textbf{误差协方差}
\(P(t) = \mathbb{E}[(x - \hat{x})(x - \hat{x})^T]\) 满足 \textbf{滤波
Riccati 方程}：

\[\dot{P} = A P + P A^T - P C^T R^{-1} C P + Q\]

\textbf{定理 229.1}（Kalman 滤波的最优性）：Kalman
滤波器是线性系统的最小方差无偏估计（BLUE）。在高斯噪声假设下，\(\hat{x}(t)\)
是状态 \(x(t)\) 在观测 \(\{y(s)\}_{s \leq t}\) 下的条件期望。

\textbf{定义 229.3}（离散时间 Kalman 滤波）：对离散时间系统
\(x_{k+1} = F x_k + B u_k + w_k\)，\(y_k = H x_k + v_k\)，Kalman
滤波由预测步和更新步交替组成：

\begin{itemize}
\tightlist
\item
  \textbf{预测步}：\(\hat{x}_{k|k-1} = F \hat{x}_{k-1|k-1} + B u_{k-1}\)，\(P_{k|k-1} = F P_{k-1|k-1} F^T + Q\)
\item
  \textbf{更新步}：\(K_k = P_{k|k-1} H^T (H P_{k|k-1} H^T + R)^{-1}\)，\(\hat{x}_{k|k} = \hat{x}_{k|k-1} + K_k (y_k - H \hat{x}_{k|k-1})\)
\end{itemize}

\subsubsection{229.2
随机最优控制}\label{ux968fux673aux6700ux4f18ux63a7ux5236}

\textbf{定义 229.4}（随机最优控制问题）：随机系统
\(dx_t = f(t, x_t, u_t) dt + \sigma(t, x_t) dW_t\)，性能指标为
\(J = \mathbb{E}[\phi(x_T) + \int_0^T L(t, x_t, u_t) dt]\)。

\textbf{定理 229.2}（随机 HJB 方程）：值函数 \(V(t, x)\) 满足

\[-\frac{\partial V}{\partial t} = \min_{u \in U} \left\{ L(t, x, u) + f^T \nabla_x V + \frac{1}{2} \operatorname{Tr}(\sigma \sigma^T \nabla_x^2 V) \right\}\]

\textbf{定理 229.3}（分离原理 / LQG 控制）：对线性随机系统（LQG：Linear
Quadratic Gaussian），最优控制 \(u = -K \hat{x}\) 分别由 LQR
的确定性最优增益 \(K\) 和 Kalman 滤波的状态估计 \(\hat{x}\) 独立决定。

\bigskip\hrule\bigskip

\subsection{第230章：鲁棒控制}\label{ux7b2c230ux7ae0ux9c81ux68d2ux63a7ux5236}

实际系统总存在模型不确定性（参数扰动、未建模动态）。鲁棒控制设计在最坏情况下保证系统稳定性和性能。

\subsubsection{230.1
小增益定理与不确定性建模}\label{ux5c0fux589eux76caux5b9aux7406ux4e0eux4e0dux786eux5b9aux6027ux5efaux6a21}

\textbf{定义 230.1}（\(\mathcal{H}_\infty\) 范数）：稳定传递函数
\(G(s)\) 的 \textbf{\(\mathcal{H}_\infty\) 范数}为

\[\|G\|_\infty = \sup_{\omega \in \mathbb{R}} \bar{\sigma}(G(j\omega))\]

其中 \(\bar{\sigma}(\cdot)\) 是最大奇异值。\(\|G\|_\infty\) 是 \(L^2\)
输入到 \(L^2\) 输出的诱导范数（最坏情况增益）。

\textbf{定理 230.1}（小增益定理 / Small Gain Theorem）：设系统 \(M\)
和不确定性 \(\Delta\) 为稳定传递函数。若
\(\|M \Delta\|_\infty < 1\)，则反馈互联系统稳定。更一般地，小增益条件是保证鲁棒稳定性的充分条件。

\textbf{定义 230.2}（结构奇异值 / \(\mu\)）：对结构化不确定性
\(\Delta = \operatorname{diag}(\Delta_1, \ldots, \Delta_k)\)，定义

\[\mu_\Delta(M) = \frac{1}{\min\{\bar{\sigma}(\Delta) : \det(I - M\Delta) = 0\}}\]

\(\mu < 1\) 等价于对所有允许的
\(\|\Delta\|_\infty \leq 1\)，系统保持稳定。

\subsubsection{\texorpdfstring{230.2 \(\mathcal{H}_\infty\)
最优控制}{230.2 \textbackslash mathcal\{H\}\_\textbackslash infty 最优控制}}\label{mathcalh_infty-ux6700ux4f18ux63a7ux5236}

\textbf{定义 230.3}（\(\mathcal{H}_\infty\) 控制问题）：考虑广义被控对象

\[\begin{bmatrix} z \\ y \end{bmatrix} = \begin{bmatrix} P_{11} & P_{12} \\ P_{21} & P_{22} \end{bmatrix} \begin{bmatrix} w \\ u \end{bmatrix}\]

其中 \(w\) 是外部扰动，\(u\) 是控制，\(z\) 是被控输出，\(y\)
是测量。寻找控制器 \(K(s)\)：\(u = K y\)，使得闭环传递函数 \(T_{zw}(s)\)
满足 \(\|T_{zw}\|_\infty < \gamma\)。

\textbf{定理 230.2}（Doyle-Glover-Khargonekar-Francis 解法 / DGKF
1989）：\(\mathcal{H}_\infty\) 次优控制问题可归结为两个代数 Riccati
方程（一个用于状态反馈，一个用于输出估计）的可解性。控制器由中心解给出。

\textbf{定理 230.3}（有界实引理 / Bounded Real
Lemma）：\(\|G\|_\infty < \gamma\) 当且仅当 Hamiltonian
矩阵没有虚轴上的特征值，等价于某线性矩阵不等式（LMI）可解。

\subsubsection{230.3
线性矩阵不等式}\label{ux7ebfux6027ux77e9ux9635ux4e0dux7b49ux5f0f}

\textbf{定义 230.4}（线性矩阵不等式 / LMI）：\textbf{线性矩阵不等式}形如
\(F(x) = F_0 + \sum_{i=1}^m x_i F_i \geq 0\)（半正定），其中 \(F_i\)
为对称矩阵。

\textbf{定理 230.4}（LMI
方法的优势）：许多控制问题（\(\mathcal{H}_\infty\)
控制、\(\mathcal{H}_2\) 控制、极点配置、\(\mu\)-综合的上界）都可转化为
LMI 可行性问题。LMI 是凸优化问题，可用内点法高效求解。

\textbf{定理 230.5}（Schur 补引理）：分块对称矩阵正定当且仅当

\[\begin{bmatrix} Q & S \\ S^T & R \end{bmatrix} > 0 \iff Q > 0,\ R - S^T Q^{-1} S > 0\]

Schur 补将非线性 Riccati 不等式转化为线性矩阵不等式。

\bigskip\hrule\bigskip

\subsection{第231章：非线性控制}\label{ux7b2c231ux7ae0ux975eux7ebfux6027ux63a7ux5236}

大多数真实系统是非线性的。非线性控制利用微分几何和 Lyapunov
方法处理非线性动态。

\subsubsection{231.1 Lyapunov
稳定性理论}\label{lyapunov-ux7a33ux5b9aux6027ux7406ux8bba}

\textbf{定理 231.1}（Lyapunov 直接法）：对系统
\(\dot{x} = f(x)\)（\(f(0) = 0\)），若存在 \(C^1\) 函数
\(V: \mathbb{R}^n \to \mathbb{R}\)，满足

\begin{itemize}
\tightlist
\item
  \(V(0) = 0\)，且 \(V(x) > 0\)（\(x \neq 0\)）（正定性）
\item
  \(\dot{V}(x) = \nabla V(x)^T f(x) < 0\)（\(x \neq 0\)）（负定性）
\end{itemize}

则原点 \(x = 0\) 是渐近稳定的。\(V(x)\) 称为 \textbf{Lyapunov 函数}。

\textbf{定理 231.2}（LaSalle 不变原理）：若 \(V(x) \geq 0\) 且
\(\dot{V}(x) \leq 0\)，则所有有界轨线收敛到最大不变集
\(M \subseteq \{x : \dot{V}(x) = 0\}\)。

\textbf{定义 231.1}（控制 Lyapunov 函数 / CLF）：对受控系统
\(\dot{x} = f(x) + g(x) u\)，光滑正定函数 \(V(x)\) 是\textbf{控制
Lyapunov 函数}，如果对所有 \(x \neq 0\)，

\[\inf_{u} \left\{ L_f V(x) + L_g V(x) \cdot u \right\} < 0\]

其中 \(L_f V = \nabla V^T f\) 和 \(L_g V = \nabla V^T g\) 是 Lie 导数。

\subsubsection{231.2 反馈线性化}\label{ux53cdux9988ux7ebfux6027ux5316}

\textbf{定义 231.2}（Lie 括号与能控性分布）：向量场 \(f\) 和 \(g\)
的\textbf{Lie 括号}为
\([f, g] = \nabla g \cdot f - \nabla f \cdot g\)。所有 \(g\) 及其迭代
Lie 括号张成的分布 \(\Delta\) 是系统的\textbf{可达分布}。

\textbf{定理 231.3}（输入-状态反馈线性化条件）：仿射非线性系统
\(\dot{x} = f(x) + g(x) u\) 可经状态变换 \(z = \Phi(x)\) 和反馈
\(u = \alpha(x) + \beta(x) v\) 化为线性系统 \(\dot{z} = A z + B v\)
的必要条件是
\(\{g, \operatorname{ad}_f g, \ldots, \operatorname{ad}_f^{n-1} g\}\) 在
\(x_0\) 附近线性独立且对合（involutive）。

\textbf{定义 231.3}（相对阶 / Relative Degree）：系统
\(\dot{x} = f(x) + g(x) u\)，\(y = h(x)\) 在 \(x_0\) 具有\textbf{相对阶}
\(r\)，如果 \(L_g L_f^{k} h(x) = 0\)（\(k = 0, \ldots, r-2\)）且
\(L_g L_f^{r-1} h(x_0) \neq 0\)。相对阶是输出被输入直接影响的延迟阶数。

\subsubsection{231.3 滑模控制}\label{ux6ed1ux6a21ux63a7ux5236}

\textbf{定义 231.4}（滑模控制 / Sliding Mode
Control）：选择\textbf{滑动流形} \(s(x) = 0\)（如
\(s = \dot{e} + \lambda e\)，其中 \(e = x - x_d\)
是跟踪误差）。控制律设计为使系统轨线在有限时间内到达并保持在滑动流形上：

\[u = u_{\text{eq}} - k \operatorname{sgn}(s)\]

其中 \(u_{\text{eq}}\) 是等效控制（\(\dot{s} = 0\)
的解），\(-k \operatorname{sgn}(s)\) 是开关项（强制到达滑动面）。

\textbf{定理 231.4}（到达条件与有限时间收敛）：若
\(s \dot{s} \leq -\eta |s|\)（\(\eta > 0\)），则 \(s \to 0\) 在有限时间
\(T \leq |s(0)|/\eta\)
内。在滑模面上的运动对匹配的不确定性/扰动具有完全鲁棒性。

\subsubsection{231.4 反步法}\label{ux53cdux6b65ux6cd5}

\textbf{定义 231.5}（反步法 /
Backstepping）：对严格反馈系统（strict-feedback form）：

\[\begin{aligned} \dot{x}_1 &= f_1(x_1) + g_1(x_1) x_2 \\ \dot{x}_2 &= f_2(x_1, x_2) + g_2(x_1, x_2) x_3 \\ &\vdots \\ \dot{x}_n &= f_n(x_1, \ldots, x_n) + g_n(x_1, \ldots, x_n) u \end{aligned}\]

反步法递归地构造 Lyapunov 函数 \(V_i(x_1, \ldots, x_i)\) 和虚拟控制
\(\alpha_i\)，最终得到真实控制 \(u\)。

\textbf{定理 231.5}（反步法的 Lyapunov 构造）：在第 \(i\) 步，令
\(z_i = x_i - \alpha_{i-1}\) 为误差变量。选择
\(V_i = V_{i-1} + \frac{1}{2} z_i^2\)，设计 \(\alpha_i\) 使得
\(\dot{V}_i < 0\)。这递归地给出全局渐近稳定的控制律。

\bigskip\hrule\bigskip

\emph{本卷在分析学（V8、V13）、微分方程（V17）和优化理论（V26）的基础上，系统建立了线性控制（可控性与可观测性）、最优控制（Pontryagin
原理与 Bellman 方程）、随机滤波（Kalman 滤波与 LQG）、鲁棒控制（H∞ 与
μ-综合）和非线性控制（反馈线性化、滑模与反步法）的理论体系。共 5 章。}

\newpage

\section{卷四十一：博弈论}\label{ux5377ux56dbux5341ux4e00ux535aux5f08ux8bba}

\begin{quote}
博弈论是研究多个理性决策者之间策略互动的数学理论，由 John von Neumann 和
Oskar Morgenstern 在《博弈论与经济行为》（1944）中奠基，John Nash 在
1950
年代建立了非合作博弈的均衡概念。博弈论不仅深刻影响了经济学（多个诺贝尔经济学奖），也是计算机科学（算法博弈论、拍卖理论）、政治学（投票理论、机制设计）和演化生物学（演化博弈论）的基础工具。本卷在概率论（V21）、优化理论（V27）和拓扑学（V11）的基础上，建立非合作博弈、合作博弈、机制设计、演化博弈和随机博弈的系统理论。
\end{quote}

\bigskip\hrule\bigskip

\subsection{第232章：非合作博弈}\label{ux7b2c232ux7ae0ux975eux5408ux4f5cux535aux5f08}

非合作博弈研究每个参与者独立选择策略以最大化自身收益的情形。John Nash
的均衡概念是非合作博弈理论的核心。

\subsubsection{232.1 策略型博弈与 Nash
均衡}\label{ux7b56ux7565ux578bux535aux5f08ux4e0e-nash-ux5747ux8861}

\textbf{定义 232.1}（策略型博弈 / Normal-Form Game）：一个 \(n\)
人\textbf{策略型博弈}
\(\Gamma = (N, \{S_i\}_{i \in N}, \{u_i\}_{i \in N})\) 由以下要素构成：
- \(N = \{1, \ldots, n\}\)：参与者集合 - \(S_i\)：参与者 \(i\)
的\textbf{纯策略}集合 -
\(u_i: S = \prod_{j \in N} S_j \to \mathbb{R}\)：参与者 \(i\)
的\textbf{支付函数}

当参与者的选择是随机的时，\textbf{混合策略}是 \(S_i\) 上的概率分布
\(\sigma_i \in \Delta(S_i)\)。

\textbf{定义 232.2}（Nash 均衡，1950）：策略剖面
\(\sigma^* = (\sigma_1^*, \ldots, \sigma_n^*)\) 是 \textbf{Nash
均衡}，如果对每个参与者 \(i\) 和每个策略 \(\sigma_i\)，

\[u_i(\sigma_i^*, \sigma_{-i}^*) \geq u_i(\sigma_i, \sigma_{-i}^*)\]

即没有参与者可以通过单方面偏离来提高自己的收益。

\textbf{定理 232.1}（Nash
的存在定理，1951）：任何有限策略型博弈至少存在一个 Nash
均衡（可能为混合策略）。

\emph{证明思路}：将均衡刻画为不动点。混合策略空间
\(\Delta = \prod_i \Delta(S_i)\)
是有限维单纯形的笛卡尔积，故为非空凸紧集。对每个
\(\sigma_{-i}\)，最优反应对应
\(B_i(\sigma_{-i}) = \arg\max_{\sigma_i \in \Delta(S_i)} u_i(\sigma_i, \sigma_{-i})\)
的图为闭集（由 \(u_i\) 在混合策略下的多线性连续性），且因期望支付对
\(\sigma_i\) 为线性函数，最优反应集为某线性函数在紧凸集上的极大点集，故
\(B_i(\sigma_{-i})\) 为非空凸集。取
\(B(\sigma) = \prod_i B_i(\sigma_{-i})\)，则 \(B: \Delta \to 2^\Delta\)
满足：值域在紧凸集 \(\Delta\) 内，图闭（上半连续），且对每个
\(\sigma\)，\(B(\sigma)\) 为非空凸集。由 Kakutani
不动点定理，\(\exists \sigma^* \in B(\sigma^*)\)，此即 Nash 均衡。∎

\textbf{例 232.1}（囚徒困境 / Prisoner's Dilemma）：两个嫌疑人各有策略
C（合作/不招供）和 D（背叛/招供）。支付矩阵为

{\def\LTcaptype{none} % do not increment counter
\begin{longtable}[]{@{}lll@{}}
\toprule\noalign{}
& C & D \\
\midrule\noalign{}
\endhead
\bottomrule\noalign{}
\endlastfoot
C & (-1, -1) & (-3, 0) \\
D & (0, -3) & (-2, -2) \\
\end{longtable}
}

唯一的 Nash 均衡是 (D, D)，但 (C, C) 是 Pareto
更优的------这是个体理性与集体理性冲突的经典例证。

\subsubsection{232.2 均衡的精炼}\label{ux5747ux8861ux7684ux7cbeux70bc}

\textbf{定义 232.3}（子博弈完美均衡 / Subgame Perfect
Equilibrium，Selten 1965）：扩展型博弈（博弈树）中策略剖面 \(\sigma\)
是\textbf{子博弈完美均衡}，如果它在每个子博弈中都是 Nash
均衡。这排除了不可信威胁（non-credible threats）。

\textbf{定义 232.4}（颤抖手完美均衡 / Trembling Hand Perfect
Equilibrium，Selten 1975）：策略剖面 \(\sigma^*\)
是\textbf{颤抖手完美均衡}，如果存在一列完全混合策略剖面
\(\sigma^k \to \sigma^*\)，使得 \(\sigma_i^*\) 对每个 \(\sigma_{-i}^k\)
都是最优反应。

\textbf{定义 232.5}（序贯均衡 / Sequential Equilibrium，Kreps-Wilson
1982）：序贯均衡要求在信息集上的信念是通过 Bayes
法则从均衡策略推导出来的（一致性），且每个参与者的策略是序贯理性的。

\subsubsection{232.3
贝叶斯博弈与不完全信息}\label{ux8d1dux53f6ux65afux535aux5f08ux4e0eux4e0dux5b8cux5168ux4fe1ux606f}

\textbf{定义 232.6}（贝叶斯博弈 / Harsanyi 1967-1968）：\textbf{Bayes
博弈}中每个参与者 \(i\) 有私有的\textbf{类型}
\(\theta_i \in \Theta_i\)，类型的联合分布 \(p(\theta)\)
是共同知识。\(i\) 的策略是函数 \(s_i: \Theta_i \to S_i\)。

\textbf{定义 232.7}（Bayes-Nash 均衡）：策略剖面 \(s^*\) 是
\textbf{Bayes-Nash 均衡}，如果对每个 \(i\) 和 \(\theta_i\)，

\[\mathbb{E}_{\theta_{-i}|\theta_i} [u_i(s_i^*(\theta_i), s_{-i}^*(\theta_{-i}), \theta)] \geq \mathbb{E}_{\theta_{-i}|\theta_i}[u_i(s_i, s_{-i}^*(\theta_{-i}), \theta)]\]

\subsubsection{232.4
相关均衡与通信均衡}\label{ux76f8ux5173ux5747ux8861ux4e0eux901aux4fe1ux5747ux8861}

\textbf{定义 232.8}（相关均衡 / Correlated Equilibrium，Aumann
1974）：概率分布 \(\mu \in \Delta(S)\) 是\textbf{相关均衡}，如果对所有
\(i\) 和所有映射 \(d_i: S_i \to S_i\)，

\[\sum_{s \in S} \mu(s) u_i(s) \geq \sum_{s \in S} \mu(s) u_i(d_i(s_i), s_{-i})\]

相关均衡允许参与者根据相关信号协调选择（比 Nash 均衡更一般）。

\textbf{定理 232.2}（Aumann
相关均衡定理，1974）：相关均衡集是凸紧集，包含所有 Nash
均衡的凸包。任何相关均衡均可通过引入有限状态的通信设备（correlation
device）------即发送给各参与者的私有随机信号------来实现，使得参与者根据接收到的信号理性选择行动构成相关均衡。

\bigskip\hrule\bigskip

\subsection{第233章：合作博弈与核}\label{ux7b2c233ux7ae0ux5408ux4f5cux535aux5f08ux4e0eux6838}

合作博弈研究参与者可以形成联盟并通过有约束力的协议分配合作收益的情形。

\subsubsection{233.1
特征函数型博弈}\label{ux7279ux5f81ux51fdux6570ux578bux535aux5f08}

\textbf{定义 233.1}（特征函数型博弈 / Coalitional Game）：\(n\)
人\textbf{特征函数型博弈} \((N, v)\) 由特征函数
\(v: 2^N \to \mathbb{R}\) 定义，\(v(S)\) 是联盟 \(S\)
可以保证获得的最大收益（\(v(\emptyset) = 0\)）。博弈是\textbf{超加性的}（superadditive），如果
\(v(S \cup T) \geq v(S) + v(T)\)（\(S \cap T = \emptyset\)）。

\textbf{定义 233.2}（转归 / Imputation）：\textbf{转归}是支付向量
\(x = (x_1, \ldots, x_n)\) 满足： -
个体理性：\(x_i \geq v(\{i\})\)（每个参与者至少获得单独行动的收益） -
集体理性 / 效率：\(\sum_{i=1}^n x_i = v(N)\)（总收益全部分配）

\subsubsection{233.2 核与稳定集}\label{ux6838ux4e0eux7a33ux5b9aux96c6}

\textbf{定义 233.3}（核 / Core，Gillies 1953）：博弈 \((N, v)\)
的\textbf{核}是满足联盟理性的转归集合：

\[C(v) = \left\{ x \in \mathbb{R}^n : \sum_{i \in N} x_i = v(N),\; \sum_{i \in S} x_i \geq v(S) \; (\forall S \subset N) \right\}\]

核中的分配不会被任何联盟通过单独行动改进。

\textbf{定理 233.1}（Bondareva-Shapley
定理，1963）：博弈的核非空当且仅当博弈是\textbf{平衡的}（balanced）：对所有平衡权
\(\{\lambda_S\}_{S \subseteq N}\)（\(\sum_{S \ni i} \lambda_S = 1\)），有
\(\sum_S \lambda_S v(S) \leq v(N)\)。

\textbf{定义 233.4}（Shapley 值，1953）：参与者 \(i\) 的 \textbf{Shapley
值}是对其边际贡献的期望：

\[\phi_i(v) = \sum_{S \subseteq N \setminus \{i\}} \frac{|S|! (n - |S| - 1)!}{n!} \left[ v(S \cup \{i\}) - v(S) \right]\]

\textbf{定理 233.2}（Shapley 值的公理化）：Shapley
值是满足以下公理的唯一解概念：效率（\(\sum_i \phi_i(v) = v(N)\)）、对称性（对等参与者获得相同支付）、虚拟参与者公理（零边际贡献的参与者获零）和可加性（\(\phi(v + w) = \phi(v) + \phi(w)\)）。

\subsubsection{233.3
核仁与谈判集}\label{ux6838ux4ec1ux4e0eux8c08ux5224ux96c6}

\textbf{定义 233.5}（超额与核仁 / Nucleolus，Schmeidler 1969）：联盟
\(S\) 对分配 \(x\) 的\textbf{超额}为
\(e(x, S) = v(S) - \sum_{i \in S} x_i\)。\textbf{核仁}是按词典序最小化超额向量（从大到小排列）的唯一分配点。

\textbf{定义 233.6}（谈判集 / Bargaining Set，Aumann-Maschler
1964）：谈判集是对异议的稳定性要求。若 \(i\) 对 \(j\)
提出异议（带替代联盟），\(j\) 能提出反异议，则 \(x\) 在谈判集中。

\bigskip\hrule\bigskip

\subsection{第234章：机制设计与拍卖}\label{ux7b2c234ux7ae0ux673aux5236ux8bbeux8ba1ux4e0eux62cdux5356}

机制设计是「反向博弈论」：给定社会目标，设计博弈规则使得参与者在追求自身利益时达到期望结果。

\subsubsection{234.1 社会选择与 Gibbard-Satterthwaite
定理}\label{ux793eux4f1aux9009ux62e9ux4e0e-gibbard-satterthwaite-ux5b9aux7406}

\textbf{定义 234.1}（社会选择函数 / SCF）：\textbf{社会选择函数}
\(f: \Theta \to X\) 将偏好剖面映射到社会结果。\(\Theta_i\) 是参与者
\(i\) 的类型空间，\(X\) 是备选结果集。

\textbf{定义 234.2}（激励相容 / Incentive Compatibility）：直接显示机制
\((f, t)\)（\(t_i\)
是转移支付）是\textbf{激励相容}的，如果真实报告类型是 Bayes-Nash
均衡：对所有的 \(i\)、\(\theta_i\)、\(\theta_i'\)，

\[\mathbb{E}_{\theta_{-i}}[u_i(f(\theta_i, \theta_{-i}), \theta_i) + t_i(\theta_i, \theta_{-i})] \geq \mathbb{E}_{\theta_{-i}}[u_i(f(\theta_i', \theta_{-i}), \theta_i) + t_i(\theta_i', \theta_{-i})]\]

\textbf{定理 234.1}（Gibbard-Satterthwaite 定理，1973/1975）：若
\(|X| \geq 3\) 且 \(f\)
的值域包含至少三个结果，则任何非独裁的社会选择函数都是可操纵的（即非激励相容的）。

\subsubsection{234.2 VCG
机制与有效分配}\label{vcg-ux673aux5236ux4e0eux6709ux6548ux5206ux914d}

\textbf{定义 234.3}（Vickrey-Clarke-Groves
机制，1961/1971/1973）：\textbf{VCG 机制}选择有效分配
\(x^*(\theta) = \arg\max_x \sum_i v_i(x, \theta_i)\)，转移支付为

\[t_i^{\text{VCG}}(\theta) = \sum_{j \neq i} v_j(x^*(\theta), \theta_j) - \max_x \sum_{j \neq i} v_j(x, \theta_j)\]

\textbf{定理 234.2}（VCG 的激励相容性）：VCG
机制是策略性的（dominant-strategy incentive
compatible），即真实报告是每个参与者的占优策略。

\textbf{定理 234.3}（Myerson-Satterthwaite
不可能定理，1983）：在双边贸易中（一个买方、一个卖方），不存在同时满足激励相容、个体理性和预算平衡且事后有效的机制。

\subsubsection{234.3
最优拍卖设计}\label{ux6700ux4f18ux62cdux5356ux8bbeux8ba1}

\textbf{定义 234.4}（Myerson
最优拍卖，1981）：卖方有一个不可分物品。买方 \(i\) 的估价 \(v_i\)
独立来自分布 \(F_i\)。\textbf{虚拟价值}为
\(\psi_i(v_i) = v_i - \frac{1 - F_i(v_i)}{f_i(v_i)}\)。最优拍卖将物品分配给虚拟价值最高的买方，保留价
\(r^* = \psi_i^{-1}(0)\)。

\textbf{定理 234.5}（Myerson
最优拍卖的推导概要）：卖方的期望收入可写为各买方虚拟剩余的期望和：

\[\mathbb{E}\left[\sum_i \psi_i(v_i) x_i(v)\right]\]

其中 \(x_i(v) \in \{0, 1\}\)
为分配规则，\(\psi_i(v_i) = v_i - \frac{1-F_i(v_i)}{f_i(v_i)}\)
为\textbf{虚拟估值函数}（virtual
valuation）。推导要点：由激励相容条件导出包络定理------买方的期望支付由其分配概率积分确定；将期望支付代入卖方目标函数后，收入仅依赖于分配规则
\(x_i(v)\) 和虚拟估值。最优分配规则为：将物品分配给
\(\max_i \psi_i(v_i)\) 的买方（若最高虚拟估值 \(<\) 卖方保留价值
\(0\)，则物品不售出）。\textbf{保留价} \(r_i^* = \psi_i^{-1}(0)\)
保证卖方不对虚拟估值为负的买方售出。当买方对称时，最优拍卖化简为含最优保留价的二价拍卖（Vickrey
1961）------这是收入等价定理的直接推论。Myerson
的理论统一了最优拍卖、最优垄断定价和非线性定价等机制设计问题。

\textbf{定理 234.4}（收入等价定理 / Revenue
Equivalence）：任何两个拍卖机制若分配函数相同且最低类型的期望支付相同，则卖方的期望收入相同。

\bigskip\hrule\bigskip

\subsection{第235章：演化博弈论}\label{ux7b2c235ux7ae0ux6f14ux5316ux535aux5f08ux8bba}

演化博弈论将博弈均衡解释为演化过程的稳定状态，将生物学的自然选择思想与博弈论结合。

\subsubsection{235.1
演化稳定策略}\label{ux6f14ux5316ux7a33ux5b9aux7b56ux7565}

\textbf{定义 235.1}（演化稳定策略 / ESS，Maynard Smith-Price
1973）：在大种群（随机配对）中，策略 \(x\)
是\textbf{演化稳定策略}，如果对任何变异策略 \(y \neq x\)，存在
\(\bar{\varepsilon} > 0\) 使得对所有
\(\varepsilon \in (0, \bar{\varepsilon})\)：

\[u(x, (1-\varepsilon)x + \varepsilon y) > u(y, (1-\varepsilon)x + \varepsilon y)\]

即 \(x\) 对自身（包括少量变异入侵后）的表现优于变异策略。

\textbf{定理 235.1}（ESS 的等价条件）：\(x\) 是 ESS 当且仅当： 1.
\((x, x)\) 是 Nash 均衡：\(u(x, x) \geq u(y, x)\)（对所有 \(y\)） 2. 若
\(u(x, x) = u(y, x)\)（\(y\) 是对 \(x\) 的备选最优反应），则
\(u(x, y) > u(y, y)\)

\subsubsection{235.2 复制者动态}\label{ux590dux5236ux8005ux52a8ux6001}

\textbf{定义 235.2}（复制者方程 / Replicator Dynamics，Taylor-Jonker
1978）：设 \(x_i\) 是种群中使用纯策略 \(i\)
的比例。\textbf{复制者方程}为

\[\dot{x}_i = x_i \left( u(i, x) - \bar{u}(x) \right)\]

其中 \(u(i, x)\) 是纯策略 \(i\) 对混合策略种群 \(x\)
的期望收益，\(\bar{u}(x) = \sum_j x_j u(j, x)\) 是平均收益。

\textbf{定理 235.2}（复制者动态的性质）： - 单纯形 \(\Delta\)
是正向不变集 - Nash 均衡是动态系统的平衡点 -
稳定的平衡点（渐近稳定的静止点）是 ESS - 被严格占优的策略最终被淘汰

\textbf{定理 235.3}（Folk Theorem of Evolutionary Game Theory）：ESS
在单纯形内部是复制者动态的渐近稳定平衡点；边界 ESS
需要额外条件（如对进入单纯形的内部轨线具有局部吸引性）才能保证渐近稳定性。反过来，渐近稳定平衡点是
Nash 均衡，但不一定是 ESS。

\subsubsection{235.3
演化博弈与学习}\label{ux6f14ux5316ux535aux5f08ux4e0eux5b66ux4e60}

\textbf{定义 235.3}（虚拟博弈 / Fictitious Play，Brown 1951）：参与者
\(i\)
假设对手使用历史经验分布，在每步选择对过去对手策略分布的最优反应。这是最简单学习模型。

\textbf{定理
235.4}（零和博弈中的虚拟博弈）：在有限二人零和博弈中，虚拟博弈的对手策略经验分布收敛到
Nash 均衡策略。

\textbf{定义 235.4}（无悔学习与 Regret
Matching）：参与者在每轮后计算各策略的累积后悔（regret），以与正后悔成比例的概率选择策略。Hart-Mas-Colell（2000）证明了无悔学习使经验分布收敛到相关均衡集。

\bigskip\hrule\bigskip

\subsection{第236章：随机博弈与重复博弈}\label{ux7b2c236ux7ae0ux968fux673aux535aux5f08ux4e0eux91cdux590dux535aux5f08}

在多期博弈中，参与者可以观察到过去的行为并据此调整策略。重复博弈研究长期的合作与冲突。

\subsubsection{236.1 重复博弈与 Folk
定理}\label{ux91cdux590dux535aux5f08ux4e0e-folk-ux5b9aux7406}

\textbf{定义 236.1}（无限重复博弈）：阶段博弈 \(G\) 无限重复，贴现因子
\(\delta \in (0, 1)\)。参与者的总收益是各阶段收益的贴现和：\(U_i = (1 - \delta) \sum_{t=0}^\infty \delta^t u_i(a^t)\)。

\textbf{定理 236.1}（Nash Folk 定理，Friedman
1971）：对无限重复博弈，任何可行且个体理性的支付向量（超过极小极大支付）都可以作为某个
\(\delta\) 充分接近 \(1\) 时的子博弈完美均衡支付。

\textbf{定理 236.2}

\textbf{定理 236.1b}（极限平均支付下的 Folk 定理 /
Aumann-Shapley-Rubinstein）：在无限重复博弈中，若采用\textbf{极限平均支付准则}（\(U_i = \liminf_{T \to \infty} \frac{1}{T} \sum_{t=0}^{T-1} u_i(a^t)\)），则任何\textbf{可行且严格个体理性}的支付向量（超过各参与者的极小极大支付）均可作为
Nash 均衡实现，且对任意 \(\varepsilon > 0\) 存在子博弈完美均衡使支付
\(\varepsilon\)-逼近该向量。证明构造：将目标支付分解为阶段博弈行动序列的凸组合，参与者相互监控对方是否遵循预定行动轨迹；一旦检测到偏离，则永久切换到极小极大惩罚阶段。极限平均准则（无贴现
\(\delta = 1\)）下惩罚成本为零------与贴现情形（\(\delta < 1\)
需足够大以确保惩罚可置信）相比，Folk 定理的结论更强。 （Neyman
定理，1985/1998）：在有限重复博弈中，若参与者有界理性（使用有限自动机实现策略），合作可作为
\(\varepsilon\)-均衡出现，即使这是 Nash Folk 定理不适用的情况。

\textbf{定义 236.2}（触发策略 / Trigger
Strategies）：\textbf{冷酷触发策略}（Grim
Trigger）：对方偏离则永远合作终止。\(\delta\)
充分大时，触发策略构成子博弈完美均衡。

\subsubsection{236.2 随机博弈}\label{ux968fux673aux535aux5f08}

\textbf{定义 236.3}（随机博弈 / Stochastic Game，Shapley
1953）：\textbf{随机博弈}由状态空间
\(\Omega\)、各状态下的阶段博弈和状态转移概率 \(p(\omega' | \omega, a)\)
组成。贴现因子为 \(\delta\)。

\textbf{定理 236.3}（Shapley
的二人零和随机博弈）：在有限状态、有限行动的二人零和随机博弈中，存在唯一的值
\(v_\delta(\omega)\)，且双方有最优稳定策略。值函数满足 Shapley 方程：

\[v(\omega) = \operatorname{Val}\left[ (1-\delta) r(\omega, a) + \delta \sum_{\omega'} p(\omega' | \omega, a) v(\omega') \right]\]

\textbf{定理 236.4}（Mertens-Neyman
定理，1981）：在有限二人零和随机博弈中，对于接近 \(1\)
的贴现因子，值函数 \(v(\omega)\) 关于 \(\delta\) 一致有界。

\subsubsection{236.3
重复博弈中的声誉效应}\label{ux91cdux590dux535aux5f08ux4e2dux7684ux58f0ux8a89ux6548ux5e94}

\textbf{定理 236.5}（Kreps-Wilson-Milgrom-Roberts
声誉定理，1982）：在有限重复囚徒困境中，若存在小概率的不理性类型（``针锋相对''型），理性参与者在博弈早期也会合作（建立声誉），直到最后几期才会背叛。

\textbf{定理 236.6}（Fudenberg-Levine
定理，1989）：单个长期参与者面对一系列短期参与者时，即使承诺类型概率极小，长期参与者在充分长的博弈中也能获得接近
Stackelberg 收益的支付。

\bigskip\hrule\bigskip

\emph{本卷在概率论（V21）、优化（V27）和拓扑学（V11）的基础上，系统建立了非合作博弈（Nash
均衡、均衡精炼、Bayes 博弈）、合作博弈（核、Shapley
值、谈判集）、机制设计（VCG、最优拍卖）、演化博弈论（ESS、复制者动态）和随机博弈的理论体系。共
5 章。}

\newpage

\section{卷四十二：数学生物学}\label{ux5377ux56dbux5341ux4e8cux6570ux5b66ux751fux7269ux5b66}

\begin{quote}
数学生物学使用微分方程、动力系统、随机过程和网络理论等数学工具研究生命现象。从古典的种群动力学（Lotka-Volterra
方程）到现代的基因组学和系统生物学，数学在理解生物学过程中发挥了不可替代的作用。
\end{quote}

\begin{quote}
\textbf{前置知识}：本卷依赖 \textbf{V17 常微分方程论}（ODE
定性理论、相平面分析、稳定性理论）和 \textbf{V26
优化理论}（线性规划、凸优化、最优控制），并在动力系统（V28）和概率论（V21）的基础上展开。建议读者先掌握
ODE 的分岔理论（V17 Ch 90-92）和 Pontryagin 极大值原理（V26 Ch
136）后再阅读本卷的传染病最优控制（Ch 238.3）和神经网络动力学（Ch
241.3）。
\end{quote}

\bigskip\hrule\bigskip

\subsection{第237章：种群动力学}\label{ux7b2c237ux7ae0ux79cdux7fa4ux52a8ux529bux5b66}

种群动力学是数学生物学最古老的分支，用微分方程和差分方程描述种群数量随时间的变化。

\subsubsection{237.1 单种群模型}\label{ux5355ux79cdux7fa4ux6a21ux578b}

\textbf{定义 237.1}（Malthus
增长模型，1798）：无限制的指数增长：\(\dot{N} = r N\)，\(N(t) = N_0 e^{rt}\)。\(r\)
是内禀增长率。

\textbf{定义 237.2}（Logistic 方程 / Verhulst
模型，1838）：有环境容量的增长：

\[\dot{N} = r N \left( 1 - \frac{N}{K} \right)\]

其中 \(K\) 是\textbf{环境容量}（carrying capacity）。解为 S
形曲线：\(N(t) = \frac{K N_0 e^{rt}}{K + N_0(e^{rt} - 1)}\)。\(N = K\)
是全局渐近稳定的平衡点。

\textbf{定义 237.3}（具有 Allee 效应的种群模型）：\textbf{Allee
效应}（种群密度过低反而不利于增长）：

\[\dot{N} = r N (N - A) \left( 1 - \frac{N}{K} \right)\]

其中 \(0 < A < K\)。\(N = 0\) 和 \(N = K\) 是稳定平衡点，\(N = A\)
是不稳定阈值。

\textbf{定义 237.4}（离散时间种群模型 / Ricker
模型）：\(N_{t+1} = N_t e^{r(1 - N_t/K)}\)。当 \(r > 2\)
时出现倍周期分岔，\(r\) 足够大时出现混沌。

\subsubsection{237.2
多种群相互作用}\label{ux591aux79cdux7fa4ux76f8ux4e92ux4f5cux7528}

\textbf{定义 237.5}（Lotka-Volterra 捕食-被捕食模型，1925-1926）：

\[\begin{cases} \dot{x} = \alpha x - \beta x y \\ \dot{y} = \gamma x y - \delta y \end{cases}\]

其中 \(x\) 是被捕食者，\(y\) 是捕食者。系统的解是周期轨道，存在首次积分
\(H(x, y) = \gamma x - \delta \ln x + \beta y - \alpha \ln y\)。

\textbf{定理 237.1}（Lotka-Volterra 模型的性质）：平衡点
\((\delta/\gamma, \alpha/\beta)\)
是中心（非双曲），所有正初始条件的轨线是周期轨道。这是保守系统，无结构稳定性。

\emph{证明概要}：构造首次积分
\(H(x,y) = \gamma x - \delta\ln x + \beta y - \alpha\ln y\)，直接验证
\(\frac{d}{dt}H(x,y) = \gamma\dot{x} - \delta\frac{\dot{x}}{x} + \beta\dot{y} - \alpha\frac{\dot{y}}{y} = 0\)。相轨线由
\(H(x,y)\)
的等值线给出，围绕平衡点形成闭曲线（一阶积分保证轨线在相平面上为闭轨，无极限环）。周期由椭圆型积分表出。∎

\textbf{定义 237.6}（竞争模型 / Competition Models）：\(n\)
个物种竞争共享资源：

\[\dot{N}_i = r_i N_i \left( 1 - \sum_{j=1}^n \alpha_{ij} N_j / K_i \right)\]

其中 \(\alpha_{ij}\) 是竞争系数（物种 \(j\) 对物种 \(i\) 的影响）。

\textbf{定理 237.2}（竞争排斥原理 / Gause
定律）：若两个物种的生态位完全重叠，其中一个最终会排斥另一个（灭绝）。更形式化地，若
\(\alpha_{12} < K_1/K_2\) 且 \(\alpha_{21} > K_2/K_1\)，物种 2 被排斥。

\subsubsection{237.3
年龄结构与积分-差分方程}\label{ux5e74ux9f84ux7ed3ux6784ux4e0eux79efux5206-ux5deeux5206ux65b9ux7a0b}

\textbf{定义 237.7}（McKendrick-von Foerster 方程 / 年龄结构方程）：

\[\frac{\partial n}{\partial t} + \frac{\partial n}{\partial a} = -\mu(a) n(t, a)\]

边界条件
\(n(t, 0) = \int_0^\infty \beta(a) n(t, a) da\)（出生过程），其中
\(\mu(a)\) 是年龄别死亡率，\(\beta(a)\) 是年龄别出生率。

\textbf{定理 237.3}（Euler-Lotka 方程）：稳定年龄分布的增长率 \(r\) 满足
\(\int_0^\infty e^{-ra} \beta(a) \ell(a) da = 1\)，其中
\(\ell(a) = \exp(-\int_0^a \mu(s) ds)\) 是存活概率。

\bigskip\hrule\bigskip

\subsection{第238章：传染病动力学}\label{ux7b2c238ux7ae0ux4f20ux67d3ux75c5ux52a8ux529bux5b66}

流行病学模型是数学生物学最成功的应用之一，对公共卫生决策（如疫苗接种策略）有直接影响。

\subsubsection{238.1 仓室模型}\label{ux4ed3ux5ba4ux6a21ux578b}

\textbf{定义 238.1}（SIR 模型 / Kermack-McKendrick
1927）：种群分为易感者 \(S\)、感染者 \(I\) 和康复者 \(R\) 三个仓室：

\[\begin{cases} \dot{S} = -\beta S I \\ \dot{I} = \beta S I - \gamma I \\ \dot{R} = \gamma I \end{cases}\]

其中 \(\beta\) 是传播率，\(\gamma\) 是康复率。

\textbf{定义 238.2}（基本再生数
\(\mathcal{R}_0\)）：\(\mathcal{R}_0 = \frac{\beta S(0)}{\gamma}\)
是一个感染者在完全易感种群中产生的新感染者的平均数目。\(\mathcal{R}_0 > 1\)
时疫情爆发，\(\mathcal{R}_0 < 1\) 时疫情消退。

\textbf{定理 238.1}（SIR 模型的阈值定理 / Kermack-McKendrick
阈值定理）： - 若 \(\mathcal{R}_0 \leq 1\)，\(I(t)\) 单调递减到
\(0\)（无疫情） - 若 \(\mathcal{R}_0 > 1\)，\(I(t)\)
先增后减（疫情爆发），最终感染者总数 \(R(\infty)\) 满足
\(S(\infty) = S(0) e^{-\mathcal{R}_0 (1 - S(\infty)/S(0))}\)

\emph{证明概要}：由 \(\dot{I} = (\beta S - \gamma)I\) 知
\(\dot{I}/I = \beta S - \gamma\)。当 \(\beta S_0 \leq \gamma\)（即
\(\mathcal{R}_0 \leq 1\)）时 \(\dot{I} \leq 0\)。若
\(\mathcal{R}_0 > 1\)，在 \(S(t)\) 降至 \(\gamma/\beta\) 之前 \(I(t)\)
递增，之后递减。由 \(dI/dS = -1 + \gamma/(\beta S)\) 积分得
\(I + S - (\gamma/\beta)\ln S = \text{const}\)，令 \(t \to \infty\) 且
\(I(\infty)=0\) 即得终态方程。∎

\textbf{定义 238.3}（SEIR 模型与扩展模型）： -
\textbf{SEIR}：增加\textbf{暴露者} \(E\)（已感染但尚未具传染性的潜伏期）
- \textbf{SIRS}：免疫是暂时的，康复者重新变为易感者 -
\textbf{SIS}：无免疫的疾病（康复后直接回到易感者仓室）

\textbf{定理 238.2}（接种阈值 / Herd Immunity Threshold）：若接种比例
\(p\) 大于 \(1 - 1/\mathcal{R}_0\)，则
\(\mathcal{R}_{\text{eff}} = \mathcal{R}_0(1 - p) < 1\)，疾病无法传播。这是群体免疫的数学基础。

\subsubsection{238.2
网络传染病模型}\label{ux7f51ux7edcux4f20ux67d3ux75c5ux6a21ux578b}

\textbf{定义 238.4}（网络上的 SIR 模型）：在接触网络 \(G = (V, E)\)
上，每个节点是 S/I/R 状态之一。传播沿边发生，速率
\(\beta\)。网络上基本再生数 \(\mathcal{R}_0\) 依赖于网络的度分布
\(P(k)\)。

\textbf{定理 238.3}（Pastor-Satorras-Vespignani
结果，2001）：在无标度网络（度分布
\(P(k) \sim k^{-\gamma}\)，\(2 < \gamma \leq 3\)）上，\(\mathcal{R}_0 \to \infty\)（当
\(N \to \infty\)），即不存在传播阈值------即使极低传播率也能维持疫情。

\subsubsection{238.3
空间传染病模型与最优控制}\label{ux7a7aux95f4ux4f20ux67d3ux75c5ux6a21ux578bux4e0eux6700ux4f18ux63a7ux5236}

\textbf{定义 238.5}（反应-扩散流行病模型）：

\[\begin{cases} \frac{\partial S}{\partial t} = d_S \nabla^2 S - \beta S I \\ \frac{\partial I}{\partial t} = d_I \nabla^2 I + \beta S I - \gamma I \end{cases}\]

\textbf{定理 238.4}（传播波速 / Fisher-KPP 型的波速）：空间 SIR
模型中传染病以速度 \(c^* = 2\sqrt{d_I (\beta S_0 - \gamma)}\) 传播（当
\(\beta S_0 > \gamma\)）。这是 Fisher-KPP 方程的行波解。

\emph{证明概要}（Fisher-KPP方程
\(\partial_t u = D\partial_{xx}u + ru(1-u/K)\)）：取行波变量
\(\xi = x - ct\)，代入得
\(D u'' + c u' + ru(1-u/K) = 0\)。化为等价一阶系统（相平面）：\(u'=v\),
\(Dv' = -cv - ru(1-u/K)\)。平衡点为 \((0,0)\)（不稳定结点）和
\((K,0)\)（鞍点）。行波解对应于连接两平衡点的异宿轨。波速
\(c \geq c_{\min}=2\sqrt{Dr}\) 时鞍-鞍连结存在，\(c < c_{\min}\)
时轨道进入负值区域不合法。∎

\textbf{定义 238.6}（传染病的最优控制）：最小化总代价（感染 +
干预）：\(J = \int_0^T [c_I I(t) + c_u u^2(t)] dt\)，受控 SIR 系统为
\(\dot{S} = -\beta S I (1 - u)\)，控制 \(u(t) \in [0, u_{\max}]\)
表示非药物干预程度。由 Pontryagin 极小值原理可求解最优干预策略。

\bigskip\hrule\bigskip

\subsection{第239章：反应-扩散系统与形态发生}\label{ux7b2c239ux7ae0ux53cdux5e94-ux6269ux6563ux7cfbux7edfux4e0eux5f62ux6001ux53d1ux751f}

Turing 在 1952 年提出了反应-扩散系统产生生物斑图的机制（形态发生 /
morphogenesis），开启了生物斑图形成的数学理论。

\subsubsection{239.1 Turing
不稳定性}\label{turing-ux4e0dux7a33ux5b9aux6027}

\textbf{定义 239.1}（反应-扩散系统）：二组分反应-扩散系统为

\[\begin{cases} \frac{\partial u}{\partial t} = D_u \nabla^2 u + f(u, v) \\ \frac{\partial v}{\partial t} = D_v \nabla^2 v + g(u, v) \end{cases}\]

其中 \(f\) 和 \(g\) 是反应项（自催化、抑制等），\(D_u, D_v > 0\)
是扩散系数。

\textbf{定义 239.2}（Turing 不稳定性，1952）：在齐次平衡点
\((u^*, v^*)\)
处，若无扩散系统是稳定的（正平衡点），但有扩散系统在某些波数下不稳定，则出现
\textbf{Turing 不稳定性}。条件是： -
\(\operatorname{Tr}(J) = f_u + g_v < 0\)（稳定性） -
\(\det(J) = f_u g_v - f_v g_u > 0\)（稳定性） -
\(D_v f_u + D_u g_v > 2 \sqrt{D_u D_v \det(J)}\)（扩散驱动的不稳定性）

\textbf{定理 239.1}（Turing 斑图的条件）：Turing
不稳定性要求「快扩散的抑制剂 + 慢扩散的激活剂」（\(D_v \gg D_u\)），且
Jacobian 矩阵满足
\(f_u > 0\)（自激活），\(g_v < 0\)（自抑制），\(f_v g_u < 0\)（交叉抑制）。

\textbf{例 239.1}（Gierer-Meinhardt 模型 / 激活-抑制模型，1972）：

\[\begin{cases} \frac{\partial a}{\partial t} = D_a \nabla^2 a + \rho_a \frac{a^2}{h} - \mu_a a \\ \frac{\partial h}{\partial t} = D_h \nabla^2 h + \rho_h a^2 - \mu_h h \end{cases}\]

在 \(D_h \gg D_a\) 时产生斑点、条纹等 Turing
斑图。这解释了动物皮毛花纹、贝壳纹理等现象。

\subsubsection{239.2 趋化性模型}\label{ux8d8bux5316ux6027ux6a21ux578b}

\textbf{定义 239.3}（Keller-Segel
趋化性模型，1970）：细胞（如黏菌）沿化学信号浓度梯度移动：

\[\begin{cases} \frac{\partial u}{\partial t} = D_u \nabla^2 u - \chi \nabla \cdot (u \nabla v) \\ \frac{\partial v}{\partial t} = D_v \nabla^2 v + \alpha u - \beta v \end{cases}\]

其中 \(u\) 是细胞密度，\(v\) 是化学吸引剂浓度，\(\chi\) 是趋化敏感度。

\textbf{定理 239.2}（Keller-Segel 模型的临界质量 / 爆破条件）：对二维
Keller-Segel 模型，存在临界质量
\(M_c = 8\pi D_u / \chi\)。若初始细胞质量 \(M_0 < M_c\)，解全局存在；若
\(M_0 > M_c\)，有限时间爆破（细胞聚集）。

\subsubsection{239.3
兴奋介质与行波}\label{ux5174ux594bux4ecbux8d28ux4e0eux884cux6ce2}

\textbf{定义 239.4}（FitzHugh-Nagumo
方程）：神经和心脏组织兴奋传播的简化模型：

\[\begin{cases} \frac{\partial u}{\partial t} = D \nabla^2 u + u(1 - u)(u - a) - v \\ \frac{\partial v}{\partial t} = \varepsilon (b u - v) \end{cases}\]

其中
\(0 < a < 1\)，\(\varepsilon \ll 1\)。系统支持行波脉冲、螺旋波和目标波。

\textbf{定理 239.3}（行波的存在性 / 快-慢分析）：在
\(\varepsilon \to 0\)
极限下，行波解由快子系统（兴奋前沿）和慢子系统（不应期）交替组成，可用奇异摄动理论分析。

\bigskip\hrule\bigskip

\subsection{第240章：系统生物学与生化网络}\label{ux7b2c240ux7ae0ux7cfbux7edfux751fux7269ux5b66ux4e0eux751fux5316ux7f51ux7edc}

系统生物学使用网络理论和动力系统方法研究细胞内分子相互作用网络的整体行为。

\subsubsection{240.1 酶动力学与 Michaelis-Menten
方程}\label{ux9176ux52a8ux529bux5b66ux4e0e-michaelis-menten-ux65b9ux7a0b}

\textbf{定义 240.1}（Michaelis-Menten 酶动力学，1913）：酶 \(E\)
催化底物 \(S\) 转化为产物 \(P\)：

\[E + S \underset{k_{-1}}{\overset{k_1}{\rightleftharpoons}} ES \overset{k_{\text{cat}}}{\longrightarrow} E + P\]

\textbf{定理 240.1}（Michaelis-Menten 方程 /
准稳态近似）：在底物远多于酶的准稳态假设下，反应速率为

\[v = \frac{V_{\max} [S]}{K_m + [S]}\]

其中
\(V_{\max} = k_{\text{cat}} [E]_T\)（最大速率），\(K_m = (k_{-1} + k_{\text{cat}}) / k_1\)（Michaelis
常数）。

\emph{证明概要}（QSSA推导）：酶-底物复合物满足
\(\frac{d[ES]}{dt} = k_1[E][S] - (k_{-1}+k_{\text{cat}})[ES]\)。准稳态假设
\(\frac{d[ES]}{dt} \approx 0\)，结合酶守恒 \([E]_T = [E] + [ES]\)，解得
\([ES] = \frac{[E]_T[S]}{K_m + [S]}\)。反应速率
\(v = k_{\text{cat}}[ES] = \frac{k_{\text{cat}}[E]_T[S]}{K_m + [S]}\)。此近似的有效性条件为
\([S] \gg [E]_T\)（初始快速过渡后）。∎

\subsubsection{240.2
基因调控网络}\label{ux57faux56e0ux8c03ux63a7ux7f51ux7edc}

\textbf{定义 240.2}（基因调控的 ODE 模型 / Goodwin 振荡器）：描述 mRNA
和蛋白质浓度变化的经典模型：

\[\begin{cases} \dot{m} = \frac{\alpha}{1 + p^n / K^n} - \gamma_m m \\ \dot{p} = \beta m - \gamma_p p \end{cases}\]

其中 \(m\) 是 mRNA 浓度，\(p\) 是蛋白质浓度，Hill 系数 \(n\)
表示合作的强度。

\textbf{定理 240.2}（Goodwin 振荡器的 Hopf
分岔）：在标准参数归一化下，当 \(n > 8\) 时（一般参数下阈值由非线性方程
\(\det(J(\lambda)) = 0\) 的 Hopf 分岔条件给出，临界 \(n\)
依赖于降解速率比 \(\gamma_m / \gamma_p\) 和合成参数
\(\alpha, \beta\)），系统通过 Hopf
分岔产生持续振荡（生物钟的数学基础）。

\textbf{定义 240.3}（布尔网络模型 / Kauffman 1969）：基因表达状态简化为
ON(1)/OFF(0)，每个基因的状态是由少数调控基因决定的布尔函数。\textbf{NK
模型}：\(N\) 个基因，每个由 \(K\) 个随机选择的输入调控。

\textbf{定理 240.3}（Kauffman 的混沌边缘）：\(K = 2\)
时系统处于「混沌边缘」（有序与混沌的边界），与真实细胞网络的动力学特征相似。

\subsubsection{240.3
代谢网络与通量平衡分析}\label{ux4ee3ux8c22ux7f51ux7edcux4e0eux901aux91cfux5e73ux8861ux5206ux6790}

\textbf{定义 240.4}（通量平衡分析 / FBA）：对代谢网络，设化学计量矩阵
\(S \in \mathbb{R}^{m \times n}\)（\(m\) 种代谢物，\(n\)
个反应）。稳态条件为 \(S \cdot v = 0\)（\(v\) 是通量向量）。FBA
是线性规划：

\[\max \; c^T v \quad \text{s.t.} \quad S v = 0, \; v_{\min} \leq v \leq v_{\max}\]

\textbf{定理
240.4}（极端通路与基本通量模式）：代谢网络的所有可行稳态通量的锥由有限条极端通路或基本通量模式生成。这些模式对应生物的生理上可实现的代谢路径。

\subsubsection{240.4
随机生化动力学}\label{ux968fux673aux751fux5316ux52a8ux529bux5b66}

\textbf{定义 240.5}（化学主方程 / Chemical Master
Equation）：分子计数低时，需用随机描述。设 \(P(\mathbf{x}, t)\) 是时间
\(t\) 各物种分子数为 \(\mathbf{x}\) 的概率。\textbf{化学主方程}为

\[\frac{dP(\mathbf{x}, t)}{dt} = \sum_j \left[ a_j(\mathbf{x} - \nu_j) P(\mathbf{x} - \nu_j, t) - a_j(\mathbf{x}) P(\mathbf{x}, t) \right]\]

其中 \(a_j(\mathbf{x})\) 是反应 \(j\) 的倾向函数，\(\nu_j\)
是化学计量变化向量。

\textbf{定理 240.5}（Gillespie 算法 / 随机模拟算法，1977）：从当前状态
\(\mathbf{x}\) 出发：(1) 生成反应发生的时间
\(\tau \sim \operatorname{Exp}(a_0(\mathbf{x}))\)（\(a_0 = \sum_j a_j\)）；(2)
选择反应 \(j\) 的概率为 \(a_j / a_0\)。Gillespie
算法精确模拟了化学主方程的轨线。

\bigskip\hrule\bigskip

\subsection{第241章：神经科学与神经动力学}\label{ux7b2c241ux7ae0ux795eux7ecfux79d1ux5b66ux4e0eux795eux7ecfux52a8ux529bux5b66}

神经科学的数学模型从 Hodgkin-Huxley
的神经元模型到神经网络理论，连接了个体神经元和认知功能。

\subsubsection{241.1 神经元模型}\label{ux795eux7ecfux5143ux6a21ux578b}

\textbf{定义 241.1}（Hodgkin-Huxley 模型，1952）：乌贼巨轴突的膜电位
\(V\) 满足

\[C_m \frac{dV}{dt} = I_{\text{ext}} - \bar{g}_{\text{Na}} m^3 h (V - V_{\text{Na}}) - \bar{g}_{\text{K}} n^4 (V - V_{\text{K}}) - g_L (V - V_L)\]

门控变量 \(m, h, n \in [0, 1]\)
遵循一阶动力学：\(\frac{dx}{dt} = \alpha_x(V)(1-x) - \beta_x(V) x\)。

\textbf{定理 241.1}（Hodgkin-Huxley 的动作电位与兴奋性）：模型中 Na+
通道的快激活和慢失活，与 K+
通道的慢激活协同产生动作电位。通过分岔分析，静息态通过鞍结分岔或 Hopf
分岔失去稳定性。

\textbf{定义 241.2}（简化神经元模型）： - \textbf{FitzHugh-Nagumo
模型}（Hodgkin-Huxley
的二维简化）：\(\dot{v} = v - v^3/3 - w + I\)，\(\dot{w} = \varepsilon (v + a - b w)\)
- \textbf{漏积分发放模型}（Leaky Integrate-and-Fire /
LIF）：\(\tau_m \dot{V} = -(V - V_{\text{rest}}) + R_m I(t)\)，当
\(V \geq V_{\text{th}}\)，发放脉冲并重置 \(V \to V_{\text{reset}}\) -
\textbf{Izhikevich 模型}（2003）：兼顾生物真实性和计算效率的二维模型

\subsubsection{241.2
神经编码与信息论}\label{ux795eux7ecfux7f16ux7801ux4e0eux4fe1ux606fux8bba}

\textbf{定义
241.3}（发放率编码与时间编码）：神经元通过脉冲序列传递信息。\textbf{发放率编码}假设信息由单位时间内的脉冲数表示；\textbf{时间编码}假设精确的脉冲时刻携带信息。

\textbf{定义 241.4}（Spike-Triggered Average / 反向相关法）：给定刺激
\(s(t)\) 和神经响应（脉冲序列 \(\rho(t)\)），\textbf{发放触发平均}为

\[\operatorname{STA}(\tau) = \frac{1}{N} \sum_{i} s(t_i - \tau)\]

其中 \(t_i\) 是脉冲时刻。STA 给出了线性滤波近似中的感受野（receptive
field）。

\textbf{定理 241.2}（互信息与神经编码效率）：刺激 \(S\) 与响应 \(R\)
的互信息 \(I(S; R)\) 衡量神经编码的效率。有效编码假说（Barlow
1961）认为感觉系统最大化互信息并最小化冗余。

\subsubsection{241.3
神经网络动力学}\label{ux795eux7ecfux7f51ux7edcux52a8ux529bux5b66}

\textbf{定义 241.5}（Hopfield 网络，1982）：\(N\)
个二值神经元（\(\pm 1\)），状态更新为
\(x_i = \operatorname{sgn}(\sum_j W_{ij} x_j - \theta_i)\)。若 \(W\)
对称（\(W_{ij} = W_{ji}\)），系统的能量函数
\(E = -\frac{1}{2} \sum_{i,j} W_{ij} x_i x_j + \sum_i \theta_i x_i\)
单调递减，收敛到局部极小值（记忆状态）。

\textbf{定理 241.3}（Hopfield 网络的记忆容量）：Hebb 学习
\(W_{ij} = \frac{1}{N} \sum_{\mu=1}^p \xi_i^\mu \xi_j^\mu\) 存储 \(p\)
个模式。\(p \leq 0.138 N\)
时存储的模式的回忆误差极小（Amit-Gutfreund-Sompolinsky 1985）。

\textbf{定义 241.6}（Wilson-Cowan 平均场方程，1972）：

\[\tau_e \frac{dE}{dt} = -E + f_e(w_{ee} E - w_{ie} I + I_e)\]
\[\tau_i \frac{dI}{dt} = -I + f_i(w_{ei} E - w_{ii} I + I_i)\]

其中 \(E\) 和 \(I\) 分别是兴奋和抑制性神经元的平均发放率，\(f_{e,i}\) 是
Sigmoid 激活函数。

\textbf{定理 241.4}（Wilson-Cowan
模型的振荡与节律）：兴奋-抑制回路通过适当参数可产生多种振荡行为（\(\alpha\)、\(\beta\)、\(\gamma\)
节律），通过 Hopf 分岔产生。

\bigskip\hrule\bigskip

\emph{本卷在微分方程（V18）、动力系统（V28）和概率论（V21）的基础上，系统建立了种群动力学（Logistic
方程、Lotka-Volterra 模型、年龄结构）、流行病学（SIR 模型与
\(\mathcal{R}_0\)、网络传播）、反应-扩散系统（Turing
不稳定性、趋化性）、系统生物学（酶动力学、基因调控、生化网络）和神经科学（神经元模型、神经编码、神经网络）的数学理论。共
5 章。}

\newpage

\section{卷四十三：辛几何与切触几何}\label{ux5377ux56dbux5341ux4e09ux8f9bux51e0ux4f55ux4e0eux5207ux89e6ux51e0ux4f55}

\begin{quote}
辛几何是经典力学（Hamilton
力学）的天然数学语言，研究偶数维光滑流形上的非退化闭
2-形式（辛形式）。切触几何是奇数维流形上的对应理论。辛几何在 1980
年代后经历了爆发式发展：Gromov
的伪全纯曲线理论（1985）引入了「刚性」概念，Gromov-Witten
不变量连接了辛几何和代数几何（镜对称 / mirror symmetry），Floer
同调为辛几何提供了无穷维 Morse
理论工具。本卷在微分几何（V12）和代数拓扑（V15）的基础上，建立辛流形、辛约化、Gromov-Witten
理论、切触几何和 Fukaya 范畴的系统理论。
\end{quote}

\bigskip\hrule\bigskip

\subsection{第242章：辛流形基础}\label{ux7b2c242ux7ae0ux8f9bux6d41ux5f62ux57faux7840}

辛几何研究带有辛形式（封闭、非退化的 2-形式）的光滑流形，是 Hamilton
力学的相空间的抽象化。

\subsubsection{242.1 辛线性代数}\label{ux8f9bux7ebfux6027ux4ee3ux6570}

\textbf{定义 242.1}（辛向量空间）：有限维实向量空间 \(V\)
上的\textbf{辛形式} \(\omega: V \times V \to \mathbb{R}\)
是双线性、反对称（\(\omega(u, v) = -\omega(v, u)\)）且非退化（\(\omega(u, v) = 0\)
对所有 \(v\) 蕴含 \(u = 0\)）的。\((V, \omega)\)
称为\textbf{辛向量空间}。

\textbf{命题 242.1}（辛向量空间的标准型 /
维数）：任何辛向量空间的维数是偶数：\(\dim V = 2n\)。存在辛基
\((e_1, \ldots, e_n, f_1, \ldots, f_n)\) 使得
\(\omega(e_i, f_j) = \delta_{ij}\)，\(\omega(e_i, e_j) = \omega(f_i, f_j) = 0\)。辛形式在辛基下的矩阵为
\(\begin{pmatrix} 0 & I \\ -I & 0 \end{pmatrix}\)。

\emph{证明概要}：设 \(\dim V = m\)，任选基下 \(\omega\) 的矩阵 \(A\)
为反对称矩阵（\(A^T = -A\)）。非退化性
\(\iff \det A \neq 0\)。对任意反对称矩阵
\(A\)，\(\det A = \det(A^T) = \det(-A) = (-1)^m \det A\)。当 \(m\)
为奇数时，\(\det A = -\det A \Rightarrow \det A = 0\)，与非退化矛盾。故
\(m\) 必为偶数 \(2n\)。辛基的存在性由归纳构造：取任意
\(e_1 \neq 0\)，由非退化性存在 \(f_1\) 使
\(\omega(e_1, f_1) = 1\)，在辛正交补 \(\{e_1, f_1\}^\omega\)
上继续归纳即得。∎

\textbf{定义 242.2}（辛群 \(\operatorname{Sp}(2n)\)）：\textbf{辛群}
\(\operatorname{Sp}(2n, \mathbb{R}) = \{A \in \operatorname{GL}(2n, \mathbb{R}) : A^T \Omega A = \Omega\}\)，其中
\(\Omega = \begin{pmatrix} 0 & I \\ -I & 0 \end{pmatrix}\)。

\subsubsection{242.2 辛流形}\label{ux8f9bux6d41ux5f62}

\textbf{定义 242.3}（辛流形）：\textbf{辛流形} \((M, \omega)\)
是光滑流形 \(M\) 配备闭（\(d\omega = 0\)）且非退化（\(\omega^n \neq 0\)
处处）的 2-形式 \(\omega \in \Omega^2(M)\)。\(\dim M = 2n\)。

\textbf{例 242.1}（辛流形的例子）： - \(\mathbb{R}^{2n}\)
的标准辛形式：\(\omega_0 = \sum_i dq_i \wedge dp_i\) - 余切丛 \(T^*Q\)
配备典范辛形式 \(\omega = -d\lambda\)，其中 Liouville 形式
\(\lambda = \sum_i p_i dq_i\) - Kähler 流形 \((M, J, g)\) 的 Kähler 形式
\(\omega(\cdot, \cdot) = g(J\cdot, \cdot)\) 是辛形式 -
余伴随轨道（Kirillov-Kostant-Souriau）：李群 \(G\)
的余伴随轨道具有自然辛结构

\textbf{定理 242.1}（Darboux 定理）：任何辛流形 \((M^{2n}, \omega)\)
局部同构于
\((\mathbb{R}^{2n}, \omega_0)\)。即辛几何无局部不变量------所有 2n
维辛流形局部看起来都一样。

\emph{证明思路}：用 Moser 道路法（Moser's path method）。设
\(p \in M\)，取 \(p\) 附近的局部坐标使
\(\omega|_p = \omega_0|_p\)（先通过线性代数找到辛基，再用指数映射拉平到一阶）。构造辛形式的线性同伦
\(\omega_t = (1-t)\omega_0 + t\omega\)，在 \(p\) 的某邻域上 \(\omega_t\)
均为闭的且 \(\omega_t|_p = \omega_0|_p\)。由 Poincare
引理，\(\omega - \omega_0 = d\alpha\)（\(\alpha\) 为 1-形式，可选取使得
\(\alpha|_p = 0\)）。目标是寻找一族微分同胚
\(\phi_t\)（\(\phi_0 = \operatorname{id}\)）满足
\(\phi_t^* \omega_t = \omega_0\)。对 \(t\) 求导得
\(\phi_t^*(\mathcal{L}_{X_t} \omega_t + \dot{\omega}_t) = 0\)，其中
\(X_t\) 是 \(\phi_t\) 的生成向量场。需解
\(\mathcal{L}_{X_t} \omega_t + \dot{\omega}_t = 0\)。利用 Cartan 公式
\(\mathcal{L}_{X_t} \omega_t = d(\iota_{X_t} \omega_t)\) 和
\(\dot{\omega}_t = d\alpha\)，方程化为
\(d(\iota_{X_t} \omega_t + \alpha) = 0\)。由于 \(\omega_t\) 非退化（在
\(p\) 的充分小邻域上），可从 \(\iota_{X_t} \omega_t = -\alpha\) 唯一解出
\(X_t\)。因 \(\alpha|_p = 0\)，\(X_t|_p = 0\)，故 \(X_t\) 的流
\(\phi_t\) 在小邻域上定义至 \(t=1\)。于是
\(\phi_1^* \omega = \phi_1^* \omega_1 = \omega_0\)，即 Darboux
坐标存在。∎

\subsubsection{242.3 Hamilton
力学与辛几何}\label{hamilton-ux529bux5b66ux4e0eux8f9bux51e0ux4f55}

\textbf{定义 242.4}（Hamilton 向量场）：对函数 \(H \in C^\infty(M)\)，其
\textbf{Hamilton 向量场} \(X_H\) 由 \(\iota_{X_H} \omega = dH\)
唯一确定。

\textbf{定义 242.5}（Poisson 括号）：\(C^\infty(M)\) 上的
\textbf{Poisson 括号}为
\(\{F, G\} = \omega(X_F, X_G) = dF(X_G)\)。Poisson 括号满足 Jacobi
恒等式，使 \((C^\infty(M), \{\cdot, \cdot\})\) 成为 Poisson 代数。

\textbf{推论 242.2}（Liouville 定理 / 相空间体积守恒）：Hamilton 流
\(\phi_H^t\) 保持辛体积形式
\(\omega^n/n!\)：\((\phi_H^t)^* \omega^n = \omega^n\)。特别是，Hamilton
相流不可压缩。

\textbf{定理 242.3}（Noether 定理的辛表述）：若 \(G\) 通过辛同胚作用在
\((M, \omega)\) 上，且作用有矩映射
\(\mu: M \to \mathfrak{g}^*\)（\(\langle \mu, \xi \rangle\) 的 Hamilton
向量场生成 \(G\) 作用），则 \(\mu\) 沿 Hamilton
流守恒：\(\mu \circ \phi_H^t = \mu\)（若 \(\{H, \mu\} = 0\)）。

\bigskip\hrule\bigskip

\subsection{第243章：辛约化与矩映射}\label{ux7b2c243ux7ae0ux8f9bux7ea6ux5316ux4e0eux77e9ux6620ux5c04}

辛约化（Marsden-Weinstein
约化）是构造新辛流形的基本方法，是经典力学中消除对称性（如守恒量）的数学抽象。

\subsubsection{243.1 矩映射与 Hamilton
群作用}\label{ux77e9ux6620ux5c04ux4e0e-hamilton-ux7fa4ux4f5cux7528}

\textbf{定义 243.1}（Hamilton 群作用 / 辛作用）：李群 \(G\) 在辛流形
\((M, \omega)\) 上的辛作用称为 \textbf{Hamilton 作用}，如果每个
\(\xi \in \mathfrak{g}\) 对应的无穷小生成元 \(X_\xi\) 是某个函数
\(H_\xi\) 的 Hamilton 向量场。

\textbf{定义 243.2}（矩映射 / Moment Map）：\(G\)-等变映射
\(\mu: M \to \mathfrak{g}^*\) 称为
\textbf{矩映射}（或动量映射），如果对所有
\(\xi \in \mathfrak{g}\)，\(\langle \mu, \xi \rangle\) 是 \(X_\xi\) 的
Hamilton 函数：

\[d\langle \mu, \xi \rangle = \iota_{X_\xi} \omega\]

\textbf{例 243.1}（余切丛的矩映射）：\(G\) 左作用在 \(Q\) 上，提升到
\(T^*Q\)。矩映射为
\(\langle \mu(q, p), \xi \rangle = p(X_\xi(q))\)（线性动量）。

\subsubsection{243.2 Marsden-Weinstein
约化}\label{marsden-weinstein-ux7ea6ux5316}

\textbf{定理 243.1}（Marsden-Weinstein-Meyer 辛约化定理，1974）：设
\(G\) 以 Hamilton 作用在 \((M, \omega)\)
上，\(\mu: M \to \mathfrak{g}^*\) 是矩映射。若 \(a \in \mathfrak{g}^*\)
是 \(\mu\) 的正则值且 \(G_a\)（\(a\) 的迷向子群）自由作用于
\(\mu^{-1}(a)\)，则商空间

\[M /\!/ G(a) = \mu^{-1}(a) / G_a\]

具有唯一的辛结构 \(\omega_{\text{red}}\)，满足
\(\pi^* \omega_{\text{red}} = \iota^* \omega\)（\(\iota: \mu^{-1}(a) \hookrightarrow M\)，\(\pi: \mu^{-1}(a) \to M /\!/ G(a)\)
是投影）。

\textbf{定理
243.2}（约化的维数）：\(\dim(M /\!/ G(a)) = \dim M - 2 \dim G_a\)。特别地，若
\(G\) 以有效自由作用在 \(\mu^{-1}(a)\) 上，则
\(\dim(M /\!/ G(a)) = \dim M - 2 \dim G\)。

\textbf{例 243.2}（辛约化的例子）： - Kepler 问题通过
\(\operatorname{SO}(3)\) 约化消除角动量对称性（二维有效问题） -
复射影空间 \(\mathbb{C}P^n = S^{2n+1} /\!/ S^1\)（使用 Hopf
纤维的辛约化） - 模空间的矩映射构造：平坦联络的模空间
\(M_{\text{flat}}\) 是辛约化的商

\subsubsection{243.3 Delzant
定理与环面簇}\label{delzant-ux5b9aux7406ux4e0eux73afux9762ux7c07}

\textbf{定义 243.3}（环面辛流形 / Toric Symplectic
Manifold）：紧致连通辛流形 \((M^{2n}, \omega)\)
是\textbf{环面的}，如果它容许 \(\mathbb{T}^n\) 的有效 Hamilton
作用（最大可能的环面维数）。

\textbf{定理 243.3}（Delzant 定理，1988）：紧致环面辛流形与 Delzant
多面体（有理多面体，满足在每个顶点处边方向形成
\(\mathbb{Z}\)-基）之间存在一一对应。矩映射的像 \(\mu(M)\) 恰好是
Delzant 多面体。

\bigskip\hrule\bigskip

\subsection{第244章：伪全纯曲线与 Gromov-Witten
理论}\label{ux7b2c244ux7ae0ux4f2aux5168ux7eafux66f2ux7ebfux4e0e-gromov-witten-ux7406ux8bba}

Gromov 在 1985
年引入了伪全纯曲线方法，革命性地改变了辛几何。这是辛几何「刚性」的源泉，也是
Gromov-Witten 不变量和镜对称的基础。

\subsubsection{244.1
近复结构与伪全纯曲线}\label{ux8fd1ux590dux7ed3ux6784ux4e0eux4f2aux5168ux7eafux66f2ux7ebf}

\textbf{定义 244.1}（与辛形式相容的近复结构）：流形 \(M\)
上的\textbf{近复结构} \(J\) 满足 \(J^2 = -\operatorname{id}\)。\(J\)
与辛形式 \(\omega\) \textbf{相容}（compatible），如果
\(g_J(v, w) = \omega(v, Jw)\) 是 Riemann 度量。\footnote{注：\textbf{tame}
  是更弱的条件------仅要求 \(\omega(v, Jv) > 0\) 对所有 \(v \neq 0\)
  成立，而不要求 \(g_J\) 是对称的。compatible 等价于 tame + \(J\) 是
  \(\omega\)-不变的（即 \(\omega(Jv, Jw) = \omega(v, w)\)）。}

\textbf{定理 244.1}（相容近复结构的存在性）：辛流形 \((M, \omega)\)
上相容近复结构的空间 \(\mathcal{J}(M, \omega)\) 是非空且可缩的。

\textbf{定义 244.2}（\(J\)-全纯曲线）：映射
\(u: (\Sigma, j) \to (M, J)\)（\(\Sigma\) 是 Riemann 面）是
\textbf{\(J\)-全纯的}，如果

\[\bar{\partial}_J u = \frac{1}{2}(du + J \circ du \circ j) = 0\]

即 \((u, J)\)-全纯曲线的 Cauchy-Riemann 方程。

\textbf{定理 244.2}（能量恒等式）：\(J\)-全纯曲线 \(u\) 的能量
\(E(u) = \frac{1}{2} \int_\Sigma |du|^2 d\operatorname{vol}_\Sigma = \int_\Sigma u^* \omega\)。因此
\(J\)-全纯曲线的面积等于辛面积（拓扑量）。

\subsubsection{244.2 Gromov
紧致性与非压缩定理}\label{gromov-ux7d27ux81f4ux6027ux4e0eux975eux538bux7f29ux5b9aux7406}

\textbf{定理 244.3}（Gromov 紧致性定理，1985）：具有一致有界能量的
\(J\)-全纯曲线的序列可以「气泡化」收敛到节点 \(J\)-全纯曲线（nodal
\(J\)-holomorphic curve），除有限个「气泡」点外在 \(\Sigma\)
上光滑收敛。

\textbf{定理 244.4}（Gromov 非压缩定理 / Non-Squeezing
Theorem）：若存在辛嵌入
\(B^{2n}(R) \hookrightarrow B^2(r) \times \mathbb{R}^{2n-2}\)（柱面），则
\(R \leq r\)。即一个大的辛球不能被「挤入」一个细的辛柱面中------这是辛刚性的里程碑。

\emph{证明思路}：使用伪全纯曲线。在 \(\mathbb{C}P^1\)
的辛嵌入中，\(J\)-全纯球的不变量限制了 \(R\) 和 \(r\) 的关系。能量
\(\geq \pi R^2\) 但面积受限于 \(\pi r^2\)。∎

\subsubsection{244.3 Gromov-Witten
不变量}\label{gromov-witten-ux4e0dux53d8ux91cf}

\textbf{定义 244.3}（Gromov-Witten 不变量）：辛流形 \((M, \omega)\) 的
\textbf{Gromov-Witten 不变量}是映射

\[\operatorname{GW}_{g, k, A}: H^*(M)^{\otimes k} \to \mathbb{Q}\]

计数（在适当意义下）亏格 \(g\)、\(k\) 个标记点且同调类为
\(A \in H_2(M; \mathbb{Z})\) 的 \(J\)-全纯曲线的数目。形式上，

\[\operatorname{GW}_{g, k, A}(\alpha_1, \ldots, \alpha_k) = \int_{[\overline{\mathcal{M}}_{g,k}(M, A)]^{\text{vir}}} \prod_{i=1}^k \operatorname{ev}_i^* \alpha_i\]

其中 \(\overline{\mathcal{M}}_{g,k}(M, A)\)
是稳定映射的模空间（紧化），\([\cdot]^{\text{vir}}\) 是虚基本类。

\textbf{定理 244.5}（GW 不变量与计数几何）：对合适的 \(\alpha_i\)（如
Poincaré 对偶于点类的上同调类），GW 不变量计算特定几何约束下的
\(J\)-全纯曲线数。这些数字是辛不变量（不依赖于 \(J\) 的选择）。

\subsubsection{244.4 量子上同调与 WDVV
方程}\label{ux91cfux5b50ux4e0aux540cux8c03ux4e0e-wdvv-ux65b9ux7a0b}

\textbf{定义 244.4}（量子上同调环 / Quantum Cohomology）：\(M\)
的\textbf{量子上同调} \(QH^*(M)\) 是 \(H^*(M; \mathbb{C})\)
作为向量空间，但乘法量子变形：

\[\alpha * \beta = \sum_{A \in H_2(M)} (\alpha * \beta)_A \cdot q^A\]

其中 \((\alpha * \beta)_A\) 由三个点的 GW
不变量给出：\(\langle \alpha, \beta, \gamma \rangle_A = \operatorname{GW}_{0,3,A}(\alpha, \beta, \gamma)\)。

\textbf{定理 244.6}（WDVV 方程 / Witten-Dijkgraaf-Verlinde-Verlinde
方程）：量子上同调的关联函数满足 WDVV 方程，这是 GW
势的完全可积系统的条件。WDVV 方程等价于量子上同调的结合性。

\bigskip\hrule\bigskip

\subsection{第245章：切触几何}\label{ux7b2c245ux7ae0ux5207ux89e6ux51e0ux4f55}

切触几何是奇数维流形上的几何，与辛几何通过「辛化」（symplectization）紧密相连。切触几何是热力学、光学和三维拓扑的自然语言。

\subsubsection{245.1 切触结构}\label{ux5207ux89e6ux7ed3ux6784}

\textbf{定义 245.1}（切触流形）：奇数维流形 \(M^{2n+1}\)
上的\textbf{切触结构} \(\xi\) 是余维 1
的极大非可积超平面场。局部地，\(\xi = \ker \alpha\)，其中
\textbf{切触形式} \(\alpha \in \Omega^1(M)\) 满足
\(\alpha \wedge (d\alpha)^n \neq 0\) 处处。

\textbf{例 245.1}（切触流形的例子）： - \(\mathbb{R}^{2n+1}\)
的标准切触形式：\(\alpha_0 = dz - \sum_{i=1}^n y_i dx_i\) - Legendre
子流形
\(\Lambda \subset T^*Q \times \mathbb{R}\)：\(\alpha = dz - \lambda\)
的积分子流形 - 给定切触流形
\((M, \alpha)\)，其\textbf{辛化}（symplectization）是辛流形
\((M \times \mathbb{R}, d(e^t \alpha))\)。反之，若辛流形 \((M, \omega)\)
满足 \(\omega = d\alpha\) 为恰当形式（如余切丛 \(T^*Q\)），则
\((M, \alpha)\) 也可视为一个切触流形 -
\(S^{2n+1} \subset \mathbb{C}^{n+1}\) 的标准切触结构（由复切空间诱导）

\textbf{定义 245.2}（Legendre 子流形）：切触流形 \((M^{2n+1}, \xi)\)
的最大维积分子流形（维数为 \(n\)）称为 \textbf{Legendre
子流形}。它们是辛几何中 Lagrangian 子流形的切触类比。

\subsubsection{245.2 切触同调}\label{ux5207ux89e6ux540cux8c03}

\textbf{定义 245.3}（Reeb 向量场）：切触形式 \(\alpha\) 的 \textbf{Reeb
向量场} \(R_\alpha\) 由 \(\iota_{R_\alpha} d\alpha = 0\) 和
\(\alpha(R_\alpha) = 1\) 唯一确定。Reeb 向量场是 Hamilton
向量场的切触类比。

\textbf{定义 245.4}（切触同调 / Contact
Homology，Eliashberg-Givental-Hofer 2000）：切触流形 \((M, \xi)\)
的\textbf{切触同调} \(HC_*(M, \xi)\) 由 Reeb
轨道的链复形生成，边界算子计数伪全纯曲线。切触同调是切触结构的不变量。

\textbf{定理 245.1}（切触几何的刚性 / Eliashberg 1989）：在维数 \(> 3\)
的闭流形上，所有过度扭转（overtwisted）切触结构是同痕的，但紧切触（tight）结构形成更丰富的分类。即切触几何将流形的切触结构分为两类：柔性的（过度扭转）和刚性的（紧的）。

\subsubsection{245.3
三维切接触拓扑}\label{ux4e09ux7ef4ux5207ux63a5ux89e6ux62d3ux6251}

\textbf{定义 245.5}（三维切触拓扑的特殊性）：维数 3 中切触结构 \(\xi\)
的平面场可以用特征叶理（characteristic foliation）和开卷分解（open book
decomposition）研究。Giroux 定理建立了开卷分解与切触结构之间的对应。

\textbf{定理 245.2}（Giroux
定理，2000）：三维闭流形上的切触结构（模同痕）与开卷分解（模正稳定化）一一对应。这是理解和分类三维切触结构的基本工具。

\bigskip\hrule\bigskip

\subsection{第246章：Fukaya
范畴与镜对称}\label{ux7b2c246ux7ae0fukaya-ux8303ux7574ux4e0eux955cux5bf9ux79f0}

Fukaya 范畴是辛流形中 Lagrangian 子流形的 \(A_\infty\)-范畴，是
Kontsevich 的同调镜对称猜想中的核心对象。

\subsubsection{246.1 Lagrangian 子流形与 Floer
同调}\label{lagrangian-ux5b50ux6d41ux5f62ux4e0e-floer-ux540cux8c03}

\textbf{定义 246.1}（Lagrangian 子流形）：辛流形 \((M^{2n}, \omega)\) 的
\(n\) 维子流形 \(L\) 是 \textbf{Lagrangian} 的，如果
\(\omega|_L = 0\)。Lagrangian 子流形是辛几何中最基本的对象。

\textbf{定义 246.2}（Lagrangian Floer 同调）：对两个横截相交的
Lagrangian 子流形 \(L_0, L_1\)，\textbf{Lagrangian Floer 同调}
\(HF(L_0, L_1)\)
由交点生成，边界算子计数连接交点的伪全纯带（strip）。\(HF(L_0, L_1)\) 在
Hamiltonian 同痕下不变。

\textbf{定理 246.1}（Floer 的 Lagrangian 交猜想 / Arnold
猜想）：在适当条件下，\(\operatorname{rank} HF(L_0, L_1) \leq \sum_k \dim H^k(L_0)\)，且若
\(L_0 \pitchfork L_1\)，则
\(|L_0 \cap L_1| \geq \sum_k \dim H^k(L_0)\)。

\subsubsection{\texorpdfstring{246.2 Fukaya
\(A_\infty\)-范畴}{246.2 Fukaya A\_\textbackslash infty-范畴}}\label{fukaya-a_infty-ux8303ux7574}

\textbf{定义 246.3}（\(A_\infty\)-范畴）：一个
\textbf{\(A_\infty\)-范畴} \(\mathcal{C}\) 由以下数据组成： - 一组对象
\(\operatorname{Ob}(\mathcal{C})\) - 对每对对象
\(X, Y\)，一个分次向量空间 \(\operatorname{Hom}_{\mathcal{C}}(X, Y)\) -
对每个 \(k \geq 1\)，\(k\) 阶乘法运算
\(\mu^k: \operatorname{Hom}(X_{k-1}, X_k) \otimes \cdots \otimes \operatorname{Hom}(X_0, X_1) \to \operatorname{Hom}(X_0, X_k)\)，次数为
\(\deg \mu^k = 2 - k\)

这些运算满足 \(A_\infty\) 关系：对所有 \(k \geq 1\)，
\[\sum_{r+s+t = k} (-1)^{r+st} \mu^{r+1+t}(\mathrm{id}^{\otimes r} \otimes \mu^s \otimes \mathrm{id}^{\otimes t}) = 0\]

其中符号规则中 \((-1)^{r+st}\) 的指数由 Koszul 符号约定确定。特别地： -
\(\mu^1 \circ \mu^1 = 0\)（\(\mu^1\) 是微分） -
\(\mu^1(\mu^2(a, b)) = \mu^2(\mu^1(a), b) + (-1)^{|a|} \mu^2(a, \mu^1(b))\)（Leibniz
法则） - \(\mu^2\)
的结合律在同伦意义下成立：\(\mu^2(\mu^2(a, b), c) - \mu^2(a, \mu^2(b, c)) = \pm \mu^1(\mu^3(a, b, c)) \pm \mu^3(\mu^1(a), b, c) \pm \mu^3(a, \mu^1(b), c) \pm \mu^3(a, b, \mu^1(c))\)

\textbf{定义 246.4}（Fukaya 范畴 / Fukaya 1993）：辛流形 \((M, \omega)\)
的 \textbf{Fukaya 范畴} \(\mathcal{F}uk(M)\) 是一个
\(A_\infty\)-范畴，其中： - 对象：合适的 Lagrangian
子流形（带额外的结构：自旋结构、分次、局部系统） -
态射空间：\(\operatorname{Hom}(L_0, L_1) = CF^*(L_0, L_1)\)（Floer
链复形），其生成元为 \(L_0\) 与 \(L_1\) 的横截交点，分次由 Maslov
指标给出 - \(A_\infty\) 运算 \(\mu^k\)：\(\mu^1\) 是 Floer
微分（计数全纯带），\(\mu^k\)（\(k \geq 2\)）由 \(k+1\)
个边界标记点的全纯 \(k\)-边形的计数给出，即全纯映射
\(u: D^2 \setminus \{z_0, \ldots, z_k\} \to M\) 的模空间在
\(\mathbb{R}\) 平移下的商

\textbf{定理 246.2}（\(A_\infty\)-关系的几何来源）：Floer 微分的平方
\(\mu^1 \circ \mu^1 = 0\) 由模空间的边界给出（气泡化和断裂）。\(\mu^2\)
结合律的偏差 \(\mu^2(\mu^2(a, b), c) - \mu^2(a, \mu^2(b, c))\) 由
\(\mu^3\) 对边界算子的同伦给出。所有 \(A_\infty\)
关系来自模空间紧化边界的几何。

\subsubsection{246.3 Kontsevich
同调镜对称}\label{kontsevich-ux540cux8c03ux955cux5bf9ux79f0}

\textbf{定义 246.5}（同调镜对称猜想 / Homological Mirror
Symmetry，Kontsevich 1994）：对镜面对 \((X, \check{X})\)（\(X\) 是
Calabi-Yau 辛流形，\(\check{X}\) 是其镜面复流形），存在范畴等价：

\[D^b \mathcal{F}uk(X) \cong D^b \operatorname{Coh}(\check{X})\]

即 \(X\) 的 Fukaya 范畴的导出范畴（辛几何）等价于镜面 \(\check{X}\)
的相干层范畴的导出范畴（复代数几何）。

\textbf{定理 246.3}（镜对称的部分证明）： -
椭圆曲线：Polishchuk-Zaslow（1998）证明了 \(T^2\) 的同调镜对称 -
四次曲面（quartic surface）：Seidel（2015）和 Sheridan（2015）证明 -
环面簇的衍生范畴等价：由 Abouzaid、Fang、Liu 等人证明 -
完全交和超曲面的镜对称：多个作者推进

\subsubsection{246.4 SYZ 猜想与
Strominger-Yau-Zaslow}\label{syz-ux731cux60f3ux4e0e-strominger-yau-zaslow}

\textbf{定义 246.6}（SYZ 猜想，Strominger-Yau-Zaslow
1996）：接近大体积极限的 Calabi-Yau 流形 \(X\) 是特殊 Lagrangian 环面
\(T^n\) 的纤维化（\(T^n\)-对偶 / T-duality）。镜面 \(\check{X}\)
通过对偶纤维构造。

\textbf{定理 246.4}（SYZ 的 T-对偶图像）：在 SYZ 图像下，\(X\) 上的
Lagrangian 子流形对应 \(\check{X}\) 上的全纯向量丛（通过纤维 \(T^n\)
上的 Fourier-Mukai 变换），\(X\) 上的伪全纯曲线对应 \(\check{X}\)
上相干层的修正复形。

\bigskip\hrule\bigskip

\emph{本卷在微分几何（V12）和代数拓扑（V15）的基础上，系统建立了辛几何（Darboux
定理、Hamilton 力学、Noether 定理）、辛约化（Marsden-Weinstein
约化、Delzant 定理）、Gromov-Witten
理论（伪全纯曲线、量子上同调、WDVV）、切触几何（Reeb
动力学、切触同调、Giroux 定理）和 Fukaya 范畴与镜对称（Lagrangian Floer
同调、Kontsevich 猜想、SYZ 猜想）的理论体系。共 5 章。}

\newpage

\section{卷四十四：分形几何与几何测度论}\label{ux5377ux56dbux5341ux56dbux5206ux5f62ux51e0ux4f55ux4e0eux51e0ux4f55ux6d4bux5ea6ux8bba}

\begin{quote}
分形几何（Fractal Geometry）由 Benoît Mandelbrot 在 1970
年代系统建立，研究处处不可微的、具有自相似结构的「粗糙」集合。几何测度论（Geometric
Measure Theory）由 De Giorgi、Federer、Fleming 等人在 1960-1970
年代发展，使用测度论工具研究 \(\mathbb{R}^n\)
中「广义曲面」的几何与拓扑性质，特别是 Plateau
问题（存在以给定曲线为边界的极小面积曲面）的严格数学解。两个领域都处理传统光滑微分几何无法描述的「奇异」几何对象，在现代分析、变分法和图像处理中有深刻应用。本卷在实分析（V9）、测度论（V9
Ch 45）和泛函分析（V14）的基础上，建立 Hausdorff
测度、分形维数、整流与电流、Plateau 问题和变分法的几何测度论基础。
\end{quote}

\bigskip\hrule\bigskip

\subsection{第247章：Hausdorff
测度与分形维数}\label{ux7b2c247ux7ae0hausdorff-ux6d4bux5ea6ux4e0eux5206ux5f62ux7ef4ux6570}

分形几何的核心概念是用非整数维数来度量集合的「粗糙度」。Hausdorff 测度和
Hausdorff 维数是最基本的工具。

\subsubsection{247.1 Hausdorff 测度}\label{hausdorff-ux6d4bux5ea6}

\textbf{定义 247.1}（\(s\)-维 Hausdorff 测度）：对 \(s \geq 0\) 和
\(A \subset \mathbb{R}^n\)，定义

\[\mathcal{H}^s(A) = \lim_{\delta \to 0} \mathcal{H}^s_\delta(A)\]

其中

\[\mathcal{H}^s_\delta(A) = \inf \left\{ \sum_{i=1}^\infty (\operatorname{diam} U_i)^s : A \subset \bigcup_i U_i,\ \operatorname{diam} U_i \leq \delta \right\}\]

\(\mathcal{H}^s\) 是 \(\mathbb{R}^n\) 上的外测度（实际上是 Borel
测度）。

\textbf{命题 247.1}（Hausdorff 测度的基本性质）： - \(\mathcal{H}^0\)
是计数测度 - \(\mathcal{H}^n\) 在 \(\mathbb{R}^n\) 上与 Lebesgue
测度相差常数因子：\(\mathcal{H}^n = c_n \mathcal{L}^n\)（\(c_n\) 是
\(n\) 维单位球的体积相关常数） - 对
\(s > n\)，\(\mathcal{H}^s(\mathbb{R}^n) = 0\) - 若 \(f\) 是 Lipschitz
映射，则
\(\mathcal{H}^s(f(A)) \leq (\operatorname{Lip}(f))^s \mathcal{H}^s(A)\)

\textbf{定理 247.1}（Hausdorff 测度的行为）：函数
\(s \mapsto \mathcal{H}^s(A)\) 在某个唯一的 \(s_0\) 处从 \(\infty\) 跳到
\(0\)。即在 \(s < s_0\) 时 \(\mathcal{H}^s(A) = \infty\)，\(s > s_0\) 时
\(\mathcal{H}^s(A) = 0\)（\(\mathcal{H}^{s_0}(A)\) 可能为
\(0\)、正有限或无穷）。

\emph{证明概要}（Hausdorff维数的基本性质）：设
\(\mathcal{H}^t(A) < \infty\)，对任意 \(s > t\) 和 \(\delta\)-覆盖
\(\{U_i\}\)，\(\sum (\operatorname{diam} U_i)^s \leq \delta^{s-t} \sum (\operatorname{diam} U_i)^t\)，取
\(\delta \to 0\) 得 \(\mathcal{H}^s(A) = 0\)。故 \(\mathcal{H}^s\)
至多在一个临界值处从 \(\infty\) 跳到
\(0\)。\textbf{可数稳定性}：\(\mathcal{H}^s(\bigcup_n A_n) \leq \sum_n \mathcal{H}^s(A_n)\)
由测度的可数次可加性直接得到。\textbf{Lipschitz不变性}：\(\mathcal{H}^s(f(A)) \leq (\operatorname{Lip}(f))^s \mathcal{H}^s(A)\)，证：每个
\(\delta\)-覆盖 \(\{U_i\}\) 在 \(f\) 下变为
\(\operatorname{Lip}(f)\delta\)-覆盖
\(\{f(U_i)\}\)，取和即得。这些性质使得Hausdorff维数成为刻画分形粗糙度的最基本工具。∎

\subsubsection{247.2 分形维数}\label{ux5206ux5f62ux7ef4ux6570}

\textbf{定义 247.2}（Hausdorff 维数）：集合 \(A \subset \mathbb{R}^n\)
的\textbf{Hausdorff 维数}为

\[\dim_H A = \sup\{s \geq 0 : \mathcal{H}^s(A) = \infty\} = \inf\{s \geq 0 : \mathcal{H}^s(A) = 0\}\]

Hausdorff 维数可以是任意非负实数。

\textbf{定义 247.3}（盒维数 / Box-Counting Dimension）：设
\(N_\delta(A)\) 是覆盖 \(A\) 所需的直径为 \(\delta\)
的集合的最少数目。\textbf{下盒维数}和\textbf{上盒维数}分别为

\[\underline{\dim}_B A = \liminf_{\delta \to 0} \frac{\log N_\delta(A)}{-\log \delta}, \quad \overline{\dim}_B A = \limsup_{\delta \to 0} \frac{\log N_\delta(A)}{-\log \delta}\]

若上下盒维数相等，则 \(\dim_B A\) 是\textbf{盒维数}（也称 Minkowski
维数）。一般
\(\dim_H A \leq \underline{\dim}_B A \leq \overline{\dim}_B A\)。

\textbf{注 247.1}（Hausdorff 维数与盒维数的关系）：不等式
\(\dim_H A \leq \underline{\dim}_B A\)
成立的原因是：盒维数使用固定大小的覆盖（所有覆盖集直径
\(\leq \delta\)），而 Hausdorff
维数允许任意尺寸的覆盖（各覆盖集直径可不同），故 Hausdorff
测度更''精细''，维数更低。严格不等式的典型例子：\(\mathbb{Q} \cap [0,1]\)
作为可数集满足 \(\dim_H = 0\)，但因稠密性其盒维数
\(\dim_B = 1\)。另一个经典例子是
\(\{1, 1/2, 1/3, \ldots\} \cup \{0\}\)：\(\dim_H = 0\) 但盒维数为
\(1/2\)。这些例子表明盒维数对''可数个点的极限密度''敏感，而 Hausdorff
维数对于可数集恒为零。对''好''的自相似分形（如 Cantor 集、Sierpiński
三角形），二者通常相等。

\textbf{例 247.1}（经典分形集的维数）： - Cantor 三分集
\(C\)：\(\dim_H C = \dim_B C = \frac{\log 2}{\log 3} \approx 0.631\) -
Sierpiński 三角形：\(\dim_H = \frac{\log 3}{\log 2} \approx 1.585\) -
Koch 雪花曲线边界：\(\dim_H = \frac{\log 4}{\log 3} \approx 1.262\) -
Menger 海绵：\(\dim_H = \frac{\log 20}{\log 3} \approx 2.727\)

\subsubsection{247.3
自相似集与迭代函数系}\label{ux81eaux76f8ux4f3cux96c6ux4e0eux8fedux4ee3ux51fdux6570ux7cfb}

\textbf{定义 247.4}（迭代函数系 /
IFS）：\textbf{迭代函数系}是一族收缩映射
\(\{f_1, \ldots, f_m\}\)（\(f_i: \mathbb{R}^n \to \mathbb{R}^n\)，\(\operatorname{Lip}(f_i) < 1\)）。IFS
的\textbf{吸引子}（或不变集）\(K\) 是满足 \(K = \bigcup_{i=1}^m f_i(K)\)
的唯一非空紧致集合。

\textbf{定理 247.2}（Hutchinson 定理，1981）：设 \(\{f_i\}\) 是相似比为
\(c_i\) 的相似压缩。若开集条件（OSC）成立（存在非空有界开集 \(V\) 使得
\(f_i(V) \subset V\) 且 \(f_i(V)\) 两两不交），则吸引子的 Hausdorff 维数
\(s\) 满足\textbf{Moran 方程}：

\[\sum_{i=1}^m c_i^s = 1\]

且 \(0 < \mathcal{H}^s(K) < \infty\)。

\emph{证明概要}（Moran-Hutchinson公式）：设 \(K\) 为IFS吸引子，则
\(K = \bigcup_i f_i(K)\)。由Hausdorff测度的性质：\(\mathcal{H}^s(K) \leq \sum c_i^s \mathcal{H}^s(K)\)。上界估计：对
\(K\) 构造由压缩迭代生成的覆盖，使用开集条件保证覆盖不重叠，得
\(\sum c_i^s \leq 1\)，故维数不超过 Moran 方程的根。下界估计：利用开集
\(V\) 上分布的质量（规范测度）和能量积分方法，证明
\(\mathcal{H}^s(K) > 0\)，从而 \(\sum c_i^s \geq 1\)。结合得 \(s\) 满足
\(\sum c_i^s = 1\)。OSC确保压缩像不重叠，使得维数计算可分离为各压缩映射的贡献之和。∎

\textbf{例 247.2}（Moran 方程的应用）： - Cantor
三分集：\(m = 2\)，\(c_1 = c_2 = 1/3\)，\(2 \cdot (1/3)^s = 1 \Rightarrow s = \log_3 2\)
- Sierpiński
三角形：\(m = 3\)，\(c_i = 1/2\)，\(3 \cdot (1/2)^s = 1 \Rightarrow s = \log_2 3\)

\bigskip\hrule\bigskip

\subsection{第248章：分形结构的自相似与多重分形}\label{ux7b2c248ux7ae0ux5206ux5f62ux7ed3ux6784ux7684ux81eaux76f8ux4f3cux4e0eux591aux91cdux5206ux5f62}

许多分形不能仅由一个维数完全描述------它们具有多层次的分形结构，需要\textbf{多重分形分析}（multifractal
analysis）来刻画。

\subsubsection{248.1
自仿射分形与随机分形}\label{ux81eaux4effux5c04ux5206ux5f62ux4e0eux968fux673aux5206ux5f62}

\textbf{定义 248.1}（自仿射分形 / Self-Affine Fractal）：由仿射压缩
\(f_i(x) = A_i x + b_i\)（\(A_i\) 是收缩矩阵，各方向收缩比不同）组成的
IFS 的吸引子。与自相似分形不同，自仿射分形的不同方向有不同的伸缩行为。

\textbf{例 248.1}（分数 Brown 运动 / FBM）：Hurst 指数 \(H \in (0, 1)\)
的分数 Brown 运动 \(B_H(t)\) 是自仿射的：\(B_H(ct)\) 与 \(c^H B_H(t)\)
同分布。\(B_H\) 的图的 Hausdorff 维数和盒维数为 \(2 - H\)。

\textbf{定义 248.2}（随机分形与 DLA）：\textbf{扩散限制聚集}（DLA /
Diffusion-Limited Aggregation，Witten-Sander 1981）通过随机粒子沿 Brown
运动轨迹黏附生长的过程产生随机分形。DLA 团簇的维数约为
\(1.71\)（二维）。

\subsubsection{248.2
多重分形分析}\label{ux591aux91cdux5206ux5f62ux5206ux6790}

\textbf{定义 248.3}（多重分形谱 / Multifractal Spectrum）：对测度
\(\mu\) 上的奇异性分析，定义局部 Hölder 指数
\(\alpha(x) = \lim_{r \to 0} \frac{\log \mu(B(x, r))}{\log r}\)。\textbf{多重分形谱}
\(f(\alpha)\) 是满足 \(\alpha(x) = \alpha\) 的点的 Hausdorff 维数。

\textbf{定义 248.4}（Rényi 维数与配分函数）：定义配分函数
\(Z_q(\delta) = \sum_i \mu(B_i)^q\)（\(B_i\) 是边长 \(\delta\)
的网格覆盖）。广义维数 \(D_q\) 满足
\(Z_q(\delta) \sim \delta^{(q-1)D_q}\)（\(\delta \to 0\)）。通过
Legendre
变换：\(f(\alpha) = \inf_q (q \alpha - \tau(q))\)，\(\tau(q) = (q-1)D_q\)。

\textbf{定理 248.1}（多重分形形式论 / Frisch-Parisi
猜想，1985）：在湍流和许多其他物理系统中，多重分形谱 \(f(\alpha)\) 是
\(\cap\) 形的，\(f(\alpha)\) 的最大值为信息维数 \(D_1\)，\(D_0\)
是测度支撑的盒维数。

\subsubsection{248.3 渗流与分形}\label{ux6e17ux6d41ux4e0eux5206ux5f62}

\textbf{定义 248.5}（渗流模型 / Percolation
Model）：在二维网格上，每条边以概率 \(p\) 被「打开」（可通行）。当 \(p\)
超过临界概率 \(p_c = 1/2\) 时，存在无穷大连通分量。

\textbf{定理 248.2}（渗流团簇的分形维数，Stanley 1977）：在 \(p = p_c\)
时，渗流团簇是分形。二维临界渗流团簇的 Hausdorff 维数为
\(91/48 \approx 1.896\)。团簇的化学维数（沿团簇测量的最短路径距离）为
\(d_{\min} \approx 1.13\)。

\bigskip\hrule\bigskip

\subsection{第249章：整流与电流}\label{ux7b2c249ux7ae0ux6574ux6d41ux4e0eux7535ux6d41}

几何测度论使用\textbf{整流}和\textbf{电流}（rectifiable sets and
currents）的概念精密地研究奇异曲面和变分问题。

\subsubsection{249.1 可整流集}\label{ux53efux6574ux6d41ux96c6}

\textbf{定义 249.1}（\(k\)-维可整流集 / Rectifiable Set）：Borel 集
\(E \subset \mathbb{R}^n\) 是\textbf{可数
\(\mathcal{H}^k\)-可整流的}，如果存在可数多个 Lipschitz 映射
\(f_i: \mathbb{R}^k \to \mathbb{R}^n\) 使得

\[\mathcal{H}^k\left(E \setminus \bigcup_i f_i(\mathbb{R}^k)\right) = 0\]

即可整流集几乎被 Lipschitz 参数化的 \(k\) 维曲面覆盖。

\textbf{定理 249.1}（可整流集的结构定理）：\(k\)-维可整流集 \(E\) 在
\(\mathcal{H}^k\)-几乎处处有 \(k\) 维近似切空间
\(\operatorname{Tan}^k(E, x)\)。即对 \(\mathcal{H}^k\)-a.e.
\(x \in E\)，

\[\lim_{r \to 0} \frac{\mathcal{H}^k(E \cap B(x, r))}{r^k} = 1\]

（在适当的旋转下，\(E\) 的放大极限是一个 \(k\) 维平面）。

\textbf{定义 249.2}（纯不可整流集 / Purely Unrectifiable Set）：Borel 集
\(E\) 是\textbf{纯不可整流的}，如果对每个 \(k\)-维可整流集
\(F\)，\(\mathcal{H}^k(E \cap F) = 0\)。Besicovitch
证明了二维中存在纯一维不可整流集。

\subsubsection{249.2
整系数整流链}\label{ux6574ux7cfbux6570ux6574ux6d41ux94fe}

\textbf{定义 249.3}（整系数 \(k\)-维整流链 / Integral Rectifiable
\(k\)-Current）：\textbf{整系数 \(k\)-维整流链} \(T\) 是分布意义下的
\(k\)-维「定向曲面」，由可整流的 \(k\) 维集合 \(E\)、整数重数
\(\theta: E \to \mathbb{Z}\) 和定向 \(k\)-向量 \(\xi\) 组成，作为一般
\(k\)-形式 \(\omega\) 上的线性泛函：

\[T(\omega) = \int_E \langle \omega(x), \xi(x) \rangle \theta(x) \, d\mathcal{H}^k(x)\]

\textbf{定义 249.4}（质量与边界 / Mass and Boundary）：整流链 \(T\)
的\textbf{质量}为
\(\mathbf{M}(T) = \int_E |\theta(x)| d\mathcal{H}^k(x)\)。\(T\)
的\textbf{边界} \(\partial T\) 是对 \((k-1)\)-形式 \(\omega\)
定义的分布：\(\partial T(\omega) = T(d\omega)\)。

\textbf{定义 249.5}（整系数平坦链 / Integral Flat
Chain）：整系数平坦链是整系数整流链和边界为整流链的整系数整流链之和的
\(L^1_{\text{loc}}\)
极限。整系数平坦链空间对于变分法非常重要（适合取极限）。

\subsubsection{249.3 等周不等式与 Sobolev
不等式}\label{ux7b49ux5468ux4e0dux7b49ux5f0fux4e0e-sobolev-ux4e0dux7b49ux5f0f}

\textbf{定理 249.2}（De Giorgi 等周不等式，1958）：对任何 \(n\)
维（\(n \geq 2\)）整系数整流链 \(T\) 满足 \(\partial T = 0\)
且紧支撑，存在 \(n+1\) 维整系数整流链 \(S\) 满足 \(\partial S = T\) 且
\[\mathbf{M}(S)^{n/(n+1)} \leq C_n \mathbf{M}(T)\] 其中 \(C_n\)
是仅依赖于 \(n\) 的常数。

更精确的等周不等式的几何测度论版本是 Plateau 问题解的基础。

\textbf{定理 249.3}（Federer-Fleming 等周不等式）：对 \(n\)
维整系数整流链 \(T\) 满足 \(\partial T = 0\) 且紧支撑，存在 \(n+1\)
维整系数整流链 \(S\) 满足 \(\partial S = T\) 且

\[\mathbf{M}(S)^{n/(n+1)} \leq \gamma \mathbf{M}(T)\]

其中 \(\gamma\) 是仅依赖于 \(n\) 的常数。

\textbf{定理 249.4}（面积公式 / Area Formula）：设
\(f: \mathbb{R}^m \to \mathbb{R}^n\) 是 Lipschitz
映射（\(m \leq n\)），则对任意 \(\mathcal{L}^m\)-可测集
\(A \subset \mathbb{R}^m\)，
\[\int_A J_m f(x) \, d\mathcal{L}^m(x) = \int_{\mathbb{R}^n} \mathcal{H}^0(A \cap f^{-1}(y)) \, d\mathcal{H}^m(y)\]
其中 \(J_m f(x) = \sqrt{\det((Df(x))^T Df(x))}\) 是 \(m\)-维 Jacobian。

\textbf{定理 249.5}（余面积公式 / Coarea Formula）：设
\(f: \mathbb{R}^n \to \mathbb{R}^m\) 是 Lipschitz
映射（\(n \geq m\)），则对任意 \(\mathcal{L}^n\)-可测集
\(A \subset \mathbb{R}^n\)，
\[\int_A J_m f(x) \, d\mathcal{L}^n(x) = \int_{\mathbb{R}^m} \mathcal{H}^{n-m}(A \cap f^{-1}(y)) \, d\mathcal{L}^m(y)\]

\emph{证明概要}（面积/余面积公式）：面积公式的证明基于线性代数中的极分解：对线性映射
\(L: \mathbb{R}^m \to \mathbb{R}^n\)，Jac(\(L\)) 是 \(L\) 将
\(m\)-维体积放大的倍数。对Lipschitz映射，利用Rademacher定理（几乎处处可微），对仿射近似局部化。通过Vitali覆盖引理将整体积分分解为局部仿射贡献的和。余面积公式是面积公式在对偶维数情形下的推广：将
\(f\) 的等高面的 \(\mathcal{H}^{n-m}\)
测度关于水平集积分。核心工具是Fubini定理的几何版本------将空间分解为水平集。这两个公式是几何测度论中最核心的积分公式，推广了经典的变量替换公式。∎

\bigskip\hrule\bigskip

\subsection{第250章：Plateau
问题与极小曲面}\label{ux7b2c250ux7ae0plateau-ux95eeux9898ux4e0eux6781ux5c0fux66f2ux9762}

Plateau
问题（求以给定闭曲线为边界的极小面积曲面）是变分法的古典问题。几何测度论用电流理论给出了第一个完全严格的解。

\subsubsection{250.1 经典 Plateau
问题}\label{ux7ecfux5178-plateau-ux95eeux9898}

\textbf{定义 250.1}（Plateau 问题 / 几何测度论形式）：给定
\(\mathbb{R}^3\) 中的（可整流）闭曲线 \(\Gamma\)（一维整系数整流链满足
\(\partial \Gamma = 0\)）。求二维整系数整流链 \(T\) 满足： -
\(\partial T = \Gamma\)（\(T\) 以 \(\Gamma\) 为边界） -
\(\mathbf{M}(T) \leq \mathbf{M}(S)\)（对所有满足 \(\partial S = \Gamma\)
的 \(S\)）

\textbf{定理 250.1}（Douglas-Radó 解法，1931/1930）：任何
\(\mathbb{R}^3\) 中的可整流 Jordan
曲线都张成某个极小面积曲面（参数化圆盘映射）。Douglas 因此获得首届
Fields 奖（1936）。

\textbf{定理 250.2}（Federer-Fleming 解法 / 几何测度论中的 Plateau
问题解，1960）：对任何一维整系数整流链 \(\Gamma\) 满足
\(\partial \Gamma = 0\) 且紧支撑，存在二维整系数整流链 \(T\)
实现极小质量。

\emph{证明思路}：取质量最小化序列
\(\{T_i\}\)。用等周不等式控制边界，用紧致性定理（Federer-Fleming
紧致性定理）提取子列收敛到整系数平坦链。由下半连续性（质量在平坦范数下是下半连续的），极限链是最小的。用正则性理论证明极限是整系数整流链。∎

\subsubsection{250.2
极小曲面的正则性}\label{ux6781ux5c0fux66f2ux9762ux7684ux6b63ux5219ux6027}

\textbf{定理 250.3}（De Giorgi
正则性定理，1961）：质量最小化的二维整系数整流链 \(T\) 在
\(\mathbb{R}^n\) 中除了孤立奇点外是光滑解析曲面。在三维中，\(T\)
是嵌入曲面（无奇点）。

\textbf{定理 250.4}（Allard 正则性定理，1972）：质量最小的整系数 \(k\)
维整流链在 \(\mathbb{R}^n\) 中具有解析的 Hausdorff 维数 \(k\)
的支撑，除一个 \(\mathcal{H}^{k-2}\) 维的闭奇异集外。

\textbf{定义 250.2}（奇异集的维数）：\(k\)
维面积最小化整系数整流链的奇异集 \(\operatorname{Sing}(T)\) 满足
\(\dim_H \operatorname{Sing}(T) \leq k - 7\)（对超曲面
\(k = n-1\)）。特别是，在 \(\mathbb{R}^8\) 以下维数中无奇异点。

\subsubsection{250.3
自由边界问题与毛细曲面}\label{ux81eaux7531ux8fb9ux754cux95eeux9898ux4e0eux6bdbux7ec6ux66f2ux9762}

\textbf{定义 250.3}（自由边界 Plateau
问题）：极小曲面不仅张成给定曲线，还要与某固定曲面
\(W\)（如容器壁）以自由边界条件相交。极小曲面与 \(W\)
沿边界以垂直角度相交。

\textbf{定理
250.5}（自由边界问题的存在性与正则性）：在适当的条件下，自由边界 Plateau
问题有质量最小化解，其在内部光滑，在自由边界处（除特定维数的奇异集）光滑。

\textbf{定义 250.4}（毛细曲面与 Gauss
自由能）：毛细曲面满足均值曲率与高度成比例的方程（Young-Laplace
方程）。Gauss 自由能泛函结合了面积和重力势能。

\bigskip\hrule\bigskip

\subsection{第251章：变分法的几何测度论}\label{ux7b2c251ux7ae0ux53d8ux5206ux6cd5ux7684ux51e0ux4f55ux6d4bux5ea6ux8bba}

几何测度论为变分法中的许多经典问题提供了严格的数学基础，包括集合的边界变分（Caccioppoli
集）和具有奇异性的最小化子。

\subsubsection{251.1 Caccioppoli
集与有限周长集}\label{caccioppoli-ux96c6ux4e0eux6709ux9650ux5468ux957fux96c6}

\textbf{定义 251.1}（有限周长集 / Caccioppoli 集）：Borel 集
\(E \subset \mathbb{R}^n\) 是\textbf{有限周长集}（Caccioppoli
集），如果其特征函数 \(\chi_E\) 是有界变差函数（\(BV\)），即分布导数
\(D\chi_E\) 是有限 Radon 测度。周长定义为
\(\operatorname{Per}(E) = \int |D\chi_E|\)。

\textbf{定理 251.1}（De Giorgi 结构定理）：有限周长集 \(E\) 的约化边界
\(\partial^* E\) 是 \(\mathcal{H}^{n-1}\)-可整流的，且
\(D\chi_E = \nu_E \mathcal{H}^{n-1} \llcorner \partial^* E\)（\(\nu_E\)
是广义外法向量）。即周长等于约化边界的 Hausdorff 面积。

\textbf{定理
251.2}（等周不等式的几何测度论证明）：\(\operatorname{Per}(E) \geq n \omega_n^{1/n} (\mathcal{L}^n(E))^{1 - 1/n}\)，等号成立当且仅当
\(E\) 是球（模去零测集）。

\subsubsection{\texorpdfstring{251.2 \(BV\)
函数与总变差}{251.2 BV 函数与总变差}}\label{bv-ux51fdux6570ux4e0eux603bux53d8ux5dee}

\textbf{定义 251.2}（有界变差函数 / \(BV\) 函数）：\(u \in L^1(\Omega)\)
属于 \(BV(\Omega)\)，如果其分布导数 \(Du\) 是有限 Radon 测度。总变差为
\(|Du|(\Omega)\)。\(BV\) 空间是 Sobolev 空间 \(W^{1,1}\)
的自然推广（允许跳跃间断）。

\textbf{定理 251.3}（\(BV\) 函数的紧致性 / \(BV\)
紧嵌入定理）：若有界开集 \(\Omega\) 具有 Lipschitz 边界且
\(\{u_k\} \subset BV(\Omega)\) 满足
\(\sup_k \|u_k\|_{BV} < \infty\)，则存在子列在 \(L^1(\Omega)\)
中强收敛到 \(u \in BV(\Omega)\)。

\textbf{定理 251.4}（\(BV\) 函数的精细结构 / Federer-Vol'pert
定理）：\(u \in BV(\Omega)\) 关于 \(\mathcal{H}^{n-1}\)-可整流的跳跃集
\(J_u\) 具有近似极限和跳跃。\(Du\) 可分解为绝对连续部分、跳跃部分和
Cantor 部分。

\emph{证明概要}（Federer-Vol'pert分解）：\(BV\)函数 \(u\) 的分布导数
\(Du\)
是Radon测度，由Radon-Nikodym定理可分解为三部分：\(Du = D^a u + D^j u + D^c u\)。绝对连续部分
\(D^a u = \nabla u \, d\mathcal{L}^n\)（\(\nabla u\)
是近似梯度）。跳跃部分
\(D^j u = (u^+ - u^-) \nu_u \mathcal{H}^{n-1} \llcorner J_u\)，其中
\(J_u\) 是跳跃集（\(\mathcal{H}^{n-1}\)-可整流），\(u^\pm\) 是沿法向
\(\nu_u\) 的近似极限。Cantor部分 \(D^c u\) 与 \(\mathcal{L}^n\)
奇异且无原子（如Cantor函数的导数）。关键工具是Blaschke选择定理和Besicovitch覆盖引理，用于建立
\(BV\) 函数的近似极限和跳跃的测度论结构。此定理完整刻画了 \(BV\)
函数的不光滑性来源。∎

\subsubsection{251.3
图像处理中的变分方法}\label{ux56feux50cfux5904ux7406ux4e2dux7684ux53d8ux5206ux65b9ux6cd5}

\textbf{定义 251.3}（Mumford-Shah 泛函 / 图像分割，1989）：

\[E(u, K) = \int_{\Omega \setminus K} |\nabla u|^2 dx + \alpha \mathcal{H}^{n-1}(K) + \beta \int_{\Omega} (u - g)^2 dx\]

其中 \(u\) 是分段光滑的分割，\(K\) 是边界跳跃集，\(g\)
是观测到的噪声图像。

\textbf{定义 251.4}（Rudin-Osher-Fatemi 模型 / 全变差去噪，1992）：

\[\min_u \left\{ \int_{\Omega} |Du| + \frac{\lambda}{2} \int_{\Omega} (u - g)^2 dx \right\}\]

全变差正则化保持图像边缘（不模糊跳跃），是现代图像处理的基石。

\textbf{定理 251.5}（ROF 模型的解的性质）：ROF 模型的解 \(u\) 是 \(g\)
的「全变差流」在尺度 \(\lambda^{-1}\)
下的结果。大尺度结构被保留，小尺度振荡被滤除。

\subsubsection{251.4 集中紧致性原理与 Gamma
收敛}\label{ux96c6ux4e2dux7d27ux81f4ux6027ux539fux7406ux4e0e-gamma-ux6536ux655b}

\textbf{定义 251.5}（Gamma 收敛 / \(\Gamma\)-收敛，De Giorgi
1970s）：泛函序列 \(F_n\) \textbf{\(\Gamma\)-收敛}到
\(F\)（\(F_n \overset{\Gamma}{\to} F\)），如果： - \textbf{下界}：对任何
\(x_n \to x\)，\(\liminf F_n(x_n) \geq F(x)\) - \textbf{上界}：对任何
\(x\)，存在收敛到 \(x\) 的序列 \(x_n\)，使得
\(\limsup F_n(x_n) \leq F(x)\)

\textbf{定理 251.6}（\(\Gamma\)-收敛与最小化子的收敛）：若
\(F_n \overset{\Gamma}{\to} F\) 且 \(x_n\) 是 \(F_n\) 的最小化子，则
\(x_n\) 的每个聚点是 \(F\) 的最小化子。

\textbf{定理 251.7}（Modica-Mortola 定理 / Allen-Cahn 泛函的
\(\Gamma\)-收敛）：Allen-Cahn 泛函

\[F_\varepsilon(u) = \int_{\Omega} \left( \frac{\varepsilon}{2} |\nabla u|^2 + \frac{1}{\varepsilon} W(u) \right) dx\]

（\(W(u) = \frac{1}{4}(1-u^2)^2\) 是双井势）\(\Gamma\)-收敛到周长泛函
\(\sigma \operatorname{Per}(\{u = 1\})\)（模去常数），其中
\(\sigma = \int_{-1}^1 \sqrt{2W(s)} ds\)。

\bigskip\hrule\bigskip

\emph{本卷在实分析（V9）和泛函分析（V14）的基础上，系统建立了 Hausdorff
测度与分形维数、自相似与多重分形分析、整流与电流（Federer
理论）、Plateau 问题（Douglas
解与正则性）和几何测度论变分法（Caccioppoli 集、BV
函数、\(\Gamma\)-收敛）的理论体系。全书以分形几何的「无限精细结构」和几何测度论的「广义曲面」收尾，展示了数学处理非光滑几何对象的强大能力。共
5 章。}

\newpage

\section{卷四十五：运筹学}\label{ux5377ux56dbux5341ux4e94ux8fd0ux7b79ux5b66}

\begin{quote}
运筹学是研究在有限资源约束下做出最优决策的数学学科，起源于二战期间的军事后勤优化。战后被广泛应用于工业生产、交通运输、供应链管理和公共服务等领域。本卷在线性规划（V27
Ch
132）和组合数学（V26）的基础上，系统建立线性规划对偶与灵敏度、整数规划、网络流、排队论和库存与供应链的数学理论。
\end{quote}

\bigskip\hrule\bigskip

\subsection{第252章：线性规划对偶与灵敏度分析}\label{ux7b2c252ux7ae0ux7ebfux6027ux89c4ux5212ux5bf9ux5076ux4e0eux7075ux654fux5ea6ux5206ux6790}

线性规划的对偶理论是运筹学最优美的数学理论之一，揭示了原问题与对偶问题之间的深刻经济含义。

\subsubsection{252.1 对偶理论}\label{ux5bf9ux5076ux7406ux8bba}

\textbf{定义 252.1}（线性规划的对偶）：原问题（标准形）：
\[\min c^T x \quad \text{s.t.} \quad Ax = b,\ x \geq 0\]

其对偶问题为：\[\max b^T y \quad \text{s.t.} \quad A^T y \leq c\]

\textbf{定理 252.1}（弱对偶定理）： 若 \(x\) 是原可行解，\(y\)
是对偶可行解，则 \(c^T x \geq b^T y\)。原目标值始终大于等于对偶目标值。

\textbf{定理 252.2}（强对偶定理）：若原问题有有限最优解
\(x^*\)，则对偶问题也有最优解 \(y^*\)，且 \(c^T x^* = b^T y^*\)。

\textbf{定理 252.3}（互补松弛条件）：\(x^*\) 和 \(y^*\)
分别为原和对偶最优解的充要条件是：
\[x_j^* (c_j - \sum_i a_{ij} y_i^*) = 0 \quad (\forall j)\]

即若某变量为正，则对应的对偶约束为紧；若对偶约束为松，则对应的原变量为零。

\textbf{定义 252.2}（影子价格 / Shadow Price）：对偶变量 \(y_i\)
的经济含义是约束 \(i\) 右端项增加一个单位时最优值的边际变化量------资源
\(i\) 的「影子价格」。

\subsubsection{252.2 灵敏度分析}\label{ux7075ux654fux5ea6ux5206ux6790}

\textbf{定义 252.3}（灵敏度分析 / Sensitivity
Analysis）：研究参数（\(c\)、\(b\)、\(A\)）变化时最优解和最优值如何变化。

\textbf{定理 252.4}（目标系数的容许变化范围）：设 \(B\)
是最优基。非基变量 \(x_j\) 的目标系数 \(c_j\) 变化 \(\Delta\)
时，最优基保持不变当且仅当
\(\bar{c}_j + \Delta \geq 0\)（极小化问题），其中 \(\bar{c}_j\)
是已约化成本。

\textbf{定理 252.5}（右端项的容许变化范围）：设
\(B^{-1} b \geq 0\)（可行性条件）。\(b\) 的分量变化允许范围为
\[\max_{i: (B^{-1})_{ik} > 0} \frac{-(B^{-1}b)_i}{(B^{-1})_{ik}} \leq \Delta b_k \leq \min_{i: (B^{-1})_{ik} < 0} \frac{-(B^{-1}b)_i}{(B^{-1})_{ik}}\]。

\subsubsection{252.3
网络规划与单纯形法的几何联系}\label{ux7f51ux7edcux89c4ux5212ux4e0eux5355ux7eafux5f62ux6cd5ux7684ux51e0ux4f55ux8054ux7cfb}

\textbf{定理 252.6}（单纯形法的几何解释）：单纯形法从多面体
\(P = \{x : Ax = b, x \geq 0\}\)
的一个顶点出发，沿下降边移动到相邻顶点，直至达到最优顶点。单纯形表对应于基矩阵
\(B\) 的逆：基本解为 \(x_B = B^{-1} b\)，非基变量为 \(0\)。

\subsubsection{252.4 非线性规划的 KKT
条件（附例）}\label{ux975eux7ebfux6027ux89c4ux5212ux7684-kkt-ux6761ux4ef6ux9644ux4f8b}

线性规划是凸优化的一类特例。对一般非线性规划，Karush-Kuhn-Tucker（KKT）条件提供了最优性的必要条件（凸情形下为充要条件）。

\textbf{定理 252.7}（KKT 条件，Karush 1939 / Kuhn-Tucker
1951）：考虑问题 \(\min f(x)\) s.t.
\(g_i(x) \leq 0\)，\(h_j(x) = 0\)。设 \(x^*\) 满足约束规范（如
LICQ），则存在 Lagrange 乘子 \(\mu_i \geq 0\)、\(\lambda_j\) 使得：

\[\nabla f(x^*) + \sum_i \mu_i \nabla g_i(x^*) + \sum_j \lambda_j \nabla h_j(x^*) = 0, \quad \mu_i g_i(x^*) = 0\]

\textbf{应用例}（带资源约束的投资组合优化）：投资 \(n\) 项资产，期望收益
\(\mathbf{r}\)，协方差矩阵 \(\Sigma\)（正定），预算约束
\(\sum_i x_i = 1\)，风险上限 \(\sigma_0^2\)。问题：

\[\min_{x} -\mathbf{r}^T x \quad \text{s.t.} \quad x^T \Sigma x \leq \sigma_0^2,\; \sum_i x_i = 1\]

KKT 条件给出
\(\mathbf{r} = 2\mu \Sigma x^* + \lambda \mathbf{1}\)，\(\mu(x^{*T}\Sigma x^* - \sigma_0^2) = 0\)。当风险约束为紧时（\(\mu > 0\)），解得
\(x^* = \frac{1}{2\mu}\Sigma^{-1}(\mathbf{r} - \lambda\mathbf{1})\)，配合预算约束和风险约束即可确定
\(\mu, \lambda\)。该解是均值-方差有效前沿的参数化表达，与 Markovitz
组合理论一致。

\bigskip\hrule\bigskip

\subsection{第253章：整数规划}\label{ux7b2c253ux7ae0ux6574ux6570ux89c4ux5212}

许多实际问题要求决策变量取整数值（人员数量、车辆数目、二选一决策），这引入了离散性------整数规划通常是
NP 难的。

\subsubsection{253.1
整数规划基础}\label{ux6574ux6570ux89c4ux5212ux57faux7840}

\textbf{定义 253.1}（整数规划 / Integer Programming /
IP）：\[\min c^T x \quad \text{s.t.} \quad Ax = b,\ x \geq 0,\ x \in \mathbb{Z}^n\]

若部分变量取整，则为\textbf{混合整数规划}（MIP）。若所有变量取值
\(\{0, 1\}\)，则为 \textbf{0-1 规划}。

\textbf{定义 253.2}（LP 松弛与对偶间隙）：放宽整数约束得到 LP 松弛。LP
松弛最优值 \(\leq\) IP
最优值（极小化问题）。两者的差距称为\textbf{对偶间隙}（duality gap）。

\textbf{定理 253.1}（全幺模矩阵 / Totally Unimodular）：若约束矩阵 \(A\)
是全幺模的（每个方子阵的行列式为 \(0, \pm 1\)），则对任意整数右端项
\(b\)，线性规划 \(Ax = b, x \geq 0\)
的所有基可行解都是整数的。典型例子：指派问题（assignment
problem）和网络流问题。

\subsubsection{253.2 Branch-and-Bound
与割平面法}\label{branch-and-bound-ux4e0eux5272ux5e73ux9762ux6cd5}

\textbf{定义 253.3}（分枝定界法 /
Branch-and-Bound）：递归地分解问题------在 LP 松弛解中选分数变量
\(x_j\)，分枝为 \(x_j \leq \lfloor x_j \rfloor\) 和
\(x_j \geq \lceil x_j \rceil\) 两个子问题。用 LP
松弛值作为下界（极小化时），剪去不可能含最优解的枝（bounding）。

\textbf{定理 253.2}（割平面法 / Cutting Plane Method / Gomory 1958）：若
LP 松弛最优解非整数，则从单纯形表生成 Gomory 割
\(-\bar{f}_j x_j - \cdots \leq -\bar{f}_0\)（\(\bar{f}_j\)
为分数部分），添加此约束割去当前分数解而不损失任何整数可行解。有限次迭代后收敛到整数最优解。

\emph{有限终止性论证}：每个 Gomory 割从当前 LP
最优解对应的单纯形表的一行导出：设基变量
\(x_{B_i} = \bar{b}_i - \sum_{j \in NB} \bar{a}_{ij} x_j\)，其中
\(\bar{b}_i = \lfloor \bar{b}_i \rfloor + \bar{f}_{i0}\)（\(0 < \bar{f}_{i0} < 1\)）。Gomory
割为 \(\sum_{j \in NB} \bar{f}_{ij} x_j \geq \bar{f}_{i0}\)（其中
\(\bar{f}_{ij} = \bar{a}_{ij} - \lfloor \bar{a}_{ij} \rfloor\)）。添加此割后，对偶单纯形法在当前基不可行但最优的前提下，经过一次对偶迭代恢复可行性并保持最优性，目标函数值不增（极小化问题）。由于整数规划的可行域是有限个整数点的集合（有界情形），而每个割严格排除当前分数
LP
最优解，结合目标函数的单调性和有限整数点集，算法必然在有限步内终止于整数最优解或证明不可行。在无界情形需结合附加的界约束以确保有限性。∎

\textbf{注
253.1}（分枝定界法的策略与定界方法）：分枝定界法的效率高度依赖分枝策略和定界方法。

\textbf{分枝策略}：常见策略包括（1）\textbf{最大不可行性}（maximum
infeasibility）：选择与整数偏离最远的分数变量 \(x_j\)（即
\(|x_j - \text{round}(x_j)|\)
最大的变量），可较快缩小搜索空间；（2）\textbf{伪费用分枝}（pseudo-cost
branching）：记录各变量分枝后目标函数的历史变化率，优先选择预计改善最大的变量；（3）\textbf{强分枝}（strong
branching）：在分枝前试探性求解若干候选分枝的 LP
松弛，选择使下界上升最多的分枝------虽每次分枝计算成本高但总节点数大幅减少，在
MIP 求解器中广泛应用。

\textbf{定界方法}：每个子节点的 LP
松弛值提供该子树的\textbf{下界}（极小化问题）。维护全局\textbf{上界}（当前最优整数解的目标值）。若某节点的下界
\(\geq\)
当前上界，则该节点可被\textbf{剪枝}（fathomed）。此外，若子问题的 LP
松弛不可行或已获得整数解，也终止该分枝。现代 MIP 求解器（如
Gurobi、CPLEX）还结合切割平面（cutting
planes）形成\textbf{分枝切割法}（branch-and-cut），在每个节点处动态添加有效不等式以收紧
LP 松弛。

\subsubsection{253.3
背包问题与动态规划}\label{ux80ccux5305ux95eeux9898ux4e0eux52a8ux6001ux89c4ux5212}

\textbf{定义 253.4}（0-1 背包问题 / Knapsack
Problem）：\[\max \sum_{i=1}^n v_i x_i \quad \text{s.t.} \quad \sum_{i} w_i x_i \leq W,\ x_i \in \{0, 1\}\]

其中 \(v_i\) 为物品价值，\(w_i\) 为重量，\(W\) 为容量。

\textbf{定理 253.3}（动态规划解背包）：定义 \(f(i, w)\) 为前 \(i\)
个物品在容量 \(w\)
下的最大价值。递推：\(f(i, w) = \max\{f(i-1, w), f(i-1, w-w_i) + v_i\}\)（若
\(w \ge w_i\)）。时间 \(O(nW)\)（伪多项式）。

\bigskip\hrule\bigskip

\subsection{第254章：网络流理论}\label{ux7b2c254ux7ae0ux7f51ux7edcux6d41ux7406ux8bba}

网络流研究网络中从源到汇的最大流和最小费用流问题，是运筹学在交通、通信和分配问题中的核心工具。

\subsubsection{254.1
最大流与最小割}\label{ux6700ux5927ux6d41ux4e0eux6700ux5c0fux5272}

\textbf{定义 254.1}（流网络 / Flow Network）：有向图
\(G = (V, E)\)，每条边 \((u, v)\) 有容量 \(c(u, v) \geq 0\)。源
\(s\)，汇 \(t\)。流 \(f: E \to \mathbb{R}_{\geq 0}\) 满足容量约束
\(0 \leq f(e) \leq c(e)\) 和守恒约束
\(\sum_{(u,v)\in E} f(u,v) = \sum_{(v,w)\in E} f(v,w)\)（对所有
\(v \neq s, t\)）。

\textbf{定义 254.2}（割 / Cut）：\(s\)-\(t\) 割
\((S, T)\)（\(s \in S, t \in T\)）的容量为
\(\sum_{u \in S, v \in T} c(u, v)\)。

\textbf{定理 254.1}（Ford-Fulkerson
最大流最小割定理，1956）：最大流的值等于最小 \(s\)-\(t\) 割的容量。

\[\max |f| = \min_{S: s \in S, t \notin S} \sum_{u \in S, v \notin S} c(u, v)\]

\emph{证明思路}：Ford-Fulkerson
算法沿增广路径增加流。当无增广路径时，从源可达的节点集形成最小割。最大流最小割定理是线性规划强对偶定理的特殊情形。∎

\textbf{定理 254.2}（Edmonds-Karp 算法，1972）：使用 BFS
找最短增广路径，Ford-Fulkerson 的复杂度为 \(O(|V| \cdot |E|^2)\)。

\subsubsection{254.2 最小费用流}\label{ux6700ux5c0fux8d39ux7528ux6d41}

\textbf{定义 254.3}（最小费用流 / Min-Cost Flow）：每条边 \((u, v)\)
有单位流费用 \(a(u, v)\)。在满足流量需求 \(d\)（从 \(s\) 到 \(t\)
的总流）的条件下最小化总费用 \(\sum_e a(e) f(e)\)。

\textbf{定理 254.3}（最小费用流的最优性条件 / 负环条件）：可行流 \(f\)
是费用最小的当且仅当残量网络 \(G_f\) 中不存在负费用有向环。

\subsubsection{254.3 匹配问题}\label{ux5339ux914dux95eeux9898}

\textbf{定义 254.4}（匹配 /
Matching）：图的匹配是两两不共享公共端点的边集合。\textbf{完美匹配}覆盖所有顶点。\textbf{最大权匹配}是权和最大的匹配。

\textbf{定理 254.4}（König 定理 /
二分图匹配）：在二分图中，最大匹配的边数 = 最小顶点覆盖的点数。

\textbf{定理 254.5}（Edmonds 花算法 / Blossom
Algorithm，1965）：一般图的最大匹配可在多项式时间（\(O(|V|^3)\)）内求得。花算法的核心是收缩奇数长的交错环（花/blossom）。

\bigskip\hrule\bigskip

\subsection{第255章：排队论}\label{ux7b2c255ux7ae0ux6392ux961fux8bba}

排队论使用随机过程理论来分析和优化服务系统中的等待现象，从电话交换机（Erlang）到计算机网络和医院管理。

\subsubsection{255.1
排队系统的基本要素}\label{ux6392ux961fux7cfbux7edfux7684ux57faux672cux8981ux7d20}

\textbf{定义 255.1}（Kendall 记号 /
\(A/B/c/K/N/D\)）：排队模型由以下要素确定： - \(A\)：到达间隔分布（\(M\)
= Markov/Poisson，\(D\) = 确定性，\(G\) = 一般分布） -
\(B\)：服务时间分布（同 \(A\)） - \(c\)：服务器数量 -
\(K\)：系统容量（默认 \(\infty\)） - \(N\)：顾客群体大小（默认
\(\infty\)） - \(D\)：服务规则（FCFS、LCFS、优先权等）

\textbf{定义 255.2}（基本量）：\(L\) = 系统平均顾客数，\(L_q\) =
队列平均等待人数，\(W\) = 系统平均逗留时间，\(W_q\) = 平均等待时间。

\textbf{定理 255.1}（Little 定律，1961）：\(L = \lambda W\)（\(\lambda\)
为有效到达率）。\(L_q = \lambda W_q\)。Little
定律对所有排队系统成立（非常一般），是排队论最基本的恒等式。

\subsubsection{255.2 Markovian
排队系统}\label{markovian-ux6392ux961fux7cfbux7edf}

\textbf{定义 255.3}（\(M/M/1\) 排队）：Poisson 到达（速率
\(\lambda\)），指数服务时间（速率 \(\mu\)），单一服务器。设
\(\rho = \lambda / \mu\)（交通强度/traffic intensity）。

\textbf{定理 255.2}（\(M/M/1\) 的平稳分布）：\(\rho < 1\)
时平稳分布存在，系统状态为几何分布：\(\pi_n = (1-\rho)\rho^n\)（\(n \geq 0\)）。由此可推导：

\[L = \frac{\rho}{1-\rho},\quad L_q = \frac{\rho^2}{1-\rho},\quad W = \frac{1}{\mu - \lambda},\quad W_q = \frac{\rho}{\mu - \lambda}\]

\textbf{定义 255.4}（\(M/M/c\) 排队 / Erlang C 公式）：\(c\) 个服务器的
\(M/M/c\) 系统的稳态队列概率为：

\[P(\text{队列长度} > 0) = \frac{(c\rho)^c}{c! (1-\rho)} \left[ \sum_{k=0}^{c-1} \frac{(c\rho)^k}{k!} + \frac{(c\rho)^c}{c! (1-\rho)} \right]^{-1}\]

（Erlang C 公式）。\(M/M/c/c\) 的损失概率由 Erlang B 公式给出。

\subsubsection{255.3 排队网络}\label{ux6392ux961fux7f51ux7edc}

\textbf{定义 255.5}（Jackson 网络 / 开放 Jackson 网络，1957/1963）：节点
\(i\) 有 \(c_i\) 个服务器，外部 Poisson 到达速率
\(\gamma_i\)，顾客以概率 \(r_{ij}\) 从节点 \(i\) 路由到节点 \(j\)（或以
\(1 - \sum_j r_{ij}\) 离开系统）。

\textbf{定理 255.3}（Jackson 定理）：在开放 Jackson
网络中，节点状态独立，平稳分布为各节点 \(M/M/c_i\)
分布的乘积形式（乘积形解 / product-form）：

\[\pi(\mathbf{n}) = \prod_i \pi_i(n_i)\]

这是排队网络乘积形解的奠基定理，意味着各节点可以独立分析。

\bigskip\hrule\bigskip

\subsection{第256章：库存论与供应链管理}\label{ux7b2c256ux7ae0ux5e93ux5b58ux8bbaux4e0eux4f9bux5e94ux94feux7ba1ux7406}

库存论研究何时订货、订多少货以平衡订货成本、存储成本和缺货成本的最优策略。

\subsubsection{256.1
确定性库存模型}\label{ux786eux5b9aux6027ux5e93ux5b58ux6a21ux578b}

\textbf{定义 256.1}（经济订货量模型 / EOQ，Harris 1913）：需求率恒定
\(D\)（单位时间），固定订货成本 \(K\)，单位时间单位存储成本 \(h\)。EOQ
是使总成本（订货 + 存储）最小的订货量：

\[Q^* = \sqrt{\frac{2KD}{h}}\]

最优订货周期 \(T^* = Q^*/D = \sqrt{2K/(hD)}\)。年总成本
\(C(Q^*) = \sqrt{2KDh}\)。

\textbf{推导}：单位时间总成本
\(C(Q) = \frac{KD}{Q} + \frac{hQ}{2}\)（订货成本：\(K\) 每次，年订货
\(D/Q\) 次；存储成本：平均库存 \(Q/2\)，年存储费 \(hQ/2\)）。令
\(C'(Q) = -\frac{KD}{Q^2} + \frac{h}{2} = 0\)，解得
\(Q^* = \sqrt{\frac{2KD}{h}}\)。二阶导
\(C''(Q^*) = \frac{2KD}{Q^{*3}} > 0\)，确为极小值。EOQ
公式揭示了固定订货成本 \(K\) 增大促使大单订货、存储成本 \(h\)
增大促使小单订货的经济直觉。

\textbf{定义 256.2}（经济生产量模型 / EPQ）：生产速率
\(P > D\)，库存以速率 \(P - D\) 积累。最优批量
\(Q^* = \sqrt{\frac{2KD}{h(1 - D/P)}}\)。

\textbf{定义 256.3}（有数量折扣的 EOQ）：供应商对大量订货提供单价折扣
\(c(Q)\)。最优策略是计算每个价格区间的 EOQ 并比较总成本。

\subsubsection{256.2
随机库存模型与报童模型}\label{ux968fux673aux5e93ux5b58ux6a21ux578bux4e0eux62a5ux7ae5ux6a21ux578b}

\textbf{定义 256.4}（报童模型 / Newsvendor Problem）：需求 \(D\)
是随机变量（分布 \(F\)）。单位进货成本 \(c\)，售价 \(p\)，残值
\(s\)（\(s < c < p\)）。最优订货量 \(Q^*\) 满足临界分位数（critical
fractile）：

\[F(Q^*) = \frac{p - c}{p - s} = \frac{\text{underage cost}}{\text{underage} + \text{overage}}\]

其中 underage cost \(= p - c\)（缺货损失），overage cost
\(= c - s\)（过剩损失）。

\textbf{定义 256.5}（\((s, S)\) 策略）：定期盘点库存，若库存水平
\(\leq s\)（再订货点），则订货到
\(S\)（目标水平）。在随机需求下，\((s, S)\)
策略对单产品库存在一定条件下是最优的。

\subsubsection{256.3
供应链协调与牛鞭效应}\label{ux4f9bux5e94ux94feux534fux8c03ux4e0eux725bux97adux6548ux5e94}

\textbf{定义 256.6}（供应链合同 / Supply Chain
Contracts）：供应链由多个独立决策主体（供应商、制造商、零售商）组成。协调机制包括：
- \textbf{回购合同}（Buyback）：供应商以约定价格回购未售出商品 -
\textbf{收益共享合同}（Revenue
Sharing）：零售商分享部分收益以换取批发价折扣 -
\textbf{数量柔性合同}（Quantity
Flexibility）：零售商可调整订购量（在约定范围内）

\textbf{定义 256.7}（牛鞭效应 / Bullwhip
Effect）：需求信息向供应链上游传递时波动逐渐放大的现象（Lee-Padmanabhan-Whang
1997）。原因包括：需求预测更新、批量订购、价格波动和短缺博弈。

\textbf{牛鞭效应的数学表述}：设下游零售商面临需求
\(D_t = d + \rho D_{t-1} + \varepsilon_t\)（AR(1)
过程，\(\varepsilon_t \sim N(0, \sigma^2)\)），采用 \((s, S)\)
策略且提前期 \(L\)。使用移动平均法预测需求时，上游供应商接收到的订单方差
\(\operatorname{Var}(O_t)\) 与下游需求方差 \(\operatorname{Var}(D_t)\)
的关系为：

\[\frac{\operatorname{Var}(O_t)}{\operatorname{Var}(D_t)} \geq 1 + \frac{2L}{p} + \frac{2L^2}{p^2}\]

其中 \(p\) 为移动平均窗口长度（Lee et al.~1997）。该不等式表明：提前期
\(L\) 越长、预测越短视（\(p\)
越小），则牛鞭效应（方差放大倍数）越显著。这是 Bullwhip 效应的定量形式。

\textbf{定理
256.1}（信息共享的价值）：在两级供应链中，下游需求信息的共享可以减小供应链总成本的方差，减少牛鞭效应。最优协调可通过两级报童模型的合同设计实现。

\bigskip\hrule\bigskip

\emph{本卷在线性规划（V27）、组合数学（V26）和概率论（V21）的基础上，系统建立了线性规划对偶与灵敏度、整数规划、网络流理论、排队论和库存与供应链管理的数学理论。共
5 章。}

\newpage

\section{卷四十六：密码学数学}\label{ux5377ux56dbux5341ux516dux5bc6ux7801ux5b66ux6570ux5b66}

\begin{quote}
密码学是研究信息保密和认证的数学基础的科学。从古典的 Caesar
密码到现代的公钥密码（RSA、椭圆曲线密码）和后量子密码（格密码），密码学始终依赖于深刻的数学理论------数论、代数、组合学和复杂性理论。本卷在数论
II（V17）、抽象代数（V8、V13）和计算理论（V31）的基础上，系统建立编码理论深化、对称密码、公钥密码、格密码与后量子密码以及密码协议的数学理论。
\end{quote}

\bigskip\hrule\bigskip

\subsection{第257章：编码理论深化}\label{ux7b2c257ux7ae0ux7f16ux7801ux7406ux8bbaux6df1ux5316}

编码理论研究在噪声信道中可靠传输信息的方法，是 Shannon
信息论（V29）的直接应用和深化。

\subsubsection{257.1
错误检测与纠错}\label{ux9519ux8befux68c0ux6d4bux4e0eux7ea0ux9519}

\textbf{定义 257.1}（线性码 / Linear Code）：\([n, k, d]\) 线性码 \(C\)
是 \(\mathbb{F}_q^n\) 的 \(k\) 维子空间，极小距离
\(d = \min_{c \neq c' \in C} \operatorname{wt}(c - c')\)。\(C\) 可纠正
\(\lfloor (d-1)/2 \rfloor\) 个错误。

\textbf{定义 257.2}（生成矩阵与校验矩阵）：\textbf{生成矩阵}
\(G \in \mathbb{F}_q^{k \times n}\)（\(C = \{mG : m \in \mathbb{F}_q^k\}\)）。\textbf{校验矩阵}
\(H \in \mathbb{F}_q^{(n-k) \times n}\)（\(C = \{x : H x^T = 0\}\)）。\(HG^T = 0\)。

\textbf{定理 257.1}（Singleton 界）：\([n, k, d]\) 码满足
\(d \leq n - k + 1\)。达到此界的码称为 \textbf{MDS
码}（最大距离可分码）。

\textbf{定理 257.2}（Hamming 界与 Gilbert-Varshamov 界）： -
\textbf{Hamming
界}（球填充界）：\(q^k \sum_{i=0}^{\lfloor (d-1)/2 \rfloor} \binom{n}{i} (q-1)^i \leq q^n\)
- \textbf{Gilbert-Varshamov 界}（存在性界）：存在 \([n, k, d]\) 码满足
\(q^{n-k} > \sum_{i=0}^{d-2} \binom{n-1}{i} (q-1)^i\)

\subsubsection{257.2
经典代数编码}\label{ux7ecfux5178ux4ee3ux6570ux7f16ux7801}

\textbf{定义 257.3}（Reed-Solomon 码 / RS 码，1960）：取 \(n \leq q-1\)
个互异点 \(\alpha_1, \ldots, \alpha_n \in \mathbb{F}_q\)。RS 码定义为
\(\{(f(\alpha_1), \ldots, f(\alpha_n)) : f \in \mathbb{F}_q[x], \deg f < k\}\)。RS
码是 MDS 码（\(d = n - k + 1\)），广泛用于 CD、DVD、QR 码和深空通信。

\textbf{定义 257.4}（BCH 码与 Reed-Muller 码）： - \textbf{BCH
码}（Bose-Chaudhuri-Hocquenghem，1960）：指定纠错能力 \(t\)
的循环码，设计距离 \(\delta = 2t + 1\) - \textbf{Reed-Muller
码}：基于布尔函数的代数构造，\(RM(r, m)\) 参数为
\([2^m, \sum_{i=0}^r \binom{m}{i}, 2^{m-r}]\)

\textbf{定理 257.3}（Berlekamp-Massey 算法，1968-1969）：BCH
码的解码可通过求解线性递推关系在 \(O(n^2)\)
时间内完成。这是代数解码的核心算法。

\subsubsection{257.3 现代编码：Turbo 码、LDPC 与 Polar
码}\label{ux73b0ux4ee3ux7f16ux7801turbo-ux7801ldpc-ux4e0e-polar-ux7801}

\textbf{定义 257.5}（Turbo 码，Berrou-Glavieux-Thitimajshima
1993）：Turbo 码通过并行级联两个卷积码并使用迭代软判决解码，性能接近
Shannon 极限（差距 \(< 0.7\) dB）。

\textbf{定义 257.6}（LDPC 码 / 低密度校验码，Gallager 1963 / MacKay-Neal
1996）：LDPC 码由稀疏校验矩阵定义，使用置信传播（belief propagation）在
Tanner 图上的迭代解码。极长码长下接近 Shannon 极限。

\textbf{定理 257.4}（Polar 码的极性化定理，Arıkan
2009）：通过信道极化（channel polarization）构造的 Polar 码，在码长
\(N \to \infty\) 时可达对称信道容量。Polar
码是首个被严格证明可在所有对称二进制输入无记忆信道下达到 Shannon
信道容量的构造性编码方案。Polar 码已被 5G 标准采纳。

\bigskip\hrule\bigskip

\subsection{第258章：对称密码}\label{ux7b2c258ux7ae0ux5bf9ux79f0ux5bc6ux7801}

对称密码（私钥密码）使用相同的密钥进行加密和解密，是通信保密的最古老和最高效的形式。

\subsubsection{258.1
密码学基础概念}\label{ux5bc6ux7801ux5b66ux57faux7840ux6982ux5ff5}

\textbf{定义 258.1}（加密方案 / Encryption
Scheme）：对称加密方案由三个算法组成： -
\(\mathsf{Gen}(1^\lambda) \to k\)：密钥生成（安全参数 \(\lambda\)） -
\(\mathsf{Enc}(k, m) \to c\)：加密（明文 \(m\)，密文 \(c\)） -
\(\mathsf{Dec}(k, c) \to m\)：解密

\textbf{定义 258.2}（Shannon 完美保密 / Perfect
Secrecy，1949）：加密方案具有完美保密性，如果对任意分布 \(M\)、任意明文
\(m\) 和任意密文 \(c\)， \[\Pr[M = m | C = c] = \Pr[M = m]\]
等价地，\(|K| \geq |M|\) 且每个密钥以等概率使用（一次性密钥包 / one-time
pad）。

\textbf{定理 258.3}（Shannon完美保密性的充要条件）：对称加密方案
\((K, M, C, E, D)\) 具有完美保密性当且仅当 \(|K| \geq |M|\)
且密钥均匀随机分布。

\emph{证明概要}（一次一密的证明）：充分性：设有均匀随机密钥 \(k\) 且
\(|K| \geq |M|\)。对任意
\((m, c)\)，\(\Pr[C=c|M=m] = \Pr[K = k: E(k,m)=c] = 1/|K|\)（或0若不存在）。由于该条件概率不依赖于
\(m\)，Bayes公式给出 \(\Pr[M=m|C=c] = \Pr[M=m]\)，即完美保密。必要性：若
\(\Pr[M=m|C=c] = \Pr[M=m]\)，则对任意 \(m_1, m_2\) 和
\(c\)，\(\Pr[C=c|M=m_1] = \Pr[C=c|M=m_2]\)。这意味着对不同明文，密文分布相同，因此必须有至少与明文空间一样大的密钥空间
\(|K| \geq |M|\)，且每个密钥被等概率选中。一次一密（One-Time
Pad）是满足此条件的典范构造，其中 \(c = m \oplus k\)。∎

\textbf{定义 258.3}（不可区分性 / IND-CPA 与
IND-CCA）：现代密码学使用计算安全性（非信息论安全）： -
\textbf{IND-CPA}（选择明文攻击下的不可区分性）：攻击者不能区分两个自选明文的加密
-
\textbf{IND-CCA}（选择密文攻击下的不可区分性）：即使有解密谕示，攻击者也不能区分加密

\subsubsection{258.2
分组密码与流密码}\label{ux5206ux7ec4ux5bc6ux7801ux4e0eux6d41ux5bc6ux7801}

\textbf{定义 258.4}（分组密码 / Block
Cipher）：\(E: \{0,1\}^k \times \{0,1\}^n \to \{0,1\}^n\)。DES（\(n=64, k=56\)）、AES（\(n=128, k=128/192/256\)）。

\textbf{定义 258.5}（AES / 高级加密标准，Daemen-Rijmen 2001）：AES
是取代-置换网络（SPN），由 SubBytes（S 盒）、ShiftRows、MixColumns 和
AddRoundKey 四步组成，迭代 10/12/14 轮。

\textbf{定义 258.6}（流密码 / Stream
Cipher）：使用伪随机生成器从种子密钥生成密钥流：\(c_i = m_i \oplus G(k)_i\)。安全需求：\(G\)
的输出在计算上与真随机不可区分。

\textbf{定理 258.1}（伪随机函数与分组密码的安全性）：若 \(F_k(\cdot)\)
是伪随机函数（PRF），则 CTR 模式（计数器模式）加密是 IND-CPA 安全的。

\subsubsection{258.3 认证与 Hash}\label{ux8ba4ux8bc1ux4e0e-hash}

\textbf{定义 258.7}（消息认证码 /
MAC）：\(\mathsf{MAC}(k, m) \to t\)（标签）。安全性要求：即使能查询任意消息的
MAC 标签（adaptive chosen-message
attack），攻击者也不能为新消息伪造有效标签。

\textbf{定义 258.8}（Hash 函数安全性质）：密码学 Hash 函数
\(H: \{0,1\}^* \to \{0,1\}^n\) 应满足： - \textbf{抗原像}（preimage
resistance）：给定 \(y\)，难找 \(x\) 使 \(H(x) = y\) -
\textbf{抗第二原像}（second preimage resistance）：给定 \(x\)，难找
\(x' \neq x\) 使 \(H(x') = H(x)\) - \textbf{抗碰撞}（collision
resistance）：难找 \(x \neq x'\) 使 \(H(x) = H(x')\)

\textbf{定理 258.2}（Merkle-Damgård 构造）：若压缩函数 \(f\)
是抗碰撞的，则由其通过 Merkle-Damgård 迭代构造的 Hash
函数也是抗碰撞的。SHA-256 基于此构造。

\bigskip\hrule\bigskip

\subsection{第259章：公钥密码}\label{ux7b2c259ux7ae0ux516cux94a5ux5bc6ux7801}

公钥密码（非对称密码）使用两个不同的密钥（公钥和私钥），解决了对称密码的密钥分发问题，是现代安全通信的基础。

\subsubsection{259.1 RSA
与因子分解}\label{rsa-ux4e0eux56e0ux5b50ux5206ux89e3}

\textbf{定义 259.1}（RSA 密码体制，Rivest-Shamir-Adleman 1978）： -
\(\mathsf{Gen}\)：选大素数
\(p, q\)，\(N = pq\)，\(\varphi(N) = (p-1)(q-1)\)。选 \(e\) 满足
\(\gcd(e, \varphi(N)) = 1\)。计算 \(d = e^{-1} \bmod \varphi(N)\)。公钥
\((N, e)\)，私钥 \(d\) - \(\mathsf{Enc}((N, e), m) = m^e \bmod N\) -
\(\mathsf{Dec}(d, c) = c^d \bmod N\)

\textbf{定理 259.1}（RSA
的正确性）：\(c^d \bmod N = m^{ed} \bmod N = m\)。

\emph{证明}：由 \(ed \equiv 1 \pmod{\varphi(N)}\)，存在整数 \(k\) 使
\(ed = k\varphi(N) + 1\)，需证
\(m^{k\varphi(N)+1} \equiv m \pmod{pq}\)。分两种情况： 1. 若
\(\gcd(m, N) = 1\)：由 Euler 定理
\(m^{\varphi(N)} \equiv 1 \pmod N\)，故
\(m^{ed} = m^{k\varphi(N)+1} \equiv m \pmod N\)。 2. 若
\(\gcd(m, N) \neq 1\)：因 \(N = pq\) 且 \(p, q\)
为素数，\(\gcd(m, N) \neq 1\) 意味着 \(p \mid m\) 或
\(q \mid m\)（或二者）。先考虑模 \(p\)：若 \(p \mid m\)，则
\(m^{ed} \equiv 0 \equiv m \pmod p\)；若 \(p \nmid m\)，由 Fermat 小定理
\(m^{p-1} \equiv 1 \pmod p\)，故
\(m^{ed} = m \cdot (m^{p-1})^{k(q-1)} \equiv m \pmod p\)。同理对模 \(q\)
也有
\(m^{ed} \equiv m \pmod q\)。由中国剩余定理，\(m^{ed} \equiv m \pmod{pq}\)。∎

\textbf{定义 259.2}（RSA 假设与整数分解）：RSA
安全性基于\textbf{大整数分解的困难性}和 \textbf{RSA 问题}（已知
\(N, e, m^e\) 求
\(m\)）。目前最佳分解算法为数域筛法（NFS），亚指数复杂度
\(L_N(1/3, c)\)。

\textbf{定理 259.2}（Wiener 攻击 / 小解密指数的 RSA）：若
\(d < N^{1/4}/3\)，可从公钥 \((N, e)\)
在多项式时间内恢复私钥（利用连分数）。

\subsubsection{259.2
离散对数与椭圆曲线密码}\label{ux79bbux6563ux5bf9ux6570ux4e0eux692dux5706ux66f2ux7ebfux5bc6ux7801}

\textbf{定义 259.3}（Diffie-Hellman 密钥交换，1976）：公开参数：素数
\(p\) 和生成元 \(g \in \mathbb{F}_p^*\)。Alice 选 \(a\)，发送
\(g^a\)；Bob 选 \(b\)，发送 \(g^b\)。共享密钥为
\(g^{ab}\)。安全性基于\textbf{计算 Diffie-Hellman（CDH）假设}。

\textbf{定理
259.4}（Diffie-Hellman安全性归约）：若决策Diffie-Hellman（DDH）假设成立------即区分
\((g^a, g^b, g^{ab})\) 与 \((g^a, g^b, g^r)\)（\(r\)
随机）在计算上不可行------则Diffie-Hellman密钥交换在窃听攻击下是语义安全的。DDH蕴含CDH，且在一般循环群模型中，CDH的困难性等价于离散对数问题（DLP）的困难性：攻击者若能从
\(g^a, g^b\) 恢复 \(g^{ab}\)，则可利用随机自归约求解离散对数。椭圆曲线群
\(E(\mathbb{F}_q)\) 上的ECDLP：已知最有效攻击（Pollard rho）复杂度为
\(O(\sqrt{n})\)（\(n\)
为群阶），与有限域上的亚指数级DLP形成对比。这是ECC使用短密钥（256位）却提供与RSA-3072同等安全性的根本原因。∎

\textbf{定义 259.4}（椭圆曲线密码 / ECC，Miller 1985 / Koblitz
1987）：以椭圆曲线群 \(E(\mathbb{F}_q)\) 代替乘法群
\(\mathbb{F}_p^*\)。椭圆曲线离散对数问题（ECDLP）的最知名攻击复杂度为指数级（\(\sqrt{|E|}\)），允许使用更短的密钥（256
位 ECC \(\approx\) 3072 位 RSA 的安全性）。

\textbf{定义 259.5}（ElGamal 加密与 ECDSA 签名）： - \textbf{ElGamal
加密}：公钥 \(h = g^x\)。加密：随机 \(r\)，输出
\((g^r, m \cdot h^r)\)。解密：\(m = (h^r)^{-x} \cdot m \cdot h^r\) -
\textbf{ECDSA}：基于椭圆曲线的数字签名算法，是 Bitcoin 和 TLS
中最广泛使用的签名方案

\textbf{定理
259.5}（ElGamal加密的IND-CPA安全性）：在DDH假设下，ElGamal加密方案在选择明文攻击下具有不可区分性（IND-CPA安全）。密文
\((g^r, m \cdot h^r)\)
中第一分量均匀随机，第二分量在DDH假设下与随机不可区分------攻击者无法区分
\(h^r = g^{xr}\) 与随机群元素。解密需要私钥
\(x\)（即离散对数），安全性完全归约到CDH/DDH。将ElGamal移植到椭圆曲线群
\(E(\mathbb{F}_q)\)
上得到EC-ElGamal，在ECDLP假设下具有同等安全性，但密钥和密文长度大幅缩短。ElGamal本身具有乘法同态性：\(E(m_1) \cdot E(m_2) = (g^{r_1+r_2}, m_1 m_2 h^{r_1+r_2}) = E(m_1 m_2)\)，这是最早的同态加密方案之一。∎

\subsubsection{259.3
同态加密与安全计算}\label{ux540cux6001ux52a0ux5bc6ux4e0eux5b89ux5168ux8ba1ux7b97}

\textbf{定义 259.6}（同态加密 / Homomorphic Encryption）：加密方案
\(\mathsf{Enc}\) 是\textbf{同态的}，如果存在运算 \(\oplus\) 使
\(\mathsf{Enc}(m_1) \oplus \mathsf{Enc}(m_2) = \mathsf{Enc}(m_1 + m_2)\)（加法同态
/ Paillier）或 \(\otimes\) 使
\(\mathsf{Enc}(m_1) \otimes \mathsf{Enc}(m_2) = \mathsf{Enc}(m_1 \cdot m_2)\)（乘法同态
/ ElGamal）。

\textbf{定理 259.3}（完全同态加密 / FHE，Gentry
2009）：存在同时支持加法和乘法同态的加密方案（基于理想格），使得可以在密文上执行任意电路计算。Gentry
的蓝图：基于格的 somewhat 同态加密 + 自举（bootstrapping）实现完全同态。

\bigskip\hrule\bigskip

\subsection{第260章：格密码与后量子密码}\label{ux7b2c260ux7ae0ux683cux5bc6ux7801ux4e0eux540eux91cfux5b50ux5bc6ux7801}

量子计算（Shor 算法，1994）可以在多项式时间分解整数和求离散对数，威胁
RSA 和 ECC
的安全。\textbf{后量子密码}（PQC）研究在量子攻击下仍然安全的密码体制。

\subsubsection{260.1 格基础}\label{ux683cux57faux7840}

\textbf{定义 260.1}（格 / Lattice）：\(\mathbb{R}^n\) 中一组线性无关向量
\(\mathbf{b}_1, \ldots, \mathbf{b}_n\) 的所有整数线性组合构成格：

\[\mathcal{L} = \left\{ \sum_{i=1}^n x_i \mathbf{b}_i : x_i \in \mathbb{Z} \right\}\]

\textbf{定义 260.2}（格中的难解问题）： -
\textbf{SVP}（最短向量问题）：求格中最短的非零向量的长度
\(\lambda_1(\mathcal{L})\) -
\textbf{CVP}（最近向量问题）：给定目标点，求最近的格点 -
\textbf{SIS}（短整数解）：给定随机矩阵 \(A\)，找短的 \(x \neq 0\) 使
\(A x = 0 \pmod q\) - \textbf{LWE}（带误差学习 / Learning With
Errors，Regev 2005）：从 \((A, As + e)\) 中恢复秘密 \(s\)（\(e\)
是小误差）

\textbf{定理 260.1}（Regev 的 LWE 安全性归约，2005/2009）：LWE
问题至少与最坏情况的格问题（GapSVP 和 SIVP）一样难。即攻破 LWE
意味着可以用量子算法近似解决所有格中的最坏情况难解问题。

\textbf{定理 260.2}（LLL 算法 / Lenstra-Lenstra-Lovász 1982）：LLL
算法在多项式时间内找到长度不超过 \(\lambda_1\) 的 \(2^{(n-1)/2}\)
倍的短向量。LLL 是格基规约的基本工具。

\subsubsection{260.2
基于格的密码体制}\label{ux57faux4e8eux683cux7684ux5bc6ux7801ux4f53ux5236}

\textbf{定义 260.3}（Regev 的 LWE 加密，2005）：公钥
\(A \in \mathbb{Z}_q^{n \times m}\) 和 \(b = s^T A + e^T\)（误差 \(e\)
是短的）。加密比特 \(m \in \{0, 1\}\) 为
\((A^T r, b^T r + m\lfloor q/2 \rfloor)\)（\(r\) 小随机向量）。解密利用
\(s^T A^T r \approx b^T r\) 恢复 \(m\)。

\textbf{定义 260.4}（NTRU，Hoffstein-Pipher-Silverman 1996 /
CRYSTALS-Kyber / ML-KEM）：NTRU 基于多项式环
\(R = \mathbb{Z}_q[x]/(x^n - 1)\) 上的格问题。CRYSTALS-Kyber（NIST PQC
标准 2024）使用 Module-LWE 的 ML-KEM。

\textbf{定义 260.5}（CRYSTALS-Dilithium / ML-DSA）：基于 Module-SIS 和
Module-LWE 的数字签名方案，已被 NIST 标准化（2024）。

\subsubsection{260.3
其他后量子密码方法}\label{ux5176ux4ed6ux540eux91cfux5b50ux5bc6ux7801ux65b9ux6cd5}

\textbf{定义 260.6}（基于 Hash 的签名 / SPHINCS+）：Merkle
树一次签名（Lamport 签名、Winternitz
签名）与少数次签名（HORS）结合的免状态 Hash 签名方案。

\textbf{定义 260.7}（基于编码的密码 / Classic McEliece）：McEliece
密码体制（1978）基于一般线性码的解码困难性（NP 完全）。使用 Goppa
码（可有效解码）作为私钥，加扰后的公开码为公钥。

\bigskip\hrule\bigskip

\subsection{第261章：密码协议}\label{ux7b2c261ux7ae0ux5bc6ux7801ux534fux8bae}

密码协议使用密码学原语实现安全多方计算、零知识证明和数字签名等高级安全目标。

\subsubsection{261.1 数字签名}\label{ux6570ux5b57ux7b7eux540d}

\textbf{定义
261.1}（数字签名方案）：\(\mathsf{Gen}(1^\lambda) \to (pk, sk)\)，\(\mathsf{Sign}(sk, m) \to \sigma\)，\(\mathsf{Verify}(pk, m, \sigma) \to \{0, 1\}\)。安全性（EUF-CMA
/
存在性不可伪造性）：即使攻击者可以自适应地查询任意消息的签名，也不能为新消息伪造有效签名。

\textbf{定理 261.1}（RSA-FDH 签名的安全性 / Full Domain
Hash）：在随机谕示模型下，RSA-FDH 签名是 EUF-CMA 安全的当且仅当 RSA
问题是难的（Bellare-Rogaway 1993/1996）。

\subsubsection{261.2 零知识证明}\label{ux96f6ux77e5ux8bc6ux8bc1ux660e}

\textbf{定义 261.2}（交互证明系统 / Interactive Proof）：证明者 \(P\)
和验证者 \(V\) 之间的协议，使得 \(P\) 可以令 \(V\) 确信陈述 \(x \in L\)
为真而不泄露其他信息。满足： - \textbf{完备性}（completeness）：若
\(x \in L\)，诚实 \(P\) 总是成功说服 \(V\) -
\textbf{可靠性}（soundness）：若
\(x \notin L\)，任何（甚至恶意的）\(P^*\) 说服 \(V\) 的概率可忽略 -
\textbf{零知识性}（zero-knowledge）：\(V\) 在交互中学到的不超过
\(x \in L\) 这个事实本身（存在模拟器可生成与真实交互不可区分的 view）

\textbf{定义 261.3}（\(\Sigma\)-协议 /
诚实验证者零知识）：三步协议（承诺 \(a\)、挑战 \(c\)、响应
\(z\)），满足特殊的可靠性（从两个接受相同 \(a\) 但不同 \(c\) 的
transcripts
可提取证据）。\(\Sigma\)-协议可转化为非交互零知识证明（Fiat-Shamir
变换）。

\textbf{定理 261.2}（Schnorr 的识别/签名协议，1989）：基于离散对数的
\(\Sigma\)-协议。\(P\) 知道 \(x\) 使
\(h = g^x\)。承诺：\(a = g^r\)（随机
\(r\)）。挑战：\(c\)。响应：\(z = r + cx\)。验证：\(g^z = a \cdot h^c\)。Fiat-Shamir
变换：\(c = H(a, m)\) 将交互证明转为签名。

\subsubsection{261.3
安全多方计算}\label{ux5b89ux5168ux591aux65b9ux8ba1ux7b97}

\textbf{定义 261.4}（安全多方计算 / Secure MPC）：\(n\)
个参与方各自有输入 \(x_i\)，联合计算 \(f(x_1, \ldots, x_n)\)
而不暴露各自的输入（除计算结果的隐含信息外）。

\textbf{定义 261.5}（秘密共享与阈值密码）： - \textbf{Shamir
秘密共享}（1979）：\(t\) 个点确定 \(t-1\) 次多项式。将秘密 \(s\) 编码为
\(f(0)=s\)，分发给 \(n\) 方 \(f(i)\)。任意 \(t\) 方可通过 Lagrange
插值恢复 \(s\)，\(t-1\) 方无法获得任何信息 -
\textbf{阈值密码}：解密/签名需要 \(\ge t\) 个参与方合作

\textbf{定理 261.3}（Yao 的乱码电路 / Garbled
Circuit，1986）：任何多项式时间可计算函数都可以被两方安全地计算（在半诚实模型下）。这是安全两方计算的基础理论。

\bigskip\hrule\bigskip

\subsection{第262章（续）：量子密码学与信息论安全}\label{ux7b2c262ux7ae0ux7eedux91cfux5b50ux5bc6ux7801ux5b66ux4e0eux4fe1ux606fux8bbaux5b89ux5168}

\subsubsection{262.1
量子密钥分发（QKD）}\label{ux91cfux5b50ux5bc6ux94a5ux5206ux53d1qkd}

量子密钥分发利用量子力学的基本原理------不确定性原理和不可克隆定理------在两个远程方之间建立信息论安全的共享密钥。

\textbf{定义 262.1}（BB84 协议 / Bennett-Brassard, 1984）：Alice
在两组建共轭基底中随机选择------计算基底 \(\{|0\rangle, |1\rangle\}\) 和
Hadamard 基底
\(\{|+\rangle = \frac{|0\rangle+|1\rangle}{\sqrt{2}}, |-\rangle = \frac{|0\rangle-|1\rangle}{\sqrt{2}}\}\)------制备单个光子态并通过量子信道发送给
Bob。Bob
随机选择测量基底。公开比对基底选择后，保留基底一致的部分作为原始密钥（sifted
key）。

\textbf{协议步骤}： 1. Alice 随机生成长度 \(\gg n\)
的随机比特串和随机基底串，制备对应量子态并发送 2. Bob
随机选择测量基底，测量每个接收到的光子并记录结果 3.
双方通过经典认证信道公开比对基底选择，丢弃基底不一致的比特 4.
随机抽样部分比特公开比对，估计量子比特错误率（QBER）以检测窃听 5. 若
QBER 低于阈值，执行信息调和（纠错）和隐私放大（通过 universal
hashing）提取最终安全密钥

\textbf{定理 262.1}（QKD 的无条件安全性 / Shor-Preskill, 2000）：BB84
协议在任意攻击（包括联合攻击和相干攻击）下的安全性可归约为
CSS（Calderbank-Shor-Steane）量子纠错码的不可区分性。窃听者 Eve
获得的信息量以 Holevo 界显式约束：

\[I(A;B) \leq S(\rho_A) = -\operatorname{Tr}(\rho_A \log \rho_A)\]

其中 \(S(\rho_A)\) 是 Alice 子系统约化密度矩阵的 von Neumann
熵。当观测到的 QBER 小于约 \(11\%\) 时，可提取正的密钥率。

\textbf{定义 262.2}（Holevo 界 / Holevo, 1973）：Bob
可从量子态中提取的关于 Alice 经典比特的信息量不超过 Alice 系统的 von
Neumann 熵。这是量子信息论中量化窃听者优势的核心工具。

\subsubsection{262.2
格密码与后量子安全}\label{ux683cux5bc6ux7801ux4e0eux540eux91cfux5b50ux5b89ux5168}

\textbf{定义 262.3}（格与格问题）：\(n\) 维\textbf{格}
\(\Lambda = \{\sum_{i=1}^n z_i \mathbf{b}_i : z_i \in \mathbb{Z}\}\)（\(\mathbf{b}_1, \ldots, \mathbf{b}_n\)
为格基）。核心困难问题： -
\textbf{SVP}（最短向量问题）：求格中长度最短的非零向量 -
\textbf{CVP}（最近向量问题）：给定目标点，求最近的格点 -
\textbf{SIVP}（最短独立向量问题）：求 \(n\) 个线性无关的短格向量

\textbf{定义 262.4}（LWE 问题 / Learning With Errors, Regev,
2005）：给定素数模数 \(q\)、维度 \(n\)、秘密向量
\(\mathbf{s} \in \mathbb{Z}_q^n\) 和误差分布 \(\chi\)，LWE 分布为

\[(\mathbf{a}, b = \langle \mathbf{a}, \mathbf{s} \rangle + e \bmod q), \quad \mathbf{a} \leftarrow \mathbb{Z}_q^n \text{ 均匀}, \quad e \leftarrow \chi\]

搜索型 LWE 问题：给定来自该分布的多项式数量样本，恢复秘密
\(\mathbf{s}\)。判定型 LWE 问题：区分 LWE 分布与均匀分布。

\textbf{定理 262.2}（Regev
量子归约，2005）：若存在高效算法以不可忽略的概率求解 LWE
问题，则存在量子算法求解最坏情况下的格上近似最短向量问题（GapSVP）和最短独立向量问题（SIVP）。这建立了
LWE 的密码学安全性------假定格上最坏情况问题对量子计算机难解，则 LWE
在平均情况下也是难解的。

\textbf{后量子密码标准化}：NIST 于 2016-2024
年的后量子密码标准化过程选定了基于格的方案： -
\textbf{KYBER}（ML-KEM）：基于 Module-LWE 的密钥封装机制 -
\textbf{Dilithium}（ML-DSA）：基于 Module-LWE 和 Module-SIS
的数字签名方案

这些方案的数学基础正是格上困难问题的（量子）归约安全性。

\subsubsection{262.3 信息论安全}\label{ux4fe1ux606fux8bbaux5b89ux5168}

\textbf{定义 262.5}（信息论安全 /
无条件安全）：加密方案是\textbf{信息论安全}的，如果其安全性不依赖于任何计算假设------即使敌手拥有无限计算能力，密文也不泄露任何关于明文的信息（除密文长度外）。形式上，\(H(M \mid C) = H(M)\)（完全保密，Shannon,
1949）。

\textbf{定理 262.3}（Shannon 完全保密定理）：对 \(n\)
比特消息，一次一密（One-Time Pad）\(c = m \oplus k\)（\(k\)
为真随机且仅用一次的 \(n\)
比特密钥）实现完全保密。反之，任何完全保密的加密方案必须满足
\(H(K) \geq H(M)\)（密钥熵不小于消息熵）。

\textbf{定义 262.6}（提取器与隐私放大 / Extractor and Privacy
Amplification）：\textbf{随机提取器}
\(\operatorname{Ext}: \{0,1\}^n \times \{0,1\}^d \to \{0,1\}^m\)
将任意具有最小熵 \(k\) 的弱随机源 \(X\) 与短的均匀种子 \(S\)
结合，输出统计上接近均匀的 \(m\) 比特。在 QKD 的隐私放大步骤中，使用
universal hashing 族（如 Toeplitz
矩阵）从部分安全的原始密钥中提取近乎完全安全的最终密钥。

\bigskip\hrule\bigskip

\emph{本卷在数论（V17）、抽象代数（V8、V13）和计算理论（V31）的基础上，系统建立了编码理论深化、对称密码、公钥密码、格密码与后量子密码以及密码协议的数学理论。共
5 章。}

\newpage

\section{卷四十七：量子信息数学}\label{ux5377ux56dbux5341ux4e03ux91cfux5b50ux4fe1ux606fux6570ux5b66}

\begin{quote}
量子信息科学是量子力学与信息论和计算机科学的交叉领域，其数学基础涉及
Hilbert
空间、算子理论、张量积和群表示论。量子纠缠是经典世界不存在的资源，量子算法（Shor、Grover）在计算复杂性上超越了经典算法。本卷在量子力学数学基础（V30
Ch 149）和函数分析（V14 Ch 73-Hilbert
空间）的基础上，系统建立量子力学信息论公理、纠缠理论、量子信道、量子算法和量子纠错码的数学理论。
\end{quote}

\bigskip\hrule\bigskip

\subsection{第262章：量子力学的信息论公理}\label{ux7b2c262ux7ae0ux91cfux5b50ux529bux5b66ux7684ux4fe1ux606fux8bbaux516cux7406}

量子力学可以用信息论的语言重新表述，将量子态视为信息载体。

\subsubsection{262.1
量子态与密度矩阵}\label{ux91cfux5b50ux6001ux4e0eux5bc6ux5ea6ux77e9ux9635}

\textbf{定义 262.1}（量子态 / Quantum
State）：物理系统的量子态由单位向量
\(|\psi\rangle \in \mathcal{H}\)（Hilbert
空间）表示（模相位：\(|\psi\rangle\) 和 \(e^{i\theta}|\psi\rangle\)
表示同一物理态）。\(\dim \mathcal{H} = d\)。

\textbf{定义 262.2}（密度矩阵 / Density Matrix）：\textbf{密度算子}
\(\rho\) 是 \(\mathcal{H}\) 上的正半定迹为 1 的算子：
\[\rho \geq 0, \quad \operatorname{Tr}(\rho) = 1\]

纯态：\(\rho = |\psi\rangle\langle\psi|\)；混合态：\(\rho = \sum_i p_i |\psi_i\rangle\langle\psi_i|\)（\(p_i > 0\)，\(\sum p_i = 1\)）。

\textbf{定义 262.3}（复合系统与张量积）：两个系统 \(\mathcal{H}_A\) 和
\(\mathcal{H}_B\) 的复合系统为
\(\mathcal{H}_A \otimes \mathcal{H}_B\)。子系统 \(\rho_A\)
的\textbf{约化密度矩阵}为
\(\rho_A = \operatorname{Tr}_B(\rho_{AB})\)（部分迹）。

\subsubsection{262.2 量子测量}\label{ux91cfux5b50ux6d4bux91cf}

\textbf{定义 262.4}（投影测量 / von Neumann 测量）：由可观测量的谱分解
\(M = \sum_m m P_m\) 定义。测量结果 \(m\) 的概率为
\(\operatorname{Tr}(\rho P_m)\)，测量后状态退相为
\(P_m \rho P_m / \operatorname{Tr}(\rho P_m)\)。

\textbf{定义 262.5}（POVM / 正算子值测度）：\textbf{POVM} 是一组正算子
\(\{E_m\}\) 满足 \(\sum_m E_m = I\)。结果 \(m\) 的概率为
\(\operatorname{Tr}(\rho E_m)\)。POVM
推广了投影测量，适用于非理想测量和信息最优获取。

\textbf{定理 262.1}（Naimark 扩张定理）：任何 POVM 可由更大 Hilbert
空间上的投影测量实现。即 POVM 测量等价于扩展系统后的投影测量。

\subsubsection{262.3 量子比特与 Bloch
球}\label{ux91cfux5b50ux6bd4ux7279ux4e0e-bloch-ux7403}

\textbf{定义 262.6}（量子比特 /
Qubit）：\(\mathcal{H} = \mathbb{C}^2\)。纯态
\(|\psi\rangle = \alpha|0\rangle + \beta|1\rangle\)（\(|\alpha|^2 + |\beta|^2 = 1\)）。任意单量子比特密度矩阵可写为

\[\rho = \frac{1}{2}(I + \mathbf{r} \cdot \boldsymbol{\sigma})\]

其中 \(\mathbf{r} \in \mathbb{R}^3\)（Bloch 向量）满足
\(|\mathbf{r}| \leq 1\)，\(\boldsymbol{\sigma} = (\sigma_x, \sigma_y, \sigma_z)\)
为 Pauli 矩阵。纯态对应 \(|\mathbf{r}| = 1\)（Bloch 球面）。

\subsubsection{262.4 量子隐形传态（Quantum
Teleportation）}\label{ux91cfux5b50ux9690ux5f62ux4f20ux6001quantum-teleportation}

量子隐形传态（Bennett et
al.~1993）利用纠缠资源在两个远程地点之间传输未知量子态，而无需实际传输物理量子比特。

\textbf{协议步骤}：Alice 欲向 Bob 传输未知态
\(|\psi\rangle = \alpha|0\rangle + \beta|1\rangle\)。双方预先共享 EPR 对
\(|\Phi^+\rangle = \frac{1}{\sqrt{2}}(|00\rangle + |11\rangle)\)。总初态为：

\[|\psi\rangle_A \otimes |\Phi^+\rangle_{AB} = (\alpha|0\rangle + \beta|1\rangle)_A \otimes \frac{1}{\sqrt{2}}(|00\rangle + |11\rangle)_{AB}\]

\textbf{步骤 1}：Alice 对持有的两个量子比特施加 CNOT
门（控制位为其未知态）和 Hadamard 门，将总态重写为 Bell 基的展开：

\[\frac{1}{2}\sum_{k=0}^{3} |\text{Bell}_k\rangle_A \otimes \sigma_k (\alpha|0\rangle + \beta|1\rangle)_B\]

其中
\(\sigma_0 = I\)、\(\sigma_1 = X\)、\(\sigma_2 = Z\)、\(\sigma_3 = XZ\)
为 Pauli 算子。

\textbf{步骤 2}：Alice 对自身两量子比特做 Bell 基测量，得到经典结果
\(k \in \{00, 01, 10, 11\}\)（2 比特经典信息）。测量后 Bob
的量子比特塌缩为 \(\sigma_k|\psi\rangle\)。

\textbf{步骤 3}：Alice 通过经典信道将 \(k\) 发送给 Bob。Bob 根据 \(k\)
施加相应的 Pauli 校正门 \(\sigma_k^\dagger\)，恢复出 \(|\psi\rangle\)。

\textbf{注}：隐形传态不违反不可克隆定理------Alice 手中的原始态在 Bell
测量中被破坏（退相），信息通过经典信道传递（不超过光速），且协议消耗了一个
EPR 对（纠缠资源）。

\bigskip\hrule\bigskip

\subsection{第263章：量子纠缠}\label{ux7b2c263ux7ae0ux91cfux5b50ux7ea0ux7f20}

量子纠缠是多个子系统之间无法用经典概率解释的量子关联，是量子信息区别于经典信息的本质特征。

\subsubsection{263.1
纠缠的定义与判定}\label{ux7ea0ux7f20ux7684ux5b9aux4e49ux4e0eux5224ux5b9a}

\textbf{定义 263.1}（可分态与纠缠态）：两体态 \(\rho_{AB}\)
是\textbf{可分的}（separable），如果可写为直积态的凸组合：
\[\rho_{AB} = \sum_i p_i \rho_A^{(i)} \otimes \rho_B^{(i)}\]

否则 \(\rho_{AB}\)
是\textbf{纠缠的}。可分态对应经典关联，纠缠态是量子独有的。

\textbf{定理 263.1}（PPT 判据 / Peres-Horodecki 判据，1996）：若
\(\rho_{AB}\) 可分，则其部分转置 \(\rho_{AB}^{T_B}\)（仅对子系统 \(B\)
取转置）是正半定的。对 \(2 \otimes 2\) 和 \(2 \otimes 3\) 系统，PPT
条件是纠缠的必要充分条件。

\textbf{定义 263.2}（纠缠目击 / Entanglement Witness）：Hermitian 算子
\(W\) 满足
\(\operatorname{Tr}(W \rho_{\text{sep}}) \geq 0\)（对所有可分态），且存在纠缠态
\(\rho_{\text{ent}}\) 使
\(\operatorname{Tr}(W \rho_{\text{ent}}) < 0\)。\(W\)「目击」了
\(\rho_{\text{ent}}\) 的纠缠性。

\subsubsection{263.2 纠缠度量}\label{ux7ea0ux7f20ux5ea6ux91cf}

\textbf{定义 263.3}（纠缠熵 / Von Neumann 熵）：对纯态
\(|\psi\rangle_{AB}\)，\textbf{纠缠熵}
\(E(|\psi\rangle) = S(\rho_A) = S(\rho_B)\)，其中 von Neumann 熵
\(S(\rho) = -\operatorname{Tr}(\rho \log \rho)\)。纠缠熵衡量了子系统混合的程度。

\textbf{定义 263.4}（形成纠缠 / Entanglement of Formation）：两体混合态
\(\rho_{AB}\) 的\textbf{形成纠缠}为
\[E_F(\rho) = \inf \left\{ \sum_i p_i S(\rho_A^{(i)}) : \rho = \sum_i p_i |\psi_i\rangle\langle\psi_i| \right\}\]

即分解为纯态的最小平均纠缠熵（为 \(2 \otimes 2\)
状态可解析计算，Wootters 1998）。

\textbf{定理 263.2}（纠缠的不可增性 / LOCC
单调性）：纠缠在\textbf{局部操作与经典通信}（LOCC）下不会增加。任何合理的纠缠度量必须是
LOCC 单调的。这是纠缠作为资源的基本性质。

\subsubsection{263.3 Bell
不等式与非局域性}\label{bell-ux4e0dux7b49ux5f0fux4e0eux975eux5c40ux57dfux6027}

\textbf{定义 263.5}（Bell-CHSH 不等式，Clauser-Horne-Shimony-Holt
1969）：对二元测量 \(A, A'\)（Alice 端）和 \(B, B'\)（Bob
端），任何经典局域隐变量理论满足
\[|\langle AB \rangle + \langle AB' \rangle + \langle A'B \rangle - \langle A'B' \rangle| \leq 2\]

量子纠缠态可违反此界，最大值可达 \(2\sqrt{2}\)（Tsirelson 界）。

\textbf{定理 263.3}（Tsirelson 界，1980）：量子力学中 CHSH
表达式的最大值为 \(2\sqrt{2}\)。这严格小于无信号条件下代数学的最大值
\(4\)（Popescu-Rohrlich
盒），表明量子力学遵循的「非局域性」比一般无信号理论弱。

\bigskip\hrule\bigskip

\subsection{第264章：量子信道与量子容量}\label{ux7b2c264ux7ae0ux91cfux5b50ux4fe1ux9053ux4e0eux91cfux5b50ux5bb9ux91cf}

量子信道描述了量子信息在噪声环境中的传输。

\subsubsection{264.1 量子信道}\label{ux91cfux5b50ux4fe1ux9053}

\textbf{定义 264.1}（量子信道 / 完全正保迹映射 / CPTP Map）：量子信道
\(\mathcal{E}: \mathcal{B}(\mathcal{H}_A) \to \mathcal{B}(\mathcal{H}_B)\)
是线性映射，满足： -
\textbf{保迹}（trace-preserving）：\(\operatorname{Tr}(\mathcal{E}(\rho)) = \operatorname{Tr}(\rho)\)
- \textbf{完全正}（completely positive / CP）：对任意扩展系统
\(R\)，\(\mathcal{E} \otimes \operatorname{id}_R\) 保持正定性

\textbf{定理 264.1}（Kraus 表示 / 算子和表示）：任何 CPTP 映射
\(\mathcal{E}\) 可表示为
\[\mathcal{E}(\rho) = \sum_{k=1}^K A_k \rho A_k^\dagger\]

其中 Kraus 算子 \(A_k\) 满足 \(\sum_k A_k^\dagger A_k = I\)。

\textbf{定义 264.2}（Stinespring 膨胀定理）：任何 CPTP 映射
\(\mathcal{E}\) 可用更大大 Hilbert 空间中的酉演化 \(U\) 和环境初态
\(|0\rangle_E\)
实现：\(\mathcal{E}(\rho) = \operatorname{Tr}_E(U (\rho \otimes |0\rangle\langle 0|_E) U^\dagger)\)。这是量子信道的纯化表示。

\subsubsection{264.2
典型量子信道}\label{ux5178ux578bux91cfux5b50ux4fe1ux9053}

\textbf{定义 264.3}（常见量子噪声信道）： -
\textbf{去极化信道}（depolarizing）：\(\mathcal{E}(\rho) = (1-p)\rho + \frac{p}{d} I\)
- \textbf{幅值阻尼信道}（amplitude
damping）：模拟自发辐射（\(|1\rangle \to |0\rangle\)） -
\textbf{相位阻尼信道}（dephasing）：消灭非对角元，模拟测量导致的信息丢失：\(\mathcal{E}(\rho) = (1-p)\rho + p \operatorname{diag}(\rho)\)

\subsubsection{264.3
量子信道容量}\label{ux91cfux5b50ux4fe1ux9053ux5bb9ux91cf}

\textbf{定义 264.4}（量子信道容量的分类）： - \textbf{经典容量}
\(C(\mathcal{E})\)：通过信道可靠传输经典信息的最大速率 -
\textbf{量子容量} \(Q(\mathcal{E})\)：通过信道可靠传输量子信息的最大速率
- \textbf{私密容量} \(P(\mathcal{E})\)：信道可用于秘密共享的容量 -
\textbf{纠缠辅助容量}
\(C_E(\mathcal{E})\)：收发双方共享无限纠缠时传输经典信息的容量

\textbf{定理 264.2}（Holevo-Schumacher-Westmoreland 定理 / HSW
定理，1996-1998）：量子信道 \(\mathcal{E}\) 的经典容量为
\[C(\mathcal{E}) = \lim_{n \to \infty} \frac{1}{n} \chi(\mathcal{E}^{\otimes n})\]

其中 \textbf{Holevo 量}
\(\chi(\mathcal{E}) = \max_{\{p_i, \rho_i\}} [S(\mathcal{E}(\sum_i p_i \rho_i)) - \sum_i p_i S(\mathcal{E}(\rho_i))]\)。

\textbf{注 264.1}（Holevo 界的简要推导）：对于发送方以概率 \(p_i\)
制备状态 \(\rho_i\)、接收方通过 POVM \(\{E_m\}\)
测量的经典通信场景，设输入随机变量为 \(X\)，测量结果为 \(Y\)。经典互信息
\(I(X:Y)\) 受限于 Holevo 界：
\[\chi = S\left(\sum_i p_i \rho_i\right) - \sum_i p_i S(\rho_i)\]
推导关键：由量子相对熵的单调性（Lindblad-Uhlmann 定理），对任意 POVM
测量，经典互信息不超过 Holevo 量；再由 HSW 定理，该上界在
\(n \to \infty\) 的意义下可由适当的编码方案渐近达到。Holevo 界的本质是
\(S(\rho) - \sum_i p_i S(\rho_i)\)
度量了''量子态集合的可区分性''------它是量子信息论中经典容量的单次使用上界，比经典
Shannon 信道容量多了 von Neumann
熵的凹性修正项。当输入为纯态且相互正交时，Holevo 量达到最大值
\(\log d\)。∎

\textbf{定理 264.3}（Lloyd-Shor-Devetak 定理 / LSD
定理）：量子信道的量子容量为相干信息的正则化极限：
\[Q(\mathcal{E}) = \lim_{n \to \infty} \frac{1}{n} I_c(\mathcal{E}^{\otimes n})\]

其中\textbf{相干信息}
\(I_c(\mathcal{E}) = \max_{\rho} [S(\mathcal{E}(\rho)) - S(\mathcal{E} \otimes \operatorname{id}_R(|\psi\rangle\langle\psi|_{\rho R}))]\)。

\bigskip\hrule\bigskip

\subsection{第265章：量子算法}\label{ux7b2c265ux7ae0ux91cfux5b50ux7b97ux6cd5}

量子算法利用叠加、干涉和纠缠实现经典计算机难以完成的计算。

\subsubsection{265.1
量子计算模型}\label{ux91cfux5b50ux8ba1ux7b97ux6a21ux578b}

\textbf{定义 265.1}（量子电路模型）：\(n\) 个量子比特初始化为
\(|0^n\rangle\)。应用酉门序列（\(H\)、\(T\)、CNOT
等通用门集合），最后测量得到结果。量子电路是经典电路对应的量子推广。

\textbf{定义 265.2}（通用量子门集合 / quantum gate universality）：
\[\{H, T, \operatorname{CNOT}\}\]

（\(H\) 为 Hadamard 门，\(T = \operatorname{diag}(1, e^{i\pi/4})\)，CNOT
为受控非门）是近似通用的------任何 \(n\) 量子比特酉变换可被
\(O(4^n \log(1/\varepsilon))\) 个门 \(\varepsilon\)-逼近（Solovay-Kitaev
定理）。

\textbf{定义 265.3}（量子 Fourier 变换 / QFT）：\(n\) 量子比特的量子
Fourier 变换 \(\operatorname{QFT}_N\)（\(N=2^n\)）：
\[|j\rangle \mapsto \frac{1}{\sqrt{N}} \sum_{k=0}^{N-1} e^{2\pi i jk / N} |k\rangle\]

QFT 可在 \(O(n^2)\) 个门内实现（经典 FFT 需 \(O(n 2^n)\)）。

\subsubsection{265.2 Shor 算法与 Grover
算法}\label{shor-ux7b97ux6cd5ux4e0e-grover-ux7b97ux6cd5}

\textbf{定理 265.1}（Shor 的因式分解算法，1994）：存在量子算法在
\(O((\log N)^3)\) 时间内分解 \(N\) 位整数（以高概率）。核心步骤：用量子
Fourier 变换求模 \(N\) 指数 \(f(x) = a^x \bmod N\)
的周期（量子周期求取），然后用数论从周期推因子。Shor
算法是量子霸权的最著名例子。

\textbf{定理 265.2}

\textbf{推导概要}（Shor 算法的核心步骤）：

\begin{enumerate}
\def\labelenumi{\arabic{enumi}.}
\item
  \textbf{约化到周期求取}：随机选取
  \(a\)（\(1 < a < N\)，\(\gcd(a, N) = 1\)）。若 \(\gcd(a, N) > 1\)
  则已完成分解。考虑函数 \(f(x) = a^x \bmod N\)，其周期 \(r\) 满足
  \(a^r \equiv 1 \pmod{N}\)。若 \(r\) 为偶数且
  \(a^{r/2} \not\equiv -1 \pmod{N}\)，则 \(\gcd(a^{r/2} \pm 1, N)\) 给出
  \(N\) 的非平凡因子。
\item
  \textbf{量子周期求取}：用 \(n = 2\lceil \log_2 N \rceil\)
  个辅助量子比特制备叠加态
  \(\frac{1}{\sqrt{M}}\sum_{x=0}^{M-1} |x\rangle|0\rangle\)（\(M = 2^n\)），Oracle
  计算
  \(f(x)\)：\(\frac{1}{\sqrt{M}}\sum_x |x\rangle|f(x)\rangle\)。对辅助寄存器做
  QFT 并测量，得到
  \(\frac{1}{\sqrt{M}}\sum_x e^{2\pi i s x / M}|x\rangle\)，其中
  \(s \approx \frac{kM}{r}\)（\(k\) 为某整数）。测量结果给出
  \(s/M \approx k/r\)。
\item
  \textbf{连分式算法}：从测量值 \(s/M\) 用连分式展开求得
  \(r\)（以高概率）。复杂度：QFT 需 \(O(n^2)\) 个量子门，连分式
  \(O(n^3)\) 经典运算，总体 \(O((\log N)^3)\)。 （Grover
  的搜索算法，1996）：在 \(N\) 个未排序元素中找标记元素，Grover 算法使用
  \(O(\sqrt{N})\) 次查询（经典需要 \(O(N)\)）。算法核心是反复应用
  \textbf{Grover 迭代}
  \(G = H^{\otimes n} (2|0\rangle\langle 0| - I) H^{\otimes n} \cdot O_f\)（\(O_f\)
  为标记元素的谕示），将状态向标记方向旋转。
\end{enumerate}

\textbf{定理 265.3}（BBBV 定理 /
量子查询下界，Bennett-Bernstein-Brassard-Vazirani 1997）：量子计算机搜索
\(N\) 个元素至少需要 \(\Omega(\sqrt{N})\) 次查询。Grover 算法是最优的。

\subsubsection{265.3
量子复杂性类}\label{ux91cfux5b50ux590dux6742ux6027ux7c7b}

\textbf{定义 265.4}（量子复杂性类）： -
\textbf{BQP}：被多项式时间量子电路以有界误差判定的语言类 -
\textbf{QMA}：量子模拟 NP------多项式时间量子验证者可验证的「量子证明」

\textbf{定理 265.4}（BQP 的位置）：BPP \(\subseteq\) \textbf{BQP}
\(\subseteq\) PSPACE（Bernstein-Vazirani 1993）。已知 BQP
包含因式分解和离散对数问题，但不知道 BQP 是否包含 NP 完全问题。

\bigskip\hrule\bigskip

\subsection{第266章：量子纠错与容错量子计算}\label{ux7b2c266ux7ae0ux91cfux5b50ux7ea0ux9519ux4e0eux5bb9ux9519ux91cfux5b50ux8ba1ux7b97}

量子纠错是量子计算实用化的关键，它克服了量子态的脆弱性和退相干。

\subsubsection{266.1
量子纠错原理}\label{ux91cfux5b50ux7ea0ux9519ux539fux7406}

\textbf{定义 266.1}（量子纠错码 / Quantum Error-Correcting
Code）：\([[n, k, d]]\) 量子码将 \(k\) 个逻辑量子比特编码为 \(n\)
个物理量子比特，可纠正任意 \(\lfloor (d-1)/2 \rfloor\)
个量子比特上的错误（包括相位错误和比特翻转错误）。

\textbf{定理 266.1}（量子纠错条件 / Knill-Laflamme
条件，1997）：一组错误算子 \(\{E_a\}\) 可被码 \(C\)（投影
\(P\)）纠正当且仅当 \[P E_a^\dagger E_b P = \lambda_{ab} P\]

对所有 \(a, b\)，\(\lambda_{ab}\) 是 Hermitian 矩阵。这是经典纠错条件
\(\mathbf{1}^T H E_a^T E_b \mathbf{1} = \lambda_{ab}\) 的量子推广。

\textbf{定义 266.2}（稳定子码 / Stabilizer Code，Gottesman
1997）：\textbf{稳定子码}由 Pauli 群的 Abel 子群 \(S\)（稳定子）的公共
+1 本征空间定义。著名的 Shor 码 \([[9, 1, 3]]\) 和 CSS 码都是稳定子码。

\textbf{定理 266.2}（CSS 构造 / Calderbank-Shor-Steane
1996）：若存在经典线性码 \(C_1 = [n, k_1, d_1]\) 和
\(C_2 = [n, k_2, d_2]\) 满足 \(C_2^\perp \subseteq C_1\)，则可构造量子
CSS 码 \([[n, k_1 - k_2, \min(d_1, d_2)]]\)。

\subsubsection{266.2
表面码与容错阈值}\label{ux8868ux9762ux7801ux4e0eux5bb9ux9519ux9608ux503c}

\textbf{定义 266.3}（表面码 / Surface Code，Kitaev 1997 / Bravyi-Kitaev
1998）：在 \(L \times L\) 二维格子上定义的稳定子码。距离
\(d = L\)，每个格点上 4 体或 2
体测量。\textbf{容错阈值}：若每个物理门的错误率
\(< p_{\text{th}} \approx 1\%\)，则增加
\(L\)（扩大码距）可使逻辑错误率指数衰减。

\textbf{定理 266.3}（阈值定理 / Threshold Theorem，Aharonov-Ben-Or /
Kitaev /
Knill-Laflamme-Zurek，1996-1997）：若物理量子门的错误率低于某个常数阈值（\(\sim 10^{-3}\)
至
\(10^{-2}\)），则原则上可以进行任意长的容错量子计算。证明方法在于使用量子纠错码的分级编码（concatenation）。

\subsubsection{266.3
拓扑量子计算}\label{ux62d3ux6251ux91cfux5b50ux8ba1ux7b97}

\textbf{定义 266.4}（anyons 与拓扑量子计算）：二维系统中的
anyons（既非玻色子也非费米子）的\textbf{编织操作}（braiding）构成量子门。非
Abelian anyons（如 Majorana 零模、Ising anyons、Fibonacci
anyons）的编织是实现拓扑保护的量子计算的途径。

\textbf{定理 266.4}（Fibonacci anyons 的通用性）：Fibonacci anyons
的编织操作可近似实现任意单量子比特门和 CNOT
门，从而是通用的拓扑量子计算模型。这是通过编织群在量子比特空间上的稠密性保证的（类似于
Solovay-Kitaev 定理）。

\bigskip\hrule\bigskip

\emph{本卷在量子力学数学基础（V30）和泛函分析（V14）的基础上，系统建立了量子态与测量公理、量子纠缠理论、量子信道与容量、量子算法和量子纠错与容错计算的数学理论。共
5 章。}

\newpage

\section{卷四十八：随机矩阵理论}\label{ux5377ux56dbux5341ux516bux968fux673aux77e9ux9635ux7406ux8bba}

\begin{quote}
随机矩阵理论（Random Matrix Theory /
RMT）研究元素为随机变量的矩阵的谱性质（特征值分布）。它起源于
Wishart（1928）在多元统计中的工作和
Wigner（1955）在核物理中的应用，现已发展成连接数论（Riemann \(\zeta\)
函数零点）、高维统计、机器学习、无线通信、量子混沌和自由概率论的统一数学框架。Wigner
半圆律是「随机矩阵的普遍性」的典范------大量随机矩阵的宏观特征值分布不依赖于微观分布细节。本卷在概率论
II（V21）、线性代数（V7）和泛函分析（V14）的基础上，建立 Wigner
半圆律、自由概率论、Tracy-Widom 分布和随机矩阵现代应用的理论体系。
\end{quote}

\bigskip\hrule\bigskip

\subsection{第267章：Wigner
半圆律}\label{ux7b2c267ux7ae0wigner-ux534aux5706ux5f8b}

Wigner 半圆律是随机矩阵理论的基本定理，断言一大类对称（或
Hermitian）随机矩阵的极限特征值分布为半圆形。

\subsubsection{267.1 经典 Wigner
矩阵}\label{ux7ecfux5178-wigner-ux77e9ux9635}

\textbf{定义 267.1}（Wigner 矩阵）：\textbf{Wigner 矩阵}
\(H = (h_{ij})_{1 \leq i,j \leq N}\) 是 \(N \times N\) 随机
Hermitian（或实对称）矩阵，满足： - 对角线元素 \(h_{ii}\) 是 i.i.d.
实随机变量（均值 \(0\)，方差 \(\sigma^2\)） - 非对角线元素
\(h_{ij}\)（\(i < j\)）是 i.i.d. 复（或实）随机变量（均值 \(0\)，方差
\(1/N\)，独立于对角线元素） - \(h_{ji} = \overline{h_{ij}}\)

\textbf{定义 267.2}（经验谱分布 / ESD）：矩阵 \(H\)
的\textbf{经验谱分布}（Empirical Spectral Distribution）为经验测度
\[\mu_H = \frac{1}{N} \sum_{i=1}^N \delta_{\lambda_i}\]

其中 \(\lambda_1 \leq \cdots \leq \lambda_N\) 是 \(H\) 的特征值。

\textbf{定理 267.1}（Wigner 半圆律，1955-1958）：设 \(H\) 是 Wigner
矩阵（矩阵元素方差适当缩放）。以概率 \(1\)（几乎必然），当
\(N \to \infty\) 时，经验谱分布 \(\mu_H\) 弱收敛到半圆律：
\[d\mu_{sc}(x) = \frac{1}{2\pi} \sqrt{4 - x^2} \, \mathbf{1}_{[-2, 2]}(x) \, dx\]

半圆律的 Stieltjes 变换为
\(s(z) = \frac{-z + \sqrt{z^2 - 4}}{2}\)（\(z \in \mathbb{C}^+\)），满足方程
\(s(z) + \frac{1}{z + s(z)} = 0\)。

\emph{证明思路}（矩法）：对 \(k\) 为偶数（奇次矩为零，由对称性），计算
\[\mathbb{E}\left[\frac{1}{N} \operatorname{Tr}(H^k)\right] = \frac{1}{N} \sum_{i_1, \ldots, i_k=1}^N \mathbb{E}[h_{i_1 i_2} h_{i_2 i_3} \cdots h_{i_k i_1}]\]
将期望展开，每个项对应一个闭合的指标序列
\((i_1, i_2, \ldots, i_k, i_1)\)（可视作边标记图）。由于矩阵元素独立且均值为零，期望非零仅当每个
\(h_{ij}\)
在乘积中出现偶数次------即指标序列中的每条边出现恰好两次，形成''配对''（pairing）。在
\(N \to \infty\)
极限下，主导贡献来自不会产生交叉的\textbf{非交叉配对}（non-crossing /
planar
pairings）------因为交叉配对导致独立指标求和时每多一个交叉损失因子
\(1/N\)。\(k\) 个顶点的非交叉配对数目恰由 Catalan 数
\(C_{k/2} = \frac{1}{k/2+1}\binom{k}{k/2}\) 给出。因此
\[\lim_{N \to \infty} \mathbb{E}\left[\frac{1}{N} \operatorname{Tr}(H^k)\right] = \begin{cases} C_{k/2}, & k \text{ 为偶数} \\ 0, & k \text{ 为奇数} \end{cases}\]
Catalan 数正是半圆律矩序列。由矩方法（moment
method），矩唯一确定紧支撑概率分布，方差估计进一步证明几乎必然收敛到半圆律。∎

\textbf{定理 267.2}（半圆律的普遍性 / Universality）：Wigner
半圆律对矩阵元素的分布不敏感------只要元素满足有限的矩条件（如四阶矩有限），极限谱分布就是半圆的。这是随机矩阵理论「普遍性」（universality）的第一个表现。

\subsubsection{267.2
矩法与自由独立}\label{ux77e9ux6cd5ux4e0eux81eaux7531ux72ecux7acb}

\textbf{定义 267.3}（Catalan 数与半圆律的矩）：半圆律的 \(2k\)
阶矩（奇阶矩为零）为\textbf{Catalan 数} \(C_k\)：
\[\int x^{2k} d\mu_{sc}(x) = C_k = \frac{1}{k+1} \binom{2k}{k}\]

\(k=1\) 时 \(C_1=1\)，\(k=2\) 时 \(C_2=2\)，\(k=3\) 时
\(C_3=5\)。Catalan 数计数了「非交叉配对」的数量。

\textbf{定理 267.3}（矩法的核心估计）：在 Wigner
矩阵中，\(\mathbb{E}[\operatorname{Tr}(H^k)]/N\)
的主导贡献来自矩阵指标序列的「非交叉配对」（non-crossing / planar
pairings），交叉配对的贡献被 \(1/N^2\) 压制。

\subsubsection{267.3 Wishart 矩阵与 Marchenko-Pastur
律}\label{wishart-ux77e9ux9635ux4e0e-marchenko-pastur-ux5f8b}

\textbf{定义 267.4}（Wishart / 样本协方差矩阵）：\(X\) 是 \(p \times n\)
矩阵，元素为 i.i.d.（均值 \(0\)，方差 \(1\)）。\textbf{样本协方差矩阵}为
\(S = \frac{1}{n} X X^T\)（\(p \times p\)）。

\textbf{定理 267.4}（Marchenko-Pastur 律，1967）：设 \(p, n \to \infty\)
且 \(p/n \to \gamma \in (0, \infty)\)。\(S\) 的经验谱分布收敛到
\textbf{Marchenko-Pastur 分布}：
\[d\mu_{MP}(x) = \frac{\sqrt{(x - a_-)(a_+ - x)}}{2\pi \gamma x} \mathbf{1}_{[a_-, a_+]}(x) dx + (1 - \gamma^{-1})_+ \delta_0(dx)\]

其中 \(a_\pm = (1 \pm \sqrt{\gamma})^2\)。

\bigskip\hrule\bigskip

\subsection{第268章：自由概率论}\label{ux7b2c268ux7ae0ux81eaux7531ux6982ux7387ux8bba}

自由概率论（Free Probability Theory）由
Voiculescu（1985-1990）创立，为随机矩阵提供了独立的非交换概率框架。

\subsubsection{268.1
非交换概率空间}\label{ux975eux4ea4ux6362ux6982ux7387ux7a7aux95f4}

\textbf{定义 268.1}（非交换概率空间 /
\(C^*\)-概率空间）：非交换概率空间是 \((\mathcal{A}, \tau)\)，其中
\(\mathcal{A}\) 是单位
\(C^*\)-代数，\(\tau: \mathcal{A} \to \mathbb{C}\)
是单位保迹的线性泛函，满足 \(\tau(1) = 1\) 且
\(\tau(a^*a) \geq 0\)（正性）。

\textbf{定义 268.2}（自由独立 / Free Independence）：\(\mathcal{A}\)
的子代数族 \(\{\mathcal{A}_i\}_{i \in I}\)
是\textbf{自由独立的}，如果对任意满足 \(\tau(a_k) = 0\) 的非交换多项式
\(a_1 \in \mathcal{A}_{i_1}, \ldots, a_k \in \mathcal{A}_{i_k}\)（相邻指标不同
\(i_1 \neq i_2 \neq \cdots \neq i_k\)），有
\(\tau(a_1 a_2 \cdots a_k) = 0\)。

\textbf{定理 268.1}（自由独立 = 独立矩阵的极限行为，Voiculescu
1991）：若 \(A_N\) 和 \(B_N\) 是独立（经典意义上的独立）的 Wigner
矩阵序列，则它们在 \(N \to \infty\)
极限下是自由独立的（在非交换概率空间中）。

\subsubsection{268.2 自由卷积与
R-变换}\label{ux81eaux7531ux5377ux79efux4e0e-r-ux53d8ux6362}

\textbf{定义 268.3}（自由加法卷积 / Free Additive Convolution
\textbf{\(\boxplus\)}）：若 \(a, b\) 自由独立，其边缘分布的加法卷积
\(\mu_a \boxplus \mu_b\) 定义为 \(a+b\) 的分布。

\textbf{定义 268.4}（R-变换 / \(R\)-transform，Voiculescu
1986）：对概率测度 \(\mu\)，定义其 Cauchy 变换
\(G_\mu(z) = \int \frac{d\mu(x)}{z - x}\)（\(z \in \mathbb{C}^+\)）。\(K_\mu(z) = G_\mu^{-1}(z)\)（函数逆）。\textbf{R-变换}为
\[R_\mu(z) = K_\mu(z) - \frac{1}{z}\]

\textbf{定理 268.2}（R-变换的线性化性质）：若
\(\mu = \mu_1 \boxplus \mu_2\)，则
\[R_\mu(z) = R_{\mu_1}(z) + R_{\mu_2}(z)\]

即\textbf{自由卷积对应
R-变换的加法}------这是自由概率论的核心结果，类比经典卷积对应特征函数的乘积。对数
Fourier 变换的角色在此被 R-变换取代。

\textbf{例 268.1}（半圆律的 R-变换）：半圆律 \(R(z) = z\)。因此 \(n\)
个自由独立的半圆律的 \(\boxplus\) 卷积仍为半圆律（方差放大 \(n\)
倍）------这解释了为什么 i.i.d. Wigner 矩阵之和的特征值分布可解析计算。

\subsubsection{268.3 自由乘法卷积与
S-变换}\label{ux81eaux7531ux4e58ux6cd5ux5377ux79efux4e0e-s-ux53d8ux6362}

\textbf{定义 268.5}（自由乘法卷积 / Free Multiplicative Convolution
\textbf{\(\boxtimes\)}）：对正算子 \(a, b\)
自由独立，\(\mu_a \boxtimes \mu_b\) 定义为
\(a^{1/2} b a^{1/2}\)（或等价的 \(ab\)）的分布。

\textbf{定义 268.6}（S-变换 / S-transform）：对具有非零均值的正测度
\(\mu\)，\(S_\mu(z)\) 满足 \(\chi_\mu(z S_\mu(z)) = z\)，其中
\(\chi_\mu\) 是矩母函数的相关函数。自由乘法卷积对应
S-变换的乘法：\(S_{\mu_1 \boxtimes \mu_2}(z) = S_{\mu_1}(z) S_{\mu_2}(z)\)。

\subsubsection{268.4 自由熵与自由 Fisher
信息}\label{ux81eaux7531ux71b5ux4e0eux81eaux7531-fisher-ux4fe1ux606f}

\textbf{定义 268.7}（自由熵 / Free Entropy，Voiculescu
1993-1998）：非交换随机变量 \(X\) 的\textbf{自由熵} \(\chi(X)\)
类比于经典 Shannon
微分熵，定义为近似矩阵模型的「微观态计数」的大偏差极限。自由熵满足自由信息论中的最大化原理：自由半圆变量最大化自由熵（给定方差约束）。

\bigskip\hrule\bigskip

\subsection{第269章：Tracy-Widom
分布与边缘统计}\label{ux7b2c269ux7ae0tracy-widom-ux5206ux5e03ux4e0eux8fb9ux7f18ux7edfux8ba1}

Tracy-Widom
分布是最大的随机矩阵特征值在适当尺度缩放下的极限分布，是随机矩阵理论中最重要的「新分布」之一。

\subsubsection{269.1 Gaussian 系综}\label{gaussian-ux7cfbux7efc}

\textbf{定义 269.1}（Gaussian 系综 / GOE、GUE、GSE）： -
\textbf{GOE}（Gaussian Orthogonal Ensemble）：实对称矩阵
\(H = H^T\)，\(h_{ii} \sim \mathcal{N}(0, 2)\)，\(h_{ij} \sim \mathcal{N}(0, 1)\)（\(i < j\)）
- \textbf{GUE}（Gaussian Unitary Ensemble）：复 Hermitian 矩阵
\(H = H^\dagger\)，\(\Re h_{ij}, \Im h_{ij} \sim \mathcal{N}(0, 1/2)\) -
\textbf{GSE}（Gaussian Symplectic Ensemble）：描述自旋-轨道耦合系统

\textbf{定理 269.1}（GUE 的联合特征值分布）：GUE
矩阵的特征值联合概率密度为

\[p(\lambda_1, \ldots, \lambda_N) = \frac{1}{Z_N} \prod_{i < j} |\lambda_i - \lambda_j|^2 \prod_{i=1}^N e^{-\lambda_i^2 / 2}\]

其中 \(|\lambda_i - \lambda_j|^2\) 是 Vandermonde
行列式的平方------表示特征值之间存在「排斥」效应（level repulsion）。

\subsubsection{269.2 最大特征值与 Tracy-Widom
分布}\label{ux6700ux5927ux7279ux5f81ux503cux4e0e-tracy-widom-ux5206ux5e03}

\textbf{定理 269.2}（Tracy-Widom 分布，1994-1996）：设
\(\lambda_{\max}\) 是 GUE（或
GOE/GSE）矩阵的最大特征值。经过适当缩放（中心化和标准化）后：

\[\frac{\lambda_{\max} - 2\sqrt{N}}{N^{-1/6}} \xrightarrow{d} F_\beta\]

其中 \(F_\beta\) 是\textbf{Tracy-Widom 分布}（\(\beta = 1\) 为
GOE，\(\beta = 2\) 为 GUE，\(\beta = 4\) 为 GSE）。\(F_2\) 由 Painlevé
II 方程的解定义：
\[F_2(s) = \exp\left( -\int_s^\infty (x - s) q^2(x) dx \right)\]

其中 \(q(x)\) 满足 \(\ddot{q} = x q + 2 q^3\)（Painlevé II），边界条件
\(q(x) \sim \operatorname{Ai}(x)\)（\(x \to \infty\)）。

\textbf{定理 269.3}（Tracy-Widom 的普遍性）：最大特征值的 Tracy-Widom
极限对矩阵元素分布不敏感------只要元素的四阶矩是有限的（Soshnikov
1999，Erdos-Yau 等人的工作）。这是随机矩阵边缘普遍性（edge
universality）的最重要结果。

\subsubsection{269.3 Tracy-Widom
的惊人出现}\label{tracy-widom-ux7684ux60caux4ebaux51faux73b0}

Tracy-Widom 分布在随机矩阵理论之外也出现： -
最长递增子序列的长度分布（Baik-Deift-Johansson 1999）：Ulam 问题的解 -
KPZ 普适类（Kardar-Parisi-Zhang）：\((1+1)\) 维随机生长界面的涨落 -
ASEP（非对称简单排它过程）中的涨落（Tracy-Widom 2007-2009） -
二维渗流和随机三角剖分的极限界面形状

\bigskip\hrule\bigskip

\subsection{第270章：随机矩阵与自由概率的应用}\label{ux7b2c270ux7ae0ux968fux673aux77e9ux9635ux4e0eux81eaux7531ux6982ux7387ux7684ux5e94ux7528}

随机矩阵理论在多个学科中有深刻应用------从数论到机器学习。

\subsubsection{\texorpdfstring{270.1 数论与 Riemann \(\zeta\)
函数}{270.1 数论与 Riemann \textbackslash zeta 函数}}\label{ux6570ux8bbaux4e0e-riemann-zeta-ux51fdux6570}

\textbf{定理 270.1}（Montgomery 的 pair correlation 猜想，1973 / Odlyzko
数值验证 1987）：Riemann \(\zeta\) 函数非平凡零点的配对相关函数（pair
correlation function）与 GUE 随机矩阵特征值的配对相关函数一致：
\[R_2(\beta, x) = 1 - \left( \frac{\sin(\pi x)}{\pi x} \right)^2\]

这是 Hilbert-Pólya 猜想的证据------Riemann \(\zeta\)
函数的零点可能是某个 Hermitian 算子的特征值。

\textbf{猜想 270.2}（Keating-Snaith
猜想，2000）：\(L\)-函数的矩与随机矩阵特征多项式的矩在 \(T \to \infty\)
下一致，预测了 Riemann \(\zeta\) 函数在临界线上
\(\log|\zeta(1/2 + it)|\) 的高斯分布性质。

\subsubsection{270.2
无线通信与高维统计}\label{ux65e0ux7ebfux901aux4fe1ux4e0eux9ad8ux7ef4ux7edfux8ba1}

\textbf{定义 270.1}（MIMO 系统的信道容量，Telatar 1999 / Foschini
1996）：在 \(n_t\) 发射天线、\(n_r\) 接收天线的 MIMO 信道
\(Y = H X + N\) 中（\(H\) 为信道矩阵），当 \(H\) 的元素为 i.i.d.
Gauss，信道的遍历容量为
\[C = \mathbb{E}_H [\log \det(I + \rho H H^\dagger)]\]

其中 \(\rho\) 是信噪比。Marchenko-Pastur 律给出了当
\(n_t, n_r \to \infty\) 时信道容量近似。

\textbf{定义 270.2}（高维协方差估计 / 随机矩阵在统计中的应用）：对样本数
\(n\) ≈ 维数 \(p\) 的高维数据，Marchenko-Pastur
律描述了样本协方差矩阵特征值的分布偏差。这对主成分分析（PCA）的可靠性有直接影响------高维下的特征值估计不能用传统低维理论。

\subsubsection{270.3
量子混沌与核物理}\label{ux91cfux5b50ux6df7ux6c8cux4e0eux6838ux7269ux7406}

\textbf{定义 270.3}（Bohigas-Giannoni-Schmit
猜想，1984）：量子混沌系统的能级统计统计（nearest-neighbor spacing
distribution）遵循适当对称类（GOE/GUE/GSE）的随机矩阵特征值间距分布（Wigner
猜忌 / Wigner surmise）：
\[P_{\text{GUE}}(s) = \frac{32}{\pi^2} s^2 e^{-4s^2 / \pi}\]

（\(s\) 为间距除以平均间距）。

\subsubsection{270.4
自由概率在算子代数中的应用}\label{ux81eaux7531ux6982ux7387ux5728ux7b97ux5b50ux4ee3ux6570ux4e2dux7684ux5e94ux7528}

\textbf{定理 270.3}（Voiculescu
的自由群因子同构问题，1990s）：使用自由概率论，Voiculescu 证明了自由群
von Neumann 代数 \(L(\mathbb{F}_n)\) 与
\(L(\mathbb{F}_m)\)（\(n \neq m\)）的非同构性在自由熵框架下可被区分。这解决了
von Neumann 代数领域的一个重大未解问题。

\bigskip\hrule\bigskip

\subsection{第271章：随机矩阵的现代发展}\label{ux7b2c271ux7ae0ux968fux673aux77e9ux9635ux7684ux73b0ux4ee3ux53d1ux5c55}

随机矩阵理论近二十年的突破集中在普遍性定理的严格证明（Erdos-Schlein-Yau
等人）和机器学习中的应用。

\subsubsection{271.1
普遍性定理的严格证明}\label{ux666eux904dux6027ux5b9aux7406ux7684ux4e25ux683cux8bc1ux660e}

\textbf{定义 271.1}（Dyson Brown 运动与局部平衡）：将 Wigner 矩阵 \(H\)
的特征值 \(\lambda_i(t)\) 视为时间演化的 Dyson Brown 运动（\(\lambda_i\)
间对数排斥力 + 噪声）。关键：Dyson Brown 运动在短时间达到局部平衡（local
equilibrium），使得特征值间距的短距离统计普遍。

\textbf{定理 271.1}（Erdos-Schlein-Yau 的 Wigner-Dyson-Mehta
普遍性，2010-2013）：对任何满足适当矩条件的广义 Wigner
矩阵，其特征值的\textbf{局部统计}（local
statistics）------包括相邻间距分布和 \(k\)-点相关函数------与 GUE/GOE
的对应统计在极限下一致。即 Tracy-Widom 边缘分布和 Wigner-Dyson
体（bulk）统计具有普遍性。

\textbf{定理 271.2}（Dyson 猜想 /
边缘隙分布，2020s）：对于矩阵边缘的最大特征值，Tracy-Widom
分布是普遍的。对于次最大特征值的隙分布也有类似的普遍性结果。

\subsubsection{271.2
随机矩阵与机器学习}\label{ux968fux673aux77e9ux9635ux4e0eux673aux5668ux5b66ux4e60}

\textbf{定义 271.2}（随机神经网络 / Random Matrix Theory of Deep
Learning）：深度神经网络的权值矩阵（随机初始化）的特征值分布影响梯度传播和训练动力学。Gaussian
权值网络的 Hessian 矩阵和 Fisher
信息矩阵的谱可用随机矩阵理论分析（Pennington-Worah
2017，Pennington-Schoenholz-Ganguli 2017-2018）。

\textbf{定理 271.3}（双下降现象 / Double Descent，Belkin-Hsu-Ma-Mandal
2019 /
随机矩阵解释）：当模型参数数超过训练样本数时，测试误差出现双下降（插值阈值后再次下降）。随机矩阵理论解释了这一现象------过参数化后的最小范数解（伪逆）比恰好在插值阈值处的解具有更小的方差。

\subsubsection{271.3 稀疏随机矩阵与稀疏
Wigner}\label{ux7a00ux758fux968fux673aux77e9ux9635ux4e0eux7a00ux758f-wigner}

\textbf{定义 271.3}（稀疏 Wigner 矩阵 / Erdős-Rényi 矩阵）：Wigner
矩阵的每个元素以概率 \(p = d/N\) 非零（\(d\) 为平均度）。当
\(d \to \infty\) 但 \(d/N \to 0\)
时，谱分布偏离半圆律，边缘出现离群特征值。

\textbf{定理 271.4}（稀疏 Wigner 谱的过渡）：稀疏 Wigner 矩阵的谱在
\(d \gg 1\) 时接近半圆，在 \(d\)
有限时呈现更尖锐的峰。存在从稀疏谱到稠密半圆谱的连续过渡。

\bigskip\hrule\bigskip

\emph{本卷在概率论 II（V21）和线性代数（V7）的基础上，系统建立了 Wigner
半圆律与普遍性、自由概率论、Tracy-Widom
分布与边缘统计、随机矩阵跨学科应用和现代发展的数学理论。共 5 章。}

\newpage

\section{卷四十九：反问题}\label{ux5377ux56dbux5341ux4e5dux53cdux95eeux9898}

\begin{quote}
反问题（Inverse
Problems）研究如何从间接观测数据中推断系统的内部参数或源。与正问题（已知参数→计算观测）不同，反问题（观测→推断参数）通常是\textbf{不适定的}（ill-posed）：解的存在性、唯一性或稳定性可能不满足。Hadamard
在 1902
年提出的适定性概念是理解反问题的基本框架。反问题在现代科学和工程中有广泛的应用：CT
扫描（Radon
变换）、地震成像、电阻抗断层成像、天气预测的数据同化、图像去模糊等。本卷在泛函分析（V14）、偏微分方程（V18）和数值分析（V23）的基础上，系统建立不适应问题理论、Radon
变换与 CT、散射反问题、Bayes 反演框架和数据同化方法的数学理论。
\end{quote}

\bigskip\hrule\bigskip

\subsection{第272章：不适定问题与正则化理论}\label{ux7b2c272ux7ae0ux4e0dux9002ux5b9aux95eeux9898ux4e0eux6b63ux5219ux5316ux7406ux8bba}

不适定问题是反问题的数学核心。正则化是将不适定问题转化为适定逼近问题的系统方法。

\subsubsection{272.1 不适定问题}\label{ux4e0dux9002ux5b9aux95eeux9898}

\textbf{定义 272.1}（Hadamard 适定性，1902）：问题
\(K f = g\)（\(K: X \to Y\) 为有界线性算子，\(X, Y\) 为 Hilbert
空间）是\textbf{适定的}，如果： 1. \textbf{存在性}：对所有
\(g \in Y\)，存在解 \(f \in X\) 2. \textbf{唯一性}：解唯一 3.
\textbf{稳定性}：\(f\) 连续依赖于 \(g\)（\(K^{-1}\) 连续，等价于
\(K^{-1}\) 有界）

否则，问题是\textbf{不适定的}。

\textbf{定理 272.1}（不适定性的典型来源）：若 \(K\) 是紧算子且 \(X\)
无穷维，则 \(K\) 的值域 \(\mathcal{R}(K)\) 不闭。这时即使 \(K\)
单射（唯一性成立），\(K^{-1}\)
也无界（稳定性不成立）。大多数反问题涉及紧正算子的求逆。

\textbf{例 272.1}（向后热方程 / Backward Heat
Equation）：\(\frac{\partial u}{\partial t} = \nabla^2 u\)，给定终值
\(u(T, x)\) 求初值
\(u(0, x)\)。正问题是耗散的（平滑化），反问题放大高频噪声------这是指数级不适定的。

\textbf{定义 272.2}（不适定度 / Degree of Ill-Posedness）：对紧算子
\(K\)，设其奇异值 \(\sigma_n \searrow 0\)。若 \(\sigma_n\)
以多项式衰减（\(\sigma_n \sim n^{-\alpha}\)），为适度不适定；若
\(\sigma_n\)
指数衰减（\(\sigma_n \sim e^{-\beta n}\)），为\textbf{严重不适定}。

\subsubsection{272.2 Tikhonov 正则化}\label{tikhonov-ux6b63ux5219ux5316}

\textbf{定义 272.3}（Tikhonov 正则化，1963）：对噪声数据
\(g^\delta\)（\(\|g^\delta - g\| \leq \delta\)），Tikhonov 正则化解
\(f_\alpha^\delta\) 最小化
\[J_\alpha(f) = \|K f - g^\delta\|^2 + \alpha \|f\|^2\]

其中 \(\alpha > 0\) 是\textbf{正则化参数}。解析解（在 Hilbert 空间下）为
\[f_\alpha^\delta = (K^* K + \alpha I)^{-1} K^* g^\delta\]

\textbf{注 272.1}（Tikhonov 正则化的统计解释）：Tikhonov 正则化与 Bayes
统计推断有深刻的对应关系。在加性 Gauss 噪声模型
\(g^\delta = K f + \eta\)（\(\eta \sim \mathcal{N}(0, \sigma^2 I)\)）下，极小化泛函
\(J_\alpha(f) = \|K f - g^\delta\|^2 + \alpha\|f\|^2\) 等价于在 \(f\)
上施加 Gauss 先验
\(\pi(f) \propto \exp(-\frac{\alpha}{2\sigma^2}\|f\|^2)\)
后求\textbf{最大后验估计}（MAP）。具体地，似然函数为
\(\pi(g^\delta|f) \propto \exp(-\frac{1}{2\sigma^2}\|Kf - g^\delta\|^2)\)，由
Bayes 定理，后验分布
\(\pi(f|g^\delta) \propto \pi(g^\delta|f)\pi(f) \propto \exp(-\frac{1}{2\sigma^2}[\|Kf - g^\delta\|^2 + \alpha\|f\|^2])\)。使后验最大即令
\(\|Kf - g^\delta\|^2 + \alpha\|f\|^2\) 最小，恰为 Tikhonov
泛函。正则化参数 \(\alpha\) 在此框架下对应先验精度与噪声方差的比值
\(\alpha = \sigma^2/\sigma_f^2\)。更一般地，\(\ell^1\) 正则化对应
Laplace 先验 \(\pi(f) \propto \exp(-\alpha\|f\|_1)\)，总变差正则化对应
Markov
随机场先验。这一统计视角统一了确定性正则化和概率反演，并为正则化参数的选择（如通过经验
Bayes 方法）提供了原则性指导。

\textbf{定理 272.2}（Tikhonov 正则化的收敛率）：若
\(f^\dagger = (K^* K)^\nu w\)（\(\|w\| \leq E\)，\(\nu > 0\)
为源条件），且 \(\alpha \sim \delta^{2/(2\nu+1)}\)，则
\[\|f_\alpha^\delta - f^\dagger\| \leq C \delta^{2\nu/(2\nu+1)}\]

源条件 \(\nu\) 是解光滑性的度量------越光滑的真解恢复越快。

\textbf{定理 272.3}（Morozov 偏差原理 / Discrepancy
Principle，1966）：选择 \(\alpha\) 使
\(\|K f_\alpha^\delta - g^\delta\| = \tau \delta\)（\(\tau > 1\)），则
\(\alpha\) 随噪声水平 \(\delta\) 自适应调整，保证
\(f_\alpha^\delta \to f^\dagger\)（\(\delta \to 0\)）。

\subsubsection{272.3
迭代正则化与稀疏正则化}\label{ux8fedux4ee3ux6b63ux5219ux5316ux4e0eux7a00ux758fux6b63ux5219ux5316}

\textbf{定义 272.4}（Landweber 迭代）：梯度下降应用于
\(\|K f - g^\delta\|^2\)：
\[f_{k+1}^\delta = f_k^\delta - \omega K^*(K f_k^\delta - g^\delta)\]

其中 \(\omega \in (0, 2/\|K\|^2)\)。早期停止（迭代次数 \(k\)
随噪声减小）起到正则化作用。Landweber 法的渐进正则化等价于
\(\alpha = 1/k\)（步长的倒数）。

\textbf{定义 272.5}（\(\ell^1\)-正则化 / 压缩感知正则化）：极小化
\(J_\alpha(f) = \|K f - g^\delta\|^2 + \alpha \|f\|_1\)（\(\|f\|_1\) 为
\(\ell^1\) 范数）。\(\ell^1\)
正则化促进稀疏解（许多分量为零），在压缩感知中有广泛应用。

\textbf{定理 272.4}（\(\ell^1\)-正则化的恢复保证 / Candès-Romberg-Tao
2006）：若 \(K\) 满足有限等距性质（RIP）且真解 \(f^\dagger\)
足够稀疏，\(\ell^1\)-极小化可以在噪声水平 \(\delta\) 下以 \(O(\delta)\)
误差恢复真解。

\bigskip\hrule\bigskip

\subsection{第273章：Radon
变换与计算机断层成像}\label{ux7b2c273ux7ae0radon-ux53d8ux6362ux4e0eux8ba1ux7b97ux673aux65adux5c42ux6210ux50cf}

X 射线 CT（计算机断层成像）的数学基础是 Radon
变换------将函数映射为它在所有超平面上的线积分。

\subsubsection{273.1 Radon 变换}\label{radon-ux53d8ux6362}

\textbf{定义 273.1}（二维 Radon 变换，Radon 1917）：对
\(f \in L^1(\mathbb{R}^2) \cap L^2(\mathbb{R}^2)\)，
\[\mathcal{R}f(\theta, s) = \int_{x \cdot \theta = s} f(x) \, d\ell = \int_{-\infty}^\infty f(s\theta + t\theta^\perp) dt\]

其中
\(\theta = (\cos\theta, \sin\theta)\)，\(s \in \mathbb{R}\)。\(\mathcal{R}f(\theta, s)\)
是 \(f\) 沿方向 \(\theta\) 垂直线（距原点 \(s\)）的线积分。

\textbf{定义 273.2}（三维 Radon 变换与 X 射线变换）： - 三维中，Radon
变换在平面（不是线）上积分 - \textbf{X 射线变换}：沿所有直线积分（在锥束
CT 中使用）

\textbf{定理 273.1}（Fourier 切片定理 / Projection-Slice
Theorem）：Radon 变换在各角度 \(\theta\) 下的一维 Fourier
变换等于原函数二维 Fourier 变换沿过原点直线 \(\xi = \rho \theta\)
的切片： \[\widehat{\mathcal{R}f}(\theta, \rho) = \hat{f}(\rho \theta)\]

这是所有 CT 重建算法的理论基石------在频域中，不同角度的投影填满整个
Fourier 空间。

\subsubsection{273.2 Radon 反演}\label{radon-ux53cdux6f14}

\textbf{定理 273.2}（Radon 逆变换公式）：\(f(x)\)
可由其所有方向的投影重建：
\[f(x) = \frac{1}{4\pi} \int_0^{2\pi} \left( \mathcal{H} \frac{\partial}{\partial s} \mathcal{R}f \right)(\theta, x \cdot \theta) \, d\theta\]

其中 \(\mathcal{H}\) 是 Hilbert 变换。即先对投影数据求 \(s\) 偏导，进行
Hilbert 变换，再做反投影（backprojection）。

\textbf{定理 273.3}（滤波反投影算法 / FBP / Filtered
Backprojection）：反演公式等价于：
\[f(x) = \frac{1}{2} \int_0^{2\pi} (h * \mathcal{R}f)(\theta, x \cdot \theta) \, d\theta\]

其中 \(h\) 是 ramp 滤波器：\(\hat{h}(\rho) = |\rho|\)。FBP 是 CT
扫描中最广泛使用的重建算法。

\textbf{定理
273.4}（有限角度问题的严重不适定性）：若投影数据仅在有限角度区间
\([\theta_{\min}, \theta_{\max}]\) 内可用，则缺失的 Fourier
空间部分造成了严重的重建伪影和不适定性。COTA（Characterization Of The
Artifacts）理论描述了有限角度重建中不可恢复的奇异性。

\subsubsection{273.3
成像模型与离散化}\label{ux6210ux50cfux6a21ux578bux4e0eux79bbux6563ux5316}

\textbf{定义 273.3}（离散 Radon 变换 / X 射线投影矩阵）：将图像离散化为
\(N \times N\) 像素网格。正向投影 \(A f = p\) 矩阵 \(A_{ij}\) 表示像素
\(j\) 对射线 \(i\) 的贡献（线长）。重建 \(f = A^{-1} p\)
通常使用迭代方法（ART、SART 或 SIRT 算法）。

\textbf{定义 273.4}（Cone-beam CT 的 Feldkamp 算法 / FDK
1984）：锥束投影的近似三维 FBP 重建公式，是当代 CBCT 系统的标准算法。

\bigskip\hrule\bigskip

\subsection{第274章：散射反问题}\label{ux7b2c274ux7ae0ux6563ux5c04ux53cdux95eeux9898}

散射反问题研究如何从波（声波、电磁波）的散射数据中恢复障碍物或不均匀介质的内部性质。这是声纳、雷达、医学超声和地球物理成像的基础。

\subsubsection{274.1 Helmholtz
方程的散射正问题}\label{helmholtz-ux65b9ux7a0bux7684ux6563ux5c04ux6b63ux95eeux9898}

\textbf{定义 274.1}（Helmholtz 方程 / 声波散射）：时间谐波声波
\(p(\mathbf{x}, t) = \Re(u(\mathbf{x}) e^{-i\omega t})\) 满足 Helmholtz
方程 \[\nabla^2 u + k^2 u = 0\]

其中 \(k = \omega/c\) 为波数。入射波 \(u^i\)
遇到障碍物（或折射率非均匀介质）产生散射波 \(u^s\)。总场
\(u = u^i + u^s\)。

\textbf{定义 274.2}（Sommerfeld
辐射条件）：散射波在无穷远处满足出射波条件：
\[\lim_{r \to \infty} r^{(n-1)/2} \left( \frac{\partial u^s}{\partial r} - i k u^s \right) = 0\]

这保证了散射波从障碍物向外辐射（物理因果性）。

\textbf{定理 274.1}（散射场的远场模式 / Far Field Pattern）：在远场
\(|\mathbf{x}| \to \infty\)，散射波的渐近行为由\textbf{远场模式}
\(u_\infty(\hat{\mathbf{x}}, \hat{\mathbf{d}}, k)\)
刻画（\(\hat{\mathbf{x}}\) 为观测方向，\(\hat{\mathbf{d}}\)
为入射方向）：
\[u^s(\mathbf{x}) = \frac{e^{ik|x|}}{|x|^{(n-1)/2}} \left( u_\infty \left( \frac{\mathbf{x}}{|x|} \right) + O\left( \frac{1}{|x|} \right) \right)\]

\subsubsection{274.2
散射反问题的唯一性与不适定性}\label{ux6563ux5c04ux53cdux95eeux9898ux7684ux552fux4e00ux6027ux4e0eux4e0dux9002ux5b9aux6027}

\textbf{定理 274.2}（障碍物散射的唯一性 / Schiffer 定理）：设
\(D_1, D_2 \subset \mathbb{R}^3\)
为两个有界开集（障碍物），具有光滑边界。考虑 Helmholtz 方程的 Dirichlet
外问题：\((\nabla^2 + k^2)u = 0\)（在
\(\mathbb{R}^3 \setminus \overline{D_i}\) 中），\(u = 0\)（在
\(\partial D_i\) 上，声软边界条件），外加 Sommerfeld
辐射条件。若两个障碍物对所有入射方向 \(\hat{\mathbf{d}} \in S^2\)
和所有波数 \(k > 0\)（或固定 \(k\) 的所有入射方向）产生的远场模式
\(u_\infty(\hat{\mathbf{x}}, \hat{\mathbf{d}}, k)\) 相同，则
\(D_1 = D_2\)。即远场模式唯一确定障碍物的形状。该定理对 Neumann
边界条件（声硬障碍物）和阻抗边界条件也有相应版本。

\textbf{定理 274.3}（折射率的唯一性 / 逆介质散射）：折射率
\(n(\mathbf{x})\)
由多频率的远场模式在多入射方向下的数据唯一确定（Nachman 1988, Novikov
1988 等）。

\textbf{定义
274.3}（散射反问题的不适定性）：散射反问题是\textbf{严重非线性且严重不适定}的------远场模式对应紧算子的值域，且散射波的高频分量随频率指数衰减（避衰波
/ evanescent waves）。

\subsubsection{274.3
线性采样法与分解法}\label{ux7ebfux6027ux91c7ux6837ux6cd5ux4e0eux5206ux89e3ux6cd5}

\textbf{定义 274.4}（线性采样法 / Linear Sampling Method，Colton-Kirsch
1996）：LSM
通过求解第一类远场积分方程（不适定的线性方程）来确定障碍物的支撑，而无需事先知道障碍物的边界条件类型。通过构造\textbf{指示函数}
\(G(\mathbf{z})\)（在障碍物内大，在障碍物外小）来成像。

\textbf{定理 274.4}（分解法 / Factorization Method，Kirsch
1998）：通过将远场算子分解为 \(F = G S G^*\)（\(G\) 有界、\(S\)
强制），障碍物的特征函数可用 \(F\)
的谱分解直接刻画。分解法是唯一不需要非线性优化的障碍物重构方法。

\subsubsection{274.5 Calderón 问题}\label{calderuxf3n-ux95eeux9898}

Calderón
问题（电阻抗断层成像的数学基础）是反问题领域中最著名的开放问题之一。

\textbf{定义 274.7}（Calderón 问题，1980）：设
\(\Omega \subset \mathbb{R}^n\)（\(n \geq 2\)）是有界 Lipschitz
区域，其内部电导率为正函数 \(\gamma(x) > 0\)。对边界 \(\partial\Omega\)
上施加的电势 \(f\)，内部电势 \(u\) 满足椭圆型方程
\(\operatorname{div}(\gamma \nabla u) = 0\)（在 \(\Omega\)
中），\(u|_{\partial\Omega} = f\)。定义
\textbf{Dirichlet-to-Neumann（DtN）映射}
\(\Lambda_\gamma: H^{1/2}(\partial\Omega) \to H^{-1/2}(\partial\Omega)\)
为
\(\Lambda_\gamma f = \gamma \frac{\partial u}{\partial\nu}|_{\partial\Omega}\)（边界法向导数，即边界电流密度）。\textbf{Calderón
问题}问：从边界测量 \(\Lambda_\gamma\) 能否唯一确定内部电导率
\(\gamma(x)\)？换言之，能否通过施加边界电势并测量边界电流来重构物体内部的电导率分布？这是医学电阻抗断层成像（EIT）和地球物理电阻率成像的核心数学问题。

\textbf{定理 274.5}（Calderón 唯一性结果，1980）：当
\(\gamma \in C^\infty\) 且 \(\gamma\)
在边界附近为常数时，\(\Lambda_\gamma\) 唯一确定 \(\gamma\)（在
\(\mathbb{R}^n\) 中，\(n \geq 2\)）。Calderón
本人使用复几何光学解（CGO解 / Complex Geometrical Optics
Solutions）给出了该定理的证明概要。

\textbf{定理 274.6}（Sylvester-Uhlmann 唯一性，1987）：对
\(\mathbb{R}^n\)（\(n \geq 3\)）中的光滑正电导率 \(\gamma\)，DtN 映射
\(\Lambda_\gamma\) 唯一确定
\(\gamma\)。证明使用指数增长解（\(\nabla^2 u + \zeta \cdot \nabla u = 0\)
形如 \(u = e^{x\cdot\zeta}(1 + \psi)\) 的复几何光学解），将问题转化为
Fourier 变换的可逆性。在 \(\mathbb{R}^2\) 中，Calderón 问题的唯一性由
Nachman（1996）使用 \(\bar{\partial}\)-方法完全解决。Calderón
问题至今仍是反问题研究中极为活跃的领域，其与隐身、超材料等前沿应用的关联使该问题持续受到关注。

\textbf{定义 274.5}（逆源问题 / Inverse Source
Problem）：从辐射场恢复源的分布。电磁/声波源 \(J(\mathbf{x})\)
产生的场满足 \(\nabla^2 u + k^2 u = J\)，从远场数据恢复
\(J\)。逆源问题一般具有严重的非唯一性------非辐射源（non-radiating
sources）产生的场在观测区域内为零。

\textbf{定义 274.6}（全波形反演 / Full Waveform Inversion，Tarantola
1984）：FWI
使用全时域波场数据（而非仅远场或初至时间），通过最小化数据残差
\(\|d_{\text{obs}} - d_{\text{sim}}(m)\|^2\)（\(m\)
为介质参数）来反演地下结构。FWI
是地震勘探中的黄金标准，涉及极大规模的非线性 PDE 约束优化。

\bigskip\hrule\bigskip

\subsection{第275章：Bayes
反演框架}\label{ux7b2c275ux7ae0bayes-ux53cdux6f14ux6846ux67b6}

Bayes
反演将反问题从确定性的正则化框架推广到概率框架，将不确定性量化系统地纳入反演过程。

\subsubsection{275.1
反问题的概率表述}\label{ux53cdux95eeux9898ux7684ux6982ux7387ux8868ux8ff0}

\textbf{定义 275.1}（反问题的 Bayes 表述）：设观测数据
\(y \in \mathbb{R}^m\)，模型参数
\(\theta \in \mathbb{R}^n\)（可能极高维）。正模型（含噪声）：
\[y = G(\theta) + \eta, \quad \eta \sim \pi_{\text{noise}}\]

由 Bayes 定理，\textbf{后验分布}为
\[\pi(\theta | y) = \frac{\pi(y | \theta) \pi(\theta)}{\pi(y)} \propto \pi(y | \theta) \pi(\theta)\]

其中 \(\pi(\theta)\)
为\textbf{先验分布}（编码物理约束和先验知识），\(\pi(y | \theta)\)
为\textbf{似然函数}（由噪声模型确定）。

\textbf{定义 275.2}（Gauss 线性反问题的后验）：若
\(G(\theta) = A \theta\) 且
\(\eta \sim \mathcal{N}(0, \Sigma_\eta)\)，先验
\(\theta \sim \mathcal{N}(m_0, \Sigma_0)\)，则后验
\(\theta | y \sim \mathcal{N}(m_{\text{post}}, \Sigma_{\text{post}})\)
其中
\[\Sigma_{\text{post}}^{-1} = A^T \Sigma_\eta^{-1} A + \Sigma_0^{-1}, \quad m_{\text{post}} = \Sigma_{\text{post}} (A^T \Sigma_\eta^{-1} y + \Sigma_0^{-1} m_0)\]

Tikhonov 正则化是 Gauss 线性反问题的 MAP（最大后验）估计在
\(\Sigma_\eta = \delta^2 I\)、\(\Sigma_0 = \alpha^{-1} I\) 下的等价值。

\subsubsection{275.2 先验建模}\label{ux5148ux9a8cux5efaux6a21}

\textbf{定义 275.3}（Markov 随机场先验 / MRF Prior）：参数场
\(\theta(\mathbf{x})\) 的空间相关结构通过 MRF 或 Gauss
过程的协方差函数编码。\textbf{Matheron 的 Matérn 协方差}族：
\[C_\nu(r) = \sigma^2 \frac{2^{1-\nu}}{\Gamma(\nu)} \left( \frac{\sqrt{2\nu} r}{\ell} \right)^\nu K_\nu\left( \frac{\sqrt{2\nu} r}{\ell} \right)\]

\(\nu\) 控制光滑度，\(\ell\) 为相关长度。\(\nu \to \infty\) 对应 Gauss
协方差（无限光滑），\(\nu = 1/2\) 对应指数协方差。

\textbf{定义 275.4}（稀疏促进先验 / Laplace 先验与 TV
先验）：\(\ell_1\)-正则化对应 Laplace 先验
\(\pi(\theta) \propto e^{-\alpha \|\theta\|_1}\)。全变差（TV）正则化对应
\(\pi(\theta) \propto e^{-\alpha \|\nabla \theta\|_1}\)（梯度场的
Laplace 先验）。

\subsubsection{275.3 高维后验采样与
MCMC}\label{ux9ad8ux7ef4ux540eux9a8cux91c7ux6837ux4e0e-mcmc}

\textbf{定义 275.5}（MCMC 用于反问题 / Markov Chain Monte
Carlo）：在高维参数空间中从后验 \(\pi(\theta | y)\) 采样。核心方法包括：
- \textbf{Metropolis-Hastings}（MH）：从提议分布中采样候选
\(\theta'\)，以概率
\(\min(1, \frac{\pi({\theta'}|y) q(\theta|\theta')}{\pi(\theta|y) q(\theta'|\theta)})\)
接受 - \textbf{Hamiltonian Monte Carlo}（HMC / 混合 Monte
Carlo）：利用后验的梯度信息（通过伴随方法高效计算），以 Hamiltonian
动力学进行提议 - \textbf{Langevin 扩散}（MALA）：使用 Langevin SDE
的离散化：\(\theta_{k+1} = \theta_k + \frac{\varepsilon^2}{2} \nabla \log \pi(\theta_k | y) + \varepsilon \xi_k\)（\(\xi_k \sim \mathcal{N}(0, I)\)）

\textbf{定理 275.1}（Cui-Law-Marzouk / dimension-independent
MCMC，2016）：使用函数空间观点（将先验和后验视为无限维测度），合理设计的
MCMC
方法的混合时间可以是\textbf{维数无关的}------即使在无限维极限下也是良定义的。

\bigskip\hrule\bigskip

\subsection{第276章：数据同化与 Kalman
反演}\label{ux7b2c276ux7ae0ux6570ux636eux540cux5316ux4e0e-kalman-ux53cdux6f14}

数据同化（Data
Assimilation）将观测数据与数学模型动态融合，是现代天气预报和气候科学的核心技术。

\subsubsection{276.1 Kalman
滤波与扩展}\label{kalman-ux6ee4ux6ce2ux4e0eux6269ux5c55}

\textbf{定义 276.1}（离散时间 Kalman 滤波 / 同化中的 Kalman
滤波）：设状态方程为 \(x_{k+1} = M_k x_k + w_k\)，观测方程为
\(y_k = H_k x_k + v_k\)（\(w_k \sim \mathcal{N}(0, Q_k)\)
为模型误差，\(v_k \sim \mathcal{N}(0, R_k)\) 为观测误差）。Kalman
滤波递归地产生最优线性无偏估计。

\textbf{定理 276.1}（Kalman 滤波 = 动态线性 Bayes 反演）：Kalman 滤波是
Gauss 线性反问题后验的递推版本。预测步计算先验
\(x_k | y_{1:k-1}\)，更新步计算后验 \(x_k | y_{1:k}\)。

\textbf{定义 276.2}（扩展 Kalman 滤波 / EKF）：对非线性模型
\(M_k(\cdot)\) 和观测 \(h_k(\cdot)\)，EKF
在每个时间步进行切向线性化（Jacobian 计算）。

\subsubsection{276.2 系综 Kalman
滤波}\label{ux7cfbux7efc-kalman-ux6ee4ux6ce2}

\textbf{定义 276.3}（系综 Kalman 滤波 / Ensemble Kalman Filter /
EnKF，Evensen 1994）：利用蒙特卡洛系综 \(\{x_k^{(i)}\}_{i=1}^N\)
近似状态分布，通过系综协方差代替 EKF 中的显式协方差矩阵传播：
\[P_k^f = \frac{1}{N-1} \sum_{i=1}^N (x_k^{f,(i)} - \bar{x}_k^f)(x_k^{f,(i)} - \bar{x}_k^f)^T\]

EnKF
避免了显式计算切向模型和高维协方差矩阵，使数据同化在高维地球系统模型中实际可行。

\textbf{定理 276.2}（EnKF 的收敛性）：当系综大小 \(N \to \infty\)，EnKF
的均值和协方差以 \(O(1/\sqrt{N})\) 的误差收敛到 Kalman 滤波的对应值。

\subsubsection{276.3 变分同化与
4D-Var}\label{ux53d8ux5206ux540cux5316ux4e0e-4d-var}

\textbf{定义 276.4}（四维变分同化 / 4D-Var，Le Dimet-Talagrand 1986,
Courtier 1994）：\textbf{4D-Var} 是最小化以下代价函数的 PDE
约束优化问题：
\[J(x_0) = \frac{1}{2}(x_0 - x_b)^T B^{-1} (x_0 - x_b) + \frac{1}{2} \sum_{k=0}^T (y_k - H_k(x_k))^T R_k^{-1} (y_k - H_k(x_k))\]

约束：\(x_{k+1} = M_k(x_k)\)。梯度通过伴随方法（adjoint
method）计算------正向运行模型存储轨迹，反向运行伴随模型传递梯度。

\textbf{定理 276.3}（4D-Var 与 Kalman 的等价）：在线性模型和 Gauss
噪声假设下，4D-Var 在分析时间窗口末尾的最优状态估计等价于 Kalman
平滑。在实际非线性系统中，4D-Var 和 EnKF 可以结合成混合方法（Hybrid
4D-EnVar）。

\subsubsection{276.4
粒子滤波与非线性同化}\label{ux7c92ux5b50ux6ee4ux6ce2ux4e0eux975eux7ebfux6027ux540cux5316}

\textbf{定义 276.5}（粒子滤波 / Particle Filter / Sequential Monte
Carlo）：在高度非线性非 Gauss 系统中，使用加权的粒子系综（而非 Gauss
近似）表示后验分布。挑战是\textbf{维数灾难}：有效粒子数随状态维数指数衰减（Snyder-Bengtsson-Bickel
2008）。

\textbf{定理
276.4}（最优提议分布与重要性采样）：使用模型动力学作为提议的最优重要性采样，粒子的权重更新为似然函数：\(w_k^{(i)} \propto w_{k-1}^{(i)} \cdot \pi(y_k | x_k^{(i)})\)。重采样步骤（重要性重采样
/ SIR）防止粒子退化的权值崩溃。

\bigskip\hrule\bigskip

\emph{本卷在泛函分析（V14）和偏微分方程（V18）的基础上，系统建立了不适定问题与正则化理论、Radon
变换与 CT 成像、散射反问题、Bayes 反演框架和数据同化方法的数学理论。共 5
章。}

\newpage

\section{卷五十：积分方程理论}\label{ux5377ux4e94ux5341ux79efux5206ux65b9ux7a0bux7406ux8bba}

\begin{quote}
积分方程是泛函分析在数学物理和工程中的重要应用领域。积分方程将未知函数置于积分号下，与微分方程形成互补------许多微分方程初边值问题可以转化为积分方程，而积分方程往往具有更好的正则性。本卷系统建立奇异积分方程、Wiener-Hopf
方程以及积分方程数值求解方法，并将在卷首补充 Fredholm 和 Volterra
积分方程的完整理论。
\end{quote}

\begin{quote}
\textbf{卷序说明}：本章节号（Ch
202-206）反映其在全书逻辑链中的位置。物理卷序上本卷位于V50，作为积分方程方向的专题补充卷。
\end{quote}

\bigskip\hrule\bigskip

\subsection{第202章：Fredholm
积分方程}\label{ux7b2c202ux7ae0fredholm-ux79efux5206ux65b9ux7a0b}

Fredholm 积分方程是积分方程理论的历史起源，Fredholm 将线性代数中 Cramer
法则的思想推广到无穷维空间，建立了 Fredholm 行列式、预解核和 Fredholm
择一定理。本章系统阐述第二类 Fredholm 积分方程的理论。

\subsubsection{202.1 Fredholm
积分方程的基本分类}\label{fredholm-ux79efux5206ux65b9ux7a0bux7684ux57faux672cux5206ux7c7b}

\textbf{定义 202.1}（Fredholm 积分方程）：\textbf{第二类 Fredholm
积分方程}为

\[\varphi(x) - \lambda \int_a^b K(x, y) \varphi(y) \, dy = f(x), \quad x \in [a, b]\]

其中 \(K(x, y) \in L^2([a,b] \times [a,b])\)
为核函数，\(\lambda \in \mathbb{C}\) 为参数，\(f\)
为已知自由项，\(\varphi\) 为未知函数。当 \(f \equiv 0\)
时称为\textbf{齐次方程}，否则称为\textbf{非齐次方程}。\textbf{第一类
Fredholm 方程}为
\(\int_a^b K(x,y) \varphi(y) \, dy = f(x)\)，由于它是严重不适定的（解对数据的微小扰动不连续依赖），本卷不予讨论。

\textbf{定义 202.2}（共轭方程）：齐次共轭方程为

\[\psi(x) - \overline{\lambda} \int_a^b \overline{K(y, x)} \psi(y) \, dy = 0\]

共轭方程在刻画原方程可解性条件时起关键作用。

\subsubsection{202.2 Fredholm
行列式与预解核}\label{fredholm-ux884cux5217ux5f0fux4e0eux9884ux89e3ux6838}

\textbf{定义 202.3}（Fredholm 行列式）：定义 Fredholm 行列式

\[D(\lambda) = \sum_{n=0}^\infty \frac{(-\lambda)^n}{n!} \int_a^b \cdots \int_a^b \det\begin{pmatrix} K(x_1, x_1) & \cdots & K(x_1, x_n) \\ \vdots & \ddots & \vdots \\ K(x_n, x_1) & \cdots & K(x_n, x_n) \end{pmatrix} dx_1 \cdots dx_n\]

\(D(\lambda)\) 是 \(\lambda\)
的整函数（在整个复平面上解析）。\(D(\lambda)\) 的零点恰为 Fredholm
积分算子的特征值。

\textbf{定义 202.4}（Fredholm 一阶子式与预解核）：定义 Fredholm 一阶子式

\[D(x, y; \lambda) = K(x, y) D(\lambda) + \sum_{n=1}^\infty \frac{(-\lambda)^n}{n!} \int_a^b \cdots \int_a^b \det\begin{pmatrix} K(x, y) & K(x, x_1) & \cdots & K(x, x_n) \\ K(x_1, y) & K(x_1, x_1) & \cdots & K(x_1, x_n) \\ \vdots & \vdots & \ddots & \vdots \\ K(x_n, y) & K(x_n, x_1) & \cdots & K(x_n, x_n) \end{pmatrix} dx_1 \cdots dx_n\]

对任意使 \(D(\lambda) \neq 0\) 的
\(\lambda\)，定义\textbf{预解核}（resolvent kernel）

\[R(x, y; \lambda) = \frac{D(x, y; \lambda)}{D(\lambda)}\]

\textbf{定理 202.1}（预解核表示的 Fredholm 解）：若
\(D(\lambda) \neq 0\)，则第二类 Fredholm 积分方程有唯一解

\[\varphi(x) = f(x) + \lambda \int_a^b R(x, y; \lambda) f(y) \, dy\]

\(R(x, y; \lambda)\) 是 \(\lambda\) 的半纯函数------在 \(\mathbb{C}\)
上除极点（特征值）外全纯。这类似于线性方程组 \(Ax = b\) 的解
\(x = A^{-1}b\)，预解核对应对 \((I - \lambda T)^{-1}\) 的积分核表示。

\emph{证明概要}（Fredholm 第一定理：\(I - \lambda T\) 可逆 \(\iff\)
\(N(I - \lambda T) = \{0\}\)，有限维类比 + Riesz-Schauder
理论）：在有限维空间 \(\mathbb{C}^n\) 中，方阵 \(A\) 可逆当且仅当
\(Ax = 0\) 仅有零解------Fredholm 理论将此推广至无穷维。Fredholm
积分算子 \(T\varphi(x) = \int_a^b K(x,y)\varphi(y)dy\) 是 \(L^2\)
上的紧算子。由 Riesz-Schauder 理论，紧算子 \(T\) 的谱 \(\sigma(T)\) 除
\(0\) 外仅由特征值构成，且每个非零特征值具有有限代数重数。预解算子
\(R_\lambda(T) = (I - \lambda T)^{-1}\)
在非特征值处存在并有界，其积分核恰为预解核
\(R(x,y;\lambda) = D(x,y;\lambda)/D(\lambda)\)，其中 \(D(\lambda)\) 是
Fredholm 行列式------它是有限维行列式 \(\det(I - \lambda A)\)
的无穷维推广。当 \(D(\lambda) \neq 0\)（即 \(\lambda\)
不是特征值）时，\(I - \lambda T\) 双射且逆有界，方程有唯一解。∎

\subsubsection{202.3 Fredholm
择一定理}\label{fredholm-ux62e9ux4e00ux5b9aux7406}

\textbf{定理 202.2}（Fredholm 择一定理）：对第二类 Fredholm 积分方程
\(\varphi - \lambda T\varphi = f\)，以下二者恰有一个成立：

\begin{enumerate}
\def\labelenumi{\arabic{enumi}.}
\tightlist
\item
  非齐次方程对任意 \(f \in L^2\) 存在唯一解（此时 \(\lambda\)
  不是特征值）；
\item
  齐次方程有非零解（此时 \(\lambda\) 是特征值）。
\end{enumerate}

在情形 2 中，非齐次方程有解当且仅当 \(f\) 与共轭齐次方程的所有解正交：

\[\int_a^b f(x) \overline{\psi(x)} \, dx = 0, \quad \forall \psi \text{ 满足 } \psi(x) = \overline{\lambda} \int_a^b \overline{K(y, x)} \psi(y) \, dy\]

此外，齐次方程及其共轭齐次方程的解空间维数有限且相等：
\[\dim \ker(I - \lambda T) = \dim \ker(I - \overline{\lambda} T^*) < \infty\]

Fredholm 择一定理是线性代数中 ``方阵 \(A\) 要么可逆，要么 \(Ax = 0\)
有非零解'' 的无穷维推广，而维数等式
\(\dim \ker = \dim \operatorname{coker}\)
是有限维线性代数中矩阵与其转置具有相同秩这一事实的无穷维类比。

\textbf{定理 202.3}（特征值的离散性）：若
\(K \in L^2([a,b] \times [a,b])\)，则 Fredholm 积分算子
\(T\varphi(x) = \int_a^b K(x,y) \varphi(y) \, dy\)
的非零特征值构成至多可数集，唯一可能的聚点是
\(0\)。每个非零特征值的几何重数（相应特征空间的维数）有限。特别地，\(T\)
是 \(L^2\) 上的紧算子。

\subsubsection{202.4 对称核的 Hilbert-Schmidt
理论}\label{ux5bf9ux79f0ux6838ux7684-hilbert-schmidt-ux7406ux8bba}

\textbf{定义 202.5}（对称核）：若 \(K(x, y) = \overline{K(y, x)}\)，则称
\(K\) 为\textbf{对称核}（实核时即为 Hermite 对称性
\(K(x,y) = K(y,x)\)）。此时积分算子
\(T\varphi(x) = \int_a^b K(x,y) \varphi(y) \, dy\) 是 \(L^2\)
上的紧自伴算子。

\textbf{定理 202.4}（Hilbert-Schmidt 展开定理）：设 \(K\) 为 \(L^2\)
对称核，\(T\) 为相应的积分算子。则存在标准正交特征函数系
\(\{\varphi_n\}\) 和实特征值
\(\{\lambda_n\}\)，\(\lambda_n \to 0\)（\(n \to \infty\)），使得

\[K(x, y) \sim \sum_{n} \frac{\varphi_n(x) \overline{\varphi_n(y)}}{\lambda_n}\]

上式在 \(L^2\) 意义下收敛（即
\(\|K(x,y) - \sum_{n=1}^N \lambda_n^{-1} \varphi_n(x) \overline{\varphi_n(y)}\|_{L^2} \to 0\)）。进一步，若
\(K\) 连续，则对任意使 \(D(\lambda) \neq 0\) 的 \(\lambda\)，解可展开为
\(f\) 的 Fourier 级数：

\[\varphi(x) = f(x) + \lambda \sum_n \frac{\langle f, \varphi_n \rangle}{\lambda_n - \lambda} \varphi_n(x)\]

\emph{证明概要}（紧自伴算子的谱定理）：紧自伴算子 \(T\) 的 Hilbert
空间谱定理确保了 \(T\) 可对角化：存在标准正交基 \(\{\varphi_n\}\)
和实特征值 \(\lambda_n \to 0\)，使得
\(T\varphi = \sum_n \lambda_n^{-1} \langle \varphi, \varphi_n \rangle \varphi_n\)。核函数
\(K(x,y)\) 可视作 \(T\) 的 Schwartz 核，从而
\(K(x,y) = \sum_n \lambda_n^{-1} \varphi_n(x) \overline{\varphi_n(y)}\)
在 \(L^2\) 意义下成立（即紧算子 Mercer
型展开）。代入积分方程并投影到各特征方向上，得
\(\langle \varphi, \varphi_n \rangle = \langle f, \varphi_n \rangle + \lambda \lambda_n^{-1} \langle \varphi, \varphi_n \rangle\)，解得
\(\langle \varphi, \varphi_n \rangle = \frac{\lambda_n}{\lambda_n - \lambda} \langle f, \varphi_n \rangle\)，经
Fourier 合成即得解的级数表达式。∎

\textbf{定理 202.5}（Mercer 定理）：若 \(K\) 是连续对称核且 \(T\)
是正算子（所有特征值 \(\lambda_n > 0\)），则级数
\(\sum_n \lambda_n^{-1} \varphi_n(x) \overline{\varphi_n(y)}\)
绝对且一致收敛到 \(K(x, y)\)。

\bigskip\hrule\bigskip

\subsection{第203章：Volterra
积分方程}\label{ux7b2c203ux7ae0volterra-ux79efux5206ux65b9ux7a0b}

Volterra 积分方程与 Fredholm 积分方程的根本区别在于积分上限是变量 \(x\)
而非固定值 \(b\)。这一差异带来了深刻的结构性变化：Volterra
算子具有无反馈（non-feedback）特性，是拟幂零算子，对任意参数 \(\lambda\)
方程均有唯一解。Volterra 方程与常微分方程初值问题有自然的等价关系。

\subsubsection{203.1 Volterra
积分方程的基本形式}\label{volterra-ux79efux5206ux65b9ux7a0bux7684ux57faux672cux5f62ux5f0f}

\textbf{定义 203.1}（Volterra 积分方程）：\textbf{第二类 Volterra
积分方程}为

\[\varphi(x) - \lambda \int_a^x K(x, y) \varphi(y) \, dy = f(x), \quad x \in [a, b]\]

其中积分上限为变量 \(x\)，这是与 Fredholm 型（上限为固定值
\(b\)）的关键区别。\(K(x, y)\) 定义在三角区域 \(a \leq y \leq x \leq b\)
上。\textbf{第一类 Volterra 方程}为
\(\int_a^x K(x,y) \varphi(y) \, dy = f(x)\)，若 \(K(x,x) \neq 0\)
可通过求导化为第二类方程。

\textbf{定义 203.2}（Volterra 算子）：定义 Volterra 积分算子

\[(T\varphi)(x) = \int_a^x K(x, y) \varphi(y) \, dy\]

\subsubsection{203.2 逐次逼近法（Neumann
级数）}\label{ux9010ux6b21ux903cux8fd1ux6cd5neumann-ux7ea7ux6570}

\textbf{定理 203.1}（Neumann 级数的收敛性）：若 \(K(x, y)\) 在
\(a \leq y \leq x \leq b\) 上连续，则对任意 \(\lambda \in \mathbb{C}\)
和任意连续函数 \(f(x)\)，第二类 Volterra
积分方程存在唯一连续解。解由一致收敛的 Neumann 级数给出：

\[\varphi(x) = f(x) + \sum_{n=1}^\infty \lambda^n \int_a^x K_n(x, y) f(y) \, dy\]

其中 \(K_1 = K\)，迭代核 \(K_n\) 由递推定义：

\[K_n(x, y) = \int_y^x K(x, z) K_{n-1}(z, y) \, dz \quad (n \geq 2)\]

\textbf{证明}：定义迭代序列 \(\varphi_0(x) = f(x)\)，

\[\varphi_{n+1}(x) = f(x) + \lambda \int_a^x K(x, y) \varphi_n(y) \, dy\]

令 \(M = \sup_{a \leq y \leq x \leq b} |K(x, y)|\)。用归纳法可证

\[|\varphi_{n+1}(x) - \varphi_n(x)| \leq |\lambda|^n M^n \frac{(x-a)^n}{n!} \|f\|_\infty\]

因此级数 \(\sum_{n=0}^\infty (\varphi_{n+1}(x) - \varphi_n(x))\)
由指数级数控制，由 Weierstrass M-判别法知该级数在 \([a,b]\)
上一致收敛。取极限 \(\varphi(x) = \lim_{n \to \infty} \varphi_n(x)\)
即为方程的解。唯一性同样由归纳估计得出。 ∎

\textbf{注记}：与 Fredholm 方程不同，Neumann 级数对任意 \(\lambda\)
均收敛------这在 Volterra 算子的拟幂零性中得到了本质解释。

\subsubsection{203.3 Volterra
方程与常微分方程的联系}\label{volterra-ux65b9ux7a0bux4e0eux5e38ux5faeux5206ux65b9ux7a0bux7684ux8054ux7cfb}

\textbf{定理 203.2}（Volterra 方程与 ODE 的等价性）：\(n\)
阶线性常微分方程的初值问题

\[y^{(n)}(x) + a_{n-1}(x) y^{(n-1)}(x) + \cdots + a_0(x) y(x) = g(x)\]

\[y(a) = y_0, \; y'(a) = y_0', \; \ldots, \; y^{(n-1)}(a) = y_0^{(n-1)}\]

等价于第二类 Volterra 积分方程

\[y(x) = F(x) + \int_a^x K(x, t) y(t) \, dt\]

其中

\[K(x, t) = -\sum_{k=1}^n a_{n-k}(x) \frac{(x-t)^{k-1}}{(k-1)!}\]

且 \(F(x)\) 由 \(g(x)\) 和初值确定。这一等价关系将 ODE
初值问题的存在唯一性理论统一到 Volterra 积分方程的框架下：ODE 的 Picard
逐次逼近法恰对应于 Volterra 方程的 Neumann 级数展开。

\textbf{例 203.1}：最简单的 Volterra 方程
\(\varphi(x) = 1 + \int_0^x \varphi(y) \, dy\) 等价于 ODE
\(\varphi'(x) = \varphi(x)\)，\(\varphi(0) = 1\)，解为
\(\varphi(x) = e^x\)。该方程的 Neumann 级数恰为指数函数的 Taylor 展开。

\subsubsection{203.4 Volterra
算子的拟幂零性}\label{volterra-ux7b97ux5b50ux7684ux62dfux5e42ux96f6ux6027}

\textbf{定理 203.3}（拟幂零性与整体可解性）：第二类 Volterra
积分方程对任意 \(\lambda \in \mathbb{C}\) 存在唯一解。等价地，Volterra
积分算子 \(T\) 在任意连续函数构成的 Banach 空间上满足谱半径
\(\rho(T) = 0\)，即 \(T\) 是\textbf{拟幂零算子}（quasinilpotent
operator）。

\textbf{证明思路}：由定理 203.1
的迭代估计，\(\|T^n\|^{1/n} \leq M \cdot \frac{(b-a)}{n^{1/n}} \to 0\)，故
\(\rho(T) = \lim_{n \to \infty} \|T^n\|^{1/n} = 0\)。 ∎

\textbf{推论}：\((I - \lambda T)^{-1}\) 是 \(\lambda\)
的整函数（在整个复平面上解析），预解式无极点。这与 Fredholm
型的预解核形成鲜明对比------Fredholm
预解核仅在特征值的补集上解析，有极点。

\bigskip\hrule\bigskip

\subsection{第204章：奇异积分方程}\label{ux7b2c204ux7ae0ux5947ux5f02ux79efux5206ux65b9ux7a0b}

奇异积分方程是积分方程理论中技术难度最高、应用最广的分支之一。其核心特征在于积分核在积分区域内部存在不可积奇异性（如
Cauchy 核 \(1/(y-x)\)），经典 Riemann 积分失效，必须借助 Cauchy
主值积分或 Hadamard
有限部分积分等广义积分概念。奇异积分方程在弹性力学（裂纹尖端应力分析）、流体力学（翼型绕流）和电磁学（天线辐射）中有重要应用，并通过
Plemelj 公式与复变函数论中的 Riemann-Hilbert
边值问题建立了深刻等价关系。

\subsubsection{204.1 Cauchy
奇异积分方程}\label{cauchy-ux5947ux5f02ux79efux5206ux65b9ux7a0b}

\textbf{定义 204.1}（Cauchy 奇异积分方程）：\textbf{Cauchy
奇异积分方程}具有形式

\[a(x) \varphi(x) + \frac{b(x)}{\pi i} \operatorname{p.v.} \int_\Gamma \frac{\varphi(y)}{y - x} \, dy = f(x)\]

其中 \(\Gamma\) 是光滑曲线（或闭曲线），\(\operatorname{p.v.}\) 表示
Cauchy 主值积分。当 \(a(x) \equiv 0\)
时称为\textbf{第一类}，\(a(x) \neq 0\) 时称为\textbf{第二类}。

\textbf{定义 204.2}（Cauchy 主值积分）：沿曲线 \(\Gamma\) 的 Cauchy
主值积分定义为

\[\operatorname{p.v.} \int_\Gamma \frac{\varphi(y)}{y - x} \, dy = \lim_{\varepsilon \to 0} \int_{\Gamma \setminus B(x, \varepsilon)} \frac{\varphi(y)}{y - x} \, dy\]

\textbf{定理 204.1}（Plemelj 公式 / Sokhotski-Plemelj 公式）：设
\(\Gamma\) 是光滑弧，\(\varphi\) 是 Hölder 连续函数。定义 Cauchy 型积分

\[\Phi(z) = \frac{1}{2\pi i} \int_\Gamma \frac{\varphi(y)}{y - z} \, dy \quad (z \notin \Gamma)\]

则 \(\Phi(z)\) 在 \(\Gamma\) 两侧有边界值，满足

\[\Phi^\pm(x) = \lim_{z \to x^\pm} \Phi(z) = \pm \frac{1}{2} \varphi(x) + \frac{1}{2\pi i} \operatorname{p.v.} \int_\Gamma \frac{\varphi(y)}{y - x} \, dy\]

因此，\(\Phi^+(x) - \Phi^-(x) = \varphi(x)\)，\(\Phi^+(x) + \Phi^-(x) = \frac{1}{\pi i} \operatorname{p.v.} \int_\Gamma \frac{\varphi(y)}{y - x} \, dy\)。

\subsubsection{204.1.1 Abel
积分方程与反演公式}\label{abel-ux79efux5206ux65b9ux7a0bux4e0eux53cdux6f14ux516cux5f0f}

\textbf{定义 204.1.1}（Abel 积分方程）：经典的 Abel
积分方程为第一类弱奇异 Volterra 方程：

\[\int_0^x \frac{\varphi(y)}{(x-y)^\alpha} \, dy = f(x), \quad 0 < \alpha < 1\]

当 \(\alpha = 1/2\) 时即为等时降落问题导出的标准 Abel 方程。

\textbf{定理 204.1.1}（Abel 反演公式）：若 \(f\) 绝对连续且
\(f(0) = 0\)，则 Abel 方程的解为

\[\varphi(x) = \frac{\sin(\pi\alpha)}{\pi} \frac{d}{dx} \int_0^x \frac{f(y)}{(x-y)^{1-\alpha}} \, dy\]

\emph{证明概要}（Laplace 变换法）：方程左端是 \(\varphi\) 与核
\(k(x) = x^{-\alpha}\) 的卷积：\((k \ast \varphi)(x) = f(x)\)。取
Laplace 变换，由卷积定理得
\(\mathcal{L}\{k\}(s) \cdot \mathcal{L}\{\varphi\}(s) = \mathcal{L}\{f\}(s)\)。已知
\(\mathcal{L}\{x^{-\alpha}\}(s) = \Gamma(1-\alpha) s^{\alpha-1}\)，故
\(\mathcal{L}\{\varphi\}(s) = \frac{s^{1-\alpha}}{\Gamma(1-\alpha)} \mathcal{L}\{f\}(s)\)。由
Laplace 变换的微分性质及
\(\mathcal{L}\{x^{\alpha-1}\}(s) = \Gamma(\alpha) s^{-\alpha}\)，取逆变换并利用
\(\Gamma(\alpha)\Gamma(1-\alpha) = \pi / \sin(\pi\alpha)\)
即得反演公式。∎

\subsubsection{204.2 Riemann-Hilbert
问题}\label{riemann-hilbert-ux95eeux9898}

\textbf{定义 204.3}（Riemann-Hilbert 问题 / 边值问题）：设 \(\Gamma\)
是简单闭曲线。求 \(\Gamma\) 外的全纯函数 \(\Phi(z)\)，满足 \(\Gamma\)
上的边界条件

\[\Phi^+(x) = G(x) \Phi^-(x) + g(x) \quad (x \in \Gamma)\]

其中 \(G(x), g(x)\) 是 \(\Gamma\) 上的给定函数（\(G(x) \neq 0\)）。

\textbf{定理 204.2}（标量 Riemann-Hilbert 问题的解）：设 \(\Gamma\)
是简单闭曲线，\(G(x) \neq 0\) 是 Hölder 连续函数。定义
\(\kappa = \frac{1}{2\pi i} [\log G]_\Gamma\)（\(G\)
的指标，index）。则： - 若 \(\kappa \geq 0\)，齐次问题（\(g = 0\)）有
\(\kappa + 1\) 个线性独立解 - 若
\(\kappa < 0\)，齐次问题只有零解，非齐次问题可解当且仅当 \(g\) 满足
\(-\kappa - 1\) 个正交条件

\textbf{定义 204.4}（Cauchy 奇异积分方程与 Riemann-Hilbert
问题的等价性）：通过 Plemelj 公式，Cauchy 奇异积分方程可转化为
Riemann-Hilbert 问题，反之亦然。具体地，方程

\[a(x) \varphi(x) + \frac{b(x)}{\pi i} \operatorname{p.v.} \int_\Gamma \frac{\varphi(y)}{y - x} \, dy = f(x)\]

等价于 \(\Phi\) 满足
\(\Phi^+(x) = \frac{a(x) - b(x)}{a(x) + b(x)} \Phi^-(x) + \frac{f(x)}{a(x) + b(x)}\)。

\subsubsection{204.3
奇异积分方程与弹性力学}\label{ux5947ux5f02ux79efux5206ux65b9ux7a0bux4e0eux5f39ux6027ux529bux5b66}

\textbf{定义
204.5}（平面弹性力学中的奇异积分方程）：平面弹性力学的基本边值问题可化为
Cauchy 奇异积分方程。位移间断和应力集中引发核的奇异性。

\textbf{定理 204.3}（Muskhelishvili
方法，1953）：平面弹性力学中的边值问题可通过复变函数方法化为 Cauchy
奇异积分方程，使用 Kolosov-Muskhelishvili
势函数。这是积分方程在力学中最重要的应用之一。

\subsubsection{204.4
弱奇异核与强奇异核}\label{ux5f31ux5947ux5f02ux6838ux4e0eux5f3aux5947ux5f02ux6838}

\textbf{定义 204.6}（弱奇异核）：核
\(K(x, y) = \frac{H(x, y)}{|x-y|^\alpha}\)（\(0 < \alpha < d\)，\(d\)
是空间维数）称为\textbf{弱奇异核}。相应的积分算子在 \(L^p\) 上是紧的。

\textbf{定义 204.7}（强奇异核 / 超奇异核）：核
\(K(x, y) = \frac{H(x, y)}{|x-y|^\alpha}\)（\(\alpha \geq d\)）称为\textbf{强奇异核}。相应的积分必须理解为
Hadamard 有限部分积分（finite part integral）。

\textbf{定义 204.8}（Hadamard 有限部分积分）：对 \(\alpha \geq 1\)，

\[\operatorname{f.p.} \int_a^b \frac{\varphi(y)}{(x-y)^\alpha} \, dy = \lim_{\varepsilon \to 0} \left[ \int_a^{x-\varepsilon} + \int_{x+\varepsilon}^b \frac{\varphi(y)}{(x-y)^\alpha} \, dy - \sum_{k=0}^{\lfloor\alpha\rfloor-1} \frac{\varphi^{(k)}(x)}{k!} \frac{1-(-1)^{\alpha-k-1}}{(\alpha-k-1) \varepsilon^{\alpha-k-1}} \right]\]

超奇异积分方程在断裂力学和边界元方法中有重要应用。

\bigskip\hrule\bigskip

\subsection{第205章：Wiener-Hopf
积分方程}\label{ux7b2c205ux7ae0wiener-hopf-ux79efux5206ux65b9ux7a0b}

Wiener-Hopf 方程是卷积型积分方程在半无穷区间上的推广，由 Wiener 和 Hopf
在 1931 年研究辐射平衡问题时引入。

\subsubsection{205.1 Wiener-Hopf
方程的基本形式}\label{wiener-hopf-ux65b9ux7a0bux7684ux57faux672cux5f62ux5f0f}

\textbf{定义 205.1}（Wiener-Hopf 积分方程）：\textbf{第一类 Wiener-Hopf
方程}为

\[\int_0^\infty K(x - y) \varphi(y) \, dy = f(x) \quad (x \geq 0)\]

\textbf{第二类 Wiener-Hopf
方程}：\(\varphi(x) - \lambda \int_0^\infty K(x - y) \varphi(y) \, dy = f(x)\)（\(x \geq 0\)）。

核 \(K\) 仅依赖于两个变量的差 \(x-y\)（卷积型），积分区间是半无穷区间
\([0, \infty)\)。这是 Wiener-Hopf 方程与经典卷积方程的区别。

\textbf{定理 205.1}（Wiener-Hopf 的分解方法）：Wiener-Hopf 方程可通过
Fourier 变换和 Wiener-Hopf 因子分解求解。关键步骤：

\begin{enumerate}
\def\labelenumi{\arabic{enumi}.}
\tightlist
\item
  将方程延拓到整个实轴：\(\varphi(x) - \lambda \int_{-\infty}^\infty K(x-y) \varphi_+(y) \, dy = f_+(x) + \psi_-(x)\)
\item
  做 Fourier
  变换：\(\hat{\varphi}_+(\xi) [1 - \lambda \hat{K}(\xi)] = \hat{f}_+(\xi) + \hat{\psi}_-(\xi)\)
\item
  将 \(1 - \lambda \hat{K}(\xi)\) 分解为
  \(G_+(\xi) G_-(\xi)\)（Wiener-Hopf 分解）
\item
  通过解析延拓和 Liouville 定理确定 \(\hat{\varphi}_+\)
\end{enumerate}

\subsubsection{205.2 Wiener-Hopf
因子分解}\label{wiener-hopf-ux56e0ux5b50ux5206ux89e3}

\textbf{定义 205.2}（Wiener-Hopf 因子分解）：设函数
\(\Phi(\xi)\)（\(\xi \in \mathbb{R}\)）非零且满足适当的正则性条件。\textbf{Wiener-Hopf
分解}是

\[\Phi(\xi) = \Phi_+(\xi) \Phi_-(\xi)\]

其中 \(\Phi_+(\xi)\) 可解析延拓到上半平面 \(\Im \xi > 0\)
且在上半平面无零点，\(\Phi_-(\xi)\) 可解析延拓到下半平面 \(\Im \xi < 0\)
且在下半平面无零点。

\textbf{定理 205.2}（Wiener-Hopf 分解的存在性）：若 \(\Phi(\xi)\) 是
Hölder 连续、非零，且 \(\Phi(\xi) \to 1\)（\(|\xi| \to \infty\)），指标
\(\kappa = \frac{1}{2\pi i} [\log \Phi]_{-\infty}^\infty\)
为有限整数，则 Wiener-Hopf 分解存在（模去常数因子）。

\textbf{定义 205.3}（Wiener-Hopf 分解的构造）：\(\Phi(\xi)\) 的
Wiener-Hopf 分解由

\[\Phi_\pm(\xi) = \exp\left(\frac{1}{2\pi i} \int_{-\infty}^\infty \frac{\log \Phi(t)}{t - \xi \mp i0} \, dt\right)\]

给出。这是 Cauchy 积分在 Wiener-Hopf 方法中的核心应用。

\subsubsection{205.3 Wiener-Hopf
方程的应用}\label{wiener-hopf-ux65b9ux7a0bux7684ux5e94ux7528}

\textbf{定理 205.3}（Milne
问题，辐射传输）：半空间中的辐射传输方程可化为 Wiener-Hopf
积分方程。Milne 问题（1921）是 Wiener-Hopf 方法的最早成功应用之一。

\textbf{例 205.1}（衍射问题）：半平面衍射（Sommerfeld 问题）可通过
Wiener-Hopf 方法严格求解。这是数学物理中 Wiener-Hopf 方法最经典的应用。

\textbf{定理 205.4}（排队论中的 Wiener-Hopf 方程）：排队论中的 Lindley
方程（等待时间分布）等价于 Wiener-Hopf 积分方程。Wiener-Hopf 方法给出了
GI/G/1 排队系统等待时间的平稳分布。

\bigskip\hrule\bigskip

\subsection{第206章：积分方程的数值方法}\label{ux7b2c206ux7ae0ux79efux5206ux65b9ux7a0bux7684ux6570ux503cux65b9ux6cd5}

积分方程的数值求解是科学计算的重要课题。与微分方程相比，积分方程离散化后通常产生满矩阵，计算成本更高，但往往具有更好的稳定性。本章讨论三类核心数值方法：投影方法（Galerkin
法和配置法）将无穷维问题投影到有限维子空间；Nyström
方法直接用数值求积公式离散积分算子；退化核逼近则通过核的低秩近似将积分方程约化为线性代数问题。此外，第一类积分方程的不适定性需要正则化技术的特殊处理。

\subsubsection{206.1 投影方法：Galerkin
与配置法}\label{ux6295ux5f71ux65b9ux6cd5galerkin-ux4e0eux914dux7f6eux6cd5}

\textbf{定义 206.1}（Galerkin 方法）：对积分方程
\((I - \lambda T) \varphi = f\)，\textbf{Galerkin 方法}在有限维子空间
\(X_n \subset X\) 中求近似解 \(\varphi_n \in X_n\)，满足

\[\langle (I - \lambda T) \varphi_n, \psi \rangle = \langle f, \psi \rangle \quad (\forall \psi \in X_n)\]

若取 \(X_n = \operatorname{span}\{\phi_1, \ldots, \phi_n\}\)，设
\(\varphi_n = \sum_{j=1}^n c_j \phi_j\)，则得到线性方程组

\[\sum_{j=1}^n c_j \left[\langle \phi_j, \phi_i \rangle - \lambda \langle T\phi_j, \phi_i \rangle\right] = \langle f, \phi_i \rangle\]

\textbf{定义 206.2}（配置法，Collocation
Method）：\textbf{配置法}在选定的配置点 \(x_1, \ldots, x_n\)
上强制方程精确成立：

\[\varphi_n(x_i) - \lambda \int_a^b K(x_i, y) \varphi_n(y) \, dy = f(x_i) \quad (i = 1, \ldots, n)\]

在有限维近似空间 \(X_n\) 中求解。

\textbf{定理 206.1}（投影方法的收敛性）：若 \(X_n\) 在 \(X\)
中稠密（\(n \to \infty\)），\(T\) 是紧算子，且 \(I - \lambda T\)
可逆，则对充分大的 \(n\)，Galerkin 近似方程可解且
\(\varphi_n \to \varphi\)（真实解）。

\subsubsection{206.2 Nyström 方法}\label{nystruxf6m-ux65b9ux6cd5}

\textbf{定义 206.3}（Nyström 方法）：\textbf{Nyström
方法}（求积法）直接用求积法则近似积分算子。设求积节点 \(\{y_j\}\) 和权
\(\{w_j\}\)，则

\[\int_a^b K(x, y) \varphi(y) \, dy \approx \sum_{j=1}^n w_j K(x, y_j) \varphi(y_j)\]

原始积分方程离散化为

\[\varphi_n(x_i) - \lambda \sum_{j=1}^n w_j K(x_i, y_j) \varphi_n(y_j) = f(x_i)\]

得到关于 \(\varphi_n(y_j)\) 的 \(n\) 维线性方程组。

\textbf{定理 206.2}（Nyström
方法的收敛性）：若求积法则收敛（\(\sum_{j=1}^n w_j g(y_j) \to \int_a^b g(y) dy\)
对所有连续 \(g\)），则 Nyström 方法得到的近似解一致收敛到真实解。

\textbf{算法 206.1}（Nyström 插值）：解出 \(\varphi_n(y_j)\) 后，Nyström
插值公式给出任意点 \(x\) 处的近似解：

\[\varphi_n(x) = f(x) + \lambda \sum_{j=1}^n w_j K(x, y_j) \varphi_n(y_j)\]

\subsubsection{206.3 退化核逼近}\label{ux9000ux5316ux6838ux903cux8fd1}

\textbf{定义 206.4}（退化核逼近）：将核 \(K(x, y)\) 近似为退化核

\[K(x, y) \approx \sum_{j=1}^n a_j(x) b_j(y)\]

则积分方程化简为有限维线性代数系统。常用的退化核逼近包括： -
\textbf{截断奇异值分解}（TSVD）：取 \(K\) 的奇异值分解的前 \(n\) 项 -
\textbf{插值逼近}：\(K(x, y) \approx \sum_{j=1}^n K(x, y_j) L_j(y)\)（\(L_j\)
是 Lagrange 基函数）

\textbf{定理 206.3}（退化核逼近的误差估计）：若
\(\|K - K_n\| \to 0\)（算子范数），则退化核逼近的解
\(\varphi_n \to \varphi\)，且
\(\|\varphi - \varphi_n\| \leq C \|K - K_n\|\)。

\subsubsection{206.4
第一类积分方程的正则化}\label{ux7b2cux4e00ux7c7bux79efux5206ux65b9ux7a0bux7684ux6b63ux5219ux5316}

\textbf{定义 206.5}（不适定问题与正则化）：第一类积分方程
\(T\varphi = f\) 是不适定的（ill-posed）：解对 \(f\)
的微小扰动不连续依赖。\textbf{Tikhonov 正则化}求解

\[\min_\varphi \|T\varphi - f\|^2 + \alpha \|\varphi\|^2\]

其中 \(\alpha > 0\) 是正则化参数。解为
\(\varphi_\alpha = (T^*T + \alpha I)^{-1} T^* f\)。

\textbf{定理 206.4}（Tikhonov 正则化的收敛性）：若真解
\(\varphi^\dagger\) 满足源条件，则当 \(\alpha \to 0\) 且 \(\alpha\)
与噪声水平适当匹配时，\(\varphi_\alpha\) 收敛到 \(\varphi^\dagger\)。

\textbf{定义 206.6}（迭代正则化）：Landweber 迭代
\(\varphi_{k+1} = \varphi_k + \omega T^*(f - T\varphi_k)\)
是另一种正则化方法，迭代次数 \(k\) 充当正则化参数。

\bigskip\hrule\bigskip

\emph{本卷建立了积分方程的核心理论框架：Fredholm 积分方程（Fredholm
行列式、Fredholm 择一定理、预解核）是泛函分析的历史起源，将线性代数中的
Cramer 法则推广到无穷维；Volterra 积分方程（逐次逼近、Neumann 级数、与
ODE 的等价性）具有''无反馈''特性，解对所有参数存在；奇异积分方程（Cauchy
核、Plemelj 公式、Riemann-Hilbert
问题）在弹性力学和流体力学中有重要应用；Wiener-Hopf
积分方程（因子分解、Fourier
变换、半无穷区间）是解决衍射问题和排队论的经典工具；数值方法（Galerkin、Nyström、配置法、正则化）为积分方程的实用求解提供了可靠算法。积分方程理论是泛函分析、微分方程和数值分析的交汇点。}

\bigskip\hrule\bigskip

\emph{本卷在已有内容基础上，共 5 章。}

\newpage

\section{卷五十一：代数 K
理论}\label{ux5377ux4e94ux5341ux4e00ux4ee3ux6570-k-ux7406ux8bba}

\begin{quote}
代数 K 理论是连接代数、拓扑、数论和代数几何的深层桥梁。从 Grothendieck
的 K0 构造到 Quillen 的高阶 K 理论，再到 Voevodsky 的母题上同调，K
理论揭示了代数结构与拓扑不变量的深刻联系。本卷系统建立低阶 K
理论（K0、K1、K2）的基本框架、Milnor K 理论、代数 K
理论与母题上同调的关联，以及 K 理论在数论、几何和物理学中的应用。
\end{quote}

\begin{quote}
\textbf{卷序说明}：本章节号（Ch
196-201）反映其在全书逻辑链中的位置。物理卷序上本卷位于V51，作为代数K理论方向的专题补充卷。
\end{quote}

\bigskip\hrule\bigskip

\subsection{K0、K1、K2：低阶K理论导引}\label{k0k1k2ux4f4eux9636kux7406ux8bbaux5bfcux5f15}

\(K_0\) 是代数 K 理论的起点，由 Grothendieck 在 1957 年研究 Riemann-Roch
定理时引入。其基本思想是将''分裂正合列的可加性''形式化为 Abel 群。

\subsubsection{196.1 Grothendieck
群的定义}\label{grothendieck-ux7fa4ux7684ux5b9aux4e49}

\textbf{定义 196.1}（环的 \(K_0\)）：设 \(R\) 为含幺环。令
\(\mathcal{P}(R)\) 为有限生成投射左 \(R\)-模的同构类全体。在
\(\mathcal{P}(R)\) 上定义加法运算 \([P] + [Q] = [P \oplus Q]\)。则
\textbf{\(K_0(R)\)} 是上述交换么半群 \((\mathcal{P}(R), \oplus)\) 的
\textbf{Grothendieck 群}（群完备化）：由形式差 \([P] - [Q]\)
生成，模去关系 \([P] - [Q] = [P'] - [Q'] \iff \exists\) 有限生成投射模
\(R\) 使得 \(P \oplus Q' \oplus R \cong P' \oplus Q \oplus R\)。

等价地，\(K_0(R)\) 由生成元 \([P]\)（\(P\) 为有限生成投射
\(R\)-模）和关系 \([P \oplus Q] = [P] + [Q]\) 生成的 Abel 群。

\textbf{例 196.1}（域和除环的 \(K_0\)）：若 \(R\)
是域（或除环），则有限生成投射模等价于有限维向量空间，分类由维数给出。因此
\(K_0(R) \cong \mathbb{Z}\)，由 \([R]\) 生成。

\textbf{例 196.2}（Dedekind 整环的 \(K_0\)）：若 \(R\) 为 Dedekind
整环（如代数整数环），则
\[K_0(R) \cong \mathbb{Z} \oplus \operatorname{Cl}(R)\] 其中
\(\operatorname{Cl}(R)\) 是 \(R\) 的理想类群。\(\mathbb{Z}\)
分量来自秩（rank），\(\operatorname{Cl}(R)\)
分量反映了非自由投射模（即非主理想）的信息。这是 \(K_0\)
与数论联系的第一个实例。

\subsubsection{\texorpdfstring{196.2 \(K_0\)
的基本性质}{196.2 K\_0 的基本性质}}\label{k_0-ux7684ux57faux672cux6027ux8d28}

\textbf{命题 196.1}（\(K_0\) 的函子性）：\(R \mapsto K_0(R)\)
是从含幺环范畴到 Abel 群范畴的协变函子（对任意环同态
\(f: R \to S\)，标量限制 \(S \otimes_R -\) 诱导
\(f_*: K_0(R) \to K_0(S)\)）。

\textbf{命题 196.2}（Morita 不变性）：若环 \(R\) 和 \(S\) 是 Morita
等价的（即它们的模范畴等价），则
\(K_0(R) \cong K_0(S)\)。特别地，\(K_0(R) \cong K_0(M_n(R))\)。

\textbf{命题 196.3}（幂等元描述）：\(K_0(R)\) 也可等价地定义为
\(M_\infty(R) = \bigcup_{n=1}^\infty M_n(R)\) 中幂等元的等价类在直和下的
Grothendieck 群。投射模 \(P\) 可表为 \(P \cong e R^n\)（\(e \in M_n(R)\)
为幂等元），因此 \(K_0(R)\) 的每个元素可写为 \([e] - [f]\)（\(e, f\)
为幂等矩阵）。

\subsubsection{\texorpdfstring{196.3 拓扑空间的
\(K^0\)}{196.3 拓扑空间的 K\^{}0}}\label{ux62d3ux6251ux7a7aux95f4ux7684-k0}

\textbf{定义 196.2}（拓扑 \(K^0\)）：紧 Hausdorff 空间 \(X\)
上复向量丛的同构类在直和下构成交换么半群。其 Grothendieck 群记为
\(K^0(X)\)。这就是 \textbf{拓扑 K 理论} 的起点（Atiyah-Hirzebruch,
1959）。

\textbf{定理 196.1}（Swan 定理）：\(X\) 紧 Hausdorff 时，\(X\)
上复向量丛的范畴等价于 \(C(X)\) 上有限生成投射模的范畴。因此
\(K_0(C(X)) \cong K^0(X)\)------代数 K 理论与拓扑 K 理论在此处交汇。

\bigskip\hrule\bigskip

\(K_1\) 由 Bass 和 Schanuel 在 1964 年引入，其定义源于一般线性群
\(GL_n(R)\) 的白头化（Whitehead
lemma）------可逆矩阵的''初等变换''生成整个换位子群。

\subsubsection{197.1 初等矩阵与 Whitehead
引理}\label{ux521dux7b49ux77e9ux9635ux4e0e-whitehead-ux5f15ux7406}

\textbf{定义 197.1}（初等矩阵）：对 \(i \neq j\) 和
\(r \in R\)，\textbf{初等矩阵} \(e_{ij}(r)\) 为单位矩阵加上第 \((i, j)\)
个非对角线位置上的 \(r\)。令 \(E_n(R)\) 为所有 \(n \times n\) 初等矩阵
\(e_{ij}(r)\) 生成的 \(GL_n(R)\) 的子群。

\textbf{引理 197.1}（初等矩阵的换位关系）：
\[[e_{ij}(r), e_{jk}(s)] = e_{ik}(rs) \quad (i \neq k)\]
\[[e_{ij}(r), e_{k\ell}(s)] = I_n \quad (j \neq k, i \neq \ell)\]

\textbf{定理 197.1}（Whitehead 引理）：对任意 \(n \geq 3\)，
\[E_n(R) = [GL_n(R), GL_n(R)] = [E_n(R), E_n(R)]\]
即初等子群同时是一般线性群和自身的换位子群。

\subsubsection{\texorpdfstring{197.2 \(K_1\)
的定义与基本性质}{197.2 K\_1 的定义与基本性质}}\label{k_1-ux7684ux5b9aux4e49ux4e0eux57faux672cux6027ux8d28}

\textbf{定义 197.2}（环的 \(K_1\)）：令
\(GL(R) = \bigcup_{n=1}^\infty GL_n(R)\)（通过标准嵌入
\(A \mapsto \begin{pmatrix} A & 0 \\ 0 & 1 \end{pmatrix}\)
取正向极限）。令 \(E(R) = \bigcup_{n=1}^\infty E_n(R)\)。定义
\[K_1(R) = GL(R) / E(R) = GL(R)^{\text{ab}}\] 即一般线性群的
\textbf{Abel 化}（\(E(R) = [GL(R), GL(R)]\) 为换位子群）。

\textbf{例 197.1}（域的 \(K_1\)）：对域 \(F\)，由 Dieudonne
行列式，\(K_1(F) \cong F^\times\)。特别地，\(K_1(\mathbb{R}) \cong \mathbb{R}^\times\)，\(K_1(\mathbb{C}) \cong \mathbb{C}^\times\)。

\textbf{例 197.2}（整数的
\(K_1\)）：\(K_1(\mathbb{Z}) \cong \{\pm 1\} \cong \mathbb{Z}/2\mathbb{Z}\)。

\textbf{命题 197.1}（行列式映射）：对交换环
\(R\)，存在行列式诱导的满同态 \(\det: K_1(R) \to R^\times\)，其核记为
\(SK_1(R)\)。一般地，\(K_1(R) \cong R^\times \oplus SK_1(R)\)
不一定成立，但存在正合列
\(1 \to SK_1(R) \to K_1(R) \to R^\times \to 1\)，该正合列在环 \(R\)
为局部环时分裂。

\textbf{命题 197.2}（Morita 不变性）：\(K_1\) 也是 Morita
不变的：\(K_1(R) \cong K_1(M_n(R))\)。

\subsubsection{\texorpdfstring{197.3 拓扑空间的
\(K^1\)}{197.3 拓扑空间的 K\^{}1}}\label{ux62d3ux6251ux7a7aux95f4ux7684-k1}

\textbf{定义 197.3}（拓扑 \(K^1\)）：\(K^1(X)\) 定义为 \(X\) 上的
\(GL_n(\mathbb{C})\)-值连续函数的同伦类的正向极限，等价于悬挂的约化
\(K^0\)：\(K^1(X) = \tilde{K}^0(\Sigma X)\)。

\bigskip\hrule\bigskip

\(K_2\) 由 Milnor 在 1970 年引入，主要动机是在 Mennicke 符号和 Moore
关于泛中心扩张的理论工作基础上，完成低阶 K 群的正合列。

\subsubsection{198.1 Steinberg
群与泛中心扩张}\label{steinberg-ux7fa4ux4e0eux6cdbux4e2dux5fc3ux6269ux5f20}

\textbf{定义 198.1}（Steinberg 群）：设 \(R\)
为含幺环（\(n \geq 3\)）。\textbf{Steinberg 群}
\(\operatorname{St}_n(R)\) 由生成元
\(x_{ij}(r)\)（\(1 \leq i \neq j \leq n, r \in R\)）定义，满足 Steinberg
关系： 1. \(x_{ij}(r) x_{ij}(s) = x_{ij}(r + s)\) 2.
\([x_{ij}(r), x_{jk}(s)] = x_{ik}(rs)\)（\(i \neq k\)） 3.
\([x_{ij}(r), x_{k\ell}(s)] = 1\)（\(j \neq k, i \neq \ell\)）

存在自然的满同态
\(\phi_n: \operatorname{St}_n(R) \twoheadrightarrow E_n(R)\)，将
\(x_{ij}(r)\) 映为 \(e_{ij}(r)\)。其核是 \(K_2(R)\) 与 \(E_n(R)\)
的泛中心扩张之间关系的关键。

\textbf{定理 198.1}（Steinberg 群的泛中心扩张，Steinberg 1962）：对
\(n \geq 3\)，\(\operatorname{St}_n(R)\) 是 \(E_n(R)\) 的
\textbf{泛中心扩张}------\(\ker(\phi_n)\) 包含在
\(\operatorname{St}_n(R)\) 的中心中，且任何 \(E_n(R)\) 的中心扩张都通过
\(\operatorname{St}_n(R)\) 分解。

\subsubsection{\texorpdfstring{198.2 \(K_2\)
的定义}{198.2 K\_2 的定义}}\label{k_2-ux7684ux5b9aux4e49}

\textbf{定义 198.2}（环的 \(K_2\)）：对 \(n \geq 3\)，定义
\[K_2(R) = \ker(\operatorname{St}_n(R) \to E_n(R))\] 该定义与
\(n \geq 3\) 的选取无关（通过稳定性定理）。\(K_2(R)\) 是 Abel 群且包含在
\(\operatorname{St}(R)\) 的中心中。

\textbf{定理 198.2}（Bass 正合列）：\(K_2\) 自然地出现在正合列中：
\[0 \to K_2(R) \to \operatorname{St}(R) \to E(R) \to 1\]
该正合列为泛中心扩张。

\textbf{例 198.1}（域上的 \(K_2\)------Matsumoto 定理, 1969）：对域
\(F\)，
\[K_2(F) \cong F^\times \otimes_{\mathbb{Z}} F^\times / \langle a \otimes (1-a) : a \in F \setminus \{0, 1\} \rangle\]
这恰好是 Milnor \(K_2^M(F)\) 的定义------Milnor \(K_2\) 与 Quillen
\(K_2\) 在域上一致。生成元 \(\{a, b\}\)（\(a, b \in F^\times\)）称为
\textbf{Steinberg 符号}，满足双乘性 \(\{a b, c\} = \{a, c\}\{b, c\}\) 和
Steinberg 关系 \(\{a, 1-a\} = 1\)（\(a \neq 0, 1\)）。

\textbf{例 198.2}（有限域的 \(K_2\)---Quillen）：若 \(\mathbb{F}_q\) 为
\(q\) 元有限域，则 \(K_2(\mathbb{F}_q) = 0\)。

\subsubsection{198.3 Milnor K
理论简介}\label{milnor-k-ux7406ux8bbaux7b80ux4ecb}

\textbf{定义 198.3}（Milnor \(K\) 群）：域 \(F\) 的 \textbf{Milnor \(K\)
群} \(K_n^M(F)\) 是分次环 \(K_*^M(F)\) 的 \(n\) 次部分，定义为
\[K_n^M(F) = (F^\times)^{\otimes_{\mathbb{Z}} n} / \langle a_1 \otimes \cdots \otimes a_n : a_i + a_{i+1} = 1 \text{ 对某 } i \rangle\]

关键性质： - \(K_0^M(F) = \mathbb{Z}\)，\(K_1^M(F) = F^\times\) -
\(K_2^M(F)\) 即上述 Matsumoto 定理的表述 - 存在符号映射
\(\{a_1, \ldots, a_n\} \in K_n^M(F)\)，满足多重线性和 Steinberg 关系

\textbf{定理 198.3}（Milnor 猜想 / Voevodsky 定理）：对特征不为 2 的域
\(F\)， \[K_n^M(F) / 2 \cong H^n_{\text{Gal}}(F, \mu_2^{\otimes n})\]
这是 Bloch-Kato 猜想的 \(\ell = 2\) 情形，由 Voevodsky 于 1996 年证明。

\bigskip\hrule\bigskip

\subsection{\texorpdfstring{第199章：Milnor \(K\) 理论与 Bloch
群}{第199章：Milnor K 理论与 Bloch 群}}\label{ux7b2c199ux7ae0milnor-k-ux7406ux8bbaux4e0e-bloch-ux7fa4}

Milnor K 理论由 Milnor 在 1970 年引入，专注于域上的 K
理论，并与二次型、Galois 上同调和 Bloch 群有深刻联系。

\subsubsection{\texorpdfstring{199.1 Milnor \(K\)
群}{199.1 Milnor K 群}}\label{milnor-k-ux7fa4}

\textbf{定义 199.1}（Milnor \(K\) 理论，域）：域 \(F\) 的 \textbf{Milnor
\(K\) 群} \(K_n^M(F)\) 是分次环 \(K_*^M(F)\) 的 \(n\) 次部分，由

\[K_*^M(F) = T(F^\times) / \langle a \otimes (1-a) : a \in F \setminus \{0, 1\} \rangle\]

定义，其中 \(T(F^\times)\) 是乘法群 \(F^\times\)
的张量代数。\(\{a_1, \ldots, a_n\}\) 记
\(a_1 \otimes \cdots \otimes a_n\) 在 \(K_n^M(F)\) 中的像（Milnor
符号）。

\textbf{命题 199.1}（Milnor \(K\) 群的低阶）： -
\(K_0^M(F) = \mathbb{Z}\) - \(K_1^M(F) = F^\times\) - \(K_2^M(F)\) 与
Quillen \(K_2(F)\) 一致（Matsumoto 定理）

\textbf{定理 199.1}（Milnor \(K\) 理论与 Galois 上同调，Bloch-Kato 猜想
/ Voevodsky 定理 2011）：对域 \(F\) 和素数
\(\ell \neq \operatorname{char}(F)\)，存在自然同构

\[K_n^M(F) / \ell \cong H^n_{\text{Gal}}(F, \mu_\ell^{\otimes n})\]

（Galois 上同调与 Milnor \(K\) 理论的联系）。这是 Milnor
猜想（\(\ell = 2\)）和 Bloch-Kato 猜想（一般 \(\ell\)）的融合，由
Voevodsky 证明（2011 年 Fields 奖）。

\subsubsection{199.2 Bloch
群与稀释子}\label{bloch-ux7fa4ux4e0eux7a00ux91caux5b50}

\textbf{定义 199.2}（Bloch 群）：域 \(F\) 的 \textbf{Bloch 群}
\(\mathcal{B}(F)\) 是由 \([x]\)（\(x \in F \setminus \{0, 1\}\)）生成的
Abel 群，模去关系

\[[x] - [y] + \left[\frac{y}{x}\right] - \left[\frac{1 - x^{-1}}{1 - y^{-1}}\right] + \left[\frac{1 - x}{1 - y}\right] = 0\]

（五项关系，five-term relation）。

\textbf{定理 199.2}（Bloch-Suslin 定理，1986）：对无限域
\(F\)，\(K_3(F)_{\text{ind}} \cong \mathcal{B}(F)\)（\(K_3\)
的不分解部分）。Bloch 群是 \(K_3\) 的精确代数描述。

\textbf{定义 199.3}（稀释子，Dilogarithm）：\textbf{Bloch-Wigner 稀释子}
\(D: \mathbb{C} \setminus \{0, 1\} \to \mathbb{R}\) 定义为

\[D(z) = \operatorname{Im}(\operatorname{Li}_2(z)) + \log|z| \cdot \arg(1 - z)\]

\(D(z)\) 是单值的且在 \(\mathbb{C}P^1\) 上连续延拓。\(D(z)\)
满足五项关系，因此诱导了 \(\mathcal{B}(\mathbb{C}) \to \mathbb{R}\)
的同态。

\textbf{定理 199.3}（Bloch 群与双对数恒等式）：Bloch
群的结构由稀释子（dilogarithm）和经典的双对数恒等式（Euler, Abel,
Kummer）控制。\(\mathcal{B}(\mathbb{C})\) 的秩为不可数无限。

\subsubsection{199.3 调节子映射}\label{ux8c03ux8282ux5b50ux6620ux5c04}

\textbf{定义 199.4}（Borel 调节子）：Borel 定义了从
\(K_{2n-1}(\mathbb{C})\) 到 \(\mathbb{R}\) 的调节子映射，其像与
\(\zeta(n)\) 相关。更一般地，对数域 \(F\)，

\[K_{2n-1}(\mathcal{O}_F) \otimes \mathbb{Q} \cong \mathbb{Q}^{d_n}\]

其中 \(d_n = \dim_{\mathbb{Q}} H^1_{\text{mot}}(F, \mathbb{Q}(n))\)。

\textbf{定理 199.4}（Beilinson 猜想，1984）：特殊 \(L\) 值 \(L^*(M, n)\)
与母题 \(M\) 的调节子映射的像（在 \(K\) 理论中）有联系。这是代数 \(K\)
理论与数论 \(L\) 函数之间最深刻的猜想。

\bigskip\hrule\bigskip

\subsection{\texorpdfstring{第200章：代数 \(K\)
理论与母题}{第200章：代数 K 理论与母题}}\label{ux7b2c200ux7ae0ux4ee3ux6570-k-ux7406ux8bbaux4e0eux6bcdux9898}

母题（motives）是 Grothendieck 在 1960
年代提出的概念，旨在统一代数簇的所有上同调理论。Voevodsky
建立了母题上同调与代数 \(K\) 理论的联系。

\subsubsection{200.1 母题范畴}\label{ux6bcdux9898ux8303ux7574}

\textbf{定义
200.1}（母题，Grothendieck）：\textbf{母题}（motive）是代数簇的''普遍上同调''对象。对光滑射影代数簇
\(X\)，其母题 \(h(X)\) 是 Abel 范畴
\(\mathcal{M}(k)\)（母题范畴）中的对象，满足对所有 Weil 上同调理论
\(H^*\) 存在实现函子 \(r_H: \mathcal{M}(k) \to \text{Vec}\) 使得
\(r_H(h(X)) = H^*(X)\)。

\textbf{定义 200.2}（Chow 母题）：\textbf{Chow 母题}范畴
\(\mathcal{M}_{\text{rat}}(k)\)
由光滑射影代数簇的对应（correspondences）构造。对象是
\((X, p, n)\)（\(X\) 是簇，\(p\) 是幂等对应，\(n\) 是 Tate
扭转）。其中： - \textbf{幂等对应}
\(p \in \operatorname{CH}^{\dim X}(X \times X)_{\mathbb{Q}}\) 满足
\(p \circ p = p\)（对应合成的意义下），作用相当于在 \(X\)
的上同调中取出某个直和分量------形式上 \(p\) 将母题 \(h(X)\)
投影到其''子母题''。 - \textbf{Tate 扭转} \(n \in \mathbb{Z}\)：对象
\((X, p, n)\) 的''母题上同调''相对于 \((X, p, 0)\) 有一个维数位移。Tate
母题 \(\mathbb{Q}(n)\) 是 \(\operatorname{Spec}(k)\) 的母题经过 \(n\) 次
Tate 扭转，其 Betti 实现给出
\((2\pi i)^n \mathbb{Q}\)，反映了母题范畴中的''维数权重''结构。Tate
扭转是定义母题范畴中 Hom 群的分次结构以及联系母题上同调与代数 \(K\)
理论谱序列的关键成分。

\textbf{定理 200.1}（Jannsen 定理，1992）：Chow 母题范畴是半单 Abel
范畴当且仅当数值等价与有理等价一致（这是 Grothendieck 标准猜想之一）。

\subsubsection{200.2 Voevodsky
的母题上同调}\label{voevodsky-ux7684ux6bcdux9898ux4e0aux540cux8c03}

\textbf{定义 200.3}（母题上同调，Voevodsky 2000）：光滑代数簇 \(X\)
的\textbf{母题上同调} \(H^{p,q}_{\text{mot}}(X, \mathbb{Z})\) 定义为
\(X\) 在母题稳定同伦范畴 \(\mathcal{SH}(k)\) 中的母题 Eilenberg-MacLane
谱 \(H\mathbb{Z}\) 的上同调。

\textbf{定理 200.2}（母题上同调与代数 \(K\) 理论，Voevodsky
2003）：对光滑代数簇 \(X\)，

\[H^{p,q}_{\text{mot}}(X, \mathbb{Z}) \cong \operatorname{CH}^q(X, 2q-p)\]

（Bloch 高阶 Chow 群）。母题上同调给出了 Bloch 高阶 Chow
群的上同调解释。

\textbf{定理
200.3}（母题谱序列，Friedlander-Suslin-Voevodsky）：存在从母题上同调到代数
\(K\) 理论的谱序列：

\[E_2^{p,q} = H^{p-q,-q}_{\text{mot}}(X, \mathbb{Z}) \Rightarrow K_{-p-q}(X)\]

这统一了代数 \(K\) 理论和母题上同调。

\subsubsection{200.3 Voevodsky 的 Fields
奖工作}\label{voevodsky-ux7684-fields-ux5956ux5de5ux4f5c}

\textbf{定理 200.4}（Milnor 猜想 / Voevodsky 定理，1996-2003）：对域
\(k\)（\(\operatorname{char}(k) \neq 2\)），Milnor \(K\) 理论模 2 与
étale 上同调之间存在自然同构：

\[K_n^M(k) / 2 \cong H^n_{\text{ét}}(k, \mathbb{Z}/2\mathbb{Z})\]

\textbf{定理 200.5}（Bloch-Kato 猜想，已证明，Voevodsky-Rost
2011）：该猜想等价于定理 199.1（Bloch-Kato 猜想的 étale
上同调版本），两者表述不同但实质相同。详见第 199 章。

\subsubsection{\texorpdfstring{200.4 代数 \(K\) 理论与 \(L\)
函数}{200.4 代数 K 理论与 L 函数}}\label{ux4ee3ux6570-k-ux7406ux8bbaux4e0e-l-ux51fdux6570}

\textbf{定理 200.6}（Lichtenbaum 猜想，Beilinson
猜想）：对光滑射影代数簇 \(X\) 在整数环 \(\mathcal{O}_F\)
上的模型，\(K\) 群的秩由 Hasse-Weil \(L\) 函数在整点处的零点阶数给出：

\[\dim_{\mathbb{Q}} K_{2n-1}(X) \otimes \mathbb{Q} = \operatorname{ord}_{s=n} L(h^{2n-1}(X), s)\]

这推广了 Borel 关于数域 \(K\) 群的结果。

\bigskip\hrule\bigskip

\subsection{\texorpdfstring{第201章：\(K\)
理论的其他方向与应用}{第201章：K 理论的其他方向与应用}}\label{ux7b2c201ux7ae0k-ux7406ux8bbaux7684ux5176ux4ed6ux65b9ux5411ux4e0eux5e94ux7528}

代数 \(K\) 理论在拓扑、几何、算子代数和物理学中有广泛的应用。

\subsubsection{\texorpdfstring{201.1 拓扑 \(K\)
理论与指标定理}{201.1 拓扑 K 理论与指标定理}}\label{ux62d3ux6251-k-ux7406ux8bbaux4e0eux6307ux6807ux5b9aux7406}

\textbf{定义 201.1}（拓扑 \(K\) 理论，Atiyah-Hirzebruch 1959）：紧
Hausdorff 空间 \(X\) 的拓扑 \(K\) 理论 \(K^0(X)\) 是 \(X\) 上复向量丛的
Grothendieck 群。\(K^{-n}(X) = \tilde{K}^0(S^n X)\)（Bott
周期性：\(K^{-n-2}(X) \cong K^{-n}(X)\)）。

\textbf{定理 201.1}（Atiyah-Singer 指标定理，1963）：椭圆算子 \(D\)
的解析指标等于拓扑指标：

\[\operatorname{Index}(D) = \int_M \operatorname{Ch}(\sigma(D)) \wedge \operatorname{Td}(M)\]

其中 \(\operatorname{Ch}\) 是 Chern 特征标（\(K\)
理论到上同调的映射），\(\operatorname{Td}\) 是 Todd 类。

\textbf{定理 201.2}（Swan 定理，1962）：\(C(X)\)
上有限生成投射模的范畴等价于 \(X\) 上复向量丛的范畴。因此
\(K_0(C(X)) \cong K^0(X)\)（代数 \(K\) 理论与拓扑 \(K\) 理论的联系）。

\subsubsection{\texorpdfstring{201.2 C*-代数 \(K\)
理论}{201.2 C*-代数 K 理论}}\label{c-ux4ee3ux6570-k-ux7406ux8bba}

\textbf{定义 201.2}（C\emph{-代数 \(K\) 理论）：C}-代数 \(A\) 的
\(K_0(A)\) 定义为 \(A\) 中投影的 Murray-von Neumann 等价类的
Grothendieck 群。\(K_1(A)\) 定义为 \(A\) 中酉元的连通分支群。

\textbf{定理 201.3}（六项正合列，\(K\) 理论的基本正合列）：对
C*-代数的短正合列 \(0 \to J \to A \to A/J \to 0\)，存在 \(K\)
理论的六项正合列：

\[\begin{CD}
K_0(J) @>>> K_0(A) @>>> K_0(A/J) \\
@AAA @. @VVV \\
K_1(A/J) @<<< K_1(A) @<<< K_1(J)
\end{CD}\]

\textbf{定理 201.4}（Bott
周期性，C\emph{-代数情形）：\(K_0(SA) \cong K_1(A)\)，\(K_1(SA) \cong K_0(A)\)（其中
\(SA = C_0(\mathbb{R}) \otimes A\) 是悬挂）。这给出了 C}-代数 \(K\)
理论的 2-周期性。

\subsubsection{\texorpdfstring{201.3 \(K\)
理论在物理学中的应用}{201.3 K 理论在物理学中的应用}}\label{k-ux7406ux8bbaux5728ux7269ux7406ux5b66ux4e2dux7684ux5e94ux7528}

\textbf{定义 201.3}（拓扑绝缘体与 \(K\) 理论）：拓扑绝缘体（topological
insulator）的拓扑相由 \(K\) 理论分类。时间反演对称的拓扑绝缘体由 \(KR\)
理论（实 \(K\) 理论带对合）分类。

\textbf{定理 201.5}（Kitaev 的周期表，2009）：自由费米子系统的拓扑相由
\(K\) 理论的 Bott 周期间接分类（``十重分类''）。\(K\)
理论是理解拓扑序和拓扑量子计算的数学基础。

\textbf{定义 201.4}（\(K\) 理论与 D-膜）：弦论中 D-膜的荷由 \(K\)
理论分类（而非整数上同调）。\(K^0(X)\) 分类 IIB 理论中
D-膜的荷，\(K^1(X)\) 分类 IIA 理论中 D-膜的荷。

\subsubsection{\texorpdfstring{201.4 Waldhausen \(K\) 理论与代数 \(K\)
理论空间的构造}{201.4 Waldhausen K 理论与代数 K 理论空间的构造}}\label{waldhausen-k-ux7406ux8bbaux4e0eux4ee3ux6570-k-ux7406ux8bbaux7a7aux95f4ux7684ux6784ux9020}

\textbf{定义 201.5}（Waldhausen \(K\) 理论，1985）：Waldhausen 将
Quillen 的 Q-构造推广到带有余纤维化（cofibration）和弱等价（weak
equivalence）的范畴（Waldhausen 范畴）。Waldhausen \(K\) 理论是代数
\(K\) 理论、代数几何 \(K\) 理论和等变稳定同伦论中最重要的统一框架。

\textbf{定理 201.6}（Waldhausen 的加性定理和纤维化定理）：Waldhausen
\(K\) 理论满足对 Waldhausen
范畴的分裂正合列和纤维化序列的加性行为，推广了 Quillen 的局部化定理。

\bigskip\hrule\bigskip

\emph{本卷建立了代数 \(K\) 理论的完整框架：\(K_0\)（投射模的
Grothendieck
群）建立了将正合列转化为群关系的代数框架；\(K_1\)（一般线性群的 Abel
化）通过 Whitehead
引理揭示了初等矩阵与换位子的深层联系；\(K_2\)（Steinberg
群与泛中心扩张）通过 Matsumoto 定理将域上的 \(K_2\) 与 Milnor
符号连接；Milnor \(K\) 理论与 Bloch 群（符号、稀释子、调节子）将 \(K\)
理论与数论 \(L\) 函数连接；母题上同调（Voevodsky
的母题稳定同伦论、Bloch-Kato 猜想）是 21 世纪数学的里程碑；拓扑 \(K\)
理论、C}-代数 \(K\) 理论和物理学应用展示了 \(K\)
理论的跨领域影响力。代数 \(K\)
理论是当代数学中连接代数、拓扑、几何、数论和物理的统一语言。*

\bigskip\hrule\bigskip

\emph{本卷在已有内容基础上，以K0/K1/K2低阶K理论导引开篇，再续以 3 章（Ch
199-201），共 4 个单元。}

\newpage

\section{卷五十二：逼近论}\label{ux5377ux4e94ux5341ux4e8cux903cux8fd1ux8bba}

\begin{quote}
逼近论研究用简单函数（多项式、有理函数、样条、小波等）逼近复杂函数的理论和方法。从
Weierstrass 逼近定理到 Chebyshev 最佳逼近，从有理逼近与 Pade
逼近到样条与小波，逼近论为数值计算、信号处理和机器学习提供了坚实的数学基础。本卷系统建立最佳逼近的存在性与唯一性理论、多项式逼近的
Jackson 与 Bernstein
定理、有理逼近、样条插值与光滑性、小波的多分辨分析框架，以及径向基函数与高维逼近方法。
\end{quote}

\begin{quote}
\textbf{卷序说明}：本章节号（Ch
208-211）反映其在全书逻辑链中的位置。物理卷序上本卷位于V52，作为逼近论方向的专题补充卷。
\end{quote}

\bigskip\hrule\bigskip

\subsection{第208章：Weierstrass 逼近定理与 Jackson
定理}\label{ux7b2c208ux7ae0weierstrass-ux903cux8fd1ux5b9aux7406ux4e0e-jackson-ux5b9aux7406}

Weierstrass
逼近定理（1885）是逼近论的奠基性结果------它断言闭区间上的任何连续函数均可被多项式一致逼近。Bernstein（1912）给出了一个构造性的概率证明，开创了用
Bernstein 多项式逼近连续函数的先河。Jackson
定理（1911）进一步建立了光滑性与逼近阶之间的定量关系。

\subsubsection{208.1 Weierstrass
第一逼近定理}\label{weierstrass-ux7b2cux4e00ux903cux8fd1ux5b9aux7406}

\textbf{定理 208.1}（Weierstrass 第一逼近定理，1885）：设
\(f \in C[a, b]\)。则对任意 \(\varepsilon > 0\)，存在多项式 \(p(x)\)
使得 \[\max_{x \in [a, b]} |f(x) - p(x)| < \varepsilon\] 即多项式空间
\(\mathcal{P}\) 在 \(C[a, b]\)（一致范数拓扑）中稠密。

\subsubsection{208.2 Bernstein
多项式与构造性证明}\label{bernstein-ux591aux9879ux5f0fux4e0eux6784ux9020ux6027ux8bc1ux660e}

\textbf{定义 208.1}（Bernstein 多项式，1912）：设
\(f \in C[0, 1]\)。\(f\) 的 \(n\) 次 \textbf{Bernstein 多项式} 定义为
\[B_n(f; x) = \sum_{k=0}^n f\left(\frac{k}{n}\right) \binom{n}{k} x^k (1-x)^{n-k}, \quad x \in [0, 1]\]

注意 \(B_n(f; x)\) 恰为二项分布 \(\operatorname{Bin}(n, x)\) 下
\(f(k/n)\) 的数学期望 \(\mathbb{E}_x[f(S_n / n)]\)，其中
\(S_n \sim \operatorname{Bin}(n, x)\)。这赋予 Bernstein
多项式自然的概率解释。

\textbf{定理 208.2}（Bernstein 多项式的收敛性）：对任意
\(f \in C[0, 1]\)，
\[\lim_{n \to \infty} \max_{x \in [0, 1]} |f(x) - B_n(f; x)| = 0\]

\emph{证明（基于概率论）}：由一致连续性，对任意
\(\varepsilon > 0\)，存在 \(\delta > 0\) 使得
\(|x - y| < \delta \implies |f(x) - f(y)| < \varepsilon\)。记
\(M = \|f\|_\infty\)。则 \[\begin{aligned}
|f(x) - B_n(f; x)| &= \left|\mathbb{E}_x\left[f(x) - f\left(\frac{S_n}{n}\right)\right]\right| \\
&\leq \mathbb{E}_x\left|f(x) - f\left(\frac{S_n}{n}\right)\right| \cdot \left[\mathbf{1}_{\{|S_n/n - x| < \delta\}} + \mathbf{1}_{\{|S_n/n - x| \geq \delta\}}\right] \\
&\leq \varepsilon + 2M \cdot \mathbb{P}_x\left(\left|\frac{S_n}{n} - x\right| \geq \delta\right)
\end{aligned}\]

由 Chebyshev
不等式，\(\mathbb{P}_x(|S_n/n - x| \geq \delta) \leq \frac{\operatorname{Var}(S_n/n)}{\delta^2} = \frac{x(1-x)}{n\delta^2} \leq \frac{1}{4n\delta^2}\)。因此
\[\max_{x \in [0, 1]} |f(x) - B_n(f; x)| \leq \varepsilon + \frac{M}{2n\delta^2}\]

取 \(n > M / (2\varepsilon \delta^2)\) 即得
\(\|f - B_n(f)\|_\infty < 2\varepsilon\)。∎

\textbf{注 208.1}（Bernstein 多项式的局限）：尽管 Bernstein
多项式给出了一致收敛的构造性证明，但其收敛速度很慢------对被逼近函数
\(f\) 即使非常光滑（如 \(f(x) = x^2\)），误差也仅为 \(O(1/n)\)（而非
\(O(1/n^2)\) 或更快）。这促使后人寻求更高效的逼近方法（如 Chebyshev
多项式插值、Jackson 定理等）。

\textbf{推论 208.1}（推广到一般区间）：对
\(f \in C[a, b]\)，通过线性变换 \(t = (x-a)/(b-a)\) 将 \([a, b]\) 映射到
\([0, 1]\)，即得 \([a, b]\) 上的 Weierstrass 逼近定理。

\subsubsection{208.3 Jackson
定理------直接定理}\label{jackson-ux5b9aux7406ux76f4ux63a5ux5b9aux7406}

Weierstrass 定理回答了''是否存在多项式逼近''的问题，Jackson
定理进一步回答了''误差随多项式次数如何衰减''的定量问题。

\textbf{定理 208.3}（Jackson 定理 / 直接定理，Jackson 1911）：设
\(f \in C^r[-1, 1]\)。则对任意正整数 \(n \geq r\)，存在次数不超过 \(n\)
的多项式 \(p_n\) 使得
\[\|f - p_n\|_\infty \leq \frac{C_r}{n^r} \cdot \omega\left(f^{(r)}; \frac{1}{n}\right)\]
其中 \(\omega(g; \delta) = \sup_{|x-y| \leq \delta} |g(x) - g(y)|\)
为连续模（modulus of continuity），\(C_r\) 为仅依赖于 \(r\) 的常数。

特别地，若 \(f^{(r)} \in \operatorname{Lip}_\alpha\)（\(\alpha\)-Holder
连续，\(0 < \alpha \leq 1\)），则
\[\|f - p_n\|_\infty \leq \frac{C_{r, \alpha}}{n^{r+\alpha}}\]

\textbf{定理 208.4}（Jackson 定理的简化形式）：设
\(f \in C^k[-1, 1]\)。则存在仅依赖于 \(k\) 的常数 \(C_k\) 使得对任意
\(n \geq k\)，
\[E_n(f) := \inf_{p \in \mathcal{P}_n} \|f - p\|_\infty \leq \frac{C_k}{n^k} \cdot \omega\left(f^{(k)}; \frac{1}{n}\right)\]

\textbf{推论 208.2}（光滑性 \(\Rightarrow\) 逼近阶）：若
\(f \in C^r[-1, 1]\) 且 \(f^{(r)}\) 满足 Lipschitz 条件，则
\(E_n(f) = O(n^{-r-1})\)。光滑性越高，多项式逼近的代数收敛阶越高。

\textbf{评注 208.1}（Jackson 定理的意义）：Jackson
定理精确刻画了''函数的可微性阶数 \(r\) 决定多项式最佳逼近的代数收敛阶
\(\geq r+1\)``。这是逼近论中\textbf{直接定理}（direct
theorem）的核心------从函数的光滑性信息推断最佳逼近的收敛速率。反方向的\textbf{逆定理}（Bernstein
/ Stechkin
逆定理）则从逼近速率反推函数的光滑性，两者共同构成完整的逼近论等价关系。

\bigskip\hrule\bigskip

\subsection{第209章：有理逼近与 Padé
逼近}\label{ux7b2c209ux7ae0ux6709ux7406ux903cux8fd1ux4e0e-paduxe9-ux903cux8fd1}

有理函数逼近通常比多项式逼近更高效，特别是在函数有极点时。Padé
逼近通过匹配 Taylor 系数构造最优有理逼近。

\subsubsection{209.1
有理逼近的基本理论}\label{ux6709ux7406ux903cux8fd1ux7684ux57faux672cux7406ux8bba}

\textbf{定义 209.1}（有理函数）：\([m/n]\) 型\textbf{有理函数}为
\(r(x) = p_m(x) / q_n(x)\)，其中
\(p_m \in \mathcal{P}_m\)，\(q_n \in \mathcal{P}_n\)（\(q_n \not\equiv 0\)）。

\textbf{定义 209.2}（有理最佳逼近）：\(f \in C[a,b]\) 在
\(\mathcal{R}_{m,n}\)（\([m/n]\)
型有理函数全体）中的\textbf{最佳一致有理逼近} \(r^*\) 满足

\[\|f - r^*\|_\infty = \inf_{r \in \mathcal{R}_{m,n}} \|f - r\|_\infty\]

\textbf{定理 209.1}（有理逼近的 Chebyshev
交替定理）：\(r^* \in \mathcal{R}_{m,n}\) 是 \(f\)
的最佳一致有理逼近当且仅当 \(f - r^*\) 在 \([a,b]\) 上至少存在
\(m+n+2 - d\) 个点的交替点集，其中
\(d = \min(m - \deg p^*, n - \deg q^*)\)。

\textbf{定理 209.2}（Newman 定理，1964）：\(|x|\) 在 \([-1, 1]\)
上的有理逼近收敛速率远快于多项式逼近：

\[\min_{r \in \mathcal{R}_{n,n}} \||x| - r(x)\|_\infty \leq 3 e^{-\sqrt{n}}\]

而多项式逼近的速率仅为
\(O(1/n)\)。这展示了有理逼近在处理非光滑函数时的巨大优势。

\subsubsection{209.2 Padé 逼近}\label{paduxe9-ux903cux8fd1}

\textbf{定义 209.3}（Padé 逼近，1892）：设
\(f(z) = \sum_{k=0}^\infty c_k z^k\) 在 \(0\) 处解析。\textbf{\([m/n]\)
型 Padé 逼近} \(r_{m,n}(z) = p_m(z) / q_n(z)\)（\(q_n(0)=1\)）满足

\[f(z) - r_{m,n}(z) = O(z^{m+n+1}) \quad (z \to 0)\]

即 Padé 逼近匹配 \(f\) 在 \(0\) 处的 Taylor 系数的前 \(m+n+1\) 项。

\textbf{定理 209.3}（Padé 逼近的构造）：Padé
逼近的系数通过求解线性方程组确定：

\[c_0 q_0 = p_0\] \[c_1 q_0 + c_0 q_1 = p_1\] \[\vdots\]

其中 \(q_j = 0\)（\(j > n\)），\(p_j = 0\)（\(j > m\)）。\(q_j\)
的确定需要求解 Hankel 矩阵方程组。

\textbf{定义 209.4}（Padé 表）：\([m/n]\) Padé 逼近的二维阵列称为
\textbf{Padé 表}。Padé 表的对角线 \([n/n]\)（对角线 Padé
逼近）和次对角线 \([n-1/n]\) 具有特别的重要性。

\textbf{例 209.1}（指数函数的 Padé 逼近）：\(e^z\) 的 \([n/n]\) Padé
逼近为

\[r_{n,n}(z) = \frac{\sum_{k=0}^n \frac{(2n-k)! n!}{(2n)! k! (n-k)!} z^k}{\sum_{k=0}^n \frac{(2n-k)! n!}{(2n)! k! (n-k)!} (-z)^k}\]

\subsubsection{209.3 Padé
逼近的收敛性}\label{paduxe9-ux903cux8fd1ux7684ux6536ux655bux6027}

\textbf{定理 209.4}（Padé 逼近的收敛性，Nuttall-Pommerenke
定理）：对大部分亚纯函数，对角线 Padé 序列 \([n/n]\) 在容量意义下收敛到
\(f\)（在极点以外的区域）。Padé 逼近自然地捕捉了函数的极点。

\textbf{定理 209.5}（Stahl
定理，1986）：对具有分支切割的代数函数，对角线 Padé
逼近的极限在容量意义下存在，且极点的极限分布由 \(f\)
的奇异性的最小容量集合给出。

\textbf{定义 209.5}（Padé 逼近与连分式）：Padé
逼近与连分式展开有自然联系。\([n/n]\) 和 \([n-1/n]\) Padé
逼近是函数连分式展开的截断。

\bigskip\hrule\bigskip

\subsection{第210章：样条与小波}\label{ux7b2c210ux7ae0ux6837ux6761ux4e0eux5c0fux6ce2}

样条是分段多项式函数，具有优良的逼近性质和计算效率。小波是时间-频率局部化的基函数，是信号处理的核心工具。

\subsubsection{210.1 样条函数}\label{ux6837ux6761ux51fdux6570}

\textbf{定义 210.1}（\(k\) 次样条）：给定节点
\(a = x_0 < x_1 < \cdots < x_N = b\)，\textbf{\(k\) 次样条函数} \(s(x)\)
满足： 1. \(s\) 在每个子区间 \([x_i, x_{i+1}]\) 上是次数 \(\leq k\)
的多项式 2. \(s \in C^{k-1}[a,b]\)（\(k-1\) 阶连续可微）

\textbf{定义
210.2}（自然三次样条）：\textbf{自然三次样条}（\(k=3\)）是满足端点条件
\(s''(a) = s''(b) = 0\) 的三次样条。

\textbf{定理 210.1}（样条插值的最优性质）：在所有满足插值条件
\(s(x_i) = f(x_i)\) 且 \(s'' \in L^2[a,b]\)
的函数中，自然三次样条插值极小化弯曲能量

\[E(s) = \int_a^b |s''(x)|^2 \, dx\]

\textbf{定理 210.2}（样条逼近的误差估计）：对
\(f \in C^4[a,b]\)，三次样条插值 \(s\) 满足

\[\|f - s\|_\infty \leq \frac{5}{384} \|f^{(4)}\|_\infty h^4, \quad \|f' - s'\|_\infty \leq \frac{1}{24} \|f^{(4)}\|_\infty h^3\]

其中 \(h = \max_i (x_{i+1} - x_i)\)。

\subsubsection{210.2 B-样条基}\label{b-ux6837ux6761ux57fa}

\textbf{定义 210.3}（B-样条，Cox-de Boor 递推，1972）：\textbf{\(k\) 次
B-样条} \(B_{i,k}(x)\) 由递推定义：

\[B_{i,0}(x) = \begin{cases} 1 & x \in [x_i, x_{i+1}) \\ 0 & \text{否则} \end{cases}\]

\[B_{i,k}(x) = \frac{x - x_i}{x_{i+k} - x_i} B_{i,k-1}(x) + \frac{x_{i+k+1} - x}{x_{i+k+1} - x_{i+1}} B_{i+1,k-1}(x)\]

\textbf{命题 210.1}（B-样条的性质）： - 局部支撑：\(B_{i,k}(x)\) 仅在
\([x_i, x_{i+k+1}]\) 上非零 - 非负性：\(B_{i,k}(x) \geq 0\) -
单位分解：\(\sum_i B_{i,k}(x) = 1\)（对所有 \(x\)） -
线性无关：\(\{B_{i,k}\}\) 是 \(k\) 次样条空间的基

\textbf{定理
210.3}（B-样条的导数）：\(B_{i,k}'(x) = \frac{k}{x_{i+k} - x_i} B_{i,k-1}(x) - \frac{k}{x_{i+k+1} - x_{i+1}} B_{i+1,k-1}(x)\)。

\subsubsection{210.3 小波分析}\label{ux5c0fux6ce2ux5206ux6790}

\textbf{定义 210.4}（多分辨分析，Mallat 1989）：\(L^2(\mathbb{R})\)
的\textbf{多分辨分析}（MRA）是一列闭子空间
\(\cdots \subset V_{-1} \subset V_0 \subset V_1 \subset \cdots\) 满足：
1. \(\bigcup V_j\) 在 \(L^2(\mathbb{R})\)
中稠密，\(\bigcap V_j = \{0\}\) 2.
\(f(x) \in V_j \iff f(2x) \in V_{j+1}\) 3.
存在缩放函数（尺度函数）\(\phi \in V_0\)，其整数平移构成 \(V_0\)
的规范正交基

\textbf{定义 210.5}（小波函数）：设 \(W_j\) 是 \(V_j\) 在 \(V_{j+1}\)
中的正交补（\(V_{j+1} = V_j \oplus W_j\)）。\textbf{小波函数}（母小波）\(\psi \in W_0\)，其整数平移构成
\(W_0\)
的规范正交基。\(\{\psi_{j,k}(x) = 2^{j/2} \psi(2^j x - k)\}_{j,k \in \mathbb{Z}}\)
构成 \(L^2(\mathbb{R})\) 的规范正交基。

\textbf{定理 210.4}（Daubechies
小波，1988）：存在具有紧支撑的任意高阶光滑小波函数。Daubechies 小波
\(\psi_N\) 的支撑为 \([0, 2N-1]\)，具有 \(N\)
个消失矩：\(\int x^k \psi_N(x) dx = 0\)（\(k = 0, \ldots, N-1\)）。

\textbf{定理 210.5}（小波的逼近性质）：具有 \(N\) 个消失矩的小波，函数
\(f\) 的小波系数的衰减速率刻画了 \(f\) 的光滑性：

\[|\langle f, \psi_{j,k} \rangle| \leq C 2^{-j(N+1/2)} \|f^{(N)}\|_\infty\]

\subsubsection{210.4
小波在信号处理中的应用}\label{ux5c0fux6ce2ux5728ux4fe1ux53f7ux5904ux7406ux4e2dux7684ux5e94ux7528}

\textbf{定义 210.6}（快速小波变换，FWT = Mallat
算法）：小波分解和重构可通过滤波器组（低通 \(h\) 和高通 \(g\)）以
\(O(n)\) 时间复杂度实现：

\[a_{j,k} = \sum_m h_{m-2k} a_{j+1,m}, \quad d_{j,k} = \sum_m g_{m-2k} a_{j+1,m}\]

\textbf{定理 210.6}（小波收缩与降噪，Donoho-Johnstone
1994）：通过阈值处理小波系数（软阈值或硬阈值），小波方法在降噪中达到近乎最优的最小最大速率。

\bigskip\hrule\bigskip

\subsection{第211章：径向基函数与高维逼近}\label{ux7b2c211ux7ae0ux5f84ux5411ux57faux51fdux6570ux4e0eux9ad8ux7ef4ux903cux8fd1}

径向基函数（RBF）是处理高维散乱数据逼近和插值的主要工具，在机器学习和偏微分方程无网格方法中有广泛应用。

\subsubsection{211.1 径向基函数}\label{ux5f84ux5411ux57faux51fdux6570}

\textbf{定义
211.1}（径向基函数）：\textbf{径向基函数}（RBF）\(\phi(x) = \phi(\|x\|)\)
仅依赖于到原点的距离。常用的 RBF 包括： -
\textbf{Gauss}：\(\phi(r) = e^{-(\varepsilon r)^2}\) -
\textbf{多二次}（Multiquadric）：\(\phi(r) = \sqrt{1 + (\varepsilon r)^2}\)
- \textbf{逆多二次}（Inverse
Multiquadric）：\(\phi(r) = 1 / \sqrt{1 + (\varepsilon r)^2}\) -
\textbf{薄板样条}（Thin Plate
Spline）：\(\phi(r) = r^2 \log r\)（\(d=2\)）

\textbf{定义 211.2}（RBF 插值）：给定散乱数据点
\(\{(x_i, f_i)\}_{i=1}^N \subset \mathbb{R}^d \times \mathbb{R}\)，RBF
插值具有形式

\[s(x) = \sum_{i=1}^N \lambda_i \phi(\|x - x_i\|) + \sum_{j=1}^M \beta_j p_j(x)\]

其中 \(p_j\) 是多项式基（保证唯一可解性）。系数 \(\lambda_i, \beta_j\)
由插值条件和矩条件 \(\sum \lambda_i p_j(x_i) = 0\) 确定。

\subsubsection{211.2 正定函数与再生核 Hilbert
空间}\label{ux6b63ux5b9aux51fdux6570ux4e0eux518dux751fux6838-hilbert-ux7a7aux95f4}

\textbf{定义 211.3}（正定函数）：函数
\(\Phi: \mathbb{R}^d \to \mathbb{R}\) 称为\textbf{正定}的，如果对任意
\(\{x_1, \ldots, x_N\} \subset \mathbb{R}^d\) 和
\(\{c_1, \ldots, c_N\} \subset \mathbb{R}\)，

\[\sum_{i,j=1}^N c_i c_j \Phi(x_i - x_j) \geq 0\]

\textbf{定理 211.1}（Bochner 定理，1933）：连续函数 \(\Phi\)
是正定的当且仅当它是有限正 Borel 测度的 Fourier
变换：\(\Phi(x) = \int_{\mathbb{R}^d} e^{i x \cdot \omega} d\mu(\omega)\)。Gauss
核是正定的，多二次核是条件正定的。

\textbf{定义 211.4}（再生核 Hilbert 空间，RKHS）：设 \(\mathcal{H}\)
是函数 \(f: \Omega \to \mathbb{R}\) 的 Hilbert 空间。\(\mathcal{H}\)
是\textbf{再生核 Hilbert 空间}，如果存在核
\(K: \Omega \times \Omega \to \mathbb{R}\) 使得对任意
\(x \in \Omega\)，\(K(\cdot, x) \in \mathcal{H}\) 且
\(f(x) = \langle f, K(\cdot, x) \rangle_{\mathcal{H}}\)（再生性质）。

\textbf{定理 211.2}（RBF 与 RKHS 的等价性）：每个正定 RBF \(\phi\)
关联一个唯一的 RKHS \(\mathcal{N}_\phi(\Omega)\)（Native 空间）。RBF
插值在此空间中是误差最小的：

\[\|f - s\|_{\mathcal{N}_\phi} \leq \|f\|_{\mathcal{N}_\phi}\]

\subsubsection{211.3 RBF
逼近的收敛性}\label{rbf-ux903cux8fd1ux7684ux6536ux655bux6027}

\textbf{定理 211.3}（RBF 插值的误差估计）：对 Gauss 核
\(\phi(r) = e^{-(\varepsilon r)^2}\)，在 Native 空间中，

\[\|f - s_{X,\varepsilon}\|_\infty \leq C e^{-c \varepsilon / h_{X,\Omega}} \|f\|_{\mathcal{N}_\phi}\]

其中 \(h_{X,\Omega}\) 是填充距离（fill
distance）。\(h_{X,\Omega} \to 0\) 时误差以谱精度收敛。

\textbf{定理 211.4}（权衡原则，Trade-off Principle）：RBF
插值在精度和稳定性之间存在权衡。小
\(\varepsilon\)（平坦核）提高精度但增大条件数，大 \(\varepsilon\)
提高稳定性但降低精度。这是 RBF 方法的核心挑战。

\textbf{定理 211.5}（平坦极限与多项式，Driscoll-Fornberg 2002）：当
\(\varepsilon \to 0\)（平坦极限），某些 RBF
插值收敛到（多元）多项式插值。这解释了 RBF 方法在平坦极限下的行为。

\subsubsection{211.4 RBF 在 PDE
中的应用}\label{rbf-ux5728-pde-ux4e2dux7684ux5e94ux7528}

\textbf{定义 211.5}（Kansa 方法，RBF 配置法，1990）：用 RBF 配置法求解
PDE：将解近似为
\(u(x) \approx \sum_{i=1}^N \lambda_i \phi(\|x - x_i\|)\)，在配置点
\(x_j\) 上强制 PDE 满足，得到关于 \(\lambda_i\) 的线性方程组。

\textbf{定理 211.6}（RBF 方法在无网格 PDE 求解中的优势）：RBF
方法自然地处理高维和不规则区域上的 PDE，无需网格生成。对于椭圆形
PDE，RBF 配置法可达谱精度。

\bigskip\hrule\bigskip

\emph{本卷建立了逼近论的核心理论框架：Weierstrass 逼近定理与 Bernstein
多项式证明建立了多项式逼近的存在性基础，Jackson
直接定理给出了光滑性与逼近阶的定量关系；有理逼近与 Padé 逼近（Newman
定理、Padé
表、收敛性）在处理非光滑函数和极点上远优于多项式；样条与小波（B-样条基、多分辨分析、Daubechies
小波、快速小波变换）提供了高效且局部化的逼近工具；径向基函数（正定核、RKHS、散乱数据逼近、无网格
PDE
方法）是处理高维和复杂几何问题的主要工具。逼近论是连接泛函分析、数值分析和信号处理的桥梁。}

\bigskip\hrule\bigskip

\emph{本卷在已有内容基础上，扩展至 4 章。}

\newpage

\section{卷五十三：统计力学数学}\label{ux5377ux4e94ux5341ux4e09ux7edfux8ba1ux529bux5b66ux6570ux5b66}

\begin{quote}
统计力学是连接微观物理定律与宏观热力学现象的桥梁，其数学基础涵盖概率论、泛函分析和数学物理方法。从
Gibbs
系综理论到相变与临界现象的数学描述，从重整化群方法到精确可解模型（如
Ising 模型的 Onsager
解），再到非平衡统计力学的现代发展，本卷系统建立统计力学的严格数学框架。
\end{quote}

\begin{quote}
\textbf{卷序说明}：本章节号（Ch
213-216）反映其在全书逻辑链中的位置。物理卷序上本卷位于V53，作为统计力学数学方向的专题补充卷。
\end{quote}

\bigskip\hrule\bigskip

\subsection{第213章：Gibbs 系综、热力学极限与相变的 Peierls
论证}\label{ux7b2c213ux7ae0gibbs-ux7cfbux7efcux70edux529bux5b66ux6781ux9650ux4e0eux76f8ux53d8ux7684-peierls-ux8bbaux8bc1}

统计力学的数学基础始于 Gibbs
系综理论，它将平衡态的热力学性质归结为概率测度（Gibbs
测度）的性质。热力学极限则是连接有限系统与宏观系统的桥梁，相变在这极限下表现为
Gibbs 测度的非唯一性。

\subsubsection{213.1 Gibbs
系综与配分函数}\label{gibbs-ux7cfbux7efcux4e0eux914dux5206ux51fdux6570}

\textbf{定义 213.1}（有限体积 Gibbs 测度 / 正则系综）：设
\(\Lambda \subset \mathbb{Z}^d\) 为有限区域，边界条件为
\(\omega \in \Omega_{\Lambda^c}\)（边界位形）。定义 \(\Lambda\) 上的
Hamilton 量
\(H_\Lambda^\omega(\sigma)\)（含相互作用能，\(\sigma \in \Omega_\Lambda\)
为 \(\Lambda\) 内的位形）。\textbf{有限体积 Gibbs 测度}（正则系综在温度
\(T = 1/(k_B \beta)\) 下）为 \(\Omega_\Lambda\) 上的概率测度：
\[\mu_{\Lambda, \beta}^\omega(\sigma) = \frac{1}{Z_{\Lambda, \beta}^\omega} \exp(-\beta H_\Lambda^\omega(\sigma))\]
其中 \textbf{配分函数}（partition function）为归一化常数
\[Z_{\Lambda, \beta}^\omega = \sum_{\sigma \in \Omega_\Lambda} \exp(-\beta H_\Lambda^\omega(\sigma))\]

该定义的核心是 \textbf{Boltzmann 权重}
\(\exp(-\beta H)\)------能量越低的位形，概率权重越大。

\textbf{定义 213.2}（自由能）：有限体积 \textbf{Helmholtz 自由能} 定义为
\[F_{\Lambda, \beta}^\omega = -\frac{1}{\beta} \log Z_{\Lambda, \beta}^\omega\]

热力学量（内能、比热、磁化等）可通过自由能对温度或外场的导数获得。

\textbf{定义 213.3}（Gibbs 测度 / DLR 方程，Dobrushin-Lanford-Ruelle
1968-1970）：无穷体积 \textbf{Gibbs 测度} \(\mu_\beta\) 是
\(\Omega = \Omega_{\mathbb{Z}^d}\) 上的概率测度，满足对任意有限区域
\(\Lambda\)，条件分布 \(\mu_\beta(\cdot \mid \omega_{\Lambda^c})\)
由有限体积 Gibbs 测度给出（DLR 方程）。Gibbs 测度的全体记为
\(\mathcal{G}(\beta)\)。

\subsubsection{213.2 热力学极限}\label{ux70edux529bux5b66ux6781ux9650}

\textbf{定义 213.4}（热力学极限）：令
\(\Lambda_n \uparrow \mathbb{Z}^d\)（如
\(\Lambda_n = [-n, n]^d \cap \mathbb{Z}^d\)）。若自由能密度
\[f(\beta) = \lim_{n \to \infty} \frac{1}{|\Lambda_n|} F_{\Lambda_n, \beta}^\omega\]
存在且与边界条件 \(\omega\)
的选取无关，则称系统具有良定义的\textbf{热力学极限}。\(f(\beta)\)
为自由能密度。

\textbf{定理 213.1}（热力学极限的存在性 / Ruelle-Fisher
定理）：对具有有限程相互作用的平移不变格点系统，热力学极限存在。具体而言，对于有限程势
\(\Phi\) 满足 \(\sum_{X \ni 0} \|\Phi_X\|_\infty < \infty\)，自由能密度
\(f(\beta)\) 是良定义的连续函数。

\subsubsection{213.3 Ising 模型与相变的 Peierls
论证}\label{ising-ux6a21ux578bux4e0eux76f8ux53d8ux7684-peierls-ux8bbaux8bc1}

\textbf{定义 213.5}（Ising 模型）：\(d\) 维 \textbf{Ising 模型}
定义在格点 \(\mathbb{Z}^d\) 上，每个格点 \(i\) 处自旋变量
\(\sigma_i \in \{-1, +1\}\)。Hamilton 量为
\[H_\Lambda(\sigma) = -J \sum_{\langle i, j \rangle \cap \Lambda \neq \emptyset} \sigma_i \sigma_j - h \sum_{i \in \Lambda} \sigma_i\]
其中 \(\langle i, j \rangle\) 表示最近邻对，\(J > 0\)
为铁磁耦合常数（相邻自旋倾向于同向），\(h\) 为外磁场。

\textbf{定理 213.2}（Peierls 论证：Ising 模型存在相变，Peierls 1936,
严格化 Griffiths 1964）：对 \(d \geq 2\) 维零场（\(h = 0\)）Ising
模型，存在临界温度 \(T_c(d) > 0\) 使得： -
\(T < T_c\)（低温）时，存在至少两个不同的平移不变 Gibbs 测度
\(\mu_\beta^+\) 和
\(\mu_\beta^-\)，分别对应于正自发磁化和负自发磁化：\(\mu_\beta^+(\sigma_0 = +1) > 0\)
且 \(\mu_\beta^+(\sigma_0 = +1) \neq \mu_\beta^-(\sigma_0 = +1)\)。 -
\(T > T_c\)（高温）时，Gibbs 测度唯一。

\emph{Peierls 论证概要}：在低温下，考虑正边界条件（边界自旋全为
\(+1\)）。对某格点 \(i\)，若 \(\sigma_i = -1\)，则其周围被 \(+1\)
自旋包围形成''Peierls 轮廓''（contour）------围绕 \(-1\)
区域的闭合路径。每条轮廓 \(\gamma\) 的能量代价为
\(2J |\gamma|\)（\(|\gamma|\) 为轮廓长度），因此轮廓出现的概率
\(\propto \exp(-2\beta J |\gamma|)\)。对轮廓的 Peierls 估计表明，只要
\(\beta J\) 足够大（即 \(T\) 足够低），长轮廓的概率被指数压制，从而
\(-1\) 自旋不能渗透整个系统------自发磁化为正。对于零场，\(+\) 和 \(-\)
边界条件给出不同的 Gibbs 测度，相变得以确立。

\textbf{推论 213.1}（\(d=1\) 无相变）：一维 Ising 模型的 Peierls
论证失效（轮廓为单个点，无能量势垒），事实上 \(d=1\) 时任意 \(T > 0\) 下
Gibbs 测度唯一------不存在有限温度相变。

\textbf{评注 213.1}：Peierls
论证是统计力学中第一个关于相变的严格数学证明思路，其核心在于利用几何-组合方法（轮廓展开）将热力学问题转化为轮廓系统的低密度展开。这一方法后来发展为
\textbf{Pirogov-Sinai 理论}（1975-1976），系统研究了低温相图的普适结构。

\bigskip\hrule\bigskip

\subsection{第214章：重整化群}\label{ux7b2c214ux7ae0ux91cdux6574ux5316ux7fa4}

重整化群（RG）是 Wilson 在 1971
年引入的处理临界现象的统一理论框架，将不同尺度上的物理联系起来。Wilson
因此获 1982 年 Nobel 物理学奖。

\subsubsection{214.1
重整化群变换}\label{ux91cdux6574ux5316ux7fa4ux53d8ux6362}

\textbf{定义 214.1}（重整化群变换，Kadanoff
块自旋变换）：\textbf{重整化群变换} \(\mathcal{R}\) 将微观 Hamilton 量
\(H\) 映射到粗粒化后的有效 Hamilton 量 \(H' = \mathcal{R}(H)\)：

\begin{enumerate}
\def\labelenumi{\arabic{enumi}.}
\tightlist
\item
  将系统分为 \(b^d\) 个自旋组成的块
\item
  每个块用块自旋变量（block spin）代替
\item
  积分掉短波涨落，得到块自旋变量的有效 Hamilton 量
\end{enumerate}

在 Hamilton 量参数空间中，\(\mathcal{R}\) 是保半群的变换。

\textbf{定义 214.2}（RG 不动点与临界曲面）：Hamilton
量参数空间中的\textbf{不动点} \(H^*\) 满足
\(\mathcal{R}(H^*) = H^*\)。临界曲面是流向不动点的 Hamilton
量的集合。不动点的稳定性由线性化 RG 变换的特征值决定。

\textbf{定理 214.1}（RG 不动点的分类）： - \textbf{无关本征方向}（特征值
\(< 1\)）：流向不动点，对应的参数是无关的 -
\textbf{相关本征方向}（特征值
\(> 1\)）：流离不动点，对应的参数是相关的（如温度 \(T - T_c\)，外场
\(h\)） - 相关本征方向的个数决定了临界指数的个数

\textbf{定理 214.2}（RG 与临界指数）：若相关本征方向的本征值为
\(\lambda_t = b^{y_t}\)，\(\lambda_h = b^{y_h}\)，则临界指数为

\[\nu = \frac{1}{y_t}, \quad \eta = d + 2 - 2y_h, \quad \alpha = 2 - \frac{d}{y_t}, \quad \beta = \frac{d - y_h}{y_t}, \quad \gamma = \frac{2y_h - d}{y_t}\]

\subsubsection{\texorpdfstring{214.2
\(\varepsilon\)-展开}{214.2 \textbackslash varepsilon-展开}}\label{varepsilon-ux5c55ux5f00}

\textbf{定理 214.3}（Wilson-Fisher 不动点与 \(\varepsilon\)-展开，Wilson
1972）：在 \(d = 4 - \varepsilon\) 维（\(\varepsilon \ll 1\)）空间中的
\(\phi^4\) 理论，RG 不动点在 \(\varepsilon\) 阶处存在（Wilson-Fisher
不动点）。临界指数可展开为 \(\varepsilon\) 的幂级数：

\[\eta = \frac{\varepsilon^2}{54} + O(\varepsilon^3), \quad \nu = \frac{1}{2} + \frac{\varepsilon}{12} + \frac{7\varepsilon^2}{162} + O(\varepsilon^3)\]

\(\varepsilon\)-展开给出了
\(d=3\)（\(\varepsilon=1\)）时临界指数的系统近似，与实验和数值模拟吻合极好。

\textbf{定理 214.4}（上临界维数）：\(d_c = 4\) 是 Ising
普适类的\textbf{上临界维数}。\(d \geq 4\)
时，平均场临界指数是精确的（Gauss 不动点稳定）。\(d < 4\)
时，涨落使系统流向 Wilson-Fisher 不动点。

\subsubsection{214.3
共形场论与临界现象}\label{ux5171ux5f62ux573aux8bbaux4e0eux4e34ux754cux73b0ux8c61}

\textbf{定理 214.5}（Polyakov
共形不变性假说，1970）：在临界点，系统在尺度变换下不变，且具有更丰富的共形对称性。\(d=2\)
维的临界系统由\textbf{共形场论}（CFT）描述。

\textbf{定理 214.6}（BPZ 定理，Belavin-Polyakov-Zamolodchikov
1984）：\(d=2\) 维共形场论由 Virasoro 代数（中心荷
\(c\)）的表示完全分类。\(c < 1\) 的酉极小模型有离散的系列：

\[c = 1 - \frac{6}{m(m+1)} \quad (m = 3, 4, 5, \ldots)\]

Ising 模型对应 \(m = 3\)（\(c = 1/2\)），三临界 Ising 对应
\(m = 4\)（\(c = 7/10\)）。

\bigskip\hrule\bigskip

\subsection{第215章：精确可解模型}\label{ux7b2c215ux7ae0ux7cbeux786eux53efux89e3ux6a21ux578b}

少数统计力学模型可以精确求解，它们提供了临界现象的非微扰理解。Ising
模型的 Onsager 解是 20 世纪理论物理的最高成就之一。

\subsubsection{215.1 二维 Ising 模型的 Onsager
解}\label{ux4e8cux7ef4-ising-ux6a21ux578bux7684-onsager-ux89e3}

\textbf{定理 215.1}（Onsager 解，1944, 精确自由能）：\(d=2\) 维零场
Ising 模型的自由能密度为

\[-\beta f = \log 2 + \frac{1}{2\pi^2} \int_0^\pi \int_0^\pi \log[\cosh(2\beta J_x) \cosh(2\beta J_y) - \sinh(2\beta J_x) \cos \omega_1 - \sinh(2\beta J_y) \cos \omega_2] \, d\omega_1 d\omega_2\]

（在各向同性情形 \(J_x = J_y = J\) 下）。

\textbf{定理 215.2}（临界温度与对数奇异性）：临界温度由
\(\sinh(2\beta_c J) = 1\)
给出（\(k_B T_c / J \approx 2.269\)）。自由能在 \(T_c\) 处有对数奇异性：

\[C \sim -\log|T - T_c| \quad (T \to T_c)\]

即 \(\alpha = 0\)（对数发散）。

\textbf{定理 215.3}（自发磁化，杨振宁 1952）：\(T < T_c\) 时，二维 Ising
模型的自发磁化为

\[m^* = (1 - \sinh^{-4}(2\beta J))^{1/8}\]

临界指数 \(\beta = 1/8\)。\(T \to T_c^-\) 时
\(m^* \sim (T_c - T)^{1/8}\)。

\textbf{方法概要}：Onsager 使用转移矩阵方法将配分函数化为
\(2^N \times 2^N\) 矩阵的迹，通过 Clifford
代数表示为自由费米子问题，使用 Jordan-Wigner 变换和 Bogoliubov
变换对角化。

\subsubsection{215.2 Bethe 拟设}\label{bethe-ux62dfux8bbe}

\textbf{定义 215.1}（Bethe 拟设，Bethe 1931）：\textbf{Bethe
拟设}是求解一维量子多体问题的方法。一维 Heisenberg 自旋链（XXX
模型，其中 XXX 指三个耦合常数相等
\(J_x = J_y = J_z\)，即各向同性情形）的 Hamilton 量为

\[H = -J \sum_{i=1}^N (\sigma_i^x \sigma_{i+1}^x + \sigma_i^y \sigma_{i+1}^y + \sigma_i^z \sigma_{i+1}^z)\]

\textbf{定理 215.4}（Bethe 拟设解，Bethe 1931）：一维 Heisenberg
反铁磁链的基态可通过 Bethe 拟设严格求解。\(N\) 个粒子的 Bethe 波函数为

\[\psi(x_1, \ldots, x_M) = \sum_{P \in S_M} A_P e^{i \sum_{j=1}^M k_{P_j} x_j}\]

其中波数 \(k_j\) 满足 \textbf{Bethe 拟设方程}：

\[e^{i k_j N} = \prod_{\ell \neq j}^M \frac{\sin(k_j - k_\ell + 2i\eta)}{\sin(k_j - k_\ell - 2i\eta)}\]

\textbf{定理 215.5}（Yang-Baxter 方程，杨振宁 1967, Baxter 1972）：Bethe
拟设可解性的代数本质是 Yang-Baxter 方程：

\[R_{12}(u) R_{13}(u+v) R_{23}(v) = R_{23}(v) R_{13}(u+v) R_{12}(u)\]

其中 \(R_{ij}\) 是作用于 \(V_i \otimes V_j\) 的 \(R\)-矩阵。Yang-Baxter
方程是量子可积性的核心代数条件。

\subsubsection{215.3
六顶角模型与八顶角模型}\label{ux516dux9876ux89d2ux6a21ux578bux4e0eux516bux9876ux89d2ux6a21ux578b}

\textbf{定义 215.2}（六顶角模型）：\textbf{六顶角模型}（6-vertex
model）是二维格点上的冰型模型，每个顶点有 6
种允许的冰规则配置。配分函数可通过 Bethe 拟设求解。

\textbf{定理 215.6}（Lieb 解，1967）：六顶角模型的自由能和极化率由 Lieb
使用 Bethe 拟设方法精确求解。残留熵（\(T=0\) 时的熵）为
\(S(0) = N k_B \frac{3}{2} \log(4/3)\)（``方冰''的残留熵）。

\textbf{定理 215.7}（Baxter 八顶角模型解，1972）：八顶角模型（8-vertex
model）是带有对角相互作用的最一般二维可解格点模型。Baxter 使用
Yang-Baxter 方程和转移矩阵的交换性，给出了自由能的参数化精确解。

\bigskip\hrule\bigskip

\subsection{第216章：非平衡统计力学}\label{ux7b2c216ux7ae0ux975eux5e73ux8861ux7edfux8ba1ux529bux5b66}

非平衡统计力学研究系统远离平衡态的行为，从 Boltzmann
方程的微观推导到涨落定理和线性响应理论。

\subsubsection{216.1 Boltzmann 方程与 H
定理}\label{boltzmann-ux65b9ux7a0bux4e0e-h-ux5b9aux7406}

\textbf{定义 216.1}（Boltzmann 方程，1872）：稀薄气体中单粒子分布函数
\(f(x, v, t)\) 满足

\[\frac{\partial f}{\partial t} + v \cdot \nabla_x f = \left(\frac{\partial f}{\partial t}\right)_{\text{coll}}\]

其中碰撞项为

\[\left(\frac{\partial f}{\partial t}\right)_{\text{coll}} = \int_{\mathbb{R}^3} \int_{S^2} [f(v') f(v_*') - f(v) f(v_*)] B(v - v_*, \omega) \, d\omega \, dv_*\]

（Boltzmann 碰撞算子）。\(B\) 是碰撞核，由分子间相互作用势决定。

\textbf{定理 216.1}（Boltzmann H 定理，1872）：定义 \textbf{H 函数}
\(H(t) = \int f(x, v, t) \log f(x, v, t) \, dx \, dv\)。则
\(dH/dt \leq 0\)，且 \(dH/dt = 0\) 当且仅当 \(f\) 是局部 Maxwell 分布：

\[f_{\text{eq}}(v) = \frac{n}{(2\pi k_B T/m)^{3/2}} \exp\left(-\frac{m|v - u|^2}{2k_B T}\right)\]

H 定理给出了热力学第二定律的微观解释。

\textbf{定理 216.2}（Cercignani 猜想，Boltzmann
方程的收敛速率）：Boltzmann 方程的解以指数速率收敛到 Maxwell 分布。对
Maxwell 分子（\(B\) 仅依赖于散射角），Desvillettes-Villani（2005）证明了
\(H(t) - H_{\text{eq}} \leq C e^{-t/\tau}\)。

\subsubsection{216.2
涨落耗散定理}\label{ux6da8ux843dux8017ux6563ux5b9aux7406}

\textbf{定理 216.3}（涨落耗散定理，Callen-Welton 1951, Kubo
1957）：线性响应函数 \(\chi(t)\) 与平衡涨落的关系：

\[\chi(t) = -\frac{1}{k_B T} \frac{d}{dt} \langle A(t) A(0) \rangle_{\text{eq}} \quad (t > 0)\]

其中 \(\langle A(t) A(0) \rangle\) 是可观测量 \(A\)
的平衡时间关联函数。涨落（关联函数）和耗散（响应函数）由同一微观机制决定。

\textbf{定理 216.4}（Green-Kubo 公式，1954,
1957）：输运系数（如扩散系数、黏度、热导率）由平衡关联函数的时间积分给出：

\[D = \frac{1}{d} \int_0^\infty \langle v(t) \cdot v(0) \rangle_{\text{eq}} \, dt\]

\textbf{定理 216.5}（Onsager
倒易关系，1931）：线性不可逆热力学中，输运系数矩阵 \(L_{ij}\)
是对称的：\(L_{ij} = L_{ji}\)。这是微观可逆性的宏观表现。

\subsubsection{216.3 涨落定理}\label{ux6da8ux843dux5b9aux7406}

\textbf{定理 216.6}（涨落定理，Evans-Cohen-Morriss 1993,
Gallavotti-Cohen 1995）：在非平衡稳态中，熵产生 \(\Sigma_t\) 满足

\[\frac{\mathbb{P}(\Sigma_t = A)}{\mathbb{P}(\Sigma_t = -A)} = e^{A t}\]

即正熵产生的概率指数大于负熵产生的概率。涨落定理将热力学第二定律提升为在涨落层面上的精确关系。

\textbf{定理 216.7}（Jarzynski 等式，1997）：对非平衡过程，功 \(W\)
的分布满足

\[\langle e^{-\beta W} \rangle = e^{-\beta \Delta F}\]

其中 \(\Delta F\) 是自由能差。Jarzynski
等式允许从非平衡功的测量中提取平衡自由能。

\subsubsection{216.4 Fokker-Planck 方程与 Langevin
动力学}\label{fokker-planck-ux65b9ux7a0bux4e0e-langevin-ux52a8ux529bux5b66}

\textbf{定义 216.2}（Fokker-Planck 方程）：\textbf{Fokker-Planck
方程}（Smoluchowski 方程）描述随机过程 \(X_t\) 的概率密度 \(p(x, t)\)
的演化：

\[\frac{\partial p}{\partial t} = -\frac{\partial}{\partial x}[\mu(x) p] + \frac{1}{2} \frac{\partial^2}{\partial x^2}[D(x) p]\]

其中 \(\mu(x)\) 是漂移系数，\(D(x)\) 是扩散系数。

\textbf{定理 216.8}（Fokker-Planck 与 Langevin 的等价性）：Langevin 方程
\(dX_t = \mu(X_t) dt + \sigma(X_t) dW_t\) 的概率密度满足 Fokker-Planck
方程（\(\mu\) 和 \(D = \sigma^2\) 的对应关系由 Stratonovich 或 Itô
解释确定）。

\textbf{定理 216.9}（Kramers 逃逸问题，1940）：Fokker-Planck
方程中，粒子从势阱中逃逸的平均首次通过时间（Arrhenius 律）为

\[\tau \sim e^{\beta \Delta E}\]

其中 \(\Delta E\) 是势垒高度。这是化学反应的过渡态理论的数学基础。

\bigskip\hrule\bigskip

\emph{本卷建立了统计力学的核心数学框架：Gibbs 系综与 DLR
理论（配分函数、Gibbs 测度、DLR
方程、热力学极限）将平衡态统计力学放在严格的概率论基础上，相变表现为
Gibbs 测度的非唯一性；相变理论（Ising 模型、Peierls
论证、临界指数、普适性）揭示了临界现象的本质规律，标度律将不同临界指数统一起来；重整化群（RG
变换、不动点、\(\varepsilon\)-展开、共形场论）是 Wilson
的伟大创造，统一了临界现象的理论描述；精确可解模型（Onsager 解、Bethe
拟设、Yang-Baxter
方程）为强关联系统提供了非微扰的精确解，是非平衡统计力学中量子可积性的核心；非平衡统计力学（Boltzmann
方程、H
定理、涨落耗散定理、涨落定理）将热力学第二定律从微观动力学中推导出来。统计力学是连接微观物理与宏观现象的数学桥梁，也是概率论和动力系统理论的重要应用领域。}

\bigskip\hrule\bigskip

\emph{本卷在已有内容基础上，扩展至 4 章。}

\newpage

\section{卷五十四：金融数学}\label{ux5377ux4e94ux5341ux56dbux91d1ux878dux6570ux5b66}

\begin{quote}
金融数学是随机分析在金融领域的重要应用。从无套利定价原理出发，经由
Black-Scholes-Merton
期权定价理论，到利率模型、信用风险模型和精算数学，本卷系统建立现代金融数学的核心理论框架。随机微积分（Ito
公式、Girsanov 定理、鞅表示定理）为金融衍生品定价提供了严格的数学基础。
\end{quote}

\begin{quote}
\textbf{卷序说明}：本章节号（Ch
167-171）反映其在全书逻辑链中的位置。物理卷序上本卷位于V54，作为金融数学方向的专题补充卷。
\end{quote}

\bigskip\hrule\bigskip

\subsection{第167章：无套利定价原理}\label{ux7b2c167ux7ae0ux65e0ux5957ux5229ux5b9aux4ef7ux539fux7406}

无套利定价是现代金融数学的第一原理------在一个有效的金融市场中，不存在''免费午餐''。这一原理将金融衍生品定价问题转化为等价鞅测度的存在问题，为
Black-Scholes 公式奠定了理论基础。

\subsubsection{167.1
金融市场的基本模型}\label{ux91d1ux878dux5e02ux573aux7684ux57faux672cux6a21ux578b}

\textbf{定义
167.1}（离散时间金融市场模型）：一个离散时间金融市场模型由以下数据构成：概率空间
\((\Omega, \mathcal{F}, P)\)，域流 \(\{\mathcal{F}_t\}_{t=0}^T\)，以及
\(d+1\) 个资产的价格过程
\(\{S_t^i\}_{t=0}^T\)（\(i = 0, 1, \ldots, d\)）。资产 \(0\)
为无风险资产（银行存款），满足 \(S_t^0 = (1+r)^t\)（\(r\)
为常数利率）。资产 \(1, \ldots, d\) 为风险资产（如股票）。

\textbf{定义
167.2}（交易策略与自融资）：一个\textbf{交易策略}（或投资组合）是一个
\(\mathbb{R}^{d+1}\)-值可料过程
\(H_t = (H_t^0, H_t^1, \ldots, H_t^d)\)，其中 \(H_t^i\) 表示 \(t\)
时刻持有资产 \(i\) 的数量。策略 \(H\) 称为\textbf{自融资的}，若
\[V_t(H) := H_t \cdot S_t = H_{t+1} \cdot S_t, \quad t = 0, 1, \ldots, T-1\]
即投资组合价值的变动完全来自资产价格变动，无外部资金注入或撤出。

\textbf{定义 167.3}（折现价格过程）：以无风险资产的折现因子 \(1/S_t^0\)
对价格进行折现，定义折现价格过程 \(\tilde{S}_t^i = S_t^i / S_t^0\)。

\subsubsection{167.2
无套利与等价鞅测度}\label{ux65e0ux5957ux5229ux4e0eux7b49ux4ef7ux9785ux6d4bux5ea6}

\textbf{定义 167.4}（套利机会）：一个自融资策略 \(H\)
称为\textbf{套利机会}，若 \(V_0(H) = 0\)，\(V_T(H) \geq 0\) a.s.，且
\(P(V_T(H) > 0) > 0\)。即零初始成本、非负最终收益，且有正收益的正概率。

\textbf{定义 167.5}（等价鞅测度）：概率测度 \(Q\)
称为\textbf{等价鞅测度}（或风险中性测度），若
\(Q \sim P\)（等价），且折现价格过程 \(\tilde{S}_t\) 在 \(Q\) 下是鞅。

\textbf{定理
167.1}（资产定价第一基本定理）：离散时间有限状态市场中，市场无套利
\(\iff\) 存在等价鞅测度 \(Q\)。

\emph{证明思路}（利用超平面分离定理）：考虑折现收益空间
\(\mathcal{K} = \{V_T(H) - V_0(H) : H \text{ 为自融资策略}\}\)。无套利条件等价于
\(\mathcal{K} \cap \mathbb{R}_{\geq 0}^{\Omega} = \{0\}\)。由凸集的超平面分离定理，存在严格正的线性泛函，经规范化后即得等价鞅测度。∎

\textbf{定理
167.2}（风险中性定价公式）：若市场无套利，则任意可复制未定权益
\(X\)（支付在时刻 \(T\) 的
\(\mathcal{F}_T\)-可测随机变量）的无套利价格为
\[V_0 = \mathbb{E}_Q\left[\frac{X}{S_T^0}\right]\] 即折现支付的
\(Q\)-期望。等价地，\(V_t = S_t^0 \cdot \mathbb{E}_Q[X / S_T^0 \mid \mathcal{F}_t]\)。

\subsubsection{167.3
完全市场与唯一鞅测度}\label{ux5b8cux5168ux5e02ux573aux4e0eux552fux4e00ux9785ux6d4bux5ea6}

\textbf{定义 167.6}（完全市场）：市场称为\textbf{完全的}，若每个未定权益
\(X\) 均可被某个自融资策略复制（对冲）：存在 \(H\) 使得 \(V_T(H) = X\)
a.s.。

\textbf{定理 167.3}（资产定价第二基本定理）：无套利市场是完全的 \(\iff\)
等价鞅测度唯一。

\textbf{命题 167.1}（完全市场的特征）：在一个无套利的有限状态市场中，设
\(D\) 为总可能状态数，\(N\) 为风险资产数。则市场完全当且仅当
\(D \leq N+1\) 且收益矩阵是满秩的。

\subsubsection{167.4
连续时间模型的推广}\label{ux8fdeux7eedux65f6ux95f4ux6a21ux578bux7684ux63a8ux5e7f}

\textbf{定义 167.7}（连续时间金融市场）：设
\((\Omega, \mathcal{F}, \{\mathcal{F}_t\}_{t \in [0,T]}, P)\)
为带域流的概率空间。无风险资产 \(dB_t = r B_t dt\)，\(B_0 = 1\)（即
\(B_t = e^{rt}\)）。风险资产（股票）满足
\[dS_t = \mu_t S_t dt + \sigma_t S_t dW_t\] 其中 \(W_t\) 为 Brown
运动（在适当的 Itô 积分定义下）。

\textbf{定理
167.4}（连续时间第一基本定理，Delbaen-Schachermayer）：市场满足
NFLVR（No Free Lunch with Vanishing Risk）\(\iff\) 存在等价局部鞅测度。

\subsubsection{167.5 Girsanov
定理------测度变换的核心工具}\label{girsanov-ux5b9aux7406ux6d4bux5ea6ux53d8ux6362ux7684ux6838ux5fc3ux5de5ux5177}

从现实测度 \(P\) 变换到风险中性测度 \(Q\) 的过程依赖 Girsanov
定理，它是随机分析中最重要的结果之一。

\textbf{定理 167.5}（Girsanov 定理 / Cameron-Martin-Girsanov 定理）：设
\(W_t\) 为概率空间 \((\Omega, \mathcal{F}, \{\mathcal{F}_t\}, P)\) 上的
\(d\) 维 Brown 运动。设 \(\theta_t\) 为 \(\mathcal{F}_t\)-适应的
\(\mathbb{R}^d\)-值过程，满足 Novikov 条件
\[\mathbb{E}_P\left[\exp\left(\frac{1}{2} \int_0^T \|\theta_t\|^2 dt\right)\right] < \infty\]
定义 \textbf{Radon-Nikodym 密度过程}（随机指数 / Doleans-Dade 指数）
\[Z_t = \exp\left(-\int_0^t \theta_s \cdot dW_s - \frac{1}{2} \int_0^t \|\theta_s\|^2 ds\right)\]
则 \(Z_t\) 是 \(P\)-鞅（\(\mathbb{E}_P[Z_T] = 1\)）。定义新测度 \(Q\)
通过 \(\frac{dQ}{dP}\big|_{\mathcal{F}_T} = Z_T\)。则在 \(Q\) 下，过程
\[\tilde{W}_t = W_t + \int_0^t \theta_s ds\] 是 \(d\) 维 Brown
运动（相对于域流 \(\{\mathcal{F}_t\}\)）。

\textbf{在 Black-Scholes 框架中的应用}：对几何 Brown 运动
\(dS_t = \mu S_t dt + \sigma S_t dW_t\)，取
\(\theta_t = (\mu - r)/\sigma\)（市场风险价格），则 Girsanov
变换将现实世界的漂移 \(\mu\) 替换为无风险利率 \(r\)：
\[dS_t = \mu S_t dt + \sigma S_t dW_t = r S_t dt + \sigma S_t d\tilde{W}_t\]
其中 \(\tilde{W}_t = W_t + \frac{\mu - r}{\sigma}t\) 是风险中性测度
\(Q\) 下的 Brown 运动。折现价格过程 \(\tilde{S}_t = e^{-rt} S_t\) 在
\(Q\) 下成为鞅：\(d\tilde{S}_t = \sigma \tilde{S}_t d\tilde{W}_t\)。

无套利定价原理是整个金融数学的逻辑起点------从这一简单而深刻的原理出发，可以推导出期权定价公式、利率期限结构理论和信用风险定价框架。

\bigskip\hrule\bigskip

\subsection{第168章：Black-Scholes
期权定价理论}\label{ux7b2c168ux7ae0black-scholes-ux671fux6743ux5b9aux4ef7ux7406ux8bba}

Black-Scholes-Merton
公式是金融数学中最重要的结果之一，它为欧式期权的定价提供了显式公式，并通过
Delta 对冲策略将期权定价与随机分析联系起来。1973 年 Black 和 Scholes
发表这一结果，Merton 同年给出了基于连续时间随机过程的严格推导（1997
年诺贝尔经济学奖）。

\subsubsection{168.1 Black-Scholes
模型与推导}\label{black-scholes-ux6a21ux578bux4e0eux63a8ux5bfc}

\textbf{假设 168.1}（Black-Scholes 模型假设）：(1) 股票价格 \(S_t\)
服从几何 Brown 运动 \(dS_t = \mu S_t dt + \sigma S_t dW_t\)，其中
\(\mu, \sigma\) 为常数；(2) 无风险利率 \(r\) 为常数；(3)
允许连续交易，无交易成本；(4) 允许卖空，资产无限可分；(5) 无红利支付。

\textbf{定义 168.1}（欧式看涨期权）：行权价 \(K\)、到期日 \(T\)
的欧式看涨期权在 \(T\) 时刻的收益为
\(C_T = (S_T - K)^+ = \max(S_T - K, 0)\)。

\textbf{定理 168.1}（Black-Scholes PDE）：在 Black-Scholes 模型中，设
\(V(t, S_t)\) 为期权的无套利价格。则 \(V(t, S)\)
满足以下抛物型偏微分方程：
\[\frac{\partial V}{\partial t} + rS\frac{\partial V}{\partial S} + \frac{1}{2}\sigma^2 S^2 \frac{\partial^2 V}{\partial S^2} - rV = 0\]
边界条件：\(V(T, S) = (S - K)^+\)（看涨）或
\(V(T, S) = (K - S)^+\)（看跌）。

\emph{推导（Delta 对冲与 Itô 公式）}：假设期权价格 \(V(t, S)\) 是
\(C^{1,2}\) 函数。构造自融资投资组合
\(\Pi_t = V(t, S_t) - \Delta_t S_t\)，其中 \(\Delta_t\)
为持有股票的数量（待定）。由自融资条件，组合价值的微分为
\(d\Pi_t = dV - \Delta_t dS_t\)。对 \(V(t, S_t)\) 应用 Itô 公式：
\[dV = \left(\frac{\partial V}{\partial t} + \mu S \frac{\partial V}{\partial S} + \frac{1}{2}\sigma^2 S^2 \frac{\partial^2 V}{\partial S^2}\right)dt + \sigma S \frac{\partial V}{\partial S} dW_t\]
代入 \(dS_t = \mu S_t dt + \sigma S_t dW_t\) 得：
\[d\Pi_t = \left(\frac{\partial V}{\partial t} + \mu S \frac{\partial V}{\partial S} + \frac{1}{2}\sigma^2 S^2 \frac{\partial^2 V}{\partial S^2} - \Delta_t \mu S\right)dt + \sigma S\left(\frac{\partial V}{\partial S} - \Delta_t\right)dW_t\]
选取 \(\Delta_t = \frac{\partial V}{\partial S}\)（Delta
对冲），则随机项 \(dW_t\) 的系数为零------组合 \(\Pi_t\)
瞬间无风险。由无套利原理，\(\Pi_t\)
必须以无风险利率增长：\(d\Pi_t = r \Pi_t dt = r(V - \Delta_t S) dt\)。令两式
\(dt\) 项系数相等：
\[\frac{\partial V}{\partial t} + \frac{1}{2}\sigma^2 S^2 \frac{\partial^2 V}{\partial S^2} = rV - rS\frac{\partial V}{\partial S}\]
整理即得 Black-Scholes PDE。此推导的关键步骤在于：(1) Itô 公式将随机过程
\(V(t, S_t)\) 的微分展开；(2) Delta 对冲消除随机波动（消除 \(dW_t\)
项）；(3) 无套利条件强制局部无风险组合的收益率等于 \(r\)。

\textbf{注 168.1}（光滑性条件的验证）：上述推导隐含地假定期权价格函数
\(V(t, S) \in C^{1,2}([0, T] \times \mathbb{R}_+)\)，即 \(V\) 对 \(t\)
一阶连续可微、对 \(S\) 二阶连续可微------这是应用 Itô 公式的必要条件。在
Black-Scholes 模型中，\(V(t, S)\) 的显式公式（定理
168.2）保证了所需的光滑性。对于一般收益函数，只要收益函数是连续且多项式增长的，则
Feynman-Kac 公式给出的解
\(V(t, S) = e^{-r(T-t)} \mathbb{E}_Q[\Phi(S_T) \mid S_t = S]\) 确实是
\(C^{1,2}\) 的，验证了 Itô
公式的适用性。此外，\(\Delta_t = \partial V / \partial S\) 的可料性（由
\(V\) 的 \(C^1\) 性质）确保了 Itô 积分的良定义性。∎

\subsubsection{168.2 Black-Scholes 公式与
Greeks}\label{black-scholes-ux516cux5f0fux4e0e-greeks}

\textbf{定理 168.2}（Black-Scholes 公式）：欧式看涨期权的价格为
\[C(t, S) = S \cdot \Phi(d_1) - K e^{-r(T-t)} \cdot \Phi(d_2)\] 其中
\[d_1 = \frac{\ln(S/K) + (r + \sigma^2/2)(T-t)}{\sigma\sqrt{T-t}}, \quad d_2 = d_1 - \sigma\sqrt{T-t}\]
且 \(\Phi(x) = \frac{1}{\sqrt{2\pi}} \int_{-\infty}^x e^{-u^2/2} du\)
为标准正态分布函数。

\emph{证明（鞅方法）}：在风险中性测度 \(Q\)
下，\(dS_t = r S_t dt + \sigma S_t d\tilde{W}_t\)（\(\tilde{W}_t\) 为
\(Q\)-Brown 运动）。解此 SDE 得
\[S_T = S_t \exp\left((r - \sigma^2/2)(T-t) + \sigma(\tilde{W}_T - \tilde{W}_t)\right)\]
计算 \(C(t, S) = e^{-r(T-t)} \mathbb{E}_Q[(S_T - K)^+ \mid S_t = S]\)
即得。∎

\textbf{定义
168.2}（Greeks------敏感度参数）：期权的风险敏感度由以下偏导数刻画： -
\textbf{Delta}:
\(\Delta = \frac{\partial V}{\partial S}\)（标的价格敏感度），看涨期权的
\(\Delta = \Phi(d_1) \in [0, 1]\) - \textbf{Gamma}:
\(\Gamma = \frac{\partial^2 V}{\partial S^2}\)（Delta
的敏感度），\(\Gamma = \frac{\Phi'(d_1)}{S\sigma\sqrt{T-t}} > 0\) -
\textbf{Theta}:
\(\Theta = \frac{\partial V}{\partial t}\)（时间衰减），通常为负 -
\textbf{Vega}:
\(\mathcal{V} = \frac{\partial V}{\partial \sigma}\)（波动率敏感度），\(\mathcal{V} = S\sqrt{T-t} \cdot \Phi'(d_1) > 0\)
- \textbf{Rho}: \(\rho = \frac{\partial V}{\partial r}\)（利率敏感度）

\textbf{命题 168.1}（Put-Call 平价）：对相同行权价和到期日的欧式期权：
\[C(t, S) - P(t, S) = S - K e^{-r(T-t)}\]

\subsubsection{168.3
隐含波动率与波动率微笑}\label{ux9690ux542bux6ce2ux52a8ux7387ux4e0eux6ce2ux52a8ux7387ux5faeux7b11}

\textbf{定义 168.3}（隐含波动率）：给定市场上期权的观测价格
\(C_{\text{market}}\)，\textbf{隐含波动率} \(\sigma_{\text{imp}}\)
是使得 Black-Scholes 公式给出的理论价格等于市场价格的波动率参数：
\[C_{\text{BS}}(t, S; K, T, r, \sigma_{\text{imp}}) = C_{\text{market}}\]

\textbf{定理
168.3}（隐含波动率的逼近公式）：在平价附近（\(S \approx K\)），隐含波动率有渐近展开：
\[\sigma_{\text{imp}}(K, T) \approx \sigma_0 \left(1 + \frac{\text{skew}}{6} \cdot \frac{\ln(K/S)}{\sigma_0 \sqrt{T}} + \cdots\right)\]

\textbf{现象
168.1}（波动率微笑与偏斜）：实际市场中，隐含波动率随行权价变化，呈现''微笑''（外汇市场）或''偏斜''（股票市场，低行权价对应高隐含波动率）。这一现象表明
Black-Scholes
模型的常数波动率假设与实数据不符，推动了随机波动率模型（Heston
模型）和局部波动率模型的发展。

\subsubsection{168.4 Black-Scholes
模型的推广与评价}\label{black-scholes-ux6a21ux578bux7684ux63a8ux5e7fux4e0eux8bc4ux4ef7}

\textbf{简要介绍}：Black-Scholes
公式的成功在于其简洁性和可操作性。然而，其假设的局限性（常数波动率、连续几何
Brown 运动）催生了丰富的后续研究：Merton 跳跃-扩散模型（1976
年）、Heston 随机波动率模型（1993 年）、Dupire 局部波动率模型、以及基于
Lévy 过程的指数 Lévy 模型。

这些推广共同构成了现代金融衍生品定价的数学基础。Black-Scholes
公式不仅是一个计算公式，更是一种思维范式------它展示了如何利用随机分析（Itô
微积分）和偏微分方程将金融问题转化为数学问题，深刻地影响了其后数十年的金融数学研究。

\bigskip\hrule\bigskip

\subsection{第169章：利率模型与固定收益}\label{ux7b2c169ux7ae0ux5229ux7387ux6a21ux578bux4e0eux56faux5b9aux6536ux76ca}

利率模型研究收益率曲线的随机演化，是固定收益衍生品定价的基础。

\subsubsection{169.1
短期利率模型}\label{ux77edux671fux5229ux7387ux6a21ux578b}

\textbf{定义 169.1}（短期利率）：\(r_t\) 是瞬时无风险利率。零息债券价格
\(P(t, T) = \mathbb{E}^{\mathbb{Q}}[\exp(-\int_t^T r_s ds) | \mathcal{F}_t]\)。

\textbf{定义 169.2}（经典短期利率模型）： - \textbf{Vasicek
模型}（1977）：\(dr_t = \kappa(\theta - r_t) dt + \sigma dW_t\)。\(r_t\)
是 Gauss 过程（利率可为负）。债券价格有显式公式 - \textbf{CIR
模型}（Cox-Ingersoll-Ross，1985）：\(dr_t = \kappa(\theta - r_t) dt + \sigma \sqrt{r_t} dW_t\)。\(r_t\)
非负（若 \(2\kappa\theta \geq \sigma^2\)），服从非中心 \(\chi^2\) 分布 -
\textbf{Hull-White
模型}（1990）：\(dr_t = \kappa(\theta_t - r_t) dt + \sigma dW_t\)（\(\theta_t\)
为时间函数，可拟合初始收益率曲线）

\textbf{定理 169.1}（Vasicek
债券定价公式）：\(P(t, T) = \exp(A(t, T) - B(t, T) r_t)\)，其中

\[B(t, T) = \frac{1 - e^{-\kappa(T-t)}}{\kappa}, \quad A(t, T) = \left(\theta - \frac{\sigma^2}{2\kappa^2}\right)(B(t, T) - (T-t)) - \frac{\sigma^2}{4\kappa} B(t, T)^2\]

\subsubsection{169.2 远期利率与 HJM
框架}\label{ux8fdcux671fux5229ux7387ux4e0e-hjm-ux6846ux67b6}

\textbf{定义 169.3}（远期利率）：\(f(t, T)\) 是在 \(t\) 时刻观测的 \(T\)
时刻的瞬时远期利率。\(P(t, T) = \exp(-\int_t^T f(t, s) ds)\)。

\textbf{定理 169.2}（Heath-Jarrow-Morton
框架，1992）：远期利率的无套利动力学为

\[df(t, T) = \alpha(t, T) dt + \sigma(t, T) dW_t\]

其中
\(\alpha(t, T) = \sigma(t, T) \int_t^T \sigma(t, s) ds\)（在风险中性测度下）。漂移项完全由波动率结构决定。

\textbf{定义 169.4}（LIBOR 市场模型 / BGM
模型，1997）：直接在离散远期利率（LIBOR）上建模，是实际应用最广泛的利率衍生品定价模型。

\subsubsection{169.3 利率衍生品}\label{ux5229ux7387ux884dux751fux54c1}

\textbf{定义 169.5}（常见利率衍生品）： -
\textbf{利率互换}（IRS）：固定利率与浮动利率的交换 -
\textbf{利率上限/下限}（Caps/Floors）：利率看涨/看跌期权组合 -
\textbf{互换期权}（Swaptions）：进入利率互换的期权 - \textbf{Bermudan
互换期权}：可在多个日期行权的互换期权

\textbf{定理 169.3}（Black 模型用于利率期权）：在 LIBOR 市场模型下，caps
的价格可分解为 caplets 的和，每个 caplet 由 Black 公式定价。

\bigskip\hrule\bigskip

\subsection{第170章：信用风险与
Copula}\label{ux7b2c170ux7ae0ux4fe1ux7528ux98ceux9669ux4e0e-copula}

信用风险研究交易对手违约导致的损失，Copula
是建模多变量相关结构的核心工具。

\subsubsection{170.1
违约风险建模}\label{ux8fddux7ea6ux98ceux9669ux5efaux6a21}

\textbf{定义 170.1}（违约时间 / 约化模型）：违约时间 \(\tau\)
是非负随机变量。生存概率
\(S(t) = \mathbb{P}(\tau > t)\)。违约强度（hazard rate）\(\lambda_t\)
满足
\(\mathbb{P}(\tau > t | \mathcal{F}_t) = \exp(-\int_0^t \lambda_s ds)\)。

\textbf{定义 170.2}（信用违约互换 / CDS）：CDS
是信用风险的主要对冲工具。保护买方定期支付保费，在违约时获得补偿（面值减去回收值）。

\textbf{定理 170.1}（CDS 定价）：公平 CDS 保费 \(s\) 满足

\[s \cdot \sum_{i=1}^n \Delta t_i P(0, t_i) S(t_i) = (1-R) \int_0^T P(0, t) (-dS(t))\]

其中 \(R\) 为回收率，\(P(0, t)\) 为无风险贴现因子。

\subsubsection{170.2 结构化模型}\label{ux7ed3ux6784ux5316ux6a21ux578b}

\textbf{定义 170.3}（Merton 模型，1974）：公司资产价值 \(V_t\) 遵循几何
Brown 运动。公司债务为零息债券面值 \(D\)（到期 \(T\)）。若
\(V_T < D\)，公司违约。股权 \(E_T = (V_T - D)^+\)（看涨期权），债务
\(B_T = \min(V_T, D) = D - (D - V_T)^+\)。

\textbf{定理 170.2}（Merton
模型中的违约概率）：\(\mathbb{P}(V_T < D) = \Phi\left( \frac{\ln(D/V_0) - (\mu - \sigma^2/2)T}{\sigma\sqrt{T}} \right)\)。信用利差随杠杆率
\(D/V_0\) 和波动率 \(\sigma\) 增加。

\subsubsection{170.3 Copula
与相关结构}\label{copula-ux4e0eux76f8ux5173ux7ed3ux6784}

\textbf{定义 170.4}（Copula）：\(n\) 维 Copula
\(C: [0, 1]^n \to [0, 1]\) 是均匀边际的联合分布函数。由 Sklar
定理，任何联合分布 \(F(x_1, \ldots, x_n)\) 可分解为
\(F(x_1, \ldots, x_n) = C(F_1(x_1), \ldots, F_n(x_n))\)。

\textbf{定理 170.3}（Sklar 定理，1959）：Copula
是连接边际分布和联合分布的唯一函数，捕捉了变量间的相关结构。

\textbf{例 170.1}（常见 Copula）： - \textbf{Gauss
Copula}：\(C(u_1, \ldots, u_n) = \Phi_\Sigma(\Phi^{-1}(u_1), \ldots, \Phi^{-1}(u_n))\)
- \textbf{t-Copula}：尾部相关性更强，适合极端事件的建模 -
\textbf{Clayton
Copula}：下尾相关，\(C(u, v) = (u^{-\theta} + v^{-\theta} - 1)^{-1/\theta}\)
- \textbf{Gumbel
Copula}：上尾相关，\(C(u, v) = \exp(-[(-\ln u)^\theta + (-\ln v)^\theta]^{1/\theta})\)

\textbf{定理
170.4}（尾部相关性）：\(\tau^L = \lim_{u \to 0} \mathbb{P}(U \leq u | V \leq u)\)，\(\tau^U = \lim_{u \to 1} \mathbb{P}(U > u | V > u)\)。Gauss
Copula 的尾部相关性为 0（在 \(\rho < 1\) 时），而 t-Copula
有正的尾部相关。

\bigskip\hrule\bigskip

\subsection{第171章：精算数学基础}\label{ux7b2c171ux7ae0ux7cbeux7b97ux6570ux5b66ux57faux7840}

精算数学是保险和养老金领域的数学基础，研究生命表、保费计算和准备金评估。

\subsubsection{171.1
生命表与生存模型}\label{ux751fux547dux8868ux4e0eux751fux5b58ux6a21ux578b}

\textbf{定义 171.1}（生存函数与死亡力）：设 \(T_x\) 是 \(x\)
岁的人的剩余寿命。生存函数
\(S_x(t) = \mathbb{P}(T_x > t)\)。\textbf{死亡力}（force of
mortality）\(\mu_{x+t} = -\frac{d}{dt} \ln S_x(t)\)。

\textbf{定义 171.2}（精算记号）： -
\({}_tp_x = \mathbb{P}(T_x > t)\)（\(x\) 岁的人活过 \(t\) 年的概率） -
\({}_tq_x = \mathbb{P}(T_x \leq t) = 1 - {}_tp_x\) -
\({}_{t|u}q_x = {}_tp_x \cdot {}_uq_{x+t}\)（\(x\) 岁的人在 \(t\) 到
\(t+u\) 年间死亡的概率）

\textbf{定理 171.1}（Gompertz
定律，1825）：死亡力近似为指数增长：\(\mu_x \approx B c^x\)（\(B > 0, c > 1\)）。Makeham
推广（1860）：\(\mu_x = A + B c^x\)。

\subsubsection{171.2
保费计算原理}\label{ux4fddux8d39ux8ba1ux7b97ux539fux7406}

\textbf{定义 171.3}（精算现值）：在利率 \(i\)（折现因子
\(v = 1/(1+i)\)）下，\(t\) 年后支付 1 的精算现值为
\({}_tE_x = v^t \cdot {}_tp_x\)。

\textbf{定义 171.4}（纯保费 / 均衡净保费）：终身寿险的趸交纯保费为
\(\bar{A}_x = \int_0^\infty v^t \cdot {}_tp_x \cdot \mu_{x+t} \, dt\)。终身年金为
\(\bar{a}_x = \int_0^\infty v^t \cdot {}_tp_x \, dt\)。

\textbf{定理 171.2}（保费等价原理 / Equivalence
Principle）：\textbf{均衡净保费}（level net premium）\(P\) 满足
\(\mathbb{E}[PV(\text{保费收入})] = \mathbb{E}[PV(\text{保险金给付})]\)。

\textbf{例 171.1}（终身寿险的年缴保费）：\(P = \bar{A}_x / \bar{a}_x\)。

\subsubsection{171.3
准备金与偿付能力}\label{ux51c6ux5907ux91d1ux4e0eux507fux4ed8ux80fdux529b}

\textbf{定义 171.5}（前瞻准备金 / Prospective Reserve）：在 \(t\)
时刻的准备金为
\(_tV = \mathbb{E}[PV(\text{未来给付}) | T_x > t] - \mathbb{E}[PV(\text{未来保费}) | T_x > t]\)。

\textbf{定理 171.3}（Thiele 微分方程）：准备金满足

\[\frac{d}{dt} {}_tV = \delta \cdot {}_tV + P_t - \mu_{x+t} \cdot (B_t - {}_tV)\]

其中 \(\delta\) 为利息力，\(P_t\) 为保费率，\(B_t\) 为死亡给付。

\subsubsection{171.4 风险理论}\label{ux98ceux9669ux7406ux8bba}

\textbf{定义 171.6}（Cramér-Lundberg 模型 /
经典风险模型）：保险公司盈余过程
\(U_t = u + ct - \sum_{i=1}^{N_t} X_i\)，其中 \(u\) 为初始资本，\(c\)
为保费率，\(N_t\) 为 Poisson 过程（索赔次数），\(X_i\) 为索赔额。

\textbf{定义
171.7}（破产概率）：\(\psi(u) = \mathbb{P}(U_t < 0 \text{ 对某个 } t > 0)\)。

\textbf{定理 171.4}（Lundberg 不等式）：\(\psi(u) \leq e^{-R u}\)，其中
\(R\)（Lundberg 调整系数）是方程
\(\lambda + cR = \lambda \mathbb{E}[e^{R X}]\) 的唯一正解。

\textbf{定理 171.5}（Cramér-Lundberg
渐近公式）：\(\psi(u) \sim C e^{-R u}\)（当
\(u \to \infty\)），在适当的矩条件下。

\bigskip\hrule\bigskip

\emph{本卷在 V21（概率论 II）和
V33（随机分析）的基础上，建立了无套利定价、Black-Scholes
期权定价、利率模型、信用风险、Copula 和精算数学的系统理论。共 5 章。}

\bigskip\hrule\bigskip

\emph{本卷在已有内容基础上，扩展至 5 章。}

\newpage

\section{卷五十五：计算数学}\label{ux5377ux4e94ux5341ux4e94ux8ba1ux7b97ux6570ux5b66}

\begin{quote}
计算数学是连接数学理论与工程实践的桥梁，它将偏微分方程等连续数学问题转化为可在计算机上求解的离散问题。本卷建立在数值分析（V23）和微分方程（V18）的基础上，系统建立有限元方法的数学理论、谱方法的逼近性质、多重网格法的收敛性分析，以及计算流体力学中的核心数值格式。计算数学的核心张力在于精度与效率的平衡------高精度方法往往计算代价高昂，而快速方法可能损失精度，算法的设计正是寻找这一平衡点的艺术。
\end{quote}

\bigskip\hrule\bigskip

\subsection{第277章：有限元方法理论基础}\label{ux7b2c277ux7ae0ux6709ux9650ux5143ux65b9ux6cd5ux7406ux8bbaux57faux7840}

有限元方法（FEM）是求解偏微分方程最强大的数值工具之一。其核心思想是将问题从无穷维函数空间的强形式转化为有限维子空间中的弱形式（变分形式），从而将
PDE 转化为代数方程组。

\subsubsection{277.1
变分原理与弱形式}\label{ux53d8ux5206ux539fux7406ux4e0eux5f31ux5f62ux5f0f}

\textbf{定义 277.1}（抽象变分问题）：设 \(V\) 为 Hilbert
空间，\(a: V \times V \to \mathbb{R}\)
为连续强制双线性形式，\(f \in V'\)。求 \(u \in V\) 使得

\[a(u, v) = f(v), \quad \forall v \in V\]

\textbf{定理 277.1}（Lax-Milgram 引理）：若 \(a\) 满足有界性
\(|a(u,v)| \leq M\|u\|\|v\|\) 和强制性
\(a(u,u) \geq \alpha\|u\|^2\)（\(\alpha > 0\)），则变分问题存在唯一解
\(u \in V\) 且 \(\|u\|_V \leq \frac{1}{\alpha}\|f\|_{V'}\)。

\textbf{例 277.1}（Poisson 方程的弱形式）：\(-\Delta u = f\) 在
\(\Omega\)，\(u|_{\partial\Omega} = 0\)。取
\(V = H_0^1(\Omega)\)，\(a(u,v) = \int_\Omega \nabla u \cdot \nabla v \, dx\)，\(f(v) = \int_\Omega f v \, dx\)。Poincaré
不等式确保强制性。

\subsubsection{277.2 有限元离散与 Céa
引理}\label{ux6709ux9650ux5143ux79bbux6563ux4e0e-cuxe9a-ux5f15ux7406}

\textbf{定义 277.2}（Galerkin 逼近）：选取有限维子空间
\(V_h \subset V\)，其中 \(h\) 为网格参数。求 \(u_h \in V_h\) 使得

\[a(u_h, v_h) = f(v_h), \quad \forall v_h \in V_h\]

\textbf{定理 277.2}（Céa 引理）：设 \(a\) 连续强制，\(u\) 和 \(u_h\)
分别为连续问题和离散问题的解。则

\[\|u - u_h\|_V \leq \frac{M}{\alpha} \inf_{v_h \in V_h} \|u - v_h\|_V\]

即有限元解的误差控制在最佳逼近误差的常数倍之内。

\textbf{定义 277.3}（Lagrange 有限元空间）：对区域 \(\Omega\) 的三角剖分
\(\mathcal{T}_h\)，定义分片多项式空间

\[V_h^k = \{v_h \in C^0(\bar{\Omega}) : v_h|_K \in P_k(K), \forall K \in \mathcal{T}_h, v_h|_{\partial\Omega} = 0\}\]

其中 \(P_k\) 为次数不超过 \(k\) 的多项式空间。

\textbf{定理 277.3}（插值误差估计）：设
\(u \in H^{k+1}(\Omega)\)，\(I_h u\) 为其在 \(V_h^k\) 中的节点插值。则

\[\|u - I_h u\|_{L^2} + h\|u - I_h u\|_{H^1} \leq C h^{k+1} |u|_{H^{k+1}}\]

\textbf{定理 277.4}（有限元误差估计）：对于 Poisson 问题，若
\(u \in H^{k+1}(\Omega) \cap H_0^1(\Omega)\)，用 \(P_k\) 有限元离散，则

\[\|u - u_h\|_{H^1} \leq C h^k \|u\|_{H^{k+1}}, \quad \|u - u_h\|_{L^2} \leq C h^{k+1} \|u\|_{H^{k+1}}\]

（\(L^2\) 估计通过 Aubin-Nitsche 对偶论证得到）。

\subsubsection{277.3 混合有限元与 inf-sup
条件}\label{ux6df7ux5408ux6709ux9650ux5143ux4e0e-inf-sup-ux6761ux4ef6}

\textbf{定义 277.4}（鞍点问题）：求 \((u, p) \in V \times Q\) 使得

\[a(u, v) + b(v, p) = f(v), \quad b(u, q) = g(q), \quad \forall (v, q) \in V \times Q\]

\textbf{定理 277.5}（Brezzi 分裂定理 / inf-sup
条件）：鞍点问题适定的充要条件是 inf-sup
条件（Ladyzhenskaya-Babuška-Brezzi 条件）：

\[\inf_{q \in Q} \sup_{v \in V} \frac{b(v, q)}{\|v\|_V \|q\|_Q} \geq \beta > 0\]

这是混合有限元（如 Stokes 问题和弹性力学问题的稳定化格式）的理论基石。

\subsubsection{277.4
自适应有限元与后验误差估计}\label{ux81eaux9002ux5e94ux6709ux9650ux5143ux4e0eux540eux9a8cux8befux5deeux4f30ux8ba1}

\textbf{定义 277.5}（后验误差估计子）：对每个单元
\(K \in \mathcal{T}_h\)，残差型后验误差估计子为

\[\eta_K^2 = h_K^2 \|f + \Delta u_h\|_{L^2(K)}^2 + \frac{1}{2} \sum_{e \subset \partial K} h_e \|[\nabla u_h \cdot \mathbf{n}_e]\|_{L^2(e)}^2\]

其中 \([\nabla u_h \cdot \mathbf{n}_e]\) 表示跨越单元边界 \(e\)
的法向梯度跳跃。

\textbf{定理 277.6}（可靠性与有效性的上下界）：存在常数 \(C_1, C_2\)（与
\(h\) 无关）使得

\[C_1 \|u - u_h\|_{H^1} \leq \left(\sum_K \eta_K^2\right)^{1/2} \leq C_2 \|u - u_h\|_{H^1}\]

后验误差估计的可靠性确保估计是误差的上界，有效性确保估计不会过高估计真实误差。基于此的
\textbf{自适应有限元方法}
对误差较大的单元进行局部加密（\(h\)-细化）或提高多项式次数（\(p\)-细化），从而以最优计算代价达到目标精度。

\bigskip\hrule\bigskip

\subsection{第278章：谱方法}\label{ux7b2c278ux7ae0ux8c31ux65b9ux6cd5}

谱方法利用全局光滑基函数（如 Fourier 级数、Chebyshev 多项式、Legendre
多项式）逼近微分方程的解，在光滑解情形下具有''指数收敛''（谱精度）的卓越性质，是模拟湍流、天气预报等需要极高精度问题的首选方法。

\subsubsection{278.1 Fourier 谱方法}\label{fourier-ux8c31ux65b9ux6cd5}

\textbf{定义 278.1}（Fourier-Galerkin 方法）：在周期区域 \([0, 2\pi]^d\)
上，取有限维子空间
\(V_N = \text{span}\{e^{ik\cdot x} : |k|_\infty \leq N\}\)。求
\(u_N \in V_N\) 使得残差与 \(V_N\) 中的检验函数正交。

\textbf{定理 278.1}（谱精度 / 指数收敛）：若 \(u\)
是解析的（周期情形下），则 Fourier 谱逼近的误差满足

\[\|u - u_N\|_{H^s} \leq C e^{-\gamma N} \|u\|_{\text{解析}}, \quad \gamma > 0, s \geq 0\]

与有限元的代数收敛 \(\mathcal{O}(h^k)\) 形成鲜明对比。

\subsubsection{278.2 Chebyshev 谱方法与 Gauss
求积}\label{chebyshev-ux8c31ux65b9ux6cd5ux4e0e-gauss-ux6c42ux79ef}

\textbf{定义 278.2}（Chebyshev 多项式）：在 \([-1, 1]\) 上定义
\(T_n(x) = \cos(n \arccos x)\)。\(\{T_n\}\) 构成加权
\(L^2_w(-1, 1)\)（\(w(x) = 1/\sqrt{1-x^2}\)）的完备正交基。

\textbf{定理 278.2}（Chebyshev-Gauss-Lobatto 求积公式）：\(N+1\) 点 CGL
求积公式 \(\int_{-1}^1 f(x) w(x) dx \approx \sum_{j=0}^N w_j f(x_j)\)
对所有次数 \(\leq 2N-1\) 的多项式精确成立。节点
\(x_j = \cos(\pi j / N)\)。

\subsubsection{278.3
伪谱方法与非线性项处理}\label{ux4f2aux8c31ux65b9ux6cd5ux4e0eux975eux7ebfux6027ux9879ux5904ux7406}

\textbf{定义
278.3}（伪谱方法）：在物理空间（配置点）和谱空间（Fourier/Chebyshev
系数）之间通过快速变换（FFT）来回转换，在物理空间计算非线性项，避免系数空间的卷积运算。

\textbf{命题 278.1}（去混淆 / 3/2 规则）：对于 Fourier 伪谱方法，若保留
\(N\) 个模，为防止二次非线性项的混淆误差，非线性计算应在 \(3N/2\)
个点上进行。

\subsubsection{278.4 谱元法（Spectral Element
Method）}\label{ux8c31ux5143ux6cd5spectral-element-method}

\textbf{定义
278.4}（谱元法）：谱元法结合了有限元的几何灵活性和谱方法的高精度------将区域分解为若干四边形/六面体单元，在每个单元上使用高阶多项式（通常
\(N = 5\) 至 \(15\)）和 Gauss-Lobatto-Legendre (GLL)
配置点。质量矩阵由于配置点与求积点重合而自然对角化。

\textbf{定理 278.3}（谱元法的误差估计）：对于 \(d\) 维椭圆问题，若解
\(u \in H^{m}(\Omega)\)，使用 \(N\) 次多项式的谱元法，\(hp\)-误差估计为

\[\|u - u_{hp}\|_{H^1} \leq C \frac{h^{\min(N+1, m)-1}}{N^{m-1}} \|u\|_{H^m}\]

当解光滑（\(m\) 很大）时，增大 \(N\)（\(p\)-细化）比减小
\(h\)（\(h\)-细化）能获得更快的收敛速度------这是谱元法相比传统有限元的核心优势。

\bigskip\hrule\bigskip

\subsection{第279章：多重网格法与预处理}\label{ux7b2c279ux7ae0ux591aux91cdux7f51ux683cux6cd5ux4e0eux9884ux5904ux7406}

多重网格法是大规模科学计算中求解离散椭圆问题的最快迭代方法之一------其巧妙之处在于利用不同尺度网格之间的信息传递，将迭代法的''光滑''收敛特性转化为''全谱''的最优收敛速度。

\subsubsection{279.1
经典迭代法的光滑性质}\label{ux7ecfux5178ux8fedux4ee3ux6cd5ux7684ux5149ux6ed1ux6027ux8d28}

\textbf{定义 279.1}（迭代法的矩阵分裂）：求解 \(Au = f\)。分裂
\(A = M - N\)（\(M\) 可逆），迭代格式
\(u^{(k+1)} = M^{-1}N u^{(k)} + M^{-1}f\)。迭代矩阵
\(G = M^{-1}N = I - M^{-1}A\)。

\begin{itemize}
\tightlist
\item
  \textbf{Jacobi 方法}：\(M = D = \text{diag}(A)\)
\item
  \textbf{Gauss-Seidel 方法}：\(M = D + L\)（下三角）
\end{itemize}

\textbf{定理 279.1}（光滑性质）：对 Gauss-Seidel 和加权 Jacobi
迭代，高频误差分量迅速衰减（每次迭代衰减约
\(\rho_{\text{high}} \ll 1\)），而低频误差分量几乎不衰减（\(\rho_{\text{low}} \approx 1\)）。因此，经典迭代法是优秀的''光滑器''。

\subsubsection{279.2
两网格法与多重网格原理}\label{ux4e24ux7f51ux683cux6cd5ux4e0eux591aux91cdux7f51ux683cux539fux7406}

\textbf{定义 279.2}（两网格法）：(1) 在前处理（光滑）阶段在细网格上执行
\(\nu_1\) 次 Gauss-Seidel 迭代消除高频误差；(2)
将残差限制到粗网格上求解粗网格校正方程；(3) 将校正延拓回细网格；(4)
在后处理阶段再执行 \(\nu_2\) 次光滑迭代。

\textbf{定理 279.2}（两网格收敛因子）：对模型问题（Poisson
方程，一致网格，双线性插值），两网格法的收敛因子满足

\[\rho_{TG} \leq \frac{C}{1 + C} < 1\]

与网格尺寸 \(h\)
无关。这意味着迭代次数不随问题规模增长而增加（数值可扩展性）。

\textbf{定理 279.3}（多重网格复杂度）：对 \(d\) 维椭圆问题，一次
V-循环多重网格迭代的计算复杂度为 \(\mathcal{O}(N)\)（\(N\)
为未知量总数），在适当条件下可达到网格无关收敛。完整求解（达到离散精度）的复杂度为
\(\mathcal{O}(N)\)，这是求解离散椭圆 PDE 的理论下限。

\subsubsection{279.3
代数多重网格（AMG）方法}\label{ux4ee3ux6570ux591aux91cdux7f51ux683camgux65b9ux6cd5}

\textbf{定义 279.3}（代数多重网格）：经典多重网格依赖几何网格层级。AMG
直接从系数矩阵 \(A\)
的代数结构出发，通过聚合或匹配算法自动构造粗网格层级，适用于复杂几何和高度非结构化网格。

\textbf{定理 279.4}（AMG 的两网格收敛理论）：若 \(A\) 是对称正定的
M-矩阵，则 AMG 的两网格收敛因子 \(\leq 1 - 1/\kappa\)，其中 \(\kappa\)
依赖于粗化和插值算子的构造质量。

\subsubsection{279.4 Krylov
子空间预处理方法}\label{krylov-ux5b50ux7a7aux95f4ux9884ux5904ux7406ux65b9ux6cd5}

\textbf{定义 279.4}（Krylov 子空间方法）：对于大型稀疏线性系统
\(Au = f\)，Krylov 子空间方法在第 \(m\) 步寻求
\(u_m \in u_0 + \mathcal{K}_m(A, r_0)\)，其中

\[\mathcal{K}_m(A, r_0) = \text{span}\{r_0, Ar_0, A^2 r_0, \ldots, A^{m-1} r_0\}\]

\textbf{CG 方法}（共轭梯度法）：当 \(A\) 对称正定时，按
\(A\)-范数实现最优逼近。收敛速率
\(\|u - u_m\|_A \leq 2\left(\frac{\sqrt{\kappa} - 1}{\sqrt{\kappa} + 1}\right)^m \|u - u_0\|_A\)，其中
\(\kappa = \kappa_2(A)\)。

\textbf{GMRES 方法}（广义极小残量法）：对非对称矩阵，最小化残差
\(\|Au_m - f\|_2\) 在 \(\mathcal{K}_m\) 上的投影。

\textbf{定理 279.5}（预处理的本质）：预处理是将 \(A\) 替换为
\(P^{-1}A\)（左预处理）或
\(AP^{-1}\)（右预处理），使得预处理后的矩阵条件数大幅减小。多重网格法作为预处理算子
\(P^{-1}\) 与 CG 或 GMRES
结合（MG-CG、MG-GMRES），是当前求解大规模稀疏线性系统的最有效策略之一。

\bigskip\hrule\bigskip

\subsection{第280章：双曲型方程与守恒律数值方法}\label{ux7b2c280ux7ae0ux53ccux66f2ux578bux65b9ux7a0bux4e0eux5b88ux6052ux5f8bux6570ux503cux65b9ux6cd5}

双曲型方程（波动方程、Euler
方程、浅水波方程等）的数值求解具有特殊的挑战------解可能出现激波等间断，对数值格式的守恒性、熵条件和
TVD（Total Variation Diminishing）性质提出严格要求。

\subsubsection{280.1
守恒律与弱解}\label{ux5b88ux6052ux5f8bux4e0eux5f31ux89e3}

\textbf{定义 280.1}（标量守恒律）：一维标量守恒律

\[\frac{\partial u}{\partial t} + \frac{\partial f(u)}{\partial x} = 0, \quad u(x, 0) = u_0(x)\]

其中 \(f\)
为通量函数。即使初始条件光滑，解也可能在有限时间内产生间断（激波），因此需要弱解概念。

\textbf{定义
280.2}（熵条件）：满足守恒律的弱解可能不唯一。物理相关的唯一解由
\textbf{熵条件} 选取：存在凸熵函数 \(U(u)\) 和熵通量
\(F(u)\)（\(F' = U' f'\)），使得熵不等式
\(\partial_t U + \partial_x F \leq 0\) 在分布意义下成立。

\subsubsection{280.2 Godunov
方法与迎风格式}\label{godunov-ux65b9ux6cd5ux4e0eux8fceux98ceux683cux5f0f}

\textbf{定义 280.3}（数值通量与守恒格式）：守恒格式具有形式

\[u_j^{n+1} = u_j^n - \frac{\Delta t}{\Delta x} (\hat{f}_{j+1/2} - \hat{f}_{j-1/2})\]

其中数值通量 \(\hat{f}_{j+1/2} = \hat{f}(u_j^n, u_{j+1}^n)\)
须满足相容性 \(\hat{f}(u, u) = f(u)\)。

\textbf{定义 280.3a}（CFL 条件 / Courant-Friedrichs-Lewy
条件，1928）：对线性对流方程 \(u_t + c u_x = 0\)
的显式差分格式，数值依赖域必须包含 PDE
的物理依赖域，否则格式必不稳定。具体地，CFL 条件要求

\[\left|c\right| \frac{\Delta t}{\Delta x} \leq 1\]

\emph{推导概要}：方程的特征线为
\(x = ct + x_0\)，\(u(x,t) = u_0(x - ct)\)------解在 \((x_j, t_{n+1})\)
的值完全由初始数据在特征线与 \(t = t_n\)
交点处的值决定。数值格式（如迎风格式）中 \(u_j^{n+1}\) 仅依赖于
\(\{u_{j-1}^n, u_j^n, u_{j+1}^n\}\)（即空间跨度至多为
\(2\Delta x\)）。若物理依赖点 \(x_j - c\Delta t\) 不在数值依赖区间
\([x_{j-1}, x_{j+1}]\) 内，即
\(|c|\Delta t > \Delta x\)，则无论数值格式设计如何精良，均无法收敛到真解------此为
CFL 条件的必要性。更一般地，对非线性问题 \(u_t + f(u)_x = 0\)，CFL
条件推广为 \(\max |f'(u)| \cdot \Delta t / \Delta x \leq 1\)。注：CFL
条件是稳定性的必要非充分条件------满足 CFL
条件不保证格式稳定（如中心差分格式），但不满足则必不稳定。∎

\textbf{定理 280.1}（Godunov
定理）：线性常系数情形下，单调守恒格式的精度最多为一阶。

\textbf{定义 280.4}（Godunov 型格式）：在每个单元界面上精确求解 Riemann
问题（初值间断的守恒律），然后对解进行单元平均。Riemann 问题的自相似解是
Godunov 方法的核心组成部分。

\subsubsection{280.3
高分辨率格式：TVD、ENO、WENO}\label{ux9ad8ux5206ux8fa8ux7387ux683cux5f0ftvdenoweno}

\textbf{定义 280.5}（TVD 格式）：数值格式称为 \textbf{TVD}（Total
Variation Diminishing），若

\[TV(u^{n+1}) \leq TV(u^n), \quad TV(u) = \sum_j |u_{j+1} - u_j|\]

TVD 条件确保无虚假振荡（Gibbs 现象被抑制）。

\textbf{定理 280.2}（Harten 引理）：若格式可写为
\(u_j^{n+1} = u_j^n - C_{j-1/2}(u_j^n - u_{j-1}^n) + D_{j+1/2}(u_{j+1}^n - u_j^n)\)
且
\(C_{j-1/2}, D_{j+1/2} \geq 0\)，\(C_{j-1/2} + D_{j+1/2} \leq 1\)，则该格式是
TVD 的。

\textbf{定义 280.6}（ENO 与 WENO 格式）：ENO（Essentially
Non-Oscillatory）格式通过自适应选择最光滑的模板构造高阶重构多项式，避免跨间断插值。WENO（Weighted
ENO）进一步通过凸组合多个模板，在光滑区域获得更高阶精度，在间断附近自动降阶------实现了高精度和无振荡的统一。

\textbf{定理 280.3}（Lax-Wendroff 定理）：若守恒格式的数值解 \(u_h\) 在
\(h \to 0\) 时几乎处处收敛到某个函数 \(u\)，且格式守恒，则 \(u\)
是守恒律的弱解。

\subsubsection{280.4 一维和二维 Euler
方程的数值求解}\label{ux4e00ux7ef4ux548cux4e8cux7ef4-euler-ux65b9ux7a0bux7684ux6570ux503cux6c42ux89e3}

\textbf{定义 280.7}（一维 Euler 方程）：理想气体的 Euler 方程为

\[\frac{\partial}{\partial t}\begin{pmatrix} \rho \\ \rho u \\ E \end{pmatrix} + \frac{\partial}{\partial x}\begin{pmatrix} \rho u \\ \rho u^2 + p \\ u(E + p) \end{pmatrix} = 0\]

其中 \(\rho\) 为密度，\(u\) 为速度，\(E\)
为总能量，\(p = (\gamma - 1)(E - \frac{1}{2}\rho u^2)\)
为压力（理想气体状态方程）。

\textbf{Riemann 问题的波系结构}：一维 Euler 方程的 Riemann
问题产生三种基本波------激波（shock）、接触间断（contact
discontinuity）和稀疏波（rarefaction wave）。精确 Riemann
求解器通过迭代求解非线性方程确定中间状态；近似 Riemann 求解器（如 Roe
求解器、HLL/HLLC 求解器）以较低计算代价获得可接受的精度。二维 Euler
方程的数值求解通常采用 \textbf{维数分裂（dimensional splitting）} 或
\textbf{非结构网格} 上的有限体积法。

\bigskip\hrule\bigskip

\subsection{第281章：计算流体力学数学基础}\label{ux7b2c281ux7ae0ux8ba1ux7b97ux6d41ux4f53ux529bux5b66ux6570ux5b66ux57faux7840}

计算流体力学（CFD）将 Navier-Stokes
方程的数值求解推向实际工程应用，涉及不可压缩流、可压缩流、湍流建模等核心问题，是计算数学最具挑战性的领域之一。

\subsubsection{281.1 Navier-Stokes
方程的数值挑战}\label{navier-stokes-ux65b9ux7a0bux7684ux6570ux503cux6311ux6218}

\textbf{定义 281.1}（不可压缩 Navier-Stokes 方程）：

\[\frac{\partial \mathbf{u}}{\partial t} + (\mathbf{u} \cdot \nabla)\mathbf{u} = -\nabla p + \nu \Delta \mathbf{u} + \mathbf{f}, \quad \nabla \cdot \mathbf{u} = 0\]

\textbf{挑战}：(1) 不可压缩约束 \(\nabla \cdot \mathbf{u} = 0\)
导致速度-压力耦合的鞍点问题------离散化必须满足 inf-sup 条件（LBB
条件）；(2) 高 Reynolds 数下对流项占优，需要稳定化处理；(3)
湍流的多尺度特性使直接数值模拟（DNS）的计算代价极高（\(\propto Re^3\)）。

\textbf{定义 281.2}（无量纲参数）：Reynolds 数 \(Re = U L / \nu\)
刻画惯性力与粘性力之比。低 \(Re\) 为层流（稳态光滑），高 \(Re\)
为湍流（非稳态混沌）。另外，Mach 数 \(Ma = U/c\)（\(c\)
为声速）判明压缩性------\(Ma < 0.3\) 可视作不可压缩。

\subsubsection{281.2
投影法与分步法}\label{ux6295ux5f71ux6cd5ux4e0eux5206ux6b65ux6cd5}

\textbf{定义 281.3}（Chorin-Temam 投影法）：(1)
预测步：忽略压力梯度求解中间速度

\[\frac{\mathbf{u}^* - \mathbf{u}^n}{\Delta t} = \nu \Delta \mathbf{u}^* - (\mathbf{u}^n \cdot \nabla)\mathbf{u}^n + \mathbf{f}\]

\begin{enumerate}
\def\labelenumi{(\arabic{enumi})}
\setcounter{enumi}{1}
\tightlist
\item
  投影步：将 \(\mathbf{u}^*\) 投影到无散度空间
  \(\mathbf{u}^{n+1} = \mathcal{P}(\mathbf{u}^*)\)，其中
  \(\mathcal{P} = I - \nabla \Delta^{-1} \nabla\cdot\)（Leray-Helmholtz
  投影），投影步需求解 Poisson 方程
  \(\Delta \phi = \nabla \cdot \mathbf{u}^*\)。
\end{enumerate}

\textbf{定理 281.1}（投影法的 Helmholtz-Hodge
分解）：对任意充分光滑的向量场 \(\mathbf{w}\)，存在唯一正交分解
\(\mathbf{w} = \mathbf{u} + \nabla q\)，其中
\(\nabla \cdot \mathbf{u} = 0\) 且
\(\mathbf{u} \cdot \mathbf{n}|_{\partial\Omega} = 0\)。投影算子
\(\mathcal{P}: \mathbf{w} \mapsto \mathbf{u}\) 是 \(L^2\) 正交投影。

\textbf{定理 281.2}（投影法的精度）：标准 Chorin
投影法对速度是一阶时间精度的。通过引入压力的时间外推（如压力校正投影法或旋转投影法），可将时间精度提升至二阶。

\textbf{注 281.1}（旋转投影法 / Rotational Projection
Method）：标准投影法的 Helmholtz-Hodge 分解中，压力项采用 \(\nabla p\)
的形式，导致边界层误差积累限制时间精度。\textbf{旋转投影法}（Timmermans-Minev-Van
De Vosse 1996, Guermond-Shen 2006）将压力梯度替换为
\(\nabla \times \nabla \times\) 形式（利用恒等式
\(\Delta = \nabla(\nabla\cdot) - \nabla \times \nabla \times\)），通过显式处置涡量边界条件，消除了标准投影法在无滑移边界上产生的法向压力梯度的人工边界层，从而将整体时间精度提升至二阶。该方法在旋转坐标系（rotational
form of Navier-Stokes）下自然导出，故名''旋转投影法''。

\subsubsection{281.3
稳定化有限元方法}\label{ux7a33ux5b9aux5316ux6709ux9650ux5143ux65b9ux6cd5}

\textbf{定义
281.4}（对流-扩散方程）：\(-\varepsilon \Delta u + \mathbf{b} \cdot \nabla u = f\)。当
Péclet 数 \(Pe = \|\mathbf{b}\|h/(2\varepsilon) > 1\) 时，标准 Galerkin
方法产生非物理振荡。需引入稳定化。

\textbf{SUPG（Streamline Upwind Petrov-Galerkin）方法}：修改检验函数
\(v \mapsto v + \tau \mathbf{b}\cdot\nabla v\)，沿流线方向增加人工扩散而不破坏精度。\(\tau\)
的选取通过误差分析确定：

\[\tau_K = \frac{h_K}{2\|\mathbf{b}\|}\left(\coth Pe_K - \frac{1}{Pe_K}\right)\]

其中 \(Pe_K = \|\mathbf{b}\| h_K / (2\varepsilon)\) 为单元 Péclet
数。该公式的含义为：当对流占优时（\(Pe_K \gg 1\)，\(\coth Pe_K \to 1\)），\(\tau_K \approx h_K / (2\|\mathbf{b}\|)\)，沿流线方向引入的稳定化量与局部网格尺寸和流速匹配；当扩散占优时（\(Pe_K \to 0\)，\(\coth Pe_K - 1/Pe_K \approx Pe_K / 3\)），\(\tau_K \approx \varepsilon h_K^2 / (6\|\mathbf{b}\|^2)\)，稳定化自动减弱以避免过度扩散。\(\tau_K\)
的这一双渐近公式确保 SUPG 在从纯对流到纯扩散的整个参数范围内都表现最优。

\textbf{GLS（Galerkin Least-Squares）方法}：在变分形式中加入最小二乘项
\(\sum_K \tau_K (\mathcal{L}v, \mathcal{L}u - f)_K\)，同时稳定对流项和压力项，适用于多物理场耦合的复杂情形。

\subsubsection{281.4 湍流模型与
LES}\label{ux6e4dux6d41ux6a21ux578bux4e0e-les}

\textbf{定义 281.5}（大涡模拟 /
LES）：不直接求解所有尺度的湍流运动，而是通过滤波操作
\(\bar{u}(\mathbf{x}) = \int G(\mathbf{x} - \boldsymbol{\xi}) u(\boldsymbol{\xi}) d\boldsymbol{\xi}\)
分离大尺度（可解析尺度）和小尺度（亚格子尺度），并对亚格子应力建立封闭模型。

\textbf{Smagorinsky 模型}：涡粘性假设

\[\tau_{ij} - \frac{1}{3}\tau_{kk}\delta_{ij} = -2\nu_t \bar{S}_{ij}, \quad \nu_t = (C_s \Delta)^2 |\bar{S}|\]

其中
\(\bar{S}_{ij} = \frac{1}{2}(\partial \bar{u}_i/\partial x_j + \partial \bar{u}_j/\partial x_i)\)
为滤波后的应变率张量，\(|\bar{S}| = \sqrt{2\bar{S}_{ij}\bar{S}_{ij}}\)，\(C_s \approx 0.1\)--\(0.2\)
为 Smagorinsky 常数，\(\Delta\) 为滤波宽度。

\textbf{RANS（Reynolds-Averaged
Navier-Stokes）方法}：将瞬时量分解为时间平均和脉动量（\(u = \bar{u} + u'\)），得到
Reynolds 平均方程------包含 Reynolds 应力张量 \(\overline{u_i' u_j'}\)
需湍流模型封闭（如 \(k\)-\(\varepsilon\) 模型、\(k\)-\(\omega\) SST
模型）。

\textbf{定理 281.3}（DNS
的计算代价）：对均匀各向同性湍流，直接数值模拟（DNS）的空间自由度
\(N \propto Re^{9/4}\)，计算代价 \(\propto Re^{3}\)。这意味着 \(Re\)
每增大 10 倍，计算量增大约 1000 倍------这解释了为何 DNS
目前仅限于低至中等 Reynolds 数的流动。

\subsubsection{结语}\label{ux7ed3ux8bed}

计算数学将纯数学的理论（泛函分析、PDE
理论、线性代数）转化为可执行的算法，是连接数学与科学计算的桥梁。从有限元的
Sobolev 空间理论到多重网格的 \(\mathcal{O}(N)\)
复杂度，从双曲守恒律的熵条件到 Navier-Stokes
的投影法------计算数学的每一步发展都展示着数学理论对实际计算的深刻指导作用。

\newpage

\section{卷五十六：图像与信号处理数学}\label{ux5377ux4e94ux5341ux516dux56feux50cfux4e0eux4fe1ux53f7ux5904ux7406ux6570ux5b66}

\begin{quote}
图像与信号处理是调和分析、变分方法和最优化理论的重要应用领域。从 Fourier
变换的频域分析到多尺度小波分解，从全变差正则化的变分去噪到压缩感知的稀疏重建，数学为信号处理提供了从建模到算法再到理论保证的完整框架。本卷建立在调和分析（V25）、泛函分析（V14）、数值分析（V23）和优化理论（V27）的基础上，系统建立小波分析的多分辨框架、变分图像处理的偏微分方程方法、压缩感知的稀疏恢复理论、稀疏建模与字典学习，以及深度学习驱动的图像重建方法。
\end{quote}

\bigskip\hrule\bigskip

\subsection{第282章：小波分析的多分辨框架}\label{ux7b2c282ux7ae0ux5c0fux6ce2ux5206ux6790ux7684ux591aux5206ux8fa8ux6846ux67b6}

小波分析是 Fourier
分析的重大发展------它在时域和频域同时具有局部化能力（``数学显微镜''），能够有效地分析信号中的瞬态特征和多尺度结构。与
Fourier
变换仅提供全局频域信息不同，小波变换揭示了信号的\textbf{多尺度时频结构}，在信号去噪、压缩和特征提取中具有不可替代的优势。

\subsubsection{282.1
连续小波变换}\label{ux8fdeux7eedux5c0fux6ce2ux53d8ux6362}

\textbf{定义 282.1}（连续小波变换）：设 \(\psi \in L^2(\mathbb{R})\)
满足\textbf{允许性条件}
\[\int_{\mathbb{R}} \frac{|\hat{\psi}(\omega)|^2}{|\omega|} d\omega < \infty\]
信号 \(f \in L^2(\mathbb{R})\) 的\textbf{连续小波变换}（CWT）定义为
\[W_f(a, b) = \frac{1}{\sqrt{|a|}} \int_{\mathbb{R}} f(t) \overline{\psi\left(\frac{t-b}{a}\right)} dt\]
其中 \(a \in \mathbb{R} \setminus \{0\}\)
为\textbf{尺度参数}（控制频率分辨率：小 \(|a|\) 对应高频细节，大 \(|a|\)
对应低频轮廓），\(b \in \mathbb{R}\)
为\textbf{平移参数}（控制时间位置）。函数 \(\psi\)
称为\textbf{母小波}（Mother Wavelet）。

\textbf{定理 282.1}（Calderon 重构公式）：若 \(\psi\)
满足允许性条件，则对任意 \(f \in L^2(\mathbb{R})\)，
\[f(t) = \frac{1}{C_\psi} \int_{-\infty}^\infty \int_{-\infty}^\infty W_f(a, b) \frac{1}{\sqrt{|a|}} \psi\left(\frac{t-b}{a}\right) \frac{da db}{a^2}\]
其中
\(C_\psi = \int_{\mathbb{R}} \frac{|\hat{\psi}(\omega)|^2}{|\omega|} d\omega\)
为允许性常数。该定理表明连续小波变换是可逆的------信号可由其小波系数完全重建。

\textbf{定义 282.2}（时频窗与 Heisenberg
不确定原理）：小波变换的时频窗为
\[\left[b + a t^* - a \Delta_\psi, b + a t^* + a \Delta_\psi\right] \times \left[\frac{\omega^*}{a} - \frac{\Delta_{\hat{\psi}}}{a}, \frac{\omega^*}{a} + \frac{\Delta_{\hat{\psi}}}{a}\right]\]
其中 \(t^*, \omega^*\) 和 \(\Delta_\psi, \Delta_{\hat{\psi}}\) 分别为
\(\psi\) 及其 Fourier 变换的中心和半径。时间窗宽度
\(\propto |a|\)，频率窗宽度
\(\propto 1/|a|\)------在低频处频率分辨率高（时间分辨率低），在高频处时间分辨率高（频率分辨率低），满足
\(\Delta t \cdot \Delta \omega \geq 1/2\)。这种\textbf{恒 Q
滤波}特性使小波变换特别适合分析具有瞬变行为的信号。

\subsubsection{282.2
多分辨分析与正交小波基}\label{ux591aux5206ux8fa8ux5206ux6790ux4e0eux6b63ux4ea4ux5c0fux6ce2ux57fa}

\textbf{定义 282.3}（多分辨分析 / MRA，Mallat-Meyer
1986-1989）：\(L^2(\mathbb{R})\)
的一个\textbf{多分辨分析}是一列嵌套闭子空间
\(\cdots \subset V_{-1} \subset V_0 \subset V_1 \subset \cdots\) 满足：

\begin{enumerate}
\def\labelenumi{\arabic{enumi}.}
\tightlist
\item
  \textbf{稠密性与分离性}：\(\overline{\bigcup_{j \in \mathbb{Z}} V_j} = L^2(\mathbb{R})\)，\(\bigcap_{j \in \mathbb{Z}} V_j = \{0\}\)
\item
  \textbf{二进尺度关系}：\(f(t) \in V_j \iff f(2t) \in V_{j+1}\)
\item
  \textbf{平移不变性}：\(f(t) \in V_0 \iff f(t - k) \in V_0, \forall k \in \mathbb{Z}\)
\item
  \textbf{Riesz 基存在性}：存在\textbf{尺度函数}（Scaling
  Function）\(\phi \in V_0\)，使得
  \(\{\phi(t - k)\}_{k \in \mathbb{Z}}\) 构成 \(V_0\) 的 Riesz
  基，即存在 \(0 < A \leq B < \infty\) 使得对所有
  \(\{c_k\} \in \ell^2(\mathbb{Z})\)，
  \[A \sum_k |c_k|^2 \leq \left\|\sum_k c_k \phi(\cdot - k)\right\|^2 \leq B \sum_k |c_k|^2\]
\end{enumerate}

MRA 的直观理解：\(V_0\) 是''基本分辨率''空间（由尺度函数生成），\(V_j\)
是分辨率逐级加倍的空间（\(j\) 越大，包含越精细的细节）。小波空间 \(W_j\)
定义为 \(V_j\) 在 \(V_{j+1}\)
中的正交补：\(V_{j+1} = V_j \oplus W_j\)。递归分解得
\(L^2(\mathbb{R}) = \overline{\bigoplus_{j \in \mathbb{Z}} W_j}\)。

\textbf{定理 282.2}（Mallat 算法---快速小波变换，1989）：由两尺度方程
\[\phi(t) = \sqrt{2} \sum_{k \in \mathbb{Z}} h_k \phi(2t - k), \quad \psi(t) = \sqrt{2} \sum_{k \in \mathbb{Z}} g_k \phi(2t - k)\]
其中 \(g_k = (-1)^k \overline{h_{1-k}}\) 为高通滤波器系数（\(h_k\)
为低通滤波器系数），可得离散信号
\(c_n^0 = \langle f, \phi_{0,n} \rangle\) 的快速分解与重构算法：

\textbf{分解}（由细到粗）：\(c_k^{j-1} = \sum_n \overline{h_{n-2k}} c_n^j\)，\(d_k^{j-1} = \sum_n \overline{g_{n-2k}} c_n^j\)

\textbf{重构}（由粗到细）：\(c_n^j = \sum_k h_{n-2k} c_k^{j-1} + \sum_k g_{n-2k} d_k^{j-1}\)

该算法的复杂度为 \(\mathcal{O}(N)\)（\(N\) 为信号长度），优于 FFT 的
\(\mathcal{O}(N \log N)\)，且在实现上仅需卷积和下采样操作。

\textbf{定义 282.4}（消失矩与正则性）：小波 \(\psi\) 具有 \textbf{\(N\)
阶消失矩}，若
\[\int_{\mathbb{R}} t^k \psi(t) dt = 0, \quad k = 0, 1, \ldots, N-1\]
消失矩越高，小波对光滑信号的压缩性能越好------信号在光滑区域的低阶
Taylor
展开项被小波自动''消除''。同时，消失矩与尺度函数的光滑性相关：高消失矩意味着
\(\hat{\phi}(\omega)\) 在 \(\omega = 0\) 处衰减更快，\(\phi\) 更光滑。

\textbf{例 282.1}（Daubechies 紧支集正交小波族，1988）：Daubechies
构造了具有紧支集和任意高消失矩的正交小波族
\(D_N\)（\(N = 1, 2, \ldots\)）。\(D_N\) 小波的性质： - \(N\)
阶消失矩：\(\int t^k \psi(t) dt = 0\)（\(k = 0, \ldots, N-1\)） -
支集长度为 \(2N-1\) - \(D_1\) 即 Haar 小波（\(N=1\)，支集 \([0,1]\)） -
\(D_4\) 是使用最广泛的 Daubechies 小波（\(N=2\)，连续但不可微）

\subsubsection{282.3
小波在信号处理中的应用}\label{ux5c0fux6ce2ux5728ux4fe1ux53f7ux5904ux7406ux4e2dux7684ux5e94ux7528-1}

\textbf{定义 282.5}（小波阈值去噪---Donoho-Johnstone VisuShrink
方法，1994）：信号去噪的三个步骤： 1.
\textbf{分解}：计算含噪信号的小波变换系数
\(w_{j,k}\)（\(j\)：尺度，\(k\)：位置） 2.
\textbf{阈值处理}：对小波系数施加非线性收缩 -
\textbf{硬阈值}：\(\hat{w} = w \cdot \mathbf{1}_{\{|w| > \lambda\}}\) -
\textbf{软阈值}：\(\hat{w} = \operatorname{sign}(w) \cdot (|w| - \lambda)_+\)
- 其中 \(\lambda = \sigma \sqrt{2 \log N}\)
为\textbf{通用阈值}（Universal Threshold），\(\sigma\)
为噪声标准差（通常由小波系数的中位数绝对偏差 MAD 估计） 3.
\textbf{重构}：通过阈值化后的小波系数逆变换重构去噪信号

\textbf{定理 282.3}（VisuShrink 的渐进最优性，Donoho-Johnstone
1994）：在适当条件下（信号属于 Besov 空间且噪声为独立
Gaussian），通用阈值去噪方法在极小极大意义下是渐进最优的------其风险
\(R(\hat{f}, f) = \mathbb{E}\|\hat{f} - f\|^2\) 以因子
\(1/\sqrt{\log N}\) 逼近理想风险（Oracle
风险，即已知哪些系数为非零时的最优风险）。软阈值在均方误差意义上略优于硬阈值，且能避免硬阈值带来的不连续性。

\textbf{定义 282.6}（SureShrink
自适应阈值）：对每个分解层使用不同的阈值，阈值通过 Stein
无偏风险估计（SURE）选择，不需要先验知道噪声水平。
\[\lambda_j = \arg\min_{\lambda} \text{SURE}(\lambda; \{w_{j,k}\}_k)\]

\textbf{应用 282.1}（JPEG2000 图像压缩标准，ISO/IEC 15444）：JPEG2000
基于离散小波变换（Cohen-Daubechies-Feauveau 9/7
双正交小波用于有损压缩，LeGall 5/3 小波用于无损压缩）和 EBCOT（Embedded
Block Coding with Optimized Truncation）编码。相比 JPEG（基于
\(8 \times 8\) 块 DCT），JPEG2000
的优势包括：在低比特率下显著减少块效应、支持渐进传输（从低分辨率到高分辨率逐级解码）、支持感兴趣区域（ROI）编码（图像中关键区域可以更高精度编码）。

\subsubsection{282.4 Besov
空间的基本定义}\label{besov-ux7a7aux95f4ux7684ux57faux672cux5b9aux4e49}

Besov 空间 \(B_{p,q}^s(\mathbb{R}^d)\) 是包含 Sobolev 空间和 Holder
空间的更一般的函数空间族，通过光滑模（modulus of smoothness）或
Littlewood-Paley 分解定义，特别适合刻画小波系数的衰减行为。

\textbf{定义 282.6}（Besov 空间------差分刻画）：设
\(s > 0\)，\(1 \leq p, q \leq \infty\)。记
\(\Delta_h^r f(x) = \sum_{k=0}^r \binom{r}{k} (-1)^{r-k} f(x + kh)\) 为
\(r\)
阶差分（\(r = \lfloor s \rfloor + 1\)）。\(f \in L^p(\mathbb{R}^d)\)
属于 \textbf{Besov 空间} \(B_{p,q}^s(\mathbb{R}^d)\)，若其半范数
\[|f|_{B_{p,q}^s} = \left( \int_0^\infty \left( t^{-s} \sup_{|h| \leq t} \|\Delta_h^r f\|_{L^p} \right)^q \frac{dt}{t} \right)^{1/q} < \infty\]
（\(q = \infty\) 时取 \(\sup\)）。Besov 范数为
\(\|f\|_{B_{p,q}^s} = \|f\|_{L^p} + |f|_{B_{p,q}^s}\)。

直观上，\(B_{p,q}^s\) 包含了''在 \(L^p\) 意义下具有 \(s\)
阶光滑性''的函数------\(s\) 越大越光滑，\(q\)
是精细指标控制光滑性衰减速度。当 \(p = q = 2\)
时，\(B_{2,2}^s = H^s\)（Sobolev 空间）；当 \(p = q = \infty\) 且 \(s\)
非整数时，\(B_{\infty,\infty}^s = C^s\)（Holder 空间）。

\textbf{在信号处理中的意义}：Besov 空间同时捕获函数的全局正则性（\(s\)
------ 光滑阶数）和局部不规则性（\(p\) ------
积分范数，控制对奇异的敏感程度）。小波系数
\(d_{j,k} = \langle f, \psi_{j,k} \rangle\) 的衰减刻画了 \(f\) 所属的
Besov 空间：\(f \in B_{p,q}^s\) 等价于
\(\|f\|_{B_{p,q}^s} \asymp \|c_{0,\cdot}\|_{\ell^p} + \|(2^{j(s+d/2-d/p)} \|d_{j,\cdot}\|_{\ell^p})_{j \geq 0}\|_{\ell^q}\)------即小波系数在不同尺度上的
\(\ell^p\) 范数以 \(2^{-js}\) 速率衰减。这一定量关系正是小波阈值去噪（如
VisuShrink）在 Besov 空间中达到极小极大最优性的理论根源。

\bigskip\hrule\bigskip

\subsection{第283章：变分图像处理与 PDE
方法}\label{ux7b2c283ux7ae0ux53d8ux5206ux56feux50cfux5904ux7406ux4e0e-pde-ux65b9ux6cd5}

变分方法和偏微分方程（PDE）为图像处理提供了一种统一的数学框架------图像处理任务被建模为\textbf{能量泛函的最小化问题}，最优解满足相应的
\textbf{Euler-Lagrange
方程}（PDE），通过数值迭代求解。这一框架的优雅之处在于：图像处理被归结为三个核心元素的设计------\textbf{数据保真项}（衡量与原始图像的偏差）、\textbf{正则项}（编码图像的先验知识，如光滑性、分段恒常性）和\textbf{数值算法}（高效求解最小化问题）。

\subsubsection{283.1
Rudin-Osher-Fatemi（ROF）全变差模型}\label{rudin-osher-fatemirofux5168ux53d8ux5deeux6a21ux578b}

\textbf{定义 283.1}（全变差 / Total Variation）：函数
\(u \in L^1(\Omega)\)（\(\Omega \subset \mathbb{R}^d\)）的\textbf{全变差}定义为
\[TV(u) = \sup\left\{\int_\Omega u \cdot \operatorname{div} \phi \, dx : \phi \in C_c^1(\Omega, \mathbb{R}^d), \|\phi(x)\|_\infty \leq 1, \forall x \in \Omega\right\}\]
若 \(u \in C^1(\bar{\Omega})\)，则
\(TV(u) = \int_\Omega |\nabla u| dx\)。全变差是 \(W^{1,1}\) 半范数在
\(BV(\Omega)\)（有界变差函数空间）上的推广------它允许函数具有跳跃间断（即边缘），这是图像处理中
TV 正则化能\textbf{保边}的根本原因。

\textbf{定理 283.1}（ROF 去噪模型 / Rudin-Osher-Fatemi,
1992）：给定含噪图像 \(f \in L^2(\Omega)\)，去噪图像 \(u\)
是以下极小化问题的解：
\[\min_{u \in BV(\Omega)} \left\{ \int_\Omega |\nabla u| dx + \frac{\lambda}{2} \int_\Omega (u - f)^2 dx \right\}\]
其中第一项 \(TV(u) = \int_\Omega |\nabla u| dx\)（对于光滑
\(u\)）为\textbf{正则项}------惩罚振荡（噪声）同时允许保留尖锐边缘（因为边缘的
TV 贡献是有限的跳跃而非无穷）；第二项
\(\frac{\lambda}{2}\|u - f\|_{L^2}^2\) 为\textbf{保真项}（Fidelity
Term）------保持去噪结果与原始图像的相似度。参数 \(\lambda > 0\)
控制去噪强度：\(\lambda\) 越大则保真项权重越大（去噪越轻），\(\lambda\)
越小则 TV 项权重越大（去噪越强）。

\textbf{定理 283.2}（ROF 模型的 Euler-Lagrange
方程与几何解释）：形式上的最优性条件为
\[-\operatorname{div}\left(\frac{\nabla u}{|\nabla u|}\right) + \lambda(u - f) = 0\]
其中第一项 \(-\operatorname{div}(\nabla u / |\nabla u|)\)
是\textbf{平均曲率算子}------它是水平集 \(\{u = c\}\)
的曲率。该式定义了\textbf{各向异性扩散}：扩散在梯度方向（法向）被抑制，在等照度线方向（切向）被允许。具体而言：
- 在 \(|\nabla u|\) 大的区域（边缘），扩散系数 \(1/|\nabla u|\)
小------扩散被抑制以\textbf{保边缘} - 在 \(|\nabla u|\)
小的区域（平滑区域），扩散系数大------扩散被促进以\textbf{去噪}

这解释了 TV 正则化能在去噪的同时保持边缘不模糊的原因。

\textbf{算法 283.1}（Chambolle-Pock 原始-对偶算法，2011）：为有效求解
ROF 模型（不可微 TV 项），引入 Legendre-Fenchel
对偶：\(TV(u) = \max_{\|p\|_\infty \leq 1} \langle u, -\operatorname{div} p \rangle\)。原始-对偶迭代：
\[p^{k+1} = \Pi_B\left(p^k + \sigma \nabla \bar{u}^k\right)\]
\[u^{k+1} = \frac{u^k + \tau \operatorname{div} p^{k+1} + \tau \lambda f}{1 + \tau \lambda}\]
\[\bar{u}^{k+1} = 2u^{k+1} - u^k\] 其中 \(\Pi_B\) 为到 \(\ell^\infty\)
单位球的投影（逐点 \(p / \max(1, |p|)\)），步长满足
\(\sigma \tau \|\nabla\|^2 < 1\)。该算法在 \(\mathcal{O}(1/N)\)
收敛速度下求解 ROF 模型，是 TV 去噪的标准实现。

\subsubsection{283.2 Mumford-Shah
分割与水平集方法}\label{mumford-shah-ux5206ux5272ux4e0eux6c34ux5e73ux96c6ux65b9ux6cd5}

\textbf{定义 283.2}（Mumford-Shah 泛函，1989）：图像分割的变分模型
\[E(u, \Gamma) = \int_{\Omega \setminus \Gamma} |\nabla u|^2 dx + \mu \cdot \mathcal{H}^{d-1}(\Gamma) + \lambda \int_\Omega (u - f)^2 dx\]
其中： - \(\Gamma \subset \Omega\)
为\textbf{分割边界}（边缘集，可能不连续但为 \((d-1)\)
维的），\(\mathcal{H}^{d-1}(\Gamma)\) 是 \(\Gamma\) 的 \((d-1)\) 维
Hausdorff 测度（在 \(\mathbb{R}^2\) 中为长度） - \(u\)
为分段光滑的分割图像（在 \(\Omega \setminus \Gamma\) 上光滑，沿
\(\Gamma\) 允许跳跃） - 第一项
\(\int_{\Omega \setminus \Gamma} |\nabla u|^2\) 鼓励分段光滑性 - 第二项
\(\mu \cdot \mathcal{H}^{d-1}(\Gamma)\) 惩罚过长的边界（控制分割的粒度）
- 第三项 \(\lambda \int (u - f)^2\) 保持与原始图像的相似度

Mumford-Shah
泛函的妙处在于它同时优化了\textbf{分割的平滑性}（在区域内部）和\textbf{边界的简洁性}（短边界优先，避免过分割）。当
\(\mu \to \infty\) 时，边界被禁止，模型退化为 \(H^1\)
正则化（纯粹光滑）；当 \(\mu \to 0\) 时，允许任意复杂的边界。

\textbf{定义 283.3}（Chan-Vese 模型 / 无边缘主动轮廓，2001）：Chan 和
Vese 提出了 Mumford-Shah 的水平集实现------用水平集函数
\(\phi: \Omega \to \mathbb{R}\) 隐式表示分割轮廓
\(\Gamma = \{x : \phi(x) = 0\}\)，区域由 \(\phi\)
的符号区分（\(\phi > 0\) 为内部，\(\phi < 0\) 为外部）。能量泛函为
\[E(c_1, c_2, \phi) = \lambda_1 \int_{\{\phi > 0\}} (f - c_1)^2 dx + \lambda_2 \int_{\{\phi < 0\}} (f - c_2)^2 dx + \mu \int_\Omega |\nabla H(\phi)| dx + \nu \int_\Omega H(\phi) dx\]
其中 \(H\) 为 Heaviside 函数（正则化版本），\(c_1, c_2\)
为两区域的平均灰度值。对应的梯度下降流（PDE 演化）为
\[\frac{\partial \phi}{\partial t} = \delta(\phi) \left[\mu \operatorname{div}\left(\frac{\nabla \phi}{|\nabla \phi|}\right) - \nu - \lambda_1 (f - c_1)^2 + \lambda_2 (f - c_2)^2\right]\]
其中 \(\delta(\phi) |\nabla \phi|\)
限制演化仅在零水平集附近活跃------保证了计算效率。

\subsubsection{283.3
图像修复与光流估计}\label{ux56feux50cfux4feeux590dux4e0eux5149ux6d41ux4f30ux8ba1}

\textbf{定义 283.4}（图像修复 / Image Inpainting）：设 \(\Omega\)
中的子集 \(D\)
为待修复的缺失区域。\textbf{Bertalmio-Sapiro-Caselles-Ballester
方法}（2000）沿等照度线（isophote）方向传播信息：
\[\frac{\partial u}{\partial t} = \nabla (\Delta u) \cdot \nabla^\perp u\]
其中 \(\nabla^\perp u = (-\partial_y u, \partial_x u)\)
为等照度线的方向向量。信息沿等照度线方向被''输送''进入缺失区域，保持图像的几何结构（边缘的连续性）。

\textbf{定义 283.5}（光流 / Horn-Schunck 模型，1981）：给定图像序列
\(I(x, y, t)\)，光流估计的变分问题为
\[\min_{\mathbf{v} = (v_1, v_2)} \left\{ \int_\Omega (\nabla I \cdot \mathbf{v} + I_t)^2 dx dy + \alpha \int_\Omega (|\nabla v_1|^2 + |\nabla v_2|^2) dx dy \right\}\]
其中第一项 \((\nabla I \cdot \mathbf{v} + I_t)^2\)
为\textbf{亮度恒常假设}（Brightness
Constancy）------假设像素的运动不改变其亮度值（光流约束方程
\(I_x v_1 + I_y v_2 + I_t = 0\)），第二项为\textbf{平滑先验}（Smoothness
Prior）------假设光流场在空间上是平滑的（即相邻像素具有相似的运动）。\(\alpha > 0\)
平衡数据项和正则项。Euler-Lagrange 方程为
\[(I_x^2 + \alpha \Delta) v_1 + I_x I_y v_2 = -I_x I_t\]
\[I_x I_y v_1 + (I_y^2 + \alpha \Delta) v_2 = -I_y I_t\]
这是一个耦合的椭圆型 PDE 系统，可通过 Gauss-Seidel 迭代求解。

\bigskip\hrule\bigskip

\subsection{第284章：压缩感知与稀疏恢复}\label{ux7b2c284ux7ae0ux538bux7f29ux611fux77e5ux4e0eux7a00ux758fux6062ux590d}

\textbf{预备定理 284.0}（Nyquist-Shannon 采样定理，Nyquist 1928 /
Shannon 1949）：设信号 \(f \in L^2(\mathbb{R})\) 是带限的，即其 Fourier
变换满足 \(\operatorname{supp} \hat{f} \subset [-B, B]\)（最大频率为
\(B\) Hz）。若以采样间隔 \(\Delta t \leq 1/(2B)\)（即采样频率
\(f_s \geq 2B\)，其中 \(2B\) 称为 Nyquist 率）对 \(f\) 均匀采样，则
\(f\) 可由其样本 \(\{f(n\Delta t)\}_{n \in \mathbb{Z}}\) 完全重建：
\[f(t) = \sum_{n=-\infty}^\infty f(n\Delta t) \cdot \operatorname{sinc}\left(\frac{t - n\Delta t}{\Delta t}\right), \quad \operatorname{sinc}(x) = \frac{\sin(\pi x)}{\pi x}\]

\emph{证明概要（利用 Poisson 求和公式）}：定义采样信号的 Dirac 梳
\(\operatorname{III}_{\Delta t}(t) = \sum_{n \in \mathbb{Z}} \delta(t - n\Delta t)\)，则采样值可写为
\(f_s(t) = f(t) \operatorname{III}_{\Delta t}(t)\)。Poisson 求和公式给出
\(\operatorname{III}_{\Delta t}\) 的 Fourier
变换：\(\widehat{\operatorname{III}_{\Delta t}}(\omega) = \frac{2\pi}{\Delta t} \sum_{k \in \mathbb{Z}} \delta(\omega - \frac{2\pi k}{\Delta t})\)。由卷积定理，采样信号的频谱为
\(\hat{f}_s(\omega) = \frac{1}{2\pi} (\hat{f} * \widehat{\operatorname{III}_{\Delta t}})(\omega) = \frac{1}{\Delta t} \sum_{k \in \mathbb{Z}} \hat{f}(\omega - \frac{2\pi k}{\Delta t})\)------即
\(\hat{f}_s\) 是 \(\hat{f}\) 以 \(2\pi/\Delta t\)
为周期的周期化拷贝的叠加。当 \(\Delta t \leq 1/(2B)\) 时，周期
\(2\pi/\Delta t \geq 4\pi B\)，各拷贝
\(\hat{f}(\omega - 2\pi k/\Delta t)\) 的支撑互不重叠。在区间
\([-2\pi B, 2\pi B]\) 上乘以矩形窗函数
\(\operatorname{rect}(\omega / (4\pi B))\) 即可恢复 \(\hat{f}\)。再取
Fourier 逆变换，矩形窗对应时域 \(\operatorname{sinc}\)
函数的卷积，即得重建公式。若 \(\Delta t > 1/(2B)\)，频谱拷贝重叠（混叠 /
aliasing），信息不可恢复。∎

压缩感知（Compressed Sensing / Compressive
Sampling）是信号处理领域的一项革命性成果（Candes-Romberg-Tao 和 Donoho,
2004-2006）------它表明，只要信号在某个基下是\textbf{稀疏的}，就能从远少于
Nyquist-Shannon
采样定理要求的测量值中\textbf{精确重建}信号。这一发现改变了数据获取的范式：采样可以与压缩\textbf{同时进行}，而非''先采样、后压缩''的传统流程。

\subsubsection{\texorpdfstring{284.1 稀疏性与 \(\ell_0\)
优化}{284.1 稀疏性与 \textbackslash ell\_0 优化}}\label{ux7a00ux758fux6027ux4e0e-ell_0-ux4f18ux5316}

\textbf{定义 284.1}（稀疏信号）：向量 \(x \in \mathbb{R}^N\) 称为
\textbf{\(s\)-稀疏}的，若
\(\|x\|_0 := |\operatorname{supp}(x)| \leq s \ll N\)。更一般地，\(x\)
在基（或框架）\(\Psi\) 下的表示 \(\alpha = \Psi x\)
可能是稀疏的（变换域稀疏性）。

\textbf{定义 284.2}（压缩感知问题）：给定\textbf{测量矩阵}
\(A \in \mathbb{R}^{m \times N}\)（\(m \ll N\)，即测量数远小于信号维数）和测量值
\(y = Ax\)（或含噪版本 \(y = Ax + e\)），从 \(y\) 中重建
\(x\)。最直接的形式是 \(\ell_0\) 极小化：
\[\min_{x \in \mathbb{R}^N} \|x\|_0 \quad \text{s.t.} \quad Ax = y\]
然而，\(\ell_0\) 极小化是 \textbf{NP-hard}
的组合优化问题，因为需要在所有 \(\binom{N}{s}\) 个可能的支撑集上搜索。

\textbf{命题 284.1}（\(\ell_0\) 恢复的唯一性条件---Spark）：若矩阵 \(A\)
的 \textbf{Spark}（最小线性相关列数）满足
\(\operatorname{spark}(A) > 2s\)，则至多存在一个 \(s\)-稀疏解满足
\(Ax = y\)。Spark 是刻画 \(\ell_0\) 唯一性的基本概念，但一般情况下计算
Spark 也是 NP-hard 的。

\subsubsection{\texorpdfstring{284.2 \(\ell_1\)
松弛与限制等距性质（RIP）}{284.2 \textbackslash ell\_1 松弛与限制等距性质（RIP）}}\label{ell_1-ux677eux5f1bux4e0eux9650ux5236ux7b49ux8dddux6027ux8d28rip}

\textbf{定理 284.1}（压缩感知基本定理 / \(\ell_1\)
松弛的合理性，Candes-Romberg-Tao 和 Donoho, 2004-2006）：若测量矩阵
\(A\) 满足适当的不相干条件（或限制等距性质），则凸的 \(\ell_1\) 极小化
\[\min_{x \in \mathbb{R}^N} \|x\|_1 \quad \text{s.t.} \quad Ax = y\]
（或含噪情形
\(\min_x \|x\|_1 \text{ s.t. } \|Ax - y\|_2 \leq \epsilon\)，称为基追踪去噪
/ BPDN）能以高概率\textbf{精确重建} \(s\)-稀疏信号 \(x\)，仅需测量数
\[m \geq C \cdot s \cdot \log(N/s)\] 其中 \(C\)
为通用常数。这一测量数远小于 Nyquist 采样率（需
\(m = \mathcal{O}(N)\)），而仅与\textbf{稀疏度 \(s\)
对数依赖}地依赖于信号维数 \(N\)。

\textbf{定义 284.3}（限制等距性质 / RIP，Candes-Tao 2005）：矩阵
\(A \in \mathbb{R}^{m \times N}\) 满足 \textbf{\(s\)
阶限制等距性质}（Restricted Isometry Property），若存在常数
\(\delta_s \in (0, 1)\) 使得对\textbf{所有} \(s\)-稀疏向量
\(x \in \mathbb{R}^N\)，
\[(1 - \delta_s)\|x\|_2^2 \leq \|Ax\|_2^2 \leq (1 + \delta_s)\|x\|_2^2\]
RIP 确保测量矩阵 \(A\) 在
\(s\)-稀疏向量的子集上近似为保距变换（isometry），从而不破坏稀疏信号中蕴含的几何结构信息。\(\delta_s\)
越小，\(A\) 在稀疏子空间上的行为越接近正交投影。

\textbf{定理 284.2}（RIP 的随机矩阵实现）：以下随机矩阵以高概率满足
\(\delta_{2s} \leq \delta\)（\(\delta\) 为任意正常数），所需测量数为
\(m = \mathcal{O}(s \log(N/s))\)： - \textbf{独立同分布 Gaussian
矩阵}：\(A_{ij} \sim \mathcal{N}(0, 1/m)\) - \textbf{Bernoulli
矩阵}：\(A_{ij} = \pm 1/\sqrt{m}\) 等概率 - \textbf{部分随机 Fourier
矩阵}：随机选取 \(m\) 行 DFT 矩阵（在 MRI 加速成像中尤为重要） -
\textbf{部分 Hadamard 矩阵}和一般的次 Gaussian 随机矩阵

\textbf{定理 284.3}（\(\ell_1\) 恢复的稳定性与稳健性 /
Candes-Romberg-Tao 2006）：若 \(A\) 满足
\(\delta_{2s} < \sqrt{2} - 1\)，则对含噪观测
\(y = Ax + e\)（\(\|e\|_2 \leq \epsilon\)），\(\ell_1\) 极小化解（BPDN）
\[\hat{x} = \arg\min \|x\|_1 \quad \text{s.t.} \quad \|Ax - y\|_2 \leq \epsilon\]
满足误差界
\[\|\hat{x} - x\|_2 \leq C_1 \epsilon + C_2 \frac{\|x - x_s\|_1}{\sqrt{s}}\]
其中 \(x_s\) 为 \(x\) 的最佳 \(s\)-项逼近（保留 \(s\)
个最大绝对值分量）。该界揭示了压缩感知的两个核心性质： -
\textbf{稳健性（Robustness）}：恢复误差与噪声水平 \(\epsilon\)
成线性关系------对测量噪声不敏感 -
\textbf{稳定性（Stability）}：恢复误差与信号的可压缩性 \(\|x - x_s\|_1\)
成线性关系------即使信号不完全稀疏但可压缩（系数快速衰减），也能近似恢复

\subsubsection{284.3 恢复算法}\label{ux6062ux590dux7b97ux6cd5-1}

\textbf{算法 284.1}（正交匹配追踪 /
OMP）：贪心算法，每次迭代选择与当前残差最相关的原子（\(A\)
的列），然后求解所选支撑上的最小二乘问题，更新残差。OMP
简单高效，每步计算复杂度为 \(\mathcal{O}(mN)\)，理论上在 RIP
条件下也有恢复保证（若 \(\delta_{s+1} < 1/(3\sqrt{s})\)）。

\textbf{算法 284.2}（迭代硬阈值 / IHT, Blumensath-Davies 2008）：对
\(\ell_0\) 约束问题的贪心迭代
\[x^{(k+1)} = H_s\left(x^{(k)} + \mu A^T(y - Ax^{(k)})\right)\] 其中
\(H_s\) 为\textbf{硬阈值算子}（保留绝对值最大的 \(s\)
个分量，其余分量置零），\(\mu\) 为步长。IHT 可视为 \(\ell_0\)
问题的近端梯度法------梯度步后投影回 \(s\)-稀疏向量集合。

\textbf{算法 284.3}（CoSaMP / 压缩采样匹配追踪，Needell-Tropp
2008）：融合了组合算法和凸优化的思想： 1.
信号代理：\(u = A^T (y - A x^{(k)})\)，选择 \(2s\) 个最大相关分量 2.
支撑合并：\(T = \operatorname{supp}(x^{(k)}) \cup \{\text{$2s$ 个最大 $|u_i|$ 的索引}\}\)
3. 最小二乘估计：\(\tilde{x}_T = A_T^\dagger y\)（在合并支撑上求解） 4.
剪枝：\(x^{(k+1)} = H_s(\tilde{x})\)（保留 \(s\) 个最大分量的支撑） 5.
残差更新：重复至收敛

CoSaMP
结合了计算效率（每次迭代仅需矩阵-向量乘法和最小二乘）和严格的恢复保证：在
RIP 条件下线性收敛到真解。

\textbf{应用 284.1}（压缩感知 MRI 加速成像）：在 MRI
中，每次测量需要较长的扫描时间。压缩感知 MRI（Lustig-Donoho-Pauly
2007）利用 MRI 图像在小波基和全变差下的稀疏性，通过随机欠采样
\(k\)-空间（相当于部分 Fourier 测量矩阵），将扫描时间减少 3-8
倍，同时保持诊断图像质量。这是压缩感知在医学成像中最成功的应用之一。

\bigskip\hrule\bigskip

\subsection{第285章：稀疏建模与字典学习}\label{ux7b2c285ux7ae0ux7a00ux758fux5efaux6a21ux4e0eux5b57ux5178ux5b66ux4e60}

压缩感知假定信号在某个\textbf{固定基}（如 Fourier
基、小波基）下是稀疏的。字典学习的核心思想更进一步------从数据本身\textbf{学习最优的稀疏表示基}（字典），使信号在该字典下具有最稀疏的表示。这种数据驱动的方法在许多图像处理任务中显著优于固定基方法。

\subsubsection{285.1
字典学习的数学模型}\label{ux5b57ux5178ux5b66ux4e60ux7684ux6570ux5b66ux6a21ux578b}

\textbf{定义 285.1}（字典学习问题）：给定训练数据
\(\{y_i\}_{i=1}^M \subset \mathbb{R}^n\)，求字典
\(D \in \mathbb{R}^{n \times K}\)（通常是\textbf{过完备}的，即
\(K > n\)）和稀疏系数矩阵
\(X = [x_1, \ldots, x_M] \in \mathbb{R}^{K \times M}\)，使得
\[\min_{D, X} \sum_{i=1}^M \|y_i - D x_i\|_2^2 \quad \text{s.t.} \quad \|x_i\|_0 \leq T_0 \text{（}i = 1, \ldots, M\text{）}, \|d_j\|_2 = 1 \text{（}j = 1, \ldots, K\text{）}\]
其中对字典原子（\(D\) 的列）的单位范数约束是为了避免
\textbf{尺度模糊性}（scale
ambiguity）------字典原子和系数的尺度可以交换（\(d \cdot \alpha = (\lambda d) \cdot (\alpha / \lambda)\)），不加约束会导致优化退化。

\textbf{算法 285.1}（K-SVD 算法 / Aharon-Elad-Bruckstein,
2006）：交替优化框架：

\textbf{阶段一 --- 稀疏编码}（固定 \(D\)，优化 \(X\)）：对每个训练样本
\(y_i\)，
\[x_i = \arg\min_x \|y_i - Dx\|_2^2 \quad \text{s.t.} \quad \|x\|_0 \leq T_0\]
可使用 OMP（正交匹配追踪）高效求解。

\textbf{阶段二 --- 字典更新}（逐原子更新）：对每个原子
\(d_k\)（\(k = 1, \ldots, K\)），定义使用该原子的样本索引集
\(\omega_k = \{i : x_{k,i} \neq 0\}\)，计算残差矩阵（除去原子 \(d_k\)
贡献后的误差）
\[E_k = Y_{\omega_k} - \sum_{j \neq k} d_j x_{j,\omega_k}\] 其中
\(Y_{\omega_k}\) 为 \(\omega_k\) 索引的样本矩阵。对 \(E_k\) 进行 SVD
分解 \(E_k = U \Sigma V^T\)，取最大奇异值对应的左奇异矢量作为 \(d_k\)
的新值，同时更新相应的系数向量。该过程保证了在每次更新中重建误差单调不增------K-SVD
因此收敛到局部极小点。

\textbf{定理 285.1}（字典辨识的理论保证，Spielman-Wang-Wright 2012
等）：在适当的不相干条件下（字典原子间角度足够大），从完备字典（\(K = n\)）下的稀疏表示样本中可以多项式时间辨识出真实字典。但对于过完备字典（\(K > n\)），严格的理论保证仍然具有挑战性------这反映了字典学习非凸优化的内在困难。

\subsubsection{285.2
稀疏表示与图像处理应用}\label{ux7a00ux758fux8868ux793aux4e0eux56feux50cfux5904ux7406ux5e94ux7528}

\textbf{定义 285.2}（图像去噪的稀疏表示框架 / Elad-Aharon
方法，2006）：将图像 \(u\)
分割为重疊的\textbf{小块}（patches）\(\{R_i u\}\)（\(R_i\)
为提取算子），去噪过程为： 1. 对每个含噪小块 \(R_i f\)
求稀疏表示：\(\hat{\alpha}_i = \arg\min_{\alpha} \|\alpha\|_0 \text{ s.t. } \|R_i f - D\alpha\|_2^2 \leq C \sigma^2\)（\(\sigma\)
为噪声水平） 2. 全局重建通过加权平均：
\[\hat{u} = \left(\lambda I + \sum_i R_i^T R_i\right)^{-1} \left(\lambda f + \sum_i R_i^T D \hat{\alpha}_i\right)\]
其中 \(\lambda\) 为 Lagrange
乘子（平衡全局保真与局部稀疏）。由于小块是重叠的，每个像素由多个小块覆盖，平均操作自然产生平滑的重建结果。

\textbf{命题 285.1}（稀疏性与自然图像统计）：自然图像的小块在过完备 DCT
或学习到的字典下呈现高度稀疏性------仅约 5-10\% 的系数非零即可实现 PSNR
\textgreater{} 30dB
的高质量图像重构。这一经验观察是稀疏编码在图像处理中成功的基础。更精确地，自然图像小块的边际分布服从重尾分布（heavy-tailed），其稀疏性反映了自然图像中主导的结构信息（边缘、纹理）是高维空间中的低维流形。

\subsubsection{285.3
卷积稀疏编码与分析模型}\label{ux5377ux79efux7a00ux758fux7f16ux7801ux4e0eux5206ux6790ux6a21ux578b}

\textbf{定义 285.3}（卷积稀疏编码 / CSC）：将图像建模为卷积和
\[u = \sum_{k=1}^K d_k * x_k\] 其中 \(\{d_k\}_{k=1}^K\)
为\textbf{卷积字典滤波器}（通常较小，如 \(11 \times 11\)），\(\{x_k\}\)
为\textbf{稀疏特征图}（与输入图像同尺寸，仅在特定位置非零）。与小块字典学习不同，卷积模型具有\textbf{平移不变性（Shift-Invariance）}------相同的特征出现在图像的不同位置时，由同一滤波器的不同位置响应自然捕获，避免了小块方法中因独立处理各小块而产生的''块状''伪影。

\textbf{定义 285.4}（分析算子模型 / Co-sparse Analysis
Model）：分析模型提供对偶视角：信号 \(u\) 在\textbf{分析算子}
\(\Omega \in \mathbb{R}^{P \times N}\)（\(P \gg N\)，过完备分析算子）下应具有稀疏表示------\(\Omega u\)
的大多数分量为零（即 \(u\) 位于 \(\Omega\)
的零空间的低余维子空间中）。例如，TV
范数是梯度算子（分析算子）\(\nabla\) 作用后 \(\ell_1\) 稀疏性的度量：
\[\|u\|_{TV} = \|\nabla u\|_1 = \sum_i |(\nabla u)_i|\]
分析稀疏模型的形式为 \(\min_u \|\Omega u\|_1 \text{ s.t. }\)
数据约束，是综合稀疏模型（\(u = D\alpha, \|\alpha\|_0\)
小）的对偶。两者关键差异：综合模型直接建模信号本身，分析模型建模信号的变换。在某些应用中（如
CT 重建），分析模型因自然编码了信号的几何先验（梯度稀疏 =
分段常值）而具有优势。

\bigskip\hrule\bigskip

\subsection{第286章：深度学习驱动的图像重建}\label{ux7b2c286ux7ae0ux6df1ux5ea6ux5b66ux4e60ux9a71ux52a8ux7684ux56feux50cfux91cdux5efa}

深度学习彻底改变了图像重建与处理领域------从卷积神经网络（CNN）的直接端到端学习，到基于物理模型和数据驱动的混合方法（如
Plug-and-Play
和深度展开），深度学习为传统数学方法注入了新的生命力，同时也带来了关于可解释性和理论保证的新挑战。

\subsubsection{286.1
卷积神经网络与图像处理}\label{ux5377ux79efux795eux7ecfux7f51ux7edcux4e0eux56feux50cfux5904ux7406}

\textbf{定义
286.1}（图像重建的深度学习方法）：将图像重建视为\textbf{映射学习}
\(f_\theta: y \mapsto x\)（从退化观测 \(y\) 到干净图像 \(x\)
的映射），通过大量配对训练数据 \(\{(y_i, x_i)\}_{i=1}^N\) 最小化重建误差
\[\min_\theta \frac{1}{N} \sum_{i=1}^N \|f_\theta(y_i) - x_i\|_2^2\]
（或更一般的感知损失、GAN 对抗损失等）。这里的 \(f_\theta\)
由深度神经网络参数化。

\textbf{代表性架构}： - \textbf{DnCNN}（Zhang et al., 2017）：残差学习
\(f_\theta(y) = y - R_\theta(y)\)（网络学习的是噪声残差而非干净图像），结合
Batch Normalization 和 ReLU 激活，在 Gaussian 去噪上取得 SOTA 性能 -
\textbf{U-Net}（Ronneberger et al.,
2015）：编码器-解码器架构加跳跃连接（skip
connections）------跳跃连接将编码器的高分辨率特征直接传递给解码器，解决了深层网络的梯度消失问题，特别适合图像分割和重建任务
- \textbf{ResNet}（He et al., 2016）：残差块
\(x_{l+1} = x_l + \mathcal{F}(x_l, W_l)\)
解决了极深网络的退化问题，使数百层的网络可以有效训练

\textbf{对比传统方法}：CNN
的优势在于从数据中自动学习最优的图像先验（\textbf{隐式正则化}），而不依赖人工设计的手工正则项（如
TV、小波稀疏等）。然而，纯数据驱动的方法缺乏理论保证------网络的黑盒性质使人们难以理解其泛化行为，且对分布外数据（OOD）表现不可预测。

\subsubsection{286.2 Plug-and-Play
与深度展开}\label{plug-and-play-ux4e0eux6df1ux5ea6ux5c55ux5f00}

\textbf{定义 286.2}（Plug-and-Play 先验 / PnP,
Venkatakrishnan-Bouman-Wohlberg 2013）：将经典优化算法（如 ADMM、HQS /
Half-Quadratic Splitting）的正则化步骤 \(\operatorname{prox}_R\)
替换为预训练的\textbf{深度去噪器} \(D_\sigma\)。以 HQS 框架为例，求解
\[\min_x \frac{1}{2}\|Ax - y\|^2 + \lambda R(x)\] 的迭代格式为：
\[x_{k+1} = D_{\sigma_k}(\hat{x}_k)\]
\[\hat{x}_{k+1} = \arg\min_x \frac{1}{2}\|Ax - y\|^2 + \frac{\mu}{2}\|x - x_{k+1}\|^2\]
其中 \(\sigma_k\)
逐步减小（相当于正则化逐渐减弱，从粗到细）。第一式是去噪步骤（由深度去噪器实现，代替了
\(\operatorname{prox}_R\)），第二式是数据保真步骤（通常有闭式解或通过共轭梯度求解）。这一框架的优美之处在于：它\textbf{解耦了物理测量模型}（数据保真项，保留了成像物理）和\textbf{图像先验}（深度去噪器，捕获了自然图像的统计规律）。

\textbf{算法 286.1}（PnP-ADMM）：更一般的 PnP 框架使用 ADMM 分裂：
\[x_{k+1} = D_\sigma(z_k - u_k) \quad \text{（正则化步骤，由去噪器实现）}\]
\[z_{k+1} = \operatorname{proj}_{\{w: \|Aw - y\| \leq \epsilon\}}(x_{k+1} + u_k) \quad \text{（数据保真投影）}\]
\[u_{k+1} = u_k + x_{k+1} - z_{k+1} \quad \text{（对偶变量更新）}\]

PnP 方法的重要理论问题：\textbf{何时收敛？} Ryu et
al.（2019）证明了当去噪器满足 Lipschitz 条件且 Lipschitz
常数足够小时，PnP-ADMM 收敛到某个不动点。实践中，先进的 CNN
去噪器（通常不满足收缩性）在 PnP 框架中仍然表现良好且经验稳定。

\textbf{定义 286.3}（深度展开 / Algorithm Unrolling, Gregor-LeCun
2010）：将迭代优化算法的每一次迭代展开为网络的一层，使算法参数成为\textbf{可学习的网络参数}：

\textbf{LISTA}（Learned ISTA）：将 \(K\) 次 ISTA（Iterative
Soft-Thresholding Algorithm）迭代
\[x^{(k+1)} = S_{\lambda/L}\left(x^{(k)} - \frac{1}{L} A^T(A x^{(k)} - y)\right)\]
展开为 \(K\) 层循环网络，其中阈值 \(\lambda\) 和权重矩阵 \(W_e, W_g\)
均可学习： \[x^{(k+1)} = S_{\theta_k}\left(W_e y + W_g x^{(k)}\right)\]

\textbf{优点}：深度展开继承了模型驱动的可解释性（每层对应于优化的一步，结构源于问题的物理）和理论框架，同时获得了数据驱动的自适应能力（从训练数据中学习最优的超参数，如步长、正则化强度）。\textbf{缺点}：展开的网络层数有限（一般
\(K \leq 10-20\)），因为计算图随 \(K\) 线性增大；网络结构固定，无法像纯
CNN 那样灵活扩展深度。

\subsubsection{286.3
生成模型与逆问题}\label{ux751fux6210ux6a21ux578bux4e0eux9006ux95eeux9898}

\textbf{定义 286.4}（基于生成模型的图像重建 / CSGM 框架,
Bora-Jalal-Price-Dimakis
2017）：不同于在像素空间中重建，生成模型方法在\textbf{潜在空间中求解}：
\[\min_{z \in \mathbb{R}^k} \|A G(z) - y\|_2^2\] 其中
\(G: \mathbb{R}^k \to \mathbb{R}^N\) 为预训练的生成器（如 GAN、VAE
的解码器），其输出分布
\(G(z)\)（\(z \sim p_z\)）近似于自然图像的流形。这种方法利用生成模型的\textbf{强大低维先验}------自然图像集被建模为
\(k\) 维潜在空间在 \(N\)
维像素空间中的嵌入（\(k \ll N\)），在潜在空间中优化自然地将重建约束在自然图像的流形上。

\textbf{定理 286.1}（生成模型压缩感知的理论保证，Bora et
al.~2017）：设生成器 \(G: \mathbb{R}^k \to \mathbb{R}^N\) 为
\(L\)-Lipschitz 的 \(d\) 层神经网络（ReLU
激活），\(A \in \mathbb{R}^{m \times N}\) 满足 S-REC（Set-Restricted
Eigenvalue Condition）且在 \(G\) 的范围内成立，则仅需
\[m = \mathcal{O}(k d \log L)\]
次测量即可保证重建误差有界。这一测量数可与压缩感知经典结果的
\(s \log N\) 对比：当
\(k \ll s \ll N\)（即潜在空间维数远小于稀疏度）时，生成模型方法需要更少的测量------这反映了生成模型比稀疏先验更强（``更低的自由度''）的先验知识。

\subsubsection{286.4 评述与展望}\label{ux8bc4ux8ff0ux4e0eux5c55ux671b}

\textbf{命题 286.1}（数学理论 vs.~数据驱动方法的核心张力）：

{\def\LTcaptype{none} % do not increment counter
\begin{longtable}[]{@{}
  >{\raggedright\arraybackslash}p{(\linewidth - 4\tabcolsep) * \real{0.1304}}
  >{\raggedright\arraybackslash}p{(\linewidth - 4\tabcolsep) * \real{0.5870}}
  >{\raggedright\arraybackslash}p{(\linewidth - 4\tabcolsep) * \real{0.2826}}@{}}
\toprule\noalign{}
\begin{minipage}[b]{\linewidth}\raggedright
维度
\end{minipage} & \begin{minipage}[b]{\linewidth}\raggedright
数学方法（PDE/变分/稀疏）
\end{minipage} & \begin{minipage}[b]{\linewidth}\raggedright
深度学习方法
\end{minipage} \\
\midrule\noalign{}
\endhead
\bottomrule\noalign{}
\endlastfoot
\textbf{可解释性} & 强------每个分量有明确的数学含义 &
弱------黑盒模型 \\
\textbf{理论保证} & 强------恢复界、收敛率 & 弱------主要为经验验证 \\
\textbf{数据需求} & 几乎不需要训练数据 & 需要大量高质量配对数据 \\
\textbf{泛化能力} & 在数学模型适用范围内容易保证 &
对分布外数据表现不可预测 \\
\textbf{对复杂模式的建模} & 受限于手工设计的先验 &
可从数据中学习任意复杂模式 \\
\textbf{计算效率} & 推理时需迭代优化 & 推理时仅需单次前向传播 \\
\end{longtable}
}

\textbf{命题 286.2}（融合的方向）：未来的发展方向在于融合两者的优势： -
\textbf{数学指导网络设计}：用 PDE 和变分原理指导网络架构设计（如等变 CNN
源于对称性原理，曲率驱动的网络层源于几何 PDE） -
\textbf{网络增强数学方法}：用深度学习从数据中学习自适应先验、加速迭代求解器
-
\textbf{理论分析数据驱动方法}：用调和分析、统计学习的工具分析深度网络的表示能力、泛化误差和收敛性（如
Neural Tangent Kernel 理论、频率原则）

\textbf{结语}：从 Fourier 到小波，从 TV
去噪到压缩感知，从字典学习到深度重建------图像与信号处理的数学基础始终围绕着三个核心问题：（1）\textbf{如何表示信号}（基的选择------Fourier
基、小波基、字典原子、深度特征）；（2）\textbf{如何建模先验知识}（正则化------\(\ell_2\)
光滑性、TV 保边性、\(\ell_1\) 稀疏性、GAN
流形）；（3）\textbf{如何高效求解}（算法设计------阈值迭代、原始-对偶、ADMM、梯度下降、反向传播）。每一种新的数学工具都将这三个问题推进到新的境界，而现代深度学习则为这一经典领域带来了前所未有的数据驱动的全新视角。

\end{document}
